\PassOptionsToPackage{hyperfootnotes=false,bookmarksnumbered=true,bookmarksopen=true,bookmarksopenlevel=0}{hyperref}
\documentclass[11pt,oneside]{book}
\usepackage{fontspec}
\usepackage[main=french,english,french]{babel}
\usepackage[leqno]{amsmath}
\usepackage{amssymb}
\usepackage{amsthm}
\usepackage{mathrsfs}
\usepackage{tikz-cd}
\usepackage{enumitem}
\usepackage[a4paper,left=25mm,right=30mm,top=24mm,bottom=26mm,headheight=14pt,marginparwidth=21mm,marginparsep=4mm]{geometry}
\usepackage{microtype}
\usepackage{marginnote}
\usepackage{xurl}
\usepackage{adjustbox}
\usepackage{needspace}
\usepackage{marginfix}
\usepackage{etoolbox}
\usepackage{lmodern}
\usepackage{array}
\usepackage{ragged2e}
\usepackage{fancyhdr}
\usepackage{mathtools}
\usepackage{booktabs}
\usepackage{derivateurs_companion}
\usetikzlibrary{decorations.pathreplacing,babel}
\frenchsetup{AutoSpacePunctuation=false}
\setlength{\emergencystretch}{2em}
\hypersetup{pdftitle={Les Dérivateurs: an integrated reader, with commentary},pdfauthor={Alexander Grothendieck},
  colorlinks=true,linkcolor=black,citecolor=black,urlcolor=blue!45!black,pdfdisplaydoctitle=true}
\makeatletter
% counters as in the chunks' own readers (article class): equations by chunk, with unique PDF anchors
\newcounter{dgrchunk}
\renewcommand{\theequation}{\arabic{equation}}
\renewcommand{\thesection}{\arabic{section}}
\renewcommand{\thefigure}{\arabic{figure}}
\renewcommand{\thetable}{\arabic{table}}
\renewcommand{\theHequation}{\arabic{dgrchunk}.\arabic{equation}}
\renewcommand{\theHsection}{\arabic{dgrchunk}.\arabic{section}}
\renewcommand{\theHfigure}{\arabic{dgrchunk}.\arabic{figure}}
\renewcommand{\theHtable}{\arabic{dgrchunk}.\arabic{table}}
% running heads
\def\ps@dgr{\let\@mkboth\@gobbletwo
  \def\@oddhead{\footnotesize\scshape\rightmark\hfil\upshape\thepage}\let\@evenhead\@oddhead
  \def\@oddfoot{}\let\@evenfoot\@oddfoot}
\pagestyle{dgr}
% the contents: wider boxes for the chapter numerals, the labels of the parts and three-digit page numbers
\renewcommand*{\@pnumwidth}{2.2em}
\renewcommand*{\@tocrmarg}{3.2em}
\renewcommand*{\l@chapter}[2]{%
  \ifnum\c@tocdepth>\m@ne
    \addpenalty{-\@highpenalty}\vskip 1.0em \@plus\p@
    \setlength\@tempdima{3.8em}%
    \begingroup
      \parindent\z@ \rightskip\@pnumwidth \parfillskip -\@pnumwidth
      \leavevmode\bfseries \advance\leftskip\@tempdima \hskip -\leftskip
      #1\nobreak\hfil\nobreak\hb@xt@\@pnumwidth{\hss #2\kern-\p@\kern\p@}\par
      \penalty\@highpenalty
    \endgroup
  \fi}
\renewcommand*{\l@section}{\@dottedtocline{1}{1.5em}{4.6em}}
\renewcommand*{\l@subsection}{\@dottedtocline{2}{6.1em}{3.2em}}
% a chapter, or a part of one, that this book does not contain: a line of the contents without page number
\newcommand*{\dgrtocmissingchapter}[2]{%
  \addpenalty{-\@highpenalty}\vskip 1.0em \@plus\p@
  \begingroup\parindent\z@\leftskip 3.8em\rightskip\@tocrmarg plus 1fil\hskip-3.8em
  \hb@xt@ 3.8em{\bfseries #1\hfil}\textit{#2}\enspace\textendash\enspace\textit{not included}\par\endgroup}
\newcommand*{\dgrtocmissing}[2]{%
  \begingroup\parindent\z@\leftskip 6.1em\rightskip\@tocrmarg plus 1fil\hskip-4.6em
  \hb@xt@ 4.6em{#1\hfil}\textit{#2}\enspace\textendash\enspace\textit{not included}\par\endgroup}
% a chapter of the work: #1 numeral (contents), #2 heading line, #3 title; without numeral the contents keep the
% numeral's place empty, and the bookmark has no number
\DeclareRobustCommand{\dgrnonumber}{\numberline{}}
\pdfstringdefDisableCommands{\def\dgrnonumber{}}
\newcommand{\dgrchapter}[3]{%
  \clearpage\thispagestyle{plain}%
  \phantomsection
  \if\relax\detokenize{#1}\relax\addcontentsline{toc}{chapter}{\protect\dgrnonumber #3}%
  \else\addcontentsline{toc}{chapter}{\protect\numberline{#1}#3}\fi
  \markright{#2}%
  \begingroup\centering
    \vspace*{2\baselineskip}{\large\scshape #2\par}\bigskip{\LARGE\bfseries #3\par}\vspace{2\baselineskip}%
  \endgroup}
% a chunk (or its part in one chapter): #1 1 = on a new page, #2 label, #3 contents line, #4 heading, #5 running head
\newcommand{\dgrpart}[5]{%
  \ifnum#1=1 \clearpage\fi
  \phantomsection\addcontentsline{toc}{section}{\protect\numberline{#2}#3}%
  \markright{#5}%
  \par\begingroup\centering\sffamily\small #4\par\endgroup
  \smallskip\hrule\bigskip}
\newcommand{\dgroverview}[2]{%
  \clearpage
  \phantomsection\addcontentsline{toc}{subsection}{Overview}%
  \markright{#2}%
  \section*{Overview \textperiodcentered\ #1}}
\newcommand{\dgrblock}[2]{\pdfbookmark[2]{#2}{#1}}
% a chunk: its group, language and counters
\newcommand{\dgrchunkbegin}{\begingroup\stepcounter{dgrchunk}\selectlanguage{french}%
  \setcounter{footnote}{0}\setcounter{equation}{0}\setcounter{section}{0}\setcounter{figure}{0}\setcounter{table}{0}}
\newcommand{\dgrchunkend}{\par\endgroup}
% the back matter
\newcommand{\dgrbackchapter}[1]{\chapter*{#1}\phantomsection\addcontentsline{toc}{chapter}{#1}\markright{#1}}
\renewcommand{\bibsection}{\dgrbackchapter{\bibname}}
\newcommand{\dgrclaim}[4]{\par\noindent\hangindent=2.5em\hangafter=1
  \textbf{ms.~#2}\enspace\textsc{#1}\enspace{\footnotesize #3}\enspace\textendash\enspace#4\par\smallskip}
\makeatother
\begin{document}
\selectlanguage{english}
\frontmatter
\begin{titlepage}
\centering
\vspace*{0.24\textheight}
{\Large Alexander Grothendieck\par}
\vspace{2.2em}
{\Huge\bfseries Les Dérivateurs\par}
\vspace{2.2em}
{\Large an integrated reader, with commentary\par}
\vfill
\end{titlepage}
\setcounter{page}{2}
\thispagestyle{empty}
\noindent The text is Alexander Grothendieck's, from the edition of \emph{Les Dérivateurs} by M.~Künzer,
J.~Malgoire and G.~Maltsiniotis (Universités de Montpellier~II et Paris~VII),
\url{https://webusers.imj-prg.fr/~georges.maltsiniotis/groth/Derivateurs.html}. The \LaTeX{} transcription of that
edition is AI-transcribed (Claude Opus~5.5, Anthropic). The commentary, the register of claims and the corrections
are AI-produced by GPT-6 Astra (OpenAI, ChatGPT Pro mode), checked and integrated by Claude (Anthropic).

\bigskip
\noindent\textbf{Rights}

\smallskip
\noindent The rights in the original work belong to its author or the author's representatives or heirs. Where we
hold rights, in the AI-produced transcription, translation, corrections and typesetting, they are dedicated to the
public domain under CC0 1.0. This edition is published for educational, non-profit purposes, and we claim no rights we
do not hold. If you hold rights in this work and object to its publication here, please open a GitHub issue and we will
take it down.

\smallskip
\begin{otherlanguage}{french}
\noindent Les droits sur l'œuvre originale appartiennent à son auteur, à ses représentants ou à ses ayants droit. Là
où nous détenons des droits, sur la transcription, la traduction, les corrections et la composition produites par IA,
ils sont dédiés au domaine public sous CC0 1.0. Cette édition est publiée à des fins pédagogiques et non lucratives, et
nous ne revendiquons aucun droit que nous ne détenons pas. Si vous détenez des droits sur cette œuvre et vous opposez
à sa publication ici, ouvrez une issue sur GitHub et nous la retirerons.
\end{otherlanguage}
\clearpage
\tableofcontents
\mainmatter
% DER_I-a: Grothendieck's text with the commentary set in, then the overview, inside a group that applies the chunk's own definitions
\dgrchapter{I}{Chapter I}{Généralités sur les (pré)dérivateurs}
\dgrpart{0}{I-a}{Part 1 of 2, manuscript pages 1--66}{I-a\enspace\textperiodcentered\enspace Chapter I, part 1 of 2\enspace\textperiodcentered\enspace manuscript pages 1--66}{Chapter I \textperiodcentered\ \textup{I-a}}
\dgrchunkbegin
% DER_I-a: the definitions of the chunk's preamble, applied inside its group
\let\printedpage\relax\newcommand{\printedpage}[1]{\marginpar{\footnotesize p.~#1}}
\let\msp\relax\newcommand{\msp}[1]{\marginpar{\footnotesize [#1]}}
\let\AMSTeX\relax\newcommand{\AMSTeX}{$\mathcal{A}\kern-.1667em\lower.5ex\hbox{$\mathcal{M}$}\kern-.125em\mathcal{S}$-\TeX}
\let\uHom\relax\newcommand{\uHom}{\underline{\mathsf{Hom}}}
\let\Ob\relax\newcommand{\Ob}{\mathsf{Ob}}
\let\Ab\relax\newcommand{\Ab}{\mathsf{Ab}}
\let\bDelta\relax\newcommand{\bDelta}{\boldsymbol{\Delta}}
\let\Dia\relax\newcommand{\Dia}{\mathsf{Dia}}
\let\Cat\relax\newcommand{\Cat}{\mathcal{C}\mathit{at}}
\let\DD\relax\newcommand{\DD}{\mathbb{D}}
\let\Hom\relax\newcommand{\Hom}{\underline{\mathsf{Hom}}}
\let\id\relax\newcommand{\id}{\mathsf{id}}
\let\Del\relax\newcommand{\Del}{\boldsymbol{\Delta}}
\let\Ens\relax\newcommand{\Ens}{\mathcal{E}ns}
\let\defeq\relax\newcommand{\defeq}{\overset{\text{déf}}{=}}
\let\D\relax\newcommand{\D}{\mathbb{D}}
\let\sHom\relax\newcommand{\sHom}{\mathsf{Hom}}
\let\DER\relax\newcommand{\DER}{\mathsf{DER}}
\let\sIm\relax\newcommand{\sIm}{\mathsf{Im}}
\let\sDia\relax\newcommand{\sDia}{\mathsf{Dia}}
\let\sOb\relax\newcommand{\sOb}{\mathsf{Ob}}
\let\sid\relax\newcommand{\sid}{\mathsf{id}}
\let\HOT\relax\newcommand{\HOT}{\mathsf{HOT}}
\let\cA\relax\newcommand{\cA}{\mathcal{A}}
\let\hot\relax\newcommand{\hot}{\mathsf{hot}}
\let\Hot\relax\newcommand{\Hot}{\mathsf{Hot}}
\let\isoto\relax\newcommand{\isoto}{\overset{\sim}{\to}}
\let\isolongto\relax\newcommand{\isolongto}{\overset{\sim}{\longrightarrow}}

% ---- from der1_p001-005.tex
% ===================== page 1 =====================
\printedpage{--} % no page number printed (title page)
% head: (none)

\begin{center}
{\Large\textsf{LES DÉRIVATEURS}}

\bigskip

\textbf{ALEXANDRE GROTHENDIECK}

\bigskip\bigskip

\textbf{Chapitre I}

\medskip

\textbf{Généralités sur les (pré)dérivateurs}
\end{center}

\vfill



\clearpage

% ===================== page 2 =====================
\printedpage{--} % no page number printed (editors' note)
% head: (none)

Ce texte a été déchiffré et transcrit en \AMSTeX{} par G. Maltsiniotis. Il a été édité par M. Künzer, J. Malgoire, et G. Maltsiniotis. La transcription est aussi fidèle que possible au manuscrit. Pour les quelques corrections évidentes, ou les rares commentaires des éditeurs, ainsi que pour la numérotation originale des pages du manuscrit, les caractères de machine à écrire \texttt{[typewriter]} entre crochets sont utilisés. Un point d'interrogation entre crochets signifie que l'on n'est pas sûr du mot qui précède.

Cette édition est provisoire. Les remarques, commentaires et corrections sont bienvenus. Envoyer un message à :

\textbf{maltsin@math.jussieu.fr}

\hfill G. Maltsiniotis

\clearpage

% ===================== page 3 =====================
\printedpage{1}
% head: A. Grothendieck : Dérivateurs, Ch. I, Édition des Universités de Montpellier II et Paris VII
% foot: Version 8 juillet 2001

\msp{page 1}

\section*{Dérivateurs}

\subsection*{\textmd{1.}\ Préliminaires sur les foncteurs images inverses des préfaisceaux sur des catégories d'indices, et leurs adjoints multiples.}

$A$ “catégorie de coefficients”,\\
$I$ “catégorie d'indices”,\\
$A(I)=\uHom(I^{\circ},A)$ préfaisceaux sur $I$ à valeurs dans $A$.

Si $f : I \to J$, on en déduit :
\[
f^{*} : A(J) \to A(I) \quad ,
\]
et on s'intéresse à l'existence des foncteurs adjoints successifs à gauche et à droite. Notation provisoire :
\[
f^{?}\ f_{!}\ f^{*}\ f_{*}\ f^{!} \quad .
\]

Voici les faits principaux :

\dgrblock{dgrb1}{Variance and the two Kan extensions}
\begin{dcomment}{exposition}{p1:Voici les faits principaux}{Variance and the two Kan extensions}
The elementary presheaf calculation fixes the conventions for everything that follows. Throughout these pages an $I$-diagram has arrows in the direction of $I^{\mathrm{op}}$. Thus a diagram on $0\to1$ is an arrow $F_1\to F_0$. We retain this convention in the commentary. If $f:I\to J$, the adjunctions are
\[
 f_!\dashv f^*\dashv f_*.
\]
They are adjunctions of ordinary categories at this stage. Their pointwise formulas, whenever the indicated limits exist, are
\[
 (f_*F)_j=\lim_{(f/j)^{\mathrm{op}}}F\circ\operatorname{pr},
 \qquad
 (f_!F)_j=\operatorname*{colim}_{(j/f)^{\mathrm{op}}}F\circ\operatorname{pr}.
\]
Here $f/j$ has objects $(i,f(i)\to j)$, whereas $j/f$ has objects $(i,j\to f(i))$. The opposite signs are essential: forgetting one exchanges the two formulas. These are the usual pointwise Kan formulas applied to $f^{\mathrm{op}}$; see \citet[\S6.2]{Riehl2016}.

For completeness, a map $G\to f_*F$ consists, at $j$, of a compatible family of maps $G_j\to F_i$, one for every $f(i)\to j$. A map $f^*G\to F$ supplies this family by composing $G_j\to G_{f(i)}\to F_i$. Conversely, the members indexed by identities recover the original map. The two constructions are inverse and natural. The left adjunction follows by reversing all arrows. This proves the existence assertion without any adjoint functor theorem; for finite categories the comma categories are finite, so finite limits or colimits suffice.

The word ``immersion'' here denotes the inclusion of a full subcategory, with ``open'' meaning a sieve and ``closed'' a cosieve. A sieve is closed under taking the source of an arrow whose target belongs to it; a cosieve is closed under taking the target. In a poset these are, respectively, downward and upward closed subsets. This agrees with the conventions forced by the displayed examples. For a fully faithful $f$, the identity at $i$ is terminal in $f/f(i)$ and initial in $f(i)/f$. The formulas therefore give $f^*f_*F\cong F\cong f^*f_!F$. Both Kan extensions are fully faithful. Calling $f^*$ a localization expresses these adjunction identities, not a claim that restriction itself is fully faithful.
\end{dcomment}


\textit{a}) Sous réserve de l'existence des $\varprojlim$ (resp. $\varinjlim$) pertinentes, $f_{*}$ (resp. $f_{!}$) existe. Par exemple s'il existe \texttt{[des]} $\varprojlim$ (resp. $\varinjlim$) \textit{finies}, alors $f_{*}$ (resp. $f_{!}$) existe si $I$, $J$ sont des catégories finies.

\textit{b}) Même si $A$ est une catégorie abélienne, avec $\varprojlim$ et $\varinjlim$ (par exemple $\Ab$), les \textit{biadjoints} n'existent pas toujours. Par exemple, dans le cas particulier où $J=e$ est la catégorie ponctuelle, et $I$ l'ensemble ordonné
\[
\begin{tikzcd}
 & a \\
0 \arrow[ur] \arrow[dr] & \\
 & b
\end{tikzcd}
\]
(dans le cas de $f^{!}$), ou l'ensemble ordonné opposé (dans le cas de $f^{?}$). Même si $f$ est une immersion ouverte, $f^{!}$ n'existe pas forcément (par exemple elle n'existe pas si $I$ est comme \texttt{[ci-]}dessus, et
\[
J=\begin{tikzcd}
 & a \arrow[dr] & \\
0 \arrow[ur] \arrow[dr] & & c \\
 & b \arrow[ur] &
\end{tikzcd}
\qquad\qquad \text{),}
\]
et si $f$ est une immersion fermée, $f^{?}$ n'existe pas forcément.

\dgrblock{dgrb2}{Why the next adjoint can fail}
\begin{dcomment}{examples}{p1:et si $f$ est une immersion fermée}{Why the next adjoint can fail}
The first counterexample is already a calculation with abelian groups. On the fork $0\to a,0\to b$, a presheaf is a cospan $A\to C\leftarrow B$, and its limit is $A\times_C B$. Consider the componentwise surjection of cospans
\[
 (\mathbb Z\xrightarrow{2}\mathbb Z\leftarrow0)
 \longrightarrow
 (\mathbb Z/2\xrightarrow{0}\mathbb Z/2\leftarrow0),
\]
whose nonzero components are reduction modulo $2$. The induced map on limits is $0\to\mathbb Z/2$, which is not surjective. A functor having a right adjoint is a left adjoint and hence preserves cokernels; in this abelian diagram category it would preserve epimorphisms. The limit functor therefore has no right adjoint. This verifies precisely the failure of $f^!$ asserted in (b).

For the opposite fork the colimit is a pushout. The componentwise injection from the span $\mathbb Z\xrightarrow{2}\mathbb Z$, $\mathbb Z\to0$, to the analogous span with $\mathbb Q$ in place of $\mathbb Z$ induces $\mathbb Z/2\to0$ on pushouts. It fails to preserve a monomorphism and cannot be a right adjoint. Thus its left adjoint $f^?$ does not exist. These examples use neither an absence of ordinary limits nor a failure of abelian structure.

The open-immersion example is the same obstruction localized at the new vertex $c$. For the inclusion of the fork into the diamond, $(f_*F)_c$ is the preceding pullback. Evaluation at $c$ preserves colimits, so a right adjoint to $f_*$ would force this pullback functor to preserve cokernels. The same surjection disproves that possibility. The dual argument treats a closed immersion and $f^?$. The point of the examples is that iterated adjunction is a property of a particular functor, not a consequence of the ambient categories having all limits and colimits.
\end{dcomment}


\textit{c}) Si $f$ est une \textit{immersion fermée}, et $A$ est \textit{ponctuée} (\textit{i.e.} admet un objet à la fois initial et final), alors $f^{!}$ existe. Si $f$ est une immersion ouverte, $f^{?}$ existe. (Toujours sous réserve
\printedpage{2}
% head: A. Grothendieck : Dérivateurs, Ch. I, Édition des Universités de Montpellier II et Paris VII
% foot: Version 8 juillet 2001
de l'existence dans $A$ des limites pertinentes). Si $f$ est une immersion, $f_{*}$ et $f_{!}$ \texttt{[sont]} pleinement fidèles \textit{i.e.} $f^{*}$ \texttt{[est]} un foncteur de localisation.

\textit{d}) Supposons seulement l'existence de l'objet final $e$ dans $A$ (plus au besoin certaines $\varprojlim$ finies). Si

\msp{page 2}
\[
i : Y \hookrightarrow S
\]
est une immersion fermée, de complémentaire $U$
\[
Y \overset{i}{\hookrightarrow} S \overset{j}{\hookleftarrow} U
\]
de sorte que $j$ soit une immersion ouverte, alors $A(Y)$ s'identifie par le foncteur $i_{*}$ à la sous-catégorie pleine de $A(S)$ formée des $F \in \Ob\ A(S)$ tels que $F|U \simeq e_{U}$ (préfaisceau final). De même, si l'objet initial $\emptyset$ existe dans $A$, et les $\varinjlim$ pertinentes, alors $A(U)$ s'identifie par $j_{!}$ à la sous-catégorie pleine de $A(S)$ formée des $F$ tels que $F|Y=\emptyset_{Y}$ (préfaisceau initial, prenant partout la valeur $\emptyset$).

Si $A$ est même \textit{ponctuée} ($e=\emptyset$), alors on a, pour tout $F\in \Ob\ A(S)$, un carré à la fois cartésien et cocartésien, fonctoriel en $F$
\[
\begin{tikzcd}
 & F \arrow[dr, "\text{adj}"] & \\
j_{!}j^{*}F \arrow[ur, "\text{adj}"] \arrow[dr] & & i_{*}i^{*}F \\
 & e \arrow[ur] &
\end{tikzcd}
\qquad\qquad \text{.}
\tag*{(1)}
\]
Dans le cas $A$ abélienne, ce diagramme prend la forme d'une suite exacte
\[
0 \longrightarrow j_{!}j^{*}F \longrightarrow F \longrightarrow i_{*}i^{*}F \longrightarrow 0 \quad .
\]

\textit{e}) Mêmes conditions pour $i$ et $j$. On suppose que $A$ a les $\varprojlim$ pertinentes, dont objet final. On a alors un carré cartésien fonctoriel en $F\in \Ob\ A(S)$ \texttt{[il faut supposer la catégorie} $A$ \texttt{ponctuée]}
\[
\begin{tikzcd}
F \arrow[d, "\text{adj}"'] & i_{*}i^{!}F \arrow[l] \arrow[d] \\
j_{*}j^{*}F & e \arrow[l]
\end{tikzcd}
\tag*{(2)}
\]
\textit{i.e.} $i^{!}F$ s'identifie par le foncteur pleinement fidèle $i_{*}$ au \textit{noyau} de $F\to j_{*}j^{*}F$. Dans le cas abélien, cela prend la forme d'une suite exacte
\[
0 \longrightarrow i_{*}i^{!}F \longrightarrow F \longrightarrow j_{*}j^{*}F \quad .
\]

\dgrblock{dgrb3}{Support, zero extension, and the exceptional inverse image}
\begin{dcomment}{exposition}{p2:Dans le cas abélien, cela prend la forme}{Support, zero extension, and the exceptional inverse image}
Let $U\subset S$ be a sieve and $Y=S\setminus U$ its complementary cosieve. A presheaf on $Y$ extends by the terminal object on $U$ under $i_*$. Indeed, $Y/u$ is empty for $u\in U$, whereas restriction to $Y$ recovers the original presheaf. Conversely, if $F|_U$ is terminal, the unit $F\to i_*i^*F$ is an isomorphism at every object. This proves the description of the essential image in (d). Dually, $j_!$ extends a presheaf on $U$ by the initial object on $Y$.

Assume now that $A$ has a zero object and the required limits. Form
\[
 K=F\times_{j_*j^*F}0.
\]
Its restriction to $U$ is zero, so $K=i_*G$ for a unique $G$ up to the canonical equivalence above. For $H\in A(Y)$, the maps $i_*H\to j_*j^*F$ are in bijection with maps $j^*i_*H=0\to j^*F$, hence form a singleton. The universal property of the pullback gives
\[
 \operatorname{Hom}(H,G)\cong\operatorname{Hom}(i_*H,K)
 \cong\operatorname{Hom}(i_*H,F).
\]
Thus $G=i^!F$, proving both the extra adjunction and square (2). This argument explains the editors' pointedness correction: a terminal object alone supplies no canonical map from it to $j_*j^*F$. For $A=\mathsf{Set}$ and $S=[1]$, the closed inclusion of $1$ has $i_*B=(B\to *)$. It does not preserve the initial object, so cannot have a right adjoint.

The dual construction uses
\[
 Q=F\amalg_{i_!i^*F}0.
\]
Now $Q|_Y=0$, hence $Q=j_!j^?F$, and the same universal-property argument proves $j^?\dashv j_!$. The first square of (d) is bicartesian for a different, simpler reason: at each point it is a square with an identity on one side and a zero object on the other. Checking at points proves both universal properties. No stability hypothesis is needed for this particular square. In an abelian category its associated short exact sequence need not split as a sequence of diagrams, even though every component is split. For $S=[1]$ and $F=(\mathbb Z\xrightarrow{1}\mathbb Z)$, the sequence is
\[
 0\longrightarrow(0\to\mathbb Z)\longrightarrow
 (\mathbb Z\xrightarrow{1}\mathbb Z)\longrightarrow
 (\mathbb Z\to0)\longrightarrow0.
\]
A splitting would require the identity map on the first copy of $\mathbb Z$ to have zero composite with the displayed identity arrow, which is impossible.
\end{dcomment}


\clearpage

% ===================== page 5 =====================
\printedpage{3}
% head: A. Grothendieck : Dérivateurs, Ch. I, Édition des Universités de Montpellier II et Paris VII
% foot: Version 8 juillet 2001

\msp{page 3}

Je me dispense d'écrire le carré cocartésien dual, quand $A$ a les $\varinjlim$ pertinentes, y compris objet initial. Il semble qu'elle n'ait pas servi en pratique. Les foncteurs adjoints et biadjoints classiquement utiles sont inclus ci-dessous (où les “pas utiles” sont entre crochets)
\[
\begin{aligned}
{}[i_{!}]\ i^{*}\ i_{*}\ i^{!} \\
[j^{?}]\ j_{!}\ j^{*}\ j_{*}
\end{aligned}
\]
\textit{i.e.} les foncteurs $i_{!}$ et $j^{?}$ n'ont pas servi classiquement.

\textit{f}) Cas des catégories $\bDelta_{i}$ (simplexes type, $i\in\mathbb{N}$), et des foncteurs entre ceux-ci (applications semi-simpliciales). Il semblerait que pour une telle $f$, si $A$ est ponctuée, il y ait les \textit{bi-adjoints} à la fois à gauche et à droite. De plus, pour $\bDelta_{i}$ et $\bDelta_{j}$ fixés quand on regarde toutes les $f:\bDelta_{i}\to\bDelta_{j}$ et toutes les $g:\bDelta_{j}\to\bDelta_{i}$, que les suites totales de foncteurs multi-adjoints pour chacun des $f^{*}$, $g^{*}$ correspondants sont toutes les mêmes \texttt{[ce n'est pas vrai]}. Ça a l'air vrai du moins si $j=i+1$ (et je l'ai vérifié pour $i=0$, $i=1$), on doit trouver une “suite d'adjonctions” de longueur $2i+7=(i+1)+(i+2)+2+2$. De façon précise, la suite comporte quatre termes “extrêmes” (les deux premiers, les deux derniers) et les $(i+1)+(i+2)$ termes intermédiaires correspondant aux foncteurs de la forme $f^{*}$ ou $g^{*}$, où $f:\bDelta_{i+1}\to\bDelta_{i}$ est une des $(i+1)$ applications de dégénérescence, et $g:\bDelta_{i}\to\bDelta_{i+1}$ une des $i+2$ applications bord \texttt{[ce n'est pas tout à fait vrai; il faut se limiter aux applications qui sont injectives ou surjectives]}.

% ---- from der1_p006-010.tex
\printedpage{4}
% head: A. Grothendieck : Dérivateurs, Ch. I, Édition des Universités de Montpellier II et Paris VII
% foot: Version 8 juillet 2001

\msp{page 4}

\dgrblock{dgrb4}{The simplex adjunctions, with their endpoints}
\begin{dcomment}{examples}{p3}{The simplex adjunctions, with their endpoints}
Write $C_n=A^{[n]^{\mathrm{op}}}$, so its objects are strings $X_n\to\cdots\to X_0$. Let $\delta^k:[n]\to[n+1]$ omit $k$, and let $\sigma^k:[n+1]\to[n]$ identify $k$ and $k+1$. The inequalities in the finite ordinals give
\[
 \delta^{k+1}\dashv\sigma^k\dashv\delta^k.
\]
Applying contravariant presheaves reverses both arrows and natural transformations; consequently the restrictions form the adjunction chain
\[
 (\delta^{n+1})^*\dashv(\sigma^n)^*\dashv
 (\delta^n)^*\dashv\cdots\dashv
 (\sigma^0)^*\dashv(\delta^0)^*.
\]
One can also check each bijection of Hom sets by writing down the compatible component maps. There are $2n+3$ terms here.

If $A$ is pointed with finite limits and colimits, two more terms occur at either end. On the left, let $L:C_n\to C_{n+1}$ prepend zero at the high end: $L(X)=(0\to X_n\to\cdots\to X_0)$. Then $L\dashv(\delta^{n+1})^*$. Define $Q:C_{n+1}\to C_n$ by
\[
 (QX)_r=X_r\amalg_{X_{n+1}}0\qquad(0\leq r\leq n).
\]
A map $X\to LY$ is exactly a compatible family $X_r\to Y_r$ killing $X_{n+1}$, so $Q\dashv L$. At the right end let $R(X)=(X_n\to\cdots\to X_0\to0)$. Then $(\delta^0)^*\dashv R$. Its right adjoint is the string $K(X)_r=X_{r+1}\times_{X_0}0$, since a map $RY\to X$ is exactly a family of maps to those fibers. The full chain has $2n+7$ terms, as suggested in the manuscript.

For $n=0$ it is particularly transparent:
\[
 \operatorname{coker}\dashv(0\to-)\dashv
 \operatorname{ev}_0\dashv\operatorname{const}\dashv
 \operatorname{ev}_1\dashv(-\to0)\dashv\operatorname{ker}.
\]
Here kernel and cokernel mean fiber and cofiber over the zero object, so abelian structure is unnecessary. This also proves the first conjecture in (f): every monotone map of finite ordinals factors into elementary faces and degeneracies, each restriction occurring far enough inside one of these chains to have two adjoints on both sides. Composing the corresponding adjunctions supplies $f^?$ and $f^!$.

The stronger statement about one common chain for \emph{all} maps between two fixed ordinals is false. Already for $[1]$ and itself, the identity restriction has only identity adjoints up to isomorphism, whereas the restriction along a constant map is not the identity on $C_1$, for example when $A=\mathsf{Ab}$. Between $[n]$ and $[n+1]$, the central chain displayed above contains exactly the injective faces and surjective degeneracies. Other monotone maps are composites, not additional entries in that chain. The editors' qualifications should therefore be retained; the nearby count is a correct statement about the specified chain, not a classification of every restriction functor.
\end{dcomment}


\section*{2. Préliminaires heuristiques sur les dérivateurs.}

On se donne une sous-catégorie pleine $\Dia$ (comme “diagrammes”) de $\Cat$. Ça peut être tout $\Cat$, ou plus petit. Mais on supposera le plus souvent que $\Dia$ contient au moins les “catégories ordonnées” (correspondant aux ensembles ordonnés, \textit{i.e.} telles que $a$) tout morphisme est un mono et $b$) tout iso est une identité \texttt{[ce n'est pas vrai : la catégorie réduite à une double flèche satisfait à ces conditions, mais n'est pas un ensemble ordonné]}). Par abus de langage, on identifie un ensemble ordonné avec la catégorie correspondante (et une catégorie au topos qu'elle définit?). On considère $\Dia$ comme une 2-\textit{catégorie}.\footnote{N.B. Il faut des conditions de stabilité pour $\Dia$, au moins la stabilité par $I \times J$ et $I \amalg J$, en fait, nettement plus \dots}

\dgrblock{dgrb5}{From derived coefficients to diagram categories}
\begin{dcomment}{context}{p4:On se donne une sous-catégorie pleine}{From derived coefficients to diagram categories}
The opening calculations make the geometric language of supports available before any homotopy theory is introduced. There is a documented reason for this starting point. In \emph{Pursuing Stacks}, \S69, Grothendieck traces his concern to the inadequacy of a single derived category for the trace and determinant formulas arising around SGA~5. The proposed replacement assigns a derived category to every diagram shape; filtered complexes and diagrams of triangles are instances of this larger construction. The demand that the same structure reappear at every shape already anticipates the shifted derivators of manuscript p.~18 \citep[\S69, pp.~1--3 of the edited extract]{Grothendieck1983}.

His letter to Thomason of 2 April 1991 returns to the same question: which intrinsic structure survives passage from models to categories of coefficients? That letter explicitly connects the renewed work to the 1983 notes and discusses both derivators and the role of small categories as homotopy models \citep[pp.~2--4 of the edited letter]{Grothendieck1991Letter}. These are statements by Grothendieck about his program; they do not establish a retrospective priority claim for every later theorem. Heller's independently formulated homotopy theories already organized homotopy categories of diagrams, homotopy Kan extensions, and pointed operations in a related axiomatic framework \citep[Chs.~II--III]{Heller1988}.

The terminology on the next pages must be read provisionally. The superscript $\circ$ on $\Dia$ reverses both functors and transformations. A transformation $\alpha:f\Rightarrow g$ induces $\alpha^*:g^*\Rightarrow f^*$ because a presheaf sends $\alpha_i:f(i)\to g(i)$ to $F(g(i))\to F(f(i))$. In modern notation the input 2-category would therefore be $\Dia^{\mathrm{coop}}$. A common modern convention uses covariant diagrams and writes $\mathscr D:\mathsf{Cat}^{\mathrm{op}}\to\mathsf{CAT}$, reversing functors but retaining the direction of 2-cells. The translation is $\mathscr D(I)=\DD(I^{\mathrm{op}})$. A left derivator in this manuscript has right Kan extensions; this is the terminology of left exactness, not that of a left adjoint. The preliminary labels $\mathrm{Der}_1$, $\mathrm{Der}_2$, and $\mathrm{Der}_3$ will also be reassigned when the systematic axioms begin.

The small correction about posets is substantive but completely local. A category is a poset precisely when it is thin and skeletal: each Hom set has at most one member, and isomorphic objects are equal. Having only monomorphisms and no nonidentity isomorphisms does not imply thinness. The parallel-pair category printed in the editors' note is a counterexample, since every composable arrow that could test cancellation is an identity. No later construction depends on the incorrect characterization.
\end{dcomment}


Un \textit{dérivateur} $\DD$ (sur $\Dia$, comme “\textit{domaine}” du dérivateur) est un 2-foncteur contravariant de $\Dia$ dans $\Cat$, satisfaisant certaines conditions. Donc
\[ \DD : \Dia^{\circ} \longrightarrow \Cat \quad . \]
Si $I \in \Ob_0 \Dia$ est une catégorie (ou un “type de diagramme”) du domaine de $\DD$, on doit donc avoir une catégorie $\DD(I)$. Les objets de $\DD(I)$ s'appelleront les \textit{objets-coefficients} (ou \textit{coefficients}) de type $\DD$, au dessus de $I$. Pour
\[ f : I \longrightarrow J \]
une 1-flèche de $\Dia$, on doit avoir un foncteur
\[ f^*_{\DD} \text{ ou } f^* : \DD(J) \longrightarrow \DD(I) \quad , \]
dit \textit{“foncteur image inverse”}. Si $F$ est un coefficient de type $\DD$ sur $J$ son “\textit{image inverse}” est donc un coefficient de type $\DD$ sur $I$. On suppose (pour simplifier!) que le 2-foncteur $\DD$ est \textit{strict}, donc qu'on a transitivité stricte
\[ (gf)^* = f^*g^* \ , \qquad \text{si} \qquad I \xrightarrow{f} J \xrightarrow{g} K \qquad \text{dans} \ \ \Dia \ . \]

\msp{page 5}

De plus, si on a deux foncteurs $f, g : I \rightrightarrows J$, et un morphisme entre eux
\[
\begin{tikzcd}[column sep=large]
I \arrow[r, shift left=2, "f"] \arrow[r, shift right=2, "g"'] \arrow[r, phantom, "{\downarrow}u"] & J
\end{tikzcd}
\quad ,
\]
on se donne un morphisme $u^*_{\DD}$ ou $u^*$ entre les foncteurs $g^*$ et $f^*$, d'où
\[
\begin{tikzcd}[column sep=large]
\DD(I) & \DD(J) \arrow[l, shift right=2, "f^*"'] \arrow[l, shift left=2, "g^*"] \arrow[l, phantom, "{\uparrow}u^*"]
\end{tikzcd}
\quad .
\]

\printedpage{5}
% head: A. Grothendieck : Dérivateurs, Ch. I, Édition des Universités de Montpellier II et Paris VII
% foot: Version 8 juillet 2001

On doit encore avoir transitivité pour une composition de morphismes de foncteurs :
\[
\begin{tikzcd}[column sep=large]
I \arrow[r, shift left=4, "f"] \arrow[r, "g"' near start] \arrow[r, shift right=4, "h"'] \arrow[r, phantom, shift left=2, "{\downarrow}u"] \arrow[r, phantom, shift right=2, "{\downarrow}v"] & J
\end{tikzcd}
\]
doit donner
\[ (vu)^* = u^*v^* \]
donc
\[
\begin{tikzcd}[column sep=large]
\DD(I) & \DD(J) \arrow[l, shift right=4, "f^*"'] \arrow[l, "g^*" near end] \arrow[l, shift left=4, "h^*"] \arrow[l, phantom, shift right=2, "{\uparrow}u^*"] \arrow[l, phantom, shift left=2, "{\uparrow}v^*"]
\end{tikzcd}
\quad .
\]
De plus, si $f$ est une identité, $f^*$ aussi
\[ (\id_I)^* = \id_{\DD(I)} \quad , \]
et de même si $u$ est une identité
\[ (\id_f)^* = \id_{f^*} \quad . \]

\msp{page 6}

J'indique tout de suite l'axiome le plus important sur un dérivateur (au-delà de la donnée d'un 2-contra foncteur strict) :

\textbf{Der}$_1$ : Si $f : I \to J$ dans $\Dia$, alors $f^*_{\DD} = f^* : \DD(J) \to \DD(I)$ a un adjoint à droite et un adjoint à gauche.

Ces foncteurs seront notés (du moins provisoirement) $f^{\DD}_*$ ou $f_*$ , et $f^{\DD}_!$ ou $f_!$ , soit
\[ f_! \, , \ f_* : \DD(I) \rightrightarrows \DD(J) \quad , \]
de sorte qu'on trouve une suite de \textit{trois} foncteurs adjoints entre $\DD(I)$ et $\DD(J)$
\[ f_! \, , \ f^*, \ f_* \quad . \]

N.B. Pour la plupart des développements on n'aura pas besoin de l'existence à la fois des $f_!$ et des $f_*$ . Un \textit{dérivateur à gauche} ne suppose que les $f_*$ (adjoints à droite des $f^*$), et l'existence de $f_!$ seulement quand $f$ est une \textit{immersion ouverte}. Dualement, un \textit{dérivateur à droite} ne suppose que les $f_!$ (adjoints à gauche des $f^*$), et $f_*$ seulement si $f$ est une \textit{immersion fermée}. Ce sont sans doute les dérivateurs à gauche qui sont les plus importants.

\msp{page 7}

\textbf{Der}$_2$ : Si $f : I \to J$ est dans $\Dia$, si c'est une immersion \textit{fermée}, le biadjoint à droite $f^!_{\DD}$ ou $f^!$ de $f^*$ existe. Si c'est une immersion ouverte, le biadjoint à gauche $f^?$ existe.\footnote{Il n'y a lieu d'exiger les biadjoints que pour un dérivateur ponctué. \textsf{HOT} n'y satisfait pas.} Dans
\printedpage{6}
% head: A. Grothendieck : Dérivateurs, Ch. I, Édition des Universités de Montpellier II et Paris VII
% foot: Version 8 juillet 2001
les deux cas, on veut que $f^*$ soit un foncteur de localisation, ce qui équivaut à dire que $f_*$ est pleinement fidèle, ou que $f_!$ est pleinement fidèle.\footnote{Il faut de plus caractériser l'image de $f_*$ si $f$ est une immersion fermée (les $F$ tels que $g^*F \overset{\sim}{\to} e_{\DD(U)}$) et celle de $f_!$ si c'est une immersion ouverte (les $F$ tels que $g^*F \simeq \emptyset_{\DD(Y)}$), où $g$ est l'immersion complémentaire.}

N.B. Dans le cas d'un dérivateur à gauche, on suppose seulement l'existence de $f^!$ si $f$ est une immersion fermée. Donc pour résumer, pour un tel dérivateur, on suppose l'existence de
\[
\left\{
\begin{aligned}
& f_* \text{ quel que soit } f : I \to J \text{ dans } \Dia \\
& f_! \text{ si } f \text{ est une immersion ouverte} \\
& f^! \text{ si } f \text{ est une immersion fermée (cas ponctué seulement).}
\end{aligned}
\right.
\]
Dualement, pour un dérivateur à droite, on suppose l'existence de
\[
\left\{
\begin{aligned}
& f_! \text{ quel que soit } f : I \to J \text{ dans } \Dia \\
& f_* \text{ si } f \text{ est une immersion fermée} \\
& f^? \text{ si } f \text{ est une immersion ouverte (cas ponctué seulement).}
\end{aligned}
\right.
\]
Je voudrais formuler des axiomes astucieux, pour le cas d'un couple
\[ Y \overset{i}{\hookrightarrow} I \overset{j}{\hookleftarrow} U \]
d'une immersion ouverte $j$ et d'une immersion fermée $i$ dans $\Dia$, complémentaires l'une de l'autre. D'abord une formulation si on suppose $\DD$ “\textit{ponctué}”, \textit{i.e.} les catégories $\DD(I)$ ponctuées, et les foncteurs $f^*$ compatibles avec les ponctuations. On veut “interpréter” le carré (1) page 2 :

\msp{page 8}

\dgrblock{dgrb6}{Coherent squares and what the axioms remember}
\begin{dcomment}{exposition}{p7}{Coherent squares and what the axioms remember}
The passage from a square in $\DD(I)$ to an object of $\DD(I\times[1]\times[1])$ is the first genuinely homotopical issue. The underlying-diagram functor discards part of the data. A commutative square in a homotopy category records equality of two homotopy classes; a coherent square also records the homotopy implementing that equality, together with whatever compatibility is required in a larger diagram. It is therefore the coherent object to which a homotopy pullback or pushout operation applies.

Put $\square=[1]\times[1]$ and let $W$ be the full subcategory obtained by deleting $(1,1)$. For presheaves, $W^{\mathrm{op}}$ is a cospan. A coherent square $X\in\DD(I\times\square)$ is cartesian if it is in the essential image of right Kan extension along $I\times W\hookrightarrow I\times\square$. Equivalently, its unit comparison with that right Kan extension is an isomorphism. The opposite corner and left Kan extension define cocartesian squares. These definitions use the universal property of an adjoint inside the whole family of diagram categories, rather than an ordinary pullback in $\DD(I)$. In the represented derivator of a complete category they reduce to ordinary cartesian and cocartesian squares.

The desire to lift diagrams does not say that every diagram in $\DD(I)$ has a preferred lift. A frequently imposed additional axiom, strongness, asks that the underlying arrow functor $\DD(I\times[1])\to\DD(I)^{[1]^{\mathrm{op}}}$ be full and essentially surjective. Even that says neither that the lift is unique nor that the functor is faithful. The special support squares have more structure: they arise from restrictions, adjunction units or counits, and zero extensions. Those operations give the coherent constructions that an arbitrary incoherent square lacks.

In the modern full derivator axioms, pointedness can be required only at $\DD(e)$. It then propagates to all diagram categories, and the exceptional adjoints required here follow; this is proved in \citet[Proposition~3.2 and Corollary~3.8]{Groth2013}. Thus the manuscript's extra-adjoint requirement was a useful way to expose the intended geometry, but is not an additional independent condition on an already pointed derivator. This does not justify imposing exceptional adjoints on the unpointed derivator of spaces: the initial and terminal spaces differ. The qualification in footnote~2 is necessary.
\end{dcomment}


\[
\begin{tikzcd}
& F \arrow[dr, "\mathrm{adj}"] & \\
j_!j^*F \arrow[ur, "\mathrm{adj}"] \arrow[dr] & & i_*i^*F \\
& e \arrow[ur] &
\end{tikzcd}
\quad ,
\tag*{(1)}
\]

\textbf{Der}$'_2$ : Pour tous $I, J \in \Dia$, on a
\[ \DD(I \amalg J) \xrightarrow[(\alpha^*, \beta^*)]{\sim} \DD(I) \times \DD(J) \quad , \]
où $\alpha, \beta$ sont les injections canoniques de $I$ et $J$ dans $I \amalg J$.

\printedpage{7}
% head: A. Grothendieck : Dérivateurs, Ch. I, Édition des Universités de Montpellier II et Paris VII
% foot: Version 8 juillet 2001

carré à la fois cartésien et cocartésien. On cherche donc un carré dans $\DD(I)$, mais mieux que ça. Soit $J = \Del_1 \times \Del_1$ le carré standard
\[
\begin{tikzcd}
& (1,0) \arrow[dl] & \\
(1,1) & & (0,0) \arrow[ul] \arrow[dl] \\
& (0,1) \arrow[ul] &
\end{tikzcd}
\]
et considérons la catégorie $I \times J$. On a
\[
\begin{aligned}
& J \longrightarrow \Hom(I, I \times J) \\
& j \longmapsto \bigl(i \longmapsto (i,j)\bigr)
\end{aligned}
\]
D'autre part, par $\DD$
\[ \Hom(I, I \times J)^{\circ} \longrightarrow \Hom(\DD(I \times J), \DD(I)) \]
d'où un composé
\[ J^{\circ} \longrightarrow \Hom(\DD(I \times J), \DD(I)) \]
ou encore
\[ \DD(I \times J) \longrightarrow \Hom(J^{\circ}, \DD(I)) \quad . \]
Ceci vaut pour tout $J \in \Ob \, \Dia$. On peut donc dire que la donnée d'un “coefficient” $\Phi$ sur une catégorie produit $I \times J$ donne naissance à un diagramme de type $J^{\circ}$ dans la catégorie $\DD(I)$ des coefficients sur $I$. La philosophie
\msp{page 9}
des dérivateurs consiste en grande partie, inversement, chaque fois qu'on a un diagramme dans une catégorie de coefficients $\DD(I)$ (sur un diagramme $I$), disons diagramme de type $K$ : $K \to \DD(I)$, de se demander s'il ne provient pas “canoniquement” d'un objet de $\DD(I \times K^{\circ})$ \textit{i.e.} d'un coefficient sur $I \times K^{\circ}$.

Ainsi pour le carré p. 8. On voudrait qu'il provienne de $\DD(I \times J)$, où $J = \Del_1 \times \Del_1$, et ceci fonctoriellement, donc on désire un foncteur “canonique”
\[ \DD(I) \longrightarrow \DD(I \times \Del_1 \times \Del_1) \]
qui en composant avec
\[ \DD(I \times \Del_1 \times \Del_1) \longrightarrow \Hom((\Del_1 \times \Del_1)^{\circ}, \DD(I)) \]
\printedpage{8}
% head: A. Grothendieck : Dérivateurs, Ch. I, Édition des Universités de Montpellier II et Paris VII
% foot: Version 8 juillet 2001
donne le foncteur décrit par le carré (1). De plus, il faut exprimer dans le contexte présent la condition “bicartésienne” du carré (1), dans le contexte de la page 7. Question analogue pour le carré (2) de la page 2 :
\[
\begin{tikzcd}
& F \arrow[dr] & \\
i_*i^!F \arrow[ur] \arrow[dr] & & j_*j^*F \\
& e \arrow[ur] &
\end{tikzcd}
\quad ,
\tag*{(2)}
\]
où il faut de plus interpréter la condition cartésienne. On va définir en fait un coefficient sur la sous-catégorie ouverte $W$ suivante de $\Del_1 \times \Del_1$
\[
\begin{tikzcd}
(1,0) & \\
& (0,0) \arrow[ul] \arrow[dl] \\
(0,1) &
\end{tikzcd}
\quad ,
\]

\msp{page 10}

plus exactement sur $W \times I$, et si $\alpha$
\[ \alpha : W \times I \longrightarrow \Del_1 \times \Del_1 \times I \]
est l'inclusion canonique, on définira le carré “cartésien” comme l'image par $\alpha_*$ de ce coefficient. Ceci nous ramène donc au problème de “globaliser” d'abord le diagramme partiel
\[
\begin{tikzcd}
F \arrow[dr] & \\
& j_*j^*F \\
e \arrow[ur] &
\end{tikzcd}
\quad .
\]
Et de même pour le carré (1) sauf qu'on est embarrassé si on choisit de le regarder surtout comme cartésien (ce qui suggère à procéder comme tantôt) ou comme \underline{co}cartésien (ce qui % verified transcription correction: the edition prints embarrassé, with double r

% ---- from der1_p011-015.tex
\printedpage{9}
% head: A. Grothendieck : Dérivateurs, Ch. I, Édition des Universités de Montpellier II et Paris VII
\noindent suggère de procéder de façon duale). Mais les deux coefficients qu’on trouvera par ces deux procédés, produits soit de
\[
\begin{tikzcd}
F \arrow[dr] & \\
 & i_*i^*F \\
e \arrow[ur] &
\end{tikzcd}
\quad ,
\]
soit de
\[
\begin{tikzcd}
 & F \\
j_!j^*F \arrow[ur] \arrow[dr] & \\
 & e
\end{tikzcd}
\quad ,
\]
devraient être \textit{canoniquement} isomorphes.

\noindent\msp{page 11}Si $\DD$ est un dérivateur, la catégorie
\[
\mathcal{A} = A_{\DD} = \DD(e) % edition reading verified: plain italic A_D is retained; this alphabet variation is not established to be an error
\]
joue un rôle important. On l’appelle la \textit{catégorie de base,} ou \textit{la catégorie fondamentale}, du dérivateur, ces objets sont les \textit{coefficients fondamentaux.} Pour $I$ quelconque dans $\Dia$, on a donc pour $F \in \DD(I)$, coefficient sur $I$ % [sic] ces objets -> ses objets, possessive intended
\[
I^\circ \longrightarrow \mathcal{A} \quad , \qquad i \longmapsto F_i
\]
un diagramme de type $I^\circ$ dans $\mathcal{A}$, fonctoriellement, \textit{i.e.}
\[
\DD(I)\longrightarrow
\uHom(I^\circ,\underset{\substack{\Vert\\\DD(e)}}{\mathcal A})
\]
% the vertical equality sign is printed under the letter \mathcal{A} of \uHom(I^\circ,\mathcal{A}), with \DD(e) beneath it
(\textit{cf.} page 8).

\noindent\textbf{Der}$_3$ : Pour tout $I \in \Dia$, ce foncteur est \textit{conservatif,}\footnote[4]{N.B. Mais ce foncteur n’a aucune envie d’être \textit{fidèle.}} \textit{i.e.} pour que
\[
F \longrightarrow G
\]
\printedpage{10}%
% head: A. Grothendieck : Dérivateurs, Ch. I, Édition des Universités de Montpellier II et Paris VII
dans $\DD(I)$ soit \texttt{[un]} isomorphisme, il faut et il suffit que pour tout $i \in I$,
\[
F_i \longrightarrow G_i
\]
le soit.

\noindent\textbf{Corollaire.} \textit{Le foncteur canonique (pour $I$, $J \in \Dia$)}
\[
\DD(I \times J) \xrightarrow{\ \theta\ } \uHom(J^\circ, \DD(I))
\]
\textit{est conservatif.}

En effet, si $F$ est dans $\DD(I \times J)$, on a pour $(i,j) \in I \times J$
\[
F_{i,j} \simeq \theta(F)(j)_i \quad .
\]

\noindent\msp{page 12}Intuitivement, on regarde un dérivateur $\DD$ comme “étant” la catégorie fondamentale $\mathcal{A}$, avec une “structure” supplémentaire, \texttt{[à]} savoir la donnée des “extensions” de $\mathcal{A}$ (comme catégorie des coefficients ponctuels) aux diagrammes $I \in \Dia$, de façon contravariante en $I$. Un coefficient $F$ dans $\DD(I) \defeq \mathcal{A}_I$ est regardé comme “étant” le diagramme $(F_i)_{i\in I}$ correspondant d’objets fondamentaux, avec une donnée supplémentaire de “globalisation”. Les $F_i$, ou “fibres” de $F$ en les $i \in I$, sont les caractéristiques locaux, du coefficient “global sur $I$” $F$. % [sic] caractéristiques locaux -> caractéristiques locales, agreement (feminine)

\dgrblock{dgrb7}{Conservative is not faithful}
\begin{dcomment}{examples}{p11}{Conservative is not faithful}
The conservativity axiom detects whether a \emph{given coherent morphism} is invertible. It does not reconstruct that morphism from its components. A small derived example makes the distinction exact. In the abelian category of arrows $X_1\to X_0$, put
\[
 P=(\mathbb Z\to0),\qquad Q=(0\to\mathbb Z),\qquad
 E=(\mathbb Z\xrightarrow{1}\mathbb Z).
\]
The nonsplit sequence
\[
 0\longrightarrow Q\longrightarrow E\longrightarrow P\longrightarrow0
\]
defines a nonzero element of $\operatorname{Ext}^1(P,Q)$, and hence a nonzero morphism $P\to Q[1]$ in the derived category of arrows. Its two evaluated components are $\mathbb Z\to0[1]$ and $0\to\mathbb Z[1]$, both zero. This is the same nonsplitting calculated in the support example, now viewed as a morphism in a derived category. It exhibits exactly the information lost by
\[
 D(\mathsf{Ab}^{[1]^{\mathrm{op}}})
 \longrightarrow D(\mathsf{Ab})^{[1]^{\mathrm{op}}}.
\]
There is no contradiction with conservativity. The nonzero extension class is not a morphism all of whose components are isomorphisms.

The corollary for partial underlying diagrams is proved in the text, not left as a gap. If a coherent morphism over $I\times J$ becomes an isomorphism at every $j$ in $\DD(I)$, evaluating again at every $i$ gives an isomorphism at each $(i,j)$. The original conservativity axiom on $I\times J$ now applies. In particular, a shifted family retains the same test for invertibility. This short argument does not need strongness, full faithfulness, a model structure, or a choice of representatives.
\end{dcomment}


Tous les dérivateurs connus sont obtenus de la façon standard suivante. On part d’une “catégorie de modèles”
\[
\mathcal{M} \quad , \quad \Sigma
\]
munie d’un ensemble de flèches, appelées “équivalences”, et satisfaisant aux axiomes pertinents pour le passage à une catégorie des fractions
\[
\mathcal{A} = \mathcal{M}\Sigma^{-1}
\]
qui sera la “catégorie fondamentale”. On définit pour tout $I$ dans $\Dia$
\[
\begin{aligned}
\mathcal{M}_I &= \uHom(I^\circ,\mathcal{M}) \\
\Sigma_I &= \{u \in \mathsf{Fl}(\mathcal{M}_I) \mid \forall i \in I, u_i \in \Sigma\} \\
\mathcal{A}_I &= \mathcal{M}_I\Sigma_I^{-1}
\end{aligned}
\]

\noindent\msp{page 13}Il est clair que $\mathcal{A}_I \defeq \DD_{\mathcal{M},\Sigma}(I)$ dépend de $I$ de façon fonctorielle contravariante, et même de façon 2-fonctorielle.

Bien entendu, pour avoir vraiment un dérivateur, il faut vérifier les axiomes. Le plus fort est l’existence des $f_!$, $f_*$. Elle est prouvée par Thomason lorsque $(\mathcal{M},\Sigma)$ provient d’une “catégorie de modèles” au sens de Quillen. Ce n’est pas clair pour moi qu’alors $(\mathcal{M}_I,\Sigma_I)$ provienne également d’une catégorie de modèles (dont par exemple les fibrations \printedpage{11}%
% head: A. Grothendieck : Dérivateurs, Ch. I, Édition des Universités de Montpellier II et Paris VII
et les cofibrations seraient définies à partir de ces données dans $\mathcal{M}$, comme tantôt pour les équivalences).

Quand on prend pour $(\mathcal{M},\Sigma)$ une catégorie de modèles pour $\mathsf{Hot}$ (par exemple les ensembles semi-simpliciaux $\mathcal{M} = \uHom(\bDelta^\circ,\Ens)$, d’où $\mathcal{M}_I \simeq \uHom(\bDelta^\circ \times I^\circ,\Ens) \simeq \uHom(\bDelta^\circ,\underbrace{\uHom(I^\circ,\Ens)}_{\widehat{I}})$, catégorie des objets semi-simpliciaux dans le topos $\widehat{I}$, on trouve une description d’un dérivateur (non additif), qui jouera un rôle particulièrement important. Je le désigne par % Resolved: exactly two closing parentheses in the formula; the source's outer prose parenthesis remains unclosed.
\[
\mathsf{HOT} \quad ,
\]
sa catégorie des coefficients fondamentaux est $\mathsf{Hot}$, la catégorie homotopique standard.

\dgrblock{dgrb8}{Localizing diagrams and deriving their Kan extensions}
\begin{dcomment}{later}{p13:sa catégorie des coefficients fondamentaux}{Localizing diagrams and deriving their Kan extensions}
For a relative category $(\mathcal M,W)$, objectwise weak equivalences are preserved by every restriction. Thus localization supplies the displayed prederivator, subject to the usual size convention ensuring that its Hom classes are sets. Naturality of the localized transformations follows from the universal property of localization. One may use a functorial localization construction to keep restriction strictly functorial. Arbitrary independent choices of localized categories instead give canonically equivalent data with coherence isomorphisms.

The difficult assertion is the existence and pointwise behavior of the derived Kan extensions. Grothendieck attributes the model-category case to Thomason here, without specifying a publication. We have not identified a published Thomason theorem with exactly the generality of this attribution. A precise later source is \citet[Theorem~6.11]{Cisinski2003}: a Quillen model category admitting small limits gives a left derivator in the presheaf convention. Applying the result to the opposite model category gives the dual assertion when small colimits exist. The proof also establishes local smallness of the diagram localizations. This settles the mathematical existence question without requiring an ordinary Quillen model structure on every diagram category.

For a combinatorial model category there is a familiar more direct route. The projective diagram model structure has objectwise fibrations and weak equivalences; the injective structure has objectwise cofibrations and weak equivalences. Their homotopy categories are the same localization, although their replacements differ. Projective replacements derive left Kan extension, while injective replacements derive right Kan extension. In general one cannot declare \emph{both} fibrations and cofibrations objectwise and expect the lifting axioms. The passage from models to presheaves is treated in a more flexible formalism by \citet[Theorem~6.22 and Corollary~6.24]{Cisinski2010}.

For the indicated model of spaces, the exponential law gives
\[
 (\mathsf{sSet})^{I^{\mathrm{op}}}
 \cong\mathsf{Set}^{\Delta^{\mathrm{op}}\times I^{\mathrm{op}}}
 \cong(\mathsf{Set}^{I^{\mathrm{op}}})^{\Delta^{\mathrm{op}}}.
\]
The weak equivalences in this construction are evaluated at each $i$. They are not a newly imposed local weak-equivalence relation for an arbitrary Grothendieck topology. Also, ``semi-simplicial'' is historical terminology here: the displayed $\Delta$ includes degeneracies and gives what is now called a simplicial set. Replacing it silently by the category of injective order maps would change the model category under discussion.
\end{dcomment}


\noindent\msp{page 14}La notion de “morphisme de Dérivateurs” est claire, c’est un morphisme entre deux 2-foncteurs stricts
\[
\Dia^\circ \longrightarrow \mathcal{C}at
\]
Donc
\[
u : \DD \longrightarrow \DD'
\]
est défini par la donnée, pour tout $I \in \Dia$, de
\[
u_I : \DD(I) \longrightarrow \DD'(I)
\]
avec commutativité (stricte?) pour $f : I \to J$
\[
\begin{tikzcd}
\DD(J) \arrow[r, "f^*_{\DD}"] \arrow[d, "u_J"'] & \DD(I) \arrow[d, "u_I"] \\
\DD'(J) \arrow[r, "f^*_{\DD'}"] & \DD'(I)
\end{tikzcd}
\quad . \tag*{(1)}
\]
On a
\[
\begin{tikzcd}[column sep=large]
\uHom(I,J)^\circ \arrow[r, "\DD_{I,J}"] \arrow[d, "\DD'_{I,J}"'] & \uHom(\DD(J),\DD(I)) \arrow[d, "\phi \mapsto u_I \circ \phi"] \\
\uHom(\DD'(J),\DD'(I)) \arrow[r, "\psi \mapsto \psi \circ u_J"'] & \uHom(\DD(J),\DD'(I))
\end{tikzcd}
\tag*{(2)}
\]
et la commutativité de (1) signifie la commutativité de (2) \textit{pour les objets}. Mais bien sûr on la veut aussi pour les flèches, \textit{i.e.} on veut la commutativité du diagramme (2).\footnote[5]{Donner l’exemple de morphisme de dérivateurs défini par un foncteur de catégories de modèles $(\mathcal{M},\Sigma) \to (\mathcal{M}',\Sigma')$.}

Sans doute on veut un peu plus sur les morphismes de dérivateurs - on veut une commutation aux $f_*$ et (?) aux $f_!$ (cette fois seulement à isomorphisme canonique près) \ldots\footnote[6]{Commutation à $f_*$ (resp. $f_!$) morphisme \textit{exact à gauche} (resp. \textit{exact à droite}). Les morphismes \textit{exacts} sont ceux à la fois exacts à gauche et à droite.}

\dgrblock{dgrb9}{Morphisms and their comparison maps}
\begin{dcomment}{exposition}{p14:Sans doute on veut un peu plus}{Morphisms and their comparison maps}
A morphism between the entire families of coefficient categories must respect restriction, including its dependence on natural transformations of indexing functors. The square (2) is the condition that is lost if one checks only the individual functors in square (1). In modern treatments one normally allows coherent invertible comparison maps
\[
 \gamma_f:f^*_{\DD'}u_J\xrightarrow{\sim}u_If^*_{\DD}
\]
rather than demanding strict equality. The coherence conditions concern identity functors, composites, and natural transformations. Modifications between such pseudonatural morphisms provide the 2-cells. A functor of relative categories preserving weak equivalences gives an example by applying it at every component and then localizing; no preservation of Kan extensions follows merely from this construction.

Adjunction does give canonical comparison maps, with precisely the orientations on manuscript pp.~15--16. For the right Kan extension the map is the composite
\[
\begin{aligned}
 u_Jf_*^{\DD}
 &\longrightarrow f_*^{\DD'}f^{*}_{\DD'}u_Jf_*^{\DD}\\
 &\xrightarrow{\ f_*^{\DD'}\gamma_f f_*^{\DD}\ }
 f_*^{\DD'}u_If^{*}_{\DD}f_*^{\DD}
 \longrightarrow f_*^{\DD'}u_I.
\end{aligned}
\]
The outer arrows are the unit in $\DD'$ and the counit in $\DD$. The left Kan comparison $f_!^{\DD'}u_I\to u_Jf_!^{\DD}$ is the adjoint of
\[
 u_I\longrightarrow u_If^*_{\DD}f_!^{\DD}
 \xrightarrow{\gamma_f^{-1}}f^*_{\DD'}u_Jf_!^{\DD}.
\]
An invertible restriction comparison need not have an invertible mate. Preservation of limits or colimits is the additional assertion that the corresponding mates are invertible. These definitions also explain which comparison is meant; an unspecified isomorphism of the two functors is not sufficient.

In the manuscript, left and right exactness at this point mean preservation of all the relevant right and left Kan extensions. Modern ``exact morphism of stable derivators'' usually refers to finite homotopy limits and colimits. Stability makes the two finite notions coincide; it does not equate preservation of arbitrary products and coproducts. For example, derived tensor product with $\mathbb Q$ is exact on $D(\mathbb Z)$ and preserves coproducts. The map
\[
 \mathbb Q\otimes\prod_{n\geq1}\mathbb Z
 \longrightarrow\prod_{n\geq1}\mathbb Q
\]
is not surjective: an element on the left has a common denominator, whereas $(1/n!)_{n\geq1}$ does not. The footnote about triangulated categories is correct in its finite setting and should not be promoted to an assertion about arbitrary Kan extensions. For model categories, a right Quillen adjoint yields a derived morphism preserving homotopy limits, with the dual statement for a left Quillen adjoint; \citet[Proposition~6.12 and its dual]{Cisinski2003} provide the relevant construction.
\end{dcomment}


\printedpage{12}
% head: A. Grothendieck : Dérivateurs, Ch. I, Édition des Universités de Montpellier II et Paris VII
\noindent\msp{page 15}Par exemple, considérons le dérivateur
\[
\mathsf{HOTAB}
\]
dont la catégorie fondamentale est $\mathsf{Hotab}$, la sous-catégorie pleine de la catégorie dérivée $\mathsf{Der}(\mathsf{Ab})$ de $\mathsf{Ab}$, formée des complexes à cohomologie nulle en degré $> 0$. On peut décrire $\mathsf{HOTAB}$ par la catégorie de modèles $\uHom(\bDelta^\circ,\mathsf{Ab})$ des groupes abéliens semi-simpliciaux, avec comme équivalences les quasi-isomorphismes au sens habituel. Le “foncteur abélianisation” canonique
\[
\mathsf{Hot} \longrightarrow \mathsf{Hotab}
\]
se prolonge en un morphisme de dérivateurs
\[
\mathsf{HOT} \longrightarrow \mathsf{HOTAB} \quad .
\]
La question de sa compatibilité aux $f_*$ ne semble pas tautologique. Elle impliquerait, par exemple, que si
\[
\begin{tikzcd}
X \arrow[d, "f"'] & X' \arrow[l, "g"'] \arrow[d, "f'"] \\
Y & Y' \arrow[l, "j"]
\end{tikzcd}
\]
est un diagramme cartésien dans $\mathcal{M} = \uHom(\bDelta^\circ,\Ens)$ avec $f$ une fibration (de Thom dans le cadre des espaces topologiques, avec $f$ par exemple fibration localement triviale) alors le diagramme de $\mathsf{Hotab}$
\[
\begin{tikzcd}
\mathbb{R}\Gamma_!(X) \arrow[d] & \mathbb{R}\Gamma_!(X') \arrow[l] \arrow[d] \\
\mathbb{R}\Gamma_!(Y) & \mathbb{R}\Gamma_!(Y') \arrow[l]
\end{tikzcd}
\]
est $\mathsf{hot}$-cartésien\footnote[7]{C’est faux déjà si $Y$ est l’espace ponctuel : trivialement le $\mathbb{R}\Gamma_!$ ne transforme pas produits en sommes! Ce foncteur est sûrement \texttt{[pas]} “exact à gauche” (comme morphisme de dérivateurs).}. Et de même pour une $\mathsf{hot}$-somme amalgamée (carrée $\mathsf{hot}$-cocartésien dans $\mathsf{Hot}$). % [sic] carrée -> carré, agreement (masculine noun)

\noindent\msp{page 16}Si c’est vrai, et pas connu de plus, ce serait quand même scandaleux!

\dgrblock{dgrb10}{Abelianization preserves homotopy colimits}
\begin{dcomment}{examples}{p15}{Abelianization preserves homotopy colimits}
The model for $\mathsf{HOTAB}$ is simplicial abelian groups. Normalization identifies these with nonnegative chain complexes, or equivalently cochain complexes in nonpositive degrees; see \citet[Theorem~14.24.3, Tag~019D]{StacksDoldKan}. Its homotopy category is $D^{\leq0}(\mathsf{Ab})$. This is an additive, pointed theory, but it is not stable: suspension need not be invertible inside the connective category. Consequently one must not infer preservation of homotopy limits from an additive appearance.

The morphism from spaces is induced by the degreewise free abelian group functor $\mathbb Z[-]$. It is left adjoint to the forgetful functor on simplicial objects. The latter preserves fibrations and trivial fibrations in the standard model structures, so this is a Quillen adjunction. Its derived left adjoint preserves homotopy colimits, including the homotopy pushouts in the last sentence of manuscript p.~15. That sentence therefore has a positive answer, although the preceding pullback assertion does not. Under normalization, this is the usual chain construction, with its coherent diagram version.

A finite discrete counterexample makes the failure of the pullback assertion explicit. Let $X=Y'=\{0,1\}$ and $Y=*$, so $X'=X\times Y'$ has four points. All objects are Kan complexes and the maps to the point are fibrations. After abelianization the square has upper left and lower right terms $\mathbb Z^2$, lower left term $\mathbb Z$, and upper right term $\mathbb Z^4$. Each map $\mathbb Z^2\to\mathbb Z$ is augmentation. Since these maps are surjective, their ordinary pullback computes their homotopy pullback in simplicial abelian groups, and it is
\[
 \{(a_0,a_1,b_0,b_1)\in\mathbb Z^4:
 a_0+a_1=b_0+b_1\}\cong\mathbb Z^3.
\]
The comparison from the chains of $X'$ sends the basis vector $e_{ij}$ to $(e_i,e_j)$. Its kernel is generated by
\[
 e_{00}-e_{01}-e_{10}+e_{11}.
\]
It is therefore not an isomorphism. This proves the correction in footnote~7 while keeping separate the valid colimit statement. Already the terminal-object test would rule out continuity, since $\mathbb Z[*]=\mathbb Z$ and the terminal simplicial abelian group is zero, but the displayed example addresses exactly the nonempty square that Grothendieck is discussing.
\end{dcomment}


En procédant comme pour les morphismes de changements de base pour les $Rf_*$, on trouve en tout cas un morphisme canonique, pour $f : I \to J$
\[
u_Jf_*^{\DD} \longrightarrow f_*^{\DD'}u_I
\]
\printedpage{13}%
% head: A. Grothendieck : Dérivateurs, Ch. I, Édition des Universités de Montpellier II et Paris VII
et de même
\[
f_!^{\DD'}u_I \longrightarrow u_Jf_!^{\DD} \quad ,
\]
et la question est de savoir si c’est des isomorphismes. Questions analogues pour les $f^!$, pour $f$ immersion fermée (ou pour $f^?$, pour $f$ immersion ouverte). Comme la notion de morphisme de dérivateurs devrait quand même “coller” à la notion bien connue de morphisme (triangulé) de catégories triangulées\footnote[8]{Non! Le cas des catégories triangulées est “misleading”, car pour celles-ci l’exactitude à gauche et à droite revient au même.}, donc être exact en un sens convenable, cela montre bien que l’on doit exiger la commutation d’un morphisme de dérivateurs aux $f_*$ et aux $f_!$.\footnote[9]{N.B. On conçoit qu’un foncteur entre dérivateurs puisse être exact à gauche ou exact à droite ou les deux.} Je devrais demander à Thomason s’il a regardé, pour un morphisme
\[
\mathcal{M} \longrightarrow \mathcal{M}'
\]
entre catégories de modèles, qui soit disons \textit{exact à gauche}, et “compatible” avec les équivalences, fibrations (et cofibrations (???)).

\noindent\textsc{Notation} : $\mathsf{DER}$ pour la 2-catégorie des \textit{dérivateurs.}

\noindent\msp{page 17}\textbf{Dualité}. Si $\DD$ est un dérivateur, on va définir un dérivateur \textit{dual} $\DD^\circ$ par
\[
\DD^\circ(I) \defeq \DD(I^\circ)^\circ \quad .
\]
Les coefficients sur $I$ de type $\DD^\circ$ sont les coefficient sur la catégorie opposée $I^\circ$ à $I$ de type $\DD$, et en tant que catégorie c’est la catégorie \textit{opposée} de celle des coefficients de type $\DD$ sur $I^\circ$. Pour définir la loi 2-fonctorielle, on note que pour $I$, $J$ dans $\Dia$ on a % [sic] les coefficient -> les coefficients, plural agreement
\[
\begin{tikzcd}[column sep=large]
\uHom(I,J)^\circ \arrow[r, "\DD_{I,J}"] \arrow[d, "\wr"', "\text{dualité}"] & \uHom(\DD(J),\DD(I)) \arrow[d, "\wr"', "\text{dualité}"] \\
\uHom(I^\circ,J^\circ) \arrow[r, "\theta"] & \uHom(\DD(J)^\circ,\DD(I)^\circ)^\circ
\end{tikzcd}
\quad ,
\]
où $\theta$ est l’unique foncteur rendant le diagramme commutatif. Alors $\theta^\circ$est un foncteur
\[
\uHom(I^\circ,J^\circ)^\circ \longrightarrow \uHom(\DD(J)^\circ,\DD(I)^\circ) \quad ,
\]
c’est ça par définition $\DD^\circ_{I,J}$ \texttt{[plutôt} $\DD^\circ_{I^\circ,J^\circ}$\texttt{]}.

Si $\DD$ est défini par $(\mathcal{M},\Sigma)$, alors $\DD^\circ$ par $(\mathcal{M}^\circ,\Sigma^\circ)$, puisque
\[
(\mathcal{M}^\circ)_I \defeq \uHom(I^\circ,\mathcal{M}^\circ) \simeq \uHom(I,\mathcal{M})^\circ \simeq (\mathcal{M}_{I^\circ})^\circ
\]
% foot (pages 9--13): Version 8 juillet 2001

% ---- from der1_p016-020.tex
\printedpage{14}
% head: A. Grothendieck : Dérivateurs, Ch. I, Édition des Universités de Montpellier II et Paris VII
% foot: Version 8 juillet 2001

et $\Sigma^\circ{}_I$ est déduit de $\Sigma_{I^\circ}$, d'où
\[
(\mathcal{M}^\circ)_I(\Sigma^\circ{}_I)^{-1} \simeq (\mathcal{M}_{I^\circ}\Sigma_{I^\circ}^{-1})^\circ = \D(I^\circ)^\circ \quad .
\]

\msp{page 18}

\dgrblock{dgrb11}{What the whole diagram theory remembers}
\begin{dcomment}{later}{p17}{What the whole diagram theory remembers}
A modern comparison makes precise the extra information carried by the variable $I$. For an $\infty$-category $\mathcal C$, define, in the convention of this manuscript,
\[
\DD_{\mathcal C}(I)=h\operatorname{Fun}(N(I)^{\mathrm{op}},\mathcal C).
\]
Here $N$ is the nerve and $h$ takes the homotopy category. One first forms coherent diagrams and then passes to homotopy; reversing these operations would give $\operatorname{Fun}(I^{\mathrm{op}},h\mathcal C)$, which is generally different. Restriction and its natural transformations supply the prederivator structure. If $\mathcal C$ has all small limits and colimits, its pointwise Kan extensions supply the two adjoints and the comma-square formula. \citet[\S5.3, especially Proposition~5.3.9 and Remark~5.3.10]{RiehlVerity2017} explain this passage using the homotopy $2$-category of an $\infty$-cosmos. Their convention uses covariant diagrams; replacing $I$ by $I^{\mathrm{op}}$ yields the formula above.

The formula also clarifies the philosophical remark on manuscript page~18. Passing to $\DD(e)=h\mathcal C$ discards most of the indexing variable. Retaining the entire prederivator can retain much more. \citet[Theorem~19 and Corollary~20]{Arlin2020} proves a fully faithful embedding using strict transformations and the specified simplicial enrichment. With pseudonatural transformations, his Theorem~29 gives a bicategorically fully faithful embedding for \emph{small} quasicategories when the domain contains all small categories. The size hypothesis, the domain, and the kind of transformation belong to the theorem; none can be silently omitted. In the model-category setting, \citet[Theorems~3.3.2 and~3.4.4]{Renaudin2009} obtains the corresponding local equivalence, and then a biequivalence with derivators of small presentation, from left proper combinatorial model categories localized at Quillen equivalences, with adjunctions as morphisms.

These results concern the entire diagram theory. They do not assert that every functor $h\mathcal C\to h\mathcal E$ admits a coherent lift, or that a lift is determined by its value at $e$. Nor does an objectwise description of a prederivator by localized diagram categories dispense with the proof of its Kan-extension axioms. The comparison is useful precisely because it keeps three questions separate: existence of coherent diagrams, existence of homotopy Kan extensions, and reconstruction of an enhanced category from all its diagram categories. Grothendieck's successive shifts $\DD_I$ belong to this richer level of information.
\end{dcomment}


\textbf{Philosophie.} Chaque fois qu'on a un foncteur
\[
\mathcal{A} \xrightarrow{u_0} \mathcal{A}'
\]
entre deux catégories qui peuvent s'interpréter naturellement comme les catégories fondamentales de deux dérivateurs $\D$, $\D'$, il se pose la question d'interpréter $u_0$ comme provenant d'un morphisme de dérivateurs
\[
\D \xrightarrow{u} \D' \quad .
\]
Celui-ci est, en pratique, toujours déterminé canoniquement, à isomorphisme canonique près. (Mais je doute pourtant que le foncteur restriction aux catégories fondamentales
\[
\uHom(\D,\D') \longrightarrow \uHom(\mathcal{A},\mathcal{A}')
\]
soit pleinement fidèle \dots\ )

\dgrblock{dgrb12}{A negative answer to the evaluation question}
\begin{dcomment}{examples}{p18:soit pleinement fidèle}{A negative answer to the evaluation question}
For unrestricted morphisms of derivators, the doubt about full faithfulness can be answered within the elementary examples already at hand. Let $A=[1]$ and let $\mathscr R(I)=\operatorname{Fun}(I^{\mathrm{op}},A)$. Since the two-element ordered set has all small meets and joins, this represented prederivator is a derivator: its Kan extensions are the ordinary pointwise ones. The represented construction is standard; compare \citet[Examples~1.2(1) and~1.6]{Groth2013}.

For any prederivator $\mathscr E$ in our convention and any $X\in\mathscr E(A^{\mathrm{op}})$, define a morphism $U_X:\mathscr R\to\mathscr E$ as follows. An object $F\in\mathscr R(I)$ is a functor $F:I^{\mathrm{op}}\to A$, so $F^{\mathrm{op}}:I\to A^{\mathrm{op}}$, and we set
\[
 (U_X)_I(F)=(F^{\mathrm{op}})^*X.
\]
If $F\Rightarrow G$ is a natural transformation, taking opposites gives $G^{\mathrm{op}}\Rightarrow F^{\mathrm{op}}$; the reversal of 2-cells by $\mathscr E$ gives the required map $(U_X)_I(F)\to(U_X)_I(G)$. The composition and restriction laws follow directly from those of $\mathscr E$. A map $m:X\to Y$ likewise induces a modification $U_m:U_X\Rightarrow U_Y$ with component $(F^{\mathrm{op}})^*m$ at $F$.

Now take $\mathscr E$ to be the unbounded derived derivator of abelian groups. Its category $\mathscr E(A^{\mathrm{op}})$ is the derived category of arrows $0\to1$. As in the preceding nonfaithfulness example, the extension
\[
 0\to(0\to\mathbb Z)\to
 (\mathbb Z\xrightarrow{1}\mathbb Z)\to(\mathbb Z\to0)\to0
\]
yields a nonzero map $m:(\mathbb Z\to0)\to(0\to\mathbb Z)[1]$ with both evaluated components zero. Therefore $(U_m)_e$ is the zero natural transformation between the base-category functors. Nevertheless $U_m$ itself is nonzero: at $I=A^{\mathrm{op}}$ and $F=\operatorname{id}_A$, its value is exactly $m$. Evaluation at $e$ is thus not faithful, and consequently not fully faithful.

This settles the question for the usual category of all pseudonatural morphisms, and already for the strictly natural morphisms just constructed. It does not assert the same counterexample for a subcategory in which every morphism must preserve all left or right Kan extensions. The manuscript is still negotiating that convention. The proof identifies the precise obstruction in the unrestricted case: an entire coherent arrow can carry a morphism which vanishes at each of its two endpoints.
\end{dcomment}


On va appliquer cette philosophie aux foncteurs $f^*$, $f_!$, $f_*$ (et éventuellement $f^!$, $f^?$), pour $f : I \to J$, donc pour les foncteurs entre $\D(I)$ et $\D(J)$. Cela nous amène donc à regarder $\D(I)$ et $\D(J)$ comme des catégories fondamentales de dérivateurs convenables.

\textbf{Dérivateur induit} par un dérivateur $\D$, sur une catégorie d'indices $I \in \Dia$. C'est un dérivateur $\D_I$ défini par
\[
\D_I(J) = \D(I \times J) \quad .
\]
Le foncteur associé à $\D_I$, pour deux catégories $J$, $J' \in \Dia$

\msp{page 19}
\[
\uHom(J,J')^\circ \xrightarrow{(\D_I)_{J,J'}} \uHom(\D_I(J'),\D_I(J)) = \uHom(\D(I\times J'),\D(I\times J))
\]
est le composé
\[
\uHom(J,J')^\circ \longrightarrow \uHom(I\times J, I\times J')^\circ \xrightarrow{\D_{I\times J,I\times J'}} \uHom(\D(I\times J'),\D(I\times J)) \quad ,
\]
où la première flèche est le foncteur
\[
f \longmapsto \id_I \times f \quad .
\]
En d'autres termes, $\D_I$ est le composé
\[
\begin{aligned}
&\Dia^\circ \longrightarrow \Dia^\circ \xrightarrow{\D} \Cat \\
&J \longmapsto I \times J \qquad .
\end{aligned}
\]
On a
\[
f^*_{\D_I} = (\id \times f)^*_{\D} \quad .
\]
Si les $f_*$ (resp. les $f_!$) existent pour $\D$, ils existent pour $\D_I$. (De même pour toute espèce de multiadjoints.) Comme $J \mapsto I \times J$ transforme immersion fermée (resp. ouverte) en itou, l'axiome \textbf{Der}$_2$ (plus précisément, chacun de ces ingrédients, \textit{cf.} annotations marginales p.~7) est stable par passage de $\D$ à $\D_I$ . De même pour \textbf{Der}$_3$. (Pour chaque nouvel axiome qu'on dégagera, il faudra vérifier la stabilité.)

\printedpage{15}
% head: A. Grothendieck : Dérivateurs, Ch. I, Édition des Universités de Montpellier II et Paris VII
% foot: Version 8 juillet 2001

\msp{page 20}

\dgrblock{dgrb13}{Duality and the change of coefficient category}
\begin{dcomment}{exposition}{p19}{Duality and the change of coefficient category}
The two opposites in $\DD^{\mathrm{op}}(I)=\DD(I^{\mathrm{op}})^{\mathrm{op}}$ have distinct roles. The first reverses the direction of the indexing arrows; the second reverses arrows between coefficient objects. Both are needed to obtain another prederivator with the same variance convention. In a model presentation this is the canonical identification
\[
 \operatorname{Fun}(I^{\mathrm{op}},\mathcal M^{\mathrm{op}})
 \cong\operatorname{Fun}(I,\mathcal M)^{\mathrm{op}},
\]
followed by localization. Localization commutes with taking opposites by its universal property: a functor inverting the selected arrows on either side is exactly the opposite of such a functor on the other side. Thus the dual model statement is proved by the displayed argument. The incorrect indices in the definition of the dual action on Hom categories are already corrected by the editors.

The shift $\DD_I(J)=\DD(I\times J)$ changes which objects count as coefficients. An entire coherent $I$-diagram now becomes a single coefficient, and a coherent $J$-diagram of such coefficients is an object over $I\times J$. At the level of a relative category this is the elementary exponential law
\[
 (\mathcal M^{I^{\mathrm{op}}})^{J^{\mathrm{op}}}
 \cong\mathcal M^{(I\times J)^{\mathrm{op}}},
\]
with exactly the same objectwise weak equivalences. This is why the shift is intrinsic to the proposed language: it permits an argument about one coefficient to be reapplied to a family with parameters.

The preliminary stability checks on p.~19 are valid for the axioms already written. Products preserve sieves and cosieves. Adjoints to $(\operatorname{id}_I\times f)^*$ are supplied by the corresponding adjoints in $\DD$, and the conservativity argument uses evaluation at each $(i,j)$. The full pointwise Kan-extension axiom still needs a separate verification. In the completed theory this is the shift theorem, \citet[Theorem~1.25]{Groth2013}; it is not an unproved additional assumption about an arbitrary family.

Likewise the candidate components $(f\times\operatorname{id}_J)_!$ and $(f\times\operatorname{id}_J)_*$ on p.~20 do not constitute a morphism of derivators merely by being written down. One must show that their comparison maps for changes of $J$ are invertible and satisfy coherence. The product base-change squares supply this, and the coherence follows from the calculus of mates. Restriction is strictly natural before choices of identifications, whereas adjoint extensions are naturally pseudonatural. The following pages are devoted to precisely these exchange maps. It is appropriate to call the present construction a sketch awaiting those checks, not an error in the definition of the shift.
\end{dcomment}


Soit
\[
f : I \longrightarrow I' \quad .
\]
Considérons entre $\D(I)$, $\D(I')$ les foncteurs
\[
\begin{tikzcd}[column sep=huge]
\D(I) \arrow[r, shift left=3ex, "f_!"] \arrow[r, shift right=3ex, "f_*"'] & \D(I') \arrow[l, "f^*"']
\end{tikzcd}
\]
(sous réserve d'existence des $f_*$ resp. $f_!$). Peut on les interpréter en termes de morphismes de dérivateurs
\[
\begin{tikzcd}[column sep=huge]
\D_I \arrow[r, shift left=3ex, "\overline{f_!}"] \arrow[r, shift right=3ex, "\overline{f_*}"'] & \D_{I'} \arrow[l, "\overline{f^*}"']
\end{tikzcd}
\qquad ?
\]
Donc il faut définir, pour $J$ dans $\Dia$, des foncteurs correspondants entre $\D_I(J)$ et $\D_{I'}(J)$, \textit{i.e.} entre $\D(I \times J)$ et $\D(I' \times J)$. Or on a bien
\[
f \times \id_J : I \times J \longrightarrow I' \times J
\]
et on posera (sous réserve d'existence des deux derniers)
\[
\left\{
\begin{aligned}
&\overline{f^*}(J) = (f \times \id_J)^*_{\D} \\[1ex]
&\overline{f_!}(J) = (f \times \id_J)^{\D}_! \\[1ex]
&\overline{f_*}(J) = (f \times \id_J)^{\D}_* \quad .
\end{aligned}
\right.
\]

\msp{page 21}

Autrement dit, on peut considérer
\[
\Dia \ni I \longmapsto \D_I \quad , \qquad \D \in \DER \quad ,
\]
comme fonctoriel en $I$ de trois façons (une contravariante, les deux autres covariantes), suivant qu'on utilise les $f^*$, les $f_!$ (s'ils existent) ou les $f_*$ (s'ils existent). On désignera par
\[
f^{\D}_! \quad , \qquad f^*_{\D} \quad , \qquad f^{\D}_*
\]
indifféremment les foncteurs entre $\D(I)$ et $\D(I')$ ou (de façon plus fine) entre $\D_I$ et $\D_{I'}$, dont les précédents se déduisent en passant aux catégories fondamentales correspondantes. De même pour $f^!$ ou pour $f^?$ sous réserve d'existence. Par exemple, si on suppose que les $f^!$ existent pour $\D$, quel que soit $f$ immersion fermée, alors on trouve pour un tel $f$ un morphisme de dérivateurs
\[
\D_I \xleftarrow{f^!_{\D}} \D_{I'} \quad .
\]
Il s'impose de regarder si tous ces morphismes de dérivateurs sont exacts, \textit{i.e.} si pour
\[
g : J \longrightarrow J'
\]
dans $\Dia$, ils commutent à $g_!$, $g_*$.

\printedpage{16}
% head: A. Grothendieck : Dérivateurs, Ch. I, Édition des Universités de Montpellier II et Paris VII
% foot: Version 8 juillet 2001

\msp{page 22 vide}

\msp{page 23}

Pour expliciter cette question, on est amenés à regarder le diagramme de catégories
\[
\begin{tikzcd}[column sep=huge, row sep=huge]
I \times J \arrow[r, "G = \id_I \times g"] \arrow[d, "F = f \times \id_J"'] & I \times J' \arrow[d, "F' = f \times \id_{J'}"] \\
I' \times J \arrow[r, "G' = \id_{I'} \times g"] & I' \times J'
\end{tikzcd}
\qquad ,
\]
qui est d'ailleurs un carré cartésien dans $\Cat$. La commutativité de $f^*_{\D}$ aux $g_*$ s'explicite alors par la commutativité dans
\[
\begin{tikzcd}[column sep=huge, row sep=huge]
\D(I \times J) \arrow[r, "G_*", ""{name=T, below, near end}] & \D(I \times J') \\
\D(I' \times J) \arrow[u, "F^*"] \arrow[r, "G'_*"] & \D(I' \times J') \arrow[u, "{F'}^*"', ""{name=R, left}]
\arrow[dashed, from=R, to=T, "\text{iso?}"]
\end{tikzcd}
\qquad .
\]
Cette commutativité sera assurée par des axiomes ultérieures sur les dérivateurs. (cf. plus % [sic] ultérieures -> ultérieurs, accord avec « axiomes »
bas page 59). Ainsi, $f^*_{\D}$ est exact à droite \texttt{[plutôt à gauche]}, et dualement, il doit être aussi exact à gauche \texttt{[plutôt à droite]}. Donc $f^*_D : \D_{I'} \to \D_I$ est \textit{exact.} Considérons $f^{\D}_*$, % [sic] f^*_D -> f^*_{\mathbb{D}}, indice imprimé en D italique au lieu de \mathbb{D}
cette fois c'est la commutativité de
\[
\begin{tikzcd}[column sep=huge, row sep=huge]
\D(I \times J) \arrow[r, "G_*"] \arrow[d, "F_*"'] & \D(I \times J') \arrow[d, "F'_*"] \\
\D(I' \times J) \arrow[r, "G'_*"] & \D(I' \times J')
\end{tikzcd}
\qquad ,
\]
c'est trivial (à isomorphisme canonique près). Ainsi, on
\msp{page 24}
trouve que $f^{\D}_* : \D_I \to \D_{I'}$ est \textit{exact à gauche}, et dualement $f^{\D}_!$ est \textit{exact à droite.} Mais est-ce que $f^{\D}_*$ est aussi exact à \textit{droite}? Cela signifie donc commutativité de
\[
\begin{tikzcd}[column sep=huge, row sep=huge]
\D(I \times J) \arrow[r, "G_!"] \arrow[d, "F_*"'] & \D(I \times J') \arrow[d, "F'_*", ""{name=R2, left}] \\
\D(I' \times J) \arrow[r, "G'_!", ""{name=B, above, near end}] & \D(I' \times J')
\arrow[dashed, from=B, to=R2]
\end{tikzcd}
\qquad .
\]

\dgrblock{dgrb14}{Three different interchange questions}
\begin{dcomment}{exposition}{p23}{Three different interchange questions}
The extra variable in a shifted derivator is a parameter. An object of
$\mathbb D(I\times J)$ is therefore a coherent $I$-family of coherent
$J$-diagrams; it is not, in general, merely a functor from $I^{\circ}$ to
$\mathbb D(J)$. Keeping the two variables coherent makes the question on these
pages well posed. There are three operations to compare: restriction in one
variable with a Kan extension in the other; two right, or two left, Kan
extensions; and a right Kan extension with a left Kan extension. Only the last
comparison introduces a new exactness condition.

Here is the pointwise mechanism for the first comparison. For a projection
in the $J$-direction, the slice occurring in its right Kan formula at
$(i,j')$ has the form
\[
 (I/i)\times(g/j').
\]
The first factor has terminal object $(i,1_i)$. Since our diagrams are
contravariant, taking cohomology in that factor amounts to evaluating at
this terminal object. Consequently the value is the right Kan formula in
the $J$-variable applied to the $i$-th fibre. The same calculation with $i$
replaced by $f(i)$ gives the restriction comparison. Der~2 promotes these
pointwise isomorphisms to an isomorphism of coherent diagrams. The dual
argument treats left Kan extensions. This is the content of
\citet[Proposition~2.5]{Groth2013}, in the convention of covariant diagrams.

For the second comparison, no pointwise calculation is necessary.
Restriction around the product square agrees in either order. The two
composites of right adjoints are therefore right adjoint to the same
functor, and uniqueness of an adjoint gives their canonical isomorphism.
The left-adjoint assertion is identical. These are the derivator forms of
Fubini's rule for limits and for colimits.

For the mixed comparison there is still a canonical map, in the direction
\[
 G'_!F_*\longrightarrow F'_*G_!,
\]
but no formal argument makes it invertible. Its construction uses a unit,
the already established restriction comparison, and an adjunction. The
asymmetry of its direction is meaningful: it is the familiar comparison
from a colimit of limits to a limit of colimits. The source's provisional
word \emph{exact} means preservation of all the Kan extensions available
in the chosen domain. This is stronger than the finite exactness used in
modern treatments of stable derivators. The distinction is needed before
reading the subsequent appeal to triangulated categories.
\end{dcomment}


\printedpage{17}
% head: A. Grothendieck : Dérivateurs, Ch. I, Édition des Universités de Montpellier II et Paris VII
% foot: Version 8 juillet 2001

On arrive bien à définir
\[
G'_!F_* \xrightarrow{?} F'_*G_!
\]
ou ce qui revient au même par adjonction
\[
F_* \xrightarrow{?} {G'}^*F'_*G_!
\]
en utilisant le morphisme de changement de base
\[
{G'}^*F'_* \xrightarrow{\overset{\alpha}{\sim}} F_*G^*
\]
qui est \texttt{[un]} \textit{isomorphisme} (on l'a vu), d'où
\[
\begin{tikzcd}[column sep=large, row sep=normal]
& {G'}^*F'_*G_! \arrow[dd, "\wr"', "\alpha * G_!"] \\
F_* \arrow[ur, "?"] \arrow[dr, "F_* * \beta"'] & \\
& F_*G^*G_!
\end{tikzcd}
\]
où $\beta : \id \to G^*G_!$ est le morphisme d'adjonction, d'où le morphisme voulu. (On peut aussi chercher ${F'}^*(G'_!F_*) \to G_!$, et utiliser la commutation $G_!F^* \overset{\sim}{\to} {F'}^*G'_!$, \textit{i.e.} l'exactitude à \textit{droite} de $f^*_{\D}$.)

\msp{page 25}

Mais ce morphisme est-il un isomorphisme? Je pense que dans le cas additif ça peut être raisonnable, mais dans le cas non additif ça n'a aucune raison d'être vrai. On voit en tout cas que si pour tout $f : I \to I'$ dans $\Dia$, $f^{\D}_* : \D_I \to \D_{I'}$ est (non seulement exact à gauche mais aussi) exact à droite, alors pour tout $g : J \to J'$ dans $\Dia$, $g^{\D}_!$ est (non seulement exact à droite mais même) exacte à gauche. % [sic] exacte -> exact, accord avec « $g^{\D}_!$ est (non seulement exact\dots) »

Prenons le dérivateur des préfaisceaux d'ensembles sur des $I$ variables, sa catégorie fondamentale est $\Ens$. C'est un dérivateur non ponctué. Considérons $I \to I'$ une immersion fermée non surjective, $f_* : \D(I) \to \D(I')$. Si $f^{\D}_*$ était exact à droite $f_*$ devrait transformer objet initial de $\D(I)$ en objet initial de $\D(I')$, sauf erreur (\textit{cf.} plus bas). Or il n'en est rien, puisque pour tout $F \in \sIm\bigl(f^{\D}_* : \D(I) \to \D(I')\bigr)$ la restriction de $F$ à $I' \setminus I$ est le faisceau final $e_U = (i \mapsto e)$, et non le faisceau initial $\emptyset_U = (i \mapsto \emptyset)$.

On va prouver ceci :

\textbf{Proposition.} \textit{Si $\D$ est un dérivateur à gauche (à droite), alors les produits finis (sommes finies) existent dans les $\D(I)$. Tout morphisme $\varphi : \D \to \D'$ de dérivateurs qui est exact à gauche (à droite) commute aux dits produits (sommes), chaque $I$ donnant $\varphi(I) : \D(I) \to \D'(I)$. En particulier, il transforme objet final (initial) en itou.}

\printedpage{18}
% head: A. Grothendieck : Dérivateurs, Ch. I, Édition des Universités de Montpellier II et Paris VII
% foot: Version 8 juillet 2001

\msp{page 26}

\dgrblock{dgrb15}{The empty shape and finite products}
\begin{dcomment}{examples}{p25}{The empty shape and finite products}
The simplest test for a claim about preservation of all homotopy colimits
is the empty diagram. If $i:Y\hookrightarrow X$ is a proper closed
inclusion of index categories, then in the represented derivator of sets
$i_*$ extends a presheaf by the terminal set outside $Y$. Thus
$(i_*\varnothing_Y)_x=\{*\}$ for $x\notin Y$, whereas
$(\varnothing_X)_x=\varnothing$. A right adjoint can preserve every limit
and still fail this first colimit test.

The finite-product proposition uses a different feature of derivators:
discrete coherent diagrams carry no hidden coherence. For a finite set
$S$, Der~1 identifies
\[
 \mathbb D(I\times S)\simeq\prod_{s\in S}\mathbb D(I).
\]
Restriction along $\pi:I\times S\to I$ becomes the diagonal. Its right
adjoint $\pi_*$ is therefore the ordinary product in the category
$\mathbb D(I)$: for $(F_s)$ and $E$ the adjunction reads
\[
 \operatorname{Hom}(E,\pi_*(F_s))
 \cong\prod_{s\in S}\operatorname{Hom}(E,F_s).
\]
This proves the assertion including $S=\varnothing$, when the right-hand
set is a singleton. The left-adjoint calculation gives coproducts. A
morphism preserving the corresponding Kan extensions preserves these
ordinary products or coproducts because it identifies their representing
adjunctions. This argument explains exactly why the empty-diagram test is
available. The phrase ``adjoint \`a gauche'' at the beginning of manuscript
page~26 is an edition-preserved slip: for products it must be a right
adjoint, as the same paragraph subsequently says.
\end{dcomment}


Dire que les produits finis indexés par l'ensemble fini $I$ existent dans la catégorie $\mathcal{A} = \D(e)$ (disons) équivaut à dire que le foncteur diagonal
\[
\mathcal{A} \xrightarrow{\Delta} \mathcal{A}^I
\]
admet un adjoint à gauche. Or si
\[
f : I \to e
\]
est la projection de $I$ (considéré comme ensemble ordonné discret) sur l'ensemble ordonné ponctuel $e$, le foncteur $\Delta$ s'identifie à $f^* : \D(e) \to \D(I)$ (grâce à l'axiome \textbf{Der}$_2$ sur le comportement du dérivateur pour une somme discrète de catégories d'indices). Donc l'existence du produit revient à celle de l'adjoint à droite $f_*$ de $f^*$, OK pour dérivateur à gauche. Et un homomorphisme exact à gauche de dérivateurs à gauche, comme il commute aux $f_*$, commute aux produits finis. Itou pour les sommes finies, et les dérivateurs à droite.

L'exemple précédent utilise le fait que le dérivateur envisagé n'est pas ponctué. Mais même pour un dérivateur ponctué $\D$, un morphisme $f^{\D}_* : \D_I \to \D_{I'}$ n'est pas forcément exact à droite. Si on prend le dérivateur associé à une catégorie $\mathcal{M}$ quelconque (avec $\Sigma = \emptyset$),
\[
\D_{\mathcal{M}}(I) = \uHom(I^\circ,\mathcal{M}) \overset{\text{déf}}{=} \mathcal{M}(I)
\]
alors le dérivateur induit
\[
(\D_{\mathcal{M}})_I = \{J \mapsto \D_{\mathcal{M}}(I \times J) = \uHom((I \times J)^\circ,\mathcal{M}) \simeq \uHom(J^\circ,\uHom(I^\circ,\mathcal{M}))\}
\]

\msp{page 27}

s'identifie à
\[
\D_{\uHom(I^\circ,\mathcal{M})} = \D_{\mathcal{M}(I)}
\]
\textit{i.e.}
\[
(\D_{\mathcal{M}})_I = \D_{\mathcal{M}(I)}
\]
N.B. Ici $\mathcal{M}$ et $\D_{\mathcal{M}}$ se déterminent mutuellement et le foncteur
\[
\mathcal{M} \longmapsto \D_{\mathcal{M}}
\]
\[
\Cat \longrightarrow \DER
\]
est 2-fidèle, \textit{i.e.} si $\mathcal{M}$, $\mathcal{M}'$ sont deux catégories
\[
\uHom(\mathcal{M},\mathcal{M}') \xrightarrow{\simeq} \sHom(\D_{\mathcal{M}},\D_{\mathcal{M}'})
\]
est une équivalence de catégories. D'autre part, on vérifie tout de suite que $u : \mathcal{M} \to \mathcal{M}'$ est exact à gauche (ou à droite) si et seulement si $\D_u : \D_{\mathcal{M}} \to \D_{\mathcal{M}'}$ l'est.

Ainsi, pour voir si un morphisme de dérivateurs de la forme
\[
\begin{tikzcd}[column sep=large, row sep=normal]
f^{\D}_* : &[-2.5em] (\D_{\mathcal{M}})_I \arrow[r] \arrow[d, equal] & (\D_{\mathcal{M}})_{I'} \arrow[d, equal] \\
& \D_{\mathcal{M}(I)} & \D_{\mathcal{M}(I')}
\end{tikzcd}
\]
est exact à \textit{droite,} il revient au même de dire

% ---- from der1_p021-025.tex
\printedpage{19}
% head: A. Grothendieck : Dérivateurs, Ch. I, Édition des Universités de Montpellier II et Paris VII
% foot: Version 8 juillet 2001
\msp{page 28}
que le foncteur pour les catégories fondamentales associées
\[ \mathcal{M}(I) \xrightarrow{f^{\mathcal{M}}_*} \mathcal{M}(I') \]
est exact à droite. Or ici c'est un foncteur image directe de préfaisceaux ordinaires, et ce foncteur n'a aucune envie d'être exact à \textit{droite.} (Même si $\mathcal{M}$ est \textit{abélienne}, d'ailleurs!) Prenons par exemple pour $f : I \hookrightarrow I'$ une immersion ouverte avec $I' = \boldsymbol{\Delta}_1 \times \boldsymbol{\Delta}_1$, $I = I' - \{c\}$, $c$ objet final de $I'$
\[
I' = \begin{tikzcd}
c & b \arrow[l] \\
a \arrow[u] & 0 \arrow[l] \arrow[u]
\end{tikzcd}
\qquad\qquad
I = \begin{tikzcd}
 & b \\
a & 0 \arrow[l] \arrow[u]
\end{tikzcd}
\]
Le foncteur $f_*$ est celui qui associe à un diagramme
\[
F = \left\{
\begin{tikzcd}
 & F_b \arrow[d] \\
F_a \arrow[r] & F_0
\end{tikzcd}
\right\}
\]
le carré cartésien complété
\[
\overline{F} = \left\{
\begin{tikzcd}
F_a \times_{F_0} F_b \arrow[r] \arrow[d] & F_b \arrow[d, "\beta"] \\
F_a \arrow[r, "\alpha"] & F_0
\end{tikzcd}
\right\}
\]
Or si $F \mapsto \overline{F}$ était exact à droite, $F \mapsto F_a \times_{F_0} F_b$ le serait aussi. Or $F_a \times_{F_0} F_b = K(F) = \mathsf{Ker}(F_a \times F_b \xrightarrow{\alpha pr_1 - \beta pr_2} F_0)$, en tant que foncteur additif en $F$, n'est pas exact à droite, par exemple il ne transforme pas nécessairement épimorphisme en épimorphisme. Prenons par exemple
\[
F' = \left\{
\begin{tikzcd}
 & F_b \arrow[d] \\
F_a \arrow[r] & 0
\end{tikzcd}
\right\}
\]
et l'homomorphisme évident $F \to F'$ qui est l'identité sur
\msp{page 29}
$F_a$ et sur $F_b$. C'est un épimorphisme et l'homomorphisme correspondant $\overline{F} \to \overline{F'}$ est l'inclusion
\[ F_a \times_{F_0} F_b \longrightarrow F_a \times F_b \]
qui n'est un épimorphisme que si $\alpha = \beta = 0$.

\dgrblock{dgrb16}{What represented prederivators remember}
\begin{dcomment}{exposition}{p27}{What represented prederivators remember}
For a category $\mathcal M$, the assignment
$\mathbb D_{\mathcal M}(I)=\operatorname{Fun}(I^{\circ},\mathcal M)$ is
always a prederivator. The existence of the limits or colimits demanded by
the diagram domain is needed before calling it a derivator. Its shift is
represented by another ordinary category because currying gives
\[
 \operatorname{Fun}((I\times J)^{\circ},\mathcal M)
 \cong\operatorname{Fun}(J^{\circ},
             \operatorname{Fun}(I^{\circ},\mathcal M)).
\]
Both sides carry the same objects, maps and composition laws after
currying; the identification commutes with restriction in $J$ and hence
identifies the prederivators, rather than only their separate values.

The fully faithful assertion about categories and represented
prederivators can also be made explicit. A functor
$u:\mathcal M\to\mathcal M'$ acts by postcomposition at every shape.
Conversely, let $\Phi:\mathbb D_{\mathcal M}\to\mathbb D_{\mathcal M'}$
be a pseudonatural morphism and put $u=\Phi_e$. Its pseudonaturality at
$i:e\to I$ identifies the $i$-th value of $\Phi_I(F)$ with $u(F_i)$.
For an arrow $i\to j$, the corresponding natural transformation between
these point functors makes these identifications commute with the
presheaf arrow $F_j\to F_i$. They therefore form a natural isomorphism
$\Phi_I(F)\cong uF$. Naturality in $F$ and in changes of shape follows
from the coherence conditions for $\Phi$. A modification between two
such morphisms is determined, by the same evaluation argument, by its
component at $e$, and every natural transformation $u\Rightarrow v$
extends by postcomposition. This proves the asserted equivalence of
hom-categories. With strict morphisms one uses strict naturality instead.

There is a useful qualification to the next assertion. Preservation of
\emph{all relevant} limits by $u$ implies preservation of right Kan
extensions because these are calculated pointwise by those limits. The
converse follows by taking Kan extensions to the point. If the domain
contains infinite categories, ordinary finite left exactness of $u$ is
not enough. The same qualification applies dually to colimits. Thus the
represented example detects a genuine obstruction to the manuscript's
strong exactness condition; it does not show that homotopical exactness
has failed to remember an elementary property of $\mathcal M$.
\end{dcomment}


\printedpage{20}
% head: A. Grothendieck : Dérivateurs, Ch. I, Édition des Universités de Montpellier II et Paris VII
% foot: Version 8 juillet 2001
\textbf{Perplexité} : On devrait pouvoir définir une notion de dérivateur (additif) “\textit{triangulé}” (qui correspondrait à l'intuition des catégories triangulées; c'est un dérivateur satisfaisant à des axiomes assez forts), de sorte que pour les morphismes entre ceux-ci l'exactitude à gauche et l'exactitude à droite soient \textit{équivalentes.}

\textbf{“Théorème de changement de base”} pour dérivateurs. Soit $\DD$ dérivateur à gauche (existence de $f_*$) et considérons
\[ f : I \longrightarrow J \qquad \text{dans } \sDia \quad , \]
soit $j_0 \in \sOb\ J$, on se propose de “calculer” le foncteur
\[ F \mapsto f_*(F)_{j_0} = \alpha^*_{j_0}(f_*(F)) \quad , \]
où
\[ \alpha_{j_0} : e \longrightarrow J \]
est l'inclusion de la catégorie ponctuelle $e$ dans $J$, de valeur $j_0$.

\msp{page 30}
Considérons le diagramme
\[
\begin{tikzcd}[row sep=large, column sep=large]
I/j_0 \arrow[r, "q"] \arrow[d, "f'"'] & I \arrow[d, "f"] \\
e \arrow[r, "p = \alpha_{j_0}"'] & J
\arrow[from=1-2, to=2-1, dashed, shorten <=14pt, shorten >=14pt, "\gamma"']
\end{tikzcd}
\quad ,
\tag*{(1)}
\]
où $I/j_0$ est la catégorie $I \times_J J/j_0$ des couples
\[ \{(i,\varphi) \mid i \in \sOb\ I,\ \varphi : f(i) \to j_0\} \]
et $q$ est le foncteur $(i,\varphi) \mapsto i$. Ce diagramme n'est pas commutatif, mais il y a un foncteur \texttt{[plutôt morphisme de foncteurs]} canonique
\[ fq \xrightarrow{\gamma} pf' \quad , \]
où $fq$ est le foncteur $(i,\varphi) \mapsto f(i)$ et $pf'$ le foncteur constant $(i,\varphi) \mapsto j_0$ et où $\gamma(i,\varphi) : f(i) \to j_0$ est défini par $\gamma(i,\varphi) = \varphi$.

\dgrblock{dgrb17}{The pullback obstruction, and the finite stable repair}
\begin{dcomment}{examples}{p29}{The pullback obstruction, and the finite stable repair}
The abelian counterexample can be made numerical. Let the coefficient
category be abelian groups and choose
\[
 F_a=F_0=\mathbb Z,\qquad F_b=0,\qquad
 \alpha=1_{\mathbb Z},\quad\beta=0.
\]
Then $K(F)=\mathbb Z\times_{\mathbb Z}0=0$. Replace the central value
$F_0$ by zero to obtain $F'$. The map $F\to F'$ is an epimorphism at
each object of the indexing cospan, but $K(F')=\mathbb Z$, and the
induced map of missing corners is $0\to\mathbb Z$. Hence the right Kan
extension fails to preserve epimorphisms and cannot be right exact even
in the ordinary finite abelian sense. Every group and every map in this
example is visible in the original three-object shape.

The ensuing ``Perplexit\'e'' points toward stability, with a necessary
change in the meaning of exactness. A pointed derivator is stable when
its coherent cartesian and cocartesian squares coincide. Between stable
derivators a morphism preserving zero objects and cartesian squares
therefore preserves cocartesian squares as well, and conversely. This
settles the finite version of the question; see
\citet[Definition~3.15 and Corollary~4.17]{Groth2013}. Additivity alone
cannot imply stability. In a represented pointed derivator the pushout
of $0\leftarrow A\to0$ is always zero, so its suspension functor is
zero. If that derivator were stable, suspension would be an equivalence,
forcing every object to be zero.

Even stability does not make arbitrary limits commute with arbitrary
colimits. Work in the derivator of complexes over a nonzero field $k$,
and regard vector spaces as complexes in degree zero. On the discrete
shape $\mathbb N\times\mathbb N$ put
\[
 X_{nm}=\begin{cases}k,&n=m,\\0,&n\ne m.\end{cases}
\]
Products and coproducts of vector spaces are exact, so here their derived
values are their ordinary values. The mixed comparison is
\[
 \bigoplus_m\prod_nX_{nm}=k^{(\mathbb N)}
 \longrightarrow
 \prod_n\bigoplus_mX_{nm}=k^{\mathbb N}.
\]
It is the inclusion of finitely supported sequences; the sequence
$(1,1,1,\ldots)$ is absent from its image. This disproves the unrestricted
interpretation of the anticipated stable exactness theorem.

The precise positive result is stronger than merely commuting a finite
product with a finite sum: in a stable derivator homotopy finite limits
commute with arbitrary homotopy colimits, and dually. A homotopy finite
category is equivalent to a finite skeletal category with no
nonidentity endomorphisms; finiteness of the underlying set of arrows
alone is insufficient. The equivalence of this commutation property
with stability is proved in
\citet[Theorem~2.14]{GrothShulman2017}. It identifies the exact boundary
between the true stable intuition and the stronger interchange asked
for on these pages.
\end{dcomment}


Plus généralement, considérons un diagramme dans $\sDia$
\[
\begin{tikzcd}[row sep=large, column sep=large]
I' \arrow[r, "q"] \arrow[d, "f'"'] & I \arrow[d, "f"] \\
J' \arrow[r, "p"'] & J
\arrow[from=1-2, to=2-1, dashed, shorten <=14pt, shorten >=14pt, "\gamma"']
\end{tikzcd}
\quad ,
\tag*{(2)}
\]
\printedpage{21}
% head: A. Grothendieck : Dérivateurs, Ch. I, Édition des Universités de Montpellier II et Paris VII
% foot: Version 8 juillet 2001
avec $\gamma : fq \to pf'$. On définit alors un \textit{morphisme de changement de base pour les images directes} relatif à ce diagramme
\[ p^* f_* \xrightarrow{c_\gamma} f'_* q^* \]
\msp{page 31}
de deux façons, qui donnent le même résultat (histoire générale de foncteurs adjoints à expliciter une bonne fois, en termes d'un diagramme de foncteurs
\[
\begin{tikzcd}[row sep=large, column sep=large]
X' & X \arrow[l, "q^*"'] \\
Y' \arrow[u, "{f'}^*"] & Y \arrow[l, "p^*"'] \arrow[u, "f^*"']
\arrow[from=2-1, to=1-2, dashed, shorten <=14pt, shorten >=6pt]
\end{tikzcd}
\]
en supposant que les foncteurs adjoints à gauche $f'_*$, $f_*$ des deux foncteurs verticaux existent). Un cas intéressant est celui où
\[ pf' = fq \quad , \qquad \gamma = \sid_{fq} \]
et où on pose $c_\gamma = c$.

Ici
\[ p^* f_* \xrightarrow{?} f'_* q^* \]
revient à
\[
\begin{tikzcd}[column sep=small]
f_* \arrow[r, "?"] & p_* f'_* q^* \simeq & (pf')_* q^* \\
 & & (fq)_* q^* \simeq (f_* q_*) q^* = f_*(q_* q^*) \quad . \arrow[u, "\gamma_* * q^*"']
\end{tikzcd}
\]
Or on a
\[ \sid \xrightarrow{\alpha} q_* q^* \qquad\qquad \text{(adjonction)} \quad , \]
d'où
\[ f_* \xrightarrow{f_* * \alpha} f_*(q_* q^*) \]
et en composant avec $\gamma * q^*$, on trouve ce qu'on veut. Et on peut au lieu de ceci chercher
\[
\begin{tikzcd}
{f'}^*(p^* f_*) \arrow[r, "?"] \arrow[d, equal] & q^* \\
(pf')^* f_* \arrow[d, "\gamma^* * f_*"] & \\
(fq)^* f_* \simeq q^*(f^* f_*) \arrow[r, "q^* * \gamma'"] & q^*
\end{tikzcd}
\quad ,
\]
où $\gamma' : f^* f_* \to \sid$ est l'adjonction.

\printedpage{22}
% head: A. Grothendieck : Dérivateurs, Ch. I, Édition des Universités de Montpellier II et Paris VII
% foot: Version 8 juillet 2001
\msp{page 32}
\textbf{Der}$_4$ : Le morphisme de changement de base pour images directes $f_*$ associé à un diagramme (1) est toujours un isomorphisme (dérivateurs à gauche). Condition duale pour dérivateurs à droite et le changement de base pour images directes $f_!$, pour le diagramme
\[
\begin{tikzcd}[row sep=large, column sep=large]
j_0 \backslash I \arrow[r] \arrow[d] & I \arrow[d, "f"] \\
e \arrow[r, "p = \alpha_{j_0}"'] & J
\arrow[from=2-1, to=1-2, dashed, shorten <=22pt, shorten >=8pt, "\gamma"]
\end{tikzcd}
\tag*{(1$'$)}
\]
associé à $f : I \to J$ dans $\sDia$, et $j_0 \in \sOb\ J$.

\dgrblock{dgrb18}{The base-change map as a mate}
\begin{dcomment}{context}{p31}{The base-change map as a mate}
The requested general adjunction calculation belongs to the calculus of
mates. Its two-categorical setting was already available in
\citet{KellyStreet1974}; a derivator-oriented account is
\citet[Appendix~A]{GrothPontoShulman2014}. Here we spell out the map in
the manuscript's variance, since reversing a single arrow changes the
statement.

Suppose $\gamma:fq\Rightarrow pf'$ is given. Contravariance on
2-cells gives
$\gamma^*:f'^*p^*\Rightarrow q^*f^*$. Write $\eta'$ for the unit of
$f'^*\dashv f'_*$ and $\varepsilon$ for the counit of
$f^*\dashv f_*$. The canonical transformation is the composite
\[
\begin{aligned}
 p^*f_*&\xrightarrow{\eta'\,p^*f_*}
       f'_*f'^*p^*f_*\\
 &\xrightarrow{f'_*\gamma^*f_*}f'_*q^*f^*f_*
   \xrightarrow{f'_*q^*\varepsilon}f'_*q^* .
\end{aligned}
\]
In particular, this construction needs only the adjoints $f_*$ and
$f'_*$; adjoints to $p^*$ and $q^*$ are convenient for the alternative
calculation, but are not part of its minimal input.

The equality of the two calculations is an equality of natural
transformations. Under the adjunction bijection
\[
 \operatorname{Nat}(p^*f_*,f'_*q^*)
 \cong\operatorname{Nat}(f'^*p^*f_*,q^*)
\]
the displayed composite is sent to
$q^*\varepsilon\circ(\gamma^*f_*)$, by the triangular identity for
$f'^*\dashv f'_*$. In the alternative construction, transpose first
across $p^*\dashv p_*$; replace $(pf')_*$ and $(fq)_*$ by the
composites of right adjoints; then transpose back across
$f'^*\dashv f'_*$. Naturality of the units moves them past
$\gamma^*$, and the unit--counit pair for $q^*\dashv q_*$ cancels by
its triangular identity. The resulting transformation is exactly
$q^*\varepsilon\circ(\gamma^*f_*)$. Since the adjunction bijection is
injective, the original two transformations agree.

A strictly commuting square supplies $\gamma=1$, but does not supply
invertibility of its mate. The role of Der~4 is to assert invertibility
for the particular comma squares that compute fibres. More general
base-change theorems will require conditions on the square. The words
``adjoints \`a gauche'' on this page are another source slip: the
functors bearing a subscript $*$ in this calculation are right adjoints.
\end{dcomment}


\textbf{Der}$_5$\footnote[10]{\underline{Totalement} canulé, devient raisonnable que pour coefficients constants sur $F$.} : Soit $f : I \to J$ une \textit{hot-équivalence stricte, i.e.} $\forall j_0 \in \sOb\ J$, $I/j_0$ est contractile. Alors $f^*$ est pleinement fidèle, \textit{i.e.} le morphisme d'adjonction
\[ \sid_{\DD(J)} \longrightarrow f_* f^* \]
est un isomorphisme ($\DD$ dérivateur à gauche). Dualement, si $f$ est une équivalence co-stricte, \textit{i.e.} $f^\circ : I^\circ \to J^\circ$ est une équivalence stricte, \textit{i.e.} $\forall j_0 \in \sOb\ J$, $j_0 \backslash I$ est contractile, alors (si $\DD$ est un dérivateur à droite) le foncteur $f^*$ est pleinement fidèle, \textit{i.e.} le morphisme d'adjonction % Resolved: the edition line-break gives co-stricte; retain that spelling.
\[ f_! f^* \longrightarrow \sid_{\DD(J)} \]
est un isomorphisme.

N.B. Donc les foncteurs $f_*$ (resp. $f_!$) sont des foncteurs de localisation.

\textbf{N.B.} Les pages 33 à 40 ont été suprimées, car canulées. % [sic] suprimées -> supprimées, spelling

\msp{page 41}
Je vais reprendre la terminologie et les axiomes relatifs aux dérivateurs.

\dgrblock{dgrb19}{The cancelled axiom and the walking arrow}
\begin{dcomment}{examples}{p32}{The cancelled axiom and the walking arrow}
The cancellation is mathematically necessary and should remain visible.
Let $Y=[1]=(0\to1)$ and let $f:e\to Y$ select $0$. Each slice $f/y$
has one object and one arrow, so $f$ satisfies the stated strict
hot-equivalence condition. In the represented derivator of abelian
groups a coefficient on $Y$ is a map $F_1\to F_0$. The restriction
$f^*$ remembers $F_0$, and its right adjoint sends $A$ to the constant
map $A\xrightarrow{1}A$. Consequently the unit is
\[
 (F_1\longrightarrow F_0)\longrightarrow
 (F_0\xrightarrow{1}F_0),
\]
whose component at $1$ is the original structure map. With
$F_1=0$ and $F_0=\mathbb Z$ it is not invertible. Contractibility of
the slice controls a constant coefficient system on that slice; it
does not turn the restriction of an arbitrary $F$ into such a system.

The failure can also be read directly as failure of full faithfulness.
The distinct endomorphisms of the diagram
$\mathbb Z\to0$ obtained from the identity and zero on $\mathbb Z$
have the same restriction to the object $0$. Thus $f^*$ is not faithful.
The dual example, using the final endpoint and left Kan extension,
disproves the dual part of the cancelled axiom.

There is a valid constant-coefficient conclusion. If
$F=p_Y^*M$, the pointwise right Kan formula gives
\[
 (f_*f^*p_Y^*M)_y
 \simeq\Gamma_{\mathbb D}(f/y,p_{f/y}^*M).
\]
Under the valid contractibility theorem the right-hand side is $M$, and
the isomorphism is precisely the component of the unit. Der~2 then
shows that the unit on $p_Y^*M$ is an isomorphism. This is the repair
anticipated in footnote~10. Full faithfulness of restriction on all
coherent diagrams is a stronger condition, usually called a
homotopical epimorphism. It must not be substituted for either
orientation of cofinality without an additional argument.
\end{dcomment}


Un \textit{prédérivateur} est un 2-foncteur
\[ \sDia^\circ \xrightarrow{\DD} \mathcal{C}at \quad , \]
où $\sDia$ est une 2-sous-catégorie pleine de $\mathcal{C}at$.\footnote[11]{$\sDia$ \textit{domaine} du dérivateur. Les objets sont les \textit{catégories d'indices} pour $\DD$. Les $\DD(I)$ sont les catégories de coefficients pour $\DD$, les objets de $\DD(I)$ les “coeficients de type $\DD$ sur $I$”. % [sic] coeficients -> coefficients, spelling
} On suppose tout au moins que $\sDia$ contient les \textit{ensembles ordonnés finis,} que $\sDia$ est stable par $\varprojlim$ finies, par sommes finies. Qu'avec toute catégorie $C$ et toute sous-catégorie strictement pleine ouverte (resp. fermée) elle contienne la sous-catégorie complémentaire.\footnote[12]{Notation $f^*$ : image inverse des coefficients.}

\printedpage{23}
% head: A. Grothendieck : Dérivateurs, Ch. I, Édition des Universités de Montpellier II et Paris VII
% foot: Version 8 juillet 2001
\msp{page 42}
On définit le prédérivateur $\DD_{\mathcal{M},\Sigma}$ associé à un système
\[ (\mathcal{M},\Sigma) \qquad\qquad \mathcal{M} \text{ Catégorie,} \quad \Sigma \subset \mathsf{Fl}(\mathcal{M}) \quad . \] % Resolved: Fl abbreviates fleches, with lowercase l; retain Fl.
Si $\Sigma = \emptyset$, on écrit $\DD_{\mathcal{M}}$.

\dgrblock{dgrb20}{Restarting the axioms without changing variance}
\begin{dcomment}{exposition}{p41}{Restarting the axioms without changing variance}
The return to definitions is a revision of the theory, not an additional
list to be imposed on top of the cancelled axioms. We keep
$\mathbb D(I)$ for contravariant coherent diagrams. In fully explicit
modern notation its source is $\Dia^{\mathrm{coop}}$, with both
1-cells and 2-cells reversed. Equivalently one can work with the
covariant-diagram convention by setting
$\mathbb D_{\mathrm{cov}}(A)=\mathbb D(A^{\circ})$. This translation
explains why modern formulas often use the opposite comma category;
compare \citet[Warning~2.2]{GrothPontoShulman2014}.

The diagram domain is part of the hypotheses. The operations used below
require products, coproducts, opposites, and the relevant comma
categories. The stated finite-limit and complement closures do not
explicitly include all these operations. Taking all small categories
in a fixed universe makes the present formulas unambiguous; when using
a smaller domain, one must require the corresponding closure conditions
before writing the dual prederivator or the pointwise Kan formulas.
No assertion about a shape outside the domain follows formally from an
axiom restricted to that domain.

For a relative category $(\mathcal M,\Sigma)$ the diagram localization
exists as a prederivator once the chosen universes allow the required
localizations. Saturating the weak equivalences does not by itself
prove the derivator axioms. On the other hand, the shift identity does
not demand a new homotopical theorem. The exponential law identifies
$\mathcal M^{(I\times J)^{\circ}}$ with
$(\mathcal M^{I^{\circ}})^{J^{\circ}}$, and identifies the two
classes of pointwise weak equivalences before localization. Localizing
therefore gives the claimed equivalence of shifted prederivators.
Functoriality in relative functors and in natural transformations is
descended through the same universal property. This is a useful place
to separate what follows from localization alone from the homotopy Kan
extension properties that require additional assumptions.
\end{dcomment}


Les prédérivateurs forment une 2- catégorie $\mathsf{PREDER}$, et
\[ (\mathcal{M},\Sigma) \longmapsto \DD_{\mathcal{M},\Sigma} \]
est un 2-foncteur
\[ \mathsf{PREMOD}
\begin{pmatrix}
\text{“Catégorie des catégories} \\
\text{de modèles” (faibles)}
\end{pmatrix}
\longrightarrow \mathsf{PREDER} \quad . \]
Pour $\DD$ donné, on définit le dérivateur induit ($I \in \sDia$)
\[ \DD_I = \big(J \longmapsto \DD(I \times J)\big) \quad , \]
dépendant de $I$ de façon contravariante, d'où un 2-foncteur
\[ \begin{aligned}
\sDia^\circ &\longrightarrow \mathsf{PREDER} \\
I &\longmapsto \DD_I
\end{aligned} \]
ou encore un bi-2-foncteur
\[ \begin{aligned}
\mathsf{PREDER} \times \sDia^\circ &\longrightarrow \mathsf{PREDER} \\
(\DD, I) &\longmapsto \DD_I \quad .
\end{aligned} \]
Si $\DD$ est défini par $(\mathcal{M},\Sigma)$, alors $\DD_I$ est défini par $(\mathcal{M}(I),\Sigma(I))$ :
\[ (\DD_{\mathcal{M},\Sigma})_I \simeq \DD_{\mathcal{M}(I),\Sigma(I)} \quad , \]
\msp{page 43}
où $\mathcal{M}(I) = \underline{\mathsf{Hom}}(I^\circ,\mathcal{M})$ (préfaisceaux sur $I$ à valeurs dans $\mathcal{M}$), et $\Sigma(I) \subset \mathsf{Fl}\ \mathcal{M}(I)$ est déduit de $\Sigma$ en regardant fibre par fibre (supposant maintenant $\Sigma$ saturé). % Resolved: the repeated arrow-class symbol is Fl, with lowercase l.

Pour tous axiomes qu'on dégagera pour certains prédérivateurs, il faudra vérifier qu'il sont stables par passage aux prédérivateurs induits $\DD_I$, et aussi comment ils se comportent par dualité (passage de $\DD$ à $\DD^\circ$). % [sic] qu'il sont -> qu'ils sont, agreement

Définition du prédérivateur dual $\DD^\circ$
\[ \DD^\circ(I) \overset{\text{déf}}{=} \DD(I^\circ)^\circ \qquad\qquad \textit{etc.} \]
\textbf{Der 1} (Premier axiome de localisation) : \textit{a}) Si $I$, $J \in \sDia$, alors
\[ \DD(I \amalg J) \longrightarrow \DD(I) \times \DD(J) \]
est une équivalence de catégories.

\textit{b}) $\DD(\emptyset)$ est la catégorie finale.

Stable par passage aux prédérivateurs induits, et passage au dual.

\textbf{Corollaire.} \textit{Soit} $(I_\alpha)_{\alpha \in A}$ \texttt{[une]} \textit{famille finie de catégories de} $\sDia$, $I = \coprod\limits_{\alpha \in A} I_\alpha$. \textit{Alors}
\[ \DD(I) \longrightarrow \prod_{\alpha \in A} \DD(I_\alpha) \]
\textit{est une équivalence de catégories.}

% ---- from der1_p026-030.tex
\printedpage{24}
% head: A. Grothendieck : Dérivateurs, Ch. I, Édition des Universités de Montpellier II et Paris VII

\noindent\textbf{Der 2} (Deuxième axiome de localisation) : Pour tout $I \in \Dia$, le foncteur canonique
\[ \D(I) \longrightarrow \underline{\mathsf{Hom}}(I^{\circ}, \mathcal{A}_{\D}) \quad , \]
$\mathcal{A}_{\D} = \D(e)$ catégorie fondamentale pour le dérivateur $\D$, est conservatif.\\
Stable par induction, et par dualité.

\noindent\msp{page 44}\texttt{[page 44]}

\dgrblock{dgrb21}{What the localization axioms do and do not say}
\begin{dcomment}{exposition}{p43}{What the localization axioms do and do not say}
Der~1 decomposes coefficients over disconnected indexing categories.
By induction on the number of summands, its binary and empty cases give
the displayed assertion for every finite family. It survives a shift
because $I\times(J\amalg K)$ is canonically
$(I\times J)\amalg(I\times K)$ and $I\times\varnothing=\varnothing$.
It survives duality because taking opposites commutes with products of
categories and with the corresponding indexing coproducts.

Der~2 is a detection statement for morphisms already present in
$\mathbb D(I)$. If such a morphism becomes invertible at every object
of $I$, it was invertible. This gives neither a lifting of every
incoherent diagram nor a lifting of every map between incoherent
diagrams. The later strongness axiom addresses a different issue and
must not be read into this one. To check Der~2 after shifting by $J$,
apply it first to the $J$-valued fibres and then to the pairs of
objects of $I\times J$. The converse direction follows because every
functor preserves isomorphisms. Duality preserves the same criterion.

The next axiom names $f_*F$ relative cohomology. Its transitivity is an
adjunction calculation: both $g_*f_*$ and $(gf)_*$ are right adjoint
to $f^*g^*=(gf)^*$. This gives a specified canonical comparison,
compatible with further composition. The existence of finite products
again follows by applying this axiom to finite discrete projections.
A right adjoint $f_*$ preserves them automatically. Preservation by
restriction $f^*$ in a one-sided derivator uses the independent-variable
base-change assertion discussed at manuscript page~23; the final
pointwise axiom supplies the missing justification. The corresponding
homological statements follow by duality.
\end{dcomment}


\noindent Un prédérivateur est appelé un \textit{dérivateur à gauche} s'il satisfait \textbf{Der 1} et \textbf{Der 2}, et de plus les axiomes suivants (dont voici le premier).

\noindent\textbf{Der 3} (g) (Existence des images directes) : Pour tout $f : I \to J$ dans $\Dia$, le foncteur $f^{*} : \D(J) \to \D(I)$ admet un \textit{adjoint à droite} $f_{*}$\,.\\
Stable par induction.

On appelle $f_{*}$ le foncteur image directe (cohomologique) pour $f$. L'objet $f_{*}(F) \in \D(J)$ (pour $F \in \D(I)$) s'appelle aussi l'objet de “\textit{cohomologie relative}” de $I$ sur $J$ (par $f$), à coefficients dans $F$. On a des isomorphismes canoniques
\[ (gf)_{*} \simeq g_{*} f_{*} \quad . \]

\noindent\textbf{Corollaire.} \textit{Les catégories $\D(I)$ ont des produits finis, les foncteurs $f^{*}$ et $f_{*}$ y commutent.}

Utilise \textbf{Der 1} pour $I = \emptyset$ et $I = e \amalg e$, et \textbf{Der 3} pour $I \xrightarrow{f} e$.

On a un axiome dual :

\noindent\textbf{Der 3} (d) : existence des adjoints à gauche $f_{!}$ des $f^{*}$, le foncteur $f_{!}$ s'appelle “image directe homologique” par $f$, et on a une notion d'\textit{homologie relative} de $I$ sur $J$ (pour $f$), à coefficients dans $F$.

Par la suite, la mention (g) ou (d) à coté de la désignation \textbf{Der} $n$ d'un axiome, % [sic] coté -> côté, missing circumflex
signifie qu'il s'applique aux dérivateurs à gauche, respectivement à droite. De même pour les énoncés, définitions \textit{etc.}

\noindent\msp{page 45}\texttt{[page 45]}

\noindent\textbf{Der 4}$_{0}$ (g)\footnote[13]{Inutile comme axiome, va devenir un théorème (résulte des autres axiomes).} : Soit $U \overset{j}{\hookrightarrow} X \overset{i}{\hookleftarrow} Y$ dans $\Dia$, immersions ouverte et fermée complémentaires l'une de l'autre. Alors le morphisme d'adjonction dans $\D(Y)$
\[ i^{*} i_{*}(F_{Y}) \longrightarrow F_{Y} \]
est un isomorphisme, donc
\[ i_{*} : \D(Y) \longrightarrow \D(X) \]
est pleinement fidèle, et $i^{*} : \D(X) \to \D(Y)$ est un foncteur de localisation. L'image essentielle de $i_{*}$ est formée des objets $F$ de $\D(X)$ tels que $j^{*}(F)$ soit l'objet final de $\D(U)$.

% foot: Version 8 juillet 2001

\printedpage{25}
% head: A. Grothendieck : Dérivateurs, Ch. I, Édition des Universités de Montpellier II et Paris VII

\noindent\textbf{Der 4}$_{0}$ (d) : Avec les notations précédentes pour $i$, $j$, le morphisme d'adjonction dans $\D(U)$
\[ F_{U} \longrightarrow j^{*} j_{!} F_{U} \]
est un isomorphisme, - $j_{!}$ est pleinement fidèle, $j^{*}$ un foncteur de localisation. L'image essentielle de $j_{!}$ est la sous-catégorie pleine de $\D(X)$ formée des $F$ tels que $i^{*}(F)$ soit objet initial de $\D(Y)$.

Il y a les deux carrés de la page 8 et 9, mais pour interpréter la condition cartésienne ou cocartésienne c'est prématuré ici.

\noindent\msp{page 46}\texttt{[page 46]}

\dgrblock{dgrb22}{Extension by an empty fibre}
\begin{dcomment}{exposition}{p45}{Extension by an empty fibre}
The proposed additional axiom $\mathrm{Der}\,4_0$ is a consequence of
the pointwise Kan formula. This is a useful calculation because it
distinguishes the special behaviour of an immersion from the false
full-faithfulness statement for arbitrary cofinal functors.

In the manuscript's convention a closed full inclusion
$i:Y\hookrightarrow X$ is a cosieve: an arrow whose source lies in
$Y$ has target in $Y$. For $y\in Y$, the comma category
$(i/i(y))$ has terminal object $(y,1)$. Since the diagram is
contravariant, its cohomology is the value at that terminal object.
Thus $(i_*E)_{i(y)}\cong E_y$, and these are the components of the
counit $i^*i_*E\to E$. Der~2 makes the counit invertible, proving
full faithfulness of $i_*$.

For $x\notin Y$, the comma category $(i/x)$ is empty: an arrow
$i(y)\to x$ would contradict the cosieve condition. Hence
$(i_*E)_x$ is terminal, the right Kan extension over an empty category.
Conversely let $F\in\mathbb D(X)$ have terminal fibres on the
complement. In the unit $F\to i_*i^*F$, the components on $Y$ are
isomorphisms by the preceding calculation, and the components off $Y$
are maps between terminal objects. Der~2 makes the whole unit
invertible, so $F$ belongs to the essential image. This proves both
assertions of $\mathrm{Der}\,4_0(g)$.

For an open inclusion $j:U\hookrightarrow X$, which is a sieve, the
left Kan formula uses $(x/j)$. Outside $U$ this is empty, and inside
$U$ its initial object becomes terminal after taking the opposite
indexing category for the colimit. Therefore $j_!$ is fully faithful
and extends by initial objects. This proves the dual assertion.
The calculation is recorded in modern covariant notation in
\citet[Proposition~1.23]{Groth2013}. Only the empty and pointwise
Kan formulas are doing work here; no separate extension-by-zero axiom
is needed. Pointedness later identifies the terminal and initial
extensions with extension by the same zero object.
\end{dcomment}


\noindent Axiomes de changement de base.

\noindent\textbf{Der 4} (g) (Fibres d'une image directe) : Soient $f : X \to Y$ dans $\Dia$, $y \in Y$, considérons le diagramme connu
\[
\begin{tikzcd}[row sep=large, column sep=large]
X/y \arrow[r, "q"] \arrow[d, "f'"'] & X \arrow[d, "f"] \\
e \arrow[r, "p = \alpha_{y}"'] & Y
\end{tikzcd}
\]
et le morphisme de changement de base correspondant
\[ f_{*}(F)_{y} \overset{\text{déf}}{=} p^{*} f_{*}(F) \longrightarrow f'_{*}(q^{*}(F))\bigl(\overset{\text{déf}}{=} \Gamma_{\D}(X/y, q^{*}(F))\bigr) \quad . \]

\noindent C'est toujours un isomorphisme.

\noindent\msp{page 47}\texttt{[page 47]}

\noindent\textbf{Der 5} (g) (Axiome de cofinalité pour la cohomologie)\footnote[14]{Il est conséquence d'énoncés plus faibles.} : Soit $\alpha : X' \to X$ dans $\Dia$ “cohomologiquement cofinal” \textit{i.e.} les $x\backslash X'$ sont contractiles ($x \in X$) \textit{i.e.} $\alpha^{\circ} : {X'}^{\circ} \to X^{\circ}$ est une hot-équivalence “localement sur $X^{\circ}$\,”. Alors $\forall F \in \D(X)$, l'homomorphisme
\[ \Gamma_{\D}(X, F) \longrightarrow \Gamma_{\D}(X', \alpha^{*} F) \]
est un isomorphisme (dans $\mathcal{A} = \D(e)$).

\noindent N.B. Si $f : X \to e$, on pose $\Gamma_{\D}(X, F) = f_{*}(F)$ (cohomologie “absolue” de $X$, à coefficients dans $F$).

\noindent Même chose\footnote[15]{Faux!} si $\alpha$ est “homologiquement cofinal”, i.e. les $X'/x$ sont contractiles \texttt{[phrase barrée]}.

% foot: Version 8 juillet 2001

\printedpage{26}
% head: A. Grothendieck : Dérivateurs, Ch. I, Édition des Universités de Montpellier II et Paris VII

\noindent\textbf{Remarque.} On devrait pouvoir renforcer considérablement cet axiome, en exigeant que
\[ \alpha_{!}(\varepsilon_{X'}) \xrightarrow{\sim} \varepsilon_{X} \qquad\qquad \text{dans} \quad \HOT(X) \quad , \]
où $\varepsilon_{X'}$, $\varepsilon_{X}$ \texttt{[sont]} les objets finaux de $\HOT(X')$, $\HOT(X)$, et où $\alpha_{!}(\varepsilon_{X'}) \to \varepsilon_{X}$ est l'homomorphisme d'adjonction, compte tenu que $\varepsilon_{X'} \simeq \alpha^{*}(\varepsilon_{X})$. On dit que $X'$ est \textit{homotopiquement trivial sur $X$}. (Et si la condition duale est satisfaite dans $\HOT(X^{\circ})$, on dit que $X'$ est cohomotopiquement trivial sur $X$.) Cette condition n'est sans doute \underline{pas} autoduale. Elle est vérifiée aussi bien si $\alpha$ est cohomologiquement cofinal, que s'il est homologiquement cofinal (hot-équivalence costricte ou stricte).

\noindent N.B. Pour $f : X \to e$, dire que $f$ est cohomologiquement cofinal, homologiquement cofinal, ou que $X$ est contractile, est équivalent.

\noindent\textbf{Corollaire 1} (g) (Théorème d'asphéricité pour les coefficients constants en cohomologie). \textit{Soit $f : X \to e$ dans} $\Dia$\textit{, $X$ contractile, $M \in \D(e)$. Alors l'homomorphisme d'adjonction}
\[ M \xrightarrow{\sim} f_{*} f^{*}(M) \quad \bigl(\overset{\text{déf}}{=} \Gamma(X, f^{*}(M))\bigr) \]
\textit{est un isomorphisme.}

\noindent\textbf{Corollaire 2} (g)\footnote[16]{\underline{Faux}, sauf pour coefficients constants.} \textit{Soit $f : X \to Y$} \texttt{[une]} \textit{hot-équivalence stricte (i.e. $\forall y \in Y$, $X/y$ contractile). Alors pour $F \in \Ob\,\D(Y)$}
\[ F \overset{\text{adj}}{\xrightarrow{\sim}} f_{*} f^{*} F \]
\texttt{[est un]} \textit{isomorphisme.} \texttt{[Ce corollaire est barré, mais pas sa démonstration.]}

\noindent\msp{page 48}\texttt{[page 48]}

\dgrblock{dgrb23}{Cofinality retains an orientation}
\begin{dcomment}{status}{p47}{Cofinality retains an orientation}
The valid cohomological cofinality condition is contractibility of
$(x/\alpha)$, whose objects are arrows $x\to\alpha(x')$.
For contravariant coefficients it preserves the homotopy limit.
Contractibility of $(\alpha/x)$ is the opposite condition and
preserves the homotopy colimit. These conditions become equivalent for
a map to the point, but not for a general map. The cofinality theorem
is part of the established derivator calculus; in covariant notation
see \citet[Corollary~3.13 and Remark~3.14]{GrothPontoShulman2014}.
The theorem's derivation for all small categories does not remove the
need to check domain hypotheses in a restricted theory.

The walking-arrow example in the commentary at manuscript page~32
settles the crossed-out sentence. For the inclusion of $0$ in
$Y=[1]$, every $(\alpha/y)$ is a point, but
$\Gamma(Y,F)=F_1$ while $\Gamma(e,\alpha^*F)=F_0$.
The comparison is the given structure map $F_1\to F_0$, which may be
$0\to\mathbb Z$. Homotopy invariance with arbitrary varying
coefficients cannot ignore its direction.

The proposed condition involving the terminal object of $\mathsf{HOT}$
can be calculated completely. In the homotopy theory of presheaves of
spaces, the pointwise left Kan formula yields
\[
 (\alpha_!\varepsilon_{X'})_x
 \simeq\operatorname*{hocolim}_{(x/\alpha)^{\circ}}*
 \simeq B(x/\alpha).
\]
For the second identification, the simplicial replacement of the
constant point has one simplex for each composable string of arrows
in the indexing category; its realization is the nerve realization.
Passing to the opposite category reverses these strings and gives the
same homotopy type. The counit in the remark has component the
augmentation $B(x/\alpha)\to *$. Der~2 therefore identifies its
invertibility with weak contractibility of every $(x/\alpha)$.
Thus this condition is a reformulation of the valid cohomological
cofinality hypothesis, not a substantial strengthening of it.

It also refutes the assertion that either orientation suffices. For
the same initial-endpoint inclusion, $(1/\alpha)$ is empty, so the
component of $\alpha_!\varepsilon$ at $1$ is the empty space and its
map to the point is not a weak equivalence. The opposite-endpoint
inclusion supplies the dual discrepancy. The remark's prediction
that the condition is not self-dual is correct, while its claimed
implication from both cofinality conditions is false.

The constant-coefficient corollary has no such problem. If
$p:X\to e$ and $NX$ is weakly contractible, the valid cohomological
cofinality theorem applied to $M$ gives
$M\cong p_*p^*M$. This is exactly the unit, rather than an
unspecified isomorphism. What fails in the following cancelled
corollary is the replacement of a genuinely varying restriction of
$F$ by a constant coefficient system on a contractible slice.
\end{dcomment}


\noindent On applique \textbf{Der 2}, \textbf{Der 4} (g), pour se ramener à regarder, pour $y_{0} \in Y$, l'homomorphisme $\alpha$ du diagramme (commutatif!) suivant :
\[
\begin{tikzcd}[row sep=large, column sep=huge]
\Gamma(\{y_{0}\}, F_{y_{0}}) = F_{y_{0}} \arrow[r, "\alpha"] & \Gamma_{\D}(X/y_{0}, f^{*}(F)) \\
\Gamma(Y/y_{0}, F) \arrow[u, "\varphi\,\wr"] \arrow[ur, "\alpha'"'] &
\end{tikzcd}
\quad ,
\]
où $\alpha'$ est déduit du foncteur induit par $f$
\[ X/y_{0} \longrightarrow Y/y_{0} \]
et où $\varphi$ est déduit du foncteur
\[ \{y_{0}\} \longrightarrow Y/y_{0} \qquad \text{(objet final)}\,. \]

% foot: Version 8 juillet 2001

\printedpage{27}
% head: A. Grothendieck : Dérivateurs, Ch. I, Édition des Universités de Montpellier II et Paris VII

\noindent Donc $\varphi$ est un isomorphisme par \textbf{Der 5} (g)\footnote[17]{\textit{cf.} corollaire 3, page 50, qu'on devrait remonter ici.} (N.B. Si $z \in Z$, d'où $\{z\} \to Z$, ce foncteur est cohomologiquement cofinal si et seulement si $z$ est objet final, homologiquement final \texttt{[plutôt cofinal]} si et seulement si $z$ est objet initial). Donc la conclusion du corollaire équivaut à dire que $\alpha'$ est un isomorphisme. Mais $X/y_{0} \to Y/y_{0}$ est aussi homologiquement cofinal, car si $u : y_{1} \to y_{0}$ est dans $Y/y_{0}$, on a un isomorphisme canonique
\[ (X/y_{0})/(y_{1}, u) \simeq X/y_{1} \]
qui est contractile. Donc par la deuxième partie de \textbf{Der 5} (g),\footnote[18]{(qui est \underline{fausse})} on gagne.

\noindent\msp{page 49}\texttt{[page 49]}

\dgrblock{dgrb24}{Locating the failed inference in the surviving proof}
\begin{dcomment}{exposition}{p48}{Locating the failed inference in the surviving proof}
Most of the surviving proof is correct. The category $Y/y_0$ has
terminal object $(y_0,1)$, and cohomology of a contravariant diagram
therefore evaluates to $F_{y_0}$. The inclusion of that terminal
object is cohomologically cofinal: its comma category over an object
is the one-element set consisting of the unique map to the terminal
object. More generally, for $z\in Z$ the comma category
$(x/\{z\})$ is the discrete set $Z(x,z)$. Its nerve is contractible
exactly when that set has one element. Hence the singleton inclusion
is cohomologically cofinal precisely when $z$ is terminal; the dual
assertion characterizes initial objects.

The displayed slice identification is also exact. An object of
$(X/y_0)/(y_1,u)$ consists of $x\in X$, a map
$v:f(x)\to y_0$, and a map $w:f(x)\to y_1$ with $uw=v$.
Forgetting $v$, which is determined by $w$, identifies this category
with $X/y_1$, on morphisms as well as objects. This proves the
\emph{homological} cofinality of $X/y_0\to Y/y_0$ under the stated
hypothesis. It does not prove cohomological cofinality. The final
sentence invokes exactly the opposite-orientation clause that was
crossed out on manuscript page~47. No additional fault in the slice
calculation is needed to account for the failure. The author and
editors explicitly mark the false premise; a modern repair retains
the calculation and changes the coefficient hypothesis or the
cofinality hypothesis, as in the preceding commentary.
\end{dcomment}


\noindent\textbf{Corollaire} 4 (g) (critère de changement de base pour images directes cohomologiques). \textit{Soit}
\[
\begin{tikzcd}[row sep=huge, column sep=huge]
X' \arrow[r, "q"] \arrow[d, "f'"'] & X \arrow[d, "f"] \arrow[dl, dashed, "\gamma"', shorten <=4ex, shorten >=1ex] \\
Y' \arrow[r, "p"'] & Y
\end{tikzcd}
\]
\textit{dans} $\Dia$ (\textit{cf.} page 30), \textit{d'où l'homomorphisme de changement de base}
\[ p^{*} f_{*}(F) \longrightarrow f'_{*} q^{*}(F) \qquad (F \in \Ob\,\D(X)) \quad . \]
\textit{Pour que ce soit un isomorphisme pour tous $F$, il suffit que pour tout $y' \in \Ob\,Y'$, le foncteur induit par} (p. 9)
\[ X'/y' \longrightarrow X/p(y') \]
\textit{soit une hot-équivalence stricte (ou costricte);\footnote[19]{\underline{Faux} pour “stricte”, ne marche \underline{que} pour “\underline{\underline{co}}stricte”.} cela signifie aussi (dans le cas strict) que pour}
\[ y'_{0} \in \Ob\,Y' \,,\ \ x_{0} \in \Ob\,X \,,\ \ u_{0} : f(x_{0}) \longrightarrow p(y'_{0}) \]
\textit{la catégorie $K(y'_{0}, x_{0}, u_{0})$ des objets de la forme}
\[ (x', v', u) \mid x' \in \Ob\,X',\ \ v' : f'(x') \to y'_{0},\ \ u : q(x') \overset{\dashleftarrow}{\longrightarrow} x_{0} \]
\textit{rendant commutatif le carré dans $Y'$}
\[
\begin{tikzcd}[row sep=huge, column sep=huge]
f(q(x')) \arrow[r, "\gamma(x')"] \arrow[d, shift right=0.5ex, "f(u)"'] & p(f'(x')) \arrow[d, "p(v')"] \\
f(x_{0}) \arrow[r, "u_{0}"'] \arrow[u, dashed, shift right=0.5ex, "\leftarrow\ \text{\scriptsize\shortstack[l]{cas de\\hot-équivalence\\costricte}}"'] & p(y'_{0})
\end{tikzcd}
\tag*{$(*)$}
\]
\textit{est contractile.}

% foot: Version 8 juillet 2001

\printedpage{28}
% head: A. Grothendieck : Dérivateurs, Ch. I, Édition des Universités de Montpellier II et Paris VII

\noindent N.B. Dans le cas de la hot équivalence \underline{co}stricte, c'est la même description de
\[ K = (x_{0}, u_{0})\backslash(X'/y'_{0}) \quad , \]
sauf que la flèche $u$ est en sens inverse $x_{0} \to q(x')$ et itou dans le diagramme commutatif $(*)$.\\
On applique \textbf{Der 2}, \textbf{Der 4} (g) et corollaire 1 \texttt{[plutôt} \textbf{Der 5} (g)\texttt{]}.

\noindent\msp{page 50}\texttt{[page 50]}

\dgrblock{dgrb25}{The factorization category controls base change}
\begin{dcomment}{exposition}{p49}{The factorization category controls base change}
The base-change criterion is a statement about the square, including its specified natural transformation. It does not assert that every cartesian square is exact, and it does not require the square to be cartesian. To see what it tests, fix $y'\in Y'$, put $y=p(y')$, and write
\[
 h_{y'}:X'/y'\longrightarrow X/y,
 \qquad (x',v:f'(x')\to y')\longmapsto
 (q(x'),p(v)\gamma_{x'}).
\]
The pointwise formula for right Kan extension identifies the value of the Beck--Chevalley map at $y'$ with restriction of cohomology along $h_{y'}$. Since the coefficients are contravariant, the appropriate cofinality test uses arrows \emph{from} an object of $X/y$ \emph{to} the image of $h_{y'}$. This is precisely the orientation retained in footnote~19. The discarded orientation is the mistake already isolated in the commentary at manuscript p.~32, ``The cancelled axiom and the walking arrow.''

For clarity, the category attached to $x\in X$ and $\alpha:f(x)\to p(y')$ has objects
\[
 (x',u:x\to q(x'),v:f'(x')\to y'),
 \qquad p(v)\gamma_{x'}f(u)=\alpha.
\]
A morphism to $(x'',u',v')$ is an arrow $t:x'\to x''$ such that $q(t)u=u'$ and $v'f'(t)=v$. Thus this is the comma category $(x,\alpha)\backslash h_{y'}$, with all compatibility conditions retained. Weak contractibility of every such category implies that restriction along $h_{y'}$ preserves homotopy limits of the contravariant diagrams in question. Applying this at every $y'$, and then using \textbf{Der 2}, proves the criterion. The words ``dans $Y'$'' before the displayed square are a slip in the edition: the square is in $Y$.

The resulting invariant is now called a \emph{homotopy exact square}. In modern covariant notation its factorization categories are those of \citet[Definition~3.7 and Theorem~3.8]{GrothPontoShulman2014}; passage to opposite indexing categories recovers exactly the arrows above. Their Theorem~3.16 also gives the converse when exactness is required for every derivator. The useful content of this later formulation is that the elementary categories of factorizations, rather than any particular presentation by model categories, determine the universal base-change property.
\end{dcomment}


\noindent\textbf{Corollaire 5} (g). \textit{Soit $f : X \to Y$ pleinement fidèle dans} $\Dia$\textit{, et $F \in \D(X)$. Alors l'homomorphisme d'adjonction}
\[ f^{*} f_{*}(F) \longrightarrow F \]
\textit{est un isomorphisme.}

En efffet, il suffit de le vérifier fibre par fibre (\textbf{Der 2}). Soit $x \in X$, $y = f(x) \in Y$, % [sic] efffet -> effet, tripled f
alors
\[ f^{*} f_{*}(F)_{x} \simeq f_{*}(F)_{y} \simeq \Gamma_{\D}(X/y, F) \]
or à cause de la pleine fidélité $X/y \simeq X/x$.

\noindent\textbf{Corollaire 3} (remonter page 48).\footnote[20]{Va résulter de \textbf{Der 6} (g) plus fort ci-dessous.} \textit{Soit $Z \in$} $\Dia$\textit{, supposons que $Z$ a un objet final $e_{Z}$. Soit $p : e \to Z$ le foncteur correspondant. Alors l'homomorphisme canonique (pour $F$ dans} $\D(Z)$\,)
\[
\begin{tikzcd}[row sep=huge]
Z \arrow[d, shift left=0.6ex, "f"] \\
e \arrow[u, shift left=0.6ex, "p"]
\end{tikzcd}
\qquad\qquad\qquad\qquad
f_{*}(F) \longrightarrow p^{*}(F)
\]
\textit{est un isomorphisme.}\footnote[21]{Axiome de la catégorie d'indices à objet final (coefficient quelconque).}

Cet homomorphisme est \texttt{[un]} cas particulier de celui-ci, pour
\[
\begin{tikzcd}[row sep=large, column sep=large]
Z' \arrow[rr, "p"] \arrow[dr, "f'"'] & & Z \arrow[dl, "f"] \\
& S &
\end{tikzcd}
\]
(en faisant $Z' = S = e$, $f' = \mathsf{id}_{e}$),
\[ f_{*}(F) \longrightarrow f'_{*}(p^{*}(F)) \quad , \qquad (\, F \in \Ob\,\D(Z)\,) \quad , \]
et c'est un cas particulier\footnote[22]{déjà envisagé dans cor. 2, p. 47} des morphismes de

% foot: Version 8 juillet 2001

% ---- from der1_p031-035.tex
\printedpage{29}
% head: A. Grothendieck : Dérivateurs, Ch. I, Édition des Universités de Montpellier II et Paris VII
% foot: Version 8 juillet 2001

\noindent\msp{page 51}
changement de base pour
\[
\begin{tikzcd}
Z' \arrow[r, "p"] \arrow[d, "f'"'] & Z \arrow[d, "f"] \\
S \arrow[r, "\id"] & S
\end{tikzcd}
\qquad .
\]

\dgrblock{dgrb26}{Fully faithful direct image and extension by a terminal object}
\begin{dcomment}{exposition}{p50}{Fully faithful direct image and extension by a terminal object}
There are two separate assertions behind the next localization calculation. First, a fully faithful inclusion retains all coefficients on its image. If $f:X\to Y$ is fully faithful, the map
\[
 X/x\longrightarrow X/f(x),\qquad
 (x',a:x'\to x)\longmapsto(x',f(a)),
\]
is an isomorphism of categories. The right member here is the comma category formed with $f$. Fullness and faithfulness give, respectively, existence and uniqueness of the inverse on objects and arrows. Since $X/x$ has terminal object $(x,1_x)$, the pointwise formula gives
\[
 (f^*f_*F)_x\simeq\Gamma_{\mathbb D}(X/x,F)\simeq F_x.
\]
The map so obtained is the counit, not merely an unspecified isomorphism between its source and target. Conservativity therefore proves that $f_*$ is fully faithful.

Second, for a cosieve $i:Y\hookrightarrow X$, right extension in the presheaf convention fills the complementary sieve $U$ with terminal objects. If $u\in U$, there can be no arrow from an object of $Y$ to $u$, so $Y/u$ is empty. The pointwise formula consequently makes $(i_*F)_u$ terminal. Conversely, let $E\in\mathbb D(X)$ be terminal on $U$. On $Y$, the unit $E\to i_*i^*E$ is invertible by the first assertion. On $U$, it is a map between terminal objects and is again invertible. The two parts cover all objects of $X$, so \textbf{Der 2} finishes the proof. This argument uses no zero object. In a pointed theory the terminal values happen also to be initial, which explains the later terminology ``extension by zero.''

These computations also explain why the argument survives the loss of ordinary diagrammatic information. We do not recover an arbitrary coherent diagram from a list of its values. We construct a specific morphism, the unit or counit, and show that this morphism is invertible at every object. That is exactly what conservativity permits. The corresponding modern statement is \citet[Proposition~1.23]{Groth2013}, with sieve and cosieve interchanged when changing the variance convention.
\end{dcomment}


\noindent\textsc{Remarque} : On retrouve le fait que $i^*i_*(F_Y) \overset{\sim}{\to} F_Y$ est \texttt{[un]} isomorphisme, pour $i$ une immersion fermée. De plus, avec les notations de \textbf{Der 4} (g), $j^*(i_*(F_Y))$ se calcule par le théorème de changement de base (corollaire 2 \texttt{[plutôt 4]} ci-dessus), qui s'applique dans toute situation commutative
\[
\begin{tikzcd}[row sep=large, column sep=large]
X' \arrow[r, "q"] \arrow[d, "f'"'] \arrow[dr, phantom, "\text{\scriptsize cart}"] & X \arrow[d, "f"] \\
Y' \arrow[r, "p"] & Y
\end{tikzcd}
\]
quand $p$ est une immersion ouverte (et $\gamma = \id$) (alors $X'/y' \to X/p(y')$ est même un isomorphisme!), et on trouve, revenant aux notations de \textbf{Der 4} (g),
\[
\begin{tikzcd}
\emptyset \arrow[r, "q"] \arrow[d, "f'"'] & Y \arrow[d, hook, "i"] \\
U \arrow[r, hook, "j"'] & X
\end{tikzcd}
\]
\textit{i.e.} $j^*(i_*(F_Y)) \simeq f'_*(\varphi)$, où $\varphi = q^*(F_Y)$ est l'unique élément de la catégorie (équivalente à une catégorie ponctuelle, par \textbf{Der 1}) $\DD(\emptyset)$. D'autre part, $f'_*$ transforme objet final en objet final, d'où $j^*i_*(F_Y)$ est objet final.

Inversement, soit $F$ dans $\DD(X)$, tel que $j^*(F) = e_{\DD(U)}$. Considérons
\[
F \xrightarrow{u} i_*i^*F \qquad\quad \text{(adjonction) ;}
\]
je dis que c'est un isomorphisme. En effet, par l'axiome \msp{page 52} de localisation \textbf{Der 2}, il suffit de prouver que $i^*(u)$ et $j^*(u)$ sont \texttt{[des]} isomorphismes. Pour $i^*(u)$ cela résulte du fait que $i^*$ soit foncteur de localisation, et pour $j^*$ du fait que $j^*(F)$ et $j^*(i_*i^*(F))$ sont des objets finaux dans $\DD(U)$.

\printedpage{30}
% head: A. Grothendieck : Dérivateurs, Ch. I, Édition des Universités de Montpellier II et Paris VII
% foot: Version 8 juillet 2001

\noindent\msp{page 53}

\dgrblock{dgrb27}{Properness as a condition on all lifts}
\begin{dcomment}{exposition}{p52}{Properness as a condition on all lifts}
The geometric terminology begins with a computational requirement. A right direct image over $y$ is initially computed on the comma category $X/y$. To compute it on the strict fibre $X_y$ instead, the inclusion $X_y\to X/y$ must be cofinal in the orientation appropriate to contravariant coefficients. Properness asserts this simultaneously for every $y$.

There is a concrete way to express the requirement without derived functors. For $x\in X$ and $\alpha:f(x)\to y$, let $L_f(x,\alpha)$ be the category whose objects are arrows
\[
 v:x\longrightarrow x_1,\qquad f(v)=\alpha,
\]
and whose morphisms from $v$ to $v'$ are arrows $t:x_1\to x_1'$ with $tv=v'$ and $f(t)=1_y$. This is a category of lifts of one fixed base arrow, with a fixed source upstairs. There is a canonical isomorphism
\[
 L_f(x,\alpha)\cong (x,\alpha)\backslash X_y.
\]
Indeed, a morphism $(x,\alpha)\to(x_1,1_y)$ in $X/y$ is exactly such a lift $v$, and a morphism between two of them is exactly the triangle defining $t$. Thus $f$ is proper precisely when every $L_f(x,\alpha)$ has weakly contractible nerve. In particular, a lift must exist, all choices must be connected, and their higher homotopy must vanish. Mere existence of lifts is insufficient.

For smoothness one fixes the \emph{target} upstairs instead. Given $x\in X$ and $\beta:y\to f(x)$, the incoming lift category consists of $v:x_0\to x$ over $\beta$ and triangles over $1_y$. Smoothness says that these categories are weakly contractible. The first display on manuscript p.~54 should consequently have target $y\backslash X$, not $y\backslash Y$; the latter is what the edition actually prints. The error is typographical, since the succeeding characterization by fibres of $X/x\to Y/f(x)$ has the correct meaning. In the modern literature this elementary equivalence is stated in \citet[Definition~3.2.1 and Proposition~3.2.3]{Maltsiniotis2020}; \citet[Theorem~1.2]{Maltsiniotis2005Asphericite} gives its formulation by base change.

The stability under strict pullback can now be proved without any cohomology. In a pullback $X'=X\times_Y Y'$, choose $x'=(x,y')$ and an arrow $\alpha':y'\to z'$ in $Y'$. A lift of $\alpha'$ out of $x'$ is exactly a lift of $p(\alpha')$ out of $x$, its second component already being prescribed. The same correspondence identifies every morphism between lifts. The lift categories are therefore isomorphic. Reversing all arrows proves the smooth case.

Several small consequences deserve to be read literally. A cosieve inclusion is proper because every allowed outgoing arrow has one lift; a sieve inclusion is smooth by the dual argument. If $f$ is proper, nonemptiness of $L_f(x,\alpha)$ shows that every arrow out of an object of its image is itself in the image. Thus the strict image is a full cosieve. Finally, $X\to e$ is both proper and smooth: the lift categories are $x\backslash X$ and $X/x$, which have initial and terminal objects. Product projections inherit both properties by strict pullback.
\end{dcomment}


\subsection*{Morphismes homotopiquement propres, homotopiquement lisses.}

\noindent Soit $f : X \to Y$ \texttt{[un]} morphisme de catégories, considérons $y_0 \in \Ob\, Y$, et le foncteur canonique
\[
X_{y_0} \xrightarrow{\varphi} X/y_0 \qquad\qquad\qquad \varphi(x) = (x, \id_{f(x)}) \tag*{(1)}
\]
($X_{y_0}$ fibre de $X$ en $y_0$). Soit $(x_0, u_0) \in \Ob\, X/y_0$ donc $x_0 \in \Ob\, X$, $u_0 : f(x_0) \to y_0$. Considérons la catégorie coinduite par $\varphi$ sur $(x_0, u_0)$
\[
(x_0, u_0)\backslash(X_{y_0}) \quad ,
\]
pour tester si $\varphi$ est une hot-équivalence costricte \textit{i.e.} homologiquement \texttt{[plutôt co\-ho\-mo\-lo\-gi\-que\-ment]} cofinal. On va l'interpréter via le foncteur induit par $f$
\[
x_0\backslash X \xrightarrow{\psi} f(x_0)\backslash Y
\]
allant de la catégorie des spécialisations de $x_0$ dans $X$ vers celle des spécialisations de $f(x_0)$ dans $Y$. On trouve alors un isomorphisme canonique de la catégorie coinduite $(x_0, u_0)\backslash X_{y_0}$ par $X_{y_0}$ au dessus de l'objet $(x_0, u_0)$ de $X/y_0$ et de la catégorie fibre $(x_0\backslash X)_{u_0}$ de $x_0\backslash X$ au dessus de $f(x_0)\backslash Y$ en $u_0$. Donc

\noindent\textbf{Proposition.} \textit{Pour que les foncteurs $\varphi : X_{y_0} \to X/y_0$ ($y_0 \in \Ob\, Y$) soient des hot-équivalences costrictes, il faut et il suffit que les fibres des foncteurs $x_0\backslash X \to f(x_0)\backslash Y$ induits par $f$ soient contractiles.}

\noindent\msp{page 54}
On dit alors que $f$ est \textit{homotopiquement propre.} La notion duale s'exprime, soit par la condition que les foncteurs induits par $f$
\[
X_{y_0} \longrightarrow y_0\backslash Y
\]
soient des hot-équivalences strictes \textit{i.e.} qu'ils soient cohomologiquement \texttt{[plutôt ho\-mo\-lo\-gi\-que\-ment]} cofinaux, soit que les fibres des morphismes induits sur les ``catégories localisées'' (catégories des générisations)
\[
X/x_0 \longrightarrow Y/f(x_0)
\]
soient contractiles. (N.B. On utilise ici qu'une catégorie $Z$ est contractile si et seulement si $Z^\circ$ l'est.) On dit alors que $f$ est \textit{homotopiquement lisse.}

\noindent\textbf{Proposition.} \textit{La notion de foncteur propre (resp. lisse) est stable par changement de base, par composition. Les isomorphismes, et plus généralement, les immersions fermées sont propres. L'image d'un morphisme propre est une sous-catégorie fermée.}\footnote[23]{Attention, une équivalence de catégories n'est pas propre! Si $X' \hookrightarrow X$ \texttt{[est]} une inclusion d'une sous-catégorie de $X$, elle n'est propre que si c'est une immersion fermée.}

\noindent N.B. Le foncteur $X \to e$ est toujours propre et lisse, car les catégories $x_0\backslash X$ et $X/x_0$ sont contractiles (ayant objet initial ou objet final). Donc tout foncteur de projection $X \times Y \xrightarrow{pr_2} Y$ est à la fois propre et lisse.

\printedpage{31}
% head: A. Grothendieck : Dérivateurs, Ch. I, Édition des Universités de Montpellier II et Paris VII
% foot: Version 8 juillet 2001

\noindent\msp{page 55}
\textsc{Démonstration}. Stabilité par changement de base :
\[
\begin{tikzcd}
X \arrow[d, "f"'] \arrow[dr, phantom, "\text{\scriptsize cart}"] & X' \arrow[l, "q"'] \arrow[d, "f'"] \\
Y & Y' \arrow[l, "p"]
\end{tikzcd}
\qquad ,
\]
d'où pour $x' \in \Ob\, X'$
\[
\begin{tikzcd}[row sep=large]
X/x \arrow[d] & X'/x' \arrow[l] \arrow[d] \\
Y/y & Y'/y' \arrow[l]
\end{tikzcd}
\qquad ,
\]
où $x$, $y'$, $y$ sont les images de $x'$ dans $X$, $Y'$, $Y$. On vérifie que ce diagramme est \textit{cartésien}, donc les fibres de $X'/x' \longrightarrow Y'/y'$ sont isomorphes à des fibres de $X/x \longrightarrow Y/y$, donc contractiles si celles-ci le sont.

\dgrblock{dgrb28}{Why an equivalence can fail properness}
\begin{dcomment}{examples}{p54}{Why an equivalence can fail properness}
Let $J$ be the groupoid with two objects $0,1$ and exactly one morphism between each ordered pair, and let $i:e\to J$ select $0$. The functor $i$ is fully faithful and essentially surjective, hence an equivalence. Nevertheless, for the unique isomorphism $\alpha:0\to1$, the outgoing lift category $L_i(*,\alpha)$ is empty. There is no object of the strict domain lying over $1$. The incoming lift category for the inverse isomorphism is empty for the same reason, so $i$ is neither proper nor smooth.

This failure is visible directly in base change. Pull $i$ back along $j:e\to J$ selecting $1$. The strict pullback is empty. For an ordinary set $M$, or a space $M$, the right Kan extension $i_*M$ has value $M$ at $1$: the comma category $e/1$ has a single object. The right direct image from the empty pullback is terminal. Hence the comparison at $1$ is
\[
 M\longrightarrow *,
\]
which fails to be invertible, for example when $M$ is a discrete two-point space. This example verifies the caution in footnote~23 and shows exactly where the strictness matters. Replacing the strict pullback by an equivalence-invariant pullback would change the problem under discussion.

For an equivalence $f:X\to Y$, the preceding argument gives a more precise criterion: $f$ is proper if and only if it is surjective on objects, and the same criterion characterizes smoothness. Necessity follows because every object of $Y$ is isomorphic to one in the image and properness supplies a lift of that isomorphism. For sufficiency, fix $\alpha:f(x)\to y$ and choose $x_1$ with $f(x_1)=y$. Full faithfulness gives a unique lift $x\to x_1$ of $\alpha$. Any two lift objects have a unique morphism between them over $1_y$, again by full faithfulness. Their lift category is therefore nonempty and equivalent to a point. The incoming version is identical. This is the criterion recorded in \citet[Remark~3.2.5]{Maltsiniotis2020}.
\end{dcomment}


Dans un cas $X \xrightarrow{f} Y$ on vérifie que pour un foncteur bilocalisé de $f$
\[
X/x \xrightarrow{\varphi=f/x} Y/y \quad ,
\]
tout ``bilocalisé''
\[
(X/x)/\xi \longrightarrow (Y/y)/\eta \qquad\qquad (\xi : x' \to x\,, \quad \eta : y' \to y)
\]
de $\varphi$ s'identifie encore à un bilocalisé de $f$
\[
X/x' \longrightarrow Y/y' \quad .
\]
Donc si $f$ est lisse, les bilocalisés $f/x$ ($x \in \Ob\, X$) le sont aussi. Notons d'autre part.

\noindent\textbf{Proposition.}\footnote[24]{\textit{cf.} amélioration dans cor. 2, page suivante.} \textit{Tout morphisme propre (ou dualement lisse) à fibres contractiles est une hot-équivalence.}

\noindent\msp{page 56}
Cela provient du fait que pour un tel foncteur propre, si $F$ est un préfaisceau, ou complexe de préfaisceaux abéliens sur $X$, la formation des $Rf_*(F)$ commute au passage aux fibres, et de même pour $R^1f_*(F)$ si $F$ est un préfaisceau en groupes (pas même abéliens).

\dgrblock{dgrb29}{Contractible fibres and the missing steps in composition}
\begin{dcomment}{exposition}{p55}{Contractible fibres and the missing steps in composition}
The assertion about contractible fibres does not need a separate theorem detecting weak equivalences by abelian and nonabelian cohomology. Suppose first that $f$ is proper and every $X_y$ is weakly contractible. By definition $X_y\to X/y$ is coaspherical, hence a weak equivalence on nerves by Quillen's Theorem~A. Therefore every $X/y$ is weakly contractible. A second application of Theorem~A, now to $f$, proves that $Nf$ is a weak equivalence. Opposite categories give the smooth case. This is a completion of the proof sketched here, using \citet[\S1, Theorem~A]{Quillen1973}, rather than a correction of the theorem.

The same argument proves the stronger assertion on manuscript p.~56. Base change along $Y/y\to Y$ and $y\backslash Y\to Y$ preserves properness or smoothness and preserves the contractibility of strict fibres. Each base has contractible nerve, since it has a terminal or an initial object. The total categories are $X/y$ and $y\backslash X$, respectively, and the fibre theorem makes each of them contractible. Thus $f$ is both aspherical and coaspherical.

Here is also a full proof of the composition assertion. Let $f:X\to Y$ and $g:Y\to Z$ be proper, fix $x\in X$, and fix $\alpha:gf(x)\to z$. There is a functor
\[
 L_{gf}(x,\alpha)\longrightarrow L_g(f(x),\alpha),
 \qquad (v:x\to x_1)\longmapsto f(v).
\]
The undercategory functor $x\backslash X\to f(x)\backslash Y$ is proper: after fixing an object $v:x\to x_1$, its outgoing lift categories identify with the corresponding outgoing lift categories of $f$ at $x_1$. Its strict fibres are contractible by the definition of properness of $f$. The displayed functor is its strict pullback along $L_g(f(x),\alpha)\to f(x)\backslash Y$. It is consequently proper with contractible fibres, so the preceding theorem makes it a weak equivalence. Its target is contractible because $g$ is proper. Its source is therefore contractible, and this is the required condition for $gf$. Duality proves composition for smooth functors. The argument is the elementary mechanism behind \citet[Propositions~3.2.6 and~3.2.10]{Maltsiniotis2020}.

There is no cancellation theorem with the other two positions in a composite. The functors $[1]\to e$ and $e\to e$ are proper, but their factor $e\to[1]$ selecting $0$ is not: the arrow $0\to1$ has no lift. The statement about sections in the text has precisely this content. It should not be read as a defect of the proof of composition.
\end{dcomment}


\noindent\textbf{Corollaire 1.} \textit{Si $f : X \to Y$ est cohomologiquement propre à fibres contractiles, alors $Y$ est contractile si et seulement si $X$ l'est.}

Grâce à ceci, on trouve que le composé de deux foncteurs propres (lisses) est propre (lisse).

\printedpage{32}
% head: A. Grothendieck : Dérivateurs, Ch. I, Édition des Universités de Montpellier II et Paris VII
% foot: Version 8 juillet 2001

\noindent N.B. Contrairement à ce qui a lieu en topologie pour la notion de propreté, si $gf$ est propre et $g$ propre, $f$ ne l'est pas nécessairement, \textit{i.e.} un $S$-morphisme de $S$-catégories $X$, $Y$ propres sur $S$ n'est pas nécessairement propre. Ainsi une section de $X/S$, pour $X$ propre sur $S$, n'est pas nécessairement propre, \textit{i.e.} n'est pas nécessairement une immersion fermée. C'est déjà faux si $S = e$ (auquel cas $f : X \to e$ est toujours propre).

\noindent\textbf{Corollaire 2.} \textit{Un morphisme propre ou lisse, à fibres contractiles est à la fois une hot-équivalence stricte et costricte.}

En effet, on applique le corollaire 1 à $X/y \to Y/y$ ou $y\backslash X \to y\backslash Y$.

\noindent\textbf{Théorème} (g) \textit{Soit une situation de changement de base dans} \textsf{Dia}
\[
\begin{tikzcd}[row sep=large, column sep=huge]
X \arrow[d, "f"'] \arrow[dr, phantom, "\text{cart}"] & X' \arrow[l, "q"'] \arrow[d, "f'"] \\
Y & Y' \arrow[l, "p"]
\end{tikzcd}
\qquad , \qquad\qquad p^*f_*(F) \xrightarrow{c} f'_*q^*(F)
\]
\textit{($F \in \DD(X)$), $\DD$ dérivateur exact à gauche. Si $f$ est (homotopiquement) propre ou $p$ homotopiquement lisse, alors $c$ est un isomorphisme.}

\noindent\msp{page 57}
La question c'est si pour tout $y' \in Y'$, posant $y = p(y')$ et considérant le foncteur induit par $q$
\[
X/y \longleftarrow X'/y'
\]
le morphisme correspondant
\[
\Gamma_{\DD}(X/y, F) \longrightarrow \Gamma_{\DD}(X'/y', q^*F)
\]
est \texttt{[un]} isomorphisme. Dans le cas $f$ propre, on considère
\[
\begin{tikzcd}[column sep=large]
X_y \arrow[d, "\sim\ \ \text{costricte}"'] & X_{y'} \arrow[l, "\sim"'] \arrow[d, "\sim\ \ \text{costricte}"] \\
X/y & X'/y' \arrow[l]
\end{tikzcd}
\]
d'où la conclusion, puisque par \textbf{Der 5} (g), dans
\[
\begin{tikzcd}[row sep=large, column sep=large]
\Gamma(X_y) \arrow[r] & \Gamma(X_{y'}) \\
\Gamma(X/y) \arrow[u] \arrow[r] & \Gamma(X'/y') \arrow[u]
\end{tikzcd}
\]
les flèches verticales et l'horizontale supérieure sont \texttt{[des]} isomorphismes, donc aussi l'horizontale inférieure.

\printedpage{33}
% head: A. Grothendieck : Dérivateurs, Ch. I, Édition des Universités de Montpellier II et Paris VII
% foot: Version 8 juillet 2001

Dans le cas $p$ lisse, considérons du coté $Y$, $Y'$ le carré de changement de base % [sic] coté -> côté, missing circumflex (cf. "Du côté" below)
\[
\begin{tikzcd}[row sep=large, column sep=large]
Y & Y' \arrow[l] \\
Y/y \arrow[u] & Y'/y' \arrow[l] \arrow[u]
\end{tikzcd}
\]

\noindent\msp{page 58}
d'où le diagramme contigu
\[
\begin{tikzcd}[column sep=small]
X \arrow[dd] & & X' \arrow[ll] \arrow[dd] & \\
& X/y \arrow[ul] & & X'/y' \arrow[ll, crossing over, "\kappa"'] \arrow[ul] \arrow[dd] \\
Y & & Y' \arrow[ll] & \\
& Y/y \arrow[ul] \arrow[from=uu, crossing over] & & Y'/y' \arrow[ll, "\pi"] \arrow[ul]
\end{tikzcd}
\qquad ,
\]
où les carrés latéraux (verticaux) sont cartésiens. On note que $\pi$ est lisse et à fibres contractiles (car $p$ est lisse) donc $\kappa$ aussi. Donc $\kappa$ est une hot-équivalence costricte, donc $\Gamma(X/y) \overset{\sim}{\longrightarrow} \Gamma(X'/y')$, on gagne encore.

\noindent N.B. Sauf erreur, on a une réciproque du théorème, \texttt{[à]} savoir : Si $f : X \to Y$ est tel que le théorème de changement de base pour $f_*(F)$ est vrai pour \textit{tout} dérivateur (il suffit sans doute de \textsf{HOTAB} et \textsf{HOT}) et pour tout $Y' \to Y$, alors $f$ est homotopiquement propre. D'autre part, si $g : Y' \to Y$ est tel que le théorème de changement de base pour $g$ est vrai pout tout $X \to Y$ et tout $F \in \DD(X)$ ($\DD$ dérivateur exact à droite quelconque), alors $g$ est homotopiquement lisse. % [sic] pout tout -> pour tout, misprint

\noindent Du côté homologique les foncteurs lisses $f$ sont exactement ceux pour lesquels $f_!$ commute à tout changement de base. Les foncteurs propres $g$ sont exactement ceux pour lesquels le changement de base par $g$ commute à $f_!$, pour tout $f : X \to Y$ et tout $\DD$ dérivateur exact à gauche.

\noindent\msp{page 59}
Ainsi, les foncteurs propres $f : X \to Y$ se caractérisent de \textit{quatre} façons remarquables
\begin{itemize}[label={}, leftmargin=0pt, itemindent=2\parindent, labelsep=0.3em]
\item[\textit{a})] $f_*$ commute à tout changement de base ($\forall$ dérivateur $\DD$ exact à gauche);
\item[\textit{b})] le changement de base $f^*$ commute à tout $g_!$ ($g : Y' \to Y$, $\DD$ dérivateur exact à droite);
\item[\textit{c})] $\forall y \in \Ob\, Y$, $X_y \to X/y$ est une hot-équivalence costricte \textit{i.e.} est un foncteur homologiquement \texttt{[plutôt co\-ho\-mo\-lo\-gi\-que\-ment]} cofinal;
\item[\textit{d})] $\forall x \in \Ob\, X$ le foncteur induit sur les colocalisés
\[
x\backslash X \longrightarrow f(x)\backslash Y
\]
\end{itemize}

\dgrblock{dgrb30}{The converse, and why spaces suffice}
\begin{dcomment}{status}{p58}{The converse, and why spaces suffice}
The conjectural converse is proved in the subsequent theory. Its exact scope is universal base change: for a fixed $f:X\to Y$, one considers every strict pullback square and every coefficient in the indicated derivator. The class of homotopy exact squares is already detected by the derivator of spaces; \citet[Theorem~3.16]{GrothPontoShulman2014} equate this test with exactness for all derivators and with weak contractibility of the factorization categories. Thus the proposed pair $\mathsf{HOTAB}$ and $\mathsf{HOT}$ can be reduced to $\mathsf{HOT}$ alone. The proper/smooth theory is developed systematically by \citet[\S3.2]{Maltsiniotis2020}. We spell out the small tests that detect the two converse assertions, because they explain why the conjecture has the stated strength.

Fix $y\in Y$ and form the fibre square, with $j:X_y\to X$ and $s:X_y\to e$. In the derivator of spaces its right comparison
\[
 y^*f_*\longrightarrow s_*j^*
\]
is invertible if and only if its left-adjoint mate
\[
 j_!s^*\longrightarrow f^*y_!
\]
is invertible. This is taking the adjoint of a transformation between right adjoints, so invertibility is preserved in both directions. Apply the left mate to the terminal space $*$. At $x\in X$ its target is the discrete space $\operatorname{Hom}_Y(f(x),y)$, because $y_!*$ is the representable presheaf $\operatorname{Hom}_Y(-,y)$. Its source is
\[
 (j_!*)_x\simeq
 \coprod_{\alpha:f(x)\to y}|N L_f(x,\alpha)|.
\]
Indeed, the left pointwise formula is the homotopy colimit of the point over $x\backslash X_y$ (with the opposite convention making no difference to its realization), and this category is the disjoint union of the lift categories indexed by $\alpha$. The mate collapses each summand to its indexing point. It is an equivalence precisely when every summand is weakly contractible. This proves properness from universal right base change, and even shows that the fibre squares alone suffice for this converse.

For the smooth converse, fix $p:Y'\to Y$ and assume that base change along $p$ commutes with every right direct image in the derivator of spaces. Choose $f:e\to Y$ selecting any $y_0$, and evaluate at $y'\in Y'$. For every space $M$, the comparison becomes
\[
 \operatorname{Map}(S,M)\longrightarrow
 \operatorname{Map}(|N(Y'_{y_0}/y')|,M),
 \qquad S=\operatorname{Hom}_Y(y_0,p(y')),
\]
where $S$ is discrete. The map underlying this comparison sends an arrow $z\to y'$ with $p(z)=y_0$ to its image in $Y$. If the displayed map is an equivalence for every $M$, Yoneda in the homotopy category of spaces says that this underlying map to $S$ is an equivalence. Its fibre over $\beta:y_0\to p(y')$ is the nerve of the incoming lift category at $y'$. All these categories are therefore weakly contractible, which is smoothness.

The converse statements thus require no conjectural detection principle left over from these pages. There is, however, a notational qualification in the printed paragraph: a comparison involving $f_*$ requires right Kan extensions, which the manuscript calls a derivator ``\`a gauche.'' The subsequent sentence concerning $f_!$ requires left Kan extensions. Where the edition switches these adjectives, the adjunctions identify the intended assertion without ambiguity. Passing to the left-adjoint mate also explains condition~(b) on manuscript p.~59: for a pullback square, properness of the right vertical functor controls
\[
 q_!f'^*\longrightarrow f^*p_!,
\]
just as it controls $p^*f_*\to f'_*q^*$. These are two adjoint forms of one exactness statement.
\end{dcomment}


% ---- from der1_p036-040.tex
\printedpage{34}
% head: A. Grothendieck : Dérivateurs, Ch. I, Édition des Universités de Montpellier II et Paris VII
% foot: Version 8 juillet 2001
\noindent est à fibres contractiles.

Autre situation standard où le changement de base marche (mais où les foncteurs qui interviennent n'ont aucune raison d'être propres ou lisses).

Soient
\[
f : I \longrightarrow I' \quad , \qquad g : J \longrightarrow J'
\]
dans $\Dia$, d'où le diagramme (comparer page 23)
\[
\begin{tikzcd}
X' = I \times J \arrow[r, "G'"] \arrow[d, "F'"'] & I \times J' = X \arrow[d, "F"] \\
Y' = I' \times J \arrow[r, "G"] & I' \times J' = Y
\end{tikzcd}
\quad .
\]
Ce diagramme est tel que pour tout dérivateur à gauche

\noindent\msp{page 60}

\dgrblock{dgrb31}{Independent changes of the two variables}
\begin{dcomment}{exposition}{p59}{Independent changes of the two variables}
The product square illustrates exactness that does not arise from properness of either original functor. Its proof is worth expanding because it also checks the correction in brackets on manuscript p.~60. Fix
\[
 x_0=(i_0,j'_0),\quad y'_0=(i'_0,j_0),\quad
 \alpha=(u'_0:f(i_0)\to i'_0,\ v'_0:j'_0\to g(j_0)).
\]
An object of the factorization category is specified by $i\in I$, $j\in J$ and arrows
\[
 a:i_0\to i,\quad r:f(i)\to i'_0,\quad
 b:j'_0\to g(j),\quad s:j\to j_0
\]
with $r f(a)=u'_0$ and $g(s)b=v'_0$. The first equation and the second are independent. The same separation holds for morphisms, giving the isomorphism
\[
 K\cong
 (i_0,u'_0)\backslash(I/i'_0)
 \ \times\ 
 (j'_0\backslash J)/(j_0,v'_0).
\]
The first factor has initial object $((i_0,u'_0),1)$; the second has terminal object $((j_0,v'_0),1)$. Their nerves are therefore contractible, as is the product. This proves the desired base-change isomorphism. The product need not itself have an initial object, which is why the second part of the edition's bracketed correction is necessary as well as the corrected formula.

The result says that a homotopy Kan extension in one variable can be performed while retaining the other variable as a parameter. Restriction in that other variable commutes with the extension. It is one reason that shifted derivators $\mathbb D^J(I)=\mathbb D(I\times J)$ provide an efficient language for parametrized constructions. It does not say that an arbitrary homotopy limit commutes with an arbitrary homotopy colimit. The latter exchange is the stronger question raised earlier in the chapter and requires additional hypotheses.
\end{dcomment}


\noindent on a le changement de base
\[
G^*F_* \simeq F'_*{G'}^* \quad ,
\]
la raison ici étant que les foncteurs induits
\[
X'/y' \longrightarrow X/y
\]
sont des hot-équivalences \underline{\underline{co}}strictes, \textit{i.e.} que les catégories $K$ (page 49) sont contractiles. En effet, si $x_0 = (i_0, j'_0) \in \Ob X$, $y'_0 = (i'_0, j_0) \in \Ob Y'$, $u_0 = (u'_0, v'_0) : F(x_0) \to G(y'_0)$, où $u'_0 : f(i_0) \to i'_0$, $v'_0 : j'_0 \to g(j_0)$ sont donnés, on vérifie que la catégorie $K$ cherchée est isomorphe simplement à $(i_0, u'_0)\backslash(I/i'_0) \times (j_0, v'_0)\backslash(J/j'_0)$ qui est contractile ayant un objet initial \texttt{[plutôt $(i_0, u'_0)\backslash(I/i'_0) \times (j'_0\backslash J)/(j_0, v'_0)$ qui est contractile étant produit d'une catégorie ayant un objet initial et d'une catégorie ayant un objet final]}. Ouf!

Je m'aperçois d'un très gros canular. Je croyais me rappeler que pour $f : X \to Y$ qui est une hot-équivalence stricte, on avait pour tout $F \in \DD(Y)$, $F \isoto f_*f^*F$ (par exemple pour $\DD = \mathsf{Der}(\mathit{Ab})$, ou $\DD = \Hot$), \textit{i.e.} $f^*$ pleinement fidèle, d'où aussi que $\Gamma_{\DD}(Y, F) \to \Gamma_{\DD}(X, f^*F)$ est \texttt{[un]} isomorphisme. Mais ce n'est vrai que si $F$ est un coefficient \textit{constant} sur $Y$, \textit{i.e.} de la forme $q^*(M)$ ($q : Y \to e$, $M \in \DD(e) = \cA_{\DD}$). C'est faux, pour $F$ général, déjà dans le cas de l'immersion $f : \{\varphi\} \to Y$, où $\varphi$ est un objet initial. Alors, $f$ est toujours une hot-équivalence stricte, mais l'homomorphisme $\Gamma(Y, F) \to \Gamma(X, f^*F)$ n'est en général pas un isomorphisme.

\clearpage
\printedpage{35}
% head: A. Grothendieck : Dérivateurs, Ch. I, Édition des Universités de Montpellier II et Paris VII
% foot: Version 8 juillet 2001
\noindent\msp{page 61}

\dgrblock{dgrb32}{A cone construction exposing the homotopy obstruction}
\begin{dcomment}{development}{p60}{A cone construction exposing the homotopy obstruction}
The lifting criterion allows an entire space of obstructions to be placed into a functor over the walking arrow. The criterion itself is established mathematics. We did not find the following explicit organization, including its universal ordinary base-change test, in the passages checked in \citet[\S1 and \S4]{Maltsiniotis2005Asphericite} and \citet[\S3]{GrothPontoShulman2014}. We record the construction and its proof as a development of the present discussion, without a claim of priority.

\textbf{Proposition (the cone test).}\ours
All categories are small relative to the fixed indexing universe, and completeness below refers to its small limits. Let $K$ be a small category and let $X=K^{\triangleleft}$ be obtained by freely adjoining an initial object $o$. Define
\[
 f:X\longrightarrow[1],\qquad f(o)=0,\quad f(k)=1\quad(k\in K).
\]
Then $f$ is smooth and a weak equivalence on nerves. It is proper if and only if $NK$ is weakly contractible. If $K$ is nonempty and connected, every strict base-change square for $f$ satisfies right Beck--Chevalley for ordinary diagrams in every complete category. If in addition $NK$ is not weakly contractible, the fibre square at $1$ fails right Beck--Chevalley for the derivator of spaces.

\emph{Proof.} For an incoming arrow into $o$, the only relevant base arrow is $1_0$, and its lift category is a point. For a target $k\in K$ and base arrow $0\to1$, the unique lift is $o\to k$, so again the lift category is a point. For base arrow $1_1$, the incoming lift category is $K/k$, with terminal object $1_k$. This proves smoothness. Both $X$ and $[1]$ have initial objects; their nerves are contractible, and every functor between them is a weak equivalence on nerves.

The outgoing lift categories are equally explicit. At $o$ over $1_0$ the lift category is a point; at $k\in K$ over $1_1$ it is $k\backslash K$, with initial object $1_k$. The remaining lift category, at $o$ over $0\to1$, is exactly $K$. Thus properness is equivalent to weak contractibility of $NK$.

For the assertion about ordinary base change, let $p:Y'\to[1]$ be arbitrary and let $X'=X\times_{[1]}Y'$. Fix $y'\in Y'$ and consider
\[
 h:X'/y'\longrightarrow X/p(y').
\]
If $p(y')=0$, its target is the one-object category consisting of $o$, and its source is $Y'/y'$, since every arrow into $y'$ must lie over $1_0$. The source has a terminal object. If $p(y')=1$, identify $X/1$ with $X$. For $k\in K$ there is an isomorphism
\[
 k\backslash h\cong(k\backslash K)\times(Y'_1/y'),
\]
where $Y'_1$ is the full subcategory of objects sent to $1$. Both factors are nonempty and connected: the first has an initial object, the second a terminal object. For the remaining object $o$, the comma category $o\backslash h$ is $X'/y'$ itself. It contains the full subcategory $K\times\{y'\}$, which is connected. Every object $(a,z\to y')$ maps to this subcategory: if $a\in K$, use $(a,y')$; if $a=o$, choose any $k\in K$ and use $o\to k$. Hence $o\backslash h$ is nonempty and connected as well.

For completeness, this connectedness criterion preserves ordinary limits of contravariant diagrams for the following direct reason. Given a cone with vertex $T$ over $Fh$, define its component at $a$ by choosing $(b,a\to h(b))$ in $a\backslash h$ and composing $T\to F(h(b))\to F(a)$. Morphisms in the comma category give equal composites, and connectedness makes the choice irrelevant. Applying this to an arrow $a\to a'$ proves compatibility of the constructed components. It is the unique cone over $F$ inducing the given cone. Thus restriction induces a bijection on cones with every vertex $T$, proving the limit comparison isomorphism and consequently ordinary right base change.

Finally take $p:e\to[1]$ to select $1$. For a constant space coefficient $M$, the comparison is
\[
 M\simeq\operatorname{Map}(|NX|,M)
 \longrightarrow\operatorname{Map}(|NK|,M).
\]
If this were an equivalence for every $M$, the map $|NK|\to *$ would be an isomorphism in the homotopy category by Yoneda, contradicting the hypothesis. This proves the last assertion.\hfill$\square$

A finite example makes the failure explicit. Take $K$ to have four objects $a_0,a_1,b_0,b_1$, with the four nonidentity arrows $a_i\to b_j$ and no others. Its nerve has four vertices, four edges, and no nondegenerate higher simplices; its realization is a circle. The category $X$ has five objects. All the ordinary base-change comparisons in the proposition are invertible. In the derivator of spaces, take $M=S^1$. The comparison for the fibre at $1$ is
\[
 S^1\longrightarrow\operatorname{Map}(S^1,S^1),
\]
sending a point to the constant map. The source has one connected component, whereas the target has components indexed by the integer degree. The map lands in degree zero and is not surjective on components. Here the distinction between connected lift categories and contractible lift categories accounts for the entire obstruction.
\end{dcomment}


\subsection*{Notions de $\mathbb{D}$-équivalence et de $\mathbb{D}$-contractibilité.}

On va donc reprendre les questions d'acyclicité et de changement de base.

\noindent\textbf{Définition 1.} Soit $\DD$ un prédérivateur, et soit
\[
\begin{tikzcd}
X \arrow[rr, "f"] \arrow[dr, "p"'] & & Y \arrow[dl, "q"] \\
& e &
\end{tikzcd}
\]
dans $\Dia$. On dit que $f$ est une $\DD$-\textit{équivalence} si le foncteur
\[
F \longmapsto f^*F \quad , \qquad \DD(Y) \longrightarrow \DD(X)
\]
induit un foncteur pleinement fidèle sur la sous-catégorie de $\DD(Y)$ formée des $\DD$-coefficients \textit{constants}, \textit{i.e.} si pour tous $M, N \in \Ob \cA = \DD(e)$ (“coefficients absolus”), on a
\[
\sHom(q^*M, q^*N) \isolongto \sHom(p^*M, p^*N) \quad . \tag*{(1)}
\]
Supposons que $p_*$ et $q_*$ existent. On posera, pour un faisceau $F$ sur $X$
\[
H^*_{\DD}(X, F) \overset{\text{déf}}{=} p_*(F) \qquad \text{(et de même } \ H^{\DD}_*(X, F) = p_!(F) \ \text{).}
\]
N.B. Et si $F = p^*M$, $M \in \cA$, on note simplement $H^*_{\DD}(X, M)$, $H^{\DD}_*(X, M)$.

\noindent L'homomorphisme (1) s'interprète par adjonction comme
\[
\begin{tikzcd}[row sep=small]
\sHom(M, q_*q^*N) \arrow[r] \arrow[d, equal] & \sHom(M, p_*p^*N) \arrow[d, equal] \\
\sHom(M, H^*_{\DD}(Y, q^*N)) \arrow[r] & \sHom(M, H^*_{\DD}(X, p^*N))
\end{tikzcd}
\]
déduit de l'homomorphisme
\[
q_*q^*N \longrightarrow p_*p^*N \quad , \qquad \textit{i.e.} \quad H^*_{\DD}(Y, N) \longrightarrow H^*_{\DD}(X, N) \quad . \tag*{(2)}
\]
Donc dire que (1) est un isomorphisme pour tous $M, N$ signifie que (2) est un isomorphisme pour tout $N$. De même, si $p_!$, $q_!$ existent cela équivaut à dire que
\[
p_!p^*N \longrightarrow q_!q^*N \quad , \qquad \textit{i.e.} \quad H^{\DD}_*(X, N) \longrightarrow H^{\DD}_*(Y, N) \tag*{$(2')$}
\]
est un isomorphisme dans $\cA$ pour tout $N \in \cA$.

\clearpage
\printedpage{36}
% head: A. Grothendieck : Dérivateurs, Ch. I, Édition des Universités de Montpellier II et Paris VII
% foot: Version 8 juillet 2001
\noindent\msp{page 62}

\dgrblock{dgrb33}{Equivalence measured by constant coefficients}
\begin{dcomment}{exposition}{p61}{Equivalence measured by constant coefficients}
The definition of $\mathbb D$-equivalence deliberately retains less information than an equivalence $\mathbb D(Y)\simeq\mathbb D(X)$. It compares only the full subcategories formed by constant coefficients, while retaining every morphism between them in the coherent diagram categories. Such morphisms can contain cohomological information not visible in the values of the constant diagrams themselves.

Write $p:X\to e$ and $q:Y\to e$, with $p=qf$. If the required right adjoints exist, adjunction gives
\[
 \operatorname{Hom}_{\mathbb D(Y)}(q^*M,q^*N)
 \cong\operatorname{Hom}_{\mathbb D(e)}(M,q_*q^*N).
\]
The analogous formula for $X$ identifies the restriction map with the map represented by $q_*q^*N\to p_*p^*N$. It is bijective for every $M$ exactly when this representing map is invertible. This is ordinary Yoneda in $\mathbb D(e)$; it does not require an enrichment. With left adjoints one instead writes
\[
 \operatorname{Hom}_{\mathbb D(Y)}(q^*M,q^*N)
 \cong\operatorname{Hom}_{\mathbb D(e)}(q_!q^*M,N)
\]
and uses contravariant Yoneda while varying $N$. This proves both characterizations on manuscript p.~62, including the weaker existence hypothesis that only the representing objects for constant coefficients are needed.

The elementary formal properties also hold already for prederivators. A natural transformation $f\Rightarrow g$ induces the identity on the restrictions of constant coefficients, because its composite with $Y\to e$ is the identity transformation. The maps on their Hom sets therefore agree; a zigzag of transformations has the same consequence. For composable functors, the maps on these sets compose. Two-out-of-three for bijections proves two-out-of-three for $\mathbb D$-equivalences. Taking the terminal category as target then gives the characterization of $\mathbb D$-contractibility by full faithfulness of the constant-diagram functor.

For $\mathbb D=\mathsf{HOT}$ the representing objects have familiar meanings:
\[
 H^*_{\mathsf{HOT}}(X,M)\simeq
 \operatorname{Map}(|NX|,M),
 \qquad
 H_*^{\mathsf{HOT}}(X,M)\simeq |NX|\times M.
\]
The first expression is the derived \emph{mapping space}, considered as an object of the homotopy category. It is not the ordinary Hom set of that category: taking only its connected components would lose higher information. This is the intended reading of the printed $\operatorname{Hom}_{\mathcal A}$ in the example. The second formula follows because the homotopy colimit of a constant space is its tensor with the classifying space of the indexing category. These two constructions are adjoint: products with $|NX|$ are left adjoint to mapping spaces out of $|NX|$. Setting $M=*$ in the homological formula proves that a $\mathsf{HOT}$-equivalence is precisely a weak equivalence of nerves; the converse follows by taking products with arbitrary $M$. The coefficient-sensitive definition thus recovers the ordinary homotopy type for this particular derivator, without asserting that every coefficient theory detects the same class of functors.
\end{dcomment}


\noindent En résumé :

\noindent\textbf{Proposition 1.} \textit{Si $p_*$ et $q_*$ existent, i.e. si $H^*_{\DD}(X, \bullet)$ et $H^*_{\DD}(Y, \bullet)$ existent (il suffit même que ça existe pour des coefficients constants), alors $f$ est une $\DD$-équivalence si et seulement si $f$ induit un isomorphisme pour les $\DD$-cohomologies à coefficients constants de $X$ et $Y$. Si $p_!$ et $q_!$ existent, i.e. si $H^{\DD}_*(X, \bullet)$ et $H^{\DD}_*(Y, \bullet)$ existent, $f$ est une $\DD$-équivalence si et seulement si $f$ induit un isomorphisme pour les $\DD$-homologies à coefficients constants de $X$ et $Y$.}

\noindent\textbf{Remarque 1.} Pour que $f$ soit \texttt{[une]} $\DD$-équivalence, il \textit{suffit} que pour tout $F$ \textit{constant} sur $Y$ le morphisme d'adjonction $F \to f_*f^*F$ soit \texttt{[un]} isomorphisme.

\noindent\textbf{Remarque 2.} Si $f$, $f' : X \to Y$ sont dans une même composante connexe de $\uHom(X, Y)$, alors $f$ est $\DD$-équivalence $\Leftrightarrow$ $f'$ $\DD$-équivalence.

\noindent\textbf{Proposition 2.} \textit{Les isomorphismes sont des $\DD$-équivalences. Si $X \overset{f}{\to} Y \overset{g}{\to} Z$, si deux parmis $f$, $g$, $gf$ sont des $\DD$-équivalences, le troisième aussi.} % [sic] parmis -> parmi, spelling

\noindent\textbf{Définition 2.} On dit que $X$ dans $\Dia$ est $\DD$-\textit{asphérique} (ou $\DD$-\textit{contractile}, ou $\DD$-\textit{trivial}) si $X \overset{p}{\to} e$ est une $\DD$-équivalence, \textit{i.e.} si $p^*$ est pleinement fidèle.

Si $p_*$ existe, cela signifie donc que pour $M \in \Ob(\cA)$
\[
M \overset{\mathrm{adj}}{\isolongto} p_*p^*M \quad , \qquad \textit{i.e.} \quad M \isolongto H^*_{\DD}(X, M) \quad ,
\]
et si $p_!$ existe, cela signifie que pour tout $M$
\[
p_!p^*M \overset{\mathrm{adj}}{\isolongto} M \quad , \qquad \textit{i.e.} \quad H^{\DD}_*(X, M) \isolongto M \quad .
\]
\textbf{Corollaire.} \textit{Si $Y$} (resp. $X$) \textit{est contractile, alors $f : X \to Y$ est une $\DD$-équivalence si et seulement si $X$} (resp. $Y$) \textit{est contractile.}

\noindent\textsc{Exemple}. Prenons $\DD = \HOT$. (On peut prendre $\Dia$ = catégorie des petites catégories). Alors, pour $M \in \cA = \Hot$ et $X$ petite catégorie, on a
\[
H^*_{\HOT}(X, M) \simeq \sHom_{\cA}(\hot(X), M) \quad ,
\]
où
\[
\hot : \mathcal{C}\mathit{at} \longrightarrow \cA = \Hot
\]
est le foncteur canonique qui associe à tout $X$ son “type d'homotopie”. On a
\[
H^{\HOT}_*(X, M) = \hot(X) \times M \quad .
\]

\noindent\msp{page 63}

\dgrblock{dgrb34}{How much constant coefficients can detect}
\begin{dcomment}{examples}{p62}{How much constant coefficients can detect}
The definition of $\DD$-equivalence invites a small calculation which makes the dependence on the coefficient theory completely explicit. This calculation concerns represented \emph{prederivators}; no existence of Kan extensions is required. We use $\pi_0I$ for the set of equivalence classes of objects under zigzags of arrows, equivalently the connected components of $NI$.

\textbf{Proposition (coefficient sensitivity of a represented theory).}
Let $A$ be a locally small category and let $\DD_A(I)=\operatorname{Fun}(I^{\mathrm{op}},A)$. For a functor $f:X\to Y$, precisely the following alternatives occur.
\begin{enumerate}[label=\textup{(\roman*)}]
\item If some set $\operatorname{Hom}_A(M,N)$ contains two different elements, then $f$ is a $\DD_A$-equivalence if and only if $\pi_0f:\pi_0X\to\pi_0Y$ is a bijection.
\item If every hom-set has at most one element and at least one hom-set is empty, then $f$ is a $\DD_A$-equivalence if and only if $X$ is empty exactly when $Y$ is empty.
\item If $A$ is empty, or if every hom-set of the nonempty category $A$ is a singleton, then every functor is a $\DD_A$-equivalence. In the latter case $A$ is equivalent to the terminal category.
\end{enumerate}

\emph{Proof.} A natural transformation between constant $A$-valued diagrams with values $M$ and $N$ consists of arrows $a_i:M\to N$. Naturality along any arrow of $I$ says that the two corresponding arrows $a_i$ are equal. Thus, canonically,
\[
\operatorname{Hom}_{\DD_A(I)}(M_I,N_I)
 \cong \operatorname{Hom}_A(M,N)^{\pi_0I}.
\]
Restriction along $f$ is precomposition with $u=\pi_0f$. If $u$ is a bijection, all these restriction maps are bijections. Conversely, suppose a hom-set $H$ contains distinct elements $a,b$. If $u$ is not surjective, choose a component of $Y$ outside its image. The constant $a$-valued function on $\pi_0Y$ and the function which takes value $b$ only at that component have the same restriction. If $u$ is not injective, choose different components $x_1,x_2$ with the same image. The function on $\pi_0X$ taking value $b$ at $x_2$ and $a$ elsewhere cannot be a restriction. Hence bijectivity on $H$-valued functions forces $u$ to be a bijection, proving~(i).

For~(ii), every hom-set is either a singleton or the empty set. Singleton-valued functions never impose a condition. For an empty hom-set, its power by a set $S$ is a singleton when $S$ is empty and is empty otherwise. Therefore the induced map is bijective exactly when $\pi_0X$ and $\pi_0Y$ have the same emptiness status. The hypothesis that an empty hom-set exists makes this condition necessary. For~(iii), there are either no pairs $(M,N)$ to test, or each tested set of transformations is a singleton. Finally, if all hom-sets in a nonempty $A$ are singletons, the unique arrows between any two objects are inverse, since the only endomorphism is the identity. This proves the stated equivalence with the terminal category.\hfill$\square$

For $A=\mathsf{Set}$, or $A=\mathsf{Ab}$, alternative~(i) applies. For the two-element ordered category $0<1$, alternative~(ii) applies: all functors between nonempty indexing categories are equivalences for this coefficient test. More strikingly, if $G$ is a nontrivial group, the functor $BG\to e$ is a $\DD_A$-equivalence for \emph{every} represented prederivator. Both categories have one connected component, so the displayed formula proves the assertion without assumptions on $A$. To see directly why spaces distinguish it, write $B=|N(BG)|$ and take the constant coefficients $M=*$ and $N=B$. The tested map identifies with
\[
 [*,B]\longrightarrow[B,B].
\]
Its image consists of the homotopy class of constant maps. The identity of $B$ is not in this image: it induces the identity on the nontrivial fundamental group $G$, whereas a constant map induces the trivial homomorphism. The identification follows from the constant-coefficient homotopy-limit formula $\operatorname{holim}_{(BG)^{\mathrm{op}}} B\simeq\operatorname{Map}(B,B)$ explained in the preceding source discussion. Thus this functor is not a $\mathsf{HOT}$-equivalence, and constant coefficients in an ordinary functor category and constant coefficients in the localized theory of spaces have quite different detecting power.

The hom-set calculation is the elementary naturality condition for constant functors; compare \citet[\S\S1.4 and~1.7]{Riehl2016}. The connected-component test in the first case is standard. We have written out the thin and empty cases as well, so that no nondegeneracy assumption on the coefficient category is hidden. This is an elementary worked consequence of naturality, not an additional axiom on a derivator.
\end{dcomment}


\noindent On en déduit que $f : X \to Y$ est une $\HOT$-équivalence (au sens de la définition 1) si et seulement si c'est une $\hot$-équivalence (ou “équivalence faible”) \textit{i.e.} si et seulement si
\[
\hot(f) : \hot(X) \isolongto \hot(Y)
\]
\printedpage{37}
% head: A. Grothendieck : Dérivateurs, Ch. I, Édition des Universités de Montpellier II et Paris VII
% foot: Version 8 juillet 2001
est un isomorphisme.

\noindent\textbf{Proposition 2.} \textit{Soit $f : X \to Y$. Conditions équivalentes (si $\DD$} \texttt{[est un]} \textit{dérivateur} à gauche). % Resolved against official supplemental edition p37: a gauche). is upright.

a) \textit{$f$ est une $\DD$-équivalence, et le reste après tout changement de base $Y' \to Y$ qui est un isomorphisme local (par exemple pour $Y/F \to Y$, $F \in \Ob \widehat{Y}$).}

b) \textit{Pour tout $y \in \Ob Y$, le foncteur induit $X/y \to Y/y$ est une $\DD$-équivalence.}

c) \textit{Pour tout $y \in \Ob Y$, $X/y$ est $\DD$-contractile.}\footnote[25]{Autres conditions équivalentes :\par
\indent d) \textit{Pour tout faisceau \underline{constant} $M_Y$ sur $Y$, on a $M_Y \isoto f_*f^*M_Y$.}\par
\indent d$'$) \textit{Comme} d) \textit{avec faisceau \underline{localement} constant.}}

\noindent N.B. On dit que $g : Z \to Y$ est un \textit{isomorphisme local} si pour tout $z \in \Ob Z$, posons $y = g(z)$, le morphisme induit sur les catégories localisées en $z$ et $y$
\[
Z/z \isolongto Y/y
\]
est un isomorphisme. Par exemple toute immersion ouverte, tout foncteur induit $Y/y \to Y$ (foncteur de localisation) sont des isomorphismes locaux.

Il en résulte a) $\Rightarrow$ b) trivialement.

b) $\Leftrightarrow$ c). Résulte du corollaire de la page précédente, si l'on sait que $Y/y$ est $\DD$-contractile\footnote[26]{Cf. cor. 1 de \textbf{Der 5} (g)}. Cela résultera de \textbf{Der 5} ci-dessous.

\noindent\msp{page 64}

c) $\Rightarrow$ a) “Passage du local au global”. Comme l'hypothèse c) est stable par changement de base par immersion ouverte, il suffit de prouver que c) implique que $f$ est une $\DD$-équivalence, \textit{i.e.} que pour tout $F = M_Y$ ($M \in \Ob \cA$), coefficient constant sur $Y$, l'homomorphisme d'adjonction
\[
M_Y \longrightarrow f_*f^*M_Y = f_*M_X
\]
est \texttt{[un]} isomorphisme. Il suffit de prouver que c'est un isomorphisme fibre par fibre (\textbf{Der 2}), donc que pour tout $y \in Y$,
\[
M_y \longrightarrow f_*(M_X)_y
\]
est un isomorphisme. Par \textbf{Der 4} (g), on a
\[
f_*(M_X)_y \isolongto \Gamma(X/y, M) \quad ,
\]
et l'hypothèse que $X/y$ est $\DD$-contractile donne ce qu'on veut.

\noindent\textbf{Proposition 2$'$.} \textit{Dualement, si $\DD$ est un dérivateur à gauche} \texttt{[plutôt à droite]}\textit{, conditions équivalentes sur $f : X \to Y$ :}

a) \textit{$f$ $\DD$-équivalence,} \texttt{[et]} \textit{le reste par tout changement de base $g : Y' \to Y$ qui est un “isomorphisme colocal” (i.e. pour tout $y' \in \Ob Y'$, le foncteur induit sur les catégories “colocalisées” $y'\backslash Y' \isolongto y\backslash Y$ (où $y = g(y')$)} \texttt{[est un isomorphisme]}\textit{).}

b) \textit{Pour tout $y \in Y$, le foncteur induit $y\backslash X \to y\backslash Y$ est une $\DD$-équivalence.}

c) \textit{Pour tout $y \in Y$, $y\backslash Y$ est $\DD$-contractile.}

\noindent Pour le prouver, on doit utiliser que les $y\backslash Y$ sont $\DD$-contractiles (cf. cor. 1 de \textbf{Der 5} (d)).

\clearpage
\printedpage{38}
% head: A. Grothendieck : Dérivateurs, Ch. I, Édition des Universités de Montpellier II et Paris VII
% foot: Version 8 juillet 2001
\noindent\msp{page 65}

\dgrblock{dgrb35}{Local isomorphisms and local constancy}
\begin{dcomment}{exposition}{p64}{Local isomorphisms and local constancy}
A local isomorphism in this text is a discrete fibration of categories. Indeed, the requirement that $Z/z\to Y/g(z)$ be an isomorphism says that every arrow $y_0\to g(z)$ has exactly one lift with target $z$. Conversely, unique lifts give the required inverse on slice categories, including their morphisms. A category of elements of a presheaf has this property: lifting an arrow into an element means taking its restriction. The projection $Y/y\to Y$ is the representable case.

This description supplies a small omitted step in the implication (c)$\Rightarrow$(a). The text mentions stability under open immersions, while (a) quantifies over all local isomorphisms. If $g:Y'\to Y$ is a discrete fibration and $X'=X\times_Y Y'$, then for every $y'\in Y'$ there is a canonical isomorphism
\[
 X'/y'\cong X/g(y').
\]
To recover the left-hand object from $(x,\alpha:f(x)\to g(y'))$, lift $\alpha$ uniquely with target $y'$. Its source $y'_0$ satisfies $g(y'_0)=f(x)$ and gives the required pair $(x,y'_0)$. Unique lifting also identifies all arrows. Hence condition (c) is stable under every base change in (a), not merely the sieve inclusions mentioned in the proof. The remainder of the argument is complete: the unit on a constant coefficient, evaluated at $y$, is
\[
 M\longrightarrow\Gamma_{\mathbb D}(X/y,M),
\]
which is invertible by $\mathbb D$-contractibility, and \textbf{Der 2} makes the unit invertible globally.

The locally constant strengthening in footnote~25 has a similarly short proof. Let $F\in\mathbb D(Y)$ have invertible structure maps. On $Y/y$, the natural transformation from the projection to the constant functor at $y$ induces, in the presheaf convention, a morphism from the constant coefficient with value $F_y$ to the restriction of $F$. Every component is invertible. By \textbf{Der 2} this is an isomorphism of coherent coefficients. Its pullback to $X/y$ is therefore constant as well, so the preceding argument applies. Conversely, a statement for every locally constant coefficient includes every constant coefficient. This proves the equivalence claimed in the footnote, rather than requiring local constancy to be an additional axiom.

In Proposition~$2'$ the displayed condition (c) must read $y\backslash X$, as duality and condition (b) require. The edition prints $y\backslash Y$. Read literally, that condition holds automatically because $y\backslash Y$ has an initial object, and cannot characterize $f$. For example, $\emptyset\to e$ fails to be an equivalence for the derivator represented by sets although the printed (c) holds. The following sentence about the contractibility of $y\backslash Y$ is correct: that is the \emph{base} category used to pass between (b) and the repaired (c).

The closure assertions on manuscript p.~65 now have transparent proofs. Pullback preserves discrete fibrations. After a local base change over $Z$, a composite of strict $\mathbb D$-equivalences remains a composite of $\mathbb D$-equivalences. The same argument and two-out-of-three prove cancellation when $f$ and $gf$ are strict. For a categorical equivalence $f$, the pullback along a discrete fibration is again an equivalence: full faithfulness follows by uniquely lifting the arrow upstairs through $f$, and essential surjectivity follows by lifting an isomorphism with prescribed target through the discrete fibration. Any prederivator preserves categorical equivalences. The assertion about equivalences therefore needs no stronger hypothesis at this point.
\end{dcomment}


\noindent\textbf{Définition 3.} Soit $\DD$ un prédérivateur (peut-être même un dérivateur à gauche ou à droite). On dit que $f$ est une $\DD$-équivalence \textit{stricte} si $f$ satisfait la condition a) de \texttt{[la]} proposition 2. On dit que c'est une $\DD$-équivalence \textit{costricte} si $f$ satisfait celle de \texttt{[la]} proposition $2'$; on dit aussi parfois que $f$ est $\DD$-cofinal. (La condition duale, \textit{i.e.} $\DD$-équivalence stricte, s'appelerait co-$\DD$-cofinal!)\footnote[27]{Meilleure terminologie : $f$ est $\DD$-\textit{asphérique} ($\DD$-équivalence stricte), $\DD$-\textit{coasphérique} ($\DD$-équivalence co-stricte).} % [sic] s'appelerait -> s'appellerait, spelling

\noindent\textbf{Proposition 3.} \textit{Soit $X \overset{f}{\to} Y \overset{g}{\to} Z$. Si $f, g$ sont des $\DD$-équivalences strictes, alors $gf$ aussi. Si $gf$ et $f$ sont des $\DD$-équivalences strictes, alors $g$ aussi. Si $\DD$} \texttt{[est un]} \textit{dérivateur à gauche ou à droite, toute équivalence} \texttt{[de catégories]} \textit{est une $\DD$-équivalence stricte.} (Car elle a un adjoint à droite.)

On utilise le fait que la notion d'isomorphisme local est stable par changement de base.

\noindent\msp{page 66}

\noindent\textbf{Proposition 4.} \textit{Soient $\DD$ un dérivateur à droite, $f : X \to Y$ un foncteur dans} $\Dia$\textit{, tel que pour $y \in Y$, $y\backslash X$ soit $\DD$-contractile (donc $f$ une $\DD$-équivalence costricte, sous réserve que $\DD$ soit un dérivateur à droite \ldots\ ). Alors pour tout $F$ dans $\DD(Y)$ (pas nécessairement à coefficients constant) on a} % [sic] constant -> constants, agreement
\[
H^*_{\DD}(Y, F) \isolongto H^*_{\DD}(X, f^*(F))
\]
\textit{est un isomorphisme. Réciproquement, si la condition est vraie pour tout $F$, alors les $y\backslash X$ sont $\DD$-contractiles.}

\noindent\textsc{Démonstration}. Il revient au même de dire que pour tout $M \in \Ob \cA = \DD(e)$, on a
\[
\sHom(M, H^*_{\DD}(Y, F)) \overset{\mathrm{iso}}{\longrightarrow} \sHom(M, H^*_{\DD}(X, f^*(F))) \ .
\]
Mais par adjonction, cette flèche s'interprète comme
\[
\sHom(M_Y, F) \longrightarrow \sHom(M_X, f^*F) \tag*{$(*)$}
\]
D'autre part, si $f_!$ existe, on a
\[
\sHom(M_X, f^*F) = \sHom(f_!(M_X), F) \ ,
\]
et l'homomorphisme $(*)$ est déduit de l'homomorphisme d'adjonction
\[
f_!(M_X) = f_!(f^*(M_Y)) \longrightarrow M_Y \ , \tag*{$(**)$}
\]
en prenant les $\sHom$ dans $F$. Donc la condition équivaut à dire que cet homomorphisme est \texttt{[un]} isomorphisme, quel que soit le coefficient constant $M_Y$ sur $Y$. Pour voir si c'est un isomorphisme, on regarde fibre par fibre (\textbf{Der 2}), or les fibres s'interprètent comme $H^{\DD}_*(y\backslash X, M_X)$ (par \textbf{Der 4} (d)) et comme $M$.


\dgrblock{dgrb36}{Cofinality and the final Yoneda argument}
\begin{dcomment}{exposition}{p66}{Cofinality and the final Yoneda argument}
The final proposition explains why coasphericity controls arbitrary coefficients although $\mathbb D$-equivalence was defined using constant ones. Its proof varies the \emph{target} coefficient while keeping the source constant. That change in the Yoneda argument is what strengthens the conclusion.

There is one existence qualification. A derivator ``\`a droite'' in the terminology of these pages supplies $f_!$, not the absolute right Kan extensions used to define $H^*_{\mathbb D}$. The cohomological statement consequently also requires the existence of $p_*$ and $q_*$, where $p:X\to e$ and $q:Y\to e$. Without them the Hom-set formulation in the proof remains meaningful and gives the correct one-sided statement. We prove that formulation first.

Suppose left Kan extensions and their pointwise formulas exist. For $M\in\mathbb D(e)$, abbreviate $p^*M$ and $q^*M$ to $M_X$ and $M_Y$. Adjunction identifies the map
\[
 \operatorname{Hom}_{\mathbb D(Y)}(M_Y,F)
 \longrightarrow
 \operatorname{Hom}_{\mathbb D(X)}(M_X,f^*F)
\]
with precomposition by the counit
\[
 \varepsilon_M:f_!M_X\longrightarrow M_Y.
\]
For fixed $M$, this is a bijection for every $F$ if and only if $\varepsilon_M$ is invertible, by Yoneda in $\mathbb D(Y)$. By conservativity and the left pointwise formula, that happens for every $M$ if and only if
\[
 H_*^{\mathbb D}(y\backslash X,M)\longrightarrow M
\]
is invertible for every $y$ and $M$. By the definition and characterization of $\mathbb D$-contractibility, this is precisely the requirement that every $y\backslash X$ be $\mathbb D$-contractible. Thus the forward implication and the converse in the Hom-set statement are both proved using only the left Kan extensions.

Now assume in addition that $p_*$ and $q_*$ exist. Adjunction identifies the same Hom comparison with
\[
 \operatorname{Hom}_{\mathbb D(e)}(M,q_*F)
 \longrightarrow
 \operatorname{Hom}_{\mathbb D(e)}(M,p_*f^*F).
\]
Varying $M$ and applying Yoneda once more makes its bijectivity equivalent to invertibility of $q_*F\to p_*f^*F$. This proves the cohomological statement exactly as intended, with its existence hypothesis made explicit. No additional repair of the algebraic proof is needed. In modern terminology this is the cofinality criterion for the relevant homotopy limits, with the comma orientation dictated throughout by the presheaf convention.
\end{dcomment}
\begin{otherlanguage}{english}
\dgroverview{I-a}{Chapter I \textperiodcentered\ \textup{I-a} \textperiodcentered\ Overview}
\selectlanguage{english}
\noindent\textbf{Chapter I, first part: from presheaf operations to homotopical base change}

These pages construct a language in which a homotopy theory can carry its own diagram categories. The point of departure is elementary: a functor of indexing categories gives restriction of presheaves, restriction has Kan-extension adjoints when the required limits exist, and special inclusions admit further adjoints when the coefficient category is pointed. What changes in passing to homotopy theory is the meaning of a diagram. A diagram in a homotopy category records only the images of objects and arrows after localization. Localizing a category of diagrams retains additional morphisms and coherence. Grothendieck proposes to organize all these localized categories together, so that a universal construction can be performed on a coherent diagram and remain compatible with further parameters.

The historical motivation is stated explicitly in \emph{Pursuing Stacks}, \S69: a single derived category does not provide an adequate setting for the trace and determinant identities encountered around SGA~5 \citep[\S69]{Grothendieck1983}. Categories of filtered complexes and diagrams of complexes supply more information, but they too must be organized into a structure stable under introducing new diagram shapes. The letter to Thomason of 2 April 1991 describes the renewed investigation of this intrinsic structure \citep[pp.~2--4]{Grothendieck1991Letter}. Heller's homotopy theories form a closely related independent axiomatic development \citep{Heller1988}. The present opening chapter is exploratory: proposed axioms are changed, conclusions are withdrawn, and the final definitions are reached through computations that reveal which hypotheses are actually needed.

\medskip\noindent\textbf{The convention that governs the calculations.}
We work in a fixed universe of small indexing categories; a restricted domain must contain every opposite, product, coproduct, and comma category used in a particular assertion. The manuscript uses presheaves,
\[
 \DD_A(I)=\operatorname{Fun}(I^{\mathrm{op}},A),
 \qquad f_!\dashv f^*\dashv f_*.
\]
Thus an object over $[1]=(0\to1)$ has an underlying arrow $F_1\to F_0$. Restriction reverses both functors and natural transformations, so the explicit modern source is $\Dia^{\mathrm{coop}}$. The manuscript calls a theory with right Kan extensions a derivator \emph{\`a gauche}, because these are the operations of left exactness. A theory with left Kan extensions is called \emph{\`a droite}. The systematic axioms starting at manuscript p.~41 replace the earlier provisional numbering. These changes of convention are substantive reading instructions: a remembered modern formula with the opposite variance can reverse the cofinality condition.

The ordinary pointwise formulas already display the relevant slices:
\[
 (f_*F)_y=\lim_{(f/y)^{\mathrm{op}}}F,
 \qquad
 (f_!F)_y=\operatorname*{colim}_{(y/f)^{\mathrm{op}}}F.
\]
For derivators these are homotopy Kan formulas, expressed by invertibility of specified base-change maps. The map, not merely an isomorphism between its source and target, is part of the assertion. This explains the care devoted to adjunction units, counits, and the two equivalent constructions of the base-change mate.

\medskip\noindent\textbf{Supports, coherent squares, and shifts.}
Open inclusions are sieves and closed inclusions are cosieves. In the presheaf convention a closed inclusion extends by terminal objects under $i_*$; an open inclusion extends by initial objects under $j_!$. Pointedness identifies the two values with zero and permits the exceptional adjoints $i^!$ and $j^?$. The commentary proves these assertions from their universal properties and works out the nonsplit short exact sequence of arrow diagrams underlying the support decomposition. It also gives the finite-ordinal adjunction chain, including its four exceptional endpoints, while retaining the edition's correction of the stronger claim about every monotone map.

The same nonsplit extension exhibits a nonzero morphism between coherent derived diagrams whose components are all zero. Consequently the underlying-diagram functor is conservative without being faithful. This is precisely the distinction the localization axiom requires. Homotopy cartesian squares are coherent squares obtained by right Kan extension from a missing-corner diagram; they need not be ordinary pullbacks in the base homotopy category. Strongness, when imposed, addresses lifting arrows and morphisms of arrows. It is a separate condition from conservativity.

The shift $\DD_I(J)=\DD(I\times J)$ turns an entire coherent $I$-diagram into one coefficient. Its elementary formal properties explain the repeated demand that each new axiom survive shifting. The completed shift theorem and product base change establish the proposed morphisms induced by $f^*$, $f_!$, and $f_*$. At manuscript p.~18 the question whether evaluation at the base category is fully faithful also has a concrete negative answer for unrestricted morphisms: the invisible extension class yields a nonzero modification between morphisms out of the represented derivator of $[1]$ whose component at the point is zero. The example does not impose continuity on these morphisms.

For a Quillen model category, existence of the homotopy Kan extensions is a theorem rather than a formal property of localization. Cisinski supplies a precise general treatment \citep[Theorem~6.11]{Cisinski2003}. In the $\infty$-categorical formulation one takes
\[
 \DD_{\mathcal C}(I)=h\operatorname{Fun}(N(I)^{\mathrm{op}},\mathcal C);
\]
the pointwise extensions of a complete and cocomplete $\infty$-category give the derivator operations \citep[\S5.3]{RiehlVerity2017}. Reconstruction results for the whole prederivator retain their hypotheses on size and morphisms. Renaudin's comparison concerns combinatorial model theories and adjunctions \citep[Theorems~3.3.2 and~3.4.4]{Renaudin2009}, while Arlin distinguishes strict transformations from pseudonatural ones \citep[Theorems~19 and~29]{Arlin2020}. These are stronger invariants than the base category alone.

\medskip\noindent\textbf{The exactness question has a finiteness boundary.}
The connective derived theory $\mathsf{HOTAB}$ is additive and pointed but is not stable. Its abelianization morphism from spaces preserves homotopy colimits. It fails to preserve homotopy pullbacks: for two discrete two-point spaces over a point, the comparison is a map from $\mathbb Z^4$ to $\mathbb Z^3$ with an explicitly displayed rank-one kernel. This separates the true pushout assertion from the false pullback assertion in the same passage.

Stability answers the later ``Perplexit\'e'' after the scope of exactness is fixed. Between stable derivators, preservation of zero objects and cartesian squares is equivalent to preservation of zero objects and cocartesian squares \citep[Corollary~4.17]{Groth2013}. Homotopy finite limits commute with homotopy colimits in the corresponding precise stable theorem \citep[Theorem~2.14]{GrothShulman2017}. Arbitrary products and coproducts still need not commute. The diagonal array of copies of a field gives the comparison $k^{(\mathbb N)}\to k^{\mathbb N}$, which omits every sequence with infinite support. The modern finite theorem therefore does not validate an unrestricted preservation claim for all Kan extensions.

\medskip\noindent\textbf{Cofinality and universal base change.}
The decisive correction in the chapter concerns the coefficient system. Contractibility of every $f/y$ does not make restriction fully faithful on all coefficients. For the inclusion of the initial endpoint into $[1]$, every such comma category is a point, but the unit on $F_1\to F_0$ contains the original structure map. The example $0\to\mathbb Z$ disproves the proposed axiom. The surviving proof of the later cancelled corollary contains a valid slice identification; its last inference uses exactly the false, opposite-orientation cofinality clause. The audit identifies this inference without inventing an additional gap in the calculation.

For right direct image in the manuscript's convention, the correct cofinality test uses arrows from an object of the target to the image of the functor. Applied pointwise to a square, it produces categories of factorizations of a fixed base arrow. Their weak contractibility is the universal homotopy-exactness criterion. Properness specializes it to the comparison between a strict fibre and a comma category; smoothness is the dual incoming-lift condition. The commentary completes the composition proof by identifying the relevant lift-category functor as a strict pullback and applying Quillen's Theorem~A \citep[\S1]{Quillen1973}.

The conjectured converse is established: spaces alone detect universal homotopy exactness, so the proposed simultaneous use of $\mathsf{HOT}$ and $\mathsf{HOTAB}$ is unnecessary \citep[Theorem~3.16]{GrothPontoShulman2014}. Proper and smooth functors are treated systematically by Maltsiniotis \citep[Theorem~1.2]{Maltsiniotis2005Asphericite}. Strict pullbacks matter here. An equivalence of categories that misses an object can fail both properness and smoothness, as the two-object groupoid example shows.

The commentators' marked cone construction makes the strength of the homotopical test explicit. Adjoin an initial object to a small category $K$, send that object to $0$, and send $K$ to $1$. The resulting functor to $[1]$ is always smooth and a weak equivalence of nerves; it is proper exactly when $NK$ is weakly contractible. If $K$ is nonempty and connected, all its strict base-change squares pass the right Kan test for every complete ordinary coefficient category. A four-object crown gives $|NK|\simeq S^1$ and a failure for space coefficients, detected by the degree components of $\operatorname{Map}(S^1,S^1)$. The construction, including its universal ordinary-coefficient assertion, is proved in full and carries no claim of priority.

\medskip\noindent\textbf{The final definition and its scope.}
The corrected notion of $\DD$-equivalence compares every morphism between constant coherent coefficients. By Yoneda it is equivalent, when the representing objects exist, to invariance of constant-coefficient homology or cohomology. Its two-out-of-three property is therefore already formal at the prederivator level. The local versions distinguish discrete fibrations from discrete opfibrations and lead to asphericity and coasphericity. The omitted stability check under all local isomorphisms is supplied by an explicit isomorphism of comma categories. Locally constant coefficients are handled by a canonical morphism on a slice and conservativity.

The last proposition turns this constant-coefficient test into a criterion for arbitrary coefficients by varying the target in a Yoneda argument. Its Hom-set formulation needs left Kan extensions only; its cohomological formulation additionally needs the absolute right Kan extensions that define $H^*$. This is an existence qualification, not a failure of the argument. The questions treated here consequently have precise answers once variance, coefficient class, finiteness, and existence of adjoints are stated. The attribution to Thomason, which the text leaves unspecified, remains a documentary uncertainty.
\end{otherlanguage}
\dgrchunkend

% DER_I-b: Grothendieck's text with the commentary set in, then the overview, inside a group that applies the chunk's own definitions
\dgrpart{1}{I-b}{Part 2 of 2, manuscript pages 67--126}{I-b\enspace\textperiodcentered\enspace Chapter I, part 2 of 2\enspace\textperiodcentered\enspace manuscript pages 67--126}{Chapter I \textperiodcentered\ \textup{I-b}}
\dgrchunkbegin
% DER_I-b: the definitions of the chunk's preamble, applied inside its group
\csname tagsleft@true\endcsname
\let\printedpage\relax\newcommand{\printedpage}[1]{\marginnote{\scriptsize\sffamily ed.\ #1}}
\let\msp\relax\newcommand{\msp}[1]{%
  \begingroup\def\thismsp{#1}\def\blankmsp{page 114 vide}%
  \ifx\thismsp\blankmsp
    \marginnote{\scriptsize\sffamily [114: vide]}[1.3\baselineskip]%
  \else\def\followingmsp{page 115}%
    \ifx\thismsp\followingmsp
      \marginnote{\scriptsize\sffamily [page 115]}[2.6\baselineskip]%
    \else\marginnote{\scriptsize\sffamily [#1]}[1.3\baselineskip]\fi
  \fi\endgroup}
\let\DD\relax\newcommand{\DD}{\ensuremath{\mathbb{D}}}
\let\cA\relax\newcommand{\cA}{\ensuremath{\mathcal{A}}}
\let\Dia\relax\newcommand{\Dia}{\ensuremath{\mathsf{Dia}}}
\let\Ob\relax\newcommand{\Ob}{\ensuremath{\mathsf{Ob}}}
\let\Hom\relax\newcommand{\Hom}{\ensuremath{\mathsf{Hom}}}
\let\hot\relax\newcommand{\hot}{\ensuremath{\mathsf{hot}}}
\let\Hot\relax\newcommand{\Hot}{\ensuremath{\mathsf{Hot}}}
\let\HOT\relax\newcommand{\HOT}{\ensuremath{\mathsf{HOT}}}
\let\isoto\relax\newcommand{\isoto}{\overset{\sim}{\to}}
\let\isolongto\relax\newcommand{\isolongto}{\overset{\sim}{\longrightarrow}}
\let\D\relax\newcommand{\D}{\ensuremath{\mathbb{D}}}
\let\Fl\relax\newcommand{\Fl}{\ensuremath{\mathsf{Fl}}}
\let\Cat\relax\newcommand{\Cat}{\ensuremath{\mathcal{C}at}}
\let\bs\relax\newcommand{\bs}{\backslash}
\let\RR\relax\newcommand{\RR}{\ensuremath{\mathbb{R}}}
\let\sfup\relax\newcommand{\sfup}[1]{\textsf{\upshape #1}}

% ---- from der1_p041-045.tex
\printedpage{39}
% head: A. Grothendieck : Dérivateurs, Ch. I, Édition des Universités de Montpellier II et Paris VII
% foot: Version 8 juillet 2001
\msp{page 67}

\noindent Donc la flèche envisagée $(**)$ s'interprète comme
\[
H_*^{\D}(y\backslash X, M_X) \longrightarrow M
\]
et c'est l'homomorphisme d'adjonction $p_!(p^*(M)) \to M$.
\[
\begin{tikzcd}
X \arrow[rr, "f"] \arrow[dr, "p"{swap}] & & Y \arrow[dl, "q"] \\
& e &
\end{tikzcd}
\]

\noindent Donc c'est un isomorphisme pour tout $M$ si et seulement si $y\backslash X$ est $\D$-contractile.

\noindent On va énoncer ce qu'on a prouvé.

\noindent\textbf{Théorème 1.} \textit{Soit}
\[
\begin{tikzcd}
X \arrow[rr, "f"] \arrow[dr, "p"{swap}] & & Y \arrow[dl, "q"] \\
& e &
\end{tikzcd}
\]
\textit{dans} $\Dia$. \textit{On suppose que} $\D$ \textit{est un prédérivateur satisfaisant} \textbf{Der 2}, \textit{que} $p_*, q_*$ \textit{et} $f_!$ \textit{existent, que} $f_!$ \textit{satisfait à} \textbf{Der 4} (d) \textit{pour le calcul des fibres. Conditions équivalentes :}

a) $f$ \textit{est une} $\D$\textit{-équivalence costricte, i.e.} (prop. 2) \textit{pour tout} $y \in \Ob\, Y$, $y\backslash X$ \textit{est} $\D$\textit{-contractile ;}

b) \textit{pour tout} $F \in \Ob\, \D(Y)$, \textit{l'homomorphisme suivant est} \texttt{[un]} \textit{isomorphisme}
\[
H^*_{\D}(Y, F) \isoto H^*_{\D}(X, f^*(F)) \ ;
\]

c) \textit{pour tout} $M \in \Ob(\mathcal{A} = \D(e))$, \textit{l'homomorphisme d'adjonction est} \texttt{[un]} \textit{isomorphisme}
\[
f_!(M_X) = f_!(f^*(M_Y)) \isoto M_Y \ .
\]

\dgrblock{dgrb37}{The conclusion reached on manuscript page 67}
\begin{dcomment}{exposition}{p67:f_!(M_X) = f_!(f^*(M_Y))}{The conclusion reached on manuscript page 67}
Manuscript pages~60--66 are commented with the first part of the chapter; its blocks ``Equivalence measured by constant coefficients'', ``How much constant coefficients can detect'', ``Local isomorphisms and local constancy'', and ``Cofinality and the final Yoneda argument'' give the definitions and the argument now being completed. The distinction needed below is quite exact. A $\D$-equivalence preserves the morphisms between constant objects; a coaspherical functor preserves cohomology with \emph{every} coefficient object on its target. Neither assertion, by itself, says that inverse image is fully faithful on the entire coefficient category.

For the theorem on manuscript page~67, the relevant test is the counit
\[
f_!p^*M\longrightarrow q^*M,\qquad p=qf.
\]
Applying $\operatorname{Hom}(-,F)$ and then the two adjunctions identifies this test with cohomological invariance for all $F$. Evaluation at $y$ identifies it instead with the augmentation of homology of $(y\backslash f)$ with constant value $M$. These are two readings of the \emph{same natural transformation}. This observation explains the appearance of homology in a theorem about cohomology; it also prevents a spurious gap from being introduced into an otherwise complete Yoneda argument. The theorem explicitly supplies the right adjoints to the terminal projections, whose existence was less explicit in the proposition on page~66.
\end{dcomment}


Je vais renuméroter les axiomes sur le prédérivateurs. % [sic] le prédérivateurs -> les prédérivateurs, article singulier devant un pluriel

\noindent\textbf{Der 1}, \textbf{Der 2} comme avant (p. 45).

\noindent\textbf{Der 3}\footnote[28]{Axiome des catégories d'indices asphériques, \textit{cf.} variantes \textbf{Der 3} a) b) c) page 97, et surtout pages 111-112.} : Soit $X$ dans $\Dia$, $p : X \to e$. Supposons $X$ contractile.\footnote[29]{N.B. Il suffit de supposer que $X$ a un objet final. Mais ce n'est pas un axiome autodual.} Alors
\[
p^* : M \longmapsto M_X \quad , \qquad p^* : \mathcal{A} = \D(e) \longrightarrow \D(X)
\]
est pleinement fidèle.

\printedpage{40}
% head: A. Grothendieck : Dérivateurs, Ch. I, Édition des Universités de Montpellier II et Paris VII
% foot: Version 8 juillet 2001
\msp{page 68}

\noindent\textbf{Der 4} (g) : existence des $f_*$.

\noindent\textbf{Der 4} (d) : existence des $f_!$.

\noindent\textbf{Der 5} (g) Calcul des fibres de $f_*(F)$.

\noindent\textbf{Der 5} (d) Calcul des fibres de $f_!(F)$.

\noindent L'ancien axiome \textbf{Der 5} (g) résulte de \textbf{Der 2}, \textbf{Der 3}, \textbf{Der 4} (g,d), \textbf{Der 5} (d), grâce au théorème précédent.

\dgrblock{dgrb38}{Keeping the axiom numbers separate}
\begin{dcomment}{exposition}{p68:grâce au théorème précédent.}{Keeping the axiom numbers separate}
Grothendieck is reorganising the assumptions after the failure diagnosed on manuscript page~60. The new numbering must be kept separate from both the preceding pages and present-day derivator terminology. Here \textbf{Der 3} asserts that a contractible indexing category has fully faithful constant-diagram functor. The two parts of \textbf{Der 4} assert existence of right and left Kan extensions; the two parts of \textbf{Der 5} give their pointwise comma-category formulas. In the now common numbering, existence is usually (Der3), and the pointwise formulas are (Der4); compare \citet[Definition~1.5]{Groth2013}. The extra strongness axiom used to lift incoherent arrows is a different matter and is not being imposed by this renumbering.

Throughout these comments we retain the manuscript's presheaf convention: a coefficient object over $X$ has underlying diagram $X^{\mathrm{op}}\to\D(e)$. Thus
\[
 (f_*F)_y\simeq H^*_{\D}(X/y,F),\qquad
 (f_!F)_y\simeq H^{\D}_*(y\backslash X,F).
\]
A modern reference using covariant diagrams reverses the comma categories. The right adjoint is still the right adjoint; the difference comes from replacing every indexing category by its opposite.

The new contractibility axiom is not an assertion that the whole category $\D(X)$ collapses to $\D(e)$. It concerns only constant objects and their morphisms. For example, over the walking arrow a coefficient is a coherent arrow, and its two values need not be isomorphic even though the walking arrow has contractible nerve. What the axiom guarantees is that taking a constant object over this shape introduces no new endomorphisms. This is precisely the amount of invariance used when the lifting categories in the following theorem are contractible.

The comment that both homological and cohomological axioms seem indispensable should consequently be read as a description of this proof strategy, not as a proved independence theorem. Modern treatments may begin with a one-sided derivator and formulate the corresponding one-sided conclusions. The argument here deliberately uses the opposite adjoint to make a universal-coefficient test accessible to pointwise calculation.
\end{dcomment}


Un aspect intéressant du théorème précédent, si on l'envisage comme énonçant l'équivalence de a) et b), c'est que cet énoncé ne fait pas appel à des notions homologiques, mais que sa démonstration y fait appel, sous une forme relativement forte (calcul des fibres de $f_!(F)$). Pour les théorèmes de changement de base pour la $\D$-cohomologie relative, \textit{i.e.} pour les $f_*^{\D}$, on a besoin aussi constamment du calcul des fibres des $f_*(F)$. Donc il semble qu'on ne peut se passer d'aucun des axiomes précédents, y compris \textbf{Der 4} (d) et \textbf{Der 5} (d), même si on ne veut s'intéresser qu'au formalisme des $f^*, f_*$ (aspect cohomologique). Par la suite, on va donc supposer ces axiomes satisfaits.

\dgrblock{dgrb39}{The category of lifts behind properness}
\begin{dcomment}{exposition}{p68:Par la suite, on va donc supposer ces axiomes satisfaits.}{The category of lifts behind properness}
There is a useful concrete form of the first two conditions of Theorem~2. Given $x\in X$ and $\alpha:f(x)\to y$, let $L_f(x,\alpha)$ have objects
\[
 h:x\longrightarrow z,\qquad f(z)=y,\qquad f(h)=\alpha.
\]
A morphism from $h$ to $h'$ is a map $k:z\to z'$ with $f(k)=1_y$ and $kh=h'$. These are strict lifting data; replacing equality by an unspecified isomorphism would change the condition. The theorem says that $f$ is $\D$-proper precisely when every $L_f(x,\alpha)$ is $\D$-contractible.

Indeed, if $i_y:X_y\to X/y$ is the canonical inclusion, an object of $X/y$ is exactly a pair $(x,\alpha)$ as above, and
\[
 ((x,\alpha)\backslash i_y)\cong L_f(x,\alpha).
\]
Condition~(a), applied to the definition of coasphericity, is therefore condition~(b). This also proves its stability under arbitrary strict base change. For a pullback $f':X\times_Y Y'\to Y'$, the lift of an arrow $\alpha':y'_0\to y'_1$ from $(x,y'_0)$ is the lift of its image in $Y$, together with the already fixed second coordinate $y'_1$. The resulting lifting categories are isomorphic.

The significance of these categories is stronger than the existence of lifts. An opfibration gives a preferred cocartesian lift, which is an initial object of $L_f(x,\alpha)$. Properness permits a whole category of choices, requiring only that constant coefficients cannot distinguish its homotopy from a point. In particular, requiring merely that each strict fibre $X_y$ be contractible misses the lifting condition altogether.

To obtain the base-change statement, first replace $Y$ by any $Y'\to Y$. At an object $y'$, the right Kan extension is calculated over $X'/y'$. By coasphericity of $X'_{y'}\to X'/y'$ and Theorem~1, restriction induces
\[
 H^*_{\D}(X'/y',F)\xrightarrow{\sim}
 H^*_{\D}(X'_{y'},F|_{X'_{y'}})
\]
for \emph{every} coefficient object $F$ on $X'/y'$. Consequently right direct image is calculated by the actual fibres. After a further base change, the actual fibre is unchanged. The two evaluations of the Beck--Chevalley comparison are therefore identified by their canonical restriction maps; \textbf{Der 2} makes the comparison an isomorphism.

Conversely, suppose the fibre formula holds after localisation $Y/y\to Y$. The category $Y/y$ has terminal object $1_y$, so evaluating the right direct image there identifies its global cohomology with that of $X/y$. The asserted fibre formula then states that restriction along $i_y$ preserves cohomology for every object of $\D(X/y)$. Theorem~1 gives the coasphericity of $i_y$. This completes the equivalence of (a), (b), (c), (c$'$), and (c$''$), without substituting constant coefficients where the argument needs arbitrary ones.
\end{dcomment}


\noindent\textbf{Théorème 2.} \textit{Soit} $f : X \to Y$ \textit{dans} $\Dia$. \textit{Conditions équivalentes :}

a) \textit{Pour tout} $y_0 \in \Ob\, Y$, \textit{le foncteur canonique}
\[
X_{y_0} \overset{\varphi}{\longrightarrow} X/y_0 \quad (\simeq X \times_Y Y/y_0) \ ,
\]
\textit{où} $X_{y_0}$ \textit{désigne la fibre en} $y_0$, \textit{est une} $\D$\textit{-équivalence co-stricte.}

b) \textit{Pour tout} $x_0 \in \Ob\, X$, \textit{le foncteur induit sur les catégories colocalisées en} $x_0$ \textit{et en} $y_0 = f(x_0)$
\[
x_0\backslash X \longrightarrow y_0\backslash Y
\]
\textit{est à fibres} $\D$\textit{-contractiles.}

\msp{page 69}

c) \textit{La formation des} $f_*F$ \textit{(}$F$ \textit{dans} $\D(X)$\textit{) commute à tout changement de base} $Y' \to Y$, \textit{et ça reste vrai après tout changement de base -- i.e. si on a} $Y'' \to Y' \to Y$, \textit{d'où}
\[
\begin{tikzcd}
X'' \arrow[r] \arrow[d, "f''"{swap}] & X' \arrow[r] \arrow[d, "f'"] & X \arrow[d, "f"] \\
Y'' \arrow[r, "p'"{swap}] & Y' \arrow[r, "p"{swap}] & Y
\end{tikzcd}
\]
\textit{le “théorème de changement de base” est vrai pour} $f'_*$ \textit{et} $p'^{*}$.

c$'$) \textit{La formation des} $f_*(F)$ \textit{(}$F \in \Ob\, \D(X)$\textit{) commute au passage aux fibres et ceci reste vrai après tout changement de base} $p : Y' \to Y$.

c$''$) \textit{Comme} c$'$) \textit{avec} $Y' \to Y$ \textit{un morphisme de localisation} $Y/y \to Y$.

d) \textit{Si} $Y'' \overset{p}{\to} Y' \to Y$, \textit{avec} $p$ \textit{coasphérique, alors} $q : X'' = X \times_Y Y'' \to X' = X \times_Y Y'$ \textit{est coasphérique.}

\printedpage{41}
% head: A. Grothendieck : Dérivateurs, Ch. I, Édition des Universités de Montpellier II et Paris VII
% foot: Version 8 juillet 2001

d$'$) \textit{Comme} d) \textit{avec} $Y'' = e$ (donc $Y'$ est catégorie avec objet final $e$).

d$''$) \textit{Comme} d) \textit{avec} $Y'' = \{1\}$, $Y' = \boldsymbol{\Delta}_1$.

.

\noindent\textsc{Démonstration}. a) $\Leftrightarrow$ b) cf. proposition page 53. Montrons que ces conditions impliquent c). Comme elles sont stables par changement de base (\textit{cf.} démonstration immédiate page 55) donc par \textbf{Der 2} on est ramené à prouver que $H^*_{\D}(X/y, F) \overset{\sim}{\to} H^*_{\D}(Y_y, F)$ $(F \in \Ob\, \D(Y))$, ce qui résulte du théorème 1.\footnote[30]{Donc on utilise \textbf{Der 5} (d).} (Cf. démonstration du théorème page 56, qui marche avec des hypothèses plus générales.)

On a c) $\Leftrightarrow$ c$'$) $\Rightarrow$ c$''$) (l'équivalence c) $\Leftrightarrow$ c$'$) par \textbf{Der 2}).

Prouvons c$''$) $\Rightarrow$ a).

\msp{page 70}

\noindent Considérons le diagramme à carrés cartésiens
\[
\begin{tikzcd}
X_{y_0} \arrow[r, "q"] \arrow[d, "p"{swap}] & X/y_0 \arrow[r] \arrow[d, "f_0"] & X \arrow[d, "f"] \\
\{y_0\} \arrow[r] & Y/y_0 \arrow[r] & Y
\end{tikzcd}
\qquad\qquad .
\]

\noindent L'hypothèse c$''$) prouve que pour tout $F$ dans $\D(X/y_0)$, on a $H^*_{\D}(X/y_0, F) \overset{\sim}{\to} H^*_{\D}(X_{y_0}, F)$. Par le théorème 1, ceci implique que $f$ est un $\D$-équivalence co-stricte. % [sic] un D-équivalence -> une D-équivalence, accord du genre

\dgrblock{dgrb40}{Completing the tests in the proper base-change theorem}
\begin{dcomment}{exposition}{p70:Par le théorème 1, ceci implique que}{Completing the tests in the proper base-change theorem}
The three conditions labelled (d) are not proved in the displayed argument. They have a short verification which also explains why a single walking arrow is a sufficient test. Write a pullback square as
\[
\begin{tikzcd}
 X''\arrow[r,"q"]\arrow[d,"f''"']&X'\arrow[d,"f'"]\\
 Y''\arrow[r,"p"']&Y'.
\end{tikzcd}
\]
If $f'$ is proper, the isomorphism $p^*f'_*\simeq f''_*q^*$ has, on passing to its left adjoint, the canonical isomorphism
\[
 q_!f''{}^*\simeq f'{}^*p_!.
\]
Here we take the left adjoints of the two \emph{composite functors}; no general assertion that every mate of every invertible cell is invertible is needed. If $p$ is coaspherical, its counit is an isomorphism on constants by Theorem~1. Applying the displayed comparison to a constant object on $Y''$ identifies the counit $q_!M_{X''}\to M_{X'}$ with its pullback. Theorem~1 now makes $q$ coaspherical. This proves (c)$\Rightarrow$(d).

Condition (d) implies its tests on inclusions of terminal objects, and these imply the test $\{1\}\to[1]$. Conversely, pull $f$ back along the functor $[1]\to Y$ determined by $\alpha:f(x)\to y$. Apply (d$''$) to the terminal-object inclusion $\{1\}\to[1]$. In the resulting coaspherical functor $X_y\to X\times_Y[1]$, the comma category under $(x,0)$ is exactly $L_f(x,\alpha)$. Thus (d$''$) implies (b). All the asserted conditions are now accounted for.

The parenthesis in (d$'$), ``donc $Y'$ est catégorie avec objet final $e$'', should be understood as specifying the terminal-object tests. A $\D$-coaspherical functor from the point need not literally select a terminal object for every possible coefficient theory. For example, take the represented derivator of the complete lattice $\{0<1\}$. Every nonempty discrete category is $\D$-contractible, since its diagonal constant-diagram functor is fully faithful. The inclusion $e\to BC_2$ is then $\D$-coaspherical: its only comma category is the discrete two-element set $C_2$. The unique object of $BC_2$ is not terminal. Reading (d$\prime$) as the restricted terminal-object test avoids attributing this false converse to the proof.

Several other slips here are already in the edition. The proof on page~69 requires $X_y$, not $Y_y$, and a coefficient on $X$, not on $Y$. At the end of the argument on page~70 the coaspherical functor is the displayed $q:X_{y_0}\to X/y_0$, not $f$. Definition~4 refers to Theorem~2, not Theorem~1. These corrections repair the typing and references; they do not alter the theorem.
\end{dcomment}


\noindent\textbf{Corollaire 1.} Énoncé dual du théorème, pour les foncteurs $X_{y_0} \to y_0\backslash X$, $X/x_0 \to Y/y_0$, et le changement de base des $j_!(F)$ \ldots

\noindent\textbf{Définition 4.} Sous les conditions équivalentes du théorème 1, on dit que $f$ est $\D$-\textit{propre.} La condition duale s'appelle $f$ $\D$-\textit{lisse.}

\noindent\textbf{Corollaire 2.} \textit{Soit} $f : X \to Y$ \textit{dans} $\Dia$. \textit{Si} $f$ \textit{est} $\D$\textit{-propre ou} $\D$\textit{-lisse, et à fibres} $\D$\textit{-contractiles, alors} $f$ \textit{est une} $\D$\textit{-équivalence stricte et costricte. Mieux,} $f$ \textit{est une} $\D$\textit{-équivalence} universelle \textit{i.e. est une} $\D$\textit{-équivalence et le reste après tout changement de base.}\footnote[31]{Se généralise : \textbf{Corollaire 2$'$.} \textit{Soit}
\[
\begin{tikzcd}[ampersand replacement=\&]
X \arrow[rr, "f"] \arrow[dr] \& \& Y \arrow[dl] \\
\& S \&
\end{tikzcd}
\]
\textit{avec} $X, Y$, $\D$\textit{-propres sur} $S$. \textit{Si} $f$ \textit{induit une} $\D$\textit{-équivalence sur les fibres c'est une} $\D$\textit{-équivalence. Itou dans le cas lisse.}\par
C'est la proposition 6 ci-dessous.}

\printedpage{42}
% head: A. Grothendieck : Dérivateurs, Ch. I, Édition des Universités de Montpellier II et Paris VII
% foot: Version 8 juillet 2001
\msp{page 71}

\noindent Il suffit de prouver que $f$ est une $\D$-équivalence. Supposons par exemple $f$ $\D$-propre. Il suffit de prouver que si $Y$ est $\D$-contractile, alors $X$ l'est (quitte à l'appliquer aux $Y/y_0$ après changement de base, pour montrer que les $X/y_0$ sont contractiles). Pour ceci, il suffit de prouver que $f$ est une $\D$-équivalence, et on sait que ceci est le cas si pour tout $F$ \textit{constant} sur $Y$, l'homomorphisme $F \overset{\text{adj}}{\longrightarrow} f_*f^*F$ est \texttt{[un]} isomorphisme. En fait, c'est vrai pour \textit{tout} $F$.

\noindent\textbf{Corollaire 3.} \textit{Si} $f : X \to Y$ \textit{est} $\D$\textit{-propre ou} $\D$\textit{-lisse, à fibres} $\D$\textit{-contractiles, alors} $f^* : \D(Y) \to \D(X)$ \textit{est pleinement fidèle, i.e. les homomorphismes d'adjonction}
\[
F \longrightarrow f_*f^*F \quad , \qquad f_!f^*F \longrightarrow F
\]
\textit{dans} $\D(Y)$ \textit{sont des isomorphismes.}

\noindent Il suffit de le voir pour $f$ propre, et cela résulte du calcul des fibres de $f_*f^*F$ \ldots

\noindent\textbf{Proposition 5.} \textit{La notion de foncteur} $\D$\textit{-propre} (resp. $\D$\textit{-lisse}) \textit{est stable par changement de base,} (trivial), \textit{et par composition} (utiliser \texttt{[le]} corollaire 2). \textit{Toute immersion fermée est} $\D$\textit{-propre, toute immersion ouverte} $\D$\textit{-lisse. Plus généralement, tout isomorphisme local est} $\D$\textit{-lisse, tout co-isomorphisme local est} $\D$\textit{-propre.}

\noindent\textbf{Proposition 6.} \textit{Soient} $X, Y$ \textit{sur} $S$, $f : X \to Y$ \textit{un} $S$\textit{-morphisme. Supposons} $X, Y$ \textit{propres sur} $S$, \textit{ou} $X$ \textit{et} $Y$ \textit{lisses sur} $S$, \textit{et que pour tout} $s \in \Ob\, S$, $f_s : X_s \to Y_s$ \textit{soit une} $\D$\textit{-équivalence, alors} $f$ \textit{aussi.}\footnote[32]{\textsc{Démonstration}. On prouvera même que sur la sous-catégorie des coefficients sur $Y$ qui proviennent de $\D(Y)$ \texttt{[plutôt de }$\D(S)$\texttt{]}, $f^* : \D(Y) \to \D(X)$ induit un foncteur pleinement fidèle. Si $f$ \texttt{[est]} propre, je l'interprète comme $q_*q^*F \overset{\sim}{\to} p_*p^*F$ pour tout $F$, et cela résulte du calcul de $p_*$, $q_*$ par passage aux fibres (propre), et de l'hypothèse sur les $f_s$. OK.}

\dgrblock{dgrb41}{Properness does not require constant fibres}
\begin{dcomment}{examples}{p71:et de l'hypothèse sur les}{Properness does not require constant fibres}
Let $A:[1]\to\Cat$ send $0$ to the point and $1$ to the discrete category $\{a,b\}$, with the transition selecting $a$. Its Grothendieck construction $X$ has three objects $u,a,b$, one nonidentity arrow $u\to a$, and projection $f:X\to[1]$ sending $u$ to $0$ and $a,b$ to $1$. The map $u\to a$ is a cocartesian lift of $0\to1$. The corresponding lift category is a point; the lift categories over identities have their identity lifts as initial objects. Thus $f$ is proper for every derivator under consideration.

For a contravariant coefficient diagram $F$ on $X$, the right direct image is
\[
 (f_*F)_0\simeq F_u,\qquad
 (f_*F)_1\simeq F_a\times F_b,
\]
and its transition is the composite $F_a\times F_b\to F_a\to F_u$. In the derivator of complexes, these are derived products, which for a finite pair agree with the ordinary finite product. For constant value $M$, the transition is the projection $M\times M\to M$. Taking $M=\mathbb Q$ in degree zero proves that it is not locally constant. The map is therefore proper but not a cohomological fibration in the definition about to be introduced.

This example separates two uses of a fibre calculation. Properness identifies a direct image with the cohomology of the fibres. A fibration condition additionally requires that these cohomologies be transported by isomorphisms along the base. There is no contradiction in a proper map having fibres with different homotopy types.

By contrast, if a proper map has $\D$-contractible fibres, then for every $F\in\D(Y)$ the component at $y$ of its unit is
\[
 F_y\longrightarrow H^*_{\D}(X_y,(F_y)_{X_y}),
\]
which is an isomorphism. Pointwise detection proves that $f^*$ is fully faithful. Since $f^*$ has both adjoints, this is equivalent both to invertibility of the unit of $f^*\dashv f_*$ and to invertibility of the counit of $f_!\dashv f^*$. This proves Corollary~3 and its smooth counterpart. It also gives universal $\D$-equivalence: properness and the fibres are preserved by base change. The claims of strict and costrict equivalence follow by applying that universal assertion to both kinds of local-isomorphism base changes.
\end{dcomment}


\msp{page 72}

\noindent\textbf{Théorème 3.} \textit{Soit}
\[
\begin{tikzcd}
X \arrow[d, "f"{swap}] & X' \arrow[l, "q"{swap}] \arrow[d, "f'"] \\
Y & Y' \arrow[l, "p"{swap}]
\end{tikzcd}
\]
\textit{cartésien dans} $\Dia$. \textit{Alors le morphisme de changement de base}
\[
p^*f_* \longrightarrow f'_*q^*
\]
\textit{est un isomorphisme si} $f$ \textit{est} $\D$\textit{-propre, ou} $p$ $\D$\textit{-lisse.}

\noindent Le cas $\D$-propre déjà prouvé. Le cas $\D$-lisse : la démonstration p. 56 est OK.

\dgrblock{dgrb42}{The geometric pattern and its categorical content}
\begin{dcomment}{context}{p72:la démonstration p. 56 est OK.}{The geometric pattern and its categorical content}
The vocabulary ``proper'' and ``smooth'' records a specific geometric analogy. In étale cohomology, proper direct image commutes with base change, and a smooth change of base gives a second base-change theorem under the appropriate torsion hypotheses. The comparison map has the same shape as the one displayed here. \citet[Introduction and Theorem~1.2]{Maltsiniotis2005Asphericite} explicitly traces the categorical definitions to this pattern in \emph{Pursuing Stacks} and identifies the corresponding SGA~4 exposés. His Theorems~4.13 and~5.4 develop the relation between lifting categories and base change in the framework of asphericity structures.

There is no implication here that a proper functor between categories is a proper morphism of spaces in the topological sense. The definition asks for a Beck--Chevalley isomorphism in a coefficient theory. The lifting categories supply its combinatorial content. The smooth half of Theorem~3 is obtained by applying that same calculation with the two adjunctions interchanged; it is not obtained merely by observing that the underlying square is a pullback. In the covariant-diagram convention, \citet[Propositions~1.24 and~1.26]{Groth2013} formulate the corresponding exactness of pullback squares for Grothendieck fibrations and opfibrations and of comma squares. Opposing all four indexing categories returns the convention of this manuscript.

For Proposition~6 one needs only to compare $q_*q^*M$ and $p_*p^*M$ over the common base $S$. Their components are constant-coefficient cohomologies of $Y_s$ and $X_s$. The hypothesis identifies those components, and \textbf{Der 2} finishes the proof. The footnote's stronger statement for coefficients pulled back from $S$ follows by the same calculation with $M$ replaced pointwise by $F_s$, followed by adjunction. This argument is valid without requiring the comparison map $X\to Y$ itself to be proper.
\end{dcomment}


\printedpage{43}
% head: A. Grothendieck : Dérivateurs, Ch. I, Édition des Universités de Montpellier II et Paris VII
% foot: Version 8 juillet 2001

\subsection*{Comparaison de $X$ et $X^{\circ}$.}

Soit $X$ \texttt{[une]} catégorie, on définit une catégorie $\Phi(X)$ et des foncteurs
\[
\begin{tikzcd}[column sep=small]
& & \Phi(X) \arrow[ddl, "p_X"{swap, name=P}] \arrow[ddr, "q_X"{name=Q}] & & \\
|[alias=L]| \text{\footnotesize\begin{tabular}{c}$\Cat$-cofibrant\\(donc propre)\end{tabular}} & & & & |[alias=R]| \text{\footnotesize\begin{tabular}{c}$\Cat$-fibrant\\(donc lisse)\\ \texttt{[plutôt }$\Cat$\texttt{-cofibrant}\\ \texttt{(donc propre)]}\end{tabular}} \\
& X & & X^{\circ} &
\arrow[from=P, to=L]
\arrow[from=Q, to=R]
\end{tikzcd}
\]
qui seront des $\hot$-équivalences, et des $\D$-équivalences si $\D$ est un dérivateur dont le domaine contient $X$ et $\Phi(X)$. On pose
\[
\Ob\, \Phi(X) = \Fl(X) \quad ,
\]
donc mêmes objets que $\mathcal{F}\ell(X) = \underline{\Hom}(\boldsymbol{\Delta}_1, X)$, mais les morphismes sont différents :

\msp{page 73}
\[
\begin{tikzcd}
Y_1 \arrow[r, "u"] & Y_0 \arrow[d, "v_0"] \\
Y'_1 \arrow[u, "v_1"] \arrow[r, "u'"{swap}] & Y'_0
\end{tikzcd}
\]
un morphisme de $u$ dans $u'$ est un couple $v_0 : Y_0 \to Y'_0$, $v_1 : Y'_1 \to Y_1$ tel que $v_0uv_1 = u'$. Ainsi source et but $s(u)$ et $b(u)$ sont l'une contravariante, l'autre covariant en $u$. On prend pour $q$ le foncteur source, pour $p$ le foncteur but. % [sic] l'autre covariant -> l'autre covariante, accord avec « l'une contravariante »

Si on applique ceci à $X^{\circ}$, on trouve
\[
\begin{tikzcd}
& \Phi(X^{\circ}) \arrow[dl, "p_{X^{\circ}}"{swap}] \arrow[dr, "q_{X^{\circ}}"] & \\
X^{\circ} & & (X^{\circ})^{\circ} = X
\end{tikzcd}
\]
On a $\Ob\, \Phi(X^{\circ}) = \Fl(X^{\circ}) \simeq \Fl(X)$
\[
\begin{tikzcd}
Y_1 \arrow[r, "u"] & Y_0 \arrow[d, "v_1"] & \text{dans}\ \ X \\
Y'_1 \arrow[u, "v_0"] \arrow[r, "u'"{swap}] & Y'_0 & \text{dans}\ \ X
\end{tikzcd}
\]
et mêmes homomorphismes au sens de $\Phi(X)$ et au sens de $\Phi(X^{\circ})$. Donc
\[
\left\{
\begin{aligned}
&\Phi(X^{\circ}) = \Phi(X) \\
&p_{X^{\circ}} = q_X \\
&q_{X^{\circ}} = p_X \quad ,
\end{aligned}
\right.
\]

% ---- from der1_p046-050.tex
\printedpage{44}
% head: A. Grothendieck : Dérivateurs, Ch. I, Édition des Universités de Montpellier II et Paris VII
le passage de $X$ à $X^\circ$ revient à échanger $p_X$ et $q_X$. La fibre de $\Phi(X)$ en $x \in \Ob X$ est $\simeq (X/x)^\circ$, celle en $y^\circ \in \Ob X^\circ$ est $\simeq y\bs X$
\[
\left\{
\begin{aligned}
&p_X^{-1}(\{x\}) \simeq (X/x)^\circ\\
&q_X^{-1}(\{y^\circ\}) \simeq y\bs X
\end{aligned}
\right.
\]

{\raggedright
N.B. : $p_X$ est $W$-lisse \texttt{[plutôt $W$-propre]} à fibres $W$-asphériques, donc dans $W^u$ \texttt{[$W^u =$ ensemble des morphismes qui sont dans $W$ et le restent après tout changement de base]}, $q_X$ est $W$-propre à fibres $W$-asphériques, donc dans $W^u$.\par}

\dgrblock{dgrb43}{The twisted-arrow category and its two projections}
\begin{dcomment}{exposition}{p73:et le restent après tout changement de base}{The twisted-arrow category and its two projections}
The category $\Phi(X)$ is the twisted-arrow category, with the orientation explicitly fixed by the equation $u'=v_0uv_1$. Keeping that equation in view avoids confusing it with the ordinary arrow category. An object $u:a\to b$ maps to $b$ under $p_X$ and to $a$, considered in $X^{\mathrm{op}}$, under $q_X$. Composition is componentwise in the appropriate orders. Replacing $u$ by $u^{\mathrm{op}}$ interchanges the two components of a twisted morphism and gives a canonical isomorphism
\[
 \Phi(X^{\mathrm{op}})\cong\Phi(X),
\]
not the opposite category on the right. The two projection formulas follow under this identification.

Both projections are opfibrations. For $p_X$, an arrow $\beta:b\to b'$ is lifted from $u$ by the twisted arrow $u\to\beta u$, with source component $1_a$ and target component $\beta$. If a twisted arrow from $u$ to $w:c\to d$ has target component $\gamma\beta$, it factors uniquely through this chosen lift, using target component $\gamma$ and the same source component. This is the cocartesian universal property. For $q_X$, an arrow $a\to a'$ in $X^{\mathrm{op}}$ is an arrow $\alpha:a'\to a$ in $X$; its cocartesian lift is $u\to u\alpha$. The same factorisation argument applies.

The fibre of $p_X$ over $b$ is $(X/b)^{\mathrm{op}}$: its arrows are the reversed triangles between maps to $b$. Its object $1_b$ is initial. The fibre of $q_X$ over $a$ is $a\backslash X$, whose object $1_a$ is initial as well. Consequently both maps are proper with contractible fibres. The preceding corollaries prove directly that they are universal $\D$-equivalences and that inverse image along either is fully faithful. This gives a shorter proof of the assertion announced on page~74; it uses exactly the corrected ``cofibrant, donc propre'' label in the edition.

The displayed slice identities are also literal categorical isomorphisms, and not merely equivalences of homotopy types. An object of $\Phi(X)/x$ is a composable pair $a\xrightarrow{u}b\xrightarrow{h}x$. It determines the arrow $(a,hu)\to(b,h)$ in $X/x$, hence an object of $\Phi(X/x)$. A twisted morphism consists of $v_1:a'\to a$, $v_0:b\to b'$ with $u'=v_0uv_1$ and $h'v_0=h$. The remaining slice equation $h'u'=huv_1$ then follows. These are precisely the morphisms in $\Phi(X/x)$. Reversing the construction is immediate. The second identity follows by applying this argument to $X^{\mathrm{op}}$ and interchanging the endpoints.

The last comma notation in the succeeding proof requires a correction. The strict-equivalence test for $q_Z:\Phi(Z)\to Z^{\mathrm{op}}$ uses $\Phi(Z)/z_0^{\mathrm{op}}$, whose additional arrow in $Z$ is $z_0\to s(u)$. It is this category that is isomorphic to $\Phi(z_0\backslash Z)$. The edition instead prints $z_0\backslash\Phi(Z)$; interpreted literally as a comma category under $q_Z$, that symbol would require $s(u)\to z_0$ in $Z$ and reverse the intended arrow. With the intended comma category restored, the proof continues correctly. If $Z$ has a terminal object, $z_0\backslash Z$ has both an initial and a terminal object. The unique arrow between these is terminal in its twisted-arrow category, exactly as Grothendieck observes.
\end{dcomment}


\msp{page 74}

On vérifie
\[
\left\{
\begin{aligned}
&\Phi(X)/x \simeq \Phi(X/x) &\qquad& x \in \Ob X\\
&\Phi(X)/x^\circ \simeq \Phi(x\bs X) &\qquad& x^\circ \in \Ob X^\circ
\end{aligned}
\right.
\]

Je dis que $p_X$, $q_X$ sont des $\D$-\textit{équivalences strictes} et \textit{costrictes.} Par les formules précédentes, cela signifie que les $\Phi(X/x)$, $\Phi(x\bs X)$ sont $\D$-contractiles. Donc que si $Z$ a un objet final, $\Phi(Z)$ est $\D$-contractile. Mais considérons
\[
\begin{tikzcd}
\Phi(Z) \arrow[dr, "q_Z"] & \\
& Z^\circ
\end{tikzcd}
\quad ;
\]
il suffit de prouver que c'est une $\D$-équivalence stricte (puisque $Z^\circ$, ayant un objet initial, est $\D$-contractil), donc que les $z_0\bs\Phi(Z)$ sont $\D$-contractiles. On a par les formules précédentes % [sic] contractil -> contractile, adjectif ailleurs écrit « contractile »
\[
z_0\bs\Phi(Z) \simeq \Phi(z_0\bs Z) = \Phi(T) \quad ,
\]
où cette fois $T$ a un objet initial $\emptyset$ et un objet final $e$. Considérons alors l'objet
\[
\emptyset \longrightarrow e
\]
de $\Phi(T)$, c'est un objet final, car pour tout $x \to y$ dans $\Phi(T)$, il y a une unique flèche
\[
\begin{tikzcd}
x \arrow[r] & y \arrow[d] \\
\emptyset \arrow[u] \arrow[r] & e
\end{tikzcd}
\]
% foot: Version 8 juillet 2001

\dgrblock{dgrb44}{Intervals in a finite ordered set}
\begin{dcomment}{examples}{p74:il y a une unique flèche}{Intervals in a finite ordered set}
For $X=[1]$, the twisted-arrow category has objects $1_0$, $u:0\to1$, and $1_1$, with precisely the nonidentity arrows
\[
 1_0\longrightarrow u\longleftarrow1_1.
\]
Thus $u$ is terminal. Under the target projection, the fibre at $0$ is the point $1_0$, whereas the fibre at $1$ is the arrow $1_1\to u$. Under the source projection, the roles are reversed. Each fibre is contractible, although the projections are not isomorphisms of categories. Their nerves subdivide the interval into two edges.

For the chain $[n]$, objects are intervals $(i,j)$ with $0\leq i\leq j\leq n$. There is exactly one morphism
\[
 (i,j)\longrightarrow(i',j')
 \quad\Longleftrightarrow\quad i'\leq i\leq j\leq j'.
\]
The largest interval $(0,n)$ is terminal. More generally this description works for comparable pairs in any poset, showing immediately that $\Phi$ preserves finite posets. This is the domain-closure observation at the end of manuscript page~75; it is not an assertion that every finite category is a finite poset.

The comparison with the opposite category is useful because it is natural. For a functor $f:X\to Y$, applying $f$ to arrows and to both components gives $\Phi(f)$, and both endpoint squares commute. Since their vertical projections are $\D$-equivalences, two-out-of-three first gives
\[
 f\in W_{\D}\ \Longleftrightarrow\ \Phi(f)\in W_{\D},
\]
and then the corresponding equivalence with $f^{\mathrm{op}}\in W_{\D}$. One must have $X^{\mathrm{op}},Y^{\mathrm{op}},\Phi(X),\Phi(Y)$ in the domain. There is no claim that $X$ and $X^{\mathrm{op}}$ are equivalent as ordinary categories. The natural span compares precisely the information detected by the chosen coefficient theory.
\end{dcomment}


\printedpage{45}
% head: A. Grothendieck : Dérivateurs, Ch. I, Édition des Universités de Montpellier II et Paris VII
\msp{page 75}

Pour résumer

\textbf{Proposition 1 :} \textit{$p_X$ et $q_X$ sont des} \textsf{hot}\textit{-équivalences strictes,} i.e. \textsf{Hot}\textit{-asphériques (donc aussi des $\D$-équivalences strictes si $\Phi(X)$ est dans le domaine de $\D$) à fibres les $(X/x)^\circ$ et $y\bs X$ respectivement (donc à fibres contractiles).}

\textbf{Corollaire 1.} \textit{Pour que $X$ soit $\D$-contractile, il faut et il suffit que $X^\circ$ le soit.}

\hspace*{\parindent}Plus généralement :

\textbf{Corollaire 2.} \textit{Soit $f : X \to Y$. Pour que $f$ soit une $\D$-équivalence, il faut et il suffit que $f^\circ$ le soit.}

\hspace*{\parindent}On utilisera que le diagramme
\[
\begin{tikzcd}
& \Phi(X) \arrow[dl, "p_X"'] \arrow[dr, "q_X"] & \\
X & & X^\circ
\end{tikzcd}
\]
est fonctoriel en $X$. On a donc
\[
\begin{tikzcd}
& \Phi(X) \arrow[rr, "\Phi(f)"] \arrow[dl, "p_X"'] & & \Phi(Y) \arrow[dl, "p_Y"] \arrow[dd, "q_Y"] \\
X \arrow[rr, "f" near end] \arrow[from=1-2, to=3-2, "q_X"' near end, crossing over] & & Y & \\
& X^\circ \arrow[rr, "f^\circ"'] & & Y^\circ
\end{tikzcd}
\]
les deux carrés commutatifs. Comme $p_X$ et $p_Y$ sont des $\D$-équivalences, $f$ est une $\D$-équivalence si et seulement si $\Phi(f)$ l'est, et de même pour $f^\circ$. Ces raisonnements supposent que si $X$ est dans le domaine de $\D$, $\Phi(X)$ aussi. C'est le cas si \Dia{} est formé des catégories ordonnées finies.
% foot: Version 8 juillet 2001

\dgrblock{dgrb45}{Twisted arrows after the opposite-category comparison}
\begin{dcomment}{later}{p75:Ces raisonnements supposent que si}{Twisted arrows after the opposite-category comparison}
The construction has become a standard indexing device for homotopy ends and coends. \citet[Example~7.2.9 and the footnote in Section~7.8]{Riehl2014} describe the twisted-arrow category and the identification of its nerve with edgewise subdivision. This supplies a simplicial interpretation of the span $X\leftarrow\Phi(X)\to X^{\mathrm{op}}$. Its realization compares the two classifying spaces without choosing an equivalence of categories.

For a derivator written with the covariant convention, \citet[Section~5]{PontoShulman2016} use the same category to define the homotopy coend of $A\in\mathbb V(X^{\mathrm{op}}\times X)$ by
\[
 \int^{X}A=(\pi_{\Phi(X)^{\mathrm{op}}})_!
       (t^{\mathrm{op}},s^{\mathrm{op}})^*A.
\]
The opposite on the indexing category is essential in this formula. Ends use the other variance. The passage from this formula to composition of homotopy profunctors explains why the two endpoint projections occur together so persistently in later derivator calculus. Grothendieck uses them here for a narrower purpose: to show that the weak-equivalence class obtained from constant coefficients is unchanged by taking opposites. The later coend formalism extends the utility of the construction; it is not an extra assumption in his proof.
\end{dcomment}


\printedpage{46}
% head: A. Grothendieck : Dérivateurs, Ch. I, Édition des Universités de Montpellier II et Paris VII
\msp{page 76}

\subsection*{$\D$-fibrations}

\textbf{Proposition 1.} \textit{Soit $f : X \to Y$ dans} \Dia\textit{, conditions équivalentes :}

\hspace*{\parindent}a) \textit{Pour tout $y \in Y$, l'inclusion $X_y \to X/y$ est une $\D$-équivalence, et pour toute flèche $y_1 \to y_0$ de $Y$, $X/y_1 \to X/y_0$ est une $\D$-équivalence.}

\hspace*{\parindent}b) \textit{Pour tout coefficient} constant \textit{$F$ sur $Y$, $f_*f^*F$ est localement constant,}\footnote[33]{Par là on entend que pour tout $u : y \to y'$ dans $Y$, $G_{y'} \to G_y$ est \texttt{[un]} isomorphisme. Sans axiomes supplémentaires sur $\D$, il n'est pas clair que cela implique que les ``restrictions'' de $G$ aux $Y/y$ ou $y\bs Y$, soient isomorphes à des coefficients constants. Voir cependant l'axiome \textbf{Der 3} a$'$), page 97, qui assure que c'est pareil.} \textit{et ``se calcule fibre par fibre''} i.e. \textit{l'homomorphisme canonique}
\[
f_*f^*(F)_y \overset{\sim}{\longrightarrow} H^*_{\D}(X_y, F|X_y)
\]
\texttt{[est un isomorphisme]}.

Ces conditions sont équivalentes, et impliquent les conditions de (b) pour tout $F$ \textit{localement} constant sur $Y$.

La première condition dans (a) signifie que pour tout $y \in Y$, et tout $M \in \Ob \mathcal{A}$,
\[
H^*_{\D}(X/y, M) \longrightarrow H^*_{\D}(X_y, M)
\]
est un isomorphisme. Or cette flèche s'interprète comme
\[
f_*(M_X)_y \longrightarrow H^*_{\D}(X_y, M)
\]
est un isomorphisme, \textit{i.e.} que $f_*(M_X)$ se calcule fibre par fibre. Comme $M_X \simeq f^*(M_Y)$, c'est la deuxième condition (b). La deuxième condition dans (a), pour $y_1 \to y_0$ signifie que pour $M$ dans $\mathcal{A}$,
\[
H^*_{\D}(X/y_0, M) \longrightarrow H^*_{\D}(X/y_1, M)
\]
est \texttt{[un]} isomorphisme, or cette flèche s'interprète comme
\[
f_*(M_X)_{y_0} \longrightarrow f_*(M_X)_{y_1} \quad ,
\]
que ce soit toujours un isomorphisme, c'est la première condition (b). Que (b) soit valable aussi pour $F$ localement constant sur $Y$ est immédiat, car c'est une condition de nature locale par rapport à $Y$.

\dgrblock{dgrb46}{Local constancy and the first fibration condition}
\begin{dcomment}{exposition}{p76:que ce soit toujours un isomorphisme,}{Local constancy and the first fibration condition}
The new fibration condition retains two separate tests. First, the canonical inclusion of an actual fibre into its comma category must be a $\D$-equivalence. Second, the transition functors between comma categories must be $\D$-equivalences. The first identifies the values of $f_*M_X$ with fibre cohomology. The second says that their transport maps are isomorphisms. Since $\D$-equivalence tests all absolute coefficients $M$, this is an equivalence of conditions, not an inference from the cohomology of a single distinguished coefficient.

The extension from constant to locally constant coefficients is valid with the usual two-dimensional functoriality and pointwise detection axioms. There is a direct proof of the local issue raised in footnote~33. If a category $A$ has terminal object $a$, the natural transformation $1_A\to ap$ induces, in the manuscript's presheaf convention, a morphism
\[
 p^*a^*F\longrightarrow F.
\]
For arrowwise locally constant $F$, every component is an isomorphism; \textbf{Der 2} makes the whole morphism invertible. If $A$ has an initial object, the corresponding natural transformation gives the comparison in the other direction, with the same conclusion. Applying these observations to $Y/y$ and $y\backslash Y$ proves that restrictions of a locally constant object are constant there. This only uses categories possessing the indicated distinguished object. The stronger assertion for \emph{every} weakly contractible category is the separate issue addressed by the later variants of \textbf{Der 3}.

Now fix an arrow $y_1\to y_0$. Pull a locally constant $F$ back to $Y/y_0$, where it is isomorphic to the constant object with value $F_{y_0}$. Pulling this isomorphism further to $X/y_0$, $X/y_1$, and the two fibres turns the desired comparison into the one already established for constant coefficients. The isomorphism $F_{y_0}\to F_{y_1}$ supplies the change of value at the other endpoint. This proves the statement for locally constant $F$, including the compatibility of transport, rather than relying only on the phrase ``de nature locale''.

For model-category derivators, a systematic treatment is \citet[Theorem~1.3 and Section~3]{Cisinski2009LocallyConstant}: weak equivalences of indexing nerves induce equivalences on the homotopy categories of locally constant diagrams. The distinction between constant and locally constant objects remains necessary when the base has nontrivial fundamental group. It is exactly this distinction that will invalidate the proposed implication on manuscript pages~80--81.
\end{dcomment}


\msp{page 77}

La proposition duale fait appel aux foncteurs
\[
y_0\bs X \longrightarrow y_1\bs X \qquad \text{et} \qquad X_y \longrightarrow y\bs X \quad ,
\]
et aux faisceaux $f_!f^*M$, pour $M$ constant, ou localement constant sur $Y$.
% foot: Version 8 juillet 2001

\printedpage{47}
% head: A. Grothendieck : Dérivateurs, Ch. I, Édition des Universités de Montpellier II et Paris VII
\textbf{Définition 1.} Si les conditions équivalentes de la proposition 1 sont satisfaites, on dit que $f$ est une $\D$-\textit{fibration cohomologique} (ou $\D$-\textit{fibration gauche}). Si les conditions duales sont satisfaites, on dit que c'est une $\D$-\textit{fibration homologique} (ou $\D$-\textit{fibration droite}).

N.B. Les $\D$-fibrations cohomologiques sont stables par tout changement de base qui est un isomorphisme local. Les $\D$-fibrations homologiques par tout changement de base qui est un co-isomorphisme local. Mais ce ne sont pas des notions stables par \textit{tout} changement de base\footnote[34]{? donner exemple!}. Mais la notion de $\D$-fibration cohomologique est stable par changement de base $\D$-lisse, et celle de $\D$-fibration homologique par changement de base $\D$-propre.

\dgrblock{dgrb47}{A two-object failure of arbitrary base change}
\begin{dcomment}{development}{p77:Mais ce ne sont pas des notions stables par}{A two-object failure of arbitrary base change}
The footnote asks for an example showing that the first fibration condition need not survive arbitrary base change. We can supply one before the more elaborate counterexamples later in the manuscript. We checked the treatments of smoothness and asphericity in \citet[Sections~1 and~4]{Maltsiniotis2005Asphericite} and of locally constant diagrams in \citet{Cisinski2009LocallyConstant}; we did not find this precise small example there. The construction below is an elementary verification for this commentary, without a claim of priority.

\textbf{A non-universal cohomological fibration.}\ours
Let $J$ be the groupoid with two objects $a,b$ and exactly one arrow between every ordered pair of objects. Let $f:[1]\to J$ send $0$ to $a$, $1$ to $b$, and the nonidentity arrow to the unique isomorphism $a\to b$. For the space derivator, $f$ is both a cohomological and a homological fibration in Definition~1, but one strict base change of $f$ has neither property.

\emph{Proof.} Each strict fibre of $f$ is a point. For either $y\in J$, an object of $[1]/y$ consists of $i\in[1]$ and an arrow $f(i)\to y$. The latter is unique. Projection therefore gives an isomorphism $[1]/y\cong[1]$. Under these isomorphisms every transition functor is the identity. The inclusion of either fibre is an endpoint inclusion $e\to[1]$, hence a weak equivalence of nerves. This proves the cohomological condition. The same argument with $y\backslash[1]$ proves the homological condition. More generally these tests succeed for any coefficient theory in which the point and the walking arrow have the usual constant-coefficient invariance.

Define $g:[1]\to J$ in the reverse order: $g(0)=b$, $g(1)=a$. Form the strict pullback
\[
 P=[1]\times_J[1]\xrightarrow{\,f'\,}[1],
\]
with $f$ on the first factor and $g$ on the second. Its objects are exactly $(1,0)$ and $(0,1)$. There is no arrow from the first to the second, because there is no arrow $1\to0$ in the first factor. There is none in the reverse direction, because there is no arrow $1\to0$ in the second factor. Thus $P$ is the discrete two-object category, with one object in each fibre of $f'$.

The comma category $P/1$ is discrete with two objects, whereas $P_1$ is a point. Hence $P_1\to P/1$ is not a weak equivalence, and $f'$ is not a cohomological fibration. Dually $0\backslash P$ is discrete with two objects while $P_0$ is a point, so the homological condition fails too. For complexes of rational vector spaces the same failure is visible in degree zero: the fibre comparison is $\mathbb Q\times\mathbb Q\to\mathbb Q$, a projection with nonzero kernel.\hfill$\square$

Both $[1]$ and $J$ in the original pair of maps are contractible, and every original strict fibre is a point. The defect is the incompatible choice of direction for a strict lift of the base isomorphism. This example therefore isolates the need for a universal condition without invoking monodromy or a failure of constant coefficients to detect it. The later counterexamples address a different, stronger assertion: whether a universal fibration preserves all $\D$-equivalences under pullback.
\end{dcomment}


\hspace*{\parindent}Je vais essayer de construire un exemple où le changement de base canule, en posant $Y = \boldsymbol{\Delta}_1$. Se donner $X$ sur $\boldsymbol{\Delta}_1$ revient à se donner une sous-catégorie ouverte $X_0$, et son complémentaire fermé $X_1$
\[
\begin{tikzcd}
& X & \\
X_1 \arrow[ur, hook] \arrow[d, dash, dashed] & & X_0 \arrow[ul, hook'] \arrow[d, dash, dashed] \\
1 & & 0 \arrow[ll]
\end{tikzcd}
\quad .
\]
Dire que $X \to Y = \boldsymbol{\Delta}_1$ est une $\D$-fibration cohomologique signifie simplement
\msp{page 78}
que les restrictions
\[
X_0 \overset{j}{\longrightarrow} X \quad , \qquad X_1 \overset{i}{\longrightarrow} X
\]
sont des $\D$-équivalences. Prenons un changement de base $Y' \to Y = \boldsymbol{\Delta}_1$, \textit{i.e.} un $Y'$ avec $Y'_0, Y'_1$
\[
\begin{tikzcd}
& Y' & \\
Y'_1 \arrow[ur, hook] \arrow[d, dash, dashed] & & Y'_0 \arrow[ul, hook'] \arrow[d, dash, dashed] \\
1 & & 0 \arrow[ll]
\end{tikzcd}
\quad .
\]
On voit que la condition $Y'$ lisse sur $\boldsymbol{\Delta}_1$ signifie que pour tout $y'_1 \in \Ob Y'_1$, la catégorie $Y'_0/y'_1$ des objets $y_0$ de $Y'_0$ au dessus de $y'_1$ (\textit{i.e.} munis de $y_0 \to y'_1$) est contractile. On voit que si $Y'_1 \neq \emptyset$ alors $Y'_0 \neq \emptyset$ -- mais un exemple avec $Y'_0 = \emptyset$, où $Y'_1 = \emptyset$, n'est pas intéressant
% foot: Version 8 juillet 2001
\printedpage{48}
% head: A. Grothendieck : Dérivateurs, Ch. I, Édition des Universités de Montpellier II et Paris VII
(car une projection $I \times J \to I$ est toujours une $\D$-fibration cohomologique). Prenons par exemple pour $Y'$
\[
\begin{tikzcd}
& a \arrow[dl] \\
b & \\
& a' \arrow[ul] \\
\underbrace{\qquad}_{\displaystyle Y'_1} & \underbrace{\qquad}_{\displaystyle Y'_0}
\end{tikzcd}
\quad ,
\]
somme amalgamée de deux copies de $Y = \boldsymbol{\Delta}_1$, donc $X' = X \times_Y Y'$ peut se voir comme somme amalgamée $X \amalg_{X_1} X$. Mais ça semble \textit{pas} être un contre-exemple. Il faudrait que $X_1 \to X'$ ne soit \textit{pas} une $\D$-équivalence.

\dgrblock{dgrb48}{Why the first doubled-arrow attempt does not work for spaces}
\begin{dcomment}{examples}{p78:somme amalgamée de deux copies de}{Why the first doubled-arrow attempt does not work for spaces}
The fork $a\to b\leftarrow a'$ is a useful test even though it does not provide the desired counterexample for the space derivator. A functor $X\to[1]$ decomposes $X$ into a sieve $X_0$ and a cosieve $X_1$. Pullback along that fork is two copies of $X$ identified along $X_1$:
\[
 X'=X\amalg_{X_1}X.
\]
Every nonidentity chain in $X'$ lies in one of the copies. Indeed, after it enters the common cosieve $X_1$, it cannot leave it to reach the other copy of $X_0$. Consequently the same pushout description holds degree by degree for nerves:
\[
 NX'=NX\amalg_{NX_1}NX.
\]
The maps $NX_1\to NX$ are inclusions of simplicial subsets. Their realizations are inclusions of CW subcomplexes. If $X\to[1]$ is a cohomological fibration for spaces, both $X_0\to X$ and $X_1\to X$ are weak homotopy equivalences. In particular the latter inclusion is a homotopy equivalence of CW complexes and a cofibration. Gluing a second copy along it preserves its homotopy type: the homotopy extension property, applied to an inverse and its homotopy, makes the inclusion of either copy in the pushout a homotopy equivalence. It follows that $|NX_1|\to|NX'|$ is a weak equivalence.

At $a$ and $a'$ the new fibre and comma category are copies of $X_0$. At $b$ they are $X_1$ and $X'$; the preceding paragraph supplies their comparison. The transport from either of the first two comma categories to the last is a copy of $X_0\to X\to X'$, hence a weak equivalence. Thus this particular base change is still a cohomological fibration. This establishes Grothendieck's expressed suspicion in the space case; it does not assert that arbitrary base changes work.
\end{dcomment}


\msp{page 79}

\textbf{Proposition 2.} \textit{Soit $f : X \to Y$ dans} \Dia. \textit{Conditions équivalentes :}

\hspace*{\parindent}a) \textit{Pour tout $Y'' \overset{p}{\to} Y' \to Y$, d'où}
\[
\begin{tikzcd}
X'' \arrow[r, "q"] \arrow[d] \arrow[dr, phantom, "\text{cart}"] & X' \arrow[r] \arrow[d] \arrow[dr, phantom, "\text{cart}"] & X \arrow[d, "f"] \\
Y'' \arrow[r, "p"'] & Y' \arrow[r] & Y
\end{tikzcd}
\quad ,
\]
\textit{si $p$ est une $\D$-équivalence,} \texttt{[alors]} \textit{$q$ aussi.}\footnote[35]{Il faut supposer $Y', Y''$ \textsf{Hot}-asphériques, cf. page 96, et supposer vérifié l'axiome Der 3 $a'$) de la page 97.}

\hspace*{\parindent}b) \textit{$f$ est une $\D$-fibration cohomologique ``universelle'',} i.e. \textit{est une fibration cohomologique, et le reste après tout changement de base.}

\hspace*{\parindent}b$'$) \textit{Condition duale de} (b) \textit{: $f$ est une $\D$-fibration homologique universelle.}\footnote[36]{b) $\Rightarrow$ a) \underline{canulée}, \textit{cf.} annotation marginale,\footnotemark[37] \textit{cf.} contre-exemple page 89. Mais ça marche pour les dérivateurs $\D_i$ ($i \in \mathbb{N}$) associés aux ``localisateurs fondamentaux'' $W_i$, \textit{i.e.} à la théorie des $i$-types d'homotopie.}\footnotetext[37]{L'implication b) $\Rightarrow$ a) pas prouvée, seulement sous forme plus faible :\par \hspace*{\parindent}a$'$) si $Y''$, $Y'$ sont $\D$-asphériques et 1-connexes, alors $X'' \to X'$ est une $\D$-équivalence. (Il suffit même de le supposer pour $Y'$, $Y''$ ayant des objets finaux (pour impliquer b) ou ayant des objets initiaux (pour impliquer b$'$).)}

\textit{a}) $\Rightarrow$ \textit{b}) : Supposons \textit{a}), soit $Y' \to Y$ changement de base, d'où
\[
\begin{tikzcd}
X' \arrow[r] \arrow[d, "f'"'] & X \arrow[d, "f"] \\
Y' \arrow[r] & Y
\end{tikzcd}
\quad .
\]
% foot: Version 8 juillet 2001

% ---- from Der-ChI-234_p051-055.tex
\printedpage{49}
% head: A. Grothendieck : Dérivateurs, Ch. I, Édition des Universités de Montpellier II et Paris VII
% foot: Version 8 juillet 2001

Pour prouver que $f'$ est une $\DD$-fibration cohomologique, on regarde au dessus de $Y'$ les foncteurs de la forme $\{y'\} \to Y'/y'$ ($y' \in \Ob\, Y'$) et $Y'/y'_1 \to Y'/y'_0$ (si $u' : y'_1 \to y'_0$ flèche de $Y'$), et on voit que les flèches dans $\Cat/Y'$ déduites de ceux-ci par le changement de base $X' \xrightarrow{f'} Y'$ restent des $\DD$-équivalences

\msp{page 80}

\[
\begin{tikzcd}[column sep=large, row sep=large]
X'_1 \arrow[r, "q"] \arrow[d, "f'_1"'] \arrow[dr, phantom, "\text{cart}"] & X'_0 \arrow[r] \arrow[d, "f'_0"] \arrow[dr, phantom, "\text{cart}"] & X' \arrow[r] \arrow[d, "f'"] \arrow[dr, phantom, "\text{cart}"] & X \arrow[d, "f"] \\
Y'_1 \arrow[r, "p"'] & Y'_0 \arrow[r] & Y' \arrow[r] & Y & .
\end{tikzcd}
\]
Comme $p$ est une $\DD$-équivalence, cela résulte bien de a).

\indent\textit{b}) $\Rightarrow$ \textit{a}) : Comme la condition \textit{b}) est évidemment stable par changement de base, on est ramené à prouver que si $f$ satisfait \textit{b}), et si on a
\[
\begin{tikzcd}[column sep=large, row sep=large]
X \arrow[d, "f"'] \arrow[dr, phantom, "\text{cart}"] & X' \arrow[l, "q"'] \arrow[d, "f'"] \\
Y & Y' \arrow[l, "p"]
\end{tikzcd}
\tag*{$(*)$}
\]
si $p$ est une $\DD$-équivalence, alors $q$ aussi, \textit{i.e.}
\[
\begin{tikzcd}[row sep=small]
H^*_{\DD}(X,M) \arrow[r, "\sim"] & H^*_{\DD}(X',M) \\
\wr| & \wr| \\
H^*_{\DD}(Y, f_*(M_X)) & H^*_{\DD}(Y', f'_*(M_{X'}))
\end{tikzcd}
\]
pour tout $M$ dans $\mathcal{A} = \DD(e)$. Mais $F = f_*(M_X)$ est localement constant et se calcule “fibre par fibre”, de même pour $F' = f'_*(M_{X'})$.

\msp{page 81}

L'équivalence \textit{a}) $\Leftrightarrow$ \textit{b}$'$) se déduit de \textit{a}) $\Leftrightarrow$ \textit{b}) par dualité (la condition \textit{a}) étant autoduale), ou en reprenant le même raisonnement pour \textit{a}) $\Rightarrow$ \textit{b}$'$), et pour \textit{b}$'$) $\Rightarrow$ \textit{a}) aussi (en remplaçant la cohomologie par l'homologie).

\dgrblock{dgrb49}{The coefficient lost in the attempted proof}
\begin{dcomment}{exposition}{p81:en remplaçant la cohomologie par l'homologie}{The coefficient lost in the attempted proof}
The proposition is an experiment whose unsuccessful part Grothendieck subsequently identifies himself. Its three conditions must therefore be read together with the marginal corrections. Condition (a) says that pullback along every base change of $f$ preserves the entire class of $\DD$-equivalences. Conditions (b) and (b$'$) say that constant coefficients have universally well-behaved cohomological or homological direct images. The first condition implies the other two. The converse requires information about coefficients which the first definition of a $\DD$-equivalence does not provide.

Here is the precise point. In a pullback square with lower horizontal map $p:Y'\to Y$, universal fibrewise calculation gives an isomorphism
\[
 p^*f_*M_X\simeq f'_*M_{X'}.
\]
To conclude that $q:X'\to X$ is a $\DD$-equivalence, one must compare the cohomology of $Y$ and $Y'$ with coefficient $F=f_*M_X$. This coefficient is locally constant. It need not be constant. The assumption that $p$ is a $\DD$-equivalence compares cohomology only with constant coefficients. There is no applicable hypothesis at the last step of the displayed argument. This is the gap, and the examples on manuscript pages 87--95 make it an actual obstruction.

The surviving implication can be proved without this step. For an object $y$ of a base category, the categories $Y/y$ and $\{y\}$ have final objects. The functors between them, and the functors $Y/y_1\to Y/y_0$ induced by arrows of $Y$, are $\DD$-equivalences under the asphericity axiom. Pulling them back along $f$ gives exactly the inclusions $X_y\to X/y$ and the transition functors between the comma categories. Condition (a), applied also after a preliminary base change, therefore implies (b). Using undercategories gives (b$'$). Nothing in this argument proves the converse.

The correction in footnote 35, restricting the test categories to weakly contractible categories, is adequate once locally constant objects on such categories are constant. The qualification is essential: replacing weak contractibility by $\DD$-acyclicity is insufficient. Nor does connectedness of a base remove monodromy. We shall return below to a two-object covering that makes this visible with set-valued coefficients alone. Thus the sentence in footnote 36 about all truncated homotopy theories cannot be retained under the usual interpretation of $\DD_i$; its case $i=0$ already fails.
\end{dcomment}


\textbf{Définition 2.} Si les conditions équivalentes de la proposition 2 sont satisfaites, on dit que $f$ est une $\DD$-\textit{fibration}.

N.B. C'est une notion \textit{à la fois} locale et “colocale” sur $Y$ (alors que fibration gauche est locale, fibration droite colocale).

\printedpage{50}
% head: A. Grothendieck : Dérivateurs, Ch. I, Édition des Universités de Montpellier II et Paris VII
% foot: Version 8 juillet 2001

\textbf{Proposition 3.} a) \textit{La notion de $\DD$-fibration est stable par changement de base et par composition.}\footnote[38]{La stabilité par composition \textit{n'est} pas claire.}

\indent b) \textit{Soient}
\[
\begin{tikzcd}[column sep=normal, row sep=large]
X \arrow[rr, "f"] \arrow[dr, "p"'] & \phantom{X} & X' \arrow[dl, "p'"] \\
& Y &
\arrow[from=1-2, to=2-2, phantom, "\text{\footnotesize comm}"]
\end{tikzcd}
\]
\textit{avec $p,p'$ des $\DD$-fibrations (il suffit que $p,p'$ soient des $\DD$-fibrations cohomologiques, ou que ce soient des $\DD$-fibrations homologiques, ou que $p,p'$ soient propres, ou $p,p'$ lisses). Si $f$ induit des $\DD$-équivalences dans les fibres, alors $f$ est une $\DD$-équivalence.}\footnote[39]{Peut se préciser considérablement dans le cas où $p$, $p'$ sont des $\DD$-fibrations, \textit{cf.} prop. 6, page 84.}

\indent\textsc{Démonstration.} \textit{a}) est immédiat sous la condition \textit{a}) (prop. 2) de la notion de $\DD$-fibration. Pour \textit{b}), supposons que $f$ induise des $\DD$-équivalences sur les fibres. On veux prouver que % [sic] On veux -> On veut, verb agreement
\[
H^*_{\DD}(X,M) \longrightarrow H^*_{\DD}(X',M)
\]
est un isomorphisme, pour tout $M$ dans $\mathcal{A}$, mais cette flèche s'identifie à
\[
H^*_{\DD}(Y, p_*(M_X)) \longrightarrow H^*_{\DD}(Y, p'_*(M_{X'}))
\]
déduite de
\[
p_*(M_X) \longrightarrow p'_*(M_{X'}) \qquad\qquad (\text{NB}\ \ M_X \simeq f^*M_{X'}\,)
\]
\msp{page 82}
et il suffit de prouver que cette flèche est un isomorphisme. Pour ceci, on regarde fibre par fibre, et on utilise le fait que “les images directes se calculent fibre par fibre”.

\dgrblock{dgrb50}{Fibration terminology during the revision}
\begin{dcomment}{exposition}{p82:les images directes se calculent fibre par fibre}{Fibration terminology during the revision}
The notation temporarily conceals two different classes of functors. If ``$\DD$-fibration'' means the pullback-preservation condition (a), closure under composition is immediate: a $\DD$-equivalence is first pulled back along one factor and then along the other. Base change stability is immediate for the same reason. If the word instead means universal cohomological fibrewise calculation, that proof is unavailable. In a composite, the intermediate coefficient is a locally constant direct image, and the available comparison controls only direct images of coefficients pulled back from the base. The reservation in footnote 38 is consequently a substantive correction to the argument, rather than a doubt about the formal stability of condition (a).

The fibrewise comparison theorem in Proposition 3(b) survives. Write $f:X\to X'$ over $Y$, with projections $p$ and $p'$. The natural cohomological comparison is
\[
 p'_*M_{X'}\longrightarrow p_*M_X,
\]
obtained by applying $p'_*$ to the unit $M_{X'}\to f_*f^*M_{X'}$. If both direct images are computed on strict fibres, its component at $y$ is the cohomological map induced by $f_y:X_y\to X'_y$. A $\DD$-equivalence on each fibre therefore makes this map invertible by the detection axiom, and applying the direct image $Y\to e$ proves the assertion. The three displayed cohomological arrows in the printed proof point in the reverse direction; the theorem does not depend on that slip. The homological proof uses the counit and has the covariant direction.

The lemma on manuscript page 85 strengthens the sufficient criterion. The property that $f_y$ is a $\DD$-equivalence propagates across every arrow of $Y$, by applying two-out-of-three to the square of comma categories. It is therefore enough to test one object in each connected component. The proposed converse, from the total-space equivalence to fibre equivalences, fails because taking global cohomology can annihilate a nonzero local system. This does not contradict pointwise detection in $\DD(Y)$: the functor of global sections is generally not conservative.
\end{dcomment}


\textbf{Question.} Soit $f : X \to Y$ une $\DD$-fibration, et soit $F \in \DD(Y)$. Alors $f_*f^*F$ se calcule-t-il par passage aux fibres? Même question pour $f_!f^*F$. Pour la première propriété, concernant $f_*$, on montre que cela marche si $F$ est de la forme $p_*(M_Y)$, où $p : Y' \to Y$ est $\DD$-\textit{propre,} et $M \in \Ob\, \mathcal{A}$. D'autre part, la sous-catégorie pleine $C$ de $\DD(Y)$ formée des $F$ pour lesquels ça marche est (sauf erreur) “stable par $\DD$-$\mathsf{hot}$-$\varprojlim$” en le sens suivant : si on a une catégorie $I$ avec objet final $e$ formant une sous-catégorie fermée (\textit{i.e.} $e \to x \Rightarrow x = e$), et si on a $\mathcal{F} \in \DD(Y \times I)$, d'où $I^\circ \to \DD(Y)$, $i \mapsto \mathcal{F}_i$ (diagramme squelettique), si les $\mathcal{F}_i$ pour $i \neq e$ sont dans $C$, et si $\mathcal{F}$ est dans l'image essentielle de $j_*$, où $j : Y \times (I \setminus \{e\}) \hookrightarrow Y \times I$ est l'inclusion, alors $\mathcal{F}_e$ est aussi dans $C$.

\dgrblock{dgrb51}{Two useful classes of coefficients}
\begin{dcomment}{exposition}{p82:Pour la première propriété, concernant}{Two useful classes of coefficients}
The question about arbitrary coefficients should be separated from the assertion actually needed in the preceding proof. A universally cohomological fibration already calculates $f_*f^*F$ on strict fibres whenever $F$ is locally constant, under the local constancy axiom. Two further operations on coefficients explain the tentative observations here.

Suppose $p:Y'\to Y$ is $\DD$-proper, $M\in\DD(e)$, and $F=p_*M_{Y'}$. Form $X'=X\times_Y Y'$, with maps $q:X'\to X$ and $f':X'\to Y'$. Proper base change gives
\[
 f^*p_*M_{Y'}\simeq q_*M_{X'},\qquad
 f_*f^*p_*M_{Y'}\simeq p_*f'_*M_{X'}.
\]
The fibre formula for $p_*$ turns evaluation at $y$ into cohomology over $Y'_y$. Universal fibrewise calculation for $f'$ identifies the restriction of $f'_*M_{X'}$ to $Y'_y$ with the direct image along
$X_y\times Y'_y\to Y'_y$. The resulting object is
$H^*_{\DD}(X_y\times Y'_y,M)$.
On the other side, $F_y=H^*_{\DD}(Y'_y,M)$, and the product interchange formula identifies
$H^*_{\DD}(X_y,(F_y)_{X_y})$ with the same object. All identifications are the canonical mates, so this proves the desired comparison, not merely an accidental isomorphism between its endpoints. The printed $p_*(M_Y)$ must here be understood as $p_*(M_{Y'})$.

There is also a formal closure argument. Fix $y$. The two sides of the fibre comparison, as functors of $F\in\DD(Y)$, preserve coherent homotopy limits: restrictions preserve limits computed in the diagram variable, and right Kan extensions compose. Their comparison is a natural transformation between these continuous functors. If it is invertible on a coherent diagram of coefficients, it is invertible on its homotopy limit, by applying the two functors and then the canonical limit comparison. In particular this proves the punctured-cone assertion in the text, provided the indicated diagram shapes belong to $\Dia$. No claim about arbitrary incoherent diagrams in the category $\DD(Y)$ is required. This distinction matters because an arbitrary diagram in a homotopy category need not have a chosen coherent lift.
\end{dcomment}


\indent On peut cependant considérer que si on veux que $f : X \to Y$ satisfasse à toutes les propriétés cohomologiques et homologiques qui s'attachent à la notion d'une “fibration” (“localement triviale”), on devrait mettre dans les conditions que $f$ est \textit{lisse}, et \textit{propre}, pour que les changements de base par $f$ commutent aux foncteurs $g_*$, $g_!$. % [sic] on veux -> on veut, verb agreement

\dgrblock{dgrb52}{The fibre formula for coefficients from the base}
\begin{dcomment}{development}{p82:pour que les changements de base par $f$ commutent aux foncteurs}{The fibre formula for coefficients from the base}
The question on this page has an affirmative answer for the universal notion of fibration retained on manuscript pages 117--118. The restriction on the coefficient is essential to the argument: it may vary arbitrarily on the base, but the coefficient on the total category is its inverse image. We give the argument because it also answers the question about trivial fibrations on page 84. Throughout this block the domain contains the comma categories, categories of factorizations, and their pullbacks used below. The usual domain of all small categories has this property; so does the domain of finite posets when the categories under consideration are finite posets.

\textbf{Proposition (fibre calculation for inverse images).}\ours
Suppose that $\D$ takes admissible disjoint unions to products, has both homotopy Kan extensions and their pointwise formulas, and detects isomorphisms pointwise. All disjoint unions indexed by the hom-sets below are assumed admissible. Let $f:X\to Y$ be a weak cohomological $\D$-fibration after every base change. Then, for every $F\in\D(Y)$ and $y\in Y$, the canonical comparison
\[
 (f_*f^*F)_y\longrightarrow
 H^*_{\D}(X_y,(F_y)_{X_y})
\]
is an isomorphism. These comparisons also imply base change for $f_*f^*F$ along an arbitrary functor into $Y$. The dual assertion holds for $f_!f^*F$ when $f$ is universally a weak homological fibration.

Fix $y$ and write $H=X/y$, $S=X_y$, $v:H\to Y$ for $(x,f(x)\to y)\mapsto f(x)$, and $s:H\to e$, $t:S\to e$ for the projections. The two functors in the comparison are the right adjoints
\[
 T=s_*v^*,\qquad R=t_*t^*y^* :\D(Y)\longrightarrow\D(e).
\]
The pointwise formula identifies $T$ with evaluation of $f_*f^*$ at $y$. Their left adjoints are respectively $v_!s^*$ and $y_!t_!t^*$. It therefore suffices to prove that the left mate of the comparison,
\[
 y_!t_!t^*\longrightarrow v_!s^*,
\]
is invertible. We can test this on $M\in\D(e)$ and then at each object $z\in Y$.

For an arrow $\gamma:z\to y$, let $E_\gamma$ be the category of factorizations
\[
 z\xrightarrow{\beta}c\xrightarrow{\alpha}y,
 \qquad \alpha\beta=\gamma.
\]
A morphism from $(c,\beta,\alpha)$ to $(c',\beta',\alpha')$ is $h:c\to c'$ with $h\beta=\beta'$ and $\alpha'h=\alpha$. This category has initial object $(z,1_z,\gamma)$ and terminal object $(y,\gamma,1_y)$. Put
\[
 K_\gamma=X\times_Y E_\gamma,
\]
where $E_\gamma\to Y$ remembers the middle object. Thus an object of $K_\gamma$ is a factorization $z\to f(x)\to y$ of $\gamma$, and its morphisms are arrows of $X$ making the two triangles commute.

The pullback $K_\gamma\to E_\gamma$ is a weak cohomological fibration by hypothesis. Its strict fibre over the terminal factorization is $X_y$. Since slicing over a terminal object changes nothing, the first clause in the definition of a weak cohomological fibration says precisely that
\[
 X_y\longrightarrow K_\gamma,
 \qquad x\longmapsto(x,\gamma,1_y),
\]
is a $\D$-equivalence. Consequently it induces an isomorphism on homology with every constant coefficient $M$. Here we use the equivalence of the homological and cohomological tests for a $\D$-equivalence established on manuscript pages 61--62; no assertion about arbitrary local systems is being substituted for it.

Now $z\backslash v$ is the disjoint union of the $K_\gamma$ over the discrete set $Y(z,y)$. Indeed, its objects give the two arrows $z\to f(x)\to y$, and every morphism preserves their composite. The pointwise formula for left Kan extension and the coproduct axiom identify the component of the mate at $z$ with
\[
 \coprod_{\gamma\in Y(z,y)}H_*^{\D}(X_y,M)
 \longrightarrow
 \coprod_{\gamma\in Y(z,y)}H_*^{\D}(K_\gamma,M).
\]
Each summand is induced by the displayed inclusion of $X_y$ and is an isomorphism. Hence the mate is an isomorphism by pointwise detection, and so is the original comparison between right adjoints. This proves the fibre formula. For a base change $a:Y'\to Y$, both sides of the base-change map at $y'\in Y'$ identify with the same cohomology of $X_{a(y')}$ with value $F_{a(y')}$. Naturality of the comparisons gives the base-change assertion. Applying the argument with the opposite indexing categories and opposite coefficient categories gives the homological assertion.\hfill$\square$

The mechanism is the relative form of the factorization-category proof of homotopy exactness in \citet[Theorem 3.8]{GrothPontoShulman2014}. The full-faithfulness consequence in the next block has an all-derivator formulation in the homotopical epimorphism criterion of \citet[Theorem 6.10]{GrothStovicek2016Trees}. We have checked these references, as well as \citet{Maltsiniotis2005Asphericite} and \citet{Cisinski2009LocallyConstant}; we did not locate this precise answer to the page-82 question in those texts. The label marks our application and proof here, not a claim of priority. Notice what the proof does not require: neither properness of $f$ nor local constancy of $F$ on $Y$.
\end{dcomment}


\printedpage{51}
% head: A. Grothendieck : Dérivateurs, Ch. I, Édition des Universités de Montpellier II et Paris VII
% foot: Version 8 juillet 2001

\msp{page 83}

Signalons :

\textbf{Proposition 4.} \textit{Soit $f : X \to Y$} propre. \textit{Pour que $f$ soit une $\DD$-fibration, il faut et il suffit que pour tout $y_1 \to y_0$ dans $Y$, le foncteur $X/y_1 \to X/y_0$ induit soit une $\DD$-équivalence} (i.e. \textit{que $f$ soit une $\DD$-fibration} cohomologique). \textit{Dualement, si $f$ est $\DD$-lisse, pour que ce soit une $\DD$-fibration, il faut et il suffit que pour tout $y_1 \to y_0$ dans $Y$, $y_0\backslash X \to y_1\backslash X$ soit une $\DD$-équivalence.}

\indent On pourrait appeler $\DD$-fibrations \textit{parfaites} les foncteurs à la fois propres et lisses qui sont des $\DD$-fibrations cohomologiques (ou ce qui revient au même des $\DD$-fibrations homologiques).

\textbf{Corollaire.} \textit{Supposons $f$ $\Cat$-cofibration (donc propre). Pour que $f$ soit une $\DD$-fibration, il faut et il suffit que les foncteurs $X_{y_0} \to X_{y_1}$ (pour $y_0 \to u_1$ dans $Y$) soient des $\DD$-équivalences. Duallement, pour $\Cat$-fibration (donc lisse).} % [sic] u_1 -> y_1, the arrow is y_0 -> y_1 ; [sic] Duallement -> Dualement, spelling

\textbf{Proposition 5.} \textit{Soit $f : X \to Y$ dans} $\Dia$. \textit{Conditions équivalentes :}
\begin{enumerate}[label={}, leftmargin=2em, labelsep=0pt]
\item a) \textit{$f$ est une $\DD$-équivalence, et le reste par tout changement de base ;}
\item b) \textit{$f$ est une $\DD$-fibration à fibres $\DD$-contractiles.}
\end{enumerate}

\textsc{Démonstration.} \textit{a}) $\Rightarrow$ \textit{b}) : Que \textit{a}) implique que $f$ soit une $\DD$-fibration est immédiat sur la définition de celle-ci (condition \textit{a}) de la proposition 2), et que les fibres sont contractiles est tautologique.

\textit{b}) $\Rightarrow$ \textit{a}) : Il suffit de montrer que \textit{b}) implique que $f$ est une $\DD$-équivalence. Mais c'est un cas particulier de la proposition 3 \textit{b}) (en y prenant $X' = Y$).

\dgrblock{dgrb53}{Properness separates the two tests}
\begin{dcomment}{exposition}{p83:Mais c'est un cas particulier de la proposition 3}{Properness separates the two tests}
For a proper functor the first half of the cohomological-fibration condition is already supplied by proper base change: the inclusion $X_y\to X/y$ is coaspherical, hence a $\DD$-equivalence, and right Kan extensions calculate on strict fibres. What remains is local constancy of $f_*M_X$ for every $M$. This remaining condition is stable under arbitrary change of base. Indeed properness identifies the new direct image with the restriction of the old one, and restriction preserves local constancy. This proves the criterion of Proposition 4 in the universal cohomological sense. Its identification with the revised bilateral fibration notion uses the local-constancy hypotheses explained at manuscript pages 97--100.

For a Grothendieck opfibration, a choice of cocartesian lifts gives transition functors $\alpha_!:X_{y_0}\to X_{y_1}$ for arrows $\alpha:y_0\to y_1$. They are defined up to canonical natural isomorphism. The map on constant-coefficient cohomology between the two strict fibres is the map induced by $\alpha_!$, with the cohomological direction reversed. Consequently local constancy says precisely that every transition functor is a $\DD$-equivalence. The criterion is independent of the chosen cleavage. The dual statement uses cartesian transport in a Grothendieck fibration.

The later adjective ``trivial'' is also justified without the false converse of Proposition 2. A universal $\DD$-equivalence has $\DD$-contractible strict fibres by pulling back to points. Its further pullbacks are again universal equivalences. It has the cohomological and homological fibration tests because the relevant maps form triangles whose other two sides are these universal equivalences. Conversely, a universal cohomological fibration with $\DD$-contractible fibres has $M_Y\to f_*M_X$ invertible at every point, hence invertible in $\DD(Y)$. Taking global cohomology proves that it is a $\DD$-equivalence, and the argument repeats after every base change. Thus Proposition 5 survives the revision. The stronger full-faithfulness question on page 84 requires the separate factorization-category argument given in its commentary.
\end{dcomment}


\msp{page 84}

Quand ces conditions sont satisfaites, on s'attend à ce que le foncteur
\[
f^* : \DD(Y) \longrightarrow \DD(X)
\]
soit pleinement fidèle, \textit{i.e.} que les morphisme d'adjonction % [sic] les morphisme -> les morphismes, plural agreement
\[
F \longrightarrow f_*f^*F \quad , \qquad f_!f^* \longrightarrow F % [sic] f_!f^* -> f_!f^*F, the object F is missing
\]
dans $\DD(Y)$ soient des isomorphismes. (C'est vrai si $f$ est $\DD$-propre ou $\DD$-lisse, cf. cor. 3, p.71.) Cela revient encore à dire que les $f_*f^*F$ se calculent par passage aux fibres, ou encore itou pour les $f_!f^*F$. Cela nous ramène à la question de la page 82.

\dgrblock{dgrb54}{Universal equivalences give fully faithful restriction}
\begin{dcomment}{development}{p84:Cela nous ramène à la question de la page 82.}{Universal equivalences give fully faithful restriction}
The expectation in the preceding paragraph follows from the fibre formula just proved. It is stronger than the assertion that $f$ preserves cohomology of constant objects, because it controls all morphisms between arbitrary coefficients on $Y$.

\textbf{Corollary (the question about trivial fibrations).}\ours
Under the domain and derivator hypotheses of the preceding development, suppose that every pullback of $f:X\to Y$ is a $\D$-equivalence. Then
\[
 f^*:\D(Y)\longrightarrow\D(X)
\]
is fully faithful. Both the unit $F\to f_*f^*F$ and the counit $f_!f^*F\to F$ are isomorphisms, for every $F\in\D(Y)$.

To prove this, first pull back over a point $y\to Y$. The hypothesis says that $X_y$ is $\D$-aspherical. Pulling back over $Y/y$ likewise shows that $X/y$ is $\D$-aspherical, since $Y/y$ has a terminal object. The inclusion $X_y\to X/y$ and all transition functors between these comma categories are therefore $\D$-equivalences by two-out-of-three. The same argument works after any base change. Thus $f$ is universally a weak cohomological fibration. At $y$, compose the adjunction unit with the comparison proved in the preceding block:
\[
 F_y\longrightarrow(f_*f^*F)_y
 \longrightarrow H^*_{\D}(X_y,(F_y)_{X_y}).
\]
The composite is the constant-coefficient unit for $X_y\to e$, which is invertible by asphericity. The second arrow is invertible by the proposition. The first arrow is therefore invertible, and pointwise detection proves that $F\to f_*f^*F$ is an isomorphism. This is equivalent to full faithfulness of $f^*$. Since $f^*$ also has the left adjoint $f_!$, full faithfulness makes the other displayed adjunction counit invertible as well.\hfill$\square$

There is also a useful direct connection with the modern criterion. For each arrow $\gamma:z\to y$, the category $E_\gamma$ of its factorizations has an initial and a terminal object. Universality makes $K_\gamma=X\times_YE_\gamma$ aspherical. These are exactly the factorization categories controlling full faithfulness. In the special case of weak equivalences of nerves, \citet[Theorem 6.10]{GrothStovicek2016Trees} therefore supplies the same conclusion simultaneously for every derivator. The argument above gives the coefficient-relative conclusion asked for in the manuscript and explains why universality, rather than the single global equivalence $X\to Y$, is the relevant hypothesis.
\end{dcomment}


\textbf{Définition 3.} On dit alors que $f$ est une $\DD$-fibration \textit{triviale}.

\textbf{Proposition 6} (précise prop. 3 page 81). \textit{Considérons un triangle commutatif dans} $\Dia$
\[
\begin{tikzcd}[column sep=small, row sep=large]
X \arrow[rr, "f"] \arrow[dr, "p"'] & & X' \arrow[dl, "p'"] \\
& S & & & .
\end{tikzcd}
\]

\printedpage{52}
% head: A. Grothendieck : Dérivateurs, Ch. I, Édition des Universités de Montpellier II et Paris VII
% foot: Version 8 juillet 2001

\indent a) \textit{Supposons que $p$ et $p'$ soient des $\DD$-fibrations gauches, ou que $p$ et $p'$ soient des $\DD$-fibrations droites. Soit $S_0$ une partie de $S$ qui rencontre toute composante connexe de $S$. Pour que $f$ soit une $\DD$-équivalence, il} suffit \textit{que pour tout $s_0 \in S_0$, $f$ induise une $\DD$-équivalence $X_{s_0} \to Y_{s_0}$.}

\indent b) \textit{Cette condition est aussi nécessaire}\footnote[40]{Pas prouvé!\footnotemark[41] Même si $X' = S$, ce n'est pas clair, même si $\DD = \mathsf{HOTAB}$, $f$ propre et lisse. Mais OK si $S$ est $\DD$-contractile.}\footnotetext[41]{\underline{Faux}, \textit{cf.} contre-exemple, p. 87.} \textit{si $p$ et $p'$ sont des $\DD$-fibrations. Donc dans ce cas, $f$ est une $\DD$-équivalence si et seulement si $f$ induit des $\DD$-équivalences fibre par fibre.}

\msp{page 85}

\textsc{Démonstration.} a) Il suffit, en vertu de la proposition 3, de montrer que l'hypothèse implique que pour \textit{tout} $s \in S$, $f$ induit une $\DD$-équivalence $f_s : X_s \to X'_s$. Cela résulte aussitôt du lemme :

\textbf{Lemme.} \textit{Si $p$ et $p'$ sont toutes deux soit des $\DD$-fibrations gauches, soit des $\DD$-fibrations droites, et si $s_1 \to s_0$ est une flèche dans $S$, alors $f_{s_1}$ est une $\DD$-équivalence si et seulement si $f_{s_0}$ l'est. (Donc l'ensemble $S' \subset \Ob\, S$ des $s \in S$ tels que $f_s$ soit une $\DD$-équivalence est une réunion de composantes connexes de $S$.)}

\indent Supposons par exemple que $p,p'$ soient des $\DD$-fibrations gauches. Considérons, pour $s \in S$, le diagramme
\[
\begin{tikzcd}[row sep=large]
X_s \arrow[r, "f_s"] \arrow[d] & X'_s \arrow[d] \\
X/s \arrow[r, "f/s"'] & X'/s & .
\end{tikzcd}
\]
Les flèches verticales sont des $\DD$-équivalences par l'hypothèse sur $p,p'$, donc $f_s$ est une $\DD$-équivalence si et seulement si $f/s$ l'est. Considérons maintenant
\[
\begin{tikzcd}[row sep=large]
X/s_1 \arrow[r, "f/s_1"] \arrow[d] & X'/s_1 \arrow[d] \\
X/s_0 \arrow[r, "f/s_0"'] & X'/s_0 & .
\end{tikzcd}
\]
Ici encore, l'hypothèse sur $p,p'$ implique que les flèches verticales sont des $\DD$-équivalences, donc $f_{s_1}$ est une $\DD$-équivalence si et seulement si $f_{s_0}$ l'est.

\msp{page 86}

Si $Y' = S$, l'énoncé (non prouvé) dit que si $f : X \to S$ est une $\DD$-fibration et une $\DD$-équivalence, alors les fibres sont $\DD$-contractiles. Ce n'est pas sûr même pour une projection $X = T \times S \to S$. Si $\eta = \mathit{hot}_0(S)$, $\xi = \mathit{hot}_0(T)$, la question se réduit à celle-ci : Si pour tout $M$ dans $\mathcal{A}$, on a : l'homomorphisme
\[
\alpha_M : \xi \times_{\DD} M \longrightarrow M
\]
\printedpage{53}
% head: A. Grothendieck : Dérivateurs, Ch. I, Édition des Universités de Montpellier II et Paris VII
% foot: Version 8 juillet 2001
satisfait
\[
\eta \times_{\DD} \alpha_M : \eta \times_{\DD} (\xi \times_{\DD} M) \xrightarrow{\sim} \eta \times_{\DD} M
\]
est un isomorphisme, alors $\alpha_M$ est lui même un isomorphisme.

\dgrblock{dgrb55}{Why a product cannot supply the counterexample}
\begin{dcomment}{development}{p86:est lui même un isomorphisme.}{Why a product cannot supply the counterexample}
The product case singled out on manuscript page 86 has a short resolution that does not use an additive coefficient theory. We give it to isolate the role of twisting in the counterexamples which follow. It is a formal retract argument; we claim no priority for it. The locally constant-coefficient results of \citet{Cisinski2009LocallyConstant} address the general homotopical framework, but we did not find this particular paragraph of the manuscript treated there.

\textbf{Cancellation for a nonempty product base.}\ours\
Assume the action of indexing homotopy types on $\DD(e)$ is functorial and that $S$ has an object. If the product projection $S\times T\to S$ is a $\DD$-equivalence, then $T\to e$ is a $\DD$-equivalence.

\emph{Proof.} Choose $s:e\to S$ and write $p:S\to e$. Then $ps=1_e$. For each coefficient $M$, let $\alpha_M:T\times_{\DD}M\to M$ be the augmentation. Functoriality gives a retract diagram in the arrow category of $\DD(e)$,
\[
 \alpha_M\longrightarrow S\times_{\DD}\alpha_M
                   \longrightarrow\alpha_M,
\]
whose composite is the identity: the first map is induced by $s$, the second by $p$. The middle arrow is an isomorphism by the product hypothesis and Fubini for homology. A retract of an isomorphism is an isomorphism. Explicitly, if an arrow $a$ is a retract of an invertible arrow $b$, with inclusions $i_0,i_1$ and retractions $r_0,r_1$, then $r_0b^{-1}i_1$ is an inverse of $a$, as the two commutative squares show. Hence every $\alpha_M$ is invertible. This is the asserted $\DD$-contractibility of $T$. \hfill$\square$

The nonempty-base hypothesis cannot be omitted: $\varnothing\times T\to\varnothing$ is an isomorphism for every $T$. The argument establishes that the proposed product example cannot work; it does not establish the analogous assertion for a bundle with monodromy. The next example exploits exactly that distinction.
\end{dcomment}


\msp{page 87}

Voici un \textit{contre-exemple}\footnote[42]{Voir un contre-exemple encore plus simple pages 93, 94. (Mais il a l'inconvenient de ne pas être à fibres connexes.) % [sic] inconvenient -> inconvénient, missing accent
} que je donne d'abord dans le contexte des espaces topologiques, au lieu de $\Cat$. Je vais construire une fibration localement triviale et non triviale
\[
\begin{tikzcd}[column sep=0.2em, row sep=large]
X \arrow[d, "f"'] & \supset S^1 & \\
S & = S^1 & \qquad ,
\end{tikzcd}
\]
de fibre $S^1$, telle que $f$ soit une $\DD$-équivalence pour tout $\DD = \mathsf{Der}(k\text{-}\mathsf{Mod})$, $k$ étant un anneau où 2 est inversible. Cela signifie que
\[
H_*(X,k) \longrightarrow H_*(S,k)
\]
est un isomorphisme modulo 2-torsion (noyau et conoyau sont annulés par une puissance de 2 - en fait par 2 ici), ou encore $H_*(X,k) \xrightarrow{\sim} H_*(S,k)$ si 2 inversible dans $k$. Mais la fibre $S^1$ n'est pas $\DD$-contractile (sauf si $k = 0$!).\footnote[43]{Exemple beaucoup plus simple encore : on prend $X = S^1$, et $X \to S$ le revêtement à deux feuillets. Si on prend le revêtement à $n$ feuillets, c'est aussi OK si $k$ anneau où $n$ est inversible.}

\indent La fibration est donnée à isomorphisme unique près, en se fixant un point $s \in S$, et en se donnant un automorphisme $u$ de la fibre $X_s$ (“monodromie”). (La classe d'isomorphie du fibré ne dépend que de la classe d'isotopie de $u$). Je prend pour $u$ la symétrie (inversant l'orientation) $z \mapsto z^{-1}$ du cercle standard $\{z \in \mathbb{C} \mid |z| = 1\}$. Cette fibration “se réalise” par des ensembles ordonés $X \to S$, $S^1$ étant représenté (sur la base, comme sur les fibres) par % [sic] Je prend -> Je prends, verb agreement ; [sic] ordonés -> ordonnés, spelling
\[
\begin{tikzcd}[column sep=normal, row sep=normal]
& & a & & \\
S \ = & 0 \arrow[ur] \arrow[dr] & & 0' \arrow[ul] \arrow[dl] & = \ S' \amalg_R S'' \quad , \\
& & b & &
\end{tikzcd}
\]
où $R$ est la catégorie discrète ayant deux éléments $a$, $b$. On définit $X$ comme $X' \amalg_Y X''$, où $X' = X|S'$, $X'' = X|S''$, $Y = X|R$.\footnote[44]{N.B. $S'$, $S''$, $R = S' \cap S''$ fermés dans $S$, $X'$, $X''$, $Y = X' \cap X''$ fermés dans $X$.}

% ---- from Der-ChI-234_p056-060.tex
\printedpage{54}
% head: A. Grothendieck : Dérivateurs, Ch. I, Édition des Universités de Montpellier II et Paris VII

\msp{page 88}
On pose
\[ X' = S' \times \bar S \qquad\qquad X'' = S'' \times \bar S \]
($\bar S$ copie conforme de $S$, cf ci-dessous) et il faut, pour recoller $X'$ et $X''$, se donner un isomorphisme entre $X'|R \simeq R \times \bar S = \bar S \amalg \bar S$ et $X''|R \simeq R \times \bar S = \bar S \amalg \bar S$ qui soit un $R$-foncteur, \textit{i.e.} qui provient de deux isomorphismes
\[
\begin{aligned}
\bar S &\xrightarrow{u_a} \bar S \\
\bar S &\xrightarrow{u_b} \bar S \qquad .
\end{aligned}
\]
On prend
\[
\begin{aligned}
u_a &= \mathsf{id}_{\bar S} \\
u_b &= \text{symétrie par rapport à l'axe } \bar 0, \bar 0' \ \textit{i.e.} % Doubt resolved: PDF leaf 18 (printed p.54) reads u_b, with plain b.
\\
&\phantom{=}\quad u_b \text{ laisse fixes } \bar 0, \bar 0' \text{ et intervertit } \bar a \text{ et } \bar b.
\end{aligned}
\]
(C'est plus clair si on désidentifie la base et la fibre, en prenant fibre
\[
\bar S \quad = \quad
\begin{tikzcd}
 & \bar a & \\
\bar 0 \arrow[ur] \arrow[dr] & & \bar 0' \arrow[ul] \arrow[dl] \\
 & \bar b &
\end{tikzcd}
\qquad ).
\]

\dgrblock{dgrb56}{The Klein bottle and its sign coefficient}
\begin{dcomment}{examples}{p88:plus clair si on désidentifie la base et la fibre}{The Klein bottle and its sign coefficient}
Take a nonzero commutative ring $k$ in which $2$ is invertible. Let $r:S^1\to S^1$ be complex conjugation and form its mapping torus
\[
 K=(S^1\times[0,1])/((z,1)\sim(r(z),0)).
\]
Projection to $[0,1]/(0\sim1)$ gives the Klein-bottle bundle $f:K\to S^1$. We use the cohomological convention in which $k[-1]$ has its copy of $k$ in degree $1$. Grothendieck's displayed $k[1]$ is explicitly annotated ``en degr\'e 1'' and should be translated accordingly, not corrected as a transcription error.

A cell structure has one vertex, two edges $a,b$, and one two-cell with attaching word $aba^{-1}b$. Here $a$ projects once around the base and $b$ runs around the fibre. The cellular complex is
\[
 0\longrightarrow k\xrightarrow{(0,2)}k a\oplus k b
                         \xrightarrow{0}k\longrightarrow0.
\]
The map to the cellular complex of the base sends $a$ to its generator and kills $b$ and the two-cell. Its kernel is the two-term complex $k\xrightarrow{2}k$, which is contractible. Consequently the projection induces an equivalence with every constant derived $k$-module coefficient, including unbounded complexes: the displayed chain contraction survives applying the appropriate Hom complex. Integrally, one instead gets $H_1(K;\mathbb Z)=\mathbb Z\oplus\mathbb Z/2$, and the kernel of the map to $H_1(S^1;\mathbb Z)$ is the indicated $2$-torsion. The fibre still has $H_1(S^1;k)=k$ and is not $\DD$-contractible.

The same computation in local-coefficient language explains why global cohomology loses the fibre. Let $L$ be the rank-one local system on the base with monodromy $-1$. The fibre cohomology gives
\[
 Rf_*k_K\simeq k_{S^1}\oplus L[-1].
\]
This splitting can be obtained from the cellular circle complex with its reflection action. Alternatively, the fixed point of the reflection gives a section and splits the degree-zero summand. The remaining summand is $L[-1]$. Global sections of $L$ are computed by $k\xrightarrow{-2}k$, in degrees $0,1$; the complex is zero in the derived category. Thus the extra summand is visible over each point and invisible after taking global sections.

The finite categorical model in the text has four objects in the base and four in each fibre. Choose $S'=\{0,a,b\}$ and $S''=\{0',a,b\}$. Their intersection is the discrete category $\{a,b\}$. Gluing the products with the four-object circle by the identity over $a$ and by reflection over $b$ produces sixteen objects. Each arrow in the base transports the fibre by an isomorphism. The projection is therefore both a Grothendieck fibration and an opfibration, and remains so after every base change. Its nerve realizes the same monodromy bundle. This supplies the categorical counterexample required by the manuscript, without appealing only to a topological analogy. Cellular calculations and the general language of local coefficients are reviewed in \citet[\S\S2.2, 3.H]{Hatcher2002}.
\end{dcomment}


Pour voir que $X$ est une $\D$-fibration sur $S$ (quelque soit $\D$) prouvons qu'il est \textit{propre}\footnote[45]{N.B. $f : X \to S$ est également \textit{lisse}, car c'est un fibré trivial (un produit) au dessus de deux ouverts $S - \{a\}$, $S - \{b\}$ qui recouvrent $S$. Mieux : $f$ est à la fois $\Cat$-fibrant et $\Cat$-cofibrant (car cette condition n'utilise que les propriétés des $X \times_S \boldsymbol{\Delta}_2 \to \boldsymbol{\Delta}_2$, pour $\boldsymbol{\Delta}_2 \to S$).} sur $S$ - on est ramenés à voir que $X' \to S'$ et $X'' \to S''$ sont propres, ce qui est trivial. Pour voir que c'est une $\D$-fibration, on est ramenés également au cas de $X' \to S'$ et $X'' \to S''$, où c'est aussi trivial. On a utilisé le : % [sic] quelque soit -> quel que soit, spelling

\textbf{Lemme.} \textit{Soient $f : X \to S$, $(S_i)_{i \in I}$ une famille de sous-catégories fermées recouvrant $S$. Pour tout $i$, soit $X_i = X \times_S S_i \xrightarrow{f_i} S_i$. Pour que $f$ soit propre, il faut et il suffit que les $f_i$ le soient. Pour que $f$ soit une $\D$-fibration propre il faut et il suffit que les $f_i$ le soient.}

\dgrblock{dgrb57}{Why the closed-cover test is sufficient}
\begin{dcomment}{exposition}{p88:Pour que $f$ soit propre, il faut et il suffit que les}{Why the closed-cover test is sufficient}
The gluing lemma has a local proof which does not require a descent theorem for arbitrary closed covers. A closed full subcategory here is a cosieve: whenever it contains the source of an arrow it contains its target. In the lifting-category criterion for properness, fix $x\in X$ and an arrow $f(x)\to s$ in the base. Choose a member $S_i$ of the cover containing $f(x)$. Its cosieve property forces $s$ and every intermediate object relevant to this lift into $S_i$. The entire category of lifts is therefore the same one computed in $X_i\to S_i$. Its $\DD$-contractibility follows from properness there. Necessity follows by base change.

Once properness has been obtained globally, the fibration condition asks whether the fibre-cohomology comparison along each arrow is invertible. Choose $S_i$ containing its source; the arrow and both endpoints lie in that subcategory. The local condition proves the comparison, and the pointwise criterion proves the global condition. This explains both assertions of the lemma. For the displayed four-object circle, the two closed subcategories $S'$ and $S''$ contain respectively $0$ and $0'$, together with both targets $a,b$, so they meet exactly these requirements. The dual open-cover argument explains smoothness.
\end{dcomment}


\msp{page 89}
Voici le petit calcul qui montre que $f : X \to S$ est une $\D$-équivalence. On fait le calcul en cohomologie.\footnote[46]{Mais dans le cas du revêtement à 2 feuillets de $S$, il n'y a pas même de calcul à faire!} On doit prouver que
\[ H^*(S,k) \longrightarrow H^*(X,k) \tag*{$(*)$} \]
\printedpage{55}
% head: A. Grothendieck : Dérivateurs, Ch. I, Édition des Universités de Montpellier II et Paris VII
est un isomorphisme (et de même pour tout coefficient constant $M$ sur $S$). Or
\[ H^*(X,k) \overset{\mathrm{def}}{=} \RR\Gamma_X(k_X) = \RR\Gamma_S(\RR f_*(k_X)) \quad , \]
et $(*)$ est induit par
\[ i : k_S \longrightarrow \RR f_* k_X \quad . \]
Or on a
\[ \RR f_* k_X \simeq \underbrace{k_S}_{= \mathsf{Im}(i)} \quad \oplus \quad \underbrace{k_S[1]}_{k_S \text{ en degré } 1}{}^{\underline{\mathrm{tordu}}} \]
savoir ``tordu'' par l'action de monodromie sur $S = S^1$, laquelle est la symétrie $\lambda \mapsto -\lambda$. On a
\[ \RR\Gamma_S(k_S[1]^{\mathrm{tordu}}) = H^1(\underbrace{\pi_1(S)}_{= \mathbb{Z}}\, , k_S^{\mathrm{tordu}}) = k_{(\mathbb{Z})} = k/(u - \mathsf{id})k = k/2k = 0 \ , \]
d'où le résultat.

Notons que si on regarde
\[ X' = X \times_S X \xrightarrow{f'} S \quad , \]
on a par Künneth relatif
\[ \RR f'_*(k_{X'}) \simeq \RR f_*(k_X) \otimes^{\mathbb{L}} \RR f_*(k_X) \simeq k_S \oplus k_S[1]^t \oplus k_S[1]^t \oplus \underbrace{k_S[1]^t \otimes k_S[1]^t}_{\simeq k_S[2] \text{ non tordu}} \ , \]
d'où
\[ H^*_{\D}(X',k) \simeq \RR\Gamma_{X'}(k_{X'}) \simeq \RR\Gamma_S(\RR f'_*(k_{X'})) \simeq \RR\Gamma_S(k_S) \oplus \RR\Gamma_S(k_S)[2] \]
et
\[ H^*_{\D}(S,k) \longrightarrow H^*_{\D}(X',k) \]
n'est \textit{pas} un isomorphisme.

\msp{page 90}
Ainsi $f : X \to S$ qui est une $\D$ fibration (propre et lisse) satisfait bien les conditions équivalentes $b$), $b'$) de la proposition 2 (page 79), mais pas la condition $a$), puisque le changement de base
\[
\begin{tikzcd}[column sep=large, row sep=large]
X' = X \times_S X \arrow[r, "{q = pr_1}"] \arrow[d, "{f' = pr_2}"'] & X \arrow[d, "f"] \\
S' = X \arrow[r, "{p = f}"'] & S
\end{tikzcd}
\]
est une $\D$-équivalence $S' \to S$, alors que $X' \to X$ ne l'est pas.

\dgrblock{dgrb58}{Self-pullback restores the missing cohomology}
\begin{dcomment}{examples}{p90:alors que $X' \to X$ ne l'est pas.}{Self-pullback restores the missing cohomology}
The self-pullback $K\times_{S^1}K$ is the mapping torus of
$r\times r:T^2\to T^2$. Its fibre has cohomology $k$ in degrees $0$ and $2$, and $k^2$ in degree $1$. The monodromy acts trivially in degrees $0,2$ and as $-1$ on both degree-one generators. The degree-two action is $(-1)(-1)=1$. The two sign summands again have acyclic global sections, whereas the degree-two summand does not. Hence
\[
 H^j(K\times_{S^1}K;k)=
 \begin{cases} k,&j=0,1,2,3,\\0,&\text{otherwise}.\end{cases}
\]
The same conclusion follows directly from the mapping-torus complex: in each fibre degree its extra differential is $1-(r\times r)^*$, which is either zero or multiplication by $2$. This calculation avoids extension ambiguities and is valid over the stated ring.

The map $K\to S^1$ is a constant-coefficient equivalence, but its pullback $K\times_{S^1}K\to K$ cannot be one, because the source has the additional nonzero groups in degrees $2,3$. Thus even a bundle which is both proper and smooth in the categorical sense does not preserve all constant-coefficient equivalences under pullback. Properness controls the base-change comparison for a direct-image functor; it does not assert that the base-change map is an equivalence for whichever weak-equivalence class one has chosen on the base.

The easier covering example in footnote 43 exhibits the same distinction. The degree-$n$ map $S^1\to S^1$ is a $k$-homology equivalence when $n$ is invertible in $k$, since it is the identity in degree $0$ and multiplication by $n$ in degree $1$. Its strict fibre has $n$ points. Its self-pullback is a disjoint union of $n$ circles, so the pulled-back projection fails already on $H_0$ for $n>1$. This example also avoids any question about a derived splitting.
\end{dcomment}


Voici un autre contre-exemple, sur le même principe mais un peu plus simple. Soit $G$ un groupe fini (par exemple $G = \mathbb{Z}/2\mathbb{Z}$), faisons le opérer sur un objet $Z$ de $\Cat$, soit $S = [G] = B_G$ la catégorie réduite à un objet $a$, avec $G = \mathsf{End}(a)$. Soit $E_G$ la catégorie (isomorphe à $(B_G)/a$)
\[ \mathsf{Ob}\ E_G = G \quad , \qquad \mathsf{Hom}_{E_G}(g,g') = \{ h \in G \mid g = g'h \} \quad , \]
composition évidente - alors $E_G$ est un groupoïde trivial (1-connexe). Le goupe $G$ y opère à gauche librement par translations, et $B_G = E_G/G$. % [sic] goupe -> groupe, misprint

\printedpage{56}
% head: A. Grothendieck : Dérivateurs, Ch. I, Édition des Universités de Montpellier II et Paris VII

\msp{page 91}
Si $G$ opère à gauche sur $Z$, on peut alors l'utiliser pour tordre $Z$, par
\[ X = E_G \times_G Z = (E_G \times Z)/\text{opérations ``diagonales'' de G} \quad . \]
Donc on a un diagramme cartésien
\[
\begin{tikzcd}
 & X' \arrow[d, equal] & X'' \arrow[d, equal] \\
X \arrow[d, "f"'] & Z \times E_G \arrow[l] \arrow[d, "pr_2"] & Z \arrow[l, dashed, "{\overset{\D}{\sim}}"'] \arrow[d] \\
B_G \arrow[d, equal] & E_G \arrow[l] \arrow[d, equal] & e \arrow[l, dashed, "{\overset{\D}{\sim}}"'] \arrow[d, equal] \\
S & S' & S''
\end{tikzcd}
\]
Comme le changement de base $S' = E_G = (B_G)/a \to B_G = S$ transforme $X$ sur $B_G$ en un produit (lequel est un $\D$-fibré propre et lisse sur $S'$), on voit que $X$ sur $B_G$ est un $\D$-fibré propre et lisse (car c'est là une condition locale sur $B_G$).

On suppose maintenant que $G$ est fini d'ordre $n > 1$, et soit $k$ un anneau où $n$ est invesible, $\D = \mathsf{Der}(k - \mathsf{Mod})$. Alors $E_G \to B_G$ est une $\D$-équivalence, ou ce qui revient au même $e \to B_G$ l'est.\footnote[47]{En fait, $e$, $E_G$, $B_G$ sont $\D$-asphériques.} Pourtant, après changement de base soit par $S' = E_G$, soit par $S'' = e$, $X' \to X$ (ou ce qui revient au même, $X'' \to X$) n'a guère de raison d'en être aussi. On a en cohomologie % [sic] invesible -> inversible, misprint
\[ \RR\Gamma_X(k_X) \simeq \RR\Gamma_{BG}(f_*(k_X)) \simeq \RR\Gamma_G(G, H^*(Z,k)) \]
\msp{page 92}
donc ($n$ étant inversible dans $k$) :
\[ H^*(X,k) = \RR\Gamma_X(k_X) \simeq H^*(Z,k)^G \longrightarrow H^*(Z,k) \quad , \]
ce n'est pas en général un isomorphisme. On peut prendre par exemple $G = \mathbb{Z}/2\mathbb{Z}$, $Z \sim S^1$ comme avant, et $G$ opérant par symétrie de $S^1$ comme précédemment.

Cet exemple est particulièrement convainquant, car c'est un exemple \textit{avec $S$, $S'$ $\D$-asphériques,} et $S' \to S$ un isomorphisme local (en fait, un revêtement principal de $S$). C'est aussi un exemple d'un $\D$-fibré lisse et propre $X \to S$, sur une base \textit{asphérique} $S$, et tel que l'inclusion d'une fibre dans $X$ % [sic] convainquant -> convaincant, spelling of the adjective
\[ X_s \longrightarrow X \]
ne soit \textit{pas} une $\D$-équivalence.

Le seul reproche qu'on puisse faire à cet exemple, en comparaison avec le précédent, c'est que les catégories $X, S$ ne sont pas des ensembles ordonnés. Et $S$ ne peut être un ensemble ordonné fini, car l'espace classifiant d'un groupe fini a de la cohomologie en toutes dimensions (sauf erreur).

\dgrblock{dgrb59}{Averaging over a finite group}
\begin{dcomment}{examples}{p92:Le seul reproche qu'on puisse faire à cet exemple}{Averaging over a finite group}
The second counterexample places the same phenomenon over a base which is already $\DD$-acyclic. Let $G$ be finite of order invertible in $k$. The category $E_G$ has one arrow between every ordered pair of objects of $G$; in particular each object is both initial and final. The quotient by left translation is $BG$. A $G$-category $Z$ gives an associated bundle
\[
 X=(E_G\times Z)/G\longrightarrow BG.
\]
An arrow in the base acts on the fibre by an isomorphism, so this is again a fibration and an opfibration with universally locally constant fibre transport. Pulling back to the unique object of $BG$ recovers $Z$.

For a $k[G]$-module $V$, the averaging operator
\[
 e_G(v)=|G|^{-1}\sum_{g\in G}gv
\]
is an idempotent with image $V^G$. More than a splitting, it proves exactness of invariants: an invariant element of a quotient can be lifted arbitrarily and its lift averaged. The same argument applies degreewise to complexes, so cohomology commutes with invariants. Thus
\[
 H^*(X;k)=H^*(Z;k)^G,\qquad H^*(BG;k)=k
\]
with the second $k$ in degree zero. The first formula is a formula for cohomology; a functorial complex-level statement uses derived sections of $Z$ with their action, followed by invariants. It does not require choosing a splitting of an arbitrary complex into its cohomology groups.

Take $G=C_2$ and let $Z$ be the four-object circle with reflection. Its cohomology is the trivial representation in degree zero and the sign representation in degree one. Therefore $X$ and $BG$ are $\DD$-acyclic, whereas their common fibre $Z$ is not. This disproves the qualification ``OK si $S$ est $\DD$-contractile'' in footnote 40 as well as the unqualified necessity assertion of Proposition 6(b). The replacement needed there is complete $\DD$-asphericity, which controls locally constant coefficients, or ordinary weak contractibility under the appropriate axiom.

An even smaller choice is the discrete regular $G$-set $Z=G$. Then $X$ is equivalent to the contractible category $E_G$, while its fibre has $|G|$ connected components. This is a wholly categorical example with finitely many objects, although its nerve is not finite-dimensional. A nontrivial finite group cannot have a finite-poset classifying space: if it did, the finite covering corresponding to its universal cover would be a finite contractible complex, so its Euler characteristic would be both $1$ and $|G|$ times an integer. The phrase ``cohomology in all dimensions'' in the manuscript should be read as the obstruction to a bounded-dimensional classifying space, not as a claim of nonvanishing in every positive degree for every coefficient ring.
\end{dcomment}


\printedpage{57}
% head: A. Grothendieck : Dérivateurs, Ch. I, Édition des Universités de Montpellier II et Paris VII

\msp{page 93}
Mais on peut s'inspirer de cet exemple pour le modifier de façon qu'il se transpose à des ensembles ordonnés. Il suffit de réaliser, pour $G$ donné :

\textit{a}) Une opération \textit{libre} de $G$ sur un objet $E$ de $\Cat$, telle que $H^*(E,k)^G \simeq k$, \textit{i.e.} $B = E/G$ soit $\D$-asphérique.

\textit{b}) Une opération de $G$ sur un objet $Z$, telle que l'inclusion
\[ H^*(Z,k)^G \longrightarrow H^*(Z,k)^G \] % [sic] H^*(Z,k)^G \longrightarrow H^*(Z,k)^G -> H^*(Z,k)^G \longrightarrow H^*(Z,k), an inclusion of the invariants (cf. p. 56)
ne soit pas \texttt{[un]} isomorphisme.

\textit{c}) Choisir pour $E$, $Z$ des ensembles ordonnés, et s'assurer que $B = E/G$ est également représenté par une catégorie ordonnée. Sauf erreur, il en sera de même pour $X = E \times Z/G$. On aura en tout les cas % [sic] en tout les cas -> en tous les cas, agreement
\[
\begin{tikzcd}[column sep=large]
 & Z \times E \arrow[d, equal] & \\
X \arrow[d] & X' \arrow[l] \arrow[d] & Z \arrow[l] \arrow[d] \\
S \arrow[d, equal] & S' \arrow[l] \arrow[d, equal] & S'' = \{s\}\, , \quad s \in E \ . \arrow[l] \\
E/G & E & \text{\footnotesize asph} \\
\text{\footnotesize asph} & \substack{\text{asph} \\ \text{si on peut}} &
\end{tikzcd}
\]
Cela donne encore un exemple avec $X$ $\D$-fibré propre et lisse, sur une base $S$ $\D$-asphérique, tel que l'inclusion d'une fibre $X_s$ dans $X$ ne soit pas une $\D$-équivalence. On peut même se débrouiller pour que $H^*(Z,k)^G \simeq k$ ce qui implique
\msp{page 94}
que $X$ est lui aussi asphérique (mais les fibres de $X \to S$ ne sont pas $\D$-asphériques). Cela donne donc un autre contre-exemple au b) de la prop. 6 (p. 84).

Voici la façon la plus économique, je crois, de réaliser la situation que j'ai dite. Topologiquement, je prend % [sic] je prend -> je prends, conjugation
\[
\begin{aligned}
G &= \{\pm 1\} \\
E &= S^2 &\qquad& \text{sphère euclidienne, sur laquelle } \{\pm 1\} \\
  &      &&      \text{opère par antipodisme } x \mapsto -x \\
B &= S^2/\{\pm 1\} && \text{plan projectif réel} \\
Z &= S^1 && \text{où G opère comme précédemment par la symétrie\,.}
\end{aligned}
\]
C'est donc un exemple très proche de celui du début (p. 87), où on pouvait aussi considérer la base ($B = S^1$) comme quotient d'un $E$ ($= S^1$) par $G = \{\pm 1\}$ opérant librement (par
\printedpage{58}
% head: A. Grothendieck : Dérivateurs, Ch. I, Édition des Universités de Montpellier II et Paris VII
l'antipodisme $z \mapsto -z$ dans $\mathbb{C} \simeq \RR^2$). Mais dans cet exemple, la base n'était pas $\D$-asphérique.

La façon la plus économique de réaliser la sphère $S^2$ par un ensemble ordonné est sans doute par sa décomposition ``globulaire'' par trois paires de cellules de dimension 0, 1, 2, donnant l'ensemble ordonné
\[
I \qquad = \qquad
\begin{tikzcd}[row sep=large, column sep=large]
S_2^- & S_2^+ \\
S_1^- \arrow[u] \arrow[ur] & S_1^+ \arrow[u] \arrow[ul] \\
S_0^- \arrow[u] \arrow[ur] & S_0^+ \arrow[u] \arrow[ul]
\end{tikzcd}
\]
et l'antipodisme échange $S_i^-$ et $S_i^+$, mais le quotient n'est pas ordonné, c'est la catégorie (importante pour la théorie des 2-catégories) :
\msp{page 95}
\[
\begin{tikzcd}[row sep=large]
S_2 \\
S_1 \arrow[u, shift left, "p_1^-"] \arrow[u, shift right, "p_1^+"'] \\
S_0 \arrow[u, shift left, "p_0^-"] \arrow[u, shift right, "p_0^+"']
\end{tikzcd}
\qquad\qquad \text{avec} \qquad\qquad
\begin{cases}
p_1^- p_0^- = p_1^- p_0^+ \\
p_1^+ p_0^- = p_1^+ p_0^+
\end{cases}
\]
Mais prenons la subdivision barycentrique de la décomposition cellulaire précédente de $S$ dont les sommets correspondent aux éléments de $\mathsf{Ob}\ I$, les simplexes aux parties totalement ordonnés de $\mathsf{Ob}\ I$. On trouve l'octaèdre, ayant 6 sommets, 12 arêtes, 8 faces, et $G$ opère librement sur l'ensemble des simplexes, donc le quotient : 3 sommets, 6 arêtes, 4 faces, décomposition cellulaire du plan projectif. (Sa décomposition par trois droites.) On présente $S^1$ par l'ensemble ordonné (comme au début) % [sic] parties totalement ordonnés -> parties totalement ordonnées, agreement
\[
\begin{tikzcd}[row sep=large, column sep=large]
S_1^- & S_1^+ \\
S_0^- \arrow[u] \arrow[ur] & S_0^+ \arrow[u] \arrow[ul]
\end{tikzcd}
\qquad\qquad \textit{i.e.} \qquad\qquad
\begin{tikzcd}
 & S_1^+ & \\
0 = S_0^- \arrow[ur] \arrow[dr] & & S_0^+ = 0' \arrow[ul] \arrow[dl] \\
 & S_1^- &
\end{tikzcd}
\]
et $G$ y opère par symétrie échangeant $S_1^-$ et $S_1^+$. Donc l'ensemble ordonné $X$ aura $4 \times 13 = 52$ sommets et c'est modeste \ldots

\dgrblock{dgrb60}{The finite poset and its residual top class}
\begin{dcomment}{development}{p95:Donc l'ensemble ordonné $X$ aura}{The finite poset and its residual top class}
The passage to finite posets deserves a complete computation. The number $52$ in the text is correct. The claimed acyclicity of the resulting total space is not. We give both verifications, retaining the printed construction and its error in the source. The transfer calculation used below is standard; see \citet[\S3.G]{Hatcher2002}. We did not find a published discussion of this particular $52$-object calculation in the sources checked.

Footnote 42 describes the later example as having disconnected fibres, whereas the explicit choice $Z=S^1$ on manuscript page 94 has connected fibres. The covering alternative of footnote 43 and a discrete regular $G$-set instead give disconnected fibres; the intended reference in footnote 42 cannot be fixed unambiguously.

\textbf{The corrected finite model.}\ours\
Let $E$ be the face poset of the boundary of an octahedron. Its vertices are $\{\pm1,\pm2,\pm3\}$; a nonempty face is a subset containing at most one element of each antipodal pair. There are $6+12+8=26$ faces. The involution $\tau$ negates every vertex. Let $Z$ be the four-object circle poset with two vertices $v_0,v_1$ below two edges $e_0,e_1$, and let $\tau$ fix the vertices and exchange the edges. Then $B=E/C_2$ and $X=(E\times Z)/C_2$ are posets with $13$ and $52$ objects. Their realizations satisfy
\[
 |NB|\simeq\mathbb{RP}^2,\qquad
 |NX|\simeq(S^2\times S^1)/((x,z)\sim(-x,\bar z)).
\]
For $2\in k^\times$, their cohomology is
\[
 H^*(B;k)=k,\qquad
 H^j(X;k)=\begin{cases}k,&j=0,3,\\0,&\text{otherwise}.\end{cases}
\]
The fibre is a circle, so its inclusion into $X$ is not a $k$-cohomology equivalence, but $X\to B$ is not a $k$-cohomology equivalence either.

\emph{Proof.} If nonempty faces $\sigma,\rho$ satisfy both $\sigma\subseteq\rho$ and $\sigma\subseteq-\rho$, then $\rho$ contains an antipodal pair, which is impossible. Hence for fixed representatives there is at most one group element transporting a possible upper face to a face containing the lower one. Quotient incidence therefore gives at most one morphism between two orbit objects. Ranks prevent nontrivial cycles. The same argument applied to the first coordinate proves the assertion for $E\times Z$. Every chain in either quotient lifts uniquely once its first object is chosen, so the nerve of the quotient is the quotient of the nerve. The nerve of $E$ is the barycentric subdivision of the octahedral sphere, and the nerve of $Z$ is a circle. Realization commutes with these finite products and these free simplicial quotients. This proves the assertions about realization and object counts.

For a finite regular covering with group $G$, pullback identifies cochains on the quotient with invariant cochains upstairs: every simplex has $|G|$ lifts, permuted freely and transitively. Averaging makes invariants exact when $|G|$ is invertible. Consequently the quotient cohomology is the invariant part of the covering-space cohomology. The antipodal map on $S^2$ acts by $-1$ on $H^2$, whereas reflection on $S^1$ acts by $-1$ on $H^1$. In $H^*(S^2\times S^1;k)$ the actions in degrees $0,1,2,3$ are respectively $+1,-1,-1,+1$. This gives exactly the displayed answer. The product of the two anti-invariant generators is invariant and survives in degree $3$. It is the term omitted by the proposed implication on manuscript pages 93--94. \hfill$\square$

There is an earlier incidence slip on the same page. For the six-object globular poset, label the two orbits of arrows between consecutive levels by the product of the endpoint signs. Composition multiplies these signs. The orbit relations are consequently
\[
 p_1^-p_0^-=p_1^+p_0^+,
 \qquad p_1^-p_0^+=p_1^+p_0^-.
\]
The displayed relations in the edition instead hold $p_1$ fixed while changing $p_0$; they do not describe this antipodal quotient. This does not affect the face-poset construction just proved, which gives an unambiguous replacement.

The broader recipe (a)--(c) also needs care. If $E$ is the discrete regular $G$-set, then $E/G$ is a point and $X=(E\times Z)/G\cong Z$. Taking a nontrivial action on $H^*(Z;k)$ satisfies the proposed invariant-cohomology conditions, yet the fibre inclusion is the identity. Thus those conditions alone do not force a counterexample. For the selected connected $E=S^2$ the fibre inclusion does fail, as calculated above, but invariants of the separate factors do not determine invariants of their tensor product.

As an independent check, the accompanying exact rational chain computation gives simplex counts $(13,36,24)$ for $NB$, with boundary ranks $(12,24)$, and $(52,340,576,288)$ for $NX$, with boundary ranks $(51,289,287)$. The resulting Betti numbers are $(1,0,0)$ and $(1,0,0,1)$. These finite calculations corroborate the proof; they are not used to infer the missing degree-three class.
\end{dcomment}


% ---- from Der-ChI-234_p061-065.tex
\printedpage{59}
% head: A. Grothendieck : Dérivateurs, Ch. I, Édition des Universités de Montpellier II et Paris VII
% foot: Version 8 juillet 2001

\msp{page 96}
Finalement, il semble qu’il ne reste plus rien de la proposition 2 (p.79) introduisant la notion de $\DD$-fibration – pas même l’équivalence des deux conditions $b$) et $b'$) ($\DD$-fibration à gauche universelle, $\DD$-fibration à droite universelle). On songe à leur relation à une condition $a$) modifiée ainsi (pour $f : X \to Y$) :

$a$) Pour tout $Y'' \to Y' \to Y$, avec $Y'$ et $Y''$ \Hot-\underline{asphériques} (et pas seulement $\DD$-asphériques) $X'' = X \times_Y Y'' \to X' = X \times_Y Y'$ est une $\DD$-équivalence. (Quand on suppose que $Y''$ et $Y'$ ont des éléments finaux, c’est la condition que $f$ soit une $\DD$-fibration à gauche universelle; quand on suppose que $Y''$ et $Y'$ ont des éléments initiaux, c’est la condition $\DD$-fibration à droite universelle.)\footnote[48]{Il faut expliciter le raisonnement. Il suffit de tester même sur $Y'' \to Y'$, où $Y''$ est \underline{ponctuelle}, et où $Y'$ a un objet final (resp. initial).}

On a évidemment $a$) $\Rightarrow$ $b$) et $a$) $\Rightarrow$ $b'$) et on aimerait que ce soient des équivalences. Voici un axiome intéressant\footnote[49]{\textit{cf.} \textbf{Der 3} $a'$) page suivante.} sur $\DD$, qui semble vérifié dans tous les cas connus, qui assure l’équivalence souhaitée. Je vais en même temps reformuler \textbf{Der 3} en trois volets (\textit{cf.} p. 67).

\msp{page 97}
\textbf{Der} 3 $a$) Si $Z \in \Dia$ est \Hot-asphérique, alors $p : Z \to e$ est une $\DD$-équivalence, \textit{i.e.} par exemple $p^* : \mathcal{A} = \DD(e) \to \DD(Z)$ est pleinement fidèle.

$b$) Soit $f : Z \to Z'$ une flèche dans \Dia. Si $f$ est une \Hot-équivalence, c’est une $\DD$-équivalence (\textit{i.e.} $f^* : \DD(Z') \to \DD(Z)$ induit une \underline{équivalence} entre les sous-catégories de $Z$, $Z'$ \texttt{[plutôt de} $\DD(Z)$\texttt{,} $\DD(Z')$\texttt{]} formées des objets constants).

$a'$) (Précise considérablement $a$). Si $Z \in \Dia$ est \Hot-asphérique, alors on a $a$), et l’image essentielle de $\DD(e) = \mathcal{A}$ dans $\DD(Z)$ est formée des objets \underline{localement} \underline{constants} de $Z$, \textit{i.e.} des objets tels que le foncteur
\[
Z^{\circ} \longrightarrow \mathcal{A} \qquad z \longmapsto F_z
\]
soit localement constant (\textit{i.e.} transforme toute flèche en isomorphisme) donc isomorphe à un foncteur constant (puisque $Z$ est \Hot-asphérique, donc 1-connexe).

\dgrblock{dgrb61}{Locally constant objects as a homotopy invariant}
\begin{dcomment}{later}{p97:transforme toute flèche en isomorphisme}{Locally constant objects as a homotopy invariant}
The additional axiom sought here became a theorem for the usual derivator framework. \citet[Corollary 3.8]{Cisinski2009LocallyConstant} proves that a weak equivalence of nerves induces an equivalence on locally constant objects for every derivator on small categories. The model-category statement is Theorem 1.3 of the same paper; Theorem 3.13 organizes locally constant objects into a derivator. Thus Der~3~b$'$ is available under those modern hypotheses, rather than remaining an unresolved extra condition. This does not retroactively prove the assertion from only the provisional axioms currently being tried in the manuscript.

The modern interpretation uses the fundamental $\infty$-groupoid $\Pi_\infty|NZ|$. For a coefficient $\infty$-category $\mathcal C$ presenting the homotopy theory, locally constant coherent diagrams are functors from this groupoid to $\mathcal C$, and the category denoted $\DD(Z)_{\mathrm{lc}}$ is their homotopy category. If $|NZ|$ is contractible, evaluation gives the desired equivalence with $\DD(e)$. Constancy of the underlying functor into $\DD(e)$ alone is weaker: it forgets higher coherence. This is the same distinction emphasized in the commentary at manuscript p.~7 (``Coherent squares and what the axioms remember''), now applied to local systems.
\end{dcomment}


Il faut bien noter que dans $a'$), il ne suffit \underline{pas} de supposer que $Z$ soit $\DD$-asphérique, au lieu de \Hot-asphérique. Dès que la $\DD$-équivalence est strictement plus faible que \Hot, il semble que $a'$) devient toujours faux, en tant que condition sur $\DD$, si on y suppose seulement que $Z$ soit $\DD$-asphérique. Par exemple, même si $\DD = \DD_i$ (théorie des $i$-types d’homotopie relative), la catégorie
\msp{page 98}
des objets localement constants sur $Z$ s’exprime en termes du $(i+1)$-type d’homotopie, mais pas en termes de son $i$-type d’homotopie seul.

\dgrblock{dgrb62}{Sets already detect the missing monodromy}
\begin{dcomment}{examples}{p98:mais pas en termes de son $i$-type d’homotopie seul.}{Sets already detect the missing monodromy}
The shift from an $i$-type to an $(i+1)$-type has a concrete beginning. For the represented derivator
$\DD_0(A)=\mathsf{Set}^{A^{\mathrm{op}}}$, natural maps between constant diagrams on a nonempty connected category are simply maps of sets. Hence $e\to BG$ is a $\DD_0$-equivalence for every group $G$. Locally constant set-valued diagrams on $BG$, however, are $G$-sets, and restriction to $e$ forgets the action. For a nontrivial $G$, this functor is not full: there are maps of sets from a one-point set to the regular $G$-set, but no equivariant such map. Thus a $0$-type records connected components, while local systems of $0$-types already record the fundamental group.

This also supplies a direct test of footnote 36. Take $G=C_2$ and the covering $f:E_G\to BG$. Its strict fibres are the discrete two-element set $G$. Since the covering is a discrete fibration and a discrete opfibration, every base change is a covering of the same kind. Right Kan extension of a constant set $M$ has fibre $M^G$, and left Kan extension has fibre $G\times M$; arrows in the base act by permutations of $G$. Both resulting diagrams are locally constant and both Kan extensions calculate on strict fibres. Therefore conditions (b) and (b$'$) hold universally for $\DD_0$.

Now pull back along $e\to BG$. The induced map of total categories is $G_{\mathrm{disc}}\to E_G$. It is not a $\DD_0$-equivalence: for a two-element coefficient set $M$, the constant-coefficient limit on the connected category $E_G$ is $M$, whereas on the discrete fibre it is $M^2$. The comparison is the diagonal map and is not bijective. Thus (b) does not imply (a), even for the derivator of $0$-types. The correct positive assertion for ordinary weak homotopy equivalences uses their invariance of locally constant coefficients, as in the preceding block; it does not replace them by $\DD_0$-equivalences. If $\Dia$ is restricted to finite posets, the connected double cover of the four-object circle gives the same counterexample. Its total poset has eight objects and its self-pullback is the disjoint union of two copies of that total poset. Both bases are connected, while the pulled-back total categories have different numbers of components.
\end{dcomment}


Il semble bien que ce soit l’axiome \textbf{Der} 3 ($a'$) qui assure l’existence d’une bonne notion (“bilatère”) de $\DD$-fibration, qui soit stable par changement de base, et qui se réduise soit
\printedpage{60}
% head: A. Grothendieck : Dérivateurs, Ch. I, Édition des Universités de Montpellier II et Paris VII
% foot: Version 8 juillet 2001
à “$\DD$-fibration à gauche universelle” soit à “$\DD$-fibration à droite universelle” et donc qui soit en plus une condition \underline{vérifiable} pratiquement. Dans le cas où $f$ est propre ou lisse, le critère est particulièrement simple, puisque il se borne à la condition que les $f_* f^* a_Y$, ou les $f_! f^* a_Y$ (pour $a \in \mathcal{A}$), soient des objets localements constants sur $Y$ (donc \underline{constants} si $Y$ est \underline{\underline{\Hot}}-asphérique, en vertu de \textbf{Der} 3 ($a'$). % [sic] localements -> localement, plural ending on the adverb; % [sic] "(donc ... ($a'$)." -> closing parenthesis after "($a'$)" missing, unbalanced parenthesis

Je vais également introduire une variante renforcée de \textbf{Der} 3 $b$) :

\textbf{Der} 3 $b'$) Soit $Z' \to Z$ une \Hot-équivalence dans \Dia. Alors $f^* : \DD(Z) \to \DD(Z')$ induit une équivalence entre les catégories des coefficients \underline{localements} \underline{constants} sur $Z$ et sur $Z'$.\footnote[50]{\textbf{Der} 3 $b'_0$) : le foncteur est fidèle\\
\hspace*{6em}$b'_1$) : le foncteur est pleinement fidèle\\
\hspace*{6em}$b'_2$) : le foncteur est une équivalence} % [sic] localements -> localement, plural ending on the adverb

\msp{page 99}
Dans certaines questions il suffit de savoir que le dit foncteur est pleinement fidèle. Cela implique que
\[
H^*_{\DD}(Z,F) \overset{\sim}{\longrightarrow} H^*_{\DD}(Z', f^*F)
\]
quand $F$ est \underline{localement} \underline{constant} sur $Z$\footnote[51]{On pourra dire que $f : Z' \to Z$ est une $\DD$-équivalence 0-\textit{complète à gauche,} et 0-\textit{complète à droite} dans le cas homologique.} (mais il n’est pas clair que cette condition implique toujours la pleine fidélité\footnote[52]{Mais c’est OK si on a des \underline{\Hom} internes dans $\DD(Z)_{\mathrm{loc.const.}}$ satisfaisant
\[
\Hom(F,G) \simeq \Hom(e_Z, \underline{\Hom}(F,G)) \quad ,
\]
où $e$ est l’élément final de $\mathcal{A}$, les \underline{\Hom} étant compatibles aux images inverses.}) et de même pour l’homologie
\[
H_*^{\DD}(Z', f^*F) \longrightarrow H_*^{\DD}(Z,F) \quad .
\]
Moyennant cette condition, on a la :

\textbf{Proposition :} \textit{On suppose} \textbf{Der} 3 \textup{b}$'$) \textit{satisfait. Soit $f : X \to Y$ dans} \Dia. \textit{Pour que $f$ soit une $\DD$-fibration, il faut et il suffit que pour $Y'' \to Y'$ sur $Y$ qui soit une} \Hot\textit{-équivalence, $X'' = X \times_Y Y'' \to X' = X \times_Y Y'$ soit une $\DD$-équivalence.}

À vrai dire, on peut énoncer quelque chose de raisonnable, sans supposer l’axiome \textbf{Der} 3 $b'$). On doit introduire la :

\textbf{Définition :} Soit $f : Z' \to Z$ dans \Dia. On dit que $f$ et une $\DD$-équivalence 1-\textit{complète} si le foncteur $f^* : \DD(Z) \to \DD(Z')$ est pleinement fidèle sur la sous-catégorie pleine de $\DD(Z)$ formée des coefficients localement constants. On dit que $f$ est une $\DD$-équivalence 2-\textit{complète} si elle induit une équivalence entre les catégories des coefficients localement constants. % [sic] f et une -> f est une, verb "est" misprinted as "et"

\msp{page 100}
On dit que $Z$ est \textit{complètement} $\DD$-\textit{asphérique}, si $Z \to e$ est une $\DD$-équivalence 2-complète.

\printedpage{61}
% head: A. Grothendieck : Dérivateurs, Ch. I, Édition des Universités de Montpellier II et Paris VII
% foot: Version 8 juillet 2001

Ainsi, l’axiome \textbf{Der} 3 $a'$) peut se reformuler, en disant que si $Z$ est \Hot-asphérique, alors $Z$ est complètement $\DD$-asphérique. Et de même, \textbf{Der} 3 $b'$) en disant que si $f : X \to Y$ dans \Dia{} est une \Hot-équivalence, alors $f$ est une $\DD$-équivalence 2-complète. Nous allons donner une variante affaiblie, de \textbf{Der} 3 $a'$), qui suffit pour donner une notion raisonnable de $\DD$-fibration :

\textbf{Der} 3 $a''$) Tout $Z$ dans \Dia{} qui a soit un objet final, soit un objet initial, est complètement $\DD$-asphérique.

\textbf{Définition.} On dit que $f : X \to Y$ dans \Dia{} est une $\DD$-fibration, si pour
\[
Y'' \longrightarrow Y' \longrightarrow Y
\]
avec $Y'' \to Y'$ une $\DD$-équivalence 1-complète (par exemple $Y'$ et $Y''$ complètement $\DD$-asphèriques), $X'' = X \times_Y Y'' \longrightarrow X' = X \times_Y Y'$ est une $\DD$-équivalence. % [sic] asphèriques -> asphériques, grave accent instead of acute

Cela équivaut à chacune des conditions $b$), $b'$) de la prop. 2 (page 79), \textit{i.e.} que $f$ soit universellement une $\DD$ fibration à gauche, ou universellement une $\DD$-fibration à droite (grâce à \textbf{Der 3} $b''$ \texttt{[plutôt} $a''$\texttt{]}). % [sic] $\DD$ fibration -> $\DD$-fibration, hyphen missing

\msp{page 101}
\textbf{Définition 1.} $f : X \to Y$ dans \Dia{} est appelé $\DD$-\textit{équivalence cohomologique} (sous-entendu : \textit{complète}) si pour tout $F \in \Ob\,\DD(Y)$, $F$ \underline{localement} \underline{constant}, on a
\[
H^*_{\DD}(Y,F) \overset{\sim}{\longrightarrow} H^*_{\DD}(X, f^*(F)) \quad .
\]
Notion duale $\DD$-\textit{équivalence homologique}
\[
H_*^{\DD}(X, f^*(F)) \overset{\sim}{\longrightarrow} H_*^{\DD}(Y,F) \quad .
\]
Si $Y$ est complètement $\DD$-asphérique, alors les coefficients localement constants sur $Y$ sont constants, et $f$ est une $\DD$-équivalence cohomologique (resp. homologique) si et seulement si c’est une $\DD$-équivalence \textit{i.e.} si $X$ est $\DD$-asphérique.\footnote[53]{L’axiome \textbf{Der} 3 $b'_1$) implique que toute \Hot-équivalence est une $\DD$-équivalence cohomologique.}

\textbf{Proposition 1} (Soritale). \textit{Soient $X \overset{f}{\to} Y \overset{g}{\to} Z$. Si $f$ et $g$ sont des $\DD$-équivalences cohomologiques, alors $gf$ aussi. Si $gf$ et $f$ sont des $\DD$-équivalences cohomologiques, alors $g$ aussi. Si $gf$ est une $\DD$-équivalence cohomologique, et si $g$ est une $\DD$-équivalence \underline{2-complète}} (donc tout $F$ localement constant sur $Y$ provient d’un $G$ localement constant sur $Z$), \textit{alors $f$ est une $\DD$-équivalence cohomologique.}

\dgrblock{dgrb63}{Complete equivalences repair the base change argument}
\begin{dcomment}{exposition}{p101:provient d’un $G$ localement constant sur $Z$}{Complete equivalences repair the base change argument}
The adjective ``complete'' records which coefficients are being tested. A $1$-complete equivalence is fully faithful on the category of locally constant objects. A $2$-complete equivalence is in addition essentially surjective there. These are levels of categorical information, not degrees of homology. In particular a $1$-complete equivalence $p:Y'\to Y$ compares both homology and cohomology with every locally constant coefficient $F$ on $Y$.

For cohomology, adjunction gives
\[
 \Hom_{\DD(e)}(M,H^*_{\DD}(Y,F))
       =\Hom_{\DD(Y)}(M_Y,F).
\]
Both arguments on the right are locally constant. Full faithfulness identifies this group of morphisms with $\Hom_{\DD(Y')}(M_{Y'},p^*F)$, naturally in $M$. Yoneda gives the asserted cohomology isomorphism. Applying the same argument to morphisms from $F$ to a constant object proves homological invariance. The internal-Hom qualification in footnote 52 is a converse criterion: if locally constant internal Homs exist and commute with pullback, cohomological invariance applied to $\underline{\Hom}(F,G)$ yields full faithfulness. Its unit must be the stated cartesian unit, or an explicitly chosen monoidal unit; the zero object of a stable theory cannot play that role merely by being terminal.

We can now complete the repaired proof. Suppose $f$ is universally cohomological, and let $p$ be $1$-complete. For a constant coefficient $M$, put $F=f_*M_X$. The coefficient $F$ is locally constant, and universal fibre calculation identifies $p^*F$ with $f'_*M_{X'}$. The preceding Yoneda argument compares the global cohomology of these two objects and proves that $X'\to X$ is a $\DD$-equivalence. This establishes the new pullback-preservation condition. Conversely, under Der~3~a$''$, every functor between categories having final objects is $2$-complete, because both locally constant categories are equivalent to $\DD(e)$ and the restriction agrees with the identity on constants. Testing the new condition on $\{y\}\to Y/y$ and on transitions between slices therefore gives universal cohomological fibre calculation. Categories with initial objects give the homological condition. These arguments prove both directions and explain why the weaker endpoint axiom suffices at this stage.

The cancellation statements on manuscript page 101 have the same logic. Composing two comparisons proves closure under composition; cancelling the comparison for $f$ proves the implication from $gf$ and $f$ to $g$. To recover $f$ from $gf$, one must first express an arbitrary locally constant coefficient on $Y$ as the pullback of a coefficient on $Z$. The $2$-complete hypothesis on $g$ supplies exactly that essential-surjectivity step, and its full faithfulness supplies the comparison for $g$. Ordinary two-out-of-three would have omitted the coefficient which must descend.
\end{dcomment}


\textbf{Définition 2.} 1°) $f : X \to Y$ dans \Dia{} est appelé une $\DD$-\textit{fibration cohomologique complète,} je dis bien \underline{complète}, si

$a$) Pour tout $y \in \Ob\, Y$, $X_y \to X/y$ est une $\DD$-équivalence cohomologique (complète).

$b$) Pour tout $y' \to y$ dans $Y$, $X/y' \to X/y$ est une $\DD$-équivalence cohomologique (complète).

\printedpage{62}
% head: A. Grothendieck : Dérivateurs, Ch. I, Édition des Universités de Montpellier II et Paris VII
% foot: Version 8 juillet 2001

\msp{page 102}
$c$) Ces conditions restent satisfaites après tout changement de base $Y' \to Y$.

2°) Notion duale de fibration homologique complète.

\textsc{Exemple.} Si $f$ est \Hot-propre, les flèches des $a$) et $b$) sont des \Hot-équivalences, donc en vertu de \textbf{Der 3} $b'_1$) (qu’il faut admettre), ce sont des $\DD$-équivalences cohomologiques. Cela signifie que $a$) et $b$) sont valables, et $c$) aussi, car la notion de flèche \Hot-propre est stable par changement de base.

\dgrblock{dgrb64}{What the completion requires}
\begin{dcomment}{exposition}{p102:car la notion de flèche \Hot-propre est stable par changement de base.}{What the completion requires}
The quantifier on the coefficients is the decisive change. A weak cohomological fibration is tested by the constant objects $M_X$. A complete cohomological fibration must work with every locally constant object of each category obtained by base change. Such an object need not descend from $X$, and need not become constant on $X$. The monodromy examples just discussed explain why that distinction cannot be suppressed. The role of completion is to make the coefficients produced by one direct image admissible for a second direct image.

We keep the variance of the manuscript: objects of $\D(X)$ behave as coherent contravariant diagrams, and the pointwise formula for a right direct image is
\[
 (f_*F)_y\simeq H^*_{\D}(X/y,F|_{X/y}).
\]
Write $\mathcal L(X)\subseteq\D(X)$ for the full subcategory of locally constant objects. Local constancy means that every transition morphism is invertible; it does not mean that an isomorphism to a constant object has been chosen. For $r:A\to B$, the cohomological condition is the isomorphism
\[
 H^*_{\D}(B,L)\longrightarrow H^*_{\D}(A,r^*L)
 \qquad(L\in\mathcal L(B)).
\]
The conditions on $X_y\to X/y$ therefore compare a true fibre with a comma category for every local system on the latter. The conditions on $X/y\to X/y'$ then control the transport of those cohomology objects along the base. Universal repetition of the first comparison is what makes the resulting fibre formula compatible with arbitrary base changes.

The example immediately preceding Proposition~2 contains a genuine slip: $\Hot$-properness gives the fibre-to-comma comparison, but does not make the comma categories locally constant as $y$ varies. Take the inclusion $\{1\}\hookrightarrow[1]$, where $[1]=\{0\to1\}$. This is a closed immersion and hence proper, but its comma categories at $0$ and $1$ are respectively empty and terminal. The transition $\varnothing\to e$ is not a weak homotopy equivalence. Thus the example becomes correct after adding the transport condition, or by assuming that $f$ is both $\Hot$-proper and a weak $\HOT$-fibration. This is precisely the hypothesis separated out on manuscript page~107 and proved usable in Proposition~6 on page~118. The printed assertion is retained in the source and is audited as false as written.
\end{dcomment}


\textbf{Proposition 2.} \textit{Soient $f : X \to Y$ $\DD$-fibration cohomologique complète, $F \in \Ob\,\DD(X)$ localement constant sur $X$. Alors $f_*(F)$ est localement constant, sa formation commute à tout changement de base, et en particulier se calcule fibre par fibre.}\footnote[54]{Il y a une réciproque \textit{cf.} prop. 3$'$.}

\textbf{Corollaire 1.} \textit{Soit un diagramme cartésien}
\[
\begin{tikzcd}
X & X' \arrow[l, "q"'] \arrow[d, "f'"] \\
Y & Y' \arrow[l, "p"']
\arrow[from=1-1, to=2-1, "f"']
\end{tikzcd}
\quad .
\]
\textit{On suppose $f$ $\DD$-fibration cohomologique et $p$ une $\DD$-équivalence cohomologique. Alors $q$ est une $\DD$-équivalence cohomologique.}

\msp{page 103}
En effet, on a pour $F$ localement constant sur $X$
\[
H^*_{\DD}(X,F) \simeq H^*_{\DD}(Y, f_*F) \overset{\sim}{\longrightarrow} H^*_{\DD}(Y', p^*f_*F) \simeq H^*_{\DD}(Y', f'_*q^*F) \simeq H^*_{\DD}(X', q^*F)
\]
qed.

\textbf{Proposition 3.} \textit{La notion de $\DD$-fibration cohomologique complète est stable par changement de base} (trivial) \textit{et par composition.}

Pour la stabilité par composition, il suffit de noter le corollaire suivant de la proposition 2.

\textbf{Corollaire 2} (de prop. 2). \textit{Soit $f : X \to Y$. Pour que $f$ soit une $\DD$-fibration cohomologique complète, il faut et il suffit que pour tout $Y'' \overset{p}{\to} Y' \to Y$ sur $Y$, avec $p$ une $\DD$-équivalence cohomologique, la flèche correspondante $X'' = X \times_Y Y'' \longrightarrow X' = X \times_Y Y'$ soit aussi une $\DD$-équivalence cohomologique.}

C’est nécessaire par le corollaire 1, compte tenu que $f' : X' = X \times_Y Y' \to Y'$ est aussi une $\DD$-fibration cohomologique. C’est suffisant, car il suffit de l’appliquer aux cas $\{y\} \to Y/\{y\}$, $Y/y \to Y/y'$, qui sont des $\DD$-équivalences cohomologiques en
\msp{page 104}
vertu de \textbf{Der} 3 $b'_1$) (ou seulement \textbf{Der} 3 $a'$)).

\printedpage{63}
% head: A. Grothendieck : Dérivateurs, Ch. I, Édition des Universités de Montpellier II et Paris VII
% foot: Version 8 juillet 2001

\textbf{Proposition 4.} \textit{On suppose que} \textbf{Der} 3 a$'$) \textit{est satisfait} (\textit{i.e.} tout $Z$ dans \Dia{} qui est \Hot-asphérique, est complètement $\DD$-asphérique). \textit{Soit $f : X \to Y$ dans} \Dia. \textit{Les conditions suivantes sont équivalentes}

1°) \textit{$f$ est une $\DD$-fibration cohomologique \underline{complète}} (conditions $a$), $b$), $c$) de \texttt{[la]} définition 2, page 101).

2°) \textit{Pour tout $F \in$} \Ob\,$\DD(X)$\textit{, localement constant, $f_*(F)$ est localement constant, se calcule fibre par fibre; et ça reste vrai après tout changement de base $Y' \to Y$ (avec $F'$ localement constant sur $X' = X \times_Y Y'$). (Alors la formation de $f_*(F)$ commute à \underline{tout} changement de base).}

3°) \textit{Pour tout $Y'' \overset{p}{\to} Y'$ sur $Y$, si $p$ est une $\DD$-équivalence cohomologique, alors $X'' = X \times_Y Y'' \longrightarrow X' = X \times_Y Y'$ aussi.}

3$'$) \textit{Comme 3, en supposant que $Y''$ est la catégorie ponctuelle, et que $Y'$ a un élément final : il faut que $X'' \to X'$ soit une $\DD$-équivalence cohomologique.}

\dgrblock{dgrb65}{The universal criterion, with the quantifiers restored}
\begin{dcomment}{exposition}{p104:la catégorie ponctuelle, et que $Y'$ a un élément final}{The universal criterion, with the quantifiers restored}
Proposition~4 is more than a convenient restatement: it supplies the missing reason that the strengthened notion composes. We give the argument, keeping the hypothesis that every weakly contractible indexing category has only constant local systems, as in \textbf{Der}~3~$a'$. The same argument will apply to Proposition~9, after the terminology is revised.

First assume the comma-category conditions of Definition~2, universally in the base. If $F\in\mathcal L(X)$, its restriction to $X/y$ is locally constant. The comparison $X_y\to X/y$ gives
\[
 (f_*F)_y\simeq H^*_{\D}(X_y,F|_{X_y}).
\]
For an arrow $y\to y'$, the functor $X/y\to X/y'$ is a cohomological equivalence, so the resulting transition $(f_*F)_{y'}\to(f_*F)_y$ is an isomorphism. Thus $f_*F$ is locally constant. In a pullback square with base map $p:Y'\to Y$ and induced maps $f',q$, evaluate the base-change map
\[
 p^*f_*F\longrightarrow f'_*q^*F
\]
at $y'$. Both sides identify with the cohomology of the same strict fibre $X_{p(y')}$ with the same restricted coefficient. The identification of the comparison with the identity is the naturality of the pointwise Kan-extension formula. Pointwise detection of isomorphisms proves base change.

Suppose now that this direct-image property holds after every base change. For a cohomological equivalence $p:Y''\to Y'$, and $L\in\mathcal L(X')$, put $G=f'_*L$. The preceding property makes $G$ locally constant and identifies $p^*G$ with $f''_*q^*L$. The canonical comparison is consequently the composite
\[
 H^*_{\D}(X',L)\simeq H^*_{\D}(Y',G)
 \longrightarrow H^*_{\D}(Y'',p^*G)
 \simeq H^*_{\D}(X'',q^*L),
\]
whose middle map is invertible. This proves the universal preservation criterion. Notice that it tests all $L$ on $X'$, rather than just coefficients pulled back from $X$.

It remains to recover the comma conditions from the apparently much smaller test in~$3'$. Every functor $e\to B$ with $B$ admitting a terminal object is a cohomological equivalence under $a'$: all local systems on $B$ are constant and $B$ is $\D$-aspherical. Apply the test to the base change $Y/y\to Y$, selecting the terminal object $1_y$. Its inverse image is $X_y$, and the total space is $X/y$. This gives the first comma comparison. For $u:y\to y'$, use $Y/y'\to Y$ and select its object $u$. We get a cohomological equivalence $X_y\to X/y'$. It factors as
\[
 X_y\longrightarrow X/y\longrightarrow X/y'.
\]
The first map and the composite are cohomological equivalences. Testing a local system on $X/y'$ in the two comparisons shows that the second map is also one. This is the valid cancellation direction from Proposition~1 on manuscript page~101; no unqualified two-out-of-three assertion for cohomological equivalences is being used. As the point test was required over every base change, all the comparisons are universal.

These arguments also justify the parenthesis in Proposition~9 that only the fibre comparison needs to be imposed universally. Applied to $Y/y'$ and to its arbitrary object $u$, that one comparison already recovers transport. They explain why the converse concerns arbitrary locally constant coefficients on the pullbacks. Merely proving the fibre formula for restrictions of coefficients from the original $X$ would leave a real gap.
\end{dcomment}


\msp{page 105}
\textbf{Définition 3.} On dit que $f : X \to Y$ dans \Dia{} est une $\DD$-\textit{fibration complète} si c’est une $\DD$-fibration cohomologique complète et une $\DD$-fibration homologique complète.

C’est stable par changement de base, par composition.

La définition 3 équivaut à ceci : pour tout $p : Y'' \to Y'$ sur $Y$, avec $Y''$ et $Y'$ \Hot-asphériques, le morphisme correspondant $X'' \to X'$ est \underline{à} \underline{la} \underline{fois} une $\DD$-équivalence cohomologique, et une $\DD$-équivalence homologique. (Et il suffit même de le supposer pour $Y'' = e$ et pour $Y'$ avec objet final ou objet initial.) On peut songer à renforcer encore cette condition en exigeant (sans regarder à gauche ni à droite) que $X'' \to X'$ soit une $\DD$-équivalence 1-complète. Mais il n’est pas clair que cette notion renforcée soit stable par composition – à moins de la renforcer encore, en exigeant que si $Y'' \to Y'$ est une $\DD$-équivalence 1-complète, alors $X'' \to X'$ l’est aussi. L’ennui avec cette définition renforcée, c’est que faute d’avoir une description cohomologique ou homologique de la notion de
\msp{page 106}
$\DD$-équivalence 1-complète, je ne me sens pas capable de donner des exemples intéressants (non totalement triviaux) de $\DD$-fibrations en ce sens très fort. Sauf les \Hot-fibrations triviales (\textit{i.e.} à fibres \Hot-asphériques) – mais ce n’est pas drôle!

Mais quand on a des produits $*$ associatifs et unitaires et des \underline{\Hom} internes correspondants dans les $\DD(X)_{\mathrm{loc.\ const.}}$, compatibles avec \texttt{[les]} images inverses, donc on aura sur $X$ si $F$ et $G$ sont localement constants dans $\DD(X)$, et si $\varepsilon_X$ désigne l’unité pour $*$ dans $\DD(X)$,
\[
\Hom(F,G) \simeq \Hom(\varepsilon_X * F, G) \simeq \Hom(\varepsilon_X, \underline{\Hom}(F,G)) \simeq \Hom_{\mathcal{A}}(\varepsilon_e, p_*\, \underline{\Hom}(F,G)) \ ,
\]
car $\varepsilon_X = p^*(\varepsilon_e)$, où $p : X \to e$, donc dans ce cas une flèche $f : X \to Y$ dans \Dia{} est une $\DD$-équivalence 1-complète si (et seulement si) c’est une $\DD$-équivalence cohomologique. Donc dans ce cas, si $f$ est une $\DD$-fibration cohomologique complète, c’est aussi une $\DD$-fibration homologique complète, donc aussi une $\DD$-fibration complète (définition 3). (Et même au sens soit-disant renforcé de tantôt : si $Y'' \to Y'$ sur $Y$ est une $\DD$-équivalence 1-complète ($\Leftrightarrow$ une $\DD$-équivalence cohomologique!) alors $X'' \to X'$ aussi. % [sic] soit-disant -> soi-disant, spelling; % [sic] "(Et même ... aussi." -> closing parenthesis missing, unbalanced parenthesis

\dgrblock{dgrb66}{What internal Hom does, and what it does not do}
\begin{dcomment}{exposition}{p106:alors $X'' \to X'$ aussi.}{What internal Hom does, and what it does not do}
The internal-Hom argument is sound, provided its compatibility hypotheses are stated at the level where they are used. Suppose that each $\mathcal L(X)$ is a closed monoidal category, that its tensor unit is $\varepsilon_X=p_X^*\varepsilon_e$, and that restriction preserves the internal Hom between locally constant objects. For $F,G\in\mathcal L(Y)$, adjunction gives
\[
 \Hom_{\D(Y)}(F,G)
 \simeq\Hom_{\D(e)}\bigl(\varepsilon_e,
                p_{Y*}\underline{\Hom}_Y(F,G)\bigr).
\]
If $f:X\to Y$ is a cohomological equivalence, apply its defining comparison to the locally constant object $\underline{\Hom}_Y(F,G)$. The assumed comparison
\[
 f^*\underline{\Hom}_Y(F,G)
 \simeq\underline{\Hom}_X(f^*F,f^*G)
\]
then identifies the displayed Hom set with $\Hom_{\D(X)}(f^*F,f^*G)$. Thus $f^*$ is fully faithful on $\mathcal L(Y)$: $f$ is 1-complete.

Conversely, full faithfulness on local systems compares the morphisms from every constant object $M_Y$ to $F\in\mathcal L(Y)$. By adjunction these are the morphisms from $M$ to $p_{Y*}F$. Yoneda gives the cohomological comparison. Comparing morphisms from $F$ to every $M_Y$, instead, gives the homological comparison. This last argument needs no tensor product. It proves the two outgoing implications from 1-completeness in the diagram on page~113, while the internal-Hom hypothesis supplies the converse only on the cohomological side.

There is no need to assume that the unit detects isomorphisms. In the forward direction the cohomological comparison is already an isomorphism before $\Hom(\varepsilon_e,-)$ is applied; in the reverse direction Yoneda uses every $M$. Nor does the argument identify 1-completeness with 2-completeness: essential surjectivity of local systems is an additional demand.

One must also retain the adjective ``cohomological''. On page~122 the printed sentence says that an unqualified $\D$-equivalence is already 1-complete in the presence of internal Homs. That stronger assertion is false. For example, with $\D=\mathsf{Der}(\mathbb Q\text{-}\mathsf{Mod})$ and $G=C_2$, the map $e\to BG$ is a constant-coefficient equivalence, but restriction sends the sign representation to a one-dimensional vector space. There are no $G$-equivariant maps from the trivial representation to the sign representation, whereas after restriction there are nonzero maps. The categories of local systems have the usual derived tensor and internal Hom. The wording on pages~106 and~113, and the diagram on page~123, supply the correct repair: cohomological equivalence, not constant-coefficient equivalence.
\end{dcomment}


% ---- from Der-ChI-234_p066-070.tex
\printedpage{64}
% head: A. Grothendieck : Dérivateurs, Ch. I, Édition des Universités de Montpellier II et Paris VII
% foot: Version 8 juillet 2001

\msp{page 107}

Il se pose la question de trouver des $f$ qui soient des $\DD$-fibrations cohomologiques, ou homologiques. Je ne sais si je serai capable d'en trouver, qui ne soient déjà des \HOT-fibrations cohomologiques, ou homologiques (cela revient maintenant au même\footnote[55]{Les \HOT-équivalences cohomologiques $=$ \HOT-équivalences homologiques $=$ \HOT-équivalences sans plus.} ), car pratiquement je ne sais guère dire qu'une flèche $Z' \to Z$ dans $\Cat$ est une $\DD$-équivalence cohomologique complète, que si cette flèche est une \HOT-équivalence (et en admettant l'axiome \textbf{Der} 3 $b'_1$), ou sa variante affaiblie cohomologique).

\dgrblock{dgrb67}{A complete fibration beyond ordinary homotopy}
\begin{dcomment}{examples}{p107:Je ne sais si je serai capable d'en trouver, qui ne soient déjà des}{A complete fibration beyond ordinary homotopy}
The request for examples that are not already $\HOT$-fibrations has a simple answer with rational coefficients. Let $G=C_2$, let $BG$ denote its one-object category, and take the Grothendieck opfibration
\[
 f:E=\int_{[1]}(BG\longrightarrow e)\longrightarrow[1].
\]
Explicitly, $E$ has objects $a$ and $t$, with $\operatorname{End}(a)=G$, $\operatorname{End}(t)=\{1\}$, one arrow $a\to t$, and no arrow $t\to a$. The unique arrow to $t$ absorbs every $g\in G$. The fibres are $BG$ and $e$, and the transport is the collapse $r:BG\to e$.

For $\D=\mathsf{Der}(\mathbb Q\text{-}\mathsf{Mod})$, the functor $r$ is a cohomological equivalence. All coefficients on its target are constant, and invariants for a finite group over $\mathbb Q$ are exact. Indeed, averaging $v\mapsto |G|^{-1}\sum_{g\in G}gv$ is a natural projection onto invariants, and averaging a lift proves that invariants preserve epimorphisms. Consequently their right derived functor on any complex is computed degreewise. On a trivial $G$-complex $M$, this gives $R\Gamma(BG,M)=M$.

More generally, an opfibration whose transport functors are cohomological equivalences is a complete cohomological fibration. Properness supplies the fibre formula for every coefficient. For a locally constant coefficient $L$ on the total category and an arrow in the base, the cocartesian arrows give an isomorphism between $L$ on the source fibre and the pullback of $L$ on the target fibre. The cohomological equivalence of transport therefore makes the direct-image transition invertible. This argument survives an arbitrary pullback, since the new transport functors are the old ones indexed by the new arrows. Proposition~4, or Proposition~9 in the revised terminology, proves the assertion.

Applying this observation to $f$ shows that it is a complete cohomological fibration. Derived tensor products and internal Homs on rational local systems satisfy the hypotheses explained at page~106, so it is a complete homological fibration as well. It is not a $\HOT$-fibration: its transport $BG\to e$ is not a weak homotopy equivalence, since $\pi_1(BG)=C_2$. Thus stronger coefficient invariance is still relative to the chosen coefficients. Reversing the transport to $e\to BG$ destroys cohomological equivalence, and is exactly the feature exploited by the composition counterexample at manuscript page~118.
\end{dcomment}


Je vais donc chercher les \HOT-fibrations, \textit{i.e.} les $f : X \to Y$ tels que pour tout $Y'' \to Y' \to Y$, avec $Y''$ et $Y'$ asphériques (sous-entendu : pour \HOT), $X'' \to X'$ est une \Hot-équivalence. (Il suffit même de supposer $Y''$ réduit à un point, $Y'$ avec objet final $\ldots$ )

Supposons par exemple $f$ propre (\textit{i.e.} \HOT-propre), donc la condition $a$) de la définition 2, p. 101 satisfaite et de même après tout changement de base. Suffit-il de vérifier de plus la condition $b$), \textit{i.e.} que
\[ X/y \longrightarrow X/y' \quad , \qquad \text{pour } y \to y' \text{ dans } Y \ , \]
sont des \Hot-équivalences? (J'entends en restant sur $Y$, sans supposer aussi d'avoir la même chose après changement de base).
\msp{page 108}
Mais cela reste-il vrai après tout changement de base? % [sic] reste-il -> reste-t-il, forme interrogative

La situation semble ici meilleur pour $f$ lisse (au lieu de propre), en interprétant les conditions $a$) et $b$) sous forme homologique, comme signifiant que $f_!(e_X)$ est un objet \underline{localement} \underline{constant} de $\mathsf{HOT}(Y)$, qui se calcule fibre par fibre. Comme $f$ lisse implique que $f_!$ commute à tout changement de base, cette hypothèse sur $f$ (que $f$ soit une \HOT-fibration cohomologique, peut-être pas complète) est stable par changement de base, donc implique que c'est une \HOT-fibration cohomologique complète. Mais cet argument utilise les propriétés générales des dérivateurs pour \HOT, que je n'ai pas encore établies - surtout pas dans le domaine $\Cat$ tout entier! Donc il est sain de me méfier. Je préfère tâter la difficulté en restant d'abord dans le cadre du formalisme \underline{cohomologique} (non commutatif), dont j'ai une certaine maîtrise. Je suppose donc $f$ propre à nouveau, et que les % [sic] meilleur -> meilleure, accord avec « situation »
\[ X_y \longrightarrow X/y \quad , \qquad X/y \longrightarrow X/y' \]
\msp{page 109}
soient des \HOT-équivalences, je veux prouver la même chose après tout changement de base, au moins que dans la situation
\[
\begin{tikzcd}
X \arrow[r, leftarrow, "q"] \arrow[d, "f"'] & X' \arrow[r, leftarrow, "q'"] \arrow[d, "f'"'] & X'' \arrow[d, "f''"] \\
Y \arrow[r, leftarrow, "p"'] & Y' \arrow[r, leftarrow, "p'"'] & Y''
\end{tikzcd}
\]
\printedpage{65}
% head: A. Grothendieck : Dérivateurs, Ch. I, Édition des Universités de Montpellier II et Paris VII
% foot: Version 8 juillet 2001
si $Y'$, $Y''$ sont asphériques, alors $q'$ est une \Hot-équivalence. Si je veux l'établir par critère cohomologique, je dois commencer par prendre sur $X'$ un faisceau d'\textit{ensembles} localement constant $F$, et prouver que $\Gamma(X',F) \overset{\sim}{\to} \Gamma(X'',{q'}^{*}F)$. Comme ce faisceau ne provient pas de $X$, je suis ennuyé pour appliquer l'hypothèse sur $f$, qui concerne à priori des faisceaux d'ensembles constants sur $X$, ou sur les $X/y$. Mais on devrait avoir ceci :

\textbf{Lemme.} \textit{Soit $f : X'' \to X'$. Alors $f$ est une} \HOT-\textit{équivalence} i.e.
\[ H^{*}_{\mathsf{HOT}}(X', a_{X'}) \xrightarrow{\ \sim\ } H^{*}_{\mathsf{HOT}}(X'', a_{X''}) \]
\textit{est} \texttt{[un]} \textit{isomorphisme pour tout $a \in \mathsf{Hot} = \mathsf{HOT}(e)$, si et seulement si $f$ est une} \Hot-\textit{équivalence,}\footnote[56]{Attention ici \texttt{[mot illisible]} la distinction entre \HOT-équivalence et \Hot-équivalence.} \textit{et alors $f$ est même une} \HOT-\textit{équivalence} $2$-\textit{complète} (j'ai besoin seulement de : \HOT-équivalence cohomologique et homologique complète) i.e. \textit{pour tout $\xi \in \DD(X')$ \underline{localement} \underline{constant} sur $X'$, on a}\footnote[57]{Est-ce vrai aussi pour les dérivateurs $\DD_i$?}
\[ H^{*}_{\mathsf{HOT}}(X',\xi) \xrightarrow{\ \sim\ } H^{*}_{\mathsf{HOT}}(X'',{q'}^{*}(\xi)) \quad . \]
\msp{page 110}

La première assertion doit résulter des formules
\[ H^{*}_{\mathsf{HOT}}(Z,a_Z) \simeq \underline{\mathsf{Hom}}_{\mathsf{Hot}}(\mathsf{hot}(Z),a) \tag*{$(*)$} \]
\[ \text{(et aussi} \qquad\qquad H_{*}^{\mathsf{HOT}}(Z,a_Z) \simeq \mathsf{hot}(Z) \times a \qquad\qquad ) \]
donc si
\[ \mathsf{hot}(X'') \xrightarrow{\ \sim\ } \mathsf{hot}(X') \tag*{$(**)$} \]
est \texttt{[un]} isomorphisme, alors
\[ H^{*}_{\mathsf{HOT}}(X', a_{X'}) \longrightarrow H^{*}_{\mathsf{HOT}}(X'', a_{X''}) \]
est \texttt{[un]} isomorphisme, puisque c'est déduit de $(*)$ en prenant $\underline{\mathsf{Hom}}_{\mathsf{Hot}}(\,.\,,a)$. Inversement, si $(*)$ est toujours \texttt{[un]} isomorphisme, exprimant les deux membres \texttt{[mot illisible]} de $\underline{\mathsf{Hom}}_{\mathsf{Hot}}(\,.\,,a)$ et prenant les $\pi_0$, on trouve que
\[ \mathsf{Hom}_{\mathsf{Hot}}(\mathsf{hot}(X'),a) \xrightarrow{\ \sim\ } \mathsf{Hom}_{\mathsf{Hot}}(\mathsf{hot}(X''),a) \]
est un isomorphisme pour tout $a$, donc $(**)$ est un isomorphisme.

Mais le fait que ça implique que $f''$ soit une \HOT-équivalence cohomologique complète, \textit{i.e.} que la relation $(*)$ reste vraie pour un $\xi \in \DD(X')$ \underline{localement} \underline{constant} sur $X'$, me paraît une propriété délicate du dérivateur \HOT, sur laquelle il me faudra revenir.

Donc le lemme est loin d'être établi, faute d'avoir seulement construit le dérivateur \HOT, sans même parler d'en établir les propriétés spéciales délicates! Mais je n'ai aucune raison de douter du lemme, il fait partie de la kyrielle des intuitions de base que je voudrais fonder solidement.

\dgrblock{dgrb68}{The lemma for spaces}
\begin{dcomment}{later}{p110:Donc le lemme est loin d'être établi, faute d'avoir seulement construit le dérivateur}{The lemma for spaces}
The lemma on pages~109--110 is proved by the modern theory of locally constant diagrams. We can separate its two assertions, which need different arguments. Put $BX=|NX|$, and let $\mathcal S$ denote the homotopy theory of spaces. In the homotopy derivator $\HOT$, the constant-diagram formulas are
\[
 H^*_{\HOT}(X,M_X)\simeq\operatorname{Map}(BX,M),
 \qquad
 H_*^{\HOT}(X,M_X)\simeq BX\times M.
\]
These are derived mapping spaces and derived products, viewed in $\Hot$. Thus a weak equivalence $BX\to BY$ induces both constant-coefficient comparisons. Conversely, if all the cohomological comparisons are isomorphisms, their sets of components give bijections
\[
 [BY,M]\longrightarrow[BX,M]
\]
for every homotopy type $M$. Yoneda in $\Hot$ makes $BX\to BY$ an isomorphism. Homologically one can simply take $M=e$. This proves the identification of the manuscript's uppercase $\HOT$-equivalences with lowercase $\Hot$-equivalences, without conflating the homotopy derivator with its value at a point.

The local-coefficient assertion has another content: all coherence in a locally constant diagram must be transportable along a weak homotopy equivalence of the base. An invariant modern formulation identifies locally constant space-valued diagrams on $X$ with space-valued functors on the homotopy groupoid $\Pi_\infty(BX)$. Restriction along an equivalence of these groupoids is an equivalence of functor categories, and remains so on passing to homotopy categories. Hence a $\Hot$-equivalence is 2-complete for $\HOT$. It preserves cohomology and homology with every local coefficient by the adjunction argument in the preceding block.

A precise derivator reference is \citet[Theorem~1.3 and Corollary~3.8]{Cisinski2009LocallyConstant}: a nerve weak equivalence induces an equivalence of categories of locally constant objects, for model categories and, in the latter result, for arbitrary Grothendieck derivators on small categories. The construction of those local systems as a derivator, with both adjoints to its inclusion into the diagram derivator, is Theorem~3.13. Thus the strongest $b$-axiom proposed below is available in the usual modern setting; it is not a currently unproved property of spaces.

The higher-categorical description is compatible with straightening and unstraightening \citep[Theorem~2.2.1.2]{Lurie2009HTT}; the homotopy limit of a space-valued local system is its space of sections \citep[Corollary~3.3.3.4]{Lurie2009HTT}. This gives a geometric reading of the two comparisons: transport the family, then compare its sections or its total homotopy colimit. It also indicates exactly what is absent from the ordinary diagram in the homotopy category: the coherence of transport around paths and higher simplices.

Footnote~57 requires a qualification. If its question asks whether a genuine $\Hot$-equivalence induces 2-completeness for the derivators of truncated homotopy types, the answer is yes, by the same invariance theorem. If it asks whether a $\D_i$-equivalence alone has the same consequence, the answer is no already for $i=0$. The example in ``Sets already detect the missing monodromy'' at manuscript page~98 has $e\to BC_2$ a constant-coefficient equivalence for sets, while restriction of $C_2$-sets is not full: a map from the one-point trivial action to the regular action exists after forgetting the action and does not exist before. Truncating the coefficients can lose information about the base that the category of local systems still retains.
\end{dcomment}


\printedpage{66}
% head: A. Grothendieck : Dérivateurs, Ch. I, Édition des Universités de Montpellier II et Paris VII
% foot: Version 8 juillet 2001

\msp{page 111}

Le lemme implique qu'une $\DD = \mathsf{HOT}$-fibration (au sens de la définition 2 page 81, revue et corrigée à la page 96-97) est en fait déjà une \HOT-fibration cohomologique et homologique complète. Or on a vu que si $f$ est \underline{propre}, $f$ est une $\DD$-fibration si et seulement si les conditions précédentes \underline{sur $Y$} (sans changement de base) sont satisfaites. Et dualement pour $f$ lisse.

\begin{flushright}
11.11.90
\end{flushright}

Je sens le besoin de revoir les notions de $\DD$-fibrations en tous genres (p. 76 ff). D'abord les diverses variantes utiles de l'axiome \textbf{Der} 3 sur les dérivateurs :

\textbf{Der} 3 $a$)\footnote[58]{\textit{cf.} aussi la série \textbf{Der 3} $c \ldots$ , page 124 ff.} Si $Z \in \mathsf{Dia}$ est \Hot-asphérique, il est $\DD$-asphérique \textit{i.e.} $\mathcal{A} = \DD(e) \to \DD(Z)$ est pleinement fidèle. % Doubt resolved: attached PDF leaf 30 (printed p.66), footnote 58 has bold Der 3.

$a'$) Comme $a$), avec en plus la caractérisation de l'image essentielle de $\mathcal{A}$ dans $\DD(Z)$, comme formée des $\DD$-coefficients localements constants. % [sic] localements -> localement, adverbe invariable

\textbf{Der} 3 $b$) Si $f : Z' \to Z$ dans \Dia{} est une \Hot-équivalence, alors c'est une $\DD$-équivalence, \textit{i.e.} le foncteur $f^{*} : \DD(Z) \to \DD(Z')$ induit une équivalence entre les sous-catégories formées des coefficients constants (ou encore un foncteur pleinement fidèle $\DD(Z)_{\text{loc. const.}} \longrightarrow \DD(Z')$).

$b'_1$) Si $f$ comme ci-dessus, alors $f$ est une $\DD$-équivalence 1-complète, \textit{i.e.} le foncteur $f^{*} : \DD(Z) \to \DD(Z')$ induit un foncteur pleinement fidèle
\msp{page 112}
entre les sous-catégories pleines formées des objets localement constants.

$b'_2$) Comme $b'_1$), mais avec $\DD$-équivalence 2-\underline{complète} au lieu de 1-complète, \textit{i.e.} $f^{*}$ induit une équivalence entre les sous-catégories formées des coefficients localement constants.

$b''_{\mathrm{coh}}$) Si $f : Z' \to Z$ dans \Dia{} est une \Hot-équivalence, alors c'est une $\DD$-équivalence cohomologique complète, \textit{i.e.} pour tout $F \in \mathsf{Ob}\ \DD(Z)$ qui est \underline{localement} \underline{constant}, on a
\[ H^{*}_{\DD}(Z,F) \xrightarrow{\ \sim\ } H^{*}_{\DD}(Z',f^{*}F) \quad . \]

$b''_{\mathrm{hom}}$) Si $f : Z' \to Z$ dans \Dia{} est une \Hot-équivalence, c'est une $\DD$-équivalence homologique complète, \textit{i.e.} pour tout $F \in \mathsf{Ob}\ \DD(Z)$ qui est \underline{localement} \underline{constant}, on a\footnote[59]{\textit{cf.} variante $\pm$ spéciale pour \HOT{} pages 124-125.}
\[ H_{*}^{\DD}(Z',f^{*}F) \xrightarrow{\ \sim\ } H_{*}^{\DD}(Z,F) \quad . \]

\printedpage{67}
% head: A. Grothendieck : Dérivateurs, Ch. I, Édition des Universités de Montpellier II et Paris VII
% foot: Version 8 juillet 2001

On a le diagramme d'implications tautologiques pour les variantes de \textbf{Der} 3 :
\[
\begin{tikzcd}
 & b'_2 \arrow[d, Rightarrow] \arrow[drr, Rightarrow] & & & \\
 & b'_1 \arrow[dl, Rightarrow] \arrow[dr, Rightarrow] \arrow[dd, Rightarrow] & & a' \arrow[ddd, Rightarrow] & \\
b''_{\mathrm{coh}} \arrow[dr, Rightarrow] & & b''_{\mathrm{hom}} \arrow[dl, Rightarrow] & & \\
 & b \arrow[drr, Rightarrow] & & & \\
 & & & a & .
\end{tikzcd}
\]
Dans tous les cas à ma connaissance, les dérivateurs satisfont la condition la plus forte \textbf{Der} 3 $b'_2$.
\msp{page 113}

\dgrblock{dgrb69}{Two families of axioms with different antecedents}
\begin{dcomment}{exposition}{p112:Dans tous les cas à ma connaissance, les dérivateurs satisfont la condition la plus forte}{Two families of axioms with different antecedents}
The $a$- and $b$-conditions should be read as assertions about the strength of the coefficient theory, not as successive definitions of the same phrase. Condition $a$ says that a weakly contractible base does not create new morphisms between constant coefficients. Condition $a'$ adds that it does not create new locally constant objects either. The $b$-conditions ask for analogous invariance under arbitrary weak homotopy equivalences, rather than just maps to a point.

Here is a direct check of the diagram. If $f^*$ is an equivalence on local systems, it is fully faithful there. Full faithfulness gives cohomological and homological invariance by testing morphisms from and to all constant objects. Each of these implies invariance for constant coefficients. Applying the latter to $Z\to e$ proves $a$. Applying the equivalence on local systems to this same map proves $a'$, because every object at $e$ is constant. Finally $a'$ contains $a$ by definition. These are the arrows drawn in the source; the diagram does not claim that full faithfulness alone gives essential surjectivity.

There is a small but consequential inconsistency in the parenthesis defining $b$ on page~111. Its first wording concerns constant coefficients, but the parenthesis writes $\D(Z)_{\mathrm{loc.const.}}$. Taken literally this is $b'_1$, the stronger condition introduced on the next line. For the displayed implication diagram and the earlier definition of a $\D$-equivalence, the subscript in that parenthesis must be read as ``constant''. We retain the edition's wording and record the distinction in the audit.

Modern homotopy invariance supplies $b'_2$ for ordinary Grothendieck derivators with all small categories as shapes. A restricted diagram domain still requires care: the statement must only quantify over shapes in that domain, and a proof that passes through simplicial resolutions or larger indexing categories cannot be imported without an extension hypothesis. This is the reason the manuscript continues to list these properties as separate assumptions. Its present calculation needs the explicitly stated hypothesis, not an unannounced change of domain.
\end{dcomment}


Je rappelle aussi la :

\textbf{Définition 1 :} Soit $f : X \to Y$ dans \Dia. On dit que $f$ est une $\DD$-\textit{équivalence cohomologique} (sous-entendu : “\textit{complète}”) si pour tout $\DD$-coefficient $F$ \underline{localement} \underline{constant} sur $Y$,
\[ H^{*}_{\DD}(Y,F) \xrightarrow{\ \sim\ } H^{*}_{\DD}(X,f^{*}F) \quad ; \]
on dit que $f$ est une $\DD$-\textit{équivalence homologique} (“\textit{complète}”) si pour tout $F$ comme \texttt{[ci-]}dessus,
\[ H_{*}^{\DD}(X,f^{*}F) \xrightarrow{\ \sim\ } H_{*}^{\DD}(Y,F) \quad . \]
On dit que $f$ est une $\DD$-\textit{équivalence} 1-\textit{complète} (resp. 2-\textit{complète}) si le foncteur
\[ f^{*} : \DD(Y)_{\text{loc. const.}} \longrightarrow \DD(X)_{\text{loc. const.}} \]
est pleinement fidèle (resp. une équivalence).
\begingroup\small
\[
\begin{tikzcd}[column sep=small]
 & \text{$\DD$-équivalence 2-complète} \arrow[d, Rightarrow] & & \\
 & \text{$\DD$-équivalence 1-complète} \arrow[dl, Rightarrow] \arrow[dr, Rightarrow] & & \\
\text{$\DD$-équivalence cohomologique} \arrow[dr, Rightarrow] & & \text{$\DD$-équivalence homologique} \arrow[dl, Rightarrow] & \\
 & \text{$\DD$-équivalence} & & .
\end{tikzcd}
\]
\endgroup
Si dans les sous-catégories $\DD(X)_{\text{loc. const.}}$ il y a une loi $*$ (associative, unitaire, commutative) et des \underline{\textsf{Hom}} internes associés, préservés par les images inverses, alors les $\DD$-équivalences cohomologiques sont déjà 1-complètes et sont donc des $\DD$-équivalences homologiques, donc dans ce cas particulier
\printedpage{68}
% head: A. Grothendieck : Dérivateurs, Ch. I, Édition des Universités de Montpellier II et Paris VII
% foot: Version 8 juillet 2001
\msp{page 114 vide}
\msp{page 115}
\[
\begin{tikzcd}
\text{$\DD$-équivalence 2-complète} \arrow[d, Rightarrow] & & \\
\text{$\DD$-équivalence 1-complète} \arrow[r, Leftrightarrow] & \text{$\DD$-équivalence cohomologique} \arrow[d, Rightarrow] & \\
 & \text{$\DD$-équivalence homologique} \arrow[d, Rightarrow] & \\
 & \text{$\DD$-équivalence} & .
\end{tikzcd}
\]
Dans $\DD = \mathsf{HOT}$, toutes ces notions devraient être équivalentes, et équivalentes à la condition encore plus faible en apparence, que $f : X \to Y$ soit une \Hot-équivalence. Mais je ne connais que \HOT{} comme dérivateur où la $\DD$-équivalence d'une flèche $f : X \to Y$ dans \Dia{} implique déjà la $\DD$-équivalence cohomologique (complète), ou la $\DD$-équivalence homologique (complète), voire même la $\DD$-équivalence 1-complète. (Je reviendrai là-dessus tantôt.)

\textbf{Proposition 1.} \textit{Soit $f : X \to Y$ dans} \Dia. \textit{Conditions équivalentes :}

a) \textit{Pour tout $y \in \mathsf{Ob}\, Y$, l'inclusion $X_y \to X/y$ est une $\DD$-équivalence, et pour toute flèche $y' \to y$ dans $Y$, $X/y' \to X/y$ est une $\DD$-équivalence.}

b) \textit{Pour tout coefficient \underline{constant} $F$ sur $Y$, $f_{*}f^{*}F$ est localement constant, et se calcule “fibre par fibre”.}

(C'est la prop. 1, p. 76).

\textbf{Proposition 1$'$.} Dual.

\msp{page 116}

\textbf{Définition 2.} Sous les conditions équivalentes de la proposition 1, on dit que $f$ est une $\DD$-\textit{fibration cohomologique faible.} Sous les conditions duales, que $f$ est une $\DD$-\textit{fibration homologique faible.}

C'est ce que dans la définition 1, p. 77 j'appelais “$\DD$-fibration cohomologique” ou “$\DD$-fibration à gauche”, et dualement. Maintenant la notion me paraît trop faible pour mériter ce nom, que je réserve pour le cas “complet”.

\textbf{Proposition 2.} \textit{Les $\DD$-fibrations cohomologiques faibles sont stables par changements de base lisses, les $\DD$-fibrations homologiques faibles stables par changements de base propres.}

Je devrais construire un exemple où $f$ est une $\DD$-fibration cohomologique faible, et ne l'est plus après changement de base.

\dgrblock{dgrb70}{A failure already among finite posets}
\begin{dcomment}{development}{p116:Je devrais construire un exemple où}{A failure already among finite posets}
The example at manuscript page~77 uses isomorphisms in the base to expose the failure of arbitrary base change. The same failure occurs for finite posets; it is therefore relevant even to the restricted domains repeatedly considered in these pages.

\textbf{A four-object example.}\ours\quad
Let $Y=[2]=\{0<1<2\}$ and let $X$ be the poset whose distinct nonidentity relations are
\[
 x<a,\qquad b<a,\qquad b<z.
\]
Define $f:X\to Y$ by $f(x)=0$, $f(a)=f(b)=1$, and $f(z)=2$. Then $f$ is both a weak cohomological and a weak homological $\HOT$-fibration. Its pullback along the long edge $[1]\to[2]$, selecting $0$ and $2$, is neither.

To prove the first assertion, the strict fibres are $\{x\}$, $\{b<a\}$, and $\{z\}$, all contractible. Since $Y$ is a poset, the overcategories $X/y$ are the subposets on the elements mapping to $\{0,\ldots,y\}$. They are respectively
\[
 \{x\},\qquad \{x\to a\leftarrow b\},\qquad X.
\]
Their nerves realize to a point, an interval subdivided into two edges, and an interval subdivided into three edges. They are all contractible. Every fibre-to-overcategory inclusion and every induced transition between these categories is consequently a weak homotopy equivalence. This is exactly the criterion for a weak cohomological fibration. Dually, the undercategories $y\backslash X$ are $X$, the poset $\{a\leftarrow b\to z\}$, and $\{z\}$; the same argument proves the weak homological condition.

In the pullback, only $x$ and $z$ remain. There is no morphism $x\to z$ in $X$: the connecting path in its nerve traverses $b\to a$ backwards. Thus the pullback category $X'$ is the discrete two-point category, mapping to the two endpoints of $[1]$. The overcategory at the terminal endpoint is all of $X'$, whereas that strict fibre is a point. Its inclusion is not an equivalence on components. The dual failure occurs at the initial endpoint. Even ordinary degree-zero cohomology detects the defect: for a nonzero coefficient field $k$, the relevant comparison is $k\times k\to k$, not an isomorphism.

This settles the request for an example on page~116 without any infinite category or nontrivial automorphism. The original weak conditions hold for every coefficient derivator in which weakly contractible shapes are completely aspherical; the failure is visible for both spaces and ordinary derived coefficients. The construction is supplied here with no priority claim. We checked the local-system and smoothness treatments cited in this commentary and did not locate this precise four-object example. In particular, the ``weakly smooth'' functors of \citet[Definition~4.1 and Proposition~4.3]{Maltsiniotis2005Asphericite} have a different, comma-asphericity condition; their terminology must not be substituted for the weak cohomological condition used here.
\end{dcomment}


% ---- from Der-ChI-234_p071-075.tex
\printedpage{69}
% head: A. Grothendieck : Dérivateurs, Ch. I, Édition des Universités de Montpellier II et Paris VII

\textbf{Proposition 3.} \textit{Soit}
\[
\begin{tikzcd}
X \arrow[rr, "f"] \arrow[dr, "p"'] & & X' \arrow[dl, "p'"] \\
& S &
\end{tikzcd}
\]
\textit{dans \Dia, avec $p$ et $p'$ des $\D$-fibrations cohomologiques faibles} (resp. \textit{des $\D$-fibrations homologiques faibles}), \textit{soit $S_0 \subset S$ tel que $\pi_0(S_0) \to \pi_0(S)$ soit surjective. Pour que $f$ soit une $\D$-équivalence, il suffit que les $f_s : X_s \to X'_s$, pour $s \in S_0$, le soient.}

\textbf{Proposition 4.} \textit{Si $f : X \to Y$ dans \Dia{} est $\D$-propre} (resp. $\D$-\textit{lisse}), \textit{alors $f$ est une $\D$-fibration cohomologique} (resp. \textit{homologique}) \textit{faible si et seulement si pour tout $F$ constant sur $Y$, $f_*f^*F$} (resp. $f_!f^*F$) \textit{est localement constant sur $Y$. Et s'il en est ainsi, pour tout $Y' \to Y$ dans \Dia, $f' : X' = X \times_Y Y' \to Y'$ est encore une $\D$-fibration cohomologique} (resp. \textit{homologique}) \textit{faible.}

Immédiat sur le critère \textit{b}) de la proposition 1.

\msp{page 117}\textbf{Proposition 5.} \textit{Soit $f : X \to Y$ dans \Dia. Considérons les conditions suivantes :}

a) \textit{Pour tout $Y'' \to Y'$ sur $Y$ dans \Dia, si $Y''$ et $Y'$ sont \Hot-asphériques, alors}
\[
X'' = X \times_Y Y'' \longrightarrow X' = X \times_Y Y'
\]
\textit{est une $\D$-équivalence.}
\[
\begin{tikzcd}
X & X' \arrow[l] & X'' \arrow[l] \\
Y \arrow[from=u] & Y' \arrow[l] \arrow[from=u] & Y'' \arrow[l, "p"] \arrow[from=u]
\end{tikzcd}
\]

a$'$) \textit{Comme} a), \textit{mais en supposant $p : Y'' \to Y'$ une \Hot-équivalence (au lieu de supposer $Y'$, $Y''$ \Hot-asphériques).}

a$_1$) \textit{Comme} a), \textit{mais en supposant (au lieu de l'asphéricité) que $Y'$ a un objet final, et $Y''$ est la catégorie unité.}

a$_2$) \textit{Comme} a), \textit{mais en supposant que $Y'$ a un objet cofinal, et $Y''$ est la catégorie unité.}

b$_1$) \textit{$f$ est une $\D$-fibration cohomologique faible, et le reste après tout changement de base.}

b$_2$) \textit{$f$ est une $\D$-fibration homologique faible, et le reste après tout changement de base.}

\printedpage{70}
% head: A. Grothendieck : Dérivateurs, Ch. I, Édition des Universités de Montpellier II et Paris VII
\textit{On a alors le diagramme d'implications tautologiques}
\[
\begin{tikzcd}[column sep=small]
& & & a' \arrow[d, Rightarrow] & & & & \\
& & & a \arrow[dlll, Rightarrow] \arrow[drrr, Rightarrow] & & & & \\
a_1 & \Longleftrightarrow & b_1 & & b_2 & \Longleftrightarrow & a_2 & .
\end{tikzcd}
\]
\textit{Lorsque $\D$ satisfait à l'axiome} \textbf{Der} 3 a$'$ (p. 111), \textit{alors toutes ces conditions, sauf} a$'$), \textit{sont équivalentes. Lorsque $\D$ satisfait à l'axiome} \textbf{Der} 3 b$'_1$ \textit{(ou seulement} \textbf{Der} 3 b$''_{\mathrm{coh}}$ ou b$''_{\mathrm{hom}}$\textit{),} \textit{alors toutes les conditions (y inclus} a$'$\textit{) sont équivalentes.}

Pour simplifier, on va par la suite supposer que l'axiome \textbf{Der} 3 $a'$) est satisfait.

\textbf{Définition 3.} Sous les conditions équivalentes $a$, $a_1$, $b_1$, $a_2$, $b_2$ de la proposition 4, on dit que $f$ est une $\D$-\textit{fibration}. % [sic] proposition 4 -> proposition 5, the conditions a, a_1, b_1, a_2, b_2 are those of proposition 5

\dgrblock{dgrb71}{Why universality identifies the two weak conditions}
\begin{dcomment}{exposition}{p117:Sous les conditions équivalentes $a$, $a_1$, $b_1$, $a_2$, $b_2$ de la proposition 4}{Why universality identifies the two weak conditions}
Proposition~5 organizes the earlier failed argument by changing its input. Under $a'$, one may pull back between weakly contractible bases. The preceding example shows why this is an additional requirement; it is not implied by the fibre criterion on the original base.

Suppose first that the weak cohomological condition holds after every base change. Over a category $B$ with a terminal object, the direct image of a constant coefficient is locally constant and hence constant by $a'$. Its cohomology over $B$ is therefore its value at any selected object $b$. The pointwise formula identifies that value with the cohomology of the strict fibre. Hence every point inclusion into such a base produces a $\D$-equivalence on total categories. Conversely, these point tests applied over $Y/y$, and then at its objects $u:y\to y'$, give the fibre and transport comparisons, exactly as in the proof of Proposition~4 above, now with constant coefficients only. This proves $a_1\Leftrightarrow b_1$.

More generally, take a pullback $p:B'\to B$ with both bases weakly contractible. Put $G=f_*M_X$ and $G'=f'_*M_{X'}$. They are locally constant. Their pointwise comparison is an isomorphism because both universal weak fibre formulas use the same strict fibre; thus $G'\simeq p^*G$. Since $G$ is constant by $a'$, the comparison of its cohomology over the two contractible bases is invertible. This proves $b_1\Rightarrow a$. The dual proof gives $b_2\Rightarrow a$, and the remaining implications are immediate. In $a_2$, the printed ``objet cofinal'' must carry the dual, initial-object meaning required by this argument.

If the stronger $b''_{\mathrm{coh}}$ is assumed, the same calculation works with an arbitrary $\Hot$-equivalence $p$: $G$ need not become constant, because $p$ preserves cohomology of every locally constant coefficient. That proves $a'$. The homological version is dual. Thus the two different uses of the prime on $a'$ must be kept apart: \textbf{Der}~3~$a'$ is an axiom about coefficients on contractible categories, whereas condition~$a'$ of Proposition~5 is a universal pullback test using all weak homotopy equivalences.
\end{dcomment}


\msp{page 118}\textbf{Proposition 6.} \textit{Supposons que $f : X \to Y$ dans \Dia{} soit $\D$-propre} (resp. $\D$-\textit{lisse}). \textit{Conditions équivalentes :}

a) \textit{$f$ est une $\D$-fibration.}

b) \textit{$f$ est une $\D$-fibration cohomologique} (resp. \textit{homologique}) \textit{faible.}

c) \textit{Pour tout $y' \to y$ dans $Y$, $X/y' \to X/y$} (resp. $y\backslash X \to y'\backslash X$) \textit{est une $\D$-équivalence.}

d) \textit{Pour tout $F$ constant sur $Y$, $f_*f^*F$} (resp. $f_!f^*F$) \textit{est localement constant sur $Y$.}

d$'$) \textit{Pour tout $F$ localement constant sur $Y$, $f_*f^*F$} (resp. $f_!f^*F$) \textit{est localement constant sur $Y$.}

Cela résulte formellement de Prop. 4 et Prop. 5.

\textbf{Corollaire.} \textit{Supposons $f : X \to Y$} \Cat\textit{-fibrant} (resp. \Cat\textit{-cofibrant}). \textit{Pour que $f$ soit une $\D$-fibration, il faut et il suffit que pour tout $u : y \to y'$ dans $Y$, $u^* : X_{y'} \to X_y$} (resp. $u_* : X_y \to X_{y'}$) \textit{soit une $\D$-équivalence.}

Il est immédiat que la notion de $\D$-fibration est stable par changement de base, mais il est douteux qu'elle le soit par composition.\footnote[60]{N.B. Si les $\D$-équivalences sont déjà des \Hot-équivalences, alors la notion coïncide avec celle de \Hot-fibration, et elle est stable par composition. Également, si les $\D$-équivalences sont au sens des $\D$-équivalences cohomologiques, ou homologiques, alors la notion coïncide avec celle des $\D$-fibrations cohomologiques (resp. homologiques) ci-dessous, et est également stable par composition.} (Ici encore, j'ai omis de faire un contre-exemple, avec des dérivateurs $\D = \sfup{Der}(k\text{-}\sfup{Mod})$ disons.) Autre lacune de la notion : il n'est pas clair, si $f$ est une $\D$-fibration et $F$ un $\D$-coefficient sur $Y$, si
\[
f_*f^*F \qquad \text{et} \qquad f_!f^*F
\]
\printedpage{71}
% head: A. Grothendieck : Dérivateurs, Ch. I, Édition des Universités de Montpellier II et Paris VII
se calculent “fibre par fibre”, comme on s'attendrait dans le cas d'un “fibré”.

Par contre, voici un résultat sympa :

\dgrblock{dgrb72}{Two fibrations whose composite is not a fibration}
\begin{dcomment}{development}{p118:Par contre, voici un résultat sympa :}{Two fibrations whose composite is not a fibration}
The doubt about composition can be settled by a small categorical example. It uses the same loss of monodromy information as the preceding counterexamples, but places that loss in a transition functor. We work with $\D=\mathsf{Der}(\mathbb Q\text{-}\mathsf{Mod})$ on all small categories and with the final definition of $\D$-fibration on page 117, which includes universality under base change.

\textbf{Proposition (failure of composition).}\ours
There are $\D$-fibrations $p:X\to Y$ and $q:Y\to[1]$ such that $qp$ is not even a weak cohomological $\D$-fibration. The first map can be chosen to be a covering functor, hence both proper and smooth; the second is a Grothendieck opfibration.

Let $G=C_2=\{1,\sigma\}$. Define $Y$ to have two objects $a,b$, with
\[
 Y(a,a)=\{1_a\},\quad Y(b,b)=G,\quad
 Y(a,b)=G,\quad Y(b,a)=\varnothing.
\]
Composition on $Y(b,b)$ is multiplication, and postcomposition of the arrow $g:a\to b$ by $h\in G$ is the arrow $hg$. Send $a$ to $0$ and $b$ to $1$. This defines $q:Y\to[1]$. It is the Grothendieck construction of the functor $e\to BG$. More explicitly, any arrow $g:a\to b$ is opcartesian: for every arrow $k:a\to b$, the unique factorization through $g$ is multiplication by $kg^{-1}$. Its endpoint fibres are $e$ and $BG$, and its transition functor is $e\to BG$.

This transition is a $\D$-equivalence. For a complex $M$ of rational vector spaces with trivial $G$-action, averaging by $(1+\sigma)/2$ makes invariants exact, and coinvariants are exact for the same reason. Thus
\[
 \mathbb R\Gamma(G,M)=M,
 \qquad \mathbb Q\mathbin{\otimes^{\mathbb L}_{\mathbb Q[G]}}M=M.
\]
These identities are natural for arbitrary, including unbounded, complexes: the trivial module is projective over $\mathbb Q[G]$, and the averaging contraction is degreewise compatible with the differential. The opfibration criterion on page 118 now makes $q$ a $\D$-fibration, with all base changes covered by that criterion.

Let $Z$ be the discrete category on the set $G$. Define a functor $L:Y\to\Cat$ by $L(a)=L(b)=Z$, and let every arrow labeled $g$ act by left translation $z\mapsto gz$. Put $X=\int_Y L$ and let $p$ be the projection. All its transport functors are isomorphisms. Every arrow of $Y$ therefore has a unique lift with either prescribed source or prescribed target. In particular $p$ is both a discrete fibration and a discrete opfibration, and the transport criterion makes it a $\D$-fibration (indeed a $\HOT$-fibration).

The composite $qp$ is an opfibration whose fibre at $0$ is the discrete two-point category $Z$. Its fibre at $1$ is the action groupoid $G\ltimes G$ for the left regular action. There is exactly one arrow between any ordered pair of objects of this groupoid, so it is equivalent to the terminal category. The transition functor is the inclusion
\[
 Z_{\mathrm{disc}}\longrightarrow G\ltimes G.
\]
On cohomology in degree zero it induces
\[
 \mathbb Q\longrightarrow\mathbb Q^2,
 \qquad r\longmapsto(r,r),
\]
which is not an isomorphism. Since $qp$ is proper, the criterion of Proposition 6 identifies this failure with failure to be a weak cohomological fibration. Equivalently, $(qp)_*\mathbb Q_X$ has values $\mathbb Q^2$ at $0$ and $\mathbb Q$ at $1$, and its transition is the displayed diagonal. This completes the construction and the proof.\hfill$\square$

The example explains why the positive answer concerning $f_*f^*F$ does not force closure under composition. The coefficient $p_*\mathbb Q_X$ on $Y$ is the permutation local system attached to the regular $G$-set. It is not pulled back from a coefficient on $[1]$: every such pullback has trivial $G$-action at $b$. The constant-coefficient test for $q$ cannot control this local system. The stronger cohomological fibrations introduced on page 121 impose exactly the missing condition. Our literature check covered the smooth-functor formalism of \citet{Maltsiniotis2005Asphericite} and the locally constant coefficient theory of \citet{Cisinski2009LocallyConstant}; we did not find this particular composition counterexample there. It is an elementary development of the mechanism already present in Grothendieck's finite-group example.
\end{dcomment}


\msp{page 119}\textbf{Proposition 7.} \textit{Soit $f : X \to Y$ dans \Dia. Conditions équivalentes :}

a) \textit{$f$ est une $\D$-équivalence “universelle”,} i.e. \textit{pour tout $Y' \to Y$ dans \Dia, $f' : X' = X \times_Y Y' \to Y'$ est une $\D$-équivalence.}

b) \textit{$f$ est une $\D$-fibration à fibres $\D$-asphériques.}

\textbf{Définition 4.} Sous les conditions équivalentes de la proposition 7, on dit que $f$ est une $\D$-fibration \textit{triviale}.

Mais attention, $f$ peut-être une $\D$-fibration et une $\D$-équivalence, sans que la $\D$-fibration soit triviale \textit{i.e.} sans que les fibres soient $\D$-asphériques, cf. p. 87 et suivantes. % [sic] peut-être -> peut être, verb "pouvoir" followed by infinitive
Et également (devrait remonter) même si $\D$ satisfait \textbf{Der} 3 $b_2$ (la plus forte des variantes de \textbf{Der} 3), donc quand \underline{toutes} les conditions de la proposition 5, y inclus $a'$), sont équivalentes, celles-ci (\textit{i.e.} que $f$ soit une $\D$-fibration) n'impliquent pourtant pas, en général, que dans la condition $a'$) on puisse se borner à $p : Y'' \to Y'$ une $\D$-équivalence (au lieu d'une \Hot-équivalence). En d'autres termes, si on a un diagramme cartésien dans \Dia
\[
\begin{tikzcd}
X \arrow[d, "f"'] & X' \arrow[l, "q"'] \arrow[d, "f'"] \\
Y & Y' \arrow[l, "p"] & ,
\end{tikzcd}
\]
avec $f$ une $\D$-fibration, et $p$ une $\D$-équivalence, il ne s'ensuit pas forcément que $q$ soit une $\D$-équivalence. (Mais ce sera pourtant le cas, “bien sûr”, si $\D = \HOT$.)

\msp{page 120}Par contre, on a :

\textbf{Proposition 8.} \textit{Soit $f : X \to Y$ dans \Dia. Conditions équivalentes :}

a) \textit{$f$ est une $\D$-fibration.}

b) \textit{Pour tout $p : Y'' \to Y'$ sur $Y$, si $p$ est une $\D$-équivalence} 1-\textit{complète} (déf. 1, p. 113), \textit{alors $q : X'' \to X'$ est une $\D$-équivalence.}
\[
\begin{tikzcd}
X \arrow[d, "f"'] & X' \arrow[l] \arrow[d, "f'"'] & X'' \arrow[l, "q"'] \arrow[d, "f''"] \\
Y & Y' \arrow[l] & Y'' \arrow[l, "p"]
\end{tikzcd}
\]
\textbf{Corollaire.} \textit{Supposons que dans \Dia, toute $\D$-équivalence soit} 1-\textit{complète.} (N.B. Je ne connais que le dérivateur \HOT{} pour satisfaire à cette condition, qui paraît quasiment caractéristique de \HOT.) \textit{Alors $f : X \to Y$ est une $\D$-fibration si et seulement si pour tout $p : Y'' \to Y'$ sur $Y$, si $p$ est une $\D$-équivalence, alors $q : X'' \to X'$ aussi. Dans ce cas exceptionnel parmi les dérivateurs, la notion de $\D$-fibration est stable par composition} (ce qui est tautologique sur le critère que je viens de donner.)

\dgrblock{dgrb73}{Universal equivalence and the 1-complete test}
\begin{dcomment}{exposition}{p120:ce qui est tautologique sur le critère que je viens de donner.}{Universal equivalence and the 1-complete test}
Two different universal assertions occur here. A universal $\D$-equivalence is a functor whose every base change is a $\D$-equivalence. A $\D$-fibration preserves certain equivalences in diagrams over its base; it need not itself be a $\D$-equivalence. Proposition~7 identifies their intersection with the condition that all strict fibres be $\D$-aspherical. For the forward implication, base change to each point gives aspherical fibres, while the two-out-of-three property applied to the two universally equivalent projections gives the fibration criterion. For the reverse implication, the universal weak fibre formula makes $f_*M_X$ a constant object $M_Y$: its unit is an isomorphism at every $y$ because $X_y$ is $\D$-aspherical. Taking global cohomology proves the equivalence, and the same argument applies after base change.

This result does not permit replacement of ``universal'' by ``ordinary''. The rational monodromy examples already show that a fibration which is a $\D$-equivalence can have a fibre that is not $\D$-aspherical. The issue is not a failure of the fibre formula, but a loss of information upon taking global cohomology of a nonconstant local system.

Proposition~8 repairs the pullback assertion by strengthening the equivalence in the base. If $p:Y''\to Y'$ is 1-complete, it preserves cohomology of every local system. For a constant $M$, universality of the weak fibration condition makes $f'_*M_{X'}$ locally constant and identifies its pullback with $f''_*M_{X''}$. Taking cohomology proves that $X''\to X'$ is a $\D$-equivalence. Conversely, apply the preservation assertion to $e\to B$ with $B$ weakly contractible. Under $a'$, restriction on local systems is an equivalence and hence 1-complete, so the point tests of Proposition~5 follow.

There is a useful distinction between the result here and the stronger definition below. Proposition~8 asks that pulling back a 1-complete equivalence produce a constant-coefficient equivalence. It does not assert that the pullback is itself 1-complete. Demanding that stronger conclusion supplies a composition-stable class by definition, but it is an extra condition unless the internal-Hom argument identifies the relevant equivalence classes.
\end{dcomment}


\printedpage{72}
% head: A. Grothendieck : Dérivateurs, Ch. I, Édition des Universités de Montpellier II et Paris VII

On va renforcer la notion de $\D$-fibration, pour la rendre plus souple, et notamment stable par composition.

\textbf{Proposition 9.} \textit{Soit $f : X \to Y$ dans \Dia. Conditions équivalentes :}

a) \textit{Pour tout $y \in Y$, $X_y \to X/y$ est une $\D$-\underline{équivalence} \underline{cohomologique}} (déf. 1, p. 113), \textit{et de même pour $X/y \to X/y'$ pour toute flèche $y \to y'$ dans $Y$. Et ces conditions sur $f$ restent valables après tout changement de base} (N.B. il suffit que la première de ces conditions le reste).

\msp{page 121}b) \textit{pour tout $F \in \Ob\ \D(X)$ \underline{localement} \underline{constant}, $f_*F$ est localement constant sur $Y$ et se calcule “fibre par fibre” -- et itou après tout changement de base.} % Doubt resolved: attached PDF leaf 36 (printed p.72) has an en dash, encoded --.

c) \textit{Pour tout $p : Y'' \to Y'$ sur $Y$, si $p$ est une $\D$-équivalence cohomologique, alors il en est de même de $q : X'' \to X'$.}

c$'$) \textit{Comme} c, \textit{mais en supposant comme hypothèse sur $p$ que $Y''$ soit la catégorie finale, et $Y'$ ait un objet final.}

\textsc{Démonstration}. Dans l'ordre ($\pm$ tautologique)
\[
\mathrm{a} \Longrightarrow \mathrm{b} \Longrightarrow \mathrm{c} \Longrightarrow \mathrm{c}' \Longrightarrow \mathrm{a} \quad .
\]

\textbf{Définition 5.} Sous les conditions équivalentes de la proposition 9, on dit que $f$ est une $\D$-\textit{fibration cohomologique}. Sous les conditions duales, on dit que $f$ est une $\D$-\textit{fibration homologique}. On dit que $f$ est une $\D$-\textit{fibration forte}, si c'est à la fois une $\D$-fibration cohomologique et une $\D$-fibration homologique.

Ces notions renforcent la notion de $\D$-fibration. Elles ont l'avantage de plus d'être \underline{stables} \underline{par} \underline{composition}, non seulement par changement de base.

\dgrblock{dgrb74}{The reason the strong notion composes}
\begin{dcomment}{exposition}{p121:Elles ont l'avantage de plus d'être}{The reason the strong notion composes}
We may now prove the composition assertion without an appeal to intuition about bundles. Let $X\xrightarrow{f}Y\xrightarrow{g}Z$ be cohomological fibrations in the revised sense. For $L\in\mathcal L(X)$, the object $f_*L$ is locally constant on $Y$, and then $g_*f_*L$ is locally constant on $Z$. The canonical composition isomorphism identifies it with $(gf)_*L$. This is the coefficient condition that was missing for the earlier notion.

For the universal base-change assertion, pull back the composite along any $Z'\to Z$, and factor the resulting rectangle into the two cartesian squares for $g$ and $f$. The comparison for $g$, applied to the locally constant object $f_*L$, is invertible. The comparison for $f$, applied to $L$, is invertible. Pasting gives the comparison for $gf$. The same reasoning works for every locally constant coefficient on the new total category, because the hypotheses on $f$ and $g$ were universal. Specializing to a point in $Z$ gives the fibre formula. Alternatively, condition~$c$ of Proposition~9 proves composition by applying two successive pullbacks to a cohomological equivalence. The two arguments express the same closure property at the levels of coefficients and indexing categories.

The dual reasoning applies to homological fibrations, and hence to strong fibrations. It does not prove that every $\D$-fibration composes: a weak fibration controls direct images of constant coefficients, whereas the first direct image in a composite is usually only locally constant. The explicit counterexample at page~118 realizes exactly this obstruction. The questions about $f_*f^*F$ and $f_!f^*F$, for an arbitrary coefficient on the base, are different again; the factorization-category argument at page~82 proves the required fibre formulas under the stated domain hypotheses. There is no conflict, because a locally constant object produced on the intermediate category need not be pulled back from the final base.

For the proposed 1-complete and 2-complete fibrations of page~122, the phrase ``not stable under composition (as far as I know)'' supplies no proof. In the compatible closed monoidal situation, the 1-complete variant reduces to the cohomological one and therefore composes. The 2-complete variant nevertheless fails to compose, even for the represented set-valued derivator: the development ``The 2-complete variant also fails to compose'' constructs two functors between finite categories from a normalized $C_2$-valued cocycle on a triangulated two-sphere. Their composite fails a permitted contractible-base test because restriction gives only trivial $C_2$-actions. This restriction is fully faithful but not essentially surjective. An internal Hom controls morphisms between objects already descended; it does not supply descent of every local system. The counterexample thus settles the 2-complete question without extending its negative conclusion to the 1-complete class.
\end{dcomment}


\msp{page 122}La notion de $\D$-fibration forte a été choisie pour coiffer les deux notions duales de $\D$-fibration cohomologique et homologique (à défaut d'une meilleur notion) tout en restant stable par composition. On pourrait exiger que pour tout % [sic] meilleur -> meilleure, agreement with "notion"
\[
Y'' \longrightarrow Y'
\]
sur $Y$, avec $Y'$, $Y''$ \Hot-asphériques, $q : X'' \to X'$ soit une $\D$-équivalence 1-complète (voire même 2-complète), appeler ça les $\D$-\textit{fibrations} 1-\textit{complètes} (resp. 2-\textit{complètes}). Ça renforce de façon assez élégante la notion de $\D$-fibration forte, mais a le grand tort de ne pas être stable par composition (que je sache!), sauf dans le cas particulier où on a des \underline{\Hom} internes sympas (cf. bas de la page 113), auquel cas la notion de $\D$-équivalence implique déjà celle de $\D$-équivalence 1-complète,
\printedpage{73}
% head: A. Grothendieck : Dérivateurs, Ch. I, Édition des Universités de Montpellier II et Paris VII
\msp{page 123}donc aussi celle de $\D$-équivalence homologique. Donc dans ce cas
\[
\begin{tikzcd}
\text{$\D$-fibration cohomologique} \arrow[r, Leftrightarrow] \arrow[d, Rightarrow] & \text{$\D$-fibration 1-complète} \\
\text{$\D$-fibration homologique} \quad , &
\end{tikzcd}
\]
et il n'y a pas à chercher plus loin que les $\D$-fibrations cohomologiques.

\dgrblock{dgrb75}{The 2-complete variant also fails to compose}
\begin{dcomment}{development}{p122:La notion de $\D$-fibration forte a été choisie}{The 2-complete variant also fails to compose}
The qualification ``que je sache'' is not a proof of failure. Nevertheless, the proposed 2-complete class really need not be closed under composition, even for a represented derivator with all small categories as shapes. We give a categorical counterexample, including the coherence responsible for the failure. The existence of the local-systems derivator in \citet[Theorem~3.13]{Cisinski2009LocallyConstant} does not supply the missing composition property. We checked the local-systems treatment just cited and did not find this particular construction there. The construction below is a verification given here, without a claim of priority.

\textbf{Proposition.}\ours\quad
Let $\D_{\mathrm{set}}(A)=\operatorname{Fun}(A^{\mathrm{op}},\mathsf{Set})$. There are functors
\[
 X\xrightarrow{p}Y\xrightarrow{q}[1]
\]
which are both 2-complete fibrations in the precise sense proposed on manuscript page~122, while $qp$ is not one. The categories $X$ and $Y$ are finite, and the functor $p$ is even a $\HOT$-fibration. Thus the failure persists in the usual domain of all small categories, with cartesian closed coefficients.

Write $\Pi_1(A)$ for the fundamental groupoid of $|NA|$, or equivalently the groupoid obtained by inverting all arrows of $A$. A locally constant set-valued presheaf factors through this groupoid. Consequently
\[
 \mathcal L_{\mathrm{set}}(A)
 \simeq\operatorname{Fun}(\Pi_1(A)^{\mathrm{op}},\mathsf{Set}).
\]
In particular, an equivalence of fundamental groupoids induces a 2-complete equivalence. We will verify this condition for the two factors, and exhibit an object which cannot descend for the composite.

Let $B$ be the poset of nonempty proper subsets of $\{0,1,2,3\}$, ordered by inclusion. Its nerve is the barycentric subdivision of the boundary of a tetrahedron, so $|NB|$ is a two-sphere. Fix the maximal chain
\[
 \sigma=(\{0\}<\{0,1\}<\{0,1,2\}).
\]
Write $C_2=\mathbb Z/2\mathbb Z$ additively. For composable arrows $u:b\to b'$ and $v:b'\to b''$, set $c(u,v)=1$ when they constitute this chain, and set $c(u,v)=0$ otherwise. In particular it is zero if either arrow is an identity. This is a normalized $C_2$-valued 2-cocycle: there is no chain of three nonidentity arrows in $B$, and the cocycle identity with an identity arrow cancels term by term.

Define $E_c$ to have the same objects as $B$, with one arrow $(u,n)$ for each arrow $u$ of $B$ and label $n\in C_2$. Put
\[
 (v,n)\circ(u,m)=(vu,m+n+c(u,v)),
 \qquad 1_b=(1_b,0).
\]
The cocycle identity is precisely associativity. The projection $E_c\to B$ is both a Grothendieck fibration and an opfibration. Indeed, after fixing a lift of one base arrow, any required factorization is obtained by uniquely solving for the $C_2$-label of the other lift. Its fibre at each object is the one-object groupoid $BC_2$, and its transport functors act identically on $C_2$. The coherence is nontrivial, although the individual transport functors are identities.

We need the fundamental group of $E_c$, not the full homotopy type of its nerve. Put $t_b=(1_b,1)$ and $h_{bb'}=(b\leq b',0)$. Its groupoid completion has these generators and the relations
\[
 t_b^2=1,\qquad t_{b'}h_{bb'}=h_{bb'}t_b,
 \qquad
 h_{b'b''}h_{bb'}=h_{bb''}t_b^{c(b,b',b'')}.
\]
Choose a spanning tree of the one-skeleton of $NB$ and use it to identify all objects. The first relations identify the $t_b$ with a single central generator $t$, with $t^2=1$. Every triangle other than $\sigma$ gives the ordinary, untwisted edge relation. These relations kill all remaining edge generators, because the complex obtained from the triangulated sphere by removing the interior of $\sigma$ is a disk, with the same one-skeleton. The remaining triangle relation then says $1=t$. This proves that $E_c$ is connected and that $\Pi_1(E_c)$ is equivalent to the terminal groupoid.

Now choose $b_0=\{0\}$ and let
\[
 q:Y=\int_{[1]}(e\xrightarrow{b_0}B)\longrightarrow[1]
\]
be the Grothendieck construction. Concretely, $Y$ consists of $B$ and one new object $a$, with
\[
 Y(a,b)=B(b_0,b),\quad Y(b,a)=\varnothing,
 \quad Y(a,a)=\{1_a\}.
\]
The functor $q$ sends $a$ to $0$ and $B$ to $1$. There is a retraction $r:Y\to B$ sending $a$ to $b_0$. If $i:B\hookrightarrow Y$ is the inclusion, then $ri=1$ and the arrow $a\to b_0$ gives a natural transformation $1_Y\to ir$. Finally define
\[
 X=Y\times_B E_c,\qquad p:X\longrightarrow Y.
\]
The retraction and its natural transformation lift to $X\to E_c$, using the arrow of $C_2$-label zero over $a\to b_0$. Naturality follows from the normalization of $c$, since this arrow maps to $1_{b_0}$ under $r$. Hence $|NX|\simeq|NE_c|$, and $X$ is simply connected. The strict fibre of $qp$ at $0$ is $BC_2$.

We now check the universal quantifiers in the definition. First, $p$ is a pullback of the opfibration $E_c\to B$; all its transports are equivalences of groupoids. The opfibration criterion already established on page~118 therefore makes $p$ a $\HOT$-fibration, after every base change. A functor between weakly contractible bases is a nerve weak equivalence. Pullback along $p$ consequently gives a nerve weak equivalence, and hence a 2-complete equivalence for sets.

There is also a direct verification of this last assertion. For any weakly contractible category $A\to Y$, the pulled-back cocycle on $A$ is a coboundary, since $H^2(|NA|;C_2)=0$. Writing it as $\delta d$ and replacing the label $n$ of an arrow $u$ by $n+d(u)$ identifies the pulled-back category with $A\times BC_2$. Choose $d$ normalized, so that this identification preserves the vertical $C_2$-loops. A comparison induced by a functor between two such bases preserves their vertical generator, and hence induces an isomorphism of their fundamental groups, both equal to $C_2$. Both total categories are connected. Thus their fundamental groupoids are equivalent, as required.

For $q$, we prove a stronger assertion. For every $A\to[1]$, with no contractibility assumption, the projection
\[
 q_A:Y_A=Y\times_{[1]}A\longrightarrow A
\]
induces an equivalence of fundamental groupoids. Its strict fibres are $e$ or $B$, and it has a section selecting $a$ over $0$ and $b_0$ over $1$. The transport functors preserve these selected objects. Within each fibre there is a unique morphism in its fundamental groupoid from the selected object to any given object, because the fibre is connected and simply connected. These fibre paths give a natural isomorphism, in $\Pi_1(Y_A)$, from the section followed by the projection to the identity. For naturality, factor an arrow into its opcartesian lift and a vertical arrow; transport of the chosen path and composition with the vertical arrow give the same path in the target fibre, by uniqueness. The section is therefore a quasi-inverse to $\Pi_1(q_A)$.

If $A'\to A$ is a functor over $[1]$ with both bases weakly contractible, the two projections just considered and the equivalence $\Pi_1(A')\to\Pi_1(A)$ show that $Y_{A'}\to Y_A$ is a 2-complete equivalence. This proves exactly the page-122 condition for $q$.

For the composite, use the permitted test $\{0\}\to[1]$. Both bases are contractible, but the upstairs comparison is
\[
 BC_2=X_0\longrightarrow X.
\]
Since $X$ is connected and simply connected, all its locally constant set-valued presheaves are constant. Restriction to $BC_2$ therefore has only trivial $C_2$-actions in its essential image. A two-element set on which $1\in C_2$ exchanges the two elements is a locally constant coefficient on $BC_2$ which does not extend to $X$. Restriction is not essentially surjective, so the comparison is not 2-complete. This proves the proposition.\hfill$\square$

The failure is precisely in descent of objects. In this final test the functor from sets to $C_2$-sets with trivial action is fully faithful; the obstruction does not concern morphisms between objects already descended. Thus the example is consistent with closure of the 1-complete class under the internal-Hom hypotheses, and explains why that argument cannot settle the 2-complete variant. Theorem~3.13 of Cisinski constructs coherent diagrams of local systems for a fixed indexing category; it does not assert that the restricted pullback tests defining the proposed class preserve essential surjectivity under a composite.
\end{dcomment}


On pourrait songer, il est vrai, à définir (dans le cas général) une notion plus forte, en exigeant que pour tout $p : Y'' \to Y'$ sur $Y$ qui soit une $\D$-équivalence 1-complète, $q : X'' \to X'$ soit aussi une $\D$-équivalence 1-complète. Cette fois-ci, la notion est stable par composition, mais elle me paraît néanmoins inutilisable (en dehors du cas qu'on vient de dire, où elle coïncide avec la notion de $\D$-fibration cohomologique, ou encore de $\D$-fibration forte), car n'étant pas en mesure de décrire les $\D$-équivalences 1-complètes, je suis aussi incapable de donner des exemples non triviaux de $\D$-fibrations “strictement 1-complètes”.

Pour vérifier qu'une $f : X \to Y$ dans \Dia{} est une $\D$-fibration cohomologique, il faut
\msp{page 124}savoir vérifier que des flèches $q : Z' \to Z$ dans \Dia{} sont des $\D$-équivalences cohomologiques, donc que $H^*_{\D}(Z,F) \overset{\sim}{\rightarrow} H^*_{\D}(Z',q^*(F))$ est \texttt{[un]} isomorphisme pour tout $F \in \Ob\ \D(Z)$ qui est localement constant sur $Z$. Or je ne sais vérifier un énoncé de cette nature (en dehors du cas où je sais d'avance que $Z$ est complètement $\D$-asphérique, \textit{i.e.} que les coefficients localement constants sur $Z$ sont en fait constant) \underline{que} si $q$ est déjà une \Hot-équivalence. % [sic] constant) -> constants), agreement with "coefficients"
C'est dire que je ne sais exhiber des $\D$-fibrations cohomologiques, que quand ce sont déjà des \HOT-fibrations cohomologiques. Je veux, dans la suite, donner des critères utilisables pour des \HOT-fibrations. Mais on va axiomatiser ces développements, en gardant un $\D$ général, sur lequel je vais faire une hypothèse si draconienne que pratiquement, je ne connais que le dérivateur \HOT{} qui y satisfait. Il s'agit de $c_1$ ci dessous.

\textbf{Der} 3 $c_1$) Toute $\D$-équivalence est une $\D$-équivalence cohomologique.

$c_2$) Toute $\D$-équivalence est une $\D$-équivalence homologique.

\msp{page 125}$c'$) Toute $\D$-équivalence est 1-complète.

$c''$) Toute $\D$-équivalence est 2-complète.

N.B. On a
\[
\begin{tikzcd}
& c'' \arrow[d, Rightarrow] & \\
& c' \arrow[dl, Rightarrow] \arrow[dr, Rightarrow] & \\
c_1 & & c_2
\end{tikzcd}
\]

% ---- from Der-ChI-234_p076-076.tex
\printedpage{74}
% head: A. Grothendieck : Dérivateurs, Ch. I, Édition des Universités de Montpellier II et Paris VII
% foot: Version 8 juillet 2001
tautologiquement. Je ne connais que \HOT{} pour satisfaire l'ensemble des deux plus faibles conditions $c_1$, $c_2$. Et en fait, sous réserve d'écrire la théorie du dérivateur \HOT, il doit être vrai qu'elle satisfait même la forme la plus forte $c''$. (Et les \HOT-équivalences doivent être simplement les \Hot-équivalences.) (Cf. le “lemme” p.~109.)

La proposition suivante est une tautologie.

\noindent\textbf{Proposition 10.} \textit{Supposons que $\D$ satisfasse la condition \textup{\textbf{Der} 3 $\mathrm{c}_1$ (p.~124),} par exemple $\D = \text{\HOT}$. Alors la notion de $\D$-fibration coïncide avec celle de $\D$-fibration cohomologique. Elle est donc stable par composition. Pour que $f : X \to Y$ soit une $\D$-fibration, il faut et il suffit que pour $p : Y'' \to Y'$ sur $Y$, si $p$ est une $\D$-équivalence, il en est de même de $q : X'' \to X'$.}

\noindent\textbf{Corollaire 1.} \textit{Si $f : X \to Y$ dans \textup{\Dia} est $\D$-propre, alors $f$ est une $\D$-fibration cohomologique si et seulement si pour toute flèche $y' \to y$ dans $Y$, $X/y' \to X/y$ est une $\D$-équivalence \textup{(i.e.} si et seulement si $f$ est une $\D$-fibration cohomologique faible\textup{).}}

Appliquer le critère $c$) de \texttt{[la]} prop.~6.

\msp{page 126}

\noindent\textbf{Corollaire 2.} \textit{Supposons que $\D$ satisfasse \textup{\textbf{Der} 3 $(\mathrm{c}_1, \mathrm{c}_2)$ (par exemple $\D = \text{\HOT}$).} Alors pour une flèche $f : X \to Y$ dans \textup{\Dia}, il revient au même que $f$ soit une $\D$-fibration, une $\D$-fibration cohomologique, ou une $\D$-fibration homologique. Si $f$ est propre \textup{(resp.} lisse\textup{),} il revient au même de dire que $f$ est une $\D$-fibration cohomologique \textup{(resp.} homologique\textup{),} faible.}

\dgrblock{dgrb76}{The closing axioms and the theorem they isolate}
\begin{dcomment}{status}{p126:Alors pour une flèche $f : X \to Y$ dans}{The closing axioms and the theorem they isolate}
The $c$-conditions have an antecedent different from that of the $b$-conditions. They start with a $\D$-equivalence, which may be much weaker than a weak homotopy equivalence of nerves. Thus the modern validity of $b'_2$ does not make the $c$-conditions automatic. Rational coefficients and the map $e\to BC_2$ already separate them: all the homotopy-invariance axioms hold, while invariance for constant coefficients does not imply invariance for the sign local system.

For spaces, the argument at page~110 proves all of $c_1,c_2,c',c''$. Propositions~5 and~9 therefore describe one class, and the criterion says that its members preserve weak homotopy equivalences under arbitrary pullback over their base. This is the intrinsic categorical fibration property being isolated here. It should not be identified with the class of fibrations in every model structure on categories, or with an assertion that the realization of the displayed functor is a literal Serre fibration. The invariant content is the comparison with homotopy fibres and the stability of the appropriate pullback squares.

The transition criterion has a classical source. In the dual form of \citet[\S1, Theorem~B and its corollary]{Quillen1973}, weak equivalences between the comma categories along every arrow identify these comma categories with homotopy fibres of the induced map of classifying spaces. For an opfibration, transport functors between strict fibres provide the corresponding test. Properness then gives the strict-fibre-to-comma comparison required in the manuscript. The resulting criterion is the content of Corollary~1 here; its dual gives Corollary~2. These statements were not awaiting a new homotopy group calculation: they were awaiting a coefficient formalism in which the calculation could be made functorial and self-dual.

The scope of the later theorems is worth stating exactly. The ordinary homotopy derivator exists on all small categories \citep[Theorem~6.11]{Cisinski2003}, and its locally constant objects have the invariance needed in the lemma. Within that setting, the question on pages~107--111 and the final conditional assertions are settled affirmatively. The weaker notions for other coefficient derivators retain the distinctions exposed by the examples in these pages. The manuscript's remark that $\HOT$ is the only example known to him is autobiographical, not a uniqueness theorem: for instance a product of $\HOT$ with another derivator satisfying $b'_2$ still detects nerve weak equivalences and satisfies $c''$.

The last statement follows directly. A map is an equivalence for the product derivator exactly when it is one for both factors. The $\HOT$ factor forces it to be a nerve weak equivalence. The $b'_2$ property of each factor then makes restriction an equivalence on their local systems, and local systems in a product are computed factorwise. This proves $c''$ for the product without identifying the coefficient category with spaces. What the axiom singles out is the strength of detection and transport, not a unique presentation of the theory.
\end{dcomment}


On voit donc que dans $\Cat$, on doit avoir une seule notion raisonnable de \HOT-fibration, qui coïncidera avec celle de \HOT-fibration cohomologique ou homologique. Mais pour que cette notion soit dès à présent disponible, il faudrait que le dérivateur \HOT{} soit déjà construit sur le domaine $\text{\Dia} = \Cat$ tout entier. Or ce n'est pas chose faite encore (et il n'a donc pas pu être question de prouver les propriétés de \HOT{} dont il vient d'être question, disons la plus forte \textbf{Der} 3 $c''$).

Néanmoins, je voudrais développer de façon directe (“élémentaire”) la notion de \Hot-fibration, sans me réferer à la théorie du dérivateur \HOT, et aux développements théoriques précédents. % [sic] réferer -> référer, missing accent on the second e


\dgrblock{dgrb77}{The later letter and the meaning of a fibration}
\begin{dcomment}{context}{p126}{The later letter and the meaning of a fibration}
There is a useful chronological comparison at this stopping point. In his letter to Thomason of 2 April 1991, Section~1, Grothendieck starts with a category equipped with a class $W$ of weak equivalences and defines its $W$-fibrations by requiring pullback to carry $W$-equivalences to $W$-equivalences \citep{Grothendieck1991Letter}. Closure under composition then follows by composing pullback functors. This is the strong universal preservation condition proposed on manuscript page~79, and it becomes the final condition in this chapter when the supplementary assumptions on $\D$ identify the appropriate equivalence classes.

The comparison explains why the word ``fibration'' must always be accompanied by its test class. A $W$-fibration defined by preservation of $W$ has a formal closure property. A $\D$-fibration defined by preservation of weak equivalences of nerves only as $\D$-equivalences has two different test classes, so that formal argument does not apply. The composition example on page~118 makes the difference explicit. In the case $\D=\HOT$, the classes agree and the difficulty disappears.

The earlier motivation in Section~69 of \emph{Pursuing Stacks} was to retain coherent diagram categories and the two Kan extensions as part of the coefficient theory \citep{Grothendieck1983}. The present pages show why that enrichment matters even before stability or triangulated categories enter: the coefficient produced by one direct image can have nontrivial transport, and the next direct image must retain it. The mathematical work of these pages is the progressive adjustment of the admissible coefficients and the admissible base changes until this operation can be iterated. It should be read with those successive adjustments intact.
\end{dcomment}
\begin{otherlanguage}{english}
\dgroverview{I-b}{Chapter I \textperiodcentered\ \textup{I-b} \textperiodcentered\ Overview}
\begin{otherlanguage}{english}
\noindent\textbf{Coefficients, monodromy, and the search for a usable fibration.}

These pages of Chapter~I begin with a correction to an overly strong use of cofinality and end with a programme for a theory of fibrations in the category of small categories. Their continuity lies in a single question: which coefficients must a comparison preserve if that comparison is to remain useful after another operation? Preservation of constant coefficients is enough to define a class of weak equivalences. It is not, in general, enough to carry out base change or to compose fibrations. The passage from constants to locally constant coherent coefficients is therefore the principal mathematical movement of the text.

These are pages~34--74 of the edition, manuscript pages~60--126. The text follows the editors' electronic edition \citep{Grothendieck1990Derivateurs}. Manuscript pages~60--66 also close the first part of the chapter, where their commentary, claims and readings are given. Printed and manuscript numbers are kept distinct throughout. No assertion in the commentary silently corrects a mathematical mistake in the French source: genuine transcription errors are corrected, whereas errors already present in the edition remain visible and are explained where they affect the argument.

We use the manuscript's convention that an object of $\D(X)$ has underlying diagram $X^{\mathrm{op}}\to\D(e)$. Thus the value of a right Kan extension $f_*F$ at $y$ is calculated on $X/y$, and the value of the left Kan extension $f_!F$ is calculated on $y\backslash X$. The symbols $H^*_{\D}$ and $H_*^{\D}$ denote objects in the coefficient category, not merely lists of cohomology or homology groups. This convention is responsible for several apparently reversed statements when one compares the manuscript with modern covariant-diagram treatments; the variance distinction is explained explicitly in \citet[Warning~2.2]{GrothPontoShulman2014}.

\medskip\noindent\textbf{From constant equivalences to base change.}

For $f:X\to Y$, with terminal projections $p$ and $q$, a $\D$-equivalence means that restriction induces bijections
\[
 \operatorname{Hom}_{\D(Y)}(q^*M,q^*N)
 \longrightarrow
 \operatorname{Hom}_{\D(X)}(p^*M,p^*N)
\]
for all $M,N\in\D(e)$. Adjunction and Yoneda identify this with invariance of constant-coefficient homology or cohomology whenever the relevant Kan extensions exist. It does not say that $f^*$ is fully faithful on all of $\D(Y)$. The distinction is already decisive for an endpoint inclusion in the walking arrow.

Theorem~1 on manuscript page~67 gives the stronger cofinality test actually needed: cohomological invariance for every coefficient on $Y$ is equivalent to asphericity of the comma categories $y\backslash X$. The proof uses the counit $f_!p^*M\to q^*M$. One adjunction reads it as a cohomology comparison; pointwise evaluation reads it as constant homology of a comma category. This explains the apparent use of homology inside a cohomological theorem without adding a missing hypothesis to a valid argument.

The proper and smooth base-change results then have a precise local content. Properness requires the categories of lifts of every arrow $f(x)\to y$ starting at $x$ to be $\D$-aspherical. The requirement involves both objects and their morphisms, rather than mere nonemptiness of the set of lifts. A Grothendieck opfibration satisfies it because an opcartesian lift is an initial object of the corresponding lift category. The dual account gives smoothness. The commentary supplies the omitted implications among the endpoint, lifting, and base-change tests in Theorem~2, while preserving the edition's mistyped objects and cross-references as source readings.

This vocabulary retains the geometry of proper and smooth base change, but defines a property of functors relative to a coefficient theory. The later theory of asphericity structures provides an appropriate comparison, rather than an identification with topological proper maps \citep[Theorems~1.2, 4.13, and~5.4]{Maltsiniotis2005Asphericite}. The intervening construction $\Phi(X)$, now called the twisted-arrow category, gives natural comparisons between $X$ and $X^{\mathrm{op}}$. Both endpoint projections are opfibrations with contractible fibres. Their universal equivalence proves that taking opposites preserves the equivalence class detected by $\D$. This is a comparison of homotopy information, not an equivalence of the original categories.

\medskip\noindent\textbf{What the counterexamples distinguish.}

From manuscript page~76 onward, the word ``fibration'' changes meaning as the argument develops. Its first version asks that the actual fibre and its comma-category model have the same constant-coefficient theory and that transport between comma categories preserve that theory. Requiring this after every base change is an additional condition. We exhibit both a two-object example and an example involving only finite posets where the first condition holds but fails after a strict base change.

The more substantial failure concerns the proposed implication that a universal fibration must preserve every $\D$-equivalence of bases. For rational complexes, a base can be cohomologically indistinguishable from a point while still carrying nontrivial local systems. Grothendieck's reflected-circle bundle over the circle makes the point concretely. The degree-one cohomology of its fibres is a sign local system on the base circle; its derived cohomology has differential multiplication by $-2$, hence is zero over $\mathbb Q$. The projection from the Klein bottle is consequently a rational constant-coefficient equivalence although its circle fibres are not rationally aspherical. In a fibre square the tensor square of the orientation local system becomes untwisted, and the extra cohomology reappears. The commentary calculates the full complexes and the comparison maps, including the degree-zero term omitted from the source's abbreviated group-cohomology calculation.

The finite-group version replaces this computation by the averaging idempotent. It also shows why the proposed extension to truncated homotopy theories needs care: already the point mapping into $BC_2$ is an equivalence for the represented set-valued derivator, whereas locally constant $C_2$-sets retain the action. The corresponding assertions about complete equivalences cannot be inferred from connectedness.

One further correction is necessary in the proposed finite-poset model. The action on $S^2$ is antipodal and the action on $S^1$ is a reflection. Their product acts trivially on the top class of degree three. Hence
\[
 H^*((S^2\times S^1)/C_2;\mathbb Q)
 \cong\mathbb Q\quad\text{in degrees }0,3,
\]
and vanishes in other degrees. The base $\mathbb{RP}^2$ is rationally aspherical, but this total space is not. The example still disproves invariance under inclusion of a fibre; it does not supply the additional aspherical-total-space claim. An exact rational chain calculation for the thirteen-object base and fifty-two-object total poset independently confirms these Betti numbers. The orbit-category relations printed before this construction also require correction, which is explained separately in the commentary.

\medskip\noindent\textbf{The final distinctions and their present interpretation.}

The later definitions separate constant-coefficient equivalence, invariance of cohomology with every locally constant coefficient, and full faithfulness or equivalence on the entire category of locally constant objects. The last two are called $1$- and $2$-complete equivalence. A closed internal-Hom formalism identifies cohomological equivalence with $1$-complete equivalence, under the stated compatibility assumptions. It does not identify ordinary constant-coefficient equivalence with either of them. The missing word ``cohomological'' in the sentence on manuscript page~122 is therefore a material slip.

The nerve-invariance assertions eventually formulated as the $b'$ variants of \textbf{Der~3} have a precise later answer. \citet[Theorem~1.3 and Corollary~3.8]{Cisinski2009LocallyConstant} proves invariance of locally constant coefficients under weak equivalences of indexing nerves, first in the model-category setting and then for derivators. In the space-valued case these coefficients can be read as coherent functors out of the fundamental infinity-groupoid. A weak equivalence changes neither that groupoid nor its coefficient categories. This establishes the strong form of the lemma about $\HOT$ that Grothendieck had deferred; it does not collapse the distinctions for a fixed linear coefficient theory. The minimal-localizer theorem places the ordinary weak-equivalence class inside the broader framework that motivated these pages \citep[Theorem~2.2.11]{Cisinski2004Minimal}.

Two developments in the commentary on these pages address unfinished questions directly. First, for a universally weak cohomological fibration, the formula for $f_*f^*F$ holds for every coefficient $F$ on the base, with no local-constancy assumption. A factorization category for each arrow $z\to y$ reduces the left mate of the desired comparison to constant homology. The dual proof applies to $f_!f^*F$. As a consequence, a universal $\D$-equivalence induces fully faithful restriction on the whole coefficient category. The required closure of the diagram domain is stated explicitly. These arguments use the factorization method underlying homotopy exactness \citep[Theorem~3.8]{GrothPontoShulman2014}, and the full-faithfulness conclusion for universal weak equivalences of nerves also follows from the modern homotopical epimorphism criterion \citep[Theorem~6.10]{GrothStovicek2016Trees}.

Second, the intermediate $\D$-fibrations are not closed under composition. Over $\mathbb Q$, take the opfibration over $[1]$ with transition $e\to BC_2$ and then the covering determined by the regular $C_2$-set. Both factors satisfy the universal constant-coefficient condition. Their composite has transition from two discrete points to a contractible action groupoid, and fails already in degree-zero cohomology. The stronger cohomological and homological fibrations on manuscript page~121 control locally constant coefficients on the total category and are closed under composition. A separate example with transition $BC_2\to e$ shows that this stronger relative theory has examples beyond $\HOT$-fibrations. The proposed $2$-complete variant also fails to compose for the represented set-valued derivator: a $C_2$-valued cocycle on the tetrahedral boundary face poset gives two functors between finite categories whose composite fails the required test. The obstructed restriction includes sets as trivial $C_2$-sets; its failure of essential surjectivity distinguishes this counterexample from the $1$-complete case, which composes under the internal-Hom hypotheses.

The mathematical status is thus stated at the level of the actual assertions. The main cofinality and base-change theorems hold with their conventions and domain hypotheses made explicit; the false implications have concrete counterexamples; several stated questions admit the supplied proofs or the cited later theorems. A failure to locate a particular construction in the sources checked is recorded as such, never converted into a claim that it is an open problem or a new result. The label \ours\ marks proofs and constructions given in the commentary; it carries no priority claim.
\end{otherlanguage}
\end{otherlanguage}
\dgrchunkend

% DER_II-a: Grothendieck's text with the commentary set in, then the overview, inside a group that applies the chunk's own definitions
\dgrchapter{II}{Chapter II}{Cofinalité (à droite et à gauche) (préliminaire à la cofinalité cohomologique)}
\dgrpart{0}{II-a}{Part 1 of 2, manuscript pages 1--53}{II-a\enspace\textperiodcentered\enspace Chapter II, part 1 of 2\enspace\textperiodcentered\enspace manuscript pages 1--53}{Chapter II \textperiodcentered\ \textup{II-a}}
\dgrchunkbegin
% DER_II-a: the definitions of the chunk's preamble, applied inside its group
\let\printedpage\relax\newcommand{\printedpage}[1]{\begingroup\def\dcpagevalue{#1}\def\dcblankpage{--}\ifx\dcpagevalue\dcblankpage\else\marginpar{\scriptsize\raggedright p.~#1}\fi\endgroup}
\def\dcmsnumber page #1\relax{#1}
\newif\ifdcskipmsp
\let\dcpairedpage\relax\newcommand{\dcpairedpage}[2]{\marginpar{\scriptsize\raggedright p.~#1\\\mbox{[#2]}}\global\dcskipmsptrue}
\let\msp\relax\newcommand{\msp}[1]{\ifdcskipmsp\global\dcskipmspfalse\else\marginpar{\scriptsize\raggedright[\dcmsnumber#1\relax]}\fi}
% printed notation: the category C' enclosed in a V-shaped wedge with the arrow name under it
\let\Vcat\relax\newcommand{\Vcat}[2]{\underset{#2}{\underset{\vee}{#1}}}
\let\Homint\relax\newcommand{\Homint}{\operatorname{Hom\,int}}
\let\sX\relax\newcommand{\sX}{\mathcal{X}} % resolved against PDF: CMSY calligraphic X
\let\X\relax\newcommand{\X}{\mathcal{X}}
\let\Hom\relax\newcommand{\Hom}{\operatorname{Hom}}
\let\Ker\relax\newcommand{\Ker}{\operatorname{Ker}}
\let\Ob\relax\newcommand{\Ob}{\operatorname{Ob}}
\let\mspin\relax\newcommand{\mspin}[1]{\texttt{[#1]}}
\let\HOT\relax\newcommand{\HOT}{\mathrm{HOT}}
\let\Cat\relax\newcommand{\Cat}{\mathrm{Cat}}
\let\Fl\relax\newcommand{\Fl}{\mathrm{Fl}}
\let\Ens\relax\newcommand{\Ens}{\mathrm{Ens}}
\let\pl\relax\newcommand{\pl}{\texttt{[plutôt}}
\let\uHom\relax\newcommand{\uHom}{\underline{\operatorname{Hom}}}
\let\bs\relax\newcommand{\bs}{\backslash}
\let\cM\relax\newcommand{\cM}{\mathcal{M}}
\let\dbl\relax\newcommand{\dbl}[2]{\mathrel{\overset{#1}{\underset{#2}{\rightrightarrows}}}}
\let\defeq\relax\newcommand{\defeq}{\overset{\text{déf}}{=}}

% ---- from Der-ChII-234_p001-005.tex
% ===================== PDF page 1 =====================
\printedpage{--}
% no page number printed (title page)
% head: (none)

\begin{center}
\vspace*{4cm}
{\Large\textsf{LES DÉRIVATEURS}}

\bigskip\bigskip

\textbf{ALEXANDRE GROTHENDIECK}

\bigskip\bigskip\bigskip

\textbf{Chapitre II}

\textbf{Cofinalité (à droite et à gauche)}

\textbf{(préliminaire à la cofinalité cohomologique)}
\end{center}

\clearpage

% ===================== PDF page 2 =====================
\printedpage{--}
% no page number printed (editorial note)
% head: (none)

\vspace*{6cm}

\noindent Ce texte a été déchiffré et transcrit en \LaTeXe{} par M. Künzer. Il a été édité par M. Künzer, J. Malgoire, et G. Maltsiniotis. La transcription est aussi fidèle que possible au manuscrit. Pour les quelques corrections évidentes, ou rares commentaires des éditeurs, ainsi que pour la numérotation originale des pages du manuscrit, les caractères de machine à écrire \texttt{[typewriter]} entre crochets sont utilisés. Un point d'interrogation entre crochets signifie que l'on n'est pas sûr du mot qui précède.

\noindent Cette édition est provisoire. Les remarques, commentaires et corrections sont bienvenus. Envoyer un message à :

\noindent\textbf{maltsin@math.jussieu.fr}

\hfill G. Maltsiniotis\hspace*{2cm}

\clearpage

% ===================== PDF page 3 =====================
\dcpairedpage{1}{1}
% head: A. Grothendieck : Dérivateurs, Ch. II, Édition des Universités de Montpellier II et Paris VII
% foot: Version 8 juillet 2001

\msp{page 1}

\section*{1\quad Cofinalité cohomologique}

Soit
\[ \alpha : C' \longrightarrow C \]
un foncteur de (petites) catégories, $\mathcal{M}$ une catégorie ayant les $\varprojlim$ pertinentes, $F$ un préfaisceau (donc contravariant en $i$) sur $C$ à valeurs dans $\mathcal{M}$. On a un homomorphisme canonique dans $\mathcal{M}$
\[ \varprojlim_{C} F \longrightarrow \varprojlim_{C'} \underbrace{\alpha^*(F)}_{F\circ\alpha} \tag*{(1)} \]
et on veut des conditions sur $\alpha$ pour que ce soit toujours un isomorphisme (\footnote{On dit alors que $\alpha$ est \textit{cofinal,} cf. SGA 4, I 8.1, ou 3.3, p. 45.}). On est ramené aussitôt au cas où $\mathcal{M} = (\mathrm{Ens})$. Disons qu'une catégorie $I$ est \textit{faiblement filtrante à gauche} si $I \neq \emptyset$ et si $\forall\, i', i'' \in \operatorname{Ob} I$, $\exists\, i \in \operatorname{Ob} I$ et des flèches
\[
\begin{tikzcd}[row sep=small, column sep=small]
 & i' \\
i \arrow[ur] \arrow[dr] & \\
 & i''.
\end{tikzcd}
\]
Cela signifie aussi que $\operatorname{Ob} I$ en tant qu'ensemble ordonné est filtrant à gauche.

\dgrblock{dgrb78}{The presheaf represented by a change of index}
\begin{dcomment}{exposition}{p1:On est ramené aussitôt au cas où}{The presheaf represented by a change of index}
Grothendieck begins again with ordinary limits because their elementary descent argument already contains the object which the cohomological theory must derive. We keep his variance throughout: a presheaf on $C$ is a functor $C^{\mathrm{op}}\to\mathrm{Set}$, and $c/\alpha$, also written $c\backslash C'$, has objects $(a,f:c\to\alpha a)$. The arrow in the first proposition is printed in the opposite direction; its intended direction is fixed both by this notation and by the proof on manuscript page~6. This is an edition slip, not a change of convention.

Put $K_\alpha(c)=\pi_0(c/\alpha)$. Precomposition with $c_0\to c_1$ defines $K_\alpha(c_1)\to K_\alpha(c_0)$, so $K_\alpha$ is a presheaf. There is a natural bijection
\[
 \lim_{C'^{\mathrm{op}}}\alpha^*F
 \;\cong\;\operatorname{Nat}(K_\alpha,F).
\]
Indeed a compatible family $(s_a)$ determines the map $[a,f]\mapsto F(f)(s_a)$. A morphism $(a,f)\to(b,g)$ in $c/\alpha$ gives $g=\alpha(h)f$, and compatibility gives $s_a=F(\alpha h)(s_b)$; hence the value is constant along every zigzag. Naturality in $c$ follows by composition. Conversely, from $t:K_\alpha\to F$, put $s_a=t_{\alpha a}[a,1_{\alpha a}]$. For $h:a\to b$, the arrow $(a,1)\to(b,\alpha h)$ and naturality of $t$ prove compatibility. These constructions are inverse.

Equivalently, $K_\alpha=\alpha_!1_{C'}$. Thus the restriction map on limits is induced by the augmentation $K_\alpha\to1_C$. Yoneda, applied in the presheaf category, proves the exact criterion:
\[
 \text{restriction preserves every presheaf limit}
 \quad\Longleftrightarrow\quad
 \pi_0(c/\alpha)=\{*\}\quad(c\in C).
\]
The nonempty, connected comma categories are what matter. A common predecessor for each pair of objects is a stronger sufficient condition; it has no equalization requirement. For coefficients in a locally small category $\mathcal M$, apply $\operatorname{Hom}_{\mathcal M}(M,-)$ for every $M$. These functors preserve existing limits and jointly reflect isomorphisms, which justifies the reduction to sets without imposing extra completeness on $\mathcal M$.

In current covariant terminology, $\alpha$ is a final functor, while $\alpha^{\mathrm{op}}$ is initial; the latter is the functor acting on the diagrams whose limits are written here. SGA~4 uses the cofinal terminology followed in the manuscript \citep[Expos\'e~I, \S8.1]{SGA4I1972}. Proposition~3.2 below returns to precisely this criterion. We shall use the representing presheaf again there, rather than repeat the elementwise argument.
\end{dcomment}


\textbf{Proposition.} \textit{Pour que} (1) \textit{soit toujours un isomorphisme} ($\alpha$ \textit{`anodin pour} $\varprojlim$ \textit{'), il suffit que pour tout} $x_0 \in \operatorname{Ob} C$\textit{, la catégorie} $x_0\backslash C'$ \textit{des couples} $(x', u)$\textit{, avec} $x' \in \operatorname{Ob} C'$\textit{,} $u : \alpha(x') \longrightarrow x$\textit{, soit faiblement filtrante à gauche} (\footnote{Cf. amélioration des remarques p. 5.}, \footnote{Cf. 3.2, p. 43-44, pour conditions nécessaires et suffisantes de cofinalité.}). % [sic] u : \alpha(x') \longrightarrow x -> x_0 in place of x, the category is x_0\backslash C'

Ceci n'implique pas la conclusion plus forte
\[ F \underbrace{\xrightarrow{\;\sim\;}}_{\text{adj.}} \alpha_*\alpha^* F \]
que le morphisme d'adjonction canonique soit un isomorphisme. Pour qu'il soit ainsi, il suffit que $\forall\, x_0 \in C$, le foncteur $\alpha/x_0 : C'/x_0 \longrightarrow C/x_0$ induit par $\alpha$ soit `anodin pour les $\varprojlim$ ', ce qui signifie aussi que pour toute flèche $u : x_1 \longrightarrow x_0$ dans $C$, la catégorie $\Vcat{C'}{u}$ formée des triples $(x', \xi, \eta)$ avec $x' \in \operatorname{Ob} C'$ et $\xi : x_1 \longrightarrow \alpha(x')$, $\eta : \alpha(x') \longrightarrow x_0$ rendant commutatif
% printed: C' inside a V-shaped wedge, with u printed below the wedge
\msp{page 2}
le diagramme
\printedpage{2}
% head: A. Grothendieck : Dérivateurs, Ch. II, Édition des Universités de Montpellier II et Paris VII
% foot: Version 8 juillet 2001
\[
\begin{tikzcd}
x_1 \arrow[rr, "u"] \arrow[dr, "\xi"'] & & x_0 \\
 & \alpha(x'), \arrow[ur, "\eta"'] &
\end{tikzcd}
\]
est faiblement filtrante à gauche.

Plus généralement, considérons
\[
\begin{tikzcd}
X' = C' \arrow[rr, "\alpha"] \arrow[dr, "f'"'] & & C = X \arrow[dl, "f"] \\
 & S &
\end{tikzcd}
\]
et posons nous la question quand l'homomorphisme fonctoriel, par rapport à un préfaisceau $F$ variable sur $X$,
\[ f_*(F) \longrightarrow f'_*(\alpha^*(F)) \]
est toujours un isomorphisme. (Le cas précédent est celui où $S = X$, $f = \mathrm{id}_X$.) On trouve qu'il suffit que $\forall\, s_0 \in \operatorname{Ob} S$, le foncteur
\[ \alpha/s_0 : X'/s_0 \longrightarrow X/s_0 \]
soit `anodin pour les $\varprojlim$ ', et pour ceci, il suffit que pour \texttt{[tout]} objet $(x_0, u_0)$ de $X/s_0$ (où
\[ f(x_0) \xrightarrow{u_0} s_0) \]
la catégorie $(x_0, u_0)\backslash(X'/s_0)$ soit faiblement filtrante à gauche.

\msp{page 3}

Cette catégorie s'interprète comme celle des triples $(x', \xi, u)$ avec
\[
\left\{
\begin{aligned}
& x' \in \operatorname{Ob} C' \\
& u : g(x') = f(\alpha(x')) \longrightarrow s_0 \\
& \xi : x_0 \longrightarrow \alpha(x')
\end{aligned}
\right.
\]
\texttt{[plutôt} $f'(x')$ \texttt{au lieu de} $g(x')$\texttt{]} de façon à rendre commutatif le diagramme
\[
\begin{tikzcd}
f(x_0) \arrow[rr, "u_0"] \arrow[dr, "f(\xi)"'] & & s_0 \\
 & f(\alpha(x')). \arrow[ur, "u"'] &
\end{tikzcd}
\]
Je voudrais dégager des variantes de ces énoncés dans le cadre cohomologique, les foncteurs $\varprojlim$ étant remplacés par des $R\Gamma$ (de préfaisceaux abéliens p.ex. - mais on pense aussi aux `cas non commutatifs', définis par des dérivateurs notamment.)

\clearpage

% ===================== PDF page 5 =====================
\printedpage{3}
% head: A. Grothendieck : Dérivateurs, Ch. II, Édition des Universités de Montpellier II et Paris VII
% foot: Version 8 juillet 2001

Si $I$ est une catégorie \textit{ordonnée finie}, la condition \texttt{[faiblement]} filtrante à gauche équivaut à l'existence d'un objet initial, ce qui implique que $I$ est contractile. Dans les énoncés précedents, il se pourrait que la bonne philosophie soit de remplacer partout `faiblement filtrante à gauche' par `contractile'. Si p.ex. $C$ est la catégorie finale, alors remplacer la condition que $C'$ soit filtrante à gauche, par l'acyclicité. C'est en tous cas en accord avec l'esprit de l'axiome $\mathrm{Der}_5$ (cas où $J$ est la catégorie finale). En effet, dans le cas où $C = e$, la condition que $\alpha$ soit anodin pour les $\varprojlim$ équivaut visiblement à dire que le foncteur $\alpha^* : \mathcal{M} \longrightarrow \underline{\mathrm{Hom}}(C'^{o}, \mathcal{M})$, associant à tout $F \in \mathcal{M}$ le préfaisceau constant de valeur $F$ % [sic] précedents -> précédents, missing accent on the second e
\msp{page 4}
sur $C'$, donne un morphisme d'adjonction
\[ \varprojlim_{e} F = F \longrightarrow \alpha_*(\alpha^*(F)) = \varprojlim_{C'} \alpha^* F \]
qui est un isomorphisme. Et l'axiome $\mathrm{Der}_5$ assure qu'il doit en être ainsi si $C'$ est contractile.

\dgrblock{dgrb79}{Why the relative problem has two endpoints}
\begin{dcomment}{exposition}{p3}{Why the relative problem has two endpoints}
The passage from global sections to a direct image introduces a second endpoint. For $f:X\to S$ and a presheaf $F$, the pointwise right Kan extension formula is
\[
 (f_*F)(s)=\lim_{(f/s)^{\mathrm{op}}}F(x),
 \qquad (f/s)=\{(x,f(x)\to s)\}.
\]
This is the ordinary comma-category formula for a Kan extension \citep[Theorem~6.2.1 and Corollary~6.2.9]{Riehl2016}. If $f'=f\alpha$, the comparison at $s$ is restriction along $\alpha/s$. Its comma category at $(x_0,u_0)$ consists exactly of factorizations
\[
 f(x_0)\xrightarrow{\,f(\xi)\,}f(\alpha a)
 \xrightarrow{\,u\,}s,
 \qquad u f(\xi)=u_0.
\]
Connectedness and nonemptiness of these categories therefore suffice. The weak cofiltering condition written here also suffices, but the phrase ``ce qui signifie aussi'' on manuscript page~1 should not turn that sufficient condition into a necessary one.

When $S=X$ and $f=1_X$, the test concerns factorizations of each arrow $x_0\to x_1$ through $\alpha a$. It is stronger than the test involving only $x_0/\alpha$. This is the same distinction already isolated in the commentary on Chapter~I, manuscript p.~49 (``The factorization category controls base change''); Proposition~4.1 will identify its ordinary categorical form. Keeping the two endpoints visible prevents an erroneous inference from invariance of global sections to full faithfulness of restriction.
\end{dcomment}


La condition filtrante dans la proposition est d'ailleurs une condition \textit{suffisante} d'anodinité, mais nullement nécessaire. Si p.ex. $C$ est la catégorie finale, alors la condition d'anodinité pour $C' \longrightarrow C$ est simplement que $C'$ soit connexe non vide, i.e. $\pi_0(C') = \{0\}$. (Condition qui visiblement néglige les $\pi_i$ supérieurs.)

Soit \fbox{$\alpha : X' \longrightarrow X$} un foncteur de petites catégories. La condition suggérée par la proposition, que les $x_0\backslash X'$ soient contractiles, signifie que $\alpha$ est une hot-équivalence co-stricte, i.e. que $\alpha^{o} : X'^{o} \longrightarrow X^{o}$ est une hot-équivalence `localement sur $X^{o}$', donc aussi globalement. Si $\alpha$
\msp{page 5}
lui-même était une hot-équivalence stricte, au lieu de $\alpha^{o}$, alors on aurait par $\mathrm{Der}_5$, que pour $F \in \mathbf{D}(X)$, $F \xrightarrow{\;\sim\;} \alpha_*\alpha^* F$ ($\mathbf{D}$ dérivateur à gauche), ce qui est bien plus fort que $\Gamma_{\mathbf{D}}(X, F) \xrightarrow[\sim]{\alpha^*} \Gamma_{\mathbf{D}}(X', \alpha^* F)$. La condition duale implique (si $\mathbf{D}$ \texttt{[est un]} dérivateur à droite) $\alpha_!\alpha^* F \xrightarrow{\;\sim\;} F$. Mais est-il raisonnable de s'attendre qu'elle implique aussi la condition `du côté cohomologie, au lieu de l'homologie',
\[ \Gamma_{\mathbf{D}}(X, F) \xrightarrow{\;\sim\;} \Gamma_{\mathbf{D}}(X', \alpha^*(F)) \ \ ? \]

Le premier cas, qui devrait emporter conviction, est celui du dérivateur DerAb, i.e. des catégories dérivées des faisceaux abéliens :

\textbf{Question.} \textit{Soit} $\alpha : X' \longrightarrow X$ \textit{une} hot-\textit{équivalence co-stricte,} $F$ \textit{un préfaisceau abélien sur} $X$\textit{. A-t-on alors} $H^i(X, F) \xrightarrow[\sim]{\alpha^*} H^i(X', \alpha^*(F))$ \textit{pour tout} $i \in \mathbf{Z}$ \textit{?}

Si $X$ admet un objet initial $\emptyset$, alors $\emptyset\backslash X' \simeq X'$, donc $X'$ contractile, donc c'est OK pour $F$ constant.

\dgrblock{dgrb80}{Cohomological cofinality and the derived weight}
\begin{dcomment}{later}{p5:donc c'est OK pour $F$ constant.}{Cohomological cofinality and the derived weight}
Grothendieck's question about abelian coefficients has an affirmative answer. Its proof can follow the argument he has just discovered, provided that the presheaf $K_\alpha$ is replaced by its chain-level predecessor. We give the argument for a commutative ring $k$; $k=\mathbb Z$ answers the stated question. Write $D(C;k)$ for the derived category of $k$-module presheaves and use $\mathbf R\operatorname{Hom}_{C;k}$ for the \emph{global} derived Hom, valued in $D(k)$.

Derived adjunction gives, naturally in $F$,
\[
 \mathbf R\Gamma(C',\alpha^*F)
 \simeq
 \mathbf R\operatorname{Hom}_{C;k}
       (\mathbf L\alpha_!k_{C'},F).
\]
For clarity, this can be constructed without assuming the disputed internal formula on manuscript page~10. Resolve $k_{C'}$ by the augmented simplicial presheaf of free modules whose degree $n$ is
\[
 P_n=\bigoplus_{a_0\to\cdots\to a_n} k\operatorname{Hom}_{C'}(-,a_0).
\]
At $a$, its augmented chain complex is the nerve complex of $a/C'$, contracted by its initial object $1_a$. Each $P_n$ is projective. Its normalized complex, in nonpositive cohomological degrees, is therefore a projective resolution. Since restriction is exact, $\alpha_!$ takes these projectives to projectives. Evaluating $\alpha_!P_\bullet$ at $c$ gives the free chain complex of $c/\alpha$. In particular
\[
 (\mathbf L\alpha_!k_{C'})(c)
 \simeq C_*(N(c/\alpha);k),
\]
where a chain degree $n$ is placed in cohomological degree $-n$. The usual formulation using the opposite comma category gives the same complex up to signed reversal.

If every $N(c/\alpha)$ is weakly contractible, its augmented complex is quasi-isomorphic to $k$. Consequently $\mathbf L\alpha_!k_{C'}\to k_C$ is an isomorphism. Applying derived Hom yields the desired comparison in every degree, and indeed for every complex $F$. This proof needs only $k$-acyclicity of these nerves; weak contractibility is the uniform condition which works for homotopy types as well. The two conditions must not be confused with mere connectedness. For example, $BG\to e$ is cofinal for ordinary presheaf limits whenever $G$ is a group, but for $G=\mathbb Z$ and the constant coefficient $\mathbb Z$ one has $H^1(BG;\mathbb Z)=\mathbb Z$.

For spaces, the corresponding weight is $c\mapsto B(c/\alpha)$ in $\operatorname{Fun}(C^{\mathrm{op}},\mathcal S)$. The same mapping-object adjunction proves cofinality; the general higher categorical criterion is \citet[Theorem~4.1.3.1, together with Proposition~4.1.1.8]{Lurie2009HTT}, applied with the appropriate opposites. Its corollary~4.1.3.3 is Quillen's Theorem~A. These results settle the question on page~5, without requiring the stronger unit or counit assertions in its preceding paragraph. Those assertions are conditional there on the manuscript's invocation of $\mathrm{Der}_5$; interpreted as consequences of comma-category asphericity under present derivator axioms, they are false. The endpoint calculation in the next block makes the distinction explicit, continuing the caution about the changing axioms in the commentary on Chapter~I, manuscript p.~43 (``What the localization axioms do and do not say'').
\end{dcomment}


\textbf{Remarque.} Je m'aperçois que dans la proposition 1, au lieu de la condition que les $x\backslash C'$ soient faiblement filtrantes à gauche, il suffirait qu'elles soient non vides et connexes, i.e.

% ---- from Der-ChII-234_p006-010.tex
\printedpage{4}
% head: A. Grothendieck : Dérivateurs, Ch. II, Édition des Universités de Montpellier II et Paris VII
% foot: Version 8 juillet 2001
0-connexes (\footnote[4]{Cela signifie aussi que $\alpha$ est $W_0$-coasphérique. Pour un traitement ab ovo, cf. prop. 3.2, page 43.}). Le fait que les $x\backslash C'$ non vide implique que $\varprojlim\limits_{C} \longrightarrow \varprojlim\limits_{C'}$ soit \textit{injectif.} Pour la sujectivité, si $\xi' \in \varprojlim\limits_{C'}$ (cas ensembliste), on définit $\xi \in \varprojlim\limits_{C}$ par les $\xi(x) \in F(x)$ à l'aide de % [sic] sujectivité -> surjectivité, lettre r omise
\msp{page 6}
\[ x \xrightarrow{u} \alpha(x') \]
en prenant $u^*(\xi(x'))$ (où $\xi(x') \in (\alpha^*(F))(x') = F(\alpha(x')))$, le point essentiel est de montrer que les $u^*(\xi(x'))$ ne dépendent que de $x$, pas du choix de $(x',u) \in x\backslash C'$. Or pour $(x',u)$ variable dans $x\backslash C'$, l'application
\[
\begin{aligned}
\varphi : \operatorname{Ob}(x\backslash C') &\quad\longrightarrow\quad F(x)\\
(x',u) &\quad\longmapsto\quad u^*(\xi(x'))
\end{aligned}
\]
est telle que s'il existe une flèche
\[ (x',u) \xrightarrow{u'} (\bar{x}',\bar{u}) \qquad \text{où} \qquad u' : x' \longrightarrow \bar{x}' \]
(donc le diagramme
\[
\begin{tikzcd}
 & \alpha(x') \arrow[dd, "\alpha(u')"] \\
x \arrow[ur, "u"] \arrow[dr, "\bar{u}"'] & \\
 & \alpha(\bar{x}')
\end{tikzcd}
\]
est commutatif), alors $\varphi(x',u) = \varphi(\bar{x}',\bar{u})$. Donc $\varphi$ est constant sur les composantes connexes de $x\backslash C'$, donc constant si $x\backslash C'$ est 0-connexe, d'où $\xi(x)$. On prouve alors aussitôt que pour $x$ variable,
\[ \xi \overset{\text{déf}}{=} (\xi(x))_{x\in\operatorname{Ob} C} \simeq \varprojlim F, \qquad \text{\texttt{[plutôt}}\ \ \in \varprojlim F, \text{\texttt{]}} \]
et que $\alpha^*(\xi) = \xi'$, cqfd.

On peut dire que pour $x_0 \in \operatorname{Ob} C_0$ fixé, si on regarde le préfaisceau \textit{constant} sur $x_0\backslash C'$ de valeur $F(x_0)$, soit $\Phi_{x_0}$, l'application
\[ (x',u) \longmapsto \xi^*(u') : x_0\backslash C' \longrightarrow F(x_0) \]
peut être regardée comme une section de ce préfaisceau sur $x\backslash C'$, défini par $\xi' \in \varprojlim\limits_{C'} \alpha^*F$, donc on a
\[ \varprojlim_{C'} \alpha^*F \xrightarrow{c_{x_0}} \Gamma(x_0\backslash C', \Phi_{x_0}) \qquad (\xleftarrow{\ \sim\ } F(x_0) \text{ si } x_0\backslash C \text{ 0-connexe}). \]
Si on a \fbox{$x_0 \xrightarrow{v} x_1$} dans $C'$, on en déduit
\[ x_1\backslash C' \xrightarrow{V} x_0\backslash C' \]
(les $x\backslash C'$ sont les fibres d'une \textit{catégorie fibrée} $C\backslash C'$ sur $C$ \dots) et on a $F(x_1) \longrightarrow F(x_0)$, d'où \fbox{$\Phi_{x_1} \longrightarrow V^*\Phi_{x_0}$} pour les faisceaux constants correspondants sur $x_0\backslash C'$ resp. $x_1\backslash C'$
\printedpage{5}
% head: A. Grothendieck : Dérivateurs, Ch. II, Édition des Universités de Montpellier II et Paris VII
% foot: Version 8 juillet 2001
\msp{page 7}
et on a commutativité dans
\[
\resizebox{\linewidth}{!}{%
\begin{tikzcd}[column sep=large,ampersand replacement=\&]
\varprojlim\limits_{C'} \alpha^*F \arrow[r, "c_{x_1}"] \arrow[d, "c_{x_0}"'] \& \Gamma(x_0\backslash C', F(x_0)) \simeq \operatorname{Hom}(\pi_0(x_0\backslash C'), F(x_0)) \arrow[d, "V^*"] \& \\
\operatorname{Hom}(\pi_0(x_1\backslash C'), F(x_1)) \simeq \Gamma(x_1\backslash C', F(x_1)) \arrow[r, "\substack{v^*\\ (F(x_1) \longrightarrow F(x_0))}"'] \& \Gamma(x_1\backslash C', F(x_0)) \simeq \operatorname{Hom}(\pi_0(x_1\backslash C'), F(x_0)) \& \\
\pi_0(x_1\backslash C') \arrow[r, "\pi_0(V)"] \arrow[d, dotted] \& \pi_0(x_0\backslash C') \arrow[d, dotted] \& \text{commutatif.} \\
F(x_1) \arrow[r, "F(v)"] \& F(x_0) \&
\end{tikzcd}
}
\]
On considère le préfaisceau d'ensembles
\[ x \longmapsto \pi_0(x\backslash C') \qquad \text{sur } C, \]
appelons-le $\pi_0(C'/C)$, on aura un morphisme canonique
\[ \boxed{\varprojlim_{C'} \alpha^*(F) \xrightarrow{\ \sim\ } \operatorname{Hom}_{\substack{\text{préfaisceaux}\\ \text{d'ensembles sur $C$}}} (\pi_0(C'/C), F),} \]
sans doute toujours \texttt{[un]} isomorphisme.

Si on travaille dans la catégorie dérivée des préfaisceaux des $k$-modules ($k$ un anneau), la formule précédente suggère conjecturalement
\[ \boxed{\mathbf{R}\Gamma(C', \alpha^*F) \simeq \mathbf{R}\underbrace{\operatorname{Hom}}_{\text{Hom internes}}\ (\underbrace{\mathbf{L}\alpha_!(k_{C'})}_{\substack{\text{homologie}\\ \text{relative de}\\ \text{$C'$ sur $C$}}}, F).} \]

\msp{page 8}
L'augmentation canonique
\[ \mathbf{L}\alpha_!(k_{C'}) \longrightarrow k_C \]
redonne, en prenant les $\mathbf{R}\mathrm{Hom}$ dans $F$,
\[ \mathbf{R}\Gamma(C,F) \longrightarrow \mathbf{R}\operatorname{Hom}(\mathbf{L}\alpha_!(k_{C'}), F) \xleftarrow{\ \sim\ } \mathbf{R}\Gamma(C', \alpha^*F) \]
l'homomorphisme canonique de `contravariance' de la cohomologie.

Si par contre on travaille avec le dérivateur non additif HOT, donc $F$ un `type d'homotopie relatif' sur $C$, on doit avoir
\[ \boxed{\mathbf{R}\Gamma_{\mathrm{HOT}}(C', \alpha^*F) \xrightarrow{\ \sim\ } \mathbf{R}\Gamma(C\,,\ \underbrace{\operatorname{Hom}}_{\substack{\text{Hom interne}\\ \text{de la catégorie des}\\ \text{types d'homotopie}\\ \text{relatifs sur $C$}}}\ (\ \alpha_!\ (\ \underbrace{e_{C'}}_{\substack{\text{objet final}\\ \text{de HOT($C'$),}\\ \text{type d'homotopie}\\ \text{relatif final}\\ \text{sur $C'$}}}\ ),\ F)).} \]

\dgrblock{dgrb81}{Global adjunction does not identify an internal Hom}
\begin{dcomment}{examples}{p8:Si par contre on travaille avec le dérivateur non additif HOT}{Global adjunction does not identify an internal Hom}
The proposed formula is useful because its global version is right. The difficulty begins only when its two sides are compared as objects over $C$. For presheaves of spaces, the correct equivalence is
\[
 \operatorname{Map}_{\mathcal P(C')}(1,\alpha^*F)
 \simeq
 \operatorname{Map}_{\mathcal P(C)}(\alpha_!1,F),
 \qquad \mathcal P(C)=\operatorname{Fun}(C^{\mathrm{op}},\mathcal S).
\]
Since $\mathcal P(C)$ is cartesian closed, the right side is also
$\mathbf R\Gamma(C,\underline{\operatorname{Hom}}(\alpha_!1,F))$.
This proves exactly the globally sectioned formula on page~8. The analogous abelian formula on page~7 takes values in $D(k)$ when ``Hom internes'' means the $k$-enriched global Hom. A sheaf-valued internal Hom would require an additional $\mathbf R\Gamma(C,-)$.

To see what global sections forget, take $C=[1]=(0\to1)$ and $i:e\to C$ with $i(*)=1$. A presheaf is an arrow $F_1\to F_0$. Direct calculation of all three adjoints gives
\[
 i^*F=F_1,\qquad i_!A=(A\xrightarrow{1}A),
 \qquad i_*A=(A\to *).
\]
Here the displayed order is always the restriction map from value at $1$ to value at $0$. Every $c/i$ is terminal, so $i$ passes even the contractible-comma criterion; $\Gamma(C,F)=F_1$ and the comparison with $\Gamma(e,i^*F)$ is an identity. Nevertheless the unit has component $F_0\to *$ at $0$, and the counit $i_!i^*F\to F$ has component $F_1\to F_0$ there. Neither is generally invertible.

Moreover, $i_!1=1_C$, hence
\[
 \underline{\operatorname{Hom}}_C(i_!1,F)=F,
 \qquad i_*i^*F=(F_1\to *).
\]
Taking $F=(\{*\}\to\{0,1\})$ disproves the internal identification. Applying global sections to both sides gives $\{*\}$, as it should. With $j(*)=0$ instead, every $j/c$ is terminal, whereas $j_*j^*F=(F_0\xrightarrow{1}F_0)$; thus the opposite comma orientation also fails to imply the universal unit assertion. The same examples work in $D(k)$, replacing $*$ by the zero module and choosing an arrow which is not an isomorphism.

The shift argument on page~9 establishes the expected Hom groups once the derived adjunction is available. Recovering an equivalence of enriched Hom objects requires the enriched adjunction itself, or a constructed comparison and a conservativity argument. It is not supplied merely by a bijection on $\pi_0$ in the space-valued case. No such extra step is needed for the explicit bar-resolution proof above.
\end{dcomment}


\printedpage{6}
% head: A. Grothendieck : Dérivateurs, Ch. II, Édition des Universités de Montpellier II et Paris VII
% foot: Version 8 juillet 2001
En fait, ces formules sont quasi-tautologiques, si on admet l'existence des foncteurs $\alpha_!$ comme adjoints \textit{à gauche} des $\alpha^*$. Dans le cas abélien, disons,
\msp{page 9}
l'isomorphie des $H^0$ des deux membres est la formule d'adjonction
\[ \operatorname{Hom}_{C'}(k_{C'}, \alpha^*F) \simeq \operatorname{Hom}_C(\alpha_!(k_{C'}), F), \]
d'où l'isomorphie pour les $H^i$ des deux membres, par décalage. De même dans le cas de HOT, l'isomorphie canonique des $\pi_0$ des deux membres est la formule d'adjonction. Pour avoir l'isomorphisme plus fin dans $\mathbf{D}(e)$, la catégorie fondamentale du dérivateur envisagé, en termes de Hom internes, on doit admettre l'existence d'objets unité (les $k_X$ dans un cas, les $e_X$ \texttt{[plutôt} $\varepsilon_X$\texttt{]} dans l'autre), compatibles avec images inverses, de telle façon \texttt{[que]} $\Homint(\varepsilon_X, \underbrace{F}_{\in\mathbf{D}_X}) \simeq F$. On aura
\[ \Gamma(C', \alpha^*F) \overset{\text{déf}}{=} f'_*(\alpha^*(F)) \simeq f_*(\alpha_*\alpha^*(F)), \]
et la question est donc de donner une interprétation de $\alpha_*\alpha^*F$ comme
\[ \alpha_*\alpha^*F \simeq \Homint(\alpha_!(\varepsilon_{C'}), F) \qquad \text{dans } \mathbf{D}(C). \]
Pour ceci, on va écrire
\[ \alpha^*(F) \simeq \Homint(\varepsilon_{C'}, F), \]
et la formule doit résulter d'une `formule d'adjonction interne'
\msp{page 10}
\[ ? \qquad\qquad \boxed{\alpha_*(\Homint_{C'}(G', \alpha^*F)) \simeq \Homint_C(\alpha_!G', F).} \]
(\textit{Formule d'adjonction `interne',} pour $\alpha_!$ et $\alpha^*$, valable en principe pour \textit{tout} $\alpha : C' \longrightarrow C$ dans Diag, tout $G'$ dans $\mathbf{D}(C')$ et tout $F$ dans $\mathbf{D}(C)$.)

\dgrblock{dgrb82}{The projection formula is an extra property}
\begin{dcomment}{exposition}{p9:et la formule doit résulter}{The projection formula is an extra property}
Grothendieck correctly reduces the proposed internal adjunction to a projection formula. What fails is the expectation that every change of indexing category has this property. In a closed symmetric monoidal setting, write $L=\alpha_!$, $R=\alpha^*$ and $U=\alpha_*$. The canonical map supplied by the adjunction $L\dashv R$ has direction
\[
 L(RG\otimes G')\longrightarrow G\otimes LG'.
\]
It is adjoint to $RG\otimes G'\to RG\otimes RLG'$, using the unit and the monoidal structure of $R$. The arrow printed on page~10 points in the inverse direction; such an inverse is available only after invertibility has been established.

For every $H$, the adjunctions give
\[
\begin{aligned}
 \operatorname{Hom}(H,U\underline{\operatorname{Hom}}(G',RF))
 &\cong\operatorname{Hom}(L(RH\otimes G'),F),\\
 \operatorname{Hom}(H,\underline{\operatorname{Hom}}(LG',F))
 &\cong\operatorname{Hom}(H\otimes LG',F).
\end{aligned}
\]
Thus the canonical internal comparison
$\underline{\operatorname{Hom}}(LG',F)\to
U\underline{\operatorname{Hom}}(G',RF)$
is invertible for all $F,G'$ exactly when the projection maps are invertible for all $H,G'$. This proves the reduction, and also locates the additional hypothesis precisely.

For cartesian products this property is called Frobenius reciprocity. An essential geometric morphism need not satisfy it. Moreover, Frobenius reciprocity itself is weaker than its validity in every slice, the latter being the condition appropriate to a locally connected geometric morphism. \citet[\S1]{GarnerStreicher2021} explicitly distinguish these conditions; they also recall the equivalence between Frobenius reciprocity and preservation of exponentials by the right adjoint of an adjunction between cartesian closed categories. These later distinctions explain why existence of both Kan extensions does not settle the question. They should not be replaced by an unrestricted ``projection formula for derivators''.
\end{dcomment}


Comment prouver cette formule ? En prenant l'opération produit $\otimes$ associé à l'opération Hom\,int, on est ramené (en prenant les foncteurs contravariants associés) à une formule
\[ G \otimes \alpha_!G' \xrightarrow{\ \sim\ } \alpha_!(\alpha^*G \otimes G'). \]
On trouve par adjonction un homomorphisme canonique, il faut voir s'il est toujours \texttt{[un]} isomorphisme.

\dgrblock{dgrb83}{The exact obstruction for a point inclusion}
\begin{dcomment}{development}{p10:On trouve par adjonction un homomorphisme canonique}{The exact obstruction for a point inclusion}
The walking-arrow example can be sharpened to measure exactly which restriction maps obstruct the projection formula. This makes the question on page~10 concrete even before a general theory of Frobenius adjunctions is introduced.

\textbf{Point-inclusion criterion.}\ours\
Let $d$ be an object of a small category $C$, and let $i:e\to C$ select $d$. For a set-valued presheaf $G$, the projection map
\[
 i_!(i^*G\times A)\longrightarrow G\times i_!A
\]
is invertible for every set $A$ if and only if $G(f):G(d)\to G(c)$ is bijective for every arrow $f:c\to d$. Consequently the formula holds for all $G,A$ if and only if every arrow with codomain $d$ is an isomorphism.

\textit{Proof.}
The left adjoint to evaluation at $d$ is $i_!A=A\times h_d$, where $h_d(c)=C(c,d)$. At $c$ the projection map is
\[
 A\times G(d)\times C(c,d)\longrightarrow
 A\times G(c)\times C(c,d),\qquad
 (a,s,f)\longmapsto(a,G(f)s,f).
\]
For $A=1$ it is a disjoint union, indexed by $f$, of the maps $G(f)$; it is bijective precisely when each of them is. This also suffices for every $A$. If every $f$ is invertible, every $G(f)$ is bijective. Conversely, suppose the formula holds for every presheaf. Taking $G=h_z$ gives a bijection $C(d,z)\to C(c,z)$ by precomposition with $f$, for every $z$. With $z=c$, obtain $g:d\to c$ such that $gf=1_c$. With $z=d$, injectivity and $(fg)f=f=1_df$ imply $fg=1_d$. Thus $f$ is invertible.\hfill$\square$

This includes the object $1$ of $[1]$, where the noninvertible incoming arrow is exactly the obstruction. It also explains why the analogous formula for a group action behaves differently: all indexing arrows are invertible. For a fixed $G$, local constancy along the arrows entering $d$ is the precise repair; no conclusion about all coefficients follows from connectedness of a comma category.

We checked the pointwise Kan-extension treatment of \citet[\S6.2]{Riehl2016} and the Frobenius discussion of \citet[\S1]{GarnerStreicher2021}, and searched for the point-evaluation formulation. We could not find this elementary criterion stated there in this form. It is recorded as our worked continuation of the calculation, with no claim of priority.
\end{dcomment}


\printedpage{7}
% head: A. Grothendieck : Dérivateurs, Ch. II, Édition des Universités de Montpellier II et Paris VII
% foot: Version 8 juillet 2001
\msp{page 11}

\dgrblock{dgrb84}{Why the argument passes through toposes}
\begin{dcomment}{context}{p10}{Why the argument passes through toposes}
The change of subject at the beginning of \S2 is part of the cofinality problem. Grothendieck wants conditions on indexing categories which survive passage to categories of coefficients and which force the relevant homotopy types to disappear. A formulation in terms of the presheaf topos can detect all objects over the category, whereas the underlying preorder remembers only whether an arrow exists.

This is a small-scale instance of the program described in \citet[\S69, especially the discussion of diagram categories on pp.~1--3 of the excerpt]{Grothendieck1983}: the derived category of coefficients is accompanied by derived categories of diagrams of every shape. The concern there comes from operations on coherent diagrams which a single triangulated category does not retain. The present calculation asks what changing the shape does to sections and their derived functors. In parallel, \citet[Chapter~II, \S1; Chapter~III, \S2]{Heller1988} axiomatizes homotopy theories and develops the computation of homotopy Kan extensions. These references explain the setting, not a historical claim that one author's argument was taken from the other.

The additional tensor language on pages~7--10 needs separate structure. Heller's Chapter~IV treats tensor and cotensor operations; the preceding projection counterexample shows why that structure cannot be inferred from the mere existence of left and right Kan extensions. We now return to a question for which the ordinary presheaf topos is already sufficient: which part of the filtering conditions is visible in the behavior of its connected objects?
\end{dcomment}


\section*{2\quad Les catégories localement filtrantes sont les catégories localement totalement $0$-connexes}

Soit $\mathcal{T}$ un topos. On dit que $\mathcal{T}$ est \textit{irréductible} si l'ensemble ordonné $\operatorname{Ouv}(\mathcal{T})$ des sous-objets de $e_{\mathcal{T}}$ a la propriété suivante : a) $\operatorname{Ouv}(\mathcal{T})$ n'est pas réduit à un élément, i.e. $\emptyset_{\mathcal{T}} \longrightarrow e_{\mathcal{T}}$ n'est pas un isomorphisme ($\mathcal{T}$ n'est pas identique à son sous-topos vide). b) Si $U'$, $U''$ sont deux objets `non vides', i.e. non initiaux, de $\operatorname{Ouv}(\mathcal{T})$, leur intersection est `non vide'. Il revient au même pour b) de dire que si $\mathcal{T}'$, $\mathcal{T}''$ sont deux sous-topos fermés distinct de $\mathcal{T}$, alors leur supremum $\mathcal{T}' \cup \mathcal{T}''$ est distinct de $\mathcal{T}$, ou encore que tout ouvert de $\mathcal{T}$ est connexe. On dit que $\mathcal{T}$ est \textit{localement irréductible} si a) $\mathcal{T}$ est localement connexe (i.e. chaque faisceau sur $\mathcal{T}$ est somme de faisceaux 0-connexes), et si b) pour tout faisceau 0-connexe $U$ sur $\mathcal{T}$, le topos induit $\mathcal{T}/U$ est irréductible (i.e. l'ensemble ordonné $\operatorname{Ouv}(U)$ des sous-objets de $U$ a la propriété b) ci-dessus ( a) signifie simplement que $U \not\simeq \emptyset_{\mathcal{T}}$, i.e. que $U$ n'est pas objet initial dans $\mathcal{T}^\wedge$ ; ou plus simplement) : tout sous-faisceau d'un faisceau connexe est connexe (\footnote[5]{\textbf{NB} Irréductible n'implique pas localement irréductible, pas plus que connexe n'implique pas localement connexe etc.}) (\footnote[6]{$\mathcal{T}$ est localement irréductible si et seulement s'il a un ensemble des générateurs $(U_\alpha)$, avec les $U_\alpha$ irréductibles (donc 0-connexes, donc $\mathcal{T}$ est localement connexe).}).

\dgrblock{dgrb85}{Irreducibility of an object and of its support}
\begin{dcomment}{exposition}{p11:tout sous-faisceau d'un faisceau connexe est connexe}{Irreducibility of an object and of its support}
Grothendieck is distinguishing two ways to forget morphisms. The space $\operatorname{Esp}(X)$ remembers downward-closed sets of objects; the topos $\widehat X=\operatorname{Fun}(X^{\mathrm{op}},\mathrm{Set})$ also remembers every category of elements. The first sees a common predecessor of two objects. The second must see common predecessors compatible with prescribed arrows.

Here an object is \emph{connected} if it has no decomposition into two noninitial disjoint summands; ``$0$-connected'' adds noninitiality. Thus the assertion that every subobject of a connected object is connected allows the initial subobject. An object $A$ is irreducible when it is noninitial and any two noninitial subobjects of $A$ meet. Equivalently, every subobject is connected: a disconnected subobject gives two disjoint nonempty subobjects, while disjoint subobjects have a disconnected union. Applied to $A=1$, this gives the definition of an irreducible topos. Applied to every nonempty connected $A$, after assuming local connectedness, it gives the local definition.

For $\widehat X$, subterminal objects are precisely sieves of objects: if $x$ belongs and $y\to x$ exists, then $y$ belongs. Any two nonempty such sieves meet exactly when each pair of objects has a common predecessor. This proves
\[
 \widehat X\text{ irreducible}
 \quad\Longleftrightarrow\quad
 X\ne\varnothing\text{ and }
 (\forall x,y)(\exists z)(z\to x, z\to y).
\]
It is a condition on the associated preorder, not on equalities of arrows. A least class in that preorder need not give an initial object: uniqueness of arrows has been discarded. If the preorder quotient is artinian, a minimal class exists, and common lower bounds force it below every other class; this proves the assertion on page~13. Finite principal lower sets suffice as well. ``Locally finite'' here must have that meaning, as indicated by the finite slices $X/x$; finiteness of intervals alone would fail, for example for the ordered set $\mathbb Z$.

The list of equivalences on pages~12--13 is provisional. Its claim about components being weakly cofiltered is too weak, and its subsequent inclusion of PS~2 is too strong. The annotations and Proposition~2.1 correct these in opposite directions. The surviving characterization is PS~1 alone. We prove it below, retaining the successive formulations as part of the text's development.
\end{dcomment}


Quand $\mathcal{T}$ est de la forme $\operatorname{Top}(T)$, où $T$ est un espace topologique, alors $\mathcal{T}$ irréductible signifie que $T$ est irréductible (car alors $\operatorname{Ouv}(\mathcal{T}) \simeq \operatorname{Ouv}(T)$). Cela implique déjà trivialement que $T$ est localement connexe (tout ouvert est somme disjoint d'ouverts connexes), donc $\operatorname{Top}(\mathcal{T})$ est localement connexe. De plus, cela implique aussi que $\operatorname{Top}(T)$ est localement irréductible (\footnote[7]{expliciter la démonstration.}). (\textbf{NB} Les objets 0-connexes % [sic] \operatorname{Top}(\mathcal{T}) -> \operatorname{Top}(T), Top s'applique à l'espace T et non au topos \mathcal{T}
\msp{page 12}
de $\operatorname{Top}(T)^\wedge$, i.e. les faisceaux d'ensembles 0-connexes sur $T$, sont ceux dont la `fibre générique' a exactement un point \dots) Aussi $\mathcal{T} = \operatorname{Top}(T)$ est localement irréductible si et seulement si $T$ \texttt{[est]} localement connexe, et les composantes connexes de $T$ sont irréductibles, i.e. si et seulement si $T = \coprod_i T_i$, les $T_i$ irréductibles.

\dgrblock{dgrb86}{The generic fibre without a generic point}
\begin{dcomment}{examples}{p12:les composantes connexes de $T$ sont irréductibles}{The generic fibre without a generic point}
The topological case supplies the intuition for the definition, but it also calls for a precise interpretation of ``generic fibre''. An irreducible space need not possess a generic point. For a sheaf $F$ on a nonempty irreducible space $T$, define instead
\[
 \gamma(F)=\mathop{\operatorname{colim}}_{\varnothing\ne U\subseteq T}F(U),
\]
where transition maps are restrictions as $U$ shrinks. Nonempty opens form a filtering system under intersection. If a generic point exists, this is its stalk; otherwise it is the stalk at the corresponding point of the soberification. The following direct argument avoids requiring that point.

For every nonempty open $U$, a section of $F(U)$ has a well-defined generic germ. Because $U$ is connected, the constant sheaf with value $\gamma(F)$ has sections $\gamma(F)$ on $U$. We obtain a map $F\to\underline{\gamma(F)}$ and a decomposition
\[
 F\cong\coprod_{g\in\gamma(F)}F_g
\]
into its fibres. Every $F_g$ is noninitial because $g$ has a representative section. If $A,B$ are two noninitial subobjects of $F_g$, choose sections $a\in A(U)$ and $b\in B(V)$ on nonempty opens. They have the same generic germ, so on some nonempty $W\subseteq U\cap V$ their restrictions agree. This supplies a section of $A\cap B$, proving that $F_g$ is irreducible.

Consequently these fibres are exactly the connected components of $F$, and $F$ is $0$-connected if and only if $\gamma(F)$ is a singleton. This proves both the assertion about the generic fibre and the missing implication in footnote~7: the sheaf topos of an irreducible space is locally irreducible. More generally, if $T$ is a disjoint union of open irreducible spaces, the conclusion follows component by component. Conversely, a locally irreducible sheaf topos is locally connected; its terminal object decomposes into connected subterminal objects, which are disjoint open subsets of $T$, and each is irreducible. This proves the topological characterization stated here.

For a general presheaf topos, no generic point of the underlying preorder can replace this argument. Parallel arrows produce different objects in slices, even when they have the same source and target. The next example shows exactly where the analogy with topological spaces ceases to suffice.
\end{dcomment}


Si $\mathcal{T}$ est de la forme $\operatorname{Top}(X)$, $X$ une catégorie, on sait que $\mathcal{T}$ est localement connexe (et même localement $\infty$-connexe). On voit d'autre part que $\mathcal{T}$ est irréductible si et seulement si $\operatorname{Esp}(X)$ (l'espace topologique associé, dont l'ensemble sous-jacent est $\operatorname{Ob} X$) est irréductible (\footnote[8]{car les ouverts de $\operatorname{Top}(X)$, ceux de $\operatorname{Esp}(X)$ et ceux de $X$ se correspondent.}) ; et on dit alors que $X$ est irréductible (\footnote[9]{Cela signifie aussi que pour $x$, $y$ dans $X$, $\exists\, z$ avec \begin{tikzcd}[ampersand replacement=\&, row sep=0.2em, column sep=1.2em, baseline=(\tikzcdmatrixname-2-1.base)] \& x \\ z \arrow[ur] \arrow[dr] \& \\ \& y \end{tikzcd}, i.e. $\operatorname{Ord}(X)$ est filtrant à gauche.}). Comme les composantes connexes de $X$ et celles de $\operatorname{Esp}(X)$ ou $\operatorname{Top}(X) = \mathcal{T}$ se correspondent, on voit donc que les conditions suivantes sont équivalentes : a) le topos $\operatorname{Top}(X)$ est localement irréductible, b)
\printedpage{8}
% head: A. Grothendieck : Dérivateurs, Ch. II, Édition des Universités de Montpellier II et Paris VII
% foot: Version 8 juillet 2001
$\operatorname{Esp}(X)$ est localement irréductible, c) les composantes connexes de $X$ sont irréductibles, d) les $X/x$ \texttt{[sont]} irréductibles. On dit alors que $X$ est localement irréductible (\footnote[10]{Faux, on a (a $\Longleftrightarrow$ d) $\Longrightarrow$ (b $\Longleftrightarrow$ c) $X$ irréductible n'implique \textit{pas} que les $X/x$ le soient, exemple $X = (0 \overset{\alpha}{\underset{\beta}{\rightrightarrows}} 1)$ $(\alpha \neq \beta)$, on a $X/1 =$ \begin{tikzcd}[ampersand replacement=\&, row sep=0.2em, column sep=1.2em, baseline=(\tikzcdmatrixname-2-2.base)] 0' \arrow[dr] \& \\ \& 1 \\ 0'' \arrow[ur] \& \end{tikzcd}, qui n'est \textit{pas} irréductible, ni localement irréductible puisqu'il est connexe. Pourtant il est étale sur $X$.}).

Cela implique donc que pour tout $X'$ étale sur $X$, avec $X'$ 0-connexe, $X'$ est irréductible.

\textit{Dire conditions équivalentes.}
\begin{enumerate}
\item[a)] $\operatorname{Top}(X)$ est localement irréductible.
\item[b)] Les $X/x$ \texttt{[sont]} irréductibles.
\item[c)] Tout objet 0-connexe de $X^\wedge$ est irréductible.
\item[d)] Les composantes connexes de $X$ sont faiblement cofiltrantes (i.e. faiblement filtrantes à gauche).
\item[e)] Condition PS $1_{\mathrm{g}}$ de SGA 4, I 2.
\end{enumerate}
Si $X$ est \texttt{[une]} catégorie, $X$ irréductible signifie aussi que $X \neq \emptyset$ et $\forall\, x, y \in \operatorname{Ob} X$, \texttt{[il]} existe \texttt{[un]} $s$ avec \begin{tikzcd}[row sep=0.2em, column sep=1.2em, baseline=(\tikzcdmatrixname-2-1.base)] & x \\ s \arrow[ur] \arrow[dr] & \\ & y \end{tikzcd}. (Cela signifie aussi que $X$ est \textit{faiblement filtrante} (\footnote[11]{Il est mieux l'appeler \textit{pseudo-filtrante,} et réservons faiblement filtrante à : localement faiblement filtrante + connexe.}) à gauche, ou faiblement cofiltrante.) $X$ localement irreductible signifie que pour $x$, $y$ \textit{appartenant à une même composante connexe,} \texttt{[il]} existe \texttt{[un]} $s$ avec \begin{tikzcd}[row sep=0.2em, column sep=1.2em, baseline=(\tikzcdmatrixname-2-1.base)] & x \\ s \arrow[ur] \arrow[dr] & \\ & y \end{tikzcd}, % [sic] irreductible -> irréductible, accent omis
\msp{page 13}
i.e. $X/[x] \cap X/[y] \neq \emptyset$ (où $[x]$ est considéré comme élément de l'ensemble ordonné $\operatorname{Ob} X$). Il revient encore au même de dire que les catégories induites $X/x$ sont irréductibles (\footnote[12]{faux, cf. annotation marginale plus haut.}).

\textbf{Exemple.} Une catégorie \textit{ayant un objet initiale,} ou seulement un `plus petit objet', est irréductible. Inversement, si l'ensemble \textit{ordonné} $\operatorname{Ord}(X)$ associé à $X$ (quotient de $\operatorname{Ob}(X)$ par la relation d'équivalence associé à son préordre) est fini, p.ex. $X$ fini, ou seulement localement fini, p.ex. $X$ localement fini (i.e. les catégories $X/x$ finies), alors si $X$ est irréductible, $\operatorname{Ob} X$ a un plus petit élément, i.e. $X$ a un plus petit objet. C'est vrai plus généralement encore, si $\operatorname{Ord}(X)$ est artinien, i.e. toute suite décroissante dans $\operatorname{Ord}(X)$ est stationnaire. % [sic] initiale -> initial, accord avec « objet »

Mais la condition intéressante est la condition topossique de \textit{locale} irréductibilité, qui équivaut : $I$ pseudo-filtrante à gauche, i.e. les composantes connexes de $I$ filtrantes, i.e. les conditions PS 1, PS 2

PS 1\enspace Tout diagramme \begin{tikzcd}[row sep=0.2em, column sep=1.2em, baseline=(\tikzcdmatrixname-2-2.base)] x \arrow[dr] & \\ & s \\ y \arrow[ur] & \end{tikzcd} se complète en \begin{tikzcd}[row sep=0.2em, column sep=1.2em, baseline=(\tikzcdmatrixname-2-1.base)] & x \arrow[dr] & \\ z \arrow[ur] \arrow[dr] & & s \\ & y \arrow[ur] & \end{tikzcd} \textit{commutatif.}

PS 2\enspace Pour toute double flèche $x \rightrightarrows y$, \texttt{[il]} existe un égalisateur $z \longrightarrow x \rightrightarrows y$.

% ---- from Der-ChII-234_p011-015.tex
\printedpage{9}
% head: A. Grothendieck : Dérivateurs, Ch. II, Édition des Universités de Montpellier II et Paris VII
% foot: Version 8 juillet 2001

\msp{page 14}

En résumé, on a la

\noindent\textbf{Proposition 2.1.} \textit{Soit $X$ un objet de} Cat. \textit{Conditions équivalentes :}
\begin{enumerate}[label={},leftmargin=3em]
\item[\textit{a)}] \textit{Le topos associé} $\operatorname{Top}(X)$ \textit{est localement irréductible, i.e. (étant localement connexe) pour tout objet $F$ de $X^{\wedge}$ et tout sous-objet $F'$ de $F$, si $F$ est connexe, $F'$ l'est.}
\item[\textit{a$'$)}] \textit{Pour tout $X'$ étale sur $X$, si $X'$ est connexe, tout ouvert de $X'$ est connexe.}
\item[\textit{a$''$)}] \textit{Tout ouvert d'une catégorie localisée $X/a$ ($a \in \operatorname{Ob} X$) est connexe (i.e. $X/a$ est une catégorie irréductible, i.e. le topos} $\operatorname{Top}(X/a)$ \textit{l'est).}
\item[\textit{b)}] \textit{Pour tout objet $a$ de $X$,} $\operatorname{Esp}(X/a)$ \textit{est un espace topologique irréductible.}
\item[\textit{c)}] \textit{Les catégories $X/a$ sont faiblement filtrantes à gauche (ou faiblement cofiltrantes), i.e. pour deux objets $\xi$, $\eta$ d'une telle catégorie,} \texttt{[il]} \textit{existe un objet $\zeta$ et deux flèches}
\begin{tikzcd}[ampersand replacement=\&, row sep=tiny, column sep=small, baseline=(\tikzcdmatrixname.center)]
 \& \xi \\
\zeta \arrow[ur] \arrow[dr] \& \\
 \& \eta
\end{tikzcd}
\textit{(i.e. l'ensemble ordonné Ord($X/a$) associé à $X/a$ est filtrant à gauche)}
\item[\textit{c$'$)}] \textit{Pour tout diagramme en traits pleins dans $X$,}
\begin{tikzcd}[ampersand replacement=\&, row sep=tiny, column sep=small, baseline=(\tikzcdmatrixname.center)]
 \& x \arrow[dr] \& \\
z \arrow[ur, dotted] \arrow[dr, dotted] \& \& s \\
 \& y \arrow[ur] \&
\end{tikzcd}
\textit{, on peut le compléter en un carré} commutatif \textit{comme indiqué. (Condition} PS 1 \textit{des catégories localement filtrantes.)}
\end{enumerate}

\noindent Ces conditions impliquent que toute composante connexe $X_{\alpha}$ de $X$ est irréductible, ou ce qui revient au même, que $\operatorname{Esp}(X_{\alpha})$ est irréductible,

\msp{page 15}

\noindent mais ne sont pas impliquées par celle-ci.

Contrexemple, la catégorie 0-connexe $0 \underset{\beta}{\overset{\alpha}{\rightrightarrows}} 1$, avec $\alpha \neq \beta$. On a $X/1 \simeq$
\begin{tikzcd}[ampersand replacement=\&, row sep=tiny, column sep=small, baseline=(\tikzcdmatrixname.center)]
0 \arrow[dr] \& \\
 \& 1 \\
0' \arrow[ur] \&
\end{tikzcd}, qui n'est \textit{pas} irréductible, puisque $\Phi \backslash \{1\} \simeq \varepsilon$ est un ouvert disconnexe \texttt{[}$\Phi = X/1$\texttt{,} $\varepsilon = \bullet\ \bullet$\texttt{]}.

J'aimerais une condition `topossique' sur $X$, équivalente à $X$ `pseudo-filtrante à gauche', i.e. les composantes connexes filtrantes, i.e. PS 1\textsuperscript{opp} (condition c$'$) ci-dessus) + PS 2\textsuperscript{opp} (condition de la double flèche), et qui de plus implique que tout objet 0-connexe de $X^{\wedge}$ (i.e. tout objet $X'$ de Cat étale sur $X$ et 0-connexe) soit $W_{\infty}$-asphérique. Et je songe à la notion de `totale asphéricité' de Pursuing Stacks.

\dgrblock{dgrb87}{Completing Proposition 2.1}
\begin{dcomment}{exposition}{p15:J'aimerais une condition}{Completing Proposition 2.1}
Proposition~2.1 isolates the part of filtering which can be expressed without equalizing parallel arrows. We give its proof because it explains both the corrected statement and the mechanism which will later be strengthened.

Recall that PS~1 says that every cospan $a\to x\leftarrow b$ admits a commuting square with some upper-left object $c$; no universal property is requested. Suppose PS~1 holds. In any category with this property, the relation ``$a$ and $b$ have a common predecessor'' is reflexive and symmetric. It is transitive: given $a\leftarrow r\to b$ and $b\leftarrow s\to d$, apply PS~1 to $r\to b\leftarrow s$ and compose the resulting arrows. Every arrow connects objects related in this sense. Hence any two objects in the same connected component have a common predecessor.

For a presheaf $F$, its category of elements $E=\int_XF\to X$ is a discrete fibration, the meaning of ``\'etale'' here. Every arrow into the image of an object has a unique lift into that object. Thus a PS~1 square in $X$ lifts to a PS~1 square in $E$: the two lifts into the common target agree by uniqueness. Therefore $E$ also satisfies PS~1. Components of $E$ correspond to coproduct components of $F$. If $F$ is $0$-connected, $E$ is nonempty and connected, so the previous paragraph gives common predecessors for every pair of its objects. Subobjects of $F$ correspond to sieves in $E$; any two nonempty such sieves meet. Thus $F$ is irreducible, and $\widehat X$ is locally irreducible.

Conversely, suppose every $0$-connected presheaf is irreducible. The representable $h_x$ is $0$-connected, since its category of elements $X/x$ has a terminal object. Given arrows $a\to x$ and $b\to x$, their generated subobjects of $h_x$ are noninitial. Their intersection is noninitial, so it contains some element represented by arrows $c\to a$ and $c\to b$ with equal composite to $x$. This is PS~1. The descriptions using $X/x$, its associated space, and open subcategories are equivalent by the preceding sieve argument. This proves all the equivalences of Proposition~2.1.

The counterexample is the category with two objects and distinct arrows $\alpha,\beta:0\rightrightarrows1$. Its objects have the common predecessor $0$, so its underlying preorder is the irreducible two-point chain. But in $X/1$, the objects $(0,\alpha)$ and $(0,\beta)$ have no common predecessor: the only endomorphism of $0$ is its identity. Their two singleton lower sets are disjoint nonempty subobjects of $h_1$. Thus $\widehat X$ is not locally irreducible. The mistaken converse on page~16 is precisely the assertion that the preorder remembers this equality of composites; the footnotes already retract it.
\end{dcomment}


\noindent\textbf{Définition 2.2.} (\footnote[13]{Cf. tableau récapitulatif pages 97, 98 et surtout 99.})\ \ Soit $X$ une catégorie.
\begin{enumerate}[label={},leftmargin=3em]
\item[a)] On dit que $X$ est \textit{faiblement filtrante} (\footnote[14]{dire plutôt \textit{pseudo-filtrante}, réserver `faiblement filtrante' à `localement faiblement filtrante' (ou \textit{irréductible}) + connexe, condition PS $1'_{\mathrm{g}}$.}) si $X \neq \emptyset$, et si pour $a, b \in \operatorname{Ob} X$, $\exists\, c$ et des flèches $c$
\begin{tikzcd}[ampersand replacement=\&, row sep=tiny, column sep=small, baseline=(\tikzcdmatrixname.center)]
 \& a \\
{} \arrow[ur, phantom] \& \\
 \& b
\end{tikzcd}
\hspace{-2.6em}% layout aid only
\begin{tikzcd}[ampersand replacement=\&, row sep=tiny, column sep=small, baseline=(\tikzcdmatrixname.center)]
 \& {} \\
{} \arrow[ur] \arrow[dr] \& \\
 \& {}
\end{tikzcd}
, i.e. $X/[a] \cap X/[b] \neq \emptyset$. (\textbf{NB} $X$ est faiblement filtrante (\footnote[15]{pseudo-filtrante.}) à gauche si et seulement si $\operatorname{Ord}(X)$ l'est, i.e. est filtrante à gauche. \textbf{ReNB} $X$ faiblement filtrante $\Longrightarrow$ $X$ 0-connexe.)

\item[b)] \printedpage{10}
% head: A. Grothendieck : Dérivateurs, Ch. II, Édition des Universités de Montpellier II et Paris VII
% foot: Version 8 juillet 2001
On dit que $X$ est \textit{localement faiblement filtrante à gauche} ou plus simplement préfiltrante à gauche si les catégories $X/x$ ($x \in \operatorname{Ob} X$) sont faiblement filtrantes à gauche, i.e. si $X$ satisfait la condition : $\forall$ diagrammes
\begin{tikzcd}[ampersand replacement=\&, row sep=tiny, column sep=small, baseline=(\tikzcdmatrixname.center)]
a \arrow[dr] \& \\
 \& x \\
b \arrow[ur] \&
\end{tikzcd}, $\exists\, c$ dans $X$ et des flèches de $c$ dans $a$, $b$ rendant \textit{commutatif} le carré
\[
\begin{tikzcd}[ampersand replacement=\&, row sep=tiny, column sep=small]
 \& a \arrow[dr] \& \\
c \arrow[ur, dotted] \arrow[dr, dotted] \& \& x \\
 \& b \arrow[ur] \&
\end{tikzcd}
\]
(\footnote[16]{Condition PS $1_{\mathrm{g}}$.}). (\textbf{NB} Cette condition implique que $\operatorname{Ord}(X)$ est localement faiblement filtrant, mais n'équivaut pas à cette dernière. Exemple, $X = 0 \underset{\beta}{\overset{\alpha}{\rightrightarrows}} 1$, $X/1 =$
\begin{tikzcd}[ampersand replacement=\&, row sep=tiny, column sep=small, baseline=(\tikzcdmatrixname.center)]
0' \arrow[dr] \& \\
 \& 1 \\
0'' \arrow[ur] \&
\end{tikzcd}
n'est pas faiblement filtrante à gauche.)

\msp{page 16}

\textbf{NB} Si $X$ localement faiblement filtrante à gauche, tout $X'$ étale sur $X$ itou.

\textbf{NB} (\footnote[17]{\textbf{Faux.}}) $X$ localement faiblement filtrante à gauche $\Longleftrightarrow$ les composantes connexes sont faiblement filtrantes à gauche, ou encore si $X$ est 0-connexe, $X$ localement faiblement filtrante à gauche $\Longleftrightarrow$ $X$ est faiblement filtrante (\footnote[18]{pseudo-filtrante.}) à gauche. On a en effet $\Longleftarrow$ tautologiquement (\footnote[19]{\textbf{faux,} ainsi $0 \underset{\beta}{\overset{\alpha}{\rightrightarrows}} 1$ est 0-connexe et faiblement filtrante à gauche, mais non pas localement faiblement filtrante à gauche car
\begin{tikzcd}[ampersand replacement=\&, row sep=tiny, column sep=small, baseline=(\tikzcdmatrixname.center)]
0 \arrow[dr, "\alpha"] \& \\
 \& 1 \\
0 \arrow[ur, "\beta"'] \&
\end{tikzcd}
ne se complète pas en un carré \textit{commutatif.}}), et $\Longrightarrow$ en montrant que si PS 1 \texttt{[est]} satisfait, alors $PS(a,b)$ défini par $X/[a] \cap X/[b] \neq \emptyset$ est une relation d'équivalence, impliquée par la relation $C(a,b)$ : $\{\operatorname{Hom}(a,b) \neq \emptyset$ \textit{ou} $\operatorname{Hom}(b,a) \neq \emptyset\}$, donc la relation d'équivalence (homotopie) engendrée par $C(a,b)$ est plus fine que $PS(a,b)$.

\item[c)] On dit que $X$ est \textit{filtrante à gauche}, ou \textit{cofiltrante,} si elle est faiblement filtrante (\footnote[20]{pseudo-filtrante.}) à gauche (donc \textit{non vide}), et satisfait (en plus de la condition PS $1'_{\mathrm{g}}$ de a)) la condition PS $2_{\mathrm{g}}$ suivante :

Pour toute double flèche $a \underset{v}{\overset{u}{\rightrightarrows}} b$, \texttt{[il]} existe \texttt{[une]} flèche $c \xrightarrow{w} a$ telle que $uw = vw$ (\footnote[21]{i.e. trois conditions a) $X \neq \emptyset$ b) PS $1'$ habituelle c) PS 2 (double flèche).}).

Il revient au même de dire que $X$ satisfait les trois conditions a$'$) $X$ est 0-connexe (plus forte que $X \neq \emptyset$), b$'$) PS 1, i.e. $X$ localement faiblement filtrante à gauche et c$'$) PS 2 comme avant.

\item[d)] On dit que $X$ est \textit{localement filtrante à gauche} si les composantes connexes sont filtrantes à gauche, i.e. si $X$ satisfait les conditions PS 1 et PS 2 (`pseudo-cofiltrante' dans la terminologie de SGA 4). \textbf{NB} Si $X$ est localement filtrante à gauche, il en est de même de tout $X'$ étale sur $X$, car c'est évident pour PS 1 qui concerne les $X/x$, et pour PS 2 cela résulte du fait que $X' \longrightarrow X$ est fidèle. Donc si $X$ est localement cofiltrante, les $X/a$ sont filtrantes, mais la réciproque est fausse.
\end{enumerate}

\dgrblock{dgrb88}{The equalization condition survives beyond every slice}
\begin{dcomment}{examples}{p16:mais la réciproque est fausse.}{The equalization condition survives beyond every slice}
The definitions now separate PS~1 from the additional condition PS~2. In the latter, ``\'egalisateur'' means an \emph{equalizing arrow}: for $u,v:a\rightrightarrows b$ there is $w:c\to a$ with $uw=vw$. It does not mean that a categorical equalizer exists. This is the usual cofiltered condition, dual to SGA~4's PS~2 \citep[Expos\'e~I, \S2.7]{SGA4I1972}.

To see why all slices can be cofiltered without the original category being locally cofiltered, take the one-object category $BG$ for a nontrivial group $G$. It satisfies PS~1: a pair $f,g$ can be completed by $1$ and $g^{-1}f$. For its unique object $*$, the slice $BG/*$ has objects $g\in G$ and exactly one arrow from $g$ to $h$, namely $h^{-1}g$. It is equivalent to the terminal category and is cofiltered. Yet distinct $u,v\in G$ cannot be equalized on the right, since $uw=vw$ and invertibility of $w$ imply $u=v$. Hence $BG$ fails PS~2.

This example also makes the topos distinction visible. Every connected $G$-set is transitive and therefore irreducible, so the presheaf topos is locally irreducible. Products of connected objects need not remain connected: for the regular right $G$-set, the diagonal right action on $G\times G$ has one orbit for each value of $hg^{-1}$. The product condition introduced next detects the missing equalization information. No defect in Proposition~2.1 is involved; its purpose was exactly to isolate PS~1.
\end{dcomment}


\printedpage{11}
% head: A. Grothendieck : Dérivateurs, Ch. II, Édition des Universités de Montpellier II et Paris VII
% foot: Version 8 juillet 2001

\noindent On a les notions duales : \textit{faiblement filtrante à droite} (ou simplement faiblement filtrante), et \textit{filtrante à droite} ou simplement filtrante. (Notions duales de a), c) ci-dessus.) Pour la notion duale de b), savoir que les $x\backslash X$ soient faiblement filtrantes, on ne peut plus l'appeler `localement faiblement filtrante', mais ce serait \textit{co}localement faiblement filtrante.

\noindent\textbf{Exemple.} $X = \Bigl(0 \xrightarrow{\gamma} 1 \underset{\beta'}{\overset{\alpha'}{\rightrightarrows}} 2 \underset{\beta}{\overset{\alpha}{\rightrightarrows}} 3\Bigr)$ avec $(\alpha\alpha' = \beta\beta') \neq (\alpha\beta' = \beta\alpha')$.

\dgrblock{dgrb89}{The finite example and its homotopy type}
\begin{dcomment}{development}{p16:avec $(\alpha\alpha' = \beta\beta')}{The finite example and its homotopy type}
The example between Definition 2.2 and Proposition 2.3 makes the loss of information in $\operatorname{Ord}(X)$ concrete. We read the displayed graph and its two relations as a presentation of $X$; in particular, no additional relation is imposed on composites with $\gamma$. Put
\[
 r=\alpha\alpha'=\beta\beta',\qquad
 s=\alpha\beta'=\beta\alpha',\qquad r\ne s.
\]
Then $\operatorname{Ord}(X)$ is the chain $0<1<2<3$. Every pair of objects has both a common predecessor and a common successor. Nevertheless, $0$ is not initial: the arrows $\alpha'\gamma,\beta'\gamma:0\rightrightarrows2$ are distinct. Likewise, $3$ is not terminal, since $r,s:1\rightrightarrows3$ are distinct.

The same calculation locates the exact failure of PS 1. A commuting square for the cospan $1\xrightarrow{\alpha'}2\xleftarrow{\beta'}1$ would require two arrows from an object into $1$. Their common source can only be $1$ or $0$, and both arrows would be respectively $1_1$ or $\gamma$. Neither choice makes the square commute. Thus the left PS 1 fails, and the pair $\alpha',\beta'$ also has no equalizing arrow on the left. Dually, a commuting cocone over $3\xleftarrow{\alpha}2\xrightarrow{\beta}3$ can only use identity arrows out of $3$, so the right PS 1 fails; the pair $\alpha,\beta$ cannot be equalized on the right either. The category is consequently neither locally cofiltered nor locally filtered, although its associated poset satisfies both conditions. This is the point of the paragraph immediately following the example.

The crossed composition relations also admit a useful homotopical interpretation which the manuscript does not develop.

\textbf{The homotopy type of the finite example.}\ours
\textit{For the category presented above, $|N X|\simeq\mathbb{R}P^2$. In particular, its contractible reachability poset does not determine even the fundamental group of its nerve.}

\textit{Proof.} Let $C$ be the full subcategory on $1,2,3$. The functor $q:X\to C$ sends $0$ to $1$, sends $\gamma$ to $1_1$, and fixes the remaining generators. It respects the defining relations. For the inclusion $j:C\hookrightarrow X$, one has $qj=1_C$ and a natural transformation $1_X\to jq$ whose component at $0$ is $\gamma$ and whose other components are identities. The resulting homotopy of nerves fixes $N C$, so $|N X|\simeq|N C|$.

Identify each set $C(i,j)$, for $1\le i<j\le3$, with $\mathbf Z/2$: label $\alpha',\alpha,r$ by $0$ and $\beta',\beta,s$ by $1$. The two displayed relations say exactly that composition adds these labels. Let $P$ be the poset with objects $(i,\varepsilon)$, where $i\in\{1,2,3\}$ and $\varepsilon\in\mathbf Z/2$, with $(i,\varepsilon)<(j,\delta)$ precisely when $i<j$. Its involution changes every $\varepsilon$ to $\varepsilon+1$. There is an invariant functor $P\to C$ sending $(i,\varepsilon)$ to $i$ and an arrow to the label $\delta-\varepsilon$. A chain of arrows in $C$ lifts uniquely after choosing the lift of its first vertex, and the two choices are exchanged by the involution. This proves, degree by degree and compatibly with faces and degeneracies,
\[
 N P/(\mathbf Z/2)\simeq N C.
\]
The order complex of $P$ is the join of three two-point sets. Sending their vertices to $\{\pm e_1,\pm e_2,\pm e_3\}$ identifies it with the boundary of the octahedron. Radial projection identifies its realization with $S^2$, and the involution with the antipodal map. Realization commutes with the quotient, giving
\[
 |N C|\simeq |N P|/(\mathbf Z/2)\cong S^2/(x\sim-x)=\mathbb{R}P^2.
\]
Since the sphere is simply connected, this double covering also gives $\pi _1(|N X|)\cong\mathbf Z/2$.\hfill$\square$

The join and antipodal descriptions are classical \citep[p.~9 and Example 1.43]{Hatcher2002}. Our addition is their explicit identification with this particular presentation. The bounded literature searches did not locate that identification; the result is marked as our calculation, without a claim of priority.
\end{dcomment}


\msp{page 17}

\noindent Seule la notion a) (faiblement filtrante (\footnote[22]{pseudo-filtrant.}) à gauche) ne dépend que de $\operatorname{Ord} X$ - mais bien sûr si $X$ est \textit{pré}ordonnée, toutes trois ne dépendent que de $\operatorname{Ord} X$. Notons que dire que ça ne dépend que de $\operatorname{Ord} X$, ou que de $\mathrm{Esp}\,\mathrm{sep}X$, revient au même, puisque les deux se déterminent mutuellement. La condition a) signifie que $\operatorname{Esp} X$, ou aussi $\operatorname{Top}(X)$, est \textit{irréductible.} On dira que $X$ est \textit{irréductible} (au lieu de pseudo-filtrante). Les trois autres notions ne peuvent s'exprimer en termes de $\operatorname{Esp} X$ ou de $\operatorname{Ord} X$ seul, mais elles s'expriment en termes du $\mathrm{Topos}\, X$. La condition b) ($X$ localement faiblement filtrante) signifie que $\operatorname{Top}(X)$ est \textit{localement irréductible,} et on dira aussi que $X$ est \textit{localement irréductible} (au lieu de localement faiblement filtrante). Ça signifie que pour tout $X'$ étale sur $X$ et 0-connexe, $X'$ est faiblement filtrante à gauche, i.e. est irréductible.

Reste à exprimer la condition filtrante elle-même en termes de $\operatorname{Top} X$. Je rappelle à ce propos :

\dgrblock{dgrb90}{Products of connected objects}
\begin{dcomment}{exposition}{p17:Reste à exprimer la condition filtrante elle-même}{Products of connected objects}
Grothendieck now asks for a condition on the topos which remembers the equalization of parallel arrows. Irreducibility alone cannot do this: it records whether two subobjects meet, whereas equalization concerns two maps with the same source. Proposition 2.3 first isolates the intermediate condition that the component functor preserve products. Throughout the following topos discussion, ``connected object'' means noninitial and connected, that is, $0$-connected; the topos is locally connected. Thus every object is the coproduct of its connected components, and $\pi _0$ is left adjoint to the constant-object functor.

Write $\mathcal E$ for the category of objects of the topos and let $S$ be the full generating subcategory of connected objects used in the text. The precise statement of 2.3(b) is
\[
 \pi _0(V\times U)\longrightarrow\pi _0(V)
 \quad\hbox{is bijective whenever $U$ is connected.}
\]
The edition prints $\pi _0(U)$ on the right, although the displayed projection has target $V$. This is a slip of the edition, retained in the source. Condition (a) implies (c) by decomposing $U$ and $V$ into connected components: products distribute over coproducts in a topos, and each resulting product is connected. Conversely, (c) applied to connected objects gives (a); the same decomposition proves the equivalence with (b).

The reduction to generators deserves the proof which the footnote requests. For $F\in\mathcal E$, let $S/F$ have objects $s\to F$, with $s\in S$. The canonical presentation
\[
 F\simeq\mathop{\rm colim}_{(s\to F)\in S/F}s
\]
holds in $\mathcal E$; it is the density presentation for the canonical site of generators. Applying $\pi _0$ gives $\pi _0(F)\simeq\pi _0(S/F)$, since every $s$ is connected. If $s\times t$ is connected for every $s,t\in S$, separate preservation of colimits by products gives
\[
 \pi _0(F\times G)
 \simeq\mathop{\rm colim}_{S/F\times S/G}\{*\}
 \simeq\pi _0(S/F)\times\pi _0(S/G).
\]
This proves (d)$\Rightarrow$(c). Equivalently, one may use the clopen-partition argument in our block on Lemma 2.9 below; this avoids any choice of a presentation.

Applying (c) to the terminal object shows that the diagonal of $\pi _0(1)$ is a bijection onto its square. Consequently $\pi _0(1)$ has at most one element. The nondegenerate case has exactly one, which explains the separate nonemptiness clause in Definition 2.4. Nothing in this argument annihilates higher homotopy groups. The inference to $W_\infty$-asphericity immediately following 2.3 is explicitly withdrawn in the manuscript's annotation; the stronger hereditary hypothesis introduced in 2.6 is what will support the cohomological argument.
\end{dcomment}


\noindent\textbf{Proposition 2.3.} \textit{Les conditions suivantes sur un topos $\sX$ sont équivalentes (quand $S$ (comme `site') désigne une partie génératrice de} $\mathrm{Fais}(\sX) = S^{\wedge}$, formée d'objets 0-connexes (\footnote[23]{donc $\sX$ est localement 0-connexe.}).\textit{)}
\begin{enumerate}[label={},leftmargin=3em]
\item[\textit{a)}] \textit{Si $U$, $V$ sont deux objets $0$-connexes de $\operatorname{Ob} \sX$, $U \times V$ est $0$-connexe.}
\item[\textit{b)}] \textit{Si $U$ est un objet $0$-connexe de $\sX^{\wedge}$, alors $U \longrightarrow e$ est `$W_0$-asphérique', i.e. pour tout $V$, la projection $V \times U \longrightarrow V$ induit $\pi_0(U \times V) \xrightarrow{\sim} \pi_0(U)$.}
\item[\textit{c)}] \textit{Si $U$, $V$} \texttt{[sont]} \textit{dans $\sX^{\wedge}$, alors}
\[
\pi_0(U \times V) \xrightarrow{\sim} \pi_0(U) \times \pi_0(V).
\]
\item[\textit{d)}] \textit{Comme a), avec $U$, $V$ dans $S$. (Mais on suppose les objets de $S$ $0$-connexes.)}
\end{enumerate}
\noindent(\footnote[24]{Il faudrait expliciter la démonstation avec celle de 2.9 ci-dessous (p.~25).}). % [sic] démonstation -> démonstration, letter r missing

\msp{page 18}

\noindent De plus, ces conditions impliquent que $\sX$ est connexe, donc 0-connexe si $\sX \neq \emptyset$. Cela implique de plus que $\sX$ est $W_{\infty}$-asphérique, cf. Pursuing Stacks §35 (Corollary p.~52) (\footnote[25]{\textbf{Faux,} loc. cit. ne dit rien de tel.}).

C'est essentiellement Proposition 1 de Pursuing Stacks §65, cf. loc. cit. p.~170-171 (\footnote[26]{Mais je m'y borne aux topos de la forme $\operatorname{Top}(A)$, avec $A$ dans Cat.}). Cf. aussi Pursuing Stacks, §35, pages 50-55.

\printedpage{12}
% head: A. Grothendieck : Dérivateurs, Ch. II, Édition des Universités de Montpellier II et Paris VII
% foot: Version 8 juillet 2001

\noindent\textbf{Définition 2.4.} On dit que $\sX$ est \textit{totalement $0$-connexe} s'il est localement connexe, non vide et satisfait les conditions équivalentes de 2.3. (NB c'est un cas particulier de la notion `totalement $W$-asphérique' de loc. cit. On a ainsi la notion de totalement $n$-connexe.)

Notons que pour un objet $X$ de Cat, dire que $\operatorname{Top}(X)$ est totalement 0-connexe revient à dire que $X \neq \emptyset$, et que pour $x, y \in \operatorname{Ob} X$, $x \times y$ dans $X^{\wedge}$ est 0-connexe, i.e. $X/x \times y$ est 0-connexe (\footnote[27]{Cela signifie aussi que $\mathrm{diag}_X : X \longrightarrow X \times X$ est co-cofinale, i.e. $X^o \longrightarrow X^o \times X^o$ est cofinale, et implique qu'il en est de même de $X \longrightarrow X^E$ pour tout ensemble \textit{fini} $E$.}). C'est une condition très forte. On dira aussi que $X$ est \textit{totalement $0$-connexe.} Notons ici

\dgrblock{dgrb91}{The ordinary sifted condition and its variance}
\begin{dcomment}{later}{p18:Notons que pour un objet}{The ordinary sifted condition and its variance}
The product condition has acquired a useful name, provided that we keep track of the opposite category. For $\widehat X=\operatorname{Fun}(X^{\mathrm{op}},\mathbf {Set})$, the left adjoint to the constant-presheaf functor is
\[
 \pi _0(F)=\mathop{\rm colim}_{X^{\mathrm{op}}}F
          =\pi _0(X/F).
\]
Thus $\widehat X$ is totally $0$-connected precisely when $X^{\mathrm{op}}$ is \emph{sifted}: its colimits in sets commute with finite products. In the ordinary categorical characterization, the indexing category is nonempty and its diagonal is final. For the conventions of these pages, this says that $X\to X\times X$ is co-cofinal. The comma category at $(x,y)$ consists of common cones $z\to x$, $z\to y$, and is exactly $X/(h_x\times h_y)$.

This characterization also explains why binary products suffice. Repeated use of finality of the diagonal gives finality of every finite nonempty diagonal; nonemptiness handles the nullary product. The terminology and its relation to algebraic theories are developed by \citet{AdamekRosickyVitale2010}. The distinction from filteredness is structural: filtered colimits preserve finite limits, while sifted colimits are required to preserve finite products. Equalizers are precisely the operation absent from the latter requirement.

There is a corresponding distinction in higher category theory, but one must not transfer a statement by changing its vocabulary alone. In the higher version the relevant comma objects must be weakly contractible, rather than merely nonempty and connected. The characterization of filtered colimits by preservation of finite limits in spaces is \citet[Proposition 5.3.3.3]{Lurie2009HTT}; sifted colimits and geometric realizations are treated in his \S5.5.8. Our present identification concerns \emph{ordinary} set-valued colimits. The higher homotopy obstruction which survives it is exhibited later in the commentary on 2.19--2.20.
\end{dcomment}


\noindent\textbf{Proposition 2.5.} \textit{Soit $\sX$ un topos. Si $\sX$ est totalement $0$-connexe, alors $\sX$ est irréductible. (En fait la conclusion vaut, tautologiquement, en supposant seulement que le produit de deux objets $0$-connexes de $\sX^{\wedge}$ est non vide, et il suffit même de le supposer pour des sous-objets de l'objet final $e_{\sX}$.) En particulier, si $\sX = \mathit{Top}(X)$, où $X$ est dans} Cat\textit{, alors $X$ est faiblement filtrante à gauche.}

Mais bien sûr, cela n'implique \textit{pas} que $\sX$ soit localement irréductible.

\msp{page 19}

\noindent Par exemple la catégorie $\boldsymbol{\Delta}$, ou toute petite catégorie stable par produits binaires, est totalement 0-connexe, mais $\boldsymbol{\Delta}$ n'est pas localement irréductible (les catégories induites $\boldsymbol{\Delta}/\boldsymbol{\Delta}^n$, pour $n > 0$, ne sont visiblement \textit{pas} irréductibles). (Mais rappelons que si $\sX = \operatorname{Top}(T)$, $T$ un espace topologique, alors $T$ irréductible implique que $\sX$ est aussi \textit{localement} irréductible. Donc $T$ totalement 0-connexe implique $\operatorname{Top}(T)$ globalement \textit{et} localement irréductible, i.e. pour tout $T'$ étale sur $T$, et 0-connexe, $T'$ est irréductible.)

\dgrblock{dgrb92}{The simplex category separates the two conditions}
\begin{dcomment}{examples}{p19:les catégories induites}{The simplex category separates the two conditions}
Grothendieck's example shows exactly where a global product condition loses the information needed locally. We spell it out because the word ``visibly'' hides two different computations. Let $\Delta$ have the nonempty finite ordinals $[n]=\{0,\ldots,n\}$ and monotone maps. For any small category $X$ with binary products, the product of two representables is representable; its category of elements has a terminal object. Such a nonempty $X$ is therefore totally $0$-connected. The separate assertion about $\Delta$ does not say that $\Delta$ has binary products: the ordered square $[1]\times[1]$ is not an ordinal.

For $\Delta$, the product $h_{[m]}\times h_{[n]}$ is the simplicial set $\Delta^m\times\Delta^n$. Its vertices are the pairs $(i,j)$. An edge joins $(i,j)$ to $(i+1,j)$ whenever the latter exists, and similarly joins $(i,j)$ to $(i,j+1)$. Hence its vertex graph is connected. Every simplex is joined in its category of elements to each of its vertices by a face map. The category of elements is consequently nonempty and connected, which verifies the generator test of 2.3 for all $m,n$.

Local irreducibility fails already over $[1]$. In $\Delta/[1]$, consider the two objects $[0]\xrightarrow{0}[1]$ and $[0]\xrightarrow{1}[1]$. A common predecessor would be a nonempty ordinal $[r]$ mapping to both copies of $[0]$, with equal composites to $[1]$. Those composites are respectively constant at $0$ and $1$, so equality is impossible. The two principal sieves have empty intersection. In particular, $\Delta/[1]$ is not irreducible, and $\Delta$ is not cofiltered. Equivalently, the parallel vertex maps $[0]\rightrightarrows[1]$ admit no equalizing map into $[0]$.

This example gives an explicit positive answer to the request in footnote 31: total $0$-connectedness does not imply local irreducibility. Its opposite also supplies an ordinary sifted category which is not filtered. Neither conclusion requires an argument about derived limits.
\end{dcomment}


\noindent\textbf{Théorème 2.6.} \textit{Soit $\sX$ un topos localement connexe, $S \subseteq \sX^{\wedge}$ une petite partie génératrice formée d'objets $0$-connexes. Conditions équivalentes (}sauf \textit{c, c$'$, qui semblent} (\footnote[28]{\textbf{sont.}}) \textit{plus faibles, cf. démonstration).)}
\begin{enumerate}[label={},leftmargin=3em]
\item[\textit{a)}] \textit{Tout topos induit $\sX/U$ sur $U$ $0$-connexe dans $\sX^{\wedge}$ est totalement $0$-connexe.}
\item[\textit{a$'$)}] \textit{Tout produit fibré $U \times_W V$ d'objets $0$-connexes dans $\sX^{\wedge}$ est $0$-connexe.}
\item[\textit{b)}] \textit{$\sX$ est localement irréductible, et ses composantes connexes sont totalement $0$-connexes.}
\item[\textit{c)}] \textit{Comme a), avec $U$ dans $S$, i.e. les $\sX/U$ ($U \in \operatorname{Ob} S$) sont totalement $0$-connexes.}
\item[\textit{c$'$)}] \textit{Comme a$'$), avec $U$, $V$, $W$ dans $S$.}
\item[\textit{d)}] \textit{$S$ est localement filtrante à gauche, i.e. $S$ satisfait les deux conditions}

PS$^*$ 1\ \ \textit{Tout diagramme}
\begin{tikzcd}[ampersand replacement=\&, row sep=tiny, column sep=small, baseline=(\tikzcdmatrixname.center)]
a \arrow[dr] \& \\
 \& x \\
b \arrow[ur] \&
\end{tikzcd}
\textit{dans $S$ s'insère dans}
\begin{tikzcd}[ampersand replacement=\&, row sep=tiny, column sep=small, baseline=(\tikzcdmatrixname.center)]
 \& a \arrow[dr] \& \\
c \arrow[ur, dotted] \arrow[dr, dotted] \& \& x \\
 \& b \arrow[ur] \&
\end{tikzcd}
\textit{commutatif.}

PS$^*$ 2\ \ \textit{Tout diagramme $a \rightrightarrows b$ s'insère dans un diagramme commutatif $c \longrightarrow a \rightrightarrows b$.}
\end{enumerate}
\noindent(\footnote[29]{NB On a équivalence de \tikz[baseline=(n.base)]\node[draw, rounded corners, inner sep=1.5pt](n){a), a$'$), b), d),}; qui impliquent c), c$'$), lesquels sont équivalents. Mais il semble (\footnotemark[30]) que c), c$'$) soient strictement plus faible. Cf. cependant les variantes e), e$'$) dans 2.12 (p.~28, 29).}). % [sic] plus faible -> plus faibles, adjective does not agree with the plural subject c), c')
\footnotetext[30]{Il s'avère. C'est bien vrai, c) et c$'$) n'impliquent \textit{pas} les autres conditions, cf. annotation marginale
% footnote 30 continues at the foot of printed page 13:
p.~28 (cas de $X = B_G$).}

\printedpage{13}
% head: A. Grothendieck : Dérivateurs, Ch. II, Édition des Universités de Montpellier II et Paris VII
% foot: Version 8 juillet 2001

\msp{page 20}

\noindent Ces conditions impliquent que les objets 0-connexes de $\sX^{\wedge}$ (en particulier les éléments de $S$) sont $\infty$-connexes.

\noindent\textbf{Corollaire 2.7} (de l'équivalence b $\Longleftrightarrow$ d). \textit{Soit $X$ un objet de} Cat. \textit{Pour que $X$ soit filtrant à gauche il faut et il suffit que $\mathit{Top}(X)$ soit totalement $0$-connexe et localement (et donc aussi globalement) irréductible (ou encore que $X$ soit totalement $0$-connexe et que pour tout $x \in \operatorname{Ob} X$, la catégorie localisée $X/x$ soit irréductible (elle est même totalement $0$-connexe), ou encore que pour deux objets $a$, $b$ dans $X$, $a \times b \in X^{\wedge}$ soit $0$-connexe, et pour}
\begin{tikzcd}[ampersand replacement=\&, row sep=tiny, column sep=small, baseline=(\tikzcdmatrixname.center)]
a \arrow[dr] \& \\
 \& x \\
b \arrow[ur] \&
\end{tikzcd}
\textit{dans $X$, $a \times_x b$ (dans $X^{\wedge}$) soit $0$-connexe)} (\footnote[31]{Exemple de $X$ totalement 0-connexe, mais non cofiltrant, i.e. non localement irréductible ? La catégorie $\boldsymbol{\Delta}$ des simplexes type est totalement 0-connexe, mais n'est pas cofiltrante car $\boldsymbol{\Delta}^0 = e \underset{\partial_1}{\overset{\partial_0}{\rightrightarrows}} \boldsymbol{\Delta}^1$ ne peut être égalisé ! Donc \texttt{[la catégorie opposée]} $\boldsymbol{\Delta}^o = I$ est telle que $I \longrightarrow I^E$ soit cofinal pour tout ensemble fini $E$, mais $I$ non filtrante.}).

\noindent\textsc{Démonstration} de 2.6. On a le diagramme d'implications $\pm$ tautologiques
\[
\begin{tikzcd}[ampersand replacement=\&, column sep=large, row sep=large]
 \& \mathrm{a} \arrow[r, Leftrightarrow] \arrow[dl, Rightarrow, "{\substack{\text{par}\\ \text{2.5}}}"'] \arrow[d, Rightarrow] \& \mathrm{a}' \arrow[d, Rightarrow] \\
\mathrm{b} \arrow[d, Rightarrow] \& \mathrm{c} \arrow[r, Leftrightarrow, "{(\Leftarrow \text{ par 2.3.})}"] \& \mathrm{c}' \arrow[dll, Rightarrow] \\
\text{PS 1.} \& \&
\end{tikzcd}
\]
Montrons quand même les deux implications c$'$) $\Longrightarrow$ PS 1 et b $\Longrightarrow$ PS 1.

\dgrblock{dgrb93}{What the hereditary condition adds}
\begin{dcomment}{exposition}{p20:Montrons quand même les deux implications}{What the hereditary condition adds}
Theorem 2.6 makes the product condition stable under passage to a connected object as base. Its surviving equivalence is
\[
 (a)\ \Longleftrightarrow\ (a')\ \Longleftrightarrow\ (b)
 \ \Longleftrightarrow\ (d),
 \qquad (c)\Longleftrightarrow(c'),
\]
with an implication from the first group to the second. Here the three objects $U,V,W$ in (a') are all connected. Connectedness of $W$ is essential: even in sets, two singleton objects mapping to different elements of a two-element set have empty pullback.

The equivalence of (a) and (a') is the same product test inside $\mathcal E/W$. Condition (a) implies local irreducibility by 2.5 and gives the assertion about components in (b). To see the decisive step (b)$\Rightarrow$(d), first use 2.8 for PS 1. Given $u,v:a\rightrightarrows b$ in $S$, both objects lie over the same connected component of $1$. By (b), $b\times b$ is connected. The equalizer is the pullback
\[
 E=a\times_{b\times b}b,
 \qquad a\longrightarrow b\times b=(u,v),\quad
 b\longrightarrow b\times b=\operatorname{diag}.
\]
Local irreducibility, applied over the connected object $b\times b$, makes $E$ noninitial. A generator mapping to $E$ supplies an equalizing arrow, which is PS 2. The text's ``c) est vérifié'' at manuscript p.~24 means (b); its reference to $\mathcal X\times\mathcal X$ should be read as an application of 2.8 in $\mathcal X$ itself.

There is a concise exactness interpretation after restricting to any connected component of the topos. Under (a'), decomposition into connected components gives, for arbitrary $U,V,W$ in that component,
\[
 \pi _0(U\times_WV)
 \simeq\pi _0(U)\times_{\pi _0(W)}\pi _0(V).
\]
Indeed, the summands of the pullback are precisely the products $U_i\times_{W_k}V_j$ for which both factors map to the same component $W_k$; each is connected by (a'). Since $\pi _0(1)=\{*\}$, $\pi _0$ preserves finite limits on this connected topos. We shall use this stronger fact in the cohomological argument. For a disconnected topos one applies it component by component; its global $\pi _0$ does not preserve the terminal object.
\end{dcomment}


\msp{page 21}

\noindent C'est contenu dans le

\noindent\textbf{Corollaire 2.8.} \textit{Soient $\sX$, $S$ comme dans 2.6. Conditions équivalentes :}
\begin{enumerate}[label={},leftmargin=3em]
\item[\textit{a)}] \textit{$\sX$ est localement irréductible, i.e. les topos induits $\sX/U$, où $U$ est $0$-connexe, sont irréductibles.}
\item[\textit{b)}] \textit{Tout produit fibré $U \times_W V$ d'objets $0$-connexes $U$, $V$ est `non vide'.}
\item[\textit{b$'$)}] \textit{Tout produit fibré $U \times_W V$, avec $U$, $V$ `non vide', $W$ $0$-connexe, est non vide.}
\item[\textit{c)}] \textit{Les $\sX/x$ avec $x \in S$ sont irréductibles.}
\item[\textit{d)}] \textit{Comme b), avec $U$, $V$, $W$ dans $S$.}
\item[\textit{e)}] \textit{$S$ est localement faiblement filtrante à gauche, i.e. satisfait PS$^*$ 1.}
\end{enumerate}
\noindent Ici les implications $\pm$ évidentes sont

% ---- from Der-ChII-234_p016-020.tex
\printedpage{14}
% head: A. Grothendieck : Dérivateurs, Ch. II, Édition des Universités de Montpellier II et Paris VII
% foot: Version 8 juillet 2001

\[
\begin{tikzcd}
\mathrm{b}' \arrow[r, Rightarrow] & \mathrm{b} \arrow[rr, Rightarrow] \arrow[d, Rightarrow] & & \mathrm{a} \arrow[d, Rightarrow] \\
 & \mathrm{d} \arrow[r, Leftrightarrow] & \mathrm{e} \arrow[r, Leftrightarrow] & \mathrm{c}.
\end{tikzcd}
\]

Prouvons les implications de la ligne du bas. Considérons d). Si $\begin{tikzcd}[cramped, row sep=0.1em, column sep=0.8em, baseline=(current bounding box.center)] a \arrow[dr] & \\ & x \\ b \arrow[ur] & \end{tikzcd}$ \texttt{[est]} dans $S$, dire que le produit fibré dans $\X$ est `non vide' signifie, puisque $S$ est un système de générateurs et en particulier \textit{non vide,} qu'il existe $c$ dans $S$ et $c\longrightarrow a\times_x b$, ou ce qui revient au même, $\begin{tikzcd}[cramped, row sep=0.1em, column sep=0.8em, baseline=(current bounding box.center)] & a \\ c \arrow[ur, dotted] \arrow[dr, dotted] & \\ & b \end{tikzcd}$, rendant commutatif le diagramme $\begin{tikzcd}[cramped, row sep=0.1em, column sep=0.8em, baseline=(current bounding box.center)] & a \arrow[dr] & \\ c \arrow[ur, dotted] \arrow[dr, dotted] & & x \\ & b \arrow[ur] & \end{tikzcd}$. Donc d) équivaut à PS 1, i.e. à e) ($\X$ localement faiblement filtrant à gauche). Considérons la condition c) : les $\X/a=(S/a)^\sim$ sont irréductibles. Cela signifie que pour deux sous-objets \textit{non vides} $U,V\subseteq a$ dans $\X^\wedge$, on a $U\cap V\neq\emptyset$, ou encore que pour deux cribles

\msp{page 22}$\underline{U}$, $\underline{V}$ de $S/x$, si $\underline{U}$, $\underline{V}$ non vides (ce qui signifie que les sous-objets de $x$ dans $\X^\wedge$ qu'ils définissent sont `non vides', ou que les objets de $S$ sont `non vides' dans $\X^\wedge$), alors $\underline{U}\cap\underline{V}$ est non vide. On peut supposer que $\underline{U}$, $\underline{V}$ est engendré chacun par une seule flèche $a\xrightarrow{u}x$, $b\xrightarrow{v}x$. Le crible intersection est formé des objets $c\longrightarrow x$ de $S/x$ tels que $c\longrightarrow x$ se factorise par un $\begin{tikzcd}[cramped, row sep=0.1em, column sep=0.8em, baseline=(current bounding box.center)] & a \\ c \arrow[ur, dotted] \arrow[dr, dotted] & \\ & b \end{tikzcd}$, donc dire que ce crible n'est pas vide, signifie que PS 1 est satisfait pour les deux flèches données $u$, $v$. Donc c) équivaut à la condition PS 1\texttt{[, i.e. à]} e.

Ainsi on a établi
\[
\mathrm{b}'\Longrightarrow\mathrm{b}\Longrightarrow\mathrm{a}\Longrightarrow(\mathrm{c}\Longleftrightarrow\mathrm{d}\Longleftrightarrow\mathrm{e}),
\]
reste à prouver que c $\Longrightarrow$ b$'$. Soit donc $\begin{tikzcd}[cramped, row sep=0.1em, column sep=0.8em, baseline=(current bounding box.center)] U \arrow[dr] & & V \arrow[dl] \\ & W & \end{tikzcd}$ un diagramme d'objets `non vides', avec $W$ 0-connexe, il faut prouver que le produit fibré est non vide. Considérant $a\longrightarrow U$ et $b\longrightarrow V$ avec $a,b\in S$, on est ramené au cas où $U=a$, $V=b$ dans $S$. Prouvons que si $a\times_W b$ est `vide', alors $W$ est disconnexe. Pour ceci, considérons la sous-catégorie $C$ de $S/W$ formée des $x\longrightarrow W$ tels que $a\times_W x$ soit `non vide'. Désignant par $w(a)$, $w(x)$ l'image de $a$, $x$ dans $W$, cela signifie que $w(a)\cap w(x)\neq\emptyset$.

\msp{page 23}Notons que tout quotient dans $\X^\wedge$ d'un objet irréductible est irréductible, donc $w(a)$, $w(x)$ sont des objets irréductibles. Cela montre que pour $x$ dans $C$ et $x'$ dans $S/W$, on a $x'$ dans $C$, i.e. $w(a)\cap w(x')\neq\emptyset$ si et seulement si $w(x)\cap w(x')\neq\emptyset$. (Lemme direct : dans un objet $W$ d'un topos $\X$, si on a \texttt{[des]} sous-objets $W_1$, $W_2$, $W_3$ irréductibles, si l'un des deux \texttt{[plutôt trois]} rencontre les deux autres, ceux-ci se rencontrent.) Cela montre que $C$ est un crible, puisque $x'\longrightarrow x$ implique $w(x')\subseteq w(x)$, donc $w(x')$ rencontre $w(x)$, donc aussi $w(a)$. Mais alors le co-crible $C'=(S/W)\backslash C$ est également un crible, car si $c$
\printedpage{15}
% head: A. Grothendieck : Dérivateurs, Ch. II, Édition des Universités de Montpellier II et Paris VII
% foot: Version 8 juillet 2001
est dans $C'$, i.e. \texttt{[si]} $w(c)$ ne rencontre pas $w(x)$, a fortiori pas tout $x'$ au-dessus de $x$, pour lequel on aura $w(x')\subseteq w(x)$. Ainsi, $S/W$ est décomposé en somme de deux ouverts $C$, $C'$, non vides puisque l'un contient $a$, l'autre $b$. Par suite, $W$ est décomposé en somme de deux sous-objets non vides, donc $W$ est non 0-connexe.

On a prouvé 2.8, donc aussi (revenant à 2.6) les implications vers PS 1 dans le diagramme page 20. Il reste à prouver
\[
\left\{
\begin{aligned}
&\mathrm{b}\Longrightarrow\mathrm{PS}\ 2,\ \mathrm{c}'\Longrightarrow\mathrm{PS}\ 2\\
&\mathrm{PS}\ 1+\mathrm{PS}\ 2\ (\text{i.e. d})\Longrightarrow\mathrm{a}'
\end{aligned}
\right.
\]
(\footnote[32]{NB finalement on ne prouve pas c$'\Longrightarrow$ PS 2.}).

\dgrblock{dgrb94}{How irreducible generators control an arbitrary slice}
\begin{dcomment}{exposition}{p23:On a prouvé 2.8, donc aussi}{How irreducible generators control an arbitrary slice}
The proof of 2.8 explains why irreducibility, although weaker than filteredness, is genuinely local in the topos. The missing hypothesis in its printed condition (b) is that $W$ be connected. The assertion used in the proof is (b'): noninitial $U$ and $V$ over a connected $W$ have noninitial pullback.

Here is the passage from the generating objects to such a $W$. Suppose every $s\in S$ is irreducible, and consider the images $W_s\subseteq W$ of all maps $s\to W$. These images cover $W$ and remain irreducible: inverse images of two nonzero subobjects under an epimorphism are nonzero, and their nonzero intersection has nonzero image. On this covering family, declare two images related when they meet. The relation is reflexive and symmetric. It is transitive because, if $A$ and $C$ both meet an irreducible $B$, the two nonzero subobjects $A\cap B$ and $B\cap C$ of $B$ meet. Thus the equivalence classes give pairwise disjoint unions of subobjects covering $W$.

These unions are a coproduct decomposition in the topos. Since $W$ is connected, there is only one class. Every pair of generating images therefore meets. Choose a generator mapping into each of the noninitial objects $U$ and $V$. Their images in $W$ meet, and epimorphisms remain epimorphisms after pullback; consequently $U\times_WV$ is noninitial. This proves the step (c)$\Rightarrow$(b') without assuming that every object of a site is representable as a presheaf on that site.

Finally, the condition that $s\times_x t$ be noninitial is equivalent to the existence of a generator $r$ mapping to that pullback. Such a map is exactly a commutative square completing $s\to x\leftarrow t$, namely PS 1. This proves the equivalence with (d) and (e). The implication (b')$\Rightarrow$(a) follows by taking two nonzero subobjects of a connected object. The remaining implications are restrictions of this same test. Grothendieck's argument is therefore sound after restoring the condition on the base; the diagram of implications is not concealing an additional homotopical assumption.
\end{dcomment}


\msp{page 24}Nous nous plaçons donc dans l'hypothèse où PS 1 est satisfait (puisque b), c$'$) et d) impliquent PS 1), donc où on est sous les conditions équivalentes de 2.8 - $\X$ est localement irréductible. Considérons alors la condition PS 2, et prouvons qu'elle est impliquée tant par b), que par c$'$), et qu'elle implique a$'$). Elle signifie que pour une double flèche $a\overset{u}{\underset{v}{\rightrightarrows}}b$ dans $S$, le conoyau \texttt{[plutôt noyau]} dans $S^\wedge$, ou, ce qui revient au même, dans $\X^\wedge=S^\sim\subseteq S^\wedge$, est `non vide'. On peut l'écrire sous la forme d'un produit fibré
\[
\begin{tikzcd}
a \arrow[r, "{\scriptstyle (u,v)}"] & b\times b \\
\operatorname{Ker}(u,v) \arrow[u] \arrow[r] & b. \arrow[u, "\mathrm{diag}"']
\end{tikzcd}
\]
Ici par hypothèse $a$ et $b$ sont des objets 0-connexes de $\X^\wedge$ (ceux de $S$ l'étant), c) est vérifié (($\X$ localement irréductible et) ses composantes connexes \texttt{[sont]} totalement 0-connexes), alors $b\times b$ est 0-connexe comme produit de deux objets 0-connexes sur une \textit{même composante connexe,} laquelle est totalement 0-connexe. Comme $\X\times\X$ est localement irréductible, il s'ensuit que le produit fibré est non vide (2.8), donc b) $\Longrightarrow$ d). Prouvons d) $\Longrightarrow$ a), i.e. si

\msp{page 25}$S$ est localement cofiltrante, alors les topos induits $\X/W$, où $W$ est 0-connexe, sont totalement 0-connexes. Ce topos induit est décrit par le site $S/W$, i.e. $S/W$ est un système de générateurs. Or si $S$ est localement filtrante, $S/W$ (étale sur $S$) l'est aussi. Donc on peut supposer $W=e$, et à prouver que si $S$ est localement cofiltrante et $\X$ 0-connexe, alors $\X$ est totalement 0-connexe. Notons que $\X$ 0-connexe équivaut à $S$ 0-connexe - il faudrait quand même l'expliciter :

\textbf{Lemme 2.9.} \textit{Soit $\X$ un topos, $S\subseteq\X^\wedge$ une sous-catégorie pleine génératrice formée d'objets} 0-connexes. \textit{Alors $\X$ est $0$-connexe si et seulement si $S$ l'est - plus généralement, l'ensemble ordonné des sous-objets à la fois ouverts et fermés de $S$ dans} Cat \textit{(ou encore des
\printedpage{16}
% head: A. Grothendieck : Dérivateurs, Ch. II, Édition des Universités de Montpellier II et Paris VII
% foot: Version 8 juillet 2001
sous-objets ouverts $=$ cribles de $S$ qui sont aussi fermés, i.e. des cocribles) est isomorphe à celui des sommandes directs de $\X$.}

Je l'ai utilisé tacitement dans la démonstration de 2.8, pour le topos induit $\X/W$. Je n'explicite pas la démonstration.

\dgrblock{dgrb95}{Connectedness in a category of generators}
\begin{dcomment}{exposition}{p25:Je n'explicite pas la démonstration}{Connectedness in a category of generators}
Lemma 2.9 is what permits the categorical and topos-theoretic arguments to meet. Its content is stronger than the statement that a generating family covers the terminal object: a partition of the generators must descend to an actual coproduct decomposition of that object.

Let $S\subseteq\mathcal E$ be full, generating, and formed of connected objects. A decomposition $1=A\amalg B$ assigns every $s\in S$ entirely to one summand, since the two pullbacks decompose $s$. Every arrow preserves this assignment. It thus gives a partition $S=S_A\amalg S_B$ into full subcategories with no arrows between them. Conversely, suppose such a partition is given. Let $A$ and $B$ be the unions in $1$ of the images of the objects in the respective parts. They cover $1$, because $S$ generates. If $A\cap B$ were noninitial, some generator would map into it. Pulling back the epimorphic families defining $A$ and $B$, and choosing a generator in a noninitial member, would produce an object $t\in S$ with maps to an object of $S_A$ and one of $S_B$. This contradicts the partition. Hence $A\cap B=0$, and $A\amalg B\simeq1$.

The two constructions are inverse: a connected generator whose image lies in $A$ cannot belong to the other part, and the images of all generators assigned to $A$ jointly cover $A$. The same argument works with any set of parts and proves
\[
 \pi _0(S)\simeq\pi _0(1_{\mathcal E}).
\]
It also applies to $S/F$, which is a full generating subcategory of $\mathcal E/F$. An object $s\to F$ is connected in the slice exactly when its domain $s$ is connected, since coproducts in a slice are computed on domains. We obtain the repeatedly used identity
\[
 \pi _0(S/F)\simeq\pi _0(F).
\]
This supplies the proof omitted at manuscript p.~25 and justifies the earlier use of the lemma in 2.8. It also explains why the category of elements, rather than only the set of objects of a site, controls connectedness.
\end{dcomment}


Ainsi (revenant à la démonstration de d) $\Longrightarrow$ c) dans 2.6.) $S$ étant 0-connexe et localement cofiltrante,

\msp{page 26}\texttt{[elle]} est en fait cofiltrante. Donc pour deux objets $a$, $b$ de $S$, $\exists\ c$ avec $c\longrightarrow a$, $c\longrightarrow b$ (condition PS 1$'$, plus forte que PS 1). Pour prouver que $\X$ est totalement 0-connexe, on est réduit à prouver par 2.3 que les $a\times b$ dans $\X^\wedge$ sont 0-connexes. Par 2.9, il revient au même de prouver qu'ils sont 0-connexes dans $S^\wedge$. Donc on est ramené à prouver le

\textbf{Lemme 2.10.} \textit{Soit $S$ une catégorie cofiltrante. Si $a$, $b$ sont deux objets de $S$, $a\times b$ est un objet} $0$\textit{-connexe de $S^\wedge$.}

On voit de suite qu'il suffit de prouver (et en fait il revient au même, compte tenu que $a\times b$ est localement irréductible) que si on a deux flèches
\[
\begin{tikzcd}
x \arrow[r] & a\times b \\
 & y, \arrow[u]
\end{tikzcd}
\]
alors $x\times_{a\times b}y\neq\emptyset$, i.e. $\exists\ z$ dans $S$ rendant commutatif
\[
\begin{tikzcd}
x \arrow[r] & a\times b \\
z \arrow[u] \arrow[r] & y. \arrow[u]
\end{tikzcd}
\]
Or par PS 1$'$ il existe bien un $z$ s'envoyant dans $x$ et dans $y$, mais les deux flèches composées
\[
z\rightrightarrows a\times b,
\]
i.e. chacune des double flèches
\[
z\rightrightarrows a \tag*{$(*)$}
\]
\[
z\rightrightarrows b \tag*{$(*')$}
\]
peuvent n'être pas égales. Mais par PS 2,

\msp{page 27}on peut trouver
\[
z'\longrightarrow z
\]
\printedpage{17}
% head: A. Grothendieck : Dérivateurs, Ch. II, Édition des Universités de Montpellier II et Paris VII
% foot: Version 8 juillet 2001
égalisant la double flèche $(*)$, donc on peut supposer que celle-ci est déjà triviale. Appliquant encore une fois PS 2, par un $z''\longrightarrow z'$, on égalise de même la double flèche $(*')$, et on gagne. On a donc prouvé, finalement,
\[
\begin{tikzcd}
\mathrm{a} \arrow[r, Leftrightarrow] \arrow[d, Rightarrow] & \mathrm{a}' \arrow[r, Rightarrow] & \mathrm{b} \arrow[r, Leftrightarrow] & \mathrm{d}\ (=\mathrm{PS}\ 1+\mathrm{PS}\ 2) \\
\mathrm{c} \arrow[r, Leftrightarrow] & \mathrm{c}',
\end{tikzcd}
\]
mais pas l'implication inverse. En d'autres termes, si les topos induits $\X/a$, pour $a$ dans $S$, sont totalement 0-connexes, je n'ai pas prouvé qu'il en est de même des $\X/W$, $W$ un objet 0-connexe \textit{quelconque} de $\X^\wedge$. Ainsi, on en est au point suivant

\dgrblock{dgrb96}{Where both equalizations are used}
\begin{dcomment}{exposition}{p27:mais pas l'implication inverse}{Where both equalizations are used}
The completion of 2.6 is a useful instance of the finite-cone argument underlying filteredness. It is not enough to find a common predecessor of two objects: one must also arrange that the resulting maps agree in both factors. This is why the proof of Lemma 2.10 invokes PS 2 twice.

Suppose $S$ is cofiltered. An object of $S/(h_a\times h_b)$ is a triple $(x,r:x\to a,s:x\to b)$. It exists by the common-predecessor axiom. Given two such triples with domains $x$ and $y$, choose $z\to x$ and $z\to y$. Equalize the two resulting maps $z\rightrightarrows a$ by a map $z'\to z$, then equalize the two maps $z'\rightrightarrows b$ by $z''\to z'$. The second operation preserves the first equality. The object with domain $z''$ is a common predecessor in $S/(h_a\times h_b)$. This proves its connectedness, as required in 2.10. In fact it is cofiltered: parallel arrows in it can be equalized after projection to $S$, and the resulting arrow still belongs to the category of elements.

Now let $S$ merely be locally cofiltered, and let $W$ be connected in the topos. The projection $S/W\to S$ is a discrete fibration. PS 1 lifts along a discrete fibration: once a common predecessor is chosen downstairs, lifting one leg determines the object upstairs, and equality of the two composites gives the other leg. PS 2 lifts for the same reason, using uniqueness of lifts. Thus $S/W$ is locally cofiltered. Lemma 2.9 says it is connected, so it is cofiltered. Applying the preceding argument in this generating category and then the generator test of 2.3 shows that $\mathcal E/W$ is totally $0$-connected. This is (d)$\Rightarrow$(a).

The proof establishes the announced equivalence of (a), (a'), (b), and (d). It does not establish (c')$\Rightarrow$PS 2, and the annotation correctly removes that assertion. Restricting the base of a pullback to a representable object has lost information: the object $b\times b$ required for the diagonal argument need not be among the chosen generators. The following group example makes this loss explicit.
\end{dcomment}


\textbf{Perplexité 2.11.} \textit{Soit $X$ dans} Cat. \textit{Si $X$ est localement filtrante à gauche, alors les produits fibrés $a\times_x b$ d'objets dans $X$ sont des objets} $0$\textit{-connexes de $X^\wedge$, et cette condition à son tour implique PS 1, i.e. que $X$ est localement irréductible. Mais en présence de cette hypothèse PS 1, ou $X$ localement irréductible, la condition sur les produits fibrés $a\times_x b$ équivaut à la suivante : pour tout diagramme} \texttt{[en]} \textit{traits pleins}

\msp{page 28}
\[
\begin{tikzcd}
c \arrow[r] & a\times_x b \\
e \arrow[u, dotted] \arrow[r, dotted] & d, \arrow[u]
\end{tikzcd}
\qquad\qquad
\textit{avec $x$, $a$, $b$, $c$, $d$ dans $X$,}
\]
\textit{on peut compléter en un diagramme} commutatif, \textit{avec $e$ dans $X$. En d'autre termes, tout diagramme commutatif à traits pleins ci-dessous se complète en y rajoutant un `objet initial' $e$} % [sic] d'autre termes -> d'autres termes, missing plural s
\[
\begin{tikzcd}[column sep=small, row sep=small]
 & & a \arrow[ddr] & \\
 & c \arrow[ur] \arrow[dddr] & & \\
e \arrow[ur, dotted] \arrow[dr, dotted] & & & x \\
 & d \arrow[uuur] \arrow[dr] & & \\
 & & b \arrow[uur] &
\end{tikzcd}
\]
(\footnote[33]{Il suffit d'exiger pour $e$ la commutativité des deux carrés $(e\ c\ b\ d)$ et $(e\ d\ a\ c)$.}).

Cette condition est donc impliquée par $X$ localement cofiltrante, et elle est satisfaite aussi par ce qui précède plus généralement si et seulement si (\footnote[34]{si et seulement si ?}) les $X/x$ sont cofiltrantes (\footnote[35]{Elle signifie que $X$ satisfait PS 1, et la condition

\textbf{PS$'$ 2} Tout diagramme commutatif
\[
a\rightrightarrows b\longrightarrow x
\]
% footnote continued on printed page 18
s'insère dans un diagramme commutatif
\[
y\longrightarrow a\rightrightarrows b\longrightarrow x.
\]
}). Mais la réciproque est-elle vraie - cette condition de totale 0-connexité des $X/x$ implique-t-elle que $X$ soit totalement 0-connexe ? Mais je n'ai pas construit d'exemple où les $X/x$
\printedpage{18}
% head: A. Grothendieck : Dérivateurs, Ch. II, Édition des Universités de Montpellier II et Paris VII
% foot: Version 8 juillet 2001
sont cofiltrantes sans que $X$ soit localement cofiltrante (ou ce qui revient au même, de $X$ 0-connexe avec les $X/x$ cofiltrantes sans que $X$ en soit cofiltrante)(\footnote[36]{Mais il y a un \textit{contre-exemple} évident, si $X$ est un groupoïde 0-connexe, les $X/x$ sont $\approx e$, donc cofiltrantes, mais $X$ n'est cofiltrant (i.e. \texttt{[mot(s) illisible(s)]}) que si $X$ est 1-connexe. P.ex. si $C$ est un groupoïde $\neq 1$, alors $X=B_G$ est tel que les $X/x$ sont cofiltrantes, mais $X$ non cofiltrant. % [sic] si $C$ est un groupoïde -> si $G$ est un groupoïde, the groupoid is then written $B_G$
}) ; ni d'exemple où les $X/x$ sont totalement 0-connexes, sans que les $X/x$ soient cofiltrantes. Mais notons le

\dgrblock{dgrb97}{A groupoid invisible to its representable slices}
\begin{dcomment}{examples}{p28:Mais la réciproque est-elle vraie}{A groupoid invisible to its representable slices}
The marginal answer to Perplexity 2.11 is decisive, and it also explains the weaker condition that remains. Let $G$ be a nontrivial group and let $BG$ be its one-object category. The objects of $BG/*$ are the elements $g\in G$, viewed as arrows to the unique object. For any $g,h$, there is exactly one arrow from $g$ to $h$, namely the element $k$ satisfying $hk=g$. Hence $BG/*$ is a nonempty indiscrete groupoid, equivalent to the terminal category, and is cofiltered.

The category $BG$ itself is not cofiltered. If $g\ne h$ are parallel arrows, an equalization $gk=hk$ would imply $g=h$ because $k$ is invertible. Its presheaf topos is the category of right $G$-sets. The representable is the regular $G$-set $G$; it is connected. Its square with diagonal action has
\[
 \pi _0(G\times G)\simeq G,
 \qquad [(a,b)]\longmapsto ba^{-1},
\]
after choosing the corresponding convention for the right action. Thus the square is not connected. This verifies directly that the topos is locally irreducible but is not totally $0$-connected, although the slice over the representable is equivalent to sets. The image of the regular object onto the terminal $G$-set is exactly the quotient missed by the generator-only test.

There is, however, no further example in which \emph{every} $X/x$ is totally $0$-connected but some $X/x$ is not cofiltered. Indeed, total $0$-connectedness of all $X/x$ makes all these categories irreducible, so $X$ is locally irreducible. For $y\to x$, the slice $(X/x)/(y\to x)$ is isomorphic to $X/y$, and is therefore irreducible. Hence $X/x$ is both totally $0$-connected and locally irreducible. Corollary 2.7, whose proof has just been completed, makes it cofiltered.

In elementary terms, the surviving weaker condition is PS 1 together with
\[
 a\mathrel{\substack{\xrightarrow{u}\\[-.5ex]\xrightarrow[v]{}}}b
 \xrightarrow{h}x,\quad hu=hv
 \quad\Longrightarrow\quad
 \exists\,t:c\to a\ \text{with }ut=vt.
\]
This is PS$'$ 2 in footnote 35. It equalizes parallel arrows only when they already become equal after a common outgoing arrow. In a groupoid that antecedent forces equality, so the condition is automatic. PS 2 has no such antecedent. The two questions in the paragraph are therefore settled differently: the first has the groupoid counterexample already supplied by Grothendieck; the second has a negative existence answer from his preceding theorem.
\end{dcomment}


\textbf{Corollaire 2.12 :} \textit{Les conditions équivalentes a), a$'$), b), d) (à l'exclusion de c), c$'$), plus faibles apparemment) du théorème 2.6}

\msp{page 29}\textit{équivalent encore à la suivante}
\begin{itemize}[leftmargin=2.5em]
\item[\textit{e)}] \textit{$\forall\ x$ dans $S$, le topos induit $\X/x$ et le topos induit $\X/x\times x$ est irréductible} (\footnote[37]{NB $\X/x$ inutile.}),
\end{itemize}
\textit{où encore à la suivante} % [sic] où encore -> ou encore, conjunction "ou" without accent
\begin{itemize}[leftmargin=2.5em]
\item[\textit{e$'$)}] \textit{$\forall\ x$ dans $S$, $S^\wedge/x$ et $S^\wedge/x{\times}x$ sont irréductibles, i.e. $S/x$ et $S/x{\times}x$ sont irréductibles} (\footnote[38]{$S/x$ inutile, car $x\xrightarrow{\mathrm{diag}}x\times x$ est une immersion ouverte non vide, donc si $x\times x$ est irréductible, $x$ l'est aussi. Mais il revient au même de dire que les $S/x$ sont irréductibles (condition PS$^*$ 1, cf. 2.8, p. 21) \textit{et} que les $x\xrightarrow{\mathrm{diag}}x\times x$ sont dominants (ce que implique en effet que $x\times x$ est irréductible). Ceci redonne immédiatement la forme PS$^*$ 2 de la condition cofiltrante. % [sic] ce que implique -> ce qui implique, relative pronoun subject
}).
\end{itemize}
C'est la forme `topologique' la plus faible en apparence que je peux trouver pour la condition $S$ localement cofiltrante, i.e. $\X$ (ou Top$(S)$, au choix) localement totalement 0-connexe.

\dgrblock{dgrb98}{Why the square of one generator detects equalizers}
\begin{dcomment}{exposition}{p29:C'est la forme `topologique' la plus faible}{Why the square of one generator detects equalizers}
Corollary 2.12 is not merely another presentation of a long list of equivalent conditions. It isolates the diagonal as the part of the product test which forces equalization. We can read its criterion entirely in the topos, without treating the chosen site as a space.

Suppose that $s\times s$ is irreducible for every $s\in S$. The diagonal $s\hookrightarrow s\times s$ is a noninitial subobject. Every noninitial subobject of an irreducible object is irreducible, so each $s$ is irreducible and 2.8 gives PS 1. For two maps $u,v:a\rightrightarrows b$ in $S$, the image of $(u,v):a\to b\times b$ is noninitial. It meets the diagonal, because $b\times b$ is irreducible. Pullback stability of epimorphisms then says that
\[
 \operatorname{Eq}(u,v)=a\times_{b\times b}b
\]
is noninitial. A generator mapping into this equalizer supplies PS 2. Conversely, PS 1 and PS 2 give local cofilteredness; the category of generators over $b\times b$ is cofiltered by the proof of 2.10, so this product is irreducible.

This also proves the comment that the separate hypothesis on $\mathcal E/s$ is redundant. The diagonal is a nonzero open subobject in the language of slice topoi, and irreducibility passes to such an open. Saying instead that $s$ is irreducible and its diagonal is dense gives the same test: density means that the diagonal meets the image of every noninitial object mapping to $s\times s$. The pair $(u,v)$ then gives the equalizer test just used. The relevant ``open immersion'' is the morphism of topoi associated with a monomorphism in $\mathcal E$; it is not an assertion that an arbitrary topological diagonal is open.

Definition 2.13 packages these equivalent conditions as local total $0$-connectedness. The generator category is allowed to vary: once the condition is expressed for all connected slices of the topos, the proof of 2.6 applies to every full generating subcategory of connected objects.
\end{dcomment}


\textbf{Définition 2.13.} On dit qu'un topos $\X$ est \textit{localement totalement} $0$\textit{-connexe} s'il satisfait aux conditions équivalentes a, a$'$, b, d du théorème 2.6, i.e. si $\X$ est localement connexe et si de plus tout topos induit 0-connexe est totalement 0-connexe - le produit fibré de deux objets 0-connexes sur un troisième est 0-connexe. Il revient au même de dire que toute (ou que une) sous-catégorie pleine génératrice $S$ de $\X^\wedge$ formée d'objets 0-connexes est localement cofiltrante (en supposant toujours $\X$ localement connexe, ou ce qui revient au même, l'existence d'un $S$ formé d'objets 0-connexes).

On n'a pas terminé encore la démonstration de 2.6 - il reste à prouver que si $\X$ (supposé localement totalement 0-connexe) est 0-connexe, il est $\infty$-connexe - ou encore que les composantes connexes de $\X$, supposé localement totalement 0-connexe, sont $\infty$-connexes, i.e. $W_\infty$-asphériques.

\msp{page 30}\textbf{Corollaire 2.14.} \textit{Notons quand même que si $\X$ est de la forme Top$(T)$, $T$ un espace topologique, alors $\X$ est localement totalement} $0$\textit{-connexe si et seulement si il est localement irréductible, i.e. si et seulement si $T$ est localement irréductible.}

% ---- from Der-ChII-234_p021-025.tex
\printedpage{19}
% head: A. Grothendieck : Dérivateurs, Ch. II, Édition des Universités de Montpellier II et Paris VII
% foot: Version 8 juillet 2001

On est ramené à prouver que cette dernière condition implique la condition d) de 2.6, et on peut prendre pour $S$ l'ensemble ordonné des ouverts connexes non vides de $S$, cela devient alors évident.

Ceci dit, il faut se rappeler que la théorie cohomologique des espaces topologiques irréductibles n'est nullement triviale - comme nous le montre FAC de \textsc{Serre}\,! D'autre part, c'est \textsc{Serre} qui m'avait fait observer que pour des coefficients \textit{constants,} la cohomologie d'un espace irréductible est triviale.

\dgrblock{dgrb99}{Irreducible spaces still have nonconstant cohomology}
\begin{dcomment}{examples}{p30:la cohomologie d'un espace irréductible est triviale}{Irreducible spaces still have nonconstant cohomology}
The reference to Serre separates two coefficient questions which must remain distinct. On an irreducible nonempty topological space $T$, every nonempty open is connected. For an abelian group $M$, the constant sheaf $\underline M$ consequently has sections $M$ on every nonempty open, with identity restriction maps. Restrictions to the empty open are also surjective. Thus $\underline M$ is flabby and
\[
 H^q(T,\underline M)=0\quad(q>0),\qquad
 H^0(T,\underline M)=M.
\]
The same argument on each irreducible component applies when the space is a disjoint union of irreducible open components, which is the local condition used in 2.14.

For a concrete contrast, take the Zariski space $T=\mathbf P^1_k$, which is irreducible, and the line bundle $\mathcal O(-2)$. On the standard affine charts with coordinate $t$ and $t^{-1}$, its degree-one \v{C}ech quotient is
\[
 \frac{k[t,t^{-1}]}{k[t]+t^{-2}k[t^{-1}]}
   =k\,[t^{-1}].
\]
The first summand removes all nonnegative powers and the second all powers at most $-2$; exactly the coefficient of $t^{-1}$ remains. The affine charts and their intersection are acyclic for this coherent sheaf, so the quotient is $H^1(\mathbf P^1_k,\mathcal O(-2))\simeq k$. The affine-cover comparison used here is \citet[no.~47, Theorem 4]{Serre1955FAC}; the displayed computation gives the particular example in full.

For 2.14 itself, choose the connected nonempty opens as generators. If the space is locally irreducible in Grothendieck's sense, each such open is irreducible. Two nonempty opens inside it meet, giving PS 1; the generator category is a poset, so parallel arrows are already equal and PS 2 imposes nothing further. The converse follows from local total $0$-connectedness implying local irreducibility. This is why the two notions coincide for spaces while the group example separates them for general topoi.
\end{dcomment}


Dans le cas actuel, prouvons déjà, si $\mathcal{X}$ est 0-connexe et localement totalement 0-connexe, que tout revêtement de l'objet final de $\mathcal{X}$ est trivial. Il suffit de voir que si $U$ est 0-connexe, $U \longrightarrow e$ est un isomorphisme. Mais considérons $U \hookrightarrow U \times U$, on sait que c'est un sommande direct, et comme $U \times U$ est connexe, c'est un isomorphisme. Donc $U \longrightarrow e$ \texttt{[est un]} monomorphisme, donc sommande direct, et comme $U \neq \emptyset$, on a $U \longrightarrow e$ isomorphisme.

Reste à prouver que pour tout faisceau \msp{page 31} abélien constant $F$ sur $\mathcal{X}$, on a $H^i(X,F)=0$ pour $i>0$. On procède par récurrence sur $i$, c'est déjà fait pour $i=1$. Supposons que ce soit prouvé pour les $j<i$ ($i$ fixé) pour tout topos \textit{totalement} 0\textit{-connexe.} Soit $\xi \in H^i(\mathcal{X},F)$, soit $U \longrightarrow e$ \texttt{[un]} épimorphisme tel que $\xi|U=0$, et considérons la suite spectrale % [sic] H^i(X,F) -> H^i(\mathcal{X},F), plain X printed where the topos \mathcal{X} is meant elsewhere
\[
E_2^{pq} = H^p(C(U^*,\mathcal{H}^q(F))) \Longrightarrow H^*(\mathcal{X},F).
\]
On a par hypothèse totalement 0-connexe
\[
\pi_0(U^q) \simeq \pi_0(U)^q, \tag*{$(*)$}
\]
et les composantes connexes de $U^q$ sont totalement 0-asphériques,
\begin{center}
\begin{tikzpicture}[scale=0.8]
\draw[->] (0,0) -- (0,5);
\node[left] at (0,5) {$q$};
\draw[->] (0,0) -- (5,0);
\node at (4.3,-0.6) {$p$};
\node[left] at (0,4) {$i$};
\foreach \x in {0,1,2,3,4} {\fill (\x,4) circle (2pt); \fill (\x,0) circle (2pt);}
\end{tikzpicture}
\end{center}
donc par hypothèse de récurrence
\[
E_1^{pq} = 0 \qquad \text{si } 0<q<i,
\]
donc on trouve une suite exacte
\[
0 \longrightarrow H^i(C(U^*,F)) \longrightarrow H^i(\mathcal{X},F) \longrightarrow H^0(\ \underbrace{H^i(U^*,F)}_{\substack{\text{complexe}\\ \text{co-semisimplicial}\\ \text{pour } i\\ \text{variable}}}\ ) = \Ker(H^i(U,F) \rightrightarrows H^i(U^2,F)).
\]

\dgrblock{dgrb100}{The hereditary hypothesis in the cohomological proof}
\begin{dcomment}{exposition}{p31:Supposons que ce soit prouvé pour les}{The hereditary hypothesis in the cohomological proof}
Grothendieck's induction proves the intended assertion after keeping its hypothesis fixed: the topos must be connected and \emph{locally} totally $0$-connected. The word ``localement'' disappears in the sentence introducing the induction at manuscript p.~31, but it is needed when the argument passes to components of powers of a covering object. Those components define connected slices of the original topos, and hence remain in the same class by 2.6. Mere total $0$-connectedness does not provide this inheritance.

First consider a connected locally constant object $L$. Its diagonal in $L\times L$ is complemented, as one checks after a cover trivializing $L$. The product is connected, so this noninitial summand is the whole product. Hence $L\to1$ is monic. Its image is also complemented, by the same local triviality; connectedness of the base and nonemptiness of $L$ give $L\simeq1$. Every covering object is therefore a coproduct of copies of $1$. In the source paragraph, $U$ must still denote a connected component of a \emph{covering}, not an arbitrary connected sheaf.

For the higher cohomology, assume the vanishing in degrees $0<q<n$ for every connected topos in this hereditary class. Given $\xi\in H^n(\mathcal E,\underline M)$, choose an epimorphism $U\to1$ killing $\xi$. Such a cover exists because a positive-degree cohomology class represented in an injective resolution becomes a boundary locally. The augmented \v{C}ech nerve gives the descent spectral sequence
\[
 E_1^{pq}=H^q(\mathcal E/U^{p+1},\underline M)
 \ \Longrightarrow\ H^{p+q}(\mathcal E,\underline M).
\]
Each $U^{p+1}$ is a coproduct of connected objects whose slices remain in the class; thus the rows $0<q<n$ vanish. The edge sequence gives
\[
 0\longrightarrow H^n(C^\bullet)
 \longrightarrow H^n(\mathcal E,\underline M)
 \longrightarrow \ker\bigl(H^n(U,\underline M)
       \rightrightarrows H^n(U^2,\underline M)\bigr).
\]
The last image of $\xi$ is zero by choice of $U$. If $I=\pi _0(U)$, product preservation identifies
\[
 C^p=\operatorname{Map}(I^{p+1},M).
\]
Choose $i_*\in I$, which exists since $U\to1$ is epic. Inserting $i_*$ as first coordinate defines
$(sc)(i_0,\ldots,i_{p-1})=c(i_*,i_0,\ldots,i_{p-1})$.
For the alternating deletion differential, direct cancellation gives $ds+sd=1$ in positive degrees, with the usual augmentation in degree zero. Thus $H^n(C^\bullet)=0$ and $\xi=0$. The ``trivial'' complex in the manuscript is this contractible augmented complex; its individual groups and differentials are not zero.
\end{dcomment}


\printedpage{20}
% head: A. Grothendieck : Dérivateurs, Ch. II, Édition des Universités de Montpellier II et Paris VII
% foot: Version 8 juillet 2001

Par hypothèse sur $U$, l'image de $\xi$ dans le troisième terme est nulle, donc $\xi$ appartient au premier terme. Mais pour $(*)$, le complexe $C(U^*,F)$ est le complexe \textit{trivial} des puissances de $I=\pi_0(U)$, à coefficients dans $M$ (si $F=M_{\mathcal{X}}$). Sa cohomologie \msp{page 32} en degrés $>0$ est nulle, donc $\xi=0$, q.e.d.

\dgrblock{dgrb101}{From connected components to the hypercover homotopy type}
\begin{dcomment}{later}{p32:en degrés $>0$ est nulle}{From connected components to the hypercover homotopy type}
The stronger homotopical conclusion of 2.6 can also be proved directly, without deducing weak contractibility from an unspecified version of a cohomological Whitehead theorem. On a connected locally totally $0$-connected topos, the component functor preserves finite limits, as shown above, and preserves epimorphisms: epimorphisms in a topos are regular, and the left adjoint $\pi _0$ preserves their coequalizer presentations. Let $U_\bullet\to1$ be a hypercover by objects of this topos. Its matching objects are finite limits. Therefore
\[
 \pi _0\bigl(M_n(U_\bullet)\bigr)
       \simeq M_n\bigl(\pi _0(U_\bullet)\bigr),
\]
and the hypercover epimorphisms become surjections of sets. The simplicial set $\pi _0(U_\bullet)$ consequently has a filler for every boundary $\partial\Delta^n\to\pi _0(U_\bullet)$, including the nonemptiness condition in degree zero. Its map to a point is a trivial Kan fibration, and it is weakly contractible.

In the hypercover construction of the homotopy type, whose classical setting is \citet[\S\S8--9]{ArtinMazur1969}, every stage is therefore weakly contractible. This proves the $W_\infty$-asphericity meant here, as well as the corresponding assertion for each connected object by passage to its slice. For a presheaf topos it says concretely that every connected category of elements is weakly contractible.

There is also a short algebraic verification of the constant-coefficient calculation. Since $\pi _0$ preserves finite limits, it induces an exact functor on abelian group objects, left adjoint to the constant-abelian-sheaf functor. The latter therefore preserves injectives. Applying it to an injective resolution of $M$, and then taking global sections, recovers the same resolution of $M$, because the topos is connected. Thus $R\Gamma(\mathcal E,\underline M)\simeq M[0]$. These proofs make explicit why the hereditary condition repairs the false inference after 2.3. They impose no claim that the topos itself is equivalent to sets; the preceding projective-line example shows that arbitrary coefficients retain substantial information.
\end{dcomment}


\noindent\textbf{Proposition 2.15.} (\footnote[39]{C'est déjà dans SGA 4, I 8.}) \textit{Soit $X$ un objet de} Cat\textit{. Conditions équivalentes :}
\begin{enumerate}[label={}, leftmargin=2em, itemsep=0pt]
\item \textit{a) $X$ est filtrant, i.e. $X$ est $0$-connexe et $X^o$ satisfait les conditions équivalentes du théorème 2.6 ($X^{o,\wedge}$ est localement totalement $0$-connexe).}
\item \textit{b) L'objet final de $X^\wedge$ est ind-représentable.}
\end{enumerate}
\noindent\textsc{Démonstration}. On sait qu'on a
\[
e_{X^\wedge} \simeq \varinjlim_{x \in X} x, \tag*{(2.15.1)}
\]
limite inductive prise dans $X^\wedge$. En effet, cela signifie aussi que pour tout objet $a$ dans $X$, on a
\[
\varinjlim_{x \in X} \ \underbrace{x(a)}_{\Hom_X(a,x)} \ \xrightarrow{\ \sim\ } e \qquad \text{(ensemble ponctuel)}
\]
Mais considérons l'objet $X' = a\backslash X \longrightarrow X$ coétale sur $X$, qui est donc Cat-cofibré sur $X$, on sait que l'on a
\[
\pi_0(X') \simeq \varinjlim_{x \in X} \pi_0(X'_x),
\]
où $X'_x = \Hom(a,x)$, donc on trouve
\[
\varinjlim_{x \in X} \Hom(a,x) = \pi_0(a\backslash X) = e
\]
car $a\backslash X$ est 0-connexe.

\dgrblock{dgrb102}{The terminal presheaf as an ind-object}
\begin{dcomment}{exposition}{p32:C'est déjà dans SGA 4}{The terminal presheaf as an ind-object}
The passage to ind-objects replaces the cone conditions by a representability statement. The general form, recorded in 2.18, is the familiar criterion from \citet[Expos\'e I, Theorem 8.3.3]{SGA4I1972}: a presheaf $F$ on a small category $X$ is a filtered colimit of representables if and only if its category of elements $X/F$ is filtered. We give the argument, since it also supplies all the details of 2.15.

For every presheaf there is the density presentation
\[
 F\simeq\mathop{\rm colim}_{(x,s)\in X/F}h_x.
\]
Objectwise, the element represented by a pair $a\to x$, $s\in F(x)$ maps to its pullback in $F(a)$; it has the representative $(a,t)$ with identity map, and the relations identify every other representative with this one. If $X/F$ is filtered, this is the required ind-presentation.

Conversely, write $F=\operatorname{colim}_{i\in I}h_{x_i}$ with $I$ filtered. The category of elements is nonempty, since the image of $1_{x_i}$ supplies an object. Two elements $s\in F(a)$ and $t\in F(b)$ have representatives at stages $i,j$; passage to a common stage $k$ supplies arrows from $(a,s)$ and $(b,t)$ to $(x_k,\overline{1}_{x_k})$. Given parallel arrows $u,v:(a,s)\rightrightarrows(b,t)$, represent $t$ by $b\to x_i$. The two composites $a\rightrightarrows x_i$ have the same image in $F(a)$. Equality in a filtered colimit of sets is witnessed at a common later stage. The resulting map $(b,t)\to(x_j,\overline{1}_{x_j})$ equalizes $u,v$. Thus $X/F$ is filtered.

For $F=1$, its category of elements is $X$ itself, giving 2.15. The formula $1=\operatorname{colim}_{x\in X}h_x$ is valid even when $X$ is not filtered: at $a$ its value is $\pi _0(a\backslash X)=\{*\}$ because $(a,1_a)$ is initial. The condition lies in the possibility of a \emph{filtered} presentation. The reindexing by a directed poset used in the manuscript is Deligne's result, \citet[Expos\'e I, Proposition 8.1.6]{SGA4I1972}; the proof above does not need this extra reindexing.
\end{dcomment}


La formule (2.15.1) montre que si \msp{page 33} $X$ est filtrant, $e_{X^\wedge}$ est limite inductive filtrante de foncteurs représentables. On sait d'ailleurs qu'on peut réindexer par une catégorie \textit{ordonnée} filtrante (cf. SGA 4, I 8.1.6 (c'est dû à P. \textsc{Deligne}\,!)).

On a prouvé a $\Longrightarrow$ b, prouvons b $\Longrightarrow$ a. Supposons donc
\[
e_{X^\wedge} = \varinjlim_{i \in I} x_i,
\]
où $I$ est un (petit) ensemble ordonné filtrant, et $I \longrightarrow X$ un foncteur. Prouvons PS 1$'$ (forme droite), i.e. pour $a, b$ dans $X$, \texttt{[il]} existe $\vcenter{\hbox{\begin{tikzcd}[row sep=0.1em, column sep=1.2em]
a \arrow[dr] & \\
& x \\
b \arrow[ur] &
\end{tikzcd}}}$. En effet, comme $\varinjlim\limits_{I} \Hom(a,x_i) = e$, il s'ensuit que $\Hom(a,x_i) \neq \emptyset$ pour $i$ grand, de même $\Hom(b,x_i) \neq \emptyset$, donc il suffit de prendre $\vcenter{\hbox{\begin{tikzcd}[row sep=0.1em, column sep=1.2em]
a \arrow[dr] & \\
& x_i \\
b \arrow[ur] &
\end{tikzcd}}}$. Prouvons PS 2$_{\mathrm{d}}$, i.e. que pour \texttt{[une]} double flèche $a \overset{u}{\underset{v}{\rightrightarrows}} b$, \texttt{[il]} existe $b \longrightarrow x$ qui égalise.

\printedpage{21}
% head: A. Grothendieck : Dérivateurs, Ch. II, Édition des Universités de Montpellier II et Paris VII
% foot: Version 8 juillet 2001

\msp{page 34}
Soit $\xi_{i_0}$ dans $\Hom(b,x_{i_0})$, ses deux images dans $\Hom(a,x_{i_0})$ ne sont pas nécessairement égales, mais le deviennent dans la limite inductive, réduite à $e$, donc deviennent égales en composant avec $x_{i_0} \longrightarrow x_i$ pour $i$ assez grand, donc le composé $b \longrightarrow x_{i_0} \longrightarrow x_i$
\[
a \rightrightarrows b \longrightarrow x_{i_0} \longrightarrow x_i
\]
égalise la double flèche.

\noindent\textbf{Corollaire 2.16.} \textit{Soit $W$ un localiseur fondamental satisfaisant l'axiome des limites} Loc (8)\textit{. Si $X$ dans} Cat \textit{est une catégorie filtrante ou cofiltrante, alors $X$ est $W$-asphérique.}

Plus généralement et agréablement :

\dgrblock{dgrb103}{Filtered categories and the minimal localizer}
\begin{dcomment}{later}{p34:Plus généralement et agréablement}{Filtered categories and the minimal localizer}
Corollary 2.16 needs only the standard contractibility of a filtered category, together with the comparison between its weak equivalences and the chosen localizer. Here is the finite-data argument behind the first assertion. In a filtered category, any finite collection of objects has maps into one object. Given finitely many arrows between those objects, successive applications of the parallel-arrow axiom modify the common target so that all the required triangles commute. Thus every finite diagram of objects and arrows occurring in a finite simplicial subset of the nerve admits a compatible cone.

Every map from a finite sphere into the realization of the nerve can, after simplicial approximation, be represented in a finite such subset. The cone makes that map nullhomotopic. The nerve is nonempty and connected, so all its homotopy groups vanish. Reversing the category does not change the homotopy type of its nerve, and gives the cofiltered assertion as well.

The manuscript formulates the conclusion for a fundamental localizer satisfying its axiom Loc(8). A later theorem permits a sharper modern statement: the nerve weak equivalences $W_\infty$ form the smallest fundamental localizer. This is \citet[Theorem 2.2.11, p.~136]{Cisinski2004Minimal}. With the standard modern definition of fundamental localizer, every filtered or cofiltered category is consequently $W$-aspherical without an additional filtered-limit axiom. This is a later strengthening of the hypothesis in the manuscript, rather than a reason to alter its historical formulation. The finite-cone argument also underlies the higher categorical characterization cited earlier: filtered colimits of spaces commute with finite limits.
\end{dcomment}


\noindent\textbf{Corollaire 2.17.} \textit{Soit $W$ comme dans 2.16. Alors toute catégorie filtrante est} totalement $W$-asphérique (\footnote[40]{cf. XII, §V pour cette notion, que j'ai beaucoup peiné à dégager et à comprendre\,!}), \textit{i.e. pour toute flèche $f : X \longrightarrow Y$ dans} Cat\textit{, l'objet $f_!(e_X)$ de $e_Y$ (défini aussi comme $y \longmapsto \pi_0(y\backslash X)$, et intervenant dans la factorisation canonique de $f$ en $ip$, avec $p$ étale et les cofibres} \texttt{[?]} \textit{de $p$ $0$-connexes)} \msp{page 35} \textit{est $W$-asphérique. En fait, $Y/f_!(e_{X^\wedge})$ est même une catégorie filtrante.}

En effet, le foncteur $f_!$ commute aux limites inductives, donc $f_! e_{X^\wedge}$ est limite inductive filtrante des $f_!(x_i) = f(x_i)$. Donc $Y/f_!(e_{X^\wedge})$ est limite inductive filtrante des $Y/f(x_i)$. Donc c'est même une catégorie filtrante, a fortiori elle est $W$-asphérique.

Donc finalement, on trouve des exemples non triviaux de catégories totalement $W$-asphériques, i.e. des exemples donc des catégories sans objet final. % [sic] donc -> dont (?), the conjunction "donc" makes no sense here

\dgrblock{dgrb104}{Left Kan extension carries the ind-presentation}
\begin{dcomment}{exposition}{p35:Donc finalement, on trouve des exemples non triviaux}{Left Kan extension carries the ind-presentation}
The reason behind 2.17 is that an ordinary left Kan extension sends a representable to a representable and preserves colimits. If $f:X\to Y$ and $X$ is filtered, then
\[
 f_!1_X\simeq\mathop{\rm colim}_{x\in X}f_!h_x
          \simeq\mathop{\rm colim}_{x\in X}h_{f(x)}.
\]
The result is an ind-representable presheaf on $Y$, so $Y/(f_!1_X)$ is filtered by the criterion just proved. This gives the precise argument for the last assertion of 2.17 and therefore its asphericity conclusion. The proof does not require $f$ to be fully faithful or cofinal.

The canonical factorization mentioned in the statement has the following form:
\[
 X\xrightarrow{p}Y/(f_!1_X)\xrightarrow{i}Y.
\]
Here $i$ is the discrete fibration associated with the presheaf, and
\[
 (f_!1_X)(y)=\pi _0(y\backslash X).
\]
The first map sends $x$ to the component represented by $(x,1_{f(x)})$. For an object $(y,c)$ of the category of elements, the comma category $(y,c)\backslash X$ is exactly the component $c$ of $y\backslash X$; it is nonempty and connected. Thus $p$ is cofinal in the terminology of these pages. The edition assigns ``\'etale'' to $p$, but it is the second map $i$ which is \'etale; the connected cofibres belong to $p$. The phrase ``objet $\cdots$ de $e_Y$'' similarly means an object of $\widehat Y$.

For the simplest example without a terminal object, take $X=Y=\mathbf N$ with its usual order and $f=1_X$. Its terminal presheaf is the increasing colimit of the representables, and its category of elements is again $\mathbf N$. This is filtered and weakly contractible, although no natural number is terminal. It supplies exactly the kind of example for which Grothendieck introduces the ind-object formulation. Corollary 2.18 is the same argument with an arbitrary presheaf in place of the terminal one.
\end{dcomment}


\noindent\textbf{Corollaire 2.18.} (\footnote[41]{déjà dans SGA 4, I 8.}) \textit{Soit $X$ dans} Cat\textit{, $F$ dans $X^\wedge$. Pour que $F$ soit ind-représentable il faut et il suffit que $X/F$ soit filtrante.}

En effet, par (2.15), $X/F$ est filtrante si et seulement si l'objet final $F$ de $X/F$ provient d'un ind-objet de $X/F$, mais cela signifie aussi que $F$ est la limite inductive dans $X^\wedge$ d'un ind-objet de $X$, i.e. que $F$ est ind-représentable.

\msp{page 36}
\noindent\textbf{Question.} \textit{Si $X$ dans} Cat \textit{est totalement $W_\infty$-asphérique, i.e. les $f_!(e_{X^\wedge})$ sont $W_\infty$-asphériques pour toute flèche $f : X \longrightarrow Y$ dans} Cat\textit{, ou encore pour toute telle flèche, avec les $y\backslash Y$ $0$-connexes ($y \in \Ob Y$), $Y$ est $W_\infty$-asphérique - cela implique-t-il que $X$ soit filtrante ?} Ça ne m'étonnerait pas - et même en supposant seulement les $f_!(e_{X^\wedge})$ 1-connexes (\footnote[42]{Ils le sont toujours (cf. XII, p. 101), c'est 2-\textit{connexe} qu'il faut lire ici\,!}) (ils sont 0-connexes, dès que $X$ l'est). Mais je ne veux pas m'attarder à présent sur cette question.

\dgrblock{dgrb105}{The quantifier in total asphericity}
\begin{dcomment}{exposition}{p36:Mais je ne veux pas m'attarder à présent sur cette question.}{The quantifier in total asphericity}
Grothendieck is asking which indexing categories remain homotopically harmless after an ordinary change of index. The adjective ``total'' concerns every possible target of a functor; it does not merely strengthen the contractibility of the nerve of the original category. It is useful to isolate the precise condition before reading the proposed answer. Write $B A=|N A|$, and write $T(A)$ for the assertion that $B E$ is weakly contractible whenever $A\longrightarrow E$ is an ordinary cofinal functor. Here and below ``cofinal'' has the orientation fixed in Section~3: each category $e\backslash A$ must be nonempty and connected.

This formulation is equivalent to the one using $f_!1_A$ on the page. Indeed, for $f:A\longrightarrow Y$, put $P=f_!1_A$, so that
\[
 P(y)=\pi_0(y\backslash A).
\]
Let $Y/P$ be its category of elements. There is a canonical factorization
\[
 A\xrightarrow{q}Y/P\longrightarrow Y,
 \qquad a\longmapsto(fa,[a,1_{fa}]).
\]
For an object $(y,c)$ of $Y/P$, the comma category $(y,c)\backslash A$ is precisely the connected component $c$ of $y\backslash A$. Thus $q$ is cofinal. Conversely, if $f$ itself is cofinal, then $P=1_Y$ and $Y/P=Y$. Consequently $T(A)$ says exactly that the category of elements of every ordinary presheaf $f_!1_A$ is weakly contractible. The occurrence of $y\backslash Y$ in the printed question is a slip: it is $y\backslash A$ that must be connected.

A filtered category satisfies $T$. Its ordinary cofinal images are filtered by Corollary~3.4 below, and their nerves are weakly contractible. The difficult direction is whether this test forces the equalization of parallel arrows. The counterexample inserted on pages~36bis--36ter withdraws the attempted negative answer; it does not prove the converse.

There is a terminological precaution for the later literature. In \citet[Definition~1.7.2 and Proposition~1.7.1]{Maltsiniotis2020}, a category is called totally aspherical when it is aspherical and its diagonal is aspherical, equivalently when products of representable presheaves are aspherical. This is a different quantifier from $T(A)$. In particular, the fact that the simplex category is a strict test category does not settle the question on this page. We retain $T(A)$ only inside this commentary, to keep the question itself unchanged.
\end{dcomment}


La réponse à la question précédente est négative, comme on voit par la :

\printedpage{22}
% head: A. Grothendieck : Dérivateurs, Ch. II, Édition des Universités de Montpellier II et Paris VII
% foot: Version 8 juillet 2001

\noindent\textbf{Proposition 2.19.} (\footnote[43]{\textbf{Faux.}}) \textit{Soit $X$ dans} Cat \textit{telle que $\mathit{diag}_X : X \longrightarrow X \times X$ soit cofinal} (\footnote[44]{cf. définition 3.3, p. 45.})\textit{, i.e. telle que $X^o$ soit totalement $0$-connexe. Alors $X$ est totalement $W_\infty$-asphérique (donc aussi totalement $W$-asphérique pour tout localiseur fondamental $W$ tel que $W \supseteq W_\infty$).}

Soit en effet un foncteur cofinal $f : X \longrightarrow Y$, il faut prouver que $Y$ est $W_\infty$-asphérique. Or on sait que totalement 0-connexe implique $W_\infty$-asphérique (\footnote[45]{C'est \textbf{faux.}}) (cf. proposition 2.3, p. 17-18), donc il suffit de voir que $Y^o$ est totalement 0-connexe, i.e. $Y \xrightarrow{d_Y} Y \times Y$ cofinal. Or on a le diagramme commutatif
\[
\begin{tikzcd}[row sep=large, column sep=large]
X \arrow[r, "d_X"] \arrow[d, "f"'] & X \times X \arrow[d, "f \times f"] \\
Y \arrow[r, "d_Y"] & Y \times Y,
\end{tikzcd}
\]
où par hypothèse $f$, donc aussi $f \times f$, et $d_X$ sont cofinaux. Donc $\varphi = (f \times f) \circ p_X$ \texttt{[plutôt} $\varphi = (f \times f) \circ d_X$\texttt{]} est cofinal, et comme $\varphi = d_Y f$ et $f$ est aussi cofinal, on conclut que $d_Y$ l'est, q.e.d.

\noindent\textbf{Corollaire.} $\boldsymbol{\Delta}$ \textit{étant totalement $0$-connexe, $\boldsymbol{\Delta}^o$ est totalement $W_\infty$-asphérique, pourtant elle n'est pas filtrante.}

\noindent\textbf{Question.} \textit{$X$ totalement $W_\infty$-asphérique $\Longrightarrow$ $X^o$ totalement $0$-connexe ? Ou du moins implique PS 1$'$ ?}

\msp{page 36bis}
\noindent\textbf{Contrexemple 2.20} à la proposition fausse 2.19.

Soit $X$ un objet de Cat, rajoutons-lui la catégorie finale $e$ comme fermé, i.e. considérons un objet $X'$ de Cat dans lequel $X$ soit un sous-objet ouvert, $e$ le fermé complémentaire :
\[
X \overset{i}{\hookrightarrow} X' \overset{j}{\hookleftarrow} e.
\]
On sait (cf. XII) qu'une telle opération est déterminée, à isomorphisme près, par la donnée d'un objet
\[
H \in \Ob X^\wedge
\]
en posant
\[
H(x) = \Hom_{X'}(x,e),
\]
la catégorie $X/e$ s'identifiant (à isomorphisme près) à $X/H$. Elle est donc 0-connexe si et seulement si l'objet $H$ de $X^\wedge$ est 0-connexe. Cela signifie aussi que le morphisme $i$ est cocofinal, i.e. que $i^o : X^o \longrightarrow X'^o$ est cofinal. Supposons maintenant $X$ $W_\infty$-asphérique et totalement 0-connexe, i.e. l'actuel Loc (6), cf. XVI. La proposition 2.19 affirme que, sous ces conditions, $X'$ (qui est lui aussi totalement 0-connexe) est même $W$-asphérique. Il s'ensuivrait donc que $i \in W_\infty$. Mais par XII, th. 4, cor. 4 (p. 149), il s'ensuivrait que
\printedpage{23}
% head: A. Grothendieck : Dérivateurs, Ch. II, Édition des Universités de Montpellier II et Paris VII
% foot: Version 8 juillet 2001
\msp{page 36ter}
$X/e$ est $W_\infty$-asphérique, i.e. $H$ $W_\infty$-asphérique. Ainsi, si $X$ est totalement 0-asphérique, tout objet 0-connexe de $X^\wedge$ serait $\infty$-connexe. Mais c'est manifestement faux, p.ex. si on prend pour $X$ une catégorie-test pour $W_\infty$, p.ex. $X = \boldsymbol{\Delta}$.

\dgrblock{dgrb106}{The retracted proposition and a visible sphere}
\begin{dcomment}{examples}{p36:p.ex. si on prend pour $X$ une catégorie-test}{The retracted proposition and a visible sphere}
The useful part of Proposition~2.19 is its elementary calculation with cofinal functors: the ordinary cofinal image of a nonempty sifted category is sifted. The unsupported step is the passage from that ordinary connectedness condition to weak contractibility. The edition marks this step and the proposition as false. The appeal in Counterexample~2.20 to Chapter~XII also needs its hypothesis restored: the reflection of a weak equivalence to the attaching category uses the strong form of $W(7\mathrm{bis})$, explicitly required in \citet[Chapter~XII, Section~V.3.4, Corollary~4, note~104]{Grothendieck1990XII}. That reflection is not valid for arbitrary weak homotopy equivalences. We can make the intended counterexample entirely explicit without using it.

The omitted reflection condition has a concrete obstruction. Choose a finite connected acyclic noncontractible simplicial complex $K$; a finite presentation complex with this homotopy type is given in \citet[Example~2.38 and Theorem~2C.5]{Hatcher2002}. Its suspension is simply connected and acyclic, hence contractible by the homology Whitehead theorem. The face posets of the two cones on $K$ are open subcategories covering the face poset of $\Sigma K$, with intersection the face poset of $K$. The two cones and their union are weakly contractible, but their intersection is not. Thus the weak equivalence from one cone into the union does not reflect to a weak equivalence from the intersection into the other cone. This contradicts the strong biconditional $W(7\mathrm{bis})$, while leaving the direct homotopy-pushout theorem intact.

Let $H=\Delta[1]/\partial\Delta[1]$, the simplicial circle, and form the category $E$ with full subcategory $\Delta$, one additional object $z$, and
\[
 E([n],z)=H_n,\qquad E(z,[n])=\varnothing,
 \qquad E(z,z)=\{1_z\}.
\]
Composition into $z$ is the simplicial action of $H$. Thus $\Delta/z=\Delta/H$, the category of simplices of $H$. Its classifying space has the homotopy type of $|H|=S^1$; this is the usual category-of-simplices realization, or the test-category equivalence for $\Delta$ \citep[Proposition~1.7.14]{Maltsiniotis2020}. Since $\Delta$ has terminal object $[0]$, its classifying space is contractible.

There is a direct description of $B E$. An $n$-simplex of $N E$ that does not lie in $N\Delta$ consists of a chain in $\Delta$ followed by an arrow into $z$, with possible repetitions of $z$. The arrow into $z$ and all its composites are exactly a simplex of $N(\Delta/H)$. These simplices glue to the mapping cone of the forgetful map $B(\Delta/H)\longrightarrow B\Delta$. Hence
\[
 B E\simeq\operatorname{hocolim}\bigl(B\Delta\longleftarrow
 B(\Delta/H)\longrightarrow *\bigr)\simeq\Sigma S^1\simeq S^2.
\]
The cone description can also be read degree by degree: adjoining the last vertex $z$ supplies the cone direction, and deleting that vertex applies the forgetful map. It therefore identifies the attaching maps, not merely the vertices of the two spaces.

The inclusion $i:\Delta\hookrightarrow E$ is ordinary initial: $\Delta/e$ has a terminal object when $e\in\Delta$, while $\Delta/z=\Delta/H$ is nonempty and connected. Thus $i^{\mathrm{op}}:\Delta^{\mathrm{op}}\longrightarrow E^{\mathrm{op}}$ is cofinal. The source $\Delta^{\mathrm{op}}$ is sifted: for $[m],[n]$, its category of common targets is opposite to $\Delta/(\Delta[m]\times\Delta[n])$, which is connected because the product of the two simplices is contractible. The target $E^{\mathrm{op}}$ is consequently sifted as well, whereas $B(E^{\mathrm{op}})\simeq S^2$. This supplies both a counterexample to the implication used in Proposition~2.19 and a direct counterexample to its asserted corollary about $\Delta^{\mathrm{op}}$.

The distinction is the same one encountered earlier in the chapter between products of connected objects and their higher homotopy. Here its failure has a specific witness, a two-dimensional sphere. It is not a failure of the composition or cancellation properties of ordinary cofinal functors.
\end{dcomment}


Cela me fait soupçonner que la conjecture
\[
X \text{ totalement } W_\infty\text{-asphérique} \Longrightarrow X \text{ filtrante}
\]
pourrait être valable.

\dgrblock{dgrb107}{A necessary directedness condition}
\begin{dcomment}{development}{p36:pourrait être valable.}{A necessary directedness condition}
Grothendieck asks, before returning to cofinality, whether his condition at least forces the existence of common successors. This weaker question admits an elementary answer. The argument below also settles his full question for preorders. We did not find this particular implication isolated in the consulted treatments of sifted categories, asphericity, or lax epimorphisms; the following is our continuation of the argument, with no assertion of priority.

\textbf{Proposition (the common-successor obstruction).}\ours
\textit{If $T(A)$ holds, then $A$ is nonempty and any two objects of $A$ admit arrows to a common object. Consequently, for a preorder $A$, the condition $T(A)$ is equivalent to filteredness.}

\textit{Proof.}
The identity of $A$ is cofinal, so $T(A)$ implies that $B A$ is weakly contractible. In particular $A$ is nonempty and connected. Suppose that objects $a,b$ have no common successor. Let $U$ be the full subcategory of objects receiving an arrow from $a$, and let $V$ be the full subcategory of objects receiving an arrow from $b$. These subcategories are disjoint and closed under passage to successors.

Let $P$ be the three-object category with precisely the nonidentity arrows
\[
 0\xrightarrow{p}1,\qquad 0\xrightarrow{q}2.
\]
Define $f:A\longrightarrow P$ by sending $U$ to $1$, $V$ to $2$, and all other objects to $0$. Closure under successors makes this a functor: an arrow cannot leave $U$ or $V$, and there are no arrows between them. The category $0\backslash A$ computed for $f$ is $A$ itself. The other two comma categories are $U$ and $V$. They are connected, since every object of $U$ receives an arrow from $a$, and every object of $V$ receives one from $b$. They need not have initial objects; the existence of those arrows is all that is used. Hence $f$ is cofinal.

We next give $P$ a cofinal image with nonzero second homology. Adjoin an object $z$ and arrows
\[
 r,s:z\rightrightarrows0,
 \qquad v:z\longrightarrow1,
 \qquad w:z\longrightarrow2,
\]
subject to $pr=ps=v$ and $qr=qs=w$, and no other nonidentity arrows. Call the resulting category $Q$. The inclusion $j:P\hookrightarrow Q$ is cofinal. For an object of $P$, the relevant comma category has an initial object. For $z$, the category $z\backslash P$ has four objects $r,s,v,w$ and four arrows
\[
 r\longrightarrow v,\quad r\longrightarrow w,
 \quad s\longrightarrow v,\quad s\longrightarrow w;
\]
it is nonempty and connected.

To see exactly what obstructs contractibility, use normalized simplicial chains of $NQ$. There are four nondegenerate two-simplices, denoted $(r,p),(s,p),(r,q),(s,q)$, and no nondegenerate simplices in higher dimensions. Their boundaries are
\[
\begin{aligned}
 \partial(r,p)&=p-v+r,&\partial(s,p)&=p-v+s,\\
 \partial(r,q)&=q-w+r,&\partial(s,q)&=q-w+s.
\end{aligned}
\]
Therefore
\[
 c=(r,p)-(s,p)-(r,q)+(s,q)
\]
is a nonzero cycle. Since there are no three-simplices, it is not a boundary. More precisely, the coefficients of $p,q,r,s,v,w$ in the boundary show that every two-cycle is an integral multiple of $c$, so $H_2(BQ;\mathbf Z)\cong\mathbf Z$. The composite $jf$ is cofinal, contradicting $T(A)$. This proves the common-successor assertion.

In a preorder, parallel arrows are already equal. Nonemptiness and common successors therefore imply filteredness. The converse follows from preservation of filteredness by ordinary cofinal images, as in Corollary~3.4. This proves the last assertion.\hfill$\square$

For general categories the argument leaves precisely the parallel-arrow equalization problem. The distinction is substantive: a one-object group category satisfies the common-successor condition for every group, while filteredness forces that group to be trivial. Our literature check did not locate a resolution of the full implication $T(A)\Rightarrow A\text{ filtered}$ under the exact quantifier used here. This is a bounded report on the sources checked, not a claim that the question is known to be open throughout the literature.
\end{dcomment}


% ---- from Der-ChII-234_p026-030.tex
\printedpage{24}
% head: A. Grothendieck : Dérivateurs, Ch. II, Édition des Universités de Montpellier II et Paris VII
% foot: Version 8 juillet 2001
\msp{page 37}

\section*{3\quad Retour sur la cofinalité}

Il est instructif de relier les questions de cofinalité (pour une flèche $u : I' \longrightarrow I$ dans Cat) au dérivateur $\HOT_{W_0}$, où $W_0 \subseteq \Fl\,\Cat$ est le plus grossier des localiseurs dans Cat satisfaisant l'axiome de connexité Loc 5, i.e.
\[
W_0 = \{ f \in \Fl\,\Cat \mid \pi_0(f) \text{ bijectif}\}. \tag*{(3.1)}
\]
On trouve alors, pour $I$ dans Cat, une équivalence canonique, 2-fonctorielle en $I$,
\[
\HOT_{W_0}(I) \xrightarrow{\;\sim\;} I^{\wedge} \defeq \uHom(I,(\Ens)), \quad \text{\pl}\ \uHom(I^o,\Ens)\,\text{\texttt{]}} \tag*{(3.2)}
\]
en d'autres termes, on a une équivalence
\[
\HOT_{W_0} \simeq \mathbf{D}_{(\Ens)}. \tag*{(3.3)}
\]
Quand $\HOT_{W_0}$ est défini par localisation à partir de la catégorie $\Cat(I) \defeq \uHom(I^o,\Cat)$, la flèche (3.2) est celle déduite, par localisation, de la flèche
\[
\begin{array}{rcl}
\pi_0^I : \uHom(I^o,\Cat) & \longrightarrow & \uHom(I^o,(\Ens)) \\[1ex]
X & \longmapsto & \pi_0 \circ X,
\end{array} \tag*{(3.4)}
\]
où
\[
\pi_0 : \Cat \longrightarrow (\Ens) \tag*{(3.5)}
\]
est le foncteur ‘composantes connexes’. Si on décrit $\HOT_{W_0}(I)$ comme localisée de $\Cat/I$, (3.2) est déduit de la flèche
\[
\left\{
\begin{array}{rcl}
\Cat/I & \longrightarrow & I^{\wedge} = \uHom(I^o,(\Ens)) \\[1ex]
X & \longmapsto & \underline{\pi_0}(X/I) = (i \longmapsto \pi_0(i\backslash X)).
\end{array}
\right. \tag*{(3.6)}
\]

\msp{page 38}
Ainsi, les foncteurs familiers, pour une flèche
\[
u : I' \longrightarrow I
\]
dans Cat
\[
\left\{
\; I'^{\wedge} \quad
\begin{array}{c}
\xrightarrow{\;u_!\;} \\[1.5ex]
\xleftarrow{\;u^*\;} \\[1.5ex]
\xrightarrow{\;u_*\;}
\end{array}
\quad I^{\wedge},
\right. \tag*{(3.7)}
\]
\printedpage{25}
% head: A. Grothendieck : Dérivateurs, Ch. II, Édition des Universités de Montpellier II et Paris VII
% foot: Version 8 juillet 2001
s'interprètent comme les foncteurs de même nom pour le dérivateur $\HOT_{W_0}$, ou plus directement, pour le dérivateur (qu'on pourrait faussement considérer comme canularesque) $\mathbf{D}_{(\Ens)}$.

\dgrblock{dgrb108}{The represented derivator at dimension zero}
\begin{dcomment}{exposition}{p37}{The represented derivator at dimension zero}
Grothendieck now uses sets to separate the genuinely categorical part of cofinality from its higher homotopical refinements. The equivalence $\HOT_{W_0}(I)\simeq\operatorname{Fun}(I^{\mathrm{op}},\mathrm{Set})$ has a short proof that explains why no hidden higher information survives. Let $d:\mathrm{Set}\longrightarrow\mathrm{Cat}$ assign the discrete category, and let $\pi_0$ be connected components. There is an adjunction $\pi_0\dashv d$, with $\pi_0d=1$. For each category $C$, the unit $C\longrightarrow d\pi_0C$ is a $W_0$-equivalence. Applied pointwise to an $I^{\mathrm{op}}$-diagram, these units give a natural objectwise $W_0$-equivalence from that diagram to its discrete component diagram. Thus the two functors become inverse after localization. Since the construction commutes with restriction, it gives the stated equivalence of diagram theories.

The description by $\mathrm{Cat}/I$ requires the relative, rather than the global, component construction. To $X\longrightarrow I$ one associates the presheaf $i\longmapsto\pi_0(i\backslash X)$. Localizing the slice merely at functors that induce a bijection on the components of their total categories would give a different category. The relevant class is detected in all these comma categories, exactly as in the relative localization recalled in the commentary on Chapter~I, manuscript p.~43 (``What the localization axioms do and do not say'').

The large table then has an economical interpretation. The words ``aspherical'' and ``coaspherical'' become ordinary initiality and finality; the proper and smooth conditions say that specified strict lift categories are nonempty and connected. Right and left Kan extensions retain their ordinary meanings. Nothing in this reduction licenses replacing connected lift categories by contractible ones. Conversely, it makes the degree-zero assertions accessible by elementary calculations with connected components, without invoking a construction of general homotopy Kan extensions.
\end{dcomment}


Cela donne un sens à la cohorte des notions cohomologiques associées à un dérivateur, ou plus particulièrement, à un localiseur fondamental $W$ sur Cat (donnant naissance au dérivateur $\HOT_W$). On les désignera par le nom de $W_0$-fourbi (au lieu de $\HOT_{W_0}$-fourbi, ou $\mathbf{D}_{(\Ens)}$-fourbi), notamment pour une flèche $u : I' \longrightarrow I$ de Cat telle la notion
\[
\left\{
\begin{array}{lll}
\text{$W_0$-équivalence} & \Bigl(\iff & \text{$\pi_0(u)$ bijectif.}\Bigr) \\
\text{$W_0$-asphérique} & \Bigl(\iff & \text{les $I'/i$ pour $i \in \Ob I$ sont} \\
 & & \text{$0$-connexes.}\Bigr) \\
\text{$W_0$-coasphérique} & \Bigl(\iff & \text{les $i\backslash I'$ pour $i \in \Ob I$ sont} \\
 & & \text{$0$-connexes.}\Bigr) \\
\text{$W_0$-propre (\footnotemark[46])} & \Bigl(\iff & \text{les inclusions $I'_i \longrightarrow I'/i$ ($i \in \Ob I$)} \\
 & & \text{sont $W_0$-coasphériques, i.e. les fibres} \\
 & & \text{des colocalisées $i'\backslash I' \longrightarrow i\backslash I$} \\
 & & \text{($i' = u(i)$) \pl\ ($i = u(i')$)]} \\
 & & \text{sont $0$-connexes.}\Bigr) \\
\text{\mspin{page 39}} & & \\
\text{$W_0$-lisse (\footnotemark[47])} & \Bigl(\iff & \text{les inclusions $I'_i \longrightarrow i\backslash I'$} \\
 & & \text{sont $W_0$-asphériques, i.e. les fibres} \\
 & & \text{des flèches localisées $I'/i' \longrightarrow I/i$} \\
 & & \text{($i = u(i')$) sont $0$-connexes.}\Bigr) \\
\text{$W_0$-parfait} & \Bigl(\iff & \text{$W_0$-propre et $W_0$-lisse.}\Bigr) \\
\text{$W_0$-équivalence universelle (\footnotemark[48])} & \Bigl(\iff & \text{$\pi_0$ bijectif et le reste par tout} \\
 & & \text{changement de base. Ce sont aussi} \\
 & & \text{les $W_0$-fibrations (ou $W_0$-fibrations} \\
 & & \text{faibles à fibres $W_0$-asphériques,} \\
 & & \text{i.e. $0$-connexes.}\Bigr) \\
\text{$W_0$-fibration} & \Bigl(\iff & \text{$u^* : \Cat/I \longrightarrow \Cat/I'$ transforme} \\
 & & \text{$W_0$-équivalences en $W_0$-équivalences.}\Bigr) \\
\text{$W_0$-fibration faible (\footnotemark[49])} & \Bigl(\iff & \text{le changement de base} \\
 & & \text{$u^* : \Cat/I \longrightarrow \Cat/I'$ transforme} \\
 & & \text{homotopismes $J_1 \longrightarrow J_2$ en} \\
 & & \text{$W_0$-équivalences - et il suffit de} \\
 & & \text{l'exiger pour l'inclusion d'un objet} \\
 & & \text{final (ou au choix initial) et pour} \\
 & & \text{l'inclusion d'un point dans $\boldsymbol{\Delta}^1$.}\Bigr)
\end{array}
\right. \tag*{(3.8)}
\]
\footnotetext[46]{Exemple : les Cat-cofibrations.}
\footnotetext[47]{Exemple : les Cat-fibrations.}
\footnotetext[48]{Exemple : les flèches $W_0$-lisses ou $W_0$-propres, à fibres 0-connexes.}
\footnotetext[49]{Exemple : les flèches $W_0$-lisses telles que $\pi_0(I'/I) = (i \longmapsto \pi_0(i\backslash I')) = u_!(e_{I'^{\wedge}})$ soit localement constant sur $I$, i.e. $\underline{\pi_0}(I'/I)$ soit un \textit{revêtement} étale de $I$, ou si $u^{\mathrm{opp}}$ satisfait cette condition.}

\printedpage{26}
% head: A. Grothendieck : Dérivateurs, Ch. II, Édition des Universités de Montpellier II et Paris VII
% foot: Version 8 juillet 2001
\texttt{[Nous avons une]} implication \textit{stricte} (et comment)
\[
\text{$W_0$-équivalence universelle} \Longrightarrow \text{$W_0$-fibration}
\]
\texttt{[et une]} implication \textit{stricte} (et comment\,!) (\footnote[50]{Mais si les fibres sont 0-connexes, les deux notions de $W_0$-fibrations sont équivalentes et signifient qu'on a une $W_0$-équivalence universelle \dots})
\[
\text{$W_0$-fibration} \Longrightarrow \text{$W_0$-fibration faible.}
\]

Signalons, comme exemple d'utilisation de ces notions, que pour un diagramme cartésien dans Cat
\[
\begin{tikzcd}
I' \arrow[d, "u"'] & J' \arrow[l, "p'"'] \arrow[d, "v"] \\
I & J, \arrow[l, "p"']
\end{tikzcd} \tag*{(3.9)}
\]
la flèche de changement de base pour $u_*$, $v_*$ pour les préfaisceaux est \texttt{[un]} isomorphisme si $u$ est $W_0$-propre, ou $p$ $W_0$-lisse, et la flèche de changement de base pour $u_!$, $v_!$ est \texttt{[un]} isomorphisme si $u$ est $W_0$-lisse ou $p$ $W_0$ propre (\footnote[51]{Sauf erreur, ca donne \textit{deux caractérisations} des morphismes $W_0$-propres, ou des morphismes $W_0$-lisses.}). Je signale à ce propos que ces énoncés, qui a priori se formulent pour les préfaisceaux à valeurs dans (Ens), restent valables % [sic] ca -> ça (in footnote 51), missing cedilla
\msp{page 40}
pour des préfaisceaux à valeurs dans une catégorie $\mathcal{M}$ quelconque (sous réserve d'existence des operations d'images directes cohomologiques ou homologiques envisagées). Cela se voit p.ex. en notant que % [sic] operations -> opérations, missing accent
\[
\text{$W_0$-équivalence} \;\Longrightarrow\; \text{$\mathbf{D}_{\mathcal{M}}$-équivalence} \tag*{(3.10)}
\]
pour toute $\mathcal{U}$-catégorie $\mathcal{M}$, d'où résulte que toute (\footnote[52]{Pas tout à fait exacte pour \textit{toute} ‘$W_0$-notion’, il faut excepter les deux notions de $W_0$-fibration forte, et faible.}) ‘$W_0$-notion’ implique la $\mathbf{D}_{\mathcal{M}}$-notion correspondante, et permet donc d'appliquer à $\mathbf{D}_{\mathcal{M}}$ les résultats le concernant. (Il y a bien sûr un argument direct, moins savant, consistant à se ramener au cas où $\mathcal{M}$ est petite, et à utiliser alors le plongement pleinement fidèle $\mathcal{M} \hookrightarrow \mathcal{M}^{\wedge}$ resp. $\mathcal{M} \hookrightarrow ((\mathcal{M}^o)^{\wedge})^o$.)

Je rappelle pour mémoire une (troisième\,?) caractérisation des flèches $u : I' \longrightarrow I$ $W_0$-propres (resp. $W_0$-lisses) : ce sont celles pour lesquelles le foncteur changement de base $u^* : \Cat/I \longrightarrow \Cat/I'$ transforme flèches $W_0$-coasphériques (resp. $W_0$-asphériques) en itou. Cela gagne en intérêt, du fait que la notion de foncteur $W_0$-coasphérique (resp. $W_0$-asphérique) s'interprète comme celle de foncteur cofinal (resp. co-cofinal), cf. plus bas.

\printedpage{27}
% head: A. Grothendieck : Dérivateurs, Ch. II, Édition des Universités de Montpellier II et Paris VII
% foot: Version 8 juillet 2001
\msp{page 41}
La notion de $W_0$-fibration apparaît comme extrêmement forte - beaucoup plus forte à certains égards que celle (disons) de $W_\infty$-fibration ou de $W_\omega$-fibration. Elle équivaut à ceci (comme on voit aussitôt) : pour toute composante connexe $I'_\alpha$ de $I'$, désignant par $I_\alpha$ la composante connexe correspondante de $I$, et pour tout $J_\alpha$ 0-connexe et au dessus de $I_\alpha$, le produit fibré $I'_\alpha \times_{I_\alpha} J_\alpha$ est encore 0-connexe. (Ça ressemble aux histoires de totale 0-connexité d'un topos\,!) Mais cela signifie aussi que $I'_\alpha \longrightarrow I_\alpha$ est une $W_0$-équivalence universelle. Cela implique que ses fibres sont 0-connexes, et la réciproque est vraie si $u : I' \longrightarrow I$ est $W_0$-propre ou $W_0$-lisse. Ainsi, la condition essentielle pour une $W_0$-fibration est celle que les fibres des flèches induites $I'_\alpha \longrightarrow I_{u(\alpha)}$ \pl\ $I'_\alpha \longrightarrow I_\alpha$\texttt{]} soient 0-connexes. En fait, on a ceci

\dgrblock{dgrb109}{The lift categories behind ordinary base change}
\begin{dcomment}{exposition}{p40}{The lift categories behind ordinary base change}
The base-change assertion is an application of the lift-category criterion already developed in the commentary on Chapter~I, manuscript pp.~68--70 (``The category of lifts behind properness'' and ``Completing the tests in the proper base-change theorem''). Its present specialization is useful because all the relevant tests are tests of connectedness. For $u:A\longrightarrow B$, properness means that, for every $a\in A$ and arrow $\beta:u(a)\longrightarrow b$, the category
\[
 L(a,\beta)=\{a\xrightarrow{h}a'\mid u(a')=b, u(h)=\beta\}
\]
of strict lifts, with arrows over $1_b$, is nonempty and connected. Smoothness is the dual condition on arrows into $u(a)$. An opfibration supplies initial objects in the displayed lift categories; a fibration supplies the dual terminal objects. Connectedness is the exact weakening that suffices for ordinary Kan extensions.

For a strict pullback square as in (3.9), the comparison of the comma categories computing a right Kan extension is final if $u$ is proper or if the horizontal base-change functor is smooth. The ordinary cofinality theorem therefore gives the right base-change isomorphism. Reversing arrows gives the left formula. The same argument works in any coefficient category in which the requisite comma limits exist, so that these Kan extensions are computed pointwise. One can apply $\operatorname{Hom}(m,-)$ to the limits, or dually $\operatorname{Hom}(-,m)$ to the colimits, and reduce to the set-valued cofinality theorem. This avoids imposing unnecessary completeness assumptions on the coefficient category.

The passage from a $W_0$-equivalence to an equivalence measured by constant coefficients has a similarly elementary explanation. A limit of a constant diagram is a product indexed by the connected components of its shape, and a colimit is the corresponding coproduct, whenever those objects exist. A bijection on components preserves both. The printed footnote correctly excludes the two fibration notions from the blanket transfer of terminology: preserving a larger class of equivalences under every base change is a stronger requirement, even when the equivalences themselves have a simpler definition. The next proposition measures this effect exactly.
\end{dcomment}


\noindent\textbf{Proposition 3.1.} \textit{Soit $u : I' \longrightarrow I$ une flèche dans} Cat.
\begin{itemize}
\item[\textit{a)}] \textit{Pour que ce soit une $W_0$-équivalence universelle, il faut et il suffit que pour toute $\boldsymbol{\Delta}^1 \longrightarrow I$, i.e. toute flèche $\alpha$ de $I$, la catégorie $I' \times_I \boldsymbol{\Delta}^1 \defeq I'|\alpha$ soit $0$-connexe, ou encore que les fibres de $I$ soient $0$-connexes, et que pour toute flèche} non identique \textit{$\alpha$ de $I$, $I'|\alpha$ soit $0$-connexe.}
\item[\textit{b)}] \textit{Pour que $u : I' \longrightarrow I$ soit une $W_0$-fibration, il faut et il suffit que pour toute composante connexe $I'_{\alpha'}$ de}
\msp{page 42}
\textit{$I'$, posant $\alpha = u(\alpha')$, la flèche induite $I'_{\alpha'} \longrightarrow I_\alpha$ soit une $W_0$-équivalence universelle (cf. critère a) pour cette notion).}
\item[\textit{c)}] \textit{Pour que ce soit une $W_0$-fibration faible, il} faut \textit{que $\forall\, i \in I$, l'inclusion $I'_i \hookrightarrow i\backslash I'$ soit une $W_0$-équivalence, i.e. $\pi_0(I'_i) \xrightarrow{\;\sim\;} \pi_0(i\backslash I')$ soit bijectif,} et \textit{que $i \longmapsto \pi_0(i\backslash I')$ soit un préfaisceau localement constant sur $I$ (i.e. $\underline{\pi_0}(I'/I)$ soit un} revêtement \textit{étale de $I$). Cette condition est aussi suffisante, si $u$ est $W_0$-lisse. De plus (pour mémoire) $u$ est une $W_0$-fibration faible si et seulement si $u^{\mathrm{opp}}$ l'est.}
\end{itemize}

\dgrblock{dgrb110}{Why arrows detect universal connectedness}
\begin{dcomment}{exposition}{p42:Cette condition est aussi suffisante}{Why arrows detect universal connectedness}
Proposition~3.1 says that universal preservation of connected components can be checked on vertices and single arrows. Its proof is a useful model for reading the much stronger componentwise condition in part~(b). We give the details because the proposition is stated, rather than proved, here.

Suppose first that $u:A\longrightarrow B$ is a universal $W_0$-equivalence. Pullback along any $[1]\longrightarrow B$ gives a connected nonempty category, since $[1]$ is connected. Conversely, assume this arrow test. Applying it to an identity arrow at $b$ shows that $A_b\times[1]$, and hence $A_b$, is connected and nonempty. Let $v:J\longrightarrow B$ be arbitrary. Every object of $J$ has a nonempty fibre in $A\times_BJ$, so the map on components is surjective. If two objects of the pullback project to objects joined by a zigzag in $J$, each edge of that zigzag has a connected pullback by hypothesis. Its endpoints can be joined through that pullback, and the connected vertex fibres join any two choices of endpoints. Concatenating these finite zigzags proves injectivity. Thus
\[
 \pi_0(A\times_BJ)\xrightarrow{\sim}\pi_0(J)
\]
for every $J$, which is the required universal equivalence. Notice that the test concerns strict pullbacks, not comma categories.

Now let $A_\lambda$ be a connected component of $A$, mapping to the component $B_\mu$. If pullback by $u$ preserves every $W_0$-equivalence, apply it to $J\amalg\coprod_{\nu\ne\mu}B_\nu\longrightarrow B$, where $J\longrightarrow B_\mu$ is any connected nonempty category over $B_\mu$. The resulting map on components forces $A_\lambda\times_{B_\mu}J$ to be connected and nonempty. Decomposing an arbitrary category over $B_\mu$ into its components then shows that $A_\lambda\longrightarrow B_\mu$ is a universal $W_0$-equivalence. Conversely, if all these component maps are universal equivalences, then for any $J\longrightarrow B$ there is a natural bijection
\[
 \pi_0(A\times_BJ)\cong
 \pi_0(A)\times_{\pi_0(B)}\pi_0(J).
\]
This formula immediately shows that pullback by $u$ preserves every bijection on connected components. It proves part~(b), including components of $B$ not met by $A$.

The necessary conditions in part~(c) come from the same tests with categorical homotopy equivalences. The inclusion of the initial object $1_b$ into $b\backslash B$ is such an equivalence; its pullback identifies $A_b\longrightarrow b\backslash A$ on components. Testing both endpoints of an arrow in $B$ makes the resulting component presheaf locally constant. When $u$ is smooth, ordinary left base change identifies the component presheaf after every base change with the restriction of this local system. Its colimit depends only on the fundamental groupoid of the indexing category, and categorical homotopy equivalences induce equivalences of those groupoids. This proves the stated sufficiency. The opposite-category assertion follows directly because categorical homotopies and $\pi_0$ are unchanged by reversal of arrows.
\end{dcomment}


La plus grande différence entre la notion de $W_0$-fibration, et celle de $W_0$-fibration faible, me paraît que les morphismes $W_0$-parfaits (i.e. $W_0$-lisses et $W_0$-propres), et a fortiori les $W$-parfaits pour un localiseur fondamental $W \subseteq W_0$, sont forcément des $W$-fibrations faibles, mais (\textit{même si $W = W_\omega$,} le plus fin de tous les localiseurs fondamentaux) ne sont pas des $W$-fibrations. L'exemple le plus frappant à cet égard est sans doute celui des revêtements étales, qui sont $W_\omega$-parfaits (donc
\msp{page 43}
des $W_\omega$-fibrations faibles et a fortiori des $W_0$-fibrations faibles), mais pas des $W_0$-fibrations\,; car dans le cas où $I'$, $I$ sont 0-connexes, il faudrait que les fibres soient 0-connexes (cf. b) ci-dessus), i.e. que $u$ soit un isomorphisme\,!

\dgrblock{dgrb111}{Monodromy explains the strictness}
\begin{dcomment}{examples}{p43:i.e. que $u$ soit un isomorphisme}{Monodromy explains the strictness}
The covering example isolates what the componentwise test excludes. This is the same monodromy issue as in the commentary on Chapter~I, manuscript p.~98 (``Sets already detect the missing monodromy''), now applied to the two precise definitions in (3.8). Let $G$ be a nontrivial group. Its translation groupoid $E G$ has object set $G$ and exactly one arrow between each ordered pair of objects. The functor
\[
 E G\longrightarrow B G
\]
sends $g\longrightarrow h$ to $hg^{-1}$. It is a discrete fibration and a discrete opfibration, so the strict lift categories are singletons and the map is proper and smooth. Both total and base categories are connected, but the strict fibre is the discrete set $G$.

The inclusion of the sole object $*\longrightarrow BG$ is a $W_0$-equivalence. Its pullback along $EG\longrightarrow BG$ is the map $G_{\mathrm{disc}}\longrightarrow EG$, whose map on components is $G\longrightarrow *$. This proves directly that the covering is not a $W_0$-fibration. It is a weak $W_0$-fibration by the smooth criterion of Proposition~3.1(c): the relative component presheaf is the locally constant $G$-set with regular action. Thus weak fibrations permit monodromy between different fibre components, whereas the strong $W_0$ condition treats each connected component of the total category separately.

There is a subscript caution in the preceding printed paragraph. The conclusion established by this example is ``not necessarily a $W_0$-fibration.'' It does not assert that coverings fail to preserve weak homotopy equivalences under pullback. In fact their nerves are simplicial coverings, and they do preserve weak homotopy equivalences. Similarly a $W_0$-perfect map gives a weak $W_0$-fibration; a $W$-perfect map gives the corresponding weak $W$-fibration. These are distinct assertions when $W\subsetneq W_0$.
\end{dcomment}


Une propriété sympathique des morphismes $W_0$-lisses (resp. $W_0$-propres), c'est que ce sont des applications \textit{ouvertes} (resp. fermées) pour les espaces topologiques sous-jacents, ce qui signifie en effet que les fibres locales, i.e. des flèches localisées en haut $I'/i' \longrightarrow I/i$ (resp. colocales, i.e. les flèches colocalisées en bas $i'\backslash I' \longrightarrow i\backslash I$), sont surjectives, i.e. à fibres non-vides (ce qui est impliqué par la condition que celles-ci soient 0-connexes).

\printedpage{28}
% head: A. Grothendieck : Dérivateurs, Ch. II, Édition des Universités de Montpellier II et Paris VII
% foot: Version 8 juillet 2001
Dans tout ça je n'ai encore guère parlé des flèches $W_0$-asphériques resp. co-asphériques, qui pourtant sont ici une motivation principale pour introduire $W_0$\,!

\noindent\textbf{Proposition 3.2.} \textit{Conditions équivalentes sur la flèche $u : I' \longrightarrow I$ dans} Cat :
\begin{itemize}
\item[a)] \textit{Pour toute $\mathcal{U}$-catégorie $\mathcal{M}$ et tout foncteur $F : I \longrightarrow \mathcal{M}$, la flèche canonique (sous réserve que les petites limites inductives soient représentables dans $\mathcal{M}$),}
\[
\varinjlim_{I'} u^*(F) \longrightarrow \varinjlim_{I} F \tag*{(3.2.1)}
\]
\msp{page 44}
\textit{est} \texttt{[un]} \textit{isomorphisme.}
\item[a$'$)] \textit{Pour toute $\mathcal{U}$-catégorie $\mathcal{M}$ où les petites limites projectives sont représentables et tout $F \in \mathcal{M}(I) = \uHom(I^o,\mathcal{M})$, i.e. tout préfaisceau sur $I$ à valeurs dans $\mathcal{M}$, la flèche canonique}
\[
\varprojlim_{I} F \longrightarrow \varprojlim_{I'} u^*(F) \tag*{(3.2.2)}
\]
\textit{est} \texttt{[un]} \textit{isomorphisme.}
\item[a$''$)] (\footnote[53]{C'est une condition sur les morphismes de topos, savoir que si $F$ \texttt{[est]} dans $\mathcal{X}^{\wedge}$, on a
\[
H^0(\mathcal{X},F) \xrightarrow{\;\sim\;} H^0(\mathcal{X}',u^*F).
\]
Il y a \texttt{[une]} notion similaire pour les $H^i$ ($i \leq n$), à dégager. Notion de morphismes $W_0$-coasphériques de topos\,! C'est dans la nature d'une condition de 0-connexité relative, duale de la notion $W_0$-asphérique, qui signifie simplement que pour $U$ 0-connexe dans $U^{\wedge}$ \pl\ $\mathcal{X}^{\wedge}$\texttt{]}, $f^*(U)$ \pl\ $u^*(U)$\texttt{]} est 0-connexe. De même, on a sans doute des notions (‘duales’ des précédentes) de $W_n$-coasphéricité (bien comprise en principe si $\mathcal{X}$ est localement $n$-asphérique) (La $W_n$-coasphéricité, étant celle évoquée plus haut.)}) \textit{Comme a$'$, avec $\mathcal{M} = (\Ens)$, i.e. $F \in \Ob I^{\wedge}$.}~(\footnote[54]{\textbf{NB} On peut introduire
\begin{itemize}
\item[a$'''$)] Comme a), mais avec $\mathcal{M} = (\Ens)$.
\item[a$^{IV}$)] Comme a), avec le foncteur canonique $I' \longrightarrow I'^{\wedge}$.
\end{itemize}}).
\item[b)] \textit{La flèche $u$ est $W_0$-coasphérique, i.e. pour tout $i \in I$, $i\backslash I'$ est $0$-connexe.}
\end{itemize}
L' équivalence de a, a$'$, a$''$ est $\pm$ tautologique en termes des définitions (cf. SGA 4, I, 8\dots). La forme a$''$ s'écrit, posant $\mathbf{D} = \HOT_{W_0}$, que $\forall\, F \in \mathbf{D}(I)$, on veut que
\[
\underbrace{H^{\bullet}_{\mathbf{D}}(I,F)}_{\defeq p_{I,*}(F)} \longrightarrow \underbrace{H^{\bullet}_{\mathbf{D}}(I',u^*F)}_{\defeq p_{I',*}(F)} \qquad \text{\pl}\quad \underbrace{H^{\bullet}_{\mathbf{D}}(I',u^*F)}_{\defeq p_{I',*}(u^*F)}\ \text{\texttt{]}} \tag*{(3.2.3)}
\]
($p_I : I \longrightarrow e$ et $p_{I'} : I' \longrightarrow e$ les foncteurs canoniques) soit \texttt{[un]} isomorphisme. On sait que cela équivaut à $u$ $\mathbf{D}$-coasphérique. La démonstration est un jeu amusant entre $u_!$, $u^*$, $u_*$ à la fois. Il faut exprimer que le Hom d'un objet $M$ de $\mathbf{D}(e)$ ($=$ (Ens) en l'occurence) dans la flèche (3.2.3) est bijective. Ici, il % [sic] occurence -> occurrence, French spelling

% ---- from Der-ChII-234_p031-035.tex
\printedpage{29}
% head: A. Grothendieck : Dérivateurs, Ch. II, Édition des Universités de Montpellier II et Paris VII
% foot: Version 8 juillet 2001

\msp{page 45}
est plus naturel de prendre $M = e$ (ensemble ponctuel) et d'interpréter la flèche (3.2.3) comme
\[
\Hom(e_I, F) \longrightarrow \Hom(e_{I'}, u^*F)
\tag*{(3.2.4)}
\]
et par adjonction de $u_!$, $u^*$ le deuxième membre s'interprète comme $\Hom(u_!(e_{I'}), F)$, et la flèche est déduite de
\[
u_!e_{I'} = u_!u^*(e_I) \xrightarrow{\text{adj.}} e_I
\tag*{(3.2.5)}
\]
en prenant le Hom dans $F$. Donc dire que (3.2.4) est \texttt{[un]} isomorphisme pour tout $F$ signifie que (3.2.5) est \texttt{[un]} isomorphisme. Passant aux valeurs des deux membres en $i$, on trouve $\pi_0(i\bs I') \xrightarrow{\sim} 1$, i.e. que $i\bs I'$ est 0-connexe, q.e.d.

\dgrblock{dgrb112}{What the ordinary cofinality theorem tests}
\begin{dcomment}{exposition}{p45:i.e. que $i\bs I'$ est 0-connexe, q.e.d.}{What the ordinary cofinality theorem tests}
Grothendieck has now proved the criterion announced on the opening page, with the sufficient directedness condition replaced by the exact connectedness condition. The proof is the representability calculation already explained in ``The presheaf represented by a change of index'': the functor $F\longmapsto\Gamma(I',u^*F)$ is represented by $u_!1_{I'}$, and
\[
 (u_!1_{I'})(i)=\pi_0(i\backslash I').
\]
Consequently the comparison on all set-valued presheaves is an isomorphism precisely when $u_!1_{I'}\longrightarrow1_I$ is an isomorphism. Applying ordinary Hom functors reduces arbitrary existing limits and colimits to this case. This is the classical theorem of SGA~4, Expos\'e~I, Section~8.1, in the variance used throughout these pages \citep{SGA4I1972}.

Two small examples distinguish its quantifiers. The inclusion $\{1\}\hookrightarrow[1]$ is cofinal, because each $i\backslash\{1\}$ is a point. For a presheaf $F(1)\longrightarrow F(0)$, its limit is therefore $F(1)$. But the restriction functor forgets $F(0)$ and is not fully faithful. On the other hand, for any nonempty connected category $C$, the projection $C\longrightarrow *$ is cofinal even when $BC$ has nontrivial homotopy groups. Ordinary cofinality is therefore stronger than a bijection on components of arbitrary source and target, weaker than full faithfulness of restriction, and weaker than homotopy cofinality in a different direction. These comparisons refer to different properties of a functor, not to alternative conventions for the same word.

Footnote~53 distinguishes two degree-zero conditions on a morphism of topoi $f:\mathcal E'\longrightarrow\mathcal E$. Its coasphericity condition preserves global sections of every object. Its asphericity condition instead requires $f^*U$ to be noninitial and connected whenever $U$ is. For locally connected topoi, the latter is equivalent to preservation of component sets: decompose $U$ into its noninitial connected components and use preservation of coproducts by $f^*$. For the morphism induced by $u:I'\longrightarrow I$, the restricted representable $u^*h_i$ has category of elements $I'/i$. Thus this asphericity condition is ordinary initiality of $u$, by the cofinality theorem with arrows reversed; the section condition is finality. The two endpoints of $[1]$ distinguish them: inclusion of $0$ is initial, while inclusion of $1$ is final.

The supplementary Yoneda test in footnote~54 should use the functor $I\longrightarrow I^{\wedge}$. Its colimit is $1_I$, whereas its restriction along $u$ has colimit $u_!1_{I'}$. The printed $I'\longrightarrow I'^{\wedge}$ has the wrong domain for the test as stated. The further suggestion in footnote~53 about higher cohomology should be read with coefficients specified: preservation of all set-valued sections, preservation of abelian cohomology through a fixed degree, and equivalence on all homotopy-valued diagrams detect different truncations. The earlier cohomological cofinality block supplies the derived comparison; the ordinary theorem here establishes its degree-zero model.
\end{dcomment}


\noindent\textbf{Définition 3.3.}
\begin{itemize}[label={}]
\item[a)] Quand les conditions équivalentes de 3.2 sont satisfaites, donc quand les $i\bs I'$ ($i \in \Ob I$) sont 0-connexes, on dit que $u : I' \longrightarrow I$ est un foncteur \textit{cofinal.} On dit que $u$ est \textit{co-cofinal} si $u^o : I'^o \longrightarrow I^o$ est cofinal, i.e. si les localisées $I'/i$ de $I'$ sur $I$ ($i \in \Ob I$) sont 0-connexes.
\item[b)] Si $I'$ est une \textit{sous-catégorie} de $I$, on dit que $I'$ \textit{est cofinale dans} $I$ si le foncteur d'inclusion est cofinal, \textit{co-cofinale} si ce foncteur est co-cofinal.
\end{itemize}

\noindent\textbf{Corollaire 3.4.} (\footnote[55]{\textbf{NB} Pour un énoncé complémentaire de cofinalité, cf. XVIII prop. 4.9.6.3 (p. 224-225, 228-242). Il faudrait un énoncé qui coiffe les deux.})
\begin{itemize}[label={}]
\item[\textit{a)}] \textit{Pour que $u$ soit cofinal, il} faut \textit{que les catégories colocalisées $i\bs I'$ soient non-vides, i.e. qu'on ait}

\msp{page 46}
\begin{itemize}[label={}]
\item[\textit{F 0)}] $\pi_0(f)$ \texttt{[plutôt $\pi_0(u)$]} \textit{est bijectif.}
\item[\textit{F 1)}] \textit{Pour tout} $i \in \Ob I$, \texttt{[il]} \textit{existe} \texttt{[un]} $i' \in I'$ \textit{et une flèche} $i \longrightarrow u(i')$ \textit{dans} $I$.
\end{itemize}
(\footnote[56]{\textbf{NB} F 1 implique déjà \texttt{[la]} surjectivité de $\pi_0$, donc sa bijectivité si $I'$ est 0-connexe - dans ce cas F 0 résulte de F 1.}).
\item[\textit{b)}] \textit{Si $I'$ est faiblement filtrante, pour que $u$ soit cofinale, il} suffit \textit{qu'elle satisfasse F 1, ainsi que la condition :}
\begin{itemize}[label={}]
\item[\textit{F 2)}] \textit{Pour toute double flèche} $i \dbl{\beta}{\gamma} u(i')$ \textit{(avec} $i \in \Ob I$, $i' \in \Ob I'$\textit{),} \texttt{[il]} \textit{existe une flèche} $i' \xrightarrow{\alpha'} j'$ \textit{dans} $I'$\textit{, de telle façon que} $u(\alpha')$ \textit{égalise} $(\beta, \gamma)$\textit{, i.e.} $u(\alpha')\beta = u(\alpha')\gamma$.
\end{itemize}
\textit{et ces conditions F 1, F 2 impliquent que les $i\bs I'$ sont faiblement filtrantes, et $I$ aussi} (\footnote[57]{Si $I'$ est localement faiblement filtrante (= préfiltrante), alors pour la cofinalité de $u$, il suffit qu'on ait F 0, F 1, F 2, et alors les $i\bs I'$ sont faiblement filtrantes, et $I$ est préfiltrante tout comme $I'$.}).

\printedpage{30}
% head: A. Grothendieck : Dérivateurs, Ch. II, Édition des Universités de Montpellier II et Paris VII
% foot: Version 8 juillet 2001
\item[\textit{c)}] \textit{Si $I'$ est filtrante, alors $u$ est cofinal} si et seulement s'\textit{il satisfait les deux conditions F1, F2, ou encore si et seulement si les $i\bs I'$ sont} filtrantes \textit{(et non seulement} 0-\textit{connexes). Et cela implique que $I'$ est également filtrante} (\footnote[58]{Si $I'$ est localement filtrante, $u$ est cofinal si et seulement si elle satisfait F 0, F 1, F 2, i.e. si et seulement si les $i\bs I'$ sont filtrantes, et alors $I$ est localement filtrante.}).
\item[\textit{d)}] \textit{Si $I$ est préfiltrante (= localement faiblement filtrante, i.e. PS 1), et $u$} pleinement fidèle, \textit{pour que $u$ soit cofinal, il} faut \textit{et il} suffit \textit{que $u$ satisfasse la condition F 1 plus haut. Et alors $I'$ est également préfiltrante, et les $i\bs I'$ sont faiblement filtrantes. Si de plus $I$ est filtrante, les $i\bs I'$ et $I'$ le sont également. (Comparer SGA 4, I 8.1.3.)}
\end{itemize}

\msp{page 47}
\textsc{Démonstration.} a) F 1 est évident, car 0-connexe $\Longrightarrow$ non vide. F 0 car $W_0$-coasphérique $\Longrightarrow$ $W_0$-équivalence, i.e. $\pi_0(u)$ bijectif.

Pour \textit{prouver} b), il suffit de prouver que F 1, F 2 (si $I'$ est faiblement filtrante) impliquent que les $i\bs I'$ sont faiblement filtrantes (a fortiori 0-connexes) et de plus que $I$ l'est également. On sait déjà qu'ils sont non vides par F 1, reste à voir que si $(i', \alpha)$, $(j', \beta)$ sont dans $i\bs I'$, alors ils sont `majorés' par un troisième objet de $i\bs I'$. Comme $I'$ est faiblement filtrante, on trouve
\begin{tikzcd}[row sep=0pt, column sep=small, baseline=(k.base)]
i' \arrow[dr, "r"] & \\
& |[alias=k]| k' \\
j' \arrow[ur, "s"'] &
\end{tikzcd}, d'où $u(r)$, $u(s)$ dans $I$ (cf. diagramme)
\[
\begin{tikzcd}[row sep=small, column sep=small]
& i' \arrow[rr, "r"] \arrow[dd, dash] & & k' \arrow[r, dotted, "t"] \arrow[dd, dash] & l' \arrow[dd, dash] \\
& & j' \arrow[ur, "s"] \arrow[dd, dash] & & \\
& u(i') \arrow[rr, "u(r)"] & & u(k') \arrow[r, dotted, "u(t)"] & u(l') \\
i \arrow[ur, "\alpha"] \arrow[rr, "\beta"'] & & u(j'), \arrow[ur, "u(s)"'] & &
\end{tikzcd}
\]
appliquant F 2 à la double flèche $u(r)\alpha, u(s)\beta : i \rightrightarrows u(k')$, on trouve $k' \xrightarrow{t} l'$ (en pointillés sur \texttt{[le]} diagramme) telle que $(u(t)u(r))\alpha = (u(t)u(s))\beta$, soit $\gamma$, ce qui montre que l'objet $(l', \gamma : i \longrightarrow u(l'))$ majore $(i', \alpha)$ et $(j', \beta)$.

La démonstration que si $I'$ est faiblement filtrante et $u$ satisfait F 1, alors $I$ est faiblement filtrante, est immédiate.

\textit{Prouvons} c). Il reste à prouver que si $I'$ est \textit{filtrante} (non seulement faiblement filtrante), alors si $u$ \texttt{[est]} cofinal, alors F 2 est vérifiée, et de plus les $i\bs I'$ et $I$ sont filtrantes (pas seulement faiblement filtrantes). Mais notons tout de suite que $i\bs I' \longrightarrow I'$ est coétale, donc $I'$ filtrante, donc localement filtrante implique que $i\bs I'$ est aussi localement filtrante\,; comme par hypothèse les $i\bs I'$ sont 0-connexes, ils sont donc filtrantes. La condition F 2 sort facilement à présent :

\printedpage{31}
% head: A. Grothendieck : Dérivateurs, Ch. II, Édition des Universités de Montpellier II et Paris VII
% foot: Version 8 juillet 2001
\msp{page 48}
\[
\begin{tikzcd}[row sep=small, column sep=small]
& i' \arrow[rr, "r"] \arrow[dd, dash] & & j' \arrow[r, dotted, "t"] \arrow[dd, dash] & k' \arrow[dd, dash] \\
& & i' \arrow[ur, "s"] \arrow[dd, dash] & & \\
& u(i') \arrow[rr, "u(r)"] & & u(j') \arrow[r, dotted, "u(t)"] & u(k') \\
i \arrow[ur, "\alpha"] \arrow[rr, "\beta"'] & & u(i'). \arrow[ur, "u(s)"'] & &
\end{tikzcd}
\]
Comme $i\bs I'$ est filtrante, donc faiblement filtrante (PS 1), on peut trouver $i' \dbl{r}{s} j'$ telles que
\[
u(r)\alpha = u(s)\beta,
\]
et comme $I'$ est filtrante, on peut trouver $t : j' \longrightarrow k'$ qui égalise la double flèche $(r, s)$ :
\[
tr = ts(\defeq \tau),
\]
alors
\[
\underbrace{u(t)u(r)}_{u(tr)}\alpha = \underbrace{u(t)u(s)}_{u(ts)}\beta \qquad\quad \text{i.e.} \qquad\quad u(\tau)\alpha = u(\tau)\beta,
\]
i.e. on a vérifié F 2.

Reste à prouver que $I$ est filtrante (si $I'$ l'est et \texttt{[si]} $u : I' \longrightarrow I$ est cofinal). On sait déjà que $I$ est faiblement filtrante par b), reste à vérifier la condition de la double flèche
\[
i \dbl{\alpha}{\beta} j
\]
(PS 2). Soit $i'$ dans $I'$ et $j \longrightarrow u(i')$, et considérons
\[
i \dbl{\gamma\alpha}{\gamma\beta} u(i'),
\]
par hypothèse F 2 $\exists\, r : i' \longrightarrow j'$ dans $I'$ telle que
\msp{page 49}
$u(r)(\gamma\alpha) = u(r)(\gamma\beta)$, donc $u(r)\gamma$ égalise $(\alpha, \beta)$, q.e.d.

La vérification de d) par AQT \texttt{[abréviation dont le sens échappe aux éditeurs]}.

\dgrblock{dgrb113}{The filtered criterion and eventual lifting}
\begin{dcomment}{development}{p49:La vérification de d) par AQT}{The filtered criterion and eventual lifting}
Corollary~3.4 converts connectedness of comma categories into the two concrete operations needed for filtered colimits: reaching the image of $u$, and equalizing parallel arrows after reaching it. The proof of parts~(b) and~(c) is present on the page. There is one harmless slip in the statement of~(c): its last assertion concerns filteredness of $I$, not the already assumed filteredness of $I'$. Part~(d) is left under the abbreviation AQT. Its proof uses full faithfulness in two different places.

One variance qualification is essential. Definition~2.2 called the completion of a cospan by a common predecessor PS~1 and explicitly called its right-hand dual ``colocalement faiblement filtrante.'' Here part~(d) and footnote~57 instead use ``pr\'efiltrante'' or ``localement faiblement filtrante'' for the right-hand condition: a span must admit a commuting square with a common successor. The proof requires that interpretation. Under the earlier left-hand reading, part~(d) would be false: the poset $I=(0\longrightarrow1,\,0\longrightarrow2)$ satisfies left PS~1, and the full inclusion $\{1,2\}\hookrightarrow I$ satisfies F~1, but its comma category under $0$ is a discrete two-point category. Thus this is a necessary clarification of the printed terminology, not a variance-neutral use of PS~1.

Assume that $I$ satisfies the right-hand Ore condition PS~1, that $u:I'\longrightarrow I$ is fully faithful, and that F~1 holds. Given objects $(a,\alpha:i\longrightarrow u a)$ and $(b,\beta:i\longrightarrow u b)$ of $i\backslash I'$, PS~1 supplies arrows $r:u a\longrightarrow j$ and $s:u b\longrightarrow j$ with $r\alpha=s\beta$. Choose $j\longrightarrow u c$ by F~1. Fullness lifts the two resulting arrows $u a\longrightarrow u c$ and $u b\longrightarrow u c$ to arrows in $I'$, and they give a common successor of the original two objects in $i\backslash I'$. F~1 also makes that comma category nonempty. It is therefore weakly filtered, hence connected, and $u$ is cofinal.

To transfer PS~1 itself to $I'$, start with a span $a\longrightarrow b$, $a\longrightarrow c$ in $I'$. Complete its image in $I$ to a commuting square, reach an object in the image of $u$ by F~1, and lift the two new arrows by fullness. Faithfulness ensures that the lifted square commutes. Thus $I'$ satisfies PS~1. If $I$ is filtered, a parallel pair $r,s:a\rightrightarrows b$ in $I'$ can likewise be equalized in $I$, followed by an arrow to $u c$ and lifted to an equalizer arrow $b\longrightarrow c$ in $I'$. Faithfulness again supplies the equality. Nonemptiness follows from F~1, so $I'$ is filtered. The same reasoning, or the discrete-opfibration projection $i\backslash I'\longrightarrow I'$, makes each comma category filtered.

For locally filtered categories, the additional condition F~0 in footnotes~57--58 ensures that objects whose images lie in the same component already lie in the same source component. The arguments can then be carried out in that component. Without this component condition, F~1 and F~2 do not prevent several disconnected copies of one indexing category from being mapped onto it. This is why the global filtered statement and its local variant have different hypotheses.

Footnote~55 asks for a statement encompassing this corollary and Chapter~XVIII. The exact later result is Lemma~4.9.6.3, part~1, manuscript p.~224, whose proof is in \S4.9.6.7. Its additional condition is eventual lifting:
\[
 \text{F~3:}\quad
 \theta:u(a)\longrightarrow u(b)
 \ \Longrightarrow\
 \exists\,v:b\longrightarrow c,\ w:a\longrightarrow c,
 \qquad u(v)\theta=u(w).
\]
For a faithful functor into a filtered category, the lemma identifies F~1 and F~3 with filteredness of the source together with cofinality \citep[Lemma~4.9.6.3(1), pp.~81--82]{Grothendieck1990XVIII}. The following finite formulation combines the two arguments, retaining only the right Ore hypothesis where that suffices. It is our proof-level synthesis of the cited passages, without a claim of priority.

\textbf{A finite comparison with eventual lifting.}\ours
\textit{If the target satisfies right PS~1, F~1 and F~3 make $u$ cofinal. If $u$ is faithful, its source also satisfies right PS~1 and is filtered whenever the target is filtered.}

Indeed, given $(a,\alpha)$ and $(b,\beta)$ under $i$, right PS~1 followed by F~1 produces $\rho_a:u(a)\to u(c)$ and $\rho_b:u(b)\to u(c)$ with $\rho_a\alpha=\rho_b\beta$. Apply F~3 first to $\rho_a$, obtaining $v:c\to d$, $w:a\to d$, and then to $u(v)\rho_b$, obtaining $t:d\to e$, $z:b\to e$. The equality
\[
 u(tw)\alpha=u(t)u(v)\rho_a\alpha
              =u(t)u(v)\rho_b\beta=u(z)\beta
\]
gives a common successor in the comma category. F~1 gives nonemptiness, hence cofinality. For a span in the source, this construction makes its two composites equal after applying $u$; faithfulness makes them equal before applying it. If the target is filtered, equalize the images of any parallel pair, reach another image object by F~1, and lift the resulting arrow eventually by F~3. Faithfulness then gives an equalizing arrow in the source. Cofinality supplies connectedness and nonemptiness, completing filteredness. Conversely, when the source is filtered and $u$ cofinal, the already proved filteredness of $u(a)\backslash I'$ gives F~3 by taking a common successor of $(a,1)$ and $(b,\theta)$.

Full faithfulness makes F~3 automatic, recovering part~(d). The source-filtered case is the F~1--F~2 criterion in part~(c). Thus the finite comparison requested in footnote~55 is explicit. Part~2 of the later lemma strengthens F~3 to simultaneous lifting for families of bounded cardinality and imposes the corresponding larger filtering hypotheses; that cardinal refinement is a separate assertion, not a consequence of the finite conditions alone.
\end{dcomment}


\printedpage{32}
% head: A. Grothendieck : Dérivateurs, Ch. II, Édition des Universités de Montpellier II et Paris VII
% foot: Version 8 juillet 2001
\section*{4\quad Foncteurs biasphériques}

Soit encore
\[
u : I' \longrightarrow I
\]
une flèche de Cat, soit $\cM$ une catégorie, d'où
\[
u^* : \cM(I) \longrightarrow \cM(I');
\tag*{(4.1)}
\]
nous cherchons des conditions sur $u$, $\cM$ pour que ce foncteur soit pleinement fidèle. Si le foncteur adjoint à gauche $u_!$ (resp. adjoint à droite $u_*$) existe, on sait qu'il revient au même de dire que le morphisme d'adjonction
\[
u_!u^*F \longrightarrow F
\tag*{(4.2)}
\]
resp.
\[
u_*u^*F \longleftarrow F
\tag*{(4.2 bis)}
\]
soit \texttt{[un]} isomorphisme pour tout $F$ dans $\cM(I)$. Disons alors que $u$ est $\cM$-connexe, ou $\cM$-biasphérique, cette terminologie sous-entendant que c'est là une propriété qui se réfère au dérivateur $\mathbf{D}_{\cM} = (S \longmapsto \cM(S))$ sur Cat défini par $\cM$. Notons que
\[
\begin{aligned}
&\text{Si $u$ est $\cM$-connexe, alors il est aussi $\cM(S)$-connexe, pour tout $S$ dans Cat,}\\
&\text{puisque $\cM(S)(I) \simeq \cM(I)(S) \defeq \uHom(S^o, \cM(I))$,}
\end{aligned}
\tag*{(4.3)}
\]

\msp{page 50}
et
\[
\text{Si $\cM' \hookrightarrow \cM$ est pleinement fidèle, alors $\cM$-connexe implique $\cM'$-connexe,}
\tag*{(4.4)}
\]
enfin, inversement,
\[
\begin{aligned}
&\text{pour que $u$ soit $\cM$-connexe, il faut et il suffit que pour toute \textit{petite}}\\
&\text{sous-catégorie pleine $\cM'$ de $\cM$, $u$ soit $\cM'$-connexe.}
\end{aligned}
\tag*{(4.5)}
\]
Comme toute petite $\cM'$ se plonge dans $\cM'^\wedge = \operatorname{Ens}(\cM')$, il résulte de (4.3), (4.4), (4.5) que
\[
\text{Si $u$ est Ens-connexe, il est $\cM$-connexe pour toute catégorie $\cM$,}
\tag*{(4.6)}
\]
car c'est le cas si $\cM$ est petite par (4.3), donc pour tout $\cM$ par (4.5).

Disons qu'un foncteur $u$ dans Cat est $\cM$-\textit{cofinal} ou $\cM$-coasphérique si l'homomorphisme fonctoriel
\[
\varprojlim_{I} F \longrightarrow \varprojlim_{I'} u^*(F)
\tag*{(4.7)}
\]
(flèche dans $\uHom(\cM(I), \cM)$ si les $\varprojlim$ de types $I$, $I'$ existent dans $\cM$, sinon flèche dans $\uHom(\cM(I), \cM^\wedge)$) est un isomorphisme. En se ramenant au cas de petites catégories $\cM$, chacune plongée dans $\cM^\wedge$, on voit par le critère (4.2 bis) que
\[
\text{$u$ est $\cM$-connexe si pour tout $i \in \Ob I$,}
\tag*{(4.8)}
\]

\printedpage{33}
% head: A. Grothendieck : Dérivateurs, Ch. II, Édition des Universités de Montpellier II et Paris VII
% foot: Version 8 juillet 2001
\msp{page 51}
le foncteur induit $I'/i \longrightarrow I/i$ est $\cM$-cofinal. La réciproque n'est pas claire, car $u$ est $\cM$-connexe si et seulement si pour tout $i \in \Ob I$ et tout $\cM$-préfaisceau $F_i$ sur $I/i$ \textit{induit par un préfaisceau $F$ sur $I$,} l'homomorphisme $\varprojlim\limits_{I/i} F_i \longrightarrow \varprojlim\limits_{I'/i} (u/i)^*(F/i)$ est \texttt{[un]} isomorphisme. On ne peut donc affirmer d'emblée que cela implique itou pour \textit{tout} $F_i$ sur $I/i$, que si on sait que $(u/i)^* : \cM(I) \longrightarrow \cM(I/i)$ est essentiellement surjectif. Il en est ainsi pour tout $i$ quand $I$ est \textit{ordonné,} car alors $I/i \longrightarrow I$ est pleinement fidèle, donc $\cM(I/i) \longrightarrow \cM(I)$ \texttt{[est]} pleinement fidèle, donc $\cM(I) \longrightarrow \cM(I/i)$ \texttt{[est]} un foncteur de localisation, en particulier essentiellement surjectif.

\msp{page 52}
Mais quand $\cM = (\text{Ens})$, cette perplexité disparaît. On trouve la

\dgrblock{dgrb114}{The missing pointwise limit hypothesis}
\begin{dcomment}{exposition}{p51}{The missing pointwise limit hypothesis}
Grothendieck isolates a real difference between testing every coefficient on a slice and testing only coefficients restricted from the ambient category. There is also a second issue here: full faithfulness for coefficients in $\mathcal M$ does not imply full faithfulness after replacing $\mathcal M$ by its presheaf completion. The sufficient implication (4.8) is valid. The converse discussed on page~51, including its asserted ordered case, requires hypotheses ensuring that the relevant pointwise Kan-extension limits exist in $\mathcal M$.

Take $I=[1]=(0\longrightarrow1)$, $I'=\{1\}$, and let $\mathcal M$ be the discrete category on two objects $m,n$. Every $\mathcal M$-valued presheaf on $I$ is constant, so restriction $\mathcal M(I)\longrightarrow\mathcal M(I')$ is an equivalence. Nevertheless, at $i=0$ the induced functor on slices is $\varnothing\longrightarrow *$. In $\mathcal M^{\wedge}\simeq\mathrm{Set}\times\mathrm{Set}$, the Yoneda image of $m$ is $(1,\varnothing)$, whereas the limit of the empty diagram is $(1,1)$. Thus the comparison
\[
 \varprojlim_{I/0}y(m)=(1,\varnothing)
 \longrightarrow\varprojlim_{I'/0}y(m)=(1,1)
\]
is not invertible. This coefficient on $I/0$ even extends to $I$. Hence the intermediate assertion that $\mathcal M$-connectedness is equivalent to the slice comparison for restricted coefficients fails without pointwise limit hypotheses.

Existence of an abstract adjoint alone is insufficient: in this example restriction is an equivalence and has both adjoints. What fails is the formula that would compute its right adjoint by an empty limit inside $\mathcal M$. Completeness of $\mathcal M$, or existence of all the particular comma limits involved, supplies the needed repair.

The extension observation for posets is independently correct. Identify $I/i$ with $A=\{x\mid x\leq i\}$. The map $r:I\longrightarrow A$ sending $x$ to itself when $x\leq i$ and to $i$ otherwise is order preserving: no arrow can enter $A$ from outside, and an arrow leaving $A$ is sent to $x\leq i$. With $j:A\hookrightarrow I$, one has $rj=1_A$, so every coefficient on $A$ extends by $r^{\mathrm{op}}$. Under the pointwise limit hypotheses this proves the ordered-case converse. It does not remove those hypotheses. Proposition~4.1 subsequently quantifies over all coefficient categories and can use sets from the start; its intended theorem is unaffected by this coefficient-specific error.
\end{dcomment}


\noindent\textbf{Proposition 4.1.} \textit{Soit $u : I' \longrightarrow I$ une flèche dans} Cat. \textit{Conditions équivalentes :}
\begin{itemize}[label={}]
\item[\textit{a)}] \textit{Pour toute catégorie $\cM$, $u^*_{\cM} : \cM(I) \longrightarrow \cM(I')$ est pleinement fidèle (i.e. $u$ est $\cM$-connexe).}
\item[\textit{a$'$)}] \textit{Pour toute catégorie $\cM$ où les limites inductives de type $I$, $I'$ existent, la flèche d'adjonction pour $F$ dans $\cM(I)$} % Transcription corrected: official printed33 reads a prime.
\[
u_!u^*F \xrightarrow{\sim} F
\]
\textit{est} \texttt{[un]} \textit{isomorphisme.}
\item[\textit{a$''$)}] \textit{Dual de a$'$), avec} \texttt{[la]} \textit{flèche d'adjonction}
\[
u_*u^* \xleftarrow{\sim} F. % [sic] u_*u^* -> u_*u^*F, the object F is missing on the left side (compare a') and (4.2 bis))
\]
\item[\textit{b) b$'$) b$''$)}] \textit{Comme a) a$'$) a$''$), avec $\cM =$} (Ens).
\item[\textit{c)}] \textit{Comme b$'$), avec $F$ représentable, i.e. de la forme $j \longmapsto$} $\Hom_I(j, i)$.
\item[\textit{d)}] \textit{Pour tout $I$ dans $I$, $I'/i \longrightarrow I/i$ est cofinal.} % [sic] Pour tout I dans I -> Pour tout i dans I, the variable of the statement is i (compare d'))
\item[\textit{d$'$)}] \textit{Pour tout $i$ dans $I$, $i\bs I' \longrightarrow i\bs I$ est cofinal.}
\item[\textit{d$''$)}] \textit{Pour toute flèche $i_0 \xrightarrow{\alpha} i_1$ dans $I$, la catégorie} $i_0\bs \underset{\alpha}{I'}/i_1 \defeq I' \times_I \underbrace{i_0\bs \underset{\alpha}{I}/i_1}_{\substack{\text{catégorie biïnduite}\\ (i_0, \alpha)\bs(I/i_1) \simeq (i_0\bs I)/(i_1, \alpha)}}$ \qquad \textit{est}

0-\textit{connexe.}
\end{itemize}

\msp{page 53}
\textsc{Démonstration.} On sait déjà, par les préliminaires, l'équivalence de a, a$'$, a$''$, b, b$'$, b$''$, soit A cette condition - nous dirons que $u$ est $W_0$-\textit{connexe} ou $W_0$-\textit{biasphérique,} si ces conditions équivalentes sont satisfaites. On a de plus l'autodualité
\[
\text{$u$ est $W_0$-connexe si et seulement si $u^o$ l'est,}
\tag*{(4.8)} % [sic] (4.8) -> another number, the label (4.8) is already used for a display on p. 32
\]
car on a évidemment
\[
\text{$u$ est $\cM$-connexe si et seulement si $u^o$ est $\cM^o$-connexe.}
\tag*{(4.9)}
\]


\dgrblock{dgrb115}{Biasphericity, lax epimorphisms, and factorization categories}
\begin{dcomment}{later}{p53}{Biasphericity, lax epimorphisms, and factorization categories}
The new question in Section~4 is whether a diagram and all its natural transformations can be recovered after restriction. The answer is governed by factorizations of arrows, just as ordinary cofinality was governed by objects of comma categories. This is the ordinary version of the factorization test encountered in the commentary on Chapter~I, manuscript p.~49 (``The factorization category controls base change''). The proof of Proposition~4.1 continues beyond these pages; we give the complete argument for its corrected statement here.

For $u:A\longrightarrow B$ and $\gamma:b_0\longrightarrow b_1$, let $\mathcal F_u(\gamma)$ have objects
\[
 (a,p,q),\qquad b_0\xrightarrow{p}u a\xrightarrow{q}b_1,
 \quad qp=\gamma,
\]
and arrows $h:a\longrightarrow a'$ satisfying $u(h)p=p'$ and $q'u(h)=q$. Then the following conditions are equivalent: restriction on presheaves is fully faithful for every coefficient category; it is fully faithful for sets; and every $\mathcal F_u(\gamma)$ is nonempty and connected.

To prove this, first work with sets, where both Kan extensions exist. Full faithfulness of $u^*$ is equivalent to the counit $u_!u^*F\longrightarrow F$ being an isomorphism. Both sides preserve colimits in $F$, and every presheaf is a colimit of representables, so it suffices to take $F=h_{b_1}=B(-,b_1)$. At $b_0$ the pointwise colimit formula gives
\[
 (u_!u^*h_{b_1})(b_0)
 \cong\coprod_{\gamma\in B(b_0,b_1)}
             \pi_0\mathcal F_u(\gamma).
\]
Indeed, the generators are pairs $b_0\longrightarrow u a\longrightarrow b_1$, and the colimit relations are exactly the zigzags induced by arrows of $\mathcal F_u(\gamma)$. Composition labels every such relation by the same $\gamma$. The counit sends a connected component to that composite. It is bijective if and only if each indicated category is nonempty and connected.

The set-valued criterion implies the version for arbitrary coefficients. It first implies full faithfulness for coefficients in every presheaf category, by applying the set-valued statement in each variable. A small category $\mathcal M$ embeds fully faithfully into its presheaf category by Yoneda; the full-faithfulness property descends to this full subcategory. For a possibly large locally small $\mathcal M$, any two diagrams on the small category $B$ take values in a small full subcategory. Applying the preceding argument there descends every natural transformation and proves uniqueness. This establishes the reductions (4.3)--(4.6), with the usual universe enlargement when necessary.

For fixed $b_1$, the comma category of the induced functor $A/b_1\longrightarrow B/b_1$ under $(b_0,\gamma)$ is $\mathcal F_u(\gamma)$. Thus condition~(d) is equivalent to the factorization criterion. Its dual says that
\[
 b_0\backslash A\longrightarrow b_0\backslash B
 \quad\text{is \emph{co-cofinal}, that is, initial.}
\]
The edition instead prints ``cofinal'' in~(d$'$). That cannot be retained as a mathematical equivalence. The inclusion $u:\{1\}\hookrightarrow[1]$ satisfies the printed condition: for $i=0$ it induces the final-vertex inclusion into $[1]$, and for $i=1$ it induces the identity of a point. Yet restriction on presheaves is not faithful: take $F(1)=\varnothing$, $F(0)=\{0,1\}$; its identity and the automorphism transposing $0,1$ become the same map after restriction. The corrected dual condition has comma categories opposite to $\mathcal F_u(\gamma)$, so it gives exactly the required connectedness. The two adjunction formulations and self-duality now follow, proving all the intended equivalences.

In later terminology these functors are \emph{lax epimorphisms}, also called absolutely dense functors. \citet{AdamekElBashirSobralVelebil2001} identify them by full faithfulness of set-valued restriction and by connected categories of factorizations; the latter characterization is restated explicitly in \citet[Theorem~4.1]{NunesSousa2022}. For derivators, \citet[Theorem~6.10]{GrothStovicek2016Trees} replace connectedness by weak contractibility of the same factorization categories and obtain \emph{homotopical epimorphisms}. Thus the arrowwise obstruction persists; the homotopy theory changes the required strength of its vanishing.

A quotient group displays the difference concretely. For a surjective homomorphism $\rho:H\twoheadrightarrow G$, restriction of $G$-sets to $H$-sets is fully faithful, because equivariance under $H$ is equivalent to equivariance under its image $G$. For any $g\in G$, the category $\mathcal F_{B\rho}(g)$ is the transitive action groupoid with object set $G$, on which $H$ acts through $\rho$; its stabilizer is $\ker\rho$. It is therefore equivalent to $B\ker\rho$. The functor $B\rho$ is an ordinary lax epimorphism, but it is a homotopical epimorphism only when $\ker\rho$ is trivial. With $\rho:\mathbf Z\twoheadrightarrow\mathbf Z/n\mathbf Z$, the obstruction has the homotopy type of a circle. The full-faithfulness statement about ordinary diagrams is correct; it simply does not recover that circle.
\end{dcomment}
\begin{otherlanguage}{english}
\dgroverview{II-a}{Chapter II \textperiodcentered\ \textup{II-a} \textperiodcentered\ Overview}
\begingroup
\selectlanguage{english}
\noindent\textbf{Cofinality, connected objects, and the recovery of diagrams.}

Chapter~II begins by asking when a change of indexing category leaves sections unchanged \citep{Grothendieck1990II}. It ends this first part with the stronger question of when restriction recovers a whole diagram and every natural transformation between diagrams. Between these questions Grothendieck develops a topos-theoretic account of filtering, examines the homotopy information lost by ordinary connectedness, and specializes the language of derivators to sets. The movement is deliberate. Ordinary cofinality is not merely an elementary preliminary to the derived theory: its proof identifies the object which must be derived, and its limitations identify the higher information which a derivator must retain.

The conventions are essential. A coefficient on $C$ is a contravariant functor, and the category denoted $c\backslash C'$ has objects $c\to\alpha(c')$. Thus a cofinal functor in these pages is final in the usual covariant terminology; its opposite is initial for the diagrams whose limits are being taken. Under this convention the first question has a representability answer. The presheaf
\[
 K_\alpha(c)=\pi_0(c\backslash C')= (\alpha_!1_{C'})(c)
\]
represents $F\mapsto\Gamma(C',\alpha^*F)$. The comparison with $\Gamma(C,F)$ is induced by $K_\alpha\to1_C$. It is an isomorphism for all $F$ exactly when every indicated comma category is nonempty and connected. This is the classical cofinality theorem of SGA~4, Expos\'e~I, \S8.1 \citep{SGA4I1972}. The initial weak cofiltering hypothesis is sufficient, and Grothendieck replaces it by the exact connectedness criterion as the argument develops.

The same proof suggests its own derived form. For a commutative ring $k$, replace the component presheaf by the derived weight $\mathbf L\alpha_!k_{C'}$. Its value at $c$ is represented by chains on $N(c\backslash C')$. Derived adjunction then gives
\[
 \mathbf R\Gamma(C',\alpha^*F)
 \simeq\mathbf R\operatorname{Hom}_{C;k}
       (\mathbf L\alpha_!k_{C'},F).
\]
Consequently weak contractibility of all the comma nerves implies cohomological cofinality, answering the question on manuscript page~5 affirmatively. For these fixed linear coefficients, their $k$-acyclicity already suffices. For spaces the weight is the corresponding presheaf of classifying spaces, and the mapping-space adjunction gives the homotopical statement. Modern higher category theory places this argument within the general criterion for cofinality \citep[Theorem~4.1.3.1]{Lurie2009HTT}. The commentary constructs the linear comparison by a bar resolution, so the answer does not depend on the internal-Hom formula subsequently proposed in the manuscript.

That formula is a genuine point of failure. The global derived Hom belongs to the category of coefficients at a point; a sheaf-valued internal Hom belongs to the diagram category. Applying global sections can identify two such objects without identifying the objects themselves. For the terminal-vertex inclusion $i:\{1\}\hookrightarrow[1]$, a presheaf is an arrow $F_1\to F_0$, and
\[
 i_!A=(A\xrightarrow{1}A),\qquad i_*A=(A\to *).
\]
The comma categories for cofinality are points, and global sections are $F_1$. Nevertheless $i_*i^*F$ need not be $F$, while $\underline{\operatorname{Hom}}(i_!1,F)=F$. This disproves the unrestricted internal identification in sets, hence already before derived complications arise. The correct internal comparison is equivalent to an additional projection-formula property, whose canonical arrow has the opposite direction from the one displayed in the manuscript. Our point-inclusion calculation identifies its obstruction exactly: for a fixed presheaf, every restriction along an arrow into the chosen object must be invertible. The formula holds for all coefficients precisely when all those incoming arrows are isomorphisms.

The turn to toposes is motivated by a different problem: how to recognize the equalization part of filtering without inspecting individual arrows. An irreducible presheaf topos records common predecessors of objects, a condition visible in the underlying preorder. Local irreducibility records common predecessors compatible with a prescribed pair of arrows, namely PS~1. It still does not supply PS~2, the equalization of arbitrary parallel arrows. The parallel-arrow category $0\rightrightarrows1$ separates the first two conditions. A nontrivial group category $BG$ separates local irreducibility from local cofilteredness: every representable slice is equivalent to a point, but cancellation prevents the equalization of distinct group elements. Several preliminary identifications in the manuscript are explicitly withdrawn by its annotations. The audit preserves this sequence of correction and distinguishes it from errors of transcription.

Total $0$-connectedness adds preservation of products by the component functor. For $\widehat X=\operatorname{Fun}(X^{\mathrm{op}},\mathrm{Set})$, it says that $X^{\mathrm{op}}$ is an ordinary sifted category. Sifted colimits preserve finite products; filtered colimits preserve finite limits, including equalizers \citep{AdamekRosickyVitale2010}. The simplex category supplies the instructive distinction. The products of its representables have connected categories of elements, but the two vertices of $[1]$ admit no common predecessor over $[1]$. Thus its presheaf topos is totally $0$-connected without being locally irreducible. Nothing about this ordinary product test alone makes all the relevant nerves weakly contractible.

The surviving theorem in 2.6 makes the product condition hereditary under connected slices. Its equivalent formulations are (a), (a'), (b), and (d); the generator-slice conditions (c) and (c') are strictly weaker. The diagonal $b\to b\times b$ explains the distinction. Once $b\times b$ is connected and the topos is locally irreducible, its pullback against $(u,v):a\to b\times b$ is noninitial and yields an equalizing arrow. Conversely, PS~1 and PS~2 lift to categories of elements, giving connected products in every connected slice. The commentary supplies the omitted generator-component lemma and the successive equalizations in this argument. It also restores connectedness of the base where the printed formulation of 2.8 omits it.

On a connected topos satisfying this hereditary condition, $\pi_0$ preserves finite limits. This makes the later cohomological argument precise. The induction must run over connected, locally totally $0$-connected toposes, since it repeatedly passes to connected components of slice toposes. The degree-zero row of the covering spectral sequence is the augmented complex of functions $I^{p+1}\to M$, where $I$ is the component set of a covering object. Inserting a chosen element of $I$ contracts that complex. The stronger homotopical conclusion follows from hypercovers: applying the exact component functor to their matching-object diagrams gives boundary fillers in every degree, so the resulting simplicial sets are weakly contractible. This connects the argument to the classical hypercover construction of homotopy types \citep[\S\S8--9]{ArtinMazur1969}. It concerns constant coefficients and shape. It does not trivialize arbitrary sheaf cohomology, as the worked calculation $H^1(\mathbf P^1_k,\mathcal O(-2))\cong k$ illustrates.

The ind-object reformulation in 2.15--2.18 returns to a classical theorem: a presheaf is a filtered colimit of representables exactly when its category of elements is filtered \citep[Expos\'e~I, Theorem~8.3.3]{SGA4I1972}. Taking the terminal presheaf recovers filteredness of the indexing category. Taking its ordinary left Kan extension along an arbitrary functor shows that the resulting category of elements is again filtered. The modern minimal-localizer theorem supplies the asphericity conclusion for every fundamental localizer, without the additional limit axiom used in the historical formulation \citep[Theorem~2.2.11]{Cisinski2004Minimal}.

The converse question on page~36 has a stronger quantifier than later uses of the phrase ``totally aspherical''. In the commentary we denote by $T(A)$ the condition that every ordinary cofinal target of $A$ have weakly contractible nerve. Filtered categories satisfy it. The attempted proof that ordinary siftedness suffices is retracted on pages~36bis--36ter; our explicit completion of the counterexample uses the simplicial circle and a one-point collage whose classifying space is $S^2$. We also prove a necessary common-successor condition for $T(A)$. Its obstruction is a four-object category with a cofinal three-object subcategory and a nonzero integral two-cycle. This settles $T(A)\Leftrightarrow A\text{ filtered}$ for preorders. For general categories the remaining issue is equalization of parallel arrows. We did not locate a resolution under this exact universal quantifier in the sources checked; this is a report of that bounded search, not a declaration of a recognized open problem. Results developed here are marked as the commentators' work and carry no claim of priority.

Section~3 shows how much of the earlier derivator terminology can be tested with sets. Pointwise localization of category-valued diagrams at component equivalences yields the represented derivator of sets. Properness and smoothness become connectedness of strict lift categories. A universal $W_0$-equivalence can be tested over vertices and arrows; a strong $W_0$-fibration is characterized component by component by those universal equivalences. Coverings explain why the weak condition is weaker: monodromy can connect different fibre components without making each fibre connected. The comparison with the commentary on Chapter~I is made at the precise lifting and monodromy arguments, rather than repeating its whole discussion.

The final transition is from preservation of global sections to full faithfulness of restriction. Proposition~4.1 characterizes the latter by nonempty connected categories of factorizations
\[
 b_0\longrightarrow u(a)\longrightarrow b_1
\]
of every arrow $b_0\to b_1$. These are the lax epimorphisms, or absolutely dense functors, of the later literature \citep{AdamekElBashirSobralVelebil2001}. The homotopical version requires weakly contractible factorization categories \citep[Theorem~6.10]{GrothStovicek2016Trees}. A quotient of groups gives a particularly concrete distinction: the factorization category has the homotopy type of the classifying space of the kernel. Two qualifications belong to the audit. The printed coslice condition in 4.1(d') must be initial, not final. Also the coefficient-specific discussion on page~51 invokes pointwise Kan-extension formulas without the necessary comma limits; a discrete coefficient category on two objects gives a counterexample even when the indexing category is ordered. The corrected theorem quantified over all coefficient categories is established by its set-valued case and Yoneda embeddings.

The historical setting is the program of coherent diagram categories described in \citet[\S69]{Grothendieck1983}, alongside the systematic homotopy Kan extensions of \citet{Heller1988}. The commentary keeps the successive formulations, counterexamples, and marginal revisions visible. It distinguishes a proof supplied here from a proof actually written in these pages, and distinguishes a settled theorem from an unsuccessful search for a later treatment.

\endgroup
\end{otherlanguage}
\dgrchunkend

\addtocontents{toc}{\protect\dgrtocmissing{II-b}{Part 2 of 2}}
% DER_III-V: Grothendieck's text with the commentary set in, then the overview, inside a group that applies the chunk's own definitions
\dgrchapter{III}{Chapter III}{Hom externes dans les dérivateurs}
\dgrpart{0}{III-V}{Manuscript pages 1--12}{III-V\enspace\textperiodcentered\enspace Chapter III\enspace\textperiodcentered\enspace manuscript pages 1--12}{Chapter III \textperiodcentered\ \textup{III-V}}
\dgrchunkbegin
% DER_III-V: the definitions of the chunk's preamble, applied inside its group
\let\printedpage\relax\newcommand{\printedpage}[1]{\marginpar{\footnotesize p.~#1}}
\let\msp\relax\newcommand{\msp}[1]{\marginpar{\footnotesize [#1]}}
\let\Hot\relax\newcommand{\Hot}{\mathrm{Hot}}
\let\D\relax\newcommand{\D}{\mathbf{D}}
\let\uHom\relax\newcommand{\uHom}{\underline{\mathrm{Hom}}}
\let\hm\relax\newcommand{\hm}{\mathrm{hom}}
\let\Hom\relax\newcommand{\Hom}{\mathrm{Hom}}
\let\homext\relax\newcommand{\homext}{\mathrm{hom}\,\mathrm{ext}}
\let\HOT\relax\newcommand{\HOT}{\mathrm{HOT}}
\let\Ob\relax\newcommand{\Ob}{\mathrm{Ob}\,}
\let\Ltens\relax\newcommand{\Ltens}{\overset{\mathbf{L}}{\otimes}}
\let\RHom\relax\newcommand{\RHom}{\mathbf{R}\mathrm{Hom}}
\let\ab\relax\newcommand{\ab}{\mathrm{ab}}
\let\defeq\relax\newcommand{\defeq}{\overset{\text{déf}}{=}}
\let\cA\relax\newcommand{\cA}{\mathcal{A}}
\let\bDelta\relax\newcommand{\bDelta}{\boldsymbol{\Delta}}
\let\id\relax\newcommand{\id}{\mathrm{id}}
\let\Fl\relax\newcommand{\Fl}{\operatorname{Fl}}
\let\uCh\relax\newcommand{\uCh}{\underline{\mathrm{Ch}}}
\let\bPhi\relax\newcommand{\bPhi}{\boldsymbol{\Phi}}
\let\uFl\relax\newcommand{\uFl}{\underline{\mathrm{Fl}}}
\let\Ch\relax\newcommand{\Ch}{\underline{\mathrm{Ch}}}
\let\card\relax\newcommand{\card}{\mathrm{card}}
\let\cM\relax\newcommand{\cM}{\mathcal{M}}
\let\Chh\relax\newcommand{\Chh}{\underline{\underline{\mathrm{Ch}}}}
\let\KK\relax\newcommand{\KK}{\underline{\underline{K}}}
\let\dfeq\relax\newcommand{\dfeq}{\overset{\text{déf}}{=}}
\let\Cathot\relax\newcommand{\Cathot}{\mathrm{Cat\,hot}}
\let\Ind\relax\newcommand{\Ind}{\operatorname{Ind}}
\let\Ens\relax\newcommand{\Ens}{\mathrm{Ens}}
\let\Cat\relax\newcommand{\Cat}{\mathrm{Cat}}
\let\uuK\relax\newcommand{\uuK}{\underline{\underline{K}}}

% ---- from Der-ChIII-567_p001-005.tex
% ===================== page 1 (title page) =====================
\printedpage{--}
% no page number printed on this page
% head: (none)

\begin{center}
{\LARGE\sffamily LES DÉRIVATEURS}

\bigskip

\textbf{ALEXANDRE GROTHENDIECK}

\bigskip\bigskip

\textbf{Chapitre III}

\medskip

\textbf{Hom externes dans les dérivateurs}
\end{center}

\clearpage

% ===================== page 2 (editorial note) =====================
\printedpage{--}
% no page number printed on this page
% head: (none)

Ce texte a été déchiffré et transcrit en \LaTeXe{} par M. Künzer. Il a été édité par M. Künzer, J. Malgoire, et G. Maltsiniotis. La transcription est aussi fidèle que possible au manuscrit. Pour les quelques corrections évidentes, ou rares commentaires des éditeurs, ainsi que pour la numérotation originale des pages du manuscrit, les caractères de machine à écrire \texttt{[typewriter]} entre crochets sont utilisés. Un point d'interrogation entre crochets signifie que l'on n'est pas sûr du mot qui précède.

Cette édition est provisoire. Les remarques, commentaires et corrections sont bienvenus. Envoyer un message à :

\textbf{maltsin@math.jussieu.fr}
\begin{flushright}
G. Maltsiniotis\qquad\qquad
\end{flushright}

\clearpage

% ===================== page 3 (printed 1) =====================
\printedpage{1}
% head: A. Grothendieck : Dérivateurs, Ch. III, Édition des Universités de Montpellier II et Paris VII
% foot: Version 2 février 2002

\msp{page 1}

\section*{1\quad Relation des dérivateurs avec la catégorie $\Hot$}

Soit $\D$ un dérivateur, $\mathcal{A}$ sa catégorie fondamentale,
\[
\mathcal{A} = \D(e).
\]
Je me propose de définir, si faire se peut, des foncteurs
\[
\left\{
\begin{array}{rccccccl}
\Hot & \times & \mathcal{A} & \xrightarrow{\pi_{\D}} & \mathcal{A} & (\xi, M) & \longmapsto & \xi \times M \ \text{ ou } \ \xi \otimes M \quad \text{(suivant les cas)}\\
\Hot^o & \times & \mathcal{A} & \xrightarrow{\varpi_{\D}} & \mathcal{A} & (\xi, M) & \longmapsto & \hm(\xi, M).
\end{array}
\right.
\tag*{(1)}
\]
Ces foncteurs doivent s'échanger par dualité : si on prend les foncteurs $\pi^o$, $\varpi^o$ sur les catégories opposées, on trouve
\[
\left\{
\begin{array}{rcccc}
\Hot & \times & \mathcal{A}^o & \xrightarrow{\varpi^o_{\D}} & \mathcal{A}^o\\
\Hot^o & \times & \mathcal{A}^o & \xrightarrow{\pi^o_{\D}} & \mathcal{A}^o,
\end{array}
\right.
\tag*{(1$'$)}
\]
et on doit avoir
\[
\left\{
\begin{array}{rcl}
\varpi^o_{\D} & = & \pi_{\D^o}\\
\pi^o_{\D} & = & \varpi_{\D^o}.
\end{array}
\right.
\]
À vrai dire, considérant (1) comme des foncteurs
\[
\begin{array}{rcl}
\Hot & \xrightarrow{\pi_{\D}} & \uHom(\mathcal{A}, \mathcal{A})\\
\Hot^o & \xrightarrow{\varpi_{\D}} & \uHom(\mathcal{A}, \mathcal{A}),
\end{array}
\]
on s'attend à ce que ces foncteurs s'étendent en
\[
\left\{
\begin{array}{rcl}
\Hot & \xrightarrow{\pi_{\D}} & \uHom(\D, \D)\\
\Hot^o & \xrightarrow{\varpi_{\D}} & \uHom(\mathcal{A}, \mathcal{A})
\end{array}
\right.
\tag*{(2)}
\]
\texttt{[plutôt deux fois} $\uHom(\D, \D)$\texttt{]}. Cela signifie, essentiellement, que pour $\xi \in \mathrm{Ob}\,\Hot$ donné, les opérations $\xi \times_{\D} M$ et $\uHom_{\D}(\xi, M)$ \texttt{[plutôt} $\hm_{\D}(\xi, M)$\texttt{]} sont définies non seulement dans $\mathcal{A} = \D(e)$, mais dans tous les $\D(I)$, et ceci de façon à commuter aux foncteurs $f^*$ :

\dgrblock{dgrb147}{Spaces acting on coefficients}
\begin{dcomment}{exposition}{p1:Ces foncteurs doivent s'échanger par dualité}{Spaces acting on coefficients}
Grothendieck is looking for the part of ordinary homotopy theory that must act on every theory of coefficients. A space should provide an instruction for assembling copies of an object, while the same space, with the variance reversed, should provide an instruction for taking their limit. This is the reason for placing the two operations together before attempting to construct an external Hom.

Write \(B X=|N X|\) for the classifying space of a small category, and temporarily assume that all the categories used below belong to the permitted diagram category. If \(p_X:X\to e\), put
\[
 T_X(M)=(p_X)_!p_X^*M,\qquad
 C_X(M)=(p_X)_*p_X^*M.
\]
These expressions are already defined by the derivator. The problem is to show that they depend on \(X\) only through \(B X\), and then to understand their compatibility as \(X\) and \(M\) vary. In the derivator of spaces they are \(B X\times M\) and the derived function space \(\operatorname{Map}(B X,M)\). In a derived category they are derived tensor and Hom constructions with the integral chains of \(B X\). The common formalism therefore sees more than the set of components of \(X\).

There are three different Hom constructions in play. The ordinary set \(\Hom_{\D(e)}(M,N)\) records morphisms after localization. The cotensor \(C_X(N)\) is another coefficient object, hence belongs to \(\D(e)\). The proposed external Hom \(\homext(M,N)\) is a homotopy type. Its intended property is
\[
 \Hom_{\D(e)}(T_XM,N)
 \cong [B X,\homext(M,N)].
\]
In a model presentation this last object is the derived mapping space. It is not the internal Hom of a monoidal category of coefficients: no multiplication of two coefficient objects has been assumed. For example, the derived category of modules over a noncommutative ring still has an external mapping space, although a tensor product of two left modules over that ring has no such general meaning.

The extension to every \(\D(I)\) is essential. With \(p:X\times I\to I\), the same formulas define \(T_{X,I}=p_!p^*\) and \(C_{X,I}=p_*p^*\). Base change along products supplies their compatibility with restriction \(f^*\). Thus the request in (2) is for morphisms of derivators, including their pseudonaturality isomorphisms, rather than unrelated endofunctors of the individual categories. The passage from \(\D(e)\) to the shifted derivator \(\D_I\) is the mechanism already used in the commentary on Chapter~I, manuscript p.~4 (\emph{From derived coefficients to diagram categories}); here its new role is to make the action itself coherent in the diagram variable.

For a represented derivator, the calculation reduces to an instructive limiting case:
\[
 T_X(M)=\coprod_{\pi_0 X}M,\qquad
 C_X(M)=\prod_{\pi_0 X}M.
\]
Indeed a cone on a constant ordinary diagram is a family of maps constant on every connected component. The richer dependence on \(B X\) belongs to the homotopy theory, not to the notation \(p_!p^*\) alone.
\end{dcomment}


\msp{page 2}
Si
\[
f : I' \longrightarrow I,
\]
\printedpage{2}
% head: A. Grothendieck : Dérivateurs, Ch. III, Édition des Universités de Montpellier II et Paris VII
% foot: Version 2 février 2002
alors pour $F \in \mathrm{Ob}\,\D(I)$, on doit avoir des isomorphismes canoniques
\[
\left\{
\begin{array}{rcl}
f^*(\xi \times_{\D} M) & \simeq & \xi \times_{\D} f^*(M)\\
f^*\hm_{\D}(\xi, M) & \simeq & \hm_{\D}(\xi, f^*(M)),
\end{array}
\right.
\tag*{(3)}
\]
avec les compatibilités habituelles pour un composé $fg$, relatif à $I'' \xrightarrow{g} I' \xrightarrow{f} I$.

Mais je me place de préférence dans $\mathcal{A}$, ce qui devrait impliquer les constructions similaires dans tous les $\D(I)$, quitte à appliquer ce qui a été dit au dérivateur $\D_I$.

L'intuition algébrique fondamentale est celle-ci, en désignant par
\[
X \longmapsto \mathrm{hot}(X) : (\mathrm{Cat}) \longrightarrow (\Hot)
\]
le foncteur de localisation de (Cat) dans (Hot), obtenu en inversant les flèches dans (Cat) qui sont des `équivalences faibles'. Supposons $X$ dans Diag, soit
\[
p : X \longrightarrow e
\]
le foncteur à valeurs dans la catégorie ponctuelle. On veut que
\[
\left\{
\begin{array}{rcccl}
\mathrm{hot}(X) \times_{\D} M & \simeq & p_!(p^*(M)) & = & H^{\D}_{\bullet}(X, M)\\
\hm_{\D}(\mathrm{hot}(X), M) & \simeq & p_*(p^*(M)) & = & H^{\bullet}_{\D}(X, M)
\end{array}
\right.
\tag*{(4)}
\]
(isomorphismes canoniques).

\msp{page 3}
Ceci donnerait un principe de définition pour les opérations cherchées. Une difficulté vient du fait qu'en général, Diag n'est pas égal à (Cat), et que les types d'homotopie provenant de Diag ne donnent pas (même à isomorphisme près) \textit{tous} les types d'homotopie. Intuitivement, on ne trouve que des types d'homotopie `assez petits'. Mais on peut localiser Diag par rapport aux équivalences faibles (il vaut mieux dire `quasi-isomorphismes'), on trouve alors une catégorie $(\Hot)_0$, et un foncteur canonique
\[
\Hot_0 \longrightarrow \Hot,
\]
si on a de la chance, ce foncteur sera pleinement fidèle, et on essayera de décrire les foncteurs (1) ou (2) en y remplaçant Hot par la sous-catégorie pleine des `petits' types d'homotopie. On peut aussi oublier (provisoirement!) Hot, et travailler avec $\Hot_0$. C'est ce qu'on va faire pour le moment.

Pour pouvoir définir les foncteurs
\[
\left\{
\begin{array}{lccccccl}
\Hot_0 & \times & \mathcal{A} & \longrightarrow & \mathcal{A} & (\xi, F) & \longmapsto & \xi \times_{\D} F\\
\Hot_0^o & \times & \mathcal{A} & \longrightarrow & \mathcal{A} & (\xi, F) & \longmapsto & \hm_{\D}(\xi, F),
\end{array}
\right.
\tag*{(5)}
\]
je note que pour $F$ \texttt{[plutôt} $M$\texttt{]} fixé, $X$ variable dans Diag, les
\printedpage{3}
% head: A. Grothendieck : Dérivateurs, Ch. III, Édition des Universités de Montpellier II et Paris VII
% foot: Version 2 février 2002
\msp{page 4}
derniers membres dans (4)
\[
H^{\D}_{\bullet}(X, M), \qquad H^{\bullet}_{\D}(X, M)
\tag*{(6)}
\]
sont bien, l'un covariant, l'autre contravariant en $X$. Donc il y a , à isomorphisme unique près, une seule paire de bifoncteurs (5) satisfaisant (bifonctoriellement) les isomorphismes (4), si et seulement si les expressions (6), considérées comme fonctorielles en $X$, transforment quasi-isomorphismes en isomorphismes. Cela équivaut donc à la validité du % [sic] "il y a ," -> "il y a,", espace imprimée avant la virgule

\textbf{Théorème} (?) (\footnote{Nouvel axiome ?}). \textit{Soit $f : X' \longrightarrow X$ dans Diag. Si $f$ est un quasi-isomorphisme, c'est une $\D$-équivalence.}

Quand $X$ est la catégorie ponctuelle, ce n'est là autre que l'axiome Der 3 (p. 67). J'ignore si c'est une conséquence des axiomes Der 1 à Der 5 (d). Si ce n'est pas le cas, il faudrait l'ajouter aux axiomes, je crois. Je n'ai pas même prouvé qu'une équivalence de catégories dans Diag (au sens des catégories!) est une $\D$-équivalence!

\dgrblock{dgrb148}{The descent theorem and its domain}
\begin{dcomment}{status}{p4:J'ignore si c'est une conséquence des axiomes}{The descent theorem and its domain}
The question marked \emph{Théorème (?)} is precisely the obstruction to making the preceding formulas into operations of homotopy types. It has a positive answer for derivators on all small categories with the usual homotopy Kan extension axioms. The scope of the indexing category matters: that answer must not be silently substituted for a theorem about every restricted \(\mathrm{Diag}\) occurring in the manuscript.

First, the claimed equivalence with descent is formal. Let \(W\) be the nerve weak equivalences in \(\mathrm{Diag}\). A bifunctor \((X,M)\mapsto T_XM\) factors through \(\mathrm{Diag}[W^{-1}]\) exactly when every natural transformation \(T_f:T_{X'}\to T_X\), \(f\in W\), is invertible. The localization universal property applies also to natural transformations, giving the factorization with its specified comparison uniquely up to unique compatible isomorphism. The same argument applies contravariantly to \(C_X\). The two conditions are equivalent because \(T_X\dashv C_X\), and the comparison transformations are mates: a natural isomorphism between left adjoints has an invertible right mate, and conversely.

The easy case that Grothendieck says he has not yet checked needs no nerve theorem. If \(f:X'\to X\) is an equivalence of categories, a quasi-inverse and its invertible unit and counit give, by the prederivator structure, an equivalence \(f^*:\D(X)\to\D(X')\). It carries \(p_X^*M\) to \(p_{X'}^*M\). Hence it induces bijections
\[
 \Hom_{\D(X)}(p_X^*M,p_X^*N)
 \longrightarrow
 \Hom_{\D(X')}(p_{X'}^*M,p_{X'}^*N),
\]
which, through the two adjunctions, imply exactly the required comparison isomorphisms. This argument is already available for a prederivator possessing the displayed adjoints.

The passage from categorical equivalences to all nerve weak equivalences is substantial. For a derivator on all small categories, its constant-coefficient equivalences form a fundamental localizer \citep[Proposition 3.7]{Cisinski2008}; \citet[Theorem 2.2.11]{Cisinski2004Minimal} proves that nerve weak equivalences form the smallest fundamental localizer. Their combination is precisely \citet[Corollary 3.10]{Cisinski2008}: every nerve weak equivalence is a derivator equivalence. The more structural answer is the same paper's Theorem 3.24 and Corollary 3.26: the derivator of spaces is freely generated under homotopy colimits by a point. Its universal action yields the tensor construction; duality yields the cotensor construction. These results settle the modern full-domain version of the question.

The universal property also identifies the operation, not merely its existence. In its cocontinuous form, evaluation at the point identifies cocontinuous morphisms from the space derivator to \(\D\) with objects of \(\D(e)\). Let \(F_M\) correspond to \(M\). The classifying space \(B X\) is the homotopy colimit of the constant point diagram over \(X\). Preservation of that homotopy colimit forces
\[
 F_M(B X)\cong (p_X)_!p_X^*M.
\]
The equivalence of categories makes this construction functorial in maps \(M\to N\), and the statement for shifted derivators supplies its diagram versions. This explains why the modern theorem answers the particular question on the page: the descended action is constrained to have Grothendieck's formula (4). It also explains why full-domain hypotheses enter the literature. The universal object is the entire space derivator; replacing its diagram domain by an arbitrary subcategory is an additional operation whose effect must be checked.

For a restricted \(\mathrm{Diag}\), two further assertions must be kept separate. The internal localization \(\Hot_0\) exists under the usual size convention, but its comparison with \(\Hot\) need not be fully faithful merely by its definition. Even if it is fully faithful, it need not contain all homotopy types. One can work perfectly well with the action of \(\Hot_0\) once descent has been established there. What one cannot do is invoke Yoneda with test objects outside that category. This is exactly the limitation that returns on manuscript pages 7--10.
\end{dcomment}


Notons la \textit{formule d'adjonction}
\[
\mathrm{Hom}(\xi \times_{\D} M, N) \xrightarrow{\sim} \mathrm{Hom}(M, \hm_D(\xi, N)).
\tag*{(7)}
\] % [sic] \hom_D (D italique maigre) -> \hom_{\mathbf{D}}, le dérivateur est noté en gras partout ailleurs
Si $\xi$ provient de $X \in \mathrm{Diag}$, alors cela s'écrit
\[
\mathrm{Hom}(p_!p^*M, N) \xrightarrow{\sim} \mathrm{Hom}(M, p_*p^*N),
\tag*{(8)}
\]
\msp{page 5}
or par adjonction les deux membres sont canoniquement isomorphes à
\[
\mathrm{Hom}_{\D(X)}(p^*M, p^*N).
\]

\textbf{Exemple.} Prenant pour $\xi$ le type d'homotopie du cercle standard (représenté p.ex. par l'ensemble ordonné
\[
\begin{tikzcd}[row sep=small, column sep=small]
 & b & \\
0 \arrow[ur] \arrow[dr] & & 0' \arrow[ul] \arrow[dl] \ , \\
 & a &
\end{tikzcd}
\]
où on convient de prendre 0 comme origine, et d'orienter de 0 vers $a$), les deux foncteurs en question sont (sauf erreur) les foncteurs suspension $\Sigma$ et `cosuspension' ou foncteur `espace des lacets' $\Omega$. Sous réserve de vérifier que $\Omega F$ coïncide bien avec ce que j'ai décrit ailleurs, à propos des carrés substantiels cartésiens dans $\mathcal{A}$.

\dgrblock{dgrb149}{What the circle actually gives}
\begin{dcomment}{examples}{p5:où on convient de prendre 0 comme origine}{What the circle actually gives}
Grothendieck wants the circle action to recover the suspension and loop constructions, thereby linking the abstract action to the elementary squares of the next chapter. The four-object poset drawn here does have the homotopy type of a circle: its nerve has four vertices and four nondegenerate edges and no higher nondegenerate simplices. The difficulty is therefore not a defective model of \(S^1\). It is the distinction between an unpointed tensor and a reduced one, which he himself identifies on page 12.

In spaces,
\[
 T_{S^1}(M)=S^1\times M,\qquad
 C_{S^1}(N)=\operatorname{Map}(S^1,N).
\]
The second space is the free loop space. Choosing \(s\in S^1\) gives the evaluation map \(\operatorname{Map}(S^1,N)\to N\); its homotopy fibre over a chosen point \(n\) is \(\Omega_nN\). A basepoint in the indexing circle alone does not turn every free loop into a loop based at \(n\).

The algebraic calculation makes the excess term completely visible. With cohomological grading and a chosen orientation,
\[
 C_*(S^1;\mathbf Z)\simeq\mathbf Z[0]\oplus\mathbf Z[1].
\]
Here \(\mathbf Z[1]\) has its nonzero term in degree \(-1\). Consequently, in \(D(k\text{-Mod})\),
\[
 T_{S^1}(M)\simeq M\oplus M[1],
 \qquad C_{S^1}(M)\simeq M\oplus M[-1].
\]
The inclusion of the basepoint and the augmentation split off \(M\). Their reduced complements are \(M[1]\) and \(M[-1]\), the suspension and loop functors of this stable derivator. Taking \(M=k\ne0\), the extra degree-zero copy of \(k\) proves that the unqualified identification on the present page is false.

In a pointed homotopy theory the corresponding operation uses the smash tensor with a based circle; its cotensor is the based mapping object. The ordinary unpointed action is recovered by adjoining a disjoint basepoint to the space first. Formula (14) is therefore correct with \(\Sigma,\Omega\) understood as the unpointed circle tensor and cotensor, but it is not then the familiar reduced suspension-loop identity. With the reduced meanings, the right side is the loop space of the pointed external Hom at its zero morphism. The explicit Kan extension construction of those reduced functors appears in \citet[Definition 3.16 and Proposition 3.17]{Groth2013}; Chapter IV is already its essential diagrammatic pattern.
\end{dcomment}


Dans $\Hot_0$ on a une opération produit, en fait ça doit être le produit catégorique dans $\Hot_0$. Mais sur les modèles il est donné par
\[
\mathrm{hot}_0(X) \times \mathrm{hot}_0(Y) \ \simeq \ \mathrm{hot}_0(X \times Y)
\tag*{(8)}
\] % [sic] (8) -> (9), le numéro (8) est déjà employé plus haut sur cette page

% ---- from Der-ChIII-567_p006-010.tex
\printedpage{4}
% head: A. Grothendieck : Dérivateurs, Ch. III, Édition des Universités de Montpellier II et Paris VII
% foot: Version 2 février 2002

\texttt{[deux fois} (8)\texttt{]}. Ceci posé, on a les isomorphismes canoniques, pour $\xi,\eta\in\Ob\Hot_0$,
\[
\left\{
\begin{aligned}
\xi\times_{\D}(\eta\times_{\D}M) &\;\simeq\; (\xi\times\eta)\times_{\D}M\\
\hm_{\D}(\xi,\hm_{\D}(\eta,M)) &\;\simeq\; \hm_{\D}(\xi\times\eta,M).
\end{aligned}
\right.
\tag*{(9)}
\]

Si $\xi$, $\eta$ sont représentés par $X$, $Y$, la première formule p.ex.
\msp{page 6}
s’établit en calculant le deuxième membre comme
\[
\varphi_!(M_{X\times Y})\simeq p_!(\;\underbrace{q'_!(M_{X\times Y})}_{\substack{\simeq\ q_!(M_Y)_X\\ \text{car $q$ est $\mathbf{D}$-lisse}}}\;)\simeq\xi\times_{\D}(\eta\times_{\D}M),
\]
O.K.
\[
\begin{tikzcd}
X \ar[d,"p"'] & X\times_e Y \ar[l,"q'"'] \ar[d,"p'"] \ar[dl,dotted,"\varphi"'] \\
e & Y \ar[l,"q"]
\end{tikzcd}
\]

J’ai aussi envie de prouver :

\dgrblock{dgrb150}{Associativity and exactness from product squares}
\begin{dcomment}{exposition}{p6:J’ai aussi envie de prouver :}{Associativity and exactness from product squares}
Grothendieck next checks that the action has the algebra expected of an action. The calculation is a Fubini argument for homotopy Kan extensions, and its force is that it works before any mapping-space enrichment has been supplied. We retain his convention that coefficients on \(I\) have underlying diagrams of shape \(I^o\).

First, the product on \(\Hot_0\) mentioned just above can be justified on the stated models. Assume that \(\mathrm{Diag}\) is closed under binary products. The nerve preserves those products, and weak equivalences of spaces are preserved by taking a product with a fixed space. Hence the product functor sends a pair of nerve weak equivalences to a nerve weak equivalence and descends to the localization. The diagonal functor descends too. Before localization the diagonal is left adjoint to the binary product; their unit and counit descend as natural transformations, and still satisfy the triangle identities. The product functor on the localization is therefore right adjoint to its diagonal, which is exactly the universal property of a categorical product. This proves the sentence beginning \emph{en fait ça doit être}, under the required closure assumption, without assuming that \(\Hot_0\to\Hot\) is fully faithful. Existence of this product should not be confused with closure under the right adjoints to products, which becomes relevant to (13).

For \(X,Y\in\mathrm{Diag}\), compose the projections \(X\times Y\to X\to e\). Pseudofunctoriality of left Kan extension identifies their composite with the extension to \(e\). Base change for the product square identifies the first intermediate extension of a constant coefficient with the constant coefficient \(T_YM\) on \(X\). Thus
\[
 T_XT_YM\cong T_{X\times Y}M.
\]
The terminal category supplies \(T_e\cong\mathrm{id}\). The analogous right Kan calculation gives \(C_XC_Y\cong C_{X\times Y}\), with the harmless interchange of \(X\) and \(Y\) supplied by the product symmetry. The associativity and unit comparisons are canonical mates. Their coherence follows by pasting product squares and using the composition constraints for Kan extension; it is not a second independent choice of associators.

The same argument proves more than preservation of finite colimits or finite limits. Fix \(X\), let \(f:I\to J\), and write \(p:X\times I\to I\), \(q:X\times J\to J\), \(f'=1_X\times f\). The product square is homotopy exact in both directions. Therefore
\[
\begin{aligned}
 f_*C_{X,I}F
 &=f_*p_*p^*F\\
 &\cong q_*f'_*p^*F\\
 &\cong q_*q^*f_*F=C_{X,J}f_*F .
\end{aligned}
\]
The dual calculation gives \(f_!T_{X,I}\cong T_{X,J}f_!\). Every Kan extension permitted by \(\mathrm{Diag}\) is covered. This is the precise content of \emph{exact à gauche} and \emph{exact à droite} here.

For completeness, the decisive product comparison can be checked directly by the pointwise formula. At \((x,j)\), the comma category computing the appropriate Kan extension is a product of the comma category for \(f\) with a slice of \(X\) having an extremal object. Homotopy Kan extension over the latter factor reduces to evaluation. The resulting isomorphism is the mate of the product square, so its pointwise verification also identifies the comparison that must be used. This is the special case of base change for Grothendieck fibrations and opfibrations formalized in \citet[Proposition 1.24]{Groth2013}.

Finally,
\[
 \Hom(T_XM,N)
 \cong\Hom_{\D(X)}(p_X^*M,p_X^*N)
 \cong\Hom(M,C_XN)
\]
proves the tensor-cotensor adjunction without representability in \(\Hot\). This separation is useful: (7), (9), and Proposition 1 are consequences of the Kan formalism once descent is available; the external Hom sought next is an additional representation problem.
\end{dcomment}


\textbf{Proposition 1.} \textit{Soit $\xi\in\Ob\Hot_0$. Alors $F\longmapsto\xi\times_{\D}F$, en tant que morphisme de dérivateurs $\D\longrightarrow\D$, est} exact à droite, \textit{et $F\longmapsto\hm_{\D}(\xi,F)$ est exact à gauche.}

On peut supposer que $\xi$ est défini par $X$ dans Diag. Soit $f:I\longrightarrow J$, soit $F$ dans $\D(I)$, il faut prouver
\[
f_*(\hm_{\D}(\xi,F))\;\simeq\;\hm_{\D}(\xi,f_*(F)).
\]
Considérons donc
\[
\begin{tikzcd}
X\times I \ar[r,"f'"] \ar[d,"{p=(x,i)\longmapsto i}"'] & X\times J \ar[d,"{q=(x,j)\longmapsto j}"] \\
I \ar[r,"f"] & J.
\end{tikzcd}
\]

On a
\[
\hm_{\D}(\xi,F)=p_*(p^*F),
\]
donc
\[
f_*(\hm_{\D}(\xi,F))=\underbrace{f_*p_*}_{(fp)_*}(p^*(F))=(qf')_*p^*(F)\simeq q_*(f'_*p^*(F)).
\]
Mais
\[
f'_*p^*(F)\;\simeq\;q^*f_*(F),
\]
formule de changement de base des $f_*$, pour le changement de base $q$, qui est lisse, car déduit de $X\longrightarrow e$, qui est lisse, par changement de base $J\longrightarrow e$. On trouve
\printedpage{5}
% head: A. Grothendieck : Dérivateurs, Ch. III, Édition des Universités de Montpellier II et Paris VII
% foot: Version 2 février 2002
\msp{page 7}
donc
\[
f_*(\hm_{\D}(\xi,F))\simeq q_*q^*(f_*(F))\defeq\hm_{\D}(\xi,f_*(F)),
\]
q.e.d. L’énoncé pour $F\longmapsto\xi\times_{\D}F$ est dual de celui qu’on vient de prouver.

Reprenons, pour $\xi\in(\mathrm{hot})_0$ \texttt{[}$=\Hot_0$\texttt{]}, la formule d’adjonction
\[
\Hom(\xi\times_{\D}M,N)\simeq\Hom(M,\hm_{\D}(\xi,N))\defeq\Hom_{\xi}(M,N).
\tag*{$(*)$}
\]

Si $\xi$ est défini par $X$ dans Diag, les deux membres de $(*)$ s’identifient à
\[
\Hom_{\D(X)}(p^*M,p^*N),
\]
où $p:X\longrightarrow e$ est le foncteur structural.

Fixons $M$, $N$, et considérons l’expression $(*)$ comme contrafoncteur
\[
(\Hot_0)^o\;\longrightarrow\;\mathrm{Ens}.
\]

Il semblerait que dans tous les cas connus, ce foncteur soit ‘représentable’ par un objet, sinon de $\Hot_0$ lui-même, du moins un objet dans Hot. Si on ne peut le définir comme objet $H$ de $\Hot_0$, le fait qu’on ait un isomorphisme canonique
\[
\Hom_{\xi}(M,N)\;\simeq\Hom_{\Hot}(\xi,H),
\tag*{(10)}
\]
fonctoriel en $\xi\in\Ob\Hot_0$, ne \texttt{[le]} définit pas forcément
\msp{page 8}
à isomorphisme unique près. En fait, quand
\[
\D=\D_{\mathcal{M},\Sigma}
\]
est défini par une catégorie de modèles $(\mathcal{M},\Sigma)$, alors il y a une définition naturelle d’une opération $\homext$ (Hom ‘externe’) à valeurs dans Hot, savoir un foncteur
\[
\left\{
\begin{aligned}
&\mathcal{A}^o\times\mathcal{A} &&\longrightarrow &&\mathrm{Hot}\\
&(M,N) &&\longmapsto &&\homext(M,N),
\end{aligned}
\right.
\tag*{(11)}
\]
où $\mathcal{A}=\mathcal{M}\Sigma^{-1}$ est la catégorie fondamentale de $\D=\D_{\mathcal{M},\Sigma}$. Il est vrai que je n’ai pas vérifié, si p. ex. $(\mathcal{M},\Sigma)$ est sous-jacente à une catégorie de modèles quillénienne, de sorte que (par \textsc{Thomason}) $\D$ est bien un dérivateur, que cette notion de $\homext$ satisfait bien à une formule de la forme (10). Ici je vais le poser [?], au besoin, ces données supplémentaires : donnée de (11), et d’un isomorphisme fonctoriel
\[
\underbrace{\Hom_{\xi}(M,N)}_{\begin{subarray}{l}=\ \Hom(\xi\times_{\D}M,N)\\ =\ \Hom(M,\hm_{\D}(\xi,N))\end{subarray}}
\quad\simeq\;\Hom_{\Hot}(\xi,\homext(M,N)).
\tag*{(12)}
\]

\dgrblock{dgrb151}{External Hom as a derived mapping space}
\begin{dcomment}{later}{p8:Il est vrai que je n’ai pas vérifié}{External Hom as a derived mapping space}
Grothendieck is now trying to recover the space of maps from the already available operations on coefficients. In a model category this space had been constructed independently: the hammock localization of \citet[Corollary 4.2, Proposition 4.4, and Corollary 4.7]{DwyerKan1980Function} supplies mapping complexes, identifies their components with localized morphisms, and compares them with the enrichment of a simplicial model category. This work predates the manuscript. Chapter V will approach the same issue through categories of zigzags.

Here is the verification of (12) in the simplicial case. Choose a cofibrant representative \(QM\) and a fibrant representative \(RN\). A simplicial set \(K\) representing \(\xi\) is cofibrant in the usual model structure. The simplicial tensor-Hom adjunction gives
\[
 \operatorname{Map}_{\mathcal M}(K\otimes QM,RN)
 \cong
 \operatorname{Map}_{\mathrm{sSet}}
       (K,\operatorname{Map}_{\mathcal M}(QM,RN)).
\]
Taking components identifies the left side with
\(\Hom_{\mathrm{Ho}(\mathcal M)}(K\otimes^{\mathbf L}M,N)\).
The right side is
\([K,\operatorname{Map}_{\mathcal M}(QM,RN)]\), because the inner mapping complex is fibrant. Moreover \(K\otimes^{\mathbf L}M\) is the homotopy colimit of the constant \(M\)-diagram indexed by the space \(K\). For \(K=B X\), the simplicial replacement of the constant \(X^o\)-diagram gives exactly that tensor; its simplices are copies of \(QM\) indexed by strings of arrows in \(X^o\). Reversal identifies its realized indexing space with \(B X\). This proves the formula with the particular tensor defined in (4).

The cotensor identity follows either from the dual model construction or from the tensor-cotensor adjunction already proved. Derived function complexes and framings give the corresponding statement for a general Quillen model category; the tensor, cotensor, and enrichment compatibilities are treated systematically in \citet[Sections 10.2--10.3]{Riehl2014}. Thus the extra datum is supplied by the model presentation and is independent, up to homotopy equivalence, of the chosen resolutions.

For a presentable \(\infty\)-category \(\mathcal C\), the same calculation has the particularly transparent form
\[
 \operatorname{Map}_{\mathcal C}
 \bigl(\operatorname*{colim}_{K}M,N\bigr)
 \simeq
 \operatorname{Map}_{\mathcal S}
 \bigl(K,\operatorname{Map}_{\mathcal C}(M,N)\bigr).
\]
It is the universal property of the colimit indexed by the space \(K\). Passing to homotopy categories gives (12); no new choice of external Hom is needed once \(\mathcal C\) is given. The universal role of space-valued presheaves is made precise in \citet[Theorem 5.1.5.6]{Lurie2009HTT}.

For a completely abstract derivator, one should distinguish an action by spaces from a closed action possessing these Hom objects. \citet[Section 1, pp.~3--4, and Definition 4.2]{GrothShulman2017} explicitly distinguish the two structures. The cited results settle the model and presentable \(\infty\)-categorical cases; they do not justify erasing Grothendieck's representability hypothesis for every restricted diagram domain considered here.
\end{dcomment}


\printedpage{6}
% head: A. Grothendieck : Dérivateurs, Ch. III, Édition des Universités de Montpellier II et Paris VII
% foot: Version 2 février 2002
\msp{page 9}
Si maintenant $\eta$ est un deuxième élément de $\Hot_0$, on a envie d’écrire
\[
\left\{
\begin{aligned}
\homext(\eta\times_{\D}M,N) &\;\simeq\;\homext(M,\hm_{\D}(\eta,N))\\
&\;\simeq\;\hm_{\HOT}(\eta,\homext(M,N)),
\end{aligned}
\right.
\tag*{(13)}
\]
où dans le deuxième membre, $\hm_{\HOT}$ est l’opération (1) dans le cas où $\D=\mathrm{HOT}$, et qui est défini aussi en termes du produit cartésien $\xi\times\eta$ dans Hot, par
\[
\Hom_{\Hot}(\xi,\hm_{\HOT}(\eta,\zeta))\;\simeq\;\Hom_{\Hot}(\xi\times\eta,\zeta),
\]
en d’autres termes, pour $\eta$ fixé, le foncteur
\[
\zeta\;\longmapsto\;\hm_{\HOT}(\eta,\zeta)
\]
est défini comme le foncteur adjoint à droite du foncteur
\[
\xi\;\longmapsto\;\xi\times\eta.
\]

Pour vérifier heuristiquement la formule (13) (vérification valable ‘en forme’ si les trois objets envisagés sont en fait définis comme objets de $\Hot_0\ \ldots$), on prend $\Hom_{\Hot}(\xi,-)$ à valeurs dans les trois membres, pour $\xi$ donné dans $\Hot_0$. Pour le premier, on trouve
\[
\begin{aligned}
\Hom_{\xi}(\eta\times_{\D}M,N) &\;=\;\Hom(\xi\times_{\D}(\eta\times_{\D}M),N)\\
&\;=\;\Hom((\xi\times\eta)\times_{\D}M,N).
\end{aligned}
\]
Pour le deuxième
\[
\begin{aligned}
\Hom_{\xi}(M,\hm_{\D}(\eta,N)) &\;=\;\Hom_{\mathcal{A}}(\xi\times_{\D}M,\hm_{\D}(\eta,N))\\
&\;\simeq\;\Hom_{\mathcal{A}}(\eta\times_{\D}(\xi\times_{\D}M),N)\\
&\;\simeq\;\Hom_{\mathcal{A}}((\eta\times\xi)\times_{\D}M,N),
\end{aligned}
\]
\msp{page 10}
et enfin pour le troisième
\[
\begin{aligned}
\Hom_{\Hot}(\xi\times\eta,\homext(M,N)) &\;\simeq\;\Hom_{\xi\times\eta}(M,N)\\
&\;\simeq\;\Hom((\xi\times\eta)\times_{\D}M,N),
\end{aligned}
\]
donc toujours la même chose ! Pas de doute, ça doit marcher.

S’il faut se donner le $\homext$ (11) comme donnée supplémentaire, avec de plus un isomorphisme trifonctoriel (12), il faut aussi se donner les isomorphismes trifonctoriels (13), \textit{‘compatibles’} avec l’isomorphisme (12) au sens du calcul précédent.

\dgrblock{dgrb152}{What the Yoneda calculation proves}
\begin{dcomment}{exposition}{p10:S’il faut se donner le}{What the Yoneda calculation proves}
Grothendieck's three calculations are intended to identify the external Hom of a tensor with the cotensor of an external Hom. The calculation is correct. His qualification about the test category is also correct, and is part of the argument rather than an unnecessary hesitation.

Suppose the external Hom represents (12), and suppose that all three candidate objects in (13) belong to \(\Hot_0\). For every \(\xi\in\Hot_0\), tensor associativity and the tensor-cotensor adjunction give natural identifications of their represented functors with
\[
 \xi\longmapsto
 \Hom_{\mathcal A}((\xi\times\eta)\times_{\D}M,N).
\]
Yoneda in \(\Hot_0\) then gives a unique isomorphism respecting these identifications. Naturality in \(\eta,M,N\) follows from the same uniqueness. The last candidate additionally requires the relevant function object to exist in \(\Hot_0\); closure under products alone does not imply that closure under function objects.

If the candidates lie only in \(\Hot\), the calculation instead identifies their \emph{restricted} represented functors. Such an identification need not determine the objects. For a concrete example, take \(\Hot_0\) to be the full subcategory of discrete spaces. For any discrete set \(S\) and any nonempty connected space \(H\), the set \([S,H]\) has one element: a map is a choice of a component for each element of \(S\), and there is only one component. Consequently the point and \(S^1\) define isomorphic functors on this test category, although they are not isomorphic homotopy types. This example isolates the exact logical failure that would result from applying unrestricted Yoneda here.

There are two legitimate completions. One may assume the restricted Yoneda functor is fully faithful on the particular objects under consideration, and then the displayed calculation suffices. Alternatively, a model or \(\infty\)-categorical construction gives the actual comparison of mapping spaces before passage to homotopy types, proving (13) directly. The manuscript follows the first course conditionally and otherwise proposes to retain the comparisons as part of the additional structure. No gap should be charged to a proof that explicitly states this restriction.
\end{dcomment}


\printedpage{7}
% head: A. Grothendieck : Dérivateurs, Ch. III, Édition des Universités de Montpellier II et Paris VII
% foot: Version 2 février 2002

Par exemple, prenant $\eta=$ type d’homotopie du cercle $S^1$, les formules (13) se spécialisent en
\[
\begin{aligned}
\homext(\Sigma M,N) &\;\simeq\;\homext(M,\Omega N)\\
&\;\simeq\;\Omega\homext(M,N)\\
&\;\defeq\;\hm_{\HOT}(S^1,\homext(M,N)).
\end{aligned}
\tag*{(14)}
\]

Mais cette formule dans Hot n’est sans doute pas tout à fait la bonne, car le bon formalisme de $\Omega$, $\Sigma$ se fait dans les catégories \textit{ponctuées,} et en particulier dans la catégorie homotopique ponctué $\Hot_{\bullet}$. % [sic] ponctué -> ponctuée, agreement with « catégorie »

Je vais expliciter un exemple, en prenant le dérivateur
\[
\mathrm{Dér}(A)
\]
‘dérivé’ d’une catégorie abélienne $A$. Donc $\D(e)=$ catégorie dérivée de $A$. Pour éviter des
\msp{page 11}
difficultés accessoires ici, je vais supposer que $A$ est la catégorie des $k$-modules, $k$ un anneau quelconque.

À tout $\xi\in\Ob\Hot$ est associé un ‘abélianisé’
\[
\xi_{\ab}\defeq H_{\bullet}(\xi)\in\underbrace{D(\mathrm{Ab})}_{\text{ou }\mathbf{D}(\mathbf{Z}\text{-Mod})}\;,\qquad\quad\text{catégorie dérivée des groupes abéliens,}
\]
objet qui n’a de cohomologie qu’en degré $\le 0$. Ceci dit, si $M$ est dans $D(k\text{-Mod})$, on aura
\[
\left\{
\begin{aligned}
\xi\times_{\D}M &\;\simeq\;\xi_{\ab}\Ltens_{\mathbf{Z}}M\\
\hm_{\D}(\xi,M) &\;\simeq\;\RHom_{\mathbf{Z}}(\xi_{\ab},M),
\end{aligned}
\right.
\tag*{(15)}
\]
formule qui est en cheville avec la formule de Künneth pour l’homologie
\[
(\xi\times\eta)_{\ab}\simeq\xi_{\ab}\Ltens_{\mathbf{Z}}\eta_{\ab},
\tag*{(16)}
\]
ce qui redonne un aspect familier aux formules (9), qui se réduisent à
\[
\left\{
\begin{aligned}
\xi_{\ab}\Ltens_{\mathbf{Z}}(\eta_{\ab}\Ltens_{\mathbf{Z}}M) &\;\simeq\;(\xi_{\ab}\Ltens_{\mathbf{Z}}\eta_{\ab})\Ltens_{\mathbf{Z}}M\\
\RHom_{\mathbf{Z}}(\xi_{\ab},\RHom_{\mathbf{Z}}(\eta_{\ab},M)) &\;\simeq\;\RHom_{\mathbf{Z}}(\xi_{\ab}\Ltens_{\mathbf{Z}}\eta_{\ab},M).\qquad\text{(\footnotemark[2])}
\end{aligned}
\right.
\]
\footnotetext[2]{Expliciter aussi la formule d’adjonction (7) (p. 4).}

D’autre part, existe-t-il bien un $\homext(M,N)$, si $M$, $N$ dans $D(k\text{-Mod})$ ? On pense à $\RHom_k(M,N)$, mais pour en faire un type d’homotopie, on doit ‘tronquer en degrés $>0$’. On pourrait donc prendre
\[
\homext(M,N)\;\simeq\;\tau_{\le 0}(\RHom_k(M,N)).
\]

\dgrblock{dgrb153}{Integral chains and the truncation of derived Hom}
\begin{dcomment}{examples}{p11:D’autre part, existe-t-il bien}{Integral chains and the truncation of derived Hom}
Grothendieck turns to complexes because they make the external Hom computable. We can verify every part of the proposed formula, provided that \emph{abélianisation} means the derived free abelian-group construction and that the truncation is the canonical one.

For a simplicial set \(K\), let
\[
 A_K=N\mathbf Z[K]
\]
be its normalized free chain complex, written cohomologically in degrees at most zero. This is the object denoted \(H_\bullet(\xi)\) in the manuscript. It is a chain complex up to quasi-isomorphism, not merely its graded homology groups. The free simplicial abelian-group functor is left adjoint to the forgetful functor. After deriving and applying the Dold--Kan equivalence, it gives
\[
 [K,U\,\mathrm{DK}(B)]
 \cong
 \Hom_{D(\mathbf Z)}(A_K,B),
 \qquad B\in D(\mathbf Z)^{\leq0}.
\]
Here \(\mathrm{DK}\) means denormalization with the change from cohomological to homological grading, and \(U\) forgets the abelian-group structure. The underlying equivalence is \citet[Theorem 14.24.3]{StacksDoldKan}. It is this derived adjunction that answers the question in footnote 3; being only an abelian-group object in the ordinary homotopy category would not by itself specify the requisite strict or coherent abelian structure.

Put \(R=\mathbf R\Hom_k(M,N)\), viewed as an object of \(D(\mathbf Z)\). For a noncommutative \(k\), this target is still correct: there is generally no left \(k\)-module structure on this Hom complex. Since \(A_K\in D(\mathbf Z)^{\leq0}\), the truncation adjunction gives
\[
\begin{aligned}
 {}[K,U\,\mathrm{DK}(\tau_{\leq0}R)]
 &\cong\Hom_{D(\mathbf Z)}(A_K,\tau_{\leq0}R)\\
 &\cong\Hom_{D(\mathbf Z)}(A_K,R)\\
 &\cong\Hom_{D(k)}(A_K\otimes_{\mathbf Z}^{\mathbf L}M,N).
\end{aligned}
\]
The final equality is the derived tensor-Hom adjunction; \(k\) acts on the factor \(M\). This proves (10) and (12) for the claimed external Hom, with the otherwise suppressed Dold--Kan and forgetful functors restored.

The truncation is the smart truncation: its degree-zero term is \(\ker(d:R^0\to R^1)\), and its negative terms are those of \(R\). Simply deleting the positive degrees would replace degree-zero cocycles by all of \(R^0\), generally giving a wrong set of components. The homotopy groups of the resulting simplicial abelian group are
\[
 \pi_n\homext(M,N)
 \cong H^{-n}R
 \cong\Hom_{D(k)}(M,N[-n]),\qquad n\geq0.
\]
For instance, with \(k=\mathbf Z\), \(M=\mathbf Z\), and \(N=\mathbf Z[r]\), \(r\geq0\), one obtains \(K(\mathbf Z,r)\). The adjunction then reads
\([K,K(\mathbf Z,r)]\cong H^r(K;\mathbf Z)\), including degree zero with its usual unreduced interpretation.

The other formula in (15) follows from the same adjunction:
\[
 T_KM=A_K\otimes_{\mathbf Z}^{\mathbf L}M,\qquad
 C_KM=\mathbf R\Hom_{\mathbf Z}(A_K,M).
\]
The Eilenberg--Zilber comparison identifies \(A_{K\times L}\) with \(A_K\otimes_{\mathbf Z}^{\mathbf L}A_L\) \citep[Section 1, theorem, pp.~200--201]{EilenbergZilber1953}; normalized free chains are bounded above complexes of free abelian groups in this grading, so the ordinary tensor already computes the derived tensor. Formula (16) is this chain-level equivalence. A formula involving only the tensor product of homology groups would lose the Tor terms and would not suffice for Grothendieck's argument.
\end{dcomment}


\printedpage{8}
% head: A. Grothendieck : Dérivateurs, Ch. III, Édition des Universités de Montpellier II et Paris VII
% foot: Version 2 février 2002
\msp{page 12}
La formule fondamentale (10) est-elle alors valable ? Le deuxième membre est en effet
\[
\Hom_{\Hot}(\xi,H)\;\simeq\;\Hom_{\mathrm{Hotab}}(\xi_{\ab},H)\qquad\quad\text{(\footnotemark[3])}
\]
\footnotetext[3]{? ? est-ce vrai ? ? ?}
car $H$ est un objet abélien de Hot (cf. le tapis sur l’abélianisation \ldots), à partir de là c’est gagné :
\[
\begin{aligned}
\Hom_{\underbrace{\scriptstyle\mathrm{Hotab}}_{\begin{subarray}{l}\approx\ D(\mathbf{Z}\text{-Mod})^{\le 0},\\ \text{sous-catégorie \textit{pleine}}\\ \text{de }D(\mathbf{Z}\text{-Mod})\end{subarray}}}(\xi_{\ab},\tau_{\le 0}\RHom_k(M,N)) &\;\simeq\;\Hom_{D(\mathbf{Z}\text{-Mod})}(\xi_{\ab},\RHom_k(M,N))\\
&\;\simeq\;\Hom_{D(k\text{-Mod})}(\xi_{\ab}\times M,N),
\end{aligned}
\]
O.K.

D’autre part, le bifoncteur
\[
(M,N)\;\longmapsto\;\homext(M,N)\;=\;\tau_{\le 0}\RHom_k(M,N)
\]
se factorise canoniquement par Hotab et a fortiori par $\Hot_{\bullet}$. Le ‘bon’ $\homext(M,N)$ serait ici plutôt celui dans Hotab, voire même dans la catégorie $D^{-}(\mathbf{Z}\text{-Mod})$, déduite de celle là ‘en rendant le foncteur translation inversible’. Il est évident d’ailleurs qu’ici le $\xi_{\ab}\times_{\D}M$, où $\xi=S^1$, n’est pas ‘la bonne’ suspension, qui doit être le décalage des degrés – $\xi_{\ab}$ ne devrait avoir de la cohomologie qu’en degré $-1$, et non aussi en degré $0$.

\dgrblock{dgrb154}{Keeping all degrees of the mapping complex}
\begin{dcomment}{status}{p12:D’autre part, le bifoncteur}{Keeping all degrees of the mapping complex}
Grothendieck's last suggestion seeks an external Hom that remembers the linear structure and the information restored by suspension. That is the correct direction, but the proposed bounded-above target requires a qualification when the coefficient category is the unbounded derived category.

The homotopy category of simplicial abelian groups is equivalent to \(D(\mathbf Z)^{\leq0}\). Allowing finite shifts of its objects produces \(D^-(\mathbf Z)\): an object in the union has cohomology bounded above, and conversely any bounded-above complex can be shifted into nonpositive degrees. This observation correctly describes the stated formal inversion of translation at the level of objects and morphisms. It does not imply that every derived Hom complex belongs to this union.

A two-term algebra already shows the difficulty. Let \(F\) be a field, let \(k=F[\epsilon]/(\epsilon^2)\), and let \(L=k/(\epsilon)\). A projective resolution of \(L\) is
\[
 \cdots\xrightarrow{\epsilon}k
 \xrightarrow{\epsilon}k
 \xrightarrow{\epsilon}k\longrightarrow L\longrightarrow0.
\]
It is exact because the kernel and image of multiplication by \(\epsilon\) are both the ideal \((\epsilon)\). Applying \(\Hom_k(-,L)\) gives one copy of \(F\) in every nonnegative cohomological degree and zero differentials. Thus
\[
 H^n\mathbf R\Hom_k(L,L)=\operatorname{Ext}^n_k(L,L)
 \cong F\qquad(n\geq0).
\]
The full derived Hom is not in \(D^-(\mathbf Z)\). The present external mapping space is nevertheless well defined: its nonpositive truncation is the discrete group \(F\). The omitted extension groups reappear as
\(\pi_0\homext(L,L[n])\), for every \(n\geq0\). This example explains why a single mapping space and the whole stable mapping object are different levels of information.

For complexes, the full enhancement is by the unbounded complex \(\mathbf R\Hom_k(M,N)\); in a general stable homotopy theory it is by a mapping spectrum. The classical differential graded formulation and its relation to triangulated categories are discussed by \citet[Sections 2--3]{Keller2006DG}. The modern theory of monoidal and enriched derivators distinguishes the unpointed, pointed, and stable actions and the existence of their right adjoints \citep[Sections 3 and 8--9]{GrothPontoShulman2014Additivity}.

The correction here concerns the generality of a tentative proposal, not a theorem whose proof has collapsed. The formulas involving \(\tau_{\leq0}\mathbf R\Hom\) are valid. The earlier circle identification is repaired by reduced chains. An additional boundedness hypothesis, such as suitable finite projective dimension and boundedness on the inputs, is needed only if one insists that the entire linear external Hom take values in \(D^-(\mathbf Z)\).
\end{dcomment}
\dgrchapter{IV}{Chapter IV}{Diagrammes substantiels}
\dgrpart{0}{III-V}{Manuscript pages 73--83}{III-V\enspace\textperiodcentered\enspace Chapter IV\enspace\textperiodcentered\enspace manuscript pages 73--83}{Chapter IV \textperiodcentered\ \textup{III-V}}
% ---- from Der-ChIV-567_p001-005.tex
\printedpage{}% no page number printed on this page
% head: (none)

\begin{center}
{\LARGE\textsf{LES DÉRIVATEURS}}

\bigskip

\textbf{ALEXANDRE GROTHENDIECK}

\bigskip\bigskip

\textbf{Chapitre IV}

\medskip

\textbf{Diagrammes substantiels}
\end{center}

\clearpage

\printedpage{}% no page number printed on this page
% head: (none)

\noindent Ce texte a été déchiffré et transcrit en \LaTeXe\ par M. Künzer. Il a été édité par M. Künzer, J. Malgoire, et G. Maltsiniotis. La transcription est aussi fidèle que possible au manuscrit. Pour les quelques corrections évidentes, ou rares commentaires des éditeurs, ainsi que pour la numérotation originale des pages du manuscrit, les caractères de machine à écrire \texttt{[typewriter]} entre crochets sont utilisés. Un point d'interrogation entre crochets signifie que l'on n'est pas sûr du mot qui précède.

\noindent Cette édition est provisoire. Les remarques, commentaires et corrections sont bienvenus. Envoyer un message à :

\noindent\textbf{maltsin@math.jussieu.fr}

\hfill G. Maltsiniotis\hspace*{2cm}

\clearpage

\printedpage{1}
% head: A. Grothendieck : Dérivateurs, Ch. IV, Édition des Universités de Montpellier II et Paris VII
% foot: Version 2 février 2002

\noindent\texttt{[Il n'y a pas de pages 1-72]}

\msp{page 73}

\subsection*{Construction de diagrammes substantiels pour un dérivateur.}

\noindent Soit $\D$ un prédérivateur, $I$ dans Diag, considérons

\noindent $\varphi_I : \mathcal{D}(I) \longrightarrow \uHom(I^o,\cA)$,\ \ où $\cA=\D(e)$ est la catégorie fondamentale des coefficients. % [sic] \mathcal{D}(I) -> \mathbf{D}(I), elsewhere the prederivator is the bold D

Les objets de $\uHom(I^o,\cA)$ s'interprètent comme des \textit{diagrammes de type $I^o$} dans $\cA$. Ce langage est surtout commode si $I$ est ordonnée. Notons, pour $F\in\mathrm{Ob}\,\D(I)$, par $(F_i)_{i\in\mathrm{Ob}\,I}$ le $I^o$-diagramme associé dans $\cA$. Par l'axiome Dér 2, le foncteur
\[
\varphi : F \longmapsto (F_i)_{i\in\mathrm{Ob}\,I}
\]
est conservatif. Mais il n'a aucune envie d'être fidèle, encore moins pleinement fidèle. Quand on connaît le diagramme $(F_i)_i$, on est loin de connaître $F$, pas même à isomorphisme non unique près, je présume. Il n'est nullement clair que si $F$, $F'$ dans $\D(I)$ sont tels que les diagrammes associés soient isomorphes, ils soient eux-mêmes isomorphes.

\dgrblock{dgrb155}{Why the skeleton is insufficient}
\begin{dcomment}{context}{p73:Il n'est nullement clair que}{Why the skeleton is insufficient}
Grothendieck's term \emph{diagramme substantiel} names the information on which the whole project depends: a diagram before its coherent relations have been discarded. The point is not that an ordinary commutative diagram has been drawn inaccurately. It is that several different homotopies and compatibilities can produce the same drawing in the homotopy category.

In the model example,
\[
 \D(I)=\mathcal M^{I^o}[(W^{I^o})^{-1}],
 \qquad
 \varphi_I:\D(I)\longrightarrow(\mathrm{Ho}\mathcal M)^{I^o}.
\]
Localization is performed on the category of whole diagrams. The target instead localizes the values first and subsequently records arrows and their equations. These operations need not commute. Axiom Der~2 says that a morphism of coherent diagrams which is an isomorphism at every vertex is invertible. It says neither that every skeletal morphism lifts nor that two lifts agree. The distinction was illustrated in the commentary on Chapter~I, manuscript p.~11 (\emph{Conservative is not faithful}); there is no need to repeat that counterexample here.

For the walking arrow one can make a useful positive statement. A \emph{strong} derivator satisfies fullness and essential surjectivity of
\[
 \D(I\times[1])\longrightarrow\D(I)^{[1]^o}
\]
for every \(I\). This is the usual modern Der~5, not an identification of the numbered axioms appearing earlier in the manuscript \citep[Definition 1.8]{Groth2013}. Fullness together with conservativity implies that two coherent arrows whose skeletal arrows are isomorphic are themselves isomorphic: lift the skeletal isomorphism by fullness, then use its invertibility at both endpoints. It does not imply a unique lift of that isomorphism. Thus Grothendieck's question about objects with isomorphic skeletons has a different answer for a single arrow in a strong derivator than for an unrestricted diagram shape.

The historical motivation is particularly close to the present construction. In \emph{Pursuing Stacks}, \citet[Section 69, pp.~1--2 of the edited extract]{Grothendieck1983} recalls the need, arising around SGA~5, to supplement a derived category by derived categories of triangles and filtered complexes. Here he seeks a procedure that constructs the supplements and their maps from Kan extension itself. Heller's parallel treatment packages homotopy theories as systems of diagram categories \citep[Introduction and Chapter I]{Heller1988}. Neither perspective asks one to manufacture a functorial cone on the bare homotopy category.

The two examples pursued below are deliberately chosen to retain a universal property throughout the construction. The loop square is determined by a zero extension followed by a right Kan extension. The unit and counit arrows are determined by a category that encodes the functor \(f\). This is why canonical coherent versions of these particular arrows can be obtained even though a general skeletal arrow has no specified canonical lift.
\end{dcomment}


Pour insister sur ce point, les objets de $\D(I)$ s'appelleront les `\textit{diagrammes substantiels} de type $I^o$ dans $\cA$' (ou $\D$-substantiel, s'il y a risque de confusion). Pour étayer l'intuition et simplifier les explications, un objet $F$ de $\D(I)$ ($\D$-coefficients sur $I$) sera représenté graphiquement par le diagramme `faible' ou `squelettique' correspondant -- on le regarde
\msp{page 74}
comme une sorte de (pâle) `squelette' du diagramme substantiel dont il provient. P. ex. si $I=\bDelta_0 : 0\longrightarrow 1$ \texttt{[plutôt $\bDelta_1$]}, on note un diagramme substantiel de type $(\bDelta_0)^o$ \texttt{[plutôt $(\bDelta_1)^o$]} par le même graphisme, que le diagramme ordinaire :
\[
F_0 \longleftarrow F_1.
\]
Si
\[
I = \begin{tikzcd}
a \arrow[r] & c \\
0 \arrow[u] \arrow[r] & b, \arrow[u]
\end{tikzcd}
\]
les diagrammes substantiels de type $I$ se notent
\[
\begin{tikzcd}
F_a \arrow[d] & F_c \arrow[l] \arrow[d] \\
F_0 & F_b. \arrow[l]
\end{tikzcd}
\]

\clearpage

\printedpage{2}
% head: A. Grothendieck : Dérivateurs, Ch. IV, Édition des Universités de Montpellier II et Paris VII
% foot: Version 2 février 2002

Il faudra bien sûr préciser chaque fois si ce graphisme désigne un diagramme squelettique, ou substantiel. Souvent on part d'un squelette, et on essaye de `l'étoffer' avec un peu de chair, pour en faire un diagramme substantiel -- si possible, de façon canonique. C'est là un point de philosophie dérivatique essentiel pour la pratique. On procédera souvent de proche en proche, en `substantifiant' des sous-diagrammes correspondants à une suite strictement croissante $I_0\subset I_1\subset\cdots\subset I_n=I$ de $I$, au mieux de la situation particulière.

\msp{page 75}

En voici un exemple typique, en supposant $\D$ un dérivateur \textit{ponctué} (les $\D(I)$ sont ponctuées, et les $f^*$ transforment objets nuls en objets nuls). Il s'agit de la définition du foncteur `cosuspension' ou `espace de lacets' dans $\cA$,
\[
\Omega : \cA \longrightarrow \cA.
\]
Il est entendu que pour $M$ dans $\cA$, $\Omega(M)$ doit s'insérer dans un carré commutatif
\[
\begin{tikzcd}
e \arrow[d] & \Omega M \arrow[l] \arrow[d] \\
M & e, \arrow[l]
\end{tikzcd}
\tag*{(C)}
\]
où $e$ désigne les éléments unités. Ce carré cependant devrait être substantiel, plus précisément, être un carré substantiel `$\D$-\textit{cartésien}'. Étant entendu qu'un objet de $\D(I)$ (carré substantiel dans $\cA$) est dit $\D$-cartésien, s'il est dans l'image essentielle de $\D(I_1)$ par $i_{1*}$, où
\[
i_1 : I_1 \hookrightarrow I,
\]
$I_1$ étant la sous-catégorie ouverte de $I$ définie par
\[
I_1 = \begin{tikzcd}
a \\
0 \arrow[u] \arrow[r] & b.
\end{tikzcd}
\]

\dgrblock{dgrb156}{Constructing the loop square without choosing homotopies}
\begin{dcomment}{exposition}{p75:En voici un exemple typique}{Constructing the loop square without choosing homotopies}
Grothendieck wants to obtain a functor \(\Omega\), not merely an object which happens to have the expected homotopy type. The two successive Kan extensions supply that functoriality. The direction of the arrows is governed by the presheaf convention: the indexing square has initial vertex \(0\) and final vertex \(c\), while its coefficient diagram has arrows pointing toward \(0\).

Let \(I_0=\{0\}\), let \(I_1=\{0,a,b\}\) with \(0\to a,0\to b\), and let \(I\) be the completed square. Denote the inclusions by \(i_0:I_0\to I_1\) and \(i_1:I_1\to I\), as in the text. Since \(I_0\) is a sieve in \(I_1\), the \emph{left} Kan extension in the presheaf convention is extension by the initial object outside \(I_0\). Pointedness identifies this initial object with zero. Hence
\[
 (i_0)_!M|_0=M,\qquad
 (i_0)_!M|_a=(i_0)_!M|_b=0.
\]
The pointwise assertion follows from empty comma categories at \(a,b\), and the assertion at \(0\) follows from full faithfulness of the inclusion. Thus the cospan \(0\to M\leftarrow0\) has been made coherent without choosing representatives or nullhomotopies.

Apply the right Kan extension \((i_1)_*\). Its value at \(c\) is the homotopy limit of that cospan. Set
\[
 \Omega M=c^*(i_1)_*(i_0)_!M.
\]
The object \((i_1)_*(i_0)_!M\) carries, at once, the two maps from \(\Omega M\) to zero and the specified coherent commutativity with the two zero maps to \(M\). It is cartesian by construction. Every morphism \(M\to N\) induces a morphism of these squares by the functoriality of the composite, so the construction really produces an endofunctor.

The stronger equivalence stated on page 77 follows with little extra work. Let \(\mathcal Q\subset\D(I)\) be the full subcategory of cartesian squares with values zero at \(a,b\). The functor
\[
 L=(i_1)_*(i_0)_!:\D(e)\longrightarrow\mathcal Q
\]
is fully faithful, since both Kan extensions are. It is essentially surjective as well. If \(Q\in\mathcal Q\), cartesianness gives \(Q\cong(i_1)_*i_1^*Q\). The counit
\((i_0)_!i_0^*i_1^*Q\to i_1^*Q\) is an isomorphism at \(0\), and at \(a,b\) both sides are zero. Der~2 therefore makes it an isomorphism. Substitution gives \(Q\cong L(Q_0)\), and restriction to \(0\) is a quasi-inverse.

The phrase \emph{unique à isomorphisme unique près} has its usual universal-property meaning: after specifying the identification of the input with \(M\), a comparison is unique when required to induce that identification. A square may have automorphisms induced by automorphisms of \(M\). Also, one should not add a requirement that the induced comparison fix an independently chosen identification of the fourth vertex. The theorem determines that vertex as part of the construction.

For complexes, the homotopy fibre of \(0\to M\) is \(M[-1]\), so this recovers the reduced loop functor, with no unwanted summand \(M\). This is the construction needed to repair the circle calculation of Chapter III.
\end{dcomment}


\msp{page 76}

La restriction du carré substantiel cherché (C) à la sous-catégorie $I_1$ est donc un diagramme substantiel de squelette
\[
\begin{tikzcd}
e \arrow[d] \\
M & e, \arrow[l]
\end{tikzcd}
\tag*{(Sq)}
\]
\printedpage{3}
% head: A. Grothendieck : Dérivateurs, Ch. IV, Édition des Universités de Montpellier II et Paris VII
% foot: Version 2 février 2002
dont il se déduit (à isomorphisme canonique près) en appliquant $i_*$ à ce dernier, de façon à le compléter en $\begin{tikzcd}
e \arrow[d] & \Omega M \arrow[l, dotted] \arrow[d, dotted] \\
M & e \arrow[l]
\end{tikzcd}$. Il reste donc à décrire une substantiation canonique du squelette (Sq), en termes de $M$. Cela nous conduit à regarder la sous-catégorie ouverte
\[
I_0 = \{0\} :\quad
\begin{tikzcd}
{} \arrow[r, dotted] & {} \\
0 \arrow[u, dotted] \arrow[r, dotted] & {} \arrow[u, dotted]
\end{tikzcd}
\]
de $I_1$ (donc aussi de $I$) réduite à l'objet initial, et l'inclusion
\[
I_0 \overset{i_0}{\hookrightarrow} I_1,
\]
et de substantifier le squelette (Sq) par $(i_0)_!(M)$.

On peut dire a posteriori que, si on connaissait le carré substantiel, il n'y a, à isomorphisme unique près, qu'une seule façon de l'obtenir comme un $i_{1*}i_{0!}(?)$, avec ? dans $\D(I_0)\simeq\cA$, car $i_{0!}$ et $i_{1*}$ étant pleinement fidèles, il en est de même de leur composé.

\msp{page 77}

En fait, on voit que le foncteur
\[
i_{1*}i_{0!} :\ \cA\simeq\D(I_0) \longrightarrow \D(I)
\]
induit une équivalence de $\cA$ avec la sous-catégorie pleine de $\D(I)$ formée des carrés $\D$-cartésiens dont les fibres en $a$ et $b$ sont nulles. (Le foncteur quasi-inverse évident est $i^*$, où $i : U_0 \hookrightarrow I$ \texttt{[plutôt $I_0 \hookrightarrow I$]}.) Cela signifie donc que dans ce cas, la donnée du squelette de nature très particulière ($F_a\simeq F_b\simeq 0_{\cA}$), plus la contrainte supplémentaire $\D$-cartésienne, sur le carré cherché, détermine celui-ci à isomorphisme unique près en termes de son squelette : il y a une substantiation non seulement `canonique' de celui-ci, mais même unique à isomorphisme unique près, à cause de la contrainte $\D$-cartésienne.

Je voudrais commencer par regarder le cas plus simple encore où $I=\bDelta_1=(0\longrightarrow 1)$, donc les squelettes envisagés sont simplement les flèches
\[
F_1 \longrightarrow F_0 \qquad \text{dans } \cA.
\]
La philosophie nous oblige, pour les flèches les plus fondamentales dans la théorie des dérivateurs $\D$ (flèches dans $\cA$, ou par extension dans une $\D(I)\simeq\cA_{\D_I}$, catégorie fondamentale du dérivateur induit $\D_I$) de se demander s'il n'y a pas

% ---- from Der-ChIV-567_p006-009.tex
\printedpage{4}
% head: A. Grothendieck : Dérivateurs, Ch. IV, Édition des Universités de Montpellier II et Paris VII
% foot: Version 2 février 2002
\msp{page 78}
une façon ‘canonique’, ‘naturelle’, ‘voulue par Dieu’, de le ‘remplumer’ en un diagramme substantiel. Les flèches les plus évidentes qui se présentent sont relatives à un foncteur
\[ f : I \longrightarrow J \]
dans Diag, et aux deux morphismes d'adjonction pour le couple $(f^*, f_*)$, savoir
\[
\begin{aligned}
G &\longrightarrow f_*f^*G && \text{dans } D(J)\\
f^*f_*F &\longrightarrow F && \text{dans } D(I). \quad (\text{\footnotemark})
\end{aligned}
\]
\footnotetext{Bien sûr, problème dual pour $(f_!, f^*)$.}

Je vais essayer de remplumer ces deux flèches, et vais généraliser en envisageant la situation ci-contre \texttt{[ci-dessous]}
\[
\begin{tikzcd}
I \arrow[d, "f"'] & F \\
J & G, \arrow[u, "u"']
\end{tikzcd}
\]
où
\[ \Hom^f(G,F) \overset{\text{déf}}{=} \Hom(f^*G,F) \overset{\text{adj.}}{\simeq} \Hom(G,f_*F). \]
Donc le même $u$ s'interprète par \textit{deux} flèches squelettiques
\[
\left\{
\begin{aligned}
f^*(G) &\longrightarrow F && \text{dans } \mathbf{D}(I)\\
G &\longrightarrow f_*(F) && \text{dans } \mathbf{D}(J).
\end{aligned}
\right.
\tag*{(1)}
\]
Les remplumer, c'est les interpréter (si faire se peut) en termes de deux objets (si possible canoniques) dans $\mathbf{D}(I \times \bDelta_1)$ et $\mathbf{D}(J \times \bDelta_1)$ respectivement.

Je vais me guider sur le cas particulier d'un dérivateur trivial
\[ \mathbf{D} = \mathbf{D}_{\mathcal{M}}, \]
où $\mathcal{M}$ est une catégorie de modèles ‘nue’ $(\Sigma = \emptyset)$.

\msp{page 79}
Pour $f$ fixé, on va construire une catégorie auxiliaire $K(f)$, de telle façon que pour toute catégorie $\mathcal{M}$, les systèmes $(G,F,u)$ d'un $\mathcal{M}$-préfaisceau $G$ sur $J$, d'un autre $F$ sur $I$, et d'un $f$-homomorphisme $G \longrightarrow F$, s'identifient (fonctoriellement) aux $\mathcal{M}$-préfaisceaux sur
\printedpage{5}
% head: A. Grothendieck : Dérivateurs, Ch. IV, Édition des Universités de Montpellier II et Paris VII
% foot: Version 2 février 2002
$K(f)$. Ceci dit, on va essayer d'interpréter les flèches (1) squelettiques, mais peu importe ici -- tous les $\mathbf{D}$-squelettes sont maintenant remplumés) en termes de la géométrie de $K(f)$. % [sic] remplumés) -> remplumés, closing parenthesis printed without any opening parenthesis

\[
\left\{
\begin{aligned}
&\Ob K(f) = \Ob K(I) \amalg \Ob K(J) \qquad \text{\texttt{[plutôt} $\Ob I \amalg \Ob J$\texttt{ ]}}\\
&\text{$I$, $J$ s'identifiant à des sous-catégories pleines de $K(f)$,}\\
&\text{il n'y a pas des flèches d'un $j \in J$ dans un $i \in I$.}\\
&\text{(Donc $I$ est ouverte dans $K(f)$, $J$ son complémentaire fermé.)}\\
&\text{Pour $i \in \Ob I$, $j \in \Ob J$, on pose $\Hom_{K(f)}(i,j) = \Hom_J(f(i),j)$.}
\end{aligned}
\right.
\]

La composition des flèches dans $K(f)$ est évidente. On va définir une rétraction canonique
\[ \bar{f} : K(f) \longrightarrow J \]
de $K(f)$ sur sa sous-catégorie fermée $J$. Donc $\bar{f}$ est connu sur $\Ob J$, et sur $\Ob I$ ce sera $f$ :
\[
\left\{
\begin{aligned}
\bar{f}(j) &= j && \text{si } j \in \Ob J\\
\bar{f}(i) &= f(i) && \text{si } i \in \Ob I.
\end{aligned}
\right.
\]
Il reste à définir, pour $x, y \in \Ob K(f)$,
\msp{page 80}
\[ \Hom_{K(f)}(x,y) \longrightarrow \Hom_J(\bar{f}(x), \bar{f}(y)). \]
Si $x, y \in \Ob J$, c'est l'identité, si $x, y \in I$, c'est $f$. Si $x \in I$, $y \in J$, on doit définir une flèche
\[ \Hom(f(x),y) \longrightarrow \Hom(f(x),y), \]
et on prend encore l'identité. On a
\[
\begin{tikzcd}
J \arrow[r, hook, "j"] \arrow[dr, equal] & K(f) \arrow[d, "\bar{f}"] & I \arrow[l, hook', "i"'] \\
& J &
\end{tikzcd}
\]
\[
\left\{
\begin{aligned}
\bar{f} \circ j &= \id_J\\
j \circ \bar{f} &= \left\{
\begin{aligned}
&j && \text{sur } J\\
&jf && \text{sur } I,
\end{aligned}
\right.
\end{aligned}
\right.
\qquad \text{i.e.} \quad
\boxed{
\begin{aligned}
j\bar{f}j &= j && \text{sur } J\\
j\bar{f}i &= jf && \text{sur } I
\end{aligned}
}\ ,
\]
et on a \texttt{[}$K = K(f)$\texttt{]}
\[
\begin{aligned}
\alpha :\ \id_K &\longrightarrow j\bar{f} && \text{i.e. homomorphisme fonctoriel}\\
x &\xrightarrow{\alpha(x)} j(\bar{f}(x)), && x \in \Ob K,
\end{aligned}
\]
\printedpage{6}
% head: A. Grothendieck : Dérivateurs, Ch. IV, Édition des Universités de Montpellier II et Paris VII
% foot: Version 2 février 2002
défini ainsi :
\[
\left\{
\begin{aligned}
&\text{Si } x \in J,\ \alpha(x) = \id_x\ :\ x \xrightarrow{\sim} j(\bar{f}(x)) = j(x) = x\\
&\text{Si } x \in I,\ \alpha(x) : \underbrace{x \longrightarrow f(x)}_{\text{dans } K(f)},\ \text{ou } \underbrace{f(x) \longrightarrow f(x)}_{\text{dans } J},\ \alpha(x) = \id_{f(x)}.
\end{aligned}
\right.
\]
On doit vérifier la fonctorialité pour $x$ variable dans $K$ \dots

\dgrblock{dgrb157}{The collage records a relative morphism}
\begin{dcomment}{exposition}{p79:La composition des flèches}{The collage records a relative morphism}
Grothendieck introduces \(K(f)\) to turn a morphism between coefficients on two different bases into an ordinary diagram on one base. This is a categorical form of the mapping cylinder, or more precisely the collage of the profunctor \(\Hom_J(f(-),-)\). The construction encodes the direction of the relative morphism in its mixed arrows.

For clarity, distinguish the copy of \(I\) from the copy of \(J\), even when they have the same objects. Internal arrows are those of \(I\) and \(J\); a mixed arrow \(x\to y\), with \(x\in I,y\in J\), is an arrow \(f(x)\to y\) in \(J\); there are no mixed arrows in the other direction. If \(a:x'\to x\) lies in \(I\), \(b:y\to y'\) lies in \(J\), and \(\beta:f(x)\to y\) represents a mixed arrow, composition is
\[
 \beta\circ a=\beta\circ f(a),\qquad
 b\circ\beta=b\beta
\]
in \(J\). Associativity follows directly from associativity in \(I,J\) and functoriality of \(f\). These formulas also prove that \(I\) is a sieve and \(J\) is its complementary cosieve.

Here is the universal property in full. Given presheaves \(F:I^o\to\mathcal M\), \(G:J^o\to\mathcal M\), and a natural transformation \(u:f^*G\to F\), define \(\Phi:K(f)^o\to\mathcal M\) to equal \(F\) and \(G\) on the two pieces. On the mixed arrow represented by \(\beta:f(x)\to y\), set
\[
 \Phi(\beta)=
 \bigl(G(y)\xrightarrow{G(\beta)}G(f(x))
       \xrightarrow{u_x}F(x)\bigr).
\]
Naturality of \(u\) proves compatibility with composition on the \(I\)-side; functoriality of \(G\) proves compatibility on the \(J\)-side. Conversely, restriction gives \(F,G\), and evaluation on the mixed arrow \(x\to f(x)\) represented by \(\mathrm{id}_{f(x)}\) gives \(u_x\). The two constructions are inverse.

A morphism between triples is a pair \(r:F\to F'\), \(s:G\to G'\) with \(r\,u=u'\,f^*s\). This equation is precisely naturality on the distinguished mixed arrows, and hence on every mixed arrow by the formula above. We have therefore obtained an isomorphism of categories, natural in \(\mathcal M\):
\[
 \mathrm{Fun}(K(f)^o,\mathcal M)
 \cong\{(F,G,u:f^*G\to F)\}.
\]
The isolated superscript \(o\) on the category of triples on manuscript page 82 is inconsistent with the preceding displayed universal property if the same variance is retained. One may either omit it in the modern formulation or dualize both sides. The printed reading is preserved in the transcription and recorded in the corrections table.

The retraction \(\bar f:K(f)\to J\) and transformation \(\alpha:\mathrm{id}\to j\bar f\) need no additional coherence data. In fact \(\bar f\dashv j\): for \(x\in K(f),y\in J\), the equality
\(\Hom_K(x,jy)=\Hom_J(\bar f x,y)\) follows from the defining cases. Its unit is precisely \(\alpha\), and its counit is the identity. This proves the naturality left for verification on page 80 and explains why the identities involving \(\bar f\) fit together.
\end{dcomment}


Soit $\Phi$ un $\mathcal{M}$-préfaisceau sur $K$. Il définit
\[
\begin{aligned}
F &= i^*(\Phi) && \in \Ob \mathbf{D}_{\mathcal{M}}(I)\\
G &= j^*(\Phi) && \in \Ob \mathbf{D}_{\mathcal{M}}(J).
\end{aligned}
\]
D'autre part $f^*(G)$ s'interprète comme $(j\bar{f})^*(\Phi)|I$, car
\[
\begin{aligned}
f^*(G)(i) &= G(\ \underbrace{f(i)}_{= \bar{f}(i) \,\in\, J}\ )\\
&= j^*(\Phi)(\bar{f}(i))\\
&= \Phi(j\bar{f}(i))\\
&= (j\bar{f})^*(\Phi)(i).
\end{aligned}
\]

\msp{page 81}
Donc
\[ f^*(G) \simeq (j\bar{f})^*(\Phi)|I. \]
Or la flèche
\[ \alpha : \id_K \longrightarrow j\bar{f} \]
donne
\[ \underbrace{\id_K(\Phi)}_{= \Phi} \xleftarrow{\Phi(\alpha)} (j\bar{f})^*(\Phi), \qquad \text{fonctoriel en } \Phi, \]
d'où, en prenant les restrictions à $I$,
\[ i^*(\Phi) \overset{\text{déf}}{=} F \xleftarrow{\ u\ } f^*G. \]

En fait, sans tellement se tortiller, on devrait pouvoir montrer que pour toute catégorie $\mathcal{M}$, la catégorie $\uHom(K^o, \mathcal{M})$ des $\mathcal{M}$-préfaisceaux sur $K$ est isomorphe à celle des triples $(F,G,u)$, avec $F \in \Ob \uHom(I^o, \mathcal{M})$, $G \in \uHom(J^o, \mathcal{M})$, et $u : f^*(G) \longrightarrow F$. Ainsi, $K$ apparaît comme solution d'un problème universel
\[ \uHom(K^o, \mathcal{M}) \longrightarrow \{F,\ G,\ u \text{ comme dessus}\}(\mathcal{M}). \]

Mais pour y préciser la flèche $\longrightarrow$ de façon conceptuelle simple, il est bon de se donner $K$ avec les structures supplémentaires \begin{tikzcd}[cramped, column sep=normal] I \arrow[r, hook, shift left, "j"] & K \arrow[l, shift left, "\bar{f}"] & I \arrow[l, hook', "i"'] \end{tikzcd}, $\bar{f}j = \id_J$, $\id_K \xrightarrow{\alpha} j\bar{f}$. Et ce système % [sic] I (first object, source of j and target of f-bar) -> J, since j : J -> K(f) and f-bar : K(f) -> J
\printedpage{7}
% head: A. Grothendieck : Dérivateurs, Ch. IV, Édition des Universités de Montpellier II et Paris VII
% foot: Version 2 février 2002
\msp{page 82}
$\{K,\ i,\ j,\ \bar{f},\ \alpha\}$ est (sauf erreur) défini à isomorphisme canonique près par la condition de représenter le foncteur $\mathcal{M} \longmapsto \{F,\ G,\ u\}^o(\mathcal{M})$ \dots

En particulier, le système défini par $G \in \Ob \mathbf{D}_{\mathcal{M}}(J)$
\[ (f^*(G),\ G,\ \id : f^*(G) \longrightarrow f^*(G)) \]
est défini sauf erreur par $\bar{f}^*(G)$ sur $K(f)$. On se propose [?] de trouver une interprétation similaire pour
\[ (F,\ f_*(F),\ f^*f_*F \longrightarrow F) \]
$(F \in \Ob \mathbf{D}(I))$. Le candidat naturel serait ici $\Phi = i_*(F)$. Admettons ce point (i.e. surtout la vérification du fait que $j^*(i_*(F)) \simeq f_*(F)$). Tout ceci montre que si $\mathbf{D}$ est à nouveau un dérivateur quelconque, on peut, par les formules précédentes, trouver un objet canonique $\Phi$ de $\mathbf{D}(K)$, en termes soit d'un $\mathbf{D}$-coefficient $F$ sur $I$ (en prenant $\Phi = i_*(F)$), soit d'un $\mathbf{D}$-coefficient $G$ sur $J$ (en prenant $\Phi = \bar{f}^*(G)$). D'autre part, on peut espérer définir canoniquement
\[
\begin{aligned}
\mathbf{D}(K) &\longrightarrow \mathbf{D}(I \times \bDelta_1)\\
\mathbf{D}(K) &\longrightarrow \mathbf{D}(J \times \bDelta_1)
\end{aligned}
\]
dont les composés avec les foncteurs ‘squelettes’
\[
\begin{aligned}
\mathbf{D}(I \times \bDelta_1) &\longrightarrow \uHom(\bDelta_1^o, \mathbf{D}(I))\\
\mathbf{D}(J \times \bDelta_1) &\longrightarrow \uHom(\bDelta_1^o, \mathbf{D}(J))
\end{aligned}
\]
\msp{page 83}
redonnent (pour $\Phi \in \Ob \mathbf{D}(K)$) les flèches squelettiques
\[
\left\{
\begin{aligned}
f^*(G) &\xrightarrow{\ u\ } F\\
G &\xrightarrow{\ v\ } f_*F
\end{aligned}
\right.
\]
associées à $(F,G,u)$. De façon précise, $f^*(G) \longrightarrow F$ se définit en termes de $\mathbf{D}$, et $\Phi \in \Ob \mathbf{D}(K)$, comme ci-dessus, et $v$ s'en déduit par adjonction.

\dgrblock{dgrb158}{The missing Kan extension comparison}
\begin{dcomment}{exposition}{p82:Le candidat naturel serait ici}{The missing Kan extension comparison}
Grothendieck's candidate \(i_*F\) is correct, and the comparison \(j^*i_*F\cong f_*F\) follows directly from the geometry of the mixed arrows. This is the point at which the collage begins to do mathematical work: a right Kan extension along an arbitrary functor is recovered from one along an inclusion.

For \(y\in J\), the comma category \((i/jy)\) has objects \(x\in I\) together with arrows \(i(x)\to j(y)\) in \(K(f)\). By definition, such an arrow is an arrow \(f(x)\to y\) in \(J\). Morphisms are the same commuting triangles in both descriptions. Hence there is a canonical isomorphism
\[
 (i/jy)\cong(f/y)
\]
commuting with the projections to \(I\), naturally in \(y\). The pointwise right Kan extension axiom in the presheaf convention computes both
\((i_*F)_{j(y)}\) and \((f_*F)_y\) as the homotopy limit of \(F\) over \((f/y)^o\). The comparison is therefore an isomorphism at every \(y\), and Der~2 makes it an isomorphism in \(\D(J)\).

There is a precise canonical map behind this pointwise calculation. The transformation
\(i\to jf\) induced by \(\alpha\) yields
\[
 f^*j^*i_*F\longrightarrow i^*i_*F\cong F.
\]
Its mate is \(j^*i_*F\to f_*F\). The comma-category identification proves that this mate is the isomorphism just computed. Consequently the displayed arrow on \(I\), after that identification, is exactly the counit \(f^*f_*F\to F\); agreement on objects is not being used as a substitute for agreement of the relevant natural transformations.

The other candidate is immediate but equally canonical. If \(\Phi=\bar f^*G\), then \(j^*\Phi=G\), \(i^*\Phi=f^*G\), and the transformation on \(I\) is the identity, since \(\bar f(\alpha_x)=\mathrm{id}_{\bar f(x)}\). These assertions hold for a derivator whenever the required indexing shapes are permitted. They use only restriction, right Kan extension, and the prederivator action on natural transformations.

In modern language the comma square expressing this argument is homotopy exact. The general calculus of such squares, with the comparison specified as a mate, is developed in \citet[Proposition 1.26]{Groth2013}; the relevant earlier convention was fixed in the commentary on Chapter~I, manuscript p.~31 (\emph{The base-change map as a mate}). The argument above completes the local proof rather than appealing to a stronger base-change statement with unchecked hypotheses.
\end{dcomment}


\dgrblock{dgrb159}{Completing both coherent adjunction arrows}
\begin{dcomment}{development}{p83:et $v$ s'en déduit par adjonction.}{Completing both coherent adjunction arrows}
Grothendieck leaves the chapter just after finding the characteristic map \(K(f)\to[1]\). His construction can be completed explicitly. The point is to use that map together with the retraction to \(J\), rather than to select a coherent representative of an arbitrary arrow in a homotopy category. We give the completion because it explains why both members of an adjoint pair can be treated by the same category.

\textbf{Collage completion.}\ours\
Let \(\D\) be a derivator in the presheaf convention, and let \(f:I\to J\). Assume its diagram category contains \(K=K(f)\), the products \(I\times[1]\), \(J\times[1]\), and the comma categories needed for the indicated right Kan extensions. There are functors
\[
 H:\D(K)\longrightarrow\D(I\times[1]),\qquad
 V:\D(K)\longrightarrow\D(J\times[1])
\]
whose skeletal arrows are, respectively, \(u:f^*G\to F\) and its adjoint \(v:G\to f_*F\), for \(F=i^*\Phi\), \(G=j^*\Phi\), and \(u\) induced by \(\alpha\).

\emph{Proof.}
Define \(h:I\times[1]\to K\) by
\[
 h(x,0)=i(x),\qquad h(x,1)=j(f(x)).
\]
On an arrow \(a:x\to x'\) at level \(0\) use \(i(a)\), and at level \(1\) use \(j(f(a))\). A mixed arrow \((x,0)\to(x',1)\), determined by \(a\), is sent to the mixed arrow represented by \(f(a)\). The definition of composition in the collage shows that this is a functor. Its vertical arrow at \(x\) is \(\alpha_{i(x)}\). Therefore \(H=h^*\) has restrictions \(F,f^*G\), and its arrow in the presheaf direction is exactly \(u:f^*G\to F\). This proves the first assertion without any Kan extension.

To construct the adjoint arrow, let \(\chi:K\to[1]\) be \(0\) on \(I\) and \(1\) on \(J\), and set
\[
 q=(\bar f,\chi):K\longrightarrow J\times[1],
 \qquad V=q_*.
\]
For \(y\in J\), the category \((q/(y,0))\) contains only objects from \(I\); its definition identifies it with \((f/y)\). Thus
\[
 (q_*\Phi)|_{J\times\{0\}}\cong f_*F.
\]
The category \((q/(y,1))\) has the terminal object given by \(j(y)\) and the identity of \((y,1)\). Indeed every object of that comma category specifies precisely one arrow to \(j(y)\), either an internal arrow of \(J\) or the mixed arrow represented by its \(J\)-component. In the presheaf convention the right Kan formula takes a homotopy limit over the opposite of this comma category, so the terminal object becomes an initial object and the limit is evaluation there. Consequently
\[
 (q_*\Phi)|_{J\times\{1\}}\cong G.
\]
These comparisons are natural, and their pointwise invertibility implies invertibility as coefficient diagrams.

It remains to identify the arrow, since its two endpoints alone would not determine it. Put \(E=q_*\Phi\), let \(t:E_1\to E_0\) be its skeletal arrow, and use the counit
\(\varepsilon:q^*E\to\Phi\).
Its restriction to \(J\) gives an isomorphism \(\lambda:E_1\to G\). Its restriction to \(I\) gives \(f^*E_0\to F\); let \(\theta:E_0\to f_*F\) be its mate. The preceding comma calculations say exactly that \(\theta\) is an isomorphism. Naturality of \(\varepsilon\) with respect to the mixed arrows \(x\to f(x)\) gives
\[
 (f^*E_0\longrightarrow F)\circ f^*t
 =
 u\circ f^*\lambda.
\]
After substituting the definitions of \(\theta\) and the counit of \(f^*\dashv f_*\), this says that
\(\theta\,t\,\lambda^{-1}:G\to f_*F\)
has adjoint \(u\). The adjunction bijection identifies it with \(v\). This proves the second assertion and the proposition. \(\square\)

Taking \(\Phi=\bar f^*G\), the first arrow is the identity on \(f^*G\), so the second is the unit \(G\to f_*f^*G\). Taking \(\Phi=i_*F\), the first is the counit \(f^*f_*F\to F\), while the second is the identity on \(f_*F\). Hence both adjunction transformations have canonical coherent realizations, natural in their coefficients. Applying the construction after a shift by \(A\) gives the corresponding statement in every \(\D(A)\); the compatibility with restriction in \(A\) follows from the product-square base change already used in Chapter III.

No strongness axiom was used. The input was the coherent \(\Phi\), and for the unit and counit that input was itself constructed by a specified Kan extension or restriction. The result does not supply a canonical section of the skeleton functor on all arrows.

The collage and the mate calculus are standard constructions. We checked the treatments of homotopy exact squares in \citet[Section 1]{Groth2013}, enriched derivator adjunctions in \citet[Sections 8--9]{GrothPontoShulman2014Additivity}, and pointwise Kan extension through modules in \citet[Section 5]{RiehlVerity2017}. We did not find in these sources the precise pair \(h^*,q_*\) written as a completion of this unfinished chapter. The marker above identifies our explicit completion and proof; it makes no claim of priority.
\end{dcomment}


L'idée naturelle ici, serait de ramener le problème de l'étoffement des flèches d'adjonction sqelettiques, au cas des immersions ouvertes ou fermées. Notons que les sous-catégories ouvertes de $K$ correspondent justement aux foncteurs $K \longrightarrow \bDelta_1$, en prenant les images inverses de l'ouvert générique. Il y a donc un unique % [sic] sqelettiques -> squelettiques, missing u (spelled squelettiques elsewhere in the text)
\[ \varphi :\ K(f) \longrightarrow \bDelta_1 \]
avec $I = \varphi^{-1}(\{0\})$. Donc \texttt{[fin du chapitre]}

\dgrblock{dgrb160}{A diagonal and its two projections}
\begin{dcomment}{examples}{p83:L'idée naturelle ici, serait de ramener}{A diagonal and its two projections}
Grothendieck's geometric construction can be seen completely in the case where \(I=\{a,b\}\) is discrete, \(J=e\), and \(f\) is the unique functor. Here \(K(f)\) is the cospan \(a\to *\leftarrow b\). Its presheaf coefficients have the opposite shape: an object \(\Phi\) is a coherent pair of arrows from \(G\) to \(F_a,F_b\).

The right Kan extension \(f_*F\) is the product \(F_a\times F_b\) in \(\D(e)\). Indeed Der~1 identifies \(\D(I)\) with \(\D(e)\times\D(e)\); the right adjoint to the diagonal restriction is the categorical product. The map \(q:K(f)\to[1]\) sends \(a,b\) to \(0\) and \(*\) to \(1\). Its right Kan extension sends the coherent pair to the single coherent arrow
\[
 G\longrightarrow F_a\times F_b
\]
whose two components are the original arrows. This proves directly, in this example, that the second construction is the adjoint of the first.

For \(\Phi=\bar f^*G\), both original arrows are identities. The output is the diagonal \(G\to G\times G\), as a specified object of \(\D([1])\). For \(\Phi=i_*F\), the value at \(*\) is \(F_a\times F_b\), and restriction along \(h\) gives the two projection arrows. The same collage has therefore produced the unit and counit of the diagonal-product adjunction, including their coherent representatives.

The empty case checks the direction of every construction. If \(I=\varnothing\), then \(K(f)=J\), \(q\) is the level-\(1\) inclusion, and \(q_*\Phi\) has value \(\Phi\) at level \(1\) and the terminal coefficient at level \(0\). Its presheaf arrow is \(\Phi\to1\), the unit of the adjunction between restriction to the empty category and its right Kan extension. Interchanging the two endpoints would give the wrong arrow even in this simplest case.
\end{dcomment}
\dgrchapter{V}{Chapter V}{Catégories de chemins et localisation}
\dgrpart{0}{III-V}{Manuscript pages 1--52}{III-V\enspace\textperiodcentered\enspace Chapter V\enspace\textperiodcentered\enspace manuscript pages 1--52}{Chapter V \textperiodcentered\ \textup{III-V}}
% ---- from Der-ChV-234_p001-005.tex
\printedpage{}
% head: (none); no page number printed on this page

\begin{center}
\vspace*{6cm}

{\LARGE\textsf{LES DÉRIVATEURS}}

\vspace{1.5cm}

{\large\textbf{ALEXANDRE GROTHENDIECK}}

\vspace{2.5cm}

{\large\textbf{Chapitre V}}

\vspace{0.4cm}

{\large\textbf{Catégories de chemins et localisation}}
\end{center}

\clearpage

\printedpage{}
% head: (none); no page number printed on this page

\vspace*{6cm}

Ce texte a été déchiffré et transcrit en \LaTeXe{} par M. Künzer. Il a été édité par M. Künzer, J. Malgoire, et G. Maltsiniotis. La transcription est aussi fidèle que possible au manuscrit. Pour les quelques corrections évidentes, ou rares commentaires des éditeurs, ainsi que pour la numérotation originale des pages du manuscrit, les caractères de machine à écrire \texttt{[typewriter]} entre crochets sont utilisés. Un point d'interrogation entre crochets signifie que l'on n'est pas sûr du mot qui précède.

Cette édition est provisoire. Les remarques, commentaires et corrections sont bienvenus. Envoyer un message à :

\textbf{maltsin@math.jussieu.fr}\\
\hspace*{\fill}G. Maltsiniotis\hspace*{2cm}

\clearpage

\printedpage{1}
% head: \textbf{A. Grothendieck :} \textit{Dérivateurs,} Ch. V, Édition des Universités de Montpellier II et Paris VII
% foot: Version 9 septembre 2001

\msp{page 1}

\section*{1\quad La catégorie des chemins}

Soit $\Phi$ la catégorie suivante :
\[
\left\{
\begin{array}{rcll}
\Ob\Phi & = & \mathbf{Z} & \\
\Fl\Phi & = & \Ob\Phi\sqcup\mathbf{Z} & \qquad (\Ob\Phi\leftrightarrow\text{identités})
\end{array}
\right.
\]
On note $\{\alpha\}$, ou simplement $\alpha$, l'objet qui correspond à $\alpha\in\mathbf{Z}$. On note $u_\alpha$ la flèche qui correspond à l'entier $\alpha\in\mathbf{Z}$.
\[
\begin{array}{rcl@{\qquad}l@{\quad}l@{\quad}l}
s(u_\alpha) & = & \alpha & \text{si } \alpha \text{ pair}, & \alpha+1 & \text{si } \alpha \text{ impair}\\[1ex]
b(u_\alpha) & = & \alpha+1 & \text{si } \alpha \text{ pair}, & \alpha-1 & \text{si } \alpha \text{ impair}
\end{array}
\]
\[
\texttt{[plutôt}\qquad b(u_\alpha) \;=\; \alpha+1 \quad \texttt{si}\ \alpha\ \texttt{pair,}\quad \alpha \quad \texttt{si}\ \alpha\ \texttt{impair]}
\]

Il n'y pas à définir de compositions, car si deux flèches sont composables, elles sont l'identité. En fait, $\Phi$ est une catégorie \textit{ordonnée}, $\{\alpha\}<\{\beta\}$ si et seulement si $\alpha$ pair, $\beta=\alpha+1$ \texttt{[plutôt} $\beta=\alpha+1$ \texttt{ou} $\beta=\alpha-1$\texttt{]}. % [sic] Il n'y pas -> Il n'y a pas, verb "a" missing

On note, pour $n\in\mathbf{N}$,
\[
\Phi_n\subseteq\Phi,\quad \Ob\Phi_n=\{\alpha\in\Ob\Phi\mid -n\le\alpha\le+n\},
\]
la sous-catégorie pleine de $\Phi$ dont les objets sont les $\{\alpha\}$ avec $-n\le\alpha\le+n$. Ainsi
\[
\Phi_1=\{-1\longleftarrow 0\longrightarrow 1\},
\]
la duale est
\[
\Psi_1=\Phi_1^{o}=\{-1\longrightarrow 0\longleftarrow 1\}.
\]
On peut considérer $\Phi$ comme somme amalgamée d'une suite ordonnée infinie (de deux côtés) de copies de $\Phi_1$, ou aussi d'une suite similaire de copies de $\Psi_1$.
\[
\begin{tikzcd}[column sep=small, execute at end picture={
\draw[decorate,decoration={brace,amplitude=4pt}] ([yshift=10pt]\tikzcdmatrixname-1-2.center) -- ([yshift=10pt]\tikzcdmatrixname-1-4.center) node[midway,above=4pt]{$\scriptstyle\Phi_1$};
\draw[decorate,decoration={brace,amplitude=4pt}] ([yshift=10pt]\tikzcdmatrixname-1-4.center) -- ([yshift=10pt]\tikzcdmatrixname-1-6.center) node[midway,above=4pt]{$\scriptstyle\Phi_1$};
\draw[decorate,decoration={brace,amplitude=4pt}] ([yshift=10pt]\tikzcdmatrixname-1-6.center) -- ([yshift=10pt]\tikzcdmatrixname-1-8.center) node[midway,above=4pt]{$\scriptstyle\Phi_1$};
\draw[decorate,decoration={brace,amplitude=4pt}] ([yshift=10pt]\tikzcdmatrixname-1-8.center) -- ([yshift=10pt]\tikzcdmatrixname-1-10.center) node[midway,above=4pt]{$\scriptstyle\Phi_1$};
\draw[decorate,decoration={brace,amplitude=4pt}] ([yshift=10pt]\tikzcdmatrixname-1-10.center) -- ([yshift=10pt]\tikzcdmatrixname-1-12.center) node[midway,above=4pt]{$\scriptstyle\Phi_1$};
\draw[decorate,decoration={brace,mirror,amplitude=4pt}] ([yshift=-10pt]\tikzcdmatrixname-1-3.center) -- ([yshift=-10pt]\tikzcdmatrixname-1-5.center) node[midway,below=4pt]{$\scriptstyle\Psi_1$};
\draw[decorate,decoration={brace,mirror,amplitude=4pt}] ([yshift=-10pt]\tikzcdmatrixname-1-5.center) -- ([yshift=-10pt]\tikzcdmatrixname-1-7.center) node[midway,below=4pt]{$\scriptstyle\Psi_1$};
\draw[decorate,decoration={brace,mirror,amplitude=4pt}] ([yshift=-10pt]\tikzcdmatrixname-1-7.center) -- ([yshift=-10pt]\tikzcdmatrixname-1-9.center) node[midway,below=4pt]{$\scriptstyle\Psi_1$};
\draw[decorate,decoration={brace,mirror,amplitude=4pt}] ([yshift=-10pt]\tikzcdmatrixname-1-9.center) -- ([yshift=-10pt]\tikzcdmatrixname-1-11.center) node[midway,below=4pt]{$\scriptstyle\Psi_1$};
}]
\cdots & {} & -4 \arrow[l] \arrow[r] & -3 & -2 \arrow[l] \arrow[r] & -1 & 0 \arrow[l] \arrow[r] & 1 & 2 \arrow[l] \arrow[r] & 3 & 4 \arrow[l] \arrow[r] & {} & \cdots
\end{tikzcd}
\]

\dgrblock{dgrb161}{A path remembers the arrows that will be inverted}
\begin{dcomment}{exposition}{p1:comme somme amalgamée d'une suite ordonnée infinie}{A path remembers the arrows that will be inverted}
Grothendieck is preparing a geometric description of localization. A morphism in a category of fractions is represented by a finite word of arrows, some traversed in their original direction and some traversed backwards. The alternation imposed here is a way of giving all such words a common shape. Inserting identity arrows makes any finite zigzag alternating; adding identities at its ends gives an interval symmetric about zero. The identities are bookkeeping, whereas the arrows traversed backwards carry the mathematical restriction imposed by the chosen class of denominators.

After the edition's correction of the target of $u_\alpha$, the category $\Phi$ is the poset on $\mathbf Z$ whose only strict relations are
\[
  2i<2i-1,\qquad 2i<2i+1.
\]
Every even vertex is a source and every odd vertex a sink. There is no composite of two nonidentity arrows. The sentence saying that two composable arrows ``sont l'identit\'e'' must therefore mean that at least one of them is an identity. Composing an identity with an arbitrary arrow supplies the elementary exception to its literal wording.

The subcategory $\Phi_n$ has $2n$ edges, not $n$. Thus the manuscript's later references to the parity of a path concern the index $n$, or half its number of edges. This distinction explains the repair on manuscript page~14. A path of index~1 has shape $\cdot\leftarrow\cdot\rightarrow\cdot$; two such paths joined end to end have the opposite orientation from $\Phi_2$. The correct assertion at the end of the present page is that $\Phi_n$, for even $n>0$, is assembled from copies of $\Psi_1=\Phi_1^o$. The printed $\Psi_n$ has the wrong orientation. We keep that reading in the source and use the corrected statement in the mathematics.

The nerve of every finite $\Phi_n$ is an ordinary subdivided interval. Nevertheless the category itself, including the directions of its arrows, matters: a functor from it to $X$ is an actual zigzag in $X$, and its natural transformations are actual commutative ladders. Passing immediately to the geometric interval would discard precisely the information needed to distinguish denominators from the other arrows.
\end{dcomment}


Les $\Phi_n$, pour $n$ impair, sont de même somme de copies de $\Phi_1$ mises bout à bout, et les $\Psi_n$, pour $n$ pair $>0$, somme de copies de $\Psi_1$.

\clearpage

\printedpage{2}
% head: \textbf{A. Grothendieck :} \textit{Dérivateurs,} Ch. V, Édition des Universités de Montpellier II et Paris VII
% foot: Version 9 septembre 2001

Les $\Phi_n$ forment un système inductif avec foncteurs de transition pleinement fidèles, dont $\Phi$ est la limite inductive
\[
\Phi=\varinjlim_n\Phi_n.
\]
Mais pour la suite on va faire des $\Phi_n$ un système \textit{projectif}, en définissant pour tout $n\le m$ une \textit{rétraction} $\varphi_{n,m}$ de $\Phi_m$ sur $\Phi_n$, par

\msp{page 2}
\[
\left\{
\begin{array}{l}
\varphi_{n,m}:\Phi_m\longrightarrow\Phi_n\subset\Phi_m\\[1ex]
\varphi_{n,m}(\{\alpha_i\})=
\left\{
\begin{array}{rcll@{\quad}l@{\ \ }r}
\alpha_i & = & i & \in\Ob\Phi_n, & \text{i.e.} & -n\le i\le n\\[1ex]
\alpha_i & = & n & & \text{si} & n\le i\le m\\[1ex]
\alpha_i & = & -n & & \text{si} & -m\le i\le -n
\end{array}
\right.
\end{array}
\right.
\]

\texttt{[plutôt}
\[
\begin{array}{rcl}
\varphi_{n,m}:\Phi_m & \longrightarrow & \Phi_n\subset\Phi_m\\[1ex]
\varphi_{n,m}(\{i\}) & = &
\left\{
\begin{array}{l@{\qquad}l@{\quad}l@{\ \ }r}
\{i\} & \texttt{si}\ i\in\Ob\Phi_n, & \texttt{i.e} & -n\le i\le n\\[1ex]
\{n\} & & \texttt{si} & n\le i\le m\\[1ex]
\{-n\} & & \texttt{si} & -m\le i\le -n\ \texttt{]}
\end{array}
\right.
\end{array}
\]
On note le système projectif d'ensembles ordonnés ainsi obtenu par
\[
\bPhi=(\Phi_n)_{n\in\mathbf{N}}\in\mathrm{Pro}(\mathrm{Cat}).
\]
Soit $X$ une catégorie, alors pour $n\in\mathbf{N}$
\[
\uCh_n(X)\overset{\text{déf}}{=}\uHom(\Phi_n,X)
\]
est la catégorie des diagrammes de $X$ de la forme
\[
A_{-n}\cdots\longleftarrow A_{-2}\xrightarrow{u_{-2}}A_{-1}\xleftarrow{u_{-1}}A_0\xrightarrow{u_0}A_1\xleftarrow{u_1}A_2\xrightarrow{u_2}A_3\longleftarrow\cdots A_n
\tag*{$(*)$}
\]
avec $2n$ flèches bout à bout $u_i$ $(-n\le i\le n)$ \texttt{[plutôt (}$-n\le i\le n-1$\texttt{)]}, deux flèches consécutives allant toujours en sens opposé,
\[
\text{soit}\ \ \cdot\longrightarrow\cdot\longleftarrow\cdot,\ \ \text{soit}\ \ \cdot\longleftarrow\cdot\longrightarrow\cdot\,.
\]
Le foncteur canonique déduit de $\varphi_{n,m}:\Phi_m\longrightarrow\Phi_n$
\[
\varphi_X^{m,n}:\uCh_n(X)\longrightarrow\uCh_m(Y),\qquad n\le m,
\]
\printedpage{3}
% head: \textbf{A. Grothendieck :} \textit{Dérivateurs,} Ch. V, Édition des Universités de Montpellier II et Paris VII
% foot: Version 9 septembre 2001
\texttt{[plutôt} $\uCh_m(X)$\texttt{]} est le `foncteur de dégénérescence', consistant à compléter un diagramme $(*)$, par $m-n$ termes $A_{-n}$ avec des identités à gauche, et par $m-n$ termes \texttt{[}$A_n$\texttt{]} avec des identités à droite, donc pour $n$ pair on trouve
\[
\underbrace{A_{-n}\cdots\xrightarrow{\text{id}}A_{-n}}_{m-n\text{ termes}}\xleftarrow{\text{id}}A_{-n}\longrightarrow A_{-(n-1)}\cdots\cdots A_{n-1}\longleftarrow A_n\xrightarrow{\text{id}}\underbrace{A_n\xleftarrow{\text{id}}A_n\longrightarrow\cdots A_n}_{m-n\text{ termes}}.
\]
C'est un foncteur pleinement fidèle (ce qui correspond au fait que $\varphi_{n,m}$ est surjectif sur les objets et sur les flèches. % [sic] flèches. -> flèches)., opening parenthesis never closed

\msp{page 3}
On pose
\[
\underbrace{\uCh_\infty(X)}_{\text{ou simplement }\uCh(X)}\qquad=\uHom(\bPhi,X)\overset{\text{déf}}{=}\varinjlim_n\uCh_n(X)
\]
(\footnote{L'ensemble des objets de cette catégorie se notera $\mathrm{Ch}_\infty(X)$ ou $\mathrm{Ch}(X)$, de même on définit les $\mathrm{Ch}_n(X)$.}) (\footnote{dans cette formule: $\bPhi$ procatégorie}). On a une inclusion pleinement fidèle
\[
\uCh_\infty(X)\lhook\joinrel\longrightarrow\uHom(\Phi,X),
\]
où la deuxième catégorie est celle des diagrammes infinis (dans les deux directions)
\[
\cdots\longleftarrow A_{-2}\xrightarrow{u_{-2}}A_{-1}\xleftarrow{u_{-1}}A_0\xrightarrow{u_0}A_1\xleftarrow{u_1}A_2\xrightarrow{u_2}\cdots
\]
et où $\uCh_\infty(X)$ s'identifie à la sous-catégorie pleine formées des objets pour lesquels les flèches $u_i$ sont des identités, pour $|i|$ assez grand. % [sic] formées -> formée, agreement with "la sous-catégorie"

On appellera $\uCh_\infty(X)$ la \textit{catégorie des chemins} dans $X$, ses objets seront les chemins dans $X$. Ceux qui proviennent de $\uCh_n(X)$ seront appelés \textit{chemins standards de longueur} $2n$. (Le terme `standard' s'applique ici à nos conventions de notations. Ainsi nous ne considérons que des chemins de longueur paire, et pour une longueur déterminée, nous choisissons l'une parmi les deux espèces raisonnables, par exemple pour les chemins de longueur 2 nous avons pris seulement
\[
A_{-1}\longleftarrow A_0\longrightarrow A_1
\]
et non pas
\[
A_{-1}\longrightarrow A_0\longleftarrow A_1.)
\]

\dgrblock{dgrb162}{The pro-interval and the eventually constant diagram}
\begin{dcomment}{exposition}{p3:ses objets seront les chemins dans}{The pro-interval and the eventually constant diagram}
Grothendieck wants paths of arbitrary finite length without making their length part of the final object. He obtains them by extending a finite diagram with identity arrows at both ends. The projective system of intervals and the inductive system of path categories describe the same operation from opposite sides.

Write $r_{n,m}:\Phi_m\to\Phi_n$ for the numerical clamping map
\[
 r_{n,m}(i)=\max(-n,\min(i,n)),\qquad n\leq m.
\]
It is a functor: an edge inside $[-n,n]$ is unchanged, and every other edge is sent to an identity. The equality $r_{n,m}r_{m,l}=r_{n,l}$ is literal, so precomposition gives a strict inductive system. Its transition functors are fully faithful, for a reason slightly more precise than mere surjectivity on objects and arrows. If $F,G:\Phi_n\to X$ and $\eta:Fr_{n,m}\to Gr_{n,m}$, naturality on each collapsed edge says that the two adjacent components of $\eta$ agree. They therefore agree throughout either constant tail. Restriction to $\Phi_n$ determines $\eta$ uniquely and yields an arbitrary natural transformation $F\to G$. This proves the required bijection on Hom sets.

Surjectivity alone would not prove this assertion for an arbitrary functor. The map from a discrete two-point category to the point is surjective on objects and arrows, while precomposition gives the diagonal $X\to X\times X$, which need not be full. What makes clamping work is the chain of edges sent to identities inside each fibre.

The filtered colimit is consequently realized as the full subcategory of $\operatorname{Fun}(\Phi,X)$ consisting of diagrams whose edges are identities outside a finite interval. A natural transformation between two such diagrams also becomes constant in both tails: choose one interval containing the nonidentity edges of both diagrams and apply naturality to the identity edges beyond it. Thus both objects and morphisms descend to a common finite stage. No assertion about commuting an arbitrary functor category with filtered colimits is being used here.

The notation $\operatorname{Fun}(\boldsymbol\Phi,X)$ is a convenient enriched version of the usual pro-object formula: it means $\mathop{\rm colim}_n\operatorname{Fun}(\Phi_n,X)$ as a category. Its set of objects is not a space of arbitrary infinite paths. The eventual identity condition is essential, and differs even from eventual invertibility.

Reflection $i\mapsto-i$ is the unique nonidentity automorphism of $\Phi_n$ when $n>0$, since an automorphism of the underlying finite path is either the identity or reflection and both preserve the alternating orientation. For $n=0$ the automorphism group is trivial. This harmless boundary case will also matter when the two endpoints coincide.
\end{dcomment}


\msp{page 4}
Je me rends compte que mes notations $u_i$ ne sont pas commodes pour la dualité, je vais les changer, en $u_i$ pour les flèches entre objets $\{i\}$ avec $i\ge0$, et $v_i$ la flèche symétrique de $u_i$ par rapport à l'origine. (Le groupe des automorphismes des $\Phi_n$ est réduit au groupe à deux éléments, le seul automorphisme non trivial étant la symétrie $\{i\}\longrightarrow\{-i\}$. Donc nos chemins seront notés désormais % [sic] {-i}. -> {-i}). , opening parenthesis "(Le groupe" not closed on this page
\[
A_{-2}\xrightarrow{v_1}A_{-1}\xleftarrow{v_0}A_0\xrightarrow{u_0}A_1\xleftarrow{u_1}A_2\longrightarrow\cdots.
\tag*{$(*)$}
\]

% ---- from Der-ChV-234_p006-010.tex
\printedpage{4}
% head: A. Grothendieck : Dérivateurs, Ch. V, Édition des Universités de Montpellier II et Paris VII
% foot: Version 9 septembre 2001

(Ce n'est pas non plus enthousiasmant \dots)

Soit
\[
\varepsilon=\{-1,+1\}
\]
l'ensemble ordonné \textit{discret,} réduit aux deux éléments $-1$, $+1$, et définissons des foncteurs d'inclusion (pleinement fidèles)
\[
\begin{aligned}
\alpha_n:\varepsilon\longrightarrow\Phi_n\quad & \alpha_n(-1)=-n\\
& \alpha_n(1)=n,
\end{aligned}
\]
ces foncteurs sont compatible avec les $\varphi_{n,m}$, d'où % [sic] compatible -> compatibles, accord avec « ces foncteurs »
\[
\alpha_\infty:\varepsilon\longrightarrow\boldsymbol{\Phi}\in\mathrm{ProCat},
\]
d'où pour tout $X$ un foncteur
\[
\kappa:\uHom(\boldsymbol{\Phi},X)\longrightarrow\uHom(\varepsilon,X)\simeq X\times X
\]
i.e.

\msp{page 5}
\[
\begin{aligned}
\uCh(X) &\longrightarrow X\times X\\
c &\longrightarrow (\sigma(x),\beta(x)).
\end{aligned}
\]

\texttt{[plutôt} $(\sigma(c),\beta(c))$\texttt{]} Pour tout chemin $c$ dans $X$ i.e. un diagramme $(*)$ (satisfaisant à la condition de dégénérescence à l'infini), on a une source ou origine du chemin, et un but ou une \textit{extrémité}
\[
\left\{
\begin{aligned}
\sigma(x) &= A_i\quad\text{pour } i\leq -\mathrm{long}(x)\\
\beta(x) &= A_i\quad\text{pour } i\geq \mathrm{long}(x)
\end{aligned}
\right.
\]
\texttt{[remplacer} $x$ \texttt{par} $c$\texttt{; pour long}$(c)$ \texttt{défini convenablement]}. La fibre du foncteur précédent en $x$, $y$ sera notée
\[
\uCh(X;x,y)\qquad\text{ou simplement }\uCh(x,y),
\]
\textit{catégorie des chemins de $x$} à \textit{$y$.}

\dgrblock{dgrb163}{Fixing the endpoints also fixes the vertical arrows}
\begin{dcomment}{examples}{p5:La fibre du foncteur précédent en}{Fixing the endpoints also fixes the vertical arrows}
Grothendieck's source--target functor allows a path to become a candidate morphism from a specified object to another. Its fibre over $(x,y)$ is a strict categorical fibre. An object is an eventually constant zigzag with tails $x$ and $y$; a morphism is a natural transformation whose two tail components are $1_x$ and $1_y$. This last condition is indispensable for the interpretation by a Hom set.

There are two different full-subcategory statements in these pages. The category of $\Sigma$-paths is full in the category of all paths, because it imposes a condition only on the objects. Likewise its fixed-endpoint fibre is full in the fixed-endpoint fibre of all paths. But a fixed-endpoint fibre is generally not full in the entire functor category. On pages~15--16 and~24 the expressions ``sous-cat\'egorie pleine'' must be read with this correction: the endpoint components of a natural transformation are required to be identities.

For a concrete test take the one-object category $BG$ of a nontrivial group $G$, and let $\Sigma$ contain only the identity. The paths representing $1$ and $g\ne1$ may be written, at index~2,
\[
 (1,1,1,1),\qquad (g,1,1,1),
\]
where the entries label the four alternating arrows. There is a natural transformation from the first diagram to the second with vertex components $(1,g,g,g,g)$. It satisfies all four naturality squares, but its final component is $g$, not the identity. Thus it is not a morphism in the strict fibre. If one used the full subcategory on objects with the prescribed endpoints, this transformation would identify $1$ with $g$ on connected components, although localization at identities leaves $BG$ unchanged.

For an ordinary arrow category, the same convention makes the fibre over $(x,y)$ discrete: a morphism from $f:x\to y$ to $g:x\to y$ with both vertical sides identities exists precisely when $f=g$, and is then unique. This is the precise assertion used on page~8.

The map $\varepsilon\to\Phi_n$ is fully faithful for $n>0$. At $n=0$ it sends both endpoint labels to the same object and is not fully faithful. It still defines the source--target functor, whose index-zero stage is the diagonal $X\to X\times X$. Consequently a path of index zero occurs in the fibre only when $x=y$, exactly as the identity construction requires.
\end{dcomment}


\printedpage{5}
% head: A. Grothendieck : Dérivateurs, Ch. V, Édition des Universités de Montpellier II et Paris VII
% foot: Version 9 septembre 2001

\section*{2\quad Calcul d'une catégorie de fractions $X\Sigma^{-1}$}

Avant de continuer le formalisme des chemins (composition etc.), je vais tester l'utilité de la notion de catégorie des chemins, pour le `calcul' des $X\Sigma^{-1}$, où
\[
\Sigma\subset\mathrm{Fl}(X),\qquad\Sigma\text{ contient les identités}
\]
est une partie autrement \textit{quelconque} de $\mathrm{Fl}(X)$. On sait que $\mathrm{Ob}\,X\Sigma^{-1}=\mathrm{Ob}\,X$, et tout revient à définir les
\[
\Hom_{X\Sigma^{-1}}(x,y),\qquad x,y\in\mathrm{Ob}\,X,
\]
\msp{page 6}
et pour $x,y,z\in\mathrm{Ob}\,X$ la composition
\[
\Hom_\Sigma(x,y)\times\Hom_\Sigma(y,z)\xrightarrow{(u,v)\mapsto vu}\Hom_\Sigma(x,z),
\]
(où on note $\Hom_\Sigma$ au lieu de $\Hom_{X\Sigma^{-1}}$, pour alléger les notations).

Grâce à $\Sigma$, on définit une sous-catégorie pleine
\[
\uCh_\Sigma(X)\hookrightarrow\uCh(X),
\]
(\footnote[3]{L'ensemble de ses objets est noté $\mathrm{Ch}(X)$. Ses éléments sont les $\Sigma$-chemins de $x$ à $y$ \texttt{[de} $x$ \texttt{à} $y$\texttt{?]}.}) par la condition sur $c\in\mathrm{Ch}(X)$ : Pour $i$ impair, \texttt{[plutôt pair]} i.e. $i-1\xleftarrow{u_i}i$ ($u_i$ est dite `rétrograde'), $c(u_i)\in\Sigma$, i.e. les `tronçons' rétrogrades dans le chemin $c$ doivent être dans $\Sigma$. On définit de même les sous-catégories pleines
\[
\uCh_\Sigma(X;x,y)=\{c\in\uCh(X;x,y)\mid c\in\uCh_\Sigma(X)\},
\]
ou encore la fibre en $(x,y)$ du foncteur source-but pour $\uCh_\Sigma(X)$
\[
\uCh_\Sigma(X)\longrightarrow X\times X.
\]

\dgrblock{dgrb164}{Why a commutative ladder gives equality after localization}
\begin{dcomment}{exposition}{p6:les `tronçons' rétrogrades dans le chemin}{Why a commutative ladder gives equality after localization}
Grothendieck places no calculus-of-fractions hypothesis on $\Sigma$: it is any set of arrows containing the identities. Its closure under composition is unnecessary for the path construction. The hypothesis concerns only the backwards horizontal arrows; the vertical arrows of a morphism of paths may be arbitrary arrows of $X$.

To see why this is sufficient, let $Q:X\to X[\Sigma^{-1}]$ be the localization and consider a square around a backwards edge,
\[
\begin{tikzcd}[column sep=large]
 A_{i-1}\arrow[d,"a_{i-1}"'] & A_i\arrow[l,"s"']\arrow[d,"a_i"]\\
 B_{i-1} & B_i\arrow[l,"t"']
\end{tikzcd}
\qquad s,t\in\Sigma.
\]
Its commutativity gives $Q(a_{i-1})Q(s)=Q(t)Q(a_i)$, hence
\[
 Q(a_i)Q(s)^{-1}=Q(t)^{-1}Q(a_{i-1}).
\]
The square therefore remains commutative after reversing its horizontal arrows. Apply this to every backwards edge and then compose the resulting forward squares. The endpoint components give a commutative square between the two composite arrows in the localization. In a strict endpoint fibre those components are identities, so the two composites are equal.

This proves the functor to the arrow category and the induced function on $\pi_0$ without assuming that any vertical arrow becomes invertible. It also explains the deliberate use of connected components: an equality may be witnessed by a finite zigzag of natural transformations rather than a single transformation. The later reparametrization argument is designed to produce just such zigzags.

The parity notation changes during the construction. It is safest to number the edges consecutively from the left whenever discussing a fixed even-index interval; its even-numbered edges are then the backwards ones. This numerical convention must not be confused with the original label $u_\alpha$ on the doubly infinite interval, or with the symmetric $u_i,v_i$ notation introduced on page~4. The invariant definition is simply that an edge traversed right-to-left belongs to $\Sigma$.
\end{dcomment}


On définit une application canonique
\[
\mathrm{Ch}_\Sigma(X,x,y)\longrightarrow\Hom_\Sigma(x,y),
\]
\msp{page 7}
en associant `a tout chemin de $x$ dans $y$ le $\Sigma$-homomorphisme de $x$ dans $y$ obtenu en \textit{inversant} (dans $X\Sigma^{-1}$) les tronçons rétrogrades (ce qui est possible, puisqu'ils sont dans $\Sigma$), de sorte qu'on trouve dans $X\Sigma^{-1}$ un `chemin direct' % [sic] `a -> à, accent imprimé comme guillemet ouvrant
\[
x=A_{-n}\longrightarrow\cdots\longrightarrow A_{-1}\longrightarrow A_0\longrightarrow A_1\longrightarrow A_2\cdots A_n=A_{n+1}=\cdots=y.
\]
et on compose tous ces morphismes.

En fait, on trouve aussitôt que l'on définit ainsi un \textit{foncteur}
\[
\uCh_\Sigma(X)\longrightarrow\uFl(X\Sigma^{-1}),
\]
\texttt{[}$\uFl$ \texttt{catégorie des flèches]} du fait que si on a un diagramme

\printedpage{6}
% head: A. Grothendieck : Dérivateurs, Ch. V, Édition des Universités de Montpellier II et Paris VII
% foot: Version 9 septembre 2001
\[
\begin{tikzcd}
A_{i-1} \arrow[d, "f_{i-1}"] & A_i \arrow[l, "u_i"'] \arrow[d, "f_i"] \\
A'_{i-1} & A'_i \arrow[l, "u'_i"']
\end{tikzcd}
\qquad\qquad (i\text{ pair})
\]
avec $u_i$, $u'_i$ dans $\Sigma$, ce diagramme, en tant que diagramme dans $X\Sigma^{-1}$, est commutatif avec flèches horizontales des \textit{isomorphismes,} et reste donc commutatif en inversant lesdites flèches. D'où un foncteur induit
\[
\uCh_\Sigma(X;x,y)\longrightarrow\uFl(\underbrace{X\Sigma^{-1}}_{X'};x,y).
\]

Mais pour toute catégorie $X'$, et des objets $x$, $y$ dans $X'$, la catégorie $\uFl(X',x,y)$ est
\msp{page 8}
une catégorie \textit{discrète,} ayant comme ensemble d'objets $\Hom_{X'}(x,y)$. Ainsi le foncteur précédant définit en fait une application canonique % [sic] précédant -> précédent, adjectif
\[
\boxed{\pi_0\uCh_\Sigma(X;x,y)\longrightarrow\Hom_\Sigma(x,y).}
\]
Je dis que cette application est toujours bijective. Pour le voir, nous allons \textit{constuire} `from scratch' la catégorie $X\Sigma^{-1}$, en prenant comme objets ceux de $X$, et en \textit{définissant} % [sic] constuire -> construire, lettre omise
\[
\boxed{\Hom_\Sigma(x,y)=\pi_0\uCh_\Sigma(X;x,y).}
\]
Il faut donc définir la composition
\[
\Hom_\Sigma(x,y)\times\Hom_\Sigma(y,z)\longrightarrow\Hom_\Sigma(x,z).
\]
Pour ceci, on va définir des foncteurs `composition de chemins'
\[
\begin{aligned}
\uCh_m(X;x,y)\times\uCh_n(X;y,z) &\longrightarrow \uCh_{m+n}(X;x,z)\\
(c,c') &\longmapsto c'\circ c,
\end{aligned}
\]
mais en nous bornant pour simplifier au cas où $m$, $n$ ont \textit{la même parité,} de telle façon qu'en `mettant bout à bout' $c$ et $c'$ (dans cet ordre), on trouve bien un chemin de type $n+m$.

\dgrblock{dgrb165}{Localization without a short-fraction hypothesis}
\begin{dcomment}{context}{p8:Je dis que cette application est toujours bijective.}{Localization without a short-fraction hypothesis}
Grothendieck now proposes to recover the ordinary localization directly from categories of paths. The basic localization by composable words and inverse symbols is classical; compare \citet[Chapter~I, pp.~6--21]{GabrielZisman1967} and the explicit word formulation in \citet[\S3.1--3.3]{DwyerKan1980Simplicial}. His additional object of study is the category of presentations of a morphism, whose connected components recover those words modulo relations.

A short fraction such as $x\leftarrow x'\to y$ need not represent every localized morphism unless one imposes extra hypotheses. Arbitrary finite zigzags do. The use of all indices thus avoids an Ore condition, a model structure, and any saturation assumption on the original set $\Sigma$. Enlarging $\Sigma$ to the subcategory it generates changes neither the ordinary localization nor the final set of components, but the unlocalized path categories can change. They should not be silently identified.

There are three tasks in the proposed proof. First, concatenation must descend from particular finite representatives to their classes. Second, the original arrows of $X$ must determine a functor to those classes, and arrows in $\Sigma$ must become invertible. Finally, every functor inverting $\Sigma$ must factor uniquely by evaluating zigzags. The first task occupies most of the combinatorics. The second is where a generalized interval lemma is temporarily assumed. We supply its relative version below and also give direct witnesses for the particular identities needed here. The universal property then follows by an explicit argument at the end of the section.
\end{dcomment}


\msp{page 9}
Cette opération est `transposée' du couple de foncteurs
\[
\begin{aligned}
\Phi_n &\xrightarrow{\alpha} \Phi_{n+m}\quad i\longmapsto (i-m)\\
\Phi_m &\xrightarrow{\beta} \Phi_{n+m}\quad i\longmapsto (i+n),
\end{aligned}
\]
\printedpage{7}
% head: A. Grothendieck : Dérivateurs, Ch. V, Édition des Universités de Montpellier II et Paris VII
% foot: Version 9 septembre 2001
où on aura
\[
\alpha(\{n\})=\beta(\{-m\})=n-m,
\]
et on aura
\[
\Phi_{n+m}\simeq\Phi_n\sqcup_e\Phi_m,
\]
i.e. $\Phi_{n+m}$ est la `somme amalgamée' dans (Cat), obtenue en `recollant' $\Phi_n$ et $\Phi_m$, par identification de l'`élément but' de $\Phi_n$ avec l'`élément source' de $\Phi_m$. Le foncteur composition de chemins induit \texttt{[dans ce qui précède, à partir du début de la page 9, il faut inverser} $m$ \texttt{et} $n$\texttt{]}
\[
\uCh_{m,\Sigma}(x,y)\times\uCh_{n,\Sigma}(y,z)\longrightarrow\uCh_{m+n,\Sigma}(x,z),
\]
d'où en passant aux $\pi_0$
\[
\Hom_{\Sigma,m}(x,y)\times\Hom_{\Sigma,n}(y,z)\longrightarrow\Hom_{\Sigma,m+n}(x,z),
\]
où on dénote, pour $a,b\in\mathrm{Ob}\,X$, par
\[
\Hom_{\Sigma,n}(a,b)\overset{\text{déf}}{=}\pi_0(\uCh_{m,\Sigma}(a,b)).
\]

\texttt{[plutôt} $\uCh_{n,\Sigma}$\texttt{]}

\msp{page 10}
Pour en déduire une application sur les $\varinjlim$
\[
\underbrace{\Hom_\Sigma(x,y)}_{=\ \varinjlim\ \Hom_{\Sigma,m}(x,y)}\quad\times\quad\underbrace{\Hom_\Sigma(y,z)}_{=\ \varinjlim\ \Hom_{\Sigma,n}(y,z)}\quad\longrightarrow\quad\underbrace{\Hom_\Sigma(x,z)}_{=\ \varinjlim\ \Hom_{\Sigma,n+m}(x,z)}
\]
il faut prouver, pour $n\leq n'$, $m\leq m'$, la commutativité dans
\[
\begin{tikzcd}
\Hom_{\Sigma,m}(x,y)\times\Hom_{\Sigma,n}(y,z) \arrow[r] \arrow[d] & \Hom_{\Sigma,m+n}(x,y) \arrow[d] \\
\Hom_{\Sigma,m'}(x,y)\times\Hom_{\Sigma,n'}(y,z) \arrow[r] & \Hom_{\Sigma,m'+n'}(x,y)
\end{tikzcd}
\]
\texttt{[plutôt} $\Hom_{\Sigma,m+n}(x,z)$\texttt{,} $\Hom_{\Sigma,m'+n'}(x,z)$\texttt{]} En d'autres termes, que pour $c:x\longrightarrow y$ et $c':y\longrightarrow z$ deux $\Sigma$-chemins composables de longueur $m$ et $n$ respectivement, le composé $c'c:x\longrightarrow z$ ne change pas, `à homotopie près' dans $\uCh_\Sigma(x,z)$, quand on remplace $c$, $c'$ par des chemins $\bar{c}$, $\bar{c}'$ de longueur $m'$, $n'$ déduits de $c$ et $c'$ par `dégénérescence aux bouts'.

Plus précisement, % [sic] précisement -> précisément, accent omis
\[
\bar{c}'\bar{c}\sim\overline{c'c},
\]
le chemin composé est homotope à celui déduit de $c'c$ par `dégénérescence aux bouts'.

Je viens de regarder pour $n=m=1$, $n'=m'=2$, et de constater qu'il n'y a pas de flèche naturelle entre $\bar{c}'\bar{c}$ et $\overline{c'c}$ - l'homotopie cherchée doit
\printedpage{8}
% head: A. Grothendieck : Dérivateurs, Ch. V, Édition des Universités de Montpellier II et Paris VII
% foot: Version 9 septembre 2001
\msp{page 11}
se définir à l'aide d'un `chemin' dans $\uCh_{\Sigma,m',n'}(X;x,z)$. Dans le cas général, on va en fait trouver une homotopie
\[
h=h_{(n,m),(n',m')}
\]
`rendant commutatif' le diagramme de foncteurs
\[
\begin{tikzcd}[column sep=large]
\mathcal{A}=\uCh_{\Sigma,m}(X;x,y)\times\uCh_{\Sigma,n}(X;y,z) \arrow[r, "*_{m,n}"] \arrow[d, "i_{m,m'}\times j_{n,n'}"] & \uCh_{\Sigma,m+n}(X;x,z) \arrow[d, "k_{m+n,m'+n'}", ""'{name=K}] \\
\uCh_{\Sigma,m'}(X;x,y)\times\uCh_{\Sigma,n'}(X;y,z) \arrow[r, "*_{m',n'}", ""{name=S}] & \uCh_{\Sigma,m'+n'}(X;x,z)=\mathcal{B},
\arrow[from=K, to=S, bend right, "h"']
\end{tikzcd}
\]
où les $i_{\alpha,\beta}$, $j_{\alpha,\beta}$, $k_{\alpha,\beta}$ sont les foncteurs de transition
\[
\uCh_{\Sigma,\alpha}(X;a,b)\longrightarrow\uCh_{\Sigma,\beta}(X;a,b)
\]
$(\alpha\leq\beta)$ pour les $(x,y)$, $(y,z)$ et $(x,z)$ (en guise de $(a,b)$) respectivement. L'homotopie $h$ sera un chemin dans la catégorie $\uHom(\mathcal{A},\mathcal{B})$, allant de
\[
k_{m+n,m'+n'}\circ{*}_{m+n}\longrightarrow{*}_{m'+n'}\circ(i_{m,m'}\times j_{n,n'}).
\]
L'existence d'une telle homotopie signifie ni plus ni moins que le diagramme carré ci-dessus est commutatif, en tant que diagramme dans la catégorie (Cat\,hot) des `catégories à homotopie près,' dont les objets sont
\msp{page 12}
les (petites) catégories, et où on définit
\[
\Hom_{\mathrm{Cat\,hot}}(\mathcal{A},\mathcal{B})\overset{\text{déf}}{=}\underbrace{\pi_0(\uHom(\mathcal{A},\mathcal{B}))}_{\begin{array}{@{}l@{}}\text{classes d'homotopie}\\ \text{de foncteurs de }\mathcal{A}\\ \text{dans }\mathcal{B}\end{array}}\,.
\]
C'est cette conclusion qui sera pour nous utile, et non la description explicite du chemin $h$. Cette existence admise (cf. ci-dessous la preuve), on voit donc que si on regarde, pour $a,b\in\mathrm{Ob}\,X$, le système des catégories $\uCh_{\Sigma,r}(X;a,b)$ $(r\in\mathbf{N})$ comme un \textit{ind-objet de} Cat\,hot\textit{,} qu'on désigne par
\[
\uCh^{\mathrm{ind}}_{\Sigma}\,(X;a,b)\in\mathrm{Ob}\,\mathrm{Ind}(\mathrm{Cat\,hot}),
\]
on aura un morphisme canonique dans Ind(Cat\,hot)
\[
\uCh^{\mathrm{ind}}_{\Sigma}\,(X;x,y)\times\uCh^{\mathrm{ind}}_{\Sigma}\,(X;y,z)\xrightarrow{\ *\ }\uCh^{\mathrm{ind}}_{\Sigma}\,(X;x,z),
\]
appelé \textit{`ind-morphisme de composition de chemins.'} On aurait aimé le déduire d'un honnête \textit{foncteur} `composition de chemins'
\[
\begin{aligned}
\uCh_\Sigma(X;x,y)\times\uCh_\Sigma(X;y,z) &\xrightarrow{\ ?\ } \uCh_\Sigma(X;x,z),\\
\text{qu'on noterait }(c,c') &\longmapsto c'\circ c\text{ ou }c'\cdot c,
\end{aligned}
\]
mais je doute qu'un tel foncteur existe,

\dgrblock{dgrb166}{What composition in an ind-category actually asserts}
\begin{dcomment}{status}{p12:On aurait aimé le déduire d'un honnête}{What composition in an ind-category actually asserts}
Grothendieck has reached a useful level of coherence before obtaining a single composition functor on the eventual path categories. For even indices $m,n$, concatenation gives an actual functor at a finite stage. Increasing either index changes its comparison with the larger-stage concatenation by a finite zigzag of natural transformations. This is exactly the relation imposed in $\mathrm{Cat\,hot}$, where
\[
 \operatorname{Hom}_{\mathrm{Cat\,hot}}(A,B)
   =\pi_0\operatorname{Fun}(A,B).
\]
Products descend to this category because one can multiply each natural transformation in a zigzag by an identity transformation.

For ind-objects the assertion is governed by the formula
\[
 \operatorname{Hom}_{\operatorname{Ind}(C)}
       ((A_i),(B_j))
   =\lim_i\mathop{\rm colim}_j\operatorname{Hom}_C(A_i,B_j).
\]
Take $A_{m,n}=\operatorname{Ch}_{\Sigma,m}(x,y)\times
\operatorname{Ch}_{\Sigma,n}(y,z)$ and $B_r=\operatorname{Ch}_{\Sigma,r}(x,z)$, with all indices even. The stagewise concatenation maps land in $B_{m+n}$. Their compatibility as elements of the displayed inverse limit is precisely what the reparametrization lemma proves. The passage to components then gives an ordinary associative composition law.

A functor between the ordinary filtered unions would be additional data. Even if such a functor exists, its restriction to one source stage need not land in a bounded target stage: the source stage can have infinitely many objects. The parenthesis on page~13 explicitly adds the necessary boundedness condition. The equivalence asserted there is therefore conditional on this extra factorization, not a formal property of every functor between filtered colimits.

We do not infer a negative theorem from Grothendieck's doubt. The strict composition demanded on his particular category with symmetric end-padding is not constructed in these pages, and we did not locate a resolution of that exact requirement in the localization sources consulted. Other choices of path presentations do support strict composition. Their existence does not by itself provide the requested functor on this fixed category. The distinction is comparable to the one already emphasized in the commentary on Chapter~I, manuscript p.~7 (``Coherent squares and what the axioms remember''): an equality after passage to a homotopy category does not specify its coherent witness.
\end{dcomment}


% ---- from Der-ChV-234_p011-015.tex
\printedpage{9}
% head: A. Grothendieck : Dérivateurs, Ch. V, Édition des Universités de Montpellier II et Paris VII

\msp{page 13}
qui par passage aux ind-objets redonne le morphisme précédent, ou, ce qui revient au même, tel que pour tout $m,n\in\mathbf{N}$ (de même parité), il existe un $N\geq m+n$ assez grand tel que l'on ait un diagramme
\[
\begin{tikzcd}
\Ch_{\Sigma,m}(X;x,y)\times\Ch_{\Sigma,n}(X;y,z) \arrow[r,"*_{m,n}"] \arrow[dd,hook] \arrow[dr] & \Ch_{\Sigma,m+n}(X;x,z) \arrow[d,hook] \\
 & \Ch_{\Sigma,N}(X;x,z) \arrow[d,hook] \\
\Ch_{\Sigma}(X;x,y)\times\Ch_{\Sigma}(X;y,z) \arrow[r,"?"] & \Ch_{\Sigma}(X;x,z)
\end{tikzcd}
\]
et où le triangle supérieur est \textit{commutatif dans} Cat\,hot. (Il faut supposer que pour $N$ assez grand, le composé
\begin{tikzcd}[cramped, sep=small]
{} \arrow[d,hook] & \\
{} \arrow[r,"?"] & {}
\end{tikzcd}
se factorise par $\Ch_{\Sigma,N}(X;x,z)$ \dots)
Pour voir que les foncteurs composition $*_{m,n}$ définissent un homomorphisme des ind-objets dans (Cat\,hot), il suffit de le prouver dans les cas particuliers $(m',n')=(m+1,n+1)$. On va faire la figure en supposant $m$, $n$ \textit{impairs.}

Je m'aperçois ici que je me suis canulé stupidement, la composition de deux chemins de longueur \textit{impaire $m$, $n$ n'est}

\msp{page 14}
\textit{pas un chemin} de longueur $m+n$, i.e. un diagramme de type $\Phi_n$, mais un diagramme de type \textit{dual} $\Phi_n^o=\Psi_n$. Par exemple, pour $m=n=1$
\[
\underbrace{A_0\longleftarrow A_1\longrightarrow A_2}_{c}=\underbrace{B_0\longleftarrow B_1\longrightarrow B_2}_{c'}
\]
est de type $\Phi_2^o$! Donc pour ne pas trop compliquer la vie, on va dans les développements précédents se borner à $m$, $n$ \textit{pairs,} donc des chemins de longueur au moins 4. C'est un peu ridicule! On est ramené à prouver la commutativité à homotopie près dans le cas où on laisse $c'$ fixé et où on envoie $\Ch_{\Sigma,m}(X;x,y)$ dans $\Ch_{\Sigma,m+2}(X;x,y)$, ou inversement $c$ fixe et $\Ch_{\Sigma,n}(X;y,z)$ dans $\Ch_{\Sigma,n+2}$ \texttt{[lire:}$\Ch_{\Sigma,n+2}(X;y,z)$\texttt{]}. Prenons le premier cas, et notons un chemin $c$ sans indices négatifs (ce n'est pas une notation bien heureuse!)
\[
A_0\longrightarrow A_1\longleftarrow A_2\cdots\longleftarrow A_{2m}
\]
(NB à cause de la condition $n$ pair, tous nos chemins commencent désormais par une flèche de $A_0$ dans $A_1$, jamais en sens inverse).

Je vois que l'explicitation de l'homotopie, sans offrir de difficulté, est très fastidieuse à expliquer. Je vais essayer de donner une démonstration conceptuelle, sans `calcul' fastidieux.

\dgrblock{dgrb167}{Even indices and reparametrization relative to the ends}
\begin{dcomment}{exposition}{p14:Je vais essayer de donner une démonstration conceptuelle}{Even indices and reparametrization relative to the ends}
Grothendieck's correction on this page is part of the proof. The earlier statement that equal parity suffices for concatenation is false for two odd indices. Translation by an odd integer reverses the source--sink parity in $\Phi$, so the proposed embeddings of the two odd intervals into $\Phi_{m+n}$ are not functors with the advertised orientation. Restricting to even indices repairs the construction and loses no path, since every finite index is bounded by an even one.

For even $m,n$ use the embeddings
\[
 \Phi_m\longrightarrow\Phi_{m+n},\quad i\longmapsto i-n,
 \qquad
 \Phi_n\longrightarrow\Phi_{m+n},\quad j\longmapsto j+m.
\]
They identify the final vertex of the first interval with the initial vertex of the second. At the join both incident arrows have the same source--sink type, so no new composite of two nonidentity arrows appears. Consequently the resulting pushout in categories is $\Phi_{m+n}$. Restricting a functor on that pushout is the same as giving the two composable diagrams with their common object literally identified. Natural transformations with identity components at the exterior endpoints concatenate as well.

To compare padding before concatenation with padding afterwards, regard the two resulting diagrams as precompositions with maps from a larger interval to a smaller one. Both maps are nondecreasing in the ordinary numerical order of the vertices and fix the two endpoint labels. Their categorical orders are still the alternating orders. This second monotonicity condition, explicitly restored in the edition's footnote on page~17, prevents a backwards source edge from travelling forwards in the smaller interval.

The boundary equation is $h\alpha'=\alpha$, with $\alpha'$ the endpoint map into the larger interval and $\alpha$ the endpoint map into the smaller one. The corrected Hom expression on page~16 fixes the order of those intervals, but its description of the restriction as $\alpha'$ still has the wrong codomain; the displayed equation is the precise interpretation.

The intermediate maps in a homotopy must have that same property. If $c$ is a $\Sigma$-path and $h$ is a numerically nondecreasing functor between alternating intervals, each backwards edge is sent either to a backwards edge or to an identity. Hence $ch$ remains a $\Sigma$-path. Every natural transformation between such $h$'s induces a ladder with identity endpoint components. The forthcoming connectedness argument therefore proves compatibility inside the required path category, with the denominator condition intact.

The reduction to simultaneous increments $(m,n)\mapsto(m+1,n+1)$ before the correction is insufficient for arbitrary independent changes of the two indices. After restriction to even indices, separate increments by two in the first or second variable generate all transitions. Grothendieck's revised formulation on this page makes exactly this repair.
\end{dcomment}


% foot: Version 9 septembre 2001

\printedpage{10}
% head: A. Grothendieck : Dérivateurs, Ch. V, Édition des Universités de Montpellier II et Paris VII

\msp{page 15}
Les deux foncteurs composés
\[
f=k\circ *_{m,n}\ ,\quad g=*_{m',n'}\circ(i\times j)
\]
dans le carré page 11, dont il faut prouver qu'ils sont homotopes dans $\uHom(\mathcal{A},\mathcal{B})$, peuvent être vus comme obtenus de la façon suivante, en termes de deux foncteurs
\[
\varphi,\psi:\Phi_{m'+n'}\rightrightarrows\Phi_{m+n}, \tag*{(1)}
\]
d'où deux foncteurs `transposés'
\[
\begin{tikzcd}
\uHom(\Phi_{m,n},X) \arrow[r,shift left,"F"] \arrow[r,shift right,"G"'] & \uHom(\Phi_{m',n'},X)
\end{tikzcd} \tag*{(2)}
\]
\texttt{[plutôt} $\Phi_{m+n}$\texttt{,} $\Phi_{m'+n'}$\texttt{]} lesquels foncteurs `induisent' $f$ et $g$. De façon précise, désignant par $a$, $b$ les objets source et but de $\Phi_{m+n}$, par $c$ l'objet qui correspond au but du chemin $c$ de longueur $m$ \texttt{[trop de} $c$\texttt{'s]}, à la source du chemin $c'$ de longueur $n$ qu'il s'agit de composer, enfin par $a_0$, $b_0$ les éléments source et but de $\Phi_{m'+n'}$, les foncteurs envisagés (1) appliquent $a_0$ dans $a$, $b_0$ dans $b$, donc définissent un diagramme commutatif
\[
\begin{tikzcd}
\Phi_{m'+n'} \arrow[r,shift left,"\varphi"] \arrow[r,shift right,"\psi"'] & \Phi_{m+n} \\
\underset{\text{\footnotesize catégorie discrète}}{\varepsilon=\{a_0,b_0\}.} \arrow[u,hook,"\alpha'"] \arrow[ur,hook,"\alpha"] &
\end{tikzcd} \tag*{(3)}
\]
Ceci dit, la catégorie but $\mathcal{B}$ pour $f$, $g$, à savoir $\Ch_{\Sigma,m'+n'}(X;x,z)$, s'interprète comme la sous-catégorie pleine de $\uHom(\Phi_{m+n},X)$ formée

\msp{page 16}
des foncteurs $c':\Phi_{m'+n'}\longrightarrow X$ dont le composé avec $\alpha'$ est le foncteur
\[
\beta:\varepsilon\longrightarrow X\qquad\qquad
\begin{aligned}
a_0&\longrightarrow x\\
b_0&\longrightarrow z,
\end{aligned} \tag*{$(*)$}
\]
et de même la catégorie source $\mathcal{A}$ s'interprète comme la sous-catégorie pleine de $\uHom(\Phi_{m+n},X)$ formée des $c:\Phi_{m+n}\longrightarrow X$ dont la restriction à $\{a,b,c\}$ est le foncteur
\[
a\longrightarrow x,\ b\longrightarrow y,\ c\longrightarrow z.
\]
\texttt{[plutôt} $b\longrightarrow z$\texttt{,} $c\longrightarrow y$\texttt{]} Elle est contenue dans la sous-catégorie pleine $\mathcal{A}'$ des foncteurs dont la restriction à $\{a,b\}$ soit $a\longrightarrow x,\ b\longrightarrow z$, i.e. dont le composé avec $\alpha$ soit le foncteur $\beta:\varepsilon\longrightarrow X$ déjà envisagé ci-dessus. Ainsi les foncteurs $f,g:\mathcal{A}\rightrightarrows\mathcal{B}$ se factorisent par $\mathcal{A}'$
% foot: Version 9 septembre 2001
\printedpage{11}
% head: A. Grothendieck : Dérivateurs, Ch. V, Édition des Universités de Montpellier II et Paris VII
\[
\begin{tikzcd}
\mathcal{A} \arrow[r,hook] \arrow[dr,shift left,"f"] \arrow[dr,shift right,"g"'] & \mathcal{A}' \arrow[d,shift right,"g'"'] \arrow[d,shift left,"f'"] \\
 & \mathcal{B},
\end{tikzcd}
\]
foncteurs déduits de $\varphi$, $\psi$ ci-dessus, donnant lieu aux triangles commutatifs (3), et du foncteur $\beta:\varepsilon\longrightarrow X$, de la façon dite. Pour prouver que $f$ et $g$ sont homotopes, il suffit de le prouver pour $f'$, $g'$. Et pour ceci, il suffit de montrer que $\varphi$ et $\psi$ sont homotopes en tant qu'objets de la catégorie $\uHom(\Phi_{m+n},\Phi_{m'+n'};\alpha)$ formée des foncteurs dont la restriction à $\varepsilon$ est $\alpha$ \texttt{[plutôt de la catégorie} $\uHom(\Phi_{m'+n'},\Phi_{m+n};\alpha')$ \texttt{formée des foncteurs dont la restriction à} $\varepsilon$ \texttt{est} $\alpha'$ \texttt{]}.

\msp{page 17}
D'ailleurs, les deux assertions (sur $f'$, $g'$, quels que soient $X$, $x$, $y$, $z$, \texttt{[plutôt} $x$\texttt{,} $z$\texttt{]} et sur $\varphi$, $\psi$, $\alpha$ \texttt{[plutôt} $\alpha'$\texttt{]}) sont en fait équivalentes comme on voit en faisant $X=\Phi_{m+n}$, $x=a$, $y=c$, $z=b$ \texttt{[plutôt} $x=a$\texttt{,} $z=b$\texttt{]}. Ceci sera contenu dans le lemme plus général. (\footnote[4]{Dans les considérations qui suivent, on considère en fait tacitement que les foncteurs envisagés sont \textit{croissants} sur les sommets pour l'ordre de leur numérotation. Donc il faudrait écrire $\uHom_{cr}$ au lieu de $\uHom$, dans les énoncés qui suivent, y compris dans le lemme de connexité ci-contre.})

\textbf{Lemme.} \textit{Soient $k$ et $l$ deux entiers. Alors la sous-catégorie pleine de $\uHom(\Phi_k,\Phi_l)$ formée des foncteurs $\Phi_k\longrightarrow\Phi_l$ appliquant source en source, but en but, est connexe.} (NB Elle est non vide si et seulement si $k\geq l$.)

On va décrire entièrement ces foncteurs (= applications croissantes) $\Phi_k\xrightarrow{\varphi}\Phi_l$, en identifiant $\Fl(\Phi_k)$ \texttt{[plutôt l'ensemble des flèches non identiques de} $\Phi_k$\texttt{]} à l'intervalle d'entiers $[1,2k]$, et en associant à $\varphi$ la partie
\[
E_\varphi\subset\Fl(\Phi_k)=[1,2k]\qquad\qquad E_\varphi=\{u\in\Fl(\Phi_k)\mid\varphi(u)\text{ une identité}\}
\]
formée des $u\in\Fl(\Phi_k)$ qui sont `contractées' par $\varphi$ en une identité. Si $E$ est un ensemble d'entiers, on y introduit la notion de composantes connexes (par la relation d'équivalence dans $E$ engendrée par la relation de succession $y=x+1$). On voit aussitôt que les composantes connnexes de $E_\varphi$ qui ne contiennent ni la première ni la dernière flèche de $\Phi_k$ sont paires, et que % [sic] connnexes -> connexes, triple n
\[
\underbrace{\card\Fl(\Phi_k)}_{2k}=\card\,(E_\varphi)+\underbrace{\card\Fl(\Phi_l)}_{2l},
\]

\dgrblock{dgrb168}{A contraction is determined by the edges it erases}
\begin{dcomment}{exposition}{p17:On va décrire entièrement ces foncteurs}{A contraction is determined by the edges it erases}
Grothendieck reduces a diagrammatic homotopy problem to finite combinatorics. Number the vertices of $\Phi_k$ from $0$ to $2k$, preserving their numerical order; the vertex numbered $t$ is the original vertex $t-k$. Number the edges from $1$ to $2k$. For a numerically nondecreasing endpoint-preserving functor $f:\Phi_k\to\Phi_l$, adjacent values differ by zero or one: the target has no nonidentity arrow joining nonadjacent vertices. Every noncollapsed edge advances the numerical position once. Therefore
\[
 f_E(t)=-l+t-\bigl|E\cap[1,t]\bigr|,
 \qquad 0\leq t\leq2k,
\]
where the values are written in the original target numbering and $E$ is the set of collapsed edges. This formula both proves uniqueness and makes the prefix and suffix observations on page~20 transparent.

When $l>0$, the admissibility conditions have a direct explanation. Exactly $2l$ edges survive, so $|E|=2(k-l)$. Between successive surviving edges, an even number of erased edges is necessary and sufficient for their alternating orientations to agree. The same condition holds at an endpoint when the initial orientations agree. If $k$ and $l$ have opposite parity, the initial and terminal erased runs must both be odd. The parity of the terminal run follows from the total cardinality together with the parity of all interior runs; the edition rightly inserts both conditions (a) and (b) at that point.

Conversely, these conditions make each consecutive pair $f_E(t-1),f_E(t)$ either equal or the endpoints of the correctly oriented target edge. They thus define a functor. It is surjective on objects and nonidentity arrows, because it starts at $-l$, ends at $l$, and advances by steps of at most one. Its only identifications are the consecutive strings of vertices joined by edges of $E$, so it is the stated quotient. This establishes necessity and sufficiency, not merely the numerical count.

For $l=0$ there is one constant functor and $E=[1,2k]$. No initial target direction exists, so the parity condition must be omitted in that case. For $k=l=0$ the same convention gives $E=\varnothing$. Endpoint-preserving functors exist exactly when $k\geq l$: necessity follows from counting surviving edges, while sufficiency follows from the clamping functor $\Phi_k\to\Phi_l$.

The footnote's restriction to numerical monotonicity is essential for the description by $E$. A general functor can retrace a target edge, and its noncollapsed edges need not number $2l$. Later we prove contractibility also for the full endpoint fibre, but that stronger statement does not make the intervening classification of contractions valid without its hypothesis.

For instance, number the source vertices of $\Phi_3$ from $0$ to $6$, and retain the target numbering $-1,0,1$ for $\Phi_1$. The row
\[
 (-1,0,-1,0,-1,0,1)
\]
defines an endpoint-preserving functor. At the odd-numbered source vertices, which are sources in this numbering, its value is $0$; every generating arrow is therefore sent to one of the two nonidentity arrows leaving $0$ in $\Phi_1$. No edge is erased. Nevertheless the difference in the indices is $3-1=2$, so the proposed cardinality would be four. This example isolates the precise role of the footnote: endpoint preservation is sufficient for the later contractibility theorem, but it is insufficient for this particular description by a set of contractions.
\end{dcomment}


\msp{page 18}
i.e.
\[
\card\,(E_\varphi)=2(k-l).
\]

% foot: Version 9 septembre 2001

\printedpage{12}
% head: A. Grothendieck : Dérivateurs, Ch. V, Édition des Universités de Montpellier II et Paris VII

Il est immédiat qu'un foncteur $\varphi$ est surjectif sur les objets et sur les flèches, donc identifie $\Phi_l$ à une catégorie quotient de $\Phi_k$ - et la relation d'équivalence dans $\Phi_k$ définie par $\varphi$ est celle obtenue en identifiant tous les objets qui sont source ou but de flèches d'une même composante connexe $I$ de $E_\varphi$
\[
\cdots\longleftarrow\ \underbrace{\longrightarrow\ \longleftarrow\ \longrightarrow\ \longleftarrow}_{I}\ \longrightarrow\cdots
\]
et les flèches dans cette composante à l'identité correspondante. Les conditions énoncées sur $E$ (indépendamment d'un $\varphi$ qui lui donnerait naissance assurent que la catégorie quotient est isomorphe (avec \textit{préservation} de source et de but) soit à $\Phi_l$, soit à la duale $\Psi_l=\Phi_l^o$, suivant la direction de la première flèche de $\bar{\Phi}_k$ (quotient de $\Phi_k$). Cette `direction initiale' est la même que celle pour $\Phi_k$ si la première flèche de $\Phi_k$ n'appartient pas à $E$, ou bien si elle y appartient et sa composante connexe dans $E$ est de cardinal pair. Dans le cas inverse ($u\in E$ et sa composante connnexe dans $E$ est de cardinal impair), la direction initiale de $\bar{\Phi}_k$ est opposée à celle de $\Phi_k$. % [sic] "(indépendamment d'un $\varphi$ qui lui donnerait naissance" -> parenthèse fermante manquante, the opening parenthesis is never closed; [sic] connnexe -> connexe, triple n

\msp{page 19}
On voit donc que $E$ doit satisfaire encore à une condition supplémentaire, pour que le $\bar{\Phi}_k$ correspondant soit isomorphe à $\Phi_l$, i.e. pour que $E$ provienne d'un foncteur $\varphi$. Je résume toutes les conditions

\begin{itemize}
\item[a)] $\card E=2(k-l)$.
\item[b)] Les composantes connexes de $E$ qui ne contiennent pas de flèche initiale ou finale sont de cardinal pair.
\item[c)] Si $\Phi_k$ et $\Phi_l$ ont même direction initiale ($k$ et $l$ de même parité), alors l'éventuelle composante de $E$ contenant la flèche initiale doit être elle aussi de cardinal pair (ce qui, à cause de a) \texttt{[plutôt a) et b)]}, équivaut à la même condition pour l'éventuelle composante connexe contenant la flèche finale de $\Phi_k$). Si $\Phi_k$ et $\Phi_l$ ont des directions initiales différentes, il faut que la flèche initiale de $\Phi_k$ appartienne à une composante connexe de cardinal \textit{impair} de $E$ (ce qui équivaut à itou pour la flèche finale).
\end{itemize}

D'autre part, le foncteur $\varphi$ est uniquement déterminé par la connaissance de $E$, car il y a un \textit{unique} isomorphisme $\bar{\Phi}_k\xrightarrow{\sim}\Phi_l$ compatible avec sources et buts, à cause du

\textbf{Lemme} (immédiat) : \textit{Un automorphisme de $\Phi_l$ qui est l'identité sur l'objet source (ou sur l'objet but) est l'identité partout.}

Les foncteurs $\varphi$, $\psi$ de tantôt correspondaient au choix, l'un de $E_\varphi=\{1,2,k-1,k\}$ \texttt{[plutôt} $E_\varphi=\{1,2,2k-1,2k\}$\texttt{]}, l'autre de $E_\psi=\{1,2,i,i+1\}$, 2 composantes connexes

\msp{page 20}
dans un cas, trois \texttt{[plutôt deux]} dans l'autre.) % [sic] "l'autre.)" -> parenthèse fermante sans parenthèse ouvrante, no matching opening parenthesis

% foot: Version 9 septembre 2001

\printedpage{13}
% head: A. Grothendieck : Dérivateurs, Ch. V, Édition des Universités de Montpellier II et Paris VII

Dans le cas général, il y a lieu d'expliciter la relation d'ordre de l'ensemble $\uHom^{!}(\Phi_k,\Phi_l)$ (! signifie : compatible avec source et but), en termes des ensembles représentatifs $E_\varphi$, $E_\psi$. Je ne vais pas expliciter cette relation d'ordre, mais seulement une propriété de celle-ci qui nous sera suffisante. Soit $\varphi_E$ le foncteur défini par une partie $E\subset[1,2k]$. Tout d'abord :

\textbf{Lemme 1} (immédiat) : \textit{Soient $E,E'\subset[1,2k]$ admissibles, et soit $j\in[1,2k]$ tel que $E\cap[1,j]=E'\cap[1,j]$. Alors $\varphi_E$ et $\varphi_{E'}$ coïncident sur le but de $u_j$ ($j$ une flèche de $\Phi_k$) et sur tous les objets de $\Phi_k$ antérieurs à celle-ci (au sens de la numérotation des objets, et non de la relation d'ordre de la catégorie $\Phi_k$ bien sûr.) Dualement, si $E\cap[j,2k]=E'\cap[j,2k]$, alors $\varphi_E$ et $\varphi_{E'}$ coïncident sur la source de $y$, et sur tous les objets de $\Phi_k$ postérieurs à celui-ci.} % [sic] la source de $y$ -> la source de $u_j$, apparently intended by duality with "le but de $u_j$"

Considérons maintenant une partie connexe $I$ de $E$ (pas forcément une composante connexe), $I=[r,s]\subset[1,2k]$ ($r$ le plus petit élément de $I$, $s$ son plus grand élément) et supposons que

\begin{itemize}
\item[a)] $\card I$ est pair (i.e. $s-r$ est pair) \texttt{[plutôt} $s-r$ \texttt{impair]}
\item[b)] $r\neq1$, et $r-1\notin E$ (i.e. $r-1\in[1,2k]\backslash E$)
\end{itemize}

et soit
\[
E'=(E\backslash\{s\})\cup\{r-1\}=E\backslash I\cup\overbrace{[r-1,s-1]}^{I'}
\]
déduit de $I$ en `translatant' sur la gauche d'un cran.

\msp{page 21}
Il est immédiat que si $E$ est admissible, $E'$ aussi (et inversement). Ceci dit, on a le

\textbf{Lemme 2.} \textit{$E$ et $E'$ étant comme ci-dessus, il existe une homomorphisme de foncteurs entre $\varphi_E$ et $\varphi_{E'}$} (je ne dis pas dans quel sens va la flèche!). % [sic] une homomorphisme -> un homomorphisme, wrong gender

\textbf{Corollaire.} Énoncé dual, quand $E'$ est déduit de $E$ en translatant à droite d'un cran un intervalle connexe de longueur paire $I=[r,s]$ de $E$, $s\neq2k$ et $s+1\notin E$.

En effet, dans ce cas $E$ est déduit de $E'$ par une opération du type explicité dans le lemme 2.

La connexité de la catégorie $\uHom^{!}(\Phi_k,\Phi_l)$ est maintenant claire. Si $k$ et $l$ ont même parité, en poussant les composantes connexes de $E$ vers la gauche, on écrit $E$ sous la forme `normale'
\[
E_0=[1,2(k-l)].
\]
Si $k$ et $l$ sont de parité inverse, on ne peut pousser toute la composante connexe de $E$ qui contient l'élément dernier $2k$, car elle est de cardinal impair. Il faut laisser $2k$ dans $E$ et pousser le reste, et on ramène $E$ à la forme normale
\[
E_0=[1,2(k-l)-1]\cup\{2k\}.
\]

\dgrblock{dgrb169}{Sliding an even block gives a natural transformation}
\begin{dcomment}{exposition}{p21:La connexité de la catégorie}{Sliding an even block gives a natural transformation}
Grothendieck's elementary move explains why different padding conventions become equivalent. The formula just obtained gives a short proof of the move, including the natural transformation he leaves implicit. Suppose $I=[r,s]\subset E$ has even cardinality, $r>1$, and $r-1\notin E$, and put $E'=(E\setminus\{s\})\cup\{r-1\}$. Then
\[
 |E'\cap[1,t]|-|E\cap[1,t]|=
 \begin{cases}1,&r-1\leq t\leq s-1,\\0,&\text{otherwise}.
 \end{cases}
\]
The map $f_E$ is constant, say equal to $a$, on the vertices $r-1,\ldots,s$, because all intervening edges belong to $E$. Thus $f_{E'}$ differs from $f_E$ only by replacing $a$ with $a-1$ at the vertices $r-1,\ldots,s-1$. These two target vertices are joined by one arrow. Its orientation determines a natural transformation in one direction between the two functors; outside the displayed interval its components are identities. Naturality is automatic in a poset once the component inequalities hold.

The endpoints are unchanged, since $r-1>0$ and $s-1<2k$. Admissibility is preserved: removing an even initial block from a collapsed run preserves the parity of its remainder, and translating that even block preserves the parity of every run it meets. The initial and terminal parity conditions are therefore unchanged. This also proves the right-moving version by reversing the move.

Repeated left-moving operations terminate: the sum of the indices in $E$ strictly decreases at each move. If the endpoint orientations agree, every collapsed run is even and all runs can be joined to the leftmost one, yielding $[1,2(k-l)]$. If they disagree, the initial and terminal runs are odd. Keep the final edge $2k$ in place and move the remaining even part of the final run, together with the interior runs, to the left. The normal form is
\[
 [1,2(k-l)-1]\cup\{2k\}.
\]
Each move is a morphism in one of the two directions, so every object is connected to the same normal form. This is a proof of connectedness, which is all the composition argument needs.

For example, with $k=4,l=2$, the set $E=\{1,2,7,8\}$ successively becomes $\{1,2,6,7\}$, $\{1,2,5,6\}$, $\{1,2,4,5\}$, and $\{1,2,3,4\}$. The associated functors form a zigzag of four natural transformations relative to the endpoints. This gives a visible instance of the homotopy comparing padding at an exterior end with padding near a join. The proof does not assert that there is always one natural transformation between the original two functors.
\end{dcomment}


\hfill cqfd!

% foot: Version 9 septembre 2001

% ---- from Der-ChV-234_p016-020.tex
\printedpage{14}
% head: A. Grothendieck : Dérivateurs, Ch. V, Édition des Universités de Montpellier II et Paris VII
Ce qui précède justifie notre raisonnement, qui donnait par composition des chemins des applications de composition bien définies
\[
\underbrace{\Hom_\Sigma(x,y)}_{=\pi_0\uCh_\Sigma(x,y)} \times \underbrace{\Hom_\Sigma(y,z)}_{=\pi_0\uCh_\Sigma(y,z)} \longrightarrow \underbrace{\Hom_\Sigma(x,z)}_{=\pi_0\underline{\textstyle\mathrm{Ch}}_\Sigma(x,z)}.
\]
% (the third "Ch" under the brace is printed in a larger size than the two others)

\msp{page 22}Je dis que cette composition est associative. Ceci résulte du fait que si on a trois chemins (d'indices \textit{pairs}!)
\[
x \xrightarrow{c} y \xrightarrow{c'} z \xrightarrow{c''} t,
\]
on a (sans avoir à effectuer des homotopies \dots) associativité stricte
\[
c''(c'c) = (c''c')c.
\]

Je dis qu'il y a des identités. On définit
\[
1_x \in \Hom_\Sigma(x,x)
\]
par le chemin d'indice 0 (qui est pair!), i.e. par le diagramme réduit à $x$, sans flèche (ou si on veut interpréter le diagramme comme un foncteur $e \longrightarrow X$, il y a la flèche identique sous-entendue). Nos raisonnements n'excluaient pas le cas $n = 0$ ou $m = 0$, et il est tautologique que pour un chemin
\[
x \xrightarrow{c} y
\]
on a
\[
c \cdot 1_x = 1_y \cdot c = c.
\]
Donc on a bien défini une catégorie, que je vais appeler $X_\Sigma$. Ja vais définir un foncteur canonique % [sic] Ja -> Je, misprint of the pronoun
\[
\theta : X \longrightarrow X_\Sigma
\]
qui sera l'identité sur les objets, il faut le définir sur les flèches
\[
\Hom_X(x,y) \longrightarrow \Hom_{X_\Sigma}(x,y) = \pi_0\mathrm{Ch}_\Sigma(x,y)
\]
\texttt{[plutôt} $\pi_0\uCh_\Sigma(x,y)$\texttt{]}. Malheureusement, une flèche
\[
x \xrightarrow{u} y
\]
\msp{page 23}ne peut s'interpréter directement comme un chemin à notre sens, il faut payer nos fuites! Mais on va à une telle flèche associer le chemin
\[
x \longrightarrow y \xleftarrow{\mathrm{id}} y \xrightarrow{\mathrm{id}} y \xleftarrow{\mathrm{id}} y,
\]
% foot: Version 9 septembre 2001
\printedpage{15}
% head: A. Grothendieck : Dérivateurs, Ch. V, Édition des Universités de Montpellier II et Paris VII
élément de $\uCh_{\Sigma,2}(x,y)$, ce qui définit l'application cherchée. Il faut vérifier la compatibilité avec les identités et sous la composition. À la flèche $1_x : x \longrightarrow x$ est associé l'élément de $\Phi_2(x,x)$
\[
x \xrightarrow{\mathrm{id}} x \xleftarrow{\mathrm{id}} x \xrightarrow{\mathrm{id}} x \xleftarrow{\mathrm{id}} x,
\]
il faut prouver que dans $\pi_0(\uCh_{\Sigma,2}(x,x))$ il devient égal à l'élément déduit de $1_x \in \Phi_{\Sigma,0}(x,x)$, par le foncteur de transition. Mais cet élément n'est autre justement que
\[
x \xrightarrow{\mathrm{id}} x \xleftarrow{\mathrm{id}} x \xrightarrow{\mathrm{id}} x \xleftarrow{\mathrm{id}} x,
\]
c'est le même (sans avoir même à passer au $\pi_0$)!

\dgrblock{dgrb170}{A denominator set need not be closed under composition}
\begin{dcomment}{examples}{p22:Je dis qu'il y a des identités.}{A denominator set need not be closed under composition}
Grothendieck's use of arbitrarily long paths is already visible in the simplest one-object example. Let $X$ have one object and endomorphism monoid $(\mathbf N,+)$, writing its generator multiplicatively as $t$. Take
\[
 \Sigma=\{1,t\}.
\]
This contains the identity but is not closed under composition: $t^2\notin\Sigma$. A backwards edge may carry $t$, but not $t^2$. For a path $c$, let $a(c)$ be the sum of the exponents on its forward edges and let $b(c)$ count the backwards edges labelled $t$. Its evaluation is determined by the integer
\[
 \nu(c)=a(c)-b(c).
\]
A commutative ladder with identity endpoints preserves this integer. Indeed, write each vertical component as $t^{h_i}$ with $h_0=h_{2n}=0$. For a forward edge of exponent $a_i$, naturality gives $h_i+a_i=a_i'+h_{i-1}$; for a backwards edge the analogous equality has the horizontal exponents reversed. Summing the signed equalities cancels all the $h_i$, leaving $\nu(c)=\nu(c')$.

Conversely, every integer is represented. A nonnegative integer $r$ is represented by the ordinary arrow $t^r$ with identity padding. For $-r<0$, concatenate $r$ backwards copies of $t$, separating them by identities, then pad to an even index. The cancellation ladders below show that a forward $t$ followed by a backwards $t$ represents the identity, in either order. Splitting a forward $t^a$ into $a$ successive copies of $t$ is allowed by the functoriality calculation for ordinary arrows. These two operations reduce any path to the representative of its signed exponent. Thus
\[
 \pi_0\operatorname{Ch}_\Sigma(X;*,*)\cong\mathbf Z,
\]
with concatenation corresponding to addition.

The resulting localization is the one-object category of the group $\mathbf Z$. Its universal property says exactly that a chosen endomorphism $T$ becomes invertible: the image of the integer $r$ must be $T^r$. Negative integers require arbitrarily many backwards edges when the denominator set is the original $\{1,t\}$. Replacing it by all powers of $t$ permits a single backwards edge to represent them, but this changes the categories of presentations. The agreement after passage to connected components is the invariant conclusion.
\end{dcomment}


Pour la compatibilité avec la composition
\[
x \xrightarrow{u} y \xrightarrow{v} z,
\]
on a d'une part
\[
x \xrightarrow{vu} z \xleftarrow{\mathrm{id}} z \xrightarrow{\mathrm{id}} z \xleftarrow{\mathrm{id}} z,
\]
d'autre part
\[
x \xrightarrow{u} y \xleftarrow{\mathrm{id}} y \xrightarrow{\mathrm{id}} y \xleftarrow{\mathrm{id}} y \xrightarrow{v} z \xleftarrow{\mathrm{id}} z \xrightarrow{\mathrm{id}} z \xleftarrow{\mathrm{id}} z.
\]
On complète le premier chemin par dégénérescence aux deux bouts, pour le comparer au deuxième en tant qu'élément de $\uCh_4(x,z)$,
\[
x \xrightarrow{\mathrm{id}} x \xleftarrow{\mathrm{id}} x \xrightarrow{vu} z \xleftarrow{\mathrm{id}} z \xrightarrow{\mathrm{id}} z \xleftarrow{\mathrm{id}} z \xrightarrow{\mathrm{id}} z \xleftarrow{\mathrm{id}} z.
\]

\dgrblock{dgrb171}{Two ordinary arrows and one comparison diagram}
\begin{dcomment}{examples}{p23:Pour la compatibilité avec la composition}{Two ordinary arrows and one comparison diagram}
Grothendieck needs to show that sending an ordinary arrow to its padded path respects composition. This particular comparison is smaller than the general interval problem and can be completed explicitly. Let $x\xrightarrow{u}y\xrightarrow{v}z$, and regard it as a functor $T:[2]\to X$ with $T(0)=x$, $T(1)=y$, and $T(2)=z$.

Number the nine vertices of $\Phi_4$ from $0$ to $8$. The concatenated path for $u$ and $v$, and the padded path for $vu$, arise from the following functors to the ordered category $[2]$:
\[
\begin{array}{c|rrrrrrrrr}
 i&0&1&2&3&4&5&6&7&8\\ \hline
 f(i)&0&1&1&1&1&2&2&2&2\\
 g(i)&0&0&0&2&2&2&2&2&2\\
 h(i)&0&0&0&1&1&2&2&2&2.
\end{array}
\]
Here $h(i)=\min(f(i),g(i))$. Every row is a functor: along each forward edge its values are nondecreasing, while the values on the two ends of every backwards edge are equal. Pointwise inequalities give natural transformations
\[
 f\longleftarrow h\longrightarrow g.
\]
Their endpoint components are identities, because all three rows take the values $0$ and $2$ at the two ends. Applying $T$ produces the desired two-step ladder in $\operatorname{Ch}_{\Sigma,4}(X;x,z)$. No nonidentity arrow is placed on a backwards edge, so the proof works for any $\Sigma$ containing the identities.

This calculation establishes $\theta(v)\theta(u)=\theta(vu)$ after passage to components. For identities the two diagrams already agree at the finite stage, as Grothendieck observes. Hence $\theta$ is a functor without assuming the general lemma. The display $\Delta_1\sqcup\Delta_1\simeq\Delta_2$ on the following manuscript page must mean amalgamation over the common object, $[1]\sqcup_{[0]}[1]\cong[2]$; the ordinary disjoint union has four objects.

The table also records why an unqualified homotopy in a functor category would be insufficient. An arbitrary map $\Phi_4\to[2]$ may send a backwards edge to a nonidentity arrow. Its image under $T$ would then require $u$, $v$, or $vu$ to lie in $\Sigma$, which has not been assumed. Our three rows avoid that problem by making every backwards edge constant.
\end{dcomment}


\msp{page 24}Il faut prouver que les deux sont homotopes. On a donc deux foncteurs
\[
f,g : \underbrace{\overbrace{\uFl(X;x,y)}^{\text{catégorie discrète}} \times \overbrace{\uFl(X;y,z)}^{\text{catégorie discrète}}}_{\mathcal{A}} \underset{g}{\overset{f}{\rightrightarrows}} \underbrace{\uCh_{4,\Sigma}(X;x,z)}_{\mathcal{B}}, \tag*{$(\ast)$}
\]
on veut prouver qu'ils sont homotopes. Pour ceci, on les regarde (comme tantôt) comme `transposés' des deux foncteurs
\[
\begin{tikzcd}[column sep=large]
\bDelta_1 \sqcup \bDelta_1 \simeq \bDelta_2 & \Phi_4 \arrow[l, shift right, "\varphi"'] \arrow[l, shift left, "\psi"] \\
& \varepsilon, \arrow[ul, "\alpha"] \arrow[u, "\alpha'"']
\end{tikzcd}
\]
vu que les catégories $\mathcal{A}$, $\mathcal{B}$ des deux membres de (*) s'interprètent comme des sous-catégories pleines de $\uHom(\bDelta_2, X)$ et $\uHom(\Phi_4, X)$ respectivement; à savoir les sous-catégories pleines formées des foncteurs dont le composé avec $\alpha$ resp. $\alpha'$ est donné par $x$ et par $z$, et de plus (pour $\mathcal{A}$) dont la valeur en l'objet médian de $\bDelta_2$ soit $y$. On peut encore remplacer $\mathcal{A}$ par la sous-catégorie pleine plus grande $\mathcal{A}'$ de $\uHom(\bDelta_2, X)$, obtenue en ignorant cette dernière condition, et il suffit de prouver que les deux foncteurs `plus gros'
% foot: Version 9 septembre 2001
\printedpage{16}
% head: A. Grothendieck : Dérivateurs, Ch. V, Édition des Universités de Montpellier II et Paris VII
\[
\mathcal{A}' \underset{g'}{\overset{f'}{\rightrightarrows}} \mathcal{B}'
\]
sont homotopes. Et pour ceci, il suffit de
\msp{page 25}voir encore que $\varphi$ et $\psi$ sont homotopes dans la catégories des foncteurs $\Phi_4 \longrightarrow \bDelta_2$ qui sont compatibles avec la source, et le but. On voit donc que notre lemme d'homotopie de tantôt n'était pas assez général. Il faudrait définir les \textit{catégories intervalles généralisés,} de longueur $n$ quelconque ($n =$ nombre de flèches bout à bout), une catégorie intervalle généralisé \textit{standard} de longueur $n$ a pour ensemble d'objets l'intervalle d'entiers % [sic] la catégories -> la catégorie, singular article with plural noun
\[
[0,n],
\]
et pour tout couple d'entiers consécutifs $i-1$, $i$ parmi les objets ($1 \leq i \leq n$), il y a exactement une flèche qui relie $i$, $i+1$ \texttt{[plutôt} $i-1$\texttt{,} $i$\texttt{]}. On affecte l'intervalle $[i,i+1]$ du signe $+$ si $i \longrightarrow i+1$ (flèche en sens direct), du signe $i$ \texttt{[plutôt} $-$\texttt{]} si $i \longleftarrow i+1$ (flèche en sens rétrograde) - donc les catégories standards sont définies, pour $n$ fixé, par un système de signes $+$ ou $-$, dont on affecte les $n$ `tronçons' consécutifs $A_i := [i-1,i]$ ($1 \leq i \leq n$) de l'intervalle d'entiers $[0,n]$. Il y en a donc exactement $2^n$.

On s'attend à ce que l'énoncé suivant soit vrai :

\dgrblock{dgrb172}{Contracting relative functor categories into an interval}
\begin{dcomment}{development}{p25:On s'attend à ce que l'énoncé suivant soit vrai}{Contracting relative functor categories into an interval}
Grothendieck's question asks whether the choices used to change a path's parametrization form a contractible category. We can prove a slightly more flexible statement that also accommodates the cancellation diagrams on page~27, where both endpoints have the same image. The proof uses the numerical order of an interval only to define minimum and maximum; the categorical order retains the prescribed directions of all its edges.

\textbf{Relative interval contraction.}\ours\ Let $J$ be the free category on an oriented path with vertices $0,\ldots,r$. Let $j:B\to A$ be a functor and $b:B\to J$ a fixed diagram. The strict fibre
\[
 \mathcal F_b=
 \{f:A\to J\mid fj=b\}
\]
of restriction, with transformations restricting to $1_b$, is either empty or categorically contractible: if $g\in\mathcal F_b$, its identity endofunctor is joined to the constant endofunctor at $g$ by a finite zigzag of natural transformations. If $A$ is itself an oriented interval and the fibre is restricted to functors nondecreasing in the numerical vertex orders, the same assertion holds whenever that restricted fibre is nonempty.

\emph{Proof.} The category $J$ is a poset. Its order relation $a\leq_J a'$ means that the unique numerical segment between $a$ and $a'$ is oriented all the way from $a$ to $a'$. For fixed $c$, each of the numerical clamping maps
\[
 a\longmapsto\min(a,c),\qquad
 a\longmapsto\max(a,c)
\]
is a functor $J\to J$: the image of a directed segment is either a point or a directed subsegment with the same orientation. Consequently numerical minimum and maximum are functors $J\times J\to J$. Indeed they are order preserving in either variable separately, and a product-order comparison factors into one comparison in each variable.

Assume that $g\in\mathcal F_b$. For each vertex $t$ of $J$ define endofunctors of $\mathcal F_b$ by pointwise operations
\[
 F_t(f)=\min\bigl(f,\max(g,t)\bigr),\qquad
 G_t(f)=\min\bigl(g,\max(f,t)\bigr),
\]
where $t$ denotes the constant functor $A\to J$. These are functors on diagrams and on natural transformations because minimum and maximum are functors. They preserve the prescribed restriction: at an object of $B$, both $f$ and $g$ have the same value $a$, and
\[
 \min(a,\max(a,t))=a.
\]
Since $J$ is a poset, the component of any resulting transformation at such an object is the identity. Thus the construction really stays inside the strict fibre.

Each generating arrow $t\to t'$ of $J$ induces natural transformations $F_t\to F_{t'}$ and $G_t\to G_{t'}$. If an edge points in the opposite direction, both induced transformations point in that opposite direction. We therefore obtain finite zigzags as $t$ runs numerically from $0$ to $r$. At the ends,
\[
 F_r=1_{\mathcal F_b},\qquad
 F_0=\min(-,g)=G_0,\qquad
 G_r=\operatorname{const}_g.
\]
Traversing the $F$ zigzag backwards and the $G$ zigzag forwards contracts the identity to $g$. Every functor in the zigzag fixes $g$ itself. On taking nerves and geometric realizations, natural transformations give homotopies, so $|N\mathcal F_b|$ is contractible as well. If $r=0$, the nonempty fibre has one object and the conclusion is immediate.

For the numerical-monotonicity variant, the pointwise minimum or maximum of two nondecreasing numerical functions is nondecreasing, and constant functions are nondecreasing. Hence all $F_t$ and $G_t$ preserve the indicated full subcategory. Their restrictions give the same contraction. This completes the proof.\hfill$\square$

The construction gives a uniform bound on the length of the contracting zigzag: there are $r$ comparisons in the $F$ family and $r$ in the $G$ family. This bound depends only on the target interval, not on the number of objects or arrows of $A$. In particular, no finite-source hypothesis was concealed in the proof. When $r=1$, the contraction has the particularly transparent form
\[
 1_{\mathcal F_b}\longleftarrow\min(-,g)
       \longrightarrow\operatorname{const}_g
\]
if the target edge points from $0$ to $1$; both arrows reverse if it points from $1$ to $0$. The fixed-boundary identity used above makes this a contraction inside the fibre even when $j$ identifies several boundary labels. That includes the two coincident endpoint labels of an excursion.

Taking $B$ to be the two endpoint labels proves both the proposed connectedness and the stronger contractibility of the generalized-interval functor category. Taking a fixed pair of equal endpoint values gives the version needed for an excursion followed by its return. For alternating intervals, the numerically monotone version strengthens the contraction-set argument on pages~17--21 and preserves the condition on backwards arrows when used for reparametrization.

The minimum--maximum preservation property is familiar in the theory of ordered digraphs; compare the definition of a Min--Max ordering in \citet[p.~1577]{HellRafiey2012}. We checked that account and the localization sources cited in this section, together with searches for relative functor categories of oriented intervals. We did not locate this fixed-restriction contraction in those sources. The result is presented as a direct continuation of the manuscript, with its full proof, and carries no claim of priority.
\end{dcomment}


\textbf{Lemme.} \textit{Soient $I$, $I'$ deux intervalles généralisés, munis d'objets source et objet but (de sorte que $I$, $I'$ soient} canoniquement \textit{isomorphes aux intervalles genéralisés standards correspondants). Soit} $\uHom^!(I,I')$ % [sic] genéralisés -> généralisés, missing accent on the first e
\msp{page 26}\textit{la catégorie des foncteurs de $I$ dans $I'$, compatibles avec source et but. Si cette catégorie est $\neq \emptyset$, elle est} 0\textit{-connexe} (et, je présume, même contractile?).

Je vais admettre provisoirement qu'il en est ainsi, du moins dans le cas particulier de $\Phi_4$ (intervalle généralisé de longueur 8) et $\bDelta_2$. Cela donne donc un foncteur
\[
\theta : X \longrightarrow X_\Sigma.
\]
On a gagné, car la propriété universelle de ce foncteur est pratiquement évidente. Mais j'ai oublié de vérifier que $\theta$ transforme éléments de $\Sigma$ en flèches inversibles. Soit donc
\[
x \xrightarrow{u} y
\]
avec $u \in \Sigma$, on va définir un chemin en sens inverse
\[
y \xrightarrow{c} x
\]
qui définira dans $X_\Sigma$ une flèche inverse de $u$. L'idée naturelle est de prendre le `chemin'
\[
y \xleftarrow{u} x
\]
% foot: Version 9 septembre 2001
\printedpage{17}
% head: A. Grothendieck : Dérivateurs, Ch. V, Édition des Universités de Montpellier II et Paris VII
de $y$ dans $x$, mais ce n'est pas un chemin au sens de nos définitions, on va le compléter en
\[
c : y \xrightarrow{\mathrm{id}} y \xleftarrow{u} x \xrightarrow{\mathrm{id}} x \xleftarrow{\mathrm{id}} x.
\]
Je dis que $c$ fait l'affaire. On doit prouver deux relations d'homotopie, de $c \cdot \bar{u}$ et de $\bar{u} \cdot c$ avec les chemins identiques en $x$ \texttt{[resp. en} $y$\texttt{]} de longueur 8 (indice 4, i.e. dans $\uCh_{\Sigma,4}(X;x,x)$ resp. $\uCh_{\Sigma,4}(X;y,y)$),
\msp{page 27}où $\bar{u}$ est le chemin associé à $u$
\[
x \xrightarrow{u} y \xleftarrow{\mathrm{id}} y \xrightarrow{\mathrm{id}} y \xleftarrow{\mathrm{id}} y.
\]

On a donc à regarder
\[
x \xrightarrow{u} y \xleftarrow{\mathrm{id}} y \xrightarrow{\mathrm{id}} y \xleftarrow{\mathrm{id}} y \xrightarrow{\mathrm{id}} y \xleftarrow{u} x \xrightarrow{\mathrm{id}} x \xleftarrow{\mathrm{id}} x
\]
et voir qu'il est homotope au chemin trivial
\[
x \xrightarrow{\mathrm{id}} x \xleftarrow{\mathrm{id}} \cdots x \xrightarrow{\mathrm{id}} x \xleftarrow{\mathrm{id}} x
\]
(des identités partout), et itou pour
\[
y \xrightarrow{\mathrm{id}} y \xleftarrow{u} x \xrightarrow{\mathrm{id}} x \xleftarrow{\mathrm{id}} x \xrightarrow{u} y \xleftarrow{\mathrm{id}} y \xrightarrow{\mathrm{id}} y \xleftarrow{\mathrm{id}} y
\]
et le chemin trivial en $y$. Cette fois, on a encore deux foncteurs
\[
\begin{tikzcd}[column sep=large]
\bDelta_1 & \Phi_4 \arrow[l, shift right, "\varphi"'] \arrow[l, shift left, "\psi"] \\
& \varepsilon, \arrow[ul, "\alpha'"] \arrow[u, "\alpha"']
\end{tikzcd}
\]
compatibles avec objets sources et objets buts, et on doit prouver qu'ils sont homotopes dans $\uHom^!(\Phi_4, \bDelta_1)$. Cela doit donc résulter du même lemme général. On a gagné - par une interprétation simple de $X\Sigma^{-1}$, et en particulier des $\Hom_{X\Sigma^{-1}}(x,y)$ et de leur composition, en termes de la géométrie des espaces de chemins :
\[
\left\{
\begin{aligned}
&\Hom_{X\Sigma^{-1}} \simeq \pi_0(\uCh_\Sigma(X;x,y)) \\ % [sic] Hom_{XΣ^{-1}} -> Hom_{XΣ^{-1}}(x,y), arguments (x,y) missing
&\text{Composition des homomorphismes dans } X\Sigma^{-1} \text{ correspond à} \\
&\text{la composition des } \Sigma\text{-chemins dans } X
\end{aligned}
\right.
\]
% foot: Version 9 septembre 2001

\dgrblock{dgrb173}{Cancellation uses the same endpoint twice}
\begin{dcomment}{examples}{p27:Cette fois, on a encore deux foncteurs}{Cancellation uses the same endpoint twice}
Grothendieck has written down the correct inverse path. The only remaining work is to make its two cancellation witnesses explicit. Let $u:x\to y$ belong to $\Sigma$, and write $T:[1]\to X$ for the corresponding functor. In the numerical ordering $0<1$, the first eight-edge excursion on the page is the image of the row
\[
 a=(0,1,1,1,1,1,0,0,0),
\]
whereas the second is the image of
\[
 b=(1,1,0,0,0,1,1,1,1).
\]
The source is $\Phi_4$, numbered from $0$ to $8$. Both rows are functors: on a forward source edge the value can increase, and on a backwards edge it can decrease. Their endpoint values are respectively $(0,0)$ and $(1,1)$.

There are pointwise natural transformations
\[
 (0,0,0,0,0,0,0,0,0)\longrightarrow a,
 \qquad
 b\longrightarrow(1,1,1,1,1,1,1,1,1).
\]
Applying $T$ gives a transformation from the constant path at $x$ to the first excursion, and one from the second excursion to the constant path at $y$. Their endpoint components are identities. Each horizontal backwards edge in these diagrams is either an identity or $u$, so it belongs to $\Sigma$. Thus both transformations are morphisms in the required $\Sigma$-path fibres, and the two cancellation identities follow immediately on $\pi_0$.

There is a small mismatch in the manuscript's invocation of the general lemma. That lemma is worded for maps taking the initial endpoint to the initial endpoint and the final endpoint to the final endpoint. These cancellation diagrams instead take both endpoints to $0$, or both to $1$. The relative-fibre theorem above includes them, but the two displayed transformations already prove exactly what is needed. This is a missing specification of the boundary condition, with an elementary repair, rather than a failure of localization.
\end{dcomment}


\dgrblock{dgrb174}{Completing the universal property}
\begin{dcomment}{exposition}{p27:par une interprétation simple}{Completing the universal property}
Grothendieck's path category now gives the promised localization. We record the complete conclusion because the construction uses less structure than many familiar fraction calculi.

\textbf{Path presentation of ordinary localization.} Let $X$ be a small category and let $\Sigma$ contain all its identity arrows. There is a canonical isomorphism of categories over $X$
\[
 X[\Sigma^{-1}]\cong X_\Sigma,
 \qquad
 \operatorname{Hom}_{X_\Sigma}(x,y)
   =\pi_0\operatorname{Ch}_\Sigma(X;x,y),
\]
where the path categories use strict endpoint fibres and composition is induced by concatenation at even indices.

\emph{Proof.} The reparametrization argument proves independence of padding. More explicitly, the maps that compare padding before and after concatenation are numerically monotone endpoint-preserving functors between alternating intervals. The connectedness proved by even-block slides supplies a zigzag between them. After precomposition with a $\Sigma$-path, every backwards edge is still a backwards edge of that path or an identity. Thus the comparison is a zigzag inside the specified path category. Morphisms of the factors concatenate to morphisms of the result, so the product descends to components. Associativity follows from strict associativity of finite concatenation, after choosing common even indices; the zero-length paths give identities.

The preceding table proves that $\theta:X\to X_\Sigma$ preserves composition, and the cancellation ladders prove that it inverts $\Sigma$. Let $F:X\to Y$ invert $\Sigma$. Evaluate each finite path in $Y$ by using $F$ on forward arrows and inverses of $F(s)$ on backwards arrows. Identity padding leaves the result unchanged. Naturality squares, with the backwards ones inverted as above, show that endpoint-fixed ladders give equal arrows of $Y$. Evaluation therefore defines maps on the proposed Hom sets. It takes concatenation to composition and constant paths to identities, so it defines a functor $\overline F:X_\Sigma\to Y$ with $\overline F\theta=F$.

To prove uniqueness, one must know that a path class is the composite of its elementary arrows and inverse arrows. Pad the original path to an even index. Replace each forward arrow by its four-edge representative under $\theta$, and each backwards arrow by the four-edge inverse representative. Concatenating these representatives gives the original list of nonidentity edges with additional identities inserted. It is the pullback of the original path by a numerically monotone contraction from a larger alternating interval. The usual end-padding is another such pullback. The two contractions belong to the same connected endpoint-preserving category, so the two paths have the same component. Hence every path class is the required composite, and any extension of $F$ must evaluate it as $\overline F$ does.

Thus $\theta$ has the defining universal property of localization. Applying the two universal properties gives inverse functors between $X_\Sigma$ and $X[\Sigma^{-1}]$ that agree with the identity on $X$. Their composites are identities by uniqueness.\hfill$\square$

The proof explains the exact status of the manuscript: its principal theorem is correct, and its proof becomes complete after the endpoint-fibre convention, the even-index repair already made by the author, and the supplied interval comparisons. No Ore condition or closure of $\Sigma$ under composition has entered.
\end{dcomment}


\dgrblock{dgrb175}{From components of path categories to mapping spaces}
\begin{dcomment}{later}{p27:composition, en termes de la géométrie des espaces de chemins}{From components of path categories to mapping spaces}
Grothendieck has recovered a set of morphisms by taking connected components of a category of presentations. The next question is whether one can retain higher homotopy information and compose it coherently. The relevant historical comparison already existed before these notes: the two papers of Dwyer and Kan appeared in 1980.

Their standard simplicial localization has the ordinary localization as its category of components \citep[\S4.2]{DwyerKan1980Simplicial}. The hammock construction gives a more diagrammatic model: its simplices are commutative ladders of variable length, with all vertical arrows and all leftward horizontal arrows in the chosen subcategory of weak equivalences. Composition is concatenation followed by reduction \citep[Definition~2.1 and Proposition~3.1]{DwyerKan1980Calculating}. These are close relatives of the present paths, but the vertical-arrow requirement is different: Grothendieck permits arbitrary vertical arrows. Therefore the formula for $\pi_0$ proved here does not itself identify the nerve of his path category with a derived mapping space.

There is also a preliminary distinction in the input data. The classical Dwyer--Kan construction starts with a subcategory $W$ containing all objects, whereas the present $\Sigma$ need contain only identities. To compare the two ordinary localizations one may take $W$ to be the subcategory generated by $\Sigma$: every functor inverting $\Sigma$ inverts its finite composites, and the converse is immediate. This proves equality of their universal properties. It does not identify the intervening categories of paths. The example with $\Sigma=\{1,t\}$ above shows concretely how adjoining composite denominators changes the allowed horizontal presentations before any homotopy comparison is proved.

The comparison becomes sharper under additional hypotheses. A homotopy calculus of fractions allows certain short zigzag categories to compute hammock mapping spaces; the exact reduction results are \citet[Proposition~6.2]{DwyerKan1980Calculating}. Their hypotheses concern the homotopy types of diagram categories, not just the existence of short ordinary fractions. Consequently a reduction that is sufficient for Hom sets need not preserve all higher homotopy groups.

One elementary check shows the difference in what is being sought. If every arrow of a category $X$ is inverted, the ordinary localization records only its groupoid completion. A higher localization retains the homotopy type of the classifying space: in the connected case the mapping spaces have the homotopy type of its loop space \citep[\S5.5]{DwyerKan1980Simplicial}. Take the poset of nonempty proper faces of a tetrahedron: its nerve is the barycentric subdivision of the boundary sphere $S^2$. Its groupoid completion is equivalent to the terminal groupoid, since the sphere is connected and simply connected, while the higher endomorphism space has the nontrivial homotopy type $\Omega S^2$. This example does not decide the homotopy type of Grothendieck's particular path category; it explains why agreement on its components is only the first comparison one would require.

The later development of simplicial and higher localizations gives a precise home for the ambition of retaining those spaces of presentations. It does not retroactively turn every category of zigzags into the same model. The admissible horizontal arrows, the vertical arrows, the allowed reparametrizations, and the endpoint conditions are part of the theorem to be proved.
\end{dcomment}


\printedpage{18}
% head: A. Grothendieck : Dérivateurs, Ch. V, Édition des Universités de Montpellier II et Paris VII
\noindent\msp{page 28}

\section*{3\quad Passage aux catégories de diagrammes}

Nous aurions dû noter la catégorie $X$ $\mathcal{M}$, car par la suite elle jouera le rôle d'une `catégorie de modèles'. On désignera par contre par $X$ les `catégories d'indices' sur lesquelles on veut construire les `catégories de diagrammes' de type $X$
\[
\mathcal{M}(X) = \uHom(X,\mathcal{M}).
\]
On se donne donc
\[
\underbrace{\Sigma}_{\text{quasi-isomorphismes}} \subseteq \mathrm{Fl}(\mathcal{M}), \qquad \Sigma \supset \text{ensemble des flèches identiques},
\]
et pour toute catégorie $X$ on définit
\[
\Sigma_X \subset \mathrm{Fl}\,\mathcal{M}(X)
\]
ainsi: si $F, G \in \mathrm{Ob}\,\mathcal{M}(X)$ ($F, G : X \rightrightarrows \mathcal{M}$), une flèche $u : F \longrightarrow G$ (homomorphisme de foncteurs) est dans $\Sigma_X$ si et seulement si $\forall x \in X$, $u(x) : F(x) \longrightarrow G(x)$ est dans $\Sigma$ (`critère de quasi-isomorphie fibre par fibre', pour l'homomophisme entre les copréfaisceaux $F$, $G$ envisagés). % [sic] homomophisme -> homomorphisme, letter r missing

\dgrblock{dgrb176}{The parameter is a diagram shape}
\begin{dcomment}{exposition}{p28:On désignera par contre}{The parameter is a diagram shape}
Grothendieck now asks the path construction to do more than compute a single localized Hom-set. He wants the same calculation after replacing the coefficient category by every category of diagrams. This is the point at which a category of paths begins to behave like an external Hom: its test objects are the indexing categories themselves. The change of letters is deliberate. In this section $\mathcal M$ supplies coefficients, whereas $X$ is allowed to vary as a parameter.

There is no model-category structure in the hypotheses. We have a category $\mathcal M$ and a class $\Sigma$ containing the identities; for the localization calculation even closure of $\Sigma$ under composition is unnecessary. One may enlarge $\Sigma$ to the subcategory it generates without changing the localization. We work in fixed universes, so that the localizations and the Hom-sets under discussion exist in the chosen larger universe. The designation ``quasi-isomorphism'' means membership in this specified class, rather than a statement about chain complexes.

Write $\mathcal M^X=\operatorname{Fun}(X,\mathcal M)$ and let $\Sigma_X$ consist of natural transformations whose components belong to $\Sigma$. Restriction along $f:X\to Y$ preserves these transformations and therefore induces
\[
 f^*: \mathcal M^Y[\Sigma_Y^{-1}]
       \longrightarrow \mathcal M^X[\Sigma_X^{-1}].
\]
The two objects of special interest are the constant diagrams $a_X,b_X$. Thus the functor to be understood is
\[
 H_{a,b}(X)=\operatorname{Hom}_{\mathcal M^X[\Sigma_X^{-1}]}(a_X,b_X),
 \qquad H_{a,b}:\Cat^{\mathrm{op}}\to\Ens.
\]
These are constant diagrams before localization. We are not replacing an arbitrary diagram by its values in $\mathcal M[\Sigma^{-1}]$; that would discard precisely the coherent information retained in Chapters~I and~II. See the commentary on Chapter~I, manuscript p.~13 (``Localizing diagrams and deriving their Kan extensions'') for that distinction.

For a fixed finite path shape $\Phi_n$, the elementary exponential law
\[
 \operatorname{Fun}(\Phi_n,\operatorname{Fun}(X,\mathcal M))
 \cong\operatorname{Fun}(X,\operatorname{Fun}(\Phi_n,\mathcal M))
\]
already contains the calculation. Its objects agree because both sides are functors $X\times\Phi_n\to\mathcal M$; its arrows agree because both sides are natural transformations of such functors. The restrictions on backward arrows are checked at each object of $X$, and the restrictions on the endpoints are also objectwise. The difficulty is therefore not the exponential law. It is the passage from one finite path shape to all path lengths, where one must retain a uniform bound.
\end{dcomment}


Soient
\[
a, b \in \mathrm{Ob}\,\mathcal{M}
\]
et considérons les foncteurs constants correspondants sur $X$
\[
a_X, b_X \in \mathrm{Ob}\,\mathcal{M}(X).
\]

\msp{page 29}On se propose d'appliquer la théorie du \S2 au `calcul' de
\[
\Hom_{\Sigma_X}(a_X, b_X),
\]
on a donc une bijection canonique
\[
\boxed{\Hom_{\Sigma_X}(a_X, b_X) \simeq \pi_0\uCh_{\Sigma_X}(\mathcal{M}(X); a_x, b_x).} % [sic] a_x, b_x -> a_X, b_X, lowercase index against a_X, b_X elsewhere
\]
On se propose d'expliciter la catégorie de chemins du second membre. On a, abstraction faite des données de $\Sigma$, $a$, $b$ dans $\mathcal{M}$, l'expression pour la catégorie de \textit{tous} les chemins
% foot: Version 9 septembre 2001

% ---- from Der-ChV-234_p021-025.tex
\printedpage{19}
% head: A. Grothendieck : Dérivateurs, Ch. V, Édition des Universités de Montpellier II et Paris VII

dans $\cM(X)$
\[
\begin{aligned}
\Ch(\cM(X)) \quad &\simeq \quad \uHom(\ \underbrace{\boldsymbol{\Phi}}_{\substack{\text{pro-catégorie}\\ (\Phi_n)_{n\in\mathbf{N}}}}\ ,\ \underbrace{\cM(X)}_{\uHom(X,M)}\ ) \\ % reading resolved: italic M is printed; category-valued Hom underlines restored.
&\dfeq \quad \varinjlim_n \uHom(\Phi_n,\cM(X)) \\
&\simeq \quad \varinjlim_n \uHom(\Phi_n\times X,\cM) \\
&\simeq \quad \varinjlim_n \uHom(X,\uHom(\Phi_n,\cM)) \\
&\dfeq \quad \varinjlim_n \uHom(X,\Ch_n(\cM)) \\
&\simeq \quad \uHom^{!}(X,\Ch(\cM)),
\end{aligned}
\]
où le signe ! signifie qu'on prend la sous-catégorie pleine des foncteurs
\[
X\longrightarrow \Ch(\cM),
\]
\msp{page 30}
qui se factorisent par une des sous-catégories pleines $\Ch_n(\cM)$. On peut dire aussi que si on considère le système inductif des catégories $\Ch_n(\cM)$ comme un ind-objet de Cat, qu'on va noter
\[
\Chh(\cM) \qquad\qquad\qquad\qquad \text{ind-objet de Cat,}
\]
on a
\[
\Ch(\cM(X)) \xrightarrow{\ \sim\ } \uHom(X,\Chh(\cM)).
\]
Si on en vient à $\Sigma$, $\Sigma_X$, on trouve aussitôt qu'un chemin dans $\cM(X)$ est un $\Sigma_X$-chemin, si et seulement si le foncteur correspondant $X\longrightarrow\Ch(\cM)$ se factorise par $\Ch_{\Sigma}(\cM)$. Ainsi on trouve
\[
\Ch_{\Sigma_X}(\cM(X))\simeq\uHom^{!}(X,\Ch_{\Sigma}(\cM)),
\]
où le signe ! a le même sens que plus haut. Si on introduit le ind-objet $\Chh_{\Sigma}(\cM)$ défini par les $\Ch_{\Sigma,n}(\cM)$, on trouve donc aussi
\[
\Ch_{\Sigma_X}(\cM(X))\xrightarrow{\ \sim\ }\uHom(X,\Chh_{\Sigma}(\cM)).
\]

\dgrblock{dgrb177}{Why the exclamation mark is necessary}
\begin{dcomment}{exposition}{p30:qui se factorisent par une}{Why the exclamation mark is necessary}
Grothendieck introduces the exclamation mark to remember that a path of diagrams has one length for the entire diagram. An object of $\operatorname{Fun}(X,\operatorname{Ch}_\Sigma(\mathcal M))$ may assign a finite path to every $x\in X$ without these lengths being bounded as $x$ varies. Such a functor need not come from a path in $\mathcal M^X$. This is a genuine distinction even before taking connected components or weak equivalences.

Let $K_n$ denote the fixed-endpoint path category at stage $n$, using the fully faithful identity-padding inclusions, and put $K_\infty=\bigcup_nK_n$. Then
\[
 \operatorname{Fun}^{!}(X,K_\infty)
   =\bigcup_n\operatorname{Fun}(X,K_n)
\]
is the full subcategory on bounded functors. Fullness deserves a verification. If $F$ and $G$ factor through $K_n$, every component of a natural transformation $F\Rightarrow G$ lies in $K_n$, because $K_n\subseteq K_\infty$ is full. Its naturality equations are the same equations in either category. Thus no additional bound is needed for arrows between already bounded objects.

The failure for an infinite source can be seen without any homotopy theory. Take $K_n$ to be the discrete category on $\{0,\ldots,n\}$ and $X$ the discrete category on $\mathbf N$. The functor $x\mapsto x$ from $X$ to $K_\infty$ lies in no $\operatorname{Fun}(X,K_n)$. It cannot even be isomorphic to a bounded functor, since the target is discrete. The ind-object $(K_n)$ records a uniform finite stage which the ordinary union has forgotten.

Connected components do commute with the filtered union used here. A component of the union has a representative object at some stage. If two objects become connected, a finite zigzag of arrows witnesses this, and all its objects and arrows occur at one sufficiently large stage. This proves directly that
\[
 \varinjlim_n\pi_0\operatorname{Fun}(X,K_n)
 \xrightarrow{\sim}
 \pi_0\!\left(\varinjlim_n\operatorname{Fun}(X,K_n)\right).
\]
Nothing in this argument permits one to replace the category inside the final parentheses by the unrestricted functor category. The exclamation mark and the double underline serve the same purpose: the first describes bounded functors, while the second retains the system of bounds as an ind-object.
\end{dcomment}


\msp{page 31}
Soit $c$ un chemin dans $\cM(X)$, et
\[
X\longrightarrow\Ch(\cM)
\]
le foncteur qu'il définit, composons-le avec les foncteurs source-but
\[
\Ch(\cM)\longrightarrow\cM\times\cM,
\]
on trouve un foncteur
\[
X\longrightarrow\cM\times\cM,
\]
i.e. un objet $(F_0,F_1)$ de $\cM(X)\times\cM(X)$. Je dis que $F_0$ et $F_1$ sont respectivement la source et le but du chemin $c$. De façon un peu plus précise, on a commutativité (stricte) dans le diagramme de foncteurs
% foot: Version 9 septembre 2001
\printedpage{20}
% head: A. Grothendieck : Dérivateurs, Ch. V, Édition des Universités de Montpellier II et Paris VII
\[
\begin{tikzcd}[column sep=large, row sep=large]
\Ch_{\Sigma_X}(\cM(X)) \arrow[r, "\sim"] \arrow[d, "{\sigma_{\cM(X)},\,\beta_{\cM(X)}}"] & \uHom^{!}(X,\Ch_{\Sigma}(\cM)) \arrow[d, "{\text{déduit de }\sigma_{\cM},\,\beta_{\cM}}"] \\
\cM(X)\times\cM(X) \arrow[r, "\sim"] & \uHom(X,\cM\times\cM),
\end{tikzcd}
\]
où les foncteurs $\sigma$ sont les foncteurs source, les foncteurs $\beta$ les foncteurs but.

Si on se donne arbitrairement les objets $F_0$, $F_1$ dans $\cM(X)$, donc un foncteur
\[
X\longrightarrow\cM\times\cM,
\]
$\Ch_{\Sigma_X}(F_0,F_1)$ s'identifie donc à la catégorie des foncteurs ‘admissibles’ (i.e. se factorisant par un $\Ch_{\Sigma,n}(\cM)$) $X\xrightarrow{\ C\ }\Ch_{\Sigma}(\cM)$ qui \textit{relèvent}
\msp{page 32 vide}
\msp{page 33}
le foncteur $(\sigma_{\cM},\beta_{\cM})$ \texttt{[plutôt} $(F_0,F_1)$\texttt{]} précédent, i.e. qui donnent lieu à un triangle commutatif (strict)
\[
\begin{tikzcd}[row sep=large]
 & \Ch_{\Sigma}(\cM) \arrow[d, "{(\sigma_{\cM},\,\beta_{\cM})}"] \\
X \arrow[ur, "C"] \arrow[r, "{(F_0,F_1)}"] & \cM\times\cM.
\end{tikzcd}
\]
Mais si les foncteurs $F_0$, $F_1$ sont constants, de valeur $a$, $b$ comme tantôt, i.e. $(F_0,F_1)$ constant de valeur $(a,b)$, alors la catégorie des foncteurs $C$ en question est tautologiquement isomorphe à celle des foncteurs de $X$ dans la fibre de $(\sigma_{\cM},\beta_{\cM})$ en $(a,b)$, d'où l'isomorphisme de catégories fondamental
\[
\boxed{\Ch_{\Sigma_X}(X;a_X,b_X)\simeq\uHom^{!}(X,\Ch_{\Sigma}(\cM;a,b)).}
\]

On va expliciter le contenu essentiel de cette formule en une

\dgrblock{dgrb178}{The precise meaning of ind-representation}
\begin{dcomment}{exposition}{p33:On va expliciter le contenu essentiel}{The precise meaning of ind-representation}
The scholium isolates what survives every arbitrary choice in the path calculation. If $\Cathot$ denotes the category with small categories as objects and
\[
 \operatorname{Hom}_{\Cathot}(X,Y)
       =\pi_0\operatorname{Fun}(X,Y),
\]
then the result says that $H_{a,b}$ is a filtered colimit of representable presheaves on $\Cathot$. The homotopy relation in this formula is generated by natural transformations, with arrows in either direction. Composition is well defined because whiskering preserves natural transformations.

For clarity, the fully faithful embedding of ind-objects in presheaves can be proved in the notation needed here. To an ind-object $K=(K_i)$ of a locally small category $C$ associate
\[
 h_K(X)=\varinjlim_i\operatorname{Hom}_C(X,K_i).
\]
If $L=(L_j)$ is another ind-object, Yoneda and the universal property of a colimit give
\[
 \operatorname{Nat}(h_K,h_L)
 \cong\varprojlim_i\operatorname{Nat}(h_{K_i},h_L)
 \cong\varprojlim_i\varinjlim_j\operatorname{Hom}_C(K_i,L_j)
 =\operatorname{Hom}_{\operatorname{Ind}(C)}(K,L).
\]
The limits and colimits occur in that order. No interchange theorem has been used. A natural isomorphism of the presheaves consequently gives a unique isomorphism of the corresponding ind-objects, compatible with the specified representations. This is the exact sense of the uniqueness statement on manuscript p.~34.

To obtain the scholium, apply the preceding path-localization theorem inside $\mathcal M^X$, then apply the exponential law with fixed endpoints, and finally commute $\pi_0$ with the filtered union. This yields
\[
 H_{a,b}(X)
  \cong\varinjlim_n\operatorname{Hom}_{\Cathot}(X,K_n)
  =h_{(K_n)}(X).
\]
The notation in the boxed formula just before the scholium has $X$ where the ambient path category should be $\mathcal M(X)$; the edition prints this slip. The preceding endpoint square determines the intended meaning unambiguously.

The result has a useful limit. It gives uniqueness in $\operatorname{Ind}(\Cathot)$, not uniqueness of the ordinary union $K_\infty$ in $\Cathot$. A colimit of a chosen sequence of actual categories contains information about the representatives and their compatibility which an ind-object in a homotopy category need not retain. The questions at the end of the chapter begin exactly at this boundary.
\end{dcomment}


\textbf{Scholie.} Soient $\cM$ une catégorie avec un ensemble $\Sigma$ de flèches de localisation, $a,b\in\mathrm{Ob}\,\cM$. Alors on peut trouver un objet $\KK=K(a,b)$, $K=(K_n)$, de Ind\,Cat (indexé par $\mathbf{N}$ mais peu importe ce fait), et un isomorphisme canonique pour toute catégorie $X$ (isomorphisme fonctoriel en $X$)
\[
\begin{aligned}
\mathrm{Hom}_{\Sigma_X}(a_X,b_X) \quad &\simeq \quad \pi_0(\uHom(X,\KK)) \\
&\dfeq \quad \pi_0(\ \varinjlim \uHom(X,K_n)) \\
&\simeq \quad \varinjlim \pi_0(\uHom(X,K_n)).
\end{aligned}
\]
% foot: Version 9 septembre 2001

\printedpage{21}
% head: A. Grothendieck : Dérivateurs, Ch. V, Édition des Universités de Montpellier II et Paris VII

\msp{page 34}
Je vais épiloguer sur ce résultat, sans me soucier pour l'instant de la façon particulière dont la ind-catégorie $K$ \texttt{[plutôt} $\KK$\texttt{]} a été obtenue en termes de $\cM$, $\Sigma$, $a$, $b$.

Comme les $\pi_0\uHom(X,K_n)=\mathrm{Hom}_{\Cathot}(X,K_n)$ ne dépendent que de $X$ et de $K_n$ en tant qu'objets de Cat\,hot, je vais interpréter $\KK$ comme un ind-objet de Cat\,hot, et écrire la formule fondamentale sous la forme
\[
\mathrm{Hom}_{\Sigma_X}(a_X,b_X)\simeq\mathrm{Hom}_{\mathrm{Ind\,Cat\,hot}}(X,\KK).
\]

Ainsi, l'énoncé du scholie implique deux choses :

\begin{enumerate}
\item[a)] Le contrafoncteur $X\longmapsto\mathrm{Hom}_{\Sigma_X}(a_X,b_X)$ sur Cat se factorise canoniquement par Cat\,hot, via un foncteur canonique
\[
\Cathot^{o}\ \xrightarrow{\ H=H^{\cM,\Sigma}_{a,b}\ }\mathrm{Ens}\ ;
\]
\item[b)] Ce dernier contrafoncteur sur Cat\,hot est \textit{ind-représentable} par un ind-objet $\KK$ de Cat\,hot.
\end{enumerate}

C'est cet ind-objet de Cat\,hot qui est défini à isomorphisme unique près par $\cM$, $\Sigma$, $a$, $b$, indépendamment de tout artifice de calcul pour expliciter les $\mathrm{Hom}_{\Sigma_X}(a_X,b_X)$.

Notons aussi la conclusion suivante, nullement évidente à priori :

\msp{page 35}
Notons d'abord (ceci étant implicite dans le scholie) que toute flèche dans Cat
\[
X\xrightarrow{\ f\ }Y
\]
définit
\[
\cM(Y)\xrightarrow{\cM(f)}\cM(X) \qquad\qquad\qquad \text{appliquant } \Sigma_Y \text{ dans } \Sigma_X,
\]
d'où
\[
\cM(Y)\Sigma_Y^{-1}\dfeq\mathbf{D}_{\cM,\Sigma}(Y)\longrightarrow\cM_X\Sigma_X^{-1}=\mathbf{D}_{\cM,\Sigma}(X),
\]
\texttt{[plutôt} $\cM(X)$\texttt{]} d'où, pour tout $a,b\in\mathrm{Ob}\,\cM$, vu que $\cM(f)$ transforme $a_Y$, $b_Y$ en $a_X$, $b_X$ :
\[
\boxed{\mathrm{Hom}_{\Sigma_Y}(a_Y,b_Y)\longrightarrow\mathrm{Hom}_{\Sigma_X}(a_X,b_X).}
\]

Ceci posé, on a :

\textbf{Corollaire :} \textit{Si $f$ est un homotopisme, alors quelque soient $a,b\in\cM$, la flèche précédente est un isomorphisme.} % [sic] quelque soient -> quels que soient, grammatical form of the locution

La question naturelle, c'est de savoir si la même conclusion est encore valable, en supposant seulement que $f$ soit un ‘\textit{quasi-isomorphisme}’ (ou ‘équivalence faible’), ici que $f$ définisse un isomorphisme dans la catégorie de fractions Hot de Cat. (À ne pas confondre avec Cat\,hot,
% foot: Version 9 septembre 2001
\printedpage{22}
% head: A. Grothendieck : Dérivateurs, Ch. V, Édition des Universités de Montpellier II et Paris VII
\msp{page 36}
nettement plus fine). S'il en était ainsi, le contrafoncteur
\[
X\longmapsto\mathrm{Hom}_{\Sigma_X}(a_X,b_X)
\]
se factoriserait non seulement par Cat\,hot, mais même par la catégorie de fractions plus grossière Hot. Et dès lors se poserait la question si ce contrafoncteur hypothétique
\[
\mathrm{Hot}^{o}\longrightarrow\mathrm{Ens}
\]
est encore ind-représentable; et de préférence (\footnote[5]{ce sera automatique si on a ind-représentabilité comme foncteur sur Hot, cf. ci-dessous.}) ind-représentable par le ind-objet image de $\KK$ dans Ind\,Hot par
\[
\mathrm{Ind}(\Cathot)\longrightarrow\mathrm{Ind}(\mathrm{Hot}),
\]
déduit du foncteur de localisation $\Cathot\longrightarrow\mathrm{Hot}$. (Car cet ind-objet est un candidat tout trouvé.)

\dgrblock{dgrb179}{Two homotopy invariance questions}
\begin{dcomment}{exposition}{p35:La question naturelle, c'est de savoir}{Two homotopy invariance questions}
Grothendieck has proved invariance under categorical homotopy equivalence, and now asks for invariance under weak equivalence of nerves. The first assertion is immediate from the ind-representation. A natural transformation $f\Rightarrow g$ induces, by precomposition, a natural transformation between the two functors on each $\operatorname{Fun}(Y,K_n)$. These functors therefore induce the same map on $\pi_0$. A finite zigzag gives the same conclusion, and a homotopy inverse to $f$ gives inverse maps on $H_{a,b}$.

The second assertion has additional content. A weak equivalence $Nf$ need not admit an inverse represented by a functor and a finite zigzag of natural transformations. Nor does the construction of $H_{a,b}$, for an arbitrary class $\Sigma$, supply a derived mapping space. Thus the corollary proves factorization through $\Cathot$; it does not prove factorization through $\Hot$. This is the same difference that appeared in the commentary on Chapter~I, manuscript p.~102 (``What the completion requires''), now applied to a particular family of coefficients.

Once that factorization is established, however, the subsequent question about ind-representability has a formal affirmative answer. It does not require the proposed right adjoint to $\Cathot\to\Hot$. More generally let $\varphi:C\to D$ be a localization, let $H=\varinjlim_i h_{K_i}$ on $C$, and suppose $H=\varphi^*H'$ for a presheaf $H'$ on $D$. Restriction $\varphi^*$ is fully faithful by the universal property of localization. Put $L=\varinjlim_i h_{\varphi K_i}$ on $D$. For every presheaf $F$ on $D$, Yoneda gives natural bijections
\[
 \operatorname{Nat}(L,F)
 \cong\varprojlim_i F(\varphi K_i)
 \cong\operatorname{Nat}(H,\varphi^*F)
 \cong\operatorname{Nat}(H',F).
\]
Yoneda, now in the presheaf category, proves $L\cong H'$. Equivalently, left Kan extension along $\varphi^{\mathrm{op}}$ sends $h_K$ to $h_{\varphi K}$, preserves the filtered colimit, and recovers a descended presheaf; this is the usual Kan-extension formalism \citep[\S\S1.1--1.4]{Riehl2014}. Thus, conditional only on weak-equivalence invariance, the image of $(K_i)$ in $\operatorname{Ind}(\Hot)$ represents the descended functor. What remains conditional on the later right adjoint is the characterization by a Thomason--Kan envelope, not this descent of ind-representability.

A useful sufficient hypothesis is that localization of the strict diagram categories presents the expected infinity-categories of diagrams. More explicitly, suppose a homotopy theory $\mathcal C$ satisfies, naturally in small $X$,
\[
 \mathcal M^X[\Sigma_X^{-1}]
   \simeq h\operatorname{Fun}(NX,\mathcal C).
\]
If $Z=\operatorname{Map}_{\mathcal C}(a,b)$, the space of natural transformations between constant diagrams is
\[
 \operatorname{Map}_{\operatorname{Fun}(NX,\mathcal C)}(a_X,b_X)
   \simeq\operatorname{Map}_{\mathcal S}(|NX|,Z).
\]
Indeed, a coherent transformation chooses a point of $Z$ at each object, a path for each arrow, and the compatible higher homotopies for composable strings; these are exactly the simplices of a map from $NX$ into $Z$. Taking $\pi_0$ proves $H_{a,b}(X)\cong[|NX|,Z]$. Under this hypothesis the functor is represented on $\Hot$, rather than merely ind-represented there.

Combinatorial simplicial model categories provide a familiar setting for this statement. The mapping space must be derived, using a cofibrant source and a fibrant target. The model-independent function-complex construction is due to Dwyer and Kan; their comparison theorem explains why the resulting homotopy type is independent of auxiliary model structures \citep[\S4]{DwyerKan1980Function}. The displayed hypothesis is stated explicitly here because it is not a consequence of an arbitrary choice of $\Sigma$.
\end{dcomment}


\dgrblock{dgrb180}{An external Hom that is an ordinary homotopy type}
\begin{dcomment}{examples}{p36:est encore ind-représentable}{An external Hom that is an ordinary homotopy type}
Grothendieck's candidate can be made completely concrete for simplicial sets with their usual weak equivalences. Take $a=\Delta^0$ and let $b=Z$ be a Kan complex. For a small category $X$, a projectively cofibrant replacement of the constant point diagram is supplied by the usual bar construction. At $x\in X$ its value is $N(X/x)$; this is contractible because $X/x$ has the terminal object $\operatorname{id}_x$. The diagram is assembled from the covariant representables by attaching simplices along their boundaries, so it is projectively cofibrant. The constant diagram $Z_X$ is projectively fibrant.

A natural transformation from this resolution to $Z_X$ is equivalent to a simplicial map from its colimit to $Z$. The colimit of the diagram $x\mapsto N(X/x)$ is $NX$: a simplex represented by a string ending in an arrow to $x$ is identified with the same string at its last vertex. Consequently the derived simplicial mapping space between the two constant diagrams is $Z^{NX}$. Since the source simplicial set is cofibrant and $Z$ is Kan,
\[
 \operatorname{Hom}_{\operatorname{Ho}(\mathbf{sSet}^{X})}
       (*_X,Z_X)=\pi_0 Z^{NX}=[NX,Z].
\]
This is the bar-construction calculation underlying the abstract assertion in the preceding block; see \citet[\S\S4.4--4.5 and 5.1]{Riehl2014} for the construction and its homotopical justification.

For $Z=K(A,m)$ with $m\geq1$ and $A$ abelian, it gives $H_{*,Z}(X)=H^m(NX;A)$, with constant coefficients. If $X$ is the one-object category attached to the cyclic group $C_r$, the familiar periodic resolution has alternating differentials $g-1$ and $1+g+\cdots+g^{r-1}$. After applying $\operatorname{Hom}_{\mathbf Z[C_r]}(-,A)$ for the trivial action, these differentials become $0$ and multiplication by $r$. Hence
\[
 H^{2q}(BC_r;A)=A/rA\quad(q\geq1),\qquad
 H^{2q+1}(BC_r;A)=A[r]\quad(q\geq0).
\]
The category $BC_r$ has finitely many objects and arrows, but its nerve has nondegenerate simplices in arbitrarily high dimensions. Thus ``finite category'' in the closing question must not be silently replaced by ``finite CW complex''. The same example also shows that successful homotopy-theoretic coefficients often produce an actually representable functor on $\Hot$. This is distinct from essential constancy in $\operatorname{Ind}(\Cathot)$, the subject of the pessimistic parenthesis on manuscript p.~46. Passing from the first assertion to the second would require precisely the right adjoint discussed next.
\end{dcomment}


\msp{page 37}
Je vais faire quelques rappels catégoriques, généralisant la situation du foncteur canonique de localisation
\[
\Cathot\xrightarrow{\ \varphi\ }\mathrm{Hot},
\]
dont on sait qu'il admet un adjoint à droite \texttt{[si un lecteur a une démonstration de ce fait, prière de contacter G. Maltsiniotis : maltsin@math.jussieu.fr]}
\[
\Cathot\xleftarrow{\ \psi\ }\mathrm{Hot}, \qquad\qquad\qquad \mathrm{Hom}_{\mathrm{Hot}}(\varphi X,Y)\simeq\mathrm{Hom}_{\Cathot}(X,\psi Y),
\]
où $\psi$ est le foncteur ‘enveloppe de Kan’ ou ‘représentant de Kan’. On va donc considérer la situation plus générale
\[
\begin{tikzcd}[column sep=large]
C \arrow[r, shift right=1, "\varphi"'] & C' \arrow[l, shift right=1, "\psi"']
\end{tikzcd}
\qquad\qquad \mathrm{Hom}_{C'}(\varphi X,X')\simeq\mathrm{Hom}_{C}(X,\psi X')
\]
i.e.
\[
\varphi,\psi \text{ couple de foncteurs adjoints.}
\]
On suppose de plus
\[
\left\{
\begin{aligned}
&\varphi \text{ un foncteur localisation, i.e. } \psi \text{ pleinement fidèle,}\\
&\text{i.e. } \varphi\psi\xrightarrow[\sim]{\text{adj}}\mathrm{id}_{C'} \quad .
\end{aligned}
\right.
\]
Soit $W=W_{\varphi}\subset\mathrm{Fl}\,C$, défini par $\varphi$ : $W_{\varphi}=\{u\in\mathrm{Fl}\,C\mid\varphi(u)\text{ isomorphisme}\}$. Voici la proposition que j'avais oubliée, et que j'ai eu du mal à reconstituer, hélas!

\dgrblock{dgrb181}{The reflection does not exist}
\begin{dcomment}{development}{p37:dont on sait qu'il admet un adjoint}{The reflection does not exist}
Grothendieck wants to convert homotopy invariance of a represented functor into a condition on its representing object. For a reflective localization this is an elementary and powerful argument, and Propositions~1 and~2 below are correct formal statements about any reflective localization. The edition questions the right adjoint to $\varphi:\mathrm{Cat}\,\mathrm{hot}\to\mathrm{Hot}$ on which their application rests. It does not exist.

\textbf{Theorem.}\ours
\textit{The functor $\varphi:\mathrm{Cat}\,\mathrm{hot}\to\mathrm{Hot}$ has no right adjoint. More precisely, there is no small category $K$ with bijections $\pi_0\operatorname{Fun}(X,K)\cong[BX,K(\mathbf Z/2,2)]$ natural in the small category $X$.}

\textit{Proof.}
A right adjoint $\psi$ would give such a $K$, namely $\psi$ of the homotopy type $Y=K(\mathbf Z/2,2)$, since morphisms $X\to K$ in $\mathrm{Cat}\,\mathrm{hot}$ are the connected components of $\operatorname{Fun}(X,K)$. Suppose $K$ is given. Taking $X=e$ shows that $\pi_0K=[\mathrm{pt},Y]$ is one point.

Let $\Omega$ be an infinite set of cardinality greater than that of the set of arrows of $K$, and $G$ the group of even permutations of $\Omega$ with finite support. It is the union of the alternating groups of the finite subsets of $\Omega$, so it is simple, and its cardinality is that of $\Omega$. A functor from the one-object category $\mathsf BG$ to $K$ is an object $k$ with a homomorphism $G\to\operatorname{Aut}(k)$. Its kernel is a normal subgroup; it is not trivial because $\operatorname{Aut}(k)$ has fewer elements than $G$; so the homomorphism is trivial. A natural transformation between two such constant functors is a single arrow of $K$. Hence $\operatorname{Fun}(\mathsf BG,K)\cong K$ and $\pi_0\operatorname{Fun}(\mathsf BG,K)$ is one point.

On the other side $[BG,K(\mathbf Z/2,2)]=H^2(G;\mathbf Z/2)$ classifies the central extensions of $G$ by $\mathbf Z/2$, and one of them does not split. Let $V$ be the real vector space with orthonormal basis $(e_\omega)_{\omega\in\Omega}$ and $\mathrm{Cl}(V)$ its Clifford algebra, with $v^2=-|v|^2$. For $\omega\neq\omega'$ the unit vector $\tau=(e_\omega-e_{\omega'})/\sqrt2$ is invertible in $\mathrm{Cl}(V)$ and $v\mapsto-\tau v\tau^{-1}$ is the reflection of $V$ that exchanges $e_\omega$ and $e_{\omega'}$. Let $\widehat G$ be the group of products of an even number of such vectors. Sending such a product to the corresponding product of transpositions is a surjective homomorphism $\widehat G\to G$; its kernel is contained in $\{\pm1\}$, the kernel of the map from the Pin group to the orthogonal group, and $-1=(\tau\tau')^2\in\widehat G$ for two orthogonal such vectors $\tau,\tau'$. So $\widehat G$ is a central extension of $G$ by $\mathbf Z/2$. The element $g=(\omega_1\,\omega_2)(\omega_3\,\omega_4)$ has order two, and its two lifts $\pm\tau\tau'$ have square $-1$; a splitting would lift $g$ to an element of order two. So $H^2(G;\mathbf Z/2)\neq0$, and the two sets cannot be in bijection.\hfill$\square$

The argument shows more than the statement: a category $K$ represents on $\mathrm{Cat}\,\mathrm{hot}$ only what homomorphisms into its automorphism groups and its arrows can see, and a group that is large and simple is invisible to it. The functors $X\mapsto[BX,Y]$ that are representable on all of $\mathrm{Cat}\,\mathrm{hot}$ are those with $Y$ a $1$-type, represented by a groupoid (the proof for a general $Y$, which replaces $G$ by large simple groups with the homology of spheres, is not reproduced here); the fundamental groupoid $X\mapsto\Pi_1X$ is a reflection of $\mathrm{Cat}\,\mathrm{hot}$ onto groupoids modulo natural isomorphism, and it inverts the functors that are bijective on $\pi_0$ and on all $\pi_1$, a larger class than the weak equivalences. Propositions~1 and~2 apply to that reflection.

Thomason's model structure gives a fibrant replacement $K\to RK$ and an equivalence between its homotopy category and $\mathrm{Hot}$ \citep[Theorem~4.9 and Corollary~5.5]{Thomason1980}. It does not identify $\pi_0\operatorname{Fun}(X,RK)$ with $[NX,NK]$ for every $X$: mapping formulas in a model category need a cofibrant source, and the cofibrant objects of that structure are posets. Below, references to an envelope or to a reflective localization of $\mathrm{Cat}\,\mathrm{hot}$ onto homotopy types are to be read as hypotheses that cannot be met; references to Thomason fibrant replacements mean the model-category construction and carry only its conclusions.
\end{dcomment}


\textbf{Proposition 1.} \textit{Soit $K\in \mathit{Ob}C$. Les conditions suivantes sur $K$ sont équivalentes :}
% foot: Version 9 septembre 2001

\printedpage{23}
% head: A. Grothendieck : Dérivateurs, Ch. V, Édition des Universités de Montpellier II et Paris VII

\begin{enumerate}
\item[(i)] \textit{Le foncteur}
\[
X\longmapsto K(X)\dfeq\mathrm{Hom}_C(X,K) \qquad\quad C^{o}\longrightarrow\mathrm{Ens}
\]
\textit{transforme flèches de $W$ (‘quasi-isomorphismes’) en bijections.}
\item[(i$'$)] \textit{Même condition pour les flèches d'adjonction}
\[
X\longrightarrow\psi\varphi X \qquad\quad (\forall X\in\mathrm{Ob}\ \mathrm{C})
\]
\item[(i$''$)] \textit{Même condition pour la flèche d'adjonction} (\footnote[6]{N.B. Il suffit même que l'élément $\mathrm{id}_K$ de $F_K(K)$ provienne de $F_K(\psi\varphi K)$ \texttt{[où} $F_K=K(-)=\mathrm{Hom}(-,K)$\texttt{]}, i.e.\par (i$'''$) Il existe une rétraction de $\psi\varphi X$ sur $X$ \texttt{[plutôt} $\psi\varphi K$ \texttt{sur} $K$\texttt{]} (NB Celle-ci sera même un isomorphisme).})
\[
K\longrightarrow\psi\varphi K.
\]
\item[(ii)] \textit{$K$ appartient à l'image essentielle de $\psi$}
\item[(ii$'$)] \textit{$K\longrightarrow\psi\varphi K$ (adjonction) est un isomorphisme.}
\item[(ii$''$)] \textit{L'homomorphisme fonctoriel en $X$}
\[
\mathrm{Hom}_C(X,K)\longrightarrow\mathrm{Hom}_{C'}(\varphi X,\varphi K)
\]
\textit{est un isomorphisme.}
\end{enumerate}

\textsc{Démonstration.} On a d'abord comme quasi-tautologies (ii$'$)$\Longrightarrow$ (ii)$\Longrightarrow$ (ii$''$)$\Longrightarrow$ (ii$'$), i.e. l'équivalence des trois dernières conditions. En effet, (ii$'$)$\Longrightarrow$ (ii) tautologique, (ii)$\Longrightarrow$ (ii$''$) car si $K=\psi(K')$ on a % reading resolved: single and double mathematical primes as printed.
\msp{page 38}
par adjonction
\[
\mathrm{Hom}_C(X,K)=\mathrm{Hom}_C(X,\psi(K'))\overset{\text{adj}}{\simeq}\mathrm{Hom}_{C'}(\varphi X,K'),
\]
or $\varphi(K)=\varphi\psi(K')\xrightarrow{\ \sim\ }K'$ (adjonction, et $\psi$ pleinement fidèle).
\[
\begin{tikzcd}[column sep=small, row sep=large]
\mathrm{Hom}_C(X,K) \arrow[rr, "\sim"] \arrow[dr, "\varphi_{X,K}"] & & \mathrm{Hom}_C(\varphi X,K') & \text{\footnotesize qui se complète en} \\ % [sic] Hom_C(\varphi X,K') -> Hom_{C'}(\varphi X,K'), \varphi X and K' are objects of C' (cf. the display above)
& \mathrm{Hom}_{C'}(\varphi X,\ \underbrace{\varphi K}_{=\varphi\psi(K')}\ ) \arrow[ur, "\sim", "\text{adj}"'] & &
\end{tikzcd}
\]
ce qui prouve que $\varphi_{X,K}$ est un isomorphisme.

(ii$''$) $\Longrightarrow$ (ii$'$). Car la flèche
% foot: Version 9 septembre 2001

% ---- from Der-ChV-234_p026-030.tex
\printedpage{24}
% head: A. Grothendieck : Dérivateurs, Ch. V, Édition des Universités de Montpellier II et Paris VII
% foot: Version 9 septembre 2001

\[
\begin{tikzcd}
\Hom_C(X,K) \arrow[r,"\alpha_X"] \arrow[dr,dashed] & \Hom_{C'}(\varphi X,\varphi K) \arrow[d,dash,"\wr"',"\mathrm{adj}"] \\
 & \Hom_C(X,\psi\varphi K)
\end{tikzcd}
\]
est déduite de la flèche d'adjonction
\[
K \longrightarrow \psi\varphi K
\]
en prenant $\Hom_C(X,-)$, donc $\alpha_X$ est un isomorphisme pour tout $X$ si et seulement si cette flèche d'adjonction est un isomorphisme. Donc
\[
\text{(ii)} \Longleftrightarrow \text{(ii$'$)} \Longleftrightarrow \text{(ii$''$)}.
\]
D'autre part,
\[
\text{(ii$''$)} \Longrightarrow \text{(i)} \Longrightarrow \text{(i$'$)} \Longrightarrow \text{(i$''$)}
\]
sont tautologiques, en tenant compte pour (i) $\Longrightarrow$ (i) \texttt{[probablement (i) $\Longrightarrow$ (i$'$)]} que les flèches d'adjonction
\[
X \xrightarrow{u} \psi\varphi X \qquad\qquad\qquad \text{dans } C
\]
sont bien dans $W$, car on a un triangle commutatif
\[
\begin{tikzcd}[column sep=large]
\varphi(X) \arrow[r,"\varphi(u)"] \arrow[dr,"\id_{\varphi(X)}"] & \varphi(\psi\varphi X) = (\varphi\psi)(\varphi X) \arrow[d,"{\begin{array}{l}\wr\ \mathrm{adj}\ \varphi\psi \longrightarrow 1,\ \text{qui est un}\\ \text{isomorphisme car } \psi \text{ est pleinement fidèle}\end{array}}"] \\
 & \varphi(X),
\end{tikzcd}
\]
donc $\varphi(u)$ est bien un isomorphisme.

Reste à prouver (i$''$) $\Longrightarrow$ (ii$'$), i.e. on suppose

\noindent\msp{page 39}%
que le foncteur $X \longrightarrow \Hom(X,K)$ transforme le morphisme d'adjonction particulier
\[
K \xrightarrow{i} \psi\varphi K
\]
en
\[
\begin{aligned}
u \quad &\longrightarrow \quad u\circ i\\
\Hom(\psi\varphi(K),K) \quad &\longrightarrow \quad \Hom(K,K),
\end{aligned}
\]
qui est une bijection. Donc $\id_K$ est dans l'image, soit $p$ un élément du premier membre qui lui donne naissance

\printedpage{25}
% head: A. Grothendieck : Dérivateurs, Ch. V, Édition des Universités de Montpellier II et Paris VII
% foot: Version 9 septembre 2001

\[
\begin{tikzcd}
K \arrow[r,hook,shift left,"i"] & \psi\,\underbrace{\varphi(K)}_{K'} \arrow[l,shift left,"p"]
\end{tikzcd}
\qquad\qquad {\scriptstyle \text{donc } p\circ i = \id_K}
\]
Considérons
\[
ip = q : \psi(K') \longrightarrow \psi(K') \qquad\qquad\qquad (\text{NB } q^2 = q),
\]
comme $\psi$ est pleinement fidèle, il existe un unique
\[
q_0 : K' \longrightarrow K' \qquad\qquad\qquad (K' \overset{\text{déf}}{=} \varphi(K))
\]
tel que l'on ait
\[
q = \psi(q_0).
\]
La relation $pi = \id_K$ donne, en appliquant $\varphi$,
\[
\varphi(p)\varphi(i) = \id,
\]
donc comme $\varphi(i)$ est un isomorphisme (c'est
\[
\begin{tikzcd}[column sep=large]
\varphi(K) \arrow[r,"\sim"] \arrow[dr,"\sim"{sloped,pos=0.45},"\id"{pos=0.6}] & \varphi(\psi\varphi)(K) = (\varphi\psi)(\varphi K) \arrow[d,"\wr"',"\mathrm{adj}"] \\
 & \varphi(K),
\end{tikzcd}
\]
rendant commutatif le triangle ci-contre), $\varphi(p)$ l'est. Donc $\varphi(q) = \varphi(i)\varphi(p)$ l'est, puisque $\varphi(i)$ et $\varphi(i)$ \texttt{[plutôt $\varphi(p)$]} le sont. Mais
\[
\varphi(q) = \varphi\psi(q_0) \quad \underbrace{\xrightarrow{\ \sim\ }}_{\mathrm{adj}\ \varphi\psi\ \xrightarrow{\sim}\ \id} \quad q_0,
\]
donc $q_0$ est un isomorphisme, donc aussi $q = \psi q_0$. Comme $q^2 = q$, cela implique $q = \id$. Donc $i$ et $p$ sont des isomorphismes inverses l'un de l'autre, ce qui établit (ii$''$) \texttt{[plutôt (ii$'$)]} et achève la démonstration.

\dgrblock{dgrb182}{A retract of a local object is local}
\begin{dcomment}{exposition}{p39:et achève la démonstration.}{A retract of a local object is local}
Grothendieck's proof of Proposition~1 identifies the precise amount of surjectivity needed: the identity of $K$ must lift across its unit. This is stronger than a convenient trick. It says that the full subcategory of local objects is closed under retracts, and explains why checking every weak equivalence can be reduced to one particular map.

Let $\varphi:C\rightleftarrows C':\psi$ be an adjunction with fully faithful right adjoint, and write $T=\psi\varphi$ and $\eta:1\to T$ for its idempotent localization monad. An object $K$ is local exactly when $\eta_K$ is invertible. Local objects invert $W_\varphi$ under $\operatorname{Hom}_C(-,K)$ by adjunction. Conversely, if precomposition with $\eta_K$ is surjective at the identity of $K$, choose $p:TK\to K$ with $p\eta_K=1_K$.

Here is a short verification of the remaining step. Since $\varphi\eta_K$ is invertible by the triangle identity, $\varphi p$ is its inverse. Thus $\varphi(\eta_Kp)=1_{\varphi TK}$. The functor $\varphi$ is faithful on the full subcategory of local objects: on that subcategory it is quasi-inverse to $\psi$. Both source and target of $\eta_Kp$ are the local object $TK$, so $\eta_Kp=1_{TK}$. Hence $\eta_K$ and $p$ are inverse isomorphisms. This proves the implication from the weakest condition in the footnote and therefore all the equivalences in the proposition.

The representability corollary is then exact. If $H=H'\varphi$, every $W_\varphi$ is sent to a bijection. When $H\cong\operatorname{Hom}_C(-,K)$, the proposition makes $K$ local, and adjunction identifies $H'$ with $\operatorname{Hom}_{C'}(-,\varphi K)$. In the other direction, a representative $K'$ for $H'$ pulls back to the representative $\psi K'$. Essential surjectivity of $\varphi$ onto $C'$ ensures that the resulting natural isomorphisms are determined by the original factorization. The repeated $X$ in the footnote's first formulation is a printed slip; the object at this step is $K$, as the edition itself corrects.
\end{dcomment}


\textbf{Corollaire :} \textit{Considérons des foncteurs $C^o \xrightarrow{H} \Ens$, et $C'^o \xrightarrow{H'} \Ens$, tels que $H \simeq H'\circ\varphi$. Pour que $H$ soit représentable, il faut et il suffit que $H'$ le soit, et dans ce cas, l'objet $K$ de $C$}

\noindent\msp{page 40}%
\textit{qui représente $H$ satisfait aux conditions équivalentes de la proposition. Si $K$, $K'$ les représentent, on a des isomorphismes canoniques $K' \simeq \varphi K$, $K \simeq \psi K'$.}

Si $H'$ est représentable par $K' \in \Ob C'$, on aura donc
\[
H'\circ\varphi(X) \simeq \Hom_{C'}(\varphi X,K') \overset{\mathrm{adj}}{\simeq} \Hom_C(X,\psi(K'))
\]
\printedpage{26}
% head: A. Grothendieck : Dérivateurs, Ch. V, Édition des Universités de Montpellier II et Paris VII
% foot: Version 9 septembre 2001
donc $H$ est représentable par $K = \psi(K')$, satisfaisant la condition (ii) de la proposition. Inversement, si $H$ est représentable par un objet $K$, alors $K$ satisfait à la condition (i), donc par (ii$''$) on a
\[
\Hom_C(X,K) \xrightarrow{\sim} \Hom_{C'}(\varphi X,\varphi K),
\]
donc le foncteur $H' : C^o \longrightarrow \Ens$ qui factorise $H$ s'explicite comme $X' \longmapsto \Hom_{C'}(X',\varphi K)$, donc est représentable par $\varphi K$. % [sic] C^o -> C'^o, H' is a functor on C'^o (cf. the Corollaire)

Passons aux ind-objets. Les foncteurs $\varphi$, $\psi$ passent aux ind-objets, d'où des foncteurs
\[
\begin{tikzcd}[column sep=large]
\Ind C \arrow[r,shift left,"\Psi = \Ind\psi"] \arrow[r,shift right,"\Phi = \Ind\varphi"'] & \Ind C'.
\end{tikzcd} % [sic] both arrows printed pointing to Ind C' -> the arrow Psi = Ind psi apparently intended from Ind C' to Ind C, psi goes from C' to C
\]
On aura encore la formule d'adjonction pour $X = (X_i)_{i\in I}$ dans $\Ind C$, $X' = (X'_{i'})_{i'\in I'}$ dans $\Ind C'$, car
\[
\begin{aligned}
\Hom(X,\Psi X') &= \varprojlim_i \varinjlim_{i'} \Hom(X_i,\psi(X'_{i'}))\\
&= \varprojlim_i \varinjlim_{i'} \Hom(\varphi X_i,X'_{i'})\\
&= \Hom(\Phi X,X').
\end{aligned}
\]

\noindent\msp{page 41}%
On prouve de même que $\Psi$ est pleinement fidèle, car si $X' = (X'_i)_{i\in I}$, $Y' = (Y'_j)_{j\in J}$ sont des objets de $\Ind C'$, on aura
\[
\begin{tikzcd}[column sep=small]
\Hom(X',Y') \arrow[d,"\Psi_{X',Y'}"] & \simeq \varprojlim\limits_i \varinjlim\limits_j & \Hom_{C'}(X'_i,Y'_j) \arrow[d,"\psi_{X'_i,Y'_j}"] \\
\Hom(\Psi X',\Psi Y') & \simeq \varprojlim\limits_i \varinjlim\limits_j & \Hom_C(\psi(X'_i),\psi(Y'_j)),
\end{tikzcd}
\]
donc $\Psi_{X',Y'}$ est un isomorphisme si les $\psi_{X'_i,Y'_j}$ le sont.

Donc la proposition 1 peut s'appliquer à la situation $\Phi$, $\Psi$ entre ind-objets. Je vais l'expliciter sous la forme la plus commode :

\textbf{Proposition 2.} \textit{Soit $K = (K_i)_{i\in I}$ un ind-objet de $C$, et soit $H$ le foncteur}
\[
\left\{
\begin{aligned}
H : X \quad &\longmapsto \quad \underbrace{\Hom_C(X,K)}_{K(X)} \overset{\text{déf}}{=} \varinjlim_i \Hom_C(X,K_i)\\
C^o \quad &\longrightarrow \quad \Ens
\end{aligned}
\right.
\]
\textit{qu'il ind-représente. Conditions équivalentes}

\begin{enumerate}[label={}]
\item[(i)] \textit{Le foncteur $H$ se factorise (à isomorphisme près) par $C'$ via $\varphi$, i.e. transforme les $u \in W_\varphi$ en isomorphismes.}
\printedpage{27}
% head: A. Grothendieck : Dérivateurs, Ch. V, Édition des Universités de Montpellier II et Paris VII
% foot: Version 9 septembre 2001
\item[(i$'$)] \textit{Même condition pour les flèches d'adjonction $X \longrightarrow \psi\varphi(X)$ (lesquelles sont dans $W$).}
\item[(i$''$)] \textit{Même condition pour la flèche d'adjonction $K \longrightarrow \psi\varphi(K)$.}
\item[(ii)] \textit{$K$ appartient à l'image essentielle de} $\Ind C'$ \textit{par $\Psi$.}
\item[(ii$'$)] $K \longrightarrow \Psi\Phi(K)$ \textit{est un} isomorphisme \textit{dans} $\Ind C$.) % [sic] ".)" -> ".", closing parenthesis without opening parenthesis
\item[(ii$''$)] \textit{L'homomorphisme fonctoriel}
\[
K(X) = \Hom_C(X,K) \longrightarrow \Hom_{C'}(\varphi(X),\varphi(K))
\]
\texttt{[plutôt} $\Hom_{\mathrm{Ind}(C)}(X,K) \longrightarrow \Hom_{\mathrm{Ind}(C')}(\Phi(X),\Phi(K))$\texttt{, en identifiant l'objet} $X$ \texttt{de} $C$ \texttt{à son image dans} $\mathrm{Ind}(C)$\texttt{]} \textit{est un isomorphisme.}
\end{enumerate}

\textsc{Démonstration.} L'équivalence de (ii), (ii$'$), (ii$''$) est la proposition 1 appliquée à $\Phi$ et $\Psi$, à cela près qu'il faut noter que (ii$''$) reste évidemment valable si on y remplace $X$ par un ind-objet $(X_j)_{j\in J}$ quelconque. On a

\noindent\msp{page 42}%
\[
\text{(ii$''$)} \Longrightarrow \text{(i)} \Longrightarrow \text{(i$'$)}
\]
tautologiquement, il reste à prouver que (i$'$) implique (i$''$) et (i$''$) (ii$'$). On ne peut pas appliquer directement la proposition 1, car l'hypothèse (i$'$) ici ne porte que sur les objets de $\Ind C$ qui proviennent de $C$, et il n'est pas évident à priori qu'elle passe aux objets quelconques. Mais cela l'est dans le cas d'un homomorphisme de ind-objets $(X_i)$, $(Y_i)$ sur le même ensemble d'indices, donné par un homomorphisme ordinaire (terme à terme) de systèmes inductifs sur $I$, $X_i \longrightarrow Y_i$: si les $X_i \longrightarrow Y_i$ sont dans $W$, l'hypothèse (i) implique que $H(Y) \longrightarrow H(X)$ est un isomorphisme, et de même pour (i$'$) pour $\underbrace{X}_{=(X_i)_i} \;\longrightarrow\; \underbrace{\Psi\Phi(X)}_{=(\psi\varphi(X_i))_i}$\,, où l'homomorphisme d'adjonction est donné terme à terme par les homomorphismes d'adjonction $X_i \longrightarrow \psi\varphi(X_i)$. Donc l'hypothèse (i$'$) ici implique l'hypothèse (i$'$) de la proposition 1 pour $\Phi$ et $\Psi$, donc aussi (ii$'$) (par la proposition 1), et on gagne. % [sic] "et (i'') (ii')" -> "et (i'') => (ii')", the implication sign or verb is missing

\dgrblock{dgrb183}{Why the argument passes to ind-objects}
\begin{dcomment}{exposition}{p42:de la proposition 1 pour}{Why the argument passes to ind-objects}
Grothendieck is careful at the one point where the extension to ind-objects is not a formal substitution. The hypothesis tests units on objects of $C$, whereas Proposition~1, applied to $\operatorname{Ind}(C)$, would test units on all ind-objects. The missing bridge is a limit, and it needs no exactness assumption.

Put $\Phi=\operatorname{Ind}(\varphi)$ and $\Psi=\operatorname{Ind}(\psi)$. For filtered systems $X=(X_i)$ and $Y'=(Y'_j)$,
\[
 \operatorname{Hom}(X,\Psi Y')
  =\varprojlim_i\varinjlim_j\operatorname{Hom}_C(X_i,\psi Y'_j)
  \cong\varprojlim_i\varinjlim_j
        \operatorname{Hom}_{C'}(\varphi X_i,Y'_j).
\]
This proves $\Phi\dashv\Psi$. Applying the same formula to two systems from $C'$ proves that $\Psi$ is fully faithful. The upper arrow in the edition's displayed pair of ind-functors points the wrong way; these types determine its direction.

For a fixed ind-object $K$, write $\widehat H(X)=\operatorname{Hom}_{\operatorname{Ind}(C)}(X,K)$. Its restriction to $C$ is the functor $H$ of the proposition. Suppose each unit $X_i\to\psi\varphi X_i$ induces a bijection under $H$. The unit of an arbitrary system $X$ is represented term by term by these maps. Therefore
\[
 \widehat H(\Psi\Phi X)=\varprojlim_i H(\psi\varphi X_i)
   \xrightarrow{\sim}\varprojlim_i H(X_i)=\widehat H(X).
\]
A limit of a natural diagram of bijections is a bijection: its inverse is obtained by taking the componentwise inverses, which remain compatible. There is no interchange of a filtered colimit with an inverse limit here. Proposition~1 now applies in $\operatorname{Ind}(C)$ and gives $K\cong\Psi\Phi K$.

Condition (i$''$) needs this extension $\widehat H$ to have a meaning, because $K$ is not in general an object of $C$. The text uses $\psi\varphi K$ as shorthand for $\Psi\Phi K$; the corrected statement is a test at the unit of an ind-object. Likewise condition (ii$''$), when tested only on $X\in C$, extends to arbitrary ind-objects by the displayed limit formula. These are not extra hypotheses or repairs to the argument. They make explicit the conventions under which its proof is complete.
\end{dcomment}


Revenons à la situation où $C = \Cathot$, $C' = \Hot$, $\varphi$ le foncteur \texttt{[de]} localisation. Si $K \in \Ob\Cathot$, alors $\psi\varphi(K)$ s'appelle \textit{l'enveloppe de Thomason-Kan} de $K$ - si $K$ est une catégorie, son enveloppe $\bar{K}$ et la flèche $K \longrightarrow \bar{K}$ ne sont définis que dans la catégorie Cat\,hot, i.e. $\bar{K}$ n'est défini qu'à `homotopisme près'. % reading resolved: no underline is printed under psi phi(K).

\noindent\msp{page 43}%
Mais on peut toujours \textit{choisir} un objet $\bar{K}$ de Cat, et une flèche
\[
i : K \longrightarrow \bar{K} \qquad\qquad\qquad \textit{dans } \Cat,
\]
\printedpage{28}
% head: A. Grothendieck : Dérivateurs, Ch. V, Édition des Universités de Montpellier II et Paris VII
% foot: Version 9 septembre 2001
pas seulement dans Cat\,hot, de telle façon de plus que $i$ soit une `\textit{cofibration} de Thomason' (et en particulier une immersion ouverte?). Il en résulte que si on a un système inductif
\[
K = (K_n)_{n\in\mathbf{N}}
\]
indexé par $\mathbf{N}$, et si on choisit pour chaque $n$ une enveloppe de Thomason-Kan
\[
K_n \overset{\alpha_n}{\hookrightarrow} \bar{K}_n
\]
de cette façon, on peut trouver des foncteurs de transition
\[
\bar{K}_n \xrightarrow{\bar{i}_n} \bar{K}_{n+1}
\]
tels que les carrés
\[
\begin{tikzcd}
K_n \arrow[r,hook,"\alpha_n"] \arrow[d,"i_n"'] & \bar{K}_n \arrow[d,"\bar{i}_n"] \\
K_{n+1} \arrow[r,hook,"\alpha_{n+1}"] & \bar{K}_{n+1}
\end{tikzcd}
\]
soient commutatifs, de telle sorte que l'on ait donc un homomorphisme de \textit{systèmes inductifs de} Cat
\[
K \longrightarrow \bar{K}
\]
et a fortiori un homomorphisme de ind-objets de Cat. Ceci posé, les conditions (i), (ii$'$) de la proposition 2 montrent que le foncteur ind-représenté par $K$ \textit{dans} $C = \Cathot$ transforme quasi-isomorphismes

\dgrblock{dgrb184}{What Thomason replacement supplies}
\begin{dcomment}{later}{p43:de cette façon, on peut trouver des foncteurs de transition}{What Thomason replacement supplies}
Grothendieck needs replacements which can be chosen together with the transition maps, so that he still has a sequence of actual categories. This part of the argument has a direct model-category realization. Given a sequence $K_n\to K_{n+1}$, choose trivial cofibrations $\alpha_n:K_n\to R_n$ with $R_n$ Thomason fibrant. The square with top map $\alpha_n$, bottom object $R_{n+1}$ and right map $R_{n+1}\to *$ admits a lift. It gives $R_n\to R_{n+1}$ extending $\alpha_{n+1}i_n$. Since the indexing category is the freely generated chain, no relations beyond these adjacent squares need to be imposed.

The parenthetical question about open immersions also has a precise answer. Thomason cofibrations are sieves, hence full inclusions with no arrows from outside to inside. A technical correction matters here: Thomason originally asserted that retracts of Dwyer maps remain Dwyer maps. \citet{Cisinski1999Dwyer} gives a counterexample and proves that every cofibration is instead a pseudo-Dwyer map; the corrected assertion retains the sieve property and the proper model structure. This repairs that specific later proof, not a gap in Grothendieck's formal Propositions~1--2.

One can also replace any sequence of fibrant objects by a sequence of fibrant objects with cofibration transitions, up to a levelwise weak equivalence. Start with $T_0=R_0$. If $g_n:T_n\to R_n$ has been constructed, factor the composite $T_n\to R_n\to R_{n+1}$ as
\[
 T_n\lhook\joinrel\longrightarrow T_{n+1}
       \xrightarrow[\sim]{g_{n+1}}R_{n+1},
\]
where the first map is a cofibration and the second a trivial fibration. Then $T_{n+1}$ is fibrant and the $g_n$ form a strict map of sequences. The transition inclusions are fully faithful. This proves a useful model-theoretic version of the rearrangement proposed on manuscript p.~46. It gives a levelwise weak equivalence and therefore an isomorphism of ind-objects in $\Hot$. Calling it an isomorphism in $\operatorname{Ind}(\Cathot)$ requires a further comparison; that conclusion is not supplied by the factorization axiom.
\end{dcomment}


\noindent\msp{page 44}%
ordinaires (équivalences faibles) en bijections, \textit{si et seulement si} l'homomorphisme précédent est un isomorphisme de ind-objets de Ind\,Cat\,hot.

Ce sera donc cela la condition naturelle à imposer aux ind-catégories de chemins
\[
\Chh_{\Sigma}(\mathcal{M};a,b) \qquad\qquad\qquad (a,b \in \mathrm{ob}\,\mathcal{M})
\]
provenant d'une `catégorie de modèles' $\mathcal{M}$, $\Sigma$.

Bien sûr, disposant d'un système inductif explicite de catégories (pas seulement d'objets de Cat\,hot)
\[
K_n = \Ch_{\Sigma,n}(\mathcal{M};a,b),
\]
il est tentant de passer à la limite inductive
\[
K_\infty = \varinjlim_n \Ch_{\Sigma,n}(\mathcal{M};a,b) \overset{\text{déf}}{=} \Ch_{\Sigma}(\mathcal{M};a,b),
\]
et ceci d'autant plus que les foncteurs de transition sont des monomorphismes et pleinement fidèles, de sorte que les $K_n$ s'identifient à une suite de sous-catégories pleines de

% ---- from Der-ChV-234_p031-034.tex
\printedpage{29}
% head: A. Grothendieck : Dérivateurs, Ch. V, Édition des Universités de Montpellier II et Paris VII
% foot: Version 9 septembre 2001
$K_\infty$, dont la réunion est $K_\infty$. Mais d'autre part je ne sais donner aucune interprétation convaincante de $K_\infty$, comme de celle du ind-objet de Cat\,hot defini par $\uuK = (K_n)_{n\in\mathbf{N}}$, qui ind-représente le foncteur remarquable % [sic] defini -> défini, missing accent
\[
H = H_{\mathcal{M};x,y} : \Cathot^{o} \longrightarrow \Ens
\]
factorisant le foncteur
\msp{page 45}
\[
X \longmapsto H_{\mathcal{M};x,y}(X) \defeq \Hom_{\mathcal{M}(X)\Sigma_X^{-1}}(x_X, y_X),
\]
par la formule donc
\[
H(X) = \varinjlim_{n} \underbrace{\pi_0\uHom(X, K_n)}_{\Hom_{\Cathot}(X,K_n)} .
\]
Je ne sais donner d'interprétation de ce type ni de la catégorie $K_\infty$ elle-même, ni de son image dans Cat\,hot, ni de son image dans Hot. Il est vrai que si je restreins le foncteur $H$ à la sous-catégorie pleine de (Cat) formée des catégories \textit{finies} (Ob\,$X$ et Fl\,$X$ finis), alors on aura
\[
\uHom^{!}(X, K_\infty) = \uHom(X, K_\infty),
\]
et l'isomorphisme précédant s'interprète comme % [sic] précédant -> précédent, adjective intended
\[
\Hom_{\mathcal{M}(X)\Sigma_X^{-1}}(a_X, b_Y) \defeq H(X) \simeq \pi_0\uHom(X, K_\infty) \defeq \Hom_{\Cathot}(X, K_\infty).
\]
\texttt{[plutôt} $(a_X, b_X)$\texttt{]} Cela prouve du moins que la connaissance de $K_\infty$ est utile. Mais cette formule, pour $X$ restreint aux catégories \textit{finies,} ne définira pas $K_\infty$ comme objet de Cat\,hot, à isomorphisme unique près dans cette catégorie - \textit{sauf} dans le cas où $K_\infty$ est elle-même homotope à une catégorie
\msp{page 46}
finie. (Je ne me rends pas bien compte si ça arrive très souvent \ldots\ Mais on ne peut fonder une théorie générale des dérivateurs et des catégories de modèles `dérivables' sur une hypothèse de ce genre.)

\dgrblock{dgrb185}{Finitely many objects already give a uniform stage}
\begin{dcomment}{development}{p45:formée des catégories}{Finitely many objects already give a uniform stage}
Grothendieck uses finite categories to remove the exclamation mark. The proof is controlled by the number of objects, because fullness has already supplied every arrow between objects that have appeared. We record the resulting refinement to make this source of finiteness explicit.

\textbf{Finite-object factorization.}\ours\quad
Let $K_0\subseteq K_1\subseteq\cdots$ be full subcategories, let $K_\infty$ be their union, and let $X$ have finitely many objects, with no restriction on its sets of arrows. Then the canonical functor
\[
 \varinjlim_n\operatorname{Fun}(X,K_n)
       \longrightarrow\operatorname{Fun}(X,K_\infty)
\]
is an isomorphism of categories. In particular, the path-localization formula on manuscript p.~45 holds for such $X$.

\textit{Proof.} Given $F:X\to K_\infty$, choose a stage containing each of its finitely many object-values, and take the largest of these stages. Every arrow of $X$ has endpoints whose images lie in that stage, so fullness supplies its image arrow there. The composition and identity equations hold because inclusion is faithful. Thus $F$ factors through this one stage. If $F,G$ are two functors represented at possibly different stages, put them in a common stage. All components of every transformation $F\Rightarrow G$ then belong to that stage by fullness; its naturality equations are unchanged. Conversely, inclusions of stages introduce no identifications between different transformations. This proves both bijectivity on objects and bijectivity on every Hom-set.\hfill$\square$

There is a version invariant under equivalence of the source. If $X$ has only finitely many isomorphism classes of objects, choose a finite-object skeleton $S$ and an equivalence $S\to X$. Precomposition gives equivalences of functor categories into every $K_n$ and into $K_\infty$. Passing to $\pi_0$ therefore proves
\[
 \varinjlim_n\pi_0\operatorname{Fun}(X,K_n)
       \cong\pi_0\operatorname{Fun}(X,K_\infty).
\]
The strict equality of bounded and unrestricted functor categories need not hold for this larger class: an actual functor may send infinitely many isomorphic source objects to different target objects appearing at unbounded stages. The statement on components is the appropriate invariant one.

The qualification at the end of manuscript p.~45 also needs its domain of comparison. Ordinary Yoneda gives uniqueness among objects of $\Cathot$ homotopy equivalent to finite categories. It does not give uniqueness among all objects merely because one candidate is finite. The same elementary fence which will make the later length obstruction visible already separates these assertions.

\textbf{Failure of unrestricted uniqueness from finite tests.}\ours\quad
Let $B$ be the poset $0\to1\leftarrow2\to3\leftarrow\cdots$, and let $B_n$ be its full subcategory on $\{0,\ldots,n\}$. The restrictions of $\operatorname{Hom}_{\Cathot}(-,B)$ and $\operatorname{Hom}_{\Cathot}(-,*)$ to finite categories are naturally isomorphic, although $B$ and $*$ are not isomorphic in $\Cathot$.

\textit{Proof.} The maps $r_j(i)=\min(i,j)$, for $0\leq j\leq n$, are endofunctors of $B_n$: they preserve every fence arrow below $j$ and collapse every arrow above it to an identity. For even $j$ there is a transformation $r_j\Rightarrow r_{j+1}$; for odd $j$ there is a transformation in the reverse direction. Thus $r_0$, the constant functor at $0$, and $r_n=1_{B_n}$ are joined by a zigzag. Postcomposition contracts $\operatorname{Fun}(X,B_n)$ for every $X$. When $X$ is finite, the preceding factorization lemma places each functor $X\to B$ in some $B_n$. It is therefore homotopic to the constant functor at $0$, and $\pi_0\operatorname{Fun}(X,B)$ is a singleton. These singleton identifications are natural in $X$.

If $1_B$ were joined to the constant functor at $0$ by $L$ transformations, evaluation at $L+1$ would give a zigzag of $L$ arrows connecting $L+1$ to $0$. Each fence arrow changes the index by at most one, which is impossible. Hence $B$ is not categorically contractible.\hfill$\square$

The nerve of $B$ realizes to an infinite ray and is contractible. This counterexample concerns $\Cathot$; the last-page obstruction in $\Hot$ will require a different example. It gives a precise reason to restrict the finite-object exception to comparisons within the finite subcategory, if that is the intended reading of the parenthesis.

We checked the general filtered-colimit and functor-category framework in \citet[Chapters~1 and 7]{Riehl2014} and the specific argument in the chapter. We did not find this exact refinement stated there. It is an elementary continuation of the printed proof, with no claim of priority. Its relevance is practical: a one-object category with infinitely many endomorphisms still satisfies the lemma, even though it is not a finite category in the author's stated sense.
\end{dcomment}


Mais supposons que le foncteur $H$ envisagé se factorise même par Hot, i.e. qu'il transforme quasi-isomorphismes en bijections. On a vu que cela équivaut à la condition qu'il puisse se représenter par un système inductif sur $\mathbf{N}$ formé de catégories de Thomason-Kan. Il doit être facile de plus, si on y tient, de s'arranger pour que les foncteurs de transition soient des monomorphismes et pleinement fidèles, pour avoir une $\varinjlim$ sympathique. Ce sera bien sûr une catégorie de Thomason-Kan. Cette catégorie est-elle déterminée par le foncteur $H : \Cathot^{o} \longrightarrow \Ens$ à isomorphisme unique près dans Cat\,hot? Bien sûr, il n'est pas question qu'elle représente le foncteur $H$, sauf \texttt{[si]} celui-ci est bel et bien représentable, i.e. si le ind-objet $\uuK$ de Cat\,hot est essentiellement constant (ce qui ne doit pratiquement jamais être le cas). Notons déjà
\printedpage{30}
% head: A. Grothendieck : Dérivateurs, Ch. V, Édition des Universités de Montpellier II et Paris VII
% foot: Version 9 septembre 2001
\msp{page 47}
que si $u$ est un homomorphimse de ind-objets \textit{dans} Cat, ou, ce qui revient maintenant au même, un homomorphisme dans Cat % [sic] homomorphimse -> homomorphisme, transposed letters
\[
K_\infty \xrightarrow{u_\infty} K'_\infty
\]
tel que $\forall n$, $u_\infty(K_n) \subseteq K'_{n'}$ pour $n' = \tau(n)$ convenable, et si on suppose que l'homomorphisme de ind-objets induit un isomorphisme pour les foncteurs ind-représentables \textit{sur} Cat\,hot (par exemple, si $\uuK$ et $\uuK'$ ind-représentent le même foncteur $H$, et l'homomorphisme de ind-objets est celui qui correspond à $\mathrm{id}_H$), alors $u_\infty$ est forcément un isomorphisme \textit{dans} Hot, donc aussi (sauf erreur) dans Cat\,hot puisque $K_\infty$, $K'_\infty$ sont des catégories de Thomason-Kan. Mais ce qui n'est pas évident, c'est l'existence d'un homomorphisme $u_\infty$ comme \texttt{[ci-]}dessus, et s'il existe, qu'il soit unique (ou du moins canonique?) comme flèche de Cat\,hot.

\dgrblock{dgrb186}{Fibrancy passes to the union}
\begin{dcomment}{exposition}{p46:Ce sera bien sûr une catégorie de Thomason-Kan.}{Fibrancy passes to the union}
Grothendieck expects the union of his replacement categories to remain a good representative of a homotopy type. For Thomason fibrancy this expectation is correct and has a finite combinatorial proof. A category $K$ is Thomason fibrant when $\operatorname{Ex}^2NK$ is a Kan complex. For a filtered system $(K_i)$, the nerve commutes with its colimit: an $n$-simplex is a finite composable string, and both the string and the finitely many equations expressing a morphism of such strings occur at one stage. Moreover
\[
 (\operatorname{Ex}^2 S)_n
   =\operatorname{Hom}_{\mathbf{sSet}}(\operatorname{Sd}^2\Delta^n,S)
\]
commutes with filtered colimits, since $\operatorname{Sd}^2\Delta^n$ is a finite simplicial set. Finally, a horn in a filtered colimit of Kan complexes lies in one stage, and has a filler at that stage. Consequently
\[
 \operatorname{Ex}^2N(\varinjlim_i K_i)
       \cong\varinjlim_i\operatorname{Ex}^2NK_i
\]
is Kan whenever all the terms are Kan. This proof even dispenses with fully faithful transition maps.

Fibrancy does not answer the next uniqueness question. There are three different conclusions available from three different kinds of information. A presheaf determines its representing ind-object in $\operatorname{Ind}(\Cathot)$. A coherent map of strict diagrams which is a weak equivalence at every stage gives a weak equivalence on the appropriate homotopy colimits. A specified categorical homotopy between maps of unions requires actual natural transformations, with a finite common length. These conclusions are not interchangeable.

In particular, an isomorphism in the ind-category of a homotopy category is only eventual equality of homotopy classes at each stage. It does not come with compatible choices of the homotopies. The claim on manuscript p.~47 that a corresponding union map must be a weak equivalence is false, even for full inclusions of finite fibrant stages with fibrant union; the example is in the commentary at ms.~49. The verified model-category replacement construction in the preceding block avoids this difficulty by giving actual coherent maps and levelwise weak equivalences. The raw categorical-homotopy strictification questions remain separately recorded in the audit.
\end{dcomment}


Pour l'existence d'un $u_\infty$, on est ramené à prouver que si on a un diagramme à flèches pleines
\[
\begin{tikzcd}
K_n \arrow[r, "i_n", "\text{cofib.}"'] \arrow[d, "u_n"] & K_{n+1} \arrow[dr, dashed] & \\
K'_{\tau(n)} \arrow[rr, "\text{transition}"'] & & K'_N ,
\end{tikzcd}
\]
et si une flèche $u_{n+1}$ est donnée `en pointillés' \textit{dans} Cat\,hot, de façon à rendre le carré commutatif,
\msp{page 48}
alors on peut la réaliser par une flèche \textit{dans} Cat qui rende le diagramme commutatif pour de bon (dans Cat). On est ramené au même problème dans une situation
\[
\begin{tikzcd}
X \arrow[r, hook, shift right, "i"'] & Y , \arrow[l, dashed, shift right, "p"']
\end{tikzcd}
\]
où $X$ est une catégorie de Thomason, $i$ une cofibration de Thomason, $p$ une rétraction de $Y$ dans $X$ \textit{dans} Cat\,hot, on veut la réaliser par une vraie rétraction (\footnote[7]{\texttt{[un ?\ \ dans la marge]}}).

En admettant que ça soit le cas, donc que tout homomorphisme des ind-objets $\uuK$, $\uuK'$ \textit{dans} Cat\,hot se représente par un homomorphisme entre les ind-objets \textit{dans} Cat, donc aussi par un foncteur $K_\infty \longrightarrow K'_\infty$, reste la question si l'homomorphisme correspondant dans Cat\,hot est unique. Donc si on a deux foncteurs $f$, $g$ qui induisent le même homomorphisme sur les foncteurs ind-représentés sur Cat\,hot, si $f$ et $g$ sont homotopes. Donc

\textbf{Question}. Deux foncteurs $f, g : K_\infty \longrightarrow K'_\infty$ tels que pour tout $n \in \mathbf{N}$, $f_n$, $g_n$ qui s'en déduisent par restriction à $K_n$ sont homotopes (et même homotopes en tant que foncteurs de $K_n$ dans un $K'_N$ convenable), sont-ils homotopes? (\footnote[8]{$K$, $K'$ catégories de Thomason, les $K_n$, $K'_n$ aussi.})

\printedpage{31}
% head: A. Grothendieck : Dérivateurs, Ch. V, Édition des Universités de Montpellier II et Paris VII
% foot: Version 9 septembre 2001
\msp{page 49}
Je présume que non - qu'il peut arriver que la longueur minimum requise $l(n)$ pour un chemin dans $\uHom(K_n, K'_{\tau(n)})$ reliant $f_n$ et $g_n$, tende vers l'infini si $n \longrightarrow \infty$. L'hypothèse que les $K'_n$ sont des catégories de Thomason y change-t-il quelque chose (deux objets homotopes dans $\uHom(X, K')$, si $K'$ est de Thomason, sont-ils reliés par une flèche?)

\dgrblock{dgrb187}{A complete instance of the length obstruction}
\begin{dcomment}{development}{p49:Je présume que non - qu'il peut arriver}{A complete instance of the length obstruction}
Grothendieck identifies a concrete mechanism by which stagewise homotopies could fail to assemble: their lengths may become unbounded. The following example realizes that mechanism with full inclusions at every stage. It deliberately leaves the additional target-fibrancy condition aside, so that the obstruction itself can be seen without a model-category comparison.

\textbf{Unbounded categorical homotopies.}\ours\quad
There are full filtrations $A_0\subset A_1\subset\cdots$ and $B_0\subset B_1\subset\cdots$, and functors $f,g:A_\infty\to B_\infty$, such that $f|_{A_n}$ and $g|_{A_n}$ are homotopic in $\operatorname{Fun}(A_n,B_n)$ for every $n$, but $f$ and $g$ are not homotopic in $\operatorname{Fun}(A_\infty,B_\infty)$.

\textit{Proof.} Take $A_n$ to be discrete on $\{0,\ldots,n\}$. Take $B_n$ to be the full subcategory on the same vertices of the infinite alternating fence
\[
 0\longrightarrow1\longleftarrow2\longrightarrow3
       \longleftarrow4\longrightarrow\cdots.
\]
There are no nonidentity composites, so this is a poset. Define $f(i)=0$ and $g(i)=i$. For $0\leq j\leq n$, put $h_j(i)=\min(i,j)$ on $A_n$. Since the source is discrete, these assignments are functors. If $j$ is even there is a transformation $h_j\Rightarrow h_{j+1}$, and if $j$ is odd there is a transformation $h_{j+1}\Rightarrow h_j$: each component is either the identity or the indicated fence arrow. Hence
\[
 h_0=f|_{A_n},\quad h_1,\quad\ldots,\quad h_n=g|_{A_n}
\]
is a categorical homotopy of length $n$, contained in $B_n$.

Suppose a zigzag of $L$ natural transformations connected $f$ and $g$ globally. Evaluate it at the object $L+1$ of $A_\infty$. Each nonidentity arrow in $B_\infty$ changes the vertex by exactly one, so a zigzag of $L$ arrows cannot connect $0$ to $L+1$. This contradiction proves the assertion; it also proves that the displayed stagewise length $n$ is minimal.\hfill$\square$

Both filtrations consist of full subcategories, and the source stages are discrete, hence Thomason fibrant. We have not shown that the target stages or their union are Thomason fibrant; this construction therefore does not answer the stronger question with that additional hypothesis. The next commentary does. It proves exactly the length mechanism suggested in the first sentence of manuscript p.~49. The argument was developed directly from that sentence. We checked the surrounding categorical-homotopy and model-category references cited in this section and did not find this exact fence formulation; it is presented as an elementary illustration, not a priority claim.
\end{dcomment}


\dgrblock{dgrb188}{The same ind-morphism, two different maps of the unions}
\begin{dcomment}{development}{p49:Je présume que non - qu'il peut arriver}{The same ind-morphism, two different maps of the unions}
With Thomason fibrancy imposed on every stage and on the unions, the questions of ms.~47--49 still have negative answers. The example is a strict form of the symmetric-group shift used at ms.~52.

\textbf{Proposition.}\ours
\textit{There are a sequence $T_0\subset T_1\subset\cdots$ of full inclusions of finite categories, each $T_n$ Thomason fibrant and isomorphic in $\mathrm{Cat}\,\mathrm{hot}$ to a group, with Thomason-fibrant union $T$, and a functor $s:T\to T$ with $s(T_n)\subseteq T_{n+1}$, such that:}
\begin{enumerate}[label=(\alph*)]
\item \textit{for every $n$, the functor $s|_{T_n}$ and the inclusion $T_n\to T_{n+1}$ are joined by a zigzag of three natural transformations, so $s$ and $\mathrm{id}_T$ induce the same endomorphism of the ind-object $(T_n)$ of $\mathrm{Cat}\,\mathrm{hot}$;}
\item \textit{$Bs$ is not a weak equivalence; in particular $s$ and $\mathrm{id}_T$ are not homotopic;}
\item \textit{there is a weak equivalence $r:T\to\mathsf BG$ onto a group, between Thomason-fibrant categories, that is not an isomorphism in $\mathrm{Cat}\,\mathrm{hot}$.}
\end{enumerate}

\textit{Proof.}
Let $G_j$ be the group of permutations of $\{0,\dots,j\}$, with $G_j\subset G_{j+1}$ by fixing $j+1$, and $G=\bigcup G_j$. Let $T_n$ have the objects $0,\dots,n$, with $\operatorname{Hom}(i,j)=G_j$ if $i\leq j$ and empty otherwise, and composition by multiplication in the group of the target. The inclusions $T_n\subset T_{n+1}$ are full. The functor $r_n:T_n\to\mathsf BG_n$ that forgets the objects and the functor $j_n:\mathsf BG_n\to T_n$ onto the object $n$ satisfy $r_nj_n=\mathrm{id}$, and the arrows $i\to n$ labelled by the unit form a natural transformation $\mathrm{id}\Rightarrow j_nr_n$. So $T_n\cong\mathsf BG_n$ in $\mathrm{Cat}\,\mathrm{hot}$.

Fibrancy. A category is Thomason fibrant when it has the right lifting property for the functors $c\,\mathrm{Sd}^2\Lambda\to c\,\mathrm{Sd}^2\Delta$ associated with the horn inclusions, and these are sieves \citep[Proposition~4.2]{Thomason1980}. Let $F:A\to T_n$ be given on such a sieve $A\subset B$. The functor $r_nF$ extends to $B$, because the nerve of a groupoid is a Kan complex. Extend $F$ by sending every object of $B$ outside $A$ to $n$ and every arrow that meets the complement of $A$ to the arrow with the label given by the extension of $r_nF$. No arrow leads from the complement into $A$, and $\operatorname{Hom}(i,n)=G_n$ for all $i$, so this is defined, and it is a functor because composition in $T_n$ multiplies labels. So $T_n$ is fibrant, and the union $T$ is fibrant by the preceding commentary on ms.~46.

The shift. For $g\in G_j$ let $\sigma(g)\in G_{j+1}$ fix $0$ and act on $\{1,\dots,j+1\}$ as $g$ does on $\{0,\dots,j\}$; let $s:T\to T$ send $i$ to $i+1$ and the arrow labelled $g$ to the arrow labelled $\sigma(g)$. On $T_n$ both $s$ and the inclusion into $T_{n+1}$ map, by the arrows labelled by the unit, to the functors with the same labels and constant value $n+1$ on objects, and these two are isomorphic by the cycle $c=(0\,1\,\cdots\,n+1)$, since $cgc^{-1}=\sigma(g)$ for $g\in G_n$. This is (a).

The functors $r_n$ are compatible and give $r:T\to\mathsf BG$; the nerve and homotopy groups commute with the union, so $Br$ is a weak equivalence and $BT$ is a $K(G,1)$. Under this identification $s$ induces $\sigma:G\to G$ on $\pi_1$, which is not surjective (its image fixes $0$). This is (b). For (c): if $r$ were invertible in $\mathrm{Cat}\,\mathrm{hot}$, there would be a functor $j:\mathsf BG\to T$ with $rj$ joined to the identity by natural transformations, that is, with $rj:G\to G$ an inner automorphism. But $j$ sends the object to some $i$ and $G$ into $\operatorname{Hom}(i,i)=G_i$, so the image of $rj$ lies in $G_i\neq G$.\hfill$\square$

Consequences for these pages. The assertion of ms.~47, that a morphism of such systems inducing an isomorphism of ind-objects induces a weak equivalence of the unions, is false: $s$ induces the identity of the ind-object. The question of ms.~48, whether two maps of unions that are homotopic on every stage inside a later stage are homotopic, has the answer no, with homotopies of length three at every stage; what fails is their compatibility, not the growth of their length as in the fence example above. The union of fibrant stages is fibrant, but the union $T$ of stages each isomorphic in $\mathrm{Cat}\,\mathrm{hot}$ to a groupoid is not isomorphic to one, by (c). And a weak equivalence between Thomason-fibrant categories need not be an isomorphism in $\mathrm{Cat}\,\mathrm{hot}$ (ms.~47, ``sauf erreur''): fibrancy is not the envelope property.
\end{dcomment}


\dgrblock{dgrb189}{Fibrancy does not reduce every homotopy to one arrow}
\begin{dcomment}{examples}{p49:y change-t-il quelque chose}{Fibrancy does not reduce every homotopy to one arrow}
The parenthetical question asks a sharper thing than a uniform bound: for a Thomason fibrant target, must two homotopic functors already be joined by a single natural transformation? A three-object category gives a negative answer, and the fibrancy can be checked directly against the generating horns.

Every category $C$ with a terminal object $t$ is Thomason fibrant. Indeed, \citet[Propositions~2.5 and 4.2]{Thomason1980} identify the lifting test with inclusions
\[
 c\operatorname{Sd}^2\Lambda^k[n]
     \lhook\joinrel\longrightarrow c\operatorname{Sd}^2\Delta[n],
\]
which are sieves. More generally, if $A\subseteq D$ is a sieve and $F:A\to C$ is a functor, extend $F$ by assigning $t$ to every object outside $A$. There are no arrows from an outside object to an object of $A$. Send arrows between outside objects to $1_t$, and each arrow from inside to outside to the unique map into $t$. Composition is respected by uniqueness of maps to the terminal object. This constructs the required lift against every displayed horn, so $C\to *$ is a Thomason fibration.

Now take $C$ to be $0\to1\leftarrow2$, whose terminal object is $1$, and take $X=*$. The functors selecting $0$ and $2$ lie in the same component of $\operatorname{Fun}(*,C)=C$, connected by the two-arrow zigzag through $1$. There is no arrow $0\to2$ or $2\to0$, so there is no single transformation in either direction between them. The target is fibrant by the preceding proof.

This answers the exact one-arrow question. It does not say that minimal lengths can grow without bound for a fixed fibrant target, and it does not establish a counterexample to the preceding question with fibrant full filtrations. A claim of bounded shortening, or of compatibility of stagewise homotopies, has different quantifiers and needs a separate argument.
\end{dcomment}


Bien sûr, on peut prétendre que les catégories $\uCh_{\Sigma,n}(\mathcal{M}; a, b)$ ont été définies `canoniquement', et leur limite inductive $\uCh_{\Sigma}(\mathcal{M}; a, b)$ aussi. Mais la définition était assez arbitraire et abracadabrante - je suis convaincu que dix mathématiciens, se proposant de développer cette même idée des catégories de chemins, dans le même but que moi, arriveraient à dix définitions différentes! Pourtant, il n'est pas exclu que l'on finisse par trouver un système transitif d'isomorphismes dans Cat\,hot, entre les catégories de chemins `infinis' (i.e. de longueur non précisée) $K_\infty = K_\infty(\mathcal{M}; a, b)$ que nous aurons construites, dans le but
\msp{page 50}
principal (disons) de représenter le foncteur $H$ \textit{sur les catégories finies.} Et que les isomorphismes obtenus seront, bien entendu, des homomorphismes dans Ind(Cat\,hot). J'ai l'impression de ne pas avoir encore axiomatisé à fond la notion même de chemins, qui après tout a un sens géométrique assez évident et important par lui-même. Donc par ce qu'il faut entendre par `\textit{une} théorie des chemins dans une catégorie'. Une telle théorie $\mathcal{T}$ associerait donc à $\mathcal{M}$, $\Sigma$, $a$, $b$ un système inductif de catégories $\Ch^{\mathcal{T}}_{n,\Sigma}(\mathcal{M}; a, b)$. Et les contraintes géométriques naturelles à imposer à de telles théories, en vue aussi de leur utilisation pour expliciter
\[
\Hom_{\mathcal{M}(X)\Sigma_X^{-1}}(a_X, b_X) \xrightarrow{\ \sim\ } \pi_0\uHom(X, \Chh^{\mathcal{T}}(\mathcal{M}; a, b)) % reading resolved: two underlines, the ind-path object.
\]
donnerait un système transitif d'isomorphismes non seulement entre les ind-objets de Cat\,hot, mais peut-être aussi entre les $\Ch^{\mathcal{T}}_{\infty,\Sigma}(\mathcal{M}; a, b)$, pour des théories differentes. Si on dispose d'un tel énoncé, la catégorie $\Ch_{\infty,\Sigma}(\mathcal{M}; a, b)$, ou plus exactement ses images dans Cat\,hot et Hot, % [sic] differentes -> différentes, missing accent
\msp{page 51}
prendraient un sens convaincant.

\dgrblock{dgrb190}{What a canonical mapping space requires}
\begin{dcomment}{later}{p51:prendraient un sens convaincant}{What a canonical mapping space requires}
Grothendieck wants a theory of paths whose result is independent of its combinatorial presentation. The preceding scholium already supplies such independence at the level it actually computes: different constructions representing the same $H_{a,b}$ give uniquely isomorphic ind-objects in $\Cathot$. To obtain a canonical homotopy type of paths, one needs information beyond the connected-component calculation.

Dwyer and Kan's simplicial localization provides the standard stronger construction. For a relative category $(\mathcal M,W)$ it gives simplicial mapping objects with the localized Hom-sets as their connected components \citep[\S2 and Proposition~3.1]{DwyerKan1980Calculating}. Their hammock construction uses diagrams of zigzags and coherent higher comparisons; in a model category its mapping objects recover the usual derived function complexes \citep[Proposition~4.4 and Corollary~4.7]{DwyerKan1980Function}. Thus composition and its higher compatibility are retained while allowing different presentations of the same homotopy theory.

There is a precise caution in comparing it with these pages. Grothendieck's $K_\infty$ is the ordinary union of full categories of fixed-endpoint paths, and the morphisms of those categories are the transformations specified in sections~1--2. The theorem proved here identifies $\pi_0$ of bounded diagram categories with an ordinary localized Hom-set. That theorem alone is not a weak-equivalence comparison between $NK_\infty$ and a Dwyer--Kan function complex. Replacing vertical morphisms, changing allowed zigzags, or using a homotopy colimit may change the higher homotopy type while leaving this $\pi_0$ calculation intact.

Accordingly, the modern theory answers the demand for an invariant mapping-space object, but does not automatically verify each proposed comparison between the particular unbounded path categories in the manuscript. In the favorable model case discussed above, the desired functor $H_{a,b}$ is represented on $\Hot$ by the derived mapping space. A direct identification of a chosen $K_\infty$ with that object is a separate theorem to be proved under stated assumptions. This distinction explains both the force of the scholium and the legitimacy of the questions which follow it.
\end{dcomment}


Il y a d'autre part un autre point de vue, qui peut-être pourrait `canonifier' les catégories $\uCh_\infty$ en tant qu'objet de Cat\,hot et surtout de Hot (on passerait à Cat\,hot en appliquant le foncteur $\psi$, enveloppe de Thomason-Kan \ldots). En effet, on est intéressé à construire, en termes de $\mathcal{M}$, $\Sigma$, un `dérivateur' (ou un prédérivateur) $\mathbf{D}_{\mathcal{M},\Sigma}$ surtout sur le `domaine' des catégories finies. L'étendre à un domaine plus large est peut-être important dans un `deuxième souffle', mais peut être considéré comme du luxe dans le premier mouvement de fondation d'une théorie des dérivateurs. Ceci dit, on peut se poser des questions du type suivant.

Soit
\[
\Hot_0 \subset \Hot
\]
\printedpage{32}
% head: A. Grothendieck : Dérivateurs, Ch. V, Édition des Universités de Montpellier II et Paris VII
% foot: Version 9 septembre 2001
la sous-catégorie pleine de Hot, définie par les catégories finies (\footnote[9]{\textit{NB} Hot$_0$ est équivalente à une petite catégorie.}) - appelons les types d'homotopie `de type fini' (ou de \textit{présentation} finie?). On a un foncteur évident
\msp{page 52}
\[
\begin{aligned}
\Hot \quad &\longrightarrow \quad \Hot_0^{\wedge} = \uHom(\Hot_0^{o}, (\Ens)) \, ,\\
X \quad &\longmapsto \quad \tilde{X} = (\xi \longmapsto \Hom_{\Hot}(\xi, X)).
\end{aligned}
\]
Ce foncteur est \textit{conservatif} \texttt{[cela semble faux]}, puisque un homomorphisme $X \longrightarrow X'$ dans Hot est un isomorphisme s'il induit des isomorphismes pour les foncteurs $\tilde{X} \longrightarrow \tilde{X}'$, et même seulement pour leurs valeurs pour les sphères $S^n$ ($n \geq 0$). (Sauf erreur. Y a-t-il des canulars, parce que le critère d'équivalence faible par les $\pi_i$ ne s'applique pas directement, en l'absence de points bases. Ça ne doit pas être très sérieux \ldots\ \texttt{[cela semble en fait sérieux]}). Mais est-il peut-être même \textit{pleinement fidèle} (\footnote[10]{\texttt{[un ?\ \ dans la marge]}})? De sorte que Hot se réaliserait comme sous-catégorie pleine de $\Hot_0^{\wedge}$, contenant les foncteurs représentables?


\dgrblock{dgrb191}{Finite unbased tests miss a symmetric-group shift}
\begin{dcomment}{examples}{p52:puisque un homomorphisme}{Finite unbased tests miss a symmetric-group shift}
Grothendieck's final question concerns a specific restricted Yoneda functor, and here a concrete counterexample settles both the conservativity assertion and the stronger full-faithfulness question. The mechanism is that every finite test sees a conjugating element, but no one element works globally. The finite-support symmetric-group example is the classical case recorded in \citet[Proposition~4.1]{MatumotoMinamiSugawara1984}; the attribution and its relation to the later general result are confirmed by \citet[Remark~2.3]{ArlinChristensen2023}. We give the elementary verification in the exact form needed for these pages.

Let $G$ be the group of permutations of $\mathbf N=\{0,1,\ldots\}$ with finite support. Define a homomorphism $s:G\to G$ by
\[
 s(\sigma)(0)=0,\qquad s(\sigma)(n+1)=\sigma(n)+1.
\]
It is injective but is not surjective, since every image permutation fixes $0$, whereas the transposition $(0\ 1)$ belongs to $G$. Therefore $Bs:BG\to BG$ is not a weak homotopy equivalence: on fundamental groups it is $s$.

Now let $P$ be any finitely generated group and $\rho:P\to G$ a homomorphism. Choose a finite generating set. The union of the supports of the images of those generators is finite, so every element of $\rho(P)$ fixes all integers greater than some $m$. Let $c=(0\ 1\ \cdots\ m+1)$. Conjugation by $c$ carries the action on $\{0,\ldots,m\}$ to the shifted action on $\{1,\ldots,m+1\}$ and fixes $0$. Directly, for every $p\in P$,
\[
 c\rho(p)c^{-1}=s(\rho(p)).
\]
Hence $s\rho$ and $\rho$ are conjugate homomorphisms.

For a connected CW complex $T$, homotopy classes of maps to $BG=K(G,1)$ correspond to homomorphisms $\pi_1T\to G$ modulo conjugation. One way to verify this description is to choose a vertex, extend a chosen group homomorphism across the $2$-skeleton, and then use the vanishing higher homotopy groups of $BG$ for existence and uniqueness of further extensions. Forgetting the chosen basepoint changes the homomorphism by conjugation. If $T$ is the realization of the nerve of a finite category, each component has finitely presented fundamental group: its finite $2$-skeleton gives generators and relations. Applying the preceding conjugacy calculation on each component therefore shows
\[
 (Bs)_*:[T,BG]\longrightarrow[T,BG]
       \quad\hbox{is the identity}
\]
for every $T\in\Hot_0$. The same holds for every sphere. This proves that the restricted Yoneda functor is not conservative.

It is also not faithful. The maps $1_{BG}$ and $Bs$ are distinct in $\Hot$, since the latter is not an equivalence and the former is. Nevertheless they induce the same natural endomorphism of the presheaf $T\mapsto[T,BG]$ on $\Hot_0$: both are the identity at each test object. Thus the proposed full faithfulness fails even at injectivity on morphisms. The two editorial warnings at the end of the chapter are justified; the missing basepoint information cannot be recovered merely by adding more finite unbased tests.
\end{dcomment}


\dgrblock{dgrb192}{The modern scope of the final obstruction}
\begin{dcomment}{status}{p52}{The modern scope of the final obstruction}
The counterexample above addresses exactly the finite-category tests in the manuscript. A much stronger theorem is now available: \citet[Theorem~2.1]{ArlinChristensen2023} prove that no set of objects of the unpointed homotopy category jointly reflects isomorphisms. Their proof enlarges the symmetric-group construction to a sufficiently large regular cardinal. Consequently the obstruction is not removed by choosing a different small family of unbased tests. The same paper distinguishes this theorem from an earlier argument of Heller whose asserted boundedness property had a gap.

The loss occurs at the passage to sets of homotopy classes. Retaining more structure changes the answer: \citet[Theorem~4.3]{ArlinChristensen2023} show that spheres jointly detect equivalences in the homotopy $2$-category of spaces, whose Hom-categories are fundamental groupoids of mapping spaces. Already the automorphisms in the mapping groupoid out of a point retain the fundamental group which a set of unbased components forgets. At the infinity-categorical level, mapping spaces provide the full Yoneda data. These results give a precise modern reason for the insistence, throughout these chapters, on keeping coherent diagrams and their transformations together.
\end{dcomment}
\begin{otherlanguage}{english}
\dgroverview{III-V}{Chapter V \textperiodcentered\ \textup{III-V} \textperiodcentered\ Overview}
These chapters study how a homotopy theory can retain operations that
ordinary homotopy classes conceal. They approach the problem from three directions.
Chapter III asks for an action of homotopy types on coefficients and for an
external Hom that represents the resulting families of morphisms. Chapter IV
constructs coherent diagrams by extending coefficients along carefully chosen
inclusions. Chapter V returns to the elementary presentation of a localization
by zigzags, then asks how much of its homotopy theory those zigzags remember.
The order matters. The external Hom is a desired object; coherent diagrams
provide a way to construct operations on it; the categories of paths are a
candidate model. What is established at each stage must be distinguished from
what the next stage would require.

These chapters have three independent manuscript numberings. Chapter III
contains manuscript pages 1--12, Chapter IV contains 73--83 (the edition
explicitly says that pages 1--72 are absent), and Chapter V contains 1--52,
with page 32 marked blank. We therefore cite manuscript locations with a
chapter numeral. In the edition, the mathematical text occupies forty-seven pages, besides
titles and editorial notices.
The inserted comments preserve the edition's corrections in square brackets.
The corrections to the transcription restore what that edition actually prints;
mathematical repairs are stated in the commentary instead of being silently
substituted for the historical text.

Chapter III begins with a concrete requirement. If $p:X\to e$ and $M$ is a
coefficient, a tensor by the homotopy type of $X$ should be $p_!p^*M$, and a
cotensor should be $p_*p^*M$. Their adjunction follows directly from the two
adjunctions involving $p^*$:
\[
 \Hom(p_!p^*M,N)\cong\Hom_{\D(X)}(p^*M,p^*N)
 \cong\Hom(M,p_*p^*N).
\]
The issue is descent in $X$: a weak homotopy equivalence of indexing categories
must induce an isomorphism on constant-coefficient homology and cohomology.
There is also a size issue. A derivator defined on a restricted class of
shapes cannot be presumed to see every homotopy type. These are separate
questions, and neither is settled by the formal adjunction alone.

The modern action of spaces on derivators gives a precise setting for the
first construction, while model presentations provide derived mapping spaces.
The distinction between an action and a closed enrichment remains essential:
the former is not, by itself, a representability theorem for the external
Hom. We relate the discussion to the universal property of the derivator of
spaces and to the later account of homotopy-theoretic module structures
\citep{Cisinski2008,GrothShulman2017}. For familiar derived categories one can
see all the operations explicitly. With cohomological grading,
\[
 K\otimes M=C_*(K;\mathbf Z)\otimes^{\mathbf L}_{\mathbf Z}M,
 \qquad
 M^K=\mathbf R\Hom_{\mathbf Z}(C_*(K;\mathbf Z),M).
\]
The mapping space from $M$ to $N$ is obtained from the good truncation
$\tau_{\leq0}\mathbf R\Hom_k(M,N)$ by the Dold--Kan construction and
forgetting the abelian structure. Thus its $n$th homotopy group at zero is
$\Hom_{D(k)}(M,N[-n])$. A positive Ext group appears as a higher homotopy
group after the corresponding shift of the target; it is not a positive
homotopy group of this unshifted mapping space merely because it is an Ext
group.

The circle calculation explains why pointedness enters the manuscript.
Unreduced chains on $S^1$ give $\mathbf Z\oplus\mathbf Z[1]$, hence
$S^1\otimes M\simeq M\oplus M[1]$ in the derived category. The suspension
is the reduced summand. Grothendieck's first identification of the tensor
with suspension is therefore corrected by his own closing observation; it
should not be reported as an unresolved defect of his final construction.
His suggestion of a bounded-above enhancement, however, requires a boundedness
hypothesis on derived Hom. An elementary example over the dual numbers shows
why arbitrary coefficients do not supply one.

Chapter IV explains why these operations must be constructed on coherent
diagrams. The underlying-diagram functor is conservative, but it discards
information about coherent compositions and homotopies. The commentary on Chapter~I
already examines this failure in detail. Here the new point is constructive:
Kan extension along a suitable sequence of inclusions supplies the coherence
that a diagram drawn in the fundamental category does not specify.
The loop object is obtained by extending a single coefficient by zero to
a punctured square and then taking its right Kan extension to the whole
square. Full faithfulness of these extensions explains both the existence
and the correctly qualified uniqueness of the resulting cartesian square.
This is the pattern developed systematically for pointed and stable
derivators in \citet{Groth2013}.

The same method makes the adjunction maps visible. For $f:I\to J$,
Grothendieck forms the collage $K(f)$ with mixed morphisms
$\Hom_{K(f)}(i,j)=\Hom_J(f(i),j)$. A presheaf on $K(f)$ is exactly a
pair of presheaves together with a map $f^*G\to F$. The elementary
description identifies the required units and counits before the derived
case is attempted. We complete the construction by giving the indexing
functors and checking the pointwise Kan-extension comparison. The benefit
is a canonical coherent realization of these particular adjunction maps.
It does not assert a canonical lift of every arrow drawn in a homotopy
category.

Chapter V begins with alternating finite paths, extended by identities at
their two ends. A backward segment is required to belong to $\Sigma$;
the connected components of the resulting category with fixed endpoints
then give the morphisms of $\mathcal M[\Sigma^{-1}]$. Here ``fixed'' includes
the morphisms: transformations in the fibre have identity components at
the two endpoints. Replacing this fibre by the full subcategory on paths
with those endpoint objects would change the construction. We give the
quotient and composition argument with this convention throughout.

The manuscript's parity correction is also substantive. Concatenating two
of the initially chosen odd-indexed standard paths reverses the standard
orientation. Passing to the cofinal subsystem of even indices repairs the
problem. Grothendieck's analysis of admissible contractions establishes the
necessary compatibility with padding. His more general interval lemma is
left as an assumption, with contractibility proposed as a strengthening.
The development supplied here proves that strengthening: for a finite
oriented path as target, every nonempty category of functors with a prescribed
restriction is categorically contractible. Numerical minimum and maximum
give explicit endofunctors and a zigzag between the identity and a constant
functor. The fixed-restriction version also covers the cancellation diagrams
whose two endpoints have the same image. We give the full argument and state
the extent of the literature search; the label \emph{ours} records responsibility
for this proof, not a claim of priority.

Passing to diagrams exposes a uniformity condition that a bare union of
path categories forgets. A path in $\mathcal M^X$ has one finite length
bound valid at every object of $X$. Consequently the natural representing
object is an ind-category $(K_n)$, with
\[
 \Hom_{\mathcal M^X[\Sigma_X^{-1}]}(a_X,b_X)
   \cong\mathop{\rm colim}_n\pi_0\operatorname{Fun}(X,K_n).
\]
This is an unconditional statement about ordinary localizations. It is
not yet an identification of $NK_\infty$ with a derived mapping space.
The later simplicial-localization theory supplies canonical mapping spaces
and comparison theorems under their stated hypotheses; the commentary keeps
those theorems separate from what this particular path model proves.

The formal propositions on reflective localizations and their ind-categories
can be preceded by a simpler affirmative result. If $H_{a,b}$ inverts nerve
weak equivalences, the image of its representing ind-object already
ind-represents the descended functor on $\Hot$. Indeed, left Kan extension
along the opposite of a localization carries representable presheaves to
representable presheaves and preserves filtered colimits; restriction along
the localization is fully faithful. This conclusion needs no right adjoint
between the original categories.

The stronger formal propositions on reflective localizations and their ind-categories
are sound when the fully faithful right adjoint is supplied. Their application
to $\mathrm{Cat\,hot}\to\Hot$ needs that additional premise. Here
$\mathrm{Cat\,hot}$ means categories and functors modulo finite zigzags of
natural transformations, whereas $\Hot$ inverts weak equivalences of nerves.
The editorial request for a proof of the right adjoint is consequently
material. Thomason's model structure does not alone establish the asserted
adjunction with this particular source category. We leave that premise
unresolved in this audit and explain the valid replacement-based conclusions
without promoting an unsuccessful search to a claim that the problem is open.

The last page admits a definite answer. Testing an unpointed homotopy type
by every finite category does not give a conservative or a faithful restricted
Yoneda functor. Let $G$ be the group of permutations of $\mathbf N$ with finite
support, and let $s:G\to G$ shift a permutation one place to the right while
fixing zero. The induced $Bs$ is not an equivalence. Every map from a finite
category's nerve into $BG$, however, uses a finitely generated fundamental
group and hence only finitely many permutation supports. On those supports
the shift is conjugation by one finite permutation. Thus $Bs$ induces the
identity on all the proposed tests. The same example separates $Bs$ from
the identity despite their identical restricted Yoneda transformations.
The modern general theorem says that no set of unpointed spaces jointly
reflects isomorphisms in $\Hot$ \citep[Theorem~2.1 and Remark~2.3]{ArlinChristensen2023}.
Retaining mapping groupoids changes this conclusion: the sphere tests do
reflect equivalences in the homotopy 2-category, as their Theorem~4.3 proves.
The obstruction on the closing page is precisely information discarded by
passing too early to sets of unpointed homotopy classes.
\end{otherlanguage}
\dgrchunkend

% DER_VI-a: Grothendieck's text with the commentary set in, then the overview, inside a group that applies the chunk's own definitions
\dgrchapter{VI}{Chapter VI}{HOT}
\dgrpart{0}{VI-a}{Part 1 of 3, manuscript pages 1--70}{VI-a\enspace\textperiodcentered\enspace Chapter VI, part 1 of 3\enspace\textperiodcentered\enspace manuscript pages 1--70}{Chapter VI \textperiodcentered\ \textup{VI-a}}
\dgrchunkbegin
% DER_VI-a: the definitions of the chunk's preamble, applied inside its group
\let\printedpage\relax\newcommand{\printedpage}[1]{\marginpar{\footnotesize p.~#1}}
\let\msp\relax\newcommand{\msp}[1]{\marginpar{\footnotesize [#1]}}
\let\deq\relax\newcommand{\deq}{\overset{\text{déf}}{=}}
\let\ul\relax\newcommand{\ul}[1]{\underline{\mathrm{#1}}}
\let\xline\relax\newcommand{\xline}[1]{\mathrel{\overset{#1}{\rule[0.6ex]{2em}{0.4pt}}}}
\let\cat\relax\newcommand{\cat}[1]{\underline{\mathrm{#1}}}
\let\circled\relax\newcommand{\circled}[1]{\tikz[baseline=(c.base)]\node[draw,circle,inner sep=1pt](c){#1};}
\let\und\relax\newcommand{\und}[1]{\underline{\mathrm{#1}}}
\let\Cat\relax\newcommand{\Cat}{\mathrm{Cat}}
\let\uu\relax\newcommand{\uu}[1]{\underline{\mathrm{#1}}}
\let\lline\relax\newcommand{\lline}{\mathrel{\relbar\joinrel\relbar\joinrel\relbar}}
\let\uc\relax\newcommand{\uc}[1]{\underline{\mathrm{#1}}}
\let\HOT\relax\newcommand{\HOT}{\mathrm{HOT}}
\let\defeq\relax\newcommand{\defeq}{\overset{\text{déf}}{=}}
\let\cercle\relax\newcommand{\cercle}[1]{\tikz[baseline=(c.base)]\node[draw,circle,inner sep=1pt](c){#1};}
\let\Ob\relax\newcommand{\Ob}{\mathrm{Ob}}
\let\Fl\relax\newcommand{\Fl}{\mathrm{Fl}}
\let\id\relax\newcommand{\id}{\mathrm{id}}
\let\cA\relax\newcommand{\cA}{\mathcal{A}}

% ---- from Der-ChVI-345_p001-005.tex
% ===================== page 1 of the PDF (title page) =====================
\begin{center}
\printedpage{[sans numéro]}% no page number printed on this page
% head: (none)

\vspace*{6cm}

{\Large\textsf{LES DÉRIVATEURS}}

\bigskip

\textbf{ALEXANDRE GROTHENDIECK}

\vspace{2cm}

\textbf{Chapitre VI}

\medskip

\textbf{HOT}
\end{center}

\clearpage

% ===================== page 2 of the PDF (editorial note) =====================
\vspace*{6cm}

\noindent\printedpage{[sans numéro]}% no page number printed on this page
% head: (none)
Ce texte a été déchiffré et transcrit en \LaTeXe{} par M. Künzer. Il a été édité par M. Künzer, J. Malgoire, et G. Maltsiniotis. La transcription est aussi fidèle que possible au manuscrit. Pour les quelques corrections évidentes, ou rares commentaires des éditeurs, ainsi que pour la numérotation originale des pages du manuscrit, les caractères de machine à écrire \texttt{[typewriter]} entre crochets sont utilisés. Un point d'interrogation entre crochets signifie que l'on n'est pas sûr du mot qui précède.

\noindent Cette édition est provisoire. Les remarques, commentaires et corrections sont bienvenus. Envoyer un message à :

\noindent\textbf{maltsin@math.jussieu.fr}

\hfill G. Maltsiniotis\hspace*{2cm}

\clearpage

% ===================== printed page 1 =====================
\noindent\printedpage{1}
% head: A. Grothendieck : Dérivateurs, Ch. VI, Édition des Universités de Montpellier II et Paris VII
% foot: Version 9 novembre 2001
\msp{page 1}

\section*{\normalsize Le dérivateur \textmd{HOT}}

\begin{flushright}
14.11.90
\end{flushright}

\begin{flushright}
(Les premières feuilles, d'il y a deux semaines, ont été larguées, dépassées par VII)
\end{flushright}

\noindent Je définis Hot à partir de la catégorie de modèles
\[
\mathcal{M} = \mathrm{Cat},\ W \subseteq \mathrm{Fl}(\mathrm{Cat}),
\]
$W$ étant en principe les Hot-équivalences. Mais je vais dans la mesure du possible axiomatiser plus, en supposant donné un ensemble $W$ plus général, appelé `localisateur fondamental' (basic localiser), p.~ex. $W = W_{\mathbf{D}}$, où $\mathbf{D}$ est un dérivateur déjà donné (mais en supposant que son domaine soit (Cat) tout entier). Nous allons noter au fur et à mesure de quelles propriétés de $W$ nous aurons besoin. Aux fins de référence ultérieure, je vais rappeler ici celles que j'ai dégagées dans Pursuing Stacks.

\noindent\textbf{L1}\ \ Saturation de $W$

\begin{itemize}[leftmargin=2.5em]
\item[a)] \textit{Les identités} \texttt{[sont]} \textit{dans $W$.}
\item[b)] \textit{Si deux parmi $gf$, $f$, $g$} \texttt{[sont]} \textit{dans $W$, le troisième aussi, $X \xrightarrow{f} Y \xrightarrow{g} Z$.}
\item[c)] \textit{Si $gf$ et $fg$ sont dans $W$, alors $f$ et $g$ aussi, $X \overset{f}{\underset{g}{\rightleftarrows}} Y$ .}
\end{itemize}

\noindent\textbf{L2}\ \ $\mathbf{\Delta}^1$ \textit{est universellement $W$-asphérique, i.e. $\mathbf{\Delta}^1 \times X \longrightarrow X$ est dans $W$ pour tout $X$ dans} Cat. Cela implique

\noindent\textbf{L2 bis}\ \ $\mathbf{\Delta}^1$ \textit{est $W$-asphérique.} (C'était L3$'$ dans Pursuing Stacks.)

\noindent\textbf{L3}\ \ \textit{Si $X$ a un objet final, $X$ est $W$-asphérique.} (\footnote[1]{Il y a un L3$'$ dual de L3.})

\noindent Je note ici \texttt{[des]} variantes renforcées.

\noindent\textbf{L3 bis}\ \ \textit{Si $X$ est} Hot-\textit{asphérique, il est $W$-asphérique.}

\noindent\textbf{L3 ter}\ \ \textit{Si $X$ est} Hot-\textit{asphérique, il est universellement $W$-asphérique, i.e. $X \times Y \xrightarrow{\mathrm{pr}_2} Y$ est dans $W$ pour tout $Y$ dans} Cat. (\footnote[2]{Suffit pour entraîner L2, L3.})

\clearpage

% ===================== printed page 2 =====================
\noindent\printedpage{2}
% head: A. Grothendieck : Dérivateurs, Ch. VI, Édition des Universités de Montpellier II et Paris VII
% foot: Version 9 novembre 2001
\textbf{L3 quater}\ \ \textit{Toute} Hot-\textit{équivalence est dans $W$.}

\noindent L3 quater $\Longrightarrow$ L3 ter $\Longrightarrow$ L3 bis $\Longrightarrow$ L3.

\noindent L3 ter $\Longrightarrow$ L2 $\Longrightarrow$ L2 bis.

\noindent\msp{page 2}

\dgrblock{dgrb161}{The axioms are being found}
\begin{dcomment}{context}{p1}{The axioms are being found}
Grothendieck wants to construct the homotopy theory of diagrams from categories themselves. The expression “catégorie de modèles” here initially means a category supplied with distinguished weak equivalences; no classes of Quillen cofibrations and fibrations have yet been specified. We write W\ensuremath{\infty} for functors inducing weak homotopy equivalences of nerves, and Q : Cat \ensuremath{\to}Cat[W \ensuremath{-}1] for ordinary localization. Categories and localizations are understood in fixed, sufficiently large universes. This avoids mistaking the construction of Q for a smallness theorem about its hom sets.

The changing list of axioms is part of the argument. In particular, the L4 of manuscript page 2 is the absolute assertion that an aspherical functor is a weak equivalence. The relative assertion for a map X \ensuremath{\to}Y over an arbitrary third category S appears only as L4 bis on page 44. Later proofs require this stronger formulation. The modern definition uses the relative statement from the outset; compare Cisinski (2004, Definition 1.1.3 and Remark 1.1.4). Similarly, the L5 on page 2 concerns descent of fibrewise equivalences, whereas the L5 on page 43 concerns recovery of fibrewise equivalences from a total equivalence. They are different assertions, with different strength. The interval [1] explains why these axioms can remember homotopy. A natural transformation F \ensuremath{\Rightarrow}G is a functor A \ensuremath{\times} [1] \ensuremath{\to}B with prescribed endpoint restrictions. L2 makes both endpoints inverse to the projection after localization. Thus a natural transformation becomes an equality between the induced maps of localized categories. This observation will repeatedly turn an elementary adjunction into a weak equivalence without requiring a model structure.

The optional L3 quater also has a precise later history. Under the modern relative axioms, the inclusion W\ensuremath{\infty}\ensuremath{\subseteq}W is a theorem, established by Cisinski (2004, Theorem 2.2.11). That theorem does not license replacing the initial absolute L4 by its relative form without explanation. In these pages the additional hypothesis is explicitly adopted at manuscript page 7, and we retain the distinction when tracing the proofs.

\end{dcomment}


\begin{itemize}[leftmargin=2.5em]
\item[\textbf{L4}] \textit{Si $f : X \longrightarrow Y$ est $W$-asphérique (i.e. les $X/y \longrightarrow Y/y$ dans $W$, i.e. (si on a L3) si les $X/y$ sont $W$-asphériques), alors $f$ aussi.} (Axiome de localisation.) (\footnote[3]{il y a un L4$'$ dual de L4.}) \texttt{[`aussi'} \texttt{= dans} $W$\texttt{]}

\item[\textbf{L5}] (Axiome de fibration). \textit{Si}
\[
\begin{tikzcd}[column sep=small]
X \arrow[rr, "f"] \arrow[dr] & & Y \arrow[dl] \\
& S &
\end{tikzcd}
\]
\textit{avec $X$ et $Y$ fibrées sur $S$ et $f$ cartésien, si pour tout $s \in S$ $f_s : X_s \longrightarrow Y_s$ est dans $W$, $f$ aussi.}
\end{itemize}

\noindent Dans Pursuing Stacks j'ai supposé toujours \fbox{L1, L3, L4} (impliquant L2) - c'étaient ces conditions qui définissaient la notion de `localisateur' (fondamental) dans Cat. (\footnote[4]{Si $W$ \texttt{[est]} un localisateur, alors les catégories contractiles sont $W$-asphériques, les homotopismes sont dans $W$, les applications homotopes sont égales dans (Cat)$W^{-1} \deq \mathrm{Hot}_W$.}) D'autre part, \fbox{L1, L2 bis, L5} impliquent L1,2,3,4,5.

\noindent\textit{On va supposer ici aussi qu'on ait au moins} L1, L3, L4 (i.e. $W$ un localisateur). En appliquant la construction de I, p.~72-\textbf{75,} on voit que $f : X \longrightarrow Y$ est dans $W$ si et seulement si $f^o : X^o \longrightarrow Y^o$ l'est. Donc L1, L3, L4 impliquent L3$'$, L4$'$ (axiomes duaux de L3, L4).

\noindent Il y avait des axiomes complémentaires

\begin{itemize}[leftmargin=2.5em]
\item[\textbf{La}] $f \in W \Longrightarrow \pi_0(f)$ \textit{bijectif.}

\item[\textbf{Lb}] \textit{Le foncteur canonique de localisation}
\[
\mathrm{Cat} \longrightarrow (\mathrm{Cat})W^{-1} \deq \mathrm{Hot}_W
\]
\textit{commute aux sommes finies.}

\item[\textbf{Lc}] \textit{Itou pour produits finis.}
\end{itemize}

\clearpage

% ===================== printed page 3 =====================
\noindent\printedpage{3}
% head: A. Grothendieck : Dérivateurs, Ch. VI, Édition des Universités de Montpellier II et Paris VII
% foot: Version 9 novembre 2001
Il faudrait quand même dégager des conditions concrètes sur $W$ pour la validité de Lb, Lc. Les conditions L3 quater et La s'écrivent respectivement
\[
W_\infty \subseteq W,\ \ W \subseteq W_0.
\]

\noindent\msp{page 3}

\dgrblock{dgrb162}{What the first fibration axiom already gives}
\begin{dcomment}{exposition}{p2}{What the first fibration axiom already gives}
Grothendieck’s short implication from L1, L2 bis and the original L5 can be expanded into a useful collection of elementary arguments. It also clarifies why that first L5 becomes dispensable after relative descent has been imposed. First apply L5 to the cartesian functor [1] \ensuremath{\times} A \ensuremath{\to}A between fibrations over A. Its fibre map is [1] \ensuremath{\to}\{\ensuremath{*}\}, so L2 follows. For the projection p : A \ensuremath{\times} [1] \ensuremath{\to}A and its sections e0, e1, two-out-of-three gives e0, e1 \ensuremath{\in}W. If H : A \ensuremath{\times} [1] \ensuremath{\to}B is a natural transformation between F and G, it follows that F \ensuremath{\in}W if and only if G \ensuremath{\in}W: both are equivalent to H \ensuremath{\in}W. Moreover, Qe0 = Qp\ensuremath{-}1 = Qe1, whence QF = QG. An adjoint pair F \ensuremath{\dashv}G has GF and FG connected to the corresponding identities by its unit and counit. Both composites belong to W, and L1(c) gives F, G \ensuremath{\in}W. Applying this to the terminal-object adjunction proves L3, and the initial-object case follows as well.

There is a particularly concrete way to obtain invariance under opposites without assuming it. Let T(A) have objects the arrows a \ensuremath{\to}b of A, with a morphism from a \ensuremath{\to}b to a\ensuremath{\prime} \ensuremath{\to}b\ensuremath{\prime} given by a\ensuremath{\prime} \ensuremath{\to}a and b \ensuremath{\to}b\ensuremath{\prime} making the evident rectangle commute. The two projections

T(A)op \ensuremath{-}\ensuremath{\to}A,

T(A)op \ensuremath{-}\ensuremath{\to}Aop are fibrations. Their fibres are respectively (a\textbackslash{}A)op and A/b, each with a terminal object. L5, applied to each projection as a cartesian functor to the identity fibration, makes both projections elements of W. Their naturality in A and two-out-of-three now imply

T(f)op \ensuremath{\in}W

f op \ensuremath{\in}W. f \ensuremath{\in}W \ensuremath{\Leftarrow}\ensuremath{\Rightarrow} \ensuremath{\Leftarrow}\ensuremath{\Rightarrow} Here the fibration assertion can be checked directly: a change of the first endpoint is made by precomposition before taking opposites, and a change of the second endpoint by postcomposition. These are the specified cartesian lifts after reversal.

Now let f : X \ensuremath{\to}Y be a map over S such that every f/s : X/s \ensuremath{\to}Y/s belongs to W. The construction \ensuremath{\Psi} on manuscript pages 17--20 produces split opfibrations with these maps as their fibre maps, and \ensuremath{\Psi}(f) preserves the chosen cocartesian lifts. After taking opposites, L5 applies, so \ensuremath{\Psi}(f) \ensuremath{\in}W. The natural maps X \ensuremath{\to}\ensuremath{\Psi}(X) and Y \ensuremath{\to}\ensuremath{\Psi}(Y ) are homotopy equivalences by the adjunction exhibited below; hence two-out-of-three gives f \ensuremath{\in}W. This proves even L4 bis, and therefore L4. Each ingredient is independent of the later, much stronger L5. The argument verifies the implication recorded here while anticipating only constructions that Grothendieck gives explicitly in the next few pages.

\end{dcomment}


\noindent Leur conjonction est l'axiome des plus raisonnables
\[
W_\infty \subseteq W \subseteq W_0.
\]
Bien sûr, l'axiome L5 a maintenant l'air un peu étriqué, i.e. faiblard. On va voir très vite de quoi on aura besoin!

\noindent On veut donc définir un dérivateur
\[
\mathrm{HOT}_W
\]
à partir de la catégorie de modèles
\[
(\mathrm{Cat}, W \subseteq \mathrm{Fl}(\mathrm{Cat})).
\]
Donc pour définir
\[
\mathrm{HOT}_W(S) \qquad\qquad\qquad (S \in \mathrm{Ob}\,\mathrm{Cat}),
\]
la récette générale est de prendre la catégorie % [sic] récette -> recette, misspelling of "recette"
\[
\mathcal{M}(S) = \mathrm{Cat}(S) = \underline{\mathrm{Hom}}(S^o, (\mathrm{Cat})),
\]
d'y introduire l'ensemble de flèches
\[
\left\{
\begin{aligned}
& W_S \subseteq \mathrm{Fl}(\mathcal{M}(S)) \\
& u \in W_S \Longleftrightarrow u_s \in W \ \ \forall\, s \in \mathrm{Ob}\,S.
\end{aligned}
\right.
\]
On doit poser
\[
\mathrm{HOT}_W(S) \deq \underline{\mathrm{Hom}}(S^o, \mathrm{Cat})W_S^{-1}.
\]
Mais la catégorie $\underline{\mathrm{Hom}}(S^o, \mathrm{Cat})$ s'interprète, à isomorphisme canonique près, comme celle
\[
\underline{\mathrm{Fibsc}}\,S
\]
des catégories $X$ Cat-fibrées sur $S$, munies d'un `scindage', i.e. d'un choix de foncteurs changement de base $u^* : X_t \longrightarrow X_s$ (pour $t \xrightarrow{u} s$ dans $S$) \texttt{[plutôt} $u^* : X_s \longrightarrow X_t$\texttt{]}, avec de plus \textit{transitivité stricte}
\[
(vu)^* = u^* v^* \qquad\qquad\qquad \text{si}\ \ t \xrightarrow{u} s \xrightarrow{v} r\ \ \text{dans}\ S
\]

\noindent\msp{page 4}

\dgrblock{dgrb163}{Products and sums survive the localization}
\begin{dcomment}{exposition}{p3}{Products and sums survive the localization}
Grothendieck asks for concrete conditions ensuring Lb and Lc because a homotopy theory should still allow objects to be assembled and compared. A useful answer is available as soon as relative L4 bis is admitted. We give the argument, including the passage from closure of weak equivalences to the universal property of products and coproducts.

Suppose W satisfies L1, L2 and L4 bis. If f : A \ensuremath{\to}B lies in W, then f \ensuremath{\times} idC lies in W for every C. View this map over C. Its comma map at c is f \ensuremath{\times} idC/c. Since C/c has a terminal object, the two projections A \ensuremath{\times} (C/c) \ensuremath{\to}A and B \ensuremath{\times} (C/c) \ensuremath{\to}B are homotopy equivalences: the terminal object defines sections and a contraction natural transformation. They are in W by the preceding block. Two-out-of-three makes the comma map a weak equivalence, and L4 bis gives the assertion. The same axiom, applied over a discrete indexing category I, shows that ` i fi \ensuremath{\in}W whenever every fi \ensuremath{\in}W. The latter closure is also Cisinski (2004, Proposition 1.1.5).

For completeness, localization of a product of relative categories is the product of their localizations. One can see this directly from the generators: (w, id) and (id, w\ensuremath{\prime}) become invertible, their commutation squares still commute after inversion, and a pair of finite zigzags can be interleaved to give a zigzag in the product. Consequently both the diagonal and the binary product functor descend. The adjunction \ensuremath{\Delta}\ensuremath{\dashv}\ensuremath{\times} descends with its original unit and counit, since those are natural transformations between relative functors and the triangle identities remain identities. Its descended right adjoint is therefore a binary product in Cat[W \ensuremath{-}1]. The same argument with ` \ensuremath{\dashv}\ensuremath{\Delta}gives binary coproducts. The adjunctions defining the initial and terminal objects descend without any closure hypothesis beyond identities belonging to W.

Thus finite products and coproducts are preserved by Q under these concrete assumptions. This is an ordinary localization argument, not a claim that arbitrary homotopy limits or colimits agree with ordinary ones. In particular, it does not supply the derived Kan extensions needed for the derivator. Those require the diagram categories whose construction begins here.

\end{dcomment}


% ---- from Der-ChVI-345_p006-010.tex
\printedpage{4}
% head: A. Grothendieck : Dérivateurs, Ch. VI, Édition des Universités de Montpellier II et Paris VII
et non seulement transitivité à isomorphisme canonique près. Ainsi on peut écrire aussi
\[
\mathrm{HOT}_W(S) \xrightarrow{\ \sim\ } (\ul{Fibsc}\,S)W_S^{-1}, \tag*{$(*)$}
\]
où cette fois $W_S$ est interprété comme
\[
W_S \subseteq \mathrm{Fl}(\ul{Fibsc}\,S),
\]
savoir
\[
W_S = \{u \in \mathrm{Fl}(\ul{Fibsc}\,S) \mid u_s \in W \text{ pour tout } s \in \mathrm{Ob}\,S\}
\]
(foncteurs strictement cartésiens, i.e. commutant aux scindages, qui induisent des $W$-équivalences entre les catégories fibres).

La flèche dans $(*)$ est bel et bien un isomorphisme de catégories, pas seulement une équivalence. Mais on va définir d'autres catégories, plus commodes à certains égards, qui seront \textit{équivalentes} (on l'espère), mais sûrement pas isomorphes, à $\mathrm{HOT}_W(S)$. (On s'occupera plus tard de la variance de $\mathrm{HOT}_W(S)$ en $S$.)

Les catégories $\ul{Hom}(S^o,\mathrm{Cat})$, $\ul{Fibsc}\,S$, ainsi que les variantes que nous allons introduire maintenant, sont en fait des 2-catégories. Mais pour décrire $\mathrm{HOT}_W(S)$, seule les structures de 1-catégories sont utilisées. On a l'impression que c'est bien % [sic] seule les structures -> seules les structures, accord du pluriel
\msp{page 5}
grossier, qu'on perd quelque chose ce faisant. Mais ce qui est perdu va se retrouver, et surabondamment, dans la structure des \underline{Hom} internes dans les $\mathrm{HOT}(S)$, si on se `rappelle' que ces \underline{Hom} internes, donnant lieu à des `types d'homotopie relatifs' $\ul{Hom}_{\mathrm{HOT}_W}(S;\xi,\eta)$ (si $\xi$, $\eta$ en sont), peuvent s'interpréter en termes de $\infty$-champs en groupoïdes sur $S$, avec toute la richesse de structure qui leur appartient. Mais ceci est une autre histoire\dots

\dgrblock{dgrb164}{A diagram written as its category of elements}
\begin{dcomment}{exposition}{p4}{A diagram written as its category of elements}
Grothendieck passes to fibrations in order to replace an externally indexed diagram by one category carrying a projection to its index. This is the same organizing construction used in descent; the relevant background is Grothendieck (1971, Exposé VI, \S{}\S{}8--9). It is useful to write out the variance because both directions of the construction will soon be needed.

For F : Sop \ensuremath{\to}Cat, its Grothendieck construction has objects (s, x) with x \ensuremath{\in}F(s). An arrow (s, x) \ensuremath{\to}(t, y) consists of u : s \ensuremath{\to}t and a : x \ensuremath{\to}F(u)y in F(s). Composition is

(v, b) \ensuremath{\circ}(u, a) = (vu, F(u)(b) \ensuremath{\circ}a).

The chosen cartesian lift of u at (t, y) is (u, id) : (s, F(u)y) \ensuremath{\to}(t, y). The strict functoriality of F makes these choices compose strictly. Conversely, a normalized splitting of a fibration gives strict pullback functors u\ensuremath{*}, with (vu)\ensuremath{*}= u\ensuremath{*}v\ensuremath{*}and identity pullbacks equal to identities. These equations, rather than the existence of unspecified isomorphisms, are what is meant by a splitting.

A natural transformation F \ensuremath{\to}G induces a functor of the total categories that preserves the specified lifts. Conversely, such a functor restricts to the components of a strict natural transformation. Therefore the weak equivalences in the diagram category correspond exactly to fibrewise weak equivalences among splitting-preserving functors. With fibrations encoded by their fibre families, this is the identification Grothendieck calls an isomorphism. If arbitrary underlying object sets for total categories are allowed, the usual construction and reconstruction are canonically isomorphic over S, and the comparison is expressed as an equivalence. No homotopical argument depends on this bookkeeping convention.

The subsequent forgetful functor to Fib S changes the morphisms: there one allows all functors over S. Forgetting a splitting is therefore not merely forgetting redundant labels. The natural comparison maps used later must be checked in the category whose morphisms they actually preserve. This is the source of the obstruction on manuscript pages 32--39, and of the strictification work beginning on page 52.

\end{dcomment}


En fait, je vois un diagramme de 13 catégories de `categories sur $S$\,', à examiner à présent % [sic] categories -> catégories, accent manquant
\[
\begin{tikzcd}[column sep=small]
 & \ul{Fibsc}_0S \arrow[dl, hook] \arrow[d, "\mathrm{fid}"] & & \ul{Cofibsc}_0S \arrow[d] \arrow[dr, hook] & \\
\ul{Fibsc}\,S \arrow[d, "\mathrm{fid}"] & \ul{Fib}_0S \arrow[dl, hook] & & \ul{Cofib}_0S \arrow[dr, hook] & \ul{Cofibsc}\,S \arrow[d, "\mathrm{fid}"] \\
\ul{Fib}\,S \arrow[dr, hook] & & \ul{Parf}\,S \arrow[dl, hook] \arrow[d, hook] \arrow[dr, hook] & & \ul{Cofib}\,S \arrow[dl, hook] \\
 & \ul{Liss}\,S \arrow[dr, hook] & \ul{Fibhot}\,S \arrow[d, hook] & \ul{Prop}\,S \arrow[dl, hook] & \\
 & & \mathrm{Cat}/S & &
\end{tikzcd}
\]
% foot: Version 9 novembre 2001

\printedpage{5}
% head: A. Grothendieck : Dérivateurs, Ch. VI, Édition des Universités de Montpellier II et Paris VII
où \underline{Fib}, \underline{Cofib} désignent les catégories Cat-fibrées, Cat-cofibrées sur $S$, \underline{Liss} et \underline{Prop} les catégories
\msp{page 6}
lisses resp. propres sur $S$, \underline{Parf} celles qui sont à la fois lisses et propres (fibrations parfaites), enfin \underline{Fibhot} les catégories Hot-fibrées sur $S$. Toutes les flèches sont des foncteurs strictement pleinement fidèles, à l'exception des deux (plutôt quatre) flèches verticales supérieures, qui sont fidèles sans plus. Je viens d'en rajouter encore 4, de catégories : $\ul{Fibsc}_0$, $\ul{Fib}_0$, $\ul{Cofibsc}_0$, $\ul{Cofib}_0$ sont les sous-catégories pleines des catégories correspondantes sans indice, formées des catégories fibrées resp. cofibrées $X$ sur $S$, pour lesquelles les foncteurs changement de base resp. cochangement de base
\[
u^* : X_t \longrightarrow X_s, \qquad\quad \text{resp.} \qquad\quad u_* : X_s \longrightarrow X_t,
\]
pour $s \xrightarrow{u} t$ dans $S$, sont des Hot-équivalences. En termes de foncteurs
\[
S^o \xrightarrow{F} \mathrm{Cat} \qquad\quad \text{resp.} \qquad\quad S \xrightarrow{F'} \mathrm{Cat},
\]
cela signifie que ce foncteur transforme flèches quelconques dans $S$ en éléments de $W$ dans Cat, i.e. qu'on a un diagramme commutatif
\[
\begin{tikzcd}[row sep=large]
S^o \arrow[r, "F"] \arrow[d, "\mathrm{can}"'] & \mathrm{Cat} \arrow[d, "\mathrm{can}"] \\
\Pi(S^o) \simeq (\Pi S)^o = S^o(\mathrm{Fl}\,S^o)^{-1} \arrow[r, "\bar F"'] & \mathrm{Hot}_W = (\mathrm{Cat})W^{-1}
\end{tikzcd}
\]
(et de même pour $F'$), i.e. qu'on a un `système
\msp{page 7}
local' $\bar F$ sur $S$, à valeurs dans $\mathrm{Hot}_W$, défini à l'aide de $F$. Je rappelle qu'on trouve également un tel système local dans le cas d'une catégorie Hot-fibrée sur $S$ (et même un système local à valeurs dans Hot lui-même, au lieu de $\mathrm{Hot}_W$ - on utilise ici tacitement que $W_\infty \subseteq W$, i.e. \fbox{L3 quater} (\footnote[5]{On va le supposer vrai dans toute la suite.})). Et il en est de même a fortiori pour les objets de $\ul{Parf}\,S$. Pour $\ul{Liss}\,S$ resp. $\ul{Prop}\,S$, on trouve simplement un foncteur
\[
S^o \xrightarrow{\bar F} \mathrm{Hot} \qquad\quad \text{resp.} \qquad\quad S \xrightarrow{\bar F'} \mathrm{Hot},
\]
d'où, en composant avec $\mathrm{Hot} \longrightarrow \mathrm{Hot}_W$, des foncteurs à valeurs dans $\mathrm{Hot}_W$.

Dans chacune de ces 14 \texttt{[13?]} catégories, on introduit un ensemble de flèches, qui serviront à localiser, et on le désignera toujours par la même notation $W_S$. Dans toutes ces catégories sauf seulement $\mathrm{Cat}/S$, $W_S$ désigne la $W$-équivalence `fibre par fibre'. Dans $\mathrm{Cat}/S$ d'autre part, ce sera la $W$-équivalence sans plus. Toutes les flèches du diagramme sont compatibles avec les ensembles localisants de flèches.
% foot: Version 9 novembre 2001

\printedpage{6}
% head: A. Grothendieck : Dérivateurs, Ch. VI, Édition des Universités de Montpellier II et Paris VII
\msp{page 8}
\[
\begin{tikzcd}[column sep=small, row sep=large]
 & \ul{Fibsc}_0S \arrow[dl, hook, "\alpha'"'] \arrow[d] & & \ul{Cofibsc}_0S \arrow[d] \arrow[dr, hook, "\beta'"] & \\
\ul{Fibsc}S \arrow[d] & \ul{Fib}_0S \arrow[dl, hook, "\alpha"'] \arrow[ddr, dotted, bend right=20] & & \ul{Cofib}_0S \arrow[dr, hook, "\beta"] \arrow[ddl, dotted, bend left=20] & \ul{Cofibsc}\,S \arrow[d] \\
\ul{Fib}S \arrow[dr, hook] & & \ul{Parf}\,S \arrow[dl, hook] \arrow[d, hook] \arrow[dr, hook] & & \ul{Cofib}\,S \arrow[dl, hook] \\
 & \ul{Liss}\,S \arrow[dr, hook] & \ul{Fibhot}\,S \arrow[d, hook] \arrow[uuul, dotted, bend left=10, "\Phi_0"] \arrow[uuur, dotted, bend right=10, "\Psi_0"'] & \ul{Prop}\,S \arrow[dl, hook] & \\
 & & \mathrm{Cat}/S \arrow[uuull, dotted, bend left=60, "\Phi"] \arrow[uuurr, dotted, bend right=60, "\Psi"'] & &
\end{tikzcd}
\]
% ?? uncertain: starting point of the dotted arrows labelled Phi_0 and Psi_0 (the dotted lines run alongside Parf S down to Fibhot S) -- Fibhot S or Parf S

\dgrblock{dgrb165}{A local system in the homotopy category misses twisting}
\begin{dcomment}{examples}{p7}{A local system in the homotopy category misses twisting}
Grothendieck’s passage from F to F\textasciimacron{} : \ensuremath{\Pi}(Sop) \ensuremath{\to}HotW records monodromy, but he also signals on page 5 that richer coherence remains. A familiar pair of bundles makes the difference visible. We use ordinary spaces to display the phenomenon; replacing spaces by their simplex categories gives categorical models with the same homotopy information.

Consider the trivial circle bundle S2 \ensuremath{\times} S1 \ensuremath{\to}S2 and the Hopf bundle (Hatcher, 2002, Examples 4.45 and 4.51)

S3 = \{(z0, z1) \ensuremath{\in}C2 : |z0|2 + |z1|2 = 1\} \ensuremath{-}\ensuremath{\to}CP1 \ensuremath{\sim}= S2, (z0, z1) 7\ensuremath{-}\ensuremath{\to}[z0 : z1]. On the chart z0 = 0, with coordinate w = z1/z0, a trivialization is

\ensuremath{\lambda}

\ensuremath{\lambda} \ensuremath{\in}S1. (w, \ensuremath{\lambda}) 7\ensuremath{-}\ensuremath{\to} 1 + |w|2 (1, w), p The other chart is analogous. Thus both maps are locally trivial bundles with fibre S1, and hence fibrations. Since the base is simply connected, their parallel transport determines the same functor from the fundamental groupoid to the homotopy category: a constant local system with value S1. Nevertheless they are not fibre homotopy equivalent. Such an equivalence would induce a weak equivalence of total spaces by the homotopy exact sequences, whereas

H1(S2 \ensuremath{\times} S1; Z) = Z.

H1(S3; Z) = 0, The missing information is the coherence of transport around two-dimensional homotopies, not an additional choice of monodromy around loops.

In contemporary language the appropriate datum is a functor from the entire homotopy type of the base, considered as an \ensuremath{\infty}-groupoid, to spaces. Passing to the fundamental groupoid of the base and the ordinary homotopy category of spaces discards part of that datum. Straightening and unstraightening provide the relevant modern comparison; Lurie (2009, Theorems 2.2.1.2 and 3.2.0.1) treat respectively space-valued and category-valued versions. Here this explains the direction of Grothendieck’s remark about \ensuremath{\infty}-groupoid stacks, without asserting that a single ordinary category HOTW (S) already determines every higher structure he envisages. The commentary on the earlier chapters, on what a diagram theory remembers, remains applicable.

There is also a small logical distinction in the definition on page 6. If the subscript 0 means that all transition functors lie in W\ensuremath{\infty}, their images lie in W once W\ensuremath{\infty}\ensuremath{\subseteq}W. The converse need not hold for a larger W. A W-locally constant diagram and an ordinarily homotopically locally constant diagram must therefore be distinguished whenever W is still variable.

\end{dcomment}


\texttt{[Dans ce diagramme, il y a aussi} $\ul{Liss}_0$ \texttt{et} $\ul{Prop}_0$ \texttt{mais sans connexion aux autres objets.]}

Je suis quasiment sûr du

\textbf{Theorème.} (\footnote[6]{C'est totalement canulé, cf. commentaires page 32 ff., solution des perplexités pages 41 ff. Voir surtout le Théorème-Scholie récapitulatif p. 70-81.}) \textit{Tous les foncteurs précédents induisent des équivalences entre les catégories localisées.} % [sic] Theorème -> Théorème, accent manquant

Comme il y a 16 foncteurs, ça a l'air d'un exercise de patience. Le principe général, c'est qu'on arrive à définir des foncteurs en sens inverse, compatibles avec les $W_S$, de telle façon que les composés soient reliées aux foncteurs identiques par des flèches qui sont dans $W_S$, de sorte que par passage aux localisés ces foncteurs induisent des (équivalences) quasi-inverses des foncteurs déduits de ceux du diagramme. Tous ces foncteurs en sens opposé résultent essentiellement par composition, de trois foncteurs bien familiers. Deux de ceux-ci sont indiqués en pointillés sur le diagramme, comme $\Phi$, $\Psi$. % [sic] exercise -> exercice, orthographe française ; % [sic] reliées -> reliés, accord avec « les composés »

\msp{page 9}
Le troisième est un foncteur canonique
\[
C : \mathrm{Cat}/S \longrightarrow \ul{Parf}\,S,
\]
qu'on va expliciter également.
% foot: Version 9 novembre 2001

\dgrblock{dgrb166}{The walking arrow detects the failed equivalence}
\begin{dcomment}{examples}{p8}{The walking arrow detects the failed equivalence}
Grothendieck’s first comparison theorem would identify all the categories in the large diagram after the specified localizations. The later marginal note withdraws it. The obstruction already occurs over the smallest base with a noninvertible arrow and explains why the distinction between fine and coarse equivalence occupies the rest of the chapter.

Let S = [1] = (0 \ensuremath{\to}1) and take W = W\ensuremath{\infty}. Consider the contravariant diagrams

F(1) = \ensuremath{\varnothing},

F(0) = \{\ensuremath{*}\}, G(0) = G(1) = \{\ensuremath{*}\}. There is a unique natural transformation F \ensuremath{\to}G. Its Grothendieck construction is the functor \{0\} ,\ensuremath{\to}[1] over S. Both source and target are split fibrations, and both total categories have contractible nerve, so this map is a coarse weak equivalence. Its fibre at 1 is \ensuremath{\varnothing}\ensuremath{\to}\{\ensuremath{*}\}, which is not a weak equivalence.

This difference persists after localization; it is not merely a difference between two presenting classes of arrows. Evaluation at 1 descends from the category of strict diagrams localized at objectwise weak equivalences to Hot. The displayed map has image \ensuremath{\varnothing}\ensuremath{\to}\{\ensuremath{*}\} there and therefore cannot become an isomorphism. Its coarse image is invertible by definition. Thus the forgetful comparison from fibrewise localization to coarse localization is not an equivalence. The same argument applies whenever W \ensuremath{\subseteq}W0, since then empty and nonempty categories cannot become isomorphic after applying \ensuremath{\pi}0.

The locally constant subcategory is also genuinely smaller. If a diagram H(1) \ensuremath{\to}H(0) has a weak equivalence as its transition and were isomorphic to F after objectwise localization, evaluation would make H(1) empty and H(0) contractible in Hot, contradicting the invertibility of its transition. Reversing all arrows gives the corresponding example for split opfibrations. This does not invalidate the formal localization lemma on the next pages: the missing hypothesis is that its comparison transformations lie in the particular weak-equivalence classes being used.

\end{dcomment}


\printedpage{7}
% head: A. Grothendieck : Dérivateurs, Ch. VI, Édition des Universités de Montpellier II et Paris VII
Le foncteur $\Phi$ (resp. $\Psi$) appliqué à $\ul{Fib}_0S$ (resp. $\ul{Cofib}_0S$), induit un foncteur
\[
\ul{Fib}_0S \longrightarrow \ul{Fibsc}_0S \qquad\quad \text{resp.} \qquad\quad \ul{Cofib}_0S \longrightarrow \ul{Cofibsc}_0S \tag*{$(*)$}
\]
(ce qu'il faudra vérifier). Cela fournit toutes les flèches `inverses' dont on a besoin, à l'exception seulement de $\alpha$, $\beta$, $\alpha'$, $\beta'$. Mais si on a prouvé, par la méthode indiquée, que toutes les flèches du diagramme induisent des équivalences sur les catégories localisées, pour prouver qu'il en est encore de même pour $\alpha$, $\beta$, $\alpha'$, $\beta'$, il suffit de le prouver pour les deux foncteurs
\[
\ul{Fibsc}_0S \longrightarrow \ul{Fib}\,S, \qquad\qquad \ul{Cofibsc}_0S \longrightarrow \ul{Cofib}\,S,
\]
donc il reste à définir les foncteurs en sens inverse
\[
\ul{Fib}\,S \longrightarrow \ul{Fibsc}_0S, \qquad\qquad \ul{Cofib}\,S \longrightarrow \ul{Cofibsc}_0S,
\]
et comme les sources sont des sous-catégories pleines de $\mathrm{Cat}/S$, il suffira de définir
\[
\Phi_0 : \mathrm{Cat}/S \longrightarrow \ul{Fibsc}_0S, \qquad\qquad \Psi_0 : \mathrm{Cat}/S \longrightarrow \ul{Cofibsc}_0S.
\]
Pour ceci, on posera
\[
\Phi_0 = \Phi r C, \qquad\qquad \Psi_0 = \Psi r C, \tag*{$(**)$}
\]
\msp{page 10}
où $C : \mathrm{Cat}/S \longrightarrow \ul{Parf}\,S$, $\Phi$ et $\Psi$ sont les 3 foncteurs `rétrogrades' inversés plus haut, et $r : \ul{Parf}(S) \hookrightarrow \mathrm{Cat}/S$ est l'inclusion. Il faut donc vérifier que $\Phi_0$, $\Psi_0$ définis par $(**)$, qui a priori ont comme catégories but $\ul{Fibsc}\,S$ et $\ul{Cofibsc}\,S$, prennent leur valeur dans les sous-catégories avec indice 0.

Je vais donner un énoncé d'équivalence unique, qui couvrira (presque) tous les 16 cas à la fois :

\textbf{Corollaire.} \textit{Considérons le diagramme}
\[
\begin{tikzcd}[column sep=small, row sep=small]
 & \ul{Parf}\,S \arrow[dd, "r"'] & \\
\ul{Fibsc}_0S \arrow[dr, "p"'] & & \ul{Cofibsc}_0S \arrow[dl, "q"] \\
 & \mathrm{Cat}/S &
\end{tikzcd}
\]
\textit{et soit $\mathcal{F}$ une sous-catégorie pleine de $\mathrm{Cat}/S$, telle que l'un des foncteurs $p$, $q$, $r$ se factorise par $\mathcal{F}$. Soit $W_{\mathcal{F}}$ induit sur $\mathcal{F}$ par $W_S$ dans $\mathrm{Cat}/S$. Alors $\mathcal{F} \hookrightarrow \mathrm{Cat}/S$ induit une équivalence sur les catégories localisées.}

\msp{page 11}
\textsc{Démonstration}. Je dis qu'il suffit de prouver, pour les trois foncteurs $p$, $q$, $r$, et pour les foncteurs en sens inverse
% foot: Version 9 novembre 2001
\printedpage{8}
% head: A. Grothendieck : Dérivateurs, Ch. VI, Édition des Universités de Montpellier II et Paris VII
\[
\begin{tikzcd}[column sep=small, row sep=small]
 & \ul{Parf}\,S & \\
\ul{Fibsc}_0S & & \ul{Cofibsc}_0S \\
 & \mathrm{Cat}/S \arrow[uu, "C"] \arrow[ul, "\Phi_0"] \arrow[ur, "\Psi_0"'] &
\end{tikzcd}
\]
que les six composés
\[
p\Phi_0, \Phi_0p, q\Psi_0, \Psi_0q, rC, Cr
\]
sont reliés aux foncteurs identiques (dans la catégorie pertinente) par une flèche qui est dans le localisateur correspondant. Pour chacun de ces trois couples de foncteurs $(p,\Phi_0)$, $(q,\Psi_0)$, $(r,C)$, cela prouvera déjà qu'en passant aux localisés, ils induisent des équivalences quasi-inverses l'une de l'autre. Mais on a mieux à l'aide du lemme formel suivant.

\textbf{Lemme.} \textit{Soient $(C,W)$, $(C',W')$ deux catégories avec localisateurs $W \subseteq \mathrm{Fl}\,C$, $W' \subseteq \mathrm{Fl}\,C'$. Soit}
\[
p : C \longrightarrow C'
\]
\textit{un foncteur compatible avec les localisateurs, et supposons qu'il existe un foncteur $\varphi : C' \longrightarrow C$ également compatible, et tel que l'on ait des flèches (à direction non précisée)}
\[
\varphi p \xline{u} id_C, \qquad\qquad p\varphi \xline{v} id_{C'},
\]
\textit{qui soient dans $W$ resp. $W'$ `argument par argument'.}

\msp{page 12}
\textit{Alors $p$ et $\varphi$ induisent des équivalences quasi-inverses l'une de l'autre en passant aux catégories localisées (elles n'ont pas de merite). De plus, pour toute sous-catégorie pleine $C'_0$ de $C'$, telle que $p$ se factorise (en $p_0 : C \longrightarrow C'_0$) par $C'_0$, si on désigne par $W'_0$ le localisateur induit dans $C'_0$ par $C'$, alors $p_0 : C \longrightarrow C'_0$ et $\varphi_0 = \varphi|C'_0 : C' \longrightarrow C$} \texttt{[plutôt} $C'_0 \longrightarrow C$\texttt{]} \textit{induisent également des équivalences quasi-inverses l'une de l'autre sur les localisés.} % [sic] merite -> mérite, accent manquant

En effet, $\varphi_0p_0 = (\varphi i)p_0 = \varphi(ip_0) = \varphi p \xline{u} \mathrm{id}_C$ (ici $i : C'_0 \hookrightarrow C'$ désigne l'inclusion), et $F \overset{\text{déf}}{=} p_0\varphi_0 = p_0\varphi i$ satisfait à
\[
iF = \underbrace{ip_0}_{p}\varphi i = (p\varphi)i.
\]
Pour définir $F \xline{v_0} \mathrm{id}_{C'_0}$, qui soit dans $W'_0$ argument par argument,
\msp{page 13}
il suffit de définir
\[
\underbrace{iF}_{=\,(p\varphi)i} \xline{\bar v_0} i,
\]
qui soit dans $W'$ argument par argument, puisque $W'_0$ induit par $W'$ et que $F \longmapsto iF$
\[
\ul{Hom}(C'_0,C'_0) \longrightarrow \ul{Hom}(C'_0,C')
\]
% foot: Version 9 novembre 2001

\dgrblock{dgrb167}{The localization lemma keeps its full force}
\begin{dcomment}{exposition}{p12}{The localization lemma keeps its full force}
Grothendieck separates a sound formal mechanism from the geometric assertions to which he hopes to apply it. The lemma is useful precisely because its hypotheses record both kinds of information: preservation of weak equivalences by the functors, and membership of the comparison maps in the specified weak-equivalence classes.

Let P : (C, W) \ensuremath{\to}(D, V ) and R : (D, V ) \ensuremath{\to}(C, W) be relative functors. Suppose natural transformations, in either direction, connect RP with idC and PR with idD, with all components respectively in W and V . Their images under localization are natural isomorphisms. The induced functors P\textasciimacron{} and R\textasciimacron{} are therefore quasi-inverse equivalences. One needs neither a calculus of fractions nor triangle identities for the original transformations. Quasi-inverse equivalences, unlike a prescribed adjunction, require only the two natural isomorphisms.

Now suppose P factors through a full subcategory i : D0 ,\ensuremath{\to}D, written P = iP0, and put V0 = V \ensuremath{\cap}Mor(D0) and R0 = Ri. The first comparison restricts to R0P0 = RP. The second restricts to a transformation between iP0R0 and i. Because both of its endpoints lie in D0, fullness gives a unique transformation between P0R0 and idD0; faithfulness verifies its naturality. Every component belongs to V0. Consequently P0 and R0 also give inverse equivalences after localization. Applying the same argument to i shows that D0[V \ensuremath{-}1 ] \ensuremath{\to}D[V \ensuremath{-}1] is an equivalence.

This explains the apparently strong corollary about an arbitrary full subcategory containing the image of a replacement functor. The weak-equivalence class on that full subcategory must be the induced one. For example, a full subcategory of Cat/S containing all perfect replacements may calculate the coarse localization. Nothing in the lemma permits changing its induced coarse class to the fibrewise class without proving their agreement. Nor does fullness apply to the passage from splitting-preserving maps to arbitrary maps over S: that functor is generally faithful but not full. These are the two points at which the first comparison diagram asks for additional mathematics.

\end{dcomment}


% ---- from Der-ChVI-345_p011-015.tex
\printedpage{9}
% head: A. Grothendieck : Dérivateurs, Ch. VI, Édition des Universités de Montpellier II et Paris VII
est pleinement fidèle ($i$ l'étant). Or il suffit de prendre
\[
\bar{v}_0 = v * i : (p\varphi)\circ i \mathrel{\relbar\joinrel\relbar} \mathrm{id}_{C'}\circ i.
\] % ?? uncertain: the arrow between (p\varphi)\circ i and id_{C'}\circ i is printed as a plain line with no visible head -- plain line / arrow

Pour résumer, il faut montrer que le composé des couples de foncteurs quasi-inverses marche dans les cinq cas suivants
\[
\begin{tikzcd}[column sep=large, row sep=huge]
\cat{Fibsc} \arrow[rrd, shift left] & \cat{Fibsc}_0 \arrow[l, hook'] \arrow[rd, shift left] & \cat{Parf} \arrow[d, shift left] & \cat{Cofibsc}_0 \arrow[r, hook] \arrow[ld, shift left] & \cat{Cofibsc} \arrow[lld, shift left] \\
& & \mathrm{Cat}/S \arrow[llu, shift left, "\Phi"] \arrow[lu, shift left, "\Phi_0"] \arrow[u, shift left, "C"] \arrow[ru, shift left, "\Psi_0"] \arrow[rru, shift left, "\Psi"] & &
\end{tikzcd}
\]

D'ailleurs par dualité, on est ramené aux trois cas de $\cat{Fibsc}$, $\cat{Fibsc}_0$, $\cat{Parf}$. (Peut-être le cas $\cat{Fibsc}_0$, avec $\Phi_0$ défini en termes de $\Phi$ et de $C$, se ramènera-t-il aux deux cas qui l'entourent, de $\cat{Fibsc}$ et $\cat{Parf}$.)

Je vois que je dois rectifier le tir, tout d'abord parce que j'ai oublié que les foncteurs
\msp{page 14}
$\Phi$, $\Psi$ que j'avais en vue ne sont \textit{pas} compatibles avec les localiseurs \texttt{[= localisateurs]} - une flèche $f : X \longrightarrow Y$ de $\mathrm{Cat}/S$ telle que la flèche correspondante dans Cat soit dans $W$, ne va pas induire une flèche dans $\cat{Fibsc}\,S$ ou $\cat{Cofibsc}\,S$ qui soit dans $W$ sur chaque fibre. Ce n'est pas grave, car on peut remplacer les foncteurs par $\Phi' = \alpha'\Phi_0$, $\Psi' = \beta'\Psi_0$, qui sauf erreur doivent marcher. Mais en vue du Lemme précédent, il y a une perplexité, car j'ai l'air d'avoir prouvé que si $\mathcal{F}$ est une sous-catégorie pleine de $\mathrm{Cat}/S$ dont je suppose seulement qu'elle contient l'image essentielle d'une des trois catégories $\cat{Fib}_0S$, $\cat{Cofib}_0S$, ou $\cat{Parf}\,S$, alors l'inclusion dans $\mathrm{Cat}/S$ induit une équivalence des catégories localisées, en prenant dans $\mathcal{F}$ comme localisateur celui induit par $\mathrm{Cat}/S$, i.e.
\begin{tikzcd}[column sep=small, row sep=small, baseline=(\tikzcdmatrixname-2-2.base)]
f : X \arrow[rr] \arrow[dr] & & Y \arrow[dl] \\
& S &
\end{tikzcd}
est dans $W_{\mathcal{F}}$ si et seulement si en tant que flèche de Cat, $f$ est dans $W$. Or ceci coïncide bien (via un théorème non trivial sur lequel je devrai revenir dans le cadre d'un localisateur fondamental général $W$) avec l'ensemble de flèches de localisation `fibre par fibre', dans le cas de $\mathcal{F} = \cat{Parf}\,S$, $\cat{Fibhot}\,S$, $\cat{Fib}_0S$, $\cat{Cofib}_0S$.

\msp{page 15}
Par contre, ce n'est nullement le cas en prenant $\mathcal{F} = \cat{Fib}\,S, \cat{Liss}\,S$ ou $\cat{Cofib}\,S$. Ayant choisi dans ces catégories comme localisateurs la notion de $W$-équivalences `fibre par fibre', j'ai donc l'air de dire que je trouve pour ces quatre catégories les \textit{mêmes} localisées (à équivalence près) que si je prenais le localisateur beaucoup plus grossier hérité par $\mathrm{Cat}/S$. Cela semble difficile à croire, et il faudra que je fasse très attention en démontrant les énoncés-clef.

Pour voir que les flèches en traits pleins du grand diagramme p.~6 sont toutes compatibles avec les localis(at)eurs \texttt{[sic]}, donc aussi les flèches composées au moyen de celles-ci, j'ai à utiliser déjà ceci (pour le cas des flèches $\cat{Liss} \longrightarrow \mathrm{Cat}/S$, $\cat{Prop} \longrightarrow \mathrm{Cat}/S$):
% foot: Version 9 novembre 2001

\printedpage{10}
% head: A. Grothendieck : Dérivateurs, Ch. VI, Édition des Universités de Montpellier II et Paris VII
\circled{1}\ \ Si $X, Y$ lisses ou propres sur $S$, et $f : X \longrightarrow Y$ un $S$-foncteur, si $f_s : X_s \longrightarrow Y_s$ est dans $W$ pour tout $s \in S$, alors $f \in W$. (\footnote[7]{OK si on suppose L3 quater et L4.}).

Cela renforce considérablement la condition L5 sur $W$, où on supposait que $f$ est un foncteur \textit{cartésien} entre catégories \textit{fibrées} sur $S$. (À dégager, pour les différentes propriétés dont on a besoin pour $W$, à quelles propriétés standard on peut les ramener. Tout sera OK si $W = W_\infty$.)

\msp{page 16}
I) \textbf{Définition de} $\Phi$, $\Psi$.

\dgrblock{dgrb168}{The precise local-to-global argument}
\begin{dcomment}{exposition}{p15}{The precise local-to-global argument}
Grothendieck needs fibrewise equivalence to imply total equivalence without requiring a map of fibrations to preserve cartesian arrows. Smoothness or properness provides the comparison between strict fibres and comma categories that makes this possible. The treatment of these notions in the commentary on Chapter~I explains the lifting categories behind those comparisons; here their role is descent for a map of total categories.

For a proper functor p : X \ensuremath{\to}S, the canonical functor Xs \ensuremath{\to}X/s is a W\ensuremath{\infty}-equivalence; for a smooth functor, the corresponding functor is Xs \ensuremath{\to}s\textbackslash{}X. These assertions, in the appropriate aspherical form, are the criterion developed in Maltsiniotis (2005, Theorem 1.2 and its dual). Suppose X, Y are proper over S and f : X \ensuremath{\to}Y is any functor over S. Naturality gives

Xs

Ys

X/s

Y/s. Under W\ensuremath{\infty}\ensuremath{\subseteq}W, the vertical maps belong to W. The two-out-of-three property implies that fs \ensuremath{\in}W if and only if f/s \ensuremath{\in}W. If this holds for every s, relative L4 bis gives f \ensuremath{\in}W. For smooth functors the proof uses the undercategories and the dual relative axiom. No cartesian or cocartesian condition on f occurs.

The argument also explains the eventual universal statement on pages 41--43. Proper and smooth functors are stable under change of base. For S\ensuremath{\prime} \ensuremath{\to}S, the strict fibre of the pulled-back map at s\ensuremath{\prime} is fs\ensuremath{\prime} = fp(s\ensuremath{\prime}). Applying the same proof over S\ensuremath{\prime} shows that a fibrewise equivalence between two proper objects, or between two smooth objects, remains a total equivalence after every base change.

The parenthetical assertion here that L3 quater and L4 suffice suppresses one hypothesis: the version of Theorem A actually used is relative L4 bis. Absolute L4 directly covers the case Y = S, since S/s has a terminal object. For general Y it is the stronger formulation that is needed in the displayed proof. Grothendieck explicitly identifies this distinction on manuscript page 44. We follow that correction rather than assigning an unproved consequence to the initial axiom.

\end{dcomment}


On va calquer la construction canonique de $\cat{Ch}(X) \longrightarrow X \times X$ (ou $\cat{Ch}_\infty(X) \longrightarrow X \times X$) factorisant $X \xrightarrow{\mathrm{diag}} X \times X$, et la factorisation canonique correspondante d'un $f : X' \longrightarrow X$ en $X' \xrightarrow{i} \bar{X}' \xrightarrow{\bar{f}} Y$ \texttt{[plutôt} $\bar{X}' \xrightarrow{\bar{f}} X$\texttt{]}. On remplacera simplement $\cat{Ch}(X)$ par la sous-catégorie $\cat{Ch}_+(X)$, laquelle n'est autre que la catégorie des flèches de $X$
\[
\cat{Fl}(X) = \cat{Hom}(\boldsymbol{\Delta}_1, X),
\]
muni des foncteurs
\[
X \xrightarrow{i_X} \cat{Fl}(X) \xrightarrow{p_X=(\sigma_X,\beta_X)} X \times X,
\]
où $i_X$ est la flèche $x \longmapsto \mathrm{id}_x$, et $\sigma_X$, $\beta_X$ sont les foncteurs source et but. La fibre de $p_X$ en $(a,b) \in \mathrm{Ob}(X \times X)$ est la catégorie discrète définie par l'ensemble $\mathrm{Hom}(a,b)$, ça dépend grandement de $(a,b)$ et pas seulement de sa composante connexe dans $X \times X$ - donc $p_X$ n'a rien d'une Hot-fibration. Il n'est en géneral ni propre ni lisse ni rien, par contre les deux flèches composées avec $\mathrm{pr}_1$ et $\mathrm{pr}_2$ % [sic] géneral -> général, missing accent
\[
\begin{tikzcd}[column sep=small]
& \cat{Fl}(X) \arrow[ld, "\sigma_X"'] \arrow[rd, "\beta_X"] & \\
X & & X
\end{tikzcd}
\]
sont sympathiques, on a en effet

\msp{page 17}
\textbf{Lemme.} \textit{On a}
\[
\sigma_X^{-1}(\{x\}) \simeq x\backslash X, \qquad\qquad \beta_X^{-1}(\{x\}) \simeq X/x.
\]
\textit{Le foncteur} $\sigma_X$ \textit{est} Cat-\textit{fibrant à fibres contractiles (donc asphériques, donc} $W$\textit{-asphériques), le foncteur} $\beta_X$ \textit{cofibrant à fibres contractiles (donc asphériques, donc} $W$\textit{-asphériques). Ce sont donc aussi des} Hot-\textit{équivalences, a fortiori des} $W$\textit{-équivalences} (\footnote[8]{En fait, ce sont même des homotopismes, cf. remarque marginale plus bas (plus générale - on retrouve le cas de $\sigma_X$ en faisant $X' = X$.)}).

\textbf{Corollaire 1.} \textit{Le foncteur} $i_X : X \longrightarrow \cat{Fl}(X)$ \textit{est une} Hot-\textit{équivalence} (puisque son composé avec la Hot-équivalence $\sigma_X$ est l'identité).
% foot: Version 9 novembre 2001

\printedpage{11}
% head: A. Grothendieck : Dérivateurs, Ch. VI, Édition des Universités de Montpellier II et Paris VII
Voici la construction générale de la factorisation de $f : X' \longrightarrow X$, correspondant au choix précédent d'une `catégorie de chemins' $\mathcal{C}(X) = \cat{Fl}(X)$
\[
\begin{tikzcd}
\bar{X}' \arrow[d, "\pi_f"] \arrow[rd] \arrow[rrd, "\bar{f}"] & & \\
X' \arrow[u, dotted, bend left=60, "i_f"] \arrow[rd, "f"'] \arrow[r, phantom, "\text{\footnotesize Cart.}"] & \cat{Fl}(X) \arrow[r, "\beta_X"'] \arrow[d, "\sigma_X"'] & X \\
& X &
\end{tikzcd}
\]
où $i_f$ est déduit de la section $i$ \texttt{[plutôt} $i_X$\texttt{]} de $\cat{Fl}(X)$ sur $X$ (\footnote[9]{\textbf{NB} $\pi_f i_f = \mathrm{id}$, et $i_f\pi_f \longrightarrow \mathrm{id}$, donc $i_f$, $\pi_f$ sont des homotopismes inverses l'un de l'autre dans Cat\,hot. Cela précise \texttt{[le]} cor.~3, page 18.}). Ainsi, $\bar{X}'$ est la catégorie des couples
\[
\{(x',u) \mid x' \in \mathrm{Ob}\,X',\ \underbrace{u}_{\in\mathrm{Fl}\,(X)} : f(x') \longrightarrow y\},
\]
le foncteur $\bar{f}$ est défini par
\[
\bar{f}(x',u) = \mathrm{but}(u) = y
\]
et $i_f$ par
\[
i_f(x') = (x', \mathrm{id}_{f(x')} : f(x') \xrightarrow{\mathrm{id}} f(x')),
\]
donc on a bien
\[
f(x') = \bar{f} i_f(x').
\]

\msp{page 18}
\textbf{Corollaire 2.} \textit{Le foncteur}
\[
\bar{f} : \bar{X}' \longrightarrow X
\]
\textit{a comme fibres les catégories $X'/x$, il est cofibrant, les foncteurs de changement de base sont les foncteurs $X'/x \longrightarrow X'/y$ associés aux flèches $x \longrightarrow y$ dans $X$.}

\dgrblock{dgrb169}{The arrow category has two different good endpoints}
\begin{dcomment}{exposition}{p17}{The arrow category has two different good endpoints}
Grothendieck introduces a path construction which remembers the direction of an arrow. Its endpoint maps are excellent replacement tools separately, although the pair of endpoints is not a homotopy-fibration model of the diagonal. The elementary adjunctions make all of these assertions transparent. Write Ar(A) = Fun([1], A), and let i(a) = ida, \ensuremath{\sigma}(a \ensuremath{\to}b) = a, \ensuremath{\beta}(a \ensuremath{\to}b) = b. A commutative square from i(x) to a \ensuremath{\to}b is determined by its arrow x \ensuremath{\to}a. A square from a \ensuremath{\to}b to i(y) is determined by b \ensuremath{\to}y. Therefore

\ensuremath{\beta} \ensuremath{\dashv}i \ensuremath{\dashv}\ensuremath{\sigma}.

In particular \ensuremath{\sigma}i = \ensuremath{\beta}i = id, and the squares associated with an arrow a \ensuremath{\to}b give natural transformations i\ensuremath{\sigma} \ensuremath{\Rightarrow}id and id \ensuremath{\Rightarrow}i\ensuremath{\beta}. Both endpoint projections and i are categorical homotopy equivalences, hence belong to every W satisfying L1 and L2.

The fibration properties can be checked without choosing anything. Pulling a \ensuremath{\to}b back along a\ensuremath{\prime} \ensuremath{\to}a means precomposing to obtain a\ensuremath{\prime} \ensuremath{\to}b; the resulting square is cartesian for \ensuremath{\sigma}. Pushing it forward along b \ensuremath{\to}b\ensuremath{\prime} means postcomposing to obtain a \ensuremath{\to}b\ensuremath{\prime}; that square is cocartesian for \ensuremath{\beta}. The fibres are a\textbackslash{}A and A/b, with initial and terminal identities respectively. This example already belongs to the original descent formalism; see Grothendieck (1971, Exposé VI, \S{}11(a)).

For the combined map (\ensuremath{\sigma}, \ensuremath{\beta}) : Ar(A) \ensuremath{\to}A \ensuremath{\times} A, however, the fibre at (a, b) is the discrete set HomA(a, b). Take A = [1]. Its fibre at (1, 0) is empty, while that at (0, 0) is a point, although the base is connected. Thus it is not a Hot-fibration. It is neither smooth nor proper: at (1, 0) its undercategory has the object arising from (1, 0) \ensuremath{\to}(1, 1), and its overcategory has the object arising from (0, 0) \ensuremath{\to}(1, 0), whereas its strict fibre is empty. In fact each of these comma categories is a point. The necessary fibre-to-comma weak-equivalence condition fails in both directions. This distinction is why the directed arrow construction gives \ensuremath{\Phi} and \ensuremath{\Psi}, while the two-sided replacement C requires the longer path category Ch\ensuremath{\infty}(S).

\end{dcomment}


En fait, on voit que cette catégorie est non seulement cofibrée, mais canoniquement \textit{scindée.}

(\footnote[10]{\textbf{Corollaire 3.} \textit{La flèche $i_f : X' \longrightarrow \bar{X}'$ est un $X$-morphisme (si $X'$ sur $X$ par $f$, $\bar{X}'$ par $\bar{f}$), i.e. $f = \bar{f}\circ i_f$, et c'est une} Hot-\textit{équivalence} (car $\pi_f : \bar{X}' \longrightarrow X'$ est Cat-fibrant, à fibres contractiles, donc une Hot-équivalence).}).

Appliquant ceci à des $f : X \longrightarrow S$, on trouve une construction
\[
\mathrm{Cat}/S \xrightarrow{\Psi} \cat{Cofibsc}\,S,
\]
c'est bel et bien fonctoriel en $X$ variant dans $\mathrm{Cat}/S$. Si
\[
\psi : X \longrightarrow Y
\]
% foot: Version 9 novembre 2001
\printedpage{12}
% head: A. Grothendieck : Dérivateurs, Ch. VI, Édition des Universités de Montpellier II et Paris VII
dans $\mathrm{Cat}/S$, alors $\psi$ induit des
\[
\psi/s : X/s \longrightarrow Y/s\ ,
\]
pour $s \in S$, et ceci de façon compatible avec les
\msp{page 19}
foncteurs `cochangement de base' associés aux $u : s \longrightarrow t$ dans $S$
\[
\begin{tikzcd}[column sep=large, row sep=large]
X/s \arrow[r, "u_*^X"] \arrow[d, "\psi/s"'] & X/t \arrow[d, "\psi/t"] \\
Y/s \arrow[r, "u_*^Y"] & Y/t \quad .
\end{tikzcd}
\]
Cela définit donc un morphisme de catégories scindées sur $S$
\[
\Psi(X) \longrightarrow \Psi(Y),
\]
c'est lui!

On a un morphisme fonctoriel en $X$
\[
X \xrightarrow{i_X} \Psi(X)
\]
dans $\mathrm{Cat}/X$, et c'est une Hot-équivalence, donc $\in W$. Sur les fibres, le morphisme $i_X$ induit
\[
X_s \xrightarrow{i_{X,s}} X/s,
\]
l'inclusion canonique.

\msp{page 20}
\textbf{Corollaire 1.} \textit{Si $X \xrightarrow{f} Y$ dans} $\mathrm{Cat}/S$ \textit{est dans $W$, alors $\Psi(X) \xrightarrow{\Psi(f)} \Psi(Y)$ aussi.} Mais on se gardera de croire que c'est dans $W_S$ au sens de $\cat{Cofib}\,S$ ou $\cat{Cofibsc}\,S$.

\dgrblock{dgrb170}{A canonical opfibration and its two projections}
\begin{dcomment}{exposition}{p19}{A canonical opfibration and its two projections}
Grothendieck’s replacement \ensuremath{\Psi} packages all the overcategories X/s into one category. Writing down its arrows shows both why it is strictly functorial and why a projection used to prove homotopy invariance need not be a morphism over the original base.

For p : X \ensuremath{\to}S, an object of \ensuremath{\Psi}(X) is (x, u : p(x) \ensuremath{\to}s). A morphism to (y, u\ensuremath{\prime} : p(y) \ensuremath{\to}t) is a pair

vu = u\ensuremath{\prime}p(a).

(a : x \ensuremath{\to}y, v : s \ensuremath{\to}t), The structural projection q sends this object to s and this morphism to v. For v : s \ensuremath{\to}t, the morphism (idx, v) : (x, u) \ensuremath{\to}(x, vu) is the specified cocartesian lift. Composition of these lifts is literal composition, so the opfibration is split and its fibre at s is X/s.

There is also the forgetful projection \ensuremath{\pi} : \ensuremath{\Psi}(X) \ensuremath{\to}X, which sends (x, u) to x, and the functor i : X \ensuremath{\to}\ensuremath{\Psi}(X) sending x to (x, idp(x)). A map from i(x) to (y, u\ensuremath{\prime}) is uniquely determined by a : x \ensuremath{\to}y, because its base arrow must be u\ensuremath{\prime}p(a). Hence i \ensuremath{\dashv}\ensuremath{\pi}, with \ensuremath{\pi}i = idX and counit

\ensuremath{-}\ensuremath{-}\ensuremath{-}\ensuremath{-}\ensuremath{\to}(x, u).

i\ensuremath{\pi}(x, u) = (x, idp(x)) This proves that i and \ensuremath{\pi} are inverse up to categorical homotopy. The identity qi = p makes i a functor over S, but p\ensuremath{\pi}(x, u) = p(x) generally differs from q(x, u) = s. Thus \ensuremath{\pi} generally is not a functor over S. The printed statement that i is a morphism in Cat/X is nonetheless true when \ensuremath{\Psi}(X) carries \ensuremath{\pi}: i is a section. Both descriptions of i are legitimate, and neither licenses using \ensuremath{\pi} as an inverse in Cat/S.

For a functor f : X \ensuremath{\to}Y over S, define \ensuremath{\Psi}(f)(x, u) = (f(x), u) and \ensuremath{\Psi}(f)(a, v) = (f(a), v). This preserves the specified cocartesian lifts and is strictly functorial in f. The maps i form a natural transformation in Cat/S. Their being in W shows, by the resulting commutative square, that \ensuremath{\Psi} preserves and reflects coarse weak equivalences. On fibres, however, its effect is f/s. This last observation specifies the stronger class for which \ensuremath{\Psi} becomes a replacement preserving fibrewise weak equivalences.

\end{dcomment}


Car sur les fibres de $\Psi(X)$, $\Psi(Y)$ on trouve les flèches
\[
f/s : X/s \longrightarrow Y/s
\]
induits par $f$, qui n'ont aucune raison d'être dans $W$. Qu'elles le soient signifie que $f : X \longrightarrow Y$ est dans $W$ `localement au dessus de S' (\footnote[11]{cf. \circled{1} page 15.}), ce qui (pour les $W$ qui nous intéressent) est beaucoup plus fort que d'être dans $W$. (Si $Y = S$, cela signifie en fait que $p : X \longrightarrow S$ est $W$-asphérique, beaucoup plus fort que $p \in W$.)

Donc pour prouver que
\[
\cat{Cofib}\,S \longrightarrow \mathrm{Cat}/S \qquad\quad \text{ou} \qquad\quad \cat{Cofibsc}\,S \longrightarrow \mathrm{Cat}/S
\]
% foot: Version 9 novembre 2001
\printedpage{13}
% head: A. Grothendieck : Dérivateurs, Ch. VI, Édition des Universités de Montpellier II et Paris VII
induisent des équivalences pour les catégories localisées, le foncteur $\Psi$ directement ne sera pas utilisable. Cependant :

\textbf{Corollaire 2.} \textit{Supposons $X$, $Y$ dans} $\cat{Prop}\,S$ \textit{ou dans} $\cat{Fibhot}\,S$\textit{, et considérons}
\[
\Psi(f) : \Psi(X) \longrightarrow \Psi(Y),
\]
\textit{pour que ce soit une $W$-équivalence dans} $\cat{Cofibsc}\,S$\textit{,}
\msp{page 21}
\textit{i.e. une $W$-équivalence fibre par fibre, il faut et il suffit que $f$ soit une $W$-équivalence fibre par fibre} (et cela doit impliquer (\footnote[12]{cf. \circled{1} page 15.}) que $f$ est une $W$-équivalence).

\dgrblock{dgrb171}{Which comma construction repairs which variance}
\begin{dcomment}{exposition}{p20}{Which comma construction repairs which variance}
Grothendieck now isolates the extra condition that makes \ensuremath{\Psi} a relative replacement for fibrewise homotopy theory. Its fibre map is exactly f/s, so the condition is that every overcategory map lie in W. When the target is S itself this says that f is W-aspherical, since S/s has a terminal object. For a general target it is a relative condition, later called fine equivalence on the left.

The interval calculation after manuscript page 8 already shows why total equivalence cannot suffice. The complete comparison of both endpoint inclusions, given alongside manuscript page 32, also separates the two comma tests. On a smooth object, the strict fibre compares to the undercategory s\textbackslash{}X; on a proper object, to the overcategory X/s. Consequently the relevant fibrewise replacement for a fibration is \ensuremath{\Phi}, and for an opfibration it is \ensuremath{\Psi}. This is the geometric reason the later fine-right equivalence belongs on the fibration side of the large diagram, while fine-left equivalence belongs on the opfibration side. The direction names refer to the comma constructions, not to the position of the objects in that diagram.

\end{dcomment}


En effet, dans le diagramme commutatif
\[
\begin{tikzcd}[column sep=large, row sep=huge]
X_s \arrow[r, "f_s"] \arrow[d] & Y_s \arrow[d] \\
X/s \arrow[r, "f/s"] & Y/s
\end{tikzcd}
\]
les flèches verticales sont dans $W$, car $X$, $Y$ \texttt{[sont]} propres sur $S$ (ou $X$, $Y$ dans $\cat{Fibhot}\,S$), donc $f_s$ est dans $W$ si et seulement si $f/s$ l'est.

\texttt{[Le foncteur]} $\Phi$ se définit de façon duale à $\Psi$, ce qui revient à échanger les rôles de $\sigma_X$, $\beta_X$ dans le diagramme page 17.

(\footnote[13]{\textbf{Corollaire 3.} \textit{Si $X \in \mathit{Ob}\,$}$\cat{Fibhot}\,S$\textit{, alors $\Psi(X)$ est dans} $\cat{Fibsc}_0S$.\par Cela signifie en effet que les $X/s \longrightarrow X/t$ sont dans $W$, or l'hypothèse sur $X$ implique que c'est dans $W_\infty$. OK.})

\textbf{II) Définition de} $C : \mathrm{Cat}/S \longrightarrow \cat{Parf}\,S$

Cette fois, c'est la factorisation des flèches $p : X \longrightarrow S$ dans Cat, en
\[
X \xrightarrow{j_{X/S}} \bar{X} \xrightarrow{\bar{p}} S,
\]
où $i_X$ \texttt{[plutôt} $j_{X/S}$\texttt{]} est une Hot-équivalence (donc dans $W$), et $\bar{p}$ parfait. Cela se construit formellement comme tantôt, en utilisant $\cat{Ch}_\infty(S)$ au lieu de $\cat{Fl}(S)$, et la factorisation
\[
S \xrightarrow{j_S} \cat{Ch}_\infty(S) \xrightarrow{p_S=(\sigma_S,\beta_S)} S \times S
\]
de $\mathrm{diag}_S : S \longrightarrow S \times S$. La situation est plus jolie, car $p_S$ est déjà `parfait', ce qui implique formellement que le $\bar{p}$ qui s'en déduit, pour factoriser $p : X \longrightarrow S$, est lui-même parfait.
% foot: Version 9 novembre 2001

% ---- from Der-ChVI-345_p016-020.tex
\printedpage{14}
% head: A. Grothendieck : Dérivateurs, Ch. VI, Édition des Universités de Montpellier II et Paris VII
% foot: Version 9 novembre 2001

\msp{page 22}
Ici encore, la construction est fonctorielle, on trouve le foncteur
\[ C : \Cat/S \longrightarrow \und{Parf}\,S \]
qu'il faut composer au foncteur en sens inverse
\[ r : \Cat/S \longleftarrow \und{Parf}\,S. \]
Tout d'abord, vérifions que $C$, tout comme $r$, est compatible avec les localiseurs. Soit $f : X \longrightarrow Y$ dans $\Cat/S$, supposons $f$ est dans $W$, prouvons $C(f)$ est dans $W_S$. On a commutativité dans
\[
\begin{tikzcd}[row sep=large, column sep=large]
X \arrow[r, "f"] \arrow[d, "j_X"'] & Y \arrow[d, "j_Y"] \\
CX \arrow[r, "Cf"] & CY
\end{tikzcd}
\]
et $j_X$, $j_Y$ sont dans $W$, donc $f$ est dans $W$ si et seulement si $Cf$ l'est. Mais dans le cas $W = W_\infty$ on sait que cela équivaut à ce que $Cf$ soit dans $W_\infty$ fibre par fibre. Admettant \circled{1} p. 15 pour $W$, il reste à s'assurer d'une réciproque dans \underline{Parf}, i.e.

\dgrblock{dgrb172}{From comma fibres to homotopy fibres}
\begin{dcomment}{later}{p21}{From comma fibres to homotopy fibres}
Grothendieck introduces the longer path category because he now wants a replacement whose transport is homotopically invertible in both directions. The guiding theorem is Quillen’s Theorem B: if every transition X/s \ensuremath{\to}X/t induced by an arrow s \ensuremath{\to}t is a weak homotopy equivalence, then N(X/s) computes the homotopy fibre of NX \ensuremath{\to}NS over s. A precise restatement appears in Cisinski (2004, Theorem 2.3.2). This is the point at which locally constant comma data become actual fibre data for the homotopy type of a map.

For an object already Hot-fibred over S, the comparison Xs \ensuremath{\to}X/s is a weak equivalence, and the transitions between these comma categories are weak equivalences. Consequently \ensuremath{\Psi}(X) is a split opfibration with homotopically invertible transition functors. The footnote’s printed target Fibsc0S is therefore a slip: the target is Cofibsc0S. Its dual statement concerns \ensuremath{\Phi} and split fibrations. The slip does not affect the argument in the adjacent square.

The functor C uses results about Ch\ensuremath{\infty}(S) established outside these pages. The input here is a natural factorization X \ensuremath{-}\ensuremath{-}\ensuremath{\to}C(X) \ensuremath{\to}S with jX \ensuremath{\in}W\ensuremath{\infty}and C(X) \ensuremath{\to}S perfect. Given this input, all subsequent uses of its functoriality are formal. For instance, the naturality square of j implies that f \ensuremath{\in}W if and only if C(f) \ensuremath{\in}W, provided W\ensuremath{\infty}\ensuremath{\subseteq}W. What is not formal is replacing the last total condition by a condition on fibres for an arbitrary larger W. That is the question on pages 22--24, and the stronger L5 introduced later addresses exactly it. The correct comparison therefore depends on which localizer is being studied, not only on the quality of the replacement object.

\end{dcomment}


\circled{2} Soit
\begin{tikzcd}[baseline=(\tikzcdmatrixname-2-2.base), row sep=small, column sep=small]
X \arrow[rr, "f"] \arrow[dr] & & Y \arrow[dl] \\
& S &
\end{tikzcd}
dans $\und{Parf}\,S$ (ou seulement dans $\und{Fibhot}\,S$). Si $f$ est dans $W$, alors les $f_s : X_s \longrightarrow Y_s$ aussi.

\msp{page 23}
Essayons de le prouver, en calquant la démonstration dans le cas $C = C_\infty$ (XI, th.~2). On peut supposer $S$ connexe (\footnote[14]{ne sert à rien}), en utilisant La (i.e. $W \subseteq W_0$) et en admettant le

\textbf{Lemme.} \textit{Soit $f : X \longrightarrow Y$ dans} Cat. \textit{Alors $f \in W$ si et seulement si $\pi_0(f)$ est bijectif, et pour toute composante connexe $X_i$ de $X$, si $Y_i$ est la composante connexe correspondante de $Y$, $f_i : X_i \longrightarrow Y_i$ induit par $f$ est dans $W$.}

Utilisons alors l'existence de
\[ S' \longrightarrow S \]
Hot-fibration avec $S'$ Hot-asphérique (\footnote[15]{\textbf{NB} Il suffit d'avoir une Hot-fibration $S' \longrightarrow S_0 \hookrightarrow S$, où $S_0$ est la composante connexe de $S$ contenant l'élément $s \in \mathrm{Ob}\,S$ envisagé.}). On peut supposer même que $S' \longrightarrow S$ est parfait. Supposons qu'on sache qu'un changement de base par de telles fibrations parfaites transforme $W$-équivalences en itou. (Ça doit être vrai pour les $W_i$, mais on a vu que c'était faux pour l'équivalence $k$-homologique, cf. I p.~87 ff.) Alors $X' \xrightarrow{f'} Y'$ sur $S'$ est
\printedpage{15}
% head: A. Grothendieck : Dérivateurs, Ch. VI, Édition des Universités de Montpellier II et Paris VII
% foot: Version 9 novembre 2001
une $W$-équivalence. Or pour $s' \in S'$, comme $\{s'\} \longrightarrow S'$ est une Hot-équivalence ($S'$ Hot-asphérique), et $X'$, $Y'$ étant Hot-fibrés sur $S$, il s'ensuit que $X'_{s'} \longrightarrow X'$, $Y'_{s'} \longrightarrow Y'$ sont des Hot-équivalences, donc aussi $X'_{s'} \xrightarrow{f'_{s'}} Y'_{s'}$, comme on voit sur le carré commutatif
\[
\begin{tikzcd}[row sep=large, column sep=large]
X'_{s'} \arrow[r, "f'_{s'}"] \arrow[d, "\wr"'] & Y'_{s'} \arrow[d, "\wr"'] \\
X' \arrow[r, "{f'}"] & Y'.
\end{tikzcd}
\]

\msp{page 24}
Comme $S' \longrightarrow S$ est surjectif, cela montre que les $f_s : X_s \longrightarrow Y_s$ sont dans $W$, qed.

\dgrblock{dgrb173}{The contractible-base argument without a component reduction}
\begin{dcomment}{exposition}{p23}{The contractible-base argument without a component reduction}
Grothendieck’s strategy is to inspect a chosen fibre after making the base contractible. Its force is easiest to see when the required base-change assumption is stated before the reduction. It then becomes clear that the componentwise lemma is unnecessary for this particular argument, as his marginal annotation already observes.

Assume W\ensuremath{\infty}\ensuremath{\subseteq}W, two-out-of-three, and the following additional property: pullback along a perfect functor preserves W-equivalences. Let f : X \ensuremath{\to}Y be a functor over S, with X, Y Hot-fibred over S, and suppose its total functor lies in W. Fix s \ensuremath{\in}S. Use the perfect replacement of the point \{s\} \ensuremath{\to}S to obtain

\ensuremath{-}\ensuremath{\to}jS\ensuremath{\prime} \ensuremath{-}\ensuremath{\to}rS, \{s\} j \ensuremath{\in}W\ensuremath{\infty}, where r is perfect and S\ensuremath{\prime} is weakly contractible. The distinguished object s\ensuremath{\prime} = j(s) lies over s. This construction is allowed even when S is disconnected; its image lies in the relevant component.

Put X\ensuremath{\prime} = X \ensuremath{\times}S S\ensuremath{\prime} and Y \ensuremath{\prime} = Y \ensuremath{\times}S S\ensuremath{\prime}. The extra pullback hypothesis gives f \ensuremath{\prime} : X\ensuremath{\prime} \ensuremath{\to}Y \ensuremath{\prime} in W. Base change preserves the Hot-fibration property. Since S\ensuremath{\prime} is weakly contractible, the inclusions of the fibres at s\ensuremath{\prime} into X\ensuremath{\prime} and Y \ensuremath{\prime} are weak homotopy equivalences. Equivalently, their homotopy-fibre squares are pulled back along the weak equivalence \{s\ensuremath{\prime}\} \ensuremath{\to}S\ensuremath{\prime}. In the commutative square

Xs

Ys

X\ensuremath{\prime}

Y \ensuremath{\prime}, three sides lie in W, so fs \ensuremath{\in}W. This proves the fibre-detection assertion under precisely the stated additional hypothesis. It does not prove fs \ensuremath{\in}W\ensuremath{\infty}from f \ensuremath{\prime} \ensuremath{\in}W.

The proof uses neither La nor a decomposition of f into connected components. Its genuinely substantial input is stability of W under these perfect base changes, later called L6. The speculation here that truncated homotopy equivalences should satisfy that stability is withdrawn in the later annotations; the obstruction is discussed at manuscript pages 43--45. The proof itself is valid as a conditional argument, and its dependence on L6 is the mathematical information to retain.

\end{dcomment}


Ainsi, (2) sera conséquence de \circled{2.} Soit $p : Y' \longrightarrow Y$ une Hot-fibration (ou même un foncteur parfait, si ça peut aider \dots). Soit $f : X \longrightarrow Y$ une $S$-flèche (dans $\Cat/S$), alors $f' : X' \longrightarrow Y'$ par changement de base. Si $f \in W$, $f' \in W'$
\[
\begin{tikzcd}[row sep=large]
X' \arrow[r, "q"] \arrow[d, "{f'}"'] \arrow[dr, phantom, "\scriptstyle\text{cart.}"] & X \arrow[d, "f"] \\
Y' \arrow[r, "p"] & Y
\end{tikzcd}
\qquad\qquad
f \in W \text{ et } p \text{ parfait} \Longrightarrow f' \in W
\]

En tous cas, il nous faut absolument \circled{2} pour pouvoir conclure que $C$ est compatible avec les localiseurs.

Considérons $rC : \Cat/S \longrightarrow \Cat/S$, on a une flèche
\[
\left\{
\begin{array}{l}
\mathrm{id}_{\Cat/S} \xrightarrow{j} rC \\[1.5ex]
j_X : X \longrightarrow CX \quad \text{la flèche canonique } i_{X/S}\ \text{\texttt{[plutôt} } j_{X/S}\text{\texttt{]}} \text{ (p.~21),}
\end{array}
\right.
\]
et $j_X$ est toujours dans $W$.

Considérons $Cr : \und{Parf}\,S \longrightarrow \und{Parf}\,S$, i.e. appliquons la construction $C$ à un $X$ parfait sur $S$. Alors on a encore
\[
\left\{
\begin{array}{l}
\mathrm{id}_{\und{Parf}\,S} \xrightarrow{j'} Cr \\[1ex]
j'_X = j_{r(X)} : X = r(X) \xrightarrow{j_X} CX.
\end{array}
\right.
\]
Ce morphisme est dans $W$, donc aussi une $W$-équivalence fibre par fibre, puisque source et but sont parfaits, donc des Hot-fibrés sur $S$.

\msp{page 25}
Cela prouve déjà que les foncteurs $r$ et $C$ induisent des équivalences quasi-inverses l'une de l'autre sur les catégories localisées. Par le lemme p.~11, cela montre que pour toute sous-catégorie pleine $\mathcal{F}$ de $\Cat/S$ contenant $\und{Parf}\,S$, l'inclusion de $\mathcal{F}$ dans $\Cat/S$ induit une
\printedpage{16}
% head: A. Grothendieck : Dérivateurs, Ch. VI, Édition des Universités de Montpellier II et Paris VII
% foot: Version 9 novembre 2001
équivalence sur les localisés, en prenant comme localiseur dans $\mathcal{F}$ le localiseur `grossier' induit par $\Cat/S$. Cela s'applique notamment aux sous-catégories pleines
\[ \und{Liss}\,S, \quad \und{Fibhot}\,S, \quad \und{Prop}\,S, \]
mais en prenant dans $\und{Liss}\,S$ et $\und{Prop}\,S$ non pas le localiseur fin `fibre par fibre', comme je l'avais annoncé à tort, mais le localiseur grossier induit par $\Cat/S$. (Pour $\und{Fibhot}\,S$, il n'y a pas de risque d'ambiguïté). Il faudrait que je me convainque que si je prends la $W$-équivalence fine (dans le cas $W = W_\infty$, disons), je ne trouve \textit{pas} les `bons' localisés, i.e. le foncteur des localisés, vers le localisé de $\Cat/S$, n'est pas une équivalence. Ceci en fait suppose que dans $\und{Fib}\,S$ et $\und{Cofib}\,S$ également, c'est l'équivalence \textit{grossière}
\msp{page 26}
qu'il convient de prendre. (Dans $\und{Fib}_0S$ et $\und{Cofib}_0S$, ou $\und{Fibsc}_0S$ et $\und{Cofibsc}_0S$, les deux équivalences sont identiques, puisque ces quatre catégories sont contenues dans $\und{Fibhot}\,S$.) La suite nous montrera s'il en est bien ainsi. Par contre, dans $\und{Fibsc}\,S$ et $\und{Cofibsc}\,S$, il n'y a pas de doute, suivant les principes généraux de construction d'un dérivateur à partir d'un $(\mathcal{M}, W)$, qu'il faut prendre l'équivalence fine, fibre par fibre. À ce sujet, il faut faire attention que les foncteurs verticaux du diagramme p.~8
\[ \und{Fibsc}\,S \longrightarrow \und{Fib}\,S, \quad \und{Cofibsc}\,S \longrightarrow \und{Cofib}\,S, \]
quoique essentiellement surjectifs et fidèles, ne sont pas pleinement fidèles. Ils le seraient, et mêmes des équivalences de catégories, si dans $\und{Fib}\,S$ resp. $\und{Cofib}\,S$ je me bornais aux $S$-foncteurs \textit{cartésiens} resp. \textit{cocartésiens}. Donc on peut interpréter, à équivalence près, $\und{Fibsc}\,S$ et $\und{Cofibsc}\,S$ comme $\und{Fib}\,S$ et $\und{Cofib}\,S$, où on aurait restreint les morphismes aux seuls $S$-foncteurs \textit{cartésiens} resp. \textit{cocartésiens.} Il semblerait donc qu'en restreignant ainsi les flèches, mais en restreignant de façon plus draconienne encore le localiseur, on trouve des catégories localisées équivalentes. % [sic] et mêmes des équivalences -> et même des équivalences, adverb "même" printed with plural s

\dgrblock{dgrb174}{What the path replacement actually proves}
\begin{dcomment}{exposition}{p24}{What the path replacement actually proves}
Grothendieck wants a category over S whose strict fibres already have the homotopy types of the homotopy fibres. The replacement jX : X \ensuremath{-}\ensuremath{\to}C(X) achieves this geometrically: jX is a weak homotopy equivalence, is natural over S, and C(X) \ensuremath{-}\ensuremath{\to}S is perfect. The difficulty is to decide which weak equivalences this replacement preserves. Two assertions must be separated, because one is formal and the other detects a substantive property of the coefficient theory.

Write W tot for maps over S whose underlying functor lies in W. Assuming W\ensuremath{\infty}\ensuremath{\subseteq}W and two-outof-three, naturality gives

f \ensuremath{\in}W tot

\ensuremath{\Leftarrow}\ensuremath{\Rightarrow}C(f) \ensuremath{\in}W tot C(f)jX = jY f, S . Thus no base-change axiom is needed to make C a functor between the coarse relative categories. If F \ensuremath{\subseteq} Cat/S is full and contains every perfect object, C factors through F. The two natural transformations j on Cat/S and its restriction to F become invertible after coarse localization. Consequently the inclusion of F induces an equivalence. This proves the assertion of manuscript page 25 under the stated elementary hypotheses, independently of the conjectural pullback assertion on the present page.

The stronger assertion is that C(f) is a W-equivalence on every fibre whenever f \ensuremath{\in}W tot S . This does require the implication later named L5 on manuscript page 43. The preceding argument with a contractible covering base explains a sufficient condition precisely. For a chosen s \ensuremath{\in}S, take a perfect S\ensuremath{\prime} \ensuremath{-}\ensuremath{\to}S with contractible total category and an object over s; it suffices to cover the component containing s. If pullback along this map preserves W, the resulting total map X\ensuremath{\prime} \ensuremath{-}\ensuremath{\to}Y \ensuremath{\prime} belongs to W. For homotopy fibrations the inclusions of the chosen fibres into X\ensuremath{\prime} and Y \ensuremath{\prime} are weak homotopy equivalences. Two-out-of-three then gives the desired fibre map in W. The final comparison gives a Wequivalence, not necessarily a weak homotopy equivalence: the sentence immediately before manuscript page 24 compresses these two conclusions.

For W = W\ensuremath{\infty}this is the familiar invariance of a homotopy fibre under equivalence of total spaces over a fixed base. For a coarser W it is an additional demand. The covering calculation at manuscript page 43 below will show its failure even when both maps to the base are discrete fibrations and discrete opfibrations. In particular the phrase “il nous faut absolument” concerns compatibility with the fibrewise localizer, not the already established compatibility with the coarse one.

\end{dcomment}


\msp{page 27}
C'est assez étrange - on va voir si c'est bien vrai!

\dgrblock{dgrb175}{Strict splittings and Cartesian maps}
\begin{dcomment}{context}{p26}{Strict splittings and Cartesian maps}
The comparison with diagram categories depends on what the morphisms preserve. The definition on manuscript page 4 uses functors preserving the chosen Cartesian lifts strictly. A Cartesian functor between split fibrations need only carry a chosen lift to some Cartesian lift; the latter is isomorphic to the chosen lift at the image object. These conditions give strict and pseudonatural transformations, respectively. They are different conditions on actual functors, before localization.

Here is a concrete failure of fullness. Put S = [1] and let G = C2 = \{1, c\}. The projection S \ensuremath{\times}BG \ensuremath{-}\ensuremath{\to}S has its constant splitting, with chosen lift (u, 1) of u : 0 \ensuremath{-}\ensuremath{\to}1. Define an endofunctor over S which fixes the two fibre groupoids and sends a cross-fibre arrow (u, g) to (u, cg). This is a functor: composition with a vertical arrow is compatible with the formula because G is abelian. Every arrow over u is Cartesian, since its group component is invertible. Our endofunctor is therefore Cartesian, but it does not preserve the splitting, since (u, 1) is sent to (u, c). Restricting the target category to Cartesian functors does not make the forgetful functor from strictly split morphisms full.

There is a separate qualification concerning objects. The functor BC4 \ensuremath{-}\ensuremath{\to}BC2 induced by reduction modulo 2 is a Grothendieck fibration: each base arrow has a lift and every arrow is Cartesian. A splitting would be a homomorphic section C2 \ensuremath{-}\ensuremath{\to}C4. Such a section cannot exist, because each lift of the nonidentity element has order 4. Hence this fibration is not isomorphic over BC2 to an object with a splitting. This is the group-extension example already explained in Grothendieck (1971, Exposé VI, \S{}9). An equivalence over the base to a split replacement is the appropriate statement; essential surjectivity in the ordinary category Cat/S would ask for an isomorphism of categories over the base.

These distinctions do not damage the intended homotopical comparison. Equivalences of fibre categories belong to W\ensuremath{\infty}, and strict replacements can be connected by such equivalences. What needs proof is precisely that replacing coherent choices by strict ones induces the correct localized comparison. Grothendieck encounters that issue directly in the failure of the proposed unit on manuscript page 30 and returns to it on pages 55--60. The underlying classical language comes from the descent formalism of SGA 1; the new task here is to control the weak equivalences during this passage. We retain the printed assertion and read it with this necessary qualification, rather than changing the category of morphisms silently.

\end{dcomment}


\subsection*{III) Construction de $\Phi_0, \Psi_0 :$ \begin{tikzcd}[ampersand replacement=\&, baseline=(\tikzcdmatrixname-2-1.base), row sep=small, column sep=normal]
\& \und{Fibsc}_0S \\
\Cat/S \arrow[ur] \arrow[dr] \& \\
\& \und{Cofibsc}_0S.
\end{tikzcd}}

On doit avoir
\[ \Phi_0(X) = \Phi r C(X) \qquad\qquad\qquad \Psi_0(X) = \Psi r C(X), \]
où
\[ r : \und{Parf}\,S \longrightarrow \Cat/S \]
est l'inclusion canonique. En tant que foncteurs à valeurs dans \underline{Fibsc} resp. \underline{Cofibsc}, cela définit déjà $\Phi_0$, $\Psi_0$ par
\[ \Phi_0 = \Phi r C \qquad\qquad\qquad \Psi_0 = \Psi r C, \]
\printedpage{17}
% head: A. Grothendieck : Dérivateurs, Ch. VI, Édition des Universités de Montpellier II et Paris VII
% foot: Version 9 novembre 2001
mais il faut voir que $\Phi_0$, $\Psi_0$ prennent leurs valeurs dans les sous-catégories pleines à indice 0. Cela équivaut donc au Lemme suivant, appliqué à $Y = rC(X)$:

\textbf{Lemme.} \textit{Soit $Y$ parfait sur $S$ (ou seulement un \textup{Hot}-fibré sur $S$). Alors $\Phi(Y)$ est dans $\und{Fibsc}_0S$, $\Psi(Y)$ dans $\und{Cofibsc}_0S$.}

Il faut prouver (pour $\Psi(Y)$) que si $s \longrightarrow t$ est une flèche dans $S$, alors
\[ \Psi(Y)_s \longrightarrow \Psi(Y)_t \]
est une $W$-équivalence. Or cette flèche s'identifie à
\[ Y/s \longrightarrow Y/t \]
et on sait bien que c'est une Hot-équi\msp{page 28}valence, par définition des Hot-fibrés, donc c'est une $W$-équivalence.

On a donc les foncteurs
\[
\begin{tikzcd}[column sep=large]
\und{Fibsc}_0S \arrow[r, shift left, "p"] & \Cat/S, \arrow[l, shift left, "\Phi_0"]
\end{tikzcd}
\]
il faut voir qu'ils satisfont aux conditions générales du lemme p.~11 (et ce sera pareil pour $q$, $\Psi_0$). On sait déjà que $p$ est compatible avec les localiseurs. Prouvons le pour $\Psi_0$. Soit
\[ f : X \longrightarrow Y \]
dans $\mathrm{Cat/S}$ une $W$-équivalence, alors $rC(f)$ l'est aussi. Il suffit de prouver ceci : % [sic] Cat/S (upright S) -> Cat/$S$, the S printed upright unlike elsewhere

Soit
\[ q : X' \longrightarrow Y' \]
un morphisme dans $\und{Fibhot}\,S$ (par exemple un morphisme de $\und{Parf}\,S$), qui soit $W$-équivalence, donc \texttt{[qui induit]} des $W$-équivalences sur les fibres. Alors $\Psi(q) : \Psi(X') \longrightarrow \Psi(Y')$ est une $W$-équivalence fibre par fibre, i.e.
\[ X'/s \longrightarrow Y'/s \]
une $W$-équivalence pour tout $s \in S$ (a été vu dans le corollaire 2, page 20).

\msp{page 29}
Considérons
\[ p\Phi_0 : \Cat/S \longrightarrow \Cat/S, \]
ça transforme donc $X$ en $\Phi C(X)$, regardé comme objet de $\Cat/S$ sans plus. On a un morphisme fonctoriel
\[ X \xrightarrow{j_X} C(X) \qquad\qquad\qquad j : \mathrm{id}_{\Cat/S} \longrightarrow C, \]
\printedpage{18}
% head: A. Grothendieck : Dérivateurs, Ch. VI, Édition des Universités de Montpellier II et Paris VII
% foot: Version 9 novembre 2001
d'où un morphisme fonctoriel
\[ \Phi(X) \xrightarrow{\Phi(j_X)} \Phi C(X) \qquad\qquad\qquad \Phi * j : \Phi \longrightarrow \Phi_0 \]
et également un morphisme fonctoriel
\[ X \xrightarrow{i_X} \Phi(X), \]
d'où en composant
\[
\begin{tikzcd}
X \arrow[r, "i_X"] \arrow[rr, bend right, "\alpha_X"'] & \Phi(X) \arrow[r, "\Phi(j_X)"] & \Phi C(X),
\end{tikzcd}
\]
i.e.
\[ \alpha : \mathrm{id}_{\Cat/S} \longrightarrow p\Phi_0. \]
Pour $X$ fixé dans $\Cat/S$, on sait que $j_X$ est dans $W$, donc aussi $\Phi(j_X)$ (quoique non fibre par fibre), cf. cor.~1 page 20. On sait aussi que $i_X$ est dans $W$, donc aussi le composé $\alpha_X$.

Considérons
\[ \Phi_0 p : \und{Fibsc}_0\,S \longrightarrow \und{Fibsc}_0\,S\ \ . \]
On part de $X$ dans $\und{Fibsc}_0\,S$, on oublie le scindage et le fait même qu'il soit Cat-fibrant
\msp{page 30}
(on a appliqué $p$), on lui applique $\Phi_0$, i.e. on regarde $\Phi CX$ comme objet de $\und{Fibsc}\,S$ (et même de $\und{Fibsc}_0\,S$). On veut définir encore
\[ X \xrightarrow{\beta_X} \Phi_0 p(X) = \Phi C(X) \qquad\quad \text{\textbf{dans} } \und{Fibsc}_0\text{\textbf{,}} \] % ?? uncertain: weight of "Fibsc_0" after bold "dans" -- regular or bold
fonctoriel en $X$. Or on a comme tantôt
\[ X \xrightarrow{j_X} C(X), \]
d'où en appliquant $\Phi$
\[ \Phi(X) \xrightarrow{\Phi(j_X)} \Phi C(X), \]
mais regardé cette fois comme flèche dans $\und{Fibsc}_0\,S$ (N.B. on a bien $\Phi(X)$ dans $\und{Fibsc}_0\,S$, car $X$, provenant de $\und{Fibsc}_0\,S$, est dans $\und{Fibhot}\,S$, et on applique la remarque `Cor.~3' page 21.) D'autre part, on a une flèche dans $\Cat/S$
\[ X \xrightarrow{i_X} \Phi(X), \]
est-elle une flèche dans $\und{Fibsc}_0(S)$, i.e. est-elle cartésienne? On voit tout de suite que non, que l'on a seulement, pour une flèche $u : s \longrightarrow t$ dans $S$, un homomorphisme de commutation (pas un isomorphisme) dans

% ---- from Der-ChVI-345_p021-025.tex
\printedpage{19}
% head: A. Grothendieck : Dérivateurs, Ch. VI, Édition des Universités de Montpellier II et Paris VII
% foot: Version 9 novembre 2001

\[
\begin{tikzcd}[row sep=large, column sep=large]
X_s \arrow[r, "u_*^X"] \arrow[d, "i_s"] & X_t \arrow[d, "i_t"] \\
X/s \arrow[r, "(\mathrm{id}_X)/u", "= u_*^{\Phi(X)}"'] \arrow[ur, bend left=30, shorten <=28pt, shorten >=12pt, "\alpha"] & X/t
\end{tikzcd}
\]

\msp{page 31}
Donc il me faut remballer! On peut considérer $i_X$ au mieux comme une flèche de $\uu{Fib}_0\, S$, et en prenant les composés on trouve une flèche dans $\uu{Fib}_0\, S$ :
\[
\begin{tikzcd}
X \arrow[r, "i_X"] \arrow[rr, bend right=25, "\beta_X(=\alpha_X)"'] & \Phi(X) \arrow[r, "\Phi(j_X)"] & \Phi C(X),
\end{tikzcd}
\]
i.e. (regardant maintenant $\Phi_0$ comme foncteur à valeurs dans $\uu{Fib}_0\, S$) un homomorphisme de foncteurs
\[
\beta : \mathrm{id}_{\uu{Fib}_0} S \longrightarrow \Phi_0 p_0, % [sic] \mathrm{id}_{\uu{Fib}_0} S -> \mathrm{id}_{\uu{Fib}_0 S}, the S is printed on the baseline, not in the subscript
\]
où $p_0 : \uu{Fib}_0 S \hookrightarrow \mathrm{Cat}/S$ est l'inclusion canonique. Ce $\beta$ est dans $W$ argument par argument, et de même pour $\alpha$ page 29, interprété cette fois comme
\[
\mathrm{id}_{\mathrm{Cat}/S} \longrightarrow p_0\Phi_0.
\]
Ce qu'on a prouvé finalement est ceci:

\dgrblock{dgrb176}{Why the obvious unit fails to be Cartesian}
\begin{dcomment}{exposition}{p30}{Why the obvious unit fails to be Cartesian}
Grothendieck has found a natural weak equivalence of total categories and now asks it to be a morphism of coherent diagrams. The request is stronger. For the fibration construction \ensuremath{\Phi}(X) the fibre at s is s\textbackslash{}X, whereas the obvious map sends x \ensuremath{\in}Xs to (x, 1s). A Cartesian lift u\ensuremath{*}y \ensuremath{-}\ensuremath{\to}y in X, lying over u : s \ensuremath{-}\ensuremath{\to}t, is sent to an arrow whose source is (u\ensuremath{*}y, 1s). The chosen Cartesian lift in \ensuremath{\Phi}(X) ending at (y, 1t) instead starts at (y, u). These two objects are connected by the lift in X, but in general they are not isomorphic in s\textbackslash{}X.

The smallest example is X = S = [1] with the identity projection. Every arrow of X is Cartesian over S. In 0\textbackslash{}[1] the two objects 10 : 0 \ensuremath{-}\ensuremath{\to}0 and u : 0 \ensuremath{-}\ensuremath{\to}1 are joined by a noninvertible arrow. Thus the image of the Cartesian arrow u under iX is not Cartesian in \ensuremath{\Phi}(X). Indeed, in the Grothendieck construction of a strict diagram, an arrow over u is Cartesian exactly when its vertical comparison component is invertible. Here that comparison component is the displayed noninvertible arrow. The obstruction remains even though both comma categories are contractible.

The square printed across manuscript pages 30--31 uses X/s, X/t and covariant symbols u\ensuremath{*}. Those are the symbols of the dual construction \ensuremath{\Psi}, while the surrounding discussion concerns \ensuremath{\Phi}. The intended argument is correct with s\textbackslash{}X, t\textbackslash{}X and contravariant transport, or after dualizing the entire discussion. This variance slip should not be turned into a purported failure of the comma construction itself.

Passing to the full category Fib0S permits the map iX, because its morphisms are arbitrary functors over S. The natural weak equivalence X \ensuremath{-}\ensuremath{\to}\ensuremath{\Phi}C(X) then proves the coarse comparison in the “théorème-remballage”. The theorem does not yet compare strict coherent diagrams by fibrewise weak equivalences. This is the exact boundary between what the argument has proved and what the next attempt must establish. The later Cartesian comparison is directed towards X, rather than away from it; its construction supplies the missing information about coherence.

\end{dcomment}


\textbf{Théorème-remballage.} \textit{Soit $\mathcal{F}$ une sous-catégorie pleine de $\mathrm{Cat}/S$ qui contient une des trois sous-catégories pleines $\uu{Fib}_0 S$, $\uu{Parf}\, S$, $\uu{Cofib}_0 S$, et munissant $\mathcal{F}$ du localiseur $W$ `grossier' (induit par celui de $\mathrm{Cat}/S$). Alors l'inclusion}
\[
\mathcal{F} \longrightarrow \mathrm{Cat}/S
\]
\textit{induit une équivalence sur les localiseurs.}

\msp{page 32}
Par contre, les catégories localisées des quatre catégories de fibrés scindés
\[
\uu{Fibsc}\, S, \quad \uu{Cofibsc}\, S, \quad \uu{Fibsc}_0 S, \quad \uu{Cofibsc}_0 S
\]
échappent à nos arguments. \textit{Et pour cause!} Il est totalement ouvert, sur $S$ connexe disons (p. ex. $S = \mathbf{\Delta}^1$), que pour un type d'homotopie relatif $\xi$ sur $S$, les types induits $\xi_s$ ($s \in S$) forment un système local sur $S$, à valeurs dans $\mathrm{Hot}_W$, comme ce serait le cas si le théorème annoncé au début était vrai. Au contraire, \textit{toute} flèche de $\mathrm{Hot}_W$ doit provenir d'un objet de $\mathrm{Hot}_W(\mathbf{\Delta}^1)$. Les localisés des 9 autres catégories du diagramme p. 8, et de toutes celles dont il est question dans le théorème qu'on vient de prouver, doivent correspondre `visiblement' aux sous-catégories pleines de $\mathrm{HOT}_W(S)$ formées des HOT-coefficients localement constants sur $S$.

Aussi, je me rends compte que c'était absurde de m'attendre à ce que $\uu{Fibsc}\, S$ et $\uu{Cofibsc}\, S$ donnent des catégories localisées canoniquement équivalentes. Car les éléments $\xi$ de la première
\printedpage{20}
% head: A. Grothendieck : Dérivateurs, Ch. VI, Édition des Universités de Montpellier II et Paris VII
% foot: Version 9 novembre 2001
\msp{page 33}
des deux catégories donnent naissance chacun à un foncteur contravariant sur $S$
\[
\bar{\xi} : S^o \longrightarrow \mathrm{Hot}_W,
\]
ceux $\eta$ de la seconde (qui n'est autre que $(\uu{Fibsc}(S^o))^o$) donnent naissance chacun à un foncteur \textit{covariant}
\[
\bar{\eta} : S \longrightarrow \mathrm{Hot}_W,
\]
ce qui n'est pas pareil! Ce n'est `pareil' que si ce sont des coefficients locaux, ce qui signifie exactement que $\xi$ est dans $\uu{Fibsc}_0 S$ resp. $\eta$ dans $\uu{Cofibsc}_0 S$. On s'attend donc bien \circled{$1^o$} que les flèches
\[
\uu{Fibsc}_0 S \longrightarrow \uu{Fib}_0 S, \quad \uu{Cofibsc}_0 S \longrightarrow \uu{Cofib}_0 S
\]
induisent des équivalences sur les localisés (chose que j'ai été incapable de prouver tantôt) (\setcounter{footnote}{15}\footnote{Les quasi-inverses étant bel et bien donnés par $\Phi_0$, $\Psi_0$ - car [?] par quoi d'autre donc ???}). Par contre, il est pire que déraisonnable de s'attendre à ce qu'il en soit de même pour les foncteurs
\[
\uu{Fibsc}_0 S \longrightarrow \uu{Fibsc}\, S, \quad \uu{Cofibsc}_0 S \longrightarrow \uu{Cofibsc}\, S,
\]
ça devrait donner par localisation les foncteurs pleinement fidèles d'inclusion, des catégories d'objets localement constants dans $\mathrm{HOT}_W(S)$ resp. dans $(\mathrm{HOT}_W(S^o))^o$. Et de même pour
\msp{page 34}
les foncteurs induits sur les localisés par
\[
(*) \qquad \uu{Fibsc}\, S \longrightarrow \uu{Fib}\, S, \quad \uu{Cofibsc}\, S \longrightarrow \uu{Cofib}\, S.
\]
Comme les localisés des catégories but de ces deux foncteurs ont tout l'air de s'identifier aux sous-catégories précédentes d'objets localement constants, il semblerait donc que : \circled{$2^o$} Les foncteurs précédents définissent des rétractions de la catégorie $\mathrm{HOT}_W(S)$ resp. $\mathrm{HOT}_W(S^o)^o$ sur la sous-catégorie des objets localement constants sur $S$, les foncteurs localisés des foncteurs d'inclusion
\[
\uu{Fibsc}_0 S \longrightarrow \uu{Fibsc}\, S, \quad \uu{Cofibsc}_0 S \longrightarrow \uu{Cofibsc}\, S
\]
étant pleinement fidèles et ayant comme images essentielles les sous-catégories en question.

\dgrblock{dgrb177}{A relative type which is not locally constant}
\begin{dcomment}{examples}{p32}{A relative type which is not locally constant}
Grothendieck’s reversal becomes transparent on the walking arrow. Let S = [1] = \{0 \ensuremath{-}\ensuremath{\to}1\}, let X be the one-object category over 0, and let Y = S with its identity projection. The inclusion

f : \{0\} \ensuremath{-}\ensuremath{\to}[1]

is a map of split fibrations over S. The corresponding contravariant diagrams are

X1 = \ensuremath{\varnothing}\ensuremath{-}\ensuremath{\to}X0 = \ensuremath{*},

Y1 = \ensuremath{*}\ensuremath{-}\ensuremath{\to}Y0 = \ensuremath{*}. Both total categories are contractible, so f \ensuremath{\in}W\ensuremath{\infty}. But the fibre at 1 is \ensuremath{\varnothing}\ensuremath{-}\ensuremath{\to}\ensuremath{*}, which is not a weak homotopy equivalence. Fibre evaluation is a functor on the localization by objectwise weak equivalences, and sends isomorphisms there to isomorphisms of homotopy types. Therefore the image of f in that localization is not invertible. Since its image in the coarse localization is invertible, the comparison between these two localizations is not an equivalence.

All four tests of manuscript page 41 can be computed here. For the overcategories,

X/0 = \ensuremath{*},

Y/0 = \ensuremath{*}, X/1 = \ensuremath{*}, Y/1 = [1], so f is a local, or left-fine, weak equivalence. For the undercategories,

1\textbackslash{}X = \ensuremath{\varnothing},

0\textbackslash{}X = \ensuremath{*}, 0\textbackslash{}Y = [1], 1\textbackslash{}Y = \ensuremath{*}. Thus it fails the colocal, or right-fine, test. Pulling back along the point 1 \ensuremath{-}\ensuremath{\to}S gives the same failure, so it is not a universal weak equivalence. The example lies inside the category of split fibrations, hence also inside the category of smooth functors. Its opposite supplies the corresponding example for opfibrations and proper functors. This will also resolve the contradictory sentence on manuscript page 47.

An object equivalent to a locally constant diagram must have isomorphic fibre homotopy types along every base arrow: invertibility is invariant under isomorphism in the homotopy category of diagrams. The first diagram above does not have this property. Accordingly, the full subcategory of locally constant objects is proper even when the classifying space of the indexing category is contractible. Contractibility of the base makes local systems constant up to equivalence; it does not force every coherent diagram to be locally constant.

This gives a definite negative answer to the question on manuscript page 25 for the fibrewise localization of smooth or fibred objects. It also explains why the initial thirteen-category comparison could not have been correct with a single choice of weak equivalences. A derived category of diagrams must retain the noninvertible transition \ensuremath{\varnothing}\ensuremath{-}\ensuremath{\to}\ensuremath{*}, whereas the total homotopy type in this example forgets it.

\end{dcomment}


Cette rétraction m'apparait comme une chose tout à fait mystérieuse. C'est lié au fait que pour le dérivateur très spécial $\mathbf{D} = \mathrm{HOT}$, quand on a un $\xi \in \mathbf{D}(S)$, on en déduit (à $\mathrm{HOT}/S$-équivalence près) un type de diagramme $X$ au dessus de $S$,
\msp{page 35}
qu'on peut regarder alors (via la construction $C$) comme définissant un type d'homotopie relatif \textit{localement constant} sur $S$.

\printedpage{21}
% head: A. Grothendieck : Dérivateurs, Ch. VI, Édition des Universités de Montpellier II et Paris VII
% foot: Version 9 novembre 2001

Donc pour élucider la situation, i.e. les relations entre les 13 catégories localisées provenant du diagramme de la page 8, il reste à regarder si les assertions hypothétiques que je viens de faire sont bien valables, savoir

\begin{enumerate}[leftmargin=*]
\item[\circled{$1^o$}] Le foncteur déduit de
\[
\uu{Fibsc}_0 S \longrightarrow \uu{Fib}_0 S
\]
par passage aux localisés est une équivalence de catégories (et on peut dans cette question remplacer $\uu{Fib}_0 S$ par n'importe quelle autre sous-catégorie pleine de $\mathrm{Cat}/S$ la contenant - cela revient au même), et

\item[\circled{$2^o$}] Le foncteur déduit de
\[
\uu{Fibsc}_0 S \longrightarrow \uu{Fibsc}\, S
\]
par localisation est pleinement fidèle. (Que son image soit exactement formée des objets localement constants du premier localisé, interprété comme $\mathrm{HOT}_W(S)$, est tautologique.)
\end{enumerate}

Il n'y a plus lieu de se soucier du côté
\msp{page 36}
cofibrés, puisque tout ce qu'on pourra en dire revient à ce que je viens d'expliciter, pour $X^o$ au lieu de $X$.

Dans la direction de \circled{$1^o$}, nos arguments précédents à coups de $\Phi_0$, $p$ nous donnent déjà un renseignement utile. Mais il sera plus commode à présent de laisser tomber $\uu{Fib}_0 S$ ou plutôt, de le remplacer par $\uu{Fib}\, S$ (ce qui est licite, on l'a dit dans \circled{$1^o$}), de sorte qu'on a maintenant un diagramme commutatif (traits pleins)
\[
\begin{tikzcd}
\uu{Fibsc}\, S \arrow[dr, "j"'] & & \uu{Fibsc}_0 S \arrow[ll, hook', "i"'] \arrow[dl, "k"'] \\
& \uu{Fib}\, S, \arrow[ur, dotted, bend right=40, "\Phi_0=\psi"'] &
\end{tikzcd}
\]
où $i$ est une inclusion (strictement) pleinement fidèle, $j$ (donc aussi $k$) fidèle. Enfin, on a un foncteur canonique $\psi$ ($= \Phi_0$) (en pointillés). Les localiseurs sur la ligne du haut sont les localiseur fins (pour $\uu{Fibsc}_0 S$ c'est pareil que la $W$-équivalence grossière), sur $\uu{Fib}\, S$ % [sic] les localiseur fins -> les localiseurs fins, missing plural s
par contre on prend la $W$-équivalence grossière. (On sait que le localisé est le même, à équivalence près, que si on se restreint à la sous-catégorie pleine $\uu{Fib}_0 S$, ce qui est assez remarquable. % [sic] "(On sait ... remarquable." -> "(On sait ... remarquable.)", the opening parenthesis is never closed

\msp{page 37}
Tous les quatre foncteurs $i, j, k, \Phi_0$ sont compatibles aux localiseurs, donc passent aux localisés
\printedpage{22}
% head: A. Grothendieck : Dérivateurs, Ch. VI, Édition des Universités de Montpellier II et Paris VII
% foot: Version 9 novembre 2001
\[
\begin{tikzcd}
\mathrm{HOT}_W(S) \arrow[dr, "\bar{j}"'] & & \mathrm{HOT}^{?}_W(S) \arrow[ll, "\bar{\imath}"'] \arrow[dl, "\bar{k}"'] \\
& \mathrm{HOT}^{\mathrm{lc}}_W(S), \arrow[ur, bend right=40, "\bar{\psi}"'] &
\end{tikzcd}
\]
où l'exposant ? dans la notation signifie que je ne sais interpréter cette catégorie, sauf si je prouve que $\bar{k}$ est une équivalence, et l'exposant lc dans la catégorie du bas dois suggerer l'intuition qu'il s'agit en tous cas d'une sorte de `types d'homotopie relatifs' (dans un sens modifié) \textit{localement constants} sur $S$. Le problème est de comprendre ce diagramme. La conjecture naturelle est % [sic] dois suggerer -> doit suggérer, verb agreement and missing accent

\begin{enumerate}[leftmargin=*]
\item[$1^o$)] $\bar{k}$ est une équivalence. (Ce qui implique que $\bar{\psi}$ est l'équivalence quasi-inverse, que $\bar{k}$, $\bar{\psi}$ permettent essentiellement d'identifier $\mathrm{HOT}^{?}_W$ et $\mathrm{HOT}^{\mathrm{lc}}_W$, de sorte que $\bar{\imath}$ et $\bar{\jmath}$ deviennent des foncteurs entre $\mathrm{HOT}_W$ et $\mathrm{HOT}^{\mathrm{lc}}_W$.)

\item[$2^o$)] $\bar{\imath}$ est pleinement fidèle (NB l'image essentielle est bien comprise de toutes façons.)
\end{enumerate}

\msp{page 38}
En direction de $1^o$), je sais pour l'instant que l'on a
\[
\bar{k}\bar{\psi} \simeq \mathrm{id}_{\mathrm{HOT}^{\mathrm{lc}}_W(S)},
\]
donc $\bar{k}$ est essentiellement surjectif (sur les objets), et induit des \textit{surjections} sur les $\mathrm{Hom}(-,-)$, tandis que $\bar{\psi}$ est injectif sur les classes d'isomorphismes d'objets, et sur les $\mathrm{Hom}(-,-)$ (i.e. \textit{fidèle}).

Si les relations entre $\mathrm{HOT}^{?}_W$ et $\mathrm{HOT}^{\mathrm{lc}}_W$ s'avèrent imbitables (p. ex. $\bar{\psi}$ pas même pleinement fidèle, i.e. $\bar{k}$ pas même fidèle), alors la catégorie $\mathrm{HOT}^{?}_W$ mérite d'être larguée dans la corbeille, et on ne gardera que
\[
\begin{tikzcd}[row sep=large]
\mathrm{HOT}_W(S) \arrow[d, shift right=1.5, "\bar{j}"'] \\
\mathrm{HOT}^{\mathrm{lc}}_W(S), \arrow[u, shift right=1.5, "\bar{\psi}'"']
\end{tikzcd}
\]
où on définit $\bar{\psi}'$ comme
\[
\bar{\psi}' = \bar{\imath}\bar{\psi} = \overline{i\psi}.
\]
Et la question (plus modeste) devient à présent :

\begin{enumerate}[leftmargin=*]
\item[$3^o$)] $\bar{\psi}'$ est-il pleinement fidèle? (L'image \textit{essentielle} est la même que celle de $\uu{Fib}_0$ ($\approx$ $\uu{Fib}$), elle est formée sauf erreur des objets localement constants - à vérifier au besoin \dots)
\end{enumerate}

\printedpage{23}
% head: A. Grothendieck : Dérivateurs, Ch. VI, Édition des Universités de Montpellier II et Paris VII
% foot: Version 9 novembre 2001

\msp{page 39}
Il est difficile d'imaginer que $\bar{k}$, ou ce qui revient au même $\bar{\psi}$, puisse être une équivalence, sans qu'il existe au moins une `$W_S$-homotopie'
\[
\mathrm{id}_{\uu{Fibsc}_0 S} \ \leftrightarrow\ \leftrightarrow^{h}\!\leftrightarrow\ \leftrightarrow\ \psi k
\]
dans la catégorie des endofoncteurs de $\uu{Fibsc}_0 S$. À vrai dire, je me rends compte maintenant que la description de $\bar{\psi}$ à partir de $\psi : X \longmapsto \Phi_0 X = \Phi C X$ est inutilement compliquée, car on avait fait le détour par $C(X)$ uniquement pour être sûr de tomber dans $\uu{Fibsc}_0 S$, pas seulement dans $\uu{Fibsc}\, S$. Mais si on laisse tomber $\uu{Fibsc}_0 S$, on s'en fout, et le foncteur $\Phi : X \longmapsto \Phi(X)$ est beaucoup plus naturel. Et si on ne laisse pas tomber, il n'y a pas d'inconvénient de remplacer la source $\uu{Fib}\, S$ par $\uu{Fib}_0 S$ (qui a même localisé, à équivalence près), et sur $\uu{Fib}_0 S$ à nouveau il n'y a pas lieu de s'encombrer de $CX$, pour le plaisir de lui appliquer $\Phi$, car $\Phi(X)$ lui-même est déjà dans $\uu{Fibsc}_0 S$.

\dgrblock{dgrb178}{What a one-sided inverse does and does not give}
\begin{dcomment}{examples}{p38}{What a one-sided inverse does and does not give}
Grothendieck has a useful piece of the comparison: a natural isomorphism k\textasciimacron{} \ensuremath{\psi}\textasciimacron{} \ensuremath{\simeq}1. It proves essential surjectivity of k\textasciimacron{} and faithfulness of \ensuremath{\psi}\textasciimacron{}, and it shows that \ensuremath{\psi}\textasciimacron{} is injective on isomorphism classes of objects. The statement about surjectivity on all morphism sets requires an additional hypothesis.

To locate the issue, let A be the discrete category on two objects a0, a1, and let B = \ensuremath{*}be terminal. There is a unique functor k : A \ensuremath{-}\ensuremath{\to}B; let \ensuremath{\psi} pick a0. Then k\ensuremath{\psi} = 1B strictly. Yet

HomA(a0, a1) = \ensuremath{\varnothing}\ensuremath{-}\ensuremath{\to}HomB(\ensuremath{*}, \ensuremath{*}) = \{1\}

is not surjective. Thus a functor with a section need not be full. In this example \ensuremath{\psi} is even fully faithful, so the parenthetical suggestion that failure of faithfulness of k is equivalent to failure of full faithfulness of \ensuremath{\psi} cannot be justified solely from the retract relation. For that converse take instead A = B \ensuremath{\coprod}BM, with B = \ensuremath{*}and a nontrivial monoid M, let k be the unique functor to \ensuremath{*}and let \ensuremath{\psi} select the first summand. Then \ensuremath{\psi} is fully faithful while k is not faithful.

The precise positive statement is easy and sufficient for the progress claimed so far. Choose

\ensuremath{-}\ensuremath{\to}1. For b, b\ensuremath{\prime} \ensuremath{\in}B and h : b \ensuremath{-}\ensuremath{\to}b\ensuremath{\prime}, the morphism \ensuremath{\psi}(h) maps under k, followed by conjugation \ensuremath{\varepsilon} : k\ensuremath{\psi} with \ensuremath{\varepsilon}, to h. Hence the map

HomA(\ensuremath{\psi}b, \ensuremath{\psi}b\ensuremath{\prime}) \ensuremath{-}\ensuremath{\to}HomB(b, b\ensuremath{\prime})

is split surjective. If \ensuremath{\psi}(h) = \ensuremath{\psi}(h\ensuremath{\prime}), applying k and conjugating gives h = h\ensuremath{\prime}, proving faithfulness. If \ensuremath{\psi}b \ensuremath{\simeq}\ensuremath{\psi}b\ensuremath{\prime}, the same argument gives b \ensuremath{\simeq}b\ensuremath{\prime}. None of these statements identifies an arbitrary a \ensuremath{\in}A with \ensuremath{\psi}ka.

A natural zigzag of weak equivalences between 1 and \ensuremath{\psi}k, sought on the following pages, would supply exactly that missing identification after localization. The full comparison theorem can therefore remain correct although the present formal inference is too strong. We record an error in this intermediate inference, with its corrected scope, rather than a gap in the later comparison once that zigzag has been constructed.

\end{dcomment}


Donc, la question est celle-ci : Si $X$ varie
\msp{page 40}
% layout: the arrows nearrow, uparrow, nwarrow are printed under each "Fibsc_0 S", pointing up to it; the arrow labelled c from X to Phi(X) is a long line with its arrowhead printed in the middle
\[
\begin{array}{l}
\text{dans } \underset{\nearrow\uparrow\nwarrow}{\uu{Fibsc}_0 S}, \quad \text{y-a-t'il une homotopie dans } \underset{\nearrow\uparrow\nwarrow}{\uu{Fibsc}_0 S}, \quad \text{`naturelle', i.e. fonctorielle en } X, \\[2ex]
\qquad\qquad X \xrightarrow{\qquad c \qquad} \Phi(X) = \ \text{Catégorie fibrée scindée de fibres les } s\backslash X
\end{array}
\]
Mais un chemin qui soit formé de flèches
\[
X = X_0 \lline X_1 \lline X_2 \cdots \lline X_n = \Phi(X)
\]
qui soient toutes cartésiennes (et de plus, bien sûr, dans $W$!). Cela signifie donc qu'on aurait des catégories fibrées scindées sur $S$, $X_1$, \dots, $X_{n-1}$, dépendant \textit{fonctoriellement} de $X$, qui relient
\[
X = X_0 \qquad \text{à} \qquad \Phi(X) = X_n
\]
i.e.
\[
X_s \qquad \text{à} \qquad s\backslash X \qquad \text{???}
\]
Je ne vois rien à l'horizon.

% ---- from Der-ChVI-345_p026-030.tex
\printedpage{24}
% head: A. Grothendieck : Dérivateurs, Ch. VI, Édition des Universités de Montpellier II et Paris VII

\msp{page 41}

\dgrblock{dgrb179}{The locally constant comparison is now a theorem}
\begin{dcomment}{status}{p40}{The locally constant comparison is now a theorem}
Grothendieck is asking for morphisms that retain coherent transport while forgetting arbitrary choices of Cartesian lifts. His unsuccessful unit is not evidence that the locally constant subcategory fails to embed fully. It identifies a defect of that particular representative of the comparison. The arguments on manuscript pages 55--60 repair the comparison, and pages 65--69 establish the relevant reflection.

The higher coherence retained by local constancy was illustrated at manuscript page 7. The later theory treats the local constancy question for a general derivator. Cisinski (2009, Proposition 3.10 and Theorem 3.13) proves that locally constant objects form a full subderivator and that its inclusion has both a left and a right adjoint. These are the published numbering; the earlier arXiv version numbers the corresponding statements 3.9 and 3.12. The theorem establishes the positive status of the full embedding and the existence of the two adjoints, without identifying all the coarse localizations of this chapter for an arbitrary W.

That last restriction matters. If W is homological or a finite homotopy truncation, a coarse Wequivalence of total spaces need not be a W-equivalence on homotopy fibres. The local constancy construction in a derivator still makes sense, but the additional identification with the coarse localization of Cat/S uses the stronger hypothesis isolated below. For W\ensuremath{\infty}the categorical path model and the modern description by coherent space-valued diagrams fit together; the commentary on manuscript pages 65--69 gives the precise reflector. This question is therefore settled in its intended homotopy-theoretic setting, while its unqualified extension to every localizer is false.

\end{dcomment}


\begin{flushright}
15.11.90
\end{flushright}

Je crois qu j'ai enfin compris, cette nuit, la situation exacte. Il y a lieu, dans Cat$/S$ et dans les autres catégories du diagramme p. 8 (s'envoyant toutes dans Cat$/S$) d'envisager quatre sortes de ‘$W$-équivalences’. % [sic] qu j'ai -> que j'ai, letter missing

\begin{itemize}
\item[a)] La \textit{$W$-équivalence ‘grossière’} : \begin{tikzcd}[baseline=(s.base), row sep=small, column sep=small]
f : X \arrow[rr] \arrow[dr] & & Y \arrow[dl] \\
& |[alias=s]| S &
\end{tikzcd} dans Cat$/S$ est une $W$-équivalence grossière si en tant que flèche de Cat, c'est une $W$-équivalence. (\footnote[17]{stable par changement de base à la fois lisse et propre, plus généralement si $S' \longrightarrow S$ est une Hot-fibration.}).

\item[b)] La \textit{$W$-équivalence ‘par fibres’} (sur $S$) : $\forall\ s \in S$, $f_s : X_s \longrightarrow Y_s$ est une $W$-équivalence. (\footnote[18]{stable par tout chagement de base $S' \longrightarrow S$}). % [sic] chagement -> changement, letter missing

\item[c)] La \textit{$W$-équivalence ‘locale sur $S$’} : $\forall\ s \in S$, $f/s : X/s \longrightarrow Y/s$ est une $W$-équivalence. Signifie que c'est une $W$-équivalence, et le reste par tout changement de base $S' \longrightarrow S$ qui est un isomorphisme local. (\footnote[19]{devrait être stable par changement de base \textit{lisse} $S' \longrightarrow S$}).

\item[c$'$)] La \textit{$W$-équivalence ‘colocale sur $S$’} : $\forall\ s \in S$, $s\backslash f : s\backslash X \longrightarrow s\backslash Y$ est une $W$-équivalence. Signifie que c'est une $W$-équivalence, et le reste par tout changement de base $S' \longrightarrow S$ qui est un co-isomorphisme local. (\footnote[20]{devrait être stable par changement de base \textit{propre} sur $S$}).
\end{itemize}

Je vais résumer dans un énoncé récapitulatif les relations entre ces quatre notions de $W$-équivalence, en pensant au cas $W = W_\infty$. Une fois la chose faite, je me préoccuperai de dégager les conditions minimales sur un $W$ général qui assurent que c'est vrai pour $W$.

\textbf{Proposition 1.} \textit{Soit} \begin{tikzcd}[baseline=(s.base), row sep=small, column sep=small]
f : X \arrow[rr] \arrow[dr] & & Y \arrow[dl] \\
& |[alias=s]| S &
\end{tikzcd} \textit{dans} Cat$/S$.

\begin{itemize}
\item[1)] \textit{La $W$-équivalence locale sur $S$, et la $W$-équivalence colocale sur $S$, impliquent la $W$-équivalence grossière :}
\[
\begin{tikzcd}
\text{c} \arrow[dr, Rightarrow] & & \text{c}' \arrow[dl, Rightarrow] \\
& \text{a} &
\end{tikzcd}
\]

\printedpage{25}
% head: A. Grothendieck : Dérivateurs, Ch. VI, Édition des Universités de Montpellier II et Paris VII

\item[2)] \textit{Si $X$, $Y$ sont propres sur $S$, ou $X$, $Y$ lisses sur $S$, ou $X$, $Y$ des} Hot-\textit{fibrés sur $S$, alors la $W$-équivalence par fibres implique que c'est une $W$-équivalence} universelle \textit{sur $S$ ($W$-équivalence, et le reste par} tout \textit{changement de base $S' \longrightarrow S$), donc dans ces trois cas, on a}
\[
\begin{tikzcd}
& \text{b} \arrow[dl, Rightarrow] \arrow[dr, Rightarrow] & \\
\text{c} \arrow[dr, Rightarrow] & & \text{c}' \arrow[dl, Rightarrow] \\
& \text{a} &
\end{tikzcd}
\]
(\footnote[21]{Il suffit $X$, $Y$ $W$-propres sur $S$, ou $W$-lisses sur $S$, ou $W$-fibrés sur $S$.}).

Mais on peut préciser chacun de ces trois cas :

\item[(2 fib)] \textit{Si $X$, $Y$ sont} Hot-\textit{fibrés sur $S$, alors la $W$-équivalence grossière de $f$ implique que c'est une $W$-équivalence par fibres, donc les quatre notions de $W$-équivalence sont équivalentes pour $f$.}
\[
\text{a} \Longleftrightarrow \text{c} \Longleftrightarrow \text{c}' \Longleftrightarrow \text{b}
\]
(\footnote[22]{a $\Longrightarrow$ b seulement moyennant L5, voir ci-dessous p. 43.})

\item[(2 lisse)] \msp{page 42}\textit{Si $X$, $Y$ sont lisse sur $S$, alors la $W$-équivalence par fibres sur $S$ équivaut à la $W$-équivalence colocale sur $S$ et implique la $W$-équivalence} universelle \textit{sur $S$, donc dans ce cas} % [sic] lisse -> lisses, plural agreement with X, Y
\[
\begin{tikzcd}
\text{b} \arrow[r, Leftrightarrow] \arrow[d, Rightarrow] & \text{c}' \\
\text{c} \arrow[d, Rightarrow] & \\
\text{a} &
\end{tikzcd}
\]

\dgrblock{dgrb180}{Four tests on a map over the base}
\begin{dcomment}{exposition}{p41}{Four tests on a map over the base}
Grothendieck now separates the total homotopy type from the two possible diagram directions. For p : X \ensuremath{-}\ensuremath{\to}S there are three kinds of fibre to consider. The strict fibre Xs remembers only objects and arrows mapping to s and 1s. The comma category X/s also allows arrows p(x) \ensuremath{-}\ensuremath{\to}s in the base; it varies covariantly with s. The comma category s\textbackslash{}X allows arrows s \ensuremath{-}\ensuremath{\to}p(x) and varies contravariantly. Thus for f : X \ensuremath{-}\ensuremath{\to}Y over S the tests are

total

f \ensuremath{\in}W,

fibrewise fs \ensuremath{\in}W (s \ensuremath{\in}S),

local, or left-fine f/s \ensuremath{\in}W (s \ensuremath{\in}S),

colocal, or right-fine s\textbackslash{}f \ensuremath{\in}W (s \ensuremath{\in}S). The word “local” here refers to these particular comma tests, not to a chosen Grothendieck topology. The names “left” and “right” are the manuscript’s conventions. They must not be interchanged merely because the fibrations are drawn on the left side of the thirteen-category diagram.

For the ordinary weak homotopy equivalences, the role of the comma tests is to detect the maps of coherent diagrams produced by \ensuremath{\Psi} and \ensuremath{\Phi}. A map can preserve the total homotopy type while changing one of those diagrams, as the walking-arrow example at manuscript page 32 proves. Conversely, strict fibres alone do not detect a total equivalence without some condition on the maps to S. For example, the functor from the discrete two-object category \{0, 1\} to [1], over [1], is an isomorphism on both strict fibres but changes the total space from two components to a contractible space. It also fails both comma tests: over 1, and under 0, respectively, it compares a discrete two-point category to a contractible one.

The assertion concerning arbitrary base change of a fibrewise equivalence is formal: the fibre of a pullback over s\ensuremath{\prime} \ensuremath{\in}S\ensuremath{\prime} is the original fibre over its image in S. The assertion concerning base change of a coarse equivalence along every homotopy fibration is not formal for a general W; it is exactly the additional stability later called L6 bis. Footnote 17 is valid for W\ensuremath{\infty}, the setting explicitly adopted for the proposition, and must not be read as a theorem about every localizer.

The local-isomorphism formulation has an elementary explanation in the usual categorical convention of discrete fibrations. If T \ensuremath{-}\ensuremath{\to}S is a discrete fibration, then

(X \ensuremath{\times}S T)/t \ensuremath{\sim}= X/p(t).

An object on the right determines the object of T above its source by the unique lifting property. Thus local equivalence is preserved by this pullback, and relative axiom L4 bis gives total equivalence afterward. Conversely the representable discrete fibrations S/s \ensuremath{-}\ensuremath{\to}S recover the comma tests themselves. Dualizing gives the colocal statement. This explains both the strength and the orientation of the local tests without equating them to arbitrary base-change invariance.

\end{dcomment}


\printedpage{26}
% head: A. Grothendieck : Dérivateurs, Ch. VI, Édition des Universités de Montpellier II et Paris VII
(\footnote[23]{il suffit même $X$, $Y$ $W$-lisses sur $S$.}).

\item[(2 prop)] \textit{Si $X$, $Y$ sont propres sur $S$, alors la $W$-équivalence par fibres sur $S$ équivaut à la $W$-équivalence locale sur $S$ :}
\[
\begin{tikzcd}
\text{b} \arrow[r, Leftrightarrow] \arrow[d, Rightarrow] & \text{c} \\
\text{c}' \arrow[d, Rightarrow] & \\
\text{a} &
\end{tikzcd}
\]
(\footnote[24]{Il suffit même $X$, $Y$ $W$-propres sur $S$.}).
\end{itemize}

\textsc{Démonstration} 1) n'est autre que L4 (\footnote[25]{Non, il y a confusion, L4 est plus faible, cf. p. 44 l'axiome L4 bis.}) et son dual L4$'$ pour un localiseur fondamental $W$, et on a vu que si on suppose L1, L3, alors L4 équivaut à L4$'$, et dans la notion de localiseur fondamental on supposera toujours que c'est vérifié. Également, nous supposons que L3 est vrai même sous la forme plus forte L3 quater (p. 1) : $W_\infty \subseteq W$, i.e. toute Hot-équivalence est une $W$-équivalence.

(2 lisse) Si $X$, $Y$ \texttt{[sont]} lisses sur $S$, alors les $X_s \longrightarrow s\backslash X$, $Y_s \longrightarrow s\backslash Y$ sont des $W_\infty$-équivalences, donc des $W$-équivalences (il suffirait même que $X$, $Y$ soient $W$-lisses$/S$), donc le diagramme commutatif
\[
\begin{tikzcd}
X_s \arrow[r, "f_s"] \arrow[d, "i_s^X"'] & Y_s \arrow[d, "i_s^Y"] \\
s\backslash X \arrow[r, "s\backslash f"] & s\backslash Y
\end{tikzcd}
\]
nous montre que $f_s$ est une $W$-équivalence si et seulement si $s\backslash f$ l'est, d'où b $\Longleftrightarrow$ c$'$. Pour voir que la $W$-équivalence par fibres est même une $W$-équivalence \textit{universelle} sur $S$, on est ramené à montrer que c'est une $W$-équivalence grossière (puisque la lissité \texttt{[est]} stable par changement de base sur $S$), et on gagne puisque c$'$ $\Longrightarrow$ a par 1).

Le cas (2 propre) est dual.

(2 fib) Si $X$, $Y$ sont des Hot-fibrés sur $S$, alors
\printedpage{27}
% head: A. Grothendieck : Dérivateurs, Ch. VI, Édition des Universités de Montpellier II et Paris VII
\msp{page 43}
les inclusions
\[
X_s \longrightarrow s\backslash X, \qquad X_s \longrightarrow X/s,
\]
et de même pour $Y$, sont des $W_\infty$-équivalences, donc des $W$-équivalences (il suffirait même que $X$, $Y$ soient des $W$-fibrations sur $S$ (\footnote[26]{comparer déf. 3, p. 117, dans (I).}), i.e. par définition que les foncteurs précédents soient des $W$-équivalences, et que la même condition soit satisfaite après tout changement de base sur $S$), l'argument dans le cas lisse resp. propre montre déjà que
\[
\begin{tikzcd}
\text{c} & \text{b} \arrow[l, Rightarrow] \arrow[r, Leftrightarrow] \arrow[d, Rightarrow] & \text{c}' \\
& \text{a,} &
\end{tikzcd}
\]
et que la $W$-équivalence par fibres (ou la $W$-équivalence locale sur $S$, resp. colocale sur $S$) implique déjà la $W$-équivalence \textit{universelle} sur $S$. (\footnote[27]{valable en supposant seulement $X$, $Y$ des $W$-fibrés sur $S$.})

\dgrblock{dgrb181}{When strict fibres compute a comma fibre}
\begin{dcomment}{exposition}{p42}{When strict fibres compute a comma fibre}
Grothendieck’s proposition is useful because it states exactly when a strict fibre is a valid homotopical substitute for one of the comma categories. The discussion of smoothness and properness in the commentary on Chapter~I, manuscript pages 52--59, supplies the geometric input; here we use that input to compare weak equivalences of objects over the same base.

If X \ensuremath{-}\ensuremath{\to}S is smooth, the canonical inclusion Xs \ensuremath{-}\ensuremath{\to}s\textbackslash{}X is a weak homotopy equivalence. When both X and Y are smooth, the commutative square

Xs

Ys

s\textbackslash{}X

s\textbackslash{}Y has vertical arrows in W if W\ensuremath{\infty}\ensuremath{\subseteq}W. Two-out-of-three proves the equivalence of the fibrewise and right-fine tests. For proper objects the same argument uses Xs \ensuremath{-}\ensuremath{\to}X/s and proves equivalence with the left-fine test. The proof requires no comparison of total spaces at this stage. The same argument works with W-smooth or W-proper objects when the vertical maps are directly known to be in W.

To obtain universality, make any change of base T \ensuremath{-}\ensuremath{\to}S. Smoothness and properness are preserved, and the fibrewise maps remain in W. The preceding comparison gives a colocal or local W-equivalence over T, respectively. Relative axiom L4 bis, or its dual, now makes the total map over T a W-equivalence. The base change was arbitrary, so the map is universally a W-equivalence. This fills the quantifier in the proof; it is not a new appeal to L6. Here base change acts on a map already known to be a fibrewise equivalence between suitably regular objects.

For homotopy fibrations both comma comparisons are weak homotopy equivalences, universally on the base. Hence the fibrewise, local and colocal conditions coincide and imply universal total equivalence. The reverse implication from a total equivalence is the extra statement later named L5. It holds for W\ensuremath{\infty} by invariance of homotopy fibres, and the next examples show why it cannot be appended for every W.

The modern extension to asphericity structures is given by Maltsiniotis (2005, Corollary 4.16 and Theorem 5.4): Grothendieck fibrations are smooth, and smoothness is characterized by the appropriate base-change property. This is the direct continuation of the present comparison argument. In highercategorical straightening one likewise distinguishes equivalence in the diagram direction from equivalence of underlying total spaces; the latter agrees with the former for locally constant fibrations, not for every right or left fibration.

\end{dcomment}


La question a $\Longrightarrow$ b est par contre nettement plus délicate, et on a vu que ce n'est nullement le cas pour tout localisateur fondamental - sans doute pas pour celui associé au dérivateur Der($\mathbf{Z}$-Mod), en tous cas pas à celui associé à Der($k$-Mod), quand 2 est inversible dans $k$. (Il suffirait même que $\exists\ n \in \mathbf{N}$, $n \geq 2$, tel que $n$ inversible dans $k$ - le contrexemple typique est $X = S^1 \xrightarrow{\ n\ } S^1 = Y = S$, le revêtement étale d'ordre $n$ du cercle.)

L'implication a $\Longrightarrow$ b est donc une condition très forte sur $W$, je vais la noter L5 (l'ancienne L5, cf. page 2, ne méritant pas de survivre, car c'est une côte mal taillée).

\textbf{L5} \textit{Soit} \begin{tikzcd}[baseline=(s.base), row sep=small, column sep=small]
X \arrow[rr, "f"] \arrow[dr] & & Y \arrow[dl] \\
& |[alias=s]| S &
\end{tikzcd} \textit{dans} Cat\textit{, avec $X$, $Y$ des} Hot-\textit{fibrés sur $S$. Si $f$ est une $W$-équivalence (‘grossière’ par rapport à $S$), alors c'est une $W$-équivalence par fibres} (\footnote[28]{Dans XII (p. 12) on voit que si $W$ satisfait à L1, L3 quater, L4 bis (cf. page suivante), L5, alors \fbox{$W = W_\infty$}. Donc c'est parfaitement illusoire, si on suppose L5 de continuer à traîner un $W$.}).

On a vu qu'à un $\varepsilon$ près (‘lemme’ p. 23), L5 est une conséquence de

\textbf{L6} \textit{Soit $f : X \longrightarrow Y$ une $W$-équivalence, $Y' \longrightarrow Y$ un morphisme parfait (i.e. lisse et propre), alors $f' : X' = X \times_Y Y' \longrightarrow Y'$ est une $W$-équivalence.}

\msp{page 44}
Notons la variante bien naturelle

\dgrblock{dgrb182}{The circle covering as a categorical counterexample}
\begin{dcomment}{examples}{p43}{The circle covering as a categorical counterexample}
Grothendieck’s covering of the circle gives a complete obstruction to L5 for homology coefficients in which its degree is invertible. It is useful to realize it inside Cat, so that no passage from topological coverings to categorical fibres is left implicit.

Fix n \ensuremath{\geq}2. Let S = BZ, the one-object category with endomorphism group Z. Let X be the action groupoid for translation of Z on Z/n. Its objects are residues a, and a morphism a \ensuremath{-}\ensuremath{\to}b is an integer m with b = a + m modulo n; composition is addition. The projection p : X \ensuremath{-}\ensuremath{\to}S sends that morphism to m. Unique lifting of each integer with specified source or target makes p both a discrete fibration and a discrete opfibration. Its strict fibre is the discrete set Z/n.

The groupoid X is connected. The automorphism group of the object 0 is nZ, so the inclusion B(nZ) \ensuremath{-}\ensuremath{\to}X is an equivalence. Under this identification, p is the homomorphism nZ ,\ensuremath{\to}Z. Thus both classifying spaces are circles, and the induced map on H1(\ensuremath{-}; k) is multiplication by n after choosing the usual generators. It is the identity on H0 and both sides have zero homology in higher degrees. If k is a nonzero commutative ring in which n is invertible, p is a k-homology equivalence. Regard it now as the map from X \ensuremath{-}\ensuremath{\to}S to S \ensuremath{-}\ensuremath{\to}1S. On strict fibres it induces

H0(\ensuremath{-}; k) : kn \ensuremath{-}\ensuremath{\to}k.

Z/n \ensuremath{-}\ensuremath{\to}\ensuremath{*}, The latter map adds coordinates and has nonzero kernel. Hence it is not a fibrewise homology equivalence, although both objects over S are perfect.

The same example disproves L6 itself, not only the fibre-detection formulation. Let E be the translation groupoid of Z acting on itself, and let E \ensuremath{-}\ensuremath{\to}BZ be the evident discrete fibration and opfibration. There is exactly one arrow between each ordered pair of objects of E, so E is contractible. The objects of X \ensuremath{\times}S E are pairs (a, z); simultaneous translation preserves a \ensuremath{-}z modulo n. Each of the resulting n components has exactly one arrow between any two of its objects and is contractible. Therefore

X \ensuremath{\times}S E \ensuremath{-}\ensuremath{\to}E

induces kn \ensuremath{-}\ensuremath{\to}k on H0 and fails to be a homology equivalence. The base-change map is perfect, as required by L6. The argument explains why replacing “homotopy fibration” by “perfect” cannot repair the homological version.

\end{dcomment}


\dgrblock{dgrb183}{An integral-coefficient version of the obstruction}
\begin{dcomment}{development}{p43}{An integral-coefficient version of the obstruction}
Grothendieck also suspects failure for integral homology, where a nontrivial finite covering of the circle no longer gives a homology equivalence. The following explicit construction settles that case. Its ingredients are standard: a semidirect product and the mapping-torus homology sequence. We checked the discussion of homological local constancy in Cisinski (2009, Scholium 1.26) and the mapping-torus calculation in Hatcher (2002, Example 2.48); we did not find this particular categorical application to the sentence here in those sources. The marking records the application and proof supplied by us, not a priority claim for the group or its homology.

Proposition (integral homology does not detect the fibre). \ours There is a perfect functor p : X \ensuremath{-}\ensuremath{\to}S which induces an isomorphism on integral homology in every degree, while its strict fibre has nonzero first integral homology. Consequently the localizer of integral homology equivalences fails L5.

Proof. Put A = Z[1/2], let h : A \ensuremath{-}\ensuremath{\to}A be multiplication by 2, and put G = A \ensuremath{\rtimes}h Z. Take X = BG, S = BZ and let p be induced by the quotient G \ensuremath{\twoheadrightarrow}Z. Since this homomorphism is surjective, every base arrow admits lifts with either prescribed source or prescribed target; all arrows in these groupoids are invertible and hence Cartesian and cocartesian. Thus p is both a Grothendieck fibration and an opfibration, and is perfect. Its strict fibre is BA.

We compute the fibre’s homology first. The group A is the filtered union of the cyclic subgroups 2\ensuremath{-}rZ. The nerve, its free simplicial chain complex and homology all commute with this filtered union: each finite tuple of group elements occurs at a finite stage, and filtered colimits of abelian groups are exact. Since BZ has the homotopy type of a circle, this gives

(i \ensuremath{\geq}2).

H0(BA; Z) = Z, H1(BA; Z) = A, Hi(BA; Z) = 0 The map induced by h acts by 2 on H1 and by the identity on H0.

The group automorphism h induces a homeomorphism of the realization of the nerve of A. Its mapping torus Th is a bundle over the circle. Its fundamental group is A \ensuremath{\rtimes}h Z after choosing the orientation of the circle, and its higher homotopy groups vanish by the homotopy sequence, since the fibre is K(A, 1). Hence Th is a classifying space for G, compatibly with the projection to BZ. The mapping-torus homology sequence has the terms

\textperiodcentered{} \textperiodcentered{} \textperiodcentered{} \ensuremath{-}\ensuremath{\to}Hi(BA)

\ensuremath{-}\ensuremath{-}\ensuremath{-}\ensuremath{\to}Hi(BA) \ensuremath{-}\ensuremath{\to}Hi(Th) \ensuremath{-}\ensuremath{\to}Hi\ensuremath{-}1(BA) \ensuremath{-}\ensuremath{\to}\textperiodcentered{} \textperiodcentered{} \textperiodcentered{} . In degree 1, 1 \ensuremath{-}h\ensuremath{*}= \ensuremath{-}1 on A is invertible; in degree 0 it is zero on Z. Exactness therefore gives Hi(Th) = 0 for i \ensuremath{\geq}2 and an isomorphism H1(Th) \ensuremath{-}\ensuremath{\to}Z. Naturality with respect to the map of fibres BA \ensuremath{-}\ensuremath{\to}\ensuremath{*}identifies this last isomorphism with the homology map induced by the projection to the circle. Both spaces are connected, so the projection is also an isomorphism on H0. This proves that p is an integral homology equivalence. Its map of fibres to the fibre of 1S is BA \ensuremath{-}\ensuremath{\to}\ensuremath{*}, which fails on H1 = A = 0. \ensuremath{\Box}

The example distinguishes absolute homology of the total category from homology after changing the base. The monodromy multiplies the fibre’s homology by 2, and 1 \ensuremath{-}2 is invertible; its contribution disappears from the total homology without disappearing from the fibre. Thus the obstruction is present even with integral coefficients, independently of inverting the degree of a finite covering.

\end{dcomment}


\textbf{L6 bis} \textit{Comme L6, mais en supposant seulement $Y' \longrightarrow Y$ une} Hot-\textit{fibration.}

\printedpage{28}
% head: A. Grothendieck : Dérivateurs, Ch. VI, Édition des Universités de Montpellier II et Paris VII

Il y a des chances que pratiquement, quand L6 (ou seulement L5) est verifié, L6 bis l'est aussi. % [sic] verifié -> vérifié, missing accent

On va introduire un axiome

\textbf{L7} \ a) \textit{Soit $(f_i : X_i \longrightarrow Y_i)_{i \in I}$ une famille de flèches dans} Cat\textit{, soit $f : \coprod X_i \longrightarrow \coprod Y_i$ la flèche correspondante sur les catégories somme. Pour que $f \in W$, il faut et il suffit que $f_i \in W$ pour tout $i \in I$.}

\hspace*{2em}b) \textit{Soit $f : X \longrightarrow Y$ dans $W$. Si $Y$ est} $0$\textit{-connexe, $X$ aussi.}

Cet axiome équivaut à la conjonction de La ($f \in W \Longrightarrow \pi_0 f$ bijectif) et de la validité du ‘Lemme’ page 23, que je remets ici pour mémoire :

\textbf{Lemme.} \textit{L'axiome L7 précédent sur le localiseur fondamental $W$ équivaut à la condition suivante : $f : X \longrightarrow Y$ est dans $W$ si et seulement si $\pi_0(f)$ est bijectif, et $f$ induit des $W$-équivalences $f_i : X_i \longrightarrow Y_i$ sur les composantes connexes.}

Il me faut revenir encore sur le 1) de la proposition, il n'est pas vrai que ‘c'est’ la condition L4 sur $W$ (que j'admets, ai-je dit), c'est L4 dans le cas particulier où $Y = S$. Réflexion faite, je le vois comme une variante renforcée de L4, qui ne semble pas devoir découler de L4 :

\textbf{L4 bis} \textit{Soit} \begin{tikzcd}[baseline=(s.base), row sep=small, column sep=small]
X \arrow[rr, "f"] \arrow[dr] & & Y \arrow[dl] \\
& |[alias=s]| S &
\end{tikzcd} \textit{dans} Cat\textit{, supposons que $f$ soit une $W$-équivalence locale sur $S$ (i.e. $f/s : X/s \longrightarrow Y/s$ une $W$-équivalence $\forall\ s \in S$), alors $f$ est une $W$-équivalence} (\footnote[29]{équivaut à l'axiome dual, en termes des $s\backslash f$}).

\msp{page 45}
Par la suite, par un ‘localiseur fondamental’, nous entendrons un ensemble
\[
W \subseteq \mathrm{Fl}\,\mathrm{Cat}
\]
satisfaisant (tout au moins) les axioms % [sic] axioms -> axiomes, French spelling
\[
\text{L1, L3 quater, L4 bis, L7,}
\]
i.e.

\dgrblock{dgrb184}{Componentwise locality and relative Theorem A}
\begin{dcomment}{exposition}{p44}{Componentwise locality and relative Theorem A}
Grothendieck separates two axiomatic matters here. The first concerns disjoint components; the second concerns descent over a base that need not be the target of the map being tested. Keeping them separate clarifies which part of the preceding argument each axiom repairs.

The component criterion is exactly equivalent to L7. To prove the nontrivial direction, let f : X \ensuremath{-}\ensuremath{\to}Y belong to W and decompose Y = ` j Yj into its nonempty connected components. Put Xj = f \ensuremath{-}1(Yj), a union of components of X. Then f = ` j fj, so L7(a) implies fj \ensuremath{\in}W for every j. L7(b) implies that every Xj is nonempty and connected. Hence \ensuremath{\pi}0(f) is bijective, and the fj are precisely the component maps. Conversely, if \ensuremath{\pi}0(f) is bijective and every component map belongs to W, L7(a) gives f \ensuremath{\in}W. The empty target creates no exception: a functor into the empty category has empty source.

To recover L7 from the criterion, consider a family (fi). The component map of ` i fi is just a component map of one fi; moreover its map on components is bijective exactly when each \ensuremath{\pi}0(fi) is bijective. The criterion therefore gives both directions of L7(a). If the target of a W-equivalence is connected, its component bijection makes the source connected, giving L7(b). This is a complete proof of the lemma, including the reflection direction in the coproduct assertion.

For the second matter, the original L4 applies to a map X \ensuremath{-}\ensuremath{\to}Y when the overcategories of its fibres over objects of Y are aspherical. In Proposition 1, however, one compares X/s and Y/s for an independent base S. These overcategories need not be aspherical; only the map between them is assumed to be a W-equivalence. Thus the printed sentence identifying the proposition directly with the original L4 does not establish it. L4 bis is exactly the required relative formulation. The dual condition uses s\textbackslash{}X \ensuremath{-}\ensuremath{\to}s\textbackslash{}Y .

The accepted modern definition of a fundamental localizer includes the relative formulation; see Cisinski (2004, \S{}1.1). Its minimal member is W\ensuremath{\infty}by that paper’s Theorem 2.2.11. This gives a precise later framework for the corrected argument. We make no claim here that the absolute and relative formulations are logically independent under every possible package of the earlier axioms; what the text has exhibited is the need for the relative assertion in this proof. Opposite-category invariance then transfers the relative assertion to its dual, using the category-of-simplices comparison discussed earlier in the chapter.

\end{dcomment}


\begin{itemize}
\item[a)] \textit{Saturation} de $W$ (conditions a) b) c) p. 1)

\item[b)] $W_\infty \subseteq W$ $[\subseteq W_0]$ (\footnote[30]{Je ne suis plus sûr s'il est vraiment utile d'imposer $W_\infty \subseteq W$. Prendre au lieu de ça l'axiome de l'objet final, soit b$'$). Si on veut a) b$'$) c) + L5, alors on aura $W \subseteq W_\infty$. \textbf{NB} a) b) c) + L5 implique $W = W_\infty$, cf. annotation page 43.}) (\footnote[31]{cf. d)}).

\item[c)] Si $f : X \longrightarrow Y$ sur $S$ est localement sur $S$ une $W$-équivalence, alors $f \in W$ (L4 bis)
\end{itemize}

% ---- from Der-ChVI-345_p031-035.tex
\printedpage{29}
% head: A. Grothendieck : Dérivateurs, Ch. VI, Édition des Universités de Montpellier II et Paris VII
% foot: Version 9 novembre 2001

\begin{enumerate}[leftmargin=*]
\item[d)] Pour que $f : X \longrightarrow Y$ soit dans $W$, il faut et il suffit que $\pi_0(f)$ soit bijectif, et que les $f_i : X_i \longrightarrow Y_i$ induits sur les composantes connexes soient dans $W$.
\end{enumerate}

Ces axiomes sont vérifiés par tous les $W$ qui je connais, et qu' on aurait envie de regarder comme localiseurs. Par contre, L5, L6, L6 bis (p. 43, 44) sont des axiomes de nature beaucoup plus spéciale. J'ai envie de dire qu'ils distinguent les localiseurs de nature proprement \textit{`homotopique'} de ceux qui sont de nature \textit{`homologique'} (ou `cohomologique'). Il faudrait que je m'assure qu'ils sont satisfaits pour les $W_i$ ($i$-équivalence d'homotopie) (\footnote[32]{C'est \textit{faux}, cf. annotation ci-dessus.}). % [sic] qui je connais -> que je connais, relative pronoun object requires "que"

Le fait que L4 bis, i.e. c) ci-dessus, implique l'axiome dual, résulte du fait que L4 bis implique ceci

\textbf{Proposition 2.} \textit{Soit $f : X \longrightarrow Y$ dans} $\Cat$\textit{, alors $f \in W$ si et seulement si $f^o : X^o \longrightarrow Y^o$ est dans $W$.}

Se voit par l'argument de I, page 75 (cor. 2) - n'utilise que L4, pas même L1, L3.

\msp{page 46}

\dgrblock{dgrb185}{Why the new fibration axiom selects full homotopy theory}
\begin{dcomment}{status}{p45}{Why the new fibration axiom selects full homotopy theory}
Grothendieck’s second L5 is much stronger than the first axiom with that label. The first passed from equivalences of fibres to an equivalence of totals. The new one attempts the reverse implication for homotopy fibrations. The annotations correctly recognize that this is incompatible with the proposed range of examples, but the displayed uniqueness claim needs the standing connectedness condition W \ensuremath{\subseteq}W0. Without it the maximal class of all functors satisfies L1, L3 quater, L4 bis and L5, so those four conditions alone do not force W = W\ensuremath{\infty}.

Here is the mechanism of the uniqueness assertion, retaining W\ensuremath{\infty}\ensuremath{\subseteq}W \ensuremath{\subseteq}W0, two-out-of-three, the new L5, and the path replacement already available in the chapter. Suppose first that T \ensuremath{-}\ensuremath{\to}\ensuremath{*} belongs to W. Its classifying space is connected by W \ensuremath{\subseteq}W0. For a chosen object t form the perfect replacement P = CT (\ensuremath{*}) \ensuremath{-}\ensuremath{\to}T of the inclusion of that object. Its total category is weakly contractible. Two-out-of-three implies P \ensuremath{-}\ensuremath{\to}T \ensuremath{\in}W. Applying L5 to this map between the homotopy fibrations P \ensuremath{-}\ensuremath{\to}T and T \ensuremath{-}\ensuremath{\to}1T shows that its fibre F \ensuremath{-}\ensuremath{\to}\ensuremath{*}belongs to W. The fibre models the based loop space of BT. Repeating the argument with F and successive fibres gives categorical models for all iterated loop spaces, each W-aspherical and hence connected. Therefore

\ensuremath{\pi}n(BT, t) = \ensuremath{\pi}0(\ensuremath{\Omega}nBT) = 0

(n \ensuremath{\geq}1). Together with connectedness this says that T is weakly contractible.

Now let f : X \ensuremath{-}\ensuremath{\to}Y belong to W. Replace it by the perfect map CY (X) \ensuremath{-}\ensuremath{\to}Y . Its total map is still in W, and L5 makes every strict fibre W-aspherical. What we have just proved makes those fibres weakly contractible. Since this is a homotopy fibration, its homotopy fibres are weakly contractible, so its total map is a weak homotopy equivalence; the replacement X \ensuremath{-}\ensuremath{\to}CY (X) is also one. Thus f \ensuremath{\in}W\ensuremath{\infty}. This proves W = W\ensuremath{\infty}. The argument uses the previously established homotopy-fibre property of the path replacement, not a new assertion that arbitrary strict fibres compute homotopy fibres. The later abstract characterization in Cisinski (2004, Theorem 2.3.6) likewise includes preservation of connected components.

One can see the failure for every finite truncation directly. Let W\ensuremath{\leq}n mean maps inducing isomorphisms on \ensuremath{\pi}i for 0 \ensuremath{\leq}i \ensuremath{\leq}n, with all basepoints understood. Choose a nonzero abelian group A and a categorical model S of K(A, n + 1). The perfect replacement P = CS(\ensuremath{*}) \ensuremath{-}\ensuremath{\to}S has contractible total space, so belongs to W\ensuremath{\leq}n, while its fibre models K(A, n) and is not W\ensuremath{\leq}n-aspherical. For n = 0 that fibre is the nontrivial discrete group A. This disproves L5. Pulling P \ensuremath{-}\ensuremath{\to}S back along itself yields a total space weakly equivalent to \ensuremath{\Omega}BS, because the base P is contractible and the map is a homotopy fibration. The projection P \ensuremath{\times}S P \ensuremath{-}\ensuremath{\to}P therefore fails to belong to W\ensuremath{\leq}n, disproving L6 even for this perfect base change. The finite truncations are homotopical theories, but they do not satisfy the unrestricted fibre-detection requirement of this L5.

\end{dcomment}


Je vais essayer de résumer ce que je crois avoir compris :

\begin{enumerate}[leftmargin=*]
\item[1)] La $W$-équivalence `par fibres' n'a d'intérêt (pour construire des catégories localisées) que dans le cas où on sait qu'elle implique déjà, soit la $W$-équivalence locale, soit la $W$-équivalence colocale (soit même la $W$-équivalence universelle) sur $S$, donc quand on suppose $X$, $Y$ pour le moins tous deux soit $W$-lisses, soit $W$-propres, soit $W$-fibrés sur $S$.

\item[2)] On pourrait appeler la $W$-équivalence locale sur $S$ la \textit{`$W$-équivalence fine à gauche'} sur $S$, et la $W$-équivalence colocale sur $S$ la \textit{`$W$-équivalence fine à droite'} sur $S$. Elles coïncident dans $\uc{Fibhot}\,S$, a fortiori dans $\uc{Parf}\,S$ et dans les quatre catégories (cf. diagramme p. 8) $\uc{Fib}_0S$, $\uc{Fibsc}_0S$, $\uc{Cofib}_0S$, $\uc{Cofibsc}_0S$ (\footnote[33]{On pourrait les appeler les six catégories \textit{internes} du diagramme.}). Par contre, elles ne coïncident pas sur les catégories \textit{bordantes} du diagramme, notamment sur $\uc{Liss}\,S$, $\uc{Prop}\,S$, $\uc{Fib}\,S$, $\uc{Cofib}\,S$, $\uc{Fibsc}\,S$, $\uc{Cofibsc}\,S$. À cet égard, c'est la \textit{$W$-équivalence fine} \textbf{droite} (en termes des $s\backslash X$) qui est intéressante sur le côté (hélas) gauche du diagramme, i.e. pour $\uc{Liss}\,S$, $\uc{Fib}\,S$, $\uc{Fibsc}\,S$, alors que c'est la \textit{$W$-équivalence fine} \textbf{gauche} (en termes des $X/s$) qui est pertinente du côté droit, pour $\uc{Prop}\,S$, $\uc{Cofib}\,S$, $\uc{Cofibsc}\,S$. On n'aura jamais à regarder simultanément les deux équivalences fines sur une même catégorie du diagramme, à la seule exception de $\Cat/S$ tout entier.
\end{enumerate}

\msp{page 47}

\begin{enumerate}[leftmargin=*]
\item[3)] La $W$-équivalence grossière, elle, est `bonne partout' - mais ça dépend pour qui. Pour les 6 catégories qui sont au dessus de $\uc{Fibhot}\,S$, i.e.
\[
\uc{Fibhot}\,S,\quad \uc{Parf}\,S,\quad \uc{Fib}_0S,\quad \uc{Fibsc}_0S,\quad \uc{Cofib}_0S,\quad \uc{Cofibsc}_0S,
\]
\printedpage{30}
% head: A. Grothendieck : Dérivateurs, Ch. VI, Édition des Universités de Montpellier II et Paris VII
% foot: Version 9 novembre 2001
elle coïncide avec toutes les autres équivalences (moyennnant l'axiome L5), pour les trois catégories bordantes supérieures gauches % [sic] moyennnant -> moyennant, triple n
\[
\uc{Liss}\,S,\quad \uc{Fib}\,S,\quad \uc{Fibsc}\,S
\]
du diagramme, elle coïncide avec la $W$-équivalence fine droite, et dualement sur
\[
\uc{Prop}\,S,\quad \uc{Cofib}\,S,\quad \uc{Cofibsc}\,S
\]
elle coïncide avec la $W$-équivalence fine gauche. Quant à $\Cat/S$ tout entier, c'est la seule des catégories du diagramme où il y ait \textit{trois} notions intéressantes différentes de $W$-équivalence, à savoir : la grossière, et les deux fines (à gauche, à droite). Sur les autres six catégories bordantes du diagramme, il y en a en principe deux : la grossière, et une spécifiée des deux fines (défendu de regarder l'autre - à ses risques et périls \ldots) : la droite du côté gauche, la gauche du côté droit. (NB Je devrais refaire le diagramme, décidément, en mettant les fibrations à droite, pas à gauche.) Quant aux six catégories internes au diagramme, dans elles il n'y a plus qu'une seule notion de $W$-équivalence (puisque sur elles les quatre notions coïncident).
\end{enumerate}

\msp{page 48 (vide)}

\msp{page 49} (\textbf{NB} il n'y a pas de page 48)

\dgrblock{dgrb186}{Reading the corrected classification}
\begin{dcomment}{exposition}{p47}{Reading the corrected classification}
Grothendieck is now trying to separate the homotopy theory of a diagram from the homotopy type of its total category. This separation is the reason for retaining three localizations. Write WS for equivalences of total categories, W d S for the condition that every s\textbackslash{}X \ensuremath{\to}s\textbackslash{}Y belongs to W, and W g S for the analogous condition on X/s \ensuremath{\to}Y/s. The superscripts record the two comma directions, not the placement of a category in the drawing. One sentence in the present paragraph must be corrected mathematically. The assertion that WS coincides with W d S on the smooth or fibred side, and dually with W g S on the proper or cofibred side, is false. The edition prints it without a negation. On the fibred side it is fibrewise equivalence that coincides with W d S; total-category equivalence remains a weaker condition. The walking-arrow example in the commentary at manuscript p. 32 proves the distinction: \{0\} ,\ensuremath{\to}[1], viewed over [1], is an equivalence of total homotopy types, but its fibres over 1 are respectively empty and nonempty. The same paragraph subsequently says that there are two notions on each border, and p. 49 uses the correct distinction.

On the six internal categories the equality with WS uses the new L5 of p. 43. This is different from the earlier axiom carrying that name. Under L5, a map between homotopy fibrations which is an equivalence on total categories is an equivalence on fibres. The ordinary replacement arguments compare all the coarse localizations without this additional axiom, and comma resolutions independently compare the fine theories. L5 is needed when we identify the coarse class on internal objects with their fine class, and thereby realize the coarse theory as a fully faithful locally constant part of the fine theory. This distinction prevents the corrected theorem from inheriting the false equality in its provisional summary.

\end{dcomment}


Je résume maintenant ce que je crois avoir compris sur les catégories localisées correspondantes. Il y en a, à équivalence près, \textit{trois} différentes seulement en tout et pour tout, quand on prend les localiseurs `admissibles' dans les catégories, comme je viens de le dire (c'étant ça en somme l'intérêt objectif des prescriptions!) De façon précise :

\begin{enumerate}[leftmargin=*]
\item[$1^o$)] La \textit{$W$-équivalence grossière} donne `la même' catégorie localisée pour les 13 catégories du diagramme, i.e. toutes les flèches de transition du diagramme induisent des équivalences pour cette $W$-équivalence grossière. Le localisé sera noté \fbox{$\HOT^{\mathrm{lc}}_W(S)$}, où l'exposant lc signifie: ($W$-types d'homotopie relatif sur $S$) \textit{`localement constants.'} (Il ne me reste qu'un petit doute pour $\uc{Fibsc}\,S$ et $\uc{Cofibsc(S)}$ (\footnote[34]{Ce doute n'est pas fondé, tout marche parfaitement!}), si l'équivalence grossière y est bel et bien raisonnable. On verra en faisant la démonstration.)

\item[$2^o$)] La \textit{$W$-équivalence fine à droite} donne le même résultat pour toutes les catégories du diagramme où elle est considérée comme admise, \textit{et où elle ne coïncide pas avec la $W$-équivalence grossière} déjà considérée, donc sur les quatre catégories bordantes gauches
\[
\Cat/S,\quad \uc{Liss}\,S,\quad \uc{Fib}\,S,\quad \uc{Fibsc}\,S
\]
(\footnote[35]{\textbf{NB} Sur les trois dernières cette équivalence est la $W$-équivalence par fibres.}). Le localisé est noté \fbox{$\HOT_W(S)$}.
\end{enumerate}

\msp{page 50}

\printedpage{31}
% head: A. Grothendieck : Dérivateurs, Ch. VI, Édition des Universités de Montpellier II et Paris VII
% foot: Version 9 novembre 2001

\begin{enumerate}[leftmargin=*]
\item[$2'^o$)] Dualement, la \textit{$W$-équivalence fine gauche} donne le même résultat sur toutes les catégories du diagramme où elle est admise, et où elle ne coïncide pas avec la grossière, i.e. sur les 4 catégories
\[
\Cat/S,\quad \uc{Prop}\,S,\quad \uc{Cofib}\,S,\quad \uc{Cofibsc}\,S.
\]
Il y a lieu de noter la catégorie des fractions commune par le sigle
\[
\HOT_W(S^o)^o \qquad \text{(ou aussi $\HOT^o_W(S)$)},
\]
puisqu'elle est isomorphe à celle notée par ce sigle.
\end{enumerate}

Si on veut des notations sans ambiguïté, désignant des catégories bien définies, et non pas seulement `à équivalence près', je choisirai
\[
\HOT_W(S) \defeq \uc{Fibsc}(S)(W^{\mathrm{d}}_S)^{-1}
\]
d'où
\[
\HOT_W(S^o) \qquad \underbrace{\simeq}_{\text{\footnotesize isomorphisme de catégories}} \qquad \uc{Cofibsc}(S)(W^{\mathrm{g}}_S)^{-1},
\]
où $W^{\mathrm{d}}_S$, $W^{\mathrm{g}}_S$ désignent les ensembles de flèches correspondants à la $W$-équivalence fine droite resp. gauche. Pour $\HOT^{\mathrm{lc}}_W(S)$, je préfère une définition aussi symétrique que possible en $S$, $S^o$, et de plus une catégorie aussi petite que possible, ce qui me conduit à choisir
\[
\HOT^{\mathrm{lc}}_W(S) \defeq (\uc{Parf}\,S)W^{-1}_S,
\]
où on a noté par $W_S$ les $W$-équivalences dans $\uc{Parf}\,S$. (Inutile cette fois de préciser lesquelles - toutes les quatre \texttt{[notions]} coïncident dans $\uc{Parf}\,S$ !)

\msp{page 51}

\dgrblock{dgrb187}{The covariant theory and the extra opposite}
\begin{dcomment}{exposition}{p50}{The covariant theory and the extra opposite}
Grothendieck wants a notation which makes the symmetry between fibrations and opfibrations immediate. There is a genuine symmetry, but the edition alternates between two incompatible formulas. With the definition fixed on manuscript p. 3, the precise identities are

[Cofibsc(S)][(W g

HOTW (S) = [Sop, Cat][W \ensuremath{-}1 S)\ensuremath{-}1] \ensuremath{\simeq}[S, Cat][W \ensuremath{-}1 obj] = HOTW (Sop). obj], Here Wobj means objectwise membership in W. There is no opposite on the value category in the last expression. The displayed isomorphism in the middle of p. 50 has this correct form; the preceding prose and the editorial correction at p. 51 insert an additional opposite. We retain those readings as part of the source, while using the comma-category models to determine the variance of every assertion.

Indeed, taking opposite total categories sends a fibration over Sop to an opfibration over S. It sends a functor X \ensuremath{\to}Y to a functor Xop \ensuremath{\to}Y op in the same direction. Thus it defines a covariant isomorphism of the categories of structured objects; it does not reverse their morphisms. The same observation justifies the correctly stated symmetry of the perfect-fibration localization at p. 51.

The extra opposite cannot be repaired by an unspecified equivalence. Take S = e and W = W\ensuremath{\infty}. The category of homotopy types is not equivalent to its opposite. Its initial object \ensuremath{\varnothing}is strict: a morphism T \ensuremath{\to}\ensuremath{\varnothing}exists only when T \ensuremath{\simeq}\ensuremath{\varnothing}. Such an equivalence would force the initial object of the opposite category, namely the one-point homotopy type, to be strict there. This would say that every map \ensuremath{*}\ensuremath{\to}T is an isomorphism, contradicted by the inclusion of one point into a two-point set. The external opposite is therefore a substantive variance slip, even though both fine localizations themselves are correctly constructed.

The notation HOTo W (S) may still be used as a label for the covariant theory, provided it is defined by the left-fine localization. It should not be interpreted here as taking the opposite of HOTW (Sop). This convention also gives the stated direction to the functor j in the diagram which follows.

\end{dcomment}


Notons l'isomorphisme
\[
\uc{Parf}(S^o) \simeq \uc{Parf}\,S
\]
(provenant du fait que $f : X \longrightarrow Y$ dans $\Cat$ est une fibration parfaite si et seulement si $f^o$ l'est). Cet isomorphisme $X \longmapsto X^o$ est compatible avec les localiseurs (car $f \in W \Longleftrightarrow f^o \in W$), d'où
\[
\HOT^{\mathrm{lc}}_W(S^o) \simeq \HOT^{\mathrm{lc}}_W(S)
\]
(isomorphisme de catégories). On aura deux foncteurs canoniques entre les trois catégories de $W$-types d'homotopie relatifs sur $S$
\[
\begin{tikzcd}[column sep=small]
(\uc{Fibsc}_0S)W_S^{-1} \arrow[rr, "\approx"] \arrow[dr, dotted, "i'"'] & & \HOT^{\mathrm{lc}}_W(S) \arrow[dl, "i"] \arrow[dr, "j"'] & & (\uc{Cofibsc}_0S)W_S^{-1} \arrow[ll, "\approx"'] \arrow[dl, dotted, "j'"] \\
& \HOT_W(S) & & \HOT^o_W(S) \defeq \HOT_W(S^o). &
\end{tikzcd}
\]
\printedpage{32}
% head: A. Grothendieck : Dérivateurs, Ch. VI, Édition des Universités de Montpellier II et Paris VII
% foot: Version 9 novembre 2001
\texttt{[plutôt} $\defeq \HOT_W(S^o)^o$\texttt{]} Ces flèches $i$, $j$ sont suggerées par les flèches en pointillés $i'$, $j'$, provenant des inclusions % [sic] suggerées -> suggérées, missing accent
\[
\uc{Fibsc}_0S \lhook\joinrel\longrightarrow \uc{Fibsc}\,S,\quad \uc{Cofibsc}_0S \lhook\joinrel\longrightarrow \uc{Cofibsc}\,S,
\]
qui définiraient $i$, $j$ à l'aide du choix de foncteurs quasi-inverses pour les flèches horizontales. Mais $i$, $j$ ne seraient définis qu'à isomorphisme (unique) près. Il sera plus joli de choisir les quasi-inverses de façon canonique, à l'aide des flèches
\[
\uc{Fibsc}_0S \xleftarrow{\Phi_0} \uc{Parf}\,S \xrightarrow{\Psi_0} \uc{Cofibsc}_0S
\]
(où $\Phi_0 = \Phi|\uc{Parf}\,S$, $\Psi_0 = \Psi|\uc{Parf}\,S$) du diagramme p. 8. La question principale
\msp{page 52}
qui reste non élucidée, c'est si ces foncteurs $i$, $j$ sont bien \textit{pleinement fidèles,} comme je m'y attends certes.

Je veux maintenant prouver les affirmations précédentes sur les localisés.

\textcircled{1} Le localisé pour la $W$-équivalence grossière. Il suffira, à un $\varepsilon$ près, d'établir les conditions du lemme p. 11, dans les trois cas suivants.
\begin{tikzcd}[baseline=(\tikzcdmatrixname-1-1.base)]
\uc{Parf}\,S \arrow[r, "r"] & \Cat/S \arrow[l, dotted, bend left, "C"]
\end{tikzcd}
(ça donnera le cas de $\uc{Fibhot}\,S$, $\uc{Liss}\,S$ et $\uc{Prop}\,S$ sur $\Cat/S$) et
\begin{tikzcd}[baseline=(\tikzcdmatrixname-1-1.base)]
\uc{Fibsc}_0S \arrow[r, "p"] & \Cat/S \arrow[l, dotted, bend left]
\end{tikzcd},
\begin{tikzcd}[baseline=(\tikzcdmatrixname-1-1.base)]
\uc{Cofibsc}_0S \arrow[r, "q"] & \Cat/S \arrow[l, dotted, bend left]
\end{tikzcd},
qui donneront respectivement les cas de
\[
\uc{Fib}_0S,\ \uc{Fib}\,S,\ \uc{Liss}\,S \qquad \text{et} \qquad \uc{Cofib}_0S,\ \uc{Cofib}\,S,\ \uc{Prop}\,S \qquad \text{sur } S.
\]
Mais l'ennui pour $p$ et $q$, c'est que les foncteurs en sens inverse que j'avais construit (et noté encore $\Phi_0$, $\Psi_0$ (mais c'étaient essentiellement $\Phi\circ r\circ C$ et $\Psi\circ r\circ c$, à ne pas confondre avec les foncteurs $\Phi_0$, $\Psi_0$ en bas de la page précédente, et marqués aussi dans le diagramme p. 8), notons les plutôt $\Phi'$, $\Psi'$, marchent bien du côté $\Cat/S$, i.e. $p\Phi'$, $q\Psi'$ reçoivent le foncteur $\mathrm{id}_{\Cat/S}$ par une $W$-flèche, mais pas du côté des catégories but - on n'arrive pas à trouver une flèche naturelle entre $\Phi'p$ resp. $\Psi'q$, et les foncteurs identiques de $\uc{Fibsc}_0S$ resp. $\uc{Cofibsc}_0S$. Je ne vois
\msp{page 53}
pas, en effet, qu'il y en ait - la difficulté, par cette voie, parait irrémédiable. Par contre, si on regarde $\Phi'$, $\Psi'$ comme foncteurs à valeurs dans $\uc{Fib}_0S$, $\uc{Cofib}_0S$ respectivement (donc dans les catégories où on n'exige pas que les flèches soient catésiennes resp. cocartésiennes), on trouve encore des $W$-morphismes % [sic] catésiennes -> cartésiennes, missing r
\[
\mathrm{id} \longrightarrow \Phi'p, \qquad \mathrm{id} \longrightarrow \psi'q, % [sic] \psi'q -> \Psi'q, lowercase psi where \Psi' is used everywhere else
\]
ce qui liquide au moins déjà le cas de toutes les catégories des diagrammes p. 8 qui sont des
\printedpage{33}
% head: A. Grothendieck : Dérivateurs, Ch. VI, Édition des Universités de Montpellier II et Paris VII
% foot: Version 9 novembre 2001
sous-catégories pleines de $\Cat/S$ (\footnote[36]{Mais c'est idiot de prendre des foncteurs si compliqués, parce qu'on veut tout ramener à $\Cat/S$. Il est de plus simple de comparer $\uc{Fibhot}\,S$ directement à $\uc{Fibsc}_0S$ et à $\uc{Cofibsc}_0S$. $\Phi_0$ et $\Psi_0$ (qui sont en fait définis sur $\uc{Fibhot}$, pas seulement sur $\uc{Parf}$), et on trouve comme ci-dessous $\mathrm{id} \xrightarrow{W} \pi_0\Phi_0$ et $\Phi_0\pi_0 \xrightarrow{W} \mathrm{id}$ (itou pour $\Psi_0$ et $\uc{Cofibsc}_0S$), où $\pi_0 : \uc{Fibsc}_0S \longrightarrow \uc{Fibhot}\,S$.}). Il ne reste à traiter que le cas des quatre catégories de structures `scindées'
\[
\uc{Fibsc}\,S,\quad \uc{Fibsc}_0S \qquad \text{et} \qquad \uc{Cofibsc}\,S,\quad \uc{Cofibsc}_0S.
\]
Bien sûr il suffit de le faire d'un côté, disons du côté des fibrés, par dualité ($S \longmapsto S^o$) cela établira le résultat symétrique. On va donc regarder d'un peu plus près le carré de foncteurs
\[
\begin{tikzcd}[column sep=huge, row sep=large]
\uc{Fibsc}\,S \arrow[d, "p"] & \uc{Fibsc}_0S \arrow[l, hook, "\alpha'"'] \arrow[d, "p_0"'] \\
\uc{Fib}\,S \arrow[u, dotted, bend left, "\varphi"] & \uc{Fib}_0S \arrow[l, hook, "\alpha"'] \arrow[u, dotted, bend right, "\varphi_0"']
\end{tikzcd}
\]
(\footnote[37]{NB Il s'agit de prouver seulement que $p$, $p_0$ induisent des équivalences pour les localisés (pour la $W$-équivalence grossière). On sait déjà qu'il en est ainsi pour $\alpha$, et ça en résultera pour $\alpha'$.}, \footnote[38]{inutile de regarder ce carré, cf. annotations marginales ci-dessus (haut de la page)}). Le foncteur
\[
\Phi : \Cat/S \longrightarrow \uc{Fibsc}\,S
\]
induit des foncteurs en sens inverse de $p$, $p_0$ sur les sous-catégories pleines $\uc{Fib}\,S$, $\uc{Fib}_0S$ de $\Cat/S$, foncteurs marqués en pointillés comme $\varphi$, $\varphi_0$.

\msp{page 54}

\dgrblock{dgrb188}{Why the free split fibration is the right comparison}
\begin{dcomment}{context}{p53}{Why the free split fibration is the right comparison}
The difficulty with the first attempt is a difficulty about morphisms, not about the existence of enough good objects. A path completion puts every object into a particularly well-behaved class, but its unit need not preserve a chosen splitting. The comma construction \ensuremath{\Phi} has a stronger universal property: it freely supplies the selected Cartesian lifts. That property gives the reverse comparison map with exactly the morphism convention required by Fibsc(S).

This construction was already part of the theory of fibred categories before the present manuscript. Gray’s treatment constructs the free split fibration, and the modern formulation identifies split fibrations and splitting-preserving functors with strict algebras for the comma-category 2-monad on Cat/S (Gray, 1966; Emmenegger et al., 2024). In the latter reference, Theorem 2.6 separates strict algebras from pseudo-algebras; Theorem 5.4 treats the two adjoints used to split an already fibred category. These are distinct operations. The present \ensuremath{\Phi} is the free construction on an arbitrary category over S.

Its importance for localization is elementary. Let U : Fibsc(S) \ensuremath{\to}Cat/S forget the splitting. There are transformations

\ensuremath{\eta} : id \ensuremath{-}\ensuremath{\to}U\ensuremath{\Phi},

\ensuremath{\gamma} : \ensuremath{\Phi}U \ensuremath{-}\ensuremath{\to}id. The first lives in Cat/S; the second lives in the category with splitting-preserving morphisms. Both become isomorphisms after the relevant localization. Consequently U and \ensuremath{\Phi} induce inverse equivalences. One does not need a comparison transformation in a category with an incorrectly enlarged or restricted class of arrows.

For coarse equivalences, this conclusion requires only that W contain categorical homotopy equivalences and satisfy two-out-of-three. The ordinary functor forgetting the initial base object gives U\ensuremath{\Phi}(X) \ensuremath{\to}X, with a section \ensuremath{\eta}X and a natural contraction. Although the forgetful map is usually not over S, the contraction proves that \ensuremath{\eta}X is an equivalence of total homotopy types. The counit is then a coarse equivalence because its composite with the unit on UY is the identity. The more informative fibrewise verification comes on pp. 58--60.

\end{dcomment}


Il faut seulement vérifier que si $X/S$ est dans $\uc{Fib}_0S$, alors $\varphi(X) = \Phi(X)$ est dans la sous-catégorie pleine $\uc{Fibsc}_0S$ de $\uc{Fibsc}\,S$. Mais on a vu que c'est le cas, pour peu seulement que $X/S$ soit dans $\uc{Fibhot}\,S$, car les morphismes de changement de base dans $\Phi(X)$, associés aux flèches $u : s \longrightarrow t$ dans $S$
\[
u^*_{\Phi_X} : t\backslash X \longrightarrow s\backslash X
\]
sont des $W_\infty$-équivalences si $X$ est Hot-fibré sur $S$. D'autre part, pour $X$ quelconque dans $\Cat/S$, on a une flèche, fonctorielle en $X$,
\[
X \longrightarrow \Phi(X),
\]
i.e. $\mathrm{id}_{\Cat/S} \longrightarrow \pi\Phi$ où $\pi : \uc{Fibsc}\,S \longrightarrow \Cat/S$ est le foncteur canonique, qui est une $W_\infty$-équivalence (même un homotopisme), donc une $W$-équivalence. A fortiori les foncteurs $p\varphi$, $p_0\varphi_0$ donnent lieu ($\uc{Fib}\,S$ et $\uc{Fib}_0S$ étant des sous-catégories pleines de $\Cat/S$) à des flèches
\[
\mathrm{id} \longrightarrow p\varphi, \qquad \mathrm{id} \longrightarrow p_0\varphi_0.
\]

% ---- from Der-ChVI-345_p036-040.tex
\printedpage{34}
% head: A. Grothendieck : Dérivateurs, Ch. VI, Édition des Universités de Montpellier II et Paris VII

Il reste à trouver une flèche (et $W$-équivalence grossière)
\[\varphi p \longrightarrow \mathrm{id},\]
\msp{page 55}
laquelle induira automatiquement
\[\varphi_0 p_0 \longrightarrow \mathrm{id}\]
puisque $\underline{\mathrm{Fibsc}}_0 S$ est \texttt{[une]} sous-catégorie pleine de $\ul{Fibsc}\,S$. En fait, on va définir $\varphi p \longrightarrow \mathrm{id}$, donc
\[\boxed{\gamma_X : \Phi(X) \longrightarrow X}\]
quand $X$ est fibrée scindée sur $X$ - une flèche \textit{cartésienne,} donc (si on veut) compatible avec les scindages de $\Phi(X)$, et de $X$. En fait, on aura une telle flèche, fonctorielle en $X$, pour $X$ variable dans $\ul{Fib}\,S$ (donc sans se borner à des $S$-morphismes cartésiens entre catégories fibrées sur $S$), mais la flèche $\longrightarrow$ étant bien sûr \textit{cartésienne.} (Ça a un sens, sans choisir de scindage sur $X$, et sans qu'il en existe forcément.) La définition de cette flèche $\gamma$ fait appel de façon essentielle à l'hypothèse que $X$ soit Cat-fibrant sur $S$. Bien sûr, composé à la flèche d'inclusion évidente % [sic] fibrée scindée sur X -> fibrée scindée sur S, the base category is S
\[X \xrightarrow{i_X} \Phi(X)\]
on va avoir
\[\gamma_X i_X = \mathrm{id}_X\]
(\footnote[39]{NB Comme on sait que $i_X \in W$, il en résultera bien $\gamma_X \in W$ comme on le veut.}). On définit $\gamma$ fibre par fibre,
\[\gamma_{X,s} : \Phi(X)_s \overset{\text{déf}}{=} s\backslash X \longrightarrow X_s,\]
donc par une \textit{rétraction} de $s\backslash X$ sur la sous-catégorie pleine $X_s$. Un objet de $s\backslash X$ est un couple $(x,u)$, $x \in \mathrm{Ob}\, X$, $u : s \longrightarrow p_X(x)$ ($p_X : X \longrightarrow S$ le morphisme structural de $X$). On pose
\[\gamma(x,u) = u_X^*(x) \in \mathrm{Ob}\, X_s,\]
\msp{page 56}
ce qui suppose donc qu'on a choisi des foncteurs image inverse $u_X^* : X_t \longrightarrow X_s$, pour la catégorie Cat-fibrée $X/S$. (Mais à isomorphisme unique près, $\gamma$ ne dépend pas de ces choix.) Je me dispense de la vérification fastidieuse que cette loi est bien ‘fonctorielle en $s$’, i.e. définit bien une flèche cartésienne
\[\gamma_X : \Phi(X) \longrightarrow X,\]
et de vérifier que celle-ci est fonctorielle en $X$.

\printedpage{35}
% head: A. Grothendieck : Dérivateurs, Ch. VI, Édition des Universités de Montpellier II et Paris VII

Peut-être est-il plus élégant de montrer que $\pi : \ul{Fibsc}\,S \longrightarrow \mathrm{Cat}/S$ et $\Phi : \mathrm{Cat}/S \longrightarrow \ul{Fibsc}\,S$ sont adjoints l'un de l'autre :
\[\boxed{\mathrm{Homcart}_S(\Phi(\underbrace{\;X\;}_{\in\mathrm{Cat}/S}),\underbrace{\;Y\;}_{\in\ul{Fibsc}}) \simeq \mathrm{Hom}_S(X,\pi Y),}\]
d'où résulteront des flèches d'adjonction
\[\Phi\pi \longrightarrow \mathrm{id}_{\ul{Fibsc}}, \qquad\qquad \pi\Phi \longleftarrow \mathrm{id}_{\ul{Fib}},\]
\texttt{[plutôt} $\pi\Phi \longleftarrow \mathrm{id}_{\mathrm{Cat}/S}$\texttt{]} et on aura gagné. Mais pour prouver cette formule d'adjonction, j'ai l'impression qu'il faut utiliser $\varphi p \longrightarrow \mathrm{id}$ \ldots\ Il y a des vérifications fastidieuses \ldots

Finalement, je vois maintenant qu'on a cinq cas à traiter, et trois si on joue sur les symétries :

\msp{page 57}

\dgrblock{dgrb189}{The free split adjunction, with its arrows specified}
\begin{dcomment}{exposition}{p56}{The free split adjunction, with its arrows specified}
Grothendieck’s proposed adjunction is the efficient way to finish the comparison, once “Cartesian” in the Hom-set on the left is read as preserving the selected Cartesian lifts. We give its construction because the same calculation resolves the naturality issue in the preceding paragraph.

For p : X \ensuremath{\to}S, an object of \ensuremath{\Phi}X is (s, x, u) with u : s \ensuremath{\to}p(x). A morphism

(s, x, u) \ensuremath{-}\ensuremath{\to}(t, y, v)

is a pair (a, f) with a : s \ensuremath{\to}t, f : x \ensuremath{\to}y, and p(f)u = va. The chosen lift along a with target (t, y, v) is

\ensuremath{-}\ensuremath{-}\ensuremath{-}\ensuremath{-}\ensuremath{\to}(t, y, v).

(s, y, va) Identity and composition laws for these chosen lifts hold strictly. The unit sends x to (p(x), x, id). Let q : Y \ensuremath{\to}S be split, and let h : X \ensuremath{\to}UY be any functor over S. Write cu,z : u\ensuremath{*}z \ensuremath{\to}z for the chosen lift. Define

bh(s, x, u) = u\ensuremath{*}h(x).

For (a, f) as above, define bh(a, f) as the unique arrow over a such that

cv,h(y) bh(a, f) = h(f)cu,h(x).

The Cartesian property of cv,h(y) gives existence and uniqueness. These same uniqueness statements prove the identity and composition laws: both possible images of a composite have the required base arrow and the same composite with the terminal chosen lift. The strict splitting laws show that a selected lift (a, id) is taken to the selected lift in Y . Thus bh is splitting-preserving.

Conversely, every splitting-preserving H : \ensuremath{\Phi}X \ensuremath{\to}Y is forced by h = H\ensuremath{\eta}X. The chosen arrow (s, x, u) \ensuremath{\to}(p(x), x, id) forces its object value to be u\ensuremath{*}h(x), and Cartesian uniqueness forces its value on every morphism. The two constructions are inverse and natural, proving the actual adjunction

Homsplit/S(\ensuremath{\Phi}X, Y ) \ensuremath{\sim}= HomS(X, UY ).

Its counit is \ensuremath{\gamma}Y (s, y, u) = u\ensuremath{*}y, and \ensuremath{\gamma}Y \ensuremath{\eta}UY = idUY . Its naturality is strict for splitting-preserving functors, which are precisely the morphisms needed on this side of the comparison. This is the free-algebra adjunction underlying the 2-monad description cited above.

For a merely cloven fibration the same formula still gives a Cartesian functor \ensuremath{\Phi}X \ensuremath{\to}X, using the coherence isomorphisms of the cleavage. Different choices yield canonically isomorphic functors compatible with their Cartesian comparison arrows. But an arbitrary functor over S need not preserve Cartesian arrows, so the formula is not strictly natural on the full category Fib(S). Nor can “Cartesian” alone imply preservation of the selected splitting. Those stronger assertions in the source are unnecessary; the strict adjunction just proved gives the comparison in Fibsc(S), and the next block explains what survives after localization for all functors over S.

\end{dcomment}


\[
\begin{tikzcd}[row sep=large]
\ul{Parf}\,S \arrow[d, "r"'] \\
\mathrm{Cat}/S \arrow[u, dotted, bend right, "C"']
\end{tikzcd}
\qquad\qquad\qquad
\left\{\begin{array}{lcl}
rC & \longleftarrow & \mathrm{id} \\
Cr & \longleftarrow & \mathrm{id}
\end{array}\right.
\]

\[
\begin{tikzcd}
\ul{Fibsc}\,S \arrow[dr, "\lambda"] & & \ul{Cofibsc}\,S \arrow[dl, "\mu"'] \\
& \mathrm{Cat}/S \arrow[ul, dotted, bend left, "\Phi"] \arrow[ur, dotted, bend right, "\Psi"'] &
\end{tikzcd}
\qquad
\begin{array}{l}
\left\{\begin{array}{lcl}
\lambda\Phi & \longleftarrow & \mathrm{id} \\
\Phi\lambda & \longrightarrow & \mathrm{id}
\end{array}\right. \\[2ex]
\left\{\begin{array}{lcl}
\mu\Psi & \longleftarrow & \mathrm{id} \\
\Psi\mu & \longrightarrow & \mathrm{id}
\end{array}\right.
\end{array}
\]

\[
\begin{tikzcd}
\underline{\mathrm{Fibsc}}_0 S \arrow[dr, "\lambda_0"] & & \underline{\mathrm{Cofibsc}}_0 S \arrow[dl, "\mu_0"'] \\
& \ul{Fibhot}\,S \arrow[ul, dotted, bend left, "\Phi_0"] \arrow[ur, dotted, bend right, "\Psi_0"'] &
\end{tikzcd}
\qquad
\begin{array}{l}
\left\{\begin{array}{lcl}
\lambda_0\Phi_0 & \longleftarrow & \mathrm{id} \\
\Phi_0\lambda_0 & \longrightarrow & \mathrm{id}
\end{array}\right. \\[2ex]
\left\{\begin{array}{lcl}
\mu_0\Psi_0 & \longleftarrow & \mathrm{id} \\
\Psi_0\mu_0 & \longrightarrow & \mathrm{id}
\end{array}\right.
\end{array}
\]

(\footnote[40]{\textbf{NB} Tous les foncteurs marqués sont compatibles avec la $W$-équivalence grossière, et des dix homomorphismes fonctoriels sont des $W$-équivalences grossières.})

NB Les quatre derniers homomorphismes de foncteurs, relatifs au troisième diagramme (avec les indices 0) sont simplement induits par ceux relatifs au diagramme qui le précède. Donc les vérifications-clef sont relatives aux deux premiers diagrammes, et par symétrie, on n'a finalement qu'à regarder les \textit{deux} couples de foncteurs $(r,C)$, $(\lambda,\Phi)$, et pour chacun de ces couples, exhiber les deux homomorphismes fonctoriels, des deux composés avec les foncteurs identiques.

\printedpage{36}
% head: A. Grothendieck : Dérivateurs, Ch. VI, Édition des Universités de Montpellier II et Paris VII

\msp{page 58}

\dgrblock{dgrb190}{Naturality recovered after localization}
\begin{dcomment}{development}{p57}{Naturality recovered after localization}
Grothendieck wants the comparison to accommodate all functors over the base, since the categories Fib(S) and Liss(S) are full subcategories of Cat/S. The raw pullback construction does not have that naturality. The localization nevertheless supplies exactly the required repair. We spell out the following elementary elaboration of his argument; the “ours” label concerns this formulation and proof, not priority. We checked the free-fibration and strictification formulations in Emmenegger et al. (2024, Sections 2 and 5). They distinguish strict and pseudo preservation of cleavages but do not state the following repair for arbitrary functors over the base. We make no broader claim about its absence from the literature.

Localized naturality of pullback retractions. \ours Suppose W\ensuremath{\infty}\ensuremath{\subseteq}W. Choose normalized cleavages on fibrations X \ensuremath{\to}S. Then the functors \ensuremath{\gamma}X : \ensuremath{\Phi}X \ensuremath{\to}X determine a natural isomorphism between the functors induced by U\ensuremath{\Phi} and the identity on the right-fine localization of the full category Fib(S). The same conclusion holds for coarse localization.

Proof. First, any natural transformation over S between functors F, G : A \ensuremath{\to}B makes their images equal in (Cat/S)[(W d S)\ensuremath{-}1]. It is represented by a functor H : A \ensuremath{\times} [1] \ensuremath{\to}B over S. The projection r : A \ensuremath{\times} [1] \ensuremath{\to}A is right-fine: for every s its comma functor is the projection

(s\textbackslash{}A) \ensuremath{\times} [1] \ensuremath{-}\ensuremath{\to}s\textbackslash{}A,

a categorical homotopy equivalence. Its two sections become the same inverse of r, hence F = Hj0 and G = Hj1 become equal. The identical proof applies to coarse localization. If A is fibred, A \ensuremath{\times} [1] is fibred, so this argument also takes place within the indicated full subcategory.

Now let f : X \ensuremath{\to}Y be any functor over S, without a Cartesian hypothesis. For (s, x, u) define a vertical arrow

\ensuremath{\theta}f,(s,x,u) : f(u\ensuremath{*}

Xx) \ensuremath{-}\ensuremath{\to}u\ensuremath{*} Y f(x) by the equation

cY

u,f(x)\ensuremath{\theta}f,(s,x,u) = f(cX u,x). Existence and uniqueness follow from the Cartesian property of the lift in Y . This defines a natural transformation f\ensuremath{\gamma}X \ensuremath{\Rightarrow}\ensuremath{\gamma}Y \ensuremath{\Phi}(f) over S. To check naturality, compose the two candidate arrows with the chosen lift to the terminal object of a morphism (a, g) in \ensuremath{\Phi}X. Both composites are f(g)f(cX u,x), u,x) = f(gcX and both have the same base arrow; Cartesian uniqueness identifies them. The first paragraph of the proof therefore gives naturality after localization.

The comparison maps are right-fine equivalences. On a fibre, \ensuremath{\gamma}X,s is right adjoint to Xs ,\ensuremath{\to}s\textbackslash{}X: an arrow y \ensuremath{\to}x over u : s \ensuremath{\to}p(x) factors uniquely through u\ensuremath{*} Xx. Every adjunction gives a homotopy equivalence of nerves. Since source and target of \ensuremath{\gamma}X are fibrations, right-fine equivalence is equivalent to fibrewise equivalence. Thus the localized natural transformation is an isomorphism. Its coarse version follows either from the same argument or from \ensuremath{\gamma}X\ensuremath{\eta}X = id and coarse invertibility of \ensuremath{\eta}X. \ensuremath{\Box}

The defect before localization is real. Take S = [1], X = S, Y = S \ensuremath{\times} [1], and f(s) = (s, s). At the object (0, 1, 0 \ensuremath{\to}1) of \ensuremath{\Phi}X, the two functors in the comparison square have values (0, 0) and (0, 1). These objects are not isomorphic. The comparison \ensuremath{\theta} is the noninvertible arrow between them. Thus the source’s strict naturality assertion cannot simply be replaced by a claim of naturality up to isomorphism; the localization, or the restriction to splitting-preserving morphisms, performs an essential part of the repair.

\end{dcomment}


\cercle{2} Les localisés pour la \textit{$W$-équivalence fine droite.} Il n'y a plus lieu, on l'a dit, de regarder les catégories internes du diagramme p. 8, où ladite équivalence coïncide avec la $W$-équivalence grossière, pour laquelle on connait déjà les localisés. Il reste à regarder les catégories
\[\mathrm{Cat}/S,\ \ \ul{Liss}\,S,\ \ \ul{Fib}\,S,\ \ \ul{Fibsc}\,S.\]
Le lemme de la page 11 nous dit qu'il suffit de montrer que dans la situation
\[\ul{Fibsc}\,S \xrightarrow{\lambda} \mathrm{Cat}/S,\]
il y a un foncteur en sens inverse, avec les propriétés specifiées. Il en résultera que toute sous-catégorie pleine $\mathcal{F}$ de $\mathrm{Cat}/S$ qui contient $\ul{Fib}\,S$, quand on la munit du localiseur de la $W$-équivalence fine droite, donne un localisé équivalent, par le foncteur induit par l'inclusion, au localisé correspondant de $\mathrm{Cat}/S$, i.e. essentiellement à $\mathrm{HOT}_W(S)$. % [sic] specifiées -> spécifiées, missing accent

Comme foncteur en sens inverse, il s'impose de prendre encore $\Phi$. Il y a à vérifier trois choses :

a) $\Phi$ \texttt{[est]} compatible avec les localiseurs, i.e. si
\begin{tikzcd}[baseline=(\tikzcdmatrixname-2-2.base), row sep=small, column sep=small]
X \arrow[rr, "f"] \arrow[dr] & & Y \arrow[dl] \\
& S &
\end{tikzcd}
sur $S$ est une $W$-équivalence fine droite, i.e. les $s\backslash X \longrightarrow s\backslash Y$ sont des $W$-équivalences, alors $\Phi(f)$ aussi. Mais ladite équivalence coïncide
\msp{page 59}
dans $\ul{Fibsc}\,S$ (comme aussi dans $\ul{Fib}\,S$ et $\ul{Liss}\,S$) avec la $W$-équivalence par fibres, et d'autre part $\Phi(f)_s$ n'est autre que la flèche précédente $s\backslash X \longrightarrow s\backslash Y$. Donc il est évident que
\[f \in \underbrace{W_S^{\mathrm{d}}}_{W\text{-équivalence droite}} \iff \Phi(f) \in \underbrace{W_S^{\mathrm{f}}}_{W\text{-équivalence par fibres}} .\]

b) Que $X \xrightarrow{i_X} \lambda\Phi(X)$ est dans $W_S^{\mathrm{d}}$, pour tout $X$ dans $\mathrm{Cat}/S$. Cela signifie donc que $i_X$ induit
\[s\backslash X \longrightarrow s\backslash\Phi(X)\]
qui est une $W$-équivalence. Ce serait O.K. s'il était clair que $X \longmapsto \Phi(X)$ commute à tout changement de base propre sur $S$ (donc en particulier à $s\backslash S \longrightarrow S$), puisqu'on sait que $X \longrightarrow \Phi(X)$ est une $W$-équivalence. C'est sans doute trop demander, mais du moins que ça commute au changement de base qui est un co-isomorphisme local (telles les $s\backslash S \longrightarrow S$).

Pour un tel changement de base $S' \xrightarrow{p} S$, vérifier que la fibre de $\Phi(X')$ en $s'$ est isomorphe à celle de $\Phi(X)$ en $s = p(s')$, je m'en contenterai. Or cette fibre est $s'\backslash X'$, déduit de $X'$ sur $S'$ par le changement de base $s'\backslash S' \longrightarrow S'$, donc déduit de $X$ par $s'\backslash S' \longrightarrow S$. Or $s'\backslash S' \longrightarrow s\backslash S$ est un isomorphisme par hypothèse
\printedpage{37}
% head: A. Grothendieck : Dérivateurs, Ch. VI, Édition des Universités de Montpellier II et Paris VII
\msp{page 60}
sur $p$ (co-isomorphisme local), donc le changement de base $s'\backslash S' \longrightarrow S$ s'identifie à $s\backslash S \longrightarrow S$, et l'image inverse est bien isomorphe à $s\backslash X$, OK.

\dgrblock{dgrb191}{The comma resolution is right-fine}
\begin{dcomment}{exposition}{p59}{The comma resolution is right-fine}
Grothendieck has reduced the fine comparison to checking that the unit and counit are fine equivalences. This is the decisive point at which the comma categories preserve the relative diagram, rather than merely its total homotopy type. A direct contraction proves the assertion for every X \ensuremath{\to}S and makes the proposed base-change verification unnecessary.

Fix s \ensuremath{\in}S. An object of s\textbackslash{}\ensuremath{\Phi}X consists of

\ensuremath{-}\ensuremath{\to}vt

\ensuremath{-}\ensuremath{\to}up(x). s Define R : s\textbackslash{}\ensuremath{\Phi}X \ensuremath{\to}s\textbackslash{}X by composition, R(v, u, x) = (x, uv). A morphism represented by (a, f) satisfies av = v\ensuremath{\prime} and p(f)u = u\ensuremath{\prime}a; it is sent to f, which satisfies p(f)uv = u\ensuremath{\prime}v\ensuremath{\prime}. Thus R is a functor. The functor J = s\textbackslash{}\ensuremath{\eta}X sends (x, w) to the factorization

\ensuremath{-}\ensuremath{\to}p(x)

\ensuremath{-}\ensuremath{\to}p(x). s We have RJ = id. The arrow (u, idx) gives a natural transformation

ids\textbackslash{}\ensuremath{\Phi}X \ensuremath{-}\ensuremath{\to}JR.

It is over s, since uv is the composite arrow of the target factorization. These identities prove that J and R are inverse up to categorical homotopy. Hence \ensuremath{\eta}X belongs to W d S whenever W contains W\ensuremath{\infty}.

For a split fibration Y , the fibre of the counit is

(y, u) 7\ensuremath{-}\ensuremath{\to}u\ensuremath{*}y.

\ensuremath{\gamma}Y,s : s\textbackslash{}Y \ensuremath{-}\ensuremath{\to}Ys, If K : Ys ,\ensuremath{\to}s\textbackslash{}Y is the inclusion, Cartesian factorization gives

Homs\textbackslash{}Y (Kz, (y, u)) \ensuremath{\sim}= HomYs(z, u\ensuremath{*}y).

Thus K \ensuremath{\dashv}\ensuremath{\gamma}Y,s, with identity unit for a normalized splitting. In particular both functors are W\ensuremath{\infty}equivalences. The counit is therefore fibrewise, and hence right-fine, an equivalence. Finally, \ensuremath{\Phi}(f)s is literally s\textbackslash{}f; this verifies compatibility of \ensuremath{\Phi} with the two localizers without an additional argument. All parts of the comparison now fit together. The unit and counit induce inverse natural isomorphisms between the localized composites. We obtain

S)\ensuremath{-}1] \ensuremath{\simeq}Fibsc(S)[W \ensuremath{-}1

(Cat/S)[(W d obj]. If F is any full subcategory of Cat/S containing the fibrations, the resolution U\ensuremath{\Phi} lands in F and its unit is a morphism of F on restriction. The same argument proves the asserted equivalence with F[(W d S)\ensuremath{-}1]. Reversing the comma directions gives the corresponding statement for opfibrations and W g S.

\end{dcomment}


c) Que de même
\[\Phi\lambda \longrightarrow \mathrm{id}\]
est non seulement dans $W_S$ (grossier), mais que c'est une $W$-équivalence par fibres. Donc on part de $X$ fibrée sur $S$. On regarde
\[\Phi(X) \longrightarrow X,\]
pour voir que c'est une $W$-équivalence par fibres, il suffit de voir qu'il en est ainsi de $X \longrightarrow \Phi(X)$,
\[X_s \longrightarrow s\backslash X\]
est une $W$-équivalence pour tout $s \in S$. Mais ceci est déjà vrai dès que $X$ est lisse sur $S$.

Comme le cas de la $W$-équivalence fine gauche se ramène au cas actuel par $S \longmapsto S^{o}$, on a donc prouvé tout ce qu'était attendu, à la seule exception de la pleine fidèlité des foncteurs $i$, $j$ de la page 51. % [sic] ce qu'était -> ce qui était, grammar; fidèlité -> fidélité, accent

Mais dans le résumé prospectif de la situation pour les localiseurs, p. 49-51, j'ai oublié d'autres foncteurs que $i$, $j$, de nature si possible encore plus triviale. Ce sont
\msp{page 61}
ceux qu' on peut décrire à partir du foncteur identique, sur une des catégories du diagramme où aient cours deux, voire trois, notions de $W$-équivalence différentes. Ce sont, je le rappelle

\dgrblock{dgrb192}{The hypotheses of the fine comparison}
\begin{dcomment}{exposition}{p60}{The hypotheses of the fine comparison}
The comparison just established has survived the earlier changes of axioms. Its proof uses the inclusion W\ensuremath{\infty}\ensuremath{\subseteq}W, the elementary saturation properties, and the smooth-fibre criterion already checked on p. 42. In particular it does not use the new L5. On split fibrations, right-fine equivalences are exactly objectwise W-equivalences of the associated strict diagrams; the analogous statement holds on split opfibrations. This is the part of the construction which remains useful for homological and truncated equivalences.

The pure coarse comparison with the internal categories also needs no new L5. The path replacement C produces homotopy fibrations, its unit X \ensuremath{\to}C(X) is a total homotopy equivalence, and naturality together with two-out-of-three makes C preserve coarse equivalences. L5 enters in the additional identification of coarse equivalences between these replacements with fibrewise, hence fine, equivalences. This is what makes C induce the functor back into the fine category used for the reflection. The later annotation correctly says that this requirement is highly restrictive. The commentary on pp. 43--45 explains why it must not be silently transferred to arbitrary homological localizers.

For W = W\ensuremath{\infty}, the modern literature gives a precise model for the coarse part: small categories, nerve weak equivalences, and smooth proper functors form a category of fibrant objects; perfect objects over a base suffice to compute the corresponding localization. See Cisinski (2009, Proposition 4.4 and Corollary 4.5). This places Grothendieck’s path replacement inside a developed homotopical framework. The explicit comma proof above is still useful, since it shows exactly which comparison remains valid when the coefficient equivalences are enlarged.

\end{dcomment}


\begin{enumerate}[label={}]
\item[$1^{o}$)] Les quatre catégories
\[(*) \qquad\qquad \ul{Fibsc}\,S \longrightarrow \ul{Fib}\,S \longrightarrow \ul{Liss}\,S \longrightarrow \mathrm{Cat}/S, \qquad\qquad W_S^{\mathrm{d}} \subseteq W\]
avec la $W$-équivalence fine droite d'une part, grossière de l'autre, et dualement
\[(*') \qquad\qquad \ul{Cofibsc}\,S \longrightarrow \ul{Cofib}\,S \longrightarrow \ul{Prop}\,S \longrightarrow \mathrm{Cat}/S, \qquad\qquad W_S^{\mathrm{g}} \subseteq W\]
pour la $W$-équivalence fine gauche d'une part, la grossière de l'autre, et de plus
\item[$2^{o}$)] $\mathrm{Cat}/S$, qui était inclue dans les deux suites précédentes, et qui est la seule catégorie du diagramme donnant lieu à trois notions de $W$-équivalence différentes % [sic] inclue -> incluse, past participle of inclure
\end{enumerate}

\[
\begin{tikzcd}[row sep=small, column sep=large]
W_S^{\mathrm{d}} \arrow[dr, hook] & \\
& W_S \\
W_S^{\mathrm{g}} \arrow[ur, hook] &
\end{tikzcd}
\qquad\qquad .
\]

\printedpage{38}
% head: A. Grothendieck : Dérivateurs, Ch. VI, Édition des Universités de Montpellier II et Paris VII

Dans chacun des quatre cas ($C$, munie de deux localiseurs $\Sigma \subseteq \Sigma'$) de la suite $(*)$, le foncteur
\[C\Sigma^{-1} \longrightarrow C\Sigma'^{-1}\]
est essentiellement ‘le même’, à cause des théorèmes de comparaison. Il suffit de regarder ce foncteur dans un seul das, p. ex. un des deux extrêmes $\ul{Fibsc}\,S$ ou $\mathrm{Cat}/S$. Même remarque pour la suite $(*')$. Donc finalement il reste deux foncteurs remarquables % [sic] das -> cas, misprint
\msp{page 62}
à examiner, qui n'ont pas l'air du tout de provenir des foncteurs du diagramme p. 8 ou de quelque composé (ça ne donne que des équivalences canoniques ou bien les foncteurs $i$, $j$ de page 51, quand on va des catégories internes au diagramme vers les catégories bordantes). Le plus joli est sans doute de les regarder à partir de $\mathrm{Cat}/S$, puisque c'est là seulement qu'on les a tous les deux à la fois
\[
\begin{tikzcd}[row sep=small, column sep=tiny]
(\mathrm{Cat}/S)(W_S^{\mathrm{d}})^{-1} \approx \mathrm{HOT}_W(S) \arrow[dr, "\rho"] & \\
& (\mathrm{Cat}/S)W_S^{-1} \approx \mathrm{HOT}_W^{\mathrm{lc}}(S). \\
(\mathrm{Cat}/S)(W_S^{\mathrm{g}})^{-1} \approx \mathrm{HOT}_W^{o}(S) \arrow[ur, "\sigma"'] &
\end{tikzcd}
\]
Par contre, les foncteurs (qu'on a envie d'appeler ‘foncteurs d'inclusion’) $i$, $j$ de la page 51, en sens inverse de $\rho$, $\sigma$, sont ceux qui proviennent, par localisation $W_S^{\mathrm{d}}$ ou $W_S^{\mathrm{g}}$ selon le cas, d'un des foncteurs du diagramme page 8 (ou d'un composé de tels foncteurs), allant d'une des six catégories internes (où toutes les $W$-équivalences envisagées coïncident) vers l'une des sept catégories bordantes, sauf quand cette dernière est $\mathrm{Cat}/S$, est déterminé déjà avec quelle $W$-équivalence, $W_S^{\mathrm{d}}$ ou $W_S^{\mathrm{g}}$, on va travailler \texttt{[structure de cette phrase?]}. La détermi-\msp{page 63}% hyphen kept as printed: the word is divided at the manuscript page marker
nation peut-être la plus jolie est par
\[(*) \qquad\qquad \ul{Parf}\,S \longrightarrow \mathrm{Cat}/S\]
qui donne le foncteur $i$ quand on travaille avec $W_S^{\mathrm{d}}$, le foncteur $j$ quand on travaille avec $W_S^{\mathrm{g}}$. On voit alors tout de suite que si on compose avec le foncteur pertinent, soit $\rho$ soit $\sigma$, défini par l'inclusion $W_S^{\mathrm{d}} \subseteq W_S$ resp. $W_S^{\mathrm{g}} \subseteq W_S$ (comme parties de $\mathrm{Fl}(\mathrm{Cat}/S)$), on trouve le foncteur
\[(**) \qquad\qquad (\ul{Parf}\,S)W_S^{-1} \xrightarrow{\approx} (\mathrm{Cat}/S)W_S^{-1}\]
qui est l'\textit{équivalence canonique} entre ces deux catégories. Ainsi le foncteur correspondant à $i$, $j$ page 51, exprimé en termes de $\mathrm{Cat}/S$ comme foncteur en sens inverse de $\rho$, $\sigma$ (p. 62), est-il le composé de l'équivalence quasi-inverse de $(**)$, laquelle s'explicite par le foncteur $C$, et du foncteur sur le localisé $W_S^{\mathrm{d}}$ ou $W_S^{\mathrm{g}}$ déduit de $(*)$, i.e.

% ---- from Der-ChVI-345_p041-045.tex
\printedpage{39}
% head: A. Grothendieck : Dérivateurs, Ch. VI, Édition des Universités de Montpellier II et Paris VII

\[
\begin{tikzcd}[column sep=large, row sep=large]
 & \HOT_W(S)\approx(\Cat/S)(W_S^{\mathrm{d}})^{-1} \\
\HOT^{\mathrm{lc}}_W(S)\approx(\Cat/S)W_S^{-1} \arrow[ur, "i"] \arrow[r, "C"] \arrow[dr, "j"'] & (\ul{Parf}\,S)W_S^{-1} \arrow[u, "\text{\footnotesize car $W_S=W_S^{\mathrm{d}}$ sur $\ul{Parf}$}"'] \arrow[d, "\text{\footnotesize car $W_S=W_S^{\mathrm{g}}$ sur $\ul{Parf}$}"] \\
 & \HOT^{o}_W(S)\approx(\Cat/S)(W_S^{\mathrm{g}})^{-1}.
\end{tikzcd}
\]

\msp{page 64}
On voit par là qu'on a des isomorphismes canoniques
\[
\rho i\simeq\id,\qquad\sigma j\simeq\id
\]
dans la catégorie des endofoncteurs de $(\Cat/S)W_S^{-1}$ ($\approx\HOT^{\mathrm{lc}}_W(S)$), isomorphisme qui s'explicite ici par le fait que $\rho i$, $\sigma j$ sont déduits du même foncteur $C$ par passage aux localisés $W_S^{-1}$ (donc en fait $\rho i=\sigma j$ !), et qu'on a une flèche fonctorielle en $X$
\[
X\xrightarrow{i_X}C(X),
\]
qui est dans $W$ pour tout $X$. Ainsi, $\rho$ et $\sigma$ sont à isomorphisme près des \textit{rétractions} des catégories $\HOT_W(S)$, $\HOT^{o}_W(S)$ sur la catégorie $\HOT^{\mathrm{lc}}_W(S)$. De plus, ces rétractions sont ce qu'on appelle des foncteurs de localisation (comme on voit sur le diagramme de la page 62). (On dit que $C\xrightarrow{F}C'$ est un \textit{foncteur de localisation} si, désignant par $\Sigma\subseteq\Fl\,C$ l'ensemble des flèches de $C$ transformées en des isomorphismes par $F$, le foncteur
\[
C\Sigma^{-1}\xrightarrow{\bar{F}}C'
\]
qui factorise $F$ est une équivalence de catégories.)

\msp{page 65}
On sait d'ailleurs que si un foncteur $f$ a un adjoint $g$ (à gauche ou à droite), alors $f$ est un foncteur \texttt{[de]} localisation si et seulement si l'adjoint $g$ est pleinement fidèle. Pour prouver que $i$, $j$ sont pleinement fidèles, il serait donc tentant de prouver qu'ils sont adjoints des foncteurs $\rho$, $\sigma$.

\dgrblock{dgrb193}{What the section of localization still has to prove}
\begin{dcomment}{exposition}{p64}{What the section of localization still has to prove}
Grothendieck has now produced the same locally constant category from both comma directions. The identity of Cat/S induces

S)\ensuremath{-}1] \ensuremath{-}\ensuremath{\to}(Cat/S)[W \ensuremath{-}1

\ensuremath{\rho} : (Cat/S)[(W d S ], and path replacement induces a functor in the reverse direction. The equation \ensuremath{\rho}i \ensuremath{\simeq}id says that the coarse category is a retract. It does not yet say that i is fully faithful; this is the distinction already encountered at manuscript p. 38.

The required additional datum is now visible. There is a natural arrow X \ensuremath{\to}C(X) in the fine category, and it becomes invertible in the coarse category. Its direction suggests a unit for \ensuremath{\rho} \ensuremath{\dashv}i, not for the opposite adjunction. The inverse coarse arrow C(X) \ensuremath{\to}X then supplies a counit. Grothendieck’s change of orientation on p. 67 is consequently forced by where the unit is invertible.

A localization with a fully faithful right adjoint is a reflection onto a full subcategory. In this case its image consists of the relative homotopy types whose transport along every base arrow is invertible. The adjective “locally constant” therefore describes an intrinsic condition on a diagram, not constancy of a chosen presentation of its total category. The path construction supplies a representative satisfying that condition and a universal map to it. The argument on the next pages establishes that universality through the two triangle identities, so that no fullness conclusion is extracted merely from a retraction.

\end{dcomment}


Pour nous y reconnaître, je vais poser
\[
\left\{
\begin{aligned}
&\cA=\Cat/S,\ \ \Sigma_0=W_S^{\mathrm{d}},\ \ \Sigma=W_S\qquad\quad(\subseteq\Fl(\cA))\\
&C:\cA\longrightarrow\cA\ \text{ le foncteur familier, qui a la vertu}\\
&\qquad\qquad\qquad\quad C(\Sigma)\subseteq\Sigma_0\quad.
\end{aligned}
\right.
\]

On a donc la situation
\printedpage{40}
% head: A. Grothendieck : Dérivateurs, Ch. VI, Édition des Universités de Montpellier II et Paris VII
\[
(*)\qquad\qquad(\cA,\Sigma_0)\ \underset{C}{\overset{\rho_0\,=\,\id_{\cA}}{\rightleftarrows}}\ (\cA,\Sigma),
\]
induisant par passage aux localisés
\[
\cA\Sigma_0^{-1}\quad\underset{\bar{C}}{\overset{\rho\,=\,\bar{\rho}_0}{\rightleftarrows}}\ \cA\Sigma^{-1},
\]
on veut vérifier que $\rho$ et $C$ \texttt{[plutôt $\bar{C}$]} sont adjoints l'un de l'autre, $C$ \texttt{[plutôt $\bar{C}$]} \textit{sûrement à droite,} ce qui signifie que l'on doit avoir des homomorphismes fonctoriels
\[
\left\{
\begin{aligned}
&\id_{\cA\Sigma^{-1}}&&\rightleftarrows&&\rho\,\bar{C}\quad(=\overline{\rho_0C})\\
&\id_{\cA\Sigma_0^{-1}}&&\rightleftarrows&&\bar{C}\rho\quad(=\overline{C\rho_0})
\end{aligned}
\right.
\]
(\footnote[41]{mais si $C$ est à \textit{droite,} les flèches ici doivent aller en sens inverse et c'est bien à quoi je dois me résoudre p.~67.}) \texttt{[les flèches en bas vont en direction corrigée, mais les flèches en haut ne sont pas effacées]} satisfaisant deux relations de compatibilité sur lesquelles on reviendra. On peut espérer trouver ces homomorphismes fonctoriels à partir de
\[
\left\{
\begin{aligned}
&\id_{\cA}&&\xrightarrow{\ i\ }&&\rho_0C&&=&&C\\
&\id_{\cA}&&\xleftarrow{\ p\ }&&C\rho_0&&=&&C
\end{aligned}
\right.
\]

\msp{page 66}
et dans notre cas on dispose bel et bien de
\[
i:\id_{\cA}\longrightarrow C\qquad\text{i.e. de }\ i_X:X\longrightarrow C(X)
\]
fonctoriel en $X\in\Ob\,\cA$ (et de plus $i_X$ est dans $\Sigma$). Il n'y a pas cependant de flèche naturelle en sens inverse $p_X$
\[
C(X)\xrightarrow{\ ?\ }X.
\]
J'avais songé à la projection $\pi_X:C(X)\longrightarrow X$, déduite de la construction $C(X)\simeq X\times_S(\ul{Ch}_\infty(S),\sigma)$, mais ce n'est malheureusement pas un $S$-morphisme. Comme on cherche une flèche dans $\cA\Sigma_0^{-1}$, on pourrait se contenter d'avoir une flèche fonctorielle en sens inverse, du moment qu'elle soit dans $\Sigma_0$, donc puisse s'inverser quand on travaille dans $\cA\Sigma_0^{-1}$. Le candidat evident est $i_X:X\longrightarrow C(X)$ dont on vient de parler. L'ennui, c'est que $i_X$ est seulement dans $\Sigma$, pas dans $\Sigma_0$ (\footnote[42]{S'il n'en était pas ainsi, le type d'homotopie relatif sur $S$ défini par $X/S$ serait localement constant (puisqu'il en est ainsi pour celui défini par $C(X)$), ce qui n'est pas toujours le cas.}). Pourtant, quand $X$ est lui-même de la forme $C(Y)$ (plus généralement, s'il est un Hot-fibré sur $S$), alors le fait que $X\longrightarrow C(X)$ soit dans $\Sigma$ implique déjà qu'il soit dans $\Sigma_0$. Notons donc ceci % [sic] evident -> évident, accent missing
\[
\left\{
\begin{array}{lrl}
i:\id_{\cA}\longrightarrow C\ \text{tel que} & 1^o) & i_X:X\longrightarrow C(X)\ \textit{dans}\ \Sigma\\[1ex]
 & & \text{pour tout }X\in\Ob\,\cA\\[1ex]
 & 2^o) & i_X\ \text{dans }\Sigma_0\text{ si }X\text{ de la}\\[1ex]
 & & \text{forme }C(Y)
\end{array}
\right.
\]

\printedpage{41}
% head: A. Grothendieck : Dérivateurs, Ch. VI, Édition des Universités de Montpellier II et Paris VII

\msp{page 67}
L'idée me vient, après des essais merdouillants \texttt{[sic]}, que $\bar{C}$ pourrait être l'adjoint à droite, et non à gauche \texttt{[c'est la version corrigée (droite et gauche échangées) cf.~p.~65]}, i.e. que les flèches cherchées pourraient aller en sens inverse
\[
\left\{
\begin{aligned}
&\id_{\cA\Sigma^{-1}}&&\xleftarrow{\ p\ }&&\rho\,\bar{C}\quad(=\overline{\rho_0C})\\
&\id_{\cA\Sigma_0^{-1}}&&\xrightarrow{\ i\ }&&\bar{C}\rho\quad(=\overline{C\rho_0})
\end{aligned}
\right.
\]
(j'inverse du coup les notations pour $p$ et $i$), donc on est amené si possible à chercher des flèches
\[
\begin{aligned}
&\id_{\cA}&&\xleftarrow{\ p\ }&&\rho_0C&&=&&C\\
&\id_{\cA}&&\xrightarrow{\ i\ }&&C\rho_0&&=&&C,
\end{aligned}
\]
donc des flèches fonctorielles en $X$
\[
X\xrightarrow{i_X}C(X)\xrightarrow{p_X}X.
\]
On tient bel et bien une $i_X$ fonctorielle - toujours la même! Pour $p_X$, on n'en a pas plus qu'avant, mais se rappelant qu'il suffit de le trouver dans $\cA\Sigma^{-1}$ (et non plus $\cA\Sigma_0^{-1}$!), il suffit d'inverser la flèche $i_X$, qui est en effet un \textit{isomorphisme} dans $\cA\Sigma^{-1}$ (mais non dans $\cA\Sigma_0^{-1}$). On posera donc
\[
p_X=i_X^{-1}:C(X)\longrightarrow X\qquad\quad\text{dans }\cA\Sigma^{-1}.
\]
Il faut prouver deux compatibilités seulement, pour les foncteurs $\bar{C}$ et $\rho$ entre $\cA\Sigma_0^{-1}$ et $\cA\Sigma^{-1}$ (et j'écrirai $C$ au lieu de $\bar{C}$) : les deux composés
\[
\left\{
\begin{array}{l}
\begin{tikzcd}[ampersand replacement=\&]
\rho \arrow[r, "\rho*i"] \arrow[rr, bend left=20, dotted, "\id"] \& \rho\,C\rho \arrow[r, "p*\rho"] \& \rho
\end{tikzcd}
\qquad\quad\text{\footnotesize (relation pour foncteurs à valeurs dans $\cA\Sigma^{-1}$)}\\[2ex]
\begin{tikzcd}[ampersand replacement=\&]
C \arrow[r, "i*C"] \arrow[rr, bend right=20, dotted, "\id"'] \& C\rho\,C \arrow[r, "C*p"] \& C
\end{tikzcd}
\qquad\quad\text{\footnotesize (relation pour foncteurs à valeurs dans $\cA\Sigma_0^{-1}$)}
\end{array}
\right.
\]
doivent être l'identité $\id_\rho$ resp. $\id_C$.

\msp{page 68}
La première relation se vérifie tautologiquement, elle revient à dire que dans $\cA\Sigma^{-1}$, on a un composé identité des
\[
X\xrightarrow{i_X}C(X)\xrightarrow{i_X^{-1}}X\qquad\quad!
\]
La seconde signifie que pour tout $X$, on a un composé identité dans $\cA\Sigma_0^{-1}$
\[
C(X)\xrightarrow{i_{C(X)}}C(C(X))\xrightarrow{C(p_X)}C(X).
\]
Or $p_X:C(X)\longrightarrow X$ avait été défini dans $A\Sigma^{-1}$ comme inverse $i_X^{-1}$ de $i_X$, d'autre part le foncteur $C$ sur $\cA$ définit par hypothèse $\cA\Sigma^{-1}\longrightarrow\cA\Sigma_0^{-1}$, plus précisément $C(\Sigma)\subseteq\Sigma_0$ % [sic] $A\Sigma^{-1}$ (italic A) -> $\mathcal{A}\Sigma^{-1}$, elsewhere script A
\printedpage{42}
% head: A. Grothendieck : Dérivateurs, Ch. VI, Édition des Universités de Montpellier II et Paris VII
ainsi $C(i_X)$ est inversible dans $\cA\Sigma_0^{-1}$, et $C(p_X)=C(i_X)^{-1}$. Donc la compatibilité revient à la relation
\[
C(i_X)^{-1}i_{C(X)}=\id\qquad\quad\textit{dans}\ \ \cA\Sigma_0^{-1},
\]
i.e.
\[
\boxed{\,i_{C(X)}=C(i_X)\,}:C(X)\longrightarrow C^2(X).
\]
En d'autre termes, il faut prouver le : % [sic] d'autre termes -> d'autres termes, plural

\dgrblock{dgrb194}{A precise reflection lemma}
\begin{dcomment}{exposition}{p67}{A precise reflection lemma}
The remaining argument is formal once the two classes of arrows and the two uses of the unit are distinguished. We state the lemma in the form needed here; it also shows that the iterated-unit equality can be established without further geometry. This is the ordinary categorical form of the recognition principle for reflective localizations; compare Lurie (2009, Proposition 5.2.7.4) for its higher categorical formulation.

Lemma. Let \ensuremath{\Sigma}0 \ensuremath{\subseteq}\ensuremath{\Sigma} be classes of arrows of a category A, and suppose the localizations below exist. Let C : A \ensuremath{\to}A and \ensuremath{\eta} : id \ensuremath{\to}C satisfy

C(\ensuremath{\Sigma}) \ensuremath{\subseteq}\ensuremath{\Sigma}0,

\ensuremath{\eta}X \ensuremath{\in}\ensuremath{\Sigma}, \ensuremath{\eta}CX \ensuremath{\in}\ensuremath{\Sigma}0 (X \ensuremath{\in}A). Then the identity and C induce functors

P : A0 = A[\ensuremath{\Sigma}\ensuremath{-}1

0 ] \ensuremath{-}\ensuremath{\to}A1 = A[\ensuremath{\Sigma}\ensuremath{-}1], R : A1 \ensuremath{-}\ensuremath{\to}A0, with P \ensuremath{\dashv}R and R fully faithful. The unit is induced by \ensuremath{\eta}, and the counit is the inverse of its image in A1.

Since C(\ensuremath{\Sigma}0) \ensuremath{\subseteq}\ensuremath{\Sigma}0, the endofunctor C and its natural transformation descend to an Proof. endofunctor T = RP and a transformation e : idA0 \ensuremath{\to}T. Naturality extends across the formal inverses because both functors invert the required arrows. At every object x, both

a = eT x : Tx \ensuremath{\to}T 2x,

b = T(ex) : Tx \ensuremath{\to}T 2x are isomorphisms: the first by the third hypothesis, the second by the first two. We claim that a = b.

Naturality at ex gives bex = aex. Apply T and cancel the isomorphism b on the right to obtain T(b) = T(a). Naturality at a gives

T(a)a = eT 2xa,

so T(a) = eT 2x. Naturality at b gives

T(b)a = eT 2xb.

Substitute the preceding equalities and cancel the isomorphism T(a). This proves a = b. Notice that we have never cancelled ex, which need not be invertible in A0.

Let \ensuremath{\epsilon} : PR \ensuremath{\to}idA1 be the inverse of the transformation induced by \ensuremath{\eta} in A1. Its first triangle identity is \ensuremath{\epsilon}P xP(ex) = idP x, true by construction. For the second, every object of A1 is represented by an object X of A. On its image under R, the composite R(\ensuremath{\epsilon}X)eRX is b\ensuremath{-}1a, which is the identity by the calculation above. Hence P \ensuremath{\dashv}R. The counit is an isomorphism, so R is fully faithful. \ensuremath{\Box}

The full image is exactly the class of x for which ex is invertible. Objects in the image satisfy this by the third hypothesis; conversely, invertibility identifies x with RPx. Applied to \ensuremath{\Sigma}0 = W d S and \ensuremath{\Sigma} = WS, this proves the desired full faithfulness of i. The left-comma version proves the corresponding assertion for j, with the covariant convention explained at p. 50. Grothendieck’s concatenation argument is a particularly transparent geometric verification of the same identity, and remains useful for understanding the replacement itself.

\end{dcomment}


\textbf{Lemme.} \textit{Soit $X$ dans $\Cat/S$, considérons les deux flèches}
\[
i_{C(X)},\ C(i_X):C(X)\rightrightarrows C(C(X)).
\]
\textit{Ces flèches sont égales, sinon au sens absolu, du moins dans $(\Cat/S)(W_S^{\mathrm{d}})^{-1}$.}

Ici $f_X:X\longrightarrow S$ est la projection standard et
\[
\begin{aligned}
C(X)\ &=\ \quad\{(x,c)\ \mid\ x\in\Ob\,X,\ c:f_X(x)\ \underbrace{\longrightarrow}_{\text{chemin dans }S}\ t\}\\
C^2(X)\ &=\ \{(x,c,c')\ \mid\ x\in\Ob\,X,\ \underbrace{f_X(x)\xrightarrow{c}t\xrightarrow{c'}t'}_{c,\ c'\text{ chemins dans }S}\},
\end{aligned}
\]
les projections dans $S$ sont données par
\[
\begin{aligned}
(x,c)&&\longmapsto&&t&&=&&\text{but }c\\
(x,c,c')&&\longmapsto&&t'&&=&&\text{but }c'
\end{aligned}
\]

\msp{page 69}
On trouve
\[
i_{C(X)}(x,c)=(x,c,1_t):
\begin{tikzcd}
x \arrow[d] & & \\
f_X(x) \arrow[r, "c"] & t \arrow[r, "1_t", "\text{chemin vide}"'] & t
\end{tikzcd}
\]
\[
C(i_X)(x,c)=(x,1_{f_X(x)},c):
\begin{tikzcd}
x \arrow[d] & & \\
f_X(x) \arrow[r, "1_{f_X(x)}", "\text{chemin vide}"'] & f_X(x) \arrow[r, "c"] & t
\end{tikzcd}
\]
Mais on a le choix entre deux théories de chemins, pour construire $C(X)$ : soit avec les $\ul{Ch}_\infty(S)$, soit avec les $\ul{Ch}(S)$. ‘Ça donne les mêmes résultats’, mais suivant les situations l'un est plus commode que l'autre. Je vais travailler avec les $\ul{Ch}(S)$, auquel cas les $1_S$ sont des unités à gauche et à droite, et on a une loi de composition pour les chemins. Cette loi nous permet de définir
\[
\begin{aligned}
C^2(X)&&\xrightarrow{\ \gamma\ }&&C(X)\\
(x,c,c')&&\longmapsto&&(x,c'\circ c),
\end{aligned}
\]
et les composés de $i_{C(X)}$ et de $C_{i(X)}$ avec cette flèche sont $\id_{C(X)}$ dans les deux cas % [sic] $C_{i(X)}$ -> $C(i_X)$, the arrow considered is C(i_X)
\printedpage{43}
% head: A. Grothendieck : Dérivateurs, Ch. VI, Édition des Universités de Montpellier II et Paris VII
\[
\begin{tikzcd}
C(X) \arrow[r, shift left] \arrow[r, shift right] \arrow[rr, bend right=15, "\id"'] & C^2(X) \arrow[r, "\gamma"] & C(X).
\end{tikzcd}
\]

Donc pour montrer que $i_{C(X)}=C(i_X)$ dans $(\Cat/S)(W_S^{\mathrm{d}})^{-1}$, il suffit de prouver que $\gamma$ est dans $W_S^{\mathrm{d}}$. Mais on sait que le foncteur C transforme $W_S$ en $W_S^{\mathrm{d}}$, donc $C(i_X)$ est dans $W_S^{\mathrm{d}}$, et comme $\gamma\circ C(i_X)=\id$, $\gamma$ est dans $W_S^{\mathrm{d}}$, OK! % [sic] foncteur C (upright roman C) -> foncteur $C$, math italic elsewhere

\msp{page 70}
J'ai envie de grouper ce qui a été prouvé dans un énoncé récapitulatif.

\dgrblock{dgrb195}{Concatenation and the two iterated units}
\begin{dcomment}{exposition}{p69}{Concatenation and the two iterated units}
Grothendieck’s path model makes the reflection tangible. An object of C(X) is an object x together with a path from p(x) to its new base point. Iterating C adds a second path. The two units place the empty path in different positions, so there is no reason for the resulting functors to be literally equal before localization.

Using the finite-path version with strict empty-path units, concatenation gives a functor over S

(x, c, c\ensuremath{\prime}) 7\ensuremath{-}\ensuremath{\to}(x, c\ensuremath{\prime}c).

\ensuremath{\mu}X : C2(X) \ensuremath{-}\ensuremath{\to}C(X), On morphisms it concatenates the two compatible path diagrams; compatibility of their endpoints makes this well-defined, and composition is respected because it is respected in each path segment. The two empty-path laws give the strict equations

\ensuremath{\mu}X\ensuremath{\eta}C(X) = idC(X),

\ensuremath{\mu}XC(\ensuremath{\eta}X) = idC(X). The relevant homotopical input is C(WS) \ensuremath{\subseteq}W d S. Since \ensuremath{\eta}X \ensuremath{\in}WS, we have C(\ensuremath{\eta}X) \ensuremath{\in}W d S. Two-out-ofthree applied to the second equation yields \ensuremath{\mu}X \ensuremath{\in}W d S. After fine localization both \ensuremath{\eta}C(X) and C(\ensuremath{\eta}X) are therefore the inverse of the same arrow \ensuremath{\mu}X, which proves their equality.

This proof does not cancel the unit X \ensuremath{\to}C(X) in the fine category. That would erase precisely the distinction the chapter has worked to establish. Instead it cancels a map between already completed objects. Nor does the argument require the associativity of concatenation: the two unit identities and two-out-of-three suffice for the triangle equation. Associativity equips C with the familiar monadic structure, but the adjunction needs only the part displayed here.

The formal lemma in the preceding block shows that the equality also follows from the stated hypotheses on pp. 65--66. Thus there is no outstanding coherence gap in this part of the proof. The choice of finite paths supplies an explicit common inverse and exhibits why adjoining paths becomes idempotent after the appropriate localization.

\end{dcomment}


\dgrblock{dgrb196}{Path completion and the modern category of local systems}
\begin{dcomment}{later}{p69}{Path completion and the modern category of local systems}
The theorem reaches its most familiar modern form when W = W\ensuremath{\infty}. Let S denote the infinitycategory of spaces, and let

\ensuremath{\ell}: N(S)op \ensuremath{-}\ensuremath{\to}\ensuremath{\Pi}\ensuremath{\infty}N(S) be groupoid completion; we identify the completion of N(S)op with that of N(S) by reversing paths. Then the fine theory models homotopy coherent contravariant space-valued diagrams, and the locally constant theory models diagrams which carry every base arrow to an equivalence:

HOTW\ensuremath{\infty}(S) \ensuremath{\simeq}Ho Fun(N(S)op, S),

HOTlc

W\ensuremath{\infty}(S) \ensuremath{\simeq}Ho Fun(\ensuremath{\Pi}\ensuremath{\infty}N(S), S). The first comparison can be justified through an actual model equivalence. Thomason’s adjunction

c Sd2 : sSet \ensuremath{\rightleftarrows}Cat : Ex2 N

is a Quillen equivalence for nerve weak equivalences (Thomason, 1980, Theorem 4.9 and the introduction). Apply both functors objectwise to Sop-diagrams with their projective model structures. The right adjoint preserves fibrations and trivial fibrations objectwise. If F is projectively cofibrant and G projectively fibrant, the criterion for an adjoint pair of maps to be weak equivalences is checked at each object using the original Quillen equivalence; projective cofibrancy implies objectwise cofibrancy. Thus the diagram adjunction is again a Quillen equivalence. The simplicial-set diagram model then presents the coherent functor category by Lurie (2009, Proposition 4.2.4.4). This explains why strict category-valued diagrams, after objectwise localization, retain coherent space-valued diagrams.

Here Ho is taken after forming the category of coherent diagrams. It is not the category of ordinary functors into Ho(S). The higher transport information discussed earlier in the chapter is retained by this order of operations.

Restriction \ensuremath{\ell}\ensuremath{*}is fully faithful by the universal property of groupoid completion. Its left adjoint is left Kan extension \ensuremath{\ell}!, and the reflector on diagrams is \ensuremath{\ell}\ensuremath{*}\ensuremath{\ell}!. This identifies the universal property of Grothendieck’s C, rather than asserting that the particular path-category construction is literally a chosen infinity-categorical Kan extension. The standard straightening comparison and the adjunction for Kan extension are given in Lurie (2009, Theorem 2.2.1.2 and Proposition 4.3.3.7); the reflection for a left proper combinatorial model category of coefficients is Cisinski (2009, Corollary 2.5).

A local system of spaces on the infinity-groupoid \ensuremath{\Pi}\ensuremath{\infty}N(S) can equivalently be viewed as a space over |N(S)|, by taking its total space and homotopy fibres. Consequently the coarse theory has the interpretation

HOTlc

W\ensuremath{\infty}(S) \ensuremath{\simeq}Ho(S/|N(S)|). For a direct model-categorical comparison see Cisinski (2009, Example 1.7, equation (1.7.1)). This is a homotopy category of a slice, and not generally the ordinary slice of a homotopy category. The distinction already matters for a base with nontrivial fundamental group, as the following example shows.

These interpretations explain what became of the full-faithfulness question: for spaces it is settled by the reflector, and more general coefficient theories have their own locally constant subcategories and reflection results. They do not justify identifying an arbitrary W-coarse localization with that subcategory. Such an identification requires the additional exactness properties discussed on pp. 43--45. The existence of a locally constant reflection and the adequacy of a particular coarse test are separate assertions.

\end{dcomment}


\dgrblock{dgrb197}{Three theories within the comparison diagram}
\begin{dcomment}{exposition}{p70:J'ai envie de grouper ce qui a été prouvé}{Three theories within the comparison diagram}
Grothendieck now wants the many presentations to express one intelligible distinction. The proof that supplies it has already been explained in the commentary at manuscript pp.~64, 67, and 69: the path replacement defines a section of the coarse localization, and the two iterated units agree after fine localization. Here the new work is to organize the consequences and identify the image of that section.

For a map $v:X\to Y$ over $S$, the three tests are
\[
\begin{array}{c|l}
 W_S&v\in W\text{ on total categories},\\
 W_S^{\mathrm d}&s\backslash X\longrightarrow s\backslash Y\in W
                    \text{ for every }s,\\
 W_S^{\mathrm g}&X/s\longrightarrow Y/s\in W
                    \text{ for every }s.
\end{array}
\]
The relative Theorem A supplies the inclusions of both fine classes in $W_S$. Thus the identity on $\Cat/S$ gives $\rho$ and $\sigma$; these are quotient functors between homotopy categories, before any assertion about adjoints. In the other direction one must first replace $X$: the functor $C$ sends a coarse equivalence to both kinds of fine equivalence. This is why the inclusion of locally constant objects is represented by a nontrivial construction on total categories.

The replacement by a split fibration does something different. Its fibre at $s$ is $s\backslash X$, so it preserves the right-fine information already present in $X$. Dually, the split opfibration has fibre $X/s$. Neither construction forces its transition maps to become equivalences. The path replacement does force this, and consequently belongs to the upper part of the diagram. Confusing these two purposes would make the distinction between $\HOT_W(S)$ and $\HOT_W^{\mathrm{lc}}(S)$ disappear before it has been proved.

Part~I is therefore a statement about coarse localization of all the presentations. Part~II identifies the appropriate fine localization along each side. Part~III says when a strict fibre is a valid test and when the two fine tests agree. The hypothesis L5 is used at precisely the passage from a coarse equivalence between homotopy fibrations to a fibrewise equivalence. It should not be silently dropped when the diagram is reused. The subsequent remarks at manuscript pp.~81--82 explicitly reopen this dependency.
\end{dcomment}


\textbf{Théorème-Scholie.} On considère le diagramme de catégories de la page 8 (\footnote[43]{On suppose sur $W$ les conditions a) b) c) de p.~45, de plus L5 (p.~43). Cf. commentaires p.~81.}).

\circled{I} Quand on munit toutes les 13 catégories du diagramme du localiseur $W_S$ ($W$-équivalence grossière, i.e. $W$-équivalence dans Cat, sans plus), \textit{tous} les foncteurs du diagramme induisent des équivalences sur les catégories localisées (\footnote[44]{De plus, si $\mathcal{F}$ est une sous-catégorie pleine de $\Cat/S$ qui contient une des deux sous-catégories $\ul{Fib}_0S$, $\ul{Cofib}_0S$ (ce qui inclut le cas des 8 catégories des diagrammes, autres que $\Cat/S$ elle-même), l'inclusion $(\mathcal{F},W_S|\mathcal{F})$ dans $(\Cat/S,W_S)$ induit une équivalence pour les localisés.}). Désignons par
\[
\HOT^{\mathrm{lc}}_W(S)
\]
‘la’ catégorie localisée commune.

\circled{II} Munissons les quatre catégories bordantes gauches
\[
(*)\qquad\qquad
\begin{tikzcd}
\Cat/S & \ul{Liss}\,S \arrow[l, hook'] & \ul{Fib}\,S \arrow[l, hook'] & \ul{Fibsc}\,S \arrow[l, hook']
\end{tikzcd}
\]
du localiseur $W_S^{\mathrm{d}}$ de la $W$-équivalence fine \textit{droite} sur $S$ ( $\begin{tikzcd}[row sep=small, column sep=small]
X \arrow[rr, "f"] \arrow[dr] & & Y \arrow[dl]\\
 & S &
\end{tikzcd}$ est dans $W_S^{\mathrm{d}}$ si et seulement si les $s\backslash X\longrightarrow s\backslash Y$ le sont \texttt{[plutôt, sss ils sont dans $W$]}, pour tout $s\in\Ob\,S$). Les foncteurs du diagramme, i.e. ceux de $(*)$, entre ces catégories, donc aussi leurs composés, induisent des équivalences entre les localisées (\footnote[45]{Plus généralement, ça marche pour toute sous-catégorie pleine de $\Cat/S$ qui contient $\ul{Fib}\,S$ (et dualement, avec $W_S^{\mathrm{g}}$ au lieu de $W_S^{\mathrm{d}}$, si $\mathcal{F}$ contient $\ul{Cofib}\,S$).}). Désignons par
\[
\HOT_W(S)
\]
‘la’ catégorie localisée commune.

Enoncé dual pour
\[
(*')\qquad\qquad
\begin{tikzcd}
\Cat/S & \ul{Prop}\,S \arrow[l, hook'] & \ul{Cofib}\,S \arrow[l, hook'] & \ul{Cofibsc}\,S, \arrow[l, hook']
\end{tikzcd}
\]

% foot (each page): Version 9 novembre 2001


\dgrblock{dgrb198}{Two computations of the reflection}
\begin{dcomment}{examples}{p70}{Two computations of the reflection}
The reflection records how the base arrows must become invertible. For the interval S = [1], a contravariant diagram is an arrow B \ensuremath{-}\ensuremath{\to}uA. Since |N[1]| is contractible, locally constant diagrams are equivalent to constant diagrams. The reflector is

\ensuremath{-}\ensuremath{\to}uA) 7\ensuremath{-}\ensuremath{\to}(A

\ensuremath{-}\ensuremath{\to}A), (B with unit given by u at 1 and the identity at 0. To verify its universal property at the coherent level, a map from B \ensuremath{\to}A to the constant diagram on K consists of maps A \ensuremath{\to}K, B \ensuremath{\to}K, and a specified homotopy relating the second to the composite through A. The space of choices of the second map together with such a homotopy is contractible over the first. Thus the mapping space is naturally equivalent to Map(A, K), exactly as required. This argument also includes empty values.

For example, the diagram \ensuremath{\varnothing}\ensuremath{\to}\ensuremath{*}becomes the constant point diagram. Its total Grothendieck construction is the category supported at 0, while that of the reflected diagram is [1]. Their total nerves are equivalent, although the unit fails to be an objectwise equivalence at 1. This is the earlier counterexample viewed as the unit of an adjunction. The inclusion of constant diagrams also has a right adjoint, sending B \ensuremath{\to}A to B: the dual mapping-space computation gives Map(K, B). The left and right adjoints are distinct constructions even for this smallest base.

For a second computation take the one-object groupoid S = BZ. Every arrow of the base is already invertible, so every coherent diagram on S is locally constant and the reflection is an equivalence. Nevertheless the coarse slice contains more morphisms than the ordinary slice of Ho(S). Consider the object z : \ensuremath{*}\ensuremath{\to}BZ given by a base point. In the infinity-categorical slice,

EndHo(S/BZ)(z) \ensuremath{\sim}= Z.

Map/BZ(z, z) \ensuremath{\simeq}\ensuremath{\Omega}BZ, Indeed, the underlying map \ensuremath{*}\ensuremath{\to}\ensuremath{*}is unique, but being a map over BZ includes a path from the chosen base point to itself. Homotopy classes of such paths are the integers. In the ordinary category Ho(S)/BZ, the same object has only its identity endomorphism, because the over-condition is imposed merely as an equality of homotopy classes.

The associated local system has fibre Z with its regular translation action. A translation-equivariant map Z \ensuremath{\to}Z is determined by the image n of 0 and is k 7\ensuremath{\to}k + n; composition adds the integers. This recovers the endomorphism calculation directly from local systems. Path completion of the point over BZ supplies the categorical analogue of its universal-cover replacement, recovering exactly the transport that the ordinary slice would forget.

\end{dcomment}
\begin{otherlanguage}{english}
\dgroverview{VI-a}{Chapter VI \textperiodcentered\ \textup{VI-a} \textperiodcentered\ Overview}
% Text recovered from pages 57--59 of the original complete reader PDF; mathematical notation flattened by extraction.
The problem of the chapter. Grothendieck starts from a category of coefficients, here the category of small categories itself, and asks how its homotopy theory should vary with a diagram shape. The value at a small category S is to be


HOTW (S) = Fun(Sop, Cat)[W \ensuremath{-}1


obj], where a natural transformation belongs to Wobj when all its components belong to W. The essential operation is to localize the category of diagrams before forgetting their coherence. A functor from Sop to the already localized category of coefficients generally retains less information. This is the same motivation that led to the diagram formulation in Grothendieck (1983, \S{}69), and that the letter to Thomason subsequently places at the centre of the derivator programme (Grothendieck, 1991). The discussions of coherence and evaluation in the commentary on Chapter~I, manuscript pages 7, 17--18 and 27, remain the background; the present chapter supplies categorical constructions rather than repeating those objections.


The first seventy manuscript pages, dated 14--15 November 1990, record the discovery of the correct comparison theorem. Thirteen categories of objects over S initially appear to offer interchangeable models. They include split fibrations, arbitrary fibrations, smooth and proper functors, homotopy fibrations, perfect functors, and all of Cat/S. The objects alone do not decide the homotopy theory. The permitted morphisms and the chosen weak equivalences matter just as much. The first theorem overlooks that distinction, is expressly withdrawn, and is replaced by a comparison between three theories. The manuscript’s revisions are part of the proof’s development, not noise to remove from it. The axioms and the four tests. The axiom list also changes. The original L4 is an absolute Theorem A: an aspherical functor is a weak equivalence. L4 bis, introduced on manuscript page 44, is the relative assertion about a map X \ensuremath{\to}Y over a third category S, tested on every map X/s \ensuremath{\to}Y/s. It is the latter formulation that supports the relative descent arguments. The initial L5 passes from fibrewise weak equivalences of Cartesian maps between fibrations to weak equivalences of total categories. The new L5 on page 43 goes in the opposite direction for maps between homotopy fibrations. Their common label must not obscure their different content. The commentary proves the initial claim that L1, L2 bis and the original L5 already imply the required elementary axioms, including relative descent, by a twisted-arrow argument for opposite invariance.


For a map f : X \ensuremath{\to}Y over S, the decisive tests are


coarse


f \ensuremath{\in}W,


fibrewise fs : Xs \ensuremath{\to}Ys belongs to W for every s,


left-fine


f/s : X/s \ensuremath{\to}Y/s belongs to W for every s,


right-fine


s\textbackslash{}f : s\textbackslash{}X \ensuremath{\to}s\textbackslash{}Y belongs to W for every s. Over a smooth object, the strict fibre compares to the undercategory; over a proper object, it compares to the overcategory. These comparisons and relative descent explain when fibrewise equivalence implies a universal equivalence of totals. The modern asphericity formulation retains precisely this mechanism (Maltsiniotis, 2005, Theorem 1.2, Corollary 4.16 and Theorem 5.4). Coarse equivalence alone does not give either fine condition on the external borders of the diagram. Over [1], the inclusion \{0\} \ensuremath{\to}[1] is a coarse equivalence of split fibrations, but its fibre at 1 is the map from the empty category to a point. Evaluation at that fibre proves that the two localizations themselves differ. The two comma resolutions. The functor \ensuremath{\Phi} replaces X \ensuremath{\to}S by the split fibration whose fibre at s is s\textbackslash{}X; \ensuremath{\Psi} gives the split opfibration with fibre X/s. Their units are categorical homotopy equivalences on total categories. More precisely, the unit of \ensuremath{\Phi} is right-fine, and the unit of \ensuremath{\Psi} is left-fine. Composition of the two arrows in a comma object gives an explicit contraction proving this stronger fact. Consequently the right-fine localization of Cat/S computes the contravariant diagram theory, while the left-fine localization computes its covariant counterpart.


The comparison with chosen splittings needs a genuine repair of the argument, but not a change in the desired theorem. A Cartesian functor need not preserve the selected Cartesian lifts strictly, and an arbitrary functor over the base need not preserve Cartesian arrows at all. The actual free split adjunction is


Homsplit/S(\ensuremath{\Phi}X, Y ) \ensuremath{\sim}= HomS(X, UY ).


Its counit supplies the comparison in the strict category. For arbitrary functors between fibrations the pullback retraction has a canonical, possibly noninvertible natural comparison. Natural transformations over S induce equality after fine localization, since the cylinder projection A \ensuremath{\times} [1] \ensuremath{\to}A is fine. This yields the natural isomorphism needed after localization. The exposition constructs the adjunction and the localized comparison in full; the historical and modern strictification references are separated from this particular elaboration (Gray, 1966; Emmenegger et al., 2024).


The covariant theory is HOTW (Sop). The additional outer opposite that appears in several passages of the edition is incorrect with the initial definition above: taking opposite total categories preserves the direction of functors between them. The corrected variance is used throughout the commentary, while both the original statements and the editorial brackets are preserved in the source. The strength of path completion. The two-sided path replacement gives a natural map X \ensuremath{\to}C(X) which is a weak homotopy equivalence, with C(X) \ensuremath{\to}S perfect. This immediately proves comparisons for the coarse localizations, without the new L5. To identify coarse and fibrewise equivalences on the internal categories, and then identify their common theory with a full subcategory of the fine theory, is stronger. That step is what the new L5 supplies.


Its restriction on W is severe. A degree-n covering of the circle is an equivalence for homology with coefficients in a ring where n is invertible, but its fibre has n components. The commentary realizes the covering as an action groupoid, so that perfectness and the strict fibre can be checked


directly. For integral homology it gives a different explicit example:


Z[1/2] \ensuremath{\rtimes}2 Z


  \ensuremath{-}\ensuremath{\to}BZ. B The fibre has first homology Z[1/2]. The mapping-torus sequence shows that the total projection is nevertheless an integral homology equivalence, because 1 \ensuremath{-}2 acts invertibly on that fibre homology. This calculation settles the tentative integral-coefficient assertion on manuscript page 43. The development label marks the supplied application and proof, with a bounded account of the literature checked; it claims no priority for the group or the homology calculation.


Finite homotopy truncations fail as well: a path fibration over K(A, n + 1) has contractible total space and fibre K(A, n). Retaining the standing condition W\ensuremath{\infty}\ensuremath{\subseteq}W \ensuremath{\subseteq}W0, the new L5 actually forces W = W\ensuremath{\infty}. Iterating path fibres shows that every W-aspherical category has connected iterated loop spaces, hence is weakly contractible; applying this to the fibres of a replacement of an arbitrary W-map proves the conclusion. The shortened annotation omitting the component condition cannot stand literally, since the maximal class of all functors is a counterexample. Cisinski’s later characterization explicitly includes the corresponding condition on components (Cisinski, 2004, Theorem 2.3.6). The reflection and its modern meaning. The final argument proves full faithfulness through an adjunction. Put


\ensuremath{\Sigma}0 = W d


A = Cat/S, S \ensuremath{\subseteq}\ensuremath{\Sigma} = WS. Path completion carries \ensuremath{\Sigma} into \ensuremath{\Sigma}0, its unit \ensuremath{\eta}X belongs to \ensuremath{\Sigma}, and \ensuremath{\eta}CX belongs to \ensuremath{\Sigma}0. These conditions give


P : A[\ensuremath{\Sigma}\ensuremath{-}1


0 ] \ensuremath{\rightleftarrows}A[\ensuremath{\Sigma}\ensuremath{-}1] : R, P \ensuremath{\dashv}R, with R fully faithful. The unit is induced by \ensuremath{\eta}, and the counit is its inverse in the coarse category. The commentary verifies both triangle identities. In particular, the equality C(\ensuremath{\eta}X) = \ensuremath{\eta}CX after fine localization follows formally from naturality and the stated invertibilities. Grothendieck’s concatenation of paths gives an additional geometric proof by providing a common inverse. There is no need to cancel \ensuremath{\eta}X in the fine category, where it need not be invertible.


For W = W\ensuremath{\infty}, the resulting interpretation is


HOTW\ensuremath{\infty}(S) \ensuremath{\simeq}Ho Fun(N(S)op, S),


HOTlc


W\ensuremath{\infty}(S) \ensuremath{\simeq}Ho Fun(\ensuremath{\Pi}\ensuremath{\infty}N(S), S) \ensuremath{\simeq}Ho(S/|N(S)|). Here S is the infinity-category of spaces, and the homotopy category is taken after forming coherent diagrams or the slice. The reflector is restriction after left Kan extension along groupoid completion. The commentary explains the comparison through Thomason’s model and straightening (Thomason, 1980; Lurie, 2009). It also computes the interval reflector and the endomorphisms of the point over BZ, which distinguish the homotopy category of the slice from the ordinary slice of the homotopy category.


The locally constant embedding and its two adjoints are now theorems in the derivator setting (Cisinski, 2009, Proposition 3.10 and Theorem 3.13). Their existence does not identify an arbitrary coarse W-localization with that subcategory. The surviving qualification is about hypotheses, not an unresolved version of the full-faithfulness problem for spaces. We also do not label the


relationship between the weakest absolute and relative axiom packages as a current open problem: this audit establishes the implications used here and records where a weaker package is insufficient for the printed proof. The construction of the path categories themselves is an input from outside these pages. The first part of the chapter ends in the opening clauses of the theorem--scholium on manuscript page 70; the continuation on pages 71--81 is commented in the second part.
\end{otherlanguage}
\dgrchunkend

% DER_VI-b: Grothendieck's text with the commentary set in, then the overview, inside a group that applies the chunk's own definitions
\dgrpart{1}{VI-b}{Part 2 of 3, manuscript pages 71--145}{VI-b\enspace\textperiodcentered\enspace Chapter VI, part 2 of 3\enspace\textperiodcentered\enspace manuscript pages 71--145}{Chapter VI \textperiodcentered\ \textup{VI-b}}
\dgrchunkbegin
% DER_VI-b: the definitions of the chunk's preamble, applied inside its group
\let\printedpage\relax\newcommand{\printedpage}[1]{\marginpar{\footnotesize p.~#1}}
\let\msp\relax\newcommand{\msp}[1]{\marginpar{\footnotesize [#1]}}
\let\HOT\relax\newcommand{\HOT}{\mathrm{HOT}}
\let\Cat\relax\newcommand{\Cat}{\mathrm{Cat}}
\let\Ob\relax\newcommand{\Ob}{\mathrm{Ob}}
\let\Fl\relax\newcommand{\Fl}{\mathrm{Fl}}
\let\id\relax\newcommand{\id}{\mathrm{id}}
\let\cA\relax\newcommand{\cA}{\mathcal{A}}
\let\ul\relax\newcommand{\ul}[1]{\underline{\mathrm{#1}}}
\let\circled\relax\newcommand{\circled}[1]{\tikz[baseline=(c.base)]\node[draw,circle,inner sep=1pt](c){#1};}
\let\Hot\relax\newcommand{\Hot}{\mathrm{Hot}}
\let\hot\relax\newcommand{\hot}{\mathrm{hot}}
\let\Hom\relax\newcommand{\Hom}{\operatorname{Hom}}
\let\cercle\relax\newcommand{\cercle}[1]{\tikz[baseline=(c.base)]\node[draw,circle,inner sep=1pt](c){#1};}
\let\ovale\relax\newcommand{\ovale}[1]{\tikz[baseline=(c.base)]\node[draw,rounded corners=0.8ex,inner sep=2pt](c){#1};}
\let\usl\relax\newcommand{\usl}[1]{\underline{\mathrm{#1}}}
\let\encadre\relax\newcommand{\encadre}[1]{\tikz[baseline=(n.base)]\node[draw,rounded corners=3pt,inner sep=1.5pt](n){#1};}
\let\cM\relax\newcommand{\cM}{\mathcal{M}}
\let\cF\relax\newcommand{\cF}{\mathcal{F}}
\let\Fib\relax\newcommand{\Fib}{\underline{\mathrm{Fib}}\,}
\let\Fibhot\relax\newcommand{\Fibhot}{\underline{\mathrm{Fibhot}}\,}
\let\Ho\relax\newcommand{\Ho}{\mathrm{Ho}}
\let\Fibu\relax\newcommand{\Fibu}{\underline{\mathrm{Fib}}\,}
\let\WL\relax\newcommand{\WL}{W_S^{\mathcal{L}}}
\let\defeq\relax\newcommand{\defeq}{\overset{\text{déf}}{=}}
\let\Ch\relax\newcommand{\Ch}{\mathrm{Ch}}
\let\pr\relax\newcommand{\pr}{\mathrm{pr}}
\let\Fibsc\relax\newcommand{\Fibsc}{\underline{\mathrm{Fibsc}}\,}
\let\Cofib\relax\newcommand{\Cofib}{\underline{\mathrm{Cofib}}\,}
\let\Cofibsc\relax\newcommand{\Cofibsc}{\underline{\mathrm{Cofibsc}}\,}
\let\Parf\relax\newcommand{\Parf}{\underline{\mathrm{Parf}}\,}
\let\Liss\relax\newcommand{\Liss}{\underline{\mathrm{Liss}}}
\let\Prop\relax\newcommand{\Prop}{\underline{\mathrm{Prop}}}
\let\M\relax\newcommand{\M}{\mathcal{M}}
\let\und\relax\newcommand{\und}[1]{\underline{\mathrm{#1}}}
\let\uHom\relax\newcommand{\uHom}{\underline{\mathrm{Hom}}}
\let\uFib\relax\newcommand{\uFib}{\underline{\mathrm{Fib}}}
\let\uFibsc\relax\newcommand{\uFibsc}{\underline{\mathrm{Fibsc}}}
\let\uCh\relax\newcommand{\uCh}{\underline{\mathrm{Ch}}}
\let\Wd\relax\newcommand{\Wd}{W^{\mathrm{d}}}

% ---- from Der-ChVI-345_p046-050.tex
\setcounter{footnote}{45}

\printedpage{44}
% head: A. Grothendieck : Dérivateurs, Ch. VI, Édition des Universités de Montpellier II et Paris VII
% foot: Version 9 novembre 2001
et les localiseurs $W_S^{\mathrm{g}}$ dans ces catégories ($f\in W_S^{\mathrm{g}}$ si et seulement si les $f/s : X/s\longrightarrow Y/s$ sont dans $W$). On trouve une `même catégorie localisée', notée
\[ \HOT^o_W(S), \]
et on a tautologiquement
\msp{page 71}
\[ \HOT^o_W(S)\approx\HOT_W(S^o). \]
\texttt{[plutôt} $(\HOT_W(S^o))^o$ \texttt{cf. p. 50]}

\ovale{III}\quad Sur les six catégories internes du diagramme, les trois notions de $W$-équivalence, en termes de $W_S$, $W_S^{\mathrm{d}}$, $W_S^{\mathrm{g}}$ coïncident, et coïncident avec la $W$-équivalence par fibres $W_S^{\mathrm{f}}$. Sur les trois catégories de $(*)$ autres que $\Cat/S$, les localiseurs $W_S^{\mathrm{f}}$, $W_S^{\mathrm{d}}$ coïncident; sur les trois catégories de $(*')$ autres que $\Cat/S$, les localiseurs $W_S^{\mathrm{f}}$, $W_S^{\mathrm{g}}$ coïncident. Mais en général (\footnote{sans doute dans tous les cas où $S$ n'est pas équivalente à une catégorie dicrète.% [sic] dicrète -> discrète, faute d'orthographe
}), sur ces six catégories, ces localiseurs $W_S^{\mathrm{f}}$ ne sont pas égaux à $W_S$, plus précisément donnent des localisées non équivalentes, elles sont donc munies chacune de \textit{deux} localiseurs essentiellement distincts, $W_S^{\mathrm{f}}$ et $W_S$. La catégorie $\Cat/S$ est la seule catégorie du diagramme sur laquelle les localiseurs $W_S$, $W_S^{\mathrm{d}}$ et $W_S^{\mathrm{g}}$ sont tous les trois essentiellment distincts. % [sic] essentiellment -> essentiellement, lettre omise
C'est la seule qui permette d'obtenir simultanément les trois catégories localisées, comme
\[
\begin{tikzcd}[column sep=small]
 & \HOT^{\mathrm{lc}}_W(S)\simeq(\Cat/S)W_S^{-1} & \\
\HOT_W(S)\simeq(\Cat/S)(W_S^{\mathrm{d}})^{-1} \arrow[ur,"\rho"] & & \HOT^o_W(S)\simeq(\Cat/S)(W_S^{\mathrm{g}})^{-1} \arrow[ul,"\sigma"']
\end{tikzcd}
\]

\dgrblock{dgrb199}{Opposite categories and opposite derivators}
\begin{dcomment}{exposition}{p71:Sur les six catégories internes du diagramme}{Opposite categories and opposite derivators}
The dual statement is meant to recover the covariant theory from the contravariant one. At the level of the ordinary localized categories used here, its variance can be checked directly. Taking opposites gives a functor
\[
 O_S:\Cat/S\longrightarrow\Cat/S^{\mathrm{op}},
 \qquad (X\to S)\longmapsto(X^{\mathrm{op}}\to S^{\mathrm{op}}),
\]
which sends $v:X\to Y$ to $v^{\mathrm{op}}:X^{\mathrm{op}}\to Y^{\mathrm{op}}$, with the same direction. Moreover
\[
 s\backslash(X^{\mathrm{op}})\cong(X/s)^{\mathrm{op}}.
\]
For the opposite-invariant localizer under consideration, $O_S$ interchanges the left-fine and right-fine classes. Its square is the identity. Consequently the localization definitions of this Scholie give
\[
 (\Cat/S)(W_S^{\mathrm g})^{-1}
       \simeq(\Cat/S^{\mathrm{op}})(W_{S^{\mathrm{op}}}^{\mathrm d})^{-1},
\]
with no opposite on the resulting category.

The editorial insertion on this page asks for an additional outer opposite. That insertion belongs to the different convention of an \emph{opposite derivator}, whose value is $\mathbb D(S^{\mathrm{op}})^{\mathrm{op}}$; it is not compatible with identifying both sides here with the displayed ordinary localizations. We preserve the edition's insertion, and distinguish the conventions rather than altering its words. See also manuscript p.~50, where the same issue first appears. The two diagram functors at manuscript p.~73 confirm the present interpretation: their values lie in $\Hot_W$, with the index category reversed, rather than in $\Hot_W^{\mathrm{op}}$.

A small check prevents a significant error later. At $S=e$ and for ordinary weak homotopy equivalences, both fine localizations are $\Ho(\mathrm{Spaces})$. This category is not equivalent to its opposite: its initial object is strict, since a map to the empty space can only have empty source, whereas its terminal object is not costrict: a map $*\to S^0$ need not be invertible. These properties are preserved by equivalences and exchanged by passage to the opposite. In particular dualizing the present comparison does not reverse the order of the two adjoints: both $\rho\dashv i$ and $\sigma\dashv j$ have the order printed in Part~IV. The later hesitation at manuscript p.~84 comes from precisely this variance distinction.
\end{dcomment}


\dgrblock{dgrb200}{The groupoid test for collapse of the comparison}
\begin{dcomment}{development}{p71:Sur les six catégories internes du diagramme}{The groupoid test for collapse of the comparison}
Grothendieck asks in footnote~46 whether the distinction between coarse and fine localization persists for every base not equivalent to a discrete category. A group with one object already shows that the proposed dividing line is too narrow. The relevant issue is the presence of noninvertible arrows, not the presence of nontrivial automorphisms.

\textbf{Proposition (the canonical comparison).}\ours
For $W=W_\infty$, the ordinary weak homotopy equivalences of categories, the canonical localization
\[
 \rho:\HOT_W(S)\longrightarrow\HOT_W^{\mathrm{lc}}(S)
\]
is an equivalence if and only if $S$ is a groupoid. The same assertion holds for $\sigma$.

We use the established interpretation of the right-fine theory as coherent space-valued presheaves and its locally constant subcategory. Equivalently, one can use Part~V and the split-fibration presentation; the proof below needs only discrete representables for the difficult implication. If $S$ is a groupoid, every functor defined on $S^{\mathrm{op}}$ sends arrows to equivalences. Hence every object is locally constant, and the fully faithful right adjoint $i$ is essentially surjective. Its left adjoint $\rho$ is therefore an equivalence.

Conversely, suppose $\rho$ is an equivalence. For each object $z$, consider the discrete presheaf $h_z(s)=\Hom_S(s,z)$ and its category of elements over $S$. It defines an object of $\HOT_W(S)$. Since this object must be locally constant, precomposition by every $u:a\to b$ induces a weak equivalence of discrete spaces
\[
 \Hom_S(b,z)\xrightarrow{-\circ u}\Hom_S(a,z).
\]
A weak equivalence between discrete spaces is a bijection. With $z=a$, surjectivity gives $v:b\to a$ with $vu=1_a$. With $z=b$, injectivity applied to $uv$ and $1_b$ gives $uv=1_b$, because both become $u$ after precomposition by $u$. Thus every arrow is invertible. Reversing the argument proves the covariant assertion.

For example, take $S=BG$ with $G\ne1$. Its unique object has automorphism group $G$, so $BG$ is not equivalent to a discrete category; nevertheless both canonical comparisons are equivalences. This disproves the printed conjecture even for the ordinary localizer. The proposition concerns the specified comparison functor: detecting a nonconstant representable proves that this functor fails to be an equivalence, without asserting an unrelated classification of possible abstract equivalences between the two categories.

This is our explicit application of Yoneda to the question in the footnote, not a claim of priority. We checked the locally constant comparison in \citet[\S\S1--2]{Cisinski2009LocallyConstant} and the presheaf interpretation in \citet[Theorem~2.2.1.2]{Lurie2009HTT}; neither is being claimed as a new result. We did not find the footnote's proposed criterion stated in the passages checked.
\end{dcomment}


Ici, $\rho$ et $\sigma$ sont des foncteurs de localisation (comme il est évident sur la définition).

Tout foncteur ou foncteur composé du diagramme page 8, allant d'une catégorie interne au diagramme vers une catégorie bordante, induit un foncteur pour les localisées pour $W_S^{\mathrm{d}}$ (quand la catégorie but fait partie des quatre gauches,
\msp{page 72}
rappelées dans $(*)$, p. 70), soit pour $W_S^{\mathrm{g}}$ (quand ces catégories font partie des quatre du côté droit, rappelées dans $(*')$), lesquels localiseurs coïncident aussi avec $W_S^{\mathrm{f}}$ sauf dans le seul cas de la catégorie $\Cat/S$. Tous les foncteurs ainsi obtenus se réduisent, à équivalence près, à deux, suivant que la catégorie localisée but est $\HOT_W(S)$ (cas de la localisation $W_S^{\mathrm{d}}$), ou qu'elle est $\HOT^o_W(S)$ (cas de la localisation $W_S^{\mathrm{g}}$). On trouve ainsi des foncteurs (en sens inverse de $\rho$ et de $\sigma$)
\printedpage{45}
% head: A. Grothendieck : Dérivateurs, Ch. VI, Édition des Universités de Montpellier II et Paris VII
% foot: Version 9 novembre 2001
\[
\begin{tikzcd}
 & \HOT^{\mathrm{lc}}_W(S) \arrow[dl,"i"'] \arrow[dr,"j"] & \\
\HOT_W(S) & & \HOT^o_W(S).
\end{tikzcd}
\]

Les foncteurs $\rho$, $\sigma$ s'interchangent par dualité $S\longmapsto S^o$, et de même pour $i$, $j$.

\ovale{IV}\quad Les foncteurs $\rho$, $i$ sont adjoints l'un de l'autre dans cet ordre, i.e. on a pour deux objets $\xi\in\Ob\HOT_W(S)$, $\eta\in\Ob\HOT^{\mathrm{lc}}_W(S)$
\[ \Hom(\rho(\xi),\eta)\simeq\Hom(\xi,i(\eta)) \]
(isomorphisme de bifoncteurs). Et dualement les foncteurs $\sigma$, $j$ sont adjoint l'un de l'autre % [sic] adjoint -> adjoints, accord du pluriel
dans cet ordre : pour $\xi\in\Ob\HOT^o_W(S)$, $\eta\in\Ob\HOT^{\mathrm{lc}}_W(S)$
\[ \Hom(\xi,j(\eta))\simeq\Hom(\sigma(\xi),\eta). \]

\msp{page 73}
Il s'ensuit notamment que \textit{les foncteurs $i$, $j$ sont pleinement fidèles.}

\cercle{V}\quad On a des foncteurs évidents
\[
\begin{aligned}
&(*) &\qquad& \HOT_W(S) &\quad& \longrightarrow &\quad& \underline{\mathrm{Hom}}(S^o,\Hot_W)\\
&(*') && \HOT^o_W(S) && \longrightarrow && \underline{\mathrm{Hom}}(S,\Hot_W).
\end{aligned}
\]
Le premier s'interprète sur les quatre catégories $(*)$ (bordantes gauches) de page 70, qui servent indifférement à décrire $\HOT_W(S)$, par % [sic] indifférement -> indifféremment, faute d'orthographe
\[ s\longmapsto\hot_W(X_s) \]
dans le cas des trois catégories distinctes de $\Cat/S$, et par
\[ s\longmapsto\hot_W(s\backslash X) \]
dans le cas de $\Cat/S$. La dépendence fonctorielle par rapport à $s$ est évidente dans le cas % [sic] dépendence -> dépendance, faute d'orthographe
de $\Cat/S$ (donnée par $s\longmapsto\hot_W(s\backslash X)$), et dans le cas de $\ul{Fibsc}\,S$, $\ul{Fib}\,S$, où on dispose, pour $u : s\longrightarrow t$ dans $S$, du changement de base
\[ u^* : X_t\longrightarrow X_s. \]
Dans le cas de $\ul{Liss}\,S$, la dépendence fonctorielle se déduit de celle de $\hot_W(s\backslash X)$, grâce % [sic] dépendence -> dépendance, faute d'orthographe
au fait que l'inclusion
\[ X_s\longrightarrow s\backslash X \]
est alors une $W$-équivalence.

On a une interprétation duale de l'application
\printedpage{46}
% head: A. Grothendieck : Dérivateurs, Ch. VI, Édition des Universités de Montpellier II et Paris VII
% foot: Version 9 novembre 2001
\msp{page 74}
$(*)'$, \texttt{[mot illisible]}
\[ s\longmapsto\hot_W(X_s)\qquad\text{ou}\qquad s\longmapsto\hot_W(s\backslash X), \]
avec la loi fonctorielle en $s$ déduite du foncteur $t\backslash X\longrightarrow s\backslash X$ \texttt{[plutôt} $X/s\longrightarrow X/t$\texttt{]} (pour $u : s\longrightarrow t$) dans le cas du $\Cat/S$, du foncteur cochangement de base dans les cas de $\ul{Cofib}\,S$ et $\ul{Cofibsc}\,S$, et de la variante pour $\Cat/S$ et de la $W$-équivalence
\[ X_s\longrightarrow X/s \]
dans le cas de $\ul{Prop}\,S$.

On trouve alors ceci : pour qu'un objet $\xi$ de $\HOT_W(S)$ (resp. $\HOT^o_W(S)$) soit dans l'image essentielle de $\HOT^{\mathrm{lc}}_W(S)$, il faut et il suffit que le foncteur associé $S^o\longrightarrow\Hot_W$ (resp. $S\longrightarrow\Hot_W$) soit un système local (\footnote{il faudra encore écrire la vérification}), i.e. transforme toute flèche de $S$ en un isomorphisme, i.e. se factorise par $(\Pi S)^o\longrightarrow\Hot_W$ resp. $\Pi S\longrightarrow\Hot_W$ (\footnote{où $\Pi S$ est le groupoïde fondamental de $S$}). Donc pour un $X$ dans $\Cat/S$, cela signifie soit que les foncteurs
\[ t\backslash X\longrightarrow s\backslash X \]
associés à des $u : s\longrightarrow t$ dans $S$ soient tous des $W$-équivalences (cas de $\HOT_W(S)$), soit que les foncteurs
\[ X/s\longrightarrow X/t \]
le soient (cas de $\HOT^o_W(S)$). (NB. ces conditions ne sont nullement équivalentes. Mais si $X$ \texttt{[est]} lisse \texttt{[sur]} $S$,
\msp{page 75}
la première implique la seconde, et l'inverse est vrai si $X$ est propre sur $S$, et dans ce cas, la condition équivaut à ce que $X$ soit dans $\ul{Fibhot}$. Enfin, si $X$ est déjà dans $\ul{Fibhot}\,S$, alors les deux conditions sont automatiquement satisfaites (\footnote{On aurait dû dans le diagramme introduire aussi $\ul{Liss}_0S$, $\ul{Prop}_0S$ !}).
\[
\begin{tikzcd}[column sep=small, row sep=small]
\ul{Fibsc}_0S \arrow[dd,hook] \arrow[drr] & & & & & & & & \ul{Cofibsc}_0S \arrow[dll] \arrow[dd,hook] \\
 & & \ul{Fib}_0S \arrow[dd,hook] \arrow[dr,hook] & & \ul{Parf}\,S \arrow[dl,hook] \arrow[dd,hook] \arrow[dr,hook] & & \ul{Cofib}_0S \arrow[dl,hook] \arrow[dd,hook] & & \\
\ul{Fibsc}\,S \arrow[drr] & & & \ul{Liss}_0S \arrow[dr,hook] \arrow[dd,hook] & & \ul{Prop}_0S \arrow[dl,hook] \arrow[dd,hook] & & & \ul{Cofibsc}\,S \arrow[dll] \\
 & & \ul{Fib}\,S \arrow[dr,hook] & & \ul{Fibhot}\,S \arrow[dd,hook] \arrow[uuullll,dotted,bend right,"\Phi_0"'] \arrow[uuurrrr,dotted,bend left,"\Psi_0"'] & & \ul{Cofib}\,S \arrow[dl,hook] & & \\
 & & & \ul{Liss}\,S \arrow[dr,hook] & & \ul{Prop}\,S \arrow[dl,hook] & & & \\
 & & & & \Cat/S \arrow[uuullll,dotted,bend left,"\Phi"] \arrow[uuurrrr,dotted,bend right,"\Psi"'] & & & &
\end{tikzcd}
\]
\printedpage{47}
% head: A. Grothendieck : Dérivateurs, Ch. VI, Édition des Universités de Montpellier II et Paris VII
% foot: Version 9 novembre 2001
(\footnote{Le diagramme de \textit{quinze} catégories remarquables.}). De plus, si deux éléments $\xi\in\Ob\HOT_W(S)$, $\xi'\in\Ob\HOT^o_W(S)$ proviennent d'un même $\eta\in\Ob\HOT^{\mathrm{lc}}_W(S)$, alors les deux foncteurs correspondants
\[ (\Pi S)^o\xrightarrow{\ \varphi\ }\Hot_W\xleftarrow{\ \varphi'\ }\Pi S \]
sont \textit{associés,} i.e. chacun est isomorphe à `l'inverse' de l'autre.
\[
\left\{
\begin{array}{lcl}
\varphi^o(s) & = & \varphi(s)\\
\varphi^o(u) & = & \varphi(u)^{-1}.
\end{array}
\right.
\]

\dgrblock{dgrb201}{Completing the test for local constancy}
\begin{dcomment}{exposition}{p74:On trouve alors ceci : pour qu'un objet}{Completing the test for local constancy}
The fibre diagram is being used here to recognize objects in an already constructed full subcategory. This is weaker than claiming that an arbitrary diagram in $\Hot_W$ reconstructs the coherent object, an issue treated in the commentary at manuscript p.~7. The recognition statement nevertheless needs the verification requested in footnote~47.

Represent $\xi$ by $X\to S$, and replace it by the split fibration $\Phi X\to S$. Its strict fibre at $s$ is $s\backslash X$, and the transition along $u:s\to t$ is precomposition
\[
 u^*:t\backslash X\longrightarrow s\backslash X.
\]
The replacement map is right-fine, so $X$ and $\Phi X$ represent the same object of $\HOT_W(S)$. If all the $u^*$ belong to $W$, this split fibration belongs to $\underline{\mathrm{Fib}}_0S$: it is a categorical fibration with weakly constant fibres. By the comparison already proved, it belongs to $\underline{\mathrm{Fibhot}}S$. Its class is therefore in the image of $i$. This proves sufficiency without selecting transport maps in an arbitrary nonsplit representative.

For necessity, choose a homotopy-fibration representative $Y\to S$ of an object in the image. Its comma-fibre transition maps are weak equivalences by the defining local-constancy property. The diagram functor sends an isomorphism in $\HOT_W(S)$ to an isomorphism of diagrams in $\Hot_W$. Whether every structural arrow of such a diagram is invertible is invariant under an isomorphism of diagrams. Hence the property holds for every representative of the same class. The covariant argument uses $\Psi X$ and the transitions $X/s\to X/t$.

To pass from ``invertible in $\Hot_W$'' to ``belongs to $W$'', one uses strong saturation. It holds for $W_\infty$ and, in the standard theory, for every basic localizer; see \citet[Proposition~4.2.4]{Cisinski2006Presheaves}. This makes explicit a logical input that the sentence ``cela signifie'' suppresses. Without saturation, the invariant statement uses the saturation of $W$, rather than arbitrary membership in the originally chosen class.

Finally, the factorization through $\Pi S$ in the text concerns only the underlying homotopy-category-valued diagram. Its existence follows from the ordinary universal property of groupoid completion. It does not replace the coherent local system by a representation of the fundamental groupoid. For space-valued coefficients, higher homotopies of the base can still act on the fibres. This leaves the argument above intact: invertibility of all arrows detects local constancy, while the object being detected still retains its coherence.
\end{dcomment}


\ovale{VI}\quad NB Dans la nouvelle version du diagramme, où j'ai rajouté des catégories $\ul{Liss}_0S$ et $\ul{Prop}_0S$ pour rendre le tracé plus transparent, les catégories appelés `bordantes' deviennent % [sic] appelés -> appelées, accord du participe
le `rez-de-chausseée', et les catégories `internes', le `premier étage' ou `étage'. Tous % [sic] rez-de-chausseée -> rez-de-chaussée, lettre en trop
les 7 carrés du diagramme sont catésiens. Toutes les flèches sont des inclusions pleinement % [sic] catésiens -> cartésiens, lettre omise
fidèles, sauf les quatre flèches horizontales extérieures des catégories de Cat-fibrés ou de Cat-cofibrés \textit{scindés}, avec morphismes exclusivement cartésiens ou cocartésiens et même compatibles avec les scindages, vers les catégories correspondantes de structures non scindées.

\dgrblock{dgrb202}{What the fifteen categories remember}
\begin{dcomment}{exposition}{p75:les 7 carrés du diagramme sont catésiens.}{What the fifteen categories remember}
The expanded diagram separates two choices which are often bundled together in contemporary notation: a property of the projection $X\to S$, and a restriction on the morphisms allowed between such projections. This separation is responsible for both its size and its usefulness.

Except for the categories of split fibrations and split opfibrations, the displayed categories are full subcategories of $\Cat/S$. A map between two fibred objects is therefore an arbitrary functor over $S$; it need not preserve Cartesian arrows. For the split categories, Grothendieck imposes preservation of the splitting, and hence preservation of the chosen Cartesian or cocartesian arrows. Their forgetful functors are not full. The commentary at manuscript p.~30 explains why the naive unit for a splitting construction need not lie in this narrower morphism category.

The seven Cartesian squares assert that the relevant conditions on objects and morphisms are simultaneous conditions. In particular, the square with $\underline{\mathrm{Parf}}S$ expresses the conjunction of the left and right properties, while the outer squares impose the locally constant condition on a split or nonsplit fibration. These assertions should be read in the ordinary category-of-categories sense in which the given subcategories and forgetful functors have been specified. They are not additional assertions that the corresponding nerves form homotopy pullback squares.

The explicit comma replacements remove the possible difficulty caused by the narrow morphism categories. Applied to a functor over $S$, the rule $s\mapsto s\backslash X$ gives a strictly natural transformation of strict $\Cat$-valued diagrams, and therefore a splitting-preserving map of their categories of elements. This supplies a functor even when the original map did not preserve Cartesian arrows. The equivalence after localization follows from the natural comparison zigzags established earlier, not from pretending that all the forgetful functors were already fully faithful. A later systematic treatment of such comma and category-of-elements comparisons is \citet[Theorem~3.8]{Maltsiniotis2005Asphericite}.
\end{dcomment}


\msp{page 76}
\[
\begin{tikzcd}[column sep=small, row sep=small]
 & & & & \ul{Fibhot}\,S \arrow[dd,hook,"\delta"] \arrow[dddllll,dotted,bend right,"\Phi_0"'] \arrow[dddrrrr,dotted,bend left,"\Psi_0"] & & & & \\
 & & & \ul{Liss}_0S \arrow[ur,hook] \arrow[dd,hook] & & \ul{Prop}_0S \arrow[ul,hook] \arrow[dd,hook] & & & \\
 & \ul{Fib}_0S \arrow[urr,hook] \arrow[dd,hook] \arrow[ddrrr,dotted,bend left] & & & \Cat/S \arrow[dd,"C"] \arrow[dddllll,dotted,bend right,"\Phi"'] \arrow[dddrrrr,dotted,bend left,"\Psi"] & & & \ul{Cofib}_0S \arrow[ull,hook] \arrow[dd,hook] \arrow[ddlll,dotted,bend right] & \\ % Reading verified against PDF leaf 9: sources Fib_0 S and Cofib_0 S, target middle Parf S.
\ul{Fibsc}_0S \arrow[ur] \arrow[dd,hook] & & & \ul{Liss}\,S \arrow[ur,hook] & & \ul{Prop}\,S \arrow[ul,hook] & & & \ul{Cofibsc}_0S \arrow[ul] \arrow[dd,hook] \\
 & \ul{Fib}\,S \arrow[urr,hook] & & & \ul{Parf}\,S \arrow[uuul] \arrow[uuur] \arrow[dd,dotted,"\mathrm{id}"] & & & \ul{Cofib}\,S \arrow[ull,hook] & \\
\ul{Fibsc}\,S \arrow[ur] & & & & & & & & \ul{Cofibsc}\,S \arrow[ul] \\
 & & & & \ul{Parf}\,S \arrow[uuul,dotted,hook] \arrow[uuur,dotted,hook] & & & &
\end{tikzcd}
\]

\dgrblock{dgrb203}{Maps over a group retain the action}
\begin{dcomment}{examples}{p75}{Maps over a group retain the action}
The common locally constant category still remembers the map to the base. A completely discrete example makes this visible without any higher coherence. Let $G=C_2=\{1,\tau\}$, let $S=BG$, and let $E=G$ with its regular right action. Its category of elements has two objects, with one arrow between any ordered pair; the arrow from $e$ to $e'$ lies over the unique $g$ with $e'=eg$. Its total category is equivalent to $e$, while its projection to $BG$ records the free transitive action.

The associated fibre diagram has one value, the discrete two-point space $E$, with the nontrivial element of $G$ acting by transposition. An endomorphism over the base is an equivariant map $E\to E$. Such a map is determined by the image of $1$, and either choice gives a map: $e\mapsto ae$ for $a\in G$. Thus there are exactly two endomorphisms. Since these diagrams have discrete values, a coherent natural transformation has no nontrivial higher homotopies; their derived mapping space is the same discrete set of two equivariant maps. Consequently the endomorphism set is also two-element in each of the equivalent coarse and fine theories over $BG$.

Forgetting the base would instead replace the total category by a point, whose endomorphism set in $\Ho(\mathrm{Spaces})$ is a singleton. The coarse theory over $BG$ therefore cannot be obtained by simply computing maps between the total homotopy types. It localizes maps \emph{over} $BG$, and the permitted zigzags continue to remember that structure.

The same computation can be made with spaces. The object is the universal covering $EG\to BG$, with contractible total space and discrete fibre $G$. Its derived space of self-maps over $BG$ is $\Omega BG\simeq G$. The ordinary space of self-maps of its contractible total space is contractible. This is the relative distinction behind the Hom formula at manuscript p.~81. It also explains why a groupoid base is not a degenerate case of the theory, even though the fine-to-coarse comparison is an equivalence for such a base.
\end{dcomment}


J'ai tracé en pointillés verts \texttt{[couleur dans l'original seulement]} les cinq foncteurs $\Phi$, $\Psi$, $\Phi_0$, $\Psi_0$, $C$ qui par passage aux localisés donnent des quasi-inverses pour toutes les équivalences de catégories énoncées dans \cercle{I} \cercle{II} du présent Scholie. De façon précise, le foncteur
\[ \Phi : \Cat/S\longrightarrow\ul{Fibsc}\,S \]
associe à toute $X$ dans $\Cat/S$ la catégorie scindée sur $S$
\[ s\longmapsto s\backslash X \]
\printedpage{48}
% head: A. Grothendieck : Dérivateurs, Ch. VI, Édition des Universités de Montpellier II et Paris VII
% foot: Version 9 novembre 2001
sur $S$. Ce foncteur, et ses restrictions à $\ul{Liss}\,S$, $\ul{Fib}\,S$ et, plus généralement, à toute sous-catégorie pleine $\mathcal{F}$ de $\Cat/S$ contenant $\ul{Fib}\,S$, est compatible avec les localiseurs $W_S$ d'une part (sur $\mathcal{F}$ et $\ul{Fibsc}\,S$), $W_S^{\mathrm{d}}$ d'autre part, et passe donc aux localisés correspondants, pour fournir le quasi-inverse cherché de
\[ (\ul{Fibsc}\,S)\Sigma^{-1}\longrightarrow\mathcal{F}\Sigma^{-1}\qquad(\Sigma=W_S\text{ ou }W_S^{\mathrm{d}}) \]

On peut aussi, pour toute sous-catégorie pleine $\mathcal{F}'$ intermédiaire
\[ \ul{Fib}\,S\subseteq\mathcal{F}'\subseteq\mathcal{F}\subseteq\Cat/S, \]
regarder $\Phi$ comme un foncteur $\mathcal{F}\longrightarrow\mathcal{F}'$ (plus précisément, prendre le composé
\[ \mathcal{F}\xrightarrow{\Phi|\mathcal{F}}\ul{Fibsc}\,S\longrightarrow\ul{Fib}\,S\hookrightarrow\ \mathcal{F}'\ ), \]
soit $\Phi_{\mathcal{F}',\mathcal{F}}$, et de ce qui vient d'être dit résulte que les deux foncteurs
\[
\begin{tikzcd}
\mathcal{F}' \arrow[r,hook] & \mathcal{F} \arrow[l,bend left,"\Phi_{\mathcal{F}',\mathcal{F}}"]
\end{tikzcd}
\]
\msp{page 77}
établissent des équivalences quasi-inverses l'une de l'autre
\[
\begin{tikzcd}
\mathcal{F}'\Sigma^{-1} \arrow[r,"\approx"] & \mathcal{F}\Sigma^{-1} \arrow[l,bend left,"\approx"',"\Phi_{\mathcal{F}',\mathcal{F}}"]
\end{tikzcd}
\qquad\qquad(\Sigma=W_S\text{ ou }W_S^{\mathrm{d}})
\]

Énoncé dual pour le foncteur
\[ \Psi : \Cat/S\longrightarrow\ul{Cofibsc}\,S, \]
associant à $X$ dans $\Cat/S$ la catégorie cofibrée scindée sur $S$
\[ s\longmapsto X/s. \]
Il donne des couples de foncteurs quasi-inverses
\[
\begin{tikzcd}
(\ul{Cofibsc}\,S)\Sigma^{-1} \arrow[r] & \mathcal{F}\Sigma^{-1} \arrow[l,bend left,"\Psi_{\mathcal{F}}"]
\end{tikzcd}
\qquad\qquad(\Sigma=W_S\text{ ou }W_S^{\mathrm{g}})
\]
pour tout sous-catégorie pleine $\mathcal{F}$ de $\Cat/S$ telle que % [sic] tout -> toute, accord avec sous-catégorie
\[ \ul{Cofib}\,S\subseteq\mathcal{F}\subseteq\Cat/S, \]
et des couples de foncteurs quasi-inverses

% ---- from Der-ChVI-345_p051-055.tex
\printedpage{49}
% head: A. Grothendieck : Dérivateurs, Ch. VI, Édition des Universités de Montpellier II et Paris VII
% foot: Version 9 novembre 2001

\[
\begin{tikzcd}
\mathcal{F}'\Sigma^{-1} \arrow[r, "\approx"] & \mathcal{F}\Sigma^{-1} \arrow[l, bend left=40, "\approx"', "\Psi_{\mathcal{F}',\mathcal{F}}"]
\end{tikzcd}
\qquad\qquad (\Sigma = W_S \text{ ou } W_S^{\mathrm{g}})
\]

pour deux sous-catégories pleines $\mathcal{F}'$, $\mathcal{F}$
\[
\usl{Cofib}\,S \subseteq \mathcal{F}' \subseteq \mathcal{F} \subseteq \mathrm{Cat}/S.
\]
Les foncteurs $\Phi$, $\Psi$ envoient la sous-catégorie pleine $\usl{Fibhot}\,S$ de ‘l’étage’, dans les sous-catégories $\usl{Fibsc}_0S$, $\usl{Cofibsc}_0S$, respectivement, donc induisent à l’étage des foncteurs $\Phi_0$, $\Psi_0$
\[
\usl{Fibsc}_0S \xleftarrow{\Phi_0} \usl{Fibhot}\,S \xrightarrow{\Psi_0} \usl{Cofibsc}_0S,
\]
jouant à l’étage exactement le même rôle que $\Phi$, $\Psi$ au rez-de-chaussée. La seule différence, c’est que à l’étage on a $W_S = W_S^{\mathrm{d}} = W_S^{\mathrm{g}}$ $(= W_S^{\mathrm{f}})$,
\msp{page 78}
donc il n’y a qu’une suite de catégories localisées à examiner. On trouve ainsi
\[
\begin{tikzcd}
(\usl{Fibsc}_0S)W_S^{-1} \arrow[r, "\approx"] & \mathcal{F}W_S^{-1} \arrow[l, bend left=40, "\approx"', "\Phi_{0,\mathcal{F}}"]
\end{tikzcd}
\]
\[
\begin{tikzcd}
\mathcal{F}'W_S^{-1} \arrow[r, "\approx"] & \mathcal{F}W_S^{-1} \arrow[l, bend left=40, "\approx"', "\Phi_{0,\mathcal{F}',\mathcal{F}}"]
\end{tikzcd}
\]
pour $\mathcal{F}$ (ou $\mathcal{F}$, $\mathcal{F}'$) sous-catégorie(s) pleine(s) de $\usl{Fibhot}\,S$ telle(s) que
\[
\usl{Fib}_0S \subseteq \mathcal{F} \subseteq \usl{Fibhot}\,S \qquad\text{resp.}\qquad \usl{Fib}_0S \subseteq \mathcal{F}' \subseteq \mathcal{F} \subseteq \usl{Fibhot}\,S,
\]
et de même du côté droit pour les sous-catégories pleines $\mathcal{F}$ ou $\mathcal{F}$, $\mathcal{F}'$ de $\usl{Fibhot}\,S$ satisfaisant
\[
\usl{Cofib}_0S \subseteq \mathcal{F} \subseteq \usl{Fibhot}\,S \qquad\text{resp.}\qquad \usl{Cofib}_0S \subseteq \mathcal{F}' \subseteq \mathcal{F} \subseteq \usl{Fibhot}\,S,
\]
NB. La catégorie $\usl{Parf}\,S$ contient à la fois $\usl{Fib}_0S$ et $\usl{Cofib}_0S$, et fait donc partie des catégories admissibles tout comme $\mathcal{F}$, ou [?] $\mathcal{F}'$.

Ainsi, les quatre foncteurs $\Phi$, $\Psi$, $\Phi_0$, $\Psi_0$ donnent tous les quasi-inverses qu’on peut souhaiter au rez-de-chaussée d’une part ($\Phi$ et $\Psi$), à l’étage de l’autre ($\Phi_0$, $\Psi_0$). Ils ramènent donc toutes les localisées pertinentes aux quatre suivantes
\[
\begin{gathered}
(\mathrm{Cat}/S)W_S^{-1},\ \ (\mathrm{Cat}/S)(W_S^{\mathrm{d}})^{-1},\ \ (\mathrm{Cat}/S)(W_S^{\mathrm{g}})^{-1} \qquad \text{(au rez-de-chaussée)}\\
(\usl{Fibhot}\,S)W_S^{-1} \qquad \text{(à l’étage).}
\end{gathered}
\]

\encadre{VII} Il reste à comparer la localisée à l’étage avec les localisées au rez-de-chaussée. Les flèches verticales descendantes de l’étage au rez-de-chaussée, sont compatibles avec $W_S^{\mathrm{d}}$ du côté gauche du diagramme, avec $W_S^{\mathrm{g}}$ du côté droit, donc elles induisent essentiellement \textit{deux} foncteurs différents
\printedpage{50}
% head: A. Grothendieck : Dérivateurs, Ch. VI, Édition des Universités de Montpellier II et Paris VII
% foot: Version 9 novembre 2001
\msp{page 79}
$i$, $j$ sur $(\usl{Fibhot}\,S)W_S^{-1}$, dans $(\mathrm{Cat}/S)(W_S^{\mathrm{d}})^{-1}$ et $(\mathrm{Cat}/S)(W_S^{\mathrm{g}})^{-1}$ :
\[
\begin{tikzcd}[column sep=small, row sep=large]
& (\usl{Fibhot}\,S)W_S^{-1} \arrow[dl, "i"'] \arrow[dr, "j"] \arrow[dd, dotted, "\approx"', "\delta"] & \\
\mathrm{HOT}_W(S) \overset{\text{déf}}{\approx} (\mathrm{Cat}/S)(W_S^{\mathrm{d}})^{-1} \arrow[dr, dotted, "\rho"'] & & (\mathrm{Cat}/S)(W_S^{\mathrm{g}})^{-1} \overset{\text{déf}}{\approx} \mathrm{HOT}^{o}_{W}(S) \arrow[dl, dotted, "\sigma"] \\
& \begin{array}{l} (\mathrm{Cat}/S)W_S^{-1} \\ \overset{\text{déf}}{\approx} \mathrm{HOT}^{\mathrm{lc}}_{W}(S) \end{array} &
\end{tikzcd}
\qquad ,
\]
de façon à faire commuter le carré complété par les pointillés, représentant les foncteurs de localisation $\rho$, $\sigma$. On a bien sûr
\[
\begin{aligned}
(*) \qquad \rho i = \sigma j : (\usl{Fibhot}\,S)W_S^{-1} \longrightarrow (\mathrm{Cat}/S)W_S^{-1}\\
\text{induit par l’inclusion}\\
\delta : \usl{Fibhot}\,S \hookrightarrow \mathrm{Cat}/S
\end{aligned}
\]
Nous avons énoncé dans \encadre{I} que cette flèche est une équivalence, de sorte que $\rho$ et $\sigma$ peuvent être interprétées, à isomorphisme de foncteurs près, comme des rétractions de $\mathrm{HOT}_W(S)$ et de $\mathrm{HOT}^{o}_{W}(S)$ sur une sous-catégorie commune, s’interprétant soit comme $(\usl{Fibhot}\,S)W_S^{-1}$, soit comme $(\mathrm{Cat}/S)W_S^{-1}$. Il reste à décrire le foncteur quasi-inverse de $(*)$, ce qui donnera en même temps les quasi-inverses de tous les foncteurs verticaux (de l’étage vers le rez-de-chaussée) dans le diagramme p. 76, plus exactement pour leurs localisés, pour $W_S$.

\dgrblock{dgrb204}{A finite comma diagram with an infinite reflection}
\begin{dcomment}{examples}{p78}{A finite comma diagram with an infinite reflection}
The comma replacement and the path replacement can have very different fibres. To see the distinction explicitly, let $S$ have two objects $a,b$ and exactly two nonidentity arrows $u,v:a\rightrightarrows b$. Its nerve realizes to a circle. Take $X=e$ over $b$.

The right comma diagram $F(s)=s\backslash X$ is discrete, with
\[
 F(b)=\{*\},\qquad F(a)=\{u,v\},\qquad
 F(u)(*)=u,\quad F(v)(*)=v.
\]
Its category of elements has three objects, two over $a$ and one over $b$, with one arrow from each object over $a$ to the object over $b$. The latter is terminal, so the total category is contractible. This split fibration is right-fine equivalent to $X$, but it is not locally constant: neither map from a singleton to the two-point set is an equivalence. Dually, $X/s$ is empty at $a$ and a singleton at $b$, since there is no arrow $b\to a$.

We can exhibit a locally constant replacement of $F$ without writing the entire category of paths. Put
\[
 L(a)=L(b)=\mathbb Z,\qquad L(u)(n)=n,
 \qquad L(v)(n)=n+1,
\]
viewing $L$ as a presheaf of discrete spaces. Define $F\to L$ by sending $*$ to $0$, $u$ to $0$, and $v$ to $1$. The two naturality squares commute. All transition maps of $L$ are bijections. Its category of elements has objects $a_n,b_n$ and arrows
\[
 a_n\longrightarrow b_n,\qquad a_{n+1}\longrightarrow b_n
 \qquad(n\in\mathbb Z).
\]
There are no nonidentity composable pairs. Its nerve is the infinite line subdivided into edges, hence is contractible. The functor from the three-object category of elements of $F$ to this infinite category is therefore a coarse weak equivalence. Both its source and target are nonempty and contractible, which verifies the assertion directly rather than by a fibrewise test.

Since $L$ is locally constant, the established reflection identifies its class with $i\rho(F)$. The fibre of this reflection is $\mathbb Z$, whereas the original comma fibres had at most two points. Transport around the generating loop of the circle is translation by one, up to the choice of orientation. Geometrically, the total locally constant object is the universal cover of the circle. The reflection must add all iterates of transport; it cannot be computed by merely declaring the two original transition maps invertible within their finite sets. This example gives a concrete role to the path construction in a place where the comparison diagram alone can conceal it.
\end{dcomment}


\msp{page 80}
C’est ce qui est obtenu à l’aide du foncteur ‘montant’ (du rez-de-chaussée vers l’étage) $C$, qu’on peut regarder comme un foncteur $\mathrm{Cat}/S \longrightarrow \usl{Fibhot}\,S$, mais qui en fait prend ses valeurs dans $\usl{Parf}\,S$. (C’est là un détail sans grande importance, me semble-t-il.)
\[
\mathrm{Cat}/S \xrightarrow{C} \usl{Parf}\,S \hookrightarrow \usl{Fibhot}\,S.
\]
Ce foncteur est compatible avec les localiseurs $W_S$ des deux côtés, et établit des équivalences (quasi-inverses des équivalences
\[
\begin{gathered}
(\usl{Parf}\,S)W_S^{-1} \xrightarrow{\sim} (\mathrm{Cat}/S)W_S^{-1} \qquad\text{et}\qquad (\usl{Fibhot}\,S)W_S^{-1} \xrightarrow{\sim} (\mathrm{Cat}/S)W_S^{-1}),\\
(\mathrm{Cat}/S)W_S^{-1} \xrightarrow{\sim} (\usl{Parf}\,S)W_S^{-1} \left( \xrightarrow{\sim} (\usl{Fibhot}\,S)W_S^{-1} \right).
\end{gathered}
\]

\printedpage{51}
% head: A. Grothendieck : Dérivateurs, Ch. VI, Édition des Universités de Montpellier II et Paris VII
% foot: Version 9 novembre 2001
Donc, quand les trois catégories HOT pour $S$ sont interprétées en termes de localisées de $\mathrm{Cat}/S$, c’est le foncteur $C$ (plus exactement le composé $\mathrm{Cat}/S \xrightarrow{C} \usl{Parf}\,S \hookrightarrow \mathrm{Cat}/S$) qui fournit à la fois le foncteur d’inclusion (!) $i$, et le foncteur d’inclusion $j$ dans $\mathrm{HOT}_W(S)$ et $\mathrm{HOT}^{o}_{W}(S)$ respectivement. Le foncteur en question
\[
C' : \mathrm{Cat}/S \longrightarrow \mathrm{Cat}/S
\]
satisfait bien
\[
C'(W_S) \subseteq W_S^{\mathrm{d}} \cap W_S^{\mathrm{g}}
\]
et définit les deux foncteurs
\[
\begin{tikzcd}[column sep=large]
& (\mathrm{Cat}/S)(W_S^{\mathrm{d}})^{-1} \approx \mathrm{HOT}_W(S) \\
\begin{array}{l} (\mathrm{Cat}/S)W_S^{-1} \\ \approx \mathrm{HOT}^{\mathrm{lc}}_{W}(S) \end{array} \arrow[ur, "i"] \arrow[dr, "j"'] & \\
& (\mathrm{Cat}/S)(W_S^{\mathrm{g}})^{-1} \approx \mathrm{HOT}^{o}_{W}(S)
\end{tikzcd}
\]
Cela nous permet d’expliciter les formules d’adjonction de \encadre{IV}, quand on y interprète $\xi$, $\eta$
\msp{page 81}
comme provenant de $X$, $Y$ dans $\mathrm{Cat}/S$ :
\[
\boxed{\mathrm{Hom}_{W_S^{\mathrm{d}}}(X,CY) \simeq \mathrm{Hom}_{W_S^{\mathrm{g}}}(X,CY) \simeq \mathrm{Hom}_{W_S}(X,Y),}
\]
où $\mathrm{Hom}_{W_S^{\mathrm{d}}}$, $\mathrm{Hom}_{W_S^{\mathrm{g}}}$ et $\mathrm{Hom}_{W_S}$ désignent respectivement les Hom dans les catégories localisées pour $W_S^{\mathrm{d}}$, $W_S^{\mathrm{g}}$ et $W_S$, i.e. dans $\mathrm{HOT}_W(S)$, $\mathrm{HOT}^{o}_{W}(S)$ et $\mathrm{HOT}^{\mathrm{lc}}_{W}(S)$ respectivement (\footnote[51]{Cette formule suggère que si $X \in \mathrm{Ob}\,\mathrm{Cat}/S$, $Y \in \mathrm{Ob}\,\usl{Fibhot}\,S$, alors
\[
\mathrm{Hom}_{W_S^{\mathrm{d}}}(X,Y) \simeq \mathrm{Hom}_{W_S^{\mathrm{g}}}(X,Y) \simeq \mathrm{Hom}_{W_S}(X,Y),
\]
mais ça résulte de la formule d’adjoinction et de $Y \longrightarrow C(Y)$ dans $W_S^{\mathrm{d}} \cap W_S^{\mathrm{g}}$ pour $Y$ dans $\usl{Fibhot}\,S$.}). % [sic] adjoinction -> adjonction, spelling (elsewhere "adjonction")

\textbf{Commentaire.} Toutes les constructions et les énoncés de \encadre{I} à \encadre{VI}, à l’exception du fait que les foncteurs descendants induisent des équivalences pour les localisés \textit{grossiers}, et des formules d’adjonction de \encadre{IV}, sont de nature quasiment soritale, et n’exigent d’autres conditions sur le localiseur fondamental $W \subseteq \mathrm{Fl}\,\mathrm{Cat}$ que L1, L3 quater, L4 bis, récapitulées comme a) b) c) de la page 45. (On n’a même pas à utiliser d).) Il est vrai qu’en l’absence de la condition L5 (page 43), on ignore si à l’étage les deux localiseurs pertinents, $W_S$ et $W_S^{\mathrm{f}}$ (égal à $W_S^{\mathrm{d}}$ du côté gauche, à $W_S^{\mathrm{g}}$ du côté droit) sont égaux. Donc à priori on trouve à l’étage \textit{deux} catégories localisées essentiellement distinctes, et non pas une seule, savoir $(\usl{Fibhot}\,S)\Sigma^{-1}$ avec $\Sigma = W_S$ ou $\Sigma = W_S^{\mathrm{f}}$ $(= W_S^{\mathrm{d}} = W_S^{\mathrm{g}})$.

\printedpage{52}
% head: A. Grothendieck : Dérivateurs, Ch. VI, Édition des Universités de Montpellier II et Paris VII
% foot: Version 9 novembre 2001
Par contre, le fait que les foncteurs de descente induisent des équivalences entre les localisés grossiers est plus profond, il exige (me semble-t-il) qu’on connaisse un foncteur ‘montant’ $C$, disons de $\mathrm{Cat}/S$ vers $\usl{Fibhot}\,S$, compatible avec $W_S$, et muni de
\[
X \xrightarrow{i_X} C(X),
\]
qui soit dans $W_S$ pour $X$ quelconque. Quand on sait que Cat est une catégorie de Quillen, même sans avoir les propriétés des foncteurs $\usl{Ch}_\infty$ ou $\usl{Ch}$ (\footnote[52]{ni même être sûr de l’existence d’une factorisation \textit{fonctorielle} de $X \xrightarrow{f} S$ en $X \xrightarrow{i_X} C(X) \xrightarrow{\bar{f}} S$ avec $i_X \in W$ et $\bar{f}$ parfait.}), on doit sans doute
\msp{page 82}
pouvoir se débrouiller. Mais la formule d’adjonction me semble rester inattendue et profonde, même en admettant les résultats de \textsc{Quillen} et \textsc{Thomason} (\footnote[53]{Voir cependant le théorème scholie 2, p. 108-113. Ce n’est finalement pas si profond qu’il me semblait, mais quasiment sorital!}). Elle parait liée ici à une propriété assez particulière du foncteur $C$, savoir le lemme p. 68
\[
\begin{gathered}
i_{C(X)} \underset{\big\uparrow}{=} C(i_X) : C(X) \longrightarrow C(C(X)),\\
\scriptstyle \text{dans } \underbrace{(\mathrm{Cat}/S)(W_S^{\mathrm{d}})^{-1}}_{\mathrm{HOT}_W(S)} \text{ et dans } \underbrace{(\mathrm{Cat}/S)(W_S^{\mathrm{g}})^{-1}}_{\mathrm{HOT}^{o}_{W}(S)}
\end{gathered}
\]
ce qui n’a pas l’air purement formel. Il serait intéressant d’examiner dans quelle mesure cette formule garde un sens dans une ‘catégorie de modèles’ $(\mathcal{M},W)$ plus ou moins quelconque, ou provenant d’un triple de Quillen $(W,\mathrm{Cof},\mathrm{Fib})$ dans $\mathcal{M}$. Il est clair que la factorisation $X \xrightarrow{f} S$ en $X \xrightarrow{i_X} C(X) \xrightarrow{\bar{f}} S$ $(i_X \in W,\ \bar{f} \in \mathrm{Parf})$, jouant le rôle des ‘fibrations’ dans le contexte Quillen, a un sens intrinsèque indépendant du choix d’une construction particulière, chose qu’on peut préciser de façon parfaite. C’est la description des localiseurs $W_S^{\mathrm{d}}$, $W_S^{\mathrm{g}}$ dans $\mathcal{M}/S$ (pour $S \in \mathrm{Ob}\,\mathcal{M}$) qui n’est pas comprise de façon générale, en dehors du cas $\mathcal{M} = \mathrm{Cat}$. En d’autre termes, je ne comprends pas comment construire des catégories $\mathrm{HOT}_{\mathcal{M},W}(S)$, $\mathrm{HOT}^{o}_{\mathcal{M},W}(S)$, même si $(\mathcal{M},W)$ provient d’un triple de Quillen, je vois seulement la description d’une catégorie $\mathrm{HOT}^{\mathrm{lc}}_{\mathcal{M},W}(S) = (\mathcal{M}/S)(W_S^{\text{grossière}})^{-1}$, et un foncteur de localisation $\mathcal{M}/S \longrightarrow \mathrm{HOT}^{\mathrm{lc}}_{\mathcal{M},W}(S)$. Il faut arriver à décrire dans $\mathcal{M}/S$ des $W$-équivalences (\textit{‘relatives à $S$’}) plus fines que la grossière donnée par $W$ sur la source $X$ sans plus. Dans Cat j’y arrive en utilisant soit la famille des isomorphismes locaux (pour décrire $W_S^{\mathrm{d}}$), soit celle des isomorphismes colocaux (pour décrire $W_S^{\mathrm{g}}$) - en tous les cas, il y a une notion de ‘localisation’ (soit sur $S$, soit sur $S^{o}$ \dots), avec % [sic] En d’autre termes -> En d’autres termes, fixed expression requires plural
\msp{page 83}
toute la richesse des intuitions liées à une telle notion. Peut-être, sans le savoir ai-je toujours traité Cat non comme une catégorie ‘nue’, mais comme un site, ou plutôt (puisqu’il
\printedpage{53}
% head: A. Grothendieck : Dérivateurs, Ch. VI, Édition des Universités de Montpellier II et Paris VII
% foot: Version 9 novembre 2001
y a l’involution $X \longmapsto X^{o}$ qui ne peut être éludée), comme un ‘bi-site’ - on y a \textit{deux} manières de localiser, soit par les $X/x$ (topologie de Zariski!) soit par les $x\backslash X$ (topologie co-Zariskienne). Il y a aussi la topologie lisse (toute proche sans doute de l’intuition de la cohomologie étale), et la topologie propre - et on s’attend que la lisse donne essentiellement les mêmes résultats que la Zariskienne (vue que la topologie étale lui est égale - il n’y a pas lieu d’en faire un plat \dots); et la topologie propre les mêmes résultats (p. ex. la même cohomologie?) que la topologie co-Zariskienne. % [sic] vue que -> vu que, conjunctional locution

\dgrblock{dgrb205}{The target controls the comparison of morphisms}
\begin{dcomment}{exposition}{p81}{The target controls the comparison of morphisms}
The formula on this page is a universal property of the replacement of the \emph{target}. It is stronger than saying that $X$ and $C(X)$ have the same total homotopy type, and it explains why homotopy fibrations are the right objects in which to receive maps.

Write $Q_{\mathrm d},Q_{\mathrm g},Q$ for the three localization functors out of $\Cat/S$. The adjunction $\rho\dashv i$ identifies
\[
 \Hom_{\HOT_W(S)}(Q_{\mathrm d}X,iQY)
       \cong\Hom_{\HOT_W^{\mathrm{lc}}(S)}(QX,QY).
\]
The chosen model for $iQY$ is $Q_{\mathrm d}CY$, and the analogous model for $jQY$ is $Q_{\mathrm g}CY$. This gives both isomorphisms printed in the box, with their naturality in $X$ and $Y$. One must retain $CY$ on the fine sides unless $Y$ is already local.

For $Y\in\underline{\mathrm{Fibhot}}S$, the map $i_Y:Y\to CY$ is both right-fine and left-fine. Postcomposition by this map identifies the corresponding Hom sets with targets $Y$ and $CY$. Combining these identifications with the box proves footnote~51. In particular, the fine and coarse localizations have the same maps \emph{into a local object}, although they can have different maps between arbitrary objects. This is the familiar defining property of a reflective full subcategory, here recovered from explicit relative-category models.

There is also a useful way to recognize the arrows inverted by $\rho$. A map $f:A\to B$ in $\HOT_W(S)$ is inverted if and only if, for every local object $L$, precomposition induces a bijection
\[
 \Hom(B,L)\longrightarrow\Hom(A,L).
\]
One implication follows from the adjunction. Conversely these bijections say that precomposition by $\rho f$ is a bijection against every object of the coarse category; Yoneda implies that $\rho f$ is invertible. The fibre diagram alone does not supply these Hom sets. The regular-action example above shows why even objects with contractible total spaces can have nontrivial relative endomorphisms.
\end{dcomment}


\dgrblock{dgrb206}{Coherence already follows from the two invertibilities}
\begin{dcomment}{exposition}{p82}{Coherence already follows from the two invertibilities}
Grothendieck is trying to decide which part of the reflection depends on the special geometry of paths. Concatenation gives an attractive proof of the equality of the two iterated units, but under the hypotheses already written at manuscript p.~66 that equality is also formal. The distinction matters for the abstraction that follows.

Let $T$ be an endofunctor of an ordinary category and $\eta:1\to T$ a natural transformation. Assume that both $T\eta_X$ and $\eta_{TX}$ are invertible for every $X$. Put $a=T\eta_X$ and $b=\eta_{TX}$. Naturality at $\eta_X$ gives $a\eta_X=b\eta_X$. Applying $T$ and cancelling the invertible map $a$ gives $Ta=Tb$. Next apply naturality to the two maps $a,b:TX\to T^2X$:
\[
 \eta_{T^2X}a=(Ta)b,
 \qquad \eta_{T^2X}b=(Tb)b.
\]
The right sides agree, and $\eta_{T^2X}$ is invertible. Hence $a=b$. Notice that the proof never cancels $\eta_X$, which need not be invertible.

Now work in $\mathcal A\Sigma_0^{-1}$. The functor $C$ induces an endofunctor $T$ because $\Sigma_0\subseteq\Sigma$ and $C(\Sigma)\subseteq\Sigma_0$. Its coaugmentation is the image of $i$. The map $Ti_X$ is invertible because $i_X\in\Sigma$, and $i_{TX}$ is invertible by the second condition at manuscript p.~66. The lemma therefore gives $i_{C(X)}=C(i_X)$ in the fine localization. The first triangle identity is already tautological, and this proves the second. The same argument applies to the left-fine localization.

This is the ordinary-category case of the standard characterization of a localization by an idempotent coaugmented endofunctor; compare \citet[Proposition~5.2.7.4]{Lurie2009HTT}. The displayed proof is included to make the precise cancellation visible. No new geometric property of $C$ is required at this step once both invertibilities are known. Establishing those invertibilities, and identifying the essential image of $C$, remain mathematical work. The source's later Theorem-Scholie~2, at manuscript pp.~108--113, puts exactly these inputs into a general replacement theorem.
\end{dcomment}


Il semble bien que dans toutes les catégories de modèles utilisées jusqu’à présent pour décrire Hot comme une localisée, on dispose de notions de localisation : (Top), $\mathbf{\Delta}^\wedge$ (ou $A^\wedge$, avec $A$ catégorie-test quelconque), (Cat); et dans Cat on dispose même de \textit{deux} telles notions. Je pense par contre aux $\infty$-groupoïdes, c’est peut-être la seule catégorie de modèles à tel point proche déjà de Hot elle-même, qu’on n’y discerne pas de notion de localisation, et que l’idée même de vouloir s’y localiser (si ce n’est en ‘localisant’ Hot lui-même) paraît quasiment
\msp{page 84}
saugrenue.

\dgrblock{dgrb207}{The extra structure hidden in a categorical base}
\begin{dcomment}{context}{p83}{The extra structure hidden in a categorical base}
The word ``bi-site'' is Grothendieck's way of isolating information that a bare presentation of a homotopy theory does not remember. The categories $S/s$ and $s\backslash S$ retain the orientation of arrows in $S$. Localizing the total category alone can destroy that orientation. Consequently a construction which is invariant under weak homotopy equivalence of the base cannot, without additional input, recover both of the directed diagram theories.


The labels also need a small correction in the surrounding prose. A local
isomorphism is a discrete fibration, modelled by a slice projection $S/s\to S$;
these test $W_S^{\mathrm g}$. A co-local isomorphism is a discrete opfibration,
modelled by $s\backslash S\to S$; these test $W_S^{\mathrm d}$. The sentence at
manuscript p.~82 interchanges the two superscripts. The conventions at pp.~83,
94, and 135 agree with the comma definitions above. The edition's wording is
retained, and the correction register records the swap. This is a notational
slip, not a change in the comparison theorem.

There is a useful historical continuity here. In \citet[\S69, pp.~1--3 of the edited extract]{Grothendieck1983}, Grothendieck describes the passage from one derived category to the family of derived categories of diagrams, relates the motivation to the SGA~5 trace formalism, and then formulates the two Kan-extension adjoints. The present passage asks for a converse kind of explanation: which geometric structure on a category of models permits its objects themselves to serve as indexing bases? For $\Cat$, the answer uses comma categories and the two orientations. It is not supplied merely by knowing that its localization presents spaces.

This also explains the distinction from a slice of a model category. For $S\in\mathcal M$, the ordinary slice model structure, when available, localizes maps by their total weak equivalences. That is the coarse theory in these notes. A category of diagrams indexed by an ordinary category $S$, with objectwise weak equivalences, has a different role: it retains the directed indexing shape. In $\Cat$ the same symbol $S$ can play both roles, which is precisely what makes their comparison possible and initially deceptive.

No Grothendieck topologies or covering axioms for the four adjectives are defined in this paragraph. The expectations about equal cohomology should therefore be recorded as a programme, rather than as a theorem about a specified pair of sites. In particular, the remark identifying ``\'etale'' and ``Zariski'' belongs to these categorical local-isomorphism intuitions; it is not an assertion that the \v{C}ech or sheaf theories of the \emph{scheme-theoretic} \`etale and Zariski topologies coincide. The precise comparison statements available in these pages are the comma and fine-localization equivalences already established.
\end{dcomment}


Cette rapide digression ne convainc à présent que la compatibilité dite pour la construction de $C(X)$ et $i_X : X \longrightarrow C(X)$, ne peut en aucun cas être purement formelle, une pure histoire de catégories de Quillen disons. Ça donnera au mieux la compatibilité dans $(\mathcal{M}/S)W_S^{-1}$, ce qui est ridicule! Mais si on prend comme localiseur $\Sigma$ dans $\mathcal{M}/S$ les isomorphismes sans plus, la relation de compatibilité va devenir fausse à tous les coups! Comme ensemble de flèches $\Sigma$ un peu plus raisonnable, qui soit encore très petit, on peut songer p. ex. à la $W$-équivalence \textit{universelle} sur $S$ (nettement plus fine encore que $W_S^{\mathrm{d}}$, $W_S^{\mathrm{g}}$ dans le cas particulier de Cat). Aurait-on encore compatibilité pour cet ensemble de flèches? Je suis aussi un peu perplexe par le fait (pas bien compris) que les formules d’adjonction pour les topologies Zariskienne et co-Zariskienne soient en sens opposé. (Ou me serai-je trompé en dualisant à la va-vite?) Mais à vrai dire, cette compatibilité à elle seule, même quand on utilise une structure de site sur $\mathcal{M}$, ne suffit pas à écrire une formule d’adjonction. On a deux catégories localisées, % [sic] ne convainc -> me convainc, the verb needs an object pronoun
\[
\mathrm{HOT}^{\mathrm{lc}}_{\mathcal{M}} = (\mathcal{M}/S)W_S^{-1}, \qquad \mathrm{HOT}_W(S) = (\mathcal{M}/S)(W_S^{\mathrm{f}})^{-1},
\]
où $W_S^{\mathrm{f}}$ désigne une $W$-équivalence ‘fine’, correspondant à la notion de localisation, plus exactement aux changements de base ‘admissibles’ qu’on a choisis, et on a un foncteur de localisation
\msp{page 85}
\[
\mathrm{Hot}_W(S) \longrightarrow \mathrm{Hot}^{\mathrm{lc}}_{W}(S),
\]
mais pas de foncteur en sens inverse. Mais si, mais si, en travaillant avec la catégorie similaire $\usl{Fibhot}\,S$, ça doit garder un sens tout au moins quand on a une bonne notion de ‘fibration’. Je vois un petit travail de dégrossissage conceptuel à faire, avec peut-être une formule d’adjonction à la clef, mais je ne vais pas m’y lancer à présent.

\dgrblock{dgrb208}{A universal version of the coherence calculation}
\begin{dcomment}{development}{p84}{A universal version of the coherence calculation}
For the categorical replacement used here, the question has an affirmative answer under the standing hypotheses of the Theorem-Scholie. Put $U_S=\{u\mid u\times_ST\in W_T\text{ for every }T\to S\}$. Part~III of the Scholie identifies $W_S$ with the fibrewise class $W_S^{\mathrm f}$ on the full subcategory of Hot-fibrant objects. These objects are stable under arbitrary strict base change, as the source states explicitly at manuscript p.~119. If $u:P\to Q$ is a coarse weak equivalence between such objects and $f:T\to S$ is arbitrary, the fibre of its pullback at $t\in T$ is precisely
\[
 (u\times_ST)_t=u_{f(t)}.
\]
It therefore belongs to $W$. The pulled-back source and target remain Hot-fibrant, so Part~III over $T$ gives $u\times_ST\in W_T^{\mathrm f}=W_T$. This proves that $U_S$ and $W_S$ agree between Hot-fibrant objects; the converse uses the identity base change.

The values of $C$ are perfect, hence Hot-fibrant. Naturality of $\eta:1\to C$ and two-out-of-three imply $C(W_S)\subseteq W_S$, and the preceding comparison strengthens this to $C(W_S)\subseteq U_S$. Likewise $\eta_{CX}$ is a coarse weak equivalence between Hot-fibrant objects and therefore lies in $U_S$. Both iterated units are now invertible in $(\Cat/S)U_S^{-1}$, so the cancellation lemma at manuscript p.~82 proves
\[
 C\eta_X=\eta_{CX}\quad\hbox{in }(\Cat/S)U_S^{-1}.
\]
The unit and counit construction below then gives the corresponding reflection. Thus the source's particular question is answered using its fibrewise comparison and base-change stability; no additional identification of strict pullbacks with homotopy pullbacks is needed for that application.

There is also a useful sufficient criterion for other replacements, formulated directly in homotopy-theoretic terms. Here its additional hypothesis is explicit.

\textbf{Proposition (universal comparison for good replacements).}\ours
Fix $S$ and $W=W_\infty$. Let $\mathcal H\subseteq\Cat/S$ be a full subcategory such that, for every $P\to S$ in $\mathcal H$ and every $T\to S$, the canonical square obtained by applying nerves to $P\times_ST$ is a homotopy pullback square. Suppose $C:\Cat/S\to\mathcal H$ has a natural map $\eta:1\to C$ whose components belong to $W_S$. Let $U_S$ consist of those maps $f$ over $S$ for which $f\times_ST$ is a weak equivalence for every $T\to S$. Then
\[
 C(W_S)\subseteq U_S,\qquad
 C\eta_X=\eta_{CX}\quad\hbox{in }(\Cat/S)U_S^{-1}.
\]
Moreover the canonical functor
\[
 (\Cat/S)U_S^{-1}\longrightarrow(\Cat/S)W_S^{-1}
\]
has a fully faithful right adjoint induced by $C$.

To prove the assertion, first consider a coarse weak equivalence $f:P\to Q$ between objects of $\mathcal H$. For any $T\to S$, its base change is represented, up to weak equivalence, by the map between homotopy pullbacks
\[
 NP\mathbin{\times^{h}_{NS}}NT
   \longrightarrow NQ\mathbin{\times^{h}_{NS}}NT.
\]
This is an equivalence because homotopy pullback preserves equivalences of diagrams. Thus $f\in U_S$. The converse holds by taking $T=S$. Consequently $U_S$ and $W_S$ agree on maps between objects of $\mathcal H$.

Naturality and two-out-of-three imply that $C$ preserves $W_S$. Its values lie in $\mathcal H$, so the preceding paragraph gives $C(W_S)\subseteq U_S$. Also $\eta_{CX}$ is a coarse equivalence between two objects of $\mathcal H$, hence belongs to $U_S$. In $(\Cat/S)U_S^{-1}$ the two iterated units are therefore invertible. The cancellation lemma at manuscript p.~82 proves their equality. The unit of the proposed adjunction is the image of $\eta$, and its counit in the coarse localization is the inverse of that image. The two triangle identities follow from this equality and from $\eta^{-1}\eta=1$. Since the counit is invertible, the right adjoint is fully faithful. This proves all assertions.

This second criterion applies only after its homotopy-pullback hypothesis has been established, for example by a relative Theorem~B argument after each base change. It does not infer that hypothesis solely from the word ``fibration'', and it does not assert that every category equipped only with weak equivalences has such a replacement. The later Theorem-Scholie~2, especially manuscript pp.~111--112, supplies the abstract comparison for any finer class agreeing on fibrant objects; the opening argument verifies that agreement for the source's universal class. Our contribution here is the explicit verification for universal equivalences by the two stated methods. We checked the localization criterion in \citet[Proposition~5.2.7.4]{Lurie2009HTT} and the source's later abstract theorem; no priority is claimed for either ingredient.
\end{dcomment}


% ---- from Der-ChVI-345_p056-060.tex
\printedpage{54}
% head: A. Grothendieck : Dérivateurs, Ch. VI, Édition des Universités de Montpellier II et Paris VII

\begin{flushright}
16.11.90
\end{flushright}

Je vais quand même essayer de faire un travail ‘de dégrossissage’, car le diagramme des
15 catégories a l'air un peut touffu, % [sic] peut -> peu, adverb "un peu" intended
finalement! Je me donne
\[
(\cM,W) \qquad W\subseteq \Fl\,\cM,
\]
un $S\in\Ob\,\cM$, d'où $\cM/S$, avec $W_S\subseteq\Fl(\cM/S)$ induit par $W$. On va poser
\[
\Ho^{\mathrm{lc}}(S)=(\cM/S)W_S^{-1},
\]
ça jouera le rôle des types d'homotopie relativement constants. Je suppose d'autre part
que l'on a distingué une classe $F\subseteq\Fl(\cM)$ de flèches, jouant le rôle de ‘fibrations’ (ce
seront les morphismes Hot-fibrants dans Cat, donnant lieu à la sous-catégorie $\Fibhot S$
de $\mathrm{Cat}/S$). Je prévois qu'il faudra
\msp{page 86}supposer, outre que $\cM$ est stable par $\varprojlim$ finie, que $F$ est stable par changement de base
quelconque, et stable par composition, qu'il contient les isomorphismes. Et surtout qu'on
a un théorème de factorisation
\[
\begin{tikzcd}
X \arrow[rr, "f"] \arrow[dr, "i_X"'] & & S \\
& \bar{X} \arrow[ur, "\bar{f}"'] &
\end{tikzcd}
\qquad\qquad\qquad
i_X\in W,\ \bar{f}\in F
\]
pout % [sic] pout -> pour, misprint
toute flèche $f$ dans $\cM$. Chez \textsc{Quillen} on exige même que $i_X$ ait la propriété de
relèvement à droite relativement à $F$, mais ça ne devrait pas être essentiel, puisque je n'ai
pas eu à utiliser ce fait dans Cat. Quand je travaille dans la catégorie $\cM/S$, où $f$ est
sous-entendu pour un objet $X$ de ladite, je noterai (par abus de langage) $C(X)$ pour $\bar{X}$,
bien qu'il n'est pas dit qu'on ait un choix canonique de $C(X)$, encore moins que $C(X)$
dépende fonctoriellement de $X$.

\dgrblock{dgrb209}{The locally constant inclusion has two adjoints}
\begin{dcomment}{status}{p85}{The locally constant inclusion has two adjoints}
The programme has a sharper established form than the reflection proved here. For a derivator $\mathbb D$ and a small category $A$, \citet[Proposition~3.10 and Theorem~3.13]{Cisinski2009LocallyConstant} proves that the full inclusion of locally constant coherent diagrams has both a left and a right adjoint, and that, for fixed $A$, these subcategories form a derivator in an additional diagram variable. This is not the assertion that the varying-base assignment $S\mapsto\HOT_W^{\mathrm{lc}}(S)$ is a derivator. These are the theorem numbers in the published 2009 article; its earlier preprint uses different numbering. The model-category construction is given there by the local structures of \S2, with the fibrant objects characterized in Proposition~2.3.

For spaces, the reason is transparent in the language of homotopy-coherent functors. Let $q:NS\to\Pi_\infty S$ be groupoid completion. Restriction
\[
 q^*:\mathrm{Fun}((\Pi_\infty S)^{\mathrm{op}},\mathrm{Spaces})
       \longrightarrow\mathrm{Fun}((NS)^{\mathrm{op}},\mathrm{Spaces})
\]
is fully faithful and identifies its image with the functors carrying all arrows to equivalences. Its left and right Kan extensions give the two adjoints, since spaces admit small homotopy colimits and limits. Passing to homotopy categories gives the corresponding statement here.

The second adjoint of this one inclusion should not be identified with the functor $j$ in the Scholie. The latter embeds the coarse category into a \emph{different}, covariant fine theory. This keeps the later general theorem distinct from Grothendieck's two parallel reflections. The vague site comparison of the preceding paragraph requires a precise choice of sites before it acquires a similarly definite present status.
\end{dcomment}


Désignons par
\[
\Fib S\subseteq\cM/S \tag*{$(*)$}
\]
la sous-catégorie de $\cM/S$ formée des $X$ dont le morphisme structural $f_X:X\longrightarrow S$ est
dans $F$. Désignons (par abus de notation) par $W_S$ le localiseur induit sur $\Fib S$ par $W_S$,
alors l'inclusion $(*)$ définit
\[
(\Fib S)W_S^{-1}\xrightarrow{\ \delta\ }(\cM/S)W_S^{-1}\overset{\text{déf}}{=}\Ho^{\mathrm{lc}}_W(S).
\]

\msp{page 87}Une première question, c'est de voir que c'est une équivalence. On voit que c'est essentiellement
surjectif grâce à l'existence des $C(X)$, $i_X$. C'est peu! On va essayer de définir
une flèche en sens inverse
\[
\bar{C}:(\cM/S)W_S^{-1}\longrightarrow(\Fib S)W_S^{-1},
\]
% foot: Version 9 novembre 2001
\printedpage{55}
% head: A. Grothendieck : Dérivateurs, Ch. VI, Édition des Universités de Montpellier II et Paris VII
i.e. un \textit{foncteur}
\[
\gamma:\cM/S\longrightarrow(\Fib S)W_S^{-1}
\]
qui transforme $W_S\subseteq\Fl(\cM/S)$ en isomorphismes. On choisira, pour $X$ dans $\cM/S$, un
$C(X)$ dans $\Fib S$ et une
\[
i_X:X\longrightarrow C(X)\ , \qquad i_X\in W_S,
\]
et on posera
\[
\gamma(X)=C(X).
\]
La difficulté, c'est de montrer que ça ne dépend pas des choix faits, à isomorphisme
canonique près, et que c'est fonctoriel en $X$ - et le deuxième implique le premier.

Pour voir que ça marche, visiblement il faut supposer quelque chose de plus. Je vois deux
hypothèses qui, l'une et l'autre, pourraient faire marcher : l'une, c'est qu'on se donne
déjà $C$ comme foncteur et $i_X$ fonctoriel en $X$ (c'est ce qui était le cas dans Cat tantôt).
L'autre, c'est que
\msp{page 88}les $i_X$ ont, par rapport à $F$, la propriété de relèvement à là Quillen. % [sic] à là -> à la, accent on article "la"

\dgrblock{dgrb210}{Objectwise replacement and functorial replacement}
\begin{dcomment}{exposition}{p87}{Objectwise replacement and functorial replacement}
The restart of 16 November isolates a familiar but consequential distinction. A weak equivalence $X\to C(X)$ with fibrant target gives a representative of each localized object; it does not by itself define a functor on localized morphisms. Grothendieck correctly stops at essential surjectivity before adding a hypothesis that controls maps.

In the functorial case the proof uses much less than a full model structure. Put $\mathcal A=\mathcal M/S$, let $\mathcal H=\underline{\mathrm{Fib}}S$ be full, and let $D:\mathcal H\hookrightarrow\mathcal A$ be its inclusion. Suppose $C:\mathcal A\to\mathcal H$ is a functor and $\eta:1_{\mathcal A}\to DC$ has components in $W_S$. For $u:X\to Y$, naturality gives
\[
 DC(u)\eta_X=\eta_Yu.
\]
If $u\in W_S$, two-out-of-three implies $C(u)\in W_S$. Hence both functors descend to the localizations. The transformation $\eta$ becomes an isomorphism $1\cong\overline D\,\overline C$. Restricting it to $\mathcal H$ gives $1\cong\overline C\,\overline D$, since fullness ensures that the same components are maps in $\mathcal H$. Thus $\overline C$ and $\overline D$ are quasi-inverse equivalences.

For any full $\mathcal F$ with $\mathcal H\subseteq\mathcal F\subseteq\mathcal A$, the same proof applies to $C|_{\mathcal F}$ and the two inclusions. This proves the intermediate-category assertion at manuscript p.~89. No lifting axiom, finite-limit hypothesis, or base-change stability is used in this functorial branch. Those properties become relevant in the other branch, where lifts have to replace a functorial choice.

The word ``fibrant'' is relative to the chosen base: the objects of $\mathcal H$ are maps $X\to S$ in the distinguished class $F$. It does not assert that $X$ is fibrant as an object of $\mathcal M$, unless $S$ is terminal or an additional argument relates the two conditions. This distinction will also matter when path objects are formed inside the slice in the following pages.
\end{dcomment}


Je commence par l'hypothèse que $C$ est déjà un foncteur. Le diagramme
\[
\begin{tikzcd}[row sep=large]
X \arrow[r, "i_X"] \arrow[d, "u"'] & CX \arrow[d, "C(u)"] \\
X' \arrow[r, "i_{X'}"] & CX',
\end{tikzcd}
\]
où $i_X,\ i_{X'}\in W$, montre que $C$ transforme $W_S$ en $W_S$, donc en des isomorphismes de
$\Ho^{\mathrm{lc}}_W(S)$. Donc on tient le foncteur cherché $\gamma$. On a en fait des foncteurs
\[
\begin{tikzcd}
\Fib S \arrow[r, hook, "\overset{\text{incl.}}{D}"] & \cM/S \arrow[l, bend left=50, "C"]
\end{tikzcd}
\]
compatibles avec les localiseurs $W_S$, induisant $\delta$ et $\gamma$, et on a
\[
\mathrm{id}_{\cM/S}\longrightarrow DC, \qquad \text{dans $W_S$ argument par argument,}
\]
car on a
\[
X\xrightarrow{\ i_X\ }DC(X)=C(X)
\]
avec $i_X\in W_S$. De même, on a
\[
\mathrm{id}_{\Fib S}\longrightarrow CD \qquad \text{dans $W_S$ argument par argument,}
\]
% foot: Version 9 novembre 2001
\printedpage{56}
% head: A. Grothendieck : Dérivateurs, Ch. VI, Édition des Universités de Montpellier II et Paris VII
avec le même
\[
X\xrightarrow{\ i_X\ }CX \qquad \text{(cette fois $X$ dans $\Fib S$).}
\]

\msp{page 89}Donc on a sans larmes les deux foncteurs $\delta$, $\gamma$ comme quasi-inverses l'un de l'autre. Et le
lemme page 11 montre même que pour \textit{toute} sous-catégorie pleine $\cF$ de $\cM/S$, telle que
\[
\Fib S\subseteq\cF\subseteq\cM/S,
\]
les inclusions induisent des équivalences sur les localisés $W_S$
\[
(\Fib S)W_S^{-1}\xrightarrow{\ \sim\ }\cF W_S^{-1}\xrightarrow{\ \sim\ }\underbrace{(\cM/S)W_S^{-1}}_{\overset{\text{déf}}{=}\Ho^{\mathrm{lc}}_W(S)}.
\]
Je vais traiter maintenant le cas de type Quillen, pour me convaincre que l'existence du
foncteur $C$, et le choix particulier qu'on peut en faire, est irrelevant. (Ce dernier point
est d'ailleurs évident a posteriori, quand on sait que $\delta$ est une équivalence, car le quasi-inverse
$\gamma$ sera alors unique à isomorphisme unique près, et le choix d'un foncteur $C$, et
d'un homomorphisme fonctoriel $i:\mathrm{id}\longrightarrow C$, n'est qu'une façon parmi d'autres d'expliciter
le quasi-inverse.)

\msp{page 90}On se place donc dans l'hypothèse quillénienne sur les $i_X$. On choisit $C(X)$, $i_X$ pour
chaque $X$ dans $\cM/S$, et on pose comme tantôt
\[
\gamma(X)=C(X) \qquad \text{comme objet de $(\cM/S)W_S^{-1}$.}
\]
Montrons que c'est fonctoriel en $X$. Soit
\[
u:X\longrightarrow X'
\]
dans $\cM/S$, considérons le diagramme anodin (commutatif)
\[
\begin{tikzcd}[row sep=large]
X \arrow[rr, "u"] \arrow[d, "i_X"'] \arrow[drr, loosely dotted, "v=i_{X'}u"] & & X' \arrow[d, "i_{X'}"] \\
C(X) \arrow[rr, dotted, "\bar{u}"] \arrow[dr] & & C(X') \arrow[dl, "\in F"] \\
& S, &
\end{tikzcd}
\]
complétons le par la flèche en gros pointillés, $v=i_{X'}u$ (rendant le triangle supérieur
commutatif). (\footnote[54]{Lemme : soit
% footnote continues on printed p. 57
\[
\begin{tikzcd}[ampersand replacement=\&, row sep=large]
X \arrow[r, "u"] \arrow[d, "i"'] \& X' \arrow[d] \\
C \arrow[r, dotted] \& C'
\end{tikzcd}
\]
dans $\cM/S$ avec $i\in\mathrm{Cof}_0$, $(C'\longrightarrow S)\in F=\mathrm{Fib}$, alors la flèche en pointillés existe dans $\cM/S$, de façon
à rendre le carré commutatif. \texttt{[cf. Lemme p. 91]}}). L'hypothèse sur $i_X$ implique que la flèche $\bar{u}$ en pointillés fins existe,
de façon à avoir encore deux triangles commutatifs. Mais ça signifie exactement qu'il y a
un $S$-\textit{morphisme} $\bar{u}$, i.e. une flèche dans $\cM/S$, rendant commutatif le carré
% foot: Version 9 novembre 2001
\printedpage{57}
% head: A. Grothendieck : Dérivateurs, Ch. VI, Édition des Universités de Montpellier II et Paris VII
\[
\begin{tikzcd}[row sep=large]
X \arrow[r, "u"] \arrow[d, "i_X"'] & X' \arrow[d, "i_{X'}"] \\
C(X) \arrow[r, "\bar{u}"] & C(X'),
\end{tikzcd}
\]
\msp{page 91}et il s'impose de prendre
\[
\gamma(u)=\ \text{image de $\bar{u}$ dans $\Fl((\Fib S)W_S^{-1})$.}
\]
Il faut prouver que cette flèche ne dépend pas du choix de $\bar{u}$. On doit dégager un Lemme,
où on désigne par $\mathrm{Cof}_0$ l'ensemble des flèches de $\cM$ qu'ont % [sic] qu'ont -> qui ont, relative pronoun
la propriété de relèvement par
rapport aux éléments de $\mathrm{Fib}=F$. ($\mathrm{Cof}_0$ joue le rôle des cofibrations triviales.)

\dgrblock{dgrb211}{What makes a comparison canonical}
\begin{dcomment}{exposition}{p89}{What makes a comparison canonical}
The desired independence of replacement is an independence with specified comparison maps. This qualification explains the otherwise too compressed phrase ``unique \`a isomorphisme unique pr\`es'' without introducing a new gap in the argument.

Suppose $\delta:\mathcal H\to\mathcal A$ is an equivalence, and $\gamma,\gamma'$ are quasi-inverses equipped with isomorphisms $\epsilon:\delta\gamma\to1$ and $\epsilon':\delta\gamma'\to1$. Full faithfulness of $\delta$ gives a unique natural transformation $\theta:\gamma\to\gamma'$ satisfying
\[
 \delta\theta=(\epsilon')^{-1}\epsilon.
\]
Its components are invertible because their images are, so this is the unique comparison compatible with the specified identifications. This is the precise meaning needed for replacement functors. If the identifications are forgotten, a quasi-inverse can have nontrivial natural automorphisms; for the identity equivalence of the one-object category $BC_2$, its natural automorphisms are the two central elements of $C_2$.

At manuscript p.~90 the chosen map $i_X$ has the \emph{left} lifting property with respect to fibrations in current terminology: it is the left vertical side of the square, while $C(X')\to S$ is the right vertical side. The printed phrase ``rel\`evement \`a droite'' on p.~86 is resolved by this diagram. The lift supplies a possible $\bar u:C(X)\to C(X')$. Its uniqueness has not been asserted. What the next lemma proves is that different lifts give the same localized arrow. Once that is proved, identity and composition follow by choosing the identity lift and the composite lift. This is how a nonfunctorial choice of fibrant objects becomes a genuinely functorial construction after localization.
\end{dcomment}


\textbf{Lemme.} \textit{Soit un diagramme en traits pleins dans $\cM/S$}
\[
\begin{tikzcd}[row sep=large]
X \arrow[r, "u"] \arrow[d, "i"'] & X' \arrow[d, "i'"] \\
C \arrow[r, dotted, "v"] & C'
\end{tikzcd}
\]
\textit{avec $i\in$ Cof$_0$, et $C'\in$ Ob$\,\Fib S$ i.e. $(C'\longrightarrow S)\in$ Fib. Alors $v$ existe, de façon à rendre
le carré commutatif. De plus, si $i\in W$ et $C\in\Ob\,\Fib S$, la classe de $v$ comme flèche de
$(\Fib S)W_S^{-1}$ en est indépendante du choix fait.}

Il reste à prouver cette indépendance,
\msp{page 92}ce qui résultera aussitôt du corollaire :

\dgrblock{dgrb212}{Choosing replacements without choosing a functor}
\begin{dcomment}{context}{p91:Il reste à prouver cette indépendance,}{Choosing replacements without choosing a functor}
Grothendieck is separating the existence of good representatives from the possibility of doing mathematics with them. A factorization of every object over $S$ gives a representative in $\Fib S$, but it does not yet identify the morphisms of the two localizations. That is the work of the next lemma. The question is narrower than the construction of a model category: we require enough lifting to extend a map into a fibrant object, and enough paths to make two extensions agree after localization.

The lifting class is
\[
 \mathrm{Cof}_0={}^{\perp}\mathrm{Fib}.
\]
Thus its maps have the \emph{left} lifting property against fibrations. Neither the notation nor the definition says that every member of this class lies in $W$; this is why the factorization hypothesis subsequently asks for $i_X\in\mathrm{Cof}_0\cap W$. The manuscript's occasional phrase ``relèvement à droite'' describes the wrong side of the lifting square, as the editors explicitly correct on p.~108. For the extension problem here, the right vertical arrow is the structural map $C'\to S$. The square with top arrow $X\to C'$ and bottom arrow $C\to S$ has a lift because $i:X\to C$ belongs to ${}^{\perp}\mathrm{Fib}$. No choice of a cleavage or of a functorial factorization is involved.

What must become canonical is the \emph{class} of the extension, not an actual extension. Once its class is independent of the lift, a composite of two chosen lifts is itself an allowable lift of the composite. Independence therefore supplies compatibility with composition. This is an instructive instance of a construction becoming functorial only after passing to a homotopy category.

In modern language the other route, hypothesis A, is a right deformation: an endofunctor $C$ with a natural weak equivalence $\mathrm{id}\to C$ and values in a full subcategory of good objects. The equivalence of localizations follows from this data, without a model structure; compare the dual of \citet[Definition~2.2.1 and Remark~2.2.2]{Riehl2014}. Hypothesis B replaces this strictly natural data by lifting and a proof of independence. Its advantage is that a model category supplies it even when we do not wish to choose a functorial factorization. Its cost is the relative path-object argument that Grothendieck now gives.
\end{dcomment}


\textbf{Corollaire.} \textit{Soit un diagramme dans $\cM/S$}
\[
\begin{tikzcd}
X \arrow[r, "i"] & C \arrow[r, shift left, "v_1"] \arrow[r, shift right, "v_2"'] & C',
\end{tikzcd}
\]
\textit{avec $C$ et $C'$ dans $\Fib S$ (i.e. $C\longrightarrow S$ et $C'\longrightarrow S$ dans Fib), et $i\in$ Cof$_0$ et $i\in W$. Si
$v_1i=v_2i$, alors $v_1$ et $v_2$ définissent la même flèche de $(\Fib S)W_S^{-1}$.}
% foot: Version 9 novembre 2001

\printedpage{58}
% head: A. Grothendieck : Dérivateurs, Ch. VI, Édition des Universités de Montpellier II et Paris VII
\textsc{Démonstration}. Soit
\[
v=(v_1,v_2):C\longrightarrow C'\times C',
\]
\texttt{[plutôt} $C'\times_S C'$\texttt{]} qui rend commutatif par hypothèse le diagramme
\[
\begin{tikzcd}[row sep=large]
X \arrow[r, "w"] \arrow[d, "i"'] & C' \arrow[d, "\mathrm{diag}_{C'}"] \\
C \arrow[r, "v"] & C'\times_S C',
\end{tikzcd}
\]
où $w=v_1i=v_2i$. Factorisons $\mathrm{diag}_{C'}$ en
\[
C'\xrightarrow{\ i'\ }\bar{C}'\xrightarrow{\ \delta'\ }C'\times_S C', \qquad \text{avec $\delta'\in\mathrm{Fib}$,}
\]
et $i'\in W$. On trouve un carré commutatif
\[
\begin{tikzcd}[row sep=large]
X \arrow[r, "w'=i'w"] \arrow[d, "i"'] & \bar{C}' \arrow[d, "\delta'"] & \\
C \arrow[r, "v"] \arrow[ur, dotted, "\bar{v}"] & C'\times_S C' \arrow[r, shift left, "\mathrm{pr}_1"] \arrow[r, shift right, "\mathrm{pr}_2"'] & C'
\end{tikzcd}
\]
avec $i\in\mathrm{Cof}_0$, $\delta'\in\mathrm{Fib}$, d'où l'existence de la flèche en pointillés, rendant les deux triangles
commutatifs. On veut prouver $\mathrm{pr}_1\circ v=\mathrm{pr}_2\circ v$ dans $(\Fib S)W_S^{-1}$, comme $v=\delta'\bar{v}$ il suffit
de prouver que l'on a $\mathrm{pr}_1\delta'=\mathrm{pr}_2\delta'$ dans $(\Fib S)W_S^{-1}$.

\msp{page 93}Il faut seulement vérifier d'abord (pour justifier l'argument) que cette assertion a un sens,
i.e. que $\bar{C}'$ est lui aussi dans $\Fib S$. Mais les hypothèses de stabilité de $F=\mathrm{Fib}$ par
changement de base et composition assurent que $1^o$) $C'\times_S C'$ est dans $\Fib S$ (qui est
stable par produit), et par suite, $2^o$) que $\bar{C}'$ l'est aussi, puisque $\delta'\in\mathrm{Fib}$. Comme $C'$ est
aussi dans $\Fib S$ et $C'\xrightarrow{\ i'\ }\bar{C}'$ est dans $W_S$, pour vérifier $\mathrm{pr}_1\delta'=\mathrm{pr}_2\delta'$ dans le localisé de
$\Fib S$, il suffit de voir que $(\mathrm{pr}_1\delta')i'=(\mathrm{pr}_2\delta')i'$, or comme $\delta'i'=\mathrm{diag}_{C'}$, on a l'égalité même
dans $\Fib S$ (sans localiser).

On a donc défini $\gamma(X)$, $\gamma(u)$ pour objets et flèches de $\cM/S$. La transitivité pour $\gamma(u)$ est
claire, et aussi $\gamma(\mathrm{id})=\mathrm{id}$. Donc on tient le foncteur
\[
\gamma:\cM/S\longrightarrow(\Fib S)W_S^{-1}.
\]
Soit maintenant $\cF$ une sous-catégorie pleine de $\cM/S$, avec
\[
\Fib S\subseteq\cF\subseteq\cM/S,
\]
considérons $\gamma$ comme foncteur à valeurs dans $\cF W_S^{-1}$. (J'ai oublié de vérifier que $\gamma$ transforme
$W_S$ dans $W_S$, mais c'est tautologique.) Ainsi on a des foncteurs
% foot: Version 9 novembre 2001

% ---- from Der-ChVI-345_p061-065.tex
\printedpage{59}
% head: A. Grothendieck : Dérivateurs, Ch. VI, Édition des Universités de Montpellier II et Paris VII
\[
\begin{tikzcd}
\mathcal{F}W_S^{-1} \arrow[r, shift left, "\delta"] & (\mathcal{M}/S)W_S^{-1}, \arrow[l, shift left, "\gamma"]
\end{tikzcd}
\]
il faut voir qu'ils sont quasi-inverse l'une de \msp{page 94} l'autre.

\dgrblock{dgrb213}{The path object makes the choices agree}
\begin{dcomment}{exposition}{p93}{The path object makes the choices agree}
Grothendieck's proof works inside the subcategory whose localization we want to understand. This point prevents a circular argument: knowing that two lifts become equal in $(\mathcal M/S)[W_S^{-1}]$ would not yet show that they become equal in $(\Fib S)[W_S^{-1}]$, since full faithfulness of the inclusion is precisely what remains to be proved.

Write $\mathcal A=\mathcal M/S$ and $\mathcal F=\Fib S$. Suppose that $v_1,v_2:C\to C'$ agree after precomposition with $i:X\to C$, where $i$ has the left lifting property against fibrations and $C,C'$ lie in $\mathcal F$. Factor the relative diagonal as
\[
 C'\xrightarrow{j}P\xrightarrow{(d_0,d_1)}C'\times_S C',
 \qquad j\in W,\quad (d_0,d_1)\in\mathrm{Fib}.
\]
The object $C'\times_S C'$ is fibrant over $S$: either projection is a pullback of $C'\to S$, and its composite with the other structural map is a fibration. The same is true of $P$. The equalities $d_0j=d_1j=\mathrm{id}_{C'}$ therefore take place in $\mathcal F$. Since $j$ is inverted in $\mathcal F[W^{-1}]$, we obtain $d_0=d_1=j^{-1}$ there.

The square with left side $i$, right side $(d_0,d_1)$, bottom side $(v_1,v_2)$, and top side $jv_1i=jv_2i$ has a lift $h:C\to P$. Consequently $v_1=d_0h=d_1h=v_2$ in $\mathcal F[W^{-1}]$. This is the whole independence argument. In fact, the additional condition $i\in W$ is not needed for this particular conclusion. It is needed when the chosen $i_X$ are used to compare an arbitrary object with its replacement.

For completeness, let $u:X\to Y$ and choose a lift $\bar u:CX\to CY$ satisfying $\bar u i_X=i_Yu$. Define $\gamma(u)=[\bar u]$. The identity of $CX$ is a permissible lift for $\mathrm{id}_X$, and $\bar v\bar u$ is a permissible lift for $vu$. Independence proves the functor laws. If $u\in W$, the equation $\bar u i_X=i_Yu$ and two-out-of-three imply $\bar u\in W$. Thus $\gamma$ descends through the localization at $W_S$.

The same squares finish the naturality checks left on p.~94. In the coarse localization they give a natural isomorphism $\mathrm{id}\xrightarrow{\sim}D\gamma$, where $D$ denotes inclusion. On any full intermediate subcategory $\mathcal F\subseteq\mathcal G\subseteq\mathcal A$, both $X$ and $CX$ belong to $\mathcal G$, so the same maps give the second natural isomorphism $\mathrm{id}\xrightarrow{\sim}\gamma D$. These two isomorphisms establish the comparison theorem without identifying arbitrary fractions by hand. The proof also records exactly where closure of $\mathrm{Fib}$ under pullback and composition is used: in constructing the path object \emph{inside} $\mathcal F$.
\end{dcomment}


Du côté de $\mathcal{M}/S$ : on doit comparer $X$ et $\delta\gamma(X) = C(X)$, c'est fait par
\[
i_X : X \longrightarrow C(X)
\]
qui est dans $W_S$, d'où l'isomorphisme dans le localisé
\[
X \simeq \delta\gamma(X) \qquad \text{dans } (\mathcal{M}/S)W_S^{-1}.
\]
Il faut vérifier que c'est fonctoriel en $X$.

Du côté de $\mathcal{F}$, comparer $X$ et $\gamma\delta(X)$, i.e. $X$ et $C(X)$. Même tabac, on choisit $i_X : X \simeq C(X) = \gamma\delta(X)$. À vérifier la fonctorialité.

Il s'agit maintenant d'introduire un localisé plus fin de $\mathcal{M}/S$. Pour ceci, on suppose donné une autre classe de morphismes, je l'appelle $\mathcal{L}$ (comme `localisation')
\[
\mathcal{L} \subseteq \mathrm{Fl}(\mathcal{M}),
\]
p. ex. dans le cas de $\mathcal{M} = \mathrm{Cat}$, $\mathcal{L}$ sera formée des isomorphismes locaux (\footnote[55]{Ou aussi, des morphismes \textit{lisses}. Ça doit donner le même résultat.}) (s'il s'agit de définir $W_S^{\mathrm{g}}$), ou des isomorphismes colocaux (\footnote[56]{Ou aussi des morphismes propres.}) (s'il s'agit de définir $W_S^{\mathrm{d}}$). On supposera que $\mathcal{L}$ contient les isomorphismes, et $\mathcal{L}$ stable par composition (\footnote[57]{Il n'en est peut-être pas besoin.}), mais non par changement de base. On définit alors
\[
\WL \subseteq \mathrm{Fl}(\mathcal{M}/S)
\]
par
\[
f \in \WL \overset{\text{déf}}{\Longleftrightarrow} (f \times_S S') \in W,\ \ \forall (S' \longrightarrow S) \in \mathcal{L}.
\]
\msp{page 95}

Comme $\mathcal{L}$ contient $\mathrm{id}_S$, on aura
\[
\WL \subseteq W_S,
\]
d'où un foncteur canonique
\[
\underbrace{(\mathcal{M}/S)(\WL)^{-1}}_{\defeq \Ho_W^{\mathcal{L}}(S)} \xrightarrow{\ \rho\ } \underbrace{(\mathcal{M}/S)(W_S)^{-1}}_{\defeq \Ho_W^{\mathrm{lc}}\mathcal{L}(S)} % [sic] \Ho_W^{\mathrm{lc}}\mathcal{L}(S) -> \Ho_W^{\mathrm{lc}}(S), cf. the definition \Ho_W^{\mathrm{lc}}(S) \defeq (\mathcal{M}/S)W_S^{-1} on p. 61
\]
% foot: Version 9 novembre 2001
\printedpage{60}
% head: A. Grothendieck : Dérivateurs, Ch. VI, Édition des Universités de Montpellier II et Paris VII
qui est un foncteur de localisation. Je voudrais définir un foncteur en sens inverse. Pour ceci, il est naturel de supposer (en analogie avec la situation de Cat)

\textbf{Hypothèse 1.} \textit{Sur $\Fibu S$, $W_S = \WL$, i.e. si $f : X \longrightarrow Y$ est dans $W$ et $X$, $Y$ dans $\Fibu S$, alors pour tout changement de base $S' \longrightarrow S$ qui est dans $\mathcal{L}$, $f' : X_{S'} \longrightarrow Y_{S'}$ est encore dans $W$.}

Dans le contexte de Quillen, cette hypothèse a une nette tendance à être vérifiée quelque soit $\mathcal{L}$, i.e. pout \textit{tout} changement de base $S' \longrightarrow S$ sans restriction. Il suffit en fait que la structure de Quillen soit `propre', plus précisément (et moins restrictivement), qu'un \texttt{[morphisme]} $\in W$ reste tel après changement de base par un morphisme $\in$ Fib. Dégageons ceci comme un Lemme. % [sic] pout -> pour, misprint

\textbf{Lemme.} \textit{L'hypothèse est vérifiée si on suppose} \msp{page 96} \textit{que $F =$ Fib et $W$ sont reliés par la condition supplémentaire (en plus de la factorisation $i_X$) : pour tout carré cartésien}
\[
\begin{tikzcd}
X \arrow[d, "f"'] & X' \arrow[l, "q"'] \arrow[d, "f'"] \\
S & S', \arrow[l, "p"']
\end{tikzcd}
\]
\textit{si $f$ est dans Fib et $p$ dans $W$, alors $q \in W$.}

\dgrblock{dgrb214}{Testing equivalences after change of base}
\begin{dcomment}{exposition}{p95}{Testing equivalences after change of base}
The class $W_S^{\mathcal L}$ is designed to remember information that a total-space equivalence can lose. It asks that an equivalence survive every observation of the base allowed by $\mathcal L$. At this stage $\mathcal L$ is a collection of tests, not a topology and not a system of weak equivalences. Its role is entirely expressed by pullback.

For a fixed base $S$, the definition can be written
\[
 W_S^{\mathcal L}
 =\bigcap_{(g:T\to S)\in\mathcal L}(g^*)^{-1}(W_T).
\]
Each inverse image contains the isomorphisms and satisfies two-out-of-three; so does their intersection. The identity test gives $W_S^{\mathcal L}\subseteq W_S$. These statements use neither closure of $\mathcal L$ under composition nor closure under base change. The doubt in footnote~57 therefore has a positive answer for this part of the construction.

Composition becomes useful when we move the base. If $g:T\to S$ lies in $\mathcal L$ and $\mathcal L$ is composition-closed, then
\[
 g^*(W_S^{\mathcal L})\subseteq W_T^{\mathcal L}.
\]
Indeed, for a further test $h:U\to T$ in $\mathcal L$, the pullback $h^*g^*u$ is canonically isomorphic to $(gh)^*u$, and $gh$ is an allowed test over $S$. This proves contravariance along the chosen class. Contravariance along an arbitrary $g$ would need the stronger condition that $gh$ be allowed whenever $h$ is allowed; it does not follow from the definition alone.

The localization functor $\rho$ forgets the extra tests. The question about a right adjoint concerns a special class of objects on which those tests no longer distinguish additional weak equivalences. Hypothesis~1 says exactly that: two fibrations over $S$ which are equivalent on total spaces remain equivalent after every allowed base change. It permits the inclusion of fibrant objects to descend to the finer localization. It does \emph{not} say that every map $X\to CX$ is a fine equivalence; that would make the reflection an equivalence and erase the very distinction being introduced.

For the categorical localizers earlier in the chapter, local and colocal tests retain opposite directions of a diagram. Under the standing fundamental-localizer hypotheses, the two suggestions in footnotes~55--56 follow from the base-change theorem discussed at p.~122 below. Every proper base change preserves $W_S^{\mathrm d}$, and hence preserves the total weak equivalence of such a map. Conversely, the coslice projections $s\backslash S\to S$ are proper: they are discrete opfibrations. Testing against all proper maps therefore includes the defining coslice tests and gives exactly $W_S^{\mathrm d}$. The dual argument identifies testing against smooth maps with $W_S^{\mathrm g}$. When properness or smoothness is taken relative to a specified $W$, the corresponding relative base-change criterion is the needed hypothesis. These conclusions do not extend to an arbitrary test class in an arbitrary relative category. The commentary on Chapter~I, manuscript pp.~52--58, and on this chapter, pp.~41--44, explains the geometric criteria; the abstract argument here needs only equality of the two localizers on $\Fib S$.
\end{dcomment}


\textsc{Démonstration}. Soit
\[
\begin{tikzcd}
X \arrow[rr, "f"] \arrow[dr, "p"'] & & Y \arrow[dl, "q"] \\
& S &
\end{tikzcd}
\]
avec $p$, $q \in \mathrm{Fib}$, $f \in W$, montrons que pour \textit{tout} changement de base $g : S' \longrightarrow S$, $f_{S'} : X_{S'} \longrightarrow Y_{S'}$ est encore dans $W$. On factorise $g$ en
\[
\begin{tikzcd}
S' \arrow[r, "i"] \arrow[rr, bend right=15, "g"'] & \bar{S}' \arrow[r, "\bar{g}"] & S,
\end{tikzcd}
\qquad\qquad i \in W,\ g \in \mathrm{Fib},
\]
\texttt{[plutôt} $\bar{g} \in$ \texttt{Fib]} et on est ramené à prouver que l'assertion est vraie dans les deux cas particuliers où $g \in \mathrm{Fib}$, ou $g \in W$. Le cas $g \in \mathrm{Fib}$ résulte aussitôt de l'hypothèse du Lemme (appliquée au carré cartésien
\[
\begin{tikzcd}
X' \arrow[r, "f'"] \arrow[d, "q"'] & Y' \arrow[d, "p"] \\
X \arrow[r, "f"] & Y,
\end{tikzcd}
\]
% foot: Version 9 novembre 2001
\printedpage{61}
% head: A. Grothendieck : Dérivateurs, Ch. VI, Édition des Universités de Montpellier II et Paris VII
où $p \in \mathrm{Fib}$.)

Dans le cas où $S' \longrightarrow S$ est dans $W$, il s'ensuit par l'hypothèse du Lemme appliquée à
\[
\begin{tikzcd}
Y \arrow[d, "\in\mathrm{Fib}"] & Y' \arrow[l, "p"'] \arrow[d] \\
S & S' \arrow[l]
\end{tikzcd}
\]
que \msp{page 97} $p : Y' \longrightarrow Y$ est $\in W$, et de même pour $q : X' \longrightarrow X$, et $f \in W$ implique encore $f' \in W$.

Revenons à l'hypothèse p. 95. On aura donc un foncteur
\[
(\Fibu S, W_S = \WL) \longrightarrow (\mathcal{M}/S, \WL),
\]
d'où pour les localisés un foncteur horizontal
\[
\begin{tikzcd}
(\Fibu S)W_S^{-1} \arrow[r, "i_0"] \arrow[d, "\approx"] & (\mathcal{M}/S)(\WL)^{-1} \defeq \Ho_W^{\mathcal{L}}(S) \arrow[dl] \\
\Ho_W^{\mathrm{lc}}(S) \defeq (\mathcal{M}/S)W_S^{-1} &
\end{tikzcd}
\]
qui factorise le foncteur vertical induit par l'inclusion $(\Fibu S, W_S) \longrightarrow (\mathcal{M}/S, W_S)$ - foncteur dont on a vu (dans les deux cas envisagés) que c'est une équivalence. En inversant ce dernier et le composant avec $i_0$, on trouve donc un foncteur
\[
\underbrace{\Ho_W^{\mathrm{lc}}(S)}_{\defeq (\mathcal{M}/S)W_S^{-1}} \quad \xrightarrow{\ i\ } \quad \underbrace{\Ho_W^{\mathcal{L}}(S)}_{\defeq (\mathcal{M}/S)(\WL)^{-1}} \qquad ,
\]
tel que
\[
\rho i \simeq \mathrm{id},
\]
donc une \textit{rétraction} de $\Ho_W^{\mathcal{L}}(S)$ sur $\Ho_W^{\mathrm{lc}}(S)$. On s'attend que $i$ soit pleinement fidèle, et plus précisément encore, \msp{page 98} que $\rho$ et $i$ soient adjoints, $i$ étant l'adjoint à droite de $\rho$ :
\[
\begin{tikzcd}
\Ho_W^{\mathrm{lc}}(S) \arrow[r, shift right, "i"'] & \Ho_W^{\mathcal{L}}(S), \arrow[l, shift right, "\rho"']
\end{tikzcd}
\]
% foot: Version 9 novembre 2001
\printedpage{62}
% head: A. Grothendieck : Dérivateurs, Ch. VI, Édition des Universités de Montpellier II et Paris VII
i.e. on cherche
\[
\begin{aligned}
i\rho &\quad \underbrace{\xleftarrow{\ i\ }}_{\textit{pas}\ \text{iso}} \quad \mathrm{id} \quad \text{dans } \Ho_W^{\mathcal{L}} \quad (\text{localisation par } \WL)\\
\rho i &\quad \underbrace{\xrightarrow{\ p\ }}_{\text{iso}} \quad \mathrm{id} \quad \text{dans } \Ho_W^{\mathrm{lc}} \quad (\text{localisation par } W_S),
\end{aligned}
\]
\texttt{[trop de} $i$\texttt{'s]} donc pour $X$ dans $\mathcal{M}/S$, on cherche
\[
\begin{aligned}
i_X :&\quad X \quad \longrightarrow \quad C(X) \quad \text{tout au moins dans } (\mathcal{M}/S)(\WL)^{-1}\\
p_X :&\quad X \quad \longleftarrow \quad C(X) \quad \text{tout au moins dans } (\mathcal{M}/S)(W_S)^{-1}.
\end{aligned}
\]
On choisira pour $i_X$ la flèche `canonique' de $X$ dans $C(X)$, définie dans $\mathcal{M}/S$ lui-même, sans avoir à localiser. Mais dans le `cas Quillen' (quand $C$ n'est pas donné comme foncteur $\mathcal{M}/S \longrightarrow \Fibu S$), il faut voir que cet $i_X$ est fonctoriel en $X$, quand on le regarde dans le localisé $(\mathcal{M}/S)W_S^{-1}$. Mais c'est contenu dans la construction de $C$ comme \textit{foncteur}
\[
\mathcal{M}/S \longrightarrow (\Fibu S)W_S^{-1}.
\]
On notera que $i_X$ est dans $W_S$ sans plus, pas dans $\WL$, donc ne définit pas en général un isomorphisme de $\Ho_W^{\mathcal{L}}$.

\dgrblock{dgrb215}{Why the properness argument works}
\begin{dcomment}{exposition}{p97}{Why the properness argument works}
Grothendieck wants an easily checked condition which guarantees that a weak equivalence between fibrations remains an equivalence after a change of base. The lemma gives a sufficient condition in his sparse axiomatic setting. Its two cases are worth keeping distinct, because they apply the same properness axiom to different squares.

Let $u:X\to Y$ be a weak equivalence over $S$, with $X\to S$ and $Y\to S$ fibrations, and let $g:T\to S$. Factor $g$ as
\[
 T\xrightarrow{a}\bar T\xrightarrow{b}S,
 \qquad a\in W,\quad b\in\mathrm{Fib}.
\]
First pull back along $b$. The map $Y\times_S\bar T\to Y$ is a fibration, being a pullback of $b$. Right properness, applied to $u$ along this fibration, says that $u_{\bar T}$ lies in $W$. The structural maps of its source and target over $\bar T$ are still fibrations.

Now pull back along $a$. Right properness applied to those two structural fibrations says that both projections
\[
 X\times_S T\longrightarrow X\times_S\bar T,
 \qquad
 Y\times_S T\longrightarrow Y\times_S\bar T
\]
lie in $W$. In their commutative square the lower horizontal arrow is $u_{\bar T}\in W$. Two-out-of-three proves $u_T\in W$. This completes the argument, with no assumption that $g$ itself be a fibration or a weak equivalence.

The sufficient condition is stronger than necessary in a full Quillen model category. For every map $g:T\to S$, the pullback functor between the slice model categories preserves fibrations and trivial fibrations. It is right adjoint to postcomposition with $g$, and is therefore a right Quillen functor. The dual of Ken Brown's lemma shows that it preserves weak equivalences between objects fibrant in the slice. Thus Hypothesis~1 holds for all base changes in any model category, without global right properness. The base $S$ need not itself be fibrant: fibrancy of $X\to S$ is fibrancy as an object of the slice model category. The corresponding statement for a category of fibrant objects is proved as Lemma~C.1.11 of \citet{RiehlVerity2022Elements}, using Lemma~C.1.10. This does not invalidate Grothendieck's proof: his preceding axioms are weaker than the full model-category axioms, so a separate sufficient condition is useful.

Right properness concerns pulling back an arbitrary weak equivalence \emph{along a fibration}. The conclusion here concerns pulling back a weak equivalence \emph{between two fibrations over the same base} along an arbitrary map. The factorization argument explains the passage from one arrangement to the other.
\end{dcomment}


On définit $p_X$ à l'aide de $i_X$, comme l'inverse de $i_X$ \textit{dans} $(\mathcal{M}/S)W_S^{-1}$, qui existe car \msp{page 99} $i_X$ est dans $W_S$, donc \texttt{[est un]} isomorphisme dans $(\mathcal{M}/S)W_S^{-1} = \Ho_W^{\mathrm{lc}}$, et $p$ est donc \texttt{[un]} isomorphisme. On retrouve donc $\rho i \simeq \mathrm{id}$, donc que $\rho$ est une \textit{rétraction} de $\Ho_W^{\mathcal{L}}$ sur $\Ho_W^{\mathrm{lc}}$.

Pour prouver l'adjonction, en procédant comme dans le cas de Cat, on est ramené à prouver que pour tout $X$,
\[
\left\{
\begin{aligned}
&C(i_X) = i_{C(X)} : C(X) \rightrightarrows C(C(X))\\
&\textit{dans } \Ho_W^{\mathcal{L}} \ (\text{localisation par } \WL).
\end{aligned}
\right.
\]
J'y étais arrivé dans Cat, en utilisant un morphisme
\[
C(C(X)) \xrightarrow{\ c\ } C(X)
\]
(essentiellement, un morphisme de composition des chemins) tel que $c \in \WL$, et que $c \circ C(i_X) = c \circ i_{C(X)}$. Ça avait l'air très délicat, car c'était à grand peine que j'étais arrivé à construire un foncteur \underline{Ch} (`chemins') raisonnable, où on puisse bel et bien composer les chemins! Je renonce pour le moment à axiomatiser les propriétés du foncteur \underline{Ch} dont j'ai besoin pour faire marcher la machine (je devrais y revenir tantôt). Je vais commencer plutôt par le `cas Quillen', avec des $i_X \in \mathrm{Cof}_0$.
% foot: Version 9 novembre 2001

\printedpage{63}
% head: A. Grothendieck : Dérivateurs, Ch. VI, Édition des Universités de Montpellier II et Paris VII
\msp{page 100}

On reprend la description du `foncteur' $C$ (à valeurs dans $(\Fibu S)W_S^{-1}$, et alors sans guillemets), on trouve que $C(i_X)$, où
\[
X \xrightarrow{u = i_X} X' = C(X)
\]
est défini simplement (dans $(\Fibu S)W_S^{-1}$) par `la' flèche $C(u)$ rendant commutatif le diagramme (dans $\Fibu S$)
\[
\begin{tikzcd}
X \arrow[r, "u\ (=i_X)"] \arrow[d, "i_X"'] & X'\ (= C(X)) \arrow[d, "i_{X'}\ (=i_{C(X)})"] \\
C(X) \arrow[r, "C(u)\ (=C(i_X))"] & C(X')\ (= C(C(X))),
\end{tikzcd}
\]
de sorte qu'on a dans $\mathcal{M}/S$ le diagramme
\[
\begin{tikzcd}
X \arrow[r, "\overbrace{i_X}^{\in\mathrm{Cof}_0\cap W}"] & \overbrace{C(X)}^{\in\mathrm{Ob}\,\Fibu S} \arrow[r, shift left, "i_{C(X)}"] \arrow[r, shift right, "C(i_X)"'] & \overbrace{C(C(X))}^{\in\mathrm{Ob}\,\Fibu S},
\end{tikzcd}
\]
commutatif, i.e. $i_{C(X)}i_X = C(i_X)i_X$. Le corollaire page 92 nous dit que les flèches $i_{C(X)}$ et $C(i_X)$ sont égales dans $(\Fibu S)W_S^{-1}$. Mais par hypothèse (p. 95), dans $\Fibu S$ on a $W_S = \WL$, donc les flèches sont égales dans $(\Fibu S)(\WL)^{-1}$, et a fortiori dans $(\mathcal{M}/S)(\WL)^{-1}$, et on gagne.

Voyons ce qu'on peut faire dans le cas où on dispose d'un \textit{foncteur} $C$. On suppose bien sûr que ce foncteur dans $\mathcal{M}/S$ \msp{page 101} n'est pas lié au $S$ particulier dans $\mathcal{M}$ qu'on a choisi comme base, que sa construction provient à la \textsc{Cartan-Serre} d'un foncteur
\[
X \longmapsto \mathrm{Ch}(X) : \mathcal{M} \longrightarrow \mathcal{M},
\]
(`objet des chemins') et de morphismes fonctoriels
\[
X \xrightarrow{\alpha_X} \mathrm{Ch}(X) \xrightarrow{p_X} X \times X, \qquad\qquad \boxed{p_X\alpha_X = \mathrm{diag}_X}
\]
(\footnote[58]{\textbf{NB} On pose $\mathrm{pr}_1 \circ p_X = \sigma_X$, $\mathrm{pr}_2 \circ p_X = \beta_X$ (flèches `source' et `but' pour l'objet $X$).}). Je présume qu'il faudra supposer surtout ceci

\dgrblock{dgrb216}{Naturality in the finer localization}
\begin{dcomment}{exposition}{p100}{Naturality in the finer localization}
The proposed unit of the reflection has to exist where fewer arrows have been inverted. This is the issue hidden by the notation $i_X$: in the coarse localization $i_X$ is invertible, whereas in the finer localization it usually is not. The construction by lifts supplies the required stronger naturality, although the sentence on p.~98 mentions only the coarse category.

Put $\mathcal A=\mathcal M/S$, $\mathcal F=\Fib S$, and write $W_0$ for either $W_S^{\mathcal L}$ or the more general fine localizer of Theorem-Scholie~2. We assume $W_0\subseteq W_S$ and equality on arrows of $\mathcal F$. For a map $u:X\to Y$, a chosen lift satisfies the strict identity
\[
 \bar u i_X=i_Yu
\]
in $\mathcal A$. Two choices of $\bar u$ agree in $\mathcal F[W_S^{-1}]$, hence also after mapping into $\mathcal A[W_0^{-1}]$, because $W_S$ and $W_0$ coincide on $\mathcal F$. The displayed identity therefore defines a natural transformation in the finer localization. Naturality for formal inverses of $W_0$-arrows follows by multiplying a naturality square by those inverses; this is precisely the universal property of localization for natural transformations.

For $u=i_X$, both $\overline{i_X}$ and $i_{CX}$ extend the map $i_{CX}i_X$. The path-object lemma identifies them in $\mathcal F[W_S^{-1}]$. This is stronger than the equality needed in $\mathcal A[W_0^{-1}]$. Let $\eta_X=[i_X]$ be the fine unit, and let the coarse counit be $[i_X]^{-1}$. One triangle identity is cancellation of $[i_X]$ with its inverse. The other is cancellation after the equality
\[
 [C(i_X)]=[i_{CX}]
\]
just proved. Consequently $\rho\dashv i$, and the counit is invertible, so $i$ is fully faithful.

No strict functor $C:\mathcal A\to\mathcal F$ has been manufactured in this argument. The expression $C(i_X)$ means a class represented by an allowable lift. The remark on p.~100 that the entire displayed square lies in $\Fib S$ is a small source slip: its upper-left object $X$ need not be fibrant. The square lies in $\mathcal M/S$; the two arrows whose classes are compared have fibrant sources and targets, which is all the proof uses.
\end{dcomment}


\textbf{Hypothèse 2.} \textit{$\alpha_X \in W$, $p_X \in$} Fib.

Pour un $X$ sur $S$, on construira $C_S(X)$ ou $C(X)$ ainsi
% foot: Version 9 novembre 2001

% ---- from Der-ChVI-345_p066-070.tex
\printedpage{64}
% head: A. Grothendieck : Dérivateurs, Ch. VI, Édition des Universités de Montpellier II et Paris VII

\[
\begin{tikzcd}[column sep=large]
 & \Ch(S) \arrow[d, "p_S"] \arrow[dl, dotted, "\beta_S"'] \arrow[dd, dotted, bend right=40, "\sigma_S"'] \arrow[dr, phantom, "\text{cart.}"]
 & C_S(X)=C(X) \arrow[l, "\pi"'] \arrow[d, "p'"] \arrow[dll, bend right=40, "\bar f_X=f_{C(X)}"'] \\
S & S\times S \arrow[l, "\pr_2^S"'] \arrow[d, "\pr_1^S"] \arrow[dr, phantom, "\text{cart.}"]
 & X\times S \arrow[l, "\pi_1=f_X\times\mathrm{id}_S"] \arrow[d, "\pr_1"] \\
 & S & X(=S'). \arrow[l, "f_X"'] \arrow[uu, dotted, bend right=40, "i_X"']
\end{tikzcd}
\]

Il est l'image inverse de $(\Ch(S),\pr_1^Sp_S=\sigma_S)$ par $X\xrightarrow{f_X}S$, mais on a explicité la factorisation de $\sigma_S$ sur le diagramme pour réinterpréter la flèche structurale $\bar f_X=f_{C(X)}$, définie comme
\[
f_{C(X)}=\beta_S\pi(=\pr_2^S\underbrace{p_S\pi}_{=\pi_1p'})
\]
par la formule
\[
f_{C(X)}=\underbrace{(\pr_2^S\pi_1)}_{\text{projection } \pr_2 \text{ de } X\times S \text{ dans } S}p'.
\]

\msp{page 102}
Or
\[
\begin{aligned}
&p'\in\mathrm{Fib}, && \text{car déduit de } p_S\in\mathrm{Fib} \text{ par changement de base } \pi_1,\\
&(\pr_2 : X\times S\longrightarrow S)\in\mathrm{Fib} && \text{pourvu que } X\longrightarrow\underbrace{e}_{\text{objet final de } \mathcal{M}}\quad \text{soit dans Fib (même raison),}
\end{aligned}
\]
donc si $X\in\Fib\,e$, on a $f_{C(X)}\in\mathrm{Fib}$. Comme c'est une propriété essentielle qu'on veut avoir pour \textit{tout} $X$ dans $\mathcal{M}/S$, on va faire l'hypothèse supplémentaire

\textbf{Hypothèse 3.} \textit{$\forall\,X\in\mathit{Ob}\,\mathcal{M}$, $X\longrightarrow e$ est dans Fib.}

Quant à
\[
i_X : X\longrightarrow C(X),
\]
il est déduit de la section $\alpha_S$ de $(\Ch(S),\sigma_S)$ sur $S$ par changement de base $X\longrightarrow S$. Notons maintenant que
\[
\begin{tikzcd}
\Ch(S) \arrow[r, shift left, "\sigma_S"] \arrow[r, shift right, "\beta_S"'] & S
\end{tikzcd}
\qquad\qquad\qquad \text{sont dans } W,
\]
car leurs composés avec $S\xrightarrow{\alpha_S}\Ch(S)$ sont l'identité, et $\alpha_S\in W$. D'autre part, $\sigma_S=\pr_1^S\circ p$ est dans Fib, comme composé de deux éléments dans Fib (NB $S\longrightarrow e$ dans Fib par Hypothèse 3). Ainsi $\sigma_S$, $\beta_S$ sont dans $\mathrm{Fib}\cap W$. On aura besoin de % [sic] \pr_1^S\circ p -> \pr_1^S\circ p_S, la flèche est notée p_S dans le diagramme et plus haut (\pr_1^Sp_S=\sigma_S)

\textbf{Hypothèse 4.} \textit{L'ensemble de flèches $\mathit{Fib}\cap W$ est stable par tout changement de base.}

% foot: Version 9 novembre 2001

\printedpage{65}
% head: A. Grothendieck : Dérivateurs, Ch. VI, Édition des Universités de Montpellier II et Paris VII

\msp{page 103}
Cela implique donc que
\[
C(X)\xrightarrow{\pr_1\pi'}X
\]
\texttt{[plutôt} $\pr_1p'$\texttt{]} est $\in W$ (étant déduit de $\sigma_S : \Ch(S)\longrightarrow S$ par changement de base), donc toute section $i_X : X\longrightarrow C(X)$ de $C(X)$ sur $X$ est dans $i_X$ \texttt{[plutôt dans} $W$\texttt{]}. On a donc bien une factorisation, fonctorielle en $X$
\[
\begin{tikzcd}
 & C(X) \arrow[dr, "\bar f_X=f_{C(X)}"] & \\
X \arrow[ur, "i_X"] \arrow[rr, "f_X"] & & S
\end{tikzcd}
\]
avec
\[
i_X\in W,\qquad\qquad f_{C(X)}\in\mathrm{Fib}.
\]
C'est tout ce qu'on avait utilisé pour établir
\[
(\Fib\,S)W_S^{-1}\xrightarrow{\sim}(\mathcal{M}/S)W_S^{-1}.
\]
Mais il s'agit maintenant, sous l'hypothèse où on introduit un $\mathcal{L}$, de prouver que pour tout $X$ dans $\mathcal{M}/S$, les deux flèches
\[
\begin{tikzcd}
C(X) \arrow[r, shift left, "C(i_X)"] \arrow[r, shift right, "i_{C(X)}"'] & C(C(X))
\end{tikzcd}
\]
sont égales dans $(\mathcal{M}/S)(W_S^{\mathcal{L}})^{-1}$. Comme par hypothèse on a $W_S^{\mathcal{L}}=W_S$ dans $\Fib\,S$ (Hypothèse 1), et que $C(X),\ C(C(X))\in\mathrm{Ob}\,\Fib\,S$,
\msp{page 104}
il suffit de prouver que les deux flèches sont égales \textit{dans} $(\Fib\,S)W_S^{-1}$. Cet énoncé est indépendant de la donnée de $\mathcal{L}$, et on a envie de le prouver.

\dgrblock{dgrb217}{The mapping-path factorization and its hypotheses}
\begin{dcomment}{exposition}{p103}{The mapping-path factorization and its hypotheses}
Grothendieck now wants a replacement with visible geometry. A point of $C_S(X)$ consists of a point of $X$ together with a path from its image to a new point of $S$. The new structural map remembers the endpoint. This changes the map to $S$ while preserving the total homotopy type, and explains why the two projections of the path object play different roles.

Assume finite limits, pullback and composition stability of $\mathrm{Fib}$, and natural maps
\[
 S\xrightarrow{\alpha_S}P(S)
 \xrightarrow{(\sigma_S,\beta_S)}S\times S
\]
with composite the diagonal, $\alpha_S\in W$, and $(\sigma_S,\beta_S)\in\mathrm{Fib}$. Define
\[
 C_S(X)=X\times_{S,\sigma_S}P(S),\qquad
 q:C_S(X)\to X,\qquad
 \bar f=\beta_S\operatorname{pr}_{P(S)}.
\]
The unit sends $x$ to $(x,\alpha_Sf(x))$, so $qi_X=\mathrm{id}_X$ and $\bar f i_X=f$. These formulas describe genuine morphisms: the generalized-element description is justified by the universal property of the pullback and the Yoneda embedding, not by an assumption that the category consists of sets.

To see that $\bar f$ is a fibration, rewrite the pullback as
\[
 C_S(X)=(X\times S)\times_{S\times S}P(S),
\]
where $X\times S\to S\times S$ is $f\times\mathrm{id}$. Its projection $p':C_S(X)\to X\times S$ is a fibration. The second projection $X\times S\to S$ is a pullback of $X\to e$, so it is a fibration if $X$ is fibrant. Hypothesis~3 makes this true for every $X$. Thus $\bar f=\operatorname{pr}_2p'$ is a fibration.

To see that $i_X$ is a weak equivalence, first note that $\sigma_S\alpha_S=\mathrm{id}$ makes $\sigma_S$ a weak equivalence. It is a fibration since $S$ is fibrant and $(\sigma_S,\beta_S)$ is a fibration. Hypothesis~4 then makes its pullback $q$ a trivial fibration. Since $qi_X=\mathrm{id}$, two-out-of-three gives $i_X\in W$. The proof uses pullback stability of \emph{trivial fibrations}; it does not require arbitrary weak equivalences to survive arbitrary pullbacks.

This is the mapping-path mechanism behind Brown's factorization lemma. Its persistence in higher category theory is especially explicit in \citet[Lemma~1.2.19]{RiehlVerity2022Elements}, where a map in an $\infty$-cosmos factors as a section of a trivial fibration followed by an isofibration. The categorical hypotheses and the meaning of equivalence differ, but the pullback calculation is the same.

The concise condition A$'$ on p.~112 must be read as the advertised \emph{strengthening of A}. In particular, it retains the conclusion $i_X\in W$ already required in A. If its path data are instead used as an independent recipe for proving A, the pullback-stability step just exhibited must be supplied, for example by Hypothesis~4. This distinguishes an abbreviated restatement from a missing hypothesis in the theorem as actually formulated.
\end{dcomment}


Peut-être y a-t-il moyen d'expliciter un langage géométrique, pour s'y reconnaître. Identifions les objets $X$ de $\mathcal{M}$ aux foncteurs $\mathcal{M}^o\longrightarrow(\mathrm{Ens})$ qu'ils définissent,
\[
X\longmapsto(T\longmapsto X(T)=\mathrm{Hom}_{\mathcal{M}}(T,X))\,.
\]
Ainsi, pout tout $T$ % [sic] pout -> pour, faute de frappe
\[
\Ch(X)(T)
\]
est un ensemble, appelé \textit{l'ensemble des chemins dans $X(T)$.} Il s'envoie (par $p_X$) fonctoriellement (en $T$) dans $(X\times X)(T)=X(T)\times X(T)$, i.e. chaque chemin $c$ dans $X(T)$ a une origine $\sigma(c)$, et un but $\beta(c)$. De plus, par $\alpha_X$, $X(T)$ s'envoie (fonctoriellement en $T$) dans l'ensemble des chemins de $X(T)$ : pour chaque élément $x$ dans $X(T)$, il y a donc un chemin (noté $1_x$ ou $\alpha_x$) bien défini de source et but $x$, qu'on appelle le chemin trivial, ou vide. Ceci dit,

% foot: Version 9 novembre 2001

\printedpage{66}
% head: A. Grothendieck : Dérivateurs, Ch. VI, Édition des Universités de Montpellier II et Paris VII

\msp{page 105}
pour $X$ dans $\mathcal{M}/S$, $C(X)$ représente le contra-foncteur suivant dans $\mathcal{M}$ : pour tout $T$ dans $\mathcal{M}$,
\[
C(X)(T)=\Bigl\{(x,c)\Bigm| x\in X(T),\ c : f_X(x)\xrightarrow{c}s\ \text{ un chemin dans } S(T), \text{ de source } f_X(x)\Bigr\}\,,
\]
où $f_X : X\longrightarrow S$ est le morphisme structural.

Les ‘calculs’ faits, il y a deux (?) jours montrent qu'on a de même une interprétation
\[
\begin{aligned}
C(C(X))(T)=\Bigl\{(x,c,c')\Bigm| x\in X(T),\ f_X(x)\xrightarrow{c}s\xrightarrow{c'}t,\\
c\text{ un chemin de source } f_X(x),\\
c'\text{ un chemin de source égale à } \beta(c)\Bigr\}
\end{aligned}
\]
et les applications $i_{C(X)}$ et $C(i_X)$ s'identifient respectivement à
\[
\begin{aligned}
(x,c)\quad&\longmapsto\quad(x,1_{f(x)},c)\quad f(x)\xrightarrow{1_x}f(x)\xrightarrow{c}s\\
(x,c)\quad&\longmapsto\quad(x,c,1_{\beta(c)})\quad f(x)\longrightarrow s\xrightarrow{1_s}s
\end{aligned}
\]
\texttt{[plutôt} $f(x)\xrightarrow{1_{f(x)}}f(x)\xrightarrow{c}s$\texttt{]}. Pour les comparer, on a une envie irrésistible d'écrire
\[
c\circ 1_x=1_s\circ c=c,
\]
mais a-t-on le droit? Il faudrait pouvoir composer les chemins, et c'est très, très difficile!

Mais on peut toujours axiomatiser. On suppose donc que pour $X$ dans $\mathcal{M}$, $\Ch(X)$
\msp{page 106}
n'est pas seulement une structure de (soi-disants) ‘chemins’ dans $X$ amorphe, seulement la donnée de la source et du but des chemins et des chemins triviaux (ces structures sur $\Ch(X)$ dépendant fonctoriellement de $X$), mais qu'on a une composition des chemins quand ceux-ci sont ‘composables’, i.e. une flèche d'un produit fibré
\[
(c,c')\longmapsto c'\circ c : (\Ch(X),\beta_X)\times_X(\Ch(X),\sigma_X)\xrightarrow{\pi_X}\Ch(X),
\]
i.e. on peut composer des chemins dans $X(T)$
\[
x\xrightarrow{c}y\xrightarrow{c'}z
\]
si et seulement si $\text{but}\,c=\text{source}\,c'$, et de plus on exige bien sûr que $\text{source}(c'\circ c)=\text{source}\,c$, $\text{but}(c'\circ c)=\text{but}\,c'$, ce qui s'exprime par un diagramme commutatif
\[
\begin{tikzcd}
(\Ch(X),\beta_X)\times_X(\Ch(X),\sigma_X) \arrow[rr] \arrow[dr] & & \Ch(X) \arrow[dl, "p_X"] \\
 & X\times X. &
\end{tikzcd}
\]

\dgrblock{dgrb218}{A groupoid in which both inserted paths are visible}
\begin{dcomment}{examples}{p105}{A groupoid in which both inserted paths are visible}
Grothendieck's element notation is most useful when it allows us to see the two maps $CX\rightrightarrows C^2X$ before localizing. A groupoid gives a completely algebraic example, including a small correction to the order of the formulas on p.~105.

Let $G$ be a group, let $S=BG$ be its one-object groupoid, and let $X=e\to BG$ select the object. Work in groupoids with equivalences as weak equivalences and isofibrations as fibrations. Take $P(S)$ to be the arrow groupoid: its objects are $g\in G$, and a morphism $g\to g'$ is a pair $(a,b)$ with
\[
 bg=g'a.
\]
Source and target send this pair to $a$ and $b$. The constant-path functor sends the unique object of $BG$ to $1\in G$ and an arrow $a$ to $(a,a)$. It is fully faithful, and every $g$ is isomorphic to $1$, so it is an equivalence. The endpoint functor $P(S)\to S\times S$ is an isofibration: prescribed endpoint isomorphisms $(a,b)$ lift by choosing $g'=bga^{-1}$.

The pullback defining $C_S(e)$ fixes the source component of a morphism to be the identity. Its objects are still $g\in G$, but between $g$ and $g'$ there is exactly one morphism, whose image in $BG$ is
\[
 b=g'g^{-1}.
\]
Thus $C_S(e)$ is a contractible groupoid. Its projection to $BG$ is an isofibration: an arrow $h$ with source the image of $g$ lifts to the unique arrow $g\to hg$. The inclusion $e\to C_S(e)$ chooses $1$ and is an equivalence. The strict fibre over the object of $BG$, on the other hand, is the discrete set $G$, since $g'g^{-1}=1$ precisely when $g'=g$. The replacement has preserved the total type while encoding the nontrivial fibre needed over $BG$.

Objects of $C_S^2(e)$ are pairs $(g,h)$, representing successive paths. The two maps are
\[
 i_{C_S(e)}(g)=(g,1),\qquad
 C_S(i_e)(g)=(1,g).
\]
The first inserts a constant path at the \emph{endpoint}; the second first replaces the original object $e$ and therefore inserts a constant path at the \emph{source}. These are the reverse of the labels assigned by ``respectivement'' in the edition. The formulas themselves are present there; their attribution is interchanged. We retain that source slip and use the correctly labelled formulas in the commentary.

Composition gives $\gamma(g,h)=hg$. It is also a functor on morphisms. The unique morphism $(g,h)\to(g',h')$ has image
\[
 d=h'g'g^{-1}h^{-1}
\]
in $BG$, which is exactly the image of the unique morphism $hg\to h'g'$ in $C_S(e)$. Therefore $\gamma i_C=\gamma C(i)=\mathrm{id}$ already before localization. Since the two insertion maps are weak equivalences between fibrant objects, they become the same inverse of $\gamma$ after localization. This example exhibits the geometric comparison which the abstract argument needs; it does not identify categorical fibrations with model-category fibrations in general.

The same example makes a noninvertible fine unit concrete. Let $W_0$ on groupoids over $BG$ consist of equivalences on each strict fibre. It is a two-out-of-three class, obtained by applying the fibre functors. It is contained in the total equivalences for maps between isofibrations, and there it agrees with the total equivalences. To obtain a class contained in the total equivalences on all objects, take its intersection with $W$; the agreement on isofibrations remains true. If $G$ is nontrivial, the unit $e\to C_S(e)$ is a total equivalence but is not an equivalence on the strict fibre: there it is the inclusion of the singleton $\{1\}$ into the discrete groupoid $G$. The fibre functor, followed by localization at equivalences of groupoids, detects that this unit cannot become invertible in the fine localization. Thus the fine and coarse localizations differ even though their equivalences between fibrant objects coincide. This is precisely the situation of the reflection theorem.
\end{dcomment}


% foot: Version 9 novembre 2001

\printedpage{67}
% head: A. Grothendieck : Dérivateurs, Ch. VI, Édition des Universités de Montpellier II et Paris VII

La pensée naturelle, c'est qu'avec cette composition et ces identités, les $X(T)$, $\Ch(T)$ deviennent les objets et les flèches d'une \textit{catégorie} - i.e. que $(X,\Ch(X),\alpha_X,p_X,\pi_X)$ définissent une structure de ‘catégorie dans $\mathcal{M}$’. Mais pour ce qu'on veut faire maintenant, on n'a sans doute pas besoin de tout, i.e.
\msp{page 107}
de l'axiome pertinent d'associativité de la composition et celui sur les unités, mais seulement que les $1_x$ soient des unités bilatères pour la composition.

En supossant donné ce supplément de structure sur le foncteur Ch, et revenant à la considération du foncteurs associé % [sic] supossant -> supposant, faute de frappe ; foncteurs -> foncteur, accord avec « du » et « associé »
\[
C : \mathcal{M}/S\longrightarrow\Fib\,S,
\]
on trouve un morphisme canonique
\[
C(C(X))\xrightarrow{\gamma_X}C(X)
\]
par
\[
(x,c,c')\longmapsto(x,c'\circ c),
\]
et l'axiome des unités bilatères nous assure aussitôt que
\[
(\ast)\qquad\qquad\gamma_X\circ i_{C(X)}=\gamma_X\circ C(i_X)=\mathrm{id}_{C(X)}.
\]
Comme $i_{C(X)}\in W$, cela prouve déjà que $\gamma_X\in W$, donc $\gamma_X\in W_S$. Donc $\gamma_X$ devient un isomorphisme dans $(\Fib\,S)W_S^{-1}$, et la relation $(\ast)$ implique bien $i_{C(X)}=C(i_X)$ dans le localisé, q.e.d.

Je vais à nouveau résumer ce que j'ai dégagé.

\msp{page 108}
\textbf{Théorème-Scholie 2.} Soit $\mathcal{M}$ une catégorie, munie de deux ensembles de flèches
\[
W,\ \mathrm{Fib}\subseteq\mathrm{Fl}(\mathcal{M}).
\]
On fait les hypothèses suivantes.

\dgrblock{dgrb219}{What composition contributes to the reflection}
\begin{dcomment}{exposition}{p107}{What composition contributes to the reflection}
Grothendieck asks for a composition of paths because it contracts the two-step replacement back to a one-step replacement. For this purpose there is no triple of paths to compare. Associativity therefore has no role in the proof: the only identities used are $c\circ1=c$ and $1\circ c=c$. This explains the deliberately weaker formulation on pp.~106--107.

With the corrected order of the two unit insertions, the formula
\[
 \gamma_X(x,c,c')=(x,c'\circ c)
\]
gives a common left inverse of $i_{CX}$ and $C(i_X)$. Two-out-of-three makes $\gamma_X$ a weak equivalence, because $\gamma_Xi_{CX}=\mathrm{id}$ and $i_{CX}\in W$. Both its source and target are fibrant over $S$, so Hypothesis~1 also makes it a fine weak equivalence. Its inverse in the fine localization is therefore both $i_{CX}$ and $C(i_X)$. The path calculation of manuscript p.~69 has now been isolated from its construction in $\Cat$.

An internal category would require associative composition in addition to the endpoint and unit identities, and its set of arrows on a test object $T$ would be $P(X)(T)$. The edition's $\mathrm{Ch}(T)$ in the paragraph spanning pp.~106--107 is a type error: $\mathrm{Ch}(T)$ is an object of $\mathcal M$, whereas the discussion concerns a set of arrows between elements of $X(T)$. The intended expression is $\mathrm{Ch}(X)(T)$.

There is also a shorter route once the comparison of localizations has been proved. That comparison can transfer the required equality from the coarse localization back to the localized fibrant subcategory. We give the argument at p.~111. It shows that the path composition is a useful geometric witness, while the formal reflection theorem itself needs only the functorial replacement of A. This strengthens the stated theorem without changing the validity of the path proof.
\end{dcomment}


\begin{enumerate}[label={}]
\item[$1^o$] $W$ contient les isomorphismes, si $f$, $g$ sont composables, si deux parmi $f$, $g$, $gf$ sont dans $W$, le troisième aussi.
\item[$2^o$] $\mathcal{M}$ est stable par $\varprojlim$ finies, Fib contient les isomorphismes, est stable par composition et par changement de base (\footnote[59]{Ces hypothèses de stabilité sur Fib ne serviront que dans II, et encore seulement dans le cas A), repris in toto sous la forme A$'$ plus bas (p.~112).}).
\end{enumerate}

Pour $S$ dans $\mathcal{M}$, \fbox{$\Fib\,S$} désigne la sous-catégorie pleine de $\mathcal{M}/S$ formée des $X\longrightarrow S$ qui sont dans Fib. On distingue deux versions différentes d'un axiome de factorisation d'une $f : X\longrightarrow S$ quelconque dans $\mathcal{M}$ en $X\xrightarrow{i}\bar X\xrightarrow{\bar f}S$, avec $i\in W$, $\bar f\in\mathrm{Fib}$.

% foot: Version 9 novembre 2001

\printedpage{68}
% head: A. Grothendieck : Dérivateurs, Ch. VI, Édition des Universités de Montpellier II et Paris VII

\begin{enumerate}[label={}]
\item[A)] On se donne, pour tout $S$ dans $\mathcal{M}$, un foncteur
\[
\mathcal{M}/S\longrightarrow\mathcal{M}/S\qquad\qquad X\longmapsto C(X)\ \ (=\bar X)
\]
et un morphisme de foncteurs en $\mathrm{X}$
\[
X\xrightarrow{i_X}C(X)
\]
avec
\[
i_X\in W,\qquad\qquad C(X)\in\mathrm{Ob}\,\Fib\,S.
\]
\item[B)] Désignant par
\[
\mathrm{Cof}_0\subseteq\mathrm{Fl}\,\mathcal{M}
\]
l'ensemble des flèches de $\mathcal{M}$ qui possèdent la propriété de relèvement à droite \texttt{[plutôt à gauche]} par rapport à Fib, on suppose ‘à la \textsc{Quillen}’, que toute flèche $f : X\longrightarrow S$ se factorise en
\[
X\xrightarrow{i_X}C(X)\xrightarrow{\bar f}S,
\]
avec
\[
i\in\mathrm{Cofib}_0\cap W,\qquad\qquad\bar f\in\mathrm{Fib}.
\]
\texttt{[où} $\mathrm{Cof}_0=\mathrm{Cofib}_0$\texttt{]}.
\end{enumerate}

\msp{page 109}
Donnons nous $S$ dans $\mathcal{M}$, considérons
\[
W_S\subseteq\mathrm{Fl}(\mathcal{M}/S)
\]
l'ensemble de flèches induit par $W$ (
\begin{tikzcd}[column sep=small, row sep=small, baseline={([yshift=.5ex]current bounding box.south)}]
f : X \arrow[rr] \arrow[dr] & & Y \arrow[dl] \\
 & S &
\end{tikzcd}
est dans $W_S$ si et seulement si il l'est en oubliant $S$), et désignons par $W_S$ (par abus de notation) l'ensemble de flèches induit par $W_S$ sur $\Fib\,S$, ou plus généralement sur toute sous-catégorie pleine $\mathcal{F}$ de $\mathcal{M}/S$, avec
\[
\Fib\,S\subseteq\mathcal{F}\subseteq\mathcal{M}/S.
\]

\dgrblock{dgrb220}{Reading the two versions of the abstract theorem}
\begin{dcomment}{exposition}{p108}{Reading the two versions of the abstract theorem}
Theorem-Scholie~2 organizes the preceding work around two independent questions. The first is whether fibrant representatives compute the coarse localization. The second is whether those representatives form a reflective subcategory of a finer localization. Keeping the two questions separate makes the scope of each hypothesis transparent.

For I under A, we need a full subcategory $\mathcal F\subseteq\mathcal A$, a functor $C:\mathcal A\to\mathcal F$, and a natural transformation $\mathrm{id}_{\mathcal A}\to C$ whose components are in $W$. Two-out-of-three shows that $C$ preserves $W$. Restricting the transformation to $\mathcal F$ gives the second comparison with the identity, and hence an equivalence of localizations. Finite limits, fibrations, and lifting properties are irrelevant to this formal argument. Their role is to provide the replacement in particular examples.

Under B, the functor only appears after localization. The left lifting property supplies its arrows, and the path-object lemma supplies independence. In the proof as written, finite limits and stability of $\mathrm{Fib}$ under pullback and composition are used at pp.~92--93 to keep the relative path object fibrant over $S$. Footnote~59 says that these stability hypotheses occur only in II, and only in A$'$. That description of the proof is incorrect. It does not affect the theorem, because the hypotheses are included in its statement.

The condition in II is equally precise:
\[
 W_0\subseteq W,
 \qquad W_0\cap\mathrm{Fl}(\mathcal F)
       =W\cap\mathrm{Fl}(\mathcal F).
\]
It is an equality on arrows between good objects, not on all arrows with one good endpoint. In particular, the replacement map $X\to CX$ need not belong to $W_0$. The resulting unit in the fine localization is generally noninvertible, while its image in the coarse localization is invertible. That asymmetry is exactly what permits a nontrivial reflection.

The term ``localement constant'' in the abstract theorem is motivated by the application to $\Cat/S$. For an arbitrary relative category $\mathcal M$, the localization $(\mathcal M/S)[W_S^{-1}]$ has been defined without a notion of transport along arrows of $S$. Its identification with local systems is an additional interpretation supplied by a model of homotopy types. The formal theorem establishes the reflection; it does not by itself define monodromy or a theory of locally constant diagrams.
\end{dcomment}


\circled{I} Le foncteur
\[
(\ast)\qquad\qquad(\Fib\,S)W_S^{-1}\xrightarrow{\sim}(\mathcal{M}/S)W_S^{-1}
\]
induit par l'inclusion est une équivalence de catégories, plus généralement pour toute $\mathcal{F}$ comme dessus
\[
(\Fib\,S)W_S^{-1}\longrightarrow\mathcal{F}W_S^{-1}\longrightarrow(\mathcal{M}/S)W_S^{-1}
\]
sont des équivalences de catégories. Un foncteur quasi-inverse est obtenu (dans les deux cas A) et B)) à l'aide du foncteur
\[
\gamma : \mathcal{M}/S\longrightarrow(\Fib\,S)W_S^{-1}
\]

% foot: Version 9 novembre 2001

% ---- from Der-ChVI-345_p071-075.tex
\setcounter{footnote}{59}

\printedpage{69}
% head: A. Grothendieck : Dérivateurs, Ch. VI, Édition des Universités de Montpellier II et Paris VII
qui est donné sur les objets par
\[\gamma(X) = C(X),\]
et qui sur les flèches se définit de façon évidente dans le cas A) (fonctorialité de $C(X)$ en $X$), et dans le cas de B) se construit ad hoc, comme la seule loi fonctorielle pour laquelle les homomorphismes
\[i_X : X \longrightarrow C(X)\]
\msp{page 110}
(interprétés comme flèches dans $(\M/S)W_S^{-1}$) soient eux-mêmes fonctoriels en $X$.

On désigne par
\[\Ho^{\mathrm{lc}}_W(S) \qquad\qquad \text{(lc comme ‘localement constant’)}\]
la catégorie des fractions $(\M/S)W_S^{-1}$, équivalente à $(\Fib S)W_S^{-1}$.

\circled{II} Donnons nous un localiseur $W_S^0$ dans $\M/S$
\[W_S^0 \subseteq W_S \subseteq \Fl(\M/S),\]
avec l'hypothèse formelle similaire à celle de $W$ (cf. 1$^o$). On suppose que

\begin{itemize}[leftmargin=4em]
\item[3$^o$)] $W_S^0 \cap \Fl(\Fib S) = W_S|\Fib S$.
\end{itemize}

(Cf. pages 94-96 pour une façon générale d'obtenir des tels localiseurs. Il faut bien voir que pour définir des ‘types d'homotopie’ relatifs sur $S$ le localiseur $W_S$ est beaucoup trop grossier, moralement il ne donne qu'une catégorie des types d'homotopie relatifs localement constants, et c'est justement ça qui est montré par l'equivalence de \circled{I} , où $\Fib S$ est visualisé comme formé des ‘fibrations localement triviales’ sur $S$.) On a alors un deuxième localisé % [sic] l'equivalence -> l'équivalence, accent manquant
\[\Ho_{W_S^0}(S) \overset{\text{déf}}{=} (\M/S)(W_S^0)^{-1},\]
et deux foncteurs
\[
\begin{tikzcd}[column sep=large]
\Ho_{W_S^0}(S) \arrow[r, shift left, "\rho"] & \Ho^{\mathrm{lc}}_W(S), \arrow[l, shift left, "i"]
\end{tikzcd}
\tag*{$(*)$}
\]
où $\rho$ est le foncteur de localisation (strict) évident, dû à l'inclusion $W_S^0 \subseteq W_S$, et où $i$ est
\msp{page 111}
essentiellement le foncteur déduit de l'inclusion
\[(\Fib S, \underbrace{W_S|\Fib S}_{=W_S^0|\Fib S}) \longrightarrow (\M/S, W_S^0)\]
par localisation \textit{au moyen de $W_S^0$}, donc le foncteur
\[(\Fib S)(W_S^0)^{-1} \longrightarrow (\M/S)(W_S^0)^{-1}, \tag*{$(**)$}\]
% foot: Version 9 novembre 2001
\printedpage{70}
% head: A. Grothendieck : Dérivateurs, Ch. VI, Édition des Universités de Montpellier II et Paris VII
compte tenu de l'équivalence $(*)$ de \circled{I} , i.e. $i$ est le composé du quasi-inverse de $(*)$, avec le foncteur $(**)$ précédent. Explicitement, il est défini par
\[i(X) = \text{classe de C(X) dans } (\M/S)(W_S^0)^{-1},\]
en tenant compte du fait que $X \longmapsto C(X)$ transforme $W_S$ en $W_S^o$, grâce à la condition 3$^o$) sur $W_S^0$. % [sic] W_S^o -> W_S^0, lettre o au lieu du chiffre 0 utilisé partout ailleurs

Ceci posé, \textit{sous l'hypothèse} B (‘Quillen’), on a : dans le couple de foncteurs $(\rho, i)$ de $(*)$, page 110, \textit{$i$ est adjoint à droite de $\rho$,} i.e. on a pour $\xi \in \mathrm{Ob}\,\Ho_{W_S^0}(S)$, $\eta \in \mathrm{Ob}\,\Ho^{\mathrm{lc}}_W(S)$
\[\Hom_{W_S^0}(\xi, i(\eta)) \simeq \Hom_{W_S}(\rho(\xi), \eta),\]
ce qui s'explicite, pour $X$, $Y$ dans $\M/S$, par
\[\boxed{\Hom_{W_S^0}(X, C(Y)) \simeq \Hom_{W_S}(X, Y)}\]
(\footnote{\textbf{NB} Le premier membre est indépendant du choix de $W_S^0$! (Pourvu seulement que l'hypothèse $W_S^0 \cap \Fl\,\Fib S = W_S \cap \Fl\,\Fib S$, en plus de $W_S^0 \subseteq W_S$, soit satisfaite.)}). Il s'ensuit que \textit{le foncteur $i$ est pleinement fidèle.}

La même formule d'adjonction, et la
\msp{page 112}
pleine fidélité de $i$, sont valables dans le cas A, quand on renforce l'hypothèse A (relative à un $S$ particulier) en une hypothèse plus restrictive et plus compréhensive, que voici.

\dgrblock{dgrb221}{Functorial replacement already gives the reflection}
\begin{dcomment}{later}{p111}{Functorial replacement already gives the reflection}
Grothendieck's construction has supplied more than the final formulation uses. Hypothesis A alone suffices for II; the additional path composition in A$'$ is unnecessary for the adjunction. The missing observation is that I has already proved full faithfulness in the coarse localization. We can use that result to establish coherence without cancelling a noninvertible unit in the fine localization.

Here is the precise argument. Let $\mathcal F\subseteq\mathcal A$ be full, let $W_0\subseteq W$ contain the isomorphisms and satisfy two-out-of-three, and assume that they agree on $\mathcal F$. Suppose given a functor $C:\mathcal A\to\mathcal F$ and a natural transformation $\eta:\mathrm{id}_{\mathcal A}\to C$ with $\eta_X\in W$. Put
\[
 \mathcal D=\mathcal A[W_0^{-1}],\qquad
 \mathcal B=\mathcal A[W^{-1}],\qquad
 \mathcal K=\mathcal F[W^{-1}].
\]
By I, the inclusion $J:\mathcal K\to\mathcal B$ is an equivalence. The inclusion into the fine localization gives $k:\mathcal K\to\mathcal D$. The functor $C$ induces $q:\mathcal D\to\mathcal K$, since it carries $W_0\subseteq W$ into $W$ between objects of $\mathcal F$. The transformation $\eta$ descends to $\mathrm{id}_{\mathcal D}\to kq$.

We first compare the iterated units in $\mathcal K$. Naturality in $\mathcal A$ gives
\[
 C(\eta_X)\eta_X=\eta_{CX}\eta_X.
\]
In $\mathcal B$, $\eta_X$ is invertible, so $C(\eta_X)=\eta_{CX}$. Both arrows are arrows of $\mathcal F$. Since $J$ is faithful, their classes are equal in $\mathcal K$ itself. There is no appeal here to the desired full faithfulness of $k$; the faithful functor is the coarse comparison $J$, whose equivalence was established before $W_0$ entered the argument.

For $Y\in\mathcal F$, the arrow $\eta_Y:Y\to CY$ is a weak equivalence in $\mathcal F$. Define a natural counit $\varepsilon:qk\to\mathrm{id}_{\mathcal K}$ by $\varepsilon_Y=[\eta_Y]^{-1}$. Naturality follows by inverting the natural isomorphism induced by $\eta|_{\mathcal F}$. The triangle identities are now explicit:
\[
 k(\varepsilon_Y)\eta_{kY}=\mathrm{id}_{kY},
 \qquad
 \varepsilon_{qX}\,q(\eta_X)
   =[\eta_{CX}]^{-1}[C(\eta_X)]
   =\mathrm{id}_{qX}.
\]
Thus $q\dashv k$. Its counit is an isomorphism, so $k$ is fully faithful. Composing on the left with the equivalence $J$ identifies $q$ with the canonical localization $\rho:\mathcal D\to\mathcal B$ and gives
\[
 \rho\dashv kJ^{-1}.
\]
The functor $kJ^{-1}$ is Grothendieck's $i$. In particular,
\[
 \Hom_{\mathcal D}(X,CY)
 \simeq \Hom_{\mathcal B}(X,Y)
\]
under A without A$'$, and for every choice of $W_0$ with the stated agreement property.

The bijection can also be written without invoking the adjunction as a black box. For $Y\in\mathcal K$, it sends $f:X\to kY$ to $\varepsilon_Yq(f):qX\to Y$. Its inverse sends $h:qX\to Y$ to $k(h)\eta_X$. Substituting one formula into the other, naturality of $\eta$ or $\varepsilon$ moves the comparison map past $f$ or $h$, and the corresponding triangle identity cancels the remaining pair. These formulas establish both naturality variables and show that the comparison with the coarse localization is the canonical one. In particular, no new choice of a map from $X$ into a fibrant target has been made.

This strengthening belongs to the general theory of derived adjunctions, rather than being a new phenomenon. Apply the deformable-adjunction theorem, \citet[Theorem~2.2.11]{Riehl2014}, to the identity adjunction between $(\mathcal A,W_0)$ and $(\mathcal A,W)$: its left functor is already homotopical, and its right functor admits the right deformation $C$, on whose image the two classes agree. The proof above is included to make the application independent of a stronger axiomatization of homotopical categories, and to show precisely how the reflection lemma of manuscript p.~67 receives its coherence condition here.
\end{dcomment}


\begin{itemize}[leftmargin=3em]
\item[A$'$)] Le foncteur $C : \M/S \longrightarrow \M/S$ et le morphisme fonctoriel $i_X : X \longrightarrow C(X)$ sont construits ‘à la \textsc{Cartan-Serre}’ à partir d'un foncteur défini sur $\M$ entier
\[\M \xrightarrow{\mathrm{Ch}} \M,\]
du type ‘chemins’, i.e. muni en plus des données suivantes
\begin{itemize}[leftmargin=3em]
\item[a)] \[X \xrightarrow{\alpha_X} \mathrm{Ch}(X) \xrightarrow{p_X = (\sigma_X, \beta_X)} X \times X\]
(structure de source et but des chemins, et de chemins triviaux), et
\item[b)] Une composition des chemins, sans exiger ici l'associativité de la composition mais seulement que les chemins triviaux soient unités bilatères pour la composition. De plus, on suppose que pout tout $X$ % [sic] pout -> pour, faute de frappe
\[\begin{cases} \alpha_X \in W,\ p_X \in \mathrm{Fib} \\ (X \longrightarrow e) \in \mathrm{Fib}. \end{cases}\]
\end{itemize}
\end{itemize}
% foot: Version 9 novembre 2001

\printedpage{71}
% head: A. Grothendieck : Dérivateurs, Ch. VI, Édition des Universités de Montpellier II et Paris VII
\msp{page 113}
\textbf{Corollaire.} Soient $X$ dans $\M/S$, $Y$ dans $\Fib S$, alors
\[\Hom_{W_S^0}(X, Y) \xrightarrow{\sim} \Hom_{W_S}(X, Y).\]
\textsc{Démonstration}. Par la formule d'adjonction (p. 111), le deuxième membre est $\Hom_{W_S^0}(X, CY)$. Mais comme $Y$ et $C(Y)$ sont dans $\Fib S$, où $W_S$ et $W_S^0$ coïncident, $Y \xrightarrow{i_Y} C(Y)$ est non seulement dans $W_S$, mais dans $W_S^0$, donc $\Hom_{W_S^0}(X, Y) \xrightarrow{\sim} \Hom_{W_S^0}(X, CY)$, d'où le résultat.

Revenons à Cat, Cat$/S$. La digression précédente sur $(\M, W, \mathrm{Fib})$ a permis de mieux situer ce qui concerne le foncteur $C$, et les relations entre les catégories localisées de
\[\Cat/S \ \ (\text{pour } W_S,\ W_S^{\mathrm{d}},\ W_S^{\mathrm{g}}) \qquad\quad \text{et} \qquad\quad \Parf S \ \ (\text{pour } W_S^{\mathrm{f}} = W_S = W_S^{\mathrm{d}} = W_S^{\mathrm{g}})\]
ainsi que pour les catégories intermédiaires
\[\Parf S \subseteq \mathcal{F} \subseteq \Cat/S\]
pleines de Cat$/S$ (sur le diagramme p. 76, il y en a cinq autres, savoir $\Fibhot S$, $\Liss\,S$, $\Liss_0 S$, $\Prop\,S$, $\Prop_0\,S$). Pour ce qui concerne les relations entre les localisés de ces sept catégories, qui se trouvent au centre du diagramme, les énoncés du théorème-scholie 1 (p. 70-81) ne donnent rien de plus que le contenu du théorème-scholie 2 (p. 108-113). Pour les établir, on a donc le choix entre les deux hypothèses A) ou B) dudit théorème. Soit utiliser un foncteur $C_S$ explicite, à l'aide d'un
\msp{page 114}
foncteur $X \longmapsto \underline{\mathrm{Ch}}(X)$ (‘chemins dans $X$’) sur Cat, donc les structures de source, but et composition pour les chemins, et les chemins triviaux (agissant comme unités bilatères pour la composition). La construction d'un tel foncteur $C$ n'a par l'air trivial et semble ne pas être connue - j'ai peiné pour y arriver! (Avec, en plus, l'associativité de la composition \ldots). Ou bien utiliser, à la \textsc{Quillen}, l'existence d'une classe de morphismes dits Q-fibrants, avec des bonnes propriétés de factorisation. Ce n'est pas non plus évident dans Cat, et c'est établi par \textsc{Thomason}, par une voie un peu compliquée. Peut-être serai-je amené à revenir dessus - mais seulement en cas de besoin! % [sic] n'a par l'air -> n'a pas l'air, faute de frappe

\dgrblock{dgrb222}{Different fine localizers can have the same fibrant targets}
\begin{dcomment}{examples}{p113}{Different fine localizers can have the same fibrant targets}
Grothendieck's corollary says more than full faithfulness on pairs of fibrant objects. The source $X$ is arbitrary: it is only the \emph{target} $Y$ that must be fibrant. The canonical map of Hom sets is bijective because $Y\to CY$ is already inverted in the finer localization, and the reflection computes maps into $CY$. This is why the Hom set in footnote~60 is independent of the chosen fine localizer.

A finite-set example isolates the statement from homotopical geometry. Let $\mathcal A$ be a skeleton of finite sets, let $W$ consist of all maps, let $\mathcal F$ consist only of a singleton $1$, and put $C(X)=1$. The unique maps $X\to1$ form a natural transformation in $W$. The coarse localization is equivalent to the terminal category: after all arrows are inverted, $X\to1$ is a natural isomorphism to the constant functor. The full subcategory $\mathcal F$ also localizes to the terminal category.

First take $W_0$ to be the bijections. Its localization is just the category of finite sets. The reflection to the terminal category has right adjoint selecting $1$, since
\[
 \Hom_{\mathcal A}(X,1)=\{*\}
\]
for every finite $X$, including the empty set. It plainly does not follow that the fine localization is terminal: for example $\Hom(1,\{0,1\})$ has two elements.

For a second fine class, let $W_1$ consist of the maps which preserve whether a finite set is empty. Equivalently, $W_1$ contains the identity of the empty set and every map between nonempty finite sets. It is the inverse image of the isomorphisms under the functor
\[
 E:\mathcal A\to[1],\qquad
 E(\varnothing)=0,\quad E(X)=1\quad(X\ne\varnothing),
\]
so it satisfies two-out-of-three. The localization at $W_1$ is equivalent to $[1]$. Indeed, the functor $j:[1]\to\mathcal A$ selecting $\varnothing$ and $1$ satisfies $Ej=\mathrm{id}$, and the natural maps $X\to jE(X)$, identity on the empty set and the unique map to $1$ otherwise, all lie in $W_1$. These exhibit quasi-inverse functors after localization.

Thus the two fine categories are $\mathrm{FinSet}$ and $[1]$, which are very different. Nevertheless, both give a singleton set of maps from any object into the distinguished fibrant target. Both $W_0$ and $W_1$ agree with $W$ on $\mathcal F$, so the corollary predicts exactly this agreement. The example also makes the direction of the result unmistakable: maps \emph{out of} the distinguished object are not fixed. In $\mathrm{FinSet}$ there are two maps $1\to\{0,1\}$; in $[1]$ their images coincide.
\end{dcomment}


Par contre, les assertions concernant les catégories $\Fib S$, $\Fibsc S$, ou $\Cofib S$, $\Cofibsc S$, et leurs variantes à indices 0 (donc les huit catégories-ailes du diagramme), et leurs relations à Cat$/S$ d'une part, à $\Fibhot S$ de l'autre (pour les catégories à indice 0), ne peuvent être contenues dans un tapis général sur des catégories de modèles $(\M, W)$, me semble-t-il, vu la nature très spéciale des quatre catégories désignées ci-dessus, de leur rôle très particulier dans le formalisme catégorique. Or tout au moins les catégories ‘scindées’ $\Fibsc S$, $\Cofibsc S$, et par suite aussi leurs variantes à indice 0, jouent un rôle essentiel pour la théorie des HOT$(S)$, vu que c'est elles qui
\msp{page 115}
interviennent dualement quand on applique la recette générale de construction d'un dérivateur à partir d'une catégorie de modèles $(\M, W)$. Donc je ne vois pas de simplification
% foot: Version 9 novembre 2001
\printedpage{72}
% head: A. Grothendieck : Dérivateurs, Ch. VI, Édition des Universités de Montpellier II et Paris VII
substantielle à apporter au théorème-scholie 1. Tout au plus pourrait-on élaguer le diagramme de catégories envisagées, en laissant tomber $\Liss\,S$, $\Liss_0 S$, $\Prop\,S$, $\Prop_0 S$, qui n'y jouent pas un rôle vraiment utile. Pour ce qui concerne les résultats énoncés, ce qui les concerne est couvert par ce qui est dit pour les sous-catégories pleines $\mathcal{F}$ de Cat$/S$ qui contiennent (disons) $\Parf S$ (ou au choix celles qui contiennent $\Fib S$ ou $\Cofib S$).

\dgrblock{dgrb223}{The model-category route in Cat}
\begin{dcomment}{later}{p114}{The model-category route in Cat}
Grothendieck invokes Thomason to explain that the abstract replacement argument is not restricted to topological or simplicial objects. Thomason's model structure has as weak equivalences the functors whose nerves are weak homotopy equivalences, and as fibrations the functors detected as Kan fibrations by $\mathrm{Ex}^2N$; its existence is Theorem~4.9 of \citet{Thomason1980}. Factorization into a trivial cofibration followed by a fibration supplies hypothesis B. The theorem concerns the ordinary homotopy theory of nerves; it should not be cited as constructing a model structure for an unspecified class $W$.

There is a historical correction relevant to the preceding properness discussion. Thomason's original proof asserted that Dwyer maps are closed under retracts and hence that every cofibration is a Dwyer map. \citet[pp.~228--229]{Cisinski1999Dwyer} gives a counterexample. His Lemma~4 and Proposition~2 on p.~231 replace Dwyer maps by pseudo-Dwyer maps and recover properness. The model-category existence theorem is unaffected. Thus the later correction repairs an auxiliary assertion in the cited literature, not the validity of Grothendieck's abstract comparison theorem.

The two uses of ``fibration'' in this chapter must remain distinct. A Grothendieck fibration over $S$ carries cartesian lifts of arrows and presents a contravariant diagram of categories. A Thomason fibration is defined by a lifting condition after applying $\mathrm{Ex}^2N$. A categorical comparison between their localized presentations is a theorem, not a consequence of the shared terminology. This is the reason the eight outer categories in the large comparison diagram do not disappear merely because one has a model structure on $\Cat$: they encode strict or split diagrammatic data which the abstract factorization axioms have not named.

The return to those categories is therefore mathematically necessary. The abstract detour has made the localizations and reflection intelligible; the split-fibration presentation will now make variation in $S$ intelligible. These are complementary uses of the same models, and neither replaces the other.
\end{dcomment}


La description des catégories du type HOT$(S)$, pour $S$ fixé, me paraît à présent parfaitement bien comprise. Il est temps de faire varier $S$, et de comprendre le caractère contravariant en $S$. On se donne donc
\[S' \xrightarrow{f} S,\]
et on se propose de décrire les foncteurs image inverse pour les trois types de catégories HOT$(S)$. Mais nous allons nous borner à présent aux $\HOT_W(S)$ et $\HOT^{\mathrm{lc}}_W(S)$, puisque tout ce qui concerne $\HOT^o_W(S) = \HOT_W(S^o)$ \texttt{[plutôt} $\HOT_W(S^o)^o$\texttt{, cf. p. 50]} s'en déduira par dualité $S \longmapsto S^o$. Donc il s'agit de décrire
\msp{page 116}
\[f^* : \HOT_W(S) \longrightarrow \HOT_W(S') \tag*{$(*)$}\]
et le foncteur induit
\[f^* : \HOT^{\mathrm{lc}}_W(S) \longrightarrow \HOT^{\mathrm{lc}}_W(S').\]
Quand on décrit les catégories HOT$(S)$, conformément à la recette générale, à l'aide de foncteurs
\[F : S^o \longrightarrow \Cat\]
(préfaisceaux en catégories sur $S$), ou, ce qui revient au même, de catégories scindées, la description de la flèche $(*)$ est évidente : elle revient à la loi
\[\left|\begin{array}{rcl}
F & \longmapsto & F \circ f^o \\[1ex]
\underline{\mathrm{Hom}}(S^o, \Cat) & \longrightarrow & \underline{\mathrm{Hom}}(S'^o, \Cat),
\end{array}\right.\]
en passant aux localisés pour la $W$-équivalence argument par argument. Dans le langage des catégories fibrées, il est bien connu que cette opération triviale sur les préfaisceaux en catégories correspond à l'opération \textit{image inverse} de catégories sur $S$ - laquelle opération passe aux scindages. Cela nous donne donc
\[f^*(\hot_{W,S}(X)) = \hot_{W,S'}(\underbrace{X \times_S S'}_{X'}),\]
pour toute $X$ fibrée scindée sur $S$. Or on sait que le scindage de $X$ ne joue aucun rôle, finalement, dans la description de la classe
\[\hot_{W,S}(X) \in \HOT_W(S)\]
% foot: Version 9 novembre 2001
\printedpage{73}
% head: A. Grothendieck : Dérivateurs, Ch. VI, Édition des Universités de Montpellier II et Paris VII
\msp{page 117}
associée à $X$. On peut oublier le scindage, et même (pour étudier les $\Hom_{\mathrm{Hot}_W(S)}(X, Y)$) travailler avec des $S$-foncteurs \textit{quelconque} entre catégories fibrées sur $S$ - pas la peine même de les supposer cartésiens, voir même compatibles avec des scindages. Cela me paraît une chose très importante pour le formalisme des $\HOT_W(S)$, et très spéciale au cas de la catégorie de modèles Cat. C'est un peu tacite dans le théorème-scholie - mais mes doutes initiaux concernant la catégorie $\underline{\mathrm{Fibsc}}$ elle-même montrent bien que je sentais qu'il y avait là quelque chose de pas évident, ni même (presque) \texttt{[?]} de vraisemblable. Et pourtant l'existence du foncteur $\Phi$ montre qu'il en est bien ainsi.

\dgrblock{dgrb224}{Pullback is precomposition in the split presentation}
\begin{dcomment}{exposition}{p116}{Pullback is precomposition in the split presentation}
Grothendieck returns to split fibrations because a change of base is completely visible in that presentation. For a strict functor $F:S^{\mathrm{op}}\to\Cat$, write $\int_SF$ for its Grothendieck construction. An object is a pair $(s,x)$ with $x\in F(s)$; an arrow $(s,x)\to(t,y)$ consists of $a:s\to t$ and a map $x\to F(a)y$ in $F(s)$. The projection to $S$ is the associated split fibration.

If $f:S'\to S$, there is a canonical isomorphism over $S'$
\[
 S'\times_S\int_SF\ \cong\ \int_{S'}(Ff^{\mathrm{op}}).
\]
On objects it sends $(s',x)$ with $x\in F(f(s'))$ to the same pair. On an arrow $a:s'\to t'$ it uses the same map $x\to F(f(a))y$. Identities and composition agree because $F$ and $f$ are functors. The chosen cartesian lifts are carried to the chosen cartesian lifts, so the isomorphism also respects the splittings.

An objectwise weak equivalence of functors remains objectwise a weak equivalence after precomposition: its component at $s'$ is the original component at $f(s')$. This proves that the operation descends to the relevant localization and gives the formula for $f^*$ on split representatives. If $g:S''\to S'$ is another functor, precomposition satisfies
\[
 (Ff^{\mathrm{op}})g^{\mathrm{op}}
   =F(fg)^{\mathrm{op}}
\]
strictly. After changing to equivalent presentations by arbitrary fibrations or by categories over $S$, these equalities are expressed by the canonical comparison isomorphisms for iterated pullbacks. The coherence comes from the same strict composition law, rather than from a new choice of replacement for each composite.

This paragraph only establishes the formula on the specified good representatives. Extending it to other objects or deciding when an arbitrary pullback computes the derived inverse image requires the base-change results on the following pages. The categorical calculation here explains why Grothendieck keeps the split models in the theory even after showing that their localized category is independent of the splitting.
\end{dcomment}


On peut donc interpréter $f^*$ en termes de $\Fib S$, $\Fib S'$, par le foncteur changement de base
\[f^* : \Fib S \longrightarrow \Fib S', \qquad\quad X \longmapsto X \times_S S'/S',\]
lequel foncteur préserve évidemment la $W$-équivalence par fibres $W_S^{\mathrm{f}}$, $W_{S'}^{\mathrm{f}}$, qu'est celle pertinente ici
\[f^*(W_S^{\mathrm{f}}) \subseteq W_{S'}^{\mathrm{f}},\]
de sorte que $f^*$ induit bien
\[(\Fib S)(W_S^{\mathrm{f}})^{-1} \longrightarrow (\Fib S')(W_S^{\mathrm{f}})^{-1}\]
\texttt{[plutôt} $(\Fib S')(W_{S'}^{\mathrm{f}})^{-1}$\texttt{]},
\msp{page 118}
i.e. le foncteur $(*)$
\[f^* : \HOT_W(S) \longrightarrow \HOT_W(S').\]
D'autre part, $\HOT_W(S)$ peut être décrit aussi comme localisé de Cat$/S$ tout entier, par rapport au localiseur $W_S^{\mathrm{d}}$, lequel sur $\Fib S$, et même sur $\Liss\,S$, coïncide avec $W_S^{\mathrm{f}}$. Mais on se bute ici à la difficulté que le foncteur image inverse
\[f^* : \Cat/S \longrightarrow \Cat/S', \qquad\quad X \longmapsto X \times_S S'/S',\]
n'est \textit{pas} compatible en général (si $S$ n'est isomorphe à la catégorie ponctuelle, pour $S'$ quelconque) avec les localiseurs $W_S^{\mathrm{d}}$, $W_{S'}^{\mathrm{d}}$. On le voit p.ex. en prenant $S' = \{s\}$, et en se rappelant que $W_S^{\mathrm{d}}$ n'implique \textit{pas}, sur Cat$/S$ tout entier, la $W$-équivalence par fibres.

\dgrblock{dgrb225}{Restriction must be calculated on a suitable model}
\begin{dcomment}{exposition}{p117}{Restriction must be calculated on a suitable model}
Grothendieck now makes the base category vary. The comparison theorem has done more than identify several presentations of one category: it provides presentations on which an operation can be performed before localization. For inverse image the useful presentation consists of fibrations, or more generally smooth categories over the base. The calculation then reduces to the familiar operation on fibres.

Write $q_S:\Cat/S\to\HOT_W(S)$ for localization at $W_S^{\mathrm d}$. If $X\to S$ is smooth, the strict fibre $X_s$ computes the same homotopy type as the comma category $s\backslash X$. For $f:S'\to S$ and $s'\in S'$, the strict fibre of $X\times_S S'$ at $s'$ is exactly $X_{f(s')}$. Consequently a map between smooth objects which is a $W$-equivalence on comma categories remains so after this pullback. The comparison with split fibrations identifies the resulting localized operation with restriction of category-valued presheaves. This proves Proposition~1, including its naturality in maps between smooth objects.

The qualification on the objects matters even when the morphism being pulled back already belongs to $W_S^{\mathrm d}$. Its membership says something about all comma fibres; it says nothing by itself about its strict fibres. Pullback along the inclusion of an object extracts a strict fibre. Thus arbitrary pullback cannot act directly on the entire relative category $(\Cat/S,W_S^{\mathrm d})$. The commentary at manuscript p.~42 (``When strict fibres compute a comma fibre'') isolates precisely the comparison required here.

There is also a useful distinction between the arrows allowed in this presentation and the arrows needed to describe transport. For calculating the category $\HOT_W(S)$, all functors over $S$ between fibrations are allowed; Cartesian functors form a smaller presentation which gives the same localization. A chosen cleavage is therefore not part of a relative homotopy type. When a natural transformation between changes of base is considered on ms.~128--130, a cleavage will nevertheless provide a convenient formula for the induced map. Independence of that auxiliary choice is a separate assertion, not an additional restriction on the morphisms of $\HOT_W(S)$.
\end{dcomment}


Ces ennuis n'ont pas lieu sûr $\Liss\,S$, où $W_S^{\mathrm{d}}$ et $W_S^{\mathrm{f}}$ coïncident. On peut donc décrire $f^*$ sur les HOT aussi à l'aide de % [sic] sûr -> sur, accent fautif
\[f^* : \Liss\,S \longrightarrow \Liss\,S', \qquad\quad X \longmapsto X \times_S S'/S',\]
en passant aux localisés
\[f^* : (\Liss\,S)(W_S^{\mathrm{f}})^{-1} \longrightarrow (\Liss\,S')(W_{S'}^{\mathrm{f}})^{-1}.\]
En d'autres termes :

\textbf{Proposition 1.} \textit{Si $X$ dans} Cat$/S$ \textit{est lisse sur $S$, alors on a un isomorphisme canonique}
\[f^*(\hot_{W,S}(X)) \simeq \hot_{W,S'}(X'), \qquad\quad \text{\textit{où} } X' = X \times_S S'/S'.\]
% foot: Version 9 novembre 2001

% ---- from Der-ChVI-345_p076-080.tex
\printedpage{74}
% head: A. Grothendieck : Dérivateurs, Ch. VI, Édition des Universités de Montpellier II et Paris VII
% foot: Version 9 novembre 2001

\msp{page 119}
Une remarque similaire s'applique à $\und{Fibhot}\,S$, où $W_S^{\mathrm{d}} = W_S^{\mathrm{f}}$ également. Le foncteur changement de base induit
\[
\und{Fibhot}\,S \longrightarrow \und{Fibhot}\,S', \qquad X \longmapsto X \times_S S'/S',
\]
compatible avec les localiseurs relatifs pertinents (tous ceux envisagés jusqu'à présent coïncident sur ces catégories, avec les hypothèses faites sur $W$), d'où
\[
f^* : (\und{Fibhot}\,S)W_S^{-1} \longrightarrow (\und{Fibhot}\,S')W_{S'}^{-1}.
\]
Les localisés cette fois ne sont pas les $\HOT_W$, mais les $\HOT_W^{\mathrm{lc}}$, donc on trouve une description sympathique de
\[
f^* : \HOT_W^{\mathrm{lc}}(S) \longrightarrow \HOT_W^{\mathrm{lc}}(S'),
\]
plus avantageuse encore qu'en termes de $\und{Liss}_0S$, $\und{Liss}_0S'$, car s'appliquant à des catégories un peu plus vastes (NB $\und{Liss}_0 \subseteq \und{Fibhot}$). Notons donc le

\textbf{Corollaire 1.} \textit{La formule de la proposition 1 (p. 118) reste valable si $X$ est supposé} Hot-\textit{fibrant sur $S$ (au lieu de lisse sur $S$).} (On trouve alors des types d'homotopie relatifs \textit{localement constants} sur $S$, $S'$.)

On a donc une description de $f^*$ qui s'applique à toutes les catégories du côté gauche
\msp{page 120}
du diagramme \texttt{[page]} 76, à la seule exception de $\Cat/S$ elle-même. Si $X$ est dans $\Cat/S$, sans être ni lisse ni Hot-fibrée sur $S$, comment trouver l'image inverse dans $\HOT_W(S')$? On a le choix entre $\Phi$ et $C$ pour améliorer la situation! On sait que $X$ définit la même classe dans $\HOT_W(S)$ que $\Phi(X)$, $C(X)$, qui sont l'une lisse (et même fibrée scindée), l'autre Hot-fibrée (et même parfaite) sur $S$. Donc

\dgrblock{dgrb226}{The two replacements do different work}
\begin{dcomment}{status}{p120:On a le choix entre $\Phi$ et $C$}{The two replacements do different work}
Grothendieck wants one formula for the inverse image of an arbitrary category over $S$. The comma replacement $\Phi_S$ supplies it. The path replacement $C_S$ supplies a formula for a different object, namely the reflection into locally constant relative types. At this point the text briefly identifies those two tasks. Both the sentence preceding Corollary~2 and its second displayed isomorphism are present in the edition; the discrepancy is not introduced by the transcription.

The distinction is already established in the source. The natural map
\[
 X\longrightarrow\Phi_S(X)
\]
belongs to $W_S^{\mathrm d}$, whereas $X\longrightarrow C_S(X)$ is generally only in $W_S$. Accordingly the valid general formula is
\[
 f^*q_S(X)\simeq q_{S'}\bigl(\Phi_S(X)\times_S S'\bigr).
\]
If $i_S:\HOT_W^{\mathrm{lc}}(S)\to\HOT_W(S)$ denotes the inclusion, the formula involving $C_S$ is
\[
 f^*i_S\rho_S q_S(X)
 \simeq q_{S'}\bigl(C_S(X)\times_S S'\bigr).
\]
Here restriction preserves locally constant objects, so that $f^*i_S$ identifies with $i_{S'}f^*_{\mathrm{lc}}$. These equations become the same only when $q_S(X)$ is already locally constant.

A small calculation makes the error decisive without repeating the examples given in the commentary at manuscript pp.~8 and~32. Take $W=W_\infty$, $S=(0\to1)$ and $X=\{0\}\to S$. The comma fibres of $X$ are a point at $0$ and the empty category at $1$. Its relative type is therefore the presheaf of spaces $(\varnothing\to *)$. The local-constant reflection is the constant point: maps from this presheaf to a locally constant arrow are determined by the value at $0$, and evaluation at $0$ identifies locally constant arrows with spaces. Pullback to $1$ consequently gives $\varnothing$ from $X$ and $*$ from $C_S(X)$. Corollary~2 cannot identify them.

This is a localized mistake in the calculation, not a defect in the comparison theorem. The later discussion on ms.~123--127 correctly distinguishes $\rho_Sq_S(X)$ from $q_S(X)$, and the construction of the adjunction on ms.~139--145 consistently uses $\Phi_S$. Keeping the printed error visible is useful: it shows exactly why the two replacement functors were needed.
\end{dcomment}


\textbf{Corollaire 2.} \textit{Si $X$ dans} $\Cat/S$ \textit{est quelconque, alors}
\[
f^*(\hot_{W,S}(X)) \simeq \hot_{W,S'}(\Phi(X) \times_S S') \simeq \hot_{W,S'}(C(X) \times_S S').
\]

\bigskip

On peut s'attendre, d'autre part, que pour certains changements de base particuliers, la formule de la proposition 1 soit valable sans restriction sur $X$ dans $\Cat/S$. Il doit revenir $\pm$ au même que le foncteur image inverse
\[
X \longmapsto X \times_S S'
\]
soit compatible avec les localiseurs $W_S^{\mathrm{d}}$, $W_{S'}^{\mathrm{d}}$. La condition de $W$-équivalence `colocale sur $S$\,' doit être vue comme une condition de nature `homologique', pour le dérivateur HOT. Elle s'associe, pour un dérivateur quelconque $\mathbf{D}$, à la condition (pour $f : X \longrightarrow Y$ dans $\Cat/S$) que les $s\backslash X \longrightarrow s\backslash Y$ sont des $\mathbf{D}$-équivalences,
\printedpage{75}
% head: A. Grothendieck : Dérivateurs, Ch. VI, Édition des Universités de Montpellier II et Paris VII
% foot: Version 9 novembre 2001
\msp{page 121}
ce qui s'écrit ici élégamment en termes d'homologie relative, par la condition que $\forall\ a \in \mathcal{A} = \mathbf{D}(e)$, la flèche induit par $f$ % [sic] induit -> induite, accord avec « la flèche »
\[
\underbrace{H_\bullet^{\mathbf{D}}(X/S, a_X)}_{\overset{\text{déf}}{=}(p_X)_!(a_X)} \longrightarrow \underbrace{H_\bullet^{\mathbf{D}}(Y/S, a_Y)}_{\overset{\text{déf}}{=}(p_Y)_!(a_Y)}
\]
soit \texttt{[un]} isomorphisme. (Utiliser le `calcul' des faisceaux d'homologie relative $(p_X)_!(a_X)$, $(p_Y)_!(a_Y)$, par la caractérisation des fibres dans Der 4 bis). Or, les changements de base $\mathbf{D}$-\textit{propres} $S' \longrightarrow S$ commutent aux $p_!$, de sorte qu'on trouve le

\textbf{Lemme.} \textit{La notion d'équivalence} $\mathbf{D}$\textit{-colocale sur $S$ dans} $\Cat/S$ \textit{est stable par tout changement de base $S' \longrightarrow S$} $\mathbf{D}$\textit{-propre. (Quel que soit le dérivateur} $\mathbf{D}$\textit{.)}

On se rappelera que les $s\backslash S \longrightarrow S$ sont justement des changements des base propres! On % [sic] rappelera -> rappellera, orthographe; des base -> de base, accord
peut dire que l'équivalence $\mathbf{D}$-colocale peut se décrire en exigeant que $f : X \longrightarrow Y$ soit transformé en $W$-équivalences par certains types de changement de base $S' \longrightarrow S$, en prenant au choix l'un des trois types

\begin{itemize}
\item[a)] Les changements de base $s\backslash S \longrightarrow S$ (ce qui redonne la définition que nous avions prise).
\item[b)] Les changements de base qui sont des co-isomorphismes locaux.
\item[c)] Les changements de base propres.
\item[d)] Les changements de base $\mathbf{D}$-propres.
\end{itemize}

On ne peut appliquer ces résultats qu'au dérivateur
\msp{page 122}
HOT, faute de l'avoir encore construit, et verifié, les axiomes pertinents. On peut poser % [sic] verifié -> vérifié, accent
simplement la

\dgrblock{dgrb227}{A homological test for colocal equivalence}
\begin{dcomment}{exposition}{p121}{A homological test for colocal equivalence}
Grothendieck introduces relative homology here because it packages all the comma-category tests into a single morphism of coherent diagrams. Let $p_X:X\to S$, $p_Y:Y\to S$, and let $u:X\to Y$ be a functor over $S$. For $a\in\mathbf D(e)$, write $a_X$ for the constant coefficient. The map of interest is
\[
 (p_X)_!a_X\longrightarrow(p_Y)_!a_Y.
\]
Its component at $s$ is the map on homotopy colimits of the constant coefficient over $s\backslash X$ and $s\backslash Y$. This is the pointwise Kan-extension formula in the contravariant convention used throughout these pages. If isomorphisms are detected by evaluation, requiring the displayed map to be invertible for every $a$ is exactly the condition that all these comma-category maps be $\mathbf D$-equivalences.

Now suppose $g:T\to S$ is $\mathbf D$-proper in the sense that the relevant base-change comparisons for homological direct image are invertible. Pulling back the displayed map gives
\[
 (p_{X_T})_!a_{X_T}\longrightarrow(p_{Y_T})_!a_{Y_T},
 \qquad X_T=X\times_S T,\quad Y_T=Y\times_S T.
\]
The first map was invertible, and restriction preserves isomorphisms; hence the new map is invertible. Evaluating at every object of $T$ gives the claimed stability of colocal equivalence. This proves the lemma without assuming that strict fibres compute comma fibres: the base-change theorem supplies that comparison at the level of relative homology.

The four classes of test maps listed in the text are genuinely nested. A representable coslice projection is a discrete opfibration, hence a co-local isomorphism; those maps are proper, and absolute properness implies $\mathbf D$-properness. Testing against the largest class therefore implies testing against the first. Conversely the lemma shows that the first condition survives the largest class, and a colocal equivalence is in particular a global $\mathbf D$-equivalence. The source's phrase ``three types'' precedes four entries and has no mathematical consequence.

The historical terminology keeps the geometrical motivation visible. Properness and smoothness were designed to distinguish the two positions in a base-change square, following the analogous cohomological theorems in algebraic geometry. For the categorical definitions, including their dual formulations by taking mates, see \citet[Definition~1.1]{Maltsiniotis2005Asphericite}. Here the lesson is particularly concrete: the general derivator theorem predicts the correct class of changes of base, but cannot yet be used to establish the axioms of the very derivator under construction.
\end{dcomment}


\textbf{Conjecture} (sûre!) : \textit{Soit $p : S' \longrightarrow S$ une flèche propre (sous-entendu :} HOT\textit{-propre) dans} Cat. \textit{Alors le foncteur changement de base}
\[
\Cat/S \longrightarrow \Cat/S', \qquad X \longmapsto X \times_S S'/S'
\]
\textit{est compatible avec les localiseurs $W_S^{\mathrm{d}}$, $W_{S'}^{\mathrm{d}}$. (Pour un foncteur $S' \longrightarrow S$ lisse, le changement de base serait compatible avec $W_S^{\mathrm{g}}$, $W_{S'}^{\mathrm{g}}$.)} (\footnote[61]{NB Cette `conjecture' est prouvé sans lemme à la page 168-169, sans avoir besoin de rien qui ne soit % [sic] prouvé -> prouvée, accord avec « conjecture »
déjà connu ici. Cf. théorème 4, cas b) p. 162.}).

Mais pout tout changement de base $S' \longrightarrow S$ (propre ou non) qui a la propriété précédente, % [sic] pout -> pour, faute de frappe
donc qui induit un foncteur
\[
\underbrace{(\Cat/S)(W_S^{\mathrm{d}})^{-1}}_{\approx\,\HOT_W(S)} \longrightarrow \underbrace{(\Cat/S')(W_{S'}^{\mathrm{d}})^{-1}}_{\approx\,\HOT_W(S')}\ ,
\]
\printedpage{76}
% head: A. Grothendieck : Dérivateurs, Ch. VI, Édition des Universités de Montpellier II et Paris VII
% foot: Version 9 novembre 2001
il est immédiat, à l'aide du théorème de comparaison (p.ex. avec $\und{Liss}\,S$, $\und{Liss}\,S'$) sur $S$ et sur $S'$, que ce foncteur n'est autre (à isomorphisme canonique près) que le foncteur $f^*$ image inverse des types d'homotopie relatifs. Donc énoncons le % [sic] énoncons -> énonçons, cédille

\textbf{Corollaire 3.} \textit{La formule de la proposition 1 reste valable sans hypothèse aucune sur $X$ dans} $\Cat/S$\textit{, pourvu que la flèche de changement de base $f : S' \longrightarrow S$ soit telle que le foncteur $X \longmapsto X \times_S S'$ de} $\Cat/S$ \textit{dans} $\Cat/S'$
\msp{page 123}
\textit{transforme $W$-équivalences colocales sur $S$ en $W$-équivalences colocales sur $S'$.} (Ce qui devrait être le cas si (et seulement si?) $f$ est $W$-propre, donc aussi si $f$ est (`absolument') propre (\footnote[62]{cf. I, pages 55-59.}).)

\dgrblock{dgrb228}{The proper base-change conjecture has an elementary proof}
\begin{dcomment}{later}{p122}{The proper base-change conjecture has an elementary proof}
The conjecture on ms.~122 is not an open question. Grothendieck himself proves it later in Chapter~VI, ms.~168--169, printed p.~103 of the complete edition; the footnote points to that proof. The argument avoids the circular appeal to a derivator whose axioms are still being established \citep[ms.~168--169]{Grothendieck1990VI}.

For clarity, isolate its mechanism. If $v:T\to U$ is a universal $W$-equivalence and $a:A\to B$ is a map over $U$, there is a pullback square
\[
\begin{tikzcd}
A\times_U T\arrow[r,"a_T"]\arrow[d] & B\times_U T\arrow[d]\\
A\arrow[r,"a"] & B.
\end{tikzcd}
\]
Both vertical maps belong to $W$. The two-out-of-three property consequently gives $a_T\in W$ if and only if $a\in W$.

Apply this with $g:S'\to S$ proper and an object $s'\in S'$, putting $s=g(s')$. The properness criterion says that
\[
 s'\backslash S'\longrightarrow s\backslash S
\]
is a universal homotopy equivalence; in the relative variant it is a universal $W$-equivalence. A map in $W_S^{\mathrm d}$ is a $W$-equivalence after pullback to $s\backslash S$. The preceding square shows that it remains so after further pullback to $s'\backslash S'$. This is exactly the definition of its pullback belonging to $W_{S'}^{\mathrm d}$. Passing to opposite categories gives the smooth, left-fine version.

Corollary~3 therefore gives the expected strict formula for proper changes of base. Its proof does not depend on choosing a preferred representative: the functor defined on all of $\Cat/S$ agrees, on the smooth presentation, with the previously constructed restriction functor. The equivalence of the two presentations then identifies their descents. The later proof establishes the sufficient condition. The converse requires the additional test calculation given below, and its formulation depends on the chosen localizer.
\end{dcomment}


On n'a travaillé jusqu'ici qu'avec les localiseurs $W_S^{\mathrm{d}}$, $W_{S'}^{\mathrm{d}}$, à l'exclusion de $W_S$, $W_{S'}$. (Sauf sur les catégories contenues dans $\und{Fibhot}\,S$, où les deux coïncident.) Quand on ne fait pas d'hypothèse draconienne sur $p : S' \longrightarrow S$, du type Hot-fibration, on n'a aucune chance pour que $X \longmapsto X \times_S S'$ transforme $W$-équivalences grossières dans itou, donc aucune chance d'attraper l'image inverse du type d'homotopie relatif `grossier' de $X$, en prenant celui de $X' = X \times_S S'$. Cela est lié au fait que pour un $S' \longrightarrow S$ général, les foncteurs de localisation
\[
\HOT_W(S) \xrightarrow{\rho_S} \HOT_W^{\mathrm{lc}}(S), \qquad \HOT_W(S') \xrightarrow{\rho_{S'}} \HOT_W^{\mathrm{lc}}(S')
\]
n'ont aucune envie de commuter au foncteur image inverse $f^*$ sur les $\HOT_W$. On peut hasarder la

\textbf{Proposition 2.} \textit{Conditions équivalentes sur le foncteur $S' \longrightarrow S$ :}

\begin{itemize}
\item[\textit{a)}] \textit{Le diagramme de foncteurs}
\[
\begin{tikzcd}
\HOT_W(S) \arrow[r, "\rho_S"] \arrow[d, "f^*"] & \HOT_W^{\mathrm{lc}}(S) \arrow[d, "f^*_{\mathrm{lc}}"] \\
\HOT_W(S') \arrow[r, "\rho_{S'}"] & \HOT_W^{\mathrm{lc}}(S')
\end{tikzcd}
\]
\textit{commute (à isomorphisme de foncteurs composés près).}

\msp{page 124}

\item[\textit{b)}] \textit{Le foncteur changment de base} % [sic] changment -> changement, faute de frappe
\[
\und{Liss}\,S \longrightarrow \und{Liss}\,S', \qquad X \longmapsto X \times_S S'/S'
\]
\textit{transforme $W_S$ en $W_{S'}$.}
\printedpage{77}
% head: A. Grothendieck : Dérivateurs, Ch. VI, Édition des Universités de Montpellier II et Paris VII
% foot: Version 9 novembre 2001
\item[\textit{b$'$)}] \textit{Énoncé analogue pour} $\und{Fib}\,S \longrightarrow \und{Fib}\,S'$\textit{.}
\item[\textit{b$''$)}] \textit{Énoncé analogue pour} $\und{Fibsc}\,S \longrightarrow \und{Fibsc}\,S'$\textit{.}
\end{itemize}

(\footnote[63]{Pas prouvée, seulement que b, b$'$, b$''$ impliquent a). Pour l'implication inverse, il faudrait savoir que toute flèche de $\mathcal{F}'$ qui devient \texttt{[un]} isomorphisme dans $\mathcal{F}'W_{S'}^{-1}$, est dans $W_{S'}$ (saturation forte de $W_{S'}$ dans $\mathcal{F}'$), chose que je ne sais pas. Mais je crois que ce doit être vrai.}, \footnote[64]{\textbf{NB} Si on suppose $p$ propre, on doit pouvoir prendre aussi $\mathcal{F} = \Cat/S$.}).

\textsc{Démonstration.} Si $\mathcal{F}$ est une des trois catégories $\und{Liss}\,S$, $\und{Fib}\,S$ ou $\und{Fibsc}\,S$, le foncteur $\rho_S$ peut se décrire comme le foncteur de localisation canonique
\[
\mathcal{F}(W_S^{\mathrm{d}})^{-1} \longrightarrow \mathcal{F}W_S^{-1}
\]
(ça marcherait pour n'importe quelle sous-catégorie pleine $\mathcal{F}$ de $\Cat/S$ qui contient $\und{Fib}\,S$). Or si b), b$'$) ou b$''$) est satisfait, on trouve un diagramme commutatif de foncteurs induits par le changement de base
\[
\begin{tikzcd}
\mathcal{F}(W_S^{\mathrm{d}})^{-1} \arrow[r, "\rho_S"] \arrow[d, "f^*_{W_S^{\mathrm{d}}}"'] & \mathcal{F}(W_S)^{-1} \arrow[d, "f^*_{W_S}"] \\
\mathcal{F}'(W_{S'}^{\mathrm{d}})^{-1} \arrow[r, "\rho_{S'}"] & \mathcal{F}'(W_{S'})^{-1},
\end{tikzcd}
\]
sous réserve, il est vrai, dans le cas $\mathcal{F} = \Cat/S$, que le foncteur image inverse soit aussi compatible avec $W_S^{\mathrm{d}}$, $W_{S'}^{\mathrm{d}}$ (ce qui est automatique si $\mathcal{F} = \und{Fib}\,S$ ou $\und{Fibsc}\,S$). On sait que $f^*_{W_S^{\mathrm{d}}}$ décrit le foncteur $f^*$ sur les types d'homotopie relatifs, et pour prouver a) comme conséquence
\msp{page 125}
de b) (renforcé), b$'$ ou b$''$, il reste à s'assurer que le foncteur $f^*_{W_S}$ décrit de même le foncteur induit par le précédent sur les types d'homotopie relatifs localement constants. Or on a les inclusions canoniques $\mathcal{F}_0 \subseteq \mathcal{F}$, $\mathcal{F}_0' \subseteq \mathcal{F}'$, et sur $\mathcal{F}_0$, $\mathcal{F}_0'$ \texttt{[on a]} $W_S = W_S^{\mathrm{d}} = W_S^{\mathrm{f}}$, et les catégories sont OK pour décrire $\HOT_W^{\mathrm{lc}}(S)$ et $\HOT_W^{\mathrm{lc}}(S')$, et le foncteur $f^*$ induit entre eux par images inverses (\footnote[65]{NB $(\Cat/S)_0 = \und{Fibhot}\,S$, pour les autres cas, c'est clair.}) (\textbf{NB} Le foncteur changement de base $X \longmapsto X \times_S S'$ est compatible avec les localiseurs sur ces catégories, car c'est trivial sous la forme $W_S^{\mathrm{f}}$.) On a alors, par images inverses $X \longmapsto X \times_S S'$, un diagramme commutatif
\[
\begin{tikzcd}
\mathcal{F}_0(W_S)^{-1} \arrow[r, "\approx"] \arrow[d, "(f^*_{W_S})_0"'] & \mathcal{F}(W_S)^{-1} \arrow[d, "f^*_{W_S}"] \\
\mathcal{F}_0'(W_{S'})^{-1} \arrow[r, "\approx"] & \mathcal{F}'(W_{S'})^{-1},
\end{tikzcd}
\]
\printedpage{78}
% head: A. Grothendieck : Dérivateurs, Ch. VI, Édition des Universités de Montpellier II et Paris VII
% foot: Version 9 novembre 2001
où les flèches horizontales sont des \textit{équivalences} de catégories, par le théorème-scholie 1 \textcircled{I}~. Cela prouve que $f^*_{W_S}$ lui aussi est bon pour calculer les images inverses de types d'homotopie relatifs.

On a prouvé que b, b$'$, b$''$ impliquent chacun a), prouvons que a) les implique. Or a) dit exactement que les images inverses $f^*$ des éléments de $\HOT_W^{\mathrm{lc}}(S)$
\msp{page 126}
sont, à isomorphisme près, obtenus par le foncteur image inverse sur $\mathcal{F}$ tout entier, en interprétant $\HOT_W^{\mathrm{lc}}(S)$ comme $\mathcal{F}W_S^{-1}$. Mais cela implique bien que $\mathcal{F} \longrightarrow \mathcal{F}'$ doit être compatible avec les localiseurs. Mais je charrie - on trouve seulement que si $u : X \longrightarrow Y$ dans $\mathcal{F}$ est dans $W_S$, alors la flèche $u' : X' \longrightarrow Y'$ donne un isomorphisme dans $\mathcal{F}'(W_{S'})^{-1}$. Je ne sais pas, d'autre part, si cela implique que $u'$ est dans $W_{S'}$.

\textbf{Corollaire 1.} \textit{Supposons que $p : S' \longrightarrow S$ soit une} Hot\textit{-fibration. Alors le foncteur image inverse $X \longmapsto X \times_S S'$ de} $\Cat/S$ \textit{dans $\mathit{Cat}/S'$ applique $W_S$ dans $W_{S'}$, donc le foncteur $p^* : \HOT_W(S) \longrightarrow \HOT_W(S')$ commute aux foncteurs de localisation $\rho_S$, $\rho_{S'}$ dans ces catégories.} (\textbf{NB} Il faut faire l'hypothèse adéquate sur $W$, savoir L6 bis page 44. C'est OK pour $W_\infty$.)

Le fait que $X \longmapsto X \times_S S'$ soit compatible avec $W_S$, $W_{S'}$ est une tautologie, compte tenu de l'hypothèse L6 bis. Il ne semble pas qu'il soit aussi compatible à $W_S^{\mathrm{d}}$, $W_{S'}^{\mathrm{d}}$ (je devrais faire un contrexemple, avec bien sûr $S'$ non propre sur $S$). Donc on est sous les conditions de b), qui implique a).

\msp{page 127}

\dgrblock{dgrb229}{Saturation supplies the missing converse}
\begin{dcomment}{exposition}{p126}{Saturation supplies the missing converse}
Grothendieck wants to recognize when restriction commutes with the reflection into locally constant objects. His forward implication is sound: preservation of coarse weak equivalences on any of the indicated smooth or fibred presentations lets the square of localizations descend. For the converse he notices precisely what has not yet been justified. An arrow becoming invertible after localization need not, for an arbitrary chosen class of arrows, have belonged to that class.

There is a short repair once $W$ is strongly saturated in $\Cat$. Let $\mathcal F$ be any category of categories over a fixed base, with any chosen class of allowed morphisms, and let $W_{\mathcal F}$ consist of those allowed morphisms whose underlying functors belong to $W$. The forgetful functor descends to
\[
 \mathcal F[W_{\mathcal F}^{-1}]\longrightarrow\Cat[W^{-1}].
\]
If an actual arrow $u$ of $\mathcal F$ becomes invertible on the left, its underlying functor becomes invertible on the right. Strong saturation gives $u\in W$. The converse is part of localization. Thus the restriction of a strongly saturated coarse class remains strongly saturated, even when $\mathcal F$ is not full.

Assume now that the square in condition~a) commutes. For $u\in W_S$ in a smooth or fibred presentation, its image under $\rho_S$ is an isomorphism. Commutativity says that $f^*u$ becomes an isomorphism in the coarse localization over $S'$. The preceding argument then gives $f^*u\in W_{S'}$. This proves each of b), b$'$), and b$''$). For $W_\infty$, the needed saturation follows immediately by applying the nerve and the homotopy category of simplicial sets. In the later theory it holds for every fundamental localizer; see \citet[Proposition~4.2.4 in the author PDF]{Cisinski2006Presheaves}. The gap belongs to the proof available at this point, not to the resulting equivalence.

Commutation itself is a real restriction. In the interval calculation above, let $f:\{1\}\to(0\to1)$ and let $F=(\varnothing\to*)$. Then
\[
\rho_{\{1\}}f^*F=\varnothing,
 \qquad f^*_{\mathrm{lc}}\rho_S F=*.
\]
This supplies a concrete obstruction without confusing it with the separate failure of arbitrary strict pullback to preserve $W_S^{\mathrm d}$. Grothendieck's Hot-fibration hypothesis is sufficient because it ensures preservation of the coarse class by pullback. Strong saturation closes his converse in Proposition~2; it does not assert that every functor satisfies its equivalent conditions.
\end{dcomment}


\dgrblock{dgrb230}{Properness and a three-object counterexample}
\begin{dcomment}{development}{p126}{Properness and a three-object counterexample}
The question on ms.~123 and the proposed counterexample on ms.~126 can both be settled concretely. The useful characterization of $W$-properness is the dual of \citet[Theorem~1.2(b)]{Maltsiniotis2005Asphericite}. For $f:A\to S$, $a\in A$, and $h:f(a)\to b$ in $S$, let $L_f(a,h)$ have objects arrows $k:a\to a'$ with $f(a')=b$ and $f(k)=h$. A morphism is an arrow $\ell:a'\to a''$ with $f(\ell)=\id_b$ and $\ell k=k'$. Then $f$ is $W$-proper exactly when every $L_f(a,h)$ is $W$-aspherical. These are the outgoing comma categories for the inclusion $A_b\to A/b$.

\textbf{Proposition (detecting properness by inverse image).}\ours
For $W=W_\infty$, strict pullback along a fixed $f:A\to S$ preserves right-fine weak equivalences on all categories over $S$ if and only if $f$ is proper. For an arbitrary fundamental localizer $W$, $W$-properness is equivalent to this preservation property holding after every change of base of $f$.

\textit{Proof.} Properness implies preservation by the preceding block, and properness is stable under base change. To prove the first converse, test the assumed preservation on the unit
\[
 \{b\}\longrightarrow\Phi_S(\{b\})=S/b.
\]
This is right-fine: its comma category at $s$ maps the discrete set $\Hom_S(s,b)$ into the category of factorizations of those arrows, and each factorization category has an initial and a terminal object.

After pullback along $f$, take the comma category at $a$. In the target, objects are triples
\[
 (a',k:a\to a',r:f(a')\to b).
\]
The composite $r f(k)$ defines a decomposition indexed by $h\in\Hom_S(f(a),b)$. The component indexed by $h$ has initial object $(a,\id_a,h)$. The corresponding source component is exactly $L_f(a,h)$. A weak equivalence of these disjoint unions to the contractible target components forces every $L_f(a,h)$ to be weakly contractible. The cited criterion gives properness.

For the general converse, suppose preservation holds after every base change. Given $h:f(a)\to b$, pull back $f$ along the functor $[1]\to S$ representing $h$. Apply the preservation hypothesis to $\{1\}\to[1]$ over $[1]$, which is the same unit test. At the object $(a,0)$ of $A\times_S[1]$, its source comma category is $L_f(a,h)$ and its target comma category has an initial object. Since the map belongs to $W$, two-out-of-three implies $L_f(a,h)\to e\in W$. The criterion again gives $W$-properness.\hfill$\square$

\textbf{Example (why arbitrary localizers require base change).}\ours
The distinction between the two formulations cannot be omitted for arbitrary fundamental localizers. Let $W_{\mathrm{ne}}$ consist of functors whose source and target are either both empty or both nonempty. This is a fundamental localizer: weak saturation and the terminal-object axiom are immediate, and the relative theorem-A axiom follows because an object of a nonempty target gives a nonempty target comma category, hence a nonempty source comma category. Take a nontrivial group $G$ and $f:e\to BG$. Over the one-object category $BG$, every comma category of $X\to BG$ is nonempty exactly when $X$ is nonempty. Right-fine $W_{\mathrm{ne}}$-equivalence therefore detects only nonemptiness. Strict pullback along $f$ does not change the object set of $X$, so it preserves that class. Yet $L_f(*,g)$ is empty for every nonidentity $g\in G$, and $f$ is not $W_{\mathrm{ne}}$-proper. This example concerns the broader localizer question, not a weakening of the stronger hypotheses used elsewhere in this chapter.

\textbf{Example (a Hot-fibration which is not proper).}\ours
Let $E$ be the three-object poset with nonidentity arrows $u\to v$ and $u\to w$, and put
\[
 p:E\longrightarrow(0\to1),\qquad p(u)=p(v)=0,\quad p(w)=1.
\]
This is a Grothendieck fibration. The lift $u\to w$ of $0\to1$ is Cartesian: for an object over $0$, its arrows to $u$ and its arrows to $w$ are both singletons when the object is $u$, and both empty when it is $v$. The two fibres are $(u\to v)$ and $\{w\}$; their transition functor chooses $u$ and is a weak equivalence. Thus $p$ lies in $\Fib_0$, hence in $\Fibhot$ under the standing comparison hypotheses.

There is also a direct verification of its particularly strong homotopical behaviour. The section choosing $u$ over $0$ and $w$ over $1$ has a natural transformation $sp\Rightarrow\id_E$ over the base, whose only nonidentity component is $u\to v$. This deformation retraction survives every strict base change. Every projection $E\times_S T\to T$ is therefore a nerve weak equivalence. Two-out-of-three in a square over $S$ shows that pullback along $p$ preserves all coarse weak equivalences.

Nevertheless $p$ is not proper, since $L_p(v,0\to1)$ is empty. More explicitly, the right-fine weak equivalence $\{1\}\to[1]$ pulls back to $\{w\}\to E$. Its comma functor at $v$ is $\varnothing\to e$, so it is not right-fine. This is precisely the counterexample requested on ms.~126. We checked the properness criterion just cited; the elementary tests and examples here are explicit completions of the manuscript's questions, without a priority claim.
\end{dcomment}


Je ne me sens pas capable d'exhiber un foncteur autre qu'une Hot-fibration (ou tout au moins une $W$-fibration), pour lequel le foncteur image inverse sur $\HOT_W$ soit compatible aux foncteurs de localisation $\rho_S$, $\rho_{S'}$. Bien sûr, si $X$ est quelconque dans $\Cat/S$, rien n'empêche de regarder son localisé pour $W_S$, i.e. $\rho_S(\hot_{W,S}(X))$, et de demander comment trouver son image inverse. Si $p : S' \longrightarrow S$ est une Hot-fibration, alors il suffit de prendre $\hot_{W,S'}(X')$, $X' = S \times_S S'$. Dans tout autre cas (j'entends, sans hypothèse sur $X/S$, plus % [sic] X' = S \times_S S' -> X' = X \times_S S', cf. la définition de X' plus haut
précisément sans supposer que $X$ est lui-même Hot-fibré sur $S$), cette formule est canulée sans espoir. On doit, pour calculer $f^*\rho_S(\hot_{W,S}(X))$, `calculer' $\rho_S(\hot_{W,S}(X))$ en `montant à l'étage', comme $\hot_{W,S}(C(X))$. On trouve alors l'image inverse $\hot_{W,S'}(C(X) \times_S S')$ (par \texttt{[le]} corollaire 1, p. 119).

Maintenant que la loi contravariante de $\HOT_W(S)$ en $S \in \mathrm{Ob}\,\Cat$ semble bien comprise, il s'agit de voir qu'on a bien un dérivateur HOT. Je ne m'attends pas à ce que les $\HOT_W^{\mathrm{lc}}(S)$
\msp{page 128}
définissent eux aussi un dérivateur, mais il faudrait m'en assurer expressément. (C'est pour l'existence de $f_!$, $f_*$ que je doit chercher.) % [sic] je doit -> je dois, conjugaison

À dire vrai, il faudrait, pour bien faire, regarder aussi la structure de 2-\textit{foncteur} de
\[
S \longmapsto \HOT_W(S) \qquad \Cat^o \longrightarrow \Cat,
\]

% ---- from Der-ChVI-345_p081-085.tex
\printedpage{79}
% head: A. Grothendieck : Dérivateurs, Ch. VI, Édition des Universités de Montpellier II et Paris VII
% foot: Version 9 novembre 2001

donc pour $S$, $S'$ données dans Cat, définir
\[ \uHom(S',S)^o \longrightarrow \uHom(\HOT_W(S),\HOT_W(S')). \]
Cela revient surtout, pour
\[ f,g : S' \rightrightarrows S \]
deux flèches de $S'$ dans $S$, et un homomorphisme de foncteurs
\[ u : f \longrightarrow g, \]
à décrire un homomorphisme de foncteurs
\[ u^* : g^* \longrightarrow f^* \]
entre $\HOT_W(S)$ et $\HOT_W(S')$. C'est d'ailleurs chose faite, dans le cadre général des $(\mathcal{M},W)$, quand on décrit $\HOT_W(S)$ et $\HOT_W(S')$ par localisation de $\mathrm{Hom}(S^o,\Cat)$ et $\mathrm{Hom}(S'^o,\Cat)$. Ces catégories s'interprètent respectivement comme $\uFibsc\,S$ et $\uFibsc\,S'$, donc comme des catégories de catégories sur $S$ resp. $S'$, on voudrait donc interpréter $u^*$ en termes de telles catégories, et tant qu'à faire, pour des catégories $X$ sur $S$ aussi générales que possible. Nous nous
\msp{page 129}
attendrons alors à devoir faire des hypothèses, soit sur $X$, soit sur $f$ et sur $g$, qui assurent que les images de $\hot_{W,S}(X)$ par $f^*$ et $g^*$ ne sont autres que $\hot_{W,S}(f^*X)$, $\hot_{W,S}(g^*X)$ respectivement, où on pose
\[ f^*(X) = X \times_S (S',f), \qquad g^*(X) = X \times_S (S',g). \]
Donc on devra sans doute se placer dans l'une des hypothèses suivantes:
\begin{enumerate}
\item[a)] $X$ est lisse sur $S$.
\item[b)] $X$ est Hot-fibrée sur $S$.
\item[c)] $f$ et $g$ sont propres (et sous réserve de la validité de la conjecture page 122).
\end{enumerate}
On s'attend alors à ce que $u^*$ soit associé à un homomorphisme
\[ u^* : g^*(X) \longrightarrow f^*(X) \qquad \textit{dans } \Cat/S' \]
défini en termes de $u$. Y a-t-il moyen de le définir bel et bien, sans supposer $X$ au moins fibrée sur $S$? Prenons $S'=e$, donc $f$, $g$ correspondant à des objets $s,t\in\Ob\,S$, et $u$ à une fléche % [sic] fléche -> flèche, wrong accent
\[ u : s \longrightarrow t, \]
les images inverses de $X/S$ par $f$, $g$ sont les fibres $X_s$, $X_t$, et on cherche
\[ u^* : X_t \longrightarrow X_s. \]
On n'a aucune chance d'en trouver canoniquement, si on ne suppose $X$ fibrée sur $S$! Donc ce ne sont pas a), b) ni c) qui peuvent nous servir, mais unique%
\printedpage{80}%
% head: A. Grothendieck : Dérivateurs, Ch. VI, Édition des Universités de Montpellier II et Paris VII
% foot: Version 9 novembre 2001
\msp{page 130}%
ment l'hypothèse $X\in\Ob\,\uFib\,S$. Donc on ne peut espérer trouver une description de $u^*$ entre $\HOT_W(S)$ et $\HOT_W(S')$ en termes d'opérations sur les catégories sur $S$, $S'$, comme suggéré à l'instant, qu'en se bornant à des `modèles' $X$ qui sont dans $\uFib\,S$, tout au moins. Et je ne doute pas en effet que ça marche sur $\uFib\,S$ pas moins bien que directement sur $\uFibsc\,S$, en termes de $\uHom(S^o,\Cat)$.

\dgrblock{dgrb231}{Local constancy and the missing comma formula}
\begin{dcomment}{development}{p128}{Local constancy and the missing comma formula}
Grothendieck asks whether the locally constant part, as the base varies, is itself a derivator. His expectation that it is not is correct for ordinary homotopy types, but the suggested obstruction to the existence of adjoints is not the obstruction. It is the pointwise formula for those adjoints that fails. We give the small calculation because it separates this question from the comparison theorem.

\textbf{Proposition (the interval obstruction).}\ours
Let $\mathcal L(S)$ be the full subcategory of the space derivator on $S$ consisting of locally constant objects, with the induced restriction functors. Then $S\mapsto\mathcal L(S)$ does not satisfy the pointwise left Kan-extension axiom.

\textit{Proof.} We use the source's presheaf convention. Let $S=(0\to1)$, let $a:\{0\}\to S$, and let $b:\{1\}\to S$. An object of $\mathcal L(S)$ is a coherent arrow $F_1\to F_0$ which is an equivalence. Evaluation at $0$ is an equivalence
\[
 a^*:\mathcal L(S)\longrightarrow\mathcal L(e)=\Hot.
\]
Indeed the constant-arrow functor is a quasi-inverse: the diagram map from $F_1\to F_0$ to $F_0\xrightarrow{\id}F_0$ has components $F_1\to F_0$ and $\id_{F_0}$, both weak equivalences. This construction is natural in $F$, before taking the homotopy category. Thus a left adjoint $a^{\mathcal L}_!$ exists and is the constant-arrow functor, up to the unique isomorphism between adjoints.

Take the one-point space $*$. We obtain
\[
 b^*a^{\mathcal L}_!(*)\simeq *.
\]
On the other hand, the comma category $1\backslash\{0\}$ of arrows $1\to a(0)$ is empty. The pointwise formula for a left Kan extension in the presheaf convention would therefore identify this value with the homotopy colimit of the empty diagram, namely $\varnothing$. The categories $\mathcal L(e)=\Hot$ and $\mathcal L(\varnothing)=e$ have the usual initial-object adjunction, so that this last value is forced. Since $\varnothing$ is not isomorphic to $*$, the required base-change map cannot be invertible.\hfill$\square$

This proof also identifies what the unrestricted direct image does: in the full space derivator $a_!(*)$ is the representable presheaf $S(-,0)$, with values $*$ at $0$ and $\varnothing$ at $1$. It is not locally constant. The locally constant left adjoint first takes this extension and then reflects it. The reflection changes the value at $1$, and the comma formula is lost.

\textbf{Corollary (commutation beyond Hot-fibrations).}\ours
The inclusion $a:\{0\}\to(0\to1)$ commutes with local-constant reflection for spaces, but is not a Hot-fibration.

\textit{Proof.} The reflection of any arrow $F_1\to F_0$ is the constant arrow on $F_0$. To see its universal property, map to a constant arrow on $K$; the space of maps is $\operatorname{Map}(\operatorname{hocolim}_{S^{\mathrm{op}}}F,K)$, and $0$ is terminal in $S^{\mathrm{op}}$, so this homotopy colimit is $F_0$. Evaluation $a^*$ at $0$ therefore commutes with reflection. Nevertheless the weak equivalence $\{1\}\to S$ pulls back along $a$ to $\varnothing\to\{0\}$, which is not a weak equivalence. Hence $a$ is not a Hot-fibration. This answers the question on ms.~127.\hfill$\square$

The later general result is stronger than is needed for the example. \citet[Corollary~3.11]{Cisinski2009LocallyConstant} proves that restriction on locally constant coefficients has both adjoints for every derivator. His Theorem~3.13 concerns a different construction: one fixes a category $A$ and varies an additional diagram category $I$, retaining arbitrary coherent $I$-diagrams of locally constant $A$-coefficients. That construction is a derivator. It does not say that the varying-base assignment $S\mapsto\mathcal L(S)$ is one.

We checked those statements, the pointwise formula in \citet[Proposition~1.26]{Groth2013}, and the localization description in Cisinski's paper. We did not find this exact interval calculation stated there. The marked result records our explicit completion of the question in these pages; it makes no claim of priority for the obstruction.
\end{dcomment}


Pour les axiomes essentiels sur les dérivateurs, je renvoie à I p. 67-68 et p. 111-112. On a vu que Der 1, \circled{Der 2} (\footnote[66]{Canulé pour \underline{Der 2}, je ne sais si c'est vrai même en supposant $W$ \textit{fortement} saturé, i.e. tel que $\hot_W(f)$ isomorphisme $\Longrightarrow f\in W$.}) sont satisfaits pour tout prédérivateur défini par un couple $(\mathcal{M},W\subseteq\mathrm{Fl}\,\mathcal{M})$, donc c'est OK ici pour $\HOT_W$. Il s'ensuit que c'est OK aussi pour $\HOT_W^{\mathrm{lc}}$, vu la caractérisation des $\HOT_W^{\mathrm{lc}}(S)$ en tant que sous-catégories de $\HOT_W(S)$. Pour Der 1, ça marche-t-il même pour des sommes quelconque de catégories, i.e. a-t-on % [sic] quelconque -> quelconques, agreement with sommes; Canulé (footnote 66) is verified authorial wording; see corrections.csv
\[ \HOT_W(\coprod_i S_i) \xrightarrow{\sim} \prod_i \HOT_W(S_i)\,? \]
On a en tous cas
\[ \uFib(\overbrace{\coprod\limits_i S_i}^{S}) \xrightarrow{\sim} \prod_i \uFib(S_i), \]
et le localiseur $W_S^{\mathrm{f}}$ dans $\uFib(S)$ s'identifie
\msp{page 131}
bel et bien au produit des localiseurs $W_{S_i}^{\mathrm{f}}$ dans les $\uFib\,S_i$. Mais dans le cas d'une famille peut-être infinie de catégories $\mathcal{F}_i$, munies de localiseurs $W_i$, si on pose
\[ \mathcal{F} = \prod\mathcal{F}_i, \qquad W = \prod W_i \subseteq \mathrm{Fl}(\mathcal{F}), \]
le foncteur canonique
\[ \mathcal{F}W^{-1} \longrightarrow \prod_i(\mathcal{F}_iW_i^{-1}) \]
est-il toujours une équivalence? Même dans \texttt{[`Même dans' barré]} le cas $I$ fini il faut (il me semble) un argument assez soigneux, soit par la propriété universelle par récurrence sur card\,$I$ (on est ramené au cas de deux foncteurs), soit à coups de catégories de chemins, en notant que
\[ \uCh_\infty^W(\mathcal{F};(s_i),(t_i)) \xrightarrow{\sim} \prod_i \uCh_\infty^{W_i}(\mathcal{F}_i;s_i,t_i), \tag*{$(*)$} \]
et en passant aux $\pi_0$, qui donne
\[ \mathrm{Hom}_W((s_i),(t_i)) \xrightarrow{\sim} \prod_i \mathrm{Hom}_{W_i}(s_i,t_i). \]
Mais la formule $(*)$, visiblement, ne marche que pour $I$ fini. Dans le cas général, on trouve la sous-catégorie $\prod_i^!$ de $\prod_i$ formée des systèmes $(c_i)_{i\in I}$ de chemins qui sont de longueur bornée.

\dgrblock{dgrb232}{A natural transformation is realized by transport}
\begin{dcomment}{exposition}{p130}{A natural transformation is realized by transport}
Grothendieck has explained restriction on objects; a prederivator also remembers how restriction varies with a natural transformation of indexing functors. For $u:f\Rightarrow g:S'\to S$, the presheaf convention requires a transformation $u^*:g^*\Rightarrow f^*$. A pullback square alone supplies no arrow in this direction. A fibration supplies it through Cartesian lifting.

Let $p:X\to S$ be a Grothendieck fibration. For every $a\in S'$ and $x\in X_{g(a)}$, choose a Cartesian arrow
\[
 \lambda_{a,x}:u_a^*x\longrightarrow x
\]
over $u_a:f(a)\to g(a)$. Send $(a,x)$ in $g^*X$ to $(a,u_a^*x)$ in $f^*X$. For an arrow $(h,v):(a,x)\to(b,y)$ in $g^*X$, the equality
\[
 g(h)u_a=u_bf(h)
\]
and the Cartesian property of $\lambda_{b,y}$ give a unique arrow $w:u_a^*x\to u_b^*y$ over $f(h)$ satisfying
\[
 \lambda_{b,y}w=v\lambda_{a,x}.
\]
Uniqueness proves preservation of identities and composition. This constructs the requested functor over $S'$.

Different choices of Cartesian lifts are related by unique vertical isomorphisms. The resulting transport functors are therefore naturally isomorphic. The same uniqueness gives the comparison between transport along a composite transformation and successive transport, together with its coherence. A splitting can make selected comparison maps identities, but the localized operation is independent of the splitting.

For an arbitrary functor $H:X\to Y$ over $S$, transport need not commute strictly with $H$: $H$ need not preserve Cartesian arrows. Factoring $H(\lambda_{a,x})$ through a chosen Cartesian lift in $Y$ gives a vertical natural transformation between the two resulting functors. A vertical natural transformation is a homotopy over $S'$: it is represented by a functor from the cylinder $X\times[1]$, whose projection to $X$ is a weak equivalence on fibres. Thus the two maps agree after localization. This supplies the naturality needed when all functors over the base, rather than just Cartesian functors, are used in the fibred presentation.

The source's negative observation is accordingly about finding a direct functor on insufficiently structured representatives. It does not obstruct the existence of the $2$-functor. That $2$-functor is already defined by precomposition on coherent presheaves; the lifting calculation explains how to recognize its $2$-cells geometrically.
\end{dcomment}


\dgrblock{dgrb233}{Strong saturation does not imply pointwise detection}
\begin{dcomment}{development}{p130}{Strong saturation does not imply pointwise detection}
Footnote~66 withdraws the assertion that every relative category defines a prederivator satisfying Der~2, even when its weak equivalences are strongly saturated. There is a distinction between two statements. If an actual natural transformation has invertible evaluations after localization, strong saturation makes each component a weak equivalence; that transformation is therefore inverted in the diagram category. An arbitrary morphism of the localized diagram category, however, need not have an actual natural transformation as representative. The following construction shows that this distinction can defeat pointwise detection. It also connects this footnote to the unbounded-length obstruction discussed on ms.~131.

\textbf{Proposition (a strongly saturated counterexample to Der~2).}\ours
There exist a small category $M$, a strongly saturated class $W$, and a countable discrete category $I$ such that
\[
 M^I[(W^I)^{-1}]\longrightarrow\prod_{i\in I}M[W^{-1}]
\]
is not conservative. In fact, the source has a morphism whose evaluations are all isomorphisms but has no morphism in the reverse direction. The category $M$ can be chosen to be a disjoint union of finite categories.

\textit{Proof: the finite categories.}
Let $R$ be the walking retraction, with objects $0,1$ and arrows
\[
 i:0\longrightarrow1,\qquad r:1\longrightarrow0,\qquad
 ri=\id_0,\qquad e=ir\ne\id_1.
\]
Its five arrows are $\id_0,\id_1,i,r,e$; in particular, its only isomorphisms are its identities. For example, this is the subcategory of sets generated by the inclusion of a singleton into a two-element set and its retraction.

Fix a positive odd integer $n$. Define a finite poset $Q_n$ with elements
\[
 x_0,\ldots,x_{n-1},\quad y_0,\ldots,y_{n-1},\quad z.
\]
For writing its order relations only, put $x_n=y_n=z$. Generate the order by the two fences
\[
 x_0\longleftarrow x_1\longrightarrow x_2\longleftarrow\cdots
 \longleftarrow x_n,
 \qquad
 y_0\longleftarrow y_1\longrightarrow y_2\longleftarrow\cdots
 \longleftarrow y_n,
\]
together with $x_j\leq y_j$ and $z\leq y_j$ for $0\leq j<n$. An arrow in these displays means an order relation from its source to its target. The generated relation is antisymmetric: no $y_j$ with $j<n$ reaches an $x_k$ or $z$, and the nonidentity arrows among the $x$'s have odd source and even target. Among the $x$'s the only nonidentity comparisons are the fence edges. In particular, $z=x_n$ reaches just $x_{n-1}$ among the other $x$'s. The full subposet on the $y$'s has initial object $y_n=z$.

Let $C_n$ be the full subcategory of the finite product category $R\times Q_n$ on the objects
\[
 X_j=(0,x_j),\quad Y_j=(1,y_j)\quad(0\leq j<n),\qquad
 X_n=(0,z),\quad Y_n=(1,z).
\]
Thus $X_n$ and $Y_n$ are distinct objects, despite having the same second coordinate. Write $F_n:C_n\to R$ for projection and set
\[
 W_n=F_n^{-1}(\operatorname{Iso}R).
\]
This class is strongly saturated. Indeed, $F_n$ inverts $W_n$ and hence factors through $C_n[W_n^{-1}]$. If an original arrow becomes invertible in that localization, its image under $F_n$ is an isomorphism, which is exactly the condition defining $W_n$. This argument also confirms that the section, the retraction and the idempotent $e$ do not inadvertently become weak equivalences.

\textit{A short morphism with a long inverse.}
There are weak arrows $c_j:Y_n\to Y_j$ given by $(\id_1,z\leq y_j)$, with $c_n=\id$, and arrows $i_j:X_j\to Y_j$ given by $(i,x_j\leq y_j)$. The product category also supplies $r_n:Y_n\to X_n$, with
\[
 r_ni_n=\id_{X_n}.
\]
For every fence edge $a:X_s\to X_t$, let $b:Y_s\to Y_t$ be the corresponding weak arrow. The following equations hold strictly:
\[
 i_ta=bi_s,\qquad bc_s=c_t.
\]
In the localization put $t_j=c_j^{-1}i_j:X_j\to Y_n$. These equations give $t_ta=t_s$. Let $p_n:X_0\to X_n$ be the morphism represented by the entire weak fence, traversed from index $0$ to index $n$. Since every fence edge is weak, $p_n$ is invertible. Transporting the preceding equation along its edges gives
\[
 i_np_n=t_0=c_0^{-1}i_0.
\]
Consequently the three-arrow word
\[
 f_n=r_nc_0^{-1}i_0:X_0\longrightarrow X_n
\]
equals $p_n$ and is an isomorphism.

Nevertheless every zigzag representing a morphism $X_n\to X_0$ has at least $n$ arrows. To prove this, assign the heights
\[
 h(X_j)=j,\qquad h(Y_j)=n\qquad(0\leq j\leq n).
\]
Every original forward arrow lowers height by at most one. For arrows between $X$ objects this follows from the fence description. Arrows from an $X$ to a $Y$ cannot lower it. Arrows between $Y$ objects preserve it. Finally, a $Y$ object can map to an $X$ object only from $Y_n$, and then the target is $X_n$ or $X_{n-1}$; these are exactly the comparisons in $Q_n$ with a $y$ source and an $x$ target. Such an arrow also lowers height by at most one.

A weak arrow has the same $R$ coordinate at its two ends. Between $Y$ objects it preserves height, and between $X$ objects its two indices differ by at most one. Its formal inverse therefore lowers height by at most one as well. Every letter in a localization zigzag, whether a forward arrow or an inverted weak arrow, lowers height by at most one. A word from height $n$ to height $0$ consequently has length at least $n$. This argument concerns every possible representative and allows all the extra endomorphisms supplied by fullness in $R\times Q_n$.

\textit{The countable diagram.}
Take
\[
 M=\coprod_{n\text{ positive odd}}C_n,\qquad
 W=\coprod_{n\text{ positive odd}}W_n,
\]
and let $I$ be the discrete category of positive odd integers. The projections $F_n$ together give $F:M\to R$, and $W=F^{-1}(\operatorname{Iso}R)$ is strongly saturated by the same argument. Define four $I$-diagrams by choosing, in the $n$th component, $X_0,Y_0,Y_n,X_n$, respectively. The families $i_0,c_0,r_n$ give a single diagram-category zigzag
\[
 \mathcal X\longrightarrow\mathcal U
 \xleftarrow{\,W^I\,}\mathcal V\longrightarrow\mathcal Z.
\]
Its localized morphism $f:\mathcal X\to\mathcal Z$ evaluates to $f_n$ and is therefore pointwise invertible.

Suppose that any morphism $\mathcal Z\to\mathcal X$ existed in $M^I[(W^I)^{-1}]$. Choose a finite zigzag representing it, of length $L$. At the coordinate indexed by any odd $n>L$, evaluation gives a zigzag $X_n\to X_0$ with at most $L$ letters. All its objects remain in $C_n$, since the disjoint union has no arrows between its components. The height bound gives $n\leq L$, a contradiction. There is thus no reverse morphism at all, and $f$ is not invertible.\hfill$\square$

The use of infinitely many coordinates is essential to this proof; it makes no assertion about the same question restricted to finite nondiscrete shapes. It refutes the unrestricted statement in footnote~66, including its strongly saturated version. It also strengthens the failure of the infinite-product comparison: that comparison can fail to reflect isomorphisms, not merely fail to be full.

For comparison, a precise sufficient hypothesis is available without a model structure. Suppose $W$ is strongly saturated and every morphism in $M^I[(W^I)^{-1}]$ has a representative
\[
 X\xleftarrow{w}X'\xrightarrow{a}Y'\xleftarrow{v}Y,
 \qquad w,v\in W^I.
\]
If its evaluations are isomorphisms, then every $a_i$ becomes invertible in $M[W^{-1}]$. Saturation gives $a_i\in W$, so the displayed word is invertible. A left or right calculus of fractions is a special case, allowing one of the outside arrows to be an identity. The same argument proves the Der~2 part of \citet[Proposition~4.2.8, pp.~85--86]{LagkasNikolos2018}; Remark~4.2.19, p.~95, explicitly discusses the role of two- and three-arrow descriptions in localization of prederivators.

The counterexample above was constructed for this commentary. The literature checks located the stated fraction criterion and the usual model-category results, but did not locate this particular finite-category construction. Its proof is included in full; the marking records our contribution to the present audit and does not assert historical priority for the general failure of localization to commute with infinite products.
\end{dcomment}


Donc la question de la validité de Der 1 sans hypothèse de finitude sur l'ensemble d'indices $I$
\printedpage{81}
% head: A. Grothendieck : Dérivateurs, Ch. VI, Édition des Universités de Montpellier II et Paris VII
% foot: Version 9 novembre 2001
\msp{page 132}
reste en suspens pour l'instant. (Je soupçonne que c'est vrai, mais il faudrait, pour le prouver, tenir une description des flèches dans $\HOT_W(S)$ par des chemins de longueur bornée, disons de longueur deux ou trois. C'est une question intéressante, mais que je réserve pour plus tard, si le besoin m'y contraint $\ldots\,$)

\dgrblock{dgrb234}{Unbounded zigzags obstruct infinite products}
\begin{dcomment}{development}{p131}{Unbounded zigzags obstruct infinite products}
Grothendieck identifies the right difficulty in commuting localization with an infinite product: a morphism in a localization is represented by a finite word, and the lengths needed in different coordinates may be unbounded. The following construction turns that observation into a counterexample, even when the localizing classes are strongly saturated.

\textbf{Proposition (a product of finite fences).}\ours
For $n\geq1$, let $P_n$ be the finite poset with objects $0,1,\ldots,2n$ and nonidentity relations
\[
 0<1>2<3>\cdots<2n-1>2n.
\]
Let $W_n$ consist of every arrow. Then each $P_n[W_n^{-1}]$ is equivalent to the terminal category, but
\[
 \left(\prod_{n\geq1}P_n\right)
 \left[\left(\prod_{n\geq1}W_n\right)^{-1}\right]
 \longrightarrow \prod_{n\geq1}P_n[W_n^{-1}]
\]
is not full. More precisely, two objects $x=(x_n)$ and $y=(y_n)$ in the category on the left can be connected by a morphism if and only if $\sup_n|x_n-y_n|<\infty$.

\textit{Proof.} The only nonidentity arrows in $P_n$ join adjacent integers; there are no composable pairs of nonidentity arrows. Inverting these arrows gives the free groupoid on the undirected path with vertices $0,\ldots,2n$. Between any two vertices there is exactly one reduced path. Hence this groupoid is equivalent to the terminal groupoid.

Put $P=\prod_nP_n$. Since every arrow of $P$ is inverted, an arrow in its localization is represented by a finite zigzag in $P$. One step in such a zigzag changes the $n$th coordinate by at most one, for every $n$. A zigzag with $k$ steps therefore implies $|x_n-y_n|\leq k$ for all $n$. This proves necessity of the bound.

Conversely, suppose $|x_n-y_n|\leq K$ for all $n$. In each fence take the unique shortest undirected path from $x_n$ to $y_n$. Its arrow directions alternate. If its first arrow has the wrong orientation for a common alternating pattern beginning with a forward arrow, prepend an identity. After the path finishes, append identities. Every coordinate now has a zigzag of the same length $K+1$ and the same pattern of directions. Taking coordinatewise products gives a zigzag in $P$. This proves sufficiency.

Take $x_n=0$ and $y_n=2n$. There is no arrow between them in $P[W^{-1}]$, while their images have exactly one arrow between them in the product of the localized categories. Thus the comparison is not full. The localizing classes are all-arrow classes, so they are strongly saturated.\hfill$\square$

The original question for $\HOT_{W_\infty}$ nevertheless has a positive answer. Diagram categories of simplicial sets carry model structures with pointwise weak equivalences. For a disjoint union of shapes, the diagram model category is the product of the diagram model categories of the components. Its cofibrant-fibrant replacements are chosen coordinatewise. Between cofibrant-fibrant objects, every homotopy-category morphism has an actual representative, and homotopy can be tested coordinatewise using coordinatewise cylinder objects. Therefore the homotopy category of this product is the product of the homotopy categories, with no bound on the number of factors. The common model-category description provides the uniform presentations absent from the fence example; compare \citet[Proposition~1.30]{Groth2013}.

We checked the finite-word construction of localization in \citet[Chapter~I, \S1]{GabrielZisman1967} and the model-category derivator result just cited. We did not find this exact fence formulation in those sources. It is an elementary elaboration of the obstruction explicitly singled out in the manuscript, not a claim that the general phenomenon is new. In particular, it does not support replacing Grothendieck's tentative general argument with an assertion that all relative categories satisfy the coproduct axiom.
\end{dcomment}


Plus importante est la question de la validité de Der 3 et de ses variantes (I, p. 111-112). Le minimum qu'on demande d'un dérivateur $\mathbf{D}$, c'est que si $S'\longrightarrow S$ dans Diag est une Hot-équivalence, ce soit une $\mathbf{D}$-équivalence, i.e. (dans le cas actuel) que le foncteur
\[ f^* : \HOT_W(S) \longrightarrow \HOT_W(S') \]
soit pleinement fidèle sur la sous-catégorie des objets \textit{constants} dans $\HOT_W(S)$. Cette condition devient plus tangible quand on dispose déjà soit des foncteurs cohomologie $H^\bullet_{\mathbf{D}}(S,-)$, $H^\bullet_{\mathbf{D}}(S',-)$, soit de l'homologie $H_\bullet^{\mathbf{D}}(S,-)$, $H_\bullet^{\mathbf{D}}(S',-)$. En termes de cette dernière, elle signifie que
\[ p_!p^*(a) \xrightarrow{\sim} p'_!p'^*(a) \qquad \text{\texttt{[est un]} isomorphisme,} \]
\msp{page 133}
pour tout $a$ dans $\mathcal{A}=\mathbf{D}(e)$. (Où $p:S\longrightarrow e$, $p':S'\longrightarrow e$ sont les morphismes structuraux sur la catégorie ponctuelle.) On a vu la formule heuristique générale
\[ p_!\underbrace{p^*(a)}_{a_S} = \hot(S)\times a, \]
où $\hot\,S$ est la classe de $S\in\Ob\,\Cat$ dans $\Hot=(\Cat)W_\infty^{-1}$, et où le produit $\times$ est une opération externe de Hot sur $\mathcal{A}=\mathbf{D}(e)$
\[ \Hot\times\mathcal{A} \longrightarrow \mathcal{A}\ . \]
Dans le cas actuel, $\mathcal{A}=\Hot$, et l'opération
\[ \Hot\times\Hot \longrightarrow \Hot \]
ne doit être autre que le foncteur produit cartésien, lequel est déduit du foncteur similaire dans Cat. Quand on dispose de ce formalisme, le résultat cherché pour un $f:S'\longrightarrow S$ se réduit simplement à dire que si $f$ est dans $W$, alors il en est de même de $f\times\mathrm{id}_X$ pour tout $X$ dans Cat (en prenant $X$ tel que $\hot(X)\simeq a$). Or là c'est un fait bien connu, résultant des axiomes mis sur $W$.

J'ai donc l'impression que je n'ai aucune chance de prouver Der 3a$'$ (p. 111), ni même Der 3a) (cas où $\mathrm{S}=\mathrm{e}$), sans avoir une maîtrise du formalisme des $p_!$. Et j'ai
\printedpage{82}
% head: A. Grothendieck : Dérivateurs, Ch. VI, Édition des Universités de Montpellier II et Paris VII
% foot: Version 9 novembre 2001
\msp{page 134}
bon espoir que cetter maîtrise va s'obtenir à bon compte, alors que la question des $p_*$ semble beaucoup plus délicate. % [sic] cetter -> cette, misspelling

\dgrblock{dgrb235}{Constant coefficients explain the turn to homology}
\begin{dcomment}{exposition}{p133}{Constant coefficients explain the turn to homology}
Grothendieck turns from inverse image to homological direct image because the latter makes homotopy invariance of constant coefficients visible. For the space derivator, the homotopy colimit of a constant space $K$ over $S^{\mathrm{op}}$ is
\[
 \operatorname{hocolim}_{S^{\mathrm{op}}}K\simeq |N(S^{\mathrm{op}})|\times K
 \simeq |NS|\times K.
\]
One way to calculate it is the simplicial replacement: in degree $n$ it is a coproduct of copies of $K$, one for every string of $n$ composable arrows of $S^{\mathrm{op}}$. Realization therefore gives the displayed product. The reversal of a string identifies the realizations of the nerves of a category and its opposite.

For $f:S'\to S$, the canonical homological comparison runs from the expression for $S'$ to the expression for $S$ and is represented by
\[
 |Nf|\times\id_K:|NS'|\times K\longrightarrow|NS|\times K.
\]
When $f\in W_\infty$, this is a weak equivalence. The isomorphism written in the source may consequently be read in either direction after localization; before inversion the functorial map has the direction just displayed. For a general coefficient derivator the corresponding operation is tensoring with a homotopy type, rather than an internal Cartesian product of coefficients.

Adjunction now relates the calculation to the desired fully faithful restriction on constant objects. For constant coefficients $a,b$,
\[
 \Hom_{\mathbf D(S)}(p^*a,p^*b)
 \simeq\Hom_{\mathbf D(e)}(p_!p^*a,b).
\]
If the homological comparison is an isomorphism for every $a$, it induces a bijection on these Hom sets for every $b$. Conversely, bijectivity for every $b$ identifies the comparison by Yoneda. This is why the direct-image formalism gives exactly the statement Grothendieck wants, rather than merely an invariant of isomorphism classes.

The axiom numbers in these pages belong to the evolving manuscript. The condition called Der~3 here should not be confused with the usual modern axiom bearing that number, which asks for both Kan adjoints. Likewise the footnote retracting the unrestricted claim about Der~2 records a genuine limitation of the argument for arbitrary relative categories. The development at manuscript p.~130 gives a counterexample even when $W$ is strongly saturated, using a countably discrete diagram shape. It does not address the restriction of that question to finite shapes. The model-category theorem establishes the pointwise detection axiom for the space derivator; it does not vindicate the cancelled general assertion.
\end{dcomment}


Je vais donc me tourner d'abord vers l'axiome Der 4$'$ d'existence de
\[ f_! : \HOT_W(S') \longrightarrow \HOT_W(S), \]
comme adjoint à gauche de $f^*$. Je vais essayer de construire ad-hoc un foncteur $f_!$, puis les deux flèches d'adjonction
\[ \begin{aligned} \mathrm{id} &\quad\longrightarrow\quad f^*f_! \\ \mathrm{id} &\quad\longleftarrow\quad f_!f^*, \end{aligned} \]
et enfin, de vérifier les deux compatibilités pertinentes.

L'idée géométrique `évidente', c'est pour $X'\in\Cat/S'$, représentant $\xi'\in\Ob\,\HOT_W(S')$, de définir $f_!(\xi')$ par
\[ f_!(\hot_{S',W}(X')) = \hot_{S,W}(X'), \]
où dans le second membre, $X'$ est regardé comme objet de $\Cat/S$, grâce à $f:S'\longrightarrow S$ et au foncteur
\[ f_! : \Cat/S' \longrightarrow \Cat/S \]
qu'il induit. Le seul problème, c'est de voir si ce foncteur est compatible avec les
\msp{page 135}
localiseurs $W_{S'}^{\mathrm{d}}$, $W_S^{\mathrm{d}}$, et sinon, s'il l'est tout au moins sur la sous-catégorie $\uFib\,S'$ ou $\uFibsc\,S'$, qui suffit pour calculer $\HOT_W(S')$ par localisation. Il faut donc voir si un $X'_1\longrightarrow X'_2$ dans $\Cat/S'$, qui est dans $W_{S'}^{\mathrm{d}}$, i.e. qui est dans $W$ et le reste par tout changement de base de la forme $s'\backslash S'\longrightarrow S'$, plus généralement pour tout co-isomorphisme local, est aussi dans $W_S^{\mathrm{d}}$, i.e. si les changements de base de la forme $s\backslash S\longrightarrow S$ le transforment en $W$-équivalences (grossières). Or les co-isomorphismes locaux sont stables par \textit{tout changement de base} (cf. I, p. 55 dans le cas dual des isomorphismes locaux), donc on gagne!
\[
\begin{tikzcd}
\overline{S'} = s\backslash S' \arrow[r] \arrow[d] & S' \arrow[d] \\
\overline{S} = s\backslash S \arrow[r] & S
\end{tikzcd}
\quad .
\]
Je vais quand même le vérifier encore. Si
\[
\begin{tikzcd}
\overline{S'} \arrow[r] \arrow[d] & S' \arrow[d] \\
\overline{S} \arrow[r] & S
\end{tikzcd}
\]
\printedpage{83}
% head: A. Grothendieck : Dérivateurs, Ch. VI, Édition des Universités de Montpellier II et Paris VII
% foot: Version 9 novembre 2001
est un carré cartésien dans Cat, on a dit dans loc. cit., essentiellement, que pour $\overline{s'}$ dans $\overline{S'}$, ayant des images $s'$, $\overline{s}$, $s$ dans $S'$, $\overline{S}$, $S$ respectivement, le carré des colocalisés
\[
\begin{tikzcd}
\overline{s'}\backslash\overline{S'} \arrow[r] \arrow[d] & s'\backslash S' \arrow[d] \\
\overline{s}\backslash\overline{S} \arrow[r] & s\backslash S
\end{tikzcd}
\]
est encore
\msp{page 136}
cartésien. Je viens encore le de revérifier à l'instant, pour les objets tout au moins. Dire que $\overline{S}\longrightarrow S$ est un isomorphisme local, signifie que les $\overline{s}\backslash\overline{S}\longrightarrow s\backslash S$ sont \texttt{[des]} isomorphismes. Mais le carré cartésien précédent montre qu'alors les $\overline{s'}\backslash\overline{S'}\longrightarrow s'\backslash S'$ sont aussi \texttt{[des]} isomorphismes, on gagne! % [sic] le de revérifier -> de le revérifier, word order

\textbf{Théorème 1.} \textit{Pour toute flèche $f:S'\longrightarrow S$ dans} Cat\textit{, le foncteur induit}
\[ f_! : \Cat/S' \longrightarrow \Cat/S \]
\textit{est compatible avec les localiseurs $(W_{S'}^{\mathrm{d}},W_S^{\mathrm{d}})$, donc dualement avec les localiseurs $(W_{S'}^{\mathrm{g}},W_S^{\mathrm{g}})$, et trivialement avec les localiseurs $(W_{S'},W_S)$. Il induit donc un diagramme commutatif de foncteurs sur les localisés}
\[
\begin{tikzcd}
\HOT_W(S') \arrow[r, "\rho_{S'}"] \arrow[d, "f_!"] & \HOT_W^{\mathrm{lc}}(S') \arrow[d, "f_!^{\mathrm{lc}}"] & \HOT_W^o(S') \arrow[l, "\sigma_{S'}"'] \arrow[d, "f_!^o"'] \\
\HOT_W(S) \arrow[r, "\rho_S"] & \HOT_W^{\mathrm{lc}}(S) & \HOT_W^o(S). \arrow[l, "\sigma_S"']
\end{tikzcd}
\]

\textbf{Définition 1.} \textit{On a défini ainsi les foncteurs verticaux $f_!$, $f_!^{\mathrm{lc}}$, $f_!^o=(f^o)_!$, appelés foncteurs image directe homologique, ou foncteurs d'intégration pour $f$ (et les types d'homotopie relatifs).}

\msp{page 137}
On a obtenu en même temps, gratuitement, la compatibilité de ces foncteurs aux foncteurs \texttt{[de]} localisation $\rho_{S'}$, $\sigma_{S'}$ pour $\HOT(S')$, et $\rho_S$, $\sigma_S$ pour $\HOT(S)$. Par contre, on ne s'attendra pas qu'ils commutent avec les foncteurs pleinement fidèles d'inclusion des $\HOT_W^{\mathrm{lc}}$ dans $\HOT_W$. En effet, si on ne suppose pas $f$ dans \underline{Fibhot} (\footnote[67]{Ce cas doit être regardé avec soin ultérieurement.}), il est tout à fait immoral de s'attendre à ce que $f_!$ dans $\HOT_W(S')$ transforme types localement constants en itou. La situation est donc exactement inverse de celle pour le foncteur $f^*$, qui était compatible avec les foncteurs d'inclusion $i_S$, $i_{S'}$ et $j_S$, $j_{S'}$ - c'était à tel point trivial qu'il n'a pas été question même de le dire, voire en faire une proposition! La situation est

\dgrblock{dgrb236}{Changing the structure map gives direct image}
\begin{dcomment}{exposition}{p136}{Changing the structure map gives direct image}
Grothendieck's proposed direct image has a particularly economical description: keep the total category and compose its map to the base. The substantial point is that this operation respects right-fine equivalence even though inverse image did not respect it on arbitrary representatives. The proof uses a class of test changes of base which is itself stable under pullback.

Let $u:X\to Y$ be in $W_{S'}^{\mathrm d}$ and let $f:S'\to S$. For an object $s\in S$, form
\[
 E=(s\backslash S)\times_S S'\longrightarrow S'.
\]
This projection is a discrete opfibration: an object is an arrow $s\to f(s')$, and an arrow out of $s'$ has a unique lift obtained by composing that arrow in $S$. Right-fine equivalences remain global $W$-equivalences after pullback along such maps. Applying this characterization to $u$ shows that
\[
 s\backslash(X/S)\longrightarrow s\backslash(Y/S)
\]
is in $W$. Since this holds for every $s$, the same underlying functor, regarded over $S$, belongs to $W_S^{\mathrm d}$.

The characterization can also be checked directly from the comma resolution. For a discrete opfibration $q:E\to S'$ and $e\in E$, unique lifting gives
\[
 e\backslash(X\times_{S'}E)\cong q(e)\backslash X.
\]
Thus the map between the comma resolutions over $E$ is a weak equivalence on fibres. The localizer's fibrewise-to-total property makes it a global weak equivalence, and the comparison from each comma resolution to its total category is a weak equivalence. Two-out-of-three gives the desired global conclusion. This uses only the comparison machinery already available.

The square of coslices considered in the source is Cartesian on morphisms as well as objects: an arrow is a pair of commuting triangles whose images in the bottom-right coslice agree. The universal property of the original pullback gives exactly one arrow in the top-left coslice with those data. Therefore the author's object check extends immediately to a proof of the asserted categorical pullback. The term ``isomorphisme local'' on ms.~136 must be read in its co-local sense, since it is outgoing coslices that occur throughout this argument.

We obtain
\[
 f_!q_{S'}(X)=q_S(X\xrightarrow{}S'\xrightarrow{f}S).
\]
Composition of structure maps gives the composition law for these functors. Their compatibility with $\rho$ and $\sigma$ is equally formal: the routes around the displayed square apply the same total-category functor and then invert a larger class. Neither observation says that direct image preserves locally constant objects. The interval example above shows that it need not.
\end{dcomment}


% ---- from Der-ChVI-345_p086-090.tex
\printedpage{84}
% head: A. Grothendieck : Dérivateurs, Ch. VI, Édition des Universités de Montpellier II et Paris VII
même encore meilleure pour le foncteur $f_!$ entre les HOT, puisque la description de ce foncteur peut se faire directement sur $\Cat/S'$ tout entier, sans avoir à se borner à des sous-catégories plus ou moins \msp{page 138}astucieuses, où le foncteur induit serait compatible avec les localiseurs ($\underline{\mathrm{Liss}}\,S$ et $\underline{\mathrm{Fibhot}}\,S$ dans le cas de $f^*$).

\dgrblock{dgrb237}{Integration and the local-constant reflection}
\begin{dcomment}{examples}{p137}{Integration and the local-constant reflection}
Grothendieck distinguishes two compatibility statements here because integration changes the support of a relative type. The inclusion $a:\{0\}\to S=(0\to1)$ already exhibits the phenomenon. The one-point object over $\{0\}$ is locally constant. Its direct image in the full relative category has comma fibres
\[
 (0\backslash\{0\})\simeq e,
 \qquad (1\backslash\{0\})=\varnothing.
\]
It is therefore not locally constant on $S$. Its reflection is the constant point, as calculated in the earlier block on ms.~128. Thus $a_!i_{\{0\}}$ and $i_Sa_!^{\mathrm{lc}}$ cannot be identified in general.

The commutation with reflection goes in the direction stated in the theorem:
\[
 \rho_S f_!\simeq f_!^{\mathrm{lc}}\rho_{S'}.
\]
There is a second explanation once the adjunctions have been established. The inclusions $i_S$ commute with restriction, and $\rho_S$ is left adjoint to $i_S$. Both functors in this display are consequently left adjoint to the same composite restriction followed by inclusion. The uniqueness of adjoints identifies them. In contrast, preserving the essential image of $i_{S'}$ under $f_!$ is a further property of $f$.

For ordinary weak equivalences, the Hot-fibration hypothesis in footnote~67 supplies precisely that property. If $s\to t$ is an arrow of $S$, the induced functor $t\backslash S'\to s\backslash S'$ is a weak equivalence. For a locally constant coefficient $F$, the pointwise value $(f_!F)(s)$ is the homotopy colimit of the pulled-back coefficient over $(s\backslash S')^{\mathrm{op}}$. A weak equivalence of indexing categories preserves this homotopy colimit with locally constant coefficients. Thus every transition of $f_!F$ is an equivalence: $f_!$ carries locally constant objects to locally constant objects, and $f_!i_{S'}\simeq i_S f_!^{\mathrm{lc}}$. This is the space-valued case of \citet[Proposition~1.24, applied to the opposite indexing convention]{Cisinski2009LocallyConstant}; the proof there makes the same pointwise homotopy-colimit comparison. It gives the affirmative answer sought in the footnote, while retaining the hypothesis that the preceding interval example shows to be substantive.

For the terminal map $p:S\to e$, composition of structure maps simply forgets the base. If a category $X\to S$ represents a coherent family of homotopy types, then $p_!q_S(X)$ is the homotopy type of the total category $X$. For a category of elements this is the homotopy-colimit formula expressed geometrically. It explains the word ``integration'': the coherent fibres, with their transport, are assembled into their total homotopy type. One should retain the word ``coherent'' here, since an ordinary diagram in the homotopy category alone does not determine that total type.
\end{dcomment}


On a de façon immédiate la transitivité de $f_!$, $f_!^{\mathrm{lc}}$, $f_!^{o}$, i.e.
\[
(gf)_! = g_!f_!, \qquad (gf)_!^{\mathrm{lc}} = g_!^{\mathrm{lc}}f_!^{\mathrm{lc}}, \qquad (gf)_!^{o} = g_!^{o}f_!^{o}.
\]

Mais on ne va pas faire bouger $S$, $S'$, $f$ pour le moment, car il faut maintenant, avant toute autre chose, essayer d'établir l'adjonction de $f_!$ et de $f^*$ (dans cet ordre).

Pour le moment, occupons-nous de $\HOT_W$ seulement - le cas de $\HOT_W^{o}$ s'en déduira en appliquant ce qu'on saura à $f^{o}$. Et on s'occupera de $\HOT_W^{\mathrm{lc}}$ plus tard.

1) \textit{L'homomorphisme}
\[
\alpha : \id_{\HOT_W(S')} \longrightarrow f^*f_!.
\]

Si $X'$ est dans $\Cat/S'$, $f_!(\hot_{S',W}(X'))$ est $\hot_{S,W}(X')$, mais on ne peut trouver son image inverse $f^*$ sur $S'$ en prenant l'image inverse catégorique. Car même si on fait des hypothèses draconiennes sur $X'$ en tant qu'objet de $\Cat/S'$ (p.ex. $X'$ fibrée sur $S'$), celles-ci seront réduites à néant quand on aura composé avec $f$, sur lequel nous ne voulons ici faire aucune hypothèse. Il \msp{page 139}faut donc d'abord appliquer le foncteur $\Phi_S$ à $X'$ sur $S$, puis prendre son image inverse sur $S'$ :
\[
f^*f_!(\hot_{S',W}(X')) = \hot_{S',W}(\Phi_S(X'/S)\times_S S').
\]
Il faut donc trouver une flèche
\[
\alpha_{X'} : X' \longrightarrow \Phi_S(X'/S)\times_S S'
\]
dans le localisé $(\Cat/S')(\Wd_{S'})^{-1}$, voire même une flèche dans $\Cat/S'$ lui-même. Mais par les propriétés d'adjonction des foncteurs $f_!$ et $f^*$ au \textit{niveau catégorique,} une telle flèche revient à une flèche
\[
f_!(X') \longrightarrow \Phi_S(X'/S)
\]
dans $\Cat/S$, i.e. à une flèche qui rend commutatif le carré
\[
\begin{tikzcd}
X' \arrow[r] \arrow[d] & \Phi_S(X'/S') \arrow[d] \\ % [sic] \Phi_S(X'/S') -> \Phi_S(X'/S), the arrow f_!(X') -> \Phi_S(X'/S) just above has X'/S
S' \arrow[r] & S.
\end{tikzcd}
\]
% foot: Version 9 novembre 2001

\printedpage{85}
% head: A. Grothendieck : Dérivateurs, Ch. VI, Édition des Universités de Montpellier II et Paris VII
Or on a une telle flèche canonique, savoir $i_{X'/S}$. Elle est fonctorielle en $X'$ sur $S$ (sans avoir même à se soucier si $X'$ sur $S$ provient d'un $X'$ sur $S'$, i.e. si la flèche structurale $X'\longrightarrow S$ se factorise par $S'\longrightarrow S$).

\leavevmode\msp{page 140}2) \textit{L'homomorphisme}
\[
\beta : \id_{\HOT_W(S)} \longleftarrow f_!f^*.
\]

\dgrblock{dgrb238}{The unit comes from the comma replacement}
\begin{dcomment}{exposition}{p139}{The unit comes from the comma replacement}
The unit must compare an object over $S'$ with the result of forgetting down to $S$ and then deriving inverse image. Forgetting the intermediate base can destroy the useful smoothness or fibration hypothesis. This is why $\Phi_S$ appears between those two operations, even if the original category was already fibred over $S'$.

It is helpful to describe the construction on objects. For $p:X\to S$, the category $\Phi_S(X)$ has objects $(s,x,v:s\to p(x))$ and is projected to $s$. A morphism from $(s,x,v)$ to $(t,y,w)$ is a pair $(h:s\to t,k:x\to y)$ satisfying $p(k)v=wh$. The unit $i_X$ sends $x$ to $(p(x),x,\id_{p(x)})$. Thus for $p':X'\to S'$ the unit of the desired adjunction is represented by
\[
 \alpha_{X'}:X'\longrightarrow\Phi_S(X'/S)\times_S S',
 \qquad x\longmapsto
 \bigl((f p'(x),x,\id_{f p'(x)}),p'(x)\bigr).
\]
The same formula on arrows uses $f p'(k)$ in the first coordinate and $p'(k)$ in the second. It is a functor over $S'$ and is natural in every functor $X'\to Y'$ over $S'$.

Equivalently, it is uniquely characterized by its two projections:
\[
 \pr_1\alpha_{X'}=i_{X'/S},
 \qquad \pr_2\alpha_{X'}=p'.
\]
This characterization is the ordinary slice adjunction between composition with $f$ and pullback along $f$, applied before localization. The target of $\alpha_{X'}$ represents $f^*f_!q_{S'}(X')$ because $\Phi_S(X'/S)$ is a split fibration over $S$. Hence the natural map descends to $\id\to f^*f_!$.

Naturality before localization is enough here, but it is important to use the correct base in the replacement. The top-right object in the source's square on ms.~139 is printed as $\Phi_S(X'/S')$. Its structure map and the preceding formula show that $\Phi_S(X'/S)$ is intended. This is a source slip, retained in the transcription and registered separately. No replacement over $S'$ can repair the loss of smoothness caused by composing its structure map with an arbitrary $f$.
\end{dcomment}


\dgrblock{dgrb239}{The replacement must be taken over the new base}
\begin{dcomment}{exposition}{p139}{The replacement must be taken over the new base}
Grothendieck now has a particularly economical candidate for direct image: a category over $S'$ is regarded as a category over $S$. The difficulty is to show that this operation has the expected universal property after localization. Its prospective right adjoint is not arbitrary strict pullback. It is pullback after a replacement which remembers arrows in the base. The next seven manuscript pages make that distinction precise enough to construct an adjunction, including its two triangle identities.

Write $Q_S=\Phi_S$ and $i_X:X\longrightarrow Q_SX$. For a category $p:X\longrightarrow S$, an object of $Q_SX$ is a triple $(s,x,a:s\longrightarrow p(x))$. A morphism from $(s,x,a)$ to $(t,y,b)$ consists of $v:s\longrightarrow t$ and $u:x\longrightarrow y$ satisfying $p(u)a=bv$. The projection to $s$ is the split fibration associated with the contravariant functor $s\longmapsto s\backslash X$. The map $i_X$ sends $x$ to $(p(x),x,\id_{p(x)})$. This recalls the comma replacement discussed at manuscript p.~20 (``Which comma construction repairs which variance''), but here its role is more demanding: it supplies the unit of a derived adjunction.

For $f:S'\longrightarrow S$, denote by $L$ the functor induced by composition with $f$, and by $R$ derived inverse image. If $[X]$ denotes $\hot_{S,W}(X)$, their formulas are
\[
 L[X']=[X'/S],\qquad R[X]=[Q_SX\times_SS'].
\]
The replacement in the second formula is over $S$. Even when $X'\longrightarrow S'$ is already a fibration, its composite with an arbitrary $f$ need not be one. Replacing over $S'$ would therefore solve the wrong problem. This is the reason for the otherwise rather insistent occurrences of $X'/S$ in the argument. The occurrence of $Q_S(X'/S')$ in the square on manuscript page 139 is a slip already present in the edition; its intended meaning is $Q_S(X'/S)$.

The unit can be written before localization. Its component is the unique $S'$-functor
\[
 a_{X'}:X'\longrightarrow Q_S(X'/S)\times_SS'
\]
whose first projection is $i_{X'/S}$ and whose second projection is the structure map of $X'$. Explicitly, it sends $x'$ to the triple with base object $f(p'(x'))$ and identity arrow, together with $p'(x')\in S'$. The construction is natural in $X'$. Passing through the localization therefore gives $\alpha:\id\longrightarrow RL$.

The counit is slightly different. If $X\longrightarrow S$ is smooth, the projection $X\times_SS'\longrightarrow X$ represents $\beta_{[X]}:LR[X]\longrightarrow[X]$. For an arbitrary representative, the formula is instead
\[
 \beta_{[X]}=[i_X]^{-1}
 [\pr_1:Q_SX\times_SS'\longrightarrow Q_SX].
\]
Both arrows in this fraction lie over $S$, and the arrow inverted is one already known to belong to $W_S^{\mathrm d}$. The forgetful functor $Q_SX\longrightarrow X$ on underlying categories does not generally lie over $S$: the two structure maps at $(s,x,a)$ have values $s$ and $p(x)$. The fraction is thus doing necessary work.

The formulas also explain why naturality survives the change of representatives. For an arrow $u:X\longrightarrow Y$ over $S$, the equalities $Q_S(u)i_X=i_Yu$ and $Q_S(u)\pr_1=\pr_1(Q_S(u)\times_SS')$ give the naturality square for $\beta$ after the $i$'s have been inverted. The construction uses the localization universal property, rather than a choice of inverses in the original category. This is precisely the amount of structure needed for the adjunction proof that follows.
\end{dcomment}


On peut décrire $\HOT_W(S)$ par $(\underline{\mathrm{Fib}}\,S)(\Wd_S)^{-1}$ ou $(\underline{\mathrm{Liss}}\,S)(\Wd_S)^{-1}$ par exemple, de sorte que $f^*$ a une description simple par
\[
f^*(\hot_{S,W}(X)) = \hot_{S',W}(X\times_S S')
\]
(pour $X$ lisse sur $S$). On a ensuite
\[
f_!f^*(\hot_{S,W}(X)) = \hot_{S,W}(X\times_S S'),
\]
mais attention, $X\times_S S'$ n'est pas a priori dans $\underline{\mathrm{Liss}}\,S$, seulement dans $\Cat/S$ sans plus. Mais cela n'empêche qu'il donne ce qu'il faut, par localisation $(\Wd_S)^{-1}$ dans $\Cat/S$. Pour définir $\beta_X$, il suffit donc de trouver une flèche
\[
\beta_X : X \longleftarrow X\times_S S'
\]
dans $\Cat/S$, et on voit bien $\pr_1^X$ (\footnote[68]{quand $\xi\in\HOT_W(S)$ est défini par $X$ dans $\Cat/S$ non nécessairement lisse, on doit trouver $\Phi_S(X)\times_S S'\longrightarrow X$ dans $(\Cat/S)(\Wd_S)^{-1}$, on prend $(i_X)^{-1}\pr_1^{\Phi_S(X)}$ (cf. page 143).}). On a donc
\[
\left\{
\begin{array}{lcl}
\alpha : \id & \longrightarrow & f^*f_! \\
\beta : \id & \longleftarrow & f_!f^*
\end{array}
\right.
\]

3) \textit{La compatibilité}
\[
\begin{tikzcd}
f_! \arrow[r, "f_!*\alpha"] \arrow[rr, bend right, "\id"'] & f_!f^*f_! \arrow[r, "\beta*f_!"] & f_!.
\end{tikzcd}
\]

On doit donc prendre $X'$ sur $S'$, et expliciter les flèches à partir de l'objet
\[
\xi' = \hot_{S',W}(X').
\]
% foot: Version 9 novembre 2001

\printedpage{86}
% head: A. Grothendieck : Dérivateurs, Ch. VI, Édition des Universités de Montpellier II et Paris VII
\leavevmode\msp{page 141}On a
\[
f_!(\xi') = \hot_{S,W}(X'/S),
\]
et pour calculer $f^*f_!(\xi')$, il faut comme tantôt prendre $\Phi_S(X')$ et on trouve
\[
f^*f_!(\xi') = \hot_{S',W}(\Phi_S(X')\times_S S').
\]
Enfin,
\[
f_!f^*f_!(\xi') = \hot_{S,W}(\Phi_S(X')\times_S S').
\]
Il faut expliciter les flèches
\[
\begin{tikzcd}[row sep=small]
\underbrace{f_!(\xi')} \arrow[r, "{f_!(\alpha_{\xi'})}"] & \underbrace{f_!f^*f_!(\xi')} \arrow[r, "{\beta_{f_!(\xi')}}"] & \underbrace{f_!(\xi')}. \\
X'/S \arrow[r, "{\alpha_{X'}}", "{\text{déduit de } i_{X'}}"'] & \Phi_S(X')\times_S S' \arrow[r, dotted, "{\beta_{X'}}"] & X'/S \\
 & \text{sur } S &
\end{tikzcd}
\tag*{$(*)$}
\]

$\alpha_{X'}$ est déduit de $i_{X'} : X'\longrightarrow\Phi_S(X')$ comme l'unique $S'$-morphisme rendant commutatif le diagramme
\[
\begin{tikzcd}
X' \arrow[r, "{\alpha_{X'}}"] \arrow[d] \arrow[rr, bend left, "{i_{X'}}"] & \Phi_S(X')\times_S S' \arrow[r, "\pr_1"] \arrow[dl, dotted] & \Phi_S(X') \arrow[d] \\ % Reading verified against PDF leaf 48: the curved i_X' arrow is solid.
S' \arrow[rr] & & S.
\end{tikzcd}
\tag*{$(**)$}
\]

Pour décrire $\beta_{f_!(\xi')}$, où $f_!(\xi') = \hot_{S,W}(X')$, il faut donner une description générale de $\beta_\xi$ qui soit valable pour \textit{tout} représentant $X$ de $\xi$, pas seulement dans le cas $X$ lisse sur $S$. Pour calculer $f^*(\xi)$, il faut donc prendre le $\hot_{S',W}$ de $\Phi(X)\times_S S'$, puis pour $f_!f^*(\xi)$, il \msp{page 142}faut prendre le $\hot_{S,W}$ du même $\Phi_S(X)\times_S S'$. Il faut définir, si on peut, une flèche
\[
\Phi_S(X)\times_S S' \longrightarrow X \qquad\qquad \text{dans } \Cat/S,
\]
ou du moins dans $(\Cat/S)(\Wd_S)^{-1}$. Mais on a
\[
i_X : X \longrightarrow \Phi_S(X)
\]
qui est dans $\Wd_S$, et qui par produit cartésien définit
\[
i_X\times_S\id_{S'} : X\times_S S' \longrightarrow \Phi_S(X)\times_S S'
\]
(en d'autres termes on a pris l'image inverse de $i_X$ par $S'\longrightarrow S$). Cette flèche est-elle encore dans $\Wd_S$? (\footnote[69]{Non, elle ne l'est pas, car si $S'=\{s\}$, cela signifierait que $X'_s\longrightarrow s\backslash X$ est dans $W$, ce qui est faux en général.}). Si oui, on a le diagramme
% foot: Version 9 novembre 2001
\printedpage{87}
% head: A. Grothendieck : Dérivateurs, Ch. VI, Édition des Universités de Montpellier II et Paris VII
\[
\begin{tikzcd}
\Phi_S(X)\times_S S' \arrow[r, dotted, "\beta_X ?"] & X \\
X\times_S S' \arrow[u, "{i_X\times_S\id_{S'}}", "{=\pi_X}"'] \arrow[ur, "\pr_1^X"] &
\end{tikzcd}
\]
et on définit la flèche cherchée $\beta_X$ \textit{dans} $(\Cat/S)(\Wd_S)^{-1}$ par la commutativité de ce triangle, donc comme
\[
\beta_X = \pr_1^X\circ\underbrace{\pi_X^{-1}}_{\text{inverse dans }(\Cat/S)(\Wd_S)^{-1}} \qquad\qquad .
\]

\leavevmode\msp{page 143}Mais ça ne marche pas - cf. annotation! On va faire alors l'inverse
\[
\begin{tikzcd}
\Phi_S(X)\times_S S' \arrow[r, "\pr_1^{\Phi_S(X)}"] \arrow[dr, dotted, "\beta_X ?"'] & \Phi_S(X) \\
 & X \arrow[u, "{i_X \qquad \text{dans } \Wd_S}"']
\end{tikzcd}
\]

\dgrblock{dgrb240}{A pullback can destroy the replacement equivalence}
\begin{dcomment}{examples}{p142}{A pullback can destroy the replacement equivalence}
Grothendieck's first attempt to construct the counit asks whether the map $i_X$ can be inverted after an arbitrary pullback. The walking arrow shows exactly where that attempt breaks. We work with ordinary weak homotopy equivalences, $W=W_\infty$, so that an empty category cannot be equivalent to a point.

Let $S=[1]=(0\longrightarrow1)$ and let $X=\{1\}\longrightarrow S$. An object of $\Phi_SX$ is an arrow of $S$ with target $1$. There are two such objects, $0\longrightarrow1$ and $\id_1$, and one nonidentity morphism between them. The projection identifies $\Phi_SX$ with $S$ over itself. Consequently
\[
 i_X:\{1\}\longrightarrow S
\]
is right-fine: for $s=0$ and $s=1$, the functor $s\backslash X\longrightarrow s\backslash S$ is a weak equivalence between contractible categories.

Now pull back along $j:\{0\}\longrightarrow S$. The resulting map is
\[
 X\times_S\{0\}=\varnothing
 \longrightarrow
 \Phi_SX\times_S\{0\}=\{0\}.
\]
It is not in $W$, nor is it a right-fine equivalence over the point. Regarded again as a map over $S$, it still fails the right-fine test at $0$. Thus the proposed inverse of $\pi_X=i_X\times_S\id_{S'}$ need not exist in either localization relevant to the attempted argument.

The calculation also distinguishes the two inverse images without any abstract machinery. Strict pullback of $X$ to $0$ is empty, whereas derived pullback of $[X]$ to $0$ is a point. The missing object records the arrow $0\longrightarrow1$. A comma fibre retains that arrow; a strict fibre does not.

There is no defect in the replacement itself. The map $i_X$ becomes an isomorphism over $S$, exactly as intended. The invalid move was to apply a functor which had not been shown to preserve that class of equivalences on all of $\Cat/S$. On manuscript page 143 Grothendieck repairs the argument himself: he first projects to $\Phi_SX$, and then inverts $i_X$ in the localization over $S$. No pullback of an arbitrary denominator is required.
\end{dcomment}


(\footnote[70]{Description de $\beta_\xi$ pour $\xi=\hot_{S,W}(X)$, avec $X$ quelconque dans $\Cat/S$, à expliciter dans définition en forme.}), on définit $\beta_X$ comme $(i_X)^{-1}\pr_1^{\Phi_S(X)}$, où $(i_X)^{-1}$ est défini dans $(\Cat/S)(\Wd_S)^{-1}$.

Pour prouver la commutativité de $(*)$ (p. 141), comme $i_{X'} : X'\longrightarrow\Phi_S(X')$ est dans $\Wd_S$, il suffit de prouver, au lieu de $\beta_{X'}\alpha_{X'} = \id_{X'}$, la relation déduite en composant avec $i_{X'} : X'\longrightarrow\Phi_S(X')$
\[
\underbrace{i_{X'}\beta_{X'}}_{\pr_1^{\Phi_S(X')}}\alpha_{X'} = i_{X'},
\]
i.e.
\[
\pr_1^{\Phi_S(X')}\circ\alpha_{X'} = i_{X'},
\]
ce qui n'est autre que la relation de définition de $\alpha_{X'}$ (cf. $(**)$ p. 141), on gagne!

4) \textit{Compatibilité}
\[
\begin{tikzcd}
f^* \arrow[r, "\alpha*f^*"] \arrow[rr, bend right, "\id"'] & f^*f_!f^* \arrow[r, "f^**\beta"] & f^*.
\end{tikzcd}
\]
% foot: Version 9 novembre 2001

\printedpage{88}
% head: A. Grothendieck : Dérivateurs, Ch. VI, Édition des Universités de Montpellier II et Paris VII
\leavevmode\msp{page 144}Partons donc d'un $\xi\in\Ob\HOT_W(S)$, qu'on peut supposer défini par $X\in\Ob\underline{\mathrm{Fib}}\,S$. Alors
\[
\begin{aligned}
f^*(\xi) &\;=\; \hot_{S',W}(X\times_S S'),\\
f_!f^*(\xi) &\;=\; \hot_{S,W}(X\times_S S'),\\
f^*f_!f^*(\xi) &\;=\; \hot_{S',W}(\Phi_S(X\times_S S')\times_S S').
\end{aligned}
\]
Il faut expliciter maintenant
\[
\underbrace{f^*(\xi)}_{\hot_{S',W}(X')} \xrightarrow{\alpha_{f^*(\xi)}} \underbrace{f^*f_!f^*(\xi)}_{\hot_{S',W}(\Phi_S(X')\times_S S')} \xrightarrow{f^*(\beta_\xi)} \underbrace{f^*(\xi)}_{\hot_{S',W}(X')} \quad ,
\]
et vérifier que le composé est l'identité. On a posé
\[
X' = X\times_S S',
\]
et le diagramme est donc aussi
\[
\begin{tikzcd}
X' \arrow[r, "{\alpha_{X'}}"] \arrow[rr, bend left, "{\id_{X'}}"] & \Phi_S(X')\times_S S' \arrow[r, "f^*(\beta_\xi)"] & X' % Reading verified against PDF leaf 50: the curved identity arrow is solid.
\end{tikzcd}
\]
dans $(\Cat/S')(\Wd_{S'})^{-1}$. Pour vérifier la commutativité, on peut composer avec
\[
X\times_S S' = X' \xrightarrow{i_X\times_S\id_{S'}} \Phi_S(X')\times_S S' \qquad\qquad \text{(où } i_X : X\longrightarrow\Phi_S(X)\text{)}
\]
qui est dans $\Wd_{S'}$ (\footnote[71]{\textbf{NB} Comme $X$ est lisse$/S$, et $\Phi_S(X)$ aussi, $X\longrightarrow\Phi_S(X)$ reste une $W$-équivalence colocale par tout changement de base.}), or on vérifie, sur la construction générale de $\beta_\xi=\beta_X$ (page précédente), que l'on a commutativité dans
\[
\begin{tikzcd}
\Phi_S(X')\times_S S' \arrow[r, "f^*(\beta_\xi)"] \arrow[dr, "\Phi_S(\pr_1^X)\times_S S'"'] & X' \arrow[d, "i_X\times_S S'"] \\
 & \Phi_S(X)\times_S S'.
\end{tikzcd}
\]

\dgrblock{dgrb241}{The counit is a fraction over the original base}
\begin{dcomment}{exposition}{p143}{The counit is a fraction over the original base}
Grothendieck's first formula for the counit is correct when the representative $X\to S$ is smooth: strict pullback represents derived inverse image, and projection
\[
 X\times_S S'\longrightarrow X
\]
represents $f_!f^*q_S(X)\to q_S(X)$. To check the first triangle identity, however, one needs the counit on $f_!q_{S'}(X')$, whose natural representative need not be smooth over $S$. The extension to arbitrary representatives is the small but essential issue being worked out on ms.~141--143.

For arbitrary $X$, put $R(X)=\Phi_S(X)\times_S S'$. The replacement map $i_X:X\to\Phi_S(X)$ is invertible in $\HOT_W(S)$, so the valid formula is
\[
 \beta_X=[i_X]^{-1}[\pr_1]:q_S(R(X))\longrightarrow q_S(X).
\]
Both arrows in this fraction are considered over $S$. Their naturality follows from the naturality of $i$ and of projection. If $X$ is smooth, the map $i_X\times_S S'$ is a right-fine equivalence over $S'$, and the natural square with the two projections identifies this general formula with the simple projection formula above.

The abandoned attempt instead inverts $i_X\times_S S'$ without a hypothesis making that pullback a weak equivalence. A different interval calculation shows the problem directly. Take $S=(0\to1)$, $X=\{1\}\to S$, and $S'=\{0\}$. The strict fibre of $X$ over $0$ is empty, while the fibre of $\Phi_S(X)$ is $0\backslash X\simeq e$. The pullback of $i_X$ is therefore $\varnothing\to e$. It is not a weak equivalence, either over $S'$ or after regarding it over $S$; its comma fibre at $0$ already fails. There is no inverse with which to form the proposed fraction.

The distinction is between an inverse in the localization over $S$, which is available, and the pullback of that inverse on arbitrary representatives, which is not represented by pulling back its denominator. The corrected roof leaves the denominator over $S$ and uses projection from a replacement on which inverse image can be computed. This is the same discipline that governs total derived functors: apply the operation to a suitable replacement, then use the comparison isomorphism in the category where it is known to exist.
\end{dcomment}


\leavevmode\msp{page 145}On est ramené à la commutativité de
\[
\begin{tikzcd}
X\times_S S' = X' \arrow[r, "{\alpha_{X'}}"] \arrow[dr, "i_X\times_S S'"', "?"] & \Phi_S(X')\times_S S' \arrow[d, "\Phi_S(\pr_1^X)\times_S S'"] \\
 & \Phi_S(X)\times_S S'
\end{tikzcd}
\tag*{$(*)$}
\]
% foot: Version 9 novembre 2001


\dgrblock{dgrb242}{The last square is naturality}
\begin{dcomment}{exposition}{p145}{The last square is naturality}
The argument has reached the last square before the second triangle identity. No new homotopical input is required. The remaining equality is the naturality of the comma-replacement map, and making that equality explicit finishes the adjunction on these pages.

For the first triangle, start with $X'\to S'$, regarded over $S$ after composition. The defining relation for the unit is
\[
 \pr_1\alpha_{X'}=i_{X'}.
\]
The general formula for the counit gives $\beta_{X'}=i_{X'}^{-1}\pr_1$ in the localization over $S$. Consequently $\beta_{X'}\alpha_{X'}=\id_{X'}$. This is the identity $\beta f_!\circ f_!\alpha=\id_{f_!}$, and its proof applies without any smoothness hypothesis on $X'$ over $S$.

For the second triangle we may represent $\xi$ by a fibration $X\to S$. Put $X'=X\times_S S'$ and $r=\pr_1:X'\to X$. Since $X$ is smooth, the counit at $\xi$ is represented by $r$. To calculate its derived inverse image on the possibly nonsmooth source $X'\to S$, replace that source by $\Phi_S(X')$. Functoriality of $\Phi_S$ then gives the relation
\[
 (i_X\times_S S')\,f^*(\beta_\xi)
 =\Phi_S(r)\times_S S'
\]
in the localization over $S'$. Here $i_X\times_S S'$ is invertible because both $X$ and $\Phi_S(X)$ are smooth over $S$.

It remains to prove
\[
 (\Phi_S(r)\times_S S')\alpha_{X'}=i_X\times_S S'.
\]
These are maps $X'\to\Phi_S(X)\times_S S'$. Their second projections both equal the structure map $X'\to S'$. Their first projections are respectively
\[
 \Phi_S(r)i_{X'}\qquad\text{and}\qquad i_Xr,
\]
which are equal by the naturality of $i:\id\to\Phi_S$. The universal property of the pullback proves the displayed equality already in $\Cat/S'$. Composing the preceding relation with $\alpha_{X'}$ and cancelling $i_X\times_S S'$ now gives
\[
 f^*(\beta_\xi)\alpha_{f^*\xi}=\id_{f^*\xi}.
\]
The two triangle identities establish $f_!\dashv f^*$.

There is a minor but consequential printed index in the intervening display on ms.~144. The map labelled $i_X\times_S\id_{S'}$ has target $\Phi_S(X)\times_S S'$, as in the next triangle, rather than $\Phi_S(X')\times_S S'$. The latter is the target of $\alpha_{X'}$ and the source of $f^*(\beta_\xi)$; identifying them would obscure the naturality square just used. The edition prints the prime, so it is preserved as a source slip. The proof given here completes this last square and makes no claim of originality.
\end{dcomment}
\begin{otherlanguage}{english}
\dgroverview{VI-b}{Chapter VI \textperiodcentered\ \textup{VI-b} \textperiodcentered\ Overview}
\begingroup
\selectlanguage{english}
\noindent\textbf{Relative homotopy types, reflection, and integration.}

These pages occupy a hinge in the construction of $\HOT$. Grothendieck has accumulated several ways of representing a homotopy type over a small category $S$: a presheaf of categories, a split or unsplit fibration, a smooth category over $S$, or simply an arbitrary functor $X\to S$. He now wants to know exactly which localizations make these presentations interchangeable. Once that question is settled, the base may vary and the categories so obtained can begin to carry the operations of a derivator. The movement from a comparison at fixed $S$ to inverse and direct images explains both the long abstract digression and the care devoted to the two triangular identities of an adjunction.

The indispensable distinction is between forgetting weak equivalences in the total category and testing them relative to the base. For a map $u:X\to Y$ over $S$, write
\[
\begin{aligned}
u\in W_S&\quad\Longleftrightarrow\quad u\in W,\\
u\in W_S^{\mathrm d}&\quad\Longleftrightarrow\quad
 (s\backslash u)\in W\quad\hbox{for every }s\in S,\\
u\in W_S^{\mathrm g}&\quad\Longleftrightarrow\quad
 (u/s)\in W\quad\hbox{for every }s\in S.
\end{aligned}
\]
Here $s\backslash X=(s\downarrow X)$ and $X/s=(X\downarrow s)$ are comma categories, not strict fibres. The two fine classes retain respectively contravariant and covariant variation along arrows of $S$. On suitably smooth or proper representatives, the corresponding comma categories may be replaced by strict fibres. On homotopically locally constant representatives, the coarse and both fine tests coincide under the additional fibration hypothesis used in this part of the manuscript. The names of the fifteen categories in the large diagram serve this comparison; they do not describe fifteen different homotopy theories.

Two replacements organize the whole argument. The comma replacement $\Phi_SX$ is a split fibration with fibre $s\backslash X$; its natural map from $X$ is a right-fine equivalence. It therefore changes the representative while retaining the relative homotopy type. The path replacement $C_SX$ has locally constant fibres and is connected to $X$ by a coarse equivalence. It retains the total homotopy type but generally changes the relative one. The distinction is expressed by a reflective localization
\[
\rho_S:\HOT_W(S)\rightleftarrows\HOT_W^{\mathrm{lc}}(S):i_S,
\qquad \rho_S\dashv i_S,
\qquad \rho_Si_S\simeq\mathrm{id}.
\]
The unit $X\to C_SX$ need not be invertible in $\HOT_W(S)$. This is why the second equivalence in Corollary~2 on manuscript page~120, which offers $C_S$ as an alternative to $\Phi_S$ for computing the inverse image of an arbitrary relative type, must be corrected. Its $C_S$ expression computes the inverse image after locally constant reflection. The commentary on the first part of Chapter~VI already explains the basic distinction; the present commentary applies it precisely where the manuscript temporarily loses it.

The opening overlap with VI-a completes the adjunction by comparing $i_{C(X)}$ with $C(i_X)$. Concatenation of paths gives a common retraction, and hence equality after localization. There is also a formal reason: for a coaugmented endofunctor, invertibility of both iterated units forces their compatibility and identifies the essential image as a reflective subcategory. This familiar recognition principle is stated in higher categorical form by \citet[Proposition~5.2.7.4]{Lurie2009HTT}. Thus the path construction supplies particularly useful representatives, but its composition law is stronger data than the adjunction itself needs.

The abstract discussion beginning on manuscript page~85 isolates precisely this point. In a category $\mathcal M$ with weak equivalences $W$ and a class of fibrations, Grothendieck compares the localization of $\mathcal M/S$ with that of its full subcategory of fibrant objects. A functorial replacement gives the equivalence directly. A replacement obtained only by lifting properties requires a separate argument: two choices of lift are compared through a factorization of a relative diagonal, and become equal after localization. The manuscript proves this comparison, rather than merely assuming that choices disappear. Theorem--Scholie~2 then introduces a finer class $W_S^0\subseteq W_S$ agreeing with $W_S$ between fibrant objects and constructs the associated reflection. Our exposition writes out the naturality checks and the independence of choices, and explains why the functorial-replacement case already yields the adjunction without the additional path-composition hypothesis $A'$. This is an application of established deformation theory, not a new localization theorem; compare \citet[Theorem~2.2.11]{Riehl2014}.

There are several distinct uses of the word fibration here. A Grothendieck fibration describes transport of objects in a category. A homotopical fibration controls comma categories and strict fibres. A fibration in a Quillen model structure is part of a lifting and factorization system. They interact, but the manuscript's abstract argument does not identify them. In particular, invoking Thomason's model structure gives a method for ordinary nerve weak equivalences on categories; it does not by itself establish every assertion for every proposed fundamental localizer. Likewise, weak saturation, meaning the formal closure axioms initially imposed on $W$, must be distinguished from strong saturation, which detects membership in $W$ by invertibility in $\mathrm{Cat}[W^{-1}]$. The latter is supplied for fundamental localizers by \citet[Proposition~4.2.4 in the cited author version]{Cisinski2006Presheaves}. It closes the explicit saturation reservation in the converse of Proposition~2 on manuscript pages~124--126.

Once the base varies, raw pullback of an arbitrary $X\to S$ is insufficient. To calculate $f^*$ one first takes a representative, such as $\Phi_SX$, on which strict pullback respects the relevant weak equivalences. Smoothness of $X\to S$ supplies one sufficient condition; properness of the change of base supplies another. The ``sure conjecture'' on manuscript page~122 is not left open in the work: Grothendieck proves it on manuscript pages~168--169 of the same chapter. The commentary separates this later proof from what has actually been established in these pages. For ordinary weak homotopy equivalences, testing the comma replacement of a point also proves the converse: preservation of right-fine equivalences by strict pullback characterizes properness. A three-object span gives the nonproper homotopy fibration requested on manuscript page~126; it preserves coarse equivalences but fails to preserve a particular fine equivalence. These computations make the difference between the two base-change conditions explicit.

A natural transformation $f\Rightarrow g$ also requires attention: it induces $g^*\Rightarrow f^*$ in the presheaf convention, but a literal functor between arbitrary strict fibres need not exist. Split fibrations make that action explicit; localization then removes the dependence on the chosen splitting.

The questions about the locally constant theory lead to a useful distinction in the modern answer. Restriction between categories of locally constant diagrams has both adjoints; this holds for arbitrary derivator coefficients by \citet[Corollary~3.11]{Cisinski2009LocallyConstant}. Nevertheless, the varying-base assignment $S\mapsto\HOT^{\mathrm{lc}}(S)$ does not form a derivator for ordinary homotopy types. The pointwise Kan extension formula already fails for a point included in the walking arrow. Fixing a base $S$ and varying an additional diagram category is a different operation: that construction does give the derivator of locally constant coefficients, as established in the same paper's Theorem~3.13. The distinction is visible without any higher machinery and prevents an attractive but incorrect explanation of Grothendieck's hesitation on manuscript page~128.

Another reservation concerns infinite disjoint unions of indexing categories. The corresponding product statement for arbitrary localizations is false. The bounded-length obstruction that Grothendieck notices can be realized by a countable family of finite fence categories: each factor becomes the terminal groupoid when all arrows are inverted, while the product has objects that no finite zigzag can connect. A refinement using finite categories inside a product with the walking retraction also answers footnote~66: strong saturation alone does not ensure that evaluation detects isomorphisms in localized diagram categories. A uniform three-arrow morphism is invertible at every coordinate, but its coordinate inverses require unbounded word lengths. The diagram shape in this counterexample is countably discrete; no finite-shape conclusion is asserted. The space derivator nevertheless satisfies both pointwise detection and the unrestricted disjoint-union axiom. Model structures and controlled fraction presentations supply the additional uniformity that the abstract questions lack.

The final pages construct integration. Postcomposition with $f:S'\to S$ sends $X'\to S'$ to the same total category over $S$, and preserves the fine weak equivalences. This yields $f_!$ before any choice of fibrant replacement. The unit uses the comma replacement over the new base; the counit for arbitrary $X$ is the localized composite
\[
\Phi_SX\times_S S'\longrightarrow\Phi_SX
  \xrightarrow{\,i_X^{-1}\,}X.
\]
The order matters. Pulling back $i_X$ and then trying to invert it can fail, because strict pullback need not preserve a fine equivalence between arbitrary representatives. Grothendieck notices and repairs that attempt on manuscript pages~142--143. At manuscript page~145 the argument has reached the last square in the second triangle; the commentary completes its verification from naturality. In the case of spaces this integration is homotopy left Kan extension, and for the structural map it recovers homotopy colimit. The present derivator formalism provides its axiomatic setting; see \citet{Groth2013}.

The modern picture may be expressed as coherent diagrams of spaces on $S^{\mathrm{op}}$, with locally constant objects corresponding to diagrams on the homotopy type of $S$. It is the entire homotopy type, not just the ordinary fundamental groupoid together with a functor into the homotopy category, that retains all the coherence. In particular, factoring the underlying diagram through $\Pi S$ detects local constancy but does not reconstruct its coherent object. The comma replacements and their categories of elements retain the additional data needed to form the homotopy colimit and to calculate its maps.
\endgroup
\end{otherlanguage}
\dgrchunkend

% DER_VI-c: Grothendieck's text with the commentary set in, then the overview, inside a group that applies the chunk's own definitions
\dgrpart{1}{VI-c}{Part 3 of 3, manuscript pages 146--204}{VI-c\enspace\textperiodcentered\enspace Chapter VI, part 3 of 3\enspace\textperiodcentered\enspace manuscript pages 146--204}{Chapter VI \textperiodcentered\ \textup{VI-c}}
\dgrchunkbegin
% DER_VI-c: the definitions of the chunk's preamble, applied inside its group
\let\printedpage\relax\newcommand{\printedpage}[1]{\marginpar{\footnotesize p.~#1}}
\let\msp\relax\newcommand{\msp}[1]{\marginpar{\footnotesize [#1]}}
\let\HOT\relax\newcommand{\HOT}{\mathrm{HOT}}
\let\hot\relax\newcommand{\hot}{\mathrm{hot}}
\let\Cat\relax\newcommand{\Cat}{\mathrm{Cat}}
\let\id\relax\newcommand{\id}{\mathrm{id}}
\let\pr\relax\newcommand{\pr}{\mathrm{pr}}
\let\Ob\relax\newcommand{\Ob}{\mathrm{Ob}}
\let\Wd\relax\newcommand{\Wd}{W^{\mathrm{d}}}
\let\Hom\relax\newcommand{\Hom}{\mathrm{Hom}}
\let\uHom\relax\newcommand{\uHom}{\underline{\mathrm{Hom}}}
\let\Liss\relax\newcommand{\Liss}{\underline{\mathrm{Liss}}\,}
\let\cA\relax\newcommand{\cA}{\mathcal{A}}
\let\dd\relax\newcommand{\dd}{\mathrm{d}}
\let\lc\relax\newcommand{\lc}{\mathrm{lc}}
\let\adj\relax\newcommand{\adj}[2]{\mathrel{\underset{#2}{\overset{#1}{\rightleftarrows}}}}
\let\multhot\relax\newcommand{\multhot}{\mathrm{multhot}}
\let\mult\relax\newcommand{\mult}{\mathrm{mult}}
\let\Hot\relax\newcommand{\Hot}{\mathrm{Hot}}
\let\D\relax\newcommand{\D}{\mathbf{D}}
\let\bs\relax\newcommand{\bs}{\backslash}
\let\ul\relax\newcommand{\ul}[1]{\underline{\mathrm{#1}}}
\let\circled\relax\newcommand{\circled}[1]{\tikz[baseline=(c.base)]\node[draw,rounded corners,inner sep=2pt](c){#1};}
\let\Fib\relax\newcommand{\Fib}{\underline{\mathrm{Fib}}}
\let\Parf\relax\newcommand{\Parf}{\underline{\mathrm{Parf}}}
\let\Prop\relax\newcommand{\Prop}{\underline{\mathrm{Prop}}}
\let\Fibsc\relax\newcommand{\Fibsc}{\underline{\mathrm{Fibsc}}}
\let\Cofibsc\relax\newcommand{\Cofibsc}{\underline{\mathrm{Cofibsc}}}
\let\lcs\relax\newcommand{\lcs}{\mathrm{lc}\,!}
\let\cercle\relax\newcommand{\cercle}[1]{\tikz[baseline=(c.base)]\node[draw,circle,inner sep=1pt](c){#1};}
\let\troisbarres\relax\newcommand{\troisbarres}{\leavevmode\llap{$|||$\quad}}
\let\Bifib\relax\newcommand{\Bifib}{\underline{\mathrm{Bifib}}}
\let\margbar\relax\newcommand{\margbar}{\marginpar{\(\|\)}}
\let\und\relax\newcommand{\und}[1]{\underline{\mathrm{#1}}}
\let\circnum\relax\newcommand{\circnum}[1]{\tikz[baseline=(c.base)]\node[draw,circle,inner sep=1pt](c){#1};}

% ---- from Der-ChVI-345_p091-095.tex
\printedpage{89}
% head: A. Grothendieck : Dérivateurs, Ch. VI, Édition des Universités de Montpellier II et Paris VII

dans $(\Cat/S')(W^{\dd}_{S'})^{-1}$. Essayons de prouver que c'est commutatif dans $\Cat/S'$, pas seulement dans le localisé par $W^{\dd}_{S'}$. On peut le regarder comme diagramme dans $\Cat/S$ et prouver la commutativité à ce titre. Cette relation de commutativité $vu = w$ se décompose alors en deux, suivant les deux facteurs $\Phi_S(X)$ et $S'$ de l'objet but $\Phi_S(X) \times_S S'$. Suivant le facteur $S'$, la relation signifie simplement que $\alpha_{X'}$ est un $S'$-morphisme, c'est OK. Pour le facteur $\Phi_S(X)$, on prolonge le diagramme $(*)$ à droite par le carré
\[
\begin{tikzcd}
X' \arrow[r, dotted, "\alpha_{X'}"] \arrow[dr, dotted] & \Phi_S(X') \times_S S' \arrow[r, "\pr_1^{\Phi_S(X')}"] \arrow[d, "\Phi_S(\pr_1 X)\times_S S'"] & \Phi_S(X') \arrow[d, "\Phi_S(\pr_1^X)"] \\
 & \Phi_S(X) \times_S S' \arrow[r, "\pr_1^{\Phi_S(X)}"] & \Phi_S(X),
\end{tikzcd}
\] % [sic] \Phi_S(\pr_1 X)\times_S S' -> \Phi_S(\pr_1^X)\times_S S', the X is printed on the line instead of as superscript of \pr_1
\msp{page 146}
et compte tenu que le composé de la première ligne horizontale est $i_{X'}$, par construction de $\alpha_{X'}$, on est ramené à la commutativité de
\[
\begin{tikzcd}
X \times_S S' = X' \arrow[r, "i_{X'}"] \arrow[d, "i_X\times_S S'"'] & \Phi_S(X') \arrow[d, "\Phi_S(\pr_1^X)"] \\
\Phi_S(X) \times_S S' \arrow[r, "\pr_1^{\Phi_S(X)}"] & \Phi_S(X).
\end{tikzcd}
\tag*{$(*)$}
\]

Or le morphisme $i_X : X \longrightarrow \Phi_S(X)$ dans $\Cat/S$ donne lieu au carré commutatif (comme toute flèche $i : X \longrightarrow Y$)
\[
\begin{tikzcd}
X \times_S S' \arrow[r, "\pr_1^X"] \arrow[d, "i_X\times_S S'"'] & X \arrow[d, "i_X"] \\
\Phi_S(X) \times_S S' \arrow[r, "\pr_1^{\Phi_S(X)}"] & \Phi_S(X),
\end{tikzcd}
\]
de sorte que le composé inférieur de $(*)$ n'est autre que celui de $(**)$, donc égal au composé supérieur de ce dernier, et la commutativité de $(*)$ revient à celle de
\[
\begin{tikzcd}
X \times_S S' = X' \arrow[r, "i_{X'}"] \arrow[d, "\pr_1^X"'] & \Phi_S(X') \arrow[d, "\Phi_S(\pr_1^X)"] \\
X \arrow[r, "i_X"] & \Phi_S(X),
\end{tikzcd}
\]
qui est évidente par fonctorialité de $Z \xrightarrow{i_Z} \Phi_S(Z)$, appliquée à la flèche $\pr_1^X : X' \longrightarrow X$ de $\Cat/S$.

% foot: Version 9 novembre 2001

\printedpage{90}
% head: A. Grothendieck : Dérivateurs, Ch. VI, Édition des Universités de Montpellier II et Paris VII

\msp{page 147}
On a ainsi prouvé que les foncteurs
\[
\HOT_W(S') \adj{f_!}{f^*} \HOT_W(S)
\tag*{$(*)$}
\]
sont adjoints l'un de l'autre, avec des morphismes d'adjonction
\[
\left\{
\begin{array}{lcl}
\id & \xrightarrow{\ \alpha\ } & f^*f_! \\[1ex]
\id & \xleftarrow{\ \beta\ } & f_!f^*
\end{array}
\right.
\]
précisés plus haut. Je présume que la même démonstration, en travaillant avec les localiseurs grossiers $W_S$, $W_{S'}$, au lieu de $W^{\dd}_S$, $W^{\dd}_{S'}$, donnerait que les foncteurs correspondants
\[
\HOT^{\lc}_W(S') \adj{f_!^{\lc}}{f^*_{\lc}} \HOT^{\lc}_W(S)
\tag*{$(**)$}
\]
sont adjoints l'un de l'autre, avec essentiellement ‘les mêmes’ morphismes d'adjonction $\alpha$, $\beta$. Je n'ai pas envie de reprendre les raisonnements précédents et préfère le \textit{déduire} de l'adjonction dans la situation $(*)$. Cela me donnera l'occasion de mieux comprendre les traits formels essentiels de la situation. Je poserai pour simplifier
\msp{page 148}
\[
\begin{array}{rcl@{\qquad\qquad}rcl}
\HOT_W(S) & = & \cA, & \HOT_W(S') & = & \cA', \\
\HOT^{\lc}_W(S) & = & \cA_0, & \HOT^{\lc}_W(S') & = & \cA'_0,
\end{array}
\]
et j'utilise les couples de foncteurs adjoints
\[
\cA \adj{\rho}{i} \cA_0, \qquad\qquad \cA' \adj{\rho'}{i'} \cA'_0,
\]
où $\rho$, $\rho'$ sont des foncteurs de localisation, i.e. $i$, $i'$ des foncteurs pleinement fidèles. Je dispose d'un couple de foncteurs adjoints
\[
\cA' \adj{f_!}{f^*} \cA,
\]
et je suppose de plus

\dgrblock{dgrb243}{The two triangles reduce to ordinary naturality}
\begin{dcomment}{exposition}{p146}{The two triangles reduce to ordinary naturality}
The adjunction proof is complete, although the second triangle occupies several diagrams. Its purpose is to turn the chosen representative formulas into a genuine universal property. We give both checks together, keeping the distinction between identities in a category of categories and identities after localization.

For $X'$ over $S'$, put $Z=\Phi_S(X'/S)$ and let $a:X'\longrightarrow Z\times_SS'$ be the unit representative. The first triangle is immediate from the defining equality $\pr_1a=i_{X'/S}$:
\[
 \beta_{L[X']}L(\alpha_{[X']})
 =[i_{X'/S}]^{-1}[\pr_1][a]=\id_{L[X']}.
\]
This calculation uses a denominator over $S$, where its invertibility was established. In particular, it does not revive the failed pullback argument of page 142.

For the second triangle we may represent an arbitrary object of $\HOT_W(S)$ by a smooth $X\longrightarrow S$. Put $X'=X\times_SS'$ and $p=\pr_1:X'\longrightarrow X$. The map
\[
 r_X=i_X\times_S\id_{S'}:
 X'\longrightarrow\Phi_SX\times_SS'
\]
is now right-fine over $S'$. Here the smoothness of both $X$ and $\Phi_SX$ is essential: right-fine equivalences between these good representatives are preserved by arbitrary base change. This is a property of the representatives, not a property of arbitrary arrows in $\Cat/S$.

Under these identifications, the description of the derived counit says
\[
 [r_X]\circ R(\beta_{[X]})
 =[\Phi_S(p)\times_S\id_{S'}].
\]
The unit $a_{X'}:X'\longrightarrow\Phi_SX'\times_SS'$ therefore satisfies
\[
 [r_X]\circ R(\beta_{[X]})\alpha_{R[X]}
 =[\Phi_S(p)\times_S\id_{S'}][a_{X'}].
\]
The right side is $[r_X]$. To see this before localization, compare its two projections. Projection to $S'$ is the structure map of $X'$. Projection to $\Phi_SX$ gives $\Phi_S(p)i_{X'}$, which equals $i_Xp$ by naturality of $i$. These are exactly the projections of $r_X$. Since $[r_X]$ is invertible, cancellation proves
\[
 R(\beta_{[X]})\alpha_{R[X]}=\id_{R[X]}.
\]

The target of $r_X$ is $\Phi_S(X)\times_SS'$. In the displayed line immediately preceding the first triangle on manuscript page 144 the edition prints $\Phi_S(X')\times_SS'$ instead. The adjacent diagram and the expression $i_X\times_S\id_{S'}$ determine the intended target unambiguously. This is a typographical slip, not a gap in the naturality argument.

We can now use $L\dashv R$ without further representative choices. Its Hom bijection sends $h:L\xi'\longrightarrow\eta$ to $R(h)\alpha_{\xi'}$, and its inverse sends $k:\xi'\longrightarrow R\eta$ to $\beta_\eta L(k)$. The two equations just proved say exactly that these transformations are inverse. They also make the construction independent of an accidental presentation of a morphism by a fraction. The result is stronger than the assertion that direct and inverse image have compatible formulas on objects.
\end{dcomment}


\dgrblock{dgrb244}{Adjunction is the organizing operation}
\begin{dcomment}{context}{p147}{Adjunction is the organizing operation}
Grothendieck's insistence on constructing both adjunction arrows belongs to the original motivation of derivators. In \emph{Pursuing Stacks}, \S69, printed pp.~1--2, he describes the diagram categories as a way to recover structure lost by a single derived category, and identifies the two adjoints to restriction as the operations replacing direct and inverse limits \citep{Grothendieck1983}. The present calculation supplies one of those adjoints in the categorical model of homotopy types.

Its immediate use is already visible. Once the unit and counit exist, a natural transformation between diagrams of indexing categories can be transported across adjunctions. This produces the comparison maps for Kan extensions, rather than merely plausible isomorphisms between their values. The commentary on Chapter~I, manuscript p.~31 (``The base-change map as a mate'') explained that formal operation; manuscript pages 138--146 have now constructed the particular adjunction to which it will be applied.

There is a further gain. An adjunction may pass to a reflective subcategory even when its two functors are not constructed there in the same manner. Direct image descends through a localization, while inverse image restricts to local objects. The next argument separates these two mechanisms. That separation will become decisive when the chapter later distinguishes $\HOT_W^{\mathrm{lc}}$ from $\HOT_W^{\mathrm{lc}!}$.
\end{dcomment}


\begin{enumerate}
\item[a)] $f_!$ est compatible avec les localiseurs de $\cA'$, $\cA$ définis par $\rho'$, $\rho$ respectivement, donc définit un foncteur (bien déterminé à isomorphisme unique près)
\[
f_!^0 : \cA'_0 \longrightarrow \cA_0,
\]
donnant donc lieu à un diagramme commutatif (à isomorphisme fonctoriel près)

% foot: Version 9 novembre 2001

\printedpage{91}
% head: A. Grothendieck : Dérivateurs, Ch. VI, Édition des Universités de Montpellier II et Paris VII
\[
\begin{tikzcd}
\cA' \arrow[r, "\rho'"] \arrow[d, "f_!"'] & \cA'_0 \arrow[d, "f_!^0"] \\
\cA \arrow[r] & \cA_0.
\end{tikzcd}
\tag*{(I)}
\]
\texttt{[plutôt} $\cA \xrightarrow{\rho} \cA_0$\texttt{]}

\item[b)] ‘Dualement’, $f^* : \cA \longrightarrow \cA'$ est compatible avec les sous-catégories pleines images essentielles des foncteurs pleinement fidèles $i$, $i'$, et induit
\msp{page 149}
alors un foncteur, défini à isomorphisme unique près :
\[
f_0^* : \cA_0 \longrightarrow \cA'_0,
\]
donnant lieu à un carré de foncteurs, commutatif à isomorphisme fonctoriel canonique près
\[
\begin{tikzcd}
\cA_0 \arrow[r, "i"] \arrow[d, "f_0^*"'] & \cA \arrow[d, "f^*"] \\
\cA'_0 \arrow[r, "i'"] & \cA'.
\end{tikzcd}
\tag*{(II)}
\]
\end{enumerate}

\textit{Sous ces conditions, le couple de foncteurs}
\[
\cA_0 \adj{f_!^0}{f_0^*} \cA'_0, \qquad \text{\texttt{[plutôt} } \cA'_0 \adj{f_!^0}{f_0^*} \cA_0, \text{\texttt{]}}
\]
\textit{est un couple de foncteurs adjoints.} (NB quand je parle de couples de foncteurs adjoints, celui écrit au-dessus de l'autre est l'adjoint à gauche.)

\textsc{Démonstration}. Soient
\[
\eta_0 \in \Ob \cA_0, \qquad\qquad \xi'_0 \in \Ob \cA'_0,
\]
il faut définir une bijection bifonctorielle
\[
\Hom_{\cA_0}(f_!^0(\xi'_0), \eta_0) \simeq \Hom_{\cA'_0}(\xi'_0, f_0^*(\eta_0))
\]
Pour cela, on va écrire $\xi'_0$ sous la forme
\[
\xi'_0 = \rho'(\xi')\,, \qquad\qquad \xi' \in \Ob \cA',
\]
il suffira d'établir la bijection bifonctorielle en remplaçant $\xi'_0$ par $\rho'(\xi')$, i.e.
\[
\Hom_{\cA_0}(f_!^0(\rho'(\xi')), \eta_0) \simeq \Hom_{\cA'_0}(\rho'(\xi'), f_0^*(\eta_0))
\]

% foot: Version 9 novembre 2001

\printedpage{92}
% head: A. Grothendieck : Dérivateurs, Ch. VI, Édition des Universités de Montpellier II et Paris VII

\msp{page 150}
(En effet, la formule d'adjonction s'interprète comme un isomorphisme de bifoncteurs, ou, ce qui revient au même, de foncteurs $\cA'_0 \mathrel{\underset{\psi}{\overset{\varphi}{\rightrightarrows}}} \cA_0^{\wedge}$, or
\[
\begin{array}{ccc}
\uHom(\cA'_0, \cA_0^{\wedge}) & \longrightarrow & \uHom(\cA', \cA_0^{\wedge}) \\[1ex]
\varphi & \longmapsto & \varphi \circ \rho'
\end{array}
\]
est pleinement fidèle, puisque $\rho'$ est un foncteur de localisation, donc il revient au même de se donner un isomorphisme $\varphi \simeq \psi$, ou $\varphi \circ \rho' \simeq \psi \circ \rho'$.) On doit donc transformer l'expression
\[
\begin{array}{cccc}
 & \Hom_{\cA_0}(\underbrace{f_!^0(\rho'(\xi'))}_{\rho f_!(\xi')}, \eta_0) & \underbrace{\simeq}_{\text{\footnotesize commutativité de (I)}} & \Hom_{\cA_0}(\rho f_!(\xi'), \eta_0) \\[3ex]
\underbrace{\simeq}_{\text{\footnotesize adjonction $\rho$, $i$}} & \Hom_{\cA}(f_!(\xi'), i(\eta_0)) & \underbrace{\simeq}_{\text{\footnotesize adjonction $f_!$, $f^*$}} & \Hom_{\cA'}(\xi', \underbrace{f^*i(\eta_0)}_{i'(f_0^*(\eta_0))}) \\[3ex]
\underbrace{\simeq}_{\text{\footnotesize commutativité de (II)}} & \Hom_{\cA'}(\xi', i'(f_0^*(\eta_0))) & \underbrace{\simeq}_{\text{\footnotesize adjonction $\rho'$, $i'$}} & \Hom_{\cA'_0}(\underbrace{\rho'(\xi')}_{\xi'_0}, f_0^*(\eta_0)),
\end{array}
\]
ce qui est la bijection annoncée.

Je résume l'essentiel des résultats obtenus sur $f_!$, $f_!^{\lc}$ :

\textbf{Théorème 2.} \textit{Soit $f : S' \longrightarrow S$ une flèche dans} $\Cat$\textit{, considérons les foncteurs}
\[
\begin{aligned}
\HOT_W(S') &\adj{f_!}{f^*} \HOT_W(S) \\
\HOT^{\lc}_W(S') &\adj{f_!^{\lc}}{f^*_{\lc}} \HOT^{\lc}_W(S),
\end{aligned}
\]
\msp{page 151}
\textit{où les foncteurs $f_!$, $f_!^{\lc}$ sont décrits dans le théorème} 1 \textit{(p. 136). Les morphismes d'adjonction}
\[
\left\{
\begin{array}{lcl}
\id & \longrightarrow & f^*f_! \\[1ex]
\id & \longleftarrow & f_!f^*
\end{array}
\right.
\]
\textit{sont décrits p. 138 à 140.}

\dgrblock{dgrb245}{Passing an adjunction through two reflections}
\begin{dcomment}{exposition}{p150}{Passing an adjunction through two reflections}
Grothendieck wants the locally constant adjunction without repeating the categorical calculations. The abstract argument is correct, and it is useful to state its hypotheses independently of any identification of local homotopy theories. Suppose
\[
 \rho:\mathcal A\rightleftarrows\mathcal A_0:i,
 \qquad
 \rho':\mathcal A'\rightleftarrows\mathcal A'_0:i'
\]
are adjunctions with $i,i'$ fully faithful, and $L:\mathcal A'\rightleftarrows\mathcal A:R$ is an adjunction. Let $\Sigma$ and $\Sigma'$ be the classes inverted by $\rho$ and $\rho'$. Assume that $L(\Sigma')\subseteq\Sigma$ and that $R$ carries the essential image of $i$ into that of $i'$. We then have functors $\bar L$ and $\bar R$, with natural isomorphisms
\[
 \bar L\rho'\simeq\rho L,
 \qquad i'\bar R\simeq Ri.
\]
The desired adjunction follows from the natural chain
\[
\begin{aligned}
 \Hom_{\mathcal A_0}(\bar L\rho' x',y)
 &\simeq\Hom_{\mathcal A_0}(\rho Lx',y)\\
 &\simeq\Hom_{\mathcal A}(Lx',iy)\\
 &\simeq\Hom_{\mathcal A'}(x',Riy)\\
 &\simeq\Hom_{\mathcal A'}(x',i'\bar Ry)\\
 &\simeq\Hom_{\mathcal A'_0}(\rho'x',\bar Ry).
\end{aligned}
\]
Every object of $\mathcal A'_0$ is isomorphic to one of the form $\rho'x'$. The resulting bijection also descends naturally in that object. Precomposition with a localization is fully faithful on functor categories: a naturality square for an inverted arrow is obtained from the original square by inversion, and then naturality holds for every composite of ordinary arrows and formal inverses. This supplies the step which cannot be replaced by essential surjectivity alone. If one wishes to avoid conventions about the notation $\mathcal A_0^{\wedge}$ in the text, the relevant object is simply the bifunctor on $(\mathcal A'_0)^{\mathrm{op}}\times\mathcal A_0$ appearing in this display.

The two compatibility assumptions actually describe the same condition. This standard observation clarifies why Grothendieck calls them dual. An object $z\in\mathcal A$ belongs to the reflective subcategory if and only if
\[
 \Hom(b,z)\longrightarrow\Hom(a,z)
\]
is bijective for every $a\longrightarrow b$ in $\Sigma$. One direction follows from adjunction. For the converse, apply the condition to the unit $z\longrightarrow i\rho z$. It gives a retraction $i\rho z\longrightarrow z$. Since $i\rho z$ is already local, the same orthogonality shows that the other composite is the identity. Thus the unit is invertible.

If $L$ preserves localization equivalences and $y$ is local, the adjunction identifies
\[
 \Hom(b',Ry)\longrightarrow\Hom(a',Ry)
\]
with $\Hom(Lb',y)\longrightarrow\Hom(La',y)$. This is a bijection for every $a'\longrightarrow b'\in\Sigma'$, so $Ry$ is local. Conversely, if $R$ preserves local objects, the same argument shows that $\rho L(a'\longrightarrow b')$ induces bijections into every object of $\mathcal A_0$. Yoneda then makes it invertible. In particular, one need not verify two unrelated conditions in applications. The basic facts about reflections and fully faithful adjoints are recorded in \citet[Lemma~4.5.11 and Definition~4.5.12]{Riehl2016}; the displayed Hom calculations give the particular descent argument needed here.

This formal lemma should be carried unchanged into the later revision of the chapter. What changes there is its application. A coarse localization of categories over $S$ does not automatically identify with the full subcategory of locally constant objects. The latter comparison requires the additional hypotheses examined on manuscript pages 171--204. Theorem~2 therefore gives the fine adjunction without that difficulty, and gives the locally constant adjunction whenever the asserted reflections and their compatibility have been established. A correct formal argument does not supply a missing identification of its input categories.
\end{dcomment}


On voudrait expliciter la formule sympathique, obtenue en passant
\[
\boxed{\;f_!f^*(\xi) = \hot_{S,W}(S'/S) \times_S \xi\;}\ ,
\]
pour ceci il faut expliciter le deuxième membre. Notons la

\textbf{Proposition 4.} \textit{Il existe un unique foncteur}
\[
(\xi, \eta) \longmapsto \xi \times_S \eta : \HOT_W(S) \times \HOT_W(S) \longrightarrow \HOT_W(S)
\]

% foot: Version 9 novembre 2001

\printedpage{93}
% head: A. Grothendieck : Dérivateurs, Ch. VI, Édition des Universités de Montpellier II et Paris VII

\textit{avec la propriété que pour $X$, $Y$ variables dans} $\Liss S$\textit{, on ait un isomorphisme bifonctoriel en $X$, $Y$}
\[
\hot_{S,W}(X) \times_S \hot_{S,W}(Y) \simeq \hot_{S,W}(X \times_S Y)
\tag*{$(*)$}
\]
(\footnote[72]{\textbf{NB} La multiplication $\times_S$ dans $\HOT_W(S)$ est associative, commutative, unitaire \textit{etc., etc.}, puisque c'est le cas dans $\Liss S$.}).

C'est formel, en utilisant seulement que sur $\Liss S$, le bifoncteur
\[
(X, Y) \longrightarrow X \times_S Y
\]
est compatible avec les localiseurs $W^{\dd}_S \times W^{\dd}_S$ et $W^{\dd}_S$. Pour le voir, on est ramené à voir que si $X \xrightarrow{u} X'$ dans $\Liss S$ est dans $W^{\dd}_S$, alors il en est de même de $X \times_S Y \xrightarrow{\ u\times_S \id_Y\ } X' \times_S Y$, pour tout
\msp{page 152}
$Y$ dans $\Liss S$. En fait, on n'utilise que la lissité de $X$, $X'$, pas celle de $Y$, et c'est immédiat, car si $X \longrightarrow X'$ est dans $W^{\dd}_S$, $X \times_S Y \longrightarrow X' \times_S Y$ est dans $W^{\dd}_Y$, donc aussi dans $W^{\dd}_S$ (cf. théorème 1 appliqué à $Y \longrightarrow S$). Cette remarque suggère en même temps :

\dgrblock{dgrb246}{Deriving the product in both variables}
\begin{dcomment}{exposition}{p151}{Deriving the product in both variables}
The product appears because Grothendieck wants to understand the endofunctor $f_!f^*$. Before identifying it with multiplication by a relative homotopy type, he has to decide which fibre products survive localization. Proposition~4 makes the safe construction: take both representatives smooth over the base.

For smooth $X,Y\longrightarrow S$, their fibre product is again smooth. If $u:X\longrightarrow X_1$ is a right-fine equivalence between smooth objects, pulling it back along an arbitrary $Y\longrightarrow S$ gives a right-fine equivalence over $Y$. Forgetting from $Y$ to $S$ preserves that class. Hence
\[
 X\times_SY\longrightarrow X_1\times_SY
\]
is right-fine over $S$. The same assertion holds in the other variable when the objects varied are smooth. Localizing separately in the two variables therefore produces a bifunctor on $\HOT_W(S)$. The uniqueness asserted in the proposition is uniqueness with its comparison isomorphisms, rather than literal equality of all choices of localized categories.

For arbitrary representatives its value is
\[
 [X]\times_S[Y]
 \simeq[\Phi_SX\times_S\Phi_SY].
\]
If $X$ is smooth, $\Phi_SX$ may be replaced by $X$ in this expression. That fact does not yet allow $\Phi_SY$ to be replaced by $Y$. The two replacements concern different maps. The first is an equivalence between smooth objects, to which the base-change statement applies. The second would require arbitrary pullback along $X\longrightarrow S$ to preserve all right-fine equivalences. Smoothness of $X$ alone does not say this.

Associativity, symmetry and the unit are transported with their coherence maps. Their diagrams commute already for strict fibre products of categories; localization preserves those equalities. The unit is the class of $S\longrightarrow S$. Equivalently, the diagonal--product adjunction and the terminal-object adjunction descend, so the multiplication is the categorical product in the localized category. In the usual homotopy theory its diagram interpretation is simply the pointwise product of space-valued presheaves.

The boxed formula preceding the proposition must consequently be read provisionally. A formula
\[
 f_!f^*(\xi)\simeq[S'/S]\times_S\xi
\]
for arbitrary $f$ would identify a construction deriving only the varying input with one deriving both factors. Grothendieck notices the discrepancy almost immediately. The following counterexample is therefore part of the proof's development, and the corrected corollary is its actual conclusion. Treating the crossed-out assertion as a theorem would obscure the main mathematical distinction these pages establish.
\end{dcomment}


\textbf{Corollaire 1. (Faux)} \textit{La formule bifonctorielle} $(*)$ \textit{reste valable pour $X$ dans} $\Liss S$\textit{, $Y$ quelconque dans} $\Cat/S$\textit{, et aussi dans le cas inverse.} (\footnote[73]{Cet énoncé est sûrement \textit{faux!} Contrexemple p. 153, cf. corollaire 1 rectifié page 154-156.})

Dans le premier cas, considérons
\[
Y \xrightarrow{i_Y} \Phi(Y)
\]
dans $W^{\dd}_S$, donc un isomorphisme dans $(\Cat/S)(W^{\dd}_S)^{-1}$, on aura donc
\[
\begin{array}{lcl}
\hot_{S,W}(X) \times_S\ \hot_{S,W}(Y) & \xrightarrow{\ \sim\ } & \hot_{S,W}(X) \times_S\ \hot_{S,W}(\Phi(Y)) \\[1ex]
 & \underbrace{\xrightarrow{\ \sim\ }}_{\text{\footnotesize car $X$ \textit{et} $\Phi(Y)$ lisses sur $S$}} & \hot_{S,W}(X \times_S \Phi(Y)),
\end{array}
\]
or je dis que
\[
X \times_S Y \xrightarrow{\ X\times_S i_Y\ } X \times_S \Phi(Y)
\]
est encore dans $W^{\dd}_S$. Il suffit de voir que c'est dans $W^{\dd}_X$ (cf. théorème 1), et pour ceci, que le changement de base lisse $X \longrightarrow S$ induit $\Cat/S \longrightarrow \Cat/X$ compatible avec les localiseurs $W^{\dd}_S$, $W^{\dd}_X$. Mais c'est \textit{faux} (prendre $X = \{s\}$).

\msp{page 153}
Montrons que les corollaire est faux. Je vais prendre $Y = \{s\}$, catégorie ponctuelle en $s$, alors $\Phi(Y)$ est la catégorie fibrée à fibres discrètes définie par le foncteur représentable $F$ défini par $s$, donc c'est $S/F$ et, en l'occurence $S/s$: % [sic] les corollaire -> le corollaire, agreement; occurence -> occurrence, spelling
\[
\Phi(\{s\}) \simeq S/s \qquad\quad \text{comme catégorie sur } S.
\]

% foot: Version 9 novembre 2001

% ---- from Der-ChVI-345_p096-100.tex
\printedpage{94}
% head: A. Grothendieck : Dérivateurs, Ch. VI, Édition des Universités de Montpellier II et Paris VII
Si p.ex. $s$ est élément final $e$ de $S$, on trouve
\[\Phi(\underset{S}{e}) \simeq S,\]
donc
\[\hot_{S,W}(X)\times_S \hot_{S,W}(e)=\hot_{S,W}(X),\]
d'autre part
\[\hot_{S,W}(X\times_S e)=\hot_{S,W}(X_e).\]

\msp{154}
Si les deux étaient isomorphes dans $\HOT_W(S)$ par la flèche canonique
\[X_e\longrightarrow X,\]
en appliquant le foncteur $p_!$ (où $p : S\longrightarrow e$), il s'ensuiverait que $X_e\longrightarrow X$ est dans $W$. Si par exemple $S=\mathbf{\Delta}^1$, cela impliquerait que pour tout $X$ lisse sur $S$, $X_1\longrightarrow X$ est dans $W$, alors que normalement la lissité implique que $X_0\longrightarrow X$ est dans $W$. Donc ça reviendrait à dire que $\Liss\,\mathbf{\Delta}^1=\Liss_0\mathbf{\Delta}^1$, d'où aussitôt $\Liss\,S=\Liss_0S$, pour tout $S$ dans Cat, ça se saurait! Le contrexemple le plus trivial est de prendre $X=\{0\}\hookrightarrow S=\mathbf{\Delta}^1$, l'inclusion de l'ouvert ‘générique’, immersion ouverte, donc lisse, la fibre $X_1$ est vide! % [sic] s'ensuiverait -> s'ensuivrait, nonstandard verb form

% marginal mark: handwritten sign in the left margin beside the following statement
\textbf{Corollaire 1 rectifié.} \textit{Si $X$ est dans} $\Cat/S$\textit{, pas nécessairement lisse sur $S$, alors le foncteur}
\[\mult_{X/S} : Y\longmapsto Y\times_S X : \Cat/S\longrightarrow \Cat/S\]
\textit{induit un foncteur}
\[\Liss\,S\longrightarrow \Cat/S\]
\textit{compatible avec les localiseurs $W_S^{\mathrm{d}}$, $W_S^{\mathrm{d}}$, donc induit un foncteur}
\[(\Liss\,S)(W_S^{\mathrm{d}})^{-1}\longrightarrow(\Cat/S)(W_S^{\mathrm{d}})^{-1},\]

\msp{155}
\textit{ou encore (compte tenu de l'équivalence induite par l'inclusion}
\[(\underline{\mathrm{Liss}}\,S)(W_S^{\mathrm{d}})^{-1}\xrightarrow{\ \sim\ }(\Cat/S)(W_S^{\mathrm{d}})^{-1}\quad\textit{),}\] % doubt resolved in DER_VI-c_corrections.csv; input note: underline of Liss -- double underline or single underline
\textit{un foncteur, encore noté} $\multhot_{X/S}$,
\[\multhot_{X/S} : \HOT_W(S)\longrightarrow\HOT_W(S).\]

\dgrblock{dgrb247}{The false corollary on every fibre}
\begin{dcomment}{examples}{p154}{The false corollary on every fibre}
Grothendieck chooses the smallest base on which the two product constructions can differ. We expand his example because it also identifies the endofunctor which replaces the false formula. Let $S=[1]=(0\longrightarrow1)$, let $X=\{0\}\longrightarrow S$, and let $Y=\{1\}\longrightarrow S$.

The inclusion of $0$ is a discrete fibration: an arrow into $0$ can only be its identity. Thus $X$ is smooth. Its comma fibres are
\[
 0\backslash X=e,\qquad1\backslash X=\varnothing.
\]
For $Y$ both comma fibres are points, and $\Phi_SY\cong S$. The relative homotopy type $[Y]$ is therefore the terminal object of $\HOT(S)$. It follows that
\[
 [X]\times_S[Y]\simeq[X].
\]
On the other hand, $X\times_SY=\varnothing$, so
\[
 [X\times_SY]=[\varnothing].
\]
Evaluation at $0$ distinguishes the two objects: it gives a point on the first side and the empty homotopy type on the second. One smooth factor has not made the strict fibre product into the derived product.

There is a useful computation for an arbitrary coefficient. Write a space-valued presheaf on $[1]$ as a map $b:B\longrightarrow A$, with values $F(1)=B$ and $F(0)=A$. It may be represented by a smooth category over $S$, or calculated in the usual space-valued diagram model. For the inclusion $j_1:\{1\}\longrightarrow S$, restriction gives $B$, and left Kan extension returns the constant diagram on $B$. Thus
\[
 j_{1!}j_1^*F=(B\xrightarrow{\id}B),
 \qquad [Y]\times_SF=F=(B\xrightarrow{b}A).
\]
The counit has components $\id_B$ at $1$ and $b$ at $0$. The candidate product formula is an equivalence precisely when $b$ is a weak equivalence. For $F=[X]$, this is the map $\varnothing\longrightarrow e$.

For comparison, the inclusion $j_0:\{0\}\longrightarrow S$ gives
\[
 j_{0!}j_0^*F=(\varnothing\longrightarrow A)
 \simeq[X]\times_SF.
\]
The difference between the endpoints has a categorical explanation. The object $0$ defines a sieve, so left Kan extension of presheaves from it is extension by the empty object. The object $1$ instead has an incoming arrow from outside its image, and that arrow transports $B$ into the value at $0$.

This also proves a sharper warning contained in the corrected corollary. The map $Y\longrightarrow S$ is a right-fine equivalence, so $[Y]\simeq[S]$. Nevertheless the two fixed-factor endofunctors are different: multiplication by the category $S/S$ is the identity, whereas $\multhot_{Y/S}$ sends $B\longrightarrow A$ to the constant diagram on $B$. Thus $\multhot_{Y/S}$ depends on the chosen category $Y$ over $S$, not only on its class in $\HOT(S)$.
\end{dcomment}


Ce foncteur satisfait donc à la relation
\[\multhot_{X/S}(\hot_{S,W}(Y))\qquad\underbrace{\ \simeq\ }_{\text{isomorphisme canonique}}\qquad\hot_{S,W}(X\times_S Y)\]
% foot: Version 9 novembre 2001
\printedpage{95}
% head: A. Grothendieck : Dérivateurs, Ch. VI, Édition des Universités de Montpellier II et Paris VII
% marginal mark: handwritten sign in the left margin beside the following two lines
pour $Y$ \textit{lisse sur $S$,} et pour $Y$ quelconque dans $\Cat/S$,
\[\multhot_{X/S}(\hot_{S,W}(Y))\qquad\underbrace{\ \simeq\ }_{\text{isomorphisme canonique}}\qquad\hot_{S,W}(X\times_S \Phi_S(Y)).\]

Quand $X$ est lui-même lisse sur $S$, ce foncteur n'est autre que la multiplication par $\hot_{S,W}(X)$ :
\[\multhot_{X/S}=\mult_{\hot_{S,W}(X)}\qquad\text{si $X$ \textit{lisse sur $S$.}}\]

% marginal mark: handwritten sign in the left margin beside the following paragraph
Mais sa valeur pour $\hot_{S,W}(Y)$, quand $Y$ n'est \textit{pas} lisse sur $S$, ne doit pas être confondue avec $\hot_{S,W}(X\times_S Y)$.

Ainsi, si $X$ et $Y$ sont dans $\Cat/S$, considérons
\[X\times_S Y\qquad\text{et}\qquad\boxed{\hot_{S,W}(X\times_S Y).}\] % the frame is printed with rounded corners

\begin{enumerate}
\item[a)] Si ni $X$ ni $Y$ ne sont supposés lisses sur $S$, cet élément ‘n'a aucune interprétation naturelle en termes d'opérations sur les catégories $\HOT_W(S)$’ (\footnote[74]{Sauf dans le cas où l'un des deux soit dans $\underline{\mathrm{Fibhot}}\,S$, cf. annotation page suivante.}).

\msp{156}

\dgrblock{dgrb248}{A precise boundary for the product formula on finite posets}
\begin{dcomment}{development}{p155}{A precise boundary for the product formula on finite posets}
The walking-arrow example suggests asking which inclusions permit the formula for every coefficient. In a finite poset the answer can be read directly from the order. The following elementary development extends Grothendieck's example to a complete criterion in that setting. We use ordinary homotopy types throughout, so $W=W_\infty$.

\textbf{Proposition (full subposets and the universal product formula).}\ours
Let $S$ be a finite poset, regarded as a category, and let $u:U\hookrightarrow S$ be the inclusion of a full subposet. In space-valued presheaves put $P_U=u_!1$, where $1$ is the terminal presheaf on $U$. The canonical map
\[
 \theta_F:u_!u^*F\longrightarrow P_U\times F
\]
is an equivalence for every $F:S^{\mathrm{op}}\longrightarrow\mathsf{Spaces}$ if and only if $U$ is a lower ideal: whenever $s\le t$ and $t\in U$, one has $s\in U$.

Here $u_!$ means homotopy left Kan extension along $u^{\mathrm{op}}$. The two components defining $\theta_F$ are $u_!u^*F\longrightarrow u_!1$, induced by the terminal map, and the counit $u_!u^*F\longrightarrow F$. This specifies the comparison whose invertibility is claimed.

\emph{Proof.} For $s\in S$, write
\[
 U_{\ge s}=\{t\in U\mid s\le t\}.
\]
The pointwise formula gives
\[
 (u_!G)(s)\simeq
 \operatorname*{hocolim}_{(U_{\ge s})^{\mathrm{op}}}G(t),
 \qquad
 P_U(s)\simeq B(U_{\ge s}),
\]
where $B$ denotes the nerve's homotopy type. Replacing a category by its opposite does not change this homotopy type.

Suppose first that $U$ is a lower ideal. If $s\notin U$, the category $U_{\ge s}$ is empty, so both $(u_!u^*F)(s)$ and $(P_U\times F)(s)$ are empty. If $s\in U$, the poset $U_{\ge s}$ has initial object $s$, and its opposite has terminal object $s$. The homotopy colimit is therefore $F(s)$; the same argument with constant value $e$ gives $P_U(s)\simeq e$. Under these identifications $\theta_F(s)$ is the identity of $F(s)$. This proves the sufficiency, including the comparison map, at every object of $S$.

For necessity, suppose that $s\notin U$ but $s\le t_0$ for some $t_0\in U$. Take the representable discrete presheaf
\[
 F=h_s,\qquad h_s(t)=\Hom_S(t,s).
\]
For every $t\in U_{\ge s}$, one has $s\le t$ and $t\ne s$. Antisymmetry then excludes $t\le s$, so $h_s(t)=\varnothing$. The diagram computing $(u_!u^*h_s)(s)$ is constantly empty, and its homotopy colimit is empty. By contrast,
\[
 (P_U\times h_s)(s)\simeq B(U_{\ge s})\times e
\]
is nonempty, since $t_0$ is an object of $U_{\ge s}$. Hence $\theta_{h_s}(s)$ is not an equivalence. This contradicts the asserted universal formula and proves necessity.\hfill$\square$

The test uses only a representable presheaf. Consequently, in this finite-poset setting the universal formula can already be tested on the collection of representables $h_s$. For a point inclusion $\{t\}\hookrightarrow S$, it holds for every coefficient exactly when $t$ is minimal. It holds for every point inclusion exactly when $S$ is discrete.

There is also a distinction from absolute weak equivalence. In the chain $0<1<2$, take $U=\{1,2\}$. Both $U$ and $S$ have contractible nerves, and $P_U\simeq1$: each nonempty upper intersection $U_{\ge s}$ has a least element. Nevertheless $u^*h_0$ is empty, so $u_!u^*h_0$ is empty while $P_U\times h_0\simeq h_0$ is not. Equivalence of the total homotopy types of the indexing categories does not remove the variation of an arbitrary coefficient diagram.

The positive sieve calculation is standard Kan-extension calculus; compare \citet[Proposition~1.23, with the indexing categories replaced by their opposites]{Groth2013}. We checked the pointwise and sieve formulations in the cited derivator literature and searched for the corresponding finite-poset converse under projection formula and Frobenius reciprocity. We did not locate this exact elementary criterion there. The proof is included as a continuation of the example, without a claim of priority or of a new general projection-formula theorem.
\end{dcomment}


\item[b)] Si $X$ et $Y$ sont lisses sur $S$, cet élément s'interprète comme
\[\hot_{S,W}(X)\times_S\hot_{S,W}(Y)\]
(\footnote[75]{Mais aussi si l'\textit{un} des deux éléments est dans $\underline{\mathrm{Fibhot}}\,S$ (et moyennant l'axiome L6 bis sur $W$).}).

\item[c)] Si l'un des deux $X$, $Y$ est lisse sur $S$, sans rien supposer sur l'autre, on a une interprétation en termes de $\mult_X$ (si $Y$ lisse), ou $\mult_Y$ (si $X$ lisse), notamment
\[\begin{aligned}
\hot_{S,W}(X\times_S Y)\ &\simeq\ \multhot_X(\hot_{S,W}(Y))\quad\text{si $Y$ lisse sur $S$}\\
&\simeq\ \multhot_Y(\hot_{S,W}(X))\quad\text{si $X$ lisse sur $S$.}
\end{aligned}\]
\end{enumerate}

\textbf{Corollaire 2.} \textit{Soit $f : S'\longrightarrow S$ une flèche dans} Cat\textit{, d'où $f_!$, $f^*$ entre} $\HOT_W(S')$ \textit{et} $\HOT_W(S)$\textit{. On a la relation}
\[f_!f^*(\xi)\simeq\multhot_{S'}(\xi),\]
\textit{donc si $S'$ est lisse sur $S$,}
\[f_!f^*(\xi)\simeq\hot_{S,W}(S')\times_S\xi.\]

Nous sommes maintenant en mesure de vérifier les axiomes Der 4$'$b, Der 5$'$ relatifs aux images directes (I, p.~67, 68), et de plus Der 3b (I, p.~111). L'axiome Der 4$'$ n'est autre que l'existence des adjoints à gauche $f_!$ des $f^*$. L'axiome Der 5$'$ permet le ‘calcul’ des fibres d'un objet $f_!(\xi)$. Il dit qu'une certaine
% foot: Version 9 novembre 2001
\printedpage{96}
% head: A. Grothendieck : Dérivateurs, Ch. VI, Édition des Universités de Montpellier II et Paris VII
\msp{157}
flèche du changement de base, relatif au diagramme commutatif dans Cat
\[\begin{tikzcd}[row sep=large]
S' \arrow[d,"f"'] & s\backslash S' \arrow[l,"q"'] \arrow[d,"f'"] \\
S & e\simeq\{s\} \arrow[l,"p"']
\end{tikzcd}\qquad(s\in\mathrm{Ob}\,S),\]
savoir
\[p^*(f_!(\xi'))\longrightarrow f'_!(q^*(\xi'))\]
est un isomorphisme (pour $\xi'\in\mathbf{D}(S')$). Soit donc
\[\xi'\in\mathrm{Ob}\,\HOT_W(S'),\qquad\xi'=\hot_{S',W}(X')\]
avec $X'$ dans $\Cat/S'$. Comme
\[f_!(\xi')=\hot_{S,W}(X'),\]
et que $X'$ n'a aucune raison d'être lisse sur $S$ (si on ne suppose pas $S'$ lui-même lisse sur $S$, auquel cas il pourrait être utile de prendre $X'$ lisse sur $S'$, donc aussi sur $S$), pour calculer $p^*(f_!(\xi'))$ il n'est pas question de prendre la fibre $X'_s$, mais il faut prendre $s\backslash X'=\Phi_S(X')_s$ :
\[p^*(f_!(\xi'))\simeq\hot_W(s\backslash X').\]
\[\begin{tikzcd}[row sep=large]
X' \arrow[d] & s\backslash X' \arrow[l] \arrow[d] \\
S' \arrow[d] & s\backslash S' \arrow[l,"q"'] \arrow[d] \\
S & s\backslash S \arrow[l]
\end{tikzcd}\]

\dgrblock{dgrb249}{Fixed-factor multiplication is a direct-image composite}
\begin{dcomment}{exposition}{p156}{Fixed-factor multiplication is a direct-image composite}
The corrected corollary does more than impose an extra hypothesis on the false formula. It isolates an operation which remains available for every category $x:X\longrightarrow S$. Its varying argument is derived, while the category $X$ is held fixed:
\[
 M_X([Y])=[X\times_S\Phi_SY].
\]
The base-change assertion used in Proposition~4 makes this independent of the chosen smooth representative of $[Y]$. The example just calculated shows why it is not generally independent of the representative of $[X]$.

This operation has a direct adjunction-theoretic description:
\[
 M_X\simeq x_!x^*.
\]
Indeed, for smooth $Y$ the inverse image $x^*[Y]$ is represented over $X$ by $Y\times_SX$. Direct image along $x$ regards the same total category as lying over $S$. Thus both sides give $[Y\times_SX]$, naturally in $Y$. Since smooth objects present the localization, the identification holds on all of $\HOT_W(S)$. Corollary~2 is the special case $X=S'$.

There is a canonical comparison with the derived product. The map $i_X:X\longrightarrow\Phi_SX$ induces, for smooth $Y$,
\[
 [X\times_SY]\longrightarrow[\Phi_SX\times_SY]
 \simeq[X]\times_S[Y].
\]
The right side derives both inputs; the left side has only derived the varying input. If $X$ is smooth, the comparison is invertible, because it is the pullback of a right-fine equivalence between smooth representatives. Without that hypothesis the comparison can fail, as it does for $X=\{1\}$ over the walking arrow. The later discussion of homotopical fibrations and properness will supply further sufficient hypotheses. Nothing in the present corollary establishes those extensions in advance.

The identity $M_X\simeq x_!x^*$ also supplies structure which is easy to miss if one sees only the fibre-product formula. It is the underlying endofunctor of a comonad. Its counit is the projection map already used in the adjunction proof, and its comultiplication is
\[
 x_!x^*\xrightarrow{x_!\alpha x^*}x_!x^*x_!x^*.
\]
The comonad laws follow from the triangle identities. When $X$ is smooth, this is multiplication by $[X]$ with the diagonal coalgebra structure coming from the cartesian product. This interpretation explains the order of the argument: the adjunction is established first, and the formula for its composite is then refined according to the geometry of the fixed factor.

Thus the three cases listed on manuscript pages 155--156 have different information content. With two smooth factors, one obtains the product of two objects of the homotopy category. With only the varying factor smooth, one obtains the value of $M_X$. For two arbitrary strict representatives, their strict fibre product need not be determined by either localized class. These are statements about the operations available from particular representatives, rather than about different possible definitions of the same product.
\end{dcomment}


Considérons le diagramme de changement de base par $s\backslash S\longrightarrow S$ ci-contre \texttt{[i.e. ci-dessus]}, à carrés cartésiens. On voit d'abord que $q^*(\xi')$ n'est autre que la classe dans $\HOT_W(s\backslash S')$ de l'image inverse de $X'$ par le changement de base $s\backslash S'\longrightarrow S'$. Pour l'assurer, il suffirait de prendre $X'$ lisse sur $S'$ (ce qui
\msp{158}
est licite) - alors la description de l'image inverse est OK pour tout changement de base. Mais on a vu que dans le cas d'un co-isomorphisme local comme $s\backslash S'\longrightarrow S'$, il n'y a pas même besoin de cet artifice - la formule est OK sans restriction de lissité. Ainsi
\[q^*(\xi')=\hot_{s\backslash S',W}(s\backslash X'),\]
% foot: Version 9 novembre 2001
\printedpage{97}
% head: A. Grothendieck : Dérivateurs, Ch. VI, Édition des Universités de Montpellier II et Paris VII
d'où
\[f'_!(q^*(\xi'))\simeq\hot_{S,W}(s\backslash X')=\hot_W(s\backslash X').\]
Donc on a bien trouvé
\[p^*f_!\simeq f'_!q^*.\]

\dgrblock{dgrb250}{The comma square and its actual comparison map}
\begin{dcomment}{exposition}{p157}{The comma square and its actual comparison map}
Grothendieck now verifies the pointwise formula for direct image. This is where the construction becomes a derivator calculation: evaluation at an object of the base is expressed by a homotopy colimit over a category of arrows. The comma replacement makes the formula almost tautological, but the comparison map and the variance deserve to be stated.

Fix $s\in S$ and write $C=s\backslash S'$. Its objects are pairs $(s',a:s\longrightarrow f(s'))$, and a morphism from $(s',a)$ to $(t',b)$ is $v:s'\longrightarrow t'$ such that $f(v)a=b$. Let $q:C\longrightarrow S'$ forget $a$, let $r:C\longrightarrow e$, and let $p:e\longrightarrow S$ select $s$. The square drawn in the text is not strictly commutative. It carries the canonical natural transformation
\[
 \tau:pr\Longrightarrow fq,
 \qquad \tau_{(s',a)}=a.
\]
The two cartesian squares in the larger diagram involving $s\backslash S$ are, by contrast, strict pullback squares. These two kinds of square should not be identified.

With the presheaf convention $\HOT(S)=\operatorname{Ho}(\mathsf{Spaces}^{S^{\mathrm{op}}})$, $\tau$ gives $q^*f^*\longrightarrow r^*p^*$. The standard left-adjoint mate is
\[
 r_!q^*\xrightarrow{r_!q^*\alpha}
 r_!q^*f^*f_!\longrightarrow
 r_!r^*p^*f_!\xrightarrow{\beta}
 p^*f_!.
\]
It points in the reverse direction to the arrow displayed on manuscript page 157. Once exactness is proved, that displayed arrow can be read as its inverse. The direction is fixed by the units, counits and structural $2$-cell; it is not a further hypothesis on $W$. For the general mate construction, see \citet[Proposition~4.3.8]{Riehl2016}.

Let $X'\longrightarrow S'$ represent $\xi'$. The right side $p^*f_!\xi'$ is represented by $s\backslash X'$. On the left, $q$ is a discrete opfibration: an arrow $v:s'\longrightarrow t'$ has the unique lift taking $(s',a)$ to $(t',f(v)a)$. The pullback rule for these co-isomorphisms permits us to compute $q^*\xi'$ by strict pullback even for arbitrary $X'$. That pullback has objects $(a:s\longrightarrow f(s'),x')$ with $p'(x')=s'$, and is canonically isomorphic to $s\backslash X'$. Applying $r_!$ then forgets its map to $C$, so
\[
 r_!q^*\xi'\simeq\hot_W(s\backslash X')
 \simeq p^*f_!\xi'.
\]

These identifications also calculate the mate. An object represented by $(a,x')$ is carried by the unit into the comma replacement with its identity arrow; the structural transformation then supplies $a$; the counit forgets the redundant intermediate object. The resulting object is precisely $(a,x')$ in $s\backslash X'$, and the same description sends a morphism to itself. Thus the mate is the identity under these categorical identifications. This proves the canonical comparison is invertible, rather than merely exhibiting an unspecified isomorphism between its source and target.

In modern diagram notation the formula reads
\[
 (f_!F)(s)\simeq
 \operatorname*{hocolim}_{(s\backslash f)^{\mathrm{op}}}F(s').
\]
The opposite is forced by the presheaf convention. Modern presentations using covariant diagrams write the corresponding over-category formula instead. The pointwise axiom and its comma-square formulation appear in \citet[Definition~1.5 and Proposition~1.26]{Groth2013}; the calculation here supplies them concretely for the categorical model.
\end{dcomment}


On peut maintenant prouver aussi Der 3b. On a vu qu'une fois qu'on dispose des adjoints à gauche $f_!$ des $f^*$, cet axiome s'énonce sous \textit{forme homologique} ainsi : Si
\[f : S'\longrightarrow S\]
est dans $W$, il faut prouver que pour tout
\[\xi\in\HOT_W(e)\overset{\text{déf}}{=}\mathrm{Hot}_W,\]
la flèche correspondante
\[(*)\qquad p'_!p'^*(\xi)\longrightarrow p_!p^*(\xi)\]
est \texttt{[un]} isomorphisme dans $\mathrm{Hot}_W$, où $p : S\longrightarrow e$, $p' : S'\longrightarrow e$ sont les flèches structurales. Or comme tout objet de Cat est lisse sur $e$, le corollaire 2 précédent permet d'identifier cette flèche $(*)$ à
\[\eta'\times_S\xi\longrightarrow\eta\times_S\xi\]
\msp{159}
déduite de
\[\hot_W(f) :\quad\underbrace{\ \eta'\ }_{\overset{\text{déf}}{=}\,\hot_W(S')}\quad\longrightarrow\quad\underbrace{\ \eta\ }_{\overset{\text{déf}}{=}\,\hot_W(S)}\quad.\]
Comme celle-ci est \texttt{[un]} isomorphisme par hypothèse, on gagne.

\dgrblock{dgrb251}{Constant coefficients detect the total homotopy type}
\begin{dcomment}{exposition}{p158}{Constant coefficients detect the total homotopy type}
Having obtained direct images, Grothendieck can express an equivalence of indexing categories by its effect on homology with constant coefficients. This assertion concerns the total homotopy type of the base. It should be kept separate from equivalence of all categories of diagrams over the base.

Let $p:S\longrightarrow e$ and represent $\xi\in\Hot_W$ by a category $T$. Since every category is smooth over a point, the preceding product computation gives
\[
 p_!p^*(\xi)\simeq\hot_W(S)\times\xi.
\]
For $f:S'\longrightarrow S$, with $p'=pf$, the induced comparison is therefore
\[
 \hot_W(S')\times\xi\longrightarrow\hot_W(S)\times\xi.
\]
The subscript $S$ on the multiplication signs in the edition's display on page 158 is a slip: both objects belong to $\Hot_W=\HOT_W(e)$, so this is the product over $e$.

If $f\in W$, its image in $\Hot_W$ is invertible; applying the product functor proves the desired comparison is invertible for every $\xi$. Conversely, if every such comparison is invertible, take $\xi=e$, the product unit. Then $\hot_W(f)$ itself is invertible. Thus the precise condition detected is membership in the strong saturation $\overline W$, the class of arrows inverted by $\Cat\longrightarrow\Hot_W$. When $W$ is already strongly saturated this is simply $f\in W$. The later restatement on manuscript page 204 makes this distinction explicit.

The walking arrow again illustrates the scope. The inclusion $\{1\}\longrightarrow[1]$ is a weak equivalence of nerves. It therefore gives the same homology of constant coefficients as $[1]\longrightarrow e$. Restriction of arbitrary presheaves along that inclusion, however, remembers only $B$ from a diagram $B\longrightarrow A$. It cannot be an equivalence of all diagram categories: the diagrams $\varnothing\longrightarrow e$ and $\varnothing\longrightarrow\varnothing$ have the same restriction but are distinguished at $0$.

The stronger statement about \emph{locally constant} diagrams is different. On $[1]$, local constancy requires $B\longrightarrow A$ to be an equivalence, so evaluation at either endpoint loses no information up to coherent equivalence. This small example explains why the next page strengthens the conclusion only on $\HOT_W^{\mathrm{lc}}$. It also explains why the axiom controlling the relation between coarse equivalence and fibrewise equivalence will enter that proof. The elementary product argument for constant coefficients does not itself establish that stronger assertion.
\end{dcomment}


Donc on a prouvé

\textbf{Théorème 3.} \textit{Le prédérivateur} HOT \textit{satisfait aux axiomes} Der 1, Der 2 (\footnote[76]{\textit{pas prouvé}, le ‘résultat’ auquel je pensais est canulé, cf. complément proposition 5, page 164.}) \textit{(pour mémoire), et} Der 3b (\footnote[77]{Et même la variante la plus forte Der 3b$'_2$ de Der 3, cf. démonstration page suivante. \textbf{NB} Pour Der 3b$'_2$, on a utilisé l'axiome L6 sur $W$ - et c'est essentiel. On voit que la validité de Der 3b$'_2$ implique L6, et même la forme plus forte L6 bis, et même si on suppose seulement Der 3b$'_1$, i.e. $f_0^*$ pleinement fidèle.}), Der 4$'$, Der 5$'$\textit{.} % [sic] canulé (footnote 76) -> annulé (?), not a French word

On a envie de prouver aussi les variantes renforcées de Der 3, la forme la plus forte étant celle-ci (Der 3b$'_2$) : si $f\in W$, alors $f^*$ induit une \textit{équivalence} de catégories
\[f_0^* : \HOT^{\mathrm{lc}}_W(S)\longrightarrow\HOT^{\mathrm{lc}}_W(S').\]

La pleine fidélité du foncteur $f_0^*$ doit résulter de la variante naturelle Der 3b$''_{\mathrm{coh}}$ de Der 3b, une fois qu'on aura développé le formalisme des Hom internes dans les $\HOT_W(S)$. Mais pour donner un sens à Der 3b$''_{\mathrm{coh}}$, on a besoin des foncteurs \textit{cohomologie,} i.e. de $p_*$ et $p'_*$ (au lieu de $p_!$, $p'_!$) pour HOT, donc c'est encore prématuré pour l'instant. La construction des $f_*$ et celle des Hom internes se fera d'ailleurs sûrement $\pm$ conjointement.
% foot: Version 9 novembre 2001

\printedpage{98}
% head: A. Grothendieck : Dérivateurs, Ch. VI, Édition des Universités de Montpellier II et Paris VII
\msp{160}
Mais je m'aperçois à l'instant qu'on trouve quasiment ‘à l'œil’ le fait que $f_0^*$ soit bien une \textit{équivalence} de catégories, et que $f_!^0$ soit l'équivalence quasi-inverse, en utilisant l'adjonction des foncteurs $f_!^0$, $f_0^*$, et se rappelant que pour prouver que ce sont des équivalences quasi-inverses l'une de l'autre, il revient au même de montrer que les morphismes d'adjonction
\[\begin{aligned}
\mathrm{id}&\quad\xrightarrow{\ \alpha_0\ }\quad f_0^*f_!^0\\
\mathrm{id}&\quad\xleftarrow{\ \beta_0\ }\quad f_!^0f_0^*
\end{aligned}\]
sont \texttt{[des]} isomorphismes. Je vais essayer de décrire ces flèches d'adjonction dans le cas général (sans supposer $f\in W$), puis de montrer qu'elles sont dans $W$ si $f$ l'est.

\dgrblock{dgrb252}{What Theorem 3 establishes at this point}
\begin{dcomment}{status}{p159}{What Theorem 3 establishes at this point}
Grothendieck has proved the left-adjoint and pointwise formulas, and has checked homological invariance for constant coefficients. These are substantial conclusions. The annotations to Theorem~3 also record two separate qualifications which must remain part of its statement.

First, the claim of Der~2 is withdrawn in footnote~76. A family of evaluations is conservative if a morphism in the localized diagram category is invertible whenever all its evaluations are invertible. Defining weak equivalences objectwise does not, by itself, prove this statement for arbitrary morphisms of a localization. One still has to control how such morphisms are represented, or compare the localization with a model in which objectwise detection can be proved. The supplement on manuscript page 164 belongs to that unfinished part of the argument. It would be inaccurate either to count the parenthetical ``pour mémoire'' as a proof or to infer that conservativity is false in ordinary homotopy theory.

For the standard model the question is settled. \citet[Theorem~6.11, including its proof of Der~2]{Cisinski2003} constructs homotopy right Kan extensions for model categories with small limits and proves the corresponding derivator axioms; his Proposition~6.13 constructs the space-valued external Hom. For combinatorial model categories, \citet[Proposition~1.30]{Groth2013} gives a strong derivator. These results justify the usual space-valued diagram interpretation used in the examples above.

For every basic localizer, Der~2 now follows from \citet[Propositions~3.1.30 and~4.2.7]{Cisinski2006}: conservativity for presheaf localizers transfers through a diagramwise equivalence using $\Delta$ (4.1.29). Thus the revised relative W$(1,2,3)$ suffice, without accessibility or W(4). This later repair preserves the historical withdrawal; it makes no assertion for arbitrary weak basic localizers.

Second, the strengthened assertion that a $W$-equivalence of bases induces an equivalence of locally constant categories uses additional control of fibres. The annotation names the earlier axiom L6, and the following pages examine exactly where it is needed. For ordinary weak equivalences, homotopy invariance of locally constant model-category diagrams is proved by \citet[Theorem~1.3]{Cisinski2009LocallyConstant}.

The final paragraph about constructing $f_*$ together with internal Hom is therefore a programme at this place in the text. The adjunction proved above supplies $f_!$, but it supplies no right adjoint to $f^*$ by formal symmetry: taking opposites of indexing categories is not the same as changing a left adjoint into a right adjoint in the same coefficient theory. The later model-category constructions give both operations in their stated settings. The localizer-by-localizer analysis in the remainder of this chapter remains necessary to determine what follows from Grothendieck's chosen axioms alone.
\end{dcomment}


\begin{enumerate}
\item[\texttt{[1)}] \texttt{Soit]} $\xi=\hot^{\mathrm{lc}}_{S,W}(X)$ dans $\HOT_W(S)$. On peut supposer $X\in\mathrm{Ob}\,\underline{\mathrm{Parf}}\,S$. Comme \texttt{[‘Comme’ barré]} $X$ lisse sur $S$, on a
\[f_0^*(\xi)=\hot^{\mathrm{lc}}_{S',W}(X\times_S S'),\qquad f_!^0f_0^*(\xi)=\hot^{\mathrm{lc}}_{S,W}(X\times_S S').\]
Donc
\[\beta_{0,\xi} : f_!^0f_0^*(\xi)\longrightarrow\xi\]
est une flèche
\[X\times_S S'\longrightarrow X\qquad\text{dans }(\Cat/S)(W_S)^{-1}\]
et je présume que c'est $\mathrm{pr}_1$ (\footnote[78]{\textbf{NB} On voit que l'axiome L6 est $\pm$ \textit{nécessaire} pour que le résultat sur $f_0^*$, $f_!^0$ soit valable.}). Est-ce dans $W$ si $f : S'\longrightarrow S$ l'est? Oui (car $X$ est parfait

\msp{161}

sur $S$), sous réserve que $W$ satisfait l'axiome L6 (cf. page 43).

\item[2)] Soit
\[\xi'=\hot^{\mathrm{lc}}_{S',W}(X')\qquad\text{dans }\HOT_W(S'),\]
on a
\[f_!^0(\xi')=\hot^{\mathrm{lc}}_{S,W}(X'),\]
mais pour expliciter $f_0^*$ de cet élément, comme $X'$ n'a aucune raison d'être un Hot-fibré sur $S$ (auquel cas il suffirait de prendre $f^*(X')\overset{\text{déf}}{=}X\times_S S'$), il faut d'abord prendre $C_S(X')$, puis appliquer $f^*$ :
\[f_0^*f_!^0(X')=\hot^{\mathrm{lc}}_{S',W}(C_S(X')\times_S S').\]

Il faut donc définir
\[X'\xrightarrow{\beta_{0,X'}}C_S(X')\times_S S'\qquad\text{dans }(\Cat/S')W_{S'}^{-1}.\]

Or on trouve une flèche canonique dans $\Cat/S'$, déduite du $S$-morphisme $i_{X'}$ : $X'\longrightarrow C_S(X')$, i.e. du carré commutatif en traits pleins
\end{enumerate}
% foot: Version 9 novembre 2001

% ---- from Der-ChVI-345_p101-105.tex
\noindent\printedpage{99}%
% head: A. Grothendieck : Dérivateurs, Ch. VI, Édition des Universités de Montpellier II et Paris VII
\[
\begin{tikzcd}
X' \arrow[rr, "i_{X'}"] \arrow[ddd, "p_{X'}"'] \arrow[dr, dotted, "\beta_{0,X'}"'] & & C_S(X') \arrow[ddd] \\
& C_S(X') \times_S S' \arrow[ur, dotted, "\pr_1"'] \arrow[ddl, dotted] & \\
& \text{\footnotesize commutatif} & \\
S' \arrow[rr, "f"] & & S,
\end{tikzcd}
\]

\begin{quote}
comme l'unique flèche pointillée $\beta_{0,X'}$ qui soit un $S'$-morphisme (triangle gauche commutatif) et qui rende commutatif le triangle supérieur.
\end{quote}

Si $f \in W$, alors $\pr_1$, déduit de $f$ par changement de base $C_S(X') \longrightarrow S$, qui est parfait, est dans $W$ par L6, on sait qu'il en est de même de $i_{X'}$, donc aussi de $\beta_{0,X'}$, on gagne.

\msp{page 162}Pour terminer les propriétés générales des foncteurs $f_!$ pour $\HOT_W$, je voudrais donner les énoncés raisonnables de changement de base. Ce serait un

\dgrblock{dgrb253}{When a change of base becomes an equivalence}
\begin{dcomment}{exposition}{p161}{When a change of base becomes an equivalence}
Grothendieck now wants to know how much of the base a locally constant coefficient can remember. The adjunction is already available; the task is to determine when its two comparison maps become invertible. This is a stronger assertion than the agreement of cohomology with constant coefficients discussed just before it. A locally constant family has transport, and changing the base must preserve the whole family if restriction is to be an equivalence.

Write $f:S'\to S$, and retain the hypotheses under which the locally constant adjunction was constructed. For a perfect representative $X\to S$, the counit is the projection
\[
 \varepsilon_X:X\times_S S'\longrightarrow X.
\]
The notation $\beta_0$ used for this map agrees with the preceding adjunction. To prove that $\varepsilon_X$ is a $W$-equivalence from the hypothesis $f\in W$, one must know that this particular change of base preserves $W$. This is the role of L6. The fact that a product with a fixed category preserves $W$ is insufficient: $X\to S$ need not be a product projection.

The unit has a different presentation because an arbitrary $X'\to S'$ need not remain a good representative after composition with $f$. Put $P=C_S(X')$. The structural map $P\to S$ is perfect, and the natural map $j:X'\to P$ is a total $W$-equivalence. The map
\[
 a:X'\longrightarrow P\times_S S',\qquad
 \operatorname{pr}_1a=j,
\]
is the unit on the chosen representative. If $f\in W$, L6 gives $\operatorname{pr}_1\in W$, and two-out-of-three gives $a\in W$. Thus both the unit and the counit are inverted. The triangle identities then make $f_!^0$ and $f_0^*$ mutually quasi-inverse equivalences. The label $\beta_{0,X'}$ in the edition is a slip: this second comparison is $\alpha_{0,X'}$. The same pages twice place $\hot^{\mathrm{lc}}$ in an ambient category without the superscript $\mathrm{lc}$; the locally constant category is intended.

The footnote about necessity also has a precise reading. Suppose that the adjunction on locally constant objects is an equivalence for every $f\in W$. Its counit then shows that $X\times_S S'\to X$ is inverted by the relevant localization for every perfect $X\to S$. Applying the total-space functor shows that it is inverted in $\Hot_W$. If $W$ is strongly saturated, it belongs to $W$. Consequently the asserted invariance forces the perfect-pullback stability used in the proof. Without strong saturation the same argument gives membership in the strong saturation of $W$, which explains why the manuscript cautiously writes ``$\pm$ n\'ecessaire''.

This argument will remain useful after the axioms are reorganized on manuscript p.~171. It identifies exactly the property needed at this point, without treating an earlier choice of notation for that property as permanent. It also distinguishes two operations which the succeeding pages repeatedly separate: replacing the total category by a better object over $S$, and testing whether that replacement is still an equivalence on each fibre.
\end{dcomment}


\noindent\textbf{Théorème 4.} \textit{Soit}
\[
\begin{tikzcd}
X \arrow[d, "f"'] & X' \arrow[l, "q"'] \arrow[d, "f'"] \\
S & S' \arrow[l, "p"']
\end{tikzcd}
\]
\textit{un diagramme cartésien dans} $\Cat$\textit{, considérons la flèche de changement de base (dans} $\underline{\Hom}(\HOT_W(X), \HOT_W(S'))$\textit{) :}
\[
p^* f_! \longleftarrow f'_! q^*.
\]
\textit{C'est un isomorphisme dans chacun des deux cas suivants.}

\begin{enumerate}
\item[\textit{a)}] \textit{$f$ est lisse.}
\item[\textit{b)}] \textit{$p$ est propre.}
\end{enumerate}

\textsc{Démonstration.} \textit{Cas a).} On teste avec $\xi$ dans $\HOT_W(X)$, on peut supposer $\xi$ de la forme
\[
\xi = \hot_{X,W}(Z) \qquad \text{avec } Z \in \mathrm{Ob}\, \underline{\mathrm{Liss}}\, X.
\]
Notons que $Z$ est alors aussi lisse sur $S$, puisque $f$ est lisse.
\[
\begin{aligned}
f_!(\xi) &= \hot_{S,W}(Z), &\qquad p^* f_!(\xi) &= \hot_{S',W}(\quad \underbrace{Z \times_S S'}_{\text{$Z'$, cf. diagramme}} \quad) \\
q^*(\xi) &= \hot_{X',W}(Z'), & f'_! q^*(\xi) &= \hot_{S',W}(Z'),
\end{aligned}
\]
OK!

\noindent\printedpage{100}%
% head: A. Grothendieck : Dérivateurs, Ch. VI, Édition des Universités de Montpellier II et Paris VII
\[
\begin{tikzcd}
Z \arrow[d] & Z' \arrow[l] \arrow[d] \\
X \arrow[d, "f"'] & X' \arrow[l, "q"] \arrow[d, "f'"] \\
S & S' \arrow[l, "p"]
\end{tikzcd}
\]
\textit{Cas b).} Il faut admettre que \msp{page 163}le foncteur changement de base par $p$ :
\[
Z \longmapsto Z \times_S S' : \Cat/S \longrightarrow \Cat/S'
\]
est compatible avec les localiseurs $W^{\mathrm{d}}_{S}$, $W^{\mathrm{d}}_{S'}$ (cf. plus bas) - c'est ça, formellemet, la propriété pertinente pour $p$ (\footnote[79]{Ça doit être \textit{équivalent} à $p$ propre. Comparer avec le corollaire 1, page 166.}). Il en résultera que % [sic] formellemet -> formellement, letter omitted
\[
p^* : \HOT_W(S) \longrightarrow \HOT_W(S')
\]
se calcule, pour un élément dans $\HOT_W(S)$ donné par un $Z$ \textit{quelconque} dans $\Cat/S$ (pas nécessairement lisse sur $S$), par
\[
p^*(\hot_{S,W}(Z)) = \hot_{S',W}(Z \times_S S').
\]
Partons alors, comme dans a), d'un élément $\xi$ de $\HOT_W(X)$ donné par un $Z$ de $\Cat/X$ (inutile cette fois de supposer $Z$ lisse), le calcul immédiat précédent établit encore la formule de commutation.

Mais il ne semble pas que la propriété de compatibilité de $X \longmapsto X \times_S S'$, avec l'équivalence colocale sur $S$ et sur $S'$, résulte de façon $\pm$ évidente de la condition que $p : S' \longrightarrow S$ soit propre (\footnote[80]{Mais si, cf. démonstration page 168. Les développements des quatre pages qui suivent sont inutiles par cette démonstration.}). J'ai envie de le déduire de fourbis généraux sur le formalisme des dérivateurs, à savoir du théorème de changement de base des $f_!^{\D}$ pour un morphisme $\D$-propre.

\msp{page 164}Sous forme duale d'un théorème de changement de base pour $f_*^{\D}$, par changement de base $\D$-lisse, c'est fait dans I, p. 56 à 58, et on doit pouvoir l'appliquer ici. Notons que dans la démonstration du théorème loc. cit., on n'utilise que L4 \texttt{[plutôt} Der 4\texttt{]}, i.e. l'existence des $f_*^{\D}$, à l'exclusion des $f_!^{\D}$ (L4$'$ \texttt{[plutôt} Der 4$'$\texttt{]}), plus L5 \texttt{[plutôt} Der 5\texttt{]} bien sûr, donc la forme duale ne fait appel à l'existence des $f_*^{\D}$, seulement des $f_!^{\D}$. On a donc une bonne chance de pouvoir l'appliquer à $\HOT_W$ (où je n'ai pas encore établi l'existence des $f_*^{\mathrm{HOT}}$). % [sic] ne fait appel -> ne fait pas appel, "pas" apparently omitted (sense and "seulement des")

\dgrblock{dgrb254}{The canonical map in the base-change theorem}
\begin{dcomment}{exposition}{p163}{The canonical map in the base-change theorem}
The purpose of Theorem~4 is to turn the concrete model for $f_!$ into a calculational rule. A pullback square of categories gives two ways of transporting a coefficient from $X$ to $S'$. An isomorphism between their values would not by itself suffice: one needs the particular comparison determined by the adjunction, since subsequent compositions of squares use its compatibility with pasting.

For the square in the text, the comparison is the composite
\[
 f'_!q^*\xrightarrow{\ f'_!q^*\eta_f\ }
 f'_!q^*f^*f_!
 \cong f'_!f'^*p^*f_!
 \xrightarrow{\ \varepsilon_{f'}p^*f_!\ }p^*f_!.
\]
Here $q^*f^*\cong f'^*p^*$ is induced by the equality $fq=pf'$. This formula fixes the orientation. In the smooth case choose $Z\to X$ smooth. Then $Z\to S$ is smooth, and all inverse images in this expression can be computed literally. The unit introduces the relevant pullback and the counit projects from it; their composite is the canonical identification between the two copies of $Z\times_S S'$. Thus the calculation on p.~162 proves that the canonical map is invertible, and not merely that its source and target happen to be isomorphic.

The proper case uses a different good feature. If literal pullback along $p$ carries $W_S^{\mathrm d}$ into $W_{S'}^{\mathrm d}$, then for every $Z\to S$ the comparison
\[
 Z\times_S S'\longrightarrow \Phi_S(Z)\times_S S'
\]
is right-fine. Hence literal pullback computes $p^*$ even for a representative which is not smooth over $S$. To finish, one may still choose $Z\to X$ smooth in order to compute $q^*$; this costs nothing, because every coefficient has such a representative. Both sides of base change are again represented by $Z\times_S S'$, and the displayed mate gives its identity after the canonical comparisons. The direct argument on pp.~168--169 proves precisely the preservation statement needed here.

A small example shows why smoothness concerns transport rather than constancy. Let $S=(0\to1)$ and take the presheaf of sets $F$ with
\[
 F(0)=\{a,b\},\qquad F(1)=\{c\},\qquad c\longmapsto a.
\]
Its category of elements $X$ has the arrow $a\to c$ and the isolated object $b$, with the evident map $f:X\to S$. This is a discrete Grothendieck fibration, hence smooth. For the terminal coefficient on $X$, the object $f_!1_X$ is represented by $X\to S$. Pullback to $0$ is the two-point category, whereas pullback to $1$ is the one-point category. Both values agree with base change along the corresponding point inclusion. The transition between the fibres is not a weak equivalence, and no local constancy has been used.

The terminology is inherited from the base-change theorems of algebraic geometry, but the hypotheses here are conditions on categories of lifts. The precise categorical characterization is recorded by \citet[Theorem~1.2]{Maltsiniotis2005Asphericite}; its proper version follows by opposites. In the modern derivator formalism, \citet[Propositions~1.24 and~1.26]{Groth2013} treats the corresponding exact squares for Grothendieck fibrations, opfibrations and comma categories. These references clarify the later language without supplying the missing Der~2 argument in the manuscript's first attempted proof.
\end{dcomment}


\noindent\printedpage{101}%
% head: A. Grothendieck : Dérivateurs, Ch. VI, Édition des Universités de Montpellier II et Paris VII
Notons d'abord

\noindent\textbf{Proposition 5.} \textit{Soit} $f : S' \longrightarrow S$ \textit{une flèche dans} $\Cat$. \textit{Alors $f$ est une} $\HOT_W$\textit{-équivalence} si et seulement si $\hot_W(f)$ \textit{est} \texttt{[un]} \textit{isomorphisme. Si $W$ est fortement saturé, cela signifie donc que $f \in W$ (et en tous cas $f \in W \Longrightarrow f$} \texttt{[est]} \textit{une} $\HOT_W$\textit{-équivalence)} (\footnote[81]{devrait remonter vers le théorème 3, p. 159.}).

C'est contenu dans la démonstration de Der 3b tantôt. Je rappelle que $f$ est une HOT-équivalence, par définition, si et seulement si pour $a, b \in \Hot_W (= \HOT_W(e))$,
\[
\Hom_{\HOT_W(S)}(a_S, b_S) \longrightarrow \Hom_{\HOT_W(S')}(a_{S'}, b_{S'}) \tag*{$(*)$}
\]
est \texttt{[un]} isomorphisme, où $a_S$ etc. désigne $p_S^*(a)$ etc. ($p_S : S \longrightarrow e$). Les formules d'adjonction de $p_{S!}$, $p_S^*$ et de $p_{S'!}$, $p_{S'}^*$ permettent d'interpréter la flèche $(*)$ comme
\[
\Hom_{\HOT_W(S)}(p_{S!}(a_S), b) \longrightarrow \Hom_{\HOT_W(S')}(p_{S'!}(a_{S'}), b),
\]
\msp{page 165}déduit par ‘transposition’ $\Hom(-, b)$ de la flèche
\[
p_{S'!}(a_{S'}) \longrightarrow p_{S!}(a_S)
\]
induite par $f$ - donc il faut prouver que celle-ci est un isomorphisme (c'est la forme que nous avions utilisé tantôt). Or cette flèche s'identifie à
\[
\eta' \times a \longrightarrow \eta \times a
\]
induite par
\[
\hot_W(f) : \underbrace{\eta'}_{= \hot_W(S')} \quad \longrightarrow \quad \underbrace{\eta}_{= \hot_W(S)} \quad ,
\]
et dire que c'est \texttt{[un]} isomorphisme quel que soit $a$ dans $\Hot_W$ revient à dire que $\eta' \longrightarrow \eta$ l'est (puisque la multiplication dans $\Hot_W$ a une élément unité \dots). % [sic] une élément -> un élément, gender agreement

Maintenant je me rappelle que pour avoir l'axiome Der 2 des dérivateurs, je devrais également supposer $W$ \textit{fortement} saturé (i.e. égal au localiseur défini par le foncteur canonique $\Cat \longrightarrow (\Cat)W^{-1} \overset{\text{déf}}{=} \Hot_W$). (C'est ce que je supposais d'emblée, en partant d'un $(\mathcal{M}, W)$ général pour définir un prédérivateur, pour avoir l'axiome Der 2.)

En appliquant le dual du théorème de I, p. 56, on trouve donc le théorème de changement de base pour $f_!$, pour changement de base $f : S' \longrightarrow S$ \textit{$W$-propre.}

% the following paragraph is marked in the left margin by a double vertical rule
\msp{page 166}\marginpar{$\Vert$}Mais la démonstration utilise sûrement Der 2 (critère d'isomorphisme par fibres, dans $\HOT_W(S)$), que je n'ai pas prouvé encore. Elle est incomplète pour le moment. Mais je peux noter dès maintenant le

\dgrblock{dgrb255}{What constant coefficients detect}
\begin{dcomment}{exposition}{p165}{What constant coefficients detect}
Grothendieck interrupts the base-change argument to identify the absolute equivalences detected by $\HOT_W$. This is a useful move: it prevents the phrase ``$\HOT_W$-equivalence'' from concealing a new, unexplained class of maps. The proof uses the adjunction just established and the already available product formula for a constant coefficient.

Let $p_S:S\to e$ and put $\eta=\hot_W(S)$. For $a\in\Hot_W$, write $a_S=p_S^*a$. The representing-category description gives a natural isomorphism
\[
 p_{S!}(a_S)\cong\eta\times a.
\]
Indeed, if $a=\hot_W(T)$, then $a_S$ is represented by $S\times T\to S$ and its left direct image is represented by $S\times T$. This also identifies the morphism induced by $f:S'\to S$: it is $\hot_W(f)\times\id_a$.

Adjunction therefore turns the defining map for a $\HOT_W$-equivalence into
\[
 \Hom_{\Hot_W}(\eta\times a,b)
 \longrightarrow
 \Hom_{\Hot_W}(\eta'\times a,b),
 \qquad \eta'=\hot_W(S').
\]
The base-dependent subscripts in the edition's displayed Hom formula are slips; both Hom sets here are in $\Hot_W$. If $\hot_W(f)$ is an isomorphism, every displayed map is a bijection. Conversely, assume these maps are bijections for all $a,b$, and take $a$ to be the terminal object $1$. The product identifications $\eta\times1\cong\eta$ and $\eta'\times1\cong\eta'$ reduce the assertion to bijectivity of
\[
 \Hom_{\Hot_W}(\eta,b)\longrightarrow\Hom_{\Hot_W}(\eta',b)
 \quad\hbox{for every }b.
\]
Yoneda gives invertibility of $\hot_W(f)$. More explicitly, taking $b=\eta'$ supplies a left inverse of $\hot_W(f)$, and injectivity for $b=\eta$ forces it to be a right inverse. This proves Proposition~5.

Strong saturation enters only when the result is translated back from the localized category to the original class $W$. The proposition without saturation says that the detected maps form $\overline W$, the class inverted by $\Cat\to\Hot_W$. If $W=\overline W$, this is exactly $W$. The proof does not require the detection of all isomorphisms in $\HOT_W(S)$ by point evaluations. In particular, it must not be counted as a proof of Der~2. That latter statement concerns arbitrary arrows of a localized category, including zigzags, whereas Proposition~5 identifies the image of a given functor of indexing categories.

The source itself recognizes the distinction on p.~166. The intended appeal to the abstract proper-base-change theorem still needs Der~2, so the proof by that route is incomplete. Nothing is wrong with Proposition~5, and nothing forces Theorem~4 to wait for a global repair of the derivator axioms: the direct proof on pp.~168--169 works before that repair. Keeping this order of dependence is the reason to retain both the interrupted argument and the author's later replacement of it.
\end{dcomment}


\noindent\printedpage{102}%
% head: A. Grothendieck : Dérivateurs, Ch. VI, Édition des Universités de Montpellier II et Paris VII
\textbf{Corollaire 1} (du théorème 4, p. 162). \textit{Si la flèche} $S' \xrightarrow{p} S$ \textit{dans} $\Cat$ \textit{donne lieu (pour tout} $f : X \longrightarrow S$ \textit{dans} $\Cat/S$\textit{) à la formule de changement de base pour}
\[
f_! : \HOT_W(X) \longrightarrow \HOT_W(S),
\]
\textit{alors le foncteur} $p^* : X \longmapsto X \times_S S'$ \textit{de} $\Cat/S$ \textit{dans} $\Cat/S'$ \textit{est compatible avec les localiseurs} $W^{\mathrm{d}}_{S}$\textit{,} $W^{\mathrm{d}}_{S'}$\textit{.} (L'implication inverse a été vue déjà, cf. pages 161-163.)

Cela résultera du

\noindent\textbf{Lemme.} \textit{Soit} $u : X' \longrightarrow X$ \textit{une flèche dans} $\Cat/S$\textit{. Pour que} $u \in W^{\mathrm{d}}_{S}$\textit{, il faut et il suffit (si $W$} \texttt{[est]} \textit{fortement saturé) que pour tout} $\xi \in \Hot_W$\textit{, la flèche correspondante}
\[
u_* : p'_!(\xi_{X'}) \longrightarrow p_!(\xi_X)
\]
\textit{soit} \texttt{[un]} \textit{isomorphisme dans} $\HOT_W(S)$\textit{.}

C'est suffisant, sous réserve de saturation forte de $W$, car si $u_*$ est \texttt{[un]} isomorphisme, passant aux fibres, on trouve par Der 5$'$
\[
H_\bullet^{\HOT_W}(s\bs X', \xi) \longrightarrow H_\bullet^{\HOT_W}(s\bs X, \xi),
\]
si c'est \texttt{[un]} isomorphisme, alors $s\bs X' \longrightarrow s\bs X$ est par définition une HOT-équivalence, ce qui implique que c'est dans $W$ (proposition 5).

\msp{page 167}C'est nécessaire, car si $\xi = \hot_W(T)$, donc $\xi_S = \hot_{S,W}(T_S)$ (où $T_S = T \times S/S$), et on a donc
\[
p'_!(\xi_{X'}) = \hot_{S,W}(X' \times T), \qquad p_!(\xi_X) = \hot_{S,W}(X \times T),
\]
et il faut prouver que si $u : X' \longrightarrow X$ est dans $W^{\mathrm{d}}_{S}$, alors $\underbrace{X' \times T}_{X' \times_S T_S} \longrightarrow \underbrace{X \times T}_{X \times_S T_S}$ aussi, ce qui signifie simplement que
\[
s\bs(X' \times T) \longrightarrow s\bs(X \times T) \tag*{$(*)$}
\]
est dans $W$ pour tout $s \in S$. Mais le changement de base $s\bs S \longrightarrow S$ commute aux produits fibrés, donc transforme $X \times T = X \times_S T_S$ en $(s\bs X) \times_{s\bs S} T_{(s\bs S)} = s\bs X \times T$, et de même pour $X'$. Donc la flèche s'identifie à
\[
(s\bs X') \times T \longrightarrow (s\bs X) \times T,
\]
déduite de
\[
s\bs X' \longrightarrow s\bs X
\]
par le foncteur $? \times T$. Comme ce foncteur transforme $W$ en $W$, on gagne.

\dgrblock{dgrb256}{A single coefficient already tests right-fine maps}
\begin{dcomment}{exposition}{p167:Comme ce foncteur transforme $W$ en $W$, on gagne.}{A single coefficient already tests right-fine maps}
The lemma on pp.~166--167 is meant to extract a statement about actual functors from a base-change isomorphism. Its family of constant coefficients is natural from the cohomological viewpoint. Once the pointwise formula for $f_!$ is known, however, the terminal coefficient already detects the maps under consideration. This observation shortens the proof and makes its independence from Der~2 explicit.

Assume $W$ is strongly saturated, and let $u:X'\to X$ be a morphism over $S$. Then
\[
 u\in W_S^{\mathrm d}
 \quad\Longleftrightarrow\quad
 \hot_{S,W}(u)\text{ is an isomorphism}.
\]
The forward implication is the definition of the localization. For the converse, apply the evaluation functor $s^*$ to the asserted isomorphism. The pointwise formula identifies its value with
\[
 \hot_W(s\backslash X')\longrightarrow\hot_W(s\backslash X).
\]
Strong saturation makes $s\backslash u$ a member of $W$ for every $s$, which is exactly the definition of $W_S^{\mathrm d}$. This argument applies evaluation to an isomorphism already known to exist; it does not infer a global isomorphism from a family of pointwise isomorphisms. Consequently it uses no unproved conservativity axiom.

Now let $p:X\to S$ and $p':X'\to S$. The image of the terminal coefficient under left direct image is $p_!1_X=\hot_{S,W}(X)$, and the map induced by $u$ is $\hot_{S,W}(u)$. Thus the preceding equivalence is precisely the manuscript's test with $\xi=1$. For a general constant coefficient represented by $T$, the induced map is represented by $u\times\id_T$. Since
\[
 s\backslash(X\times T)\cong(s\backslash X)\times T,
\]
product stability of $W$ proves the implication from $u\in W_S^{\mathrm d}$ to the test for every $\xi$. The larger family of tests is therefore equivalent to the single one, rather than an additional hypothesis.

This gives Corollary~1 with its quantifiers visible. Suppose that base change along $p:S'\to S$ is valid for left direct image along every $X\to S$. If $u:X_1\to X_2$ belongs to $W_S^{\mathrm d}$, its map of pushed-forward terminal coefficients is invertible. Pulling that map back along $p$ preserves invertibility. Naturality of base change identifies the resulting map with the left direct image of the terminal-coefficient map induced by
\[
 u':X_1\times_S S'\longrightarrow X_2\times_S S'.
\]
The single-coefficient test over $S'$ implies $u'\in W_{S'}^{\mathrm d}$. Strong saturation is used at this last conversion to membership in $W$; without it the argument establishes the corresponding statement for the class detected in the localization.

This is a useful instance of the distinction between a criterion for a representative and a criterion for an arbitrary localized morphism. The former is enough for the present base-change theorem. The latter is the broader problem of Der~2, which remains separate.
\end{dcomment}


\noindent\printedpage{103}%
% head: A. Grothendieck : Dérivateurs, Ch. VI, Édition des Universités de Montpellier II et Paris VII
\msp{page 168}Finalement, je m'aperçois que j'ai eu un blocage stupide, pour la question que le changement de base propre soit compatible à la $W$-équivalence colocale. Soit
\[
\begin{tikzcd}[column sep=small]
X \arrow[rr, "f"] \arrow[dr] & & Y \arrow[dl] & \\
& S & & S', \arrow[ll, "g"']
\end{tikzcd}
\]
avec $f \in W^{\mathrm{d}}_{S}$ et $g$ propre, il faut voir que $f' : X' \longrightarrow Y'$ déduit de $f$ par changement de base $S' \longrightarrow S$ est encore dans $W^{\mathrm{d}}_{S'}$. Cela signifie que tout changement de base de la forme $s'\bs S' \longrightarrow S'$ le transforme en $\in W$. Mais soit $s = g(s')$, on a le carré commutatif
\[
\begin{tikzcd}
S & S' \arrow[l] \\
\underbrace{s\bs S}_{= Z} \arrow[u] & \underbrace{s'\bs S'}_{= Z'} \arrow[l] \arrow[u] \ ,
\end{tikzcd}
\]
et il suffit de voir que pour toute flèche $\varphi : \mathcal{X} \longrightarrow \mathcal{Y}$ dans $\Cat/Z$ qui est dans $W$, son image inverse par changement de base $Z' \longrightarrow Z$ est itou. (En appliquant ceci au cas $\mathcal{X} = X \times_S (s\bs S)$, $\mathcal{Y} = Y \times_S (s\bs S)$, $\varphi = s\bs f$.) Or ceci résulte du fait que, par \msp{page 169}hypothèse de propreté, $Z' \longrightarrow Z$ est propre et à fibres Hot-asphériques, donc c'est une Hot-équivalence universelle. Il suffirait même que $g$ soit $W$-propre, pour impliquer que $Z' \longrightarrow Z$ est une $W$-équivalence universelle, i.e. qui reste telle après tout changement de base. Or on a le

\noindent\textbf{Lemme.} Si $g : Z' \longrightarrow Z$ est une $W$-équivalence universelle, alors le foncteur
\[
g^* : X \longmapsto X \times_Z Z' : \Cat/Z \longrightarrow \Cat/Z'
\]
est compatible avec les localiseurs $W_Z$, $W_{Z'}$, i.e. si $f : X \longrightarrow Y$ dans $\Cat/Z$ est dans $W$, de même $f' : X' \longrightarrow Y'$ dans $\Cat/Z'$.
\[
\begin{tikzcd}
X \arrow[dr, "f"'] \arrow[ddr] & & X' \arrow[ll, "g_X"'] \arrow[dr, "f'"] \arrow[ddr] & \\
& Y \arrow[d] & & Y' \arrow[ll, "g_Y"', crossing over] \arrow[d] \\
& Z & & Z' \arrow[ll, "g"']
\end{tikzcd}
\]
Cela est évident sur le carré supérieur du diagramme ci-dessus, où $g_X$, $g_Y$ sont dans $W$ par l'hypthèse sur $g$, donc $f$ est dans $W$ si et seulement si $f'$ l'est, cqfd. % [sic] hypthèse -> hypothèse, letter omitted

\msp{page 170 manque}\texttt{[page 170 manque]}

% ---- from Der-ChVI-345_p106-110.tex
\setcounter{footnote}{81}

\printedpage{104}
% head: A. Grothendieck : Dérivateurs, Ch. VI, Édition des Universités de Montpellier II et Paris VII
% foot: Version 9 novembre 2001

\msp{page 171}

\dgrblock{dgrb257}{Proper base change without a pointwise conservativity axiom}
\begin{dcomment}{exposition}{p169}{Proper base change without a pointwise conservativity axiom}
The short proof on pp.~168--169 succeeds because it turns a relative question into an absolute one over a suitable comma category. Properness supplies a comparison of bases which remains a $W$-equivalence after every pullback. Two-out-of-three can then transport the desired equivalence across that comparison.

We first isolate the elementary part. Let $v:T'\to T$ be a universal $W$-equivalence, meaning that for every $A\to T$ the projection $A\times_TT'\to A$ belongs to $W$. If $h:A\to B$ is a map over $T$, the commutative square
\[
\begin{tikzcd}
 A\times_TT' \arrow[r,"h'"]\arrow[d,"v_A"']&
 B\times_TT'\arrow[d,"v_B"]\\
 A\arrow[r,"h"']&B
\end{tikzcd}
\]
has both vertical arrows in $W$. If $h\in W$, then $hv_A\in W$ and $v_Bh'=hv_A$ forces $h'\in W$. The reverse implication is the same calculation in the other direction. Literal pullback along a universal equivalence therefore preserves and reflects $W$. No localization, no adjunction and no pointwise detection is involved in this lemma.

Take now $g:S'\to S$ proper, or more generally $W$-proper, and fix $s'\in S'$ with $s=g(s')$. Properness gives a universal $W$-equivalence
\[
 v:s'\backslash S'\longrightarrow s\backslash S.
\]
Its fibre over an arrow $s\to t$ is the category of its lifts starting at $s'$. The defining asphericity conditions for properness, together with their stability after base change, give the stated universality. This is the geometric content of the proof; it is not a consequence of the existence of initial objects in the two comma categories alone. Two contractible categories can be connected by a functor which fails to be a universal equivalence.

For $u:X\to Y$ over $S$, form $u_s:s\backslash X\to s\backslash Y$, a map over $s\backslash S$. If $u\in W_S^{\mathrm d}$, then $u_s\in W$. Pullback along $v$ identifies $u_s$ with
\[
 s'\backslash(X\times_SS')\longrightarrow
 s'\backslash(Y\times_SS').
\]
The elementary lemma makes this map a member of $W$. Since $s'$ was arbitrary, literal pullback along $g$ preserves right-fine equivalences. This completes the proper case of Theorem~4 by the model calculation explained above.

For a discrete opfibration, the role of the lift categories is especially transparent. An arrow $g(s')\to t$ has a unique lift with prescribed source $s'$. Hence the induced functor on undercategories is an isomorphism; the preceding comparison needs no homotopy argument at all. For an ordinary Grothendieck opfibration, a chosen cocartesian lift is initial in the corresponding category of lifts, so its nerve is contractible. The general proper case replaces this initial object by the weaker asphericity condition required by $W$.

The manuscript calls the earlier obstruction a ``blocage''. Mathematically, the change is a change of test objects: the universal comparison over $s\backslash S$ carries exactly the information needed for right-fine equivalence. The abandoned four-page route is not a counterexample to proper base change, nor does it expose a defect in the construction of $f_!$.
\end{dcomment}


\subsection*{Nouvelle version de la théorie de $\mathrm{HOT}_W$.}

Je me suis aperçu que l'hypothèse $W \supseteq W_\infty$ (L3 quater, p. 1), que j'avais fini par mettre sur $W$, est artificielle, et que l'hypothèse L5 (p. 43), voire L6, L6 bis (p. 43-44), est inutile pour une grande partie de la théorie de $\mathrm{HOT}_W$. J'ai été ainsi amené à revoir la théorie (i.e. ce que j'en ai fait jusqu'à présent) sous les hypothèses suivantes sur $W$ (cf. XII) :

\begin{itemize}[leftmargin=4em]
\item[W(1)] (\textit{Saturation} de $W$) : comme avant (L1 page 1)
\item[W(2)] (localisation) : c'est l'ancien L4 bis (\footnote{i.e. le L5 de la page 2. L'axiome est repris (sous une forme voisine et équivalente) au bas de la page 43.}).
\item[W(3)] (objet final) : Si $X$ a un objet final, $X$ est $W$-asphérique (ancien L3).
\end{itemize}

Ces trois axiomes sont suffisants pout tout ce que j'avais fait dans Pursuing Stacks, et dans une grande partie de ce qui a été fait ici sur HOT. Parfois cependant il faut l'ancien L 5, que je réénonce ici (comparer XII, théorème \texttt{[numéro du théorème manque]}) : % [sic] pout -> pour, misprint

\dgrblock{dgrb258}{Axioms for the localization and axioms for its fibres}
\begin{dcomment}{exposition}{p171:Ces trois axiomes sont suffisants}{Axioms for the localization and axioms for its fibres}
Grothendieck's new beginning separates the construction of a homotopy theory from the stronger rules which make its total spaces determine its fibres. Much of the preceding work survives under three localizer axioms. The extra fibration axiom enters at identifiable comparison maps, and the next ten pages are an effort to locate those maps rather than to replace the entire construction.

The localization axiom W(2) is a relative Theorem~A statement. For a triangle $X\to Y\to S$, it says that membership of $X/s\to Y/s$ in $W$ for every $s$ implies membership of $X\to Y$ in $W$. Together with the terminal-object axiom and the saturation rules, it gives the familiar local-to-global mechanism. The weakened W(2~bis) keeps only the case in which all the relevant comma categories are $W$-aspherical. The distinction matters when the minimum class $W_\omega$ is introduced: it is defined using this weaker axiom, not by simply repeating the complete list W(1,2,3).

There is a second distinction in the word ``fibration''. The ordinary category $\Fib S$ consists of Grothendieck fibrations over $S$. The category $\Fib_W S$ in the revised diagram consists of weak $W$-fibrations and is defined by a homotopical condition. One equivalent test is this: for every categorically contractible $T\to S$ and every $t\in T$, the inclusion
\[
 X_t\longrightarrow X\times_ST
\]
belongs to $W$. This is the formulation of Chapter~XII, Theorem~3, manuscript pp.~51--52, to which Grothendieck refers here \citep{Grothendieck1990XIIa}. The word ``contractible'' in this test cannot be replaced silently by ``$W$-aspherical''. Doing that asks for a stronger condition, whose relation to the weak condition is precisely part of the fibration-axiom problem.

For an ordinary Grothendieck fibration, the weak condition has a concrete interpretation: each transport functor $u^*:X_{s'}\to X_s$ must belong to $W$. Thus an ordinary fibration can be smooth without being weakly $W$-fibred, as in the two-fibre example above. Conversely, a weak $W$-fibration need not possess chosen cartesian lifts. The comma replacement $\Phi$ is what supplies an ordinary split fibration representing its homotopical data.

W(4b) compares two tests for a morphism between such weakly fibred objects. The global test asks whether $X\to Y$ lies in $W$ as a functor of total categories. The fibre test asks whether each $X_s\to Y_s$ lies in $W$. The implication from the fibre test to the global test already follows from W(2), after comparing fibres with comma fibres. W(4b) asserts the converse. It is therefore a cancellation rule for total spaces over a fixed base, not merely a restatement of two-out-of-three.

The first version of W(4b$'$) restricts that rule to $W_\omega$-perfect objects. This does prove equality of the two tests inside that smaller category. It does not yet compare an arbitrary weakly fibred $X$ with its perfect replacement $C(X)$, because $X$ need not belong to the smaller category. The footnote retracting the claimed sufficiency points to exactly this missing scope in the available argument. The retraction does not establish that the implication is false. With the later minimality and filtered-colimit results explained below, the alternating replacement on pp.~192--195 supplies the missing fine unit; the original restricted condition does then imply W(4b). The commentary at p.~195 gives the transfer argument. Later a different condition receives the same label W(4b$'$); the two formulations must be kept separate throughout the audit.
\end{dcomment}


\begin{itemize}[leftmargin=4em]
\item[W(4b)] Soient $X$, $Y$ $W$-fibrés sur $S$, $f : X \longrightarrow Y$ un $S$-morphisme. Si $f \in W$, alors $\forall s \in \mathrm{Ob}\, S$, $f_s : X_s \longrightarrow Y_s$ est dans $W$.
\end{itemize}

Pour la notion de $W$-fibré (\textit{faible}), j'ai été long à trouver enfin \textit{la} bonne notion, bonne pour un $W$ ne satisfaisant que $\mathrm{W}(1,2,3)$ (et même seulement $\mathrm{W}(1, 2 \text{ bis}, 3)$, où W(2 bis) est l'ancien L4). Je renvoie pour cela à (XII, théorème 3, p. 51), où la notion est bien cernée à coups de conditions équivalentes. En fait, pour les applications principales à $\mathrm{HOT}_W$, vues jusqu'à présent, il suffit (\footnote{C'est finalement une erreur, semble-t-il à présent. Je crois que l'axiome W(4b$'$) ne sert finalement à strictement rien (si ce n'est en fait à entretenir des idées fausses.)}) d'une variante affaiblie de W(4b) :

\begin{itemize}[leftmargin=4em]
\item[W(4b$'$)] Comme W(4b), mais en exigeant de plus que $X$, $Y$ soient dans $\ul{Parf}_\omega\, S$.
\end{itemize}

\msp{page 172}

Pour la notion de morphismes $W$-propres, $W$-lisses, $W$-parfaits dans Cat, je renvoie à XII, partie III (p. 22-52), et les notations $\ul{Liss}_W S$, $\ul{Prop}_W S$, $\ul{Parf}_W S$ (sous-catégories de $\mathrm{Cat}/S$), et $\mathrm{Parf}_W \subseteq \mathrm{Fl}(\mathrm{Cat})$. Dans l'énoncé de W(4b$'$) on a posé
\[
\ul{Parf}_\omega S = \ul{Parf}_{W_\omega} S, \quad
\begin{array}[t]{l}
\text{où } W_\omega \subseteq \mathrm{Fl}(\mathrm{Cat}) \text{ est la plus petite partie de } \mathrm{Fl}(\mathrm{Cat})\\
\text{satisfaisant } \mathrm{W}(1, 2\mathrm{bis}, 3).
\end{array}
\]

Je note de même
\[
W'_\omega \subseteq \mathrm{Fl}(\mathrm{Cat}), \qquad \text{plus petite partie de } \mathrm{Fl}(\mathrm{Cat}) \text{ satifaisant } \mathrm{W}(1, 2, 3, 4\mathrm{b}). % [sic] satifaisant -> satisfaisant, misprint
\]

\printedpage{105}
% head: A. Grothendieck : Dérivateurs, Ch. VI, Édition des Universités de Montpellier II et Paris VII
% foot: Version 9 novembre 2001

L'hypothèse $W \supseteq W_\infty$ que nous avions fini par imposer sur $W$, et qui est apparu artificielle au cours de la réflexion du XII, a avantage à être remplacée par % [sic] apparu -> apparue, agreement with « l'hypothèse »
\[
W \supseteq W'_\omega \qquad (\text{ou par } W \supseteq \overline{W'_\omega}, \text{ où } \overline{W'_\omega} \text{ est le saturé fort de } W'_\omega)
\]
(\footnote{$\overline{W'_\omega}$ est le plus petit sous-ensemble \textit{fortement saturé} de $\mathrm{Fl}(\mathrm{Cat})$, satisfaisant W(4b), et il satisfait même la condition W$'$(4a), variante la plus forte des axiomes du type W(4).}). Mais nous ne l'imposons pas comme condition impérative, au même titre que $\mathrm{W}(1,2,3)$, pas plus que W(4b). Le rôle que j'ai fait jouer à $W_\infty$ dans les semaines passées, semble revenir plus naturellement à $W'_\omega$ ou $\overline{W'_\omega}$. Il en est ainsi notamment aussi, il me semble, dans l'énoncé des conditions du type Der 3 sur un dérivateur (I, p. 111-112), qu'il faudrait revoir dans cette lumière. Mais tout dépend, si (comme j'en ai l'impression) on peut faire vraiment de l'algèbre homotopique substantielle, avec notamment les suites exactes longues de fibration et de suspension, pour $\mathrm{HOT}_W$, avec des $W$ qui sont \textit{strictement} contenus dans $W_\infty$. Avant de même être assuré, il est trop tôt encore pour me prononcer sur la `bonne' formulation de Der 3,

\dgrblock{dgrb259}{The proposed theory below ordinary weak equivalence}
\begin{dcomment}{later}{p172:Le rôle que j'ai fait jouer}{The proposed theory below ordinary weak equivalence}
Grothendieck is trying to remove an externally imposed reference to classical homotopy. His question about a localizer strictly smaller than $W_\infty$ concerns the strength of the axioms themselves: perhaps their formal consequences already contain more classical homotopy theory than he has yet extracted.

For the modern basic-localizer axioms, this question is settled. \citet[Definition~1.1.3 and Theorem~2.2.11]{Cisinski2004Minimal} proves that every basic localizer contains the nerve weak equivalences $W_\infty$. Thus no strictly smaller basic localizer satisfies that complete package. The weaker list defining $W_\omega$ requires a further theorem. \citet[Theorem~1.10 and Remark~1.11]{Cisinski2009LocallyConstant} proves the corresponding minimality among weak basic localizers, using the absolute Theorem~A condition. It follows that $W_\omega=W_\infty$ as well. To check the saturation hypothesis, if $ri=\id$ and $ir\in W$, the manuscript's condition on the two composites $ri,ir$ gives $i,r\in W$; this is the retract clause in the weak-basic-localizer formulation. W(2~bis) and W(3) supply its other two clauses. The two axiom lists remain different even though their minimum classes now have the same description.

This result changes the status of the proposed example, but it leaves the reason for the revision intact. Removing an explicit hypothesis $W_\infty\subseteq W$ can still clarify which finite arguments use only the localizer axioms and which use a separate fibre-detection condition. Indeed, a class larger than $W_\infty$ can satisfy the basic axioms while failing W(4b). The connected-components localizer used below gives an elementary instance. The ordinary theory is minimal in the order of classes of equivalences; further localization is entirely possible, and it can forget information that a fibrewise test retains.

The same minimum is obtained after adjoining W(4b): $W'_\omega=W_\infty$. Indeed, Quillen's Theorem~B identifies the fibres of a weak $W_\infty$-fibration with its homotopy fibres \citep[Theorem~2.3.2]{Cisinski2004Minimal}. A total weak equivalence over a fixed base therefore induces weak equivalences on the fibres. Thus $W_\infty$ satisfies W(4b) and is an admissible candidate for $W'_\omega$; basic minimality gives the reverse inclusion. Its strong saturation is the same class. The homotopy-pullback description also supplies the strongest base-change property intended in footnote~84.

This identifies the minimum within the standing W(1,2,3) framework. The inclusion $W'_\omega\subseteq W$ does not, by itself, imply that $W$ inherits every closure property appearing in the defining list: the notion of weak $W$-fibration changes when $W$ changes. The fibre-cancellation axiom is not simply an upward-closed property of subsets of $\mathrm{Fl}(\Cat)$. This explains why Grothendieck treats the suggested containment and the assertion W(4b) as different conditions in the paragraph itself.

By \citet[Corollary~4.2.22]{Cisinski2006}, W(6~bis) is automatic for basic localizers; the theorem assumes relative W(2).
\end{dcomment}


\msp{page 173}

qui est l'axiome mettant en relation le dérivateur $\mathbf{D}$ avec $W_\infty$, ou avec quelque autre localiseur fondamental.

Je vais faire à présent un résumé de la théorie de $\mathrm{HOT}_W$, revue dans l'optique que je viens de dire : on suppose pour $W$ \textit{les conditions} $\mathrm{W}(1,2,3)$ \textit{sans plus,} sauf mention explicite.

\circled{1} \textbf{Localiseurs principaux} sur $\mathrm{Cat}/S$ et sur les principales catégories de $S$-catégories (cf. diagramme p. 76).

Ils seront tous induits par les trois localiseurs familiers
\[
\begin{tikzcd}[column sep=small]
W_S^{\mathrm{d}} \arrow[dr, hook] & & W_S^{\mathrm{g}} \arrow[dl, hook'] \\
& W_S &
\end{tikzcd}
\]
sur $\mathrm{Cat}/S$. Il y a lieu aussi de regarder $W_S^{\mathrm{f}}$, mais jamais sur $\mathrm{Cat}/S$ tout entier, mais sur d'autres catégories du diagramme, et toujours dans des cas où $W_S^{\mathrm{f}}$ coïncide soit avec $W_S^{\mathrm{d}}$, soit avec $W_S^{\mathrm{g}}$. Mais le point de vue $W_S^{\mathrm{f}}$ est souvent plus directement commode, et a l'avantage de plus d'être symétrique.

\circled{a} Dans $\ul{Liss}_W S$, on a

\printedpage{106}
% head: A. Grothendieck : Dérivateurs, Ch. VI, Édition des Universités de Montpellier II et Paris VII
% foot: Version 9 novembre 2001

\[
\begin{tikzcd}[row sep=large, column sep=0pt]
W_S^{\mathrm{f}} \arrow[d, Rightarrow] & {} = W_S^{\mathrm{d}} = W_S^{\mathrm{univ}} \\
W_S^{\mathrm{g}} \arrow[d, Rightarrow] & \\
W_S &
\end{tikzcd}
\]

Dualement,

\circled{b} dans $\ul{Prop}_W S$, on a
\[
\begin{tikzcd}[row sep=large, column sep=0pt]
W_S^{\mathrm{f}} \arrow[d, Rightarrow] & {} = W_S^{\mathrm{g}} = W_S^{\mathrm{univ}} \\
W_S^{\mathrm{d}} \arrow[d, Rightarrow] & \\
W_S, &
\end{tikzcd}
\]

\msp{page 174}

enfin,

\circled{c} dans $\ul{Fib}_W S$, on a
\[
\begin{tikzcd}[row sep=large]
W_S^{\mathrm{f}} \arrow[r, Leftrightarrow] \arrow[d, Rightarrow] & W_S^{\mathrm{d}} \arrow[r, Leftrightarrow] & W_S^{\mathrm{g}} \arrow[r, Leftrightarrow] & W_S^{\mathrm{univ}} \\
W_S. & & &
\end{tikzcd}
\]

Cette implication $W_S^{\mathrm{f}} \Longrightarrow W_S$ est stricte si on ne fait pas sur $W$ l'axiome W(4b), dont l'utilité ici est exactement d'assurer
\[
\left|\;
\begin{array}{l}
W_S^{\mathrm{f}} \Longleftrightarrow W_S \ \text{ dans } \ul{Fib}_W S \ \textit{si}\ \mathrm{W}(4\mathrm{b})\ \texttt{[est]}\ \text{satisfait}\\[1.5ex]
\text{et alors les cinq localiseurs envisagés ci-dessus}\\[1.5ex]
\text{dans } \ul{Fib}_W S \text{ sont tous égaux.}
\end{array}
\right.
\]

Les relations entre [?] localiseurs dans $\ul{Liss}_W S$, $\ul{Prop}_W S$, $\ul{Fib}_W S$ se transportent automatiquement aux sous-catégories de celles-ci, plus généralement à toutes celles qui s'envoient
\printedpage{107}
% head: A. Grothendieck : Dérivateurs, Ch. VI, Édition des Universités de Montpellier II et Paris VII
% foot: Version 9 novembre 2001
par un foncteur d'inclusion dans une de celles-ci. Dans le diagramme p. 76, toutes les catégories, à la seule exception de $\mathrm{Cat}/S$ tout entier, sont dans ce cas.

\dgrblock{dgrb260}{Five tests and the role of good objects over the base}
\begin{dcomment}{exposition}{p174:Les relations entre [?] localiseurs}{Five tests and the role of good objects over the base}
The diagram of localizers is the bookkeeping on which the revised theory depends. Grothendieck is distinguishing the information carried by a strict fibre, a comma fibre, and a total category. To read its equalities, every test must refer to the same map $u:X\to Y$ over the same base $S$.

The five conditions are
\[
\begin{aligned}
 u\in W_S&\iff u\in W\text{ on total categories},\\
 u\in W_S^{\mathrm f}&\iff u_s:X_s\to Y_s\in W\quad(s\in S),\\
 u\in W_S^{\mathrm d}&\iff s\backslash X\to s\backslash Y\in W\quad(s\in S),\\
 u\in W_S^{\mathrm g}&\iff X/s\to Y/s\in W\quad(s\in S),\\
 u\in W_S^{\mathrm{univ}}&\iff
 X\times_ST\to Y\times_ST\in W\quad(T\to S).
\end{aligned}
\]
In the last line universality is relative to $S$. It is not the assertion that $u:X\to Y$ remains a $W$-equivalence after every map into its total target $Y$. The difference is already visible when $S=e$.

If $X$ and $Y$ are $W$-smooth over $S$, the inclusions of strict fibres into under-comma fibres give a commutative square
\[
\begin{tikzcd}
 X_s\arrow[r,"u_s"]\arrow[d]&Y_s\arrow[d]\\
 s\backslash X\arrow[r,"s\backslash u"']&s\backslash Y
\end{tikzcd}
\]
whose vertical arrows lie in $W$. Two-out-of-three identifies the fibre and right-fine tests. Smoothness survives change of base. Thus the fibre test, applied after every $T\to S$, gives a right-fine map over $T$ and hence a total $W$-equivalence by W(2). This proves $W_S^{\mathrm f}\subseteq W_S^{\mathrm{univ}}$ on smooth objects. The reverse inclusion follows by taking $T=\{s\}$. Finally, taking $T=S/s$ gives $W_S^{\mathrm{univ}}\subseteq W_S^{\mathrm g}$, and the relative localizer axiom gives $W_S^{\mathrm g}\subseteq W_S$. These are exactly the first chain of implications in the text. Opposites give the proper chain.

For weak $W$-fibrations, the defining universal comparisons apply to both the initial and the terminal comma categories. Hence both vertical comparisons, for $s\backslash X$ and for $X/s$, are in $W$, and remain so after changing the base. The same proof identifies all four fine or universal tests. W(4b) adds the reverse implication from the total test. Its statement is equivalent to equality of all five classes on $\Fib_W S$ for every $S$.

A two-object example explains the restriction to good objects. Let $S=(0\to1)$, let $X$ be the discrete category on $0,1$ mapped identically on objects to $S$, and let $Y=S$ with its identity structural map. The functor $u:X\to Y$ is the identity on objects. Every strict fibre map is the identity of a point, so $u\in W_S^{\mathrm f}$ for ordinary weak equivalences. But $X$ has two connected components and $Y$ has one, so $u\notin W_S$. At $0$, the under-comma category $0\backslash X$ consists of two isolated objects, whereas $0\backslash Y$ is an interval. The right-fine test detects exactly the missing arrow. This is why the text avoids using the strict-fibre localizer on all of $\Cat/S$.

The phrase that the inclusion is strict ``si on ne fait pas'' W(4b) should be read as a statement about what the basic axioms do not force. For a fixed $W$ and a fixed $S$, absence of an additional assumption is not a proof that the classes differ. The exact universal assertion is the equivalence with W(4b); the example below supplies a localizer and a base where the difference is real.
\end{dcomment}


Pour les résultats principaux développés jusqu'à présent sur $\mathrm{HOT}_W$, quand l'axiome W(4b) est utilisé, la forme faible W(4b$'$) suffit. On a
\[
\left|\;
\begin{array}{l}
W_S^{\mathrm{f}} \Longleftrightarrow W_S \ \text{ dans } \ul{Parf}_\omega S \ \textit{si}\ \mathrm{W}(4\mathrm{b}')\ \texttt{[est]}\ \text{satisfait}\\[1.5ex]
\text{et alors les cinq localiseurs } W_S^{\mathrm{f}},\, W_S^{\mathrm{d}},\, W_S^{\mathrm{g}},\, W_S^{\mathrm{univ}},\, W_S\\[1.5ex]
\text{sont égaux dans } \ul{Parf}_\omega S.
\end{array}
\right.
\]

\circled{2} \textbf{Principales catégories de catégories sur $S$.}

Ce sont, à des modifications d'éclairage près, celles du diagramme p. 76. Les catégories $\ul{Fibsc}\,S$, $\ul{Fib}\,S$, $\ul{Cofibsc}\,S$, $\ul{Cofib}\,S$ et leurs variantes à indice 0, sont inchangées - leur rôle est d'être aussi petites

\msp{page 175}

que possible pour remplir leur fonction (de description de $\mathrm{HOT}_W(S)$, $\mathrm{HOT}^{o}_W(S)$, $\mathrm{HOT}^{\mathrm{lc}}_W(S)$, ou $\mathrm{HOT}^{\mathrm{lc}!}_W(S)$ ci-dessous). La catégorie $\ul{Parf}\,S$ du diagramme remplissait une fonction du même type, c'est pourquoi nous la remplaçons à présent (en surenchérissant) par $\ul{Parf}_\omega S$. Les autres catégories du diagramme sont, en plus de $\mathrm{Cat}/S$, $\ul{Liss}\,S$ et $\ul{Prop}\,S$ et leurs variantes à indice 0, ainsi que $\ul{Fibhot}\,S$. (Toutes ces catégories, ainsi que $\ul{Parf}\,S$, étaient donc construites par référence à $W_\infty$, non pas à $W$.) Leur rôle est d'être aussi grandes que possible pour donner lieu à certaines propiétés et relations. Or il se trouve qu'ici la référence à $W_\infty$ était artificielle, et il y a lieu de mettre plutôt les catégories $\ul{Liss}_W S$, $\ul{Prop}_W S$, $\ul{Fib}_W S$, et en plus les variantes $\ul{Liss}_{0,W} S$, $\ul{Prop}_{0,W} S$. Ainsi le diagramme devient % [sic] propiétés -> propriétés, misprint
\[
\begin{tikzcd}[column sep=tiny, row sep=small]
 & & & & \ul{Fib}_W S \arrow[dd, hook, "\delta"] \arrow[dddllll, dotted, bend right=35, "\Phi_0"'] \arrow[dddrrrr, dotted, bend left=35, "\Psi_0"] & & & & \\
 & & \ul{Liss}_{0,W} S \arrow[urr, hook] \arrow[dd, hook] & & & & \ul{Prop}_{0,W} S \arrow[ull, hook'] \arrow[dd, hook] & & \\
 & \ul{Fib}_{0,W} S \arrow[ur, hook] \arrow[dd, hook] & & & \mathrm{Cat}/S \arrow[dd, "C"] \arrow[dddllll, dotted, bend right=30, "\Phi"] \arrow[dddrrrr, dotted, bend left=30, "\Psi"'] & & & \ul{Cofib}_{0,W} S \arrow[ul, hook'] \arrow[dd, hook] & \\
\ul{Fibsc}_{0,W} S \arrow[ur] \arrow[dd, hook] & & \ul{Liss}_W S \arrow[urr, hook] & & & & \ul{Prop}_W S \arrow[ull, hook'] & & \ul{Cofibsc}_{0,W} S \arrow[ul] \arrow[dd, hook] \\
 & \ul{Fib}\,S \arrow[ur, hook] & & & \ul{Parf}_\omega S \arrow[uuull, hook] \arrow[uuurr, hook'] & & & \ul{Cofib}\,S \arrow[ul, hook'] & \\
\ul{Fibsc}\,S \arrow[ur] & & & & & & & & \ul{Cofibsc}\,S \arrow[ul]
\end{tikzcd}
\]
(\footnote{NB la nouvelle notation serait plutôt $\ul{Fibf}_W S$, car il s'agit de $W$-fibrations qu'à présent j'appelle `faibles' (et définies seulement dans Cat, à l'opposé des $W$-fibrations, définissables dans toute catégorie de modèles).}) (tous les carrés \texttt{[sont]} cartésiens, les foncteurs notés $\hookrightarrow$ sont pleinement fidèles),

\printedpage{108}
% head: A. Grothendieck : Dérivateurs, Ch. VI, Édition des Universités de Montpellier II et Paris VII
% foot: Version 9 novembre 2001

\msp{page 176}

où j'ai encore marqué en vert \texttt{[dans l'original seulement]} les cinq \textit{`foncteurs d'inversion'}
\[
\Phi,\ \Psi,\ \Phi_0,\ \Psi_0,\ C,
\]
qui fournissent des quasi-inverses des foncteurs sur les catégories localisées déduits par les \textit{foncteurs d'inclusion.}

\circled{3} \textbf{Les quatre types de catégories} $\mathrm{HOT}_W$\textbf{.}

Sur chacune des catégories de l'etage, il y a \textit{deux} localiseurs à regarder, % [sic] l'etage -> l'étage, missing accent (printed « l'étage » below)
\[
\begin{tikzcd}[row sep=huge, column sep=0pt]
W_S^{\mathrm{f}} \arrow[d, Rightarrow] & {} = W_S^{\mathrm{g}} = W_S^{\mathrm{d}} = W_S^{\mathrm{univ}} \\
W_S, &
\end{tikzcd}
\]
donnant lieu à deux localisés distincts, notés
\[
\mathrm{HOT}^{\mathrm{lc}}_W(S) \qquad \text{et} \qquad \mathrm{HOT}^{\mathrm{lc}!}_W(S),
\]
il se trouve que (grâce à $\Phi_0$ et $\Psi_0$) ces localisés `sont les mêmes' pour les 8 catégories de l'étage, à l'exception de $\ul{Parf}_\omega S$; plus exactement, les foncteurs d'inclusion (à l'exception peut-être de ceux ayant pour source $\ul{Parf}_\omega S$) induisent des équivalences sur les deux types de localisés, les équivalences quasi-inverses étant obtenues par $\Phi_0$ (aile gauche) ou $\Psi_0$ (aile droite). Quand on localise \textit{pour} $W_S$ (à l'exclusion de $W_S^{\mathrm{f}}$), on trouve aussi des quasi-inverses par $C|\ul{Liss}_{0,W} S$ resp. $C|\ul{Prop}_{0,W} S$, pour les deux inclusions
\[
\begin{tikzcd}
\ul{Liss}_{0,W} S & \ul{Parf}_\omega S \arrow[l, hook'] \arrow[r, hook] & \ul{Prop}_{0,W} S
\end{tikzcd}
\]
(\footnote{cf. plus bas, ce qu'on peut dire pour les localisés pour $W_S^{\mathrm{f}}$.}). Par raison de symétrie, le plus joli semble de définir
\[
\left\{
\begin{array}{lcl}
\mathrm{HOT}^{\mathrm{lc}}_W(S) & = & (\ul{Fib}_W S)(W_S^{\mathrm{f}})^{-1}\\
\mathrm{HOT}^{\mathrm{lc}!}_W(S) & = & (\ul{Fib}_W S)(W_S)^{-1},
\end{array}
\right.
\]
et on a un foncteur de localisation
\[
\mathrm{HOT}^{\mathrm{lc}}_W(S) \xrightarrow{\ \lambda\ } \mathrm{HOT}^{\mathrm{lc}!}_W(S)
\]
induit par l'inclusion
\[
W_S^{\mathrm{f}} \subseteq W_S.
\]

\dgrblock{dgrb261}{Why four localizations appear}
\begin{dcomment}{exposition}{p176:induit par l'inclusion}{Why four localizations appear}
Grothendieck's large diagram organizes choices of representatives. The smaller categories make constructions explicit; the larger ones make stability properties easier to state. The object represented should not depend on a convenient replacement, but that independence must be proved for the particular class of arrows one inverts.

There are two first choices. Localizing all of $\Cat/S$ by the under-comma test gives $\HOT_W(S)$; using the over-comma test gives $\HOT_W^o(S)$. The replacements $\Phi$ and $\Psi$ make these two categories accessible through smooth and proper models respectively. Their symmetry comes from taking opposite categories, not from an identification of the two comma tests on arbitrary objects over $S$.

The weakly fibred objects make a second pair of choices possible. On $\Fib_W S$, all the fine tests agree, so one may form
\[
 \mathcal L_S=\Fib_W S[(W_S^{\mathrm f})^{-1}].
\]
This retains the homotopy type of each fibre together with its transport. One may instead invert every map of weakly fibred objects whose total functor belongs to $W$:
\[
 \mathcal G_S=\Fib_W S[W_S^{-1}].
\]
The notation in the text is $\mathcal L_S=\HOT_W^{\mathrm{lc}}(S)$ and $\mathcal G_S=\HOT_W^{\mathrm{lc}!}(S)$. The exclamation mark records that this second construction may identify more arrows. It does not add an extra adjunction or an additional kind of pointwise constancy.

Since $W_S^{\mathrm f}\subseteq W_S$ on weakly fibred objects, the identity functor induces $\lambda:\mathcal L_S\to\mathcal G_S$. This is a localization by its universal property. Namely, a functor out of $\mathcal L_S$ factors through $\lambda$ precisely when it inverts the images of all total $W$-equivalences between weakly fibred objects. No full faithfulness of $\lambda$ follows from this description.

The replacement $C$ explains why the coarse category also has the presentation
\[
 \mathcal G_S\simeq(\Cat/S)[W_S^{-1}].
\]
It sends every object into $\Parf_\omega S\subseteq\Fib_W S$, and its natural unit is a total $W$-equivalence. Therefore both composites of the inclusion and replacement become naturally isomorphic to the identity after coarse localization. The same proof applies to a full intermediate category containing the image of $C$. For fine localization, however, the proof requires a fibrewise unit; the coarse unit cannot be used in its place.

For $\Phi_0$ and $\Psi_0$ the necessary finer comparison is available on $\Fib_W S$. This is why the two wings of the upper diagram still describe $\mathcal L_S$, whereas the perfect vertex is temporarily excluded from that assertion. The pattern is not an arbitrary list of exceptions. Each exception records a precise natural transformation whose known class is too coarse for the desired localization.

The diamond on p.~178 consequently says $\rho i=\lambda=\sigma j$, up to the canonical identifications of presentations. Calling $i$ and $j$ inclusions expresses the intended role of locally constant objects inside the two directional theories. Their full faithfulness still requires an argument. The reflection formulas established after W(4b) is imposed provide such an argument; the shape of the diamond alone does not.
\end{dcomment}


% ---- from Der-ChVI-345_p111-115.tex
\printedpage{109}
% head: A. Grothendieck : Dérivateurs, Ch. VI, Édition des Universités de Montpellier II et Paris VII

\noindent\msp{page 177}
Pour toute sous-catégorie pleine
\[
\mathcal{F} \subseteq \Fib_W S
\]
qui contient l’une des trois sous-catégories
\[
\Fib_{0,W}S,\ \Parf_\omega S,\ \Prop_{0,W}S,
\]
l’inclusion de $\mathcal{F}$ dans $\Fib_W S$ définit des équivalences
\[
\mathcal{F}(W_S)^{-1} \approx \HOT^{\lcs}_W(S),
\]
et si $\mathcal{F}$ contient l’une des deux sous-catégories $\Fib_{0,W}S$, $\Prop_{0,W}S$ (à l’exclusion de $\Parf_\omega S$), on trouve aussi
\[
\mathcal{F}(W^{\mathrm{f}}_S)^{-1} \approx \HOT^{\mathrm{lc}}_W(S).
\]
Le rez-de-chaussée contribue deux autres variantes de HOT, en localisant soit par $W^{\mathrm{d}}_S$ (aile gauche), soit par $W^{\mathrm{g}}_S$ (aile droite). Les foncteurs d’inclusion définissent des équivalences pour les localisées soit pour $W^{\mathrm{d}}_S$ (aile gauche), soit pour $W^{\mathrm{g}}_S$ (aile droite). Le plus joli (symétrie) semble de définir ces deux localisés par $\Cat/S$
\[
\HOT_W(S) = (\Cat/S)(W^{\mathrm{d}}_S)^{-1}, \qquad \HOT^{o}_W(S) = (\Cat/S)(W^{\mathrm{g}}_S)^{-1}.
\]
Les quasi-inverses pour les équivalences de catégories localisées, induites par les foncteurs d’inclusion, sont donnés soit par $\Phi$ (aile gauche), soit par $\Psi$ (aile droite).

Mais $\Phi$ et $\Psi$ fournissent des quasi-inverses également quand on localise pour $W_S$, la $W$-équivalence grossière au lieu de la fine - pareil qu’à l’étage. On trouve donc au rez-de-chaussée un (unique) troisième localisé, qu’on peut décrire plus joliment par
\[
(\Cat/S)(W_S)^{-1},
\]
mais il est équivalent à $(\Fib_W S)W_S^{-1}$ par le
\msp{page 178}
foncteur induit par l’inclusion. (Un quasi-inverse étant fourni par $C$.) C’est donc $\HOT^{\lcs}_W(S)$. On trouve ainsi, par le diagramme commutatif de catégories avec localiseurs
\[
\begin{tikzcd}[column sep=small]
 & (\Fib_W S, W^{\mathrm{f}}_S) \arrow[dl] \arrow[dd] \arrow[dr] & \\
(\Cat/S, W^{\mathrm{d}}_S) \arrow[ddr] & & (\Cat/S, W^{\mathrm{g}}_S) \arrow[ddl] \\
 & (\Fib_W S, W_S) \arrow[d] & \\
 & (\Cat/S, W_S) &
\end{tikzcd}
\]
% foot: Version 9 novembre 2001
\printedpage{110}
% head: A. Grothendieck : Dérivateurs, Ch. VI, Édition des Universités de Montpellier II et Paris VII
(où la flèche verticale courte induit une équivalence sur les localisés), le diagramme commutatif entre les catégories $\HOT_W$ liées à $S$ :
\[
\begin{tikzcd}[column sep=small]
 & \HOT^{\mathrm{lc}}_W(S) \arrow[dl, "i"'] \arrow[dd, "\lambda"'] \arrow[dr, "j"] & \\
\HOT_W(S) \arrow[dr, "\rho"'] & & \HOT^{o}_W(S)(\simeq \HOT_W(S^{o})) \arrow[dl, "\sigma"] \\
 & \HOT^{\lcs}_W(S), &
\end{tikzcd}
\tag*{$(\ast)$}
\]
où $\rho$, $\sigma$, $\lambda$ sont des foncteurs de localisation, et où $i$, $j$ sont ressentis comme des foncteurs d’inclusion. Ce diagramme $(\ast)$ semble épuiser les foncteurs naturels qui relient entre elles les quatre catégories envisagées.

\noindent\msp{page 179}
Les propriétés des foncteurs $\Phi$, $\Psi$, $C$ qui ont été utilisées, pour fournir tous les quasi-inverses voulus de foncteurs (entre catégories localisées) induits par les foncteurs d’inclusion du diagramme, sont les suivantes.

\dgrblock{dgrb262}{An explicit failure of faithfulness for the coarse comparison}
\begin{dcomment}{development}{p178}{An explicit failure of faithfulness for the coarse comparison}
Grothendieck's four-category diagram raises a question about arrows, not only about which objects become isomorphic. The coarse localization can erase distinct automorphisms of a locally constant family. We make that loss explicit for the connected-components localizer. The covering obstruction discussed at manuscript p.~43 already distinguished a total equivalence from a fibre equivalence; the calculation here follows the comparison functor $\lambda$ itself and determines the localization of a whole class of examples.

Let $W_0$ be the class of functors inducing bijections on connected components. It satisfies the basic localizer axioms. Saturation and the terminal-object condition follow directly from the functor $\pi_0:\Cat\to\mathrm{Set}$. For the relative axiom, given $X\to S$, there is a natural bijection
\[
 \pi_0X\cong\mathop{\mathrm{colim}}_{s\in S}\pi_0(X/s).
\]
To see this without a homotopy-colimit theorem, send the class of an object $(x,p(x)\to s)$ to the component of $x$. Every component occurs by using $(x,\id_{p(x)})$. The relations induced by morphisms in a slice, and by arrows of $S$, preserve that component. Conversely, every object $(x,p(x)\to s)$ is identified in the colimit with $(x,\id_{p(x)})$, and every arrow $x\to x'$ gives the corresponding identification through the slice over $p(x')$. These observations give inverse maps. A pointwise bijection of the slice-component diagrams therefore induces a bijection on total components, proving W(2).

\textbf{Proposition (orbit collapse).}\ours\quad
Let $G$ be a nontrivial group and $S=BG$. For $W=W_0$, the functor
\[
 \lambda:\HOT_{W_0}^{\mathrm{lc}}(BG)
 \longrightarrow\HOT_{W_0}^{\mathrm{lc}!}(BG)
\]
is not faithful. On the full category of discrete fibrations over $BG$ associated with right $G$-sets, its fine and coarse source presentations are respectively the category of right $G$-sets and the category of sets; the induced functor is the orbit functor.

\emph{Proof.} For a right $G$-set $E$, let $\int E\to BG$ be its category of elements. Its objects are elements of $E$, and a morphism over $g\in G$ with target $e'$ has source $e'g$. Unique lifting, and the invertibility of $g$, make the projection both a discrete fibration and a discrete opfibration. The categories of prescribed lifts are singletons, so these objects are $W_\omega$-perfect and in particular weakly $W_0$-fibred. Functors over $BG$ between them correspond exactly to equivariant maps of $G$-sets.

The strict fibre of $\int E$ is the discrete set $E$, while its set of total connected components is $E/G$. It follows that a map induced by $h:E\to F$ is fine precisely when $h$ is a bijection, and coarse precisely when $E/G\to F/G$ is a bijection. A bijective equivariant map is already an isomorphism. Hence fine localization of this category changes nothing.

For the coarse localization, write $O(E)=E/G$, and give every set $A$ the trivial $G$-action $T(A)$. The natural equivariant map
\[
 q_E:E\longrightarrow T(O(E))
\]
induces the identity on orbit sets. Thus every $q_E$ is inverted by coarse localization, while $OT=\id_{\mathrm{Set}}$. The induced functors $O$ and $T$ are consequently quasi-inverse equivalences after that localization: $TO$ is naturally isomorphic to the identity through the inverted $q_E$. Under this equivalence the fine-to-coarse functor is exactly $O$.

It remains to verify that the loss of faithfulness persists in the full categories of the proposition, rather than only in the restricted example category. Take $E=G$ with its regular right action. For each $h\in G$, left multiplication $L_h(g)=hg$ is equivariant and defines an endomorphism of $\int G$ over $BG$. Evaluation on the strict fibre followed by $\pi_0$ is a functor
\[
 \Fib_{W_0}(BG)\longrightarrow\mathrm{Set}
\]
which inverts every fine equivalence. It therefore factors through $\HOT_{W_0}^{\mathrm{lc}}(BG)$. Its values on the $L_h$ are distinct set maps, so these endomorphisms remain distinct in the fine localization. On the other hand, $q_G:G\to\{*\}$ is coarse, because both orbit sets are singletons. The equalities $q_GL_h=q_G$ and invertibility of $q_G$ in the coarse localization imply that every $\lambda(L_h)$ is the identity. Thus $\lambda$ is not faithful.\hfill$\square$

This also gives a completely explicit failure of W(4b), even on the perfect objects used in its first weakened variant: $\int G\to BG$ is a total $W_0$-equivalence, whereas its fibre map $G\to\{*\}$ is not. The proposition does not assert that the full coarse category is always a category of sets. The equivalence with sets was proved for the displayed category of discrete fibrations; nonfaithfulness in the full category was established separately by fibre evaluation.

We checked the localizer discussion in \citep{Cisinski2004Minimal} and the locally constant formalism in \citep{Cisinski2009LocallyConstant}. We did not find this exact calculation of the manuscript's $\lambda$ there. The argument is an elementary development of familiar action-groupoid and localization constructions, with no claim of priority.
\end{dcomment}


\cercle{a} On a un homomorphisme fonctoriel dans $\Cat/S$
\[
X \xrightarrow{i_X} \Phi(X), \qquad \underline{i_X \in W^{\mathrm{d}}_S}.
\]
Quand $X$ est dans $\Fib_W S$, $\Phi(X)$ est dans $\Fibsc_0 S$ \texttt{[plutôt }$\Fibsc_{0,W}S$\texttt{, cf. diag. p. 175]}, donc $i_X \in W^{\mathrm{f}}_S$ ($= W^{\mathrm{d}}_S$ sur $\Fib_W S$). Ainsi $\Phi$ induit le foncteur $\Phi_0$ du diagramme.

Dualement, on a
\[
X \xrightarrow{i_X} \Psi(X), \qquad \underline{i_X \in W^{\mathrm{g}}_S},
\]
et quand $X$ est dans $\Fib_W S$, $\Psi(X)$ est dans $\Cofibsc_0 S$ \texttt{[plutôt }$\Cofibsc_{0,W}S$\texttt{]}, d’où le foncteur $\Psi_0$.

\cercle{b} On a un homomorphisme fonctoriel
\[
\left|
\begin{aligned}
& X \xrightarrow{i_X} C(X), \qquad \underline{i_X \in W_S} \\
& \text{et } C(X) \in \Ob\,\Parf_\omega S.
\end{aligned}
\right.
\]

\troisbarres Mais on fera attention qu’en l’absence de l’axiome W(4b) sur $W$, on ne peut \textit{pas} affirmer $i_X \in W^{\mathrm{f}}_S$, même si $X$ est dans $\Fib_W S$ (voire, dans $\Parf_\omega S$), seul cas (de toutes façons) où on est en droit de s’y attendre. Je note à ce propos le diagramme \textit{cartésien,} pour $f : X \longrightarrow Y$ dans $\Cat/S$
% foot: Version 9 novembre 2001
\printedpage{111}
% head: A. Grothendieck : Dérivateurs, Ch. VI, Édition des Universités de Montpellier II et Paris VII
\[
\begin{tikzcd}
C(Y)_s \arrow[d, "\pi^{s}_{Y}"'] & C(X)_s \arrow[l, "C(f)_s"'] \arrow[d, "\pi^{s}_{X}"] \\
Y & X, \arrow[l, "f"']
\end{tikzcd}
\qquad \text{\footnotesize flèches verticales dans $\Parf_\omega \subseteq \Fib_W$}
\]
qui montre par exemple que
\msp{page 180}
l’on ne peut pas s’attendre que $f \in W_S \implies C(f) \in W^{\mathrm{f}}_S$ sans disposer de l’axiome W(4a$'$) \texttt{[qu’est-ce que c’est?]} (ou d’un axiome de force comparable, tel W(4b)).

Ce sont ces faits négatifs (en l’absence de W(4b)) qui sont cause que je ne peux rien dire de particulier sur les foncteurs induits par \textit{localisation fine} $(W^{\mathrm{f}}_S)^{-1}$ des deux foncteurs d’inclusion
\[
\Liss_0 S \longleftarrow \Parf_\omega S \longrightarrow \Prop_0 S,
\]
ou, ce qui revient au même, de celui induit par le composé commun de ceux-là avec les inclusions dans $\Fib_W S$
\[
\begin{tikzcd}[column sep=tiny, row sep=small]
 & \Parf_\omega S \arrow[dl] \arrow[dd] \arrow[dr] & \\
\Liss_0 S \arrow[dr] & & \Prop_0 S \arrow[dl] \\
 & \Fib_W S, &
\end{tikzcd}
\]
en d’autres termes, je ne peux rien dire sur les propriétés de fidélité du foncteur
\[
(\Parf_\omega S)(W^{\mathrm{f}}_S)^{-1} \longrightarrow (\Fib_W S)(W^{\mathrm{f}}_S)^{-1},
\]
et encore moins que ce soit une équivalence de catégories. Donc, en l’absence de W(4b), pour les sous-catégories pleines $\mathcal{F}$ de $\Fib_W S$ qui contiennent $\Parf_\omega S$, sans contenir ni $\Liss_0 S$ ni $\Prop_0 S$, il se pourrait a priori que les localisées
\[
\mathcal{F}(W^{\mathrm{f}}_S)^{-1}
\]
soient toutes différentes. Je pense notamment à $\Parf_W S$, aussi intéressante en elle-même, après tout, que $\Parf_\omega S$. Donc en plus des quatre localisées du diagramme $(\ast)$, p. 178, il y a toute une série infinie ‘tératologique’ de localisées
\msp{page 181}
qui refusent de se laisser ‘boucler’ \texttt{[?]} par $\Phi_0$ ou $\Psi_0$, et pour lesquelles même $C$ n’est d’aucun secours, faute de l’axime W(4b). % [sic] axime -> axiome, missing letter

\dgrblock{dgrb263}{Which replacement actually gives a quasi-inverse}
\begin{dcomment}{exposition}{p180}{Which replacement actually gives a quasi-inverse}
Grothendieck now tests the tools used to draw the preceding diagram. The decisive information is carried by the natural maps into the replacements. A replacement whose value belongs to a desirable subcategory gives an equivalence of localizations only when its comparison map lies in the localizer used on both sides.

For $\Phi$ the comparison $X\to\Phi(X)$ is right-fine on all objects. If $X$ is weakly $W$-fibred, its value is a split locally constant fibration, and the equivalence of fine tests on weakly fibred objects makes the comparison fibrewise. The same reasoning for $\Psi$ uses the left-fine comparison and the proper side. These statements explain why the upper wings remain equivalent after fine localization. They are applications of the localization comparison lemma already discussed in the commentary at manuscript p.~12, rather than a new assertion that localization preserves arbitrary full embeddings. At the split vertices one also uses the free split-fibration adjunction discussed in the commentary at manuscript pp.~53--59; the original unit need not preserve the chosen splitting.

The map $X\to C(X)$ comes with a weaker recorded guarantee: it belongs to $W_S$. Its target is perfect, but perfection of the target alone does not make the comparison fine. The Cartesian square displayed on p.~179 isolates the issue even more sharply. For a total $W$-equivalence $u:X\to Y$, the map $C(u)_s$ is obtained by pulling $u$ back along the perfect map $C(Y)_s\to Y$. To conclude that $C(u)_s\in W$ one needs a suitable preservation theorem for that pullback. Two-out-of-three for the total units only proves $C(u)\in W_S$; it supplies no missing fibre comparison.

There are two different questions here. One asks whether $C$ turns coarse equivalences into fine ones. The other asks whether the inclusion of perfect objects induces an equivalence after fine localization. Failure of the first proposed proof does not establish a negative answer to the second question. Another replacement, or an additional limit argument, could still produce the required fine comparisons. This is exactly what the alternating construction on pp.~192--197 is designed to investigate.

Accordingly, the ``t\'eratologique'' family of localizations is a possibility recorded during the investigation, not a theorem that infinitely many pairwise inequivalent categories have been constructed. The source supplies no family of counterexamples on these pages. Modernly, the alternating replacement settles this particular question: the results explained at p.~172 make its values $W_\omega$-perfect and its units fine on every weakly $W$-fibred input. Every full intermediate category containing $\Parf_\omega S$ therefore has the same fine localization as $\Fib_W S$, by the common replacement and its units. The proof is given at p.~195. This does not identify fine and coarse localization: the connected-components counterexample still applies to $\lambda$.
\end{dcomment}


\subsection*{\cercle{4} Cas où l’axiome \textmd{W(4b)} est satisfait.}
% foot: Version 9 novembre 2001

\printedpage{112}
% head: A. Grothendieck : Dérivateurs, Ch. VI, Édition des Universités de Montpellier II et Paris VII
Alors les perplexités que je viens de dire disparaissent. On a aussitôt, en vertu de W(4b) :
\[
|||\qquad
\begin{array}{@{}r@{}}
\text{Si } X \in \Ob\,\Fib_W S, \text{ alors } i_X : X \longrightarrow C(X) \text{ est dans } W^{\mathrm{f}}_S \\
(\textit{en supposant l’axiome } \text{W(4b)})
\end{array}
\]
et $C$ fournit des quasi-inverses pour les foncteurs ‘d’inclusion’ entre localisés fins (pour $W^{\mathrm{f}}_S$) envisagés à la page précédente. Il en fournit \texttt{[‘Il en fournit’ barré]} un également pour le foncteur $\lambda$ envisagé plus haut ($(\ast)$ page 178)
\[
\underbrace{(\Fib_W S)(W^{\mathrm{f}}_S)^{-1}}_{\overset{\text{déf}}{=}\HOT^{\mathrm{lc}}_W(S)} \longrightarrow \underbrace{(\Cat/S)(W_S)^{-1}}_{\simeq \HOT^{\lcs}_W(S)},
\]
lequel est donc une équivalence. Une autre façon, plus évidente, de dire la même chose, c’est que sur chacune des catégories de l’étage les deux localiseurs $W^{\mathrm{f}}_S$ et $W_S$ sont désormais identiques (puisqu’elles sont toutes contenues dans $\Fib_W S$, où elles \texttt{[plutôt ils]} le sont justement par l’axiome W(4b)), donc le foncteur $\lambda$ entre ces deux localisés (définition originale de $\lambda$, dans le cas $\mathcal{F} = \Fib_W S$, cf. page 176) est même le \textit{foncteur identique,} les deux catégories localisées étant \textit{égales :}
\[
\HOT^{\mathrm{lc}}_W(S) = \HOT^{\lcs}_W(S) \qquad\qquad !
\]
Ainsi le diagramme $(\ast)$ p. 178 prend la forme
\[
\begin{tikzcd}[column sep=small]
 & \HOT^{\mathrm{lc}}_W(S) \arrow[dl, shift left=1.2, "i"] \arrow[dr, shift right=1.2, "j"'] & \\
\HOT_W(S) \arrow[ur, shift left=1.2, "\rho"] & & \HOT^{o}_W(S). \arrow[ul, shift right=1.2, "\sigma"']
\end{tikzcd}
\tag*{$(\ast)$}
\]

\dgrblock{dgrb264}{The fibre axiom closes the comparison}
\begin{dcomment}{exposition}{p181:Ainsi le diagramme $(\ast)$ p. 178 prend la forme}{The fibre axiom closes the comparison}
The force of W(4b) is now visible in one line. For a weakly $W$-fibred object $X$, both $X$ and $C(X)$ are in $\Fib_W S$, and the natural map $i_X:X\to C(X)$ is a total $W$-equivalence. W(4b) makes it a fibrewise $W$-equivalence. This is the missing comparison in the preceding block.

It is useful to check that this gives a functor on the desired localizations, not only good objects. If $u:X\to Y$ is fine, naturality gives
\[
 C(u)i_X=i_Yu.
\]
On each fibre, $i_X$ and $i_Y$ are $W$-equivalences. Two-out-of-three therefore makes $C(u)$ fine. The functor $C$ thus descends to fine localizations. On the full category of weakly fibred objects and on every indicated intermediate full subcategory containing its perfect image, the two composites with inclusion are naturally isomorphic to the identity through the same units. This proves the claimed equivalences.

For $\lambda$ an even shorter argument is available. The two presentations are localizations of the very same category $\Fib_W S$, and W(4b) identifies their two classes of denominators. Choosing the same localization construction on each side makes $\lambda$ literally the identity. Without that choice, the invariant conclusion is its canonical equivalence. Grothendieck's exclamation mark after the displayed equality expresses this formal simplification, not an additional theorem about particular models of localization.

The reduction from four categories to three consequently has an exact hypothesis. The connected-components example above shows that it cannot be asserted for every localizer satisfying only the basic axioms. The next step is to identify the remaining maps as adjoints and to prove that the locally constant category sits fully faithfully inside the two directional theories. Those assertions require the mapping-property calculation on pp.~182--183; they do not follow from drawing a triangle in place of the diamond.
\end{dcomment}


\noindent\msp{page 182}
Quand on interprète non seulement les deux catégories latérales $\HOT_W(S)$, $\HOT^{o}_W(S)$, mais aussi celle du milieu, comme des localisées de $\Cat/S$, pour les localiseurs $W^{\mathrm{d}}_S$, $W^{\mathrm{g}}_S$ et $W_S$ respectivement, les foncteurs $\rho$, $\sigma$ sont ceux déduits des deux inclusions
\[
\begin{tikzcd}[column sep=small, row sep=small]
W^{\mathrm{d}}_S \arrow[dr, hook] & & W^{\mathrm{g}}_S \arrow[dl, hook'] \\
 & W_S &
\end{tikzcd}
\qquad \text{\footnotesize comme parties de Fl(Cat/$S$),}
\]
tandis que les foncteurs $i$, $j$ sont déduits, par passage aux localisés, d’un seul et même foncteur
\[
\Cat/S \xrightarrow{C} \Cat/S,
\]
ce foncteur ayant la vertu d’appliquer $W_S$ dans $W^{\mathrm{d}}_S \cap W^{\mathrm{g}}_S$,
\[
f \in W_S \overset{\text{moyennant W(4b)}}{\Longrightarrow} C(f) \in W^{\mathrm{d}}_S \cap W^{\mathrm{g}}_S \qquad (\text{et même } \in W^{\mathrm{univ}}_S).
\]
% foot: Version 9 novembre 2001

\printedpage{113}
% head: A. Grothendieck : Dérivateurs, Ch. VI, Édition des Universités de Montpellier II et Paris VII
(Et pour cause, puisque $C$ se factorise par $\Fib_W S$ (et même par $\Parf_\omega S$), où $W_S$ est identique à $W^{\mathrm{d}}_S$, $W^{\mathrm{g}}_S$, $W^{\mathrm{univ}}_S$. Que $f \in W_S$ implique $C(f) \in W$ est indépendant de l’axiome W(4b), et se voit sur le diagramme
\[
\begin{tikzcd}
X \arrow[r, "f"] \arrow[d, "i_X"'] & Y \arrow[d, "i_Y"'] \\
C(X) \arrow[r, "C(f)"] & C(Y),
\end{tikzcd}
\]
où les flèches $i_X$, $i_Y$ sont dans $W_S$.)

Ceci dit, il est vrai que les couples
\[
(\rho, i),\ (\sigma, j)
\]
sont des couples de foncteurs adjoints, donc ($\rho$ et $\sigma$ étant des foncteurs de localisation) que $i$ et $j$ sont \textit{pleinement fidèles.} La formule d’adjonction peut se lire comme
\[
\boxed{\Hom_{W^{\mathrm{g}}_S}(X, CY) \simeq \Hom_{W^{\mathrm{d}}_S}(X, CY) \simeq \Hom_{W_S}(X, Y)}
\tag*{$(\ast)$}
\]
\msp{page 183}
pour $X$, $Y$ dans $\Cat/S$ sans plus. Un $\Hom_\Sigma$ indique le Hom dans $(\Cat/S)\Sigma^{-1}$.

\noindent\textbf{Corollaire.} \textit{Si $X$ dans} $\Cat/S$\textit{, $Y$ dans} $\Fib_W S$\textit{, alors}
\[
\begin{tikzcd}[column sep=small]
\Hom_{W^{\mathrm{g}}_S}(X, Y) \arrow[dr, "\sim"] & & \Hom_{W^{\mathrm{d}}_S}(X, Y) \arrow[dl, "\sim"'] \\
 & \Hom_{W_S}(X, Y) &
\end{tikzcd}
\]
\textit{sont bijectifs.}

En effet, $Y$ dans $\Fib_W(S)$ implique que $Y \xrightarrow{i_Y} C(Y)$ est dans $W^{\mathrm{univ}}_S \subseteq W^{\mathrm{d}}_S \cap W^{\mathrm{g}}_S$, donc
\[
\Hom_{W^{\mathrm{g}}_S}(X, Y) \xrightarrow{\sim} \Hom_{W^{\mathrm{g}}_S}(X, CY) \underbrace{\xrightarrow{\sim}}_{\text{adjonction}} \Hom_{W_S}(X, Y),
\]
et itou pour $\Hom_{W^{\mathrm{d}}_S}(X, Y)$.

\subsection*{\cercle{5} Perplexités sur les variantes affaiblies de l’axiome \textmd{W(4b)}.}

Que reste-t-il de toutes ces belles choses (plus que trois catégories HOT, au lieu de $4+\infty$!) quand on ne suppose plus que W(4b$'$) au lieu de W(4b), i.e. l’égalité de $W_S$ et $W^{\mathrm{f}}_S$ dans
% foot: Version 9 novembre 2001

% ---- from Der-ChVI-345_p116-120.tex
\printedpage{114}
% head: A. Grothendieck : Dérivateurs, Ch. VI, Édition des Universités de Montpellier II et Paris VII
% foot: Version 9 novembre 2001

$\Parf_\omega S$? Finalement, il me semble (contrairement à ce que j'avais cru voir) qu'il ne reste presque rien (\footnote{\textbf{NB} Si $W$ satisfait l'axiome des limites, alors W(4b$'$) implique W(4b), cf. théorème récapitulatif dans XVI 4.4.13 (p. 58-62), notamment $2^{o}$). La condition W(4b) est la condition B de cet énoncé, et W(4b$'$) la condition D. (Elle est impliquée par W(4a) $=$ C, ce que je soupçonnais nullement en écrivant XII \dots)}). On ne sait alors toujours pas, pour $X$ dans $\Fib_W S$, si
\[ i_X : X \longrightarrow C(X) \]
n'est dans $W^{\mathrm{f}}_S$, et non seulement dans $W_S$; on ne le sait que pour $X \in \Parf_\omega S$, mais dans ce cas on n'en a pas besoin, car l'utilité du foncteur $C(X)$ et de l'inclusion $i_X : X \longrightarrow C(X)$, c'est ‘d'améliorer’ $X$ et de le rendre parfait, voire même $W$-parfait. Mais s'il l'est déjà, c'est plus la peine. Il y aura, comme
\msp{page 184}
avant, 4 localisées essentiellement distinctes $\HOT_W$, $\HOT^{o}_W$, $\HOT^{\mathrm{lc}}_W$, $\HOT^{\mathrm{lc}\,!}_W$, plus une infinité d'autres $\mathcal{F}(W^{\mathrm{f}}_S)^{-1}$ pour $\mathcal{F}$ entre $\Parf_\omega S$ et $\Fib_W S$ (cf. page 180), avec le seul bénéfice qu'on peut maintenant exclure le cas $\mathcal{F} = \Parf_\omega S$, car dans ce cas $\mathcal{F}(W^{\mathrm{f}}_S)^{-1} = \mathcal{F}(W_S)^{-1} = \HOT^{\mathrm{lc}\,!}_W(S)$ fait partie des quatre canoniques. La chose décisive, c'est qu'on n'a toujours aucun moyen de prouver que
\[ \lambda : \HOT^{\mathrm{lc}}_W(S) \longrightarrow \HOT^{\mathrm{lc}\,!}_W(S) \]
est une équivalence, et je doute que ce soit le cas. Il faudrait regarder avec $W = W_{\mathbf{D}}$ la $\mathbf{D}$-équivalence, où $\mathbf{D}$ est un dérivateur Der($k$-Mod), $k$ un anneau donné (p.ex. $k = \mathbf{Z}$, ou $k$ un corps $\mathbf{Z}/p\mathbf{Z}$ ou $\mathbf{Q}$ \dots). Je ne vais pas le faire maintenant cependant. (Il est vrai que l'axiome W(4b$'$) n'est pas non plus satisfait dans ce cas \dots)

\dgrblock{dgrb265}{Local objects absorb the extra denominators}
\begin{dcomment}{exposition}{p183}{Local objects absorb the extra denominators}
Grothendieck has reached the point where the three remaining homotopy categories can be compared by their maps, rather than by successive descriptions of their objects. This is the purpose of the boxed formula. Its force is that the source $X$ is arbitrary: local constancy is imposed on the target, where it supplies the extension property needed to invert the additional arrows.

Write $\mathcal H_d(S)=(\Cat/S)[(W_S^{\mathrm d})^{-1}]$ and $\mathcal H_c(S)=(\Cat/S)[W_S^{-1}]$. Under W(4b), the functor $C$ represents a fully faithful right adjoint $i$ to the localization $\rho:\mathcal H_d(S)\to\mathcal H_c(S)$. For $X,Y$ in $\Cat/S$, the adjunction gives
\[
 \Hom_{\mathcal H_d(S)}(X,CY)
 \simeq \Hom_{\mathcal H_c(S)}(X,Y).
\]
This formula is natural in both variables. In particular, it identifies the map induced by localization, after the canonical identification of $CY$ with $Y$ in $\mathcal H_c(S)$; it is not merely an abstract bijection between two sets of the same size. The dual localization gives the other bijection in the text.

To see what changes when $Y$ is already $W$-fibred, use the comparison $i_Y:Y\to CY$. It belongs to $W_S^{\mathrm f}$, and on the category of $W$-fibred objects this is also $W_S^{\mathrm d}$ and $W_S^{\mathrm g}$. Thus postcomposition by $i_Y$ is a bijection on Hom sets in either lateral category. Composing with the displayed adjunction proves the corollary exactly as stated. Notice that no replacement of the source is required. An arbitrary object can map into a local object, and this map already incorporates all the relations that the coarse localization would introduce.

There is a useful intrinsic reformulation. Let $v:A\to B$ be a coarse weak equivalence, and regard it as a map in $\mathcal H_d(S)$. For every $W$-fibred $Y$, precomposition induces a bijection
\[
 v^*:\Hom_{\mathcal H_d(S)}(B,Y)
       \longrightarrow \Hom_{\mathcal H_d(S)}(A,Y).
\]
Indeed the two sides are the corresponding Hom sets in $\mathcal H_c(S)$, where $v$ is invertible. This is the categorical meaning of being local for the additional denominators. It is this property, rather than the mere existence of a map $Y\to CY$, that makes a reflection useful. The distinction was prepared in the commentary at manuscript pp.~64--69; here it becomes an operational formula for computing morphisms.

The hypothesis W(4b) enters at a precise place. Naturality of $X\to CX$ and two-out-of-three show, without W(4b), that $C$ preserves total $W$-equivalences. To conclude that $C(v)$ is also right-fine, one must pass from a total equivalence between fibred objects to equivalences on their fibres. That is exactly what W(4b) supplies. The following pages investigate whether a smaller family of test objects can supply the same passage.
\end{dcomment}


Si dans l'énoncé de l'axiome, on remplaçait $\Parf_\omega S$ par $\Parf_W S$, rien d'essentiel ne serait changé - tout au plus je remplacerais $\Parf_\omega S$ par $\Parf_W S$ dans le brillant diagramme p. 175. Finalement, je suis à présent convaincu qu'il faut tout ou rien : ou bien \textit{quatre} catégories $\HOT(S)$ pour $W$, en se contentant de W$(1,2,3)$, ou bien rien que trois, i.e. pas de $\HOT^{\mathrm{lc}\,!}_W(S)$ - et pour ça, il ne semble pas qu'on puisse se contenter de moins de W(4b).

\msp{page 185}
Il suffirait, il est vrai, de supposer que $W_S = W^{\mathrm{f}}_S$ sur l'une quelconque des sept catégories de l'étage autres que $\Parf_\omega S$. Cela impliquerait alors que dans les six autres, on aura $\overline{W_S} = \overline{W^{\mathrm{f}}_S}$ (mêmes saturés forts dans $\mathrm{Fl}(\mathcal{F})$). Un des deux les plus faibles des axiomes en question serait : % doubt resolved in DER_VI-c_corrections.csv; input note: line under "étage" -- read as the overline of $\overline{W^{\mathrm{f}}_S}$ on the next line, not an underline of "étage"

\dgrblock{dgrb266}{The later annotation changes the question}
\begin{dcomment}{status}{p184}{The later annotation changes the question}
Grothendieck is asking how much fibre detection is needed to identify the two locally constant theories. The note added to this passage belongs to a later stage of the investigation. In Chapter XVI, Theorem 4.4.13, especially part~$2^\circ$, he proves an equivalence of the relevant fibration conditions when the localizers concerned are stable under sequential colimits. Its conditions $B$, $C$, and $D$ are the later counterparts of the fibre-detection condition, the base-change condition, and the restriction to perfect objects discussed here. The proof uses the alternating replacement of manuscript pages~192--195, together with a factorization argument \citep[\S4.4, especially Corollaries~4.4.10--4.4.12 and Theorem~4.4.13]{Grothendieck1991XVI}.

Thus the apparent weakness of a perfect-object test must be read together with the limit hypothesis. The test cannot be used on an arbitrary weak fibration until a suitable replacement has been constructed and its unit has been proved fine. These are precisely the tasks undertaken later in the chapter. The annotation supplies a conditional resolution in its historical setting. The modern automatic-colimit result discussed after manuscript p.~195 supplies the limit hypothesis for the standing basic localizers, while the restricted fibre-detection test remains necessary. The references to ``XVI~4.9'' in notes~95--96 should be located through the statements: in the current electronic edition the relevant replacement results occur in \S4.4.
\end{dcomment}


\begin{quote}
Si $f : X \longrightarrow Y$ est un foncteur \textit{cartésien} de catégories Cat-fibrées et $W$-fibrées (\footnote{‘$W$-fibrées’ signifiant ici que pour $u \in \mathrm{Fl}\,S$, les foncteurs $u^*$ de changement de base $X_{s'} \longrightarrow X_s$, $Y_{s'} \longrightarrow Y_s$ sont dans $W$.}) sur une autre $S$, si $f \in W$, alors $f_s \in W$ pour tout $s \in S$,
\end{quote}

(et l'autre est la version duale, pour les foncteurs cocartésiens de catégories cofibrées). Sauf erreur, cet axiome doit déjà impliquer que $\overline{W}$ satisfait la version W(4a) la plus forte
\printedpage{115}
% head: A. Grothendieck : Dérivateurs, Ch. VI, Édition des Universités de Montpellier II et Paris VII
% foot: Version 9 novembre 2001
de toutes (a fortiori W(4b)) des axiomes W(4). Cet axiome me paraît donc à présent beaucoup plus raisonnable que W(4b$'$), et ceci d'autant plus qu'il se réfère directemment à la description de $\HOT^{\mathrm{lc}}_W(S)$ en termes de foncteurs (ou de ‘pseudofoncteurs’, si on ne veut supposer les catégories scindées et les foncteurs compatibles aux scindages) $S^{o} \longrightarrow \mathrm{Cat}$. % [sic] directemment -> directement, spelling
D'ailleurs, bien qu'il n'en ait pas l'air, cet axiome \textit{est} autodual, puisque la version duale, pour $S$ fixée, est la version directe pour $S^{o}$! Cela rend cet axiome encore plus appétissant. Donc pour formuler W(4a$'$), j'ai envie à présent de laisser tomber $\Parf_\omega$, et de faire pareil. Si on garde deux versions W(4a$'$) (l'ancienne) et W(4a$''$) (la nouvelle), l'une et l'autre ont la vertu d'impliquer W(4b). Et W(4b$'$) (nouvelle version) doit impliquer W(4a) (version ultraforte), sinon pour $W$, du moins pour $\overline{W}$.

Mes idées étaient tenacement [?] embrouillées
\msp{page 186}
jusqu'à aujourd'hui même, car j'avais l'idée que, grâce au mirifique foncteur $C$, la catégorie $\Parf_\omega S$, comme les sept autres catégories de l'étage, pourrait servir à définir $\HOT^{\mathrm{lc}}_W(S)$. Or il n'en est rien, semble-t-il, à moins de supposeer l'axiome W(4b), ou la variante W(4b$'$) (nouvelle version). L'ancienne version de W(4b$'$) s'avère parfaitement inutilisable. % [sic] supposeer -> supposer, spelling

\dgrblock{dgrb267}{A smaller test family must still contain replacements}
\begin{dcomment}{exposition}{p185}{A smaller test family must still contain replacements}
Grothendieck is trying to economize the fibre-detection axiom without losing the reflection. The difficulty is structural. Testing total equivalences only between perfect objects says nothing about the map which first makes a given fibred object perfect, unless that map is already known to be a fibrewise equivalence. This is why the first version of W(4b$'$) does not immediately do the work expected of it.

The issue can be stated without the changing notation. Let $\mathcal F$ be a category of objects over $S$, and let $E\subseteq V$ be two classes of arrows, thought of as fibrewise and total equivalences. Suppose $\mathcal P\subseteq\mathcal F$ is a full subcategory on which $E$ and $V$ agree. This equality alone does not identify $\mathcal F[E^{-1}]$ with $\mathcal F[V^{-1}]$. What is additionally needed is a functor $R:\mathcal F\to\mathcal P$ with a natural map $\eta:\mathrm{id}\to R$ whose components belong to $E$. Once such a functor exists, the conclusion follows, and in fact the argument proves more. For $u:X\to Y$ in $V$, naturality and two-out-of-three imply $R(u)\in V$. Since $RX,RY$ lie in $\mathcal P$, it follows that $R(u)\in E$. In the square
\[
\begin{tikzcd}
 X\arrow[r,"u"]\arrow[d,"\eta_X"'] &Y\arrow[d,"\eta_Y"]\\
 RX\arrow[r,"R(u)"']&RY
\end{tikzcd}
\]
the other three arrows therefore lie in $E$, so $u$ does also. Thus $E=V$ throughout $\mathcal F$, provided these are the actual two-out-of-three classes under consideration; if only invertibility in the localization was assumed, the conclusion is equality of their strong saturations.

In the present application, $\mathcal F=\Fib_W S$, $E=W_S^{\mathrm f}$, and $V=W_S$. The old perfect-object test does not initially provide the required fine unit. By contrast, a test on split fibrations with $W$-equivalent transport does: the functor $\Phi$ takes a $W$-fibred object into that class, and its unit is fibrewise in $W$. The infinite alternating construction below provides another such replacement, now into $\Parf_W S$, assuming the sequential-colimit axiom.

This explains both the frustration and the eventual simplification. There are two different statements to prove: enough good objects exist, and coarse equivalences between these objects are fine. The first is a replacement statement, the second a detection statement. Neither follows merely from writing down a localization. The proposition on manuscript p.~196 combines them; the omitted adjunction check on p.~197 can then be completed by transferring the already established reflector, as we shall do there.
\end{dcomment}


À ce propos, je n'ai toujours pas répondu clairement pourtant, pas un ‘oui’ ou un ‘non’, à la question :

\begin{quote}
\textbf{Question :} \textit{les inclusions}
\[
\begin{tikzcd}
 & \Liss_{0,W}S \\
\Parf_W S \arrow[ur] \arrow[dr] & \\
 & \Prop_{0,W}S
\end{tikzcd}
\]
\textit{induisent-elles, même en l'absence de W(4b) (ou de ses variantes affaiblies) des équivalences pour les localisées fines (par $W^{\mathrm{f}}_S$)? Même question pour} $\Parf_\omega S$ \textit{au lieu de} $\Parf_W S$ (\footnote{OK en présence de l'axiome W(6bis) sur les $\varinjlim$ , cf. p. 192 ff.}). \textit{Même question pour}
\[ \Bifib\, S \subseteq \Parf_\omega S, \]
\textit{la catégorie des catégories sur $S$ qui sont à la fois} Cat-\textit{fibrées et} Cat-\textit{cofibrées.}
\end{quote}

Je ne me rends pas compte, même en admettant (cette fois) l'axiome W(4a), si $\Bifib\, S$ est assez grosse pour pouvoir décrire tous le types d'homotopie relatifs localement constants sur $S$, i.e. si % [sic] tous le types -> tous les types, missing letter
\[ (\Bifib\, S)(W^{\mathrm{f}}_S)^{-1} \longrightarrow (\Fib_W S)(W^{\mathrm{f}}_S)^{-1} \]
est au moins essentiellement surjectif, sinon une équivalence de catégories. Il manque un foncteur
\[ \mathrm{Cat}/S \longrightarrow \Bifib\, S \qquad \text{ou du moins} \qquad \underbrace{\Fib_W S}_{\text{ou $\Fib_\omega S$\,?}} \longrightarrow \Bifib\, S, \]
\printedpage{116}
% head: A. Grothendieck : Dérivateurs, Ch. VI, Édition des Universités de Montpellier II et Paris VII
% foot: Version 9 novembre 2001
qui puisse jouer le rôle de $\Phi$ et de $\Psi$ \textit{simultané\msp{page 187}ment.} Je n'ai pas même vérifié que $\Phi$ ne transforme \textit{pas} objets de $\Prop_W S$ ou de $\Prop_{0,W}(S)$ en objets $W$-propres (donc $W$-parfaits, étant de toutes façons $W$-lisse et même $W_\omega$-lisse) sur $S$ (\footnote{Voir réponse à cette intéressante question page 192.}). Si c'était pourtant vrai sur $\Prop_W S$, ça montrerait que l'image essentielle de $\Prop_W S$ dans $\HOT_W(S)$ ($= (\mathrm{Cat}/S)(W^{\mathrm{d}}_S)^{-1}$) est contenu dans $\HOT^{\mathrm{lc}}_W(S)$. % [sic] contenu -> contenue, agreement with "l'image"
Prenant $S = \mathbf{\Delta}^1$, cela signifierait que si dans $X$ on a une sous-catégorie fermée $X_1$ telle que $X_1 \overset{i}{\hookrightarrow} X$ soit $W_\omega$-asphérique, alors $1\backslash X \longrightarrow 0\backslash X$ est dans $W$ - or $0\backslash X \simeq X$, $1\backslash X \simeq X_1$, la flèche est $X_1 \longrightarrow X$, qui est bien dans $W$ étant $W$-asphérique! C'est étrange et quasiment immoral (\footnote{On observe plus bas que c'est ainsi chaque fois que $S$ a un objet final.}). Mais à force de gribouillis je viens de me convaincre que même si $X$ est cofibrée sur $S$ (donc aussi $W$-propre), $\Phi(X)$ n'est pas $W$-propre (ce qui impliquerait $W$-parfait), ni
\msp{page 188}
même un $W$-fibré, ce qui signifierait que
\[ s'\backslash X \longrightarrow s\backslash X \qquad \text{associé à } s \longrightarrow s' \text{ dans } S \]
serait toujours dans $W$. Posant $Z' = s'\backslash S$, $Z = s\backslash S$, cela signifierait que si on a un \textit{co-isomorphisme local}
\[ i : \underbrace{Z'}_{\simeq z_0\backslash Z} \longrightarrow Z, \]
où $Z'$, $Z$ ont des éléments initiaux, et si $X$ est cofibrée sur $Z$, $X'$ son image inverse sur $Z'$, alors $X' \longrightarrow X$ est dans $W$. En fait, si $Z$ a aussi un élément \textit{final} $e$ (comme c'était le cas tantôt pour $\mathbf{\Delta}^1$), alors $Z' \simeq z_0\backslash Z$ en a aussi en $e'$, qui s'envoie sur $e$, et on aura le diagramme commutatif
\[
\begin{tikzcd}
X' \arrow[r, "j"] & X \\
X'_{e'} \arrow[u, "W"] \arrow[r, "\sim"] & X_e, \arrow[u, "W"']
\end{tikzcd}
\]
d'où $j \in W$. Donc pour faire un contrexemple,
\msp{page 189}
il faut prendre $Z$ sans élément final.

\dgrblock{dgrb268}{What the bifibration question asks}
\begin{dcomment}{status}{p186}{What the bifibration question asks}
Grothendieck now asks for a replacement which has both kinds of categorical lifting at once. The question is stronger than the existence of a $W$-perfect replacement. A Grothendieck bifibration has actual Cartesian and coCartesian lifts. Smoothness and properness ask for suitable categories of lifts to be $W$-aspherical. The latter condition can hold without a distinguished universal lift, so an equivalence using $\Parf_W S$ does not by itself answer the question about $\Bifib S$.

One case admits a complete answer with the constructions already at hand. Suppose the base $S$ is a groupoid. Every Grothendieck fibration $p:E\to S$ is then also an opfibration. To prove this, let $u:s\to t$ and $x\in E_s$. Choose a Cartesian lift of $u^{-1}$ with codomain $x$. Its domain lies over $t$, and the lift is an isomorphism because its base arrow is invertible: Cartesian lifting of the inverse gives a two-sided inverse by the uniqueness in the Cartesian universal property. The inverse of this lift is therefore a coCartesian arrow over $u$ with domain $x$. This supplies the required coCartesian lifts for all $u$ and $x$.

Consequently $\Phi(X)\to S$ is a bifibration for every $X\to S$. If $X$ is $W$-fibred, its unit $X\to\Phi(X)$ is in $W_S^{\mathrm f}$. The functor $\Phi$ preserves fine equivalences between such objects: in each fibre it gives a comma functor, and the two fibre-inclusion equivalences identify this with the original fibre map. The same unit, restricted to bifibrations, stays inside that full subcategory. Passing to the fine localizations therefore gives mutually inverse equivalences
\[
 (\Bifib S)[(W_S^{\mathrm f})^{-1}]
 \simeq (\Fib_W S)[(W_S^{\mathrm f})^{-1}]
 \qquad (S\text{ a groupoid}).
\]
No colimit axiom is used. This includes discrete bases, where the assertion is simply the independent comparison on every fibre.

For an arbitrary base, the proof breaks at exactly the attempted inversion of $u$. A fibration classified by $F:S^{\mathrm{op}}\to\Cat$ is a bifibration when its transport functors have the appropriate left adjoints, with the coherent structure coming from the liftings. Knowing only that these transport functors lie in $W$ does not produce such adjoints. The sequential construction below proves the fine-localization comparison for $\Parf_W S$; it does not equip its output with Cartesian and coCartesian lifts, and therefore does not establish the general $\Bifib S$ comparison.

A relevant categorical construction is now available: \citet[Theorem~1.17]{ClarkeSchererZeilberger2025} constructs the free cloven bifibration on any functor $p:X\to S$, with a natural unit $X\to\mathrm{Bif}(p)$. Its universal property uses functors preserving the chosen lifts. Propositions~1.23--1.24 make the unit fully faithful and give it a left inverse on total categories. Those conclusions do not make the unit a fibrewise $W$-equivalence on weakly $W$-fibred inputs, which is the missing assertion for the present application. The antecedent path construction in \citet[\S4, pp.~22--23]{Lamarche2013Path} explicitly distinguishes its own weak equivalences from ordinary nerve weak equivalences. Thus these free constructions cannot be cited as a verified homotopical replacement for this question, even at $W=W_\infty$.

We did not find the precise arbitrary-base, arbitrary-localizer bifibration comparison in the literature checked. This is a limit of the present audit, not an assertion that the problem is open. The modern theory of locally constant coefficients gives a homotopy invariant category of such coefficients; for example, \citet[Theorem~1.3 and Corollary~1.28]{Cisinski2009LocallyConstant} proves invariance under nerve weak equivalences and its detection by locally constant coefficients. That theorem does not assert the strict bifibration presentation asked for here.
\end{dcomment}


Mais prenons $Z$ \texttt{[un]} ensemble ordonné, avec un objet initial $0$, mais sans objet final, et prenons pour $z_0$ un élément \textit{maximal} de $Z$, de sorte que $Z' = z_0\backslash Z$ soit isomorphe à $\{z_0\}$. Il faut donner dans Cat un exemple où l'inclusion de $X_{z_0}$ dans $X$ n'est pas dans $W$. On va prendre
\printedpage{117}
% head: A. Grothendieck : Dérivateurs, Ch. VI, Édition des Universités de Montpellier II et Paris VII
% foot: Version 9 novembre 2001
\[
\begin{tikzcd}
 & & a \\
Z = & 0 \arrow[ur] \arrow[dr] & \\
 & & b = z_0,
\end{tikzcd}
\]
donc la donnée de $X$ cofibrée coscindée sur $Z$ équivaut à un diagramme dans Cat
\[
\begin{tikzcd}
 & X_a \\
X_0 \arrow[ur, "\alpha"] \arrow[dr, "\beta"] & \\
 & X_b.
\end{tikzcd}
\]
Je pense que sous des conditions assez fortes (p.ex. si $\alpha$, $\beta$ sont des immersions ouvertes, ou des immersions fermées - à vérifier!), la catégorie $X$ est isomorphe dans Hot à la somme amalgamée dans Cat
\[ X_a{\textstyle\coprod_{X_0}}X_b, \]
et il n'y \texttt{[a]} aucune raison (sauf si on suppose $\alpha \in W$, et encore \dots) que $X_b \longrightarrow X_a \coprod_{X_0} X_b$ soit dans $W$. Prenons $X_0 = X_b = e$, $\beta = \mathrm{id}_e$, donc $\alpha$ est donné par la donnée d'un élément $t$ de $\mathrm{Ob}(X_a = T)$. La somme amalgamée est $T$, qui n'est pas isomorphe à $e$ dans $\mathrm{Hot}_W$,
\msp{page 190}
si $T$ n'a pas été prise $W$-asphérique. Mais je vais essayer de montrer directement que l'inclusion
\[ T = X_a \longrightarrow X \]
est $\in W$. Pour ceci je regarde la sous-catégorie ouverte $Z''$ de $Z$, complémentaire de la sous-catégorie pleine fermée $Z'$, et je note que $X'' = X|Z''$ étant propre sur $Z''$ et $a$ étant objet final de $Z''$, on a
\[ X_a \longrightarrow X'' \qquad \text{dans } W, \]
il reste à voir qu'il en est de même de $X'' \longrightarrow X$. Or $X''$ est déduite de $X$ en ajoutant un unique élément $\xi$ (celui de $X_b$), lequel ne s'envoie dans aucun élément de $X''$, et qui ne reçoit aucun élément de $X_a$ - par contre pour l'unique élément $\xi_0$ de $X_0$ il y a une unique flèche de celui-ci dans $\xi$. On a donc une rétraction de $X$ sur $X''$, envoyant $\xi$ sur $\xi_0$ et la flèche précédente sur $\mathrm{id}_{\xi_0}$. Est-ce un homotopisme? Il suffit de \texttt{[montrer]} qu'en tant qu'endomorphisme $\varphi$ de $X$, c'est $\sim \mathrm{id}_X$. Or on a bien $\mathrm{id}_X \longleftarrow \varphi$, dont la valeur sur $X''$ est l'identité, et la valeur sur $\xi$ est l'unique $\varphi(\xi) = \xi_0 \longrightarrow \xi$. C'est un peu bref, mais ça doit être vrai. (De toutes façons, $X$ est la réunion de deux ouverts $X''$, et de l'ouvert $X|\{0,b\} = X'''$, avec $X''' \simeq \mathbf{\Delta}^1$, et $X'' \cap X''' \simeq e$, donc au moins pour la cohomologie à coefficients localement constants de $X$ je suis à l'aise par Mayer-Vietoris, et sais bien que c'est celle du morceau $X''$, donc celle de $X_0$ \texttt{[ou $X_a$, indice illisible]} nullement celle de $e$!)

\printedpage{118}
% head: A. Grothendieck : Dérivateurs, Ch. VI, Édition des Universités de Montpellier II et Paris VII
% foot: Version 9 novembre 2001
\msp{page 191}
\margbar On a donc vu au moins ceci : Si $X$ est cofibrée donc $W$-propre sur $S$, il ne s'ensuit nullement que $\Phi(X)$ soit encore $W$-propre sur $S$ (\footnote{voir résultat plus précis page suivante!}), ni que ce soit un $W$-fibré. Mais que se passe-t-il si $X \in \Prop_{0,W}S$, par exemple s'il provient de $\Cofibsc_0 S$? Alors $X$ est dans $\Fib_W S$, donc on sait déjà que $\Phi(X)$ est dans $\Fibsc_{0,W}S$, donc ce n'est pas par là qu'on fera un contrexemple. Mais considérons encore $X$ cofibrée sur
\[
\begin{tikzcd}
 & & a \\
S = & 0 \arrow[ur] \arrow[dr] & \\
 & & b,
\end{tikzcd}
\]
donnée par le diagramme
\[
\begin{tikzcd}
 & X_a \\
X_0 \arrow[ur, "u"] \arrow[dr, "v"] & \\
 & X_b
\end{tikzcd}
\]
dans Cat, en supposant cette fois $u, v \in W$ ($X \in \Cofibsc_{0,W}S$). Alors $\Phi(X)$ est la catégorie fibrée scindée sur $S$ définie par le diagramme contravariant sur $S$
\[
\begin{tikzcd}
 & X_a \arrow[dl, hook, "i"'] \\
X_0 & \\
 & X_b, \arrow[ul, hook, "j"']
\end{tikzcd}
\]
où $i$, $j$ sont les inclusions (qui sont des immersions fermées). Dire que $\mathcal{X} = \Phi(X)$ est propre sur $S$ signifie aussi que ses restrictions au dessus des sous-catégories ouvertes $\{0,a\}$ et $\{0,b\}$ de $S$ le sont (car toute flèche $\mathbf{\Delta}^1 \longrightarrow S$ se factorise par l'une ou l'autre). Cela
\msp{page 192}
signifie donc aussi que l'inclusion des fibres fermées
\[ X_a \longrightarrow X, \qquad X_b \longrightarrow X \]
sont coasphériques. Mais $X_a \longrightarrow X$ ne peut être coasphérique que si $X_b = \emptyset$, car pour $x \in X_b \subseteq X$, $x\backslash X_a$ (ainsi que $X_a/x$) est vide, donc \textit{pas} $W$-asphérique (si on ne suppose $W = \mathrm{Fl}(\mathrm{Cat})$, cas que nous excluons). Donc $\Phi(X)$ n'est propre sur $S$ \textit{que} si $X_a$ et $X_b = \emptyset$ (résultat valable pour \textit{toute} $X$ sur $S$, qu'elle soit cofibrée sur $S$ ou $W$-propre, ou non). Si on veut $u, v \in W$, cela implique donc que $X_0 = \emptyset$, donc on a prouvé :

\dgrblock{dgrb269}{A branching base defeats one-sided completion}
\begin{dcomment}{examples}{p190}{A branching base defeats one-sided completion}
Grothendieck needs a counterexample which separates transport in one direction from transport in the other. A base with a terminal object hides this distinction: an opfibration over such a base has the homotopy type of its terminal fibre. The branching poset $S=(0\to a,\,0\to b)$ removes that shortcut.

Take a pointed category $(T,t)$, and form the split opfibration classified by
\[
 e\xrightarrow{t}T,\qquad e\xrightarrow{\mathrm{id}}e.
\]
Its total category $X$ has the objects of $T$, together with objects $\xi_0$ over $0$ and $\xi_b$ over $b$. The additional arrows are a single arrow $\xi_0\to\xi_b$ and, for every $z\in T$, arrows $\xi_0\to z$ indexed by $\Hom_T(t,z)$. Composition with arrows of $T$ is the evident one. There are no arrows from $T$ to $\xi_0$ or $\xi_b$, and none from $\xi_b$ to the other objects.

Let $X''$ be the full subcategory obtained by deleting $\xi_b$. There is a retraction $r:X\to X''$ which sends $\xi_b$ to $\xi_0$ and the arrow $\xi_0\to\xi_b$ to an identity. If $j:X''\hookrightarrow X$ denotes the inclusion, then $rj=\mathrm{id}$ and there is a natural transformation $jr\Rightarrow\mathrm{id}_X$, equal to the identity on $X''$ and to $\xi_0\to\xi_b$ at $\xi_b$. Naturality needs checking only on the new arrow, where both composites are that arrow. Thus $j$ is a categorical homotopy equivalence and belongs to every localizer in use here.

There is a second retraction $r_a:X''\to T$, sending $\xi_0$ to $t$ and fixing $T$. In fact $r_a$ is left adjoint to the inclusion $j_a:T\hookrightarrow X''$: the only nontrivial adjunction formula is
\[
 \Hom_T(r_a\xi_0,z)=\Hom_T(t,z)
              =\Hom_{X''}(\xi_0,j_a z).
\]
Hence $T\to X''\to X$ belongs to $W$. The fibre over $b$ is nevertheless $e$. If $T\to e$ is not in $W$, the map $X_b\to X$ is not in $W$, by two-out-of-three. Since $\Phi(X)_0\simeq X$ and $\Phi(X)_b\simeq e$, transport along $0\to b$ in the completed fibration is not a $W$-equivalence. This proves that $\Phi(X)$ need not even be $W$-fibred, although $X\to S$ is an opfibration and hence proper.

For the ordinary weak-homotopy localizer, a completely finite choice is $T=\{t,t'\}$ discrete. Then $X$ is the disjoint union of the three-object poset $\xi_b\leftarrow\xi_0\to t$ and the isolated object $t'$. Its nerve has two contractible components, whereas $X_b$ has one. No homology calculation or pushout theorem is required. For a general $W$, the argument requires a pointed $T$ which is not $W$-aspherical; the nondegeneracy assumption $W\ne\mathrm{Fl}(\Cat)$ alone should not be substituted for that requirement in this particular example.

The parenthetical pushout question has a positive answer under the suggested symmetric hypotheses. If $A\hookrightarrow B$ and $A\hookrightarrow C$ are both full sieves, their pushout $P$ contains no arrows from $B\setminus A$ to $C\setminus A$, or conversely. Every composable chain lies in one side, so $NP=NB\cup_{NA}NC$. The nerve inclusions are monomorphisms, and this union is the homotopy pushout of the two inclusions. Equivalently one can check the canonical functor from the total category of $B\leftarrow A\to C$ by comma categories: for $p\in B\setminus A$, retract the copy of $A/p$ lying on the $C$-side onto the middle copy. The resulting category has terminal object $p$ on the $B$-side. The case $p\in C\setminus A$ is symmetric, and for $p\in A$ the comma category contracts to the three copies of $p$ joined by a span. Theorem A therefore proves the comparison is in $W$. For two cosieves, apply the dual argument. The special example above is simpler still, since the explicit retractions determine the homotopy type for every allowed $W$.
\end{dcomment}


% ---- from Der-ChVI-345_p121-125.tex
\printedpage{119}
% head: A. Grothendieck : Dérivateurs, Ch. VI, Édition des Universités de Montpellier II et Paris VII
% foot: Version 9 novembre 2001

% the following paragraph is marked by a triple vertical rule in the left margin
\noindent\llap{$|||$\quad}Soit $X \in \mathrm{Cat}/S$, $S = \begin{tikzcd}[row sep=tiny, column sep=small] & a \\ 0 \arrow[ur] \arrow[dr] & \\ & b \end{tikzcd}$. Si $\Phi(X)$ est propre sur $S$, alors $X = \emptyset$ (si on suppose $W \neq \mathrm{Fl}(\mathrm{Cat})$).

Mais chemin faisant est apparue une idée bizarre. Considérons le diagramme infini, pour $X$ dans $\mathrm{Cat}/S$ :
\[
X \longrightarrow \Phi(X) \longrightarrow \Psi\Phi(X) \longrightarrow \Phi\Psi\Phi(X) \longrightarrow \cdots
\]
Si $X$ est quelconque, ces flèches sont alternativement dans $W_S^{\mathrm{d}}$ et $W_S^{\mathrm{g}}$, donc leurs composées sont dans $W_S$ sans plus. Mais si $X$ est dans $\und{Fib}_W S$, comme cette catégorie est \textit{stable} par $\Phi$, $\Psi$, et que sur $\und{Fib}_W S$ $W_S^{\mathrm{d}} = W_S^{\mathrm{g}} = W_S^{\mathrm{f}} = W_S^{\mathrm{univ}}$, on voit que dans ce cas les flèches sont toutes dans $W_S^{\mathrm{f}}$. On a envie maintenant de poser
\[
\Phi_\infty(X) = \varinjlim \text{ de tout ça}
\]
\msp{page 193}
et on définirait symétriquement
\[
\Psi_\infty(X) = \varinjlim (X \longrightarrow \Psi(X) \longrightarrow \Phi\Psi(X) \longrightarrow \cdots).
\]

\dgrblock{dgrb270}{The obstruction persists on any base with two maximal branches}
\begin{dcomment}{development}{p192}{The obstruction persists on any base with two maximal branches}
Grothendieck strengthens the preceding counterexample: over the three-object branching poset, a single application of $\Phi$ cannot be proper unless the original category is empty. The obstruction lies in the absence of arrows between the two maximal branches. We can isolate that reason and extend the statement without introducing additional homotopical machinery.

\textbf{Branching obstruction.}\ours\enspace
Let $S$ be a poset with a least element $0$ and two distinct maximal elements $a,b$. Let $W$ be a localizer satisfying the standing axioms and suppose that $\varnothing\to e$ does not belong to $W$. For every $X\to S$, if $\Phi_S(X)\to S$ is $W$-proper, then $X=\varnothing$.

\emph{Proof.} The fibre diagram of $\Phi_S(X)$ is $s\mapsto s\backslash X$. Since $0$ is least, $0\backslash X$ identifies with $X$. Since $a$ and $b$ are maximal, $a\backslash X=X_a$ and $b\backslash X=X_b$. The two transport functors to the fibre at $0$ are therefore the inclusions $X_a\hookrightarrow X$ and $X_b\hookrightarrow X$.

Properness forces each of these inclusions to be coaspherical. Here is the precise test being used. Pull the fibration back along the arrow $0\to a$. For an object $x$ of its fibre over $0$, the category of arrows from $x$ to the fibre over $a$ is the comma category $x\backslash X_a$. Properness requires it to be $W$-aspherical. This is the elementary lifting-category criterion for properness, applied to this one arrow; the same statement holds for $b$.

If $x\in X$ existed, the category $x\backslash X_a$ would be $W$-aspherical and hence nonempty. Choose one of its objects, an arrow $x\to x_a$ with $x_a\in X_a$. Properness along $0\to b$, applied now to the object $x_a$ of $X$, would require $x_a\backslash X_b$ to be $W$-aspherical. But this category is empty: an arrow $x_a\to x_b$ would project to an arrow $a\to b$, impossible for distinct maximal elements of a poset. This contradicts the hypothesis on the empty category. Thus $X$ has no objects. Conversely, the empty category over $S$ is proper, since all the lifting tests with a specified source object are vacuous.\hfill$\square$

For the standing localizers, the nondegeneracy condition in the text implies the displayed hypothesis on $\varnothing\to e$. Indeed, if $\varnothing\to e$ were in $W$, product stability would put $\varnothing\to A$ in $W$ for every category $A$. Two-out-of-three would then put every functor in $W$.

This proof also fills the final step in the boxed assertion on manuscript p.~192. Showing that $X_a=X_b=\varnothing$ is not yet, by itself, a proof that $X_0=\varnothing$ for an arbitrary $X\to S$. One must use coasphericity once more: $\varnothing\to X$ cannot be coaspherical when $X$ has an object. The argument above incorporates this step directly. The preceding contravariant diagram prints $X_0$ at its centre; its mathematical value is $X=0\backslash X$, as both the following inclusions and the proof require. That printed slip is retained in the source and recorded in the correction ledger.

We checked the lifting criteria in \citet[Proposition~4.3, Theorem~4.13, and Example~4.17]{Maltsiniotis2005Asphericite}; we did not find this particular branching formulation there. The proposition is an elementary extension of the present manuscript argument, with no claim of priority.
\end{dcomment}


Pour que ce soit utilisable, il faut poser l'axiome \underline{W(6 bis)} (comparer XII, pages 14, 15). La limite inductive d'une suite de morphismes $f_i : X_i \longrightarrow Y_i$, pour un morphisme de systèmes inductifs dans Cat indexés par $\mathbf{N}$, est dans $W$ si les $f_i$ le sont.

Cet axiome est voisin de W(6) de loc. cit., mais n'est impliqué par lui que moyennant W(4b). Il est plus fort dans le cas de systèmes inductifs indexés par $\mathbf{N}$ (donc aussi, plus généralement, de systèmes inductifs sur des ensembles ordonnés filtrants ayant un ensemble cofinal dénombrable), mais ne dit rien sur les systèmes inductifs sur des ensembles ordonnés d'indices plus généraux.

Cet axiome implique aussi qu'une limite inductive filtrante de morphismes $f_i$ est $W$-lisse, resp. $W$-propre, si les $f_i$ le sont. On trouve alors que $\Phi_\infty(X)$ est à la fois $W$-lisse (comme limite inductive des $\Phi(X)$, $\Phi\Psi\Phi(X)$, \ldots, $\Phi(\Psi\Phi)^n(X)$, lesquelles sont $W$-lisses sur $S$) et $W$-propre (comme limite inductive des $\Psi\Phi(X)$, \ldots, $(\Psi\Phi)^n(X)$ \ldots). D'autre part, que $X \longrightarrow \Phi_\infty(X)$ est dans $W_S^{\mathrm{f}}$ - car $W_S^{\mathrm{f}}$ est stable par limites inductives filtrantes dénombrables dans $\mathrm{Cat}/S$. On a donc trouvé un foncteur
\msp{page 194}
\[
\Phi_\infty : \und{Fib}_W S \longrightarrow \und{Parf}_W S \subseteq \und{Fib}_W S
\]
et un morphisme fonctoriel
\[
i_X : X \longrightarrow \Phi_\infty(X), \qquad \boxed{i_X \in W_S^{\mathrm{f}}} \qquad \text{si } X \in \und{Fib}_W S.
\]
% the box around i_X \in W_S^f is printed with rounded corners

Ce foncteur permet alors, en l'absence de W(4b), de faire ce que le foncteur $C$ refusait absolument. Il est vrai que $C$ a l'avantage d'être défini sur $\mathrm{Cat}/S$ tout entier, et d'y être
\printedpage{120}
% head: A. Grothendieck : Dérivateurs, Ch. VI, Édition des Universités de Montpellier II et Paris VII
% foot: Version 9 novembre 2001
raisonnable et même fort utile, quand on a W(4b). Mais $\Phi_\infty$ est, lui aussi, défini sur $\mathrm{Cat}/S$ tout entier, ainsi que $i_X$
\[
\begin{aligned}
&\Phi_\infty : \mathrm{Cat}/S \longrightarrow \und{Parf}_W S\\
&i : \mathrm{id} \longrightarrow \Phi_\infty \qquad \text{i.e.}\ \begin{cases} i_X : X \longrightarrow \Phi_\infty(X)\\ i_X \in W_S, \end{cases}
\end{aligned}
\]
avec le seul handicap (inévitable) que si $X$ n'est pas dans $\und{Fib}_W S$, on ne peut dire mieux que $i_X \in W_S$, à l'exclusion de $i_X \in W_S^{\mathrm{f}}$. Mais pour $C(X)$ c'était forcément pareil. Donc on a l'impression que $\Phi_\infty$, au choix, $\Psi_\infty$, bonnet blanc, blanc bonnet, joue en présence de W(6 bis) le même rôle que C en présence de W(4b). Tout au plus (en présence des deux axiomes à la fois) le foncteur $C$ a-t-il l'avantage esthétique de prendre ses valeurs dans $\und{Parf}_\omega S$, au lieu de $\und{Parf}_W S$ seulement (\footnote[93]{\textit{En fait} $\Phi_\infty = C$, si $C$ est défini via des chemins infinis d'un seul côté simplement. Si on définit $C$ via les chemins \textit{finis} de tous types $c$, alors il prend en effet ses valeurs dans $\und{Parf}_\omega S$.}). Mais il s'impose de % [sic] "que C en présence" (C in roman type) -> $C$, italic everywhere else
\msp{page 195}
se rendre compte à présent que $\Phi_\infty$ a un avantage exactement du même ordre sur $C$, en introduisant
\[
W''_\omega = \text{plus petite partie } \Sigma \text{ de } \mathrm{Fl}(\mathrm{Cat}) \text{ satisfaisant } \mathrm{W}(1, 2\text{ bis}, 3, 6\text{ bis})
\]
et noter que
\[
\Phi_\infty(X) \in \mathrm{Ob}\,\und{Parf}_{W''_\omega} S, \qquad i_x \in W''_\omega % [sic] i_x -> i_X, the morphism is i_X everywhere else
\]
- mais là je m'aperçois qu'on a
\[
W_\omega \subseteq W''_\omega,
\]
donc $C$ est quand même un peu meilleur à cet égard \ldots

\dgrblock{dgrb271}{Why alternating completion works at the limit}
\begin{dcomment}{exposition}{p194}{Why alternating completion works at the limit}
The failure of one-sided completion suggests Grothendieck's next move. Instead of asking $\Phi$ to preserve properness, he restores properness with $\Psi$, then restores smoothness with $\Phi$, and continues. No finite stage has to possess both properties. At the limit, the odd and even stages are cofinal and can establish the two properties separately.

Put $X_0=X$, $X_1=\Phi X$, $X_2=\Psi\Phi X$, and continue alternately, with the canonical unit at each step. Let $R X=\varinjlim_n X_n$. W(6 bis) means that an objectwise $W$-equivalence between two sequential diagrams induces a $W$-equivalence on their ordinary colimits in $\Cat$. Its quantifier is sequential. The unqualified words ``limite inductive filtrante'' in the next paragraph must be read with this restriction; an arbitrary filtered indexing category is not covered by the axiom as written.

The passage from this axiom to smoothness and properness of $RX\to S$ can be made explicit by the lifting-category criterion. For a functor $p:E\to S$, a base arrow $u:s\to t$, and $y\in E_t$, let $L_p(u,y)$ have objects arrows $v:x\to y$ with $p(v)=u$ and morphisms commuting triangles whose arrow between sources lies over $\mathrm{id}_s$. Smoothness is equivalent to $L_p(u,y)$ being $W$-aspherical for every $u,y$. Properness is the dual condition on arrows with a specified source. These are the familiar fibre categories of $E/y\to S/t$ and their duals.

Suppose first that $E=\varinjlim_n E_n$ over the fixed base $S$, with every $E_n\to S$ smooth. An object $y$ of $E$ is represented by $y_n\in E_n$ at some stage. After restricting to that tail of the sequence, there is a canonical isomorphism
\[
 L_p(u,y)\simeq\varinjlim_{m\ge n}L_{p_m}(u,y_m).
\]
To verify it, represent an arrow to $y$ at a sufficiently late stage; its image in $S$ is already the required fixed arrow $u$. A morphism of the lifting category is represented by one further arrow and one equality of composites, and that equality holds at a later stage. The same argument detects equality of two such morphisms. Thus the displayed functor is bijective on objects and arrows after taking the colimit. Every stage maps to $e$ by a member of $W$, so W(6 bis) makes the colimit $W$-aspherical. This proves sequential stability of smoothness. Reversing arrows proves stability of properness.

Now the odd subsequence $X_1,X_3,\ldots$ consists of smooth categories over $S$, and the even subsequence $X_2,X_4,\ldots$ consists of proper ones. Both have colimit $RX$, because each is cofinal in the original sequence. Hence $RX$ is $W$-perfect. Every transition $X_n\to X_{n+1}$ belongs to $W_S$, so every composite $X\to X_n$ does. Apply W(6 bis) to the natural transformation from the constant sequence at $X$ to $(X_n)$: its colimit is $X\to RX$, proving that this map lies in $W_S$ for arbitrary $X$.

When $X$ is $W$-fibred, all stages remain $W$-fibred and the transitions are fibrewise $W$-equivalences. Fibres over a fixed $s$ commute with these colimits: their objects and arrows are defined by fixed equations in $S$, and any equality in a filtered colimit is witnessed at a later stage. Applying W(6 bis) in each fibre yields
\[
 X_s\longrightarrow (RX)_s\quad\text{in }W
 \qquad\text{for every }s\in S.
\]
Thus the same construction supplies both kinds of unit advertised in the text: a total equivalence for arbitrary objects and a fine equivalence for fibred objects. It is the latter statement that was missing for the earlier proposed perfect-object replacement. Nothing in this argument asserts that $RX\to S$ is a Grothendieck bifibration: a sequential colimit of lifting categories known to be aspherical does not supply actual universal lifts.
\end{dcomment}


Quand on dispose de $\Phi_\infty$ avec les propriétés qu'on vient d'expliciter, on trouve que l'inclusion
\[
\und{Parf}_W S \longrightarrow \und{Fib}_W(S)
\]
induit une équivalence non seulement pour les localisés pour $W_S$, mais aussi pour les localisés fins
\[
(\und{Parf}_W S)(W_S^{\mathrm{f}})^{-1} \longrightarrow (\und{Fib}_W(S))(W_S^{\mathrm{f}})^{-1} \overset{\text{déf}}{=} \mathrm{HOT}^{\mathrm{lc}}_W(S)
\]
(\footnote[94]{On peut y remplacer même $\und{Parf}_W S$ par $\und{Parf}_{W''_\omega} S$, \textit{ou} par toute autre sous-catégorie pleine de $\und{Fib}_W S$ contenant $\und{Parf}_{W''_\omega} S$.}). Mais on ne peut pour autant réduire le nombre des localisés utiles à trois, en l'absence de W(4b). Tout au plus, ayant un foncteur $\Phi_\infty$ qui est plus simple que $C$, en l'absence de W(4b), du fait qu'on sait de toutes façons que
\[
i_X : X \longrightarrow \Phi_\infty(X)
\]
est dans $W_S^{\mathrm{f}}$ si $X$ est dans $\und{Fib}_W S$, peut on
\printedpage{121}
% head: A. Grothendieck : Dérivateurs, Ch. VI, Édition des Universités de Montpellier II et Paris VII
% foot: Version 9 novembre 2001
\msp{page 196}
espérer s'en tirer avec la variante affaiblie de W(4b) obtenue en se bornant dans l'énoncé de W(4b) (p. 171) au cas où $X$, $Y$ sont dans $\und{Parf}_{W''_\omega} S$. Cette fois, ça suffit en effet à impliquer que le foncteur
\[
\lambda : (\und{Fib}_W S)(W_S^{\mathrm{f}})^{-1} \longrightarrow (\und{Fib}_W S)(W_S)^{-1}
\]
est une équivalence de catégories, puisque on peut le ‘calculer’ aussi bien en utilisant $\und{Parf}_{W''_\omega} S$ au lieu de $\und{Fib}_W S$, or l'hypothèse W(4b($\mathrm{W}''_\omega$)) implique que sur $\und{Parf}_{W''_\omega} S$ les deux localiseurs $W_S^{\mathrm{f}}$, $W_S$ sont \textit{égaux.}

\dgrblock{dgrb272}{Fine replacement gives an equivalence of localizations}
\begin{dcomment}{exposition}{p195}{Fine replacement gives an equivalence of localizations}
Grothendieck can now return to the question on manuscript p.~186 with an actual construction in hand. The essential point is stronger than the existence of a perfect object with the right total homotopy type: for a fibred input, the comparison is an equivalence in every fibre. This is what allows the replacement to control morphisms as well as objects.

Today W(6 bis) is automatic: every basic localizer is closed under filtered colimits \citep[Corollary~4.2.22]{Cisinski2006}. This applies to the standing relative W$(1,2,3)$, without accessibility; we do not extend it here to arbitrary weak basic localizers.

Let $R=\Phi_\infty$, let $\mathcal F=\Fib_W S$, and let $\mathcal P=\Parf_W S$. Under W(6 bis), the unit $\eta_X:X\to RX$ is fine for every $X\in\mathcal F$. The functor $R$ preserves fine equivalences on $\mathcal F$. Indeed, if $u:X\to Y$ is fine, then in each fibre the equality
\[
 (Ru)_s(\eta_X)_s=(\eta_Y)_s u_s
\]
and two-out-of-three show that $(Ru)_s$ lies in $W$. Thus both the inclusion $J:\mathcal P\hookrightarrow\mathcal F$ and $R$ induce functors after fine localization. The units give natural isomorphisms
\[
 \mathrm{id}\xrightarrow{\sim}\overline J\,\overline R
 \quad\text{on }\mathcal F[(W_S^{\mathrm f})^{-1}],
 \qquad
 \mathrm{id}\xrightarrow{\sim}\overline R\,\overline J
 \quad\text{on }\mathcal P[(W_S^{\mathrm f})^{-1}].
\]
For the second assertion, both endpoints of the unit belong to $\mathcal P$, and the subcategory is full. These two natural isomorphisms prove that the induced functors are quasi-inverse equivalences. In particular, full faithfulness is part of the result; there is no separate assumption about roofs or a calculus of fractions.

The same proof works with any full intermediate subcategory containing the image of $R$. The more precise claim in footnote~94 follows because the construction already lands in $\Parf_{W''_\omega}S$. One must keep the two minimal classes distinct at this stage of the manuscript. The inclusion $W_\omega\subseteq W''_\omega$ does not, on its own, put the output in $\Parf_\omega S$. It puts it in the class named in the footnote.

The minimality comparison above does, however, allow us to settle the stronger annotation in footnote~89. It identifies $W_\omega$ with the nerve weak-homotopy equivalences $W_\infty$. This class is stable under sequential colimits. Indeed, the nerve commutes with filtered colimits because each finite ordinal category $[n]$ is finitely presented. Weak equivalences of simplicial sets are stable under such colimits: apply $\operatorname{Ex}^{\infty}$, which commutes with them since every subdivided simplex is finite; a filtered colimit of Kan complexes is Kan, and maps and homotopies from finite simplicial spheres are witnessed at finite stages, giving the required isomorphisms on all homotopy groups and on components. Thus $W_\infty$ satisfies W(6 bis). Minimality gives
\[
 W_\omega\subseteq W''_\omega\subseteq W_\infty=W_\omega.
\]
Consequently the replacement can indeed be taken in $\Parf_\omega S$, with the fine unit just proved. This conclusion uses the modern identification of the minimal class; it was not a consequence of the single inclusion written on manuscript p.~195.

The original restricted W(4b$'$) now suffices as well. Given a total $W$-equivalence $u:X\to Y$ between weakly fibred objects, naturality of the unit makes $Ru$ a total $W$-equivalence between $\Parf_\omega$-objects. The restricted hypothesis makes $Ru$ fine, and the two fine units then make $u$ fine by two-out-of-three in each fibre. This proves full W(4b), and therefore the fine-to-coarse comparison. It explains precisely how the historical retraction acquires a positive answer once the later colimit and minimality theorems are available.

This replacement equivalence still does not identify the fine and coarse localizations. It allows either of them to be computed on a smaller category; it does not say that the two classes of denominators agree there. That remaining assertion is exactly the restricted W(4b) hypothesis in the following proposition. Grothendieck's distinction between replacing objects and comparing equivalences is therefore preserved even after the successful infinite construction.
\end{dcomment}


Il faut une conclusion pratique utilisable de toutes ces perplexités et observations. Essayons

\textbf{Proposition.} \textit{Supposons que l'on soit sous l'une des conditions suivantes.}

\begin{enumerate}[label={}]
\item[$1^o$)] \textit{L'axiome} W(4b) \textit{est satisfait dans le cas de catégories $X$, $Y$ qui sont dans} $\und{Fibsc}_{0,W} S$ (\footnote[95]{mais c'est \textit{équivalent} à W(4b)! Cf. XVI 4.9. loc. cit.}) \textit{(cette condition sur $W$ est en fait autoduale).}

\item[$2^o$)] \textit{L'axiome} W(4b) \textit{est satisfait dans le cas de catégories $X$, $Y$ qui sont dans} $\und{Parf}_{W''_\omega} S$\textit{, où $W''_\omega$ est le plus petit ensemble $\Sigma \subseteq$ Fl}(Cat) \textit{satisfaisant les conditions} $\mathrm{W}(1, 2\text{ bis}, 3, 6\text{ bis})$. \textit{(On peut aussi se contenter de la condition, plus forte mais plus compréhensible, que} W(4b) \textit{est OK si $X$, $Y$ sont dans} $\und{Parf}_W S$ (\footnote[96]{Mais cela implique aussi W(4b), cf. loc. cit.}).\textit{) De plus, on suppose que $W$ satisfait} W(6 bis).
\end{enumerate}

\textit{Sous ces conditions, le foncteur correspondant [?] $\lambda$}
\[
\mathrm{HOT}^{\mathrm{lc}}_W(S) \longrightarrow \mathrm{HOT}^{\mathrm{lc!}}_W(S)
\]
\textit{est une équivalence de catégories (cf. diagramme p. 178), donc $(\rho, i)$, $(\sigma, j)$ peuvent être regardés comme des couples de foncteurs entre deux catégories (cf. $(*)$ p. 181).} Ce sont des couples de foncteurs adjoints, \textit{donc $i$, $j$ sont pleinement fidèles, et on a la formule d'adjonction $(*)$, page 182.}

\msp{page 197}
Pour prouver l'adjonction de $(\rho, i)$ (ou dualement de $(\sigma, j)$), on exprime encore les deux catégories en présence $\mathrm{HOT}_W(S)$, $\mathrm{HOT}^{\mathrm{lc}}_W(S)$ comme $(\mathrm{Cat}/S)(W_S^{\mathrm{d}})^{-1}$, $(\mathrm{Cat}/S)(W_S)^{-1}$, le foncteur $\rho$ étant la localisation canonique allant du premier vers le second, et le foncteur $i$ étant donné par un foncteur d'inversion
\[
\gamma : (\mathrm{Cat}/S, W_S) \longrightarrow (\mathrm{Cat}/S, \overline{W_S^{\mathrm{d}}})
\]
qui est $C$ dans le cas $1^o$, $\Phi_\infty$ dans le cas $2^o$. Les morphismes d'adjonction
\[
\rho i \longrightarrow \mathrm{id}, \qquad i\rho \longleftarrow \mathrm{id}
\]
doivent se décrire alors comme dans le cas où on suppose tout bonnement W(4b), et les compatibilités se voient de même : on est ramené à prouver que deux flèches
\printedpage{122}
% head: A. Grothendieck : Dérivateurs, Ch. VI, Édition des Universités de Montpellier II et Paris VII
% foot: Version 9 novembre 2001
\[
\begin{tikzcd}
\gamma(X) \arrow[r, shift left, "i_{\gamma(X)}"] \arrow[r, shift right, "\gamma(i_X)"'] & \gamma(\gamma(X))
\end{tikzcd}
\]
sont égales dans $(\mathrm{Cat}/S)(W_S^{\mathrm{d}})^{-1}$. La vérification dans le cas $1^o$), où $\gamma = C$, ne doit pas être changée. Dans le cas $2^o$) où $\gamma = \Phi_\infty$, je ne l'ai pas fait. Je ne le ferai que sous la pression des besoins.

\subsection*{\circnum{6} Loi contravariante de $\mathrm{HOT}_W(S)$, $\mathrm{HOT}^{\mathrm{lc}}_W(S)$.}

Il y a pour
\[
f : S' \longrightarrow S
\]
un homomorphisme du diagramme
\[
\adjustbox{max width=\linewidth}{$\begin{tikzcd}[column sep=small]
& \mathrm{HOT}^{\mathrm{lc}}_W(S') \arrow[dl] \arrow[dr] & & \text{\texttt{[Plutôt}} & & \mathrm{HOT}^{\mathrm{lc}}_W(S) \arrow[dl] \arrow[dr] & \\
\mathrm{HOT}_W(S') & & \mathrm{HOT}^{o}_W(S') & & \mathrm{HOT}_W(S) & & \mathrm{HOT}^{o}_W(S)\ \text{\texttt{]}}
\end{tikzcd}$}
\]
\msp{page 198}
dans le diagramme similaire pour $S$ \texttt{[plutôt} $S'$\texttt{]}, par des foncteurs $f^*$ qui sont définis respectivement sur
\[
\begin{aligned}
\und{Liss}_W S &\quad \text{pour}\quad \mathrm{HOT}_W(S) & \quad\longrightarrow\quad & \mathrm{HOT}_W(S')\\
\und{Prop}_W S &\quad \text{pour}\quad \mathrm{HOT}^{o}_W(S) & \quad\longrightarrow\quad & \mathrm{HOT}^{o}_W(S')\\
\und{Fib}_W S &\quad \text{pour}\quad \mathrm{HOT}^{\mathrm{lc}}_W(S) & \quad\longrightarrow\quad & \mathrm{HOT}^{\mathrm{lc}}_W(S'),
\end{aligned}
\]
par l'intermédiare des foncteurs % [sic] intermédiare -> intermédiaire, missing i
\[
f^* : X \longmapsto X \times_S S' : \left\{
\begin{aligned}
\und{Liss}_W S &\quad\longrightarrow\quad \und{Liss}_W S'\\
\und{Prop}_W S &\quad\longrightarrow\quad \und{Prop}_W S'\\
\und{Fib}_W S &\quad\longrightarrow\quad \und{Fib}_W S',
\end{aligned}
\right.
\]
lequel foncteur est compatible avec les localiseurs pretinents, qui sont $W_S^{\mathrm{f}}$ dans tous les cas. Par contre, en l'absence d'hypothèses draconiennes sur $W$ ou sur $S$, on ne peut dire que $f^*$ soit compatible avec $W_S$, $W_{S'}$ sur une quelconque des catégories du diagramme p. 175, permettant de définir $\mathrm{HOT}^{\mathrm{lc!}}_W(S)$, de façon à nous donner un foncteur % [sic] pretinents -> pertinents, transposed letters
\[
\mathrm{HOT}^{\mathrm{lc!}}_W(S) \longrightarrow \mathrm{HOT}^{\mathrm{lc!}}_W(S').
\]
En l'absence d'axiomes sur $W$, le seul cas qui me soit connu est celui où $f$ est \textit{une $W$-équivalence universelle,} i.e. une $W$-fibration à fibres $W$-asphériques, auquel cas il est évident que pour \textit{toute} flèche $u : X \longrightarrow Y$ dans $\mathrm{Cat}/S$, $u \in W_S$ si et seulement si $(u' : X' \longrightarrow Y') \in W_{S'}$. On trouve alors un homomorphisme du diagramme complet $(*)$ p. 176 \texttt{[plutôt p. 178]} pour $S$, dans le diagramme correspondant pour $S'$.

\dgrblock{dgrb273}{Completing the omitted compatibility of the two units}
\begin{dcomment}{exposition}{p197}{Completing the omitted compatibility of the two units}
Grothendieck has proved enough to compare the localizations, but deliberately leaves the triangle check for $\Phi_\infty$ unfinished. There is a short way to complete it which uses the fine replacement already constructed. It also explains the later annotation that the proposed weakened hypotheses recover the full fibre-detection axiom.

First consider any $S$-map $u:X\to Y$ between $W$-fibred categories, with $u\in W$. Under condition~$1^o$, put $R=\Phi$; under condition~$2^o$, put $R=\Phi_\infty$. In either case $\eta_X$ and $\eta_Y$ are fine equivalences. Naturality and two-out-of-three show $Ru\in W$. In the first case $RX,RY$ are split fibrations with $W$-equivalent transports; in the second they lie in $\Parf_{W''_\omega}S$. The restricted fibre-detection hypothesis therefore gives $Ru\in W_S^{\mathrm f}$. Applying two-out-of-three in each fibre to
\[
 (Ru)_s(\eta_X)_s=(\eta_Y)_s u_s
\]
gives $u_s\in W$. Thus each displayed hypothesis implies W(4b) itself. In the first case, it is enough to impose the test only on Cartesian maps, since $\Phi(u)$ is Cartesian. This is the replacement argument formalized later in \citet[Lemma~4.4.9 and Corollary~4.4.10]{Grothendieck1991XVI}; the calculation here uses only the constructions already present in Chapter~VI.

Under condition~$1^o$, the established reflector $C$ already proves the proposition. We now work under condition~$2^o$, retaining W(6 bis), and compare $R=\Phi_\infty$ with $C$. Write $\alpha_X:X\to CX$ for the unit of $C$. Work in $\mathcal H_d(S)$, where the established reflection gives $C\simeq i\rho$. Let $R=\Phi_\infty$ and retain $\eta_X:X\to RX$ for its unit. The object $RX$ is $W$-fibred, so $\alpha_{RX}$ is an isomorphism in $\mathcal H_d(S)$. Also $\eta_X\in W_S$, hence $C(\eta_X)$ is right-fine under W(4b). We therefore have a natural isomorphism
\[
 \theta_X=(\alpha_{RX})^{-1}C(\eta_X):CX\xrightarrow{\sim}RX.
\]
Naturality of $\alpha$ gives the exact identity
\[
 C(\eta_X)\alpha_X=\alpha_{RX}\eta_X,
 \qquad\text{hence}\qquad \theta_X\alpha_X=\eta_X.
\]
Thus $R$ represents the same reflection, with its given unit. This is stronger than an unspecified isomorphism between the two endofunctors: it identifies the units, which are what determine the adjunction.

It remains to extract the equality Grothendieck asks for. Since $R^2X$ is a local object, the universal property of the reflection says that precomposition by $\eta_X$ is a bijection
\[
 \Hom_{\mathcal H_d(S)}(RX,R^2X)
       \xrightarrow{\sim}
 \Hom_{\mathcal H_d(S)}(X,R^2X).
\]
The two maps $\eta_{RX}$ and $R(\eta_X)$ have the same image under this bijection, because naturality of $\eta$ gives
\[
 \eta_{RX}\eta_X=R(\eta_X)\eta_X.
\]
Therefore $\eta_{RX}=R(\eta_X)$ in $\mathcal H_d(S)$, as required. Both maps are isomorphisms: the first because $RX$ is fibred, the second because $R$ sends coarse equivalences to fine ones after the fibre-detection argument. Inverting this common map supplies the multiplication of the idempotent reflector, and the triangle identities follow.

The missing verification is consequently a genuine omission in the written proof, but no additional axiom is needed to repair it. The repair is not the generally invalid cancellation of $\eta_X$ from an equality of composites; $\eta_X$ need not be invertible in the lateral category. Cancellation is justified here by its proven universal property against the local target $R^2X$.
\end{dcomment}


Si on suppose W(4b), alors $W_S = W_S^{\mathrm{f}}$ sur $\und{Fib}_W S$, donc sur cette catégorie-là le foncteur $f^*$ est compatible (sans mérite!) avec les localiseurs $W_S$, $W_{S'}$, mais ce n'est pas nouveau vraiment
\printedpage{123}
% head: A. Grothendieck : Dérivateurs, Ch. VI, Édition des Universités de Montpellier II et Paris VII
% foot: Version 9 novembre 2001
\msp{page 199}
par rapport à ce qui a été dit plus haut, pour $\und{Fib}_W S$, $\und{Fib}_W S'$, et pour les $\mathrm{HOT}^{\mathrm{lc}}_W$.

Mais si on suppose W(4a), et $f \in \und{Fib}_W$, alors $f^*$ sur $\mathrm{Cat}/S$ tout entier est compatible avec les localiseurs $W_S$, $W_{S'}$, et on trouve un calcul de $f^*$ sur $\mathrm{HOT}^{\mathrm{lc!}}_W(S) \approx \mathrm{HOT}^{\mathrm{lc}}_W(S)$ par
\[
f^*(\mathrm{hot}^{\mathrm{lc}}_{S,W}(X)) = \mathrm{hot}^{\mathrm{lc}}_{S',W}(X \times_S S').
\]

Donc c'est là un cas, exceptionellement, où les foncteurs $f^*$ pour les $\mathrm{HOT}_W$ sont compatibles non seulement aux foncteurs d'inclusion $i$, $j$ (de la page 176 \texttt{[plutôt p. 178]}), mais aussi aux foncteurs de localisation $\rho$, $\sigma$. % [sic] exceptionellement -> exceptionnellement, French spelling has double n

\textbf{Oubli :} Si $p : S' \longrightarrow S$ est $W$-propre, alors $p^* : \mathrm{Cat}/S \longrightarrow \mathrm{Cat}/S'$ est compatible avec $W_S^{\mathrm{d}}$, $W_{S'}^{\mathrm{d}}$, et $p^* : \mathrm{HOT}_W(S) \longrightarrow \mathrm{HOT}_W(S')$ peut se ‘calculer’ par la formule
\[
p^*(\mathrm{hot}_{S,W}(X)) = \mathrm{hot}_{S',W}(X \times_S S')
\]
pour \textit{tout} $X$ dans $\mathrm{Cat}/S$.

\subsection*{\circnum{7} Produit dans $\mathrm{HOT}_W(S)$.}

Le foncteur
\[
(X, Y) \longmapsto X \times_S Y : \und{Liss}_W S \times \und{Liss}_W S \longrightarrow \und{Liss}_W S
\]
est compatible avec les localiseurs $W_S^{\mathrm{f}} \times W_S^{\mathrm{f}}$ et $W_S^{\mathrm{f}}$, donc passe aux localisés en
\[
\left\{
\begin{aligned}
\mathrm{HOT}_W(S) \times \mathrm{HOT}_W(S) &\quad\longrightarrow\quad \mathrm{HOT}_W(S)\\
(\xi, \eta) &\quad\longmapsto\quad \xi \times_S \eta
\end{aligned}
\right.
\]
(\footnote[97]{\textbf{NB} Si $X$, $Y$ sont \textit{quelconques} dans $\mathrm{Cat}/S$, on a $\mathrm{hot}_{S,W}(X) \times_S \mathrm{hot}_{S,W}(Y) = \mathrm{hot}_{S,W}(\Phi_S(X) \times_S \Phi_S(Y))$. Si $X$ ou $Y$ est lisse sur $S$, on peut y remplacer $\Phi(X)$ resp. $\Phi(Y)$ par $X$ resp. $Y$.}). Ce produit est associatif, unitaire, commute avec toutes les compatibilités habituelles, comme il résulte du fait que c'est comme ça dans $\und{Liss}_S$. L'élément unité est la classe de $S \xrightarrow{\mathrm{id}} S$. La sous-catégorie $\mathrm{HOT}^{\mathrm{lc}}_W(S)$ est stable par le produit, qui induit
\msp{page 200}
\[
\mathrm{HOT}^{\mathrm{lc}}_W(S) \times \mathrm{HOT}^{\mathrm{lc}}_W(S) \longrightarrow \mathrm{HOT}^{\mathrm{lc}}_W(S),
\]
lui-même induit par le foncteur produit $(X, Y) \longrightarrow X \times_S Y$ dans $\und{Fib}_W S$.

\dgrblock{dgrb274}{Restriction depends on the chosen equivalences}
\begin{dcomment}{exposition}{p199}{Restriction depends on the chosen equivalences}
Grothendieck's three inverse-image functors are built from the same categorical pullback, but they are defined on different useful representatives. This is why the functorial diagram is more informative than an undifferentiated assertion that pullback preserves weak equivalences.

For a smooth representative $X\to S$, its pullback along $f:S'\to S$ is again smooth, and fibrewise equivalences remain fibrewise equivalences. Since fine and right-fine equivalences agree on smooth objects, this constructs $f^*$ on $\HOT_W$. The dual argument uses proper representatives for $\HOT_W^o$. On $\Fib_W S$, the same pullback preserves the class of objects and the fine equivalences, and therefore constructs $f^*$ on $\HOT_W^{\mathrm{lc}}$. In each case the composition and identity constraints are those of categorical pullback, passed through localization. Their coherence is inherited from the coherent pullback identifications before localization.

For the coarse localization, this argument is unavailable until total equivalences are controlled. If $f$ is universally in $W$, every projection $X\times_S S'\to X$ belongs to $W$. Given $u:X\to Y$, the natural square with these projections proves
\[
 u\in W\quad\Longleftrightarrow\quad
 u\times_S S'\in W
\]
by two-out-of-three. This proves both preservation and reflection of total equivalences; no good-representative restriction is needed. Under W(4a), the same preservation statement holds when $f$ is $W$-fibred, by the defining universal preservation property of that axiom. It then computes the coarse restriction directly on every $X\to S$.

There is a second, differently placed improvement. If $p:S'\to S$ is proper, pullback preserves right-fine equivalences on all of $\Cat/S$. The displayed formula for $p^*$ on an arbitrary representative is therefore legitimate. The condition is on the map changing the base, whereas the first construction required smoothness of the object being pulled back. Either condition can make the ordinary pullback compute the desired derived inverse image, but the two arguments test different arrows.

These distinctions also explain the compatibility claims. Restriction on fine representatives commutes with the inclusions of locally constant objects into the lateral theories. Commutation with the coarse localization requires preservation of the coarse denominators, and therefore the extra hypotheses just discussed. A diagram of functors cannot be promoted from one assertion to the other merely because the object formulas look identical.
\end{dcomment}


Si $X$ est quelconque dans $\mathrm{Cat}/S$, le foncteur
\[
\mathrm{mult}_{X/S} : Y \longmapsto X \times_S Y : \und{Liss}_W S \longrightarrow \mathrm{Cat}/S
\]
est compatible avec les localiseurs $W_S^{\mathrm{d}}$, car ce foncteur peut s'écrire comme le composé du foncteur correspondent \textit{à valeurs dans $\und{Liss}_W X$,} lequel est compatible avec les localiseurs % [sic] correspondent -> correspondant, adjective intended

% ---- from Der-ChVI-345_p126-129.tex
\setcounter{footnote}{97}

\printedpage{124}
% head: A. Grothendieck : Dérivateurs, Ch. VI, Édition des Universités de Montpellier II et Paris VII
% foot: Version 9 novembre 2001

$W_S^{\mathrm{d}}$, $W_S^{\mathrm{d}}$, et le foncteur ‘oubli de $X$’ de $\underline{\mathrm{Liss}}_W X$ (même de $\mathrm{Cat}/X$) dans $\mathrm{Cat}/S$, lequel est compatible avec $W_X^{\mathrm{d}}$ et $W_S^{\mathrm{d}}$. On trouve ainsi un foncteur % [sic] W_S^d, W_S^d -> apparently W_S^d, W_X^d (or W_X^d, W_S^d), the same symbol is printed twice
\[
\left\{\begin{aligned}
&\mathrm{mult}_{X/S} : \mathrm{HOT}_W(S) \longrightarrow \mathrm{HOT}_W(S)\\
&\mathrm{mult}_{X/S}(\mathrm{hot}_{S,W}(Y)) \overset{\text{déf}}{=} \mathrm{hot}_{S,W}(X\times_S Y) \qquad \text{si $Y$ \textit{lisse} sur $S$}
\end{aligned}\right.
\]
(\footnote{Cette formule est également valable si $Y$ est dans $\underline{\mathrm{Fib}}_W\, S$, cf. $(*)$, p. 202.}), puisque $\mathrm{HOT}_W(S)$ peut s'obtenir comme localisé tant de $\underline{\mathrm{Liss}}_W S$ que de $\mathrm{Cat}/S$, pour $W_S^{\mathrm{d}}$. Quand $X$ est lui-même $W$-lisse sur $S$, ce foncteur n'est autre que
\[
\left\{\begin{aligned}
&\mathrm{mult}_\xi : \eta \longmapsto \xi\times_S \eta,\\
&\text{où }\ \xi = \mathrm{hot}_{S,W}(X).
\end{aligned}\right.
\]

Mais il n'en est pas de même si $X$ n'est pas lisse, i.e. il faut distinguer dans ce cas, même pour $Y$ $W$-lisse sur $S$,
\[
\mathrm{hot}_{S,W}(X\times_S Y) \qquad\text{et}\qquad \mathrm{hot}_{S,W}(X)\times_S \mathrm{hot}_{S,W}(Y).
\]

Le deuxième terme peut s'expliciter comme $\mathrm{hot}_{S,W}(\Phi(X)\times_S Y)$ qui reçoit $\mathrm{hot}_{S,W}(X\times_S Y)$, i.e.
\msp{page 201}
on a un homomorphisme canonique
\[
\underbrace{\mathrm{hot}_{S,W}(X\times_S Y)}_{=\,\mathrm{mult}_{X/S}(Y)} \longrightarrow \underbrace{\mathrm{hot}_{S,W}(X)\times_S \mathrm{hot}_{S,W}(Y)}_{=\,\mathrm{hot}_{S,W}(\Phi(X)\times_S Y)},
\]
induit par
\[
X\times_S Y \longrightarrow \Phi(X)\times_S Y,
\]
mais cette dernière flèche n'a aucune raison d'être dans $W_S^{\mathrm{d}}$, et n'induit donc pas nécessairement un isomorphisme dans $\mathrm{HOT}_W(S)$. (Contrexemple page 154.)

\dgrblock{dgrb275}{Multiplication by a category over the base}
\begin{dcomment}{exposition}{p200}{Multiplication by a category over the base}
Grothendieck introduces $\mathrm{mult}_{X/S}$ because an arbitrary category over the base still acts on relative homotopy types, even when its strict fibre product does not compute the internal product. The action remembers the map $X\to S$ before replacement.

Let $p:X\to S$. For a smooth representative $Y\to S$, pullback computes $p^*[Y]$ as $[X\times_S Y]$ over $X$. Postcomposition then gives
\[
 \mathrm{mult}_{X/S}=p_!p^*.
\]
This identity already explains why the functor is well defined for arbitrary $X$: the right-fine equivalences are respected in the two stages where Grothendieck factors the construction. The repeated $W_S^{\mathrm d}$ at the start of printed p.~124 is a source slip; the pullback stage has source class $W_S^{\mathrm d}$ and target class $W_X^{\mathrm d}$.

By contrast, the internal product is obtained by replacing both factors by smooth categories:
\[
 [X]\times_S[Y]=[\Phi X\times_S\Phi Y].
\]
The map $X\to\Phi X$ induces a comparison from the action to multiplication by $[X]$. For smooth $Y$, it is represented by $X\times_S Y\to\Phi X\times_S Y$. Its failure to be right-fine in general is the counterexample already analyzed on manuscript pp.~153--156; there is no need to repeat it here. What changes in the present discussion is the proposed hypothesis that $X$ is $W$-fibred.

The argument given on pp.~201--202 correctly proves the action formula when $X$ is fibred. Indeed the unit $X\to\Phi X$ is fine; after arbitrary base change to $Y$ it is fine between $W$-fibred objects over $Y$, hence right-fine there. Forgetting the intermediate base preserves right-fine equivalences. This proves that the action of $Y/S$ on $[X]$ can be computed by the strict fibre product even for arbitrary $Y$.

There remains an apparent asymmetry: identifying that value with $[X]\times_S[Y]$ also requires replacing $Y$. Grothendieck proves this when $X$ is proper and suspects that the extra condition is necessary. The next development supplies the missing comparison for a $W$-fibred $X$ by using the transports in its split-fibration replacement. The argument will identify the canonical comparison map itself, so it also gives the natural compatibility needed for products and inverse images.
\end{dcomment}


Il faut regarder à part le cas \textit{où $X$ est dans} $\underline{\mathrm{Fib}}_W S$, j'ai envie de dire que dans ce cas on a
\[
\mathrm{hot}_{S,W}(X\times_S Y) \xrightarrow{\ \sim\ } \underbrace{\mathrm{hot}_{S,W}(X)\times_S \mathrm{hot}_{S,W}(Y)}_{=\,\mathrm{hot}_{S,W}(\Phi(X)\times_S\, \Phi(Y))} \tag*{$(*)$} % doubt resolved in DER_VI-c_corrections.csv; input note: under the brace "hot" is printed larger than "(Phi(X)x_S Phi(Y))", which stands at subscript level -- argument of hot_{S,W} or part of its subscript
\]
(\footnote{Cette formule n'est établie qu'en supposant ou bien $Y$ $W$-lisse (en plus de $X$ $W$-fibré sur $S$), ou bien $X$ $W$-propre, p.ex. $X$ $W$-parfait.}) pour \textit{tout} $Y$ dans $\mathrm{Cat}/S$ (pas nécessairement $W$-lisse), a fortiori que dans ce cas
\[
\left\{\begin{aligned}
&\mathrm{mult}_{X/S} = \mathrm{mult}_\xi \qquad \text{dans } \mathrm{HOT}_W(S)\\
&\text{où } \xi = \mathrm{hot}_{S/W}(X/S) \text{ (si $X$ dans $\underline{\mathrm{Fib}}_W S$).} % [sic] hot_{S/W} -> hot_{S,W}, the notation used everywhere else
\end{aligned}\right. \tag*{$(**)$}
\]

\printedpage{125}
% head: A. Grothendieck : Dérivateurs, Ch. VI, Édition des Universités de Montpellier II et Paris VII
% foot: Version 9 novembre 2001

Dans le cas où $Y$ est lisse sur $S$ (ce qui suffit pour établir $(**)$), il suffit de se rappeler que $X\longrightarrow \Phi(X)$ est dans $W_S^{\mathrm{f}}$ (résulte du fait que $X$ est $W$-fibré sur $S$ (\footnote{il suffit même $W$-prélisse sur $S$}), donc $X_s \longrightarrow s\backslash X$ est dans $W$, $\forall\ s\in S$), donc $\underbrace{X\times_S Y}_{X_Y} \longrightarrow \underbrace{\Phi(X)\times_S Y}_{\Phi(X)_Y}$ est aussi dans $W_Y^{\mathrm{f}}$, donc (comme $X_Y$, $\Phi(X)_Y$ sont $W$-fibrés sur $Y$) dans $W_Y^{\mathrm{d}}$, donc aussi dans $W_S^{\mathrm{d}}$. En fait, même sans supposer $Y$ lisse, ceci montre que l'on a
\msp{page 202}
\[
\mathrm{hot}_{S,W}(X\times_S Y) \xrightarrow{\ \sim\ } \mathrm{mult}_{Y/S}(\mathrm{hot}_{S,W}(X)) \qquad \text{si $X$ est dans $\underline{\mathrm{Fib}}_W S$.} \tag*{$(*)$}
\]

Pour avoir $(*)$ pour $Y$ quelconque, on peut par ce qui précède supposer $X$ lisse (quitte à le remplacer par $\Phi(X)$), donc $X$ dans $\underline{\mathrm{Liss}}_{0,W}S$. Il faut voir que
\[
X\times_S Y \longrightarrow X\times_S \Phi(Y)
\]
est encore dans $W_S^{\mathrm{d}}$. Or si $X$ est \textit{$W$-propre} sur $S$, le changement de base $X\longrightarrow S$ est compatible avec $W_S^{\mathrm{d}}$, $W_X^{\mathrm{d}}$ (sur $\mathrm{Cat}/S$ tout entier), et on gagne. (Ça doit être faux sans hypothèse de $W$-propreté sur $X$, \textit{ou} de $W$-lissité sur $Y$ \dots)

Quand on a
\[
f : S' \longrightarrow S,
\]
alors
\[
f^* : \mathrm{HOT}_W(S) \longrightarrow \mathrm{HOT}_W(S')
\]
est compatible avec les structures de produit.

\subsection*{\protect\cercle{8} Les foncteurs $f_!$ : théorème d'adjonction, propriétés du dérivateur $\mathrm{HOT}_W$.}

Soit
\[
f : S' \longrightarrow S
\]
dans Cat, alors le foncteur induit
\[
f_! : \mathrm{Cat}/S' \longrightarrow \mathrm{Cat}/S
\]
est compatible avec les localiseurs $(W_{S'}^{\mathrm{d}}, W_S^{\mathrm{d}})$, ou $(W_{S'}^{\mathrm{g}}, W_S^{\mathrm{g}})$, ou $(W_{S'}, W_S)$. On trouve donc un foncteur induit pour les localisées, et \texttt{[un]} diagramme commutatif
\msp{page 203}
\[
\begin{tikzcd}[column sep=small]
\mathrm{HOT}_W(S') \arrow[rrd, "\rho'"] \arrow[dd, "f_!"'] & & & & \mathrm{HOT}^{o}_W(S') \arrow[lld, "\sigma'"'] \arrow[dd, "f^{o}_!"] \\
& & \mathrm{HOT}^{\mathrm{lc}\,!}_W(S') \arrow[dd, "f^{\mathrm{lc}\,!}_!"] & & \\
\mathrm{HOT}_W(S) \arrow[rrd, "\rho"] & & & & \mathrm{HOT}^{o}_W(S). \arrow[lld, "\sigma"'] \\
& & \mathrm{HOT}^{\mathrm{lc}\,!}_W(S) & &
\end{tikzcd}
\]

\dgrblock{dgrb276}{A fibred factor suffices for the product formula}
\begin{dcomment}{development}{p202}{A fibred factor suffices for the product formula}
Grothendieck's desired formula on manuscript p.~201 is stronger than the proof immediately following it. The remaining difficulty can be resolved by calculating over a comma category: the arrow from its fixed starting object trivializes the transport of a locally constant fibration. This gives the unrestricted second factor that the manuscript leaves undecided.

\textbf{Product comparison for a fibred factor.}\ours\enspace
Assume W$(1,2,3)$. If $X\to S$ is $W$-fibred in the weak, universally stable sense used in this chapter, then for every $Y\to S$ the canonical map
\[
 [X\times_S Y]\longrightarrow [X]\times_S[Y]
 \quad\text{in }\HOT_W(S)
\]
is an isomorphism. No properness assumption on $X$ or smoothness assumption on $Y$ is required.

\emph{Proof.} We first treat a split Grothendieck fibration $X=\int_S F$ with $F:S^{\mathrm{op}}\to\Cat$ and $F(u)\in W$ for every arrow $u$ of $S$. Fix $s\in S$ and put $B=s\backslash Y$. An object of $B$ is $(y,a)$, with $a:s\to q(y)$, where $q:Y\to S$ is the structure map. The projection
\[
 E=s\backslash(X\times_S Y)\longrightarrow B
\]
is the split fibration classified by $(y,a)\mapsto F(qy)$. A morphism $(y,a)\to(y',a')$ is an arrow $h:y\to y'$ with $qh\,a=a'$, and its transport is $F(qh)$.

The maps $F(a):F(qy)\to F(s)$ form a natural transformation to the constant diagram at $F(s)$, because
\[
 F(a)F(qh)=F(qh\,a)=F(a').
\]
They therefore define a Cartesian functor
\[
 T:E\longrightarrow F(s)\times B,
 \qquad (x,y,a)\longmapsto (F(a)x,(y,a)).
\]
Each fibre map of $T$ belongs to $W$. Hence $T\in W$: a fibrewise $W$-equivalence between Grothendieck fibrations is right-fine, as seen by the fibre-to-comma adjunctions, and the relative localization axiom sends right-fine equivalences to total ones. This implication uses the fibrewise-to-total direction; it does not invoke W(4b).

We check that this proves the canonical comparison, not just an equivalence between its endpoints. Since $X$ is smooth, $[X]\times_S[Y]$ is represented by $X\times_S\Phi Y$. Let
\[
 k_s:s\backslash(X\times_S Y)
       \longrightarrow s\backslash(X\times_S\Phi Y)
\]
be induced by the unit of $Y$. The target is a comma category of a split fibration, so pullback to its fibre at $s$ defines a functor
\[
 r_s:s\backslash(X\times_S\Phi Y)
          \longrightarrow F(s)\times(s\backslash Y).
\]
The inclusion of that fibre is left adjoint to $r_s$: an arrow from a fibre object to an object over $t$, lying over $c:s\to t$, is uniquely an arrow in the fibre to its Cartesian pullback. Thus $r_s\in W$. On an object $(x,y,a)$ of the source of $k_s$, the unit of $Y$ introduces the identity arrow at $q(y)$, and $r_s$ pulls back along $a$. Consequently $r_s k_s=T$, on objects and on morphisms. Since $r_s,T\in W$, two-out-of-three gives $k_s\in W$. This holds for every $s$, proving
\[
 X\times_S Y\longrightarrow X\times_S\Phi Y
       \quad\text{belongs to }W_S^{\mathrm d}.
\]

For a general $W$-fibred $X$, its replacement $\Phi X$ is a split fibration with $W$-equivalent transports. The unit $X\to\Phi X$ is fine. After pullback to any $Y$, it is fine between $W$-fibred objects over $Y$, because this class of objects is stable under arbitrary base change. Fine and right-fine equivalences agree on that class, and forgetting $Y$ preserves right-fine equivalences. Hence
\[
 X\times_S Y\longrightarrow\Phi X\times_S Y
       \quad\text{belongs to }W_S^{\mathrm d}.
\]
Apply the split case to $\Phi X$. Composing the two maps gives a right-fine equivalence from $X\times_S Y$ to $\Phi X\times_S\Phi Y$, whose image is exactly the canonical product comparison. This proves the assertion.\hfill$\square$

The universal stability used in the last paragraph is part of the weak-fibration notion, not a hidden strengthening of it; see \citet[Chapter~XII, Part~III, Theorem~3, manuscript pp.~51--52]{Grothendieck1990XIIa}. The construction also shows naturality in both $X$ and $Y$, since each map was induced by the functorial comma replacement or by its Cartesian transport.

In particular, the proposed general formula $(*)$ on manuscript p.~201 is valid under its stated fibred hypothesis. The suspicion at the end of p.~202 that properness of $X$ or smoothness of $Y$ is necessary is not borne out. We checked the comma-category and base-change treatments in \citet[\S\S4--5]{Maltsiniotis2005Asphericite} and the locally constant framework in \citet{Cisinski2009LocallyConstant}; we did not locate this exact formulation there. The proof is included as a continuation of the manuscript's argument, without a claim that the result is new.
\end{dcomment}


\printedpage{126}
% head: A. Grothendieck : Dérivateurs, Ch. VI, Édition des Universités de Montpellier II et Paris VII
% foot: Version 9 novembre 2001

On n'a pas de foncteur $f_!^{\mathrm{lc}}$ au niveau des $\mathrm{HOT}_W^{\mathrm{lc}}$, sauf si on fait une hypothèse, telle \underline{W(4b)}, qui assure que $\mathrm{HOT}_W^{\mathrm{lc}\,!} \approx \mathrm{HOT}_W^{\mathrm{lc}}$.

Le théorème principal, que le foncteur $f_!$ est adjoint à gauche de $f^*$, est valable \textit{sous la seule hypothèse} W$(1,2,3)$ (\footnote{i.e. avec ces nouvelles notations Loc$(1,2,3\ \text{bis})$ (cf. XVI).}). Donc dualement pour $f_!^{o}$, $f_{o}^*$. Mais comme on n'a pas de foncteur $f_{\mathrm{lc}\,!}^*$, la question d'adjonction ne se pose pas pour $f$ général, et sans hypothèse W(4b) sur $W$. Elle ne se pose, apparemment, que pour $f$ un $W$-fibré trivial (cf. fin de \cercle{6}, p. 198), et dans ce cas l'adjonction ne doit pas faire de difficultés. Mais tant que je n'ai pas de sentiment bien net au sujet de la signification des $\mathrm{HOT}_W^{\mathrm{lc}\,!}$, ce n'est pas là un résultat bien bandant \texttt{[sic]}.

L'adjonction de $f_!$, $f^*$ pour $\mathrm{HOT}_W$ permet déjà de montrer que le prédérivateur $\mathrm{HOT}_W$
\msp{page 204}
satisfait les conditions Der 1, Der $4'$, Der $5'$ des dérivateurs. De plus, on a ce qu'on peut appeler la propriété Der 3b pour $\mathrm{HOT}_W$ en tant que $W$-dérivateur, i.e.

\begin{quote}
\textit{Si $f\in W$, alors $f$ est une} $\mathrm{HOT}_W$\textit{-équivalence} (\footnote{On le prouve en interprétant la notion de $\mathrm{HOT}_W$-équivalence par la $\mathrm{HOT}_W$\textit{-homologie.} Cela revient finalement à ceci : Si $f : S' \longrightarrow S$ est dans $W$, alors $\forall\, T$, $S'\times T \longrightarrow S\times T$ aussi!}) (moyennant les seuls axiomes W$(1,2,3)$ sur $W$). Plus précisément, $f$ est une $\mathrm{HOT}_W$-équivalence si et seulement si $\mathrm{hot}_W(f)$ est \texttt{[un]} isomorphisme, i.e. $f\in \overline{W}$, le saturé fort de $W$.
\end{quote}

Si $W$ \textit{satisfait} W(4b)\textit{,} alors on a également un couple de foncteurs adjoints
\[
\begin{tikzcd}
\mathrm{HOT}_W^{\mathrm{lc}}(S') \arrow[r, shift left, "f_!^{\mathrm{lc}}"] & \mathrm{HOT}_W^{\mathrm{lc}}(S) \arrow[l, shift left, "f^*_{\mathrm{lc}}"]
\end{tikzcd}
\]
pout toute flèche $S'\longrightarrow S$ dans Cat. On voit aussi dans ce cas que % [sic] pout -> pour, misprint

\begin{quote}
% printed with a vertical rule at the left of the block
\begin{tabular}{|p{0.9\linewidth}}
si $f : S' \longrightarrow S$ est dans $W$, alors $f^*_{\mathrm{lc}}$ et $f_!^{\mathrm{lc}}$ entre $\mathrm{HOT}_W^{\mathrm{lc}}(S)$ et $\mathrm{HOT}_W^{\mathrm{lc}}(S')$ sont des équivalences quasi-inverses l'une de l'autre (moyennant W(4b)),
\end{tabular}
\end{quote}

ce qui est l'axiome Der 3b$_2$ ‘version $W$-dérivateurs’. On ne peut s'attendre à mieux, du côté Der 3!

\textbf{Oubli :} Soit $f : S' \longrightarrow S$, on a alors $\boxed{f_!f^*} : \mathrm{HOT}_W(S) \longrightarrow \mathrm{HOT}_W(S)$ est égal à $\mathrm{mult}_{S'/S}$ (cf. p. 200). Il est donc égal à $\mathrm{mult}_\xi$, où $\xi = \mathrm{hot}_{S,W}(S')$, si on suppose ou bien $S'$ $W$-lisse sur $S$, ou bien $S'$ $W$-fibré sur $S$ (cf. $(**)$, p. 201).

\subsection*{\protect\cercle{9} Théorème de changement de base pour $f_!$ dans $\mathrm{HOT}_W$.}

Si on a
\printedpage{127}
% head: A. Grothendieck : Dérivateurs, Ch. VI, Édition des Universités de Montpellier II et Paris VII
% foot: Version 9 novembre 2001
\[
\begin{tikzcd}
X \arrow[d, "f"'] & X' \arrow[l, "q"'] \arrow[d, "f'"] \\
S & S', \arrow[l, "g"']
\end{tikzcd}
\]
avec $f$ $W$-lisse \textit{ou} $g$ $W$-propre \texttt{[la suite manque]}


\dgrblock{dgrb277}{The final adjunction and the surviving base-change fragment}
\begin{dcomment}{later}{p204}{The final adjunction and the surviving base-change fragment}
Grothendieck closes the revision by separating the constructions that need only the three standing axioms from the stronger locally constant comparison. Postcomposition $f_!$ respects the three relative equivalence classes, and the earlier unit and counit prove $f_!\dashv f^*$ for $\HOT_W$ without W(4b). The locally constant adjunction requires the additional comparison that identifies the fine theory with the coarse localization.

For $f\in W$ under W(4b), the locally constant unit and counit are isomorphisms; hence $f_!^{\mathrm{lc}}$ and $f^*_{\mathrm{lc}}$ are inverse equivalences. This assertion concerns locally constant coefficients. It does not say that restriction on all coherent diagrams is an equivalence whenever the indexing categories have the same homotopy type. For instance, the terminal category and the walking arrow have equivalent nerves, whereas the latter still indexes a genuine coherent morphism. The distinction is already built into the separate lateral and locally constant categories.

The constant-coefficient claim has an elementary calculation. If $p:S\to e$ and $T\in\Cat$, the constant coefficient represented by $T$ is represented over $S$ by $S\times T$. Its homotopy colimit is therefore represented by $S\times T$ as an absolute category. A functor $f:S'\to S$ acts by $f\times\mathrm{id}_T$. Product stability gives the implication $f\in W\Rightarrow f\times T\in W$. Conversely, testing $T=e$ detects whether $[f]$ is invertible in $\Cat[W^{-1}]$. This explains the appearance of the strong saturation $\overline W$ in the final statement, rather than an assertion that every chosen class $W$ was strongly saturated from the outset.

The last paragraph of the source displays a square and the hypotheses ``$f$ smooth or $g$ proper'', then stops. Its missing continuation is a source lacuna, not an unproved conclusion which we may attribute to Grothendieck. If the square is Cartesian, the expected statement can be read against the base-change result proved earlier in the chapter: the canonical mate
\[
 f'_!q^*\longrightarrow g^*f_!
\]
is invertible for such a square. The surviving fragment does not itself repeat the Cartesian hypothesis; it is necessary in this reconstruction. In the smooth-$f$ case, represent the coefficient by $Z\to X$ smooth over $X$. Its composite to $S$ is smooth, so both sides are represented by the same category $Z\times_S S'$ over $S'$. In the proper-$g$ case, ordinary pullback along $g$ computes inverse image for every representative, and its pullback $q$ is proper as well; the same equality of representatives proves the formula. The map is the standard mate built from the unit and counit, so the identification agrees with the canonical base-change comparison rather than merely producing some isomorphism between the functors.

This is the categorical form of the base-change mechanism isolated in \citet[Proposition~5.3, Theorem~5.4, and Remark~5.5]{Maltsiniotis2005Asphericite}, with the opposite-category convention adjusted to the present $s\backslash X$ model. The surviving text is left unchanged at its stopping point; the preceding calculation belongs to the commentary.
\end{dcomment}
\begin{otherlanguage}{english}
\dgroverview{VI-c}{Chapter VI \textperiodcentered\ \textup{VI-c} \textperiodcentered\ Overview}
\selectlanguage{english}
\noindent\textbf{Chapter VI, manuscript pages 138--204: operations, fibres, and the choice of equivalences.}

These pages bring the construction of $\HOT_W$ to a point where it can support a calculus of operations \citep{Grothendieck1990VI}. The principal achievement is an adjunction. For a functor $f:S'\to S$, composition of structural maps already gives a direct-image functor on categories over the base. After localization, this becomes the left adjoint $f_!$ of the inverse-image functor $f^*$. The work lies in making the inverse image homotopically meaningful, identifying the unit and counit, and proving that their two composites are identities. The ensuing discussion of products, base change, and locally constant objects repeatedly returns to the choice of a concrete representative and to the changes of representative that survive the intended operation.

The notation needs a fixed orientation. We follow the manuscript's presheaf convention. A split category fibred over $S$ represents a functor $S^{\mathrm{op}}\to\Cat$. For $X\to S$, the comma replacement $\Phi_S(X)$ has fibre $s\backslash X$ at $s$; its unit $i_X:X\to\Phi_S(X)$ is a right-fine equivalence. Consequently
\[
 f^*\hot_{S,W}(X)
 \simeq\hot_{S',W}\bigl(\Phi_S(X)\times_S S'\bigr),
 \qquad
 f_!\hot_{S',W}(X')=\hot_{S,W}(X').
\]
The second formula remembers the change of structural map. It does not assert that the representative remains fibred or smooth over the new base. This is the reason the first formula cannot generally be replaced by strict pullback of $X$.

The adjunction proof on pages~138--146 is unusually revealing. The unit is the $S'$-map induced by $i_{X'}$ and the structural map $X'\to S'$. For the counit one first tries to invert the pullback of $i_X$; the attempt fails because that pullback need not belong to the relevant class of equivalences. The correct formula is
\[
 \beta_X=(i_X)^{-1}\,\pr_1^{\Phi_S(X)}.
\]
Its inverse is taken in the localization over $S$, where $i_X$ is known to be invertible. One triangle identity is then the definition of the unit, and the other is naturality of $i$. Our commentary writes out this reduction, with the domains and codomains kept visible. The later passage to locally constant objects is a formal adjunction argument between reflective subcategories, once the reflection and the appropriate commutation squares are available.

\medskip\noindent\textbf{Products and the operation attached to a representative.}

The next issue is more subtle than the construction of a multiplication law might suggest. Two smooth representatives have a strict product over $S$ that computes the product in $\HOT_W(S)$. With arbitrary representatives one must use
\[
 \hot_{S,W}(X)\times_S\hot_{S,W}(Y)
 \simeq\hot_{S,W}\bigl(\Phi_S(X)\times_S\Phi_S(Y)\bigr).
\]
The first corollary on page~152 claims that one smooth factor suffices to omit both replacements. Grothendieck retracts it immediately. Over the walking arrow $0\to1$, the smooth inclusion of $0$ and the terminal-point inclusion of $1$ have empty strict intersection, whereas replacing the latter by its comma fibration changes the calculation. This small example determines the correct scope of every later product formula.

What remains for an arbitrary $X\to S$ is the endofunctor $\operatorname{multhot}_{X/S}$, obtained by multiplying smooth representatives by this particular $X$. It gives the unconditional identity
\[
 f_!f^*\simeq\operatorname{multhot}_{S'/S}.
\]
An identification with multiplication by the localized object $\hot_{S,W}(S')$ requires a further argument. The distinction concerns functorial dependence on the first representative, rather than a failure of associativity or of the usual coherence of products. We return to it at the end of the chapter, where a proof using the Grothendieck construction extends the calculation for weak $W$-fibrations. The proof and its precise assumptions are given in the development block, independently of an appeal to a general projection formula.

The formula for a left Kan extension at an object is obtained in the same concrete way. Its value at $s$ is represented by $s\backslash X'$, rather than by the strict fibre $X'_s$. The comma square carries its canonical natural transformation; it is not in general a strictly commutative square. With that transformation and the presheaf variance made explicit, the base-change mate has the standard orientation, and the computation proves it invertible. This supplies the pointwise left-Kan-extension axiom. The smooth/proper base-change theorem then follows by choosing representatives on which strict pullback is legitimate. The initial detour through a general derivator theorem is incomplete because it invokes pointwise detection of isomorphisms. The direct argument on pages~168--169 removes that dependence.

\medskip\noindent\textbf{Why four localizations appear.}

The revision beginning on page~171 is structural. Grothendieck separates the basic localizer axioms from the additional principle that a total equivalence between weak fibrations is an equivalence on every fibre. For a map $u:X\to Y$ over $S$, the classes used in the discussion are
\[
\begin{array}{ll}
 W_S &: u\in W\text{ on the total categories},\\
 W_S^{\mathrm f} &: u_s\in W\text{ for every strict fibre},\\
 W_S^{\mathrm d} &: s\backslash u\in W\text{ for every }s,\\
 W_S^{\mathrm g} &: u/s\in W\text{ for every }s.
\end{array}
\]
On weak fibrations, the fine conditions agree; their implication to total equivalence is available under the basic axioms. The converse is W(4b). Without it, the fine localization of weak fibrations and their coarse localization cannot be identified merely because both admit useful replacements. This yields $\HOT_W$, its opposite-variance counterpart, $\HOT_W^{\mathrm{lc}}$, and $\HOT_W^{\mathrm{lc}!}$, with a localization $\lambda$ between the last two.

The path functor $C$ exposes the difficulty. Its unit is a total equivalence and its values are perfect over the base. These properties give coarse equivalences of presentations. They do not, by themselves, prove that the unit is fine on a weak fibration. W(4b) supplies exactly that step, identifies the two locally constant theories, and allows the reflection formulas to be read as genuine adjunctions with fully faithful right adjoints. The audit keeps this hypothesis attached to each use of that conclusion. A calculation with action groupoids makes the distinction visible on morphisms: coarse localization can identify maps that remain distinct in the fine theory.

This is also the place to separate historical uncertainty from present knowledge. The minimality theorem for the usual homotopy localizer, in the modern basic-localizer framework, is due to \citet{Cisinski2004Minimal}. It answers the proposed search for a smaller basic localizer; it does not identify every generalized coefficient theory with ordinary homotopy theory. The manuscript's varying axioms and its changing use of ``fibration'' must still be translated before a later theorem is applied.

\medskip\noindent\textbf{Alternating replacement and the questions it answers.}

One application of $\Phi$ cannot produce all the desired geometry. The fork-shaped base with two incomparable terminal branches gives a particularly strong obstruction: for a nondegenerate localizer, the comma replacement of a nonempty category over that base cannot be proper. The commentary verifies the categorical retraction in the preceding example, repairs the central object of the printed diagram, and gives the missing last step in the emptiness argument. A separately marked extension treats a class of partially ordered bases by the same mechanism. These arguments concern properness of a particular replacement, not the nonexistence of any satisfactory presentation.

Grothendieck then alternates the two comma constructions and takes the colimit of the resulting sequence. The odd and even subsequences respectively provide smooth and proper presentations. The proof uses W(6 bis): closure of $W$ under colimits of sequences of arrows. Its historical formulation does not assert arbitrary filtered closure. The limiting object is perfect, and its unit is fine when the original object is a weak fibration. This gives a fine equivalence between the weak-fibration and perfect-object presentations without using W(4b). It does not identify fine and coarse equivalences.

The later theorem of \citet[Corollary~4.2.22]{Cisinski2006} makes W(6 bis) automatic for basic localizers. With the identification $W_\omega=W'_\omega=W''_\omega=W_\infty$ established in the commentary, the replacement therefore answers the perfect-presentation question affirmatively under the standing axioms. The same fine units transfer the original restricted W(4b$'$) test to all weak fibrations. These conclusions concern the relative localization axiom W(2); they do not silently extend a colimit theorem to every weaker axiom system considered in the manuscript.

The practical proposition on pages~196--197 combines this construction with restricted fibre detection. Its last adjunction verification is expressly postponed. We supply it by first transporting fibre detection along the functorial replacement and then comparing the resulting unit with the established $C$-reflection. The later Chapter XVI confirms the role of replacement in relating the fibration axioms; the brief status block locates its theorem and keeps its historical hypotheses visible. No priority claim is attached to the elementary arguments supplied here.

\medskip\noindent\textbf{The reach of the construction.}

For the ordinary homotopy theory of spaces, the contemporary description of locally constant diagrams uses homotopy Kan extensions and localization of diagram categories; \citet{Cisinski2009LocallyConstant} gives a precise treatment. The model-category construction of derivators supplies a standard framework \citep{Cisinski2003,Groth2013}. More generally, Der~2 holds for every basic localizer \citep[Propositions~3.1.30 and~4.2.7]{Cisinski2006}. The commentary after manuscript p.~159 explains the comparison which permits that conclusion. The unrestricted strict-bifibration presentation over an arbitrary base remains outside what we have established: a free bifibration alone does not provide a fine equivalence from the given weak fibration. Over groupoid bases we give the complete comparison directly.

The surviving final recap deliberately lists Der~1, Der~$4'$, Der~$5'$, and the appropriate W-relative form of Der~3b. Der~2 has been withdrawn from the earlier list pending a proof; its later resolution does not supply a missing argument to the historical text. The audit also preserves the missing manuscript page~170 and the incomplete final base-change paragraph. What can be proved from the surviving material is supplied in commentary; the missing text is not reconstructed.
\end{otherlanguage}
\dgrchunkend

% DER_VII-a: Grothendieck's text with the commentary set in, then the overview, inside a group that applies the chunk's own definitions
\dgrchapter{VII}{Chapter VII}{Catégories de chemins (2)}
\dgrpart{0}{VII-a}{Part 1 of 3, manuscript pages 1--57}{VII-a\enspace\textperiodcentered\enspace Chapter VII, part 1 of 3\enspace\textperiodcentered\enspace manuscript pages 1--57}{Chapter VII \textperiodcentered\ \textup{VII-a}}
\dgrchunkbegin
% DER_VII-a: the definitions of the chunk's preamble, applied inside its group
\let\printedpage\relax\newcommand{\printedpage}[1]{\if\relax\detokenize{#1}\relax\else\marginpar{\raggedright\footnotesize p.~#1}\fi}
\let\msp\relax\newcommand{\msp}[1]{\marginpar{\raggedright\footnotesize [#1]}}
\let\Fltr\relax\newcommand{\Fltr}{\operatorname{Fl}\operatorname{tr}}
\let\Ob\relax\newcommand{\Ob}{\operatorname{Ob}}
\let\bDelta\relax\newcommand{\bDelta}{\boldsymbol{\Delta}}
\let\bD\relax\newcommand{\bD}{\boldsymbol{\Delta}}
\let\iddots\relax\newcommand{\iddots}{\mathinner{\mkern1mu\raise1pt\hbox{.}\mkern2mu\raise4pt\hbox{.}\mkern2mu\raise7pt\hbox{.}\mkern1mu}}
\let\Ch\relax\newcommand{\Ch}{\underline{\mathrm{Ch}}}
\let\uHom\relax\newcommand{\uHom}{\underline{\mathrm{Hom}}}
\let\uFl\relax\newcommand{\uFl}{\underline{\mathrm{Fl}}}
\let\opp\relax\newcommand{\opp}{\mathrm{opp}}
\let\Hom\relax\newcommand{\Hom}{\mathrm{Hom}}
\let\uCh\relax\newcommand{\uCh}{\underline{\mathrm{Ch}}}
\let\Ty\relax\newcommand{\Ty}{\mathrm{Ty}}
\let\uTy\relax\newcommand{\uTy}{\underline{\mathrm{Ty}}}
\let\Cat\relax\newcommand{\Cat}{\mathrm{Cat}}
\let\circled\relax\newcommand{\circled}[1]{\text{\textcircled{$\scriptstyle #1$}}}
\let\ul\relax\newcommand{\ul}[1]{\underline{\mathrm{#1}}}
\let\red\relax\newcommand{\red}{\text{\upshape réd}}
\let\Tr\relax\newcommand{\Tr}{\mathrm{Tr}}
\let\Cr\relax\newcommand{\Cr}{\mathrm{Cr}}
\let\lng\relax\newcommand{\lng}{\operatorname{long}}
\let\card\relax\newcommand{\card}{\operatorname{card}}
\let\uTyp\relax\newcommand{\uTyp}{\underline{\mathrm{Typ}}}
\let\ovalbox\relax\newcommand{\ovalbox}[1]{\tikz[baseline=(X.base)]\node[draw,rounded corners=6pt,inner sep=2pt](X){#1};}

% ---- from Der-ChVII-456_p001-005.tex
% ===================== page 1 =====================
\printedpage{}
% head: (none; title page, no page number printed)

\begin{center}
{\Large\textsf{LES DÉRIVATEURS}}

\bigskip

\textbf{ALEXANDRE GROTHENDIECK}

\bigskip\bigskip

\textbf{Chapitre VII}

\medskip

\textbf{Catégories de chemins (2)}
\end{center}

\newpage

% ===================== page 2 =====================
\printedpage{}
% head: (none; no page number printed)

Ce texte a été déchiffré et transcrit en \LaTeXe\ par M. Künzer. Il a été édité par M. Künzer, J. Malgoire, et G. Maltsiniotis. La transcription est aussi fidèle que possible au manuscrit. Pour les quelques corrections évidentes, ou rares commentaires des éditeurs, ainsi que pour la numérotation originale des pages du manuscrit, les caractères de machine à écrire \texttt{[typewriter]} entre crochets sont utilisés. Un point d'interrogation entre crochets signifie que l'on n'est pas sûr du mot qui précède.

Cette édition est provisoire. Les remarques, commentaires et corrections sont bienvenus. Envoyer un message à :

\textbf{maltsin@math.jussieu.fr}

\hfill G. Maltsiniotis\hspace*{3em}

\newpage

% ===================== page 3 =====================
\printedpage{1}
% head: A. Grothendieck : Dérivateurs, Ch. VII, Édition des Universités de Montpellier II et Paris VII
% foot: Version 22 décembre 2001

\msp{page 1}
Comme le formalisme ‘simplifié’ des chemins que j'ai voulu dégager à la va-vite s'est révélé d'une lourdeur insupportable, de nature à le rendre quasiment inutilisable, je vais essayer de dégager un formalisme souple, collant à l'intuition, au prix seulement d'un effort conceptuel supplémentaire.

\section*{1\quad Intervalles et leur variantes} % [sic] leur variantes -> leurs variantes, plural agreement

On appelle \textit{catégorie intervalle standard} une catégorie définie par un ensemble ordonné, dont l'ensemble des élements est un intervalle d'entiers $[0,n]$ ($n\in\mathbf{N}$, on n'exclut pas le cas $n=0$, i.e. de la catégorie ponctuelle $\{0\}$) : % [sic] élements -> éléments, missing accent
\[
\Ob I=[0,n]\subseteq\mathbf{N},
\]
et satisfaisant l'hypothèse suivante :

\begin{enumerate}
\item[a)] pour toute paire de deux sommets consécutifs $i-1$, $i$ de $I$ ($1\leq i\leq n$), il y a une flèche qui les relie, donc soit $i\longrightarrow i+1$, soit $i\longleftarrow i+1$ \texttt{[plutôt} $i-1\longrightarrow i$\texttt{,} $i-1\longleftarrow i$\texttt{]} (cette flèche est notée $\underline{u}_i$,
\[
(i-1)\xrightarrow{\underline{u}_i}i\qquad\quad\text{ou}\qquad\quad(i-1)\xleftarrow{\underline{u}_i}i\ \ ).
\]
\item[b)] La catégorie $I$ est engendrée par l'ensemble de ses flèches $\underline{u}_i$ (appelées \textit{flèches de transition}).
\end{enumerate}

La flèche de transition est dite \textit{directe} si elle va de $i-1$ à $i$, \textit{rétrograde} dans le sens inverse. Si on affecte l'intervalle $[i-1,i]$ du signe $+$ ou $-$ suivant que $\underline{u}_i$ est directe ou rétrograde, on voit donc que les
\msp{page 2}
intervalles standard de longueur donnée $n$ correspondent biunivoquement aux systèmes de signes $+$ ou $-$ sur les $n$ intervalles de transition $[i-1,i]$ ($1\leq i\leq n$), il y en a donc $2^n$. Le sommet $0$ est l'\textit{origine} ou la \textit{source} de l'intervalle, le sommet $n$ l'\textit{extrémité} ou le \textit{but} de l'intervalle.

\dgrblock{dgrb278}{A combinatorial interval carries two orders}
\begin{dcomment}{exposition}{p1}{A combinatorial interval carries two orders}
Grothendieck begins again because a path should be easy to concatenate and easy to change without losing its endpoints. The extra conceptual effort announced here consists in making its shape an object of study. Before forming a category of all paths, he isolates the finite categories which can serve as their domains. The commentary on Chapter~VI, manuscript p.~69 (``Concatenation and the two iterated units'') discusses the preceding path construction; the present chapter reorganizes its combinatorics.

A standard interval is the free category on an oriented path graph. Write its transition signs as
\[
 \tau=(\varepsilon _1,\ldots,\varepsilon _n),\qquad
 \varepsilon_i\in\{+,-\}.
\]
There is one object at each integer from $0$ to $n$. A positive sign supplies the generating arrow $i-1\to i$ and a negative sign supplies $i\to i-1$. The free category is thin: between two fixed vertices there is at most one directed path, since the underlying undirected graph has no cycle. Thus the description as a category associated with a partially ordered set introduces no additional relations. Conversely, the two conditions in the text reconstruct precisely this free category.

For later use we record the complete hom-set calculation. If $i<j$, then
\[
\begin{split}
 I_\tau(i,j)\ne\varnothing
   &\quad\Longleftrightarrow\quad
   \varepsilon_{i+1}=\cdots=\varepsilon_j=+,\\
 I_\tau(j,i)\ne\varnothing
   &\quad\Longleftrightarrow\quad
   \varepsilon_{i+1}=\cdots=\varepsilon_j=-.
\end{split}
\]
Every nonempty hom-set contains exactly one arrow. These formulas prove both the classification by sign words and the count $2^n$. They also explain why a functor out of an interval will be specified freely by its values on the transition arrows.

The numerical order and the categorical partial order have different jobs. The first tells us in which order the vertices are visited; the second tells us which composites exist. For the shape $0\to1\leftarrow2$, the designated source of the interval is $0$ and its designated target is $2$, although neither statement says that $0$ is initial or $2$ terminal. A path of this shape travels along the first arrow and against the second. Confusing these two orders would already make the later distinction between an opposite path and an opposite category disappear.

Length counts transitions. The interval of length zero has one vertex and no transition, and the interval of length $n$ has $n+1$ vertices. We keep this convention throughout; it will identify the isolated equality $n=\operatorname{card}I$ on ms.~13 as a printed slip rather than a competing convention.
\end{dcomment}


Un intervalle standard est \textit{orienté,} du fait qu'on à la relation d'ordre \textit{total} naturelle sur l'ensemble de ses sommets, savoir celle des entiers dans $[0,n]$. Quand on parlera de la relation d'ordre sur $I$, qu'un sommet $i$ en \textit{‘précède’} ou \textit{‘suit’} un autre $j$ ($i\leq j$ ou $i\geq j$ respectivement), c'est de cette relation d'ordre qu'on voudra parler. Celle associée à la structure de catégorie ne sera jamais regardée comme une relation d'ordre. Il y a également une relation d'ordre total naturelle sur l'ensemble des flèches de transition $\underline{u}_i$ ($1\leq i\leq n$), donnée par la relation d'ordre des entiers $i$, $1\leq i\leq n$, qui les indexent. On dira qu'une flèche $i$ \texttt{[plutôt} $\underline{u}_i$\texttt{]} de transition \textit{précède} ou \textit{suit} une autre \texttt{[}$\underline{u}_j$\texttt{]}, ou qu'elle est \textit{successeur immédiat} ou \textit{prédécesseur immédiat} d'une autre (si $i=j+1$ ou $i=j-1$). % [sic] qu'on à -> qu'on a, verb; [sic] un sommet i en ‘précède’ -> un sommet i ‘précède’, superfluous "en"

% Reader layout: continuous text; edition coordinate retained below.
% ===================== page 4 =====================
\printedpage{2}
% head: A. Grothendieck : Dérivateurs, Ch. VII, Édition des Universités de Montpellier II et Paris VII
% foot: Version 22 décembre 2001

Si on a un ensemble d'entiers
\[
E\subseteq\mathbf{Z},
\]
on considérera dans $E$ la relation d'équivalence engendrée par la relation entre deux éléments $(i,j)$ de $E$ : $j$ est successeur immédiat de $i$ \textit{dans} $\mathbf{Z}$ (i.e. $j=i+1$). Deux éléments distinct $i$, $j$ de $E$ sont congrus pour cette relation si et seulement si l'intervalle d'entiers qu'ils déterminent est $\subseteq E$, donc % [sic] distinct -> distincts, plural agreement
\[
\text{si } i<j,\ i \text{ et } j \text{ sont \textit{‘connectés’} dans } E \text{ si et seulement si } [i,j]\subseteq E.
\]

\msp{page 3}
Quand on parlera des ‘composantes connexes’ de $E$, il s'agira toujours des classes suivant cette relation d'équivalence.

Nous identifierons l'ensemble des flèches de transition $\Fltr(I)$ de $I$ à l'intervalle d'entiers $[1,n]$, donc à une partie de $\mathbf{Z}$. Si on a une partie $E\subseteq\Fltr(I)$, on y a donc la notion de flèches de transition ‘connectées’ dans $E$, et la notion de connexité pour $E$. Si $i<j$, $\underline{u}_i$ et $\underline{u}_j$ sont connectées si et seulement si toutes les flèches de transition intermédiaires entre $\underline{u}_i$ et $\underline{u}_j$ sont dans $E$. Et $E$ est connexe si et seulement s'il est soit vide, soit formée d'une suite de flèches de transition dont chacune est successeur immédiat de celle qui la précède, % [sic] soit formée -> soit formé, agreement with "il"
\[
E=\{\underline{u}_i,\underline{u}_{i+1}\ldots,\underline{u}_{i+l-1}\}, % [sic] \underline{u}_{i+1}\ldots -> \underline{u}_{i+1},\ldots, missing comma
\]
où $l$ est le cardinal de $E$. Une telle $E$ est appelée une \textit{chaîne} de flèches de transition, et $l$ est sa longueur. Une chaîne est dite \textit{composable} si le composé $\underline{u}_{i+l-1}\circ\cdots\circ\underline{u}_{i+1}\circ\underline{u}_i$ \texttt{[ou} $\underline{u}_i\circ\underline{u}_{i+1}\circ\ldots\circ\underline{u}_{i+l-1}$\texttt{]} a un sens, i.e. si elles sont toutes de même sens, soit directe, soit rétrograde.

On désigne par $\Fltr^+(I)$, $\Fltr^-(I)$ l'ensemble des flèches de transition directes, ou rétrogrades, respectivement. Leurs composantes connexes ne sont autres que les chaînes composables maximales. Ainsi, les chaînes composables maximales forment une partition de l'ensemble des flèches de transition, en intervalles d'entiers mutuellement disjoints. Les chaînes forment elles-mêmes un ensemble totalement ordonné
\[
C_1,\ C_2,\ \ldots\ C_r\qquad\quad(1\leq r\leq n).
\]
L'entier $r$ s'appelle la \textit{‘longueur utile’} de l'intervalle
\msp{page 4}
$I$, c'est donc le cardinal de l'ensemble de ses chaînes composables maximales.

Deux chaînes composables maximales \texttt{[consécutives]} sont toujours de signe opposé, si l'une est directe l'autre est rétrograde, et inversement. La première chaîne composable maximale $C_1$ (disons simplement \textit{‘chaîne structurale’}) a le même signe que $\underline{u}_1$; \texttt{[on]} dit que la \textit{direction initiale} de $I$ est directe, resp. rétrograde, s'il en est ainsi de $\underline{u}_1$, ou ce qui revient au même, de $C_1$. De même on définit la \textit{direction finale} de $I$ comme celle de $\underline{u}_n$, ou encore celle de $C_r$. Direction initiale et finale sont égales si la longueur utile $r$ est
\printedpage{3}
% head: A. Grothendieck : Dérivateurs, Ch. VII, Édition des Universités de Montpellier II et Paris VII
% foot: Version 22 décembre 2001
impaire, inégales si elle est paire. (\textbf{NB} Ce n'est que pour $n=0$, i.e. quand $I=\{0\}$ est la catégorie ponctuelle, que direction initiale et direction finale ne sont pas définies.)

Les chaînes composables correspondent aux sous-catégories pleines de $I$ qui sont isomorphes à une catégorie $\bDelta_i$ (où $i$ est la longueur de la chaîne). Notons que $\bDelta_i$ n'a pas d'autre automorphisme que l'identité, donc les chaînes maximales de longueur donnée $i\geq 1$ correspondent exactement aux foncteurs
\[
\bDelta_i\lhook\joinrel\longrightarrow\ I
\]
qui sont injectifs sur les sommets. Un tel foncteur est pleinement fidèle, et un monomorphisme. Ce foncteur définit une application strictement croissante ou strictement décroissante sur
\msp{page 5}
les sommets, suivant que la chaîne est directe ou rétrograde.

\dgrblock{dgrb279}{Runs, simplices, and the meaning of maximality}
\begin{dcomment}{exposition}{p4}{Runs, simplices, and the meaning of maximality}
Grothendieck next separates the places where a path can be composed from the places where its direction turns. A maximal constant-sign run of the word $\tau$ is his structural chain. If the run lengths are $l_1,\ldots,l_r$, then
\[
 n=l_1+\cdots+l_r,
\]
and the sign of the $a$th run is $(-1)^{a-1}\varepsilon_1$, when signs are written as $\pm1$. This proves the parity rule for initial and final directions. The useful length $r$ records the number of changes of compositional regime, rather than the number of arrows which happen to have been written down.

The paragraph identifying chains with simplex subcategories needs a small but substantive separation. A chain was defined on ms.~3 to be a contiguous set of transitions. An object-injective functor $\Delta_k\to I$ may skip vertices. For example,
\[
 \Delta_1\longrightarrow\Delta_2,
 \qquad 0\longmapsto0,\quad1\longmapsto2,
\]
is fully faithful, but its image is not the domain of one transition of $\Delta_2$; nor does it describe its maximal structural chain, whose length is two. The words ``chaînes maximales'' and the unrestricted injectivity condition both occur in the edition. They therefore remain in the transcription, with their correction supplied here.

The exact statement is useful in its own right. For $k>0$, an injective functor from $\Delta_k$ picks $k+1$ vertices in a single structural interval, in their categorical order. Indeed its first and last image vertices are connected by an arrow, so every vertex between them in the numerical order belongs to the same constant-sign run. All intermediate image vertices lie between these two in that run's categorical order. Conversely, any such selection determines an injective functor. It is fully faithful because the target is thin and its order on those vertices is exactly the simplex order. Its object function is strictly increasing in the external orientation for a positive run and strictly decreasing for a negative run.

To recover a chain in the earlier sense, impose numerical contiguity of the image. To recover a maximal structural chain, require the image to be the whole structural interval. These distinctions do not disturb the intended decomposition into runs. They identify which extra condition the printed correspondence omitted, while preserving its correct full-faithfulness assertion.
\end{dcomment}


Notons que tout simplexe type $\bDelta_i$, $i\geq 0$, peut être vu comme un intervalle, de longueur $i$, de longueur utile $1$ (si $i\geq 1$, $0$ si $i=0$). Son origine est $0$, son extrémité est $i$. Mais appelons une catégorie quelconque une \textit{catégorie intervalle} si elle est isomorphe à un intervalle standard. Notons ici

\textbf{Proposition 1.} \textit{Soit $I$ un intervalle standard de longueur $n\geq 1$.}

\begin{enumerate}
\item[\textit{a)}] \textit{Soit $i\in I$. Pour que $i\in\{0,n\}$, i.e. pour que $i$ soit un des deux} ‘bouts’ \textit{(ou} ‘extrémités’\textit{) (source ou but) de $I$, il faut et il suffit que l'une des deux catégories $I/i$ ou $i\backslash I$ soit la catégorie ponctuelle et que l'autre soit isomorphe à un} $\bDelta_j$ (\footnote{où $j$ est la longueur de l'unique chaîne structurale qui ‘touche’ $i$ (i.e. telle que $i$ soit dans la sous-catégorie pleine définie par ladite chaîne).}). \textit{(Caractérisation intrinsèque des bouts d'un intervalle.) Pour que $i$ soit la source, il faut et il suffit que $i\backslash I$} \texttt{[phrase incomplète]}
\item[\textit{b)}] \textit{Tout foncteur $I\longrightarrow I$ qui fixe les deux extrémités de $I$ est l'identité.}
\end{enumerate}

\textbf{Corollaire.} \textit{Le groupe des automorphismes de $I$ a au plus deux éléments. S'il en a deux, l'automorphisme non trivial échange les deux extrémités. (On dit alors que l'intervalle est} symétrique.\textit{)}

\dgrblock{dgrb280}{The endpoint determines an orientation, not an endofunctor}
\begin{dcomment}{exposition}{p5:Le groupe des automorphismes}{The endpoint determines an orientation, not an endofunctor}
The purpose of Proposition~1 is to make orientation intrinsic enough to transport it through isomorphisms. Its first assertion is correct. At an endpoint, one of the two slices is a point and the other is the linearly ordered category running along the unique adjacent structural chain. An interior vertex at which the signs agree has both a nonidentity incoming and a nonidentity outgoing arrow, so neither slice is a point. At a turning vertex, one slice is a point, but the other has two nonempty branches meeting at their common source or common target. In the source case it has two incomparable maximal objects, and in the target case two incomparable minimal objects. It cannot be a simplex category. This gives the stated characterization of the two endpoints.

The second assertion is false with the word ``foncteur'' which the edition prints. On $\Delta_2$ the monotone map
\[
 0\longmapsto0,\qquad1\longmapsto0,\qquad2\longmapsto2
\]
defines a nonidentity endofunctor fixing both endpoints. The arrow $1\to2$ is sent to the composite $0\to2$; there is no requirement that a functor preserve a generating transition as a generating transition. This is exactly the flexibility which the later category of types is intended to accommodate.

What the subsequent argument proves is that an \emph{automorphism} fixing the endpoints is the identity. Here is a direct proof which also supplies the uniqueness statements used on ms.~7--9. A nonidentity arrow of an interval is indecomposable into two nonidentity arrows if and only if it is a transition: the displayed hom-set formula proves this, by factoring every longer directed segment at an intermediate vertex. Isomorphisms preserve indecomposability. The graph of indecomposable arrows, after forgetting directions, is the ordinary path with vertices $0,\ldots,n$. A graph automorphism fixing $0$ fixes its unique neighbour, then that neighbour's still unvisited neighbour, and so on. It fixes every vertex and hence every arrow. Every graph automorphism is consequently the identity or vertex reversal. The categorical automorphism group is a subgroup of this two-element group.

This is also the intrinsic reconstruction of the external total order: choose one endpoint, call it $s_0$, and successively take the unique neighbour not already visited. Grothendieck's slice argument implements the same reconstruction in categorical terms. In its inductive paragraph the slices must be taken in the remaining full tail $I'$ at $s_i$; the printed $i\backslash I$ and $I/i$ have dropped that qualification. Neither reconstruction is preserved by arbitrary functors, which may collapse or lengthen transitions.

The unfinished sentence about characterizing the \emph{source} of the interval should not be completed by guessing. The underlying category determines the pair of endpoints, but a symmetric interval has an automorphism interchanging them. Selecting one of them is precisely the additional orientation. After the definition of a reduced interval on ms.~12 we prove a sharp version of the stronger rigidity assertion: all endpoint-fixing endofunctors are identities exactly for reduced intervals.
\end{dcomment}


\msp{page 6}
NB Si $i$ n'est pas un bout, et si néanmoins l'une des deux catégories $i\backslash I$, $I/i$ est ponctuelle, disons $I/i$, alors au voisinage immédiat de $i$, $I$ a la forme
\[
\begin{tikzcd}[row sep=0pt]
\bullet & \bullet \arrow[l, "\underline{u}_i"'] \arrow[r, "\underline{u}_{i+1}"] & \bullet \\
i-1 & i & i+1
\end{tikzcd}
\]
($\underline{u}_i$ rétrograde, $\underline{u}_{i+1}$ directe - si c'était $i\backslash I$ la ponctuelle, alors $\underline{u}_i$ serait directe, $\underline{u}_{i+1}$ rétrograde - en tous les cas ces deux flèches sont de sens opposé, i.e. non composables). On trouve donc que $i\backslash I$ est une catégorie de la forme

% ---- from Der-ChVII-456_p006-010.tex
\printedpage{4}
% head: A. Grothendieck : Dérivateurs, Ch. VII, Édition des Universités de Montpellier II et Paris VII
% foot: Version 22 décembre 2001

\[
\begin{tikzcd}[column sep=small, row sep=small]
 & & & i \arrow[dl] \arrow[dr] & & & \\
 & & i-1 & & i+1 & & \\
 & \iddots \arrow[dl] & & & & \ddots \arrow[dr] & \\
 i-r & & & & & & i+r, % [sic] i+r -> i+s, s est la longueur de la chaîne contenant u_{i+1}
\end{tikzcd}
\]

où $r$ est la longueur de la chaîne structurale contenant $\underline{u}_i$, $s$ celle de la chaîne structurale contenant $\underline{u}_{i+1}$. Cette catégorie est une somme amalgamée $\bD_r \amalg \bD_s$ (recollées suivant leur origine), mais n'est pas isomorphe à une $\bD_j$, p.~ex. faute d'avoir un objet final.

Cette observation prouve a). Pour prouver b), il suffit de prouver que l'ordre total sur $I$ est déterminé de façon intrinsèque, par le choix de \textit{l'origine} $s_0$ de $I$ (parmi ces % [sic] ces -> ses, possessif
deux extrémités). Supposons que l'on ait déjà déterminé intrinsèquement les sommets $s_0$, $s_1$, $\ldots$, $s_i$ ($i \geq 0$). Comment déterminer $s_{i+1}$? (Quand $i < n$.) Soit $I'$ la sous-catégorie pleine de $I$
\msp{page 7}
formée des $j \in I$ qui ne sont pas parmi les $s_j$ \texttt{[plutôt} $s_k$\texttt{]} déjà construits, sauf le dernier $s_i$ (donc $s_i \in I'$, $\operatorname{card} I' \geq 2$). Alors $I'$ est isomorphe à un intervalle standard, avec $s_i$ correspondant à l'origine. Une des catégories $i\backslash I$, $I/i$ est ponctuelle, l'autre isomorphe à une catégorie $\bD_j$, et de plus, $i$ est objet initial ou final de $\bD_j$. Si c'est l'objet initial $0$, on prend pour $s_{i+1}$ le successeur immédiat de $i$ dans $\bD_j$. Si c'est l'objet final $j$, on prend pour $s_{i+1}$ le prédécesseur immédiat $j-1$.

Si on était parti du but $n$ de $I$, au lieu de partir de la source, on aurait trouvé la relation d'ordre total opposée. Affectant l'élément but du numéro d'ordre $0$, on trouve une renumération de $I$, $i$ devenant $n-i$, et les flèches de transition qui étaient directes pour la relation d'ordre initiale, deviennent rétrogrades dans la relation d'ordre opposée, et inversement. On voit ainsi :

\textbf{Proposition 2.} \textit{Si $I$ est un intervalle standard, il existe un unique couple $(I',\theta)$, où $I'$ est un intervalle standard, et $\theta$ un isomorphisme $I \xrightarrow{\sim} I'$ qui transforme le but de $I$ dans la source de $I'$ (et la source de $I$ en le but de $I'$).}

Cet intervalle $I'$ est appelé l'intervalle \textit{opposé} de $I$. Il faut le voir, intrinsèquement, comme la `même' categorie $I$, % [sic] categorie -> catégorie, accent manquant
mais avec une `orientation' opposée.

\msp{page 8}

Une façon plus claire de voir ceci est la suivante. Soit $J$ une catégorie intervalle, i.e. isomorphe à un intervalle standard. Alors il y a exactement (si $I \not\simeq e$) \textit{deux} couples $(I,\theta)$ d'un intervalle standard $I$ et d'un isomorphisme $\theta : I \xrightarrow{\sim} J$. Ces deux isomorphismes induisent sur $J$ deux ordres totaux, opposés l'un de l'autre. Chaque couple est déterminé de façon unique par ledit ordre total, qu'on appelera % [sic] appelera -> appellera, orthographe
une \textit{orientation} de l'intervalle $J$. Ainsi $J$ a deux orientations, et chacune détermine un \textit{isomorphisme orienté} (i.e. respectant
\printedpage{5}
% head: A. Grothendieck : Dérivateurs, Ch. VII, Édition des Universités de Montpellier II et Paris VII
% foot: Version 22 décembre 2001
les orientations) de $J$ avec un intervalle standard (\footnote[2]{\textbf{NB} On parlera de l'intervalle \textit{orienté opposé} $I^{\mathrm{opp}}$ d'un intervalle \textit{orienté} $I$, qu'il faut se garder de confondre avec la catégorie opposée $I^o$. Ici $I^{\mathrm{opp}}$ a même catégorie sous-jacente que $I$, mais avec une orientation opposée.}). Ces deux intervalles sont les mêmes si et seulement s'il existe un automorphisme de $J$ (nécessairement unique) qui échange ses deux orientations, ou ce qui revient au même, qui échange entre elles les deux extrémités. (Ainsi les catégories opposées l'une de l'autre
\[
\begin{tikzcd}[column sep=small, row sep=small]
 & \bullet \arrow[dl] \arrow[dr] & & & & \bullet & \\
 \bullet & & \bullet & & \bullet \arrow[ur] & & \bullet \arrow[ul]
\end{tikzcd}
\]
sont chacunes % [sic] chacunes -> chacune, orthographe
des intervalles symétriques.) Choisir une orientation revient aussi `a % [sic] `a -> à, accent imprimé comme guillemet
choisir l'origine parmi les deux extrémités.

Notons que la \textit{catégorie opposée} à un intervalle est un intervalle, et la catégorie opposée à un intervalle standard est l'intervalle standard qui s'en déduit en changeant tous les signes sur les intervalles de transition $[i-1,i]$ ($1 \leq i \leq n$). On se gardera de confondre la propriété $I \simeq I^o$, avec la propriété de symétrie.

\msp{page 9}

Ainsi les deux catégories ($\Psi_1$ et $\Phi_1$) ci-dessus sont symétriques, mais opposées l'une de l'autre sans être isomorphes. D'autre part, les $\bD_i$ ne sont pas symétriques, mais chacune isomorphe à la catégorie opposée.

Si $I$ est isomorphe à la catégorie opposée $I^o$, et si $\theta$ est un tel isomorphisme (nécessairement unique), alors il est immédiat que $\theta$ renverse l'orientation (sinon ce serait l'identité sur les sommets, et les flèches de transition seraient à la fois directes et indirectes), donc $\theta(0) = n$, $\theta(1) = n-1$, et on veut que $\underline{u}_1$ et $\underline{u}_n$ sont de même signe. Mais si $I$ est symétrique, $\underline{u}_1$ et $\underline{u}_n$ sont de signe opposé. Donc $I$ ne peut être à la fois symétrique et isomorphe à l'intervalle opposé \texttt{[plutôt à la catégorie opposée; exception :} $I = \{0\}$\texttt{]}.

Pour toute catégorie $\bD_i$, celle-ci étant un intervalle, il y a (si $i > 0$) \textit{deux orientations.} En tant qu'intervalles orientés, il y a $\bD_i^+$ (avec l'orientation ordinaire, donné par l'ordre ordinaire sur $\operatorname{Ob} \bD_i = [0,i]$), et $\bD_i^-$ avec l'orientation opposée.

\printedpage{6}
% head: A. Grothendieck : Dérivateurs, Ch. VII, Édition des Universités de Montpellier II et Paris VII
% foot: Version 22 décembre 2001

\msp{page 10}

\dgrblock{dgrb281}{Reversal and categorical duality in a sign word}
\begin{dcomment}{examples}{p9}{Reversal and categorical duality in a sign word}
Grothendieck's two uses of ``opposite'' describe operations which can be written down independently. This makes the examples of symmetry and self-duality entirely explicit. For a sign word $\tau=(\varepsilon_1,\ldots,\varepsilon_n)$, put
\[
\begin{aligned}
 R\tau&=(-\varepsilon_n,\ldots,-\varepsilon_1),\\
 D\tau&=(-\varepsilon_1,\ldots,-\varepsilon_n).
\end{aligned}
\]
The first is the standard model of the same category with its external orientation reversed; the standardizing isomorphism sends $i$ to $n-i$. The second is the categorical opposite with the original numerical orientation. The two involutions commute, and $RD\tau$ is the word read backwards without changing its signs.

For $n>0$, the interval is symmetric precisely when
\[
 \varepsilon_i=-\varepsilon_{n+1-i}\qquad(1\leq i\leq n).
\]
Necessity follows because a nonidentity automorphism must be reversal, by the preceding proof; sufficiency follows by checking each transition. This condition forces $n$ to be even. Likewise $I\cong I^o$ precisely when
\[
 \varepsilon_i=\varepsilon_{n+1-i}\qquad(1\leq i\leq n).
\]
The isomorphism must reverse the numerical order: an order-preserving bijection of the vertex line would be the identity and would have to reverse a nonidentity transition without changing its endpoints. These two sign conditions cannot hold together for a nonempty word.

For instance, $(+,-)$ and $(-,+)$ satisfy the first condition. They give the cospan and span, each symmetric, but the two categories are not isomorphic: one has a terminal object and two minimal objects, the other an initial object and two maximal objects. The constant word $(+,+,+)$ satisfies the second condition and gives the self-dual category $\Delta_3$ with no nonidentity automorphism. A less uniform self-dual example is $(+,-,+)$; its middle sink and source are exchanged under the anti-isomorphism. For a reduced alternating word, symmetry holds at positive even length and self-duality at odd length.

The point is separate: it has a unique orientation and only the identity automorphism, and both operations leave it unchanged. Keeping this case explicit avoids construing its coincident endpoints as two different choices.
\end{dcomment}


\section*{2\quad Sous-catégories pleines d'un intervalle}

Soit $I$ une catégorie intervalle orientée, de longueur $n \geq 1$. Nous identifierons $\operatorname{Ob} I$ à l'ensemble des entiers $[0,n]$. Pour une chaîne $\underline{u}_i$, $\underline{u}_{i+1}$, $\ldots$, $\underline{u}_{i+l-1}$ de flèches de transition dans $I$, on appelle \textit{domaine} de la chaîne, l'ensemble des sources et de buts de ces éléments. C'est donc un intervalle d'entiers $[i-1, i+l-1]$, de longueur $l$ égale à celle de la chaîne, dont l'origine est déterminé % [sic] déterminé -> déterminée, accord avec l'origine
par la première flèche $\underline{u}_i$ de la chaîne (c'est sa source, si $\underline{u}_i$ est directe, et son but si $\underline{u}_i$ est rétrograde), et dont l'extrémité est déterminée par la dernière $\underline{u}_{i+l-1}$. Les intervalles d'entiers (vus comme des sommets de $I$) correspondant aux chaînes structurales sont appelés les \textit{intervalles structuraux} ou \textit{intervalles critiques} de $I$. Contrairement aux chaînes structurales, ces intervalles ne forment pas une partition de $I$, puisque deux intervalles structuraux peuvent avoir une extrémité commune. Si les chaînes structurales sont
\[
C_1,\ C_2,\ \ldots\ ,\ C_r \qquad (r = \text{ longueur utile de } I)
\]
de longueurs
\[
l_1,\ l_2,\ \ldots\ ,\ l_r,
\]
alors les intervalles structuraux sont
\[
\underbrace{[0,l_1]}_{I_1},\ \underbrace{[l_1,l_1+l_2]}_{I_2},\ \ldots\ ,\ \underbrace{[l_1+\cdots+l_{r-1},l_1+\cdots+l_{r-1}+l_r]}_{I_r},
\]
une suite d'intervalles d'entiers dont deux consécutifs ont une extrémité commune, et tels que si $i < j$, chaque élément de $I$ % [sic] I -> I_i, comparaison des éléments de I_i et de I_j
soit $\leq$ chaque élément de $I_j$. On appelle \textit{sommet critique} de $I$ les extrémités des intervalles structuraux.

\textbf{Proposition 3.} \textit{Soit $s \in I$, soit $I' = I/s$ (sous-catégorie pleine de $I$ des $s'$ tels qu'il existe $s' \longrightarrow s$), $I'' = I/s$ \textup{\texttt{[plutôt} $I'' = s\backslash I$\texttt{]}}, $I'_0 = (I/s)\backslash\{s\}$, $I''_0 = (s\backslash I)\backslash\{s\}$} (\footnote[3]{i.e. qu'on ait : il n'y a pas de flèches entre un élément de $I'_0$ et un élément de $I''_0$.}). Pour que $s$ soit critique, il faut et il suffit que $I\backslash\{s\}$ soit catégorie somme de $I'_0$ et $I''_0$.

\msp{page 11}

\dgrblock{dgrb282}{Critical vertices and the correct deletion test}
\begin{dcomment}{exposition}{p10}{Critical vertices and the correct deletion test}
Grothendieck wants to retain exactly the vertices which a path cannot bypass by composing arrows. These are the endpoints of the structural intervals. If their positions are
\[
 0=c_0<c_1<\cdots<c_r=n,
\]
then $[c_{a-1},c_a]$ is one constant-sign run, and each interior $c_a$ is a change of sign. In the categorical partial order, the critical vertices are exactly the minimal and maximal vertices. This observation explains the two slices in Proposition~3, but it also exposes an error in its global formulation.

Let
\[
 A=(I/s)\setminus\{s\},\qquad
 B=(s\backslash I)\setminus\{s\}.
\]
These are respectively the strict predecessors and strict successors of $s$ in the categorical order. They do not in general exhaust the other vertices. In $0\to1\leftarrow2$, take $s=0$. It is critical; $A$ is empty and $B$ consists of $1$, whereas $I\setminus\{0\}$ still contains $2$. Thus the printed equality with all of $I\setminus\{s\}$ cannot be right. The edition prints the same reading.

The footnote gives the valid local test: there is no arrow between $A$ and $B$. If $s$ is critical, it is minimal or maximal and one of these categories is empty. If $s$ is not critical, its two adjacent transitions have the same sign. One adjacent vertex is a strict predecessor and the other a strict successor; their composite through $s$ is an arrow between $A$ and $B$. Equivalently, the full subcategory on the vertices comparable with $s$, after removing $s$, is the coproduct $A\amalg B$ exactly when $s$ is critical.

For the application to subintervals on the next page, the appropriate \emph{global} test uses the numerical order. Suppose $0<s<n$, and let $L$ and $H$ be the full subcategories on $[0,s-1]$ and $[s+1,n]$. Then
\[
 I\setminus\{s\}=L\amalg H
 \quad\Longleftrightarrow\quad s\text{ is critical}.
\]
If $s$ is a turning vertex, an arrow from one side to the other would traverse two transitions of opposite signs, which is impossible. If it is not a turning vertex, the two neighbouring vertices are connected by their composite through $s$, an arrow still present in the full subcategory obtained by deleting $s$. This corrected statement proves the separation property used in Proposition~5. At a numerical endpoint, deleting it leaves a connected interval; its criticality is already accounted for by its being an endpoint.
\end{dcomment}


\textbf{Proposition 4.} \textit{Soient $i,j$ dans $I$, avec $i < j$. Conditions équivalentes :}
\begin{enumerate}
\item[\textit{a)}] \textit{Il existe une flèche entre $i$ et $j$ (soit $i \longrightarrow j$, soit $j \longrightarrow i$).}
\item[\textit{b)}] \textit{$i$ et $j$ appartiennent à un même intervalle structural (nécessairement unique).}
\item[\textit{c)}] \textit{La sous-catégorie pleine de $I$ définie par l'ensemble d'objets $[i,j]$ est isomorphe à une $\bD_\alpha$ (on aura $\alpha = j-i$).}
\end{enumerate}

\printedpage{7}
% head: A. Grothendieck : Dérivateurs, Ch. VII, Édition des Universités de Montpellier II et Paris VII
% foot: Version 22 décembre 2001
On dira alors que $i$ et $j$ sont \textit{directement connectés} dans $I$. (Et on utilisera la même terminologie dans une catégorie quelconque : $i$ et $j$ sont directement connectés si et seulement si $\operatorname{Hom}_I(i,j)$ \textit{ou} $\operatorname{Hom}_I(j,i)$ est non vide.)

\textbf{Proposition 5.} \textit{Soit $J \subseteq I$ une partie de $I$, considérons la sous-catégorie pleine de $I$ qu'elle définit. Alors :}
\begin{enumerate}
\item[\textit{a)}] \textit{Les composantes connexes $J_\alpha$ de $J$ sont des intervalles et l'ordre induit sur un $J_\alpha$ par celui de $I$ est une orientation (dite \textup{orientation induite}) de $J_\alpha$.}
\item[\textit{b)}] \textit{Supposons $\operatorname{card} J > 1$. Soit $i$ le plus petit élément de $J$, $j$ le plus grand élément, donc $i < j$. Pour que $J$ lui-même soit un intervalle, il faut et il suffit que tout sommet critique dans $[i,j]$ soit élément de $J$.}
\end{enumerate}
\nopagebreak[4]
(\footnote[4]{On appelle \textit{sous-intervalle} de $I$ une sous-categorie % [sic] sous-categorie -> sous-catégorie, accent manquant
pleine connexe, i.e. qui est elle-même un intervalle. Notons d'ailleurs que \textit{toute} sous-catégorie $I'$ de $I$ qui est un intervalle, est nécessairement pleine, cf. plus bas \texttt{[paragraphe]} des foncteurs entre intervalles.}).

\textbf{NB} Si $s \in [i,j]$ est un sommet critique de $I$ qui n'est pas dans $J$, alors par la proposition 3 il disconnecterait $J$, i.e. $J$ serait somme directe de deux sous-catégories $J'$ et $J''$ non vides ($i \in J'$, $j \in J''$), donc $J$ ne serait pas un intervalle.

\msp{page 12}

\dgrblock{dgrb283}{Connected full subcategories and the vertices they must retain}
\begin{dcomment}{examples}{p11}{Connected full subcategories and the vertices they must retain}
Grothendieck is preparing a reduction which composes consecutive arrows without identifying a change of direction. The distinction between a full subcategory and an arbitrary subcategory matters here. A full subcategory retains every composite whose endpoints remain, even when an intermediate vertex has been removed. Thus numerical gaps do not necessarily disconnect it.

Let $J=\{j_0<\cdots<j_m\}$ be a nonempty subset of the vertices of $I$, with its full induced category. Two successive chosen vertices $j_a,j_{a+1}$ are joined by an arrow if and only if the signs between them are constant. By Proposition~4 this is equivalent to there being no critical vertex strictly between them. Whenever this condition fails, there can be no arrow from any chosen vertex at or before $j_a$ to any chosen vertex at or after $j_{a+1}$, in either direction: such an arrow would cross the same change of sign. The connected components are therefore the maximal consecutive blocks of the list $(j_0,\ldots,j_m)$ in which successive entries are directly connected.

Within one such block, make a generating transition between each successive pair, in the direction of its arrow in $I$. The free category on this oriented line is exactly the induced category. To verify fullness, an arrow between nonconsecutive retained vertices exists in $I$ precisely when the signs on the whole intervening segment agree. That is also precisely the condition that the successive arrows of the block compose. This proves Proposition~5(a), including the claim about the induced orientation. It proves (b) as well: the whole chosen list is connected exactly when every ambient critical vertex between its extreme vertices was retained.

For a concrete calculation take
\[
 I:\quad0\longrightarrow1\longrightarrow2
       \longleftarrow3\longleftarrow4\longrightarrow5.
\]
Its critical vertices are $0,2,4,5$. The full category on $\{0,2,4,5\}$ is the reduced interval
\[
 0\longrightarrow2\longleftarrow4\longrightarrow5.
\]
Deleting $1$ or $3$ has merely replaced a pair of arrows by its composite. The full category on $\{0,1,4,5\}$ has two components, $\{0,1\}$ and $\{4,5\}$, because the omitted turning vertex $2$ separates them. The least subinterval containing $1$ and $5$ is $\{1,2,4,5\}$: the endpoints must remain, the two intervening turns must remain, and the preceding calculation shows that nothing else is required. This is the corollary on ms.~12.

The last sentence of footnote~4 asserts more than fullness permits: an interval subcategory need not be full. Inside $\Delta_2$, keep all three objects and the arrows $0\to1$ and $0\to2$, but omit $1\to2$. This is a subcategory, because its two nonidentity arrows are not composable. Abstractly it is the span interval with vertex order $1,0,2$. It is not full in $\Delta_2$. Thus fullness must remain part of the definition of a subinterval; it cannot be inferred solely from the abstract shape of the subcategory.

Finally, for a full subinterval with at least two vertices, an interior retained vertex is a turn in the induced shape exactly when it was an ambient critical vertex. This proves Proposition~6. The reduced interval is consequently the smallest full connected part of $I$ which contains its endpoints, and its length is the useful length of $I$.
\end{dcomment}


\textbf{Corollaire.} \textit{L'ensemble des sommets critiques de $I$ est aussi l'ensemble des sommets du plus petit sous-intervalle de $I$ contenant les extrémités. Plus généralement, si $i < j$ sont deux sommets de $I$, le plus petit sous-intervalle de $I$ contenant $i$, $j$ est formé de $i$, $j$ et des sommets critiques $s$ tels que $i < s < j$.}

\textbf{Définition.} Un intervalle est dit \textit{réduit} (ou \textit{irrédondant}) si tous ses sommets sont critiques, ou ce qui revient au même, si ses chaînes structurales sont de longueur 1.

Un intervalle orienté \textit{réduit} est connu, à isomorphisme près, quand on connaît
\begin{enumerate}
\item[a)] sa longueur $\mathrm{n}$,
\item[b)] sa direction initiale $+$ (directe) ou $-$ (rétrograde).
\end{enumerate}

\textbf{NB} Cette direction est la même que la direction finale si $n$ est impair, elle est opposée si $n$ est pair.

Ainsi, suivant les deux cas de direction initiale, on trouve
\[
\begin{aligned}
0 \longrightarrow 1 \longleftarrow 2 \longrightarrow \cdots \qquad\qquad \cdots\ n & \qquad (\text{Avec } & n-1 \longrightarrow n & \quad \text{si } n \text{ impair,}\\
& & n-1 \longleftarrow n & \quad \text{si } n \text{ pair.)}\\
0 \longleftarrow 1 \longrightarrow 2 \longleftarrow \cdots \qquad\qquad \cdots\ n & \qquad (\text{Avec } & n-1 \longleftarrow n & \quad \text{si } n \text{ impair,}\\
& & n-1 \longrightarrow n & \quad \text{si } n \text{ pair.)}
\end{aligned}
\]

\printedpage{8}
% head: A. Grothendieck : Dérivateurs, Ch. VII, Édition des Universités de Montpellier II et Paris VII
% foot: Version 22 décembre 2001

\textbf{Proposition 6.} \textit{Un sous-intervalle $J$, ayant $i$ et $j$ ($i < j$) comme plus petit et plus grand sommet, est réduit si et seulement si tout $s \in J$ distinct de $i$ et de $j$ est critique.}

\textbf{Corollaire.} \textit{La sous-catégorie pleine de $I$ formée des sommets critiques est un sous-intervalle. On l'appelle l'intervalle réduit associé à $I$, noté $I_{\textup{réd}}$.}

% ---- from Der-ChVII-456_p011-015.tex
\printedpage{9}
% head: A. Grothendieck : Dérivateurs, Ch. VII, Édition des Universités de Montpellier II et Paris VII
% foot: Version 22 décembre 2001

\msp{page 13}

\dgrblock{dgrb284}{The critical-vertex inclusion has two adjoints}
\begin{dcomment}{development}{p12}{The critical-vertex inclusion has two adjoints}
The reduced interval retains every turn, but Grothendieck has not yet described how to return to it from an arbitrary vertex. There are two particularly useful choices, one following the categorical source of a run and one following its categorical target. They relate reduction to the endpoint-preserving transformations which a path calculus needs.

\textbf{Two canonical collapses.}\ours{}
Let $j:R=I_{\mathrm{red}}\hookrightarrow I$ be the full inclusion of the critical vertices of an oriented interval. There are retractions $r_s,r_t:I\to R$ with
\[
 r_t\dashv j\dashv r_s,\qquad
 r_sj=r_tj=\mathrm{id}_R.
\]
Both retractions fix the endpoints. At a noncritical vertex $x$, $r_s(x)$ is the categorical source of its structural interval and $r_t(x)$ is that interval's categorical target. At a critical vertex both are the identity. For the point we set $R=I$ and take all three functors to be the identity.

\emph{Proof.} A noncritical vertex lies strictly inside a unique structural interval. On a positive run $[a,b]$, the source-collapse sends every interior vertex to $a$, leaves $b$ at $b$, and is a monotone map of its categorical chain. On a negative run the source is $b$ and the same argument applies with the categorical order reversed. The target-collapse is obtained by exchanging source and target. The run-wise maps agree at shared critical vertices. Every arrow belongs to one run, so these formulas give functors on the whole interval.

To verify the adjunctions rather than merely exhibit homotopies, let $q$ be a critical vertex. If $x$ is itself critical, fullness of $j$ gives the required hom-set bijections. If $x$ is interior to a run, the only critical vertices comparable with $x$ are its source $a$ and target $b$, with $a\to x\to b$ in the categorical order. Indeed an arrow extending beyond either endpoint of that run would have to pass a change of sign. Consequently
\[
\begin{split}
 I(jq,x)\ne\varnothing&\ \Longleftrightarrow\ q=a
     \ \Longleftrightarrow\ R(q,r_sx)\ne\varnothing,\\
 I(x,jq)\ne\varnothing&\ \Longleftrightarrow\ q=b
     \ \Longleftrightarrow\ R(r_tx,q)\ne\varnothing.
\end{split}
\]
For the middle equivalences we use that a source endpoint is a minimal vertex of $I$ and a target endpoint a maximal vertex; hence the same holds in $R$. All these hom-sets are empty or singletons, so the bijections are unique and natural. They establish $j\dashv r_s$ and $r_t\dashv j$. The two retraction identities follow directly from the definition. The counit and unit are the unique arrows
\[
 j r_s(x)\longrightarrow x,\qquad x\longrightarrow j r_t(x),
\]
which are identities whenever $x$ is critical. This also verifies the triangle identities directly.\hfill$\square$

These natural transformations are categorical deformations relative to all critical vertices, in particular relative to both endpoints. Applying the nerve, a natural transformation $F\Rightarrow G$ supplies a simplicial homotopy through the functor $I\times\Delta_1\to I$ which it defines. Thus $Nj$ has homotopy inverses $Nr_s$ and $Nr_t$. Nothing in this argument asserts that the interval categories themselves are equivalent: a skeletal interval with an interior noncritical vertex has strictly more objects than its reduction.

The effect on diagrams gives a more informative comparison. For every category $X$, precomposition yields
\[
 r_s^*\dashv j^*\dashv r_t^*,\qquad
 j^*r_s^*=j^*r_t^*=\mathrm{id}_{\operatorname{Fun}(R,X)}.
\]
This reversal of adjunctions under precomposition is the standard calculation of \citet[Proposition~4.3.7]{Riehl2016}; it can also be read directly from the two displayed natural transformations by applying each diagram $I\to X$. The source extension inserts identity arrows before the final arrow of each directed run, and the target extension inserts them after its first arrow. Restriction to $R$ composes each run.

For $I=\Delta_2$, $R=\{0,2\}$ and $X=\Delta_1$, consider the two diagrams with vertex values
\[
 F=(0,0,1),\qquad G=(0,1,1).
\]
They have the same restriction $(0,1)$ to $R$. There is no natural transformation $G\to F$, since its middle component would have to be $1\to0$ in $\Delta_1$, whereas their restrictions have an identity map. Hence $j^*$ is not full and is not an equivalence. If $X$ is a groupoid, every component of both comparison transformations is invertible, and $j^*$ is an equivalence, with either section as an inverse. Reduction therefore loses factorization data in general, but loses no such data when all arrows of the target are already invertible.

The general precomposition-adjunction statement is known and cited above. The critical-vertex reduction itself also appears in \citet[Lemma~4.16]{ArakawaCnossen2026}. For the additional explicit adjoints and the sharp rigidity criterion below, we checked that construction, the local interval arguments here, and the adjunction calculus in \citet[\S4.3]{Riehl2016}; targeted searches for critical-vertex reductions and their adjoints did not locate these particular formulations. We supply the extra deductions as elementary developments of the present definitions, without a priority claim.
\end{dcomment}


\dgrblock{dgrb285}{The exact scope of endpoint rigidity}
\begin{dcomment}{development}{p12}{The exact scope of endpoint rigidity}
The false generality of Proposition~1(b) suggests asking precisely which part of an interval allows an endpoint-preserving functor to move a vertex. The preceding collapses give one implication, while a finite displacement argument gives the other.

\textbf{Endpoint rigidity criterion.}\ours{}
For a finite oriented interval $I$, the following are equivalent: $I$ is reduced; every endofunctor of $I$ fixing its two designated endpoints is the identity. We include the point as a reduced interval.

\emph{Proof.} Suppose $I$ is reduced, of length $n>0$, and identify its vertices with $[0,n]$. A nonidentity arrow of $I$ is then a single transition: two consecutive transitions have opposite signs and cannot compose. If $f:I\to I$ is any functor, the image of a transition is therefore either an identity or a transition. Its object function satisfies
\[
 -1\leq f(i)-f(i-1)\leq1\qquad(1\leq i\leq n).
\]
If $f(0)=0$ and $f(n)=n$, summation gives
\[
 n=f(n)-f(0)=\sum_{i=1}^n\bigl(f(i)-f(i-1)\bigr)\leq n.
\]
Equality forces every summand to be $1$. Thus $f(i)=i$ for all $i$, and thinness forces $f$ to be the identity on arrows as well. We have made no monotonicity assumption on $f$; it is a consequence of the equality. The point case is immediate.

Conversely, if $I$ is not reduced, a structural interval has an interior vertex. The endofunctor $j r_s$ of the preceding result fixes every critical vertex and hence both endpoints, but sends that interior vertex to the source of its run. It is not the identity.\hfill$\square$

The obstruction is therefore exactly a nontrivial factorization inside a constant-sign run. Automorphisms cannot change such a factorization because they preserve the graph of indecomposable arrows. General functors can replace it by an identity followed by a longer composite. This is why the automorphism proof on ms.~6--7 is sound while its printed extension to all endofunctors fails.
\end{dcomment}


\section*{3\quad Chemins. Catégorie de chemins d'un type donné $\tau$.}

Soient $I$ un intervalle orienté, et $X$ une catégorie. On appelle \textit{chemin de type $I$ dans $X$} tout foncteur $I \longrightarrow X$, et \textit{catégorie de chemins de type $I$ dans $X$} la catégorie
\[
\Ch_I(X) = \uHom(I,X).
\]
Comme $I$ est \textit{canoniquement} isomorphe à une catégorie intervalle standard $I_0$, on peut identifier les chemins de type $I$, au de type $I_0$. Le type $\tau$ d'un chemin le longueur $n$ $(= \operatorname{card} I)$ dans $X$ peut donc être décrit par une suite de $n$ signes $+$ ou $-$ % [sic] au de type -> à ceux de type, word(s) missing; [sic] le longueur -> de longueur, wrong preposition
\[
\tau_1,\ \tau_2,\ \dots\ ,\ \tau_n \qquad (\tau_i \in \{-1,+1\} \text{ pour } 1 \leq i \leq n).
\]
À cause du fait que $I$ est engendré pas ses flèches de transition, un foncteur $c : I \longrightarrow X$ est déterminé quand on connaît % [sic] pas ses flèches -> par ses flèches, misprint

\begin{enumerate}
\item[a)] Les $c(i)$ $(i \in I)$, i.e. une suite de $n+1$ objets de $X$ :
\[
a_0,\ a_1,\ \dots\ ,\ a_n\ (\text{\footnotemark[5]});
\]
\footnotetext[5]{$a_0$ est dit la \textit{source} du chemin, $a_n$ son but.}

\item[b)] Les $c(\underline{u}_i)$ (où les $\underline{u}_i$ sont les flèches de transition dans $I$), i.e. pour tout $i \in [1,n]$ une flèche $u_i$ dans $X$ entre $a_{i-1}$ et $a_i$, allant de $a_{i-1}$ dans $a_i$ si $\underline{u}_i$ est directe $(\tau_i = +1)$, et de $a_i$ dans $a_{i-1}$ si $\underline{u}_i$ est rétrograde $(\tau_i = -1)$,
\end{enumerate}

et à toute telle donnée a), b) correspond un foncteur. Ainsi on pourra identifier les chemins de type $I$ \msp{page 14} dans $X$ aux diagrammes du type précisé dans $X$. Et les morphismes entre chemins correspondent de même aux morphismes entre diagrammes $a$ et $a'$, une flèche $a \xrightarrow{\alpha} a'$ étant donc une suite de $n+1$ flèches
\[
\alpha_i : a_i \longrightarrow a'_i, \qquad i \in [0,n],
\]
satisfaisant la condition de commutativité pour les $n$ carrés
\[
\begin{tikzcd}
a_{i-1} \arrow[r, "u_i"] \arrow[d, "\alpha_{i-1}"'] & a_i \arrow[d, "\alpha_i"] \\
a'_{i-1} \arrow[r, "u'_i"] & a'_i
\end{tikzcd}
\qquad\qquad \text{(cas $\underline{u}_i$ directe)}
\]

\printedpage{10}
% head: A. Grothendieck : Dérivateurs, Ch. VII, Édition des Universités de Montpellier II et Paris VII
% foot: Version 22 décembre 2001
ou
\[
\begin{tikzcd}
a_{i-1} \arrow[d, "\alpha_{i-1}"'] & a_i \arrow[l, "u_i"'] \arrow[d, "\alpha_i"] \\
a'_{i-1} & a'_i. \arrow[l, "u'_i"']
\end{tikzcd}
\qquad\qquad \text{(cas $\underline{u}_i$ rétrograde)}
\]
Si on a une flèche entre deux sommets quelconque $i$, $j$ dans $I$, avec $i < j$ et $j$ non successeur immédiat de $i$, alors $i$, $j$ appartiennent à un même intervalle critique, et dans $I$ on aura % [sic] quelconque -> quelconques, agreement
\[
i \xrightarrow{\underline{u}_{i+1}} i+1 \xrightarrow{\underline{u}_{i+2}}\ \cdots\ j-1 \xrightarrow{\underline{u}_j} j \qquad \text{une chaîne de longueur } j-i \geq 2 % doubt resolved: PDF p. 12 (edition p. 10) clearly underlines u; the existing reading is retained.
\]
ou
\[
i \longleftarrow i+1 \longleftarrow\ \ \cdots\ j-1 \longleftarrow j,
\]
et la flèche correspondante entre $a_i$ et $a_j$ s'obtient en composant les flèches $u_i$ correspondantes.

\msp{page 15}

\dgrblock{dgrb286}{A fixed path type is a freely generated diagram}
\begin{dcomment}{exposition}{p14}{A fixed path type is a freely generated diagram}
Grothendieck now uses the work on intervals to make path data and path morphisms completely explicit. Fixing the type does not impose any compatibility among the arrows of the path beyond their sources and targets. This is the benefit of having chosen a free path category as the domain.

For example, let the type be $(+,+,-)$. A diagram consists of
\[
 a_0\xrightarrow{u_1}a_1\xrightarrow{u_2}a_2
       \xleftarrow{u_3}a_3.
\]
The only nonidentity arrow of the domain which is not a transition is $0\to2$, and its image must be $u_2u_1$. Given a second diagram $a'_0\to a'_1\to a'_2\leftarrow a'_3$, a map between them consists of four arrows $\alpha_i:a_i\to a'_i$ satisfying
\[
 \alpha_1u_1=u'_1\alpha_0,\qquad
 \alpha_2u_2=u'_2\alpha_1,\qquad
 \alpha_2u_3=u'_3\alpha_3.
\]
Naturality for the longer arrow follows by calculation:
\[
 \alpha_2(u_2u_1)
 =(u'_2\alpha_1)u_1
 =u'_2(u'_1\alpha_0)
 =(u'_2u'_1)\alpha_0.
\]
Every longer constant-sign run is handled by the same induction, with the index order reversed when its arrows are retrograde. Thus the displayed transition squares are both necessary and sufficient for a natural transformation. Identities and composition are taken vertex by vertex; composing commutative squares verifies that the result is again a morphism of paths.

This also gives a useful exact model of the fixed-type category. Let $\operatorname{Arr}(X)=\operatorname{Fun}(\Delta_1,X)$, with ordinary domain and codomain functors $d,c$. For a positive transition its ordered endpoint pair is $(d,c)$; for a negative one it is $(c,d)$. Then $\operatorname{Fun}(I_\tau,X)$ is the iterated strict fibre product of $n$ copies of $\operatorname{Arr}(X)$, matching the second ordered endpoint of one to the first ordered endpoint of the next. On objects this identifies the shared vertices. On arrows it identifies the corresponding vertical components of the naturality squares. The two descriptions are therefore isomorphic as categories, not merely bijective on objects. At length zero the category is $X$ itself.

The construction retains all factorizations and all commuting ladders between them. It has not inverted the retrograde arrows in $X$ and has not identified paths related by a ladder. Those operations would require additional choices or a localization. Here a retrograde transition simply says that the arrow of $X$ points against the chosen order of traversal. Keeping this distinction explicit will matter when the category of types starts to vary.
\end{dcomment}


\textbf{Exemples.}

\begin{enumerate}
\item[1.] $I = \bDelta_0 = e$, $\Ch_I(X) \simeq X$. Les chemins de longueur 0 s'identifient aux objets $a$ de $X$. L'unique élément $\mathrm{id}_e$ de $\mathrm{Fl}(I)$ donne $\mathrm{id}_a$ dans $X$.

\item[2.] $I = \bDelta_1^+$, $\Ch_I(X) \simeq \uFl(X)$. Les chemins de type $\bDelta_1^+$ sont ceux de longueur 1 et \textit{directs,} i.e. donnés par une flèche
\[
a_0 \longrightarrow a_1
\]
dont la source est aussi la source $a_0$ du chemin. Par contre, les chemins de type $\bDelta_1^-$ sont donnés par des diagrammes
\[
a_0 \longleftarrow a_1,
\]
i.e. par des flèches dont le but est considéré comme source du chemin. Bien sûr, les catégories $\Ch_{\bDelta_1^+}(X)$ et $\Ch_{\bDelta_1^-}(X)$ sont canoniquement isomorphes, mais par un isomorphisme qui \textit{échange source et but du chemin.} Nous nous garderons donc bien de confondre les deux chemins précédents.

\msp{page 16}

\item[3.] $\bDelta_2^+$, $\bDelta_2^-$, correspondant aux diagrammes
\[
a_0 \longrightarrow a_1 \longrightarrow a_2
\]
resp.
\[
a_0 \longleftarrow a_1 \longleftarrow a_2,
\]
on peut faire une remarque similaire que pour $\bDelta_1^+$, $\bDelta_1^-$.

\printedpage{11}
% head: A. Grothendieck : Dérivateurs, Ch. VII, Édition des Universités de Montpellier II et Paris VII
% foot: Version 22 décembre 2001

\item[4.] $\Phi_1$, $\Psi_1$, ce sont respectivement les intervalles
\[
\Phi_1 :\ 0 \xleftarrow{\underline{u}_0} 1 \xrightarrow{\underline{u}_1} 2 \qquad\qquad \text{ou} \qquad\qquad
\begin{tikzcd}[row sep=small, column sep=small]
 & 0 \\
1 \arrow[ur] \arrow[dr] & \\
 & 2
\end{tikzcd}
\qquad , % doubt resolved: PDF p. 13 (edition p. 11) clearly underlines u in both labels; the existing reading is retained.
\]
et
\[
\Psi_1 :\ 0 \xrightarrow{\underline{u}_0} 1 \xleftarrow{\underline{u}_1} 2 \qquad\qquad \text{ou} \qquad\qquad
\begin{tikzcd}[row sep=small, column sep=small]
 & 0 \arrow[dl] \\
1 & \\
 & 2 \arrow[ul]
\end{tikzcd}
\qquad .
\]
Ils correspondent respectivement aux diagrammes
\[
a_0 \xleftarrow{u_0} a_1 \xrightarrow{u_1} a_2
\]
et
\[
a_0 \longrightarrow a_1 \longleftarrow a_2.
\]
Les catégories $\Phi_1$, $\Psi_1$ sont opposées l'une de l'autre, et ne sont pas isomorphes. Aussi n'y a-t-il en général aucun isomorphisme entre $\Ch_{\Phi_1}(X)$ et $\Ch_{\Psi_1}(X)$.
\end{enumerate}

\dgrblock{dgrb287}{Zigzags before taking a localization}
\begin{dcomment}{context}{p16:sont opposées l'une de l'autre}{Zigzags before taking a localization}
Grothendieck's paths belong to the same categorical language as fractions, while retaining information which a category of fractions has already identified. In \citet[Chapter~I, \S1.1]{GabrielZisman1967}, localization adjoins inverses to a specified collection of arrows and imposes the relations needed for functoriality and inversion. A span $x\leftarrow z\to y$ can consequently represent a morphism from $x$ to $y$ when its left arrow is among those to be inverted. The fixed-type path category written here makes no such selection: the objects and every component of a natural transformation are ordinary categorical data.

The two-dimensional comparison is equally concrete. The reduced hammocks of \citet[Definition~2.1]{DwyerKan1980Calculating} consist of commutative ladders of zigzags, but their vertical arrows and their left-pointing arrows are required to lie in a specified subcategory of weak equivalences; identity columns are removed and neighbouring columns of the same direction are composed. Their Proposition~3.1 identifies connected components of the resulting mapping spaces with ordinary localized morphisms. Grothendieck's full functor category of one fixed interval is therefore an ingredient of a comparable calculus, not by itself a hammock mapping space. The later commentary on the varying category of types explains the precise relation; here the historical comparison tells us why it is useful to retain ladders rather than only strings of arrows.
\end{dcomment}


Soit $I$ un intervalle orienté, et considérons l'\textit{intervalle opposé $I^{\opp}$}, ayant même catégorie sous-jacente mais l'orientation opposé. Comme la définition de la catégorie des chemins du type $I$ ne fait pas intervenir l'orientation de $I$, on a donc un isomorphisme canonique % [sic] l'orientation opposé -> l'orientation opposée, agreement
\[
\Ch_I(X) \xrightarrow{\sim} \Ch_{I^{\opp}}(X).
\]

\msp{page 17}

\dgrblock{dgrb288}{A span and a cospan give different diagram categories}
\begin{dcomment}{examples}{p16}{A span and a cospan give different diagram categories}
Grothendieck emphasizes that the direction of a turn changes the fixed-type category. The difference can be seen without any localization or homological algebra. We give a finite example which also distinguishes equivalence of categories, not only strict isomorphism.

Let $X$ be the three-object poset
\[
             a\longleftarrow b\longrightarrow c.
\]
A span diagram in $X$ is a triple $(x_0,x_1,x_2)$ with $x_1\leq_X x_0$ and $x_1\leq_X x_2$. If its middle value is $b$, each endpoint can independently be $a,b,$ or $c$, giving nine diagrams. If the middle value is $a$, both endpoints must be $a$; if it is $c$, both must be $c$. Thus
\[
 \bigl|\operatorname{Ob}\operatorname{Fun}(\Phi_1,X)\bigr|
 =3^2+1^2+1^2=11.
\]
A cospan diagram has $x_0\leq_X x_1$ and $x_2\leq_X x_1$. For middle value $b$ there is one possibility. For middle value $a$, both endpoints independently lie in $\{b,a\}$, giving four; for middle value $c$, both independently lie in $\{b,c\}$, again giving four. Hence
\[
 \bigl|\operatorname{Ob}\operatorname{Fun}(\Psi_1,X)\bigr|
 =1^2+2^2+2^2=9.
\]
Morphisms in both functor categories are completely determined: there is one from $(x_0,x_1,x_2)$ to $(y_0,y_1,y_2)$ exactly when $x_i\leq_Xy_i$ for each $i$. The naturality squares then commute automatically because $X$ is thin. These categories are posets and have no distinct isomorphic objects. An equivalence between posets is an isomorphism, since essential surjectivity is actual surjectivity and full faithfulness reflects their order. Their different object counts therefore rule out an equivalence.

Orientation reversal is a separate operation. For the span, interchanging $x_0$ and $x_2$ gives another span and an automorphism of its diagram category. It leaves the middle vertex and both arrow directions of the underlying span shape intact, while exchanging the two endpoint evaluations. Passing from the span to the cospan reverses the categorical arrows at the turn. The preceding count detects exactly that change.
\end{dcomment}


\dgrblock{dgrb289}{When a change of orientation becomes a stable equivalence}
\begin{dcomment}{later}{p16}{When a change of orientation becomes a stable equivalence}
The span--cospan distinction remains real for the ordinary functor categories which Grothendieck has just defined. A later theorem explains why analogous diagram categories can nevertheless become equivalent after one passes to a stable homotopy theory. For a finite oriented tree $Q$, \citet[Theorem~9.19 and Corollary~9.23, arXiv v3]{GrothStovicek2016Trees} construct reflection equivalences and deduce strong stable equivalence between any two orientations of the same underlying tree. In particular, for a stable derivator $\mathbb D$,
\[
 \mathbb D^{\Phi_1}\simeq\mathbb D^{\Psi_1}.
\]
These are equivalences of shifted derivators, natural with respect to exact morphisms. They are constructed using stable operations such as fibres, cofibres, and biproducts; they are not induced by an isomorphism of the span and cospan categories. The finite example above therefore presents no obstruction. It concerns arbitrary ordinary targets, whereas the reflection theorem uses stability to transform the values and arrows of the diagram. This gives a precise later setting in which changing the signs of an interval preserves a diagram theory even though it does not preserve its indexing category.
\end{dcomment}


Néanmoins on se gardera de confondre un chemin de type $I$ avec le chemin opposé de type $I^{\opp}$, la source de l'un étant le but de l'autre et inversement. C'est dire qu'un chemin doit être vu comme \textit{plus} qu'un simple foncteur
\[
I \longrightarrow X,
\]
il faut de plus se donner l'\textit{orientation de $I$} comme faisant partie de la structure de chemin. Donc la définition formelle raisonnable, qui colle à l'intuition du chemin comme allant d'une \textit{source} bien définie vers un \textit{but} également bien défini, c'est qu'un chemin (défini en termes d'une catégorie intervalle $I$) est la donné d'un couple $(\omega,c)$, où $\omega$ est une orientation de $I$, et $c$ un foncteur $c : I \longrightarrow X$. Quand on regarde la catégorie des chemins d'un `type' déterminé, ce \textit{type} est $(I,\omega)$, et non $I$. C'est quand $\omega$ est fixé une bonne fois que cette catégorie s'identifie impunément à $\uHom(I,X)$ sans plus. % [sic] la donné -> la donnée, missing final e

Il faut voir cette catégorie de chemins comme munie de la structure essentielle
\[
\begin{tikzcd}
\Ch_{(I,\omega)}(X) \arrow[d, "{(\sigma,\beta)}"] \\
X \times X
\end{tikzcd}
\qquad
\begin{array}{l}
c \longmapsto (\sigma(c),\beta(c)) \\
\text{ou $\sigma(c)$ est la \textit{source,}} \\
\qquad \text{$\beta(c)$ le \textit{but} de $c$.}
\end{array} % [sic] ou sigma(c) -> où sigma(c), missing accent
\]

\printedpage{12}
% head: A. Grothendieck : Dérivateurs, Ch. VII, Édition des Universités de Montpellier II et Paris VII
% foot: Version 22 décembre 2001

\msp{page 18}

\dgrblock{dgrb290}{The endpoints are part of the fixed-type object}
\begin{dcomment}{exposition}{p17}{The endpoints are part of the fixed-type object}
Grothendieck's insistence on orientation is most transparent over the category of endpoint pairs. The category attached to an oriented type $(I,\omega)$ is not merely $\operatorname{Fun}(I,X)$; it comes with the functor
\[
 (\sigma,\beta):\operatorname{Fun}(I,X)\longrightarrow X\times X.
\]
On a natural transformation $\alpha:c\Rightarrow d$, its value is $(\alpha_0,\alpha_n)$. Thus a morphism of paths may move the endpoints, and the displayed functor records exactly how. Its strict fibre over $(x,y)$ consists of paths with endpoints $x,y$ and natural transformations which are the identity at both endpoints. This fibre is a concrete category of fixed-endpoint ladders.

Let $\omega'$ be the reverse orientation on the same underlying interval category. The identity on that underlying category identifies the two ordinary functor categories; under this identification the endpoint pair changes from $(c(0),c(n))$ to $(c(n),c(0))$. Consequently the comparison is a functor over the transposition $s(x,y)=(y,x)$. If $x\ne y$, it carries the fibre over $(x,y)$ to the fibre over $(y,x)$, not to itself. This explains why equality of the underlying functor categories does not erase the distinction between a path and its reverse.

The endpoint map $\alpha:\{0,1\}_{\mathrm{disc}}\to I$ packages the same structure as restriction. For positive length its two images are distinct. For length zero it is the fold onto the sole vertex, so the word ``inclusion'' on the next page needs that qualification; restriction is then the diagonal $X\to X\times X$. All the formulas still hold. The subsequent discussion of symmetric types compares this identity-on-underlying-diagram reversal with the additional comparison induced by a reflection of the shape. Their difference belongs to the structure over $X\times X$, which is why the base category has been kept visible.
\end{dcomment}


Ce foncteur s'identifie au foncteur de composition
\[
c \longmapsto c \circ \alpha
\]
où $\alpha$ est le foncteur d'inclusion
\[
\left\{
\begin{array}{lll}
\varepsilon \xrightarrow{\alpha} I & & \\[1ex]
\varepsilon = \{0,1\}, & \multicolumn{2}{l}{\text{catégorie discrète à deux éléments,}} \\[1ex]
\alpha \text{ donné par} & \alpha(0) = & \text{source de } I \\[1ex]
 & \alpha(1) = & \text{but de } I.
\end{array}
\right.
\]
Si $\omega'$ est l'orientation opposée de $I$, on a un diagramme commutatif
\[
\begin{tikzcd}
\Ch_{(I,\omega)}(X) \arrow[r, "\varphi", "\sim"'] \arrow[d, "{(\sigma,\beta)}"] & \Ch_{(I,\omega')}(X) \arrow[d, "{(\sigma',\beta')}"] \\
X \times X \arrow[r, "s"] & X \times X,
\end{tikzcd}
\]
où $s$ est le foncteur de symétrie
\[
s(x,y) = (y,x)
\]
(\footnote[6]{Le foncteur $\varphi$ est celui qui associe à tout chemin $c$, de type $\tau$, le \textit{chemin opposé $c^{\opp}$} ou $c^{-1}$, de type $\tau^{\opp}$. Si $\tau = \tau^{\opp}$, c'est donc un chemin \textit{de même type,} et $\varphi$ est un \textit{automorphisme} de $\Ch_\tau(X)$. Cet automorphisme (fonctoriel en $X$) est déduit d'un automorphisme $\theta$ de $I$, cf. ci-dessous.}). Donc l'isomorphisme $\varphi$ ne respecte pas les flèches structurales `source-but' $(\sigma,\beta)$.

Il y a cependant un embarras quand il existe un isomorphisme d'intervalles \textit{orientés}
\[
\theta : (I,\omega) \xrightarrow{\sim} (I,\omega'),
\]
i.e. quand $I$ est symétrique. Ce n'est pas le cas pour les $\bDelta_i$ $(i \geq 1)$ (ces catégories n'ont pas d'autres automorphismes que l'identité), mais c'est le cas pour les catégories $\Phi_1$, $\Psi_1$ de l'exemple $4^o$, plus généralement \msp{page 19} pour les catégories réduites de longueur \textit{paire.} Comme chaque fois qu'on a un isomorphisme $\theta$ d'intervalles \textit{orientés,} on trouve un isomorphisme entre les catégories de chemins de type $(I,\omega)$ et $(I,\omega)^{\opp} = (I,\omega')$, \textit{compatible} cette fois avec les foncteurs $(\sigma,\beta)$ :
\[
\begin{tikzcd}
\Ch_{(I,\omega)}(X) \arrow[rr, "\Psi_\theta", "\sim"'] \arrow[dr, "{(\sigma,\beta)}"'] & & \Ch_{(I,\omega')}(X) \arrow[dl, "{(\sigma',\beta')}"] \\
 & X \times X. &
\end{tikzcd}
\]

\printedpage{13}
% head: A. Grothendieck : Dérivateurs, Ch. VII, Édition des Universités de Montpellier II et Paris VII
% foot: Version 22 décembre 2001

\msp{page 20}

Il est bon de regarder un \textit{type} de chemins comme une donné non catégorique, mais combinatoire, une simple suite $\tau$ de $n$ signes $\tau_i$ $(1 \leq i \leq n)$ $+$ ou $-$ ($n =$ longueur des chemins envisagés). Cette donné combinatoire définit un intervalle orienté standard, $I_\tau$. Quand on a un intervalle orienté quelconque $(I,\omega)$, il définit un type $\tau$, et on l'identifie canoniquement au type $I_\tau$. Il vaut mieux écrire $\Ch_\tau(X)$, plutôt que $\Ch_{(I,\omega)}(X)$, pour bien indiquer que pour $X$ donné, c'est là une catégorie totalement `épinglée', sans automorphismes liés à la nature particulière de $\tau$ (à des automorphismes de $\tau$!). % [sic] une donné -> une donnée, missing final e; [sic] Cette donné -> Cette donnée, missing final e

On peut préciser ce point ainsi. Soit $\tau'$ un autre type d'intervalles, et proposons-nous de déterminer des `lois' qui permettent à tout chemin de type $\tau$ dans une catégorie arbitraire $X$, d'associer un chemin de type $\tau'$ de cette même catégorie, et ceci de façon `fonctorielle', i.e. compatible avec les applications horizontales
\[
\begin{tikzcd}
\mathrm{Ch}_\tau(X) \arrow[r, "\mathrm{Ch}_\tau(f)"] \arrow[d, "L_X"'] & \mathrm{Ch}_\tau(X') \arrow[d, "L_{X'}"] \\
\mathrm{Ch}_{\tau'}(X) \arrow[r, "\mathrm{Ch}_{\tau'}(f)"] & \mathrm{Ch}_{\tau'}(X')
\end{tikzcd}
\]
associées à un foncteur $f : X \longrightarrow X'$.

\msp{page 21}

\dgrblock{dgrb291}{Operations on paths are maps of their shapes}
\begin{dcomment}{exposition}{p20}{Operations on paths are maps of their shapes}
Grothendieck wants to do the combinatorics once, in a category small enough that every operation can be seen. A path in an arbitrary category carries many labels, but an operation valid for every category cannot depend on accidental properties of those labels. Its entire content is already present in the universal path. The passage from an oriented interval to its sign word makes this universal object unambiguous. It does not assert that a particular category of paths has no automorphisms; it removes the ambiguity coming from presenting the same unoriented interval in two different ways.

Write $P_I(X)=\Hom_{\Cat}(I,X)$ for the \emph{set} of paths of shape $I$, keeping $X^I=\uHom(I,X)$ for their category. The distinction between these two expressions matters in the displayed Yoneda argument. If
\[
 L_X:P_I(X)\longrightarrow P_J(X)
\]
is natural in all functors $X\to Y$, put $\varphi=L_I(\mathrm{id}_I):J\to I$. Applying naturality to the functor $c:I\to X$ gives
\[
 L_X(c)=c\circ L_I(\mathrm{id}_I)=c\varphi.
\]
Thus every such operation is precomposition with exactly one $\varphi$, and every $\varphi$ defines such an operation. This is the covariant Yoneda lemma, with the representing objects appearing in the reverse order; see \citet[Theorem~2.2.4 and Corollary~2.2.8]{Riehl2016}. No hypothesis about interval categories is needed for this argument.

There is then a canonical operation on morphisms of paths. If $v:c\Rightarrow d$, set $(L_Xv)_j=v_{\varphi(j)}$. For an arrow $a:j\to k$ of $J$, its naturality square is the naturality square of $v$ at $\varphi(a)$. Identities and composites are preserved component by component. Consequently the operation lifts to the functor
\[
 \varphi^*:X^I\longrightarrow X^J.
\]
This proves the extension asserted on manuscript page~23, without assuming a separate coherence theorem. Natural transformations between changes of shape have an equally concrete meaning: $\alpha:\varphi\Rightarrow\psi$ gives $c\alpha:c\varphi\Rightarrow c\psi$, naturally in $c$ and $X$.

The special advantage of intervals now enters. They are posets, so a functor into an interval is determined by its object function, and a natural transformation between two such functors exists precisely when the corresponding pointwise arrows exist. If it exists, it is unique. This uniqueness makes every required diagram of \emph{shape transformations} commute. It does not make the categories $X^I$ thin when $X$ is an arbitrary category: those categories still remember all morphisms of $X$.

For fixed endpoints the same universal-path argument is even more economical. Suppose $I$ and $J$ have specified sources and targets. The operation $c\mapsto c\varphi$ preserves both endpoints for every $X$ and $c$ exactly when $\varphi(0_J)=0_I$ and $\varphi(1_J)=1_I$. Necessity follows by taking $X=I$ and $c=\mathrm{id}_I$; sufficiency is direct evaluation. This is the reason the first category of all changes of shape will soon have to be restricted.
\end{dcomment}


Par exemple, prenant le type $\bDelta_2^+$ (i.e. $+1,+1$), correspondant aux chemins
\[
a_0 \xrightarrow{u_1} a_1 \xrightarrow{u_2} a_2,
\]
on peut associer à un tel chemin le chemin de type \texttt{[}$\bDelta_1^+$\texttt{]}
\[
a_0 \xrightarrow{u_2u_1} a_2,
\]
ou aussi le chemin de type $\bDelta_2^+$
\[
a_0 \xrightarrow{u_2u_1} a_2 \xrightarrow{\mathrm{id}} a_2,
\]
etc. Comme le foncteur
\[
\left\{
\begin{array}{cccl}
X & \longmapsto & \mathrm{Ch}_\tau(X) & \qquad (= \mathrm{Ob}\, \Ch_\tau(X)) \\[1ex]
\mathrm{Cat} & \longrightarrow & (\mathrm{Ens}) &
\end{array}
\right.
\]
est \textit{représenté} par l'objet $I_\tau$ de $X$ \texttt{[plutôt de} Cat\texttt{]}, le sorite des foncteurs représentables nous dit que de telles lois de formation de chemins de type $\tau$ à partir des chemins de type $\tau'$ \texttt{[plutôt} $\tau'$ \texttt{à partir des chemins de type} $\tau$\texttt{]} correspondent biunivoquement aux flèches dans Cat
\[
I_{\tau'} \xrightarrow{\varphi} I_\tau\ ,
\]

% ---- from Der-ChVII-456_p016-020.tex
\printedpage{14}
% head: A. Grothendieck : Dérivateurs, Ch. VII, Édition des Universités de Montpellier II et Paris VII
% foot: Version 22 décembre 2001

en prenant simplement la ‘loi’ correspondante à $\varphi$
\[
c \longmapsto c\circ\varphi \qquad \text{pour } c : I_\tau \longrightarrow X.
\]
De façon plus précise, l'application naturelle
\[
\underbrace{\Hom_{\Cat}(I_{\tau'},I_\tau)}_{=\mathrm{Ob}\,\uHom(I_{\tau'},I_\tau)} \longrightarrow \Hom(\mathrm{Ch}_\tau,\mathrm{Ch}_{\tau'}),
\]
\msp{page 22}
où le deuxième Hom est considéré comme un Hom d'objets dans $\Cat^{\vee}=\uHom(\Cat,\mathrm{Ens})$.
\par\Needspace{70mm}\noindent
Ainsi, dans les deux exemples plus haut, où $\tau=\bDelta_2^+$ et $\tau'=\bDelta_1^+$ ou $\bDelta_2^+$, les foncteurs en sens inverse sont
\[
\begin{tikzcd}
\bDelta_2^+\ \Big\{ & 0 \arrow[r] & 1 \arrow[r] & 2 \\
\bDelta_1^+\ \Big\{ & 0 \arrow[u, mapsto] \arrow[rr] & & 1 \arrow[u, mapsto]
\end{tikzcd}
\]
et
\[
\begin{tikzcd}
\bDelta_2^+\ \Big\{ & 0 \arrow[r] & 1 \arrow[r] & 2 \\
\bDelta_2^+\ \Big\{ & 0 \arrow[u, mapsto] \arrow[r] & 1 \arrow[ur, mapsto] \arrow[r] & 2\,. \arrow[u, mapsto]
\end{tikzcd}
\]
C'est là un principe général pour ramener toutes les ‘opérations’ $\pm$ canoniques sur des chemins dans des catégories arbitraires, et des énoncés les concernant, à des questions où $X$ disparaît et où on travaille dans les catégories intervalles elles-mêmes. Un des avantages techniques, c'est que ce sont des catégories fort petites, donnant lieu de plus à une intuition topologique sûre (qui risque de se volatiliser par moments en travaillant dans des catégories de chemins), et aussi que dans de telles catégories, donc aussi dans les catégories de foncteurs entre elles, il y a toujours au plus
\msp{page 23}
\textit{une seule flèche} allant d'un objet donné vers un autre. En d'autres termes, dans de telles catégories tous les diagrammes sont commutatifs.

Bien sûr, on voit aussi que si on a une loi de formation $L$ de chemins de type $\tau'$ à partir de chemins de type $\tau$ (avec la condition de fonctorialité par rapport à la catégorie $X$ où on prend les chemins), cette loi est nécessairement \textit{fonctorielle par rapport aux chemins :} si on a un homomorphisme
\[
c \xrightarrow{v} c' \qquad \text{dans } \uCh_\tau(X),
\]
\printedpage{15}
% head: A. Grothendieck : Dérivateurs, Ch. VII, Édition des Universités de Montpellier II et Paris VII
% foot: Version 22 décembre 2001
alors la loi $L$ ‘se prolonge’ à de tels $v$ pour donner
\[
L(c) \xrightarrow{L(v)} L(c'),
\]
avec les relations habituelles
\[
L(\mathrm{id}_C)=\mathrm{id}_{L(c)}, \qquad L(v'v)=L(v')L(v). % [sic] \mathrm{id}_C -> \mathrm{id}_c, capital C printed while the chemin is c (cf. L(c))
\]
Cela résulte de la description de $L$ comme une opération $L_\varphi$ en termes d'un foncteur bien déterminé
\[
I_{\tau'} \xrightarrow{\varphi} I_\tau,
\]
lequel à son tour donne pour toute $X$ un \textit{foncteur}
\[
\left\{
\begin{array}{rcl}
\uCh_\tau(X) & \longrightarrow & \uCh_{\tau'}(X)\\[1ex]
c & \longmapsto & c\circ\varphi
\end{array}
\right.
\]
plus précis que simplement un \textit{application} $\mathrm{Ch}_\tau(X)\longrightarrow\mathrm{Ch}_{\tau'}(X)$ entre les ensembles d'objets correspondants. % [sic] un application -> une application, gender agreement

\dgrblock{dgrb292}{Composition and insertion of an identity}
\begin{dcomment}{examples}{p23:entre les ensembles d'objets correspondants.}{Composition and insertion of an identity}
The two examples on manuscript pages~21--22 illustrate how a familiar operation becomes a finite calculation before any coefficient category is chosen. Take $I=[2]=(0\to1\to2)$ and $J=[1]$. The unique endpoint-preserving functor $\varphi:J\to I$ sends $0$ to $0$ and $1$ to $2$. A path $c=(a_0\xrightarrow{u}a_1\xrightarrow{v}a_2)$ therefore becomes the arrow $vu:a_0\to a_2$.

A morphism from $c$ to $d=(b_0\xrightarrow{u'}b_1\xrightarrow{v'}b_2)$ is a triple $(h_0,h_1,h_2)$ satisfying $h_1u=u'h_0$ and $h_2v=v'h_1$. The induced morphism between the composite arrows is the square with sides $h_0,h_2$. Its commutativity follows from the explicit computation
\[
 h_2vu=v'h_1u=v'u'h_0.
\]
The middle component has done its work and can be forgotten. The reversed variance is essential: an operation which reduces the length of a path is represented here by a functor from the shorter interval to the longer one.

For the second operation let $J=I=[2]$ and take the nondecreasing object function $\psi=(0,2,2)$. Then $\psi^*c=(a_0\xrightarrow{vu}a_2\xrightarrow{1}a_2)$, and $\psi^*h=(h_0,h_2,h_2)$. The first naturality square is the square just computed; the second is $h_2\,1=1\,h_2$. Since $\psi\psi=\psi$, the resulting operation on paths is strictly idempotent. This equality is checked on the three objects of $[2]$, rather than on all categories and all pairs of composable arrows.

Endpoint preservation alone has not made such maps rigid. On $[2]$ the three nondecreasing endpoint-preserving endofunctors have object functions $(0,0,2)$, $(0,1,2)$, and $(0,2,2)$. They respectively place the composite in the second edge, leave the path unchanged, and place the composite in the first edge. This also supplies a concrete reminder of the correction to Proposition~1(b): fixing the endpoints does not force an arbitrary endofunctor of a nonreduced interval to be the identity. The rigidity there belongs to isomorphisms, while the present theory deliberately needs the noninvertible maps.
\end{dcomment}


\msp{page 24}
On a le choix maintenant entre

\begin{enumerate}[label=\alph*)]
\item Une étude plus ou moins systématique des foncteurs entre intervalles. Je préfère attendre pour cela d'avoir besoin de ceci et de cela, et de ne faire que le strict minimum pour les constructions que j'ai en vue.
\item Une description de la compositions des chemins. % [sic] la compositions -> la composition, number agreement
\item Un formalisme des catégories de chemins \textit{de type non précisé.}
\end{enumerate}

Comme mon objectif est surtout de décrire de telles catégories $\uCh(X)$ (sans indice $\tau$!), avec des bonnes propriétés, je commencerai par c), regarderai ensuite b) comme un test crucial pour l'utilisabilité des notions obtenues, et enfin m'occuperai de a) quand je saurai de quoi j'aurai en besoin.

\printedpage{16}
% head: A. Grothendieck : Dérivateurs, Ch. VII, Édition des Universités de Montpellier II et Paris VII
% foot: Version 22 décembre 2001

\section*{4\quad Catégorie de chemins $\uCh(X)$}

Je veux ‘inclure’ toutes les catégories $\uCh_\tau(X)$ de chemins de tous les types qu'on veut, dans une même catégorie $\uCh(X)$. Rappelons
\[
\uCh_\tau(X)=\uHom(I_\tau,X),
\]
où $I_\tau$ est l'intervalle (orienté) standard de type $\tau$.

L'ensemble de tous les types d'intervalles
\msp{page 25}
peut être identifié au monoïde libre Ty à deux générateurs $\Delta^+$ et $\Delta^-$, dont les éléments sont les suites finies quelconques (eventuellement vides) d'objets alternativement égales à $\Delta^+$ et à $\Delta^-$. Si le premier élément est $\Delta^+$, on a un type de chemin à direction initiale directe, rétrograde dans le cas contraire. La décomposition de $I_\tau$, ou de $\mathrm{Fl}\,(I_\tau)$, en intervalles critiques correspond à l'écriture réduite d'un mot comme % [sic] eventuellement -> éventuellement, missing accent
\[
(\Delta^+)^{l_1}(\Delta^-)^{l_2}\ \cdots\ (\Delta^?)^{l_r} \qquad (\Delta^?=\Delta^+ \text{ si } r \text{ impair, } \Delta^?=\Delta^- \text{ si } r \text{ pair}),
\]
ou comme
\[
(\Delta^-)^{l_1}(\Delta^+)^{l_2}\ \cdots\ (\Delta^?)^{l_r} \qquad (\Delta^?=\Delta^- \text{ si } r \text{ impair, } \Delta^?=\Delta^+ \text{ si } r \text{ pair}).
\]
Les $l_i$ sont les longueurs des intervalles critiques successifs (ou des chaines structurales successives), leur nombre est la longueur réduite de l'intervalle.

On va considérer Ty comme l'ensemble des objets d'une catégorie $\uTy$ :
\[
\left\{
\begin{array}{l}
\mathrm{Ob}\,\uTy=\Ty=\ \text{monoïde libre à deux générateurs } \Delta^+,\ \Delta^- \text{ précédents.}\\[1ex]
\text{Si } \tau,\tau'\in\Ty,\ \Hom(\tau,\tau')=\Hom_{\Cat}(I_\tau,I_{\tau'}).
\end{array}
\right.
\]

(\footnote[7]{Définition un peu canulée, il faut se borner aux homomorphismes de $I_\tau$ dans $I_{\tau'}$ compatibles avec source et but - cf. p.~29.}). On notera qu'on a non seulement un ensemble $\Hom_{\Cat}(I_\tau,I_{\tau'})$, mais même une catégorie $\uHom(I_\tau,I_{\tau'})$, ce qui revient ici à dire que $\Hom(I_\tau,I_{\tau'})$ (i.e. $\Hom(\tau,\tau')$) peut être vu comme un ensemble \textit{ordonné.} Mais pour l'instant nous allons ignorer cette relation d'ordre.

\msp{page 26}
La composition des morphismes est la composition des foncteurs entre catégories $I_\tau$. Ainsi on a un foncteur \textit{pleinement fidèle} (\footnote[8]{faux, cf. annotation marginale page précédente - on ne prend pas \textit{tous} les foncteurs dans Cat entre $I_\tau$ et $I_{\tau'}$.}) canonique
\[
\tau \longmapsto I_\tau : \uTy \xrightarrow{\rho} (\Cat) \qquad \text{(\footnotemark[9]),}
\]
\footnotetext[9]{foncteur \textit{réalisation} d'un type (combinatoire) d'intervalles.}% footnote text printed at the foot of p. 17
\printedpage{17}
% head: A. Grothendieck : Dérivateurs, Ch. VII, Édition des Universités de Montpellier II et Paris VII
% foot: Version 22 décembre 2001
par quoi la catégorie $I_\tau$ \texttt{[plutôt} $\uTy$\texttt{]} pourrait s'identifier à une sous-catégorie pleine de (Cat) (celle des $I_\tau$). Mais je préfère ne pas identifier $\tau\in\Ty$ (donnée combinatoire) avec $I_\tau$ (donnée catégorique).

\dgrblock{dgrb293}{Three successive categories of types}
\begin{dcomment}{exposition}{p25}{Three successive categories of types}
Grothendieck is testing which changes of parameter should count as morphisms of paths. The changes of definition in these pages are purposeful, but a notation retained through all three stages can obscure which fullness assertion remains true. We distinguish the stages here, without altering the source's notation. All three categories have the finite sign words as objects, including the empty word $e$.

Let $T_{\mathrm{all}}$ have every functor $I_\tau\to I_\upsilon$ as a morphism; let $T_{\partial}$ retain just the functors preserving the two chosen endpoints; and let $T_{\uparrow}$ retain those which also preserve the external, linear order of vertices. Thus
\[
 T_{\uparrow}\subseteq T_{\partial}\subseteq T_{\mathrm{all}}.
\]
The category eventually denoted by $\uTy$ in these pages is $T_{\uparrow}$. On page~25 it is initially $T_{\mathrm{all}}$, while pages~29--31 pass through $T_{\partial}$. For $T_{\mathrm{all}}$, realization in $\Cat$ is fully faithful by definition. For $T_{\partial}$, realization is fully faithful into $\varepsilon\backslash\Cat$, where $\varepsilon$ is the discrete category on two objects; forgetting $\varepsilon$ loses fullness. For $T_{\uparrow}$ even the embedding into $\varepsilon\backslash\Cat$ is generally not full. The footnotes retracting fullness on page~26 and the assertion of fullness on page~31 are therefore compatible when read at their respective stages.

Here is a small witness to the last distinction. From the alternating interval $I_{+-+}$ to $I_+$ there is a functor with object function
\[
 (0,1,2,3)\longmapsto(0,1,0,1).
\]
Every forward transition goes to $0\to1$ and the backward transition, whose categorical direction is from its right vertex to its left vertex, also goes to $0\to1$. Both endpoints are preserved, but the vertex function is not nondecreasing. Thus it belongs to $T_{\partial}$ and not to $T_{\uparrow}$. The monotonicity condition excludes this backtracking.

The sign words themselves form a free monoid on $+$ and $-$. Arbitrary adjacent repetitions are allowed; the alternating feature is the alternation of signs between \emph{maximal constant runs}. A nonempty word has a unique expression $\varepsilon_1^{l_1}\cdots\varepsilon_r^{l_r}$ with $l_i>0$ and $\varepsilon_{i+1}=-\varepsilon_i$. Its length is $\sum l_i$, its useful length is $r$, and its critical intervals are these runs with their endpoint vertices restored. The empty word has length and useful length zero and no initial sign. This last exception should remain in force whenever the text speaks of the first generator.

In the undercategory description, $\varepsilon\to I_e$ identifies the two objects of $\varepsilon$. It is a structure map, not an embedding; the hooked arrows drawn on page~29 should not be read as asserting injectivity in this degenerate case. Keeping $e$ is indispensable: it will be the unit for concatenation.
\end{dcomment}


Soit $X$ une (petite?) catégorie. On a un contrafoncteur
\[
\left\{
\begin{array}{rcl}
S & \longmapsto & \uHom(S,X)\\[1ex]
(\Cat)^o & \longrightarrow & (\Cat)
\end{array}
\right.
\]
défini par $X$, d'où en composant avec $\rho$ un contrafoncteur
\[
\uTy^{o} \xrightarrow{\mathrm{CH}_X} \Cat.
\]
C'est le contrafoncteur
\[
\left\{
\begin{array}{l}
\tau \longmapsto \uHom(I_\tau,X) \overset{\text{déf}}{=} \uCh_\tau(X),\\[1ex]
\text{i.e. } \mathrm{CH}_X(\tau)=\uCh_\tau(X).
\end{array}
\right.
\]
Mais un contrafoncteur sur une petite catégorie $\uTy$ donne lieu à une \textit{catégorie fibrée} (et même scindée) sur celle-ci.

\msp{page 27}
Je vais désigner cette catégorie fibrée par $\uCh(X)$. (Provisoirement - si elle s'avère trop grosse, je la restreindrai \dots) Ainsi on a
\[
\begin{tikzcd}[row sep=large]
\uCh(X) \arrow[d, "\pi_X"] \\
\uTy\ .
\end{tikzcd}
\]
La fibre de $\uCh(X)$ en $\tau\in\mathrm{Ob}\,\uTy$ est la catégorie $\uCh_\tau(X)$ des chemins de type $\tau$ dans $X$. Cela nous dit donc quelles sont les flèches dans $\uCh(X)$ au dessus des flèches identiques dans $\uTy$. Mais que sont les flèches plus générales? On se donne une flèche dans $\uTy$
\[
\varphi : \tau' \longrightarrow \tau, \qquad \text{i.e. } I_{\tau'} \xrightarrow{\varphi} I_\tau,
\]
quelles sont les flèches de $\uCh(X)$ au dessus de $\varphi$? Par le sorite général des catégories fibrées, si
\[
\begin{array}{lclcl}
c\in\uCh(X)_\tau & = & \uCh_\tau(X) & \overset{\text{déf}}{=} & \uHom(I_\tau,X)\\[1ex]
c'\in\uCh(X)_{\tau'} & = & \uCh_{\tau'}(X) & \overset{\text{déf}}{=} & \uHom(I_{\tau'},X)
\end{array}
\]
on a
\[
\boxed{\Hom_\varphi(c',c)\simeq\Hom(c',\underbrace{\varphi^*(c)}_{=c\circ\varphi}).}
\]
Donc : les homomorphismes du
\printedpage{18}
% head: A. Grothendieck : Dérivateurs, Ch. VII, Édition des Universités de Montpellier II et Paris VII
% foot: Version 22 décembre 2001
\msp{page 28}
chemin $c'$ de type $\tau'$ dans $X$, dans le chemin $c$ de type $\tau$, ‘au dessus’ du morphisme $\varphi$ de $\tau'$ dans $\tau$, i.e. du foncteur
\[
\varphi : I_{\tau'} \longrightarrow I_\tau,
\]
est un homomorphisme $v$ de $c'$ dans $c\circ\varphi$, i.e.
\[
\begin{tikzcd}[row sep=large]
I_{\tau'} \arrow[r, "\varphi"] \arrow[dr, "c'"', ""{name=A}] & I_\tau \arrow[d, "c", ""'{name=B}] \\
& X.
\arrow[from=A, to=B, bend left=20, shorten <=6pt, shorten >=2pt, "\circled{v}"]
\end{tikzcd}
\]
La composition des flèches est évidente :
\[
\begin{tikzcd}[row sep=large]
I_{\tau''} \arrow[r, "\varphi'"] \arrow[drr, "c''"', ""{name=C}] & I_{\tau'} \arrow[r, "\varphi"] \arrow[dr, "c'"', ""{name=A}] & I_\tau \arrow[d, "c", ""'{name=B}] \\
& & X,
\arrow[from=A, to=B, bend left=20, shorten <=6pt, shorten >=2pt, "v"]
\arrow[from=C, to=A, shorten <=4pt, shorten >=4pt, "v'"]
\end{tikzcd}
\]
on a
\[
\begin{array}{rccl}
v : & c' & \longrightarrow & c\circ\varphi, \quad \text{d'où } c'\circ\varphi' \xrightarrow{v*\varphi'} \underbrace{(c\circ\varphi)\circ\varphi'}_{=c\circ(\varphi\varphi')}\\[3ex]
v' : & c'' & \longrightarrow & c'\circ\varphi'
\end{array}
\]
et
\[
v''=(v*\varphi')\circ v', \qquad \text{donc} \qquad
\begin{tikzcd}
c'' \arrow[r, "v'"] \arrow[rr, bend right=12, "\circled{v''}"'] & c'\circ\varphi' \arrow[r, "v*\varphi'"] & c\circ(\varphi\varphi').
\end{tikzcd}
\]
Mais a-t-on des foncteurs source et but
\[
\uCh(X) \overset{\sigma}{\underset{\beta}{\rightrightarrows}} X \qquad
\begin{array}{ccl}
c & \longmapsto & \text{source}(c)\\[1ex]
c & \longmapsto & \text{but}(c)
\end{array}
\qquad ?
\]
Il faudrait donc qu'un homomorphisme $v$ comme dessus
\[
v : c' \longrightarrow c\circ\varphi
\]
définisse
\[
\begin{array}{lrcl}
\sigma(c) : & \text{source}(c') & \longrightarrow & \text{source}(c)\\[1ex]
\beta(c) : & \text{but}(c') & \longrightarrow & \text{but}(c).
\end{array}
\]

\dgrblock{dgrb294}{The split fibration of all paths}
\begin{dcomment}{context}{p27}{The split fibration of all paths}
Grothendieck now uses the machinery of descent to assemble paths of different shapes. The construction is the one in \citet[Expos\'e~VI, \S\S8--9]{GrothendieckSGA1}: a strict contravariant functor to categories supplies a category with a chosen, strictly composable family of Cartesian lifts. The commentary on Chapter~VI, manuscript pages~26 and~53--56, discussed splittings and their universal constructions. Here the splitting needs no choices, because it is given by literal precomposition.

For any of the three shape categories $T$, put $F_X(\tau)=X^{I_\tau}$. An object of its Grothendieck construction is $(\tau,c)$, and a morphism
\[
 (\tau',c')\longrightarrow(\tau,c)
\]
is a pair $(\varphi,v)$ with $\varphi:\tau'\to\tau$ and $v:c'\Rightarrow c\varphi$. If $(\varphi',v'):(\tau'',c'')\to(\tau',c')$, their composite is
\[
 (\varphi,v)(\varphi',v')=
 \bigl(\varphi\varphi',(v\varphi')\circ v'\bigr).
\]
At $i\in I_{\tau''}$ the second component is $v_{\varphi'(i)}v'_i$. Threefold composition is associative because both parenthesizations have component $v_{\varphi'\varphi''(i)}v'_{\varphi''(i)}v''_i$ and the same shape map. The identity is $(1_\tau,1_c)$. These checks explain the word ``\emph{\'evidente}'' in the source without suppressing the variance.

The chosen lift of $\varphi:\tau'\to\tau$ at $(\tau,c)$ is
\[
 (\varphi,1_{c\varphi}):(\tau',c\varphi)\longrightarrow(\tau,c).
\]
To verify its Cartesian property, let a map from $(\omega,d)$ to $(\tau,c)$ have shape $\varphi\theta$ and transformation $w:d\Rightarrow c\varphi\theta$. There is exactly one factor over $\theta$ through the displayed lift, namely $(\theta,w)$. The chosen lifts compose on the nose, since $(c\varphi)\theta=c(\varphi\theta)$. More generally $(\varphi,v)$ is Cartesian exactly when $v$ is invertible: it factors as the vertical arrow $(1,v)$ followed by this chosen lift, and a vertical Cartesian arrow is an isomorphism.

For $T=T_{\partial}$ or $T_{\uparrow}$ the endpoint functors send $(\varphi,v)$ to $v_{0_{\tau'}}$ and $v_{1_{\tau'}}$. Their targets are $c(0_\tau)$ and $c(1_\tau)$ because $\varphi$ preserves endpoints. The formula for composition immediately gives functoriality. Likewise a functor $H:X\to Y$ acts by $(\tau,c)\mapsto(\tau,Hc)$ and $(\varphi,v)\mapsto(\varphi,Hv)$, preserving the chosen Cartesian lifts. Restricting the permitted shape maps is literally base change of this split fibration, which proves the fiber-product assertion on page~30.

Size causes no mathematical obstruction here. Fix universes $\mathcal U\in\mathcal V$. The set of finite words and their finite interval functors is $\mathcal U$-small. If $X$ is a $\mathcal V$-small category, each $X^{I_\tau}$ and the total construction are $\mathcal V$-small. One may therefore take $X=\Cat_{\mathcal U}$ after working in a sufficiently larger universe. This supplies the convention which the source deliberately leaves implicit.
\end{dcomment}


% ---- from Der-ChVII-456_p021-025.tex
\printedpage{19}
% head: A. Grothendieck : Dérivateurs, Ch. VII, Édition des Universités de Montpellier II et Paris VII
% foot: Version 22 décembre 2001

\msp{page 29}
Mais comme
\[
\left\{
\begin{array}{l}
\sigma(c)=c(0_\tau),\ \beta(c)=c(1_\tau)\\[1ex]
\text{où } 0_\tau,\ 1_\tau \text{ sont les objets source et but de l'intervalle orienté standard } I_\tau\ ,
\end{array}
\right.
\]
et de même pour $\tau'$, $v$ n'induira bien
\[
\underbrace{\sigma(c')}_{=c'(0_{\tau'})} \longrightarrow c(\varphi(0_{\tau'})) \overset{?}{=} \underbrace{\sigma(c)}_{c(0_\tau)},
\]
et itou pour $\beta(c')$ et $\beta(c)$, que si on impose
\[
\boxed{
\left\{
\begin{array}{lcl}
\varphi(0_{\tau'}) & = & 0_\tau\\
\varphi(1_{\tau'}) & = & 1_\tau,
\end{array}
\right.
}
\]
i.e. que le foncteur envisagé $\varphi$ soit \textit{compatible avec la source et le but.} C'est ce visiblement qu'il fallait imposer dans la définition de % [sic] C'est ce visiblement -> C'est visiblement ce, word order
\[
\mathrm{Hom}_{\ul{Ty}}(\tau',\tau)=\underline{\mathrm{Hom}}^{\circled{int}}(I_{\tau'},I_\tau)\quad\text{-}
\]
il faut prendre les foncteurs $\varphi$ respectant les structures d'intervalle, i.e. rendant commutatif le diagramme
\[
\begin{tikzcd}[column sep=small]
I_{\tau'} \arrow[rr] & & I_\tau \\
& \varepsilon=\{0,1\} \arrow[ul, hook', "i_{\tau'}"] \arrow[ur, hook, "i_\tau"'] &
\end{tikzcd}
\]
définissant l'orientation de $I_{\tau'}$ et de $I_\tau$ (\footnote[10]{En fait, il faut sans doute restreindre encore plus les morphismes en exigeant qu'il soient \textit{croissants} sur les objets - cf. page 31.}). Cette rectification faite, j'ai bon espoir que la catégorie $\ul{Ch}(X)$ qu'on vient de construire est `la bonne'. On pourrait désigner par $\ul{CH}(X)$ la catégorie plus grande, sur la catégorie $\ul{Typ}$ ayant mêmes objets que $\ul{Ty}$, mais plus de morphismes : à savoir % [sic] (in footnote 10) qu'il soient -> qu'ils soient, missing plural s
\msp{page 30}
\[
\begin{array}{lcl}
\mathrm{Hom}_{\ul{Typ}}(\tau,\tau') & = & \mathrm{Hom}_{\mathrm{Cat}}(I_\tau,I_{\tau'})\\[0.5ex]
\mathrm{Hom}_{\ul{Ty}}(\tau,\tau') & = & \mathrm{Hom}^{\mathrm{int}}_{\mathrm{Cat}}(I_\tau,I_{\tau'})\\[0.5ex]
 & & \text{(homomorphismes respectant la structure d'intervalle).}
\end{array}
\]
On aura
\[
\ul{Ch}(X)\simeq\ul{CH}(X)\times_{\ul{Typ}}\ul{Ty},
\]
\printedpage{20}
% head: A. Grothendieck : Dérivateurs, Ch. VII, Édition des Universités de Montpellier II et Paris VII
% foot: Version 22 décembre 2001
i.e. $\ul{Ch}(X)$ est la \textit{restriction} de la catégorie fibrée $\ul{CH}(X)$ sur $\ul{Typ}$, à la sous-catégorie (non-pleine) $\ul{Ty}$ de $\ul{Typ}$.

\dgrblock{dgrb295}{Changing shape need not have a direct image}
\begin{dcomment}{examples}{p28}{Changing shape need not have a direct image}
The projection to the category of types is a fibration because paths can always be restricted. Extending a path to a new shape asks a different question, and generally requires colimits in its coefficient category. This is the usual distinction between a fibration and a bifibration; compare \citet[Expos\'e~VI, \S10 and \S11(f)]{GrothendieckSGA1}.

Consider the interval $J=I_{-+}=(0\leftarrow1\to2)$ and its unique collapse $q:J\to e$. This collapse belongs even to $T_{\uparrow}$. Its inverse-image functor is the diagonal
\[
 q^*:X\longrightarrow X^J,
\]
sending an object to its constant span. If an opCartesian lift over $q$ existed at every span $c$, its target $q_!c$ would represent the functor
\[
 x\longmapsto\operatorname{Nat}(c,\mathrm{const}_x).
\]
Indeed every map from $c$ over $q$ to a constant path would factor uniquely through the opCartesian lift. This is precisely the universal property of the pushout of the span $c(0)\leftarrow c(1)\to c(2)$.

Take $X=J$ itself and $c=1_J$. There is no cocone on this diagram: neither $0$ nor $2$ maps to the other, and neither maps to $1$. Thus $q_!c$ does not exist, and the path projection is not an opfibration in this example. This does not reveal a gap in Grothendieck's assertion, which is only that the construction is split fibrational. It identifies precisely what would be added by a stronger assertion.

For $X=\mathrm{Set}$ the same calculation has a complete answer. Let $A\leftarrow C\to B$ be a span. Then $q_!c$ is the quotient of $A\amalg B$ by the equivalence relation generated by identifying the two images of each $z\in C$. A map from this quotient to $S$ is exactly a pair $A\to S$, $B\to S$ which agrees on $C$, so it is exactly a natural transformation $c\Rightarrow\mathrm{const}_S$. For a general change of shape $\varphi$ and a cocomplete $X$, left Kan extension supplies the corresponding left adjoint to $\varphi^*$. Such extra adjoints come from $X$; the split fibration itself is available without them.
\end{dcomment}


\dgrblock{dgrb296}{How much endpoint preservation is forced}
\begin{dcomment}{development}{p29}{How much endpoint preservation is forced}
The correction on page~29 gives exactly what is needed for endpoint maps obtained by taking components of a natural transformation. There is a slightly larger possibility if one allows those components to be followed by canonical arrows inside the target interval. Describing it clarifies why endpoint preservation is a structural choice, and why concatenation will justify the stricter choice.

\textbf{Endpoint evaluation criterion.}\ours{}
Let $T$ be a subcategory of $T_{\mathrm{all}}$ with all sign words as objects. Its path construction admits functors $\sigma_X,\beta_X:\int F_X\to X$, natural in $X$, restricting in each fiber to evaluation at the specified source and target, if and only if every $\varphi:\tau'\to\tau$ in $T$ admits arrows
\[
 a_\varphi:\varphi(0_{\tau'})\longrightarrow0_\tau,
 \qquad
 b_\varphi:\varphi(1_{\tau'})\longrightarrow1_\tau
 \quad\text{in }I_\tau.
\]
When they exist, these endpoint functors are unique. The permitted functors form a category $T_{\mathrm{ev}}$ containing $T_{\partial}$, and it is the largest such subcategory of $T_{\mathrm{all}}$.

\emph{Proof.} For necessity take $X=I_\tau$ and the universal path $c=1_{I_\tau}$. The Cartesian arrow
\[
 (\varphi,1_\varphi):(\tau',\varphi)\longrightarrow(\tau,1_{I_\tau})
\]
must be sent by $\sigma_X$ to an arrow $\varphi(0_{\tau'})\to0_\tau$, and by $\beta_X$ to an arrow $\varphi(1_{\tau'})\to1_\tau$. Each is unique when it exists, because an interval is thin. For composable maps $\tau''\xrightarrow{\psi}\tau'\xrightarrow{\varphi}\tau$, the arrows
\[
 a_\varphi\circ\varphi(a_\psi),\qquad
 b_\varphi\circ\varphi(b_\psi)
\]
have precisely the endpoints required for $\varphi\psi$. Identities satisfy the requirement as well. Hence the permitted functors form a category.

To construct the endpoint maps, define on $(\varphi,v):c'\to c$
\[
 \sigma_X(\varphi,v)=c(a_\varphi)\circ v_{0_{\tau'}},
 \qquad
 \beta_X(\varphi,v)=c(b_\varphi)\circ v_{1_{\tau'}}.
\]
For a second morphism $(\psi,w):c''\to c'$, naturality of $v$ at $a_\psi$ gives
\[
 v_{0_{\tau'}}c'(a_\psi)
   =c\varphi(a_\psi)v_{\psi(0_{\tau''})}.
\]
After composing with $c(a_\varphi)$ and $w_{0_{\tau''}}$, this says that the product of the two source arrows is
\[
 c(a_{\varphi\psi})
 v_{\psi(0_{\tau''})}w_{0_{\tau''}},
\]
which is the source arrow assigned to the composite. The same argument with $b_\psi$ proves the target identity. Identity morphisms and naturality in $X$ follow directly from the formulas.

Finally, every $(\varphi,v)$ factors as a vertical morphism $(1,v)$ followed by a chosen Cartesian arrow. The endpoint functors are prescribed on vertical arrows. Naturality in $X$ applies to $c:I_\tau\to X$ and forces their values on the Cartesian arrow to be $c(a_\varphi)$ and $c(b_\varphi)$. Thus the formulas are forced, proving uniqueness and maximality.\hfill$\square$

For example, the functor $e\to I_+$ selecting $0$ belongs to $T_{\mathrm{ev}}$: its source comparison is an identity and its target comparison is $0\to1$. It does not preserve both endpoints. The functor selecting $1$ does not belong to $T_{\mathrm{ev}}$, since there is no arrow $1\to0$. This already prevents endpoint evaluation on $T_{\mathrm{all}}$.

The enlargement does not solve the concatenation problem. Concatenate two copies of the admissible map $e\to I_+$ selecting $0$. The common vertex of $e*e=e$ would have to map both to vertex $0$ of the first target interval and to vertex $1$, the starting vertex of the second interval in $I_{++}$. The two prescribed maps therefore do not glue. Strict endpoint preservation makes precisely these images agree. The theorem is an elementary specialization of the general description of functors out of a split fibration in \citet[Expos\'e~VI, Proposition~12.1]{GrothendieckSGA1}. We checked that description and the Yoneda treatment cited above; we did not find this particular maximal class of interval maps stated there. The marker records this explicit deduction and proof, without a claim of priority.
\end{dcomment}


\textbf{NB} Les catégories fibrées $\ul{CH}(X)$, $\ul{Ch}(X)$ sur $\ul{Typ}$, $\ul{Ty}$ dépendent fonctoriellement de $X$, variant dans (Cat), comme il se doit. Comme de juste, on s'y intéresse surtout quand $X$ \textit{n'est pas} dans (Cat), i.e. quand $X$ est une `grosse' catégorie (telle que (Cat) elle-même). Je me dispense d'expliciter ici les questions pertinentes d'Univers.

Deux tests essentiels, pour s'assurer si la notion est bonne
\begin{enumerate}
\item[a)] La théorie des $\mathcal{M}\Sigma^{-1}$.
\item[b)] La composition des chemins doit bien marcher.
\end{enumerate}

\msp{page 31}
On a le diagramme de foncteurs canoniques
\[
\begin{tikzcd}
\ul{Ty} \arrow[r, hook] \arrow[d, hook] & \varepsilon\backslash\mathrm{Cat} \arrow[d] \\
\ul{Typ} \arrow[r, hook] & \mathrm{Cat},
\end{tikzcd}
\]
où les flèches horizontales sont pleinement fidèles.

\dgrblock{dgrb297}{The fraction test and hammock localization}
\begin{dcomment}{status}{p30}{The fraction test and hammock localization}
Grothendieck's two tests concern different levels of the construction. Concatenation should be an operation on actual paths and their morphisms; localization should arise after paths are subjected to specified identifications or homotopical comparisons. The category $\uCh(X)$ constructed here has not yet been asserted to satisfy the universal property of a localization.

For a category $\mathcal M$ and a specified class $\Sigma$ of arrows, ordinary localization is already constructed by \citet[Chapter~I, \S1.1 and Lemma~1.2]{GabrielZisman1967}. Its arrows are composable words of forward arrows and formal inverses of arrows in $\Sigma$, modulo the category relations and inverse relations. A functor out of the quotient is consequently equivalent to a functor from $\mathcal M$ sending $\Sigma$ to isomorphisms. This existence statement does not require the Ore hypotheses needed for a calculus with uniformly short fractions.

The homotopical comparison already available before these notes is quite specific. In \citet[Definition~2.1 and Propositions~2.2 and~3.1]{DwyerKan1980Calculating}, a hammock is a finite grid of zigzags, with backward horizontal arrows and vertical arrows required to be weak equivalences. Identity columns are removed and adjacent columns pointing in the same direction are composed. Their mapping simplicial sets are equivalent to the standard simplicial localization, and their sets of components recover ordinary localization:
\[
 \pi_0L^H(\mathcal M,\Sigma)(x,y)
   \cong\Hom_{\mathcal M[\Sigma^{-1}]}(x,y).
\]
Thus there is a proven construction in which changing the length and subdivision of zigzags is accompanied by the higher comparison data needed for mapping spaces.

This is a comparison of mathematical mechanisms, not an identification of the constructions. The present $\uCh(X)$ allows arbitrary vertical natural transformations, has not specified which backward arrows may be inverted, and uses a single ordinary category of types. The hammock conditions, the fixed endpoints, and the higher simplices must all be supplied before a localization theorem can be invoked. The source's demand that the fraction theory ``work'' is a design requirement at this point, not a research conjecture left open in these pages. Its exact implementation in Grothendieck's subsequent text should be checked on its own definitions.
\end{dcomment}


\printedpage{21}
% head: A. Grothendieck : Dérivateurs, Ch. VII, Édition des Universités de Montpellier II et Paris VII
% foot: Version 22 décembre 2001

\section*{5\quad Propriétés des morphismes d'intervalles}

Je me suis plus ou moins convaincu que pour se prêter à des calculs et à des raisonnements simples, il faut encore plus restreindre les morphismes dans la définition de la catégorie $\ul{Ty}$. Je vais exiger que les morphismes $\tau\longrightarrow\tau'$ soient des foncteurs
\[
w : I_\tau\longrightarrow I_{\tau'}
\]
qui sont compatibles avec la structure d'\textit{intervalles orientés} de $I_\tau$, $I_{\tau'}$, i.e. qu'ils transforment source en source, but en but, mais de plus qu'ils soient \textit{croissants} sur les objets (pour la relation d'ordre de progression des chemins). C'est une restriction qui géométriquement me paraît assez naturelle. (Elle est similaire à la restriction analogue, dans la définition pré-catégorique des morphismes entre $\bDelta_i$ \ldots\,) C'est cette hypothèse qui permet de faire une théorie de structure assez simple des morphismes d'intervalles. Mais pour l'instant je vais distinguer `morphismes' et `morphismes croissants'.

\dgrblock{dgrb298}{The same type category in later zigzag constructions}
\begin{dcomment}{later}{p31:Je vais exiger que les morphismes}{The same type category in later zigzag constructions}
Grothendieck's final restriction on changes of shape has a precise counterpart in the later literature on localization. The correspondence can be checked directly, before making any assertion about homotopy theory. Given a word $\tau$, mark the backward transitions of $I_\tau$ and their composites, together with the identities. This makes $I_\tau$ a relative category. An increasing endpoint-preserving functor automatically respects this marking: a nonconstant image of a backward transition is a composite of backward transitions. Conversely, forgetting the marking of an increasing endpoint-preserving relative functor gives exactly a morphism of $\uTy$. The objects, morphisms, and concatenation are therefore unchanged by this decoration.

Under $+\leftrightarrow A$ and $-\leftrightarrow W^{-1}$, this identifies $\uTy$ with the category $Z$ of \citet[Definition~3.5, Remark~3.6, and footnotes~8--10, pp.~341--342]{MazelGee2018Hammocks}. Their weak-equivalence labels restrict which diagrams in a coefficient category are allowed; they do not impose an additional restriction on these monotone shape maps. The empty word remains the monoidal unit. This is an identification of the finite indexing categories, rather than an inference from a resemblance between pictures.

The image factorization implicit in the manuscript also explains a useful structure. If $f:I_\tau\to I_\upsilon$ is increasing, the full subcategory on its image is an interval. Factor $f$ through its standard oriented model:
\[
 I_\tau\mathrel{\mathop{\longrightarrow}^{q}}I_\omega
 \mathrel{\mathop{\lhook\joinrel\longrightarrow}^{j}}I_\upsilon.
\]
The first map is surjective on vertices and the second injective. Their uniqueness follows from the unique increasing enumeration of the image. Nonidentity surjections lower the length, and nonidentity injections raise it. Thus length is a Reedy degree. The reduced interval gives the smallest image of an endpoint-preserving injection: Proposition~5 says that every such image contains all critical vertices. If the interval is not reduced, its reduced inclusion is consequently initial among the proper injections into it; for a reduced interval there are none. These statements follow from the full-subcategory argument already given, without a new hypothesis on coefficients.

This Reedy structure and this exact use of critical vertices occur in \citet[Example~4.12 and Lemma~4.16]{ArakawaCnossen2026}, where the category is denoted $\mathrm{Zig}_{\partial}$. Their work is an August~2026 preprint. We cite these known applications, rather than present the Reedy interpretation as an original result. What remains to be added to the present manuscript is the relative and homotopical diagram theory over this category of shapes.
\end{dcomment}


\msp{page 32}
Soit
\[
f : I_\tau\longrightarrow I_{\tau'}
\]
un morphisme d'intervalles orientés. Considérons les sous-catégories pleines, isomorphes à des $\Delta_i^+$ ou $\Delta_i^-$ comme intervalles orientés, définis par les intervalles critiques de $I_\tau$. Ils forment une suite $\Sigma_\tau$ de sous-catégories de $I_\tau$
\[
\Sigma_\tau=(\bDelta^{\varepsilon_1}_{l_1}(\tau),\bDelta^{\varepsilon_2}_{l_2}(\tau),\ldots,\bDelta^{\varepsilon_r}_{l_r}(\tau))
\]
dont $I_\tau$ est la somme amalgamée (on reviendra sur ce point plus \texttt{[bas]}), $r$ étant la longueur utile de $\tau$. On les appelle les \textit{sous-intervalles critiques} de $I_\tau$. Ainsi $f$ se décompose en une suite de morphismes
\[
f_i : \bDelta^{\varepsilon_i}_{l_i}\longrightarrow I_{\tau'},\qquad i\in[1,r],
\]
tels que
\begin{enumerate}
\item[a)] pour $i\in[1,r-1]$, le but de $f_i$ soit égal à la source de $f_{i+1}$.

De plus,
\item[b)] $f_1$ transforme source en source, et $f_r$ transforme but en but.
\end{enumerate}

Inversement, quand on se donne des $f_i$ satisfaisant ces deux conditions, ils définissent un morphisme d'intervalles $f$ unique. (Et si la condition b) n'est pas exigée, on trouve une description des morphismes de catégories $I_\tau\longrightarrow I_{\tau'}$.)

\msp{page 33}
On va d'abord regarder un foncteur quelconque
\[
f : \bDelta^{\varepsilon}_i\longrightarrow I_{\tau'}\qquad(\varepsilon\in\{\pm1\})\quad .
\]

\dgrblock{dgrb299}{A functor cannot cross a turn in one run}
\begin{dcomment}{exposition}{p32}{A functor cannot cross a turn in one run}
Grothendieck begins the structure theorem by cutting an interval where its direction changes. Each piece is a chain, and the image of a chain remains comparable even when its length changes. This elementary observation is what allows a map of long zigzags to be described by a finite amount of ordered data.

For a nonempty word $\tau$ let
\[
 0=s_0<s_1<\cdots<s_r=n
\]
be its critical vertices, and let $K_i=[s_{i-1},s_i]$ be its maximal constant-sign runs, regarded as full subcategories. The category $I_\tau$ is obtained by gluing the $K_i$ at their common endpoint vertices. To see the universal property directly, choose functors $f_i:K_i\to Y$ agreeing at shared endpoints. They determine the image of every generating transition of $I_\tau$. Each nonidentity arrow of $I_\tau$ is the unique composite of the transitions along a single run, so its image is forced to be the corresponding composite in $Y$. The only composition equations are associativity along those runs, and these hold in $Y$. Thus the $f_i$ extend uniquely to a functor. For $Y=I_{\upsilon}$, imposing the values of $f_1(s_0)$ and $f_r(s_r)$ additionally imposes endpoint preservation, exactly as on page~32.

Now let $K$ be a critical run of the source and suppose $f|K$ is nonconstant. The image vertices are pairwise comparable in $I_{\upsilon}$: any two source vertices are joined by an arrow in one direction, whose image joins their images. Let $a<b$ be the smallest and largest distinct image vertices in the external order. Their comparability means that all target transitions between $a$ and $b$ have one sign. They therefore lie in a unique maximal run $J'$, and so do all the image vertices between them. This proves both the existence and uniqueness in Proposition~7.

If the source run has sign $\varepsilon$ and $J'$ has sign $\eta$, the object function of $f|K$ is nondecreasing when $\varepsilon=\eta$ and nonincreasing when $\varepsilon=-\eta$. A nonconstant function which is nondecreasing in the external order cannot be of both sorts. Hence an increasing functor sends every nonconstant source run into a target run of the same sign. This proves the sign assertion as well; no endpoint hypothesis has been used.

The preceding paragraph of the source overstates the constant case. A vertex strictly inside a long target run belongs to one critical interval; an interior \emph{critical} vertex belongs to two; an endpoint of a positive-length interval belongs to one. For instance, the constant map with value $1$ in $I_{++}=[2]$ factors through its sole critical interval, although $1$ is not an endpoint. This is a printed slip, retained in the text, and it does not affect Proposition~7 because that proposition assumes nonconstancy. There is also a harmless zero-length exception to the footnote's comparable-subset lemma: $I_e$ has one vertex and no critical runs. A nonconstant map cannot have this target, so the proposition itself remains valid without an added exception.
\end{dcomment}


\printedpage{22}
% head: A. Grothendieck : Dérivateurs, Ch. VII, Édition des Universités de Montpellier II et Paris VII
% foot: Version 22 décembre 2001
Supposons pour fixer les idées que $\varepsilon=+1$, i.e. que le premier membre soit un intervalle critique direct \ldots{} On trouve que $f$ se factorise par un sous-intervalle critique de $I_{\tau'}$, soit $\bDelta^{\varepsilon'}_j$. Il est donc défini par une application croissante de $\bDelta^{\varepsilon}_i$ dans $\bDelta^{\varepsilon'}_j$ (si $\varepsilon$, $\varepsilon'$ \texttt{[sont]} de même signe), une application décroissante dans le sens inverse (\footnote[11]{\textbf{Lemme.} \textit{Si $I_{\tau'}$ est un intervalle, toute partie $J$ de $I_{\tau'}$ telle que pour deux éléments de $J$, il existe toujours une flèche de l'un dans l'autre, est contenue dans un intervalle critique (unique si} $\operatorname{card} J\geq 2$\textit{).}}). Mais si on exige que $f$ soit croissant, alors il faut que \textit{ou bien} $f$ soit constant (auquel cas, si sa valeur n'est pas une extrémité de $I_{\tau'}$, il y a exactement \textit{deux} intervalles critiques de $I_{\tau'}$ par lesquels se factorise $f$) \textit{ou bien} $f$ non constant et $\varepsilon=\varepsilon'$ (et en tous cas, si $f$ non constant, alors l'intervalle critique par lequel se factorise $f$ est unique). Pour résumer :

\textbf{Proposition 7.} \textit{Soient $I_\tau$, $I_{\tau'}$ deux intervalles orientés,}
\[
f : I_\tau\longrightarrow I_{\tau'}
\]
\textit{un foncteur. Alors pour tout intervalle critique $J\subseteq I_\tau$ tel que $f|J$ soit non constant, il existe un unique intervalle critique $J'\subseteq I_{\tau'}$ tel que $f|J$ se factorise par $J'$, i.e. tel que $f(J)\subseteq J'$. Si $f$ est croissant, alors $J$ et $J'$ sont de même signe,}
\msp{page 34}
\textit{i.e. tous deux directs, ou tous deux rétrogrades.}

Ainsi, si on regarde la suite $s_0,s_1,\ldots,s_n$ des sommets critiques de $\tau$, alors
\[
t_0=f(s_0),\ t_1=f(s_1),\ \ldots\ ,\ t_n=f(s_n)
\]
forment une suite de sommets de $I_{\tau'}$, telle que deux consécutifs soient reliés par une flèche de $I_{\tau'}$. Donc la sous-catégorie pleine de $I_{\tau'}$ définie par ces $t_i$ est connexe, c'est donc un sous-intervalle. Si $f$ est croissant, ce sous-intervalle est $[t_0,t_n]$. Si $f$ est un morphisme d'intervalles, il contient tous les sommets critiques de $I_{\tau'}$. Dans tous les cas, prenant $t_i$, $t_j$ $(i<j)$ et appliquant ce qui précède à $f|J$, où $J$ est le sous-intervalle de $I$ d'origine $s_i$, d'extrémité $s_j$, on voit que tout sommet critique de $I_{\tau'}$ compris entre $t_i$ et $t_j$ est parmi les $t_k$ $(i\leq k\leq j)$. Pour résumer :

\textbf{Corollaire.} \textit{Soient $s_i\leq s_j$ deux sommets critiques de $I_{\tau'}$. Alors tout sommet critique de $I_{\tau'}$ compris entre $f(s_i)=t_i$ et $f(s_j)=t_j$, est l'image d'un sommet critique $s_k$ compris entre $s_i$ et $s_j$ (i.e. $i\leq k\leq j$). En particulier, si $f$ est un morphisme d'intervalles,} % [sic] sommets critiques de I_{\tau'} (first occurrence) -> de I_\tau, the s_i are the critical vertices of \tau
\[
f(I_{\tau\,\red})\supseteq I_{\tau'\,\red}
\]
\textit{(où $I_{\tau\,\red}$ désigne l'ensemble des sommets critiques de $I_\tau$ ou aussi le sous-intervalle réduit qu'il définit).}

\dgrblock{dgrb300}{Critical vertices are crossed, not necessarily preserved}
\begin{dcomment}{exposition}{p34:sous-intervalle réduit qu'il définit}{Critical vertices are crossed, not necessarily preserved}
The next step records which turns in the target must be encountered by a map. This is weaker than preservation of all critical vertices, and the distinction is needed when the target is not reduced. It also prevents an ordinary interval of integers from being confused with a connected full subcategory of an interval category.

Let $s_0,\ldots,s_r$ be the critical vertices of the source and put $t_i=f(s_i)$. Consecutive $t_i$ are comparable, since they are images of the endpoints of a source run. Thus the full subcategory on the \emph{set} $\{t_0,\ldots,t_r\}$ is connected. Suppose $t_i$ and $t_j$ lie on opposite sides of an interior critical vertex $z$ of the target. If no $t_k$, $i\leq k\leq j$, were equal to $z$, some consecutive pair in this finite sequence would lie strictly on opposite sides of $z$. Such a pair cannot be comparable: the two transitions adjoining an interior critical vertex have opposite signs. This contradiction proves that $z=t_k$ for some $k$ in the indicated range. The argument permits backtracking and repetitions; it does not require monotonicity.

If $f$ preserves the two endpoints, take $i=0$ and $j=r$. Every target critical vertex is consequently the image of a source critical vertex. This is the content of the corollary, whose first occurrence of $I_{\tau'}$ should be read as $I_\tau$ when it describes the $s_i$. The stronger inclusion in the other direction is false. In the displayed example $I_{+-+}\to I_{++}$, the object function is
\[
 (0,1,2,3)\longmapsto(0,1,1,2).
\]
All four source vertices are critical. Their image contains the two target endpoints and also its noncritical vertex $1$; the middle source transition collapses. The example therefore gives exactly the one-sided containment asserted in the text.

There is a second printed overstatement just before the corollary. Even for increasing $f$, the full subcategory on the $t_i$ need not have \emph{every} integer vertex between $t_0$ and $t_r$. Take the identity of $I_{++}$. Its source critical vertices are $0,2$, so their image is $\{0,2\}$, a full subcategory isomorphic to $[1]$, whereas the ambient interval $[0,2]$ has three vertices. The correct statement is that the image subcategory has least and greatest vertices $t_0,t_r$, its order-convex hull is $[t_0,t_r]$, and it contains every ambient critical vertex in that hull. This is all that the argument uses.

One can now see why the image of $\Cr(\tau,f)\to\Cr(\upsilon)$ is an interval in the ordered set of target runs. Join successive numbers $t_i$ by real line segments, obtaining a continuous walk in $[0,m]$. Each nonconstant segment is contained in the unique target run attached to $K_i$. If the walk reaches two distinct target runs, it must traverse every intervening run, and each such traversal occurs on at least one nonconstant segment. Conversely, a target run not encountered by the walk cannot be attached to any $K_i$. The same argument shows that every target run meeting the segment between $f(0)$ and $f(n)$ in a positive-length segment is encountered. In particular, endpoint preservation makes the run map surjective. This completes the proofs of Corollary~2(a)--(c), including the parts for maps which backtrack.
\end{dcomment}


\msp{page 35}
On fera attention que même si $f$ est croissant, on \textit{n'a pas} en général $f(I_{\tau\,\red})\subseteq I_{\tau'\,\red}$, i.e. $f$ ne transforme pas forcément sommets critiques en sommets critiques, Exemple minimal :
\printedpage{23}
% head: A. Grothendieck : Dérivateurs, Ch. VII, Édition des Universités de Montpellier II et Paris VII
% foot: Version 22 décembre 2001

\noindent
\begin{minipage}[c]{0.6\textwidth}
\[
\begin{tikzcd}
\Phi_3^+\ : & s_0 \arrow[r] \arrow[d, mapsto] & s_1 \arrow[d, mapsto] & s_2 \arrow[l] \arrow[r] \arrow[dl, mapsto] & s_3 \arrow[dl, mapsto] \\
\bDelta_2^+\ : & {} \arrow[r] & {} \arrow[r] & {} &
\end{tikzcd}
\]
\end{minipage}\hfill
\begin{minipage}[c]{0.35\textwidth}
$\Phi_3^+$ est réduite, les deux sommets critiques transformés en le sommet non critique 1.
\end{minipage}

Désignons par $\mathrm{Cr}(\tau)$ l'ensembles des intervalles critiques de $I_\tau$, en correspondance biunivoque (respectant l'ordre total) avec l'ensemble des flèches de transition de $I_{\tau\,\red}$, % [sic] l'ensembles -> l'ensemble, spurious plural s
\[
\mathrm{Cr}(\tau)\simeq\mathrm{Fl}\,\mathrm{tr}\,(I_{\tau\,\red}).
\]
Soit, pour un foncteur
\[
f : I_\tau\longrightarrow I_{\tau'},
\]
$\mathrm{Cr}(\tau,f)$ l'ensemble des intervalles critiques de $I_\tau$ sur lesquels $f$ ne soit pas constant. On a alors une application
\[
\mathrm{Cr}(\tau,f)\longrightarrow\mathrm{Cr}(\tau') \tag*{$(*)$} % label (*) printed at the left margin
\]
induite par $f$ via la proposition 7, associant à tout $J\in\mathrm{Cr}(\tau,f)$ l'unique $J'\in\mathrm{Cr}(\tau')$ tel que $f(J)\subseteq J'$. On a de plus

\textbf{Corollaire 2.}
\begin{enumerate}
\item[\textit{a)}] \textit{L'image de $(*)$ est un segment de l'ensemble totalement ordonné} $\mathrm{Cr}(\tau')$\textit{, i.e. si deux intervalles critiques $J'$, $J''$ dans $I'_{\tau'}$ sont dans l'image, il en est de même de tous} % [sic] I'_{\tau'} -> I_{\tau'}, spurious prime
\msp{page 36}
\textit{ceux qui sont compris entre eux.}
\item[\textit{b)}] \textit{Cette image contient aussi l'intervalle critique de $I_{\tau'}$ (il y en a au plus un) qui contient $f(0_\tau)$ (resp. $f(1_\tau)$) et dont l'intersection avec le segment $[f(0_\tau),f(1_\tau)]$ est de cardinal $\geq 2$, ainsi que tous les intervalles critiques contenus dans ledit segment.}
\item[\textit{c)}] \textit{Si $f$ est un morphisme d'intervalles orientés, alors l'application induite}
\[
\mathrm{Cr}(f):\mathrm{Cr}(\tau,f)\longrightarrow\mathrm{Cr}(\tau')
\]
\textit{est surjective. De façon plus précise, soit $J'\subseteq$} $\mathrm{Cr}(\tau')$\textit{, et soit $J=f^{-1}(J')$. Alors} $\mathrm{Cr}(f)^{-1}(\{J'\})$ \textit{est formé des intervalles critiques dans} $\mathrm{Cr}(\tau,f)$ contenus dans $J$. \textit{On a de plus,} si $f$ est croissant, % [sic] J' \subseteq Cr(\tau') -> J' \in Cr(\tau'), J' is an element of Cr(\tau')
\[
\mathrm{long}(J')=\sum_{K\in\mathrm{Cr}(f)^{-1}(\{J'\})}\pi_{\mathrm{n}}(K,f),
\]
\textit{où $\pi_{\mathrm{n}}(K,f)$ est la `partie' (numérique) de $f|K$, i.e. la longueur de l'intervalle $[f(0_K),f(1_K)]$. (}\textbf{NB} \textit{Si $K$ est un intervalle critique de $I$ tel que $K\notin Cr(\tau,f)$, i.e. $f|K$ est constant, on posera $\pi_{\mathrm{n}}(K,f)=0$. Donc la formule précédente se lit aussi}
\[
\mathrm{long}(J')=\sum\limits_{\underbrace{K\in\mathrm{Cr}(I_\tau),\ K\subseteq J}_{\text{(i.e. }K\,\in\,\mathrm{Cr}(f^{-1}(J'))\text{)}}}\pi_{\mathrm{n}}(K,f).\ \textit{)}
\]
\end{enumerate}

% ---- from Der-ChVII-456_p026-030.tex
\printedpage{24}
% head: A. Grothendieck : Dérivateurs, Ch. VII, Édition des Universités de Montpellier II et Paris VII

\textbf{Corollaire 3.} \textit{Si $f$ est un morphisme} croissant \textit{d'intervalles orientés, alors pour tout $J' \in \Cr(\tau')$, on a}
\[
(0 \leq)\ \card \Cr(f)^{-1}(\{J\}) \longrightarrow \lng(J'),
\] % [sic] \longrightarrow -> \leq, the next sentence speaks of equality; and {J} -> {J'}, as in Cr(f)^{-1}({J'}) below
\msp{page 37}
\textit{l'égalité n'étant obtenue que si les parties de $f$ sur les intervalles critiques de $I_\tau$ qui sont au dessus de $J'$ sont toutes égales à $0$ ou $1$ (donc si et seulement si sur les $J \in \Cr(f)^{-1}(\{J'\})$, elles sont égales à $1$).}

\dgrblock{dgrb301}{The length budget over a target run}
\begin{dcomment}{exposition}{p36}{The length budget over a target run}
Monotonicity turns the preceding coverage statement into an additive formula. A map may pause, but it cannot spend the same portion of the target twice. This is the numerical mechanism behind Corollaries~3--5.

Let $f:I_\tau\to I_\upsilon$ be increasing and preserve endpoints. For a source critical run $K_i=[s_{i-1},s_i]$, write
\[
 d_i=f(s_i)-f(s_{i-1})\geq0.
\]
When $d_i>0$, Proposition~7 associates to $K_i$ the unique target run containing its image; this is $\Cr(f)(K_i)$. Fix a target run $J'=[a,b]$. The nonconstant segments $[f(s_{i-1}),f(s_i)]$ attached to $J'$ cover $[a,b]$, meet only at consecutive endpoints, and are encountered in their natural order. Coverage follows because $f(0)=0$, $f(n)=m$ and every intervening critical vertex is attained; disjointness of interiors follows from the monotonicity of the sequence $f(s_i)$. Their lengths therefore add to $b-a$:
\[
 b-a=\sum_{\substack{K_i\in\Cr(\tau,f)\\\Cr(f)(K_i)=J'}}d_i.
\]
This is the formula for $\pi_{\mathrm n}(K_i,f)$ in the source. It is valid even if $f$ skips noncritical target vertices, because length measures the whole segment spanned by the image, not the number of image vertices.

Every summand displayed is a positive integer. Consequently, if $k_{J'}=|\Cr(f)^{-1}(J')|$, then
\[
 1\leq k_{J'}\leq b-a,
\]
and equality on the right holds precisely when every positive $d_i$ in the sum is $1$. The printed arrow in Corollary~3 is a typesetting slip for $\leq$, and its unprimed $J$ in the singleton should be $J'$. The subsequent discussion of equality gives the intended reading independently of our proof.

For a first calculation take $I_{+-+}\to I_{++}$ with vertex function $(0,1,1,2)$. Its three source runs have lengths of image $1,0,1$; the single target run has length $2$. Thus $k_{J'}=2$ and the upper bound is attained. For the identity $I_{++}\to I_{++}$, the sole source run has image length $2$, so $k_{J'}=1<2$. Both maps are increasing and preserve endpoints. What is counted is the number of nonconstant \emph{runs}, not the number of nonidentity generating arrows.

The alternative summation on page~36 may harmlessly include every global source run contained in $f^{-1}(J')$: the added terms have zero length. Its parenthetical identification of these runs with all critical runs of the preimage is, however, too literal. A preimage can contain a truncated piece of a neighboring global run on which $f$ is constant. For example, take
\[
 f:I_{++--}\longrightarrow I_{+-},\qquad
 (0,1,2,3,4)\longmapsto(0,1,1,1,2).
\]
For $J'=[0,1]$, its preimage has vertices $0,1,2,3$ and critical runs $[0,2]$ and $[2,3]$. Only $[0,2]$ is a full global source run contained in the preimage; $[2,3]$ is a proper part of the global run $[2,4]$. The extra piece is constant with value $1$, hence contributes zero. The equality of sums survives, but equality of the indexing sets does not. One should either keep the global-run indexing or explicitly allow such zero-contribution boundary pieces.
\end{dcomment}


\textbf{Corollaire 4.} \textit{Supposons de plus $I'$} \texttt{[}$= I_{\tau'}$\texttt{]} \textit{réduit. Alors}
\[
\mathit{Cr}(f) : \mathit{Cr}(\tau, f) \longrightarrow \mathit{Cr}(\tau')\ (\simeq \mathrm{Fl}\,\mathrm{tr}(\tau'))
\]
\textit{est} bijective.

Donc dans ce cas, $\tau'$ réduit, on obtient une application en sens inverse
\[
\begin{tikzcd}[column sep=small]
\Tr(\tau') \simeq \Cr(\tau') \arrow[rr, "\Cr^*(f)"] \arrow[dr, "\Cr(f)^{-1}"'] & & \Cr(\tau) \\
 & \Cr(\tau, f), \arrow[ur, hook, "\text{inclusion}"'] &
\end{tikzcd}
\]
application qui transforme l'intervalle critique $J'$ $(= \{i - 1, i\})$ en l'\textit{unique} intervalle critique $J$ de $I_\tau$ tel que
\[
f(J) = J'.
\]

\msp{page 38}
Mais $J \simeq \bD_i^{\varepsilon}$, et on a donc un morphisme \textit{surjectif}
\[
f_0 : \bD_i^{\varepsilon} \longrightarrow J' \simeq \bD_1^{\varepsilon}.
\]
On voit tout de suite qu'il existe alors une \textit{unique} flèche de transition $u$ de $\bD_i^{\varepsilon}$, telle que $f(u) \neq \mathrm{id}$. Donc on trouve le corollaire suivant, plus précis que le corollaire 4 :

\textbf{Corollaire 5.} \textit{Soit $f : I_\tau \longrightarrow I_{\tau'}$ un morphisme croissant d'intervalles orientés, avec $I_{\tau'}$ réduit, et soit $\Tr(\tau, f)$ l'ensemble des flèches de transition $\underline{u}_i$ de $I_\tau$ telle que $f(\underline{u}_i) \neq \mathrm{id}$, i.e. telles que $\pi(u_i, f) = 1$. Alors $f$ induit une application} % [sic] telle -> telles, agreement with "flèches" (cf. "telles que" just after)
\[
\Tr(\tau, f) \longrightarrow \Tr(\tau')
\]
\textit{qui est bijective, et (bien sûr) croissante, donc un isomorphisme d'ensembles ordonnés.}

Ainsi l'application $\Tr(\tau') \longrightarrow \Cr(\tau, f)$ de tantôt se précise par une application bijective
\[
\Tr(f)^* : \Tr(\tau') \longrightarrow \Tr(\tau, f),
\]
reliée à $\Cr(f)^*$ plus haut par le triangle commutatif
\[
\begin{tikzcd}
 & \Tr(\tau, f) \arrow[dd, "\wr\ \alpha"] \\
\Tr(\tau') \arrow[ur, "\Tr(f)^*"] \arrow[dr, "\Cr(f)^*"'] & \\
 & \Cr(\tau, f),
\end{tikzcd}
\]
% foot: Version 22 décembre 2001
\printedpage{25}
% head: A. Grothendieck : Dérivateurs, Ch. VII, Édition des Universités de Montpellier II et Paris VII
\msp{page 39}
l'application $\alpha$ associant à toute $u_i \in \Tr(\tau, f)$ l'unique intervalle critique qui ‘contienne’ $u_i$ (i.e. tel que $u_i$ soit une flèche dudit intervalle, ou encore, fasse partie de la ‘chaîne critique’ associée \dots).

\dgrblock{dgrb302}{A reduced target remembers the surviving transitions}
\begin{dcomment}{examples}{p38}{A reduced target remembers the surviving transitions}
Grothendieck next specializes the length formula to a target with no redundant vertices. Every target run then has length one. There is exactly one nonconstant source critical run over it, and that source run can cross its two target vertices only once. This converts a numerical constraint into the combinatorial invariant used in Proposition~8.

More explicitly, on a source run $K=[s,t]$ mapping onto the two vertices $a,a+1$, the nondecreasing vertex function has the form
\[
 f(s)=\cdots=f(i-1)=a,
 \qquad
 f(i)=\cdots=f(t)=a+1
\]
for exactly one $i$ with $s<i\leq t$. Hence exactly one transition in $K$ has nonidentity image. The run and its target have the same sign, so this transition maps to the target's unique transition in its categorical direction, including when that direction is backward. Performing this construction for all target runs gives the bijection of Corollary~5. It is increasing because the source runs and their crossings occur in the same external order as the target runs. The triangle with $\alpha$ commutes because $\alpha$ forgets which transition was selected and remembers its containing source run.

Take $\tau=++--++$ and $\upsilon=+-+$. Choose transitions $2,4,6$ in the source. The associated vertex map is
\[
 (0,1,2,3,4,5,6)\longmapsto(0,0,1,1,2,2,3).
\]
The first source run maps across the first forward target edge at transition~2; the second maps across the backward target edge at transition~4; the third crosses the final forward edge at transition~6. The other three transitions become identities. If a path of type $\upsilon$ in $X$ is
\[
 x_0\xrightarrow{a}x_1\xleftarrow{b}x_2\xrightarrow{c}x_3,
\]
restriction along this map is the path
\[
 x_0\xrightarrow{1}x_0\xrightarrow{a}x_1
 \xleftarrow{1}x_1\xleftarrow{b}x_2
 \xrightarrow{1}x_2\xrightarrow{c}x_3.
\]
In each of the three source runs either of its two transitions could instead be the surviving one. These choices give $2^3=8$ different maps to this reduced target, all explicitly determined by their three crossing positions. The general theorem on the next pages will allow arbitrary sign words and express this determination as a signed-subsequence bijection.

When the target is $e$, there are no target runs and no nondegenerate transitions. The unique endpoint-preserving map $I_\tau\to I_e$ collapses every transition. Both sides of the bijection in Corollary~5 are empty, so the assertion includes this case. When the source is $e$ and the target has positive length, no endpoint-preserving map exists. These two degenerate cases explain why the later subset description needs the empty word as well as the alternating nonempty words.
\end{dcomment}


On a mieux encore :

\textbf{Proposition 8.} \textit{Soient $\tau, \tau' \in \mathrm{Ty}$, $\tau'$} réduit. \textit{Pour tout morphisme croissant d'intervalles orientés}
\[
f : I_\tau \longrightarrow I_{\tau'}
\]
\textit{(on va dire simplement que $f$ est une flèche dans $\uTy$, et écrire $f \in \Hom_{\uTy}(I_\tau, I'_\tau)$), associons lui la partie $\Tr(\tau, f)$ des flèches de transition de $I_\tau$ non dégénérées pour $f$. Ces flèches forment un ensemble totalement ordonné affecté de signes $+$ ou $-$ (caractère direct ou rétrograde des flèches), et $f$ induit une bijection d'ensembles totalement ordonnés} % [sic] I'_\tau -> I_{\tau'}, the target of f is I_{\tau'}
\[
\underbrace{E_f}_{\subseteq \Tr(\tau)} = \Tr(\tau, f) \simeq \Tr(\tau'),
\]
\textit{et même un isomorphisme d'ensembles totalement ordonnés ‘affectés de signes’. Ainsi, à toute $f$ on associe une partie $E_f$ de $\Tr(\tau)$, et telle que son ‘type’ soit égal à $\tau'$. Désignons par}
\[
\mathfrak{P}_{\tau'}(\Tr(\tau)) \subseteq \mathfrak{P}(\Tr(\tau))
\]
\textit{l'ensemble des parties $E$ de $\Tr(\tau)$ dont le type soit $\tau'$.}

\msp{page 40}
\textit{Alors l'application}
\[
f \longmapsto \Tr(\tau, f) : \Hom_{\uTy}(\tau, \tau') \longrightarrow \mathfrak{P}_{\tau'}(\Tr(\tau))
\]
\textit{est} bijective \textit{: se donner un morphisme $f : \tau \longrightarrow \tau'$ dans $\uTy$,} revient au même \textit{que de se donner une suite strictement croissante de flèches de transition}
\[
\underline{u}_{i_1},\ \underline{u}_{i_2},\ \dots\ ,\underline{u}_{i_{n'}} \qquad (n' = \lng \tau')
\]
\textit{dans $I_\tau$ $(i_1 < i_2 < \cdots < i_{n'})$, telle que la suite des signatures correspondante}
\[
(\varepsilon_{i_1}, \varepsilon_{i_2}, \dots, \varepsilon_{i_{n'}}) = \tau'
\]
\textit{soit égale à $\tau'$.}

\dgrblock{dgrb303}{Choosing the transitions that survive}
\begin{dcomment}{exposition}{p40:se donner une suite strictement croissante}{Choosing the transitions that survive}
The reduced target allows Grothendieck to replace a functor by a selection of edges. This is the first form of the normal form that will make concatenation manageable: each edge of the target has one identifiable ancestor in the source. The hypothesis that the target is reduced is doing exactly the work needed here.

Write the source signature as $\tau=(\varepsilon_1,\ldots,\varepsilon_n)$ and the reduced target as $\tau'=(\varepsilon'_1,\ldots,\varepsilon'_m)$. An increasing endpoint-preserving functor has a nondecreasing object map
\[
 f:[0,n]\longrightarrow[0,m],\qquad f(0)=0,\quad f(n)=m.
\]
The image of the $i$th transition is either an identity or a nonidentity arrow joining $f(i-1)$ to $f(i)$. In a reduced interval every nonidentity arrow is a transition: two successive transitions have opposite signs and cannot compose. Consequently
\[
 f(i)-f(i-1)\in\{0,1\}.
\]
Exactly $m$ differences are $1$. If their indices are $i_1<\cdots<i_m$, preservation of arrow direction gives $\varepsilon_{i_j}=\varepsilon'_j$. Conversely such a signed subsequence defines
\[
 f(k)=\#\{j\mid i_j\leq k\}.
\]
On a selected transition this map uses the corresponding target transition; on every other transition it uses an identity. These choices extend uniquely to a functor because the source is the free category on its oriented linear graph. They preserve the two endpoints and recover the selected subsequence. This proves Proposition~8, including its empty-target case: the empty subsequence gives the unique functor to the point.

For example, from the word $++-+$ to the reduced word $+-+$ there are exactly two morphisms. The negative transition must be the third source edge, the last positive transition must be the fourth, and the first positive transition may be either the first or the second. Their object maps are respectively
\[
 (0,1,1,2,3),\qquad (0,0,1,2,3).
\]
The two paths retain different occurrences of the first positive step. Thus the type records positions, not merely the numbers of positive and negative transitions.

For a single-edge target the same argument gives the stated representability of positive and negative transitions. A positive edge at position $i$ determines the map to $\Delta_+$ which is $0$ before that edge and $1$ after it; a negative edge determines the analogous map to $\Delta_-$. Naturality remains to be explained because an edge can be stretched over several target edges. The next construction supplies exactly that explanation. When the target is not reduced, the selection alone is insufficient: the unique morphism $\Delta_+\to\Delta_+^2$ stretches one edge across two, so the size of each surviving image must also be retained.
\end{dcomment}


\textbf{Corollaire 1.} \textit{Pour tout $\tau$, on a des bijections canoniques}
\[
\left\{
\begin{array}{lcl}
\Tr^+(\tau) & \simeq & \Hom_{\uTy}(\tau, \Delta^+) \\
\Tr^-(\tau) & \simeq & \Hom_{\uTy}(\tau, \Delta^-),
\end{array}
\right.
\]
fonctorielles en $\tau$ \textit{(variant dans} $\uTy$\textit{).}
% foot: Version 22 décembre 2001

\printedpage{26}
% head: A. Grothendieck : Dérivateurs, Ch. VII, Édition des Universités de Montpellier II et Paris VII

Pour donner un sens au corollaire, il faut définir
\[
\tau \longmapsto \Tr^+(\tau), \qquad \tau \longmapsto \Tr^-(\tau)
\]
comme \textit{foncteurs} contravariant en $\tau$, donc pour un morphisme
\[
f : I_{\tau_2} \longrightarrow I_{\tau_1}
\]
dans $\uTy$, définir des applications
\msp{page 41}
\[
\begin{array}{rcl}
\Tr^+(f)^* : \Tr^+(\tau_1) & \longrightarrow & \Tr^+(\tau_2) \\
\Tr^-(f)^* : \Tr^-(\tau_1) & \longrightarrow & \Tr^-(\tau_2),
\end{array}
\]
en d'autres termes, une application
\[
\Tr(f)^* : \Tr(\tau_1) \quad \longrightarrow \quad \Tr(\tau_2)
\]
‘compatible avec les signatures’. Dans le corollaire 5 plus haut, on a défini une telle application seulement pour $\tau'$ réduit \texttt{[}$\tau_1 = \tau'$\texttt{,} $\tau_2 = \tau$\texttt{]} (\footnote[12]{Ce qui suit ne fait pas intervenir les signatures - c'est un sorite sur les applications croissantes entre ensembles finis totalement ordonnés.}). Pour une flèche $u$ de $I_\tau$, appelons \textit{partie ensembliste} de $f$ sur $u$, où de $u$ pour $f$, la partie ensembliste de la flèche $f(u)$ dans $I_{\tau'}$, i.e. le segment $[f(i-1), f(i)] \subseteq I_{\tau'}$ (ici $u = \underline{u}_i$). Soit % [sic] où -> ou, alternative "de f sur u, ou de u pour f"
\[
\underline{u}_{i_1},\ \dots\ ,\underline{u}_{i_N} \qquad \text{ou } (v_1, v_2, \dots, v_n)
\] % [sic] v_n -> v_N, the sequence has N elements (cf. v_N below)
la suite croissante des éléments de $\Tr(\tau, f)$, on voit alors par récurrence sur $i$ que pour $i \in [1, N-1]$, on a
\[
\begin{array}{rcl}
f(\mathrm{but}(v_i)) & = & f(\mathrm{source}(v_{i+1}))\ \ (= t_i), \\
f(\mathrm{source}(v_1)) & = & 0_{\tau'}, \\
f(\mathrm{but}(v_N)) & = & 1_{\tau'}.
\end{array}
\]
Considérons la suite
\[
\{\overbrace{0_{\tau'}}^{= t_0}, t_1, t_2, \dots, t_{N-1}, \overbrace{1_{\tau'}}^{= t_N}\},
\]
elle est définie par le sous-ensemble correspondant de $I_{\tau'}$, qu'on appelera l'ensemble des \textit{sommets critiques de $I_{\tau'}$ pour $f$.} C'est une suite contenant $0_{\tau'}$, $1_{\tau'}$, et telle que deux éléments consécutifs sont reliés dans $I_{\tau'}$ par une flèche, ce qui équivaut à dire qu'elle contient les sommets critiques (intrinsèques) de $I_{\tau'}$. D'autre part, pour tout intervalle de transition $[i-1, i]$ de $I_{\tau'}$ $(i \in [1, n'])$, % [sic] appelera -> appellera, French spelling
% foot: Version 22 décembre 2001
\printedpage{27}
% head: A. Grothendieck : Dérivateurs, Ch. VII, Édition des Universités de Montpellier II et Paris VII
\msp{page 42}
il existe un et un seul intervalle de transition $[t_{j-1}, t_j]$ pour cette suite $(j \in [1, N])$, telle qu'on ait $[i-1, i] = [t_{j-1}, t_j]$. C'est dire qu'il existe une et une seule $v_j \in \Tr(\tau, f)$ telle que la partie de $f$ sur $v_j$ (savoir justement $[j-1, j]$) contienne (celle de) $\underline{u}_i$. Pour résumer :

\textbf{Corollaire 2.} \textit{Soit $f : \tau \longrightarrow \tau'$ dans $\uTy$. Alors pour toute $u \in \Tr(\tau') = \Tr(I_{\tau'})$, il existe une et une seule $v \in \Tr(\tau)$ telle que la partie de $f$ sur $v$ contienne $u$. De plus, les signatures de $u$ et de $v$ sont égales.}

C'est ainsi qu'on obtient donc une application
\[
\Tr(f)^* : \Tr(\tau') \quad \longrightarrow \quad \Tr(\tau) \qquad (\text{se factorisant par } \Tr(\tau, f))
\]
induisant
\[
\begin{array}{rcl}
\Tr^+(f)^* : \Tr^+(\tau') & \longrightarrow & \Tr^+(\tau) \\
\Tr^-(f)^* : \Tr^-(\tau') & \longrightarrow & \Tr^-(\tau).
\end{array}
\]
La transitivité pour un composé
\[
\tau \xrightarrow{f} \tau' \xrightarrow{g} \tau''
\]
ou
\[
I_\tau \xrightarrow{f} I_{\tau'} \xrightarrow{g} I_{\tau''}
\]
est immédiate. Si $u$, $u'$, $u''$ sont des flèches de transition dans $I_\tau$, $I_{\tau'}$, $I_{\tau''}$ respectivement, telles que
\[
\begin{array}{rcll}
\pi(g, u') & \supseteq & \pi(u''), & \qquad \text{i.e. } u' = \Tr(g)^*(u'') \\
\pi(f, u) & \supseteq & \pi(u'), & \qquad \text{i.e. } u = \Tr(f)^*(u'),
\end{array}
\]
alors on a
\[
\pi(gf, u) \ \supseteq\ \pi(u''), \qquad \text{i.e. } u = \Tr(gf)^*(u'').
\]

\msp{page 43}
Si les flèches en question ont comme partie intrinsèque $[i-1, i]$, $[i'-1, i']$, $[i''-1, i'']$ respectivement, l'hypothèse est
\[
\begin{array}{rclrcl}
g(i'-1) & \leq & i''-1, & g(i') & \geq & i'' \\
f(i-1) & \leq & i'-1, & f(i) & \geq & i'
\end{array}
\]
et comme $g$ est croissant, on a déduit
\[
\begin{array}{rcll}
gf(i-1) & \leq & g(i'-1) & (\leq i''-1) \\
gf(i) & \geq & g(i') & (\geq i''),
\end{array}
\]
\begin{flushright}
c.q.f.d.
\end{flushright}
% foot: Version 22 décembre 2001

\dgrblock{dgrb304}{A stretched edge and the containment convention}
\begin{dcomment}{examples}{p42}{A stretched edge and the containment convention}
Grothendieck is assigning every target transition to the unique source transition whose image crosses it. One small example fixes both the meaning of this assignment and two printed slips immediately before Corollary~2.

Let $f:\Delta_+\to\Delta_+^2$ have object map $f(0)=0$, $f(1)=2$. There is one surviving source transition, so $N=1$, $t_0=0$ and $t_1=2$. The two target transitions have underlying vertex segments $[0,1]$ and $[1,2]$. Both are contained in $[t_0,t_1]=[0,2]$, and neither equals it. Thus the printed relation
\[
 [i-1,i]=[t_{j-1},t_j]
\]
must be read mathematically as $[i-1,i]\subseteq[t_{j-1},t_j]$. The parenthesis identifying the part of $f$ on $v_j$ as $[j-1,j]$ has the same problem: that part is $[t_{j-1},t_j]$. Both readings occur in the edition; they are retained in the source and repaired here. Corollary~2 itself uses containment and is correct.

The intervals $[t_{j-1},t_j]$ may meet at a vertex, but their sets of transitions are disjoint. Since $t_0=0$, $t_N=m$ and the $t_j$ are strictly increasing, these transition sets partition $\{1,\ldots,m\}$. A transition cannot lie in two of them: its two endpoints would then have to belong to their intersection, which contains at most one vertex. This proves uniqueness without assuming that the target is reduced. Functoriality forces every edge in the block assigned to $v_j$ to have the sign of $v_j$.

The words ``source'' and ``but'' in the preceding display concerning the $v_j$ must here be read as the first and last vertices in the traversal order. On a retrograde transition they are the reverse of its categorical domain and codomain. Writing $f(i_j-1)$ and $f(i_j)$ avoids this ambiguity. If the target is the point there are no target transitions and the pullback map of transition sets is the unique map from the empty set; no first or last surviving transition should be invoked.
\end{dcomment}


\printedpage{28}
% head: A. Grothendieck : Dérivateurs, Ch. VII, Édition des Universités de Montpellier II et Paris VII

Ainsi, le corollaire 1 plus haut obtient un sens précis et il est démontré du même coup, vu que les correspondances du corollaire 1 sont données par
\[
\underbrace{f}_{\in \Hom_{\uTy}(\tau, \Delta^{\varepsilon})} \qquad \longmapsto \Tr(f)^*(\underline{u}),
\]
où $\underline{u}$ est l'unique flèche de transition de $\Delta^{\varepsilon}$.

\dgrblock{dgrb305}{Edges reverse the direction of a map}
\begin{dcomment}{exposition}{p43:Ainsi, le corollaire 1 plus haut}{Edges reverse the direction of a map}
The contravariant transition map contains more information than its notation initially suggests. It determines the whole functor, and gives a particularly economical modern description of the final category of types. This observation explains why the later calculations require only ordered lists and nonnegative integers.

Let $\mathcal O_{\pm}$ have as objects finite linearly ordered sets equipped with a sign at each element, including the empty set, and as arrows the nondecreasing maps which preserve signs. We use the skeletal version with underlying sets $[1,n]$. For $f:\tau\to\tau'$ in $\uTy$, where $\tau$ and $\tau'$ have lengths $n$ and $m$, define
\[
 q_f:[1,m]\longrightarrow[1,n],\qquad
 q_f(j)=\min\{i\mid j\leq f(i)\}.
\]
The endpoint conditions ensure that this minimum exists when $m>0$; its defining inequalities are
\[
 f(q_f(j)-1)<j\leq f(q_f(j)).
\]
This is Grothendieck's $\Tr(f)^*$. It is nondecreasing, and it preserves signs because the image of a source transition is a monochromatic string of target transitions of that same sign.

Conversely, suppose $q:[1,m]\to[1,n]$ is nondecreasing and sign-preserving. Put
\[
 F_q(i)=\#q^{-1}([1,i]),\qquad 0\leq i\leq n.
\]
The fibre $q^{-1}(i)$ is a consecutive block, possibly empty, whose target signs are all $\varepsilon_i$. Hence the consecutive vertices $F_q(i-1)$ and $F_q(i)$ are joined in the target by an arrow of the required direction, or coincide when the fibre is empty. Sending the $i$th source transition to this arrow defines a unique functor. Its object map is nondecreasing, begins at $0$, and ends at $m$. The inequalities above show $q_{F_q}=q$, while counting the target edges crossed by the first $i$ source edges gives $F_{q_f}(i)=f(i)$. These constructions are inverse, including the empty cases. There is no map from the point to a nonpoint interval in $\uTy$, just as there is no function from a nonempty edge set to the empty one.

If $g:\tau'\to\tau''$, a target edge of $\tau''$ has a unique ancestor in $\tau'$ and then a unique ancestor in $\tau$. Equivalently, applying the displayed inequalities twice gives
\[
 q_{gf}=q_fq_g.
\]
Identities induce identities. We therefore have an isomorphism of skeletal categories
\[
 \uTy\ \cong\ \mathcal O_{\pm}^{\mathrm{op}}.
\]
The reversal of variance is essential: a functor stretches a source edge into a target block, whereas $q_f$ sends each member of that block back to its source edge. The signs can be forgotten for the construction on ordered sets, as Grothendieck observes, but must be restored to say when the resulting vertex map is a functor between the oriented interval categories.

This identification also establishes the naturality promised in Corollary~1. If an edge of sign $\varepsilon$ represents $h:\tau'\to\Delta_\varepsilon$, the edge representing $hf$ is $q_f$ of the edge representing $h$. It is precisely ordinary precomposition in the representable presheaf $\Hom_{\uTy}(-,\Delta_\varepsilon)$. The underlying single-sign calculation is the familiar ordinal description of the simplicial category; see \citet[Appendix~D.1, pp.~713--715]{AguiarMahajan2010}. The signed formulation follows by retaining the labels throughout the inverse constructions above.
\end{dcomment}


On va utiliser ces sorites pour donner, pour tout $\tau \in \uTy$, une description combinatoire simple des catégories $\uTy/\tau$ et $\tau\backslash\uTy$, ou du moins de leurs ensembles d'objets.

\textbf{1$^o$ Description de} $\Ob\,(\tau\backslash\uTy)$\textbf{.}

À $f \in \Ob(\tau\backslash\uTy)$, i.e.
\[
f : \tau \longrightarrow \tau' \qquad \text{i.e. } f : I_\tau \longrightarrow I_{\tau'}
\]
dans $\uTy$, est associé $E = E_f = \Tr(\tau, f) \subseteq \Tr(\tau)$, l'ensemble des flèches de transition critiques (i.e. non dégénérées) pour $f$. À chaque $u \in E$, associons l'entier $\lng(\pi(f, u))$,
\msp{page 44}
la \textit{partie} de $f$ sur $u$. Alors $(\tau', f)$ est déterminé de façon unique par le couple
\[
(E, \pi), \qquad E \in \mathfrak{P}(\Tr(\tau)), \qquad \underbrace{\pi}_{\text{pondération}} : E \ \longrightarrow\ \underbrace{\mathbf{N}^*}_{\substack{= \text{ensemble des} \\ \text{entiers} \geq 1}} \quad .
\]
Ou si on préfère, par l'application $\pi = \pi_f$
\[
\Tr(\tau) \xrightarrow{\pi} \mathbf{N}, \qquad u \longmapsto \pi(u) = \pi(u, f) = \lng(\pi_{\tau'}(f(u)))
\]
(qui redonne $E$ comme
\[
E = \pi^{-1}(\mathbf{N}^*) \qquad ).
\]
Cette application n'est soumise à aucune restriction,
\[
\Ob\,(\tau\backslash\uTy) \ \xrightarrow[\sim]{\text{bij.}}\ \Hom_{\mathrm{Ens}}(\Tr(\tau), \mathbf{N}).
\]
On récupère $\tau'$ à partir de $\pi$ ainsi. On prend $E = \pi^{-1}(\mathbf{N}^*)$ (flèches de transition non dégénérées pour $\pi$), $E$ est totalement ordonné par l'ordre induit par celui de $\Tr(\tau)$, il est de plus signaturé \texttt{[sic]} par la signaturation \texttt{[sic]} de $\Tr(\tau)$, soit
\[
(\varepsilon_1, n_1),\ \dots\ ,\ (\varepsilon_r, n_r)
\]
la suite des signatures et pondérations, de la suite des éléments de $E$ $(\varepsilon_i \in \{\pm 1\}$, $n_i \in \mathbf{N}^*$ pour $i \in [1, r])$. Alors,
\[
\tau' = \Delta_{\varepsilon_1}^{n_1} \Delta_{\varepsilon_2}^{n_2} \cdots \Delta_{\varepsilon_r}^{n_r}, \qquad r = \card E,
\]
dans le monoïde libre $\mathrm{Ty}$, où on a désigné par $\Delta_+$, $\Delta_-$ (au lieu de $\Delta^+$, $\Delta^-$) les deux générateurs. L'application
\[
I_\tau \ \longrightarrow\ I_{\tau'}
\]
% foot: Version 22 décembre 2001

% ---- from Der-ChVII-456_p031-035.tex
\printedpage{29}
% head: A. Grothendieck : Dérivateurs, Ch. VII, Édition des Universités de Montpellier II et Paris VII
\msp{page 45}est décrite de façon unique en disant qu'elle est \textit{constante} sur chacune des `composantes connexes' (au sens de la numérotation des sommets de $I_\tau$) de l'ensemble $E' = E'(f) \subseteq I_\tau = [0,n]$ des sommets qui sont sources ou buts de flèches de transition dégénérées pour $f$, i.e. dans $\mathrm{Tr}(\tau)\backslash E$. Et sur le $\alpha$-ème élément de $E$ ($\alpha \in [1, r = \operatorname{card} E]$), d'indice $i_\alpha$ dans $\mathrm{Tr}(\tau)$, $f$ induit le foncteur
\[
\boldsymbol{\Delta}_{\varepsilon_\alpha} \simeq \underbrace{J_{i_\alpha}}_{\substack{\text{intervalle}\\ \text{de transition}\\ \text{d'indice } i_\alpha\\ \text{dans } I_\tau}} \longrightarrow I_{\tau'}
\]
donné par \textit{l'unique} \underline{Ty}-\textit{flèche} dans $I_\tau$ de $J_{i_\alpha}$ dans le $\alpha$-ème facteur, de type $\Delta^{n_\alpha}_{\varepsilon_\alpha}$, de $\tau'$, donné par $(*)$
\[
\Delta^{n_\alpha}_{\varepsilon_\alpha} \simeq \text{type de } \boldsymbol{\Delta}^{\varepsilon_\alpha}_{n_\alpha}\ \ \ldots
\]

\dgrblock{dgrb306}{Weight vectors and composition by block sums}
\begin{dcomment}{examples}{p44}{Weight vectors and composition by block sums}
The outgoing classification now has coordinates which are useful for actual calculations. For a fixed source $\tau=(\varepsilon_1,\ldots,\varepsilon_n)$, put
\[
 a_i=f(i)-f(i-1)=\#q_f^{-1}(i),\qquad
 A_i=\sum_{h=1}^{i}a_h,\quad A_0=0.
\]
Every $a_i$ is a nonnegative integer. The target is the word
\[
 \tau'=\varepsilon_1^{a_1}\varepsilon_2^{a_2}\cdots\varepsilon_n^{a_n},
 \qquad f(i)=A_i,
\]
where exponent zero means omission. These equations prove the claimed unrestricted parametrization of $\Ob(\tau\backslash\uTy)$: they construct an outgoing morphism from every weight vector and recover the vector from every outgoing morphism. When adjacent source signs agree, their target blocks may merge into one critical interval, but their individual lengths remain part of the morphism.

Take the source $+-++-$ and the weights $(2,0,1,3,0)$. Its target is $++++++$, and its vertex map is
\[
 (0,2,2,3,6,6).
\]
The pullback on target transitions is $(1,1,3,4,4,4)$. Both negative source edges have been collapsed. The two consecutive positive source edges at positions $3$ and $4$ survive separately, even though their target images lie in the same positive run.

To see composition, give the six target edges the further weights
\[
 (b_1,\ldots,b_6)=(1,0,2,1,0,3).
\]
The second map has vertex values $(0,1,1,3,4,4,7)$ and target $+^7$. The composite has weights $(1,0,2,4,0)$, since its $i$th weight is
\[
 c_i=\sum_{j=A_{i-1}+1}^{A_i}b_j.
\]
Thus its vertex map is $(0,1,1,3,7,7)$, which is also obtained by direct substitution into the two displayed object maps. An empty block contributes zero; a block of length greater than one contributes a sum. This is the numerical content of composition in Section~6.

The formula gives a useful direct verification of associativity. Given a third weight vector, first replacing every $b_j$ by the sum over its descendants and then summing over the $a_i$-block is the same finite sum as first merging the $b_j$-blocks and then summing over the merged block. The order and boundaries of the summands are unchanged. No cancellation, choice of representative, or equivalence relation enters this calculation.
\end{dcomment}


\subsection*{2$^o$ Détermination de $\mathrm{Ob}\,(\uTy/\tau')$.}

Si $(\tau, f) \in \mathrm{Ob}\,(\uTy/\tau')$, donc
\[
f : I_\tau \longrightarrow I_{\tau'},
\]
on lui associe tout d'abord la partie
\[
\mathrm{Cr}(\tau', f) = E'_f = E' \subseteq I_{\tau'}
\]
formée des sommets (dits `critiques pour $f$', cf. page 41) de $I_{\tau'}$ qui sont source ou but d'une flèche $f(u)$, où $u \in \mathrm{Tr}(\tau, f)$, i.e. $f(u)$ n'est pas dégénérée (\footnote[13]{dégager un lemme $f$ croissante inutile.}). $E'$ contient $f(I_{\tau\,\text{réd}})$ (i.e. les images par $f$ des sommets critiques de $I_\tau$),
\msp{page 46}mais peut être plus grand. (P.ex. si $\tau = \tau'$, $f = \mathrm{id}$, $E' = I_{\tau'}$, qui est donc strictement plus grand que $I_{\tau'\,\text{réd}}$ si $\tau'$ n'est pas réduit.) Pour tout $i \in E'$, considérons $f^{-1}(\{i\}) \subseteq I_\tau$. Comme $f$ \texttt{[est]} croissant, c'est un intervalle de l'ensemble totalement ordonné $I_\tau$, donc c'est un sous-intervalle de la catégorie intervalle $I_\tau$. Comme tel, il a une signature $\tau_i \in \Ty$. On a donc une application
\[
\eta_f \text{ ou } \eta : E' \longrightarrow \Ty.
\]
Je dis que $(\tau, f) \in \mathrm{Ob}\,(\uTy/\tau')$ est déterminé de façon unique par ce couple
\[
(E', \eta) \begin{cases} E' \in \mathfrak{P}(I_{\tau'}) = \mathfrak{P}([0,n']) \\[1ex] \eta : E' \longrightarrow \Ty\ . \end{cases}
\]
De plus, celui-ci est soumis à la seule restriction que
\[
E' \supseteq \mathrm{Cr}(\tau') = \text{ensemble des sommets critiques de } I_{\tau'}
\]
\printedpage{30}
% head: A. Grothendieck : Dérivateurs, Ch. VII, Édition des Universités de Montpellier II et Paris VII
(où ce qui revient au même, que $E'$ contient $0_{\tau'}$, $1_{\tau'}$ et que deux éléments consécutifs de $E'$ sont reliés par une flèche de $I_{\tau'}$) : % [sic] où -> ou, « ou ce qui revient au même »
\[
\left\{
\begin{array}{lcl}
\mathrm{Ob}\,(\uTy/\tau') & \simeq & \text{ensemble des couples } (E',\eta), \text{ avec} \\[1ex]
 & & E' \subseteq [0, \underbrace{\ n'\ }_{=\,\mathrm{long}\,\tau'}], \quad E' \supseteq \mathrm{Cr}(\tau'), \\[1ex]
 & & \eta : E' \longrightarrow \Ty\ .
\end{array}
\right.
\]
Soit
\[
0 = t_0 < t_1 < \ \cdots\ < t_r = n'
\]
la suite des éléments de $E'$, soit pour \ovalbox{$i \in [1,N]$}
\[
\left\|
\begin{array}{l}
l_i = t_i - t_{i-1} \\[1ex]
\varepsilon_i \quad \text{signature de } [t_{i-1}, t_i] \text{ dans } I_{\tau'}\ ,
\end{array}
\right.
\]
\msp{page 47}et pour \ovalbox{$i \in [0,N]$}\,, soit
\[
\|\ \ \tau_i = \eta(t_i).
\]
Alors le type $\tau$ est donné par
\[
\tau = \tau_0 \Delta_{\varepsilon_1} \tau_1 \Delta_{\varepsilon_2} \cdots \Delta_{\varepsilon_N} \tau_N, \tag*{$(*)$}
\]
\texttt{[de]} longueur $N + \sum_{i=0}^N \underbrace{\mathrm{long}(\tau_i)}_{\lambda_i \geq 0}$.

\dgrblock{dgrb307}{Fibres are components of the collapsed-edge graph}
\begin{dcomment}{exposition}{p45}{Fibres are components of the collapsed-edge graph}
Grothendieck wants a geometric reconstruction of the same weight vector. The intended reconstruction is correct once the graph of collapsed transitions is kept, but the printed description in terms of a set of vertices loses information that is needed to find its components.

Consider $\tau=+++$ with weights $(0,1,0)$. The corresponding map to $+$ is
\[
 f(0)=f(1)=0,\qquad f(2)=f(3)=1.
\]
The vertices incident to a collapsed transition form the whole set $\{0,1,2,3\}$. As a subset of the integers this set is connected. Consequently the assertion that $f$ is constant on its integer-connected components would make $f$ constant everywhere, contrary to the displayed functor. This is a printed mathematical slip, not a defective proof of the weight classification.

Instead, take the graph with all source vertices and only those edges $[i-1,i]$ for which $a_i=0$. Its components are consecutive strings of zero-weight edges, together with any isolated vertices. In the example they are $\{0,1\}$ and $\{2,3\}$. Since
\[
 f(j)-f(i)=\sum_{h=i+1}^{j}a_h\qquad(i\leq j),
\]
nonnegativity shows that two vertices have the same image if and only if every intervening edge has weight zero. The graph components are therefore exactly the nonempty vertex fibres of $f$. On each surviving transition the endpoint values are determined by the cumulative sum of the earlier weights. This proves both existence and uniqueness of the proposed geometric reconstruction after the correction.

The same observation produces a canonical factorization. First collapse the zero-weight edges, leaving the signed subword consisting of the active source edges; then stretch its successive edges by their positive weights. The first vertex map is surjective, the second strictly increasing, and both preserve the endpoints. Uniqueness follows either from the fibre components or from the unique image factorization of $q_f$ in the category of signed ordered sets. This factorization will let us distinguish a combinatorial contraction from a contraction which has an inverse on one side.
\end{dcomment}


\subsection*{3$^o$ Description de $\mathrm{Fl}(\uTy)$.}

Si
\[
f : \tau \longrightarrow \tau' \qquad \texttt{[est]} \ \text{dans } \uTy\ ,
\]
on lui associe comme tantôt un entier \underline{$N \geq 0$} ($N = \operatorname{card}\mathrm{Cr}(\tau', f){-}1$, $N = 0$ si et seulement si $\tau' = 1$, type de l'intervalle ponctuel), des applications
\[
\left\{
\begin{array}{lcl@{\qquad\qquad}ccl}
{[1,N]} & \longrightarrow & [+1,-1] \times \mathbf{N}^*, & i & \longmapsto & (\varepsilon_i, l_i) \\[1ex]
{[0,N]} & \longrightarrow & \Ty & i & \longmapsto & \tau_i
\end{array}
\right.
\]
(\footnote[14]{\textbf{NB} Si $N = 0$, $[1,N] = \emptyset$, et il n'y a pas de $\varepsilon_i$, $l_i$; l'application $\underbrace{[0,N]}_{=\,\{0\}} \longrightarrow \Ty$ se réduit à la donnée de $\tau_0 \in \Ty$.}). La flèche $f$ dans $\uTy$ est determinée par ces données. Le type $\tau$ est donné par la formule $(*)$ ci-dessus, le type $\tau'$ par % [sic] determinée -> déterminée, accent manquant
\[
\tau' = \Delta^{l_1}_{\varepsilon_1} \Delta^{l_2}_{\varepsilon_2}\ \cdots\ \Delta^{l_N}_{\varepsilon_N}, \qquad \text{longueur } n' = \textstyle\sum l_i, \tag*{$(*')$}
\]
\printedpage{31}
% head: A. Grothendieck : Dérivateurs, Ch. VII, Édition des Universités de Montpellier II et Paris VII
et l'application croissante de $I_\tau = [0, N + \sum \mathrm{long}(\tau_i)]$ dans $I_{\tau'} = [0, \sum l_i]$ est donnée ainsi :
\[
\begin{array}{l@{\qquad}cl}
\text{constante sur } [0, 0+\lambda_0] & \longmapsto & 0 \\[1ex]
\text{constante sur } [\lambda_0+1, \lambda_0+1+\lambda_1] & \longmapsto & l_1 \\[1ex]
\text{constante sur } [\lambda_0+\lambda_1+2, \lambda_0+\lambda_1+2+\lambda_2] & \longmapsto & l_2 \\[1ex]
\ldots & & \\[1ex]
\text{constante sur } [\lambda_0+\lambda_1+\cdots+\lambda_i+(i+1), \lambda_0+\cdots+\lambda_i+\lambda_{i+1}+(i+1)] & \longmapsto & l_i.
\end{array}
\]

\printedpage{32}
% head: A. Grothendieck : Dérivateurs, Ch. VII, Édition des Universités de Montpellier II et Paris VII
\msp{page 48}

\dgrblock{dgrb308}{Reading the morphism from its target}
\begin{dcomment}{exposition}{p47}{Reading the morphism from its target}
The incoming classification organizes a morphism by the places where its nonconstant pieces meet. It is the same normal form viewed from the target: choose the visited vertices, and at each visited vertex record the entire interval that has been collapsed there. This description is useful because the intervening noncollapsed transitions are then forced.

Let $f:I_\tau\to I_{\tau'}$ be increasing and endpoint-preserving, with target length $m$. The uniform definition is
\[
 E'=\operatorname{im}(f)\subseteq[0,m].
\]
When $m>0$, it agrees with the printed definition by endpoints of nondegenerate images. Indeed, take a vertex value $t$ assumed by $f$. If not every source vertex has value $t$, connectedness of the underlying source path supplies an edge from the nonempty set $f^{-1}(t)$ to its complement. That edge is nondegenerate and one of its image endpoints is $t$. Conversely every such endpoint is a value of $f$. This argument does not require monotonicity, and supplies the elementary lemma requested in footnote~13, provided the functor is nonconstant.

For $m=0$ all transition images are identities, so the printed endpoint definition would instead give $E'=\varnothing$. The later claim that $E'$ contains the endpoint, and the definition $N=\#E'-1$, then fail. Taking $E'=\operatorname{im}(f)=\{0\}$ repairs the nullary case and agrees with the explicit description of $N=0$ on manuscript page~47. This is an omitted boundary convention; it does not change the nonconstant construction.

Write the distinct elements of $E'$ as
\[
 0=t_0<t_1<\cdots<t_N=m.
\]
Every two consecutive values are joined by a target arrow: the source transition between their two adjacent nonempty fibres has exactly these image endpoints. Hence the segment from $t_{j-1}$ to $t_j$ is monochromatic, with sign $\varepsilon_j$. Equivalently, $E'$ contains the intrinsic critical vertices of the target, namely its endpoints and every place where the signs change. The notation $\mathrm{Cr}(\tau')$ has changed meaning at this point in the manuscript: here it denotes these vertices, whereas earlier it denoted critical intervals. Keeping these two sets distinct removes a potential confusion in the formula.

The fibre over $t_j$ is a nonempty consecutive interval of vertices; let its signature be $\tau_j$ and its length $\lambda_j$. Reading the source from left to right gives exactly
\[
 \tau=\tau_0\Delta_{\varepsilon_1}\tau_1
        \Delta_{\varepsilon_2}\cdots
        \Delta_{\varepsilon_N}\tau_N.
\]
Conversely, choose any subset $E'$ containing the critical vertices and any words $\tau_0,\ldots,\tau_N$. Form this source word. Collapse the $\tau_j$-block to $t_j$, and map the transition between consecutive blocks to the target arrow joining $t_{j-1}$ and $t_j$, in the direction prescribed by $\varepsilon_j$. This defines an increasing endpoint-preserving functor; its image is exactly $E'$, and its fibres have the prescribed words. Each arrow is uniquely determined by its endpoints in an interval category, so there is no additional gluing datum. The two constructions are inverse. For a point target they reduce to the arbitrary word $\tau_0$ and its unique map to the point.

For use in computations, set
\[
 l_j=t_j-t_{j-1},\qquad
 B_j=j+\sum_{h<j}\lambda_h.
\]
The correct full formula is
\[
 f(k)=t_j=\sum_{h=1}^{j}l_h
 \quad\text{for } B_j\leq k\leq B_j+\lambda_j,
 \qquad 0\leq j\leq N.
\]
The source's explicit table gives $l_2$ on the third fibre, where the value must be $l_1+l_2$; its final row similarly needs the corresponding cumulative sum. The formula above repairs all rows at once. For example, take target $++-$, visited vertices $E'=\{0,2,3\}$, and fibre words $\tau_0=-$, $\tau_1=+$, $\tau_2$ empty. Then the source is $-++-$ and the object map is $(0,0,2,2,3)$. The first positive transition crosses the two target positive edges, the next is collapsed, and the final negative transition crosses the final target edge. The two jumps have sizes $2$ and $1$, while the final vertex value is their sum $3$.
\end{dcomment}


\dgrblock{dgrb309}{Which contractions admit a section}
\begin{dcomment}{development}{p47}{Which contractions admit a section}
The normal form makes it possible to ask precisely when a deleted collection of transitions can be restored functorially. This is useful for the later path formalism: a surjection of underlying vertex sets need not be a split epimorphism in the category of signed types. The signs can obstruct every possible section.

\textbf{Section criterion for a contraction.}\ours{}
Let $\tau=(\varepsilon_1,\ldots,\varepsilon_n)$ and let $E=\{e_1<\cdots<e_k\}$ be a subset of its transition positions. Let
\[
 p_E:\tau\longrightarrow\tau|_E
\]
be the contraction with weights $1$ at the positions in $E$ and $0$ elsewhere. If $k=0$, it has a section exactly when $n=0$. If $k>0$, it has a section exactly under the following conditions. Every sign before $e_1$ equals $\varepsilon_{e_1}$, and every sign after $e_k$ equals $\varepsilon_{e_k}$. In each open gap $(e_j,e_{j+1})$, the sequence of signs is a possibly empty string of $\varepsilon_{e_j}$ followed by a possibly empty string of $\varepsilon_{e_{j+1}}$. When these conditions hold, the number of sections is
\[
 \prod_{\substack{1\leq j<k\\
                  \varepsilon_{e_j}=\varepsilon_{e_{j+1}}}}
       (e_{j+1}-e_j).
\]
An empty product is $1$.

\emph{Proof.} Under the signed-edge anti-isomorphism, $p_E$ corresponds to the inclusion $E\hookrightarrow[1,n]$. A section of $p_E$ therefore corresponds exactly to a nondecreasing sign-preserving retraction
\[
 r:[1,n]\longrightarrow E,\qquad r(e_j)=e_j.
\]
If $E$ is empty such a map exists exactly when its domain is empty. Assume henceforth $k>0$. Monotonicity forces every index before $e_1$ to map to $e_1$, and every index after $e_k$ to map to $e_k$. This proves the two exterior conditions, and these parts of $r$ are unique.

For $e_j<i<e_{j+1}$, monotonicity forces
$e_j\leq r(i)\leq e_{j+1}$. Since there is no element of $E$ strictly between them, $r(i)$ is one of these two endpoints. Moreover the indices mapped left must all precede the indices mapped right. There is therefore a single cut in the gap. Sign preservation is equivalent to the signs before the cut being $\varepsilon_{e_j}$ and those after it being $\varepsilon_{e_{j+1}}$. This proves necessity of the gap condition. If the endpoint signs differ, an admissible cut is unique. If they agree, every sign in the gap must have that common value and all $d+1$ cuts are admissible, where $d=e_{j+1}-e_j-1$.

Conversely, choose an admissible cut independently in every gap, use the forced exterior values, and fix every $e_j$. The resulting map is nondecreasing across each gap and its boundary, hence globally nondecreasing; it preserves signs by construction and retracts onto $E$. These choices describe every retraction exactly once. Multiplying their independent numbers gives the displayed product. The signed-edge correspondence turns each retraction back into a unique section of $p_E$, completing the proof.

For $\tau=+-$ and $E=\{1\}$ the trailing negative sign violates the exterior condition. Thus the surjective vertex map $+ -\to+$ has no section. For $\tau=+++$ and $E=\{1,3\}$ the gap has one omitted positive edge, and there are exactly two sections; their vertex maps $[0,2]\to[0,3]$ are $(0,1,3)$ and $(0,2,3)$. This distinguishes the collapse of an edge from a simplicial degeneracy with an available section.

We checked the ordinal description and comonoid interpretation in \citet[Appendix~D]{AguiarMahajan2010}, the zig-zag and Reedy structures of \citet[Definition~1.1 and Example~4.12]{ArakawaCnossen2026}, and searched for signed or coloured ordinal retractions and sections of Grothendieck's interval maps. We could not find this exact interval-form criterion and count in those sources. The result is an elementary consequence of the normal form, recorded here as a proved continuation of the calculation, without a claim of priority.
\end{dcomment}


\section*{6\quad Composition dans $\uTy$}

On a une loi de composition (partout définie) dans $\mathrm{Ob}\,(\uTy) = \Ty$, c'est celle du monoïde libre à deux générateurs $\Delta_+$, $\Delta_-$. Cette loi de composition s'interprète, au niveau catégorique (et non plus seulement combinatoire) comme une somme amalgamée
\[
I_{\tau\tau'} = I_\tau \amalg_{(e, \underline{1}_\tau, \underline{0}_{\tau'})} I_{\tau'},
\]
où $e$ est la catégorie unité, et où $\underline{1}_\tau$, $\underline{0}_{\tau'}$ représentent maintenant les morphismes $e \longrightarrow I_\tau$, $e \longrightarrow I_{\tau'}$ , ayant comme image l'élément but et l'élément source respectivement de l'intervalle orienté envisagé. On a donc un carré cocartésien dans (Cat)
\[
\begin{tikzcd}[column sep=small, row sep=small]
 & e \arrow[dl, hook', "\underline{1}_\tau"'] \arrow[dr, hook, "\underline{0}_{\tau'}"] & \\
I_\tau \arrow[dr, hook, "i_{\tau,\tau'}"'] & \text{\scriptsize cocart.} & I_{\tau'} \arrow[dl, hook', "j_{\tau,\tau'}"] \\
 & I_{\tau\tau'}, &
\end{tikzcd}
\]
où toutes les flèches sont des foncteurs mono et pleinement fidèles. Que ce diagramme (évidemment commutatif) soit bien cocartésien dans (Cat) revient, par définition, à dire que pour toute $X$ dans Cat, le diagramme transposé
\[
\begin{tikzcd}[column sep=small, row sep=small]
 & X & \\
\uCh_\tau(X) \arrow[ur, "\beta_\tau" near end, "\text{\scriptsize (but)}" near start] & & \uCh_{\tau'}(X) \arrow[ul, "\sigma_{\tau'}"' near end, "\text{\scriptsize (source)}"' near start] \\
 & \uCh_{\tau\tau'}(X), \arrow[ul] \arrow[ur] &
\end{tikzcd}
\]

\texttt{[est cartésien]}

\msp{page 49}ce qui revient à la tautologie suivante qu'un chemin $C$ dans $X$ de type $\tau\tau'$ `n'est autre chose' qu'un couple $(c,c')$ d'un chemin $c$ de type $\tau$ et d'un chemin $c'$ de type $\tau'$ `mis bout à bout' (dans cet ordre), i.e. tels que
\[
\mathrm{but}(c) = \mathrm{source}(c')
\]
(et on a alors
\[
\begin{array}{rcll}
\mathrm{source}(C) & = & \mathrm{source}(c) & \\[1ex]
\mathrm{but}(C) & = & \mathrm{but}(c') & \qquad\quad ).
\end{array}
\]

\printedpage{33}
% head: A. Grothendieck : Dérivateurs, Ch. VII, Édition des Universités de Montpellier II et Paris VII
Il se pose ici le sempiternel problème si on va noter ce chemin $c{\cdot}c'$, ce qui semble à première vue s'imposer, ou plutôt $c' \circ c$, regardant le chemin $c$ de $x = \mathrm{source}(c)$ vers $y = \mathrm{but}(c)$ comme une sorte de morphisme de $x$ dans $y$ (cf. plus bas le formalisme des catégories de fractions), $c'$ comme un morphisme de $y = \mathrm{source}(c')$ vers $z = \mathrm{but}(c')$, et en interprétant le composé $C$ comme le composé de ces deux morphismes $c$, $c'$ (dans cet ordre du point de vue des opérations à effectuer), qu'on note traditionellement $c' \circ c$ ou $c' \cdot c$, et non $c \circ c'$ ou $c \cdot c'$, hélas! C'est cette notation que nous allons suivre plus bas. % [sic] traditionellement -> traditionnellement, orthographe

Mais pour l'instant on ne va pas se soucier des chemins dans des catégories quelconques $X$, mais revenir au formalisme des $\tau$ et des $I_\tau$ eux-mêmes. L'interprétation catégorique de $\tau\tau'$ comme une somme amalgamée nous assure que cette opération de composition est fonctorielle :

\msp{page 50}
\[
\tau_1 \xrightarrow{\ f\ } \tau_2, \quad \tau'_1 \xrightarrow{\ f'\ } \tau'_2 \qquad\quad \text{donne } \tau_1\tau'_1 \xrightarrow{f*f'} \tau_2\tau'_2
\]
(on n'écrira pas $f \cdot f'$ pour la flèche notée $f * f'$, à cause de la confusion possible avec la composition des flèches dans $\uTy$). On a transitivité de façon évidente : si de plus $\tau_2 \xrightarrow{\ g\ } \tau_3$, $\tau'_2 \xrightarrow{\ g'\ } \tau'_3$, \texttt{[on a]}
\[
(gf) * (g'f') = (g * g') \circ (f * f').
\]
Donc on a un bifoncteur de composition
\[
\uTy \times \uTy \longrightarrow \uTy \qquad\quad (c,c') \longmapsto c \cdot c' \text{ ou } c * c'.
\]
(En fait, on a aussi un bifoncteur
\[
\uTyp \times \uTyp \longrightarrow \uTyp,
\]
mais on préfère oublier maintenant la catégorie \underline{Typ}, qui donne un formalisme nettement moins simple.)

\dgrblock{dgrb310}{Concatenation is a pushout with chosen endpoints}
\begin{dcomment}{context}{p49}{Concatenation is a pushout with chosen endpoints}
Grothendieck's reason for introducing types now becomes visible: the concatenation of words should encode the actual gluing of paths, including their transformations. An interval is a free category on an oriented linear graph. Gluing the last vertex of the first graph to the first vertex of the second produces the graph of the concatenated word, and the free-category construction preserves this pushout. Therefore
\[
 I_{\tau\tau'}=I_\tau\amalg_e I_{\tau'}
\]
is a pushout in $\Cat$ with the specified endpoint maps. Any directed path between two vertices of one summand stays in that summand: entering the other would require crossing the gluing vertex twice in an underlying linear graph. Thus the inclusions are fully faithful as well as injective on objects.

The corresponding assertion about path categories is also literal. Restriction gives an isomorphism
\[
 \operatorname{Fun}(I_{\tau\tau'},X)
 \cong
 \operatorname{Fun}(I_\tau,X)
 \times_X\operatorname{Fun}(I_{\tau'},X),
\]
where the fibre product uses evaluation at the two endpoints being identified. On objects, two functors agree at the common vertex exactly when their two paths concatenate. On morphisms, two natural transformations glue exactly when their components at that vertex are equal. Naturality along each generating edge then gives naturality on every composite. This verifies the categorical pullback, rather than only the corresponding assertion about sets of paths.

The convention on the order of composition now has a concrete meaning. The word $\tau\tau'$ is read first through $\tau$ and then through $\tau'$, whereas the corresponding categorical composition is written $c'\circ c$. This harmless reversal of written order should be fixed before discussing fractions. It is not an opposite-category operation and does not reverse any transition.

The endpoint data are indispensable. The parenthetical assertion on manuscript page~50 that concatenation also gives a bifunctor on $\uTyp$, as defined on page~30 using \emph{all} functors, is false. The one-arrow categories $I_+$ and $I_-$ are isomorphic in $\uTyp$. A bifunctor whose object operation is word concatenation would carry this isomorphism, together with $\mathrm{id}_{I_+}$, to an isomorphism between $I_{++}$ and $I_{-+}$. The former is a three-object chain; the latter has one initial object and two incomparable maximal objects. They are not isomorphic. Thus the claimed bifunctor cannot exist on that larger category. It does exist on endpoint-preserving functors, even without monotonicity, because their restrictions agree at the glued vertex. The final category $\uTy$ satisfies this requirement and remains closed under concatenation. This correction isolates precisely the structure that the path formalism needs.
\end{dcomment}


Cette composition est \textit{strictement} associative et \textit{strictement unitaire} (l'unité étant bien sûr l'élément unité 1 du monoïde libre Ty, correspondant à la catégorie ponctuelle $\{0\}$), et pas seulement `à isomorphisme près'. Et pour cause, puisque dans $\uTy$ les isomorphismes sont des identités. \textit{Donc} $\uTy$ \textit{est un monoïde dans la catégorie} (Cat) des petites catégories!

\msp{page 49$'$ (il y a un décalage dans la numérotation)}Cela implique que non seulement $\mathrm{Ob}\,\uTy = \Ty$, mais aussi $\mathrm{Fl}(\uTy)$, et plus généralement tous les ensembles de la forme $\mathrm{Hom}_{\mathrm{Cat}}(X, \uTy)$ sont des monoïdes, et ceci fonctoriellement en $X$, par exemple les composantes $\mathrm{Hom}_{\mathrm{Cat}}(\boldsymbol{\Delta}_i, \uTy)$ de l'ensemble semi-simplicial `nerf' de $\uTy$,
\[
\mathrm{Nerf}_i(\uTy) = \mathrm{Hom}_{\mathrm{Cat}}(\boldsymbol{\Delta}_i, \uTy),
\]
sont des monoïdes, et les applications semi-simpliciales entre elles sont des homomorphismes de monoïdes : le nerf de $\uTy$ est un \textit{monoïde semi-simplicial.}

% ---- from Der-ChVII-456_p036-040.tex
\printedpage{34}
% head: A. Grothendieck : Dérivateurs, Ch. VII, Édition des Universités de Montpellier II et Paris VII
On va décrire ces monoïdes $\mathrm{Nerf}_i(\uTy)$, et constater que ce sont tous des monoïdes \textit{libres,} tout comme $\mathrm{Nerf}_0(\uTy) = \mathrm{Ob}\,\uTy = \mathrm{Ty}$. Les applications semi-simpliciales entre ces monoïdes sont donc connus quand on les connaît sur les \textit{générateurs} de ces monoïdes libres (lesquels générateurs sont aussi les éléments indécomposables, donc définis intrinsèquement en termes de la structure multiplicative). % [sic] connus -> connues, agreement with « applications »

Pour $\mathrm{Nerf}_0(\uTy) = \mathrm{Ty}$, c'est une chose vue, les générateurs sont
\[
\Delta_+,\Delta_-,
\]
correspondant respectivement aux intervalles orientés
\[
\begin{array}{rccc}
\boldsymbol{\Delta}_+ : & 0 & \longrightarrow & 1\\[4pt]
\boldsymbol{\Delta}_- : & 0 & \longleftarrow & 1.
\end{array}
\]

\msp{page 50$'$}
Pour $\mathrm{Nerf}_1(\uTy) = \mathrm{Fl}(\uTy)$, i.e. pour les flèches, le résultat $3^o$) du paragraphe précédent (détermination de $\mathrm{Fl}(\uTy)$) s'interprète aussi en disant que $\mathrm{Fl}(\uTy)$ est le monoïde libre engendré par un système de générateurs infini
\[
\mathfrak{f}_{\varepsilon,n}\ , \qquad\quad (\varepsilon,n)\in\{+1,-1\}\times\mathbf{N},
\]
indexé par l'ensemble produit $\{+1,-1\}\times\mathbf{N}$. Ici $\mathfrak{f}_{\varepsilon,n}$ désigne l'unique flèche dans $\uTy$
\[
\mathfrak{f}_{\varepsilon,l} : \Delta_\varepsilon \longrightarrow \Delta^l_\varepsilon
\]
(où bien sûr l'exposant $l$ désigne une exponentation pour la loi de composition dans $\uTy$
\[
\Delta^l_\varepsilon = \underbrace{\Delta_\varepsilon\circ\Delta_\varepsilon\circ\ \cdots\ \circ\Delta_\varepsilon}_{l\text{ copies}},
\]
et non la puissance cartésienne dans Cat, où dans $\uTy$ (où elle n'existe que si $l\in\{0,1\}$)). % [sic] où dans -> ou dans, conjunction intended
Dans la description de $\mathrm{Fl}(\uTy)$ dans le paragraphe précédent, on a écrit le `terme général' du monoïde libre en groupant ensemble les groupes de foncteurs successifs de la forme $f_{+,0}$, $f_{-,0}$, i.e. corréspondant aux deux flèches de dégénérescence indécomposables (pour $*$) % [sic] corréspondant -> correspondant, accent
\[
\left\{
\begin{array}{lccc}
f_{+,0} : & \Delta_+ & \longrightarrow & e\\[4pt]
f_{-,0} : & \Delta_- & \longrightarrow & e
\end{array}
\right.
\]
(\footnote[15]{\textbf{NB} Le `mot' vide de $\mathrm{Fl}(\uTy)$ correspond à la flèche $\mathrm{id}_1$, où $1$ est l'élément unité de Ty, réalisant l'`intervalle' ponctuel.}). Le sous-monoïde libre de $\mathrm{Fl}(\uTy)$ engendré par ces deux éléments est le monoïde des flèches de dégénérescence (pure) quelconques
% foot: Version 22 décembre 2001
\printedpage{35}
% head: A. Grothendieck : Dérivateurs, Ch. VII, Édition des Universités de Montpellier II et Paris VII
\msp{page 51}
dans $\uTy$ :
\[
f_\tau : I_\tau \longrightarrow e \qquad\quad (\tau\in\mathrm{Ty}),
\]
et on a
\[
f_\tau * f_{\tau'} = f_{\tau\tau'},
\]
i.e. l'application
\[
\tau \longmapsto f_\tau :\quad \underbrace{\mathrm{Ob}\,(\uTy)}_{=\,\mathrm{Ty}\,=\,\mathrm{Nerf}_0(\uTy)} \quad\longrightarrow\quad \underbrace{\mathrm{Fl}(\uTy)}_{\mathrm{Nerf}_1(\uTy)}
\]
est un homomorphisme de monoïdes, transformant $\Delta_\varepsilon$ ($\varepsilon\in\{+1,-1\}$) dans $\mathfrak{f}_{\varepsilon,0}$. Ce n'est pas l'application de dégénérenscence de $\mathrm{Nerf}_0$ dans $\mathrm{Nerf}_1$, i.e. l'application $\tau\longmapsto\mathrm{id}_\tau$. Celle-ci est donnée par % [sic] dégénérenscence -> dégénérescence, spelling
\[
\begin{array}{ccc}
\mathrm{Ob}\,(\uTy) & \longrightarrow & \mathrm{Fl}(\uTy)\\[4pt]
\tau & \longmapsto & \mathrm{id}_\tau
\end{array}
\ :\
\left|\
\begin{array}{rcccl}
 & \Delta_\varepsilon & \longmapsto & \mathfrak{f}_{\varepsilon,1} & \qquad(\varepsilon\in\{+1,-1\})\\[4pt]
\text{i.e.} & \Delta_+ & \longmapsto & \mathfrak{f}_{+,1} & \\[4pt]
 & \Delta_- & \longmapsto & \mathfrak{f}_{-,1}\ . &
\end{array}
\right.
\]
Quant aux homomorphismes source et but
\[
\begin{tikzcd}[column sep=large]
\mathrm{Fl}(\uTy) \arrow[r, shift left, "\sigma\text{ (source)}"] \arrow[r, shift right, "\beta\text{ (but)}"'] & \mathrm{Ob}\,(\uTy),
\end{tikzcd}
\]
ils sont évidents sur les générateurs
\[
\left\{
\begin{array}{lcl}
\sigma(\mathfrak{f}_{\varepsilon,l}) & = & \Delta_\varepsilon\\[4pt]
\beta(\mathfrak{f}_{\varepsilon,l}) & = & \Delta^l_\varepsilon.
\end{array}
\right.
\]

\dgrblock{dgrb311}{The type category is generated by two comonoids}
\begin{dcomment}{later}{p50}{The type category is generated by two comonoids}
The strict monoidal structure places Grothendieck's calculation beside a classical universal property of the simplicial category. In the single-sign case, endpoint-preserving maps of finite intervals are opposite to monotone maps of finite ordered edge sets. The latter category, with ordinal sum, contains the universal monoid; its opposite contains the universal counital coassociative comonoid. An explicit account, with the weak-composition coordinates used here, is \citet[Appendix~D.1--D.2, especially Proposition~D.2, pp.~713--716]{AguiarMahajan2010}, which traces the universal property to Mac Lane's Proposition~VII.5.1. The two signs simply keep two independent copies of this structure in the same ordered word.

In $\uTy$ put
\[
 \delta_\varepsilon=\mathfrak f_{\varepsilon,2}:
      \Delta_\varepsilon\longrightarrow\Delta_\varepsilon^2,
 \qquad
 e_\varepsilon=\mathfrak f_{\varepsilon,0}:
      \Delta_\varepsilon\longrightarrow1,
 \qquad \varepsilon\in\{+,-\}.
\]
Both ways of applying $\delta_\varepsilon$ twice give the unique endpoint-preserving map from $\Delta_\varepsilon$ to $\Delta_\varepsilon^3$. Applying $e_\varepsilon$ to either output of $\delta_\varepsilon$ gives the identity, since its surviving weight is $1$. Thus each $\Delta_\varepsilon$ is a counital coassociative comonoid. There is no specified relation between the positive and negative comonoids.

More precisely, a strict monoidal functor from $\uTy$ to a strict monoidal category $(\mathcal C,\otimes,\mathbf1)$ is uniquely specified by a pair of counital coassociative comonoids $C_+$ and $C_-$ in $\mathcal C$. Here is the construction, which also proves the assertion. For either comonoid, let $\delta^{(a)}:C\to C^{\otimes a}$ denote the iterated comultiplication, with $\delta^{(0)}$ its counit and $\delta^{(1)}=\mathrm{id}$. Coassociativity and the counit laws give, for every list $b_1,\ldots,b_a$ of nonnegative integers,
\[
 \bigl(\delta^{(b_1)}\otimes\cdots\otimes\delta^{(b_a)}\bigr)
       \delta^{(a)}
 =\delta^{(b_1+\cdots+b_a)}.
\]
For completeness, the binary identity
$(\delta^{(p)}\otimes\delta^{(q)})\delta^{(2)}
 =\delta^{(p+q)}$ follows by induction on positive $p,q$ using coassociativity; the cases $p=0$ or $q=0$ are the counit identities, including $p=q=0$. Applying this binary identity successively to the first $a-1$ factors and the last factor proves the displayed formula by induction on $a$. The case $a=0$ says that the identity of the unit composed with the counit is the counit.

Send the word $\varepsilon_1\cdots\varepsilon_n$ to
$C_{\varepsilon_1}\otimes\cdots\otimes C_{\varepsilon_n}$ and send its outgoing morphism of weights $(a_1,\ldots,a_n)$ to
\[
 \delta_{\varepsilon_1}^{(a_1)}\otimes\cdots\otimes
 \delta_{\varepsilon_n}^{(a_n)}.
\]
The block-sum formula for composition and the preceding identity prove that this assignment is a functor. The all-one vector gives identities; the empty word gives the unit; concatenation of vectors gives the tensor product of maps. It is therefore strict monoidal. Conversely every outgoing morphism is the tensor product of its one-edge factors, and each such factor is an iterated comultiplication or counit. A strict monoidal functor must have the displayed values, establishing uniqueness.

This is a reinterpretation of the known ordinal/comonoid universal property, not a new homotopy theorem. It explains the two apparently different generators in the manuscript: $\mathfrak f_{\varepsilon,0}$ erases a transition, whereas $\mathfrak f_{\varepsilon,2}$ subdivides one. Their iterates produce every allowed reparametrization. No braiding is present: the order of the source transitions is part of the structure and cannot be permuted.
\end{dcomment}


Je vais décrire maintenant $\mathrm{Nerf}_2(\uTy)$, i.e. les systèmes
\[
\tau \xrightarrow{\ f\ } \tau' \xrightarrow{\ g\ } \tau''.
\]

\msp{page 52}
Se donner $f\in\mathrm{Fl}(\uTy)$ revient à se donner une suite ordonnée de générateurs $f_{\varepsilon,l}$, donc une suite
\[
\begin{aligned}
&\mathfrak{f}_{\varepsilon_1,l_1},\ \mathfrak{f}_{\varepsilon_2,l_2},\ \ \dots\ \ ,\mathfrak{f}_{\varepsilon_n,l_n}\\
&((\varepsilon_i,l_i))_{1\leq i\leq n}\in[\{\pm1\}\times\mathbf{N}]^n.
\end{aligned}
\]
On a alors
\[
\begin{aligned}
x &= \Delta_{\varepsilon_1}\Delta_{\varepsilon_2}\cdots\Delta_{\varepsilon_n} &\qquad\qquad& \mathrm{long}\,x = n\\
y &= \Delta^{l_1}_{\varepsilon_1}\Delta^{l_2}_{\varepsilon_2}\cdots\Delta^{l_n}_{\varepsilon_n} && \mathrm{long}\,y = \textstyle\sum_{i\in[1,n]} l_i.
\end{aligned}
\]
% foot: Version 22 décembre 2001

\printedpage{36}
% head: A. Grothendieck : Dérivateurs, Ch. VII, Édition des Universités de Montpellier II et Paris VII
Se donner $q : y\longrightarrow z$, i.e. un objet de $y\backslash\uTy$, revient à se donner une application de l'ensemble des $m=\sum l_i$ flèches de transition de $y$ dans $\mathbf{N}$. En fait, on a % [sic] q -> g, the map is called g below (« la donnée de g »)
\[
\mathrm{Tr}(y) = \mathrm{Tr}(\Delta^{l_1}_{\varepsilon_1})\amalg\mathrm{Tr}(\Delta^{l_2}_{\varepsilon_2})\amalg\ \cdots\ \amalg\mathrm{Tr}(\Delta^{l_n}_{\varepsilon_n})
\]
(`ensemble somme', i.e. réunion disjointe), donc
\[
\mathrm{Hom}_{\mathrm{Ens}}(\mathrm{Tr}(y)) \simeq \prod_{i\in[1,n]}\mathrm{Hom}_{\mathrm{Ens}}(\underbrace{\mathrm{Tr}(\Delta^{l_i}_{\varepsilon_i})}_{[1,l_i]},\mathbf{N}), % [sic] Hom_Ens(Tr(y)) -> Hom_Ens(Tr(y),N), second argument missing
\]
donc la donnée de $g$ équivaut à la donnée d'un système de $n$ applications
\[
\left\{
\begin{array}{l}
\pi_i : [1,l_i]\longrightarrow\mathbf{N} \qquad\quad (\text{pondération de }\Delta^{l_i}_{\varepsilon_i})\\[6pt]
\text{pour } i\in[1,n].
\end{array}
\right.
\]
Mais pour $i\in[1,n]$, considérons l'élément de $\mathrm{Nerf}_2$
\[
\Delta_{\varepsilon_1}\xrightarrow{\mathfrak{f}_{\varepsilon_i,l_i}}\Delta^{l_i}_{\varepsilon_i}\xrightarrow{g_{\varepsilon,\pi_i}}\underbrace{\Delta^{\pi_i(1)}_{\varepsilon_i}\Delta^{\pi_i(2)}_{\varepsilon_i}\cdots\Delta^{\pi_i(l_i)}_{\varepsilon_i}}_{\Delta_{\varepsilon_i}^{\Sigma_j\pi_i(j)\overset{\text{déf}}{=}\mathrm{long}(\pi_i)}}, % [sic] \Delta_{\varepsilon_1} -> \Delta_{\varepsilon_i}, and g_{\varepsilon,\pi_i} -> g_{\varepsilon_i,\pi_i}, index i intended
\tag*{(1)}
\]

\msp{page 53}
où $g_{\varepsilon_i,\pi_i}$ est la flèche définie en termes de $\pi_i$ au paragraphe précédent. Mais ce n'est autre aussi que
\[
g_{\varepsilon_i,\pi_i} = \mathfrak{f}_{\varepsilon_i,\pi_i(1)}*\mathfrak{f}_{\varepsilon_i,\pi_i(2)}*\cdots*\mathfrak{f}_{\varepsilon_i,\pi_i(l_i)}.
\tag*{(2)}
\]
L'élément de $\mathrm{Nerf}_2$ envisagé ne dépend que de
\[
(\varepsilon_i,l_i,\pi_i), \qquad\quad \text{où } (\varepsilon_i,l_i)\in\{\pm1\}\times\mathbf{N},\ \pi_i : [1,l_i]\longrightarrow\mathbf{N}.
\]
(\textbf{NB} Si $l_i=0$, $[1,l_i]=[1,0]=\emptyset$ par convention, et $\pi_i$ est l'unique application $\emptyset\longrightarrow\mathbf{N}$, donc la donnée $(\varepsilon_i,l_i,\pi_i)$ revient alors à $(\varepsilon_i,l_i)$.)

Si
\[
(\varepsilon,l,\pi)\ \text{ comme dessus,} \qquad\quad (\varepsilon,l)\in\{\pm1\}\times\mathbf{N},\ \pi : [1,l]\longrightarrow\mathbf{N},
\]
désignons par
\[
\mathfrak{f}_{\varepsilon,l,\pi}
\]
l'élément de $\mathrm{Nerf}_2$ donné par
\[
\Delta_\varepsilon\xrightarrow{\mathfrak{f}_{\varepsilon,l}}\Delta^l_\varepsilon\ \xrightarrow{\quad\mathfrak{f}_{\varepsilon,\pi(1)}\ \ *\ \cdots\ *\ \ \mathfrak{f}_{\varepsilon,\pi(l)}\quad}\Delta^{\pi(1)}_\varepsilon\Delta^{\pi(2)}_\varepsilon\cdots\Delta^{\pi(l)}_\varepsilon.
\]
On peut formuler la construction précédente en disant que tout élément de $\mathrm{Nerf}_2$ peut se décomposer de façon unique en un produit
\[
(f,g) = \mathfrak{f}_{\varepsilon_1,l_1,\pi_1}*\mathfrak{f}_{\varepsilon_2,l_2,\pi_2}*\cdots*\mathfrak{f}_{\varepsilon_n,l_n\pi_n}, % [sic] l_n\pi_n -> l_n,\pi_n, comma missing in index
\]
ce qui prouve que $\mathrm{Nerf}_2(\uTy)$ est le monoïde libre sur les générateurs précédents $\mathfrak{f}_{\varepsilon,l,\pi}$. (\textbf{NB} à la suite vide correspond
\[
1\xrightarrow{\mathrm{id}_1}1\xrightarrow{\mathrm{id}_1}1 \qquad\qquad .)
\]
% foot: Version 22 décembre 2001

\printedpage{37}
% head: A. Grothendieck : Dérivateurs, Ch. VII, Édition des Universités de Montpellier II et Paris VII
\msp{page 54}
Si l'objet de $\mathrm{Nerf}_2$ est
\[
x\xrightarrow{\ f\ }y\xrightarrow{\ g\ }z,
\]
la construction précédente nous dit, en termes de la donnée combinatoire, i.e. de la suite de données
\[
(\varepsilon_1,l_1,\pi_1),\ \dots\ ,\ (\varepsilon_n,l_n,\pi_n),
\]
quels sont $x$, $y$, $z$ : chacun d'eux est le produit dans $\mathrm{Nerf}_0=\mathrm{Ty}$ des $n$ facteurs $x_i$, $y_i$ ou $z_i$ respectivement, pour $i\in[1,n]$, donnés dans la définition des $\mathfrak{f}_{\varepsilon_i,l_i,\pi_i}$, i.e. de
\[
\Delta_{\varepsilon_i},\ \Delta^{l_i}_{\varepsilon_i},\ \text{ et }\Delta^{\pi_i(1)}_{\varepsilon_i}\Delta^{\pi_i(2)}_{\varepsilon_i}\cdots\Delta^{\pi_i(l_i)}_{\varepsilon_i}=\Delta^{\mathrm{long}(\pi_i)}_{\varepsilon_i}
\]
(où $\mathrm{long}(\pi_i)=\sum_{j\in[1,l_i]}\pi_i(j)$). Donc
\[
\left\{
\begin{array}{lcl}
x & = & \prod^{\mathrm{Ty}}_{i\in[1,n]}\Delta_{\varepsilon_i}\\[4pt]
y & = & \prod^{\mathrm{Ty}}_{i\in[1,n]}\Delta^{l_i}_{\varepsilon_i}\\[4pt]
z & = & \prod^{\mathrm{Ty}}_{i\in[1,n]}\Delta^{\mathrm{long}(\pi_i)}_{\varepsilon_i}
\end{array}
\right.
\]
(\footnote[16]{attention, les produits doivent être pris en respectant l'ordre des facteurs donné par la notation indicielle en $i$.}), et on trouve de même $f$, $g$ comme des produits (au sens $*$) de $n$ facteurs $f_i$ resp. $g_i$ ($i\in[1,n]$) explicités dans les formules (1), (2) (p. 52, 53) :
\[
\begin{array}{lclccl}
f_i & = & \mathfrak{f}_{\varepsilon_i,l_i} : & \Delta_{\varepsilon_i} & \longrightarrow & \Delta^{l_i}_{\varepsilon_i}\\[4pt]
g_i & = & g_{\varepsilon_i,l_i} : & \Delta^{l_i}_{\varepsilon_i} & \longrightarrow & \Delta^{\mathrm{long}(\pi_i)}_{\varepsilon_i}, % [sic] g_{\varepsilon_i,l_i} -> g_{\varepsilon_i,\pi_i}, notation of (1), (2)
\end{array}
\]
$g_i$ lui-même étant représenté comme un $*$-produit de $l_i$ facteurs indécomposables, de sorte que $f$ et $g$ apparaissent sous forme
\msp{page 55}
de $*$-produits de $n$, resp. de $\sum_{i\in[1,n]}l_i$ facteurs de la forme $\mathfrak{f}_{\varepsilon,\nu}$. Pour terminer de tenir en mains l'objet semi-simplicial tronqué $(\mathrm{Nerf}_i(\uTy))_{0\leq i\leq 2}$, il reste à décrire les deux applications de dégénérescence $\mathrm{Nerf}_1\longrightarrow\mathrm{Nerf}_2$ (cas où $f$ ou $g$ est une flèche identique) et surtout l'application de \textit{composition} des flèches
\[
\mathrm{Nerf}_2(\uTy)\xrightarrow{\ \gamma\ }\mathrm{Nerf}_1(\uTy)=\mathrm{Fl}(\uTy).
\]
Il suffit de donner $\gamma$ sur les générateurs, et on a tautologiquement, par définition des $\mathfrak{f}_{\varepsilon,l,\pi}$
\[
\gamma(\mathfrak{f}_{\varepsilon,l,\pi})=\mathfrak{f}_{\varepsilon,\mathrm{long}(\pi)}.
\]
(\textbf{NB} Il n'existe au plus qu'un seul morphisme de $\Delta_\varepsilon$ (de $\Delta_+$ ou de $\Delta_-$) dans un intervalle $I_\tau$. Il existe si et seulement si $\tau$ est de la forme $\Delta^l_\varepsilon$, $l\in\mathbf{N}$.)
% foot: Version 22 décembre 2001

\dgrblock{dgrb312}{Free monoids in degrees one and two}
\begin{dcomment}{exposition}{p53}{Free monoids in degrees one and two}
Grothendieck uses the nerve to keep composition together with concatenation. A string of composable reparametrizations can still be cut along the original source edges. The resulting factors are independent, and this is the reason the monoid in each nerve degree is free.

In degree one, a morphism with source word $\varepsilon_1\cdots\varepsilon_n$ and weights $(a_1,\ldots,a_n)$ has the unique decomposition
\[
 f=\mathfrak f_{\varepsilon_1,a_1}*\cdots*
       \mathfrak f_{\varepsilon_n,a_n}.
\]
The factors are indexed by source edges, even when their weights vanish. A generator $\mathfrak f_{\varepsilon,0}$ is therefore not the unit of this free monoid: its source has one edge. The monoid unit is the empty factorization, namely the identity of the point. Nor is $\mathfrak f_{\varepsilon,1}$ the monoid unit; it is the identity of a one-edge object and is indecomposable under concatenation. This separates the categorical identity operation from the unit of the tensor product.

Now consider $\tau\xrightarrow f\tau'\xrightarrow g\tau''$. The $i$th original edge produces an ordered block of $a_i$ edges after $f$. Let $\pi_i:[1,a_i]\to\mathbb N$ record the weights that $g$ assigns to that block. All descendants have the original sign $\varepsilon_i$. Consequently the pair has the unique factorization
\[
 (f,g)=\prod_{i=1}^{n,*}\mathfrak f_{\varepsilon_i,a_i,\pi_i}.
\]
Existence comes from cutting every intermediate word into its descendant blocks; uniqueness comes from the fixed order of the original edges and the unique weight vector of each restriction. The nullary case $a_i=0$ has one empty function $\pi_i$, so it contributes a legitimate generator whose two targets are the point. For $n=0$ the whole pair is the identity string at the point and is the monoid unit. This proves the claimed freeness in degree two without discarding extinct blocks.

Writing $N_r=N_r(\uTy)$, the three faces of a degree-two generator are
\[
\begin{aligned}
 d_2(\mathfrak f_{\varepsilon,a,\pi})
    &=\mathfrak f_{\varepsilon,a},\\
 d_1(\mathfrak f_{\varepsilon,a,\pi})
    &=\mathfrak f_{\varepsilon,\sum_{j=1}^{a}\pi(j)},\\
 d_0(\mathfrak f_{\varepsilon,a,\pi})
    &=\mathfrak f_{\varepsilon,\pi(1)}*\cdots*
        \mathfrak f_{\varepsilon,\pi(a)}.
\end{aligned}
\]
The last line is the empty product when $a=0$. The two degeneracies of a degree-one generator insert the identity before or after its map:
\[
 s_0(\mathfrak f_{\varepsilon,a})
       =\mathfrak f_{\varepsilon,1,(a)},\qquad
 s_1(\mathfrak f_{\varepsilon,a})
       =\mathfrak f_{\varepsilon,a,1_a}.
\]
Here $(a)$ is the one-entry weight vector and $1_a$ the all-one vector of length $a$. These formulas recover precisely the composition and degeneracies described on manuscript pages~54--56, while displaying the other two faces as well.

The word ``semi-simplicial'' in this edition includes degeneracies; in current terminology the nerve is a simplicial monoid. The faces and degeneracies are monoid homomorphisms because concatenation is a bifunctor and is strictly associative and unital. This statement concerns the whole degreewise free monoid. Its distinguished basis need not itself be a simplicial set: $d_0$ can send one generator to a product of several generators, or to the unit.
\end{dcomment}


\printedpage{38}
% head: A. Grothendieck : Dérivateurs, Ch. VII, Édition des Universités de Montpellier II et Paris VII
Le cas $g=\mathrm{id}$ correspond au cas où
\[
\pi_i(j)=1 \qquad\quad \forall j\in[1,l_i],
\]
les $\mathfrak{f}_{\varepsilon_i,l_i,\pi_i}$ sont alors détérminés en fonction des $\mathfrak{f}_{\varepsilon_i,l_i}$, et l'application de dégénérescence est % [sic] détérminés -> déterminés, accent
\[
\mathfrak{f}_{\varepsilon,l}\longmapsto\mathfrak{f}_{\varepsilon,l,1_l} \qquad\quad (\text{où } 1_l \text{ est la fonction constante de valeur 1 sur } [1,l]).
\]
Le $f=\mathrm{id}$ est celui où les $l_i$ sont égaux à 1. Alors chaque $\pi_i : [1,1]\longrightarrow\mathbf{N}$ peut être vu comme un élément de $\mathbf{N}$, la donnée de $(\varepsilon_i,l_i,\pi_i)$ revient à celle des $(\varepsilon_i,n_i)\in\{\pm1\}\times\mathbf{N}$, i.e. la donnée % [sic] Le f = id -> Le cas f = id, word « cas » missing
\msp{page 56}
d'une suite de $n$ générateurs parmi les $\mathfrak{f}_{\varepsilon,l}$. L'application de dégénérescence $\mathrm{Nerf}_1\longrightarrow\mathrm{Nerf}_2$ est ici donnée sur les générateurs par
\[
\mathfrak{f}_{\varepsilon,l}\longmapsto\mathfrak{f}_{\varepsilon,1,l}.
\]
Heureusement, on n'a plus à prouver l'associativité de la composition.

\textbf{NB} La simplicité de ces descriptions renforce mon sentiment que la catégorie $\uTy$ définie à présent, avec les \textit{morphismes d'intervalles orientés croissants} $I_\tau\longrightarrow I_{\tau'}$ comme flèches de $\uTy$, est bien `la bonne'. On va déjà voir si elle donne une description sympa des catégories de fractions $X\Sigma^{-1}$, sans tortillements comme précédemment.

Mais on a promis de déterminer $\mathrm{Nerf}_r(\uTy)$, pour tout $r\in\mathbf{N}$. Le générateur général de $\mathrm{Nerf}_r$ est de la forme (pour $\varepsilon\in\{\pm1\}$ fixé)
\[
\underbrace{\Delta_\varepsilon}_{=\,\Delta^{l_0}_\varepsilon\text{avec } l_0=1}\xrightarrow{\ f_1\ }\Delta^{l_1}_\varepsilon\xrightarrow{\ f_2\ }\Delta^{l_2}_\varepsilon\longrightarrow\ \cdots\ \xrightarrow{\ f_r\ }\Delta^{l_r}_\varepsilon,
\]
où les $f_i : \Delta^{l_{i-1}}_\varepsilon\longrightarrow\Delta^{l_i}_\varepsilon$, $i\in[1,r]$, sont donnés par des applications
\[
\pi_{i-1} : [1,l_{i-1}]\longrightarrow\mathbf{N} \qquad\quad \text{telles que }\ \underbrace{\mathrm{long}(\pi_{i-1})}_{\overset{\text{déf}}{=}\ \Sigma_{j\in[1,l_i]}\pi_i(j)}\ = l_i. % [sic] \Sigma_{j\in[1,l_i]}\pi_i(j) -> \Sigma_{j\in[1,l_{i-1}]}\pi_{i-1}(j), indices shifted with respect to long(\pi_{i-1})
\]
(NB $\pi_0 : [1,1]=\{1\}\longrightarrow\mathbf{N}$ est identifié à $\mathrm{long}(\pi_0)=\pi_0(1)=l_1$.)

\msp{page 57}
Pour ne pas faire dépendre la définition du domaine des $\pi_i$ ($i\in[0,r-1]$) des entiers $l_i$, on peut aussi considérer les $\pi_i$ comme des applications
\[
\mathbf{N}^*\xrightarrow{\ \pi\ }\mathbf{N} \qquad\quad (\mathbf{N}^*=\{i\in\mathbf{N}\mid i\geq 1\})
\]
à supports finis (i.e. tels que $\pi(i)=0$ pour $i$ grand), et exiger de façon précise
\[
\begin{aligned}
&\mathrm{supp}\,\pi_0\ \subseteq\ \{1\}\\
&\mathrm{supp}\,\pi_1\ \subseteq\ [1,\mathrm{long}(\pi_0)]\\
&\qquad\ \cdots\\
&\boxed{\mathrm{supp}\,\pi_i\ \subseteq\ [1,\mathrm{long}(\pi_{i-1})]}, \qquad\quad \text{pour }\ \boxed{i\in[0,r-1]}\quad .
\end{aligned}
\tag*{$(**)$}
\]
% foot: Version 22 décembre 2001


\dgrblock{dgrb313}{All nerve degrees as ordered forests}
\begin{dcomment}{exposition}{p57}{All nerve degrees as ordered forests}
The final assertion of manuscript page~57 is that the same decomposition works for arbitrarily long strings of reparametrizations. It is useful to complete the argument now, since it explains both the support conditions at the bottom of the page and the one face map that behaves differently from the others.

Fix $r\geq0$. Call an $r$-step string \emph{rooted} when its initial type has one transition. Every intermediate type is then a power of the same sign, and the string has the form
\[
 \Delta_\varepsilon^{l_0}\longrightarrow
 \Delta_\varepsilon^{l_1}\longrightarrow\cdots\longrightarrow
 \Delta_\varepsilon^{l_r},\qquad l_0=1.
\]
For $0\leq a<r$ its $a$th weight vector is
\[
 \pi_a:[1,l_a]\longrightarrow\mathbb N,
 \qquad l_{a+1}=\sum_{j=1}^{l_a}\pi_a(j).
\]
One can regard this as an ordered rooted tree with $r$ levels of edges: a transition at level $a$ has $\pi_a(j)$ ordered descendants at level $a+1$. A vertex may have no descendants. Once a whole level is empty, all subsequent levels are empty because there is no morphism from the point to a nonpoint type.

For a general string, let its initial word have transitions $1,\ldots,n$. At the first stage the descendants of each original transition form a consecutive monochromatic block. At the next stage the descendants of all transitions in one such block again form a consecutive block: concatenate their individual image blocks, allowing zero lengths. Induction gives such a block at every level. Restricting all maps to the descendants of a fixed original edge produces a rooted string, and concatenating the $n$ rooted strings reconstructs the original string. Each boundary is determined by cumulative weights, so this factorization is unique. A rooted string cannot be a product of two nonunit strings, since the lengths of their initial words would be positive integers summing to one. Conversely a string whose initial word has at least two edges splits after its first edge. Thus rooted strings are exactly the indecomposables, and
\[
 N_r(\uTy)\cong\text{the free monoid on rooted $r$-step strings}.
\]
For $r=0$ there are exactly two generators, the two signs. For an empty initial word there is only the all-identity string at the point, which is the monoid unit. These boundary cases establish the assertion for every $r$.

The simplicial operators can now be read directly on these forests. An inner face removes one level and replaces the two successive weight vectors by the sums over consecutive descendant blocks. The last face forgets the last level. A degeneracy repeats a level, inserting weight $1$ at each of its transitions. These operations preserve a single root. The first face removes the original root level: a rooted string with $l_1$ first-generation transitions becomes the concatenation of their $l_1$ rooted descendant strings. This is the higher-dimensional version of the formula for $d_0$ above. Since these are exactly restriction and composition of actual functors $[r]\to\uTy$, they satisfy all simplicial identities; the unique forest decomposition identifies their formulas with those categorical operators.

For $r>0$, to use fixed-domain functions as on manuscript page~57, extend each $\pi_a$ by zero to $\mathbb N^*$. The exact conditions are
\[
 \operatorname{supp}(\pi_0)\subseteq\{1\},\qquad
 \operatorname{supp}(\pi_a)\subseteq
       [1,\operatorname{long}(\pi_{a-1})]
       \quad(1\leq a<r).
\]
The source's final boxed range starts at $a=0$, although $\pi_{-1}$ is undefined; it must begin at $1$. The preceding display defining $\operatorname{long}(\pi_{i-1})$ likewise needs the indices $j\in[1,l_{i-1}]$ and $\pi_{i-1}(j)$. These printed indexing slips do not affect the argument once the successive domains are kept in view.

One should not infer from this result that the simplicial monoid is freely generated by a simplicial set of indecomposables: the first-face calculation rules that out for this particular basis. Nor does the algebraic freeness make its classifying space complicated. The point is terminal in $\uTy$, since every weight vector admits the unique further collapse to zero. Appending that terminal object gives the usual contraction of the nerve. The free presentations describe how all reparametrizations fit together, while this terminal object describes their ordinary homotopy type.
\end{dcomment}


\dgrblock{dgrb314}{What the shape calculus does for localization}
\begin{dcomment}{status}{p57}{What the shape calculus does for localization}
The free monoids just constructed organize subdivision and concatenation. To obtain mapping spaces, a second ingredient is needed: comparisons of labelled paths must carry the specified weak-equivalence information. For a relative category $(C,W)$ and fixed $x,y$, let $Z_\tau(C,W)_{x,y}$ be the category of paths of type $\tau$ from $x$ to $y$, with backward arrows in $W$, and natural transformations whose components lie in $W$ and whose endpoint components are identities. Precomposition makes these categories contravariant in $\uTy$. Their nerves, or their groupoid completions, retain comparisons that the set of path components alone would discard.

The analogous construction for relative infinity-categories is developed by \citet[Definitions~3.9 and~3.17]{MazelGee2018Hammocks}. Under a homotopical three-arrow calculus, his Theorem~4.4 identifies the geometric realizations of its hammock mapping objects with the groupoid completions of three-arrow zigzag categories. This is an additional hypothesis, not a consequence of the finite normal form in the manuscript. Independently, \citet[Theorem~3.8]{MazelGee2019Rezk} proves that completing the relative Rezk nerve presents localization. That universal property supplies the criterion with which a proposed path model must be compared.

For the general mapping-space comparison, \citet[Theorem~1.9 and Corollary~1.11]{ArakawaCnossen2026} prove in their preprint that
\[
 \operatorname{Map}_{C[W^{-1}]}(x,y)
 \simeq \operatorname*{colim}_{\tau\in\uTy^{\mathrm{op}}}
          \bigl|Z_\tau(C,W)_{x,y}\bigr|,
\]
where the colimit is taken in spaces and $|{-}|$ is groupoid completion. The displayed formula uses the ordinary relative categories just defined; the passage from their strict endpoint fibres to the invariant formula is supplied by Convention~4.2, Definition~4.3, Theorem~4.7, and Proposition~4.8(b) of the same paper. The general infinity-categorical statement uses the corresponding infinity-categories of relative diagrams and their endpoint fibres. Their Corollary~5.7 supplies the localization identification for Mazel-Gee's hammock construction. These references distinguish the established Rezk-nerve theorem from the more recent explicit comparison.

The comparison makes the scope of the present pages clear. They provide the shapes, their functorial changes, and strict concatenation. They do not yet choose $W$ or impose the conditions on vertical arrows used above. No new open conjecture is required to finish their finite combinatorics. The subsequent localization construction still has to be checked with its own morphisms and universal property; the contractibility of an interval, by itself, does not perform that check.
\end{dcomment}
\begin{otherlanguage}{english}
\dgroverview{VII-a}{Chapter VII \textperiodcentered\ \textup{VII-a} \textperiodcentered\ Overview}
\selectlanguage{english}
\noindent\textbf{Chapter VII, manuscript pages 1--57: the shapes of paths.}

The opening of Chapter VII is a restart. Grothendieck has found the earlier path formalism too cumbersome to serve as a foundation for localization, and he now separates a path from the finite combinatorial object on which it is parametrized. The distinction gives these pages their unity. First one understands intervals independently of any coefficient category; next one regards every operation on paths as a change of interval; finally one assembles all types and makes their concatenation strictly functorial. The later theory of fractions is a test of the construction, not yet a theorem established in these pages. The commentary on localization and path completion in Chapter~VI, especially at manuscript pages~64--70, supplies the background; the present chapter changes the finite apparatus on which such arguments are to rest.

An interval of length $n$ is the free category on an oriented path with vertices $0,\ldots,n$. Its signature is a word $\tau=(\varepsilon_1,\ldots,\varepsilon_n)$ in two signs. There are two orders to remember. The progression order is the linear order of the vertices; the categorical order records the existence of arrows. A path proceeds from vertex~$0$ to vertex~$n$ even when some of its arrows point backwards. Reversing this progression order changes the signature to $(-\varepsilon_n,\ldots,-\varepsilon_1)$, whereas taking the opposite category changes it to $(-\varepsilon_1,\ldots,-\varepsilon_n)$. Their distinction is already needed to understand why a span is symmetric as an interval but is not isomorphic to its opposite cospan.

The maximal constant runs of the signature are the structural chains. Their endpoints are the critical vertices; retaining just these vertices gives the reduced interval. These notions organize all full subcategories of an interval. A nonempty full subcategory has interval components, and a subset with extreme vertices $i<j$ is connected precisely when it retains every critical vertex between them. The proof is elementary and decisive: two consecutive retained vertices can be joined by an arrow exactly when no change of sign has been omitted between them. Thus reduction remembers the turns while suppressing subdivisions inside a composable run. It does not mean that arbitrary diagrams on the interval may forget the intermediate objects.

Several literal statements near the beginning need repair. Proposition~1(b) asserts that every endpoint-fixing endofunctor is the identity. On $[2]$, the function $(0,1,2)\mapsto(0,0,2)$ is a counterexample. The proof on the following pages establishes rigidity of endpoint-fixing \emph{isomorphisms}, which is sufficient for the classification of orientations. In fact, the stronger endofunctor assertion holds exactly for reduced intervals. Proposition~3's deletion criterion also needs the vertices incomparable with the deleted vertex to be separated from its incoming and outgoing links. Finally, an interval-shaped subcategory need not be full. All these readings occur in the edition and remain in the source text. The accompanying blocks give counterexamples and prove the repaired assertions, so that subsequent correct arguments do not inherit an unnecessary uncertainty.

Reduction has a useful additional interpretation. For the full inclusion $j:I_{\mathrm{red}}\hookrightarrow I$, send an interior vertex of each run to that run's categorical target or categorical source, and fix all critical vertices. The resulting maps satisfy
\[
 p_{\mathrm{target}}\dashv j\dashv p_{\mathrm{source}}.
\]
The commentary proves these adjunctions by comparing the relevant hom-sets. They give categorical deformation retractions and explain exactly why reduction is harmless for groupoid-valued paths. For a general category $X$, restriction $X^I\to X^{I_{\mathrm{red}}}$ need not be an equivalence. This is the first place where the geometry of the shape and the information in a labelled diagram must be kept distinct.

Once the shape is fixed, a path of type $\tau$ in $X$ is a functor $I_\tau\to X$, and its morphisms are natural transformations. The category is $X^{I_\tau}$; its set of objects is the represented functor $\Hom_{\Cat}(I_\tau,X)$. Yoneda then says that every operation on paths natural in the coefficient category is precomposition by a unique functor between their shapes, in the reverse direction. The proof evaluates the operation on the identity path in $I_\tau$. This explains both composition of successive arrows and insertion of identities, and supplies their action on morphisms without an additional coherence choice. The manuscript's distinction between the set and the category of paths must therefore survive the typography.

The category of types changes twice while Grothendieck tests it. Initially it permits all functors between intervals; then it retains functors preserving the two endpoints; finally it also requires their vertex functions to be nondecreasing in the progression order. Only the first version is fully faithful in $\Cat$. The second is fully faithful in the category of categories under the discrete two-point category. The final version, denoted $\uTy$, is generally not full there either. These are successive definitions, not contradictory assertions about a fixed category. The final restriction preserves directed progress through a path and excludes changes of parameter which backtrack.

For any chosen version $T$, the strict contravariant functor $\tau\mapsto X^{I_\tau}$ gives a split Grothendieck fibration over $T$. A morphism from $(\tau',c')$ to $(\tau,c)$ is a pair
\[
 \varphi:\tau'\to\tau,
 \qquad v:c'\Rightarrow c\varphi.
\]
Composition is whiskering followed by composition of transformations. This is the machinery of \citet[Expos\'e~VI, \S\S8--10]{GrothendieckSGA1}, specialized to restriction along finite shapes. Endpoint-preserving shape maps make evaluation at the two ends into functors on the total category. The commentary also determines the maximal enlargement for which natural endpoint evaluations can still exist, allowing comparison arrows at the ends. That enlargement illustrates why the subsequent requirement of strict gluing is stronger. No colimits in $X$ are needed for the fibration; extending paths in the opposite direction would require Kan extensions, and already a collapsed span asks for a pushout.

The main calculation of the second half has a particularly short modern form. Give the transition set $\Tr(\tau)=\{1,\ldots,n\}$ its order and its signs. If $f:I_\tau\to I_\upsilon$ is a morphism of the final type category, every target transition belongs to the image interval of exactly one source transition. This defines a monotone, sign-preserving map
\[
 f^*:\Tr(\upsilon)\longrightarrow\Tr(\tau).
\]
Conversely, given such a map $h$, recover the vertex function by
\[
 f(i)=\#\{j:h(j)\leq i\}.
\]
These assignments are inverse and reverse composition. Thus the type category is opposite to the category of finite signed linear orders with monotone sign-preserving maps. Its many normal forms are descriptions of this one correspondence.

For a reduced target, $f^*$ is injective, so a map is specified by a signed subsequence of the source word. In general, a map out of a fixed source is specified by nonnegative integer weights: an edge of sign $\varepsilon$ and weight $a$ becomes a run of $a$ edges of sign $\varepsilon$. Weight zero collapses that edge. A map into a fixed target is instead described by its image vertices and by the signatures of the fibers over them. The latter description must use cumulative sums for the target coordinates, and collapsed fibers are components of the graph of collapsed edges, not components of the union of their endpoint vertices with every ambient adjacency restored. The point target is treated separately. These corrections preserve the two classification results.

Concatenation glues the final vertex of one interval to the initial vertex of the next. On standard words and endpoint-preserving maps it is strictly associative and unital, and makes $\uTy$ a strict monoidal category. It cannot extend with this object operation to the category allowing \emph{all} interval functors: there $+$ and $-$ are isomorphic, whereas $++$ and $-+$ are not. The restriction on morphisms is therefore essential to the source's own construction. Through the signed-edge description, $\uTy$ is the free strict monoidal category on two counital coassociative comonoids. The one-generator background is the opposite of the augmented simplex category; see \citet[Appendix~D, Proposition~D.2]{AguiarMahajan2010}. The two-generator version follows by ordered concatenation, without a symmetry interchanging the letters.

The nerve calculation is now transparent. A chain of $r$ composable type maps splits uniquely according to the edges of its initial word. Each factor begins at a single signed edge; its later stages are successive ordered expansions of that edge, allowing complete disappearance. Consequently every $N_r(\uTy)$ is a free monoid. The face deleting the initial object can send a generator to several generators, or to the empty word, so this assertion does not identify the nerve with the free simplicial monoid on its degreewise generators. The commentary gives every face and degeneracy by the corresponding operation on ordered forests and resolves the final support-indexing slip.

The later comparison can be made at the level of the indexing category itself. Marking the negative transitions identifies $\uTy$ with the zigzag type category in \citet[Definition~3.5]{MazelGee2018Hammocks}. The closing status block distinguishes this finite identification from the additional construction of localization mapping spaces, including the recent comparison in \citet[Theorem~1.9 and Corollary~1.11]{ArakawaCnossen2026}. The development blocks prove explicit consequences left unworked here: the reduction adjunctions and sharp rigidity criterion, the endpoint evaluation criterion, and a criterion with a counting formula for sections of contractions. Their markers record our proofs and limited literature searches, not a claim of priority. The unresolved task at manuscript page~57 is the promised construction of fractions with its precise morphisms; the finite combinatorics considered in these pages has been settled in the commentary.
\end{otherlanguage}
\dgrchunkend

% DER_VII-b: Grothendieck's text with the commentary set in, then the overview, inside a group that applies the chunk's own definitions
\dgrpart{1}{VII-b}{Part 2 of 3, manuscript pages 58--117}{VII-b\enspace\textperiodcentered\enspace Chapter VII, part 2 of 3\enspace\textperiodcentered\enspace manuscript pages 58--117}{Chapter VII \textperiodcentered\ \textup{VII-b}}
\dgrchunkbegin
% DER_VII-b: the definitions of the chunk's preamble, applied inside its group
\providecommand\IfDocumentMetadataT[1]{}
\let\printedpage\relax\newcommand{\printedpage}[1]{\marginpar{\footnotesize p.~#1}}
\let\msp\relax\newcommand{\msp}[1]{\marginpar{\vspace*{10pt}\footnotesize [#1]}}
\let\uTy\relax\newcommand{\uTy}{\underline{\mathrm{Ty}}}
\let\Ch\relax\newcommand{\Ch}{\underline{\mathrm{Ch}}}
\let\Ty\relax\newcommand{\Ty}{\underline{\mathrm{Ty}}}
\let\Nerf\relax\newcommand{\Nerf}{\mathrm{Nerf}}
\let\Fl\relax\newcommand{\Fl}{\mathrm{Fl}}
\let\Hom\relax\newcommand{\Hom}{\mathrm{Hom}}
\let\ch\relax\newcommand{\ch}{\mathrm{ch}}
\let\hot\relax\newcommand{\hot}{\mathrm{hot}}
\let\id\relax\newcommand{\id}{\mathrm{id}}
\let\CH\relax\newcommand{\CH}{\underline{\mathrm{CH}}}
\let\red\relax\newcommand{\red}{\mathrm{r\acute{e}d}}
\let\uHom\relax\newcommand{\uHom}{\underline{\mathrm{Hom}}}
\let\ul\relax\newcommand{\ul}[1]{\underline{#1}}
\let\triv\relax\newcommand{\triv}{\mathrm{triv}}
\let\uCh\relax\newcommand{\uCh}{\underline{\mathrm{Ch}}}
\let\Cat\relax\newcommand{\Cat}{\mathrm{Cat}}
\let\Cathot\relax\newcommand{\Cathot}{\mathrm{Cat\,hot}}
\let\Hot\relax\newcommand{\Hot}{\mathrm{Hot}}
\let\Ob\relax\newcommand{\Ob}{\mathrm{Ob}\,}
\let\sHom\relax\newcommand{\sHom}{\underline{\mathcal{H}\mathit{om}}}
\let\To\relax\newcommand{\To}{T_{\mathrm{o}}}
\let\circled\relax\newcommand{\circled}[1]{\tikz[baseline=(c.base)]\node[draw,circle,inner sep=1pt](c){#1};}
\let\Hotz\relax\newcommand{\Hotz}{\mathrm{Hot}^{\wedge}_{0}}
\let\CNo\relax\newcommand{\CNo}{C_{\mathbf{N}^{o}}}
\let\lab\relax\newcommand{\lab}[2][l]{\text{\footnotesize\begin{tabular}{@{}#1@{}}#2\end{tabular}}}

% ---- from Der-ChVII-456_p041-045.tex
\printedpage{39}
% head: A. Grothendieck : Dérivateurs, Ch. VII, Édition des Universités de Montpellier II et Paris VII
D'autre part, posant $l_i = \mathrm{long}(\pi_{i-1})$ (pour $i \in [1,r]$),
\[
f_i : \Delta_\varepsilon^{l_{i-1}} \longrightarrow \Delta_\varepsilon^{l_i\ (=\mathrm{long}(\pi_{i-1}))}
\]
est défini comme
\[
f_i = \mathfrak{f}_{\varepsilon,\pi_{i-1}(1)} * \mathfrak{f}_{\varepsilon,\pi_{i-1}(2)} * \cdots * \mathfrak{f}_{\varepsilon,\pi_{i-1}(l_{i-1})},
\]
de sorte que l'élément $(*)$ de $\Nerf_r(\Ty)$ est déterminé de façon unique par la donnée des $\varepsilon \in \{\pm 1\}$ et de
\[
\pi_\bullet = (\pi_i)_{i\in[0,r-1]} \in (\mathrm{Hom}_{\mathrm{Ens}}(\mathbf{N}^*,\mathbf{N}))^r
\]
satisfaisant les conditions $(**)$. Soit $F_{\varepsilon,\pi_\bullet}$ cet élément de $\Nerf_r(\Ty)$, et soit $N_r$ l'ensemble des $\pi_\bullet$ qu'on vient d'expliciter. Alors
\[
(F_{\varepsilon,\pi_\bullet})_{(\varepsilon,\pi_\bullet)\in\{\pm1\}\times N_r}
\]
est un système de générateurs libre du monoïde $\Nerf_r(\Ty)$.

Si $l_i = 0$ (i.e. \fbox{$\pi_{i-1} = 0$}), on doit avoir $\mathrm{supp}(\pi_{i+1}) = \emptyset$, i.e. $\pi_{i+1} = 0$, d'où $l_{i+1} = 0$, et de même pour les $l_j$ et $\pi_j$ suivants. Géométriquement l'hypothèse s'exprime par la condition que le but
\msp{page 58}
de $f_i$ est l'élément 1 de Ty (correspondant à l'intervalle ponctuel), et comme $\mathrm{Hom}(1,\tau) = \emptyset$ pour $\tau \neq 1$, cela signifie que les flèches qui suivent $f_i$ doivent être toutes égales à la flèche $\mathrm{id}_1 : 1 \longrightarrow 1$.

\dgrblock{dgrb315}{Degreewise free is a precise assertion}
\begin{dcomment}{exposition}{p57}{Degreewise free is a precise assertion}
Grothendieck has now made the combinatorics of a changing path shape completely explicit. The commentary at manuscript pages~53--57 already explains the free generators in every degree and counts refinement stages. One additional distinction helps us use that calculation correctly: freeness in each degree does not mean that the intrinsic generators themselves form a simplicial set.

To see the distinction, take the generator of degree two represented by
\[
 \Delta_+\longrightarrow\Delta_+^2
       \xrightarrow{\mathrm{id}}\Delta_+^2.
\]
Its first face, which forgets the initial object and first arrow, is
$\mathrm{id}_{\Delta_+^2}=f_{+,1}*f_{+,1}$. This is a product of two nonunit elements in the degree-one monoid. It is therefore not a generator. In the notation of the text, the general formula is
\[
 d_0(f_{\varepsilon,l,\pi})
    =f_{\varepsilon,\pi(1)}*\cdots*f_{\varepsilon,\pi(l)},
\]
where the empty product for $l=0$ is the unit. An isomorphism between free monoids preserves their intrinsic indecomposables. Consequently $N_\bullet\underline{\mathrm{Ty}}$ cannot be the degreewise free monoid on a simplicial set: such a construction would send each generator under every face to a generator. The monoid structure maps are determined by generators, as Grothendieck says; their images can be arbitrary words.

There is also a useful boundary case in his finite-support description. The condition on $\pi_0$ is separate, namely $\operatorname{supp}\pi_0\subseteq\{1\}$. The recursive condition must then run over $1\leq i\leq r-1$, since $\pi_{-1}$ has not been defined. If the length at stage $i$ is zero, the next weighting is $\pi_i=0$, and every later stage remains the empty interval. The edition instead writes $\pi_{i+1}$ at this step. Both index slips are retained in the source and registered in the corrections table. They change neither the free decomposition nor its geometric meaning: a path shape which has disappeared cannot reappear under an endpoint-preserving interval map.
\end{dcomment}


\textbf{Remarque.} Dans la description combinatoire de $\Ty$ et de $\Nerf_\bullet(\Ty)$, l'ensemble des `signatures' $\{\pm1\}$ joue un rôle assez bidon. Il peut être remplacé par n'importe quel ensemble $\Sigma$ (dit `ensemble des signatures'), et on construit alors une catégorie monoïdale (ou un monoïde catégorique) ayant comme ensemble d'objets le monoïde libre engendré par $\Sigma$, interprétés comme des intervalles d'entiers $[0,n]$, munis de `signatures' dans $\Sigma$ pour chaque intervalle de transition, ou ce qui revient au même, des `graphes' orientés $\Sigma$-pondérés de type combinatoire un intervalle $\boldsymbol{\Delta}_i$, ce qui revient à la donnée d'une suite (ordonnée) de $n$ éléments $\varepsilon_1,\ \varepsilon_2,\ \ldots\ ,\ \varepsilon_n \in \Sigma$, i.e. d'un `mot' $\varepsilon_1\varepsilon_2\ldots\varepsilon_n$, qu'on préfère écrire $\Delta_{\varepsilon_1}\Delta_{\varepsilon_2}\cdots\Delta_{\varepsilon_n}$. On peut définir de façon évidente les homomorphismes \textit{croissants} de tels graphes pondérés orientés (sous-entendu : avec l'élément source et l'élément but dans la structure, à respecter par les morphismes), sans avoir besoin d'une interprétation catégorique. La fonctorialité de l'opération de composition des intervalles pondérés, pour cette notion de morphisme, est évidente. Le
\msp{page 59}
choix $\Sigma = \{\pm1\}$ est à ma connaissance le seul utile, à cause de son lien avec la notion de chemins et le calcul de $\pi_0$ d'un objet de $X$ \texttt{[plutôt d'un objet} $X$ \texttt{de} Cat\texttt{]}, etc.
% foot: Version 22 décembre 2001

\dgrblock{dgrb316}{The localization test}
\begin{dcomment}{exposition}{p59:le seul utile, à cause de son lien}{The localization test}
Grothendieck is about to test his category of types against a construction whose universal property is already settled: the category obtained by formally inverting specified arrows. This is a demanding first test because literal concatenation gives too much information. It remembers the length of a zigzag, every inserted identity, and every subdivision of a composite. A localization must identify presentations which differ in these ways, while retaining their common source and target.

Let $P(X)$ denote the category whose arrows are literal finite zigzags in $X$. Length is additive under composition. Thus an invertible arrow of $P(X)$ must have length zero. Moreover the assignment $u\mapsto c(u)$ from arrows of $X$ to one-edge paths does not preserve composition: $c(vu)$ has one edge and $c(v)c(u)$ has two. It does not preserve identities either, since $c(\mathrm{id}_a)$ has length one. These observations concern the literal path construction uniformly in $X$. For particular categories there can of course be useful functors in one direction: if $X$ is a groupoid, every zigzag has an effect in $X$, and evaluation is a functor $P(X)\to X$.

The category $\underline{\mathrm{Ch}}(a,b)$ supplies the missing identifications as morphisms. A type map replaces a string of composable arrows by its composite or inserts identity arrows through pullback. A natural transformation compares two paths of the same type by a ladder of commuting squares. Since the endpoints are fixed, the components at $a$ and $b$ must be identities. Passing to $\pi_0$ identifies objects joined by a finite zigzag of these morphisms. A single directed natural transformation is enough for this purpose; it need not be invertible.

The desired category therefore has
\[
 \operatorname{Hom}(a,b)=\pi_0\underline{\mathrm{Ch}}(a,b).
\]
Its composition comes from concatenation before taking components. This is legitimate because $\pi_0(A\times B)=\pi_0(A)\times\pi_0(B)$ for categories: a zigzag in a product projects to the factors, while zigzags in the factors can be traversed one coordinate at a time. Associativity and units are then inherited, rather than imposed as further equivalence relations.

The relative version allows only a chosen class $\Sigma$ of arrows to be traversed backwards. Its aim is the formula
\[
 \pi_0\underline{\mathrm{Ch}}_\Sigma(a,b)
       \cong \operatorname{Hom}_{X[\Sigma^{-1}]}(a,b).
\]
At this stage the assertion concerns ordinary hom-sets. The homotopy type of the whole path category is a separate, stronger question, and the later pages of the chapter turn precisely to the additional choices and coherence needed for it.
\end{dcomment}


\printedpage{40}
% head: A. Grothendieck : Dérivateurs, Ch. VII, Édition des Universités de Montpellier II et Paris VII
\section*{7\quad Explicitation des catégories de fractions $X\Sigma^{-1}$ en termes des $\Ch(X;a,b)$}

Si $X$ est une catégorie, on désigne par $\Ch(X)$ la catégorie des chemins dans $X$ (cf. §4), et pour $a,b \in \mathrm{Ob}\,X$, par
\[
\Ch(X;a,b) \qquad \text{ou simplement } \Ch(a,b)
\]
celle des chemins de $a$ vers $b$, i.e. de source $a$, but $b$. L'intuition géométrique (qu'on explicitera), c'est qu'un tel chemin $c$ définit un chemin `direct' de $a$ vers $b$ (i.e. un chemin de type $\tau'$ `direct', i.e. toutes les flèches de transition sont directes, i.e. $\tau'$ de la forme $\Delta_+^n$), de même longueur que le chemin initial $c$, obtenu en inversant toutes les flèches rétrogrades dans le chemin $c$. Cette opération est effectivement possible dans $X$ lui-même si et seulement si lesdites flèches sont inversibles. On obtient alors une suite composable de flèches, dont le composé total (qu'on peut appeler `l'effet utile' du chemin)
\msp{page 60}
est une flèche de $a$ vers $b$. Dans le cas général, les chemins de $a$ vers $b$ peuvent être vus comme une sorte de `flèches généralisées' allant de $a$ vers $b$, dont la composition serait la composition évidente, et les identités les `chemins vides' relatifs aux $a \in \mathrm{Ob}\,X$. On trouve d'ailleurs bien ainsi une catégorie, ayant mêmes objets que $X$, mais où les
\[
\mathrm{Hom}_{\mathrm{Ch}}(a,b) = \mathrm{Ch}(a,b) = \mathrm{Ob}\,\Ch(a,b)
\]
sont prohibitivement gros. Cela s'exprime notamment par le fait que les seules flèches inversibles dans $X_{\mathrm{Ch}}$ sont les identités ! On n'a pas de foncteur naturel entre $X_{\mathrm{Ch}}$ et $X$ qui soit l'identité sur les objets. Car à un chemin $a \xrightarrow{c} b$ on ne peut associer de $X$-morphismes $a \longrightarrow b$ (sauf si $c$ \texttt{[est]} de type $\Delta_+$), et à un $X$-morphisme $a \xrightarrow{u} b$ on peut associer un chemin $c(u)$ de type $\Delta_+$, mais si $a \xrightarrow{u} b \xrightarrow{v} c$, $c(vu)$ est encore de type $\Delta_+$ alors que $c(v)c(u)$ est de type $\Delta_+^2$, ils ne sont pas égaux. Par contre, il y a une flèche naturelle entre ces deux éléments de $\Ch(a,c)$, flèche au dessus de l'unique homomorphisme
\[
\Delta_+ \xrightarrow{\mathfrak{f}_{+,2}} \Delta_+^2.
\]
\msp{page 61}
On a alors
\[
\mathfrak{f}_{+,2}^*(\underbrace{a \xrightarrow{u} b \xrightarrow{v} c}_{c(v)c(u)}) = (\underbrace{a \xrightarrow{vu} c}_{c(vu)}),
\]
d'où une flèche canonique
\[
c(vu) \longrightarrow c(v)c(u) \qquad \text{dans } \Ch(a,c).
\]
Cela montre que $c(vu)$ et $c(v)c(u)$ sont tout au moins dans une même composante connexe de $\Ch(a,c)$, ou (comme on dit encore) qu'ils sont `homotopes'. Cela suggère ici de modifier la définition de $\mathrm{Hom}_{\mathrm{Ch}}(a,b)$ en
\[
\pi(a,b) = \pi_0(\Ch(a,b)).
\]
% foot: Version 22 décembre 2001

\printedpage{41}
% head: A. Grothendieck : Dérivateurs, Ch. VII, Édition des Universités de Montpellier II et Paris VII
Le foncteur de composition des chemins
\[
\begin{array}{ccccc}
\Ch(a,b) & \times & \Ch(b,c) & \longrightarrow & \Ch(a,c)\\
(c & , & c') & \longmapsto & c' \circ c
\end{array}
\]
\texttt{[il y a trop de} $c$\texttt{]} (\footnote[17]{Ce foncteur est déduit des foncteurs
\[
\Ch_\tau(a,b) \times \Ch_{\tau'}(b,c) \longrightarrow \Ch_{\tau\tau'}(a,c)
\]
évidents, cf. page 77.}) induit une application sur les $\pi_0$, ce qui donne une loi de composition
\[
\pi(a,b) \times \pi(b,c) \longrightarrow \pi(a,c)
\]
sur les flèches. Il est immédiat qu'elle est associative (la composition des chemins l'étant), et unitaire. Elle admet comme flèche identique en $x$ soit le `lacet vide en $x$', lequel est (directement) homotope dans $\Ch(x,x)$ au chemin `dégénéré' $x \xrightarrow{\mathrm{id}} x$
\msp{page 62}
de type $\Delta_+$, ou aussi au chemin dégénéré de n'importe quel type $\tau$, disons
\[
x \xleftarrow{\mathrm{id}} x \xrightarrow{\mathrm{id}} x \xrightarrow{\mathrm{id}} x \xleftarrow{\mathrm{id}} x \xleftarrow{\mathrm{id}} x,
\]
un tel chemin n'étant autre que $p_\tau^*(x)$, où
\[
p_\tau : \tau \longrightarrow 1
\]
est le morphisme de dégénérescence (l'unique morphisme dans $\Ty$ vers 1, qui est objet final), d'où
\[
p_\tau^* : \Ch_1(x,x) = \{x\} \longrightarrow \Ch_\tau(x,x) \subseteq \Ch(x,x).
\]
La catégorie $\Pi(X)$ ainsi obtenu doit s'appeler le \textit{groupoïde fondamental} de $X$. En effet, % [sic] obtenu -> obtenue, agreement with « La catégorie »

\dgrblock{dgrb317}{A span and its fraction}
\begin{dcomment}{examples}{p61}{A span and its fraction}
The construction becomes concrete as soon as reversing one arrow creates a morphism which the original category did not contain. Consider the category $X$ generated by the span
\[
 a\xleftarrow{s}x\xrightarrow{t}b
\]
with no nonidentity arrows besides $s$ and $t$, and put
$\Sigma=\{\mathrm{id}_a,\mathrm{id}_x,\mathrm{id}_b,s\}$. The zigzag from $a$ to $b$ is a $\Sigma$-path although $\operatorname{Hom}_X(a,b)$ is empty. Its effect after localization is the fraction $ts^{-1}$.

We can compute the entire localization. It has the same three objects and the following hom-sets, with rows denoting sources:
\[
\begin{array}{c|ccc}
 &a&x&b\\ \hline
 a&\{\mathrm{id}_a\}&\{s^{-1}\}&\{ts^{-1}\}\\
 x&\{s\}&\{\mathrm{id}_x\}&\{t\}\\
 b&\varnothing&\varnothing&\{\mathrm{id}_b\}.
\end{array}
\]
Composition is determined by $s^{-1}s=\mathrm{id}_x$,
$ss^{-1}=\mathrm{id}_a$, and $(ts^{-1})s=t$. For any functor $F:X\to Y$ with $F(s)$ invertible, there is exactly one extension to this category: it sends $s^{-1}$ to $F(s)^{-1}$ and $ts^{-1}$ to $F(t)F(s)^{-1}$. This verifies the universal property directly.

The same calculation can be read inside the path categories. A path starting at $b$ can only stay at $b$, because $t$ is not permitted as a backward arrow. A nonconstant path ending at $b$ reaches it by a final occurrence of $t$; before that it moves between $a$ and $x$ using $s$ in either direction. Adjacent traversals of $s$ in opposite directions cancel by the homotopies on manuscript page 63. After removing identity edges, only one of the entries in the displayed table remains. This proves directly that the connected components of the relative path categories give exactly those hom-sets.

The two objects $a$ and $x$ are isomorphic in the result. Choosing $a$ as their representative identifies a skeleton with the walking arrow $a\to b$. Thus localization has changed the hom-set from $a$ to $b$, while it has introduced no path from $b$ to $a$. This is the elementary reason that a relative localization is generally a category rather than a groupoid.
\end{dcomment}


\begin{enumerate}[label=\alph*)]
\item toutes les flèches de $\Pi(X)$ sont inversibles, i.e. $\Pi(X)$ est bien un groupoïde, et
\item le foncteur canonique
\[
X \longrightarrow \Pi(X)
\]
qui est l'identité sur les objets, et qui est défini sur les flèches par $u \longmapsto c(u)$ comme ci-dessus, est \textit{universel} (dans Cat) pour les foncteurs (i.e. les Cat-morphismes) de $X$ vers des \textit{groupoïdes.} (Plus généralement encore, il est universel pour les foncteurs $X \xrightarrow{F} Y$ tels que $F$ transforme toute flèche de $X$ en \texttt{[une]} flèche inversible.)
\end{enumerate}
% foot: Version 22 décembre 2001

\printedpage{42}
% head: A. Grothendieck : Dérivateurs, Ch. VII, Édition des Universités de Montpellier II et Paris VII
\msp{page 63}
Pour prouver que toute flèche dans $\Pi(X)$ est inversible, comme tout chemin non vide est composé de chemins `élémentaires' de type $\Delta_+$ ou $\Delta_-$, il suffit de prouver que ces chemins sont inversibles. En fait, ils se correspondent biunivoquement,
\[
\Ch_{\Delta_+}(a,b) \simeq \Ch_{\Delta_-}(b,a),
\]
en associant au chemin de $a$ vers $b$
\[
c(u) : a \xrightarrow{u} b \qquad \text{de type } \Delta_+
\]
le chemin
\[
c(u)^o : b \xleftarrow{u} a \qquad \text{de type } \Delta_-
\]
de $b$ vers $a$. Je dis que
\[
\begin{array}{rcl}
c(u)^o c(u) & \sim & 1_a\\
c(u) c(u)^o & \sim & 1_b
\end{array}
\]
\[
\left.
\begin{tikzcd}
 & a \arrow[r, "u"] & b & a \arrow[l, "u"'] \\
 {} & a \arrow[u, "\mathrm{id}"'] \arrow[r, "\mathrm{id}"] & a \arrow[u, "u"'] & a \arrow[u, "\mathrm{id}"'] \arrow[l, "\mathrm{id}"'] \\
 {} \arrow[u] & & a &
\end{tikzcd}
\right\}\
\begin{array}{l}
\text{on a un homomorphisme dans } \Ch_{\Delta_+\Delta_-}(a,a)\\
\text{de } 1_a \text{ dans } c(u)^o c(u).
\end{array}
\]
\[
\left.
\begin{tikzcd}
 & b & a \arrow[l, "u"'] \arrow[r, "u"] & b \\
 {} & a \arrow[u, "u"'] & a \arrow[u, "\mathrm{id}"'] \arrow[l, "\mathrm{id}"'] \arrow[r, "\mathrm{id}"] & a \arrow[u, "u"'] \\
 {} \arrow[u] & & a &
\end{tikzcd}
\right\}\
\begin{array}{l}
\text{itou dans } \Ch_{\Delta_-\Delta_+}(b,b),\\
\text{de } 1_b \longrightarrow c(u)c(u)^0 % [sic] c(u)^0 -> c(u)^o, digit 0 printed instead of the letter o used elsewhere
\end{array}
\]
\texttt{[plutôt} $1_b \longleftarrow c(u)c(u)^0$\texttt{, viz.} % [sic] c(u)^0 -> c(u)^o, digit 0 printed instead of the letter o used elsewhere
\[
\begin{tikzcd}
b \arrow[d, "\mathrm{id}"] & a \arrow[l, "u"'] \arrow[r, "u"] \arrow[d, "u"] & b \arrow[d, "\mathrm{id}"] \\
b & b \arrow[l, "\mathrm{id}"'] \arrow[r, "\mathrm{id}"] & b \ \texttt{].}
\end{tikzcd}
\]
% foot: Version 22 décembre 2001

\printedpage{43}
% head: A. Grothendieck : Dérivateurs, Ch. VII, Édition des Universités de Montpellier II et Paris VII
\msp{page 64}
Notons que, plus généralement, pour tout chemin
\[
c : a \longrightarrow b, \qquad \text{de type } \tau,
\]
on définit un \textit{chemin opposé}
\[
c^o : b \longrightarrow a \qquad \text{de type } \tau^o,
\]
où
\[
\tau \longmapsto \tau^o : \mathrm{Ty} \longrightarrow \mathrm{Ty}
\]
est l'antiïnvolution du monoide libre Ty qui échange les deux générateurs $\Delta_-$, $\Delta_+$, et qui est donc donné sur l'élément général par
\[
(\Delta_{\varepsilon_1}\Delta_{\varepsilon_2}\cdots\Delta_{\varepsilon_n})^o = \Delta_{-\varepsilon_n}\Delta_{-\varepsilon_{n-1}}\cdots\Delta_{-\varepsilon_1}.
\]
Par récurrence sur la longueur de $c$ (\footnote[18]{inutile, résulte directement de la formule $(*)$ \texttt{[de la]} page suivante, en décomposant $c$ en composés de chemins `élémentaires' de type $\Delta_+$ ou $\Delta_-$.}), en utilisant le résultat précédent, que $c(u)$, $c(u)^o$ sont inverses l'un de l'autre dans $\Pi(X)$, on trouve
\[
\left.
\begin{array}{rcl}
c^o c & = & \mathrm{id}_a\\
c c^o & = & \mathrm{id}_b
\end{array}
\right\}\text{dans } \Pi(X).
\]

\dgrblock{dgrb318}{Two cancellation homotopies with fixed endpoints}
\begin{dcomment}{exposition}{p63}{Two cancellation homotopies with fixed endpoints}
Grothendieck wants reversal of a single edge to give an inverse after taking components. The proof requires two different directions of natural transformation, because the two loops have different shapes. Keeping those directions visible also explains the editors' correction on this page.

For $u:a\to b$, the loop at $a$ is the diagram
$L_a=(a\xrightarrow{u}b\xleftarrow{u}a)$. Let $K_a$ be the constant $+-$ diagram at $a$. The components
\[
 (\mathrm{id}_a,u,\mathrm{id}_a):K_a\longrightarrow L_a
\]
form a natural transformation. Both endpoint components are identities, so this is a morphism in the fixed-endpoint fibre
$\underline{\mathrm{Ch}}_{+-}(a,a)$. The collapse of the type $+-$ to the empty type gives a cartesian arrow from $K_a$ to the empty path at $a$. We have therefore joined $L_a$ to the empty path by a zigzag in the total fixed-endpoint category.

For the other loop, $L_b=(b\xleftarrow{u}a\xrightarrow{u}b)$, the appropriate constant diagram is $K_b$ at $b$. This time the transformation goes in the other direction:
\[
 (\mathrm{id}_b,u,\mathrm{id}_b):L_b\longrightarrow K_b.
\]
It is a morphism in $\underline{\mathrm{Ch}}_{-+}(b,b)$, and $K_b$ again maps to the empty path through the type collapse. The uncorrected diagram in the edition has a constant diagram at $a$ below $L_b$. Its endpoint components are $u$, so it is not a morphism in this fixed-endpoint fibre. The bracketed correction supplies exactly the downward transformation just described.

Two conventions are doing work here. First, a constant diagram of nonempty type is an object of its fibre, whereas the empty path belongs to the fibre over the unit type. When the source writes $1_a$ or $1_b$ in a fibre, it is temporarily using the constant diagram as a representative of the unit component. Second, connectedness forgets the direction of the arrows used to connect objects. Nothing requires both cancellation homotopies to point the same way.

The relative argument is identical when $u\in\Sigma$: every nonidentity backward edge in these diagrams is then an allowed denominator, and all remaining backward edges are identities. This proves that $\kappa(u)$ is invertible in $\Pi_\Sigma(X)$. There is no condition on a vertical component of a path transformation beyond naturality and the identity endpoint condition.
\end{dcomment}


\textbf{Exemple.} Si
\[
\begin{array}{rcl}
\tau & = & \Delta_+\Delta_-\Delta_+^2\Delta_-\\
\tau^o & = & \Delta_+\Delta_-^2\Delta_+\Delta_-,
\end{array}
\]
et si on a un chemin
\[
c : a_0 \xrightarrow{u_1} a_1 \xleftarrow{u_2} a_2 \xrightarrow{u_3} a_3 \xrightarrow{u_4} a_4 \xleftarrow{u_5} a_5
\]
de type $\tau$, on en déduit le chemin
\[
c^o : a_5 \xrightarrow{u_5} a_4 \xleftarrow{u_4} a_3 \xleftarrow{u_3} a_2 \xrightarrow{u_2} a_1 \xleftarrow{u_1} a_0
\]
de type $\tau^o$. L'opération $c \longmapsto c^o$ satisfait en
\msp{page 65}
effet \texttt{[pour]}
\[
\begin{array}{ccccc}
a & \xrightarrow{\gamma} & b & \xrightarrow{\gamma'} & c\\
c & \xrightarrow{\gamma'^o} & b & \xrightarrow{\gamma^o} & a
\end{array}
\]
la relation
\[
(\gamma'\gamma)^o = \gamma^o\gamma'^o, \tag*{$(*)$}
\]
% foot: Version 22 décembre 2001

\dgrblock{dgrb319}{A circle in a category with two arrows}
\begin{dcomment}{examples}{p64}{A circle in a category with two arrows}
Grothendieck's reversal operation removes backtracking without destroying genuine loops. A small example distinguishes the two phenomena. Let $X$ have two objects $a,b$ and two distinct nonidentity arrows $f,g:a\to b$, and invert every arrow. The zigzag
\[
 a\xrightarrow{f}b\xleftarrow{g}a
\]
represents the loop $h=g^{-1}f$ at $a$.

All morphisms in the resulting groupoid have unique expressions of the following forms, with $n\in\mathbf Z$:
\[
\begin{array}{c|cc}
 &a&b\\ \hline
 a&h^n&gh^n\\
 b&h^ng^{-1}&gh^ng^{-1}.
\end{array}
\]
For existence, use $g$ as a chosen path from $a$ to $b$ and replace $f$ by $gh$. After cancelling adjacent $g^{-1}g$ and $gg^{-1}$, a word has exactly the indicated form. For uniqueness, consider the functor from $X$ to the one-object group $\mathbf Z$ which sends $f$ to $1$ and $g$ to $0$. Since its target is a groupoid, it extends to the groupoid completion, and it sends $h^n$ to $n$. Thus different integers cannot give equal arrows. Composition in the table adds the integer exponents, with the displayed $g$ and $g^{-1}$ providing the changes of endpoint.

It follows that $\operatorname{Aut}_{\Pi(X)}(a)\cong\mathbf Z$ and
$\pi_0\underline{\mathrm{Ch}}(a,a)\cong\mathbf Z$. The nerve supplies a geometric check: the only nondegenerate positive-dimensional simplices of $NX$ are the two edges $f,g$. There are no composable pairs of nonidentity arrows. Its realization is therefore the graph consisting of two vertices joined by two edges, which is a circle.

If instead $\Sigma$ consists only of identities, every backward edge is an identity and the effect of a path is already defined in $X$. The relative localization is then $X$ itself. In particular the two arrows $f,g$ remain distinct and there is no nonidentity loop at $a$. The difference lies entirely in which arrows may be traversed backwards, not in a choice of representative for the same unrestricted path space.
\end{dcomment}


% ---- from Der-ChVII-456_p046-050.tex
\printedpage{44}
% head: A. Grothendieck : Dérivateurs, Ch. VII, Édition des Universités de Montpellier II et Paris VII
% foot: Version 22 décembre 2001

ce qui est en accord avec l'intuition que l'opération $\gamma \longmapsto \gamma^o$ joue le rôle d'un passage à l'inverse (`à homotopie près'), et on a la formule
\[
(\gamma'\gamma)^{-1} = \gamma^{-1}\gamma'^{-1}
\]
pour deux flèches inversibles $\gamma$, $\gamma'$ composables (dans l'ordre $\gamma'$, $\gamma$).

Pour la propriété universelle de
\[
X \longrightarrow \Pi(X)
\]
(qui permettra d'écrire
\[
\Pi(X) \simeq X\Sigma^{-1}, \qquad \text{avec } \Sigma = \Fl(X) \ \ ),
\]
autant l'établir tout de suite dans le cadre général où on se donne une partie
\[
\Sigma \subseteq \Fl(X),
\]
et on définit, pour $\tau \in \Ty$,
\[
\Ch_{\Sigma,\tau}(X) \subseteq \Ch_{\tau}(X)
\]
comme la sous-catégorie pleine de $\Ch_{\tau}(X)$ formée des chemins dans $X$ dont \textit{toutes les flèches rétrogrades}

\msp{page 66}

sont dans $X$ \texttt{[plutôt dans $\Sigma$]}. On notera que pour qu'un chemin $c$ et son opposé $c^o$ soient tous les deux des $\Sigma$-\textit{chemins,} il faut et il suffit que \textit{toutes} les flèches du chemin $c$ soient dans $\Sigma$ (on dira alors que $c$ est un $\Sigma$-\textit{chemin stricte}), puisque les flèches rétrogrades de $c^o$ sont les flèches directes de $c$.

Pour que les sous-catégories pleines $\Ch_{\tau,\Sigma}(X)$ soient stables par les foncteurs images inverses
\[
f^* : \Ch_{\tau'}(X) \longrightarrow \Ch_{\tau}(X)
\]
associés aux flèches
\[
f : \tau \longrightarrow \tau'
\]
dans $\uTy$, ou simplement pour qu'y ait stabilité pour les deux foncteurs $f^*$ associés à un $f : \Delta_- \longrightarrow 1$ et $\Delta_- \longrightarrow \Delta_-^2$, il faut et il suffit que $\Sigma$ soit stable par composition et contienne toutes les flèches identiques % [sic] qu'y ait -> qu'il y ait, missing word
\[
\left\{
\begin{aligned}
& x \in \mathrm{Ob}\,X \ \Longrightarrow \ \id_x \in \Sigma \\
& u, v \in \Sigma, \ \ u,\, v \text{ composables} \Longrightarrow uv \in \Sigma.
\end{aligned}
\right.
\]
En effet, les flèches dans un chemin $f^*(c') \in \Ch_{\tau}(X)$ sont soit des flèches identiques (qui correspondent aux flèches de transition de $I_\tau$ dégénérées pour $f$), soit des flèches parmi celles de $\tau$, ou des composés d'une suite de flèches consécutives composables de $\tau'$ (qui correspondent aux flèches de transition dans $I_{\tau'}$ incluses dans la partie d'une flèche de transition de $I_\tau$), avec conservation des signatures.

\printedpage{45}
% head: A. Grothendieck : Dérivateurs, Ch. VII, Édition des Universités de Montpellier II et Paris VII
% foot: Version 22 décembre 2001

\msp{page 67}

\dgrblock{dgrb320}{Why denominators must form a subcategory}
\begin{dcomment}{exposition}{p66}{Why denominators must form a subcategory}
Grothendieck imposes exactly the closure needed to let every interval map act on relative paths. A backward source edge may cover several consecutive backward target edges, or it may collapse to a vertex. Pullback must therefore allow composition of denominators and insertion of identity denominators. These are the two tests singled out in the text.

The map $\Delta_-\to1$ pulls the empty path at $x$ back to the one-edge path $x\xleftarrow{\mathrm{id}_x}x$. Stability forces $\mathrm{id}_x\in\Sigma$ for every $x$. Next, for composable denominators $v:z\to y$ and $u:y\to x$, the backward path
\[
 x\xleftarrow{u}y\xleftarrow{v}z
\]
pulls back along $\Delta_-\to\Delta_-^2$ to
$x\xleftarrow{uv}z$. Thus stability forces $uv\in\Sigma$.
Conversely, assume these two closure properties. The image of a source edge under an endpoint-preserving increasing type map is either a vertex or a consecutive monochromatic block. A backward source edge consequently becomes either an identity or a composite of backward target arrows. Both belong to $\Sigma$. This proves stability under every type map, including maps with both collapses and subdivisions.

The condition says that $\Sigma$ is the arrow set of a wide subcategory of $X$. It does not require an Ore condition, saturation, or a two-out-of-three property. If we begin with an arbitrary collection $S$ of arrows to invert, we may first take its closure under identities and finite composition. A functor inverts this closure exactly when it inverts $S$, so the intended universal localization is unchanged. Saturation is a later operation: it adds every arrow whose image in the localization happens to be invertible.

The adjective ``full'' in $\underline{\mathrm{Ch}}_{\Sigma,\tau}(X)$ deserves equal attention. The restriction is on backward \emph{horizontal} arrows of the path. Given two such paths, every natural transformation between them remains available, with identity components at prescribed endpoints. Its internal vertical arrows need not lie in $\Sigma$. This fullness is consistent with the ordinary localization theorem, as the calculation at manuscript page 68 will show. It must be remembered when making a higher-homotopical comparison with constructions which restrict vertical arrows as well.

Finally, reversal of an individual path remains allowed precisely when its forward arrows also belong to $\Sigma$. Closure of the entire path system under reversal therefore forces $\Sigma=\operatorname{Fl}(X)$. This statement concerns literal permitted paths; it does not say that a localization can be a groupoid only in that case. Arrows outside $\Sigma$ may already be invertible in $X$, or become invertible after localization.
\end{dcomment}


On suppose ces conditions satisfaites, de sorte que les $\Ch_{\tau,\Sigma}(X)$ forment les fibres d'une sous-catégorie fibrée pleine de $\Ch(X)$ sur $\uTy$, qu'on dénotera par
\[
\Ch_{\Sigma}(X) \hookrightarrow \Ch(X).
\]
Il est évident que cette sous-catégorie est stable par l'opération de composition des chemins (mais non pas, en général, on l'a dit, par passage de $c$ à $c^o$ (il faudrait pour cela que $\Sigma$ contienne \textit{toutes} les flèches de $X$, puisque pour toute $u : a \longrightarrow b$, qui est donc dans $\Ch_{\Delta_+,\Sigma}(X)$, $u^o$ donné par $b \xleftarrow{\ u\ } a$ dans $\Ch_{\Delta_-}(X)$ devrait être un $\Sigma$-chemin, d'où $u \in \Sigma$)).

On va définir une catégorie canonique
\[
\Pi_\Sigma(X),
\]
ayant même objets que $X$, et un foncteur canonique
\[
X \longrightarrow \Pi_\Sigma(X)
\]
qui est l'identité sur les objets, et qui transforme les $u \in \Sigma$ en flèches inversibles. Pour cela on calque la description du groupoïde fondamental (mais $\Pi_\Sigma(X)$ ne sera pas en général un groupoïde), en posant
\[
\Hom_\Sigma(x,y) = \pi_0(\Ch_\Sigma(x,y)),
\]

\msp{page 68}

et en définissant la composition des `$\Sigma$-flèches' en termes de la composition des chemins. L'argument qui prouvait, dans le cas de $\Pi(X)$, que \textit{toute} flèche de $X$, vue comme chemin $c$ de type $\Delta_+$, était inversible dans $\Pi(X)$, ayant comme inverse $c^o$ de type $\Delta_-$, est encore valable, pourvu que $c^o$ soit encore bien un $\Sigma$-chemin, ce qui exige que $u \in \Sigma$.

On va l'interpréter en termes d'un foncteur canonique
\[
\kappa : X \longrightarrow \Pi_\Sigma(X) \qquad \text{(\footnotemark[19])}
\tag*{$(*)$}
\]
\footnotetext[19]{Sur les flèches, $\kappa$ est défini, pour $u : a \longrightarrow b$, par $\kappa(u) =$ classe d'homotopie de $\ch(u)$ dans $\Ch_\Sigma(a,b)$, où $\ch(u)$ est le chemin de type $\Delta_+$ de $a$ vers $b$ défini par $u$.}
qui se décrit encore comme dans le cas du groupoïde fondamental, en associant à toute flèche $u$ de $X$ la classe d'homotopie du chemin $c(u)$ qu'elle définit. On voit donc que ce foncteur canonique transforme les $u \in \Sigma$ en flèches inversibles.

Quand les flèches de $\Sigma$ sont déjà inversibles dans $X$, je dis que le foncteur $(*)$ est un isomorphisme, i.e. que pour tout $a, b \in \mathrm{Ob}\,X$,
\[
\alpha : \Hom_X(a,b) \longrightarrow \pi_0\Ch_\Sigma(a,b), \qquad u \longmapsto \underbrace{\ch_\Sigma(u)}_{= \ch(u)}
\]
est bijectif. Pour le voir, on va définir
\[
\beta : \pi_0(\Ch_\Sigma(a,b)) \longrightarrow \Hom_X(a,b),
\]
qui sera inverse de $\beta$. On définit % [sic] inverse de $\beta$ -> inverse de $\alpha$, $\beta$ is meant to be the inverse of $\alpha$

\printedpage{46}
% head: A. Grothendieck : Dérivateurs, Ch. VII, Édition des Universités de Montpellier II et Paris VII
% foot: Version 22 décembre 2001

\msp{page 69}

\dgrblock{dgrb321}{The effect of a zigzag descends to components}
\begin{dcomment}{exposition}{p68}{The effect of a zigzag descends to components}
Grothendieck now uses evaluation in a category where the denominators are already invertible. Before checking that evaluation is inverse to $\alpha$, we should verify that it is defined on a connected component of the path category. This is a short verification, but it explains why the internal vertical arrows of a path transformation need not be invertible.

The useful statement is slightly more general. Let $F:X\to Y$ send every arrow of $\Sigma$ to an isomorphism. For a $\Sigma$-path
$c=(x_0,\ldots,x_n)$, let
$e_i:F(x_{i-1})\to F(x_i)$ be $F(u_i)$ if its $i$th edge is forward and $F(u_i)^{-1}$ if that edge is backward. Define
\[
 E_F(c)=e_n\cdots e_1,
 \qquad E_F(\text{empty path at }a)=\mathrm{id}_{F(a)}.
\]
This formula is defined solely by the inversion hypothesis.

First compare paths $c,c'$ of the same type by a natural transformation with components $\eta_i:F(x_i)\to F(x'_i)$. For a forward edge the commuting square reads
$\eta_i e_i=e'_i\eta_{i-1}$. For a backward edge the original square is
\[
 \eta_{i-1}F(u_i)=F(u'_i)\eta_i.
\]
Multiplication by the inverses of $F(u_i)$ and $F(u'_i)$ gives exactly the same formula
\[
 \eta_i e_i=e'_i\eta_{i-1}.
\]
Successively composing these equations yields
\[
 \eta_n E_F(c)=E_F(c')\eta_0.
\]
At fixed endpoints, $\eta_0$ and $\eta_n$ are identities. Hence the two effects agree. No inverse of any $\eta_i$ has been taken.

We must also compare different types. If $q:\tau\to\tau'$ and $c'$ is a path of type $\tau'$, then each edge of $q^*c'$ is a composite over the corresponding consecutive block of edges of $c'$, or an identity for an empty block. A forward block has the same effect as its composite. A backward block does too, because inversion reverses the order of composition. The blocks partition the target edges in their traversal order, so
\[
 E_F(q^*c')=E_F(c').
\]
A general morphism in the total path category factors into a morphism within a fibre followed by a cartesian morphism of this form. We have therefore proved invariance under every morphism and hence under every connecting zigzag. Evaluation defines a map on $\pi_0\underline{\mathrm{Ch}}_\Sigma(a,b)$.

Taking $F=\mathrm{id}_X$ gives the missing descent check for $\beta$ in the text. The identity $\beta\alpha=\mathrm{id}$ is immediate. For the other identity, a backward traversal of an invertible $u\in\Sigma$ represents the inverse of $\kappa(u)$; this is also $\kappa(u^{-1})$, since $\kappa$ is a functor. Replacing each backward edge in this way gives the all-forward path $c^\natural$, and repeated composition identifies its class with the elementary path of its composite. These equalities take place in $\Pi_\Sigma(X)$. The omitted subscripts in equations (1) and (2) on manuscript page 69 should not lead us to infer equality in a relative localization merely from equality in its groupoid completion.
\end{dcomment}


$\beta(c)$, pour une $\Sigma$-flèche de type $\tau$, le longueur $n$, en inversant toutes les flèches rétrogrades (ce qui est possibles, puisqu'elles sont dans $\Sigma$), de façon à obtenir un chemin $c^\natural$ de type $\Delta_+^n$, et en prenant le composé (dans $X$) des flèches de ce chemin $c^\natural$. Il est évident alors que $\beta\alpha = \id$, il faut prouver que l'on a aussi $\alpha\beta = \id$, i.e. que dans $\Pi_\Sigma(X)$, les flèches $u = \int c^\natural$ et $c$ sont égales, i.e. $u$ et $c$ sont homotopes dans $\Ch_\Sigma(a,b)$. Considérons le chemin $c^\natural$, de type $\Delta_+^n$ : % [sic] le longueur -> la longueur, gender; possibles -> possible, number
\[
a = a_0 \xrightarrow{u_1} a_1 \xrightarrow{u_2} a_2 \longrightarrow \ \cdots \ \longrightarrow a_{n-1} \xrightarrow{u_n} a_n = b,
\]
et soit $u_i$ ($i \in [1,n]$) une flèche de ce chemin, correspondant à une flèche de transition $\underline{u}_i$ \textit{rétrograde} de $\tau$, donc par construction
\[
u_i = v_i^{-1} \qquad \text{dans } X\text{, avec } v_i \in \Sigma
\]
la flèche d'indice $i$ dans le chemin initial $c$ de type $\tau$. Si $\underline{u}_i$ par contre est directe, on aura
\[
u_i = v_i.
\]
(NB On désigne par $v_i$ la $i$-flèche dans le chemin $c$.) Or on a
\[
c = c_n * c_{n-1} * \cdots * c_1 \qquad \text{dans } \Pi_\Sigma(X)
\tag*{$(*)$}
\]
et le facteur $c_i$ est soit la classe d'homotopie de $\ch(v_i) = \ch(u_i)$ (cas où la flèche de transition $\underline{u}_i$ de $\tau$ est directe), soit celle du chemin \textit{opposé} au chemin $\ch(v_i) : a_i \xrightarrow{v_i} a_{i-1}$ de type $\Delta_+$, donc le chemin $\ch(v_i)^o : a_{i-1} \xleftarrow{v_i} a_i$ de type $\Delta_-$. Je dis que
\[
\boxed{\ \ch(v_i)^o = \ch(u_i) \qquad \text{\textit{dans} } \Pi(X),\ }
\tag*{(1)}
\]
ce qui prouve que dans $\Pi_\Sigma(X)$, $c$ est égal au composé des $\ch(u_i)$ ($i \in [1,n]$), donc à $c^\natural$. Il faut prouver de plus que
\[
c^\natural = \ch(u) \qquad (u = u_nu_{n-1}\cdots u_1) \qquad \text{dans } \Pi(X).
\tag*{(2)}
\]

\msp{page 70}

\dgrblock{dgrb322}{Relations before higher homotopies}
\begin{dcomment}{context}{p69}{Relations before higher homotopies}
Grothendieck's path construction is testing the same universal problem as the classical category of fractions, while retaining the witnesses for its elementary relations. In \citet[Chapter I, \S1.1 and Lemma 1.2]{GabrielZisman1967}, localization is constructed from paths with formal reverse arrows by imposing composition, identity, and inverse relations. Lemma 1.2 also gives the universal property on functor categories and natural transformations. Here composition and identity relations come from type changes; the inverse relations come from the two fixed-endpoint ladders. The evaluation calculation above shows that arbitrary path ladders respect these relations. The universal property on the following pages will prove that taking components introduces no additional identifications.

This is also the appropriate scale on which to read the connection with the derivator programme. In \citet[\S69, pp.~1--4 of the separately edited text]{Grothendieck1983}, the diagramwise construction
\[
 I\longmapsto
 \operatorname{Fun}(I,M)[W_I^{-1}]
\]
retains a whole family of localized diagram categories, where $W_I$ consists of pointwise weak equivalences. The present calculation makes morphisms in such ordinary localizations accessible by diagrams of arrows and commuting squares. It does not by itself supply the homotopy Kan extensions or the higher comparisons required for a derivator.

The chronological comparison with Dwyer and Kan is particularly instructive, since their simplicial localization papers appeared in 1980, before these manuscript pages. Their hammock construction restricts every vertical arrow to the chosen subcategory of weak equivalences, as well as every backward horizontal arrow; its component category is the ordinary localization, and it is weakly equivalent to their standard simplicial localization \citep[Definition 2.1 and Propositions 2.2, 3.1]{DwyerKan1980Calculating}. Grothendieck's full relative path fibres allow arbitrary internal vertical arrows. Agreement of component categories is therefore a common starting point, not an automatic comparison theorem for the entire nerves. The stronger assertions later in the chapter must be justified with the actual path category and its actual indexing maps.
\end{dcomment}


Donc il y a deux vérifications à faire (où je suppose à nouveau $\Sigma$ quelconque) :

\begin{enumerate}[label={}]
\item[$1^o$)] Soit $u \in \Sigma$, $u : a \longrightarrow b$, supposons $\Sigma$ inversible dans $X$. Alors
\[
\ch(u)^o = \ch(u^{-1}) \qquad \text{dans } \Pi_\Sigma(X),
\]
i.e. les deux chemins de $b$ vers
\[
\begin{array}{cccl}
b & \xleftarrow{\ u\ } & a & \text{de type } \Delta_- \\
b & \xrightarrow{u^{-1}} & a & \text{de type } \Delta_+
\end{array}
\]
\printedpage{47}
% head: A. Grothendieck : Dérivateurs, Ch. VII, Édition des Universités de Montpellier II et Paris VII
% foot: Version 22 décembre 2001
sont homotopes dans $\Ch_\Sigma(b,a)$. Or on sait déjà, comme $\kappa : X \longrightarrow \Pi_\Sigma(X)$ est un foncteur, que
\[
\hot\,\ch(u^{-1}) = (\hot\,\ch(u))^{-1} \qquad \text{dans } \Pi_\Sigma(X)
\]
où $\hot(\bullet)$ désigne la classe d'homotopie d'une chaine (ici classe dans $\Ch_\Sigma(b,a)$ et dans $\Ch_\Sigma(a,b)$ respectivement). Donc l'énoncé à prouver c'est que $\hot(\ch(u)^o)$ est égal à l'\textit{inverse} de $\hot(\ch(u))$ dans $\Pi_\Sigma(X)$. C'est donc ce qu'on a prouvé pour \textit{tout} $u \in \Fl(X)$, dans le cas de $\Pi(X)$. L'argument doit marcher, du moment qu'il a un sens, i.e. que $\ch(u)^o$ est bel et bien un $\Sigma$-chemin, i.e. $u \in \Sigma$.

\item[$2^o$)] La formule (2) signifie que pour tout chemin $c$ de type $\Delta_+^n$ (donc tautologiquement un \msp{page 71} $\Sigma$-chemin), allant de $a$ vers $b$, sa classe dans $\Hom_\Sigma(a,b)$ n'est autre que celle de $\ch(u)$, où $u$ est le composé $u_nu_{n-1}\cdots u_1$. Mais on a
\[
c = c_1 * c_2 * \cdots * c_n,
\]
d'où dans $\Pi_\Sigma(X)$
\[
\hot(c) = \hot(c_n)\hot(c_{n-1})\cdots\hot(c_1),
\]
or $\hot(c_i) = \hot\,\ch(u_i) = \kappa(u_i)$, donc
\[
\begin{aligned}
\hot(c) \ &= \ \kappa(u_n)\kappa(u_{n-1})\cdots\kappa(u_1) \\
&= \ \kappa(\underbrace{u_nu_{n-1}\cdots u_1}_{u}),
\end{aligned}
\]
puisque $\kappa$ est un foncteur, donc
\[
\hot(c) = \hot\,\ch(u),
\]

\hfill c.q.f.d.
\end{enumerate}

Notons maintenant que si $(X', \Sigma')$ est un deuxième couple d'une catégorie $X'$ et d'un ensemble \textit{`admissible'} de flèches de $X'$, et si
\[
f : X \longrightarrow X'
\]
est un foncteur tel que
\[
f(\Sigma) \subseteq \Sigma',
\]
alors il existe un \textit{unique} foncteur
\[
\bar{f} = \Pi_\Sigma(f) : \Pi_\Sigma(X) \longrightarrow \Pi_{\Sigma'}(X')
\]
qui `prolonge' $f$, i.e. tel que le carré
\printedpage{48}
% head: A. Grothendieck : Dérivateurs, Ch. VII, Édition des Universités de Montpellier II et Paris VII
% foot: Version 22 décembre 2001
\msp{page 72}
\[
\begin{tikzcd}
X \arrow[r, "\kappa_X"] \arrow[d, "{\scriptstyle f}"] & \Pi_\Sigma(X) \arrow[d, "{\scriptstyle \bar{f} = \Pi_\Sigma(f)}"] \\
X' \arrow[r, "\kappa_{X'}"] & \Pi_{\Sigma'}(X')
\end{tikzcd}
\]
soit commutatif. L'unicité résulte du fait que sur les objets $\bar{f}$ doit être donné par $f$, et que sur les flèche la commutativité du carré implique que $\bar{f}$ est connu sur les flèches de la forme $\kappa_X(u)$, ainsi donc que sur les inverses de telles flèches, quand ces inverses sont inversibles, donc en particulier sur les $\kappa_X(u)^{-1}$ pour $u \in \Sigma$. Or toute flèche de $\Pi_\Sigma(X)$, provenant d'un chemin, est composé d'une suite de flèches de cette forme. Pour l'existence de $\bar{f}$, on note que $f$ induit des foncteurs % [sic] les flèche -> les flèches, plural
\[
\Ch_\tau(X; a, b) \longrightarrow \Ch_\tau(X'; f(a), f(b))
\]
qui commutent aux opérations $\varphi^*$ sur les chemins associées aux foncteurs
\[
\varphi : \tau \longrightarrow \tau',
\]
d'où un homomorphisme de catégories fibrées sur $\uTy$
\[
f_{a,b} : \Ch(X; a, b) \longrightarrow \Ch(X'; f(a), f(b)).
\]

\msp{page 73}

D'où une application sur les $\pi_0$, i.e.
\[
\bar{f}_{a,b} : \Hom_\Sigma(a,b) \longrightarrow \Hom_\Sigma(f(a), f(b)),
\]
donnée par
\[
\bar{f}_{a,b}(\hot(c)) = \hot(f_{a,b}(c)).
\]
On a évidemment, pour deux chemins $\gamma : a \longrightarrow b$ et $\gamma' : b \longrightarrow c$,
\[
f_{a,c}(\gamma\gamma') = f_{b,c}(\gamma')f_{a,b}(\gamma),
\]
ce qui prouve que $\bar{f}$ est compatible avec la composition. Il est trivial qu'il est compatible avec les identités (car $f_{a,a}$ transforme le chemin vide en $a$ en le chemin vide en $f(a)$), donc on a un foncteur $\bar{f}$, et il rend bel et bien le carré commutatif.

On peut maintenant prouver la propriété universelle de $\Pi_\Sigma(X)$ et de
\[
X \xrightarrow{\ \kappa\ } \Pi_\Sigma(X),
\]
savoir que tout foncteur $f : X \longrightarrow X'$ tel que $u \in \Sigma \Longrightarrow f(u)$ inversible, se factorise \textit{de façon unique} par $\Pi_\Sigma(X)$. L'unicité se voit comme tantôt (\footnote[20]{C'est inutile, cf. ci-dessous.}). Pour l'existence, on note que $f$ applique $\Sigma$ dans l'ensemble $\Sigma'$

% ---- from Der-ChVII-456_p051-055.tex
\setcounter{footnote}{20}

\printedpage{49}
% head: A. Grothendieck : Dérivateurs, Ch. VII, Édition des Universités de Montpellier II et Paris VII
\msp{page 74}
des flèches inversibles de $X'$ (lequel ensemble satisfait bien aux deux conditions d'admissibilité). Donc par ce qui précède on a un diagramme commutatif
\[
\begin{tikzcd}
X \ar[r,"\kappa"] \ar[d,"f"'] & \Pi_{\Sigma}(X) \ar[d,"\bar f = \Pi_{\Sigma}(f)"] \ar[dl,dotted,"\tilde f"'] \\
X' \ar[r,"\kappa'","\sim"'] & \Pi_{\Sigma'}(X'),
\end{tikzcd}
\]
où $\bar f$ est uniquement déterminé. Comme $\kappa'$ est un isomorphisme (comme on a vu tantôt, puisque $\Sigma'$ \texttt{[est]} formé des flèches inversibles), l'existence et l'unicité de $\bar f$ équivaut à la conclusion annoncée,

\hfill c.q.f.d.

\textbf{Encore un sorite} sur les classes d'homotopie d'éléments d'une catégorie $X$. De façon générale, on voit que $X$ et $\Pi_{\Sigma}(X) = X\Sigma^{-1}$ ont mêmes composantes connexes, car si
\msp{page 75}
$a$ et $b$ dans $X$ sont reliés par une flèche de $X\Sigma^{-1}$, i.e. par un $\Sigma$-chemin, alors ils sont dans une même composante connexe de $X$. Mais l'inverse est également vrai, comme on voit p.~ex. en prenant $\Sigma = \mathrm{Fl}(X)$, et en disant que dans le \textit{groupoïde} $\Pi(X)$, où deux sommets sont dans une même composante \texttt{[connexe]} si et seulement si $\exists\ x \longrightarrow y$ (relation qui est non seulement transitive et reflexive, mais aussi symétrique dans un groupoïde). Un chemin
\[
a \xrightarrow{\quad c \quad} b
\]
sera appelé aussi une \textit{homotopie} de $a$ vers $b$. On dit que $a$ et $b$ sont \textit{homotopes} s'il existe une homotopie $c$ de $a$ vers $b$, auquel cas il existe aussi une homotopie `inverse' $c^{o}$ de $b$ vers $a$. Cela signifie aussi que $a$ et $b$ sont dans la même composante connexe de $X$, i.e. définissent le même élément de $\pi_0(X)$.

\dgrblock{dgrb323}{The universal property includes natural transformations}
\begin{dcomment}{exposition}{p74}{The universal property includes natural transformations}
Grothendieck has now brought the path construction back to the familiar category of fractions. The work with arbitrary lengths and their changes of parameter pays for a universal property requiring only that the denominators contain identities and be closed under composition. No condition allowing fractions to be shortened to spans has entered the argument. This is the general localization of \citet[Chapter~I, \S1.1 and Lemma~1.2]{GabrielZisman1967}; a calculus of left or right fractions is additional structure that can simplify its representatives.

Write $Q:X\longrightarrow L=\Pi_\Sigma(X)$. A functor $F:X\longrightarrow Y$ carrying $\Sigma$ into isomorphisms must send a zigzag
\[
a=x_0\xrightarrow{u_1}x_1\xleftarrow{s_2}x_2
 \xrightarrow{u_3}\cdots\xleftarrow{s_n}x_n=b
\]
to the corresponding composite of the $F(u_i)$ and $F(s_j)^{-1}$, in their order of traversal. The preceding discussion of the evaluation map proves that this prescription is constant on components of the path category. It therefore defines $\overline F:L\longrightarrow Y$. Every arrow of $L$ is represented by a zigzag, so these prescribed values also prove uniqueness. The factorization through $\Pi_{\mathrm{Iso}(Y)}(Y)\simeq Y$ used in the text packages precisely this argument.

There is a useful categorical refinement, already part of the classical universal property. Precomposition induces an isomorphism
\[
 Q^*:\operatorname{Fun}(L,Y)\xrightarrow{\ \sim\ }
       \operatorname{Fun}_{\Sigma}(X,Y),
\]
where the right side is the full subcategory of $\operatorname{Fun}(X,Y)$ on the functors inverting $\Sigma$. The statement concerns all natural transformations, not merely natural isomorphisms. To see the extra assertion directly, let $\eta:F\Rightarrow G$ be a natural transformation between such functors. Since $Q$ is the identity on objects, the only possible extension has components $\overline\eta_x=\eta_x$. Naturality for $s:x\to y$ gives
\[
 \eta_yF(s)=G(s)\eta_x,
 \qquad
 \eta_xF(s)^{-1}=G(s)^{-1}\eta_y.
\]
The second equality supplies naturality for the formal inverse of $s$. Naturality then passes to every composite representing an arrow of $L$. These components define the required transformation, and uniqueness follows again from the identity on objects. This also explains why relative functors induce functors on localizations without further choices: both their action on arrows and their action on transformations are determined before any presentation of a fraction is chosen.

For the present chapter, this refinement matters because a path whose vertices are themselves functors is a diagram of natural transformations. The later passage to diagram categories will repeatedly use the same distinction between determining objects and controlling the morphisms between them. The localization established here already supplies both at the ordinary categorical level.
\end{dcomment}


L'homotopie sera dite \textit{simple} si $c$ est de type $\Delta_{+}$ ou $\Delta_{-}$, i.e. si elle est donnée par une flèche $u : a \longrightarrow b$ ou $v : b \longrightarrow a$. Si on cherche une homotopie entre $a$ et $b$, on essaiera de la trouver la plus `courte' possible. Si on a une homotopie $c$ de longueur $l$ de $a$ vers $b$, alors en remplaçant $\tau$ par $\tau_{\red}$, donc $I_{\tau}$ par $I_{\tau\,\red}$, le sous-intervalle réduit
\msp{page 76}
correspondant, on trouve une homotopie de $a$ à $b$ de type $\tau_{\red}$, dont la longueur est égale à la \textit{`longueur utile'} du type $\tau$ de $c$ (ce qui justifie la terminologie). Notons que si
\[
\tau = \Delta_{\varepsilon}^{l_1}\Delta_{-\varepsilon}^{l_2}\Delta_{\varepsilon}^{l_3}\cdots\Delta_{(-1)^{r}\varepsilon}^{l_r},
\]
alors
\[
\tau_{\red} = \Delta_{\varepsilon}\Delta_{-\varepsilon}\Delta_{\varepsilon}\cdots\Delta_{(-1)^{r}\varepsilon},
\]
% foot: Version 22 décembre 2001
\printedpage{50}
% head: A. Grothendieck : Dérivateurs, Ch. VII, Édition des Universités de Montpellier II et Paris VII
\texttt{[plutôt deux fois} $(-1)^{r-1}$\texttt{]} où $r$ est la longueur utile de $\tau$ (la longueur pas utile étant $n = \sum\limits_{i\in[1,r]} l_i$).

Quand on part d'une homotopie entre \textit{foncteurs}
\[
X \longrightarrow Y,
\]
il est entendu qu'il s'agit d'homotopies dans la catégorie $\uHom(X,Y)$. (Ou, suivant le contexte, dans une sous-catégorie précisée de $\uHom(X,Y)$.)

\msp{page 77}
À la page 61, j'ai invoqué l'existence d'un \textit{foncteur} de \textit{`composition des chemins'}
\[
\Ch_{\Sigma}(X;a,b)\times\Ch_{\Sigma}(X;b,c) \ \longrightarrow\ \Ch_{\Sigma}(X;a,c),
\]
qui sur la fibre du premier membre au dessus de $(\tau,\tau') \in \operatorname{Ob}(\Ty\times\Ty)$, serait donné par
\[
\Ch_{\Sigma,\tau}(X;a,b)\times\Ch_{\Sigma,\tau'}(X;b,c) \ \longrightarrow\ \Ch_{\Sigma,\tau\tau'}(X;a,c), \tag*{$(*)$}
\]
la composition des chemins de types précisés $\tau$, $\tau'$. Qu'on ait là un \textit{foncteur,} pas seulement une application sur les ensembles d'objets, est assez évident, et peut s'exprimer par exemple à l'aide du diagramme d'amalga\texttt{[ma]}tion
\[
\begin{tikzcd}[column sep=small]
 & \varepsilon = \{0,1\} \ar[dl,hook'] \ar[dr,hook] & \\
I_{\tau} \ar[dr,hook] & & I_{\tau'} \ar[dl,hook'] \\
 & I_{\tau\tau'} &
\end{tikzcd}
\]
en le transformant par $\uHom(\cdot,X)$ (au lieu de le transformer par $\mathrm{Hom}_{\mathrm{Cat}}(\cdot,X)$ seulement), du moins dans le cas particulier où $\Sigma = \mathrm{Fl}(X)$ (qui était celui de la page 61). Mais le passage aux flèches est assez évident - pour deux composés $c'_1\circ c_1$ et $c'_2\circ c_2$, de
\[
\begin{tikzcd}[column sep=small]
a \ar[rr,"c_1"] \ar[d,equal] & {} \ar[d,"w"'] & b \ar[rr,"c'_1"] \ar[d,equal] & {} \ar[d,"w'"] & c \ar[d,equal] \\
a \ar[rr,"c_2"] & {} & b \ar[rr,"c'_2"] & {} & c
\end{tikzcd}
\]
\texttt{[on]} voit que les homomorphismes $w''$ du premier dans le second, dans $\Ch_{\Sigma}(X;a,c)$, qui sont l'identité sur le sommet d'indice $n+1$ des deux chemins (où $n = \operatorname{long}\tau$), correspondant à $1_{\tau}$ et à $0_{\tau'}$, `sont' exactement les couples
% foot: Version 22 décembre 2001
\printedpage{51}
% head: A. Grothendieck : Dérivateurs, Ch. VII, Édition des Universités de Montpellier II et Paris VII
\msp{page 78}
formés d'un homomorphisme $w : c_1 \longrightarrow c_2$ et $w' : c'_1 \longrightarrow c'_2$ dans $\Ch_{\Sigma,\tau}(a,b)$ et $\Ch_{\Sigma,\tau'}(b,c)$ respectivement, et si $w''$, $w$, $w'$ se correspondent ainsi, on posera
\[
w'' = w * w'.
\]
Il y a une transitivité évidente à vérifier pour voir que ça définit un \textit{foncteur} $(*)$. On peut exprimer la situation à l'aide du diagramme
\[
\begin{tikzcd}[column sep=small]
 & \Ty\times\Ty \ar[rrr,"\pi\ (\text{composition})"] \ar[dl,"p_1"'] \ar[dr,"p_2"] & & & \Ty \\
\Ty & & \Ty, & &
\end{tikzcd}
\]
où $p_1$, $p_2$ sont les deux projections. On considère sur $\Ty\times\Ty$ les trois catégories fibrés % [sic] fibrés -> fibrées, agreement with "catégories"
\[
p_1^*(\CH_{\Sigma,\tau}(a,b)),\ \ p_2^*(\CH_{\Sigma,\tau'}(b,c)),\ \ \pi^*(\Ch_{\Sigma,\tau\tau'}(a,c))
\]
\texttt{[plutôt} $\pi^*(\CH_{\Sigma,\tau\tau'}(a,c)$\texttt{]}, dont les fibres en $(\tau,\tau')$ sont respectivement % [sic] \pi^*(\CH_{\Sigma,\tau\tau'}(a,c) -> \pi^*(\CH_{\Sigma,\tau\tau'}(a,c)), one closing parenthesis missing
\[
\Ch_{\Sigma,\tau}(a,b),\ \ \Ch_{\Sigma,\tau'}(b,c),\ \ \Ch_{\Sigma,\tau\tau'}(a,c).
\]
(On écrit $\CH$ au lieu de $\Ch$, pour indiquer que l'on considère les catégories envisagées comme des catégories (fibrées) \textit{au dessus de} $(\Ty)$ (\footnote{Mais il s'agit pour $p_1^*$, $p_2^*$ et $\pi^*$ de produits cartésiens ordinaires dans Cat, il n'y a pas à avoir peur et à écrire des $\CH$ \ldots}).) Au dessus de $\Ty\times\Ty$, on a un homomorphisme (cartésien) de catégories fibrées
\[
p_1^*(\CH_{\tau}(a,b))\times_{\Ty\times\Ty} p_2^*(\CH_{\tau}(b,c)) \ \longrightarrow\ \pi^*(\CH_{\tau}(a,c)) \tag*{$(**)$}
\]
\texttt{[plutôt} $p_1^*(\CH_{\tau}(a,b))\times_{\Ty\times\Ty} p_2^*(\CH_{\tau'}(b,c)) \ \longrightarrow\ \pi^*(\CH_{\tau\tau'}(a,c))$\texttt{]},
\msp{page 79}
car c'est OK sur les fibres ($(*)$ page 77), et si
\[
(\varphi,\varphi') : (\tau_1,\tau'_1) \longrightarrow (\tau,\tau') \qquad (\varphi : \tau_1 \longrightarrow \tau,\ \varphi' : \tau'_1 \longrightarrow \tau')
\]
est une flèche dans $\Ty\times\Ty$, on en déduit un diagramme commutatif
\[
\begin{tikzcd}
\Ch_{\Sigma,\tau}(a,b)\times\Ch_{\Sigma,\tau'}(b,c) \ar[r,"\text{comp.}"] \ar[d,"\varphi^*_{a,b}\times(\varphi'_{b,c})^*"] & \Ch_{\Sigma,\tau}(a,c) \ar[d,"(\varphi*\varphi')^*_{a,c}"] \\
\Ch_{\Sigma,\tau_1}(a,b)\times\Ch_{\Sigma,\tau'_1}(b,c) \ar[r,"\text{comp.}"] & \Ch_{\Sigma,\tau_1\tau'_1}(a,c)
\end{tikzcd}
\]
% foot: Version 22 décembre 2001
\printedpage{52}
% head: A. Grothendieck : Dérivateurs, Ch. VII, Édition des Universités de Montpellier II et Paris VII
\texttt{[plutôt} $\Ch_{\Sigma,\tau\tau'}(a,c)$\texttt{]}, i.e. la composition des chemins d'extrémités fixés est compatible avec les images inverses de chemins par flèches entre les types de chemins.

\dgrblock{dgrb324}{Connectedness does not impose an orientation}
\begin{dcomment}{examples}{p76}{Connectedness does not impose an orientation}
Grothendieck next isolates the least information carried by his paths: whether two objects belong to one component. For any admissible $\Sigma$, the map $Q:X\to X[\Sigma^{-1}]$ induces a bijection on connected components. Indeed, every arrow of $X$ supplies an edge between its two objects in the localization, so a zigzag in $X$ remains a zigzag there. Conversely, an arrow of the localization is represented by a $\Sigma$-zigzag in $X$; forgetting which of its arrows were designated denominators gives an ordinary zigzag in $X$. Applying this observation to the finitely many arrows of a zigzag in the localization proves the converse on components.

The orientation must be kept separate from this argument. Let $X=[1]$, with its single nonidentity arrow $0\to1$, and let $\Sigma$ contain only identities. Then $X[\Sigma^{-1}]=X$ and $0,1$ belong to the same component, although there is no arrow $1\to0$ in the localization. Thus the phrase about the converse on manuscript page~75 is read as a statement about connected components. It would be false if read as the assertion that connectedness supplies a directed $\Sigma$-path in every prescribed direction. When $\Sigma=\mathrm{Fl}(X)$, reversal is available, and the groupoid formulation does permit that stronger assertion.

Reduction of a type has an equally concrete meaning. A string
\[
 a_0\longrightarrow a_1\longrightarrow a_2
 \longleftarrow a_3\longleftarrow a_4\longrightarrow a_5
\]
can be replaced by
\[
 a_0\longrightarrow a_2\longleftarrow a_4\longrightarrow a_5
\]
by composing the arrows in each maximal run. The signs of the three remaining edges are $+,-,+$. For $r$ runs beginning with $\varepsilon$, their signs are $(-1)^{j-1}\varepsilon$, $1\leq j\leq r$; this is the reason for the edition's correction of the last exponent on pages~76--77. If the original backward arrows belong to $\Sigma$, closure under composition ensures that the reduced backward arrows still do. The same construction in $\operatorname{Fun}(A,B)$ composes successive natural transformations of the same direction. Here a homotopy between functors means a finite zigzag of natural transformations; its endpoints need not be naturally isomorphic.
\end{dcomment}


Mais l'homomorphisme cartésien $(**)$ de catégories fibrées sur $\Ty\times\Ty$ n'est pas tout à fait la même chose qu'un anodin foncteur
\[
\Ch_{\Sigma}(a,b)\times\Ch_{\Sigma}(b,c) \ \longrightarrow\ \Ch_{\Sigma}(a,c). \tag*{$(*)$}
\]
La formule $(*)$ définit en tout cas une application sur les ensembles d'objets, puisque les ensembles d'objets de chacune des trois catégories envisagées se décompose en somme d'ensembles $\Ch_{\Sigma,\tau}(\bullet,\bullet)$, où $\tau$ parcourt $\mathrm{Ty} = \operatorname{Ob}\Ty$. Mais cela se prolonge-t-il de façon naturelle aux morphismes?

\msp{page 80}
Mais voilà que je me rappelle que ce fameux produit cartésien $p_1^*(\text{fourbi})\times p_2^*(\text{machin})$ n'est autre que le produit cartésien tout bête
\[
\Ch_{\Sigma}(a,b)\times\Ch_{\Sigma}(b,c),
\]
envoyé sur $\Ty\times\Ty$ par le produit cartésien des deux foncteurs de projection
\[
\begin{aligned}
p_{a,b} : \Ch_{\Sigma}(a,b) &\ \longrightarrow\ \ \Ty\\
p_{b,c} : \Ch_{\Sigma}(b,c) &\ \longrightarrow\ \ \Ty,
\end{aligned}
\]
d'où
\[
\begin{tikzcd}
\Ch_{\Sigma}(a,b)\times\Ch_{\Sigma}(b,c) \ar[d,"p_{a,b}\times p_{b,c}"] \\
\Ty\times\Ty
\end{tikzcd}
\]
\msp{page 81}
qui est bel et bien fibrant (comme produit cartésien de deux foncteurs fibrants) et a les fibres qu'il faut! On trouve donc bien, par le sorite général sur les foncteurs fibrants (ici sur $\Ty\times\Ty$) un foncteur de catégories fibrées sur $\Ty\times\Ty$
\[
\Ch_{\Sigma}(a,b)\times\Ch_{\Sigma}(b,c) \ \longrightarrow\ \qquad \underbrace{\pi^*(\CH_{\Sigma}(a,c))}_{=\,(\Ty\times\Ty)\times_{\Ty}\Ch_{\Sigma}(a,c)\ \ \text{(\footnotemark)}} \qquad .
\]
\footnotetext{produit fibré par $\pi : \Ty\times\Ty \longrightarrow \Ty$.}
Comme on a d'autre part un foncteur du dernier membre (produit fibré) vers le facteur $\Ch_{\Sigma}(a,c)$, on gagne. (On l'aurait retrouvé par le raisonnement direct que je m'apprêtais de faire, et que j'ai biffé - mais je suis plus heureux comme ça!)

% foot: Version 22 décembre 2001
\printedpage{53}
% head: A. Grothendieck : Dérivateurs, Ch. VII, Édition des Universités de Montpellier II et Paris VII
Pour bien faire, il faudrait prouver l'\textit{associativité} des foncteurs de composition, pas seulement au niveau des objets, par des arguments comme ceci, mais ça revient simplement à une relation du type
\[
(w''*w')*w = w''*(w'*w)
\]
pour un système de trois flèches dans
\[
\Ch_{\Sigma}(a,b),\ \Ch_{\Sigma}(b,c)\ \ \text{et}\ \ \Ch_{\Sigma}(c,d)
\]
respectivement. Ça ne peut être que trivialement
\msp{page 82}
vrai, mais fastidieux à expliciter.

\dgrblock{dgrb325}{The total composition and its associativity}
\begin{dcomment}{exposition}{p81}{The total composition and its associativity}
Grothendieck returns to composition because the localization argument uses a functor on the total path categories, including morphisms that change the type. The commentary at manuscript pages~48--50 describes the fixed-type operation. The additional verification here is that this operation respects the arrows of the Grothendieck construction. Writing these arrows explicitly also settles the associativity left to the reader on pages~81--82.

For fixed endpoints, put $P_{a,b}(\tau)=\Ch_{\Sigma,\tau}(X;a,b)$. An arrow in its total category,
\[
 (\tau,c)\longrightarrow(\sigma,d),
\]
is a pair $(\varphi,\eta)$ with $\varphi:\tau\to\sigma$ and an endpoint-fixed natural transformation $\eta:c\Rightarrow d\varphi$. If a second arrow is $(\psi,\theta):(\sigma,d)\to(\rho,e)$, its composite with the first is
\[
 (\psi\varphi,(\theta\varphi)\eta).
\]
This is the concrete form of the fibred-category formalism used in the text.

To concatenate two such arrows, concatenate the type maps and glue their transformations at the common endpoint:
\[
 ((\varphi,\eta),(\varphi',\eta'))
       \longmapsto(\varphi*\varphi',\eta*\eta').
\]
Both components at the joining object are the identity. Hence the two definitions there agree, and the naturality square for every edge comes from exactly one of the two given naturality squares. The resulting objects remain $\Sigma$-paths because concatenation does not change an edge or its direction. A vertical component of a transformation is unrestricted: $P_{a,b}(\tau)$ is a full subcategory on its allowed objects, with the endpoint condition imposed on transformations.

Compatibility with composition is checked at vertices. On the first interval the two possible composites have component $(\theta\varphi)_i\eta_i$; on the second they have the corresponding primed component; at the joining vertex both have the identity. They are therefore equal. The base maps agree by functoriality of concatenation in $\underline{\mathrm{Ty}}$. Thus the formula defines the desired functor on total categories. Equivalently,
\[
 p_1^*\Ch_\Sigma(a,b)\times_{\underline{\mathrm{Ty}}\times\underline{\mathrm{Ty}}}
 p_2^*\Ch_\Sigma(b,c)
 \cong \Ch_\Sigma(a,b)\times\Ch_\Sigma(b,c),
\]
and the cartesian functor to the pullback along concatenation can be followed by the projection to $\Ch_\Sigma(a,c)$. This is the recognition made on page~80.

For associativity, number the vertices of three paths of lengths $n,m,l$ by successive blocks in $[0,n+m+l]$. The first block keeps its indices, the second is shifted by $n$, and the third by $n+m$. Either parenthesization produces those same indices, objects, and arrows. The same statement holds for the components of a glued transformation and for the underlying type map. Consequently both parenthesizations agree as functors, and the length-zero path is a strict unit. With these standard representatives, no later strictification is required.

One diagrammatic slip deserves explicit repair. The amalgamation on page~77 is
\[
 I_\tau\amalg_{\{1_\tau=0_{\tau'}\}}I_{\tau'},
\]
a pushout over a singleton category. The displayed $\varepsilon=\{0,1\}$ would glue two pairs of endpoints. Already for two forward one-edge intervals, that different pushout has two parallel arrows between two objects, rather than the two successive arrows of $I_{\Delta_+^2}$. The joining vertex has index $n$ in the numbering $0,\ldots,n+m$; the printed $n+1$ uses the wrong index. Neither slip affects the correct componentwise construction.
\end{dcomment}


Les sueurs froides sont terminées!

En résumé, le formalisme des $\Ch_{\Sigma}(X)$ et des $\Ch_{\Sigma}(X;a,b)$ développé jusqu'à présent donne satisfaction complète sur les points suivants :

\begin{enumerate}
\item[a)] Bon formalisme de composition des chemins (qui était dégueulasse avant).

\item[b)] Traitement élégant de la construction des catégories de fractions $\mathcal{M}\Sigma^{-1}$. Mais je m'attends que ce sera moins joli qu'avant, quand il s'agira de remplacer un $(\mathcal{M},\Sigma)$ par des $(\ \underbrace{\mathcal{M}(X)}_{=\,\uHom(X,\mathcal{M})}\ ,\Sigma_X)$ ($X$ une catégorie d'indices) \texttt{[plutôt} $\uHom(X^{o},\mathcal{M})$\texttt{]}. C'est pour y pallier qu'il me faudra développer des énoncés de comparaison, entre la catégorie $\Ch_{\Sigma}$ construite ici, et d'autres dont le type dans Hot se prête aux calculs de façon plus directe, comme dans l'état premier de la question. De plus, on sera sûrement obligé d'en venir à des Ind-catégories, pour pouvoir suivre la variance de $\mathcal{M}(X)\Sigma_X^{-1}$ pour $X$ variable.

\item[c)] Bon formalisme des `homotopies' et des compositions des homotopies (ce qui est une autre façon de parler des chemins et de leur composition).
\end{enumerate}
% foot: Version 22 décembre 2001

% ---- from Der-ChVII-456_p056-060.tex
\printedpage{54}
% head: A. Grothendieck : Dérivateurs, Ch. VII, Édition des Universités de Montpellier II et Paris VII
% foot: Version 22 décembre 2001

\msp{page 83}

\dgrblock{dgrb326}{Why the category of diagrams returns}
\begin{dcomment}{context}{p82}{Why the category of diagrams returns}
The reservation at the end of this section brings the calculation back to the purpose of a derivator. Localizing $\mathcal M$ answers the question at one point. To vary coefficients over an indexing category $A$, one must also localize
\[
 \mathcal M(A)=\operatorname{Fun}(A^{\mathrm{op}},\mathcal M),
 \qquad
 \Sigma_A=\{\eta\mid\eta_a\in\Sigma\text{ for every }a\in A\}.
\]
The opposite in the edition's correction records the presheaf convention. It is a variance convention, not an alteration of the localization theorem: the latter applies to either functor category once its convention is fixed.

There is a small but decisive issue of uniformity. An object of a path category for $\mathcal M(A)$ is a finite zigzag of natural transformations with one common word of signs over all $a\in A$. An arbitrary family of zigzags in $\mathcal M$ may use a different word, or require an unbounded length, as $a$ varies. Thus a description that works perfectly for individual objects still has work to do before it describes localization naturally in the parameter $A$. The ind-objects considered later organize precisely the freedom to enlarge a finite path type while retaining its functorial meaning.

This interpretation agrees with Grothendieck's subsequent account in his letter to Thomason of 2 April 1991. In \S2, printed page~9, he explicitly defines the associated prederivator by $A\mapsto\mathcal M(A)[\Sigma_A^{-1}]$, with presheaf diagrams and objectwise denominators; on printed page~8 he emphasizes the transformations between restriction functors as part of the structure \citep{Grothendieck1991Letter}. The present concern about types is a concrete instance of that requirement. It asks for controlled variation of the entire diagram theory, beyond the universal property of a single localization.
\end{dcomment}


\section*{8\quad Le lemme d'homotopie pour les morphismes d'intervalles}

C'est le suivant :

\textbf{Lemme 1.} \textit{Soient $\tau,\tau'\in\mathrm{Ty}$, et $f,g:\tau\rightrightarrows\tau'$ deux morphismes dans $\uTy$. Alors $f$, $g$ sont homotopes dans la catégorie $\uHom\,\underline{\mathrm{int}}(I_\tau,I_{\tau'})$ des morphismes croissants d'intervalles orientés de $I_\tau$ dans $I_{\tau'}$.}

Cette catégorie pourrait licitement se désigner par $\uHom_{\uTy}(\tau,\tau')$, ou $\uHom(\tau,\tau')$. C'est une sous-catégorie de la catégorie $\uHom(I_\tau,I_{\tau'})$ de \textit{tous} les foncteurs de $I_\tau$ dans $I_{\tau'}$, et une sous-catégorie pleine de la catégorie $\uHom^\varepsilon(I_\tau,I_{\tau'})$ des foncteurs $f$ qui sont compatibles avec la donnée des sources et des buts, i.e. qui rendent commutatif le diagramme
\[
\begin{tikzcd}[column sep=small]
 & \varepsilon=\{0,1\} \arrow[dl, hook', "{(\sigma_\tau,\beta_\tau)}"'] \arrow[dr, hook, "{(\sigma_{\tau'},\beta_{\tau'})}"] & \\
I_\tau \arrow[rr, "f"] & & I_{\tau'}.
\end{tikzcd}
\]

Si $f$, $g$ sont dans cette catégorie (pas forcément morphismes croissants d'ailleurs), une flèche de $f$ dans $g$ est donc la donnée, pour tout $i\in I_\tau$, d'une flèche $f(i)\xrightarrow{w_i}g(i)$ dans $I_{\tau'}$, avec les deux conditions :

\begin{enumerate}[label={}, leftmargin=*]
\item[a)] Si $i\in\{0_\tau,1_\tau\}$, $w_i=\mathrm{id}$, et
\item[b)] le carré commutatif habituel pour une flèche $i\longrightarrow j$ dans $I_\tau$. (Il suffit d'avoir cette compatibilité pour les flèches de transition $\ul{u}_i:i-1\longrightarrow i$ dans $I_\tau$ ($i\in[1,\operatorname{long}\tau]$).)
\end{enumerate}

Une autre façon de décrire cette catégorie $\uHom^\varepsilon$, c'est de dire que c'est la catégorie
\msp{page 84}
\[
\underline{\mathrm{Ch}}_\tau(I_{\tau'};0_{\tau'},1_{\tau'})\qquad(\simeq\uHom^\varepsilon(I_\tau,I_{\tau'})),
\]
mais $\uHom_{\uTy}(\tau,\tau')$ est la sous-catégorie pleine de cette catégorie formée des \textit{chemins croissants} dans $I_{\tau'}$. Le lemme d'homotopie signifie aussi que cette catégorie $\uHom_{\uTy}(\tau,\tau')$ est toujours \textit{connexe.} Elle n'est d'ailleurs pas forcément 0-connexe, car elle peut être vide. Notons ici, pour plus tard, la

\textbf{Proposition 1 :} \textit{Soient $\tau,\tau'\in\mathrm{Ty}$, soient $r$, $r'$ leurs longueurs utiles. Pour que $\mathrm{Hom}_{\uTy}(\tau,\tau')$ soit non vide, il faut et il suffit que l'une ou l'autre des deux conditions suivantes soit satisfaite :}

\begin{enumerate}[label={}, leftmargin=*]
\item[\textit{a)}] $r\geq r'+1$

\printedpage{55}
% head: A. Grothendieck : Dérivateurs, Ch. VII, Édition des Universités de Montpellier II et Paris VII
% foot: Version 22 décembre 2001

\item[\textit{b)}] \textit{$r=r'$, et (si $r\neq0$) $\tau$ et $\tau'$ ont même direction initiale (de sorte que $\tau_{\red}=\tau'_{\red}$).}
\end{enumerate}

(Démonstration $\pm$ immédiate.)

Notons que si $\tau$, $\tau'$ n'ont pas même direction initiale, alors tout morphisme $\tau\longrightarrow\tau'$ est forcément dégénéré sur le premier intervalle critique de $I_\tau$. On aura donc $\tau=\Delta_\varepsilon^{l_0}\cdot\tau_0$ (où $l_0$ est la longueur du premier intervalle critique, et où $\tau_0$ a \texttt{[la]} direction initiale opposée de celle de $\tau$), et un isomorphisme
\[
\uHom(\tau,\tau')\simeq\uHom(\tau_0,\tau'),
\]
où cette fois $\tau_0$, $\tau'$ ont même direction initiale. On supposera donc désormais que \textit{la direction initiale de $\tau$, $\tau'$ est la même.} La condition $\uHom(\tau,\tau')\neq\emptyset$ équivaut maintenant à $\underline{r\geq r'}$, ce que nous pouvons supposer. Je vais alors décrire un élément canonique
\msp{page 85}
$F_0$ de $\uHom(\tau,\tau')$, ainsi. Soient
\[
I_1,I_2,\dots,I_r
\]
les intervalles critiques de $I_\tau$, et
\[
I'_1,I'_2,\dots,I'_{r'}\qquad(r'\leq r)
\]
ceux de $I_{\tau'}$. Le foncteur $F_0$ proviendra de morphismes d'intervalles orientés (comme
\[
\left\{
\begin{array}{ccc}
I_1 & \xrightarrow{f_1} & I'_1\\[4pt]
I_2 & \xrightarrow{f_2} & I'_2\\[4pt]
 & \dots & \\[4pt]
I_{r'} & \xrightarrow{f_{r'}} & I'_{r'}
\end{array}
\right.\qquad)
\]
et sera donc forcément trivial sur les intervalles critiques $I_\alpha$ pour $r'<\alpha\leq r$ (il y en a si et seulement si $r>r'$). Comme $I_\tau$, $I_{\tau'}$ ont même direction initiale, il s'ensuit que pour $\alpha\in[1,r']$, $I_\alpha$ et $I'_\alpha$ ont même signature (à savoir $\varepsilon_\alpha=(-1)^{\alpha+1}\varepsilon_1$, où $\varepsilon_1$ est la signature commune de $I_1$ et $I'_1$). Donc
\[
I_\alpha=\Delta^{l_\alpha}_{\varepsilon_\alpha},\qquad I'_\alpha=\Delta^{l'_\alpha}_{\varepsilon_\alpha},
\]
et on définit $f_\alpha$ comme
\[
f_\alpha=\mathfrak{f}_{\varepsilon_\alpha,l'_\alpha}*\triv_\alpha,
\]
où
\[
\mathfrak{f}_{\varepsilon_\alpha,l'_\alpha}:\Delta_{\varepsilon_\alpha}\longrightarrow\Delta^{l'_\alpha}_{\varepsilon_\alpha}
\]
et où
\[
\triv_\alpha:\Delta^{l_{\alpha-1}}_{\varepsilon_\alpha}\longrightarrow1. % [sic] l_{\alpha-1} -> l_\alpha-1, the exponent is printed as l with subscript \alpha-1, while I_\alpha = \Delta^{l_\alpha} = \Delta * \Delta^{l_\alpha-1}
\]
En d'autres termes, pour $\varepsilon=+1$ (p. ex.) on doit définir un morphisme des intervalles type $\underbrace{\mathbf{\Delta}^+_{l_\alpha}}_{=[0,l_\alpha]}\longrightarrow\underbrace{\mathbf{\Delta}^+_{l'_\alpha}}_{=[0,l'_\alpha]}$, et on pose (avec les numérotations internes à $\mathbf{\Delta}_{l_\alpha}$, $\mathbf{\Delta}_{l'_\alpha}$)
\[
f_\alpha(0)=0,\qquad f_\alpha(i)=l'_\alpha\ \text{ pour }i\in[1,l_\alpha].
\]

\dgrblock{dgrb327}{The two orders in the interval lemma}
\begin{dcomment}{exposition}{p84}{The two orders in the interval lemma}
Grothendieck now wants the choice of a type map to disappear after homotopy. The statement is about a particularly constrained mapping category, and its constraints explain both the nonemptiness criterion and the proof. Number the source and target vertices by $[0,n]$ and $[0,m]$, and write their edge signs as $\varepsilon_i$ and $\delta_j$. An object is a nondecreasing function $F:[0,n]\to[0,m]$ satisfying $F(0)=0$, $F(n)=m$, and
\[
 F(i-1)<F(i)\quad\Longrightarrow\quad
 \delta_{F(i-1)+1}=\cdots=\delta_{F(i)}=\varepsilon_i.
\]
The condition says that one source edge can traverse a run of target edges only when their directions agree. It cannot cross a change of direction in the target. An arrow $F\to G$ is instead defined by the categorical order of the target: for every $i$, there must be an arrow $F(i)\to G(i)$ in $I_{\tau'}$. Such arrows, when present, are unique, and their naturality is automatic because the target is a poset. In particular the ordering of maps is generally different from pointwise numerical comparison.

Let $r,r'$ count the maximal constant-sign runs. Each target run must be traversed using at least one source edge of the same sign. Distinct target runs require distinct source runs in the same order, because their signs alternate and a source run cannot pass a target turning vertex. This proves $r\geq r'$ when the initial signs agree. If the initial signs disagree, the whole first source run is forced to be constant at the initial target vertex, and the necessary inequality is $r-1\geq r'$. Conversely, under either inequality, map the first edge of each of the earliest matching source runs across the corresponding whole target run, collapse its remaining edges, and collapse every unused source run at the final target vertex. This constructs an admissible map and proves Proposition~1. If the target has no edges, there is exactly one map from every source; if the source has no edges and the target has an edge, there is none.

These observations also clarify the word ``connected'' in the text. Its assertion is that any two objects are connected by a finite zigzag of transformations; it allows the mapping category to be empty. For a nonempty mapping category, the lemma asserts that its nerve is connected. Grothendieck then asks whether a stronger contraction exists. The following example separates that question from the proposed terminal-object argument, and the development after it gives a uniform contraction.
\end{dcomment}


\printedpage{56}
% head: A. Grothendieck : Dérivateurs, Ch. VII, Édition des Universités de Montpellier II et Paris VII
% foot: Version 22 décembre 2001

\msp{page 86}
L'idée donc, c'est d'être dans les dispositions d'un voyageur qui, tenu à un plan de voyage par étapes décrit par $\tau$, voudrait arriver aussi vite que possible à $1_{\tau'}$, en partant de $0_{\tau'}$. Il va dès la première étape $\Delta_\varepsilon\subseteq I_1$, se porter jusqu'au bout de $I'_1$, regrettant ne pas pouvoir aller plus loin (à cause du plan de voyage). Il doit rester inactif jusqu'à ce qu'il soit arrivé au moment marqué par l'extrémité finale de $I_1$ (si celle-ci est au delà de celle de sa première flèche de transition, i.e. si $l_1>1$), puis il fait son deuxième grand pas jusqu'à l'extrémité terminale de $I'_2$ etc. Il peut espérer, à chaque moment de son voyage, être en avance (ou du moins pas en retard) sur les autres voyageurs, qui auraient suivi de tactiques de progression différentes. Si cet espoir était justifié, cela signifierait que $F_0$ serait un plus grand élément de $\uHom(\tau,\tau')$ (\textbf{NB} toutes ces catégories $\uHom$ entre intervalles sont ‘ordonnées’, i.e. il y a au plus une flèche entre deux objets), i.e. un objet final. Cela impliquerait que $\uHom(\tau,\tau')$ serait même contractile - chose peut-être vraie, mais sûrement pas de façon si simpliste je crains. Je me suis d'ailleurs canulé : $F_0$ est sans doute bien un ‘plus grand’ élément, mais pour l'ordre défini par la numérotation des sommets de $I_{\tau'}$ (donc $F\leq G$ si et seulement si $F(i)\leq G(i)$ $\forall i\in I_\tau$), et nullement pour l'ordre défini par la structure de catégorie de $\uHom(\tau,\tau')$.

\msp{page 87}
L'idée pour prouver la connexité de $\uHom(\tau,\tau')$, c'est de ramener tout $F\in\mathrm{Hom}(\tau,\tau')$, par une suite d'homotopies simples, à la forme normale. Rappelons que $F$ est donné par une suite croissante de flèches de transition
\[
\ul{u}_{i_1},\ \ul{u}_{i_2},\ \dots\ ,\ \ul{u}_{i_s},\qquad i_1<i_2<\cdots<i_s,
\]
et une suite d'entiers $\geq1$
\[
\lambda_1,\ \lambda_2,\ \dots\ ,\ \lambda_s
\]
avec la condition que
\[
\Delta^{\lambda_1}_{\varepsilon_{i_1}}\Delta^{\lambda_2}_{\varepsilon_{i_2}}\cdots\Delta^{\lambda_s}_{\varepsilon_{i_s}}=\tau'. \tag*{$(*)$}
\]
Les $\ul{u}_{i_\alpha}$ sont les flèches de transition de $I_\tau$ \textit{critiques} pour $F$, i.e. non dégénérées pour $F$, et $\lambda_\alpha$ est la \textit{partie} de $\ul{u}_{i_\alpha}$ pour $F$. Le foncteur normalisé (‘pressé’) $F_0$ est distingué par les conditions suivantes a) et b) \texttt{[et c)]}. Tout d'abord :

\dgrblock{dgrb328}{A quickest map that is neither initial nor terminal}
\begin{dcomment}{examples}{p86}{A quickest map that is neither initial nor terminal}
The normal form reaches each target turning vertex as early as the source permits. It maximizes numerical progress, but the categorical direction reverses whenever the target changes direction. A four-object mapping category makes the distinction visible without any large diagram.

Take
\[
 \tau=\Delta_+^2\Delta_-^2,
 \qquad \tau'=\Delta_+\Delta_-.
\]
The target is $0\to1\leftarrow2$. Every increasing endpoint-preserving functor has vertex list
\[
 F_{a,b}=(0,a,1,b,2),\qquad a\in\{0,1\},\quad b\in\{1,2\}.
\]
Indeed the middle source vertex must map to the target turning vertex $1$: it cannot lie below $1$, since the remaining negative edges could not reach $2$, and it cannot lie above $1$, since the initial positive edges could not reach it. The choices of $a,b$ are then unrestricted in the displayed ranges. The exact arrow relation is
\[
 F_{a,b}\longrightarrow F_{a',b'}
 \quad\Longleftrightarrow\quad a\leq a'\ \text{ and }\ b\geq b',
\]
where the inequalities on the right are numerical. The reversal of the second inequality is forced by $2\to1$ in the target.

The quickest functor is $F_0=F_{1,2}$. It is numerically greater than or equal to every other functor, but the terminal object of this mapping category is $F_{1,1}$ and its initial object is $F_{0,2}$. In particular $F_{1,1}$ has no arrow to $F_0$, and $F_0$ has no arrow to $F_{0,2}$. The objection on manuscript page~86 is therefore exact, even in this small case.

Nevertheless a contraction to $F_0$ is immediate once we normalize one target run at a time. Define $R(F_{a,b})=F_{1,b}$. There is a natural transformation $\mathrm{id}\Rightarrow R$, and another $\mathrm{const}_{F_0}\Rightarrow R$. Hence
\[
 \mathrm{id}\ \Longrightarrow\ R\ \Longleftarrow\ \mathrm{const}_{F_0}.
\]
This is the pattern of the general argument: each step advances across only one constant-sign target run, so all its components have one common categorical direction. The entire numerical maximum need not define a single transformation.
\end{dcomment}


\dgrblock{dgrb329}{A uniform contraction of the interval mapping category}
\begin{dcomment}{development}{p86}{A uniform contraction of the interval mapping category}
Grothendieck's proposed normal form can be made into a contraction by advancing across one target run at a time. This answers the question on page~86 for the increasing endpoint-preserving mapping categories used in this section. We did not find this particular argument in the literature checked for the present commentary: the source's \S8 proves connectedness, while the institutional abstract of \citet{FernandesKoppitzMusunthia2019} concerns order-preserving self-transformations of a fence, without these simultaneous endpoint and linear-order constraints. We checked that abstract and metadata; its manuscript could not be retrieved. The result below is an elementary supplement, without a claim of priority.

\textbf{Proposition (run-by-run contraction).}\ours\
\textit{Let $P(\tau,\tau')$ be the category of nondecreasing endpoint-preserving functors $I_\tau\to I_{\tau'}$, with endpoint-fixed natural transformations. If $P(\tau,\tau')$ is nonempty, its identity functor is connected to the constant functor at the quickest map by a finite zigzag of natural transformations. In particular its nerve is contractible. One can use at most $r'$ transformations, where $r'$ is the number of target runs.}

\textit{Proof.} We induct on $r'$. When the target has no edges, every functor is the same constant functor and $P(\tau,\tau')$ is the terminal category. Assume that the target is nontrivial. The source is then nontrivial too. If its initial sign differs from the target's, every map is constant on the entire first source run, as proved above. Every transformation is the identity at those vertices. Restriction past that run is consequently an isomorphism of mapping categories. We may use this isomorphism to assume that the first source and target runs have the same sign $\varepsilon$.

Let the first source run have vertices $0,\ldots,a$, and the first target run have vertices $0,\ldots,b$, where $a,b\geq1$. The numerical clamping map
\[
 h:[0,m]\longrightarrow[0,m],\qquad h(j)=\max\{j,b\},
\]
is a functor of the target interval category. On its first run it collapses every arrow to an identity at $b$; on every later edge it is the identity. These checks on generating edges prove functoriality, including at the turning vertex $b$.

For each $F\in P(\tau,\tau')$, define
\[
 (RF)(0)=0,\qquad (RF)(i)=h(F(i))\quad(1\leq i\leq n).
\]
The first source run cannot pass the first target turning vertex. Thus $F(i)\leq b$ for $0\leq i\leq a$, and $(RF)(i)=b$ for $1\leq i\leq a$. The new first edge is an allowed traversal from $0$ to $b$ of sign $\varepsilon$; all later edges are obtained by applying the functor $h$ to their old images. This proves that $RF$ is a functor. Its vertex map remains nondecreasing, it fixes both endpoints, and hence it is again an object of $P(\tau,\tau')$.

For a transformation $F\Rightarrow G$, apply $h$ to its components at vertices $i\geq1$ and use the identity at vertex $0$. The resulting components define a transformation $RF\Rightarrow RG$. Indeed both functors are already valid interval functors, their target is thin, and each component exists by functoriality of $h$. Identities and composites are preserved componentwise. We have therefore constructed an endofunctor $R$ of $P(\tau,\tau')$.

The comparison with the identity has one uniform direction. If $\varepsilon=+$, the unique arrows $j\to h(j)$ give $\mathrm{id}_P\Rightarrow R$; if $\varepsilon=-$, the arrows $h(j)\to j$ give $R\Rightarrow\mathrm{id}_P$. At vertex $0$ the component is the identity, and at the final vertex it is also the identity. Naturality in the source interval and in $F$ follows from thinness of the target and of the mapping category. Only values in the first target run have changed, which is why this comparison exists despite the failure of a global numerical-maximum argument.

Let $Q\subseteq P$ be the full subcategory of maps with values $0,b,\ldots,b$ on the first source run. The endofunctor $R$ takes values in $Q$ and restricts to the identity there. Delete the first source run, retain its last vertex as the new initial vertex, and do the same for the target. Restriction identifies $Q$ with the endpoint-preserving mapping category of these two remaining intervals. To verify this identification on morphisms, observe that every component over the deleted prefix is forced to be an identity; restriction and the prescribed extension are inverse functors. The remaining target has $r'-1$ runs. Since $Q$ contains $RF$ for every $F$, it is nonempty, so induction applies.

Compose the inductive zigzag on $Q$ with $R:P\to Q$ and the inclusion $Q\hookrightarrow P$. Adjoining the single comparison between $\mathrm{id}_P$ and $R$ gives a zigzag to the constant functor at the map that crosses the first target run on the first available source edge and then follows the inductively quickest map. This is exactly Grothendieck's normal form, after the forced initial deletion when necessary. The number of comparisons increases by one, hence is at most $r'$. A natural transformation induces a simplicial homotopy after taking nerves, so this categorical contraction contracts the nerve as asserted.\hfill$\square$

The proposition supplies more than a path joining two chosen maps: its choices of intermediate functors are simultaneous and natural in the object of the mapping category. It concerns the increasing subcategory. No conclusion about all endpoint-preserving functors, without numerical monotonicity, is being used here.
\end{dcomment}


\begin{enumerate}[label={}, leftmargin=*]
\item[a)] La suite des signes $\varepsilon_{i_1},\varepsilon_{i_2},\dots,\varepsilon_{i_s}$ est ‘alternée’, i.e. $\varepsilon_{i_\alpha}\neq\varepsilon_{i_{\alpha+1}}$ ($\alpha\in[1,s-1]$). (Cela traduit le fait que le voyageur pressé ne fait pas en plusieurs pas ce qu'il pourrait faire en un seul.)

Comme le premier de ces signes est forcément $\varepsilon'_1=\varepsilon_1$, cela signifie que la suite des signes est la même que celle des $\varepsilon_1$, $\varepsilon_2$, \dots, $\varepsilon_{r'}$ (à la longueur près). Mais mieux, que l'expression $(*)$ de $\tau'$ est son unique expression réduite, donc que $s=r'$ (longueur utile de $\tau'$), et que les $\lambda_i$ sont les longueurs $l'_i$ des intervalles critiques de $I_{\tau'}$.

\printedpage{57}
% head: A. Grothendieck : Dérivateurs, Ch. VII, Édition des Universités de Montpellier II et Paris VII
% foot: Version 22 décembre 2001

\msp{page 88}
Soit donc $\alpha\in[1,s=r']$, et considérons la flèche de transition $F$-critique $\ul{u}_{i_\alpha}$, soit $j_\alpha$ l'indice marquant l'intervalle critique de $I_\tau$ contenant la partie de $\ul{u}_{i_\alpha}$ dans $I_\tau$,
\[
\ul{u}_{i_\alpha}\in\mathrm{Fl}(I_{j_\alpha}),
\]
alors $F$ induit des morphismes (croissants) d'intervalles orientés
\[
F_\alpha=F|I_{j_\alpha}:I_{j_\alpha}\longrightarrow I'_\alpha
\]
et est trivial sur les autres intervalles critiques (ceux pas de la forme $I_{j_\alpha}$). La donnée des $F_\alpha$ détermine donc $F$ sur les $I_{j_\alpha}$, donc sur \textit{tous} les $I_j$, donc partout (par l'opération $*$ de composition par recollement). De plus la condition a) assure que chaque $F_\alpha$ est de nature assez particulière : une seule flèche de transition critique $\ul{u}_{i_\alpha}$. De plus, pour qu'on ait $F=F_0$, il faudra :

\item[b)] $\forall\alpha\in[1,r']$, $\ul{u}_{i_\alpha}$ est la \textit{première} flèche de transition de l'intervalle critique $I_{j_\alpha}$ qui la contient.

Et enfin :

\item[c)] La suite des intervalles critiques de $I_\tau$
\[
I_{j_1},\ I_{j_2},\ \dots\ ,\ I_{j_{r'}}\qquad(\text{de signes alternés }\varepsilon_1,\ -\varepsilon_1,\ \varepsilon_1,\ \ldots)
\]
associés aux flèches $F$-critiques, n'est autre que celle
\[
I_1,\ I_2,\ \dots\ ,\ I_{r'}
\]
des $r'$ premiers intervalles critiques de $I_\tau$, i.e.
\[
(j_1,j_2,\dots,j_{r'})=(1,2,\dots,r').
\]
\end{enumerate}

\msp{page 89}
Ceci vu, l'idée naturelle pour ‘déformer’ $F$ en $F_0$, c'est de trouver une suite d'homotopies qui transforment $F$ de façon à lui faire acquérir successivement les caractéristiques a), b), c) du voyageur pressé. Au besoin, pour aller plus vite au but, on utilisera l'hypothèse de récurrence sur $n=\operatorname{long}\tau$, que le théorème de connexité est déjà prouvé quand $I$ est de longueur $\leq n-1$. (NB Ce théorème est trivial pour $n\leq1$, car alors les catégories $\uHom(\tau,\tau')$ sont toutes ou bien vides, ou bien ponctuelles.)

\msp{page 90}
1) On veut assurer a). Pour ceci, on utilise l'hypothèse de récurrence, pour remplacer un groupe de \texttt{[mot illisible]} flèches critiques consécutives de même signe, par une seule flèche critique. Cette méthode ne foire que si les flèches critiques sont toutes de même signe $\varepsilon$ et incluent $\ul{u}_1$ et $\ul{u}_n$. Mais l'hypothèse de récurrence appliquée à $I=[1,n-1]\subseteq I_\tau$ nous ramène au cas où il n'y a que deux flèches critiques $\ul{u}_1$ et $\ul{u}_n$. Si le premier intervalle critique n'est pas réduit à $\Delta_\varepsilon$, i.e. soit de longueur 2, alors l'hypothèse de récurrence appliquée à $I=[1,n]\subseteq I_\tau$ nous ramène au cas où les deux flèches critiques sont $\ul{u}_1$ et $\ul{u}_2$, et on est
\printedpage{58}
% head: A. Grothendieck : Dérivateurs, Ch. VII, Édition des Universités de Montpellier II et Paris VII
% foot: Version 22 décembre 2001
\msp{page 91}
alors ramené au cas où $\boxed{I=\Delta^2_\varepsilon,\ \tau'=\Delta^2_\varepsilon.}$ Sinon, au cas où $\boxed{\tau=\Delta_\varepsilon\Delta^{l_2}_{-\varepsilon}\Delta_\varepsilon,\ \tau'=\Delta^2_\varepsilon,}$ et où les deux flèches critiques sont encore la première et la dernière, et $\tau'=\Delta^2_\varepsilon$. Ce sont là les deux cas irréductibles, qu'il faudra traiter à part.

Ces cas admis, on a donc moyen pour définir $F$ de façon à satisfaire a). Ensuite on veut :

2) Assurer que la suite des intervalles critiques $I_{j_1}$, $I_{j_2}$, \dots, $I_{j_{r'}}$ qui contient les flèches critiques $\ul{u}_{i_1}$, \dots, $\ul{u}_{i_{r'}}$, soit la suite $I_1$, \dots, $I_{r'}$. C'est encore immédiat par l'hypothèse de récurrence, le seul cas irréductible étant celui où $r'=1$ (une seule flèche critique), et où il s'agit de faire descendre cette flèche critique jusqu'à $\ul{u}_1$. Si $I_1$ est de longueur $>1$, alors on fait descendre en $\ul{u}_2$ par hypothèse de récurrence, et le cas irréductible est celui où $\boxed{\tau=\Delta^2_\varepsilon,\ \tau'=\Delta_\varepsilon,}$ une flèche critique $\ul{u}_2$, la faire descendre en $\ul{u}_1$. Si par contre $I_1$ \texttt{[est]} de longueur 1, le cas irréductible est $\boxed{\tau=\Delta_\varepsilon\Delta^{l_2}_{-\varepsilon}\Delta_\varepsilon,\ \tau'=\Delta_\varepsilon,}$ une flèche critique $\ul{u}_{l_2+2}$, la faire descendre en $\ul{u}_1$.

Sous réserve de traiter les deux nouveaux cas irréductibles, on peut déformer $F$ en un foncteur qui satisfait à a) et à c). Reste à normaliser chacun des $F_\alpha:I_\alpha\longrightarrow I'_\alpha$ ($\alpha\in[1,r']$). Donc on est ramené à $\Delta^n_\varepsilon\longrightarrow\Delta_\varepsilon$ : une seule flèche critique $\ul{u}_\alpha$, $\alpha\in[1,n]$, à faire descendre en $\ul{u}_1$. Si $n\geq2$, par hypothèse de récurrence elle descend en $\ul{u}_2$, cas supposé traité en $2^o$). Donc
\msp{page 92}
on gagne, sous réserve de traiter les 4 cas irréductibles.

A) Cas $\tau=\Delta^2_\varepsilon$, $\tau'=\Delta^2_\varepsilon$ ou $\Delta_\varepsilon$.

La catégorie $\uHom(\Delta^2_\varepsilon,\Delta_\varepsilon)$ a deux éléments, les morphismes $\Delta^2_\varepsilon\longrightarrow\Delta_\varepsilon$ correspondent aux flèches de transition de $\Delta^2_\varepsilon$ ($=\mathbf{\Delta}^+_2$ ou $\mathbf{\Delta}^-_2$). Dans le cas $\varepsilon=+1$, ce sont les deux applications de $\{0,1,2\}$ dans $\{0,1\}$, qui appliquent 0 en 0, 2 en 1, et 1 soit en 1 ($F_0$), soit en 1 \texttt{[plutôt en 0]} ($F_1$).
\[
\begin{tikzcd}
0 \arrow[r] & 1 \arrow[r] & 2
\end{tikzcd}
\]
\[
\begin{tikzcd}
0 \arrow[r] & 1
\end{tikzcd}
\]

On a évidemment une flèche $F_1\longrightarrow F_0$ (le voyageur pressé arrive plus vite \ldots). Similaire pour $\varepsilon=-1$,
\[
\uHom(\Delta^2_\varepsilon,\Delta_\varepsilon)\simeq\mathbf{\Delta}^1,\qquad\text{catégorie 0-connexe.}
\]

La catégorie $\uHom(\Delta^2_\varepsilon,\Delta^2_\varepsilon)$ a trois éléments, s'identifiant à l'ensemble des parties non vide de l'ensemble $\mathrm{Tr}(\tau)$ à deux éléments des flèches de transition $\ul{u}_1$, $\ul{u}_2$ de $\tau$ : il y a $F_{\ul{u}_1}=F_0$, $F_{\ul{u}_2}$, $F_{\ul{u}_1,\ul{u}_2}$. Si $\varepsilon=+1$, il s'agit des applications croissantes (sans plus) de $\{0,1,2\}$ dans $\{0,1,2\}$, appliquant 0 en 0, 2 en 2, et 1 en ce qu'elle voudra. % [sic] parties non vide -> parties non vides, agreement

% ---- from Der-ChVII-456_p061-065.tex
\printedpage{59}
% head: A. Grothendieck : Dérivateurs, Ch. VII, Édition des Universités de Montpellier II et Paris VII
% foot: Version 22 décembre 2001

\begin{tabular}{lll}
Donc & $1\longmapsto 2$ & correspond à \quad $F_0$\\[1ex]
 & $1\longmapsto 1$ & correspond à \quad $F_1$\\[1ex]
 & $1\longmapsto 0$ & correspond à \quad $F_2$
\end{tabular}

et on a bien sur $F_2\longrightarrow F_1\longrightarrow F_0$, % [sic] bien sur -> bien sûr, accent missing
\[
\uHom(\Delta^2_\varepsilon,\Delta^2_\varepsilon)\simeq\boldsymbol{\Delta}^2,\qquad\quad\text{catégorie 0-connexe.}
\]

Dans ces deux cas, deux éléments de $\mathrm{Hom}(\tau,\tau')$ sont reliés par une homotopie \textit{simple.}

\msp{page 93}

B) Cas $\tau=\Delta_\varepsilon\Delta^l_{-\varepsilon}\Delta_\varepsilon$, $\tau'=\Delta^2_\varepsilon$ ou $\Delta_\varepsilon$.

Notons que l'on a, pour $\tau,\tau'\in\uTy$ quelconque, un isomorphisme de catégories
\[
\uHom(\tau,\tau')\longrightarrow\uHom(\tau'^o,\tau^o),
\]
exprimant simplement que
\[
\tau\longmapsto\tau^o
\]
est un \texttt{[2]}-foncteur de $\uTy$ dans $\uTy$, et même un automorphisme de cette \texttt{[2]}-catégorie. Donc dire que $\uHom(\tau,\tau')$ est connexe revient à dire que $\uHom(\tau^o,\tau'^o)$ l'est. Donc on peut supposer ci-dessus que $\varepsilon=+1$ (et pareil dans A!).
\[
\begin{aligned}
\tau :\quad & \longrightarrow\ \underbrace{\longleftarrow\ \longleftarrow\ \cdots\ \longleftarrow}_{l}\ \longrightarrow\\
\tau' :\quad & \longrightarrow\ \longrightarrow\quad\text{ou}\quad\longrightarrow\ .
\end{aligned}
\]

Mais tous les morphismes $\tau\longrightarrow\tau'$ sont triviaux sur l'intervalle critique $I_2$ (rétrograde), ce qui implique que
\[
\uHom(\tau,\tau')\xleftarrow{\ \sim\ }\uHom(\Delta^2_\varepsilon,\tau'),
\]
via le foncteur de contraction canonique
\[
I_\tau\longrightarrow I_{\Delta^2_\varepsilon}=\boldsymbol{\Delta}^2.
\]

Or dans A on a vu que $\uHom(\Delta^2_\varepsilon,\tau')$ est connexe (pour $\tau'$ de la forme $\Delta_\varepsilon$ ou $\Delta^2_\varepsilon$).

On a gagné!

\newpage
\printedpage{60}
% head: A. Grothendieck : Dérivateurs, Ch. VII, Édition des Universités de Montpellier II et Paris VII
% foot: Version 22 décembre 2001

\msp{page 94}

\dgrblock{dgrb330}{The four small cases and the completed proof}
\begin{dcomment}{exposition}{p93}{The four small cases and the completed proof}
Grothendieck's induction reduces the normalization to two monotone cases and two cases with one intervening run of the opposite sign. The local calculations are sound; the reversal formula printed at the start of case~B needs its arguments restored to their original order. Reversal of the numbered interval gives
\[
 P(\tau,\tau')\cong P(\tau^o,\tau'^o),
 \qquad F\longmapsto r_{\tau'}F r_\tau^{-1},
\]
where $r_\tau:I_\tau\to I_{\tau^o}$ reverses the vertex numbering. It is a covariant isomorphism of the underlying categories, since the reversed word also changes every sign. The displayed $P(\tau'^o,\tau^o)$ in the edition would reverse the direction of a functor and is not this operation. The next sentence in the source already has the correct order.

For $\tau=\Delta_+^2$ and target $\Delta_+^k$ with $k=1$ or $2$, a map is determined by $F(1)=j$, $0\leq j\leq k$. There is a unique transformation from $j$ to $j'$ exactly when $j\leq j'$. Thus the two mapping categories are the chains $[1]$ and $[2]$. With negative signs the order reverses, producing isomorphic chains. In particular any two objects are joined by one simple homotopy. These are the calculations on pages~92--93, including the edition's correction of the repeated value $1$ in the first example.

For $\tau=\Delta_+\Delta_-^l\Delta_+$ and a positive target, every edge of the negative middle run must collapse: a nondecreasing map cannot carry a negative source edge to a nonconstant positive target arrow. It therefore factors uniquely through the contraction $I_\tau\to I_{\Delta_+^2}$. The factorization also holds for transformations, because all components over the collapsed middle run coincide. This proves the displayed isomorphism of mapping categories and reduces the remaining cases to the preceding chains. Two small corrections clarify the longer induction on page~90. A first source run that is not a singleton has length at least two; the printed ``length two'' omits that qualification. More substantially, to concentrate all activity except the final edge, the inductive subinterval must be $[0,n-1]$, not the printed $[1,n-1]$. The latter fixes both interior boundary values, and cannot eliminate a nonzero traversal between them: the identity on a positive three-edge interval already demonstrates the problem. On the prefix $[0,n-1]$, induction instead leaves its first edge active and its remaining edges constant, while the final edge is unchanged. The following use of $[1,n]$ then moves the tail activity to the second edge when the first run has length at least two. With these corrections the stated reductions have their intended meaning. The run-by-run contraction above gives a full alternative proof of their connectedness conclusion.
\end{dcomment}


\dgrblock{dgrb331}{The path model and simplicial localization}
\begin{dcomment}{later}{p93}{The path model and simplicial localization}
The interval lemma makes changes of parameter homotopically harmless in a precise local sense. It is natural to compare the resulting path constructions with Dwyer--Kan localization, but the categories being compared must first be specified. In the present text, $\Ch_{\Sigma,\tau}(X;a,b)$ is full on the diagrams whose backward edges lie in $\Sigma$; its interior vertical transformations are arbitrary arrows of $X$. In the hammock construction, the vertical arrows are required to lie in the chosen weak-equivalence subcategory as well. These prescriptions agree when $\Sigma=\mathrm{Fl}(X)$, and can differ for a proper subcategory of denominators.

\citet[Definition~2.1, Proposition~2.2 and Proposition~3.1]{DwyerKan1980Calculating} construct the hammock mapping simplicial sets, compare them with simplicial localization, and identify their components with ordinary localization. Their definition also permits strict composition by concatenation followed by reduction. This gives a useful historical comparison: their work appeared in 1980, before these notes of 1990. It supplies an established construction retaining homotopical information about fractions, but its comparison theorem does not identify the nerve of the particular full path category in this manuscript without an additional argument.

Here the elementary result is already valuable on its own. It says that a category of allowed parameter maps carries a contraction whenever it has an object. The next section transports natural transformations between such maps through precomposition on path diagrams. That step explains why Grothendieck needed the homotopies to remain inside the increasing endpoint-preserving category, rather than merely to exist between maps of the underlying realizations.
\end{dcomment}


\section*{9\quad Théorème d'homotopie pour les catégories de chemins}

Soient $X$ une catégorie, $\Sigma\subseteq\mathrm{Fl}(X)$ \texttt{[un]} ensemble admissible de flèches dans $X$, $a,b\in\mathrm{Ob}\,X$. Considérons pour $\tau\in\mathrm{Ty}$ variable, les catégories
\[
\uCh_{\Sigma,\tau}(X),\qquad\quad\uCh_{\Sigma,\tau}(X;a,b),
\]
qui forment les fibres des catégories Cat-fibrées
\[
\begin{tikzcd}[column sep=small]
\uCh_\Sigma(X) \arrow[dr] & & \uCh_\Sigma(X;a,b) \arrow[ll, "\text{incl.}"'] \arrow[dl]\\
 & \uTy &
\end{tikzcd}
\]
sur la catégorie $\uTy$. Ainsi, pour $\tau,\tau'\in\mathrm{Ty}$, et pour toute
\[
\varphi:\tau\longrightarrow\tau',
\]
on a un carré commutatif d'homomorphismes
\[
\begin{tikzcd}
\uCh_{\Sigma,\tau'}(X) \arrow[d, "\varphi^*_X"'] & \uCh_{\Sigma,\tau'}(X;a,b) \arrow[l, "\text{incl.}"'] \arrow[d, "\varphi^*_{a,b}"]\\
\uCh_{\Sigma,\tau}(X) & \uCh_{\Sigma,\tau}(X;a,b), \arrow[l, "\text{incl.}"']
\end{tikzcd}
\]
avec transitivité pour $\varphi$ variable (composition $\psi\circ\varphi$ de flèches dans $\uTy$).

En fait, pour $\varphi$ variable dans $\uHom(\tau,\tau')$, les foncteurs images inverses $\varphi^*_X$, $\varphi^*_{a,b}$ dépendent \textit{fonctoriellement} de $\varphi$; i.e. on a des \textit{foncteurs}
\[
\begin{tikzcd}[row sep=small]
 & \uHom(\uCh_{\Sigma,\tau'}(X),\uCh_{\Sigma,\tau}(X))\\
\uHom(\tau,\tau') \arrow[ur] \arrow[dr] & \\
 & \uHom(\uCh_{\Sigma,\tau'}(X;a,b),\uCh_{\Sigma,\tau}(X;a,b))
\end{tikzcd}
\tag*{$(*)$}
\]

\newpage
\printedpage{61}
% head: A. Grothendieck : Dérivateurs, Ch. VII, Édition des Universités de Montpellier II et Paris VII
% foot: Version 22 décembre 2001

\msp{page 95}

\dgrblock{dgrb332}{Why a homotopy of types acts on every path}
\begin{dcomment}{exposition}{p94}{Why a homotopy of types acts on every path}
Grothendieck now uses the interval calculation to remove a choice from the theory: the operation induced on paths should depend on the source and target types, rather than on the particular subdivision or collapse chosen between them. The relevant assertion concerns natural transformations, not merely the induced maps of connected components. This stronger form is what will allow a diagram of categories to descend to a diagram in $\Cathot$.

Write $P_\tau=\uCh_{\Sigma,\tau}(X)$ and let $\mathcal T(\tau,\tau')$ be the category of endpoint-preserving type maps used in the preceding section. If $\eta:\varphi\Rightarrow\psi$ is a morphism of this category and $F:I_{\tau'}\to X$ is a permitted path, the transformation
\[
 F\eta:F\varphi\Longrightarrow F\psi
\]
is a morphism in $P_\tau$. Indeed, both its objects are permitted paths because $\varphi$ and $\psi$ carry backward arrows to backward arrows or identities. The fibre $P_\tau$ is the \emph{full} subcategory of the functor category on these objects: the components of $F\eta$ need not belong to $\Sigma$. Naturality in $F$ is precisely the interchange law for natural transformations. Thus whiskering gives a functor
\[
 \mathcal T(\tau,\tau')\longrightarrow\uHom(P_{\tau'},P_\tau),
 \qquad \varphi\longmapsto\varphi^*.
\]
This is the elementary two-dimensional mechanism behind the theorem.

The endpoint condition does a second piece of work. For a transformation of type maps its two endpoint components are identities, so $(F\eta)_0=\id_{F(0)}$ and $(F\eta)_1=\id_{F(1)}$. The displayed functor therefore lands in the category of functors and vertical transformations over $X\times X$. Restriction to the strict fibre at $(a,b)$ consequently preserves the entire homotopy. An arbitrary homotopy of functors over an unspecified base would not have this property. The words ``au dessus de $X$'' on manuscript page~97 should here be understood as ``over $X\times X$'', as the following diagram makes explicit.

Finally, connectedness of $\mathcal T(\tau,\tau')$, proved by the interval homotopy lemma, supplies a finite zigzag from $\varphi$ to $\psi$. Applying the whiskering functor supplies the required finite zigzag of natural transformations. Taking nerves turns each transformation into a simplicial homotopy. No assertion that the individual pullback functors are equivalences is made: the theorem compares parallel operations, and this is exactly what the next construction needs.
\end{dcomment}


Cela provient de l'inclusion
\[
\uHom(\tau,\tau')\hookrightarrow\uHom^{\varepsilon}(I_\tau,I_{\tau'})
\]
et de la description des catégories $\uCh_\tau$, $\uCh_{\tau'}$ comme catégories de foncteurs, telle
\[
\uCh_{\Sigma,\tau}(X)\simeq\uHom_{\Sigma_\tau,\Sigma}(I_\tau,X),
\]
où $\Sigma_\tau$ désigne l'ensemble des flèches dans $I_\tau$ qui sont soit inversibles, soit \textit{rétrogrades.} (Cet ensemble de flèches est admissible, car stable par composition.) On désigne par $\uHom_{\Sigma_\tau,\Sigma}$ la sous-catégorie pleine de $\uHom$ formée des foncteurs qui appliquent $\Sigma_\tau$ dans $\Sigma$. Or de façon générale, si
\[
(I,\Sigma_I),\ \ (I',\Sigma_{I'}),\ \ (X,\Sigma)
\]
sont trois catégories, munies d'ensembles admissibles de flèches, considérant les foncteurs $\varphi:I\longrightarrow I'$ avec $\varphi(\Sigma_I)\subseteq\Sigma_{I'}$, alors le foncteur familier
\[
\uHom(I,I')\longrightarrow\uHom(\uHom(I',X),\uHom(I,X))
\]
(transposé du foncteur de composition
\[
\begin{matrix}
\uHom(I,I') & \times & \uHom(I',X) & \longrightarrow & \uHom(I,X) & \\
(\varphi & , & F') & \longmapsto & F'\circ\varphi & ),
\end{matrix}
\]
se généralise en un foncteur
\[
\uHom_{\Sigma_I,\Sigma_{I'}}(I,I')\longrightarrow\uHom(\uHom_{\Sigma_{I'},\Sigma}(I',X),\uHom_{\Sigma_I,\Sigma}(I,X)),
\tag*{$(*)$}
\]
transposé du foncteur de composition
\[
\uHom_{\Sigma_I,\Sigma_{I'}}(I,I')\times\uHom_{\Sigma_{I'},\Sigma}(I',X)\longrightarrow\uHom_{\Sigma_I,\Sigma}(I,X)
\]
induit par $(*)$.

\msp{page 96}

Comme toute
\[
\varphi:\tau\longrightarrow\tau'
\]
induit
\[
\varphi:I_\tau\longrightarrow I_{\tau'},
\]
transformant flèches rétrogrades en flèches identiques ou rétrogrades, i.e.
\[
\varphi(\Sigma_\tau)\subseteq\Sigma_{\tau'},
\]
\printedpage{62}
% head: A. Grothendieck : Dérivateurs, Ch. VII, Édition des Universités de Montpellier II et Paris VII
% foot: Version 22 décembre 2001
on trouve le premier des deux foncteurs $(*)$ annoncés page 94, qui est associé par transposition au foncteur $\gamma$
\[
\begin{tikzcd}[column sep=large]
\uHom(\tau,\tau')\times\ \underbrace{\uCh_{\Sigma,\tau'}(X)}_{\simeq\,\uHom_{\Sigma_{\tau'},\Sigma}(I_{\tau'},X)} \arrow[r, "\gamma"] \arrow[d, "\text{can.}"'] & \uCh_{\Sigma,\tau}(X) \arrow[d, equal]\\
\uHom_{\Sigma_\tau,\Sigma_{\tau'}}(I_\tau,I_{\tau'})\times\uHom_{\Sigma_{\tau'},\Sigma}(I_{\tau'},X) \arrow[r, dotted, "\gamma_0"] & \uHom_{\Sigma_\tau,\Sigma}(I_\tau,X)
\end{tikzcd}
\tag*{$(*)$}
\]
déduit du foncteur composition $\gamma_0$ dans ce diagramme : $\gamma=\gamma_0\circ\mathrm{can}.$ D'autre part, ce foncteur de composition $\gamma$ induit visiblement un foncteur sur les sous-catégories de chemins à source et but fixés
\[
\uHom(\tau,\tau')\times\uCh_{\Sigma,\tau'}(X;a,b)\longrightarrow\uCh_{\Sigma,\tau}(X;a,b),
\]
d'où la deuxième des deux flèches annoncées dans $(*)$ p. 94.

Cela montre donc que si
\[
\varphi_1,\varphi_2:\tau\rightrightarrows\tau'
\]
sont deux morphismes de $\tau$ dans $\tau'$ homotopes dans $\uHom(\tau,\tau')$, alors les foncteurs induits entre catégories de $\Sigma$-chemins ou entre catégories de $\Sigma$-chemins à extrémités précisés : % [sic] précisés -> précisées, agreement with extrémités

\msp{page 97}
\[
\begin{aligned}
\varphi^*_{1\,X},\varphi^*_{2\,X} : && \uCh_{\Sigma,\tau'}(X) &\quad\rightrightarrows\quad\uCh_{\Sigma,\tau}(X),\\
(\varphi_1)^*_{a,b},(\varphi_2)^*_{a,b} : && \uCh_{\Sigma,\tau'}(X;a,b) &\quad\rightrightarrows\quad\uCh_{\Sigma,\tau}(X;a,b)
\end{aligned}
\]
sont également homotopes.

Or le lemme d'homotopie nous assure que l'hypothèse $\varphi_1\sim\varphi_2$ est \textit{toujours} satisfaite. Donc :

\textbf{Théorème d'homotopie sur les opérations universelles sur les chemins.} \textit{Soient $X$, $\Sigma$, $a$, $b$ comme dessus, et $\tau,\tau'\in\mathrm{Ty}$ telles que $\mathrm{Hom}(\tau,\tau')\neq\emptyset$. Alors les foncteurs $\varphi^*_X$, $\varphi^*_{X;a,b}$ induits par un $\varphi\in\mathrm{Hom}(\tau,\tau')$}
\[
\left\{
\begin{aligned}
\varphi^*_X : && \uCh_{\Sigma,\tau'}(X) &\quad\longrightarrow\quad\uCh_{\Sigma,\tau}(X)\\
\varphi^*_{X;a,b} : && \uCh_{\Sigma,\tau'}(X;a,b) &\quad\longrightarrow\quad\uCh_{\Sigma,\tau}(X;a,b)
\end{aligned}
\right.
\]
\textit{ne dépendent pas de $\varphi$,} à homotopie près.

On peut préciser cet énoncé, de telle façon que l'énoncé pour les $\varphi^*_{X;a,b}$ apparaisse, comme il devrait, comme corollaire pour celui sur $\varphi^*_X$. Pour ceci, on va regarder les catégories de chemins dans $X$ comme catégories \textit{au dessus de $X$,} par les foncteurs `source et but' :
\[
\begin{tikzcd}[column sep=small]
\uCh_{\Sigma,\tau}(X) \arrow[dr, "(\sigma_\tau{,}\beta_\tau)"'] & & \uCh_{\Sigma,\tau'}(X) \arrow[dl, "(\sigma_{\tau'}{,}\beta_{\tau'})"]\\
 & X\times X &
\end{tikzcd}
\]

\newpage
\printedpage{63}
% head: A. Grothendieck : Dérivateurs, Ch. VII, Édition des Universités de Montpellier II et Paris VII
% foot: Version 22 décembre 2001

\msp{page 98}

Alors le foncteur $(*)$ (p. 94) se factorise par la sous-catégorie $\uHom_{\mathrm{Cat}/X\times X}(\cdot,\cdot)$ de $\uHom_{\mathrm{Cat}}(\cdot,\cdot)$ formée des foncteurs qui commutent aux foncteurs de projection dans $X\times X$, et les homomorphismes entre tels foncteurs qui s'effectuent par des flèches `verticales' de la catégorie but (i.e. par des flèches au dessus de flèches identiques de la catégorie base $X\times X$) :
\[
\uHom(\tau,\tau')\longrightarrow\uHom_{\mathrm{Cat}/X\times X}(\uCh_{\Sigma,\tau'}(X),\uCh_{\Sigma,\tau}(X)).
\]

Donc dans l'énoncé du théorème d'homotopie, on peut préciser que l'homotopie se fait dans cette sous-catégorie, i.e. s'effectue par une suite d'homotopies simples qui sont des morphismes de foncteurs entre catégories au dessus de $X\times X$. Cela implique alors l'énoncé d'homotopie sur les $\varphi^*_{X;a,b}$.

% ---- from Der-ChVII-456_p066-070.tex
\leavevmode\printedpage{64}
% head: A. Grothendieck : Dérivateurs, Ch. VII, Édition des Universités de Montpellier II et Paris VII
\msp{page 99}

\section*{10\quad Question de comparaison pour les catégories \textmd{Cat}-fibrées satisfaisant au ‘théorème d'homotopie' (base filtrante)}

Nous allons généraliser maintenant au cas d'une catégorie (Cat)-fibrée sur une autre
\[
\begin{tikzcd}
C \arrow[d, "\pi"'] \\
T,
\end{tikzcd}
\]
en pensant surtout au cas $T = \uTy$, $C = \uCh_{\Sigma}(X;a,b)$. L'hypothèse essentielle est la suivante :

\noindent\begin{tabular}{@{}p{0.3\textwidth}p{0.64\textwidth}@{}}
(Hypothèse d'homotopie) & Si $t,t' \in \Ob T$ sont tels que $\Hom(t,t') \neq \emptyset$, alors les foncteurs $\varphi^{*} : C_{t'} \longrightarrow C_{t}$ associés aux $\varphi \in \Hom(t,t')$ sont homotopes.
\end{tabular}

Ils définissent donc une seule et même flèche dans Cat\,hot. Je désigne par $T_{\mathrm{ord}}$ ou $\To$ l'ensemble $\Ob(T)$, avec comme relation d'ordre celle déduite de la structure de catégorie :
\[
t \leq t' \qquad\Longleftrightarrow\qquad \Hom_{T}(t,t') \neq \emptyset,
\]
de sorte qu'on a un foncteur
\[
T \longrightarrow \To
\]
qui est \textit{universel} pour les foncteurs de $T$ dans des catégories ‘ordonnées' (i.e. où toute flèche est monomorphique \texttt{[non, cf. }$\bullet\rightrightarrows\bullet$\texttt{; plutôt, où toujours }$\#\Hom(a,b) \leq 1$ \texttt{(admettant des isomorphismes non identiques)]}) ou pour les foncteurs $F$ sur $T$ tels que $F$ soit constant sur tout $\Hom(t,t')$. L'hypothèse d'homotopie peut s'exprimer en disant que le foncteur composé
\[
\begin{tikzcd}
T^{o} \arrow[r] & \Cat \arrow[r] & \Cathot \\
t \arrow[r, mapsto] & C_{t} &
\end{tikzcd}
\]
\msp{page 100}
se factorise (évidemment de façon unique) en un foncteur
\[
H_{C} : \To^{o} \longrightarrow \Cathot,
\]
i.e. qu'on a un diagramme commutatif
\printedpage{65}
% head: A. Grothendieck : Dérivateurs, Ch. VII, Édition des Universités de Montpellier II et Paris VII
\[
\begin{tikzcd}
T^{o} \arrow[r, "t\,\longmapsto\, C_{t}"] \arrow[d] & \Cat \arrow[d] \\
\To^{o} \arrow[r, "H_{C}"] & \Cathot.
\end{tikzcd}
\]

\dgrblock{dgrb333}{The preorder remembers existence, not multiplicity}
\begin{dcomment}{exposition}{p99}{The preorder remembers existence, not multiplicity}
Grothendieck's purpose is to retain the common value of all operations between two fibres while discarding their multiplicity. Let $q:T\to P$ be the identity on objects, with
\[
 P(t,t')=\begin{cases}\{*\},&T(t,t')\ne\varnothing,\\
 \varnothing,&T(t,t')=\varnothing.\end{cases}
\]
The category $P$ is a preorder; distinct objects may be mutually comparable. Its universal property is among thin categories, whose hom-sets have at most one element. Having every arrow monic is a different condition: the category with two parallel arrows already separates the two notions, as the editors' correction observes.

Let $F:T^{\mathrm{op}}\to\Cat$ describe the fibration. In this discussion
\[
 \Cathot(A,B)=\pi_0\uHom(A,B)
\]
means functors modulo finite zigzags of natural transformations. It is not the localization $\Hot$ of categories at weak equivalences of nerves. Natural transformations become equalities in $\Cathot$, and hence also in $\Hot$, but not every weak equivalence of nerves has a functorial homotopy inverse. The distinction matters throughout the comparison problem.

The homotopy hypothesis says exactly that the composite of $F$ with $\Cat\to\Cathot$ factors through $P^{\mathrm{op}}$. If $P$ is nonempty and any two objects have a common lower bound, $P^{\mathrm{op}}$ is filtered and the resulting diagram defines an ind-object $H_C$. Here nonemptiness is part of the usual convention for ``filtrant''. The displayed pairwise condition in the manuscript must be read with that convention; otherwise the empty base would satisfy it vacuously and would have no cofinal sequence of the proposed form.

One must keep this modest hypothesis distinct from $T$ being cofiltered. A cofiltered category must also admit a common precomposition equalizing any pair of parallel arrows. Nothing of that kind is assumed here. The usual definition and its role in finite-limit arguments are recorded in \citet[Categories, Definition~4.19.1 and Lemma~4.19.2]{StacksProject}. For $\Ty$, concatenation supplies common lower bounds because either factor can be collapsed to the empty type. It does not identify the different ways to perform a collapse. The homotopy theorem, rather than strict equalization, is what allows those different ways to be forgotten in $H_C$.
\end{dcomment}


Nous allons supposer de plus :

\noindent\begin{tabular}{@{}p{0.36\textwidth}p{0.58\textwidth}@{}}
(Hypothèse de filtration) & L'ensemble ordonné $\To$ associé à $T$ est \textit{filtrant} décroissant (filtrant à gauche), i.e. si $t,t' \in \Ob T$, il existe $t'' \in \Ob T$ et des flèches $t'' \longrightarrow t$, $t'' \longrightarrow t'$.
\end{tabular}

Cette hypothèse est bien satisfaite dans $\uTy$, car pour $\tau,\tau' \in \Ob \uTy = \mathrm{Ty}$, il suffit de prendre $\tau\tau'$, qui s'envoie dans $\tau$ par $\mathrm{id}_{\tau} * p_{\tau'}$, dans $\tau'$ par $p_{\tau} * \mathrm{id}_{\tau'}$, où $p_{\tau} : \tau \longrightarrow 1$ est le morphisme de $\tau$ dans l'objet final.

Ainsi, $H_{C}$ définit un ind-objet de Cat\,hot, que nous allons désigner par $\underline{\underline{C}}$. Si on suppose de plus :

\noindent\begin{tabular}{@{}p{0.36\textwidth}p{0.58\textwidth}@{}}
(Hypothèse de dénombrabilité, toujours satisfaite si $\mathrm{card}\,\To \leq \aleph_{0}$) & $\To$ admet un sous-ensemble \textit{dénombrable} cofinal à gauche.
\end{tabular}

alors cet ind-objet de Cat\,hot peut être défini, de beaucoup de façons, en termes
\msp{page 101}
d'un système inductif dans Cat (pas seulement dans Cat\,hot)
\[
(C_{n})_{n\in\mathbf{N}},
\]
d'où on récupère l'ind-objet $\underline{\underline{C}}$ dans Cat\,hot en passant au système inductif correspondant dans Cat\,hot, puis à l'ind-objet associé.

Pour construire le système inductif $C_{n}$, il suffit de prendre une chaîne infinie
\[
t_{0} \longleftarrow t_{1} \longleftarrow t_{2} \longleftarrow \;\cdots\; \longleftarrow t_{n} \longleftarrow t_{n+1} \longleftarrow \;\cdots
\]
de morphismes dans $T$, telle que l'ensemble des $t_{i}$ soit cofinal à gauche dans $\To$. On va supposer de plus $C$ \textit{strictement fibrée} sur $T$ (i.e. provenant d'un véritable foncteur $T^{o} \longrightarrow \Cat$, i.e. on peut choisir les foncteurs images inverses
\[
\varphi^{*} : C_{t'} \longrightarrow C_{t}
\]
de façon transitive stricte, pas seulement transitive à isomorphisme près) - cas auquel on peut toujours se ramener d'ailleurs par l'astuce de \textsc{Giraud}. Cette hypothèse est d'ailleurs vérifiée dans le cas $\uCh_{\Sigma}(X;a,b)$ sur $\uTy$. Ainsi, les
\[
(C_{t_{i}})_{i\in\mathbf{N}}
\]
forment un système inductif de catégories indexées par $\mathbf{N}$, le morphisme de transition
\[
C_{n} \longrightarrow C_{n'} \qquad \text{i.e. } C_{t_{n}} \longrightarrow C_{t_{n'}}
\]
\printedpage{66}
% head: A. Grothendieck : Dérivateurs, Ch. VII, Édition des Universités de Montpellier II et Paris VII
\msp{page 102}
étant le foncteur image inverse $\varphi^{*}_{n,n'}$, où
\[
\varphi_{n,n'} : t_{n'} \longrightarrow t_{n}
\]
est le composé des $t_{n'} \longrightarrow t_{n'-1} \longrightarrow \;\cdots\; \longrightarrow t_{n}$. (Si on n'avait pas transitivité stricte, on définirait $\varphi^{*}_{n,n'}$ comme le composé des foncteurs image inverse par des flèches de transition successives $t_{n'} \longrightarrow t_{n'-1} \longrightarrow \;\cdots\; \longrightarrow t_{n+1} \longrightarrow t_{n}$, ce ne serait $\varphi^{*}_{n,n'}$ qu'à isomorphisme canonique près.)

\dgrblock{dgrb334}{Making a sequence strict}
\begin{dcomment}{context}{p101}{Making a sequence strict}
Grothendieck wants an actual category as a sequential colimit, so he temporarily makes the reindexing functors strictly composable. This is a familiar issue from the theory of fibred categories: a cleavage supplies canonical comparison isomorphisms, whereas a splitting makes those comparisons identities. The distinction is already developed in \citet[Expos\'e~VI, \S\S8--9]{SGA1}. The replacement invoked here is an equivalence of fibred categories; it is not a claim that a chosen cleavage on the original total category can always be made split without changing that category.

There is a concrete construction behind the reference to Giraud. For $t\in T$, let $\widetilde F(t)$ be the category of cartesian functors $T/t\to C$ over $T$, with vertical natural transformations. Precomposition along $T/s\to T/t$, induced by $s\to t$, is strictly associative. Evaluation at the terminal object $\id_t$ gives an equivalence $\widetilde F(t)\simeq C_t$: a cartesian functor is reconstructed from its value there by cartesian lifts, and any morphism between those values extends uniquely to a vertical transformation. Thus $\widetilde F$ is a strict replacement of the original pseudofunctor. The right splitting and its equivalence property are treated explicitly in \citet[Theorem~5.4 and Remark~5.5]{EmmeneggerMesitiRosoliniStreicher2024}; this gives a precise modern reference for the ``astuce de Giraud''.

Countability has a different role. Enumerate a cofinal family $s_0,s_1,\ldots$ in $P$. Choose $t_0$ below $s_0$, then choose $t_{n+1}$ below both $t_n$ and $s_{n+1}$, together with an arrow $t_{n+1}\to t_n$. The resulting chain is cofinal in the preorder. Pullback produces
\[
 C_{t_0}\longrightarrow C_{t_1}\longrightarrow C_{t_2}\longrightarrow\cdots.
\]
Even without replacing the entire fibration, this particular sequence can be made strict simply by defining its long transition maps as composites of its adjacent ones. There are no competing paths in the freely generated chain to compare. The arrow $C_{t_n}\to C_{t_{n-1}}$ printed on manuscript page~102 runs in the wrong direction; the intended transition is $C_{t_{n-1}}\to C_{t_n}$.
\end{dcomment}


Soit
\[
C_{\infty} = \varinjlim_{n} C_{n} = \varinjlim_{n} C_{t_{n}}
\]
la catégorie $\varinjlim$ . (Elle sera surtout sympathique si on a pu choisir le système $(t_{n})_{n\in\mathbf{N}}$ de telle façon que les foncteurs de transition
\[
C_{t_{n}} \longrightarrow C_{t_{n-1}}
\]
soient pleinement fidèles, et peut-être même \texttt{[des]} monomorphismes, de telle façon que les $C_{n}$ s'identifient à des sous-catégories pleines de leur limite inductive $C_{\infty}$.) Je voudrais comprendre les relations précises entre les trois genres d'objets suivants :

\begin{itemize}[leftmargin=2.5em]
\item[\circled{a}] $C$ elle-même (comme objet dans Cat, pas de $\Cat/T$), et surtout les objets de Cat\,hot et de Hot qui lui correspondent.

\item[\circled{b}] \msp{page 103}$C_{\infty}$, et les objets de Cat\,hot et de Hot qui lui correspondent. Dans quelle mesure ces deniers % [sic] deniers -> derniers, missing letter
sont-ils indépendant % [sic] indépendant -> indépendants, agreement with the plural subject
du choix du système projectif cofinal à gauche $(t_{n})_{n\in\mathbf{N}}$ dans $T$? Quel est le rôle de $C_{\infty}$? Est-elle canoniquement isomorphe à $C$ dans Cat\,hot, ou dans Hot? (Auquel cas la question de son ‘rôle' me semblerait résolue!)

\item[\circled{c}] Le ind-objet % [sic] Le ind-objet -> L'ind-objet, missing elision
$H_{C}$ dans Cat\,hot, et l'ind-objet de Hot qu'il définit.

Je note que l'ind-objet $H_{C}$, et aussi les catégories $C_{\infty}$ construites comme dessus, ne changent pas quand on remplace $T$ par une sous-catégorie $T'$ cofinale à gauche, ou aussi par une catégorie de la forme $T/t_{0}$ ($t_{0} \in \Ob T$), qui donne lieu à l'ensemble ordonnée % [sic] ordonnée -> ordonné, agreement with "ensemble"
associé
\[
(T/t)_{\mathrm{o}} = \To/t = \{t \in \To \mid t \leq t_{0}\}
\]
\texttt{[plutôt} $(T/t_{0})_{\mathrm{o}} = \To/t_{0} = \{t \in \To \mid t \leq t_{0}\}$\texttt{]}. Si la relation entre les objets dans a), et ceux dans b) et dans c), est très étroite, on en concluerait % [sic] concluerait -> conclurait, spelling of the conditional
que le ‘type' de $C$ dans Cat\,hot, ou dans Hot suivant le cas, ne change pas en remplaçant $C$ par
\printedpage{67}
% head: A. Grothendieck : Dérivateurs, Ch. VII, Édition des Universités de Montpellier II et Paris VII
\msp{page 104}
$C \times_{T} T'$ ou par $C/t_{0} \simeq C \times_{T} T/t_{0}$. Bien mieux, elle ne changerait pas en remplaçant $T$ par $(\mathbf{N})^{o}$, la catégorie (ordonnée) opposée à $\mathbf{N}$ (ayant donc un unique objet final), et la catégorie fibrée $C$ sur $T$ par la catégorie $C_{\underline{t}}$ induite par $C$ sur $(\mathbf{N})^{o}$, via le foncteur
\[
\underline{t} : (\mathbf{N})^{o} \longrightarrow T
\]
défini par le système projectif $\underline{t} = (t_{n})_{n\in\mathbf{N}}$ dans $T$. Cela signifierait donc que pour comprendre le type homotopique (au sens strict, i.e. dans Cat\,hot, ou large, i.e. dans Hot) de $C$, on peut le ‘nettoyer' de façon draconienne, la base $T$ étant vue comme une sorte de gangue un peu vaseuse, ou du moins pléthorique.

\dgrblock{dgrb335}{A slice does not have the same preorder}
\begin{dcomment}{development}{p103}{A slice does not have the same preorder}
Grothendieck proposes restricting the construction to types below a chosen type. This is sound for a lower ideal of the preorder, but a slice category retains the actual arrow to the chosen object. The equality of preorders printed here therefore needs a correction that affects the proposed argument.

\textbf{A slice obstruction.}\ours\quad There is a finite category $T$ whose associated preorder is cofiltered, while the associated preorder of $T/t_0$ is not cofiltered.

\emph{Proof.} Take two objects $s,t_0$, their identities, and exactly two further arrows $f,g:s\rightrightarrows t_0$. The preorder $P$ has $s<t_0$, so $s$ is a lower bound for every pair. The slice $T/t_0$ has three objects
\[
 (s,f),\qquad (s,g),\qquad(t_0,\id_{t_0}),
\]
with one arrow from each of the first two objects to the last. An arrow between the first two would have to be $\id_s$ and would force $f=g$. There is no fourth object that could map to both. Its preorder is consequently not cofiltered, although the lower ideal $P/t_0=P$ is. This proves the assertion.\hfill$\square$

The example also prevents interpreting the displayed equality as an equivalence after removing duplicate isomorphic objects: the two incomparable objects remain distinct in the skeletal preorder. It does not dispute the connectedness of the slice, which has a terminal object, nor the separate construction of $H_C$ over $P^{\mathrm{op}}$.

There are two precise repairs. If only the ind-object is wanted, use the full subcategory $T_{\le t_0}\subseteq T$ on objects admitting some arrow to $t_0$. Its preorder is exactly the lower ideal of $P$. Its inclusion is cofinal on associated preorders: given $t$, choose a common lower bound of $t$ and $t_0$. It is cofiltered as a preorder for the same reason. Restriction therefore leaves $H_C$ unchanged. The actual slice $T/t_0$ can instead be used when $T$ is cofiltered in the categorical sense. Given $u:a\to t_0$ and $v:b\to t_0$, first choose $r\to a,b$, then precompose once more to equalize the resulting two arrows $r\rightrightarrows t_0$. This supplies a common predecessor in the slice. Parallel arrows in the slice can likewise be equalized in $T$, and the resulting equality still lies over $t_0$. Thus the slice is then cofiltered.

We checked the standard filtering results in \citet[Categories, \S4.19]{StacksProject} and the homotopy-colimit comparison discussed below; we did not find this particular correction to manuscript page~103 stated there. The example and the two repairs are recorded as a local development of the argument, without a claim of priority.
\end{dcomment}


\item[\circled{d}] De toutes façons, si $T'$ est une autre catégorie filtrante à gauche, et
\[
\alpha : T' \longrightarrow T
\]
un foncteur cofinal à gauche, de sorte que $C$ sur $T$, et $C' = C \times_{T} T'$ sur $T'$, définissent essentiellement le ‘même' objet de Ind\,Cat\,hot, et que la catégorie $C'_{\infty}$ pour $C'$ sur $T'$ fait partie des catégories $C_{\infty}$ pour $C$ sur $T$ - on peut se demander
\msp{page 105}
si le foncteur canonique
\[
C' \;\longrightarrow\; C
\]
est également toujours un isomorphisme dans Cat\,hot, ou sinon, du moins dans Hot.
\end{itemize}

Un intérêt particulier s'attache à $\pi_{0}(C)$, et de plus à sa variance pour $C$ sur $T$ variable de la façon suivante. Si $S$ est dans Cat, on regarde les catégories
\[
\uHom_{\Cat}(S,C_{t}) \qquad (t \in T),
\]
on trouve un foncteur
\[
t \longmapsto \uHom_{\Cat}(S,C_{t}), \qquad T^{o} \longrightarrow (\Cat),
\]
d'où une catégorie fibrée sur $T$, que je vais désigner par
\[
\sHom_{T}(S_{T},C)
\]
(Sauf erreur, elle s'interprète comme la catégorie fibrée sur $T$ correspondante au champ sur $\mathrm{Top}(T) = T^{\wedge}$ du morphisme de $T$-champs de $S_{T}$ dans $C$, où $S_{T} = S \times T$ (considérée comme catégorie fibrée sur $T$ à fibre essentiellement constante $S$) dans $C$. Ici j'identifie catégorie fibrée sur $T$, et champs sur le topos $T^{\wedge}$.) Il est immédiat que la propriété d'homotopie (p. 99) est stable par passage de $C$ à $\sHom_{T}(S_{T},C)$. On s'intéresse à la variation de
\[
\pi_{0}(\sHom_{T}(S_{T},C))
\]
\printedpage{68}
% head: A. Grothendieck : Dérivateurs, Ch. VII, Édition des Universités de Montpellier II et Paris VII
\msp{page 106}
comme contrafoncteur en $S$,
\[
H : S \longmapsto \pi_{0}\sHom_{T}(S_{T},C) \qquad (\Cat)^{o} \longrightarrow (\mathrm{Ens}),
\]
et à une éventuelle expression de ce contrafoncteur $H$ en termes de l'ind-objet $H_{C}$ de Cat\,hot, par la formule hypothétique
\[
H(S) \simeq \Hom_{\mathrm{Ind}(\Cathot)}(\gamma(S),H_{C}),
\]
où
\[
\gamma : \Cat \longrightarrow \Cathot
\]
est le foncteur canonique. Si on avait un tel isomorphisme canonique fonctoriel en $S$, cela donnerait une description (à isomorphisme unique près) de l'ind-objet $H_{C}$ en termes du foncteur précédent $H$ sur Cat.

Il me vient un doute ici, en vue de l'exemple $T = \uTy$, $C = \uCh_{\Sigma}(X;a,b)$. Les trois types d'objets a) b) c) ont le même $\pi_{0}$, et il est sûr que ‘c'est le bon'. D'autre part, pour une catégorie paramètre variable $S$, il est sûr aussi que le bon foncteur en $S$ est
\[
S \longmapsto \pi_{0}(\uHom_{T}(S_{T},C)) \simeq \Hom_{\mathrm{Ind}\,\Cathot}(S,H_{C}),
\]
isomorphe aussi, \textit{si $S$ est finie,} à
\[
\Hom_{\Cathot}(S,C_{\infty}),
\]
ce qui montre que, du moins pour des arguments $S \in (\Cat)$ qui sont des catégories \textit{finies}, les objets
\msp{page 107}
$C_{\infty}$ et $H_{C}$ sont ‘bons', il nous permettent % [sic] il nous permettent -> ils nous permettent, agreement with the plural subject
de décrire un foncteur important qui domine la situation envisagée. Moyennant une hypothèse très forte, mais satisfaite dans les cas importants, on peut dans ces deux formules remplacer Cat\,hot par Hot.

\dgrblock{dgrb336}{The functor represented by the ind-object}
\begin{dcomment}{exposition}{p106}{The functor represented by the ind-object}
Grothendieck's parameter category $S$ restores information that the connected components of a single fibre cannot carry. The proposed formula for $H(S)$ is valid under the hypotheses already imposed, and its proof explains both the ind-object and the restriction to finite $S$ in the next formula.

For any strict diagram $A:T^{\mathrm{op}}\to\Cat$, the connected components of its fibred total category are
\[
 \pi_0\left(\int_T A\right)
 \cong\mathop{\mathrm{colim}}_{T^{\mathrm{op}}}\pi_0 A(t).
\]
Indeed, the total category is generated by vertical arrows and cartesian transport arrows. Vertical arrows impose exactly the component relation inside each $A(t)$; a cartesian arrow imposes exactly the relation identifying a class with its restriction along an arrow of $T$. Those are the generators and relations defining the colimit of sets, with no further relation. Apply this observation to $A(t)=\uHom(S,C_t)$. A zigzag between two reindexing functors induces a zigzag after postcomposition with functors from $S$, so the diagram of component sets factors through $P^{\mathrm{op}}$. We obtain, naturally in $S$,
\[
 \begin{split}
 H(S)&=\pi_0\sHom_T(S_T,C)\\
 &\cong\mathop{\mathrm{colim}}_{P^{\mathrm{op}}}
       \Cathot(S,C_t)
 =\Hom_{\operatorname{Ind}(\Cathot)}(S,H_C).
 \end{split}
\]
The last equality is the defining formula for maps from a constant object into an ind-object. The usual fully faithful embedding of ind-objects into set-valued presheaves then gives the asserted uniqueness of $H_C$. Setting $S=e$ proves the equality of the three component sets in the text.

To compare with $C_\infty$, suppose $S$ has finitely many objects and arrows. A functor $S\to\mathop{\mathrm{colim}}_n C_n$ involves finitely many objects, arrows, and composition equations. Represent its objects and arrows at a common stage, then move to a later stage where all the finitely many required equations hold. It therefore descends to some $C_n$. A natural transformation between two such functors also consists of finitely many arrows and naturality equations and descends in the same way. Finally, if two descended functors lie in the same component, a witnessing homotopy is a \emph{finite} zigzag of such transformations; the whole zigzag descends to one later stage. This proves the bijection
\[
 \mathop{\mathrm{colim}}_n\Cathot(S,C_n)
 \xrightarrow{\ \sim\ }\Cathot(S,C_\infty).
\]
It is an instance of the finite-presentation argument underlying filtered colimits, not a fibrancy assertion.

The finiteness cannot simply be omitted. Let $C_n$ be the discrete category on $\{0,\ldots,n\}$ and let $S$ be the discrete category on $\mathbf N$. Then $C_\infty$ is discrete on $\mathbf N$. The left side consists of functions $\mathbf N\to\mathbf N$ with bounded image; the right side consists of all functions. The identity function lies only in the latter. The ind-object retains the information that a whole $S$-family must occur at one stage, while the ordinary colimit allows the stage to vary with the individual element of $S$.
\end{dcomment}


Par contre, sauf pour le cas où $S$ est réduit à un point et où $\Hom_{\Cathot}$ ou $\Hom_{\Hot}$ de $S$ dans $C$ n'est autre que $\pi_{0}(C)$ ($\simeq \pi_{0}(C_{\infty}) \simeq \pi_{0}(H_{C})$), donc effectivement ‘bon', je ne sais pas si
\[
\underbrace{\Hom_{\Cathot}(S,C)}_{=\,\pi_{0}\uHom(S,C)} \qquad \text{ou} \qquad \Hom_{\Hot}(S,C)
\]
est lié de façon raisonnable aux ensembles précédents (déterminés par $H_{C}$, resp. par $C_{\infty}$), p. ex. s'ils lui sont isomorphes de façon naturelle, et s'ils ont une signification géométrico-topologique raisonnable. Je commence à suspecter que \textit{non} - donc qu'on aurait ici un cas où une relation de type très général concernant un $\pi_{0}$, ici donc
\[
\pi_{0}(C) \simeq \pi_{0}(C_{\infty}) \simeq \pi_{0}(H_{C}),
\]
ne se prolonge pas en un isomorphisme sur les $\pi_{i}$ supérieurs, et en une propriété de Hot-équivalence.

% ---- from Der-ChVII-456_p071-075.tex
\printedpage{69}
% head: A. Grothendieck : Dérivateurs, Ch. VII, Édition des Universités de Montpellier II et Paris VII

\msp{page 108}
Je ne trouve d'ailleurs aucune flèche directe entre $C$ et l'un des objets $C_{\infty}$, $H_{C}$; tout ce que je vois dans ce sens, c'est
\[
\begin{tikzcd}[column sep=large, row sep=huge]
\CNo \arrow[r, "{\substack{\text{foncteur de}\\ \text{localisation}}}"] \arrow[d] & C_{\infty} \\
C, &
\end{tikzcd}
\]
où la flèche verticale est la flèche de projection de $\CNo = C \times_{T} \mathbf{N}^{o}$ (avec $\mathbf{N}^{o} \xrightarrow{\ t\ } T$). Ceci n'établit une relation topologique étroite entre $C$ (en tant qu'objet de Cat\,hot ou de Hot simplement) et $C_{\infty}$, que si les deux flèches précédentes sont tout au moins des Hot-équivalences. Et le cas-type (problématique) où la question de la relation entre $C$ et $C_{\infty}$ se pose, est donc celui où $T = \mathbf{N}^{o}$, i.e. celui d'un système inductif quelconque de (petites) catégories, indexé par $\mathbf{N}$, et interprété comme défini par une catégorie fibrée $C = \CNo$ sur $\mathbf{N}^{o}$. Dans ce cas particulier, on a du moins une flèche canonique
\[
C \longrightarrow C_{\infty}
\]
sur laquelle on peut se poser la question si c'est
\msp{page 109}
une Hot-équivalence. Sinon, cela prouvera que dans le cas général, le Hot-type de $C$ (et a fortiori sa classe d'isomorphie dans Cat\,hot) n'a pas de sens géométrique intéressant. C'est à dire que $C$ est intéressante uniquement \textit{en tant que catégorie fibrée sur $T$}, permettant notamment de construire l'ind-objet $H_{T}$ de Cat\,hot ou de Hot, mais pas en tant que `catégorie nue'.

\dgrblock{dgrb337}{A finite category need not be a finite homotopy probe}
\begin{dcomment}{examples}{p107}{A finite category need not be a finite homotopy probe}
Grothendieck explicitly reserves the passage from $\Cathot$ to $\Hot$ for an additional hypothesis. The reservation is necessary. The finite-presentation proof above counts objects, arrows, and equations in a category. Compactness in the homotopy category of spaces concerns the whole nerve, and a finite category can have nondegenerate simplices in arbitrarily high degrees.

Let $S=\mathbf B(\mathbf Z/2)$, which has one object and two arrows, and put $K=NS$, $K_n=\operatorname{sk}_n K$. Let $D_n=\Delta/K_n$ be the category of simplices of $K_n$. Every simplex of $K$ belongs to one skeleton, so
\[
 D_\infty=\mathop{\mathrm{colim}}_nD_n=\Delta/K.
\]
The simplex-category resolution gives $|ND_n|\simeq|K_n|$ and $|ND_\infty|\simeq|K|$, naturally for the inclusions. The precise last-vertex statement for an arbitrary simplicial set is \citet[Corollary~11.9.11.5]{Kerodon2026Simplex}. Thus the filtered colimit has the homotopy type of $B(\mathbf Z/2)$, while each stage has an $n$-dimensional realization.

The identity of $B(\mathbf Z/2)$ cannot factor in $\Hot$ through any $|K_n|$. Such a factorization would make its identity map on cohomology factor through a group vanishing in degrees greater than $n$. But
\[
 H^*(B(\mathbf Z/2);\mathbf F_2)\cong\mathbf F_2[u],\qquad |u|=1.
\]
For completeness, the normalized group cochains with trivial $\mathbf F_2$ coefficients have one basis element in each degree: a normalized cochain is determined by its value on the string consisting entirely of the nonidentity group element. Its differential is zero, since the two end terms cancel and every interior product gives the identity, on which a normalized cochain vanishes. The degree-one generator and its cup powers take value one on these strings. This proves the displayed calculation and the obstruction.

Consequently the canonical map
\[
 \mathop{\mathrm{colim}}_n\Hom_{\Hot}(S,D_n)
 \longrightarrow\Hom_{\Hot}(S,D_\infty)
\]
is not surjective, although $S$ is finite as a category. There is no contradiction with the preceding $\Cathot$ formula. Finite CW homotopy types are compact probes for filtered homotopy colimits, but the nerve of this finite group is not one. This is the point at which the two meanings of ``finite'' must be separated when reading the notation $\Hotz$ later on these pages.
\end{dcomment}


\dgrblock{dgrb338}{The sequential total category is a homotopy colimit}
\begin{dcomment}{later}{p108}{The sequential total category is a homotopy colimit}
Grothendieck has reduced one part of the comparison to a single sequence. That part has a uniform answer: the canonical functor from its fibred total category to its ordinary filtered colimit is a weak equivalence of nerves. No full-faithfulness assumption on the transition functors is needed.

Write $E=\int_{\mathbf N^{\mathrm{op}}} C_n$ and $D=\mathop{\mathrm{colim}}_n C_n$. With the variance of this manuscript, an arrow $(m,y)\to(n,x)$ over $m\to n$, where $m\ge n$, is an arrow $y\to F_{n,m}x$ in $C_m$. The canonical functor $q:E\to D$ sends it to the corresponding arrow in $D$. Cartesian transport arrows become identities. The localization description in the text is understood in the usual categorical sense: after inverting those arrows, $E$ becomes \emph{equivalent} to $D$. Literal equality of its object set with that of $D$ is neither needed nor generally true.

The homotopy-colimit theorem of \citet[Theorem~1.2]{Thomason1979}, in its dual form for a fibration, identifies $|NE|$ with $\operatorname{hocolim}_n|NC_n|$. The convention for a contravariant diagram and the comparison map are displayed explicitly in \citet[\S3, equation~(9)]{Cegarra2020}. Two elementary filtered-colimit facts complete the argument. First,
\[
 N D\cong\mathop{\mathrm{colim}}_n NC_n:
\]
an $r$-simplex is a finite string of composable arrows, and finitely many representatives and equalities can all be realized at one stage. Second, for a filtered diagram of simplicial sets $K_n$, the augmentation
\[
 \operatorname{hocolim}_n K_n\longrightarrow
 \mathop{\mathrm{colim}}_n K_n
\]
is a weak equivalence. One way to see the second fact is to use the simplicial replacement. In each fixed simplicial degree, the fibre over a simplex of the colimit is the category of its representatives at the different stages. Any two representatives eventually agree, and finitely many such agreements hold simultaneously at a later stage. This category is nonempty and filtered, hence its nerve is contractible: every finite diagram admits a cone, which contracts each finite simplicial sphere. Thus the augmentation is a weak equivalence in each row of the bisimplicial replacement, and diagonal realization preserves these rowwise weak equivalences. The composite is the map induced by $q$.

This identifies the genuine issue in the span on manuscript page~108. The horizontal arrow is always a weak equivalence. Whether the vertical arrow is one depends on how the chosen sequence sits inside $T$, including the homotopy carried by its parallel arrows. Passing to the preorder has deliberately forgotten that information.
\end{dcomment}


\dgrblock{dgrb339}{The total category and the telescope}
\begin{dcomment}{exposition}{p108:sur laquelle on peut se poser la question}{The total category and the telescope}
Grothendieck wants a path category that carries composition and at the same time has the intended homotopy type. The two objects in this opening comparison solve different parts of that problem. The total category remembers how the types of paths vary; the filtered colimit forgets the stages at which paths were represented. We write $BC=|NC|$, and call a functor a weak equivalence when its nerve is a weak homotopy equivalence. This is the meaning of an isomorphism in $\Hot$. It is stronger than a bijection on connected components and different from an isomorphism in $\mathrm{Cat\,hot}$, where morphisms are connected components of ordinary functor categories.

Here is a proof of the assertion announced on the following page. Let $F_0\to F_1\to\cdots$ be a strict sequence of small categories, let $E\to\mathbf N^o$ be its fibred Grothendieck construction, and let $L=\mathop{\rm colim}_nF_n$. The canonical functor $q:E\to L$ is a weak equivalence. For clarity, a morphism $(n,x)\to(m,y)$ in $E$, with $n\geq m$, is a morphism $x\to F_{m,n}y$ in $F_n$.

To prove the assertion, fix $z\in L$ and consider the comma category $q/z$. Any finite collection of its objects, arrows, and commuting triangles can be represented at one sufficiently large stage $r$: filtered colimits of categories are computed from finite diagrams and their eventual relations. We may also represent $z$ at stage $r$. Transporting the finitely many objects to that stage gives a diagram in $F_r/z_r$. Its projection to $q/z$ has a natural transformation to the original diagram, supplied by the cartesian arrows in $E$. It has a natural transformation to the constant diagram $(r,z_r,1)$, supplied by the arrows to $z_r$. Thus every finite simplicial diagram in $N(q/z)$ contracts after passing to a stage. The same argument with a chosen vertex proves connectedness and vanishing of all homotopy groups; the comma category is nonempty because $z$ has a representative. Quillen's Theorem A now gives the assertion \citep[\S1, Theorem A]{Quillen1973}. This argument does not require the transition functors to be faithful or injective on objects.

The second leg of the displayed roof is a separate matter. Take a nontrivial group $G$, let $T=BG$ be its one-object category, and take $C=T$ over itself. The preorder obtained by forgetting the distinction between parallel arrows is a point. A sequence constant on the identity arrow is therefore cofinal in that preorder, whereas the comparison is $*\to BG$. Its fundamental group is $1\to G$, so it is not a weak equivalence. This proves exactly the limitation noticed on manuscript page~109: cofinality in the preorder of objects does not account for the automorphisms of the indexing category. It does not predict failure for the particular category of interval types, which has additional structure.
\end{dcomment}


Mais je viens de me convaincre que $\CNo \longrightarrow C_{\infty}$ est bien une Hot-équivalence. Donc la flèche $\CNo \longrightarrow C$ définit une flèche \textit{dans} Hot
\[
C_{\infty} \longrightarrow C
\]
(en composant avec l'inverse dans Hot de la précédente). Est-ce un isomorphisme? Donc pour cette question, on est ramené à la question pour $\CNo \longrightarrow C$. Mais si on prend $C = T$, cela pose déjà la même question pour $\mathbf{N}^{o} \longrightarrow T$, donc la question si $T$ est asphérique, sous la seule hypothèse que $T_{\mathrm{o}}$ soit filtrante à gauche et dénombrable à l'infini (à gauche). Mais bien sûr c'est faux (p. ex. si $T$ n'a qu'un seul objet, il n'est pas pour autant asphérique !).

Cela me convainc donc qu'à défaut
% foot: Version 22 décembre 2001
\printedpage{70}
% head: A. Grothendieck : Dérivateurs, Ch. VII, Édition des Universités de Montpellier II et Paris VII
\msp{page 110}
d'hypothèse supplémentaire draconienne sur $T$, la catégorie $C$ `nue' n'offre pas d'intérêt particulier, à part son $\pi_{0}$ (qui est donné aussi par $C_{\infty}$, et par $H_{C}$). Il n'est pas dit, certes, que \underline{Ty} ne satisfasse pas à ces `conditions draconiennes'. Mais comme $C$ `nue' ne s'insère pas dans un formalisme sympathique (p.ex. on n'a pas d'interprétation géométrique de $\mathrm{Hom}_{\mathrm{Cat\,hot}}(S,C)$, ou de $\mathrm{Hom}_{\mathrm{Hot}}(S,C)$), je crois que par la suite je peux en faire abstraction. Exit a) !

\dgrblock{dgrb340}{What a one-object base can retain}
\begin{dcomment}{examples}{p109}{What a one-object base can retain}
Grothendieck's one-object example can be made completely explicit, and it shows exactly which information is absent from the ind-object. Let $G$ be a nontrivial discrete group and let $T=\mathbf BG$ be the category with one object and endomorphism group $G$. Take $C=T$ with its identity projection. This is a split fibration with terminal fibre and identity reindexing functors. Its homotopy hypothesis holds strictly; its associated preorder is the one-point category, which is cofiltered and has a countable cofinal family.

Choose the constant chain at that object, with identity transition arrows. Every $C_n$ and $C_\infty$ is the terminal category, and $H_C$ is the constant terminal ind-object. The total category remains $\mathbf BG$. Consequently the comparison in $\Hot$ is
\[
 *\longrightarrow BG,
\]
where $BG=|N\mathbf BG|$. Its source has trivial fundamental group and its target has fundamental group $G$. It is a bijection on components and is not a weak equivalence. More generally the constant fibration with fibre $D$ has total category $T\times D$, so it retains the product $BG\times|ND|$ while its sequential colimit is $D$.

There is also a direct answer to the mapping-set question of manuscript page~107. Set $G=\mathbf Z/2$ and $S=\mathbf BG$. Since the fibres of $C\to T$ are terminal, $H(S)$ is a singleton. The identity and the trivial functor $S\to C$ are not homotopic in $\Cathot$: natural transformations between functors of one-object groupoids express conjugacy of the corresponding homomorphisms, and these two homomorphisms are not conjugate. They remain distinct in $\Hot$, because their maps on $\pi_1$ differ. Thus neither of the proposed mapping sets into the bare total category recovers $H(S)$ in this generality.

The example concerns the abstract hypotheses of \S10. It does not decide the special path-category comparison over $\Ty$. That case has additional geometric structure, used later in the chapter. Nor does it make $C$ topologically meaningless: $C$ remembers the base monodromy, which is precisely the datum suppressed by the thin reflection. The warning is against identifying two different invariants.
\end{dcomment}


\begin{flushright}
{\footnotesize 5.11.90}
\end{flushright}

J'ai quand même passé la journée entière hier à essayer de comprendre la relation de $C$ et de son type dans Hot, à celui de $C_{\infty}$. En effet, comme on a finalement une flèche dans Hot
\[
C_{\infty} \longrightarrow C,
\]
déduite du diagramme dans (Hot)
\[
\begin{tikzcd}[column sep=huge, row sep=huge]
\CNo \arrow[r, "\sim"] \arrow[d] & C_{\infty} \\
C, &
\end{tikzcd}
\]
de prouver que $C_{\infty} \longrightarrow C$ est un isomorphisme dans Hot
\msp{page 111}
semble la façon la plus simple, et peut-être unique, de prouver que l'image de $C_{\infty}$ dans Hot ne dépend pas, à isomorphisme unique près, des choix faits pour construire $C_{\infty}$. Sans compter que dans le cas motivant, $C = \Ch_{\Sigma}(X)$ ou $\Ch_{\Sigma}(X;a,b)$, $T = \Ty$, c'est la catégorie des chemins $C$ qui donne lieu à une loi de composition des chemins impeccable, fonctorielle, associative et tout, alors qu'il n'en est rien pour $C_{\infty}$. On ne trouve pas une loi de composition
\[
C_{\infty} \times C_{\infty} \longrightarrow C_{\infty}
\]
dans Cat, ni même dans Cat\,hot ou dans Hot - ou plutôt, si on arrive pourtant à construire une flèche $C_{\infty} \times C_{\infty} \longrightarrow C_{\infty}$ dans Cat, c'est encore laborieusement à coups de choix $\pm$ arbitraires, cran par cran, et sans savoir si ce qu'on obtient a une propriété d'unicité, ne serait ce que dans Hot; on ne trouve une interprétation géométrique tangible de cette composition que dans $\Hotz$ (où on a une véritable loi de groupoïde), ou dans $H_{T}$ en tant qu'objet de Ind\,Cat\,hot
% foot: Version 22 décembre 2001
\printedpage{71}
% head: A. Grothendieck : Dérivateurs, Ch. VII, Édition des Universités de Montpellier II et Paris VII
\msp{page 112}
ou de Ind\,Hot. (Donc la limite inductive dans $\Hotz$ n'est autre que l'élément de $\Hotz$ défini par $C_{\infty}$.)

Mes efforts pour prouver que
\[
C_{\infty} \longrightarrow C
\]
était un isomorphisme dans Hot, ou ne serait-ce qu'un 1-isomorphisme (i.e. un isomorphisme non seulement pour les $\pi_{0}$, ce qui est toujours vrai, mais aussi pour les $\pi_{1}$ de toutes les composantes), ont échoués. J'ai trouvé qu'il suffisait que le foncteur % [sic] ont échoués -> ont échoué, accord fautif du participe
\[
T \longrightarrow T_{\mathrm{o}}
\]
soit ou bien lisse et coasphérique, ou bien propre et asphérique - donc tout est OK si $T = T_{\mathrm{o}}$, i.e. si $T$ est ordonnée. Mais dans le cas $T = \Ty$, le foncteur $T \longrightarrow T_{\mathrm{o}}$ n'est ni propre, ni lisse (pas même 0-propre ou 0-lisse).

Finalement j'en suis arrivé à la conviction que dans le cas qui nous occupe ($T = \Ty$, $C = \Ch_{\Sigma}(X;a,b)$), la flèche $C_{\infty} \longrightarrow C$ n'est en général \textit{pas} un isomorphisme dans Hot. En effet, dans le cas $a = b$, si cela \texttt{[était vrai]},
\msp{page 113}
on aurait une description du Hot-groupe $\Omega(X,a) = \Ch(X,a)$ (\footnote[23]{On fait ici $\Sigma = \mathrm{Fl}(X)$.}) en termes d'une catégorie $C$ munie d'une loi de composition \textit{associative} (et unitaire) \textit{dans} Cat, redonnant la loi de groupe de $\Omega$ en passant de (Cat) à Hot. Prenant le nerf de $C$, on trouverait un monoïde semi-simplicial qui, dans Hot, donnerait la loi de groupe de $\Omega(X,a)$. Or je me rappelle (\footnote[24]{C'est un erreur - ça doit être pour d'\textit{autres} objets groupe de Hot que les groupes de lacets. J'ai vu dans Quillen que $\Omega(X,x)$ doit se décrire par un groupe semi-simplicial (quand $X$ est un ensemble semi-simplicial, $x$ un sommet). Donc le problème si $C_{\infty} \simeq C$ dans Hot reste entier.}) qu'on m'a expliqué qu'en général, la loi de groupe de $\Omega(X,a)$ ne pourrait pas être obtenue par une loi \textit{associative} sur un ensemble semi-simplicial. Ainsi, l'insuccès de mes efforts semble tenir, non à ma maladresse, mais à la nature des choses. % [sic] C'est un erreur (note 24) -> C'est une erreur, genre

\dgrblock{dgrb341}{The cofinality condition that controls total categories}
\begin{dcomment}{exposition}{p112}{The cofinality condition that controls total categories}
Grothendieck is seeking a hypothesis that allows the base to be replaced without changing the total homotopy type. The appropriate general test is a condition on comma categories, stronger than cofinality of the associated preorders. It gives a precise sufficient criterion without requiring every reindexing functor between fibres to be a weak equivalence.

Let $\alpha:J\to T$ be a functor. If every comma category $\alpha/t$ has contractible nerve, then for every Grothendieck fibration $C\to T$ the projection
\[
 C\times_T J\longrightarrow C
\]
is a weak equivalence of nerves. To prove this, fix $x\in C_t$ and consider the comma category of the displayed projection over $x$. Its objects are triples consisting of $j\in J$, $y\in C_{\alpha j}$, and an arrow $y\to x$ over some $u:\alpha j\to t$. Forgetting $y$ gives a functor to $\alpha/t$. Cartesian lifting supplies a right adjoint: send $(j,u)$ to the cartesian arrow $u^*x\to x$. The adjunction follows from the unique factorization of an arrow through a cartesian lift. Hence these two comma categories have equivalent nerves. The dual of Quillen's Theorem~A now applies; see \citet[Corollary~4.1.3.3]{Lurie2009HTT}, applied to the opposite functor, for a precise modern formulation.

The condition is also necessary if the conclusion is required for \emph{every} fibration over $T$. Use the discrete fibration $T/t\to T$. Its total category has a terminal object, and its pullback is $\alpha/t$. Universal invariance of total homotopy types therefore forces $\alpha/t$ to be contractible. This elementary converse explains why a criterion expressed only in terms of which pairs of objects admit arrows cannot suffice.

For a cofiltered preorder $P$ and a decreasing cofinal sequence $t_n$, the category $\alpha/t$ is the nonempty tail of indices for which $t_n\le t$. It has a terminal object, the least such index. Thus the criterion proves the positive assertion in the text when $T=P$. For $T=\mathbf BG$ and the constant identity sequence, by contrast, $\alpha/t$ is a disjoint union of copies of $\mathbf N^{\mathrm{op}}$ indexed by $G$, so it is not contractible when $G\ne1$. These are the same two cases that the preceding examples separate. This comma criterion characterizes invariance for every fibration. The properness and smoothness conditions in the source instead exploit its hypothesis on parallel pullbacks, as the following proof shows; they need not imply the universal comma criterion.
\end{dcomment}


\dgrblock{dgrb342}{Why the smooth and proper criteria suffice}
\begin{dcomment}{later}{p112}{Why the smooth and proper criteria suffice}
The two sufficient conditions asserted here do follow from the earlier homotopy hypothesis. They are conditions on the base together with this particular class of coefficients; they need not make the chosen sequence cofinal for every fibration. The distinction matters because the hypothesis on parallel pullbacks supplies exactly the local constancy needed in the proof.

Write $q:T\to P=T_{\mathrm o}$, let $\alpha:\mathbf N^{\mathrm{op}}\to T$ be the chosen sequence, and replace the fibration by a strict diagram $F:T^{\mathrm{op}}\to\Cat$ as above. Put $K=NF$, regarded as a diagram of spaces, and $k=q^{\mathrm{op}}$. The fibre $E_p=q^{-1}(p)$ has the single object $p$ and endomorphism monoid $\operatorname{End}_T(p)$. Every endomorphism $u:p\to p$ is parallel to $\id_p$. Hence $F(u)$ is homotopic to the identity, so its nerve is a weak equivalence. Thus $K|_{E_p^{\mathrm{op}}}$ is a locally constant diagram: every one of its structure maps is a weak equivalence.

Use smoothness and properness for the localizer of weak equivalences of nerves. The fibre criterion of \citet[Theorem~1.2(b), printed p.~4]{Maltsiniotis2005Asphericite} says that smoothness makes
\[
 j_p:E_p\longrightarrow p\backslash q
\]
aspherical, meaning that every $j_p/z$ has contractible nerve. By the smooth--proper duality on pp.~4--5, properness makes
\[
 j_p:E_p\longrightarrow q/p
\]
coaspherical, meaning that every $z\backslash j_p$ has contractible nerve. In the first case coasphericity of $q$ says that $p\backslash q$ is contractible; in the second case asphericity of $q$ says that $q/p$ is contractible. Quillen's Theorem~A and its dual therefore show that $E_p$ is contractible in either case.

We shall use the following precise local-constancy fact. If $I$ has contractible nerve and $L:I\to\mathrm{sSet}$ sends every arrow to a weak equivalence, then for every $i\in I$ the canonical maps
\[
 L(i)\longrightarrow\operatorname{hocolim}_I L,
 \qquad
 \operatorname{holim}_I L\longrightarrow L(i)
\]
are weak equivalences. The first is \citet[Corollary~1.18, printed p.~599]{Cisinski2009LocallyConstant}; the second is its dual, applied to the opposite index category and opposite model category. The pointwise homotopy Kan-extension formula is Proposition~1.14 of the same paper, with the right-extension formula obtained by the same duality. These are statements about coherent diagrams, so no choice of mutually compatible homotopies between parallel pullbacks is required.

\emph{Smooth and coaspherical.} Let $G=\mathbf Lk_!K$ be the left homotopy Kan extension. Its value at $p$ is
\[
 G(p)\simeq
 \operatorname{hocolim}_{(p\backslash q)^{\mathrm{op}}}K
 \simeq\operatorname{hocolim}_{E_p^{\mathrm{op}}}K.
\]
Here the second equivalence follows because $j_p^{\mathrm{op}}$ is homotopy final: its under-comma categories are the opposites of the contractible categories $j_p/z$. The local-constancy fact shows that the adjunction unit $K\to k^*G$ is an objectwise weak equivalence. Let $v=(q\alpha)^{\mathrm{op}}:\mathbf N\to P^{\mathrm{op}}$. Cofinality of the chosen sequence in the preorder makes $v$ homotopy final: for $p\in P$, the category $p/v$ is the nonempty tail of indices for which $t_n\le p$, and it has a least element. Consequently
\[
 \operatorname{hocolim}_{\mathbf N}K\alpha^{\mathrm{op}}
 \xrightarrow{\ \sim\ }
 \operatorname{hocolim}_{\mathbf N}v^*G
 \xrightarrow{\ \sim\ }
 \operatorname{hocolim}_{P^{\mathrm{op}}}G
 \simeq\operatorname{hocolim}_{T^{\mathrm{op}}}K.
\]
The last equivalence is composition of left homotopy Kan extensions. By the unit and its adjunction identity, this composite is the canonical comparison of homotopy colimits.

\emph{Proper and aspherical.} Instead take $G=\mathbf Rk_*K$. The pointwise formula and the proper fibre criterion give
\[
 G(p)\simeq
 \operatorname{holim}_{(q/p)^{\mathrm{op}}}K
 \simeq\operatorname{holim}_{E_p^{\mathrm{op}}}K
 \xrightarrow{\ \sim\ }K(p).
\]
Now $j_p^{\mathrm{op}}$ is homotopy initial, which justifies restricting the homotopy limit, and the final arrow is the local-constancy equivalence. These maps are the components of the counit $k^*G\to K$, so this counit is an objectwise weak equivalence. Asphericity of $q$ also makes $k$ homotopy final, since $p/k\cong(q/p)^{\mathrm{op}}$. Both $k$ and $v=k\alpha^{\mathrm{op}}$ therefore give homotopy-colimit equivalences to $\operatorname{hocolim}_{P^{\mathrm{op}}}G$. It follows by two-out-of-three that
\[
 \operatorname{hocolim}_{\mathbf N}\alpha^{\mathrm{op}*}k^*G
 \longrightarrow\operatorname{hocolim}_{T^{\mathrm{op}}}k^*G
\]
is a weak equivalence. The counit transfers this conclusion to $K$.

Finally, the natural fibred form of Thomason's homotopy-colimit theorem identifies the comparison just proved with $C\times_T\mathbf N^{\mathrm{op}}\to C$. The sequential comparison already proved identifies its source with $C_\infty$ in $\Hot$. Thus each of the source's two conditions indeed makes its canonical $C_\infty\to C$ an isomorphism in $\Hot$.

This proof also identifies the exact limit of the statement. It used local constancy on $E_p^{\mathrm{op}}$, supplied by the homotopy hypothesis on $F$. For example, let $T$ be the one-object category of the monoid $\{1,0\}$ with absorbing zero. The natural transformation from the identity functor to the constant functor whose component is $0$ contracts its nerve. The map $q:T\to *$ is both smooth and proper, by the fibre criterion, and both aspherical and coaspherical. But for the constant identity sequence $\alpha$, the category $\alpha/t$ is a disjoint union of two copies of $\mathbf N^{\mathrm{op}}$. Thus these conditions do not imply the universal comma criterion for arbitrary fibrations; the coefficient hypothesis cannot be dropped from the argument.
\end{dcomment}


\dgrblock{dgrb343}{Two disconnected lifting categories}
\begin{dcomment}{examples}{p112}{Two disconnected lifting categories}
The assertion about the projection to the preorder has a small, concrete explanation.  Passing to the preorder remembers that a change of type exists and forgets how many such changes there are.  For smoothness or properness, those forgotten choices must be connected by commuting triangles of lifts.  Natural transformations between interval functors do not by themselves provide such triangles.

Write $q:\Ty\to T_o$ for the projection, with $T_o$ retaining the same objects and having one arrow whenever the corresponding Hom set in $\Ty$ is nonempty.  A necessary condition for $0$-smoothness is the following connected lifting property.  Given $a\in\Ty$, $b\in T_o$, and an arrow $b\to q(a)$, consider the category of arrows $x\to a$ with $q(x)=b$ lying over that arrow; morphisms are commuting triangles over $\mathrm{id}_b$.  This category must be nonempty and connected.  The dual condition, with fixed source, is necessary for $0$-properness.  These are the $W_0$ specializations of the smoothness criterion and its dual, where $W_0$ consists of functors inducing bijections on connected components \citep[Theorem~1.2(b), p.~4]{Maltsiniotis2005Asphericite}.

For smoothness take $b=+-+$ and $a=+$.  The interval of type $b$ has four vertices $0,1,2,3$, and the interval of type $a$ has two.  An increasing endpoint-preserving functor $b\to a$ is forced to collapse the middle, negative edge.  There are exactly two such functors, with vertex lists
\[
 (0,0,0,1)\qquad\hbox{and}\qquad(0,1,1,1).
\]
On the other hand, the only endpoint-preserving increasing endofunctor of $+-+$ is its identity.  To see this, its first internal value must be at most $1$, since a nonconstant image of the first positive edge cannot cross the following negative edge.  Its second internal value must be at least $2$, by the same argument at the last positive edge.  The middle negative edge can then be preserved only when these two values are $1$ and $2$.  As $q$ is the identity on objects, the strict fibre over $b$ has this one object and only its identity arrow.  The lifting category therefore has precisely two objects and only their identities.  It is disconnected, proving that $q$ is not $0$-smooth.

For properness take the fixed source $a=++$ and target $b=+$.  There are exactly two endpoint-preserving increasing functors $a\to b$, with vertex lists
\[
 (0,0,1)\qquad\hbox{and}\qquad(0,1,1).
\]
The only endpoint-preserving increasing endofunctor of $+$ is the identity.  The fixed-source lifting category is consequently another discrete category with two objects.  Thus $q$ is not $0$-proper.

Both examples are compatible with the earlier contractibility theorem for nonempty categories of interval functors.  In that theorem the morphisms are endpoint-fixed natural transformations.  Here they are factorizations through endofunctors of a fixed type.  The distinction explains why connectedness of every mapping category cannot supply the lifting property required on this page.
\end{dcomment}


\dgrblock{dgrb344}{What the preorder forgets about a lift}
\begin{dcomment}{examples}{p112}{What the preorder forgets about a lift}
Grothendieck's proposed sufficient criterion fails for the category of types for a concrete reason: its preorder identifies distinct changes of parameter which cannot be connected inside a lifting category. The definition on manuscript page~99 retains the same objects in $\Ty_o$ and replaces every nonempty set of morphisms by a singleton. It does not identify different objects whose images happen to be isomorphic in that preorder.

Take
\[
 S=(0\longrightarrow1\longleftarrow2\longrightarrow3),
 \qquad T=(0\longrightarrow1).
\]
An endpoint-preserving monotone type map $S\to T$ changes its vertex value from $0$ to $1$ at exactly one of the two forward edges. There are precisely two such maps, with vertex lists
\[
                    (0,1,1,1),\qquad(0,0,0,1).
\]
Both $S$ and $T$ have only the identity endomorphism in $\Ty$. For $S$, a monotone endomorphism can advance by at most one vertex on each edge, since its alternating target has no nonidentity composite of two consecutive edges. Its three increments must sum to three because the endpoints are fixed, so each increment is one. The assertion for $T$ is immediate.

Let $p:\Ty\to\Ty_o$ and consider the unique base arrow $pS\to pT$. An incoming lift with target $T$ has source exactly $S$, and its morphisms are endomorphisms of $S$ commuting with the lifted arrow. Its lifting category is therefore the discrete category on the two displayed maps. The outgoing lifting category with source $S$ is likewise discrete on those maps, now because its morphisms are endomorphisms of $T$. Neither category is connected. The incoming and outgoing lift criteria consequently show that $p$ is neither $0$-smooth nor $0$-proper, and hence neither smooth nor proper. This verifies the source's assertion; failure of a sufficient criterion still does not imply failure of the geometric comparison itself.
\end{dcomment}


Je vais résumer les faits positifs auxquels je suis arrivé dans le diagramme suivant, qui a lieu dans la situation générale $C \longrightarrow T$ envisagée au début du présent paragraphe :

\dgrblock{dgrb345}{The obstruction withdrawn in the footnote}
\begin{dcomment}{context}{p113:Je vais résumer les faits positifs}{The obstruction withdrawn in the footnote}
Grothendieck is trying to explain why an associative composition is easy on all finite paths but elusive on a chosen infinite refinement. The proposed explanation on this page is withdrawn in footnote~24, and the withdrawal is mathematically decisive. A based loop space does admit a strictly associative model. One elementary topological model consists of Moore loops: a loop is a pair $(\ell,\gamma)$, where $\ell\geq0$, $\gamma:[0,\ell]\to X$, and both endpoints are the basepoint. Concatenation adds lengths and joins the maps without rescaling their domains. Its associativity is an equality, and the loop of length zero is a strict unit. Reversal is an inverse up to homotopy, which is exactly the group law required in $\Hot$; a strict monoid model does not require strict inverses.

The simplicial statement is stronger. For a reduced simplicial set $K$, Kan's loop-group construction gives a simplicial group $GK$ and a weak equivalence $K\to\overline W GK$. Consequently $|GK|$ models the based loop space of $|K|$. We checked the precise reducedness hypothesis and the unit equivalence in \citet[Theorem~21]{Stevenson2012}, which gives a modern proof of Kan's theorem. The word ``semi-simplicial'' in this part of the manuscript includes the simplicial formalism with degeneracies; it must not be read as a claim about an arbitrary semisimplicial object without them.

These constructions settle the alleged obstruction, but do not by themselves identify either of Grothendieck's particular categories with the model. That identification is the work still to be done. For $\Sigma=\mathrm{Fl}(X)$, section~13 of the present fragment eventually compares the all-types category with $\Ch_\infty$ and restores the strict concatenation model. The footnote should therefore be read as a correction to an inference, not as an incidental qualification of a theorem. The unsuccessful search described in the surrounding prose is part of the discovery of the later proof.
\end{dcomment}


\msp{page 114}
% foot: Version 22 décembre 2001
\printedpage{72}
% head: A. Grothendieck : Dérivateurs, Ch. VII, Édition des Universités de Montpellier II et Paris VII
\[
\begin{tikzcd}[column sep=small, row sep=huge]
 & & C_{t'} \arrow[d, "{u^{*}\quad\lab{(pour $u : t \longrightarrow t'$\\ dans $T$)}}"] \arrow[ldd, dotted] \arrow[rdd, dotted] & \\
 & & C_{t} \arrow[ld, swap, "{\lab{flèche définie dans\\ Ind\,Cat\,hot\\ (\texttt{[ou dans]}\\ IndHot,\\ Ind\,($\Hotz$))}}"] \arrow[rd, "{\lab[r]{flèche dans Cat\\ (donc dans\\ Cat\,hot,\\ Hot,\\ $\Hotz$)}}"] & \\
\varinjlim C_{t} = & |[draw, circle]| H_{T} \arrow[rd, swap, "{\lab{flèche\\ définie dans\\ Ind\,Cat\,hot\\ ou dans Ind\,Hot\\ ou dans Ind\,$\Hotz$}}"] & \lab{diagramme\\ commutatif\\ dans Ind\,Hot\\ ou Ind\,($\Hotz$)} & C \\
\lab{dans Ind\,Cat\,hot\\ (ou dans Ind\,Hot\\ ou dans Ind\,($\Hotz$))} & & |[draw, circle]| C_{\infty} \arrow[ru, swap, "{\lab[r]{flèche\\ définie\\ dans Hot,\\ donc dans $\Hotz$,\\ en général pas un\\ isomorphisme, même\\ dans $\Hotz$}}"] & \\
 & & = \varinjlim H_{T} \text{ dans } \Hotz &
\end{tikzcd}
\]

\dgrblock{dgrb346}{Strict loop multiplication and the later comparison}
\begin{dcomment}{status}{p113}{Strict loop multiplication and the later comparison}
Grothendieck wants to retain the strictly associative composition already constructed on the total category of paths.  His hesitation here concerns whether that useful algebraic structure can live on a correct model of the loop space.  Footnote~24 retracts the proposed obstruction, and the retraction is mathematically decisive.  A multiplication that is associative only in the homotopy category need not come equipped with the coherent data needed for a strict model.  A loop space, however, already has those data; its origin as a space of loops supplies much more than an isolated group object in the homotopy category.

There is a precise classical replacement.  For a reduced simplicial set $K$, Kan constructs a simplicial group $GK$ and a principal $GK$-bundle over $K$ with weakly contractible total space \citep[Definition~6.3 and Theorem~7.1]{Kan1958Combinatorial}.  After passage to a reduced Kan model, this identifies $GK$ with the based loop homotopy type.  Its multiplication and inverses are strict in every simplicial degree.  Thus the existence of a strictly associative simplicial model cannot disprove the comparison proposed here.  This assertion concerns the existence of a model, rather than a claim that a particular category of paths is that model; the latter still requires a comparison argument.

The same manuscript subsequently supplies that argument for unrestricted paths.  In Chapter~VII, \S13, manuscript pages~157--170, Grothendieck returns to the case $\Sigma=\Fl(X)$ and obtains the comparison through the two projections from the sequential Grothendieck construction.  The endpoint functors and their fibres are essential; a weak equivalence of total categories alone would not justify a statement about their strict fibres.  The elementary lemma that constant paths induce a homotopy equivalence $X\to\operatorname{Fun}(I_\tau,X)$ is left without a written proof at manuscript page~170.  The doubts recorded on the present page therefore belong to the progress of the argument, not to its final conclusion \citep[Ch.~VII, \S13]{GrothendieckDerivateursVII}.

That remaining lemma has a direct proof.  Delete the last vertex of a nontrivial oriented interval, retract it to its neighbour, and let $j$ denote the inclusion of the remaining interval.  The retraction $r$ is a functor, $rj$ is the identity, and either $jr\Rightarrow\mathrm{id}$ or $\mathrm{id}\Rightarrow jr$ is a natural transformation, according to the direction of the deleted edge.  Iterating gives a finite zigzag between the identity of $I_\tau$ and its constant functor at the initial vertex.  Precomposition gives the corresponding zigzag between the identity of $\operatorname{Fun}(I_\tau,X)$ and the composite of evaluation at that vertex followed by the constant-path functor.  The other composite is strictly the identity of $X$.  Thus constant paths give a categorical homotopy equivalence, and hence a weak equivalence of nerves.

For the audit this distinction has two consequences.  The inference from an alleged impossibility of strict loop multiplication is false as written, and is explicitly withdrawn by the author.  The unrestricted path comparison must be recorded with its later positive treatment, whereas the unrestricted \emph{general-fibration} comparison of \S10 still has the counterexamples discussed above.  These statements have different hypotheses.  Neither footnote~24 nor the later path argument removes the base-category obstruction for an arbitrary $C\to T$.
\end{dcomment}


\texttt{[Les annotations dans ce diagramme concernent les flèches en traits pleins.]}

Les objets importants sont $H_{T}$ et $C_{\infty}$, alors que $C$ (sur $T$) et les $C_{t}$ prennent un rôle technique intermédiaire, pour construire $H_{T}$ et $C_{\infty}$. On peut considérer $C_{\infty}$ comme une catégorie (objet de (Cat)), mais dont la construction dépend de nombreux arbitraires, à l'inverse de $H_{T}$ dans Ind\,Cat\,hot. C'est l'objet de $\Hotz$ déduit de $C_{\infty}$ qui a un sens canonique, et peut être décrit comme la limite inductive de $H_{T}$ comme objet de Ind\,$\Hotz$, i.e. comme représentant le même foncteur sur $\mathrm{Hot}_{0}$ que $H_{T}$.

\msp{page 115}
Dans le cas géométrique $T = \Ty$, et $C$ de la forme $\Ch_{\Sigma}(X;a,b)$, donnant lieu à $C_{\infty}(a,b)$, on \textit{n'a pas} de définition naturelle d'un bifoncteur de composition des chemins
\[
C_{\infty}(a,b) \times C_{\infty}(b,c) \ \longrightarrow\ C_{\infty}(a,c).
\]
On peut définir, certes, des accouplements
\[
\Ch_{\Sigma,\tau}(a,b) \times \Ch_{\Sigma,\tau'}(b,c) \ \longrightarrow\ \Ch_{\Sigma,\tau\tau'}(a,c),
\]
ou, avec les notations relatives à la donnée de
\[
\begin{aligned}
&\underline{t} : \mathbf{N}^{o} \longrightarrow \Ty,\\
&\Ch_{n}(a,b) \times \Ch_{m}(b,c) \longrightarrow \Ch_{n+m}(a,c),
\end{aligned}
\]
% foot: Version 22 décembre 2001
\printedpage{73}
% head: A. Grothendieck : Dérivateurs, Ch. VII, Édition des Universités de Montpellier II et Paris VII
mais ces accouplements \textit{ne sont pas} compatibles avec les morphismes de transition des trois systèmes inductifs de catégories $(\Ch_{n})_{n\in\mathbf{N}}$, ils ne le sont \textit{qu'à homotopie près.} Cela suffit cependant pour avoir des accouplements parfaits
\[
H_{T}(a,b) \times H_{T}(b,c) \ \longrightarrow\ H_{T}(a,c),
\]
avec associativité, et (si $\Sigma = \mathrm{Fl}(X)$ tout entier) existence de \textit{vrais} inverses, i.e. de flèches (dans Ind\,Cat\,hot) de `passage au chemin opposé'

\msp{page 116}
\[
H_{T}(a,b) \ \longrightarrow\ H_{T}(b,a), \qquad c \longmapsto c^{o}, \tag*{$(*)$}
\]
donnant lieu au diagramme commutatif exprimant que c'est bien un inverse pour la composition. Une autre façon de le dire, c'est que pour tout objet $Y$ de Cat\,hot (ou de Ind\,Cat\,hot), les ensembles
\[
\mathrm{Hom}_{\mathrm{Cat\,hot}}(Y, H_{T}(a,b)) \overset{\text{déf}}{=} H_{T}^{Y}(a,b)
\]
forment, pour $a,b \in X$ variables, les flèches d'un \textit{groupoïde} (\footnote[25]{Si $\Sigma = \mathrm{Fl}(X)$. Si $\Sigma$ est plus petit, au lieu d'un groupoïde, on trouve une \textit{catégorie} sans plus.}) $\mathrm{Ch}_{\Sigma}^{Y}(X)$ ayant pour ensemble de sommets $\mathrm{Ob}\,X$, et $H_{T}^{Y}(a,b)$ étant formé des flèches de $a$ vers $b$. L'application
\[
H_{T}^{Y}(a,b) \longrightarrow H_{T}^{Y}(b,a), \qquad c \longmapsto c^{o},
\]
déduite du foncteur $(*)$ plus haut n'est autre que l'application de passage à l'inverse. Cette catégorie n'est d'ailleurs autre que celle déduit de % [sic] celle déduit -> celle déduite, accord
\[
X(Y)\Sigma_{Y}^{-1}, \qquad \text{où } \left\{
\begin{aligned}
&X(Y) = \underline{\mathrm{Hom}}_{\mathrm{Cat}}(Y,X)\\
&\Sigma_{Y} \subseteq \mathrm{Fl}(X(Y)),\\
&\text{défini en termes de } \Sigma \text{ `argument par argument'},
\end{aligned}
\right.
\]
à l'aide du foncteur
\[
X \longrightarrow X(Y), \qquad a \longmapsto a_{Y},
\]
\msp{page 117}
associant à tout $a \in X$ le foncteur \textit{constant} sur $Y$ de valeur $a$, en prenant
\[
H_{T}(a,b) \simeq \mathrm{Hom}_{X(Y)\Sigma_{Y}^{-1}}(a_{Y}, b_{Y}) \ .
\]

\dgrblock{dgrb347}{The stage of a concatenation}
\begin{dcomment}{development}{p115}{The stage of a concatenation}
Grothendieck's ind-object accommodates concatenation because a sufficiently refined type can dominate any two given types and their concatenation. The notation $n+m$ is convenient for a multiplicative standard sequence, but does not follow from an arbitrary cofinal sequence alone. The following small calculation locates the missing convention.

\textbf{A slow cofinal sequence.}\ours\quad There is a cofinal decreasing sequence of types $t_n$ for which no type map $t_{n+m}\to t_nt_m$ exists for some $n,m$.

\emph{Proof.} Put $w=\Delta_+\Delta_-$ and
\[
 t_n=w^{k(n)},\qquad k(n)=\lfloor\log_2(n+1)\rfloor.
\]
Use identity maps when $k(n+1)=k(n)$ and collapse the final copies of $w$ otherwise. This gives $t_{n+1}\to t_n$. It is cofinal: every finite signed word is obtained from a sufficiently long alternating word by deleting unwanted edges and expanding the selected edges into its maximal constant-sign runs. For $n=m=3$, however, $t_6=w^2$, whereas $t_3t_3=w^4$. A type map from a word to a reduced alternating target must use at least one source edge for every target edge, since the image of a single source edge is a constant-sign run. Four source edges cannot cover the eight edges of $w^4$. Thus there is no such map.\hfill$\square$

The obstruction is visible on path categories themselves. Take $X=I_{w^4}$, with endpoints $a,b$, and let $c$ be its middle vertex. The inclusions of its two four-edge halves give objects of $\uCh_{t_3}(X;a,c)$ and $\uCh_{t_3}(X;c,b)$. The category $\uCh_{t_6}(X;a,b)$ is empty: every nonidentity arrow of this alternating interval joins adjacent vertices, so four transitions cannot cover the graph distance eight between its endpoints. Hence a pairing with target stage $6$ cannot exist in this example, even with $\Sigma=\mathrm{Fl}(X)$.

For a general chain the correct target is a chosen $k=k(n,m)$ admitting $t_k\to t_nt_m$. Cofinality guarantees such a $k$, and passage to a larger stage or a different chosen map makes no difference in $\Cathot$, by the homotopy theorem. Equivalently, one may choose the standard sequence $t_n=w^n$, where $t_nt_m=t_{n+m}$. We checked the type-map description in the preceding part of this chapter and the usual filtered-colimit formalism; we did not find this particular indexing correction recorded in the literature. Its purpose is to make the intended ind-construction exact, not to alter that construction.
\end{dcomment}


\dgrblock{dgrb348}{Parameters turn the ind-object back into localization}
\begin{dcomment}{exposition}{p116}{Parameters turn the ind-object back into localization}
Grothendieck ends this part by explaining why the ind-object is geometrically useful: its values on parameter categories recover ordinary localization of diagram categories. The assertion has a direct proof, and it supplies associativity and inverses at the same time.

For an actual small category $Y$, write $X(Y)=\uHom(Y,X)$ and let $\Sigma_Y$ consist of natural transformations whose components belong to $\Sigma$. It contains identities and is closed under composition. A functor
\[
 Y\longrightarrow\uCh_{\Sigma,\tau}(X;a,b)
\]
is the same thing as a functor $Y\times I_\tau\to X$ whose endpoints are the constant functors $a_Y,b_Y$, with identity endpoint components on morphisms in $Y$, and whose backward arrows belong pointwise to $\Sigma$. By exchanging the two variables, it is exactly a $\Sigma_Y$-path of type $\tau$ in $X(Y)$ from $a_Y$ to $b_Y$. This identification respects natural transformations and every type-reindexing functor. The two total categories over $\Ty$ are therefore isomorphic.

Taking connected components and applying the localization theorem of \S7 gives
\[
 \begin{split}
 H_T^Y(a,b)
 &:=\Hom_{\operatorname{Ind}(\Cathot)}(Y,H_T(a,b))\\
 &\cong\Hom_{X(Y)[\Sigma_Y^{-1}]}(a_Y,b_Y).
 \end{split}
\]
This corrects two notational slips in the printed formulas: the ambient category on manuscript page~116 is $\operatorname{Ind}(\Cathot)$, and the left side on page~117 needs the superscript $Y$. No finiteness of $Y$ is required for this formula. The uniform type of a $Y$-family is retained by the ind-object; replacing it by $C_\infty$ is the separate step for which finite presentation was needed.

Concatenation on either side is the same operation on functors $Y\times I_\tau\to X$, so this bijection respects composition and units. It follows that the sets $H_T^Y(a,b)$ are the hom-sets of the full subcategory of $X(Y)[\Sigma_Y^{-1}]$ on the constant diagrams. They form a groupoid when $\Sigma$ is all arrows: the reversed path represents the inverse, and the backtracking contractions of \S7 are natural in $Y$. Yoneda in the ind-category turns these pointwise identities into the associativity, unit, and inverse identities of the ind-valued path system. Finite products of ind-objects are represented by the corresponding product diagrams, so the pairings have exactly the domain stated in the manuscript.

If $Y$ itself is an ind-object, the abstract hom-set on the first line still makes sense. For a presentation $Y=(Y_i)$ it is $\lim_i H_T^{Y_i}(a,b)$. The expression $\uHom(Y,X)$ on the second line is only the stated functor category when $Y$ is an actual small category; it should not be silently applied to a formal ind-object. The natural extension is the displayed inverse limit on source indices.
\end{dcomment}


\dgrblock{dgrb349}{What the ind-object represents}
\begin{dcomment}{exposition}{p117:associant à tout $a \in X$}{What the ind-object represents}
The ind-object retains a uniform bound on the complexity of a family of paths. This is why it can encode localization before a single category has acquired a satisfactory geometric interpretation. With $[Y,D]_{\mathrm{cat}}=\pi_0\underline{\mathrm{Hom}}(Y,D)$, the formula is
\[
 H_T^Y(a,b)=\mathop{\rm colim}_n[Y,C_n(a,b)]_{\mathrm{cat}}
 \cong\mathrm{Hom}_{X(Y)[\Sigma_Y^{-1}]}(a_Y,b_Y).
\]
Here $\Sigma_Y$ consists of natural transformations whose components belong to $\Sigma$, and the right side uses the constant diagrams. A morphism on the right is represented by a finite zigzag of natural transformations. Its finite length is a bound shared by every object of $Y$; it therefore gives a map into one stage on the left. Refinement and the elementary relations of localization give the identifications in the colimit. This is the mechanism already discussed in the commentary on Chapter~V, manuscript pages~30--33, concerning the parameter category and ind-representation. The formula printed on page~117 omits the superscript $Y$ from $H_T^Y$; without it an ind-object is being equated to a set.

For finite $Y$, the natural map
\[
 \mathop{\rm colim}_n[Y,C_n]_{\mathrm{cat}}
       \longrightarrow [Y,\mathop{\rm colim}_n C_n]_{\mathrm{cat}}
\]
is bijective. Indeed a functor uses finitely many objects, arrows, and relations, all of which lift to one stage. A finite zigzag of natural transformations uses finitely many more data, so equality also occurs at one stage. Finite presentation, rather than any topological fibrancy, is doing the work. For infinite $Y$, a functor to the colimit can have an unbounded range of stages, and this argument stops.

Composition of localized zigzags is associative because juxtaposition of finite strings is associative. When all arrows of $X$ are inverted, reversal supplies the inverse string. Hence the family of sets $H_T^Y(a,b)$ is a groupoid with object set $\mathrm{Ob}(X)$. For a general $\Sigma$ it is only a category. Passing from this set-valued representation to mapping \emph{spaces} is an additional step: the notation does not identify the ordinary nerve of every functor category with a derived mapping object.
\end{dcomment}


Dans le même cas géométrique, je me suis convaincu d'ailleurs que $C_{\infty}$ était, en tant qu'objet de Hot (pas seulement de $\Hotz$) \textit{unique à} Hot-\textit{isomorphisme près,} mais (chose étrange) peut-être pas à Hot-isomorphisme \textit{unique} près. De façon générale, revenant aux
% foot: Version 22 décembre 2001


\dgrblock{dgrb350}{From path components to mapping spaces}
\begin{dcomment}{later}{p117}{From path components to mapping spaces}
The last formula brings the discussion back to its original purpose: families of paths should describe morphisms after localization, and should continue to do so when an indexing category is introduced.  At the level of sets of components, the preceding arguments give a complete description.  To retain higher homotopies, a modern treatment keeps the diagram of zigzag categories together with the coherences in its changes of shape.

One precise formulation is the hammock formula.  For a relative infinity-category $(\mathcal C,\mathcal W)$, let $\mathrm{Zig}_{Z}(\mathcal C,\mathcal W)_{x,y}$ denote the infinity-category of zigzags of shape $Z$ with fixed endpoints and pointwise weak-equivalence transformations.  Write $|{-}|$ for groupoid completion, which for an ordinary category is represented by the realization of its nerve.  Then, with colimit taken in spaces,
\[
 \operatorname{Map}_{\mathcal C[\mathcal W^{-1}]}(x,y)
 \simeq
 \mathop{\operatorname{colim}}_{Z\in\mathrm{Zig}_{\partial}^{\mathrm{op}}}
 \bigl|\mathrm{Zig}_{Z}(\mathcal C,\mathcal W)_{x,y}\bigr|.
\]
Here $\mathrm{Zig}_{\partial}$ has endpoint-preserving changes of shape.  Arakawa and Cnossen prove this formulation in their 2026 preprint, Definition~1.7, Theorem~1.9 and Corollary~1.11; Remark~1.10 identifies composition with concatenation \citep{ArakawaCnossen2026}.  Their \S4 compares the construction with Dwyer--Kan's original hammock localization.

The notation must not obscure a difference in definitions.  Grothendieck's $\Ch_{\Sigma,\tau}$ on these pages is a \emph{full} subcategory of the functor category: membership constrains the backward arrows of an object, while morphisms are arbitrary natural transformations with the specified endpoint condition.  The hammock construction just quoted also constrains the vertical components of those transformations.  A comparison of the two models for general $\Sigma$ is an additional theorem, not a consequence of recognizing the same words of signs.  When all arrows are marked this difference disappears, and the later comparison within Chapter~VII is directly relevant.  This qualification is particularly useful when passing from $X$ to $X(Y)$: it prevents an equality of localized morphism \emph{sets} from silently being promoted to an identification of mapping \emph{spaces}.
\end{dcomment}


\dgrblock{dgrb351}{Independence of the sequential homotopy type}
\begin{dcomment}{development}{p117}{Independence of the sequential homotopy type}
Grothendieck's final distinction between an isomorphism and a uniquely determined isomorphism can be made precise without solving the comparison with the total category $C$.  The two questions separate: a bad base can obstruct $C_\infty\to C$ while all the sequential models still have the same weak homotopy type.  The following argument gives the latter conclusion under the abstract hypotheses of \S10.  Its ingredients are standard telescope arguments; we include their application here because we did not locate this exact reachability-preorder statement in the sources checked.  For the usual sequential homotopy-colimit construction see \citet[\S14.9, Proposition~14.10 and Corollary~14.11]{DuggerHocolim}.

\textbf{Proposition (independence of choices).}\ours\quad
Let $T$ be a nonempty small category and let $F:T^{\mathrm{op}}\to\Cat$ be a strict functor.  Suppose that $F(u)$ and $F(v)$ are joined by a finite zigzag of natural transformations whenever $u,v$ are parallel arrows of $T$.  Let
\[
 t_0\longleftarrow t_1\longleftarrow\cdots,
 \qquad s_0\longleftarrow s_1\longleftarrow\cdots
\]
be two chains whose objects are cofinal to the left in the reachability preorder of $T$.  Then
\[
 \varinjlim_n F(t_n)\quad\hbox{and}\quad\varinjlim_m F(s_m)
\]
have isomorphic images in $\Hot$.  The assertion includes arbitrary choices of the arrows in either chain.  It does not assert a preferred isomorphism.

\textit{Proof.}
We first record precisely which properties of sequential homotopy types will be used.  For a sequence of small categories $A_n$, the nerve commutes with its filtered colimit.  Indeed, a simplex is a finite string of objects and arrows subject to finitely many incidence equations; these data, and equalities between two presentations of them, occur at some common stage.  Hence
\[
 N\bigl(\varinjlim_n A_n\bigr)=\varinjlim_n N A_n.
\]
For a sequence of simplicial sets $K_n$, the canonical map
\[
 \operatorname{Tel}(|K_n|)\longrightarrow
       \bigl|\varinjlim_n K_n\bigr|
\]
is a weak equivalence.  Here the telescope glues the upper end of
$|K_n|\times[0,1]$ to the next space by its transition map.  To check this assertion on homotopy groups, replace simplicial sets by $\operatorname{Ex}^{\infty}K_n$.  This replacement commutes with filtered colimits: $\operatorname{Sd}\Delta^r$ is finite for each $r$, so $\operatorname{Ex}$ commutes with them, and the sequential iteration does also.  A filtered colimit of Kan complexes is Kan, because every horn is finite.  Maps and homotopies defining a homotopy group use finite simplicial data, so homotopy groups, with a basepoint represented at some stage, commute with this filtered colimit.  On the telescope side a map or a homotopy from a compact sphere or cylinder has bounded height.  Each bounded initial telescope retracts to its final space.  Its homotopy groups therefore give the same direct limits.  The canonical map induces these identifications, including the statement about components, and is a weak equivalence.

We also need invariance when each transition map is replaced by a homotopic map.  Suppose $f_n,g_n:A_n\to A_{n+1}$ are continuous maps with homotopies $H_n$ from $g_n$ to $f_n$.  A map from $\operatorname{Tel}(f)$ to $\operatorname{Tel}(g)$ is the identity on every level $A_n$.  On the cylinder for $f_n$, it first traverses the cylinder for $g_n$, then follows $H_n$ in level $A_{n+1}$.  More explicitly it sends $(x,r)$, for $0\leq r\leq\tfrac12$, to $(x,2r)$ in the $g_n$ cylinder, and for $\tfrac12\leq r\leq1$ to $H_n(x,2r-1)$ in level $A_{n+1}$.  The endpoint at $r=1$ is $f_n(x)$, so these maps glue.  Reversing the homotopies gives a map back.  Their composites are homotopic to the identities, fixing the levels, by cancelling a path immediately followed by its reverse and reparametrizing the cylinders.  These cancellations are defined separately on each cylinder and agree at its ends; they consequently glue on the entire telescope.  No compatibility between different $H_n$ is required.  Finally, passing to a cofinal subsequence preserves the weak homotopy type of a telescope, by the same bounded-height calculation of homotopy groups.

We now compare the chains.  Choose indices, increasing without bound, and arrows in $T$ to obtain an interleaving
\[
 t_{n_0}\longleftarrow s_{m_0}\longleftarrow
 t_{n_1}\longleftarrow s_{m_1}\longleftarrow\cdots.
\]
At each step this is possible by cofinality of the opposite chain; after finding an index, moving farther down that chain allows the indices to increase.  Apply $F$ and take nerves and realizations.  The even subsequence has objects $|NF(t_{n_i})|$.  Its bonding map is obtained from the composite of two of the newly chosen arrows.  That composite and the original transition $t_{n_{i+1}}\to t_{n_i}$ are parallel arrows of $T$.  By hypothesis their pullbacks are connected by natural transformations; taking nerves turns these into a homotopy of continuous maps.  The preceding telescope construction therefore identifies the homotopy type of this even subsequence with that of the original $t$-sequence.  The odd subsequence has the corresponding property for the $s$-sequence.  Both are cofinal subsequences of the same interleaved telescope.  Combining these comparisons with the telescope-to-colimit weak equivalences proves the proposition.\hfill$\square$

\end{dcomment}
\begin{otherlanguage}{english}
\dgroverview{VII-b}{Chapter VII \textperiodcentered\ \textup{VII-b} \textperiodcentered\ Overview}
\selectlanguage{english}
\noindent\textbf{Paths, localization, and the information in an indexing category}



\medskip
These pages occupy the middle of Chapter~VII.  They begin with the end of the combinatorial analysis of path types, establish the ordinary category of fractions by means of path categories, and then examine whether a category containing all shapes of paths has the same homotopy type as a sequentially enlarged model.  Their central concern is a compatibility problem.  Concatenation is completely satisfactory before taking limits: it is functorial, associative and unital.  The sequential models make other calculations accessible, but their transition maps need not respect that composition strictly.  Grothendieck is trying to preserve both advantages in one construction.

The category $\underline{\mathrm{Ty}}$ consists of finite oriented intervals, equipped with their spatial order and their two endpoints.  Its arrows preserve that order, the categorical direction of each traversed edge, and the endpoints.  Concatenation makes it a strict monoidal category.  The opening calculation describes its nerve degree by degree as a free monoid: a composable string of type maps splits uniquely over the edges of its initial interval.  The data attached to one initial edge are successive finite weightings recording how many edges replace it at each stage.  The commentary on the first part of the chapter explains this calculation through manuscript page~57.  The commentary here adds the distinction needed for its later use: a face can send an indecomposable generator to a product.  Degreewise freeness is therefore not freeness on a simplicial set.

For an ordinary category $X$, a path of type $\tau$ is a functor $I_\tau\to X$.  Between paths of the same type one keeps natural transformations, and between different types one also keeps pullback along interval maps.  These categories assemble into a category fibred over $\underline{\mathrm{Ty}}$.  Fixing endpoints $a,b$ means taking the strict fibre of the endpoint functor: the components of a natural transformation at both ends must be identities.  This requirement is responsible for several of the correct diagrams in the text.  Dropping it would change the hom-sets obtained after passing to components.

If $\Sigma$ is a wide subcategory of $X$, we allow backward edges only when they lie in $\Sigma$, and write $\Ch_\Sigma(X;a,b)$ for the resulting total fixed-endpoint path category.  Closure of $\Sigma$ under identities and composition is exactly what makes this construction stable under every type map.  A source edge can collapse, or cover a consecutive block of target edges; the two possibilities require precisely these two closure conditions.  No calculus-of-fractions hypothesis is needed at this stage.  The fibres are full subcategories of functor categories, so their internal vertical transformations remain unrestricted.

The first decisive result is
\[
 \operatorname{Hom}_{X[\Sigma^{-1}]}(a,b)
      \cong \pi_0\Ch_\Sigma(X;a,b).
\]
Concatenation induces the composition on the right.  A one-edge path followed by its reversal is connected to a constant path through a ladder with identity endpoint components, and the constant path is connected to the empty path by a type collapse.  Thus every arrow of $\Sigma$ becomes invertible.  For a functor $F:X\to Y$ which inverts $\Sigma$, evaluation of a zigzag in $Y$ is unchanged under either a ladder or a type change.  In a ladder, the equations for its successive squares telescope to
\[
 \eta_n E_F(c)=E_F(c')\eta_0;
\]
the endpoint identities give equality of effects.  The argument uses inverses of the backward horizontal edges but never inverses of the internal vertical components.  It supplies the omitted well-definedness check for the map $\beta$ on manuscript page~68.  The resulting evaluation has the required unique factorization property.  This is the classical localization problem of \citet[Chapter~I, \S1]{GabrielZisman1967}, with the witnesses for its relations retained as categorical data.

Two examples make the theorem tangible.  Inverting one arm of a span creates the fraction across the span and produces a category with a walking-arrow skeleton.  Inverting both arrows of the category $a\rightrightarrows b$ produces a connected groupoid with automorphism group $\mathbf Z$; the nerve of the original category is a circle.  The complete calculations distinguish cancellation of backtracking from the survival of a genuine loop.  They also show why a relative localization need not be a groupoid.

The next task is to establish that concatenation really acts on all morphisms of the total path categories.  A ladder can be glued to a second ladder because their common endpoint components are identities.  A change of type can be glued to a second change of type because the common endpoint has a specified image.  Writing a total morphism as a type map together with a vertical transformation makes the functoriality and strict associativity transparent.  The commentary gives this calculation explicitly.  It also repairs an edition diagram which glues over the two-element set $\{0,1\}$ where the construction requires a single identified endpoint.

Sections~8 and~9 then explain why the particular interval map inducing a pullback ceases to matter after categorical homotopy.  The relevant homotopy relation joins functors by finite zigzags of natural transformations.  Grothendieck proves that every nonempty category of endpoint-preserving interval maps is connected, and asks whether it is contractible.  We give a complete categorical contraction.  After a forced initial collapse when the first signs disagree, the first target run can be filled by the first source edge.  Clamping the remaining vertex images to the end of that run defines an idempotent functor and a natural transformation in the direction determined by the run's sign.  The retract is a smaller interval-mapping category.  Induction contracts it to a point.  This is marked as a development of the present commentary, without a claim of priority.

Precomposition transports these natural transformations to the path categories.  Because the endpoint components remain identities, the homotopy takes place over $X\times X$ and restricts to every fixed-endpoint fibre.  In modern notation, the diagram of fibres becomes constant on parallel arrows after passage to $\Cathot$, the category whose morphisms are categorical-homotopy classes of functors.  This is the bridge to the abstract discussion of a fibration $C\to T$ in \S10.

Here three objects must be distinguished.  There is the total category $C$, the ind-object determined by its fibres in $\Cathot$, and an ordinary sequential colimit $C_\infty$ formed after choosing a chain cofinal in the reachability preorder $T_{\mathrm o}$.  The source assumes that this preorder is cofiltered.  It does not assume that parallel arrows of $T$ can be equalized.  Forgetting that difference would conceal the principal obstruction: for the constant point-valued diagram over a one-object group $BG$, the ind-object and every sequential model are points, while the total category is $BG$.

There is a further failure on manuscript page~103.  The reachability preorder of a slice $T/t_0$ need not be the lower ideal of $t_0$ in $T_{\mathrm o}$.  The category with two parallel arrows already disproves that assertion.  The lower ideal may still be used by restricting to the full subcategory of objects which admit an arrow to $t_0$, but this is only cofinality on reachability preorders.  Comparison of total categories requires the actual comma categories.  The commentary formulates and proves the corresponding sufficient-and-necessary condition for base change to preserve the homotopy type of every Grothendieck fibration, using Quillen's Theorem~A \citep[\S1]{Quillen1973}.  The source's separate properness and smoothness criteria are also proved under its hypothesis on parallel pullbacks.  Contractibility of the strict fibres of $T\to T_{\mathrm o}$, together with local constancy of the coefficients there, allows homotopy Kan extension to descend the diagram to the preorder.  This argument explains why those sufficient criteria can hold without the universal comma condition \citep{Maltsiniotis2005Asphericite,Cisinski2009LocallyConstant}.

The sequential colimit itself is less problematic.  The nerve preserves filtered colimits, and for a sequence its homotopy colimit maps by a weak equivalence to its ordinary colimit.  We prove that two cofinal chains give isomorphic objects of $\Hot$ under the source's homotopy hypothesis: interleave cofinal subsequences, compare their bonding maps by homotopies, and compare the resulting telescopes.  This answers the existence part of the independence question.  The proof neither chooses a preferred isomorphism nor promotes it to a categorical homotopy equivalence in $\Cathot$.

Finite probes explain part of the remaining caution.  A finite category is finitely presented for ordinary categorical filtered colimits, which gives the asserted formula for $\Cathot$-morphisms into $C_\infty$.  Its nerve need not have the homotopy type of a finite CW complex.  The one-object category of the group of order two is a concrete counterexample to unrestricted compactness after passage to $\Hot$: its mod-two cohomology persists in every degree, whereas each finite-dimensional skeleton of its nerve has bounded cohomological dimension.  Replacing $\Cathot$ by $\Hot$ therefore needs a separate argument.  The mapping-set formulas at the end of manuscript page~117 are audited with their parameter $Y$ and their ambient ind-category restored in the commentary.

The proposed obstruction from strict loop multiplication on manuscript page~113 is already withdrawn in footnote~24.  Kan's loop-group theorem supplies strict simplicial group models \citep[Theorem~7.1]{Kan1958Combinatorial}.  For unrestricted paths, the specific comparison is taken up positively later in this same chapter, \S13; it is not a current open problem.  The constant-path lemma left unwritten at manuscript page~170 is supplied here by successively retracting the end vertices of an oriented interval.  Modern hammock localization gives a separate, precise mapping-space theory.  Its requirement that vertical arrows also be weak equivalences must be checked when comparing it with Grothendieck's full path fibres \citep{DwyerKan1980Calculating,ArakawaCnossen2026}.
\end{otherlanguage}
\dgrchunkend

% DER_VII-c: Grothendieck's text with the commentary set in, then the overview, inside a group that applies the chunk's own definitions
\dgrpart{1}{VII-c}{Part 3 of 3, manuscript pages 118--171}{VII-c\enspace\textperiodcentered\enspace Chapter VII, part 3 of 3\enspace\textperiodcentered\enspace manuscript pages 118--171}{Chapter VII \textperiodcentered\ \textup{VII-c}}
\dgrchunkbegin
% DER_VII-c: the definitions of the chunk's preamble, applied inside its group
\providecommand\IfDocumentMetadataT[1]{}
\let\printedpage\relax\newcommand{\printedpage}[1]{\marginpar{\footnotesize p.~#1}}
\let\msp\relax\newcommand{\msp}[1]{\marginpar{\footnotesize [#1]}}
\let\Hotz\relax\newcommand{\Hotz}{\mathrm{Hot}^{\wedge}_{0}}
\let\Ch\relax\newcommand{\Ch}{\underline{\mathrm{Ch}}}
\let\Ty\relax\newcommand{\Ty}{\underline{\mathrm{Ty}}}
\let\CNo\relax\newcommand{\CNo}{C_{\mathbf{N}^{o}}}
\let\lab\relax\newcommand{\lab}[2][l]{\text{\footnotesize\begin{tabular}{@{}#1@{}}#2\end{tabular}}}
\let\Hot\relax\newcommand{\Hot}{\mathrm{Hot}}
\let\Hc\relax\newcommand{\Hc}{\underline{\mathrm{Hom\,cof}}}
\providecommand{\No}{}
\providecommand\No{}\renewcommand{\No}{\mathbf{N}^o}
\let\uHom\relax\newcommand{\uHom}{\underline{\mathrm{Hom}}}
\let\comp\relax\newcommand{\comp}{\operatorname{comp}}
\let\uCh\relax\newcommand{\uCh}{\underline{\mathrm{Ch}}}
\let\hot\relax\newcommand{\hot}{\mathrm{hot}}
\let\dashto\relax\newcommand{\dashto}[1]{\mathrel{\overset{#1}{\rule[0.6ex]{2em}{0.4pt}}}}
\let\CCh\relax\newcommand{\CCh}{\underline{\underline{\mathrm{Ch}}}}
\let\Xb\relax\newcommand{\Xb}{\overline{X'}}
\let\Xbb\relax\newcommand{\Xbb}{\overline{\overline{X'}}}
\let\cC\relax\newcommand{\cC}{\mathcal{C}}
\let\cX\relax\newcommand{\cX}{\mathcal{X}}
\let\NN\relax\newcommand{\NN}{\mathbf{N}^{o}}
\let\Ob\relax\newcommand{\Ob}{\operatorname{Ob}}
\let\CN\relax\newcommand{\CN}{\mathcal{C}_{\mathbf{N}^o}}

% ---- from Der-ChVII-456_p076-080.tex
\printedpage{74}
% head: A. Grothendieck : Dérivateurs, Ch. VII, Édition des Universités de Montpellier II et Paris VII
conditions du début du paragraphe, $C_\infty$ doit s'écrire $C_\infty(\underline{t})$, elle dépend du choix du foncteur cofinal
\[
\underline{t} : \No \longrightarrow T.
\]
Quand $\underline{t}$ varie dans la catégorie
\[
\Hc(\No, T),
\]
les $C_\infty(\underline{t})$ varient de façon contravariante en $\underline{t}$ (donc définissent une catégorie fibrée sur $\Hc(\No, T)$). De plus, le foncteur
\[
C_\infty(\underline{t'}) \longrightarrow C_\infty(\underline{t})
\]
associé à un morphisme de systèmes projectifs
\[
\underline{t} \longrightarrow \underline{t'}
\]
est un \textit{isomorphisme} dans Hot. Donc le
\msp{page 118}
foncteur
\[
\underline{t} \longmapsto C_\infty(\underline{t}), \qquad \Hc(\No, T) \longrightarrow (\Hot)
\]
se factorise par le groupoïde fondamental
\[
\Pi(\Hc(\No, T)) \longrightarrow (\Hot).
\]
Quand ce groupoïde, i.e. quand la catégorie $\Hc$ elle-même est connexe, on peut donc conclure que tous les $C_\infty(\underline{t})$ sont isomorphes entre eux dans Hot. Pour conclure même à un système \textit{transitif} d'isomorphismes entre les $C_\infty(\underline{t})$, i.e. que $C_\infty(\underline{t})$ est ‘unique à isomorphisme \textit{unique} près’, il faudrait savoir de plus que le groupoïde fondamental est équivalent à la catégorie ponctuelle (groupe fondamental trivial), ou, comme on dit, que $\Hc$ soit \textit{simplement connexe,} ou 1-connexe.

La connexité de $\Hc$ est assurée quand la catégorie $T$ a (comme c'est le cas pour $\underline{\mathrm{Ty}}$) une loi de composition $(t,t') \longmapsto t * t'$, avec des morphismes fonctoriels en $t, t'$
\[
\begin{tikzcd}[row sep=small]
 & t \\
t * t' \arrow[ru] \arrow[rd] & \\
 & t'
\end{tikzcd}
\]
(cette dernière condition étant automatique
% foot: Version 22 décembre 2001
\printedpage{75}
% head: A. Grothendieck : Dérivateurs, Ch. VII, Édition des Universités de Montpellier II et Paris VII
\msp{page 119}
si la loi de composition admet l'objet final de $T$, supposé existant, comme unité bilatère). Dans ce cas, $\underline{t} = (t_n)_n$ et $\underline{t'} = (t'_n)_n$ sont reliés entre eux par $\underline{t} * \underline{t'} = (t_n * t'_n)_n$, puisqu'on a un morphisme de systèmes projectifs
\[
\begin{tikzcd}[column sep=small]
 & \underline{t} * \underline{t'} \arrow[ld] \arrow[rd] & \\
\underline{t} & & \underline{t'}.
\end{tikzcd}
\]
Je ne sais d'ailleurs pas si la flèche dans Hot
\[
C_\infty(\underline{t}) \longrightarrow C_\infty(\underline{t'})
\]
associée à $\underline{t'} \longrightarrow \underline{t}$ dans $\Hc$ ne dépend que de $\underline{t}$, $\underline{t'}$, et pas du choix de $\underline{t'} \longrightarrow \underline{t}$. Ce qui est clair, c'est que son image dans $\Hotz$ est unique, c'est l'identité quand on identifie $C_\infty(\underline{t})$ et $C_\infty(\underline{t'})$ au seul et même élément de $\Hotz$ défini par $\varinjlim H_T$.

Et pourtant, je m'aperçois à l'instant que, sous l'hypothèse de l'existence de la loi de composition et de son unité bilatère, et admettant de plus que la loi en question est associative à isomorphisme près, \textit{on a transitivité} dans les isomorphismes (canoniques!) entre les $C_\infty(\underline{t})$ qu'on vient d'indiquer !

\msp{page 120}
(Sans avoir pour autant à calculer le groupe fondamental de $\Hc(\No, T)$!) Si on a trois éléments $\underline{t}$, $\underline{t'}$, $\underline{t''}$ de $\Hc$, il suffit de regarder le diagramme déduit du diagramme commutatif suivant
\[
\begin{tikzcd}[column sep=small]
 & & \underline{t} * \underline{t'} * \underline{t''} \arrow[ld] \arrow[rd] & & \\
 & \underline{t} * \underline{t'} \arrow[ld] \arrow[rd] & & \underline{t'} * \underline{t''} \arrow[ld] \arrow[rd] & \\
\underline{t} & & \underline{t'} & & \underline{t''}
\end{tikzcd}
\]
en lui appliquant le contrafoncteur $\underline{s} \longrightarrow C_\infty(\underline{s})$ de $\Hc$ dans Hot.

Donc dans le cas des $C(a,b) = \underline{\mathrm{Ch}}_\Sigma(X; a, b)$, je n'ai plus aucun scrupule à considérer la catégorie $C_\infty(a,b)$ comme un objet \textit{canonique} de Hot, relié à l'ind-objet canonique $H_T(a,b)$ dans Hot par un homomorphisme canonique
\[
C_\infty(a,b) \longleftarrow H_T(a,b), \tag*{$(*)$}
\]
qui dans $\Hotz$ donne un isomorphisme
\[
C_\infty(a,b) \simeq \varinjlim H_T(a,b), \tag*{$(**)$}
\]
i.e. la flèche $(*)$ induit, pour $\xi \in \mathrm{Ob}\,\Hot$, une application ensembliste
\[
\mathrm{Hom}_{\Hot}(\xi, C_\infty(a,b)) \longleftarrow \underbrace{\mathrm{Hom}_{\mathrm{Ind}\,\Hot}(\xi, H_T(a,b))}_{\underset{n}{\varinjlim}\, \mathrm{Hom}_{\Hot}(\xi, C_n(a,b))},
\]
qui est \textit{bijective} si $\xi$ est dans Hot.
% foot: Version 22 décembre 2001

\printedpage{76}
% head: A. Grothendieck : Dérivateurs, Ch. VII, Édition des Universités de Montpellier II et Paris VII
\msp{page 121}
Bien entendu, les flèches
\[
C_\infty(\underline{t'}) \longrightarrow C_\infty(\underline{t}) \tag*{$(*)$}
\]
dans Hot, associés aux flèches $\underline{t} \longrightarrow \underline{t'}$ dans $\Hc(\No, T)$, sont en fait des flèches dans % [sic] associés -> associées, accord avec « les flèches »
(Cat) elle-même - et même des foncteurs pleinement fidèles, dans le cas ‘géométrique’ et pour les deux flèches canoniques
\[
\begin{tikzcd}[column sep=small]
 & \underline{t} * \underline{t'} \arrow[ld] \arrow[rd] & \\
\underline{t} & & \underline{t'}.
\end{tikzcd}
\]
Elles sont isomorphiques \textit{dans} Hot, parce qu'elles définissent un homomorphisme sur les ind-objets associés de Cat\,hot, donc aussi pour les ind-objets associés dans Hot et dans $\Hotz$, d'où résulte que
\[
\mathrm{Hom}_{\Hot}(Y, C_\infty(\underline{t'})) \longrightarrow \mathrm{Hom}_{\Hot}(Y, C_\infty(\underline{t}))
\]
associé à $(*)$ est une bijection si $Y$ \textit{est finie} (\footnote[26]{Ce qui suffit à assurer que $(*)$ est \texttt{[un]} isomorphisme dans Hot : $\Hot \longrightarrow \Hotz$ est conservatif!}), cette flèche s'identifie alors à
\[
\underset{n}{\varinjlim}\, \mathrm{Hom}_{\Hot}(Y, C_{t'_n}) \longrightarrow \underset{n}{\varinjlim}\, \mathrm{Hom}_{\Hot}(Y, C_{t_n}).
\]
(Car $Z \longmapsto \mathrm{Hom}_{\Hot}(Y, Z)$ commute aux limites inductives filtrantes, si $Y$ est une catégorie \textit{finie.} Cela provient du fait qu'une $\varinjlim$ filtrante de catégories de Thomason-Kan est itou \dots) Mais j'ignore par contre si $(*)$ est \texttt{[un]} isomorphisme dans Cat\,hot. Seulement que pour une catégorie \textit{finie} $Y$, elle induit
\[
\mathrm{Hom}_{\mathrm{Cat\,hot}}(Y, C_\infty(\underline{t'})) \xrightarrow{\ \sim\ } \mathrm{Hom}_{\mathrm{Cat\,hot}}(Y, C_\infty(\underline{t})),
\]
les deux membres s'identifiant alors à
\msp{page 122}
\[
H_T(Y)_{\mathrm{Cat\,hot}} = \mathrm{Hom}_{\mathrm{Ind\,Cat\,hot}}(Y, H_T).
\]
Pour résumer ces lumières et perplexités sur les $C_\infty(\underline{t})$ :

\dgrblock{dgrb384}{Contracting the category of choices}
\begin{dcomment}{development}{p120}{Contracting the category of choices}
Grothendieck seeks a coherent comparison between different cofinal sequences, so that changing a path parameter does not change the resulting operation in $\Hot$. His three-sequence diagram supplies this comparison, but the natural projections already imply a stronger statement about the space of choices.

\textbf{Contraction of the refinement category.}\ours\quad
Let $T$ be a category with a bifunctor $*$ and natural maps $t*t'\to t$ and $t*t'\to t'$. Let $\mathcal H$ be the nonempty category of sequences $\mathbf N^o\to T$ cofinal in the preorder of objects. Then $N\mathcal H$ is contractible. Consequently every functor $\mathcal H^o\to\mathcal A$ taking all arrows to isomorphisms has a unique system of comparison isomorphisms induced by arrows and their inverses; these comparisons depend only on their two endpoints and compose transitively.

\textit{Proof.} Fix $s\in\mathcal H$. Define $Q(t)=t*s$ pointwise, on both objects and morphisms of sequences. For every $a\in T$ there is an index $n$ and a map $t_n\to a$. Composing $t_n*s_n\to t_n\to a$ shows that $Q(t)$ is still cofinal. The two natural projections give natural transformations
\[
 \mathrm{id}_{\mathcal H}\ \longleftarrow\ Q\ \longrightarrow\ \mathrm{const}_s.
\]
The nerve of a natural transformation gives a homotopy of the associated maps on realizations: explicitly, a chain of arrows and a cut in that chain use the first functor before the cut, the transformation at the cut, and the second functor afterwards. Thus the identity of $B\mathcal H$ is homotopic to the constant map at $s$, proving contractibility.

A functor taking all arrows to isomorphisms factors through the localization of $\mathcal H$ at all its arrows, its fundamental groupoid. Contractibility makes this groupoid equivalent to a point: between any two of its objects there is exactly one morphism. Applying the functor to that morphism proves the uniqueness and transitivity assertions, including independence from the particular arrow of sequences used. The opposite variance changes the direction of the comparisons but not the argument.\hfill$\square$

Associativity of $*$ and a terminal unit are unnecessary for this proof once the two natural projections are given. In the interval application, concatenation has these projections by collapsing the other interval, and a cofinal sequence exists. The proposition isolates the formal part of the argument. It does \emph{not} prove that the actual functors $C_\infty(t')\to C_\infty(t)$ are weak equivalences; that premise requires the geometric comparison or a valid relative argument. In particular the conservativity assertion in footnote~26 cannot supply it in the stated generality. We checked the general nerve and cofinality mechanisms in Quillen's \S1 and the smooth-functor formulation of \citet[Theorem~1.2]{Maltsiniotis2005Asphericite}; we did not locate this explicit application to Grothendieck's category of sequences in those sources. The result is an elementary continuation of his construction, without a claim of priority.
\end{dcomment}


\dgrblock{dgrb385}{Two different finiteness arguments}
\begin{dcomment}{exposition}{p121}{Two different finiteness arguments}
The comparison on this page must keep finite categorical data separate from a finite homotopy type. The preceding block on ind-representation proves a filtered-colimit formula in $\mathrm{Cat\,hot}$ whenever the source category is finitely presented. In $\Hot$, the corresponding safe hypothesis is that the classifying space of the source has the homotopy type of a finite CW complex, or more generally is compact for the filtered homotopy colimits in question. A finite poset has a finite nerve and meets this hypothesis. But Grothendieck explicitly defines a finite category by finiteness of both its objects and its arrows in Chapter~V, manuscript page~45, and uses these categories to define $\Hot_0$ on pages~51--52 \citep{Grothendieck1990V}. With that convention, the assertion is false.

For a concrete distinction, let $Y$ be the one-object category of the group $\mathbf Z/2$. It is finite, but $BY\simeq\mathbf{RP}^{\infty}$. Consider the increasing skeleta $\mathbf{RP}^n$. The identity of $\mathbf{RP}^{\infty}$ cannot factor, up to homotopy, through one skeleton: such a factorization would make the identity on $H^k(\mathbf{RP}^{\infty};\mathbf F_2)$ factor through zero for every $k>n$, whereas these cohomology groups are nonzero. Thus
\[
 \mathop{\rm colim}_n[BY,\mathbf{RP}^n]
       \longrightarrow[BY,\mathbf{RP}^{\infty}]
\]
is not surjective. The same example can be expressed entirely in categories by using compatible face-poset models for the finite skeleta, or the categories of simplices of their simplicial models. It is the size of the nerve in all dimensions, not merely the number of arrows, that matters here.

For a finite simplicial source $K$, the valid filtered-colimit argument is direct. Replace the targets by a functorial Kan replacement built using finite horn data. Such a replacement commutes with filtered colimits. Maps from $K$ and simplicial homotopies between them involve finitely many simplices, so they appear at one stage. Finite CW sources follow by a finite simplicial model up to homotopy. The statement about filtered colimits of Thomason-fibrant categories is also valid, since the relevant horn domains are finite after applying the fixed subdivisions. It does not establish compactness for $B(\mathbf Z/2)$.

Accordingly, the literal last sentence of page~120 omits even the intended subscript zero; restoring that subscript would still leave the compactness error. A valid replacement restricts to homotopy-finite tests and thereby changes the stated test category. This repair does not alter the elementary finite-presentation calculation in $\mathrm{Cat\,hot}$.
\end{dcomment}


\dgrblock{dgrb386}{The unpointed detection claim}
\begin{dcomment}{status}{p121}{The unpointed detection claim}
Grothendieck uses finite tests here to recognize a weak equivalence. In the unpointed homotopy category, the assertion that $\Hot\to\widehat{\Hot_0}$ is conservative is false. The editors already flag the claim in Chapter~V, manuscript page~52, and the commentary there gives the symmetric-group shift example; we do not repeat it. The general theorem is stronger: no set of spaces jointly detects all isomorphisms in the unpointed homotopy category \citep[Theorem~2.1]{ArlinChristensen2023}. Their Theorem~4.3 shows why retaining one more level of data helps: spheres do detect equivalences in the homotopy 2-category, whose morphism categories retain homotopy classes of homotopies.

This defect in the inference does not disprove the particular path comparisons. They can be proved with based or relative information, or, in the all-arrows situation of section~13, by comparing the explicit total path categories over $X\times X$. An audit must separate that geometric theorem from the invalid general detection principle used on this page. For based spaces, all basepoints and their homotopy groups give the usual Whitehead test. For free homotopy classes, the fundamental-group information lost to conjugacy must first be restored.
\end{dcomment}


\dgrblock{dgrb387}{A comparison that fails under the general hypotheses}
\begin{dcomment}{development}{p121}{A comparison that fails under the general hypotheses}
The preceding distinction can be tested inside the exact fibred setting of section~10, rather than only in the abstract homotopy category. Grothendieck assumes there that parallel reindexing functors are homotopic, that the preorder of the base is cofiltered and has a countable coinitial subset, and that a strict cleavage has been chosen. These assumptions produce the same ind-object, but they do not force every map between the resulting colimits to be a weak equivalence.

\textbf{A fibred symmetric-group example.}\ours\quad
There is a strict fibred category $C\to T$ satisfying these assumptions, with $T$ also admitting an associative product with terminal unit and natural projections, for which a morphism of cofinal sequences induces an isomorphism in $\mathrm{Ind}\,\mathrm{Cat\,hot}$ and a noninvertible map in $\Hot$.

\textit{Proof.} Let $\mathrm{FI}$ be the skeleton of finite sets and injections, with objects $[n]=\{0,\ldots,n-1\}$, and set $T=\mathrm{FI}^o$. Write $\mathbf B G$ for the one-object category of a group $G$. Define a strict functor
\[
 F:T^o=\mathrm{FI}\longrightarrow\mathrm{Cat},\qquad
 F([n])=\mathbf B\mathfrak S_n.
\]
For an injection $j:[n]\hookrightarrow[m]$, send a permutation $\sigma$ to the permutation $j\sigma j^{-1}$ on $j([n])$, extended by the identity on its complement. Extension along a composite of injections agrees exactly with successive extension, so $F$ is a functor. Its Grothendieck construction is strictly fibred over $T$.

If $j,k:[n]\hookrightarrow[m]$ are parallel injections, choose $h\in\mathfrak S_m$ with $h j=k$, by extending the resulting bijection of their images to the complements. Then
$F(k)(\sigma)=hF(j)(\sigma)h^{-1}$ for every $\sigma$. Conjugation by $h$ is a natural isomorphism between the associated one-object functors. Thus the homotopy hypothesis is satisfied. The preorder $T_o$ is $\mathbf N^o$. Disjoint union provides the required associative operation and the empty set is its terminal unit in $T$; the two inclusions in $\mathrm{FI}$ become its natural projections.

Take $t'_n=[n]$ and $t_n=[n+1]$, with the opposite of the usual initial-segment inclusions as transitions. Both sequences are cofinal. The injections
$\delta_n:[n]\hookrightarrow[n+1]$, $i\mapsto i+1$, commute with those inclusions and therefore define a morphism $t\to t'$ in $\mathcal H$. The two filtered colimits of the fibres identify with $\mathbf B G$, where $G$ is the group of finite-support permutations of $\mathbf N$. The induced functor is the homomorphism
\[
 s:G\longrightarrow G,\qquad
 (s\sigma)(0)=0,\quad(s\sigma)(i+1)=\sigma(i)+1.
\]
It is not surjective, since the transposition $(0\ 1)$ is absent from its image. Consequently $Bs$ does not induce an isomorphism on fundamental groups and is not a weak equivalence.

At each finite stage, however, $\delta_n$ and the standard inclusion are parallel injections, so their induced functors are naturally isomorphic. In $\mathrm{Ind}\,\mathrm{Cat\,hot}$ the comparison is therefore the ordinary cofinal-shift isomorphism. This proves all assertions.\hfill$\square$

The earlier Chapter~V counterexample supplies the shift mechanism; the additional content here is its realization by a strict fibred category satisfying the hypotheses actually used on pages~117--123. We checked those hypotheses against the editors' complete Chapter~VII, manuscript pages~99--101, and the detection theorem of \citet{ArlinChristensen2023}; we did not find this indexed realization in those sources. No priority is claimed. The example does not satisfy the stronger full-faithfulness features of the geometric padding maps. It leaves the later all-arrows path theorem intact, while disproving the general inference that an ind-isomorphism of homotopy classes must remain an isomorphism after these ordinary colimits are formed.
\end{dcomment}


\begin{enumerate}
\item[a)] À tout $\underline{t} = (t_n)_{n \in \No} \in \Hc(\No, T)$ est associé un système inductif, indexé par $\mathbf{N}$, de catégories $C_n(\underline{t})$ $(= C_{t_n})$, et leur limite inductive
\[
C_\infty(\underline{t}).
\]
% foot: Version 22 décembre 2001

\printedpage{77}
% head: A. Grothendieck : Dérivateurs, Ch. VII, Édition des Universités de Montpellier II et Paris VII
\item[b)] Ce système inductif, et aussi sa limite est \texttt{[un]} foncteur contravariant en $\underline{t}$. Si $\underline{t}, \underline{t'} \in \mathrm{Ob}\,\Hc$, alors on peut comparer $C_\infty(\underline{t})$, $C_\infty(\underline{t'})$ et les systèmes inductifs correspondants, grâce à
\[
\begin{tikzcd}[column sep=small]
 & \underline{t} * \underline{t'} \arrow[ld] \arrow[rd] & \\
\underline{t} & & \underline{t'}.
\end{tikzcd}
\]

\item[c)] L'objet
\[
(C_n(\underline{t}))_{n \in \mathbf{N}} \in \mathrm{Ob\,Ind\,Cat\,hot}
\]
défini par le système inductif $(C_n(\underline{t}))_{n \in \mathbf{N}}$ de Cat est (à isomorphisme unique près) indépendant du choix de $\underline{t}$, soit $H_C$ ou $H$. Les homomorphismes
\[
C_\infty(\underline{t'}) \xrightarrow{u^*} C_\infty(\underline{t}) \qquad \text{associé à } u : \underline{t} \longrightarrow \underline{t'} \tag*{$(*)$}
\]
sont compatibles avec cette donnée, i.e. on a commutativité dans
\[
\begin{tikzcd}
\underset{\mathrm{Ind\,Cat\,hot}}{\varinjlim} C_*(\underline{t'}) \arrow[rr] & & \underset{\mathrm{Ind\,Cat\,hot}}{\varinjlim} C_*(\underline{t}) \\
 & H \arrow[lu, "\sim"] \arrow[ru, "\sim"'] &
\end{tikzcd}
\]

\msp{page 123}
\item[d)] Il s'ensuit que les homomorphismes $(*)$ sont des \textit{isomorphismes} dans Hot, donc tous les $C_\infty(\underline{t})$ sont isomorphes \textit{dans} Hot. L'astuce des $t * t'$ (dans le cas où on a une loi $t * t'$ associative et admettant l'objet final comme objet unité) donne même des isomorphismes \textit{canoniques}
\[
C_\infty(\underline{t'}) \longrightarrow C_\infty(\underline{t}) \qquad \text{dans Hot,}
\]
et c'est là un système transitif d'isomorphismes. Donc on a dans Hot un objet canonique
\[
C_\infty \in \mathrm{Ob}\,\Hot,
\]
dont la relation à $H \in \mathrm{Ob\,Ind\,Hot}$ est déduite de c), c'est une flèche canonique
\[
H \longrightarrow C_\infty \qquad \text{dans Ind\,Hot (\footnotemark[27]).}
\]
\footnotetext[27]{\textbf{NB} $H$ n'est pratiquement jamais essentiellement constant, donc la flèche ci-contre \texttt{[n'est]} pas un isomorphisme dans Ind\,Hot.}

Cette flèche induit un isomorphisme pour les objets correspondants de $\Hotz$; mais en regardant $H$ non comme ind-objet de $\Hotz$, mais prenant le foncteur $H_0$ induit sur $\Hot_0$ par le foncteur ind-représentable $H$ sur Hot. Ce foncteur induit est ‘représentable’ par l'objet $C_\infty$ de la catégorie plus grande
\[
H_0 \simeq C_\infty \qquad \text{dans } \Hotz.
\]

(Si le foncteur canonique $\Hot \longrightarrow \Hotz$ est pleinement fidèle, alors cette formule décrit $C_\infty$ comme objet de Hot, à isomorphisme unique près, en termes de $H_0$, donc aussi en termes de $H \in \mathrm{Ob\,Ind\,Cat\,hot}$.)
% foot: Version 22 décembre 2001

\printedpage{78}
% head: A. Grothendieck : Dérivateurs, Ch. VII, Édition des Universités de Montpellier II et Paris VII
\msp{page 124}
\item[e)] J'ignore par contre si les flèches
\[
C_\infty(\underline{t'}) \longrightarrow C_\infty(\underline{t})
\]
sont des isomorphismes dans Cat\,hot, bien qu'elles soient définies en termes d'homomorphismes de systèmes inductifs de catégories, qui sont des isomorphismes dans Ind\,Cat\,hot. Il en est ainsi, même pour les flèches $C_\infty(\underline{t}), C_\infty(\underline{t'}) \rightrightarrows C_\infty(\underline{t} * \underline{t'})$ utilisées pour comparer $C_\infty(\underline{t})$ et $C_\infty(\underline{t'})$ canoniquement. Ainsi, je ne vois pas comment définir des flèches en sens inverse dans Cat, qui seraient candidats à être des isomorphismes inverses dans Cat\,hot.

\dgrblock{dgrb388}{A colimit need not be the same ind-object}
\begin{dcomment}{examples}{p123}{A colimit need not be the same ind-object}
Grothendieck's warning in footnote~27 concerns the test objects allowed, not a defect of the colimit. It is already visible for sets. Put $D_n=\{0,\ldots,n\}$, viewed as discrete categories, and $D=\mathbf N$. For every finite discrete $K$ one has
\[
 \mathop{\rm colim}_n\mathrm{Hom}(K,D_n)=\mathrm{Hom}(K,D),
\]
because the image of a finite set is bounded. The map from the ind-object $(D_n)_n$ to the constant ind-object $D$ is nevertheless not an isomorphism. Applying maps from the constant object $\mathbf N$, its source consists of bounded-image functions $\mathbf N\to\mathbf N$, while the target also contains the identity. These computations are unchanged in $\Hot$: maps between discrete spaces are functions, and no homotopy identifies different functions.

Thus restriction to the finite-test category can identify two presheaves that differ on larger tests. The ind-object records how a family is bounded; the colimit records its pointwise values after the bound has been forgotten. This is the precise content behind ``pratiquement jamais essentiellement constant.'' It is not a universal prohibition: a sequence that is eventually constant is a constant ind-object, and more generally one can have an ind-isomorphism to a constant object without literal eventual constancy.
\end{dcomment}


\item[f)] Dans le cas géométrique ($T = \underline{\mathrm{Ty}}$, situation décrite en termes de $X$, $\Sigma$, et $a, b \in \mathrm{Ob}\,X$, d'où des $C_\infty(\underline{t}; a, b)$, $C_\infty(\overline{a, b})$, $H(a,b)$ etc.), on \textit{n'a pas} d'opération de composition de chemins au niveau des \textit{catégories} $C_\infty(\underline{t}; a, b)$ et $C_\infty(\underline{t}; b, c)$, mais seulement au niveau des objets $H(a,b)$, $H(b,c)$ de Ind\,Cat\,hot ou de Ind\,Hot. Suis-je seulement sûr de savoir définir \textit{dans} Hot
\[
C_\infty(a,b) \times C_\infty(b,c) \longrightarrow C_\infty(a,c) \ \ ?
\]

Sûrement, il n'y a pas de problèmes pour les images dans $\Hotz$, puisque là ce sont les mêmes que celles de $H(a,b)$, $H(b,c)$, $H(a,c)$ respectivement, et on a des accouplements parfaits, associatifs, avec unité, et avec inverses là ou il doit y % [sic] là ou -> là où, accent omis sur le relatif
\msp{page 125}
en avoir. Mais pas pour les $C_\infty(a,b)$ eux-mêmes?
\end{enumerate}

C'est un peu vexant qu'avec tout ce travail il semblerait qu'on n'ait pas même été foutu de définir l'analogue dans Cat du classique espace des lacets $\Omega(X, x)$ comme un groupe \textit{dans} Hot. (Que fait \textsc{Thomason}?) Mais voilà une idée pensable [?]. Pour définir \textit{quand-même}
\[
C_\infty(a,b) \times C_\infty(b,c) \longrightarrow C_\infty(a,c)
\]
dans Hot, on va décrire les deux foncteurs $C_\infty(a,b)$ et $C_\infty(b,c)$ par des systèmes projectifs différents $\underline{t}$, $\underline{t'}$. On peut l'expliciter en partant d'un $\tau_0$ de longueur utile $\geq 2$, et en prenant dans les deux cas
\[
t_n = t'_n = \tau_0^n,
\]
mais des morphismes de transition différents
\[
\begin{aligned}
(t_{n+p} \longrightarrow t_n) &= \pi_p * \mathrm{id}_{\tau_0^n} \\
(t'_{n+p} \longrightarrow t'_n) &= \mathrm{id}_{\tau_0^n} * \pi_p,
\end{aligned}
\]
où $\pi_p : \tau_0^p \longrightarrow 1$ est la dégénérescence totale. On définira $C_\infty(a,c)$ en termes de $\underline{t''} = \underline{t} * \underline{t'}$, donc
\[
\left\{
\begin{aligned}
t''_n &= \tau_0^{2n} \\
(t''_{n+p} \longrightarrow t''_n) &= \pi_p * \mathrm{id}_{\tau_0^{2n}} * \pi_p.
\end{aligned}
\right.
\]
% foot: Version 22 décembre 2001

% ---- from Der-ChVII-456_p081-085.tex
\printedpage{79}
% head: A. Grothendieck : Dérivateurs, Ch. VII, Édition des Universités de Montpellier II et Paris VII
% foot: Version 22 décembre 2001

Cette fois, les foncteurs de composition
\[ C_n(\underline{t};a,b)\times C_n(\underline{t'};b,c) \longrightarrow C_n(\underline{t''};a,c), \]
sont compatibles avec les foncteurs transition pour $n$ variable, et on trouve
\msp{page 126}
\[ C_\infty(\underline{t};a,b)\times C_\infty(\underline{t'};b,c) \longrightarrow C_\infty(\underline{t''};a,c), \]
donc une flèche dans Hot
\[ \comp_{\tau_0} : C_\infty\times C_\infty \longrightarrow C_\infty. \]
Il faut voir qu'elle ne dépend pas du choix de $\tau_0$, mais si on a $\tau_0$ et $\tau'_0$, on les compare par $\tau_0\times\tau'_0 \begin{array}{l}\longrightarrow \tau_0\\ \longrightarrow \tau'_0\end{array}$, il suffit donc de voir que s'il existe une flèche $\tau_0\longrightarrow\tau'_0$, alors $\comp_{\tau_0} = \comp_{\tau'_0}$. Mais une telle flèche induit des flèches $\underline{t}_{\tau_0}\longrightarrow\underline{t}_{\tau'_0}$, $\underline{t'}_{\tau_0}\longrightarrow\underline{t'}_{\tau'_0}$, $\underline{t''}_{\tau_0}\longrightarrow\underline{t''}_{\tau'_0}$, d'où un diagramme dans Cat
\[
\begin{tikzcd}
C_\infty(\underline{t}_{\tau'_0};a,b)\times C_\infty(\underline{t'}_{\tau'_0};b,c) \arrow[r, "{\comp_{\tau'_0}}"] \arrow[d] & C_\infty(\underline{t''}_{\tau'_0};a,c) \arrow[d] \\
C_\infty(\underline{t}_{\tau_0};a,b)\times C_\infty(\underline{t'}_{\tau_0};b,c) \arrow[r, "{\comp_{\tau_0}}"] & C_\infty(\underline{t''}_{\tau_0};a,c),
\end{tikzcd}
\]
dont on vérifie tout de suite qu'il est commutatif. Or ce diagramme peut se réécrire dans Hot comme
\[
\begin{tikzcd}
C_\infty(a,b)\times C_\infty(b,c) \arrow[r, "{\comp_{\tau'_0}}"] \arrow[d, "{\mathrm{id}}"] & C_\infty(a,c) \arrow[d, "{\mathrm{id}}"] \\
C_\infty(a,b)\times C_\infty(b,c) \arrow[r, "{\comp_{\tau_0}}"] & C_\infty(a,c),
\end{tikzcd}
\]
\msp{page 127}
d'où $\comp_{\tau_0} = \comp_{\tau'_0}$.

\dgrblock{dgrb389}{Padding at the two outer ends}
\begin{dcomment}{exposition}{p126}{Padding at the two outer ends}
Grothendieck's construction now solves the binary-composition problem by choosing the locations where new constant segments are inserted. A common stage in each factor is harmless because the diagonal $\mathbf N\to\mathbf N\times\mathbf N$ is cofinal. The location of the padding is what makes the squares commute.

Write a path in the first system as $c_n$ and one in the second as $d_n$. Under $n\mapsto n+p$, the first becomes a string of $p$ constant blocks followed by $c_n$, while the second becomes $d_n$ followed by $p$ constant blocks. Concatenation gives
\[
 (1_a^{,p}*c_n)*(d_n*1_c^{,p})
       =1_a^{,p}*(c_n*d_n)*1_c^{,p},
\]
which is exactly the transition prescribed for the target system. The equality holds for the vertical transformations as well as for objects, because concatenation is a bifunctor. It therefore induces a functor on the filtered colimits. Products of categories commute with these filtered colimits, again because objects, arrows, and their defining equations are finite data.

A morphism $\tau_0\to\tau'_0$ gives a commutative square before passage to homotopy; once the refinement maps have been established to be weak equivalences and coherently identified, the square gives equality of the two binary operations in $\Hot$. To compare unrelated types, one can use $\tau_0*\tau'_0$ and its two collapse maps. The printed $\tau_0\times\tau'_0$ on this page cannot literally mean their categorical product as interval types: the product of two nontrivial intervals generally has a square rather than a linear underlying shape. Concatenation supplies the common refinement required by the argument.

The remaining associativity issue is real at the level of this proof. A third path has two joining positions, and the two binary constructions use different padding conventions at those positions. One must compare the resulting maps with coherent reparametrizations, or compare with a single strictly associative all-types model. Equality after restriction to finite tests cannot be substituted for that comparison. In the all-arrows case, the comparison established in section~13 is precisely the route back to the strict concatenation which motivated the discussion on page~111.
\end{dcomment}


Je recule devant la vraie difficulté - celle de prouver l'associativité de cette loi de composition. Elle devient associative dans $\Hot^{\wedge}_{0}$ - mais pour que deux flèches dans Hot soient égales (ici donc
\[ C_\infty(a,b)\times C_\infty(b,c)\times C_\infty(c,d) \rightrightarrows C_\infty(a,d)\quad ), \]
suffit-il qu'il en soit ainsi de leurs images dans $\Hot^{\wedge}_{0}$ - en d'autres termes, le foncteur
\[ \Hot\longrightarrow\Hot^{\wedge}_{0}, \]
est-il (en plus d'être conservatif) au moins fidèle (sinon pleinement fidèle)?

\clearpage
\printedpage{80}
% head: A. Grothendieck : Dérivateurs, Ch. VII, Édition des Universités de Montpellier II et Paris VII
% foot: Version 22 décembre 2001

Il y a aussi que cette loi de composition est compatible avec celle des $H(a,b)$ dans Ind\,Hot, pour les flèches canoniques
\[ H(a,b)\longrightarrow C_\infty(a,b) \qquad \text{dans Ind\,Hot}, \]
si on saurait que ce sont des épimorphismes, et le restent par produits cartésiens, on aurait gagné.

Bien sûr, on saura que c'est associatif, une fois qu'on saura que la composititon définie ici est compatible avec la composition des chemins dans le cadre semi-simplicial. Mais comment fait-on dans ce cadre, pour composer les chemins de longueur non bornée? N'y a-t-il pas là exactement la même difficulté? À moins de s'en % [sic] composititon -> composition, misspelling % Reading resolved against the supplied PDF; see corrections.csv.
\msp{page 128}
tirer en se ramenant tout de suite, pour la théorie des opérations $\Omega$, aux complexes de Kan, et aux chemins de type $\Delta_+$. Mais pour les composer \textit{fonctoriellement,} il me semble y avoir toujours le même problème - on trouvera des chemins de type $\Delta_+\Delta_+ = \boldsymbol{\Delta}_2^+$, et non pas $\Delta_+$, mais on peut alors composer les homotopies directes (revient à s'induire sur $\Delta_+\simeq\{0,2\}\subseteq\boldsymbol{\Delta}_2^+$). Or j'ai fait le formalisme des catégories de chemins sans me soucier de conditions du type Kan. Je présume que dans le cadre semi-simplicial, les $\uHom(\boldsymbol{\Delta}_1,X)$ ne donnent \textit{pas} le bon objet des chemins, pour $X$ pas de Kan. (Car sinon, il semblerait que pour $X=X_0$ à composantes \textit{finies,} le calcul des $\pi_i$ devrait être fini, ce qui n'est pourtant pas le cas, m'a-t-on toujours assuré \dots)

\dgrblock{dgrb390}{A map invisible at every finite stage}
\begin{dcomment}{examples}{p127}{A map invisible at every finite stage}
Grothendieck asks whether finite tests determine a morphism and whether the ind-object maps epimorphically onto its colimit. A small telescope gives a negative answer to both questions. We give the calculation because it exhibits the obstruction without relying on an unspecified phantom-map example.

Let $T$ be the mapping telescope of the degree-two maps
\[
 S^1\xrightarrow{2}S^1\xrightarrow{2}S^1\xrightarrow{2}\cdots.
\]
Its finite subtelescopes $T_n$ retract onto circles. Its fundamental group is $\mathbf Z[1/2]$, and its higher homotopy groups vanish: a map or nullhomotopy from a compact sphere lies in a finite subtelescope. Thus $T$ is a $K(\mathbf Z[1/2],1)$. After contracting the ray of vertices, a cellular model has one vertex, one loop $e_n$ for each $n\geq0$, and one two-cell with boundary $e_n-2e_{n+1}$. Cellular cohomology consequently gives
\[
 H^2(T;\mathbf Z)=\operatorname{coker}\left(
 \prod_{n\geq0}\mathbf Z\longrightarrow\prod_{n\geq0}\mathbf Z,
 \quad(b_n)\longmapsto(b_n-2b_{n+1})\right).
\]
This is the standard telescope computation of \citet[Examples~3F.7 and~3F.9]{Hatcher2002}.

To see nonvanishing directly, take $a_{2j}=1$ and $a_{2j+1}=0$. If $a_n=b_n-2b_{n+1}$, iteration would give
\[
 b_0=\sum_{j=0}^{N}2^ja_j+2^{N+1}b_{N+1}
 \qquad(N\geq0).
\]
In the ring of $2$-adic integers, this forces
$b_0=1+4+16+\cdots=-1/3$, impossible for an integer $b_0$. Hence $a$ determines a nonzero cohomology class. By the representation of integral degree-two cohomology, it determines an essential map
\[
 f:T\longrightarrow K(\mathbf Z,2)
\]
\citep[Theorem~4.57]{Hatcher2002}. Each restriction to $T_n$ is nullhomotopic, since $H^2(T_n;\mathbf Z)=0$. Every map from a finite CW complex into $T$ has image in a finite subtelescope up to homotopy, so all finite homotopy tests identify $f$ with the constant map.

Even if the test category includes nerves of all categories with finitely many objects and arrows, this example remains invisible. Each connected component of such a nerve has finitely generated fundamental group. Any homomorphism from that group to $\mathbf Z[1/2]$ has cyclic image, because finitely many denominators have a common power of two. The associated map to $T=K(\mathbf Z[1/2],1)$ therefore factors up to homotopy through $K(\mathbf Z,1)=S^1$, or through a point. Its composite with $f$ is nullhomotopic because $[S^1,K(\mathbf Z,2)]=0$.

Finally let $H=(T_n)_n$ in $\mathrm{Ind}\,\Hot$. The two maps $f,0:T\rightrightarrows K(\mathbf Z,2)$ become equal after precomposition with $H\to T$, since
\[
 \mathrm{Hom}_{\mathrm{Ind}\,\Hot}(H,K(\mathbf Z,2))
      =\mathop{\rm lim}_n[T_n,K(\mathbf Z,2)]=\{0\}.
\]
They are distinct before precomposition. Hence $H\to T$ is not an epimorphism. This example refutes the proposed general inference; it does not contradict associativity of the particular path composition, which requires its own geometric proof.
\end{dcomment}


Je présume pourtant, quand on prend $\Sigma=\mathrm{Fl}(X)$, que les objets $\Ch_\infty(X;a,b)$ et $\Ch_\infty(\bar{X};a,b)$ sont isomorphes dans Hot, quand on immerge $X$ dans une enveloppe de Thomason-Kan $\bar{X}$.

\clearpage
\printedpage{81}
% head: A. Grothendieck : Dérivateurs, Ch. VII, Édition des Universités de Montpellier II et Paris VII
% foot: Version 22 décembre 2001

\msp{page 129}

\dgrblock{dgrb391}{Edges and the fibres of the path object}
\begin{dcomment}{exposition}{p128}{Edges and the fibres of the path object}
Grothendieck wants the categorical construction to work before fibrant replacement. The distinction needed here concerns the endpoint fibres. For a Kan complex $K$, the map $K^{\Delta^1}\to K\times K$ is a Kan fibration, and its fibre over two vertices models their space of paths. For an arbitrary simplicial set, the same strict fibre need not compute a homotopy fibre. This does not mean that the total space $K^{\Delta^1}$ has the wrong homotopy type: evaluation at an endpoint has the constant-path section and the simplicial interval contraction gives a homotopy from their composite to the identity.

A finite example makes the endpoint problem explicit. Let $X$ be the category with objects $0,1,2$ and nonidentity arrows $0\to1\leftarrow2$. Its classifying space is contractible, so the topological space of paths from $0$ to $2$ is nonempty and contractible. But there is no morphism $0\to2$. The fibre of the functor category $\underline{\mathrm{Hom}}([1],X)\to X\times X$ over $(0,2)$ is therefore empty. The two-edge zigzag already supplies the missing path, and the refinement formalism is designed to supply all the further homotopies between such representatives.

There is a related distinction in the discussion of composition. Two composable edges of a simplicial set specify the spine $\Delta^1\cup_{\Delta^0}\Delta^1$, not an entire $\Delta^2$. For a Kan complex a filler produces a third edge, with higher fillers controlling the dependence on choices. In an ordinary category the interval $[2]$ already includes the composite arrow; this explains why the categorical notation for two direct segments can conceal the simplicial filling step. A simplicial loop group, a Moore-path model, or a coherent localization retains the data needed to compare repeated compositions. None requires an algorithm computing all homotopy groups of a finite simplicial set.

The last conjecture on this page is correct for the all-arrows path object: a weak equivalence of categories induces weak equivalences on the corresponding derived path spaces, with their endpoints carried along. The justification for identifying Grothendieck's specific categories with those spaces belongs to the geometric argument in sections~12--13, not to a calculation with one-edge paths.
\end{dcomment}


\section*{11\quad Rapide retour sur les $(\mathcal{M},\Sigma)$, dans l'optique des dérivateurs}

\dgrblock{dgrb392}{What an external Hom must represent}
\begin{dcomment}{exposition}{p129:Rapide retour sur les}{What an external Hom must represent}
Grothendieck now asks what his paths would have to accomplish in order to serve a theory of derived diagrams. A path from $a$ to $b$ is no longer only a point of a topological mapping space. It must also describe a morphism between the constant diagrams $a_Y$ and $b_Y$, after objectwise localization of $Y$-diagrams. This is the reason for testing the path construction against every parameter category $Y$.

Write
\[
 \mathbb D(Y)=\operatorname{Fun}(Y,\mathcal M)[\Sigma_Y^{-1}],
 \qquad
 F_{a,b}(Y)=\operatorname{Hom}_{\mathbb D(Y)}(a_Y,b_Y).
\]
The conventions on variance can be reversed together without changing this discussion. The first displayed condition says precisely that $F_{a,b}$ factors through the localization of categories at nerve weak equivalences. It is a necessary invariance property for the proposed representation by a space. It is not by itself a definition of a derivator: existence and pointwise computation of homotopy Kan extensions are additional requirements. For combinatorial model categories these requirements hold by \citet[Proposition~1.30]{Groth2013}. For their locally constant diagrams, the invariance asserted here is a theorem of \citet[Theorem~1.3]{Cisinski2009LocallyConstant}; constant diagrams are particular locally constant diagrams, and the subcategory in that theorem is full.

The contemporary form of the intended answer is especially transparent when a relative category presents an $\infty$-category $\mathcal C$ and its localized strict diagram categories present $\operatorname{Fun}(NY,\mathcal C)$. With $K=\operatorname{Map}_{\mathcal C}(a,b)$ one then has
\[
 F_{a,b}(Y)\cong
 \pi_0\operatorname{Map}_{\mathcal S}(|NY|,K).
\]
Indeed, a coherent transformation between the two constant diagrams is a coherent $Y$-indexed family of points of $K$; this is the limit of the constant $K$-diagram, hence the mapping space from the groupoid completion $|NY|$ to $K$. Taking components gives the formula. The diagram-presentation hypothesis is substantive: one cannot replace it by the existence of an arbitrary admissible set $\Sigma$. The comparison between strict model-category diagrams and coherent diagrams is established, under the relevant model-category hypotheses, in \citet[Proposition~4.2.4.4]{Lurie2009HTT}.

Here the source means finite sets of objects and arrows by a finite category \citep[ms.~45]{Grothendieck1990V}; such a category need not have a finite nerve. Consequently the earlier compactness argument does not establish the comparison for every finite category. The conditional assertion under the additional locality of $H$ is not disproved merely by that observation.

This also clarifies the finite-stage issue. The ind-object $H(a,b)$ records a bound on the length of a zigzag that works simultaneously over the entire parameter. Passing to $\Ch_\infty$ can forget this uniformity. It does not follow that modern derived mapping spaces must themselves be ind-objects: in the good diagram presentations just described, one space $K$ represents the functor for every small $Y$. What fails for a general pair $(\mathcal M,\Sigma)$ may be the identification of its localized strict diagrams with coherent diagrams, rather than representability of the latter. The example below separates these two issues with a six-object category.
\end{dcomment}


Quand on s'intéresse au formalisme des $\Ch_\Sigma(X;a,b)$ en vue d'un \textit{dérivateur} $\mathbf{D}_{X,\Sigma}$ qu'on voudrait définir en termes de $X$, $\Sigma$, il semblerait que le prédérivateur qu'on obtient bel et bien ne mérite le nom de dérivateur que lorsqu'on suppose que pour toute flèche
\[ Y\longrightarrow Y' \]
dans Cat qui est une Hot-équivalence, les flèches correspondantes
\[ \mathrm{Hom}_{\mathrm{Ind\,Cat\,hot}}(Y',H(a,b)) \longrightarrow \mathrm{Hom}_{\mathrm{Ind\,Cat\,hot}}(Y,H(a,b)), \]
ou encore
\[ \varinjlim_n \mathrm{Hom}_{\mathrm{Cat\,hot}}(Y',\Ch_{n,\Sigma}(a,b)) \longrightarrow \varinjlim_n \mathrm{Hom}_{\mathrm{Cat\,hot}}(Y,\Ch_{n,\Sigma}(a,b)) \]
sont des bijections. J'ai regardé dans V ce que ça signifie pour l'objet $H(a,b) = (\Ch_{n,\Sigma}(a,b))_n$ de Ind\,Cat\,hot. C'est cette condition seulement qui assure, du point de vue des dérivateurs, une interprétation convaincante de $\Ch_{\infty,\Sigma}(a,b)$,
\[ \Ch_{\infty,\Sigma}(a,b) \;=\; \varinjlim_n \Ch_{n,\Sigma}(a,b) \]
en tant qu'objet de Hot, ou plus exactement de $\Hot^{\wedge}_{0}$, car on a alors pour $Y$ une catégorie \textit{finie}
\[ \underbrace{\mathrm{Hom}_{\mathrm{Cat\,hot}}(Y,H(a,b))}_{\simeq\,\mathrm{Hom}_{X(Y)\Sigma_Y^{-1}}(a_X,b_X)} \;\simeq\; \mathrm{Hom}_{\Hot}(Y,\Ch_{\infty,\Sigma}(a,b)), \]
\msp{page 130}
car je veux quand-même que le type dans Hot soit le bon, qu'il correspond bien à l'homotopie des espaces de chemins. Donc on voudrait qu'un morphisme
\[ X\xrightarrow{\;f\;}X' \]
dans (Cat), qui est une Hot-équivalence, définisse
\[ \Ch_\infty(X;a,b) \longrightarrow \Ch_\infty(X';f(a),f(b)), \]
qui soient eux aussi des Hot-équivalences. D'autre part, c'est évidemment compatible avec la brillante composition des chemins, que j'ai réussi à définir entre ces objets de Hot. Donc (sous réserve d'avoir vérifié ce que je viens de dire) la vérification de l'associativité se ramenerait au cas où $X$ est de Thomason-Kan. Mais dans ce cas, il y a des chances que, tout comme dans le cas semi-simplicial, les $\Ch_\infty(X;a,b)$ peuvent se définir (comme objets de (Hot)) comme étant simplement les chemins de type $\Delta_+$ de $a$ à $b$, qu'on compose alors de la façon triviale. Mais non, ces catégories sont discrètes! % [sic] ramenerait -> ramènerait, missing accent

Je renonce pour l'instant à poursuivre plus la question d'associativité. Demander à \textsc{Thomason} pour les propriétés du foncteur $\Hot\longrightarrow\Hot^{\wedge}_{0}$.

\clearpage
\printedpage{82}
% head: A. Grothendieck : Dérivateurs, Ch. VII, Édition des Universités de Montpellier II et Paris VII
% foot: Version 22 décembre 2001

\msp{page 131}
Cela signifie, à très peu de choses près, qu'on a une interprétation de $\Ch_{\infty,\Sigma}(a,b)$, en tant qu'objet de Hot, comme un `Hom externe' $\mathrm{Hom\,ext}(a,b)$ à valeurs dans Hot, défini sur la catégorie $\mathcal{M}=X$ à l'aide du seul ensemble admissible de flèches $\Sigma\subseteq\mathrm{Fl}(\mathcal{M})$, opération satisfaisant la condition
\[ \boxed{\mathrm{Hom}_{\Hot}(Y,\mathrm{Hom\,ext}(a,b)) \;\simeq\; \mathrm{Hom}_{\mathcal{M}(Y)\Sigma_Y^{-1}}(a_X,b_X)} \tag*{$(*)$} \]
(qui peut s'interpréter aussi, soit comme
\[ \mathrm{Hom}_{\mathcal{M}}(p_!(a_Y),b) \]
si $p_! : \underbrace{\mathcal{M}(Y)\Sigma_Y^{-1}}_{\mathbf{D}(Y)} \longrightarrow \mathcal{M}$ existe, soit comme
\[ \mathrm{Hom}_{\mathcal{M}}(a,p_*(b_Y)) \]
si $p_* : \underbrace{\mathcal{M}(Y)\Sigma_Y^{-1}}_{\mathbf{D}(Y)} \longrightarrow \underbrace{\mathcal{M}}_{=\,\mathbf{D}(*)}$ existe); avec le seul petit ennui que cette formule n'est valable que pour $Y$ \textit{finie.} Si on veut exprimer le second membre pour $Y$ variable quelconque dans Cat, il ne suffit pas d'un objet de Hot ($\mathrm{Hom\,ext}(a,b)$) dans le premier membre, mais il faut l'ind-objet $H(a,b) = (\Ch_n(a,b))_{n\in\mathbf{N}}$ dans Hot. Si donc on tient à travailler dans un `domaine' Diag du dérivateur $\mathbf{D}=\mathbf{D}_{\mathcal{M},\Sigma}$ aussi grand que possible, en se bornant à des catégories finies,
\msp{page 132}
et qu'on veut qu'une formule $(*)$ reste valable pour toute $Y$ dans le domaine, il faut `en toute rigueur' définir le $\mathrm{Hom\,ext}(a,b)$ non comme objet de Hot, mais de Ind\,Hot. Il se décrirait donc alors comme
\[ \mathrm{Hom\,ext}(a,b) = (\Ch_{n,\Sigma}(a,b))_n \qquad \text{(ind-objet dans Hot)}\ , \]
et l'objet $\Ch_{\infty,\Sigma}$ dans Hot en est une sorte de limite inductive.

\dgrblock{dgrb393}{The original category and its localized value}
\begin{dcomment}{context}{p131}{The original category and its localized value}
Grothendieck wants the same external Hom to explain both homology and cohomology with constant coefficients. This is an adjunction calculation once the ambient categories have been fixed. Put $p:Y\to e$. If the relevant adjoints exist, then
\[
\begin{aligned}
 \operatorname{Hom}_{\mathbb D(Y)}(p^*a,p^*b)
 &\cong\operatorname{Hom}_{\mathbb D(e)}(p_!p^*a,b)\\
 &\cong\operatorname{Hom}_{\mathbb D(e)}(a,p_*p^*b).
\end{aligned}
\]
Here $\mathbb D(e)=\mathcal M[\Sigma^{-1}]$, not the original $\mathcal M$ unless the notation has been changed. Likewise the $a_X,b_X$ of the displayed formula mean $a_Y,b_Y$. These small notational adjustments matter because an adjoint to restriction of localized diagrams is a derived operation; an adjoint before localization need not compute it.

The historical motivation is unusually explicit. In the retrospective in \emph{Pursuing Stacks}, \citet[\S69, pp.~1--2]{Grothendieck1983} relates the inadequacy of an isolated derived category to the trace questions arising from SGA~5, and proposes retaining derived categories of all diagram shapes. The present pages carry that programme into a concrete question: can the dependence on the parameter shape itself be represented homotopically? The adjunction formulas explain why a positive answer belongs simultaneously to the homological and the cohomological sides of the theory. The subsequent calculations concern the geometric adequacy of the proposed representative, rather than another formal manipulation of these adjunctions.
\end{dcomment}


Je me suis demandé si la condition préalable, savoir que les foncteurs $H(a,b)$ sur Cat transforment Hot-équivalences $Y'\longrightarrow Y$ en bijections, était bel et bien une (forte) restriction. Je me suis convaincu que c'est bien le cas, déjà dans le cas particulier où $\Sigma=\mathrm{Fl}(X)$. Il me semble que lorsque les composantes connexes de $\Omega(X,a)$ ne sont pas contractiles, ladite condition n'est pratiquement jamais satisfaite. J'ai pensé p.~ex. au cas où $X$ a le type d'homotopie (dans Hot) d'une 2-sphère, de sorte que son $\Omega(X,a)=\Ch_\infty(X,a,a)$ doit avoir un $\pi_1\simeq\mathbf{Z}$. Prenons $Y$ ayant le type dans Hot d'un cercle, disons
\[
\begin{tikzcd}
 & & a & & \\
Y\colon & 0 \arrow[ur] \arrow[dr] & & 0' \arrow[ul] \arrow[dl] & , \\
 & & b & &
\end{tikzcd}
\]

\clearpage
\printedpage{83}
% head: A. Grothendieck : Dérivateurs, Ch. VII, Édition des Universités de Montpellier II et Paris VII
% foot: Version 22 décembre 2001

\msp{page 133}
donc les foncteurs
\[ f : Y\longrightarrow \Ch_\infty(X;\bullet,\bullet) \overset{\text{déf}}{=} \Omega(X,\bullet) \]
correspondent à \textit{certains} `lacets' de $\Omega(X,\bullet)$ (mais lacets à origine non précisée), ils ont un type d'homotopie au sens topologique (`nombre de tours' autour de la sphère $X$) donné par un entier $n\in\mathbf{Z}$ (le `degré' du lacet $\lambda$). L'intuition évidente, c'est que l'ensemble des degrés $\delta$ qu'on peut obtenir ainsi, avec une catégorie \textit{finie} donnée $Y$ qui représente le cercle $S^1\in(\Hot)$, est fini - $|\delta|$ majoré sûrement par le nombre de ces flèches (donc au moins par 4 dans le cas envisagé, chaque flèche de $Y$ faisant faire \textit{au plus} un tour autour du pseudo-cercle $\Omega(S^2,\bullet)$) (\footnote[28]{probablement on ne trouve tout au plus que les degrés $\delta$ entre $-2$ et $+2$.}). D'autre part, quand on représente $S^1$ par des catégories $Y_n$ arbitrairement grandes (disons déduites des intervalles $\Phi_n=(\Delta_+\Delta_-)^n$ en identifiant la source et le but), on doit trouver des degrés arbitrairement élevés. Or toute $Y_n$ s'envoie dans $Y(=Y_2)$ par une flèche
\[ Y_n = Y' \longrightarrow Y = Y_2 \]
\msp{page 134}
qui est de degré 1, donc qui est une Hot-équivalence. Mais l'application correspondante
\[ \pi_0\uHom(Y,\Omega(X,\bullet)) \longrightarrow \pi_0\uHom(Y',\Omega(X,\bullet)) \]
n'est pas bijective. Ce genre d'argument me semble devoir marcher chaque fois que l'ensemble des groupes d'homotopie $\pi_j(X,x)$, pour $j\geq 2$ (donc l'ensemble des groupes d'homotopie $\pi_i$ de $\Omega(X,x)$, pour $i\geq 1$) est \textit{infini,} en prenant pour $Y$, $Y'$ différents représentants $Y_n$ de la sphère $S^i$, disons les $Y_n$ déduites de $\Phi_n$ $(n\geq 2)$ en prenant $\Phi_n^i = \Phi_n\times\Phi_n\times\cdots\times\Phi_n$ (qui est une sorte de $i$-cube combinatoire), et en identifiant son `bord' à un seul point. Je présume que le nombre d'élements dans $\pi_i(\Omega(X,x))$ qu'on peut obtenir à partir de flèches $Y_n\longrightarrow\Omega(X,x)$ est limité par la taille de $Y_n$, mais on peut sûrement les obtenir chacun, en prenant $n$ assez grand. Si donc $\xi\in\pi_i(\Omega)$ échappe à $Y_{n_0}$ ($n_0=0$, disons), on prend un $n'$ assez grand pour que $\xi$ n'échappe pas à $Y_{n'}$, et on choisit % [sic] d'élements -> d'éléments, missing accent
\[ Y_{n'} \longrightarrow Y_n\ , \]
\texttt{[}$n=n_0$\texttt{?]}

\msp{page 135}
qui soit une Hot-équivalence (p.~ex. à l'aide d'une rétraction de $\Phi_{n'}$ sur $\Phi_n$, compatible avec sources et buts).

\dgrblock{dgrb394}{A fixed circle already sees every degree}
\begin{dcomment}{development}{p134}{A fixed circle already sees every degree}
Grothendieck's intuition that localization of strict diagrams can depend on the chosen category presenting a homotopy type is correct. The proposed explanation by a bound on the degree of a four-arrow circle is not. A natural transformation between long paths contains an arbitrarily wide commuting ladder. Its contribution to a two-dimensional homotopy class is not bounded by the number of arrows in the parameter category. We give a finite example that both disproves the suggested bound and proves the desired failure of invariance.

\textbf{A suspension calculation.}\ours\ Let $Y$ be the four-object poset with minimal objects $r_0,r_1$, maximal objects $q_0,q_1$, and all four relations $r_i<q_j$. Let $X$ be obtained by adjoining two incomparable objects $s,n$, each greater than every object of $Y$. Take all arrows of $X$ as weak equivalences. Then $|NY|$ is a circle and $|NX|$ is its suspension, hence a $2$-sphere. For the constant diagram $s_Y$ the group
\[
 G_Y=\operatorname{Aut}_{\Pi\operatorname{Fun}(Y,X)}(s_Y)
\]
contains an element of infinite order. Moreover, for one fixed four-arrow $Y$, functors $Y\to\Ch_\infty(X;s,s)$ realize every integer in $\pi_1\Omega|NX|$.

\emph{Proof.} Denote the inclusion by $j:Y\hookrightarrow X$, and choose a vertex $y_0\in Y$. Since all objects of $Y$ lie below both suspension vertices, there is a loop of natural transformations
\[
 s_Y\ \longleftarrow\ j\ \longrightarrow\ n_Y
       \ \longleftarrow\ (y_0)_Y\ \longrightarrow\ s_Y.       \tag{B1}
\]
This is a loop in the nerve of the functor category, based at $s_Y$. Under evaluation it gives a continuous family of loops in $|NX|$: at $y$ the loop travels from $s$ to $n$ along the meridian through $y$, and returns along the fixed meridian through $y_0$.

To check its class, let $\gamma$ be an oriented fundamental $1$-cycle of $|NY|$. The two varying portions of the family sweep out the two cones on $\gamma$. The fixed return meridian contributes no $2$-chain. Thus the image of the fundamental class of $|NY|\times S^1$ is, with a choice of orientation,
\[
                    n*\gamma-s*\gamma,
\]
which is the fundamental cycle of the suspension $|NX|$. In particular its degree is $1$. Equivalently, collapse the fixed meridian in the suspension: the adjoint of the family becomes the usual suspension map of the identity of $|NY|$. The collapse of that contractible arc induces an isomorphism on $H_2$, so this gives the same calculation. Consequently (B1) is nontrivial and its $k$-fold power has degree $k$. A null-homotopy in the functor category would give a null-homotopy after evaluation and degree zero. All these powers are therefore distinct in $G_Y$.

The same construction takes place in one finite path stage. For
$\rho=\Delta_-\Delta_+\Delta_-\Delta_+$ define
\[
 F(y)=\bigl(s\longleftarrow y\longrightarrow n
                    \longleftarrow y_0\longrightarrow s\bigr).
\]
For an arrow $y\to y'$ use that arrow in the varying position and identities in the other four positions. All squares commute because $X$ is a poset. This defines a functor $F:Y\to\Ch_\rho(X;s,s)$. Concatenating $k$ copies pointwise gives a functor into $\Ch_{\rho^k}(X;s,s)$ of degree $k$. Insert identity arrows at the two ends to fit this alternating word into a sufficiently long symmetric stage $t_N$; this changes only parametrization, not the degree. Reversal gives negative $k$, and the constant family gives zero. This proves the assertion for $\Ch_\infty$, without assuming that its concatenation has already been made into a globally defined bifunctor.\hfill$\square$

The example also supplies a genuine change-of-parameter failure. Write $B\mathbb Z$ for the one-object category with automorphism group $\mathbb Z$. There is a functor $u:Y\to B\mathbb Z$ sending one of the four edges to $1\in\mathbb Z$ and the other three to $0$. Choose the sign so that the induced map on the circle has degree $1$. Its nerve is a weak equivalence: both classifying spaces are $K(\mathbb Z,1)$ and the induced map on fundamental groups is an isomorphism. Every functor $B\mathbb Z\to X$ is constant, since a poset has only identity automorphisms; every natural transformation between such functors is exactly one arrow of $X$. Therefore
\[
 \operatorname{Fun}(B\mathbb Z,X)\cong X,
 \qquad
 \operatorname{Aut}_{\Pi\operatorname{Fun}(B\mathbb Z,X)}(s)=1.
\]
Restriction along $u$ maps this trivial group to $G_Y$, which contains all the powers of (B1). It is not surjective. By the localization interpretation of $H(s,s)$ used on ms.~129, this is precisely a failure of the required invariance of $H(s,s)$.

We checked the relevant diagram-localization and fibrancy results in \citet{Cisinski2009LocallyConstant}, \citet{Thomason1980}, and \citet{Raptis2010}. We did not find this finite suspension calculation presented there. It is included as a direct completion of Grothendieck's proposed example, without a claim of priority. The failure proved here concerns localized \emph{strict} diagrams in this relative category. It is consistent with the homotopy invariance of coherent diagrams explained above.
\end{dcomment}


Ceci attire fortement l'attention sur le fait que si $Y$, $C$ sont deux catégories, la catégorie $\uHom(Y,C)$ \textit{n'est pas} un représentant de $\uHom_{\Hot}(\mathrm{hot}(Y),\mathrm{hot}(C))$, on a seulement une flèche canonique
\[ \mathrm{hot}(\uHom(Y,C)) \longrightarrow \uHom_{\Hot}(\mathrm{hot}(Y),\mathrm{hot}(C)), \tag*{$(*)$} \]

% ---- from Der-ChVII-456_p086-090.tex
\printedpage{84}
% head: A. Grothendieck : Dérivateurs, Ch. VII, Édition des Universités de Montpellier II et Paris VII
% foot: Version 22 décembre 2001

i.e.
\[
\underbrace{\hot(\uHom(Y,C))\times\hot(Y)}_{\hot(\uHom(Y,C)\times Y)}\longrightarrow\hot(C),
\]
laquelle flèche est simplement déduite de
\[
\uHom(Y,C)\times Y\longrightarrow C.
\]
Mais la flèche $(*)$ n'a aucune envie d'être \texttt{[un]} isomorphisme. C'est parce qu'elle ne l'est déjà pas dans le cas particulier $Y=\mathbf{\Delta}^1$, essentiellement, qui nous oblige, pour faire une théorie des chemins dans les catégories qui soit topologiquement raisonnable, de ne pas nous borner à $\mathbf{\Delta}^1$ comme type de chemins, mais à prendre des intervalles composés $\tau\in\mathrm{Ty}$.

\dgrblock{dgrb395}{The absolute arrow category and its endpoint fibres}
\begin{dcomment}{examples}{p135:Mais la flèche $(*)$ n'a aucune envie}{The absolute arrow category and its endpoint fibres}
Grothendieck is trying to isolate the defect that forced him to lengthen paths. The defect is real, but the specific absolute comparison on this page does not exhibit it. For every category $C$, put $A=\operatorname{Fun}([1],C)$, let $c:C\to A$ send an object to its identity arrow, and let $e_0:A\to C$ evaluate at the source. Then $e_0c=\mathrm{id}_C$. For an arrow $f:x\to y$, the square with vertical arrows $\mathrm{id}_x$ and $f$ defines a natural transformation
\[
                    ce_0\Longrightarrow\mathrm{id}_A.
\]
After taking nerves, this is a homotopy from $Bce_0$ to the identity. Hence $Be_0$ is a homotopy equivalence. Evaluation at $0$ also gives a homotopy equivalence
$\operatorname{Map}(|N[1]|,BC)\to BC$, since the interval contracts to $0$. The canonical comparison
\[
 B\operatorname{Fun}([1],C)
       \longrightarrow\operatorname{Map}(|N[1]|,BC)
\]
commutes with these evaluations and is consequently a weak equivalence. No fibrancy assumption on $C$ is involved.

What cannot be inferred is the corresponding assertion on \emph{strict} endpoint fibres. The fibre over $(a,b)$ of $A\to C\times C$ is the discrete category $\operatorname{Hom}_C(a,b)$: a square with identity arrows at its two ends can only relate an arrow to itself. The homotopy fibre of the topological path fibration is instead the full path space from $a$ to $b$. The fork example discussed at ms.~128 already shows that these can differ even on components. Thus the motivation for alternating paths survives, while the sentence claiming failure of the absolute $Y=[1]$ comparison must be corrected.

There is also a direct example of absolute failure for another shape. Let $P$ be any poset with $BP\simeq S^1$, and take $Y=B\mathbb Z$. As above, $\operatorname{Fun}(B\mathbb Z,P)\cong P$. Its classifying space is connected, whereas $\operatorname{Map}(S^1,S^1)$ has one component for every integer degree. The comparison misses every nonconstant degree. This pinpoints the missing derived dependence on the source, without confusing the total arrow category with its fibres.
\end{dcomment}


\msp{page 136}

Par contre, $(*)$ est \texttt{[un]} isomorphisme quand $C$ est une catégorie de Thomason-Kan. Il suffit de voir que pour tout $Z$ dans (Cat), le Hom de $\hot(Z)$ dans les deux membres est isomorphe par la flèche déduite de $(*)$. Or si $C$ est de Thomason-Kan, de même $\uHom(Y,C)$ (immédiat sur la définition de Thomason-Kan en termes de prolongements de
\begin{tikzcd}[baseline=(\tikzcdmatrixname-1-1.base), row sep=small]
A_0 \arrow[r, hook, "i"] \arrow[dr] & A \arrow[d, dotted] \\
& C
\end{tikzcd}
$\ldots$(\footnote[29]{un peu vif! Il faut savoir que les applications cofibrantes $i$ de \textsc{Thomason} le restent après multiplication cartésienne par $Y$ \ldots})), donc
\[
\mathrm{Hom}_{\mathrm{Hot}}(\hot(Z),\hot(\uHom(Y,C)))\xleftarrow{\ \sim\ }\pi_0(\underbrace{\uHom_{\mathrm{Cat}}(Z,\uHom(Y,C))}_{\uHom_{\mathrm{Cat}}(Z\times Y,C)})\simeq\pi_0\uHom(Z\times Y,C), % Reading resolved against the supplied PDF; see corrections.csv.
\]
et
\[
\begin{aligned}
&\mathrm{Hom}_{\mathrm{Hot}}(\hot(Z),\uHom_{\mathrm{Hot}}(\hot(Y),\hot(C)))\simeq\mathrm{Hom}_{\mathrm{Hot}}(\underbrace{\hot(Z)\times\hot(Y)}_{\hot(Z\times Y)},\hot(C))\\
&\qquad\underbrace{\xleftarrow{\ \sim\ }}_{\text{car $C$ est de Thomason-Kan}}\qquad\quad\pi_0(\uHom_{\mathrm{Cat}}(Z\times Y,C)),
\end{aligned}
\]
O.K.

On voit donc que les catégories $\uCh_{\infty,\Sigma}(X;a,b)$ n'ont aucune tendance à être des catégories de Thomason-Kan. Pourtant, son type dans Hot doit être `le bon', sans avoir d'abord à remplacer $X$ par une enveloppe de Thomason-Kan $\bar{X}$.

\dgrblock{dgrb396}{A fibrant target does not derive every source}
\begin{dcomment}{development}{p136:On voit donc que les catégories}{A fibrant target does not derive every source}
Grothendieck's next attempt is to impose a Kan condition on the target. His footnote identifies one missing assertion about products of cofibrations, but the displayed proof also replaces arbitrary maps in the localized category by ordinary functors from an arbitrary source. Thomason fibrancy alone does not permit that replacement. The following example makes the issue visible already on components.

\textbf{A fibrant-poset obstruction.}\ours\ There are a Thomason-fibrant category $P$ and a finite category $Y$ for which
\[
 B\operatorname{Fun}(Y,P)\longrightarrow
                 \operatorname{Map}(BY,BP)
\]
is not a weak equivalence. Here ``finite'' means finitely many objects and arrows, not a finite nerve.

\emph{Proof.} Let $G=\mathbb Z/2$ and $Y=BG$. By \citet[Theorem~2.6]{Raptis2010}, Thomason's model structure restricts to posets, with the same weak equivalences and fibrations. Start with a poset representing $BG$, for example the twofold subdivision construction $c\operatorname{Sd}^2N(BG)$, and take its fibrant replacement within posets. We obtain a poset $P$ that is fibrant as a category and satisfies $BP\simeq BG$. Every functor $BG\to P$ sends the single object to some $p$ and every element of $G$ to $\mathrm{id}_p$. Natural transformations are exactly the relations $p\leq p'$. Hence
\[
                        \operatorname{Fun}(BG,P)\cong P.
\]
The source of the comparison is therefore connected. Its image consists of constant maps $BG\to BP$. But the target has at least two components: a constant map induces the zero homomorphism on $\pi_1$, whereas a homotopy equivalence $BG\to BP$ induces an isomorphism. These maps cannot be homotopic, even without specified basepoints, since conjugation cannot turn a trivial homomorphism into an isomorphism. The comparison fails on $\pi_0$.\hfill$\square$

This also supplies an immediate test of the first isomorphism asserted in the proof on ms.~136: take its $Z=e$. The claimed bijection would identify the connected set $\pi_0\operatorname{Fun}(Y,P)$ with $[BY,BP]$, which has the two distinct classes just exhibited. One cannot repair that step merely by proving that the functor category is fibrant. In Thomason's model structure the source has to be handled as well; in particular every cofibrant category is a poset \citep[Proposition~5.7]{Thomason1980}, so this $BG$ is not cofibrant.

The literal product claim in footnote~29 would also be too strong for arbitrary cofibrations: multiplying the cofibration $\varnothing\to e$ by $Y$ would make every category cofibrant. The next development also rules out preservation of all trivial cofibrations under product with arbitrary $Y$: the lifting adjunction would otherwise force every such functor category to be fibrant.

A safe model for the desired derived mapping space is the simplicial function complex
\[
              \operatorname{Map}_{\mathrm{sSet}}(NY,RNP),
\]
where $RNP$ is a Kan replacement. Every simplicial set is cofibrant, and this construction derives the target in the cartesian model category of simplicial sets. The ordinary nerve identity
$N\operatorname{Fun}(Y,P)=(NP)^{NY}$ explains precisely what the uncorrected formula omits. We checked the model structure and the poset replacement in the cited papers; we did not find this particular $B(\mathbb Z/2)$ application there. The result is recorded as an elementary application, not a priority claim.
\end{dcomment}


\printedpage{85}
% head: A. Grothendieck : Dérivateurs, Ch. VII, Édition des Universités de Montpellier II et Paris VII
% foot: Version 22 décembre 2001

\msp{page 137}

\dgrblock{dgrb397}{Fibrancy is not preserved by strict diagram categories}
\begin{dcomment}{development}{p136}{Fibrancy is not preserved by strict diagram categories}
The failure of the preceding mapping-space formula does not logically decide the separate assertion that $\operatorname{Fun}(Y,C)$ is fibrant whenever $C$ is. That assertion also fails, but it needs a different example. An equivariant construction makes the obstruction precise: ordinary horn filling need not respect an action on the diagram.

\textbf{A finite parameter with a nonfibrant functor category.}\ours\quad
There are a countable Thomason-fibrant category $E$ and a finite category $Y$ such that $\operatorname{Fun}(Y,E)$ is not Thomason-fibrant.

\emph{Proof.} Let $D$ be the four-object crown used above, so $BD\simeq S^1$. It is not Thomason-fibrant. Indeed, $\operatorname{Ex}^2ND$ has finitely many $1$-simplices: they are the functors from the finite poset $c\operatorname{Sd}^2\Delta^1$ to $D$. In a Kan complex every element of a based fundamental group is represented by one loop edge, since horn fillers compose and invert edge paths. If $\operatorname{Ex}^2ND$ were Kan, its fundamental group would therefore be finite, contrary to its being isomorphic to $\pi_1(BD)=\mathbb Z$.

Put $G=\mathbb Z/2$ and give $D$ the trivial action. We first construct a $G$-poset $Q$ whose underlying category is fibrant and whose fixed subposet is exactly $D$. Start with $Q_0=D$. At step $r$, for every generating Thomason horn inclusion
\[
 A=c\operatorname{Sd}^2\Lambda_i^n
       \ \longrightarrow\ B=c\operatorname{Sd}^2\Delta^n
\]
and every ordinary functor $F:A\to Q_r$, attach the free $G$-cell $G\times B$ along $G\times A$, using the equivariant attaching map $(g,a)\mapsto gF(a)$. Here $G$ in a product denotes the discrete set of its two elements. Take all these pushouts simultaneously and put $Q=\varinjlim_rQ_r$.

These horn inclusions are Dwyer inclusions of finite posets. The pushout description in \citet[Lemmas~4.7--4.8]{MayZakharevichStephan2017} shows that the old poset stays full and that the new objects are exactly the disjoint sets $G\times(B\setminus A)$. All new objects have free $G$-action. Thus $Q^G=D$, with its original order. Every horn map into $Q$ factors through some finite stage, because its domain is finite, and it obtains a filler at the next stage. Consequently $Q$ is Thomason-fibrant. There are only countably many finite horn shapes and, at each countable stage, countably many maps from them. Hence $Q$ is countable.

Define the transport category $E=Q\rtimes G$ as follows. Its objects are those of $Q$, and an arrow $x\to y$ labelled $g\in G$ exists exactly when $gx\leq y$; composition multiplies labels. There is a projection $E\to BG$. We claim that $E$ is Thomason-fibrant. Given a horn diagram $F:A\to E$, write $g_{ab}$ for the label of an arrow $a\to b$. Since $BA$ is contractible, its fundamental groupoid is equivalent to a point. The projected functor $A\to BG$ is therefore naturally isomorphic to a constant functor: choose labels $h_a$ with
\[
                         h_b g_{ab}=h_a.
\]
Replacing $F(a)$ by $h_aF(a)$ turns every arrow label into $1$ and gives a functor $A\to Q$. Extend it to $B\to Q$ by fibrancy of $Q$. Extend the assignments $h_a$ arbitrarily to the remaining objects of $B$, and undo the replacement: at $b$ use $h_b^{-1}F'(b)$, and label $b\to c$ by $h_c^{-1}h_b$. The inequalities defining $F'$ make this a functor $B\to E$, and it restricts exactly to the original horn diagram. This proves the claim.

Finally set $Y=BG$. A functor $BG\to E$ induces a homomorphism $\phi:G\to G$, either trivial or the identity. The trivial case is the identity action at an arbitrary object of $E$. In the identity case, an endomorphism labelled by the nontrivial $g$ requires $gx\leq x$; applying $g$ gives the reverse inequality, so $gx=x$. Its objects are therefore precisely $D=Q^G$. A natural transformation labelled by $h$ must satisfy $h\phi(g)=\psi(g)h$. Since $G$ is abelian, this forces $\phi=\psi$. For the identity case the remaining condition is just $x\leq y$, and every $h\in G$ is allowed. We obtain an actual decomposition
\[
             \operatorname{Fun}(BG,E)
                      \cong E\amalg(D\times BG).
\]
The category $D$ is a retract of this coproduct: project the second summand onto $D$, send the first constantly to a chosen object of $D$, and use identity labels for the section. A retract of a fibrant category is fibrant. Since $D$ is not fibrant, neither is $\operatorname{Fun}(BG,E)$.\hfill$\square$

Thus both claims on ms.~136 fail: arbitrary strict diagram categories need not preserve Thomason fibrancy, and even when a particular functor category is fibrant, it need not represent the derived mapping space. The equivariant cell and fixed-point ingredients above are standard and are cited at their use. We did not locate this application to strict diagram categories in the primary sources checked, including the published equivariant-poset paper; the argument is supplied here without a claim of priority.
\end{dcomment}


\section*{12\quad Retour sur les `catégories de chemins' dans \textmd{Cat}, comme catégorie de modèles pour \textmd{Hot}}

Je vois exactement deux applications intéressantes du formalisme des chemins dans les catégories. L'une, qui m'a surtout forcé la main, est pour pouvoir construire des Hom\,ext, Hom \textit{externes} à valeurs dans Hot, pour tout prédérivateur $\mathbf{D}$ défini par un couple $(\mathcal{M},\Sigma)$ (et je ne connais pas de prédérivateur qui ne soit défini ainsi).

L'autre application est pour avoir dans (Cat) une \textit{construction fonctorielle,} en termes d'un objet $X$ quelconque ou d'une flèche $X\longrightarrow Y$, de l'analogue du formalisme de \textsc{Cartan-Serre} sur les espaces de chemins. Dans ce cadre, il est entendu qu'on prend toujours $\Sigma_X=\mathrm{Fl}(X)$, ensemble de \textit{toutes} les flèches de $X$. Pour construire fonctoriellement en $X$ dans (Cat)
\[
\begin{tikzcd}[row sep=huge]
\uCh_\infty(X) \arrow[d] \\
X\times X,
\end{tikzcd}
\]
il faut avoir fixé un système projectif
\[
\underline{t}=(t_n)_{n\in\mathbf{N}}
\]
\msp{page 138}
dans \underline{Ty} qui soit cofinal dans \underline{Ty}, ce qui signifie simplement que parmi les $t_n$ il y en a de longueur utile arbitrairement grande (et comme
\[
n\longmapsto\mathrm{long.ut.}(t_n)
\]
est croissante, cela implique que long.ut.$(t_n)$ tend vers $+\infty$ pour $n\longrightarrow+\infty$). Les faits principaux à vérifier ensuite (si faire se peut) sont les propriétés essentielles des flèches dans le diagramme
\printedpage{86}
% head: A. Grothendieck : Dérivateurs, Ch. VII, Édition des Universités de Montpellier II et Paris VII
% foot: Version 22 décembre 2001
\[
\begin{tikzcd}[column sep=large, row sep=large]
& X \arrow[d, "\tau"] & \\
& \uCh_\infty(X) \arrow[d, "p"] \arrow[ddl, "\sigma"'] \arrow[ddr, "\beta"] & \\
& X\times X \arrow[dl, "\mathrm{pr}_1"] \arrow[dr, "\mathrm{pr}_2"'] & \\
X & & X\ ,
\end{tikzcd}
\qquad\qquad
\left|\begin{array}{l}
\sigma\tau=\beta\tau=\mathrm{id}_X\\
p\tau=\mathrm{diag}_X
\end{array}\right.
\]
et plus généralement, pour une flèche
\[
X'\xrightarrow{\ f\ }X
\]
dans Cat, des flèches dans le diagramme correspondant
\[
\begin{tikzcd}
& \overline{X'}=C(f) \arrow[dl] \arrow[dr, "g"] \arrow[ddrr, bend left=30, "\bar{f}"] & & \\
X' \arrow[ur, dotted, bend left=40, "\alpha"] \arrow[dr, "f"'] & |[draw, rounded corners]| \text{cart.} & \uCh_\infty(X) \arrow[dl, "\sigma"] \arrow[dr, "\beta"] & \\
& X & & X,
\end{tikzcd}
\qquad
\bar{f}\alpha=(\beta g)\alpha=f,
\]
et plus particulièrement de $\beta g$. La flèche
\msp{page 139}
en pointillés $\alpha$ est la section de $\overline{X'}$ sur $X'$ déduite de la section $\tau$ de $\uCh_\infty(X)$ sur $X$ qui réalise la décomposition de $f$ en
\[
X'\xrightarrow{\ \alpha\ }\overline{X'}\xrightarrow{\ \bar{f}\ }X,
\]
où $\bar{f}$ devrait être une flèche \textit{`fibrante'} et une Hot-équivalence, $\alpha$ une flèche \textit{`cofibrante'} au sens de \textsc{Thomason}. (Mais s'il en est ainsi, est à voir!) Voici pour le moment les propriétés essentielles que je voudrais vérifier.

\dgrblock{dgrb398}{The two maps in a mapping-path factorization}
\begin{dcomment}{exposition}{p139:Voici pour le moment les propriétés essentielles}{The two maps in a mapping-path factorization}
Grothendieck turns from representing an external Hom to replacing a map by its homotopy fibres. This change of purpose makes the endpoints decisive. Put $E=\Ch_\infty(X)$, with source $\sigma$, target $\beta$, and constant-path section $i$. Then
\[
 C(f)=X'\times_{X,\sigma}E,
 \quad r:C(f)\longrightarrow X',
 \quad \bar f=\beta\operatorname{pr}_E,
 \quad \alpha(x')=(x',i(fx')).
\]
The strict identities are $r\alpha=\mathrm{id}$ and $\bar f\alpha=f$. Geometrically, $r$ forgets a path whose initial point is prescribed; its fibre should be contractible. The other map remembers the final point; its fibre should be the homotopy fibre of $f$. Thus the weak equivalence is $\alpha$, whereas $\bar f$ is a weak equivalence exactly when $f$ is. The sentence on ms.~139 that attaches the unconditional weak-equivalence assertion to $\bar f$ is contradicted by item~$3^o$ immediately below it.

Once $r$ is known to be a weak equivalence, the rest of this assertion uses only two-out-of-three. Since $r\alpha=\mathrm{id}$, the section $\alpha$ is a weak equivalence. Since $\bar f\alpha=f$, the maps $f$ and $\bar f$ are weak equivalences simultaneously. Neither inference says that $\alpha$ is a Thomason cofibration, or that $\bar f$ is a Thomason fibration. Smoothness and properness concern the homotopy exactness of base change, while Thomason's chosen fibrations are tested by $\operatorname{Ex}^2N$. Their definitions must not be interchanged.

The cofibration expectation actually fails for the unaltered constant-path inclusion. Take $f=\mathrm{id}_X$ with $X=[1]$. A nonconstant alternating path from $0$ to $1$, padded by identity arrows, admits a natural transformation to the constant path at $1$: at a vertex labelled $0$ use $0\to1$, and at a vertex labelled $1$ use its identity. Hence the image of $\alpha:X\to E$ is not a sieve. Every Thomason cofibration is a pseudo-Dwyer map and therefore a sieve \citep[\S2, following Proposition~2.4]{Raptis2010}. Thus $\alpha$ need not be a Thomason cofibration. This does not damage the intended homotopy-fibre replacement; it limits the proposed identification with a particular model-category factorization.

The expectation that $\bar f$ is always a Thomason fibration also fails directly. Take $X=e$. Since $\Ch_\infty(e)=e$ strictly, the construction gives $C(f)=X'$ and $\bar f=f:X'\to e$. Choosing $X'=D$, the nonfibrant crown of the preceding development, gives a counterexample. The homotopical lifting properties proved later do not impose Thomason fibrancy on an arbitrary category mapping to $e$.

The properties in items~$1^o$--$3^o$ are requirements at this stage of the manuscript. Their proofs begin only after the concrete symmetric model is fixed on ms.~145. We will use them in that order, rather than reading the list as an already established theorem for an unspecified projective system.
\end{dcomment}


\begin{enumerate}[label={}, leftmargin=*]
\item[1$^o$)] $p$ est lisse et propre, et est une HOT-fibration.

\item[2$^o$)] Donc il en est de même de $\sigma=\mathrm{pr}_1p$ et $\beta=\mathrm{pr}_2p$. De plus, ce sont des HOT-fibrations \textit{triviales,} i.e. à fibres asphériques.

\printedpage{87}
% head: A. Grothendieck : Dérivateurs, Ch. VII, Édition des Universités de Montpellier II et Paris VII
% foot: Version 22 décembre 2001

\item[3$^o$)] Plus généralement (pour ce qui concerne $\beta$), $\overline{X'}\xrightarrow{\ \bar{f}\ }X$ devrait encore être propre et lisse et une HOT-fibration, mais bien sûr en général pas triviale. (Elle devrait être triviale si et seulement si $f$ est une Hot-équivalence.) (\footnote[30]{résulte formellement de 1$^o$ et 2$^o$, cf. page 170.}).
\end{enumerate}

Un mot sur la `dualité'. Les énoncés sur $\sigma$ (resp. $\beta$) se traduisent en énoncés sur $\beta$ (resp. $\sigma$) par passage de $\underline{t}$ à $\underline{t}^o=(t^o_n)$, où $\tau\longmapsto\tau^o$ est le foncteur dans \underline{Ty} `passage à l'intervalle opposé'. On a en effet % Reading resolved against the supplied PDF; see corrections.csv.

\msp{page 140}

un isomorphisme fonctoriel en $X$ (passage au chemin opposé)
\[
\boxed{\begin{array}{ccc}
c & \longmapsto & c^o\\[4pt]
\uCh^{\underline{t}}_\infty(X) & \longrightarrow & \uCh^{\underline{t}^o}_\infty(X)
\end{array}}
\qquad\qquad
\text{(\footnotemark[31])}
\]
\footnotetext[31]{cf. page 143 pour le passage de $\underline{t}$ à $\underline{t}^-$.}
qui `échange source et but'. Bien sûr, on a intérêt à prendre $\underline{t}$ tel qu'on ait $\underline{t}=\underline{t}^o$ (c'est ce que j'avais fait, par un sain instinct, dès le début dans V), de sorte qu'on a un foncteur involutif canonique
\[
\begin{array}{ccc}
c & \longmapsto & c^o\\[4pt]
\uCh_\infty(X) & \longrightarrow & \uCh_\infty(X)
\end{array}
\]
(je laisse tomber à nouveau l'exposant $\underline{t}$), qui `échange source et but'. Donc on trouve, à isomorphisme canonique près, la `même' factorisation de $f:X'\longrightarrow X$ en $\bar{f}\circ\alpha$, si on fait comme dans page 38 pour définir $\overline{X'}$ et $\alpha$, ou si dans ce diagramme on intervertit les rôles de $\sigma$ et de $\beta$. % [sic] page 38 -> page 138, the diagram defining \overline{X'} and \alpha stands on manuscript page 138 (apparently)

\dgrblock{dgrb399}{Two different reversals}
\begin{dcomment}{exposition}{p140:Donc on trouve, à isomorphisme canonique près}{Two different reversals}
Grothendieck needs two symmetries because two different operations are being compared: exchanging the endpoints of a path, and replacing the ambient category by its opposite. For a signed word $\tau=(\varepsilon_1,\ldots,\varepsilon_r)$, with $\varepsilon_i\in\{+,-\}$, write
\[
 \tau^o=(-\varepsilon_r,\ldots,-\varepsilon_1),
 \qquad
 \tau^-=(-\varepsilon_1,\ldots,-\varepsilon_r).
\]
The first operation reverses the ordered interval. It is an anti-automorphism of concatenation and interchanges source and target without changing $X$. The second reverses every categorical arrow while preserving the ordering of the vertices. It is an automorphism of concatenation and is the operation accompanying $X\mapsto X^{\mathrm{op}}$.

At a fixed stage the corresponding identities are obtained by reindexing diagrams and natural transformations:
\[
 \Ch_\tau(X)\cong\Ch_{\tau^o}(X),
 \qquad
 \Ch_\tau(X^{\mathrm{op}})
          \cong\Ch_{\tau^-}(X)^{\mathrm{op}}.
\]
The first identity is understood over the transposition of $X\times X$; the second is over $(X\times X)^{\mathrm{op}}$. Passing through the specified transition maps gives the identities for projective systems. For $\tau_0=\Delta_+\Delta_-$ one has $\tau_0^o=\tau_0$ but $\tau_0^-=\Delta_-\Delta_+$. This explains why path reversal becomes a strict involution in the symmetric model, while categorical duality initially changes its parity. The translation introduced on ms.~147 supplies the remaining comparison. The commentary at manuscript p.~9 already distinguishes these operations; here their role is to transfer the endpoint and base-change assertions.
\end{dcomment}


Enfin :

\begin{enumerate}[label={}, leftmargin=*]
\item[4$^o$)] Si
\[
c:a\longrightarrow b
\]
est un chemin quelconque dans $X$, on voudrait en conclure un isommorphisme dans Hot (de `transport par $c$') % [sic] isommorphisme -> isomorphisme, doubled m
\[
\overline{X'}_a\xrightarrow[\sim]{\bar{c}}\overline{X'}_b,
\]
avec bien sûr transitivité pour la composition,

\msp{page 141}

ce qui donnera aussi
\[
\overline{c^{\mathrm{opp}}}=\overline{c}^{-1}
\]
(en admettant que les chemins triviaux $a\xrightarrow{\mathrm{id}}a$ donnent l'identité $\overline{X'}_a\xrightarrow[\sim]{\mathrm{id}}\overline{X'}_a$ !) (\footnotemark[32]\textsuperscript{,}\footnotemark[33]). Bien sûr, $\overline{c}$ ne doit dépendre que de la \textit{classe d'homotopie} de $\overline{c}$ (dans % [sic] de $\overline{c}$ (dans -> de $c$ (dans, the homotopy class is apparently that of the path c
\footnotetext[32]{Surtout pas se borner au cas où $c$ serait un des $t_n$!}
\footnotetext[33]{décrire quand-même l'effet sur les $\pi_0$!}
\printedpage{88}
% head: A. Grothendieck : Dérivateurs, Ch. VII, Édition des Universités de Montpellier II et Paris VII
% foot: Version 22 décembre 2001
$\uCh_\tau(X;a,b)$, ou même dans $C=\uCh(X;a,b)$, catégorie de tous les chemins dans $X$, de $a$ vers $b$, de type non précisé). En d'autres termes, on doit avoir un foncteur
\[
\begin{array}{rcl}
\Pi(X) & \longrightarrow & \mathrm{Hot}\\[4pt]
a & \longmapsto & \hot(\overline{X'}_a)
\end{array}
\]
du groupoïde fondamental de $X$ dans Hot. À vrai dire, ceci n'a rien à voir avec la situation particulière étudiée ici - ça doit marcher chaque fois qu'on a une HOT-fibration, ici $\overline{X'}\longrightarrow X$, mais aussi $\uCh_\infty(X)\longrightarrow X\times X$ (\footnote[34]{Ici aussi, il faudrait décrire l'effet du transport $\overline{c}$ sur les $\pi_0$ des fibres.}). Et pour une $\mathbf{D}$-fibration $X'\longrightarrow X$ ($\mathbf{D}$ un dérivateur quelconque), et $a\in\mathbf{D}(e)$, on doit avoir de même deux foncteurs
\[
\Pi(X)\rightrightarrows\mathbf{D}(e)
\]
correspondant à
\[
\left\{\begin{array}{lcl}
x & \longmapsto & H^\bullet_{\mathbf{D}}(X'_x,\bullet)\\[4pt]
x & \longmapsto & H_\bullet^{\mathbf{D}}(X'_x,\bullet).
\end{array}\right.
\]

\msp{page 142}

Le cas qui nous occupe ici serait celui où $\mathbf{D}=\mathrm{HOT}$, où on prend $H^{\mathbf{D}}_\bullet$, et $a=$ objet final de $\mathbf{D}(e)=\mathrm{Hot}$.)

\item[5$^o$)] Si 1$^o$, 2$^o$, 3$^o$ sont vérifiés, on est moralement sûr que la construction des $\uCh_\infty(X)$ est bel et bien l'équivalent exact du formalisme de \textsc{Cartan-Serre}. Pour en avoir le c\oe ur totalement net, i.e. raccrocher explicitement cette théorie aux constructions classiques, il faudrait prouver un \textit{énoncé de comparaison} avec la théorie semi-simpliciale, qui dirait essentiellement ceci : Si on désigne, pour $X$ dans Cat, par $\tilde{X}$ l'objet associé dans $\mathbf{\Delta}^\wedge$ (le `nerf' de $X$), alors la flèche
\[
\uCh_\infty(X)\tilde{}\longrightarrow\underbrace{(X\times X)\tilde{}}_{\simeq\,\tilde{X}\times\tilde{X}}
\]
dans $\mathbf{\Delta}^\wedge/(X\times X)\tilde{}$, est $W$-équivalente sur $(X\times X)\tilde{}$ à la flèche
\[
\uCh(\tilde{X})\longrightarrow\tilde{X}\times\tilde{X}
\]
de la théorie semi-simpliciale des chemins (à supposer qu'il y ait une telle flèche classique définie pour \textit{tout} objet $X$ de $\mathbf{\Delta}^\wedge$). Ou qu'on ait tout au moins des isomorphismes canoniques dans Hot
\[
\hot_{\mathrm{Cat}}(\uCh_\infty(X;a,b))\ \simeq\ \hot_{\mathbf{\Delta}^\wedge}(\uCh(\tilde{X};\tilde{a},\tilde{b})).
\]
\end{enumerate}

\dgrblock{dgrb400}{Transport and the comparison with ordinary path spaces}
\begin{dcomment}{later}{p142:Ou qu'on ait tout au moins des isomorphismes canoniques}{Transport and the comparison with ordinary path spaces}
Grothendieck wants the replacement to remember how its fibres vary, and his marginal request to compute $\pi_0$ is the right concrete test. For the ordinary mapping-path construction, a point in the homotopy fibre at $a$ is a pair $(x',\gamma)$ with $\gamma:f(x')\leadsto a$. Transport along $c:a\leadsto b$ sends it to
\[
                       (x',\gamma*c).
\]
Homotopic paths give homotopic maps of fibres; concatenation and reversal give composition and inverse in the homotopy category. This produces the functor from the fundamental groupoid requested in item~$4^o$. It is a truncation of the coherent family of fibres. A functor into $\Hot$ alone need not retain all that family's higher twisting, a distinction already illustrated in the commentary on Chapter~VI, manuscript pp.~7 and~74.

For a connected pointed map $f:(X',x'_0)\to(X,a)$, this transport has an explicit component formula. Set $G=\pi_1(X,a)$ and $H=\operatorname{im}\pi_1(f)$. Then
\[
                   \pi_0\operatorname{hofib}_a(f)\cong H\backslash G.
\]
To see the formula, join $x'_0$ to $x'$ by a path $\delta$. The pair $(x',\gamma)$ determines the loop $f\delta*\gamma$. A different $\delta$ multiplies its class on the left by an element of $H$. Conversely, a homotopy of pairs gives precisely this relation, by tracing the path of their first coordinates; any such relation is realized by that construction. A loop $c$ at $a$ therefore acts by right multiplication $Hg\mapsto Hg[c]$. For the inclusion $2\mathbb Z\subset\mathbb Z$, viewed as a map of one-object groupoids, there are two components and a generator of the base interchanges them. This answers the request about $\pi_0$ without choosing representatives in a categorical lifting category.

There is a functorial simplicial path object for every simplicial set $K$, not only for a Kan complex. Choose a functorial Kan replacement $K\to RK$ and put
\[
 P(K)=(K\times K)\times_{RK\times RK}(RK)^{\Delta^1}.
\]
Its endpoint map is a Kan fibration. The constant-path map $K\to P(K)$ is a weak equivalence: right properness first identifies $P(K)\to(RK)^{\Delta^1}$ as a weak equivalence, and evaluation identifies the latter with $RK$. Apply this with $K=NX$. Once $i:X\to\Ch_\infty(X)$ has been proved a weak equivalence, the span
\[
             N\Ch_\infty(X)\ \longleftarrow\ NX
                              \ \longrightarrow\ P(NX)
\]
is already a zigzag of weak equivalences over $NX\times NX$. This gives the first form of comparison in item~$5^o$. Identifying the \emph{strict} endpoint fibres with ordinary path spaces still requires the homotopy-fibre property of the categorical endpoint map; a weak equivalence of total spaces over a base does not by itself identify strict fibres.

For a general coefficient $A\in\mathbb D(e)$, the analogous homology transport is covariant; cohomology transport uses pullback along the inverse fibre equivalence. These constructions require the homotopy base-change isomorphisms appropriate to the stated $\mathbb D$-fibration.

The later theory of locally constant diagrams makes the distinction between a coherent family and its monodromy systematic \citep[Theorem~1.3 and the discussion preceding it]{Cisinski2009LocallyConstant}. The language of straightening provides a higher-categorical model of the family itself \citep[Theorem~2.2.1.2]{Lurie2009HTT}. These results explain the final form of the construction; the elementary fibre comparison remains the work that this chapter undertakes.
\end{dcomment}


% ---- from Der-ChVII-456_p091-095.tex
\printedpage{89}
% head: A. Grothendieck : Dérivateurs, Ch. VII, Édition des Universités de Montpellier II et Paris VII
% foot: Version 22 décembre 2001

\msp{page 143}
Dans $1^o$ et $3^o$, on veut prouver que certaines flèches ($p$ et $\bar{f}$) sont à la fois \textit{propre} et \textit{lisse} - ce qui est vraiment idéal ! Je dis qu'il suffit à chaque fois de prouver l'un des deux, que l'autre s'en déduit (comme `cas particulier') par dualité. Pour ceci, notons
\[
\tau \longmapsto \tau^{-}
\]
l'automorphisme (\textit{pas anti}-automorphisme) de $\Ty$ qui consiste à changer de signe toutes les flèches de transition - c'est donc, au niveau des objets, l'unique automorphisme du monoïde libre Ty, qui échange $\Delta_+$ et $\Delta_-$. Au système projectif $\underline{t}$ dans $\Ty$ correspond le système projectif $\underline{t}^-$. Ceci dit, on a un isomorphisme de catégories
\[
\boxed{\Ch^{\underline{t}}_{\infty}(X^o) \ \simeq\ (\Ch^{\underline{t}^-}_{\infty}(X))^o} \tag*{$(^{35})$} % Reading resolved against the supplied PDF; see corrections.csv.
\]
\footnotetext[35]{cf. page 140 pour le passage de $\underline{t}$ à $\underline{t}^o$.}
et le diagramme commutatif
\[
\begin{tikzcd}
\Ch^{\underline{t}}_{\infty}(X^o) \arrow[r, "\sim"] \arrow[d, "p^{\underline{t}}_{X^o}"] & \Ch^{\underline{t}^-}_{\infty}(X)^o \arrow[d, "(p^{\underline{t}^-}_{X})^o"] \\
X^o \times X^o \arrow[r, equal] & (X \times X)^o.
\end{tikzcd}
\]
Ainsi, dire que $p^{\underline{t}}_{X}$ \texttt{[plutôt} $p^{\underline{t}^-}_{X}$\texttt{]} est lisse, i.e. que $(p^{\underline{t}}_{X})^o$ \texttt{[plutôt} $(p^{\underline{t}^-}_{X})^o$\texttt{]} est propre, équivaut à dire que $p^{\underline{t}}_{X^o}$ est propre.

\msp{page 144}
Il nous reste donc, à ce sujet, que deux faits à vérifier : la lissité de $p : \Ch_{\infty}(X) \longrightarrow X \times X$, et celle de $\bar{f} : \overline{X'} \longrightarrow X$.

a) Lissité de $p : \Ch_{\infty}(X) \longrightarrow X \times X$. On se donne une flèche
\[
w : (a,b) \longrightarrow (a',b'), \qquad \text{i.e.} \quad \left\{ \begin{array}{lccc} u : & a & \longrightarrow & a' \\[1ex] v : & b & \longrightarrow & b', \end{array} \right.
\]
sur la base, un objet $c'$ au dessus du but, i.e.
\[
c' \in \Ch_{\infty}(X; a', b'),
\]
et on regarde la catégorie des flèches dans $\Ch_{\infty}(X)$, de but $c'$, qui relèvent $w$,
\[
\begin{tikzcd}
\underset{\text{\scriptsize variable}}{c} \arrow[r, "\varphi"] \arrow[d, dash] & \underset{\text{\scriptsize fixe}}{c'} \arrow[d, dash] \\
(a,b) \arrow[r, "w"] & (a',b'),
\end{tikzcd}
\]
\printedpage{90}
% head: A. Grothendieck : Dérivateurs, Ch. VII, Édition des Universités de Montpellier II et Paris VII
% foot: Version 22 décembre 2001
i.e. la catégorie des couples $(c,\varphi)$, où
\[
\left\{ \begin{array}{l}
c \in \Ch_{\infty}(X; a, b) \\[1ex]
\varphi : c \longrightarrow c' \ \ (\in \Ch_{\infty}(X; a', b')) \\[1ex]
\text{une flèche induisant } \begin{array}[t]{l} u \ \text{ sur les sources} \\ v \ \text{ sur les buts,} \end{array}
\end{array} \right.
\]
\[
\begin{tikzcd}
a \arrow[rrrr, dotted, dash, "c"] \arrow[d] & {} \arrow[d, dotted] & {} \arrow[d, dotted] & {} \arrow[d, dotted] & b \arrow[d] \\
a' \arrow[rrrr, dash, "c'"] & {} & {} & {} & b'.
\end{tikzcd}
\]

\msp{page 145}
\begin{flushright}
{\footnotesize 9.11.90}
\end{flushright}

\dgrblock{dgrb401}{Why a category of lifts is the right test}
\begin{dcomment}{exposition}{p144}{Why a category of lifts is the right test}
Grothendieck now replaces the geometric word ``smooth'' by a calculation that can be carried out on paths. Fix an arrow $w:b\to b'$ in a base $B$ and an object $e'$ above $b'$. For $p:E\to B$, the relevant category has objects $(e,\varphi)$ with $p(e)=b$ and $\varphi:e\to e'$ satisfying $p(\varphi)=w$. A morphism from $(e,\varphi)$ to $(e_1,\varphi_1)$ is a vertical map $h:e\to e_1$ with $\varphi_1h=\varphi$. Thus one must control both existence of lifts and the homotopy type of the space of choices.

This is exactly a comma category for the inclusion $E_b\to b\backslash E$. Requiring it to be contractible for every $(e',w)$ says that this inclusion is aspherical. The equivalence with smoothness is the criterion of \citet[Theorem~1.2]{Maltsiniotis2005Asphericite}. In the present case $B=X\times X$ and $w=(u,v)$. The annotation placing $\varphi$ inside $\Ch_\infty(X;a',b')$ cannot be literal: its source has endpoints $(a,b)$, so it is a morphism of the total path category over $(u,v)$.

The next construction is designed to make these comma categories manageable. At a finite interval whose two endpoint arrows point into the interior, a lift is obtained by changing the two endpoints and composing with $u$ and $v$. This is a Cartesian lift, hence a terminal object in the corresponding category of lifts. To deal with a path already lying in some stage, enlarge its support to a stage with the required endpoint parity. The filtered union then accommodates every finite diagram in the lifting category. This is why the two-index system introduced next is useful: it permits the same infinite path category to be approached by finite stages suited separately to incoming and outgoing lifts.

The proof of properness uses the opposite variance, as the duality formula on ms.~143 indicates. It does not follow from the mere existence of one lift, nor from a statement about the homotopy type of the total path category. The full calculation is carried out on ms.~150--153.
\end{dcomment}


Pour la question précédente, relative aux flèches dans le diagramme
\[
\begin{tikzcd}
& X \arrow[d, "i"] & \\
& \Ch_{\infty}(X) \arrow[d, "p"] \arrow[ddl, "\sigma"'] \arrow[ddr, "\beta"] & \\
& X \times X \arrow[dl, "\mathrm{pr}_1"] \arrow[dr, "\mathrm{pr}_2"'] & \\
X & & X,
\end{tikzcd}
\]
il faut préciser absolument la déscription de $\Ch_{\infty}(X)$ \textit{dans} Cat, donc il faut préciser le système projectif $\underline{t} = (t_n)_{n \in \mathbf{N}}$. Nous allons prendre en effet % [sic] déscription -> description, misspelling
\[
t_n = (\underbrace{\Delta_+ \Delta_-}_{\tau_0})^{2n},
\]
avec les morphismes de transition décrits par ailleurs
\[
\underbrace{t_{n+p}}_{=\tau_0^{2n+2p}} \ \longrightarrow \ \underbrace{t_n}_{\tau_0^{n}}, % [sic] \tau_0^{n} -> \tau_0^{2n}, since t_n = (\Delta_+\Delta_-)^{2n} = \tau_0^{2n}
\]
\printedpage{91}
% head: A. Grothendieck : Dérivateurs, Ch. VII, Édition des Universités de Montpellier II et Paris VII
% foot: Version 22 décembre 2001
étant donnés par $\delta_p * \mathrm{id}_{\tau_o^{2n}} * \delta_p$, où $\delta_p : \tau_0^p \longrightarrow 1$ est la dégénérescence de $\tau_0^p$. Ainsi on peut considérer $\Ch_{\infty}(X)$ comme étant la catégorie des chemins biïnfinis dans $X$ % [sic] \tau_o^{2n} -> \tau_0^{2n}, letter o printed in the index instead of 0
\[
\cdots \longrightarrow a_{-5} \longleftarrow \overbrace{a_{-4} \longrightarrow a_{-3} \xleftarrow{u_{-2}} \overbrace{a_{-2} \xrightarrow{u_{-1}} a_{-1} \xleftarrow{u_0} \overbrace{a_0}^{\Ch_0} \xrightarrow{u_1} a_1 \xleftarrow{u_2} a_2}^{\Ch_1} \xrightarrow{u_3} a_3 \longleftarrow a_4}^{\Ch_2} \longrightarrow a_5 \longleftarrow \cdots,
\]
avec les $u_i$ des identités pour $|i|$ assez grand.

La sous-catégorie (pleine)
\[
\Ch^{W}_{t_n}(X) \ = \ \Ch^{W}_{(\Delta_+\Delta_-)^{2n}}
\]
s'identifie à celle pour laquelle on a
\begin{quote}
$u_i$ une identité pour $i \notin [-2i+1, 2i]$, i.e. les seuls $u_i$ `significatifs' sont les $4i$ flèches $u_{-2i+1}$, $u_{-2i+2}$, $\ldots$, $u_{2i}$. % [sic] i (in [-2i+1,2i], 4i, u_{-2i+1}, u_{-2i+2}, u_{2i}) -> n, the index of t_n is n
\end{quote}

Mais il sera plus commode de regarder $\Ch_{\infty}(X)$ comme la limite inductive d'une
\msp{page 146}
famille à deux indices de sous-catégories pleines
\[
\Ch_{n,m}(X), \qquad n, m \in \mathbf{Z},
\]
en définissant cette catégorie comme \textit{vide} si $n > m$, et pour $n \leq m$ comme formée des objets
\[
\left\{ \begin{array}{l}
(a_i, u_i)_{i \in \mathbf{Z}} \in \Ch^{W}_{n,m}(X) \iff u_i \text{ une identité pour } i \notin [n, m-1] \\[1ex]
\overset{\mathrm{id}}{\cdots}\, a_{n-1} \dashto{\mathrm{id}} \underbrace{a_n \dashto{u_n} a_{n+1} \cdots \dashto{} a_{m-1} \dashto{u_{m-1}} a_m}_{\text{\scriptsize `partie utile' du chemin}} \dashto{\mathrm{id}} a_{m+1} \overset{\mathrm{id}}{\cdots}\,.
\end{array} \right.
\]
Donc pour $n = m$, on trouve la sous-catégorie pleine isomorphe à $X$, formée des chemins infinis `totalement dégénérés', i.e. tels que \textit{toutes} les $u_i$ soient des identités. (Cette sous-catégorie de $\Ch_{\infty}$ ne dépend donc pas du choix de $n = m$.)
\[
\xrightarrow{\mathrm{id}} \underbrace{a_{-1}}_{=\, a} \xleftarrow{\mathrm{id}} \underbrace{a_0}_{=\, a} \xrightarrow{\mathrm{id}} \underbrace{a_1}_{=\, a} \xleftarrow{\mathrm{id}}
\]
Tous les sommets d'un tel chemin sont égaux à un même objet de $X$, et les flèches de transition entre eux sont identiques.

La structure de $\Ch^{W}_{n,m}(X) \ \simeq \ \Ch^{W}_{\tau_{n,m}}(X)$ ne dépend que de la longueur $m-n$ de $\tau_{n,m}$ et de la \textit{parité} de $n$. Si $n$ est pair, la direction initiale de $\tau_{n,m}$ est positive, et
\[
\tau_{n,m} = \underbrace{\Delta_+ \Delta_- \cdots}_{m-n \text{ facteurs en alternance}} = \left\{ \begin{array}{ll}
(\Delta_+ \Delta_-)^{\frac{m-n}{2}} & \text{si } m \text{ pair} \\[1ex]
(\Delta_+ \Delta_-)^{\frac{m-n}{2}-1} \Delta_+ & \text{si } m \text{ impair}
\end{array} \right.
\]
\texttt{[plutôt} $\frac{m-n-1}{2}$ \texttt{au lieu de} $\frac{m-n}{2} - 1$\texttt{]}, tandis que pour $n$ impair c'est juste l'inverse
\printedpage{92}
% head: A. Grothendieck : Dérivateurs, Ch. VII, Édition des Universités de Montpellier II et Paris VII
% foot: Version 22 décembre 2001
\msp{page 147}
\[
\tau_{n,m} = \underbrace{\Delta_- \Delta_+ \cdots}_{m-n \text{ facteurs alternés}} = \left\{ \begin{array}{ll}
(\Delta_- \Delta_+)^{\frac{m-n}{2}} & \text{si } m \text{ impair} \\[1ex]
(\Delta_- \Delta_+)^{\frac{m-n}{2}-1} \Delta_- & \text{si } m \text{ pair}
\end{array} \right.
\]
\texttt{[plutôt} $\frac{m-n-1}{2}$ \texttt{au lieu de} $\frac{m-n}{2} - 1$\texttt{]}. La direction initiale est conforme à la parité de $n$, i.e. égale à $(-1)^n$, tandis que la direction terminale est opposée à la parité de $m$, i.e. égale à $(-1)^m$ \texttt{[plutôt} $(-1)^{m+1}$\texttt{]}.

\dgrblock{dgrb402}{Finite support and endpoint parity}
\begin{dcomment}{exposition}{p147:La direction initiale est conforme}{Finite support and endpoint parity}
Grothendieck's two-sided sequence makes the eventual constancy of a path concrete. Its advantage is that support may be enlarged independently at the two ends. To avoid the shift in edge numbering between ms.~145 and ms.~146, write $e_i$ for the edge joining $a_i$ and $a_{i+1}$, directed forwards for even $i$ and backwards for odd $i$. Thus $e_i$ is the $u_{i+1}$ of the large diagram on ms.~145. A path supported on $[n,m]$ has $e_i$ equal to an identity outside $n\leq i<m$.

For $n<m$, the source-end sign is $(-1)^n$ and the target-end sign is $(-1)^{m-1}=(-1)^{m+1}$. The four cases are
\[
\begin{array}{c|c|c|c}
 n\bmod2&m\bmod2&\hbox{initial sign}&\hbox{terminal sign}\\ \hline
 0&0&+&-\\
 0&1&+&+\\
 1&0&-&-\\
 1&1&-&+
\end{array}
\]
For an odd length $m-n$, the number of full two-letter blocks is $(m-n-1)/2$; this is the correction already supplied by the editors in both displayed formulas. When $n=m$ there are no endpoint arrows and the subcategory is the constant-path copy of $X$; the sign table is not intended for that degenerate case.

The inclusion for $n'\leq n\leq m\leq m'$ extends a diagram by identity arrows. It is fully faithful, since a natural transformation is forced to be constant on each identity tail. Every eventually constant path and every natural transformation between two such paths belongs to a sufficiently large common support interval. Hence the filtered colimit of these full subcategories is the asserted category of two-sided paths, on morphisms as well as objects. The even-even intervals form a cofinal family of support intervals, and so do the odd-odd intervals. The first family has both endpoint arrows pointing into the interval and is adapted to Cartesian lifts; the second has both pointing out and is adapted to coCartesian lifts. Cofinality here is the ordinary directedness of support intervals, not an assertion that a fixed inclusion of categories is a weak equivalence.
\end{dcomment}


Je note que l'on a $\tau_0 = (\tau_0)^o$, d'où
\[
\underline{t} = \underline{t}^o, \qquad \text{mais } \underline{t} \neq \underline{t}^-,
\]
et pourtant on a des isomorphismes fonctoriels en $X$
\[
\Ch_{\underline{t}}(X) \ \simeq \ \Ch_{\underline{t}^-}(X)
\]
respectant source et but, à cause d'un isomorphisme
\[
\underline{t} \ \simeq \ \underline{t}^- \qquad \text{dans Pro}(\Ty).
\]
On peut dire p. ex. que le \textit{foncteur translation} $T^p$ d'amplitude $p$ (pour $p \in \mathbf{Z}$) est défini sur $\Ch^{W}_{\underline{t}}(X)$, par
\[
\left\{ \begin{array}{l}
T^p(((a_i),(u_i))_{i \in \mathbf{Z}}) = ((b_j),(v_j))_{j \in \mathbf{Z}} \\[1ex]
\text{avec} \left\{ \begin{array}{lcl} b_j & = & a_{j-p} \\[1ex] v_j & = & u_{j-p} \end{array} \right.
\end{array} \right.
\]
et que
\[
\begin{aligned}
T^p &: \left\{ \begin{array}{lcl} \Ch_{t}(X) & \xrightarrow{\sim} & \Ch_{-t}(X) \\[1ex] \Ch_{-t}(X) & \xrightarrow{\sim} & \Ch_{t}(X) \end{array} \right\} \ \text{si } p \text{ impair} \\
T^p &: \left\{ \begin{array}{lcl} \Ch_{t}(X) & \xrightarrow{\sim} & \Ch_{t}(X) \\[1ex] \Ch_{-t}(X) & \xrightarrow{\sim} & \Ch_{-t}(X) \end{array} \right\} \ \text{si } p \text{ pair.}
\end{aligned} % Reading resolved against the supplied PDF; see corrections.csv.
\]

\msp{page 148}
En d'autres termes, on peut formellement regarder la catégorie somme
\[
\Ch_{\pm}(X) = \Ch_{+}(X) \amalg \Ch_{-}(X)
\]
et l'opération translation $T$ comme un automorphisme qui échange les deux morceaux, et $T^p$ n'est autre que son itéré $p$-ème, qui échange les morceaux $+$ et $-$ si $p$ impair, qui les stabilise si $p$ pair.

\printedpage{93}
% head: A. Grothendieck : Dérivateurs, Ch. VII, Édition des Universités de Montpellier II et Paris VII
% foot: Version 22 décembre 2001

Il y a aussi l'automorphisme involutif de passage au chemin opposé, qui échange source et but, et que je définis (conformément à l'isomorphisme canonique identique
\[
\underline{t} \xrightarrow{\sim} \underline{t}^o \quad )
\]
comme la symétrie $S$
\[
\left\{ \begin{array}{l}
S((a_i),(u_i)) = ((b_j),(u_j)) \ \ \text{avec} \\[1ex]
\begin{array}{lcl} b_j & = & a_{-j} \\[1ex] u_j & = & u_{-j+1}. \end{array}
\end{array} \right.
\]
Si je le regarde comme opérant sur $\Ch^{W}_{\pm}(X)$, il est lié au foncteur $T$ par
\[
\left\{ \begin{array}{rcl}
S^2 & = & \mathrm{id} \\[1ex]
STS^{-1} & = & T^{-1},
\end{array} \right.
\]
de sorte que c'est le groupe semidirect de $\mathbf{Z}/2\mathbf{Z}$ par $\mathbf{Z}$ (groupe affine de dimension 1 de l'anneau $\mathbf{Z}$) qui opère sur $\Ch_{\pm}$.

\dgrblock{dgrb403}{The dihedral action on a two-sided path}
\begin{dcomment}{examples}{p148:de sorte que c'est le groupe semidirect de}{The dihedral action on a two-sided path}
Grothendieck introduces translations to compare the two parity choices without changing either endpoint. In the edge convention of the preceding block, translation and reversal are
\[
 (T^pa)_j=a_{j-p},\quad (T^pe)_j=e_{j-p},
 \qquad
 (Sa)_j=a_{-j},\quad (Se)_j=e_{-j-1}.
\]
These formulas apply equally to the components of a natural transformation, so they define functors, not merely permutations of the objects. Since paths are eventually constant, translation leaves both eventual values unchanged; reversal exchanges them. Odd translation reverses parity, whereas even translation preserves it. Reversal preserves each parity convention because reversing the direction of an edge also reverses the parity of its index.

The group relations are identities of functors. On vertices,
\[
 (S^2a)_j=a_j,
 \qquad
 (STSa)_j=a_{j+1}=(T^{-1}a)_j,
\]
and substitution in the edge formula gives the same equalities on edges and natural transformations. The acting group is therefore
\[
 D_\infty=\mathbb Z\rtimes\mathbb Z/2,
 \qquad (\epsilon,k):j\longmapsto\epsilon j+k,
 \quad \epsilon\in\{1,-1\}.
\]
The integer subgroup is the translation subgroup and reflection acts on it by inversion. This specifies the semidirect-product convention independently of the order of the groups in the French phrase.

For example, a path with one nonidentity edge $e_0:x\to y$ is carried by $T^2$ to the same edge at position $2$, with two additional identity positions on one side and two fewer on the other. Both have source $x$ and target $y$, but they are different objects of the strict path category. The formulas above prove an isomorphism of categories and the dihedral relations. They do not by themselves prove that $T^2$ is the identity after passage to $\Hot$. This is the specific coherence question taken up on the next page. It should be resolved by a homotopy or a comparison of the relevant diagrams, not by treating coordinate translation as literal equality.
\end{dcomment}


\msp{page 149}
On est embarassé de trouver un automorphisme $T^2$ de la situation $\Ch_{\infty}(X) \longrightarrow X \times X$, sans qu'il soit évident que dans Hot, ce foncteur induise l'identité! Je ne vois aucun moyen de relier un chemin $c$ et $T^2 c$ par une flèche dans $\Ch_{\infty}$, ni même pas un chemin `canonique'. Mais je me dis que cet automorphisme est décrit par un automorphisme, non pas du système projectif % [sic] embarassé -> embarrassé, misspelling
\[
(\tau_{n,m})_{n,m \in \mathbf{Z}^o \times \mathbf{Z},\ n \leq m}
\]
(où $\mathbf{Z}^o$ désigne l'ensemble ordonné \textit{opposé} de $\mathbf{Z}$), mais du \textit{pro-objet} de $\Ty$ qu'il définit (qui est bien sûr `le même' pour tous les systèmes projectifs cofinaux à gauche dans $\Ty$). Il en résulte déjà que $T^2$ induit l'identité sur l'ind-objet de Cat\,hot, $\underline{\underline{\mathrm{Ch}}}_{\infty}(X)$, et a fortiori sur l'ind-objet correspondant de Hot, ou de $\mathrm{Hot}^{\wedge}_{0}$, ou enfin pour l'objet de $\mathrm{Hot}^{\wedge}_{0}$ lui-même. De même pour les $\Ch_{\infty}(X; a, b)$. Il s'ensuit, revenant à $\Ch_{\infty}(X)$ lui-même et aux $\Ch_{\infty}(X; a, b)$, que l'automorphisme induit par $T^2$ est l'identité dans $\mathrm{Hot}^{\wedge}_{0}$. J'espère que $\mathrm{Hot} \longrightarrow \mathrm{Hot}^{\wedge}_{0}$ est
\msp{page 150}
tout au moins fidèle, de sorte que ça prouvera que $T^2$ induit l'identité dans Hot également.

Je reviens aux questions $1^o$ à $4^o$ d'il y a quelques jours.

Pour $1^o$), il faut prouver que
\[
p : \Ch_{\infty}(X) \longrightarrow X \times X
\]
est lisse et propre et une Hot-fibration. Par dualité, $p$ propre résultera de $p$ lisse. Mais $\Ch_{\infty}(X)$ est $\varinjlim$ des $\Ch_{n,m}(X)$, avec $n \longrightarrow -\infty$, $m \longrightarrow +\infty$, et où on peut se borner à $n$ et $m$ \textit{pairs.} Or on voit que pour $n$, $m$ pairs, i.e. $\tau_{n,m}$ à flèche initiale directe, à flèche finale

% ---- from Der-ChVII-456_p096-100.tex
\printedpage{94}
% head: A. Grothendieck : Dérivateurs, Ch. VII, Édition des Universités de Montpellier II et Paris VII
rétrograde, $\Ch_{n,m}(X) \simeq \Ch_{\tau_{n,m}}(X)$ est Cat-\textit{fibrée} sur $X \times X$, donc lisse (i.e. HOT-lisse), et cela passe aux $\varinjlim$ (\footnote[36]{\textbf{NB} $\Ch_\infty(X)$ n'est ni fibrée ni cofibrée sur $X \times X$ - car les foncteurs d'inclusion entre les $\Ch_{n,m}$ ne sont pas cartésiens resp. cocartésiens.}).

Ainsi $\Ch_\infty(X) \xrightarrow{p} X \times X$ est à la fois propre et lisse (\footnote[37]{12.11. Mais ça implique déjà que c'est une fibration !}). Dire que c'est une HOT-fibration, revient p. ex. à dire que pour
\[
(a,b) \xrightarrow{(u,v)} (a',b')
\]
dans $X \times X$, la flèche correspondante
\[
\Ch_\infty(X)/(a,b) \longrightarrow \Ch_\infty(X)/(a',b')
\]
est une Hot-équivalence. Or, comme $p$ est propre, les inclusions des fibres en $(a,b)$ et en $(a',b')$ dans l'un et l'autre membre sont des Hot-équivalences :
\msp{page 151}
\[
\begin{tikzcd}[row sep=large]
\Ch_\infty(X)/(a,b) \arrow[r, "\rho"] & \Ch_\infty(X)/(a',b') \\
\Ch_\infty(X;a,b) \arrow[u, hook, "\wr"] \arrow[r, dotted, "\varphi"] & \Ch_\infty(X;a',b') \arrow[u, hook, "\wr"]
\end{tikzcd}
\tag*{$(*)$}
\]
Ainsi on trouve une flèche en pointillés dans Hot, dont il faut prouver que c'est une Hot-équivalence. Mais en considérant
\[
u : a \longrightarrow a' \qquad v : b \longrightarrow b'
\]
comme des \textit{chemins} dans $X$, et
\[
u^o : a' \longleftarrow a
\]
comme un chemin rétrograde de $a'$ à $a$, la théorie de variance générale (dans IX) nous donne bel et bien un isomorphisme dans (Hot) :
\[
\Ch_\infty(X;a,b) \longrightarrow \Ch_\infty(X;a',b'),
\]
on a envie de dire que c'est $\varphi$, i.e. qu'en l'insérant à la place de $\varphi$, le carré ci-dessus devient commutatif (ce qui prouverait ce qu'on veut). Je pense que si on travaille avec les ind-objets $\CCh_\infty(X)/(a,b)$, $\CCh_\infty(X;a,b)$ etc., pour lesquels on a des accouplements qui s'explicitent agréablement aux crans finis, et donnent naissance aux $\varphi$, cela devrait être facile (quoiqu'un peu fastidieux). Cela ne prouverait pas pour autant la commutativité de $(*)$ dans Hot, mais
% foot: Version 22 décembre 2001
\printedpage{95}
% head: A. Grothendieck : Dérivateurs, Ch. VII, Édition des Universités de Montpellier II et Paris VII
\msp{page 152}
seulement dans $\Hotz$, ce qui n'est peut-être pas suffisant. Mais cela suffit pour prouver que la flèche $\rho$ dans $(*)$ devient \texttt{[un]} isomorphisme dans $\Hotz$, donc elle est \texttt{[un]} isomorphisme.

\dgrblock{dgrb404}{Why both parities give the endpoint fibration}
\begin{dcomment}{exposition}{p150}{Why both parities give the endpoint fibration}
Grothendieck wants the categorical path object to carry the same geometric information as an ordinary path fibration, before any fibrant replacement of $X$. The two ends of a finite zigzag already explain how this can happen. For a positive-length interval, when both endpoints are sources in the interval category, changing their values contravariantly only composes the two outer arrows; this gives a Grothendieck fibration over $X\times X$. When both endpoints are sinks, the analogous operation is covariant and gives an opfibration. The two parity choices, with positive length imposed, are cofinal among the finite supports of an eventually constant path. The zero-length case is the diagonal functor and must not be used in the fibration claim.

The passage to the filtered union deserves a proof, since the inclusion functors need not preserve Cartesian arrows. For a functor $q:E\to B$, an object $e\in E$, and an arrow $h:b\to q(e)$, put
\[
 L_q(e,h)=\{(d,\gamma):q(d)=b,\ \gamma:d\to e,\ q(\gamma)=h\},
\]
with morphisms over $1_b$ commuting with $\gamma$. Smoothness means that every such category is weakly contractible; properness has the dual outgoing-lift test. These are the fibre-to-comma criteria of \citet[Theorem~1.2 and its dual]{Maltsiniotis2005Asphericite}. Suppose $E=\varinjlim_i E_i$ over a filtered category, with all functors over $B$. Once $e$ is represented at one stage, its lift category is the filtered colimit of the corresponding lift categories at later stages. Every object, arrow, and equation defining a finite simplex occurs at a stage, and every equality holds at some later stage. Nerves preserve these filtered colimits. A filtered colimit of weakly contractible simplicial sets is weakly contractible: after $\operatorname{Ex}^{\infty}$, which also commutes with filtered colimits, the assertions about based homotopy groups involve only finitely many simplices. Consequently smoothness and properness pass separately to this filtered colimit. Preserving Cartesian arrows is unnecessary.

It follows that $\Ch_\infty(X)\to X\times X$ is smooth and proper. We can also justify directly the later marginal observation that these two conditions imply the required homotopy-fibre property. For $h:b\to b'$ in $B$, let $M_h$ be the category of arrows $d\to d'$ in $E$ over $h$, with its two projections to $E_b$ and $E_{b'}$. Properness and smoothness show that these projections are weak equivalences. For example, the comma category under an object $d_0\in E_b$ retracts, by precomposition, onto its outgoing-lift category; the dual argument uses postcomposition. Quillen's Theorem~A applies. The arrows defining $M_h$ then give a natural transformation between its two maps into $b\backslash E$, the second obtained by precomposition from $b'\backslash E$. Since smoothness identifies each strict fibre with its under-comma category, precomposition
\[
 b'\backslash E\longrightarrow b\backslash E
\]
is a weak equivalence. Quillen's Theorem~B identifies the strict fibres with homotopy fibres \citep[\S1, Theorem~B]{Quillen1973}. This proof also survives every base change. It establishes what the manuscript needs without the finite-test conservativity invoked on the following pages.
\end{dcomment}


\dgrblock{dgrb405}{A homotopy fibration without Cartesian lifts}
\begin{dcomment}{examples}{p151}{A homotopy fibration without Cartesian lifts}
The warning in footnote~36 can be seen in the walking-arrow category $X=[1]$. It explains why a contractible category of choices is the appropriate conclusion, even when no best choice exists.

Consider the constant eventual path $c=1$ and the base arrow $(0\to1,1_1):(0,1)\to(1,1)$. An incoming lift is precisely an eventual alternating path $d$ from $0$ to $1$: its natural transformation to the constant path $1$ is unique. If the endpoint map were a Grothendieck fibration, this lifting category would have a terminal object. Its arrows are pointwise inequalities between the $0$--$1$ diagrams. For any candidate $d$, choose a direct edge far enough to the left that $d$ is still constantly $0$ there. Let $e$ change from $0$ to $1$ across that edge and remain $1$ afterwards. This is an allowed alternating path, because all its other edges are identities. At a vertex immediately after the chosen edge, $e$ is $1$ and $d$ is $0$, so there is no natural transformation $e\to d$. No terminal lift exists.

Dually, take the constant path $0$ and the outgoing base arrow $(1_0,0\to1)$. A cocartesian lift would give an initial object among paths from $0$ to $1$. For any candidate $d$, choose a single-step path whose direct transition is farther right than the point where $d$ has become constantly $1$. There is no arrow from $d$ to this later-step path. Thus no initial lift exists either. The incoming and outgoing lift categories are nevertheless weakly contractible by the filtered-stage proof. This is the exact distinction between the two homotopical lifting properties proved here and an ordinary fibration or opfibration.
\end{dcomment}


2$^o$) Il faut montrer que les HOT-fibrations
\[
\Ch_\infty(X) \overset{\sigma}{\underset{\beta}{\rightrightarrows}} X
\]
sont triviales, ou ce qui revient au même, qu'elles sont à fibres asphériques. Il suffit de le voir pour $\sigma$, donc que
\[
\Ch_\infty(X;a,?)
\]
est asphérique. Mais cette catégorie est la limite inductive des $\Ch_n(X;a,?)$, i.e. des
\[
\underline{\mathrm{Hom}}^{\text{source fixée}}(\tau^n, X).
\]
Mais comme l'intervalle $I_{\tau^n}$ se rétracte par déformation sur son origine, il est immédiat que cette catégorie \underline{Hom} est même contractile.

3$^o$) Considérons
\[
\begin{tikzcd}[column sep=small, row sep=large]
 & \Xb = C(f) \arrow[dl] \arrow[dr, "g"] \arrow[ddrr, bend left, "\bar f"] & & \\
X' \arrow[ur, dotted, bend left, "\alpha"] \arrow[dr, "f"'] & \underline{\text{cart.}} & \Ch_\infty(X) \arrow[dl, "\sigma"'] \arrow[dr, "\beta"] & \\
 & X & & X,
\end{tikzcd}
\]
\msp{page 153}
je voudrais prouver que $\bar f$ est une \textbf{D}-fibration propre et lisse (\footnote[38]{Ça résulte formellement de ce qui a été déjà vu, cf. pages 170-171.}).

\dgrblock{dgrb406}{The contraction with one endpoint fixed}
\begin{dcomment}{exposition}{p152}{The contraction with one endpoint fixed}
The next task is to make a path starting at $a$ disappear while keeping its starting point fixed. This is stronger than knowing merely that the interval has contractible nerve, because an arbitrary weak equivalence of interval categories need not induce a weak equivalence of strict functor categories. Here the interval has an explicit categorical deformation, and that is enough.

Write its vertices as $0,\ldots,r$, in their geometric order, with either orientation on each edge. For $0\leq k\leq r$, define $r_k(j)=\min(j,k)$. Each $r_k:I\to I$ is a functor. Successive $r_k,r_{k+1}$ are connected by a natural transformation in the direction of the edge between $k$ and $k+1$: on the tail $j\geq k+1$ use that edge, and elsewhere use identities. Naturality follows on the generating edges, hence on all arrows. Every $r_k$ fixes $0$, while $r_0$ is constant and $r_r=1_I$. Precomposition therefore joins the identity of $\operatorname{Fun}(I,X)_{a,?}$ to the constant functor at the constant path $a_I$, through a finite zigzag of natural transformations which is relative to the source. Its nerve is contractible.

Filtered colimits of these one-endpoint categories are weakly contractible by the finite simplicial argument just given. Thus both $\sigma$ and $\beta$ have weakly contractible strict fibres. Since they are smooth and proper, their pullbacks along arbitrary functors are weak equivalences. In the mapping-path construction, this universal statement is precisely what allows the category $X'$ to replace $X$ at the source endpoint. No supposition that $X$ itself is a Thomason fibrant category enters this argument.
\end{dcomment}


$\Xb$ est la catégorie des couples
\[
(x',c) \ \in \ \mathrm{Ob}\,X' \times \mathrm{Ch}_\infty(X)
\]
tels que
\[
\sigma(c) = f(x')
\]
et
\[
\bar f(x',c) = \beta(c), \qquad \text{but de } c.
\]
La fibre de $\Xb$, $\bar f$ en $x_0 \in X$ est donc la catégorie des couples $(x',c) \in \mathrm{Ob}\,X' \times \mathrm{Ch}_\infty(X)$, tels que
\[
\sigma(c) = f(x'), \qquad \beta(c) = x_0.
\]
% foot: Version 22 décembre 2001

\printedpage{96}
% head: A. Grothendieck : Dérivateurs, Ch. VII, Édition des Universités de Montpellier II et Paris VII
On peut regarder encore $\Xb$ comme limite inductive des sous-catégories $\Xb_{n,m}$, correspondante aux couples $(x',c)$ avec $c \in \mathrm{Ch}_{n,m}(X)$. On trouve que pour $m$ pair (direction terminale de $\tau_{n,m}$ rétrograde) les $\Xb_{n,m}$ sont Cat-fibrées sur $X$, pour $m$ impair (direction terminale de $\tau_{n,m}$ directe) elles sont Cat-cofibrées. Donc on conclut encore comme tantôt, $\Xb$ étant (au choix) limite inductive filtrante de sous-catégories lisses sur $X$, ou propres sur $X$. % [sic] correspondante -> correspondantes, agreement with « sous-catégories »

\msp{page 154}
Il faut prouver enfin que $\Xb \xrightarrow{\bar f} X$ est une \textbf{D}-fibration (\footnote[39]{Cela résulte déjà du fait que c'est propre et lisse!}), ce qui revient p. ex. à dire que pour une flèche
\[
a \xrightarrow{v} b
\]
dans $X$, la flèche induite
\[
\Xb/a \xrightarrow{\bar v} \Xb/b
\]
est une Hot-équivalence. On sait déjà comme tantôt que l'on a un diagramme à flèches verticales des Hot-équivalences
\[
\begin{tikzcd}[row sep=large]
\Xb/a \arrow[r] & \Xb/b \\
\Xb_a \arrow[u, hook, "\wr"] \arrow[r, dotted] & \Xb_b, \arrow[u, hook, "\wr"]
\end{tikzcd}
\]
et tout revient à montrer qu'une certaine flèche entre $\Xb_a$ et $\Xb_b$ dans Hot est \texttt{[un]} isomorphisme. J'ai envie de définir directement une telle flèche dans Hot, voire dans Cat. Or pour les $m$ impairs, les $\Xb_{n,m}$ sont Cat-cofibrées sur $X$. Pour $\Xb_{-m,m}$ ($m$ impair) on trouve donc
\[
(\Xb_{-m,m})_a \xrightarrow{v_*} (\Xb_{-m,m})_b,
\]
\msp{page 155}
l'ennui est qu'elles ne sont pas compatibles avec les morphismes de transition pour $m$ variable. Elles le sont cependant \textit{à homotopie près} sûrement, donc on trouvera au moins une flèche d'ind-objets dans Cat\,hot
\[
\Xbb_a \xrightarrow{v_*} \Xbb_b,
\]
ce qui par passage à $\Hotz$ donne quand-même une flèche
\[
\Xb_a \xrightarrow{v_*} \Xb_b.
\]
En travaillant avec les $m$ impairs, on trouve de même des foncteurs en sens inverses % [sic] m impairs -> m pairs, the Cat-fibrées (giving v^*) are stated above for m pair
\[
(\Xb_{-m,m})_a \xleftarrow{v^*} (\Xb_{-m,m})_b,
\]
% foot: Version 22 décembre 2001
\printedpage{97}
% head: A. Grothendieck : Dérivateurs, Ch. VII, Édition des Universités de Montpellier II et Paris VII
qui ne se recollent pas pour $m$ variable, mais qui se recollent à homotopie près, pour donner une flèche dans Ind\,Cat\,hot
\[
\Xbb_a \xleftarrow{v^*} \Xbb_b,
\]
d'où aussi
\[
\Xb_a \xleftarrow{v^*} \Xb_b \qquad \text{dans } \Hotz.
\]
D'autre part, on doit pouvoir montrer soritalement,
\msp{page 156}
cran par cran, que $v_*v^*$ et $v^*v_*$ dans Ind\,Cat\,hot sont égales à l'identité, donc $v_*$ et $v^*$ sont des isomorphismes inverses l'un de l'autre. On doit vérifier d'autre part soritalement que
\[
\begin{tikzcd}[row sep=large]
\Xbb/a \arrow[r, "\bar v"] & \Xbb/b \\
\Xbb_a \arrow[u] \arrow[r, "v_*"] & \Xbb_b, \arrow[u]
\end{tikzcd}
\]
dans Ind\,Cat\,hot est commutatif. Passant au diagramme correspondant dans $\Hotz$, en prenant les limites inductives, on trouve un carré commutatif dans $\Hotz$
\[
\begin{tikzcd}[row sep=large]
\Xb/a \arrow[r, "\bar v"] & \Xb/b \\
\Xb_a \arrow[u, "\wr"'] \arrow[r, "v_*", "\sim"'] & \Xb_b, \arrow[u, "\wr"']
\end{tikzcd}
\]
d'où résulte que $\bar v$ est \texttt{[un]} isomorphisme dans $\Hotz$, donc aussi dans Hot.

\dgrblock{dgrb407}{Evrard's categorical cocylinder}
\begin{dcomment}{later}{p153}{Evrard's categorical cocylinder}
Grothendieck is seeking an explicit factorization of an arbitrary functor by adjoining paths. The same purpose already appears in \citet[\S I.1]{Evrard1975}, whose variable-length zigzags carry both natural transformations and maps which insert identity segments. The indexing conventions differ from those used here; an identification of the two categories would therefore require a comparison functor, not just the observation that their objects look alike.

A precise modern theorem is \citet[Theorem~2.8]{Shoikhet2016}: the categorical cocylinder gives a functorial factorization by a weak equivalence followed by a homotopy-fibred functor, and there is a dual construction. Its relation with the manuscript is the role of the second map: after applying the nerve, ordinary fibres compute homotopy fibres. Shoikhet's Proposition~1.14 formulates a stronger, base-change-stable version through simplicial sharp maps, under the prefibration hypotheses stated there. We use Quillen's comma-category theorem directly when those strict prefibration hypotheses have not been established. These distinctions matter for the total category $\Ch(X)$, whose type-changing arrows will obstruct properness below, even though its fixed-endpoint categories still have the correct homotopy types.
\end{dcomment}


Pour vérifier que la HOT-fibration $\bar f$ est triviale si et seulement si $f$ est une HOT-équivalence, il faudrait avoir développé dans Cat le formalisme de la suite exacte des lacets - on est tout près visiblement !
% foot: Version 22 décembre 2001

\printedpage{98}
% head: A. Grothendieck : Dérivateurs, Ch. VII, Édition des Universités de Montpellier II et Paris VII
\msp{page 157}

\dgrblock{dgrb408}{When the replacement is trivial}
\begin{dcomment}{exposition}{p156}{When the replacement is trivial}
The desired test for triviality is already formal once the mapping-path construction has been established. Put $r:C(f)\to X'$ for the projection which forgets the path. It is the pullback of the universally weakly equivalent source map $\sigma$, so $r$ is a weak equivalence. Its section $i_f$ is consequently a weak equivalence as well. The equality $\bar f i_f=f$ gives
\[
 \bar f\text{ is a weak equivalence}\quad\Longleftrightarrow\quad
 f\text{ is a weak equivalence}
\]
by two-out-of-three. If triviality is expressed instead by weakly contractible fibres, the homotopy-fibre theorem gives the same criterion. One can see both directions by replacing the realized map by a Serre fibration: a weak equivalence has contractible homotopy fibres, and a fibration with such fibres induces isomorphisms on every based homotopy group and a bijection on components.

Thus an independently developed long exact sequence for categorical loop objects is not needed for this particular conclusion. Nor should the preceding finite-test inference bear its proof. The claimed conservativity of the unbased restricted Yoneda functor is false in the generality used in the manuscript; the commentary on the finite-test obstruction in Chapter~V, manuscript p.~52, explains this. The two-out-of-three argument, together with the actual endpoint-fibration theorem, proves exactly the assertion required here.
\end{dcomment}


\begin{flushright}
\footnotesize 13.11.90
\end{flushright}

\section*{13\quad Où on prouve enfin que $\Ch(X;-,-)$ est aussi bon (voire meilleur) que $\Ch_\infty(X;-,-)$.}

Je crois finalement que les catégories $\Ch(X)$, $\Ch(X;a,b)$ (sans indice $\infty$!) sont `bonnes', que les flèches canoniques dans Hot
\[
\left\{
\begin{aligned}
\Ch_\infty(X) \quad &\xrightarrow{\sim} \quad \Ch(X)\\
\Ch_\infty(X;a,b) \quad &\xrightarrow{\sim} \quad \Ch(X;a,b)
\end{aligned}
\right.
\]
sont des isomorphismes. (La première est même \texttt{[un]} isomorphisme dans $\mathrm{HOT}(X \times X)$, ce qui implique la deuxième relation, en passant aux fibres \dots) Je crois que je suis à présent (après XI) en mesure de le prouver. Cela résultera des deux résultats-clef suivants :

\begin{enumerate}
\item[1)] Le foncteur canonique
\[
X \longrightarrow \Ch(X)
\]
(associant à tout $x \in \mathrm{Ob}\,X$ le chemin vide de source et but $X$) est une Hot-équivalence, et : % [sic] but $X$ -> but $x$, the object is $x \in \mathrm{Ob}\,X$
\item[2)] Le foncteur source-but
\[
\Ch(X) \longrightarrow X \times X
\]
est \textit{lisse} et \textit{propre} (donc une `Hot-fibration parfaite').
\end{enumerate}

Cela doit donner ce qu'on cherche. Comme la $X \times X$-flèche $\Ch_\infty(X) \longrightarrow \Ch(X)$ (et les flèches déduites $\Ch_\infty(X;a,b) \longrightarrow \Ch(X;a,b)$) a été définie par l'intermédiaire de la catégorie
\[
\mathcal{C}_{\mathbf{N}^o} = \mathbf{N}^o \times_{\Ty} \mathcal{C}
\qquad
\begin{aligned}
&\text{(où $\mathcal{C} = \Ch(X)$ \texttt{[est]} considérée comme catégorie fibrée sur $\Ty$,}\\
&\text{et $\mathbf{N}^o \longrightarrow \Ty$ est le système projectif des $\underbrace{(\Delta_+\Delta_-)^{2n}}_{t_n} \dots$)}
\end{aligned}
\]
\msp{page 158}
à l'aide des deux flèches canoniques
\[
\begin{tikzcd}[row sep=huge, column sep=large]
\mathcal{C}_{\mathbf{N}^o} \arrow[r] \arrow[d] & \Ch_\infty(X) \\
\Ch(X), &
\end{tikzcd}
\]
% foot: Version 22 décembre 2001

\dgrblock{dgrb409}{A properness claim which the empty path disproves}
\begin{dcomment}{development}{p157}{A properness claim which the empty path disproves}
The return to paths of arbitrary type is motivated by their strictly associative concatenation. Their geometric usefulness is real, but the proposed proof gives them a stronger lifting property than they possess. We isolate the error before using the comparison diagram.

\textbf{The empty-path obstruction.}\ours\quad
For $X=[1]=(0\xrightarrow{w}1)$, neither
\[
 \Ch(X)\longrightarrow X\times X
 \qquad\text{nor}\qquad
 \mathcal C_{\mathbf N^o}\longrightarrow X\times X
\]
is proper, with the morphisms defined in this chapter.

\emph{Proof.} Let $z$ be the empty path at $0$, lying above $(0,0)$, and consider the base arrow $(1_0,w):(0,0)\to(0,1)$. Properness requires the category of its outgoing lifts with source $z$ to be weakly contractible. In $\Ch(X)$, a morphism from a path of type $e$ to a path of type $\tau$ lies over an endpoint-preserving type map $e\to\tau$. The single vertex of $e$ would have to map to both endpoints of $\tau$. Hence $\tau=e$. But an empty path has equal endpoints and cannot lie above $(0,1)$. The lift category is empty. For $\mathcal C_{\mathbf N^o}$ the same argument is even more immediate: an arrow from stage $0$ can only end at stage $0$, because its image in $\mathbf N^o$ has source $0$. Again its target must have equal endpoints. This proves both failures.\hfill$\square$

The opposite of a total Grothendieck construction reverses the arrows in its indexing category as well as the arrows inside its fibres. Reversing a path or replacing $X$ by $X^o$ does not perform that operation on $\Ty$ or on $\mathbf N^o$. This is the missing distinction in the deductions of properness on pp.~161 and 169. The same lifting test is the dual of \citet[Theorem~1.2]{Maltsiniotis2005Asphericite}.

This counterexample does not contradict the comparison of fixed-endpoint homotopy types which opens the section. That conclusion has a repair: smoothness of $\Ch(X)$, together with explicit cancellation of a path and its reverse, gives the homotopy-fibre theorem by Quillen~B. The intermediary $\mathcal C_{\mathbf N^o}$ can then be compared with $\Ch_\infty(X)$ using its filtered-colimit construction. We give the full argument below. We did not find this empty-path counterexample or this precise repair of the present pages in the primary sources checked, namely Evrard, Shoikhet, and Maltsiniotis; the mathematical ingredients of the repair are classical.
\end{dcomment}


% ---- from Der-ChVII-456_p101-105.tex
\printedpage{99}
% head: A. Grothendieck : Dérivateurs, Ch. VII, Édition des Universités de Montpellier II et Paris VII
on sera obligé de prouver aussi les résultats qui correspondent à 1) 2) pour $\cC_{\NN}$, i.e.

\begin{itemize}[label={}]
\item[1$'$)]
\[
\begin{tikzcd}
X \ar[r] \ar[d,"\wr"'] & \cC_{\NN} \\
\cC_{t_0} = \cC_0 &
\end{tikzcd}
\]
(l'inclusion canonique de $\cC_0$, fibre de $\cC_{\NN}$ (fibrée sur $\mathbf{N}^o$) en $0 \in \mathbf{N}^o$) est une Hot-équivalence, et :

\item[2$'$)] $\cC_{\NN} \longrightarrow X \times X$ est propre et lisse.
\end{itemize}

Moyennant ces quatre résultats techniques, que je vais prouver tantôt (et $1'$ et $2'$ ne font aucune surprise pour moi), on gagne. Considérons en effet le diagramme
\[
\begin{tikzcd}
& X \ar[dl,"\alpha_1"'] \ar[d,"\alpha_0"] \ar[dr,"\alpha_2"] & \\
\Ch(X) \ar[dr,"p_1"'] & \cC_{\NN} \ar[l,"\pi_1"'] \ar[r,"\pi_2"] \ar[d,"p_0"] & \Ch_{\infty}(X) \ar[dl,"p_2"] \\
& X \times X. &
\end{tikzcd}
\]

Par 1) et $1'$) (concernant $\alpha_1$ et $\alpha_0$) et le résultat
\msp{page 159}
similaire déjà connu pour $\alpha_2$, ces trois foncteurs sont des Hot-équivalences. Donc il en est de même de $\pi_1$, $\pi_2$. Comme $p_1$, $p_0$, $p_2$ sont des Hot-fibrations (étant propres et lisses), il s'ensuit que $\pi_1$ et $\pi_2$ induisent des Hot-équivalences sur les fibres
\[
\Ch(X;a,b) \quad \xleftarrow{\ \pi_1(a,b)\ } \quad \cC_{\NN}(a,b) \quad \xrightarrow{\ \pi_2(a,b)\ } \quad \Ch_{\infty}(X;a,b).
\]
Je l'avais vérifié (sans mal) déjà pour $\pi_2$ et les $\pi_2(a,b)$, ce qui avait permis, en inversant $\pi_2$ dans Hot, de définir \textit{dans} Hot
\[
\left\{
\begin{aligned}
\Ch_{\infty}(X) &\quad\longrightarrow\quad \Ch(X) \\
\Ch_{\infty}(X;a,b) &\quad\longrightarrow\quad \Ch(X;a,b),
\end{aligned}
\right.
\]
et dire que ce sont des isomorphismes signifie justement que $\pi_1$ et les $\pi_1(a,b)$ sont des isomorphismes.

Je ne crois pas que $\cC_{\NN}$ ait en lui-même d'autre interêt, que de permettre de comparer % [sic] interêt -> intérêt, missing acute accent
les $X \times X$-catégories $\Ch_{\infty}(X)$ et $\Ch(X)$, à l'aide du diagramme précédent. Par contre, $\Ch_{\infty}(X)$ et $\Ch(X)$ m'aparaissent comme les deux ingrédients essentiels du formalisme des % [sic] m'aparaissent -> m'apparaissent, spelling
chemins dans Cat.
% foot: Version 22 décembre 2001

\printedpage{100}
% head: A. Grothendieck : Dérivateurs, Ch. VII, Édition des Universités de Montpellier II et Paris VII
\msp{page 160}
Je vais prouver 1, 2, $1'$, $2'$.

1) Considérons le diagramme
\[
\begin{tikzcd}
X = (X_{\Ty})_e \ar[rr,hook,"i"] \ar[dr,dotted] & & X_{\Ty} = X \times \Ty \ar[rr,"j"] \ar[dr] & & \Ch(X) \ar[dl] \\
& e \ar[rr,dotted] & & \Ty\ , &
\end{tikzcd}
\]
où $X_{\Ty}$ et $\Ch(X)$ sont considérées comme catégories Cat-fibrées sur $\Ty$, et le foncteur $j$ est déduit des foncteurs canoniques
\[
X \longrightarrow \Ch(X)_{\tau} \overset{\text{déf}}{=} \Ch_{\tau}(X),
\]
associant à $x \in \Ob X$ le chemin $p_{\tau}^{*}(x)$, où $p_{\tau} : \tau \longrightarrow e$ est la projection de $\tau$ vers l'élément unité (et élément final) $e = 1_{\Ty}$ de $\Ty$, et $\Ch_{e}(X)$ est identifiée à $X$. Il faut vérifier le

\textbf{Lemme.} \textit{Pour tout $\tau \in \Ty$, le foncteur précédent est une} Hot-\textit{équivalence (et même un homotopisme).}

Mais on sait que pour deux catégories $\cX$, $\cC$ lisses sur une autre, $T$, tout $T$-morphisme de $\cX$ dans $\cC$ qui induit des Hot-équivalences sur les fibres, est une Hot-équivalence. Ainsi $j$ est
\msp{page 161}
est une Hot-équivalence. D'autre part, $\Ty$, ayant un objet final, est asphérique, donc % [sic] est ... est -> est, word repeated across the manuscript page break
$X \times \Ty \xrightarrow{\mathrm{pr}_1} X$ est une Hot-équivalence, d'où résulte aussitôt que $i$ est une Hot-équivalence. Donc de même $ji$, or $ji = \alpha_1$, q.e.d.

\dgrblock{dgrb410}{Completing the constant-path lemma}
\begin{dcomment}{exposition}{p160}{Completing the constant-path lemma}
The first comparison is a statement about total homotopy types, and here Grothendieck's argument works. Its useful feature is that it does not ask the chosen sequence of interval types to be homotopy cofinal in $\Ty$.

For every finite oriented interval $I$, evaluation at the initial vertex and the constant-diagram functor give
\[
 X\xrightarrow{c}\operatorname{Fun}(I,X)
 \xrightarrow{\operatorname{ev}_0}X,
 \qquad \operatorname{ev}_0c=1_X.
\]
The truncation functors $r_k(j)=\min(j,k)$ constructed above join $c\operatorname{ev}_0$ to the identity by a finite zigzag of natural transformations. This proves the lemma promised here and left unwritten on p.~170: $c$ is a categorical homotopy equivalence, not just a weak equivalence of nerves. The argument applies to every sign pattern and to the zero-length interval.

These constant-diagram functors commute strictly with precomposition by type maps. They therefore assemble to the displayed Cartesian functor
\[
 j:X\times\Ty\longrightarrow\Ch(X)
\]
between Grothendieck fibrations over $\Ty$, and each of its fibre functors is a weak equivalence. The standard fibrewise comparison theorem for Grothendieck constructions makes $j$ a weak equivalence. This is also the objectwise invariance behind Thomason's homotopy-colimit theorem; the fibred convention is recalled explicitly in \citet[\S1.5]{Shoikhet2016}. Since $\Ty$ has terminal object $e$, the inclusion $X\to X\times\Ty$ at $e$ is a weak equivalence. Their composite is the empty-path section $X\to\Ch(X)$.

Exactly the same proof over $\mathbf N^o$ proves $X\to\mathcal C_{\mathbf N^o}$ is a weak equivalence. For $\Ch_\infty(X)$ the constant-diagram functors assemble to a levelwise weak equivalence of filtered diagrams; invariance of weak equivalences under filtered colimits gives the corresponding assertion. Consequently the two comparison arrows out of $\mathcal C_{\mathbf N^o}$ are weak equivalences on total categories, by two-out-of-three. What remains is genuinely relative: a weak equivalence over $X\times X$ need not be a weak equivalence on its strict fibres until those fibres have been identified as homotopy fibres.
\end{dcomment}


2) Prouvons p. ex. la lissité de $\Ch(X) \longrightarrow X \times X$, la propriété s'en déduira en appliquant ce résultat à $X^{\mathrm{opp}}$. Soit donc
\[
(a',b') \ \longrightarrow\ (a,b), \qquad \text{i.e. } a' \xrightarrow{u} a,\ \ b' \xrightarrow{v} b,
\]
une flèche dans $X \times X$, et $c$ un objet de $\Ch(X)$ au dessus de $(a,b)$ :
\[
\begin{tikzcd}[column sep=large]
a' \ar[r,dotted,"c'","\scriptstyle\rightarrow"'] \ar[d,"u"'] & b' \ar[d,"v"] \\
a \ar[r,"c","\scriptstyle\rightarrow"'] & b.
\end{tikzcd}
\]

On regarde la catégorie $\cC(c;u,v)$ des couples $(c',w)$ qui relèvent $(c,(u,v))$, où donc $c'$ est un chemin (de type non précisé) de $a'$ vers $b'$, et
\[
w : c' \longrightarrow c
\]
une flèche dans $\Ch(X)$ qui sur les objets source induise $u$, et sur les objets buts induise $v$ (ainsi $c'$ est un chemin de $a'$ vers $b'$). Il faut
% foot: Version 22 décembre 2001
\printedpage{101}
% head: A. Grothendieck : Dérivateurs, Ch. VII, Édition des Universités de Montpellier II et Paris VII
\msp{page 162}
prouver que cette catégorie $\cC(c;u,v)$ est asphérique. Pour cela, on regarde le chemin composé
\[
\tilde{c} = \mathrm{ch}(u)\ c\ \mathrm{ch}(v)^{o},
\]
où $\mathrm{ch}(u)$ est le chemin (direct) de $a'$ vers $a$ défini par $u : a' \longrightarrow a$, et $\mathrm{ch}(v)^{o}$ désigne donc le chemin (rétrograde) de $b$ à $b'$ défini par $v : b' \longrightarrow b$. Si $c$ est de type $\tau$, $\tilde{c}$ est donc de type $\Delta_{+}\tau\Delta_{-}$. On va comparer $\cC(c;u,v)$ et la catégorie $\Ch(X)/\tilde{c}$, à l'aide de deux foncteurs
\[
\begin{tikzcd}
\cC(c;u,v) \ar[r,shift left,"\varphi"] & \Ch(X;a',b')/\tilde{c}\ , \ar[l,shift left,"\psi"]
\end{tikzcd}
\]
dont on prouvera qu'ils définissent des flèches inverses l'une de l'autre dans Hot. Comme $\Ch(X)/\tilde{c}$ est contractile (ayant un objet final), donc asphérique, il en est de même de $\cC(c;u,v)$, et on gagnera.

On définit $\varphi$ par $\boxed{\varphi(c',w') = (\tilde{c}',\tilde{w}'),}$ où ($c'$ étant du type $\tau'$) $\tilde{c}'$ est le chemin de type $\Delta_{+}\tau'\Delta_{-}$ déduit de $c'$ comme image inverse par la dégénérescence $\Delta_{+}\tau'\Delta_{-} \longrightarrow \tau'$, i.e. on prolonge $c'$ d'un `cran $\Delta_{+}$' $a' \xrightarrow{\mathrm{id}} a'$ à gauche et d'un `cran $\Delta_{-}$' $b' \xleftarrow{\mathrm{id}} b'$ à droite, cf. figure.
\[
\overbrace{\underbrace{\begin{tikzcd}[ampersand replacement=\&]
a' \ar[r,"\mathrm{id}"] \ar[d,"\mathrm{id}"'] \ar[dr,phantom,"\ast"] \& a' \ar[rr,dotted,no head,"c'"] \ar[d,"u"'] \& {} \ar[d,dotted,"w'"] \& b' \ar[d,"v"'] \ar[dr,phantom,"\ast"] \& b' \ar[l,"\mathrm{id}"'] \ar[d,"\mathrm{id}"] \\
a' \ar[r,"u"] \& a' \ar[rr,no head,"c'"'] \& {} \& b' \& b' \ar[l,"v"'] % [sic] a' ... b' with label c' in the bottom row -> a ... b with label c, the bottom row of the figure is the path c~ = ch(u) c ch(v)^o
\end{tikzcd}}_{\tilde{c}}}^{\tilde{c}'}
\]

\msp{page 163}
On peut dire que
\[
\tilde{c}' = a'_{\Delta_{+}} \circ c' \circ b'_{\Delta_{-}},
\]
où de façon générale, si $x \in X$ et $\tau \in \mathrm{Ty}$, $x_{\tau}$ désigne le `chemin constant' de valeur $x$, type $\tau$. On a d'autre part des morphismes évidents
\[
a'_{\Delta_{+}} \longrightarrow \mathrm{ch}(u), \qquad b'_{\Delta_{-}} \longrightarrow \mathrm{ch}(v)^{\mathrm{opp}},
\]
comme montré [?] par les flèches verticales $(\mathrm{id}_{a'},u)$ et $(v,\mathrm{id}_{b'})$ dans les deux carrés marqués $\ast$ du diagramme, flèches qui sont l'une dans $\Ch_{\Delta_{+}}(X)$, l'autre dans $\Ch_{\Delta_{-}}(X)$. Ceci dit, on pose
\[
\tilde{w}' : \tilde{c}' \longrightarrow \tilde{c}, \qquad \tilde{w} = \lambda(u) \ast w' \ast \mu(v), % [sic] \tilde{w} -> \tilde{w}', the morphism just named is \tilde{w}'
\]
où
\[
\left\{
\begin{aligned}
\lambda(u) :\ \ a'_{\Delta_{+}} &\quad\longrightarrow\quad \mathrm{ch}(u) \\
\mu(v) :\ \ b'_{\Delta_{-}} &\quad\longrightarrow\quad \mathrm{ch}(v)^{\mathrm{opp}}
\end{aligned}
\right.
\]
% foot: Version 22 décembre 2001
\printedpage{102}
% head: A. Grothendieck : Dérivateurs, Ch. VII, Édition des Universités de Montpellier II et Paris VII
(cf. figure). Le couple
\[
(\tilde{c}',\tilde{w}') \in \Ob \Ch(X)/\tilde{c}
\]
ainsi défini dépend fonctoriellement de $(c',w')$ variable dans $\cC(c;u,v)$ : Si $(c'_1,w'_1)$ est un deuxième objet de cette catégorie, une flèche du premier dans le second est une flèche
\[
c' \xrightarrow{\alpha} c'_1
\]
dans $\Ch(X;a',b')$, telle qu'on ait commutativité dans
\[
\begin{tikzcd}
c' \ar[rr,"\alpha"] \ar[dr,"w'"'] & & c'_1 \ar[dl,"w'_1"] \\
& c, &
\end{tikzcd}
\qquad\qquad w' = w'_1\alpha.
\]

Une telle flèche $\alpha$ définit
\[
\tilde{\alpha} = a'_{\Delta_{+}} \ast \alpha \ast b'_{\Delta_{-}} : \tilde{c}' \longrightarrow \tilde{c}'_1
\]
dans $\Ch(X;a',b')$, et on vérifie
\msp{page 164}
aussitôt que l'on a commutativité dans
\[
\begin{tikzcd}
\tilde{c}' \ar[rr,"\tilde{\alpha}"] \ar[dr,"\tilde{w}'"'] & & \tilde{c}'_1 \ar[dl,"\tilde{w}'_1"] \\
& \tilde{c}, &
\end{tikzcd}
\qquad\qquad \text{i.e. } \tilde{w}' = \tilde{w}'_1\tilde{\alpha},
\]
i.e. que $\tilde{\alpha}$ est bien une flèche dans $\Ch(X;a',b')/\tilde{c}$. De même, que $\alpha \longmapsto \tilde{\alpha}$ définit une dépendance fonctorielle, i.e. $(\beta\alpha)\tilde{} = \tilde{\beta}\tilde{\alpha}$, est évident sur la définition de $\tilde{\alpha}$. Donc, on a bien défini % Reading resolved against the supplied PDF; see corrections.csv.
\[
\cC(c;u,v) \ \xrightarrow{\ \varphi\ }\ \Ch(X;a',b')/\tilde{c},
\]
définissons maintenant
\[
\cC(c;u,v) \ \xleftarrow{\ \psi\ }\ \Ch(X;a',b')/\tilde{c}.
\]
Soit $(c'',w'')$ un objet de $\Ch(X;a',b')/\tilde{c}$,
\[
\overbrace{\underbrace{\begin{tikzcd}[ampersand replacement=\&]
a'_0 = a' \ar[r] \ar[d,"\mathrm{id}"'] \& a'_1 \ar[rr,no head,"\rightarrow"] \ar[d,"w''_1"'] \& {} \ar[d,dotted,"w''"] \& a'_n \ar[d,"w''_n"'] \& b' \ar[l] \ar[d,equal] \\
a' \ar[r,"u"] \& a \ar[rr,no head] \& {} \& b \& b', \ar[l,"v"']
\end{tikzcd}}_{c}}^{c''}
\]
% foot: Version 22 décembre 2001

\printedpage{103}
% head: A. Grothendieck : Dérivateurs, Ch. VII, Édition des Universités de Montpellier II et Paris VII
donc
\[
w'' : c'' \longrightarrow \tilde{c}, \qquad\quad \text{flèche dans } \Ch(X;a',b'),
\]
notons qu'on a une flèche de rétraction évidente
\[
\tilde{c} \xrightarrow{\rho_c} c \qquad\quad \text{associée à } \Delta_{+}\tau\Delta_{-} \longrightarrow \tau \text{ (rétraction),}
\]
induisant $u$ et $v$ sur la source et le but :
\[
\begin{tikzcd}
a' \ar[r,"u"] \ar[d,"u"] & a \ar[rrr,no head,"c"] \ar[d,"\mathrm{id}"] & {} \ar[d,"\mathrm{id}"] & {} \ar[d,"\mathrm{id}"] & b \ar[d,"\mathrm{id}"] & b' \ar[l,"v"'] \ar[d,"v"] & \Big\}\ \tilde{c} & \in \Ch_{\Delta_{+}\tau\Delta_{-}}(a',b') \\
a \ar[r,"\mathrm{id}"] & a \ar[rrr,no head,"c"] & {} & {} & b & b' \ar[l,"v"'] & \Big\}\ \delta^{*}_{\tau}(c) & \in \Ch_{\Delta_{+}\tau\Delta_{-}}(a,b) \\
& a \ar[rrr,no head,"c"] & & & b & & \Big\}\ c & \in \Ch_{\tau}(a,b).
\end{tikzcd}
\]

\msp{page 165}
Donc on définit un objet
\[
(c',w') = \psi(c'',w'')\ \in \Ob \cC(c;u,v)
\]
par
\[
c' = c'', \qquad\quad w' = \rho_c w'',
\]
et c'est fonctoriel en $(c'',w'')$.

\dgrblock{dgrb411}{The lifting category is contracted relative to its endpoints}
\begin{dcomment}{exposition}{p164}{The lifting category is contracted relative to its endpoints}
The long diagram calculation has a simple purpose: to replace the category of all incoming lifts of $(u,v)$ by a slice with a terminal object. It is a proof of smoothness, and no properness assertion is needed to make it work.

Fix $u:a'\to a$, $v:b'\to b$, and a path $c:a\rightsquigarrow b$. Write $A=\Ch(X;a',b')$ and $\widetilde c=\mathrm{ch}(u)c\mathrm{ch}(v)^o$. Its slice $A/\widetilde c$ has terminal object $1_{\widetilde c}$. The functor $\varphi$ pads an incoming lift $w':c'\to c$ with a constant edge at each end. Its new map to $\widetilde c$ has the components $(1_{a'},u)$ on the left and $(v,1_{b'})$ on the right. A morphism $\alpha:c'\to c'_1$ is padded by identities, so the equality $w'_1\alpha=w'$ becomes $\widetilde w'_1\widetilde\alpha=\widetilde w'$. Both functoriality and the slice condition follow from the interchange law for concatenation of natural transformations.

To recover a lift from a slice object, there is an arrow $\rho_c:\widetilde c\to c$ over the collapse $\Delta_+\tau\Delta_-\to\tau$, with endpoint components $u,v$. Define $\psi(c'',w'')=(c'',\rho_cw'')$. The identity which makes the argument relative is
\[
 \rho_c\widetilde w'=w'\rho_{c'}.
\]
It holds on the central interval by the definition of $\widetilde w'$, and on the two added edges by the two commutative endpoint squares. The maps $\rho_{c'}$ are natural because collapsing constant collars commutes with every endpoint-preserving type map. They give a natural transformation $\psi\varphi\Rightarrow1$ on the lifting category.

On the other side, both $\varphi\psi$ and the identity map naturally to the constant endofunctor with value $1_{\widetilde c}$, because the slice has a terminal object. Thus their nerves are homotopic. The two functors are homotopy inverses, and the lifting category is contractible. This proves smoothness of $\Ch(X)\to X\times X$ completely. The wrong primes in the first padding diagram and the wrong final vertex in the diagram of $\rho_c$ are printed errors; the corrected endpoint descriptions above specify the maps without ambiguity.
\end{dcomment}


Le composé $\varphi\psi$ étant un endomorphisme de la catégorie contractile $\Ch(X;a',b')/\tilde{c}$, est homotope à l'identité. Il faut vérifier qu'il en est de même de $\psi\varphi$, endomorphisme de $\cC(c;u,v)$. Ce foncteur associe à $(c',w')$, $(\tilde{c}',\overline{w'})$, où $\tilde{c}'$ est défini comme plus haut, $\tilde{c}' = a'_{\Delta_{+}} \circ c' \circ b'_{\Delta_{-}}$, et où $\overline{w'}$ est le composé
\[
\underbrace{\tilde{c}'}_{= \delta^{\prime *}_{\tau'}(c')}\ \ \underbrace{\xrightarrow{\ \rho_{c'}\ }}_{\substack{\text{flèche}\\ \text{canonique}}}\ \ c' \xrightarrow{w'} c.
\]

On a donc
\[
(\psi\varphi)(c',w') = (\tilde{c}',\overline{w'}) \xrightarrow{\rho_{c'}} (c',w'),
\]
c'est bien un morphisme dans $\cC(c;u,v)$, i.e. on a commutativité dans
\[
\begin{tikzcd}
\tilde{c}' \ar[rr,"\rho_{c'}"] \ar[dr,"\overline{w'}"'] & & c' \ar[dl,"w'"] \\
& c &
\end{tikzcd}
\qquad\qquad \text{(flèche de } \Ch(X;a',b')\text{).}
\]
% foot: Version 22 décembre 2001

% ---- from Der-ChVII-456_p106-109.tex
\printedpage{104}
% head: A. Grothendieck : Dérivateurs, Ch. VII, Édition des Universités de Montpellier II et Paris VII
\msp{page 166}

Ainsi, le foncteur $\psi\varphi$ est bien homotope au foncteur identique. On a gagné !

Passons à $1'$, $2'$) pour $\CN$. L'assertion $1'$) se démontre exactement de la même manière que 1) pour \underline{Ch}$(X)$, en regardant $\CN$ comme fibrée sur $\mathbf{N}^o$ (qui a un objet final 0), et en utilisant le diagramme
\[
\begin{tikzcd}
X = (X_{\mathbf{N}^o})_e \arrow[rr, "i'"] & & X_{\mathbf{N}^o} = X \times \mathbf{N}^o \arrow[rr, "j'"] \arrow[dr] & & \mathcal{C}_{\mathbf{N}^o} \arrow[dl] \\
 & & & \mathbf{N}^o \ , &
\end{tikzcd}
\]
où $j'$ est (cartésien et) une Hot-équivalence sur chaque fibre d'après le lemme, donc une Hot-équivalence globalement, et $i'$ est une Hot-équivalence car $\mathbf{N}^o$ est asphérique (ayant un objet final). Donc $j'i'$ est asphérique, q.e.d.

Il reste à prouver que
\[ \CN \longrightarrow X \times X \]
est propre et lisse. Ici encore,
\msp{page 167}
c'est exactement le même argument que tantôt, je présume. Il y a cependant une complication technique. Nous avions pris, pour définir
\[ \underline{t} : \mathbf{N}^o \longrightarrow \underline{\mathrm{Ty}}, \qquad \underline{t} = (t_n)_n, \]
le choix
\[ t_n = (\Delta_+\Delta_-)^{2n}, \]
ou si on préfère
\[ t_n = (\Delta_+\Delta_-)^n, \]
en tous cas les $t_n$ sont à direction initiale $+$, terminale $-$, et ce sont des types réduits. Si on prend $\Delta_+t_n\Delta_-$, on trouve un type non réduit, qui n'est pas parmi les $t_n$. Il ne servirait à rien de changer la définition de $t_n$ en $t_n = (\Delta_-\Delta_+)^{\text{quelque chose}}$, car cette fois $\Delta_+t_n\Delta_-$ sera bien réduit, mais n'aura pas les bons signes initial et terminal. Donc pour définir $\tilde{c}$ en termes de $c$, comme objet de $\CN$, il faudra, si $c$ est de type $\tau = t_n$, prendre $\tilde{c}$ de type $\Delta_+\Delta_-\tau\Delta_+\Delta_- = t_{n+1}$ (en prenant $t_{n+1} = (\Delta_+\Delta_-)^{2n}$), et de même pour définir $\tilde{c}'$ en termes de $c'$. Les définitions pertinentes du passage de $(c',w')$ à $(\tilde{c}',\tilde{w}')$ sont décrites dans la figure : % [sic] t_{n+1} = (\Delta_+\Delta_-)^{2n} -> t_{n+1} = (\Delta_+\Delta_-)^{2n+2}, exponent inconsistent with t_n = (\Delta_+\Delta_-)^{2n}

\dgrblock{dgrb412}{Recovering the path-space theorem without properness}
\begin{dcomment}{development}{p166}{Recovering the path-space theorem without properness}
The smoothness argument supplies enough control once we also know that transport along an arrow of $X\times X$ is invertible up to homotopy. Full path categories are particularly convenient for this second step, because their concatenation is a genuine functor.

\textbf{Transport and homotopy fibres for all path types.}\ours\quad
Let $E=\Ch(X)$ with endpoint functor $p:E\to B=X\times X$. For every $(a,b)$, the natural square with upper left corner $\Ch(X;a,b)$ and lower left corner $\{(a,b)\}$ is homotopy Cartesian. In particular,
\[
 B\Ch(X;a,b)\simeq
 \operatorname{hofib}_{(a,b)}(BX\xrightarrow{\Delta}BX\times BX),
\]
where $B(-)$ denotes the classifying space, so the right side is the ordinary path space from $a$ to $b$ up to weak equivalence.

\emph{Proof.} For $h=(u,v):(a',b')\to(a,b)$ define
\[
 S_h:\Ch(X;a,b)\longrightarrow\Ch(X;a',b'),
 \qquad c\longmapsto\mathrm{ch}(u)c\mathrm{ch}(v)^o.
\]
The opposite paths define a candidate inverse. Its composites insert the two backtracking paths associated to $u$ and $v$. For an arrow $u:x\to y$, cancellation is provided by the explicit diagrams
\[
 \text{empty}_x\ \longleftarrow\ x_{\Delta_+\Delta_-}
 \ \longrightarrow\ (x\xrightarrow{u}y\xleftarrow{u}x),
\]
\[
 (y\xleftarrow{u}x\xrightarrow{u}y)
 \ \longrightarrow\ y_{\Delta_-\Delta_+}
 \ \longrightarrow\ \text{empty}_y.
\]
The first rightward arrow has components $(1_x,u,1_x)$; the second display's first arrow has components $(1_y,u,1_y)$. The remaining arrows collapse constant intervals. They are morphisms with fixed endpoints. Whiskering with an arbitrary path makes these diagrams natural in that path. Cancelling each of the two inserted backtracks consequently gives homotopies from both composites to the identity. Hence $S_h$ is a homotopy equivalence.

Smoothness identifies the inclusion $i_z:E_z\to z\backslash E$ as a weak equivalence for every $z\in B$. The arrows $\rho_c:S_h(c)\to c$ over $h$ give a natural transformation
\[
 i_{(a',b')}S_h\ \Longrightarrow\ h^*i_{(a,b)}
\]
in $(a',b')\backslash E$. The three known weak equivalences in this square force $h^*:(a,b)\backslash E\to(a',b')\backslash E$ to be a weak equivalence. Quillen's Theorem~B now identifies $E_z$ with the homotopy fibre of $p$. Finally the constant-path section $X\to E$ is a weak equivalence over $X\times X$, so it identifies that homotopy fibre with the homotopy fibre of the diagonal. The latter is equivalent to the space of paths from $a$ to $b$: in the standard homotopy-fibre model, join the two paths from the chosen intermediate point, reversing the first.\hfill$\square$

The use of Quillen~B is classical \citep[\S1]{Quillen1973}; the explicit application above repairs the proof in these pages. We have checked the categorical cocylinder treatments of \citet{Evrard1975} and \citet{Shoikhet2016}, but did not locate this precise correction of Grothendieck's total path category there. No claim of priority is attached to the argument.
\end{dcomment}


\msp{page 168}
% foot: Version 22 décembre 2001

\printedpage{105}
% head: A. Grothendieck : Dérivateurs, Ch. VII, Édition des Universités de Montpellier II et Paris VII
\[
\tilde{w}'\Big\downarrow
\left\{
\overbrace{\underbrace{
\begin{tikzcd}[ampersand replacement=\&]
a' \arrow[r, "\mathrm{id}"] \arrow[d, "\mathrm{id}"'] \& a' \arrow[d, "u"'] \& a' \arrow[l, "\mathrm{id}"'] \arrow[d, "u"'] \arrow[rrrr, dotted, no head, "c'", "\longrightarrow"'] \& {} \arrow[d, "w'"] \& {} \arrow[d, "w'"] \& {} \arrow[d, "w'"] \& b' \arrow[r, "\mathrm{id}"] \arrow[d, "v"'] \& b' \arrow[d, "v"] \& b' \arrow[l, "\mathrm{id}"'] \arrow[d, "\mathrm{id}"] \\
a' \arrow[r, "u"] \& a \& a \arrow[l, "\mathrm{id}"'] \arrow[rrrr, no head, "c", "\longrightarrow"'] \& {} \& {} \& {} \& b \arrow[r, "\mathrm{id}"] \& b \& b'. \arrow[l, "v"']
\end{tikzcd}
}_{\tilde{c}}}^{\tilde{c}'}
\right\}
\tilde{w}'
\]

Ainsi
\[
\left\{
\begin{array}{lcl ll}
\tilde{c} & = & \xi_u \circ c \circ \eta_v & \text{où} &
\left\{
\begin{array}{lccccc}
\xi_u : & a' & \xrightarrow{u} & a & \xleftarrow{\mathrm{id}} & a \\[1ex]
\eta_v : & b & \xrightarrow{\mathrm{id}} & b & \xleftarrow{v} & b'
\end{array}
\right\}
\text{ type } \Delta_+\Delta_- \\[4ex]
\tilde{c}' & = & a'_{\Delta_+\Delta_-} \circ c' \circ b'_{\Delta_+\Delta_-} & & \\[4ex]
\tilde{w}' & = & \lambda(u) * w * \mu(v) & \text{avec} & % [sic] w -> w', the homomorphism of the pair (c',w')
\left\{
\begin{array}{lccc}
\lambda(u) : & a'_{\Delta_+\Delta_-} & \longrightarrow & \xi_u \\[1ex]
\lambda(v) : & b'_{\Delta_+\Delta_-} & \longrightarrow & \eta_v % [sic] \lambda(v) -> \mu(v), the formula above has \mu(v)
\end{array}
\right\}
\text{ comme dans \texttt{[la]} figure.}
\end{array}
\right.
\]
On trouve ainsi, mutatis mutandis,
\[ \CN(c;u,v) \xrightarrow{\varphi} \CN(a',b')/\tilde{c}, \qquad \varphi(c',w') = (\tilde{c}',\tilde{w}'), \]
et on définit le foncteur
\[ \CN(c;u,v) \xleftarrow{\psi} \CN(a',b')/\tilde{c} \]
ainsi : à $(c'',w'')$ dans $\CN(a',b')/\tilde{c}$, donc un chemin $c''$ de type $t_n = (\Delta_+\Delta_-)^{2n}$ de $a'$ à $b'$ et un homomorphisme $w'' : c'' \longrightarrow \tilde{c}$ à extrémités fixes, on associe $(c',w')$ dans $\CN(c;u,v)$, avec
\[ c' = c'', \qquad w' = \rho_c w'', \]
où
\[ \rho_c : \tilde{c} \longrightarrow c \]
est défini de façon similaire que tantôt, par
\msp{page 169}
\[
\begin{tikzcd}
a' \arrow[r, "u"] \arrow[d, "u"'] & a \arrow[d, "\mathrm{id}"] & a \arrow[l, "\mathrm{id}"'] \arrow[d, "\mathrm{id}"] \arrow[rr, zigzag] & {} \arrow[d, "\mathrm{id}"'] & b \arrow[r, "\mathrm{id}"] \arrow[d, "\mathrm{id}"] & b \arrow[d, "\mathrm{id}"] & b' \arrow[l, "v"'] \arrow[d, "v"] & \tilde{c},\ \text{ type } t_{n+1}. \\
a \arrow[r, "\mathrm{id}"] & a & a \arrow[l, "\mathrm{id}"'] \arrow[rr, zigzag, "c"] & {} & b \arrow[r, "\mathrm{id}"] & b & b' \arrow[l, "\mathrm{id}"'] & \varphi^*_{n,n+1}(c).
\end{tikzcd}
\]

On doit seulement vérifier que l'endofoncteur $\psi\varphi$ de $\CN(c;u,v)$ est homotope à l'identité. Mais ce foncteur transforme $(c',w')$ en $(\tilde{c}' = a'_{\Delta_+\Delta_-}c'b'_{\Delta_+\Delta_-}, \overline{w'})$, où $\overline{w'}$ est encore le composé
% foot: Version 22 décembre 2001
\printedpage{106}
% head: A. Grothendieck : Dérivateurs, Ch. VII, Édition des Universités de Montpellier II et Paris VII
\[
\begin{tikzcd}
\tilde{c}' \arrow[r, "\text{can.}"] \arrow[rr, bend right=20, "\overline{w'}"'] & c' \arrow[r, "w'"] & c.
\end{tikzcd}
\]

On a donc
\[ \psi\varphi \longrightarrow \mathrm{id}, \]
$\psi\varphi$ est homotope au foncteur identique.

Cela prouve la lissité de $\CN \longrightarrow X \times X$. Pour la propreté, il faut prouver la lissité de le cas de $X^o$, et du système projectif $\underline{t}^{\mathrm{opp}} \simeq (t_n^{\mathrm{opp}})_n$. Mais on a $\underline{t}^{\mathrm{opp}} = \underline{t}$, et il suffit d'appliquer ce qu'on a déjà prouvé au cas de $X^o$, pour voir que $\CN \longrightarrow X \times X$ est propre. % [sic] de le cas -> dans le cas, grammatical

Il ne reste plus que le lemme à prouver. Ça doit être essentiellement
\msp{page 170}
trivial, et je renonce à rédiger une démonstration.

\dgrblock{dgrb413}{The sequential intermediary and the repaired comparison}
\begin{dcomment}{development}{p169}{The sequential intermediary and the repaired comparison}
Grothendieck's intermediary remembers a finite stage as part of an object. That memory explains both its value for comparison and the limitation of the padding argument just attempted. Put $D=\Delta_+\Delta_-$, $t_n=D^{2n}$, and denote the specified symmetric-collapse transition by $\delta_{m,n}:t_m\to t_n$. An arrow of $\mathcal C_{\mathbf N^o}$ over $m\to n$ has exactly this underlying type map. Padding both ends of an arbitrary such arrow produces
\[
 1_D*\delta_{m,n}*1_D:t_{m+1}\longrightarrow t_{n+1},
\]
which is generally different from $\delta_{m+1,n+1}$. The former retains the outermost blocks and collapses blocks immediately inside them; the latter collapses the outermost blocks. Thus the formula on p.~168 need not even define the promised object of the slice of $\mathcal C_{\mathbf N^o}$. This is a restriction of the chosen indexing arrows, independent of the corrected exponents and primes.

\textbf{Smoothness of the sequential intermediary.}\ours\quad
Nevertheless $\mathcal C_{\mathbf N^o}\to X\times X$ is smooth.

\emph{Proof.} Fix a target path $c$ of stage $n$ and an incoming base arrow $h:(a',b')\to(a,b)$. Its lift category is fibred over the tail $\{m\geq n\}^o$. The fibre at $m$ is the incoming-lift category of $h$ into $\delta_{m,n}^*c$ for the fixed-type endpoint functor
\[
 \operatorname{Fun}(t_m,X)\longrightarrow X\times X.
\]
For $m>0$ this functor is a Grothendieck fibration: both endpoints of $t_m$ are sources, and a Cartesian lift is obtained by changing their objects and composing the first and last arrows with the specified endpoint maps. The corresponding lift category has that Cartesian lift as a terminal object. The only exception is $m=0$, possible when $n=0$. Its fibre is empty unless $a'=b'$ and the two components of $h$ coincide; when these equalities hold it is a singleton. If it is empty, omit stage $0$; otherwise retain it. In either case every fibre over the remaining tail is contractible, and that tail has a terminal object. The fibrewise comparison theorem for Grothendieck fibrations, applied to the projection to this tail, shows that the total lift category is weakly contractible. This is the smoothness criterion.\hfill$\square$

This argument proves only smoothness. Properness is excluded by the empty-path example above. Likewise the word ``asph\'erique'' at the end of the proof of $(1')$ on p.~166 is stronger than what was proved: for a positive-stage object $c$, the comma category of $X\to\mathcal C_{\mathbf N^o}$ over $c$ is empty, since there is no base arrow $0\to n$ for $n>0$. The section is a weak equivalence, as required.

\textbf{Comparison on every fixed-endpoint fibre.}\ours\quad
The canonical span
\[
 \Ch(X;a,b)\ \longleftarrow\
 \mathcal C_{\mathbf N^o}(a,b)\ \longrightarrow\
 \Ch_\infty(X;a,b)
\]
consists of weak equivalences.

\emph{Proof.} Write $E_1=\Ch(X)$, $E_0=\mathcal C_{\mathbf N^o}$, $E_2=\Ch_\infty(X)$, all over $B=X\times X$, and $\pi_i:E_0\to E_i$. The constant-path lemma proves both $\pi_i$ are weak equivalences on total categories. The filtered Grothendieck-to-colimit comparison, established in the commentary at manuscript pages~108--109, proves $\pi_2:E_{0,z}\to E_{2,z}$ is a weak equivalence on each strict fibre: apply that comparison to the sequence of fixed-endpoint categories itself.

By the endpoint theorem for $E_2$, its strict fibre is its homotopy fibre over $z$. The natural diagram
\[
\begin{tikzcd}[column sep=large]
 E_{0,z}\arrow[r,"\pi_2"]\arrow[d] & E_{2,z}\arrow[d,"\sim"]\\
 \operatorname{hofib}_z(BE_0\to BB)\arrow[r,"\sim"] &
 \operatorname{hofib}_z(BE_2\to BB)
\end{tikzcd}
\]
is understood after applying $B$ to the top row. Its lower arrow is a weak equivalence because the maps on total spaces are weak equivalences over the same base. The left vertical map is therefore a weak equivalence. The preceding block proves the same homotopy-fibre description for $E_1$. Repeating the square with $\pi_1$ shows that $E_{0,z}\to E_{1,z}$ is a weak equivalence.\hfill$\square$

The conclusion of the section is thus recovered, although its claimed properness and its particular sequential padding construction cannot be retained. The repair uses based homotopy fibres, never a conservativity assertion about unbased finite tests. The literature check is the one stated in the preceding development block; these are proofs supplied here for the specified models.
\end{dcomment}


\dgrblock{dgrb414}{Even translation acts trivially after localization}
\begin{dcomment}{development}{p169}{Even translation acts trivially after localization}
The comparison now also settles the question about $T^2$ left on pp.~149--150. Its resolution uses a visible change of finite support, rather than equality after restriction to finite test objects.

\textbf{Translation invariance.}\ours\quad
The even translation of an eventually constant alternating path induces the identity in $\Hot$, both on $\Ch_\infty(X)$ and on each fixed-endpoint category.

\emph{Proof.} Keep $D=\Delta_+\Delta_-$ and $t_n=D^{2n}$. Define $R$ on the sequential Grothendieck category by sending a path of stage $n$ to stage $n+1$, adding two constant $D$-blocks at its right end. The asymmetric collapse $q_n:t_{n+1}\to t_n$ defining this operation commutes with the symmetric transitions:
\[
 q_n\delta_{m+1,n+1}=\delta_{m,n}q_m\qquad(m\geq n).
\]
Both sides collapse the same number of left and right blocks. Hence $R$ is a functor. On the eventual bi-infinite realization, with the centred indexing of these pages this operation is $U=T^{-2}$, and $\pi_2R=U\pi_2$.

The maps $q_n$ are permitted in $\Ty$. They give a natural transformation $\pi_1R\Rightarrow\pi_1$ into the full path category, whose component is the Cartesian arrow from the padded path to the original path. By the comparison just proved, $k=\pi_1\pi_2^{-1}$ is an isomorphism in $\Hot$. Therefore
\[
 kU\pi_2=\pi_1R=\pi_1=k\pi_2,
\]
and cancellation gives $U=1$. The same construction fixes the endpoints, so the argument applies to every fibre. Taking inverses gives $T^2=1$ as well.\hfill$\square$

We did not find this finite-support proof in the checked literature. It is a consequence of the repaired comparison, with no additional general theorem about the restricted Yoneda functor.
\end{dcomment}


\dgrblock{dgrb415}{Why the comparison preserves composition}
\begin{dcomment}{development}{p169}{Why the comparison preserves composition}
Grothendieck's original reason for retaining all finite path types was their strictly associative concatenation. The comparison of endpoint fibres should therefore identify that operation with ordinary path concatenation, not merely identify the underlying homotopy types. The constant-path sections determine this compatibility once composable pairs are treated as carefully as individual paths.

\textbf{Compatibility with path concatenation.}\ours\quad
Put $E=\Ch(X)$ and $K=BX$. The equivalences
\[
 u_{a,b}:B\Ch(X;a,b)\longrightarrow P_{a,b}K
 \quad\hbox{in }\Hot,
\]
obtained from the constant-path diagonal factorization, can be chosen canonically in the localized slice category and satisfy
\[
 u_{a,c}\,B(\mathord{*})
   =\mathord{*}_{\mathrm{top}}(u_{a,b}\times u_{b,c})
 \quad\hbox{in }\Hot.
\]
Here $P_{a,b}K$ is the ordinary space of paths from $a$ to $b$, and the products have the indicated endpoint fibres as their domain. In particular the equivalence on loops respects the usual multiplication in the homotopy category.

\emph{Proof.} Let $D=E\times_{X}E$, using the target of the first path and the source of the second. Its projection to $\Ty\times\Ty$ is the fibred Grothendieck construction with fibre
$\operatorname{Fun}(\tau*\upsilon,X)$ at $(\tau,\upsilon)$: matching the two endpoint values is exactly the condition for gluing the diagrams on the two intervals. The constant-diagram functor from $X$ to this fibre is a categorical homotopy equivalence, by the truncation homotopies proved at manuscript page~160. These functors commute with the type maps. The same fibrewise comparison used there, and the terminal object $(e,e)$ of $\Ty\times\Ty$, show that the constant-pair section $j:X\to D$ is a weak equivalence.

We also need the relative assertion for $D$. Its three-endpoint functor $r:D\to X^3$ factors as
\[
 D\longrightarrow E\times X\longrightarrow X^3,
 \qquad
 (c,d)\longmapsto(c,\beta d)
       \longmapsto(\sigma c,\beta c,\beta d).
\]
The first functor is the base change of $E\to X^2$ along $(c,z)\mapsto(\beta c,z)$; the second is its product with $1_X$. Thus $r$ is smooth, by stability under base change and composition \citep[Proposition~4.7]{Maltsiniotis2005Asphericite}. For a base arrow $(u,v,w):(a',b',c')\to(a,b,c)$ the functor on strict fibres is
\[
 S_{(u,v)}\times S_{(v,w)}.
\]
Both factors are homotopy equivalences by the explicit backtrack cancellations at page~166. The two collapse arrows $\rho$ agree on their common middle endpoint, with component $v$, so together give a natural arrow over $(u,v,w)$. Smoothness identifies each strict fibre with its under-comma category. The natural arrow then shows that every transition between those comma categories is a weak equivalence, exactly as in the proof for $E$. Quillen's Theorem~B applies. Consequently
\[
 B\bigl(\Ch(X;a,b)\times\Ch(X;b,c)\bigr)
     \longrightarrow\operatorname{hofib}_{(a,b,c)}(BD\to K^3)
\]
is a weak equivalence.

Now use ordinary spaces, with compactly generated mapping spaces. Put $P=K^{[0,1]}$ and $Q=P\times_K P$, and let $s:K\to P$ and $t:K\to Q$ be their constant-path sections. Both are weak equivalences: $P$ and $Q$ are the spaces of maps from the intervals $[0,1]$ and $[0,2]$, respectively, and evaluation at the initial point is a deformation retraction onto constants. Their endpoint maps, to $K^2$ and $K^3$, are fibrations. Write $i:K\to BE$ for the constant-path section and still write $j:K\to BD$ for its two-path version. In the homotopy categories of spaces over these endpoint bases, define
\[
 u=s\,i^{-1}:BE\longrightarrow P,
 \qquad
 v=t\,j^{-1}:BD\longrightarrow Q.
\]
These are the unique comparison isomorphisms whose composites with the constant sections are $s$ and $t$.

All maps in the remaining argument can be regarded over $K^3$. For concatenation the target is the pullback of $P\to K^2$ along $(x,y,z)\mapsto(x,z)$; for either projection it is the corresponding pullback along $(x,y,z)\mapsto(x,y)$ or $(y,z)$. These pullbacks are just products with the unused copy of $K$, so they preserve weak equivalences and the displayed localized comparisons. Both composites comparing categorical concatenation with topological concatenation send $j$ to the same constant path. Since $j$ is invertible in this slice homotopy category, cancellation gives
\[
                u\,B(\mathord{*})=\mathord{*}_{\mathrm{top}}\,v.
\]
The identical argument with either path projection shows that $v$ commutes with both projections under $u$.

Take the derived fibre over $(a,b,c)$. The relative assertion for $D$, the endpoint theorem for $E$, and the ordinary path fibrations identify these derived fibres with the stated strict path spaces and their products. The two projection identities imply that the induced map of the product fibres is $u_{a,b}\times u_{b,c}$, since a map to a product in $\Hot$ is determined by its two projections. The concatenation identity is therefore the required formula. Constants give the corresponding unit identity.\hfill$\square$

This argument supplies the compatibility omitted by a bare comparison of homotopy types. It uses Quillen's relative fibre theorem and localization at weak equivalences, not a conservativity or faithfulness assertion about finite tests. We checked the path constructions in \citet[Section~I.1]{Evrard1975} and the diagonal and cocylinder comparisons in \citet[Sections~1.5 and~2.1, Theorem~2.8]{Shoikhet2016}; we did not find this particular comparison through the category of composable pairs there. The ingredients are classical, and no priority is claimed.
\end{dcomment}


Je vais maintenant montrer `formellement' que la factorisation `canonique' d'une flèche de Cat $f : X' \longrightarrow X$ en
\[ X' \xrightarrow{i_f} \overline{X'} \xrightarrow{\bar{f}} X, \qquad i_f \text{ Hot-équivalence,} \]
avec (on l'espérait) $\bar{f}$ propre et lisse, à partir d'une factorisation similaire pour la flèche diagonale $X \xrightarrow{\ \ \mathrm{diag}_X\ \ } X \times X$ en
\[
\begin{tikzcd}
X \arrow[d, "i_X\ \ \text{\footnotesize Hot-équivalence}"] \\
\mathcal{C}_X\ \ \text{`chemins'} \arrow[d, "p_X"] \\
X \times X
\end{tikzcd}
\]
est chose formelle. (Cf. page 139 pour les notations et \texttt{[la]} problématique.) On regarde le diagramme
\[
\begin{tikzcd}
 & \mathcal{C}_X \arrow[d, "p_X"] \arrow[dl, "\beta_X"'] \arrow[dr, phantom, "\text{cart.}"] & C(f) \arrow[l, "\mathrm{pr}'_2"'] \arrow[d, "p_{X,f}"] \arrow[dll, bend right=50, "\bar{f}"'] & [-1.5em] \simeq X' \times_X (C_X, \sigma = \mathrm{pr}_1 \circ p_X) \\ % Reading resolved against the supplied PDF; see corrections.csv.
X & X \times X \arrow[l, "\mathrm{pr}_2^X"] \arrow[d, "\mathrm{pr}_1^X"'] \arrow[dr, phantom, "\text{cart.}"] & X' \times X \arrow[l, "f \times \mathrm{id}_X = f'"] \arrow[d, "\mathrm{pr}_1"] & \\
 & X & X' \arrow[l, "f"'] \arrow[uu, bend right=60, dotted, "i_f"'] &
\end{tikzcd}
\]
% foot: Version 22 décembre 2001
\printedpage{107}
% head: A. Grothendieck : Dérivateurs, Ch. VII, Édition des Universités de Montpellier II et Paris VII
(\footnote[40]{NB $C(f)$ est défini comme l'image inverse sur $X'$ du fibré $C_X$ sur $X$ (par $\sigma_X = \mathrm{pr}_1^X \circ p_X$) par $f : X' \longrightarrow X$.}). Comme $p_X$ est propre et lisse, i.e. `parfait', son image inverse sur $X' \times X$, savoir $p_{X,f}$, l'est aussi; d'autre part $\mathrm{pr}_2^X \circ f'$ est la
\msp{page 171}
projection $\mathrm{pr}_2$ de $X' \times X$ sur $X$, donc c'est aussi un morphisme parfait (puisque $X' \longrightarrow e$ est parfait, donc il en est de même du composé
\[ (\mathrm{pr}_2^X \underbrace{f') \circ p_{X,f}} = \underbrace{\mathrm{pr}_2^X(p_X}\, \mathrm{pr}'_2) = \beta_X \mathrm{pr}'_2 \overset{\text{déf}}{=} \bar{f} \ \ ). \]

\dgrblock{dgrb416}{The formal diagonal factorization, with its hypothesis retained}
\begin{dcomment}{exposition}{p170}{The formal diagonal factorization, with its hypothesis retained}
The closing argument extracts the reusable part of the path construction. Suppose a category $C_X$ is given with a factorization
\[
 X\xrightarrow{i_X}C_X\xrightarrow{p_X}X\times X,
 \qquad p_Xi_X=\Delta_X,
\]
where $i_X$ is a weak equivalence and $p_X$ is smooth and proper. Put $\sigma_X=\operatorname{pr}_1p_X$ and $\beta_X=\operatorname{pr}_2p_X$. Both are smooth and proper, because these classes are closed under composition and base change; a projection is a base change of a map to the terminal category. The identity $\sigma_Xi_X=1_X$ implies that $\sigma_X$ is a weak equivalence. Its strict fibres are homotopy fibres by the endpoint argument, hence weakly contractible. The lift-category criterion and Quillen~A then show that every base change of $\sigma_X$ is a weak equivalence.

For $f:X'\to X$, form the ordinary pullback $C(f)=X'\times_{X,\sigma_X}C_X$. Its projection $r:C(f)\to X'$ is a weak equivalence. The section $i_f$ induced by $i_X$ satisfies $ri_f=1_{X'}$, so $i_f$ is a weak equivalence. The map to the other endpoint factors as
\[
 C(f)\longrightarrow X'\times X\longrightarrow X.
\]
The first arrow is the base change of $p_X$, and the second is a projection. Their composite $\bar f$ is therefore smooth and proper. The equation $\bar f i_f=f$ proves the requested factorization. This is exactly the formal argument of the manuscript, with the universal weak-equivalence step justified.

It applies to $C_X=\Ch_\infty(X)$, whose two lifting properties were proved using the two cofinal parity systems. It cannot be applied with the asserted hypothesis to $C_X=\Ch(X)$ or $\mathcal C_{\mathbf N^o}$: those endpoint maps fail properness. Their weak equivalence over $X\times X$ does not transfer a strict comma-category lifting condition. The formal result and the comparison theorem therefore coexist without identifying their hypotheses.
\end{dcomment}


Je rappelle que la section $i_f$ de $C(f)$ sur $X'$ est définie à partir de la section $i_X$ de $(\mathcal{C}_X,\sigma_X)$ sur $X$. C'est une Hot-équivalence, car $\sigma_X = \mathrm{pr}_1^X p_X$ est parfait à fibres asphériques, donc universellement une Hot-équivalence, donc $\mathrm{pr}_1 p_{X,f} : C(f) \longrightarrow X'$ est aussi une Hot-équivalence, et toute section en est donc une aussi.
% foot: Version 22 décembre 2001


\dgrblock{dgrb417}{A mapping-path fibre for a homomorphism of groups}
\begin{dcomment}{examples}{p171}{A mapping-path fibre for a homomorphism of groups}
The factorization records more than the kernel of a functor on components. For a homomorphism $\phi:H\to G$, write $\mathbf B H$ and $\mathbf B G$ for their one-object groupoids and consider the functor $\mathbf B\phi$ between them. The homotopy fibre has a particularly small categorical model, the action groupoid $F_\phi$ with objects $g\in G$ and arrows
\[
 h:g\longrightarrow g'
 \quad\Longleftrightarrow\quad
 g'\phi(h)=g.
\]
If $h:g\to g'$ and $k:g'\to g''$, their composite is $kh$, since $g''\phi(kh)=g$. All arrows are invertible. This is the over-comma category $(\mathbf B\phi)/*$, where $*$ is the unique object of $\mathbf B G$. Every arrow in the base groupoid acts on it by an isomorphism, so Quillen~B identifies $BF_\phi$ with the homotopy fibre.

Two objects lie in the same component precisely when they are in the same coset $g\operatorname{im}\phi$. An automorphism of any $g$ is an $h$ with $\phi(h)=1$, hence the isotropy group is $\ker\phi$. Choosing one representative in each coset gives an equivalence of groupoids
\[
 F_\phi\simeq\coprod_{G/\operatorname{im}\phi}\mathbf B(\ker\phi).
\]
Thus $\pi_0$ detects the failure of surjectivity and the fundamental group of each component detects the failure of injectivity. For multiplication by a nonzero integer $m$ on $\mathbf Z$, the fibre has $|m|$ contractible components. For the map $\mathbf Z\to1$, it is a circle up to homotopy. The corresponding fibre of Grothendieck's mapping-path factorization has exactly these homotopy types. In particular it is weakly contractible precisely when $\phi$ is an isomorphism, which illustrates the triviality criterion at p.~156 without computing any individual categorical path.
\end{dcomment}
\begin{otherlanguage}{english}
\dgroverview{VII-c}{Chapter VII \textperiodcentered\ \textup{VII-c} \textperiodcentered\ Overview}
\begingroup
\selectlanguage{english}
\noindent\textbf{Paths, comparison, and the information carried by a model.}

These pages bring the long construction of categorical paths to its decisive
test.  A path should do more than give a word in arrows and their inverses.  Its
category should have the homotopy type of a space of paths, should vary with
both endpoints, and should turn an arbitrary functor into a map whose strict
fibres compute homotopy fibres.  Grothendieck wants these properties before
making any fibrant replacement of the category in which the paths lie.  The
problem is therefore one of comparison between concrete models, with their
actual arrows and their actual endpoint maps
\citep[ms.~108--171]{GrothendieckDerivateursVII}.

There are three models in play.  The category $\Ch(X)$ retains a finite interval
type as part of a path, and permits changes of type.  Concatenation is strictly
associative there.  A chosen sequence of longer alternating intervals gives
categories $C_n$ and a filtered colimit $C_\infty$; identity arrows added at the
ends become identifications in this colimit.  Between them stands the
Grothendieck construction $C_{\mathbf N^o}$, which retains the stage and the
transition morphisms.  Passing the diagram $(C_n)$ through homotopy first gives
yet another object, the ind-object $H_T$.  None of these operations should be
silently substituted for another.  The opening pages, from manuscript page~108, are largely
an investigation of what is lost by doing so.

For a filtered diagram the comparison from its Grothendieck construction to
its ordinary colimit is a weak equivalence.  Nerves commute with filtered
colimits, and filtered homotopy colimits of simplicial sets map by weak
equivalences to ordinary colimits.  This explains one side of the comparison
span on manuscript page~108.  The other side depends on the indexing category
and on the diagram.  Grothendieck correctly observes that cofinality in the
preorder of objects does not make the original indexing category contractible:
a one-object group already distinguishes the two statements.  The appropriate
general machinery is the homotopy-colimit theorem for the Grothendieck
construction and the comma-category criteria for homotopy cofinality
\citep{Thomason1979,Quillen1973}.

The proposed independence of the chosen sequence is more delicate than the
text acknowledges.  An isomorphism of ind-objects in the ordinary homotopy
category does not control the compatibility of the homotopies used at
successive stages.  Our finite-injection example makes this precise within the
general hypotheses of manuscript pages~99--101.  To a finite set we associate
the one-object category of its permutation group.  Injections act by extending
permutations by the identity.  Parallel injections induce naturally isomorphic
functors, yet a morphism between two cofinal sequences induces, on the colimit,
the non-surjective shift of the finitary permutation group.  Thus the general
comparison assertion on pages~117--121 is false.  This does not invalidate the
geometric path comparison proved later by a different argument.

The same distinction clarifies the questions about
$\Hot\longrightarrow\Hot_0^{\wedge}$.  In this notation the subscript $0$
refers to a chosen category of finite tests; it is not a truncation to
homotopy $0$-types.  Restricting unbased representable functors to finite tests
is neither generally faithful nor generally conservative.  The locally inner
permutation-group example in the commentary on Chapter~V, manuscript page~52,
already treats conservativity.  Here a telescope of degree-two maps of circles supplies an
explicit failure of faithfulness: an essential map to $K(\mathbf Z,2)$ is
invisible on every finite stage.  Its obstruction is calculated in the
commentary, rather than attributed vaguely to a failure of coherence.  The
cellular and representability tools are classical
\citep[\S3F and Theorem~4.57]{Hatcher2002}.  In addition, a category with
finitely many arrows need not have a homotopy-finite nerve.  Compactness of a
test in the homotopy category requires its own hypothesis.  These qualifications
matter whenever a displayed bijection is used to infer an isomorphism or an
equality of maps.

\medskip\noindent\textbf{External Hom and the role of the parameter category.}

Section~11 returns to the reason that paths entered the theory of derivators.
For a relative category $(\mathcal M,\Sigma)$ one can always form the
prederivator
\[
 A\longmapsto\mathcal M^A[\Sigma_A^{-1}],
\]
where $\Sigma_A$ is defined objectwise.  Grothendieck asks when the morphisms
between constant diagrams are represented by a homotopy type, with the expected
dependence on the homotopy type of $A$.  This is a substantive condition on a
relative category.  It does not follow merely from the existence of an
ordinary localization, and the existence of a prederivator is not the
existence of all the Kan extensions required of a derivator.

The discussion of a categorical two-sphere on pages~132--135 identifies a real
problem but gives an incorrect explanation.  A fixed finite circle-shaped
parameter can already support arbitrarily high degrees when the length of the
target path is allowed to increase: concatenate the same family of loops
repeatedly.  We work this out for the six-object suspension of the four-object
circle poset.  The same example gives a genuine failure of homotopy invariance
by comparing that circle poset with the one-object category $B\mathbf Z$.
Every functor from a group to a poset is constant; the resulting localization
has lost the family that the circle-shaped parameter sees.  The example
separates the valid concern about parameter categories from the proposed
bound on degrees.

There is a second distinction at page~135.  The total arrow category
$\operatorname{Fun}([1],C)$ has the homotopy type of $C$, for every $C$.
Evaluation and the constant-diagram functor give the needed adjunction and
natural transformation.  What fails for an arbitrary category is the claim
that its strict fixed-endpoint arrow sets are derived path spaces.  Likewise,
Thomason fibrancy of the target alone does not turn every ordinary functor
category into a derived internal Hom.  The commentary supplies a fibrant-poset
counterexample and distinguishes a derived mapping-space construction from
this incorrect assertion.  The separate claim that ordinary functor categories
preserve Thomason fibrancy is also false.  An explicit equivariant horn-cell
construction gives a fibrant category $E$ for which
$\operatorname{Fun}(B(\mathbf Z/2),E)$ has a nonfibrant retract; the proof uses
the fixed-point pushout formulas of
\citet[Lemmas~4.7--4.8]{MayZakharevichStephan2017}.
Simplicial localization provides a systematic
external Hom for a relative category, while additional model-theoretic
hypotheses justify comparisons with categories of ordinary diagrams
\citep{DwyerKan1980Simplicial,DwyerKan1980Calculating,Groth2013,Cisinski2009LocallyConstant,Lurie2009HTT}.

\medskip\noindent\textbf{What the path comparison establishes.}

Sections~12 and~13 concern $\Sigma=\operatorname{Mor}(X)$.  Alternating paths
of unbounded length are the essential construction.  By choosing the parity
of their end arrows one obtains finite-stage Grothendieck fibrations or
opfibrations over the endpoint category.  The categories of lifts pass to
filtered unions.  This proves smoothness and properness of
$\Ch_\infty(X)\to X\times X$, although the union need not itself be a
Grothendieck fibration or opfibration.  The lift-category argument, not the
existence of compatible Cartesian arrows at all stages, is what survives the
passage to the union.  The terminology comes from base-change phenomena:
modern asphericity theory makes the categorical tests precise
\citep{Maltsiniotis2005Asphericite}.

For the category of all finite paths the source proves smoothness by a useful
explicit construction.  Given endpoint arrows and a path $c$, insert those
arrows at the two ends to form $\widetilde c$.  The category of lifts is
compared with the slice over $\widetilde c$, which has a terminal object.
Padding, retraction, and their naturality give the contraction.  This is a
sound argument once the notation in the diagrams is read correctly.  Its
claimed dual, however, does not prove properness.  An empty path admits no
outgoing change to a nonempty interval type.  Over $X=[1]$, changing only one
endpoint therefore gives an empty category of outgoing lifts.  The same
obstruction applies to $C_{\mathbf N^o}$.  The assertion that both these
endpoint projections are proper must remain false in the audit.

The desired comparison nevertheless survives.  For $\Ch(X)$, concatenating an
endpoint arrow, or its reversed path, defines transport between fibres.
The two composites are connected to the identity by the explicit
backtracking-cancellation zigzags.  Combining these equivalences with the
proved smoothness gives the hypotheses of the dual of Quillen's Theorem~B.
Hence strict endpoint fibres compute homotopy fibres.  The known total and
fibrewise equivalences
$C_{\mathbf N^o}\to\Ch_\infty(X)$ transfer the same property to the middle
model.  Finally the total equivalence
$C_{\mathbf N^o}\to\Ch(X)$ becomes a fibrewise equivalence by comparison of
homotopy fibres.  This repairs the concluding argument without assigning a
false properness property to either total category.

It follows that the two geometric path models have equivalent endpoint
fibres.  Strict concatenation in the category of finite paths then furnishes
the expected composition in $\Hot$, including the loop-space multiplication.
The diagonal factorization and the mapping-path construction on pages~170--171
work with the perfect model $\Ch_\infty(X)$; their formal proof is expanded in
the commentary.  Evrard's earlier categorical cocylinder and Shoikhet's
refinement give closely related factorization theorems, with precisely stated
homotopy-fibration conclusions.  They should not be cited as asserting that
Grothendieck's particular maps are Thomason fibrations or that the inclusion
of constant paths is a Thomason cofibration \citep{Evrard1975,Shoikhet2016}.

The larger programme is continuous with \emph{Pursuing Stacks}, \S69: homotopy
theory is to retain the diagram categories needed to form and compare fibres,
cofibres, and their operations.  The letter to Thomason of 2~April~1991
explicitly returns to unbounded categorical paths and to constructions that
avoid a preliminary Kan replacement
\citep{Grothendieck1983,Grothendieck1991Letter}.  Here the useful boundary is
precise.  The geometric comparison and mapping-path factorization can be
proved; the general finite-test and ordinary-internal-Hom assertions require
correction.  The developments marked \ours\ give the explicit counterexamples
and repairs used in the commentary.  They record arguments given here,
with the literature checks stated in their blocks, and make no claim of
priority.



\endgroup
\end{otherlanguage}
\dgrchunkend

% DER_VIII-X: Grothendieck's text with the commentary set in, then the overview, inside a group that applies the chunk's own definitions
\dgrchapter{VIII}{Chapter VIII}{\texorpdfstring{1-types d'homotopie relatifs : leur intégration\,\dots}{1-types d'homotopie relatifs : leur intégration}}
\dgrpart{0}{VIII-X}{Manuscript pages 1--9}{VIII-X\enspace\textperiodcentered\enspace Chapter VIII\enspace\textperiodcentered\enspace manuscript pages 1--9}{Chapter VIII \textperiodcentered\ \textup{VIII-X}}
\dgrchunkbegin
% DER_VIII-X: the definitions of the chunk's preamble, applied inside its group
\let\printedpage\relax\newcommand{\printedpage}[1]{\marginpar{\footnotesize p.~#1}}
\let\msp\relax\newcommand{\msp}[1]{\marginpar{\footnotesize [#1]}}
\let\Grp\relax\newcommand{\Grp}{\underline{\mathrm{Grp}}}
\let\uHom\relax\newcommand{\uHom}{\underline{\mathrm{Hom}}}
\let\uGrp\relax\newcommand{\uGrp}{\underline{\mathrm{Grp}}}
\let\uuGrp\relax\newcommand{\uuGrp}{\underline{\underline{\mathrm{Grp}}}}
\let\fl\relax\newcommand{\fl}{\longrightarrow}
\let\Ch\relax\newcommand{\Ch}{\mathrm{Ch}}
\let\Hot\relax\newcommand{\Hot}{\mathrm{Hot}}
\let\Hom\relax\newcommand{\Hom}{\mathrm{Hom}}
\let\HotO\relax\newcommand{\HotO}{\mathrm{Hot}_0^{\wedge}}
\let\MOP\relax\newcommand{\MOP}{\mathcal{M}_{0\,\Pi}}
\let\M\relax\newcommand{\M}{\mathcal{M}}
\let\uCh\relax\newcommand{\uCh}{\underline{\mathrm{Ch}}}
\let\Homext\relax\newcommand{\Homext}{\mathrm{Hom\,ext}}
\let\Lotimes\relax\newcommand{\Lotimes}{\overset{\mathbf{L}}{\otimes}}
\let\cercle\relax\newcommand{\cercle}[1]{\tikz[baseline=(c.base)]\node[draw,circle,inner sep=1pt](c){#1};}
\let\Chu\relax\newcommand{\Chu}{\underline{\mathrm{Ch}}}
\let\D\relax\newcommand{\D}{\mathbf{D}}
\let\Cat\relax\newcommand{\Cat}{\mathrm{Cat}}
\let\Fl\relax\newcommand{\Fl}{\operatorname{Fl}}
\let\Ob\relax\newcommand{\Ob}{\operatorname{Ob}}
\let\bDelta\relax\newcommand{\bDelta}{\boldsymbol{\Delta}}

% ---- from Der-ChVIII-567_p001-005.tex
% ===================== page 1 =====================
\printedpage{}
% head: (aucun ; page de titre non numérotée)

\begin{center}
{\Large\textsf{LES DÉRIVATEURS}}

\bigskip

\textbf{ALEXANDRE GROTHENDIECK}

\bigskip\bigskip

\textbf{Chapitre VIII}

\medskip

\textbf{1-types d'homotopie relatifs : leur intégration \dots}
\end{center}

\clearpage

% ===================== page 2 =====================
\printedpage{}
% head: (aucun ; page non numérotée)

Ce texte a été déchiffré et transcrit en \LaTeXe{} par M. Künzer. Il a été édité par M. Künzer, J. Malgoire, et G. Maltsiniotis. La transcription est aussi fidèle que possible au manuscrit. Pour les quelques corrections évidentes, ou rares commentaires des éditeurs, ainsi que pour la numérotation originale des pages du manuscrit, les caractères de machine à écrire \texttt{[typewriter]} entre crochets sont utilisés. Un point d'interrogation entre crochets signifie que l'on n'est pas sûr du mot qui précède.

Cette édition est provisoire. Les remarques, commentaires et corrections sont bienvenus. Envoyer un message à :

\textbf{maltsin@math.jussieu.fr}

\hfill G. Maltsiniotis\hspace*{2cm}

\clearpage

% ===================== page 3 =====================
\printedpage{1}
% head: A. Grothendieck : Dérivateurs, Ch. VIII, Édition des Universités de Montpellier II et Paris VII
% foot: Version 2 février 2002

\msp{page 1}

\section*{1\quad Intégration des $1$-types d'homotopie relatifs sur une catégorie.}

Les types d'homotopie relatifs sur une (petite) catégorie $S$ sont définis par les \textit{champs} (en catégories) sur $S^{\wedge}$, qui correspondent essentiellement aux contrafoncteurs $S^o \longrightarrow (\mathrm{Cat})$, ou encore aux catégories $X$ sur $S$ qui sont Cat-fibrées sur $S$. (Dans $(\mathrm{Cat}/S)_{\mathrm{fib}}$, on inverse les flèches qui sont des Hot-équivalences sur les fibres.) D'autre part, les 1-types d'homotopie relatifs sur $S$ doivent se décrire par des \textit{champs en groupoïdes} sur $S$. Si on a un type d'homotopie relatif, défini par
\[ X \longrightarrow S, \qquad\qquad \text{catégorie Cat-fibrée sur } S, \]
ou par
\[
\begin{aligned}
S^o \quad &\longrightarrow \quad \mathrm{Cat}\\
s \quad &\longmapsto \quad X_s\ ,
\end{aligned}
\]
le 1-type d'homotopie associé est défini par
\[ s \longmapsto \Pi(X_s), \qquad\qquad S^o \longrightarrow (\mathrm{Grp}), \qquad\qquad \text{catégorie des groupoïdes,} \]
où $\Pi(X_s)$ est le groupoïde fondamental de $X_s$. On note la catégorie fibrée correspondante sur $S$ par \fbox{$\Pi(X/S)$}. Il est vrai que si
\[ X \ \longrightarrow\ X' \]
est un homomorphisme de catégories fibrées sur $S$ qui induise une Hot-équivalence sur les fibres, elle induit une équivalence fibrée
\[ \Pi(X/S) \ \longrightarrow\ \Pi(X'/S), \]
mais non un isomorphisme. Il y aurait donc du travail à faire pour définir un foncteur
\[
\underbrace{\mathrm{Hot}(S)}_{\text{\footnotesize\begin{tabular}{c}catégorie des\\ types d'homotopie\\ relatifs sur $S$\end{tabular}}}
\quad\xrightarrow{\ \xi\,\longmapsto\,\Pi(\xi/S)\ }\quad
\underbrace{\Grp(S)}_{\text{\footnotesize\begin{tabular}{c}catégorie des\\ groupoïdes\\ fibrés sur $S$\\ (ou catégories\\ des champs en\\ groupoïdes\\ sur $S^{\wedge}$)\end{tabular}}}
\quad .
\]

Mais ce n'est pas pour l'instant mon propos. Plutôt, de donner une formule de transitivité en homologie non commutative, analogue à $R(gf)_* \simeq Rg_*Rf_*$ en cohomologie.

\dgrblock{dgrb418}{Why the target is a 2-category}
\begin{dcomment}{exposition}{p1:Mais ce n'est pas pour l'instant mon propos.}{Why the target is a 2-category}
Grothendieck wants integration to retain the automorphisms of a homotopy type, and not merely its connected components. A groupoid is exactly the amount of structure needed for a homotopy $1$-type. Its objects represent points, its arrows represent homotopy classes of paths, and the nerve has no homotopy groups above degree one. For a category $C$, the groupoid $\Pi C$ is its localization at \emph{all} arrows. In particular, it is generally very different from the maximal subgroupoid of $C$.

The relative version must remember transport as well. We write $\mathsf{Gpd}(S)$ for the $2$-category of categories fibred in groupoids over $S$, functors over $S$, and natural transformations over $S$. Every arrow of a category fibred in groupoids is cartesian; consequently every functor over the base between such categories preserves cartesian arrows, and every $2$-arrow is invertible. A choice of cleavage gives a pseudofunctor $S^{\mathrm{op}}\to\mathsf{Gpd}$, including the isomorphisms comparing successive inverse images. This is the language of \citet[Expos\'e VI, \S\S6--8]{SGA1}. No topology imposing additional descent has been specified here: ``champs'' is being used through the equivalence with fibred categories on the presheaf topos.

Taking fundamental groupoids fibre by fibre is compatible with these comparison isomorphisms. It therefore gives a relative truncation, rather than a mere collection of unrelated groupoids. A weak equivalence of fibres becomes an equivalence of groupoids, which explains Grothendieck's hesitation about a functor to an ordinary category of actual groupoids: localization requires its inverted arrows to become \emph{isomorphisms}. One can either pass to the homotopy category of $\mathsf{Gpd}(S)$, whose morphisms are isomorphism classes of cartesian functors, or retain the entire $2$-category. In modern language the underlying operation is the pointwise truncation
\[
 \operatorname{Fun}(NS^{\mathrm{op}},\mathcal S)
 \longrightarrow
 \operatorname{Fun}(NS^{\mathrm{op}},\mathcal S_{\leq1}).
\]
Here $\mathcal S$ denotes spaces and $\mathcal S_{\leq1}$ their full subcategory of $1$-types. The composition law for integration will be an equivalence with coherent comparison isomorphisms, precisely the form in which it can be useful again under a further change of base.

The terminology continues the programme of \citet[\S69, especially printed pp.~3--5]{Grothendieck1983}: integration is the left adjoint to change of the indexing category. The present calculation isolates a particularly concrete part of that programme, where ordinary categories and their groupoid completions suffice to construct the adjoint.
\end{dcomment}


\clearpage

% ===================== page 4 =====================
\printedpage{2}
% head: A. Grothendieck : Dérivateurs, Ch. VIII, Édition des Universités de Montpellier II et Paris VII
% foot: Version 2 février 2002

\msp{page 2}

Si
\[ f : S \longrightarrow T \]
est un foncteur, on va définir un foncteur
\[ f_! :\ \Grp(S) \ \longrightarrow\ \Grp(T) \]
(image directe de 1-types d'homotopie relatifs). Mais notons d'abord que si
\[ p :\ X \ \longrightarrow\ S \]
est une $S$-catégorie \textit{quelconque} (pas nécessairement Cat-fibrée), on lui associe néanmoins un type d'homotopie relatif, via la catégorie fibrée $\tilde{X}^{(S)}$, ou $\tilde{X}$, ayant comme fibres les $s\backslash X$ (en notant que $s' \longrightarrow s$ dans $S$ définit bien $s\backslash X \longrightarrow s'\backslash X$) (\footnote{Ayant $\tilde{X}$, on définit $\Pi(X/S) \overset{\text{déf}}{=} \Pi(\tilde{X}/S)$.}). On a un $S$-foncteur canonique
\[ X \xrightarrow{\ i\ } \tilde{X}, \qquad\qquad x \longmapsto (x, p(x), \underbrace{p(x)}_{=\,\text{‘}\,s\,\text{’}} \xrightarrow{\ \mathrm{id}\ } p(x)). \]
C'est une $S$-Hot-équivalence si $X$ était déjà fibrée. En fait, sur chaque fibre c'est un homotopisme car on définit alors $\tilde{X}_{s_0} \xrightarrow{\ j\ } X_{s_0}$ de façon évidente, de façon que
\[
X_{s_0} \underset{j}{\overset{i}{\rightleftarrows}} \tilde{X}_{s_0}, \qquad ji \simeq \mathrm{id}, \qquad ij \underbrace{\xrightarrow{\text{homotopie simple}}}_{\text{\footnotesize\begin{tabular}{l}NB cette flèche\\ n'est pas \texttt{[un]}\\ isomorphisme\\ en général.\end{tabular}}} \mathrm{id}.
\]
Revenons alors à $f : S \longrightarrow T$, et considérons un groupoïde fibré sur $S$,
\[ G \ \longrightarrow\ S, \]
ce n'est pas une catégorie fibrée sur $T$, mais on peut lui associer
\[ t \longmapsto t\backslash G, \qquad\qquad T^o \longrightarrow (\mathrm{Cat}), \]
d'où une catégorie fibrée sur $T$, dont on prend le $\Pi$ relatif. Mais il est plus simple de prendre directement
\[ t \ \longmapsto\ \Pi(t\backslash G), \]
qui définit une catégorie fibrée en groupoïdes sur $T$.

\dgrblock{dgrb419}{The comma category computes integration}
\begin{dcomment}{exposition}{p2:qui définit une catégorie fibrée en groupoïdes sur}{The comma category computes integration}
The comma construction makes the variance visible. For $p:X\to S$, an object of $s\backslash_S X$ is a pair $(x,\alpha:s\to p(x))$; an arrow to $(x',\alpha')$ is an arrow $v:x\to x'$ with $p(v)\alpha=\alpha'$. Precomposing $\alpha$ gives the contravariant dependence on $s$. Thus the displayed formula is meaningful even when $p$ has no cartesian lifts.

When $p$ is fibred, those lifts explain the comparison with the actual fibre. Choose a normalized cleavage and send $(x,\alpha)$ to $\alpha^*x\in X_s$. If $v$ is a comma arrow, its composite with the chosen cartesian arrow $\alpha^*x\to x$ factors uniquely through $\alpha'^*x'\to x'$. This defines a functor $j_s:s\backslash_S X\to X_s$. Its composite with the inclusion $i_s$ is the identity, and the cartesian arrows give $i_sj_s\Rightarrow\mathrm{id}$. Applying the nerve yields a homotopy equivalence. The last natural transformation need not be an isomorphism, exactly as the manuscript warns; after applying $\Pi$, it becomes one. This also proves the asserted equivalence $H\simeq\Pi(H/S)$ when $H$ is fibred in groupoids.

For $f:S\to T$ and $G\to S$, the same construction gives
\[
 (f_!G)_t=\Pi(t\backslash_T G).
\]
The subscript on the comma is essential: its objects use arrows $t\to f(s)$ in $T$. Equivalently, if $F:S^{\mathrm{op}}\to\mathsf{Gpd}$ describes $G$, then
\[
 (f_!F)(t)\simeq
 \tau_{\leq1}\operatorname*{colim}_{(s,\alpha)\in(t\backslash f)^{\mathrm{op}}}F(s),
\]
where the colimit on the right is the homotopy colimit in spaces. The total category is the concrete way to form this colimit before truncation. In particular, the $s\backslash G$ printed in the unit on manuscript page~4 must be read as $f(s)\backslash_T G$; this is a notational slip in the edition, not an error introduced by the transcription.

For orientation, let $S=[0\to1]$, and let $F$ be the diagram $F(1)\to F(0)$. Integration over $S$ gives $F(0)$ up to equivalence, whereas its groupoid of sections is $F(1)$. Indeed, the construction above sends an object over $1$ to its inverse image over $0$, and its cartesian arrow becomes invertible in $\Pi G$. Sections are instead determined by their value over $1$. The two endpoints already separate homology from cohomology.
\end{dcomment}


\clearpage

% ===================== page 5 =====================
\printedpage{3}
% head: A. Grothendieck : Dérivateurs, Ch. VIII, Édition des Universités de Montpellier II et Paris VII
% foot: Version 2 février 2002

\msp{page 3}

Ainsi on trouve bien un foncteur (en fait, même un 2-foncteur de 2-catégories)
\[ f_! :\ \Grp(S) \ \longrightarrow\ \Grp(T). \]
Si
\[ u :\ G \ \longrightarrow\ G' \]
est un homomorphisme de groupoïdes fibrés sur $S$, qui induise des Cat-équivalences sur les fibres, alors
\[ f_!(u) :\ f_!(G) \ \longrightarrow\ f_!(G') \]
est aussi une Cat-équivalence sur les fibres. (Cela est formel, à partir du fait que $f_!$ est un 2-foncteur, compte tenu que si on a une équivalence sur les fibres, alors il y a un quasi-inverse -- c'est une histoire générale de catégories fibrées \dots)

Soit maintenant
\[ p :\ X \ \longrightarrow\ S, \]
et prenons le cas $G = \Pi(X/S)$. Je dis qu'on a alors une $T$-équivalence
\[ \boxed{\ f_!(\Pi(X/S)) \ \simeq\ \Pi(X/T),\ } \tag*{$(*)$} \]
en d'autres termes, pour connaître $\Pi(X/T)$ (à $T$-équivalence près), il suffit de connaître $\Pi(X/S)$, et de lui appliquer le foncteur $f_!$ d'intégration des 1-types d'homotopie relatifs, pour $f : S \longrightarrow T$.

Notons que l'on a trivialement
\[ \Pi(X/S) \ =\ p_!(e_X), \]
où $e_X$ désigne le $X$-groupoïde fibré trivial (toutes les fibres ponctuelles). Donc la formule précédente s'écrit comme une formule de transitivité
\[ f_!(p_!(\Gamma)) \ \simeq\ (fp)_!(\Gamma), \]
\msp{page 4}
où $\Gamma$ est le groupoïde fibré trivial sur $X$. Cela suggère de généraliser la formule $(*)$ en
\[ \boxed{\ (fp)_!(\Gamma) \ \simeq\ f_!(p_!(\Gamma))\ } \]
pour \textit{tout} groupoïde fibré sur $X$, dans une situation
\[ X \xrightarrow{\ p\ } S \xrightarrow{\ f\ } T. \]

Ce sera formel, si on établit une propriété de 2-adjonction pour le foncteur $f_!$ :

% ---- from Der-ChVIII-567_p006-010.tex
\printedpage{4}
% head: A. Grothendieck : Dérivateurs, Ch. VIII, Édition des Universités de Montpellier II et Paris VII
% foot: Version 2 février 2002

\textbf{Théorème.} \textit{Soit $f : S \fl T$ une flèche dans} (Cat)\textit{, $G$ un champ en groupoïdes sur $S$, $H$ un champ en groupoïdes sur $T$. On a alors une équivalence de catégories canonique}
\[
\boxed{\ \uHom_{\uGrp(S)}(G, f^*H) \ \simeq\ \uHom_{\uGrp(T)}(f_!G, H).\ }
\]

Les flèches d'adjonction
\[
\left\{
\begin{array}{lclcrcl}
f_!f^*H & \fl & H & \qquad\qquad & H & \in & \mathrm{Ob}\,\uGrp(T)\\[4pt]
f^*f_!G & \longleftarrow & G & \qquad\qquad & G & \in & \mathrm{Ob}\,\uGrp(S)
\end{array}
\right.
\]
devraient se définir sans problème. Pour la deuxième, il faut définir pour $s \in S$
\[
G_s \ \fl\ (f^*(f_!(G)))_s = (f_!(G))_{f(s)} \overset{\text{déf}}{=} \Pi(s\backslash G),
\]
or on a
\[
G_s \fl s\backslash G,
\]
d'où
\[
\underbrace{\Pi(G_s)}_{\begin{subarray}{l} =\ G_s\\ \text{car } G_s\\ \text{est un}\\ \text{groupoïde}\end{subarray}} \ \fl \Pi(s\backslash G),
\]
O.K. C'est fonctoriel en $s$ \dots

\msp{page 5}
Pour la flèche $f_!f^*H \fl H$ sur $T$, il faut trouver pour tout $t \in \mathrm{Ob}\,T$, et fonctoriellement en $t$,
\[
\underbrace{(f_!f^*H)_t}_{\overset{\text{déf}}{=}\Pi(t\backslash f^*(H))} \quad \fl \quad \underbrace{H_t}_{=\,\Pi(H_t)} \ ,
\]
il suffirait donc de définir
\[
t\backslash f^*(H) \ \fl\ H_t,
\]
mais on a le diagramme cartésien
\[
\begin{tikzcd}[column sep=large, row sep=large]
f^*H = H \times_S T \arrow[r] \arrow[d] & H \arrow[d] \\
S \arrow[r, "f"] & T \ni t
\end{tikzcd}
\]
\texttt{[plutôt} $H \times_T S$\texttt{]} induisant
\[
f^*H \ \fl\ H,
\]
\printedpage{5}
% head: A. Grothendieck : Dérivateurs, Ch. VIII, Édition des Universités de Montpellier II et Paris VII
% foot: Version 2 février 2002
d'où
\[
t\backslash f^*H \ \fl\ t\backslash H \qquad\quad \text{(au lieu de s'envoyer dans $H_t$),}
\]
mais d'autre part, comme $H$ est fibrée sur $T$, donc lisse sur $T$, l'inclusion
\[
H_t \fl t\backslash H
\]
est une Hot-équivalence. Ou disons plutôt : $f^*H \fl H$ est un morphisme de catégories sur $T$, induisant
\[
\underbrace{\Pi(f^*H/T)}_{\overset{\text{déf}}{=}\ f_!(f^*H)} \ \fl\ \Pi(H/T),
\]

\msp{page 6}
or le $T$-foncteur canonique
\[
H \fl \Pi(H/T) \qquad\quad (= \text{catégorie sur } T \text{ définie par } t \longmapsto \Pi(t\backslash H))
\]
est une équivalence, parce que $H$ est un fibré en groupoïdes sur $T$. Donc on peut trouver un $T$-quasi-inverse
\[
\Pi(H/T) \fl H
\]
et en composant on trouve
\[
f_!(f^*H) \ \fl\ \Pi(H/T) \ \fl\ H
\]
la flèche cherchée.

Ayant les 2-flèches d'adjonction, je m'arrête là dans les vérifications !

\dgrblock{dgrb420}{Completing the adjunction and its transitivity}
\begin{dcomment}{exposition}{p6:Ayant les 2-flèches d'adjonction}{Completing the adjunction and its transitivity}
The proposed unit and counit are the right ones. To finish their verification without making incompatible choices of quasi-inverses, it is useful to prove a slightly more general universal property. For any category $X\to T$, put $R_TX=\Pi(X/T)$. If $H\to T$ is fibred in groupoids, restriction along $X\to R_TX$ induces an equivalence
\[
 \operatorname{Hom}_T(R_TX,H)\simeq\operatorname{Fun}_T(X,H).
 \tag{VIII.1}
\]
Both sides are groupoids of functors and transformations over $T$.

Here is a direct construction of the inverse. Given $A:X\to H$ over $T$, choose cartesian lifts in $H$ and send $(x,\alpha:t\to p(x))$ to $\alpha^*A(x)$. For a comma arrow $v$, the cartesian factorization of $A(v)$ gives an arrow between these two objects in $H_t$. All such arrows are invertible, so the functor extends uniquely across the localization $t\backslash_T X\to\Pi(t\backslash_T X)$. The comparison isomorphisms of the cleavage supply its dependence on $t$; their cocycle equation follows from the uniqueness of cartesian factorization. A transformation $A\Rightarrow B$ extends by applying $\alpha^*$ to its component at $x$. Thus this construction is a functor on the Hom-groupoids, not just a bijection between classes of objects.

Restricting the extension to $(x,\mathrm{id}_{p(x)})$ recovers $A$. Conversely, an over-$T$ functor $B:R_TX\to H$ sends the cartesian arrow from $(x,\alpha)$ to $(x,\mathrm{id}_{p(x)})$ to a cartesian arrow in $H$. Cartesian factorization therefore gives a natural isomorphism between $B(x,\alpha)$ and $\alpha^*B(x,\mathrm{id})$. This identifies $B$ with the extension of its restriction, naturally also in transformations. It proves (VIII.1), including all the coherence needed for an adjunction.

Apply (VIII.1) to the total category of $G$ viewed over $T$. An over-$T$ functor $G\to H$ is exactly an over-$S$ functor $G\to f^*H$. We obtain the claimed equivalence of Hom-groupoids. The unit sends $g\in G_s$ to the class of $(g,\mathrm{id}_{f(s)})$. The counit sends the class of $((s,h),\alpha:t\to f(s))$ to $\alpha^*h$. The first triangle evaluates the latter formula at an identity arrow. The second uses the canonical cartesian arrow just described, so is the natural isomorphism between extension after restriction and the original functor. Their compatibility is the uniqueness assertion in the same cartesian factorization. Choosing different cleavages changes these formulas by a unique compatible comparison; it does not change the biadjunction.

Transitivity now follows at the same level of precision. For $X\xrightarrow{p}S\xrightarrow{f}T$,
\[
\begin{aligned}
 \operatorname{Hom}_T(f_!p_!\Gamma,H)
 &\simeq\operatorname{Hom}_S(p_!\Gamma,f^*H)\\
 &\simeq\operatorname{Hom}_X(\Gamma,p^*f^*H)\\
 &\simeq\operatorname{Hom}_T((fp)_!\Gamma,H).
\end{aligned}
\]
The bicategorical Yoneda argument identifies the two representing objects. For three composable maps, both parenthesizations represent the same composite restriction; the representing equivalences consequently satisfy the associativity constraint. Taking $\Gamma=e_X$ gives $p_!e_X=\Pi(X/S)$ and the first boxed transitivity formula. What the source leaves unchecked is therefore a coherence verification, not a missing hypothesis in the theorem.
\end{dcomment}


\printedpage{6}
% head: A. Grothendieck : Dérivateurs, Ch. VIII, Édition des Universités de Montpellier II et Paris VII
% foot: Version 2 février 2002

\section*{2\quad Questions de changement de base}

\[
\begin{tikzcd}[column sep=large, row sep=large]
S & S' \arrow[l, "q"'] \arrow[d, "f'"] \\
T \arrow[from=1-1, "f"'] & T' \arrow[l, "p"]
\arrow[from=2-2, to=1-1, bend right=50, shorten <=18pt, shorten >=18pt, "\gamma"']
\end{tikzcd}
\qquad\qquad\qquad
p^*f_!(G) \longleftarrow f'_!q^*(G)
\tag*{$(*)$}
\]

La flèche est définie, pourvu qu'on ait commutativité ou une flèche de commutativité $\gamma : pf' \fl fq$ (pas nécessairement \texttt{[un]} isomorphisme). Par adjonction, il revient au même de donner
\[
\begin{tikzcd}[row sep=large]
p_!f'_!q^*G \arrow[r] \arrow[d, dash, "\wr"', "\text{transitivité}"] & f_!(G) \\
(pf')_!q^*G \arrow[d, "\gamma"] & \\
(fq)_!q^*G \arrow[d, dash, "\text{transitivité}\ \wr"'] & \\
f_!(q_!q^*)G. \arrow[uuur, dashed, "\text{car } q_!q^*G \longrightarrow G"'] &
\end{tikzcd}
\]

On se demande quand la flèche est \texttt{[un]} isomorphisme.

\msp{page 7}
Le théorème suivant doit être vrai, et résulte de mes fourbis généraux.

% vertical bar printed in the left margin beside the following theorem
\textbf{Théorème.} \textit{Si le carré $(*)$ est} cartésien \textit{dans} (Cat)\textit{, $\gamma$ un isomorphisme, et si on suppose que $f$ soit} 1\textit{-lisse, ou $p$} 1\textit{-propre, alors le $T'$-foncteur de changement de base est une équivalence de $T'$-catégories.}

Les ‘fourbis’ généraux seraient en termes d'un dérivateur $\mathbf{D}$ ad hoc, celui des 1-types d'homotopie relatifs. J'y reviendrai par la suite.

\dgrblock{dgrb421}{Base change for truncated homotopy types}
\begin{dcomment}{later}{p7:Les ‘fourbis’ généraux seraient}{Base change for truncated homotopy types}
Grothendieck's next step asks whether integration can be performed after restricting the base. The natural transformation $\gamma:pf'\Rightarrow fq$ fixes the direction of the comparison: contravariance gives $q^*f^*\Rightarrow f'^*p^*$, whose mate is $f'_!q^*\Rightarrow p^*f_!$. At $t'$, its effect before groupoid completion is the evident functor sending a lift $\delta:t'\to f'(s')$ to the composite $p(t')\xrightarrow{p\delta}pf'(s')\xrightarrow{\gamma_{s'}}fq(s')$.

The theorem is valid for the derivator of $1$-types. More generally, set
\[
 \mathbb D_n(I)=h\operatorname{Fun}(NI^{\mathrm{op}},\mathcal S_{\leq n}),
 \qquad n\geq-1.
\]
Its left Kan extensions are homotopy left Kan extensions followed by $n$-truncation; its right Kan extensions are homotopy limits. The truncation is an accessible reflector, and the truncated category remains presentable \citep[Proposition~5.5.6.18 and Remark~5.5.6.21]{Lurie2009HTT}, so these Kan extensions exist for every small indexing category. An $n$-aspherical category is one whose $n$-truncation is a point, with the case $n=-1$ meaning nonemptiness. Smoothness and properness must be taken relative to this theory. They are conditions on the relevant comma categories and their changes of base, not merely on the homotopy types of the strict fibres.

A precise later formulation is \citet[Theorem~1.2, Corollary~4.16, Example~4.17, and \S5]{Maltsiniotis2005Asphericite}. Theorem~1.2 identifies smoothness with the asphericity of the fibre-to-comma inclusions; \S5 treats the corresponding base-change morphism. With his covariant convention replaced by the contravariant convention here, its two alternatives are exactly that the integrating map $f$ is smooth or that the restricting map $p$ is proper. Applying the statement to the localizer detected by fundamental groupoids gives equivalences of groupoids at every $t'$, hence an equivalence of fibred groupoids. No information above the first homotopy group is needed for this particular comparison.

This distinction between a correct integration theorem and an insufficient local-constancy criterion will matter on manuscript page~9. Truncating the \emph{values} of a diagram does not truncate all the coherence by which the values are assembled over its base.
\end{dcomment}


On se pose la question quand le 2-foncteur associé à $f : S \fl T$
\[
f^* :\ \uuGrp(T) \ \fl\ \uuGrp(S)
\]
\printedpage{7}
% head: A. Grothendieck : Dérivateurs, Ch. VIII, Édition des Universités de Montpellier II et Paris VII
% foot: Version 2 février 2002
est 0-fidèle, 1-fidèle, 2-fidèle (\footnote[2]{La $i$-fidélité de $f^*$ ($i \in \{0,1,2\}$) équivaut à celle des foncteurs (obtenus en prenant $G = e_T$)
\[
\underline{\Gamma}(H/T) \ \fl\ \underline{\Gamma}(f^*H/S).
\]}). Si $G$, $H$ sont des fibrés en groupoïdes sur $T$, on considère
\[
\uHom_{\uuGrp(T)}(G,H) \ \fl\ \uHom_{\uuGrp(S)}(f^*G, f^*H),
\]
et suivant que ce foncteur est toujours fidèle, ou pleinement fidèle, ou une équivalence de catégories, on dit que $f^*$ est 0-, 1-, ou 2-fidèle. Dans le cas 2-fidèle, cela implique que si $K$ est un fibré en groupoïdes sur $S$, s'il est équivalent à un $f^*G$, alors $G$ est déterminé à $T$-équivalence près, elle même déterminée à isomorphisme unique près.

\dgrblock{dgrb422}{Sections alone do not test full faithfulness}
\begin{dcomment}{examples}{p7:elle même déterminée à isomorphisme unique près.}{Sections alone do not test full faithfulness}
The three faithfulness conditions distinguish the lifting of transformations, of functors, and of their isomorphisms. The footnote proposes testing them using section groupoids alone. As an assertion for an arbitrary $f$, that reduction is false.

Let $T=[0\to1]$, $S=\{1\}$, and let $f$ be the inclusion. For every contravariant groupoid diagram $H=(H_1\to H_0)$, evaluation at $1$ is an equivalence $\Gamma(T,H)\simeq H_1$: a section consists of $h_1$, $h_0$, and an isomorphism $h_0\simeq u^*h_1$, and the choices of the latter pair form a contractible groupoid. Thus
\[
 \Gamma(T,H)\longrightarrow\Gamma(S,f^*H)
\]
is an equivalence for every $H$.

Nevertheless, take $G_1=H_1=\varnothing$, $G_0=*$, and $H_0=B(\mathbf Z/2)$. The only structure maps from the empty groupoid exist uniquely. A functor $G\to H$ is now just a functor $*\to B(\mathbf Z/2)$, and its natural automorphism group is $\mathbf Z/2$. Consequently
\[
 \operatorname{Hom}_T(G,H)\simeq B(\mathbf Z/2)
 \longrightarrow
 \operatorname{Hom}_S(f^*G,f^*H)\simeq *
\]
is not faithful. The proposed section test therefore fails even for $0$-faithfulness, despite satisfying the strongest of its own three tests.

There is a useful exact replacement. If $\pi_G:G\to T$ denotes the \emph{total category} of the first coefficient, then
\[
 \operatorname{Hom}_T(G,H)
 \simeq\Gamma(G,\pi_G^*H).
\]
Indeed, a section of $G\times_T H\to G$ is precisely a functor $G\to H$ over $T$, and the same identification holds for transformations. After restriction, the base becomes $S\times_T G$. Testing sections on these changed bases, with the pulled-back coefficients, is equivalent to testing the required Hom-groupoids. The missing base change cannot in general be suppressed by invoking an internal Hom: arbitrary restriction of presheaves does not preserve that internal Hom.
\end{dcomment}


\dgrblock{dgrb423}{Factorizations detect all three levels}
\begin{dcomment}{development}{p7:elle même déterminée à isomorphisme unique près.}{Factorizations detect all three levels}
Grothendieck's faithfulness question can be answered without imposing smoothness or properness. The appropriate test keeps both endpoints of an arrow of $T$, as did the factorization categories discussed in the commentary on Chapter~II, manuscript page~53. In the untruncated setting, \citet[arXiv version, Definition~6.9 and Theorem~6.10]{GrothStovicek2016Trees} give the corresponding criterion for homotopical epimorphisms. We spell out its elementary three-level refinement for groupoids. We did not find this combined formulation in the sources checked; the argument uses standard density and adjunction, and carries no claim of priority.

For an arrow $\beta:t\to t'$ of $T$, let $\mathcal F_f(\beta)$ have objects
\[
 (s,a:t\to f(s),b:f(s)\to t'),\qquad ba=\beta,
\]
and arrows $v:s\to s'$ satisfying $f(v)a=a'$ and $b'f(v)=b$.

\textbf{Proposition (the three tests).}\ours\quad
For $i\in\{0,1,2\}$, restriction $f^*:\mathsf{Gpd}(T)\to\mathsf{Gpd}(S)$ is $i$-faithful if and only if every $\mathcal F_f(\beta)$ is $(i-1)$-connected. Explicitly, these conditions are nonemptiness for $i=0$, nonempty connectedness for $i=1$, and $\Pi\mathcal F_f(\beta)\simeq*$ for $i=2$.

\emph{Proof.} The counit makes the factorization categories appear. Write $h_{t'}=\operatorname{Hom}_T(-,t')$, viewed as a discrete groupoid-valued presheaf. The comma formula gives, at $t$,
\[
 (f_!f^*h_{t'})_t
 \simeq\coprod_{\beta:t\to t'}\Pi\mathcal F_f(\beta)
 \longrightarrow \operatorname{Hom}_T(t,t').
 \tag{VIII.3}
\]
The counit maps every summand to its indexing arrow. The biadjunction identifies restriction on $\operatorname{Hom}_T(G,H)$ with precomposition by $\varepsilon_G:f_!f^*G\to G$.

To prove necessity at a fixed $t$, use a coefficient $H=t_*K$, where $t_*$ is the right biadjoint to evaluation at $t$ and $K$ is an arbitrary groupoid. The adjunction reduces the test for $G=h_{t'}$ to precomposition along (VIII.3), with values in $K$. If a summand is empty, take $K=B(\mathbf Z/2)$: an automorphism at that index is invisible after restriction, so faithfulness fails. If a summand has two connected components, the same $K$ allows different automorphisms on those components, which cannot descend from the single index; full faithfulness fails. If a summand is connected but has nontrivial fundamental group, take $K$ to be that summand's groupoid. Its identity functor, extended by constant functors on the other summands, is not isomorphic to a functor factoring through the discrete target. Essential surjectivity fails. These observations give precisely the three necessary conditions.

For sufficiency, first keep $G=h_{t'}$ and allow arbitrary $H$. Under the nonemptiness condition, a modification between two maps $h_{t'}\to H$ is determined by its pullback along $\varepsilon_{h_{t'}}$: at every $t$ and every $\beta$, there is an object over $\beta$ on which to evaluate it. Under connectedness, a modification between the pulled-back maps has equal components on all objects over a fixed $\beta$, since the two maps send every arrow of that groupoid to an identity. It therefore defines a unique component at $\beta$. These components are natural in $t$: reindex an object lying over $\beta$, and use naturality of the given modification. Thus the modification descends uniquely, proving full faithfulness. In the third case, (VIII.3) is itself an equivalence at every $t$, so precomposition is an equivalence of Hom-groupoids.

To pass from representables to an arbitrary $G$, one needs their coherent density, not just generation of connected components. Index the representables by the total category of $G$, sending $(t',g)$ to $h_{t'}$. Their homotopy colimit in groupoid-valued presheaves, evaluated at $t$, is $\Pi(t\backslash_T G)$, hence is equivalent to $G_t$ by the cartesian-lift argument already proved. Thus $G$ is this homotopy colimit. Restriction preserves it, because such colimits are computed pointwise. Taking maps into $H$, respectively into $f^*H$, changes these colimits into homotopy limits of the corresponding Hom-groupoids. Faithful functors remain faithful under these limits, because a modification is determined by its components. Fully faithful functors remain fully faithful, because the components lift uniquely and their compatibility follows by faithfulness. Equivalences remain equivalences by the universal property of a homotopy limit. Equivalently, the homotopy fibres of faithful and fully faithful functors of groupoids are, respectively, $0$-truncated and $(-1)$-truncated spaces; these conditions are preserved by limits. The representable tests therefore imply the required assertion for every $G,H$. This completes the proof.

At level zero the condition has a particularly concrete form: every $t$ is a retract of some $f(s)$. Indeed, apply nonemptiness to $\mathrm{id}_t$; conversely such a retraction supplies a factorization of every arrow with target $t$. For $f:Z\to e$, the sole factorization category is $Z$. The three tests become, respectively, that $Z$ is nonempty, connected, or $1$-connected. Even the strongest condition concerns maps between constant coefficients; it leaves room for the nonconstant diagram over $S^2$ constructed below.
\end{dcomment}


\msp{page 8}
Soit $i \in \mathbf{N}$. On dit que $f$ est $i$-asphérique si $\forall t \in T$, $Z = S/t$ est $i$-asphérique, i.e.
\[
\pi_0(Z),\ \dots,\ \pi_i(Z) \qquad\quad \text{triviaux.}
\tag*{$(*)$}
\]

On définit aussi $(-1)$-asphérique : les $S/t$ sont non-vides. De même pour la notion de $i$-coasphérique.

\dgrblock{dgrb424}{The valid descent assertions in Theorem 3}
\begin{dcomment}{exposition}{p8:De même pour la notion}{The valid descent assertions in Theorem 3}
Theorem~3 contains a useful distinction: full faithfulness of restriction describes how a coefficient already on $T$ is recovered, whereas its essential image describes which coefficients on $S$ descend. The hypotheses in parts~1 and~2 do give these conclusions. The assertion in part~3 needs an additional degree of connectivity.

We give the argument for the valid assertions. Put $n=i-1$. Under hypothesis~a, the inclusions $S_t\to S/t$ are $n$-coaspherical; under hypothesis~b, the inclusions $S_t\to t\backslash S$ are $n$-aspherical. This is the fibre-to-comma criterion for properness or smoothness, respectively. Combining it with the stated asphericity or coasphericity shows that every strict fibre $S_t$ is $n$-connected. The same criterion says that the categories of lifts of an arrow of $T$, with its source fixed in the proper case or its target fixed in the smooth case, are $n$-connected.

For $i=0$, every fibre is nonempty. A natural transformation between two functors $G\to H$ over $T$ is determined by its components over each object of $T$; choosing one object of $S_t$ detects each component. Restriction is therefore faithful on every Hom-groupoid.

For $i=1$, every fibre is connected. Suppose a transformation between the restrictions of two functors $G\to H$ has been given. At $g\in G_t$, its component calculated using $s\in S_t$ is independent of $s$, because naturality along the vertical arrows of $S_t$, and then along zigzags of such arrows, identifies the components. These components are natural within $G_t$. To verify naturality across an arrow $a:t\to t'$, choose a lift of $a$ with one endpoint fixed; its existence follows from the lift-category criterion. Naturality upstairs is the required naturality downstairs. The descended transformation is unique because every fibre is nonempty. This proves full faithfulness on Hom-groupoids.

For $i=2$, base change computes the relevant adjunction on each strict fibre. Under hypothesis~b, if $H$ is on $T$, the counit has the fibrewise form
\[
 (f_!f^*H)_t\simeq\Pi(S_t)\times H_t
 \longrightarrow H_t.
\]
Under hypothesis~a, use the right biadjoint $f_*$ instead: its unit is
\[
 H_t\longrightarrow(f_*f^*H)_t
 \simeq\operatorname{Fun}(\Pi(S_t),H_t).
\]
The right biadjoint is the groupoid of cartesian sections over the comma category $S/t$; the proper base-change formula replaces that comma category by $S_t$. Since $\Pi(S_t)\simeq*$, the displayed arrows are equivalences. The adjunction equivalences on Hom-groupoids now prove that $f^*$ is $2$-faithful. Notice that this argument uses constant coefficients \emph{on each fibre}; it does not assert that every locally constant coefficient there is constant.

The essential-image statement follows by applying the other adjunction arrow. Under hypothesis~b, its restriction to a fibre is the unit
\[
 G|_{S_t}\longrightarrow
 \operatorname{const}_{S_t}\!\left(\Pi\!\int_{S_t}G\right).
\]
If $G|_{S_t}$ is equivalent to a constant groupoid $K$, integration gives $\Pi(S_t)\times K\simeq K$, so this unit is an equivalence. Thus $G\to f^*f_!G$ is fibrewise an equivalence, and $G$ descends. Under hypothesis~a, use $f^*f_*G\to G$ and the same argument with sections: $\operatorname{Fun}(\Pi(S_t),K)\simeq K$. Necessity is immediate, because a coefficient pulled back from $T$ restricts to a constant coefficient on each strict fibre. This proves part~2 with precisely its printed fibrewise-\emph{constant} condition. Its replacement by equivalence of transport functors requires a separate argument and is where the next difficulty lies.
\end{dcomment}


En fait, pour $i \in [-1,+\infty[ \ \subseteq \mathbf{Z}$, on définit un dérivateur $\mathbf{D}_i$ des $i$-types d'homotopie relatifs (cf. plus bas). D'où les notions d'objet $\mathbf{D}_i$-asphérique ($=$ $i$-asphérique), de $\mathbf{D}_i$-équivalence, de foncteur $\mathbf{D}_i$-asphérique et $\mathbf{D}_i$-coasphérique, de foncteur $\mathbf{D}_i$-lisse et $\mathbf{D}_i$-propre, de $\mathbf{D}_i$-fibration et de $\mathbf{D}_i$-fibration \textit{triviale} (i.e. à fibres $\mathbf{D}_i$-asphériques). On doit avoir le

\textbf{Théorème 3.} \textit{Soit $f : S \fl T$ dans} (Cat)\textit{, $i \in \{0,1,2\}$. Supposons l'une des conditions suivantes satisfaite :}

\begin{itemize}[leftmargin=3em]
\item[\textit{a)}] \textit{$f$ est $(i-1)$-propre, et $(i-1)$-asphérique (s'exprime par la condition que certaines catégories, telles les $S/t$, soient $(i-1)$-asphériques).}
\item[\textit{b)}] \textit{$f$ est $(i-1)$-lisse, et $(i-1)$-coasphérique.}
\end{itemize}

\textit{Alors :}

\begin{itemize}[leftmargin=3em]
\item[\textit{1\textsuperscript{o})}] \textit{Le foncteur $f^*$ image inverse}
\[
f^* : \uGrp(T) \ \fl\ \uGrp(S)
\]
\textit{est $i$-fidèle.}
\item[\textit{2\textsuperscript{o})}] \textit{Dans le cas $i = 2$ ($f$ pleinement fidèle), on}
\printedpage{8}
% head: A. Grothendieck : Dérivateurs, Ch. VIII, Édition des Universités de Montpellier II et Paris VII
% foot: Version 2 février 2002
\msp{page 9}
\textit{a de plus ceci : si $G \in \mathrm{Ob}\,\uGrp(S)$ est un fibré en groupoïdes sur $S$, alors $G$ est équivalent à un fibré en groupoïdes de la forme $f^*(H)$ si et seulement si pour tout $t \in T$, la restriction $G_t$ de $G$ à $S_t$ est équivalente à un fibré en groupoïdes constant sur $S_t$. (}\textbf{NB} \textit{La condition a) ou b) implique déjà que $S_t$ est} 1\textit{-connexe, i.e. que son groupoïde fondamental est équivalent à $e$.)}
\item[\textit{3\textsuperscript{o})}] \textit{Soit $Z$ un objet} 1\textit{-connexe dans} (Cat)\textit{, $G$ un fibré en groupoïdes sur $Z$. Pour que $G$ soit $Z$-équivalent à un groupoïde constant sur $Z$, il faut et il suffit que pour toute flèche $u$ dans $Z$, le morphisme de groupoïdes $u^*$ induit sur les fibres de $G$ soit une $Z$-équivalence} (\footnote[3]{On dit alors que $G$ est \textit{essentiellement localement constant} sur $Z$.}).
\end{itemize}

\textbf{Corollaire.} \textit{La condition nécessaire et suffisante de 2\textsuperscript{o}) s'écrit aussi ainsi : pour toute flèche $u$ de $S$ ‘verticale’ sur $T$, i.e. telle que $f(u)$ soit une flèche identique, le foncteur $u^*$ induit par $u$ sur les fibres de $G$ soit une équivalence de groupoïdes.}

\dgrblock{dgrb425}{A finite obstruction in degree two}
\begin{dcomment}{examples}{p9:le foncteur $u^*$ induit par}{A finite obstruction in degree two}
Grothendieck wants a practical descent test that can be checked on the arrows of the fibres. The obstruction to this last simplification already appears over a finite poset, with every fibre the same two-element automorphism group. It makes the distinction from the preceding valid argument quite explicit. The general coherence issue was isolated in the commentary on Chapter~VI, manuscript pages~7 and~74; here it produces a counterexample to a stated theorem.

Let $Z$ be the poset of nonempty proper subsets of $\{0,1,2,3\}$, ordered by inclusion. It has fourteen objects. Its nerve is the barycentric subdivision of the boundary of a tetrahedron, so $|NZ|\cong S^2$ and $\Pi Z\simeq*$. Write $A=\mathbf F_2$. On the strictly increasing triples in $Z$, define
\[
 c(\sigma,\tau,\upsilon)=
 \begin{cases}
 1,&(\sigma,\tau,\upsilon)=(\{0\},\{0,1\},\{0,1,2\}),\\
 0,&\text{otherwise}.
 \end{cases}
\]
Extend $c$ by zero to triples with a repeated object. There is no strictly increasing chain of four objects, so this normalized $2$-cochain is a cocycle.

It is useful to verify directly that it is not a coboundary. Let $z$ be the sum, with coefficients in $\mathbf F_2$, of the twenty-four flags vertex $\subset$ edge $\subset$ face of the tetrahedral boundary. Every edge of this order complex occurs in exactly two such flags: a vertex--edge pair extends to two faces, a vertex--face pair has two intervening edges, and an edge--face pair has two possible vertices. Hence $\partial z=0$, whereas $\langle c,z\rangle=1$. If $c=\delta b$, this pairing would equal $\langle b,\partial z\rangle=0$, a contradiction.

Now define a category $E_c$ over $Z$. There is one object over each $\sigma$, and for $\sigma\subseteq\tau$ the arrows over that inclusion are labelled by $A$; there are no other arrows. The identity has label zero, and composition is
\[
 (\tau\subseteq\upsilon,b)\circ(\sigma\subseteq\tau,a)
 =\bigl(\sigma\subseteq\upsilon,
        a+b+c(\sigma,\tau,\upsilon)\bigr).
 \tag{VIII.2}
\]
The cocycle equation is exactly associativity; normalization gives the unit equations. Each fibre is $BA$. Every arrow is cartesian: composition with a fixed arrow carries the set of lifts of any preceding inclusion bijectively onto the set of lifts of its composite, since in (VIII.2) it adds a fixed element of $A$. Thus $E_c\to Z$ is a category fibred in groupoids. Choosing the lift labelled zero for each inclusion makes every inverse-image functor $BA\to BA$ the identity. Only the comparison isomorphisms for successive inverse images are nontrivial, and they are encoded by $c$. Passing to an equivalent split fibration gives the same counterexample in terms of a strict groupoid-valued functor; the obstruction does not arise from a choice of presentation.

Suppose that $E_c$ were equivalent over $Z$ to a constant fibred groupoid. Its fibre would have to be equivalent to $BA$, so the constant fibred groupoid would admit a section. An equivalence over $Z$ and its quasi-inverse transport that section to a section of $E_c$. Such a section chooses a label $b(\sigma,\tau)$ on each inclusion; its functoriality says
\[
 b(\sigma,\upsilon)
 =b(\sigma,\tau)+b(\tau,\upsilon)+c(\sigma,\tau,\upsilon).
\]
This is precisely $\delta b=c$, which has just been excluded. Therefore $E_c$ is not constant, although $Z$ is $1$-connected and every transport functor is an equivalence. This disproves part~3.

It also disproves the corollary under its own hypotheses. Take $f:Z\to e$. A functor to $e$ is both a Grothendieck fibration and an opfibration, hence smooth and proper for the localizer of $1$-types. Its sole fibre is the $1$-connected category $Z$, so hypotheses~a and~b for $i=2$ both hold. Every arrow is vertical and satisfies the proposed transport test, but $E_c$ is not in the image of $f^*$. Part~2 itself remains valid: its actual condition demands that $E_c|_Z$ be constant, which it is not.
\end{dcomment}


\dgrblock{dgrb426}{The corrected local-constancy criterion}
\begin{dcomment}{status}{p9:le foncteur $u^*$ induit par}{The corrected local-constancy criterion}
The missing datum is classical $2$-monodromy. A locally constant diagram with fibre a groupoid $K$ determines a classifying map
\[
 |NZ|\longrightarrow B\operatorname{Aut}(K),
\]
where $\operatorname{Aut}(K)$ is the groupoid of self-equivalences of $K$ and natural isomorphisms. The target is a $2$-type: its first homotopy group records equivalence classes of self-equivalences, and its second is $\operatorname{Aut}(\mathrm{id}_K)$. A simply connected base removes the first kind of monodromy but can still support the second. In the finite example it is the nonzero class in $H^2(S^2;\mathbf F_2)$.

A precise classification is given by \citet[arXiv version, Theorem~2.2.5 and \S2.3]{PoleselloWaschkies2005}; their Corollary~2.2.4 gives constancy over a $2$-connected base. Applied to the CW complex $|NZ|$, it proves the corrected statement: if $Z$ is $2$-connected, a groupoid diagram is equivalent to a constant one if and only if all its transport functors are equivalences. Alternatively, over a $1$-connected base one must require its remaining $2$-monodromy to vanish. This is a known obstruction, not a new conjecture or a priority claim. Strengthening the corollary by requiring $2$-connected strict fibres, while retaining the hypotheses of part~2, makes the vertical-arrow criterion sufficient by applying this corrected statement separately to every fibre.
\end{dcomment}
\dgrchapter{IX}{Chapter IX}{\texorpdfstring{Retour sur les catégories de fractions $\mathcal{M}W^{-1}$, comparaison avec Quillen}{Retour sur les catégories de fractions MW-1, comparaison avec Quillen}}
\dgrpart{0}{VIII-X}{Manuscript pages 1--25}{VIII-X\enspace\textperiodcentered\enspace Chapter IX\enspace\textperiodcentered\enspace manuscript pages 1--25}{Chapter IX \textperiodcentered\ \textup{VIII-X}}
% ---- from Der-ChIX-678_p001-005.tex
\printedpage{[non numérotée]}
% head: (aucun titre courant imprimé)

\begin{center}
{\LARGE\textsf{LES DÉRIVATEURS}}

\bigskip

\textbf{ALEXANDRE GROTHENDIECK}

\bigskip\bigskip

\textbf{Chapitre IX}

\textbf{Retour sur les catégories de fractions $\boldsymbol{\mathcal{M}W^{-1}}$,\\
comparaison avec Quillen}
\end{center}

\clearpage

\printedpage{[non numérotée]}
% head: (aucun titre courant imprimé)

Ce texte a été déchiffré et transcrit en \LaTeXe{} par M. Künzer. Il a été édité par M. Künzer, J. Malgoire, et G. Maltsiniotis. La transcription est aussi fidèle que possible au manuscrit. Pour les quelques corrections évidentes, ou rares commentaires des éditeurs, ainsi que pour la numérotation originale des pages du manuscrit, les caractères de machine à écrire \texttt{[typewriter]} entre crochets sont utilisés. Un point d'interrogation entre crochets signifie que l'on n'est pas sûr du mot qui précède.

Cette édition est provisoire. Les remarques, commentaires et corrections sont bienvenus. Envoyer un message à :

\textbf{maltsin@math.jussieu.fr}

\hfill G. Maltsiniotis\hspace*{2cm}

\clearpage

\printedpage{1}
% head: A. Grothendieck : Dérivateurs, Ch. IX, Édition des Universités de Montpellier II et Paris VII
% foot: Version 12 mars 2002

\msp{page 1}

\section*{1\quad Fonctorialité de $\Ch^{\Sigma}_{\infty}(a,b)$ en $a$, $b$ dans $\mathcal{M}\Sigma^{-1}$}

On se donne une catégorie $\mathcal{M}$ avec \texttt{[un]} ensemble admissible de flèches $\Sigma \subseteq \mathrm{Fl}(\mathcal{M})$, d'où (en prenant p.~ex. $\underline{t} = (t_n)$, avec $t_n = (\underbrace{\Delta_+\Delta_-}_{\tau_0})^{2n}$, et $\underbrace{t_{n+1}}_{=\,\tau_0^{2n+2}} \longrightarrow \underbrace{t_n}_{=\,\tau_0^{2n}}$ défini comme $p * \mathrm{id}_{\tau_0^{2n}} * p$, où $p : \tau_0 \longrightarrow 1$ la dégénérescence.) On a alors une théorie des $\underline{\Ch}_{\Sigma,\infty}(\mathcal{M},a,b)$, que je note simplement $\Ch_\infty(a,b)$ en tant qu'objet de Hot. Des accouplements
\[
\Ch_\infty(a,b) \times \Ch_\infty(b,c) \longrightarrow \Ch_\infty(a,c) \qquad \text{dans Hot}
\]
avec (nous allons l'admettre) \textit{associativité} dans Hot. Cela définit donc des accouplements sur les $\pi_0(\Ch_\infty(a,b)) \overset{\text{déf}}{=} \pi_0(a,b)$, qui n'est autre que la composition des flèches dans $\mathcal{M}\Sigma^{-1}$. D'autre part, on a pour $a,b \in \mathrm{Ob}\,\mathcal{M}$
\[
\Hom_{\mathcal{M}}(a,b) \longrightarrow \pi_0(\Ch_\infty(a,b)) = \pi_0(a,b),
\]
compatible avec les compositions. Mais on peut l'ignorer pour l'instant, et regarder $\mathrm{Ob}\,\mathcal{M}$, ou simplement $\mathcal{M}_0$, avec la seule structure
\[
\left\{
\begin{array}{ccccccc}
\mathcal{M}_0 \times \mathcal{M}_0 & \longrightarrow & \Hot & \qquad & (a,b) & \longmapsto & \Ch_\infty(a,b) \\[1ex]
\mathcal{M}_0 \times \mathcal{M}_0 \times \mathcal{M}_0 & \longrightarrow & \mathrm{Fl}\,\Hot & \qquad & (a,b,c) & \longmapsto & (\Ch_\infty(a,b) \times \Ch_\infty(b,c) \underbrace{\longrightarrow}_{\substack{\text{flèche}\\ \text{dans Hot}}} \Ch_\infty(a,c))
\end{array}
\right.
\]
avec condition \texttt{[d']}associativité, et aussi une condition d'unités :
\[
\Ch_\infty(a,a) \text{ est muni d'un } 1_a \in \pi_0\Ch_\infty(a,a) = \pi_0(a,a),
\]
jouant le rôle d'une `identité bilatère', dans le sens suivant. Notons que les accouplements
\[
\begin{array}{ccccc}
\Ch_\infty(a,b) & \times & \Ch_\infty(b,c) & \longrightarrow & \Ch_\infty(a,c) \\[1ex]
P & \times & Q & \longrightarrow & R
\end{array}
\]
\msp{page 2}
définissent
\[
\begin{array}{ccc}
P & \longrightarrow & \underline{\Hom}_{\Hot}(Q,R) \\[1ex]
Q & \longrightarrow & \underline{\Hom}_{\Hot}(P,R),
\end{array}
\]
d'où en passant aux $\pi_0$
\[
\begin{array}{ccc}
\pi_0(P) & \longrightarrow & \Hom_{\Hot}(Q,R) \\[1ex]
\pi_0(Q) & \longrightarrow & \Hom_{\Hot}(P,R),
\end{array}
\]
\printedpage{2}
% head: A. Grothendieck : Dérivateurs, Ch. IX, Édition des Universités de Montpellier II et Paris VII
% foot: Version 12 mars 2002
en d'autres termes,
\[
\begin{array}{ccc}
\pi_0(a,b) & \longrightarrow & \Hom_{\Hot}(\Ch_\infty(b,c),\Ch_\infty(a,c)) \\[1ex]
\pi_0(b,c) & \longrightarrow & \Hom_{\Hot}(\Ch_\infty(a,b),\Ch_\infty(a,c)).
\end{array}
\]

\dgrblock{dgrb427}{Composition before functoriality}
\begin{dcomment}{exposition}{p1:avec (nous allons l'admettre)}{Composition before functoriality}
Grothendieck wants a morphism in the localized category to act on an entire space of paths, not merely on its set of components. This is why he temporarily forgets how the path categories were constructed. Once composition and its units are available, change of either endpoint is a formal operation. The hypothesis being made is substantial: associativity is assumed for the particular spaces $\Ch_\infty(a,b)$. The abstract argument that follows does not establish that hypothesis.

Here is the argument with all its variance visible. Write $C(a,b)$ for the proposed Hom object in a monoidal category $(\mathcal H,\otimes,I)$, and suppose that
\[
 m_{abc}:C(a,b)\otimes C(b,c)\longrightarrow C(a,c),
 \qquad j_a:I\longrightarrow C(a,a)
\]
satisfy associativity and the two unit equations, with the associator and unitors of $\mathcal H$ understood. This is enrichment, with the factors ordered in the direction in which a path is traversed; reversing the tensor convention gives the usual notation of \citet[\S\S1.1--1.3]{Kelly1982}. Symmetry is not required. In particular, the Cartesian signs in the general displays on manuscript page 7 should be read as the specified monoidal product $*$.

The ordinary category $C_0$ has
\[
 C_0(a,b)=\operatorname{Hom}_{\mathcal H}(I,C(a,b)).
\]
For $u:I\to C(a,b)$ and $v:I\to C(b,c)$ its composition is
\[
 I\xrightarrow{\sim}I\otimes I
   \xrightarrow{u\otimes v}C(a,b)\otimes C(b,c)
   \xrightarrow{m_{abc}}C(a,c).
\]
The enriched associativity and unit equations give the ordinary ones after precomposition with $I$. Thus the $j_a$ become the identities of $C_0$. Their uniqueness follows already in the endomorphism monoid: if $j$ and $j'$ are both two-sided units, then $j=j\circ j'=j'$. Existence as units on the Hom objects is a separate assumption; an element that is a unit only after applying $\operatorname{Hom}(I,-)$ need not satisfy the stronger enriched equations.

To change endpoints, let $u:a'\to a$ and $v:b\to b'$ belong to $C_0$. Insert the corresponding maps out of $I$ on the two sides of $C(a,b)$ and compose twice. This gives
\[
 C(u,v):C(a,b)\longrightarrow C(a',b').
\]
For composable $u':a''\to a'$ and $v':b'\to b''$, the two composites $C(u',v')C(u,v)$ and $C(uu',v'v)$ are the two parenthesizations of the same fivefold composition, with $u',u$ on the left and $v,v'$ on the right. Associativity identifies them. The unit equations identify $C(1_a,1_b)$ with the identity. Inserting only the left or only the right element also shows that the two actions commute. We have therefore obtained an ordinary bifunctor
\[
 C_0^{\mathrm{op}}\times C_0\longrightarrow\mathcal H.
\]
Applying $\operatorname{Hom}_{\mathcal H}(I,-)$ gives the ordinary Hom bifunctor, including its action on morphisms, rather than only objectwise bijections.

For $\mathcal H=\Hot$, $I$ is a point and $\operatorname{Hom}(I,-)$ is $\pi_0$. One can obtain the two endpoint actions by currying in the Cartesian closed homotopy category, as Grothendieck does on page 2, but the insertion argument shows that closedness is unnecessary. This economy matters when he passes to ind-objects: it is the monoidal structure and its coherence that the construction uses.
\end{dcomment}


On en déduit en particulier pour $a,b \in \mathcal{M}_0$
\[
\left\{
\begin{array}{cccc}
1) & \pi_0(a,a) & \longrightarrow & \Hom_{\Hot}(\Ch_\infty(a,b),\Ch_\infty(a,b)) \\[1ex]
2) & \pi_0(a,a) & \longrightarrow & \Hom_{\Hot}(\Ch_\infty(b,a),\Ch_\infty(b,a)),
\end{array}
\right.
\tag*{$(*)$}
\]
et un objet $\alpha_a$ de $\pi_0(a,a)$ est unité bilatère si ses deux images sont les homomorphismes identiques dans $\Ch_\infty(a,b)$ resp. $\Ch_\infty(b,a)$. Passant aux $\pi_0$ de ceux-ci, cela implique que $\alpha_a$ est unité bilatère pour la composition des $\pi_0(x,y)$, relativement à $b$. Si on veut que ce soit le cas pour \textit{tout $b$,} cela signifie (prenant $b = a$) que $\alpha_a$ est un objet unité bilatère pour la structure de composition associative (donc de monoïde) dans $\pi_0(a,a)$, donc que $\pi_0(a,a)$ est un monoïde \textit{unitaire,} ayant $\alpha_a$ pour objet unité. L'élément $\alpha_a$ est donc déterminé de façon unique. Et on veut que $\alpha_a$ existe
\msp{page 3}
pour tout $a \in \mathcal{M}_0$, donc $\mathcal{M}_0$, muni des $\pi_0(a,b)$ et de leur composition, devient une \textit{catégorie,} soit $\mathcal{M}_\Pi$; et les flèches identiques de celle-ci doivent jouer le rôle d'unité bilatère également pour les `Hom externes' $\Ch_\infty(a,b)$, à valeurs dans Hot.

On peut dire que pour $a = a_0$ fixé, les $\Ch_\infty(a_0,b)$ ($b$ variable) forment les valeurs d'un \textit{foncteur} covariant
\[
\mathcal{M}_\Pi \longrightarrow \Hot \qquad b \longmapsto \Ch_\infty(a_0,b),
\]
en associant à
\[
b \xrightarrow{\ u\ } b', \qquad u \in \pi_0(b,b'),
\]
la flèche
\[
\Ch_\infty(a,b) \longrightarrow \Ch_\infty(b') % [sic] \Ch_\infty(b') -> \Ch_\infty(a_0,b'), argument manquant (cf. la note des éditeurs qui suit)
\]
\texttt{[plutôt} $\Ch_\infty(a_0,b) \longrightarrow \Ch_\infty(a_0,b')$\texttt{]} correspondante $*_1$ [?]. Que ce soit bien là une loi fonctorielle doit résulter formellement de l'associativité des accouplements sur les $\Ch_\infty$. De même, on a une loi fonctorielle, pour $b$ fixé,
\[
\mathcal{M}^o_\Pi \longrightarrow \Hot, \qquad a \longmapsto \Ch_\infty(a,b),
\]
grâce à l'associativité. En fait, l'associativité impliquera aussi qu'on a un \textit{bifoncteur}
\[
\begin{array}{ccc}
\mathcal{M}^o_\Pi \times \mathcal{M}_\Pi & \longrightarrow & \Hot \\[1ex]
(a\,,\,b) & \longmapsto & \Ch_\infty(a,b),
\end{array}
\]
\printedpage{3}
% head: A. Grothendieck : Dérivateurs, Ch. IX, Édition des Universités de Montpellier II et Paris VII
% foot: Version 12 mars 2002
\msp{page 4}
car si
\[
a' \xrightarrow{\ u\ } a, \qquad b \xrightarrow{\ v\ } b',
\]
on a le diagramme
\[
\begin{tikzcd}[column sep=large, row sep=large]
\Ch_\infty(a,b) \arrow[r, "\substack{\text{comp. avec}\\ v}"] \arrow[d, "\text{comp. avec } u"'] & \Ch_\infty(a,b') \arrow[d, "\text{comp. avec } u"] \\
\Ch_\infty(a',b) \arrow[r, "\substack{\text{comp. avec}\\ v}"] & \Ch_\infty(a',b'),
\end{tikzcd}
\]
qui est commutatif grâce à l'associativité pour
\[
\Ch_\infty(a',a),\ \Ch_\infty(a,b),\ \Ch_\infty(b,b').
\]

\textbf{Corollaire} de cette dépendance bifonctorielle : \textit{Si $u : a' \longrightarrow a$, $v : b \longrightarrow b'$ sont dans $\mathcal{M}_\Pi$, alors la flèche correspondante}
\[
\Ch_\infty(a,b) \longrightarrow \Ch_\infty(a',b') \qquad \textit{dans } \Hot
\]
\textit{est un} isomorphisme.

\dgrblock{dgrb428}{An invertibility hypothesis already restored on page 8}
\begin{dcomment}{examples}{p4:de cette dépendance bifonctorielle}{An invertibility hypothesis already restored on page 8}
The intended application is that weakly equivalent endpoints have the same Hom object. For this application $u$ and $v$ are invertible in the localization. The corollary as printed has omitted this condition; manuscript page 8 states the correct version explicitly.

A discrete example detects the omission without any homotopy theory. Enrich the ordered category $0<1$ in sets, regarded as discrete spaces, by its usual Hom sets. Change the first variable by the noninvertible map $u:0\to1$ and fix the second variable at $0$. The induced map is
\[
 C(1,0)=\varnothing\longrightarrow C(0,0)=\{1_0\}.
\]
This is the endpoint action of $u:0\to1$ with the second variable fixed at $0$, and is not an isomorphism. If $u$ and $v$ are invertible, functoriality supplies the inverse $C(u^{-1},v^{-1})$, proving the repaired corollary.

The following display also omits the opposite on the first factor: its domain is $\mathcal M_\Pi^{\mathrm{op}}\times\mathcal M_\Pi$. Both omissions occur in the edition, so they are retained in the source and entered separately in the audit. They do not invalidate the preceding construction of the bifunctor.
\end{dcomment}


Les mêmes arguments devraient marcher si au lieu des $\Hom\,\mathrm{ext}$ à valeurs dans Hot, on les prenait dans Ind\,Hot, ou dans Cat\,hot, ou dans Ind\,Cat\,hot -- je pense bien sûr aux $\mathrm{CH}_\infty(a,b)$ (dans la théorie de la localisation) dans Ind\,Cat\,hot, qui peuvent être regardés comme provenant d'un bifoncteur (où $\mathcal{M}_\Pi = \mathcal{M}\Sigma^{-1}$)
\[
\mathcal{M}_\Pi \times \mathcal{M}_\Pi \longrightarrow \mathrm{Ind}\,\mathrm{Cat}\,\mathrm{hot}.
\]

\msp{page 5}
Voici l'axiomatisation définitive (?). On remplace Hot par une catégorie

\begin{itemize}[label={$\mathcal{H}$,}]
\item catégorie avec produit $*$ associatif, commutatif (?) (\footnote{On ne doit pas avoir besoin de la commutativité.}), unitaire, et toutes les compatibilités habituelles au besoin,
\end{itemize}

l'objet unité pour $*$ sera noté 1, et on notera
\[
\left\{
\begin{array}{ccc}
\pi_0(\xi) & \overset{\text{déf}}{=} & \Hom_{\mathcal{H}}(1,\xi) \\[1ex]
\pi_0 : \mathcal{H} & \longrightarrow & \mathrm{Ens}.
\end{array}
\right.
\]

Si on a
\[
P * Q \xrightarrow{\ u\ } R,
\]

% ---- from Der-ChIX-678_p006-010.tex
\printedpage{4}
% head: A. Grothendieck : Dérivateurs, Ch. IX, Édition des Universités de Montpellier II et Paris VII
% foot: Version 12 mars 2002

on en déduit
\[
\left\{
\begin{array}{lcl}
\pi_0(P) & \xrightarrow{c_u^{\mathrm{l}}} & \mathrm{Hom}_{\mathcal{H}}(Q,R)\\
\pi_0(Q) & \xrightarrow{c_u^{\mathrm{r}}} & \mathrm{Hom}_{\mathcal{H}}(P,R)
\end{array}
\right.
\]
ainsi :
\[
\begin{array}{rlcl}
\text{Si } \alpha\in\pi_0(P),\ \text{i.e. } 1\xrightarrow{\ \alpha\ }P,\ \text{alors} & \multicolumn{3}{l}{\alpha*\mathrm{id}_Q : Q\longrightarrow P*Q}\\
\text{et} & c^{\mathrm{l}}_u(\alpha) & = & u\circ(\alpha*\mathrm{id}_Q)\\
\text{et dualement} & c^{\mathrm{r}}_u(\beta) & = & u\circ(\mathrm{id}_P*\beta).
\end{array}
\]
Si on a trois objets $P_0$, $P_1$, $P_2$ dans $\mathcal{H}$, des accouplements
\[
\left\{
\begin{array}{rcl}
P_0*P_1 & \xrightarrow{u_1} & Q_0\\
P_1*P_2 & \xrightarrow{u_2} & Q_1\\
P_0*Q_1 & \xrightarrow{u_3} & R\\
Q_0*P_2 & \xrightarrow{u_4} & R
\end{array}
\right.
\]
\[
\begin{tikzcd}
a \arrow[r, "P_0"'] \arrow[rr, bend left=20, "Q_0"] \arrow[rrr, bend left=35, "R"] & b \arrow[r, "P_1"'] \arrow[rr, bend right=20, "Q_1"'] & c \arrow[r, "P_2"'] & d
\end{tikzcd}
\]

\msp{page 6}
avec une relation d'associativité, disant que
\[
\begin{tikzcd}
P_0*P_1*P_2 \arrow[r] \arrow[d] & Q_0*P_2 \arrow[d]\\
P_0*Q_1 \arrow[r] & R
\end{tikzcd}
\]
est commutatif, alors on en déduit deux relations de transitivité et une de commutation, pour les opérations associées aux éléments de $\pi_0(P_0)$, $\pi_0(P_1)$, $\pi_0(P_2)$, qu'on n'explicitera pas, sauf l'une :
\[
\begin{array}{rrrcl}
\multicolumn{5}{l}{\text{Soit } \beta\in P_1,\ \gamma\in P_2,}\\
\text{d'où par } \beta : & c_{u_1}(\beta) : & P_0 & \longrightarrow & Q_0,\\
\text{et par } \gamma : & c_{u_4}(\gamma) : & Q_0 & \longrightarrow & R,\\
\text{et par } \beta *_{u_2}\gamma\in\pi_0(Q_1) : & c_{u_3}(\beta\,\texttt{[}{*_{u_2}}\texttt{]}\gamma) : & P_0 & \longrightarrow & R,
\end{array}
\]

\printedpage{5}
% head: A. Grothendieck : Dérivateurs, Ch. IX, Édition des Universités de Montpellier II et Paris VII
% foot: Version 12 mars 2002
\texttt{[plutôt} $\beta\in\pi_0(P_1)$\texttt{,} $\gamma\in\pi_0(P_2)$\texttt{]}, et bien on a
\[
\boxed{\;c(\beta *_{u_2}\gamma) \;=\; c(\gamma)c(\beta).\;}
\]
Il faut encore définir $\beta *_{u_2}\gamma$, via l'accouplement $P_1*P_2\xrightarrow{u_2}Q_1$, on a en effet
\[
\begin{array}{ccc}
\pi_0(P_1)\times\pi_0(P_2) & \longrightarrow & \pi_0(P_1*P_2)\\
(\beta\,,\,\gamma) & \longmapsto & \beta*\gamma,
\end{array}
\]
et on compose avec
\[
\pi_0(u_2) : P_1*P_2 \longrightarrow Q_1
\]
\texttt{[plutôt} $\pi_0(u_2) : \pi_0(P_1*P_2) \longrightarrow \pi_0(Q_1)$\texttt{]}.

\msp{page 7}
On se donne un ensemble $\mathcal{M}_0$, et
\[
\left\{
\begin{array}{rclccrcl}
\mathcal{M}_0\times\mathcal{M}_0 & \xrightarrow{\mathrm{Ch}} & \mathcal{H} & \text{noté} & & (a,b) & \longmapsto & \mathrm{Ch}(a,b)\ \text{ou}\ \mathrm{Hom\,ext}_{\mathcal{H}}(a,b)\\
\mathcal{M}_0\times\mathcal{M}_0\times\mathcal{M}_0 & \xrightarrow{\text{comp.}} & \mathrm{Fl}\,\mathcal{H} & & & (a,b,c) & \longmapsto & \bigl(\mathrm{Ch}(a,b)\times\mathrm{Ch}(b,c)\longrightarrow\mathrm{Ch}(a,c)\bigr)
\end{array}
\right.
\]
avec deux axiomes

\begin{enumerate}
\item[a)] Associativité pour les compositions, pour $(a,b,c,d)\in\mathcal{M}_0^4$
\[
\begin{tikzcd}
\mathrm{Ch}(a,b)\times\mathrm{Ch}(b,c)\times\mathrm{Ch}(c,d) \arrow[r] \arrow[d] & \mathrm{Ch}(a,c)\times\mathrm{Ch}(c,d) \arrow[d]\\
\mathrm{Ch}(a,b)\times\mathrm{Ch}(b,d) \arrow[r] & \mathrm{Ch}(a,d).
\end{tikzcd}
\]
\item[b)] Existence d'objets unité bilatères
\[
1_a\in\pi_0(\mathrm{Ch}(a,a)) \qquad \text{pour tout } a\in\mathcal{M}_0
\]
\end{enumerate}

Alors les $\pi_0(\mathrm{Ch}(a,b))$ et leur composition définissent sur $\mathcal{M}_0$ comme ensemble d'objets une structure de \textit{catégorie} $\MOP$, ayant les $1_a$ comme identités, et on trouve que
\[
(a,b)\longmapsto\mathrm{Ch}(a,b)
\]
peut s'interpréter à partir d'un \textit{bifoncteur} canonique
\[
\mathcal{M}^o_{0\,\Pi}\times\MOP \longrightarrow \mathcal{H},
\]
dont le composé avec
\[
\mathcal{H}\xrightarrow{\pi_0}\mathrm{Ens}
\]
n'est autre que le bifoncteur $(a,b)\longrightarrow\mathrm{Hom}_{\MOP}(a,b)$,
\printedpage{6}
% head: A. Grothendieck : Dérivateurs, Ch. IX, Édition des Universités de Montpellier II et Paris VII
% foot: Version 12 mars 2002
\msp{page 8}
\[
\mathrm{Hom}_{\MOP}(a,b) \;\simeq\; \pi_0(\mathrm{Ch}(a,b)).
\]
Cette relation est la \textit{définition} de $\underline{\mathrm{Hom}}_{\MOP}$, mais ici la relation est entendue comme une % Reading resolved against PDF leaf 18: a single underline; see corrections.csv.
bijection \textit{fonctorielle} en $a$, $b$.

On trouve notamment que si
\[
u : a'\longrightarrow a, \qquad v : b\longrightarrow b'
\]
sont des flèches \textit{inversibles} dans $\MOP$, alors
\[
\mathrm{Ch}(a,b)\longrightarrow\mathrm{Ch}(a',b')
\]
correspondant dans $\mathcal{H}$ est \texttt{[un]} \textit{isomorphisme.}

On appliquera tout ceci surtout aux deux situations déduites d'une situation $(\mathcal{M},\Sigma)$, avec
\[
\begin{array}{rlcllcl}
1^o) & \mathcal{H} & = & \mathrm{Hot}, & \mathrm{Ch}(a,b) & = & \mathrm{Ch}^{\Sigma}_{\infty}(\mathcal{M};a,b)\\
2^o) & \mathcal{H} & = & \mathrm{Ind\,Cat\,hot}, & \mathrm{Ch}(a,b) & = & \mathrm{CH}^{\Sigma}_{\infty}(\mathcal{M};a,b)\ \text{(\footnotemark[2])}.
\end{array}
\]
\footnotetext[2]{$*$ dans Hot et Ind\,Cat\,hot est le produit cartésien sans plus.}
Dans les deux cas, on a donc un foncteur sur $(\mathcal{M}\Sigma^{-1})^o\times(\mathcal{M}\Sigma^{-1})$, à valeurs soit dans Hot, soit dans Ind\,Cat\,hot. (\textbf{NB} Les deux structures de catégorie sur $\mathcal{M}_0$ déduites des % [sic] des -> de, « déduites des ces deux foncteurs »
ces deux foncteurs sont canoniquement isomorphes.)
\[
\begin{tikzcd}[column sep=small]
 & \mathrm{Hot} & (a,b)\longmapsto\mathrm{Ch}^{\Sigma}_{\infty}(a,b)\\
(\mathcal{M}\Sigma^{-1})^o\times(\mathcal{M}\Sigma^{-1}) \arrow[ur] \arrow[dr] & & \\
 & \mathrm{Ind\,Cat\,hot} & (a,b)\longmapsto\mathrm{CH}^{\Sigma}_{\infty}(a,b)
\end{tikzcd}
\]

\msp{page 9}
Comme les foncteurs canoniques
\[
\begin{tikzcd}
\mathrm{Ind\,Cat\,hot} \arrow[r, "\mathrm{Ind}\,\varphi"] & \mathrm{Ind\,Hot} \arrow[r, "\mathrm{Ind}\,\psi"] & \mathrm{Ind\,Hot}_0^{\wedge} \arrow[r, "\varinjlim"] & \mathrm{Hot}_0^{\wedge},\\
\mathrm{Cat\,hot} \arrow[u, hook] \arrow[r, "\varphi"] & \mathrm{Hot} \arrow[u, hook] \arrow[r, "\psi"] & \mathrm{Hot}_0^{\wedge} \arrow[u, hook] \arrow[ur, "\mathrm{id}"'] &
\end{tikzcd}
\]
commutant aux produits finis, toute théorie de `Hom externes' à valeurs dans une de ces catégories, en définit une à valeurs dans chacune de celles dans laquelle elle s'envoie.
\printedpage{7}
% head: A. Grothendieck : Dérivateurs, Ch. IX, Édition des Universités de Montpellier II et Paris VII
% foot: Version 12 mars 2002
D'ailleurs, ces foncteurs ne sont pas tous pleinement fidèles (c'est en tous cas faux pour $\varphi$ et $\mathrm{Ind}\,\varphi$, et j'ignore ce qui en est de $\psi$ et $\mathrm{Ind}\,\psi$ -- les flèches verticales par contre sont pleinement fidèles (\footnote[3]{et également, $\mathrm{Ind\,Hot}_0^{\wedge}\xrightarrow{\varinjlim}\mathrm{Hot}_0^{\wedge}$ n'est \textit{pas} pleinement fidèle.}).) Mais ils induisent des bijections pour les $\pi_0$. Donc la catégorie $\MOP$ déduite d'une théorie de Hom externes à valeurs dans une de ces catégoires, et celle % [sic] catégoires -> catégories, faute d'orthographe
déduite dans une autre où elle s'envoie, est la même à isomorphisme canonique près. Pour ce qui est de cette définition de $\MOP$, rien n'est changé donc quand on regarde les $\mathrm{CH}_{\infty}$ comme des objets de Ind\,Hot, ou seulement de $\mathrm{Ind\,Hot}_0^{\wedge}$, ou même seulement de $\HotO$, ou quand on regarde les $\mathrm{Ch}_{\infty}$ comme des objets
\msp{page 10}
dans $\HotO$. Mais bien sûr, la donnée des accouplements de composition perd chaque fois en précision (sauf peut-être en passant de Hot à $\HotO$, si $\psi$ est pleinement fidèle \dots\,). Quand on prend les images des $\mathrm{CH}^{\Sigma}_{\infty}(a,b)$ et $\mathrm{Ch}^{\Sigma}_{\infty}(a,b)$ dans $\HotO$ (la dernière catégorie du diagramme), on trouve des objets canoniquement isomorphes (fonctoriellement en $a$, $b$ dans $\MOP\simeq\mathcal{M}\Sigma^{-1}$). Quand on les regarde comme étant respectivement dans Ind\,Hot et dans $\mathrm{Hot}\hookrightarrow\mathrm{Ind\,Hot}$, on trouve un homomorphisme fonctoriel
\[
\mathrm{CH}^{\Sigma}_{\infty}(a,b)\longrightarrow\mathrm{Ch}^{\Sigma}_{\infty}(a,b) \qquad\quad \text{dans Ind\,Hot ou dans } \mathrm{Ind\,Hot}_0^{\wedge} \text{ au choix.} \tag*{$(*)$}
\]
La $\varinjlim$ de $\mathrm{CH}^{\Sigma}_{\infty}(a,b)$ dans Hot n'existe pas en général (sauf si l'ind-objet est essentiellement constant), mais elle existe dans $\HotO$, et l'isomorphisme canonique
\[
\varinjlim_{\HotO}\mathrm{CH}^{\Sigma}_{\infty}(a,b)\xrightarrow{\ \sim\ }\mathrm{Ch}^{\Sigma}_{\infty}(a,b) \qquad\quad \text{dans } \HotO \tag*{$(**)$}
\]
est la flèche déduite de $(*)$ dans $\mathrm{Ind\,Hot}_0^{\wedge}$.

\dgrblock{dgrb429}{What the finite tests forget}
\begin{dcomment}{status}{p9:D'ailleurs, ces foncteurs ne sont pas tous pleinement fidèles}{What the finite tests forget}
Grothendieck is trying to recover a space of morphisms from its values on a manageable collection of diagram shapes. The attraction is clear: if those values determined both spaces and maps between them, a difficult coherence argument could be replaced by calculations on finite diagrams. The notation needs its earlier definition. In Chapter V, manuscript pages 51--52, $\Hot_0$ is the full subcategory of homotopy types represented by categories with finitely many objects and arrows \citep[Chapter V, printed pp.~31--32]{GrothendieckDerivateursV}. It is not the subcategory of finite CW complexes. For example, the one-object category of a finite nontrivial group is finite, whereas its nerve need not have finite dimension.

Choose a small skeleton $\mathcal A$ of $\Hot_0$. The functor in question is
\[
 \psi(X)(K)=[K,X],\qquad K\in\mathcal A.
\]
It preserves products, and evaluation at the one-point category recovers $\pi_0X$. These useful facts survive the negative answer to the full-faithfulness question. The commentary on Chapter~V, manuscript p.~52, discusses the obstruction (``Finite unbased tests miss a symmetric-group shift'' and ``The modern scope of the final obstruction''). The general theorem is that no set of unpointed homotopy types jointly detects isomorphisms in $\Hot$ \citep[Theorem~2.1]{ArlinChristensen2023}. Applied to $\mathcal A$, this says that $\psi$ is not even conservative, and hence is not fully faithful.

There is also a separate answer to the faithfulness question, which is what the hoped-for associativity argument actually needs. The presheaf category on $\mathcal A$ has a faithful set-valued functor
\[
 F\longmapsto\coprod_{K\in\operatorname{Ob}\mathcal A}F(K):
\]
a natural transformation is determined by its components, which this disjoint union retains. If $\psi$ were faithful, the composite would make $\Hot$ concrete. Freyd's nonconcreteness theorem rules this out; his argument already applies to a full subcategory of simply connected, finite-dimensional CW complexes, so forgetting basepoints does not evade the obstruction \citep[pp.~1 and 5 of the 2004 reprint]{Freyd1970}. Thus $\psi$ is not faithful either. These are statements with the usual fixed-universe interpretation of small spaces and small sets; restricting the allowed spaces to a particular small category changes the question.

Consequently, agreement after passage to $\Hot_0^{\wedge}$ cannot in general prove equality of the two composites in $\Hot$. It follows as well, by restriction to constant ind-objects, that $\operatorname{Ind}\psi$ is not fully faithful. None of this invalidates the underlying category of morphisms obtained by taking $\pi_0$; it identifies exactly the additional coherence that manuscript page 13 still has to supply.
\end{dcomment}


\dgrblock{dgrb430}{Three meanings of a limiting space}
\begin{dcomment}{examples}{p10:dans Hot n'existe pas en général}{Three meanings of a limiting space}
Grothendieck keeps an ind-object because it records the stage at which a family of paths can be represented. Passing to its union, or to a colimit of its homotopy types, asks a different question. The phrase ``n'existe pas en général'' should be read as the absence of a general existence theorem here. Essential constancy is a sufficient condition, not a necessary condition for an ordinary colimit in $\Hot$.

For a complete elementary example, let $E_n=\{0,\ldots,n\}$, with the discrete topology and the evident inclusions, and let $E=\mathbf N$. For every space $Y$ there are natural bijections
\[
 [E,Y]=(\pi_0Y)^{\mathbf N}
       \simeq\varprojlim_n(\pi_0Y)^{E_n}
       =\varprojlim_n[E_n,Y].
\]
Hence $E$ is the ordinary colimit in $\Hot$. Nevertheless, the ind-object $(E_n)$ is not essentially constant. If it were isomorphic to a constant object $Z$, its functor of maps to constant objects would imply $Z\simeq E$. An inverse morphism from the constant $E$ to $(E_n)$ factors through some $E_n$, by the defining formula for maps in an ind-category. Its composite with the cocone would make the identity of $\mathbf N$ factor through a finite set on $\pi_0$, which is impossible.

The same example explains the footnote about the colimit functor on presheaves. Let $F_n$ be the constant presheaf with value $E_n$ on $\mathcal A$, and $F$ the constant presheaf with value $\mathbf N$. Then $\operatorname{colim}F_n=F$. The identity of $F$ has no lift to a morphism from the constant ind-object $F$ to $(F_n)$: such a lift would factor through one finite $F_n$, again contradicting evaluation at any object of $\mathcal A$. Thus
\[
 \operatorname{Ind}(\widehat{\mathcal A})
       \longrightarrow\widehat{\mathcal A},\qquad
 (F_n)\longmapsto\operatorname{colim}F_n,
\]
is not full.

A further qualification is needed for the displayed comparison with $\Ch_\infty$. A filtered homotopy colimit $K=\operatorname{hocolim}K_n$ is tested by a \emph{compact} homotopy type $A$ through
\[
 \operatorname{colim}_n[A,K_n]\simeq[A,K].
\]
Finite-category homotopy types are not all compact. Take $A=BC_2=\mathbf{RP}^{\infty}$ and $K_n=\mathbf{RP}^{n}$. Their union is $A$, but the identity of $A$ cannot factor, even up to homotopy, through any $K_n$: it would kill the nonzero mod-$2$ cohomology in degrees above $n$ \citep[Theorem~3.19]{Hatcher2002}. Thus the finite-category definition of $\Hot_0$ alone does not justify the comparison in manuscript page 10. For the particular path system, one must either prove the required factorization on these test categories by its special categorical structure, or impose an appropriate finite-test/compactness hypothesis. We do not replace that missing model-specific argument by a general assertion about filtered homotopy colimits.
\end{dcomment}


\msp{page 11}
Dans les deux cas CH et Ch (qu'on va désigner par le symbole commun $C_{\infty}(a,b)$), la fonctorialité de cet objet $C_{\infty}(a,b)$ (même au sens le plus faible, comme objet de $\HotO$) par rapport à $a$ et $b$ (même au sens le plus faible, comme foncteur $\mathcal{M}^o\times\mathcal{M}\longrightarrow\HotO$ et pas même $(\mathcal{M}\Sigma^{-1})^o\times(\mathcal{M}\Sigma^{-1})\longrightarrow\HotO$) n'est pas évidente du tout sur la définition directe en termes du système inductif explicite dans Cat
\[
C_{\infty}(a,b) \;=\; (\mathrm{Ch}_{\tau_0^{2n}}(a,b))_{n\in\mathbf{N}},
\]
car une flèche $a'\longrightarrow a$ ou $b\longrightarrow b'$ dans $\mathcal{M}$, et même qu'elle serait dans $\Sigma$, ne définit aucunement une flèche dans Cat
\[
\begin{array}{rcl}
C_{\infty}(a,b) & \longrightarrow & C_{\infty}(a',b)\\
\text{ou}\ \ C_{\infty}(a,b) & \longrightarrow & C_{\infty}(a,b').
\end{array}
\]
Cette difficulté rappelle celle qu'avait à surmonter \textsc{Quillen}, dans Homotopical Algebra, pour affronter la différence entre `left homotopies' et `right homotopies'. Je me rappelle comment
\printedpage{8}
% head: A. Grothendieck : Dérivateurs, Ch. IX, Édition des Universités de Montpellier II et Paris VII
% foot: Version 12 mars 2002
\msp{page 12}
on peut néanmoins définir de telles flèches, comme résultant d'une flèche (beaucoup plus riche) de `composition' de chemins entre $a'$ et $a$, et $a$ et $b$, ou entre $a$ et $b$, et $b$ et $b'$, soit en général donc l'accouplement
\[
C_{\infty}(a,b)\times C_{\infty}(b,c)\longrightarrow C_{\infty}(a,c),
\]
dont la définition directe n'est pas moins problématique, et ne peut en tous cas \textit{pas} se donner comme une flèche dans Cat lui-même, ou dans Ind\,Cat. Interprétant p.~ex. $\mathrm{Ch}_{\infty}(a,b)$ en termes de chemins \textit{infinis} (des deux cotés) de $a$ à $b$, en un sens évident (mais % [sic] cotés -> côtés, accent omis
`essentiellement finis' des deux cotés), on peut définir aussi % [sic] cotés -> côtés, accent omis
\[
\mathrm{Ch}^{-}_{\infty}(a,b), \qquad \mathrm{Ch}^{+}_{\infty}(a,b)
\]
comme les sous-catégories pleines formées respectivement de chemins qui sont totalement dégénérés en degrés $\geq 0$ (i.e. $a_0\longleftarrow a_1\longrightarrow a_2\longleftarrow a_3\cdots$ sont tous des identités), ce qui revient aux chemins infinis seulement à gauche, i.e. s'arrêtant à $a_0$, et symétriquement pour $\mathrm{Ch}^{+}_{\infty}(a,b)$. Même tabac pour les $\mathrm{CH}^{+}$, $\mathrm{CH}^{-}$.

\msp{page 13}
Ceci dit, on a ceci (revenant à la notion $C_{\infty}$ ambivalente) :

\begin{enumerate}
\item[a)] les inclusions pleinement fidèles
\[
C^{-}_{\infty}(a,b)\hookrightarrow C_{\infty}(a,b)\hookleftarrow C^{+}_{\infty}(a,b)
\]
sont des isomorphismes dans Ind\,Cat\,hot resp. dans Hot, et
\item[b)] On a bel et bien une composition des chemins
\[
C^{-}_{\infty}(a,b)\times C^{+}_{\infty}(b,c)\longrightarrow C_{\infty}(a,c).
\]
\end{enumerate}

L'\textit{associativité} dans le cas du CH n'offre d'ailleurs pas de difficulté, à cause de la possibilité de se ramener aux
\[
C_m(a,b)\times C_n(b,c)\times C_p(c,d)\rightrightarrows C_N(a,d)
\]
pour $N$ grand, et prouver l'egalité comme flèches dans Cat\,hot -- ce qui résultera du lemme % [sic] egalité -> égalité, accent omis
d'homotopie sur les intervalles dans Cat. Dans le cas du $\mathrm{Ch}_{\infty}$ par contre je n'ai pas trouvé de démonstration. (Faute seulement d'avoir essayé sérieusement de m'y mettre? Je n'ai pas trop envie, vu qu'il se pourrait que $\mathrm{Hot}\longrightarrow\HotO$ soit pleinement fidèle, auquel cas le travail devient inutile \dots\,)

\dgrblock{dgrb431}{The path model and the localization it is meant to represent}
\begin{dcomment}{later}{p13:dans le cas du CH n'offre}{The path model and the localization it is meant to represent}
Grothendieck is trying to make concatenation independent of an arbitrary placement of the middle endpoint. Moving one path to the left and the other to the right produces an actual functor; the homotopy equivalences from these one-sided path categories then give composition in $\Hot$. For the ind-object $\mathrm{CH}_\infty$, every finite collection of paths and reparametrizations can be handled at one larger stage. His interval homotopies therefore give the asserted associativity in $\mathrm{Ind\,Cat\,hot}$. For the realized $\Ch_\infty$, the text explicitly leaves the associativity proof unfinished. An eventual equality at every finite stage is not, by itself, a specified compatible homotopy on the whole realization.

The commentary on Chapter~V, manuscript p.~12 (``What composition in an ind-category actually asserts'') explains the relevant ind-Hom formula. Here it gives the exact associativity argument. Put $\mathcal A=\mathrm{Cat\,hot}$ and write $C_m(a,b)$ for the finite stages. Morphisms from the triple product into $\mathrm{CH}_\infty(a,d)$ are computed by
\[
 \varprojlim_{m,n,p}\ \varinjlim_N
 \operatorname{Hom}_{\mathcal A}\bigl(
 C_m(a,b)\times C_n(b,c)\times C_p(c,d),C_N(a,d)\bigr).
\]
For each fixed $(m,n,p)$, concatenate the three finite words before returning to the chosen standard shape. Parenthesization does not change this concatenated word; it changes only the finite reparametrization and the inserted identity padding. At a common larger stage the interval-homotopy lemma identifies those reparametrizations; its relative endpoint condition is the one explained in ``Even indices and reparametrization relative to the ends'' and ``Sliding an even block gives a natural transformation'' in the commentary on Chapter~V, manuscript pp.~14 and~21. Hence the two representatives have the same class in the displayed filtered colimit for every $(m,n,p)$, which proves equality in the inverse limit. No choice of these finite homotopies compatible as $(m,n,p)$ varies is required for this equality of ind-morphisms. Concatenating an identity path similarly adds only identity padding, so the same argument proves the strong unit equations for $\mathrm{CH}_\infty$. This is the precise advantage of the ind-object in the present argument.

The general target of this construction was already available in a different model. The hammock localization of \citet[\S2; Propositions 2.2 and 3.1]{DwyerKan1980Calculating} is a simplicial category whose composition concatenates and reduces hammocks; its category of components is $\mathcal M[W^{-1}]$. Thus it supplies coherent mapping spaces, with strictly associative composition at that model level, for a relative category. These statements concern the hammock construction. They do not identify its simplicial sets with Grothendieck's fixed-type path categories, whose transformations and stabilization conventions must be compared separately.

This distinction continues the commentary on Chapter~V, manuscript p.~27 (``From components of path categories to mapping spaces''). Equality of $\pi_0$, agreement on finite tests, and a weak equivalence of the mapping spaces are three different assertions. We have not located in the sources checked a comparison theorem stated for exactly the bilateral stabilization and unrestricted natural transformations used here. The model-specific comparison, including compatibility with concatenation, is consequently left identified as an unfinished verification in these pages, not advertised as a new open problem in homotopy theory. The abstract bifunctor construction above is complete once its stated associativity and unit hypotheses hold.
\end{dcomment}


% ---- from Der-ChIX-678_p011-015.tex
\printedpage{9}
% head: A. Grothendieck : Dérivateurs, Ch. IX, Édition des Universités de Montpellier II et Paris VII
% foot: Version 12 mars 2002

\msp{page 14}

\section*{2\quad Comparaison avec Quillen}

Il se pose un certain nombre de questions sur ce foncteur $(\M\Sigma^{-1})^o \times (\M\Sigma^{-1}) \longrightarrow \Hot$, et sa relation aux constructions classiques, notamment à la théorie de \textsc{Quillen}. Je vais maintenant noter
\[
\begin{aligned}
\Ch^{\Sigma}_{\infty}(a,b) \;&=\; \Homext_{\Hot}(a,b) && \text{ou simplement}\\
&\phantom{\;=\;} \Homext(a,b).
\end{aligned}
\]

1$^o$). Soient $f,g : a \rightrightarrows b$ dans $\M$, alors (si $\M$ est une `model category') \textsc{Quilllen} définit un ensemble % [sic] Quilllen -> Quillen, three l's printed
\[
\pi_1(a,b;f,g) \qquad \text{ou} \qquad \pi_1(f,g)
\]
des classes d'homotopie de $f$ à $g$, de telle façon que $\mathrm{Hom}_{\M}(a,b)$ devient l'ensemble des objets d'un groupoïde, dont les flèches sont les éléments de $\pi_1(f,g)$, pour une notion de composition naturelle. Je suspecte que l'on a
\[
\pi_1(f,g) \;=\; \mathrm{Hom}_{\underbrace{\Pi(\uCh^{\Sigma}_{\infty}(a,b))}_{\substack{\text{groupoïde}\\ \text{fodamental}\\ \text{de } \uCh^{\Sigma}_{\infty}(a,b)}}}(f,g), % [sic] fodamental -> fondamental, missing n
\]
quand $f$ et $g$ sont considérés comme des `chemins infinis' (en insérant chaque $f : a \longrightarrow b$ dans
\[
\cdots \underbrace{a_{-2}}_{=a} \xleftarrow{\ \mathrm{id}\ } \underbrace{a_{-1}}_{=a} \xrightarrow{\ f\ } \underbrace{a_{0}}_{=b} \xleftarrow{\ \mathrm{id}\ } \underbrace{a_{1}}_{=b} \longrightarrow \cdots\ \ )
\]

\dgrblock{dgrb432}{Homotopies as arrows in a mapping groupoid}
\begin{dcomment}{exposition}{p14:Je suspecte que l'on a}{Homotopies as arrows in a mapping groupoid}
Grothendieck's question asks for the first layer of the mapping space, where homotopies themselves become morphisms. Its modern formulation begins with a derived mapping space $\operatorname{Map}^{h}(a,b)$, and takes its fundamental groupoid. The objects corresponding to actual maps $a\to b$ form a full subgroupoid. They need not meet every component until the usual cofibrancy and fibrancy hypotheses, or replacements, have been imposed.

For a cofibrant object $a$ and a fibrant object $b$ in a simplicial model category, $\operatorname{Map}(a,b)$ is a Kan complex. A simplicial path from $f$ to $g$ represents a homotopy; two such paths give the same arrow of its fundamental groupoid precisely when they are homotopic relative to their endpoints. Concatenation in the groupoid is defined by filling an inner two-dimensional horn. Different fillers give the same class because the Kan conditions provide the requisite relative homotopy. Constant paths give identities and reversed paths give inverses. Consequently
\[
 \operatorname{Hom}_{\Pi\operatorname{Map}(a,b)}(f,g)
\]
is the exact object Grothendieck wants, not the set of equivalence classes of maps $a\to b$.

Quillen's original construction has precisely the cofibrant-source and fibrant-target qualifications: see Chapter I, \S2, Proposition 1.2.1 in the checked retypesetting; Chapter II, \S2 identifies this groupoid with the fundamental groupoid of the simplicial mapping complex \citep{Quillen1967}. For a general model category, \citet[Propositions 4.4--4.5 and 4.8]{DwyerKan1980Function} compare resolutions with the hammock mapping spaces and with the simplicial enrichment when it exists. These results justify the same description using $\operatorname{Map}^{h}$ after replacements. The identification with $\Pi\Ch^W_\infty(a,b)$ still requires the particular path-model comparison just discussed. The source's conjectural equality should therefore be stated as a natural equivalence of groupoids after that comparison, not as a literal equality of two chosen constructions.
\end{dcomment}


\msp{page 15}

Donc la définition des $\pi_1(f,g)$ (si c'est bien ça) se prolonge de façon naturelle aux chemins infinis (essentiellement finis) arbitraires.

2$^o$) Supposons que $\M$ admette un objet à la fois initial et final $0_{\M}$ ($\M$ `ponctuée'). Je dis qu'alors toutes les catégories $\uCh^{\Sigma}_{\infty}(a,b)$ sont ponctuées de façon naturelle : à savoir, par le chemin infini de $a$ à $b$ défini (comme plus haut) à partir de $f = 0_{a,b} : a \longrightarrow b$, l'application nulle de $a$ à $b$. On a envie alors de regarder les $\uCh^{\Sigma}_{\infty}(a,b)$ comme définissant des éléments dans le $\Hot_{\bullet}$ (`Hot ponctué'). C'est plus fin que de dire qu'on se donne une classe neutre
\[
e \;\longrightarrow\; \Ch(a,b) \qquad \text{dans Hot,}
\]
car le foncteur de $\Hot_{\bullet}$ dans Hot (ou dans $e\backslash\Hot$) n'est pas fidèle. Le problème, c'est de définir les flèches
\[
\Ch(a,b)\times\Ch(b,c) \;\longrightarrow\; \Ch(a,c)
\]
dans $\Hot_{\bullet}$, et pas seulement dans Hot, et avec associativité (et les unités) dans $\Hot_{\bullet}$ également. Si on applique la définition avec $\Ch^{-}(a,b)$ et $\Ch^{+}(b,c)$, le composé des éléments nuls $a \xrightarrow{\ 0\ } b$ de $\Ch^{-}(a,b)$ et $b \xleftarrow{\ \mathrm{id}\ } b \xrightarrow{\ 0\ } c$ de $\Ch^{+}(b,c)$
\printedpage{10}
% head: A. Grothendieck : Dérivateurs, Ch. IX, Édition des Universités de Montpellier II et Paris VII
% foot: Version 12 mars 2002
\msp{page 16}
est l'élément de $\Ch_{\infty}(a,c)$ que voici
\[
\begin{tikzcd}[row sep=huge]
 & a_{-2} & a_{-1} \arrow[l] \arrow[r] & a_{0} \arrow[d, equal] & a_{1} \arrow[l] \arrow[r] & a_{2} & \cdots \\
{}_{\mathrm{id}}\ \cdots & a & a \arrow[l, "\mathrm{id}"'] \arrow[r, "0"'] & b & b \arrow[l, "\mathrm{id}"'] \arrow[r, "0"'] & c & \cdots_{\mathrm{id}} \arrow[l, "\mathrm{id}"']
\end{tikzcd}
\]
ce n'est \textit{pas} l'élément nul de $\Ch_{\infty}(a,c)$ ! Mais on a des flèches canoniques
\[
e^{-}_{a,b}\circ e^{+}_{b,c} \;\;\rightleftarrows\;\; e_{a,c} \qquad \text{dans } \Ch_{\infty}(a,c)
\]
qui sont l'identité sur tous les sommets sauf celui d'indice 0. Il y a un sorite à développer, comme quoi une paire
\[
(f : K_{\bullet} \longrightarrow L_{\bullet}, c)
\]
d'un foncteur entre catégories ponctuées, et d'un chemin $c : f(e_{K_{\bullet}}) \longrightarrow e_{L_{\bullet}}$, définit dans Hot une flèche canonique
\[
\mathrm{hot}_{\bullet}(K_{\bullet}) \;\longrightarrow\; \mathrm{hot}_{\bullet}(L_{\bullet})
\]
avec transitivité si on a
\[
(f' : L_{\bullet} \longrightarrow M_{\bullet}, c' : f'(e_{L_{\bullet}}) \longrightarrow e_{M_{\bullet}}),
\]

\dgrblock{dgrb433}{The path needed to preserve a basepoint}
\begin{dcomment}{exposition}{p16:Il y a un sorite à développer}{The path needed to preserve a basepoint}
Grothendieck wants the zero map to remain a specified point throughout composition, because its loop group will define suspension and loops. A point chosen in the unpointed homotopy category does not contain the required information. For example, based maps $S^1\to BG$ are classified by elements of a group $G$, whereas free homotopy identifies conjugate elements. For $G=S_3$, the two distinct transpositions $(12)$ and $(13)$ give different based classes and the same free class. This proves the asserted failure of faithfulness of the forgetful functor.

The proposed extra datum is exactly the right one at the level of spaces. For pointed CW spaces $(K,k)$ and $(L,l)$, a map $f:K\to L$ together with a path $c$ from $f(k)$ to $l$ defines a morphism in the pointed homotopy category. To see this without suppressing the path, attach an interval to $K$ at $k$ and point the resulting space $K^+$ at its free endpoint. Define $F:K^+\to L$ by $f$ on $K$ and by $c$ on the interval. The collapse $r:K^+\to K$ is a pointed homotopy equivalence. The required morphism is $[F][r]^{-1}$. The same construction applies to nerves of pointed categories after geometric realization, because a vertex inclusion is a cofibration.

For a second pair $(g,d)$, the composite is represented by $(gf,(gc)*d)$. Indeed, the space obtained by attaching two successive intervals to $K$ maps to the target by $f$, then the first path followed under $g$, then $d$. Collapsing the subdivision of the two intervals identifies this representative with the single attached interval carrying $(gc)*d$. Repeating this argument gives associativity: the increasing piecewise-linear reparametrizations of a concatenated interval are homotopic relative to the two endpoints. Homotopic choices of $c$ relative endpoints give the same pointed morphism. Arbitrary paths with the same endpoints need not do so; their difference can act through the fundamental group.

The specific claim immediately before this construction is stronger and is false as written, but it has a simple canonical repair. Denote the displayed concatenation by $P$ and the standard zero path from $a$ to $c$ by $Q$. Then $P$ has $b$ at vertices $0,1$, while $Q$ has $c$ at both. A transformation cannot be the identity at every vertex except $0$: the backward identity edge between $1$ and $0$ forces its components at those two vertices to agree. Instead define $\alpha:P\to Q$ by
\[
 \alpha_0=\alpha_1=0_{b,c},\qquad
 \alpha_i=\mathrm{id}\quad(i\ne0,1).
\]
The square at $-1\to0$ commutes because a composite with a zero map is zero; the square at $1\to0$ compares two identity edges; and the square at $1\to2$ says $\mathrm{id}_c\,0_{b,c}=0_{b,c}$. Every other square is an identity square. Since transformations of paths are unrestricted, this is an admissible canonical comparison for every $W$ containing identities. It supplies the needed basepoint path without a reparametrization or an additional choice. A reverse arrow with zero components would force $\mathrm{id}_c=0_{c,c}$ at the last square and is not available in general. This repairs the single basepoint comparison; associativity and the other coherence conditions for the particular pointed composition still require verification.
\end{dcomment}


\msp{page 17}

d'où
\[
f'' = f'\circ f : K_{\bullet} \longrightarrow M_{\bullet}, \qquad \text{et}\ \ c'' : f''(e_{K_{\bullet}}) \longrightarrow e_{M_{\bullet}}), % [sic] closing parenthesis without opening parenthesis -> (f'' = f'\circ f : ..., c'' : ...), unbalanced parenthesis
\]
définit ainsi
\[
\begin{tikzcd}
f'(f(e_{K_{\bullet}})) \arrow[r, "f'(c)"] \arrow[rr, bend right=12, "{c''\ \text{par définition}}"'] & f'(e_{L_{\bullet}}) \arrow[r, "c'"] & e_{M_{\bullet}}
\end{tikzcd}
\]

\noindent\llap{?\quad}Cela nous donne donc bel et bien des accouplements dans $\Hot_{\bullet}$ (et aussi dans $\mathrm{Cat\,hot}_{\bullet}$?). Il est difficile de croire qu'ils pourraient ne pas être associatifs! D'ailleurs $\Hot_{\bullet}$, avec le produit cartésien ordinaire (\footnote[4]{Non, il faut prendre le produit $\xi_{\bullet} \divideontimes \eta_{\bullet}$ ($\xi_{\bullet}\times\eta_{\bullet}$, avec $\xi_{\bullet}\vee\eta_{\bullet}$ contracté en $e$).}), satisfait à l'axiomatique de tantôt pour $\mathcal{H}$, et $\Hot_{\bullet} \longrightarrow \Hot$ est compatible avec ce produit (\footnote[5]{Non, mais on s'en tire quand même \ldots}), et avec $\pi_0$. % Reading resolved against PDF leaf 22: \divideontimes, meaning smash; see corrections.csv.

\dgrblock{dgrb434}{Why the monoidal unit changes in the pointed theory}
\begin{dcomment}{exposition}{p17:Il est difficile de croire}{Why the monoidal unit changes in the pointed theory}
The two footnotes identify a genuine change of monoidal structure. In pointed spaces the Cartesian unit is the one-point space, and
\[
 \operatorname{Hom}_{\Hot_\bullet}(*,X)=\{*\}.
\]
It cannot recover the components of $X$, so the Cartesian structure does not instantiate the earlier convention $\pi_0(X)=\operatorname{Hom}(I,X)$. With the smash product, the unit is $S^0$, and
\[
 \operatorname{Hom}_{\Hot_\bullet}(S^0,X)=\pi_0(X).
\]
Thus the corrected enrichment uses smash composition
\[
 C(a,b)\wedge C(b,c)\longrightarrow C(a,c).
\]
The identity is the pointed map $S^0\to C(a,a)$ carrying the nonbasepoint to $1_a$ and the basepoint to zero. Composition must kill either zero input, not merely the pair of zero inputs. That is the additional requirement hidden by a single basepoint comparison for a Cartesian product. The standard pointed simplicial model of derived mapping spaces has this absorbing-zero property, hence its composition factors through the smash product. This is the familiar monoidal example in \citet[\S1.1]{Kelly1982}.

Forgetting basepoints is not strong monoidal for smash and Cartesian products. Nevertheless there is a natural quotient map
\[
 U(X)\times U(Y)\longrightarrow U(X\wedge Y)
\]
and a unit map $*\to U(S^0)$ selecting its nonbasepoint. These make $U$ lax monoidal: both associativity composites collapse exactly those tuples with a zero coordinate, and both unit composites leave the other coordinate unchanged. Composing these maps with smash composition recovers associative unpointed composition. The footnote ``on s'en tire quand m\^eme'' can therefore be made precise without asserting preservation of products. The symbol in the edition is the crossed product $\divideontimes$, explained there as collapsing the wedge, and is faithfully transcribed; its meaning here is the smash product.
\end{dcomment}


Si tout ceci marche, on aurait donc une théorie des $\Homext_{\bullet}$ à valeurs dans $\Hot_{\bullet}$; et sans doute aussi pour les $\mathrm{CH}_{\infty}$ comme objets de $\mathrm{Ind\,Hot}_{\bullet}$, et non seulement de $\mathrm{Ind\,Hot}$. Mais
\printedpage{11}
% head: A. Grothendieck : Dérivateurs, Ch. IX, Édition des Universités de Montpellier II et Paris VII
% foot: Version 12 mars 2002
de toutes façons on ne peut en général laisser tomber Cat hot, qui reste importante -- donc il faudrait avoir des objets $\mathrm{CH}_{\infty}(a,b)$ dans $\mathrm{Ind\,Cat\,hot}_{\bullet}$, pour avoir une côte bien taillée.

Les $\pi_0(C(a,b))$ seront à présent des ensembles ponctués, et ce sera par l'image de l'application nulle $a \xrightarrow{\ 0\ } b$, i.e. le composé
\[
a \longrightarrow 0 \longrightarrow b
\]

\msp{page 18}

qu'il faut composer avec l'elément déduit du chemin tronqué % [sic] l'elément -> l'élément, missing accent
\[
a \longrightarrow \underbrace{0}_{=a_0} \longleftarrow b \qquad \ldots
\]

Il faudrait montrer que 0 est aussi un objet neutre de $\M\Sigma^{-1}$, i.e. que
\[
\pi_0(a,0) = \pi_0(0,a) = \{0\} \qquad \text{pour tout } a \in \mathrm{Ob}\,\M
\]
et interpréter l'élément nul dans $\pi_0(a,b)$ comme l'application nulle dans $\M\Sigma^{-1}$. De façon générale, il faut des conditions qui assurent (sans que $\M$ soit forcément ponctuée) que
\[
\M \xrightarrow{\ \gamma\ } \M\Sigma^{-1}
\]
commute aux sommes finies, aux produits finis (donc $\gamma(\varphi_{\M}) = \varphi_{\M\Sigma^{-1}}$, resp. $\gamma(e_{\M}) = e_{\M\Sigma^{-1}}$) (\footnote[6]{à faire.}).

\dgrblock{dgrb435}{Zero objects and finite sums after localization}
\begin{dcomment}{exposition}{p18:Il faudrait montrer que 0 est aussi}{Zero objects and finite sums after localization}
Grothendieck needs the distinguished component to be an actual zero morphism of the localized category. At the level of ordinary localization this requires no additional compatibility condition on $W$. If $z$ is initial in $\mathcal M$, write $i_x:z\to x$ for its unique maps. For every ordinary arrow $u:x\to y$ we have $ui_x=i_y$. If $u$ is inverted, the same equality gives $u^{-1}i_y=i_x$ in the localization. An induction along any finite zigzag starting at $z$ therefore identifies it with $i_x$. Thus $z$ remains initial. The reversed argument preserves terminal objects. In particular a zero object remains zero, and $a\to0\to b$ is the zero morphism in $\mathcal M[W^{-1}]$. This is also the nullary case of descent of adjunctions through localization \citep[Lemma 1.1.6]{Krause2021}.

The displayed path $a\to0\leftarrow b$ needs a different caution: its backward edge belongs to $W$ only when $b\to0$ does. The two forward zero maps $a\to0\to b$ always give the required ordinary zero morphism; inserting backward identity edges puts this composite into the permitted zigzag shapes. The preservation of the zero object therefore does not depend on the inadmissible shortened path or on the unresolved pointed composition.

For binary products there is a useful sufficient condition. Suppose $\mathcal M$ has finite products and $W$ contains identities and is closed under their finite products. The product functor and the diagonal then descend to the localizations, as do the unit and counit of $\Delta\dashv\prod$. Naturality for inverse generators follows by multiplying the original naturality squares by the inverses of their $W$-edges; the triangle identities remain equalities after localization. Finally
\[
 (\mathcal M^n)[(W^n)^{-1}]
 \simeq (\mathcal M[W^{-1}])^n
 \qquad(n<\infty):
\]
a tuple of finite zigzags is represented by one tuple-shaped zigzag after padding the finitely many words by identities, and the relations are coordinatewise. Consequently the descended diagonal still has the descended product as right adjoint. This proves that localization preserves finite products. Reversing arrows gives the corresponding statement for finite coproducts. One can weaken closure by requiring only that the relevant products or coproducts of $W$-arrows become invertible.

For an additive category, closure of $W$ under finite direct sums is already enough for the ordinary localized category to be additive; a five-lemma hypothesis is unnecessary for this conclusion \citep[Lemma 2.2.1]{Krause2021}. Indeed, products and coproducts survive and the original canonical comparison from their sum to their product remains an isomorphism, giving biproducts. Addition is therefore defined by the usual diagonal--biproduct--codiagonal formula. The morphism $-1_x$ remains an inverse for the addition on $x$, and its naturality extends over inverted arrows, so every Hom monoid is a group. A central $k$-action similarly extends. These conclusions concern ordinary Hom groups; constructing the proposed complexes of higher Hom information is the further problem on the next page.
\end{dcomment}


Supposons enfin que $\M$ soit \textit{additive,} et même $k$-linéaire ($k$ un anneau fixé), i.e. on a de plus un homomorphisme d'anneaux
\[
k \;\longrightarrow\; \mathrm{End}(\mathrm{id}_{\M}).
\]

On doit supposer quelque compatibilité de $\Sigma$ à la structure additive, p. ex. que la somme directe de deux flèches dans $\Sigma$ est dans $\Sigma$, et sans doute quelque chose de plus fort, du genre `lemme des cinq' pour $\Sigma$. (Mais on ne suppose pas forcément $\M$ abélienne, donc on aurait dû mal à formuler un lemme des cinq?) % [sic] dû mal -> du mal, accent on du

\msp{page 19}

Mais peut-être ici faut-il un axiome plus subtil, p.ex. (si $\underbrace{\Sigma}_{\text{suspension}}$ et $\Omega$ existent (\footnote[7]{cf. plus bas ce que ça signifie.})) que l'ensemble des morphismes d'adjonction $\Sigma\Omega \longrightarrow \mathrm{id}$, $\mathrm{id} \longrightarrow \Omega\Sigma$ soit \texttt{[formé d']}isomorphismes, ce qui implique que $\Omega$ resp. $\Sigma$ est pleinement fidèle. En tous cas, l'attente naturelle ici, c'est que $\Ch^{\Sigma}_{\infty}(a,b)$ soit non seulement dans $\Hot_{\bullet}$, mais même dans la catégorie $D^{\leq 0}(k\text{-Mod})$, sous-catégorie pleine de la catégorie dérivée $D^{-}(k\text{-Mod})$ formée des $K$ tels que $H^{i}(K) = 0$ si $i > 0$. On peut aussi la noter $D_{\geq 0}(k\text{-Mod})$, l'interprétant en
\printedpage{12}
% head: A. Grothendieck : Dérivateurs, Ch. IX, Édition des Universités de Montpellier II et Paris VII
% foot: Version 12 mars 2002
termes de complexes de `chaînes' plutôt que de cochaînes. Ici les accouplements doivent correspondre à l'opération $* = \Lotimes_k$ dans $D_{\geq 0}(k\text{-Mod})$, admettant $k$ comme objet unité, et
\[
\pi_0(K_{\bullet}) \;\simeq\; H_0(K_{\bullet}).
\]
Donc on doit avoir
\[
\Homext(a,b) \Lotimes_k \Homext(b,c) \;\longrightarrow\; \Homext(a,c),
\]
où cette fois Hom ext désigne l'objet de $D_{\geq 0}(k\text{-Mod})$, au lieu d'un objet de Hot seulement, ou de $\Hot_{\bullet}$. On utilise les foncteurs
\[
D_{\geq 0}(k\text{-Mod}) \longrightarrow \Hot_{\bullet} \longrightarrow \Hot
\]
\msp{page 20}
pour faire le lien avec la théorie ponctuée dans $\Hot_{\bullet}$, ou \texttt{[la]} théorie générale dans Hot. Si on désigne par
\[
\gamma :\ D_{\geq 0}(k\text{-Mod}) \;\longrightarrow\; \Hot_{\bullet}
\]
le premier de ces foncteurs, il ne commute pas aux produits $*$ ($\Lotimes_k$ et $\times$ respectivement), mais on a, pour $P_{\bullet}$, $Q_{\bullet}$ variable dans $D_{\geq 0}(k\text{-Mod})$, une flèche fonctorielle canonique % [sic] variable -> variables, agreement with P, Q
\[
\gamma(P_{\bullet} \Lotimes_k Q_{\bullet}) \xleftarrow{\ \mathit{can.}\ } \gamma(P_{\bullet})\times\gamma(Q_{\bullet})
\]
`$k$-bilinéaire', ce qui montre que de donner des accouplements
\[
P_{\bullet} * Q_{\bullet} \xrightarrow{\ c_k\ } R_{\bullet},
\]
d'où
\[
\gamma(P_{\bullet} * Q_{\bullet}) \xrightarrow{\ \gamma(c_k)\ } \gamma(R_{\bullet}), \tag*{$(*)$}
\]
est l'un et l'autre plus fin que de se donner seulement
\[
\gamma(P_{\bullet}) * \gamma(Q_{\bullet}) \xrightarrow{\ c_{\bullet}\ } \gamma(R_{\bullet}), \tag*{$(*')$}
\]
vu qu'une donnée $(*)$ impliquera une donne $(*')$ par % [sic] donne -> donnée, missing letters
\[
\begin{tikzcd}
 & \gamma(P_{\bullet} * Q_{\bullet}) \arrow[dr, "\gamma(c)"] & \\
\gamma(P_{\bullet}) * \gamma(Q_{\bullet}) \arrow[ur, "\text{can.}"] \arrow[rr, "c_{\bullet}"] & & \gamma(R_{\bullet}).
\end{tikzcd}
\]
Ceci vu, on s'attend bien à avoir une théorie \textit{forte} des accouplements
\printedpage{13}
% head: A. Grothendieck : Dérivateurs, Ch. IX, Édition des Universités de Montpellier II et Paris VII
% foot: Version 12 mars 2002
\msp{page 21}
\[
\Homext_k(a,b) \Lotimes_k \Homext_k(b,c) \;\longrightarrow\; \Homext_k(a,c),
\]
d'où on déduira les accouplements plus faibles dans $\Hot_{\bullet}$ (au lieu de $D_{\geq 0}(k\text{-Mod})$)
\[
\Homext_{\Hot_{\bullet}}(a,b) \times \Homext_{\Hot_{\bullet}}(b,c) \;\longrightarrow\; \Homext_{\Hot_{\bullet}}(a,c).
\]

\dgrblock{dgrb436}{A sufficient differential graded setting}
\begin{dcomment}{later}{p19:l'attente naturelle ici}{A sufficient differential graded setting}
Grothendieck now asks for more than additive Hom groups: he wants composition on their derived complexes before taking $H_0$. A precise sufficient setting is a differential graded model, for example complexes of modules with quasi-isomorphisms. Choose one cofibrant replacement $Qa\to a$ for each object, in the projective model structure; all complexes are fibrant. The complexes $\operatorname{Hom}^\bullet(Qa,Qb)$ then compose literally by differential graded composition, since the middle representative is the same in the two factors. This is the enhancement described in \citet[\S\S2.2 and 3.2]{Keller2006DG}. The claim here is about this specified setting. The text does not state a set of hypotheses on an arbitrary additive relative category that would construct such an enhancement, and the existence of an adjoint equivalence $\Sigma\dashv\Omega$ alone does not give that construction.

For a commutative ring $k$, put
\[
 C(a,b)=\tau_{\leq0}\operatorname{Hom}^\bullet_k(Qa,Qb)
       \simeq\tau_{\leq0}\mathbf R\operatorname{Hom}_k(a,b).
\]
The good truncation retains the negative degrees and the cycles in degree zero. It belongs to $D^{\leq0}(k)$, or equivalently to $D_{\geq0}(k)$ after reversing degrees. Differential graded composition gives
\[
 C(a,b)\mathbin{\overset{\mathbf L}{\otimes}}_k C(b,c)
       \longrightarrow C(a,c),
\]
using the Koszul interchange to put the factors into the usual composition order. Composing two cycles of degree zero gives a cycle, and the sum of two nonpositive degrees is nonpositive, so composition restricts to the good truncations. Replacing the two input complexes by projective resolutions provides the indicated derived tensor map. Its associativity is inherited from differential graded composition and the coherent tensor comparison; its unit is $k\to C(a,a)$ carrying $1$ to the identity. Since $k$ is projective over itself,
\[
 \operatorname{Hom}_{D(k)}(k,C(a,b))=H^0C(a,b)
 =\operatorname{Hom}_{D(k)}(a,b).
\]

The unqualified ``ring'' in the source needs attention at exactly this tensor product. Two left modules over a noncommutative ring do not have the displayed tensor product with values in left modules. One may assume $k$ commutative, as above, or use the appropriate bimodule category. Moreover, the image of $k\to\operatorname{End}(\mathrm{id}_{\mathcal M})$ is central: any two natural endomorphisms commute by applying the naturality of one to a component of the other. Thus the scalar action already factors through a commutative image.

The connective complex produces a pointed space by Dold--Kan and forgetting the simplicial module structure. The convention is
\[
 \pi_n\operatorname{Map}^{h}(a,b)
   =H^{-n}\mathbf R\operatorname{Hom}_k(a,b)
   =\operatorname{Hom}_{D(k)}(a[n],b),\qquad n\geq0,
\]
with the first equality at the zero basepoint for $n>0$. The higher positive Ext groups of two degree-zero modules are not positive homotopy groups of this mapping space. This is the grading distinction explained in the commentary on Chapter~III, manuscript p.~12 (``Keeping all degrees of the mapping complex'').

For a concrete calculation, take $a=k[0]$ and $b=k[1]$, with cohomological convention $K[1]^i=K^{i+1}$. The source is projective, so the ordinary Hom complex already computes derived Hom and is $k[1]$, concentrated in degree $-1$. Dold--Kan gives $K(k,1)$; there is one component, its fundamental group is the additive group of $k$, and every higher homotopy group vanishes. On the other hand $\operatorname{Hom}_{D(k)}(k[1],k[1])=k$, exactly the group predicted by suspension representability. This example makes visible the information that would disappear if one kept only the ordinary zero map $k\to k[1]$.
\end{dcomment}


\dgrblock{dgrb437}{The bilinear comparison is lax monoidal}
\begin{dcomment}{exposition}{p20:une flèche fonctorielle canonique}{The bilinear comparison is lax monoidal}
Grothendieck correctly asks for a comparison map in the direction needed to transport composition. It need not be an isomorphism. Let $P$ and $Q$ be nonnegative chain complexes and choose nonnegative projective resolutions $P'\to P$ and $Q'\to Q$. Their Dold--Kan simplicial modules $A=\Gamma P'$ and $B=\Gamma Q'$ have a degreewise map of pointed sets
\[
 U(A)\times U(B)\longrightarrow U(A\otimes_k B),
 \qquad (a,b)\longmapsto a\otimes b.
\]
The shuffle equivalence identifies the normalized complex of $A\otimes_k B$ with $P'\otimes_k Q'$, which computes $P\otimes_k^{\mathbf L}Q$. This yields the source's comparison in $\Hot_\bullet$. It factors through the smash product because either zero input is sent to zero. For coherence, use the unital associative Alexander--Whitney comparison $\Gamma P'\otimes_k\Gamma Q'\to\Gamma(P'\otimes_k Q')$; its normalized map is inverse to shuffle up to chain homotopy. The derived construction is independent of the chosen projective replacements. Dold--Kan itself is Theorem 14.24.3 of \citet{StacksDoldKan}; the tensor comparisons, their identities, and the induced monoidal structure on $\Gamma$ are checked in \citet[\S2.3, especially equation (2.8)]{SchwedeShipley2003}.

Already with $P=Q=k[0]$, the map on components is
\[
 k\times k\longrightarrow k,\qquad (x,y)\longmapsto xy,
\]
which is not injective for any nonzero ring, since $(0,0)$ and $(0,1)$ have the same image. It nonetheless carries a tensor pairing to a pointed or unpointed pairing by composition. Thus ``more information'' means that the bilinear factorization is specified, not that the forgetful functor is fully faithful or that its monoidal comparison is invertible. The displayed Cartesian product is legitimate for this weaker pairing; to apply the enriched unit convention in pointed spaces one uses its smash factorization, as explained above.
\end{dcomment}


3$^o$) Revenant au cas $\M$ ponctuée sans plus, \textsc{Quillen} définit, pour $a$, $b$ dans $\M$,
\[
\pi_1(a,b) \;=\; \pi_1(a,b\,;0_a,0_b)
\]
comme bifoncteur sur $(\M\Sigma^{-1})^o\times(\M\Sigma^{-1})$ (avec les notations de 1$^o$)). Si ce que j'avance sans preuve dans 1$^o$) est vrai, il doit s'ensuivre qu'on a un isomorphisme canonique
\[
\pi_1(a,b) \;\simeq\; \pi_1(\Ch^{\Sigma}_{\infty}(a,b)),
\]
où $\Ch^{\Sigma}_{\infty}(a,b)$ est vu comme objet de $\Hot_{\bullet}$, donc son $\pi_1$ est défini. On a en tous cas, sans aucune référence à des cofibrations et des fibrations, un bifoncteur à valeurs dans Groupes, composé dans
\[
(\M\Sigma^{-1})^o\times(\M\Sigma^{-1}) \xrightarrow{\ \ (a,b)\longmapsto \Ch^{\Sigma}_{\infty}(a,b)\ \ } \Hot_{\bullet} \xrightarrow{\ \pi_1\ } (\text{Groupes}).
\]
Si pour $a$ fixé, $b$ variable,
\[
b \;\longmapsto\; \pi_1(a,b)
\]
est représentable, on note l'objet qui le représente par $\Sigma(a)$ (\textit{suspension} -- ne pas confondre ce $\Sigma$ avec $\Sigma \subseteq \mathrm{Fl}(\M)$, qu'il vaut mieux maintenant noter $W$, comme `weak equivalence'). Donc

\msp{page 22}

\[
\mathrm{Hom}(\Sigma a,b) \;\simeq\; \pi_1(a,b),
\]
et $\Sigma a$ dépend fonctoriellement de $a$, sous réserve d'existence. S'il existe toujours, on a
\[
\Sigma :\ \M W^{-1} \;\longrightarrow\; \M W^{-1}.
\]
On définit de même, sous réserve d'existence,
\[
\Omega :\ \M W^{-1} \;\longrightarrow\; \M W^{-1}
\]
par
\[
\mathrm{Hom}(a,\Omega b) \;\simeq\; \pi_1(a,b),
\]
et si tous les deux existent, on a la formule d'adjonction
\[
\mathrm{Hom}(\Sigma a,b) \;\simeq\; \mathrm{Hom}(a,\Omega b) \qquad (\simeq \pi_1\Homext_{\Hot_{\bullet}}(a,b)).
\]

\dgrblock{dgrb438}{Representability and its higher counterpart}
\begin{dcomment}{exposition}{p22:et si tous les deux existent}{Representability and its higher counterpart}
The representability argument is complete once the pointed mapping bifunctor has been constructed. Let $\mathcal C=\mathcal M[W^{-1}]$ and write $F(a,b)=\pi_1\operatorname{Map}^{h}(a,b)$ at zero. If $F(a,-)$ is represented by $\Sigma a$, then a morphism $a\to a'$ induces $F(a',-)\to F(a,-)$; Yoneda gives $\Sigma a\to\Sigma a'$. Uniqueness in Yoneda proves the functor laws. The same argument in the other variable constructs $\Omega$, and the two representing bijections give $\Sigma\dashv\Omega$. The appearance of $0_a,0_b$ in the source's page-21 expression must be read as two copies of $0_{a,b}$, since both endpoints of a path lie in the same mapping space.

Representability of $\pi_1$ alone supplies only a statement on components of loop spaces. To obtain the displayed equivalences of whole mapping spaces one requires homotopy limits and colimits in the enhanced category. In a pointed category with derived mapping spaces and a homotopy pullback defining $\Omega b=0\times_b^h0$, the defining universal property gives
\[
 \operatorname{Map}^{h}(a,\Omega b)
  \simeq *\times_{\operatorname{Map}^{h}(a,b)}^h*
  =\Omega\operatorname{Map}^{h}(a,b).
\]
Dually, the homotopy pushout $\Sigma a=0\amalg_a^h0$ gives the other equivalence. Both respect loop multiplication because concatenation is induced by the same pointed path construction. This is the additional structure behind the source's test~(c), rather than a consequence of an arbitrary bijection on $\pi_0$.

Quillen proves the group-valued bifunctor and its two representabilities for a pointed model category in Chapter I, \S2; the checked retypesetting labels the result Theorem 1.2 \citep{Quillen1967}. The same construction is built into pointed derivators, where the coherent square, rather than a square in the underlying homotopy category, defines suspension and loops \citep[Definition 3.16 and Proposition 3.17]{Groth2013}. The unit being invertible makes $\Sigma$ fully faithful; the counit being invertible makes $\Omega$ fully faithful. If both are invertible, they form an equivalence. Requiring both is the stability condition, not a prerequisite for the existence of the two functors in a pointed model theory.

The following paragraph also uses existence of the appropriate finite products and coproducts. With those in place, a natural multiplication on $\operatorname{Hom}(-,\Omega b)$ is, by Yoneda, a map $\Omega b\times\Omega b\to\Omega b$; the group equations follow by the same argument. Dually $\Sigma a$ is a cogroup. In the enhanced setting, a derived mapping-space functor carries homotopy products in its second variable, and homotopy coproducts in its first, to products of spaces. These are the products and coproducts that justify the source's compatibility with group objects. Existence of ordinary products in the localized category, without a comparison to the enhanced products, would not by itself prove the mapping-space statement. On a pointed derivator $\Omega$ is right adjoint to $\Sigma$ as a morphism of derivators, so it preserves homotopy right Kan extensions, and $\Sigma$ preserves homotopy left Kan extensions; this supplies the left and right exactness anticipated in footnote 8.
\end{dcomment}


% ---- from Der-ChIX-678_p016-018.tex
\printedpage{14}
% head: A. Grothendieck : Dérivateurs, Ch. IX, Édition des Universités de Montpellier II et Paris VII
% foot: Version 12 mars 2002

L'existence de
\[
\begin{aligned}
\Omega :\ & \mathcal{M}W^{-1} \ \longrightarrow\ \mathcal{M}W^{-1}, \quad \text{ou de}\\
\Sigma :\ & \mathcal{M}W^{-1} \ \longrightarrow\ \mathcal{M}W^{-1},
\end{aligned}
\]
voire même des deux à la fois, apparaît comme une hypothèse \cercle{a} extrêmement forte sur $(\mathcal{M}, W)$. Cette condition \textit{plus} celle \cercle{b} que les $\Chu_{\infty}(a,b)$ soient homotopes à des catégories de Thomason-Kan, ou (ce qui revient au même, à peu de choses près?) que $Y \longmapsto \mathrm{Hom}_{\mathcal{M}(Y)W_Y^{-1}}(a_Y, b_Y)$ transforme Hot-équivalences de Cat en bijections (ne serait-ce que pour $Y$ fixé \dots), semble le test crucial pour savoir si `$(\mathcal{M}, \Sigma)$ \texttt{[plutôt} $(\mathcal{M}, W)$\texttt{]} se prête à une théorie homotopique `dans $\mathcal{M}$' '. Je doute que

\msp{page 23}

ces conditions puissent être satisfaites pour $\mathcal{M}$ finie, sauf dans le cas trivial où $\mathcal{M} = \emptyset$, où $\mathcal{M}\Sigma^{-1}$ \texttt{[plutôt} $\mathcal{M}W^{-1}$\texttt{]} est équivalente à la catégorie ponctuelle.

\dgrblock{dgrb439}{A finite test of the proposed axioms}
\begin{dcomment}{development}{p23:ces conditions puissent être satisfaites}{A finite test of the proposed axioms}
Grothendieck suspects that his tests already force an infinite supply of objects and operations. This would be natural for a model of a nontrivial stable homotopy theory, but the tests actually written permit a finite pointed category. The following calculation isolates the missing requirement. We checked the cited localization and model-category sources and the commentary on the earlier chapters; we did not find this particular example discussed as a response to this passage. Its role is diagnostic, without a claim of priority.

\textbf{A finite relative category satisfying tests (a) and (b).}\ours
Let $\mathcal E$ have objects $0,x$, with $0$ a zero object, $\operatorname{End}(x)=\{1_x,z\}$, and $z$ the composite $x\to0\to x$, so $z^2=z\ne1_x$. All other Hom sets are singletons. This is the full subcategory of pointed sets on a singleton and a two-element pointed set. Take $W$ to consist of identities. Then $\mathcal E[W^{-1}]=\mathcal E$, which is neither empty nor equivalent to a point.

We first compute the actual path category. With all backward edges identities, a bilateral eventually constant path is an eventually constant forward chain from $a$ to $b$, with duplicated vertices at the backward steps. Composition of its finite nonidentity part defines a functor
\[
 q:\Ch_\infty^W(a,b)\longrightarrow
       \operatorname{Hom}_{\mathcal E}(a,b)_{\mathrm{disc}}.
\]
Naturality with fixed endpoints makes the composite invariant under morphisms of paths. Choose the section $s$ that inserts a morphism at the forward edge $-1\to0$, with identities elsewhere. Define a third functor $R$ by replacing the part strictly before $0$ by the constant value $a$, retaining the part from $0$ onward, and using the composite $a\to a_0$ at the joining edge. The eventual identity tails make that composite well-defined. There is a natural transformation $R\to\mathrm{id}$: its early components are the composites $a\to a_i$, and its later components are identities. There is also $R\to sq$: its early components are identities and its later components are the composites $a_i\to b$. All the squares commute by composition of forward chains; duplicated backward vertices have identical components. Thus
\[
 \mathrm{id}\longleftarrow R\longrightarrow sq,
 \qquad qs=\mathrm{id}.
\]
This proves a categorical homotopy equivalence to the discrete Hom category. In particular every component is contractible, every $\pi_1$ is zero, and the path categories are homotopy equivalent to discrete categories, which are Thomason--Kan fibrant. This verifies the first formulation of test~(b).

The second formulation is equally direct. Since $W_Y$ contains only identities, no localization changes $\mathcal E(Y)$. A natural transformation between the constant diagrams $a_Y,b_Y$ assigns one arrow $a\to b$ to each connected component of $Y$, because naturality along every arrow forces equality. Hence
\[
 \operatorname{Hom}_{\mathcal E(Y)}(a_Y,b_Y)
 \cong\operatorname{Hom}_{\mathcal E}(a,b)^{\pi_0Y}.
\]
Every Hot-equivalence induces a bijection here. Finally $\Sigma a=0=\Omega a$ represents the identically trivial group-valued bifunctor $\pi_1(a,b)$ in both variables. Test~(a) holds, as do the mapping-space equivalences of test~(c): all their terms are points. The nontrivial endomorphisms of $x$ survive in $\pi_0$, so this is not the terminal exception.

The finite-limit hypothesis introduced on the next page would rule this example out for an elementary reason. If a category has finitely many arrows, say $N$, and all finite powers, then two distinct maps $a\to b$ would give at least $2^n$ maps $a\to b^n$ for every $n$, contradicting $2^n>N$. Thus every Hom set has at most one element. If the category is pointed, the zero maps make every Hom set nonempty, so it is equivalent to the terminal category. This is the finite form of the usual small-completeness argument; see \citet[Exercise 6.3.23(d)]{Leinster2014}. The distinction is therefore exact: the stated tests do not imply the finiteness obstruction, while a finite pointed category with finite products is already trivial before localization.
\end{dcomment}


Un troisième test \cercle{c} qui me paraît important (mais peut-être ça résulte du reste), c'est les formules dans $\Hot_\bullet$
\[
\left\{
\begin{aligned}
\Homext_\bullet(\Sigma a, b) \ &\simeq\ \Homext_\bullet(a, \Omega b)\\
&\simeq\ \Omega\Homext_\bullet(a, b),
\end{aligned}
\right.
\tag*{$(*)$}
\]
où le $\Homext_\bullet$ désigne le Hom externe à valeurs dans $\Hot_\bullet$. Il faut le voir séparément (sous condition d'existence soit de $\Omega$, soit de $\Sigma$) comme deux formules `duales' l'une de l'autre,
\[
\left\{
\begin{aligned}
\Homext_\bullet(a, \Omega b) \ &\simeq\ \Omega\Homext_\bullet(a, b)\\
\Homext_\bullet(\Sigma a, b) \ &\simeq\ \Omega\Homext_\bullet(a, b),
\end{aligned}
\right.
\tag*{$(**)$}
\]
isomorphismes de \textit{groupes} dans $\Hot_\bullet$. (\textbf{NB} $\Omega b$ est un objet groupe de $\mathcal{M}W^{-1}$, car il représente un contrafoncteur à valeurs dans Gr, et de même $\Sigma a$ est un cogroupe de $\mathcal{M}W^{-1}$. Le foncteur $y \longmapsto \Homext(a, y)$ doit commuter aux produits finis, et donc transformer objets groupes en itou. Et de même $x \longmapsto \Homext(x, b)$ doit transformer sommes finies en produits finis, donc cogroupes en groupes (\footnote[8]{Le foncteur $\Omega$ doit provenir d'un endomorphisme du dérivateur $\mathbf{D} = \mathbf{D}_{\mathcal{M},W}$, et celui-ci doit être exact à gauche, et dualement pour $\Sigma$.}).)

Je me suis convaincu que les isomorphismes $(**)$ n'ont de chance d'exister que si on a la condition préalable \cercle{b}.

\newpage

\printedpage{15}
% head: A. Grothendieck : Dérivateurs, Ch. IX, Édition des Universités de Montpellier II et Paris VII
% foot: Version 12 mars 2002

\msp{page 24}

\section*{3\quad Essai de construction du foncteur $\Omega$}

Il n'y a sans doute guère de chance de prouver quelque chose sur le foncteur
\[
\Omega :\ \mathcal{M}W^{-1} \longrightarrow \mathcal{M}W^{-1},
\]
sans tenir en mains une \textit{construction} plus ou moins explicite de celui-ci. Je vais essayer, en m'inspirant de \textsc{Quillen}, de dégager des conditions générales sur $(\mathcal{M}, W)$ (plus générales que celles de \textsc{Quillen}) qui assurent l'existence de $\Omega$, et permettent de le construire. L'idée géométrique remonte à \textsc{Serre-Cartan}. Il s'agit, pour un objet $S$ de $\mathcal{M}$, de construire un objet $\mathrm{Ch}(S)$ \textit{dans $\mathcal{M}$,} jouant le rôle d'un `espace de chemins', et qui doit s'inclure dans un diagramme de factorisation de $\mathrm{diag}_S : S \longrightarrow S \times S$ (on suppose que dans $\mathcal{M}$ les $\varprojlim$ finies, donc les produits finis, existent)
\[
\begin{tikzcd}
S \arrow[d, "i_S"] \arrow[dd, bend right=40, "\mathrm{diag}_S"'] & \\
\mathrm{Ch}(S) \arrow[d, "p_S"] & \\
S \times S & e. \arrow[l]
\end{tikzcd}
\]
La propriété essentielle, c'est que $p$ soit une \textit{`fibration'} (dans un sens convenable, à dégager).

\dgrblock{dgrb440}{Which factorization computes loops}
\begin{dcomment}{exposition}{p24:La propriété essentielle}{Which factorization computes loops}
Grothendieck returns to an actual path object because representability must eventually be proved by a construction. In a pointed model category, choose a fibrant replacement of $S$ and a factorization
\[
 S\xrightarrow{\sim}P S\xrightarrow{(d_0,d_1)}S\times S
\]
of its diagonal, with the first map a weak equivalence and the second a fibration. The fibre over $(0,0)$ computes $\Omega S$. Fibrancy is part of the useful hypothesis: it makes the endpoint maps fibrations to a fibrant object, so their pullbacks compute the intended homotopy fibres. The first map may be chosen to be a trivial cofibration, but its weak-equivalence property is essential. At this point the source is still seeking its axioms, so this is a completion of the proposed construction, not a correction of a theorem claimed to have been proved.

Merely asking for the second map to be a fibration, even adding a cofibration for the first, would be insufficient. In pointed simplicial sets take a nontrivial Kan complex $S=K(A,1)$, let $P S=S\times S$, and let the second map be the identity. The diagonal is a monomorphism, hence a cofibration, and $S$ is fibrant. Nevertheless the indicated fibre is a point, whereas the true loop space has components $A$. What fails is precisely the omitted weak equivalence $S\to P S$.

The representation can be read directly from mapping spaces. In a pointed simplicial model category, let $T$ be cofibrant. Applying $\operatorname{Map}(T,-)$ to the displayed factorization gives a path object for the Kan complex $K=\operatorname{Map}(T,S)$: its endpoint map is a fibration by the simplicial model axiom and its first map is a weak equivalence between fibrant objects. Mapping out of $T$ preserves the pullback that defines the fibre. Therefore
\[
 \operatorname{Map}(T,\Omega S)\simeq\Omega K,
 \qquad
 [T,\Omega S]\cong\pi_1\operatorname{Map}(T,S).
\]
The latter bijection is natural in both variables after replacements. It is independent of the chosen path object because all such objects compute the same homotopy pullback $0\times_S^h0$. For arbitrary model categories the corresponding statement is Quillen's path-object argument, using homotopy classes of left and right homotopies; the checked version gives it as equation (5) in the proof of Theorem 1.2 \citep[Chapter I, \S2]{Quillen1967}.

To replace the mapping space by $\Ch_\infty^W(T,S)$ in the final display of the manuscript requires its comparison with the derived mapping space. A path object in $\mathcal M$ and an interval of zigzags in $\mathrm{Cat}$ play complementary roles; they are not interchangeable definitions. The audit leaves the specific comparison visibly conditional while giving the completed argument in the model setting.
\end{dcomment}


\dgrblock{dgrb441}{Path objects before the full model structure}
\begin{dcomment}{context}{p24:sans tenir en mains une}{Path objects before the full model structure}
Grothendieck is looking for precisely the axioms needed to construct loops. This route had already been isolated by \citet[\S1, pp.~420--421]{Brown1973}, whose categories of fibrant objects retain weak equivalences, fibrations, and path objects without specifying cofibrations. Brown requires pullbacks of fibrations and of trivial fibrations to exist and remain of the same kind; every object is fibrant. His path-object axiom factors the diagonal by a weak equivalence followed by a fibration. These are sufficient structural data for the factorization lemma proved there, which replaces an arbitrary map by a weak equivalence followed by a fibration.

The comparison clarifies the two tentative choices on the next manuscript page. Fibrancy of $S$ supplies part of the working setting, while the essential property of $S\to\mathrm{Ch}(S)$ is its being a weak equivalence. A cofibration condition on that arrow is not needed for Brown's axioms. We use this established precedent to specify a sufficient framework, without suggesting that the manuscript cites Brown or that these axioms alone identify its particular bilateral path categories with derived mapping spaces.
\end{dcomment}


\msp{page 25}

À titre accessoire il peut être utile (peut-être) que $i$ soit une `cofibration' (\footnote[9]{et peut-être que $S \longrightarrow e$ soit également une fibration ?}). On définit $\Omega(S)$ comme la fibre de $\mathrm{Ch}(S)$ sur $S{\times}S$ (par rapport à l'unique morphisme $e = e_{\mathcal{M}} \longrightarrow S\times S$, où $e_{\mathcal{M}}$ est l'objet nul de $\mathcal{M}$).
\[
\begin{tikzcd}[column sep=large]
\mathrm{Ch}(S) \arrow[d] & \Omega S \arrow[l] \arrow[d]\\
S \times S & e. \arrow[l]
\end{tikzcd}
\]

\newpage

\printedpage{16}
% head: A. Grothendieck : Dérivateurs, Ch. IX, Édition des Universités de Montpellier II et Paris VII
% foot: Version 12 mars 2002

C'est un objet de $\mathcal{M}$, qu'on regardera surtout comme un objet de $\mathcal{M}W^{-1}$. Il faut s'assurer qu'on a une bijection fonctorielle en $T$ dans $\mathcal{M}W^{-1}$
\[
\underbrace{\pi_1\Chu^W_\infty(T, S)}_{\overset{\text{déf}}{=}\,\pi_1(T,S)} \ \simeq\ \underbrace{\mathrm{Hom}_{\mathcal{M}W^{-1}}(T, \Omega S)}_{\overset{\text{déf}}{=}\,\pi_0\Chu^W_\infty(T,\Omega S)}.
\]
Pour cela, il s'impose de comparer
\[
\Chu^W_\infty(T, \Omega S) \qquad\quad \text{et} \qquad\quad \underbrace{\Omega_{\Hot_\bullet}(\Chu^W_\infty(T, S))}_{\overset{\substack{\text{`par définition'}\\ \text{pour $\Omega$ de $\mathrm{Hot}_\bullet$}}}{=}\ \Chu_\infty(\mathrm{Ch}^W_\infty(T,S);0,0)} \qquad ,
\]
puisqu'il s'agit de construire une bijection entre leurs $\pi_0$.

\dgrblock{dgrb442}{The loop object of a complex, with its contraction}
\begin{dcomment}{examples}{p25:puisqu'il s'agit de construire une bijection}{The loop object of a complex, with its contraction}
For complexes the proposed construction can be written down without any choices, and its fibre displays the sign in the loop shift. Use cochain complexes of $k$-modules with the model structure whose weak equivalences are quasi-isomorphisms and whose fibrations are degreewise surjections. Every complex is fibrant. For a complex $B$, define
\[
 (P B)^n=B^n\oplus B^n\oplus B^{n-1},\qquad
 d(x,y,h)=(dx,dy,y-x-dh).
\]
The third component of $d^2$ is $dy-dx-d(y-x-dh)=0$, so this is a complex. The maps
\[
 i(b)=(b,b,0),\qquad p(x,y,h)=(x,y),\qquad r(x,y,h)=x
\]
are chain maps, $pi$ is the diagonal, and $p$ is a degreewise surjection. To verify the required weak equivalence, define the degree $-1$ homotopy $H(x,y,h)=(0,h,0)$. Direct calculation gives
\[
 (dH+Hd)(x,y,h)=(0,y-x,h)
               =(\mathrm{id}-ir)(x,y,h),\qquad ri=\mathrm{id}.
\]
Thus $i$ is a chain homotopy equivalence. Its fibre over zero under $p$ is the complex with terms $B^{n-1}$ and differential $-d$, namely $B[-1]$. We have obtained $\Omega B=B[-1]$ and, dually, $\Sigma B=B[1]$.

For a cofibrant replacement $Q T\to T$, the corresponding Hom complex satisfies
\[
 H^0\operatorname{Hom}^\bullet(QT,B[-1])
   =H^{-1}\operatorname{Hom}^\bullet(QT,B)
   =\pi_1\operatorname{Map}^{h}(T,B),
\]
which verifies the source's required bijection in this example. More than the bijection is visible: applying the good nonpositive truncation and Dold--Kan to these Hom complexes identifies the whole mapping space into $B[-1]$ with the loop of the mapping space into $B$. All maps used above commute with chain maps $B\to B'$, so the construction is functorial before localization. This is exactly the kind of explicit factorization that the closing paragraph asks us to find.
\end{dcomment}
\dgrchapter{X}{Chapter X}{\texorpdfstring{Comparaison de Cat avec $\boldsymbol{\Delta}^{\wedge}$}{Comparaison de Cat avec Δ∧}}
\dgrpart{0}{VIII-X}{Manuscript pages 1--5}{VIII-X\enspace\textperiodcentered\enspace Chapter X\enspace\textperiodcentered\enspace manuscript pages 1--5}{Chapter X \textperiodcentered\ \textup{VIII-X}}
% ---- from Der-ChX-789_p001-005.tex
% ===================== page 1 (p001) =====================
\printedpage{}
% no page number printed (title page)
% head: (none)

\begin{center}
{\LARGE \textsf{LES DÉRIVATEURS}}

\bigskip

{\large \textbf{ALEXANDRE GROTHENDIECK}}

\vspace{3em}

{\large \textbf{Chapitre X}}

\medskip

{\large \textbf{Comparaison de} Cat \textbf{avec} $\bDelta^{\wedge}$}
\end{center}

\clearpage

% ===================== page 2 (p002) =====================
\printedpage{}
% no page number printed (editorial note)
% head: (none)

Ce texte a été déchiffré et transcrit en \LaTeXe{} par M. Künzer. Il a été édité par M. Künzer, J. Malgoire, et G. Maltsiniotis. La transcription est aussi fidèle que possible au manuscrit. Pour les quelques corrections évidentes, ou rares commentaires des éditeurs, ainsi que pour la numérotation originale des pages du manuscrit, les caractères de machine à écrire \texttt{[typewriter]} entre crochets sont utilisés. Un point d’interrogation entre crochets signifie que l’on n’est pas sûr du mot qui précède.

Cette édition est provisoire. Les remarques, commentaires et corrections sont bienvenus. Envoyer un message à :

\textbf{maltsin@math.jussieu.fr}

\hfill G. Maltsiniotis\hspace*{3em}

\clearpage

% ===================== page 3 (p003), printed 1 =====================
\printedpage{1}
% head: A. Grothendieck : Dérivateurs, Ch. X, Édition des Universités de Montpellier II et Paris VII
% foot: Version 9 avril 2002

\msp{page 1}

\section*{1\quad Rappel des notions de D-asphéricité, et lien avec le ‘localisateur fondamental $W_{\mathbf{D}}$’}

Pour un dérivateur $\D$, il s’agit des notions suivantes (cf. I) :

\dgrblock{dgrb443}{The homotopy shapes seen by a coefficient theory}
\begin{dcomment}{exposition}{p1:il s’agit des notions suivantes}{The homotopy shapes seen by a coefficient theory}
Grothendieck now returns to the part of a derivator that can be expressed without choosing a particular category of coefficients. The relevant equivalences concern \emph{constant} coefficients. Write $p_A:A\to e$ and $K_A=p_A^*K$. Then $u:A\to B$ belongs to $W_{\D}$ precisely when, for every $K,L\in\D(e)$, restriction gives a bijection
\[
 \Hom_{\D(B)}(K_B,L_B)\ \longrightarrow\
 \Hom_{\D(A)}(K_A,L_A).
\]
This recalls the distinction explained in the commentary on Chapter~I, manuscript p.~61 (``Equivalence measured by constant coefficients''). It is not the assertion that $u^*:\D(B)\to\D(A)$ is an equivalence on arbitrary coherent diagrams. The distinction is what makes $W_{\D}$ a theory of homotopy \emph{shapes}, while $\D$ still records the coefficients carried by those shapes.

The two comma categories retain different orientations. An object of $A/b$ is an arrow $u(a)\to b$; an object of $b\backslash A$ is an arrow $b\to u(a)$. Accordingly, the fibre of $A/a\to B/u(a)$ over $\beta:b\to u(a)$ consists of factorizations of $\beta$ by arrows $a'\to a$ lying over it. Its asphericity is the lifting criterion for smoothness, with the opposite criterion for properness. These are the categorical forms of the two base-change conditions, not assertions that the strict fibres $A_b$ alone control base change. The later formulation is \citet[Theorem~1.2 and its dual]{Maltsiniotis2005Asphericite}.
\end{dcomment}


\begin{itemize}
\item[(i)] \textit{Objet} $\D$-\textit{asphérique} de Cat (i.e. $Z \longrightarrow e$ est une $\D$-équivalence). NB Je préfère dire ‘$\D$-asphérique’ au lieu de ‘$\D$-contractile’, qui prête parfois à confusion.

\item[\textcircled{ii}] $\D$-équivalence pour une flèche $f : X \longrightarrow Y$ dans Cat.

\item[(iii)] \textit{Foncteur} $f : X \longrightarrow Y$ $\D$-\textit{asphérique} (i.e. les $X/y$ sont $\D$-asphériques) et $\D$-\textit{coasphérique} (les $y\backslash X$ sont $\D$-asphériques). C’est ce que j’avais appelé ‘$\D$-équivalence stricte’ (ou ‘costricte’), ou foncteur homologiquement $\D$-cofinal (ou cohomologiquement $\D$-cofinal), mais la terminologie ‘asphérique’ me semble la plus simple et la meilleure. (Pêchée dans PS, Pursuing Stacks.)

\item[(iv)] \textit{Foncteur} $f : X \longrightarrow Y$ $\D$-\textit{lisse,} $\D$-\textit{propre} (i.e. tel que les fibres des foncteurs induits $X/x \longrightarrow Y/f(x)$, respectivement $x\backslash X \longrightarrow f(x)\backslash Y$, soient $\D$-asphériques).

\item[(v)] $\D$-\textit{fibration} $f : X \longrightarrow Y$ (notion autoduale), i.e. telle que pour $Y'' \xrightarrow{u} Y'$ au dessus de $Y$, si $Y''$, $Y'$ \texttt{[sont]} Hot-asphériques, alors $u_X : X'' \longrightarrow X'$ dans
\[
\begin{tikzcd}
X \arrow[d] & X' \arrow[l] \arrow[d] & X'' \arrow[l] \arrow[d] \\
Y & Y' \arrow[l] & Y'' \arrow[l]
\end{tikzcd}
\]
est une $\D$-équivalence (\footnote{On dira aussi que $f$ est $\D$-\textit{fibrant.} (\textbf{NB} Il n’y a pas de notion ‘$\D$-cofibrant’ à ma connaissance.)}). La $\D$-\textit{fibration} est dite \textit{triviale} (ou \textit{asphérique}) si et seulement si ses fibres sont asphériques, ou encore si et seulement si $f$ est asphérique, ou si et seulement si $f$ est coaspérique. Cela signifie aussi que $f$ est une $\D$-équivalence \textit{universelle,} % [sic] coaspérique -> coasphérique, letter h missing
\msp{page 2}
i.e. que pour tout $Y' \longrightarrow Y$, $X' = X \times_Y Y' \longrightarrow Y'$ est une $\D$-équivalence.

\item[(vi)] On a aussi des variantes affaiblies de la notion de $\D$-fibrations, savoir les ‘$\D$-\textit{fibrations à gauche’,} ou les $f$ ‘\textit{cohomologiquement} $\D$-\textit{fibrants’,} et de même les ‘$\D$-\textit{fibrations à droite’,} i.e. les $f$ ‘\textit{homologiquement} $\D$-\textit{fibrants’} (\footnote{Nouvelle terminologie : fibrations cohomologiques (resp. homologiques) \textit{faibles,} cf. I, pages 111-126 pour une révision corrigée des diverses notions de ‘$\D$-fibrations’.}). La notion ‘gauche’
\printedpage{2}
% head: A. Grothendieck : Dérivateurs, Ch. X, Édition des Universités de Montpellier II et Paris VII
% foot: Version 9 avril 2002
ou cohomologique pour $f : X \longrightarrow Y$ signifie que les foncteurs $X_y \hookrightarrow X/y$ ($y \in \Ob Y$) et $X/y \longrightarrow X/y'$ (pour $(u : y \longrightarrow y') \in \Fl Y$) sont des $\D$-équivalences. Et dualement avec les $X_y \hookrightarrow y\backslash X$, $y'\backslash X \longrightarrow y\backslash X$ pour la notion ‘droite’ ou ‘homologique’. Les $\D$-fibrations ne sont autres que les $\D$-fibrations gauches (cohomologiques) universelles ou les $\D$-fibrations droites (homologiques) universelles.
\end{itemize}

\dgrblock{dgrb444}{Why triviality is a universal condition}
\begin{dcomment}{exposition}{p2:fibrations gauches (cohomologiques) universelles}{Why triviality is a universal condition}
The point of the fibration condition is to make the fibre calculations survive a change of base. The qualifier ``universal'' in the next sentence therefore carries real mathematical content; the earlier distinctions are reviewed in the commentary on Chapter~I, manuscript p.~104 (``The universal criterion, with the quantifiers restored'').

Here is a useful direct verification of the asserted triviality criterion. Suppose $f:X\to Y$ is a $\D$-fibration and every $X_y$ is $W_{\D}$-aspherical. After any change of base $B\to Y$, the same hypotheses hold. For $b\in B$, apply the definition to $e\to B/b$, selecting its final object. Both bases are ordinarily contractible, so
\[
 (X\times_Y B)_b\longrightarrow (X\times_Y B)/b
\]
lies in $W_{\D}$. Hence every category on the right is $W_{\D}$-aspherical. The asphericity theorem gives $X\times_Y B\to B\in W_{\D}$. Conversely, a universal $W_{\D}$-equivalence has aspherical fibres by base change to points. It is a $\D$-fibration: in its defining square the two projections to $Y''$ and $Y'$ lie in $W_{\D}$, and $Y''\to Y'$ does also because both bases are ordinarily contractible; two-out-of-three gives the required arrow between pullbacks. Applying the same comparison with $Y/y$ and $y\backslash Y$ proves the two equivalent comma-category criteria. The ordinary contractibility of the test bases, specified in the manuscript, has not been replaced by $W_{\D}$-asphericity.
\end{dcomment}


Finalement, toute cette kyrielle de notions se déduit de la seule notion de $\D$-équivalence (ii). Si je désigne par
\[
W_{\D} \subseteq \Fl(\Cat)
\]
l’ensemble de ces $\D$-équivalences, on trouve ce que dans Pursuing Stacks j’ai appelé un ‘localisateur’ (localiseur) dans Cat (\footnote{Mais attention, je suppose ici que le domaine Diag de $\D$ est Cat tout entier!}). Dans Pursuing Stacks c’était une telle partie
\[
W \subseteq \Fl(\Cat),
\]
et non la donnée d’un dérivateur, qui
\msp{page 3}
était pris comme la donnée de départ fondamental pour faire de l’algèbre homotopique en tous genres. J’avais d’ailleurs développé une kyrielle de notions du même genre (i) à (vi), en partant de cette donnée $W$. Il se pourrait que les énoncés pour les notions que j’ai développé dans I, en termes d’un dérivateur, soient valables plus généralement en termes d’un localisateur $W$, dans la mesure du moins où ils gardent un sens. (Donc pour l’instant je mets à part les énoncés proprement cohomologiques ou homologiques, qui font directement appel à $\D$.) Encore qu’il se pourrait que moyennant des conditions à expliciter sur $W$, qui impliquent que $(\Cat, W)$ définissent bel et bien un dérivateur $\D_{\Cat,W} = \D$, on arrive encore à interpréter $W$ comme un $W_{\D}$, de sorte que les énoncés cohomologiques et homologiques gardent un sens.

\dgrblock{dgrb445}{Realizing a localizer by a derivator}
\begin{dcomment}{later}{p3:les énoncés cohomologiques et homologiques gardent un sens.}{Realizing a localizer by a derivator}
Grothendieck asks for a coefficient theory realizing the proposed weak equivalences. A precise answer is available under an accessibility hypothesis. Fix a universe for small diagram shapes and a larger universe for coefficient categories, and call a fundamental localizer accessible when it is generated by a small set of functors. \citet[Theorem~5.2.15]{Cisinski2006} constructs a cofibrantly generated, left proper model structure on $\Cat$ with weak equivalences $W$. The model-category construction of derivators then gives
\[
 \D_W(A)=\Cat^{A^o}\bigl[(W^{A^o})^{-1}\bigr]
\]
with both homotopy Kan extensions \citep[Theorem~6.11 and its dual]{Cisinski2003}. The comparison with simplicial presheaves and the recovery of $W$ from constant-coefficient tests are supplied by \citet[Lemma~4.2.13, Theorem~4.2.15, and Corollary~4.2.18]{Cisinski2006}. Thus $W_{\D_W}=W$. Accessibility is the hypothesis of this realization statement; it is not implicit in the definition of an arbitrary derivator.

Realization does not determine a pre-existing coefficient theory. For a complete and cocomplete category $\mathcal C$, the represented derivator satisfies
\[
 \Hom_{\mathcal C^{A^o}}(K_A,L_A)
   =\operatorname{Map}\bigl(\pi_0 A,\Hom_{\mathcal C}(K,L)\bigr).
\]
Indeed, naturality says exactly that the component of a natural transformation is constant along every zigzag in $A$. For $\mathcal C=\mathrm{Set}$ or $\mathrm{Ab}$, a Hom set with at least two elements detects whether $\pi_0(u)$ is bijective: a missing component prevents injectivity of precomposition, and two identified components prevent surjectivity. Both derivators consequently give the same localizer $W_0$ of component bijections. Their categories of coefficients are not equivalent, since $\mathrm{Ab}$ has a zero object and $\mathrm{Set}$ does not. The equality $W_{\D_W}=W$ answers the question on these pages without identifying the original $\D$ with $\D_W$.
\end{dcomment}


J’ai envie d’abord de regarder dans Pursuing Stacks les axiomes que j’ai dégagé pour les localisateurs, et vérifier s’ils sont tous vérifiés pour
\msp{page 4}
$W_{\D}$. Ces conditions L 1 à L 5 sont énoncées dans Pursuing Stacks, §64 (page 164).

\textbf{L 1}\enspace Saturation de $W_D$. O.K.

Cela résulte formellement du fait qu’on a un foncteur
\[
\Cat \xrightarrow{\varphi} C \qquad (= \Cat^o)
\]
tel que pour $u \in \Fl(\Cat)$, $u \in W_{\D} \iff \varphi(u)$ isomorphisme. (Prendre pour $\varphi(X)$ la sous-catégorie de $\D(X)$ formée des $\D$-coefficients \textit{constants} sur $X$).

\dgrblock{dgrb446}{Making the saturation argument literal}
\begin{dcomment}{exposition}{p4:isomorphisme. (Prendre pour}{Making the saturation argument literal}
Grothendieck's argument is sound once the category of constants is represented with its coefficient labels retained. Define $C_{\D}(A)$ to have objects $K\in\D(e)$ and morphisms
\[
 C_{\D}(A)(K,L)=\Hom_{\D(A)}(K_A,L_A).
\]
Restriction is the identity on these object labels. Thus $u\in W_{\D}$ if and only if $C_{\D}(u)$ is an \emph{isomorphism} of categories, whereas the essentially full subcategory of constant objects gives the equivalent formulation using equivalences of categories. This avoids choosing representatives for coincident or isomorphic constant diagrams.

The construction proves strong saturation. If the image of $u$ becomes invertible in $\Cat[W_{\D}^{-1}]$, the factorization of the contravariant functor $C_{\D}$ through that localization makes $C_{\D}(u)$ invertible; hence $u\in W_{\D}$. In \emph{Pursuing Stacks}, \S64, ``saturation'' is the weaker list of identities, two-out-of-three, and the two-sided-composite condition. The present argument proves more than that list requires \citep[\S64, manuscript pp.~163--165]{Grothendieck1983}.
\end{dcomment}


\textbf{L 2}\enspace “$\bDelta^1$ est $W$-asphérique sur $e$”, i.e. pour tout $X$ dans Cat, $X \times \bDelta^1 \longrightarrow X$ est dans $W$.

Cela signifie aussi que $\bDelta^1 \longrightarrow e$ est une $\D$-fibration triviale, ce qui est bien le cas.

\clearpage

% ===================== page 5 (p005), printed 3 =====================
\printedpage{3}
% head: A. Grothendieck : Dérivateurs, Ch. X, Édition des Universités de Montpellier II et Paris VII
% foot: Version 9 avril 2002

\textbf{L 3}\enspace Si $X \in \Ob\Cat$ a un objet final, alors $X \longrightarrow e$ est dans $W$.

En fait, nous exigeons pour les dérivateurs que ceci soit vrai en supposant seulement $X$ asphérique (au sens absolu), et on en conclut même que $X \longrightarrow e$ est $W$-asphérique, i.e. $X \times Y \xrightarrow{\mathrm{pr}_2} Y$ est dans $W$ pour tout $Y$ dans Cat (\footnote{En fait L 2 $\Longrightarrow$ L 3.}).

\textbf{L 4}\enspace Si $f : X \longrightarrow Y$ est tel que $X/y \longrightarrow Y/y$ sont dans $W$, alors $f$ est dans $W$.

C’est OK.

\msp{page 5}
\textbf{L 5}\enspace Si $X$, $Y$ sont Cat-fibrées sur $S$, et $X \longrightarrow Y$ est un $S$-morphisme \textit{cartésien,} induisant sur les fibres des $X_s \longrightarrow Y_s$ qui sont dans $W$, alors $f$ est dans $W$.

Pour $W_{\D}$, c’est la proposition 6, page 71 de I.

Ces axiomes L 1 à L 5 (équivalent à L 1, L 3$'$ ($\bDelta^1 \longrightarrow e$ est dans $W$) et L 5) sont dans la nature d’axiomes de stabilité, assurant que $W$ est assez grand. Il y a une autre condition qui limite $W$ à ne pas être trop grand :

\dgrblock{dgrb447}{Completing the reduction to three axioms}
\begin{dcomment}{exposition}{p5:assurant que $W$ est assez grand.}{Completing the reduction to three axioms}
The shortened axiom list expresses a useful economy: integration of families, together with one interval, already forces the basic homotopy calculus. The reduction is sketched in \emph{Pursuing Stacks}, \S64. We spell it out, using the twisted-arrow construction of Chapter~I rather than its mapping-cone presentation; compare the commentary on Chapter~I, manuscript p.~73 (``The twisted-arrow category and its two projections'').

Assume L~1, L~3$'$, and L~5. Applying L~5 to the constant fibrations $A\times[1]\to A$ gives L~2. Their two sections belong to $W$ by two-out-of-three. A natural transformation $F\Rightarrow G$, viewed as a functor $A\times[1]\to B$, therefore shows that $F\in W$ if and only if $G\in W$. A category with a final or initial object retracts to that object, with its other composite naturally homotopic to the identity. Saturation gives L~3 and its opposite.

To pass between the two comma orientations, let $T(A)$ have objects arrows $a\to b$ and morphisms from $h:a\to b$ to $h':a'\to b'$ pairs
\[
 r:a\to a',\qquad s:b'\to b,\qquad s h'r=h.
\]
Both projections $T(A)\to A$ and $T(A)\to A^o$ are Grothendieck fibrations. Their fibres are $(a\backslash A)^o$ and $A/b$, respectively, and have final objects. L~5 consequently puts both projections in $W$. The construction is natural in $A$, so its two commutative squares and two-out-of-three give
\[
 u\in W\quad\Longleftrightarrow\quad u^o\in W.
\]
This also establishes the final assertion L~d).

Next let $u:A\to B$ have $W$-aspherical categories $b\backslash A$. The Grothendieck construction
\[
 E=\int_{b\in B}(b\backslash A)\longrightarrow B
\]
is a fibration with these fibres, hence lies in $W$ by L~5. Its objects are triples $(b,a,\alpha:b\to u(a))$. The forgetful functor $r:E\to A$ has the section $i(a)=(u(a),a,1)$, satisfies $ri=1$, and has a natural transformation $1_E\Rightarrow ir$ with component $(\alpha,1_a)$. Thus $r$ and $i$ belong to $W$, and $u$ does also. Opposite invariance proves the analogous assertion for $A/b$. Under the hypothesis of L~4, the final object of $B/b$ makes every $A/b$ aspherical by two-out-of-three, so L~4 follows. Conversely, L~1--L~5 include L~1 and L~5 and imply L~3$'$ by L~3. This proves the claimed equivalence of axiom systems. For the modern relative form, suppose $A\to B$ lies over $C$ and every $A/c\to B/c$ is in $W$. Opposite invariance gives objectwise equivalences $(A/c)^o\to(B/c)^o$. These diagrams define fibrations over $C^o$, so L~5 applies to their Grothendieck constructions; their forgetful functors are deformation retractions to $A^o$ and $B^o$. Thus $u^o$, and hence $u$, belongs to $W$. This yields the relative Theorem~A axiom of \citet[Definition~1.1.3]{Cisinski2004Minimal}. No model structure is used.
\end{dcomment}


\textbf{L a)}\enspace Si $f \in W$, alors $\pi_0(f)$ est bijectif.

Ceci n’est pas satisfait pour tous les $W_{\D}$ – p. ex. si $\D$ est le dérivateur \textit{nul,} alors $W_{\D} = \Fl\Cat$! Mais pour les dérivateurs utiles en pratique, cela semble toujours satisfait.

\dgrblock{dgrb448}{A nonzero coefficient theory that forgets components}
\begin{dcomment}{examples}{p5:Mais pour les dérivateurs utiles en pratique}{A nonzero coefficient theory that forgets components}
The additional condition L~a) restricts the amount of homotopy theory the coefficients may forget. Even a nonzero derivator need not satisfy it. Take the complete lattice $\mathbf 2=\{0<1\}$ and its represented derivator
\[
 \D_{\mathbf 2}(A)=\operatorname{Fun}(A^o,\mathbf 2).
\]
For nonempty $A$, the Hom sets between the two constant diagrams are the Hom sets of $\mathbf 2$: all are singletons except $\Hom(1_A,0_A)=\varnothing$. For $A=\varnothing$ there is one diagram and one morphism. Therefore
\[
 u:A\to B\in W_{\D_{\mathbf 2}}
 \quad\Longleftrightarrow\quad
 (A=\varnothing\ \Longleftrightarrow\ B=\varnothing).
\]
In particular the fold $e\amalg e\to e$ belongs to this localizer, although its map on components is not injective. This is the usual coarse localizer $W_{-1}$ \citep[Examples~3.3.5--3.3.6]{Cisinski2006}. The terminal coefficient category gives the still larger localizer of all functors, exactly as Grothendieck observes. By contrast, coefficients with a Hom set of at least two elements detect components in the represented examples calculated above. His remark about useful examples is a comment on the intended coefficients, not an additional derivator axiom.
\end{dcomment}


\textbf{L b), c)}\enspace Le foncteur
\[
\Cat \longrightarrow \Cat\, W^{-1}
\]
commute aux sommes finies et aux produits finis (\footnote{Ça marche si $(\Cat, W)$ provient d’un triple de Quillen (avec des Q-fibrations, Q-cofibrations, Q-équivalences faibles).}).

À examiner ce que ça signifie en termes concrets sur $W$ – je n’ai pas vérifié que c’est O.K. pour les $W_{\D}$.

\dgrblock{dgrb449}{Finite products and sums already descend}
\begin{dcomment}{status}{p5:À examiner ce que ça signifie en termes concrets}{Finite products and sums already descend}
The questions L~b) and L~c) have a formal answer under the axioms already retained. If $u:A\to B$ is in $W$, the map of constant fibrations over any $C$ gives
\[
 u\times1_C:A\times C\longrightarrow B\times C\quad\in W
\]
by L~5. Applying L~5 over a discrete indexing category likewise shows that coproducts of arrows in $W$ belong to $W$. These closure properties are also \citet[Propositions~1.1.5 and 1.1.7]{Cisinski2004Minimal}.

Closure alone should be connected to the universal properties of the claimed products and sums. Put $H=\Cat[W^{-1}]$. The localization of $\Cat\times\Cat$ at the componentwise weak equivalences is $H\times H$: a pair of zigzags can be interleaved by holding the other coordinate fixed, and the commutation squares between the two coordinates remain valid after inversion. The diagonal and the product functor are homotopical, so their adjunction, including its unit and counit, descends to $H$. Thus the image of $A\times B$ represents the product of the images of $A$ and $B$. Descending the adjunction between coproduct and diagonal proves the corresponding coproduct assertion. The nullary adjunctions give the initial and terminal objects as well. Consequently the canonical localization preserves finite sums and finite products. A Quillen structure is unnecessary here; the explicit closure argument supplies the homotopy invariance needed for the \emph{ordinary} products and sums appearing in the question.
\end{dcomment}


\textbf{L d)}\enspace Si $f : X \longrightarrow Y$ dans $W$, alors $f^o : X^o \longrightarrow Y^o$ aussi.

C’est le corollaire 2, page 75 de I.
\begin{otherlanguage}{english}
\dgroverview{VIII-X}{Chapter X \textperiodcentered\ \textup{VIII-X} \textperiodcentered\ Overview}
\selectlanguage{english}
\noindent\textbf{Integration, spaces of morphisms, and the choice of weak equivalences.}

These three chapters investigate how much of homotopy theory can be reconstructed from operations on ordinary categories, provided that one keeps the right forms of coherence. Chapter VIII isolates integration for relative $1$-types. Chapter IX returns to the path categories introduced earlier, first to extract an abstract theory of composition and then to compare it with Quillen's constructions. The opening of Chapter X changes the point of departure: instead of beginning with a derivator and its coefficients, Grothendieck asks whether a class of weak equivalences of categories can itself generate the theory. The movement is from a concrete adjunction, through the structure of its spaces of morphisms, to the axioms that make such constructions possible.

These pages comprise all of Chapters VIII and IX and the opening of Chapter X, twenty-seven pages of the edition. Manuscript references must always include the chapter: they run through VIII.1--9, IX.1--25, and X.1--5. Our audit follows those separate sequences. Original statements, including the edition's retained slips and editorial interventions, remain distinguishable from the commentary. The two doubtful readings are settled against the edition; neither licenses a mathematical change to Grothendieck's text.

\medskip\noindent\textbf{Relative groupoids and a genuine biadjunction.}

The decisive construction of Chapter VIII is simple enough to write down on an object of the base. For $f:S\to T$ and a category $G$ fibred in groupoids over $S$, set
\[
 (f_!G)_t=\Pi(t\backslash_TG).
\]
The comma category records a point of $G$ together with an arrow from $t$ to its image in $T$; groupoid completion turns this record into a homotopy $1$-type. Its use is not limited to an already fibred category. For any $X\to T$, the same construction gives a relative groupoid $R_TX$, and extension along $X\to R_TX$ is universal for functors into categories fibred in groupoids. Cartesian lifts construct the extension, while their uniqueness supplies the coherence.

This universal property completes the adjunction verification left unfinished in the manuscript. It proves an equivalence of Hom-groupoids, not only a bijection on isomorphism classes. Both triangle comparisons arise from the same extension-and-restriction construction. Transitivity of integration is then the uniqueness of the left biadjoint to composite restriction. The point of the $2$-categorical formulation is thus practical: the automorphisms that would have been discarded by taking components are precisely what make the comparison usable in a further composition. The base-change theorem has the corresponding later expression in the theory of asphericity structures and smooth functors \citep[\S5]{Maltsiniotis2005Asphericite}.

The final descent statements require a more careful separation. The proper or smooth hypotheses of Theorem 3 support its assertions about faithfulness and about the essential image described by coefficients that are \emph{constant on each strict fibre}. They do not justify replacing this condition by the assertion that every vertical transport is an equivalence. A locally constant groupoid family can carry degree-two coherence even when its base is simply connected. We give a finite example over the fourteen-element face poset of the boundary of a tetrahedron. Every fibre is $B(\mathbf Z/2)$ and every transport is the identity, but a nonzero normalized $2$-cocycle prevents the total category from admitting a section. Consequently it cannot be equivalent to a constant family. This disproves part 3 and the subsequent corollary as printed, while leaving the preceding fibrewise-constant criterion intact.

The repair is expressed by higher monodromy: equivalence classes of locally constant groupoid families with a fixed fibre are controlled by maps to its automorphism $2$-type. A $1$-connected base removes ordinary monodromy but can retain the degree-two class; $2$-connectivity removes both \citep[Corollary~2.2.4 and Theorem~2.2.5, arXiv version]{PoleselloWaschkies2005}. A separate, shorter counterexample shows that section groupoids alone do not test all the proposed faithfulness conditions. The correct section tests must also be made after passing to the total category of the first coefficient. Both corrections concern the quantifiers and the coherence in descent, rather than the integration construction itself.

The general faithfulness question also admits an exact answer without smoothness or properness. For each arrow of $T$, form its category of factorizations through objects in the image of $f$. Restriction is faithful on transformations, fully faithful on Hom-groupoids, or an equivalence on Hom-groupoids precisely when these factorization categories are respectively nonempty, connected, or $1$-connected. We prove this three-level criterion by the counit on representables and coherent density, alongside the known untruncated homotopical-epimorphism criterion \citep[Theorem~6.10, arXiv version]{GrothStovicek2016Trees}.

\medskip\noindent\textbf{What composition gives, and what it presupposes.}

Chapter IX begins with a strong idea: composition of paths should explain endpoint functoriality. In a monoidal category $\mathcal H$, Hom objects $C(a,b)$ with associative composition and coherent units determine an ordinary category with morphisms $\operatorname{Hom}_{\mathcal H}(I,C(a,b))$. Inserting these elements into composition produces the bifunctor $C_0^{\mathrm{op}}\times C_0\to\mathcal H$. The proof is formal once the hypotheses are in place. Neither symmetry nor an internal Hom is needed. This is the familiar passage from an enriched category to its underlying ordinary category \citep[\S\S1.1--1.3]{Kelly1982}.

There are two limits to what this argument proves in the passage under review. First, the endpoint maps give isomorphisms only when those endpoint maps are invertible. The early corollary omits that hypothesis, and a discrete two-object example disproves the wording; the manuscript itself restores it on page 8. Second, the abstract argument assumes associative composition on the particular objects $\Ch_\infty(a,b)$. It does not construct that composition. Grothendieck sketches the finite-stage argument for the ind-object $\mathrm{CH}_\infty$ and explicitly leaves the corresponding realized-space verification unfinished.

The possible shortcut through finite tests does not solve this difficulty. The earlier definition of $\Hot_0$ uses homotopy types of finite categories, not just finite CW complexes. The resulting restricted Yoneda functor is not conservative, by the general theorem on small families of unpointed tests, and it is not faithful, by nonconcreteness of the homotopy category \citep[Theorem~2.1]{ArlinChristensen2023} and \citet[pp.~1, 5 of the reprint]{Freyd1970}. Thus equality after passing to those presheaves cannot replace a compatible homotopy on the original mapping spaces. The commentary also distinguishes an ind-object, an ordinary colimit in $\Hot$, and a homotopy colimit. An explicit increasing sequence of finite discrete sets has an ordinary colimit without being essentially constant as an ind-object. A finite-category test such as $BC_2$, on the other hand, is not compact with respect to all filtered homotopy colimits. These examples delimit the arguments that the comparison maps actually require.

Modern localization supplies the intended kind of mapping-space theory. The hammock localization has coherent mapping spaces and a component category equal to the ordinary localization \citep[\S2]{DwyerKan1980Calculating}. This identifies the modern object to which Grothendieck's construction should be compared; it does not by itself prove a comparison for his exact bilateral stabilization and transformation conventions. We keep that model-specific verification explicit. It is not described as a new open problem in the general theory of localization.

\medskip\noindent\textbf{Pointed and linear refinements.}

The later part of Chapter IX asks that the zero map become a specified basepoint in every mapping space. A basepoint in the unpointed homotopy category is insufficient: based homotopies remember data lost under free homotopy. The commentary constructs the pointed morphism associated to a map and a chosen path between basepoints, including its composition law. It also explains the significance of the manuscript's two corrective footnotes. The Cartesian unit in pointed spaces has only one map to every object; the unit that recovers all components is $S^0$ for the smash product. Composition must therefore annihilate either zero input. Forgetting basepoints carries this structure to unpointed composition through a lax monoidal comparison.

Some questions at the level of ordinary localization have particularly direct answers. Initial and terminal objects survive localization without extra hypotheses, by naturality of their unique arrows along both ordinary and inverted generators. Suitable closure under finite products or sums lets the corresponding adjunctions descend. For an additive category, stability of the denominator class under finite direct sums already yields an additive ordinary localization \citep[Lemma~2.2.1]{Krause2021}. The stronger request for complexes of morphisms requires an enhanced setting. We treat complexes of modules explicitly, with connective truncations of derived Hom, derived tensor composition, and the Dold--Kan comparison. This makes the grading visible: positive homotopy groups of the mapping space are negative cohomology groups of the Hom complex.

The suspension and loop objects are then representability statements. Representing the fundamental-group bifunctor yields an adjunction; the equivalences of whole mapping spaces arise from homotopy pushouts and pullbacks. An explicit path complex verifies the loop construction with its differential, its contraction, and its fibre. We also develop a finite counterexample to the proposed force of tests (a) and (b): the full category of pointed sets on one and two elements, with only identities inverted, satisfies both tests and even the displayed mapping-space identities. Its loop and suspension functors are zero, while its ordinary category is nontrivial. The missing strength is therefore not finiteness of the category alone. This calculation is marked as the commentators' development, without a priority claim.

\medskip\noindent\textbf{Localizers and the information retained by coefficients.}

The opening of Chapter X extracts a class $W_{\D}$ by testing maps of indexing categories against constant coefficients. It must be distinguished from the condition that restriction be an equivalence on all coherent diagrams. The localizer axioms express the stability of those constant-coefficient tests under interval homotopy, comma constructions, and integration of fibrations. We complete the reduction from the five displayed axioms to the shorter list, including the role of the twisted-arrow category in proving invariance under opposites. The finite-product and finite-sum questions are then settled by explicit closure arguments and descent of the product and coproduct adjunctions. No model structure is needed for these formal consequences.

Realization is a further theorem. For an accessible fundamental localizer, the later model structures and derivator constructions give a coefficient theory that realizes its weak equivalences \citep[Theorems~4.2.15 and 5.2.15]{Cisinski2006}; \citep[Theorem~6.11 and its dual]{Cisinski2003}. This does not reconstruct an arbitrary original derivator from $W_{\D}$. The represented derivators of sets and of abelian groups have the same component-detecting localizer, although their coefficient categories are inequivalent. The represented truth-value lattice gives an equally concrete nonzero example whose localizer detects only emptiness, disproving any attempt to impose component detection on all nonzero derivators. These examples show precisely what Chapter X has isolated: the homotopy shapes visible to a theory, while its further algebra of coefficients remains additional structure.
\end{otherlanguage}
\dgrchunkend

% DER_XI-XIIa: Grothendieck's text with the commentary set in, then the overview, inside a group that applies the chunk's own definitions
\dgrchapter{XI}{Chapter XI}{Hot-fibrations, foncteurs propres et foncteurs lisses etc. (dans Cat)}
\dgrpart{0}{XI-XIIa}{Manuscript pages 127--136}{XI-XIIa\enspace\textperiodcentered\enspace Chapter XI\enspace\textperiodcentered\enspace manuscript pages 127--136}{Chapter XI \textperiodcentered\ \textup{XI-XIIa}}
\dgrchunkbegin
% DER_XI-XIIa: the definitions of the chunk's preamble, applied inside its group
\let\printedpage\relax\newcommand{\printedpage}[1]{\marginpar{\footnotesize p.~#1}}
\let\msp\relax\newcommand{\msp}[1]{\marginpar{\footnotesize [#1]}}
\let\Hot\relax\newcommand{\Hot}{\mathrm{Hot}}
\let\Cat\relax\newcommand{\Cat}{\mathrm{Cat}}
\let\Ob\relax\newcommand{\Ob}{\operatorname{Ob}}
\let\Ch\relax\newcommand{\Ch}{\underline{\mathrm{Ch}}}
\let\Dl\relax\newcommand{\Dl}{\boldsymbol{\Delta}}
\let\fl\relax\newcommand{\fl}{\longrightarrow}
\let\hot\relax\newcommand{\hot}{\mathrm{hot}}
\let\uHom\relax\newcommand{\uHom}{\underline{\mathrm{Hom}}}
\let\id\relax\newcommand{\id}{\mathrm{id}}
\let\pr\relax\newcommand{\pr}{\mathrm{pr}}
\let\Fl\relax\newcommand{\Fl}{\mathrm{Fl}}
\let\card\relax\newcommand{\card}{\operatorname{card}}
\let\uFib\relax\newcommand{\uFib}{\underline{\mathrm{Fib}}}
\let\HOT\relax\newcommand{\HOT}{\mathrm{HOT}}
\let\HOTlc\relax\newcommand{\HOTlc}{\mathrm{HOT}^{\mathrm{lc}}_W}
\let\Hom\relax\newcommand{\Hom}{\mathrm{Hom}}
\let\Wd\relax\newcommand{\Wd}{W^{\mathrm{d}}_S}
\let\Wf\relax\newcommand{\Wf}{W^{\mathrm{f}}_S}
\let\dashedbox\relax\newcommand{\dashedbox}[1]{\tikz[baseline=(n.base)]\node[draw,dashed,inner sep=2pt](n){#1};}
\let\Fib\relax\newcommand{\Fib}{\mathrm{Fib}}
\let\Parf\relax\newcommand{\Parf}{\mathrm{Parf}}

% ---- from Der-ChXI-890_p001-005.tex
% ===================== page 1 =====================
\printedpage{--}
% head: (none) -- title page, no page number printed

\begin{center}
\vspace*{5cm}
{\LARGE\textsf{LES DÉRIVATEURS}}

\vspace{1.5em}
{\large\textbf{ALEXANDRE GROTHENDIECK}}

\vspace{3em}
{\large\textbf{Chapitre XI}}

\vspace{0.5em}
{\large $\Hot$\textbf{-fibrations, foncteurs propres et foncteurs lisses etc.}\\
\textbf{(dans} $\Cat$\textbf{)}}
\end{center}

\clearpage

% ===================== page 2 =====================
\printedpage{--}
% head: (none) -- no page number printed

\vspace*{6cm}
\noindent Ce texte a été déchiffré et transcrit en \LaTeXe\ par M. Künzer. Il a été édité par M. Künzer, J. Malgoire, et G. Maltsiniotis. La transcription est aussi fidèle que possible au manuscrit. Pour les quelques corrections évidentes, ou rares commentaires des éditeurs, ainsi que pour la numérotation originale des pages du manuscrit, les caractères de machine à écrire \texttt{[typewriter]} entre crochets sont utilisés. Un point d'interrogation entre crochets signifie que l'on n'est pas sûr du mot qui précède.

\noindent Cette édition est provisoire. Les remarques, commentaires et corrections sont bienvenus. Envoyer un message à :

\noindent\textbf{maltsin@math.jussieu.fr}

\hfill G. Maltsiniotis\hspace*{2.5em}

\clearpage

% ===================== page 3 =====================
\printedpage{1}
% head: A. Grothendieck : Dérivateurs, Ch. XI, Édition des Universités de Montpellier II et Paris VII
% foot: Version 28 mai 2002

\noindent\texttt{[Il n'y a pas de pages 1-126]}

\msp{page 127}

\section*{1\quad $\Hot$-fibrations dans $\Cat$}

\noindent\textbf{Théorème 1.} \textit{Soit $f : X \longrightarrow Y$ dans} $\Cat$\textit{. Conditions équivalentes :}

\dgrblock{dgrb450}{The universal meaning of a Hot-fibration}
\begin{dcomment}{exposition}{p127:Conditions équivalentes :}{The universal meaning of a Hot-fibration}
Grothendieck wants the ordinary pullback of categories to calculate a homotopy pullback whenever one of its sides belongs to a recognisable class. This is what makes a fibre a homotopically meaningful object, rather than an accident of a presentation. Write \(BC=|NC|\), where \(N\) is the nerve, and let \(W_\infty\) consist of the functors inducing weak homotopy equivalences of classifying spaces. Condition a) says that
\[
 f^*:\Cat/Y\longrightarrow\Cat/X,\qquad
 (T\longrightarrow Y)\longmapsto(X\times_Y T\longrightarrow X),
\]
preserves \(W_\infty\). The quantifier ranges over maps between \emph{arbitrary} objects of \(\Cat/Y\), and not only over weak equivalences with target \(Y\). It therefore includes every further change of base.

In the terminology of \citet[\S2, Proposition 2.7]{Rezk1998Sharp}, these are precisely the sharp maps in the Thomason model category. The identification uses the definition of sharpness and the right properness of that model category; it does not identify this class with the smaller class of its model fibrations. Every strict pullback square along such a map is homotopy cartesian. This interpretation also explains why composition and arbitrary base change will be available without choosing liftings.

There is a small but decisive quantifier in a') and a''). The point \(e\longrightarrow Y'\) is arbitrary; it need not select the final, respectively initial, object whose existence is assumed. Over the interval \(0\longrightarrow1\), both endpoint inclusions must therefore become weak equivalences after pullback. Asking only for the distinguished final endpoint, for every category having a final object, leads instead to the later notion of preproperness. Asking only over intervals is weaker again. These distinctions become the subject of Chapter XII, and they should not be erased by replacing every occurrence of “fibration” with a single modern definition.
\end{dcomment}


\begin{itemize}[leftmargin=3em]
\item[\textit{a)}] \textit{Pour tout $Y'' \xrightarrow{p} Y'$ sur $Y$, d'où $X'' = X \times_Y Y'' \xrightarrow{q} X' = X \times_Y Y'$, si $p$ est une} $\Hot$\textit{-équivalence, $q$ l'est aussi.}

\item[\textit{a$'$)}] \textit{Comme a), avec $Y'' = e$, $Y'$ ayant un objet final.}

\item[\textit{a$''$)}] \textit{Comme a), avec $Y'' = e$, $Y'$ ayant un objet initial} (\footnote{garder a$''$) et b$''$), conditions duales de a$'$), b$'$) \texttt{[a$''$) et b$''$) étaient barrés]}.})\textit{.}

\item[\textit{b$'$)}] \textit{Pour tout $y \in \Ob Y$, $X_y \longrightarrow X/y$ est une} $\Hot$\textit{-équivalence, pour tout $y \longrightarrow y'$ dans $Y$, $X/y \longrightarrow X/y'$ est une} $\Hot$\textit{-équivalence; enfin ces conditions sont encore vérifiées pour $f' : X' \longrightarrow Y'$, pour tout changement de base $Y' \longrightarrow Y$ dans} $\Cat$\textit{. (En fait, il suffit de le demander pour la première des deux conditions, concernant $X_y \longrightarrow X/y$.)}

\item[\textit{b$''$)}] \textit{Dual de b$'$), avec $X_y \longrightarrow y\backslash X$ et $y'\backslash X \longrightarrow y\backslash X$.}

\item[\textit{c)}] \textit{Pour tout faisceau d'ensembles localement constant $F$ sur $X$, $f_*(F)$ est localement constant et se calcule `fibre par fibre', de même pour $R^1 f_*(F)$ si $F$ est un faisceau en groupes, et pour les $R^i f_*(F)$ si $F$ est un faisceau abélien. De plus, on veut que ces conditions restent satisfaites après tout changement de base.} (\footnote{\textit{c$'$) (version homologique). 1$^o$) en termes d'images directes de groupoïdes localement constants, et 2$^o$) en termes des $R^i f_!(L)$} \texttt{[plutôt $L_i f_!(L)$]}\textit{, pour $L$ faisceau abélien localement constant.}})
\end{itemize}

\noindent\textsc{Démonstration}. On a les implications $\pm$ tautologiques
\[
\mathrm{a} \implies \mathrm{a}' \implies \mathrm{b}' \implies \mathrm{c},
\]
et il reste à prouver que c $\implies$ a. Comme a est autoduale, étant équivalente à a$'$, b$'$, elle est aussi équivalente aux conditions duales de celles-ci, a$''$ et b$''$ (cf. cor. 2, page 75). On est ramené à prouver que si c) est satisfait, et si on a
\[
\begin{tikzcd}[row sep=large, column sep=large]
X \arrow[d, "f"'] & X' \arrow[l, "q"'] \arrow[d, "f'"] \\
Y & Y' \arrow[l, "p"]
\end{tikzcd}
\]
\printedpage{2}
% head: A. Grothendieck : Dérivateurs, Ch. XI, Édition des Universités de Montpellier II et Paris VII
% foot: Version 28 mai 2002
cartésien dans $\Cat$, si $p$ est une $\Hot$-équivalence, $q$ aussi. On utilise le critère d'Artin-Mazur.

\dgrblock{dgrb451}{Comma categories connect the geometric and coefficient tests}
\begin{dcomment}{exposition}{p127:On utilise le critère d'Artin-Mazur.}{Comma categories connect the geometric and coefficient tests}
The comma categories in b') serve two purposes at once: they express a pointwise right Kan extension, and they supply canonical transition maps. We use \(X/y=(f\downarrow y)\), with objects \((x,f(x)\longrightarrow y)\). Its strict fibre \(X_y\) is the full subcategory on identity arrows. A morphism \(u:y\longrightarrow z\) gives a genuine functor \(u_*:X/y\longrightarrow X/z\), so the coherence needed for transport is already present before localization.

Here is the short argument hidden in a') \(\Rightarrow\) b'). Apply a') over \(Y/y\), which has final object \((y,1_y)\). The inclusion of the strict fibre over that object is \(X_y\longrightarrow X/y\). Next work over \(Y/z\) and select its object \((y,u)\). The resulting weak equivalence \(X_y\longrightarrow X/z\) is the composite of \(X_y\longrightarrow X/y\) and \(u_*\). Two-out-of-three proves that \(u_*\in W_\infty\). The same argument works after every base change.

The parenthesis in b') retains the transition condition on the original base; it says that only the fibre-to-comma condition needs to be imposed universally. Indeed, after \(v:Y'\longrightarrow Y\), the comparison
\[
 X'/y'\longrightarrow X/v(y')
\]
is a weak equivalence: both its source and target receive a weak equivalence from the same strict fibre \(X_{v(y')}\). Comparing this square for an arrow of \(Y'\) with the corresponding arrow of \(Y\) then proves the required transition condition on \(Y'\).

For a locally constant coefficient \(F\), the value of the derived direct image at \(y\) is computed on \(X/y\). Thus the preceding weak equivalences identify it with the cohomology of \(X_y\), and identify its transition maps with isomorphisms. In degrees zero and one this means sections and torsors; for abelian coefficients it means every derived degree. The use of local coefficients is essential: untwisted integral homology alone cannot recognise the fundamental group or its action on higher information. The homological form of the categorical Whitehead criterion is stated explicitly in \citet[Proposition 2.4]{Thomason1980}; the cohomological formulation used here is the criterion to which Grothendieck gives the name Artin--Mazur.
\end{dcomment}


% ===================== page 4 (printed 2) =====================

\noindent 1$^o$) Si $F$ est un faisceau d'ensembles localement constant sur $X$, il faut prouver que $H^0(X,F) \xrightarrow{\sim} H^0(X', q^*F)$.

\msp{page 128}

\noindent Mais
\[
H^0(X,F) \simeq H^0(Y,\ \underbrace{f_*(F)}_{\substack{\text{localement constant}\\ \text{par hypothèse c)}}}\ )
\qquad
\underbrace{\xrightarrow{\qquad\sim\qquad}}_{\substack{\text{car } p\\ \text{Hot-équivalence}}}
\qquad
H^0(Y', p^* f_*(F)),
\]
mais
\[
p^* f_*(F) \xrightarrow{\ \sim\ } f'_*(q^*(F)),
\]
car $f_*(F)$ et $f'_*(q^*F)$ se calculent `fibre par fibre' par hypothèse. Donc
\[
H^0(X,F) \simeq H^0(Y', f'_*(q^*(F))) \simeq H^0(X', q^*(F)),
\]

\hfill q.e.d.

\noindent 2$^o$) Si $F$ est un faiscau en groupes localement constant sur $X$, il faut prouver % [sic] faiscau -> faisceau, misspelling
\[
H^1(X,F) \xrightarrow{\ \sim\ } H^1(X', q^*F).
\]
Le fait que c'est injectif résulte déjà de 1$^o$), qui implique que le foncteur
\[
\begin{array}{rcl}
F\text{-torseurs sur } X & \longrightarrow & q^*(F)\text{-torseurs sur } X' \\
P & \longmapsto & f^*P % [sic] f^*P -> q^*P, the functor is the inverse image by q from X to X'
\end{array}
\]
est pleinement fidèle, par l'argument bien connu utilisant les faisceaux $\mathcal{H}$om. Il faut voir que ce foncteur est essentiellement surjectif.

\msp{page 129}

\noindent On a un homomorphisme de suites exactes d'ensembles ponctués
\[
\begin{tikzcd}[column sep=large, row sep=large,
  every matrix/.append style={name=m},
  execute at end picture={
    \coordinate (lb) at ([xshift=0.4em]m-2-3.east);
    \coordinate (lt) at ([xshift=1.8em]m-2-3.east);
    \draw (lt |- m-1-3.north) -- (lb |- m-1-3.north) -- (lb |- m-2-3.south) -- (lt |- m-2-3.south);
    \coordinate (rb) at ([xshift=0.6em]m-2-4.east);
    \coordinate (rt) at ([xshift=-0.6em]m-2-4.east);
    \draw (rt |- m-1-3.north) -- (rb |- m-1-3.north) -- (rb |- m-2-3.south) -- (rt |- m-2-3.south);
  }]
H^1(Y, f_*F) \arrow[r, "u"] \arrow[d, "\alpha"] & H^1(X,F) \arrow[r, "v"] \arrow[d, "\beta"] & \overbrace{H^0(Y, R^1 f_*(F))}^{\ni\, \eta} \arrow[r, "w"] \arrow[d, "\gamma"] & H^2(Y, f_*F) \arrow[d, "\delta"] \\
H^1(Y', f'_*F) \arrow[r, "u'"] & \underbrace{H^1(X', F')}_{\ni\, \xi'} \arrow[r, "v'"] & \underbrace{H^0(Y', R^1 f'_*(F'))}_{\ni\, \eta'} \arrow[r, "w'"] & H^2(Y', f'_*F') % [sic] H^1(Y', f'_*F) -> H^1(Y', f'_*F'), the other entries of the lower row are in F'
\end{tikzcd}
\]
(\footnote{où on a posé $F' = q^*(F)$.}), où les faisceaux $f_*F$, $R^1 f_*(F)$ sont localement constants par hypothèse et se calculent fibre par fibre, les faisceaux $f'_*(F')$ et $R^1 f'_*(F')$ aussi, qui implique que ces derniers sont
\printedpage{3}
% head: A. Grothendieck : Dérivateurs, Ch. XI, Édition des Universités de Montpellier II et Paris VII
% foot: Version 28 mai 2002
canoniquement isomorphes (par les flèches de changement de base) aux images inverses de leurs homologues $f_*(F)$, $R^1 f_*(F)$. Comme $p$ est une $\Hot$-équivalence, cela implique que les flèches $\alpha$, $\gamma$ sont \texttt{[des]} isomorphismes. Donc si $\xi' \in H^1(X', F')$, son image $\eta'$ dans $H^0(Y', R^1 f'_*(F'))$ provient d'un $\eta \in H^0(Y, R^1 f_*(F))$ bien déterminé. Il faut voir que cet $\eta$ provient d'un $\xi \in H^1(X,F)$. Ceci (si je me rappelle bien) est possible si et seulement si une certaine obstruction $w(\eta)$ dans $H^2(Y, f_*(F))$ est dans le sous-ensemble neutre du $H^2$. \textbf{NB} Le champ sur $Y^\wedge$ des $F$-torseurs sur les $X^\wedge/U$ qui remontent $\eta$ est une $\mathcal{F} = f_*F$-gerbe $\mathcal{G}_\eta$, et il faut voir qu'elle est neutre. Mais il y a mieux à faire :

\dgrblock{dgrb452}{Torsor descent and the band of the obstruction}
\begin{dcomment}{exposition}{p129:Mais il y a mieux à faire}{Torsor descent and the band of the obstruction}
The difficult part of the coefficient argument is essential surjectivity for torsors. Full faithfulness is already secure: if \(P,Q\) are \(F\)-torsors, the sheaf \(\mathcal I=\underline{\operatorname{Isom}}_F(P,Q)\) is locally constant, and
\[
 \operatorname{Hom}_F(P,Q)=H^0(X,\mathcal I).
\]
An equivariant map between two torsors is invertible, and the degree-zero calculation identifies these Hom sets after pullback.

The displayed obstruction sequence should, however, be read as a provisional sketch. For a general nonabelian \(F\), the gerbe of lifts of a section \(\eta\) of \(R^1f_*F\) need not have the fixed band of \(f_*F\). Its band is obtained from automorphisms of local representatives of \(\eta\), hence involves inner forms of \(F\). For a direct check, take \(f:BC_2\to e\), constant \(F=S_3\), and the torsor represented by the inclusion of a transposition subgroup \(C_2\hookrightarrow S_3\). Its automorphism group is \(C_{S_3}(C_2)=C_2\), whereas \(f_*F=S_3\). The component of this torsor therefore cannot have the fixed band \(S_3\). Nor does nonabelian \(H^2\) behave as a group with a distinguished zero: neutrality is a condition on a gerbe. The notation \(G\) in the next paragraph is not defined. These are reasons to replace that portion of the argument, not reasons to reject the theorem. The classical framework is the theory of bands and gerbes in \citet[Chapter IV]{Giraud1971}; we give a direct repair that avoids choosing an obstruction class.

For each \(y\), let \(\mathcal T(y)\) be the groupoid of torsors under the restriction of \(F\) to \(X/y\). Restriction along \(X/y\longrightarrow X/z\) makes \(\mathcal T\) a groupoid-valued contravariant diagram on \(Y\). We claim that every transition functor is an equivalence. Essential surjectivity follows from the local constancy of \(R^1f_*F\): its value at \(y\) is precisely the set of isomorphism classes of \(\mathcal T(y)\). To prove full faithfulness for a transition \(u:y\longrightarrow z\), take \(P,Q\in\mathcal T(z)\). Their isomorphism sheaf is a locally constant set-valued coefficient on \(X/z\). Apply the universally quantified degree-zero hypothesis to the base change
\[
 f_z:X/z\longrightarrow Y/z.
\]
The values at \((z,1_z)\) and \((y,u)\) are respectively
\(\operatorname{Hom}(P,Q)\) and
\(\operatorname{Hom}(u^*P,u^*Q)\).
Local constancy makes the restriction map bijective. This is the place where coefficients on a base-changed total category, rather than only coefficients originally chosen on \(X\), are needed.

The same hypotheses identify \(\mathcal T(y)\) with the groupoid of \(F|_{X_y}\)-torsors. On isomorphism classes this is the fibre formula for \(R^1(f_y)_*F\), applied over \(Y/y\). On Hom sets it is the fibre formula for \((f_y)_*\mathcal I\), for the isomorphism sheaf just used. Thus the identification is an equivalence of groupoids, not merely a bijection on \(H^1\).

Now let \(p:Y'\longrightarrow Y\) be a weak equivalence, and form the cartesian square of the source. The natural comparison from \(p^*\mathcal T\) to the torsor diagram \(\mathcal T'\) on \(Y'\) is an objectwise equivalence: both values at \(y'\) restrict equivalently to the torsors on the same strict fibre \(X_{p(y')}\). Finally,
\[
 \operatorname{Tors}_F(X)
 \simeq \operatorname*{holim}_{y\in Y^{\mathrm{op}}}\mathcal T(y).
\]
This is ordinary descent for torsors, expressed on the representable objects \(Y/y\). Explicitly, a compatible family of torsors on these slices determines its torsor over \(x\in X\) by evaluation at \((x,1_{f(x)})\); the comparison maps along \(f(a)\), for \(a:x\longrightarrow x'\), give its restriction maps. The descent cocycle is exactly their composition law. Pulling a torsor on \(X\) to all slices gives the inverse construction.

Homotopy limits of locally constant groupoid diagrams are invariant under a weak equivalence of indexing categories. One can express this as invariance of sections under equivalence of fundamental \(2\)-groupoids, or obtain it by applying the homotopy-invariance theorem for locally constant model-category diagrams to groupoids and their derived right adjoint of constant diagrams; see \citet[Theorem 1.3]{Cisinski2009LocallyConstant}. Consequently
\[
 \operatorname{Tors}_F(X)\longrightarrow
 \operatorname{Tors}_{q^*F}(X')
\]
is an equivalence. This proves the missing essential surjectivity and the required bijection on \(H^1\). It also supplies the local-constancy argument for the lifting gerbe that the source compresses into a parenthesis.
\end{dcomment}


% ===================== page 5 (printed 3) =====================

\msp{page 130}

\noindent On note que cette gerbe $\mathcal{G}$ est localement constante (puisqu'elle correspond à un lien $f_*(G)$ localement constant) et que $\mathcal{G}'$ est son image inverse. Mais comme $p$ est une $\Hot$-équivalence, il s'ensuit que le foncteur
\[
\underline{\Gamma}(Y, \mathcal{G}) \longrightarrow \underline{\Gamma}(Y', \mathcal{G}')
\]
est une équivalence de catégories, donc $\xi'$ provient (à isomorphisme près) d'un unique $\xi$ de $\underline{\Gamma}(Y, \mathcal{G})$, q.e.d.

\noindent 3$^o$) Soit $F$ faisceau abélien localement constant sur $X$, il faut voir \texttt{[que]}
\[
H^i(X,F) \xrightarrow{\ \sim\ } H^i(X', F') \qquad (F' = q^*F), \tag*{$(*)$}
\]
pour tout $i \in \mathbf{N}$. Mais les deux termes sont l'aboutissement des suites spectrales de Leray, de termes initiaux
\[
\begin{aligned}
E_2^{pq} &\simeq H^p(Y, R^q f_*(F)) \\
E^{\prime pq}_2 &\simeq H^p(Y', R^q f'_*(F')),
\end{aligned}
\]
et $(*)$ est associé à un homomorphisme de suites spectrales, qui pour les termes initiaux est
\[
H^p(Y, R^q f_*(F)) \xrightarrow{\sim} H^p(Y', R^q f'_*(F)), \tag*{$(**)$} % [sic] R^q f'_*(F) -> R^q f'_*(F'), cf. the E'_2 term just above
\]
associé à
\[
p^*(R^q f_*(F)) \longrightarrow R^q f'_*(F'). \tag*{$(*{*}*)$}
\]
Mais par hypothèse les $R^q f_*(F)$, $R^q f'_*(F')$ sont localement constants et se calculent fibre par fibre,

\msp{page 131}

\noindent d'où il résulte que les homomorphismes de changement de base $(*{*}*)$ sont \texttt{[des]} isomorphismes, et comme il s'agit de faisceaux localement constants et que $p$ est une $\Hot$-équivalence, il s'ensuit que $(**)$ est \texttt{[un]} isomorphisme, donc aussi $(*)$ sur les aboutissements, q.e.d.

\dgrblock{dgrb453}{The Leray comparison has finite diagonals}
\begin{dcomment}{exposition}{p131:donc aussi $(*)$ sur les aboutissements, q.e.d.}{The Leray comparison has finite diagonals}
Grothendieck can now return to abelian coefficients, where the formal argument is both shorter and complete. For each \(q\), the canonical comparison
\[
 p^*R^qf_*F\longrightarrow R^qf'_*q^*F
\]
is an isomorphism because its values are the cohomology groups of identical strict fibres. Both sides are locally constant. Since \(p\) is a weak equivalence, it therefore induces an isomorphism on every \(E_2^{p,q}\) term of the Leray spectral sequences.

There is no convergence obstruction here. These are first-quadrant cohomological spectral sequences, and in total degree \(n\) only the \(n+1\) pairs \(p+q=n\) occur. Comparison gives isomorphic associated graded groups of finite filtrations of \(H^n(X,F)\) and \(H^n(X',q^*F)\); induction through those filtrations gives the asserted isomorphism. No boundedness of the cohomological dimension of \(X\) or \(Y\) is required. The missing prime on \(F\) in the printed \(E_2\)-comparison is an edition slip, retained in the transcription and recorded separately. Together with the torsor argument, this completes the coefficient criterion.
\end{dcomment}


\noindent\textbf{Définition 1.} Sous les conditions équivalentes du théorème 1, on dit que $f$ est une $\Hot$-\textit{fibration.}

% ---- from Der-ChXI-890_p006-008.tex
\printedpage{4}
% head: A. Grothendieck : Dérivateurs, Ch. XI, Édition des Universités de Montpellier II et Paris VII
% foot: Version 28 mai 2002

\textbf{Proposition 1.} \textit{Cette notion est stable par changement de base (trivial), par composition (formel sur a)). Pour que $f$ soit une \textup{Hot}-fibration, il faut et il suffit que $f$ \textup{\texttt{[plutôt}} $f^{o}$\textup{\texttt{]}} le soit.}

\dgrblock{dgrb454}{From strict fibres to homotopy fibres}
\begin{dcomment}{later}{p131:Cette notion est stable par}{From strict fibres to homotopy fibres}
The stability assertions now have a geometric interpretation. For composable maps \(X\longrightarrow Y\longrightarrow Z\), pulling a weak equivalence back along the composite is the succession of the two pullbacks. Opposite categories preserve strict pullbacks, and reversing a nerve simplex identifies the classifying spaces of \(C\) and \(C^{\mathrm{op}}\). Thus both composition and self-duality follow directly from a), without another coefficient argument.

The strict fibre also calculates the homotopy fibre. By b'), the categories \(X/y\) vary by weak equivalences along all arrows of \(Y\), and \(X_y\longrightarrow X/y\) is a weak equivalence. The dual form of Quillen's Theorem B therefore identifies
\[
 BX_y \simeq \operatorname{hofib}_y(BX\longrightarrow BY).
\]
The exact citation is \citet[\S1, Theorem B and its corollary]{Quillen1973}; its hypotheses concern the comma-category transition functors. Merely knowing the homotopy types of the strict fibres, without the comparison to comma categories, would not suffice.

This perspective explains the force of the theorem while leaving room for Grothendieck's chosen constructions. \citet[Proposition 4.4]{Cisinski2009LocallyConstant} later proves that weak equivalences of nerves together with the maps that are both smooth and proper form a category of fibrant objects, with functorial factorization. These smooth-and-proper maps are a particularly useful subclass of the present Hot-fibrations. The statement does not say that all Hot-fibrations are smooth and proper. Their advantages are the explicit categories of lifts and the simultaneous control of left and right homotopy Kan extensions.
\end{dcomment}


\textbf{Proposition 2.} \textit{Soit $f : X \longrightarrow S$ une \textup{Hot}-fibration. Alors $s \longmapsto \mathrm{hot}(X_s)$ est un système local sur $S$, à valeurs dans \textup{Hot}. Donc les $X_s$ pour $s$ dans une même composante connexe de $S$ sont \textup{Hot}-isomorphes entre eux (mais pas canoniquement, si cette composante n'est pas \textup{1}-connexe).}

\dgrblock{dgrb455}{A group extension and its monodromy}
\begin{dcomment}{examples}{p131:système local sur $S$, à valeurs dans}{A group extension and its monodromy}
The local system of fibre homotopy types already contains nontrivial information in a finite group example. Regard a group as its one-object category. If
\[
 1\longrightarrow K\longrightarrow G
 \xrightarrow{\varphi}Q\longrightarrow1
\]
is exact, the functor from the one-object groupoid of \(G\) to that of \(Q\) is an isofibration. Its nerve is a Kan fibration. Pullbacks of weak equivalences along Kan fibrations are weak equivalences, and the nerve preserves pullbacks, so this functor is a Hot-fibration. Its strict fibre is the one-object groupoid \(K\).

The comma category over the unique object of \(Q\) has objects \(q\in Q\). A morphism from \(q\) to \(q'\) is an element \(g\in G\) satisfying \(q=q'\varphi(g)\), with a harmless reversal if the other convention for composition is chosen. Surjectivity of \(\varphi\) makes this a connected groupoid; the automorphism group of any object is \(K\). Thus the inclusion of the strict fibre is an equivalence of categories. If \(\varphi\) is not surjective, its connected components are the corresponding cosets of \(\varphi(G)\), and that inclusion is not a weak equivalence. This computes precisely why an arbitrary homomorphism of groups need not be a Hot-fibration.

For the nontrivial monodromy, take \(G=\mathfrak S_3\), \(K=\mathfrak A_3=C_3\), and \(Q=C_2\), with \(\varphi\) the sign. A lift of the nonidentity element of \(C_2\) is a transposition. Conjugation by it sends a three-cycle to its inverse. A different lift changes this automorphism by an inner automorphism of \(K\); inner automorphisms induce naturally isomorphic functors on the one-object groupoid and hence the same map in \(\Hot\). The action on the fibre is consequently well defined there. It is nontrivial, since inversion acts nontrivially on \(\pi_1(BC_3)=C_3\).

With constant coefficient \(\mathbb F_3\) on \(B\mathfrak S_3\), the degree-one fibre cohomology is
\[
 H^1(C_3;\mathbb F_3)
 =\operatorname{Hom}(C_3,(\mathbb F_3,+))
 \cong\mathbb F_3,
\]
and the base generator acts by \(-1\). Its invariant subgroup is zero. The low-degree Leray sequence gives \(H^1(\mathfrak S_3;\mathbb F_3)=0\): the other term \(H^1(C_2;\mathbb F_3)\) is zero, and positive-degree cohomology of \(C_2\) with these coefficients vanishes by the averaging operator \(\frac12(1+\tau)\). Independently, the same answer follows from \(\mathfrak S_3^{\mathrm{ab}}=C_2\). The fibre cohomology is locally constant on the base, but it is not a constant coefficient with trivial monodromy. This is the distinction that the printed qualifier about noncanonical identifications is meant to preserve; a nonsimply connected base can also carry trivial monodromy, as a product shows.
\end{dcomment}


% (a vertical change bar is printed in the left margin beside Théorème 2)
\textbf{Théorème 2.} {\itshape Soit
\begin{tikzcd}[ampersand replacement=\&, baseline=(\tikzcdmatrixname-1-1.base), column sep=small, row sep=small]
X' \arrow[rr, "f"] \arrow[dr, "p'"'] \& \& X \arrow[dl, "p"] \\
\& S \&
\end{tikzcd}
un diagramme commutatif dans \textup{Cat}, avec $p$, $p'$ des \textup{Hot}-fibrations, et soit $S_0 \subseteq S$ une partie de $S$ telle que $S_0 \longrightarrow \pi_0(S)$ soit surjectif. Pour que $f$ soit une \textup{Hot}-équivalence, il faut et il suffit que les $f_s : X_s \longrightarrow Y_s$ $(s \in S_0)$ le soient.} % [sic] f_s : X_s -> Y_s -> f_s : X'_s -> X_s, the diagram has X' and X (cf. f_s : X'_s -> X_s p. 6)

a) La suffisance se prouve comme tantôt, avec le critère de Artin-Mazur.

\msp{page 132}

Notons que l'hypothèse sur $S_0$ assure que si $F$, $F'$ sont deux faisceaux d'ensembles localement constants sur $S$, et $u : F \longrightarrow F'$ un homomorphisme tel que $u_s : F_s \xrightarrow{\ \sim\ } F'_s$ pour $s \in S$, alors $u$ est \texttt{[un]} isomorphisme. % [sic] s \in S -> s \in S_0, otherwise the hypothesis on S_0 plays no role

1$^{o}$) Soit $F$ un faisceau d'ensembles localement constant sur $X$, $F'$ son image inverse sur $X'$, il faut prouver
\[
H^0(X,F) \ \xrightarrow{\ \sim\ }\ H^0(X',F').
\]
Or
\[
H^0(X,F) \simeq H^0(Y, \underbrace{f_*(F)}_{\substack{\text{localement}\\ \text{constant sur } S}}) \ \longrightarrow\ H^0(Y, f'_*(F')) \simeq H^0(X',F'),
\] % [sic] H^0(Y, f_*(F)), H^0(Y, f'_*(F')) -> H^0(S, p_*(F)), H^0(S, p'_*(F')), the brace says "localement constant sur S" and (*) uses p_*, p'_*
il faut prouver que
\[
p_*(F) \ \xrightarrow{\ \sim\ } p'_*(F') \tag*{$(*)$}
\]
est \texttt{[un]} isomorphisme. Or par hypothèse sur $p$, $p'$, les deux faisceaux sont localement constants et se calculent fibre par fibre, d'où résulte que $(*)$ est un isomorphisme pour les fibres en $s \in S_0$, donc un isomorphisme partout. OK.

2$^{o}$) Il en résulte déjà, si $F$ est un faisceau en \textit{groupes,} que le foncteur
\[
(F\text{-torseurs sur } X) \ \longrightarrow\ (F\text{-torseurs sur } X')
\]
est pleinement fidèle, donc $H^1(X,F) \longrightarrow H^1(X',F')$ \texttt{[est]} injectif; il reste à prouver que ce foncteur est essentiellement surjectif. À cause de la pleine fidélité, c'est \textit{local sur $S$} (\textbf{NB} On voit directement que l'hypothèse fait sur les $s \in S_0$ implique la même hypothèse pour tout $s \in S$, en utilisant le fait que $p$, $p'$ sont des fibrations) (\footnote{Les hypothèses faites sont bel et bien locales sur $S$ (du moins si $S_0 = S$).}). Donc on est ramené au cas où $S = S/s_0$ a un objet final $s_0$.

\dgrblock{dgrb456}{Testing one fibre in each component}
\begin{dcomment}{exposition}{p132:est pleinement fidèle, donc}{Testing one fibre in each component}
Theorem 2 is the fibrewise Whitehead theorem for these categorical fibrations. Its fibre map is \(f_s:X'_s\longrightarrow X_s\); the displayed \(X_s\longrightarrow Y_s\) is an edition slip. Transport is natural with respect to \(f\), so if \(f_s\) is a weak equivalence at one point of a component then it is a weak equivalence at every point of that component. The choice of \(S_0\) therefore loses no information.

For sufficiency, the source compares direct images of local coefficients. Its degree-zero comparison should be written
\[
 H^0(X,F)=H^0(S,p_*F)\longrightarrow
 H^0(S,p'_*F')=H^0(X',F').
\]
Both direct images are locally constant, and their values on \(S_0\) are isomorphic by the fibre hypothesis. A morphism between locally constant diagrams is an isomorphism everywhere if it is so at one object of each component: propagate the property along an unoriented chain of arrows, whose transition maps are invertible. This proves degree zero and, by the isomorphism-sheaf argument, full faithfulness on torsors.

The local reduction for essential surjectivity is legitimate once it is applied to these groupoids of torsors, rather than only their sets of isomorphism classes. On each representable \(S/s\), both total categories still satisfy the hypotheses. The torsor diagrams constructed above are locally constant there, and \(S/s\) has a final object. Their groupoids of global sections are thus equivalent to their values at that object. Fibrewise equivalence supplies local preimages of any torsor on \(X'\). Full faithfulness uniquely lifts the overlap isomorphisms between these preimages; their cocycle follows from its validity after pullback and faithfulness. Effective descent for torsors gives the global preimage. This justifies both the phrase “local on \(S\)” and the degree-one edge identifications on manuscript page 133. The abelian comparison follows from Leray as before.

For necessity there is a shorter modern proof, complementary to the printed path-category argument. The preceding block identifies both strict fibres with homotopy fibres. An equivalence \(BX'\longrightarrow BX\) over \(BS\) induces equivalences on all homotopy pullbacks along \(\{s\}\longrightarrow BS\), so every \(Bf_s\) is a weak equivalence. Conversely, maps of spaces over a fixed base are detected on their homotopy fibres: on each base component the long exact homotopy sequences give the higher homotopy isomorphisms, and the action of \(\pi_1(BS,s)\) on the components of the fibre identifies the components of the total space with its orbit set. A fibrewise equivalence preserves that action and those orbit sets. This also treats disconnected or empty fibres, which a statement only about based groups would overlook.
\end{dcomment}


\printedpage{5}
% head: A. Grothendieck : Dérivateurs, Ch. XI, Édition des Universités de Montpellier II et Paris VII
% foot: Version 28 mai 2002

\msp{page 133}

On a alors
\[
\begin{aligned}
H^1(X,F) &\ \simeq\ H^0(S, R^1p_*(F)) \\
H^1(X',F') &\ \simeq\ H^0(S, R^1p'_*(F')),
\end{aligned}
\]
et on voit encore que
\[
R^1p_*(F) \ \longrightarrow\ R^1p'_*(F')
\]
est \texttt{[un]} isomorphisme (car il se calcule fibre par fibre, et on utilise l'hypothèse sur les $f_s$), donc induit un isomorphisme
\[
H^1(X,F) \ \longrightarrow\ H^1(X',F').
\]

3$^{o}$) Si $F$ est même abélien, on doit prouver
\[
H^i(X,F) \ \xrightarrow{\ \sim\ }\ H^i(X',F')
\]
pour tout $i$, et cela résulte encore de la suite spectrale de Leray pour $(p,F)$ et pour $(p',F')$.

b) Inversement, supposons que $f$ soit une Hot-équivalence, et prouvons que les $f_s : X'_s \longrightarrow X_s$ sont des Hot-équivalences. Je vais utiliser le

% (a vertical change bar is printed in the left margin beside the Lemme)
\textbf{Lemme.} \textit{Soit $S \in \Ob\Cat$, $S$ \textup{0}-connexe. Alors il existe un $\overline{S} \in \Ob\Cat$, $\overline{S}$ \textup{Hot}-asphérique, et $\overline{S} \longrightarrow S$ une \textup{Hot}-fibration.}

Voyons comment ceci implique le théorème 2 (partie réciproque). On peut supposer $S$ 0-connexe, et on fait le changement de base par $\overline{S} \longrightarrow S$

\msp{page 134}

comme dans le lemme :
\[
\begin{tikzcd}[column sep=small, row sep=small]
X' \arrow[rr, "f"] \arrow[dr, "p'"'] & & X \arrow[dl, "p"] & & & & \\
& S & & & \overline{X'} \arrow[rr, "\overline{f}"] \arrow[dr, "\overline{p'}"'] & & \overline{X} \arrow[dl, "\overline{p}"] \\
& & & & & \overline{S} \arrow[ullll, "\varphi"] &
\end{tikzcd}
\quad .
\]

Alors $\overline{p'}$ et $\overline{p}$ sont des Hot-fibrations, et comme $\varphi$ est une Hot-fibration, le changement de base par $\varphi$ transforme la Hot-équivalence $f$ en une Hot-équivalence $\overline{f}$. Considérons alors $\overline{s} \in \Ob\overline{S}$, et le diagramme
\[
\begin{tikzcd}
\overline{X'} \arrow[r, "\overline{f}"] & \overline{X} \\
\overline{X'}_{\overline{s}} \arrow[u, hook, "i'"] \arrow[r, "\overline{f}_{\overline{s}}"] & \overline{X}_{\overline{s}} \arrow[u, hook, "i"']
\end{tikzcd}
\qquad .
\]

\printedpage{6}
% head: A. Grothendieck : Dérivateurs, Ch. XI, Édition des Universités de Montpellier II et Paris VII
% foot: Version 28 mai 2002

Comme $\{\overline{s}\} \longrightarrow \overline{S}$ est une Hot-équivalence, la flèche $i$ qui s'en déduit par le changement de base $\overline{p} : \overline{X} \longrightarrow \overline{S}$ l'est aussi, puisque $\overline{p}$ est une Hot-fibration, et de même pour $i'$. Ainsi dans le carré commutatif précédent, $\overline{f}$, $i$ et $i'$ sont des Hot-équivalences, donc aussi $\overline{f}_{\overline{s}}$. Mais si $s = \varphi(\overline{s})$, $\overline{f}_{\overline{s}}$ s'identifie à $f_s : X'_s \longrightarrow X_s$. D'ailleurs, il est immédiat que $\varphi$, étant une Hot-fibration

\msp{page 135}

à source non-vide, et à but connexe, est surjective. (\textbf{NB} toutes les fibres d'une même composante connexe de $S$ sont Hot-équivalentes, donc si l'une est non-vide, toutes le sont.) Cela achève la démonstration, modulo le lemme.

\dgrblock{dgrb457}{A contractible replacement without a derivator}
\begin{dcomment}{exposition}{p135:Cela achève la démonstration, modulo le lemme.}{A contractible replacement without a derivator}
The lemma is not asserting the existence of a contractible covering space: a general connected homotopy type has no such covering. It asks for a fibration whose fibre retains the entire loop homotopy type. The source's based path category is designed to do exactly this. The dependency it leaves open at the end of the chapter concerns the proof that this specific categorical path construction is both proper and smooth.

For the existential lemma alone, Thomason's theory already gives an independent proof. Choose \(s_0\in S\) and factor its inclusion in the Thomason model category as
\[
 e\xrightarrow{\ i\ }E\xrightarrow{\ p\ }S,
 \qquad i\text{ a trivial cofibration},\quad
 p\text{ a model fibration}.
\]
Then \(BE\) is weakly contractible. Since \(\operatorname{Ex}^2N\) preserves pullbacks and sends \(p\) and all its base changes to Kan fibrations, right properness of simplicial sets proves that \(p\) is a Hot-fibration in the sense of a). The existence of the factorization is \citet[Theorem 4.9]{Thomason1980}; its weak equivalences and fibrations are Definitions 2.1--2.2, and properness is Corollary 5.5.

The map \(p\) is surjective on objects when \(S\) is nonempty and connected. The object \(i(*)\) supplies one nonempty fibre. Pull back over an arrow of \(S\); both endpoint fibres have the homotopy type of that pullback, so one is nonempty if and only if the other is. A zigzag in \(S\) reaches every object of its connected component. This is the elementary justification of the surjectivity used on manuscript page 135. It gives the printed proof of the converse in Theorem 2 without invoking any derivator: pull back along \(E\longrightarrow S\), use that the inclusions of its fibres are weak equivalences, and apply two-out-of-three.

This replacement does not diminish the purpose of Chapter XII. The abstract model factorization proves existence, whereas Grothendieck needs constructions with explicit categories of lifts, duality, and control of both direct images. The later smooth-and-proper factorization quoted above addresses that stronger requirement.
\end{dcomment}


Il reste à prouver celui-ci. Comme $S = \emptyset$ \texttt{[plutôt} $S \neq \emptyset$\texttt{]}, je peux choisir une ‘origine’
\[
s_0 \in S,
\]
et regarder la catégorie
\[
\Ch_{\infty}(S; s_0, -)
\]
des chemins infinis dans $S$, d'origine $s_0$, fibre de $\Ch_{\infty}(S)$ par le foncteur-source
\[
\sigma : \Ch_{\infty}(S) \ \longrightarrow\ S.
\]
On sait que les fibres de $\sigma$ sont Hot-asphériques (VII), et la catégorie construite ici n'est autre que (avec les notations de VII, p.~152) que la catégorie $\overline{S'}$ qui entre dans la factorisation canonique de
\[
f : \{s_0\} = S' \ \longrightarrow\ S
\]
en
\[
\{s_0\} = S' \ \underbrace{\longrightarrow}_{\text{Hot-équivalence}}\ \overline{S'} \ \xrightarrow{\ \overline{f}\ }\ S,
\]
où ici $\overline{f}$ s'interprète simplement comme le foncteur but (i.e.\ le foncteur induit sur

\msp{page 136}

$\Ch_{\infty}(S; s_0, -)$ par le foncteur but
\[
\beta : \Ch_{\infty}(S) \ \longrightarrow\ S \ \ ).
\]
On a prouvé dans loc.\ cit.\ que $\overline{f}$ est une Hot-fibration (‘propre’ et ‘lisse’), mais dans un contexte qui semblait faire appel tacitement à la théorie du dérivateur HOT. Donc pour mettre les résultats ‘topologiques’ de loc.\ cit.\ sur une base solide, il faut développer les notions pertinentes de foncteur propre, foncteur lisse, sans référence à une théorie HOT. C'est ce que j'avais d'ailleurs fait déjà en 1983, du temps de Pursuing Stacks, mais sans rédaction en forme à ce sujet. Il me faut donc à présent revenir sur ces notions, de façon un peu systématique. D'ailleurs, j'aurais dû mal à construire un seul foncteur Hot-fibrant (sauf les exemples triviaux, tels les isomorphismes, ou les produits $Y \times Z \xrightarrow{\mathrm{pr}_1} Y$), si je ne disposais d'une théorie des foncteurs lisses et des foncteurs propres (\footnote{J'exagère un peu quand même -- après tout il doit être assez évident que les foncteurs de Thomason-Kan sont des Hot-fibrations (mais il y en a encore un peu trop peu !).}). % [sic] dû mal -> du mal, "avoir du mal à" (partitive article, no accent)

\dgrblock{dgrb458}{Why the next chapter starts again}
\begin{dcomment}{context}{p136:Il me faut donc à présent revenir sur ces notions}{Why the next chapter starts again}
Grothendieck's return to proper and smooth functors is an effort to remove a dependency, not a rejection of the homotopy theory just obtained. A path construction had been used to produce fibrations while some of its asserted properties still appealed to the derivator that those fibrations were meant to construct. Categories of lifts offer a direct combinatorial foundation for those properties.

The historical background is visible in \emph{Pursuing Stacks}, \S69: the coefficient theory is to remember categories of coherent diagrams together with both adjoints of restriction, rather than merely a single homotopy category \citep{Grothendieck1983}. The letter to Thomason of 2 April 1991 makes a related distinction explicit: the weak equivalences carry the intrinsic homotopy theory, while classes of fibrations and cofibrations make its operations calculable \citep[\S1]{Grothendieck1991Letter}. We refer to the commentary on Chapter~I, manuscript p.~126 (“The later letter and the meaning of a fibration”), for the broader historical discussion.

The next pages isolate two different issues. One is to describe the maps along which categorical pullback behaves well. The other is to decide how much of the usual homotopy theory is forced by the proposed axioms on \(W\). The revisions of terminology in Chapter XII are necessary because an axiom suitable for the full homotopy type can be too strong for a localization that deliberately forgets higher homotopy groups.
\end{dcomment}
\dgrchapter{XII}{Chapter XII}{\texorpdfstring{Caractérisation de $W_\infty$. Foncteurs $W$-propres, $W$-lisses etc. Sommes amalgamées et carrés $W$-cocartésiens dans Cat}{Caractérisation de W∞. Foncteurs W-propres, W-lisses etc. Sommes amalgamées et carrés W-cocartésiens dans Cat}}
\dgrpart{0}{XI-XIIa}{Sections 1--4, manuscript pages 1--58}{XI-XIIa\enspace\textperiodcentered\enspace Chapter XII, sections 1--4\enspace\textperiodcentered\enspace manuscript pages 1--58}{Chapter XII \textperiodcentered\ \textup{XI-XIIa}}
% ---- from Der-ChXIIa-012_p001-005.tex
\printedpage{--}
% head: (none; no running head and no page number printed on this page)

\begin{center}
{\Large \textsf{LES DÉRIVATEURS}}

\bigskip

\textbf{ALEXANDRE GROTHENDIECK}

\bigskip\bigskip

\textbf{Chapitre XII}

\medskip

\textbf{Caractérisation de $W_\infty$. Foncteurs $W$-propres,\\
$W$-lisses etc. Sommes amalgamées et carrés\\
$W$-cocartésiens dans $\Cat$.}
\end{center}

\clearpage

\printedpage{--}
% head: (none; no running head and no page number printed on this page)

Ce texte a été déchiffré et transcrit en \LaTeXe{} par M. Künzer. Il a été édité par M. Künzer, J. Malgoire, et G. Maltsiniotis. La transcription est aussi fidèle que possible au manuscrit. Pour les quelques corrections évidentes, ou rares commentaires des éditeurs, ainsi que pour la numérotation originale des pages du manuscrit, les caractères de machine à écrire \texttt{[typewriter]} entre crochets sont utilisés. Un point d'interrogation entre crochets signifie que l'on n'est pas sûr du mot qui précède.

Cette édition est provisoire. Les remarques, commentaires et corrections sont bienvenus. Envoyer un message à :

\textbf{maltsin@math.jussieu.fr}

\hfill G. Maltsiniotis\hspace*{2em}

\clearpage

\printedpage{1}
% head: A. Grothendieck : Dérivateurs, Ch. XII, Édition des Universités de Montpellier II et Paris VII
% foot: Version 7 décembre 2002

\msp{page 1}

\hfill 19.11.90

\section*{I\quad Caractérisation de $W_\infty$.}

Je voudrais caractériser $W_\infty$ en tant que partie de $\mathrm{Fl}(\Cat)$, par ses propriétés comme localiseur, et sans faire appel à la connaissance de notions utilisant déjà $W_\infty$, comme la propreté etc. Je voudrais tout au moins des propriétés sur un localiseur $W$, familières pour $W_\infty$, qui soient si fortes qu'elles impliquent $W \subseteq W_\infty$. Bien sûr, dans les démonstrations je me permets d'utiliser tout ce que je sais sur $W_\infty$.

Bien sûr, il faudra au moins une forme de saturation, peut-être très faible, je présume le suivant suffira (\footnote{(non, j'ai dû le renforcer en 2 [?]).}).

\begin{itemize}[leftmargin=*]
\item[\textbf{W(1)}] (Axiome de saturation.) a) Les identités sont dans $W$, b) si $f$, $g$ \texttt{[sont]} composables et deux parmi $f$, $g$, $gf$ sont dans $W$, le troisième aussi, c) si $fg \in W$, $gf \in W$, alors $f, g \in W$.
\end{itemize}

Une propriété particulièrement cruciale, qui déjà distingue $W_\infty$ de l'ensemble des homotopismes (notion trop stricte à beaucoup d'égards), est l'axiome que j'ai appelé L4 bis :

\dgrblock{dgrb459}{A contractible nerve without a categorical contraction}
\begin{dcomment}{examples}{p1:Une propriété particulièrement cruciale, qui déjà distingue}{A contractible nerve without a categorical contraction}
Grothendieck's insistence that homotopisms are too restrictive has a concrete explanation. Let \(P\) be the ordered set with objects \(0,1,2,\ldots\) and nonidentity relations
\[
  0<1>2<3>4<5>\cdots .
\]
There are no composable pairs of nonidentity arrows. Its nerve therefore realizes to the half-line subdivided at the integers; in particular \(P\to e\) belongs to \(W_\infty\).

Nevertheless \(P\to e\) is not a homotopism in the sense used on the next manuscript page. Suppose that the identity of \(P\) could be joined to the constant functor at \(c\) by a zigzag of \(m\) natural transformations. Evaluating the zigzag at an integer \(n\) would give a path of at most \(m\) edges from \(n\) to \(c\) in the underlying undirected graph of \(P\). Every arrow changes the integer by at most one. Thus \(|n-c|\leq m\) for every \(n\), which is impossible. A homotopy inverse to \(P\to e\) would select precisely such a constant functor, so no homotopy inverse exists.

This example separates two finite conditions that are easily conflated. Each individual vertex can be connected to \(0\) by a finite path. A natural-transformation contraction would require one finite zigzag working simultaneously at every vertex. The localization axioms allow the classifying space to forget that uniform bound.
\end{dcomment}


\begin{itemize}[leftmargin=*]
\item[\textbf{W(2)}] (Axiome de localisation.) Soit $f : X \fl Y$ une flèche dans $\Cat/S$, telle que $\forall\ s \in \Ob S$, $f_s : X/s \fl Y/s$ soit dans $W$ (on dira que $f$ est dans $W$ `localement sur $S$\,'). Alors $f \in W$.
\end{itemize}

\dgrblock{dgrb460}{The two directions of the characterization}
\begin{dcomment}{exposition}{p1:soit dans $W$ (on dira que $f$ est dans $W$}{The two directions of the characterization}
Grothendieck wants the equivalences of ordinary homotopy theory to emerge from operations on categories. There are two logically different assertions to establish. The first says that the proposed axioms cannot identify too many maps: \(W\subseteq W_\infty\). The second says that a sufficiently coherent theory cannot identify too few: \(W_\infty\subseteq W\). This section attacks the first assertion by repeatedly taking fibres. The later minimal-localizer theorem supplies the second. Keeping these directions separate makes the successive changes of axioms intelligible.

Write \(BA=|NA|\) for the classifying space of a small category, and let \(e\) be the terminal category. Here \(W_\infty\) consists of the functors \(u\) for which \(Bu\) is a weak homotopy equivalence. The expression \(W\)-aspherical means only \(A\to e\in W\). It does not mean that \(A\) has vanishing higher homotopy groups until a comparison with \(W_\infty\) has been established. In particular, it is not the usual topological convention under which a connected \(K(G,1)\) is called aspherical.

The local-to-global axiom W(2) tests a map \(u:X\to Y\) over \(S\) on the comma categories \(X/s\to Y/s\). These are not the strict fibres \(X_s\to Y_s\). An object of \(X/s\) is a pair \((x,p(x)\to s)\), so the test retains all ways of approaching \(s\). W(2 bis) is the absolute theorem-A form: contractibility, in the sense of \(W\), of every \(X/y\) implies \(u\in W\). The distinction matters: the source explicitly uses the relative form for Proposition~3, whereas the loop argument and the final nondegeneracy repair need only the absolute form.

The commentary on Chapter~VI, manuscript pp.~44--45, explained why a fibre axiom begins to select ordinary homotopy theory. The present pages give the mechanism: any nonzero homotopy group of a \(W\)-aspherical category can be lowered, by looping, to a forbidden nontrivial set of components.
\end{dcomment}


\texttt{[Les axiomes W(1) etc. sont appelés aussi} \textcircled{1} \texttt{etc. dans l'original.]}

Dans le cas $Y = S$, cela donne

\textbf{Corollaire :} \textit{Si $f : X \fl Y$ est $W$-asphérique, i.e. $y\backslash X \fl y\backslash Y$} \texttt{[plutôt $X/y \fl Y/y$]} \textit{dans $W$ pour tout $y \in \Ob Y$, alors $f \in W$} (\footnote{On peut considérer la validité du corollaire comme une variante affaiblie W(2 bis) de W(2) ci-dessus. Dans la proposition 1, on n'utilise que W(2 bis), forme faible de W(2).}).

\textbf{Proposition 1.} \textit{Si $W$ satisfait} W(1) \textit{et} W(2)\textit{, les propriétés suivantes sur $W$ sont équivalentes :}

\begin{itemize}[leftmargin=*]
\item[\textit{a)}] \textit{Pour tout $X$ dans $\Cat$, $X \times \Dl^1 \xrightarrow{\pr_1^X} X$ est dans $W$ (i.e. $\Dl^1 \fl e$ est `universellement dans $W$\,').}

\item[\textit{b)}] \textit{Si $X$ dans $\Cat$ a un objet final, alors $X$ est $W$-asphérique, i.e. $X \fl e$ est dans $W$.}

\item[\textit{c)}] \textit{Tout homotopisme est dans $W$.}
\end{itemize}

\clearpage

\printedpage{2}
% head: A. Grothendieck : Dérivateurs, Ch. XII, Édition des Universités de Montpellier II et Paris VII
% foot: Version 7 décembre 2002

\textsc{Démonstration.} b) $\Longrightarrow$ a) (\footnote{n'utilise W(1) que sous forme faible, sans la condition de saturation c). NB Seul b) $\Longrightarrow$ c) utilise W(2 bis), les implications a) $\Longrightarrow$ c) $\Longrightarrow$ b) sont valables sous la seule hypothèse W(1).}), car pour voir que $\pr_1 \in W$, en appliquant le corollaire ci-dessus, il faut prouver que $\forall\ x \in X$
\[ (X \times \Dl^1)/x \fl X/x \tag*{$(*)$} \]
est dans $W$. Or
\[ (X \times \Dl^1)/x \simeq X/x \times \Dl^1, \]
\[
\begin{tikzcd}
\Dl^1 \arrow[d] & X \times \Dl^1 \arrow[l] \arrow[d] & X/x \times \Dl^1 \arrow[l] \arrow[d] \\
e & X \arrow[l] & X/x \arrow[l]
\end{tikzcd}
\]
\texttt{[la flèche verticale à droite est ajoutée]} et les deux catégories $X/x$, $X/x \times \Dl^1$ sont $W$-asphériques car elles ont un objet final, donc $(*)$ est dans $W$ par saturation faible W(1).

a) $\Longrightarrow$ c). Je rappelle que
\[ f : X \fl Y \]
\msp{page 2}
est un homotopisme s'il existe
\[ g : Y \fl X \]
telle que
\[ gf \sim \id_X, \quad fg \sim \id_Y, \tag*{$(*)$} \]
où la relation $(\sim)$ désigne l'homotopie dans $\uHom(X,X)$ et $\uHom(Y,Y)$ respectivement (définie par chemins allant d'un élément à l'autre). Prouvons le

\textbf{Lemme.} \textit{Supposons que $W$ satisfait} W(1) \textit{et le} a) \textit{de la proposition, i.e. $X \times \Dl^1 \fl X$ dans $W$ pour tout $X$ dans} $\Cat$. \textit{Alors si $X$, $Y$} \texttt{[sont]} \textit{dans} $\Cat$ \textit{et $f, g$ dans} $\uHom(X,Y)$ \textit{sont homotopes, on a}
\[ \hot_W(f) = \hot_W(g). \]
\textit{De plus, $f \in W$ si et seulement si $g \in W$.}

\dgrblock{dgrb461}{Saturation at the point where it is used}
\begin{dcomment}{exposition}{p2:si et seulement si $g \in W$.}{Saturation at the point where it is used}
The proof is carefully arranged so that equality after localization is not mistaken for membership in \(W\). For a natural transformation \(f\Rightarrow g\), let \(F:X\times[1]\to Y\) be its corresponding functor and let \(p:X\times[1]\to X\) be the projection. If \(p\in W\), its sections \(i_0,i_1\) lie in \(W\) by two-out-of-three. Since \(pi_0=pi_1=\id_X\), their images in \(\Cat[W^{-1}]\) are both \(p^{-1}\). Thus \(f=Fi_0\) and \(g=Fi_1\) become equal there. Independently, two-out-of-three gives
\[
 f\in W\quad\Longleftrightarrow\quad F\in W
 \quad\Longleftrightarrow\quad g\in W.
\]
Iteration handles a finite zigzag of transformations.

Now suppose \(f\) has a categorical homotopy inverse \(g\). The preceding argument puts both \(gf\) and \(fg\) in \(W\). The separate clause W(1c) then puts \(f\) and \(g\) in \(W\). This last step would not follow merely from the existence of their inverse images in the localization unless strong saturation had already been proved.

Three usages of saturation should therefore remain distinct. Two-out-of-three is W(1b); the two-cycle condition is W(1c); strong saturation says that every map inverted by the localization already belongs to \(W\). The modern weak-saturation axiom used for fundamental localizers asks, in addition to identities and two-out-of-three, only that \(ri=\id\) and \(ir\in W\) imply \(r\in W\) \citep[Definition~1.1.1]{Cisinski2004Minimal}. W(1c) implies this condition. Consequently the later results apply to the source's axioms without silently inserting strong saturation into this elementary proof.
\end{dcomment}


Il suffit de le prouver quand il existe un homomorphisme $f \fl g$. Cela peut s'interpréter par l'existence d'une flèche
\[ X \times \Dl^1 \xrightarrow{F} Y \]
dont les restrictions à $X \times \{0\}$ et $X \times \{1\}$ respectivement soient $f$, $g$. Pour prouver que $\hot_W(f) = \hot_W(g)$, il suffit de prouver que
\[ i_0, i_1 : X \fl X \times \Dl^1 \qquad x \longmapsto (x,0),\ x \longmapsto (x,1) \]
\printedpage{3}
% head: A. Grothendieck : Dérivateurs, Ch. XII, Édition des Universités de Montpellier II et Paris VII
% foot: Version 7 décembre 2002
sont égales dans $\Hot_W$, puisque
\[ f = F \circ i_0, \quad g = F \circ i_1. \]
Or cela résulte du fait que ce sont des sections de $X \times \Dl^1$ sur $X$, et que dans $\Hot_W$ la projection $X \times \Dl^1 \xrightarrow{p} X$ est \texttt{[un]} isomorphisme (donc $i_0 = i_1 = p^{-1}$ dans $\Hot_W$). De plus, le diagramme
\[
\begin{tikzcd}
X \arrow[r, "i_0"] \arrow[dr, "f"'] & X \times \Dl^1 \arrow[d, "F"] \\
& Y,
\end{tikzcd}
\]
où $i_0 \in W$, montre que $f \in W$ si et seulement si $F \in W$, et de même $g \in W$ si et seulement si $F \in W$, donc $f \in W \Longleftrightarrow g \in W$.

Revenons à $(*)$, on voit que dans $\Hot_W$ $g \circ f = \id_X$, $f \circ g = \id_Y$, donc $f$ et $g$ sont des $\Hot_W$-isomorphismes inverses l'un de l'autre. D'autre part, $gf \in W$, $fg \in W$, puisque $\id_X, \id_Y \in W$. Donc on gagne, si on a formulé W(1) de façon assez forte. On a prouvé donc a) $\Longrightarrow$ c).

c) $\Longrightarrow$ b) est trivial, puisque si $X$ a un objet final, $X \fl e$ est un homotopisme.

\msp{page 3}

On doit bien sûr par la suite supposer que $W$ satisfait aux conditions équivalentes de la proposition. On prendra p. ex. la forme

\begin{itemize}[leftmargin=*]
\item[\textbf{W(3)}] (Axiome de l'élément final.) Si $X$ a un objet final, alors $X$ est $W$-asphérique.
\end{itemize}

\dgrblock{dgrb462}{Terminal objects, initial objects, and duality}
\begin{dcomment}{exposition}{p3:alors $X$ est $W$-asphérique}{Terminal objects, initial objects, and duality}
The terminal-object axiom gives a usable supply of trivial categories. If \(t\) is terminal in \(A\), the natural transformation \(\id_A\Rightarrow \operatorname{const}_t\) makes \(A\to e\) a homotopism; an initial object gives the transformation in the reverse direction. Proposition~1 therefore derives initial-object asphericity from the same axioms. No second axiom is being assumed.

For a functor \(u:A\to B\), the category \(B/b\) has the terminal object \((b,\id_b)\). Hence two-out-of-three identifies the two formulations
\[
 A/b\longrightarrow B/b\in W
 \quad\Longleftrightarrow\quad A/b\longrightarrow e\in W .
\]
This is why the corrected slice orientation in the source's first corollary is essential: it determines whether the subsequent notion is asphericity or coasphericity.

Proposition~2 next turns this one-sided language into a self-dual theory. The bridge is the twisted-arrow category \(\Phi(A)\), with comparison functors to \(A\) and \(A^{\mathrm{op}}\). Their comma categories have explicit contractions, so absolute theorem A puts both comparisons in \(W\). Naturality then compares \(u\), \(\Phi(u)\), and \(u^{\mathrm{op}}\) by two-out-of-three. The construction and its variance were treated in the commentary on Chapter~I, manuscript pp.~73--75, and are not repeated here. What is new in the present argument is its economy of hypotheses: the comparisons use W(2 bis), not the full relative W(2). Applying the same construction to \(u^{\mathrm{op}}\) gives the reverse implication, so \(u\in W\) if and only if \(u^{\mathrm{op}}\in W\).
\end{dcomment}


\textbf{Corollaire 1.} \textit{Pour que $f : X \fl Y$ soit $W$-asphérique, il faut et il suffit que les $X/y$ soient $W$-asphériques} ($y \in \Ob Y$).

Car pour que
\begin{tikzcd}[baseline=(\tikzcdmatrixname-1-1.base), column sep=small, row sep=small]
X/y \arrow[rr] \arrow[dr] & & Y/y \arrow[dl] \\
& e &
\end{tikzcd}
soit dans $W$, comme $Y/y$ est $W$-asphérique (objet final), il faut et il suffit que $X/y$ le soit aussi.

\textbf{Corollaire 2.} \textit{Si $X$ dans} $\Cat$ \textit{a un objet initial, il est $W$-asphérique.} (En effet, $X \fl e$ est un homotopisme, et on applique prop. 1.) (\footnote{C'est donc l'énoncé dual de l'axiome W(3) de l'objet final.})

Je mettrai pour mémoire

\textbf{Corollaire 3.} \textit{Tout homotopisme dans} $\Cat$ \textit{est dans $W$. Si $f, g : X \rightrightarrows Y$ dans} $\Cat$ \textit{sont homotopes, alors elles sont égales dans} $\Hot_W$\textit{, et $f$ est dans $W$ si et seulement si $g$ l'est.}

\textbf{Proposition 2.} Soit $f : X \fl Y$ dans $\Cat$. Si $f \in W$, $f^\circ \in W^\circ$ \texttt{[plutôt $f^\circ \in W$].}

% ---- from Der-ChXIIa-012_p006-010.tex
\printedpage{4}
% head: A. Grothendieck : Dérivateurs, Ch. XII, Édition des Universités de Montpellier II et Paris VII
On utilise la démonstration de I, p. 75, à l'aide de la catégorie nommée $\Phi(X)$ ($\mathrm{Ob}\,\Phi(X) = \mathrm{Fl}(X)$, mais avec une notion de flèche dans $\Phi(X)$ un peu particulière). On trouve que $f \in W \iff \Phi(f) \in W$, et de même $f^\circ \in W \iff \Phi(f) \in W$.

NB. Dans la démonstration de la proposition 2, comme dans celle de la proposition 1 (pour l'implication b) $\Longrightarrow$ a)), on n'a utilisé la propriété de localisation W(2) que sous la forme faible du corollaire, i.e. sous la forme $f$ $W$-asphérique $\Longrightarrow f \in W$.

\textbf{Corollaire} (de la prop. 2) : \textit{Soit $W \subseteq \mathrm{Fl}(\mathrm{Cat})$ satisfaisant l'axiome} W(1) (saturation). \textit{Alors la conjonction des axiomes} W(2), W(3) \textit{équivaut à celle des axiomes duaux} W($2'$), W($3'$) (\footnote[5]{Itou si on remplace dans cet énoncé W(2) par W(2 bis).}).

Il suffit de prouver que 1,2,3 $\Longrightarrow$ $2'$,$3'$.

\msp{page 4}

Pour $2'$ \texttt{[plutôt} $3'$\texttt{]}, c'est le corollaire 2 ci-dessus. Pour $3'$ \texttt{[plutôt} $2'$\texttt{]}, cela résulte aussitôt de la proposition 2. (NB Cet énoncé resterait valable, si on aurait pris W(2) sous la forme faible seulement, cf. le NB page précédente.)

Je dois encore mettre un axiome de la forme L5, L6 (cf. VI, p. 43-44). Mais l'énoncé de ces axiomes utilise la notion de Hot-fibrés, ou de foncteurs parfaits, i.e. fait appel à la connaissance de $W_\infty$, ce que j'aimerais éviter (si je peux). Je dois donc introduire la notion de $W$-fibration, sur le modèle de ce que j'ai fait dans le contexte des dérivateurs, dans I (cf. p. 111-126). On introduit la notion de \textbf{D}-fibration à la page 117 (prop. 5), par \texttt{[des]} conditions équivalentes, dont certaines ont un sens en termes de la donnée de $W$ sans plus. Soit donc
\[ f : X \longrightarrow Y \]
dans Cat et considérons les conditions

\dgrblock{dgrb463}{What the four fibration tests quantify over}
\begin{dcomment}{exposition}{p4:Je dois encore mettre un axiome de la forme}{What the four fibration tests quantify over}
Grothendieck is looking for a notion of family whose strict pullbacks compute homotopy-invariant changes of base. His strongest condition says that
\[
  f^*:\Cat/Y\longrightarrow\Cat/X
\]
preserves the underlying \(W\)-maps. The intermediate condition tests this only between \(W\)-aspherical bases. The left-hand test asks that, for every \(v:T\to Y\) whose domain \(T\) has a terminal object and every object \(t\in T\), the inclusion of the strict fibre
\[
   X_{v(t)}\longrightarrow X\times_Y T
\]
belong to \(W\). The chosen \(t\) need not be the terminal object. The right-hand test uses bases with an initial object. This universal quantification over \(T\to Y\) distinguishes the definition from an isolated comparison \(X_y\to X/y\).

The arrows in the source's diagram follow at once because a terminal or initial category is \(W\)-aspherical. Their converses express additional mathematics. For \(W_\infty\), Chapter~XI supplies them. For a general \(W\), the step from a one-sided test to invariance under every \(W\)-map is precisely the proposed axiom W(4). One should not import it from the formalism of a model category: an arbitrary categorical functor is not a fibration of a model structure merely because it satisfies one of these tests.

The eventual terminology is settled only at XII, manuscript pp.~51--52. The early adjective \emph{ultraforte} names preservation of all \(W\)-maps, and it is useful to retain that descriptive meaning while reading these pages. The later \(W\)-fibration is defined using categorical homotopisms in the base. In a truncated theory the distinction survives: the path family over an \((n+1)\)-sphere passes the ordinary homotopy tests but does not preserve \(W_n\)-equivalences. The complete computation is given below at manuscript pp.~16--17.
\end{dcomment}


\begin{enumerate}
\item[a)] Si $Y'' \longrightarrow Y'$ dans $\mathrm{Cat}/Y$ est dans $W$, alors $X'' \longrightarrow X'$ aussi (le changement de base par $X \longrightarrow Y$ transforme $W$ dans $W$).
\item[a$'$)] Comme a), mais avec $Y''$, $Y'$ $W$-asphériques.
\item[a$_1$)] Comme a), dans le cas particulier où $Y'$ a un objet final et $Y''$ est la catégorie ponctuelle.
\item[a$_2$)] Dual de a$_1$), avec objet initial de $Y'$.
\end{enumerate}

On a donc
\[
\begin{tikzcd}[row sep=large, column sep=1.4em]
 & & \mathrm{a} \arrow[d, Rightarrow] & \text{$W$-fibration ultraforte} & \\
 & & \mathrm{a}' \arrow[dl, Rightarrow] \arrow[dr, Rightarrow] & \text{$W$-fibration forte} & \\
\text{$W$-fibration gauche} & \mathrm{a}_1 & & \mathrm{a}_2 & \text{$W$-fibration droite}
\end{tikzcd}
\]
% foot: Version 7 décembre 2002
\printedpage{5}
% head: A. Grothendieck : Dérivateurs, Ch. XII, Édition des Universités de Montpellier II et Paris VII
et dans le cas $W = W_\infty$ toutes ces conditions sont équivalentes. J'ai envie \texttt{[de]} prendre la plus forte de ces conditions, mais je dois me donner des moyens pour trouver des $W$-fibrations en un sens aussi fort. Par exemple j'aurais envie de savoir que si $X$
\msp{page 5}
est dans Cat, alors $\Ch_\infty(X) \longrightarrow X \times X$ ou $\Ch(X) \longrightarrow X \times X$ est une $W$-fibration. Dans les deux cas, et (disons) pour $W_\infty$, je le prouve en prouvant que $p$ est à la fois lisse, et propre. Il faudrait donc revoir les notions de lissité et de propreté dans le contexte de la donnée d'un $W$.

Si je me rappelle bien, la propriété a$'$) \texttt{[plutôt a}$_1$\texttt{)]} était quelque chose d'assez aisément vérifiable, il suffisait de supposer que les
\[ X_y \longrightarrow X/y \qquad \text{et} \qquad X/y \longrightarrow X/y' \]
(pour $y \in \mathrm{Ob}\,Y$ resp. $u : y \longrightarrow y'$ dans $\mathrm{Fl}\,Y$) sont dans $W$, et que ça reste vrai par tout changement de base - ce qui revient à la condition a$_1$) sur $f$. On pourrait donc formuler une condition sur $W$ ainsi :

\begin{enumerate}[leftmargin=*]
\item[\textbf{W(4)}] (Axiome des $W$-fibrations (\footnote[6]{Voir variante affaiblie suffisante, dans le NB p. 10. Voir encore [?] de variantes pages 13, 14 ci-dessous.}).) Soit $f : X \longrightarrow Y$ dans Cat satisfaisant la condition de $W$-fibration a$_1$) ci-dessus (p. 4), alors elle satisfait même à a), i.e. le changement de base par $f$ transforme flèches dans $W$ en flèches dans $W$ (\footnote[7]{i.e. une $W$-fibration à gauche est une fibration ultraforte.}).
\end{enumerate}

Grâce à la proposition 2, il s'ensuit que la condition duale a$_2$) sur $f$ implique la condition de $W$-fibration forte, donc l'axiome W(4), en apparence asymétrique, signifie aussi que les quatre conditions de fibration sur $f$ sont équivalentes, il est donc en fait symétrique.

\textbf{NB} Je ne vois pas, sans axiomes supplémentaires sur $W$, comment passer de a$_1$ à a$'$ seulement. La difficulté, c'est de montrer, moyennant a$_1$ seulement, que si $Y$ est $W$-asphérique, l'inclusion $X_y \longrightarrow X$ (où $y \in \mathrm{Ob}\,Y$) est dans $W$. Dans le contexte des dérivateurs \textbf{D}, cela se voyait par voie \textbf{D}-cohomologique.

Cet axiome très fort suffit à impliquer que si $X$, $Y$ sont $W$-fibrés sur $S$, et si $f : X \longrightarrow Y$ est un $S$-morphisme qui est dans $W$, alors $f_s : X_s \longrightarrow Y_s$ est dans $W$ ($\forall\ s \in S$). Montrons déjà au moins l'inverse.

\msp{page 6}

\textbf{Proposition 3.} \textit{Soient $X$, $Y$ dans} Cat$/S$\textit{, $f : X \longrightarrow Y$ une $S$-flèche. Supposons ($X$, $Y$ $W$-fibrés) que $\forall\ s \in S$, $f_s : X_s \longrightarrow Y_s$ soit dans $W$, alors $f \in W$} (\footnote[8]{\textit{Sans} axiome W(4). C'est la première fois qu'on a besoin de W(2) sous sa forme forte, à l'exclusion de W(2 bis).}).

Il s'impose d'utiliser l'axiome de localisation, donc de regarder les $f/s$ induits par $f$
% foot: Version 7 décembre 2002
\printedpage{6}
% head: A. Grothendieck : Dérivateurs, Ch. XII, Édition des Universités de Montpellier II et Paris VII
\[
\begin{tikzcd}[row sep=large]
X/s \arrow[r, "f/s"'] & Y/s \\
X_s \arrow[u, "i_s^X"] \arrow[r, "f_s"] & Y_s \arrow[u, "i_s^Y"'] \quad .
\end{tikzcd}
\]

\dgrblock{dgrb464}{From strict fibres to the total category}
\begin{dcomment}{exposition}{p6:Il s'impose d'utiliser l'axiome de localisation}{From strict fibres to the total category}
Proposition~3 is a descent statement, and the comma categories supply the descent data. Suppose \(p:X\to S\) and \(q:Y\to S\) satisfy the left fibration test and \(u:X\to Y\) lies over \(S\). Pulling \(p\) back to \(S/s\), with terminal object \((s,\id_s)\), shows that the canonical inclusion \(X_s\to X/s\) is in \(W\); similarly for \(Y\). The square in the text therefore gives
\[
 u_s\in W\quad\Longleftrightarrow\quad u/s\in W .
\]
If this holds at every \(s\), the relative axiom W(2) gives \(u\in W\). The dual proof applies when both maps satisfy the right test. The maps must satisfy a common variance; a mixed pair is not covered by this argument.

There is also a direct reason why one representative per connected component suffices. For an arrow \(a:s\to t\), the left fibration condition makes \(X/s\to X/t\) and \(Y/s\to Y/t\) weak equivalences. Comparing the two horizontal maps in the resulting square shows that \(u_s\in W\) exactly when \(u_t\in W\). Propagate this statement along a finite zigzag of arrows joining the chosen representative to any other object.

None of these steps uses W(4). That distinction is the point of the proposition: fibrewise equivalences imply a total equivalence by local-to-global descent, while a total equivalence need not be detectable on fibres without a stronger hypothesis on the theory. The path-space example below will show the failure sharply.
\end{dcomment}


La condition a$_1$ suffit à impliquer que $i_s^X$ et $i_s^Y$ sont dans $W$, donc $f/s$ est dans $W$ si et seulement si $f_s$ l'est. On n'a pas eu à utiliser l'axiome W(4), et il suffisait de supposer que $X$, $Y$ soient fibrés ‘$W$-cohomologiques’ (ou $W$-fibrés à gauche) sur $S$, ou dualement, des $W$-fibrés à droite.

De plus :

\textbf{Corollaire} (\footnote[9]{Sans axiome W(4).}). \textit{Soit $S_0 \subseteq \mathrm{Ob}\,S$ une partie de $S$ telle que $S_0 \longrightarrow \pi_0(S)$ surjectif. Alors dans la proposition 3 il suffit de supposer $f_s \in W$ pour $s \in S_0$.}

Il suffit de voir que si $u : s \longrightarrow s'$ est une flèche dans $S$, alors $f_s \in W \iff f_{s'} \in W$. On utilise le fait que $f_s \in W \iff f/s \in W$, et de même pour $f_{s'}$. Donc à voir que $f/s \in W \iff f/s' \in W$. Or on a
\[
\begin{tikzcd}[row sep=large]
X/s \arrow[r, "f/s"] \arrow[d, "X/u"'] & Y/s \arrow[d, "Y/u"] \\
X/s' \arrow[r, "f/s'"] & Y/s' \quad ,
\end{tikzcd}
\]
où $X/u$ et $Y/u$ sont dans $W$, d'où le résultat.

Mais revenons à la réciproque : si $f \in W$, alors a-t-on $f_s \in W_S$? \texttt{[plutôt} $f_s \in W$\texttt{?]} Il y a un argument dans ce sens dans VI, p. 23, que je vais reprendre. L'idée est celle-ci. On considère la composante connexe $Y_0$ de $Y$ \texttt{[plutôt} $S_0$ \texttt{de} $S$\texttt{]} contenant $s$, et on construit un
\msp{page 7}
$W$-fibré (au \textit{sens fort} a))
\[ S' \longrightarrow S_0 \subseteq S \]
avec $S'$ \textit{$W$-asphérique.} Comme $X \longrightarrow Y$ sur $S$ est dans $W$, il en sera donc de même de $X' \longrightarrow Y'$ sur $S'$. Comme les fibres de $X'$, $Y'$ en $s'$ s'identifient à celles de $X$, $Y$ en l'image de $s'$ dans $S$, on est ramené au cas où $S$ est $W$-asphérique. Mais alors on a
\[
\begin{tikzcd}[row sep=large]
X \arrow[r, "f"] & Y \\
X_s \arrow[u, "i_s^X"] \arrow[r, "f_s"] & Y_s \arrow[u, "i_s^Y"'] \quad ,
\end{tikzcd}
\]
% foot: Version 7 décembre 2002
\printedpage{7}
% head: A. Grothendieck : Dérivateurs, Ch. XII, Édition des Universités de Montpellier II et Paris VII
où $i_s^X$, $i_s^Y$ sont dans $W$, car déduites de $\{s\} \longrightarrow S$ (qui est dans $W$, $S$ étant $W$-asphérique) par les changements de base $W$-fibrants $X \longrightarrow S$, $Y \longrightarrow S$. Donc $f \in W \implies f_s \in W$.

NB On n'a utilisé la propriété de fibration pour $X$, $Y$ qu'au sens a$'$) (\footnote[10]{Et même, dans a$'$) on pourrait exiger que $Y'' = e$, $Y'$ $W_\omega$-asphérique, où $W_\omega$ \texttt{[est]} défini page 14, $W_\omega \subseteq W \cap W_\infty \ldots$}).

Il faut donc prouver, si on peut, le

\textbf{Lemme} (?) (\footnote[11]{Est prouvé moyennant l'axiome W(4) (+ des petites vérifications ultérieures).}). \textit{Soit $S$ un objet} $0$\textit{-connexe de} Cat\textit{. Alors il existe un objet $W$-asphérique $S'$ de} Cat\textit{, et une $W$-fibration (au sens} fort \textit{de a), page 4) $S' \longrightarrow S$.}

\dgrblock{dgrb465}{A path cover makes the converse local}
\begin{dcomment}{exposition}{p7:une $W$-fibration (au sens}{A path cover makes the converse local}
The converse fibre test needs a cover on which the base becomes aspherical without changing the homotopy type of the fibres. Fix \(s_0\) in a connected component of \(S\), and use the endpoint map
\[
 \beta:\Ch(S;s_0,-)\longrightarrow S .
\]
Its source is intended to model based paths with free endpoint. Every object in the chosen component is the endpoint of a finite zigzag starting at \(s_0\), so this cover reaches each object of that component.

Assume the two properties under construction: \(E=\Ch(S;s_0,-)\) is \(W\)-aspherical and \(\beta\) is ultraforte. If \(u:X\to Y\) over \(S\) is in \(W\), then the pullback \(u_E:X_E\to Y_E\) is in \(W\). For any \(e_0\in E\), the point map \(e\to E\) is in \(W\), by two-out-of-three with \(E\to e\). The strong-aspherical-base condition for \(X_E\) and \(Y_E\) now places both fibre inclusions in \(W\). Their comparison square gives \(u_{\beta(e_0)}\in W\).

This formulation avoids an unnoticed component restriction: one pulls back \(u\) along the actual composite \(E\to S\). The inclusion \(S_0\to S\) of the selected component is a pullback of the point inclusion \(e\to\pi_0S\). A functor to a discrete category is perfect, so this inclusion is perfect as well; W(4) makes it preserve \(W\) under base change. The composite of this inclusion with the perfect path cover has the required property. No prior assertion that restriction of an arbitrary \(W\)-map to one component remains in \(W\) is needed here.

The dependency that remains is real but explicitly acknowledged: asphericity of \(E\) must be proved before the converse fibre theorem is used. Deriving it from the desired conclusion about fibres of \(\Ch(S)\to S\) would be circular. The next paragraph changes the base from \(S\) to the category of path types precisely to avoid that circle.
\end{dcomment}


L'idée naturelle, c'est de prendre, pour $s_0 \in \mathrm{Ob}\,S$ fixé,
\[ S' = \Ch(S; s_0, -) \qquad \text{ou} \qquad S' = \Ch_\infty(S; s_0, -), \]
et la flèche
\[ \beta : c \longmapsto \text{but } c : S' \longrightarrow S. \]
On sait que $S'$ est $W_\infty$-aphérique et que $\beta$ est $W_\infty$-fibrante (et même parfaite). Il faut voir que $S'$ est aussi $W$-asphérique, et $\beta$ $W$-fibrante. Il faut donc revoir la démonstration dans le cas de $W_\infty$, et prouver qu'elle se transforme à $W$, % [sic] aphérique -> asphérique, lettre s omise
\msp{page 8}
moyennant les axiomes qu'il faut sur $W$. J'ai envie de regarder d'abord le choix $S' = \Ch(S; s_0, -)$, sans passage à une limite inductive qui m'obligerait à mettre un axiome de limites sur $W$. Les propriétés homotopiques délicates et essentielles de $\Ch(S)$ sont établies dans VII, p. 157-169, savoir que dans le diagramme
\[ S \xrightarrow{i_S} \Ch(S) \xrightarrow{p_S} S \times S, \]
$i_S \in W_\infty$ et $p_S$ est une $W_\infty$-fibration (en fait, \textit{propre} et \textit{lisse} pour $W_\infty$) (\footnote[12]{Dans l'énoncé des ‘propriétés essentielles’ de $\Ch(S)$, on peut travailler avec $W_\omega$, la plus petite partie de Fl(Cat) satisfaisant les axiomes W(1), W(2 bis), W(3) (cf. page 14), donc $W_\omega \subseteq W \cap W_\infty$.}). Le fait que $i_S \in W$ \texttt{[plutôt} $i_S \in W_\infty$\texttt{]} implique que les foncteurs source et but
\[ \begin{tikzcd} \Ch(S) \arrow[r, shift left, "\sigma"] \arrow[r, shift right, "\beta"'] & S \end{tikzcd} \qquad \sigma = \mathrm{pr}_1 \circ p_S,\ \beta = \mathrm{pr}_2 \circ p_S \]
sont dans $W_\infty$ (puisque $\sigma i_S = \beta i_S = \mathrm{id}_S$). Comme ce sont des Hot-fibrations, il en résulte que les fibres $\Ch(S; s, -)$ resp. $\Ch(S; -, s)$ sont Hot-asphériques. Mais même en admettant que $\sigma$, $\beta$ sont dans $W$ et sont des $W$-fibrations, je ne peux en conclure tout de suite que les fibres sont $W$-asphériques - il y aurait cercle vicieux! Il me faut donc une démonstration directe.

Considérons le foncteur
\[
\begin{tikzcd}[row sep=huge]
\Ch(S; s, -) \arrow[d, "p"'] \\
\underline{Ty} \quad ,
\end{tikzcd}
\]
% foot: Version 7 décembre 2002
\printedpage{8}
% head: A. Grothendieck : Dérivateurs, Ch. XII, Édition des Universités de Montpellier II et Paris VII
c'est un foncteur Cat-fibrant, donc $W$-lisse (je devrai développer la notion de $W$-lissité (\footnote[13]{$f : X \longrightarrow Y$ est $W$-lisse si les $X_y \longrightarrow y\backslash X$ sont $W$-asphériques.}) avec soin, j'admettrai pour le moment les propriétés formelles familières, et les justifierai par la suite). À cause de la $W$-lissité, pour prouver $p \in W$, il suffit de
\msp{page 9}
prouver que les fibres de $p$ sont $W$-asphériques (\footnote[14]{On utilise ici W(2 bis), pas $W(2)$.}). Or les fibres sont les $Z_\tau = \Ch_\tau(S; s, -)$, qui sont des catégories \textit{contractiles} (i.e. $Z_\tau \longrightarrow e$ est un homotopisme), donc $W$-asphériques par W(3) (cf. cor. 3, page 3).

\dgrblock{dgrb466}{Why finite path types avoid a limiting axiom}
\begin{dcomment}{exposition}{p9:prouver que les fibres de $p$ sont}{Why finite path types avoid a limiting axiom}
The projection to path types separates combinatorics from homotopy. For a fixed finite oriented interval \(\tau\), the category \(\Ch_\tau(S;s,-)\) consists of its diagrams in \(S\) whose initial endpoint is \(s\). Replace the last vertex by its predecessor and the last edge by the identity. This defines a functor on the diagram category; the last arrow of each diagram, with identities at the other vertices, defines a natural transformation to or from the identity functor, according to the orientation of the edge. Its image is the subcategory of diagrams with final edge an identity, canonically identified with the diagram category for the shortened interval. Repeat there and include back into the original category. Induction on the number of edges gives a finite zigzag from the identity to the constant-path functor at \(s\). This supplies the fixed-type contraction which Chapter~VII states and later leaves unwritten \citep[manuscript pp.~152, 160, 170]{GrothendieckDerivateursVII}.

The total category over \(\underline{Ty}\) is a Grothendieck fibration. Therefore it is \(W\)-smooth: every required category of lifts with prescribed target has a terminal cartesian lift. Aspherical strict fibres then give a \(W\)-equivalence to \(\underline{Ty}\), by the dual of absolute theorem A. The unit path type is terminal in \(\underline{Ty}\), so W(3) makes this base \(W\)-aspherical as well \citep[manuscript pp.~160--161]{GrothendieckDerivateursVII}. A claim that the fibres alone make the total category aspherical would omit this base hypothesis. Thus the path cover is built from categorical contractions and theorem A, before invoking any stability of filtered colimits.

The same economy explains the lifting calculation for
\(\Ch(S)\to S\times S\). For endpoint arrows \(u:a'\to a\), \(v:b'\to b\) and a path \(c:a\leadsto b\), form the concatenation \(\widetilde c=\operatorname{ch}(u)c\operatorname{ch}(v)^{\mathrm{op}}\). The category of lifts is compared with \(\Ch(S;a',b')/\widetilde c\), which has a terminal object. Chapter~VII constructs both comparison functors and links their composites to the identities by natural transformations \citep[manuscript pp.~161--166]{GrothendieckDerivateursVII}. These are actual categorical homotopies, so W(1)--W(3) put the comparisons in \(W\). Applying the construction to the opposite category supplies the other variance. This establishes properness and smoothness independently of W(4). Only afterwards does W(4) upgrade the family to preservation of arbitrary \(W\)-maps.

The lifting-category comparison is a proof by reference to Chapter~VII; the audit checked the cited construction and retains that dependency. The mapping-path factorization then follows from base-change and composition of perfect maps, with its first map a section of a universally aspherical source projection \citep[manuscript pp.~170--171]{GrothendieckDerivateursVII}. The distinction between finite zigzags and the infinite completion \(\Ch_\infty\), already developed in the commentary on Chapter~VII, manuscript pp.~157--170, explains why the present proof need not assume W(6).
\end{dcomment}


D'autre part, si on admet que $\Ch(S) \xrightarrow{p_S} S{\times}S$ est $W$-lisse (et $W$-propre), donc $\Ch(S) \xrightarrow{\sigma_S} S$ aussi, étant à fibres $W$-asphériques comme on vient de voir, il s'ensuit que $\sigma_S \in W$, d'où $S \xrightarrow{i_S} \Ch(S)$ aussi dans $W$ (puisque $\sigma_S \circ i_S = \mathrm{id}_S$). Donc il reste à prouver seulement que
\[ \Ch(S) \xrightarrow{p_S} S \times S \tag*{$(*)$} \]
est $W$-propre et $W$-lisse. Il en résultera que
\[ \Ch(S; s_0, -) \xrightarrow{\beta} S \]
est lui aussi $W$-propre et $W$-lisse, car déduit de $p_S$ par le changement de base
\[ S \longrightarrow S \times S \qquad s \longmapsto (s_0, s). \]

Il faut ensuite voir que $W$-lisse et $W$-propre implique la propriété de $W$-fibration \textit{sous sa forme la plus forte} a) (alors que je prévois qu'a priori on ne l'aura que sous les formes faibles a$_1$, a$_2$ - pas même sous la forme a$'$). C'est ici donc que l'axiome W(4) des $W$-fibrations sera essentiel (\footnote[15]{sous la forme affaiblie W(4a$''$) (p. 13).})!

Voici l'argument pour la lissité de $(*)$, dans loc. cit. p. 161 ff. Il faut prouver que pour $c \in \mathrm{Ob}\,\Ch(S)$, et une flèche $u \in \mathrm{Fl}(S \times S)$ \texttt{[plutôt} $(u,v) \in$ \texttt{Fl(}$S \times S$\texttt{)]}, une certaine catégorie de ‘relèvements de $u$ de source $c$ ’ \texttt{[plutôt ‘relèvements de} $(u,v)$ \texttt{de but} $c$ \texttt{’]}, $\mathcal{C} = \mathcal{C}(c,u,v)$, est $W$-asphérique (\footnote[16]{En fait, on prouve que $\mathcal{C}(c,u,v)$ est \textit{contractile}, donc elle est $W_\omega$-asphérique.}). Pour ceci, on la compare à une catégorie $\mathcal{C}'$ à objet final, donc contractile et a fortiori $W$-asphérique, à l'aide de deux foncteurs
\[ \begin{tikzcd} \mathcal{C} \arrow[r, shift left, "\varphi"] & \mathcal{C}', \arrow[l, shift left, "\psi"] \end{tikzcd} \]
et des homotopies $\varphi\psi \sim \mathrm{id}_{\mathcal{C}'}$ (sans mérite, $\mathcal{C}'$ étant contractile)
\msp{page 10}
et $\psi\varphi \sim \mathrm{Id}_{\mathcal{C}}$ (on trouve même une flèche $\psi\varphi \longrightarrow \mathrm{id}_{\mathcal{C}}$). Ainsi $\varphi$ est un homotopisme, donc une $W$-équivalence (cor. 3, p. 3), donc $\mathcal{C}$ est $W$-asphérique, $\mathcal{C}'$ l'étant. La propreté de $p_S$ équivaut à la lissité de $p_{(S^\circ)}$, elle est donc prouvée du même coup.
% foot: Version 7 décembre 2002

% ---- from Der-ChXIIa-012_p011-015.tex
\printedpage{9}
% head: A. Grothendieck : Dérivateurs, Ch. XII, Édition des Universités de Montpellier II et Paris VII

Ainsi, le lemme de la page 7 est prouvé, moyennant l'axiome W(4) qui assure qu'une flèche qui est à la fois $W$-lisse et $W$-propre, satisfait à la condition très forte a) des $W$-fibrations.

\textbf{NB} En fait, il suffirait, au lieu de l'axiome W(4) sous sa forme forte, de supposer qu'une $f : X \longrightarrow Y$ qui est à la fois `homotopiquement lisse' et `homotopiquement propre', i.e. les fibres des $X/x \longrightarrow Y/y$ et des $x\backslash X \longrightarrow y\backslash Y$ sont contractiles (donc $W$-asphériques pour tout $W$ satisfaisant W(1,2,3)), soit $W$-fibration forte. Énonçons quand même ce que nous avons prouvé.

\textbf{Proposition 4.} \textit{Si $W$ satisfait} W(1,2,3) (\footnote[17]{Pour W(2), il le suffit sous la forme W(2 bis).}) (nous le supposons toujours par la suite) \textit{plus} W(4) (sous une forme un peu affaiblie suffisante, cf. remarque précédente), \textit{on a ceci: Soit $f : X \longrightarrow Y$ une flèche dans} Cat$/S$\textit{, avec $X$, $Y$ des $W$-fibrés sur $S$. Si $f \in W$, alors $f_s : X_s \longrightarrow Y_s$ est dans $W$ pour tout $s \in S$.}

Je me demande si j'en sais assez pour pouvoir prouver
\[ W \subseteq W_\infty . \]
Il faut quand même éliminer des cas canularesques, comme $W = \mathrm{Fl}(\mathrm{Cat})$, donc mettre l'axiome

\begin{itemize}[leftmargin=4em, labelwidth=3em, labelsep=0.6em]
\item[W(5)] (Axiome de connexité (\footnote[18]{Voir à ce sujet le corollaire page 22.}).) Si $f : X \longrightarrow Y$ est dans $W$, $\pi_0(f) : \pi_0(X) \longrightarrow \pi_0(Y)$ \texttt{[est]} bijectif.
\end{itemize}

\msp{page 11}Je factorise alors $f$ en (\footnote[19]{À nouveau, dans ces arguments on a intérêt à travailler avec $W_\omega$.})
\[ X \xrightarrow{i} \bar{X} \xrightarrow{\bar{f}} Y, \]
avec $\bar{f}$ une $W_\infty$-fibration et une $W$-fibration (\footnote[20]{en fait, $\bar{f}$ est à la fois $W_\infty$-parfait (i.e. $W_\infty$-lisse et propre) et $W$-parfait (en fait, $W_\omega$-parfait).}), $i \in W \cap W_\infty$ - cela marche finallement, % [sic] finallement -> finalement, spelling
par la construction de VII, p. 170-171 (de \textsc{Cartan-Serre}), à partir des propiétés de % [sic] propiétés -> propriétés, missing letter
$Y \longrightarrow \Ch(Y) \longrightarrow Y \times Y$ sans plus. Comme $\bar{f} i = f \in W$ et $i \in W$, on en déduit $\bar{f} \in W$. Comme $i \in W_\infty$, on en conclut $f \in W_\infty \iff \bar{f} \in W_\infty$. Donc pour prouver que $f$ est dans $W_\infty$, on est ramené au cas où $f$ est une $W$-fibration et une $W_\infty$-fibration. Mais alors la proposition 4 nous montre que les fibres de $f$ sont $W$-asphériques (\footnote[21]{on n'a utilisé ici que cette propriété sur les fibres, comme consequence de W(4). (On utilise W(4) sous sa forme la plus faible W(4b$''$), cf. pages 13, 14).}), et il suffit de voir qu'elles sont $W_\infty$-asphériques (proposition 3 (\footnote[22]{on n'utilise la proposition 3 que dans le cas $Y = S$, pour lequel W(2 bis) suffit.}) appliquée à $W_\infty$). Ainsi on est ramené au cas où $Y = e$, i.e. à prouver que si $X$ dans Cat est $W$-asphérique, il est $W_\infty$-asphérique.
% [sic] (footnote 21) consequence -> conséquence, missing accent

Soit $x_0 \in \mathrm{Ob}\, X$, considérons comme tantôt
\[ E = \Ch(X, x_0, -) \xrightarrow{p} X, \]
% foot: Version 7 décembre 2002
\printedpage{10}
% head: A. Grothendieck : Dérivateurs, Ch. XII, Édition des Universités de Montpellier II et Paris VII
on sait que $E$ est $W \cap W_\infty$-asphérique, et que c'est une fibration pour $W$ et pour $W_\infty$. Comme $X$ est $W$-asphérique (comme $E$), $p \in W$, donc à nouveau par la proposition 4 (\footnote[23]{en fait, la condition W(4b$''$) sur $W$.}), les fibres de $p$ sont $W$-asphériques, en particulier (par W(5)) 0-connexes. On a donc prouvé le

\msp{page 12}\textbf{Lemme.} (\footnote[24]{Utilise W(1, 2 bis, 3, 4b$''$) sur $W$.}) \textit{Soit $(X, x_0)$ dans} Cat$_\bullet$\textit{. Si $X$ est $W$-asphérique, alors $\Omega(X, x_0) \overset{\text{déf}}{=} \Ch(X; x_0, x_0)$ l'est aussi} (moyennant les axiomes W(1,2,3,4) sur $W$)\textit{. Donc (par itération) les $\Omega^n(X, x_0)$ sont $W$-asphériques} ($\forall\, n \in \mathbf{N}$).

\textbf{Corollaire.} \textit{Si de plus $W$ satisfait l'axiome} W(5) \textit{de connexité, alors $X$ est $W_\infty$-asphérique.}

En effet, `on sait' que cela signifie que tous les $\Omega^n(X, x_0)$ sont 0-connexes. Or l'axiome de connexité W(5) assure que $W$-asphérique implique 0-connexe.

\dgrblock{dgrb467}{The loop comparison that closes Theorem 1}
\begin{dcomment}{exposition}{p12:que cela signifie que tous les}{The loop comparison that closes Theorem 1}
Grothendieck has isolated the decisive argument: a \(W\)-aspherical category remains \(W\)-aspherical after forming its based loop category. To obtain the theorem, these loops must detect the ordinary homotopy groups. This is the comparison left open in the running text, and it has a short proof once the ordinary fibration properties already used are available.

Let \(E=\Ch(X;x,-)\) and \(F=\Ch(X;x,x)\). The endpoint map \(E\to X\) is a \(W_\infty\)-fibration. Its strict fibre over \(x\) maps by a weak equivalence to the corresponding comma category, and the comma categories vary by weak equivalences along arrows of \(X\). Quillen's theorem B therefore identifies \(BF\) with the homotopy fibre of \(BE\to BX\) \citep[Theorem~B and its corollary, pp.~97--98]{Quillen1973}. Since \(BE\) is weakly contractible, this gives a pointed equivalence
\[
 B\Ch(X;x,x)\simeq\Omega_x BX.
\]
At the next step choose the constant loop as basepoint. Repeating the same argument gives
\[
 B\Omega^r(X,x)\simeq\Omega^r_x BX,\qquad
 \pi_0\Omega^r(X,x)\cong\pi_r(BX,x)
 \quad(r\geq1),
\]
with the right side regarded as its underlying pointed set where only connectivity is used.

Now W(4b\(^{\prime\prime}\)) applies to the perfect path map \(E\to X\): both total space and base are \(W\)-aspherical, so their map is in \(W\), and its fibre is \(W\)-aspherical. Induction gives this for every \(\Omega^r(X,x)\). W(5) makes each of them nonempty and connected, hence all positive homotopy groups of \(BX\) vanish; W(5) applied to \(X\to e\) gives connectivity in degree zero. Thus \(X\to e\in W_\infty\).

Finally, for \(u:A\to B\in W\), the canonical perfect replacement \(A\to\bar A\to B\) has first map in both \(W\) and \(W_\infty\). The second map has \(W\)-aspherical fibres by W(4b\(^{\prime\prime}\)); the preceding argument makes those fibres weakly contractible. Its ordinary homotopy-fibre sequence makes it a weak equivalence, and two-out-of-three gives \(u\in W_\infty\). The conclusion is proved; the source presents it with the loop comparison still deferred.
\end{dcomment}


Mais le `on sait' est un tout petit peu léger (\footnote[25]{C'est prouvé fort simplement dans le corollaire du `theorème 1', page 190.}). Pour le justifier, il faudrait, ou bien % [sic] (footnote 25) theorème -> théorème, missing accent
que je fasse le lien avec la théorie 1/2-simpliciale classique, et prouver que les $\Omega^n$ que je construis sont `les bons', i.e. que leur nerf redonne bien (à isomorphisme près dans Hot) les $\Omega^n$ classiques (??). En particulier, que le $\pi_0(\Omega^n(X, x_0))$ est bien le $\pi_n(X, x_0)$ `de tout le monde'. Ou bien que je prouve directement que la 0-connexité des $\Omega^n(X, x_0)$ équivaut à la $W_\infty$-asphéricité de $X$, définie cohomologiquement à la \textsc{Artin-Mazur}. C'est donc un travail qu'il me faudra faire absolument. Mais ce n'est plus qu'un simple travail de routine, et je peux énoncer dès maintenant le

\textbf{Théorème 1.} (\footnote[26]{Il suffit même de W(4) sous une forme affaiblie, avec des foncteurs `fibrants' dans un sens `absolu'. Cf. le NB p. 10 - et surtout, cf. pages 13, 14.},\footnote[27]{NB Au lieu de W(5), il suffit de supposer $W \neq \mathrm{Fl}(\mathrm{Cat})$, cf. th. 2, p. 20.},\footnote[28]{Les hypothèses utiles sur $W$ sont finalement W(1), W(2 bis), W(3), W(4b$''$), et $W \neq \mathrm{Fl}(\mathrm{Cat})$.}) \textit{Soit $W \subseteq$} Fl(Cat) \textit{satisfaisant les conditions} W(1,2,3,4,5) (axiome de saturation, de localisation, de l'objet final, de fibration, de connexité)\textit{. Alors}
\[ W \subseteq W_\infty . \]

\textbf{Corollaire.} \textit{On peut caractériser $W_\infty$ comme la plus grande partie de} Fl(Cat) \textit{satisfaisant aux}
\msp{page 13}\textit{conditions} W(1) \textit{à} W(5) \textit{précédentes.}

Bien sûr, on aimerait un axiome 6 plus intéressant, qui joint aux précédents assure que $W = W_\infty$ (\footnote[29]{Cf. axiome W(6) page suivante.}). Ou est-ce que ces propriétés assurent déjà que $W = W_\infty$? Pour le voir, il faudrait peut-être que j'en sache plus sur la structure des objets $W_\infty$-asphériques? Si c'était faux, je songe à un axiome de `descente homotopique', qui impliquerait qu'un objet $W_\infty$-asphérique est aussi $W$-asphérique.

\dgrblock{dgrb468}{The conjectural minimum is now known}
\begin{dcomment}{status}{p13:Bien sûr, on aimerait un axiome 6 plus intéressant}{The conjectural minimum is now known}
The question on this page has a positive answer without the proposed additional axiom W(6). Cisinski proved that \(W_\infty\) is the smallest fundamental localizer \citep[Theorem~2.2.11]{Cisinski2004Minimal}. The refinement to the absolute theorem-A axioms identifies it also as the smallest \emph{weak} fundamental localizer \citep[Theorem~6.1.18]{Cisinski2006}. The source's W(1c) implies the weak-saturation condition in that theorem. Hence W(1), W(2 bis), and W(3) already give \(W_\infty\subseteq W\).

Combining this inclusion with Theorem~1 gives \(W=W_\infty\) under W(4b\(^{\prime\prime}\)) and W(5). The nondegeneracy argument at manuscript pp.~20--22 replaces W(5) by \(W\ne\Fl(\Cat)\). Accordingly, the conjecture at manuscript p.~15 is no longer open, and adding closure under retracts is unnecessary.

There is a local notational change to observe. On manuscript p.~13, \(W_\omega\) is introduced as the intersection of classes satisfying W(1)--W(5). On the next page it denotes the minimum satisfying only the basic localizer axioms, preferably W(2 bis). These are different definitions during the exploration. The later theorem makes both classes equal to \(W_\infty\), but that equality must not be used to conceal their different roles in the original argument.

The proof strategy here also reappears in Cisinski's characterization by a theorem-B property: a fundamental localizer preserving components and making every locally constant equivalence aspherical must be \(W_\infty\) \citep[Theorem~2.3.6]{Cisinski2004Minimal}. Its loop-space argument is the modern counterpart of the preceding manuscript proof.
\end{dcomment}


% foot: Version 7 décembre 2002
\printedpage{11}
% head: A. Grothendieck : Dérivateurs, Ch. XII, Édition des Universités de Montpellier II et Paris VII

Je note dès maintenant que si des $W_i$ ($i \in I$) satisfont aux axiomes W(1) à W(5), il en est de même de $\bigcap_i W_i$ (\footnote[30]{C'est vrai pour tout ensemble d'axiomes dans la nature d'axiomes de stabilité.}). Donc il y a un plus petit ensemble de flèches $W_\omega$ satisfaisant à ces conditions. S'il n'est égal à $W_\infty$, i.e. si les axiomes W(1) à W(5) ne caractérisent déjà $W_\infty$, ce serait un localiseur vraiment très intéressant.

Je me rends compte que je n'ai utilisé qu'une forme assez faible de l'axiome W(4) - essentiellement, la validité de la proposition 4. Je vais résumer différentes variantes de l'axiome de fibration.

\textbf{W(4)}

\begin{itemize}[leftmargin=4em, labelwidth=2.5em, labelsep=0.6em, align=right]
\item[a$_1$)] Si $f : X \longrightarrow Y$ dans Cat est une $W$-fibration gauche, c'est une $W$-fibration ultraforte.

\item[a$_2$)] Si $f : X \longrightarrow Y$ est une $W$-fibration droite, c'est une $W$-fibration ultraforte.

\item[a$'$)] Si $f : X \longrightarrow Y$ est $W$-parfait, i.e. à la fois $W$-propre et $W$-lisse (\footnote[31]{ce qui implique que c'est une $W$-fibration gauche \textit{et} droite.}), alors $f$ est une $W$-fibration ultraforte.

\item[a$''$)] Si $f : X \longrightarrow Y$ est à la fois homotopiquement propre et homotopiquement lisse (cf. NB page 10), $f$ est une $W$-fibration ultraforte. (Il vaut mieux dire : Si $f$ est $W_\omega$-parfait, c'est une $W$-fibration ultraforte.)

\item[b)] Soit $f : X \longrightarrow Y$ une flèche dans Cat$/S$, avec $X$ et $Y$ des $W$-fibrés sur $S$. Si $f \in W$, alors $\forall\, s \in S$, $f_s : X_s \longrightarrow Y_s$ est \texttt{[dans]} $W$.

\item[b$'$)] Soit $f : X \longrightarrow Y$ une flèche de Cat$/S$ qui soit $W$-parfaite et $W_\infty$-parfaite, alors $f \in W$ implique que les fibres de $f$ sont $W$-asphériques. (Il vaut mieux dire : Si $f$ est $W_\omega$-parfait, il est une $W$-fibration ultraforte.) \texttt{[Plutôt : Soit $f : X \longrightarrow Y$ une flèche dans \textrm{Cat}$/S$ telle que $X \longrightarrow S$ et $Y \longrightarrow S$ soient $W_\omega$-parfaites. Si $f \in W$, alors $\forall\, s \in S$, $f_s : X_s \longrightarrow Y_s$ est dans $W$.]}

\item[b$''$)] \msp{page 14}Soit $W_\omega \subseteq W \cap W_\infty$ la plus petite partie de Fl(Cat) satisfaisant les conditions W(1,2,3) (\footnote[32]{Il vaut mieux prendre W(2 bis) au lieu de W(2), pour trouver $W_\omega$ la plus petite possible!}) (axiomes de saturation, de localisation, de l'objet final). Soit $f : X \longrightarrow Y$ une flèche dans Cat qui soit $W_\omega$-parfaite. Si $f \in W$, alors les fibres de $f$ sont $W_\omega$-asphériques \texttt{[plutôt $W$-asphériques]}.
\end{itemize}

Entre ces propriétés de $W$, on a les implications tautologiques (sauf a$''$) $\Longrightarrow$ b), qui est essentiellement la proposition 4)

% foot: Version 7 décembre 2002
\printedpage{12}
% head: A. Grothendieck : Dérivateurs, Ch. XII, Édition des Universités de Montpellier II et Paris VII

\[
\begin{tikzcd}[row sep=huge, column sep=0pt]
\mathrm{a} \arrow[d, Rightarrow] & {}= (\mathrm{a}_1 \Longleftrightarrow \mathrm{a}_2) & \hspace{2.5cm} \\
\mathrm{a}' \arrow[d, Rightarrow] & & \smash{\left.\rule[-1.8cm]{0pt}{3.6cm}\right\}\ \text{\begin{tabular}{l}En fait,\\ toutes ces\\ conditions sont\\ équivalentes\end{tabular}}} \\
\mathrm{a}'' \arrow[d, Rightarrow] & & \\
\mathrm{b} \arrow[d, Rightarrow] & & \\
\mathrm{b}' \arrow[d, Rightarrow] & & \\
\mathrm{b}'' & &
\end{tikzcd}
\]

(\footnote[33]{Je reviens sur les diverses variantes de W(4) plus bas, p. 53-56.}).

Eh bien, dans le lemme page 12, il suffisait de supposer que $W$ satisfait à la version la plus faible b$''$) de l'axiome des fibrations (sept versions!). En effet, les fibrations $\bar{X} \longrightarrow Y$ et $E \longrightarrow X$ utilisées dans la démonstration du lemme sont $W_\omega$-parfaites, et il suffisait de savoir que si elles sont dans $W$, alors les fibres sont $W$-asphériques.

\dgrblock{dgrb469}{What the weakest fibre axiom actually needs}
\begin{dcomment}{exposition}{p14:Eh bien, dans le lemme page 12}{What the weakest fibre axiom actually needs}
The list of seven variants is an attempt to reduce the amount of fibration theory needed for the characterization. Its least demanding member, W(4b\(^{\prime\prime}\)), asks only this: a \(W\)-equivalence which is \(W_\omega\)-perfect has \(W\)-aspherical strict fibres. The conclusion is relative to \(W\), as the editors correct; it does not assert that the fibres are already \(W_\omega\)-aspherical.

For the proof of Theorem~1 this is enough because the maps to which the axiom is applied have a special origin. The replacement \(\bar A\to B\) and the based-path projection \(\Ch(X;x,-)\to X\) are \(W_\omega\)-perfect before any fibration axiom for \(W\) is assumed. Two-out-of-three puts these maps in \(W\), and W(4b\(^{\prime\prime}\)) then supplies exactly the asphericity required. Arbitrary \(W\)-fibrations need not be considered.

The distinction between preservation and detection also explains the middle variants. W(4a\(^{\prime\prime}\)) requires base change along an absolute perfect map to preserve \(W\). W(4b\(^{\prime}\)) concerns a map \emph{between} two absolute perfect families over the same base. W(4b\(^{\prime\prime}\)) concerns one perfect map and compares it to the identity family on its target. The editorial replacement of the entire statement of b\(^{\prime}\) restores this intended progression.

The brace printed on this page spans the three a-conditions; it should not be read as a proof of every converse in the seven-term chain. Section~IV returns to those converses and to the saturation needed for the localization arguments. In the modern setting, minimality and the nondegeneracy criterion below identify the two possibilities as \(W_\infty\) and the class of all maps, so these different rigidity axioms ultimately select the same possibilities. That later conclusion is stronger than the argument supplied by the diagram itself.
\end{dcomment}


Mais je ne pense pas que même la version la plus forte W(4a) de l'axiome des fibrations, jointe aux trois axiomes précédents et à W(5) (connexité), implique $W = W_\infty$. Il faudrait tout au moins y joindre l'axiome

\begin{itemize}[leftmargin=4em, labelwidth=3em, labelsep=0.6em]
\item[\textbf{W(6)}] (Axiome des limites filtrantes.) Soit $X = \varinjlim\limits_{i} X_i$ une limite inductive filtrante dans Cat. Si les $X_i$ sont $W$-asphériques, $X$ aussi.
\end{itemize}

Ceci devrait être équivalent à la propriété, en apparence plus forte, en termes de flèches de $W$ :

\textbf{Proposition 5.} \textit{Supposons qu'on ait} W(6) (en plus de W(1) à W(4) (\footnote[34]{W(4) sous la forme la plus faible W(4b$''$) .})),
\msp{page 15}\textit{et soit}
\[ (f_i)_{i\in I} : (X_i)_{i\in I} \longrightarrow (Y_i)_{i\in I} \]
% foot: Version 7 décembre 2002
\printedpage{13}
% head: A. Grothendieck : Dérivateurs, Ch. XII, Édition des Universités de Montpellier II et Paris VII
\textit{un homomorphisme de systèmes inductifs dans} Cat\textit{, indexé par un ensemble ordonné filtrant, soient $X = \varinjlim X_i$, $Y = \varinjlim Y_i$, $f = \varinjlim f_i$, $f : X \longrightarrow Y$. Si $f_i$ est dans $W$ pour tout $i$, alors $f \in W$.}

On factorise chaque $f_i$ en
\[ X_i \xrightarrow{\alpha_i} \bar{X}_i \xrightarrow{\bar{f}_i} Y_i, \qquad \alpha_i \in W_\omega,\ \bar{f}_i \in \mathrm{Parf}\, W_\omega \]
(\footnote[35]{$W_\omega$ introduit dans la condition W(4b$''$), page 14.}) suivant le système breveté. La construction étant fonctorielle, les $\bar{X}_i$ forment eux aussi un système inductif. De plus, comme le foncteur $S \longrightarrow \Ch(S)$ qui donne naissance à la factorisation canonique des morphismes, commute aux $\varinjlim$ filtrantes (immédiat), il s'ensuit que, si on pose
\[ \bar{X} = \varinjlim_{i} \bar{X}_i \]
et qu'on regarde les flèches
\[ X \xrightarrow{\alpha} \bar{X} \xrightarrow{\bar{f}} Y, \qquad \alpha = \varinjlim \alpha_i,\ \bar{f} = \varinjlim \bar{f}_i, \]
avec $\bar{f}\alpha = \bar{f}$, cette factorisation de $f$ est bien sa factorisation canonique. Donc $\alpha \in W$, % [sic] $\bar{f}\alpha = \bar{f}$ -> $\bar{f}\alpha = f$, the composite should be f
$\bar{f} \in \mathrm{Parf}\, W_\omega$. Ceci vu, en vertu de prop. 3, $f \in W$, i.e. $\bar{f} \in W$, est vrai si les fibres de $\bar{f}$ sont $W$-asphériques. Mais les fibres de $\bar{f}$ s'interprètent comme limites inductives filtrantes de fibres des $\bar{f}_i$, lesquelles sont $W$-asphériques par la condition W(4b$''$) des fibrations. Donc l'axiome W(6) donne ce qu'il faut.

\dgrblock{dgrb470}{Filtered fibres require a moving basepoint}
\begin{dcomment}{exposition}{p15:La construction étant fonctorielle}{Filtered fibres require a moving basepoint}
Proposition~5 turns an axiom about aspherical objects into an axiom about maps by making the fibres of those maps explicit. Two finite-presentability assertions are implicit in the sketch, and both can be checked directly.

First, a path of a fixed finite type is a functor from a finite category, and a morphism between two such paths is specified by finitely many arrows and equations. If \(S=\varinjlim_i S_i\) is a filtered colimit, any path or morphism of paths in \(S\) is represented at one stage, and an equality between representatives holds at some later common stage. Thus \(\Ch(-)\), including its source and target maps, commutes with filtered colimits. Since filtered colimits of sets commute with finite limits, the finite pullbacks defining the mapping-path replacement commute with them as well.

Second, a fibre over \(y\in Y=\varinjlim_iY_i\) is not the colimit of fibres over an unspecified fixed \(y\). Choose a representative \(y_i\in Y_i\), and let \(y_j\) be its image for \(j\geq i\). There is a canonical isomorphism
\[
 (\bar X)_y\ \cong\
 \varinjlim_{j\geq i}(\bar X_j)_{y_j}.
\]
For objects, a representative \(x_j\) mapping to \(y\) has image equal to \(y_k\) at a sufficiently late stage \(k\), so it becomes an object of the indicated fibre. For arrows, the condition that the image be the identity of \(y\) is likewise witnessed at a later stage. The same finite-stage argument proves injectivity and compatibility with composition. Passing to \(j\geq i\) is cofinal and remains filtered.

Each map \(\bar f_j\) belongs to \(W\), since \(f_j\) and its replacement inclusion do. W(4b\(^{\prime\prime}\)) makes its fibres \(W\)-aspherical, W(6) makes their displayed colimit \(W\)-aspherical, and theorem A for a proper or smooth map gives \(\bar f\in W\). Finally \(\alpha\in W_\omega\subseteq W\), so \(f=\bar f\alpha\in W\). The printed equation \(\bar f\alpha=\bar f\) is a slip retained in the source.

For fundamental localizers in the modern sense, filtered-colimit stability is already a theorem, with no accessibility assumption \citep[Corollary~4.2.22]{Cisinski2006}. The calculation here remains useful because it shows exactly how a fibre argument reduces that stability to the object form.
\end{dcomment}


\textbf{Conjecture} (hardie? sotte?). \textit{Si $W \subseteq$} Fl(Cat) \textit{satisfait les axiomes} W(1) \textit{à} W(6) \textit{(avec} W(4) \textit{sous sa forme la plus faible} W(4b$''$) (\footnote[36]{et peut-être W(2) sous la forme plus faible W(2 bis).})\textit{), alors $W = W_\infty$.}

\msp{page 16}Mais il est bien possible qu'il faille renforcer un peu W(1) (saturation), en exigeant que tout rétracte d'un $f \in W$ soit dans $W$.

\dgrblock{dgrb471}{The proposed retract axiom is unnecessary}
\begin{dcomment}{status}{p16:Mais il est bien possible qu'il faille renforcer}{The proposed retract axiom is unnecessary}
The possible need for an extra retract axiom is resolved by the later theory. Every fundamental localizer with the relative theorem-A axiom W(2) is strongly saturated and consequently closed under retracts \citep[Proposition~4.2.4]{Cisinski2006}. In the weaker setting using W(2 bis), the specific equality conjectured here already follows from the weak minimal-localizer theorem and the fibre argument, so no independent retract hypothesis is required for that conclusion either. This does not assert strong saturation for every weak fundamental localizer before the fibre axiom is imposed.

This does not retroactively make every appeal to an isomorphism in \(\Hot_W\) valid under two-out-of-three alone. The elementary proofs above deliberately establish membership in \(W\) directly. The later strong-saturation theorem explains when such arguments may be shortened; it should be cited at the point where that stronger conclusion is used.
\end{dcomment}


\textbf{Remarque.} Considérons, pour $n \in \mathbf{N}$, l'ensemble $W_n \subseteq$ Fl(Cat) des `$n$-équivalences d'homotopie' dans Cat :
\[ f \in W_n \iff \pi_0(f) \text{ bijectif, et } \pi_i(f, x) : \pi_i(X, x) \longrightarrow \pi_i(Y, f(x)) \text{ bijectif pour } 0 < i \leq n. \]

Je croyais que ça satisfaisait à tout ce qu'on peut demander à un localiseur, sans avoir pris la peine de l'écrire avec soin. Ma ça ne peut satisfaire à W(1) jusqu'à W(4b$''$), puisqu'on % [sic] Ma -> Mais, word truncated
n'a pas $W_n \subseteq W_\infty$!

\dgrblock{dgrb472}{Truncated equivalences really are localizers}
\begin{dcomment}{exposition}{p16:n'a pas $W_n \subseteq W_\infty$!}{Truncated equivalences really are localizers}
The class \(W_n\) is not a defective homotopy theory. It is the ordinary theory observed through the first \(n\) homotopy groups. More precisely,
\[
  u\in W_n\quad\Longleftrightarrow\quad
  \tau_{\leq n}Bu\text{ is an equivalence},
\]
where \(\tau_{\leq n}\) is Postnikov truncation. This functor is left adjoint to the inclusion of \(n\)-types \citep[Proposition~5.5.6.18]{Lurie2009HTT}. Its universal property gives a direct modern proof of the assertions the source recalls without repeating.

Identities, two-out-of-three, and the two-cycle condition hold because they hold for equivalences of \(n\)-types. Strong saturation follows too: the functor \(\tau_{\leq n}B\), regarded as taking values in the homotopy category of \(n\)-types, factors through \(\Cat[W_n^{-1}]\). Hence any map inverted there induces an equivalence of \(n\)-types and therefore was in \(W_n\) already. Categories with terminal objects have contractible nerves, giving W(3).

For the relative axiom W(2), let \(X\to S\) be any functor. The simplicial bar construction on the covariant diagram \(s\mapsto N(X/s)\) has bisimplices consisting of
\[
 x_0\to\cdots\to x_q,\qquad
 p(x_q)\to s_0\to\cdots\to s_r.
\]
For a fixed \(q\)-simplex of \(NX\), its \(r\)-direction is the nerve of \(p(x_q)\backslash S\), which has an initial object. Realizing this direction therefore gives a weak equivalence of the diagonal bar construction with \(NX\). In other words,
\[
 BX\simeq\operatorname*{hocolim}_{s\in S}B(X/s).
\]
This is a homotopy colimit of spaces, retaining the arrows of \(S\), rather than a colimit of their homotopy categories. If \(X/s\to Y/s\) is in \(W_n\) for every \(s\), apply \(\tau_{\leq n}\) to the displayed colimits. As a left adjoint, truncation preserves them as colimits in the category of \(n\)-types. The resulting objectwise equivalence induces an equivalence of colimits, proving \(X\to Y\in W_n\).

Finally, the nerve commutes with filtered colimits because each finite chain of arrows and each relation is witnessed at a finite stage. A Kan replacement such as \(\operatorname{Ex}^{\infty}\) also commutes with filtered colimits: each simplex of \(\operatorname{Ex}\) is represented by a map from a finite subdivision, and its sequential iteration commutes with colimits. Homotopy groups of Kan complexes are computed by finite simplicial spheres and finite homotopies. They consequently commute with filtered colimits, with basepoints represented at sufficiently late stages. This proves W(6) for \(W_n\). What fails is the fibre axiom, as the next computation shows.
\end{dcomment}


W(1) est trivial sur la définition (et on a sans larmes même la \textit{saturation forte} de $W_i$, % [sic] $W_i$ -> $W_n$ (also in the following lines), the remark concerns W_n
i.e. si $f$ dans Cat est \texttt{[un]} isomorphisme dans $\mathrm{Hot}_{W_i} = (\mathrm{Cat})W_i^{-1}$, alors $f \in W_i$). Je crois me rappeler avoir vérifié W(2), par la même méthode que pour $W_\infty$, en utilisant un critère cohomologique à la \textsc{Artin-Mazur}, pour exprimer $f \in W_i$, et en utilisant les suites spectrales de Leray pour $X \longrightarrow S$ et $Y \longrightarrow S$. La condition W(3) est trivialement satisfaite puisque $W_n \supseteq W_\infty$ et qu'elle l'est pour $W_\infty$. Donc ça canule là où je m'y
% foot: Version 7 décembre 2002

% ---- from Der-ChXIIa-012_p016-020.tex
\printedpage{14}
% head: A. Grothendieck : Dérivateurs, Ch. XII, Édition des Universités de Montpellier II et Paris VII
attendais le moins : l'axiome des fibrations. Prenons la forme faible : Soit $f : X \longrightarrow Y$ une $W_\omega$-fibration et $f \in W_i$, les fibres $X_s$ sont elles $W_i$-asphériques (\footnote[37]{$\omega > \infty$, $W_\omega \subset W_\infty \subset W_n$})? On peut supposer $X$ et $Y$ 0-connexes, soit $x \in X_s$, la suite exacte d'homotopie donne
\[
\pi_{i+1}(X) \xrightarrow[\text{si } i \leq n-1]{\text{iso}} \pi_{i+1}(Y) \longrightarrow \pi_i(X_s) \longrightarrow \pi_i(X) \xrightarrow[\text{si } i \leq n]{\text{iso}} \pi_i(Y),
\]
donc
\[
\pi_i(X_s) = 1 \qquad \text{si } 0 \leq i \leq n-1,
\]
mais pour $i = n$ on trouve\msp{page 17}
\[
\pi_n(X_s) = \mathrm{Coker}\left(\pi_{n+1}(X) \longrightarrow \pi_{n+1}(Y)\right),
\]
il n'y a aucune raison que ce soit nul, comme ça devrait si $X_s \longrightarrow e$ était dans $W_n$.

\dgrblock{dgrb473}{The missing group occurs one degree higher}
\begin{dcomment}{examples}{p17:il n'y a aucune raison que ce soit nul}{The missing group occurs one degree higher}
Grothendieck's exact sequence identifies the precise loss of information. For a connected homotopy fibre \(F\) of an \(n\)-equivalence \(E\to B\), with \(n\geq1\), exactness gives
\[
 \pi_iF=0\quad(0\leq i<n),\qquad
 \pi_nF\cong
 \operatorname{coker}\bigl(\pi_{n+1}E\to\pi_{n+1}B\bigr).
\]
At \(n=1\), the cokernel is still meaningful: the connecting homomorphism from the abelian group \(\pi_2B\) surjects onto \(\pi_1F\), so the latter group is abelian in this situation. The long exact sequence and its low-dimensional interpretation are recalled in \citet[Theorem~4.41]{Hatcher2002}.

For an example entirely within \(\Cat\), take \(S\) to be the poset of nonempty proper faces of an \((n+2)\)-simplex. Then \(BS\) is the barycentric subdivision of its boundary, hence has the homotopy type of \(S^{n+1}\). Choose a vertex \(s\), and use
\[
 E=\Ch(S;s,-)\xrightarrow{\beta}S .
\]
The total space \(BE\) is contractible, and \(\beta\) is \(W_\omega\)-perfect by the construction recalled on manuscript pp.~8--11. Since the homotopy groups of the sphere vanish through degree \(n\), \(\beta\in W_n\). Its fibre \(F=\Ch(S;s,s)\), however, satisfies
\[
 BF\simeq\Omega S^{n+1},\qquad \pi_n BF\cong\mathbb Z.
\]
Thus \(F\to e\notin W_n\), contradicting W(4b\(^{\prime\prime}\)). The same square shows that \(\beta\) does not preserve all \(W_n\)-equivalences under base change: \(e\to S\) is in \(W_n\), whereas its pullback \(F\to E\) is not.

For \(n=0\), the sphere example remains valid, with \(\pi_0\Omega S^1\cong\mathbb Z\), but the displayed \emph{group cokernel} in the source requires correction. In general \(\pi_0F\) is the pointed coset set of the image of \(\pi_1E\) in \(\pi_1B\), not its quotient by the normal closure. Take \(G=S_3\), \(H=\langle(12)\rangle\), and the action groupoid \(E=(G/H)\mathbin{/\mkern-6mu/}G\). The projection \(E\to BG\) is both a Grothendieck fibration and an opfibration; its strict fibre is the discrete three-element set \(G/H\). Both \(E\simeq BH\) and \(BG\) are connected, so the projection is in \(W_0\). Yet the group cokernel of \(H\to G\) is trivial, since the conjugates of \((12)\) generate \(S_3\). The fibre has three components. This is a genuine degree-zero qualification, not a defect in the argument for positive \(n\).
\end{dcomment}


Il est vrai que parfois on renforce justement la condition pour $W_n$, disons qu'on définit $W'_n \subseteq W_n$ par
\[
f \in W'_n \iff \left\{
\begin{array}{ll}
\text{a)} & \pi_0(f) \text{ bijectif,} \\[1ex]
\text{b)} & \pi_i(f,x) : \pi_i(X,x) \longrightarrow \pi_i(Y,f(x)) \text{ bijectif pour } 0 < i \leq n \text{ et } x \in X, \\[1ex]
\text{c)} & \pi_{n+1}(f,x) : \pi_{n+1}(X,x) \longrightarrow \pi_{n+1}(Y,f(x)) \ \textit{surjectif,}
\end{array}
\right.
\]
c'est à dire aussi
\[
f \in W'_n \iff \left\{
\begin{array}{l}
\text{les `fibres homotopiques' } Z \text{ de } X \longrightarrow Y \text{ (i.e. les fibres de } \bar{X} \xrightarrow{\bar{f}} Y\text{)} \\
\text{sont } n\text{-connexes, i.e. 0-connexe et } \pi_i(Z_y,z) = 0 \text{ pour } 1 \leq i \leq n.
\end{array}
\right.
\]
Alors $X$ dans Cat sera $W'_n$-asphérique si et seulement s'il est $W_n$-asphérique, i.e. $n$-connexe. Donc par la forme 2 de la définition de $W'_n$, l'axiome W(4b$''$) est satisfait. Il faut croire que ça pèche ailleurs! Peut-être déjà la saturation.

\dgrblock{dgrb474}{Connectivity repairs fibres but loses cancellation}
\begin{dcomment}{examples}{p17:Alors $X$ dans Cat sera}{Connectivity repairs fibres but loses cancellation}
The strengthened class \(W'_n\) measures a connectedness condition on a map: its homotopy fibres are \(n\)-connected. The long exact sequence makes this equivalent to bijectivity on \(\pi_0,\ldots,\pi_n\) and surjectivity on \(\pi_{n+1}\). For the terminal map \(X\to e\), the last surjectivity condition is automatic, so \(W_n\) and \(W'_n\) have exactly the same aspherical objects.

The failure of two-out-of-three can be seen with the same finite category \(S\) modeling \(S^{n+1}\). Let \(a:e\to S\) select a vertex and let \(b:S\to e\) be the terminal functor. Then \(ba=\id_e\in W'_n\). Also \(b\in W'_n\): it is an isomorphism on homotopy groups through degree \(n\), and the map \(\mathbb Z=\pi_{n+1}S\to0\) is surjective. But \(a\notin W'_n\), because \(0\to\mathbb Z\) is not surjective. This works also for \(n=0\), using the circle.

Thus the problematic cancellation is exactly
\[
 b,\ ba\in W'_n\quad\nRightarrow\quad a\in W'_n.
\]
The source's informal argument about surjective homomorphisms is correct; the example supplies actual functors and verifies every homotopy group involved. A class of connected maps can be stable under many geometric operations without being the class inverted by a localization. The modern theory treats connectedness as one side of a factorization system; it is not required to satisfy the cancellation law for equivalences.
\end{dcomment}


\hspace*{3em}Si $gf$ et $g$ sont dans $W'_i$, il n'est pas forcément de même de $f$,

car si $\bar{g}\bar{f}$ (et donc $\bar{g}$) est surjectif, il n'en est par forcément de même pour $\bar{f}$! Donc $W'_n$ ne mérite pas le nom de `localiseur'. % [sic] par forcément -> pas forcément, misprint

Par contre, $W_n$ me paraît excellent, si on ne lui demande pas l'impossible - de satisfaire à un `axiome de fibration' qui semble bien (à peu de choses près) caractériser le seul $W_\infty$ parmi les `bons' localiseurs. En tous cas l'axiome de connexité W(5) est satisfait, et l'axiome des limites filtrantes W(6) aussi. Par ailleurs W(1) (et même la saturation forte), W(2), W(3) sont satisfaits par tout $W_{\mathbf{D}}$ associé à un dérivateur - ce qui confirme à nouveau leur importance. Les axiomes W(4) (l'axiome `draconien' des fibrations, $\pm$ spécial à $W_\infty$), W(5) (connéxité) et % [sic] connéxité -> connexité, misspelling (written connexité elsewhere)
\printedpage{15}
% head: A. Grothendieck : Dérivateurs, Ch. XII, Édition des Universités de Montpellier II et Paris VII
\msp{page 18}W(6) (limites) m'apparaissent comme de nature plus spéciale, je ne songerais à imposer aucun d'eux pour l'usage terminologique de `localiseur fondamental'.

\newpage
\printedpage{16}
% head: A. Grothendieck : Dérivateurs, Ch. XII, Édition des Universités de Montpellier II et Paris VII
\section*{II\quad Perplexités sur les axiomes et la terminologie.}

Dans I et la plupart du temps, en travaillant avec les dérivateurs, j'ai supposé \fbox{$\ W_{\mathbf{D}} \supseteq W_\infty\ $}, hypothèse qui m'a parue essentielle (à tort (\footnote[38]{à \textit{tort} il s'avère maintenant.}) ou à raison). On voit donc que si $W_{\mathbf{D}}$ satisfait à l'axiome de fibrations W(4), même sous sa forme la plus faible W(4b$''$), il est \textit{égal} à $W_\infty$. Cela montre bien que cet axiome doit être vu comme pratiquement une \textit{caractérisation} de $W_\infty$ (parmi les localiseurs fondamentaux qui contiennent $W_\infty$ -- et ces localisateurs me paraissent les plus importants, et peut-être les seuls vraiment utiles).

La théorie des \textbf{D}-fibrations dans I, p. 111-126, a été faite dans l'esprit qu'on vient de dire, en supposant $W_{\mathbf{D}} \supseteq W_\infty$. La notion de \textbf{D}-fibrations à laquelle je m'étais finalement arreté était par suite différente des notions envisagées ici dans un esprit différent (page 4). Elle revient à tester la condition a) (trop forte presque toujours, sauf pour $W = W_\infty$), caractérisant les $W$-fibrations ultrafortes, par des $Y'' \longrightarrow Y'$, où $Y''$, $Y'$ sont tous deux $W_\infty$-asphériques (ce qui est plus exigeant que $W$-asphérique). C'est seulement quand $Y''$ et $Y'$ sont $W_\infty$-asphériques (au lieu de $W$-asphériques), qu'on demande que $X'' \longrightarrow X'$ soit dans $W$. On trouve donc\msp{page 19} une notion moins exigeante que celle de $W$-fibration forte. Si on dénote ici cette condition par b), on trouve (sous réserve que $W_\infty \subseteq W$), avec les notations de la page 4 (\footnote[39]{Cf. la bonne notion de $W$-fibration dans un contexte où on ne suppose pas $W_\infty \subseteq W$, dans th. 3 et déf. 3, p. 51, 52.}) % [sic] arreté -> arrêté, missing circumflex
% In the print, a curved arrow starts at "b" in the diagram below and points up to "b)" in the line "par b), on trouve ..." above; it is drawn here to the empty node at the top left.
\[
\begin{tikzcd}[column sep=small, row sep=large]
{} & & & & \\
& & \mathrm{a} \arrow[d, Rightarrow] & & W\text{-fibration ultraforte} \\
& & \mathrm{a}' \arrow[d, Rightarrow] & & W\text{-fibration forte} \\
& & \mathrm{b} \arrow[dl, Rightarrow] \arrow[dr, Rightarrow] \arrow[uuull, bend left=30] & & W\text{-fibration} \\
W\text{-fibration gauche} & \mathrm{a}_1 & & \mathrm{a}_2 & W\text{-fibration droite.}
\end{tikzcd}
\]
a$_1$), a$_2$) étaient aussi appelées \textbf{D}-fibrations cohomologique resp. \textbf{D}-fibrations homologiques (\footnote[40]{Non, c'est une notion encore plus forte, à coups de coefficients localement constants, cf. I, p. 121.}). Quand \textbf{D} satisfait Der 3a$'$ (I, p. 111) sur les \textbf{D}-coefficients localement constants sur les $X$ $W_\infty$-asphériques, alors les notions b), a$_1$), a$_2$) sont équivalentes. (Je présume que c'est le cas pour tous les \textbf{D} et tous les $W$ vraiment intéressants.) Par contre, j'ai l'impression que la notion de $W$-fibration forte, voire ultraforte, n'est guère utilisable que dans le cas $W = W_\infty$, où elle coïncide avec les précédentes. C'est dire qu'on peut probablement oublier la terminologie, et se rappeler par contre que dans le cas $W = W_\infty$, les % [sic] cohomologique -> cohomologiques, agreement with "D-fibrations"
\printedpage{17}
% head: A. Grothendieck : Dérivateurs, Ch. XII, Édition des Universités de Montpellier II et Paris VII
$W$-fibrations ont des propriétés d'une force extrême, s'exprimant par les deux propriétés dégagées dans W(4a) et W(4b).

\dgrblock{dgrb475}{Changing the test bases changes the theory}
\begin{dcomment}{exposition}{p19:Par contre, j'ai l'impression que la notion}{Changing the test bases changes the theory}
The revised hierarchy separates the homotopy type of the indexing category from the coefficients used to observe it. A \(W_\infty\)-aspherical category is an absolutely contractible test shape. A \(W\)-aspherical category may instead be invisible only after truncation or after a chosen cohomological localization. Demanding preservation over the latter shapes is stronger when \(W_\infty\subseteq W\).

For \(W_n\), the sphere category used above is \(W_n\)-aspherical but not \(W_\infty\)-aspherical. It is therefore an admissible base in the early strong test \(a'\), and the path family fails that test; it is not an admissible base in the new condition \(b\). This is the concrete reason that the revised notion can accommodate truncated homotopy theories. The source's footnote~40 also corrects a different confusion: a condition tested using every locally constant coefficient is stronger than a condition on the underlying localizer alone. The two distinctions should not be merged.

The opening inference of section~II needs the nondegeneracy hypothesis that Theorem~2 will make explicit. From \(W_{\mathbf D}\supseteq W_\infty\) and W(4), one cannot conclude \(W_{\mathbf D}=W_\infty\) without excluding \(W_{\mathbf D}=\Fl(\Cat)\). The terminal derivator, whose value on every indexing category is the terminal category, gives \(W_{\mathbf D}=\Fl(\Cat)\); all the fibration implications then hold trivially. With \(W_{\mathbf D}\ne\Fl(\Cat)\), the corrected inference follows from Theorems~1 and~2. This is the degenerate exception already recognized at manuscript p.~10, rather than a new obstruction.
\end{dcomment}


\msp{page 20}Dans la théorie de $\mathrm{HOT}_W$, influencé par la réflexion générale sur les dérivateurs et l'axiome Der 3b, j'ai travaillé sous l'hypothèse $W \supseteq W_\infty$, mais en utilisant un axiome de fibration L5 $=$ W(4b) qui, il s'avère, implique finalement $W = W_\infty$. J'ai l'impression que je n'ai jusqu'à présent (propriétés de $f_!$ pour HOT) pas utilisé $W \supseteq W_\infty$, donc ça peut être utile, comme `ça ne coûte pas plus cher' (et oblige d'être plus attentif à ce qu'on utilise), de laisser tomber $W \supseteq W_\infty$ (impliquant $W = W_\infty$), et de mettre par contre sur $W$ des axiomes aussi forts (notamment de fibration) que les propriétés de $W_\infty$, `sans se priver en rien'. Je supposerai donc W(1), W(2), W(3), W(4a) (la forme la plus forte de l'axiome de fibration, cf. pages 13, 14) - c'est tout ce qui a été utile jusqu'à présent. Si on y joint W(5) (axiome de connexité) ça implique $W \subseteq W_\infty$. Mais pour l'instant je ne crois pas en avoir un besoin.

\dgrblock{dgrb476}{A fibre criterion that avoids component restriction}
\begin{dcomment}{development}{p20:Mais pour l'instant je ne crois pas en avoir un besoin.}{A fibre criterion that avoids component restriction}
Grothendieck wants to deduce preservation of components from the nontriviality of the theory. The displayed proof begins by restricting a \(W\)-map to each component of its target; footnote~42 correctly observes that this invokes a stronger fibre-detection statement than W(4b\(^{\prime\prime}\)). The theorem can be proved with the weaker axiom. The following formulation isolates the repair. We did not find this exact formulation in the sources checked for this audit; it is an elementary reorganization of the fibre argument, with no claim of priority.

\textbf{Proposition (two-step nondegeneracy criterion).}\ours\
Let \(V\) be a class of functors containing the categorical homotopisms and satisfying two-out-of-three and absolute theorem A: if every \(A/b\to e\) lies in \(V\), then \(A\to B\) lies in \(V\). Suppose also that for every \(u:A\to B\) in \(V\) and every \(b\in B\), its ordinary homotopy fibre has the weak homotopy type of a category \(F_b\) with \(F_b\to e\in V\). Then either every map of \(V\) induces a bijection on components, or \(V=\Fl(\Cat)\).

\emph{Proof.}
We first record two consequences of the stated axioms. If \(\emptyset\to e\in V\), then every \(\emptyset\to C\) belongs to \(V\) by absolute theorem A, since all its comma categories are empty. Two-out-of-three applied to
\(\emptyset\to C\to e\) gives \(C\to e\in V\). Applying it once more to \(C\to D\to e\) gives every functor \(C\to D\in V\).

If \(D\to e\in V\), then \(D\times D\to e\in V\). Indeed, for \(d\in D\), the comma category of the second projection is
\((D\times D)/d\cong D\times(D/d)\).
The category \(D/d\) has a terminal object, so the projection \(D\times(D/d)\to D\) is a categorical homotopism: its section at that terminal object has composite naturally joined to the identity. Hence \((D\times D)/d\to e\in V\). Absolute theorem A gives \(D\times D\to D\in V\), and composition gives the assertion. When \(D\) is empty the assertion is immediate.

Now take \(u\in V\) for which \(\pi_0u\) is not bijective. If a component of \(B\) is missed, choose \(b\) in it. The homotopy fibre of \(Bu\) at \(b\) is empty: a point in that fibre would supply a path from an image point to \(b\). Consequently its categorical model \(F_b\) is empty as well, because a category with an object has nonempty realization. The preceding paragraph gives \(V=\Fl(\Cat)\).

Suppose instead that two distinct components of \(A\) map into the same component of \(B\). Choose one object in each and paths from their images to some \(b\). These give two points of the homotopy fibre which cannot lie in the same component, since projection to \(BA\) would connect the chosen objects. Its categorical model \(D=F_b\) is therefore disconnected and \(V\)-aspherical. The product argument makes \(D\times D\) \(V\)-aspherical, so its diagonal \(d:D\to D\times D\) belongs to \(V\). Choose objects \(d_0,d_1\) in different components of \(D\). The homotopy fibre of \(Bd\) at \((d_0,d_1)\) is empty: it would otherwise provide a point connected to both \(d_0\) and \(d_1\). Apply the fibre hypothesis to this diagonal. It again gives \(\emptyset\to e\in V\), and hence \(V=\Fl(\Cat)\). This proves the proposition.

For the source's class \(W\), the hypothesis about homotopy fibres follows from its canonical factorization \(A\to\bar A\to B\). The first arrow belongs to \(W_\omega\subseteq W\cap W_\infty\); the second is \(W_\omega\)-perfect and lies in \(W\), so W(4b\(^{\prime\prime}\)) makes its strict fibres \(W\)-aspherical. Ordinary theorem B identifies those fibres with the homotopy fibres of the original map. The proposition therefore proves Theorem~2 using W(1), W(2 bis), W(3), and W(4b\(^{\prime\prime}\)), including the weakened hypotheses asserted in footnote~28 of Theorem~1.
\end{dcomment}


Mais méditant sur la question, je crois avoir trouvé le

\textbf{Théorème 2} (!) (\footnote[41]{Cf. pour un résultat plus fort XV 9.12, p. 133 bis. \textbf{NB} Le résultat est plus général, mais pas vraiment plus fort, car il ne donne pas $W = \Fl(\Cat)$.}) : \textit{Soit} $W \subseteq \Fl(\Cat)$ \textit{satisfaisant} W(1,2,3,4b$''$)\textit{. Si} $W$ \textit{ne satisfait pas à} W(5)\textit{, alors} $W = \Fl(\Cat)$.

\textsc{Démonstration}. Soit $f : X \longrightarrow Y$ dans Cat, tel que $\pi_0(f)$ ne soit pas bijectif. Soit $I$ la catégorie discrète ayant pour ensemble d'objets l'ensemble $\pi_0(Y)$, on a donc un diagramme commutatif
\[
\begin{tikzcd}[column sep=small]
X \arrow[rr, "f"] \arrow[dr, "p"'] & & Y \arrow[dl, "q"] \\
& I, &
\end{tikzcd}
\]
d'autre part un foncteur d'une catégorie dans une catégorie discrète est toujours $W_\omega$-parfait, donc $p$, $q$ sont des $W$-fibrations. Comme\msp{page 21} $f \in W$, il s'ensuit $f_i \in W$ pour tout $i \in I$ (\footnote[42]{On utilise W(4) sous forme plus forte que W(4b$''$)!}).

Si $\pi_0(f)$ n'était pas surjective, prenant $i$ qui n'est pas dans l'image, on aurait $X_i = \emptyset \longrightarrow Y_i \neq \emptyset$, qui serait dans $W$. Soit donc
\[
f : X = \emptyset \longrightarrow Y
\]
dans $W$. Alors $\forall$ $Z$ dans Cat, $f \times \mathrm{id}_Z : \underbrace{X \times Z}_{= \emptyset} \longrightarrow Y \times Z$ est dans $W$ (car $Z$ est
\printedpage{18}
% head: A. Grothendieck : Dérivateurs, Ch. XII, Édition des Universités de Montpellier II et Paris VII
$W$-fibrée sur $e$), donc $\emptyset \longrightarrow Y \times Z$ est dans $W$. Donc aussi la projection
$\begin{tikzcd}[column sep=small, row sep=small, baseline=(\tikzcdmatrixname-1-1.base)]
Y \times Z \arrow[rr] & & Y. \\
& \emptyset \arrow[ul] \arrow[ur] &
\end{tikzcd}$
Comme cette projection est une $W$-fibration, il s'ensuit par W(4b$''$) que ses fibres sont $W$-asphériques. Comme $Y \neq \emptyset$, il y a une telle fibre, elle est $\simeq Z$. Donc $Z$ est $W$-asphérique. Donc on trouve que \textit{toute} catégorie $Z$ dans Cat est $W$-asphérique. Donc toute flèche $Z \longrightarrow Z'$ dans Cat est dans $W$, q.e.d.

Supposons maintenant $\pi_0(f)$ surjective, mais non injective. Il y aurait une fibre de $\pi_0(X) \longrightarrow \pi_0(Y) = I$ qui serait de cardinal $> 1$. En prenant le $f_i$ correspondant, on trouve $X_i \longrightarrow Y_i$ avec $Y_i$ 0-connexe et $\card \pi_0(X_i) \geq 2$. Donc on est ramené au cas
\[
f : X \longrightarrow Y \qquad\qquad \left\{
\begin{array}{l}
f \in W \\[1ex]
Y \text{ 0-connexe} \\[1ex]
\card \pi_0(X) \geq 2.
\end{array}
\right.
\]
Considérons la factorisation canonique
\[
X \xrightarrow{i_f} \bar{X} \xrightarrow{\bar{f}} Y,
\]
avec $i_f \in W$, donc $\bar{f} \in W$. Donc par W(4b$''$) les fibres de $\bar{f}$ sont $W$-asphériques. Mais on voit tout de suite, en décomposant $\bar{X}$ suivant les composantes connexes $X_i$ de $X$, que ces fibres ont un $\pi_0$ de cardinalité $\geq \card \pi_0(X) \geq 2$. Donc on est ramené au cas où $Y = e$, i.e.
\[
X \ W\text{-asphérique}, \qquad \card \pi_0(X) \geq 2.
\]
Considérons $X \times X$, qui est aussi $W$-asphérique, et l'application diagonale\msp{page 22}
\[
d : X \longrightarrow X \times X = Y,
\]
qui est donc dans $W$. Mais $\pi_0(d)$ s'identifie à l'application diagonale $\pi_0(X) \longrightarrow \pi_0(X) \times \pi_0(X)$, qui n'est pas surjective puisque $\card \pi_0(X) \geq 2$. Donc pour la première partie, $W = \Fl(\Cat)$.

Donc sous forme plus positive :

\textbf{Corollaire.} \textit{Soit} $W \subseteq \Fl(\Cat)$ \textit{qui satisfait à} W(1,2,3,4b$''$)\textit{, et telle que} $W \neq \Fl(\Cat)$. \textit{Alors $W$ satisfait l'axiome de connexité} W(5)\textit{, et par suite (th. 1, page 12)} $W \subseteq W_\infty$.

% ---- from Der-ChXIIa-012_p021-025.tex
\printedpage{19}
% head: A. Grothendieck : Dérivateurs, Ch. XII, Édition des Universités de Montpellier II et Paris VII
% foot: Version 7 décembre 2002

\section*{III\quad Classes remarquables de foncteurs liés à $W$.}

Finalement, j'appelerai \textit{‘localiseur fondamental’} une partie $W$ de Fl(Cat) satisfaisant les trois conditions suivantes : % [sic] appelerai -> appellerai, spelling

\begin{enumerate}
\item[(i)] Saturation W(1) (cf. rappel p. 1).
\item[(ii)] Axiome de l'élément final W(3) (page 3).
\item[(iii)] La variante affaiblie W(2 bis) de l'axiome de localisation W(2) donnée dans le corollaire page 1 (anciennement c'était L4$'$ dans Pursuing Stacks).
\end{enumerate}

Quand j'aurai besoin de plus sur $W$, je le dirai explicitement. Ces trois axiomes ont suffi pour pratiquement tout ce que j'ai fait avec $W$ dans Pursuing Stacks, la seule exception étant un usage (si je me rappelle) unique, quelque part, d'une variante affaiblie de W(2) (pour les $S$-foncteurs \textit{cartésiens} entre $S$-catégories \textit{fibrées} - si sur toute fibre il est dans $W$, il est lui-même dans $W$).

\dgrblock{dgrb477}{The localizer retained for the geometric argument}
\begin{dcomment}{exposition}{p22:Quand j'aurai besoin de plus sur}{The localizer retained for the geometric argument}
Grothendieck now separates the part of his theory that depends on a notion of weak equivalence from the much stronger requirements imposed on fibrations. This is the point of retaining only W(1), W(3), and W(2 bis). The distinction is consequential: a map can have a well-behaved theory of lifting along arrows without preserving every weak equivalence under pullback. Throughout this section, a category is \emph{$W$-aspherical} when its map to the terminal category belongs to $W$. For $W=W_\infty$, this means that its nerve has weakly contractible realization. The word does not mean that only its higher homotopy groups vanish; connectedness and trivial fundamental group are included.

The notation for comma categories fixes the direction of every assertion that follows. For $f:X\to Y$, an object of $X/y$ is a pair $(x,\alpha:f(x)\to y)$, whereas an object of $y\backslash X$ is $(x,\alpha:y\to f(x))$. The strict fibre $X_y$ contains the objects mapping \emph{equal} to $y$ and only those arrows mapping to $1_y$. Even when several objects of $Y$ are isomorphic to $y$, they do not enter this strict fibre. These three categories solve different problems. The comma categories allow a comparison with $y$; the strict fibre requires equality over $y$.

An aspherical functor has all $X/y$ aspherical. W(2 bis) then says that it belongs to $W$. This is the absolute form of Quillen's Theorem A. W(2), by contrast, tests a map $X\to Y$ over an arbitrary common base $S$ through all $X/s\to Y/s$. The formally stronger axiom will be named again at the corollary on manuscript page 30; it is not imported into the preceding elementary proofs. In later terminology, the distinction is between the absolute and relative Theorem A axioms: \citet[\S1.1]{Maltsiniotis2012Exact} calls the version with the absolute axiom a \emph{weak} fundamental localizer and reserves ``fundamental localizer'' for the relative version. \citet[Theorem 6.1.18]{Cisinski2006} proves that the minimal weak fundamental localizer is $W_\infty$. That minimality theorem does not assert that every weak fundamental localizer satisfies the relative axiom. The manuscript's terminology is still settling here.

The next propositions convert these axioms into geometric tests. Their economy comes from arranging a comparison so that a category with a terminal or initial object appears, and then using two-out-of-three. No model structure on $\Cat$, and no derived direct-image functor, is needed for that argument.
\end{dcomment}


On suppose donc par la suite donné un localiseur fondamental
\[
W \subseteq \mathrm{Fl}(\mathrm{Cat}).
\]

\textbf{Proposition 1.} \textit{Soit $f : X \longrightarrow Y$ dans} \textup{Cat}. \textit{Conditions équivalentes :}

\begin{enumerate}
\item[\textit{a)}] \textit{Pour tout isomorphisme local $Y' \longrightarrow Y$, $f' : X' \longrightarrow Y'$ est $\in W$.}
\item[\textit{b)}] \textit{Même condition, dans le cas particulier $Y' = Y/y \longrightarrow Y$, donc $\forall\ y \in Y$, $X/y \simeq X \times_Y (Y/y)$ est $W$-asphérique.}
\end{enumerate}

\msp{page 23}\textbf{Définition 1.} Sous ces conditions équivalentes a), b) de la proposition 1, on dit que $f$ est \textit{$W$-asphérique.} On dit \texttt{[que]} $f$ est \textit{$W$-coasphérique} si $f^{\circ}$ est $W$-asphérique, i.e. si les $y\backslash X$ sont $W$-asphériques (pour $y \in \Ob Y$).

\textsc{Démonstration.} Il suffit de prouver b) $\Longrightarrow$ a). Par définition d'un isomorphisme local, on voit que la condition b) est stable par changement de base $Y' \longrightarrow Y$, avec $Y'$ un isomorphisme local (puisque les $Y'/y' \xrightarrow{\sim} Y/y$, donc $X'/y' \xrightarrow{\sim} X/y$). Il suffit donc de voir que si $f$ satisfait à b), alors $f \in W$. Mais c'est l'axiome W(2 bis).

\begin{center}
\begin{tikzcd}[row sep=large]
X \arrow[d] & & & \\
Y & & Y' \arrow[ll] & \\
& Y/y \arrow[ul] & & Y'/y' \arrow[ll] \arrow[ul]
\end{tikzcd}
\end{center}

\printedpage{20}
% head: A. Grothendieck : Dérivateurs, Ch. XII, Édition des Universités de Montpellier II et Paris VII
% foot: Version 7 décembre 2002

\textbf{Remarque.} Si $f$ a un adjoint à droite $g$, alors pour tout $y \in Y$, $X/y \simeq X/g(y)$, donc $X/y$ a un objet final et est donc $W$-asphérique. Donc si $f$ a un adjoint à droite (à gauche), alors $f$ est $W$-asphérique ($W$-coasphérique).

Mais même si $f$ n'a pas d'adjoint à droite, il y a quand même un adjoint généralisé
\[
g : Y \longrightarrow X^{\wedge}
\]
et on a
\[
X/y \simeq X/g(y) \qquad \text{(pour l'inclusion pleinement fidèle $X \hookrightarrow X^{\wedge}$),}
\]
donc $\mathrm{Top}(X/y)$ s'identifie au topos induit $\mathrm{Top}(X)/g(y)$. La condition $f$ $W$-asphérique signifie donc que les faisceaux $g(y)$ dans $\mathrm{Top}(X)$ sont
\msp{page 24}
$W$-asphériques. Pour de nombreuses autres formulations équivalentes, je renvoie à Pursuing Stacks.

\dgrblock{dgrb478}{The generalized right adjoint is a presheaf of arrows}
\begin{dcomment}{exposition}{p23:Mais même si $f$ n'a pas d'adjoint à droite}{The generalized right adjoint is a presheaf of arrows}
The passage through presheaves explains why the comma categories are the correct substitutes for an adjoint. Define
\[
 g(y)(x)=\operatorname{Hom}_Y(f(x),y),\qquad
 g:Y\longrightarrow [X^{\mathrm{op}},\mathrm{Set}].
\]
Precomposition makes $g(y)$ a presheaf on $X$, and postcomposition makes $g$ covariant in $y$. Its category of elements has objects $(x,\alpha)$ and arrows $h:x\to x'$ satisfying $\alpha=\alpha'f(h)$. It is therefore exactly $X/y$. If $f\dashv r$, the adjunction identifies $g(y)$ with the representable presheaf $X(-,r(y))$; the category of elements then becomes $X/r(y)$ and has terminal object $1_{r(y)}$. The asphericity of a functor admitting a right adjoint is thus the representable case of the construction in the text.

The asserted identification of topoi can also be made explicit. For any presheaf $P$ on $X$, a presheaf $A\to P$ determines on its category of elements the presheaf
\[
 (x,p)\longmapsto\{a\in A(x):a\mapsto p\}.
\]
Conversely, from a presheaf $B$ on $\int_X P$, recover $A(x)$ as the disjoint union of $B(x,p)$ over all $p\in P(x)$. Restriction along $h:x\to x'$ uses the arrow $(x,P(h)p')\to(x',p')$. These operations are inverse up to their canonical isomorphisms and yield
\[
 [ (\textstyle\int_X P)^{\mathrm{op}},\mathrm{Set}]
 \simeq [X^{\mathrm{op}},\mathrm{Set}]/P.
\]
Taking $P=g(y)$ proves the statement about $\operatorname{Top}(X/y)$. Its use here is interpretive: the failure of an actual right adjoint is measured by the homotopy type of a presheaf of possible arrows, rather than by the absence of a distinguished representing object.

The commentary on Chapter~I, manuscript pages~64--66, discusses local constancy and the Yoneda argument. Here the additional point is that an adjoint is sufficient but far from necessary: it supplies a terminal element of every comma category, while $W$-asphericity requires only their homotopical triviality.
\end{dcomment}


\textbf{Proposition 2.} \textit{Soit $X \xrightarrow{f} Y \xrightarrow{g} Z$ dans} \textup{Cat}. \textit{Si $f$ et $g$ sont $W$-asphériques ($W$-coasphériques), de même $gf$. Si $gf$ et $f$ sont $W$-asphériques ($W$-coasphériques), $g$ aussi.}

\textbf{Rappel.} Un morphisme $p : S' \longrightarrow S$ dans Cat est un \textit{isomorphisme local} si $\forall\ s' \in \Ob S'$, le morphisme induit
\[
p_{s'} : S'/s' \xrightarrow{\sim} S/s \qquad (s = p(s'))
\]
est un isomorphisme : $p$ induit un isomorphisme sur les ‘localisées’ de $S'$ et $S$ en $s'$ et $s = p(s')$ respectivement. Il est immédiat que cette notion est stable par composition, que si $gf$ et $g$ sont des isomorphismes locaux, alors $f$ l'est aussi, que tout isomorphisme est un isomorphisme local. De plus, les isomorphismes locaux sont \textit{stables par changement de base,} ce qui résulte du fait (trivial et important) que si on a un carré cartésien dans Cat

\begin{center}
\begin{tikzcd}[row sep=large]
X & X' \arrow[l, "q"'] \arrow[d, "f'"] \\
S \arrow[from=u, "f"'] & S', \arrow[l, "p"']
\end{tikzcd}
\end{center}

et si $x' \in \Ob X'$ a pour image $x$, $s'$, $s$ dans $X$, $S'$, $S$, alors le diagramme

\begin{center}
\begin{tikzcd}[row sep=large]
X/x \arrow[d, "f/x"'] & X'/x' \arrow[l, "q/x'"'] \arrow[d, "f'/x'"] \\
S/s & S'/s' \arrow[l, "p/s'"']
\end{tikzcd}
\end{center}

est également cartésien. (Donc si $p/s'$ est \texttt{[un]} isomorphisme, $q/x'$ aussi.)

\dgrblock{dgrb479}{Local isomorphisms and the scope of cancellation}
\begin{dcomment}{exposition}{p24:est également cartésien.}{Local isomorphisms and the scope of cancellation}
The ``local isomorphisms'' in this argument are discrete fibrations. Indeed, an isomorphism $S'/s'\to S/p(s')$ says that every arrow ending at $p(s')$ has exactly one lift ending at $s'$, and that its factorization triangles lift uniquely. Conversely, that unique-lifting property gives the displayed isomorphism of slice categories. The categories of elements of presheaves supply the basic examples. In particular, $Y/y\to Y$ is such a map.

This description makes Proposition 1 transparent. If every $X/y$ is $W$-aspherical and $q:Y'\to Y$ is a local isomorphism, then, for $y'\in Y'$ with image $y$, strict pullback gives
\[
 (X\times_Y Y')/y'\cong X/y.
\]
Theorem A applied over $Y'$ puts $X\times_Y Y'\to Y'$ in $W$. Conversely, the particular base changes $Y/y\to Y$ recover the same condition because $Y/y$ has a terminal object.

The cancellation in Proposition 2 uses precisely this class of allowed base changes. Given $X\xrightarrow fY\xrightarrow gZ$ and a local isomorphism $Z'\to Z$, form the three pullbacks $X',Y',Z'$. Its pullback $Y'\to Y$ is again a local isomorphism. Thus asphericity of $f$ gives $f'\in W$. If $g$ is also aspherical, $g'\in W$ and hence $g'f'\in W$. If instead $gf$ is aspherical, then $g'f'\in W$ and two-out-of-three gives $g'\in W$. Proposition 1 now gives the two assertions. The opposite-category argument proves the coaspherical forms.

There is no cancellation statement with $g$ and $gf$ implying asphericity of $f$. For example, let $f$ select the terminal object of $[1]$, and let $g:[1]\to e$. Both $g$ and $gf$ are aspherical, whereas the comma category of $f$ over $0$ is empty. For any $W$ with $\varnothing\to e\notin W$, $f$ is not aspherical. This elementary directional distinction is exactly what the local-isomorphism proof preserves.
\end{dcomment}


\printedpage{21}
% head: A. Grothendieck : Dérivateurs, Ch. XII, Édition des Universités de Montpellier II et Paris VII
% foot: Version 7 décembre 2002

\msp{page 25}Revenons à la proposition 2. La stabilité de la notion de morphisme asphérique par composition résulte aussitôt de la forme a) de la condition. De même pour la condition que si $gf$ et $f$ sont $W$-asphériques, $g$ aussi - il faut montrer que si $Z' \longrightarrow Z$ est \texttt{[un]} isomorphisme local, alors $g'$ est dans $W$, or l'hypothèse implique que $g'f'$ et $f'$ le sont.

\begin{center}
\begin{tikzcd}[row sep=large]
X \arrow[d, "f"'] & X' \arrow[l, "r"'] \arrow[d, "f'"] \\
Y \arrow[d, "g"'] & Y' \arrow[l, "q"'] \arrow[d, "g'"] \\
Z & Z' \arrow[l, "p"']
\end{tikzcd}
\end{center}

\textbf{Corollaire 1.} \textit{L'image inverse d'un $f : X \longrightarrow Y$ qui est $W$-asphérique (resp. $W$-coasphérique) par un isomorphisme local (resp. par un coïsomorphisme local) est $W$-asphérique ($W$-coasphérique).}

\begin{center}
\begin{tikzcd}[row sep=huge]
X \arrow[d] & X' \arrow[l] \arrow[d] & X'' \arrow[l] \arrow[d] \\
Y & Y' \arrow[l, "\text{iso local}"] & Y'' \arrow[l, "\text{iso local}"]
\end{tikzcd}
\end{center}

Résulte formellement de la définition par a) et du fait qu'un composé de deux isomorphismes locaux (resp. coïsomorphismes locaux) est itou.

\textbf{NB} Bien sûr, la notion de foncteur $W$-asphérique n'est pas stable par changement de base quelconque. Par exemple, l'inclusion $X = \{y_0\} \longrightarrow Y$ est $W$-asphérique si et seulement si les catégories discrètes $\mathrm{Hom}(y_0, y)$ sont $W$-asphériques, i.e. (sous réserve que l'on ait W(5), axiome de connexité) si et seulement si les $\mathrm{Hom}(y_0, y)$ sont de cardinal 1, i.e. $y_0$ \textit{objet initial} de $Y$. Mais si $Y \neq \{y_0\}$, prenant $y$ dans $Y$, $y \neq y_0$, et posant $Y' = \{y\} \longrightarrow Y$, on a $X' = X \times_Y Y' = \emptyset$ et $\emptyset \longrightarrow e$ n'est pas asphérique!

\msp{page 26}Par contre, on s'attend que les morphismes $W$-asphériques soient stables par changements de base \textit{$W$-lisses} (qui généralisent les isomorphismes locaux), et les morphismes $W$-coasphériques par changements de base $W$-propres. Cf. plus bas.

\dgrblock{dgrb480}{Properness asks for contractible categories of lifts}
\begin{dcomment}{exposition}{p26:Par contre, on s'attend que les morphismes}{Properness asks for contractible categories of lifts}
The anticipated base-change theorem needs a condition on the map along which one pulls back. Grothendieck finds it by inspecting every possible lift of an arrow. For $f:X\to Y$, an object $x\in X$, and $u:f(x)\to y$, write $L_f(x,u)$ for the category whose objects are arrows $a:x\to z$ with $f(a)=u$ and whose morphisms are arrows $h:z\to z'$ over $1_y$ satisfying $ha=a'$. Its morphisms matter as much as its objects: merely asking for at least one lift would lose the homotopy information.

Let $j_y:X_y\to X/y$ send $z$ to $(z,1_y)$. The object $(x,u)$ of $X/y$ gives an equality of categories
\[
 (x,u)\backslash j_y\cong L_f(x,u).
\]
The same category is the strict fibre over $u$ of $x\backslash X\to f(x)\backslash Y$. These identifications prove the equivalence of conditions (a), (b), and (d) in Proposition 3. The edition's intermediate list ``a) b) c)'' on manuscript page 27 should read ``a) b) d)''; the remaining argument establishes (c).

To see that argument without suppressing an endpoint, pull $f$ back along the arrow $u:[1]\to Y$. Write $E=X\times_Y[1]$. For an object in $E_1$, the comma category into $E_1\hookrightarrow E$ has an initial object. For an object $(x,0)$ in $E_0$, that comma category is $L_f(x,u)$. Consequently $E_1\hookrightarrow E$ is coaspherical exactly when all these lift categories are aspherical. This is why a test over the walking arrow is sufficient for \emph{properness}: the test asks for coasphericity of the inclusion, not merely that its total homotopy type be an equivalence.

The dual incoming-lift category consists of arrows $a:z\to x$ over a fixed $u:y\to f(x)$. Its morphisms are vertical factorizations. Requiring these categories to be $W$-aspherical defines smoothness. The terminology records two directions of the same lifting geometry. The commentary on Chapter~I, manuscript pages~52 and~68, explains the coefficient-theoretic origin; here Grothendieck obtains the criterion directly from $W$.
\end{dcomment}


\textbf{Corollaire 2.} \textit{Soit $f : X \longrightarrow Y$ dans} \textup{Cat} \textit{une} \textup{équivalence} \textit{de catégories. Alors $f$ est $W$-asphérique et $W$-coasphérique.}

En effet, $f$ a un adjoint à gauche et un adjoint à droite, et on utilise la remarque page 23.

\textbf{Proposition 3.} \textit{Soit $f : X \longrightarrow Y$ dans} \textup{Cat}. \textit{Conditions équivalentes :}

\printedpage{22}
% head: A. Grothendieck : Dérivateurs, Ch. XII, Édition des Universités de Montpellier II et Paris VII
% foot: Version 7 décembre 2002

\begin{enumerate}
\item[\textit{a)}] \textit{Pour tout $y \in \Ob Y$, le foncteur d'inclusion}
\[
X_y \longrightarrow X/y \tag*{$(*)$}
\]
\textit{est $W$-}co\textit{asphérique.}
\item[\textit{b)}] \textit{Pour tout $x_0 \in X$, le morphisme induit pour les colocalisées}
\[
x_0\backslash X \longrightarrow y_0\backslash Y \qquad \textit{(où $y_0 = f(x_0)$)}
\]
\textit{est à fibres $W$-asphériques.}
\item[\textit{c)}] \textit{Pour tout morphisme $\mathbf{\Delta}^1 = Y' \longrightarrow Y$, considérons $X' = X \times_Y Y'$ sur $\mathbf{\Delta}^1$, l'inclusion $X'_1 \longrightarrow X'$ de la fibre fermée de $X'$ (fibre sur le point fermé de $Y'$) est $W$-coasphérique.}
\item[\textit{d)}] \textit{Pour toute flèche $u_0 : y_0 \longrightarrow y_1$ dans $Y$, et tout $x_0$ dans $X_{y_0}$, la catégorie $X(x_0;u_0)$ formée des flèches $u' : x_0 \longrightarrow x$ dans $X$ de source $x_0$, et qui relèvent $u$} \texttt{[plutôt} $u_0$\texttt{]}\textit{, est $W$-asphérique.}
\end{enumerate}

\msp{page 27}\textsc{Démonstration.} Pour expliciter la condition a), il faut expliciter la catégorie $\xi_0\backslash(X_y)$, pour un objet quelconque $\xi_0 = (x_0, u_0)$ de $X/y$, $x_0 \in \Ob X$,
\[
u_0 : y_0 \longrightarrow y_1 = y \quad \text{dans $Y$, avec $y_0 = f(x_0)$.}
\]

Cette catégorie $\xi_0\backslash X_{y_1}$ s'explicite alors comme la catégorie des $x_1$ dans $X_{y_1}$, \textit{munis} d'une flèche de $\xi_0$ dans l'image $\xi_1 = (x_1, \mathrm{id}_{f(x_1)} = \mathrm{id}_{y_1})$ de $x_1$ dans $X/y_1$ dans $\xi_0$ \texttt{[‘dans} $\xi_0$\texttt{’ à effacer]}, i.e. d'une flèche $u$ de $x_0$ dans $x_1$, telle que $f(u)$ donne lieu au triangle commutatif dans $Y$, i.e. $f(u) = u_0$. % AUDIT resolved: regular CMTT12 editorial bracket; see corrections.csv, PDF leaf 32

\begin{center}
\begin{tikzcd}
x_1 \arrow[dd, dash] & & \\
& & x_0 \arrow[ull, "u"'] \arrow[dd, dash] \\
y_1 & y_1 \arrow[l, "\mathrm{id}"'] & \\
& & y_0 \arrow[ul, "u_0"'] \arrow[ull, "f(u)"]
\end{tikzcd}
\end{center}

Donc c'est la catégorie notée $X(x_0;u_0)$ dans d). On voit tout de suite qu'elle est aussi isomorphe à la fibre de $x_0\backslash X \longrightarrow y_0\backslash Y$ en l'élément $u_0$ de $y_0\backslash Y$, dont il est question dans b). D'où l'équivalence de a) b) c). D'autre part, la donnée de $u_0$ dans $Y$ équivaut à celle d'un foncteur $\mathbf{\Delta}^1 \longrightarrow Y$, comme dans c), et pour exprimer alors que $\underbrace{X'_1}_{\simeq X_{y_1}} \longrightarrow X'$ est coasphérique, il suffit d'exprimer que les $x'\backslash X'_1$ ($x' \in \Ob X'$) sont $W$-asphériques. Or c'est trivial si $x' \in \Ob X'_1$ (objet initial), donc il suffit de le vérifier pour $x'$ dans $X'\backslash X'_1 = X'_0 \simeq X_{y_0}$. Alors $x'$ s'identifie à un élément $x_0$ de $X_{y_0}$, et on vérifie tout de suite que $x_0\backslash X'_1$ est isomorphe à la catégorie $X(x_0;u_0)$ de d). Cela achève la démonstration.

\printedpage{23}
% head: A. Grothendieck : Dérivateurs, Ch. XII, Édition des Universités de Montpellier II et Paris VII
% foot: Version 7 décembre 2002

\msp{page 28}\textbf{Définition 2.} Sous les conditions de la proposition 3, on dit que $f$ est \textit{$W$-propre.} On dit que $f$ est \textit{$W$-lisse} si $f^{\circ}$ est $W$-propre, i.e. si les morphismes d'inclusion $(*)$ \texttt{[plutôt} $X_y \longrightarrow y\backslash X$\texttt{]} sont $W$-asphériques etc. (conditions duales de celles de la proposition 3).

\textbf{Remarque.} La condition a) de la proposition 3 est satisfaite en particulier si l'inclusion $(*)$ a un adjoint à gauche (cf. remarque page 23). Cela signifie que les catégories $X(x_0;u)$ ont un objet initial. Cela signifie aussi que la catégorie $X$ sur $Y$ est Cat-cofibrée sur $Y$ au sens faible, i.e. qu'il existe, pour $u_0 : y_0 \longrightarrow y_1$ flèche dans $Y$, et $x_0$ dans $X_{y_0}$, un objet $u_{0*}(x_0)$ dans $X_{y_1}$ qui représente le foncteur covariant $X_{y_1} \longrightarrow (\mathrm{Ens})$, $x_1 \longmapsto \mathrm{Hom}_u(x_0, x_1)$. Ainsi, si $X$ est Cat-cofibrée au sens faible sur $Y$, et a fortiori si elle est Cat-cofibrée, alors $f : X \longrightarrow Y$ est $W$-propre. Dualement, si $f$ est faiblement Cat-fibrée, $f$ est \texttt{[}$W$\texttt{-]}lisse.

\dgrblock{dgrb481}{Representable lifts and their homotopical relaxation}
\begin{dcomment}{examples}{p28:La condition a) de la proposition 3}{Representable lifts and their homotopical relaxation}
Grothendieck first places ordinary fibred categories inside the new theory. Fix $x\in X$ and $u:f(x)\to y$. An initial object $a:x\to u_*x$ of $L_f(x,u)$ gives, for every $z\in X_y$, a bijection
\[
 \operatorname{Hom}_{X_y}(u_*x,z)
 \xrightarrow{\ \sim\ }
 \{b:x\to z:f(b)=u\},\qquad h\longmapsto ha.
\]
Conversely, such a representing arrow is initial in $L_f(x,u)$. Choosing it for all $(x,u)$ is the weak cofibred condition used here. It only asserts this factorization property for arrows over the fixed arrow $u$; the stronger condition on composition is part of the usual opfibration formalism. Either condition gives properness, since initial categories are $W$-aspherical. The dual argument gives smoothness for prefibrations.

A useful calculation shows how much weaker asphericity can be than representability. Let $C$ be the poset with three objects $b,d,e$ and relations $b<d$, $b<e$. Form its left cone $X=C^{\lhd}$ by adjoining an initial object $a$, and let $p:X\to[1]$ send $a$ to $0$ and $C$ to $1$. The outgoing-lift category from $a$ over $0\to1$ is $C$ itself. It is contractible because $b$ is initial. This example is already a weak opfibration. To remove representability while retaining contractibility, replace $C$ by the four-object fence
\[
 b<d,\qquad e<d,\qquad e<c.
\]
Its nerve is the path $b-d-e-c$, hence contractible, but it has no initial object. In the corresponding cone, $L_p(a,0\to1)\cong C$ is therefore aspherical without being representable by a universal lift. For every identity in the base, the lift category is a coslice in the appropriate fibre and has an initial object. There are no other base arrows to test. Thus this $p$ is $W_\infty$-proper but is not weakly opfibred.

The dual cone gives a smooth map without the prefibration property. \citet[Proposition 4.3 and Example 4.17]{Maltsiniotis2005Asphericite} places these criteria in the broader language of asphericity structures; for a structure coming from a fundamental localizer, the resulting weak smoothness criterion recovers $W$-smoothness. For an arbitrary asphericity structure those notions need not agree.
\end{dcomment}


\textbf{Proposition 4.} (\footnote[43]{Dire aussi : Si $X \xrightarrow{f} Y \xrightarrow{g} Z$, si $gf$ et $f$ $W$-lisses, alors $g$ est lisse en les points de $\mathrm{Im}\, f$, idem pour propre.}) \textit{La notion de morphisme $W$-propre est stable par composition, par changement de base. Les coisomorphismes locaux sont propres (trivial sur b)), a fortiori les isomorphismes sont propres, les immersions fermées sont propres.}

\textit{Énoncé dual pour les morphismes lisses : stabilité par changement de base, par composition, les isomorphismes locaux sont lisses, notamment les immersions ouvertes.}

\msp{page 29}La stabilité par changement de base est immédiate sur la forme c). La stabilité par composition \texttt{[est]} moins triviale :

\begin{center}
\begin{tikzcd}[row sep=large]
{\scriptstyle x} &[-1.5em] X \arrow[d, "f"] & x\backslash X \arrow[l] \arrow[d] & \bar{X} \arrow[l] \arrow[d] \\
{\scriptstyle y} & Y \arrow[d, "g"] & y\backslash Y \arrow[l] \arrow[d] & \bar{Y} \arrow[l] \arrow[d] \\
{\scriptstyle z} & Z & z\backslash Z \arrow[l] & \{u\} \arrow[l]
\end{tikzcd}
\end{center}

Dans le diagramme à carrés cartésiens ci-contre \texttt{[plutôt dont les deux carrés de droite sont cartésiens]} (déduit de $x \in \Ob X$, d'où $y = f(x)$, $z = g(y) = gf(x)$, et de $u \in \Ob z\backslash Z$), il faut prouver que $\bar{X}$ est $W$-asphérique. On sait que $\bar{Y}$ l'est ($g$ propre) et que $\bar{X} \longrightarrow \bar{Y}$ est à fibres $W$-asphériques. Il suffit donc d'utiliser la proposition 5 ci-dessous, et la remarque immédiate % AUDIT resolved: regular CMTT12 editorial bracket; see corrections.csv, PDF leaf 33

\textbf{Lemme.} \textit{Si $f$ est propre, il en est de même de $x\backslash f : x\backslash X \longrightarrow y\backslash Y$, pour tout $x$ dans $X$, $y = f(x)$.}

% ---- from Der-ChXIIa-012_p026-030.tex
\printedpage{24}
% head: A. Grothendieck : Dérivateurs, Ch. XII, Édition des Universités de Montpellier II et Paris VII
% foot: Version 7 décembre 2002
En effet, si $u'$ est dans $x\backslash X$ et $v'$ dans $y\backslash Y$, et si $x'$, $y'$ sont leurs images dans $X$, $Y$ respectivement, alors
\[
u'\backslash(x\backslash X)\simeq x'\backslash X,\qquad v'\backslash(y\backslash Y)\simeq y'\backslash Y,
\]
et la flèche entre les deux s'identifie à la flèche de localisation
\[
x'\backslash X\longrightarrow y'\backslash Y.
\]

\dgrblock{dgrb482}{The composition proof has no circular dependence}
\begin{dcomment}{exposition}{p29:et la flèche entre les deux s'identifie à la flèche de localisation}{The composition proof has no circular dependence}
The proof of composition proceeds by factoring a category of lifts through another category of lifts. It uses Proposition 5, but Proposition 5 itself needs only the definition of properness and the already established base-change stability. Thus its appearance after Proposition 4 does not create a circular proof.

First prove the part of Proposition 5 that will be used. If $q:A\to B$ is proper and every strict fibre $A_b$ is aspherical, then $A_b\to A/b$ is coaspherical and hence belongs to $W$. Two-out-of-three makes $A/b$ aspherical, and W(2 bis) gives $q\in W$. The same reasoning applies after every base change, because properness and the asphericity of the fibres persist. The smooth version follows by taking opposites.

Now let $X\xrightarrow fY\xrightarrow gZ$ be proper. Fix $x\in X$ and $u:gf(x)\to z$. There is a functor
\[
 L_{gf}(x,u)\longrightarrow L_g(f(x),u),\qquad a\longmapsto f(a).
\]
It is the base change of $x\backslash f:x\backslash X\to f(x)\backslash Y$ along the fibre of $f(x)\backslash g$ over $u$. The coslice functor $x\backslash f$ is proper: taking the coslice at an arrow $x\to x'$ reduces its incoming data to $x'\backslash X\to f(x')\backslash Y$, precisely the original lifting test. Its displayed base change is therefore proper. The fibre over a lift $v:f(x)\to y'$ is $L_f(x,v)$, which is aspherical. Proposition 5 makes the displayed functor a $W$-equivalence. Since its target is aspherical by properness of $g$, its source is aspherical too. This proves properness of $gf$.

The same diagram proves the observation in footnote 43. If $f$ and $gf$ are proper, its source is aspherical, and the displayed proper functor still has aspherical fibres. Its target is therefore aspherical. This proves the lifting condition for $g$ at every object $f(x)$ in the strict image of $f$; it does not assert anything about objects outside that image. Duality proves the smooth statement. These arguments require W(2 bis), not the full relative axiom W(2).
\end{dcomment}


\textbf{Proposition 5.} \textit{Soit $f : X\longrightarrow Y$ dans $\Cat$, supposons $f$ $W$-propre ou $W$-lisse, et à fibres $W$-asphériques, alors $f\in W$ (et même $f$ est} \textup{universellement} \textit{dans $W$, i.e. il est dans $W$ et le reste par tout changement de base).}

\leavevmode\msp{page 30}Comme les hypothèses sont stables par changement de base, il suffit de prouver $f\in W$. Supposons p.~ex. $f$ propre, prouvons que $f$ est $W$-\textit{asphérique}, i.e. que les $X/y$ ($y\in\Ob\,Y$) sont $W$-asphériques. Mais $X_y\longrightarrow X/y$, étant $W$-coasphérique par hypothèse de propreté, est dans $W$ (par W(2 bis)), et comme $X_y$ est $W$-asphérique par hypothèse, $X/y$ aussi. Mais par W(2 bis), $f$ $W$-asphérique $\Longrightarrow f\in W$. Le cas $f$ lisse est dual.

\textbf{Remarque.} Ainsi, si $f$ est propre, s'il est à fibres $W$-asphériques il est `universellement $W$-asphérique', i.e. est et reste $W$-asphérique par tout changement de base. Inversement, si $f$ est universellement dans $W$ (pas forcément propre), alors en particulier ses fibres sont $W$-asphériques. Mais ça n'implique bien sûr pas que $f$ soit propre, ou lisse (\footnote[44]{Mais `universellement dans $W$' implique déjà `universellement $W$-asphérique', puisque la notion de $W$-asphéricité est définie par \texttt{[la]} proposition 1 a).}).

\dgrblock{dgrb483}{A universal equivalence with neither lifting property}
\begin{dcomment}{examples}{p30:Mais ça n'implique bien sûr pas}{A universal equivalence with neither lifting property}
The qualification after Proposition 5 has a small model which makes the role of the lifting condition visible. Let $P$ be the poset with objects $u,v,w,z$ and exactly the strict relations
\[
 u<v,\qquad u<z,\qquad w<z.
\]
Define $p:P\to[1]$ by $p(u)=p(v)=0$ and $p(w)=p(z)=1$. Both fibres are intervals, and the nerve of $P$ is the path $v-u-z-w$. More is true than the contractibility of these three nerves: $p$ has a contraction over its base.

The section $s:[1]\to P$ selects $u$ and $z$. Let $R:P\to P$ send $v$ to $u$ and fix $u,w,z$. Each displayed relation is preserved, so $R$ is a functor over $[1]$. The pointwise inequalities give natural transformations over $[1]$
\[
 \operatorname{id}_P\longleftarrow R\longrightarrow sp.
\]
For the first transformation use $u\to v$ at $v$; for the second use $w\to z$ at $w$; all remaining nonidentity components are absent. Since $ps=\operatorname{id}_{[1]}$, these transformations exhibit $p$ as a homotopy equivalence of categories over $[1]$. After any base change $T\to[1]$, the section and both transformations pull back, so $P\times_{[1]}T\to T$ is still a homotopy equivalence. It therefore belongs to every $W$ considered in this section.

Nevertheless $p$ is not $W_\infty$-proper: the arrow $0\to1$ has no outgoing lift from $v$, so its lift category is empty. Nor is it $W_\infty$-smooth: the same arrow has no incoming lift ending at $w$. Thus universal membership in $W$ does not imply either geometric lifting property. The same failure holds for any $W$ for which the empty category is not aspherical.

This also verifies the terminology in footnote 44. A universal $W$-equivalence is universally aspherical because, after any base change and any further local-isomorphism test, one is still taking a base change of the original map. The four-object example is universally aspherical and universally coaspherical without being proper or smooth. The latter properties inspect the lift category at each particular object, which a relative contraction may move.
\end{dcomment}


\textbf{Corollaire.} \textit{Si $W$ satisfait} \textup{W(2) (pas seulement W(2 bis))}\textit{, alors pour un $S$-morphisme $X\longrightarrow Y$ de catégories qui sont toutes deux $W$-propres sur $S$, ou toutes deux $W$-lisses, si $f_s : X_s\longrightarrow Y_s$ est dans $W$ pour tout $s\in S$, alors $f\in W$.}

\leavevmode\msp{page 31}Supposons p.~ex. $f$ propre, considérons, pour $s\in\Ob\,S$
\[
\begin{tikzcd}[row sep=large, column sep=large]
X_s \arrow[r, "f_s"] \arrow[d, "i_s^X"'] & Y_s \arrow[d, "i_s^Y"] \\
X/s \arrow[r, "f/s"] & Y/s,
\end{tikzcd}
\]
comme $X$, $Y$ \texttt{[sont]} propres, les flèches verticales sont dans $W$ (il suffirait de beaucoup moins que la propreté \dots), donc $f_s\in W\Longrightarrow f/s\in W$. Ainsi, `$f$ est dans $W$ localement sur $S$', ce qui implique $f\in W$ par W(2).

\dgrblock{dgrb484}{Fibrewise comparison uses the structural maps}
\begin{dcomment}{exposition}{p31:les flèches verticales sont dans $W$}{Fibrewise comparison uses the structural maps}
The corollary compares two families over $S$; it does not assume that the comparison $f:X\to Y$ is itself proper. The phrase ``Supposons p.~ex.\ $f$ propre'' at the start of the proof is an edition slip. The case being considered is that both structural maps $X\to S$ and $Y\to S$ are proper. This is confirmed by the next sentence and by the two vertical arrows in the displayed square.

In fact, the proof needs only
\[
 X_s\longrightarrow X/s\in W,
 \qquad Y_s\longrightarrow Y/s\in W
 \quad\text{for every }s\in S.
\]
If the fibre map $f_s$ belongs to $W$, the commutative square and two-out-of-three imply $f/s\in W$. The remaining step from all these slice comparisons to $f\in W$ is exactly the relative axiom W(2). Replacing all slices by coslices gives the smooth case, using the opposite form of the same axiom.

This distinction will persist when properness is weakened to preproperness. The geometric hypothesis supplies comparisons between fibres and comma categories; the relative localization axiom assembles the comma comparisons over the whole base. There is no appeal to the converse assertion that a weak equivalence between total categories must be a weak equivalence on fibres. That substantially stronger assertion belongs to the fibration discussion later in this chapter.
\end{dcomment}


\printedpage{25}
% head: A. Grothendieck : Dérivateurs, Ch. XII, Édition des Universités de Montpellier II et Paris VII
% foot: Version 7 décembre 2002
\textbf{Proposition 6.} \textit{Soit}
\[
\begin{tikzcd}[row sep=large, column sep=large]
X & X' \arrow[l, "q"'] \arrow[d, "f'"] \\
S \arrow[from=u, "f"'] & S', \arrow[l, "p"]
\end{tikzcd}
\]
\textit{cartésien dans} $\Cat$. \textit{Si $f$} \texttt{[est]} \textit{$W$-asphérique, $p$ $W$-lisse, alors $f'$ est $W$-asphérique.}

Il faut prouver \texttt{[que]} $X'\times_{S'}S'/s'$ est $W$-asphérique ($s'\in\Ob\,S'$). Soit $s=p(s')$, on a
\[
\begin{tikzcd}[column sep=small]
S & & S' \arrow[ll, "p"'] & \\
& S/s \arrow[ul] & & S'/s', \arrow[ul] \arrow[ll, "p_0"]
\end{tikzcd}
\]
avec $p_0$ $W$-lisse (lemme p.~28 \texttt{[plutôt 29]}) et à fibres $W$-asphériques, donc une $W$-équivalence universelle (prop.~5). Comme l'image inverse de $X$ sur $S/s$ est $W$-asphérique (hypothèse pour $f$), il \texttt{[s'ensuit]}

\leavevmode\msp{page 32}que dans le carré cartésien ci-contre \texttt{[ci-dessous]}
\[
(*)\qquad\qquad
\begin{tikzcd}[row sep=large, column sep=large]
X/s \arrow[d] & X'/s' \arrow[l, "W"'] \arrow[d] \\
S/s & S'/s' \arrow[l, "W"']
\end{tikzcd}
\]
les flèches horizontales sont dans $W$, et la première flèche verticale l'est par hypothèse sur $f$, donc la seconde aussi, donc $X'/s'$ \texttt{[est]} $W$-asphérique.

\dgrblock{dgrb485}{Why smooth base change preserves asphericity}
\begin{dcomment}{exposition}{p32:les flèches horizontales sont dans $W$, et la première flèche verticale}{Why smooth base change preserves asphericity}
Grothendieck's anticipated stability theorem now follows from the same categories of lifts, with the two variables interchanged. For a smooth $p:S'\to S$ and $s'\in S'$, put $s=p(s')$ and consider
\[
 p_{s'}:S'/s'\longrightarrow S/s.
\]
The dual of the coslice lemma makes $p_{s'}$ smooth. Its strict fibre over an arrow $v:t\to s$ is the category of incoming lifts of $v$ ending at $s'$. Smoothness of $p$ says that this category is aspherical. Proposition 5 therefore makes $p_{s'}$ a universal $W$-equivalence.

For an arbitrary $X\to S$, the square
\[
\begin{tikzcd}
 (X\times_S S')/s'\ar[r]\ar[d]&X/s\ar[d]\\
 S'/s'\ar[r,"p_{s'}"']&S/s
\end{tikzcd}
\]
is a strict pullback. An object of either description of its upper-left corner is a pair $(x,t')$ together with an arrow $t'\to s'$ and the equality $f(x)=p(t')$; the arrow conditions agree as well. Its top map belongs to $W$, independently of any assumption on $X$. If $X\to S$ is aspherical, $X/s$ is aspherical, so the top map makes the upper-left category aspherical. This proves Proposition 6.

For the corollary, apply this same universal equivalence once to $X$ and once to $Y$. If $f:X\to Y$ over $S$ satisfies $f/s\in W$, the commutative comparison square with horizontal arrows in $W$ gives $f'/s'\in W$. Thus smooth pullback preserves $W_S^{\mathrm g}$, the class detected on all over-categories. Opposite categories give the preservation of $W_S^{\mathrm d}$ by proper pullback. Neither argument needs the extra inference $W_S^{\mathrm g}\subseteq W$; accordingly the proof works with W(2 bis) alone.
\end{dcomment}


\textbf{Corollaire} (généralise \texttt{[la]} prop.~6, obtenue en faisant $Y=S$). \textit{Soit $p : S'\longrightarrow S$} \texttt{[un]} \textit{morphisme $W$-lisse dans} $\Cat$. \textit{Alors le foncteur changement de base}
\[
\Cat/S\longrightarrow\Cat/S'
\]
\textit{est compatible avec les localiseurs $W^{\mathrm{g}}_S$, $W^{\mathrm{g}}_{S'}$, i.e. transforme $f : X\longrightarrow Y$ qui est dans $W$} localement \textit{sur $S$, dans $f' : X'\longrightarrow Y'$ dans $W$ localement sur $S'$.}

Cette fois, on regarde le diagramme similaire à $(*)$
\[
\begin{tikzcd}[column sep=small]
X/s \arrow[dd] \arrow[dr, "f/s"] & & X'/s' \arrow[ll, "W"'] \arrow[dd] \arrow[dr, "f'/s'"] & \\
& Y/s \arrow[dl] & & Y'/s' \arrow[ll, "W"', crossing over] \arrow[dl] \\
S/s & & S'/s', \arrow[ll, "W"'] &
\end{tikzcd}
\]

\printedpage{26}
% head: A. Grothendieck : Dérivateurs, Ch. XII, Édition des Universités de Montpellier II et Paris VII
% foot: Version 7 décembre 2002
où les flèches horizontales sont dans $W$, donc $f/s\in W$ (hypothèse) implique $f'/s'\in W$.

\leavevmode\msp{page 33}Dualement, si $p$ est $W$-propre, le foncteur image inverse est compatible avec les localiseurs $W^{\mathrm{d}}_S$, $W^{\mathrm{d}}_{S'}$.

\dgrblock{dgrb486}{The passage from lifting categories to exact squares}
\begin{dcomment}{later}{p33:Dualement, si $p$ est $W$-propre}{The passage from lifting categories to exact squares}
The geometry just proved is the ingredient that later formulations package as exactness of squares. The term ``proper'' is justified by a universal base-change property; ``smooth'' describes the corresponding property when it is the map of bases that varies. The connection with the geometric proper and smooth base-change theorems, already discussed in the commentary on I.72, is made explicit in \citet[Introduction and Theorem 1.2]{Maltsiniotis2005Asphericite}.

The precise categorical statement appears in \citet[Propositions 2.12--2.13]{Maltsiniotis2012Exact}: a cartesian square with proper right side or smooth bottom side is $W$-exact, and the universal square condition characterizes properness. Its derivator interpretation is given by Theorem 4.32 of that paper, where cohomological base change is characterized through the appropriate weak exactness condition. These assertions explain why the lifting categories in the present text are sufficient; they are the categories appearing in pointwise base-change comparisons.

There is a convention to keep visible when comparing modern derivator texts. Grothendieck's coefficients here are contravariant diagrams. If a derivator is written with covariant diagrams instead, opposite categories exchange fibrations and opfibrations and interchange the two directions of the statement. In that convention \citet[Propositions 1.24 and 1.26]{Groth2013} proves base change for cartesian squares with a Grothendieck fibration on the right or an opfibration at the bottom, and identifies this property with the pointwise Kan-extension axiom. The apparent reversal has this variance explanation.

These developments do not identify every strict pullback with a homotopy pullback. An exact square controls a specified comparison of Kan extensions. Here the point is that properness or smoothness makes the relevant comma-category comparison cofinal in the required direction. The construction can therefore be transported to other coefficient theories without choosing path objects in each one separately.
\end{dcomment}


Si $f : X\longrightarrow S$ est $W$-propre, on va définir un foncteur canonique
\[
\begin{aligned}
\hot_*^{X/S,W} : S &\quad\longrightarrow\quad \Hot_W\overset{\text{déf}}{=}(\Cat)W^{-1}\\
\text{donné sur les objets par}\qquad s &\quad\longmapsto\quad \hot_W(X_s).
\end{aligned}
\]
Il faut définir, pour une flèche $u$ dans $S$, $u : s\longrightarrow s'$, une flèche
\[
\hot_W(X_s)\longrightarrow\hot_W(X_{s'})
\]
dans $\Hot_W$ - flèche que nous pouvons noter $\hot_*^{X/S,W}(u)$. On la définit par le diagramme
\[
\begin{tikzcd}[row sep=large, column sep=large]
X_s \arrow[r, dotted] \arrow[d, "i_s^X"] & X_{s'} \arrow[d, "i_{s'}^X"] \\
X/s \arrow[r, "X/u"] & X/s',
\end{tikzcd}
\]
où les flèches verticales sont $W$-coasphériques, donc dans $W$. Ce sont donc des isomorphismes dans $\Hot_W$, d'où l'unique flèche en pointillés dans $\Hot_W$, rendant le carré commutatif. Sur cette définition il est immédiat que
\[
u\longmapsto\hot_W(u)
\]
est compatible avec composition, et unités.

Ce foncteur est donc défini en supposant

\leavevmode\msp{page 34}seulement que les
\[
X_y\longrightarrow X/y
\]
sont dans $W$, ce qui est nettement moins fort que propre. Si on suppose que cette hypothèse est satisfaite `universellement' (i.e. reste satisfaite après tout changement de base $S'\longrightarrow S$) (\footnote[45]{On dira que $f$ est $W$-\textit{prépropre.}}), alors on a donc, pour tout $S'$ dans $\Cat/S$, un foncteur
\[
\hot_*^{X'/S',W} : S'\longrightarrow\Hot_W,
\]
et on vérifie tout de suite que ce n'est autre que le composé dans

\printedpage{27}
% head: A. Grothendieck : Dérivateurs, Ch. XII, Édition des Universités de Montpellier II et Paris VII
% foot: Version 7 décembre 2002
\[
(*)\qquad\qquad
\begin{tikzcd}[column sep=small]
S' \arrow[rr, "p"] \arrow[dr, "\hot_W^{X'/S'}"'] & & S \arrow[dl, "\hot_W^{X/S}"] \\
& \Hot_W &
\end{tikzcd}
\]
\texttt{[plutôt} $\hot_*^{X'/S',W}$ \texttt{resp.} $\hot_*^{X/S,W}$\texttt{]}. On a des observations duales si $X$ est lisse sur $S$, plus généralement si les inclusions
\[
X_s\longrightarrow s\backslash X
\]
sont dans $W$. On trouve alors un contrafoncteur
\[
\left\{
\begin{aligned}
S^\circ &\quad\longrightarrow\quad \Hot_W\\
s &\quad\longmapsto\quad \hot_W(X_s)\\
u &\quad\longmapsto\quad \hot^*_{X/S,W}(u)\ \text{ou}\ \hot^*_{X/S}(u)
\end{aligned}
\right.
\]

\leavevmode\msp{page 35}et une compatibilité similaire à $(*)$, si l'hypothèse est satisfaite \textit{universellement} sur $S$, p.~ex. $f$ $W$-lisse.

\dgrblock{dgrb487}{Transport obtained by conjugating the comma functor}
\begin{dcomment}{exposition}{p35:et une compatibilité similaire à}{Transport obtained by conjugating the comma functor}
To construct transport, Grothendieck uses a functor that already exists before localization. Let $Q:\Cat\to\Hot_W$ be the localization. The assignment
\[
 s\longmapsto X/s,\qquad
 (u:s\to t)\longmapsto(X/u):(x,\alpha)\longmapsto(x,u\alpha)
\]
is a strict functor $S\to\Cat$. If $i_s:X_s\to X/s$ belongs to $W$ for every $s$, define
\[
 \lambda(u)=Q(i_t)^{-1}Q(X/u)Q(i_s).
\]
The existence of the inverse is the only use of the hypothesis. For $s\xrightarrow u t\xrightarrow v r$, multiplication cancels $Q(i_t)$ against its inverse and gives
\[
 \lambda(v)\lambda(u)
 =Q(i_r)^{-1}Q(X/v)Q(X/u)Q(i_s)
 =\lambda(vu).
\]
The identity case follows in the same way. Thus the word ``functor'' is justified by the comma construction, not by a choice of lifts in the strict fibres.

For base change $q:S'\to S$, the canonical maps $X'/s'\to X/q(s')$ form a natural transformation of these comma diagrams, and the square with $X'_{s'}\cong X_{q(s')}$ and the two inclusions commutes. If both inclusions belong to $W$, two-out-of-three puts the comma comparison in $W$ and the displayed formula gives precisely the compatibility asserted in the text. Requiring the inclusion hypothesis after every base change is the definition of $W$-preproperness in footnote 45. Opposite categories give $W$-presmoothness, called $W$-prélissité in the source, and a contravariant transport functor.

Properness is a practical sufficient condition because it says more: each $i_s$ is coaspherical. Preproperness keeps the universal weak-equivalence conclusion and discards its local lifting explanation. Neither should be confused with the condition on the two endpoint inclusions over individual arrows considered on manuscript page 36. As the development below shows, those individual comparisons need not obey the displayed composition law.
\end{dcomment}


\textbf{Proposition 7.} \textit{Supposons $f : X\longrightarrow S$ à la fois $W$-propre et $W$-lisse, plus généralement que les morphismes}
\[
X_s\longrightarrow X/s,\qquad X_s\longrightarrow s\backslash X\qquad(s\in S)
\]
\textit{soient dans $W$, et que cette hypothèse reste verifiée après tout changement de base. Alors les foncteurs correspondants} % [sic] verifiée -> vérifiée, accent missing
\[
\begin{aligned}
\lambda &\;=\;\hot_*^{X/S,W} : &&\quad S \quad\longrightarrow\quad \Hot_W\\
\lambda^\circ &\;=\;\hot^*_{X/S,W} : &&\quad S^\circ \quad\longrightarrow\quad \Hot_W
\end{aligned}
\]
\textit{transforment toute flèche $u$ de $S$ en} \texttt{[un]} \textit{isomorphisme de} $\Hot_W$\textit{, et on a}
\[
\lambda^\circ(u)=\lambda(u)^{-1}\qquad\textit{dans}\quad\Hot_W,
\]
\textit{i.e. $\lambda(u)$ et $\lambda^c irc(u)$ sont inverses l'un de l'autre.} % [sic] λ^c irc(u) -> λ°(u), command \circ misprinted as superscript c followed by "irc"

\textsc{Démonstration}. La donnée de $u : s_0\longrightarrow s_1$ équivaut à celle de
\[
S'=\mathbf{\Delta}^1\longrightarrow S,
\]
d'où
\[
X'=X\times_S S'\qquad\text{sur } S'=\mathbf{\Delta}^1,
\]
avec les deux fibres
\[
X'_0\simeq X_{s_0},\qquad X'_1\simeq X_{s_1}.
\]
L'hypothèse assurant l'existence de $\lambda$ s'exprime par la condition que
\[
X'_1\longrightarrow X'/1\simeq X'
\]
\printedpage{28}
% head: A. Grothendieck : Dérivateurs, Ch. XII, Édition des Universités de Montpellier II et Paris VII
% foot: Version 7 décembre 2002
\msp{page 36}est dans $W$, de sorte que le diagramme
\[
X'_0\xrightarrow{i_0}X'\xleftarrow[W]{i_1}X'_1
\]
définit une flèche $\lambda(u) : X'_0\longrightarrow X'_1$ dans $\Hot_W$, savoir $\lambda(u)=i_1^{-1}i_0$. L'hypothèse qui assure l'existence de $\lambda^\circ$ s'exprime par la condition symétrique $i_0\in W$, qui permet de définir la flèche $\lambda^\circ(u) : X'_1\longrightarrow X'_0$ par $\lambda^\circ(u)=i_0^{-1}i_1$. Quand les deux conditions sont satisfaites simultanément, on voit donc que $\lambda(u)$, $\lambda^\circ(u)$ sont inverses l'un de l'autre.

\dgrblock{dgrb488}{Two directions of transport give monodromy}
\begin{dcomment}{examples}{p36:Quand les deux conditions sont satisfaites simultanément}{Two directions of transport give monodromy}
Proposition 7 says that the two canonical transports become inverse when both comma comparison hypotheses hold universally. Pulling back over $u:[1]\to S$ gives $E=X\times_S[1]$ and inclusions $i_0:X_{s_0}\to E$, $i_1:X_{s_1}\to E$. Preproperness makes $i_1$ a $W$-equivalence, and presmoothness makes $i_0$ one. The already constructed functors therefore satisfy
\[
 \lambda(u)=Q(i_1)^{-1}Q(i_0),\qquad
 \lambda^{\mathrm{op}}(u)=Q(i_0)^{-1}Q(i_1).
\]
These formulas prove the claimed inverse relation without an additional saturation assumption. The covariant functor factors through the ordinary groupoid completion $S[\mathrm{Mor}(S)^{-1}]$. It records transport around loops, or monodromy, at the level of the homotopy category.

For a concrete example, let $C_2=\{1,\sigma\}$ act on the discrete set $F=\{0,1\}$ by interchange. Regard $BC_2$ as the one-object groupoid. The action groupoid $E$ has objects $0,1$ and an arrow labelled $g$ from $x$ to $gx$. The projection $E\to BC_2$ is both a fibration and an opfibration: every base arrow has a unique incoming and a unique outgoing lift once the appropriate endpoint is specified. Its strict fibre is $F$. The transport along $\sigma$ interchanges the two points, in both directions, since $\sigma^{-1}=\sigma$.

The action is free and transitive, so there is exactly one arrow between each ordered pair of objects of $E$. Either object is consequently both initial and terminal, and $BE$ is contractible. The base has fundamental group $C_2$, while the fibre consists of two points. This is the usual universal two-sheeted covering expressed entirely by categories. It illustrates why the fibre functor and its monodromy carry information that the homotopy type of the total category alone does not retain. It also illustrates the scope of Proposition 7: it gives a functor into $\Hot_W$, not a claim that ordinary homotopy-category monodromy classifies all homotopy-coherent families.
\end{dcomment}


\textbf{Remarques.} Pour définir $\lambda$ `universellement', il n'était pas indispensable de considérer les $X/s$ et les inclusions $X_s\longrightarrow X/s$, au lieu de cela il suffisait de regarder les changements de base $S'=\mathbf{\Delta}^1\xrightarrow{u}S$, et de faire l'hypothèse que pour tout tel changement de base, l'inclusion $X'_1\xrightarrow{i_1}X'=X_{(\mathbf{\Delta}^1,u)}$ est dans $W$. Cette hypothèse est automatiquement stable par changement de base arbitraire. (La $W$-propreté de $f : X\longrightarrow S$ signifie que $X'_1\longrightarrow X'$ est non seulement dans $W$, mais même $W$-coasphérique.) Dualement, pour la définition `universelle' de $\lambda^\circ$, il suffirait que pour tout changement de base $\mathbf{\Delta}^1\longrightarrow S$ comme dessus, l'inclusion $X'_0\xrightarrow{i_0}X'$ est dans $W$.

\dgrblock{dgrb489}{An eight-object obstruction to edge transport}
\begin{dcomment}{development}{p36:Cette hypothèse est automatiquement stable par changement de base arbitraire.}{An eight-object obstruction to edge transport}
Grothendieck proposes to replace the universal comma comparisons by endpoint comparisons over the walking arrow. Those tests do define a morphism for each arrow of the base, but they do not by themselves make these morphisms functorial. The distinction is already visible on connected components, and also answers the question about pseudofibrations on manuscript page 40.

\textbf{Proposition (edge tests do not force coherence).}\ours
There exists a functor $p:P\to[2]$ such that, for every $u:[1]\to[2]$, both endpoint inclusions into $P\times_{[2]}[1]$ belong to $W_\infty$, but the resulting edge transports fail the composition law. Moreover, $p$ does not satisfy the fibration test over contractible bases.

\emph{Proof.} Indices will be taken modulo $2$. Let $P$ have eight objects
\[
 a_0,a_1,b_0,b_1,d_0,d_1,c_0,c_1,
\]
and exactly the following nonidentity order relations:
\[
 a_i<d_i,\qquad b_i<d_i,\qquad
 b_i<c_i,\qquad a_i<c_{1-i}\quad(i=0,1).
\]
Every relation starts at an $a_i$ or a $b_i$ and ends at a $d_i$ or a $c_i$. Thus no two strict relations are composable, and the list already defines a transitive partial order. Define $p(a_i)=0$, $p(b_i)=p(d_i)=1$, and $p(c_i)=2$. All listed relations preserve this order, so $p$ is a functor.

The fibre $P_0$ is the discrete pair $\{a_0,a_1\}$, the fibre $P_1$ is the disjoint union of the two intervals $b_i<d_i$, and $P_2$ is the discrete pair $\{c_0,c_1\}$. For the nonidentity base arrow $0\to1$, the pullback is the disjoint union of the two posets $a_i<d_i>b_i$. Each component has terminal object $d_i$, and each endpoint fibre meets it in a nonempty contractible subcategory. Both endpoint inclusions are therefore weak homotopy equivalences. For $1\to2$, the two components are $d_i>b_i<c_i$; each has initial object $b_i$, and the same conclusion follows. For $0\to2$, the pullback consists of the two intervals $a_i<c_{1-i}$, so again both endpoint inclusions are weak equivalences. A constant functor $[1]\to[2]$ has pullback $P_j\times[1]$; its endpoint inclusions are categorical homotopy equivalences. These six functors exhaust all $[1]\to[2]$, proving every required edge test.

Label the two connected components of each fibre by $i$. The span through the $0,1$ pullback identifies component $i$ with component $i$; so does the span through the $1,2$ pullback. The $0,2$ pullback instead identifies $i$ with $1-i$. Thus the maps induced on $\pi_0$ satisfy
\[
 \pi_0\lambda(0\to1)=1,\qquad
 \pi_0\lambda(1\to2)=1,\qquad
 \pi_0\lambda(0\to2)=(0\ 1).
\]
Since $\pi_0$ sends $W_\infty$ to bijections, it factors through $\Hot$, and these unequal functions prove that the proposed transports cannot form a functor into $\Hot$.

For completeness, the nerve of $P$ has eight vertices, eight nondegenerate edges, and no higher nondegenerate simplices. Those edges form exactly the cycle
\[
 a_0-d_0-b_0-c_0-a_1-d_1-b_1-c_1-a_0.
\]
Its realization is therefore a circle. Each fibre has two connected components, whereas $P$ is connected. Consequently none of the inclusions $P_j\to P$ is in $W_\infty$, although the base $[2]$ has both an initial and a terminal object. This disproves the stronger fibration condition over contractible bases and completes the proof.

The counterexample also satisfies the edge tests for $W_0$, the localizer detected by connected components, and has the same failure of composition there. The issue is therefore not a higher-dimensional obstruction inside the fibres. The missing information is how transports for composable arrows agree across a triangle in the base. The universal comma hypotheses used in Proposition 7 supply precisely the strict diagram from which that agreement was obtained. We checked the lift-category and smoothness treatments in \citet[\S4]{Maltsiniotis2005Asphericite} and \citet[\S\S1--2]{Maltsiniotis2012Exact}; we did not find this finite example there. This is a development of the manuscript's question, with no claim of priority.
\end{dcomment}


\leavevmode\msp{page 37}La notion la plus faible raisonnable de $W$-fibré, me semble être celle où on suppose les deux choses à la fois. Les hypothèses qu'on a faites jusqu'à présent, à l'exclusion de la forme forte de $W$-propreté ou de $W$-lissité, reviennent toujours à supposer que le \textit{changement de base par $f : X\longrightarrow S$} transforme certaines flèches dans $\Cat/S$ en \texttt{[élément de]} $W$ : Pour définir $\lambda$, ces flèches sont
\[
\text{soit } \{s\}\longrightarrow S/s,\qquad \text{soit } \{1\}\longrightarrow\mathbf{\Delta}^1(\longrightarrow S),
\]
pour définir $\lambda^\circ$, ce sont
\[
\text{soit } \{s\}\longrightarrow s\backslash S,\qquad \text{soit } \{0\}\longrightarrow\mathbf{\Delta}^1(\longrightarrow S),
\]
le `soit' suivant la définition qu'on a adoptée. Dans ces exemples, il s'agit toujours de l'inclusion de l'objet final d'une catégorie $S'$, dans ladite, ou inversement de l'objet initial. Cela montre d'ailleurs (si on exige l'universalité), que la première hypothèse (assurant $\lambda$) est \textit{exactement} cela (cas de l'inclusion d'un objet final $s'$ - puisque pour un tel objet, on a $S'/s'\simeq S'$), tandis que la deuxième hypothèse exige moins, seulement pour le cas de $\mathbf{\Delta}^1$. Et dualement pour $\lambda'$ \texttt{[plutôt $\lambda^\circ$]}. Mais il convient de noter que pratiquement, la façon la plus commode, et de loin,

\leavevmode\msp{page 38}pour vérifier de telles propriétés, c'est soit par la $W$-propreté de $f$, soit par la $W$-lissité - ce qui ramène à vérifier que certaines catégories bien précisées sont $W$-asphériques, sans avoir à se soucier de changements de base et de vérifier que les flèches construites en termes de changements de base arbitraires sont dans $W$. Et l'emploi de la $W$-propreté et de la $W$-lissité est d'autant plus sympathique, que ces notions sont stables par composition,

% ---- from Der-ChXIIa-012_p031-035.tex
\printedpage{29}
% head: A. Grothendieck : Dérivateurs, Ch. XII, Édition des Universités de Montpellier II et Paris VII
alors que toutes les notions dérivées, ainsi que les notions de fibrations à la clef, ont une fâcheuse tendance à ne \textit{pas} être stable par composition.

On pourrait appeler pseudo-$W$-propre resp. quasi-$W$-propre (un peu mieux que pseudo) ce qu'on vient de dégager - j'entends des ricanements! Et si on veut les deux à la fois - $\lambda$ et $\lambda^{\circ}$ associés comme on a dit? Je vois le même principe de description avec des $S''\fl S'$ sur $S$, d'exiger que $X''\fl X'$ soit transformé en des \texttt{[éléments de]} $W$. Avec les choix suivants qui ont l'air sympathique

\begin{itemize}
\item[a$_1$)] \msp{page 39}$Y'$ a un objet final, $Y''=e$. La définition revient donc à la condition que les
\[ X_s\fl X/s, \qquad X/s'\fl X/s \]
(pour $s\in\Ob S$ où $u : s'\fl s$ dans $\mathrm{S}$) soient dans $W$, et que ça reste comme ça après tout changement de base $S'\fl S$.

\item[a$_2$)] $Y'$ a un objet initial, $Y''=e$, c'est la condition duale à coups de
\[ X_s\fl s\backslash X, \qquad s\backslash X\fl s'\backslash X. \]

\item[a$_0$)] $Y'=\boldsymbol{\Delta}^1$, $Y''=e$ (\footnote[46]{NB a$_1$), a$_2$) sont notés pareil p. 4, mais a$_0$) a été oublié.}).

\item[b)] $Y'$ est contractile, $Y''$ aussi (ou au choix, $Y'$ contractile, $Y''=e$)
\end{itemize}

donc la condition revient à ceci : pour toute $Y'$ sur $Y$ \textit{contractile} (notion indépendante de $W$, c'est ça son principal avantage) les inclusions des fibres de $X'=X\times_Y Y'$ sur $Y'$ dans $X'$ sont \texttt{[dans]} $W$.

\dgrblock{dgrb490}{The test objects explain the hierarchy}
\begin{dcomment}{exposition}{p39:donc la condition revient à ceci}{The test objects explain the hierarchy}
The proposed definitions differ by the category that is allowed to serve as a test base. For condition $a_1$, a base $T$ with terminal object $t$ is tested at \emph{every} object $r\in T$, not only at $t$. Testing the terminal object gives preproperness. Once it is available, the unique arrow $r\to t$ identifies the comparison $X_r\to X_T$ with transport to the terminal fibre. Requiring that comparison to belong to $W$ adds the demand that all such transports be weak equivalences. The dual reading gives $a_2$.

Condition $a_0$ restricts the base to $[1]$ and allows either endpoint. The preceding eight-object example shows that it cannot be substituted for $a_1$ or $a_2$, even for ordinary weak homotopy equivalences. In particular, existence of individually invertible edge roofs is not a substitute for a functor of fibres. This is the reason for keeping the preproper or presmooth hypothesis in the equivalences proved later in the section.

For condition $b$, the equivalence of the two tests stated in the text has a short proof. Suppose all fibre inclusions are in $W$ over every categorically contractible test base. If $v:T''\to T'$ is a map between two such bases, choose an object $t''$; contractibility ensures that an object exists. In the pullback square, the fibre at $t''$ identifies with the fibre at $v(t'')$. The two inclusions from that common fibre into the total categories both belong to $W$, so their commutative triangle puts the total pullback map in $W$ by two-out-of-three. The converse is the special case $T''=e$.

Condition $b'$ allows $W$-aspherical bases instead. Condition $c$ allows every map in $W$ as a change of base. Its stronger quantifier is substantial: it asks for preservation of a whole class of weak equivalences, whereas the earlier conditions test selected contractible shapes. The terminology at this point is provisional, but these quantifiers are exact and should be retained when comparing later definitions.
\end{dcomment}


Je sens que même pour des $W$ générales, b) reste du domaine du vérifiable, tandis qu'on mollit si on passe à l'extrême en généralité dans ce sens :

\begin{itemize}
\item[b$'$)] Comme b), mais en supposant seulement $Y'$ $W$-asphérique (\footnote[47]{C'est essentiellement le a$'$) de p. 4}).
\end{itemize}

Enfin (la condition a) de la page 4)

\begin{itemize}
\item[c)] Si $Y''\fl Y'$ est \texttt{[dans]} $W$ sans plus, $X''\fl X'$ aussi.
\end{itemize}
% foot: Version 7 décembre 2002

\printedpage{30}
% head: A. Grothendieck : Dérivateurs, Ch. XII, Édition des Universités de Montpellier II et Paris VII
\msp{page 40}La seule de ces notions qui ait la moindre envie d'être stable par composition, c'est la plus forte c). Par contre, la notion de flèche $W$-parfaite (cf. déf. 3) est stable spontanément par composition :

\textbf{Définition 3.} On dit que $f : X\fl S$ est $W$-\textit{parfaite} si elle est à la fois $W$-propre et $W$-lisse.

\dgrblock{dgrb491}{The interval test and the full fibration condition}
\begin{dcomment}{exposition}{p40:à la fois $W$-propre et $W$-lisse.}{The interval test and the full fibration condition}
Grothendieck is trying to retain the mobility of fibres without asking that every weak equivalence survive every pullback. The latter requirement is exceptionally strong when $W$ forgets some homotopy groups. The useful intermediate condition asks instead that a categorically contractible base make the family constant up to $W$. Here ``contractible'' means that its identity functor is joined to a constant functor by a finite zigzag of natural transformations. This is the notion of homotopy already used in section~I; it must not silently be replaced by weak contractibility of the nerve.

The weakest test in the diagram examines only $[1]\to S$. It gives, for each arrow $u:s\to t$, a span
\[
 X_s\longrightarrow X\times_S[1]\longleftarrow X_t
\]
whose arrows belong to $W$. Inverting them produces an isomorphism between the two fibres in $\Hot_W$. What it does not supply is the compatibility of this isomorphism with composition of arrows of $S$. The development ``An eight-object obstruction to edge transport'' above gives a finite poset over $[2]$ for which all these endpoint tests are weak equivalences, whereas the transports along $0\to1\to2$ and $0\to2$ disagree already on connected components. Its total category has the homotopy type of a circle and each fibre has two components. It therefore answers negatively the question about $a_0$ raised in the next paragraph.

Proposition~8 adds precisely the missing global control. Universal preproperness identifies a fibre with the appropriate comma category; the comma categories carry genuinely functorial transport, because postcomposition in $S$ is an actual functor. This explains why preproperness together with the interval test succeeds where the interval test alone fails. Perfectness is a convenient sufficient condition, since it supplies both universal comma comparisons. It is stronger than the fibration condition being isolated here.
\end{dcomment}


Dans le cas $W=W_\infty$, toutes les variantes de la page précédente, à l'exclusion seulement de a$_0$), sont équivalentes. Ce serait intéressant de dégager des conditions minimales sur $W$ pour assurer qu'il en soit ainsi. (Et il se pose la question si dans le cas de $W_\infty$, a$_0$) suffit pour impliquer le reste - je ne m'y attends guère \dots)

Dans le cas général, on a le diagramme d'implications (et encore je me suis limité)
\[
\begin{tikzcd}[column sep=small]
 & & \mathrm{c} \arrow[d, Rightarrow] & & \text{$W$-fibration ultraforte} \\
 & & \mathrm{b}' \arrow[dd, Rightarrow] & & \text{$W$-fibration $W$-forte} \\
 & \text{$W$-parfait} \arrow[dr, Rightarrow] & & & \\
 & & \mathrm{b} \arrow[dl, Rightarrow] \arrow[dr, Rightarrow] & & \text{$W$-fibration hot-forte (\footnotemark[48])} \\
\text{$W$-fibration gauche} & \mathrm{a}_1 \arrow[dr, Rightarrow] & & \mathrm{a}_2 \arrow[dl, Rightarrow] & \text{$W$-fibration droite} \\
 & & \mathrm{a}_0 & & \text{$W$-pseudofibration}
\end{tikzcd}
\]
\footnotetext[48]{Il semblerait que ce soit ça, finalement, la ‘bonne notion’ de $W$-fibration, \textit{équivalente} à a$_1$ et a$_2$, lesquelles seraient équivalentes entre elles. Cf. prop. 9, p. 48.}

\texttt{[‘$W$-parfait $\Longrightarrow$ b’ encerclé]}.

\msp{page 41}\textbf{Proposition 8.} \textit{Soit $f : X\fl S$ dans} Cat. \textit{Les conditions suivantes sont équivalentes.}

\begin{itemize}
\item[\textit{a)}] \textit{$f$ est une $W$-fibration gauche, i.e. la condition} a$_1$ \textit{de la page 39 est vérifiée.}

\item[\textit{b)}] \textit{$f$ est $W$-prépropre, et $\forall\; s_0\fl s_1$ dans $S$, $X/s_0\fl X/s_1$ est dans $W$.}

\item[\textit{c)}] \textit{$f$ est $W$-prépropre, et une $W$-pseudofibration, i.e. la condition a$_0$ de la page 39 est vérifiée.} % AUDIT resolved: CMTI12 text-italic a with CMR8 subscript; see corrections.csv, PDF leaf 40, ms41
\end{itemize}

(\footnote[49]{\textbf{Corollaire.} \textit{Quand $f$ est} Cat-\textit{cofibrant, alors ces conditions équivalent aussi à ceci : $\forall\; u : s_0\fl s_1$ dans $S$, $u_* : X_{s_0}\fl X_{s_1}$ est dans $W$.}}).
% foot: Version 7 décembre 2002

\printedpage{31}
% head: A. Grothendieck : Dérivateurs, Ch. XII, Édition des Universités de Montpellier II et Paris VII
\textbf{NB} Ces conditions sont satisfaites si $f$ est en même temps $W$-prélisse, cf. prop. 7.

\dgrblock{dgrb492}{Completing the terminal-object test}
\begin{dcomment}{exposition}{p41:Ces conditions sont satisfaites si}{Completing the terminal-object test}
Grothendieck wants a condition that can be checked on the original base, while remaining valid after every change of base. We spell out the argument because the compressed reference to $Y/s$ and $Y/s'$ hides the universal quantifier.

Write $E_u=X\times_S[1]$ for $u:s_0\to s_1$. If the terminal-object test holds, apply it first to $T/t\to T\to S$, for an arbitrary $T\to S$ and $t\in T$. Its terminal fibre is $X'_t$ and its total category is $X'/t$, where $X'=X\times_ST$. Thus $X'_t\to X'/t$ belongs to $W$: this is universal preproperness. Next apply the same test to $S/s_1$. Both the object $(s_1,1)$ and the object $(s_0,u)$ give fibre inclusions into $X/s_1$ in $W$. Since the latter factors as
\[
 X_{s_0}\longrightarrow X/s_0\longrightarrow X/s_1,
\]
two-out-of-three gives $X/s_0\to X/s_1\in W$. This supplies the missing detail in $a)\Rightarrow b)$.

Assume now universal preproperness. The canonical functor $E_u\to X/s_1$ restricts on the terminal fibre to $X_{s_1}\to X/s_1$. Both this inclusion and $X_{s_1}\to E_u$ belong to $W$, so $E_u\to X/s_1$ belongs to $W$. Its restriction to the initial fibre is the composite displayed above. Consequently
\[
 X_{s_0}\to E_u\in W
 \quad\Longleftrightarrow\quad
 X/s_0\to X/s_1\in W.
\]
The other endpoint inclusion is already in $W$. This proves $b)\Longleftrightarrow c)$, including the asserted stability under arbitrary base change, since preproperness and the interval condition both have that stability by definition.

Finally, after any base change to a category $T$ with terminal object $t$, the inclusion of the fibre over $s\in T$ factors through $X'/s\to X'/t=X'$. Both factors belong to $W$, proving $a)$. The dual argument uses initial objects and the categories $s\backslash X$. For an ordinary opfibration the inclusions $X_s\to X/s$ have adjoints and are homotopy equivalences; its postcomposition maps identify, up to these comparisons, with the pushforward functors $u_*$. This proves the corollary in footnote~49 with its stated variance.
\end{dcomment}


\textsc{Démonstration}. On a a) $\Longrightarrow$ b) trivialement (faire $Y''=Y/s$, $Y''=Y/s'$ \texttt{[lien avec l'assertion? Cf. p. 4.]}).

b) $\Longleftrightarrow$ c) par le raisonnement de la proposition 7, car en termes de $X'=X_u=X\times_S(\boldsymbol{\Delta}^1,u)$, la condition b) pour $u : s_0\fl s_1$ donné s'interprète dans le diagramme commutatif
\[
\begin{tikzcd}[column sep=small, row sep=small]
X'_0 \arrow[dr, "i'_0\in W"] \arrow[dd, "\wr"] & & X'_1 \arrow[dr, "i'_1\in W"] \arrow[dd, "\wr"] & \\
 & \underset{(\simeq X'_0)}{X'/0} \arrow[rr, "\varphi'", crossing over] \arrow[dd, "p_0"] & & X'/1\ (\simeq X') \arrow[dd, "p_1"] \\
X_{s_0} \arrow[dr, "i_0\in W"'] & & X_{s_1} \arrow[dr, "i_1\in W"] & \\
 & X/s_0 \arrow[rr, "\varphi"] & & X/s_1
\end{tikzcd}
\]

\texttt{[il y a quelques $W$ redondants dans le diagramme original]} où $i_0$, $i_1$, $i'_0$, $i'_1$ sont dans $W$ à cause de l'hypothèse de $W$-prépropreté (stable par changement de base).

Donc $p_0$, $p_1$ sont \texttt{[dans]} $W$, donc $\varphi\in W$ si et seulement si $\varphi'\in W$, or $\varphi'\in W$ est la condition a$_0$.

\msp{page 42}Prouvons enfin que b), c) impliquent a). Comme c) (donc aussi b)) est stable par changement de base, il suffit de montrer ceci : Si $S$ a un objet final $e$, alors pour tout $s\in\Ob S$, l'inclusion $X_s\xrightarrow{\ i\ }X$ est dans $W$. Or on a le diagramme commutatif
\[
\begin{tikzcd}
X_s \arrow[r, "i_s\in W"] \arrow[dr, "i"'] & X/s \arrow[d, "u\,\in\,W\text{ par c) }\texttt{[plutôt par b)]}"] \\
 & X/e=X,
\end{tikzcd}
\]
qui montre bien $i=ui_s\in W$.

Dualement, si $f$ est $W$-prélisse, alors la condition a$_0$ ($f$ une $W$-pseudofibration) est équivalente à a$_2$ ($f$ une $W$-fibration droite), et signifie aussi que les
\[ s_1\backslash X\fl s_0\backslash X \]
sont dans $W$. Mais la notion de ‘$W$-fibration’ sur laquelle on tombe (si on ne se contente du minable a$_0$!) est différente dans les deux cas duaux l'un de l'autre. On aimerait dans les deux cas trouver des $W$-fibrations ‘hot-fortes’ (condition b), tout au moins en mettant le paquet (s'il le faut) - i.e. pour $f$ $W$-parfait. Ça revient donc à ceci.

\textbf{Question :} \textit{Soient $X\xrightarrow{\ f\ }S$, avec $f$ $W$-parfait, $S$ contractile, $s\in\Ob S$. L'inclusion $X_s\fl X$, est-elle dans $W$?}
% foot: Version 7 décembre 2002

\printedpage{32}
% head: A. Grothendieck : Dérivateurs, Ch. XII, Édition des Universités de Montpellier II et Paris VII
\msp{page 43}On va supposer seulement que $f$ est $W$-lisse et une $W$-fibration gauche \texttt{[plutôt que $f$ est simplement une $W$-fibration à droite, i.e que $f$ est $W$-prélisse et une $W$-pseudofibration]} (cf. prop. 8). On va se ramener au cas où $X$ est même fibrant sur $S$, en introduisant $\Phi_S(X)$ \texttt{[catégorie fibrée définie par le foncteur $S^{\circ}\fl\Cat$, $s\longmapsto s\backslash X$]} et % [sic] i.e que -> i.e. que, missing period of the abbreviation
\[ X\xrightarrow{\ i_X\ }\Phi_S(X). \]

\dgrblock{dgrb493}{The split replacement and its actual comparison maps}
\begin{dcomment}{exposition}{p43:On va se ramener au cas}{The split replacement and its actual comparison maps}
Grothendieck now replaces a family by a split fibration so that transport can be written before localization. The replacement is legitimate for a more concrete reason than an appeal to a previously constructed derivator.

An object of $\Phi_S(X)$ is a triple $(s,x,a)$ with $a:s\to p(x)$. A morphism to $(t,y,b)$ is a pair $(u:s\to t,v:x\to y)$ satisfying $p(v)a=bu$. The structural projection sends this object to $s$. Its fibre at $s$ is $s\backslash X$, and an arrow $u:s\to t$ acts by precomposition. This is a strict functor $S^{\mathrm{op}}\to\Cat$, hence a split Grothendieck fibration.

The comparison $i_X$ sends $x$ to $(p(x),x,1)$. The forgetful functor $r:\Phi_S(X)\to X$ satisfies $ri_X=1_X$, and the arrows $(a,1_x)$ define a natural transformation $1_{\Phi_S(X)}\Rightarrow i_Xr$. In fact $r$ is left adjoint to $i_X$. Thus $i_X$ is a categorical homotopy equivalence and belongs to every $W$ under consideration. On the fibre at $s$ it is exactly
\[
 X_s\longrightarrow s\backslash X.
\]
For a universally prelisse $X/S$, these fibre comparisons belong to $W$ after every base change. The comparison of total categories and the comparison of individual fibres are therefore two separate, verified assertions.

If $X/S$ is also a pseudofibration, the dual of Proposition~8 says that every precomposition $t\backslash X\to s\backslash X$ is in $W$. Consequently $\Phi_S(X)$ is a split fibration with transport in $W$. In a square comparing a fibre inclusion for $X$ with the corresponding inclusion for $\Phi_S(X)$, both horizontal arrows are in $W$; two-out-of-three transfers the desired conclusion. No saturation argument is needed for this step. The familiar free-fibration construction and its variance were discussed in the commentary on Chapter~VI, manuscript pp.~53--60; the additional point here is the universal fibre comparison that permits its use over an arbitrary new base.
\end{dcomment}


L'ennui, c'est que j'ai fait la théorie de $\HOT_W$ jusqu'à maintenant avec l'hypothèse (qui s'avère idiote) que $W\supseteq W_\infty$ - je serai donc obligé de revoir rapidement sous quelles hypothèses plus raisonnables sur $W$ les résultats familiers sont valables. Par exemple j'aimerais que $i_X\in W$ (\footnote[50]{en fait, $i_X\in(W_\omega)^{\mathrm{d}}_S\subseteq W^{\mathrm{d}}_S\subseteq W_S$.}), de sorte qu'on ait
\[
\begin{tikzcd}[column sep=large, row sep=large]
X_s \arrow[r, "i_s\in W", "\text{car $X$ $W$-lisse}"'] \arrow[d, "j^X_s"'] & \Phi_S(X)_s=s\backslash X \arrow[d, "j_s^{\Phi_S(X)}"] \\
X \arrow[r, "i_X\in W"] & \Phi_S(X),
\end{tikzcd}
\]
on voit donc que pour prouver que l'inclusion $j^X_s$ est \texttt{[dans]} $W$, il suffit de prouver la même chose pour $\Phi_S(X)$, qui est fibrant (pas seulement $W$-lisse). De plus, l'hypothèse pseudo-fibration sur $X$ passe sûrement à $\Phi_S(X)$ (je réserve les deux vérifications à plus tard (\footnote[51]{O.K.})). On est donc ramené au cas de
\[ X\in\Ob\uFib_0S, \qquad S\text{ contractile,} \]
à prouver que l'inclusion
\[ X_s\fl X \]
\msp{page 44}est dans $W$. Je vais prouver ceci :

\textbf{Lemme.} \textit{Soient $S$ dans} Cat, \textit{$X\in\Ob\uFib_0S$, et considérons, pour un $S'$ dans} Cat,
\[
\begin{aligned}
f &\quad\longmapsto\quad f^*(X/S)=X\times_S(S',f)\\
\uHom(S',S)^{\circ} &\quad\longrightarrow\quad \uFib_0S'.
\end{aligned}
\]
\textit{(Ce foncteur est défini grâce au fait que $X$ est fibrant sur $S$, et peut s'interpréter donc, à $S$-équivalence près, en termes d'un foncteur $S^{\circ}\fl$} Cat \textit{\dots\ (\footnote[52]{\textbf{NB} $\Phi_S(X)$ était même déjà scindée, définie par le foncteur $s\longmapsto s\backslash X$.})). Les flèches de $\uHom(S',S)$ sont transformées en flèches $\in W^{\mathrm{d}}_{S'}$.}
% foot: Version 7 décembre 2002

\printedpage{33}
% head: A. Grothendieck : Dérivateurs, Ch. XII, Édition des Universités de Montpellier II et Paris VII
En effet, en termes de foncteurs
\[ S'^{\circ}\fl\Cat, \]
si $X$ est définie par
\[ S^{\circ}\xrightarrow{\ \mathscr{X}\ }\Cat, \]
alors $f^*(X)$ est définie par $\mathscr{X}\circ f$, et la flèche
\[ \mathscr{X}\circ f\fl\mathscr{X}\circ g \]
définie par $u : f\fl g$ n'est autre que $\mathscr{X}*u$. Sa valeur en $s'\in\Ob S'$ est obtenue en prenant
\[ u(s') : f(s')\fl g(s'), \]
d'où par $\mathscr{X}$
\[ \mathscr{X}(u(s')) : \underbrace{\mathscr{X}(g(s'))}_{(\mathscr{X}\circ g)(s')}\fl\underbrace{\mathscr{X}(f(s'))}_{(\mathscr{X}\circ f)(s')} \]
(\footnote[53]{à dégager un énoncé à part.}). Mais l'hypothèse $X\in\uFib_0S$ signifie que les $\mathscr{X}(v)$, $v\in\mathrm{Fl}(S)$, sont tous dans $W$. Donc
\msp{page 45}$\mathscr{X}$ est dans $W$ fibre par fibre, donc dans $W^{\mathrm{d}}_{S'}$.

\dgrblock{dgrb494}{Transport along a natural transformation}
\begin{dcomment}{exposition}{p45:est dans $W$ fibre par fibre}{Transport along a natural transformation}
The purpose of the lemma is to turn a homotopy of indexing functors into an actual zigzag of families. For this purpose it is useful to work with the split fibration just constructed. Let $F:S^{\mathrm{op}}\to\Cat$ be its classifying functor, with every $F(u)$ in $W$, and let $\alpha:f\Rightarrow g$ be a natural transformation of functors $T\to S$.

Contravariance gives the strict natural transformation
\[
 F(g(-))\longrightarrow F(f(-)),\qquad
 t\longmapsto F(\alpha_t).
\]
Its Grothendieck construction is a Cartesian functor $g^*X\to f^*X$ over $T$. Each component belongs to $W$. The relative localization axiom, applied in the dual comma convention, implies that this total functor belongs to $W$ as well. Explicitly, the inclusion of the fibre of a Grothendieck fibration into its under-comma category has an adjoint, so the componentwise assertion gives an equivalence on those comma categories; the dual of W(2) then gives the total equivalence.

The first displayed arrow in the source proof points in the opposite direction from the correctly written component immediately below it. It is an edition slip: the transformation is $F\circ g^{\mathrm{op}}\Rightarrow F\circ f^{\mathrm{op}}$. Composition is nevertheless completely controlled: for $\beta:g\Rightarrow h$, the identity
\[
 F(\beta_t\alpha_t)=F(\alpha_t)F(\beta_t)
\]
identifies the transport for $\beta\alpha$ with the composite of the two transports, in the contravariant order. Units are identities. After localization every such transport is invertible, and one may use its inverse when writing the covariant version of the functor in the next paragraph.

The factorization through the groupoid completion of $\underline{\mathrm{Hom}}(T,S)$ follows from its universal property. Connected functors therefore induce isomorphic pullbacks, but a connected component does not choose a particular path, and different paths may induce different isomorphisms. This is the precise reason for the parenthetical qualification ``non canoniquement''.
\end{dcomment}


\textbf{Corollaire 1.} \textit{Le foncteur du lemme, composé avec}
\[ \uFib_0S'\fl\HOT^{\mathrm{lc}}_W(S')\overset{\text{déf}}{=}(\uFib_0S')(W^{\mathrm{d}}_{S'})^{-1} \]
\textit{se factorise par le groupoïde fondamental de} $\uHom$ \textit{:}
\[
\begin{tikzcd}[row sep=huge]
\uHom(S',S)^{\circ} \arrow[r] \arrow[d] & \uFib_0S' \arrow[d] \\
\textstyle\prod(\uHom(S',S))^{\circ} \arrow[r] & \HOT^{\mathrm{lc}}_W(S').
\end{tikzcd}
\]
\textit{Si $f$, $g$ sont dans une même composante connexe de $\uHom(S',S)$, alors $f^*X$ et $g^*X$ sont isomorphes (non canoniquement, en général) dans} $\HOT^{\mathrm{lc}}_W(S')$.

\dgrblock{dgrb495}{A loop can act on the fibre even when the functor is fixed}
\begin{dcomment}{examples}{p45:isomorphes (non canoniquement, en général) dans}{A loop can act on the fibre even when the functor is fixed}
Grothendieck's qualification about noncanonical isomorphisms concerns genuine monodromy. It can be seen without higher homotopy groups. Let $S=BC_2$, the one-object category with automorphism group $C_2=\{1,\tau\}$, and let $F:S^{\mathrm{op}}\to\Cat$ take its value in the discrete category $D=\{0,1\}$, with $F(\tau)$ exchanging the two objects. All transports are isomorphisms.

There is only one functor $e\to BC_2$, but it has two natural automorphisms. Pullback along this functor is the fibre $D$. The identity natural transformation induces $1_D$, whereas the transformation $\tau$ induces the transposition of $D$. These are different even after localization at $W_\infty$, because they induce different maps on $\pi_0(D)$. Thus ``the two indexing functors are homotopic'' cannot single out an isomorphism between their pullbacks.

The total category $\int F$ is equally explicit. Its objects are $0,1$, and a morphism $d\to d'$ is an element $g\in C_2$ with $d=F(g)(d')$. There is exactly one such element for each ordered pair. The total category is therefore equivalent to $e$, whereas the fibre has two components. Its map to $BC_2$ is the categorical model of the universal two-sheeted covering. The inclusion of the fibre is not a weak equivalence; the base is not contractible. This example also shows why the word ``contractile'' in the next argument is a mathematical hypothesis, not a convenient choice of coordinates.
\end{dcomment}


Il est plus astucieux de ne pas fixer $X$ et de considérer
\[ \uHom(S',S)\fl\uHom(\HOT^{\mathrm{lc}}_W(S),\HOT^{\mathrm{lc}}_W(S')) \tag*{$(*)$} \]
correspondant à
\[
\begin{aligned}
\uHom(S',S) &\quad\times\quad \HOT^{\mathrm{lc}}_W(S) &&\quad\longrightarrow\quad &&\HOT^{\mathrm{lc}}_W(S')\\
(f &\quad,\quad\ \xi) &&\quad\longmapsto\quad && f^*(\xi).
\end{aligned}
\]
Le lemme dira que les flèches de $S'\fl S$ sont transformées par $(*)$ en isomorphismes, d'où
\[ \prod\uHom(S',S)\fl\uHom(\HOT^{\mathrm{lc}}_W(S),\HOT^{\mathrm{lc}}_W(S')), \]
% foot: Version 7 décembre 2002

% ---- from Der-ChXIIa-012_p036-040.tex
\printedpage{34}
% head: A. Grothendieck : Dérivateurs, Ch. XII, Édition des Universités de Montpellier II et Paris VII
% foot: Version 7 décembre 2002

\msp{page 46}
donc si $f,g : S' \rightrightarrows S$ sont deux foncteurs homotopes, alors
\[ f^* \simeq g^* : \HOTlc(S) \rightrightarrows \HOTlc(S'). \]

Il s'ensuit aussitôt ceci

\textbf{Corollaire 2.} \textit{Si $f$ est un homotopisme, $f^* : \HOTlc(S) \longrightarrow \HOTlc(S')$ est une équivalence de catégories.}

En effet, prenant $g : S \longrightarrow S'$ telle que
\[ fg \sim \id_S, \qquad gf \sim \id_{S'}, \]
on trouve que
\[ \begin{aligned}
g^*f^* \simeq (\id_S)^* &= \id_{\HOTlc(S)} \\
f^*g^* \simeq (\id_{S'})^* &= \id_{\HOTlc(S')},
\end{aligned} \]
donc on trouve que (quitte à ajuster ces isomorphismes fonctoriels) $f^*$ et $g^*$ sont des équivalences quasi-inverses l'une de l'autre.

J'ai l'impression (\footnote[54]{O.K.}) que tout ça ne doit pas utiliser plus que les axiomes standard sur $W$ comme localiseur fondamental, tout au plus W(2) au lieu de la forme faible W(2 bis).

Revenons à $X$ lisse sur $S$ contractile. Le corollaire 2 implique que
\[ \begin{aligned}
\Hot_W &\longrightarrow \HOTlc(S) \\
\xi &\longmapsto p^*_S(\xi),
\end{aligned} \]
où $p_S : S \longrightarrow e$, est une équivalence. Donc dans la catégorie $\HOTlc(S)$, $X$ est isomorphe \msp{page 47} à un Cat-fibré de la forme $Z \times S \xrightarrow{\mathrm{pr}_2} S$, avec $Z$ dans Cat. Or pour $X_1 = Z \times S \longrightarrow S$, il est immédiat que $X_{1\,s} \longrightarrow X_1$ est dans $W$, car un produit cartésien de flèches \texttt{[dans]} $W$ (savoir $\id_Z$ et $s \longrightarrow S$) est dans $W$. (En fait, l'inclusion est même un homotopisme!) Notons maintenant que si $X_0$, $X_1$ sont deux objets dans $\underline{\mathrm{Liss}}_0S$, tels qu'il existe une $S$-flèche $X_0 \longrightarrow X_1$ qui soit dans $\Wd$, alors $X_{0\,s} \longrightarrow X_0$ est dans $W$ si et seulement si $X_{1\,s} \longrightarrow X_1$ l'est. En effet, on a
\[ \begin{tikzcd}
X_{0\,s} \arrow[r, "\in W"] \arrow[d, "i_s^{X_0}"'] & X_{1\,s} \arrow[d, "i_s^{X_1}"] & \\
X_0 \arrow[r, "\in W", ""{name=A, below}] & X_1 & \\
 & & \text{\footnotesize\begin{tabular}{@{}l@{}}ici on utilise sans doute\\ W(2), non seulement W(2 bis),\end{tabular}}
\arrow[from=A, to=3-3, bend right]
\end{tikzcd} \]
\printedpage{35}
% head: A. Grothendieck : Dérivateurs, Ch. XII, Édition des Universités de Montpellier II et Paris VII
% foot: Version 7 décembre 2002
qui prouve le résultat.

\dgrblock{dgrb496}{The finite zigzag that makes contraction effective}
\begin{dcomment}{exposition}{p47:qui prouve le résultat.}{The finite zigzag that makes contraction effective}
Grothendieck correctly notices, immediately below, that an isomorphism in a localization does not by itself provide a zigzag consisting entirely of the specified weak equivalences. The contraction supplies such a zigzag directly.

Suppose $F:S^{\mathrm{op}}\to\Cat$ sends arrows into $W$, and choose a finite zigzag
\[
 1_S=\varphi_0\;\Longleftrightarrow\;\varphi_1
 \;\Longleftrightarrow\;\cdots
 \;\Longleftrightarrow\;\varphi_m=c_t
\]
of natural transformations, with the displayed double arrows indicating a choice of one direction at each step. Set $X_j=\int(F\circ\varphi_j^{\mathrm{op}})$ over $S$. Transport along the chosen transformation gives, in the appropriate direction, a functor $X_j\to X_{j+1}$ over $S$ which is in $W$ on every fibre. By the argument of the preceding block it is also in $W$ on total categories.

For every $s\in S$ the resulting square
\[
\begin{tikzcd}
 (X_j)_s\arrow[r]\arrow[d] & (X_{j+1})_s\arrow[d]\\
 X_j\arrow[r] & X_{j+1}
\end{tikzcd}
\]
has horizontal arrows in $W$; if the chosen natural transformation points backwards, reverse the horizontal arrows. Two-out-of-three says that one vertical arrow belongs to $W$ exactly when the other does. At the last stage the family is $F(t)\times S\to S$. The map $e\to S$ selecting $s$ is a categorical homotopy equivalence: a contraction to $t$, evaluated at $s$, supplies a zigzag between $s$ and $t$, which joins the constant functors $c_s$ and $c_t$. Multiplying this homotopy by $F(t)$ shows that
\[
 F(t)\cong F(t)\times\{s\}\longrightarrow F(t)\times S
\]
is a categorical homotopy equivalence, hence is in $W$. Moving back through the finite zigzag proves $X_s\to X\in W$.

The replacement comparison at p.~43 transfers this conclusion to every prelisse pseudofibration; the dual argument treats preproper pseudofibrations. This proves the question of p.~42 and the contractible-base part of Proposition~9 under W(1,2,3). It uses the explicit contraction, componentwise transport, the relative localization axiom, and two-out-of-three. Strong saturation is not needed for this particular conclusion.
\end{dcomment}


Il est vrai que le fait que $X = X_0$ et $X_1 = Z \times S$ soient isomorphes dans $\HOTlc(S)$ n'implique pas, par lui-même, l'existence d'une \msp{page 48} chaîne d'isomorphismes provenant de flèches dans $\Wd$. Mais en fait, si dans l'énoncé du corollaire 2 ci-dessus on compare $X$ dans $\underline{\mathrm{Fib}}_0S$ à $g^*f^*(X)$, grâce à une homotopie $\id_S \sim fg$, on trouve bel et bien une suite de foncteurs $\varphi_i : S \longrightarrow S$, avec $\varphi_0 = \id_S$, $\varphi_N = fg$, et $\varphi_i$ et $\varphi_{i+1}$ reliés par une flèche dans $\underline{\mathrm{Hom}}(S,S)$, donc on a une suite d'objets $X_i = \varphi_i^*(X)$ de $\underline{\mathrm{Fib}}_0S$ (avec $X_0 = X$, $X_N = g^*f^*X$), et $X_i$ et $X_{i+1}$ reliés par une flèche \textit{dans} $\underline{\mathrm{Fib}}_0S$ (pas seulement dans $\HOTlc(S)$), laquelle flèche est dans $\Wd = \Wf$. Donc on gagne, sous réserve de vérifications de routine diverses (\footnote[55]{O.K.}). Donc modulo celles-ci, on aura

\textbf{Proposition 9} (\footnote[56]{C'est ok en supposant que $W$ est \textit{fortement} saturé, i.e. $\hot_W(f)$ isomorphisme $\Longrightarrow f \in W$.},\footnote[57]{Donc $W$-fibration à gauche $\Longleftrightarrow$ $W$-fibration à droite $\Longleftrightarrow$ $W$-fibration.}). \textit{Supposons que $X$ dans $\mathrm{Cat}/S$ soit une $W$-fibration à gauche (cf. prop. 8). Alors c'est même une $W$-fibration (au sens de b), page 40, où c'est appelé `$W$-fibration \textup{hot}-forte'), en d'autres termes, si $S'' \longrightarrow S'$ est sur $S$ avec $S''$, $S'$ contractiles, alors $X'' \longrightarrow X'$ est dans $W$ (\footnote[58]{on trouve seulement que $\hot_W(X'') \longrightarrow \hot_W(X')$ \texttt{[est un]} isomorphisme.}). On a sans doute mieux : si $S'' \longrightarrow S'$ est un homotopisme, alors $X'' \longrightarrow X'$ est dans $W$.}

\dgrblock{dgrb497}{What the localization argument proves}
\begin{dcomment}{exposition}{p48:On a sans doute mieux}{What the localization argument proves}
Grothendieck next asks for invariance under an arbitrary categorical homotopy equivalence of bases. Here the argument changes: he identifies the map of total categories with a map of homotopy colimits and applies an adjunction. Let $v:T\to S$ be a homotopy equivalence and let $\xi$ be a locally constant coefficient represented by the family. The contraction argument for the two composites of a homotopy inverse shows that restriction
\[
 v^*:\HOT_W^{\mathrm{lc}}(S)\longrightarrow\HOT_W^{\mathrm{lc}}(T)
\]
is an equivalence. Once the comparison with strict pullback and the adjunction $p_!\dashv p^*$ have been established, the displayed Hom calculation on p.~50 proves that the actual map of total categories becomes invertible in $\Hot_W$.

This is an assertion about the specified map, rather than merely an abstract isomorphism between its source and target. The last inference to membership in $W$ requires strong saturation. The footnote makes that distinction correctly. Later theory proves that every fundamental localizer in the modern, stronger sense is strongly saturated; see the status discussion at p.~16 and \citep[Proposition~4.2.4]{Cisinski2006}. Here the stronger sense includes the relative axiom W(2), not merely W(2 bis) in the terminology fixed at the beginning of section~III. Theorem~3 explicitly assumes W(1,2,3), so that later theorem removes its additional saturation hypothesis. We are not asserting this consequence for an arbitrary weak fundamental localizer satisfying only W(2 bis).

The finite-zigzag argument in the preceding block already proves the contractible-base tests directly in $W$. Thus it is consistent for Theorem~3 to state those equivalences under W(1,2,3), while presenting its more general assertion $d)$ first in the localized category.
\end{dcomment}


\msp{page 49}
Prouvons (si nous pouvons) ce complément sympathique. Je précise que j'ai eu à faire l'hypothèse de $W$-prélissité (où de $W$-prépropreté) pour être sûr que l'image inverse $X \times_S S'$, par $f' : S' \longrightarrow S$, coïncide bien avec $f^*$ sur les $\HOTlc$ - autrement je n'ai aucune chance. % [sic] où -> ou, the conjunction "or" is meant
`On sait' (tout est à revoir!) que c'est O.K. dans les trois cas $X$ $W$-lisse, $X$ $W$-propre, $X$ une $W$-fibration (mais en quel sens???). La vérification à faire (quand on travaille avec les $\Wd$) est celle-ci :
\[ \begin{aligned}
&X \xrightarrow{i_X} \Phi_S(X) = Y \text{ est dans } \Wd \text{ (sans mérite), et} \\
&\text{le reste par tout changement de base.}
\end{aligned} \tag*{$(*)$} \]

Mais pour des $X, Y \in \mathrm{Ob}\,\mathrm{Cat}/S$ qui sont $W$-prélisse, la $\Wd$-équivalence coïncide avec $\Wf$, laquelle est préservée par changement de base, ainsi que l'hypothèse sur $X$, $Y$. Donc on est tranquille.

On est ramené à prouver que si $X$ est $W$-prélisse sur $S$ et une $W$-pseudofibration, et si $S' \longrightarrow S$ est un homotopisme, alors $X' \longrightarrow X$ est dans $W$. Par l'argument de tantôt utilisant $X \longrightarrow \Phi_S(X)$, on se ramène au cas où $X$ est fibrée scindée.

\printedpage{36}
% head: A. Grothendieck : Dérivateurs, Ch. XII, Édition des Universités de Montpellier II et Paris VII
% foot: Version 7 décembre 2002

\msp{page 50}
\[ \begin{tikzcd}
X' \arrow[rr] \arrow[d] & & X \arrow[d] \\
S' \arrow[rr, "f"] \arrow[dr, "p'"'] & & S \arrow[dl, "p"] \\
 & e &
\end{tikzcd} \]

La flèche $X' = f^*(X) \longrightarrow X$ s'interprète dans $\Hot_W$ comme
\[ p'_!(f^*(\xi)) \longrightarrow p_!(\xi) \ , \qquad \text{où } \xi = \hot_{S,W}(X). \tag*{$(*)$} \]

Je dis que cette flèche est \texttt{[un]} isomorphisme, i.e. que pour tout $\eta \in \Hot_W$, la flèche transposée
\[ \Hom_W(p_!(\xi),\eta) \longrightarrow \Hom_W(p'_!(f^*(\xi)),\eta) \tag*{$(**)$} \]
est \texttt{[une]} bijection. Admettant la formule d'adjonction, la flèche $(**)$ s'interprète comme
\[ \Hom_W(\xi,\underbrace{p^*(\eta)}_{\eta_S}) \longrightarrow \Hom_W(f^*(\xi),\underbrace{p'^*(\eta)}_{f^*(\eta_S)}), \]
i.e.
\[ \Hom_W(\xi,\eta_S) \longrightarrow \Hom_W(f^*(\xi),f^*(\eta_S)), \]
que ce soit \texttt{[un]} isomorphisme résulte donc de $f^*$ pleinement fidèle sur les $\HOTlc$. Or on a vu que $f^*$ est une équivalence (cor. 2, page 46).

\msp{page 51}
Je vais reprendre les propositions 9 et 8, un peu mal foutues.

\textbf{Théorème 3.} \textit{On suppose que $W$ satisfait à \textup{W(1,2,3)}. Soit $f : X \longrightarrow S$. Considérons les conditions suivantes.}

\begin{itemize}[leftmargin=4em]
\item[\textit{a)}] \textit{$\forall\ S'' \longrightarrow S'$ sur $S$, avec $S''$, $S'$} contractiles, \textit{$X'' \longrightarrow X'$ est dans $W$.}

\item[\textit{a}$'$\textit{)}] \textit{$\forall\ S'$ sur $S$, avec $S'$ contractile, et $s' \in S'$, l'inclusion} % AUDIT resolved: the label is a prime, not d prime; see corrections.csv, PDF leaf 46
\[ X_{s'} \hookrightarrow X' = X \times_S S' \]
\textit{est dans $W$.}

\item[\textit{a}$_1$\textit{)}] \textit{Comme a$'$), mais en supposant que $S'$ a un objet final (au lieu de $S'$ contractile).} % AUDIT resolved: the reference is a prime; see corrections.csv, PDF leaf 46

\item[\textit{a}$_2$\textit{)}] \textit{Comme a$'$), mais en supposant que $S'$ a un objet initial.} % AUDIT resolved: the reference is a prime; see corrections.csv, PDF leaf 46

\printedpage{37}
% head: A. Grothendieck : Dérivateurs, Ch. XII, Édition des Universités de Montpellier II et Paris VII
% foot: Version 7 décembre 2002

\item[\textit{b}$_1$\textit{)}] \textit{$f$ est $W$-}prépropre \textit{(i.e. a$'$) est satisfait quand $s'$ est un objet final de $S'$, sans autre condition sur $S'$ dans $\mathrm{Cat}/S$), et de plus pour toute flèche $s_0 \longrightarrow s_1$ dans $S$, la flèche induite $X/s_0 \longrightarrow X/s_1$ est dans $W$ (i.e. a$'$ \textup{\texttt{[plutôt a]}}) est satisfait pour $Y/s_0 \longrightarrow Y/s_1$ \textup{\texttt{[plutôt} $S/s_0 \longrightarrow S/s_1$\texttt{]}}).}

\item[\textit{b}$_2$\textit{)}] \textit{$f$ est $W$-}prélisse \textit{(i.e. a$'$) est satisfait quand $s'$ est un objet initial de $S'$, sans autre condition sur $S'$), et de plus les flèches $s_0\backslash X \longrightarrow s_1\backslash X$ \textup{\texttt{[plutôt} $s_1\backslash X \longrightarrow s_0\backslash X$\texttt{]}} (pour $s_0 \longrightarrow s_1$ dans $S$) sont dans $W$, (i.e. a$'$ \textup{\texttt{[plutôt a]}}) satisfait pour $s_1\backslash Y \longrightarrow s_0\backslash Y$ \textup{\texttt{[plutôt} $s_1\backslash S \longrightarrow s_0\backslash S$\texttt{]}}).}

\item[\textit{c}$_1$\textit{)}] \textit{$f$ est $W$-prépropre et une $W$-pseudofibration, i.e. la condition a$'$) est satisfaite dans les deux cas suivants: 1$^\circ$) $s'$ est un objet final de $S'$, 2$^\circ$) $S' = \mathbf{\Delta}^1$.}

\item[\textit{c}$_2$\textit{)}] \textit{$f$ est $W$-prélisse et une $W$-pseudofibration, i.e. a$'$) dans les deux cas: 1$^\circ$) $s'$ est objet initial de $S'$, 2$^\circ$) $S' = \mathbf{\Delta}^1$.}
\end{itemize}

\textit{Toutes ces 8 conditions sont équivalentes. Elles impliquent la suivante :}

\begin{itemize}[leftmargin=4em]
\item[\textit{d)}] \textit{Soit $S'' \xrightarrow{u} S'$ un homotopisme dans $\mathrm{Cat}/S$, et considérons $u_X : X'' \longrightarrow X'$. Alors $\hot_W(u_X)$ est \textup{\texttt{[un]}} isomorphisme. (Donc si $W$ est fortement saturé, $u_X \in W$. Donc sous cette condition sur $W$, d) équivaut aux conditions précédentes.)}
\end{itemize}

\msp{page 52}
\textbf{Définition 4.} On dit que $f : X \longrightarrow Y$ dans Cat est une \textit{$W$-fibration}, si elle satisfait les conditions équivalentes a) à c$_2$) de la proposition 9.

\dgrblock{dgrb498}{Reading the eight equivalent conditions}
\begin{dcomment}{exposition}{p52:si elle satisfait les conditions équivalentes}{Reading the eight equivalent conditions}
Grothendieck has now isolated a usable fibration notion. Its eight formulations organize two independent requirements: strict fibres must compute comma fibres universally, and the latter must vary by weak equivalences. The reference to ``proposition~9'' in Definition~4 means the immediately preceding Theorem~3; Proposition~9 did not contain eight labelled conditions.

Here is the logical structure, including the point at which the contractible-base argument is used. Condition $a)$ implies $a')$ by taking $S''=e$. Conversely, for a map $u:T\to U$ between contractible categories over $S$, choose $t\in T$. The triangle of inclusions of the same strict fibre into $X_T$ and $X_U$ commutes. Both inclusions are in $W$ by $a')$, so $X_T\to X_U\in W$. Contractible categories are nonempty; there is no missing empty-base case.

The implications $a')\Rightarrow a_1),a_2)$ follow because categories with terminal or initial objects admit categorical contractions. Proposition~8 gives $a_1)\Longleftrightarrow b_1)\Longleftrightarrow c_1)$ and its dual gives $a_2)\Longleftrightarrow b_2)\Longleftrightarrow c_2)$. Starting with either $c_1)$ or $c_2)$, use the corresponding split replacement and the finite transport zigzag to obtain the fibre inclusion over every contractible base. This proves the reverse implication to $a')$ and therefore all eight equivalences.

For an ordinary Grothendieck fibration, the comma comparisons are automatic and the remaining test is simply that every pullback functor $u^*$ lies in $W$; the opfibred statement uses $u_*$. For a family that is both preproper and prelisse, the two interval endpoint inclusions already lie in $W$, so the pseudofibration test is automatic. These observations prove the three corollaries. They do not identify a $W$-fibration with a strict Grothendieck fibration, nor with a perfect functor: the first supplies strict Cartesian lifts, and the second supplies stronger asphericity of lift categories.
\end{dcomment}


\textbf{Corollaire 1.} \textit{Si $f$ est $W$-propre ou $W$-lisse, pour que $f$ soit une $W$-fibration, il faut et il suffit que $f$ soit une $W$-pseudofibration, i.e. que pour $S' = \mathbf{\Delta}^1 \longrightarrow S$, les inclusions des deux fibres}
\[ X'_0 \xrightarrow{i_0} X' \xleftarrow{i_1} X'_1 \]
\textit{de $X' = X \times_S S'$ soient dans $W$.}

\textbf{Corollaire 2.} \textit{Supposons $f$ \textup{Cat}-fibrant. Alors $f$ est une $W$-fibration si et seulement si pout toute flèche $u : s \longrightarrow s'$ dans $S$, $u^* : X_{s'} \longrightarrow X_s$ est dans $W$. Énoncé dual si $f$ est \textup{Cat}-cofibrant.} % [sic] pout -> pour, misprint

\textbf{Corollaire 3.} \textit{Supposons $f$ $W$-prépropre et $W$-prélisse (par exemple $f$ $W$-parfait). Alors $f$ est une $W$-fibration.}

\textbf{Proposition 10.} \textit{Soit $f : X \longrightarrow Y$ une flèche dans $\underline{\mathrm{Fib}}_WS$ (\footnote[59]{$X$ et $Y$ tous deux prépropres ou tous deux prélisses suffit.}) (sous-catégorie pleine de $\mathrm{Cat}/S$ formée des $W$-fibrés sur $S$). Si $\forall\ s \in S$, $f_s : X_s \longrightarrow Y_s$ est dans $W$, alors $f \in W$.}

\dgrblock{dgrb499}{Fibrewise equivalence gives equivalence of total categories}
\begin{dcomment}{exposition}{p52:est dans $W$, alors $f \in W$.}{Fibrewise equivalence gives equivalence of total categories}
Proposition~10 is a comparison theorem in one direction. It says that equivalent fibres give equivalent totals when the structural maps make those fibres compute the relevant comma categories. It does not yet assert that an equivalence of totals can be detected fibre by fibre; that converse is the separate axiom W(4b).

This argument uses the full relative axiom W(2), as assumed in Theorem~3; the weaker standing definition of section~III alone is not the stated justification. Suppose first that $X$ and $Y$ are prelisse over $S$, and that $h:X\to Y$ is in $W$ on every fibre. For each $s$ the commutative square
\[
\begin{tikzcd}
 X_s\arrow[r,"h_s"]\arrow[d] &Y_s\arrow[d]\\
 s\backslash X\arrow[r] &s\backslash Y
\end{tikzcd}
\]
has its upper and vertical arrows in $W$. Hence its lower arrow is in $W$. The dual form of W(2) gives $h\in W$. If the structural maps are preproper, use $X/s$ and $Y/s$ instead. This proves exactly the strengthening in footnote~59, without any pseudofibration hypothesis.

The source's replacement proof packages the same argument differently. The maps $i_X,i_Y$ are homotopy equivalences on total categories and are fibrewise equivalences under prelissity. Their split replacements are smooth over $S$, so the corollary to Proposition~5 applies to the induced map. Referring to that corollary, rather than requiring the map between the replacements itself to be smooth, is the correct reading of the final sentence.
\end{dcomment}


Considérons
\printedpage{38}
% head: A. Grothendieck : Dérivateurs, Ch. XII, Édition des Universités de Montpellier II et Paris VII
% foot: Version 7 décembre 2002
\[ \begin{tikzcd}
X \arrow[rr, "f"] \arrow[dr, "i_X"'] & & Y \arrow[dr, "i_Y"] & \\
 & \Phi_S(X) \arrow[rr, "\Phi_S(f)"'] & & \Phi_S(Y),
\end{tikzcd} \]
où $i_X, i_Y \in W$, donc $\Phi_S(f) \in W \Rightarrow f \in W$. Or le fait que $X$ \texttt{[est]} $W$-fibrée sur $S$ implique que $i_X \in \Wf$ (il suffit même que $X$ soit prépropre \texttt{[plutôt prélisse]} sur $S$), itou pour $Y$, $i_Y$. Donc l'hypothèse $f \in \Wf$ implique $\Phi_S(f) \in \Wf$, donc $\Phi_S(f) \in W$ par \texttt{[la]} proposition 5 (page 29).

% ---- from Der-ChXIIa-012_p041-044.tex
\printedpage{39}
% head: A. Grothendieck : Dérivateurs, Ch. XII, Édition des Universités de Montpellier II et Paris VII

\msp{page 53}

\section*{IV\quad Sur l'axiome W(4) et ses variantes.}

(\footnote[60]{Cf. aussi partie VI (HOT$_W$), p. 186-196}). J'ai passé des jours à essayer de prouver que toutes ces variantes sont équivalentes. Je vais donner une liste de celles, parmi ces variantes, qui avec le recul me paraissent encore avoir une utilité. (Avec des notations différentes de celles des pages 13-14.)

\dgrblock{dgrb500}{The two directions of a fibration axiom}
\begin{dcomment}{exposition}{p53:J'ai passé des jours à essayer}{The two directions of a fibration axiom}
Grothendieck separates the two properties that ordinary homotopy theory makes look inseparable. The first concerns a change of base: a weak equivalence $T\to S$, pulled back along a suitable family $X\to S$, should remain a weak equivalence. The second concerns a map between families over one fixed base: if its total map is a weak equivalence, each fibre map should be one. These are W(4a) and W(4b), respectively. W(4c) is the special case of the second in which the second family is $S\to S$ itself.

The primes restrict the objects on which these properties are tested. In W(4a') and W(4c') the map is perfect for the minimal localizer; in W(4b') the structural families are split and the comparison preserves the splittings; in W(4b'') the structural families are $W$-perfect. These restrictions are not interchangeable merely by renaming a fibration. In particular, the argument of section~I used a stronger base-change condition than the eventual definition on p.~52. The editors' objection to the arrow $a'\Rightarrow b$ records a genuine dependence on that change of definition.

There is also a change of lettering. The weak perfect-fibre axiom called W(4b'') in section~I is now W(4c'). Keeping its content fixed is safer than carrying its old label into a later argument. The subsequent proof establishes one implication through locally constant coefficients; it does not prove every arrow that the surrounding discussion anticipates.
\end{dcomment}


\begin{itemize}[leftmargin=5em,labelsep=0.5em]
\item[\fbox{W(4a)}] Toute $W$-fibration (cf. th. 3, p. 51, déf. 3, p. 52) est ultraforte, ce qui revient à ceci: pour
\[
\begin{tikzcd}
X & X' \arrow[l,"q"'] \arrow[d,"f'"] \\
Y \arrow[from=u,"f"'] & Y' \arrow[l,"p"']
\end{tikzcd}
\]
cartésien dans Cat, si $f \in \underline{\Fib}_W$, $p \in W$, alors $q \in W$.

\item[W(4a$'$)] Comme a), en supposant $f \in \Parf_\omega$, i.e. $f$ $W_\omega$-parfait (donc une $W_\omega$-fibration), où $W_\omega \subseteq W$ est la plus petits partie $\Sigma$ de Fl(Cat) satisfaisant W(1, 2 bis, 3). % [sic] petits -> petite, accord avec « partie »

\item[\fbox{W(4b)}] Soit un diagramme commutatif
$\begin{tikzcd}[ampersand replacement=\&, column sep=small, row sep=small, baseline=(X.base)]
|[alias=X]| X \arrow[rr,"f"] \arrow[dr,"p"'] \& \& Y \arrow[dl,"q"] \\
\& S \&
\end{tikzcd}$
dans Cat, avec $p, q \in \Fib_W$, $f \in W$.
Alors $\forall s \in S$, $f_s : X_s \longrightarrow Y_s$ \texttt{[est dans} $W$\texttt{]}.

\item[W(4b$'$)] Comme b), en supposant \texttt{[de plus]} $X$ et $Y$ fibrés scindés sur $S$, $f$ compatible avec \texttt{[les]} scindages (donc foncteur cartésien).

\item[W(4b$''$)] Comme b), en supposant $p$, $q$ $W$\textit{-parfaits.}

\item[\fbox{W(4c)}] Soit $f : X \longrightarrow Y$, avec $f \in \Fib_W \cap W$, alors les fibres de $f$ sont $W$-asphériques.

\item[\dashedbox{W(4c$'$)}] Comme c), en supposant $f \in \Parf_\omega \cap W$. (C'est l'ancien W(4b$''$), dépouillé du $W_\infty$-parasite.)
\end{itemize}

\msp{page 54}

J'ai cerclé les notions \texttt{[encadrées ici]} qui ont été vraiment utiles. Pratiquement, on peut parier que c'est la notion la plus forte, W(4a), qui est la seule importante. Voici les implications tautologiques (à l'exception de W(4a$'$) $\Longrightarrow$ W(4b), qui est démontré dans la démonstration de la proposition 4 (pages 6-10) (\footnote[61]{Cette démonstration utilise une notion plus forte de $W$-fibration que celle du théorème 3, p. 51.})).

\printedpage{40}
% head: A. Grothendieck : Dérivateurs, Ch. XII, Édition des Universités de Montpellier II et Paris VII

\[
\begin{tikzcd}[column sep=small]
& \mathrm{a} \arrow[d,Rightarrow] & \\
& \mathrm{a}' \arrow[dd,Rightarrow,"{\text{\rule[0.6ex]{2cm}{0.4pt}\ \begin{tabular}[t]{@{}l@{}}implication douteuse,\\ il faut affaiblir\\ un peu b), en\\ renforçant la\\ notion de\\ $W$-fibration\\ utilisée\\ pour énoncer b)\end{tabular}}}"] & \\
& & \\
& \mathrm{b} \arrow[dl,Rightarrow] \arrow[dr,Rightarrow] \arrow[dd,Rightarrow] & \\
\mathrm{b}' & & \mathrm{b}'' \\
& \mathrm{c} \arrow[d,Rightarrow] & \\
& \mathrm{c}' &
\end{tikzcd}
\]

J'ai surtout sué pour montrer que c$'$) $\Longrightarrow$ c), ou c) $\Longrightarrow$ b), sans y arriver. L'effort m'a permis de me débarrasser de diverses idées fausses, et de les remplacer par des justes, cf. plus bas. Il y a quand même un résultat positif qui surnage :

\textbf{Proposition 1.} \textit{La condition b) (et sans doute même b$'$) et b$''$)} (\footnote[62]{et en supposant dans le cas de b$''$, que $W$ satisfait la condition des limites inductives W(6 bis), cf. prop. 5 du `corollaire' de la, page 14.})\textit{, mais je ne vais pas le prouver ici \dots) implique sans doute a) pour $\bar{W}$, le saturé fort de $W$. En tous cas, si $W$ est fortement saturé, les conditions a), a$'$), b) (et sûrement aussi b$'$) et b$''$)) sont équivalentes.}

\dgrblock{dgrb501}{The later resolution of the weakest fibre axiom}
\begin{dcomment}{status}{p54:En tous cas, si $W$ est fortement saturé}{The later resolution of the weakest fibre axiom}
Grothendieck's reservations here belong to the search for a proof from a deliberately small list of axioms. The minimal weak-localizer theorem gives $W_\omega=W_\infty$ even with W(2 bis), by \citet[Theorem~6.1.18]{Cisinski2006}. The proof in section~I, completed by the nondegeneracy argument at p.~20, then shows that W(4c') forces
\[
 W=W_\infty\quad\hbox{or}\quad W=\operatorname{Fl}(\Cat).
\]
The second possibility must be retained: all the fibration axioms hold vacuously when every functor is a weak equivalence. This argument uses the perfect replacement and iterated homotopy-fibre construction, together with the minimality theorem; it does not require filtered-colimit stability or strong saturation of an arbitrary weak localizer. Those properties follow after the displayed classification. For the nontrivial alternative, ordinary homotopy theory gives the full pullback and fibrewise comparison statements. A related published characterization under the full relative-localization hypothesis is \citet[Theorem~6.4.22 and Remark~6.4.23]{Cisinski2006}, formulated using locally constant functors.

The split comparison test $b')$ also detects the weakest perfect-fibre test. If $E\to B$ is $W_\infty$-perfect and belongs to $W$, compare the split replacements $\Phi_B(E)\to\Phi_B(B)$. Its total map is in $W$, since both replacement units are. Their transports are weak equivalences, and the comparison is strictly Cartesian. Thus $b')$ makes each $b\backslash E\to b\backslash B$ a $W$-equivalence. The perfect fibre comparison $E_b\to b\backslash E$ and the initial object of $b\backslash B$ then show that $E_b$ is $W$-aspherical. The implication from the perfect pullback test $a')$ uses the mapping-path comparison already audited in section~I, rather than the disputed reuse of an older fibration definition.

Accordingly the weak perfect test, the split test, and the full tests have the same eventual alternatives. This retrospective conclusion does not turn the unfinished implications in the manuscript into proofs written there. In the claim register their original qualifications remain conjectural or sketched, with the later resolution recorded separately.
\end{dcomment}


\textsc{Démonstration} de b) $\Longrightarrow$ a), quand $W$ est fortement saturé. En fait, moyennant b) (et sans même supposer $W$ fortement saturé), nous allons prouver que si on a
\[
\begin{tikzcd}
X & X' \arrow[l,"q"'] \arrow[d,"f'"] \\
Y \arrow[from=u,"f"'] & Y' \arrow[l,"p"']
\end{tikzcd}
\]
cartésien dans Cat, avec
\msp{page 55}
$f$ une $W$-fibration, et $p \in W$, alors $\underline{q \in \bar{W}}$, i.e. $\hot_W(q)$ \texttt{[est un]} isomorphisme. Pour ceci, nous interprétons la flèche $q$ dans $\Hot_W$ de la façon suivante. Soit
\[
\xi = \hot_{Y,W}(X/Y), \qquad \xi' = p^*(\xi),
\]
\printedpage{41}
% head: A. Grothendieck : Dérivateurs, Ch. XII, Édition des Universités de Montpellier II et Paris VII
qui dit que $q$ s'identifie à la flèche canonique
\[
H_\bullet^{\HOT_W}(Y',\xi') \longrightarrow H_\bullet^{\HOT_W}(Y,\xi) \qquad \text{dans } \Hot_W = \HOT_W(e), \tag*{$(*)$}
\]
où on note de plus que $\xi \in \HOT_W(Y)$ est un coefficient \textit{localement constant} sur $Y$.

Établissons d'abord cette interprétation de $q$. Notons que comme $X \in \underline{\Fib}_W Y$, l'image inverse $\xi' = p^*(\xi)$ n'est autre que $\hot_{Y',W}(X' = X \times_Y Y')$. D'autre part, si $g : Y \longrightarrow e$, $g' : Y' \longrightarrow e$ sont les flèches structurales, les deux membres de $(*)$ s'interprètent resp. comme $g'_!(\xi')$ et $g_!(\xi)$ (les $g'_!$ et $g_!$ au sens du prédérivateur $\HOT_W$). Or le calcul de $g'_!$, $g_!$ nous donne que les deux membres s'identifient resp. à
\[
\hot_W(X'), \qquad \hot_W(X),
\]
et nous admettons que la flèche
\[
\hot_W(X') \longrightarrow \hot_W(X)
\]
n'est autre que $\hot_W(q)$. D'ailleurs, comme $X$ est $W$-fibrée sur $S$, par définition pratiquement, $\xi$ est dans $\HOT^{\mathrm{lc}}_W(S)$. Donc il nous suffira de prouver le

\msp{page 56}

\textbf{Corollaire.} \textit{Supposons que $W$ satisfait à W(4b), et soit $p : S' \longrightarrow S$ dans $W$, $\xi \in \HOT^{\mathrm{lc}}_W(S)$, $\xi' = p^*(\xi)$. Alors}
\[
H_\bullet(S',\xi') \xrightarrow{\ \sim\ } H_\bullet(S,\xi)
\]
\textit{est un isomorphisme dans} $\Hot_W$.

\dgrblock{dgrb502}{Why fully faithful restriction controls integration}
\begin{dcomment}{exposition}{p56:est un isomorphisme dans}{Why fully faithful restriction controls integration}
Grothendieck's proof is driven by the fact that a homotopy colimit is characterized by its maps into constant coefficients. We can state the argument in a form which makes every required comparison visible. Let $p:T\to S$, and suppose that restriction is fully faithful on locally constant coefficients. Let $q:S\to e$ and $r:T\to e$ be the structural maps, and suppose the adjunctions $q_!\dashv q^*$ and $r_!\dashv r^*$ are available. For a locally constant $\xi$, there is a canonical comparison
\[
 c:r_!p^*\xi\longrightarrow q_!\xi.
\]
For every coefficient $A$ at the point, adjunction identifies precomposition by $c$ with
\[
 \operatorname{Hom}(\xi,q^*A)
 \longrightarrow
 \operatorname{Hom}(p^*\xi,p^*q^*A).
\]
Both $\xi$ and $q^*A$ are locally constant, so full faithfulness makes this a bijection. Yoneda implies that $c$ is an isomorphism. Essential surjectivity of restriction is stronger than this argument requires.

The application to a literal pullback $X_T\to X$ needs two further identifications: strict pullback must represent $p^*\xi$, and the chosen integration functors must send these representatives to their total categories, with $c$ represented by the actual projection. Prelissity supplies the first identification; the construction of $q_!$ supplies the second. The source acknowledges these comparisons instead of proving them anew. The same formal Hom argument appears in \citet[Proposition~1.10]{Cisinski2009LocallyConstant}, and its Theorem~1.3 supplies ordinary weak-homotopy invariance for model-category-valued locally constant diagrams. These later results clarify the argument without making W(4b) an unmentioned consequence of W(1,2,3).
\end{dcomment}


\textsc{Démonstration}. Il suffit de prouver que pour tout $\eta \in \mathrm{Ob}\,\Hot_W$ on a une bijection pour les $\Hom(-,\eta)$. Par la formule d'adjonction sur $S$, $S'$, cette application entre les Hom dans $\Hot_W$ s'identifie \texttt{[à]}
\[
\Hom_{\HOT_W(S)}(\xi,\eta_S) \longrightarrow \Hom_{\HOT_W(S')}(\xi',\eta_{S'}),
\]
où $\eta_S = p^*(\eta)$, $\eta_{S'} = p'^*(\eta)$ \texttt{[plutôt} $\eta_S = g^*(\eta)$\texttt{,} $\eta_{S'} = g'^*(\eta)$\texttt{, où} $g : S \longrightarrow e$\texttt{,} $g' : S' \longrightarrow e$ \texttt{sont les flèches structurales]}, donc $\eta_{S'} = p^*(\eta_S)$, et la flèche $(*)$ est la flèche sur les Hom induite par le foncteur
\[
p^*_{\mathrm{lc}} : \HOT^{\mathrm{lc}}_W(S) \longrightarrow \HOT^{\mathrm{lc}}_W(S').
\]
Or moyennant l'axiome W(4b) (que nous n'avons pas utilisé jusqu'à présent), ledit foncteur est une \textit{équivalence de catégories} si $p \in W$ (cf. VI, p. 204).

\hfill q.e.d.

De peur de les oublier, je vais maintenant faire une liste des petites questions que j'ai élucidées chemin faisant en gribouillant.

\dgrblock{dgrb503}{Why preservation of Cartesian arrows matters}
\begin{dcomment}{examples}{p56:De peur de les oublier}{Why preservation of Cartesian arrows matters}
Grothendieck's first observation separates fibrewise smoothness from smoothness of the total functor. A small example makes the lost information explicit. Let $C$ be the one-object category with endomorphism monoid $(\mathbf N,+)$, and put $X=Y=C\times[1]$, both projected to $[1]$. Define an endofunctor $h:X\to Y$ over $[1]$ to be the identity on both fibres and to send a morphism crossing from $0$ to $1$, labelled $n\in\mathbf N$, to the morphism labelled $n+1$.

This is a functor. A composite crosses from $0$ to $1$ at most once, so adding $1$ commutes with precomposition and postcomposition by fibre morphisms. Each fibre functor is the identity and hence is smooth and proper. Both structural projections are products, hence split fibrations. But $h$ does not preserve Cartesian arrows: the Cartesian crossing labelled $0$ is sent to the non-Cartesian crossing labelled $1$.

To test smoothness of $h$, take the target object over $1$ and the arrow in $Y$ from the object over $0$ labelled $0$. A lift with this prescribed image would be a crossing morphism labelled $n$ in $X$ with $n+1=0$. There is none. Its lifting category is empty, so $h$ is not $W_\infty$-smooth. This is a fully explicit instance of the failure claimed in item~1. If the functor preserves the splitting, an arbitrary arrow over $u:s\to t$ factors through the chosen Cartesian lift, and its lift category is the corresponding lift category of the fibre functor at $s$. That identification proves the positive assertion in item~1.
\end{dcomment}


\begin{itemize}[leftmargin=5em,labelsep=0.5em]
\item[1°)] Si $X$, $Y$ sont $W$-lisses sur $S$, et $f : X \longrightarrow Y$ un $S$-morphisme tel que $f_s : X_s \longrightarrow Y_s$ soit $W$-lisse pour tout $s \in \mathrm{Ob}\,S$, il ne s'ensuit \textit{pas} que $f$ soit $W$-lisse. Mais c'est pourtant vrai quand $X$ et $Y$ sont fibrées scindées sur $S$ (ce
\printedpage{42}
% head: A. Grothendieck : Dérivateurs, Ch. XII, Édition des Universités de Montpellier II et Paris VII
\msp{page 57}
qui implique déjà leur $W$-lissité sur $S$), et $f$ compatible avec les scindages et a fortiori cartésien. Énoncé dual pour les catégories cofibrées scindées et les foncteurs strictement cartésiens entre elles qui sont $W$-propres fibre par fibre : ils sont $W$-propres.

\item[2°)] J'étais intéressé, dans un cas de catégories fibrées scindées $X$, $Y$ sur $S$ et d'un foncteur strictement cartésien $f : X \longrightarrow Y$, supposé $W$-parfait fibre par fibre, de conclure que $f$ soit aussi $W$-parfait. (C'était un cas où de plus on sait que $X$, $Y$ sont aussi des $W$-fibrés sur $S$, i.e. les fonteurs changement de base $X_{s'} \longrightarrow X_s$, $Y_{s'} \longrightarrow Y_s$ sont dans $W$. Mais ça n'a pas beaucoup servi \dots) La question qui restait, en 1°), était celle de la $W$\textit{-propreté} de $f$. Mais c'est sans espoir (même pour $Y = S$), en vertu de la brillante % [sic] fonteurs -> foncteurs, lettre omise

\textbf{Proposition 2.} \textit{Supposons que} $S = 0$
$\begin{tikzcd}[ampersand replacement=\&, row sep=tiny, column sep=small, baseline=(Z.base)]
\& a \\
|[alias=Z]| \phantom{0} \arrow[ur] \arrow[dr] \& \\
\& b
\end{tikzcd}$
\textit{, et soit $Z$ une catégorie sur $S$. Alors $\Phi(Z)$ est propre sur $S$ si et seulement si $Z = \emptyset$.}

\dgrblock{dgrb504}{The fork obstructs outgoing lifts}
\begin{dcomment}{exposition}{p57:Alors $\Phi(Z)$ est propre sur $S$}{The fork obstructs outgoing lifts}
The fork is being used for a second purpose here: a free split fibration supplies incoming Cartesian lifts, while properness requires the appropriate outgoing lift categories to be aspherical. The commentary on Chapter~VI, manuscript pp.~190--192, gives the full fork calculation and the support obstruction; we refer to it rather than repeat the calculation.

For the present application, one may see the obstruction at an object $(0,z,a)$ of $\Phi_S(Z)$ whose label in $S$ lies on one branch. The arrow from $0$ to the other branch has no compatible outgoing lift from that object: it would require a morphism in the fork from the first terminal branch to the other. If $Z$ has an object over either branch, this construction applies. If every object of the nonempty category $Z$ lies over $0$, then an outgoing lift towards either terminal branch is again impossible, because no object of $Z$ can label its target. These two cases exhaust the possibilities. In each case the relevant lift category is empty.

Thus the assertion is correct for ordinary properness, and for $W$-properness whenever $W$ is nontrivial: the empty category is then not $W$-aspherical. If $W=\operatorname{Fl}(\Cat)$, every lift category, including the empty one, is $W$-aspherical and every functor is $W$-proper. That degenerate case must be excluded from the relative reading of ``if and only if $Z=\varnothing$''. Fibrewise perfectness of the structural map to $S$ carries no information about this outgoing obstruction, since every map from a category to the point is perfect.
\end{dcomment}


Prenons alors dans 2°) ci-dessus $X = \Phi(Z)$, $Y = S$, $f : \Phi(Y) \longrightarrow S$ \texttt{[plutôt} $f : \Phi(Z) \longrightarrow S$\texttt{]} le morphisme structural (donc les $f_s$ sont $W$-parfaits), alors $f$ n'est propre que si $X = \emptyset$. Pourtant, si on prend $Z$ $W$-fibrée sur $S$, $\Phi(Z)$ l'est aussi, et on a un contre-exemple à 2°).
\end{itemize}

Mais on s'est quand même convaincu du résultat suivant :

\textbf{Proposition 3.} \textit{Soient $X$, $Y$ catégories fibrées scindées sur $S$, $f : X \longrightarrow Y$ $S$-foncteur strictement cartésien, tel que}
\msp{page 58}
\textit{pour tout $s \in S$, $f_s : X_s \longrightarrow Y_s$ soit} $W$-fibrant. \textit{Alors,} si $W$ satisfait W(4a), $f$ \textit{est $W$-fibrant.}

\dgrblock{dgrb505}{A missing hypothesis in the last proposition}
\begin{dcomment}{development}{p58:si $W$ satisfait W(4a), $f$}{A missing hypothesis in the last proposition}
Grothendieck is trying to lift a fibrewise fibration statement through the Grothendieck construction. The preceding paragraph mentions an additional assumption that both structural families are $W$-fibrations over $S$. That assumption is absent from Proposition~3 as printed. The literal statement is false, even for ordinary weak equivalences.

\textbf{A strictly Cartesian counterexample.}\ours\quad
Take $S=[1]$. Let $F,G:S^{\mathrm{op}}\to\Cat$ be diagrams of discrete categories
\[
 F(0)=F(1)=e,\qquad
 G(0)=e,\quad G(1)=\{a,b\},
\]
with the unique transition maps. Let $\eta_0$ be the identity and let $\eta_1:e\to\{a,b\}$ select $a$. These components form a strict natural transformation. Its Grothendieck construction is the inclusion
\[
 h:\{o\to a\}\longrightarrow\{o\to a,\ o\to b\}.
\]
Both source and target are split fibrations over $[1]$, and $h$ preserves the chosen Cartesian lifts: the only nonidentity chosen lift in the source is sent to $o\to a$. Each component of $\eta$ is perfect. Indeed every functor between discrete categories is both a Grothendieck fibration and an opfibration; the only arrows to lift in its target are identities. It is consequently a $W_\infty$-fibration.

Pull back $h$ along the arrow $o\to b$ in the target. The resulting category has one object, lying over $0\in[1]$. Its fibres are $e$ at $0$ and $\varnothing$ at $1$. The second endpoint inclusion is $\varnothing\to e$, which is not a weak equivalence. Thus $h$ is not even a $W_\infty$-pseudofibration. This proves the counterexample. Notice that the target structural family is not a $W_\infty$-fibration: its transport is $\{a,b\}\to e$. This explains exactly how the extra assumption mentioned in the preceding discussion would exclude the example.

Here is a sufficient repair for the fibrewise-perfect situation that motivated the proposition. It states explicitly the two comparison properties used in the proof rather than hiding them in an appeal to ordinary intuition.

\textbf{A sufficient fibrewise-smooth repair.}\ours\quad
Assume W(4a) and W(4b). Let $F,G:S^{\mathrm{op}}\to\Cat$ send every arrow into $W$, and let $\eta:F\Rightarrow G$ be a strict natural transformation. Suppose every $\eta_s$ is both $W$-smooth and a $W$-fibration; fibrewise perfectness is sufficient. Then $h=\int\eta$ is $W$-smooth and a $W$-fibration.

First, an arrow in $\int G$ from $(s,y)$ to the image of $(t,x)$ is given by $u:s\to t$ and $\beta:y\to\eta_sF(u)x$. Its lift category for $h$ is exactly the lift category of $\beta$ for $\eta_s$. This proves smoothness, because the latter category is $W$-aspherical.

Next fix $u:s\to t$. Form
\[
 P=G(t)\times_{G(s)}F(s).
\]
The projection $P\to F(s)$ belongs to $W$ by W(4a), since it is the pullback of $G(u)\in W$ along the $W$-fibration $\eta_s$. The naturality map $F(t)\to P$ composes with this projection to $F(u)\in W$, so it also belongs to $W$. Both $F(t)\to G(t)$ and $P\to G(t)$ are $W$-fibrations. W(4b) therefore shows that for every $z\in G(t)$ the induced relative-fibre transport
\[
 F(t)_z\longrightarrow F(s)_{G(u)z}
\]
belongs to $W$.

To finish, factor an arbitrary arrow of $\int G$ into an arrow within a fibre followed by a chosen Cartesian arrow. For an arrow within $G(s)$, transport on the under-comma categories of $h$ is in $W$ because $\eta_s$ is a $W$-fibration. The comparison with its total under-comma category is a homotopy equivalence: an object above $t$ with base arrow $u:s\to t$ is sent to its Cartesian pullback $F(u)x$ above $s$, giving a right adjoint to the inclusion of the fibre under-comma category. For a chosen Cartesian arrow, the displayed relative-fibre transport and smoothness identify the corresponding under-comma map as an element of $W$. Composition now gives the assertion for every arrow of $\int G$. The dual form of Proposition~8, together with smoothness, proves that $h$ is a $W$-fibration.

This repair covers the preceding discussion when its structural families and fibrewise-perfect hypotheses are retained. It does not claim that adding only the omitted structural-fibration hypothesis proves the entire printed formulation. We checked the later treatments of smooth functors and locally constant families cited in this commentary and did not locate this exact small counterexample or this explicit sufficient repair there. These are deductions given here, without a claim of priority.
\end{dcomment}


(J'espérais avoir que $f$ est $W$-fibrant `à l'\oe{}il', pour prouver que W(4c) implique W(4b). Mais j'étais obligé de supposer W(4a) \dots)


\dgrblock{dgrb506}{The coherence condition in a later gluing theorem}
\begin{dcomment}{later}{p58:J'espérais avoir que $f$}{The coherence condition in a later gluing theorem}
The last failed shortcut has a useful later analogue. \citet[Theorem~7.1]{Rezk1998Sharp} proves that the homotopy colimit of a diagram of sharp maps is sharp when every naturality square is homotopy Cartesian. This square condition is the coherence missing from a bare list of fibrewise fibration assertions. In the counterexample just given, the square of discrete fibres
\[
\begin{tikzcd}
 e\arrow[r]\arrow[d] & e\arrow[d]\\
 \{a,b\}\arrow[r] & e
\end{tikzcd}
\]
is not homotopy Cartesian: the homotopy pullback of the other three corners has two points.

This is a comparison of mathematical mechanisms, not an identification of all the models. Thomason's comparison between the nerve of a Grothendieck construction and a homotopy colimit is a weak equivalence; such a comparison alone does not transport the property that every \emph{strict} pullback is homotopy Cartesian. To deduce a statement about the literal functor of Proposition~3, the comparison must also control the pullbacks that occur in its fibration test. That is why the source's dependence on W(4a), and the omitted control of structural transport, cannot be removed by a slogan about homotopy colimits.
\end{dcomment}
\begin{otherlanguage}{english}
\dgroverview{XI-XIIa}{Chapter XII \textperiodcentered\ \textup{XI-XIIa} \textperiodcentered\ Overview}
\begin{otherlanguage}{english}
\noindent\textbf{Fibrations, localizers, and the coherence of fibres}

\medskip
The pages assembled here turn a useful geometric intuition into a problem of foundations. A family of categories should behave like a fibration when its strict fibres calculate homotopy fibres and vary coherently over the base. Grothendieck first characterizes that behaviour for ordinary weak homotopy equivalences. He then asks which axioms on an abstract class \(W\) of functors permit the same constructions, and discovers that several apparently natural requirements already force the whole ordinary homotopy theory. The resulting revisions are part of the argument: the word “fibration” changes its tests as the desired scope of the theory becomes clearer.

These pages comprise Chapter XI, manuscript pp.~127--136, and Chapter XII, manuscript pp.~1--58, through the end of section~IV. These are separate manuscript sequences. The edition explicitly says that Chapter XI has no pp.~1--126; their absence is not a gap in the text. We write \(BC=|NC|\) for the classifying space of a category and \(W_\infty\) for the functors whose nerves are weak homotopy equivalences.

\medskip
\noindent\textbf{The geometric theorem and its coefficient proof.}
Chapter XI begins with a property of the actual pullback functor:
\[
 f^*:\Cat/Y\longrightarrow\Cat/X
 \quad\hbox{preserves }W_\infty.
\]
This is a universal condition on changes of base, not a condition merely on the homotopy types of the individual fibres. Theorem~1 identifies it with universal comma-category tests and with base-change and local-constancy statements for direct images of locally constant coefficients. In the terminology later used for model categories, these are sharp maps for the Thomason weak equivalences \citep[\S2]{Rezk1998Sharp}. Quillen's Theorem B explains the geometry of the comma tests: when their transition functors are weak equivalences, the strict fibre, compared with the appropriate comma category, represents the homotopy fibre \citep[\S1]{Quillen1973}.

Most of the coefficient argument is complete as written. The degree-zero comparison and the first-quadrant Leray comparison require no new finiteness hypothesis. The nonabelian passage needs more care. A component of a torsor groupoid has the automorphism group of a representative torsor; this group can be an inner form of the coefficient group, and its direct image is not generally the fixed group printed in the proposed obstruction sequence. We therefore replace that sequence by an argument with the groupoids of torsors on \(X/y\). The universal degree-zero hypothesis gives full faithfulness of their transition functors; the degree-one hypothesis gives essential surjectivity. The same two hypotheses identify their values with the torsor groupoids on strict fibres. Descent reconstructs global torsors from this locally constant diagram, and homotopy invariance of such diagrams proves the required comparison \citep[Theorem~1.3]{Cisinski2009LocallyConstant}.

Theorem~2 then detects an equivalence between two such families on one fibre in each connected component of the base. The source's local reduction for torsors becomes precise through the same descent argument. Its converse uses a contractible replacement of the base. We give an independent construction by factoring a point in Thomason's model category \citep[Theorem~4.9]{Thomason1980}, and also check the specific path-category dependency that Chapter XII recalls from Chapter VII. This separates the existence of a suitable replacement from the stronger task of constructing it explicitly with both properness and smoothness.

\medskip
\noindent\textbf{What the axioms force.}
Chapter XII starts with identities and cancellation, a theorem-A principle, and the asphericity of categories with a terminal object. Two forms of theorem A occur. W(2) is relative: equivalences on comma categories over every object of a third category imply equivalence of the total functor. W(2 bis) is absolute: aspherical comma categories imply that a functor is an equivalence. The manuscript deliberately uses the latter in several arguments. Modern “fundamental localizer” normally includes the former. We keep the distinction throughout the audit.

The minimum conjectured in these pages is now known. Cisinski identifies \(W_\infty\) as the smallest fundamental localizer \citep[Theorem~2.2.11]{Cisinski2004Minimal}, and as the smallest weak fundamental localizer with the absolute theorem-A axiom \citep[Theorem~6.1.18]{Cisinski2006}. Strong saturation and stability under filtered colimits are also established for full fundamental localizers; those results are not assertions about every class satisfying only the weaker axiom. An inference from “invertible in the localization” to “belongs to \(W\)” is justified at the point where its applicable hypothesis is available.

The fibration axioms add a rigidity that the basic localizer axioms do not possess. Their weakest perfect-fibre form, combined with nondegeneracy, makes homotopy fibres of \(W\)-maps \(W\)-aspherical. Iterating the loop construction then detects all homotopy groups. We repair the source's nondegeneracy argument without its unjustified restriction to a component: a missed component produces an empty homotopy fibre, while two identified components produce a disconnected aspherical fibre, whose diagonal in its square again has an empty homotopy fibre. Asphericity of the empty category collapses the theory. Thus the alternatives are \(W_\infty\) and the class of every functor. The second alternative cannot be discarded from the opening assertion about derivators; the terminal derivator realizes it.

The truncation examples explain why a more flexible fibration notion is needed. The class \(W_n\) retains homotopy groups only through degree \(n\), and is a genuine localizer. But the based-path family of an \((n+1)\)-sphere has an invisible total space and an invisible base while its fibre retains a nonzero \(n\)-th homotopy group, or a nontrivial component set when \(n=0\). Hence the strongest fibre axioms fail. In degree zero, the exact sequence must be read in pointed sets and with the appropriate group action: a cokernel of groups does not describe the component set in general.

\medskip
\noindent\textbf{Lift categories supply the missing structure.}
Section~III develops proper and smooth functors through categories of lifts with prescribed source or target. Those categories turn universal comma comparisons into workable tests, prove stability under base change and composition, and explain why Grothendieck (op)fibrations furnish examples. They also support the passage from fibrewise equivalence to equivalence of total categories. The later formulations by Maltsiniotis describe the corresponding exact squares, while the derivator account explains the associated base-change transformations \citep{Maltsiniotis2005Asphericite,Maltsiniotis2012Exact,Groth2013}.

The proposed weakening to endpoint equivalences over every arrow loses a necessary compatibility. We give a poset with eight objects over \([2]\). Its fibres have two contractible components, every interval pullback has both endpoint inclusions in \(W_\infty\), and transport along \(0\to1\) and \(1\to2\) preserves the component labels. Transport along \(0\to2\) exchanges them. Its total nerve is an eight-cycle. Consequently the edge transports do not compose, and the pseudofibration question raised on manuscript p.~40 has a negative answer. The complete proof is independent of the accompanying finite computation.

The successful replacement uses universal preproperness or prelissity in addition to the interval test. Comma categories then provide actual functorial transport. A split replacement converts this into a diagram of categories whose arrows lie in \(W\). For a categorically contractible base, a finite zigzag of natural transformations gives an explicit finite zigzag of fibrewise and total \(W\)-equivalences. This proves the contractible-base portion of Theorem~3 directly, without extracting a zigzag from an abstract isomorphism in a localization. Here “contractible” has the manuscript's categorical meaning; a weakly contractible nerve need not admit such a finite contraction.

\medskip
\noindent\textbf{The last proposition and the limits of the repair.}
Section~IV returns to the relations among the fibration axioms, then proposes a fibrewise gluing statement. Its final Proposition~3 is false literally as printed. Over \([1]\), take the strictly Cartesian inclusion of the split fibration \(\{o\to a\}\) into \(\{o\to a,o\to b\}\). Every fibre map is a functor between discrete categories and is perfect. Pulling the total functor back along \(o\to b\) gives endpoint fibres \(e\) and \(\varnothing\); it is not even a pseudofibration. The target structural family is not a \(W_\infty\)-fibration, precisely the extra condition mentioned in the preceding discussion and omitted from the proposition.

We prove a sufficient repair for that preceding, fibrewise-perfect situation: the structural transports lie in \(W\), the fibre maps are smooth \(W\)-fibrations, and W(4a) and W(4b) control the relevant pullback and fibre comparison. We do not claim that adding only the omitted structural condition repairs the full printed statement. Rezk's later gluing theorem has an analogous additional requirement that naturality squares be homotopy Cartesian \citep[Theorem~7.1]{Rezk1998Sharp}; passing between models does not by itself remove the need to control strict pullbacks.

The marked developments give these deductions with complete proofs and bounded accounts of the literature examined, without a claim of priority. The registers distinguish what the source proves, sketches, assumes, or conjectures from the status supplied by the commentary. They preserve edition slips in the French text and apply only corrections of the transcription. Questions not decided under their exact weaker hypotheses are identified as limitations of this audit, rather than declared open in the literature.

\end{otherlanguage}
\end{otherlanguage}
\dgrchunkend

% DER_XIIb-a: Grothendieck's text with the commentary set in, then the overview, inside a group that applies the chunk's own definitions
\dgrpart{1}{XIIb-a}{Section 5, part 1 of 2, manuscript pages 59--122}{XIIb-a\enspace\textperiodcentered\enspace Chapter XII, section 5, part 1 of 2\enspace\textperiodcentered\enspace manuscript pages 59--122}{Chapter XII \textperiodcentered\ \textup{XIIb-a}}
\dgrchunkbegin
% DER_XIIb-a: the definitions of the chunk's preamble, applied inside its group
\let\printedpage\relax\newcommand{\printedpage}[1]{\marginpar{\footnotesize p.~#1}}
\let\msp\relax\newcommand{\msp}[1]{\marginpar{\footnotesize [#1]}}
\let\uHom\relax\newcommand{\uHom}{\underline{\mathrm{Hom}}}
\let\Hom\relax\newcommand{\Hom}{\mathrm{Hom}}
\let\Cat\relax\newcommand{\Cat}{\mathrm{Cat}}
\let\Ens\relax\newcommand{\Ens}{\mathrm{Ens}}
\let\id\relax\newcommand{\id}{\mathrm{id}}
\let\Ob\relax\newcommand{\Ob}{\mathrm{Ob}\,}
\let\circled\relax\newcommand{\circled}[1]{\tikz[baseline=(c.base)]\node[draw,circle,inner sep=1pt](c){#1};}
\let\cercle\relax\newcommand{\cercle}[1]{\tikz[baseline=(c.base)]\node[draw,circle,inner sep=1pt](c){#1};}
\let\defeq\relax\newcommand{\defeq}{\overset{\text{déf}}{=}}
\let\upi\relax\newcommand{\upi}{\underline{\pi}}
\let\uCh\relax\newcommand{\uCh}{\underline{\mathrm{Ch}}}
\let\bDelta\relax\newcommand{\bDelta}{\boldsymbol{\Delta}}
\let\hra\relax\newcommand{\hra}{\lhook\joinrel\longrightarrow}

% Reader layout additions only; all edition slips in the body are retained.

% ---- from Der-ChXIIb-012_p001-005.tex
\printedpage{43}
% head: A. Grothendieck : Dérivateurs, Ch. XII, Édition des Universités de Montpellier II et Paris VII ; foot: Version 7 décembre 2002

\msp{page 59}

\section*{V\quad Sommes amalgamées et carrés $W$-cocartésiens dans $\mathrm{Cat}$}

\subsection*{1\quad Généralités sur $U \overset{i}{\hookrightarrow} X \overset{j}{\hookleftarrow} T$, et les sommes amalgamées.}

\subsubsection*{1.1\quad Données de recollement entre un ouvert et un fermé.}

Soit $X$ dans Cat, $U$ un ouvert dans $X$, $T$ le complémentaire fermé. (\textbf{NB} Les ouverts correspondent bien aux ouverts de $\mathrm{Ob}\,X$, topologisé par sa relation de préordre naturelle $x \leq y \iff \mathrm{Hom}_X(x,y) \neq \emptyset$, de même pour les fermés. Les topologies associées aux relations de préordre sont celles où toute intersection d'ouverts est encore ouverte, i.e. toute réunion de fermés est encore fermée. Donc les ouverts d'une telle topologie sont les fermés d'une topologie duale, correspondant à la relation de préordre opposée. La relation de préordre associée à la topologie est
\[
x \leq y \qquad\qquad \iff \qquad\qquad x \in \overline{\{y\}},
\]
la topologie associée au préordre sur $I$ est celle pour laquelle les ouverts sont les réunions d'ensembles de la forme $I/x$ (lesquels forment une base d'ouverts), ou encore les \textit{`cribles'} ($x \in U$, $y \leq x \Longrightarrow y \in U$), les fermés sont les réunions d'ensembles $x\backslash I$, c'est-à-dire les cocribles. Pour toute partie $A$ de $I$, il y a un plus petit ouvert contenant $A$, qui est aussi son plus petit voisinage, c'est la réunion des $I/x$ pour $x \in A$. Et dualement il y a un plus grand fermé contenu dans $A$, dont le complémentaire est le plus petit ouvert contenant $B = I\backslash A$, donc $\bigcup_{x\in B} I/x$. L'adhérence de $A$ est la réunion des adhérences $x\backslash I$ de ses points.)

\msp{page 60}

\dgrblock{dgrb507}{An open–closed cut remembers a profunctor}
\begin{dcomment}{exposition}{p59}{An open--closed cut remembers a profunctor}
Grothendieck begins with a cut of a category which permits arrows to cross in only one direction. This is the reason that the subsequent pushout calculation is manageable. An open subcategory $U\subset X$ is a full sieve: if an arrow has its target in $U$, its source belongs to $U$. The complementary full subcategory $T$ is therefore a cosieve, and there are no arrows from $T$ to $U$. Once the compositions internal to $U$ and $T$ are known, the remaining information is precisely the action of these compositions on the sets of arrows from $U$ to $T$.

There is a reversal in the displayed specialization convention on this page. With $x\leq y$ meaning that $\Hom_X(x,y)$ is nonempty and with the opens taken to be downsets, one has
\[
 \overline{\{x\}}=\{y\mid x\leq y\},\qquad
 x\leq y\ \Longleftrightarrow\ y\in\overline{\{x\}}.
\]
The edition really prints $x\in\overline{\{y\}}$. That formula is incompatible with its following description of opens as sieves; it is a source slip, not a transcription error. For the walking arrow $0\to1$, the opens are $\varnothing,\{0\},\{0,1\}$ and the closure of $\{0\}$ is the whole space. All the subsequent formulas for closure and boundary use the downset convention and are correct with it. In particular, the smallest open containing $A$ is its downward closure, while its ordinary closure is its upward closure. No antisymmetry is required: categorically isomorphic objects, and more generally objects connected by arrows in both directions, cannot be distinguished by these opens.

For completeness, the reconstruction in Proposition~1 can be written without a choice of compositions. Given
\[
 H:U^{\mathrm{op}}\times T\longrightarrow\Ens,
\]
form a category $\operatorname{Coll}(H)$ whose objects are the tagged disjoint union of those of $U$ and $T$, whose hom sets within each part are the original ones, whose hom sets from $T$ to $U$ are empty, and whose hom sets from $u\in U$ to $t\in T$ are $H(u,t)$. Define composition with $a:u'\to u$ by $h\mapsto H(a,1)(h)$ and composition with $b:t\to t'$ by $h\mapsto H(1,b)(h)$. The two operations commute by bifunctoriality. Every composable string contains at most one arrow crossing from $U$ to $T$, so these identities, together with associativity in the two pieces, prove associativity in the collage. The identities follow from the identity laws for $H$.

A functor preserving the two parts consists of functors $f_U,f_T$ and a natural transformation $H\to H'(f_U^{\mathrm{op}}\times f_T)$. The naturality squares are exactly the equations expressing preservation of composition involving a crossing arrow. This proves full faithfulness of the correspondence of Proposition~1; the collage construction proves essential surjectivity. With tagged objects there is a specified inverse up to canonical isomorphism, which explains the editors' preference for an equivalence over an unqualified isomorphism of categories. Today $H$ is called a profunctor and $\operatorname{Coll}(H)$ its collage. The point of the terminology is that a category can be assembled from its internal categories and a bimodule of arrows between them.
\end{dcomment}


La situation $(X,U,T)$ se reconstitue, à isomorphisme unique près, par la connaissance de

\begin{enumerate}[label={}]
\item[$1^\circ$)] Les deux catégories $U$, $T$.
\item[$2^\circ$)] Le bifoncteur $\displaystyle\left\{\begin{aligned}
&H : U^\circ \times T \longrightarrow \mathrm{Ens}\\[1ex]
&H(u,t) = \mathrm{Hom}_X(u,t).
\end{aligned}\right.$
\end{enumerate}

De plus, ces données $1^\circ$) $2^\circ$) en sont soumises à aucune restriction. On a, plus exactement % [sic] en sont -> ne sont, negation required by the sense

\textbf{Proposition 1.} \textit{Un isomorphisme} \texttt{[plutôt une équivalence de catégories]} \textit{entre la catégorie $\mathcal{T}$ dont les objets sont les triples $(X,U,T)$, avec $X \in \mathrm{Ob}\,\mathrm{Cat}$, $U \in \mathrm{Ouv}(X)$, $T \in \mathrm{Ferm}(X)$, $U$ et $T$ complémentaires l'un de l'autre, les morphismes de $(X,U,T)$ dans $(X',U',T')$ étant les foncteurs $f : X \longrightarrow X'$ qui induisent $U \longrightarrow U'$, $T \longrightarrow T'$ (d'où
\printedpage{44}
% head: A. Grothendieck : Dérivateurs, Ch. XII, Édition des Universités de Montpellier II et Paris VII ; foot: Version 7 décembre 2002
$U = f^{-1}(U')$, $T = f^{-1}(T')$), avec la catégorie $\mathcal{T}'$ des triples $(U,T,H)$, de deux catégories $U$, $T$ et d'un bifoncteur}
\[
H : U^\circ \times T \longrightarrow \mathrm{Ens}
\]
\hspace*{2em}\textit{(ou encore un préfaiscau $U^\circ \longrightarrow T^\vee$ sur $U$, donc un objet de} $\uHom(U^\circ,T^\vee)$\textit{)),} % [sic] préfaiscau -> préfaisceau, misspelling

\textit{un morphisme de $(U,T,H)$ dans $(U',T',H')$ étant un triple}
\[
(f_U, f_T, u) \qquad\qquad
\left\{\begin{array}{l}
\left.\begin{array}{llll}
f_U : & U & \longrightarrow & U' \\[1ex]
f_T : & T & \longrightarrow & T'
\end{array}\right\}\text{des foncteurs}\\[2ex]
u : H \longrightarrow H' \circ (f_U^\circ \times f_T),\\
\text{i.e. } H(u,t) \longrightarrow H'(f_U(u), f_T(t))\\
\text{un homomorphisme fonctoriel}\\
\text{(flèche dans } (U \times T^\circ)^\wedge \simeq \uHom(U^\circ,T^\vee)\text{)),}
\end{array}\right.
\]
\textit{composition des flèches évidente.}

Si $U \hookrightarrow X \hookleftarrow T$ est un triple, et
\[
f_{U'} : U' \longrightarrow U,\ f_{T'} : T' \longrightarrow T
\]
des flèches dans Cat, il y a une façon privilégiée de `recoller' $U'$ et $T'$, de façon à ce que l'objet $(X',U',T')$ s'envoie dans $(X,U,T)$ en induisant

\msp{page 61}

$f_{U'}$, $f_{T'}$:
\[
\begin{tikzcd}[column sep=large, row sep=large]
U' \arrow[r, hook, "i'"] \arrow[d, "f_{U'}"'] & X' \arrow[d, "f"] & T' \arrow[l, hook', "j'"'] \arrow[d, "f_{T'}"] \\
U \arrow[r, hook, "i"] & X & T, \arrow[l, hook', "j"']
\end{tikzcd}
\]
elle s'obtient à l'aide des bifoncteurs
\[
\left\{\begin{aligned}
&H'(u',t') \overset{\text{déf}}{=} H(f_{U'}(u'), f_{T'}(t'))\\
&\text{i.e. } H' = H \circ (f_{U'}^\circ \times f_{T'}).
\end{aligned}\right.
\]
Ce triple, i.e. cet objet de $\mathcal{T}$, représente le foncteur suivant sur $\mathcal{T}$ :
\[
(X'',U'',T'') \longmapsto \mathrm{Hom}_{\mathrm{Cat}}(U'',U) \times \mathrm{Hom}_{\mathrm{Cat}}(T'',T),
\]
i.e. on a pour tout objet $(X'',U'',T'')$ de $\mathcal{T}$ une bijection canonique, fonctorielle en cet objet
\[
\mathrm{Hom}_{\mathcal{T}}((X'',U'',T''),(X',U',T')) \simeq \mathrm{Hom}_{\mathrm{Cat}}(U'',U) \times \mathrm{Hom}_{\mathrm{Cat}}(T'',T). \tag*{$(*)$}
\]

\printedpage{45}
% head: A. Grothendieck : Dérivateurs, Ch. XII, Édition des Universités de Montpellier II et Paris VII ; foot: Version 7 décembre 2002

Nul doute qu'on a même
\[
\uHom_{\mathcal{T}}((X'',U'',T''),(X',U',T')) \simeq \uHom_{\mathrm{Cat}}(U'',U) \times \uHom_{\mathrm{Cat}}(T'',T), \tag*{$(**)$}
\]
tout aussi fonctoriel, dont $(*)$ se déduit en passant aux ensembles d'objets. On dit que le triple $(X',U',T')$ est le triple \textit{induit} par le triple $(X,U,T)$, via les foncteurs $f_{U'}$, $f_{T'}$.

\textbf{NB} Si $f_{U'}$, $f_{T'}$ sont pleinement fidèles, alors il en est de même de $f : X' \longrightarrow X$. Nous utiliserons cette opération de recollement privilégié surtout quand $U' = U$, $f_{U'} = \mathrm{id}_U$, ou inversement $T' = T$, $f_{T'} = \mathrm{id}_T$.

\subsubsection*{1.2\quad Relation avec le recollement des topos.}

La situation d'un triple
\[
(\mathcal{X},\mathcal{U},\mathcal{T})
\]
d'un topos $\mathcal{X}$, d'un ouvert $\mathcal{U}$ de $\mathcal{X}$, et d'un

\msp{page 62}

\dgrblock{dgrb508}{The induced collage has a relative universal property}
\begin{dcomment}{exposition}{p61}{The induced collage has a relative universal property}
The induced triple is meant to express change of the two categories on which the gluing data live. Its universal property has to remember a map to the original triple. Without that map the displayed formulas $(*)$ and $(**)$ are false; both formulas are reproduced faithfully from the edition.

Write $\mathcal T$ for the category of triples and $P:\mathcal T\to\Cat\times\Cat$ for $(X,U,T)\mapsto(U,T)$. Let $b=(f_{U'},f_{T'}):(U',T')\to(U,T)$ and let $X'$ be the collage of $b^*H$. For any triple $Y=(Y,V,S)$ the correct bijection is
\[
\begin{split}
 \Hom_{\mathcal T}(Y,X')\simeq
 &\Hom_{\mathcal T}(Y,X)\mathop{\times}_{
 \Hom(V,U)\times\Hom(S,T)}\\[-2pt]
 &\hspace{18mm}\bigl(\Hom(V,U')\times\Hom(S,T')\bigr).
\end{split}
\]
On the right one chooses a morphism $g:Y\to X$ and functors $a:V\to U'$, $c:S\to T'$ such that $g|_V=f_{U'}a$ and $g|_S=f_{T'}c$. The map on mixed hom sets of $g$ then has type
\[
 H_Y\longrightarrow H(f_{U'}a-,f_{T'}c-)
                 =(b^*H)(a-,c-).
\]
It defines a unique lift $Y\to X'$, and restriction of a lift recovers exactly these data. Thus $X'\to X$ is a cartesian arrow of $P$, not an object representing arbitrary pairs of functors to $U$ and $T$.

The same argument works on natural transformations. A transformation between two lifts is a pair of transformations on $V,S$ satisfying the naturality equations for every mixed arrow, and these are precisely the equations in the corresponding strict pullback of functor categories. This is the correctly typed version of $(**)$. One may express the statement in the language of a split fibration of gluing data, using the literal restriction operation on bifunctors; the corresponding assertion for unlabelled collages is invariant under the canonical isomorphisms.

The missing datum cannot be supplied by merely inserting primes on the right of the printed formulas. Take $U=T=U'=T'=e$, both change-of-base functors the identity, and $H=\{a,b\}$. Then $X'=X$ is the category with two parallel arrows from the open object to the closed object. There are two part-preserving functors from the walking arrow to $X'$, whereas there is only one pair of functors from its two singleton parts to $U',T'$. The choice of crossing arrow is precisely the information lost by the printed right-hand side. The following observation about full faithfulness is unaffected: on the two diagonal sorts of hom sets it follows from $f_{U'},f_{T'}$, and on mixed hom sets it follows from the defining equality $H'=b^*H$.
\end{dcomment}


sous-topos fermé complémentaire, est toute similaire (\textbf{NB} Si $\mathcal{X} = \mathrm{Top}(X)$, avec $X$ dans Cat, alors $\mathrm{Ouv}(\mathcal{X}) \simeq \mathrm{Ouv}(X)$, puisque les ouverts correspondent aux sous-faisceaux du faisceau final sur $\mathcal{X}$, i.e. du préfaisceau final sur $X$. Cette bijection est compatible avec le passage aux complémentaires.) La théorie des topos dit que les couples $(\mathcal{X},\mathcal{U},\mathcal{T})$ \texttt{[plutôt triples]} se reconstituent, à équivalence définie à isomorphisme unique près, à l'aide d'un triple
\[
\left\{\begin{aligned}
&\mathcal{U},\ \mathcal{T} \text{ deux topos}\\[1ex]
&\varphi : \mathrm{Fais}\,(\mathcal{U}) \longrightarrow \mathcal{T} \text{ un foncteur exact à gauche (`accessible'),}
\end{aligned}\right.
\]
ce foncteur $\varphi$ étant obtenu, quand on part de
\[
\mathcal{U} \xrightarrow{\;i\;} \mathcal{X} \xleftarrow{\;j\;} \mathcal{T},
\]
comme
\[
\varphi = j^*i_*.
\]
De plus, le triple $(\mathcal{U},\mathcal{T},\varphi)$ peut être choisi arbitrairement, avec la seule restriction (pratiquement anodine, car toujours satisfaite pour tous les foncteurs naturels qu'on rencontre) que $\varphi$ soit (en plus d'être exact à gauche) `accessible'.

Dans le cas actuel $\mathcal{X} = \mathrm{Top}(X)$, $\mathcal{U} = \mathrm{Top}(U)$, $\mathcal{T} = \mathrm{Top}(T)$, le foncteur $\varphi = j^*i_*$ est non seulement exact à gauche, mais même il commute aux $\varprojlim$ quelconques, puisqu'il en est ainsi (toujours) de $i_*$, et aussi (parce que le $j$ topossique est induit par le $j : T \longrightarrow X$ catégorique) de $j^*$. En fait, on
\printedpage{46}
% head: A. Grothendieck : Dérivateurs, Ch. XII, Édition des Universités de Montpellier II et Paris VII ; foot: Version 7 décembre 2002
\msp{page 63}
trouve ceci :

\dgrblock{dgrb509}{The presheaf case is already present in SGA 4}
\begin{dcomment}{context}{p62}{The presheaf case is already present in SGA 4}
Grothendieck's comparison with toposes recovers a particularly exact antecedent. In \citet[Expos\'e~IV, Theorem~9.5.4]{SGA4I1972}, an open subtopos and its closed complement reconstruct the ambient topos through the accessible left exact gluing functor. Exercise~9.5.10 of the same expos\'e then treats presheaf toposes: its parts~(a)--(c) identify opens with sieves, introduce the cross-hom bifunctor, and characterize the presheaf case by preservation of products by the gluing functor. Thus Proposition~2 is a return to a construction already formulated there, now serving a calculation of categorical pushouts and their homotopy properties. The letters $i,j$ are interchanged between the SGA convention and the present pages.

The category of sheaves in the glued topos is the category of triples
\[
 (F,G,\theta),\qquad F\in\mathcal U,\quad G\in\mathcal T,
 \quad\theta:G\longrightarrow\varphi(F).
\]
The comparison arrow records how a section on the closed side approaches the open side. For $X=\operatorname{Coll}(H)$, a presheaf $E$ on $X$ restricts to $F=E|_U$ and $G=E|_T$. Its values on the mixed arrows are maps
\[
 G(t)\times H(u,t)\longrightarrow F(u)
\]
natural in $u,t$. Transposing in the first variable produces $G(t)\to\Hom_{U^\wedge}(H(-,t),F)$, exactly the required gluing arrow. Conversely these maps reconstruct the restriction maps of $E$ on all arrows; their naturality supplies the mixed composition laws. This directly explains why the ordinary categorical collage and the topos gluing construction agree. The comparison is not an analogy between two shapes of diagrams: both constructions encode the same mixed restriction maps.
\end{dcomment}


\textbf{Proposition 2 :} \textit{Considérons une situation de recollement de deux topos $\mathcal{U}$, $\mathcal{T}$ par $\varphi : \mathrm{Fais}\,\mathcal{U} \longrightarrow \mathrm{Fais}\,\mathcal{T}$, foncteur exact à gauche. Supposons que $\mathcal{U} = \mathrm{Top}(U)$, $\mathcal{T} = \mathrm{Top}(T)$, avec $U$ et $T$ dans $\mathrm{Cat}$. Pour que le topos recollé $\mathcal{X}$ soit équivalent à un topos $\mathrm{Top}(X)$, avec $X$ dans $\mathrm{Cat}$, il faut et il suffit que le foncteur de recollement $\varphi$ commute aux $\varprojlim$ quelconques, ou ce qui revient au même, qu'il ait un adjoint à gauche :}
\[
\psi : \underbrace{\mathrm{Fais}\,\mathcal{T}}_{= T^\wedge} \longrightarrow \underbrace{\mathrm{Fais}\,\mathcal{U}}_{= U^\wedge}.
\]

On s'attend dès lors que la donnée de recollement $\varphi$, ou ce qui revient au même le foncteur $\psi$ ayant un adjoint à droite, i.e. commutant aux $\varinjlim$ quelconques, revient à celle d'une donnée de recollement $H$
\[
H : U^\circ \times T \longrightarrow \mathrm{Ens}
\]
pour $U$ et $T$. Mais on sait bien que les foncteurs $\psi : T^\wedge \longrightarrow \mathcal{C}$ dans une catégorie $\mathcal{C}$ avec $\varinjlim$ quelconques, qui commutent auxdites $\varinjlim$, forment une catégorie équivalente (par le foncteur $\psi \longmapsto \psi|T(\subseteq T^\wedge)$) à celle des foncteurs
\[
\psi_0 : T \longrightarrow \mathcal{C} \quad\ (= \underbrace{U^\wedge}_{\uHom(U^\circ,\mathrm{Ens})} \text{ en l'occurrence}),
\]
\msp{page 64}
et la donnée d'un tel foncteur $\psi_0$ revient bien à celle d'un $H$. Donc

\textbf{Corollaire.} \textit{La catégorie des données de recollement topossique $\varphi : U^\wedge \longrightarrow T^\wedge$ entre} $\mathrm{Top}(U)$, $\mathrm{Top}(T)$\textit{, telles que le topos recollé soit un `topos} paradisiaque\textit{', i.e. équivalent à un topos de la forme $X^\wedge$ (avec $X$ dans} Cat\textit{), est équivalente à celle des données de recollement $H : U^\circ \times T \longrightarrow$} Ens \textit{entre $U$ et $T$. La donnée topossique $\varphi$ associée à $H$ est le foncteur adjoint à droite de $T^\wedge \longrightarrow U^\wedge$ déduit de $T \xrightarrow{H} U^\wedge$ par prolongement commutant aux $\varinjlim$, ou encore elle est donnée par}
\[
\varphi(\underbrace{F}_{\in U^\wedge})(t) \simeq \mathrm{Hom}_{U^\wedge}(H(t),F),
\]
\textit{où dans le second membre $H$ est interprété comme un foncteur $T \longrightarrow U^\wedge$. Le topos recollé par $\varphi$ est canoniquement équivalent à} $\mathrm{Top}(X)$\textit{, où $X$ est l'objet de} Cat \textit{obtenu par recollement à l'aide de $H$.}

Du point de vue de la théorie des faisceaux, cela donne l'interprétation de $H$, comme exprimant le foncteur
\[
\varphi = j^*i_* : U^\wedge \longrightarrow T^\wedge
\]
par la formule précédente.

\printedpage{47}
% head: A. Grothendieck : Dérivateurs, Ch. XII, Édition des Universités de Montpellier II et Paris VII ; foot: Version 7 décembre 2002

\msp{page 65}

\dgrblock{dgrb510}{The adjunction which characterizes the presheaf gluing}
\begin{dcomment}{exposition}{p64}{The adjunction which characterizes the presheaf gluing}
The extra preservation of limits isolates those gluing functors that come from actual sets of arrows. Let $H_t=H(-,t)\in U^\wedge$. Its colimit-preserving extension from the Yoneda image of $T$ is
\[
 \psi(G)=\int^{t\in T}G(t)\times H_t,
 \qquad
 \varphi(F)(t)=\Hom_{U^\wedge}(H_t,F).
\]
Here the coend is an ordinary colimit in $U^\wedge$. A map $\psi(G)\to F$ is a family of maps $G(t)\times H_t\to F$ compatible with the arrows of $T$. By the ordinary set-valued adjunction between product and hom, this is exactly a natural transformation $G\to\varphi(F)$. This proves $\psi\dashv\varphi$ and the displayed formula of the corollary. It also proves that $\varphi$ preserves all small limits.

Conversely, suppose that an accessible gluing functor $\varphi:U^\wedge\to T^\wedge$ preserves all small limits. The adjoint functor theorem for locally presentable categories supplies a left adjoint $\psi$. Set $H_t=\psi(y_Tt)$. Every presheaf is the colimit of its representables indexed by its category of elements; since $\psi$ is a left adjoint, this determines $\psi$ from $H$. The preceding adjunction then identifies its right adjoint with $\varphi$. Applying the explicit construction of the preceding block identifies the glued topos with $\operatorname{Coll}(H)^\wedge$. This proves the sufficient direction of Proposition~2 with the size assumption made explicit. The accessibility assumption belongs to the ambient gluing problem; it should not be silently discarded in an abstract claim that every limit-preserving functor has a left adjoint.

For the converse direction of the proposition, let the glued topos be equivalent to $C^\wedge$. Its specified open subtopos corresponds under this equivalence to a subterminal presheaf. Such a presheaf takes the value either $\varnothing$ or a singleton at each object, and its support is a sieve $C_U\subset C$. The complementary full cosieve $C_T$ presents the closed complement. Restriction to $C_T$ preserves all limits, as does the right adjoint to restriction to $C_U$. Their composite is the transported gluing functor, so $\varphi$ preserves all limits. The argument uses the specified open and closed pieces but does not require the given presentations $U,T$ to coincide literally with $C_U,C_T$. Equivalences of presheaf categories preserve the limiting operations involved. This is precisely the content of the presheaf case of the SGA exercise just cited.

There is a variance convention in the word ``category'' in the corollary. For fixed $U,T$, a transformation $H\to K$ induces $\psi_H\to\psi_K$, but it induces a transformation $\varphi_K\to\varphi_H$ on the right adjoints. Thus the asserted equivalence of categories of gluing data uses reversed transformations on the right-adjoint side; with ordinary pointwise transformations it is an antiequivalence. The object classification and every displayed objectwise formula remain unchanged.

A useful test is $U=T=e$ and $H=S$ a set. The collage has two objects and one crossing arrow for each element of $S$. The functors are $\psi(G)=G\times S$ and $\varphi(F)=F^S$. A presheaf on the collage is a pair of sets $F,G$ with an $S$-indexed family of maps $G\to F$. If $S=\varnothing$, there is no compatibility datum and the presheaf topos is the product of two copies of $\Ens$. If $S$ is a singleton, it is the category of arrows of sets. These limiting cases also verify the direction $G\to\varphi(F)$ of the gluing arrow.
\end{dcomment}


\subsubsection*{1.3\quad Flèches dans $\mathrm{Cat}$, de source ou but un objet recollé $(X,U,T)$.}

Quand $(X,U,T)$ est donné par une donnée de recollement
\[
U,\ T,\ H : U^\circ \times T \longrightarrow \mathrm{Ens},
\]
les flèches
\[
f : X' \longrightarrow X
\]
dans Cat se décrivent immédiatement ainsi : la donnée de $f$ définit d'abord
\[
U' = f^{-1}(U), \qquad\qquad T' = f^{-1}(T),
\]
couple d'un ouvert de $X'$ et du fermé complémentaire, de sorte que cette donnée équivaut à celle de

\begin{enumerate}[label={}]
\item[$1^\circ$)] un couple $(U',T')$ dans $X'$
\item[$2^\circ$)] un morphisme (dans $\mathcal{T}$) de $(X',U',T')$ dans $(X,U,T)$, ce qui équivaut par 1.1 à un triple
\[
\left\{\begin{array}{rccl}
f_{U'} : & U' & \longrightarrow & U\\[1ex]
f_{T'} : & T' & \longrightarrow & T\\[1ex]
u : & H_{U',T'} & \longrightarrow & H \circ (f_{U'}^\circ \times f_{T'}),
\end{array}\right.
\]
où on a posé
\[
\left\{\begin{aligned}
&H_{U',T'} : U^{\prime\circ} \times T' \longrightarrow \mathrm{Ens}\\[1ex]
&H_{U',T'}(u',t') = \mathrm{Hom}_{X'}(u',t'),
\end{aligned}\right.
\]
i.e. $H_{U',T'}$ est la donnée de recollement qui conrrespond au triple $(X',U',T')$ dans $\mathcal{T}$. % [sic] conrrespond -> correspond, misspelling
\end{enumerate}

\msp{page 66}

Nous serons plus intéressés par la situation inverse, i.e. par la question de déterminer les foncteurs
\[
\underbrace{X}_{= X_{U,T,H}} \quad \longrightarrow X'
\]
de $X$, obtenu par $(U,T,H)$, dans un objet quelconque de Cat. On trouve tautologiquement

\textbf{Proposition 3.} $\mathrm{Hom}_{\mathrm{Cat}}(X,X')$ \textit{est en correspondance bijective canonique, fonctorielle en $X'$ dans} Cat\textit{, avec l'ensemble des triples}

% ---- from Der-ChXIIb-012_p006-010.tex
\printedpage{48}
% head: A. Grothendieck : Dérivateurs, Ch. XII, Édition des Universités de Montpellier II et Paris VII
% foot: Version 7 décembre 2002

\[
\left\{
\begin{aligned}
&\left.\begin{aligned} f_U &: U \longrightarrow X' \\ f_T &: T \longrightarrow X' \end{aligned}\right\}\text{flèches dans Cat}\\
&u : H \longrightarrow H_{X'}(f_U^{\circ}\times f_T),\ \text{flèche dans } \uHom(U^{\circ}\times T,\Ens),
\end{aligned}
\right.
\tag*{\texttt{[}$(*)$\texttt{]}}
\]
\textit{où pour tout $X'$ dans} Cat\textit{, on a posé}
\[
H_{X'} = ((x,y)\longmapsto \Hom_{X'}(x,y)) : X'^{\circ}\times X' \longrightarrow \Ens.
\]

Sûrement on aura même un isomorphisme de catégories
\[
\uHom(X,X') \simeq \text{catégorie des triples } (f_U,f_T,u) \text{ comme dans } (*).
\]

\subsection*{1.4\quad Construction de sommes amalgamées à partir d'une immersion ouverte où fermée.} % [sic] où fermée -> ou fermée, conjonction « ou » et non adverbe « où »

Pour une construction d'une sorte de somme amalgamée, en termes d'un triple $(X,U,T)$ $+\ \ldots$, \msp{page 67} qui soit symétrique en $U$, $T$, j'ai envie de regarder d'abord la situation ‘mixte’, d'un diagramme en traits pleins
\[
\begin{tikzcd}
U \arrow[r, hook, "i"] \arrow[d, "f_U"'] & X \arrow[d, dotted, "f"] & T \arrow[l, hook', "j"'] \arrow[d, "f_T"] \\
U' \arrow[r, hook, dotted, "i'"] & X' & T', \arrow[l, hook', dotted, "j'"']
\end{tikzcd}
\tag*{$(*)$}
\]
qu'on aimerait compléter ‘canoniquement’, si faire se peut, en un morphisme de triples (i.e. en une flèche dans $\mathcal{T}$). L'idée naturelle, c'est que le $X'$ qui s'y insère serait la limite inductive, dans Cat, du diagramme initial. Si $T'=T$, $f_T=\id_T$, cela signifie donc que $X'$ s'identifie à la somme amalgamée $X\amalg_U U'$, et inversement si $U'=U$, $f_U=\id_U$.

\dgrblock{dgrb511}{Functors out of a collage include the mixed naturality equation}
\begin{dcomment}{exposition}{p66}{Functors out of a collage include the mixed naturality equation}
The next construction will be recognized by mapping out of it. Proposition~3 is therefore the correct place to retain natural transformations, even though its first formulation only mentions sets of functors. A functor $\operatorname{Coll}(H)\to Y$ consists of $a:U\to Y$, $b:T\to Y$ and maps
\[
 \eta_{u,t}:H(u,t)\longrightarrow\Hom_Y(a(u),b(t))
\]
compatible with precomposition in $U$ and postcomposition in $T$. The argument constructing the collage proves both existence and uniqueness of the functor from these data.

A morphism from $(a,b,\eta)$ to $(\bar a,\bar b,\bar\eta)$ is a pair of natural transformations $r:a\Rightarrow\bar a$, $s:b\Rightarrow\bar b$ satisfying, for every $h\in H(u,t)$,
\[
 s_t\,\eta_{u,t}(h)=\bar\eta_{u,t}(h)\,r_u.
\]
The components $r_u$ and $s_t$ define a natural transformation on the collage because its naturality is already known on arrows within either part, while this equation is its naturality on a crossing arrow. Conversely restriction recovers the pair and the equation. Composition and identities are componentwise. This gives the claimed isomorphism of functor categories, rather than merely an equivalence on their objects.

For a functor into a collage the procedure is different in presentation but just as exact. Pulling back the open part specifies a sieve in the source, and its complement specifies the cosieve; these are forced by the functor. The residual data are then a morphism of profunctors. This is why Proposition~3 does not require the images of $a$ and $b$ to be disjoint: the target is arbitrary, and a functor out of the collage may identify objects from the two parts. The universal property of the forthcoming ordinary pushout must allow precisely these functors.
\end{dcomment}


Désignons par $X'$ la somme amalgamée \texttt{[plutôt la limite inductive]}, et explicitons, pour $X''$ quelconque dans Cat, l'ensemble des flèches $g$ de $X'$ dans $X''$, comme foncteur en $X''$. C'est l'ensemble des triples
\[
(g_X, g_{U'}, g_{T'}) \qquad
\left\{\begin{aligned} g_X &: X \longrightarrow X'' \\ g_{U'} &: U' \longrightarrow X'' \\ g_{T'} &: T' \longrightarrow X''\end{aligned}\right.
\qquad \text{avec} \quad
\begin{aligned} g_X\circ i &= g_{U'} f_U \\ g_X \circ j &= g_{T'} f_T. \end{aligned}
\tag*{$(*)$}
\]
\printedpage{49}
% head: A. Grothendieck : Dérivateurs, Ch. XII, Édition des Universités de Montpellier II et Paris VII
% foot: Version 7 décembre 2002
\[
\begin{tikzcd}
U \arrow[r, hook, "i"] \arrow[d, "f_U"'] & X \arrow[dd, "g_X"'] & T \arrow[l, hook', "j"'] \arrow[d, "f_T"] \\
U' \arrow[dr, "g_{U'}"'] & & T' \arrow[dl, "g_{T'}"] \\
& X'' &
\end{tikzcd}
\]

Notons que pour un tel triple, $g_X$ est partiellement connu déjà en termes de $g_{U'}$, $g_{T'}$, par ses restrictions $g_X\circ i$ et $g_X\circ j$ à $U$ et $T$ respectivement, par les formules
\[
g_U = g_{U'}f_U, \qquad g_T = g_{T'}f_T.
\]

\msp{page 68} Donc la donnée d'un triple $(*)$ équivaut à celle d'un triple
\[
(g_{U'}, g_{T'}, g_X) \qquad
\left\{\begin{aligned}
&g_X \text{ une flèche } X\longrightarrow X'' \text{ correspondant à}\\
&g_U = g_{U'}\circ f_U \text{ et } g_T = g_{T'}\circ f_T \text{ déjà donnés (via } g_{U'}\text{, } g_{T'}\text{),}
\end{aligned}\right.
\]
donc en vertu de 1.3. cela correspond à
\[
(g_{U'}, g_{T'}, u) \qquad
\left\{\begin{aligned}
g_{U'} &: U' \longrightarrow X''\\
g_{T'} &: T' \longrightarrow X''\\
u &: H \longrightarrow \underbrace{H_{X''}\left((g_{U'}\circ f_U)^{\circ}\times(g_{T'}\circ f_T)\right)}_{(u,t)\longmapsto \Hom_{X''}(g_{U'}f_U(u),g_{T'}f_T(t))}.
\end{aligned}\right.
\tag*{$(*)$}
\]
On va interpréter la donnée de $u$, en termes d'un foncteur
\[
F^{*} = (H'\longmapsto H'\circ F) : \uHom(U'^{\circ}\times T',\Ens) \longrightarrow \uHom(U^{\circ}\times T,\Ens)
\tag*{$(**)$}
\]
où
\[
F : U^{\circ}\times T \longrightarrow U'^{\circ}\times T' \qquad \text{est donné par } F = f_U^{\circ}\times f_T.
\]
On a un objet $H''=H''_{g_{U'},g_{T'}}$ de la catégorie source de $(**)$ déterminé en termes du couple $(g_{U'},g_{T'})$, par
\[
H''(u',t') = \Hom_{X''}(g_{U'}(u'),g_{T'}(t')),
\]
et le bifoncteur but de la flèche $u$ dans $(*)$ s'interprète comme
\[
F^{*}(H'') = H''\circ F : (u,t)\longmapsto \underbrace{H''(f_U(u),f_T(t))}_{=\Hom(g_{U'}f_U(u),g_{T'}f_T(t))} \quad .
\]
Ainsi, $u$ s'interprète comme une flèche
\[
u : H \longrightarrow F^{*}(H'').
\]
Or le foncteur $F^{*}$ a un adjoint à gauche $F_{!}$, i.e. on a
\[
\Hom(H,F^{*}(H'')) \simeq \Hom(F_{!}(H),H'').
\tag*{$(*{*}*)$}
\]

\printedpage{50}
% head: A. Grothendieck : Dérivateurs, Ch. XII, Édition des Universités de Montpellier II et Paris VII
% foot: Version 7 décembre 2002
\msp{page 69} avec
\[
F_{!} : \uHom(U^{\circ}\times T,\Ens) \longrightarrow \uHom(U'^{\circ}\times T',\Ens).
\]
Posons
\[
H' = F_{!}(H) : U'^{\circ}\times T' \longrightarrow \Ens,
\]
c'est un bifoncteur bien déterminée par la seule donnée du diagramme initial $(*)$ p. 67, ou ce qui revient au même, en termes de % [sic] déterminée -> déterminé, accord avec « bifoncteur » (masculin)
\[
(U,T,H,f_U : U\longrightarrow U', f_T : T\longrightarrow T').
\]

\dgrblock{dgrb512}{Transporting the crossing arrows by a Kan extension}
\begin{dcomment}{exposition}{p68}{Transporting the crossing arrows by a Kan extension}
The two categories $U,T$ can now be changed independently. Their new internal arrows are already present in $U',T'$. What must be supplied are all composites crossing from the first new category to the second, subject to the relations forced by old compositions. This is what the left Kan extension supplies.

Put $F=f_U^{\mathrm{op}}\times f_T$. The ordinary pointwise formula can be written in a form that displays every generator:
\[
 H'(u',t')=
 \int^{u\in U,\,t\in T}
 \Hom_{U'}(u',f_Uu)\times H(u,t)
                  \times\Hom_{T'}(f_Tt,t').
\]
The notation means the quotient of the disjoint union of these products by the relations for composition in $U$ and $T$. An element is represented by a triple
\[
 u'\xrightarrow{\alpha}f_Uu,\qquad
 h\in H(u,t),\qquad f_Tt\xrightarrow{\beta}t'.
\]
For $a:u_0\to u_1$ in $U$ and $b:t_0\to t_1$ in $T$, the defining relations are
\[
 [f_U(a)\alpha,h,\beta]=[\alpha,H(a,1)h,\beta],
 \qquad
 [\alpha,H(1,b)h,\beta]=[\alpha,h,\beta f_T(b)],
\]
whenever the displayed terms are typed. This is the coend formula for $\operatorname{Lan}_F H$; the equivalent comma-category formula is the pointwise Kan-extension theorem \citep[Theorem~6.2.1]{Riehl2016}.

To prove the pushout assertion directly, form $C=\operatorname{Coll}(H')$ and send a crossing arrow $h$ of $X$ to $[1,h,1]$. The relations make this a functor $X\to C$, and the inclusions of $U',T'$ agree with it on $U,T$. Now take compatible functors $X\to Y$, $U'\to Y$, $T'\to Y$. On the generator above there is only one possible extension: it must take the value $\beta\,h\,\alpha$, after applying the respective functors. The two relations hold because the original functors preserve composition. Thus there is a unique functor $C\to Y$ inducing all three given functors. This proves the strict colimit universal property, and also identifies its canonical map from $X$.

Since $C$ was constructed as a collage, $U'$ is a full sieve and $T'$ its full complementary cosieve. Their internal hom sets have not been enlarged by the quotient: no string can leave the open part and subsequently return to it. This supplies the structural content of Theorem~1 beyond existence of a colimit in $\Cat$. It also clarifies the distinction from the earlier induced triple: restriction of $H$ gave a cartesian lift, whereas this left Kan extension gives the universal forward transport. The labels ``pullback'' and ``pushout'' refer to these different universal properties, not to a formal substitution of primes in the same formula.
\end{dcomment}


Et on voit que la donnée d'une flèche $X'\longrightarrow X''$, $X''$ objet variable de Cat, équivaut à celle de
\[
g_{U'} : U'\longrightarrow X'',\ g_{T'} : T'\longrightarrow X'',\ u : H'\longrightarrow H'' = H_{X''}(g_{U'}^{\circ}\times g_{T'}),
\]
c'est à dire, en vertu de 1.3., à celle d'une flèche
\[
X'_0 \longrightarrow X'',
\]
où $X'_0$ désigne l'objet de Cat déduit de $U'$, $T'$ par recollement au moyen de $H'$. Cela montre donc que la limite inductive cherchée n'est autre que $X'_0$. On a ainsi trouvé le

\textbf{Théorème 1.} \textit{La limite inductive $X'$ dans} Cat \textit{du diagramme} $(*)$ \textit{page 67, défini en termes d'un triple $(X,U,T)$ dans $\mathcal{T}$ et de flèches $f_U : U\longrightarrow U'$, $f_T : T\longrightarrow T'$ dans} Cat\textit{, existe (sans mérite)} \msp{page 70} \textit{et les flèches canoniques}
\[
U' \overset{i'}{\hookrightarrow} X' \overset{j'}{\hookleftarrow} T'
\]
\textit{sont des immersions ouvertes et fermées complémentaires l'une de l'autre. La donnée de recollement qui définit $X'$ via $U'$, $T'$ est donnée par}
\[
H' = (\underbrace{f_U^{\circ}\times f_T}_{F})_{!}(H) : U'^{\circ}\times T' \longrightarrow \Ens,
\]
\textit{où $H$ est la donnée de recollement qui définit $X$, i.e. $H : U^{\circ}\times T\longrightarrow$} Ens \textit{est le bifoncteur $(u,t)\longmapsto \Hom_X(u,t)$.}

Ainsi on a, pour $u'$ dans $U'$ et $t'$ dans $T'$
\[
\Hom_{X'}(u',t') = F_{!}(H)(u',t'),
\]
et il reste à expliciter $F_{!}(H)$. Ici $F=f_U^{\circ}\times f_T$ est un foncteur covariant sur $U^{\circ}\times T$, on peut l'interpréter comme un préfaisceau sur $U\times T^{\circ}$, et écrire dès lors
\[
\underbrace{F_{!}(H)(u',t')}_{\Hom_{X'}(u',t')} = \varinjlim_{(u',t')\backslash U\times T^{\circ}} \underbrace{\Hom_X(u,t)}_{H(u,t)} \ .
\tag*{$(*)$}
\]

\printedpage{51}
% head: A. Grothendieck : Dérivateurs, Ch. XII, Édition des Universités de Montpellier II et Paris VII
% foot: Version 7 décembre 2002
La catégorie $(u',t')\backslash U\times T^{\circ}$ sur laquelle on prend la $\varinjlim$ s'explicite comme un produit
\[
(u',t')\backslash U\times T^{\circ} \simeq (u'\backslash U)\times(T/t')^{\circ}.
\]
Les applications
\[
\Hom_X(u,t)\longrightarrow \Hom_{X'}(u',t') \qquad (=F_{!}(H)(u',t'))
\]
\msp{page 71} déduites de cette formule ne sont autres que les suivantes : $u$, $t$ sont données comme objets de $u'\backslash U$, $T/t'$, i.e. avec des flèches
\[
u'\xrightarrow{\alpha} f_U(u), \qquad f_T(t)\xrightarrow{\beta} t',
\]
d'où
\[
\Hom_X(u,t) \xrightarrow{\text{via } f=f_X:X\to X'} \Hom_{X'}(\underbrace{f(u)}_{f_U(u)},\underbrace{f(t)}_{f_T(t)}) \xrightarrow{\text{via } \alpha,\ \beta} \Hom_{X'}(u',t').
\]

Notons pour mémoire que pour tout objet $Z$ de Cat, et tout préfaisceau $G$ sur $Z$, on a la formule
\[
\varinjlim_{Z} G \simeq \pi_0(Z/G),
\]
où $Z/G$ s'interprète en termes de l'immersion canonique $Z\longrightarrow Z^{\wedge}$ $(G\in \Ob Z^{\wedge})$. C'est donc une catégorie fibrée sur $Z$ à fibres discrètes, et on a en fait un isomorphisme entre la sous-catégorie pleine de $\Cat/Z$ formée de telles catégories fibrées à fibres discrètes, avec $Z^{\wedge}$. Dans ce dictionnaire parfait entre préfaisceaux d'ensembles et catégories fibrées à fibres discrètes, le foncteur $\varinjlim$ sur $Z^{\wedge}$ correspond simplement au foncteur $\pi_0$ sur la catégorie des objets au dessus de $Z$ dans Cat, lui-même induit par le foncteur (indépendant de $Z$)
\[
\Cat \xrightarrow{\pi_0} \Ens.
\]

\msp{page 72} On peut donc réécrire la formule fondamentale $(*)$ p. 70, en termes du foncteur $H$ définissant le recollement $X$ de $U$ et de $T$, par
\[
\boxed{\underbrace{\Hom_{X'}(u',t')}_{H'(u',t')} \simeq \pi_0(\underbrace{(u',t')\backslash U\times T^{\circ}/H}_{((u'\backslash U)\times(T/t')^{\circ})/H}).}
\tag*{$(*)$}
\]

Par la suite, nous n'aurons sans doute qu'à utiliser le cas où on a soit $T'=T$, $f_T=\id_T$, soit l'inverse. Nous travaillerons surtout dans le premier cas, en déduisant ce dont nous aurons besoin pour le cas dual, par principe de dualité (automorphisme $X\longmapsto X^{\circ}$ de la
\printedpage{52}
% head: A. Grothendieck : Dérivateurs, Ch. XII, Édition des Universités de Montpellier II et Paris VII
% foot: Version 7 décembre 2002
catégorie Cat). Dans ce cas, dans la formule fondamentale $(*)$ l'élément $t'$ de $T'$ s'écrit comme un élément $t$ de $T$, et la formule se retranscrit comme
\[
\boxed{\Hom_{X'}(u',t) \simeq \pi_0(\underbrace{u'\backslash U/t}_{\text{catégorie des triples } (u,\alpha,\beta'),\ \left\{\begin{aligned} &u\in \Ob U\\ &\alpha : u\longrightarrow t\\ &\beta' : u'\longrightarrow f_U(u)\end{aligned}\right.}).}
\tag*{$(**)$}
\]

C'est la formule que j'avais trouvée directement dans ce cas, sans me soucier encore de faire varier $T$ en même temps que $U$. Je voudrais faire le lien entre les deux formules $(*)$ et $(**)$, par voie topologique, donc établir une bijection canonique, fonctorielle etc.
\[
\pi_0(u'\backslash U/t) \xrightarrow{\sim} \pi_0(((u'\backslash U)\times(T/t)^{\circ})/H).
\]

\msp{page 73} On a sûrement un foncteur
\[
u'\backslash U/t \longrightarrow ((u'\backslash U)\times(T/t)^{\circ})/H
\]
qui associe à un triple
\[
(u,\alpha,\beta') : u\in \Ob U,\ \alpha : u\longrightarrow t,\ \beta' : u'\longrightarrow f_U(u)
\]
le triple $(\bar u,\bar t,\xi)$ $(\xi\in H(\bar u,\bar t))$ dans le second membre, avec $\bar u$ dans $u'\backslash U$, $\bar t$ dans $T/t$, défini ainsi :
\[
\bar u = (u,\beta'),\ \bar t = (t,\id_t), \underbrace{\xi}_{=\alpha} \in H(\bar u,\bar t) = \Hom(u,t).
\]
En somme
\[
u'\backslash U/t = \text{fibre de } ((u'\backslash U)\times(T/t)^{\circ})/H \text{ en l'objet initial } \bar t \text{ de } (T/t)^{\circ}
\]
par la projection naturelle, flèche composée
\[
\begin{tikzcd}
((u'\backslash U)\times(T/t)^{\circ})/H \arrow[d] \\
(u'\backslash U)\times(T/t)^{\circ} \arrow[d, "\pi"] \\
(T/t)^{\circ}.
\end{tikzcd}
\]

\dgrblock{dgrb513}{Why connected components give the ordinary hom set}
\begin{dcomment}{exposition}{p72}{Why connected components give the ordinary hom set}
Grothendieck wants a topological description of a purely ordinary quotient. For a presheaf $G:Z^{\mathrm{op}}\to\Ens$, the category of elements $Z/G$ has objects $(z,g)$ and arrows
\[
 (z,g)\longrightarrow(z',g')\quad\text{given by}\quad
 a:z\longrightarrow z',\qquad G(a)(g')=g.
\]
The colimit of $G$ is the disjoint union of its values modulo the identifications $G(a)(g')\sim g'$. These are exactly the identifications imposed on objects when one passes to connected components of $Z/G$. Hence
\[
 \mathop{\mathrm{colim}}_{Z^{\mathrm{op}}}G\simeq\pi_0(Z/G).
\]
The source writes $\varinjlim_Z G$ for this presheaf colimit. The opposite belongs in the modern indexing category; it is not a correction to the numerical value of the formula.

The same convention explains the indexing category in the preceding page. If $H$ is regarded as a covariant functor on $U^{\mathrm{op}}\times T$, its ordinary left Kan extension is indexed by
\[
 F/(u',t')\simeq (u'\backslash f_U)^{\mathrm{op}}\times(f_T/t').
\]
Equivalently, regard $H$ as a presheaf on $U\times T^{\mathrm{op}}$ and take its presheaf colimit over the opposite category
\[
 (u'\backslash f_U)\times(f_T/t')^{\mathrm{op}}.
\]
This is the convention used in the edition. Its sentence ``Ici $F=\cdots$ ... on peut l'interpr\'eter comme un pr\'efaisceau'' contains a harmless but real letter slip: it is $H$, not the indexing functor $F$, that is the set-valued presheaf.

When $T'=T$, the resulting objects can be compressed to factorizations
\[
 (u,\alpha:u'\to p(u),\gamma:u\to t).
\]
An arrow from $(u,\alpha,\gamma)$ to $(v,\alpha',\gamma')$ is $a:u\to v$ with $p(a)\alpha=\alpha'$ and $\gamma'a=\gamma$. Its connected component is the crossing arrow of the pushout represented by $\gamma\alpha$. Thus the factorization category keeps track of all presentations of a new arrow and all relations between those presentations. Taking $\pi_0$ intentionally forgets its higher homotopy. The later asphericity conditions ask precisely when this forgetting is harmless for the homotopy problem, and do not follow from the existence of the ordinary colimit alone.
\end{dcomment}


Il faut voir que l'inclusion de la fibre donne une bijection sur les $\pi_0$, ou ce qui revient au même, que l'inclusion de la fibre $u'\backslash U$ \textit{de} $\pi$ en l'objet initial $(t,\id_t)$ de $(T/t)^{\circ}$ donne une bijection sur les $\varinjlim$ du foncteur $H$. Donc l'inclusion a intérêt à être asphérique, ou du moins co-cofinalle. % [sic] co-cofinalle -> co-cofinale, orthographe

% ---- from Der-ChXIIb-012_p011-015.tex
% ===================== page 11 =====================
\printedpage{53}
% head: A. Grothendieck : Dérivateurs, Ch. XII, Édition des Universités de Montpellier II et Paris VII

\msp{page 74}
Mais l'inclusion d'un objet initial est toujours asphérique, et le reste par produit avec un objet quelconque (i.e. $u'\backslash U$) de Cat (et même par tout changement de base lisse \dots), d'où la conclusion.

Dans la suite, nous nous intéressons surtout à la situation de sommes amalgamées (où on change de notation un peu, à présent)
\[
\begin{tikzcd}
U \arrow[r, hook, "i"] \arrow[d, "p"'] & X \arrow[d, "q"'] & T \arrow[l, hook', dotted, "j"'] \arrow[d, dotted] \\
U' \arrow[r, hook, "i'"] & X' & T' \simeq T. \arrow[l, hook', dotted]
\end{tikzcd}
\tag*{$(*)$}
\]

Le fermé $T'$ complémentaire de $U'$ dans $X'$ est alors canoniquement isomorphe au fermé $T$ complémentaire de $U$ dans $X$, l'isomorphisme étant induit par $q$. Souvent nous identifierons $T'$ et $T$, ce qui est licite en regardant $X'$ comme déduit du recollement de $U'$ et de $t$ \texttt{[plutôt }$T$\texttt{]}, par le bifoncteur $U'^{\circ}\times T\longrightarrow \mathrm{Ens}$ explicité dans la formule $(**)$ p.~72. Nous nous intéressons surtout à dégager des conditions qui assurent les deux implications suivantes

\begin{itemize}
\item[a)] $i\in W\Longrightarrow i'\in W$ (et même $\Longleftrightarrow$).
\item[b)] $p\in W\Longrightarrow q\in W$ (et même $\Longleftrightarrow$).
\end{itemize}

\msp{page 75}
\textbf{Proposition 4.} \textit{Soit $t\in\Ob T$. Pour que $t\in\Ob\bar{U}$, il faut et il suffit que $q(t)\in\Ob\overline{U'}$. En d'autres termes, le `bord' $\dot{T}=T\cap\bar{U}$ de $T$ dans $X$ est `le même' que son bord dans $X'$.}

\dgrblock{dgrb514}{The one-sided formula has an explicit deformation retraction}
\begin{dcomment}{exposition}{p74}{The one-sided formula has an explicit deformation retraction}
The comparison of the two factorization formulas admits a direct argument which also verifies Grothendieck's cofinality claim with its variance intact. Write $E$ for the category of elements occurring in the general formula when $T'=T$. Its objects are
\[
 (u,\alpha:u'\to p(u),v,\beta:v\to t,h:u\to v).
\]
An arrow to $(u_1,\alpha_1,v_1,\beta_1,h_1)$ consists of
\[
 a:u\to u_1,\quad b:v_1\to v,\qquad
 p(a)\alpha=\alpha_1,\quad \beta b=\beta_1,\quad
 h=b h_1a.
\]
Let $D=u'\backslash U/t$ be the category in the boxed one-sided formula. Define
\[
 R:E\longrightarrow D,\qquad
 R(u,\alpha,v,\beta,h)=(u,\alpha,\beta h).
\]
On arrows $R$ takes $(a,b)$ to $a$; the equality
$\beta_1h_1a=\beta b h_1a=\beta h$ proves that this is an arrow of $D$. The inclusion $J:D\to E$ takes $(u,\alpha,\gamma)$ to $(u,\alpha,t,1_t,\gamma)$, and $RJ=1_D$.

For each object of $E$, the pair $(1_u,\beta)$ gives an arrow
\[
 JR(u,\alpha,v,\beta,h)
       \longrightarrow(u,\alpha,v,\beta,h).
\]
The equations above verify naturality, so $JR\Rightarrow1_E$. Consequently $J$ and $R$ induce inverse homotopy equivalences on nerves, and in particular inverse bijections on $\pi_0$. This proves more than the displayed comparison requires. The construction is natural in $u'$ and $t$, because precomposition of $\alpha$ and postcomposition of $\beta$ commute with $\beta h$.

The reference to the initial object $(t,1_t)$ of $(T/t)^{\mathrm{op}}$ is therefore sound. An initial object is cofinal for the colimit of a \emph{presheaf}; in covariant colimit notation one has passed to the opposite and it becomes terminal. The proof does not assert that an initial object computes arbitrary covariant colimits. This is the same variance discipline needed in the commentary on Chapter~I, manuscript p.~47 (``Cofinality retains an orientation''), but here the explicit retraction identifies the whole category of presentations rather than only invoking the abstract cofinality criterion.
\end{dcomment}


\textsc{Démonstration}. Dire que $t\in\Ob\bar{U}$ signifie aussi
\[
U/t\neq\emptyset
\]
et implique, par continuité, que $q(t)\in\overline{q(U)}\subseteq\overline{U'}$. Inversement, si $q(t)\in\overline{U'}$, i.e. s'il existe $u'\in U'$ tel que
\[
\Hom(u',q(t))(=\pi_0(u'\backslash U/t))\neq\emptyset,
\]
il en résulte que $u'\backslash U/t\neq\emptyset$, a fortiori $U/t$ (qui reçoit $u'\backslash U/t)\neq\emptyset$, q.e.d.

\textbf{Proposition 5.} \textit{Soit}
\[
\begin{tikzcd}
Y \arrow[r, "i"] \arrow[d, "p"'] & X \arrow[d, "q"] \\
Y' \arrow[r, "i'"] & X'
\end{tikzcd}
\]
% foot: Version 7 décembre 2002
% ===================== page 12 =====================
\printedpage{54}
% head: A. Grothendieck : Dérivateurs, Ch. XII, Édition des Universités de Montpellier II et Paris VII
\textit{un carré cocartésien dans \textup{Cat}. Si $i$ est une immersion ouverte (resp. fermée), alors $i'$ aussi, et le carré est aussi cartésien. Idem pour $p$ et $q$. Si $i$ et $p$ sont des immersions ouvertes (\footnote{\textbf{NB} Nous utilisons \textit{`immersion ouverte'}, \textit{`immersion fermée'} au \textit{sens strict}\,: c'est un \textit{monomorphisme} pleinement fidèle dont l'image est une sous-catégorie \textit{pleine,} soit ouverte, soit fermée.}) (resp. \dots), alors la situation est isomorphe à celle d'un objet $X'$ de \textup{Cat}, de deux ouverts (resp. fermés) $X$, $Y'$ de $X$ ayant pour réunion $X'$, pour intersection} % [sic] de $X$ -> de $X'$, X and Y' are opens of X' (their union is X')
\textit{$Y$. Inversement, dans une telle situationm le carré correspondant est cocartésien, i.e. $X$} % [sic] situationm -> situation, (comma), typing error
\texttt{[plutôt }$X'$\texttt{]} \textit{s'identifie à la somme amalgamée}
\[
X'\simeq X\amalg_{Y}Y'.
\]

Donc aussi on a la proposition 15, page 158, et ses corollaires.
% foot: Version 7 décembre 2002

% ===================== page 13 =====================
\printedpage{55}
% head: A. Grothendieck : Dérivateurs, Ch. XII, Édition des Universités de Montpellier II et Paris VII

\section*{2\quad Sommes amalgamées et $W$-équivalences.}

\subsection*{2.1\quad Axiome du carré cocartésien et axiome du cube (ou `Lemme des cinq').}

Nous supposons à présent donné un localiseur fondamental
\[
W\subseteq\mathrm{Fl}(\mathrm{Cat}),
\]
donc $W$ satisfait à W(1, 2 bis, 3). Nous ne supposons pas pour le moment pour les autres axiomes standard W(4) (axiome de fibration), W(5) (axiome de connexité), W(6) (axiome des limites inductives). Mais nous allons introduire à présent deux axiomes de stabilité supplémentaire, de nature différente, concernant les sommes amalgamées.

\begin{description}[leftmargin=5em, style=nextline]
\item[W(7 bis)] (Axiome du carré cocartésien.) Soient $X$ un objet de Cat, $U$, $U'$ deux ouverts de $X$ tels que $U\cup U'=X$, considérons le diagramme cocartésien correspondant (cf. prop. 5 en marge de la page \texttt{[ci-dessus]}) :

\msp{page 76}
\[
\begin{tikzcd}[column sep=small]
& U\cap U' \arrow[dl, hook', "i'_0"'] \arrow[dr, hook, "i_0"] & \\
U \arrow[dr, hook, "i"'] & & U' \arrow[dl, hook', "i'"] \\
& X &
\end{tikzcd}
\tag*{$(*)$}
\]

\dgrblock{dgrb515}{Orbits of crossing arrows and the information discarded}
\begin{dcomment}{examples}{p75}{Orbits of crossing arrows and the information discarded}
The pushout formula can be seen in a small calculation before it is used homotopically. Let $U$ be the discrete category on $u_0,u_1$, let $T=e$ with object $t$, and let $H(u_i,t)$ be a singleton. Thus $X$ is the span $u_0\to t\leftarrow u_1$. Collapse $U$ to $U'=e$ by $p$. The category $u'\backslash U/t$ has exactly two objects and no nonidentity arrows, so the pushout $X'=X\amalg_Ue$ has two objects and \emph{two} distinct parallel arrows $u'\rightrightarrows t$. Identifying the two objects of $U$ does not identify arrows which have no relation between them. Its nerve realizes to a circle, whereas the nerve of $X$ is contractible. Here $p$ is not a weak homotopy equivalence, so this is an illustration of the formula, not a counterexample to any asserted stability under a weak equivalence.

A group action makes the quotient and its lost homotopy particularly explicit. Put $U=BG$, let $T=e$, and take a right $G$-set $S$ as the presheaf $H(-,t)$. On pushing $BG$ to $e$, the mixed hom set becomes
\[
 \Hom_{X'}(u',t)=S/G.
\]
Indeed the category of elements has objects $s\in S$ and an arrow labelled $g$ from $s$ to $s'$ when $s=s'g$. All arrows are invertible, and two objects lie in the same connected component exactly when they lie in the same orbit. If $S=K\backslash G$ is a transitive right action, the automorphism group of the object $K$ is $K$. Inclusion of that object is fully faithful and essentially surjective, hence the category of elements is equivalent to $BK$.

Thus a single crossing arrow in the quotient may stand for a factorization category with homotopy type $BK$. Choosing $K$ nontrivial shows that a singleton ordinary colimit says nothing about the contractibility of its category of presentations. If $S=G$ with the free right action, the stabilizer is trivial and the category of elements is contractible. If $S$ is the singleton with trivial action, it is $BG$. Both have the same orbit set. The distinction is exactly the one that motivates the total asphericity conditions later in this section.
\end{dcomment}


Alors
\[
i_0\in W\qquad\Longleftrightarrow\qquad i\in W
\]
(et symétriquement
\[
i'_0\in W\qquad\Longleftrightarrow\qquad i'\in W).
\]
\end{description}

\textbf{NB} Par dualité, cet axiome est équivalent à l'axiome similaire, avec des \textit{fermés} $T$, $T'$ de $X$ de réunion $X$.

\textbf{Corollaire.} \textit{Supposons \textup{W(7 bis)}. Considérons les conditions suivantes sur le carré précédent.}

\begin{itemize}
\item[\textit{a)}] \textit{$U\cap U'$ $W$-asphérique.}
\item[\textit{b)}] \textit{$U$, $U'$ $W$-asphériques.}
\item[\textit{c)}] \textit{$X$ $W$-asphérique.}
\end{itemize}

\textit{Alors (a et b) implique c, et (b et c) implique a.}

\textbf{NB} Moralement, il faudrait aussi que (c et a) implique b, i.e. si $U\cap U'$ et $X$ sont $W$-asphériques, il en est de même de $U$ et de $U'$. Mais il ne semble pas que l'axiome du carré cocartésien permette de le prouver. Par contre, l'axiome qui suit impliquera le corollaire précédent sous sa forme complète.
% foot: Version 7 décembre 2002

% ===================== page 14 =====================
\printedpage{56}
% head: A. Grothendieck : Dérivateurs, Ch. XII, Édition des Universités de Montpellier II et Paris VII

\begin{description}[leftmargin=5em, style=sameline]
\item[W(7)] (`Axiome du cube', ou `Lemme des cinq'.) Soit un diagramme commutatif cubique
\[
\begin{tikzcd}[row sep=small, column sep=small]
U_1\cap U'_1 \arrow[rr] \arrow[dr] \arrow[dd, "f_{U_1\cap U'_1}"'] & & U_1 \arrow[dr] \arrow[dd, "f_{U_1}"{swap,pos=.25}] & \\
& U'_1 \arrow[rr, crossing over] & & X_1 \arrow[dd, "f"] \\
U_2\cap U'_2 \arrow[rr] \arrow[dr] & & U_2 \arrow[dr] & \\
& U'_2 \arrow[rr] \arrow[from=uu, crossing over, "f_{U'_1}"{pos=.25}] & & X_2
\end{tikzcd}
\]
représentant un morphisme (flèches verticales) d'un

\msp{page 77}
carré cocartésien\quad $\vcenter{\hbox{\begin{tikzcd}
U_1\cap U'_1 \arrow[r] \arrow[d] & U_1 \arrow[d] \\
U'_1 \arrow[r] & X_1
\end{tikzcd}}}$\quad du type envisagé dans W(7 bis) ($X_1$ réunion de deux ouverts $U_1$, $U'_1$) dans un autre (à indices 2). Considérons les conditions suivantes.

\dgrblock{dgrb516}{Ordinary covers yield homotopy pushouts, with one direction of excision}
\begin{dcomment}{exposition}{p76}{Ordinary covers yield homotopy pushouts, with one direction of excision}
Grothendieck's new axioms try to turn a cover by two categorical opens into an excision principle. The underlying geometric fact is especially clean. Suppose $X=U\cup V$ with $U,V$ full sieves and $A=U\cap V$. A simplex of $NX$ is a composable chain $x_0\to\cdots\to x_n$. Its last vertex lies in $U$ or in $V$; since these are sieves, all earlier vertices then lie in the same one. Hence
\[
 NX=NU\cup_{NA}NV
\]
is an actual pushout of simplicial sets. The two maps out of $NA$ are monomorphisms. On geometric realization they are inclusions of CW subcomplexes, so this pushout also computes the homotopy pushout. For cosieves the same proof uses the first vertex of the chain.

This proves the direct implication of W(7 bis) for $W_\infty$: if $A\to V$ is a weak equivalence then $U\to X$ is a weak equivalence. It also proves the direct cube assertion $(a\text{ and }b)\Rightarrow c$: an objectwise weak equivalence of the two spans induces a weak equivalence of their homotopy pushouts. This invariance is the property of a derived colimit. It may be checked using the double mapping-cylinder model, or by mapping into a fibrant space, where the universal mapping space is the homotopy pullback of the three mapping spaces of the span.

The reverse implication requested in W(7 bis), and the two additional implications requested in W(7), are strictly stronger assertions. The formal corollaries drawn here from the proposed axioms are valid as conditional statements, as is Proposition~6's implication W(7)$\Rightarrow$W(7 bis). However, Theorem~2 on manuscript page~78 is false for the full printed axiom. A homotopy pushout can destroy the fundamental group of its intersection. The next block gives a finite categorical counterexample, with the ordinary and homotopy pushouts agreeing throughout. This is not a failure of the nerve to preserve the relevant pushout, and is not repaired by using a derived pushout instead.
\end{dcomment}


\begin{itemize}
\item[a)] $f_{U_1\cap U'_1}\in W$.
\item[b)] $f_{U_1}$ et $f_{U'_1}$ sont \texttt{[dans]} $W$.
\item[c)] $f\in W$.
\end{itemize}

Alors deux parmi ces conditions impliquent la troisième. En d'autres termes, quand l'une de ces conditions est satisfaite, les deux autres sont équivalentes entre elles.
\end{description}

\textbf{NB} Cet axiome est équivalent à l'axiome dual, en termes d'objets $X_1$, $X_2$ représentés chacun comme réunion de deux fermés.

\textbf{Corollaire.} \textit{Supposons \textup{W(7)} satisfait. Alors le corollaire à \textup{W(7 bis)} (axiome qu'on n'a pas à utiliser ici), sous sa forme complétée (cf. le NB qui le suit) est valable : deux parmi les conditions a), b), c) impliquent la troisième.}

Il suffit de regarder l'unique morphisme du diagramme $(*)$, p.~76, dans le diagramme trivial\quad $\vcenter{\hbox{\begin{tikzcd}[column sep=small]
& e \arrow[dl] \arrow[dr] & \\
e \arrow[dr] & & e, \arrow[dl] \\
& e &
\end{tikzcd}}}$\quad correspondant au cas où $X$ est la catégorie finale $e$, et $U$, $U'$ égaux à $X$. Alors le corollaire n'est autre que W(7) appliqué à ce cas particulier.

Cet argument fait sentir que W(7) est sans doute plus fort que W(7 bis). Incité par ce sentiment, je m'aperçois avec plaisir qu'il est fondé :
% foot: Version 7 décembre 2002

% ===================== page 15 =====================
\printedpage{57}
% head: A. Grothendieck : Dérivateurs, Ch. XII, Édition des Universités de Montpellier II et Paris VII

\msp{page 78}
\textbf{Proposition 6.} \textit{L'axiome \textup{W(7)} (lemme des cinq) implique \textup{W(7 bis)} (axiome du carré cocartésien).}

On regarde le diagramme cubique suivant
\[
\begin{tikzcd}[row sep=small, column sep=small]
U\cap U' \arrow[rr, "i_0"] \arrow[dr, "i'_0"] \arrow[dd, "f_{U\cap U'}=i_0"'] & & U' \arrow[dr, "i'"] \arrow[dd, "f_{U'}=\mathrm{id}"{swap,pos=.25}] & \\
& U \arrow[rr, crossing over, "i"] & & X \arrow[dd, "f=\mathrm{id}_X"] \\
U' \arrow[rr, "\mathrm{id}"] \arrow[dr, "i'"] & & U' \arrow[dr, "i'"] & \\
& X \arrow[rr, "\mathrm{id}"] \arrow[from=uu, crossing over, "f_U=i"{pos=.25}] & & X,
\end{tikzcd}
\]
du carré cocartésien donné dans le carré cocartésien déduit de $X$ et de ses deux ouverts $X$, $U'$, d'intersection $U'$. Ici les deux flèches verticales $f$ et $f_{U'}$ sont des isomorphismes, donc \texttt{[dans]} $W$. Avec les notations de W(7), on a donc c) ($f\in W$), donc a) et b) sont équivalentes, i.e. $i_0\in W$ et $i\in W$ sont équivalentes.

\hfill q.e.d.

\textbf{Théorème 2.} \textit{Le localiseur fondamental $W_\infty$ satisfait à l'axiome \textup{W(7)}.}

C'est sûrement une conséquence du critère cohomologique d'\textsc{Artin-Mazur}, compte tenu de la suite exacte de Mayer-Vietoris pour l'un et l'autre des deux carrés cocartésiens, et les variantes non commutatives de Mayer-Vietoris. Je n'ai pas écrit la vérification en forme, et devrais sans doute réparer cette omission.

\dgrblock{dgrb517}{A finite counterexample to the asserted reverse excision}
\begin{dcomment}{development}{p78:Je n'ai pas écrit la vérification en forme}{A finite counterexample to the asserted reverse excision}
Grothendieck hopes that a five-lemma argument will recover the homotopy type of an intersection from the two pieces and their union. The obstruction can be isolated using the familiar acyclic noncontractible complex of \citet[Example~2.38]{Hatcher2002}. We give the categorical consequence in full; the input space and its properties are classical. We checked that example, the finite simplicial replacement theorem, the subsequent distinction between the direct and reverse axioms in Chapter~XVI, and the model-independent homotopy-pushout argument. We did not find this exact finite-poset formulation attached to Theorem~2 in the sources checked; no priority is claimed.

\textbf{Proposition (finite failure of W(7 bis) and W(7)).}\ours
\emph{There is a finite poset $X$, covered by two full sieves $U,V$, such that $NU,NV,NX$ are contractible but $N(U\cap V)$ is not weakly contractible. Consequently $W_\infty$ does not satisfy either printed axiom.}

\emph{Proof.} Start with a connected CW complex $K_0$ with one vertex, two oriented edges $a,b$, and two two-cells attached along
\[
 a^5b^{-3},\qquad b^3(ab)^{-2}.
\]
Its cellular differential from dimension two to dimension one is
\[
 d_2=\begin{pmatrix}5&-2\\-3&1\end{pmatrix}.
\]
Its determinant is $-1$, so $d_2$ is an isomorphism. The differential $d_1$ is zero and there are no higher cells. It follows that $\widetilde H_*(K_0;\mathbb Z)=0$.

Here is an explicit verification that its fundamental group is nontrivial. With permutations composed from right to left, send
\[
 a\longmapsto(1\ 4\ 3\ 2\ 5),\qquad
 b\longmapsto(1\ 2\ 3).
\]
The product $ab$ is $(1\ 5)(3\ 4)$, so the images satisfy $a^5=b^3=(ab)^2=1$ and define a homomorphism from the presented group to $A_5$. The image of $a$ is nonidentity, proving nontriviality. We need neither a computation of the full group nor asphericity of the presentation complex.

By the finite case of \citet[Theorem~2C.5]{Hatcher2002}, choose a finite simplicial complex $K$ homotopy equivalent to $K_0$. Let $A$ be the poset of its nonempty simplices ordered by inclusion. Its nerve is the barycentric subdivision of $K$, so $|NA|$ is connected and acyclic, and has nontrivial fundamental group. Adjoin two incomparable elements $u,v$ above every element of $A$, and call the resulting finite poset $X$. Put
\[
 U=A\cup\{u\},\qquad V=A\cup\{v\}.
\]
Both are downsets of $X$, hence full sieves, and $U\cap V=A$. They have terminal objects $u$ and $v$, respectively. Their nerves are the two cones on $NA$. A chain in $X$ cannot contain both $u$ and $v$, and therefore $|NX|$ is the unreduced suspension of $|NA|$.

The suspension of a connected space is simply connected, by van Kampen applied to its two thickened cones. Mayer--Vietoris gives
\[
 \widetilde H_n(|NX|;\mathbb Z)
       \cong\widetilde H_{n-1}(|NA|;\mathbb Z)=0.
\]
Thus $|NX|$ is a simply connected acyclic CW complex. It is contractible by the homological Whitehead theorem \citep[Corollary~4.33]{Hatcher2002}. The inclusion $U\to X$ is consequently a weak homotopy equivalence. The inclusion $A\to V$ is not, since its source has nontrivial fundamental group and its target is contractible. This contradicts the reverse implication of W(7 bis). The map from the cover square of $X$ to the constant square on $e$ satisfies conditions $b,c$ of W(7) but not $a$. It therefore contradicts W(7) as well.\hfill$\square$

The proposed cohomological proof fails at a specific point. Local systems on the intersection need not extend to either piece or to the union. In the example, the nontrivial representation of $\pi_1(NA)$ in $A_5$ cannot extend across either cone. Mayer--Vietoris for coefficients restricted from the ambient contractible space only sees constant coefficients on the intersection, for which its acyclicity suppresses the obstruction. The nonabelian comparison required by a weak-equivalence criterion cannot be recovered from that restricted coefficient class. This is the same general difficulty as in the commentary on Chapter~I, manuscript p.~81 (``The coefficient lost in the attempted proof''), now realized in a finite categorical cover.
\end{dcomment}


\msp{page 79}
\textbf{Remarque.} La version duale (au sens de \textsc{Quillen}) de l'axiome W(7 bis) du carré cocartésien, i.e. l'axiome correspondant pour les carrés cartésiens, n'est pas valable pour $W_\infty$. L'énoncé pertinent serait le suivant : soit
\[
\begin{tikzcd}
X & X' \arrow[l, "q"'] \arrow[d, "f'"] \\
S & S' \arrow[l, "p"']
\arrow[from=1-1, to=2-1, "f"']
\end{tikzcd}
\]
un carré cartésien, avec $f$ une \textit{$W$-fibration.} Alors $f\in W\Longleftrightarrow f'\in W$. Si $W$ satisfait à W(4c) (ce qui est le cas pour $W_\infty$), alors $f\in W$ implique bien que ses fibres sont $W$-asphériques, donc que c'est une $W$-équivalence universelle, i.e. $f'\in W\ \forall\, p:S'\longrightarrow S$. Mais pour l'implication inverse $f'\in W\Longrightarrow f\in W$, il faut supposer visiblement que
\[
\pi_0(p):\pi_0(S')\longrightarrow\pi_0(S)
\]
% foot: Version 7 décembre 2002

% ---- from Der-ChXIIb-012_p016-020.tex
\printedpage{58}
% head: A. Grothendieck : Dérivateurs, Ch. XII, Édition des Universités de Montpellier II et Paris VII
% foot: Version 7 décembre 2002

est \textit{surjectif.} (P. ex. si $S'=\emptyset$, on a $X'=\emptyset$, donc $f'\in W$, sans que ça implique $f\in W$!) Le fait que dans le cas cocartésien de W(7 bis), $i_0\in W\Longrightarrow i\in W$ doit être considéré comme ‘allant de soi’ - autrement le carré ne mérite pas, du point de vue $W$-homotopique, le nom de ‘cocartésien’. Par contre, l’implication inverse $i\in W\Longrightarrow i_0\in W$ m’apparaît comme un trait très particulier du dérivatuer HOT$_W$ (quand il le possède), et notamment % [sic] dérivatuer -> dérivateur, misprint
de HOT $=$ HOT$_{W_\infty}$. Le début
\msp{page 80}
de la remarque signifie simplement que le dérivateur dual HOT$^\circ$ ne possède pas cette même propriété. Je dois renvoyer à ce sujet aux deux axiomes pour les dérivateurs correspondants à W(7 bis), W(7), ce que je ferai sans doute sous peu, de façon à donner un sens plus précis aux présents commentaires.

\dgrblock{dgrb518}{What survives of the cube principle}
\begin{dcomment}{status}{p79}{What survives of the cube principle}
The distinction between direct and reverse excision is already visible in Grothendieck's later reorganization. Chapter~XVI separates the cube conditions into Loc(6a), Loc(6b), and Loc(6c); its Commentary~3.4 distinguishes the direct theory of homotopy cofibres from the extra reversibility property \citep[\S3, manuscript pp.~16--27]{GrothendieckDerivateursXVI}. For the pages at hand the reliable assertion for ordinary homotopy types is the direct one: weak equivalences of all three entries of a span give a weak equivalence of its homotopy pushout. A reverse five lemma cannot be added without hypotheses that exclude the preceding example. Theorem~2 has no repair with its stated hypotheses; its direct part is a theorem.

The pullback remark makes a different, correct observation. Base change need not meet every connected component of the base. Take $S=\{0,1\}$ discrete, and the discrete fibration $X\to S$ with one point over $0$ and two points over $1$. Pulling back along $\{0\}\to S$ gives an isomorphism, while $X\to S$ is not even bijective on connected components. This is a nonempty version of the empty-base example printed in the source.

Under the intended fibration hypothesis, contractibility of homotopy fibres is the correct test. If base change meets every component and the pullback is a weak equivalence, its homotopy fibres are contractible. Fibre transport along paths then shows that every homotopy fibre of the original fibration is contractible, so the original map is a weak equivalence. This argument explains the appearance of surjectivity on $\pi_0$; it must be applied where the source's W-fibration axioms really identify categorical fibres with homotopy fibres. Merely calling an arbitrary functor a fibration does not supply that identification. The corresponding discussion of categorical fibres was developed in the commentary on Chapter~VI, manuscript p.~45 (``Why the new fibration axiom selects full homotopy theory'').

The final contrast with the dual derivator should therefore be read in two parts. Its warning about detection after pullback is justified. Its attribution of unrestricted reverse excision to $\mathrm{HOT}$ depends on the false Theorem~2. Both qualifications are needed before applying the proposed axioms to the less special pushouts that follow.
\end{dcomment}


\dgrblock{dgrb519}{What the symmetric cube rule can detect}
\begin{dcomment}{development}{p79}{What the symmetric cube rule can detect}
The failed implication in Theorem~2 has a structural source. A homotopy pushout can erase information in its intersection. It is useful to determine whether passing to a bounded homotopy theory removes that obstruction, and to compare the answer with the homological argument suggested in the manuscript. The following finite family gives a uniform negative answer to the first question. We did not locate this precise family of tests in the source, its Appendix, or the later papers checked; the observation is elementary and is offered without a priority claim.

\textbf{Finite truncation obstruction.}\ours\quad
For an integer $n\geq0$, let $W_{\leq n}$ consist of functors whose nerves induce a bijection on components and isomorphisms on all homotopy groups in degrees $1,\ldots,n$, at every basepoint. There is a finite poset covered by two sieves for which the reverse implication of W(7 bis) fails for $W_{\leq n}$.

\emph{Proof.} Let $A_n$ be the poset of nonempty proper subsets of a set with $n+2$ elements, ordered by inclusion. A nondegenerate simplex of $NA_n$ is a strictly increasing chain of such subsets. Thus $NA_n$ is the barycentric subdivision of the boundary of the $(n+1)$-simplex, and $|NA_n|$ is homeomorphic to $S^n$. Adjoin two distinct elements $+$ and $-$, each greater than every element of $A_n$, with no order relation between them. Call the resulting poset $D_n$, and put $U_+=A_n\cup\{+\}$ and $U_-=A_n\cup\{-\}$. These are sieves, their intersection is $A_n$, and each has a greatest element. Their nerves are cones on $NA_n$. Since a chain of $D_n$ can contain at most one of $+$ and $-$, their union is exactly $ND_n$, whose realization is the suspension of $|NA_n|$, hence $S^{n+1}$.

It follows that $U_+\hookrightarrow D_n$ belongs to $W_{\leq n}$: both spaces are connected, and all their homotopy groups in degrees at most $n$ vanish. On the other hand, $A_n\hookrightarrow U_-$ does not belong to $W_{\leq n}$. For $n=0$ it identifies two components. For $n\geq1$ it kills the nonzero group $\pi_n(S^n)\simeq\mathbb Z$. This is precisely the forbidden reverse implication in the cover square.\hfill$\square$

The example for $n=1$ has eight objects in the union: six objects forming the barycentric subdivision of a triangle boundary and the two cone points. These tests do not replace the acyclic example at p.~78, since the latter is needed for untruncated weak equivalences. Together they show that the issue is not repaired by merely stopping the homotopy groups at some finite degree.

For comparison, consider the class of functors inducing isomorphisms on nerve homology with coefficients in a fixed ring $R$. The full symmetric cube rule \emph{does} hold for that class. Here is the complete chain-level reason. If $X=U\cup U'$ is a union of two sieves and $A=U\cap U'$, every chain of composable arrows in $X$ belongs to one of the two subcategories: choose a subcategory containing its last vertex and use the sieve condition. The same statement holds for nondegenerate chains. Consequently the normalized chain complexes have the natural short exact sequence
\[
0\longrightarrow C_*(NA;R)
 \xrightarrow{(\iota,-\iota')}
 C_*(NU;R)\oplus C_*(NU';R)
 \xrightarrow{\jmath+\jmath'} C_*(NX;R)
 \longrightarrow0.
\]
Injectivity, exactness in the middle, and surjectivity can each be checked on the basis of chains: those in the intersection occur twice, all others once. For a morphism between two such covers, take the mapping cones of the three vertical chain maps. The result is again a short exact sequence of complexes. Its long exact homology sequence implies that if any two of these cones are acyclic, then so is the third. The middle cone is the direct sum of the cones for the maps on $U$ and $U'$, so its acyclicity is equivalent to the acyclicity of both summands. This proves all three implications requested in W(7). For cosieve covers, the first vertex supplies the same proof.

The manuscript's proposed Mayer--Vietoris justification therefore has an exact domain of validity: it works for fixed-coefficient homology equivalences. The passage from those equivalences to weak homotopy equivalences is the step that cannot be supplied. The acyclic intersection in the preceding counterexample makes that distinction visible without any ambiguity about the existence or computation of the pushout.
\end{dcomment}


\subsection*{2.2 Application aux carrés cocartésiens déduits d’une immersion ouverte ou fermée.}

Les axiomes W(7), W(7 bis) (du cube, et du carré cocartésien) concernent des carrés cocartésiens de nature très particulière. Nous allons voir à présent comment en déduire des énoncés du même type pour des carrés cocartésiens beaucoup plus généraux, y compris pour des carrés qui n’ont rien de cocartésien (dans Cat) - et que nous appelerons % [sic] appelerons -> appellerons, misspelling
‘homotopiquement cocartésiens’ (relativement à $W$), ou simplement $W$-cocartésiens. (Il y aura également des carrés $W$-cartésiens un jour \dots) Dans le cas où $W=W_\infty$, je les appelerai aussi Hot-cocartésiens. % [sic] appelerai -> appellerai, misspelling

Je commencerai par le cas particulier d’un carré cocartésien associé à une immersion ouverte (diagramme ($*$), p. 74) :
\[
(*)\qquad\qquad
\begin{tikzcd}
U \arrow[r, hook, "i"] \arrow[d, "p"'] & X \arrow[d, "q"'] & T \arrow[l, hook', dotted, "j"'] \arrow[d, dotted] \\
U' \arrow[r, hook, "i'"] & X' & T'\simeq T. \arrow[l, hook', dotted, "j'"']
\end{tikzcd}
\]
\msp{page 81}

Le grand traquenard ici qu’il faut arriver à exorciser, c’est qu’il peut arriver que $i\in W$, sans qu’on ait $i'\in W$, et aussi $p\in W$ sans qu’on ait $q\in W$. Je m’en étais convaincu en 1983 en écrivant Pursuing Stacks (durant les semaines de ‘scratchwork’), et n’ai pas envie de rechercher un contre-exemple ici, s’il ne se présente de lui-même (ce qu’il fera sans doute).

Je vais donner un cas important où ces ennuis disparaissent :

\dgrblock{dgrb520}{Separating the boundary from the ambient category}
\begin{dcomment}{exposition}{p81:Je vais donner un cas important où ces ennuis disparaissent}{Separating the boundary from the ambient category}
Grothendieck wants to isolate the part of a pushout which can change its homotopy type. The complement $T$ is carried isomorphically to $T'$, whereas the arrows entering it from $U$ are changed by a Kan extension. The useful intermediate category is therefore the closure $\bar U$, which consists of $U$ and the objects of $T$ reached by at least one arrow from $U$. Write $A=\bar U\cap T$. We then have a square
\[
\begin{tikzcd}
A\arrow[r]\arrow[d]&T\arrow[d]\\
\bar U\arrow[r]&X
\end{tikzcd}
\]
of cosieve inclusions. It is both Cartesian and cocartesian. Every arrow of $X$ lies in one of these two full subcategories: if its source lies in $\bar U$, its target does too; otherwise its source lies in $T$, and its target lies in $T$. The same assertion for composable strings shows that the nerve is the union of the two corresponding simplicial subsets. Its realization is consequently a homotopy pushout, since inclusions of simplicial subsets realize as cofibrations.

The formula for the pushed-out gluing presheaf shows that an object $t\in T$ receives an arrow from $U'$ after pushout exactly when it received an arrow from $U$ before pushout. Indeed, a nonempty Kan-extension value is represented by an original arrow and a morphism in $U'$; conversely, an original arrow supplies such a value at the image of its source. Thus $A\to A'$ is an isomorphism. All the movement has been confined to $\bar U\to\overline{U'}$.

This explains the architecture of Proposition~6. If $U\to\bar U$ and $U'\to\overline{U'}$ are weak equivalences, two-out-of-three replaces the question about $p$ by the question about $\bar p$. The resulting map of closed-cover squares is the identity on $A$ and $T$. The ordinary homotopy gluing lemma therefore proves, for $W_\infty$, the useful implication
\[
p\in W_\infty\quad\Longrightarrow\quad q\in W_\infty.
\]
This argument requires only the usual forward gluing statement.

The biconditionals in the printed proposition have stronger hypotheses. They explicitly invoke W(7~bis) and W(7), with their reverse implications as formulated on pp.~76--77. The counterexample discussed at p.~78 prevents us from treating those axioms as properties of $W_\infty$. Proposition~6 remains a valid formal deduction under its stated axioms; its proof is not a proof that ordinary weak homotopy equivalences can be detected backwards from every such pushout. This distinction will matter when the one-point test is proposed on p.~85.
\end{dcomment}


\textbf{Proposition 6 :} \textit{Considérons les adhérences $\bar U$ et $\overline{U'}$ de $U$ et $U'$ dans $X$, $X'$ respectivement, et supposons que les inclusions}
\[
U\xrightarrow{i_0}\bar U,\qquad U'\xrightarrow{i'_0}\overline{U'}
\]
\printedpage{59}
% head: A. Grothendieck : Dérivateurs, Ch. XII, Édition des Universités de Montpellier II et Paris VII
% foot: Version 7 décembre 2002
\textit{soient dans $W$} (\footnote[64]{cf. th. 3 p. 107 pour des conditions générales assurant que cette hypothèse est satisfaite.}). \textit{Alors}
\[
\left\{
\begin{array}{llll}
\circled{\textit{a}}\ i\in W & \iff & i'\in W & \textit{si on a }\mathrm{W}(7\text{ bis})\\
\circled{\textit{b}}\ p\in W & \iff & q\in W & \textit{si on a }\mathrm{W}(7).
\end{array}
\right.
\]

\textsc{Démonstration}. Considérons d’abord le cas \circled{a} de $i$, $i'$. Soit
\[
k:\bar U\hookrightarrow X
\]
l’inclusion. Le triangle
\begin{tikzcd}
U \arrow[r, "i_0"] \arrow[rr, bend right=8, "i"'] & \bar U \arrow[r, "k"] & X
\end{tikzcd}
, où $i_0\in W$, montre que
\[
i\in W\iff k\in W.
\]

D’autre part, $X$ est la réunion des fermés $\bar U$, $T$, donc on a un carré cocartésien dans Cat
\[
(*)\qquad\qquad
\begin{tikzcd}
\bar U\cap T=\dot T \arrow[r, "k_0"] \arrow[d, "j_0"'] & T \arrow[d, "j"] \\
\bar U \arrow[r, "k"] & X.
\end{tikzcd}
\]

L’axiome W(7 bis) (sous sa forme duale (i.e. appliqué au carré formé par les catégories opposées \dots)) implique
\msp{page 82}
que
\[
k\in W\iff k_0\in W,
\]
d’où
\[
i\in W\iff (k_0:\dot T\longrightarrow T)\in W.
\]
De même on trouve, en appliqunt ceci à $(X',U')$ au lieu de $(X,U)$ % [sic] appliqunt -> appliquant, misprint
\[
i'\in W\iff (k'_0:\dot T'\longrightarrow T')\in W.
\]
Or on a $T\xrightarrow{\sim}T'$ (isomorphisme), et par la proposition 4 cet isomorphisme transforme $\dot T$ en $\dot T'$. Donc
\[
k_0\in W\iff k'_0\in W\qquad\quad\text{(trivialement)},
\]
d’où a).

Prouvons b). Le carré commutatif induit par ($*$), p. 80,
\printedpage{60}
% head: A. Grothendieck : Dérivateurs, Ch. XII, Édition des Universités de Montpellier II et Paris VII
% foot: Version 7 décembre 2002
\[
\begin{tikzcd}
U \arrow[r, hook, "i_0"] \arrow[d, "p"'] & \bar U \arrow[d, "\bar p"] \\
U' \arrow[r, hook, "i'_0"] & \overline{U'},
\end{tikzcd}
\]
où $i_0,i'_0\in W$ par hypothèse, montre que
\[
p\in W\iff\bar p\in W,
\]
donc tout revient à prouver que
\[
\bar p\in W\overset{?}{\iff}q\in W.
\]
Mais $q$ induit un homomorphisme du carré ($*$) déduit de $(X,\bar U,T)$ (avec $X=\bar U\cup T$), dans le carré similaire pour $X'$
\[
\begin{tikzcd}
\overline{U'}\cap T'=\dot T' \arrow[r, "k'_0"] \arrow[d, "j'_0"'] & T \arrow[d, "j'"] \\ % [sic] T -> T', the square for X' should have the corner T' (although T'\simeq T)
\overline{U'} \arrow[r, "k'"] & X',
\end{tikzcd}
\]
les flèches verticales étant notées $q_{\dot T}$, $q_T$, $\bar p$, $q$.
\msp{page 83}
On sait que $q_{\dot T}$, $q_T$ sont des isomorphismes, donc dans $W$, donc le ‘lemme des cinq’ W(7) nous assure que $\bar p\in W\iff q\in W$, q.e.d.

\textbf{Commentaire.} Cette proposition nous amène à chercher des conditions sur un ouvert $U$ d’un $X$, qui assurent que l’inclusion $i_0:U\longrightarrow\bar U$ soit \texttt{[dans]} $W$, et \textit{que cette propriété soit conservée par tout cochangement de base $U\longrightarrow U'$.} De plus, nous recherchons des conditions qui soient de nature locale sur $\bar U$, plus exactement sur $\bar U\cap T=\dot U=\dot T$. Nous pouvons motiver ce souci (a posteriori) en disant que nous voudrions que la propriété $i_0\in W$ soit conservée, non seulement par cochangement de base quelconque $U\longrightarrow U'$, mais aussi, de plus, par passage de $X$ à toute sous-catégorie pleine $X_1$ contenant $U$, donnant donc naissance à une situation
\[
U\overset{i_Y}{\hookrightarrow}X_1,\qquad U\xrightarrow{(i_Y)_0}\bar U^{X_1}=\bar U\cap X_1
\]
\texttt{[}$Y=X_1$\texttt{?]}, similaire à la situation de départ. Plus généralement encore, nous voudrions qu’elle soit conservée en passant du triple $(X,U,T)$ à un triple “induit” $(X_1,U,T_1)$, par une flèche quelconque $g:T_1\longrightarrow T$ de but $T$ dans Cat, donnant naissance au couple de flèches
\[
U\xrightarrow{\mathrm{id}}U,\ g:T_1\longrightarrow T
\]
(cf. pages 60, 61). Le cas précédent est celui où
\printedpage{61}
% head: A. Grothendieck : Dérivateurs, Ch. XII, Édition des Universités de Montpellier II et Paris VII
% foot: Version 7 décembre 2002
\msp{page 84}
$T_1$ est une sous-catégorie pleine de $T$, $g$ la flèche d’inclusion, Notons à ce sujet la proposition soritale % [sic] inclusion, Notons -> inclusion. Notons, comma before a capital letter

\textbf{Proposition 7.} \textit{Conditions équivalentes sur l’ouvert $U$ d’un objet $X$ de} Cat (moyennant W(7 bis)) :

\begin{enumerate}[label={}, leftmargin=*]
\item[\textit{a)}] \textit{L’inclusion $i_0:U\longrightarrow\bar U$ est $W$-asphérique, i.e. pour tout $t\in\dot U=\bar U\backslash U$, la catégorie $U/t$ est $W$-asphérique.}
\item[\textit{b)}] \textit{L’inclusion $i_0:U\longrightarrow\bar U$ est dans $W$, et cette condition reste satisfaite, quand on remplace $(X,U)$ par $(X_1,U)$, où $(X_1,U)$ est déduit de $(X,U)$ au moyen d’une flèche $T_1\longrightarrow T(=X\backslash U)$ dans} Cat\textit{, comme rappelé ci-dessus.}
\item[\textit{b$'$)}] \textit{Comme b), dans le cas où $T_1$ est la catégorie ponctuelle $e$, et où de plus son image est dans $\dot U$.}
\end{enumerate}

\textsc{Démonstration}. Il est évident que si dans b), b$'$) on remplace la condition “$i_0\in W$” par la condition plus forte ‘$i_0$ est $W$-asphérique’, les quatre \texttt{[plutôt trois]} conditions a) à b$'$) sont équivalentes. Il reste à prouver que la condition que les $U/t$ soient $W$-asphériques ($t\in\bar U$ - mais c’est trivial pour $t\in U$, donc il suffit de l’exprimer pour $t\in\bar U\backslash U=\dot U$) est déjà impliquée par la plus faible b$'$ et on devine qu’on a ceci :

\msp{page 85}

\textbf{Corollaire.} (Moyennant W(7 bis).) \textit{Soit $X$ dans} Cat\textit{, $U$ un ouvert de $X$ tel que $X\backslash U=\overline{\{t\}}$ soit la catégorie finale. Alors l’inclusion $U\xrightarrow{i}X$ est dans $W$ si et seulement si elle est $W$-asphérique, i.e. si et seulement si $U/t$ est $W$-asphérique (où $t$ est l’unique élément de $T$).}

Nous démontrons ultérieurement ce corollaire (\footnote[65]{cf. cor. 4 de la proposition 15, page 149.}). La proposition 7 ne nous servira pas d’ici là - elle n’aura qu’un rôle purement esthétique à jouer, pour justifier notre propos délibéré de chercher des conditions sur $i:U\longrightarrow X$ (pour assurer que $U\longrightarrow\bar U$ est dans $W$ et le reste, quand on remplace $(X,U)$ par d’autres couples $(X_1,U_1)$ qui s’en déduisent par certaines opérations \dots) qui
\msp{page 86}
soient de \textit{nature locale} sur $\bar U$ (en fait, en les points de $\dot U$), et qui en tout état de cause entrainent la \textit{$W$-asphéricité} de $U\longrightarrow\bar U$.

\dgrblock{dgrb521}{The one-point test and the direction of Theorem A}
\begin{dcomment}{exposition}{p85}{The one-point test and the direction of Theorem A}
The purpose of Proposition~7 is to turn a condition on the whole inclusion into a condition at each new point. For a sieve $i:U\hookrightarrow X$, the comma category $U/u$ has final object $(u,\mathrm{id}_u)$ whenever $u\in U$. All possible failures of asphericity are therefore concentrated at the complementary objects. Replacing $T$ by the point selecting $t$ preserves precisely the presheaf $u\mapsto\Hom_X(u,t)$, and hence precisely the category $U/t$. This proves the equivalence of the three versions in which every occurrence of ``belongs to $W$'' is replaced by ``is $W$-aspherical.'' The remaining implication is exactly the deferred corollary, not a consequence of the definition.

For ordinary weak homotopy equivalences, the corollary's reverse implication is false. To see its geometry without altering the source, take a connected acyclic category $A$ whose nerve is not weakly contractible, as in the counterexample at p.~78. Form $U$ by adjoining a final object $u_+$ to $A$. Form $X$ by adjoining instead two distinct final cone points $u_+$ and $t$, with no morphisms between them. Then $U\hookrightarrow X$ is a sieve and its complement is the one-object category $\{t\}$. Every object of $A$ has its unique arrow to $t$, while $u_+$ has none, so
\[
U/t\simeq A.
\]
On the other hand, $BU$ is contractible and $BX\simeq\Sigma BA$. The suspension of a connected acyclic CW complex is simply connected and acyclic, hence weakly contractible by the Hurewicz and Whitehead theorems. Thus $U\to X$ is a weak homotopy equivalence although it is not aspherical.

The missing distinction is between the homotopy type of an attachment and the homotopy type of its attaching category. A cone attachment can erase a nontrivial perfect fundamental group without leaving any homology behind. Theorem~A supplies the implication from contractible comma categories to weak equivalence; it does not supply the converse. In modern language this is also the distinction between a weak equivalence of classifying spaces and a homotopy initial functor: the comma categories here are $U/t=(i\downarrow t)$. Compare the dual of \citet[Theorem~4.1.3.1 and Corollary~4.1.3.3]{Lurie2009HTT}, applied to $i^{\mathrm{op}}$. The source has retained W(7~bis) in the corollary, and the audit retains it too. Its subsequent use for $W_\infty$ needs an independent argument in each application.
\end{dcomment}


J’avais espéré que cette condition d’asphéricité (suggérée par la notion ‘d’immersion de Dwyer’) serait stable par cochangement de base quelconque $U\longrightarrow U'$, ce qui aurait assuré en même temps qu’après une telle opération, l’immersion ouverte $i':U'\longrightarrow X'$ obtenue satisfait encore à la condition que $i'_0:U'\longrightarrow\overline{U'}$ est dans $W$. Nous verrons que cette attente n’est pas fondée - la condition de $W$-asphéricité n’est \textit{pas} stable par cochangement de base, et elle ne suffit \textit{pas} même à assurer qu’après un cochangement de base, $i'_0:U'\longrightarrow\overline{U'}$ soit
\printedpage{62}
% head: A. Grothendieck : Dérivateurs, Ch. XII, Édition des Universités de Montpellier II et Paris VII
% foot: Version 7 décembre 2002
dans $W$. Il faudra une notion plus forte, qui est suggérée par le cas où $T=\{t\}$ est la catégorie finale et contenue dans $\bar U$ (cas auquel nous avons essayé de nous ramener dans la proposition 7). Dans ce cas, on l’a dit (corollaire précédent), la condition que $i:U\longrightarrow\bar U=X$ soit dans $W$, équivaut à la $W$-asphéricité de $i$, i.e. à la $W$-asphéricité de $U/t$, et de même pour $i':U'\longrightarrow\overline{U'}=X'$ : elle est dans $W$ si et seulement si $U'/t'$ est $W$-asphérique, où
\msp{page 87}
$t'$ est l’unique élément de $T'(\simeq T)=X'\backslash U'$. Donc il faut des conditions sur $U/t$ qui assurent que non seulement $U/t$, mais aussi toutes les catégories similaires $U'/t$ \texttt{[plutôt} $U'/t'$\texttt{]}, soient $W$-asphériques. Dans le cas général où $T$ est quelconque, c’est cette condition-là que nous allons imposer aux $U/t$ ($t\in\Ob T$), et il s’avérera que cette fois-ci, elles sont stables \texttt{[plutôt elle est stable]} par tout cochangement de base $U\xrightarrow{p}U'$, i.e. que les $U'/t'$ ($t'\in\Ob T'$) satisferont alors à la même condition. C’est cette condition que nous allons à présent reformuler.

\dgrblock{dgrb522}{A boundary presheaf and all its direct images}
\begin{dcomment}{exposition}{p87:C’est cette condition que nous allons à présent reformuler.}{A boundary presheaf and all its direct images}
Grothendieck now strengthens the boundary condition in the way dictated by the actual pushout formula. The extra quantifier is the essential feature. Given $H:U^{\mathrm{op}}\to\mathbf{Set}$, let $\int_U H=U/H$ denote its category of elements. Its objects are pairs $(u,h)$ with $h\in H(u)$, and a morphism $(u,h)\to(v,k)$ is an arrow $a:u\to v$ satisfying $H(a)(k)=h$. If $X$ is obtained by adjoining a complementary point $t$ with $\Hom_X(u,t)=H(u)$, then $U/t=\int_U H$ literally records all ways to enter that point.

For $p:U\to U'$, the new boundary is $p_!H$, where $p_!$ is the ordinary left Kan extension along $p^{\mathrm{op}}$. It is useful to calculate it through the category of elements, rather than to lose the variance in an informal colimit. With $\alpha:\int_UH\to U$ the projection,
\[
(p_!H)(u')\cong\pi_0\bigl(u'\downarrow p\alpha\bigr).
\]
An object on the right is $(u,h,\beta:u'\to p(u))$. Its connected component is the generator of the Kan extension represented by that triple. A map $a:u\to v$ with $H(a)(k)=h$ identifies the generator $(u,h,\beta)$ with $(v,k,p(a)\beta)$; these are exactly the relations in the pointwise colimit. This gives both the formula and its functoriality in $u'$.

The ordinary asphericity of $H$ only tests $p=\mathrm{id}_U$. Total asphericity requires the entire category of elements of every $p_!H$ to be aspherical. The operation takes connected components before the resulting presheaf is assembled over $U'$. It can therefore destroy homotopy information inside each comma category while creating nontrivial topology in the assembled category. The finite example at p.~100 will make this effect visible in degree two.

The definition has a deliberate stability property. If $H$ is totally $W$-aspherical and $p:U\to U'$ is any functor, then $p_!H$ is totally $W$-aspherical: for every $q:U'\to U''$ there is a canonical isomorphism $q_!p_!H\cong(qp)_!H$. Thus the strengthened condition is preserved by the very operation against which it was designed to protect the boundary. This is the formal reason the later definition of a $W$-Dwyer immersion uses these presheaves.

The editors' warning about terminology is necessary. The earlier test-category usage of ``totally aspherical'' concerns products or diagonal functors. Here the quantifier runs over all outgoing functors and their set-valued left Kan extensions. The commentary on Chapter~II, manuscript p.~36 (``The quantifier in total asphericity'') already treats this quantifier; the present pages apply it to a concrete gluing problem. Likewise the filtered examples should be read together with the commentary on Chapter~II, manuscript pp.~32--35, rather than as a new assertion that contractibility alone controls all direct images.
\end{dcomment}


\subsection*{2.3 Objets totalement $W$-asphériques de \textnormal{Cat}, et morphismes totalement $W$-asphériques.}

Reprenons la situation
\[
U\overset{i}{\hookrightarrow}X\overset{j}{\hookleftarrow}T,\qquad\qquad\text{avec }T=\{t\}\text{, catégorie finale, }t\in\bar U\text{, i.e. }X=\bar U,
\]
et celle d’un cochangement de base $U\longrightarrow U'$, donnant naissance à un carré cocartésien dans Cat
\[
(*)\qquad\qquad
\begin{tikzcd}
U \arrow[r, hook, "i"] \arrow[d, "p"'] & X \arrow[d, "q"'] & T=\{t\} \arrow[l, hook', dotted, "j"'] \arrow[d, dotted, "\mathrm{iso}"] \\
U' \arrow[r, hook, "i'"] & X' & T'=\{t'\}. \arrow[l, hook', dotted, "j'"']
\end{tikzcd}
\]
La donnée de recollement $H$ entre $U$ et $T$, $H:\underbrace{U^\circ\times T}_{\simeq U^\circ}\longrightarrow\mathrm{Ens}$, peut s’interpréter comme un préfaisceau sur $U$, $H\in\Ob U^\wedge:U^\circ\longrightarrow\mathrm{Ens}$,
\msp{page 88}
et on a
\[
U/t\simeq U/H,
\]
où $U/H$ est interprété en termes du plongement canonique
\[
U\hookrightarrow U^\wedge.
\]
De même, le recollement entre $U'$ et $T'$ sera donné par un préfaisceau
\[
H'\in\Ob U^{\prime\wedge}:U^{\prime\circ}\longrightarrow\mathrm{Ens},
\]

% ---- from Der-ChXIIb-012_p021-025.tex
\printedpage{63}
% head: A. Grothendieck : Dérivateurs, Ch. XII, Édition des Universités de Montpellier II et Paris VII
et on aura
\[ U'/t' \simeq U/H'. \]
La relation entre $H'$ et $H$ a été explicitée au §1, comme
\[ H' = p_!H, \qquad \text{où } p_! : U^{\wedge} \longrightarrow U'^{\wedge} \]
est l'adjoint à gauche de $p^* : U'^{\wedge} \longrightarrow U^{\wedge}$, donné par
\[ p^*F' = F' \circ p. \]
On a
\[ p_!(H)(u') \simeq \pi_0(u'\backslash U/H). \]

\textbf{Définition 1.} \cercle{a} Soit $U$ un objet de Cat, $H \in \Ob U^{\wedge}$ un préfaisceau sur $U$. On dit que $H$ est \textit{totalement $W$-asphérique} si pour toute flèche $p : U \longrightarrow U'$ de Cat, $p_!(H)$ dans $U'^{\wedge}$ est un objet $W$-asphérique, i.e. $U'/p_!H$ est $W$-asphérique. \cercle{b} On dit qu'un objet $V$ de \msp{page 89} Cat est \textit{totalement $W$-asphérique,} si le foncteur final $e_V$ sur $V$ ($v \longmapsto \{\emptyset\}$) est totalement $W$-asphérique, i.e. si pour toute flèche $f : V \longrightarrow V'$ dans Cat, le préfaisceau
\[ \pi_0(f) \defeq f_!(e_V) \defeq (v' \longmapsto \pi_0(v'\backslash V)) \]
sur $V'$ est $W$-asphérique, i.e. si la catégorie
\[ \upi_0(f) \ \text{ ou } \ \upi_0(V/V') \defeq V'/\pi_0(f) \]
est $W$-asphérique. \texttt{[Cette notion est différente de la notion de totale asphéricité de Pursuing Stacks.]}

Notons que la définition des foncteurs $p_!$ (ou $f_!$) comme adjoints, et leur `calcul' en termes de $\pi_0$ de certaines catégories, ne dépend pas du localiseur $W$. Donnons quelques précisions à ce sujet. On a un morphisme d'adjonction dans $V^{\wedge}$
\[ f^*f_!(F) \longleftarrow F, \]
d'où
\[ V/F \longrightarrow V/f^*f_!(F) \simeq V \times_{V'} (V'/f_!(F)), \tag*{$(*)$} \]
le dernier isomorphisme provenant du fait que l'on a commutativité, à isomorphisme canonique près, dans le diagramme de foncteurs
\[
\begin{tikzcd}[column sep=14em, row sep=large]
V'^{\wedge} \arrow[r, "f^*=(F'\longmapsto F'\circ f)"] \arrow[d, "F'\longmapsto V'/F'"'] & V^{\wedge} \arrow[d, "F\longmapsto F/V"] \\
\Cat/V' \arrow[r, "f^*_{\Cat}=(W'\longmapsto W'\times_{V'}V)"] & \Cat/V,
\end{tikzcd}
\]
% foot: Version 7 décembre 2002
\printedpage{64}
% head: A. Grothendieck : Dérivateurs, Ch. XII, Édition des Universités de Montpellier II et Paris VII
\msp{page 90}
en d'autres termes, l'opérateur $f^*$ sur les préfaisceaux s'interprète, en termes géométriques, comme l'opération $W' \longmapsto W' \times_{V'} V$ de changement de base par $V \longrightarrow V'$, sur les catégories au dessus de $V'$ qui sont fibrées à fibres discrètes (propriété qui se conserve par changement de base).

Ceci dit, la donnée d'une flèche $(*)$ équivaut, par la propriété universelle du changement de base, à celle d'une flèche (canonique en l'occurrence)
\[ f/F : V/F \longrightarrow V'/f_!(F) \]
rendant commutatif le carré
\[
\begin{tikzcd}[column sep=large, row sep=huge]
V/F \arrow[r, "\text{oubli de }F", "\substack{\text{ou foncteur}\\ \text{source}}"'] \arrow[d, "f/F"'] & V \arrow[d, "f"] \\
V'/(F' = f_!(F)) \arrow[r, "\text{oubli de }F'"] & V'.
\end{tikzcd}
\tag*{$(*)$}
\]
Dans le cas qui nous intéresse surtout ici, $F = e_V$, $V/F \longrightarrow V$ est un isomorphisme, et on obtient une factorisation canonique de foncteurs \msp{page 91}
\[
\begin{tikzcd}[row sep=large]
 & V \arrow[dl, "\tilde{f}=f/e_V"'] \arrow[dd, "f"] \\
\begin{array}{r} \upi_0(V/V') \\ = V'/f_!(e_V) \end{array} \arrow[dr, "\text{can.}"'] & \\
 & V'.
\end{tikzcd}
\tag*{$(**)$}
\]
Explicitons le $V'$-morphisme $\tilde{f}$ sur les fibres,
\[ \tilde{f}_{v'} : V_{v'} \longrightarrow \upi_0(V/V')_{v'} \defeq \pi_0(v'\backslash V), \tag*{$(***)$} \]
où dans le dernier membre, on a identifié l'ensemble $\pi_0(v'\backslash V)$ à la catégorie discrète qu'il définit. Pour expliciter cette flèche $(***)$, notons qu'on a une inclusion canonique
\[ V_{v'} \longrightarrow v'\backslash V \]
et pour toute catégorie $W$, un morphisme canonique (d'adjonction) $W \longrightarrow \pi_0(W)$, d'où un composé
\[ V_{v'} \longrightarrow v'\backslash V \longrightarrow \pi_0(v'\backslash V), \]
qui n'est autre que $\tilde{f}_{v'}$ dans $(***)$.

Il doit être facile de montrer que la factorisation $(**)$ de $f$ est caractérisée, à isomorphisme unique près, en tant que
% foot: Version 7 décembre 2002
\printedpage{65}
% head: A. Grothendieck : Dérivateurs, Ch. XII, Édition des Universités de Montpellier II et Paris VII
\msp{page 92}
factorisation dans $\Cat/V'$
\[
\begin{tikzcd}[column sep=small]
 & \tilde{V} \arrow[dr, "\pi"] & \\
V \arrow[ur, "\tilde{f}"] \arrow[rr, "f"'] & & V',
\end{tikzcd}
\]
par les propriétés suivantes :
\begin{enumerate}[label={}, leftmargin=*]
\item[a)] $\pi$ est fibrée à fibres discrètes, i.e. un isomorphisme local (i.e. $\tilde{V}$ est isomorphe dans $\Cat/V'$ à une catégorie $V/F'$, avec $F'$ dans $V'^{\wedge}$). % [sic] V/F' -> V'/F', une catégorie au-dessus de V' définie par F' dans V'^∧ est de la forme V'/F'
\item[b)] $\tilde{f}_!(e_V) = e_{\tilde{V}}$, i.e. pout tout $\tilde{v} \in \Ob \tilde{V}$, la catégorie $\tilde{v}\backslash V$ est 0-connexe : les colocalisés de $V$ au dessus de $\tilde{V}$ sont 0-connexes. % [sic] pout -> pour, faute de frappe
\end{enumerate}

Notons ici, en passant, que pour \textit{tout} morphisme $\pi : W' \longrightarrow W$ dans Cat qui soit fibré à fibres discrètes, i.e. tel que $W'$ soit isomorphe dans $\Cat/W$ à une catégorie de la forme $W/F$ ($F \in \Ob W^{\wedge}$), on a
\[ \pi_!(e_{W'}) = F, \qquad \text{i.e. } W' \underbrace{\ \simeq\ }_{\text{dans } \Cat/W} W/\pi_!(e_{W'}). \tag*{$(*)$} \]
Il faut vérifier que le foncteur
\[ w \longmapsto \pi_0(\underbrace{w\backslash(W/F)}_{\defeq w\backslash W/F}) \]
est canoniquement isomorphe à $F$. Mais l'inclusion
\[ \underbrace{W'_w}_{= F(w) \text{ (catégorie discrète)}} \qquad \longrightarrow w\backslash W' \]
\msp{page 93}
définit bien comme tantôt
\[ F(w) = W'_w \longrightarrow \pi_0(w\backslash W'), \]
il faut prouver que c'est bijectif. Or, $W'$ étant $W_\infty$-lisse sur $W$, on sait que l'inclusion $W'_w \hookrightarrow w\backslash W'$ est dans $W_\infty$, a fortiori induit un isomorphisme pour les $\pi_0$.

Ceci posé, si on a une factorisation de $f : V \longrightarrow V'$ en $f = \pi\tilde{f}$ ayant les propriétés a) b) de la page précédente, on voit déjà que l'on a
\[ f_!(e_V) \simeq \pi_!(\, \underbrace{\tilde{f}_!(e_V)}_{\simeq e_{\tilde{V}} \text{ par b)}} \,) \simeq \pi_!(e_{\tilde{V}}), \]
ce qui déjà détermine $\tilde{V}$ à isomorphisme canonique près (grâce à a) comme
\[ \tilde{V} \simeq V'/f_!(e_V) \defeq \upi_0(V'/V), \] % [sic] \underline{\pi}_0(V'/V) -> \underline{\pi}_0(V/V'), notation employée partout ailleurs (Définition 1, (**))
% foot: Version 7 décembre 2002
\printedpage{66}
% head: A. Grothendieck : Dérivateurs, Ch. XII, Édition des Universités de Montpellier II et Paris VII
et il ne doit pas être bien dur de se convaincre que $\tilde{f}$ n'est autre que la flèche canonique de la page 91.

\textbf{Proposition 8.} \textit{Pour qu'un objet $H$ de $U^{\wedge}$ soit totalement $W$-asphérique, il suffit que l'objet $U/H$ de} Cat \textit{soit totalement $W$-asphérique.}

\msp{page 94}
Supposons en effet cette dernière condition satisfaite et prouvons que $H$ est totalement $W$-asphérique, i.e. que pour tout $p : U \longrightarrow U'$, $U'/p_!(H)$ soit $W$-asphérique.
\[
\begin{tikzcd}[row sep=large]
V \defeq U/H \arrow[r, "\alpha"] \arrow[dr, "f"'] & U \arrow[d, "p"] \\
 & U'
\end{tikzcd}
\]
Soit $f$ le foncteur composé du diagramme ci-contre \texttt{[ci-dessus]}, on a $H \simeq \alpha_!(e_V)$ (par $(*)$ p. 92), d'où
\[ p_!(H) \simeq p_!\alpha_!(e_V) \simeq (p\alpha)_!(e_V) \simeq f_!(e_V), \]
qui est $W$-asphérique par hypothèse sur $V = U/H$.

\dgrblock{dgrb523}{The connected-component factorization, with its proof}
\begin{dcomment}{exposition}{p93}{The connected-component factorization, with its proof}
Grothendieck's relative $\underline\pi_0$ turns a map of categories into a genuine family of sets over its target. This is the right construction for understanding what ordinary Kan extension forgets. Let $f:A\to B$, put
\[
F(b)=\pi_0(b\downarrow f),\qquad E=B/F,
\]
and let $\pi:E\to B$ be the projection. Precomposition with a morphism $b'\to b$ gives the restriction map $F(b)\to F(b')$, so $F$ is a presheaf. The factorization map
\[
j:A\longrightarrow E,\qquad
j(a)=\bigl(f(a),[(a,\mathrm{id}_{f(a)})]\bigr)
\]
is functorial: a morphism $a\to a'$ identifies the evident two objects in $f(a)\downarrow f$. We have $f=\pi j$.

To prove the missing connectedness assertion, fix $e=(b,c)\in E$. An object of $e\downarrow j$ is an $a\in A$ together with $\beta:b\to f(a)$ such that the component of $(a,\beta)$ is $c$. Its morphisms are exactly the morphisms of $b\downarrow f$ between those objects. Thus
\[
e\downarrow j\ \cong\ (b\downarrow f)_c,
\]
the component indexed by $c$. It is nonempty and connected by definition. In the usual ordinary-category terminology, $j$ is a \emph{final functor}. No higher contractibility is being asserted.

We also need to recover a discrete fibration from this procedure. Suppose $\pi:E\to B$ is the category of elements of a presheaf $F$. For $(e,\beta:b\to\pi e)$ in $b\downarrow\pi$, the unique lift of $\beta$ into $e$ has source $(b,F(\beta)(e))$. This assignment gives a functor to the discrete fibre $E_b$. Its composite with the fibre inclusion is naturally connected to the identity by the lifted arrows, and its restriction to the fibre is the identity. Consequently
\[
\pi_0(b\downarrow\pi)\cong F(b),\qquad
\pi_!1_E\cong F.
\]
The same argument proves a homotopy equivalence between the comma category and the discrete fibre, which explains the appeal to smoothness on p.~93 without requiring that theory here.

Finally, suppose $f=\rho k$ is another factorization with $\rho:D\to B$ discrete and $k:A\to D$ final. By the Kan-extension formula $k_!1_A=1_D$, and hence
\[
f_!1_A\cong\rho_!k_!1_A\cong\rho_!1_D.
\]
The preceding recovery calculation identifies $D$ over $B$ with $B/F=E$. This identification also respects the first factor. Indeed, for any presheaf $G$ on $B$, a functor $A\to B/G$ over $B$ is the same as a natural transformation $1_A\to f^*G$, and adjunction identifies these with natural transformations $f_!1_A\to G$. For $G=F$, the canonical first factor corresponds to $\mathrm{id}_F$; the identification obtained from $\rho_!1_D\cong F$ carries $k$ to that same map. This proves uniqueness with its compatibility, rather than only determining the middle category as an abstract object.

The result is the discrete-fibration form of the comprehensive factorization of \citet{StreetWalters1973}. A complete direct proof in the opposite variance is in \citet[\S4.7, pp.~73--75, especially (4.77)--(4.79)]{Kelly1982}; Kelly also identifies an invalid Kan-extension inference in the original proof and replaces it by the component calculation used above. The theorem predates these manuscript pages. Its role here is to distinguish the final first factor from a homotopy final map, whose comma categories would have to be weakly contractible.
\end{dcomment}


\dgrblock{dgrb524}{From comprehensive factorization to homotopy families}
\begin{dcomment}{later}{p93}{From comprehensive factorization to homotopy families}
The geometric interpretation of $\underline\pi_0$ has continued to be useful independently of derivators. \citet[\S1.3, Lemma~1.6, and Theorem~1.7]{BergerKaufmann2017} place the two ordinary comprehensive factorizations in a general theory of comprehension schemes: presheaves supply discrete fibrations, and the unit of comprehension supplies their universal approximation. Their Proposition~1.10 and Remark~1.11 explain the associated base-change and Frobenius properties. These are structural statements about the factorization, not claims that its first factor is a weak homotopy equivalence.

For the comparison with higher categories, the decisive change is to keep the homotopy type of each comma category instead of only its set of components. The homotopy-finality criterion requires $B(b\downarrow f)$ to be contractible for every $b$, as in \citet[Theorem~4.1.3.1]{Lurie2009HTT}. Here $j$ only makes these spaces connected. Thus comprehensive factorization is a useful discrete approximation, but it cannot replace a construction of a homotopy Kan extension. The examples below measure precisely what is lost when this approximation is used before forming a total category.
\end{dcomment}


La réciproque de \texttt{[la]} prop. 8 est-elle valable, i.e. si $H$ est totalement $W$-asphérique, en est-il de même de $V = U/H$? La difficulté provient du fait que cela signifie que pour \textit{toute} flèche $f : V \longrightarrow V'$ dans Cat de source $V$, on ait $f_!(e_V)$ $W$-asphérique, alors que l'hypothèse équivaut à le postuler seulement pour les flèches qui se factorisent par $U$. Notons que si on ne fait pas cette hypothèse sur $f$, mais qu'on suppose au lieu de cela que $f$ s'insère dans un carré commutatif \msp{page 95}
\[
\begin{tikzcd}[row sep=large]
V \arrow[r, "\alpha"] \arrow[d, "f"'] & U \arrow[d, "p"] \\
V' \arrow[r, "\alpha'"] & U',
\end{tikzcd}
\]
avec $\alpha'$ lui aussi un isomorphisme local, alors on peut conclure que $\upi_0(f) = \upi_0(V/V')$ est bien asphérique. En effet, on a un isomorphisme canonique dans $\Cat/U'$,
\[ \upi_0(\alpha'f) \simeq \upi_0(f), \qquad \text{i.e. } \ \upi_0(V/V') \simeq \upi_0(V/U'), \]
du fait p. ex. que la factorisation canonique de $f$
% foot: Version 7 décembre 2002
\printedpage{67}
% head: A. Grothendieck : Dérivateurs, Ch. XII, Édition des Universités de Montpellier II et Paris VII
\[
\begin{tikzcd}[row sep=normal]
 & V \arrow[dl, "\tilde{f}"'] \arrow[dd, "f"] & \\
\upi_0(V/V') \arrow[dr, "\pi"'] & & \\
 & V' \arrow[dr, dotted, "\alpha'"'] & \\
 & & U'
\end{tikzcd}
\]
fournit une factorisation du même type pour $\alpha'f$, en
\[ \alpha'f = (\alpha'\pi) \circ \tilde{f}. \]
Ceci dit, on peut appliquer l'hypothèse à $p$ et on trouve bien la $W$-asphéricité de $\upi_0(f) \simeq \upi_0(\alpha'f) = \upi_0(p\alpha)$.

\msp{page 96}
L'idée naturelle, pour une réciproque à \texttt{[la]} prop. 8, serait de prendre pour $U'$ la somme amalgamée de $U$ et $V'$ sur $V$, ce qui marchera pourvu qu'on sache que $\alpha'$ est bel et bien un isomorphisme local tout comme $\alpha$. C'est bien le cas p. ex. si $V \longrightarrow U$ est une immersion ouverte, i.e. $H$ un sous-préfaisceau du préfaisceau final. Mais je crains que ce ne soit pas vrai en général - ça a l'air un peu trop beau pour être vrai! C'est en effet déjà faux dans le cas
\[
\begin{tikzcd}[row sep=huge]
V = E_G \arrow[r, "\alpha\ \text{can}"] \arrow[d, "f"'] & B_G = U \\
V' = B_{G'}, &
\end{tikzcd}
\]
dans le cas p. ex. où $G = \{\pm 1\}$, groupe à deux éléments, et où l'élément $\alpha$ de $G'$, image par $f$ de l'une des deux flèches non identiques de $E_G$ est $\neq 1$ \texttt{[plutôt de carré} $\neq 1$\texttt{]}. Si $H'$ est le sous-groupe invariant engendré par $\alpha$ \texttt{[plutôt par} $\alpha^2$\texttt{]} dans $G'$, la somme amalgamée s'identifie à $B_{G'/H'}$, or $B_{G'} \longrightarrow B_{G'/H'}$ n'est \textit{pas} un isomorphisme local (car $(B_{G'})/t_0 \longrightarrow (B_{G'/H'})/t'_0$ n'est pas bijectif sur les objets, s'identifiant à $G' \overset{q}{\longrightarrow} G'/H' \ \ldots$).

Mais d'autre part la condition de totale $W$-asphéricité sur un objet $H$ de $U^{\wedge}$ paraît déjà à tel point exigeante, que je me vois mal con\msp{page 97}struire un exemple où elle est satisfaite, sans que $U/H$ elle-même ne soit totalement $W$-asphérique.

\dgrblock{dgrb525}{The converse for subterminal presheaves and its obstruction}
\begin{dcomment}{exposition}{p96}{The converse for subterminal presheaves and its obstruction}
Proposition~8 compares two quantifiers. Write $V=U/H$ and $\alpha:V\to U$. Total asphericity of $V$ tests every outgoing functor $f:V\to C$. Total asphericity of $H$ tests the special family $p\alpha$, with $p:U\to B$. The proposition follows because $\alpha_!1_V=H$, but reversing it requires a way of placing an arbitrary $f$ in that special family without changing its comprehensive middle category.

The square criterion on p.~95 does this whenever one can complete
\[
\begin{tikzcd}
V\arrow[r,"\alpha"]\arrow[d,"f"']&U\arrow[d,"p"]\\
C\arrow[r,"\beta"']&B
\end{tikzcd}
\]
with $\beta$ a discrete fibration. To justify the assertion completely, factor $f$ as a final functor $V\to E$ followed by a discrete fibration $E\to C$. Since discrete fibrations compose, the same $E$ is the comprehensive middle category of $\beta f=p\alpha$. Its category is therefore $W$-aspherical when $H$ is totally $W$-aspherical.

If $H$ is subterminal, $\alpha$ is a sieve inclusion. Form the ordinary categorical pushout $B=U\amalg_V C$. Theorem~1 preserves the sieve inclusion on the other side, so $\beta:C\to B$ is again a sieve, in particular a discrete fibration. The preceding argument applies to every $f:V\to C$. This supplies the full proof of the source's subterminal case: for such $H$, the total asphericity of $H$ and of $U/H$ are equivalent.

The groupoid example identifies why the same proof cannot be extended merely by replacing ``sieve'' with ``discrete fibration.'' Let $G=C_2$ and let $E_G$ be the indiscrete groupoid on two objects. Its projection to $BG$ sends both nonidentity arrows to the unique nontrivial element $g$. A functor $E_G\to BG'$ sends one of these arrows to an element $a\in G'$ and the inverse arrow to $a^{-1}$. In a cocone on this span, the two objects must become equal, and compatibility with the projection to $BG$ forces
\[
a=a^{-1},\quad\text{equivalently }a^2=1.
\]
There are no additional relations. The pushout is therefore
\[
B\bigl(G'/\langle\!\langle a^2\rangle\!\rangle\bigr).
\]
For example, with $G'=\mathbb Z$ and $a$ its generator, the induced map is $B\mathbb Z\to BC_2$. A discrete fibration between one-object groupoids must lift every arrow uniquely into the sole object, so its map of morphism sets must be bijective. This quotient map is not. The bracketed corrections in the edition, replacing $a$ by $a^2$, are thus mathematically necessary.

This counterexample defeats the proposed pushout method; it does not by itself refute the general converse of Proposition~8. That converse is false: a presheaf with three elements on a monoid with two elements is given in the commentary at ms.~97. The subterminal case is proved, and the obstruction to the broader argument is explicit.
\end{dcomment}


L'existence d'objets totalement $W$-asphériques dans Cat n'est pas tellement évidente a priori. Qu'il y en ait un bon paquet est assuré par la
% foot: Version 7 décembre 2002

% ---- from Der-ChXIIb-012_p026-030.tex
\setcounter{footnote}{65}

\printedpage{68}
% head: A. Grothendieck : Dérivateurs, Ch. XII, Édition des Universités de Montpellier II et Paris VII

\noindent\textbf{Proposition 9.}

\begin{enumerate}[label=\textit{\alph*)}]
\item \textit{Tout objet $V$ de} Cat \textit{ayant un objet final est totalement $W$-asphérique - donc tout préfaisceau d'ensembles $H$} représentable \textit{sur une catégorie $U$ (la représentabilité signifiant justement que $U/H$ a un objet final) est totalement $W$-asphérique.}

\item (\footnote{!! Sans doute \textit{faux} tel quel, cf. rectification page suivante.}) \textit{(Supposons} W(7 bis) \textit{satisfait.) Soient $V$ dans} Cat, \textit{$V'$, $V''$ deux ouverts de réunion $V$. Si $V'$, $V''$, $V'\cap V''$ sont totalement $W$-asphériques, il en est de même de $V$.}
\end{enumerate}

\noindent\textsc{Démonstration} de a). Soit $V$ avec objet final et $f : V \longrightarrow V'$, d'où la factorisation
\[
\begin{tikzcd}[column sep=small, row sep=small]
 & \tilde{V} = \upi_0(V/V') \arrow[dr, "\pi"'] & \\
V \arrow[ur, "\tilde{f}"] \arrow[rr, "f"'] & & V'
\end{tikzcd}
\]
de $f$. On démontre sans mal (\footnote{Voici une démonstration simple : On a pour toute $V$ dans Cat $e_V = \varinjlim_{v \text{ dans } V} v$ (limite inductive du foncteur d'inclusion de $V$ dans $V^\wedge$), car $\forall\, x$ dans $V$
\[
\varinjlim_{v \text{ dans } V} \Hom(x,v) \simeq \pi_0(x\backslash V) = e.
\]
Comme $f_!$ commute aux $\varinjlim$ , on a donc $f_!(e_V) = \varinjlim_{v \text{ dans } V} f(v)$. Si $v_0$ est un objet final de $V$, cela montre que $f_!(e_V) = f(v_0)$, q.e.d.}) (histoire de cofinalité) ceci :

\noindent\textbf{Corollaire.} \textit{L'image par $\tilde{f}$ d'un objet final de $V$ est un objet final de $\tilde{V} = \upi_0(V/V')$.}

Comme un objet de Cat avec objet final est $W$-asphérique par W(3), cela établit a).

Prouvons b). J'ai envie de dire que si on a $f : V \longrightarrow W$ dans Cat, alors $\upi_0(V/W)$ est réunion de deux ouverts isomorphes à $\upi_0(V'/W)$, $\upi_0(V''/W)$, dont l'intersection est isomorphe à $\upi_0(V'\cap V''/W)$. Cela signifierait ni plus ni
\msp{page 98}
moins que ceci : $f_!$ transforme les monomorphismes $V' \longrightarrow V$, $V'' \longrightarrow V$ (interprétés comme flèches dans $V^\wedge$, de but $e_V$) en monomorphismes, et le produit fibré $V'\cap V'' = V' \times_V V''$ en le produit fibré correspondant dans $W^\wedge$. Mais le foncteur $f_!$ commute aux $\varinjlim$ quelconques (étant un adjoint à gauche), et je ne vois aucune raison a priori pour qu'il commute (disons) aux produits fibrés (ce qui impliquerait aussi qu'il transforme monomorphismes en itou). Si je me rappelle bien, c'est ultrafaux en général. Donc je doute fort à présent que l'énoncé b) soit vrai, sans hypothèse supplémentaire. Pour assurer la monomorphie, on songe à imposer l'existence de \textit{rétractions} de $V$ sur $V'$, $V''$ et de $V'$ ou $V''$ sur $V'\cap V''$.

\dgrblock{dgrb526}{Total asphericity tested on final images}
\begin{dcomment}{development}{p97}{Total asphericity tested on final images}
Grothendieck's search becomes more transparent if we express the quantifier through the first half of comprehensive factorization. The following elementary consequence is our formulation of what the construction already permits; it is not a claim of priority. We checked the factorization accounts of Street--Walters, Kelly, and Berger--Kaufmann, which establish the factorization but do not state this characterization of the manuscript's total-asphericity class.

\noindent\textbf{Final-image criterion.}\ours\quad
For any localizer $W$ for which asphericity is invariant under isomorphism, a small category $A$ is totally $W$-aspherical if and only if every ordinary final functor $A\to C$ has $W$-aspherical target. Moreover, the class of totally $W$-aspherical categories is closed under ordinary final functors.

\emph{Proof.} If $u:A\to C$ is final, then $u_!1_A=1_C$. Hence the category of elements used in the definition of total asphericity is $C/1_C\cong C$, proving the necessity. Conversely, factor an arbitrary $f:A\to B$ comprehensively as
\[
A\xrightarrow{j}E=B/(f_!1_A)\xrightarrow{\pi}B.
\]
The first factor is final, so the assumed condition says that $E$ is $W$-aspherical. This is exactly the required test for $f$.

For the closure assertion, suppose $A$ is totally $W$-aspherical and $u:A\to C$ is final. For every $g:C\to D$,
\[
g_!1_C\cong g_!u_!1_A\cong(gu)_!1_A.
\]
Its category of elements is $W$-aspherical by the hypothesis on $A$. Thus $C$ is totally $W$-aspherical.\hfill$\square$

There is an equivalent extremal description which is sometimes useful for constructing examples. Among subclasses of the $W$-aspherical categories which are closed under taking targets of ordinary final functors, the totally $W$-aspherical categories form the largest one. They themselves have that closure property, and they are aspherical by taking $f=\mathrm{id}$. Conversely, if $A$ belongs to any such subclass, every comprehensive middle category of an outgoing map from $A$ belongs to it, hence is aspherical. The criterion just proved applies.

The result also separates two kinds of positive examples. If $A$ has final object $a_0$, then for every $f:A\to B$,
\[
f_!1_A\cong B(-,f(a_0)),\qquad
B/(f_!1_A)\cong B/f(a_0),
\]
whose final object is $(f(a_0),\mathrm{id})$. This is Proposition~9(a) in an explicit form. Filtered categories provide the further examples already explained in the commentary on Chapter~II, manuscript pp.~32--35, through filtered colimits of representables. For instance $(\mathbb N,\leq)$ has no final object, so the conjecture in footnote~73 cannot survive even in ordinary homotopy theory. The new criterion explains what both examples must control: all final images, not just their own classifying spaces.
\end{dcomment}


\dgrblock{dgrb527}{A total presheaf whose category of elements is not total}
\begin{dcomment}{development}{p97}{A total presheaf whose category of elements is not total}
Grothendieck expects that Proposition~8 has no converse outside the subterminal case, and proposes to look for a totally $W$-aspherical presheaf $H$ whose category of elements is not totally $W$-aspherical. Such a presheaf exists on the smallest category that is not a preorder and not a group. We found no source for the example; it is given without a claim of priority.

\textbf{Proposition (the converse of Proposition~8 fails).}\ours
\textit{Let $U$ be the category with one object $*$ and $\operatorname{End}(*)=\{1,e\}$, $e^2=e$. Let $H$ be the presheaf on $U$ with $H(*)=\{i,h_1,h_2\}$, on which $e$ acts by the constant map with value $i$. Then $H$ is totally $W$-aspherical for every fundamental localizer $W$. The category $U/H$ is not totally $W$-aspherical for $W=W_\infty$, for the homology equivalences $W_R$ with coefficients in a nonzero commutative ring $R$, and for the $n$-equivalences with $n\geq2$.}

\textit{Proof.}
Let $R=U(-,*)$ be the representable presheaf. Its elements are $1$ and $e$, and $e$ acts on them by $x\mapsto xe$, the constant map with value $e$. The endomorphism $\varepsilon:x\mapsto ex$ of $R$ is idempotent with image the subpresheaf $Q=\{e\}$. The presheaf $H$ is the pushout $R\amalg_QR$ of two copies of $R$ along $Q$: the element $e$ is $i$, and the two copies of $1$ are $h_1,h_2$.

Let $p:U\to C$ be any functor to a small category, and put $c=p(*)$, $d=p(e)$, so that $d^2=d$. The functor $p_!$ preserves colimits, sends $R$ to $C(-,c)$, and sends $\varepsilon$ to the endomorphism $\alpha\mapsto d\alpha$. Every functor preserves the splitting of an idempotent. Hence $p_!Q$ is the subpresheaf
\[
 Q_d(v)=\{\alpha:v\to c\mid d\alpha=\alpha\}\subseteq C(v,c),
\]
with its inclusion, and
\[
 p_!H\cong C(-,c)\amalg_{Q_d}C(-,c).
\]
This pushout is formed objectwise along injections. The category of elements $E'=C/p_!H$ is therefore the union of two full sieves $S_1,S_2$, each isomorphic to $C/c$, with $S_1\cap S_2=C/Q_d$. A chain of arrows of $E'$ lies in the sieve that contains its last object, so $NE'=NS_1\cup NS_2$ and $NS_1\cap NS_2=N(C/Q_d)$. Each $S_j$ has the terminal object $(c,\mathrm{id}_c)$. The category $C/Q_d$ has the object $(c,d)$, and the arrows $\alpha:(v,\alpha)\to(c,d)$ are the components of a natural transformation from the identity functor to the constant functor at $(c,d)$: for an arrow $g:(v,\alpha)\to(w,\beta)$ one has $\beta g=\alpha$. Thus $BS_1$, $BS_2$ and their intersection are contractible, and so is $BE'$. Since $p$ was arbitrary, $H$ is totally $W_\infty$-aspherical. Every fundamental localizer contains $W_\infty$ \citep[Th\'eor\`eme~2.2.11]{Cisinski2004Minimal}, which gives the first assertion.

The category $E=U/H$ has the three objects $i,h_1,h_2$ and, besides identities, the arrows
\[
 e_i:i\to i,\qquad u_1:i\to h_1,\qquad u_2:i\to h_2,
 \qquad e_i^2=e_i,\quad u_je_i=u_j.
\]
Let $q:E\to\Psi$ be the functor to the outgoing fork $\Psi=(a\leftarrow z\to b)$ of the commentary at ms.~100 which sends $i,h_1,h_2$ to $z,a,b$ and $e_i$ to the identity. It is bijective on objects and full. For an object $y$ of $\Psi$ with preimage $\tilde y$, every object $(x,\,y\to q(x))$ of $y\backslash E$ receives an arrow from $(\tilde y,\mathrm{id}_y)$, obtained by lifting $y\to q(x)$. So $q$ is final. The inclusion $f:\Psi\to C$ into the four-object category of the same commentary is final and $BC\simeq S^2$. The composite $fq:E\to C$ is a final functor whose target is not aspherical for any of the localizers named. By the final-image criterion above, $U/H$ is not totally aspherical for them.\hfill$\square$

The objects $h_1,h_2$ of $U/H$ have no common successor; this is the property that fails, and it cannot be seen by functors that factor through $U$. For the localizers $W_0$ and $W_1$ (bijections on $\pi_0$; bijections on $\pi_0$ and isomorphisms on all $\pi_1$) the converse of Proposition~8 does hold. A final functor induces a bijection on connected components, and a final image of a $1$-connected category is $1$-connected (the argument with coverings at ms.~101). If $H$ is totally $W_0$-aspherical, the test $p=\mathrm{id}$ shows that $U/H$ is nonempty and connected, and then so is every final image of $U/H$; for $W_1$ the same reasoning applies with $1$-connectedness.
\end{dcomment}


Je viens de vérifier la chose, et donne le
% foot: Version 7 décembre 2002
\printedpage{69}
% head: A. Grothendieck : Dérivateurs, Ch. XII, Édition des Universités de Montpellier II et Paris VII
\textbf{Corollaire 1} (rectification de \texttt{[la]} prop. 9, partie b) (\footnote{NB Il y a un énoncé légèrement plus général, en partant d'une catégorie $U$ et d'un diagramme
$\begin{tikzcd}[ampersand replacement=\&, row sep=small]
V'\cap V'' \arrow[r] \arrow[d] \& V' \arrow[d] \\
V'' \arrow[r] \& V
\end{tikzcd}$
\textit{dans} $U^\wedge$ : si $V'$, $V''$, $V'\cap V''$ sont totalement $W$-asphériques, il en est de même de $V$.}) : \textit{L'énoncé est valable si on suppose qu'il existe une} rétraction \textit{$r'_0$ de $V''$ sur $V'\cap V''$, et de même une rétraction $r''_0$ de $V'$ sur $V'\cap V''$.} (\footnote{Corollaire rectifié faux également - les rétractions ne sont pas des flèches de $V^\wedge$, et ne passent dans $p_!$.}) % [resolved: see DER_XIIb-a_corrections.csv]
\[
\begin{tikzcd}
V'\cap V'' \arrow[r, "i''_0"'] \arrow[d, "i'_0"] & V' \arrow[d, "i'"] \arrow[l, dotted, bend right=50, "r''_0"'] \\
V'' \arrow[r, "i''"] \arrow[u, dotted, bend left=50, "r'_0"] & V \arrow[u, dotted, bend right=50, "r'"'] \arrow[l, dotted, bend left=50, "r''"]
\end{tikzcd}
\]

(\textbf{NB} Nous verrons plus tard qu'il suffit de l'existence d'\textit{une} de ces rétractions pour assurer la conclusion.)

\textsc{Démonstration}. Comme $V$ est somme amalgamée de $V'$ et $V''$, il existe une rétraction $r'$ de $V$ sur $V'$ (donnée par $\mathrm{id}_{V'}$ sur $V'$, par $i''_0r'_0$ sur $V''$) qui induit $r'_0$ sur $V''$, et de même il existe
\msp{page 99}
une rétraction $r''$ de $V$ sur $V''$, induisant $r''_0$ sur $V'$. L'existence de ces rétractions assure déjà que pour tout foncteur $f : V \longrightarrow W$, le diagramme transformé de
\[
\begin{tikzcd}[row sep=large]
V'\cap V'' \arrow[r] \arrow[d] & V' \arrow[d] \\
V'' \arrow[r] & V
\end{tikzcd}
\]
(intérprété comme diagramme dans $V^\wedge$) est formé de monomorphismes (\footnote{\textbf{faux,} déjà dans le cas où $V = \Phi$, $V'$, $V''$ ses deux ouverts canoniques isomorphes à $\bDelta^1$, et $V'{\cap}V'' = e$.}) % [sic] intérprété -> interprété, misspelt accent
\[
\begin{tikzcd}[row sep=large]
f_!(V'\cap V'') \arrow[r, hook] \arrow[d, hook] & f_!(V') \arrow[d, hook] \\
f_!(V'') \arrow[r, hook] & f_!(V),
\end{tikzcd}
\]
et de plus que $f_!(V') \cap f_!(V'')$ (qui contient $f_!(V'\cap V'')$) lui est en fait \textit{égal.} Il en résulte que l'on a un diagramme à la fois cartésien et cocartésien d'immersions ouvertes dans Cat
% foot: Version 7 décembre 2002
\printedpage{70}
% head: A. Grothendieck : Dérivateurs, Ch. XII, Édition des Universités de Montpellier II et Paris VII
\[
\begin{tikzcd}[row sep=large]
W/f_!(V'\cap V'') \arrow[r, hook] \arrow[d, hook] & W/f_!(V') \arrow[d, hook] \\
W/f_!(V'') \arrow[r, hook] & W/\Bigl(f_!(V) \overset{\text{déf}}{=} f_!(e_V)\Bigr),
\end{tikzcd}
\]
lequel s'identifie aussi à
\[
\begin{tikzcd}[row sep=large]
\upi_0(V'\cap V''/W) \arrow[r, hook] \arrow[d, hook] & \upi_0(V'/W) \arrow[d, hook] \\
\upi_0(V''/W) \arrow[r, hook] & \underline{\pi_0}(V/W).
\end{tikzcd}
\]

\dgrblock{dgrb528}{Why a categorical retraction cannot be sent through this Kan extension}
\begin{dcomment}{exposition}{p98}{Why a categorical retraction cannot be sent through this Kan extension}
Grothendieck tries to protect a union from collapsing under $f_!$ by splitting its inclusions. Splitting is indeed preserved by every functor, but the retraction must be a morphism in the category on which that functor acts. Here $f_!$ acts on $V^{\wedge}$, whereas the retractions just constructed are functors between full subcategories of $V$.

For a sieve $S\subseteq V$, the associated subterminal presheaf $1_S\subseteq1_V$ takes the value one on objects of $S$ and the empty set elsewhere. If $S$ omits an object of $V$, there cannot be a map $1_V\to1_S$: at that object it would have to be a function from a singleton to the empty set. Thus a proper sieve, even if it admits a categorical retraction, does not define a split subobject of the terminal presheaf. A functor $r:V\to S$ moves the base object; a morphism in $V^{\wedge}$ must stay over it. This is exactly the defect identified by the source's corrective footnote.

What would suffice is a genuinely different hypothesis. Given a union $H=H_1\cup H_2$ in a presheaf category, suppose that for every $p$ the images of its three inclusion maps remain monomorphisms and the square remains Cartesian. Since $p_!$ preserves the pushout, the resulting categories of elements then form a union of two sieves with the expected intersection. The ordinary forward homotopy gluing lemma proves total $W_\infty$-asphericity of $H$ when the other three presheaves are totally $W_\infty$-aspherical. This is the precise exactness condition the proposed proof needs. Categorical retractions provide none of it.
\end{dcomment}


\msp{page 100}
L'hypothèse sur $V'$, $V''$, $V'\cap V''$ implique que les trois sommets du carré autres que $\upi_0(V/W)$ sont $W$-asphériques, donc par W(7 bis) il en est de même de ce dernier, q.e.d.

\textbf{Corollaire 2} (\footnote{Plus généralement, si $I$ est une catégorie \textit{ordonnée finie,} ayant un objet initial, alors}). \textit{Moyennant} W(7 bis)\textit{, la catégorie} $\Psi = \begin{tikzcd}[ampersand replacement=\&, row sep=small, column sep=small, baseline=(B.base)]
 \& a \\
|[alias=B]| 0 \arrow[ur] \arrow[dr] \& \\
 \& b
\end{tikzcd}$ \textit{(ou $0 \longleftarrow 1 \longrightarrow 2$) est totalement $W$-asphérique} (\footnote{\textbf{Faux,} $\Psi$ n'est \textit{pas} totalement $W$-asphérique.}).

On applique le corollaire 1 à $\Psi$ et à ses deux sous-catégories ouvertes, $\Psi/a$, $\Psi/b$ (isomorphes à $\bDelta^1$), et à leur intersection $\{0\} \simeq \bDelta^0 = e$. Comme ces trois catégories ont un objet final (contrairement à $\Psi$ elle-même, qui n'a qu'un objet initial \dots), elles sont totalement $W$-asphériques. Le corollaire s'applique, l'existence des rétractions $r'_0$, $r''_0$ étant ici évidente.

C'est surtout la partie b) (rectifiée) de la page 9 qui assure l'existence de catégories totalement $W$-asphériques autres que celles ayant un objet final (\footnote{Je conjecture à présent que ce sont les seules! (\footnotemark)}\footnotetext{Faux, il y a les catégories filtrantes (si $W$ satisfait l'axiome des limites).}). Il est vrai qu'en sens inverse il faudrait quand même s'assurer qu'un objet $X$ $W$-asphérique de Cat n'est pas \textit{toujours} totalement $W$-asphérique! Le contre-exemple n'est pas tellement évident finalement -
\msp{page 101}
j'ai passé des heures avant d'en trouver un et d'y voir un peu clair! Il suffit en somme d'exhiber un exemple d'une flèche dans Cat
\[
f : X \longrightarrow Y
\]
telle que
% foot: Version 7 décembre 2002
\printedpage{71}
% head: A. Grothendieck : Dérivateurs, Ch. XII, Édition des Universités de Montpellier II et Paris VII
\begin{enumerate}[label=\alph*)]
\item $\underline{\pi_0}(f) \overset{\text{déf}}{=} \upi_0(X/Y) \xrightarrow{\ \sim\ } Y$, i.e. les $\pi_0(y\backslash X)$ \texttt{[plutôt les} $y\backslash X$\texttt{]} sont 0-connexes (pour $y \in \Ob Y$),
\end{enumerate}

\dgrblock{dgrb529}{The finite fork counterexample, computed completely}
\begin{dcomment}{examples}{p100}{The finite fork counterexample, computed completely}
Grothendieck eventually finds a very small witness to the failure of the preceding gluing claims. The construction is in his Appendix to Chapter~XII, manuscript pp.~10--12, where the final corrective footnote identifies the homotopy type as $S^2$ \citep{GrothendieckDerivateursXIIAppendix}. We work it out here because it explains at once the failure of total asphericity and the failed use of retractions.

Let $C$ have four objects $x,z,a,b$ and the nonidentity morphisms generated by
\[
\begin{tikzcd}[column sep=large]
x\arrow[r,shift left=.5ex,"\alpha"]\arrow[r,shift right=.5ex,"\beta"']&z\arrow[r,"r"]\arrow[d,"s"']&a\\
&b&
\end{tikzcd}
\qquad r\alpha=r\beta,\quad s\alpha=s\beta.
\]
Besides the four displayed arrows there are just the two composites $x\to a$ and $x\to b$. Let $A$ be its full subcategory on $z,a,b$, the outgoing fork $\Psi$, and let $f:A\hookrightarrow C$ be the inclusion. The category $A$ has initial object $z$, so $BA$ is contractible.

We first compute $f_!1_A$. At $a$ and $b$ the outgoing comma category is a point. At $z$ it has initial object $(z,\mathrm{id})$. At $x$ it has four objects
\[
(z,\alpha),\ (z,\beta),\ (a,r\alpha),\ (b,s\alpha),
\]
and one arrow from each of the first two to each of the last two. Its nerve is a four-edge circle. In particular it is nonempty and connected. Thus every $c\downarrow f$ is connected, and
\[
f_!1_A=1_C,\qquad \underline\pi_0(A/C)=C.
\]
The map $f$ is final in the ordinary sense, although its comma category at $x$ is not contractible.

The homotopy type of $C$ can be read without a spectral sequence. Let $K$ be the full subcategory on $x,z$, with its two parallel arrows. Its nerve realizes as a circle. The full subcategories $C_a$ on $x,z,a$ and $C_b$ on $x,z,b$ each adjoin a final object to $K$. Their nerves are therefore cones on $NK$. They cover $NC$, with intersection $NK$, because there are no arrows between $a$ and $b$. Consequently
\[
BC=BC_a\cup_{BK}BC_b\cong\Sigma BK\simeq S^2.
\]
Here the union description is literal at the level of nerves. Hence $A$ is not totally $W_\infty$-aspherical: its direct image along $f$ has category of elements $C$, with $H_2(BC;\mathbb Z)\cong\mathbb Z$.

Now return to the proposed retraction repair. Inside $A$, put $A_a=\{z\to a\}$ and $A_b=\{z\to b\}$. Both have final objects, their intersection is $\{z\}$, and both admit retractions to that intersection. Proposition~9(a) therefore gives total asphericity of all three pieces. Nevertheless their union is $A$, which we have just shown is not total. The false step can be inspected at a single value. Write $y_c=C(-,c)$. Since left Kan extension carries a representable to the corresponding representable,
\[
f_!(1_{\{z\}})=y_z,\quad
f_!(1_{A_a})=y_a,\quad
f_!(1_{A_b})=y_b.
\]
The equality $1_{A_a}=A(-,a)$ follows because $a$ receives exactly one arrow from $z$ and none from $b$, and similarly for $b$. Evaluating the first inclusion after Kan extension at $x$ yields
\[
\{\alpha,\beta\}=C(x,z)\longrightarrow C(x,a)=\{r\alpha=r\beta\}.
\]
It is not injective. The other leg has the same defect. Their pushout is the singleton $1_C(x)$, exactly as preservation of colimits predicts. The putative intersection square has failed before any homotopy argument is applied.

The same finite construction also explains the two kinds of pushout failure announced on p.~81. Adjoin a final object to $A$, giving $A^{\triangleright}$. Pushing out $A\subset A^{\triangleright}$ along $f$ adjoins a final object to $C$, because $f_!1_A=1_C$. The upper inclusion is an aspherical weak equivalence; the lower inclusion $C\to C^{\triangleright}$ is not, since $BC\simeq S^2$.

To make the other leg a weak equivalence as well, replace the target of $f$ by $C^{\triangleright}$, with final point $c_+$. The composite $p:A\to C^{\triangleright}$ is a weak equivalence between contractible categories. Its direct image of $1_A$ is the subterminal presheaf supported on $C$: it is one on every object of $C$ and empty at $c_+$. The resulting pushout $P$ is $C$ with two distinct final cone points, $c_+$ and the newly attached point. Hence
\[
BP\simeq\Sigma BC\simeq S^3.
\]
Both $A\to A^{\triangleright}$ and $p$ are weak equivalences, while neither map from $C^{\triangleright}$ or $A^{\triangleright}$ to $P$ is. This proof uses the explicit nerve, so it does not depend on the invalid reverse implication of the one-point corollary on p.~85.
\end{dcomment}


\dgrblock{dgrb530}{Testing the retracted assertions under their stated axiom}
\begin{dcomment}{development}{p100}{Testing the retracted assertions under their stated axiom}
The fork example proves that ordinary weak contractibility is insufficient for total asphericity. There is a logical point still to check: several of the source's retracted conclusions explicitly assume W(7 bis), whereas $W_\infty$ does not satisfy that axiom. To refute those conditional conclusions themselves, rather than only their proposed proof, we need an example in a weak equivalence class for which the axiom holds. Fixed-coefficient homology provides such a test. The homological tools are standard; the application below is our explicit check of these conditional assertions, with no claim of priority.

\textbf{A homological counterexample with W(7).}\ours\quad
Let $R$ be a nonzero commutative ring, and let $W_R$ consist of functors inducing isomorphisms on all homology groups of their nerves with coefficients in $R$. Then $W_R$ is a fundamental localizer satisfying the full W(7). Nevertheless a union of two totally $W_R$-aspherical sieves with totally $W_R$-aspherical intersection need not be totally $W_R$-aspherical, even when both pieces retract onto their intersection.

\emph{Verification of the test class.} Identities, two-out-of-three, and closure under retracts follow by applying homology. A category with a terminal object has contractible nerve and hence maps to the terminal category by an arrow in $W_R$. We spell out the relative Theorem~A axiom, since that is the remaining condition needed here. Suppose $A\xrightarrow{u}B\to C$ is a triangle and every induced map $A/c\to B/c$ is an $R$-homology equivalence. For the composite $f:A\to C$, form the bisimplicial set with $(p,q)$-simplices
\[
 \bigl(a_0\to\cdots\to a_p,\quad
       f(a_p)\to c_0\to\cdots\to c_q\bigr).
\]
Deleting the last $a$-vertex composes the crossing arrow with the image of that last arrow under $f$; deleting the first $c$-vertex composes on its other side. These rules, with the ordinary remaining faces and identity degeneracies, make the two simplicial directions commute. Take free $R$-modules, normalize in each direction, and use the usual sign in the total differential.

For each fixed $a$-chain, the second direction is the nerve of $f(a_p)\backslash C$, a category with initial object $\id_{f(a_p)}$. Its positive homology vanishes and its degree-zero homology is $R$. The spectral sequence that first takes this homology therefore identifies the homology of the total complex with $H_*(NA;R)$. In the other direction, for each fixed chain $c_0\to\cdots\to c_q$, the first simplicial direction is $N(A/c_0)$. Consequently the map induced by $u$ gives an isomorphism on that first homology page by the hypothesis on the comma categories. The comparison of the first-quadrant spectral sequences gives an isomorphism of total homology, and thus $H_*(NA;R)\cong H_*(NB;R)$. Convergence is sufficient here without any finiteness assumption: there are only finitely many bidegrees in any given total degree. This proves the relative localizer axiom. Finally the normalized-chain argument after ms.~79 proves W(7) for precisely this class $W_R$.

\emph{The counterexample.} Use the finite functor $f:\Psi\to C$ computed in the preceding block, with $\Psi=(z\to a,z\to b)$. It is ordinary final, so $f_!e_\Psi=e_C$. But $BC\simeq S^2$, and
\[
 H_2(BC;R)=R\ne0.
\]
Hence $\Psi$ is not totally $W_R$-aspherical. Its two full sieves $\{z\to a\}$ and $\{z\to b\}$, and their intersection $\{z\}$, each have a terminal object. Proposition~9(a), whose proof does not use W(7), makes all three totally $W_R$-aspherical. The constant functors from the two arrows to $\{z\}$ are retractions of the inclusions of the intersection. Thus both the union assertion and its proposed retraction repair fail, and the conclusion asserting totality of the outgoing fork fails as well. All this occurs in a fundamental localizer satisfying W(7), so the conditional qualifications in those statements do not save them.\hfill$\square$

The same calculation keeps two different defects apart. The proof's use of left Kan extension to preserve a monomorphism is invalid before any weak equivalence class is chosen. The example above additionally shows that the resulting total-asphericity conclusion can be false after all its stated homotopical axioms have been imposed.
\end{dcomment}


et telle que de plus

\begin{enumerate}[label=\alph*), start=2]
\item $X$ est $W$-asphérique, et $Y$ non.
\end{enumerate}

% a vertical bar is printed in the left margin beside the next paragraph (lines "La difficulté ..." to "... le foncteur")
La difficulté est la suivante, quand on cherche des contre-exemples `idiots' : c'est que l'hypothèse a) implique déjà, quand $X$ est $W_\infty$-asphérique (ou seulement 1-connexe) que $Y$ est 1-connexe : le foncteur
\[
Y' \longmapsto X' = X \times_Y Y'
\]
de la catégorie des revêtements étales $Y'$ de $Y$ vers celle des revêtements étales $X'$ de $X$ est pleinement fidèle (vérification immédiate), d'où le résultat annoncé. Donc un contre-exemple doit faire appel à l'existence d'un $\pi_i(Y,y)$ non nul, avec $i \geq 2$. Cela implique
\msp{page 102}
qu'il n'y a pas de contre-exemple avec des catégories de la forme $Y = Z/F$ (déduites d'un foncteur $q : X \longrightarrow Z$ en prenant $F = q_!(e_X) = \pi_0(X/Z)$), où $Z$ est une catégorie telle que $\bDelta^1$ (?), $\Phi$, $\Psi$, ou $\begin{tikzcd}[ampersand replacement=\&, sep=small, baseline=(M.base)]
 \& {} \& \\
{} \arrow[ur] \arrow[dr] \& |[alias=M]| \phantom{x} \& {} \arrow[ul] \arrow[dl] \\
 \& {} \&
\end{tikzcd}$ (ayant le type d'homotopie du cercle) etc., car les composantes connexes de telles catégories $Y$ sont des $K(\pi,1)$ - les $\pi_i$ supérieures ($i \geq 2$) sont nuls.

\dgrblock{dgrb531}{Why the first obstruction occurs above the fundamental group}
\begin{dcomment}{exposition}{p101}{Why the first obstruction occurs above the fundamental group}
Grothendieck's assertion about coverings explains why a counterexample with a circle as target cannot work. Let $f:A\to B$ be final, so every $b\downarrow f$ is nonempty and connected. Let $L,M:B\to\mathbf{Set}$ be local systems, meaning that every arrow is sent to a bijection. Given a natural transformation $\theta:Lf\to Mf$, choose any object $(a,\beta:b\to f(a))$ of $b\downarrow f$ and define
\[
\bar\theta_b=M(\beta)^{-1}\theta_a L(\beta).
\]
If $h:a\to a'$ is a morphism in that comma category, then $f(h)\beta=\beta'$ and the naturality of $\theta$ identifies the expressions obtained from $(a,\beta)$ and $(a',\beta')$. Connectedness makes the expression independent of all choices. For $\gamma:b'\to b$, the object $(a,\beta\gamma)$ computes $\bar\theta_{b'}$, proving naturality in $b$. Its restriction to $A$ is $\theta$, using $(a,\mathrm{id})$, and the same formula proves uniqueness. Thus restriction along $f$ is fully faithful on local systems.

Finality also preserves the set of connected components: it preserves the colimit of the constant singleton diagram, which is $\pi_0$. Suppose now that $BA$ is simply connected. Then $BB$ is connected, and every pulled-back local system on $A$ is constant. Fix a fibre of a local system $L$ on $B$. Each element of that fibre determines a section of the pulled-back system; full faithfulness, applied to maps from the singleton local system, makes it descend uniquely to a global section of $L$. Every monodromy action is therefore trivial. Applying this to the regular action of $\pi_1(BB)$ shows that the group itself is trivial. This supplies the omitted verification on p.~101.

The four-object example above is consistent with this restriction: $BC\simeq S^2$ is simply connected but not contractible. Its critical comma category is a circle, and taking $\pi_0$ discards exactly that loop before rebuilding the total category. By contrast, if the target category is a discrete fibration over one of the height-one categories drawn on p.~102, it has no composable nonidentity arrows. Its nerve is a graph. Each connected component is a $K(\pi,1)$, so simple connectivity would force it to be contractible. The source's exclusion of these examples is thus correct for the displayed bases; it is not a statement about arbitrary categories merely sharing their homotopy type.

In the modern cofinality theorem, connectedness of the comma categories is strengthened to their weak contractibility \citep[Theorem~4.1.3.1]{Lurie2009HTT}. That stronger condition controls all homotopy colimits and implies weak equivalence of classifying spaces. The ordinary finality used here controls ordinary colimits and already constrains coverings, but it leaves the higher homotopy obstruction exposed.
\end{dcomment}


Cette difficulté vue, le contre-exemple devient évident. Je pars d'un objet 1-connexe \textit{quelconque} $Y$ de $X$ \texttt{[plutôt de} Cat\texttt{]}, et je prends $y_0 \in \Ob Y$,
\[
X = \uCh(Y; y_0, -) \xrightarrow{\ f\ } Y.
\]
On sait que $X$ est $W_\infty$-asphérique, et que $f$ est $W_\infty$-parfait, donc les inclusions
\[
X_y \longrightarrow y\backslash X
\]
sont des $W_\infty$-équivalences. L'hypothèse a) que les $y\backslash X$ soient 0-connexes équivaut donc à la même hypothèse sur les fibres $X_y$, ou
\msp{page 103}
encore sur une seule fibre ($f$ étant parfait, donc une $W_\infty$-fibration, et $Y$ 0-connexe), soit
\[
X_{y_0} \simeq \Omega(Y, y_0).
\]
Sa 0-connexité équivaut à l'hypothèse de 1-connexité de $Y$ (on n'a utilisé encore que celle de 0-connexité). Donc pour un contre-exemple, il suffit de prendre $Y$ 1-connexe et non $\infty$-connexe. Le cas le plus simple est celui où $Y$ a le type d'homotopie de la sphère $S^2$, réalisé p. ex. par la catégorie ordonnée ci-contre \texttt{[ci-dessous]} à 6 sommets.
% foot: Version 7 décembre 2002

\printedpage{72}
% head: A. Grothendieck : Dérivateurs, Ch. XII, Édition des Universités de Montpellier II et Paris VII
\[
\begin{tikzcd}[row sep=small]
e_2^+ \bullet & \bullet\, e_2^- \\
e_1^+ \bullet \arrow[ur] & \bullet\, e_1^- \arrow[ul] \\
e_0^+ \bullet \arrow[ur] & \bullet\, e_0^- \arrow[ul]
\end{tikzcd}
\]
\texttt{[de plus,} $e_0^+ \leq e_1^+ \leq e_2^+$\texttt{,} $e_0^- \leq e_1^- \leq e_2^-$\texttt{]} (\footnote{C'est en somme cet exemple (à homotopie topologique près) que je réalise (sans d'abord m'en apercevoir dans l'Appendice p. 10-12, mais avec des catégories finies).}). Je résume le résultat obtenu.

\textbf{Proposition 10.} \textit{Soit $Y$ une petite catégorie} 1\textit{-connexe ($y_0 \in \Ob Y$). Alors} $\uCh(Y; y_0, -)$ \textit{est $W_\infty$-asphérique, mais n'est pas totalement $W_\infty$-asphérique si $Y$ elle-même n'est $\infty$-connexe, i.e. $W_\infty$-asphérique.} (\footnote{Mais un exemple plus percutant est $\Psi = \begin{tikzcd}[ampersand replacement=\&, row sep=small, column sep=small, baseline=(A.base)]
 \& b \\
|[alias=A]| a \arrow[ur] \arrow[dr] \& \\
 \& c
\end{tikzcd}$ , qui est $W_\infty$-asphérique, mais pas totalement!}).

On peut se demander si $Y$ $W_\infty$-asphérique implique que $\uCh(Y; y_0, -)$ est totalement $W_\infty$-asphérique. J'en doute fort!

\msp{page 104}
Revenons maintenant à la situation du début de la section, diagramme $(*)$, p.~87,
\[
\begin{tikzcd}
U \arrow[r, hook, "i"] \arrow[d, "p"'] & X \arrow[d, "q"'] & T = \{t\} \arrow[l, hook', dotted, "j"'] \arrow[d, dotted, "\mathrm{iso}"] \\
U' \arrow[r, hook, "i'"] & X' & T' = \{t'\}. \arrow[l, hook', dotted, "j'"']
\end{tikzcd}
\tag*{$(*)$}
\]
Nous pouvons donner maintenant un exemple où $i$ est dans $W$ (et même $W$-asphérique) sans que $i'$ le soit :
\[
i \in W, \qquad i' \notin W. \tag*{$(**)$}
\]
On peut même supposer que $t$ soit objet final de $X$ (ce qui signifie que l'objet $H$ de $U^\wedge$ est isomorphe à $e_U$, l'objet final), et $t'$ objet final de $X'$ (i.e. $p_!(e_U) = e_{U'}$, i.e. les $u'\backslash U$ sont 0-asphériques). Dans ce cas, comme $X$, $X'$ sont $W$-asphériques, les relations $(**)$ équivalent à
\[
U \text{ est } W\text{-asphérique, et } U' \text{ non,} \tag*{$(***)$}
\]
ce qui est justement la situation que nous venons de construire (p. 101-103).

\dgrblock{dgrb532}{The information lost by taking components}
\begin{dcomment}{exposition}{p103}{The information lost by taking components}
Grothendieck has found the obstruction that the preceding attempts were missing: an ordinary direct image of a set-valued presheaf remembers connected components of comma categories, whereas a weak equivalence of nerves must preserve all homotopy groups. The path construction makes the difference particularly visible. Put $P=\underline{\mathrm{Ch}}(Y;y_0,-)$ and let $f:P\to Y$ be its endpoint functor. The earlier path theory supplies the two facts used here: $BP$ is contractible and the homotopy fibre of $Bf$ over $y$ is a based path space from $y_0$ to $y$. We write $BA=|NA|$ for the classifying space of a category. The comparison of strict fibres with comma categories is the additional consequence of the perfectness hypothesis; contractibility of $BP$ alone would not supply that comparison.

For connected $Y$, the components of each such path space form a torsor for $\pi_1(BY,y_0)$. Thus simple connectivity gives
\[
 (f_!e_P)(y)=\pi_0(y\backslash P)=\{*\}.
\]
Consequently $f_!e_P=e_Y$, and the category of elements of this direct image is $Y$ itself. If $BY$ is not contractible, this single functor $f$ disproves total $W_\infty$-asphericity of $P$. This argument uses precisely the universal quantifier in Definition~1: the homotopy type of $P$ cannot determine how all its ordinary direct images behave.

The six-element ordered set displayed in the text is a useful finite model of the required target. Its elements occur in three levels $\{e_k^+,e_k^-\}$, $k=0,1,2$, and every element of a lower level precedes every element of a higher level. Its order complex is therefore
\[
 S^0*S^0*S^0\cong S^2.
\]
There are six vertices, twelve edges, and eight triangular faces. Each simplex chooses at most one vertex in each level, so these are exactly the faces of the boundary of an octahedron. The two chains supplied in the editors' brackets are essential: the diagonal arrows alone would not specify this ordered set.

The finite fork construction in the commentary at ms.~100 supplies the same failure without a path category: a contractible $\Psi$ admits a functor $\Psi\to C$ with connected undercategories and $BC\simeq S^2$. We shall use it below. The question about $\underline{\mathrm{Ch}}(Y;y_0,-)$ when $BY$ is contractible is a different question. Proposition~10 gives no implication in that case, and a replacement of the path category by some equivalent contractible category would not answer it. The next commentary answers it, negatively, with this same fork.
\end{dcomment}


\dgrblock{dgrb533}{The path category over a contractible base need not be total}
\begin{dcomment}{development}{p103}{The path category over a contractible base need not be total}
The question that closes this page, whether $\underline{\mathrm{Ch}}(Y;y_0,-)$ is totally $W_\infty$-aspherical when $Y$ is $W_\infty$-aspherical, has a negative answer, and the witness is the category that Grothendieck names in the footnote to Proposition~10.

\textbf{Proposition.}\ours
\textit{Let $Y=\Psi$ be the outgoing fork $a\leftarrow z\to b$ and $y_0=z$. Then $BY$ is contractible, and $P=\underline{\mathrm{Ch}}(Y;z,-)$ is $W_\infty$-aspherical and not totally $W_\infty$-aspherical. It is not totally $W_R$-aspherical for any nonzero commutative ring $R$ either.}

\textit{Proof.}
The object $z$ is initial, so $BY$ is contractible; in particular $Y$ is $1$-connected. Proposition~10 applies: $P$ is $W_\infty$-aspherical, and, as recalled in the preceding commentary, the endpoint functor $f:P\to Y$ satisfies $f_!e_P=e_Y$, that is, every $y\backslash P$ is nonempty and connected. So $f$ is final. The inclusion $g:\Psi\to C$ of the commentary at ms.~100 is final, and $BC\simeq S^2$. Final functors compose, because $(gf)_!e_P=g_!f_!e_P=g_!e_Y=e_C$. The category of elements of $(gf)_!e_P$ is $C$, which is neither $W_\infty$-aspherical nor $W_R$-aspherical.\hfill$\square$

The proof uses only two facts: the endpoint functor of the path category over a $1$-connected base is final, and total asphericity passes to final images. Hence $\underline{\mathrm{Ch}}(Y;y_0,-)$ is totally $W_\infty$-aspherical only if $Y$ is. Since total asphericity of $Y$ is far stronger than asphericity (for a category with finitely many objects it means that $Y$ is filtered; see the commentary on Chapter~II, ms.~36), the path construction does not produce totally aspherical categories from aspherical ones.
\end{dcomment}


On aimerait aussi trouver un exemple où dans $(*)$ on ait $p \in W$, $q \notin W$.
% foot: Version 7 décembre 2002

% ---- from Der-ChXIIb-012_p031-035.tex
\setcounter{footnote}{76}

\printedpage{73}
% head: A. Grothendieck : Dérivateurs, Ch. XII, Édition des Universités de Montpellier II et Paris VII
% foot: Version 7 décembre 2002

\msp{page 105}

Je prendrai encore $(X,U,\{t\})$ tel que $t$ soit objet final de $X$, mais si on veut $q\notin W$, on ne peut supposer $t'$ élément final de $X'$. Au lieu de cela, je vais supposer

\begin{enumerate}[label=\alph*)]
\item $U$, $U'$ (et aussi $X$, ayant objet final $t$) $W$-asphériques, donc $i\in W$ (et même $i$ $W$-asphérique), $p\in W$, mais
\item $X'$ non $W$-asphérique, donc $i',q\notin W$.
\end{enumerate}

Pour ceci, notons que
\[ U'/t'\simeq U'/p_!(e_U)\overset{\text{déf}}{=}\underline{\pi}_0(U/U'). \]
Le corollaire de la proposition 7 (cf. p. 85), non encore démontré, nous dit que $i'\in W$ (i.e. ici, $X'$ $W$-asphérique) si et seulement si $i'$ est $W$-asphérique, i.e. $U'/t'$ est $W$-asphérique. Nous allons l'admettre ici. Tout revient donc à trouver une flèche
\[ p:U\longrightarrow U' \]
dans Cat avec les propriétés suivantes.

\begin{enumerate}[label=\alph*)]
\item $U$ et $U'$ sont $W$-asphériques,
\item $\underline{\pi}_0(U/U')$ n'est pas $W$-asphérique.
\end{enumerate}

Nous allons reprendre l'exemple précédent
\[ f:X\longrightarrow Y \]
(p. 101-103), en prenant
\[ U=X,\qquad \underline{\pi}_0(U/U')=Y, \]

\msp{page 106}

donc tout revient, à partir de l'exemple donné $X\xrightarrow{f}Y$, à construire une flèche
\[ \pi:Y\longrightarrow U' \tag*{$(*)$} \]
dans Cat qui satisfasse les conditions

\begin{enumerate}[label=\alph*)]
\item $\pi$ est un isomorphisme local,
\item $U'$ est $W$-asphérique.
\end{enumerate}

On prendra alors
\[ p=\pi f. \]
Mais par la façon dont $X$ a été construit à partir d'un $Y$ 1-connexe et non $\infty$-connexe, par ailleurs quelconque, on est ramené finalement à trouver $(*)$ satisfaisant a), b), et de plus

\begin{enumerate}[label=\alph*), start=3]
\item \printedpage{74}
% head: A. Grothendieck : Dérivateurs, Ch. XII, Édition des Universités de Montpellier II et Paris VII
% foot: Version 7 décembre 2002
$Y$ est 1-connexe, mais non $\infty$-connexe.
\end{enumerate}

(NB On suppose maintenant $W=W_\infty$.) Il revient au même de trouver

% printed with a vertical rule at the left of this block
\begin{quote}
$U'$ objet \textit{asphérique} de Cat, $F\in\operatorname{Ob}U'^\wedge$ tels que $U'/F$ soit 1-connexe et non $\infty$-connexe.
\end{quote}

On songe par exemple à prendre simplement une sous-catégorie ouverte $U'_0$ de $U'$, qui soit 1-connexe et non $\infty$-connexe. Partons d'une telle catégorie $U'_0$ - par exemple une réalisation homotopique dans Cat, de la sphère $S^2$. On l'immerge comme ouvert dans un $U'$ asphérique, en lui ajoutant un objet final. On a gagné.

\msp{page 107}

Faisons un énoncé en forme.

\textbf{Proposition 11.} \textit{On peut trouver une situation} $(*)$ (p. 124 \texttt{[plutôt 104]}) \textit{avec}

\begin{enumerate}[label=\textit{\alph*)}]
\item \textit{$t$ objet final de $X$.}
\item \textit{$U$, $U'$ (et $X$) $W_\infty$-asphériques, donc $i,p\in W_\infty$ (et même $i$ $W_\infty$-asphérique).}
\item \textit{$X'$ non $W_\infty$-asphérique, donc $i',q\notin W_\infty$.}
\end{enumerate}

Il suffit pour cela de partir d'un objet 1-connexe $Y$ de Cat, de lui adjoindre un objet final $u'_0$, ce qui donne $U'$, de prendre $U=\underline{\mathrm{Ch}}(Y;y_0,-)$, $p:U\longrightarrow U'$ le composé $U\xrightarrow{\text{but}}Y\underset{\text{can.}}{\overset{\text{incl.}}{\hra}}U'$, et $X$ déduit de $U$ en lui ajoutant un objet final $t$. On a alors $U'/t'\simeq Y$, d'où
\[ X'\ W_\infty\text{-asphérique (i.e. }i'\in W_\infty)\Longleftrightarrow Y\ W_\infty\text{-asphérique.} \]
On a le contre-exemple cherché en prenant $Y$ (1-connexe mais) non $\infty$-connexe.

\dgrblock{dgrb534}{Two cones and the two failures of pushout invariance}
\begin{dcomment}{examples}{p107:On a le contre-exemple cherché en prenant}{Two cones and the two failures of pushout invariance}
Grothendieck now converts the failure of total asphericity into failures of ordinary categorical pushouts. The two constructions deserve to be separated because they test different edges of the square. We can make both finite by using the functor $f:\Psi\to C$ constructed at ms.~100. Here $\Psi=(z\to a,z\to b)$ has an initial object, $BC\simeq S^2$, and every category $c\backslash\Psi$ is connected and nonempty. Hence
\[
 f_!e_\Psi=e_C.
\]
The crucial point is that this is an equality of ordinary presheaves; $f$ is not a weak equivalence.

For the first square take $U=\Psi$, $U'=C$, and let $X=\Psi^\triangleright$ be obtained by adjoining a new terminal object $t$. The gluing weight for $U\subset X$ is $e_\Psi$. The pushout formula proved in \S1 says that $X'=X\amalg_\Psi C$ has gluing weight $f_!e_\Psi=e_C$. Thus $X'=C^\triangleright$, with its new terminal object $t'$. Both $X$ and $X'$ are contractible. The inclusion $i:\Psi\to\Psi^\triangleright$ is a weak equivalence and is even $W_\infty$-aspherical: its slices over old objects have terminal objects and its slice over $t$ is $\Psi$. But $i':C\to C^\triangleright$ is not a weak equivalence, because its source has the homotopy type of $S^2$. This proves the assertion at ms.~104 without any converse to the gluing lemma.

For the second square first adjoin a terminal object $s$ to $C$, and write $U'=C^\triangleright$. Let $j:C\hookrightarrow U'$ be the resulting sieve and put $p=jf$. Now $U$, $U'$ and $X=U^\triangleright$ are all contractible, so both $i$ and $p$ are weak equivalences. Since $j$ is a sieve, extension by $j_!$ sends $e_C$ to the subterminal presheaf which is a singleton on $C$ and empty at $s$. Indeed, $c\backslash C$ has an initial object for $c\in C$, whereas $s\backslash C$ is empty. Therefore
\[
 p_!e_U=j_!e_C,
\qquad
 (p_!e_U)(c)=\{*\},\quad (p_!e_U)(s)=\varnothing.
\]
The pushout $X'=X\amalg_U U'$ is consequently obtained from $C$ by adjoining two distinct objects $s,t'$ above every object of $C$, with no arrow between $s$ and $t'$. It is the union of two terminal cones on $C$.

This description computes the homotopy type directly. Every composable string in $X'$ lies in one of the two cones, because a string cannot contain both of their incomparable apexes. Hence
\[
 NX'=N(C^\triangleright_s)
       \amalg_{NC}N(C^\triangleright_{t'}).
\]
The two inclusions of $NC$ are simplicial monomorphisms, so this ordinary union is a homotopy pushout. Its realization is the unreduced suspension of $BC$, and therefore
\[
 BX'\simeq\Sigma BC\simeq S^3.
\]
Neither $i':U'\to X'$ nor $q:X\to X'$ is a weak equivalence. This proves Proposition~11 at ms.~107 and shows exactly how an arbitrary sieve can destroy pushout invariance even when the other three categories are contractible.

The text invokes the converse in Proposition~7 to justify the second construction. That converse depends on the stronger gluing axiom discussed at ms.~78--85 and cannot be used as an unconditional theorem for $W_\infty$. The suspension computation repairs this use of it. More generally, for the simply connected $Y$ of the manuscript, $\Sigma BY$ is contractible exactly when $BY$ is contractible: suspension shifts reduced integral homology, and a simply connected acyclic CW complex is contractible by the Hurewicz and Whitehead theorems. The hypothesis of simple connectivity is doing real work in this last equivalence.
\end{dcomment}


\subsection*{2.4\quad Application aux immersions ouvertes dwyériennes.}

Donnons quand même l'enoncé \textit{positif} auquel nous en avions. % [sic] l'enoncé -> l'énoncé, missing accent

\textbf{Théorème 3.} \textit{Soient $(X,U,T)$ un triple $\in\mathcal{T}$ ($U$ ouvert de $X$ dans} Cat\textit{, $T=X\backslash U$), et pour toute flèche $p:U\longrightarrow U'$ flèche de} Cat\textit{, le carré cocartésien dans} Cat

\[
\begin{tikzcd}
U \arrow[r, hook, "i"] \arrow[d, "p"'] & X \arrow[d, "q"'] & T \arrow[l, dotted, hook'] \arrow[d, dotted, "\text{iso}"] \\
U' \arrow[r, hook, "i'"] & X' & T'. \arrow[l, dotted, hook']
\end{tikzcd}
\]

\printedpage{75}
% head: A. Grothendieck : Dérivateurs, Ch. XII, Édition des Universités de Montpellier II et Paris VII
% foot: Version 7 décembre 2002

\msp{page 108}

\textit{Supposons que pour tout $t\in\dot{T}=\dot{U}$, le préfaisceau d'ensembles}
\[ i^*(t)=(u\longmapsto\operatorname{Hom}(u,t))\in U^\wedge \]
\textit{soit totalement $W$-asphérique (déf. 1, p. 88)} (\footnote{\textbf{NB} On dit sous ces conditions que $i_0:U\longrightarrow\bar{U}$ est \textit{strictement $W$-asphérique} (cf. déf. 2 page suivante).}) \textit{- p. ex. que la catégorie $U/t$ soit totalement $W$-asphérique (prop. 8, page 93). (NB le cas où les $U/t$ ont un objet final, cf. prop. 9 a), i.e. celui où l'inclusion $U\longrightarrow\bar{U}$ admet un adjoint à droite, est celui envisagé par Dwyer, celui des immersions ouvertes de Dwyer} (\footnote{Je ne connais pas d'autre cas où th. 3 s'applique - faute de connaître de catégorie totalement $W$-asphérique qui n'ait un objet final, ni d'objet $F$ d'une catégorie $U^\wedge$ qui soit totalement $W$-asphérique sans être représentable, i.e. $U/F$ a un objet final. À présent, je sais que les catégories \textit{filtrantes} sont totalement $W$-asphériques si $W$ satisfait l'axiome des limites inductives.})\textit{.) Alors pour toute $p:U\longrightarrow U'$, l'immersion ouverte}
\[ i'_0:U'\longrightarrow\overline{U'} \]
\textit{est $W$-asphérique (et même, tout comme $i_0$, strictement $W$-asphérique), et a fortiori elle est} \texttt{[dans]} \textit{$W$. Par suite (prop. 6} \texttt{[p. 81]}\textit{) on a ceci.}

\[
\begin{aligned}
&\text{\textcircled{a}}\qquad && i\in W && \Longleftrightarrow\quad && i'\in W\quad && \textit{moyennant } \mathrm{W}(7\text{ bis}).\\
&\text{\textcircled{b}}\qquad && p\in W && \Longleftrightarrow\quad && q\in W\quad && \textit{moyennant } \mathrm{W}(7).\ (\footnotemark)
\end{aligned}
\]
\footnotetext{Démonstration page 112.}

\dgrblock{dgrb535}{A condition on weights that survives every direct image}
\begin{dcomment}{exposition}{p108:Démonstration page 112.}{A condition on weights that survives every direct image}
After the counterexamples, Grothendieck seeks a hypothesis expressed in the actual gluing data. For a sieve $i:U\hookrightarrow X$, write $\overline U$ for the cosieve generated by $U$, and put $\dot T=\overline U\setminus U$. An object $t$ belongs to $\dot T$ precisely when at least one arrow $u\to t$ exists with $u\in U$. Its gluing weight is
\[
 H_t=\Hom_X(i(-),t)\in\widehat U.
\]
Only these objects need testing. At $u\in U$ the weight is representable and therefore totally aspherical; outside $\overline U$ it is empty and no assertion of asphericity is intended.

The hypothesis of Theorem~3 is that every $H_t$, $t\in\dot T$, remains aspherical after every ordinary left Kan extension. A pushout along $p:U\to U'$ replaces this weight by $p_!H_t$. Its total asphericity follows without a gluing axiom: for a second functor $r:U'\to V$,
\[
 r_!(p_!H_t)\cong(rp)_!H_t,
\]
and the right-hand side is aspherical by the original hypothesis. This calculation explains both the terminology ``strictly aspherical'' for $U\to\overline U$ and its stability under arbitrary pushout.

The support is preserved as well. If $H_t$ is nonempty, a representative $h\in H_t(u)$ produces a class at $p(u)$ in $p_!H_t$; a colimit of sets cannot delete that class. If $H_t$ is empty, its direct image is empty. Thus the boundary objects of the two pushouts identify with the same $\dot T$, and the preceding test covers exactly the new boundary. It follows that
\[
 i'_0:U'\longrightarrow\overline{U'}
\]
is strictly $W$-aspherical and hence $W$-aspherical. The claim at the start of the proof on ms.~112 that $i'$ itself is aspherical must read $i'_0$. A $W$-dwyérienne inclusion need not be a weak equivalence: the empty sieve into a nonempty category already demonstrates the distinction.

The two equivalences printed as (a) and (b) have a different logical status from this weight calculation. They use respectively W(7 bis) and W(7), exactly as stated. Those additional axioms must remain visible; they are not automatic properties of weak homotopy equivalences of nerves. The modern Dwyer pushout theorem gives the useful forward implications in its own hypotheses, as discussed in the later-development block. The intrinsic statement here, valid before imposing either gluing axiom, is preservation of the condition on all the weights.
\end{dcomment}


\dgrblock{dgrb536}{A boundary with no last object}
\begin{dcomment}{examples}{p108:Démonstration page 112.}{A boundary with no last object}
The footnote about filtered categories points beyond the original adjoint criterion. The commentary on Chapter~II, manuscript pp.~35--36 (``Left Kan extension carries the ind-presentation'' and ``The quantifier in total asphericity''), explains why filtered presentations make this universal condition tractable. Here their role can be seen in one concrete inclusion.

Let $U$ be the ordered set $\mathbb N$ with its usual order, and let $X=\mathbb N\cup\{\infty\}$ with $n<\infty$ for every $n$. The inclusion $i:U\hookrightarrow X$ is a sieve and $\overline U=X$. For an old object $n$, the slice $U/n$ has terminal object $n$. For the new object $\infty$, the slice is $U$ itself. For any $p:\mathbb N\to B$, one has
\[
 F:=p_!e_{\mathbb N}\cong\varinjlim_n\Hom_B(-,p(n)).
\]
Write $p_{nm}:p(n)\to p(m)$ for the transition map when $n\leq m$, and let $E=B/F$. This category is nonempty, since it contains $(p(0),[\id_{p(0)}])$. Given finitely many objects $(b_j,x_j)$ of $E$, choose representatives $g_j:b_j\to p(n_j)$ of their elements. For $m\geq n_j$ for every $j$, the arrows $p_{n_jm}g_j$ give morphisms from all these objects to $(p(m),[\id_{p(m)}])$.

For parallel morphisms $h,k:(b,x)\rightrightarrows(c,y)$ in $E$, choose a representative $g:c\to p(n)$ of $y$. The defining condition on these morphisms gives $[gh]=x=[gk]$ in $F(b)$. Equality in this filtered colimit means that, for some $m\geq n$,
\[
 p_{nm}gh=p_{nm}gk.
\]
Thus $p_{nm}g:(c,y)\to(p(m),[\id_{p(m)}])$ equalizes $h$ and $k$. These are exactly the nonemptiness, finite-cocone, and parallel-arrow tests for a filtered category. A filtered category has weakly contractible nerve: a cocone on the finite data of a simplicial sphere map contracts that map. Thus every such $E$ is $W_\infty$-aspherical.

It follows that $i$ is totally $W_\infty$-dwyérienne. It is also a couniversal weak equivalence, without using either additional gluing axiom. After any pushout its slice over the new object is the filtered category of elements just computed, while its slices over the old objects have terminal objects. Quillen's Theorem~A therefore applies to every pushout inclusion. It is not a Dwyer immersion in the original adjoint sense. A right adjoint $r:X\to U$ to $i$ would require
\[
 n\leq r(\infty)\quad\text{for every }n\in\mathbb N,
\]
which is impossible. This distinguishes the two classes even though both $\mathbb N$ and $\mathbb N^\triangleright$ are contractible.

The negative answer also holds within the full gluing framework assumed in the source's question. Let $W_{\mathbb Z}$ denote integral-homology equivalences of nerves. Every pushout just considered is a weak homotopy equivalence and therefore an integral-homology equivalence, so $i$ is couniversally a $W_{\mathbb Z}$-equivalence. Every tested category of elements is already weakly contractible, so $i$ is also totally $W_{\mathbb Z}$-dwyérienne. The development at ms.~100, \emph{Testing the retracted assertions under their stated axiom}, verifies that $W_{\mathbb Z}$ is a fundamental localizer satisfying full W(7). The failure of the right-adjoint condition is independent of the choice of weak equivalences. Thus the proposed recovery of an original Dwyer immersion fails even with W(7) imposed. The distinct converse from an arbitrary homotopical cofibration to Grothendieck's condition on every boundary weight is not established by this example.
\end{dcomment}


\textbf{Corollaire 1.} \textit{Supposons} W(7)\textit{. Si $i:U\hra X$ est une immersion ouverte telle que $i_0:U\hra\bar{U}$ soit totalement} \texttt{[il suffit strictement]} \textit{$W$-asphérique, alors $i:U\longrightarrow X$ est une $W$-cofibration (`ultraforte'), par quoi j'entends ceci : pour toute flèche $U'\xrightarrow{f}U''$ dans $U\backslash$}Cat\textit{,}

\msp{page 109}

\textit{telle que $f\in W$, la flèche $g:X'\longrightarrow X''$ qui s'en déduit par cochangement de base à partir de l'objet $X$ de $U\backslash$}Cat\textit{, est aussi dans $W$.}

\textsc{Démonstration du corollaire} : on applique le théorème 3 à l'immersion ouverte $U'\overset{i'}{\hra}X'$ déduite de $U\overset{i}{\hra}X$ par le cochangement de base $U\xrightarrow{f}U'$, en notant de plus que l'hypothèse du théorème 3 sur $i$ est préservée en passant de $i$ à $i'$.

Il faut quand même introduire une terminologie adéquate :

\textbf{Définition 2 :}

\begin{enumerate}[label=\alph*)]
\item Une flèche $f:X\longrightarrow Y$ est dite \textit{strictement $W$-asphérique} si pour toute $y\in\operatorname{Ob}Y$, $f^*(y)=(x\longmapsto\operatorname{Hom}_Y(f(x),y))$ est un objet totalement $W$-asphérique de $X^\wedge$ (cf. déf. 1, page 88). Elle est dite \textit{totalement $W$-asphérique} si pour tout $y\in\operatorname{Ob}Y$, $X/y$ ($\simeq X/f^*(y)$) est un objet totalement $W$-asphérique de Cat.

\item \printedpage{76}
% head: A. Grothendieck : Dérivateurs, Ch. XII, Édition des Universités de Montpellier II et Paris VII
% foot: Version 7 décembre 2002
Une immersion ouverte $i:U\longrightarrow X$ est dite \textit{$W$-dwyérienne} (resp. \textit{totalement} \texttt{[$W$-]} \textit{dwyérienne}) si l'immersion induite $i_0:U\longrightarrow\bar{U}$ est strictement (resp. totalement) $W$-asphérique. Une immersion fermée $j:T\longrightarrow X$ est dite \textit{$W$-dwyérienne} (resp. \textit{totalement $W$-dwyérienne}) si et seulement si $j^\circ:T^\circ\longrightarrow X^\circ$ est une immersion ouverte $W$-dwyérienne (resp totalement $W$-dwyérienne).

\item Une flèche $i:U\longrightarrow X$ de Cat (pas nécessairement une immersion)

\msp{page 110}

est appelée une \textit{$W$-cofibration,} si pour toute flèche $f:U'\longrightarrow U''$ dans $U\backslash$Cat, la flèche $g:X'\longrightarrow X''$ qui s'en déduit par le cochangement de base par $X$ (objet de $U\backslash$Cat) est dans $W$ quand $f$ l'est.
\end{enumerate}

\textbf{NB} Cette notion de $W$-cofibration est stable par composition et par cochangement de base, et les isomorphismes de Cat sont des $W$-cofibrations. Si $f$ est une $W$-cofibration, $f^\circ$ l'est et inversement.

(Le corollaire s'énonce en disant qu'une immersion ouverte ou fermée qui est $W$-dwyérienne est une $W$-cofibration.)

\begin{enumerate}[label=\alph*), start=4]
\item Une flèche $i:U\longrightarrow X$ dans Cat est appelée une \textit{$W$-cofibration triviale}, ou une \textit{$W$-équivalence couniverselle,} si pour toute $U'$ dans $U\backslash$Cat, i.e. toute flèche $p:U\longrightarrow U'$ dans Cat, la flèche déduite de $i$,
\[ i':U'\longrightarrow X'=U'\amalg_U X'\quad\texttt{[plutôt }U'\amalg_U X\texttt{]}\ \ , \]
est dans $W$. (NB Il est immédiat que cela implique bien que $i$ est une $W$-cofibration.)
\end{enumerate}

(Cette notion de $W$-cofibration triviale est elle aussi stable par composition et par cochangement de base, et les isomorphismes dans Cat en sont.)

\msp{page 111}

\dgrblock{dgrb537}{What the cofibration condition quantifies over}
\begin{dcomment}{exposition}{p110}{What the cofibration condition quantifies over}
Grothendieck's definition concerns the entire undercategory, and this is why he calls the condition ``ultraforte''. For $i:U\to X$, let
\[
 i_!:U\backslash\Cat\longrightarrow X\backslash\Cat,
 \qquad (U\to A)\longmapsto(X\to X\amalg_U A).
\]
The condition is that this functor preserve the weak equivalences of underlying categories. In particular it tests a weak equivalence $A\to B$ after an arbitrary preliminary change of base $U\to A$. Merely requiring preservation of weak equivalences whose source is $U$ would be weaker. The modern terminology and its relation to model-category cofibrations are treated in the later-development block.

The elementary closure assertions follow from the universal properties of pushouts. If $U\to X\to Z$ are two such maps, pushout along the composite is the composite of their pushout functors, up to the canonical isomorphism, so it preserves weak equivalences. If $U\to V$ is arbitrary and $i':V\to X\amalg_U V$ is the pushout of $i$, then for $A\to B$ under $V$ its two pushouts along $i'$ are canonically the pushouts along $i$ of the same map considered under $U$. This proves stability under change of base. Isomorphisms give equivalences of the undercategories. Passage to opposite categories proves the duality assertion provided that $W$ is invariant under opposites, as in the standing localizer framework.

The adjective ``trivial'' is defined by a separate universal property: every pushout $A\to X\amalg_U A$ of $i$ belongs to $W$. This implies the preceding preservation condition by two-out-of-three. Indeed the natural square
\[
\begin{tikzcd}
A\arrow[r]\arrow[d,"f"']&X\amalg_U A\arrow[d,"g"]\\
B\arrow[r]&X\amalg_U B
\end{tikzcd}
\]
has both horizontal maps in $W$; when $f\in W$, commutativity and two-out-of-three give $g\in W$. Stability of these universal weak equivalences under composition and pushout follows by factoring each new pushout into two pushouts of the original maps.

Without an additional hypothesis one should not replace this definition by ``a $W$-cofibration which belongs to $W$''. Corollary~2 establishes that implication for the $W$-dwyérienne sieves considered there, using Theorem~3 and the stated gluing assumptions. This is the reason for retaining the full definition even though its name resembles the terminology of a model category.
\end{dcomment}


\dgrblock{dgrb538}{The later theory of homotopy cofibrations}
\begin{dcomment}{later}{p110}{The later theory of homotopy cofibrations}
Grothendieck wants a property of an attachment that is visible from weak equivalences alone. This is exactly the property now called an \emph{$h$-cofibration}: the functor of cobase change on undercategories preserves weak equivalences. The terminology and a general model-categorical treatment appear in \citet[Definition~1.1]{BataninBerger2017}, who explicitly acknowledge the antecedent in Grothendieck's manuscripts. In a left proper model category their Proposition~1.6 and Corollary~1.9 identify this condition with the requirement that every pushout along the map compute a homotopy pushout. Their Lemma~1.10 also proves that an $h$-cofibration which is a weak equivalence is couniversal. These are precise later replacements for the useful direction of the present gluing programme; none requires the reverse implication of W(7 bis). The same undercategory definition is retained as ``$W$-cofibration'' by \citet[Definition~2.1 and Lemma~2.4]{ManettiMeazzini2019}.

The definition here must be kept distinct from the cofibrations of a chosen Quillen model structure. For example, for every small category $C$, the map $\varnothing\to C$ is a $W_\infty$-cofibration: its cobase-change functor is $A\mapsto A\amalg C$, and the nerve of a disjoint union is a disjoint union of nerves. A weak equivalence on the first summand and the identity on the second induce a weak equivalence on the union, component by component. Thus every category is cofibrant for Grothendieck's \emph{relative} notion. This statement does not assert that every category is cofibrant in Thomason's model structure. The distinction explains why a satisfactory description of these attachments can be broader than the particular generating cofibrations used to construct a model category.
\end{dcomment}


\textbf{Corollaire 2.} (Moyennant W(7).) \textit{Soit $i:U\longrightarrow X$ une immersion ouverte $W$-dwyérienne. Pour que $i$ soit une $W$-cofibration triviale, il faut et il suffit que $i\in W$, ou ce qui revient au même} (cf. page 82), \textit{que $\dot{T}\hra T$ soit dans $W$.}

C'est évidemment nécessaire, et la suffisance résulte du théorème 3 a).

\textbf{Questions ouvertes.}

\begin{enumerate}[label=\alph*)]
\item J'ignore si une immersion ouverte $i$ qui est une $W$-cofibration, voire une $W$-cofibration triviale, est forcément $W$-dwyérienne (\footnote{et peut-être même une immersion de Dwyer?}). Mais par le corollaire de la proposition 7 (p. 85), je sais que c'est le cas si on suppose $T$ isomorphe à la catégorie finale $e$. Il s'ensuit toute au moins ceci : Pour que non seulement $i:U\longrightarrow X$, mais aussi toutes les immersions ouvertes $i_1:U\longrightarrow X_1$ qui s'en déduisent par le procédé de prop. 7 b) au moyen de flèches quelconques $T_1\longrightarrow T$ de but $T$ dans Cat, soient des $W$-cofibrations, il faut et il suffit que $i$ soit dwyérienne. % [sic] toute au moins -> tout au moins, agreement error

\item \printedpage{77}
% head: A. Grothendieck : Dérivateurs, Ch. XII, Édition des Universités de Montpellier II et Paris VII
% foot: Version 7 décembre 2002
Y-a-t-il une bonne notion d' immersion dwyérienne, applicable à toutes les immersions

\msp{page 112}

\dgrblock{dgrb539}{The open questions of this page}
\begin{dcomment}{development}{p111}{The open questions of this page}
Question~a) asks whether a sieve that is a $W$-cofibration, or even a trivial $W$-cofibration, must be $W$-dwy\'erienne, and asserts this when the complement $T$ is the final category. Three small examples separate the notions. In each of them $U=\{u\}$ is a single object, so a functor $p:U\to A$ is the choice of an object $c=p(u)$ of $A$.

\textbf{Proposition (a trivial cofibration that is not dwy\'erienne).}\ours
\textit{Let $X$ have objects $u,t,s$ and nonidentity arrows $\alpha,\beta:u\to t$, $\gamma:t\to s$, $\delta:u\to s$, with $\gamma\alpha=\gamma\beta=\delta$. The sieve $U=\{u\}\hookrightarrow X$ is a trivial $W$-cofibration for every fundamental localizer $W$. It is not $W$-dwy\'erienne when $W$ is $W_\infty$, a homology localizer $W_R$ with $R\neq0$, or the class $W_n$ of $n$-equivalences for some $n\geq0$.}

\textit{Proof.}
The presheaf $v\mapsto\operatorname{Hom}_X(v,t)$ on $U$ has two elements, so $U/t$ is a discrete category with two objects; the test $p=\mathrm{id}$ shows that this presheaf is not totally $W$-aspherical for the localizers named. Let $p:U\to A$ be arbitrary, $c=p(u)$, and $P=X\amalg_UA$. Its objects are those of $A$ and $t,s$; for $v\in A$,
\[
 P(v,t)=A(v,c)\times\{\alpha,\beta\},\qquad P(v,s)=A(v,c),\qquad P(t,s)=\{\gamma\},
\]
composition with $\gamma$ forgets the second coordinate, and there are no arrows from $t$ or $s$ to $A$. Let $j:A\to P$ be the inclusion and $r:P\to A$ the functor that is the identity on $A$, sends $t,s$ to $c$, projects the two crossing sets to $A(v,c)$ and sends $\gamma$ to $\mathrm{id}_c$. Let $F:P\to P$ be the identity on $A$, send $t$ and $s$ to $s$, each crossing arrow to its composite into $s$, and $\gamma$ to $\mathrm{id}_s$. There are natural transformations
\[
 \mathrm{id}_P\Longrightarrow F\Longleftarrow jr,
\]
both the identity on $A$; the first has components $\gamma$ at $t$ and $\mathrm{id}_s$ at $s$, the second has at $t$ and at $s$ the arrow $c\to s$ corresponding to $\mathrm{id}_c$. Naturality at a crossing arrow $(g,\alpha)$ or $(g,\beta)$ is the relation $\gamma\alpha=\gamma\beta=\delta$. Since $rj=\mathrm{id}_A$, the map $Bj$ is a homotopy equivalence. So every cobase change of $U\to X$ is in $W_\infty\subseteq W$.\hfill$\square$

This answers the first sentence of question~a) negatively; the footnote's stronger hope fails with it, since a Dwyer immersion is dwy\'erienne.

\textbf{Proposition (the case $T\simeq e$).}\ours
\textit{Let $X_0$ have objects $u,t$ and exactly two nonidentity arrows $\alpha,\beta:u\to t$, and let $U=\{u\}$, $T=\{t\}$. For $W=W_\infty$ and for $W=W_R$, the sieve $U\hookrightarrow X_0$ and every sieve $U\hookrightarrow X_1$ deduced from it by a change of complement $T_1\to T$ (Proposition~7~b)) is a $W$-cofibration, and none of them with $T_1\neq\varnothing$ is $W$-dwy\'erienne.}

\textit{Proof.}
The category $X_1$ has the objects $u$ and those of $T_1$, the arrows of $T_1$, and two arrows $\alpha_x,\beta_x:u\to x$ for every $x\in T_1$, compatible with the arrows of $T_1$. Each crossing presheaf has two elements, so no $U/x$ is connected. For $p:U\to A$ with $c=p(u)$, the pushout $P_A=X_1\amalg_UA$ has $P_A(v,x)=A(v,c)\times\{\alpha,\beta\}$ for $v\in A$, $x\in T_1$. A chain of arrows of $P_A$ either lies in $A$, or lies in $T_1$, or consists of a chain in $A$, a crossing arrow $(g,\lambda)$ with $\lambda\in\{\alpha,\beta\}$, and a chain in $T_1$; the chain in $A$ together with $g$ is a chain of $A/c$. Write $J=N(A/c)*NT_1$ for the join, and $J_\alpha\cup J_\beta$ for two copies of $J$ glued along $NT_1$. Then $NP_A$ is the pushout of simplicial sets
\[
 NA\longleftarrow N(A/c)\amalg N(A/c)\hookrightarrow J_\alpha\cup J_\beta ,
\]
in which the left map forgets the arrow to $c$ and the right map is a monomorphism; it is therefore a homotopy pushout. A functor $h:A\to A'$ under $U$ induces a map of these diagrams. The categories $A/c$ and $A'/h(c)$ have terminal objects, so their nerves and the joins $J$ are contractible, and the maps on $N(A/c)\amalg N(A/c)$ and on $J_\alpha\cup J_\beta$ are weak equivalences whatever $h$ is. If $h\in W_\infty$, the map $NP_A\to NP_{A'}$ is then a weak equivalence. For $W_R$ one uses the Mayer--Vietoris sequences of the same pushouts.\hfill$\square$

So the assertion ``je sais que c'est le cas si on suppose $T$ isomorphe \`a la cat\'egorie finale $e$'' is false for $W$-cofibrations, and with it the equivalence stated at the end of question~a), also under the axiom W(7), which the localizers $W_R$ satisfy. The corollary of Proposition~7 concerns membership in $W$ of one inclusion, not the cofibration property. What it does give is the statement for \emph{trivial} cofibrations and homology:

\textbf{Proposition.}\ours
\textit{Let $U\hookrightarrow X$ be a sieve whose complement is one object $t$ with only its identity, and suppose that $H=\operatorname{Hom}_X(-,t)|_U$ is not empty. For a nonzero commutative ring $R$, the inclusion is a trivial $W_R$-cofibration if and only if $H$ is totally $W_R$-aspherical, that is, if and only if it is $W_R$-dwy\'erienne.}

\textit{Proof.}
After a cobase change $p:U\to A$ the complement is still $\{t\}$ and the crossing presheaf is $p_!H$. A chain of arrows of $X\amalg_UA$ that does not lie in $A$ is a chain of $A/p_!H$ followed by $t$, so $N(X\amalg_UA)$ is $NA$ with the cone on $N(A/p_!H)$ attached along the forgetful map, and $B(X\amalg_UA)$ is the mapping cone of $B(A/p_!H)\to BA$. The inclusion of $BA$ is an $R$-homology isomorphism exactly when the reduced $R$-homology of $A/p_!H$ vanishes, that is, when $A/p_!H$ has the $R$-homology of a point. Quantifying over $p$ gives the statement.\hfill$\square$

For $W_\infty$ the same computation shows that a one-point sieve is a trivial $W_\infty$-cofibration if and only if every $A/p_!H$ is nonempty and integrally acyclic and the cone attachment does not change fundamental groups. Whether this forces $H$ to be totally $W_\infty$-aspherical is not decided here: a counterexample would be a presheaf $H$ all of whose direct images have acyclic categories of elements, one of them with a nontrivial (necessarily perfect) fundamental group.
\end{dcomment}


et incluant le cas des immersions ouvertes et des immersions fermées?
\end{enumerate}

\textsc{Démonstration} du théorème 3 (p. 107). Il suffit de prouver que si $i$ est $W$-dwyérienne (i.e. $i_0$ strictement $W$-asphérique). alors $i'$ l'est également - car il en résultera que $i'$ est $W$-asphérique et a fortiori $\in W$. Il faut donc vérifier que pour tout $t'\in\dot{U}'=\dot{T}$, l'élément $i^*(t')=(u'\longmapsto\operatorname{Hom}_{X'}(u',t'))$ de $U'^\wedge$ est totalement $W$-asphérique. Or $t'$ provient d'un unique élément $t$ de $\dot{T}$, et on sait qu'on a % [sic] asphérique). alors -> asphérique), alors, period before lower-case word

\[ i'^*(t)=p_!(i^*(t))\quad\texttt{[plutôt }i'^*(t')=p_!(i^*(t))\texttt{]},\qquad\text{où }\ p_!:U^\wedge\longrightarrow U'^\wedge. \]

Par hypothèse, $F=i^*(t)$ est un élément totalement $W$-asphérique de $U^\wedge$, ce qui signifie que pour toute flèche $p:U\longrightarrow U'$ dans Cat de souce $U$, $p_!(F)$ est un élément $W$-asphérique de $U'$. En fait, on a le lemme immédiat (par transitivité $(p'p)_!\simeq p'_!p_!$). % [sic] souce -> source, missing letter

\textbf{Lemme.} \textit{Soit $U$ dans} Cat\textit{, et $F$ dans $U^\wedge$. Si $F$ est totalement $W$-asphérique (resp. si $U/F$ est totalement $W$-asphérique), alors pour toute flèche $p:U\longrightarrow U'$ dans} Cat \textit{de source $U$, $F'=p_!(F)$ est un objet totalement $W$-asphérique de $U'^\wedge$ (resp.}

\msp{page 113}

\textit{$U'/F'$ est un objet totalement $W$-asphérique de} Cat\textit{.)}

\textbf{Corollaire 3} (du théorème 3). \textit{Si $i:U\longrightarrow X$ est une immersion ouverte ou fermée $W$-dwyérienne (resp. totalement $W$-dwyérienne), alors pour toute flèche $p:U\longrightarrow U'$ dans} Cat\textit{, $i':U'\longrightarrow X'=U'\amalg_U X'$} \texttt{[plutôt $U'\amalg_U X$]} \textit{l'est aussi.}

(NB Le cas non respé est essentiellement le théorème 3. Il ne dépend d'aucun axiome supplémentaire sur $W$, mais n'a d'intérêt qu'en vue des équivalences \textcircled{a} et \textcircled{b} (p. 108) qui en résultent, et qui sont valables moyennant W(7 bis) resp. W(7) seulement.)

\dgrblock{dgrb540}{The stronger clause uses comprehensive factorization}
\begin{dcomment}{exposition}{p113:Le cas non respé est essentiellement le théorème 3.}{The stronger clause uses comprehensive factorization}
The lemma proves more than stability of a condition on a presheaf. Its parenthesized assertion says that the category of elements after direct image is itself totally aspherical, so that one may test it by arbitrary functors leaving that category. This does not follow merely by repeating the formula for a second direct image on the original base. The canonical comprehensive factorization supplies the missing step.

Write $V=U/F$, let $a:V\to U$ be the projection, and set $b=pa:V\to U'$. The identity $a_!e_V=F$ gives
\[
 b_!e_V=p_!F=F'.
\]
Let $E=U'/F'$. The canonical functor $v:V\to E$ appearing in the comprehensive factorization of $b$ is final in the ordinary sense: for every $e\in E$, the category $e\backslash V$ is connected and nonempty. Equivalently,
\[
 v_!e_V=e_E.
\]
This is the factorization established at ms.~91--93; it involves connectedness, not contractibility.

Now suppose that $V$ is totally $W$-aspherical, and let $r:E\to Z$ be any functor. Transitivity of ordinary left Kan extension gives
\[
 r_!e_E\cong r_!v_!e_V\cong(rv)_!e_V.
\]
The last presheaf is $W$-aspherical by totality of $V$, so $E$ is totally $W$-aspherical. This proves the parenthesized clause with its full quantifier. The unparenthesized clause is simpler: for $s:U'\to Z$, one has $s_!F'\cong(sp)_!F$, to which the definition applies directly.

The same argument yields a useful formal principle. If $v:V\to E$ is an ordinary final functor and $V$ is totally $W$-aspherical, then $E$ is totally $W$-aspherical. The converse is not being asserted. In particular, an arbitrary final functor need not be a weak homotopy equivalence; the finite example at ms.~100 makes that failure explicit. Here totality of the source is precisely what permits the final functor to be used in every subsequent direct-image calculation.

Corollary~3 follows by applying the two clauses to the gluing weights of each boundary object, with the support calculation already explained after Theorem~3. This proof establishes stability under pushout. It supplies no composition theorem for arbitrary strictly or totally $W$-aspherical functors: in a composite one pulls back a weight of the second functor, whereas the lemma controls left Kan extension. The distinction identifies the mathematical content of the doubt expressed on ms.~113--114. Composition is treated in the commentary at ms.~114: both classes of functors are stable under it, by a different argument.
\end{dcomment}


\textbf{Commentaire de motivation :} les immersions ouvertes ou fermées $W$-dwyériennes, ou totalement $W$-dwyériennes, sont stables par cochangement de base (d'après le corollaire précédent), mais n'ont pas du tout l'air de l'être par composition. À cet égard, les `immersions de Dwyer' (où on suppose l'existence de l'adjoint à droite de $i_0:U\longrightarrow\bar{U}$, dans le cas d'une immersion ouverte) se comportent mieux, étant stables par composition. (\textbf{NB} Il ne semble pas

\msp{page 114}

non plus clair que les flèches strictement $W$-asphériques, ou les totalement $W$-asphériques, soient stables par composition - sorry!) Ma motivation pour m'échiner sur la notion $W$-dwyérienne est double.

\begin{enumerate}[label=\alph*)]
\item J'aimerais avoir des conditions à la fois compréhensibles, et aussi générales que possible, sur une immersion (ouverte p. ex), pour que ce soit une $W$-cofibration, propriété visiblement importante (toute comme la notion duale de $W$-fibration). % [sic] toute comme -> tout comme, agreement error
\end{enumerate}

% ---- from Der-ChXIIb-012_p036-040.tex
\printedpage{78}
% head: A. Grothendieck : Dérivateurs, Ch. XII, Édition des Universités de Montpellier II et Paris VII
\begin{enumerate}
\item[b)] J'ai envie, pour des applications ultérieures, de disposer d'une classe d'immersions ouvertes ayant les propriétés suivantes :
\begin{enumerate}
\item[1°)] Ce sont des $W$-cofibrations, et
\item[2°)] Si $U \overset{i}{\hookrightarrow} X$ est dans la classe, et si $\varepsilon \xrightarrow{j} I$ est l'inclusion ‘source-but’ de la catégorie discrète $\{0,1\}$ dans l'intervalle
\[
\Phi_n = 0 \longrightarrow 1 \longleftarrow 2 \longrightarrow \dots 2n-1 \longleftarrow 2n
\]
$(j(0)=0$, $j(1)=2n)$, $n\geq 2$, l'inclusion $\alpha$ de $(U\times I)\cup(X\times\varepsilon)$ dans $X\times I$ est dans la classe.
\end{enumerate}
Or cette dernière propriété de stabilité n'est \textit{pas} vérifiée par les immersions de Dwyer. (Prendre $i : \{0\} \longrightarrow \boldsymbol{\Delta}^1 = \{0,1\}$, l'inclusion $\alpha$ n'est pas immersion de Dwyer.)
\end{enumerate}

\msp{page 115}

\dgrblock{dgrb541}{The interval test fails because a slice has two maxima}
\begin{dcomment}{examples}{p114}{The interval test fails because a slice has two maxima}
The interval condition asks for a class of admissible inclusions that can also attach both ends of a homotopy. Its categorical form is a pushout product. The original Dwyer condition is too rigid for this purpose, and the example in the text can be checked at a single object.

Take $X=[1]=(0<1)$, $U=\{0\}$, and
\[
 I=\Phi_n=(0\to1\leftarrow2\to\cdots\leftarrow2n),
 \qquad n\geq2.
\]
Let $A=(\{0\}\times I)\cup([1]\times\{0,2n\})$, considered as a full subposet of $[1]\times I$. The inclusion $A\hookrightarrow[1]\times I$ is a sieve: both end vertices of $I$ are sources, and the entire lower row is already in $A$. Its generated cosieve is the whole product, since every upper-row object $(1,k)$ receives an arrow from $(0,k)$.

If this were an original Dwyer immersion, the inclusion into that cosieve would admit a right adjoint. For an inclusion of ordered sets, this means that for every object $x$ of the larger poset the lower set $A/x$ has a greatest element. Test $x=(1,1)$. Its lower set inside $A$ contains exactly
\[
 (0,0),\quad(0,1),\quad(0,2),\quad(1,0).
\]
The strict relations are
\[
 (0,0)<(0,1),\qquad
 (0,2)<(0,1),\qquad
 (0,0)<(1,0).
\]
Thus $(0,1)$ and $(1,0)$ are distinct incomparable maximal elements, and neither is above the other. There is no greatest element. This proves the failure of the adjoint condition for every $n\geq2$.

The calculation also explains why ordinary contractibility of this slice is insufficient to recover the missing adjoint: its nerve is an interval with three edges, hence contractible. An adjoint selects a universal object, a much stronger requirement. Whether another class provides the desired interval stability is a structural question about the interaction of pushout and homotopy, rather than a reason to weaken \emph{terminal object} to \emph{contractible} in this particular criterion. The subsequent use of integration replaces a difficult ordinary pushout by a category designed to retain the homotopy of the entire diagram.
\end{dcomment}


\dgrblock{dgrb542}{A terminal fibre controls all the ordinary direct images}
\begin{dcomment}{exposition}{p114}{A terminal fibre controls all the ordinary direct images}
The last subsection returns to a situation in which the apparent loss of higher information can be controlled. For a $W$-proper functor over a category with a terminal object, the terminal-fibre inclusion is $W$-coaspherical. Under W(5) it is ordinary final, which is the property needed for total asphericity. Making this explicit shortens the proof and explains the editors' repeated suggestion to replace $W$-properness by $W_0$-properness.

Let $f:X\to I$ be $W$-proper, let $i_0$ be terminal in $I$, and write $k:X_{i_0}\hookrightarrow X$. The properness criterion says that $k$ is $W$-coaspherical: every undercategory $x\backslash X_{i_0}$ is $W$-aspherical. If W(5) holds, a $W$-aspherical category is nonempty and connected. Hence $k$ is an ordinary final functor, and the pointwise formula gives
\[
 (k_!e_{X_{i_0}})(x)=\pi_0(x\backslash X_{i_0})=\{*\}.
\]
For every $p:X\to Y$ it follows that
\[
 (pk)_!e_{X_{i_0}}
 \cong p_!k_!e_{X_{i_0}}
 \cong p_!e_X.                                      \tag{C1}
\]
If $X_{i_0}$ is totally $W$-aspherical, the left side of (C1) is an aspherical presheaf on $Y$ for every $p$, proving Proposition~11 of ms.~115. This proposition has the same number as the counterexample on ms.~107; the manuscript page is needed to distinguish them. The proof uses no W(7) or W(7 bis).

The general component formula in Proposition~12 explains how (C1) fits into a family. $W$-properness together with W(5) identifies $\pi_0(X_i)$ with $\pi_0(X/i)$, naturally in $i$, and therefore supplies the transition maps even when $f$ is not an opfibration. For an arbitrary functor $f$, there is a canonical bijection
\[
 \varinjlim_{i\in I}\pi_0(X/i)\ \cong\ \pi_0(X).        \tag{C2}
\]
One can see it without choosing a cleavage. A representative $(x,\beta:f(x)\to i)$ in $X/i$ has the same image in the colimit as $(x,\id_{f(x)})$, by the transition along $\beta$. A morphism $x\to x'$ identifies the latter representatives through the induced object of $X/f(x')$. Thus the colimit imposes exactly the zigzag relation defining components of $X$. Conversely each morphism in a comma category, and every transition between comma categories, preserves that zigzag component. These two observations construct mutually inverse maps in (C2).

Replacing $\pi_0(X/i)$ by $\pi_0(X_i)$ proves Proposition~12. The dual proof printed in the manuscript is equally valid, but after passing to presheaves its colimit is indexed by $I^{\mathrm{op}}$; the printed $I$ at ms.~117 suppresses this reversal. The precise proper/smooth criteria, their duality and their stability under composition and base change are given in \citet[Theorem~1.2 and Proposition~4.7]{Maltsiniotis2005Asphericite}. The result here concerns ordinary components and ordinary Kan extensions, so the connectivity hypothesis cannot be omitted just because the discussion takes place in a homotopical setting.
\end{dcomment}


\dgrblock{dgrb543}{Dwyer attachments and higher categorical pushouts}
\begin{dcomment}{later}{p114}{Dwyer attachments and higher categorical pushouts}
The adjoint condition belongs to the Dwyer maps of \citet[Definition~4.1 and Proposition~4.3]{Thomason1980}. One subsequent correction matters when comparing them with model-category cofibrations. Classical Dwyer maps are not closed under retracts, and some Thomason cofibrations are not Dwyer maps. \citet[Definition~1, Lemma~4 and Proposition~2]{Cisinski1999Dwyer} repairs the relevant argument with pseudo-Dwyer maps: these are retract-stable and include every Thomason cofibration. Thus the manuscript's search for a broader attachment condition addresses a distinction that cannot be removed by simply identifying the original Dwyer maps with all cofibrations.

The desired stability properties concern whether an ordinary attachment retains the information needed after passage to a homotopy theory. For classical Dwyer maps, a stronger result is now available. \citet[Theorem~1.5]{HackneyOzornovaRiehlRovelli2024} prove that, for a pushout $D=B\amalg_A C$ along a Dwyer inclusion $A\hookrightarrow B$, the comparison
\[
 NB\amalg_{NA}NC\longrightarrow ND
\]
is a weak \emph{categorical} equivalence in the Joyal model structure. Their general Theorem~1.6, in its expanded statement on p.~2179, treats discretely flat cofibrations. These are defined by requiring the corresponding functor of discrete simplicial categories to preserve Bergner weak equivalences under cobase change, together with injectivity on objects. Proposition~3.3 supplies this property for Dwyer maps; Example~3.5 includes co-Dwyer maps and pseudo-Dwyer maps. This gives one modern class that accommodates both orientations requested on p.~111.

This higher categorical assertion concerns the comparison between the ordinary categorical pushout and a pushout in the category of $(\infty,1)$-categories. The weak homotopy type of a category discards all noninvertible-arrow information, whereas its image among $(\infty,1)$-categories retains that information. Consequently the two ambient weak equivalence classes must be specified before comparing definitions of flatness. The theorem just cited does not identify discretely flat inclusions with the manuscript's $W$-Dwyer immersions, nor does it by itself establish the particular interval-product stability requested here. Those further identifications would require additional arguments. The concrete calculations in the preceding blocks keep the original question at the level at which it is posed.
\end{dcomment}


\dgrblock{dgrb544}{Composition of aspherical functors and of dwyerienne immersions}
\begin{dcomment}{development}{p114}{Composition of aspherical functors and of dwy\'erienne immersions}
Grothendieck doubts that strictly or totally $W$-aspherical functors, and the two classes of dwy\'erienne immersions, are stable under composition. For functors the doubt is unfounded; for immersions it is justified for some localizers and not for others. The statements below were not found in the sources consulted; no priority is claimed.

\textbf{Theorem (functors compose).}\ours
\textit{Let $W$ be a fundamental localizer.}
\begin{enumerate}[label=(\alph*)]
\item \textit{If $f:X\to Y$ is strictly $W$-aspherical and $H$ is a totally $W$-aspherical presheaf on $Y$, then $f^*H$ is a totally $W$-aspherical presheaf on $X$. Consequently strictly $W$-aspherical functors are stable under composition.}
\item \textit{Totally $W$-aspherical functors are stable under composition.}
\item \textit{If $B$ is a totally $W$-aspherical category and $D:B\to\mathrm{Cat}$ a functor with totally $W$-aspherical values, the cofibred category $\int_BD$ is totally $W$-aspherical.}
\end{enumerate}

\textit{Proof.}
(a) Let $p:X\to X'$ be arbitrary. For $y\in Y$ put $Q_y=p_!\operatorname{Hom}_Y(f(-),y)$, a presheaf on $X'$, covariant in $y$. Each $Q_y$ is totally $W$-aspherical, by the hypothesis on $f$ and the lemma of ms.~112. Let $Z$ be the category with full subcategories $X'$ and $Y$, with $\operatorname{Hom}_Z(x',y)=Q_y(x')$ and no arrows from $Y$ to $X'$; let $j:Y\hookrightarrow Z$ be the inclusion and $F=j_!H$. Then $F|_Y=H$ and
\[
 F|_{X'}=\int^{y\in Y}Q_y\times H(y)\cong p_!\Bigl(\int^{y}\operatorname{Hom}_Y(f(-),y)\times H(y)\Bigr)=p_!f^*H=:K .
\]
Since $H$ is total, $Z/F$ is $W$-aspherical. The inclusion $l:X'/K\hookrightarrow Z/F$ is a sieve. Its comma category over an object $(x',k)$ of $X'/K$ has a terminal object. Its comma category over $(y,h)$, $h\in H(y)$, has as objects the pairs $(x',q)$ with $q\in Q_y(x')$ (the element of $K(x')$ is then the class of $(q,h)$) and as arrows the arrows of $X'$ that respect $q$; it is $X'/Q_y$, which is $W$-aspherical. By the relative form of Theorem~A that defines a fundamental localizer, $l\in W$, and so $X'/K$ is $W$-aspherical. As $p$ was arbitrary, $f^*H$ is total. If $g:Y\to Z_0$ is also strictly aspherical, apply this to $H=\operatorname{Hom}_{Z_0}(g(-),z)$, whose restriction along $f$ is $\operatorname{Hom}_{Z_0}(gf(-),z)$.

(b) Let $f:X\to Y$ and $g:Y\to Z_0$ be totally aspherical and $z\in Z_0$. The functor $f_z:X/z\to Y/z$, $(x,\gamma)\mapsto(f(x),\gamma)$, has comma category over $(y,\beta)$ canonically isomorphic to $X/y$, because the arrow to $z$ is determined as $\beta g(h)$ by $h:f(x)\to y$. So $f_z$ is totally aspherical, hence strictly aspherical by Proposition~8. The final presheaf of $Y/z$ is total because $g$ is totally aspherical; by (a) its restriction, the final presheaf of $X/z$, is total. That is, $X/z$ is a totally $W$-aspherical category.

(c) Put $E=\int_BD$, with objects $(b,x)$, $x\in D(b)$, and projection $\pi:E\to B$. For $b\in B$ the functor $D(b)\to E/b$, $x\mapsto(b,x,\mathrm{id}_b)$, has the left adjoint $(b',x,\beta)\mapsto D(\beta)(x)$. A right adjoint is final. So $E/b$ is a final image of $D(b)$ and is total; $\pi$ is a totally aspherical functor, as is $B\to e$, and (b) applies.\hfill$\square$

\textbf{Proposition (immersions: two counterexamples).}\ours
\textit{For $W=W_0$ and for $W=W_1$ neither the $W$-dwy\'erienne nor the totally $W$-dwy\'erienne sieves are stable under composition.}

\textit{Proof.}
Let $V$ be the discrete category on two objects for $W_0$, and the category with two objects and two parallel arrows for $W_1$; $BV$ is $S^0$, respectively $S^1$. Let $B=V^{\triangleleft}$ be $V$ with an initial object $o$ adjoined. Let $X$ consist of $B$ and of a second copy $V_0$ of $V$, with
\[
 X(v_0,b)=B(v,b),\qquad X(b,v_0)=\varnothing\qquad(v\in V,\ b\in B),
\]
and composition induced by $B$. The sieve $V_0\hookrightarrow X$ generates the cosieve $V_0\cup V$, on which the inclusion of $V_0$ has the right adjoint $v\mapsto v_0$; it is a Dwyer immersion. The functor $X\to B$ forgetting the copy is left adjoint to the inclusion $B\subset X$, so $BX\simeq BB$ is contractible. A final functor is bijective on $\pi_0$ and surjective on $\pi_1$; hence $X$ is totally $W_0$- and totally $W_1$-aspherical. Let $Y=X^{\triangleright}$ be $X$ with a terminal object $t$ adjoined. The sieve $X\hookrightarrow Y$ is totally $W$-dwy\'erienne, its only boundary category being $X/t=X$. For the composite sieve $V_0\hookrightarrow Y$ the boundary category at $t$ is $V_0/t\simeq V$ and the boundary presheaf is the final presheaf of $V_0$; neither is $W$-aspherical.\hfill$\square$

The targets $Y$ have six objects and $15$, respectively $21$, arrows.

\textbf{Theorem (total immersions for $W_\infty$).}\ours
\textit{For $W=W_\infty$ the totally $W$-dwy\'erienne sieves are stable under composition.}

\textit{Proof.}
We use the criterion of the commentary on Chapter~II, ms.~36: a category $A$ is totally $W_\infty$-aspherical if and only if it is nonempty and connected and every connected component of the category of elements of every functor $A\to\mathrm{Set}$ is weakly contractible; and then those components are themselves totally $W_\infty$-aspherical and any two objects of $A$ have a common successor. A nonempty cosieve $S$ of such an $A$ is the category of elements of the functor with value a point on $S$ and empty elsewhere; it is connected, because two of its objects have a common successor in $A$, which lies in $S$. So $S$ is totally $W_\infty$-aspherical.

Let $U\subset X\subset Y$ be totally $W_\infty$-dwy\'erienne sieves and $y\in Y$ with $U/y\neq\varnothing$. Then $X/y$ is total (it has a terminal object if $y\in X$). The objects $(x,\beta)$ of $X/y$ with $U/x\neq\varnothing$ form a nonempty cosieve $S_y$, total by the preceding paragraph, and the functor $(x,\beta)\mapsto U/x$ on $S_y$ has total values. By part~(c) of the theorem its cofibred category $G$ is total. The functor $G\to U/y$ that composes $u\to x\to y$ has the left adjoint $(u,\gamma)\mapsto((u,\gamma),(u,\mathrm{id}_u))$, so it is final, and $U/y$ is total.\hfill$\square$

The same proof works for the homology localizers $W_R$, with the homological form of the criterion ($R$-acyclic in place of weakly contractible). For the strict class, the $W_R$-dwy\'erienne sieves are also stable under composition: the restriction of a totally $W_R$-aspherical presheaf to a cosieve on which it is nonempty is totally $W_R$-aspherical, after which part~(a) applies; the proof of the restriction statement is not reproduced here. For $W_\infty$-dwy\'erienne sieves the question is undecided. It comes down to the question left open in the commentary at ms.~111: does a presheaf all of whose direct images have integrally acyclic categories of elements have only weakly contractible ones?
\end{dcomment}


\dgrblock{dgrb545}{The endpoint product is not a cofibration}
\begin{dcomment}{development}{p114}{The endpoint product is not a cofibration}
The class asked for in b) does not exist for $W_\infty$: the inclusion $\alpha$ built from the simplest sieve is not a $W_\infty$-cofibration at all. This is stronger than the observation that $\alpha$ is not a Dwyer immersion.

\textbf{Theorem.}\ours
\textit{Let $n\geq2$, $I=\Phi_n$, $\varepsilon=\{0,2n\}\subset I$, and let $\alpha:A\hookrightarrow P$ be the inclusion of $A=(\{0\}\times I)\cup([1]\times\varepsilon)$ in $P=[1]\times I$, deduced from $\{0\}\hookrightarrow[1]$. Let $C$ be the four-object category of the commentary at ms.~100, of the homotopy type of $S^2$.}
\begin{enumerate}[label=(\alph*)]
\item \textit{If $W$ is a fundamental localizer for which a category of the homotopy type of $S^3$ is not $W$-aspherical (for instance $W_\infty$, $W_R$ with $R\neq0$, or the $m$-equivalences with $m\geq3$), then $\alpha$ is not a $W$-cofibration.}
\item \textit{If $C$ is not $W$-aspherical (which includes the $2$-equivalences), then $\alpha$, which belongs to $W$, is not a trivial $W$-cofibration.}
\end{enumerate}

\textit{Proof.}
Write $l_k=(0,k)$ and $u_k=(1,k)$. The objects of $A$ are $l_0,\dots,l_{2n},u_0,u_{2n}$; its nonidentity arrows are those of the lower copy of $I$ and $l_0\to u_0$, $l_{2n}\to u_{2n}$, and no two of them are composable. Both $A$ and $P$ have contractible classifying spaces, so $\alpha\in W_\infty\subseteq W$.

Recall that $C$ has objects $x,z,a,b$ and arrows $\alpha',\beta':x\to z$, $r:z\to a$, $s:z\to b$ with $r\alpha'=r\beta'$, $s\alpha'=s\beta'$. Let $D$ be either $C$ or the category $C^{\triangleright}$ obtained by adjoining a terminal object $\omega$. Define $p:A\to D$ by
\[
 p(l_0)=z,\qquad p(u_0)=b,\qquad p(l_k)=a\ (1\leq k\leq2n),\qquad p(u_{2n})=a,
\]
sending $l_0\to l_1$ to $r$, $l_0\to u_0$ to $s$ and all other arrows to $\mathrm{id}_a$. Let $Q=P\amalg_AD$. Since $A$ is a sieve in $P$, the category $Q$ consists of $D$, of the objects $u_1,\dots,u_{2n-1}$ with the arrows of the upper copy of $I$ between them, and of the crossing sets $Q(v,u_k)=(p_!H_k)(v)$, where $H_k=\operatorname{Hom}_P(-,u_k)|_A$ (Theorem~1).

The presheaf $H_1$ is the final presheaf of the sieve $A_1=\{l_0,l_1,l_2,u_0\}$ of $A$ and is empty elsewhere. So $(p_!H_1)(v)=\pi_0(v\backslash A_1)$, the comma category being formed with the restriction $A_1\to D$ of $p$. For $v=z$ it is $A_1$ itself, a tree; for $v=a$ it is $l_2\to l_1$; for $v=b$ it is the point $u_0$; for $v=x$ it has the five objects $(l_0,\alpha'),(l_0,\beta'),(l_1,r\alpha'),(l_2,r\alpha'),(u_0,s\alpha')$ and is connected; for $v=\omega$ it is empty. Hence the object $u_1$ receives exactly one arrow from each object of $C$ and none from $\omega$. For $2\leq k\leq2n-1$ the category $A/u_k$ is nonempty and connected and $p$ is constant on it with value $a$, so $p_!H_k=D(-,a)$; the arrows between $u_2,\dots,u_{2n-1}$ induce identities of $D(-,a)$, and $u_2\to u_1$ induces the unique map $D(-,a)\to p_!H_1$.

Let $R\subset Q$ be the full subcategory on $D$ and $u_1$. The functor $\rho:Q\to R$ that is the identity on $R$, sends $u_k$ to $a$ for $k\geq2$, is the identity on the crossing sets $D(v,a)$, sends the arrows between these $u_k$ to $\mathrm{id}_a$ and $u_2\to u_1$ to the arrow $a\to u_1$, satisfies $\rho|_R=\mathrm{id}_R$; and the arrows $a\to u_k$ corresponding to $\mathrm{id}_a$ are the components of a natural transformation from $\rho$, followed by the inclusion of $R$, to $\mathrm{id}_Q$. So $BQ\simeq BR$.

(a) Take $D=C^{\triangleright}$. Then $p$ is a functor between categories with contractible classifying spaces, so $p\in W_\infty\subseteq W$; it is a morphism of $A\backslash\mathrm{Cat}$ from $\mathrm{id}_A$ to $p$, and its image under cobase change along $\alpha$ is $P\to Q$. The category $R$ is $C$ with the two terminal objects $\omega$ and $u_1$ adjoined, with no arrow between them: its nerve is the union of two cones on $NC$, so $BR\simeq\Sigma BC\simeq S^3$. Thus $P$ is $W$-aspherical and $Q$ is not, and $P\to Q$ is not in $W$.

(b) Take $D=C$. Then $R$ is $C$ with one terminal object adjoined and $BQ\simeq BR$ is contractible. The cobase change $C\to Q$ of $\alpha$ along $p$ goes from a category that is not $W$-aspherical to one that is, so it is not in $W$.\hfill$\square$

For $n=2$ the category $Q$ of (a) has $8$ objects and $30$ arrows. Part~(b) also excludes the cofibration property whenever a $W$-cofibration that belongs to $W$ is a trivial $W$-cofibration. For $W_\infty$ this is \citet[Lemma~1.10]{BataninBerger2017}, applied to Thomason's model structure, which is left proper; there it adds nothing to (a).

\textbf{Corollary (no class of the kind requested).}\ours
\textit{Let $W$ be as in (a) and $n\geq2$. If a class $\mathcal K$ of sieves consists of $W$-cofibrations and contains, with $i:U\hookrightarrow X$, the inclusion of $(U\times I)\cup(X\times\varepsilon)$ in $X\times I$, then every member of $\mathcal K$ is an isomorphism. The same holds for $W$ as in (b) and classes of trivial $W$-cofibrations.}

\textit{Proof.}
$W$-cofibrations, and trivial $W$-cofibrations, are stable under retracts in the category of arrows, because a retract of a functor in $W$ is in $W$. Let $i:U\hookrightarrow X$ be in $\mathcal K$. If there is an arrow $u\to x$ with $u\in U$, $x\notin U$, then $\{0\}\hookrightarrow[1]$ is a retract of $i$ (by the functor $X\to[1]$ with fibre $U$ over $0$, and the functor $[1]\to X$ choosing that arrow), the endpoint product $\alpha$ of $\{0\}\hookrightarrow[1]$ is a retract of that of $i$, and the theorem is contradicted. If there is no such arrow, then $X=U\amalg V$; if $V\neq\varnothing$, the endpoint product of $i$ is the inclusion of $(U\times I)\amalg(V\times\varepsilon)$ in $(U\amalg V)\times I$, which is in $\mathcal K$ and has the crossing arrow $(v,0)\to(v,1)$, so the first case applies to it.\hfill$\square$

At the other extreme, for $W_0$ every functor is a $W_0$-cofibration, because $\pi_0:\mathrm{Cat}\to\mathrm{Set}$ preserves pushouts, and all sieves form a class with the two properties. The same is true for $W_1$ and for every fundamental localizer containing $W_1$: the pushout along a functor injective on objects has the same fundamental groupoid as the cofibred category of the span; this proof is not reproduced here. By \citet[Proposition~9.2.17]{Cisinski2006} a fundamental localizer contains $W_1$ exactly when a category of the homotopy type of $S^2$ is aspherical for it. So for the localizers of homotopy theory proper the endpoint product has to be replaced by a homotopy-invariant construction, which is what the integration of the following pages provides.
\end{dcomment}


\subsection*{2.5\quad Complément : critère d'asphéricité totale d'une catégorie $W$-propre sur une autre.}

\dgrblock{dgrb546}{Why the interval test exceeds Dwyer maps}
\begin{dcomment}{examples}{p115:Complément : critère d'asphéricité totale}{Why the interval test exceeds Dwyer maps}
The preceding paragraph asks for enough cofibrations to build homotopies with prescribed endpoints.  Its product with the zigzag interval is therefore a structural test, rather than an incidental counterexample.  We can make the obstruction completely visible in a finite ordered set.

Write $P=[1]\times\Phi_n$ and
\[
 A=(\{0\}\times\Phi_n)\cup([1]\times\{0,2n\}).
\]
Here $[1]=\{0<1\}$, and the even vertices of $\Phi_n$ map to their adjacent odd vertices.  The inclusion $A\subset P$ is a sieve: a predecessor of an object in its bottom row is again in that row, and a predecessor of either of its two top endpoints is already in $A$.  Every object $(1,k)$ receives a map from $(0,k)$, so the cosieve generated by $A$ is all of $P$.

For a full inclusion of ordered sets, the existence of a right adjoint is particularly concrete.  If $r$ is right adjoint to $A\hookrightarrow P$, then $r(p)$ must be the greatest element of $\{a\in A:a\leq p\}$.  At $p=(1,1)$, this set has the two incomparable maximal elements $(1,0)$ and $(0,1)$; for $n\geq2$ the other top endpoint cannot lie below $p$.  There is consequently no greatest element, hence no right adjoint.  This is exactly the obstruction in the definition of a Dwyer map, since the relevant cosieve is already $P$.  The original inclusion $\{0\}\hookrightarrow[1]$ is Dwyer, with right adjoint the constant functor to $0$.  Thus the endpoint product operation genuinely leaves that class.

The definition and generated-cosieve criterion are those of \citet[Definition~4.1 and Lemma~5.2]{Thomason1980}.  There is a second, distinct limitation in the history of this class: retracts of Dwyer maps need not be Dwyer.  \citet[pp.~228--231, especially Proposition~2]{Cisinski1999Dwyer} gives the counterexample and repairs the properness argument for Thomason's model structure using pseudo-Dwyer maps.  The interval calculation above concerns a different closure operation.  Both facts explain why Grothendieck seeks a larger homotopically defined class rather than treating the adjoint criterion as the final definition of a cofibration.
\end{dcomment}


Soit
\[
f : X \longrightarrow I
\]
une catégorie propre sur une catégorie d'indices $I$. Si $I$ a un objet final $i_0$, alors $X/i_0$ s'identifie à $X$, et par suite l'inclusion
\[
X_{i_0} \longrightarrow X
\]
de la fibre de $X$ en $i_0$ est $W$-asphérique \texttt{[plutôt }$W$\texttt{-coasphérique]}, et a fortiori est dans $W$. Donc $X$ est asphérique si et seulement si $X_{i_0}$ l'est. Nous voulons étendre ce résultat en un critère de \textit{totale} $W$-asphéricité, de façon précise prouver la

\textbf{Proposition 11 :} \textit{Soit $f : X \longrightarrow I$ $W$-propre, et soit $i_0$ un objet final de $I$. Si $X_{i_0}$ est totalement $W$-asphérique} (déf. 1, page 88), \textit{il en est de même de $X$.} (Sous réserve que $W$ satisfasse l'axiome de connexité W(5).)

La démonstration sera donnée plus bas, après quelques préliminaires ayant un intérêt indépendant. Tout d'abord, rappelons que si $X$ est $W$-propre sur $I$, alors il donne naissance à un foncteur
\[
\begin{aligned}
I \quad &\longrightarrow \quad \mathrm{Hot}_W\\
i \quad &\longmapsto \quad \mathrm{hot}_W(X_i)(\simeq \mathrm{hot}_W(X/i));
\end{aligned}
\]
en le composant avec le foncteur
\[
\pi_0 : \mathrm{Hot}_W \longrightarrow \mathrm{Ens}
\]
(sous réserve que l'axiome de connexité W(5) (p. 10) soit satisfait)
% foot: Version 7 décembre 2002
\printedpage{79}
% head: A. Grothendieck : Dérivateurs, Ch. XII, Édition des Universités de Montpellier II et Paris VII
\msp{page 116}
on trouve un foncteur
\[
i \longmapsto \underbrace{\pi_0(X_i)}_{\simeq\, \pi_0(X/i)} : I \longrightarrow \mathrm{Ens}.
\]
Ceci posé, on a le résultat sorital

\dgrblock{dgrb547}{Components of a proper family}
\begin{dcomment}{exposition}{p115:La démonstration sera donnée plus bas}{Components of a proper family}
Grothendieck wants to transfer a condition on every ordinary left Kan extension from the last fibre to the whole category.  A weak equivalence of categories alone would not accomplish this: the functor $p_!$ here is the underived left Kan extension between categories of set-valued presheaves.  The useful consequence of properness is that it controls the components of all the relevant comma categories.

We use $\widehat Y=\operatorname{Fun}(Y^{\mathrm{op}},\mathrm{Set})$ and write $1_X$ for the terminal presheaf.  For $p:X\to Y$ one has
\[
 (p_!1_X)(y)=\pi_0(y\backslash X),
\]
where an object of $y\backslash X$ is $(x,a:y\to p(x))$.  In this chapter, ``totally $W$-aspherical'' means that $Y/(p_!1_X)$ is $W$-aspherical for every $Y$ and every $p:X\to Y$.  This is the definition on manuscript pp.~88--89; the editors explicitly distinguish it from the earlier use of total asphericity in \emph{Pursuing Stacks} \citep[Definition~1]{Grothendieck1990XII}.  It is this presheaf condition that Proposition~11 transfers.

The calculation of components has a direct proof that also fixes its variance.  For any $f:X\to I$, without a properness hypothesis, the categories $X/i$ form a covariant diagram by postcomposition.  There is a canonical bijection
\[
 \operatorname*{colim}_{i\in I}\pi_0(X/i)\ \longrightarrow\ \pi_0(X).
 \tag{A1}
\]
Indeed, every object $x$ occurs as $(x,\mathrm{id}_{f(x)})$ in $X/f(x)$, so the map is surjective.  The class of an arbitrary $(x,a:f(x)\to i)$ is identified in the colimit with the class of $(x,\mathrm{id})$, by the transition functor associated to $a$.  If $v:x\to x'$ is a morphism of $X$, the two resulting classes are equal: move the first object along $f(v)$ into $X/f(x')$, where $v$ becomes a morphism to $(x',\mathrm{id})$.  These are precisely the relations defining $\pi_0(X)$.  This proves injectivity as well.

If $f$ is $W_0$-proper, the fibre inclusion $X_i\to X/i$ induces a bijection on components.  Transporting the covariant structure through these bijections supplies the functor $i\mapsto\pi_0(X_i)$ and turns (A1) into Proposition~12.  There is no choice of transport functors on the fibres in this argument.  Under W(5), the inclusion $W\subseteq W_0$ makes $W$-properness sufficient.  The two footnotes recommending $W_0$-properness identify the exact hypothesis needed for this calculation.

For the relative formula, apply the same argument to $y\backslash X\to I$.  Its fibre over $i$ is $y\backslash X_i$.  The projection $y\backslash Y\to Y$ is a discrete opfibration, hence proper; properness is preserved by base change and composition.  Consequently
\[
 (p_!1_X)(y)\cong
 \operatorname*{colim}_{i\in I}(p_{i!}1_{X_i})(y).
 \tag{A2}
\]
Precomposition by a map $y'\to y$ commutes with every functor used to construct this bijection.  Thus (A2) is an isomorphism of presheaves, not merely a collection of unrelated bijections.  This supplies both the naturality and the universal property asserted in Corollary~1.
\end{dcomment}


\textbf{Proposition 12 :} \textit{Soit $f : X \longrightarrow I$ $W$-propre, et supposons que $W$ satisfasse l'axiome de connexité} W(5) (sinon, la $W$-propreté ne nous sera d'aucune utilité \dots) (\footnote[81]{Il vaut mieux dire : Soit $f$ $W_0$-propre, où $W_0$ est le localiseur de la 0-connexité \dots}). \textit{On a alors une bijection canonique}
\[
\pi_0(X) \xleftarrow{\ \sim\ } \varinjlim_{I} \pi_0(X_i), \tag*{$(*)$}
\]
\textit{où la flèche est induite par les applications canoniques $\pi_0(X_i) \longrightarrow \pi_0(X)$.}

\textsc{Démonstration.} Passant à $f^\circ : X^\circ \longrightarrow I^\circ$, on est ramené à prouver le même résultat pour $X$ \textit{$W$-lisse} sur $I$, avec la seule différence que le foncteur
\[
i \longmapsto \pi_0(X_i)
\]
sur $I$ sera contravariant, i.e. un préfaisceau, ce qui offre l'avantage de l'intuition attachée au formalisme des préfaisceaux. Soit
\[
g : I \longrightarrow e, \qquad \text{d'où } gf = p : X \longrightarrow e,
\]
il est immédiat, par le calcul standard
\[
p_!(F) = \pi_0(X/F) \qquad (F \in \mathrm{Ob}\, X^\wedge)
\]
(où on a identifié $e^\wedge$ à (Ens)), qu'on a
\[
\pi_0(X) = p_!(e_X) = g_!(f_!(e_X)).
\]
Or le préfaisceau
\[
i \longmapsto \pi_0(X_i) \qquad (\simeq \pi_0(i\backslash X), \text{ car } X\ W\text{-lisse sur } I)
\]
\msp{page 117}
n'est autre que $f_!(e_X)$. D'autre part, pour un préfaisceau quelconque $G$ sur $I$, on a
\[
g_!(G) \simeq \pi_0(I/G) \simeq \varinjlim_{I} G(i),
\]
et appliquant ceci à $f_!(e_X)$, on trouve la formule $(*)$ de la proposition 12.

Revenons à la situation avec $f$ \textit{$W$-propre,} et considérons un foncteur
\[
p : X \longrightarrow Y,
\]
nous sommes intéressés à ‘calculer’
\[
p_!(e_X) = (y \longmapsto \pi_0(y\backslash X)) \in \mathrm{Ob}\, Y^\wedge.
\]
Considérons pour ceci le diagramme suivant
% foot: Version 7 décembre 2002
\printedpage{80}
% head: A. Grothendieck : Dérivateurs, Ch. XII, Édition des Universités de Montpellier II et Paris VII
\[
\begin{tikzcd}
y\backslash X_i \arrow[r] \arrow[d] & y\backslash X \arrow[r] \arrow[d] & y\backslash Y \arrow[d] \\
X_i \arrow[r] \arrow[d] & X \arrow[r, "p"] \arrow[d, "f"] & Y \\
\underset{\simeq\, e}{\{i\}} \arrow[r] & I\rlap{\quad ,} &
\end{tikzcd}
\]
où les trois carrés sont cartésiens. Comme $y\backslash Y \longrightarrow Y$ est $W$-propre, il en est de même de $y\backslash X \longrightarrow X$ par changement de base, donc du composé $y\backslash X \longrightarrow X \longrightarrow I$ par transitivité. Le diagramme montre d'autre part que la fibre de $y\backslash X$ en $i \in \mathrm{Ob}\, I$ s'identifie à $y\backslash X_i$.

\msp{page 118}
La proposition 12 nous donne donc la formule
\[
\pi_0(y\backslash X) = \varinjlim_{I} \pi_0(y\backslash X_i),
\]
laquelle formule ne peut s'empêcher d'être fonctorielle en $y$, tout comme celle de prop. 12 l'est dans $\underline{\mathrm{Prop}}_W I$. Or pour $y$ variable dans $Y$, le préfaisceau sur $Y$
\[
y \longmapsto \pi_0(y\backslash X_i)
\]
n'est autre que $p_{i\,!}(e_{X_i})$, où
\[
p_i : X_i \longrightarrow Y
\]
est le composé $X_i \longrightarrow X \xrightarrow{p} Y$ (induit par $p$). On a ainsi prouvé

\textbf{Corollaire 1} (de prop. 12). \textit{Soit $f : X \longrightarrow I$ $W$-propre} (\footnote[82]{dire $W_0$-propre \dots}), \textit{et $p : X \longrightarrow Y$ un foncteur quelconque. Soit, pour tout $i \in \mathrm{Ob}\, I$, $p_i : X_i \longrightarrow Y$ le foncteur induit sur la fibre $X_i$ de $X$ en $i$. On a alors une loi} fonctorielle \textit{naturelle}
\[
i \longmapsto p_{i\,!}(e_{X_i}) : I \longrightarrow Y^\wedge,
\]
\textit{et on a la formule}
\[
p_!(e_X) \xleftarrow{\ \sim\ } \varinjlim_{I} p_{i\,!}(e_{X_i}), \tag*{$(*)$}
\]
\textit{où la flèche est induite par les morphismes canoniques}
\[
p_{i\,!}(e_{X_i}) \longrightarrow p_!(e_X)
\]
% foot: Version 7 décembre 2002
\printedpage{81}
% head: A. Grothendieck : Dérivateurs, Ch. XII, Édition des Universités de Montpellier II et Paris VII
\msp{page 119}
\textit{induites par les inclusions $X_i \longrightarrow X$.}

La loi fonctorielle de $p_{i\,!}(e_{X_i})$ en $i$ a été explicitée en cours de démonstration.

Notons que lorsque $X$ sur $I$ est cofibrée et scindée, i.e. provient d'un foncteur
\[
i \longmapsto X_i, \quad I \longrightarrow \mathrm{Cat},
\]
alors la loi fonctorielle de $p_{i\,!}(e_{X_i})$ en $i$ s'explicite de façon particulièrement agréable : à une flèche
\[
\alpha : i \longrightarrow j, \qquad \text{d'où } \alpha_* : X_i \longrightarrow X_j,
\]
est associé
\[
p_!(\alpha) : p_{i\,!}(e_{X_i}) \longrightarrow p_{j\,!}(e_{X_j}),
\]
en notant qu'on a un diagramme
\[
\begin{tikzcd}
X_i \arrow[rr, "\alpha_*"] \arrow[dr, "p_i"', ""{name=A}] & & X_j \arrow[dl, "p_j", ""'{name=B}] \\
& Y &
\arrow[from=A, to=B, bend left, "\lambda_\alpha"] % [resolved: see DER_XIIb-a_corrections.csv]
\end{tikzcd}
\]
avec une flèche de commutativité
\[
\lambda_\alpha : p_i \longrightarrow p_j\alpha_*
\]
(pas \texttt{[un]} isomorphisme en général), d'où
\[
p_{i\,!}(e_{X_i}) \xrightarrow{(\lambda_\alpha)_!} p_{j\,!}\underbrace{\alpha_{*\,!}(e_{X_i})}_{\longrightarrow\, e_{X_j}} \longrightarrow p_{j\,!}(e_{X_j}),
\]
plus de doute c'est lui! (\footnote[83]{Cette flèche \texttt{[la flèche }$(\lambda_\alpha)_!$\texttt{]} devrait être explicitée avec soin!}).

\dgrblock{dgrb548}{The comparison map requested in footnote 83}
\begin{dcomment}{exposition}{p119:plus de doute c'est lui!}{The comparison map requested in footnote 83}
The point of the split-opfibration calculation is to see the transition map before passing to components.  This also answers the request in footnote~83 without relying on a convention for mates of adjunctions.

Let $\alpha:i\to j$, let $\alpha_*:X_i\to X_j$ be the chosen transport, and let $\lambda_x:p_i(x)\to p_j(\alpha_*x)$ be the image under $p$ of the chosen cocartesian arrow over $\alpha$.  For every $y\in Y$ there is a functor
\[
 y\backslash X_i\longrightarrow y\backslash X_j,
 \qquad (x,a)\longmapsto(\alpha_*x,\lambda_x a).
 \tag{A3}
\]
A morphism $v:(x,a)\to(x',a')$ is sent to $\alpha_*v$.  This is a morphism of the target comma category because
\[
 p_j(\alpha_*v)\lambda_x a
   =\lambda_{x'}p_i(v)a=\lambda_{x'}a'.
\]
Passing to components gives the displayed map $p_!(\alpha)$ in the manuscript.  It is natural in $y$ because both routes for a map $b:y'\to y$ replace $a$ by $ab$.  For two composable base arrows, the split identities and the cocartesian composition law identify the two composites of (A3).  Hence these maps form the asserted covariant diagram of presheaves.

To recover the factorization written by Grothendieck, first keep the intermediate object $\alpha_{*!}1_{X_i}$ and apply the natural transformation of left Kan extensions induced by $\lambda:p_i\Rightarrow p_j\alpha_*$.  Then use the counit-induced map $\alpha_{*!}1_{X_i}\to1_{X_j}$.  The first arrow is the mate of $(p_j\alpha_*)^*\Rightarrow p_i^*$, whose component on a presheaf $F$ is $F(\lambda_x)$.  Formula (A3) proves that this mate has the stated direction.  An invertible $\lambda$ is unnecessary.

If $I$ has terminal object $i_0$, the colimit in (A2) is its value at $i_0$.  We therefore obtain, naturally for every $p:X\to Y$,
\[
 p_!1_X\cong p_{i_0!}1_{X_{i_0}}.
\]
If $X_{i_0}$ is totally $W$-aspherical, the category of elements of the right-hand presheaf is $W$-aspherical, which proves Proposition~11.  The printed $p_{i_0!}(X_{i_0})$ on the next page is a source slip: the argument must be $1_{X_{i_0}}$.  Taking $X=Z\times I$ proves Corollary~3, since the product projection is a split opfibration with final fibre $Z$.  The argument in fact requires only $W_0$-properness of $f$ together with the stated total $W$-asphericity of its final fibre.
\end{dcomment}


\msp{page 120}
\textbf{Corollaire 2 :} \textit{Sous les conditions du corollaire 1, supposons que $I$ ait un objet final $i_0$. Alors}
\[
(p_{i_0})_!(e_{X_{i_0}}) \longrightarrow p_!(e_X)
\]
\textit{est un isomorphisme dans $Y^\wedge$.}

\dgrblock{dgrb549}{The transition map on representatives}
\begin{dcomment}{exposition}{p119}{The transition map on representatives}
The footnote asks that the arrow induced by $\lambda_\alpha$ be written out. This matters because the triangle is only laxly commutative: replacing its natural transformation by an equality would erase the data that moves a fibre over $i$ to a fibre over $j$.

Suppose $X\to I$ is the split opfibration associated to a strict functor $D:I\to\Cat$, and write $X_i=D(i)$. An arrow $\alpha:i\to j$ gives $\alpha_*:X_i\to X_j$ and chosen opcartesian arrows $x\to\alpha_*x$ in $X$. Applying $p:X\to Y$ gives a natural transformation
\[
 \lambda_\alpha:p_i\longrightarrow p_j\alpha_*.
\]
For $y\in Y$, an object of $y\backslash X_i$ is a pair $(x,\beta:y\to p_i(x))$. Define
\[
 L_{\alpha,y}(x,\beta)
   =\bigl(\alpha_*x,\lambda_{\alpha,x}\beta\bigr)
     \quad\text{in }y\backslash X_j.                 \tag{C3}
\]
If $h:x\to x'$ is a morphism in the first comma category, so that $p_i(h)\beta=\beta'$, its image is $\alpha_*h$. Naturality of $\lambda_\alpha$ gives
\[
 p_j(\alpha_*h)\lambda_{\alpha,x}\beta
 =\lambda_{\alpha,x'}p_i(h)\beta
 =\lambda_{\alpha,x'}\beta',
\]
which verifies the comma-category condition. Hence (C3) is a functor. Taking connected components yields the value at $y$ of the required presheaf map
\[
 (p_i)_!e_{X_i}\longrightarrow(p_j)_!e_{X_j}.
\]
Precomposition with a map $y'\to y$ commutes strictly with (C3), proving naturality in $y$.

This also identifies the two arrows in the manuscript's displayed composite. First keep the object $x$ in $X_i$ and replace $\beta$ by $\lambda_{\alpha,x}\beta$. That is the map induced by $\lambda_\alpha$ from $(p_i)_!e_{X_i}$ to $(p_j\alpha_*)_!e_{X_i}$. Next apply $\alpha_*$ to $x$; under transitivity of Kan extension this is the counit map
\[
 (p_j)_!(\alpha_*)_!e_{X_i}\longrightarrow(p_j)_!e_{X_j}.
\]
Their composite is exactly (C3), so its direction is fixed without an appeal to an unspecified mate convention.

Functoriality in $i$ follows from the split identities. For $\alpha:i\to j$ and $\gamma:j\to k$ one has
\[
 (\gamma\alpha)_*=\gamma_*\alpha_*,
 \qquad
 \lambda_{\gamma\alpha,x}
 =\lambda_{\gamma,\alpha_*x}\lambda_{\alpha,x}.
\]
Substitution in (C3) gives $L_{\gamma,y}L_{\alpha,y}=L_{\gamma\alpha,y}$; identity arrows act identically. Thus the transition maps constitute an actual functor $I\to\widehat Y$, not merely objectwise maps in a localization. With a non-split opfibration the comparison isomorphisms of a cleavage make the same formulas agree after taking components.

For a general $W$-proper functor, still assuming W(5), the manuscript obtains these transitions through the comma-category construction. The projection $y\backslash Y\to Y$ is a discrete opfibration, hence proper; base change and composition make $y\backslash X\to I$ proper. Its fibre over $i$ is $y\backslash X_i$. Proposition~12 applied there gives
\[
 \pi_0(y\backslash X)
   \cong\varinjlim_{i\in I}\pi_0(y\backslash X_i).
\]
The construction is natural in $y$ and so proves the presheaf isomorphism of Corollary~1. If $I$ has a terminal object, the colimit is its value at that object, recovering (C1). In the product case $X=Z\times I$, the terminal fibre is $Z$, which proves Corollary~3 as well.
\end{dcomment}


D'où aussitôt la proposition 11 : si $X_{i_0}$ est totalement $W$-asphérique, alors $X$ l'est, car pour tout $p : X \longrightarrow Y$, $p_!(e_X) \simeq p_{i_0\,!}(X_{i_0})$, et ce dernier élément de $Y^\wedge$ est $W$-asphérique % [sic] p_{i_0 !}(X_{i_0}) -> p_{i_0 !}(e_{X_{i_0}}), elsewhere the argument of p_{i!} is e_{X_i}
par hypothèse sur $X_{i_0}$.

Notons le cas particulier :

\textbf{Corollaire 3.} \textit{Soit $Z$ un objet totalement $W$-asphérique dans} Cat. \textit{Alors son produit $Z\times I$ par toute $I$ dans} Cat \textit{ayant un objet final est aussi totalement $W$-asphérique.}
% foot: Version 7 décembre 2002
\printedpage{82}
% head: A. Grothendieck : Dérivateurs, Ch. XII, Édition des Universités de Montpellier II et Paris VII

\dgrblock{dgrb550}{Total asphericity depends on more than homotopy type}
\begin{dcomment}{development}{p120:ayant un objet final est aussi totalement}{Total asphericity depends on more than homotopy type}
The transfer result just proved has a stronger content than invariance of the classifying space.  The finite construction developed below, in ``A strict cocone can repair the wrong pushout'', makes this distinction explicit.

\textbf{Corollary (failure of homotopy invariance).}\ours
Total $W_\infty$-asphericity in the sense of this chapter is not invariant under weak homotopy equivalences of categories.

\emph{Proof.}
The cited construction supplies finite face posets $E,S$ and a functor $f:E\to S$ such that $BE$ is contractible, $BS\simeq S^2$, and every $E/s$ is nonempty and connected.  Their construction and all these properties are proved there.  Set $p=f^{\mathrm{op}}:E^{\mathrm{op}}\to S^{\mathrm{op}}$.  For each $s\in S$,
\[
 (p_!1_{E^{\mathrm{op}}})(s)
   =\pi_0(s\backslash E^{\mathrm{op}})
   =\pi_0((E/s)^{\mathrm{op}})=\{*\}.
\]
Thus $p_!1_{E^{\mathrm{op}}}$ is the terminal presheaf on $S^{\mathrm{op}}$, whose category of elements is $S^{\mathrm{op}}$.  Its classifying space is a two-sphere, so $E^{\mathrm{op}}$ is not totally $W_\infty$-aspherical.  Nevertheless $E^{\mathrm{op}}\to e$ is a weak homotopy equivalence.  The category $e$ is totally $W_\infty$-aspherical: along any functor selecting $y\in Y$, its terminal presheaf is sent by left Kan extension to the representable $Y(-,y)$, and the category of elements $Y/y$ has a terminal object.\hfill$\square$

This corollary is a consequence of the explicit finite calculation below, not an assertion that the underlying topological phenomenon is new.  It explains why Proposition~11 uses properness to identify ordinary Kan extensions rather than inferring total asphericity merely from $X_{i_0}\to X\in W$.
\end{dcomment}


\msp{page 121}

\dgrblock{dgrb551}{A span is remembered by its transition arrows}
\begin{dcomment}{exposition}{p120}{A span is remembered by its transition arrows}
The next construction turns an arbitrary span into a category in which the gluing instructions remain visible as morphisms.  The ordinary pushout identifies their images.  Integration keeps those images separate and inserts comparison arrows between them.  This small change is what retains the homotopy information of the gluing.

Write $\Phi=(1\to0\leftarrow2)$ and let the span $X_1\xleftarrow{\alpha_1}X_0\xrightarrow{\alpha_2}X_2$ be the functor $\Phi^{\mathrm{op}}\to\Cat$.  Its fibred Grothendieck construction $\mathcal X$ has the disjoint union of the three sets of objects.  Within a fibre the morphisms are unchanged.  The only morphisms between distinct fibres go from a side fibre to the middle fibre, and are
\[
 \operatorname{Hom}_{\mathcal X}(a,x)
     =\operatorname{Hom}_{X_k}(a,\alpha_kx),
 \qquad a\in X_k,\ x\in X_0,\ k=1,2.
 \tag{A4}
\]
Composition on the source side is ordinary composition in $X_k$; composition with $v:x\to x'$ on the target side uses $\alpha_k(v)$.  These rules satisfy associativity because $\alpha_k$ is a functor.  They also show directly that $\mathcal X\to\Phi$ is a split fibration.

The identity of $\alpha_kx$ determines a morphism
\[
 i_k\alpha_k(x)\longrightarrow i_0(x).
\]
Naturality follows from the composition rule in (A4).  These two transformations are the missing sides of the apparently noncommutative square on manuscript p.~123.  Their sources and targets are usually different objects, so the square is not strictly commutative.  Natural transformations give homotopies on nerves; for every fundamental localizer the resulting functors become equal after localization.  The assertion that the square commutes in $\mathrm{Hot}_W$ is therefore correct, but does not replace the two transformations that constructed it.

This is the variance of the \emph{smooth} integral used in this section.  An opfibred construction of the covariant span places the transition arrows in the opposite direction.  Grothendieck compares those two constructions later in \S3.3.  Reversing only the base while leaving the fibres untouched is not an innocuous notational change.

For $W=W_\infty$, realization of the nerve of $\mathcal X$ is a homotopy pushout of $BX_1\leftarrow BX_0\to BX_2$, where $BC=|NC|$.  The closed-piece calculation below gives a direct proof in this particular shape.  The general Grothendieck-construction theorem is due to \citet[Theorem~1.2]{Thomason1979}; a primary modern proof and an explicit comparison map are given in \citet[\S3, equation~(9) and Theorem~3]{Cegarra2020}.  The manuscript expresses the corresponding fact intrinsically as a left Kan extension in the category of coherent diagrams.  Its footnote~84 correctly retains the stronger localization hypothesis needed by the preceding construction of that adjoint.
\end{dcomment}


\section*{3\quad Sommes amalgamées et intégration\quad \begin{tikzcd}[ampersand replacement=\&, row sep=small, column sep=small, baseline=(\tikzcdmatrixname-2-1.base)]
\& X_1 \\
\int X_0 \arrow[ur] \arrow[dr] \& \\
\& X_2
\end{tikzcd}}

\subsection*{3.1\quad Réduction canonique au cas d'un diagramme d'immersions fermées}

Considérons un diagramme
\[
\begin{tikzcd}[row sep=small, column sep=small]
& X_1 \\
X_0 \arrow[ur, "\alpha_1"] \arrow[dr, "\alpha_2"'] & \\
& X_2
\end{tikzcd}
\tag*{$(*)$}
\]
dans Cat. Nous pouvons le regarder comme un contrafoncteur
\[
\Phi^\circ \longrightarrow \mathrm{Cat},
\]
où
\[
\Phi = \left( \vcenter{\hbox{\begin{tikzcd}[row sep=small, column sep=small]
& 1 \arrow[dl] \\
0 & \\
& 2 \arrow[ul]
\end{tikzcd}}} \right),
\]
donc comme définissant une catégorie fibrée (scindée) $\mathscr{X}$ sur $\Phi$. Il est parfois commode de représenter cette catégorie $\mathscr{X}$ par le ‘signe d'integration’ % [sic] integration -> intégration, accent missing in the print
\[
\begin{tikzcd}[row sep=small, column sep=small]
& X_1 \\
\int X_0 \arrow[ur] \arrow[dr] & \\
& X_2,
\end{tikzcd}
\]
et la regarder comme une sorte de ‘somme amalgamée homotopique’ de $X_1$ et $X_2$ sous $X_0$, ou comme une ‘limite inductive homotopique’ du diagramme de départ $(*)$.

\msp{page 122}
Il faut entendre par là que (pour $W \subseteq \mathrm{Fl}(\mathrm{Cat})$ donné) l'image $\mathrm{hot}_W(\mathscr{X})$ de $\mathscr{X}$ dans $\mathrm{Hot}_W$ a un ‘sens homotopique’ remarquable en termes de $(*)$, de façon précise, le diagramme $(*)$ définit un objet $\xi$ de $\mathrm{HOT}_W(\Phi)$ (qui est une catégorie de fractions de la catégorie de tous les diagrammes du type $(*)$ \dots), et si
\[
\pi : \Phi \longrightarrow e
\]
désigne la projection absolue de $\Phi$, on a
\[
\pi_!(\xi) = \mathrm{hot}_W(\mathscr{X}),
\]
% foot: Version 7 décembre 2002


\dgrblock{dgrb552}{Integration as a calculus of gluing}
\begin{dcomment}{context}{p121}{Integration as a calculus of gluing}
Grothendieck's integral is designed to make a family and its transition maps into one geometric object. Grothendieck's discussion of integration in Chapter~VI, manuscript pages~81 and~146, gives the general background. Here the additional issue is quite specific: the total category must retain enough information to decide whether a given cocone is a homotopy pushout.

In \emph{Pursuing Stacks}, \S69, \citet{Grothendieck1983} connects diagram categories to the trace problems of SGA~5, then describes integration by a total cofibred category. He also singles out presheaf categories as a setting where actual amalgamated sums are available. These two concerns meet in the present pages. The integral supplies the homotopical comparison; the last theorem shows that a pushout of presheaves along a monomorphism gives an actual categorical pushout with the required homotopy type. The intervening analysis of open immersions is the bridge between the two constructions.

The letter to Thomason of 2 April 1991 gives a closely related explanation of why one diagram category cannot suffice: a derived category ``ne vient jamais seule'' \citep[\S2]{Grothendieck1991Letter}. In \S1 of the same letter, Grothendieck places weak equivalences and localization before the auxiliary classes used in model structures. This explains the order of the present argument. He constructs a comparison map first and asks whether it belongs to $W$; cofibration criteria are then ways of proving that assertion.

The source is also a record of the distinctions being worked out. A natural transformation provides a homotopy without becoming an equality of functors; a weak equivalence of total categories need not preserve every ordinary Kan extension; a gluing theorem need not reflect equivalences backwards. The proofs below keep the transition arrows, the presheaf adjunctions, and the excision hypotheses visible at precisely the steps where each is used.
\end{dcomment}


\dgrblock{dgrb553}{Integration retains the homotopy of the span}
\begin{dcomment}{exposition}{p122}{Integration retains the homotopy of the span}
The shift to integration changes what is being glued. An ordinary pushout identifies the common category immediately. The integrated category instead retains a copy of that category and arrows connecting it to both sides. These arrows provide the homotopies needed for a derived colimit.

For the span $X_1\xleftarrow{\alpha_1}X_0\xrightarrow{\alpha_2}X_2$, the manuscript uses the indexing category $\Phi=(1\to0\leftarrow2)$ and a contravariant functor on it. The associated fibred category $\mathscr X$ has objects $(k,x)$ with $x\in X_k$. It contains each $X_k$ fully, and its additional morphisms are
\[
 \Hom_{\mathscr X}((k,x),(0,z))
   =\Hom_{X_k}(x,\alpha_kz),\qquad k=1,2.
\]
There are no arrows between the $1$- and $2$-fibres, nor from the $0$-fibre to either of them. Composition with an arrow in $X_k$ is precomposition; composition with an arrow $z\to z'$ in $X_0$ is postcomposition by its image under $\alpha_k$. These rules give the split Grothendieck construction with precisely the variance used by Grothendieck. This continues the discussion of variance in Chapter~VI, manuscript p.~81, rather than replacing it by a covariant convention without explanation.

For this span one can prove the homotopy-pushout assertion directly. Let $M_k$ be the full subcategory of $\mathscr X$ on its $0$- and $k$-fibres. It is a categorical mapping cylinder for $\alpha_k$. The functor
\[
 r_k:M_k\longrightarrow X_k
\]
is the identity on $X_k$ and sends $z\in X_0$ to $\alpha_kz$. If $s_k:X_k\hookrightarrow M_k$ is the inclusion, then $r_ks_k=\id$, and the canonical arrow $\alpha_kz\to z$ in $M_k$ gives a natural transformation $s_kr_k\to\id_{M_k}$, equal to the identity on $X_k$. Thus $Br_k$ is a homotopy equivalence, and its restriction to $BX_0$ is $B\alpha_k$.

Every composable string in $\mathscr X$ lies in $M_1$ or $M_2$, and their intersection is exactly $X_0$. Therefore
\[
 N\mathscr X=NM_1\amalg_{NX_0}NM_2.
\]
Both maps out of $NX_0$ are monomorphisms, so this is a homotopy pushout of simplicial sets. The two homotopy equivalences $Br_k$, together with the identity on $BX_0$, now identify its homotopy type with
\[
 B\mathscr X\simeq
 \operatorname{hocolim}\bigl(BX_1\leftarrow BX_0\to BX_2\bigr).
\]
This is the span instance of \citet[Theorem~1.2]{Thomason1979}, whose statement is also recorded by \citet[Corollary~26.9]{ChacholskiScherer2001}; the direct proof records the orientation and the connecting morphisms needed in these pages.

For example, take $X_0$ to be two discrete objects and $X_1=X_2=e$. The ordinary pushout is $e$. The integrated category has four objects and four nonidentity arrows: each of the two side objects maps to each of the two middle objects. Its nerve is a square with no two-dimensional nondegenerate simplices, hence has realization $S^1$. The integration remembers the circle formed by the two independent identifications; the ordinary pushout has forgotten it.

The final formula $\pi_!(\xi)=\mathrm{hot}_W(\mathscr X)$ expresses this construction in the derivator already under discussion. It is a canonical isomorphism in the localized category, rather than a literal equality of chosen representatives. For $W_\infty$ the preceding argument gives its concrete meaning. For another $W$ the formula presupposes the standing hypotheses under which the diagram localization and its homotopy Kan extension are available; an arbitrary class of arrows would not by itself justify that notation. The advertised reduction to a diagram of closed immersions begins here; it is not completed on manuscript page~122.
\end{dcomment}
\begin{otherlanguage}{english}
\dgroverview{XIIb-a}{Chapter XII \textperiodcentered\ \textup{XIIb-a} \textperiodcentered\ Overview}
\begingroup
\selectlanguage{english}
\noindent\textbf{Attachments, ordinary direct images, and homotopy.}

These pages, Chapter~XII, section~V, manuscript pp.~59--122, begin with a concrete problem: how to attach one category to another along a full subcategory, and how to recognize when the attachment respects weak equivalences. The first part answers the categorical question completely. The second discovers how much more is needed for the homotopical question. Along the way Grothendieck introduces total asphericity, isolates the Dwyer condition, and reaches the Grothendieck construction as a replacement for an ordinary pushout. The text is exploratory in a mathematically consequential sense. Some tentative arguments are explicitly retracted; one proposed theorem survives in the text even though its asserted reverse implication is false. The commentary keeps these different situations separate.

An ``open'' subcategory $U\subseteq X$ is a sieve: an arrow whose target lies in $U$ has its source there as well. Its complementary full subcategory $T$ is a cosieve. There are no arrows from $T$ to $U$, so the information not already contained in these two categories is exactly the bifunctor
\[
 H:U^{\mathrm{op}}\times T\longrightarrow\mathrm{Set},
 \qquad H(u,t)=\Hom_X(u,t).
\]
Its two functorial actions encode composition on either side of a crossing arrow. Conversely any such bifunctor constructs a category on the disjoint union of the object sets. This is the collage of a profunctor. Maps and natural transformations out of the collage have equally explicit descriptions. The formulas on manuscript p.~61 must be read relatively: changing $U$ and $T$ by pullback of $H$ gives a cartesian lifting over the category of pairs of categories, not the unqualified representability asserted in the printed displays. The corrected universal property is proved in the commentary.

This description also explains the relation with open--closed gluing of topoi. A presheaf on the collage is a pair of presheaves on $U$ and $T$, together with the restriction data across $H$. The gluing functor is
\[
 \varphi(F)(t)=\Hom_{\widehat U}(H(-,t),F).
\]
It preserves limits and is right adjoint to the cocontinuous extension of $t\mapsto H(-,t)$. The converse recovers the collage from an accessible gluing functor preserving small limits. This is already the construction of SGA~4, Expos\'e~IV, Exercise~9.5.10; the comparison is exact, including the characterization of the resulting topos as a presheaf topos \citep{SGA4I1972}. Natural transformations between right adjoints reverse the direction of transformations between their left adjoints, a convention that must be retained when stating the category-level equivalence. The historical connection helps explain why the categorical attachment problem first appears here in the language of an open and its closed complement.

The advantage of the collage is that a pushout can be computed by an ordinary Kan extension. Given $p:U\to U'$ and $r:T\to T'$, the new crossing bifunctor is
\[
 H'=\operatorname{Lan}_{p^{\mathrm{op}}\times r}H.
\]
Its pointwise coend description records a crossing arrow together with paths in the new categories to its two ends, modulo the relations imposed by composition. When only $U$ changes, this becomes
\[
 \Hom_{X'}(u',t)\cong\pi_0(u'\backslash U/t),
 \qquad X'=X\amalg_U U'.
\]
The category on the right remembers factorizations through $U$; the ordinary pushout retains only their connected components. This formula proves that $U'$ is again a sieve, that the complement remains $T$, and that the support of the gluing weight is unchanged. It is also the source of the later difficulty: taking components does not preserve the higher homotopy of a factorization category.

\medskip
\Needspace{5\baselineskip}
\noindent\textbf{The gluing theorem that survives, and the reverse implication that fails.}

If a category is the union of two sieves, every simplex of its nerve lies in one of them. Its nerve is therefore the simplicial union along the intersection, and the square is a homotopy pushout. Consequently weak equivalences on the intersection and on both pieces induce a weak equivalence on the union. This is the valid gluing direction. It does not imply the symmetric rule W(7), which asks that any two among equivalence on the intersection, equivalence on both pieces, and equivalence on the union imply the third.

Theorem~2 on manuscript p.~78 asserts that this stronger rule holds for $W_\infty$. A complete counterexample uses a connected acyclic noncontractible finite complex. Take a finite poset $A$ with that homotopy type and adjoin two incomparable terminal cone points. The two one-cone subcategories are sieves, each contractible, and their union has the homotopy type of the suspension of $BA$. The suspension is simply connected and acyclic, hence contractible, whereas $BA$ remains noncontractible. Thus the map from a cone to the union is a weak equivalence, but the opposite map from the intersection to the other cone is not. The counterexample refutes both the reverse part of W(7 bis) and the full W(7); it does not merely exhibit a missing proof.

The commentary gives the classical two-cell input explicitly, verifies its homology and a nontrivial fundamental-group quotient, and passes to a finite simplicial model and its face poset \citep[Example~2.38 and Theorem~2C.5]{Hatcher2002}. A further development supplies finite tests against every bounded homotopy-group version of the axiom. It also explains why the proposed Mayer--Vietoris argument succeeds for fixed-coefficient homology equivalences: the three mapping-cone complexes form a short exact sequence. The failure occurs in passing from that homological statement to the asserted unstable homotopical converse.

This finding changes how the subsequent pages must be read. Assertions deduced under W(7) or W(7 bis) remain conditional assertions. Their application to $W_\infty$ cannot use Theorem~2. In particular, the proposed one-point converse to asphericity on pp.~84--85 is not an unconditional theorem about weak homotopy equivalences. The audit records the actual hypotheses and, where possible, supplies a direct proof of a useful consequence that avoids the false reverse principle.

\medskip
\Needspace{5\baselineskip}
\noindent\textbf{Total asphericity measures ordinary images.}

For a presheaf $F$ on $U$, total $W$-asphericity means that the category of elements of $p_!F$ is $W$-aspherical for every functor $p:U\to U'$. These are ordinary set-valued Kan extensions. For a category $V$, the definition applies to its terminal presheaf. It is substantially stronger than $V\to e\in W$, and it is not interchangeable with other uses of ``totally aspherical'' in the literature on test categories.

The relevant ordinary factorization sends $f:V\to V'$ to
\[
 V\longrightarrow V'/f_!e_V\longrightarrow V'.
\]
The second map is a discrete fibration; the first is final in the ordinary categorical sense. Its undercategories are nonempty and connected, rather than necessarily contractible. This is the comprehensive factorization. The commentary proves its universal property and relates it to the treatment of Street--Walters in Kelly's account \citep[\S4.7]{Kelly1982}. It yields a useful formulation developed here: a category is totally $W$-aspherical exactly when every target of an ordinary final functor out of it is $W$-aspherical. In particular the class of totally aspherical categories is closed under ordinary final images. This reformulation identifies the quantified condition without confusing ordinary finality with homotopy finality.

Grothendieck's attempted stability under unions fails because a left Kan extension need not preserve the monomorphisms describing a cover. His attempted use of retractions does not repair that defect: a categorical retraction does not supply a splitting of the associated inclusion of presheaves. A finite example makes both issues explicit. A final functor from the contractible outgoing three-object fork lands in a four-object category whose classifying space is $S^2$. Its relevant comma category is a square, connected but not contractible. Two distinct representatives of a Kan-extension value become equal after adjoining one branch. The construction was located in Grothendieck's Appendix~XII and is credited there. It also produces the pushout failures on pp.~104--107; adjoining a further cone gives a classifying space $S^3$, so their validity can be checked independently of the disputed one-point converse.

The converse from totality of a presheaf to totality of its category of elements is established here for subterminal presheaves, as the source proposes. A general answer was not located in the primary literature checked. The same limitation applies to certain composition and classification questions posed later in these pages. An unsuccessful search is not turned into a claim of present-day openness.

\medskip
\Needspace{5\baselineskip}
\noindent\textbf{Good attachments and the passage to integration.}

A $W$-Dwyer sieve asks that every nonempty crossing weight be totally $W$-aspherical. Ordinary Kan-extension transitivity proves that this condition is preserved by pushout. The stronger version, expressed through totality of the categories of elements themselves, uses the final first map of the comprehensive factorization. Classical Dwyer immersions are the representable case: the inclusion of $U$ into the cosieve it generates admits a right adjoint. Filtered examples show why total asphericity is broader than the existence of a terminal object. The interval-product calculation on p.~114, however, shows directly why the classical adjoint condition is too rigid for every intended homotopy attachment.

The manuscript's $W$-cofibrations have since been studied as $h$-cofibrations, defined by preservation of weak equivalences under cobase change on undercategories. In left proper model categories they characterize attachments whose pushouts compute homotopy pushouts; the theory also recovers the equivalence between trivial $h$-cofibrations and couniversal weak equivalences \citep[Proposition~1.6, Corollary~1.9 and Lemma~1.10]{BataninBerger2017}. This is a precise later form of the useful gluing principle. For classical Dwyer maps, the nerve comparison is now known to be a weak categorical equivalence, with a generalization to discretely flat cofibrations \citep[Theorems~1.5--1.6]{HackneyOzornovaRiehlRovelli2024}. These later notions are compared with the definitions here without identifying their different ambient weak equivalences.

The final properness criterion is independent of the failed cube axiom. Under the necessary connectivity hypothesis, the inclusion of a terminal fibre is ordinary final. Its terminal presheaf therefore has direct image equal to the terminal presheaf of the total category. Composing Kan extensions proves that total asphericity of that fibre implies total asphericity of the total category. The commentary gives the general component-colimit formula and writes out the natural transformation requested in the footnote on p.~119, including its action on representatives and its compatibility with composition.

The text then begins a new construction. Integrating a span retains its middle category and the arrows connecting it to both sides. For two points mapping to a point on each side, the ordinary pushout is a point, whereas the integrated category has classifying space a circle. Two categorical mapping cylinders give a direct proof that the integrated span models the homotopy pushout. The final displayed identity $\pi_!\xi\cong\mathrm{hot}_W(\mathscr X)$ states the intended derivator interpretation. The announced reduction to closed immersions continues after that identity, beyond manuscript page~122.

\endgroup
\end{otherlanguage}
\dgrchunkend

% DER_XIIb-b: Grothendieck's text with the commentary set in, then the overview, inside a group that applies the chunk's own definitions
\dgrpart{1}{XIIb-b}{Section 5, part 2 of 2, manuscript pages 123--169}{XIIb-b\enspace\textperiodcentered\enspace Chapter XII, section 5, part 2 of 2\enspace\textperiodcentered\enspace manuscript pages 123--169}{Chapter XII \textperiodcentered\ \textup{XIIb-b}}
\dgrchunkbegin
% DER_XIIb-b: the definitions of the chunk's preamble, applied inside its group
\providecommand\IfDocumentMetadataT[1]{\IfDocumentMetadataTF{#1}{}}
\let\printedpage\relax\newcommand{\printedpage}[1]{\marginpar{\footnotesize p.~#1}}
\let\msp\relax\newcommand{\msp}[1]{\marginpar{\footnotesize [#1]}}
\let\Hot\relax\newcommand{\Hot}{\mathrm{Hot}}
\let\HOT\relax\newcommand{\HOT}{\mathrm{HOT}}
\let\Cat\relax\newcommand{\Cat}{\mathrm{Cat}}
\let\cX\relax\newcommand{\cX}{\mathcal{X}}
\let\sX\relax\newcommand{\sX}{\mathscr{X}}
\let\tX\relax\newcommand{\tX}{\tilde{\mathscr{X}}}
\let\Hom\relax\newcommand{\Hom}{\mathrm{Hom}}
\let\Fl\relax\newcommand{\Fl}{\mathrm{Fl}}
\let\tsX\relax\newcommand{\tsX}{\tilde{\mathscr{X}}}
\let\circled\relax\newcommand{\circled}[1]{\tikz[baseline=(char.base)]{\node[shape=circle,draw,inner sep=1pt] (char) {#1};}}
\let\uHom\relax\newcommand{\uHom}{\underline{\mathrm{Hom}}}
\let\tld\relax\newcommand{\tld}{^{\tilde{}}}
\let\hto\relax\newcommand{\hto}[1][]{\overset{#1}{\lhook\joinrel\longrightarrow}}
% Transcription repair: complete the truncated macro; the printed integral diagrams are unchanged.
\let\intdiag\relax\newcommand{\intdiag}[3]{\int\begin{tikzcd}[row sep=small,column sep=small,ampersand replacement=\&,baseline=(\tikzcdmatrixname-2-1.base)] \& #2 \\ #1 \arrow[ur]\arrow[dr] \& \\ \& #3 \end{tikzcd}}
\let\Ob\relax\newcommand{\Ob}{\operatorname{Ob}}
\let\Fibdisc\relax\newcommand{\Fibdisc}{\underline{\mathrm{Fib\ disc}}}
\let\cercle\relax\newcommand{\cercle}[1]{\tikz[baseline=(c.base)]\node[draw,circle,inner sep=1pt](c){#1};}

% ---- from Der-ChXIIb-012_p041-045.tex
\printedpage{83}
% head: A. Grothendieck : Dérivateurs, Ch. XII, Édition des Universités de Montpellier II et Paris VII
d'après le calcul connu du foncteur $\pi_!$, adjoint à gauche du foncteur $\pi^*$
\[
\HOT_W(e) \overset{\text{déf}}{=} \Hot_W \longrightarrow \HOT_W(\Phi)
\]
(\footnote[84]{Cela suppose qu'on a W(2), non seulement W(2 bis).}). Or pour tout objet $I$ dans Cat, de projection absolue $p_I : I \longrightarrow e$, le foncteur
\[
p_{I\,!} : \HOT_W(I) \longrightarrow \HOT_W(e) = \Hot_W
\]
doit être regardé comme l'équivalent `homotopique' du foncteur limite inductive $\varinjlim\limits_{I}$ des diagrammes de type $I^{\circ}$ dans la catégorie $\Hot_W$ - lesquelles limites inductives \textit{n'existent pas} dans $\Hot_W$ (sauf
\msp{page 123}
pour des catégories très particulières - les seules que je connaisse sont les catégories discrètes (qu'il faut peut-être même supposer finies en plus)).

Notons qu'on a un diagramme de foncteurs canoniques
\[
\begin{tikzcd}[column sep=tiny, row sep=small, ampersand replacement=\&]
 \& X_0 \arrow[dl, "\alpha_1"'] \arrow[dr, "\alpha_2"] \& \\
X_1 \arrow[dr, "i_1"'] \& {\shortstack{\scriptsize non\\ \scriptsize commuta-\\ \scriptsize tif!}} \& X_2 \arrow[dl, "i_2"] \\
 \& \cX, \&
\end{tikzcd}
\tag*{$(*)$}
\]
lequel diagramme n'a rien de commutatif dans Cat. \textit{Il est cependant commutatif dans} $\Hot_W$ (quel que soit le localiseur fondamental $W$), et même dans Cat\,hot, car si on désigne par
\[
i_0 : X_0 \longrightarrow \cX
\]
l'inclusion canonique, on a des homomorphismes canoniques
\[
i_1\alpha_1 \longrightarrow i_0 \longleftarrow i_2\alpha_2,
\]
ce qui montre que $i_1\alpha_1$ et $i_2\alpha_2$ sont homotopes, donc définissent la même flèche dans $\Hot_W$.

\msp{page 124}
Considérons dans $\Phi$ les deux fermés isomorphes à $\boldsymbol{\Delta}^1$
\[
\begin{aligned}
F_1 &= 1\backslash\Phi &&= \{0,1\}\\
F_2 &= 2\backslash\Phi &&= \{0,2\},
\end{aligned}
\]
on a
\[
F_1\cap F_2 = 0\backslash\Phi = \{0\}, \qquad F_1\cup F_2 = \Phi.
\]

\dgrblock{dgrb460}{A circle hidden by an ordinary pushout}
\begin{dcomment}{examples}{p123}{A circle hidden by an ordinary pushout}
The distinction between a homotopy colimit and a colimit in the homotopy category is already visible in a span of finite discrete categories.  Take $X_0=\{a,b\}$ and $X_1=X_2=e$, with both structure maps the unique functor.  Call the two side objects $l,r$.  Formula (A4) gives precisely four nonidentity arrows:
\[
 l\to a,\quad l\to b,\quad r\to a,\quad r\to b.
\]
No two nonidentity arrows are composable.  The nondegenerate nerve is consequently a graph with four vertices and four edges, namely a four-cycle.  Its realization is a circle.  One can compute this without any homotopy-colimit theorem: choose three of the edges as a spanning tree and collapse that tree; the remaining edge is one loop.  Thus $\pi_1(B\mathcal X)\cong\mathbb Z$, $H_1(B\mathcal X;\mathbb Z)\cong\mathbb Z$, and all reduced homology in other degrees vanishes.

The ordinary pushout $e\amalg_{\{a,b\}}e$ in $\Cat$ is $e$.  Indeed, a strictly compatible pair of functors from its two copies of $e$ to $Y$ chooses a single object of $Y$; both points of $X_0$ impose the same equality.  The canonical comparison $\mathcal X\to e$ therefore collapses the circle and is not in $W_\infty$.  This example will fail Definition~3 exactly where it should.

The ordinary pushout of the same diagram also exists in the homotopy category of spaces and is a point.  For a space $T$, a homotopy class $e\to T$ is a component of $T$.  Compatibility of two such classes along the discrete two-point space forces them to be the same component, so the set of compatible pairs is $\pi_0(T)$, represented by $e$.  A homotopy pushout instead remembers two paths implementing the two compatibilities; their difference produces the circle.  Thus even existence of the ordinary colimit does not give the desired answer.

The parenthetical uncertainty on manuscript p.~123 about infinite discrete indexing categories should not be carried into the modern statement: the homotopy category of spaces has all small coproducts, represented by disjoint unions of CW models.  What fails is a general ordinary-colimit calculus that recovers coherent homotopy colimits.  The four-object calculation isolates that failure without making an unnecessary nonexistence assertion.
\end{dcomment}


\printedpage{84}
% head: A. Grothendieck : Dérivateurs, Ch. XII, Édition des Universités de Montpellier II et Paris VII
Soient
\[
\begin{aligned}
1\backslash\cX &= \cX\times_{\Phi} 1\backslash\Phi &&= \cX\times_{\Phi} F_1\\
2\backslash\cX &= \cX\times_{\Phi} 2\backslash\Phi &&= \cX\times_{\Phi} F_2
\end{aligned}
\]
les images inverses de $F_1$, $F_2$ dans $\cX$, on a donc
\[
1\backslash\cX\cap 2\backslash\cX = X_0, \qquad 1\backslash\cX\cup 2\backslash\cX = \cX,
\]
donc le diagramme
\[
\begin{tikzcd}[column sep=small, row sep=small, ampersand replacement=\&]
 \& X_0 \arrow[dl, hook', "\alpha'_1"'] \arrow[dr, hook, "\alpha'_2"] \& \\
1\backslash\cX \arrow[dr, hook, "i'_1"'] \& \& 2\backslash\cX \arrow[dl, hook', "i'_2"] \\
 \& \cX, \&
\end{tikzcd}
\tag*{$(**)$}
\]
déduit du diagramme similaire
\[
\begin{tikzcd}[column sep=small, row sep=small, ampersand replacement=\&]
 \& \{0\} \arrow[dl, hook'] \arrow[dr, hook] \& \\
F_1 \arrow[dr, hook] \& \& F_2 \arrow[dl, hook'] \\
 \& \Phi, \&
\end{tikzcd}
\]
est (comme ce dernier) un carré d'immersions fermées à la fois cartésien et cocartésien. D'autre part, on a un homomorphisme
\msp{page 125}
canonique du diagramme $(*)$ dans $(**)$, qui est l'identité sur les sommets $X_0$ et $\cX$, et qui sur $X_1$, $X_2$ est donné par les inclusions canoniques
\[
\underbrace{X_1}_{\text{\scriptsize $\cX_1$ fibre de $\cX$ en 1}} \longrightarrow 1\backslash\cX, \qquad
\underbrace{X_2}_{\text{\scriptsize $\cX_2$ fibre de $\cX$ en 2}} \longrightarrow 2\backslash\cX.
\]
Ces inclusions, on le sait, sont $W$-asphériques (elles ont même un adjoint à droite), a fortiori elles sont dans $W$. Donc le diagramme non commutatif $(*)$ est isomorphe, en tant que diagramme dans $\Hot_W$, au diagramme commutatif cartésien et cocartésien $(**)$ formé d'immersions fermées. Comme $\cX$ est bel et bien somme amalgamée de ses deux fermées % [sic] fermées -> fermés, cf. « les deux fermés » p. 83
$1\backslash\cX$, $2\backslash\cX$ sous leur intersection $X_0$, cela étaye l'intuition qui voit $\cX$ comme une `somme amalgamée homotopique'. Donnons tout de suite l'application suivante :

\dgrblock{dgrb461}{Closed pieces give a genuine gluing diagram}
\begin{dcomment}{exposition}{p125:cela étaye l'intuition qui voit}{Closed pieces give a genuine gluing diagram}
Grothendieck now replaces the informal square by a strict cover, so that an excision axiom can actually be applied.  The construction is a categorical version of replacing the two legs of a span by mapping cylinders.

Let $C_k$ be the full subcategory of $\mathcal X$ on $X_k\amalg X_0$.  This is $k\backslash\mathcal X$ in the manuscript.  It is a cosieve: an arrow starting in its side fibre can only remain there or enter $X_0$, and arrows starting in $X_0$ remain in $X_0$.  We have
\[
 C_1\cap C_2=X_0,\qquad C_1\cup C_2=\mathcal X,
 \qquad \mathcal X=C_1\amalg_{X_0}C_2.
 \tag{A5}
\]
The last equality holds in $\Cat$, not just on objects.  A string of composable arrows cannot pass from one side fibre to the other; it can enter the common fibre but cannot leave it.  Giving a functor out of $\mathcal X$ is therefore exactly giving functors out of the two pieces that agree on $X_0$.  The intersection statement also makes the square cartesian.

The inclusion $j_k:X_k\to C_k$ has right adjoint $r_k$, defined by $r_k(a)=a$ on $X_k$ and $r_k(x)=\alpha_kx$ on $X_0$.  On the mixed morphisms it uses (A4).  Indeed, for every $a\in X_k$ and $z\in C_k$,
\[
 \operatorname{Hom}_{C_k}(j_ka,z)
       =\operatorname{Hom}_{X_k}(a,r_kz).
\]
The unit is the identity and the counit is the canonical transition arrow on the middle fibre.  Thus $Bj_k$ and $Br_k$ are mutually inverse homotopy equivalences, and both functors belong to every fundamental localizer.

For spaces, (A5) proves the promised homotopy-pushout calculation.  Every simplex of $N\mathcal X$ lies in $NC_1$ or $NC_2$, and their intersection is $NX_0$.  Hence it is the pushout of these nerves.  The two inclusions of $NX_0$ are simplicial monomorphisms, whose realizations are CW-subcomplex inclusions, so this ordinary union computes the homotopy pushout of $BC_1\leftarrow BX_0\to BC_2$.  The retractions $r_k$ give a strictly commuting map of this span to $BX_1\leftarrow BX_0\to BX_2$, because $r_k|_{X_0}=\alpha_k$; all three components are homotopy equivalences.  Homotopy invariance of the double mapping cylinder proves the result.

The phrase ``homomorphism of diagrams'' on p.~125 must retain this distinction: the inclusions from the original side fibres form a comparison with specified natural transformations, and become a comparison of commutative squares in $\Hot_W$.  The retractions just written give a strict map of spans in the other direction.  No strict commutativity of the original square is being asserted or needed.  For a general $W$, the same replacement reduces Proposition~13 to precisely the closed-cover form of the stated axiom W(7 bis).
\end{dcomment}


\textbf{Proposition 13.} \textit{Soit un diagramme} $\begin{tikzcd}[column sep=small, row sep=tiny, ampersand replacement=\&, baseline=(\tikzcdmatrixname-2-1.base)]
 \& X_1 \\
X_0 \arrow[ur, "\alpha_1"] \arrow[dr, "\alpha_2"'] \& \\
 \& X_2
\end{tikzcd}$ \textit{dans} Cat\textit{, d'où le diagramme} $(*)$\textit{, p. 123,}

\printedpage{85}
% head: A. Grothendieck : Dérivateurs, Ch. XII, Édition des Universités de Montpellier II et Paris VII
\[
\begin{tikzcd}[column sep=small, row sep=small, ampersand replacement=\&]
 \& X_0 \arrow[dl, "\alpha_1"'] \arrow[dr, "\alpha_2"] \& \\
X_1 \arrow[dr, "i_1"'] \& \& X_2 \arrow[dl, "i_2"] \\
 \& \cX, \&
\end{tikzcd}
\]
\msp{page 126}
\textit{commutatif dans} $\Hot_W$ \textit{(mais non dans} Cat\textit{, sauf si $X_0 = \emptyset$). Supposons que $W$ satisfait} W(7 bis) (axiome du carré cocartésien, cf. p. 75). \textit{Alors on a}
\[
\left\{
\begin{aligned}
\alpha_1\in W &\iff i_2\in W\\
\alpha_2\in W &\iff i_1\in W. \text{ (\footnotemark[85])}
\end{aligned}
\right.
\]
\footnotetext[85]{\textbf{NB} Pour l'implication $\Longrightarrow$ on n'a besoin que de la forme directe de W(7 bis), i.e. (avec les notations de XVI) de Loc(6$'$a).}
\textsc{Démonstration}. Utilisant l'homomorphisme précédent de $(*)$ dans $(**)$, on voit que chacune des flèches $\alpha_1$, $\alpha_2$, $i_1$, $i_2$ est dans $W$, si et seulement si la flèche correspondante $\alpha'_1$, $\alpha'_2$, $i'_1$ ou $i'_2$ l'est. Donc les équivalences ci-dessus équivalent à
\[
\begin{aligned}
\alpha'_1\in W &\iff i'_2\in W\\
\alpha'_2\in W &\iff i'_1\in W,
\end{aligned}
\]
lesquelles sont vraies en vertu de W(7 bis) (formulation duale) appliqué à $\cX$ et aux deux fermés $1\backslash\cX$, $2\backslash\cX$ de $\cX$.

\dgrblock{dgrb462}{Why converse excision is an additional axiom}
\begin{dcomment}{development}{p126:lesquelles sont vraies en vertu de W(7 bis)}{Why converse excision is an additional axiom}
Proposition~13 is correctly proved under its stated hypothesis.  The strength of that hypothesis nevertheless matters: the displayed equivalences contain converse excision, whereas the usual invariance of a homotopy pushout under weak equivalences gives only the forward implication.  A suspension makes the difference precise.

\textbf{Proposition (the acyclic test for converse excision).}\ours
Let $W$ contain the weak homotopy equivalences of small categories and satisfy two-out-of-three.  Suppose that, whenever $C=C_1\cup C_2$ is a union of two full cosieves with intersection $A$, the implication
\[
 C_2\longrightarrow C\text{ belongs to }W
 \quad\Longrightarrow\quad A\longrightarrow C_1\text{ belongs to }W
 \tag{A6}
\]
holds.  Then every category $A$ whose classifying space is connected and integrally acyclic is $W$-aspherical.

\emph{Proof.}
Adjoin to $A$ two distinct initial vertices $l,r$, with one arrow from each vertex to every object of $A$, and no arrows between $l$ and $r$.  Call the resulting category $C$ and its two one-vertex enlargements $C_1,C_2$.  The full subcategories $C_1,C_2$ are cosieves, and their intersection is $A$.  Their classifying spaces are cones on $BA$; the classifying space of $C$ is their union along $BA$, namely the unreduced suspension $\Sigma BA$.

Since $BA$ is connected, van Kampen applied to the two slightly enlarged cones makes $\Sigma BA$ simply connected.  The reduced Mayer--Vietoris sequence gives $\widetilde H_n(\Sigma BA;\mathbb Z)\cong\widetilde H_{n-1}(BA;\mathbb Z)=0$ for every $n$.  A simply connected acyclic CW complex is contractible: if it had a first nonzero homotopy group, the Hurewicz theorem would identify it with a nonzero homology group; with all homotopy groups zero, Whitehead's theorem applies.  Equivalently, use \citet[Corollary~4.33]{Hatcher2002}.  Thus $C_2\to C$ is a weak homotopy equivalence, and lies in $W$.  By (A6), $A\to C_1$ lies in $W$.  Since $C_1\to e$ also lies in $W$, two-out-of-three gives $A\to e\in W$.\hfill$\square$

For completeness, an explicit finite CW complex shows why this conclusion is too strong for $W_\infty$.  Let $K$ have one vertex, three oriented edges $x,y,z$, and three two-cells with relators
\[
 x^2(xyz)^{-1},\qquad y^3(xyz)^{-1},\qquad z^5(xyz)^{-1}.
\]
Its cellular boundary $C_2(K;\mathbb Z)\to C_1(K;\mathbb Z)$ has, in these bases, the matrix
\[
 M=\begin{pmatrix}1&-1&-1\\-1&2&-1\\-1&-1&4\end{pmatrix},
 \qquad \det M=-1.
\]
It is an isomorphism.  The other reduced cellular groups are zero, so $K$ is acyclic.  Its fundamental group is nontrivial: with permutation multiplication read from right to left, send
\[
 x\longmapsto(1\,2)(3\,4),\quad
 y\longmapsto(1\,3\,5),\quad
 z\longmapsto(1\,2\,5\,3\,4).
\]
These permutations have orders $2,3,5$ respectively and product $xyz=1$, so they satisfy the three relators and define a homomorphism with nontrivial image in $A_5$.  A finite simplicial complex homotopy equivalent to $K$ exists by \citet[Theorem~2C.5]{Hatcher2002}; let $A$ be its poset of nonempty faces.  The nerve of $A$ is its barycentric subdivision.  Thus $BA$ is connected and acyclic but has nontrivial fundamental group.

Apply the construction in the proof to this finite poset.  Then $C_1$, $C_2$, and $C$ are weakly contractible, while $A$ is not.  Hence $C_2\to C$ belongs to $W_\infty$ but $A\to C_1$ does not.  This disproves full W(7 bis) for $W_\infty$.  Mapping this closed-cover square to the trivial square of one-point categories also disproves the implication $b\mathbin{\&}c\Rightarrow a$ in full W(7).  The same example disproves the corresponding converses in Proposition~13 and the corollaries if their extra hypotheses are deleted.  It does not contradict the propositions as actually stated.

The topology behind the example is classical; compare the acyclic presentation complex in \citet[Example~2.38]{Hatcher2002}.  We checked that example, the manuscript's definitions of W(7 bis) and W(7), and the standard nerve and suspension arguments.  We did not find this exact formulation as a test of these manuscript axioms in the references inspected.  The proposition records the deduction needed here, without a claim of priority.  In particular the assertion earlier in the chapter that $W_\infty$ satisfies full W(7) cannot be justified by a Mayer--Vietoris argument alone: local systems on the contractible union restrict to constant systems on $A$, and do not test its nontrivial fundamental group.

There is a precise positive form of the five-lemma intuition.  For a closed cover, every nondegenerate nerve simplex belongs to one of the two pieces, so normalized integral chains form a natural short exact sequence
\[
 0\longrightarrow C_*(NA)
  \longrightarrow C_*(NC_1)\oplus C_*(NC_2)
  \longrightarrow C_*(NC)\longrightarrow0.
 \tag{A7}
\]
The first map has signs $(+,-)$ and the second adds the two inclusions.  Exactness follows separately on each simplex, including degree zero.  Given a map of two covers, if any two of the maps on the three terms of (A7) are quasi-isomorphisms, the third is a quasi-isomorphism by the associated long exact homology sequences.  For the middle term this is equivalent to being a homology isomorphism on each piece.  Thus all three implications of the cube axiom hold for integral homology equivalences of closed covers.  The passage from this homological statement to weak homotopy equivalences is exactly the passage that the counterexample prevents.
\end{dcomment}


\Needspace{13\baselineskip}
\textbf{Corollaire :} \textit{Considérons} deux \textit{diagrammes de type} $\Phi^{\circ}$\textit{, et un homomorphisme} $f = (f_0, f_1, f_2)$ \textit{de l'un dans l'autre (cf. diagramme en traits pleins ci-contre} \texttt{[ci-dessous]}\textit{):}
\[
\begin{tikzcd}[column sep=small, row sep=small, ampersand replacement=\&]
 \& X_1 \arrow[dd, "f_1"'] \arrow[drr, dotted, "i_1"] \& \& \\
X_0 \arrow[ur, "\alpha_1"] \arrow[drr, "\alpha_2"] \arrow[dd, "f_0"'] \& \& \& \cX \arrow[dd, dotted, "f_{\cX}"] \\
 \& X'_1 \arrow[drr, dotted, "i'_1"] \& X_2 \arrow[ur, dotted, "i_2"] \arrow[dd, "f_2"'] \& \\
X'_0 \arrow[ur, "\alpha'_1"] \arrow[drr, "\alpha'_2"] \& \& \& \cX', \\
 \& \& X'_2 \arrow[ur, dotted, "i'_2"] \&
\end{tikzcd}
\]
\textit{d'où un homomorphisme}
\msp{page 127}
\[
f_{\cX} = \int f : \quad \cX = \int
\begin{tikzcd}[column sep=small, row sep=tiny, ampersand replacement=\&, baseline=(\tikzcdmatrixname-2-1.base)]
 \& X_1 \\
X_0 \arrow[ur] \arrow[dr] \& \\
 \& X_2
\end{tikzcd}
\qquad\longrightarrow\qquad
\cX' = \int
\begin{tikzcd}[column sep=small, row sep=tiny, ampersand replacement=\&, baseline=(\tikzcdmatrixname-2-1.base)]
 \& X'_1 \\
X'_0 \arrow[ur] \arrow[dr] \& \\
 \& X'_2,
\end{tikzcd}
\]

\printedpage{86}
% head: A. Grothendieck : Dérivateurs, Ch. XII, Édition des Universités de Montpellier II et Paris VII
\textit{s'insérant dans un diagramme cubique (cf. diagramme précédent complété par des pointillés). Considérons les conditions suivantes.}

\begin{itemize}
\item[\textit{a)}] $f_0\in W$\textit{.}
\item[\textit{b)}] $f_1$ \textit{et} $f_2$ \textit{sont} \texttt{[dans]} $W$\textit{.}
\item[\textit{c)}] $f_{\cX} = \int(f_0, f_1, f_2)\in W$\textit{.}
\end{itemize}

\textit{Supposons l'axiome} W(7) (`lemme des 5', cf. page 76) \textit{satisfait. Alors deux parmi les conditions a), b), c) impliquent la troisième}
\[
\left\{
\begin{aligned}
a \mathbin{\&} b &\quad\Longrightarrow\quad c\\
b \mathbin{\&} c &\quad\Longrightarrow\quad a\\
c \mathbin{\&} a &\quad\Longrightarrow\quad b.
\end{aligned}
\right.
\]

En effet, on a un homomorphisme du diagramme cubique ci-dessus dans le diagramme similaire, où $X_1$, $X_2$, $X'_1$, $X'_2$ sont remplacés respectivement pas les fermés % [sic] pas -> par, misprint
$1\backslash\cX$, $2\backslash\cX$ de $\cX$, et $1\backslash\cX'$, $2\backslash\cX'$ de $\cX'$, les deux carrés supérieur et inférieur étant des inclusions d'ouverts \texttt{[plutôt de fermés]} de $\cX$ resp. de $\cX'$. Ainsi, les conditions a), b), c) sont équivalentes respectivement aux conditions similaires pour le nouveau diagramme cubique
\msp{page 128}
(a) et c) n'ayant d'ailleurs pas changé, et b) donnant
\begin{itemize}
\item[b$'$)] $f'_1 : \cX/1 \longrightarrow \cX'/1$\ \ et\ \ $f'_2 : \cX/2 \longrightarrow \cX'/2$\ \ sont dans $W$.
\end{itemize}
\texttt{[plutôt} $f'_1 : 1\backslash\cX \longrightarrow 1\backslash\cX'$ \texttt{et} $f'_2 : 2\backslash\cX \longrightarrow 2\backslash\cX'$\texttt{]}). Alors la conclusion du corollaire est équivalente à W(7) (formulation duale) dans le cas du nouveau diagramme cubique, q.e.d.

\subsection*{3.2\quad Comparaison de \ $\displaystyle\int
\begin{tikzcd}[column sep=small, row sep=tiny, ampersand replacement=\&, baseline=(\tikzcdmatrixname-2-1.base)]
 \& X_1 \\
X_0 \arrow[ur] \arrow[dr] \& \\
 \& X_2
\end{tikzcd}$ \ avec une somme amalgamée ordinaire.\\
Carrés $W$-homotopiquement cocartésiens.}

La catégorie intégrée $\cX = \displaystyle\int
\begin{tikzcd}[column sep=small, row sep=tiny, ampersand replacement=\&, baseline=(\tikzcdmatrixname-2-1.base)]
 \& X_1 \\
X_0 \arrow[ur] \arrow[dr] \& \\
 \& X_2
\end{tikzcd}$ est solution du problème universel suivant, variante 2-catégorique de celui qui définit la limite inductive des diagrammes. Il s'agit, pour $Y$ objet variable dans Cat, de considérer les systèmes de trois foncteurs $f_0$, $f_1$, $f_2$ et des deux homomorphismes fonctoriels $\lambda_1$, $\lambda_2$ :
\printedpage{87}
% head: A. Grothendieck : Dérivateurs, Ch. XII, Édition des Universités de Montpellier II et Paris VII
\[
\begin{tikzcd}[column sep=large, row sep=large, ampersand replacement=\&]
 \& X_0 \arrow[dl, "\alpha_1"'] \arrow[dr, "\alpha_2"] \arrow[dd, "f_0", ""{name=M, pos=0.62}] \& \\
X_1 \arrow[dr, "f_1"', ""{name=A, pos=0.45}] \& \& X_2 \arrow[dl, "f_2", ""{name=B, pos=0.45}] \\
 \& Y, \&
\arrow[from=A, to=M, bend left, shorten=3pt, "\lambda_1"]
\arrow[from=B, to=M, bend right, shorten=3pt, "\lambda_2"']
\end{tikzcd}
\tag*{$(*)$}
\]
\[
\begin{aligned}
\lambda_1 :\quad & f_1\alpha_1 \quad\longrightarrow\quad f_0\\
\lambda_2 :\quad & f_2\alpha_2 \quad\longrightarrow\quad f_0.
\end{aligned}
\]

\msp{page 129}
Se donner une flèche de Cat
\[
\cX \longrightarrow Y
\]
`revient au même' que de se donner un système $(f_0, f_1, f_2, \lambda_1, \lambda_2)$ comme dessus, et ceci `fonctoriellement en $Y$'. C'est là le cas particulier de la détermination générale des foncteurs de la catégorie fibrée `intégrée' d'un contrafoncteur
\[
I^{\circ} \longrightarrow \Cat, \qquad
\begin{aligned}
i &\quad\longmapsto\quad X_i\\
u &\quad\longmapsto\quad u^*
\end{aligned}
\]
vers une catégorie $Y$, en termes d'un système de foncteurs
\[
X_i \xrightarrow{f_i} Y,
\]
et d'homomorphismes fonctoriels
\[
f_i\circ u^* \xrightarrow{\lambda_u} f_j
\]
pour $u : i \longrightarrow j$, d'où $u^* : X_j \longrightarrow X_i$, avec des conditions de transitivité pour des flèches composées, dont on n'a pas à tenir compte dans le cas actuel, vu qu'il n'y a pas de flèches composées à regarder dans $I = \Phi$ (sauf des composées avec une flèche identique, qu'on peut se dispenser de regarder ici \dots).

\dgrblock{dgrb463}{The universal property keeps its two-cells}
\begin{dcomment}{exposition}{p129:vu qu'il n'y a pas de flèches composées à regarder}{The universal property keeps its two-cells}
The universal property on these pages explains why integration is a useful replacement for an ordinary colimit.  It represents functors equipped with comparison transformations.  Those transformations are part of the input and are allowed to be noninvertible.

We can prove the assertion directly from (A4).  Suppose given $f_i:X_i\to Y$ and $\lambda_k:f_k\alpha_k\Rightarrow f_0$ for $k=1,2$.  Define $F:\mathcal X\to Y$ to be $f_i$ on each fibre.  A mixed morphism from $a\in X_k$ to $x\in X_0$, represented by $h:a\to\alpha_kx$, must be sent to
\[
 f_k(a)\xrightarrow{f_k(h)}f_k(\alpha_kx)
       \xrightarrow{\lambda_{k,x}}f_0(x).
 \tag{A8}
\]
This rule respects precomposition in $X_k$.  It respects postcomposition by $v:x\to x'$ because naturality says
\[
 f_0(v)\lambda_{k,x}=\lambda_{k,x'}f_k(\alpha_kv).
\]
There are no composites of two mixed arrows, so these checks prove functoriality.  Conversely, restrict $F$ to the fibres and apply it to the mixed arrows represented by identities.  This recovers the $f_i$ and the $\lambda_k$, and every mixed arrow factors as an internal arrow followed by one of these identities.  The constructions are inverse.

The assertion also holds at the level of categories of functors.  A natural transformation $F\Rightarrow G$ corresponds to transformations $\theta_i:f_i\Rightarrow g_i$ such that, for $k=1,2$,
\[
 \theta_0\lambda_k=\mu_k(\theta_k\alpha_k),
\]
where $\mu_k$ are the comparison transformations for $G$.  Necessity is naturality on the canonical mixed arrows; sufficiency follows from their factorization just used.  Thus the universal property is a natural isomorphism of categories with the category of these lax cocones.  The word ``lax'' here is fixed by the displayed direction $f_k\alpha_k\to f_0$, so it carries no hidden variance convention.

For a general strict functor $I^{\mathrm{op}}\to\Cat$, the additional condition is explicit.  If $u:i\to j$ and $v:j\to k$, with maps $u^*:X_j\to X_i$ and $v^*:X_k\to X_j$, then
\[
 \lambda_{vu}=\lambda_v\circ(\lambda_u v^*),
 \qquad \lambda_{\mathrm{id}_i}=\mathrm{id}_{f_i}.
\]
For the two-legged shape these conditions have no nonidentity instances.  This is why the construction can be given here without a separate coherence argument.

When the original square is strictly commutative, set $f_0=f_k\alpha_k$ and take both $\lambda_k$ to be identities.  Formula (A8) then defines $\varphi_Q$ on every morphism and shows immediately that $\varphi_Qi_i=f_i$.  Every chosen cartesian transition arrow is sent to an identity.  This supplies the requested construction in footnote~87; the comma-category calculations announced there are carried out in the later lemmas on pp.~152--154.  Definition~3 asks whether this particular comparison belongs to $W$.  A bare commutative square in $\Hot_W$ would not determine the comparison, because it has forgotten the transformations from which it was made.
\end{dcomment}


\msp{page 130}
Si en particulier on a un diagramme \textit{commutatif} dans Cat

% ---- from Der-ChXIIb-012_p046-050.tex
\printedpage{88}
% head: A. Grothendieck : Dérivateurs, Ch. XII, Édition des Universités de Montpellier II et Paris VII
\[
Q \quad = \qquad
\begin{tikzcd}
 & X_0 \arrow[dl, "\alpha_1"'] \arrow[dr, "\alpha_2"] & \\
X_1 \arrow[dr, "f_1"'] & & X_2 \arrow[dl, "f_2"] \\
 & Y, &
\end{tikzcd}
\]
on peut le compléter par
\[
f_0 \overset{\text{déf}}{=} f_1\alpha_1 = f_2\alpha_2
\]
en un diagramme
\[
\begin{tikzcd}[row sep=large, column sep=large]
 & X_0 \arrow[dl, "\alpha_1"'] \arrow[dr, "\alpha_2"] \arrow[dd, "f_0", ""'{name=F0, pos=0.65}] & \\
X_1 \arrow[dr, "f_1"', ""{name=F1, pos=0.4}] & & X_2 \arrow[dl, "f_2", ""'{name=F2, pos=0.4}] \\
 & Y, &
\arrow[from=F1, to=F0, bend left, "\mathrm{id}", shorten <=2pt, shorten >=2pt]
\arrow[from=F2, to=F0, bend right, "\mathrm{id}"', shorten <=2pt, shorten >=2pt]
\end{tikzcd}
\]
ce qui montre qu'on a une donnée du type $(*)$ ci-dessus, définissant donc une flèche
\[
\varphi_Q : \quad \sX = \int
\begin{tikzcd}[row sep=small, column sep=small]
 & X_1 \\
X_0 \arrow[ur] \arrow[dr] & \\
 & X_2
\end{tikzcd}
\qquad \longrightarrow \qquad Y
\]
qui donne lieu à un diagramme cubique

\msp{page 131}
\[
\begin{tikzcd}[row sep=small, column sep=small]
 & X_1 \arrow[dd, "\mathrm{id}"'] \arrow[rrd, "i_1", ""'{name=I1}] & & \\
X_0 \arrow[ru, "\alpha_1"] \arrow[rrr, dotted, "i_0", ""'{name=I0, pos=0.8}] \arrow[rrd, "\alpha_2"] \arrow[dd, "\mathrm{id}"'] & & & \sX \arrow[dd, "\varphi"] \\
 & X_1 \arrow[rrd, "f_1"] & X_2 \arrow[ru, "i_2"', ""{name=I2}] \arrow[dd, "\mathrm{id}"'] & \\
X_0 \arrow[ru, "\alpha_1"] \arrow[rrr, dotted, "f_0"] \arrow[rrd, "\alpha_2"] & & & Y, \\
 & & X_2 \arrow[ru, "f_2"'] &
\arrow[from=I1, to=I0, "\lambda_1"', shorten <=1pt, shorten >=1pt]
\arrow[from=I2, to=I0, bend left, "\lambda_2", shorten <=1pt, shorten >=1pt] % Resolved from PDF leaf 11: a solid curve, unlike the genuinely dotted i_0 line.
\end{tikzcd}
\]
% foot: Version 7 décembre 2002

\printedpage{89}
% head: A. Grothendieck : Dérivateurs, Ch. XII, Édition des Universités de Montpellier II et Paris VII
commutatif sauf le carré supérieur, i.e. on a
\[
\varphi i_1 = f_1, \qquad \varphi i_2 = f_2, \qquad \text{et de plus} \quad \varphi i_0 = f_0,
\]
$\varphi$ transformant de plus les homomorphismes fonctoriels
\[
\lambda_1 : i_1\alpha_1 \longrightarrow i_0, \qquad \lambda_2 : i_2\alpha_2 \longrightarrow i_0
\]
en les homomorphismes fonctoriels identiques pertinents.

\setcounter{footnote}{85}
\textbf{Définition 3.} On dit que le diagramme commutatif
\[
\begin{tikzcd}[row sep=small, column sep=small]
 & X_0 \arrow[dl, "\alpha_1"'] \arrow[dr, "\alpha_2"] & \\
X_1 \arrow[dr, "f_1"] & & X_2 \arrow[dl, "f_2"'] \\
 & Y &
\end{tikzcd}
\tag*{$(*)$}
\]
est \textit{$W$-homotopiquement cocartésien} (\footnote{ou simplement $W$-cocartésien, ou qu'il définit $Y$ comme somme amalgamée $W$-homotopique de etc.}) si l'homomorphisme canonique
\[
\sX = \int
\begin{tikzcd}[row sep=small, column sep=small]
 & X_1 \\
X_0 \arrow[ur] \arrow[dr] & \\
 & X_2
\end{tikzcd}
\qquad \xrightarrow{\ \varphi\ } \qquad Y
\]
qu'il définit est dans $W$. (\footnote{Introduire $\varphi$ dans une proposition en forme, et donner ses propriétés essentielles : flèches strictement cartésiennes sont transformées en flèches \textit{identiques} $+$ lemmes 1 et 2 des p. 152, 154.})

\msp{page 132}
On a alors le :

\textbf{Corollaire 1.} \textit{Supposons} W(7 bis)\textit{. Si le carré} $(*)$ \textit{est $W$-homotopiquement cocartésien, alors}
\[
\left\{
\begin{aligned}
\alpha_1 \in W &\quad\iff\quad f_2 \in W \\
\alpha_2 \in W &\quad\iff\quad f_1 \in W.
\end{aligned}
\right.
\]

Cela résulte aussitôt du diagramme cubique de la page précédente, où toutes les flèches verticales sont dans $W$, et de l'énoncé similaire concernant la face supérieure du cube, donné dans la proposition 13, p. 125.

\Needspace{13\baselineskip}
\textbf{Corollaire 2.} \textit{Supposons} W(7)\textit{. Considérons un diagramme cubique}
% foot: Version 7 décembre 2002

\printedpage{90}
% head: A. Grothendieck : Dérivateurs, Ch. XII, Édition des Universités de Montpellier II et Paris VII
\[
\begin{tikzcd}[row sep=small, column sep=small]
 & X_1 \arrow[dd, "p_1"'] \arrow[rrd, "f_1"] & & \\
X_0 \arrow[ru, "\alpha_1"] \arrow[rrd, "\alpha_2"] \arrow[dd, "p_0"'] & & & Y \arrow[dd, "p"] \\
 & X'_1 \arrow[rrd, "f'_1"] & X_2 \arrow[ru, "f_2"] \arrow[dd, "p_2"'] & \\
X'_0 \arrow[ru, "\alpha'_1"] \arrow[rrd, "\alpha'_2"] & & & Y', \\
 & & X'_2 \arrow[ru, "f'_2"] &
\end{tikzcd}
\]
\textit{où les deux faces inférieure et supérieure du cube sont des carrés $W$-homotopiquement cocartésiens. Considérons les conditions}

\msp{page 133}
\begin{enumerate}[label={}]
\item[\textit{a)}] $p_0 \in W$\textit{.}
\item[\textit{b)}] $p_1$ \textit{et} $p_2 \in W$\textit{.}
\item[\textit{c)}] $p \in W$\textit{.}
\end{enumerate}

\textit{Alors deux de ces conditions impliquent la troisième.}

Résulte aussitôt de la définition 3, et du résultat similaire (cor. page 126) pour le diagramme cubique similaire, avec $Y$, $Y'$ remplacées par $\sX = \int
\begin{tikzcd}[row sep=small, column sep=small]
 & X_1 \\
X_0 \arrow[ur] \arrow[dr] & \\
 & X_2
\end{tikzcd}$
et $\sX' = \int
\begin{tikzcd}[row sep=small, column sep=small]
 & X'_1 \\
X'_0 \arrow[ur] \arrow[dr] & \\
 & X'_2
\end{tikzcd}$\;.

\section*{3.3\quad Intégration propre versus intégration lisse. Énoncé d'équivalence entre les deux points de vue. Commentaires sur le cas des limites $W$-homotopiques quelconques.}

La définition 3, et les corollaires précédents, n'ont d'intérêt que dans la mesure où on dispose de critères maniables, permettant d'affirmer qu'un carré commutatif est bel et bien homotopiquement cocartésien. C'est ce qu'on va essayer de faire tantôt. Il y a aussi une perplexité, concernant le rôle privilégié joué, dans la théorie développée ici, par les immersions fermées, à l'exclusion des immersions ouvertes. Nous avons en somme ramené, du
% foot: Version 7 décembre 2002
\printedpage{91}
% head: A. Grothendieck : Dérivateurs, Ch. XII, Édition des Universités de Montpellier II et Paris VII
\msp{page 134}
point de vue homotopique, les diagrammes quelconques $\begin{tikzcd}[row sep=small, column sep=small]
 & X_1 \\
X_0 \arrow[ur] \arrow[dr] & \\
 & X_2
\end{tikzcd}$ dans Cat, au cas d'un diagramme d'immersions fermées, où la somme amalgamée prend des allures familières. Mais pourquoi pas aussi bien le ramener au cas d'immersions ouvertes? On considérerait alors le diagramme de départ comme définissant une catégorie \textit{cofibrée} sur
\[
\Phi^{\circ} = \Psi = \
\begin{tikzcd}[row sep=small, column sep=small]
 & 1 \\
0 \arrow[ur] \arrow[dr] & \\
 & 2,
\end{tikzcd}
\]
soit $\tX$. On a dans $\Psi$ les deux ouverts
\[
\underbrace{\Psi/1}_{=\{0,1\}} = U_1, \qquad \underbrace{\Psi/2}_{=\{0,2\}} = U_2, \qquad U_1 \cap U_2 = \{0\},
\]
d'où dans $\tX$ deux ouverts images inverses
\[
\tX/1, \qquad \tX/2, \qquad \text{avec } \tX/1 \cap \tX/2 = X_0,
\]
d'où un diagramme commutatif d'inclusions ouvertes
\[
\begin{tikzcd}
 & X_0 \arrow[dl, "\alpha'_1"'] \arrow[dr, "\alpha'_2"] & \\
\tX/1 \arrow[dr, "i'_1"'] & & \tX/2 \arrow[dl, "i'_2"] \\
 & \tX, &
\end{tikzcd}
\tag*{$(*)$}
\]

\dgrblock{dgrb464}{The other orientation of integration}
\begin{dcomment}{exposition}{p133:La définition 3, et les corollaires précédents}{The other orientation of integration}
Grothendieck wants the definition of a homotopy pushout to belong to the span itself, without retaining the auxiliary choice of open or closed inclusions used to construct it. The question is quite concrete. The smooth integral introduced above has transition arrows pointing from the two outer fibres towards the central fibre. The proper integral has transition arrows in the other direction. Neither construction is obtained by merely changing the names of the objects.

Write $A=X_0$, $B=X_1$, $C=X_2$, and let $a:A\to B$ and $b:A\to C$ be the given functors. The proper integral $P$ has the objects of $A$, $B$, and $C$, retained as three disjoint sets. In addition to the arrows within those categories, its mixed hom-sets are
\[
 P(x,y)=B(a(x),y),\qquad
 P(x,z)=C(b(x),z)
 \quad(x\in A,\ y\in B,\ z\in C).
\]
There are no mixed arrows in the opposite directions, or between $B$ and $C$. Composition is composition in $B$ and $C$, with $a$ and $b$ applied when necessary. Thus a functor $P\to Y$ consists of functors $f_A,f_B,f_C$ and transformations
\[
 f_A\longrightarrow f_Ba,\qquad f_A\longrightarrow f_Cb.
\]
This is the reverse of the transformation data described in ``A span is remembered by its transition arrows''. For a strictly commuting cocone these transformations are identities, so both constructions have a canonical comparison with $Y$.

The full subcategory $P_1$ on $A$ and $B$ is a sieve in $P$, as is the corresponding $P_2$ on $A$ and $C$. Their intersection is $A$, and their union is $P$. The inclusion $j:B\to P_1$ has a left adjoint $r$: it is the identity on $B$ and sends $x\in A$ to $a(x)$. The identity
\[
 P_1(x,jy)=B(r(x),y)
\]
proves the adjunction, including its mixed case. Consequently $j$ is a weak equivalence for every fundamental localizer. The same argument applies to $C\to P_2$. This explains the parallel with the closed-cover replacement of the smooth integral: the two replacements have the same three vertex homotopy types, but their transition arrows have opposite orientations.

The equivalences asserted on manuscript page 135 inherit the hypotheses of their earlier counterparts. In particular, the reverse implication from an outer inclusion being in $W$ to the corresponding original leg being in $W$ uses the strong form of W(7 bis). The cube statement uses W(7). These are the additional excision assumptions analysed in ``Why converse excision is an additional axiom''. The comparison between the two integrals, which is the object of Proposition 14, does not require either converse-excision principle.
\end{dcomment}


\msp{page 135}
lequel reçoit le diagramme non commutatif
\[
\begin{tikzcd}[row sep=large, column sep=large]
 & X_0 \arrow[dl, "\alpha_1"'] \arrow[dr, "\alpha_2"] \arrow[dd, dotted, "i_0", ""'{name=J0, pos=0.6}] & \\
X_1 \arrow[dr, "i_1"', ""{name=J1, pos=0.4}] & & X_2 \arrow[dl, "i_2", ""'{name=J2, pos=0.4}] \\
 & \tX, &
\arrow[from=J0, to=J1, dotted, bend right, "\tilde{\lambda}_1"', shorten <=2pt, shorten >=2pt]
\arrow[from=J0, to=J2, dotted, bend left, "\tilde{\lambda}_2", shorten <=2pt, shorten >=2pt]
\end{tikzcd}
\]
% foot: Version 7 décembre 2002
\printedpage{92}
% head: A. Grothendieck : Dérivateurs, Ch. XII, Édition des Universités de Montpellier II et Paris VII
les flèches
\[
X_1 \longrightarrow \tX/1, \qquad X_2 \longrightarrow \tX/2
\]
étant dans $W$ (et même $W$-coasphériques - elles ont même un adjoint à gauche).

Donc on trouve, comme précédemment
\[
\left\{
\begin{aligned}
\alpha_1 \in W &\quad\iff\quad i_2 \in W \\
\alpha_2 \in W &\quad\iff\quad i_1 \in W,
\end{aligned}
\right.
\]
et de même on peut formuler un ‘lemme des cinq’ comme le corollaire page 125 \texttt{[plutôt 126]}, mais avec $\tX$, $\tX'$ au lieu de $\sX$, $\sX'$.

De plus, $\tX$ est la solution d'un problème universel dans Cat, en termes de données
\[
\begin{tikzcd}[row sep=large, column sep=large]
 & X_0 \arrow[dl, "\alpha_1"'] \arrow[dr, "\alpha_2"] \arrow[dd, "f_0", ""'{name=G0, pos=0.6}] & \\
X_1 \arrow[dr, "f_1"', ""{name=G1, pos=0.4}] & & X_2 \arrow[dl, "f_2", ""'{name=G2, pos=0.4}] \\
 & Y, &
\arrow[from=G0, to=G1, bend right, "\tilde{\lambda}_1"', shorten <=2pt, shorten >=2pt]
\arrow[from=G0, to=G2, bend left, "\tilde{\lambda}_2", shorten <=2pt, shorten >=2pt]
\end{tikzcd}
\]

\msp{page 136}
la seule différence avec celui formulé à la page 128 étant la direction opposée des flèches $\lambda_1$, $\lambda_2$. Pour tout carré commutatif
\[
\begin{tikzcd}
 & X_0 \arrow[dl, "\alpha_1"'] \arrow[dr, "\alpha_2"] & \\
X_1 \arrow[dr, "f_1"'] & & X_2 \arrow[dl, "f_2"] \\
 & Y, &
\end{tikzcd}
\tag*{$(*)$}
\]
on trouve encore une flèche canonique
\[
\tilde{\varphi} : \tX \longrightarrow Y.
\]

Il serait tout aussi raisonnable de dire que le carré $(*)$ est $W$-homotopiquement cocartésien, lorsque $\tilde{\varphi}$ est \texttt{[dans]} $W$, que d'adopter la définition 3, page 131 - on trouverait exactement les mêmes corollaires. Les deux définitions seraient elles équivalentes?? % [sic] seraient elles -> seraient-elles, hyphen of the inverted subject missing
% foot: Version 7 décembre 2002

% ---- from Der-ChXIIb-012_p051-055.tex
\printedpage{93}
% head: A. Grothendieck : Dérivateurs, Ch. XII, Édition des Universités de Montpellier II et Paris VII
% foot: Version 7 décembre 2002

Je note que $\tsX$ étant cofibrée sur $\Psi=\Phi^\circ$, $\tsX^\circ$ est fibrée sur $\Psi^\circ=\Phi$. De façon précise, on aura
\[
\tsX^\circ = \int
\begin{tikzcd}[row sep=small, column sep=small]
 & (X_1)^\circ \\
(X_0)^\circ \arrow[ur, "\alpha_1^\circ"] \arrow[dr, "\alpha_2^\circ"'] & \\
 & (X_2)^\circ,
\end{tikzcd}
\]
\msp{page 137}
et il s'agit donc de comparer les deux intégrales
\[
\sX = \int
\begin{tikzcd}[row sep=small, column sep=small]
 & X_1 \\
X_0 \arrow[ur, "\alpha_1"] \arrow[dr, "\alpha_2"'] & \\
 & X_2
\end{tikzcd}
, \qquad
\tsX^\circ = \int
\begin{tikzcd}[row sep=small, column sep=small]
 & (X_1)^\circ \\
(X_0)^\circ \arrow[ur, "\alpha_1^\circ"] \arrow[dr, "\alpha_2^\circ"'] & \\
 & (X_2)^\circ
\end{tikzcd}
,
\]
notamment comme objets de $\HOT_W(\Phi)$. Pour cela, je reprends la construction du foncteur
\[
X \longmapsto \Phi(X) : \Cat \longrightarrow \Cat
\]
de I, p. 72-75, mais je préfère ici noter ce foncteur $\Sigma(X)$ (car $\Phi$ désigne déjà une petite catégorie $\begin{tikzcd}[row sep=tiny, column sep=small]
 & 1 \arrow[dl] \\
0 & \\
 & 2 \arrow[ul]
\end{tikzcd}$). On utilisera les homomorphismes fonctoriels en $X$ :
\[
\begin{tikzcd}[column sep=small]
 & \Sigma(X) \arrow[dl, "p_X"'] \arrow[dr, "q_X"] & \\
X & & X^\circ,
\end{tikzcd}
\]
qui sont $W$-asphériques, donc $\in W$. De plus, $\Sigma(X^\circ)=\Sigma(X)$, $p_{X^\circ}=q_X$, $q_{X^\circ}=p_X$. Ainsi, si
\[
\begin{aligned}
X_* \quad &= \qquad
\begin{tikzcd}[ampersand replacement=\&, row sep=small, column sep=small]
 \& X_1 \\
X_0 \arrow[ur, "\alpha_1"] \arrow[dr, "\alpha_2"'] \& \\
 \& X_2
\end{tikzcd}
\qquad \in \quad \Hom(\Phi^\circ,\Cat)\\
X_*^\circ \quad &= \qquad
\begin{tikzcd}[ampersand replacement=\&, row sep=small, column sep=small]
 \& (X_1)^\circ \\
(X_0)^\circ \arrow[ur, "\alpha_1^\circ"] \arrow[dr, "\alpha_2^\circ"'] \& \\
 \& (X_2)^\circ
\end{tikzcd}
\qquad \in \quad \Hom(\Phi^\circ,\Cat),
\end{aligned}
\]
\msp{page 138}
on trouve
\[
\Sigma(X_*) =
\left(
\begin{tikzcd}[row sep=small, column sep=small]
 & \Sigma(X_1) \\
\Sigma(X_0) \arrow[ur, "\Sigma(\alpha_1)"] \arrow[dr, "\Sigma(\alpha_2)"'] & \\
 & \Sigma(X_2)
\end{tikzcd}
\right)
\]
et les flèches
\printedpage{94}
% head: A. Grothendieck : Dérivateurs, Ch. XII, Édition des Universités de Montpellier II et Paris VII
% foot: Version 7 décembre 2002
\[
\begin{tikzcd}[column sep=small]
 & \Sigma(X_*) \arrow[dl, "p_{X_*}"'] \arrow[dr, "q_{X_*}"] & \\
X_* & & X_*^\circ,
\end{tikzcd}
\]
dans $\underline{\Hom}(\Phi^\circ,\Cat)$. Ces flèches sont dans $W$ terme à terme, donc les homomorphismes dans $\Cat$
\[
\begin{tikzcd}[column sep=small]
 & \int\Sigma(X_*) \arrow[dl, "\int p_{X_*}"'] \arrow[dr, "\int q_{X_*}"] & \\
\sX=\int X_* & & \int X_*^\circ=\tsX^\circ
\end{tikzcd}
\]
sont dans $W$.

\dgrblock{dgrb465}{A common resolution of a category and its opposite}
\begin{dcomment}{exposition}{p137}{A common resolution of a category and its opposite}
Grothendieck's comparison uses a resolution that reverses the category without reversing the direction of the functors between categories. This distinction is responsible for the utility of $\Sigma$. Taking opposites is a covariant operation on $\Cat$: a functor $u:A\to B$ gives $u^\circ:A^\circ\to B^\circ$ in the same direction. Applied objectwise, it therefore preserves the shape of the span.

Here is a concrete convention for $\Sigma(C)$. Its objects are arrows $h:x\to y$ of $C$. An arrow from $h:x\to y$ to $h':x'\to y'$ is a pair
\[
 r:x\to x',\qquad s:y'\to y,
 \qquad h=s h'r.
\]
The projections $p_C(h)=x$ and $q_C(h)=y$ take values in $C$ and $C^\circ$, respectively. This is the opposite of one common convention for the twisted-arrow category. Reversing every displayed arrow identifies $\Sigma(C^\circ)$ with $\Sigma(C)$ and interchanges $p$ and $q$. The equality written in the manuscript may be read as this specified canonical identification.

For completeness, the asphericity of the projections can be seen from their fibres. The projection $p_C$ is a split fibration: over $r:x\to x'$ the inverse-image functor sends $h':x'\to y'$ to $h'r:x\to y'$. Its fibre over $x$ is $(x\backslash C)^\circ$, which has the final object $\mathrm{id}_x$. To test asphericity, pull back along $C/c\to C$. The resulting category $\Sigma(C)/c$ is again split fibred, now over $C/c$, with those same aspherical fibres. Since $C/c$ has a final object, the fibre theorem gives that $\Sigma(C)/c$ is $W$-aspherical. Thus $p_C$ is $W$-aspherical. Applying the identical argument to $C^\circ$ proves the assertion for $q_C$.

The importance of the resolution is its strict functoriality. For a diagram $F:I\to\Cat$, it supplies a diagram of natural transformations
\[
 F\ \xleftarrow{p_F}\ \Sigma F\ \xrightarrow{q_F}\ F^\circ
\]
whose components lie in $W$. Integrating the two transformations again gives weak equivalences. At the level of total categories this compares the smooth integral of $F$ with the smooth integral of $F^\circ$, the latter being canonically the opposite of the proper integral of $F$. Applying the same resolution once more to that total category removes the final opposite.

For a strict cocone $Q$ into $Y$, all these maps commute with the comparison maps into $Y$, $\Sigma(Y)$, or $Y^\circ$. The calculation on manuscript page 140 is therefore a calculation in an arrow category, not merely a list of equivalences between their sources. In each commuting square, both horizontal maps belong to $W$; two-out-of-three gives that one vertical comparison belongs to $W$ precisely when the other does. This proves the equivalence of the four conditions in Proposition 14(b), and the same zigzags prove the stronger assertions about coherent diagrams in parts (c) and (d).

This proof is also useful for its scope. It compares the images after inverting the chosen weak equivalences; it does not assert that $C$ and $C^\circ$ are equivalent categories. The fork $0\to1$, $0\to2$ already separates those statements, since it has an initial object and no final object, while its opposite has a final object and no initial object.
\end{dcomment}


Finalement, l'innocent carré commutatif
\[
Q \quad = \qquad
\begin{tikzcd}[row sep=small, column sep=small]
 & X_0 \arrow[dl] \arrow[dr] & \\
X_1 \arrow[dr] & & X_2 \arrow[dl] \\
 & Y &
\end{tikzcd}
\tag*{$(*)$}
\]
(et le diagramme correspondant
\[
\widehat{Q} \quad = \qquad
\begin{tikzcd}[row sep=small, column sep=small]
 & X_0 \arrow[dl] \arrow[dr] & \\
X_1 \arrow[dr] & & X_2 \arrow[dl] \\
 & \sX=\int X_* \arrow[d, "\varphi_Q"] & \\
 & Y &
\end{tikzcd}
\qquad\qquad )
\tag*{$(**)$}
\]
donne naissance, en le transformant
\msp{page 139}
par $X\longmapsto X^\circ$ et par $X\longmapsto\Sigma(X)$, et en tenant compte des homomorphismes fonctoriels en $X$ $\begin{tikzcd}[row sep=small, column sep=small]
 & \Sigma(X) \arrow[dl] \arrow[dr] & \\
X & & X^\circ
\end{tikzcd}$, à trois diagrammes du type $(**)$, i.e. modelés sur la catégorie opposée de
\[
\Gamma \quad = \qquad
\begin{tikzcd}[row sep=small, column sep=small]
 & 0 & \\
1 \arrow[ur, "\alpha_1"] & & 2 \arrow[ul, "\alpha_2"'] \\
 & 3 \arrow[ul, "i_1"] \arrow[ur, "i_2"'] & \\
 & 4 \arrow[u, "\varphi"'] &
\end{tikzcd}
\qquad
\text{\parbox[c]{0.35\textwidth}{\footnotesize (i.e. $\boldsymbol{\Delta}^1\times\boldsymbol{\Delta}^1$ (\footnotemark[88]), auquelle on a rajouté un objet initial)}} % [sic] auquelle -> à laquelle, grammatical form (catégorie is feminine)
\]
\footnotetext[88]{faux, car le carré n'est pas commutatif}
\printedpage{95}
% head: A. Grothendieck : Dérivateurs, Ch. XII, Édition des Universités de Montpellier II et Paris VII
% foot: Version 7 décembre 2002
(\footnote[89]{On a $\alpha_1i_1\neq\alpha_2i_2$, mais la relation $\alpha_1i_1\varphi=\alpha_2i_2\varphi$.}), savoir
\[
\widehat{Q}, \qquad \widehat{\Sigma(Q)}, \qquad \widehat{Q^\circ},
\]
et deux homomorphismes de diagrammes (\texttt{[dans]} $\underline{\Hom}(\Gamma^\circ,\Cat)$)
\[
\begin{tikzcd}[column sep=small]
 & \widehat{\Sigma(Q)} \arrow[dl] \arrow[dr] & \\
\widehat{Q} & & \widehat{Q^\circ}
\end{tikzcd}
\qquad
\text{(déduit de }
\begin{tikzcd}[column sep=small]
 & \Sigma(Q) \arrow[dl] \arrow[dr] & \\
Q & & Q^\circ
\end{tikzcd}
\text{),}
\tag*{$(***)$}
\]
lesquels sont dans $W$ terme à terme (i.e. sont dans $W_\Gamma\subseteq\Fl(\underline{\Hom}(\Gamma^\circ,\Cat))$). Explicitement, en les sommets $i=0,1,2$, resp. $4$ de $\Gamma$, $(***)$ n'est autre que
\[
D_i =
\left(
\begin{tikzcd}[column sep=small]
 & \Sigma(X_i) \arrow[dl, "p_{X_i}"'] \arrow[dr, "q_{X_i}"] & \\
X_i & & X_i^\circ
\end{tikzcd}
\right)
\qquad \text{resp.} \qquad
\begin{tikzcd}[column sep=small]
 & \Sigma(Y) \arrow[dl, "p_Y"'] \arrow[dr, "q_Y"] & \\
Y & & Y^\circ,
\end{tikzcd}
\]
en le sommet 3 c'est
\[
\begin{tikzcd}[column sep=small]
 & \int\Sigma(X_*) \arrow[dl, "\int p_{X_*}"'] \arrow[dr, "\int q_{X_*}"] & \\
\sX=\int X_* & & \int X_*^\circ=\tsX^\circ,
\end{tikzcd}
\]
obtenu par intégration des flèches
\[
\begin{tikzcd}[column sep=small]
 & \Sigma(X_*) \arrow[dl, "p_{X_*}"'] \arrow[dr, "q_{X_*}"] & \\
X_* & & X_*^\circ.
\end{tikzcd}
\]
\msp{page 140}
Regardant la restriction des deux morphismes de diagrammes $(***)$ aux parties des trois diagrammes de type $\Gamma^\circ$ au dessus de la queue $\{3,4\}$ de $\Gamma$ on trouve le diagramme en traits pleins que voici
\[
\begin{tikzcd}[column sep=3em, row sep=normal]
 & \int\Sigma(X_*) \arrow[dl, "W"'] \arrow[dr, "W"] \arrow[dd, "\varphi_{\Sigma(Q)}"'] & & \Sigma\tsX \arrow[dl, dotted, "W"'] \arrow[dr, dotted, "W"] \arrow[dd, dotted, "\Sigma\tilde{\varphi}_Q"] & \\
\sX \arrow[dd, "\varphi_Q"'] & & \tsX^\circ \arrow[dd, "\varphi_{Q^\circ}=(\tilde{\varphi}_Q)^\circ"'] & & \tsX \arrow[dd, dotted, "\tilde{\varphi}_Q"] \\
 & \Sigma(Y) \arrow[dl, "W"'] \arrow[dr, "W"] & & \Sigma(Y) \arrow[dl, dotted, "W"'] \arrow[dr, dotted, "W"] & \\
Y & & Y^\circ & & Y
\end{tikzcd}
\]
\printedpage{96}
% head: A. Grothendieck : Dérivateurs, Ch. XII, Édition des Universités de Montpellier II et Paris VII
% foot: Version 7 décembre 2002
qu j'ai complété par un deuxième volet en pointillés. Dans ce diagramme de quatre carrés accolés, les côtés horizontaux des carrés représentent des flèches de $\Cat$ qui sont dans $W$. On trouve ainsi : % [sic] qu j'ai -> que j'ai, missing letter
\[
\varphi_Q\in W \iff \varphi_{\Sigma Q}\in W \iff (\varphi_{Q^\circ}=\tilde{\varphi}_Q^\circ)\in W \iff \Sigma\tilde{\varphi}_Q\in W \iff \tilde{\varphi}_Q\in W.
\]
On trouve aussi que les cinq flèches verticales sont isomorphes, en tant que flèches de $\Hot_W$. Pour résumer :

\dgrblock{dgrb466}{The indexing category must retain both paths}
\begin{dcomment}{exposition}{p139}{The indexing category must retain both paths}
Grothendieck keeps the entire comparison diagram because the two routes from $X_0$ to the integral become equal only after applying the map to $Y$. The correction in footnotes 88 and 89 is essential to that purpose. Before reaching $Y$, the routes even land in different fibres of the integral. They cannot be made equal by calling the drawing a commutative square.

Let $J=\Gamma^\circ$. A precise presentation has vertices $0,1,2,3,4$, arrows
\[
 0\xrightarrow{a_1}1\xrightarrow{j_1}3\xrightarrow{v}4,
 \qquad
 0\xrightarrow{a_2}2\xrightarrow{j_2}3,
\]
and the single relation $v j_1a_1=v j_2a_2$. The two arrows $j_1a_1,j_2a_2:0\to3$ remain distinct. Thus $4$ is terminal, but the full subcategory on $0,1,2,3$ is the free diamond, not the ordered square $[1]\times[1]$. This is exactly the relation stated by footnote 89, with the orientation reversed.

To include the dotted central inclusion mentioned in footnote 91, add a generator $r:0\to3$ and the relation $vr=vj_1a_1$. The natural transformations connecting $j_1a_1$ and $j_2a_2$ with $r$ are additional two-dimensional data; they are not relations in this ordinary indexing category. Every resolution used in Proposition 14 also preserves $r$, so adding it causes no difficulty for the objectwise weak-equivalence argument. It simply records more of the construction.
\end{dcomment}


\textbf{Proposition 14.} \textit{Considérons un diagramme (carré commutatif)}
\[
Q \quad = \qquad
\begin{tikzcd}[row sep=small, column sep=small]
 & X_0 \arrow[dl] \arrow[dr] & \\
X_1 \arrow[dr] & & X_2 \arrow[dl] \\
 & Y &
\end{tikzcd}
\]
\textit{dans} $\Cat$\textit{, et le diagramme `opposé' déduit par} $X\longmapsto X^\circ$
\[
Q^\circ \quad = \qquad
\begin{tikzcd}[row sep=small, column sep=small]
 & X_0^\circ \arrow[dl] \arrow[dr] & \\
X_1^\circ \arrow[dr] & & X_2^\circ \arrow[dl] \\
 & Y^\circ, &
\end{tikzcd}
\]
\msp{page 141}
\textit{et les flèches correspondantes}
\[
\left\{
\begin{array}{rlclcl}
\varphi_Q : & \sX_Q & = &
\left(\displaystyle\int
\begin{tikzcd}[ampersand replacement=\&, row sep=small, column sep=small]
 \& X_1 \\
X_0 \arrow[ur] \arrow[dr] \& \\
 \& X_2
\end{tikzcd}
\right)
& \longrightarrow & Y \\[2ex]
\varphi_{Q^\circ} : & \sX_{Q^\circ} & = &
\left(\displaystyle\int
\begin{tikzcd}[ampersand replacement=\&, row sep=small, column sep=small]
 \& X_1^\circ \\
X_0^\circ \arrow[ur] \arrow[dr] \& \\
 \& X_2^\circ
\end{tikzcd}
\right)
& \longrightarrow & Y^\circ.
\end{array}
\right.
\]
\textit{Enfin, considérons aussi les flèches $\tilde{\varphi}_Q$, $\tilde{\varphi}_{Q^\circ}$ canoniques, relative à l'intégrale propre $\int_{\mathrm{pr}}$} % [sic] relative -> relatives, agreement with « les flèches »
\printedpage{97}
% head: A. Grothendieck : Dérivateurs, Ch. XII, Édition des Universités de Montpellier II et Paris VII
% foot: Version 7 décembre 2002
\textit{(au lieu de l'intégrale lisse $\int=\int_{\mathrm{lis}}$)}
\[
\left\{
\begin{array}{rlclcl}
\tilde{\varphi}_Q : & \tsX_Q & = &
\left(\displaystyle\int_{\mathrm{pr}}
\begin{tikzcd}[ampersand replacement=\&, row sep=small, column sep=small]
 \& X_1 \\
X_0 \arrow[ur] \arrow[dr] \& \\
 \& X_2
\end{tikzcd}
\right)
& \longrightarrow & Y \\[2ex]
\tilde{\varphi}_{Q^\circ} : & \tsX_{Q^\circ} & = &
\left(\displaystyle\int_{\mathrm{pr}}
\begin{tikzcd}[ampersand replacement=\&, row sep=small, column sep=small]
 \& X_1^\circ \\
X_0^\circ \arrow[ur] \arrow[dr] \& \\
 \& X_2^\circ
\end{tikzcd}
\right)
& \longrightarrow & Y^\circ.
\end{array}
\right.
\]
\textit{On a alors ce qui suit :}

\begin{enumerate}[label={}, leftmargin=2.5em]\itshape
\item[\circled{\textit{a}}] $\varphi_{Q^\circ}\simeq(\tilde{\varphi}_Q)^\circ$, $\varphi_Q\simeq(\tilde{\varphi}_{Q^\circ})^\circ$ (d'où résulte déjà que $\left\{\begin{array}{lcl}\varphi_Q\in W & \Longleftrightarrow & \tilde{\varphi}_{Q^\circ}\in W\\ \varphi_{Q^\circ}\in W & \Longleftrightarrow & \tilde{\varphi}_Q\in W\end{array}\right.$ (\footnote[90]{\upshape et que les flèches $\varphi_Q$ et $\tilde{\varphi}_{Q^\circ}$ d'une part, $\varphi_{Q^\circ}$ et $\tilde{\varphi}_Q$ d'autre part, sont isomorphes comme flèches dans $\Hot_W$.})).

\item[\circled{\textit{b}}] Les relations
\[
\varphi_Q\in W, \qquad \varphi_{Q^\circ}\in W, \qquad \tilde{\varphi}_Q\in W, \qquad \tilde{\varphi}_{Q^\circ}\in W
\]
sont équivalentes, en d'autres termes : pour que le carré soit $W$-homotopiquement cocartésien (version intégrale lisse), il faut et il suffit que le carré opposé $Q^\circ$ le soit, et ces conditions équivalent aussi aux conditions correspondantes en termes des intégrales propres (au lieu des intégrales lisses).

\item[\circled{\textit{c}}] Les quatre flèches $\varphi_Q$, $\varphi_{Q^\circ}$, $\tilde{\varphi}_Q$, $\tilde{\varphi}_{Q^\circ}$ de $\Cat$ en tant que flèches dans $\Hot_W$, sont isomorphes dans $\underline{\Fl}(\Hot_W)$. Plus précisément, en tant que flèches de $\Cat$,
\msp{page 142}
elles définissent des objets isomorphes de $\HOT_W(\boldsymbol{\Delta}^1)$.

\item[\circled{\textit{d}}] Plus précisément encore, les quatre diagrammes dans $\Cat$ de type $\Gamma^\circ$, où
\[
\Gamma = \;
\begin{tikzcd}[row sep=tiny, column sep=small]
 & 1 \arrow[dl, "\alpha_1"'] & & \\
0 & & 3 \arrow[ul, "i_1"'] \arrow[dl, "i_2"] & 4, \arrow[l, "\varphi"'] \\
 & 2 \arrow[ul, "\alpha_2"] & &
\end{tikzcd}
\]
$\alpha_1i_1\varphi=\alpha_2i_2\varphi$, dans lesquels s'insèrent les quatre flèches qu'on vient d'envisager, savoir
\[
\widehat{Q} = 
\begin{tikzcd}[row sep=tiny, column sep=small]
 & X_1 \arrow[dr, "i_1"] & & \\
X_0 \arrow[ur, "\alpha_1"] \arrow[dr, "\alpha_2"'] \arrow[rr, dotted] & & \sX_Q \arrow[r, "\varphi_Q"] & Y, \\
 & X_2 \arrow[ur, "i_2"'] & &
\end{tikzcd}
\]
\end{enumerate}

% ---- from Der-ChXIIb-012_p056-060.tex
\setcounter{footnote}{90}

\printedpage{98}
% head: A. Grothendieck : Dérivateurs, Ch. XII, Édition des Universités de Montpellier II et Paris VII
% foot: Version 7 décembre 2002
\[
\begin{tikzcd}[row sep=small, column sep=small]
 & X_1^\circ \arrow[dr,"i^{\prime}_{1}"] & & \\
\widehat{Q^\circ} = X_0^\circ \arrow[ur,"\alpha_1^\circ"] \arrow[dr,"\alpha_2^\circ"'] \arrow[rr,dotted] & & \mathcal{X}_{Q^\circ} \arrow[r] & Y^\circ, \\
 & X_2^\circ \arrow[ur,"i^{\prime}_{2}"'] & &
\end{tikzcd}
\]
\[
\begin{tikzcd}[row sep=small, column sep=small]
 & X_1 \arrow[dr] & & \\
\widehat{Q}\tld = X_0 \arrow[ur] \arrow[dr] \arrow[rr,dotted] & & \tilde{\mathcal{X}}_{Q} \arrow[r] & Y, \\
 & X_2 \arrow[ur] & &
\end{tikzcd}
\]
\[
\begin{tikzcd}[row sep=small, column sep=small]
 & X_1^\circ \arrow[dr] & & \\
\widehat{Q^\circ}\tld = X_0^\circ \arrow[ur] \arrow[dr] \arrow[rr,dotted] & & \tilde{\mathcal{X}}_{Q^\circ} \arrow[r] & Y^\circ, \\
 & X_2^\circ \arrow[ur] & &
\end{tikzcd}
\]
(\footnote{\textbf{NB} On devrait encore compléter ces diagrammes (non commutatifs) par les flèches en pointillés \dots}) \textit{sont isomorphes dans} $\Hot_W(\Gamma)$\textit{. De façon plus précise encore, on a les isomorphismes dans} $\uHom(\Gamma^\circ,\Cat)$
\[
(\widehat{Q}\tld)^\circ \simeq \widehat{Q^\circ},
\]
\textit{d'où en l'appliquant à $Q^\circ$ au lieu de $Q$,}
\[
(\widehat{Q^\circ}\tld)^\circ \simeq \widehat{Q},
\]
\textit{ce qui implique (via les isomorphismes canoniques valables dans} $\HOT_W(S)$ \textit{pour tout objet $S$ dans} $\Cat$ \textit{et tout} $X_* \in \mathrm{Ob}\,\uHom(S,\Cat)$\textit{,} $X_* \simeq X_*^\circ$ (\footnote{via le foncteur $\Sigma$ utilisé plus haut \dots})\textit{), qu'on a les isomorphismes canoniques}
\[
\left.\begin{aligned}
\widehat{Q} &\simeq \widehat{Q^\circ}\tld \\
\widehat{Q^\circ} &\simeq \widehat{Q}\tld
\end{aligned}\right\}
\qquad \text{\textit{dans} } \HOT_W(\Gamma),
\]
\textit{mais aussi directement}
\[
\widehat{Q} \simeq \widehat{Q^\circ}, \qquad \widehat{Q}\tld \simeq \widehat{Q^\circ}\tld \qquad \text{\textit{dans} } \HOT_W(\Gamma)
\]
\msp{page 143}\textit{(et il reste à vérifier qu'on a ainsi un système} transitif \textit{d'isomorphismes canoniques dans} $\HOT_W(\Gamma)$\textit{, entre les quatre objets envisagés, ce qui doit être essentiellement sorital \dots).}

\dgrblock{dgrb467}{The four comparisons form a coherent system}
\begin{dcomment}{development}{p143:et il reste à vérifier qu'on a ainsi un système}{The four comparisons form a coherent system}
Grothendieck leaves a small but real task after Proposition 14: the canonical comparisons should compose transitively. Without this check, choosing a route through the four constructions could leave an unspecified automorphism. The following argument supplies the check directly in the localized category of diagrams. We did not find this particular four-comparison verification stated separately in the literature consulted; it is a formal completion of the construction on these pages, with no claim of priority.

\textbf{Coherence of the four integration comparisons.}\ours\
Let $Q$ be a strictly commuting cocone on a small category-valued diagram, and use the same expanded indexing category $J$ for the integral, the cocone, and their structural inclusions. In the localization
\[
 \mathcal L_J=\operatorname{Fun}(J,\Cat)[W_J^{-1}],
\]
the comparisons furnished by $\Sigma$ between the smooth and proper integrals of $Q$ and of $Q^\circ$ form a transitive system of isomorphisms. The assertion includes the arrows from the integral to the cocone vertex and the optional central inclusions.

\emph{Proof.} Write $O$ for objectwise passage to opposite categories. It preserves $W_J$, so it acts on $\mathcal L_J$. The resolution of the preceding block gives a natural isomorphism
\[
 \eta_D:D\xrightarrow{\sim}OD,
 \qquad \eta_D=q_Dp_D^{-1}.
\]
Because $\Sigma(OD)$ is identified with $\Sigma(D)$ with its projections exchanged,
\[
 \eta_{OD}=\eta_D^{-1}.
\]
Naturality holds first for actual maps of diagrams and then for their inverses in the localization; hence it holds for every arrow in $\mathcal L_J$. Applying it to $\eta_D$ gives
\[
 O(\eta_D)\eta_D=\eta_{OD}\eta_D,
 \qquad O(\eta_D)=\eta_D^{-1}.
\]
Thus the opposite operation and its canonical comparison already satisfy the required involution identity.

Now put $A=\widehat Q_{\mathrm{lis}}$ and $B=\widehat{Q^\circ}_{\mathrm{lis}}$. Integrating the objectwise resolution of the original diagram, while using $\Sigma(Y)$ at the cocone vertex, gives
\[
 \alpha_Q:A\xrightarrow{\sim}B.
\]
The same resolution for $Q^\circ$ interchanges its two projections, including at the integrated vertex. Consequently $\alpha_{Q^\circ}=\alpha_Q^{-1}$. By the strict descriptions of the two Grothendieck constructions, the other two objects are canonically
\[
 C=\widehat Q_{\mathrm{pr}}=OB,
 \qquad D=\widehat{Q^\circ}_{\mathrm{pr}}=OA.
\]
Proper integration is the composite of objectwise opposite, smooth integration, and objectwise opposite again. Since $O(\eta_Q)=\eta_Q^{-1}$ also holds in the localized category of input cocones, applying this composite to the resolution shows that its direct proper comparison from $C$ to $D$ is $O(\alpha_Q^{-1})$. Naturality of $\eta$ for $\alpha_Q$ says exactly
\[
 \eta_B\alpha_Q=O(\alpha_Q)\eta_A,
\]
or, equivalently,
\[
 O(\alpha_Q^{-1})\eta_B\alpha_Q=\eta_A.
\]
This is the commutativity of the square formed by the four constructions. Together with the involution identities, it shows that the composite around the square, in either direction, is the identity. Every route between two vertices can be reduced by cancelling inverse adjacent edges and using this square identity. The resulting comparison is independent of the route, and its composites are transitive. All equalities take place in $\mathcal L_J$, so they also hold after restricting to the comparison arrow or to any other subdiagram.\hfill$\square$

The argument is unchanged for an arbitrary index category. For $F:I\to\Cat$, one may define $J$ by freely adjoining a vertex $v$ and arrows $r_i:i\to v$, then a terminal vertex $y$, so that the composites $i\to v\to y$ satisfy the relations of the given strict cocone. The arrows $r_j u$ and $r_i$ are not identified for $u:i\to j$; their images in the integral are linked by transition transformations. With this presentation, the arbitrary-shape claim on manuscript pages 143--144 follows from the same proof. The phrase ``without commutation relations'' describes the first, free adjunction of $v$; only the subsequent vertex $y$ is terminal. Keeping these two operations distinct prevents the indexing error already corrected on page 139.
\end{dcomment}


\subsection*{Remarques.}

\begin{enumerate}
\item[1)] Ce qu'on vient de dire dans le cas des diagrammes dans Cat de type $\Phi^\circ = \Psi$, a l'air de garder un sens pour les diagrammes de type quelconque
\[
X_* : S^\circ \longrightarrow \Cat, \qquad s \longmapsto X_s
\]
\printedpage{99}
% head: A. Grothendieck : Dérivateurs, Ch. XII, Édition des Universités de Montpellier II et Paris VII
% foot: Version 7 décembre 2002
($S$ un objet de Cat). On peut l'intégrer de façon lisse, $\int X_* = \int_{\mathrm{lis}} X_*$, obtenant une catégorie fibrée (scindée) sur $S$, ou de façon propre, $\int_{\mathrm{pr}} X_*$, obtenant une catégorie cofibrée scindée sur $S^\circ$. On peut également considérer $X_*^\circ$, composé de $X_*$ avec l'automorphisme $Y \longmapsto Y^\circ$ dans Cat, et l'intégrer de façon soit propre, soit lisse. Dans tous ces cas, on trouve un diagramme de type $\widehat{S}^\circ$, où $\widehat{S}$ est déduit de $S$ en ajoutant un élément qui s'envoie par une flèche ‘canonique’ (donnée) dans chacun des autres (\textit{sans} relations de commutation!) Un
\msp{page 144}\textit{système inductif} de morphismes dans Cat
\[
f \qquad X_s \xrightarrow{f_s} Y
\]
donne alors un morphisme
\[
\int_{\mathrm{lis}} X_* \xrightarrow{\varphi_Q} Y,
\]
où $Q$ désigne le diagramme, de type $\widehat{S}^\circ$, formé des $X_s$, de $Y$, des $u^* : X_t \longrightarrow X_s$ (pour $u : s \longrightarrow t$ dans $S$) et des $f_s$. De même on aura
\[
\int_{\mathrm{prop}} X_* \xrightarrow{\tilde{\varphi}_Q} Y,
\]
et si
\[
\Gamma = (\widehat{S} + \text{objet initial}),
\]
on trouve des diagrammes $\widehat{Q}$, $\widehat{Q}\tld$ de type $\Gamma^\circ$ dans Cat. Sauf erreur, les relations entre $\widehat{Q}$, $\widehat{Q}\tld$, $\widehat{Q^\circ}$, $\widehat{Q^\circ}\tld$ sont les mêmes que celles explicitées dans la proposition 14 ci-dessus. Je ne vais pas m'arrêter là-dessus avec le soin requis, sauf en cas de besoin. On dira qu'un diagramme $Q$ définit $Y$ comme $\varinjlim\limits_{I}$ $W$-homotopique \texttt{[où} $I = S^\circ$\texttt{]} si $\varphi_Q \in W$. Il revient au même de dire que $\varphi_{Q^\circ} \in W$, ou que $\tilde{\varphi}_Q \in W$ ou que $\tilde{\varphi}_{Q^\circ} \in W$ \dots
\msp{page 145}

\dgrblock{dgrb468}{The comparison map in later formalisms}
\begin{dcomment}{later}{p144}{The comparison map in later formalisms}
The definition by a specified map from an integral has a direct later interpretation. For ordinary homotopy types, put
\[
 \mathbb D(J)=\operatorname{Fun}(J,\Cat)[W_{\infty,J}^{-1}]
             =\HOT_{W_\infty}(J^{\mathrm{op}}).
\]
The opposite in the last expression corrects for the manuscript's presheaf convention. A square is an object of $\mathbb D(\square)$, where $\square=[1]\times[1]$. Let $i:\ulcorner\hookrightarrow\square$ omit the terminal vertex. In the derivator formulation a square $Q$ is cocartesian when the counit
\[
                   i_!i^*Q\longrightarrow Q
\]
is an isomorphism \citep[Definition~3.9]{Groth2013}. At the three vertices of $\ulcorner$ its components are already isomorphisms. At the last vertex, the component is the homotopy-colimit comparison represented here by $\varphi_Q$. This identifies Definition~3 with the later coherent-square definition. It also explains why a square merely drawn in $\mathbb D(e)$ is insufficient: it need not specify the coherent object to which $i_!$ applies.

The composition and cancellation theorem for adjacent cocartesian squares is \citet[Proposition~3.13]{Groth2013}. Its hypotheses concern squares inside one coherent rectangle. They do not assert that an arbitrary map between homotopy pushouts reflects an equivalence on the intersection from equivalences on the pieces and the union. The latter is the stronger cube assertion tested by the acyclic example above.

There is a second distinction when the fibres are retained as categories. A Cartesian fibration classified by a diagram of $\infty$-categories models its colimit only after its Cartesian transition edges have been inverted; see \citet[Corollary~3.3.4.3 and Warning~3.3.4.4]{Lurie2009HTT}. Passing further to the groupoid completion inverts every edge. The corresponding space-valued statement is Proposition~3.3.4.5 of the same reference. This is compatible with the classifying-space interpretation of the manuscript's integral, while preserving the stronger categorical information before localization.

The constant terminal-category diagram over $[1]$ makes the distinction concrete. Its ordinary total category is $[1]$. The colimit of the diagram in categories is $e$, and its colimit in $\mathrm{Cat}_\infty$ is also $e$. The unique nonidentity arrow of $[1]$ is a transition arrow; inverting it gives a category equivalent to $e$. Before that inversion, $[1]$ and $e$ are not equivalent categories: the two objects of $[1]$ are not isomorphic. Their classifying spaces are nevertheless both contractible. Thus the three assertions---a total category, a categorical colimit, and a colimit of homotopy types---agree only after the indicated operations. The present text is computing the last of them and retaining a strict representative for its map to the proposed target.
\end{dcomment}


\item[2)] Le cas ‘universel’, pour
\begin{tikzcd}[baseline=(\tikzcdmatrixname-2-1.base), row sep=small, column sep=small]
 & X_1 \\
X_0 \arrow[ur] \arrow[dr] & \\
 & X_2
\end{tikzcd}\par\smallskip
% Reader layout: separate the two inline diagrams without changing the wording.
donné, d'un carré commutatif $Q = {}$%
\begin{tikzcd}[baseline=(\tikzcdmatrixname-2-1.base), row sep=tiny, column sep=small]
 & X_1 \arrow[dr] & \\
X_0 \arrow[ur] \arrow[dr] & & Y \\
 & X_2 \arrow[ur] &
\end{tikzcd}
dans Cat, donnant naissance à
\[
\begin{tikzcd}[row sep=small, column sep=small]
 & & & X_1 & & & \\
\varphi_Q : &[1em] \left(\vrule height 6ex depth 6ex width 0pt\right. &[-1.5em] \int X_0 \arrow[ur] \arrow[dr] & & \left.\vrule height 6ex depth 6ex width 0pt\right) &[2em] \longrightarrow &[2em] Y, \\
 & & & X_2 & & &
\end{tikzcd}
\]
est celui de $Y = X_1 \amalg_{X_0} X_2$, donnant naissace à % [sic] naissace -> naissance, letter n omitted
\printedpage{100}
% head: A. Grothendieck : Dérivateurs, Ch. XII, Édition des Universités de Montpellier II et Paris VII
% foot: Version 7 décembre 2002
\[
\begin{tikzcd}[row sep=small, column sep=small]
 & X_1 \arrow[dr] & \\
Q_0 = Q(X_*) = X_0 \arrow[ur] \arrow[dr] & & X_1 \amalg_{X_0} X_2. \\
 & X_2 \arrow[ur] &
\end{tikzcd}
\]

Le cas général s'y ramène, car on aura un diagramme commutatif
\[
\begin{tikzcd}[row sep=small, column sep=small]
 & & X_1 & & & \\
\left(\vrule height 6ex depth 6ex width 0pt\right. &[-1.5em] \int X_0 \arrow[ur] \arrow[dr] & & \left.\vrule height 6ex depth 6ex width 0pt\right) \arrow[rr,"{\varphi_{Q_0}=\varphi_{Q(X_*)}}"] \arrow[ddrr,"\varphi_Q"] &[4em] \hphantom{MMMM} &[4em] X_1 \amalg_{X_0} X_2 \arrow[dd,"\alpha_Q"] \\
 & & X_2 & & & \\
 & & & & & Y.
\end{tikzcd}
\]
\end{enumerate}

Mais il doit y avoir des cas intéressants où $\varphi_{Q_0}$ n'est \textit{pas} dans $W$, i.e. le carré Cat-cocartésien
\begin{tikzcd}[baseline=(\tikzcdmatrixname-2-1.base), row sep=tiny, column sep=small]
 & X_1 \arrow[dr] & \\
X_0 \arrow[ur] \arrow[dr] & & X_1 \amalg_{X_0} X_2 \\
 & X_2 \arrow[ur] &
\end{tikzcd}
n'est pas $W$-homotopiquement cocartésien, et où pourtant on peut trouver un carré commutatif $Q$
\msp{page 146}complétant $X_*$, de telle façon que ce carré soit $W$-homotopiquement commutatif \texttt{[plutôt cocartésien]}, i.e. que $\varphi_Q = \alpha_Q \varphi_{Q_0}$ soit dans $W$. C'est d'autant plus plausible que dans Cat, les sommes amalgamées ont tendance à ne pas avoir des bonnes propriétés homotopiques, même si l'une des flèches $\alpha_1$, $\alpha_2$ est \texttt{[un]} monomorphisme (voire une immersion ouverte), contrairement à ce qui se passe dans $\boldsymbol{\Delta}^\wedge$, et plus généralement dans les catégories de la forme $Z^\wedge$ ($Z$ dans Cat).

\dgrblock{dgrb469}{A strict cocone can repair the wrong pushout}
\begin{dcomment}{development}{p146:de telle façon que ce carré soit $W$-homotopiquement}{A strict cocone can repair the wrong pushout}
Grothendieck asks for more than a strict pushout that fails to be a homotopy pushout. He asks whether its failure can disappear after mapping that pushout to a different target, while retaining the original span. The following finite example answers his question affirmatively for $W_\infty$, even when one leg is an open immersion. The sources checked on preservation of pushouts, including the Dwyer-map criterion and its refinement in \citet[Theorems~1.5--1.6, arXiv v3]{HackneyOzornovaRiehlRovelli2024}, do not supply this particular construction. We give a complete argument rather than a priority claim.

\textbf{A finite repair example.}\ours\
There exist finite posets $E\subset B$, with $E$ a sieve in $B$, such that $BE$ and $BB$ are contractible, while the categorical pushout $P=B\amalg_E e$ has classifying space $S^3$. Consequently the canonical comparison
\[
 c:\int_{\mathrm{lis}}(B\leftarrow E\to e)\longrightarrow P
\]
is not in $W_\infty$, but its composite with $P\to e$ is in $W_\infty$. Here $BD=|ND|$ denotes the classifying space of a category $D$.

\emph{Proof.} We first construct a map from a contractible face poset to a spherical one whose comma categories are connected. Let $K$ be the finite simplicial complex with six vertices $a,b,c,A,B,C$ and maximal faces
\[
 abc,\quad abB,\quad aBA,\quad bcC,\quad bCB,\quad caA,\quad cAC.
\]
It is the disk obtained by placing the triangle $abc$ inside the triangle $ABC$ and triangulating the surrounding annulus with the remaining six triangles. In particular it is contractible. This description can also be checked without a drawing: successive free-face collapses remove
\[
\begin{gathered}
(AB,aAB),\ (AC,cAC),\ (aA,acA),\ (A,cA),\ (BC,bBC),\\
(aB,abB),\ (B,bB),\ (bC,bcC),\ (C,cC),\\
(ab,abc),\ (a,ac),\ (b,bc),
\end{gathered}
\]
leaving the vertex $c$. In each pair the first simplex is a free codimension-one face of the second at the moment it is removed.

Let $L$ be the boundary of the tetrahedron on vertices $a,b,c,n$. The vertex map fixing $a,b,c$ and sending $A,B,C$ to $n$ is simplicial from $K$ to $L$: the seven triangles have images $abc$, $abn$, $an$, $bcn$, $bn$, $can$, $cn$. Let $E$ and $S$ be the posets of nonempty faces of $K$ and $L$, and let $f:E\to S$ send a face to its image. Thus $BE$ is a subdivided disk and $BS$ a subdivided two-sphere.

For every face $s$ of $L$, the comma category $E/s$ is the face poset of the subcomplex consisting of faces whose image is contained in $s$. If $n\notin s$, this subcomplex is the simplex on $s$, so it is connected and nonempty. If $n\in s$, it contains the entire outer boundary $AB\cup BC\cup CA$. Every inner vertex allowed by $s$ attaches to that boundary by a radial edge: $a$ by $aA$, $b$ by $bB$, and $c$ by $cC$. Every remaining simplex contains vertices of this same connected subcomplex. It follows that
\[
 \pi_0(E/s)=\{*\}\quad\text{for every }s\in S.
\]
Only connectedness is claimed: for $s=\{n\}$ the comma category is the face poset of the outer boundary circle.

Adjoin an initial object $o$ to $S$ and call the resulting poset $T$. Define a poset $B$ on the disjoint union of $E$ and $T$. Keep their internal orders and add
\[
 e\leq s\quad\Longleftrightarrow\quad f(e)\leq s
 \qquad(e\in E,\ s\in S).
\]
There are no relations from $T$ to $E$, and none from $E$ to $o$. Transitivity follows from monotonicity of $f$. The subposet $E$ is a sieve. The functor $r:B\to T$ which is the identity on $T$ and sends $e$ to $f(e)$ is left adjoint to the inclusion $j:T\to B$, because
\[
 B(b,jt)=T(r(b),t)
\]
for every $b\in B$, $t\in T$, including the empty hom-sets when $t=o$ and $b\in E$. Hence $Bj$ is a homotopy equivalence. The initial object $o$ makes $BT$, and therefore $BB$, contractible.

It remains to compute the categorical pushout, since a topological quotient would give the wrong answer. Collapsing the sieve $E$ to one object $*$ introduces no arrows from $T$ to $*$ and does not change arrows within $T$. Any arrow from $*$ to $s\in S$ is represented by a relation $e\leq s$ in $B$. The relations after collapsing $E$ identify two representatives exactly along zigzags of comparable elements in $E/s$. Thus
\[
 \operatorname{Hom}_{P}(*,s)=\pi_0(E/s)=\{*\},
 \qquad \operatorname{Hom}_{P}(*,o)=\emptyset.
\]
For clarity, this description also verifies the universal property: a functor out of $B$ that is constant on $E$ assigns the same arrow to any two representatives in one such connected component, and therefore factors uniquely through the displayed category. Conversely every functor from that category supplies exactly such compatible data on $B$ and $e$. Composition is induced by the maps $E/s\to E/s'$ for $s\leq s'$.

Consequently $P$ is the poset with two incomparable minimal objects $o,*$, both lying below every object of $S$. Its order complex is exactly the simplicial join
\[
 NP=N\{o,*\}*NS,
 \qquad BP\cong S^0*|\operatorname{sd}L|\cong S^3.
\]
The three categories in the original span $B\leftarrow E\to e$ are contractible. Their smooth integral is a split fibration over the contractible indexing fork, with contractible fibres, and hence is contractible. Its canonical map $c$ cannot be a weak equivalence to $P$, whose third homology is $\mathbb Z$. On the other hand the unique functor from the integral to $e$ is a weak equivalence, and it is precisely the composite $P\to e$ after $c$. This proves the assertion.\hfill$\square$

The construction has $25$ objects in $E$, $14$ in $S$, $40$ in $B$, and $16$ in $P$. The accompanying finite computation checks the face lists, the elementary collapses, all fourteen connected comma categories, and the homology of $NP$; the join calculation above identifies its entire homotopy type. The mechanism is visible in the proof: the ordinary categorical pushout keeps only $\pi_0(E/s)$ in each mixed hom-set. It replaces a family containing a circular comma category by a family of singletons, and the resulting total category acquires the additional sphere. Sending that sphere to a point supplies the strict cocone envisaged on manuscript pages 145--146.
\end{dcomment}


\subsection*{3.4\quad Critères pour les carrés homotopiquement cocartésiens.}

Les critères que j'ai dégagé sont tous contenus dans l'énoncé suivant :

\textbf{Théorème 4} (\footnote{\textbf{NB} Dans ce théorème, et les corollaires 1, 2, 3, on n'a besoin que de la \textit{forme faible} de W(7 bis), savoir $i_0 \in W \Longrightarrow i \in W$ (cf. p. 79) (Loc(6$'$a) avec les notations de XVI).},\footnote{\textbf{NB} Au total, le théorème (partie directe) est valable sans hypothèse W(7 bis) ou W(7 bis faible) sur $W$, et il en est de même de ses corollaires 1 à 3 (à l'exclusion des corollaires 4, 5, qui eux utilisent W(7 bis) dans toute sa force).}). \textit{Considérons un diagramme carré commutatif dans} Cat
\[
Q \quad = \quad
\begin{tikzcd}
X_0 \arrow[r,"\alpha_1"] \arrow[d,"\alpha_2"'] \arrow[dr,dotted,"f_0"] & X_1 \arrow[d,"f_1"] \\
X_2 \arrow[r,"f_2"'] & Y,
\end{tikzcd}
\tag*{$(*)$}
\]
\printedpage{101}
% head: A. Grothendieck : Dérivateurs, Ch. XII, Édition des Universités de Montpellier II et Paris VII
% foot: Version 7 décembre 2002
\textit{et considérons le morphisme canonique}
\[
\begin{tikzcd}[row sep=small, column sep=small]
 & & X_1 & & & \\
\varphi_Q : \quad \mathcal{X} = \Bigl(\int &[-1.2em] X_0 \arrow[ur,"\alpha_1"] \arrow[dr,"\alpha_2"'] & & \Bigr) &[2em] \longrightarrow &[2em] Y, \\
 & & X_2 & & &
\end{tikzcd}
\]
\msp{page 147}\textit{l'intégrale étant l'intégrale} lisse. \textit{(Il y a un énoncé dual concernant l'intégrale propre de}\newline
\begin{tikzcd}[baseline=(\tikzcdmatrixname-2-1.base), row sep=small, column sep=small]
 & X_1 \\
X_0 \arrow[ur] \arrow[dr] & \\
 & X_2
\end{tikzcd}
\textit{.)}

\begin{enumerate}
\item[\textit{a)}] \textit{Pour que $\varphi_Q$ soit $W$-lisse, il faut et il suffit que $f_1$, $f_2$ et $f_0 = f_1\alpha_1 = f_2\alpha_2$ le soient.} (Moyennant W(7 bis) (\footnote{Pour la partie directe ($f_0$, $f_1$, $f_2$ lisses $\Longrightarrow \varphi_Q$ lisse), on n'a besoin que de W(7 bis) (\footnotemark), forme faible. Ce n'est que pour la réciproque (qui ne m'a encore jamais servie) qu'on a besoin de W(7 bis), dans toute sa force.}\footnotetext{On n'a pas même besoin de W(7 bis) faible, cf. annotation page 156. On utilise une hypothèse restrictive sur $W$ dans la réciproque seulement ($\varphi_Q$ lisse $\Longrightarrow f_0$, $f_1$, $f_2$ lisses).})).

\item[\textit{b)}] \textit{La fibre de $\varphi_Q$ en un élément $y$ de $Y$ est isomorphe à}
\[
\begin{tikzcd}[row sep=small, column sep=small]
 & & X_{1\,y} & \\
\int &[-1.2em] X_{0\,y} \arrow[ur] \arrow[dr] & &[2em] (\simeq \mathcal{X}_y). \\
 & & X_{2\,y} &
\end{tikzcd}
\tag*{$(**)$}
\]
\end{enumerate}

\textbf{Corollaire 1.} (\footnote{Valable sans hypothèse sur $W$ (autre que Loc(1,2,3 bis)).}) \textit{Supposons $f_1$, $f_2$ et $f_1\alpha_1 = f_2\alpha_2$ $W$-lisses, et les catégories intégrées $(**)$ $W$-asphériques pour tout $y \in \mathrm{Ob}\,Y$. (NB Il s'agit encore de l'intégrale lisse, mais peu importe au fond, car on a vu (cf. partie supérieure du diagramme p. 140, qui implique que $\mathcal{X}$} \texttt{[est]} \textit{$W$-asphérique si et seulement si $\tilde{X}$ l'est) que l'une est $W$-asphérique si et seulement si l'autre l'est.) Alors $\varphi_Q \in W$, i.e. le carré $(*)$ est $W$-homotopiquement cartésien} \texttt{[plutôt cocartésien]}\textit{. Par dualité, même conclusion si on suppose $f_1$, $f_2$, $f_1\alpha_1 = f_2\alpha_2$} propres \textit{au lieu de lisses. (Car $Q$ $W$-homotopiquement trivial} \texttt{[plutôt cocartésien]}\textit{, si et seulement si $Q^\circ$ l'est, par} \texttt{[la]} \textit{prop. 14.)}

\msp{page 148}En effet, un morphisme dans Cat qui est $W$-lisse et à fibres $W$-asphériques est dans $W$ (et c'est même une $W$-équivalence universelle \dots).

\dgrblock{dgrb470}{Smoothness and the two different uses of a fibre}
\begin{dcomment}{exposition}{p147}{Smoothness and the two different uses of a fibre}
Grothendieck now turns the comparison into a criterion that can be checked on a given square. There are two separate operations in Theorem 4. The fibre formula is an identity of categories: imposing that the comparison map has the value $y$ imposes the same condition in each of its three pieces. Smoothness is a local lifting condition, tested on slices at objects of the integral. The latter condition is what permits a strict fibre to determine the homotopy-theoretic conclusion.

It is useful to read the annotations as a refinement of the hypotheses, rather than as several competing versions of the theorem. If $f_0,f_1,f_2$ are $W$-smooth, then $\varphi_Q$ is $W$-smooth for every fundamental localizer. The reverse implication for $f_0$ is asserted under the strong form of W(7 bis). No extra excision axiom is required for the fibre identity
\[
 (\mathcal X_Q)_y\cong
 \int_{\mathrm{lis}}(X_{1y}\leftarrow X_{0y}\to X_{2y}).
\]
The calculations on manuscript pages 152--157 establish these assertions; their local content is explained below at the two lemmas. In particular, the annotations do not justify removing the strong hypothesis from the converse.

The last step of Corollary 1 deserves to be made explicit. If $u:E\to Y$ is $W$-smooth, the fibre inclusion $E_y\to y\backslash E$ is $W$-aspherical. If every $E_y$ is $W$-aspherical, two-out-of-three gives that every $y\backslash E$ is $W$-aspherical. The dual of the fundamental Theorem A then yields $u\in W$. This proof is stable under arbitrary base change: smoothness is stable under base change and the new strict fibres are the old ones. Hence $u$ is in fact a universal $W$-equivalence. The same argument with properness uses $E/y$ instead. These are the fibre criteria formalized in \citet[Theorem~1.2]{Maltsiniotis2005Asphericite}.

Consequently the forward criterion in Corollary 1 has precisely the strength needed later: smoothness of all three structure maps, together with asphericity of the integrated fibres, proves the square homotopy cocartesian, and even proves the comparison a universal weak equivalence. The qualifications in the first two footnotes to Theorem 4 correctly anticipate this result.
\end{dcomment}


\textbf{Corollaire 2}. \textit{On suppose que $W$ satisfait} W(7 bis) (axiome du carré cocartésien) (\footnote{inutile.}).
\printedpage{102}
% head: A. Grothendieck : Dérivateurs, Ch. XII, Édition des Universités de Montpellier II et Paris VII
% foot: Version 7 décembre 2002
\textit{Supposons le carré $Q$ cocartésien dans} Cat (\footnote{NB On n'utilise pas vraiment que $Q$ soit cocartésien, plutôt que $Q$ est \textit{cartésien,} et $f_2$ (donc aussi $\alpha_1$) une immersion ouverte (resp. fermée), et que pour $y \in Y\backslash X_2$, $X_{1\,y}$ soit $W$-asphérique. (Dans le cas où $Q$ est cocartésien, $X_{1\,y}$ est $\simeq e$.)})\textit{, et que $\alpha_1$ soit une immersion ouverte (resp. fermée) et $f_1$ lisse (resp. propre). Alors $Q$ est $W$-homotopiquement cocartésien.} (\footnote{On a besoin ici de W(7 bis), forme directe, i.e. de Loc(6$'$a) (notation de XVI).})

Il suffit d'établir le résultat non respé, l'autre en résultera par dualité (puisque $Q$ est $W$-cocartésien si et seulement si $Q^\circ$ l'est). Par hypothèse $f_1$ est $W$-lisse, d'autre part $\alpha_1$ et $f_2$ sont des immersions ouvertes, donc $W$-lisses, donc $f_1\alpha_1$ est $W$-lisse aussi. (D'ailleurs $\alpha_2$ aussi, car le carré est aussi cartésien, donc $\alpha_2$ déduit de $f_1$ (qui est $W$-lisse par hypothèse) par changement de base $X_2 \longrightarrow Y$.) Donc $\varphi_Q$ est \texttt{[}$W$\texttt{-]}lisse, et il suffit de voir que ses fibres sont $W$-asphériques. Le calcul $(**)$ donne, pour $y$ dans $Y\backslash X_2$,
\[
\begin{tikzcd}[row sep=small, column sep=small]
 & & e & \\
\mathcal{X}_y \simeq \quad \int &[-1.2em] \emptyset \arrow[ur] \arrow[dr] & &[1em] \simeq e, \qquad \text{O.K.} \\
 & & \emptyset &
\end{tikzcd}
\]
Et pour $y$ dans $X_2$
\[
\begin{tikzcd}[row sep=small, column sep=small]
 & & X_{1\,y} & \\
\mathcal{X}_y \simeq \int &[-1.2em] X_{0\,y} \arrow[ur,"\sim"] \arrow[dr] & &[2em] . \\
 & & e &
\end{tikzcd}
\]
\msp{page 149}Comme $X_{0\,y} \longrightarrow X_{1\,y}$ est un isomorphisme, donc est dans $W$, il s'ensuit que $e \longrightarrow \mathcal{X}_y$ l'est aussi (prop. 13, p. 125), moyennant W(7 bis), donc $\mathcal{X}_y$ est $W$-asphérique, O.K. (\footnote{On n'a besoin que de la forme directe de W(7 bis), i.e. (notation de XVI) de Loc(6$'$a).}).

\textbf{Corollaire 3.} (Moyennant W(7 bis) (\footnote{inutile.}).) \textit{Soit $X$ dans} Cat\textit{, $U_1$, $U_2$ deux ouverts de $X$ tels que $X = U_1 \cup U_2$, $U_0 = U_1 \cap U_2$ l'ouvert intersection. Alors le carré d'immersions ouvertes}
\[
\begin{tikzcd}[row sep=large, column sep=large]
U_0 \arrow[r,hook] \arrow[d,hook] & U_1 \arrow[d,hook] \\
U_2 \arrow[r,hook] & X
\end{tikzcd}
\]
\textit{est $W$-cocartésien. (Énoncé dual pour deux fermés $T_1$, $T_2$ tels que $T_1 \cup T_2 = X$, avec $T_0 = T_1 \cap T_2$.)}

On applique le corollaire 2. On sait que le carré est cocartésien dans Cat (\footnote{On n'a pas même besoin de le savoir, cf. annotation marginale du corollaire 2.}), de plus $U_1 \lhook\joinrel\longrightarrow X$ est bien $W$-lisse, étant une immersion ouverte.

% ---- from Der-ChXIIb-012_p061-065.tex
\printedpage{103}
% head: A. Grothendieck : Dérivateurs, Ch. XII, Édition des Universités de Montpellier II et Paris VII

\noindent\textbf{Corollaire 4.} (Moyennant W(7 bis) (\footnote[104]{Ici on a besoin de W(7 bis) forme forte, i.e. $i_0\in W \iff i\in W$ (cf. annotation th. 4).}).) \textit{Soit $U \hto[i] X$ une immersion ouverte telle que $T=X\backslash U$ soit isomorphe à la catégorie ponctuelle $e$. Pour que $i\in W$, il faut et il suffit que $i$ soit $W$-asphérique, i.e. que $U/t$ (où $t$ est l'unique objet de $T$) soit $W$-asphérique.}

Il suffit bien sûr de prouver que
\[
i\in W \Longrightarrow U/t \ W\text{-asphérique.}
\]

\msp{page 150}
Considérons le carré cartésien
\[
\begin{tikzcd}[row sep=large,column sep=large]
U/t \arrow[r, hook, "i/t"] \arrow[d, "p"'] & X/t \arrow[d, "q"] \\
U \arrow[r, hook, "i"] & X,
\end{tikzcd}
\]
nous allons prouver qu'il est $W$-cocartésien. Moyennant W(7 bis), il s'ensuivra que $i\in W$ implique $i/t\in W$, et comme $X/t$ est $W$-asphérique (ayant un objet final), cela signifie que $U/t$ est $W$-asphérique, ce qu'on veut prouver.

\dgrblock{dgrb471}{Why the cover criterion needs no converse excision}
\begin{dcomment}{examples}{p149}{Why the cover criterion needs no converse excision}
Grothendieck's cover criterion is meant to provide a large supply of actual homotopy pushouts. The annotations correctly remove W(7 bis) from Corollaries 2 and 3, although the displayed proof still invokes it. The missing simplification is elementary: one of the maps in the only nontrivial fibre is an isomorphism, and that special case admits a direct contraction.

For a category $A$, consider
\[
 E=\int_{\mathrm{lis}}(A\xleftarrow{\mathrm{id}}A\longrightarrow e).
\]
Keep the central copy $A_0$ and the outer point $*$, and call their full subcategory $K$. There is a unique arrow $*\to a$ for every $a\in A_0$, so $*$ is initial in $K$. Define $R:E\to K$ to be the identity on $K$ and to send the outer copy $a\in A_1$ to its central copy in $A_0$. On mixed arrows it uses the corresponding arrow of $A$. The canonical transition arrows $a_1\to a_0$ give a natural transformation $\mathrm{id}_E\to jR$, where $j:K\to E$ is the inclusion, and $Rj=\mathrm{id}_K$. Thus $j$ and $R$ are homotopy inverse after taking nerves. Since $K$ has an initial object, $E$ is $W$-aspherical for every fundamental localizer. This proof also covers $A=\emptyset$.

In Corollary 2, a point outside $X_2$ has integrated fibre $e$. At a point of $X_2$, cartesianness identifies $X_{0y}\to X_{1y}$ with an isomorphism, and the fibre is of the form just contracted. The three maps to $Y$ are smooth: $f_1$ by hypothesis, $f_2$ because it is a sieve inclusion, and $f_0=f_1\alpha_1$ by composition. The forward part of Theorem 4 therefore proves the conclusion without W(7 bis). This also establishes the stronger version recorded in footnote 99: strict cocartesianness may be replaced by cartesianness, the sieve condition on $f_2$, smoothness of $f_1$, and $W$-asphericity of $X_{1y}$ outside $X_2$. In that version the exterior fibres need not be singletons.

For a cover $X=U_1\cup U_2$ by sieves, there is a second, particularly transparent argument in ordinary homotopy theory. Every simplex $x_0\to\cdots\to x_n$ in $NX$ lies in a member of the cover containing $x_n$, because a sieve contains every predecessor of its objects. Therefore
\[
 NX=NU_1\cup_{N(U_1\cap U_2)}NU_2
\]
is an equality of simplicial sets. The two inclusions are monomorphisms, so its realization is a pushout along subcomplex inclusions and computes the homotopy pushout. For cosieves one uses the first vertex $x_0$ instead. This calculation explains why categorical open and closed covers are especially well adapted to excision.

As a worked example, let $X$ be the poset of nonempty proper faces of a tetrahedron. Let $U_1$ consist of the faces of one facet, and let $U_2$ consist of the faces of the other three facets. Both are downward closed, and their union is $X$. Their intersection is the face poset of the boundary of the first facet. The order complexes are barycentric subdivisions: $BU_1$ is one triangle, $BU_2$ is the disk made of the other three triangles, and $B(U_1\cap U_2)$ is their common boundary circle. The integral therefore has homotopy type
\[
 D^2\mathop{\cup}^{h}_{S^1}D^2\simeq S^2,
\]
and the canonical comparison identifies this sphere with $BX$, the subdivided tetrahedral boundary. The comparison is determined by the actual inclusions of face posets; it is not merely an abstract equivalence between two spaces known to be spheres.
\end{dcomment}


Il suffit, en vertu du corollaire 2, de prouver que le carré est cocartésien dans Cat, puisque $i/t$ est une immersion ouverte ($i$ l'étant), et $X/t\longrightarrow X$ est un isomorphisme local, donc $W$-lisse. On voit que $(X/t)\backslash(U/t)\simeq e$ (catégorie ponctuelle), d'ailleurs l'unique élément de $(X/t)\backslash(U/t)$ est l'élément final $t_0 : t\xrightarrow{\mathrm{id}_t} t$ de $X/t$. Ainsi $H_{t_0}\simeq e_{U/t}$ (préfaisceau final), et $p_{!}(H_{t_0})\simeq H'_t \overset{\text{déf}}{=} (u\longmapsto \mathrm{Hom}_X(u,t))$, ce qui caractérise la somme amalgamée dans Cat
\dots

Le corollaire 4 n'est autre que le corollaire non démontré de la proposition 7 (p. 85).

Voici une variante du corollaire 4.

\noindent\textbf{Corollaire 5.} (Moyennant W(7 bis) (\footnote[105]{forme forte.}).) \textit{Soit $i : U \hto X$ une immersion ouverte, soit $t_0\in\mathrm{Ob}\,T$ un point \textup{fermé} de $T=X\backslash U$, tel que $\{t_0\}$}
\msp{page 151}
\textit{soit isomorphe à la catégorie ponctuelle, i.e. $\mathrm{Hom}_X(t_0,t_0)=\{\mathrm{id}\}$. Supposons que pour tout $t\in T\backslash\{t_0\}$, $U/t$ soit $W$-asphérique. Pour que l'on ait $i\in W$, il faut et il suffit que $i$ soit asphérique, i.e. que $U/t_0$ soit également $W$-asphérique.}

Il faut prouver que $i\in W$ implique que $U/t_0$ est $W$-asphérique. Considérons l'inclusion
\[
i_0 : U\longrightarrow U' \overset{\text{déf}}{=} X\backslash\{t_0\},
\]
l'hypothèse sur les $U/t$, pour $t\in T\backslash\{t_0\}$, signifie que $i_0$ est $W$-asphérique. Donc $i_0$ est dans $W$, et comme le composé $i : U\xrightarrow{i_0} U'\xrightarrow{i'} X$ est dans $W$, cela implique que
% foot: Version 7 décembre 2002
\printedpage{104}
% head: A. Grothendieck : Dérivateurs, Ch. XII, Édition des Universités de Montpellier II et Paris VII
$i' : U'\longrightarrow X$ est dans $W$. Nous sommes maintenant dans la situation du corollaire 4, qui implique que $i'$ est $W$-asphérique. Donc $i=i'i_0$ est $W$-asphérique, comme composé de deux morphismes $W$-asphériques, q.e.d.

\dgrblock{dgrb472}{A single missing point and the role of the strong axiom}
\begin{dcomment}{exposition}{p151:comme composé de deux morphismes $W$-asphériques}{A single missing point and the role of the strong axiom}
Grothendieck uses the preceding criterion to finish a result deferred earlier in the chapter. When only one object is missing, the slice over that object records every possible attachment to the open subcategory. The square on manuscript page 150 is consequently both a strict pushout and a homotopy pushout. The remaining implication is a reflection statement, and this is where the strong form of W(7 bis) enters.

Here is the strict pushout assertion in detail. Suppose $U\subset X$ is a sieve, its complement consists of one object $t$, and $\operatorname{End}_X(t)=\{\mathrm{id}_t\}$. Let $F:U\to Z$ and $G:X/t\to Z$ agree on $U/t$ after projection to $U$. Define $H:X\to Z$ on $U$ to be $F$, and set $H(t)=G(\mathrm{id}_t)$. For an arrow $a:u\to t$, define $H(a)$ to be the image under $G$ of the canonical arrow from $(u,a)$ to $(t,\mathrm{id}_t)$ in $X/t$. The source of this image is $F(u)$ by the agreement assumption. If $b:u'\to u$ lies in $U$, the factorization of $(u',ab)\to(t,\mathrm{id}_t)$ through $(u,a)$ gives
\[
 H(ab)=H(a)F(b).
\]
There are no arrows from $t$ into $U$, since $U$ is a sieve, and the only endomorphism at $t$ is its identity. Thus these rules define a functor, and they determine it uniquely. This proves the pushout universal property of the slice square, without any appeal to the presheaf notation $H_{t_0}$ used in the manuscript.

The projection $X/t\to X$ is a discrete fibration, hence smooth. Corollary 2 proves that this square is homotopy cocartesian without an extra excision hypothesis. If the lower inclusion $U\to X$ is in $W$, however, deducing that the upper inclusion $U/t\to X/t$ is in $W$ is exactly a converse-excision step. Under strong W(7 bis), Proposition 13 supplies it. The target $X/t$ has a final object, so $U/t$ is then $W$-aspherical. Conversely, if $U/t$ is $W$-aspherical, the inclusion $U\to X$ is $W$-aspherical: its slices at objects of $U$ have final objects, and its only remaining slice is $U/t$. Theorem A gives the weak equivalence.

The reduction in Corollary 5 is now precise. Because $t_0$ is closed in $X\setminus U$ and $U$ is a sieve, $t_0$ has no arrow to any other object of $X$. Hence $U'=X\setminus\{t_0\}$ is again a sieve. The hypotheses make $U\to U'$ aspherical. If $U\to X$ is in $W$, two-out-of-three makes $U'\to X$ a weak equivalence; Corollary 4 makes it aspherical. Composition then makes $U\to X$ aspherical, which supplies the remaining slice condition at $t_0$. The argument is complete under its stated strong hypothesis. It should not be read as a general converse to Theorem A for ordinary weak homotopy equivalences.
\end{dcomment}


\noindent\textsc{Démonstration de th. 4.}

\noindent\circled{a} Critère de lissité de $\varphi_Q : \mathcal{X} = \intdiag{X_0}{X_1}{X_2} \longrightarrow Y$.

Comme $\{1\}$, $\{2\}$ sont des ouverts de $\Phi = \begin{tikzcd}[row sep=small,column sep=small,baseline=(\tikzcdmatrixname-2-1.base)]
 & 1 \arrow[dl] \\
 0 & \\
 & 2 \arrow[ul]
\end{tikzcd}$, $i_1 : X_1\longrightarrow\mathcal{X}$ et $i_2 : X_2\longrightarrow\mathcal{X}$ sont des immersions ouvertes. Comme $\varphi_Q\circ i_1 = f_1$,
\msp{page 152}
$\varphi_Q\circ i_2=f_2$, on voit déjà que
\[
f_1 \text{ et } f_2 \text{ lisses} \iff \varphi_Q \text{ lisse sur l'ouvert réunion de } i_1(X_1) \text{ et } i_2(X_2)\text{, i.e. sur } \mathcal{X}\backslash i_0(X_0).
\]
Il reste à exprimer la lissité de $\varphi_Q$ en les points $x_0$ de $i_0(X_0)$. Soient $x_1$, $x_2$, $y$ les images de $x_0$ dans $X_1$, $X_2$, $Y$ respectivement. Il faut exprimer que les fibres de
\[
(\varphi_Q)/x_0 : \mathcal{X}/x_0\longrightarrow Y/y \tag*{$(*)$}
\]
sont asphériques. Nous allons pour cela expliciter $\mathcal{X}/x_0$ et la flèche $(*)$ en termes du carré commutatif
\[
(**)\qquad\qquad
\begin{tikzcd}[row sep=large,column sep=large]
X_0/x_0 \arrow[r, "\alpha_1/x_0"] \arrow[d, "\alpha_2/x_0"'] & X_1/x_1 \arrow[d, "f_1/x_1"] \\
X_2/x_2 \arrow[r, "f_2/x_2"] & Y/y,
\end{tikzcd}
\]
induit par $Q$. Je dis

\noindent\textbf{Lemme 1.} \textit{Soit un carré commutatif}
\[
Q \quad = \quad
\begin{tikzcd}[row sep=large,column sep=large]
X_0 \arrow[r, "\alpha_1"] \arrow[d, "\alpha_2"'] & X_1 \arrow[d, "f_1"] \\
X_2 \arrow[r, "f_2"] & Y
\end{tikzcd}
\]
\textit{dans} Cat, \textit{$x_0\in\mathrm{Ob}\,X_0$, $x_1$, $x_2$, $y$ ses images dans $X_1$, $X_2$, $Y$ respectivement, $Q/x_0$ le carré $(**)$ ci-dessus, induit par $Q$ sur les localisées $X_0/x_0$ etc., $\mathcal{X} = \intdiag{X_0}{X_1}{X_2}$. On a alors un}
% foot: Version 7 décembre 2002
\printedpage{105}
% head: A. Grothendieck : Dérivateurs, Ch. XII, Édition des Universités de Montpellier II et Paris VII
\msp{page 153}
\textit{isomorphisme canonique}
\[
\mathcal{X}/x_0 \quad\simeq\quad \intdiag{X_0/x_0}{X_1/x_1}{X_2/x_2} ,
\]
\textit{et la flèche $(\varphi_Q)/x_0$, localisée de $\varphi_Q : \mathcal{X}\longrightarrow Y$ en $x_0$, s'identifie à la flèche $\varphi_{(Q/x_0)}$ :}
\[
\varphi_{(Q/x_0)} : \quad \intdiag{X_0/x_0}{X_1/x_1}{X_2/x_2} \qquad\longrightarrow\qquad Y/y.
\]

\noindent\textsc{Démonstration} du Lemme. On sait que
\[
\mathcal{X}/x_0\longrightarrow\mathcal{X}
\]
est Cat-fibrée (scindée) (à fibres discrètes), et comme
\[
\mathcal{X}\longrightarrow\Phi
\]
est aussi Cat-fibrée (scindée), il s'ensuit qu'il en est de même de la flèche composée
\[
\mathcal{X}/x_0\longrightarrow\Phi.
\]
Donc pour `calculer' $\mathcal{X}/x_0$, en tant que catégorie sur $\Phi$, il faut expliciter ses trois fibres, et les deux foncteurs de changement de base qui les relient. On trouve aussitôt que ces fibres sont canoniquement isomorphes à
\[
X_0/x_0, \qquad X_1/x_1, \qquad X_2/x_2,
\]
et on admettra (!) que les foncteurs changement de base
\msp{page 154}
\[
X_0/x_0\longrightarrow X_1/x_1, \qquad X_0/x_0\longrightarrow X_2/x_2
\]
(définis à isomorphisme unique près) ne sont autres que les localisés $\alpha_1/x_0$ et $\alpha_2/x_0$ de $\alpha_1 : X_0\longrightarrow X_1$ et $\alpha_2 : X_0\longrightarrow X_2$. Cela montre donc qu'on a un isomorphisme
\[
\mathcal{X}/x_0 \quad\simeq\quad \intdiag{X_0/x_0}{X_1/x_1}{X_2/x_2} ,
\]
% foot: Version 7 décembre 2002
\printedpage{106}
% head: A. Grothendieck : Dérivateurs, Ch. XII, Édition des Universités de Montpellier II et Paris VII
et l'on admettra aussi que la flèche $(\varphi_Q)/x_0 : \mathcal{X}/x_0\longrightarrow Y/y$ s'identifie à $\varphi_{(Q/x_0)}$.

Pour exprimer la $W$-lissité de $\varphi_Q$ en $x_0\in\mathrm{Ob}\,X_0\subseteq\mathrm{Ob}\,\mathcal{X}$, on est donc ramené à exprimer la $W$-asphéricité des fibres de
\[
\varphi_{(Q/x_0)} : \quad \intdiag{X_0/x_0}{X_1/x_1}{X_2/x_2} \qquad\longrightarrow\qquad Y/y.
\]
Pour expliciter ces fibres, on a besoin d'un autre lemme auxiliaire

\dgrblock{dgrb473}{The slice of the integral, with its arrows}
\begin{dcomment}{exposition}{p154:Pour expliciter ces fibres, on a besoin d'un autre lemme auxiliaire}{The slice of the integral, with its arrows}
Grothendieck is reducing smoothness to a calculation at one object. The reduction succeeds because the smooth integral remembers the arrows into the central stratum, rather than merely the three categories as disjoint pieces. We make this calculation explicit; it supplies both identifications which the proof asks us to admit.

Write $E$ for the smooth integral of $X_0\xrightarrow{\alpha_i}X_i$, $i=1,2$. Its objects are the tagged objects of the three categories. Arrows within a stratum are the original arrows, and the additional hom-sets are
\[
 \Hom_E((i,b),(0,a))=\Hom_{X_i}(b,\alpha_i a),\qquad i=1,2.
\]
There are no arrows in the reverse direction or between the two outer strata. Composition is composition in $X_i$, using $\alpha_i$ when a central arrow is involved. Thus the cartesian arrow over $i\to0$ has the direction
\[
 (i,\alpha_i a)\longrightarrow(0,a).
\]
The arrows printed in the opposite direction on manuscript page 155 are an edition slip, not a different convention for this integral. The functor $\varphi_Q:E\to Y$ sends an additional arrow represented by $h:b\to\alpha_i a$ to $f_i(h)$, whose target is $f_i\alpha_i(a)=f_0(a)$. In particular it sends the cartesian arrow just displayed to an identity. This is the property which the source notices it had left implicit.

Fix $x_0\in X_0$ and put $x_i=\alpha_i x_0$. An object of $E/x_0$ in the central stratum is precisely an arrow $g:a\to x_0$ in $X_0$. An object in stratum $i$ is an arrow $h:b\to x_i$ in $X_i$. A cross-stratum morphism from $h$ to $g$ is an arrow $k:b\to\alpha_i a$ satisfying
\[
       (\alpha_i g)k=h.
\]
This is exactly a morphism from $h$ to the image of $g$ in $X_i/x_i$. Within each stratum we have the ordinary slice morphisms. These descriptions identify objects, hom-sets, identities, and compositions, giving the asserted canonical isomorphism
\[
 E/x_0\cong\int_{\mathrm{li}}
       (X_0/x_0\longrightarrow X_1/x_1,\ X_0/x_0\longrightarrow X_2/x_2).
\]
Under it the functor to $Y/y$, $y=f_0x_0$, sends $g$ to $f_0g$ and $h$ to $f_i h$; on a cross-stratum arrow it is $f_i k$. This proves the compatibility with $\varphi_{Q/x_0}$ as well as the isomorphism of the source categories.

The same description proves Lemma 2 without a further coherence argument. The strict fibre over $y$ consists of objects mapping to $y$ and arrows mapping to $1_y$. A cross-stratum arrow belongs to this fibre exactly when its representing arrow in $X_i$ does. Consequently
\[
 E_y\cong\int_{\mathrm{li}}(X_{0y}\longrightarrow X_{1y},\ X_{0y}\longrightarrow X_{2y}).
\]
Strict fibres need not be full subcategories. For instance, the strict fibre of the identity functor of a one-object group category over its object is the terminal category, whereas the full subcategory on that object is the whole group category. Keeping the identity condition on arrows is essential in the present calculation.

For an arrow $u:y'\to y$, set
\[
 A_i(u)=(X_i/x_i)_u.
\]
Its objects are arrows $a\to x_i$ whose image is exactly $u$; its morphisms are commuting triangles whose remaining edge maps to $1_{y'}$. The preceding two identifications now give, with no choices,
\[
       (E/x_0)_u\cong\int_{\mathrm{li}}
          (A_0(u)\longrightarrow A_1(u),\ A_0(u)\longrightarrow A_2(u)).
\]
This is the precise local input to the smoothness criterion. For the later systematic treatment, the locality of smoothness under passage to slices is Proposition 4.10 of \citet{Maltsiniotis2005Asphericite}; the argument here retains the additional explicit description needed for this particular integral.
\end{dcomment}


\noindent\textbf{Lemme 2} (\footnote[106]{NB C'est le b) du th. 4.}). \textit{Soit un carré commutatif dans} Cat
\[
Q \quad = \quad
\left(
\begin{tikzcd}[row sep=large,column sep=large]
X_0 \arrow[r, "\alpha_1"] \arrow[d, "\alpha_2"'] & X_1 \arrow[d, "f_1"] \\
X_2 \arrow[r, "f_2"] & Y
\end{tikzcd}
\right),
\]
\textit{d'où}
\[
\varphi_Q : \quad \left( \intdiag{X_0}{X_1}{X_2} \right) \qquad\longrightarrow\qquad Y,
\]

\msp{page 155}
\textit{et soit $y\in\mathrm{Ob}\,Y$. Alors la fibre de $\varphi_Q$ en $y$}
\[
\mathcal{X}_y = \varphi_Q^{-1}(\ \underbrace{\{y\}}_{\text{catégorie ponctuelle}}\ )
\]
\textit{est donnée par (comme catégorie au dessus de $\Phi$)}
\[
\mathcal{X}_y \quad\simeq\quad \intdiag{X_{0\,y}}{X_{1\,y}}{X_{2\,y}} .
\]

\noindent\textsc{Démonstration.} Il est immédiat qu'on a
\[
\mathcal{X}_y\cap\underbrace{\mathcal{X}_s}_{X_s}\simeq(X_s)_y \qquad\quad \text{pour } s\in\mathrm{Ob}\,\Phi=\{0,1,2\},
\]
% foot: Version 7 décembre 2002
\printedpage{107}
% head: A. Grothendieck : Dérivateurs, Ch. XII, Édition des Universités de Montpellier II et Paris VII
et (le hasard fait bien les choses) les foncteurs changement de base
\[
X_0\longrightarrow X_1, \qquad X_0\longrightarrow X_2
\]
envoient la sous-catégorie (non pleine) $X_{0\,y}$ de $X_0$ dans les sous-catégories (non pleines) $X_{1\,y}$ et $X_{2\,y}$ de $X_1$ et $X_2$ respectivement.
\[
\begin{tikzcd}[row sep=large,column sep=large]
X_{s\,y} \arrow[r] \arrow[d] \arrow[dr, phantom, "\scriptstyle\text{cart.}"] & \mathcal{X}_y \arrow[r] \arrow[d] \arrow[dr, phantom, "\scriptstyle\text{cart.}"] & \{y\} \arrow[d] \\
X_s \arrow[r] \arrow[d] \arrow[dr, phantom, "\scriptstyle\text{cart.}"] & \mathcal{X} \arrow[r] \arrow[d] & Y \\
\{s\} \arrow[r] & \Phi &
\end{tikzcd}
\]
On a donc une sous-catégorie fibrée de $\mathcal{X}$, sous réserve de vérifier que les flèches strictement cartésiennes
\[
x_0\longrightarrow\alpha_1(x_0), \qquad x_0\longrightarrow\alpha_2(x_0)
\]
dans $\mathcal{X}$, pour $x_0$ dans $X_{0\,y}$, sont dans $\mathcal{X}_y$, i.e. sont au-dessus de la flèche identique de $y$. Or c'est là la propriété essentielle
\msp{page 156}
du foncteur $\varphi_Q : \mathcal{X}\longrightarrow Y$, qu'on a oublié d'expliciter dans 3.2. Ça a vraiment l'air de marcher (comme sur des roulettes!)

Appliquons le Lemme 2 au diagramme carré $Q/x_0$ et à $\varphi_{(Q/x_0)}$ (page 154), on voit donc que les fibres à examiner en les $u\in\mathrm{Ob}\,Y/y$ ($u : y'\longrightarrow y$) peuvent s'écrire comme les intégrales
\[
\mathcal{X}(x_0,u) \quad\overset{\text{déf}}{=}\quad \intdiag{(X_0/x_0)_u}{(X_1/x_1)_u}{(X_2/x_2)_u} .
\]
Utilisons maintenant le fait que $f_1$, $f_2$ sont $W$-lisses (exprimant que $\varphi_Q$ est $W$-lisse dans $\mathcal{X}\backslash X_0$), donc que $(X_1/x_1)_u$ et $(X_2/x_2)_u$ sont $W$-asphériques. Dans le carré $W$-cocartésien
\[
\begin{tikzcd}[row sep=large,column sep=small]
 & (X_1/x_1)_u \arrow[dr, "{i_1(x_0,u)}"] & \\
(X_0/x_0)_u \arrow[ur, "{\alpha_1(x_0,u)}"] \arrow[dr, "{\alpha_2(x_0,u)}"'] & & (\mathcal{X}/x_0)_u\simeq\mathcal{X}(x_0,u) \\
 & (X_2/x_2)_u \arrow[ur, "{i_2(x_0,u)}"'] &
\end{tikzcd}
\]
% foot: Version 7 décembre 2002

% ---- from Der-ChXIIb-012_p066-070.tex
\printedpage{108}
% head: A. Grothendieck : Dérivateurs, Ch. XII, Édition des Universités de Montpellier II et Paris VII
(\footnote[107]{attention, ce n'est pas un carré commutatif. Dire plutôt que $(\mathcal{X}/x_0)_u \simeq \int \begin{tikzcd}[ampersand replacement=\&, row sep=0.4em, column sep=0.4em, baseline=(\tikzcdmatrixname-2-1.base)] \& {} \\ {} \ar[ur]\ar[dr] \& \\ \& {} \end{tikzcd}$ est lisse sur $\Phi$ à fibres $W$-asphériques (si $f_0$, $f_1$, $f_2$ sont lisses), donc $W$-asphérique puisque $\Phi$ l'est.}), il faut exprimer que $(\mathcal{X}/x_0)_u$ est $W$-asphérique, sachant que $(X_1/x_1)_u$ et $(X_2/x_2)_u$ le sont. Donc cela équivaut à % Resolved from PDF leaf 31: the original diagram has exactly these two unlabelled arrows and no vertex labels.
\[
i_1(x_0,u) \in W \qquad (\text{ou } i_2(x_0,u) \in W),
\]
donc aussi (moyennant W(7 bis)) à
\msp{page 157}
\[
\alpha_1(x_0,u) \in W, \qquad \text{i.e. } (X_0/x_0)_u \ W\text{-asphérique.}
\]
Dire que c'est vrai pour tout $u$ dans $Y/y$, signifie que $f_0 : X_0 \longrightarrow Y$ est lisse en $u$ \texttt{[plutôt en $x_0$]}.

Donc on a prouvé a), sous la forme plus précise :

\dgrblock{dgrb474}{Repairing the noncommuting fibre square}
\begin{dcomment}{exposition}{p157:Donc on a prouvé a), sous la forme plus précise :}{Repairing the noncommuting fibre square}
The aim of the calculation on page 156 is to reconstruct smoothness at the central object from smoothness on the three pieces. The displayed diamond cannot itself be used as a commutative square: its two composites land in different outer strata of the integral. They are related to the inclusion of the central stratum by cartesian natural transformations. Footnote 107 correctly withdraws that use of the square.

Here is the direct argument, with its hypotheses visible. Suppose $f_i$ is $W$-smooth at $x_i$ for $i=0,1,2$. For every $u:y'\to y$, all three categories $A_i(u)$ from the preceding block are $W$-aspherical. Their smooth integral $L_u$ is a split fibration over $\Phi=(1\to0\leftarrow2)$, hence is $W$-smooth. For each $s\in\Phi$, smoothness identifies its strict fibre, up to a $W$-equivalence, with the comma category $s\backslash L_u$. These comma categories are therefore $W$-aspherical. The dual form of the relative Theorem A gives $L_u\to\Phi$ in $W$. Since $0$ is final in $\Phi$, $L_u$ is $W$-aspherical. The slice calculation identifies $L_u$ with $(E/x_0)_u$, proving smoothness of $\varphi_Q$ at $x_0$. No pushout reflection axiom has entered.

For the converse, suppose the two outer maps are smooth at $x_1,x_2$, and retain the \emph{strong} form of W(7 bis) stipulated in the source. Then $A_1(u)$ and $A_2(u)$ are $W$-aspherical. Proposition 13, applied to their integral and its closed mapping-cylinder replacement, gives
\[
 A_1(u)\longrightarrow L_u\text{ in }W
 \quad\Longleftrightarrow\quad
 A_0(u)\longrightarrow A_2(u)\text{ in }W.
\]
The left side is equivalent to asphericity of $L_u$, and the right side to asphericity of $A_0(u)$, by two-out-of-three. This proves the asserted equivalence. It applies Proposition 13 to the correctly constructed integral; it does not promote the original noncommuting diamond to a strict square. The reverse implication is a reflection statement, stronger than the preservation statement used in the direct proof. The counterexample to strong reflection discussed earlier in the commentary on this chapter explains why it cannot be transferred without its hypothesis to ordinary weak homotopy equivalences.

At an outer object $x_i$, every arrow with target $x_i$ already belongs to $X_i$, since that stratum is open. Hence $E/x_i=X_i/x_i$, and local smoothness is exactly that of $f_i$. At a central object the hypotheses on both outer objects remain necessary for the converse equivalence just established. The final sentence after Corollary 6 drops them, as the editorial annotation observes; the following finite example measures precisely what is lost.
\end{dcomment}


\textbf{Corollaire 6.} \textit{Soit $x \in \Ob \mathcal{X}$ (où $\mathcal{X} = \int \begin{tikzcd}[ampersand replacement=\&, row sep=tiny, column sep=small, baseline=(\tikzcdmatrixname-2-1.base)] \& X_1 \\ X_0 \ar[ur]\ar[dr] \& \\ \& X_2 \end{tikzcd}$.) On veut exprimer la $W$-lissité de}
\[
\varphi_Q : \mathcal{X} \longrightarrow Y
\]
\textit{en $x$.}

\begin{enumerate}[label={}, leftmargin=2.5em]
\item[\textit{a)}] \textit{Si $x \in \Ob X_i$ ($i \in \{1,2\}$), $\varphi_Q$ est $W$-lisse en $x$ si et seulement si $f_i : X_i \longrightarrow Y$ l'est.}
\item[\textit{b)}] \textit{Supposons} \textup{W(7 bis)} (\footnote[108]{Inutile pour la partie directe, en supposant \texttt{[phrase incomplète]}.})\textit{. Soit $x_0 \in \Ob X_0$, et supposons $f_1$, $f_2$ lisses en $x_1$, $x_2$. Alors $\varphi_Q$ est $W$-lisse en $x_0$ si et seulement si $f_0 : X_0 \longrightarrow Y$ l'est.}
\end{enumerate}

\textit{Donc} (moyennant W(7 bis)) \textit{dans les trois cas $x \in \Ob X_i$ ($i \in \{0,1,2\}$) on trouve que $\varphi_Q$ est $W$-lisse en $x$ si et seulement si $f_i : X_i \longrightarrow Y$ l'est} \texttt{[pas vraiment pour $i=0$]}.

\cercle{b} Quant à b) du théorème 4, c'est le lemme 2 ci-dessus. Donc tout est prouvé!

\msp{page 158}

\dgrblock{dgrb475}{A finite witness for the missing outer hypothesis}
\begin{dcomment}{development}{p157}{A finite witness for the missing outer hypothesis}
Grothendieck has isolated a useful local criterion, but the final formulation would make the middle map alone control the whole neighbourhood. We give a small witness showing why the two outer hypotheses must stay in the statement. We did not find this particular finite witness in the smoothness discussion of \citet{Maltsiniotis2005Asphericite} or the integration discussion of \citet[\S3.3]{Cisinski2006}; it is an elementary application of their criteria, not a claim of priority.

\textbf{Proposition (failure of the unconditional central criterion).}\ours
Let $W$ be a fundamental localizer satisfying W(5), so that every $W$-equivalence induces a bijection on connected components. There is a commutative square of finite posets for which $f_0$ is smooth and $f_2$ is smooth, but the smooth integral map $\varphi_Q$ is not $W$-smooth at an object of the central stratum.

\textit{Proof.} Let $Y=X_0=X_2=[1]=(0\to1)$, with $f_0=f_2=\alpha_2=1$. Let $X_1$ be the poset with objects $a,b,c$ and only nonidentity relations $a<c$ and $b<c$. Define
\[
 f_1(a)=f_1(b)=0,\qquad f_1(c)=1,
 \qquad \alpha_1(0)=a,\quad\alpha_1(1)=c.
\]
Both nonidentity arrows of $X_1$ map to $u:0\to1$. These formulas define functors and give $f_1\alpha_1=1_Y$, so the required square commutes. Identity functors are smooth; this proves the assertions about $f_0$ and $f_2$.

Take $x_0=1$, $x_1=c$, $x_2=1$ and the arrow $u:0\to1$ in $Y$. The three strict lift categories are
\[
 A_0(u)=e,\qquad A_1(u)=\{a\to c,\ b\to c\}_{\mathrm{disc}},
 \qquad A_2(u)=e.
\]
The first map chooses $a\to c$, and the second is the unique map between terminal categories. Their smooth integral has four objects: a central object $z$, outer objects $a_*,b_*$ in stratum 1, and $d_*$ in stratum 2. Its only nonidentity arrows are
\[
                 a_*\longrightarrow z\longleftarrow d_*.
\]
The object $b_*$ is isolated. Thus its connected-component set has two elements. By W(5) it is not $W$-aspherical. The slice identification gives this category as $(E/x_0)_u$, which violates the defining local smoothness test.\hfill$\square$

The failure occurs in a missing lift relation at the first outer stratum. It involves no question about strictification and no exotic homotopy type. In Corollary 6(b), the assumption that $f_1$ is smooth at $x_1$ excludes exactly this additional component.
\end{dcomment}


\subsection*{3.5\quad Applications aux sommes amalgamées dans une catégorie $X^{\wedge}$ (avec $X \in \Ob \Cat$).}

D'abord une caractérisation des sommes amalgamées impliquant une immersion ouverte :

\textbf{Proposition 15.} \textit{Soit}
% foot: Version 7 décembre 2002
\printedpage{109}
% head: A. Grothendieck : Dérivateurs, Ch. XII, Édition des Universités de Montpellier II et Paris VII
\[
Q : \qquad
\begin{tikzcd}
U \ar[r, "i"] \ar[d, "p"'] & X \ar[d, "q"] \\
U' \ar[r, "i'"] & X'
\end{tikzcd}
\]
\textit{un carré commutatif dans} $\Cat$\textit{. Pour que ce carré soit cocartésien dans} $\Cat$\textit{, avec $i$ (donc $i'$) une immersion ouverte, il faut et il suffit que les conditions suivantes soient satisfaites :}

\begin{enumerate}[label={}, leftmargin=2.5em]
\item[\textit{a)}] \textit{$Q$ est cartésien.}
\item[\textit{b)}] \textit{$i'$ est une immersion ouverte (donc $i$ aussi).}
\item[\textit{c)}] \textit{Si $T' = X' \backslash U'$, et $T = X \times_{X'} T' = X \backslash U$, $T \longrightarrow T'$ est un isomorphisme.}
\item[\textit{d)}] \textit{Pour tout $t \in T$, si $t'$ est son image dans $T'$, le foncteur induit par $p$}
\[
p/t : \underbrace{U/t}_{= V_t} \longrightarrow \underbrace{U'/t'}_{= V'_{t'}}
\]
\textit{satisfait à $(p/t)_!(e_{U/t}) \simeq e_{U'/t'}$, ou ce qui revient au même, les `fibres colocales' de $p/t$, $v' \backslash V_t$, avec $v'$ dans $V'_{t'}$, sont} $0$\textit{-connexes.}
\end{enumerate}

\msp{page 159}
Démonstration immédiate - ça devrait plutôt remonter au §1.

\dgrblock{dgrb476}{Gluing the morphisms across an open immersion}
\begin{dcomment}{exposition}{p158}{Gluing the morphisms across an open immersion}
Grothendieck now changes scale: the homotopical assertion will rest on an ordinary categorical pushout. An open immersion $U\hookrightarrow X$ is a full sieve. Its complementary full subcategory $T$ is a cosieve; there are no arrows from $T$ to $U$. Thus $X$ is specified by $U$, $T$, and the bifunctor of crossing arrows
\[
       H:U^{\mathrm{op}}\times T\longrightarrow\mathrm{Set},
       \qquad H(u,t)=\Hom_X(u,t).
\]
Precomposition and postcomposition encode every composition involving a crossing arrow. This description explains the four conditions in Proposition 15: the first three identify the two pieces, while the fourth identifies the arrows joining them.

Assume first that $Q$ is cartesian, that $U'\hookrightarrow X'$ is open, and that the induced map $T\to T'$ is an isomorphism. Identify $T$ with $T'$ by this specified map. For $t\in T$, let $H_t$ and $H'_t$ be the presheaves $\Hom_X(-,t)$ and $\Hom_{X'}(-,t)$ on $U$ and $U'$. The square supplies a natural comparison
\[
                  p_!H_t\longrightarrow H'_t.
\]
The universal category obtained by pushing $U\hookrightarrow X$ along $p$ has pieces $U'$ and $T$ and crossing-arrow presheaf $p_!H_t$. Indeed, a functor from this category to $Z$ is a pair of functors on $U'$ and $T$, together with compatible images of crossing arrows; by the defining adjunction of $p_!$, these data are exactly a functor $X\to Z$ and a functor $U'\to Z$ agreeing on $U$. Composition with arrows of $T$ is preserved because $H_t$ varies functorially with $t$. This proves the universal property, not merely an objectwise description. Consequently $Q$ is a pushout precisely when all the displayed comparisons are isomorphisms.

We can test these comparisons by explicit factorization categories. An element of $(p_!H_t)(u')$ is represented by a pair
\[
       b:u'\longrightarrow p(u),\qquad a:u\longrightarrow t,
\]
modulo the relation generated by morphisms $h:u\to v$ with $a=a_vh$ and $p(h)b=b_v$. Its image is the composite $q(a)b:u'\to t$ in $X'$. Fix such a composite $a':u'\to t$. The category of its representatives and these morphisms is exactly
\[
        (u',a')\backslash(U/t),
\]
where $U/t\to U'/t$ is induced by $p$. Thus the fibre of the comparison over $a'$ is the set of connected components of this category. The comparison is bijective for every $u'$ if and only if every such category is nonempty and connected. This proves condition (d), including the equivalence with $(p/t)_!e=e$.

Conversely, the explicit universal category just constructed has $U'$ as a full sieve, has complement exactly $T$, and its inverse image of $U'$ is exactly $U$. Hence a pushout along $U\hookrightarrow X$ necessarily satisfies (a)--(c); its crossing-arrow comparison is an isomorphism, so it also satisfies (d). This completes both directions of Proposition 15. The argument is entirely about ordinary categories and ordinary left Kan extension.

There is a small example showing why the last condition cannot be omitted. Take $U=U'=e$, $X$ the discrete category on $u,t$, and $X'=(u\to t)$, with $q$ the identity on objects and $p=1$. The square is cartesian; both open subcategories are $\{u\}$, and the complements are the same terminal category. But $U/t$ is empty whereas $U'/t=e$. The missing crossing arrow makes the factorization category empty, so (d) fails. The pushout along the identity $p$ is $X$, not $X'$. The condition about connected factorization categories is exactly the information which the object sets and the two strata do not contain.
\end{dcomment}


\textbf{Corollaire 1.} \textit{Soit $Q$ un carré} cartésien (cf. page 158), \textit{avec $i'$ une} immersion ouverte \textit{(donc $i$ aussi), et $q$} $W$-lisse \textit{(donc $p$ aussi). On suppose que $W$ satisfait l'axiome} W(5) \textit{de connexité} (\footnote[109]{Donc il suffit $q$ $W_0$-lisse, où $W_0 = \{f \in \Fl \Cat \mid \pi_0(f) \text{ bijectif}\}$.})\textit{. Soit $T' = X' \backslash U'$, $T = X \backslash U$, supposons que $T' \longrightarrow T$ soit un isomorphisme. Alors $Q$ est cocartésien.}

Il suffit de prouver la condition d) de prop. 15. Considérons le carré
\[
\begin{tikzcd}
U/t \ar[r, hook, "i/t"] \ar[d, "p/t"'] & X/t \ar[d, "q/t"] \\
U'/t' \ar[r, hook, "i'/t'"] & X'/t',
\end{tikzcd}
\tag*{$(*)$}
\]
je dis qu'il est \textit{cartésien.} Cela résulte de la seule hypothèse que $Q$ soit cartésien, sans hypothèse de lissité sur $q$ ni que $i'$ soit une immersion ouverte, et se voit p. ex. en regardant le diagramme cubique :
% foot: Version 7 décembre 2002
\printedpage{110}
% head: A. Grothendieck : Dérivateurs, Ch. XII, Édition des Universités de Montpellier II et Paris VII
\[
\begin{tikzcd}[row sep=small, column sep=small]
& U/t \ar[rr] \ar[dl] \ar[dd] & & X/t \ar[dl] \ar[dd] \\
U'/t' \ar[rr, crossing over] \ar[dd] & & X'/t' & \\
& U \ar[rr] \ar[dl] & & X \ar[dl] \\
U' \ar[rr] & & X' \ar[from=uu, crossing over] &
\end{tikzcd}
\]
(\footnote[110]{NB Ici $t$ est un élément quelconque de $\Ob X$, $t'$ son image dans $X'$.}). Dans ce diagramme, la face antérieure et postérieure du cube est cartésienne, par définition de $U/t$ et de $U'/t'$
\msp{page 160}
respectivement, et la base l'est par hypothèse. Il en résulte que la face supérieure du cube est également un carré cartésien.

Ceci vu, commme $q$ est lisse, il en est de même de $q/t$ dans le carré cartésien $(*)$, donc aussi de $p/t$ par changement de base. D'autre part, la lissité de $q$ implique que $q/t$ est à fibres $W$-asphériques, donc par changement de base il en est de même de $p/t$ (qui est donc une $W$-équivalence universelle, i.e. une $W$-fibration triviale). De plus, à cause de la lissité de $p/t$, pour $v'$ dans $U'/t'$, l'inclusion % [sic] commme -> comme, repeated letter
\[
(U/t)_{v'} \longrightarrow v' \backslash (U/t)
\]
est dans $W$, donc (par W(5)) définit une bijection sur les $\pi_0$. Comme $(U/t)_{v'}$ est $W$-asphérique par ce qui précède, il s'ensuit (par W(5)) que son $\pi_0$ est de cardinal 1, donc il en est de même de $v' \backslash (U/t)$, i.e. celui-ci est 0-connexe. Or c'est là la condition d) de la proposition 15, qu'il fallait établir.

\msp{page 161}
\textbf{Corollaire 2.} \textit{Soit $Q$ un carré cartésien} (p. 158), \textit{avec $i'$ (donc aussi $i$) une immersion ouverte, et $q$ (donc aussi $p$) un} isomorphisme local. \textit{Soit $T = X \backslash U$, $T' = X' \backslash U'$ (donc $T \simeq X \times_{X'} T'$), et supposons $T' \longrightarrow T$ isomorphisme. Alors $Q$ est cocartésien.}

C'est un cas particulier du corollaire 1, mais sans référence explicite à un localiseur $W$. On peut pendre $W = W_{\infty}$, ou $W = W_0$, pour se ramener au corollaire 1. Mais il est plus raisonnable de dire que dans le carré cartésien $(*)$, l'hypothèse que $q$ soit un isomorphisme local implique que $q/t$ est un isomorphisme, donc $p/t$ l'est aussi. Il en résulte bien que les $v' \backslash (U/t)$ sont 0-connexes, cette catégorie étant isomorphe à $v' \backslash (U'/t')$, donc a un objet initial. % [sic] pendre -> prendre, missing letter

\dgrblock{dgrb477}{From smoothness to the connected factorization test}
\begin{dcomment}{exposition}{p161:C'est un cas particulier du corollaire 1, mais sans référence explicite à un localiseur $W$.}{From smoothness to the connected factorization test}
The two corollaries of Proposition 15 replace its explicit factorization condition by familiar geometric hypotheses. Their force is that a condition on the ambient map $q$ controls every crossing arrow simultaneously.

For a cartesian square $Q$ and an object $t\in X$, the square of slices displayed on page 159 is cartesian. An object of its pullback is an arrow $x\to t$ in $X$ together with a chosen lift of its image-source to $U'$; the cartesian property of $Q$ identifies that choice with an object of $U$ lying above $x$. The same description of morphisms gives the isomorphism with $U/t$. No smoothness is needed for this step.

If $q$ is $W$-smooth, then $q/t:X/t\to X'/t'$ is $W$-smooth and has $W$-aspherical strict fibres. The first assertion is locality of smoothness on slices; the second is its defining lift-category test at $t$. Both statements are inherited by the pullback $p/t$. For $v'\in U'/t'$, smoothness gives a $W$-equivalence
\[
                (U/t)_{v'}\longrightarrow v'\backslash(U/t).
\]
Its source is $W$-aspherical, so its target is too. W(5) converts this into the required statement that the target has exactly one connected component. Proposition 15 then proves that $Q$ is cocartesian. Only the connectedness part of smoothness is used: as footnote 109 says, $W_0$-smoothness is sufficient.

For a local isomorphism $q$, the entire slice map $q/t$ is an isomorphism, and so is $p/t$. In that case every factorization category $v'\backslash(U/t)$ has an initial object. This gives the second corollary without any localizer. The later formulation that local isomorphisms are smooth is \citet[Proposition 4.9]{Maltsiniotis2005Asphericite}, but the direct slice argument also explains the stronger ordinary conclusion needed here.

The arrows on complements in the statements of both corollaries are printed $T'\to T$. The map induced by the square is $T\to T'$, as in Proposition 15(c) and in the application on page 166. The condition is that this induced map be an isomorphism. Its inverse then exists, but no inverse map is part of the initial data. We retain the printed direction in the text and record the correction in the register.
\end{dcomment}


Soit maintenant $X$ un objet de Cat, et considérons un diagramme cocartésien dans $X^{\wedge}$ :
\[
\begin{tikzcd}
G \ar[r, "i"] \ar[d, "p"'] & F \ar[d, "q"] \\
G' \ar[r, "i'"] & F',
\end{tikzcd}
\tag*{$(*)$}
\]
% foot: Version 7 décembre 2002
\printedpage{111}
% head: A. Grothendieck : Dérivateurs, Ch. XII, Édition des Universités de Montpellier II et Paris VII
\msp{page 162}
où $i$ est un monomorphisme. On sait que le caractère cocartésien du carré se vérifie `argument par argument' par le caractère cocartésien de
\[
\begin{tikzcd}
G(x) \ar[r, "i(x)"] \ar[d, "p(x)"'] & F(x) \ar[d, "q(x)"] \\
G'(x) \ar[r, "i'(x)"] & F'(x),
\end{tikzcd}
\tag*{$(*)$}
\]
dans Ens, $\forall\, x$ dans $X$. Or un tel carré cocartésien dans Ens, avec $i(x)$ \texttt{[un]} monomorphisme, est aussi cartésien. Cela implique que le carré $Q$ dans $X^{\wedge}$ est également cartésien.

Considérons, d'autre part, le foncteur canonique
\[
\varphi_X : F \longmapsto X/F : X^{\wedge} \longrightarrow \mathrm{Cat}.
\]
Il vaut mieux, d'ailleurs, le regarder comme un foncteur
\[
\varphi_X : X^{\wedge} \longrightarrow \mathrm{Cat}/X,
\]
et je rappelle, pour mémoire, que ce foncteur est pleinement fidèle, et que son image essentielle est formée des catégories $X'$ sur $X$ telles que
\[
X' \longrightarrow X
\]
\msp{page 163}
soit un isomorphisme local, ou ce qui revient au même, que $X'$ soit \textit{fibrée} sur $X$ à fibres discrètes :
\[
\varphi_X : X^{\wedge} \xrightarrow{\ \sim\ } \underbrace{\Fibdisc\, X}_{\substack{\text{catégories fibrées}\\ \text{à fibres discrètes}}} \ .
\]
Notons que toute flèche $X' \longrightarrow X''$ dans $\Fibdisc\, X$ est un isomorphisme local (tautologique sur $\begin{tikzcd}[ampersand replacement=\&, row sep=small, column sep=tiny, baseline=(\tikzcdmatrixname-1-1.base)] X'/x' \ar[rr] \ar[dr, "\text{iso.}"'] \& \& X''/x'' \ar[dl, "\text{iso.}"] \\ \& X/x \& \end{tikzcd}$), ce qui montre que si on a un diagramme
\[
\begin{tikzcd}[row sep=small, column sep=small]
X_1 \ar[dr] & & X_2 \ar[dl] \\
& X_0 &
\end{tikzcd}
\]
dans $\Fibdisc\, X$, le produit fibré $X_3 = X_1 \times_{X_0} X_2$ dans Cat est tel que $X_3 \longrightarrow X_2$ est un isomorphisme local, donc aussi $X_3 \longrightarrow X$. C'est-à-dire $\Fibdisc(X)$ est stable par produits fibrés. Ceci implique que les produits fibrés existent dans $\Fibdisc\, X$, et que ce sont les produits fibrés dans Cat, ou dans Cat$/X$.
% foot: Version 7 décembre 2002

\printedpage{112}
% head: A. Grothendieck : Dérivateurs, Ch. XII, Édition des Universités de Montpellier II et Paris VII
\msp{page 164}
De même pour les produits finis dans $\Fibdisc\, X$, ce sont les produits finis dans Cat$/X$. Bref

\textbf{Lemme 1.} \textit{La sous-catégorie pleine} $\Fibdisc\, X$ \textit{de} $\Cat/X$ \textit{est stable par} $\varprojlim$ \textit{finies, donc les} $\varprojlim$ \textit{existent dans} $\Fibdisc\, X$ \textit{et le foncteur d'inclusion}
\[
\Fibdisc\, X \longrightarrow \Cat/X
\]
\textit{y commute, et en particulier transforme carrés cartésiens de} $\Fibdisc\, X$ \textit{en carrés cartésiens de} $\Cat/X$\textit{, ou ce qui revient au même, de} $\Cat$\textit{. De plus, en vertu de l'équivalence}
\[
\varphi_X : X^{\wedge} \xrightarrow{\ \sim\ } \Fibdisc\, X \hookrightarrow \Cat/X,
\]
\textit{le foncteur $\varphi_X$ de $X^{\wedge}$ dans} $\Cat/X$ \textit{est exacte à gauche, en particulier il transforme carrés cartésiens de $X^{\wedge}$ en carrés cartésiens de} $\Cat/X$\textit{, ou encore de} $\Cat$\textit{.} % [sic] exacte -> exact, agreement with "le foncteur"

\msp{page 165}
Notons aussi qu'étant exact à gauche, le foncteur $\varphi_X$ de $X^{\wedge}$ dans Cat$/X$ transforme monomorphismes en monomorphismes. D'autre part :

\dgrblock{dgrb478}{The exactness of the category of elements}
\begin{dcomment}{later}{p164}{The exactness of the category of elements}
Grothendieck's finite-limit lemma is the first part of a stronger known statement. The relative category-of-elements functor
\[
             j_X:\widehat X\longrightarrow\Cat/X,
             \qquad F\longmapsto (X/F\to X),
\]
preserves all small limits and colimits. This is Proposition 3.2.12 of \citet{Cisinski2006}, in the published version, page 141. The absolute functor $i_X:F\mapsto X/F$ preserves colimits and pullbacks, and has the right adjoint stated in Corollary 3.2.13, page 142. We spell out the mechanism because it will also explain the pushout in Theorem 5.

An object of $X/F$ is $(x,a)$ with $a\in F(x)$. An arrow $(x,a)\to(y,b)$ is an arrow $u:x\to y$ such that $F(u)b=a$. Thus an $n$-simplex of its nerve is uniquely determined by a chain in $X$ and the element at its last vertex:
\[
 N(X/F)_n\cong
       \coprod_{x_0\to x_1\to\cdots\to x_n} F(x_n).
\]
All earlier elements are recovered by restriction. The final face uses the last restriction map of $F$; the other face and degeneracy maps change the underlying chain. The displayed bijections therefore define the nerve functor, not just its degreewise cardinalities.

Colimits of presheaves are evaluated pointwise, and coproducts commute with colimits. It follows that $N(X/F)$ preserves every small colimit in $F$. A cocone of nerves whose colimit is again a nerve is also a colimit cocone in categories: apply the fully faithful nerve to maps into an arbitrary target category. This proves the colimit assertion for $i_X$, and the cocones retain their maps to $X$, proving it for $j_X$ as well. For limits over $X$, work over the fixed simplicial set $NX$. The fibre in degree $n$ above a specified chain is $F(x_n)$, so limits are again calculated pointwise. Since the nerve creates limits, this proves the limit assertion for $j_X$.

For completeness, the absolute right adjoint is
\[
       R_X(C)(x)=\operatorname{Fun}(X/x,C),
\]
where the right side is the set of functors. For $u:x\to y$, restriction is precomposition with $X/x\to X/y$. A functor $H:X/F\to C$ assigns to $a\in F(x)$ the composite
\[
 X/x\longrightarrow X/F\xrightarrow{H}C,
 \qquad (z\xrightarrow{v}x)\longmapsto(z,F(v)a).
\]
These assignments form a map $F\to R_X(C)$. Conversely, given such a map, evaluate its functor at $1_x$ to obtain the image of $(x,a)$. An arrow $u:(x,a)\to(y,b)$ is sent to the arrow in the assigned functor on $X/y$ from $u$ to $1_y$. Naturality gives the correct source and ensures preservation of identities and compositions. The two constructions are inverse. This proves the adjunction directly.

The relative target matters for products: $X/(F\times G)$ is the fibre product $(X/F)\times_X(X/G)$. It is generally not the ordinary product of the two total categories. The absolute functor sends the terminal presheaf to $X$, not to the terminal category. The manuscript correctly formulates finite exactness in $\Cat/X$.
\end{dcomment}


\textbf{Lemme 2.} \textit{Les monomorphismes qui sont des isomorphismes locaux sont exactement les immersions ouvertes. Donc dans} $\Fibdisc\, X$\textit{, les monomorphismes sont les immersions ouvertes, et le foncteur $\varphi_X : F \longmapsto X/F$ transforme monomorphismes en immersions ouvertes.}

Revenons alors au diagramme cocartésien $(*)$ (p. 161), et au diagramme correspondant dans Cat (et même dans $\Fibdisc\, X$)
\[
\begin{tikzcd}
X/G \ar[r, hook, "X/i"] \ar[d, "X/p"'] & X/F \ar[d, "X/q"] & T \ar[l, hook', dotted] \ar[d, dotted] \\
X/G' \ar[r, hook, "X/i'"] & X/F' & T'. \ar[l, hook', dotted]
\end{tikzcd}
\]
Comme $(*)$ est cartésien, le carré précédent est cartésien, par ce qui précède. Je veux prouver qu'il est également cocartésien
\msp{page 166}
(dans Cat$/X$, ou dans Cat, cela revient au même), et pour ceci j'applique le corollaire 2 de la proposition 15 (p. 161), compte tenu que $X/q$ est un isomorphisme local, et $X/i'$ (étant un isomorphisme local \texttt{[et un]} monomorphisme) une immersion ouverte. Il suffit de prouver que, $T'$ et $T$ étant les sous-catégories fermées complémentaires des ouverts $X/G'$ de $X/F'$ et de $X/G$ de $X/F$, que $T \longrightarrow T'$ est \texttt{[un]} isomorphisme.

{\small \texttt{[Maintenant, $F$ et $G$, et $F'$ et $G'$, mais pas $p$ et $q$, échangent leur rôles.]}} % [sic] leur rôles -> leurs rôles, agreement

La vérification est immédiate : $T$ est la catégorie des flèches $x \longrightarrow G$ dans $X^{\wedge}$ (de source $x$ dans $X$) qui ne se factorisent pas par $G$ \texttt{[plutôt par $F$]}, $T'$ celle des flèches $x \longrightarrow G'$
% foot: Version 7 décembre 2002

% ---- from Der-ChXIIb-012_p071-072.tex
\printedpage{113}
% head: A. Grothendieck : Dérivateurs, Ch. XII, Édition des Universités de Montpellier II et Paris VII
qui ne se factorisent pas par $F'$, le foncteur $T \longrightarrow T'$ étant le foncteur composition de $x \longrightarrow G$ avec $G \xrightarrow{q} G'$
\[
\begin{tikzcd}[row sep=large]
F \arrow[r, hook] \arrow[d] & G \arrow[d] \\
F' \arrow[r, hook] & G'
\end{tikzcd}
\]
(NB Le composé ne se factorise pas par $F'$, car autrement, à cause du caractère cartésien de $Q$, $x \longrightarrow G$ se factoriserait par $F$.) Mais pour $x$ donné, la donnée de $x \longrightarrow G$ ne se factorisant pas par $F$ est celle d’un élément \msp{page 167} de $G(x)\backslash F(x)$, et à cause du caractère cocartésien de $Q$,
\[
(G(x)\backslash F(x)) \longrightarrow (G'(x)\backslash F'(x))
\]
est bijectif, i.e. ‘il revient au même’ de se donner $x \longrightarrow G$ ne se factorisant pas par $F$, ou $x \longrightarrow G'$ ne se factorisant pas par $F'$.

\dgrblock{dgrb479}{Why the complement calculation is enough}
\begin{dcomment}{exposition}{p167:est bijectif, i.e. ‘il revient au même’ de se donner}{Why the complement calculation is enough}
Grothendieck wants to transport a pushout of presheaves to a pushout of categories, where the homotopical criterion is already available. We use throughout the notation of the final theorem, $F\hookrightarrow G$ and $F'\hookrightarrow G'$. The source changes the letters immediately before this calculation; the editorial note marks that change, and no mathematical reversal is intended.

At each $x\in X$, the pushout of sets has the concrete description
\[
 G'(x)\cong F'(x)\amalg\bigl(G(x)\setminus F(x)\bigr).
\]
The map $q_x$ sends an element of $F(x)$ through $p_x$ and leaves the complementary elements distinct. Hence $F'(x)\to G'(x)$ is injective and its inverse image in $G(x)$ is exactly $F(x)$. This proves simultaneously that $i'$ is a monomorphism and that the square of presheaves is cartesian. It also gives a bijection between the two objectwise complements.

A monomorphism $F\hookrightarrow G$ induces a full sieve $X/F\hookrightarrow X/G$: if $(y,b)$ belongs to $X/F$ and $u:(x,a)\to(y,b)$ is an arrow of $X/G$, then $a=G(u)b$ belongs to $F(x)$. Conversely, a monomorphism which is a local isomorphism is a full sieve. To see this, lift uniquely every arrow into an object in its image using the isomorphism on slices; injectivity on objects and arrows then identifies the lift with the given arrow. This is Lemma 2 with its categorical content made explicit.

It remains to check that the bijection of complementary objects respects all arrows in both directions. Let $(x,a)$ and $(y,b)$ be outside $X/F$. An arrow between their images in $X/G'$ has an underlying arrow $u:x\to y$ satisfying
\[
       G'(u)q_y(b)=q_x(a).
\]
By naturality its left side is $q_xG(u)b$. The element $G(u)b$ cannot belong to $F(x)$, since then its image would belong to $F'(x)$, whereas $q_x(a)$ does not. Both $G(u)b$ and $a$ are therefore in the complement where $q_x$ is injective. They are equal. This gives precisely one arrow in the original complement, with the same underlying $u$. Identities and compositions are preserved because both complement functors lie over $X$. The functor $T\to T'$ is thus an isomorphism of categories, not merely a bijection on objects. Proposition 15, Corollary 2, now proves the categorical pushout.

The stronger exactness result in the previous block gives an alternative: $i_X$ preserves this pushout even without the monomorphism hypothesis. The monomorphism is needed for the homotopical conclusion, and for the simultaneous pullback assertion. Indeed, the nerve formula sends the square to a simplicial pushout along a monomorphism. Such a pushout computes the homotopy pushout of the nerves. Together with the Grothendieck-construction comparison, this proves that its canonical smooth-integral map is in $W_\infty$, and hence in every fundamental localizer, by \citet[Corollary 4.2.19]{Cisinski2006}. Thus the W(7 bis) proviso in the homotopy-pushout clause of Theorem 5 can be removed. It must still be retained for the \emph{reflection} implications in the following corollary.
\end{dcomment}


Ainsi on a prouvé le

\textbf{Théorème 5.} \textit{Soit $X$ dans} Cat\textit{, et}
\[
Q : \qquad
\begin{tikzcd}[row sep=large, baseline=(current bounding box.center)]
F \arrow[r, hook, "i"] \arrow[d, "p"'] & G \arrow[d, "q"] \\
F' \arrow[r, hook, "i'"] & G'
\end{tikzcd}
\]
\textit{un carré cocartésien dans $X^\wedge$, avec $i$} \texttt{[un]} \textit{monomorphisme. Alors $i'$ est} \texttt{[un]} \textit{monomorphisme, le carré $Q$ est aussi cartésien, et le carré correspondant dans} Cat$/X$ \textit{ou dans} Cat
\[
\begin{tikzcd}[row sep=large]
X/F \arrow[r, hook, "X/i"] \arrow[d, "X/p"'] & X/G \arrow[d, "X/q"] \\
X/F' \arrow[r, hook, "X/i'"] & X/G'
\end{tikzcd}
\]
\msp{page 168} \textit{est à la fois cartésien et cocartésien. De plus, si $W$ est un localiseur fondamental dans} Cat\textit{, satisfaisant} W(7 bis) (en plus de W(1,2,3))\textit{, alors le carré $X/Q$ dans} Cat \textit{est $W$-cocartésien.}

\printedpage{114}
% head: A. Grothendieck : Dérivateurs, Ch. XII, Édition des Universités de Montpellier II et Paris VII
Ce dernier énoncé résulte de ce qui précédait et du corollaire 2 du théorème 4, compte tenu que $X/q$ est un isomorphisme local, donc est $W$-lisse.

\dgrblock{dgrb480}{An attachment over the walking arrow}
\begin{dcomment}{examples}{p168:Ce dernier énoncé résulte de ce qui précédait}{An attachment over the walking arrow}
Theorem 5 can be seen without any large category of simplices. Let $X=[1]=(0\to1)$. A presheaf on $X$ is a map of sets $F(1)\to F(0)$. Take $G(1)=\{r,s\}$, $G(0)=\{a,b,c\}$, with restrictions $r\mapsto a$, $s\mapsto b$. Let $F$ be the subpresheaf with $F(1)=\{r\}$ and $F(0)=\{a,b\}$, and let $F'=e$ be the terminal presheaf. The unique map $F\to e$ gives the pushout
\[
 G'(1)=\{*,s\},\qquad G'(0)=\{*,c\},
 \qquad *\mapsto*,\quad s\mapsto*.
\]

The category $X/G$ has five objects and two nonidentity arrows,
\[
 (0,a)\longrightarrow(1,r),\qquad
 (0,b)\longrightarrow(1,s),
\]
with $(0,c)$ isolated. Its open subcategory $X/F$ has three objects: the first arrow and the isolated object $(0,b)$. The category $X/e$ is $[1]$. In the pushout, $(1,r)$ is identified with $(1,*)$, and both $(0,a)$ and $(0,b)$ with $(0,*)$. The resulting category has exactly the two arrows
\[
 (1,*)\longleftarrow(0,*)\longrightarrow(1,s)
\]
and the isolated object $(0,c)$. This is $X/G'$. There are no further nonidentity compositions to add.

The two complementary closed subcategories are both the discrete category on $(1,s)$ and $(0,c)$. The new attachment changes the crossing arrow into $(1,s)$ but leaves the complement itself unchanged, exactly as Proposition 15 predicts. Their nerves give a homotopy pushout. On connected components the four corners have respectively $2,3,1,2$ components, and the last set is the pushout of the first three component sets. This also checks concretely that the word ``cartesian'' for the presheaf square does not mean that either vertical map is an equivalence: the left vertical map merges two components, and the right one merges the corresponding two components after attachment.
\end{dcomment}


\textbf{Corollaire :} (\footnote[111]{\textbf{NB} Si on veut seulement $i \in W_X \Longrightarrow i' \in W_X$, $p \in W_X \Longrightarrow q \in W_X$, il suffit que l’on ait W(7 bis) \textit{faible}, i.e. $i_0 \in W \Longrightarrow i \in W$ dans l’énoncé de W(7 bis), p. 76 (Loc(6$'$a) avec les notations de XVI).}) \textit{On a} (moyennant W(7 bis)) \textit{les équivalences}
\[
\begin{aligned}
i \in W_X \quad &\Longleftrightarrow \quad i' \in W_X \\
p \in W_X \quad &\Longleftrightarrow \quad q \in W_X,
\end{aligned}
\]
\textit{où}
\[
W_X \subseteq \mathit{Fl}(X^\wedge)
\]
\textit{est défini par (cf. Pursuing Stacks \ldots)}
\[
u \in W_X \qquad\qquad \Longleftrightarrow \qquad\qquad X/u : X/F \longrightarrow X/G \quad \textit{est dans } W.
\]

Ainsi, dans une catégorie $X^\wedge$, les sommes amalgamées impliquant un monomorphisme se comportent de façon parfaite du point de vue homotopique, plus exactement pour les localiseurs $W_X$, pourvu \msp{page 169} seulement que $W$ satisfasse à W(7 bis). Notamment, les monomorphismes dans $X^\wedge$ sont des ‘$W_X$-cofibrations’, en un sens similaire de celui de la définition 2 c) (pages 99, 100 \texttt{[plutôt 109, 110]}).

\texttt{[page 170 n’existe pas.]}


\dgrblock{dgrb481}{From attachment stability to model structures}
\begin{dcomment}{status}{p169}{From attachment stability to model structures}
The final assertion concerns a class of cofibrations in the sense developed in this chapter: pushing a $W_X$-equivalence out along a monomorphism again gives a $W_X$-equivalence. This preservation property survives for ordinary homotopy types. The two reverse implications displayed in the corollary are different assertions and depend on the strong W(7 bis) hypothesis. The earlier failure of pushout reflection applies here as well after representing spaces by simplicial presheaves.

It would go beyond the text to infer that, for every small $X$, the class $W_X$ is the weak-equivalence class of a model structure whose cofibrations are all monomorphisms. Already $X=e$ shows the missing condition. For $W=W_\infty$, a presheaf is a set and $X/F$ its discrete category, so $W_X$ consists of bijections. If a model structure on sets had these weak equivalences and all monomorphisms as cofibrations, every trivial fibration would be a bijection. Factoring a noninjective map as a cofibration followed by a trivial fibration would make it injective, a contradiction. Attachment stability alone does not supply the factorization axiom.

The later theory identifies the additional hypotheses. For an accessible fundamental localizer and a $W$-local test category, the induced presheaf equivalences admit a cofibrantly generated model structure with monomorphisms as cofibrations; this is \citet[Corollary 4.2.18]{Cisinski2006}. In the ordinary case, the local test condition is characterized by the requirement that maps with the right lifting property against all monomorphisms be weak equivalences; see \citet[Theorem 4.1.19]{Cisinski2006}. The general existence theorem for accessible presheaf localizers, Theorem 1.4.3 of the same work, supplies the factorizations. The present calculation establishes the attachment mechanism on which that construction depends, while the test-category condition ensures that it models the intended homotopy theory.
\end{dcomment}
\begin{otherlanguage}{english}
\dgroverview{XIIb-b}{Chapter XII \textperiodcentered\ \textup{XIIb-b} \textperiodcentered\ Overview}
\selectlanguage{english}
\noindent\textbf{Chapter XII, section 5, manuscript pages 115--169.}
\medskip

These pages develop a concrete calculus for deciding when an ordinary gluing diagram of categories represents a homotopy gluing. Its central object is a comparison map. Given a strictly commuting square with initial corner $X_0$, adjacent corners $X_1,X_2$, and terminal corner $Y$, Grothendieck constructs a category $\mathcal X_Q$ from the three-category span and a canonical functor
\[
                    \varphi_Q:\mathcal X_Q\longrightarrow Y.
\]
The square is $W$-homotopically cocartesian when this particular functor belongs to the chosen weak-equivalence class $W$. The definition remembers the cocone, not only its four vertex homotopy types. The work that follows explains why the construction is independent of its two orientations, how smoothness can test the comparison on fibres, and why monomorphisms of presheaves provide a reliable supply of such squares.

The opening argument concerns a different but related exactness problem. Total $W$-asphericity here is the requirement that, for every $p:X\to Y$, the category of elements of the ordinary left Kan extension $p_!1_X$ be $W$-aspherical. It is the definition of this chapter, explicitly distinguished by the editors from the earlier use of the same phrase in \emph{Pursuing Stacks}. $W_0$-properness, or $W$-properness under W(5), allows Grothendieck to compute
\[
 p_!1_X\cong\operatorname*{colim}_{i\in I}p_{i!}1_{X_i}
\]
for a family $X\to I$. A terminal object of $I$ reduces this colimit to the last fibre. The commentary supplies the generators-and-relations proof on components and the precise functor on comma categories requested in footnote~83. It also proves that total $W_\infty$-asphericity is not invariant under weak homotopy equivalence. The properness hypothesis controls ordinary Kan extensions; equivalence of classifying spaces alone does not do that work.

The smooth integral of $X_1\leftarrow X_0\to X_2$ keeps the three sets of objects disjoint. Its mixed arrows run from the outer fibres to the middle fibre, with
\[
 \operatorname{Hom}_{\mathcal X_Q}(a,x)
       =\operatorname{Hom}_{X_i}(a,\alpha_i x)
       \qquad(a\in X_i,\ x\in X_0).
\]
Its universal property therefore includes natural transformations $f_i\alpha_i\Rightarrow f_0$. They need not be invertible. Strict commutativity of the given square supplies identity transformations and hence the canonical $\varphi_Q$. Dividing the integral into two full closed pieces gives a genuine categorical pushout and pullback. Each piece retracts onto its outer fibre by an adjunction. The nerve of this closed cover proves directly that the total category represents the homotopy pushout of classifying spaces. The general theorem is Thomason's; the bibliography distinguishes verification of the original metadata from verification of an independent primary proof \citep{Thomason1979,Cegarra2020}.

Even the smallest examples explain why this construction is needed. For the span $e\leftarrow\{a,b\}\to e$, the integral is a four-object category whose nerve is a circle. Its ordinary categorical pushout is a point. The ordinary pushout of the same diagram in the homotopy category is also a point: it retains equality of components but loses the two paths implementing compatibility. The integral retains those paths as transition arrows. This is the concrete content, on these pages, of the insistence on forming diagram categories before localization.

The excision hypotheses require particular care. Proposition~13 and its cube corollary are valid conditional assertions under the full W(7 bis) and W(7) printed in the manuscript. Those axioms include reflection statements stronger than the usual preservation of weak equivalences by homotopy pushout. They fail for $W_\infty$. The commentary proves this with a connected acyclic finite complex having nontrivial fundamental group. Its suspension is contractible, so the union of its two cones and both cones are contractible although their intersection is not. A map of this closed-cover square to the point square violates the converse cube implication. The usual forward gluing theorem survives. A normalized-chain calculation also explains exactly why the complete three-way implication works for integral homology equivalences. The earlier chapter's claim about full W(7) for $W_\infty$ is therefore not a valid justification of the converses in the present portion.

The proper integral reverses the transition arrows. Grothendieck compares the two constructions by the functorial resolution
\[
                         C\longleftarrow\Sigma(C)\longrightarrow C^{\mathrm{op}}.
\]
Both projections are aspherical, and the same comparison can be integrated. Proposition~14 proves agreement of the four versions obtained from smooth or proper integration and from the span or its opposite. The commentary writes out the chosen twisted-arrow convention and completes the deferred transitivity verification. The natural comparison with the opposite operation is involutive; its naturality makes the square of four integration comparisons commute. The indexing category must keep its two paths distinct until they reach $Y$. Replacing that free diamond by a commutative square would erase precisely the information the proof is using.

The question on manuscript pages 145--146 asks whether a bad ordinary pushout can become a good homotopy cocone after mapping to another target. We answer it with finite posets. A collapsible disk with seven triangles maps to the tetrahedral two-sphere so that all associated comma categories are nonempty and connected. Its face poset $E$ embeds as a sieve in a contractible mapping-cylinder poset $B$. The categorical pushout
\[
                         P=B\amalg_E e
\]
has two incomparable minimal vertices below the face poset of that sphere. Its nerve is their join with the subdivided sphere, hence has classifying space $S^3$. The span $B\leftarrow E\to e$ has contractible integral. Consequently the comparison to $P$ is not a weak equivalence, whereas its composite with $P\to e$ is. The proof includes the categorical pushout universal property; a separate exact finite calculation checks the collapses, comma connectivity, and rational homology. This is a worked development, with no claim of priority. The same disk-to-sphere map supplies the promised counterexample to homotopy invariance of total asphericity.

Theorem~4 turns from construction to recognition. Its forward assertion says that smoothness of the three maps $f_i$ implies smoothness of $\varphi_Q$. The reverse assertion for the middle map retains strong W(7 bis). Two slice calculations identify the strict fibres relevant to local smoothness with integrals of lift categories. Footnote~107 correctly observes that the apparent square in the original proof is not strictly commutative. The repair uses the split fibration of the integral over the fork and its aspherical fibres. This proves the forward direction without the strong excision axiom. An explicit contraction removes the same unnecessary hypothesis from the open-cover corollaries. The final unconditional sentence after Corollary~6, however, omits essential outer smoothness assumptions; a four-object lift category with an isolated component gives a finite counterexample.

The last part returns to ordinary category theory. A full sieve $U\subset X$ and its complement $T$ are glued by the presheaves $\operatorname{Hom}_X(-,t)$ on $U$. Proposition~15 says that a square along $U\subset X$ is a pushout exactly when it identifies the complementary categories and transports these crossing-arrow presheaves by left Kan extension. The connected factorization categories in its statement are the fibres of that comparison of presheaves. This gives a full proof of the criterion and explains why an agreement on objects and strata would be insufficient.

For a presheaf $F$ on $X$, the formula
\[
 N(X/F)_n=\coprod_{x_0\to\cdots\to x_n}F(x_n)
\]
then makes the final gluing theorem transparent. It proves preservation of colimits by the category-of-elements construction; the relative construction into $\Cat/X$ also preserves all small limits, the known strengthening recorded by \citet[Proposition~3.2.12 and Corollary~3.2.13]{Cisinski2006}. A pushout of presheaves along a monomorphism becomes a simplicial pushout along a monomorphism. Thus the positive homotopy-pushout clause of Theorem~5 holds for every fundamental localizer. Its subsequent reflection implications still require the strong hypothesis. Model structures require further conditions, supplied later by accessibility and local test categories; attachment stability alone is insufficient.

The later derivator definition identifies $\varphi_Q$ with the terminal component of the coherent left-Kan-extension counit \citep[Definition~3.9]{Groth2013}. The comparison with $\infty$-categories requires another localization: a total category models the categorical colimit after its transport arrows are inverted, while classifying spaces invert all arrows \citep[Corollary~3.3.4.3]{Lurie2009HTT}. These formulations clarify the target of the present calculations without replacing their elementary categorical proofs. The equivalence of the two integration conventions and the strict-cocone question are resolved, as are the specified coherence and local-smoothness gaps. The opening request for a sufficiently large class of cofibrations closed under the stated interval operation is not settled here in its full generality. We record that limitation without declaring it a present-day open problem.
\end{otherlanguage}
\dgrchunkend

% DER_XIIc-a: Grothendieck's text with the commentary set in, then the overview, inside a group that applies the chunk's own definitions
\dgrpart{1}{XIIc-a}{Section 6, part 1 of 2, manuscript pages 171--234}{XIIc-a\enspace\textperiodcentered\enspace Chapter XII, section 6, part 1 of 2\enspace\textperiodcentered\enspace manuscript pages 171--234}{Chapter XII \textperiodcentered\ \textup{XIIc-a}}
\dgrchunkbegin
% DER_XIIc-a: the definitions of the chunk's preamble, applied inside its group
\let\printedpage\relax\newcommand{\printedpage}[1]{{\reversemarginpar\marginnote{\raggedleft\scriptsize ed.~#1}}}
\let\msp\relax\newcommand{\msp}[1]{\marginpar{\raggedright\scriptsize ms.~#1}}
\let\Ch\relax\newcommand{\Ch}{\underline{\mathrm{Ch}}}
\let\Hotp\relax\newcommand{\Hotp}[1]{\mathrm{Hot}_{\bullet\,#1}}
\let\hotp\relax\newcommand{\hotp}[1]{\mathrm{hot}_{\bullet\,#1}}
\let\Ob\relax\newcommand{\Ob}{\mathrm{Ob}\,}
\let\HomEns\relax\newcommand{\HomEns}{\mathrm{Hom}_{\mathrm{Ens}}}
\let\utpi\relax\newcommand{\utpi}{\underset{\sim}{\pi}}
\let\Fib\relax\newcommand{\Fib}{\underline{\mathrm{Fib}}}
\let\Fl\relax\newcommand{\Fl}{\underline{\mathrm{Fl}}}
\let\Hom\relax\newcommand{\Hom}{\underline{\mathrm{Hom}}}
\let\Cat\relax\newcommand{\Cat}{\mathrm{Cat}}
\let\HOT\relax\newcommand{\HOT}{\mathrm{HOT}}
\let\Hot\relax\newcommand{\Hot}{\mathrm{Hot}}
\let\hot\relax\newcommand{\hot}{\mathrm{hot}}
\let\id\relax\newcommand{\id}{\mathrm{id}}
\let\cM\relax\newcommand{\cM}{\mathcal{M}}
\let\cF\relax\newcommand{\cF}{\mathcal{F}}
\let\Ho\relax\newcommand{\Ho}{\mathrm{Ho}}
\let\ho\relax\newcommand{\ho}{\mathrm{ho}}
\let\Wu\relax\newcommand{\Wu}{W^{\mathrm{u}}}
\let\lc\relax\newcommand{\lc}{\mathrm{lc}}
\let\circled\relax\newcommand{\circled}[1]{\tikz[baseline=(c.base)]\node[draw,circle,inner sep=1pt](c){#1};}
\let\deff\relax\newcommand{\deff}{\overset{\text{déf}}{=}}

\selectlanguage{french}


% ---- from Der-ChXIIc-30-03-03_p001-005.tex
\printedpage{115}
% head: A. Grothendieck : Dérivateurs, Ch. XII, Édition des Universités de Montpellier II et Paris VII
% foot: Version 30 mars 2003

\msp{page 171}

\section*{VI\quad Critère de $W_\infty$-asphéricité pour les $\Omega^n(X,x)$.}

\subsection*{1\quad Cogitations sur $W_n$.}

On va définir, pour $n \in \mathbf{N}$,
\[
W_n \subseteq \mathrm{Fl}(\mathrm{Cat})
\]
par les conditions suivantes (où $f : X \longrightarrow Y$ est une flèche dans Cat).

\noindent\begin{tabular}{@{}l c l >{\RaggedRight\arraybackslash}p{0.64\textwidth}@{}}
$f \in W_0$ & $\Longleftrightarrow$ & a) & $\pi_0(f)$ bijectif.\\
 & $\Longleftrightarrow$ & b) & Pour tout préfaisceau d'ensembles $F$ \textit{constant} sur $Y$, $H^0(Y,F) \longrightarrow H^0(X,f^*F)$ est bijectif.\\
 & $\Longleftrightarrow$ & c) & Le foncteur $F \longmapsto f^*F$ des préfaisceaux d'ensembles \textit{constants} sur $Y$ vers les faisceaux d'ensembles constants sur $X$ est pleinement fidèle.
\end{tabular}

\texttt{[Les termes de faisceau ou de préfaisceau sont utilisés ici indifféremment.]}

\textbf{NB} Si $F_0$ est un ensemble, $F_{0\,Y}$ le faisceau d'ensembles constant défini par $F$ sur $Y$, on a
\[
\begin{aligned}
H^0(Y,F_{0\,Y}) &\simeq \HomEns(\pi_0(Y),F_0) && \text{(fonctoriel en $Y$)}\\
\mathrm{Hom}_{Y^\wedge}(F_{0\,Y},G_{0\,Y}) &\simeq \HomEns(\pi_0(Y),\HomEns(F_0,G_0)) && \text{(fonctoriel en $Y$),}
\end{aligned}
\]
ce qui implique que a) $\Longrightarrow$ c) $\underset{\substack{\text{en prenant}\\ \mathrm{Hom}_{Y^\wedge}(e_Y,F)}}{\overset{\text{tautologiquement}}{\Longrightarrow}}$ b) $\Longrightarrow$ a).

\noindent\begin{tabular}{@{}l c l >{\RaggedRight\arraybackslash}p{0.64\textwidth}@{}}
$f \in W_1$ & $\Longleftrightarrow$ & a) & $\pi_0(f)$ bijectif, et pour tout $x \in \Ob X$, $\pi_1(f,x) : \pi_1(X,x) \longrightarrow \pi_1(Y,f(x))$ est bijectif. (NB Il suffit de le vérifier pour \textit{un} $x$ dans chaque composante connexe de $X$.)\\
 & $\Longleftrightarrow$ & a$'$) & $\Pi(X) \xrightarrow{\Pi f} \Pi(Y)$ est une équivalence de catégories.\\
 & $\Longleftrightarrow$ & b) & $H^0(Y,F) \longrightarrow H^0(X,f^*F)$ bijectif pour tout faisceau d'ensembles \textit{localement constant} sur $Y$, et $H^1(Y,F) \longrightarrow H^1(X,f^*F)$ bijectif si de plus $F$ est un faisceau en groupes.\\
 & $\Longleftrightarrow$ & b$'$) & Comme b), mais pour le $H^1$ on se borne à $F$ \textit{constant}, et on exige $H^1(f)$ \textit{surjectif} seulement.\\
 & $\Longleftrightarrow$ & c) & Le foncteur $F \longmapsto f^*F$ de la catégorie des faisceaux d'ensembles localement constants sur $Y$ vers celle des faisceaux localement constants sur $X$ est une équivalence de catégories.
\end{tabular}

\clearpage

\printedpage{116}
% head: A. Grothendieck : Dérivateurs, Ch. XII, Édition des Universités de Montpellier II et Paris VII
% foot: Version 30 mars 2003

\msp{page 172}

\dgrblock{dgrb539}{What the first two tests measure}
\begin{dcomment}{exposition}{p171}{What the first two tests measure}
Grothendieck is trying to recognize finite homotopy information without leaving the category of small categories. The two initial cases explain why locally constant coefficients, rather than arbitrary presheaves, are the appropriate witnesses. Write $BX=|NX|$ for the classifying space of $X$. A locally constant set-valued presheaf is a functor on the fundamental groupoid $\Pi(BX)$, with the variance reversed if one retains the presheaf convention. In this block all changes of variance are absorbed by inversion in the groupoid. The word ``constant'' on p.~171 means globally constant, whereas ``locally constant'' permits monodromy. For a disconnected category these are substantially different conditions.

For a set $A$, a section of the constant presheaf $A_X$ chooses one element of $A$ on each connected component. Thus $H^0(X,A_X)=\operatorname{Map}(\pi_0X,A)$. A natural transformation $A_X\to B_X$ similarly chooses one map $A\to B$ per component. These formulas prove all three tests for $W_0$. In the converse direction, if precomposition with $\pi_0f$ is bijective for every target set, Yoneda gives that $\pi_0f$ is a bijection. The empty category is included: the map from a two-element set to its set of sections on the empty category cannot be bijective.

There is a useful elementary explanation of the apparently weak condition b$'$) for $W_1$. After using constant sets to identify components, work with connected $X,Y$ and put $\alpha:H=\pi_1(X,x)\to G=\pi_1(Y,f(x))$. Sections of the local system corresponding to a $G$-set $A$ are $A^G$, and sections after pullback are $A^H$. Their equality for every $A$ forces $\alpha$ to be surjective. Indeed, let $K=\alpha(H)$ and take the transitive $G$-set $G/K$. The coset $K$ is fixed by $H$; if every $H$-fixed element is $G$-fixed, then $K=G$. Constant nonabelian $H^1$ has the familiar description
\[
 H^1(X,L_X)=\operatorname{Hom}(H,L)/L,
\]
where $L$ acts by conjugation. Surjectivity for $L=H$ says that the class of the identity of $H$ is represented by $\rho\alpha$ for some $\rho:G\to H$. Conjugating $\rho$ makes $\rho\alpha=1_H$, so $\alpha$ is injective as well. If there are several components, choose the identity representation on the chosen component and trivial representations on the others. This proves b$'$)$\Rightarrow$a), including the componentwise assertion.

An equivalence of fundamental groupoids conversely identifies every category of local systems, its global sections, and its torsors. This proves a)$\Rightarrow$b), c). To see directly that c) also forces the groupoid equivalence, full faithfulness recovers global sections as maps from the terminal object and hence forces $\alpha$ to be surjective by the same coset argument; essential surjectivity makes the regular $H$-set the restriction of a $G$-set, so $\ker\alpha$ acts trivially on a faithful $H$-set and is trivial. The torsor interpretation is the first step of the comparison between cohomology and homotopy types recalled in \citet[Preface, pp.~vii--viii]{Lurie2009HTT}. Nothing here requires a choice of a universal covering category: the fundamental groupoid already contains precisely the needed information.
\end{dcomment}


\textbf{NB} Le groupe $\pi_1(X,x)$ est défini via la théorie du groupoïde fondamental. On sait qu'il est en bijection canonique avec $\pi_0(\Omega(X,x))$, où $\Omega(X,x) = \Ch(X;x,x)$ (ce qui permet d'avoir une loi de composition monoïdale sur $\Omega(X,x)$, donnant la loi de groupe sur le $\pi_0$), ou au choix $\Omega(X,x) = \Ch_\infty(X;x,x)$ (plus commode dans certaines situations, mais ne donnant pas directement la loi de composition dans $\pi_0(\Omega(X,x))$).

Je vais poser, pour $n \in \mathbf{N}$,
\[
\pi_n(X,x) = \pi_0(\Omega^n(X,x)),
\]
$\Omega^n$ désignant l'itéré $n$-ième du foncteur $(X,x) \longmapsto \Omega(X,x)$ de la catégorie $\mathrm{Cat}_\bullet$ des objets ponctués de Cat dans elle-même. Je préfère ici travailler avec la théorie \underline{Ch} au lieu de $\Ch_\infty$. Tôt ou tard il faudra donner un énoncé de comparaison \dots

\textbf{Théorème 1} (\footnote[112]{Sans doute \textit{faux}. On doit avoir a) $\Longrightarrow$ b), mais b) $\Longrightarrow$ a) \texttt{[est]} sans doute faux déjà pour $n = 2$ (sauf le cas $X$, $Y$ 1-connexes). Cf. énoncé rectifié \texttt{[p. 178]}.}). \textit{Soit $f : X \longrightarrow Y$ dans} Cat\textit{, et soit $n \in \mathbf{N}$, $n \geq 2$, $S$ une partie de} $\Ob X$ \textit{telle que le composé $S \longrightarrow \Ob X \longrightarrow \pi_0(X)$ soit} \texttt{[un]} \textit{épimorphisme. Les conditions suivantes sont équivalentes.}

\begin{itemize}
\item[\textit{a)}] \textit{$\pi_i(f,x)$ isomorphe pour tout $x \in \Ob X$ et tout $i \in [0,n]$.}

\item[\textit{a$'$)}] \textit{Comme a), mais avec $x \in S$.}

\item[\textit{b)}] \textit{On a}
\[
\begin{array}{rcll}
H^0(Y,F) & \xrightarrow{\ \sim\ } & H^0(X,f^*F) & \textit{pour $F$ faisceau d'ensembles}\\
 & & & \textit{localement constant sur $Y$.}\\
H^1(Y,F) & \xrightarrow{\ \sim\ } & H^1(X,f^*F) & \textit{si de plus $F$ faisceau en groupes.}\\
H^i(Y,F) & \xrightarrow{\ \sim\ } & H^i(X,f^*F) & \textit{si de plus $F$ faisceau abélien, $i \in [2,n]$.}
\end{array}
\]
\end{itemize}

\msp{page 173}

\dgrblock{dgrb540}{Loops, finite types, and the scope of the proposed theorem}
\begin{dcomment}{exposition}{p172}{Loops, finite types, and the scope of the proposed theorem}
The iterated path construction is being asked to do two jobs. First it supplies the groups $\pi_i(X,x)$ intrinsically in $\Cat$; subsequently it must show that these groups are the usual homotopy invariants of $BX$. The later identification of the minimal fundamental localizer with ordinary weak homotopy equivalences removes the distinction between $W_\omega$ and $W_\infty$ in the modern interpretation; see \citet[Theorem~2.2.11]{Cisinski2004Minimal}. It should not be used to erase Grothendieck's effort to obtain the result from the categorical path construction itself.

Theorem~1 contains three separate assertions. The equivalence a)$\Leftrightarrow$a$'$) is the basepoint transport statement established below. The implication a)$\Rightarrow$b) is correct, and the argument on pp.~180--185 supplies its proof after the twisted-coefficient evaluation map is made explicit. The converse b)$\Rightarrow$a) is false in general. The first failure, however, is in degree $3$, not in degree $2$ as footnote~112 suggests. The degree-two case survives even with nontrivial fundamental group. We prove this distinction below, because it materially changes the interpretation of the ensuing ``rectification''.

It is also necessary to distinguish an equivalence of $n$-types from a map with $n$-connected homotopy fibres. For the former, $\pi_0f$ is bijective and $\pi_if$ is an isomorphism for $1\leq i\leq n$. For the latter one needs, in addition, surjectivity on $\pi_{n+1}$. Grothendieck identifies this additional condition correctly on p.~179, but calls the fibre condition the ``good localizer'' on p.~187. That last sentence cannot stand as a definition of a class of weak equivalences: the fibre condition fails the two-out-of-three axiom. The distinction does disappear for a map to a point, which explains why the argument of Theorem~2 is unaffected.
\end{dcomment}


\textbf{NB} L'équivalence de a) et a$'$) devrait être triviale, mais pour cela il faut avoir défini les $\pi_n(X,x)$, pour $x$ variable dans $\Ob X$, comme système local sur $X$, chose que je n'ai pas faite encore. Je vais essayer de le faire dans un courant d'air. Tout d'abord, les $\Omega(X,x)$ sont les fibres d'une $W_\omega$-fibration parfaite (\footnote[113]{\textbf{NB} $W_\omega$ est le plus petit localiseur de Cat (satisfaisant W(1, 2 bis, 3)).}), s'insérant dans le carré cartésien
\[
\begin{tikzcd}
\Ch(X) \arrow[d, "{p_X=(\sigma_X,\beta_X)}"'] & \Omega(X,-) \arrow[l] \arrow[d] \\
X \times X & X. \arrow[l, "\mathrm{diag}_X"'] \arrow[ul, dotted, "i_X"']
\end{tikzcd}
\]

\clearpage

\printedpage{117}
% head: A. Grothendieck : Dérivateurs, Ch. XII, Édition des Universités de Montpellier II et Paris VII
% foot: Version 30 mars 2003

On sait alors que l'on a un foncteur
\[
\begin{aligned}
\Pi X &\longrightarrow \mathrm{Hot}_{W_\omega}\\
x &\longmapsto \mathrm{hot}_{W_\omega}(\Omega(X,x)).
\end{aligned}
\tag*{$(*)$}
\]
Or on a, pour $i \geq 1$,
\[
\pi_i(X,x) = \pi_{i-1}(\ \underbrace{\Omega(X,x)}_{\text{objet ponctué}}\ ).
\]
Cela donne le système local cherché des $\pi_i$, si seulement on avait un foncteur un peu plus précis que $(*)$, savoir
\[
\Pi X \longrightarrow \Hotp{W_\omega} \qquad \text{(catégorie $W_\omega$-homotopique \textit{ponctuée})}
\]
(\footnote[114]{Mais il faut aussi vérifier que les $\pi_n(X,x)$ ($n \geq 1$) sont des foncteurs sur $\Hotp{W_\infty}$, i.e. transforment $W_\infty$-équivalences en itou. Cela résulte de W(4a) pour $W_\infty$, qui implique que $(X,x) \mapsto \Omega(X,x)$ se factorise par $\Hotp{W_\infty}$ (en tant que foncteur à valeurs dans $\Hotp{W_\infty}$).}). Or la flèche $i_X : X \longrightarrow \Ch X$ relevant $p_X$ donne une section canonique de $\Omega(X,-)$ sur $X$ -
\msp{page 174}
celle justement qui associe à tout $x \in X$ le point marqué $1_x$ de $\Omega(X;x,x)$ (`lacet vide en $x$'). La question est donc celle-ci :

\begin{quote}
On a une $W$-fibration $\Omega \longrightarrow X$, avec une section donné $\sigma$. Cela définit-il un foncteur canonique % [sic] donné -> donnée, accord avec « section »
\[
X \longrightarrow \Hotp{W},
\]
transformant toute flèche en \texttt{[un]} isomorphisme?
\end{quote}

L'idée naturelle est de prendre
\[
x \longmapsto \hotp{W}(\Omega/x)
\]
en ponctuant $\Omega/x$ par le point marqué $\sigma(x)$ de $\Omega_x \hookrightarrow \Omega/x$. L'ennui, c'est que les morphismes de transition $\Omega/x \longrightarrow \Omega/x'$ (pour $x \xrightarrow{u} x'$ dans $X$) ne sont pas compatibles avec les points marqués. Mais on a une flèche (savoir $\sigma(u)$) allant
\msp{page 175}
de l'image $\sigma(x)/x'$ du point marqué de $\Omega/x$ vers le point marqué $\sigma(x')/x'$ de $\Omega/x'$. Ceci doit suffire pour définir la flèche
\[
(\Omega/x,\sigma(x)/x) \longrightarrow (\Omega/x',\sigma(x')/x')
\]
comme flèche dans $\Hotp{W}$, avec transitivité pour une composée $x \xrightarrow{u} x' \xrightarrow{v} x''$ dans $X$. Je vais pour l'instant me borner
\printedpage{118}
% head: A. Grothendieck : Dérivateurs, Ch. XII, Édition des Universités de Montpellier II et Paris VII
% foot: Version 30 mars 2003
\msp{page 176}
à une indication, quitte à faire des vérifications circonstanciées par la suite. Au besoin, j'utiliserai dès maintenant que pour $i \in \mathbf{N}$, $i \geq 1$,
\[
\pi_i(f,x) : \pi_i(X,x) \longrightarrow \pi_i(Y,f(x))
\]
peuvent s'interpréter comme provenant d'un homomorphisme de systèmes locaux
\[
\utpi_i(X,-) \longrightarrow f^*(\utpi_i(Y,-))
\]
sur $X$. La condition a) du théorème 1 s'interprète en disant que cet homomorphisme est un \textit{isomorphisme} pour $i \in [1,n]$, et la condition a$'$) se borne alors à rappeler qu'il suffit de tester qu'il en est ainsi pour \textit{un} point de chaque composante connexe de $X$.

\dgrblock{dgrb541}{Making pointed transport explicit}
\begin{dcomment}{development}{p175}{Making pointed transport explicit}
Grothendieck needs to transport a marked object along an arrow even when the underlying transition functor does not preserve that object strictly. The discrepancy is useful data: the arrow supplied by the section tells us exactly which change of basepoint to make. The following construction fills in the indicated verification using only arrow categories and localization. We checked the natural-homotopy and adjunction lemmas in \citet[\S1.1, especially Proposition~1.1.9 and Corollary~1.1.10]{Cisinski2004Minimal}; we did not locate a discussion of this particular pointed comma-category formula in the consulted sources. The result is an elementary completion of the manuscript's construction, with no assertion of priority.

\textbf{Pointed transport lemma.}\ours\ Suppose $W$ contains every functor admitting a left or a right adjoint and satisfies two-out-of-three. In the localization of pointed categories at functors whose underlying functors lie in $W$, an arrow $a:d\to e$ of a category $D$ defines a canonical isomorphism
\[
 t_a:(D,d)\longrightarrow(D,e).
\]
It is natural in functors $D\to D'$, satisfies $t_{1_d}=1$, and satisfies $t_{ba}=t_bt_a$ for composable arrows $a,b$.

\textit{Proof.} Regard $a$ as the marked object of $D^{[1]}$. The two evaluation functors give a pointed span
\[
 (D,d)\xleftarrow{\ e_0\ }(D^{[1]},a)
 \xrightarrow{\ e_1\ }(D,e).
\]
The constant-diagram functor is left adjoint to $e_0$ and right adjoint to $e_1$, so both evaluations belong to $W$. Define $t_a=[e_1][e_0]^{-1}$. For the identity arrow the constant-diagram functor is a pointed common section, proving the unit assertion. For a composable pair, mark $D^{[2]}$ by the diagram $d\xrightarrow a e\xrightarrow b h$. Evaluation at every vertex belongs to $W$. For the endpoints this again follows from adjunctions. For the middle vertex use the natural transformation $e_0\Rightarrow e_1$ and the following elementary consequence of the assumptions: naturally related functors belong to $W$ simultaneously. Indeed, represent a natural transformation by $H:A\times[1]\to B$. The projection $A\times[1]\to A$ has an adjoint and lies in $W$; each endpoint section lies in $W$ by two-out-of-three. Thus $H i_0\in W$ if and only if $H\in W$, if and only if $H i_1\in W$. Restriction to the three edges commutes with vertex evaluation and is therefore in $W$ by two-out-of-three. On this common pointed source,
\[
 [e_2][e_0]^{-1}
 =\bigl([e_2][e_1]^{-1}\bigr)
  \bigl([e_1][e_0]^{-1}\bigr).
\]
This is the composition assertion. Applying a functor to arrow diagrams commutes strictly with all evaluations and restrictions, proving naturality.\hfill$\square$

Now let $q:E\to X$ be the fibration under discussion, with section $\sigma$. For $u:x\to x'$ write $u_*:E/x\to E/x'$ for postcomposition. Its marked source maps to $u_*\sigma(x)$, and the section gives an arrow
\[
 a_u:u_*\sigma(x)\longrightarrow\sigma(x')
\]
in $E/x'$. Define the pointed transition to be $t_{a_u}[u_*]$. For $v:x'\to x''$, the equality $a_{vu}=a_v\,v_*(a_u)$, naturality of $t$, and the composition formula prove the functor law. When the comma transitions belong to $W$, these maps are isomorphisms and the functor factors through $\Pi X$. The pointed inclusions $E_x\to E/x$ identify it with the family of strict fibres whenever those inclusions belong to $W$. Applied to the loop fibration, this gives the local systems of higher homotopy groups and proves the equivalence of the one-basepoint and every-basepoint formulations. It gives ordinary transport in the pointed homotopy category; a coherent family of pointed spaces contains further information, which is not claimed to follow from this ordinary functor alone.
\end{dcomment}


L'énoncé du théorème 1 garde un sens pour $n = 1$, et dans ce cas a été `rappelé' précédemment (et est $\pm$ tautologique). Cela va nous permettre de procéder par récurrence sur $n$, en supposant que le théorème est prouvé pour $n-1$. Cela signifie que l'on peut supposer que les conditions a) \textit{et} b) sont satisfaites par $f$ pour $i \in [1,n-1]$, et se borner à prouver que, sous ces conditions, la condition a) (ou a$'$), pour $i = n$, équivaut à la condition b) pour $i = n$.

\msp{page 177}

Donc il faut prouver simplement ceci :

\textbf{Lemme} (\footnote[115]{Sans doute faux.}). \textit{Supposons que $\pi_i(f,x)$} \texttt{[est un]} \textit{isomorphisme pour $i \in [0,n-1]$ et b) satisfait pour $i \in [0,n-1]$, où $n \in \mathbf{N}$, $n \geq 2$, est fixé. Alors les conditions suivantes sont équivalentes :}

\begin{itemize}
\item[\textit{a)}] $\pi_n(f,x) : \pi_n(X,x) \longrightarrow \pi_n(Y,f(x))$ \texttt{[est un]} \textit{isomorphisme pour tout $x \in \Ob X$.}

\item[\textit{b)}] $H^n(Y,F) \xrightarrow{\sim} H^n(X,f^*F)$ \texttt{[est un]} \textit{isomorphisme pour tout faisceau abélien localement constant sur $Y$.}
\end{itemize}

Notons que $f$ induit une flèche dans $\mathrm{Cat}_\bullet$,
\[
\Omega(f,x) : \Omega(X,x) \longrightarrow \Omega(Y,f(x)),
\]
d'où
\[
\begin{tikzcd}[row sep=small]
\pi_i(\Omega(X,x)) \arrow[r] \arrow[d, dash, "\wr"'] & \pi_i(\Omega(Y,f(x))) \arrow[d, dash, "\wr"'] \\
\pi_{i+1}(X,x) \arrow[r] & \pi_{i+1}(Y,f(x)),
\end{tikzcd}
\]
s'identifiant à $\pi_{i+1}(f,x)$ :
\[
\pi_i(\Omega(f,x)) \simeq \pi_{i+1}(f,x).
\]


\printedpage{119}
% head: A. Grothendieck : Dérivateurs, Ch. XII, Édition des Universités de Montpellier II et Paris VII
% foot: Version 30 mars 2003

L'hypothèse a) du lemme implique que $\pi_i(\Omega(X,x)) \xrightarrow{\sim} \pi_i(\Omega(Y,y))$ pour $i \in [0,n-1]$. Il s'agit ici des $\pi_i$ basés en les points unités. On devra sans doute prouver que, du fait qu'il s'agit d'un morphisme de monoïdes dans Cat (satisfaisant certaines conditions?) ça reste vrai pour les $\pi_i$ en \textit{tout} point $\xi$ de $\Omega(X,n)$. Par l'hypothèse de récurrence, on saura donc que $\Omega(X,x) \longrightarrow \Omega(Y,y)$ % [sic] \Omega(X,n) -> \Omega(X,x), le point de base est x
\msp{page 178}
induit un isomorphisme pour la cohomologie à coefficients localement constants, jusqu'en dimension $n-1$ inclus.

Finalement, je me suis $\pm$ convaincu que l'énoncé du théorème 1 est faux tel quel. Voici un énoncé que j'éspère vrai, et que je vais essayer de prouver dans un esprit `utilitaire', en utilisant tout ce que je sais en homotopie et cohomologie (et pas seulement ce que j'ai développé `from scratch' dans le modélisateur Cat). % [sic] j'éspère -> j'espère, orthographe

\textbf{Théorème 1 rectifié.} \textit{Considérons, pour un entier $n \geq 2$ donné, les conditions suivantes sur une flèche $f : X \longrightarrow Y$ dans} Cat \textit{(ou une application continue d'espaces topologiques, ou un morphisme dans $\mathbf{\Delta}^\wedge$ etc. etc.) :}

\begin{itemize}
\item[$a_n$)] \textit{Si $Z$ est une `fibre homotopique' de $f$ (NB il y en a essentiellement une par composante connexe de $Y$), alors $Z$ est $n$-connexe, i.e.}
\msp{page 179}
\textit{(p. ex) $\pi_i(Z) = 0$ pour $i \in [0,n]$.}

\dgrblock{dgrb542}{The first genuine failure and the role of loop multiplication}
\begin{dcomment}{examples}{p178}{The first genuine failure and the role of loop multiplication}
The withdrawn induction fails because ordinary cohomology in degree $n$ need not see all of $\pi_n$ when lower homotopy is already present. A concrete example makes the failure exact. Take
\[
 f:S^2=\mathbf{CP}^{1}\longrightarrow\mathbf{CP}^{\infty}=K(\mathbf Z,2).
\]
Both spaces are simply connected. Their integral cellular chain complexes have $\mathbf Z$ in degree $0$ and degree $2$, and no groups in degrees $1$ and $3$; $\mathbf{CP}^{\infty}$ has additional even-dimensional cells starting in degree $4$. Thus for every abelian group $M$, the map induces isomorphisms on $H^i(-;M)$ for $0\leq i\leq3$: the nonzero groups in this range are $H^0=M$ and $H^2=M$. Nonabelian $H^1$ is trivial on both spaces and every local system is constant. Yet
\[
 \pi_3(S^2)=\mathbf Z\longrightarrow
 \pi_3(K(\mathbf Z,2))=0
\]
is not an isomorphism. The source group is detected by the Hopf fibration, and the target calculation follows from $S^1\to S^\infty\to\mathbf{CP}^{\infty}$; both computations appear in \citet[Examples~4.50--4.51]{Hatcher2002}. This disproves Theorem~1 and the lemma on p.~177 for $n=3$, even with simply connected source and target. Applying the category-of-simplices construction to singular simplicial models gives an actual functor of small categories with the same failure; the nerve/category-of-simplices equivalence is recorded in \citet[Lemma~2.1.15]{Cisinski2004Minimal}.

The proposed passage from the unit component of a loop space to all components is nevertheless sound. If $M$ is a topological monoid with $\pi_0M$ a group, multiplication by any $m\in M$ is a homotopy equivalence. Choose $m'$ representing its inverse component. Paths from $mm'$ and $m'm$ to the unit show that multiplication by $m'$ is a homotopy inverse, using associativity. A homomorphism of such monoids commutes with these translations. Hence its effect on homotopy groups in every component is determined by its effect at the unit, together with the bijection on components. Realization of a monoidal category has this property when its monoid of components is a group, as it is for Grothendieck's loop category. The monoidal argument repairs the specific verification requested at the top of p.~178; it does not repair the false cohomological converse used by the abandoned induction.
\end{dcomment}


\item[$b_n$)] \textit{$\pi_0(f)$ est bijectif, et $\forall\ x \in X$, $\pi_i(f,x) : \pi_i(X,x) \longrightarrow \pi_i(Y,f(x))$ est} \texttt{[un]} \textit{isomorphisme pour $i \in [1,n]$. (}\textbf{NB} \textit{Il suffit de le dire pour} \textup{un} \textit{point $x$ de chaque composante connexe de $X$.)}

\item[$c_n$)] $H^i(Y,F) \longrightarrow H^i(X,f^*F)$ \texttt{[est un]} \textit{isomorphisme pour $F$ faisceau localement constant sur $Y$ (avec $F$ abélien si $i \geq 2$, $F$ faisceau en groupes si $i = 1$).}
\end{itemize}

\textit{On a alors}
\[
a_n \ \Longrightarrow\ b_n \ \Longrightarrow \text{(\footnotemark[116])}\ c_n \ \Longrightarrow \text{(\footnotemark[117])}\ a_{n-1}
\]
\footnotetext[116]{si $X$, $Y$ 1-connexes, c'est vraiment prouvé, sinon, il y a des vérifications à faire (qui ont l'air soritales \dots).}%
\footnotetext[117]{\textbf{NB} Même si $X$, $Y$ 1-connexes, il n'est pas clair que $\mathrm{c}_n \Longrightarrow \mathrm{b}_n$, car il faut montrer que $\pi_n(Z) \longrightarrow \pi_n(X)$ est nul, et on trouve seulement que le composé avec $\pi_n(X) \longrightarrow H_n(X)$ est nul. Je présume qu'il y a des contre-exemples.}%
\textit{(et je présume que chacune de ces implications est stricte).}

\textsc{Démonstration}. $\mathrm{a}_n \Longrightarrow \mathrm{b}_n$ est conséquence immédiate de la suite exacte d'homotopie pour une Hot-fibration $\bar{X} \xrightarrow{\bar{f}} Y$ qui factorise $X \xrightarrow{f} Y$ en $f = \bar{f} \circ i$, $i \in W_\infty$. En fait, $\mathrm{a}_n$ \textit{équivaut} à $\mathrm{b}_n$ + la surjectivité des $\pi_{n+1}(X,x) \longrightarrow \pi_{n+1}(Y,f(x))$, laquelle bien sûr ne résulte \textit{pas} de $\mathrm{b}_n$, i.e. de la bijectivité des $\pi_i(f,x)$ pour $i \leq n$ (prendre des $K(\pi,n+1)$). Donc l'implication est bien stricte.

% ---- from Der-ChXIIc-30-03-03_p006-010.tex
\printedpage{120}
% head: A. Grothendieck : Dérivateurs, Ch. XII, Édition des Universités de Montpellier II et Paris VII

\msp{page 180}
$\mathrm{b}_n \Longrightarrow \mathrm{c}_n$. On peut supposer que $f$ est une Hot-fibration (et même Hot-parfait si ça nous chante \dots). On sait déjà que pour
\[
H^i(Y,F) \longrightarrow H^i(X,f^*F) \tag*{$(*)$}
\]
tout est O.K. si $i \in \{0,1\}$. Donc on peut supposer $F$ abélien. Si ça nous avance, on peut supposer, par récurrence, (que l’implication $\mathrm{b}_{n-1} \Longrightarrow \mathrm{c}_{n-1}$ est connue, donc) que $(*)$ est \texttt{[un]} isomorphisme pour $i<n$. Reste à le voir pour $i=n$. On est ramené bien sûr, comme $\pi_0(f)$ est bijectif, au cas où $X$, $Y$ sont 0-connexes, donc il y a essentiellement une seule fibre $Z$. La suite exacte d’homotopie nous apprend que
\[
\pi_i(Z,z)=0 \qquad \text{pour } i \le n-1 \ (Z \text{ est } (n-1)\text{-connexe})
\]
et de plus
\[
\pi_n(Z,z) \longrightarrow \pi_n(X,z) \qquad \text{est nul}
\]
(car $\pi_n(X,z) \longrightarrow \pi_n(Y,y)$ \texttt{[est un]} isomorphisme donc \texttt{[un]} monomorphisme), donc
\[
\pi_{n+1}(Y,y) \longrightarrow \pi_n(Z,z) \qquad \text{\texttt{[est un]} épimorphisme.}
\]
Comme $Z$ est $(n-1)$-connexe (donc aussi 1-connexe), on a
\[
H^i(Z,G)=0\ , \qquad \text{si } 0<i<n\ , \tag*{$(*)$}
\]
pour $G$ localement constant (donc constant) sur $Z$.

\dgrblock{dgrb543}{Three conditions and the extra degree}
\begin{dcomment}{exposition}{p179}{Three conditions and the extra degree}
The corrected theorem separates fibre connectivity, homotopy isomorphisms, and cohomological isomorphisms. This separation is the main conceptual gain of the first section. In condition $c_n$ the printed text has omitted the range of $i$; the intended range is $0\leq i\leq n$, as both proof and notation require. We use singular local coefficients on classifying spaces, or on CW replacements when discussing spaces. A statement about arbitrary topological spaces with their ordinary sheaf cohomology would require additional hypotheses and is not supplied by the argument here.

For a homotopy fibre $Z$ over $y=f(x)$, the exact homotopy sequence gives
\[
 \pi_{n+1}(X,x)\longrightarrow\pi_{n+1}(Y,y)
 \longrightarrow\pi_n(Z,z)\longrightarrow
 \pi_n(X,x)\longrightarrow\pi_n(Y,y).
\]
Together with the preceding terms, this proves
\[
 a_n\quad\Longleftrightarrow\quad
 b_n\ \text{and}\ 
 \pi_{n+1}(X,x)\twoheadrightarrow\pi_{n+1}(Y,y)
\]
for one basepoint in every component, with nonempty connected fibres understood. In contemporary terminology $a_n$ says that $f$ is $(n+1)$-connected when a $k$-connected map is defined to have $(k-1)$-connected fibres. Authors also use a convention shifted by one; specifying the fibres avoids that ambiguity.

The map $*\to K(A,n+1)$, for $A\ne0$, satisfies $b_n$ but not $a_n$, since its homotopy fibre is $K(A,n)$. At the other end, $S^n\to *$ has $(n-1)$-connected fibre, but fails $c_n$: with $M=\mathbf Z$, its degree-$n$ cohomology map is $0\to\mathbf Z$. These examples prove strictness of the first and last arrows in the chain. The middle arrow is strict for $n=3$ by the preceding example, but it is an equivalence for $n=2$. Thus the conjecture that \emph{each} implication is strict for every $n\geq2$ is false. This does not affect the implication chain itself.

The lower boundary in the fibre statement matters. A bijection on components and isomorphisms on $\pi_i$ through degree $n$ yield only $(n-1)$-connected fibres; the remaining group $\pi_nZ$ can be the quotient of $\pi_{n+1}Y$ by the image of $\pi_{n+1}X$. It is precisely this group that enters the next low-degree exact sequence. The proof succeeds by showing that its restriction map from the cohomology of $X$ is zero, rather than by claiming that the group itself vanishes.
\end{dcomment}


\msp{page 181}
Écrivons la suite spectrale de Leray pour $f : X \longrightarrow Y$ et $f^*F = F_X$ :
\[
H^*(X,f^*F) \Longleftarrow E_2^{pq} = H^p(Y,H^q(Z,f^*F)),
\]
où ‘$H^q(Z)$’ doit être vu comme un système local sur $Y$. Donc par $(*)$ on a
\[
E_2^{pq}=0 \qquad \text{si } 0<q<n, \qquad
\begin{tikzpicture}[baseline=0pt, x=0.6cm, y=0.6cm]
\draw[->] (0,0) -- (0,4);
\draw[->] (0,0) -- (6,0);
\node[left] at (0,3) {$n$};
\foreach \i in {0,...,4} { \fill (\i,3) circle (2.5pt); \fill (\i,0) circle (2.5pt); }
\node at (1,1.5) {$0$};
\node at (3,1.5) {$0$};
\end{tikzpicture}
\qquad ,
\]
donc on trouve la suite exacte
\[
0 \longrightarrow H^n(Y,F) \longrightarrow H^n(X,F_X) \xrightarrow{\ \alpha\ } H^0(Y,\ \underbrace{H^n(Z,\overbrace{F_Z}^{M_Z})}_{\begin{array}{l}\scriptstyle = \ \mathrm{Hom}(\pi_n(Z,z),M)\\ \scriptstyle \text{par \texttt{[le]} th. de Hurewicz}\end{array}}\ ) \longrightarrow H^{n+1}(Y,F).
\]

\printedpage{121}
% head: A. Grothendieck : Dérivateurs, Ch. XII, Édition des Universités de Montpellier II et Paris VII

Ici je fixe un point $x$ de $X$, $Z$ désigne la fibre de $f$ en $y=f(x)$, et $M$ celle de $F$ en $y$, $M=F_y$. Le groupe $\pi_1(Y)$ y opère, d’ailleurs $\pi \overset{\text{déf}}{=} \pi_1(X,x) \simeq \pi_1(Y,y)$ aussi. Le $H^0(Y,H^n(Z,F_Z))$ peut s’interpréter, sauf erreur, comme
\[
\mathrm{Hom}_\pi(\pi_n(Z,x),M),
\]
en tenant compte de l’opération de $\pi=\pi_1(X,x)$ sur $\pi_n(Z,x)$ également. (À vérifier.) Il faut interpréter la flèche
\[
H^n(X,F_X) \longrightarrow \mathrm{Hom}_\pi(\pi_n(Z,x),M),
\]
via l’application composée \texttt{[où $z=x$, vu comme point de $Z$]}
\[
\pi_n(Z,z) \longrightarrow \pi_n(X,x) \longrightarrow H_n(X,\mathbf{Z}). \tag*{$(*)$}
\]
Si $F$ est \textit{constant} sur $Y$ (p. ex. si $Y$ \texttt{[est]} 1-connexe), donc de valeur $M$, on a les flèches canoniques
\[
H^n(X,F_X) = H^n(X,M) \longrightarrow \mathrm{Hom}(H_n(X,\mathbf{Z}),M) \longrightarrow \mathrm{Hom}(\pi_n(Z,z),M)\text{(\footnotemark[118])} \tag*{$(*')$}
\]
\footnotetext[118]{grâce à la composition avec $(*)$.}
\texttt{[lire : }$\mathrm{Hom}_\pi$\texttt{]},

\msp{page 182}
et sûrement (en dehors de toute hypothèse de type $\mathrm{b}_n$ ou autre) l’application canonique directe ‘de localisation sur $Y$’, suivi de \dots,
\[
H^n(X,F_X) \longrightarrow H^0(Y,H^n(Z,F_Z)) \longrightarrow \mathrm{Hom}_\pi(\pi_n(Z),M) \tag*{$(**)$}
\]
n’est autre que la composée des applications canoniques ci-dessus. Dans le cas actuel où on suppose $\mathrm{b}_n$, la deuxième flèche dans $(**)$ est \texttt{[un]} isomorphisme, donc la composée ‘n’est autre’ que la flèche de localisation elle-même. D’autre part $\pi_n(Z,z) \longrightarrow \pi_n(X,x)$ dans $(*)$ est nul, donc aussi $\pi_n(Z,z) \longrightarrow H_n(X,\mathbf{Z})$, donc aussi la flèche composée $(*')$ qui s’en déduit. Revenons alors à la suite exacte de Leray (page précédente), on trouve que $H^n(Y,F) \xrightarrow{\ \sim\ } H^n(X,F_X)$.

\dgrblock{dgrb544}{The exact sequence behind the cohomological implication}
\begin{dcomment}{exposition}{p181}{The exact sequence behind the cohomological implication}
The use of a path replacement now lets Grothendieck import the classical fibration argument into his categorical model. Assume $b_n$, reduce to a connected component, and replace $f$ by a fibration with fibre $Z$. Then $Z$ is $(n-1)$-connected, and the homotopy sequence also says that $\pi_nZ\to\pi_nX$ is zero. For an abelian local system $F$ on $Y$ with fibre $M$ at $y$, the restriction $F|_Z$ is constant. The cohomological Serre--Leray spectral sequence has
\[
 E_2^{p,q}=H^p(Y;\mathcal H^q(Z;M))
 \quad\Longrightarrow\quad H^{p+q}(X;f^*F).
\]
Here $\mathcal H^q(Z;M)$ includes both the transport of the fibre and that of $M$; replacing it by an untwisted constant group would lose the content of the argument. The historical spectral sequence is \citet[Chapter~II, \S8, Proposition~6]{Serre1951}; for the indexing and filtration in the untwisted presentation see \citet[Theorem~1.14]{HatcherSS2004}.

Since $H^q(Z;M)=0$ for $0<q<n$, the only contributions in total degrees at most $n$ are the bottom row and $E_2^{0,n}$. No differential reaches the bottom row before $d_{n+1}$ from $(0,n)$ to $(n+1,0)$. Reading the filtration therefore gives the natural exact sequence
\[
\begin{split}
0\longrightarrow H^n(Y;F)&\longrightarrow H^n(X;f^*F)
 \xrightarrow{\alpha}\operatorname{Hom}_{\mathbf Z[\pi]}
       (\pi_nZ,M)\\
 &\longrightarrow H^{n+1}(Y;F)
 \longrightarrow H^{n+1}(X;f^*F),
\end{split}
\]
where $\pi=\pi_1Y=\pi_1X$. Only exactness through the penultimate group is needed. The identification of $E_2^{0,n}$ uses the natural Hurewicz isomorphism $\pi_nZ\cong H_nZ$ and the universal coefficient theorem, followed by taking monodromy invariants. The fibre is simply connected, so no ambiguity from inner automorphisms survives in its degree-$n$ group. The map $\alpha$ is restriction to a fibre, viewed as an invariant homomorphism.

There is an efficient chain-level way to finish the twisted case. If $T$ is connected with fundamental group $\pi$, a local system $G$ with fibre $M$ is represented by a left $R=\mathbf Z[\pi]$-module. A cocycle representing $c\in H^n(T;G)$ is an $R$-linear map
\[
 z:C_n(\widetilde T)\longrightarrow M,
 \qquad z\partial=0.
\]
Evaluation on cycles gives a homomorphism $H_n(\widetilde T)\to M$; changing $z$ by a coboundary has zero effect on cycles. Composing with the equivariant Hurewicz homomorphism produces
\[
 \varepsilon_G:H^n(T;G)\longrightarrow
 \operatorname{Hom}_R(\pi_n(T,t),M),\qquad n\geq2.
\]
This is the desired spherical evaluation. Its additivity follows from the homomorphism property of Hurewicz, or from the pinch map of the sphere; its equivariance follows directly from $R$-linearity. Naturality follows by lifting a pointed map to universal covers and applying the chain map. The local-coefficient chain model is described in \citet[\S3H, pp.~327--334]{Hatcher2002}; the verification above supplies the particular application the manuscript leaves open.

For $j:Z\to X$, choose the lift $Z\to\widetilde X$ corresponding to the marked point. Because $Z$ is simply connected and $(n-1)$-connected, naturality of evaluation and Hurewicz identify $\alpha$ with
\[
 H^n(X;f^*F)\xrightarrow{\varepsilon_{f^*F}}
 \operatorname{Hom}_R(\pi_nX,M)
 \xrightarrow{j^*}\operatorname{Hom}_R(\pi_nZ,M).
\]
The last map is zero because $j_*:\pi_nZ\to\pi_nX$ is zero. Hence $\alpha=0$, and exactness proves that $H^n(Y;F)\to H^n(X;f^*F)$ is an isomorphism. In lower degrees it is already an isomorphism by the zero rows of the spectral sequence. Degrees $0$ and nonabelian $1$ are controlled by the fundamental groupoid. This completes $b_n\Rightarrow c_n$ for all the coefficients actually stated, without an assumption of simple connectivity.
\end{dcomment}


Comment récupérer le cas où $F$ n’est \textit{pas} constant? Dans le cas $n=2$, cela marchait, en passant aux revêtements universels $\tilde{X}$, $\tilde{Y}$ de $X$, $Y$ (par rapport aux points $x$, $y$ choisis), et en écrivant les suites spectrales
\[
\begin{tikzcd}[row sep=small, column sep=small]
H^*(Y,F) \arrow[d] & \Longleftarrow & E_2^{pq} \arrow[d] & = & H^p(\pi,H^q(\tilde{Y},F_{\tilde{Y}})) \\
H^*(X,F_X) & \Longleftarrow & E^{\prime pq}_2 & = & H^p(\pi,H^q(\tilde{X},F_{\tilde{X}})).
\end{tikzcd}
\]
L’hypothèse $\mathrm{b}_n$ est conservée en passant de $f : X \longrightarrow Y$ à $\tilde{f} = \tilde{f}_x : \tilde{X} \longrightarrow \tilde{Y}$, donc on sait que
\printedpage{122}
% head: A. Grothendieck : Dérivateurs, Ch. XII, Édition des Universités de Montpellier II et Paris VII
\msp{page 183}
\[
E_2^{pq} \xrightarrow{\ \sim\ } E^{\prime pq}_2 \qquad \text{\texttt{[est un]} isomorphisme si } q \le n. \tag*{$(*)$}
\]
Dans le cas $n=2$, utilisant que
\[
E_2^{p1} = E^{\prime p1}_2 = 0 \qquad \text{car } H^1(\tilde{X}) = H^1(\tilde{Y}) = 0, \qquad
\begin{tikzpicture}[baseline=0pt, x=0.5cm, y=0.5cm]
\draw[->] (0,0) -- (0,2.8);
\draw[->] (0,0) -- (5,0);
\node[above left] at (0,2.6) {$q$};
\node[left] at (0,2) {$2$};
\foreach \i in {0,...,3} { \fill (\i,2) circle (2.5pt); \fill (\i,0) circle (2.5pt); }
\node at (0,1) {$0$};
\node at (1,1) {$0$};
\node at (2,1) {$0$};
\end{tikzpicture}
\qquad ,
\]
on trouve deux suites exactes
\[
\begin{tikzcd}
0 \arrow[r] & H^2(\pi,H^0(\tilde{Y})) \arrow[r] \arrow[d, "\text{iso}"] & H^2(Y) \arrow[r] \arrow[d] & H^0(\pi,H^2(\tilde{Y})) \arrow[r] \arrow[d, "\text{iso}"] & H^3(\pi,H^0(\tilde{Y})) \arrow[d, "\text{iso}"] \\
0 \arrow[r] & H^2(\pi,H^0(\tilde{X})) \arrow[r] & H^2(X) \arrow[r] & H^0(\pi,H^2(\tilde{X})) \arrow[r] & H^3(\pi,H^0(\tilde{X})),
\end{tikzcd}
\]
et un homomorphisme de l’une dans l’autre, qui permet de conclure par le lemme des 5 que $H^2(Y,F) \longrightarrow H^2(X,F_X)$ est \texttt{[un]} isomorphisme. Mais pour $n$ quelconque, la relation $(*)$ sur les termes $E_2$ (pour $q \le n$) ne me semble pas impliquer la relation similaire pour les $E_n$ supérieures, et de là pour les termes $E_\infty$, ce qui permettrait de conclure.

Je reviens donc à la suite exacte p. 181, avec $F$ faisceau \textit{tordu.} J’ai envie de trouver une application canonique
\[
H^n(X,G) \longrightarrow \overbrace{\mathrm{Hom}_\pi(\pi_n(X,x),G_x)}^{\text{\scriptsize Homs \textit{additifs} bien sûr, i.e. de $\mathbf{Z}[\pi]$-modules}} \qquad ,
\]
où $G$ est un faisceau abélien localement constant quelconque sur $X$, et $\pi = \pi_1(X,x)$. Si j’interprète
\msp{page 184}
$\pi_n(X,x)$ par des classes d’homotopie d’applications ponctuées
\[
(S^n,x_0) \xrightarrow{\ \varphi\ } (X,x)
\]
j’ai bien, pour une telle application $\varphi$,
\[
\varphi^* : H^n(X,G) \longrightarrow H^n(S^n,\ \overbrace{\varphi^*G}^{\simeq\,(G_x)_{S^n}}\ ) \simeq G_x
\]
(tenant compte que $S^n$ est 1-connexe, et que sa $n$-cohomologie à valeurs dans un faisceau constant $M$ est $M$.) Sûrement l’application
\[
H^n(X,G) \longrightarrow G_x
\]
\printedpage{123}
% head: A. Grothendieck : Dérivateurs, Ch. XII, Édition des Universités de Montpellier II et Paris VII
ainsi définie ne dépend que de la classe d’homotopie ponctuée de $\varphi$, d’où
\[
\pi_n(X,x) \times H^n(X,G) \longrightarrow G_x,
\]
et il faut de plus prouver a) que c’est compatible avec les opérations de $\pi = \pi_1(X,x)$, et b) que c’est additif en $\pi_n(X,x)$, pour avoir l’application cherchée
\[
H^n(X,G) \longrightarrow \mathrm{Hom}_\pi(\pi_n(X,x),G_x).
\]
Si j’admets ceci, il n’y a qu’un pas pour admettre aussi que la flèche qui nous intéresse vraiment (dans la suite exacte p. 181)
\msp{page 185}
\[
H^n(X,F_X) \longrightarrow \mathrm{Hom}_\pi(\pi_n(Z,z),\underbrace{M}_{=\,F_y})
\]
n’est autre que la composée de la flèche précédente (page précédente) pour $G = F_X$, avec
\[
\mathrm{Hom}_\pi(\pi_n(X,x),M) \longrightarrow \mathrm{Hom}_\pi(\pi_n(Z,x),M)
\]
déduite de
\[
\pi_n(Z,x) \longrightarrow \pi_n(X,x).
\]
Or cette flèche est nulle, par l’hypothèse $\mathrm{b}_n$. Donc la flèche $\alpha$ dans la suite exacte page 181 est nulle, donc $H^n(Y,F) \longrightarrow H^n(X,F_X)$ est \texttt{[un]} isomorphisme, q.e.d. (mais il y a pas mal de travail de vérification à faire \dots).

\dgrblock{dgrb545}{Completing the universal-cover comparison}
\begin{dcomment}{development}{p183}{Completing the universal-cover comparison}
Grothendieck briefly abandons the universal-cover method because he does not see how its comparison of $E_2$ rows controls later pages. The method can be completed directly at the level of complexes. It is useful to record the argument, since the needed range is finite and one should not invoke an unrestricted comparison theorem without checking its boundary. We use the ordinary truncation of complexes and the local-coefficient complex in \citet[\S3H]{Hatcher2002}. We did not find this particular completion of p.~183 in the consulted literature; its ingredients are standard homological algebra.

\textbf{Finite-range descent lemma.}\ours\ Let $\pi$ be a group and let $u:C\to D$ be a map of nonnegative cochain complexes of $\mathbf Z[\pi]$-modules which induces isomorphisms on cohomology in degrees $0\leq q\leq n$. Then
\[
 H^k(R\Gamma(\pi,C))\longrightarrow
 H^k(R\Gamma(\pi,D))
\]
is an isomorphism for $0\leq k\leq n$.

\textit{Proof.} The smart truncation $\tau_{\leq n}C$ agrees with $C$ below $n$, has $\ker(d:C^n\to C^{n+1})$ in degree $n$, and is zero above $n$. The natural inclusion fits into the truncation triangle
\[
 \tau_{\leq n}C\longrightarrow C\longrightarrow\tau_{\geq n+1}C.
\]
Derived invariants carry a complex whose cohomology vanishes below $n+1$ to one with the same lower bound. Indeed, choose an injective resolution starting in degree $n+1$, or inspect the first-quadrant spectral sequence with terms $H^p(\pi,H^q(-))$: all its potentially nonzero terms have $p\geq0$ and $q\geq n+1$. Applying $R\Gamma$ to the triangle and taking its long exact cohomology sequence shows that
\[
 H^k(R\Gamma(\pi,\tau_{\leq n}C))
 \xrightarrow{\sim}H^k(R\Gamma(\pi,C)),\qquad k\leq n,
\]
because the adjacent terms of degrees $k-1$ and $k$ for $\tau_{\geq n+1}C$ both vanish. The same holds for $D$. Finally $\tau_{\leq n}u$ is a quasi-isomorphism, and derived invariants preserve quasi-isomorphisms. The two truncation comparisons prove the assertion.\hfill$\square$

Apply the lemma to the equivariant cochain complexes of the universal covers with constant fibre $M$ and its given $\pi$-action. Derived invariants compute the cohomology with the corresponding local coefficients on the quotient: one may use a free resolution of $\mathbf Z$ over $\mathbf Z[\pi]$ and the usual total double complex. Indeed, each group of covering chains is free over $\mathbf Z[\pi]$, so each cochain term is a product of coinduced modules and is acyclic for the invariants functor. Its ordinary invariant complex is the local-coefficient cochain complex. The simply connected case of $b_n\Rightarrow c_n$ gives the required cohomology isomorphisms of the cover complexes through degree $n$, and the lemma gives the desired twisted comparison on $X$ and $Y$. Thus the alternative proposed on p.~182 does work. The spherical-evaluation proof remains closer to the geometric reason the edge map vanishes.
\end{dcomment}


$\mathrm{c}_n \Longrightarrow \mathrm{a}_{n-1}$. On peut procéder par récurrence sur $n$. (NB L’enoncé a un sens même pour $n=1$, où il est prouvé, et même $\mathrm{c}_1 \Longrightarrow \mathrm{b}_1$, et $\mathrm{b}_1 \Longrightarrow \mathrm{a}_0$, plus généralement $\mathrm{b}_n \Longrightarrow \mathrm{a}_{n-1}$, immédiat par \texttt{[la]} suite exacte d’homotopie.) Donc on sait déjà, en plus de l’hypothèse $\mathrm{c}_n$, qu’on a $\mathrm{a}_{n-2}$, i.e. $Z$ est $(n-2)$-connexe. Il faut prouver qu’il est même $(n-1)$-connexe, i.e. $\pi_{n-1}(Z)=0$. Le cas $n=2$ demande % [sic] L’enoncé -> L’énoncé, missing accent on é
\msp{page 186}
peut-être une attention particulière, car on ne sait pas même encore que $Z$ est 1-connexe. Mais on sait que $\mathrm{c}_2$ implique $\mathrm{c}_1$, qui équivaut \texttt{[?]} à $\mathrm{b}_1$, donc on a
\[
\pi_1(X,x) \longrightarrow \pi_1(Y,y) \qquad \text{isomorphisme,}
\]
donc
\[
\pi_1(Z,x) \longrightarrow \pi_1(X,x) \qquad \text{\texttt{[est]} \textit{nul,}}
\]
donc
\[
\pi_2(Y,y) \longrightarrow \pi_1(Z,x) \qquad \text{\texttt{[est un]} épimorphisme,}
\]
ce qui implique tout au moins que $\pi_1(Z,x)$ est \textit{abélien.} Il faut le voir comme un $\pi$-module, et prouver sa nullité comme tel (moyennant $\mathrm{c}_2$).

\dgrblock{dgrb546}{Recovering the first missing homotopy group}
\begin{dcomment}{exposition}{p185}{Recovering the first missing homotopy group}
For the converse direction Grothendieck changes the use of the spectral sequence: instead of annihilating the edge map, he chooses coefficients in which the first possible fibre group can test itself. This is why all local systems must be admitted. A single constant coefficient group need not detect a nontrivial module with monodromy.

Assume $c_n$. The degree-$0$ and degree-$1$ conditions give $\pi_0X=\pi_0Y$ and identify $\pi_1X$ with $\pi_1Y=\pi$. Induction gives that the homotopy fibre $Z$ is $(n-2)$-connected. For $n\geq3$ its first possible nonzero group $A=\pi_{n-1}Z$ is abelian and Hurewicz identifies it with $H_{n-1}Z$. For $n=2$ one must first justify the same assertion differently: exactness gives a surjection $\pi_2Y\twoheadrightarrow\pi_1Z$, since $\pi_1X\to\pi_1Y$ is an isomorphism. Thus $\pi_1Z$ is abelian, and its abelianization map to $H_1Z$ is already an isomorphism. In both cases transport makes $A$ a $\pi$-module; the possible ambiguity in degree one is inner conjugation in an abelian group and therefore vanishes.

For the local system determined by a $\pi$-module $M$, the first nonzero fibre row gives
\[
\begin{split}
0\longrightarrow H^{n-1}(Y;M)&\longrightarrow H^{n-1}(X;f^*M)
 \longrightarrow\operatorname{Hom}_{\mathbf Z[\pi]}(A,M)\\
 &\longrightarrow H^n(Y;M)\longrightarrow H^n(X;f^*M).
\end{split}
\]
For $n=2$, the degree-one groups in this display use abelian coefficients, a case included in $c_n$. The first comparison is surjective and the last is injective. Exactness therefore forces $\operatorname{Hom}_{\mathbf Z[\pi]}(A,M)=0$. Every $\pi$-module occurs as a coefficient system on $Y$, so take $M=A$. The identity of $A$ must vanish, and hence $A=0$. This proves $c_n\Rightarrow a_{n-1}$, with the endpoint case checked rather than hidden under an appeal to Hurewicz.

This proof uses less than full $c_n$ at its final step: after the lower-degree induction, surjectivity in degree $n-1$ and injectivity in degree $n$ suffice. The full hypothesis is a convenient uniform formulation, not an optimal list of separate tests. The source has found the correct mechanism; the qualifications about coefficients and the degree-one fibre group complete its sketch.
\end{dcomment}


\printedpage{124}
% head: A. Grothendieck : Dérivateurs, Ch. XII, Édition des Universités de Montpellier II et Paris VII

Passons au cas général. On a pour $F$ localement constant abélien sur $Y$
\[
H^*(X,F_X) \Longleftarrow \underbrace{E_2^{pq} = H^p(Y,H^q(Z,F_Z))}_{\begin{array}{l}\scriptstyle = 0 \text{ si } 0<q<n-1\\ \scriptstyle \text{car } Z \text{ est } (n-2)\text{-connexe}\end{array}}, \qquad
\begin{tikzpicture}[baseline=0pt, x=0.6cm, y=0.6cm]
\draw[->] (0,0) -- (0,5);
\draw[->] (0,0) -- (6,0);
\node[left] at (0,4) {$n-1$};
\foreach \i in {0,...,4} { \fill (\i,4) circle (2.5pt); \fill (\i,0) circle (2.5pt); }
\node at (1,2) {$0$};
\node at (3,2) {$0$};
\end{tikzpicture}
\qquad ,
\]
D’où
\[
0 \longrightarrow H^{n-1}(Y,F) \longrightarrow H^{n-1}(X,F_X) \longrightarrow \underbrace{H^0(Y,H^{n-1}(Z,F_Z))}_{\begin{array}{l}\scriptstyle \mathrm{Hom}_\pi(\pi_{n-1}(Z,x),F_x)\\ \scriptstyle \text{par Hurewicz}\end{array}} \longrightarrow H^n(Y,F) \longrightarrow H^n(X,F_X).
\]
L’hypothèse $\mathrm{c}_n$ implique qu’on a
\msp{page 187}
\[
\mathrm{Hom}_\pi(\pi_{n-1}(Z,x),\underbrace{M}_{=\,F_x}) = 0
\]
et ceci \textit{quel que soit le $\pi$-module $M$,} puisque $\pi = \pi_1(X,x) \xrightarrow{\ \sim\ } \pi_1(Y,y)$. Cela implique (en prenant p. ex. $M = \pi_{n-1}(Z,x)$ !) que $\pi_{n-1}(Z,x)=0$. (Mais en faisant attention, dans le cas $n=2$, qu’on sait déjà que $\pi_1(Z,x)$ est abélien).

Donc le théorème rectifie page 178 a bien l’air vrai. Il est un peu triste, puisque dans chacune des implications % [sic] rectifie -> rectifié, past participle intended
\[
\mathrm{a}_n \ \Longrightarrow\ \mathrm{b}_n \ \Longrightarrow\ \mathrm{c}_n \ \Longrightarrow\ \mathrm{a}_{n-1}
\]
on perd un chouia, et en fin de compte on sort avec la magnifique implication
\[
\mathrm{a}_n \ \Longrightarrow\ \mathrm{a}_{n-1}.
\]
De toutes façons, il est $\pm$ clair que le “bon” localiseur $W_n$, celui qui doit définir les “$n$-types d’homotopie” comme objets de $\mathrm{Hot}_{W_n}$, est défini par la condition $\mathrm{a}_n$.

\section*{2\quad Critère de $W_n$-asphéricité.}

Je vais essayer de faire mieux quand $Y$ est réduit à un point. Cette fois, je

% ---- from Der-ChXIIc-30-03-03_p011-015.tex
\printedpage{125}
% head: A. Grothendieck : Dérivateurs, Ch. XII, Édition des Universités de Montpellier II et Paris VII
% foot: Version 30 mars 2003

\msp{page 188}
veux être strict, et donner (si je peux) des démonstrations complètes dans le cadre où je me place à présent, celui du modéliseur Cat, et en n'utilisant que ce que je sais y démontrer directement à présent, sans me référer à la théorie semi-simpliciale ou topologique.

\dgrblock{dgrb547}{Degree two and the localizer correction}
\begin{dcomment}{development}{p187}{Degree two and the localizer correction}
Two brief remarks in the manuscript deserve a more precise disposition. The suggestion of a degree-two counterexample is stronger than the examples support, and the choice of the fibre condition as the localizer confuses a useful generating class with its saturation. Both can be settled without changing the successful implication arguments.

\textbf{The degree-two converse.}\ours\ For maps of classifying spaces, the condition $c_2$ implies $b_2$. Consequently $b_2$ and $c_2$ are equivalent without a simple-connectivity hypothesis.

\textit{Proof.} The degree-one argument already identifies fundamental groupoids, so work on one component and write $R=\mathbf Z[\pi]$. Let
\[
 I=\operatorname{Hom}_{\mathbf Z}(R,\mathbf Q/\mathbf Z),
 \qquad (r\varphi)(s)=\varphi(sr).
\]
This is a left $R$-module and hence a permitted local system on $Y$. Evaluation at $1\in R$ gives a natural adjunction isomorphism
\[
 \operatorname{Hom}_R(C_*(\widetilde Y),I)
 \cong\operatorname{Hom}_{\mathbf Z}
                  (C_*(\widetilde Y),\mathbf Q/\mathbf Z),
\]
and likewise for $X$. Its inverse sends a homomorphism $\lambda$ to the map $c\mapsto(r\mapsto\lambda(rc))$. Since $\mathbf Q/\mathbf Z$ is injective as an abelian group, evaluation on homology gives
\[
 H^2(Y;I)\cong
 \operatorname{Hom}_{\mathbf Z}(H_2\widetilde Y,\mathbf Q/\mathbf Z).
\]
The isomorphism is natural in the lifted map $\widetilde f$. Hypothesis $c_2$ therefore says that the dual of $H_2(\widetilde f)$ is an isomorphism. Dualization into $\mathbf Q/\mathbf Z$ is exact and detects nonzero abelian groups: a nonzero element generates a cyclic subgroup with a nonzero homomorphism to $\mathbf Q/\mathbf Z$, and injectivity extends this homomorphism to the whole group. Applying this observation to kernel and cokernel shows that $H_2(\widetilde f)$ itself is an isomorphism. The universal covers are simply connected, so the natural Hurewicz isomorphisms $\pi_2\widetilde X\cong H_2\widetilde X$ and $\pi_2\widetilde Y\cong H_2\widetilde Y$ prove that $\pi_2f$ is an isomorphism.\hfill$\square$

The algebraic ingredients are the standard injective-cogenerator argument, checked in \citet[Tags 01D8 and 08WE]{StacksInjective2026}, together with \citet[Theorem~4.32]{Hatcher2002}. We did not find a source addressing footnote~112 specifically; this is our application of those familiar tools to its question. It is not a new form of the Hurewicz theorem. The same argument shows that the degree-$j$ part of $c_n$ detects the map on $H_j$ of universal covers. For $j=2$ these homology groups are exactly $\pi_2$; for $j=3$ the Hopf class in the earlier example explains why the conclusion changes.

Now consider the last sentence of p.~187. If $A\ne0$, the maps
\[
 *\xrightarrow{s}K(A,n+1)\xrightarrow{r}*
\]
satisfy $rs=1$. The fibre of $r$ is $n$-connected and the fibre of $rs$ is contractible, but the fibre of $s$ is $K(A,n)$, which is not $n$-connected. Thus the class defined by $a_n$ fails two-out-of-three. The class defined by $b_n$, equivalently the maps made into equivalences by $n$th Postnikov truncation, is the correct weak-equivalence class for $n$-types. The modern truncation adjunction is given in \citet[Proposition~5.5.6.18]{Lurie2009HTT}.

There is a charitable precise reading of Grothendieck's intuition.

\textbf{Saturation of the fibre class.}\ours\ In spaces, $a_n$ \emph{generates} the same localization as $b_n$.

\textit{Proof.} Every Postnikov map $X\to P_nX$ has $n$-connected fibres: it is an isomorphism on $\pi_i$ through $n$, and the target has zero $\pi_{n+1}$. Given a $b_n$-map $f:X\to Y$, its natural square with $P_nf:P_nX\to P_nY$ has vertical maps in $a_n$, and its lower map is a weak homotopy equivalence, hence also lies in $a_n$. Two-out-of-three puts $f$ in the saturation of $a_n$. Conversely, $b_n$ is saturated and contains $a_n$. This proves the equality.\hfill$\square$

The proof is a direct application of the standard truncation adjunction just cited; the formulation addresses the specific localizer sentence, without a priority claim. It leaves the literal definition on p.~187 marked as false. For $X\to *$, the two conditions agree already before saturation, so the next theorem needs no change of indexing.
\end{dcomment}


\textbf{Théorème 2} (?). \textit{Soit $X$ dans} Cat, \textit{$n \in \mathbf{N}$ un entier $\geq 1$. Alors les conditions suivantes sont équivalentes :}

\begin{itemize}
\item[\textit{$a_n$)}] \textit{$X$ est $0$-connexe, et si $x$ est dans $X$, $\pi_i(X,x) = 0$ pour $i \in [1,n]$ (condition qui ne dépend pas du choix de $x$), i.e. les $\Omega^i(X,x)$ pour $i \in [1,n]$ sont $0$-connexes.}
\item[\textit{$b_n$)}] \textit{Pour tout groupe $G$, $G \longrightarrow H^0(X,G_X)$ est un isomorphisme, et $H^1(X,G_X) = 0$. De plus, si $G$ est abélien, on a aussi $H^i(X,G_X) = 0$, pour $i \in [1,n]$.}
\end{itemize}

\textsc{Démonstration}. Pour $n = 1$ cette équivalence est $\pm$ tautologique et a été rappelée (sous une forme légèrement différente) au début (p.~71). On procède par récurrence. On peut donc supposer que $n \geq 2$, que $X$ satisfait $\mathrm{a}_{n-1}$, $\mathrm{b}_{n-1}$, et prouver qu'alors $\mathrm{a}_n$, $\mathrm{b}_n$ sont équivalentes, i.e. $\pi_n(X,x) = 0$ équivaut à $H^n(X,M) = 0$ pour tout groupe abélien $M$.

\msp{page 189}
On utilise la $W_\infty$-fibration (parfaite)
\[
\begin{tikzcd}[row sep=large]
E = \Ch(X;x,-) \arrow[d] & E\ W_\infty\text{-asphérique} \\
X &
\end{tikzcd}
\]
de fibre $Z$ en $x$ égale à $\Omega(X,x)$.

\dgrblock{dgrb548}{Why the path fibration is the right final test}
\begin{dcomment}{context}{p188}{Why the path fibration is the right final test}
Grothendieck now asks the categorical path construction to reproduce an argument whose classical power came from the contractible total space of based paths. Serre's thesis developed precisely this exchange between the homology of a space, that of its loop space, and their homotopy groups; see \citet[Chapter~IV, \S5, Proposition~7]{Serre1951}. The point of repeating the argument in $\Cat$ is not a new calculation of the groups. It is to show that the categorical paths already carry enough geometry to recognize the ordinary weak homotopy type, before appealing to a comparison with topological paths.

For $E=\Ch(X;x,-)\to X$, the total category is $W_\infty$-aspherical and the fibre is $Z=\Omega(X,x)$. Under the induction hypothesis, both base and fibre have long vanishing ranges. The spectral sequence then has only its bottom row and one possible first fibre row. The vanishing of the total cohomology forces those two rows to communicate by transgression. This is the reason the proof over a point is sharper than the preceding relative theorem: there is no independent cohomology of a nontrivial target that can absorb the first obstruction.

The condition in degree zero should be read literally for every group. Taking a nontrivial group $G$, the diagonal map $G\to G^{\pi_0X}$ is bijective only when $\pi_0X$ has one element. It excludes the empty category as well. Once $X$ is connected, the condition $H^1(X,G_X)=*$ for every group $G$ forces $\pi_1X=1$: choose $G=\pi_1X$ and the conjugacy class of its identity representation. This is why abelian coefficients alone would not prove the initial case; a perfect nontrivial fundamental group could otherwise survive.
\end{dcomment}


Supposons $\mathrm{a}_n$) pour $X$, donc $a_{n-1}$) pour $\Omega(X,x)$. Par hypothèse de récurrence, on a donc
\[
H^i(Z,M) = 0 \qquad \text{pour } i \in [1,n-1],
\]
$M$ groupe abélien. Donc dans la suite spectrale de Leray,
\[
H^*(E,M) \Longleftarrow E_2^{pq} = H^p(X,H^q(Z,M))
\]
on a
\[
E_2^{pq} = 0 \qquad \text{si } q \in [1,n-1], \qquad
\begin{tikzpicture}[baseline=0]
\draw[->] (0,0) -- (2.2,0);
\draw[->] (0,0) -- (0,2.1);
\foreach \x in {0,0.35,0.7,1.05,1.4} {
  \fill (\x,0) circle (1.5pt);
  \fill (\x,1.4) circle (1.5pt);
}
\node[left] at (-0.05,1.4) {$n$};
\node at (0.35,0.7) {$0$};
\node at (1.05,0.7) {$0$};
\end{tikzpicture}
\quad .
\]

\printedpage{126}
% head: A. Grothendieck : Dérivateurs, Ch. XII, Édition des Universités de Montpellier II et Paris VII
% foot: Version 30 mars 2003

On a donc la suite exacte
\[
0 \longrightarrow H^n(X,M) \longrightarrow H^n(E,M)
\]
et comme $E$ est $W_\infty$-asphérique, on a $H^n(E,M) = 0$, donc $H^n(X,M) = 0$.

Inversement, supposons $\mathrm{b}_n$), prouvons $\mathrm{a}_n$). On sait déjà qu'on a $\mathrm{a}_{n-1}$), donc $\mathrm{a}_{n-2}$ pour $\Omega(X,x)$, ce qui implique
\[
H^i(Z,M) = 0 \qquad \text{pour } i \in [1,n-2],
\]
donc
\[
E_2^{pq} = 0 \qquad \text{si } q \in [1,n-2].
\]
On a donc la suite exacte

\msp{page 190}
\[
0 \longrightarrow \underbrace{H^{n-1}(X,M)}_{=\,0} \longrightarrow \underbrace{H^{n-1}(E,M)}_{=\,0} \longrightarrow H^0(X,H^{n-1}(Z,M)) \longrightarrow \underbrace{H^n(X,M)}_{=\,0} \longrightarrow H^n(E,M),
\]
donc
\[
H^0(X,\ \underbrace{H^{n-1}(Z,M)}_{\begin{array}{l}\text{système local}\\ \textit{constant}\text{ sur } X,\\ \text{car } X \text{ est 1-connexe}\end{array}}\ ) = 0,
\]
donc
\[
H^{n-1}(Z,M) = 0,
\]
ce qui, par hypothèse de récurrence, implique $\pi_{n-1}(Z,z) = 0$, donc $\pi_n(X,x) = 0$, q.e.d.

\textbf{Corollaire.} \textit{Pour que $X$ soit $W_\infty$-asphérique, i.e. pour qu'on ait $G \xrightarrow{\ \sim\ } H^0(X,G)$, $H^1(X,G) = 0$} \texttt{[pour]} \textit{tout groupe} \texttt{[}$G$\texttt{]}\textit{, et $H^i(X,G) = 0$ si $i \geq 2$ pour $G$ abélien, il faut et il suffit que} $1^\circ$) \textit{$X$ soit} $0$\textit{-connexe, et ($x$ étant un point de $X$)} $2^\circ$) \textit{les $\Omega^i(X,x)$} \texttt{[soient]} $0$\textit{-connexes pour $i \geq 1$.}

\msp{page 191}

\dgrblock{dgrb549}{The loop criterion, including its degree-two endpoint}
\begin{dcomment}{exposition}{p190}{The loop criterion, including its degree-two endpoint}
Theorem~2 is proved by the displayed induction, with one small verification at its lower endpoint. Suppose first that $X$ is $n$-connected. Then $Z=\Omega X$ is $(n-1)$-connected, so the induction hypothesis annihilates $H^q(Z;M)$ for $1\leq q<n$. In the path-fibration spectral sequence this gives an injection
\[
 H^n(X;M)\hookrightarrow H^n(E;M)=0.
\]
The lower-degree groups have already vanished, and the degree-zero and torsor conditions follow from connectedness and simple connectivity. This proves the forward implication exactly in the form stated.

Conversely, assume the cohomological conditions through degree $n$. The induction hypothesis makes $X$ $(n-1)$-connected and $Z$ $(n-2)$-connected. For $n\geq2$, $X$ is simply connected, so the fibre-cohomology local systems on $X$ are constant. The low-degree exact sequence for $E\to X$ gives
\[
 H^0(X;\mathcal H^{n-1}(Z;M))=0,
 \qquad\text{hence}\qquad H^{n-1}(Z;M)=0
\]
for every abelian $M$. If $n\geq3$, $Z$ is simply connected. Its nonabelian degree-one conditions and all its lower abelian cohomology conditions are already satisfied; the induction hypothesis applied to $Z$ therefore gives $\pi_{n-1}Z=0$.

For $n=2$, one cannot immediately apply that same sentence: $Z$ is only known to be connected, and the displayed computation gives vanishing of $H^1(Z;M)$ only for abelian groups $M$. The loop multiplication supplies the missing information. The two compositions of loops in the connected loop space, concatenation of based loops and the multiplication induced from the loop-space multiplication, have the same unit and satisfy the interchange law. The Eckmann--Hilton argument makes $\pi_1Z$ abelian. Thus
\[
 H^1(Z;M)=\operatorname{Hom}(\pi_1Z,M),
\]
and choosing $M=\pi_1Z$ forces $\pi_1Z=0$. Equivalently, this is the familiar fact that $\pi_1\Omega X=\pi_2X$ is abelian. It closes the induction without requiring an additional theorem about nonabelian cohomology of $Z$.

The corollary follows by allowing every finite $n$. Connectedness of $X$ together with connectedness of every iterated loop category says that all homotopy groups of $BX$ vanish. Since a realization of a nerve has CW type, Whitehead's theorem makes $BX$ weakly equivalent, indeed homotopy equivalent, to a point; see \citet[Theorem~4.5]{Hatcher2002}. Here ``aspherical'' means weakly contractible, not merely the modern usage $\pi_i=0$ for $i>1$. For example $B\mathbf Z\simeq S^1$ has no higher homotopy groups, but its first loop category has components indexed by $\mathbf Z$, so it fails this criterion immediately.

The later minimal-localizer theorem identifies this test with Grothendieck's original axiomatic weak equivalences. Its local characterization also makes the loop method visible in the final proof: \citet[Theorem~2.3.6, especially 2.3.6.3--2.3.6.4]{Cisinski2004Minimal} passes from a weakly aspherical object to its loops and then uses vanishing of homotopy groups. Thus the categorical construction has reached its intended target. The distinction between different choices of categorical paths has become a question of canonical comparison, which is exactly the subject of the following pages.
\end{dcomment}


\section*{3\quad Sur la philosophie des fibres homotopiques.}

Le critère d'asphéricité (pour $W_\infty$) que je viens d'obtenir, grâce aux foncteur $\Ch(X)$ (ou au choix, $\Ch_\infty(X)$) et ses dérivés $\Omega(X,x)$ et les itérés, enlève tout doute (s'il pouvait en rester encore) que les `espaces de chemins' que j'ai construit (assez péniblement!) sont bel et bien `les bons'. Je voudrais maintenant me rendre compte de façon la plus précise possible de leur `canonicité', dans quelle mesure ces constructions sont indépendant d'un choix (p.ex. celui entre la théorie $\Ch$, et $\Ch_\infty$). % [sic] aux foncteur -> au foncteur, number agreement

Tout revient, finalement, à des questions de factorisation d'une flèche quelconque
\[
f : X \longrightarrow Y
\]
\printedpage{127}
% head: A. Grothendieck : Dérivateurs, Ch. XII, Édition des Universités de Montpellier II et Paris VII
% foot: Version 30 mars 2003
de Cat en
\[
X \xrightarrow{\ i_f\ } \bar{X} \xrightarrow{\ \bar{f}\ } Y,
\]
avec
\[
i_f \in W, \qquad \bar{f} \in W\text{-fib}
\]
(le fait que $\bar{f}$ puisse même être choisie $W_\omega$-parfait, ne me paraissant pas essentiel ici).

C'est bien de cela qu'il s'agit pour construire

\msp{page 192}
un objet `chemins dans $X$'
\[
\Ch(X),
\]
en factorisant l'application diagonale
\[
X \xrightarrow{\ \mathrm{diag}_X\ } X \times X,
\]
et aussi quand on construit $\Omega(X,x)$ comme fibre en $x$ d'une $W$-fibration
\[
E(X,x) \longrightarrow (X,x)
\]
qui factorise l'inclusion
\[
\{x\} \longrightarrow X
\]
(donc $E(X,x)$ étant $W$-asphérique).

\dgrblock{dgrb550}{From a model for paths to an intrinsic operation}
\begin{dcomment}{context}{p191}{From a model for paths to an intrinsic operation}
The question of canonicity belongs to Grothendieck's original reason for introducing derivators. In \emph{Pursuing Stacks}, \S69, he recalls the inadequacy of a single derived category for some of the diagrammatic constructions arising from the SGA~5 trace formalism. He proposes retaining derived categories for all diagram shapes, with their variance in the indexing category. The two three-object shapes for spans and cospans already encode homotopy fibres and cofibres \citep[\S69, printed pp.~1--3]{Grothendieck1983}. The move on these pages, from a particular category of paths to a property of localized slices, is a local instance of that larger programme.

The letter to Thomason of 2 April 1991 makes the distinction especially explicit. Its first numbered discussion treats $W$ as the essential homotopical datum and fibrations or cofibrations as additional means of proving existence and calculating operations; it then describes the intrinsic class of pullback-preserving $W$-fibrations \citep[point~1, pp.~3--4]{Grothendieck1991Letter}. This is a retrospective clarification of the position being explored here, not evidence that every provisional axiom of the present pages has already been checked. In particular, the retract qualification discussed below remains necessary under the literal two-out-of-three assumptions. The calculation of paths serves the intrinsic operation by providing enough useful representatives on which the operation becomes ordinary pullback.
\end{dcomment}


On peut être plus ou moins exigeant, et regarder, soit pour un `but' $Y$ (ou plutôt un $S$) fixé dans Cat, une factorisation fonctorielle
\[
C : \Cat/S \longrightarrow \Fib_W S \hookrightarrow \Cat/S
\]
du morphisme structural $X \longrightarrow S$ d'un objet, en
\[
X \xrightarrow{\ i_X\ } C(X) \longrightarrow S,
\]
i.e. on se donne le foncteur $C$, \textit{plus}

\msp{page 193}
le morphisme fonctoriel
\[
i_C : \mathrm{id}_{\Cat/S} \longrightarrow C
\]
avec
\[
(i_X : X \longrightarrow C(X)) \in W \qquad \forall\ X \text{ dans } \Cat/S,
\]
soit être encore plus exigeant et faire varier en même temps la source et le but de $f : X \longrightarrow Y (= S)$, et demander un foncteur
\[
\begin{aligned}
C : \Fl(\Cat) \quad &\longrightarrow \quad \Cat \\
f \quad &\longmapsto \quad C(f),
\end{aligned}
\]
\printedpage{128}
% head: A. Grothendieck : Dérivateurs, Ch. XII, Édition des Universités de Montpellier II et Paris VII
% foot: Version 30 mars 2003
avec des morphismes fonctoriels en $f$
\[
\underbrace{s(f)}_{=\,X} \xrightarrow{\ i_f\ } \underbrace{C(f)}_{=\,\bar{X}} \xrightarrow{\ p_f\ } \underbrace{b(f)}_{=\,Y}
\]
($i$, $p$ des flèches dans $\Hom(\Fl(\Cat),\Cat)$), avec
\[
i_f \in W, \qquad p_f \in W\text{-fib}
\]
pour toute flèche $f$ dans Cat.

\dgrblock{dgrb551}{A replacement must remember its comparison map}
\begin{dcomment}{exposition}{p192}{A replacement must remember its comparison map}
Grothendieck now asks what the lengthy construction of paths has bought him beyond a convenient calculation. The preceding connectivity criterion identifies its homotopical effect; it does not yet identify the construction independently of its presentation. The change of viewpoint is to regard a path category as a factorization of a structural map. A factorization of the diagonal gives paths, and a factorization of a point inclusion gives based paths. Once the second factor has the required base-change property, its ordinary fibres can serve as homotopy fibres. The same question can therefore be asked without mentioning intervals or zigzags.

For a fixed base put $\mathcal A=\Cat/S$, and let $\mathcal F\subseteq\mathcal A$ be the full subcategory of the permitted fibrations over $S$. The data are a functor $C:\mathcal A\to\mathcal F$ and a natural transformation $i:\id_{\mathcal A}\to jC$, where $j$ is inclusion and every $i_X$ belongs to $W_S$. The transformation is part of the data. An arbitrary choice of an equivalent object for every $X$ would neither provide the arrows $C(f)$ nor specify which comparison is to be called canonical. Naturality supplies both controls at once. Indeed,
\[
 C(f)i_X=i_Yf
\]
and two-out-of-three show that $f\in W_S$ if and only if $C(f)\in W_S$. Thus the homotopicality of $C$ is a consequence of its coaugmentation, not an extra axiom.

There are two levels of functoriality in the text. A replacement on $\mathcal A$ controls variation of the object over one base. A replacement on the arrow category also controls maps between bases, although it need not commute with ordinary pullbacks. The latter distinction matters: a functorial factorization of arrows does not, merely by being functorial, identify $C(X\times_S T)$ with $C(X)\times_S T$. Derived base change will be obtained from a localization argument. It is not asserted to be an identity between these raw constructions.

This is the same elementary geometry that underlies Brown's mapping-path factorization. In a category of fibrant objects, a factorization of the diagonal of $Y$ yields a factorization of $f:X\to Y$ by pulling that path object back along $f$. Brown requires no functorial choice of path objects for this argument \citep[Part I, \S1, pp.~420--422]{Brown1973}. Grothendieck deliberately asks for a more explicit choice here so that comparison maps can be written in the original category. The two approaches answer different parts of the same problem: existence of the derived construction and a convenient formula for computing it.
\end{dcomment}


\msp{page 194}
Dans les deux points de vue, dans quel sens a-t-on \textit{unicité} d'une telle construction?

Prenons d'abord le premier point de vue, factorisation des flèches à but fixé $S$, donc \texttt{[c'est une]} question dans $\Cat/S$. J'ai oublié une propriété importante, mais elle résulte de ce qui a été dit : le foncteur
\[
\begin{tikzcd}[column sep=small]
X \longmapsto C(X) : & \Cat/S \arrow[rr] \arrow[dr, dotted] & & \Cat/S \\
& & \Fib_W S \arrow[ur, dotted, hook] &
\end{tikzcd}
\]
\textit{transforme} \texttt{[les]} \textit{flèches} \texttt{[de]} \textit{$W$ en itou}, comme on voit sur le diagramme commutatif
\[
\begin{tikzcd}[row sep=large, column sep=large]
X \arrow[r, "i_X"] \arrow[d, "f"'] & C(X) \arrow[d, "C(f)"] \\
Y \arrow[r, "i_Y"] & C(Y),
\end{tikzcd}
\]
où $i_X$, $i_Y$ sont dans $W$, donc
\[
f \in W \qquad \Longleftrightarrow \qquad C(f) \in W.
\]
Soit
\[
(C',i')
\]
un deuxième foncteur, $\Cat/S \longrightarrow \Cat/S$,

\msp{page 195}
muni d'un morphisme
\[
i' : \mathrm{id}_{\Cat/S} \longrightarrow C'
\]
qui est dans $W$ argument par argument, et tel que $C'$ se factorise par $\Fib_W S$. On peut comparer $C$ et $C'$ moyennant $C \circ C'$, car
\[
\begin{tikzcd}[column sep=small]
& C \circ C' & \\
C = C \circ \mathrm{id} \arrow[ur, "C\ast i'"] & & C' = \mathrm{id} \circ C', \arrow[ul, "i\ast C'"']
\end{tikzcd}
\]
\printedpage{129}
% head: A. Grothendieck : Dérivateurs, Ch. XII, Édition des Universités de Montpellier II et Paris VII
% foot: Version 30 mars 2003
i.e. on a, pour tout $X$ dans $\Cat/S$,
\[
\begin{tikzcd}[column sep=small]
& C(C'(X)) & \\
C(X) \arrow[ur, "C(i'_X)"] & & C'(X), \arrow[ul, "i_{C'(X)}"']
\end{tikzcd}
\]
le deux flèches étant dans $W$ (celle de gauche parce que $C(W) \subseteq W$). Ce sont des flèches dans $\Fib_W S$, et si on suppose que $W$ satisfait W(4b), elles induisent des $W$-équivalences en chaque fibre. Cela montre donc déjà % [sic] le deux flèches -> les deux flèches, typo

\msp{page 196}
qu'à isomorphisme (peut-être non-canonique) près, le foncteur composé défini par $C$,
\[
\begin{tikzcd}
X \arrow[r, mapsto] & C(X) \arrow[r, mapsto] & (s \longmapsto \mathrm{hot}_W(C(X)_s)) \\
\Cat/S \arrow[r] \arrow[rr, bend right, "\varphi_C"'] & \Fib_W S \arrow[r] & \Hom(S^\circ,\mathrm{Hot}_W),
\end{tikzcd}
\]
ne dépend pas du choix particulier de $C$, i.e. deux tels foncteurs `fibres homotopiques' (correspondant à deux choix $C$, $C'$) sont toujours isomorphes. Mais au lieu d'utiliser $C \circ C'$ pour définir un tel isomorphisme, on aurai pu utiliser $C' \circ C$, ou quelque autre procédé. On aimerait comprendre s'il y a unicité véritable, une caractérisation intrinsèque de ces isomorphismes $\varphi_C \xrightarrow{\ \sim\ } \varphi_{C'}$, indépendamment de telle construction ou de telle autre. % [sic] aurai pu -> aurait pu, typo

Je présume que $C \circ C'$ est lui aussi muni de
\[
\mathrm{id}_{\Cat/S} \longrightarrow C \circ C',
\]

\msp{page 197}
et même a priori de deux façons possible
\[
\begin{tikzcd}[column sep=large, row sep=large]
& C(C'(X)) & \\
C(X) \arrow[ur, "C(i'_X)"] & & C'(X) \arrow[ul, "i_{C'(X)}"'] \\
& X, \arrow[ul, "i_X"] \arrow[ur, "i'_X"'] \arrow[uu, dotted] &
\end{tikzcd}
\]
soit par $C(i'_X) \circ i_X$, soit par $i_{C(X)} \circ i'_X$, heureusement elles sont égales, car pour toute flèche $X \xrightarrow{\ i'\ } X'$ dans $\Cat/S$, le carré % [sic] i_{C(X)} -> i_{C'(X)}, the diagram above has i_{C'(X)}

\dgrblock{dgrb552}{Two replacements and one normalized comparison}
\begin{dcomment}{exposition}{p196}{Two replacements and one normalized comparison}
The composite replacement is a useful device because it compares two constructions without asking either to map directly to the other. Write the two coaugmentations as $i:\id\to C$ and $i':\id\to C'$. For every $X$ the diagram gives
\[
 C(X)\xrightarrow{C(i'_X)}CC'(X)
       \xleftarrow{i_{C'(X)}}C'(X).
\]
Both arrows are weak equivalences: the first because $C$ preserves $W_S$, and the second because it is a component of $i$. Their passage to fibrewise equivalences uses the additional fibre axiom invoked in the text. This last step must be retained. A weak equivalence of total objects over a base need not induce weak equivalences of their ordinary fibres when the objects have not been suitably replaced.

There is already a precise answer to the question about the order $CC'$ versus $C'C$. Let $q:\mathcal A\to\mathcal A[W_S^{-1}]$ be localization. The required comparison between the two localized endofunctors is
\[
 \theta_{C,C'}(X)=q(i'_X)q(i_X)^{-1}.
\]
It is the unique comparison compatible with the specified maps from $q(X)$. Naturality of $i$, applied to $i'_X$, gives
\[
 C(i'_X)i_X=i_{C'(X)}i'_X.
\]
After localization this proves
\[
 q(i_{C'(X)})^{-1}q(C(i'_X))
       =q(i'_X)q(i_X)^{-1}.
\]
The common-refinement construction therefore gives exactly the normalized comparison. Reversing the order of the replacements gives the same answer. For three replacements the cocycle identity is literal:
\[
 \theta_{C',C''}\theta_{C,C'}=\theta_{C,C''}.
\]
The apparent ambiguity disappears as soon as compatibility with the coaugmentations is made explicit.

When this calculation is transported to the category of fibres, one uses the full faithfulness of the localization comparison proved below. In that category the formula is not an instruction to apply an ordinary fibre functor to $i_X$: its source $X$ need not be fibrant. Instead, it characterizes the comparison after passing through the equivalence of localized categories. This is the categorical content of the word ``canonical'' here. It does not say that the underlying functors admit no other automorphisms after their specified comparison maps have been forgotten.
\end{dcomment}


% ---- from Der-ChXIIc-30-03-03_p016-020.tex
\printedpage{130}
% head: A. Grothendieck : Dérivateurs, Ch. XII, Édition des Universités de Montpellier II et Paris VII

\[
\begin{tikzcd}[column sep=small, row sep=small]
 & C(X') & \\
C(X) \arrow[ur, "C(i')"] & & X' \arrow[ul, "i_{X'}"'] \\
 & X \arrow[ul, "i_X"] \arrow[ur, "i'"'] &
\end{tikzcd}
\]

est commutatif (par fonctorialité de $i_X : X \longrightarrow C(X)$), et on applique ceci au cas de $X' = C'(X)$, $i' = i'_X$.

Cette remarque nous encourage à faire un saut dans l'abstraction, et d'introduire la catégorie\msp{page 198}
\[
\mathcal{C} = \mathcal{C}_S = \begin{cases}
\text{catégorie des couples } (C,i), \text{ avec} \\
C : \Cat/S \longrightarrow \Cat/S \text{ se factorisant par } \Fib_W S, \\
i : \id_{\Cat/S} \longrightarrow C \text{ tel que } i_X \in W \ \forall\ X \text{ dans } \Cat/S,
\end{cases}
\tag*{$(*)$}
\]
une flèche
\[
(C,i) \longrightarrow (C',i')
\]
étant une flèche
\[
C \xrightarrow{u} C' \qquad \text{(dans } \Hom(\Cat/S, \Cat/S))
\]
rendant commutatif le carré \texttt{[lire : le diagramme]}
\[
\begin{tikzcd}[column sep=small]
 & \id_{\Cat/S} \arrow[dl, "i"'] \arrow[dr, "i'"] & \\
C \arrow[rr, "u"] & & C'.
\end{tikzcd}
\]

En d'autres termes, $\mathcal{C}$ est défini comme une sous-catégiorie pleine de % [sic] sous-catégiorie -> sous-catégorie, misspelled word
\[
\id_{\Cat/S}\backslash\Hom(\Cat/S, \Cat/S),
\]
à savoir celle dont les objets $(C,i)$ satisfont les deux conditions $(*)$ plus haut (impliquant l'une et l'autre le localiseur $W$).

\msp{page 199}Ces conditions impliquent ceci

\begin{enumerate}
\item[a)] $\forall\ X$ dans Cat/$S$, $(u_X : C(X) \longrightarrow C'(X)) \in W$, puisque
\[
\begin{tikzcd}[column sep=small]
 & X \arrow[dl, "i_X\in W"'] \arrow[dr, "i'_X\in W"] & \\
C(X) \arrow[rr, "u_X"] & & C'(X),
\end{tikzcd}
\]
\printedpage{131}
% head: A. Grothendieck : Dérivateurs, Ch. XII, Édition des Universités de Montpellier II et Paris VII
et

\item[b)] si $W$ satisfait W(4b), de sorte que $u_X \in W^{\mathrm{f}}_S$, alors ce définit un \textit{isomorphisme} entre $\varphi_C(X)$ et $\varphi_{C'}(X)$ % [sic] ce définit -> ceci définit, missing word
\[
\begin{tikzcd}[row sep=small]
s \arrow[r, mapsto] & \hot_W(C(X)_s) \arrow[d, "\wr\, u_{X,s}"] \\
s \arrow[r, mapsto] & \hot_W(C'(X)_s).
\end{tikzcd}
\]
\end{enumerate}

Mieux encore, dans ce cas, le foncteur composé
\[
\mathcal{C} \times \Cat/S \longrightarrow \Fib_W S \longrightarrow \HOT^{\mathrm{lc}}_W(S)
\]
donne naissance à
\[
\mathcal{C} \longrightarrow \Hom(\Cat/S, \HOT^{\mathrm{lc}}_W(S))
\]
qui transforme toute flèche de $\mathcal{C}$ en \texttt{[un]} \textit{isomorphisme,} donc se factorise par le groupoïde fondamental de $\mathcal{C}$,
\[
\Pi\mathcal{C} \longrightarrow \Hom(\Cat/S, \HOT^{\mathrm{lc}}_W(S)).
\]

\msp{page 200}Quant au foncteur envisagé tantôt
\[
\begin{aligned}
\mathcal{C} &\quad\longrightarrow\quad \underbrace{\Hom\,(\Cat/S, \Hom(S^\circ, \Hot_W))}_{\simeq\Hom((\Cat/S)\times S^\circ, \Hot_W)} \\
C &\quad\longmapsto\quad ((X,s) \longmapsto \hot_W(C(X)_s))\,,
\end{aligned}
\]
il se déduit du précédent en composant avec
\[
\begin{tikzcd}[column sep=small]
\HOT_W(S) \arrow[rr] \arrow[dr, dotted] & & \Hom(S^\circ, \Hot_W), \\
 & \Hom(\Pi S^\circ, \Hot_W) \arrow[ur, dotted] &
\end{tikzcd}
\]
plus exactement, avec le foncteur qui s'en déduit
\[
\Hom(\Cat/S, \HOT_W(S)) \longrightarrow \Hom(\Cat/S, \Hom(S^\circ, \Hot_W)).
\]

On voit poindre ici un principe plausible de canonicité d'isomorphismes, en ce qui concerne les foncteurs
\[
C_W : \Cat/S \longrightarrow \HOT^{\mathrm{lc}}_W(S)
\]
définis par les $C \in \mathrm{Ob}\,\mathcal{C}$, par la relation
\[
\Pi\mathcal{C} \approx \ \text{catégorie ponctuelle, i.e. } \mathcal{C} \text{ est 1-connexe.}
\]

De plus, les arguments utilisés jusqu'ici sont de nature à tel point formelle, que les aspects spécifiques du Modéliseur Cat
\printedpage{132}
% head: A. Grothendieck : Dérivateurs, Ch. XII, Édition des Universités de Montpellier II et Paris VII
\msp{page 201}semblent n'y jouer pratiquement aucun rôle, ou un rôle très effacé.

\dgrblock{dgrb553}{The category of choices is contractible}
\begin{dcomment}{exposition}{p200}{The category of choices is contractible}
Grothendieck's category $\mathcal C$ packages the compatibility just used. Its arrows are natural transformations of replacement functors that commute with their coaugmentations. Every such arrow is pointwise in $W_S$, by two-out-of-three. Hence any construction which sends these weak equivalences to isomorphisms factors through $\Pi\mathcal C$. The remaining question is whether this groupoid has exactly one arrow between any two objects. The natural common refinements give more: provided a replacement exists, the nerve of $\mathcal C$ is contractible.

Here is a direct proof. Fix $(R,r)\in\mathcal C$ and define
\[
 T(C,i)=\bigl(CR,(Cr)i\bigr).
\]
The displayed unit is also $(iR)r$, by naturality of $i$. The composite $CR$ takes values in $\mathcal F$, and its coaugmentation is a weak equivalence. Thus $T$ is an endofunctor of $\mathcal C$. On an arrow $u:C\to D$ put $T(u)=uR$. There are natural transformations
\[
 \id_{\mathcal C}\xrightarrow{\alpha}T
       \xleftarrow{\beta}\operatorname{const}_{(R,r)},
 \qquad \alpha_C=Cr,\quad\beta_C=iR.
\]
Both are arrows under the identity functor because $(Cr)i=(iR)r$. To check their naturality, use, respectively, naturality of $u$ at $r_X$ and the equation $ui=j$ for a map $(C,i)\to(D,j)$. These are equalities before localization.

A natural transformation $F\to G$ gives a functor $\mathcal D\times[1]\to\mathcal E$, hence a homotopy between the maps of nerves induced by $F$ and $G$; this is the elementary nerve calculation used in \citet[\S1, Proposition~2]{Quillen1973}. The two transformations above consequently homotope the identity of $|N\mathcal C|$ to a constant map. With the usual change of universe if necessary, $|N\mathcal C|$ is contractible. In particular, $\Pi\mathcal C$ is equivalent to the terminal groupoid. Any two zigzags of comparisons between the same replacements induce the same isomorphism in a target in which all arrows of $\mathcal C$ become invertible.

This argument also locates the imprecision in the word \emph{filtrante} on the next page. A common target for two objects is only one of the conditions for a filtered category; parallel arrows must also become equal after postcomposition. The common-refinement square does not establish that second condition. Nor do the natural maps into a binary product, even with all the displayed coherence, establish it in general. For example, the category of nonempty finite ordered sets and order-preserving injections has ordered concatenation and natural inclusions from both factors. It is not filtered: the two injections from a singleton into a two-element ordered set remain distinct after every injection. Nevertheless its nerve is contractible by the same argument. This example refutes the inference being used, not a separately proved assertion that this particular $\mathcal C$ fails to be filtered. Contractibility supplies precisely the canonicity needed here, so no strict equalization claim is required.
\end{dcomment}


Le carré de la page 197 montre surtout que la catégorie $\mathcal{C}$ est filtrante - $(C,i)$, $(C,i')$ s'envoient l'un et l'autre dans $(C'',i'')$, avec $C'' = C \circ C'$, et $i'' = (C * i') \circ i = (i * C') \circ i'$. Il y a même une loi de composition dans $\mathcal{C}$, nullement commutative ni unitaire c'est sûr, mais sûrement associative, et (si cette composition est notée $*$) avec des morphismes fonctoriels % [sic] (C,i') -> (C',i'), the second couple of the following formula is (C',i')
\[
\xi \xrightarrow{\alpha_{\xi,\xi'}} \xi * \xi' \xleftarrow{\beta_{\xi,\xi'}} \xi'
\]
qui assurent le caractère filtrant croissant de $\mathcal{C}$. Ces flèches semblent remplacer en quelque sorte l'unité bilatère absente, et on s'attend à ce qu'elles aient les mêmes excellentes propriétés que si on avait une telle unité, concernant les composés doubles $(\xi * \xi') * \xi'' \simeq \xi * (\xi' * \xi'')$,

\msp{page 202}
\[
(\xi * \xi') * \xi'' \xrightarrow{\sim} \xi * (\xi' * \xi'').
\]

Il y a deux façons d'envoyer un produit binaire
\[
\xi * \xi', \qquad \xi' * \xi'', \qquad \xi * \xi''
\]
dans le produit ternaire, et bien elles doivent dans chacun des trois cas, être égales. Et de même il y a, ceci admis, encore deux façons d'envoyer $\xi$ dans le produit ternaire, en factorisant soit par $\xi * \xi'$, soit par $\xi * \xi''$, et de même pour les deux autres foncteurs $\xi'$, $\xi''$; et là encore, ces deux façons doivent être égales, dans chacun de ces trois cas - donc $3 + 3 = 6$ égalités en tout à écrire et à vérifier - on a vu pire!

La philosophie calculatoire devrait être celle-ci : Si on a un produit $n$-uple
\[
\xi_1 * \xi_2 * \cdots * \xi_n
\]
(où je me dispense de mettre des parenthèses \dots), pour tout produit $m$-uple partiel $(m \in [1, n-1])$\msp{page 203}
\[
\xi_{i_1} * \xi_{i_2} * \cdots * \xi_{i_m} \qquad\qquad (i_1 < i_2 < \cdots < i_m),
\]
en utilisant à tour de bras les isomorphismes d'associativité et les $\alpha_{\xi,\xi'}$, $\beta_{\xi,\xi'}$, il n'y a pourtant qu'une seule façon de l'envoyer dans le produit $n$-uple. J'ai bien l'impression que ça doit être vrai. (Est-ce que ça implique déjà que $\mathcal{C}$ est 1-connexe?)

En fait, il y a bien une unité bilatère, mais elle n'est pas dans $\mathcal{C}$, mais dans la catégorie plus grande
\[
\id_{\Cat/S}\backslash\Hom(\Cat/S, \Cat/S),
\]
qui elle, comme toute catégorie
\[
\id_{\mathcal{X}}\backslash\Hom(\mathcal{X}, \mathcal{X}),
\]
\printedpage{133}
% head: A. Grothendieck : Dérivateurs, Ch. XII, Édition des Universités de Montpellier II et Paris VII
est une catégorie monoïdale, i.e. avec produit associatif \textit{et} unitaire. Les flèches $\alpha_{\xi,\xi'}$, $\beta_{\xi,\xi'}$ sont les flèches qui résultent de la présence de cette unité (vu que $\mathcal{C}$ est une sous-catégorie pleine), et alors il doit devenir clair qu'on a les compatibilités dites.

\msp{page 204}Mais même une catégorie monoïdale, est elle 1-connexe? Par exemple $B_G$, pour $G$ un groupe disons, est bien une catégorie monoïdale et même ‘groupale’, une $G$-catégorie (tout ce qu'il y a de strict), et pourtant elle n'est \textit{pas} 1-connexe!

\dgrblock{dgrb554}{Insertion of units accounts for all six equalities}
\begin{dcomment}{exposition}{p203}{Insertion of units accounts for all six equalities}
The associativity discussion is an effort to make the comparison procedure usable without remembering a particular sequence of refinements. It can be completed inside the ordinary monoidal category of endofunctors. An object of the coslice $\id_{\mathcal A}\backslash\operatorname{Fun}(\mathcal A,\mathcal A)$ is a pointed endofunctor $(C,i)$. Set
\[
 (C,i)*(D,j)=(CD,(Cj)i)=(CD,(iD)j).
\]
The equality is naturality of $i$. Horizontal composition of natural transformations makes this a bifunctor. Its associativity is strict when functor composition is taken strictly, and its unit is $(\id,\id)$. That unit is also initial in the coslice: the unique arrow to $(C,i)$ is $i$ itself. Although the unit need not lie in the replacement subcategory, its arrows supply $\alpha$ and $\beta$ there.

For an ordered list $(C_1,i_1),\ldots,(C_n,i_n)$ and a subset $J\subseteq\{1,\ldots,n\}$, insert an identity endofunctor in every position not in $J$. The map from this partial composite to $C_1\cdots C_n$ is the horizontal composite of the identity transformations in the retained positions and the transformations $i_k:\id\to C_k$ in the omitted positions. This gives a specified map for every inclusion of ordered subsets. Its compatibility with a second insertion follows from the interchange law for horizontal and vertical composition of natural transformations. In components that law is just the naturality square which permits two adjacent insertions to change order. Any two orders of insertion differ by adjacent interchanges, so they give the same map. This proves the three binary-to-ternary and the three unary-to-ternary equalities requested in the manuscript, and proves the general ordered-subword assertion at the same time.

The concern raised using $B_G$ has a valid lesson but needs a restriction on $G$. The one-object category with automorphism group $G$ is a strict monoidal category of this form only when $G$ is abelian. Indeed, the tensor product and composition share the same unit and obey interchange, so
\[
 a\otimes b=(a\circ1)\otimes(1\circ b)
 =(a\otimes1)\circ(1\otimes b)=a\circ b;
\]
interchanging the two factorizations gives $a\circ b=b\circ a$. Conversely, an abelian group defines such a strict monoidal category by using its group operation for tensor product. Taking $G=\mathbf Z$ therefore gives a valid example with nontrivial fundamental group. What contracts the replacement category is the natural pair of maps into $C*R$, not monoidality alone. The source's broad assertion about an arbitrary group is left as printed and is registered separately from the repaired example.
\end{dcomment}


Bien sûr, on peut court-circuiter la question de la 1-connexité de $\mathcal{C}$, en notant qu'après tout on a déjà a priori un objet privilégié dans
\[
\Hom(\Cat/S, \HOT^{\mathrm{lc}}_W(S)),
\]
i.e. un foncteur canonique
\[
\Cat/S \longrightarrow \HOT^{\mathrm{lc}}_W(S),
\]
savoir le foncteur composé
\[
\Cat/S \underbrace{\longrightarrow}_{\substack{\text{localisation}\\ \text{par } W_S}} (\Cat/S)W_S^{-1} \longrightarrow \underbrace{\HOT^{\mathrm{lc}}_W(S)}_{\overset{\text{déf}}{=} (\Fib_W S)(W^{\mathrm{f}}_S)^{-1}} \quad ,
\tag*{$(*)$}
\]
où la deuxième flèche est définie comme le quasi-inverse de
\[
(\Fib_W S)(W^{\mathrm{f}}_S)^{-1} \longrightarrow (\Cat/S)W_S^{-1}
\tag*{$(**)$}
\]

\msp{page 205}déduit de l'inclusion
\[
(\Fib_W S, W^{\mathrm{f}}_S) \hookrightarrow (\Cat/S, W_S).
\]

Le fait que $(*)$ soit une équivalence, se prouve justement par le fait de l'existence d'un foncteur $C : \Cat/S \longrightarrow \Cat/S$ et d'un $i : \id_{\Cat/S} \longrightarrow C$, ayant les propriétés qu'on sait, i.e. du fait que $C \neq \emptyset$. De plus, la donnée d'un objet $C$ de $\mathcal{C}$ donne un moyen commode d'expliciter le quasi-inverse dudit foncteur. Mais ce quasi-inverse a une existence canonique et indépendante du choix d'un $C$, i.e. d'un mode de calcul. Le fait que le foncteur composé $\varphi$ dans $(*)$ puisse se ‘relever’ (à isomorphisme près) en un foncteur % [sic] C \neq \emptyset -> \mathcal{C} \neq \emptyset, the non-emptiness concerns the category \mathcal{C}
\[
C : \Cat/S \longrightarrow \Fib_W S
\]
entre $\varphi(X) = \hot_{W,S}(X)$ et $\hot_{W,S}(C(X))$ et que l'isomorphisme puisse s'expliciter, lui aussi, par une flèche $X$-fonctorielle
\[
i_X : X \longrightarrow C(X), \qquad i_X \in W,
\]
dans Cat/$S$ lui-même (et non seulement dans le localisé $(\Cat/S)W_S^{-1}$),
\printedpage{134}
% head: A. Grothendieck : Dérivateurs, Ch. XII, Édition des Universités de Montpellier II et Paris VII
\msp{page 206}est une vertu particulière du \texttt{[mot illisible]}, qui ne doit pas pour autant faire oublier le caractère intrinsèque des flèches $(*)$, indépendamment de tout mode de calcul.

\dgrblock{dgrb555}{Why the inclusion becomes an equivalence}
\begin{dcomment}{exposition}{p205}{Why the inclusion becomes an equivalence}
Grothendieck's second answer to the canonicity problem identifies the category being computed. We give the localization argument explicitly, since its hypotheses explain both its strength and its limits. Let $(\mathcal A,W)$ be a category with weak equivalences satisfying two-out-of-three, let $j:\mathcal F\hookrightarrow\mathcal A$ be full, and suppose $C:\mathcal A\to\mathcal F$ is equipped with $i:\id_{\mathcal A}\to jC$ whose components lie in $W$. Give $\mathcal F$ the induced weak equivalences. Then
\[
 \overline j:\mathcal F[W^{-1}]\longrightarrow\mathcal A[W^{-1}]
\]
is an equivalence. Indeed, naturality and two-out-of-three show that $C$ preserves weak equivalences, so it induces $\overline C$. The transformation $i$ induces an isomorphism $\id\simeq\overline j\,\overline C$. Its restriction to $\mathcal F$ has all its components in $\mathcal F$, because that subcategory is full, and gives $\id\simeq\overline C\,\overline j$. These two natural isomorphisms prove the assertion. This is the same replacement mechanism used earlier in the commentary, notably on Chapter~VI, manuscript p.~112; here it proves an equivalence with the entire localized slice.

For the particular functor $(**)$ the weak equivalences on $\Fib_W S$ have been written $W_S^{\mathrm f}$ rather than as the induced $W_S$. One therefore also uses their agreement on this subcategory: the fibre axiom supplies the implication from total to fibrewise equivalence, while the converse is the relevant local-to-global property. Nonemptiness of $\mathcal C$ alone, without this compatibility of weak-equivalence classes, is not the whole statement. Once the compatibility is available, the preceding proof applies exactly as written.

The printed sentence on manuscript p.~205 calls $(*)$ an equivalence. The equivalence proved by the argument is $(**)$; $(*)$ starts in the unlocalized category and includes a localization functor. This is an edition-level slip, not a defect in the replacement argument. For example, even over the terminal base, localization identifies distinct natural-homotopic maps of categories, so the original localization functor cannot generally be fully faithful.

There is also a precise uniqueness statement for the inverse. If $J:\mathcal B\to\mathcal H$ is an equivalence, and $P,P':\mathcal H\to\mathcal B$ are supplied with counits $\epsilon:JP\xrightarrow{\sim}\id$ and $\epsilon':JP'\xrightarrow{\sim}\id$, full faithfulness of $J$ gives a unique natural isomorphism $\vartheta:P\to P'$ satisfying
\[
 \epsilon'\,J\vartheta=\epsilon.
\]
Its image under $J$ is $(\epsilon')^{-1}\epsilon$, which proves existence, uniqueness, naturality, and the cocycle law. Thus the intrinsic object is the inverse equipped with its comparison to the given equivalence. A formula using paths is one realization of that object. In the more general setting of derivable categories, \citet[Proposition~1.8]{Cisinski2010} obtains the analogous equivalence between a category and its fibrant subcategory even when the chosen replacements are organized by a cleavage rather than by a strict functor into the original category.
\end{dcomment}


{\reversemarginpar\marginnote{$\|$}}% double vertical bar printed in the left margin beside the first lines of this paragraph
La clef de la situation, et des questions de canonicité, me semble être dans le fait que $(**)$ est une équivalence de catégories - et il faut avouer que, tout familier et quasiment banal qu'il soit finalement devenu pour moi, c'est là un fait tout à fait remarquable et nullement évident a priori. (Mais sans doute peut-on l'avoir aussi par les fourbis généraux à la \textsc{Quillen-Thomason}, sans se battre avec un formalisme de chemins \dots) Les ‘isomorphismes canoniques’ qu'on peut établir entre une ‘théorie $C$ ’ et une ‘théorie $C'$ ’ (via deux constructions différentes, p. ex., de catégories de chemins, ou via l'approche de \textsc{Quillen-Thomason}, sans chemins) doivent toujours revenir à ceux

\msp{page 207}qu'on peu définir via la comparaison de la ‘théorie $C$ ’ d'une part, la ‘théorie $C'$ ’ de l'autre, avec ce que nous donne directement la flèche composée $(*)$ (p. 204). Finalement, il me semble que \textit{l'existence} d'une théorie des ‘fibres homotopiques’ (ou ‘$W$-homotopiques’) d'une flèche $X \longrightarrow S$ dans Cat, résulte de l'équivalence $(**)$ (p. 204), et cette théorie pourrait sans doute dans une certaine mesure se développer à partir de là, sans choix d'un foncteur $C$ ni d'une structure de catégorie de modèles au sens de \textsc{Quillen}. % [sic] qu'on peu définir -> qu'on peut définir, verb form

\dgrblock{dgrb556}{A homotopy fibre which an ordinary kernel misses}
\begin{dcomment}{examples}{p206}{A homotopy fibre which an ordinary kernel misses}
The replacement construction is already visible in elementary homological algebra, where its comparison maps and the information lost by taking a raw fibre can both be written down. This gives a different example from the groupoid calculations in the commentary on earlier chapters. Let $C,D$ be cochain complexes of abelian groups, with differential of degree $1$, and let $f:C\to D$ be a cochain map. Take $W$ to be the quasi-isomorphisms. Degreewise surjections preserve quasi-isomorphisms under pullback: pulling such a surjection $E\to D$ back along a map $A\to B$ over $D$ gives a map between the short exact sequences with common kernel $\ker(E\to D)$ and quotient map $A\to B$. The long exact cohomology sequences prove the assertion. Thus these surjections belong to the class of $W$-fibrations used below.

Define a complex
\[
 E(f)^n=C^n\oplus D^{n-1}\oplus D^n,
 \qquad
 d(c,h,y)=(d_Cc,\ y-f(c)-d_Dh,\ d_Dy).
\]
The equation $d^2=0$ follows from $d_Df=fd_C$. The maps
\[
 j(c)=(c,0,f(c)),\qquad p(c,h,y)=y,\qquad r(c,h,y)=c
\]
are cochain maps, $pj=f$, and $rj=\id_C$. The degree $-1$ map
\[
 H(c,h,y)=(0,0,h)
\]
satisfies $dH+Hd=\id_{E(f)}-jr$. Consequently $j$ is a cochain-homotopy equivalence, in particular a quasi-isomorphism, while $p$ is degreewise surjective. We have obtained the required factorization $C\xrightarrow{j}E(f)\xrightarrow{p}D$ with its comparison map explicit.

The strict fibre of $p$ over the zero map $0\to D$ is
\[
 F(f)^n=C^n\oplus D^{n-1},\qquad
 d(c,h)=(d_Cc,-f(c)-d_Dh).
\]
With the convention $\operatorname{Cone}(f)^n=D^n\oplus C^{n+1}$ and differential $(y,c)\mapsto(d_Dy+f(c),-d_Cc)$, this is canonically isomorphic to $\operatorname{Cone}(f)[-1]$ after reordering the summands. The sign is forced by the shift; it is not an additional convention about homotopy fibres. The short exact sequence $0\to D[-1]\to F(f)\to C\to0$ gives the familiar exact sequence relating the cohomology of the fibre, $C$, and $D$.

For a concrete calculation, let $C=D=\mathbf Z[0]$ and let $f$ be multiplication by an integer $m\geq2$. The ordinary kernel is zero. The homotopy fibre is the two-term cochain complex
\[
 \mathbf Z\xrightarrow{-m}\mathbf Z
\]
in degrees $0$ and $1$. It has $H^0=0$ and $H^1=\mathbf Z/m\mathbf Z$. The replacement has retained the failure of surjectivity which the ordinary kernel discards. This example is classical: Brown lists complexes in an abelian category, with quasi-isomorphisms and surjections, among his categories of fibrant objects \citep[Part I, \S1, Example~3, p.~421]{Brown1973}. The calculation explains why Grothendieck insists on replacing the structural map before taking its fibre.
\end{dcomment}


\section*{4\quad Axiomatisation des ‘fibres $W$-homotopiques’.}

La reflexion de hier me fait revenir sur l'axiomatique dans des ‘catégories de modèles’ % [sic] reflexion -> réflexion, missing accent
\[
(\mathcal{M}, W)
\]
(\footnote[119]{$W$ satisfait l'axiome de saturation faible : les isomorphismes \texttt{[sont]} dans $W$, et si deux parmi $f$, $g$, $gf$ \texttt{[sont]} dans $W$, le troisième aussi.}) plus générales que Cat, dans l'esprit de VI (cf. théorème-scholie 2, p. 108-113). Pour le moment, on va supposer seulement $\mathcal{M}$ stable par produits fibrés (on n'a pas besoin d'objet final). On introduit l'ensemble
\[
\mathrm{Fib}_W \subseteq \mathrm{Fl}(\mathcal{M})
\]

\msp{page 208}des \textit{$W$-fibrations} (sous-entendu ‘ultrafortes’!) par la condition

\begin{quote}
$f : X \longrightarrow Y$ est dans $\mathrm{Fib}_W$ si et seulement si le foncteur changement de base par $X$,
\[
\begin{aligned}
Y' &\quad\longmapsto\quad X' = X \times_Y Y' \\
\mathcal{M}/Y &\quad\longrightarrow\quad \mathcal{M}/X
\end{aligned}
\]
\end{quote}

% ---- from Der-ChXIIc-30-03-03_p021-025.tex
\printedpage{135}
% head: A. Grothendieck : Dérivateurs, Ch. XII, Édition des Universités de Montpellier II et Paris VII

\begin{quote}
transforme $W_X$ dans $W_Y$, où $W_X \subseteq \Fl(\cM/X)$ \texttt{[est]} défini par

$u \in W_X \iff \sigma(U) \in W$ \ ($\sigma : \cM/X \longrightarrow \cM$ le foncteur ‘source’, ou ‘oubli de $X$’). % [sic] \sigma(U) -> \sigma(u), the arrow is called u
\end{quote}

(\footnote[120]{\textbf{NB} Donner ici la proposition: $\mathrm{Fib}_W$ \textit{stable par composition et par changement de base, et contient les isomorphismes, stable par facteurs directs.}}). Posons, pour tout $S$ dans $\cM$,
\[
\Ho^!_W(S) = (\cM/S)W_S^{-1},
\]
qui joue le rôle d'une catégorie de $W$-types d'homotopie relatifs \textit{localement constants} sur $S$. On aimerait décrire une loi contravariante en $S$ pour $\Ho^!_W(S)$, i.e. pour
\[
f : S' \longrightarrow S
\]
flèche quelconque dans $\cM$, définir
\[
f^*_W : \Ho^!_W(S) \xrightarrow{?} \Ho^!_W(S'). \tag*{$(*)$}
\]
La difficulté est, bien sûr, que le foncteur changement de base
\[
f^* : \cM/S \longrightarrow \cM/S'
\]
n'est \textit{pas} en général compatible avec les \msp{page 209} localiseurs $W_S$, $W_{S'}$. La notion de $W$-fibration sert donc en premier lieu à dégager la condition la plus évidente sur une flèche $f : S' \longrightarrow S$ dans $\cM$, pour que le foncteur changement de base soit compatible avec les localiseurs $W_S$, $W_{S'}$ (c'est \textit{ça} la définition), donc induit $(*)$ sur les localisées.

\dgrblock{dgrb557}{The intrinsic fibrations are sharp maps}
\begin{dcomment}{later}{p208}{The intrinsic fibrations are sharp maps}
The purpose of the new definition is to extract from $(\mathcal M,W)$ the largest class of maps along which ordinary pullback is homotopical. For $f:X\to Y$ the correctly typed condition is
\[
 f^*:(\mathcal M/Y,W_Y)\longrightarrow(\mathcal M/X,W_X).
\]
The edition reverses $W_X$ and $W_Y$ in the sentence following the displayed pullback functor; the direction of that functor fixes the intended reading. Unpacked, the definition says that whenever $a:A\to B$ is in $W$ and is a map over $Y$, the map
\[
 X\times_Y A\longrightarrow X\times_Y B
\]
is again in $W$. The quantification over all objects over $Y$ is essential. Testing only pullbacks of weak equivalences whose target is exactly $Y$ would omit part of the requirement.

In a model category, this is precisely the condition called \emph{sharp} by \citet[\S2]{Rezk1998Sharp}. His Proposition~2.2 says that all model fibrations are sharp exactly when the model category is right proper. His Proposition~2.7 further identifies sharp maps, in a right proper model category, with those maps all of whose ordinary pullback squares are homotopy cartesian. These hypotheses distinguish three statements which can otherwise look indistinguishable: a map can be a fibration for a chosen model structure, preserve weak equivalences by pullback, or compute all homotopy pullbacks strictly. The first is extra structure, the second depends only on $(\mathcal M,W)$, and the equivalence with the third needs the stated properness hypothesis. Rezk attributes the sharp-map idea to Hopkins; we are comparing equivalent conditions, not asserting a historical line of transmission from this manuscript.

The basic stability properties require little machinery. If $f:X\to Y$ and $g:Y\to Z$ preserve weak equivalences by pullback, so does $gf$, because $(gf)^*$ is canonically $f^*g^*$. To see stability under a pullback $f':X\times_Y T\to T$, regard a weak equivalence over $T$ also as a map over $Y$; its pullback by $f'$ is the same underlying map as its pullback by $f$. Isomorphisms plainly have the property. The statement about retracts in footnote~120 is different: it needs retract closure of $W$. Under that hypothesis, pulling back a retract diagram makes the map to be tested a retract of a map in $W$, proving the result. Under only the weak saturation explicitly listed in footnote~119, the retract assertion fails, as the following finite example shows.
\end{dcomment}


Le \textit{‘problème des fibres homotopiques’} est justement de définir $(*)$ pour \textit{toute} flèche dans $\cM$, pas seulement pour les $W$-fibrations, de tellle façon que les $\Ho^!_W(S)$ forment, pour $S$ % [sic] tellle -> telle, letter doubled
variable, les fibres d'une catégorie fibrée remarquable sur $\cM$. Quand $\cM$ admet un objet final, de sorte que l'on peut regarder les flèches
\[
p : e \longrightarrow S \qquad (p \text{ comme ‘point’})
\]
dans $\cM$ comme des ‘points’ de $S$, la ‘fibre $W$-homotopique’ d'un objet quelconque $X/S$ de $\cM/S$ en ledit point sera par définition l'élément de
\[
\Ho_W = \Ho_W(\cM) \simeq \Ho^!_W(e)
\]
défini par
\[
p^*(\ho^!_{W/S}(X/S)),
\]
\printedpage{136}
% head: A. Grothendieck : Dérivateurs, Ch. XII, Édition des Universités de Montpellier II et Paris VII
\msp{page 210}
où
\[
\ho^!_{W/S} : \cM/S \longrightarrow \Ho^!_W(S) \overset{\text{déf}}{=} (\cM/S)(W_S)^{-1}
\]
est le foncteur de localisation canonique.

\dgrblock{dgrb558}{The retract condition needs a hypothesis}
\begin{dcomment}{development}{p209}{The retract condition needs a hypothesis}
The minimal axioms are being separated here so that the fibre theory will apply beyond familiar model categories. A small counterexample helps to keep that economy honest. We use the usual meaning of a direct factor of a morphism, namely a retract in the arrow category. The observation below addresses only the retract clause of footnote~120; composition, pullback, and isomorphism stability remain valid under the assumptions stated in the text.

\noindent\textbf{Finite-set counterexample to the retract clause.}\ours\ 
Let $\mathcal M$ be the category of finite sets, and let $W$ consist of all maps $A\to B$ for which $|A|\equiv|B|\pmod2$. Then $\mathcal M$ has pullbacks and $W$ contains all isomorphisms and satisfies two-out-of-three. Nevertheless the class $\mathrm{Fib}_W$ defined on manuscript p.~208 is not closed under retracts.

\emph{Proof.} The axioms for $W$ follow by applying equality in $\mathbf Z/2\mathbf Z$ to the cardinalities of the three objects in a composable pair. Set $S=\{0,1\}$. Let $E$ have one element over $0$ and three over $1$, and write $p:E\to S$ for its structural map. For every finite set $T\to S$,
\[
 |E\times_S T|=|T_0|+3|T_1|\equiv|T|\pmod2.
\]
Thus if $A\to B$ over $S$ belongs to $W$, its pullback along $p$ also belongs to $W$. This proves $p\in\mathrm{Fib}_W$.

Let $E'$ have one element over $0$ and two over $1$, and let $p':E'\to S$. Include $E'$ into $E$ by choosing two of the three elements over $1$. Define a retraction $E\to E'$ which is the identity on the chosen elements and sends the remaining element to either of the two elements over $1$. Together with the identity of $S$, these maps exhibit $p'$ as a retract of $p$ in the arrow category. But the map $\varnothing\to S$, considered as a map from the empty object to the identity object of $\mathcal M/S$, lies in $W$, since both cardinalities are even. Its pullback along $p'$ is $\varnothing\to E'$, whose target has cardinality $3$. This map is not in $W$. Hence $p'\notin\mathrm{Fib}_W$.\hfill$\square$

This counterexample also makes clear what repairs the statement. One may assume $W$ is retract closed, as in a model category; strong saturation is more than sufficient. Alternatively, one can omit the retract clause, which is not needed for the formal base-change construction on the following pages. We checked the axioms in Brown's Part~I, Cisinski's \S1, and Rezk's sharp-map discussion; those sources do not provide this particular example. The argument is supplied here as a local clarification of the manuscript, without a claim of priority.
\end{dcomment}


Pour tout $S$ dans $\cM$, soit
\[
\Fib_W S \subseteq \cM/S
\]
la sous-catégorie pleine formée des objets $X/S$ de $\cM/S$ tels que $X \longrightarrow S$ soit $W$-fibrant. On va définir une variante
\[
\Ho^{\lc}_W(S)
\]
de $\Ho^!_W(S)$, comme
\[
\Ho^{\lc}_W(S) \overset{\text{déf}}{=} (\Fib_W S)(\Wu_S)^{-1},
\]
où

\begin{quote}
$\Wu_S \subseteq \Fl(\Fib_W S)$ est le localiseur formé des flèches $u : X \longrightarrow Y$ qui sont ‘universellement sur $S$ dans $W$’, i.e. telles que pour toute flèche $S' \longrightarrow S$ de $\cM$ de but $S$, la flèche déduite par changement de base,
\[
u' : X' = X \times_S S' \longrightarrow Y' = Y \times_S S'
\]
est dans $W_{S'}$.
\end{quote}

On a ‘fait ce qu'il fallait’ pour que $\Ho^{\lc}_W(S)$ soit contravariant en $S$ pour des flèches \textit{quelconques} $f : S' \longrightarrow S$ de $\cM$, elles définissent
\[
f^{*\,\lc}_W : \Ho^{\lc}_W(S) \longrightarrow \Ho^{\lc}_W(S'), \tag*{$(**)$}
\]
\msp{page 211}
puisque les $W$-fibrations sont stables par changement de base, et que par suite $f^* : \cM/S \longrightarrow \cM/S'$ induit
\[
f^*_{\Fib_W} : \Fib_W S \longrightarrow \Fib_W S'
\]
qui, par définition de $\Wu_S$ et $\Wu_{S'}$, est compatible avec les localiseurs
\[
f^*_{\Fib_W} : (\Fib_W S, \Wu_S) \longrightarrow (\Fib_W S', \Wu_{S'}),
\]
et passe donc aux catégories des fractions en une flèche $(*)$.

On peut se demander pourquoi ne pas être plus généreux, et définir $\Ho^{\lc}_W(S)$ comme $(\cM/S)(\Wu_S)^{-1}$, en définissant $\Wu_S$ dans $\cM/S$ de la même façon que tantôt dans $\Fib_W S$ - ce qui ne coûte pas plus cher. Une première réponse est d'expérience le localiseur $\Wu_S$ dans $\cM/S$ tout entier s'avère pratiquement inutilisable, d'une part parce que l'appartenance à $\Wu_S$ (pour une flèche entre objets de $\cM/S$ qui ne sont pas très particuliers, p.ex. $W$-fibrés sur $S$) est invérifiable (sauf pour des isomorphismes et des choses comme ça \dots), et que d'autre part la catégorie de fractions est prohibitivement
\printedpage{137}
% head: A. Grothendieck : Dérivateurs, Ch. XII, Édition des Universités de Montpellier II et Paris VII
\msp{page 212}
grosse et ne peut servir (semble-t-il) à rien. Une autre réponse, c'est qu'on veut que la catégorie obtenue réponde bien à l'intuition de types d'homotopie relatifs \textit{localement constants} sur $S$ - et on veut donc que pour un changement de base quelconque $f : S' \longrightarrow S$ par un

\dgrblock{dgrb559}{Two localizations with different immediate advantages}
\begin{dcomment}{exposition}{p211}{Two localizations with different immediate advantages}
Grothendieck keeps two constructions in view because each makes one half of the desired theory immediate. Write
\[
 \mathcal H_S=(\mathcal M/S)[W_S^{-1}],\qquad
 \mathcal L_S=(\mathrm{Fib}_W S)[(W_S^{\mathrm u})^{-1}].
\]
The first accepts every object over $S$. Its ordinary inverse-image functor is initially available only for a base map in $\mathrm{Fib}_W$. The second accepts only objects whose structural map is a $W$-fibration, but its inverse images exist for every base map. The universal quantifier defining $W_S^{\mathrm u}$ is precisely what makes this second assertion formal.

For $f:T\to S$, ordinary pullback sends $\mathrm{Fib}_W S$ into $\mathrm{Fib}_W T$, as just proved. If $u\in W_S^{\mathrm u}$ and $g:U\to T$, the iterated pullback $g^*f^*u$ is isomorphic to $(fg)^*u$, hence belongs to $W$. Therefore $f^*u\in W_T^{\mathrm u}$. The universal property of localization supplies a functor
\[
 P_f:\mathcal L_S\longrightarrow\mathcal L_T.
\]
The canonical isomorphisms between iterated ordinary pullbacks survive localization and give $P_gP_f\simeq P_{fg}$. Their associativity follows from the universal property of the pullback: both routes identify an iterated pullback with the same object, carrying the same projections. Thus the result is naturally a pseudofunctor, or equivalently a fibred category. No strictly functorial global choice of pullbacks is needed.

One can also see exactly what must be checked to define the fibre at a generalized point. Fix $X\to S$ in $\mathrm{Fib}_W S$. Sending $T\to S$ to the absolute homotopy type of $X\times_S T$ is a covariant functor on $\mathcal M/S$; a map $T'\to T$ induces the corresponding map between total pullbacks. If that map of generalized points is in $W_S$, its image is in $W$ by the defining property of $X\to S$. The functor consequently inverts weak equivalences between generalized points. This assertion is compatible with replacing $X$ by a morphism in $W_S^{\mathrm u}$, since evaluation at every $T$ then gives a weak equivalence. It therefore makes sense for all $\xi\in\mathcal L_S$, not only for a chosen representative.

This explains the two separate roles of the subcategory and the localizer. Restricting to $\mathrm{Fib}_W S$ controls variation of the generalized point by a weak equivalence. Restricting the arrows to $W_S^{\mathrm u}$ controls variation of the representative by a weak equivalence after any base change. Omitting either condition loses one of these immediate guarantees. The manuscript's remark that the unrestricted universal localization is too large is a methodological judgment; it is not a theorem saying that every such localization is mathematically useless. What has been proved is the precise pair of invariance properties above.
\end{dcomment}


\begin{quote}
$S'$ dans $\cM/S$, d'où un foncteur composé
\[
\begin{tikzcd}[column sep=small]
\cF_S(\Wu_S)^{-1} \arrow[r, "f^{*\,\lc}_W"] \arrow[rr, dotted, bend right=15, "\varphi_{S'}"'] & \cF_{S'}(\Wu_{S'})^{-1} \arrow[r] & \Ho_W
\end{tikzcd}
\]

(\footnote[121]{$\cF$ désigne ici une sous-catégorie fibrée de $\cM/\cM = \underline{\Fl}(\cM)$ sur $\cM$, substitut éventuel de celle formée de $\Fib_W S$ ou de $\cM/S$ tout entier.}) déduit des foncteurs de catégories avec localiseurs
\[
(\cF_S, \Wu_S) \xrightarrow{f^*} (\cF_{S'}, \Wu_{S'}) \xrightarrow{\text{oubli de } S'} (\cM, W)
\]
\end{quote}

(foncteur composé $\varphi_{S'}$ qui peut être vu comme le foncteur \textit{‘fibre homotopique en l'objet $S'$ sur $S$’}, cet objet étant visualisé comme un ‘point’ de $S$ ‘à valeurs dans $S'$ ’ suivant le yoga fonctoriel général \dots), le foncteur
\[
\begin{aligned}
S' &\longmapsto \varphi_{S'}(\xi) \qquad \text{pour } \xi \text{ dans } \cF_S(\Wu_S)^{-1}\\
\cM/S &\longrightarrow \Ho_W
\end{aligned}
\]
transforme flèches dans $W_S$ en isomorphismes : la ‘fibre $W$-homotopique’ en $S'/S$ d'un objet de $\cF_S(\Wu_S)^{-1}$ (vu comme $W$-type d'homotopie \textit{localement constant} sur $S$) doit être ‘la même’, si $S'$ remplacé par $S''$ avec $S'' \longrightarrow S'$ donné et dans $W_S$. C'est pour cela qu'il est raisonnable de prendre $\cF_S$ ‘assez petit’ (savoir $\Fib_W S$, qui est \msp{page 213} exactement ajusté pour cela), plutôt que de prendre $\cM/S$ tout entier.

On peut donc dire que la catégorie $\Fib_W S$, et le localiseur $\Wu_S$ ont été construits ad-hoc, le plus simplement possible, pour donner lieu à un localisé $\Ho^{\lc}_W(S)$ qui réponde à l'intuition de $W$-types d'homotopie relatifs localement constants des deux façons suivantes :

\begin{enumerate}[label={}]
\item[1°)] Il y a spontanément une loi contravariante en $S$, pour des flèches quelconques $S' \longrightarrow S$ dans Cat (on a défini $\Wu_S$ juste exprès en conséquence).
\item[2°)] que la notion de ‘fibre $W$-homotopique en $S'/S$’ d'un objet $\xi$ de $\Ho^{\lc}_W$ satisfasse à la condition précédente, i.e. une flèche $S'' \longrightarrow S'$ dans $\cM/S$ qui est dans $W_S$ définisse un \textit{isomorphisme} (dans $\Ho_W$) des fibres homotopiques
\[
\underbrace{\xi_{S''}}_{\varphi_{S''}(\xi)} \longrightarrow \underbrace{\xi_{S'}}_{\varphi_{S'}(\xi)}
\]
(et on a choisi $\cF_S = \Fib_W(S)$ juste exprès pour ça).
\end{enumerate}

\printedpage{138}
% head: A. Grothendieck : Dérivateurs, Ch. XII, Édition des Universités de Montpellier II et Paris VII

On a donc une excellente notion de fibres $W$-homotopiques pour des $X$ $W$-fibrés sur $S$ - cette fois excellente pour des changements de base \textit{quelconque} - et il n'y a aucun mérite à ça puisqu'on a tout défini exprès pour ça! % [sic] quelconque -> quelconques, agreement with "des changements"

\msp{page 214}
Pour le dire autrement, on a défini \textit{deux} notions de $W$-types d'homotopie relatifs sur $S$, l'une de façon à donner un tel type d'homotopie $\ho^!_{W,S}(X)$ pour \textit{tout} objet de $\cM/S$, mais il n'y a a priori une bonne notion de changement de base que pour les changements de base $S' \longrightarrow S$ qui sont $W$-fibrés, l'autre, qui correspond plus directement à l'intuition d'un $W$-type d'homotopie localement constant sur $S$, n'associe un tel type $\ho^{\lc}_{W,S}(X)$ qu'aux $X$ de $\cM/S$ qui sont $W$-fibrés sur $S$, en revanche cette notion admet une bonne notion d'image inverse et de ‘fibre $W$-homotopique’ pour des changements de base quelconques.

Une ‘bonne’ théorie des fibres $W$-homotopiques doit réunir en une seule les deux notions précédentes, de sorte que les avantages des deux s'ajoutent, et les désavantages s'éliminent! Elle s'obtient à l'aide d'un axiome draconien sur la situation (alors que jusqu'à maintenant on n'a rien utilisé sur $W \subseteq \Fl(\cM)$, si ce n'est tout au plus la saturation faible, en plus de l'existence des produits fibrés) :

\msp{page 215}
\circled{a} (\textit{Axiome de la fibre homotopique.}) Pour tout objet $S$ dans $\cM$, le foncteur canonique
\[
\Ho^{\lc}_W(S) \longrightarrow \Ho^!_W(S)
\]
déduit de l'inclusion
\[
(\Fib_W(S), \Wu_S) \longrightarrow (\cM/S, W_S)
\]
est une \textit{équivalence de catégories.}

Cet axiome implique que dans $\Fib_W S$, $\Wu_S$ et le localiseur induit par $W_S$ (que je désigne encore par $W_S$ par abus de notation) ont même saturé fort. En fait, il est naturel de compléter l'axiome précédent par le suivant.

\circled{b} (Axiome adjoint de l'axiome de la fibre homotopique.) Sur $\Fib_W S$, les localiseurs $\Wu_S$ et $W_S$ coïncident, i.e. on a ceci :

\begin{quote}
Soient $X$, $Y$ dans $\Fib_W S$, $X \xrightarrow{u} Y$ un $S$-morphisme tel que $u \in W_S$, alors $u \in \Wu_S$ (i.e. le morphisme $u' : X' \longrightarrow Y'$ déduit par tout changement de base $S' \longrightarrow S$ est encore dans $W$).
\end{quote}

\msp{page 216}
Je présume que a) sera rarement satisfait sans que b) le soit (et inversement d'ailleurs), et que la démonstration de a) et b) se fera pratiquement en une seule haleine.

\dgrblock{dgrb560}{Transporting base change and separating the two axioms}
\begin{dcomment}{exposition}{p215}{Transporting base change and separating the two axioms}
The comparison axiom now combines the two advantages. Let
\[
 J_S:\mathcal L_S\longrightarrow\mathcal H_S
\]
be induced by inclusion. Axiom~a says that $J_S$ is an equivalence. Choose an inverse $R_S$ together with the unit and counit of an adjoint equivalence. For an arbitrary $f:T\to S$, define derived inverse image by
\[
 f^!_W=J_TP_fR_S:\mathcal H_S\longrightarrow\mathcal H_T.
\]
This is the operation denoted $f^*_W$ in the manuscript; the temporary exclamation mark here only distinguishes it from raw pullback. For composable $U\xrightarrow{g}T\xrightarrow{f}S$, the comparison
\[
 J_UP_gR_TJ_TP_fR_S\xrightarrow{\sim}
 J_UP_gP_fR_S\xrightarrow{\sim}J_UP_{fg}R_S
\]
uses $R_TJ_T\simeq\id$ and the ordinary pullback comparison. The triangle identities of the chosen adjoint equivalences and the associativity of pullback give the unit and pentagon identities. Equivalently, these identities can be checked after the fully faithful equivalences $J_S$, where they reduce to those same identities. A different choice of inverses gives the unique normalized comparison described above, so the resulting fibred theory is independent of that choice up to the corresponding coherent equivalence.

Axiom~a also has a useful immediate consequence which is weaker than axiom~b. If $u$ between $W$-fibrant objects belongs to $W_S$, then $J_S$ sends its image in $\mathcal L_S$ to an isomorphism. An equivalence reflects isomorphisms, so $u$ belongs to the strong saturation of $W_S^{\mathrm u}$ on $\mathrm{Fib}_W S$. Since $W_S^{\mathrm u}\subseteq W_S$, the two classes have the same strong saturation there. This is exactly the intermediate observation in the text. It does not, under bare two-out-of-three, identify the two original classes of arrows.

There is a clean case in which axiom~b follows. Suppose $W$ itself is strongly saturated: a morphism belongs to $W$ exactly when it is inverted by $\mathcal M\to\mathcal M[W^{-1}]$. Then $W_S^{\mathrm u}$ is strongly saturated in $\mathrm{Fib}_W S$. Indeed, if $u$ becomes invertible after localization at $W_S^{\mathrm u}$, every evaluation functor
\[
 X\longmapsto X\times_ST\longmapsto
 \mathcal M[W^{-1}]
\]
for $T\to S$ sends $u$ to an isomorphism. Strong saturation gives $u_T\in W$ for every $T$, which is the definition of $u\in W_S^{\mathrm u}$. Thus axiom~a implies axiom~b under this extra hypothesis. Without it, the statement in the manuscript should retain its saturation formulation; the author's expectation that the two axioms will usually be verified together is not a formal converse.

The construction gives a derived base-change theory, but it does not by itself assert all the axioms of a derivator on arbitrary diagram categories. Those require further control of diagrams, pointwise detection, and homotopy Kan extensions. The distinction preserves the considerable content of the comparison axiom without crediting it with statements that are proved elsewhere.
\end{dcomment}


Si $\cM = \mathrm{Cat}$, $W$ un localiseur fondamental, alors la notion adaptée ici de $W$-fibration (lire : ‘ultraforte’, dans la terminologie spécifique développée dans Cat) est à tel point forte, qu'elle implique déjà pratiquement b) (moyennant W(4a$'$), il est vrai, cf. p. 53), et
\printedpage{139}
% head: A. Grothendieck : Dérivateurs, Ch. XII, Édition des Universités de Montpellier II et Paris VII
a) n'est guère vérifiable que si on a W(4a) (\footnote[122]{Non, il me semble que W(4b) doit suffire \dots}) (loc. cit.). Donc ici c'est ‘l'axiome des fibrations’ sur $W$, sous sa forme forte W(4a), qui est la clef de voûte d'une bonne théorie des fibres homotopiques.

Une telle théorie (revenant au cas général) revient à dire, pour l'essentiel, qu'il y a ‘assez’ de $W$-fibrations (‘ultrafortes’) pour pouvoir décrire \textit{tout} objet d'un $\cM/S$, à $W_S$-équivalence près, comme un objet $W$-fibrant sur $S$. Dans les cas à ma connaissance, on a de plus ceci : on \msp{page 217} a pour tout $S$ dans $\cM$ une sous-catégorie $\cF_{0\,S}$ de $\Fib_W S$ (je pense par exemple à $\underline{\mathrm{Parf}}_W S$, voire $\underline{\mathrm{Parf}}_{W_\omega} S$, dans le cas où $\cM = \mathrm{Cat}$), telle que pour toute sous-catégorie pleine $\cF$ de $\cM/S$ telle que
\[
\cF_{0\,S} \subseteq \cF \subseteq \cM/S
\]
(a fortiori, pour une $\cF$ pleine telle que $\Fib_W S \subseteq \cF \subseteq \cM/S$), les inclusions induisent des équivalences de catégories
\[
\cF_{0\,S}(W_S)^{-1} \xrightarrow{\sim} \cF(W_S)^{-1} \xrightarrow{\sim} \underbrace{(\cM/S)(W_S)^{-1}}_{\overset{\text{déf}}{=} \mathrm{Hot}^!_W(S)}. % [sic] Hot -> Ho, elsewhere the notation is Ho^!_W(S)
\]
Le choix de $\cF_{0\,S}$ ne semble pas pouvoir se fixer de façon ‘la meilleure possible’, une fois pour toutes dans le contexte axiomatique adopté ici. C'est un élément de structure supplémentaire dans $\cM$ (en plus du localiseur $W$), au même titre que la donnée des classes de $Q$-fibrations et $Q$-cofibrations, à la \textsc{Quillen}. J'ignore si les axiomes a), b) ci-dessus (axiome de la fibre homotopique) impliquent déjà \msp{page 218} la variante de \circled{a} :

\dgrblock{dgrb561}{What the weaker model formalisms establish}
\begin{dcomment}{later}{p216}{What the weaker model formalisms establish}
The new axioms are close to Brown's economy of assumptions, but they organize it differently. Brown begins with specified fibrations and weak equivalences on a category whose objects are fibrant, and assumes the path-object and pullback axioms needed for a calculus of homotopy. Here the fibrations are extracted from $W$ itself by asking for homotopical base change; the existence of enough suitable representatives is imposed afterwards. One should therefore compare the actual hypotheses before using a theorem about one formalism in the other.

Cisinski's derivable categories give a precise later setting for this programme. Definition~1.1 separates weak equivalences, fibrations, and factorization toward fibrant objects. Proposition~1.8 proves that restricting to fibrant objects does not change the homotopy category. Theorem~2.21 then constructs a weak left derivator on finite direct categories, with the additional diagram axiom stated there \citep{Cisinski2010}. This is a useful progression for reading the present section: comparison of a single localization comes first; coherent diagram calculus requires further work. None of these results says that an arbitrary category equipped only with a two-out-of-three class automatically has homotopy fibres.

For ordinary spaces or simplicial sets with weak homotopy equivalences, a right proper model presentation supplies a familiar way to meet the requirements. Model fibrations are sharp, and ordinary pullback of a sharp map computes homotopy pullback. Weak equivalences between such representatives remain weak equivalences after derived base change. These are the roles of properness and replacement in the usual theory, rather than unexplained special properties of a path formula. In the generality of the manuscript, the corresponding assertions are the explicit axioms a and b, unless a chosen factorization theorem proves them.

The fibrewise ordinary diagram $s\mapsto\mathrm{hot}_W(X_s)$ still forgets coherence. In manuscript p.~200 the functor through $\Pi S^\circ$ is being used for the locally constant subcategory: the missing superscript $\mathrm{lc}$ on the displayed $\mathrm{HOT}_W(S)$ should be understood in that context. It is not a factorization of every varying relative type through the fundamental groupoid. The commentary on earlier chapters already treats the difference between an ordinary diagram in a homotopy category and a coherent family (notably on Chapter~VI, manuscript pp.~7 and 74). The equivalence of localized slice categories here concerns the coherent family; evaluation in ordinary homotopy types is one invariant of it.
\end{dcomment}


\circled{c} Pour toute sous-catégorie pleine $\cF$ de $\cM/S$ telle que
\[
\Fib_W S \subseteq \cF \subseteq \cM/S
\]
les inclusions induisent des équivalences
\[
\underbrace{(\Fib_W S)(W_S)^{-1}}_{\begin{subarray}{l} = \ \Ho^{\lc}_W(S)\\ \text{moyennant a)}\end{subarray}} \xrightarrow{\sim} \cF(W_S)^{-1} \xrightarrow{\sim} \underbrace{(\cM/S)(W_S)^{-1}}_{=\Ho^!_W(S)}.
\]

Mais je présume que pratiquement, quand on pourra vérifier \circled{a}, on pourra vérifier \circled{c} également (\footnote[123]{Cf. page 240 ff. Il faudrait ici développer un peu plus le formalisme de $\Ho^{\lc}_W(S)$ et les fibres homotopiques, sous l'hypothèse \circled{a}\,\circled{b} \ plus haut.}).

% ---- from Der-ChXIIc-30-03-03_p026-030.tex
\printedpage{140}
% head: A. Grothendieck : Dérivateurs, Ch. XII, Édition des Universités de Montpellier II et Paris VII

\section*{5 Axiomatique des $W$-types d'homotopie relatifs non localement constants $\mathrm{Ho}^{\mathcal{L}}_W(S)$.}

On n'a rien dit jusqu'à ce point d'éventuels $W$-types d'homotopie relatifs (sur un objet $S$ de $\mathcal{M}$) qui ne seraient \textit{pas} ``localement constants''. Il en était question pourtant dans loc. cit. (VI, théorème-scholie 2, partie II du théorème pages 109-113.) Je profite de l'occasion pour revenir là dessus. Il faut introduire, dans chaque $\mathcal{M}/S$, un deuxième localiseur
\[
W^0_S \subseteq \mathrm{Fl}(\mathcal{M}/S), \qquad W^0_S \subseteq W_S
\]
(plus fin que $W_S$), ce qui permettra de définir
\msp{page 219}
\[
\mathrm{Ho}_W(S) \deff (\mathcal{M}/S)(W^0_S)^{-1}.
\]
Qu'il n'y ait pas de choix \textit{canonique} qui s'impose ici, déduit de la seule structure $(\mathcal{M},W)$, est claire sur l'exemple de $(\mathrm{Cat},W)$, auquel cas on trouve \textit{deux} choix symétriques, transformés l'un en l'autre par l'involution
\[
X \longmapsto X^\circ : (\mathrm{Cat},W) \xrightarrow{\sim} (\mathrm{Cat},W),
\]
à savoir $W^{\mathrm{g}}_S$ et $W^{\mathrm{d}}_S$. Le choix de l'un ou de l'autre, disons (pour fixer les idées) de $W^{\mathrm{g}}_S$, est associé à une notion d'\textit{isomorphisme local} (ou de `coisomorphisme local' dans le cas de $W^{\mathrm{d}}_S$) que je voudrais axiomatiser. Il semble bien que pour les principales catégories de modèles pour Hot (et pour ses variantes), on dispose de façon assez naturelle d'une telle notion d'isomorphisme local (sans l'embarras, le plus souvent, d'un automorphisme involutif fournissant une notion symétrique \dots).

\dgrblock{dgrb562}{Intermediate categories and the extra information in a replacement}
\begin{dcomment}{status}{p218}{Intermediate categories and the extra information in a replacement}
The last question of this section asks for more than essential surjectivity. It asks whether one may admit any collection of additional representatives without changing the localized category. A functorial replacement proves exactly this strengthened assertion, and it is useful to isolate the short proof. Let $\mathcal F_0\subseteq\mathcal A$ be full, and suppose a coaugmented replacement $C:\mathcal A\to\mathcal F_0$, $i:\id\to C$, has components in $W$. For every full intermediate category
\[
 \mathcal F_0\subseteq\mathcal E\subseteq\mathcal A,
\]
restricting $C$ and $i$ to $\mathcal E$ gives a replacement into $\mathcal F_0$. The localization proof on p.~205 then applies both to $\mathcal F_0\subseteq\mathcal E$ and to $\mathcal F_0\subseteq\mathcal A$. Consequently
\[
 \mathcal F_0[W^{-1}]\xrightarrow{\sim}
 \mathcal E[W^{-1}]\xrightarrow{\sim}
 \mathcal A[W^{-1}].
\]
Fullness ensures that the components of the restricted coaugmentation really are arrows of $\mathcal E$; no closure of $\mathcal E$ under limits or factorizations is needed. When $\mathcal F_0=\mathrm{Fib}_W S$ and axiom~b identifies the weak-equivalence classes, this is axiom~c in the text. If the replacement lands in a smaller full category of perfect objects, the stronger assertion preceding c follows as well.

The point is that a replacement supplies comparisons inside every such intermediate category. An abstract equivalence after localization has forgotten where its representing zigzags live. One therefore cannot justify c merely by factoring an equivalence of categories through the intermediate localization: a composite of two functors can be an equivalence without either factor being an equivalence. The source recognizes precisely this issue when it distinguishes axioms a and b from its question about c.

Our conclusion about the source's unanswered implication is deliberately narrower than a general claim of openness. The implication from the bare axioms a and b to c is not proved in these pages, and the literature checked here does not settle that exact formulation. The argument above establishes c under the explicit functorial-replacement hypothesis already available in the categorical construction. Grothendieck returns to the issue at manuscript p.~240 and following in this same chapter, with factorization hypotheses and a qualification on p.~241 that every object should be $W$-fibrant over the terminal object; that later discussion is commented with those pages. Brown's and Cisinski's theorems provide other sufficient structures, not a proof that the extra structure follows from a and b alone. Thus the practical expectation on p.~218 has a complete proof in the replacement setting, while its unrestricted logical formulation is kept separate.
\end{dcomment}


On supposera donc, comme dans loc. cit. p. 94-96, donné une classe
\[
\mathcal{L} \subseteq \mathrm{Fl}(\mathcal{M})
\]
de morphismes, jounant le rôle de `flèches de % [sic] jounant -> jouant, faute de frappe
\msp{page 220}
localisation' ou `d'isomorphismes locaux', satisfaisant les conditions suivantes.

\dgrblock{dgrb563}{What the finer relative theory remembers}
\begin{dcomment}{exposition}{p219}{What the finer relative theory remembers}
Grothendieck now wants a family to retain the direction in which its fibres vary. The preceding fibre theory has deliberately made transport invertible. A general coefficient system on a category need not have this property: a specialization can create a component, identify two components, or alter a higher homotopy group. Localizing the total category over the base at all total-space equivalences loses such information. The additional class $W^0_S$ is introduced to preserve it.

There are consequently two distinct localizations to keep in view. The category $(\Cat/S)[W_S^{-1}]$ describes a homotopy type equipped with a map to the homotopy type of $S$. In the ordinary homotopy theory, its derived fibres form a locally constant coherent family over $|NS|$. A directional diagram, by contrast, is a functor from $NS$, or from $NS^{\mathrm{op}}$ according to the chosen variance, to spaces. Its transition maps need not be equivalences. It is the localization of the whole diagram category at objectwise equivalences that records this datum; a diagram in the ordinary category of homotopy types need not retain its coherence. The distinction is the one already made in the commentary on Chapter VI, pp.~64--74, and not a new objection to the fibre construction.

The notation $W^{\mathrm g}_S$ has a concrete test inherited from Chapter VI, manuscript p.~41: a map $u:X\to Y$ over $S$ belongs to this class precisely when
\[
 X/s\longrightarrow Y/s\quad\hbox{belongs to }W
 \qquad(s\in\Ob S).
\]
Here $X/s=X\times_S(S/s)$. Replacing these over-categories by the under-categories $s\backslash X$ gives $W^{\mathrm d}_S$. Taking opposite categories exchanges the two tests, together with the base $S$ and its opposite. It does not give an equivalence between the two theories over a fixed, unmodified $S$. The phrase that neither choice is canonically imposed should be read as a statement about the absence of a preferred orientation in the specified data, not as a theorem that no invariant construction of any kind is possible \citep[Chapter VI, manuscript pp.~41--42]{Grothendieck1990VI}.

The axioms on $\mathcal L$ make the test intrinsic to a chosen class of base changes. Since $\id_S$ is an admissible test, membership in $W^{\mathcal L}_S$ implies membership in $W_S$. Testing against every base change is stronger, giving $W^{\mathrm u}_S\subseteq W^{\mathcal L}_S$. All three agree on the good locally constant representatives by the preceding fibre axiom. Their agreement there says nothing about their agreement on the whole slice: the whole purpose of the finer theory is to retain objects outside that subcategory.
\end{dcomment}


\begin{enumerate}
\item[$(*)$] $\mathcal{L}$ contient les isomorphismes, est stable par composition, et par changement de base.
\end{enumerate}

On pose, pour $S$ dans $\mathcal{M}$,
\[
\mathcal{L}_S \subseteq \mathcal{M}/S \qquad \text{sous-catégorie pleine des } S'/S \text{ telle que } S' \longrightarrow S \text{ \texttt{[soit]} dans } \mathcal{L}.
\]
On définit alors pour tout $S$ dans $\mathcal{M}$
\[
\begin{aligned}
&W^{\mathcal{L}}_S \subseteq \mathrm{Fl}(\mathcal{M}/S)\\
&(X \xrightarrow{u} Y) \in W^{\mathcal{L}}_S \iff \forall\, S' \in \mathcal{L}_S,\ (u_{S'} : X \times_S S' \longrightarrow Y \times_S S') \text{ est dans } W.
\end{aligned}
\]
% foot: Version 30 mars 2003

\printedpage{141}
% head: A. Grothendieck : Dérivateurs, Ch. XII, Édition des Universités de Montpellier II et Paris VII

C'est donc une variante du localiseur $W^{\mathrm{u}}_S$, où maintenant on ne prend comme changements de base que les flèches `$\mathcal{L}$-localisantes', i.e. $\in \mathcal{L}$. On a donc
\[
W^{\mathrm{u}}_S \subseteq W^{\mathcal{L}}_S \subseteq W_S
\]
en tant que parties de $\mathrm{Fl}(\mathcal{M}/S)$, et de même dans toute sous-catégorie pleine $\mathcal{F}$ de $\mathcal{M}/S$. Dans le cas de $\underline{\mathrm{Fib}}_W S$, moyennant l'axiome \circled{b} ci-dessus ($W^{\mathrm{u}}_S = W_S$ dans $\underline{\mathrm{Fib}}_W S$), les trois localiseurs coïncident.

On pose
\[
\mathrm{Ho}^{\mathcal{L}}_W(S) \text{ (ou } \mathrm{Ho}_W(S)\text{)} \ \deff (\mathcal{M}/S)(W^{\mathcal{L}}_S)^{-1},
\]
et on a donné dans loc. cit. (partie II du théorème-scholie) un résultat intéressant sur cet invariant, pour $S$ fixé. Mon propos
\msp{page 221}
présent est de voir ce qui se passe pour $S$ variable. Visiblement, une notion de $W$-type d'homotopie relatif sur un objet quelconque $S$ de $\mathcal{M}$ n'est satisfaisante, que si elle permet des images inverses par des flèches \textit{quelconques}
\[
f : S' \longrightarrow S
\]
dans $\mathcal{M}$, qui doivent définir absolument
\[
f^{*\mathcal{L}}_W \text{ ou } f^*_W : \mathrm{Ho}^{\mathcal{L}}_W(S) \longrightarrow \mathrm{Ho}^{\mathcal{L}}_W(S').
\]
On est donc confronté à la même perplexité que tantôt - et ce n'est \textit{pas,} cette fois, l'introduction de $\underline{\mathrm{Fib}}_W(S)$ qui va sauver la situation, car
\[
(\underline{\mathrm{Fib}}_W S)(\underbrace{W^{\mathcal{L}}_S}_{\substack{= W_S \\ = W^{\mathrm{u}}_S}})^{-1} \longrightarrow (\mathcal{M}/S)(W^{\mathcal{L}}_S)^{-1}
\]
n'a aucune envie, cette fois, d'être une équivalence de catégories, mais bien plutôt une inclusion (pleinement fidèle) des $W$-types d'homotopie relatifs constants sur $S$, dans les $W$-types d'homotopie quelconques. Il faut visiblement un axiome qui remplace l'axiome `de la fibre homotopique'.

\msp{page 222}

\dgrblock{dgrb564}{An endomorphism which becomes invertible only after changing the fibre}
\begin{dcomment}{examples}{p221}{An endomorphism which becomes invertible only after changing the fibre}
A single irreversible transition already explains why the new theory is needed. Consider the category $B\mathbf N$ with one object and endomorphism monoid $(\mathbf N,+)$. Let a covariant diagram of spaces assign a circle to its object and the degree-two map $d:S^1\to S^1$ to its generating endomorphism. This diagram cannot be locally constant: multiplication by $2$ on $\pi_1(S^1)=\mathbf Z$ is not invertible. Objectwise weak equivalences of diagrams preserve this obstruction.

Its homotopy colimit is the mapping torus
\[
 T_d=(S^1\times[0,1])/((z,1)\sim(d(z),0)).
\]
One can obtain this model directly from the bar construction for the free monoid on one generator: the higher simplices express compositions of that generator, and collapse onto the mapping-cylinder description. Van Kampen, using the circle in the fibre and the loop round the parameter as generators, gives
\[
 \pi_1(T_d)=\langle a,t\mid tat^{-1}=a^2\rangle.
\]
The direction chosen for $t$ fixes the displayed convention. This group is nonabelian: the transformations $a(x)=x+1$ and $t(x)=2x$ of the real line satisfy the relation and do not commute. The cellular chain calculation for a mapping torus is the cone of $1-d_*$ on the chains of the fibre. On $H_1(S^1)$ this is multiplication by $-1$, and on $H_0(S^1)$ it is zero. The resulting exact sequence yields
\[
 H_0(T_d;\mathbf Z)=\mathbf Z,\qquad
 H_1(T_d;\mathbf Z)=\mathbf Z,\qquad H_q(T_d;\mathbf Z)=0\quad(q\geq2).
\]
Thus the total space has the homology of the base circle although it retains a substantial nonabelian fundamental group.

Now regard the projection $T_d\to S^1$ solely as a map of homotopy types. Its homotopy fibre is the infinite cyclic cover of the mapping torus. Cutting that cover at the integral parameter values exhibits it as the telescope of successive degree-two maps. Its fundamental group is therefore
\[
 \mathop{\mathrm{colim}}(\mathbf Z\xrightarrow{2}\mathbf Z\xrightarrow{2}\cdots)
       =\mathbf Z[1/2],
\]
and its higher homotopy groups vanish: every map from a sphere has compact image in a finite portion of the telescope, which retracts onto a circle. Transport around the base now acts invertibly, by multiplication by $2$ or its inverse according to orientation. The passage to locally constant relative homotopy types has changed the fibre from $K(\mathbf Z,1)$ to $K(\mathbf Z[1/2],1)$.

This example can be made entirely categorical. Use the one-object groupoid $B\mathbf Z$ and the functor induced by $m\mapsto2m$, then form its Grothendieck construction over $B\mathbf N$. Its nerve has the indicated homotopy-colimit type by the homotopy-colimit theorem \citep[Theorem 1.2]{Thomason1979}; the precise statement and natural comparison are reproduced by \citet[\S2.3, Theorem 2.6]{Rubin2020Elmendorf}. The example extends the earlier walking-arrow tests: it shows, through an explicit group calculation, the information that is lost when all transport is forced to be invertible.
\end{dcomment}


Dans le cas de Cat, la classe $\mathcal{L}$ des isomorphismes locaux (disons) permet de définir, comme ci-dessus, le localiseur $W^{\mathcal{L}}_S = W^{\mathrm{g}}_S$ dans les $\mathrm{Cat}/S$, comme formé des $u : X \longrightarrow Y$ sur $S$ qui sont dans $W$ et le restent par tout changement de base $S' \longrightarrow S$ qui est dans $\mathcal{L}$. Mais il y a lieu de \textit{saturer} $\mathcal{L}$, en regardant la classe $\mathcal{L}^\wedge$ des flèches $f : S \longrightarrow S'$ de Cat telles que le changement de base
\[
f^* : \mathrm{Cat}/S \longrightarrow \mathrm{Cat}/S'
\]
envoie $W^{\mathrm{g}}_S$ dans $W^{\mathrm{g}}_{S'}$. Ce $\mathcal{L}^\wedge$ est tautologiquement stable par composition, et contient les isomorphismes (mieux, contient $\mathcal{L}$ = ensemble des isomorphismes locaux). Il n'est pas
% foot: Version 30 mars 2003
\printedpage{142}
% head: A. Grothendieck : Dérivateurs, Ch. XII, Édition des Universités de Montpellier II et Paris VII
clair, par la définition donnée ici, qu'il soit stable par changement de base - on va renforcer la définition en conséquence, en exigeant, pour que $f : S' \longrightarrow S$ soit dans $\mathcal{L}^\wedge$, que pour tout $T$ sur $S$, considérons $f_T : T' = S' \times_S T \longrightarrow T$, $f^*_T : \mathrm{Cat}/T \longrightarrow \mathrm{Cat}/T'$ soit compatible avec les localiseurs $W^{\mathrm{g}}_T$, $W^{\mathrm{g}}_{T'}$. Alors $\mathcal{L}^\wedge$ est (tout comme $\mathcal{L}$) stable par composition et par changement de base, et de plus il contient $\mathcal{L}$. De plus, pour tout $S$, on a $W^{\mathcal{L}}_S = W^{\mathcal{L}^\wedge}_S$. Ce petit sorite pourrait se faire bien sûr pour tout modéliseur $(\mathcal{M},W)$. Mais la chose intéressante ici, c'est que $\mathcal{L}^\wedge$ contient l'ensemble $\mathrm{Liss}_W$ des flèches $W$-lisses. (J'ignore s'il est égal.) On aurait donc pu
\msp{page 223}
aussi bien partir de l'ensemble de flèches $\mathrm{Liss}_W$, pour définir les localiseurs $W^{\mathrm{g}}_S$ comme
\[
W^{\mathrm{g}}_S = W^{\mathrm{Liss}_W}_S \qquad \text{au lieu de } W^{\mathrm{g}}_S = W^{\mathrm{Isoloc}}_S,
\]
ce qui est équivalent.

\dgrblock{dgrb565}{The maximal class of harmless base changes}
\begin{dcomment}{development}{p222}{The maximal class of harmless base changes}
Grothendieck strengthens his first definition of $\mathcal L^\wedge$ because a useful class of changes of base must remain useful after it is pulled back. The strengthened definition does more than supply that stability. It identifies a largest class of tests which leaves the relative weak equivalences unchanged.

\textbf{Proposition (universal saturation of the tests).}\ours\
Let $\mathcal M$ have pullbacks, let $W$ be a class invariant under isomorphisms of arrows, and let $\mathcal L$ contain all isomorphisms and be closed under composition and pullback. Write $E_S=W^{\mathcal L}_S$. Define $\widehat{\mathcal L}$ to contain $f:S'\to S$ when, for every $T\to S$, pullback along $f_T:S'\times_ST\to T$ sends $E_T$ into $E_{S'\times_ST}$. Then $\widehat{\mathcal L}$ contains $\mathcal L$, is closed under composition and pullback, and
\[
 W^{\widehat{\mathcal L}}_S=E_S,\qquad
 \widehat{\widehat{\mathcal L}}=\widehat{\mathcal L}.
\]
Moreover, if $\mathcal K$ has the same closure properties and $W^{\mathcal K}_S=E_S$ for every $S$, then $\mathcal K\subseteq\widehat{\mathcal L}$.

\textit{Proof.} First take $f$ in $\mathcal L$ and $u$ in $E_T$. To test $f_T^*u$ at $g:R\to S'\times_ST$ in $\mathcal L$, pull back $u$ along the composite $f_Tg$. Both factors belong to $\mathcal L$, so that composite is an admissible test and its pullback of $u$ belongs to $W$. Thus $\mathcal L\subseteq\widehat{\mathcal L}$.

For the pullback of a member $f$ of $\widehat{\mathcal L}$, every subsequent test is already among the base changes in its definition. This proves pullback stability. For composable members $f:S'\to S$ and $g:S''\to S'$, each base change of $fg$ factors as the corresponding base change of $f$ followed by one of $g$. Their pullback functors preserve the respective $E$ classes in succession. This proves composition stability; identities satisfy the same condition immediately.

Since $\mathcal L\subseteq\widehat{\mathcal L}$, testing against $\widehat{\mathcal L}$ can only shrink the resulting class, so $W^{\widehat{\mathcal L}}_S\subseteq E_S$. Conversely, for $u\in E_S$ and $f:S'\to S$ in $\widehat{\mathcal L}$, the defining condition with $T=S$ gives $f^*u\in E_{S'}\subseteq W_{S'}$. Hence $u$ passes every $\widehat{\mathcal L}$ test. This proves equality. Repeating the construction consequently uses exactly the same family $E_S$, so gives exactly the same class $\widehat{\mathcal L}$.

Finally, take $f\in\mathcal K$ and any $T\to S$. If $u\in E_T=W^{\mathcal K}_T$, its pullback along $f_T$ passes every $\mathcal K$ test because such a test followed by $f_T$ is again in $\mathcal K$. The pulled-back map therefore belongs to $W^{\mathcal K}_{S'\times_ST}=E_{S'\times_ST}$. This is precisely $f\in\widehat{\mathcal L}$.\hfill$\square$

The proof does not use a model structure, a factorization axiom, or even two-out-of-three for $W$. Invariance under arrow isomorphisms only ensures independence of the choices of pullbacks and their canonical associativity isomorphisms. If $W$ is weakly saturated, each $E_S$ is weakly saturated because it is an intersection of inverse images of $W$ under functors. We give the maximality and idempotence explicitly as a development of the source's sorite; we did not locate this exact formulation in the checked accounts of smooth functors and derived base change. The label makes no claim of priority. The variance of the first displayed $f$ on manuscript p.~222 is reversed in the edition; the strengthened definition, and the argument above, use $f:S'\to S$.
\end{dcomment}


D'autre part, en travaillant avec les localisés $\mathrm{HOT}^\circ_W(S) = (\mathrm{Cat}/S)(W^{\mathrm{g}}_S)^{-1}$, le rôle de $\underline{\mathrm{Fib}}_W S$ pour la théorie des $\mathrm{HOT}^{\mathrm{lc}}_W(S)$ est joué ici par $\underline{\mathrm{Prop}}_W(S)$. On voit donc que dans la théorie des $W$-types d'homotopie relatifs non localement constants, mode $\mathrm{HOT}^\circ_W$, les ensembles de flèches $\mathrm{Liss}_W$ et $\mathrm{Prop}_W$ jouent un rôle en quelque sorte complémentaire : les flèches $f : S' \longrightarrow S$ de $\mathrm{Liss}_W$ sont (pratiquement) celles qui permettent directement de définir
\[
f^{*\circ}_W : \mathrm{HOT}^\circ_W(S) \longrightarrow \mathrm{HOT}^\circ_W(S')
\]
comme induit par
\[
f^* : (\mathrm{Cat}/S, W^{\mathrm{g}}_S) \longrightarrow (\mathrm{Cat}/S', W^{\mathrm{g}}_{S'}),
\]
tandis que les flèches de $\mathrm{Liss}_W$ permettent de définir pour tout $S$ dans Cat la sous-catégorie pleine
\[
\underline{\mathrm{Prop}}_W S \subseteq \mathrm{Cat}/S,
\]
donnant lieu à une équivalence de catégories
\[
\left\{
\begin{aligned}
&(\underline{\mathrm{Prop}}_W S)(W^{\mathrm{u}}_S)^{-1} \xrightarrow{\sim} \underbrace{(\mathrm{Cat}/S)(W^{\mathrm{g}}_S)^{-1}}_{= \mathrm{HOT}^0_W(S)} % ?? uncertain: exposant de HOT sous l'accolade -- 0 (zéro) ou \circ
\\
&\text{avec}\\
&W^{\mathrm{u}}_S = W^{\mathrm{g}}_S \text{ sur } \underline{\mathrm{Prop}}_W S
\end{aligned}
\right.
\]
(ce qui visiblement est le substitut annoncé de ``l'axiome de la fibre $W$-homotopique''). Cette
\msp{page 224}
équivalence, qui permet de décrire $\mathrm{Hot}^\circ_W(S)$ en termes de $\underline{\mathrm{Prop}}_W S \subseteq \mathrm{Cat}/S$ et d'un localiseur $W^{\mathrm{u}}_S$ \textit{stable par} \textbf{tout} \textit{changement de base,} permet d'autre part de définir la loi contravariante attendue pour $\mathrm{HOT}^\circ_W(S)$.
% foot: Version 30 mars 2003

\dgrblock{dgrb566}{The saturation question has an affirmative answer in Cat}
\begin{dcomment}{later}{p223}{The saturation question has an affirmative answer in Cat}
The question in parentheses on p.~222 can be answered using the later characterization of smooth functors. With Grothendieck's localizer $W^{\mathrm g}_S$ and the universal definition of $\mathcal L^\wedge$, one has
\[
 \mathcal L^\wedge=\mathrm{Liss}_W.
\]
This is an implication of an established theorem, not a new smoothness criterion claimed by the commentators. A functor is $W$-smooth exactly when all its base changes preserve $W$-aspherical functors \citep[Theorem 1.2(c)]{Maltsiniotis2005Asphericite}.

Here is the precise translation. A functor $v:A\to T$, regarded as a morphism from $A/T$ to the terminal object $T/T$ of the slice, belongs to $W^{\mathrm g}_T$ exactly when $A/t\to T/t$ belongs to $W$ for every $t$. Since $T/t$ has a terminal object, this says exactly that every $A/t$ is $W$-aspherical. Thus it says that $v$ is a $W$-aspherical functor. If $f:S'\to S$ belongs to $\mathcal L^\wedge$, apply the preservation condition after every $T\to S$ to these particular arrows $v:A\to T$. Their pullbacks are again $W$-aspherical. The cited characterization implies that $f$ is $W$-smooth. The opposite inclusion is the smooth base-change property used in the source. This proves the equality, including its universal quantifier.

There is an editorial distinction on this page that should not obscure the argument. The sentence introducing the full subcategory $\underline{\mathrm{Prop}}_W S$ repeats ``Liss'' where its role requires ``Prop''. The underbrace subsequently prints $\mathrm{HOT}^{0}_W(S)$, with a zero, although the theory throughout the paragraph is $\mathrm{HOT}^{\circ}_W(S)$. Both readings are retained and recorded; neither denotes a third relative theory.
\end{dcomment}


\printedpage{143}
% head: A. Grothendieck : Dérivateurs, Ch. XII, Édition des Universités de Montpellier II et Paris VII

D'autre part, quand on échange le rôle de $\mathrm{Liss}_W$ et $\mathrm{Prop}_W$, on trouve la théorie complémentaire du $\mathrm{HOT}_W(S)$,
\[
\left\{
\begin{aligned}
&\mathrm{HOT}_W(S) \deff (\mathrm{Cat}/S)(W^{\mathrm{d}}_S)^{-1} \xleftarrow{\sim} (\underline{\mathrm{Liss}}_W S)(W^{\mathrm{u}}_S)^{-1},\\
&\text{sur } \underline{\mathrm{Liss}}_W S,\ W^{\mathrm{d}}_S = W^{\mathrm{u}}_S,
\end{aligned}
\right.
\]
la description de $\mathrm{HOT}_W(S)$ en termes de $\underline{\mathrm{Liss}}_W S$ et de $W^{\mathrm{u}}_S$ permet d'en définir le caractère contravariant en $S$, pour \textit{toute} flèche $f : S' \longrightarrow S$ dans Cat. Tandis que si on veut pouvoir passer aux catégories de fractions sur $\mathrm{Cat}/S$ et $\mathrm{Cat}/S'$ tout entier (via $W^{\mathrm{d}}_S$, $W^{\mathrm{d}}_{S'}$), il faut se limiter aux changements de base $W$-propres.

Du coup, me vient la perplexité si cette complémentaire, voire cette dualité, entre deux ensembles de flèches, les `lisses' et les `propres', ne serait pas un fait général qui se retrouvait tout au moins dans les catégories de modèles proche de la description de HOT ou d'un $\mathrm{Hot}_W$.

\msp{page 225}

Il se pourrait bien que dans $\boldsymbol{\Delta}^\wedge$ on ait une dualité parfaite entre ces deux espèces de notions (il faudra que j'examine la chose), peut-être même dans toute catégorie $\mathcal{A}^\wedge$ (mais dans $\boldsymbol{\Delta}^\wedge$ il y a un automorphisme involutif $K \longmapsto K^\circ$, correspondant au $X \longmapsto X^\circ$ dans Cat \dots). Et après tout, même dans le contexte des espaces topologiques, les notions de propreté d'une part, d'asphéricité locale d'autre part, pour une flèche $X \longrightarrow Y$, sont bien définies, et pourraient bien s'ajuster comme le font les notions correspondantes dans Cat, et (qui sait) donner lieu à deux versions complémentaires, ou `duales' l'une de l'autre, d'une notion de type d'homotopie relatif.

N'en ayant pas le cœur net, je suis un peu indécis comment présenter une description axiomatique de la théorie des $\mathrm{Ho}^{\mathcal{L}}_W(S)$, pour une catégorie de modèles générale $(\mathcal{M},W)$. Dois-je me donner d'emblée
\[
\mathrm{Liss}, \mathrm{Prop} \subseteq \mathrm{Fl}(\mathcal{M}),
\]
stables par composition, par changement de base, contenant les isomorphismes, et tels que (peut-être)
\[
\mathrm{Liss} \cap \mathrm{Prop} \subseteq \mathrm{Fib}_W,
\]
\msp{page 226}
mais surtout ils permettent de définir, pour $S$ de $\mathcal{M}$
\[
\left\{
\begin{array}{lcccl}
W^{\mathrm{g}}_S & \deff & W^{\mathrm{Liss}}_S & \subseteq & \mathcal{F}(\mathcal{M}/S)\\
W^{\mathrm{d}}_S & \deff & W^{\mathrm{Prop}}_S & \subseteq & \mathcal{F}(\mathcal{M}/S)
\end{array}
\right.
\]
avec les propriétés suivantes :

\dgrblock{dgrb567}{Smoothness and properness depend on a coefficient theory}
\begin{dcomment}{context}{p225}{Smoothness and properness depend on a coefficient theory}
Grothendieck's proposed pair of classes is meant to encode two ways in which a square admits base change. The geometric names come from proper and smooth base-change theorems in étale cohomology, but their categorical definitions concern the invertibility of specific comparison transformations. This is why exchanging the two classes in the text also exchanges the direction of the relative theory. The comparison with geometry is a source of the formalism; it does not identify a categorical proper functor with a proper continuous map.

Once a coherent indexed category $\mathscr E:S\mapsto\mathscr E(S)$ is given, an adjunction $f_!\dashv f^*$ produces a base-change mate for each pullback square, and an adjunction $f^*\dashv f_*$ produces its other version. Requiring the relevant mate to be invertible after all further base changes defines a stable class. Thus a meaningful duality has three pieces of data: the coefficient theory, its two adjunctions where they exist, and the chosen variance. Merely knowing which arrows of an ambient category are weak equivalences does not supply all these pieces.

A modern formulation makes precisely this dependence explicit. \citet[Definition 9 and Proposition 10]{AnelWeinberger2024} associate smooth and proper classes to an indexed category over a base with finite limits; these are the largest classes along which the corresponding indexed sums and products exist with stable Beck--Chevalley compatibility. Their examples in Sections~3.2--3.3 include categorical and topological versions. This gives a general setting for the complementarity sought here. It is not a proof that every proposed $(\mathcal M,W,\mathcal L)$ satisfies the localization equivalences of p.~226.

The involution on simplicial sets mentioned by Grothendieck is entirely concrete. Reverse the order on each finite ordinal: for $\alpha:[m]\to[n]$, set $\alpha^{\mathrm{rev}}(i)=n-\alpha(m-i)$. These maps define an involutive automorphism of $\Delta$, hence, by precomposition, an involution $K\mapsto K^{\mathrm{op}}$ of simplicial sets. The identity $N(C^{\mathrm{op}})=(NC)^{\mathrm{op}}$ is immediate from a string of composable arrows. The corresponding assertion for a general presheaf category requires an appropriate duality of its indexing category; none is part of the present hypotheses. Consequently the speculation about all presheaf categories should remain a programme with specified additional structure, rather than an unqualified theorem or a single well-posed open conjecture.
\end{dcomment}


\begin{enumerate}
\item[{\circled{a}}] Si $f : S' \longrightarrow S$ est dans Prop (resp. Liss), alors $f^* : \mathcal{M}/S \longrightarrow \mathcal{M}/S'$ transforme $W^{\mathrm{d}}_S$ dans $W^{\mathrm{d}}_{S'}$ (resp. $W^{\mathrm{g}}_S$ dans $W^{\mathrm{g}}_{S'}$),
\end{enumerate}
% foot: Version 30 mars 2003
\printedpage{144}
% head: A. Grothendieck : Dérivateurs, Ch. XII, Édition des Universités de Montpellier II et Paris VII
\noindent ce qui est une pure tautologie, et ne concerne pas encore les relations de Prop avec Liss, mais par contre :

\begin{enumerate}
\item[{\circled{b}}]
\[
\left\{
\begin{aligned}
&(\underline{\mathrm{Liss}}S)(W^{\mathrm{u}}_S)^{-1} \xrightarrow{\sim} (\mathcal{M}/S)(W^{\mathrm{d}}_S)^{-1} \text{ \texttt{[est]} une équivalence}\\
&\text{et } W^{\mathrm{u}}_S = W^{\mathrm{d}}_S \text{ sur } \underline{\mathrm{Liss}}S,
\end{aligned}
\right.
\]
et dualement
\[
\left\{
\begin{aligned}
&(\underline{\mathrm{Prop}}S)(W^{\mathrm{u}}_S)^{-1} \xrightarrow{\sim} (\mathcal{M}/S)(W^{\mathrm{g}}_S)^{-1} \text{ \texttt{[est]} une équivalence}\\
&\text{et } W^{\mathrm{u}}_S = W^{\mathrm{g}}_S \text{ sur } \underline{\mathrm{Prop}}S.
\end{aligned}
\right.
\]
\end{enumerate}

De plus, il y a lieu d'énoncer \circled{c} deux propriétés d'adjonction (duales l'une de l'autre), pour lesquelles je renvoie à VI, théorème-scholie 2, p. 108-113. Il se pose ici la question subsidiaire de simplifier les conditions de validité de la formule d'adjonction, un peu hétéroclites dans loc. cit.

Il se pose d'autre part la question de faire la théorie des lois covariantes $f_!$ et $f_*$ pour $\mathrm{Ho}^{\mathrm{lc}}_W(S)$ et pour $\mathrm{Ho}^{\mathrm{Prop}}_W(S)$, $\mathrm{Ho}^{\mathrm{Liss}}_W(S)$.

\msp{page 227}

\dgrblock{dgrb568}{Postcomposition already gives a direct image}
\begin{dcomment}{development}{p226}{Postcomposition already gives a direct image}
Grothendieck postpones the direct images because the hard part of his construction has been the inverse image. One part of his expectation follows from the elementary structure of the test classes alone. It is useful to isolate it before bringing in a theory of derived adjunctions.

\textbf{Proposition (descent of postcomposition).}\ours\
Under the hypotheses of the universal-saturation proposition, every $f:S'\to S$ induces a functor
\[
 f_\#: (\mathcal M/S')[ (W^{\mathcal L}_{S'})^{-1}]
       \longrightarrow(\mathcal M/S)[(W^{\mathcal L}_S)^{-1}]
\]
by postcomposition with $f$. These functors compose canonically. If $f$ belongs to $\widehat{\mathcal L}$, ordinary pullback also descends, and the ordinary adjunction $f_\#\dashv f^*$ descends with its unit and counit.

\textit{Proof.} Take $u:X\to Y$ over $S'$ in $W^{\mathcal L}_{S'}$. For a test $T\to S$ in $\mathcal L$, its pullback $T'=S'\times_ST\to S'$ again belongs to $\mathcal L$. The underlying arrow obtained by pulling back $f_\#u$ to $T$ is canonically the underlying arrow obtained by pulling back $u$ to $T'$. It belongs to $W$. Thus postcomposition sends the required weak equivalences to weak equivalences and factors through localization. The equalities of ordinary postcomposition give the composition isomorphisms, with their associativity.

When $f\in\widehat{\mathcal L}$, its pullback functor is homotopical by definition. Both functors in the ordinary slice adjunction therefore descend. A natural transformation between relative functors descends as well: its naturality equation for an inverted arrow remains valid after that arrow is inverted. Apply this observation to the unit and counit. Their two triangle identities are equalities before localization and remain equalities afterwards, proving the adjunction.\hfill$\square$

This argument establishes the direct-image functor for every $f$, without a functorial replacement. It does not by itself establish an adjunction to an inverse image obtained indirectly from proper or smooth representatives when raw $f^*$ is not homotopical. In that case one must compare the replacement with the ordinary slice adjunction. The source's proposed axioms are designed to make this possible; equality of the three localizers on a full subcategory alone is not the comparison. The earlier Chapter VI commentary, at manuscript pp.~137--145, supplies that comparison in its categorical setting. The proposition above records the additional generality that follows directly from the present definition, without claiming a new priority for descent of adjunctions.
\end{dcomment}


J'ai bon espoir que le foncteur $f_!$ peut être décrit, et son adjonction avec $f^*$ établie, à aussi bon compte dans ce contexte général que dans Cat. Mais je n'ai pas envie de m'y arrêter à présent.

\section*{6 Cofibre $W$-homotopique dans $\mathrm{Cat}$.}

Par contre, j'ai envie de regarder ici sous quelles conditions sur le localiseur fondamental $W$, il existe dans $(\mathrm{Cat},W)$ une théorie duale de celle de la fibre $W$-homotopique.

La question-clef est, bien sûr, celle d'avoir une bonne prise sur les $W$-cofibrations. Or je ne connais jusqu'à présent \textit{que} les immersions ouvertes de Dwyer, les immersions fermées de Dwyer, et ce qui s'en déduit par composition etc. Le travail de \textsc{Thomason} peut motiver l'espoir que la connaissance de ces immersions-là, qui semblent ici jouer un rôle similaire (dual) à celui des morphismes $W$-parfaits, suffira pour faire une bonne théorie des `cofibres $W$-homotopiques'.

% trois traits verticaux imprimés dans la marge gauche de ce paragraphe
{\reversemarginpar\marginnote{\textbar\textbar\textbar}}Je suppose dans toute cette section que $W$ satisfait W(7 bis) (axiome des carrés cocartésiens), pour avoir des bonnes propriétés des immersions ouvertes de Dwyer (\footnote[124]{\begin{enumerate}
\item[a)] Ce sont des $W$-cofibrations (résulte de b) et c)).
\item[b)] Si
\end{enumerate}}), et plus généralement, pour les $W$-cofibrations. % ?? uncertain: coupure en fin de ligne "co-cartésiens" -- cocartésiens ou co-cartésiens
% foot: Version 30 mars 2003

% ---- from Der-ChXIIc-30-03-03_p031-035.tex
\printedpage{145}
% head: A. Grothendieck : Dérivateurs, Ch. XII, Édition des Universités de Montpellier II et Paris VII
% foot: Version 30 mars 2003

\msp{page 228}
Il se pose donc la question de la factorisation duale d'une flèche $f : X \longrightarrow Y$ de Cat en
\[
\begin{tikzcd}[column sep=huge]
X \arrow[r, hook, "W\text{-cofibration}", "i_f"'] & \tilde{Y} \arrow[r, "\in W", "p_f"'] & Y.
\end{tikzcd}
\]

\dgrblock{dgrb569}{Why a longer categorical interval is needed}
\begin{dcomment}{exposition}{p227}{Why a longer categorical interval is needed}
The passage from fibres to cofibres begins with a manageable class of maps along which strict gluing has the correct homotopical effect. For an inclusion $A\hookrightarrow B$, the Dwyer condition provides exactly such control. In its original form it says that $A$ is a full sieve, and that its inclusion in the cosieve $Z\subseteq B$ generated by $A$ has a right adjoint $r:Z\to A$ restricting to the identity of $A$. The counit $ir\to\id_Z$ is a categorical deformation of the neighbourhood $Z$ onto $A$. This accounts for the usefulness of Dwyer maps in gluing \citep[Definition 4.1 and Proposition 4.3]{Thomason1980}.

For $I=\Phi_n$, the endpoints $0$ and $2n$ form a sieve. Their generated cosieve consists of the four vertices $0,1,2n-1,2n$ when $n\geq2$. It is the disjoint union of two arrows, $0\to1$ and $2n\to2n-1$. Send the first arrow to the endpoint $0$ and the second to $2n$. This is the required right adjoint. For $n=1$, the generated cosieve is the whole cospan $0\to1\leftarrow2$. A right adjoint to the inclusion of the two endpoints would have to assign to $1$ an endpoint receiving all arrows from either endpoint, which is impossible. The lower bound $n\geq2$ therefore has an exact purpose: the two deformation neighbourhoods must not meet.

Taking a product with $X$ preserves the sieve and cosieve conditions and the adjunction, now with right adjoint $\id_X\times r$. The nerve of $\Phi_n$ realizes to a subdivided interval, so $X\times I\to X$ is an ordinary weak equivalence. The stronger description as $W_\omega$-perfect records the corresponding universal statement in the categorical theory. Thus the factorization of the codiagonal has both the disjoint endpoint inclusion needed for the mapping cylinder and the weakly contractible cylinder direction needed for its retraction.

The more general fibred base change asserted in the source has a similarly explicit proof. Pull a Dwyer inclusion $A\hookrightarrow B$ back along a Grothendieck fibration $E\to B$. Over its generated cosieve $Z$, lift each counit arrow $r(b)\to b$ cartesianly to an arrow $r'(e)\to e$. Choose identity lifts over $A$. The cartesian universal property defines $r'$ on morphisms, with functoriality following from uniqueness. For an object $a'$ over $A$, an arrow $a'\to e$ factors uniquely through that cartesian lift: its image in $B$ first factors uniquely through the counit by the original adjunction. Hence $r'$ is right adjoint to the pulled-back inclusion. Moreover every object over $Z$ receives such a lifted arrow from an object over $A$, so the pulled-back cosieve is exactly the generated one. This proves the asserted stability and explains why an arbitrary base change, without a lifting condition, was not claimed.

The terminology must remain precise. Grothendieck's $W$-cofibration is a map whose cobase changes preserve the relevant $W$-equivalences. It is a homotopical gluing condition, not the definition by a lifting property in Thomason's model structure. Dwyer maps supply examples; one should not identify the two classes. Nor should one assert that every Thomason cofibration is a Dwyer map: \citet[pp.~227--231]{Cisinski1999Dwyer} gave a counterexample to closure of the original Dwyer maps under retracts and repaired the properness argument. That correction leaves the explicit mapping-cylinder construction below intact.
\end{dcomment}


Les résultats de \textsc{Thomason} assurent que ça existe, avec $i$ une immersion ouverte de Dwyer, et $p$ une $W$-fibration - du moins dans le cas $W = W_\infty$. Je vais essayer de le voir directement, en calquant dualement la construction de \textsc{Cartan-Serre}.

\dgrblock{dgrb570}{An interval with room at both ends}
\begin{dcomment}{exposition}{p228:assurent que ça existe, avec}{An interval with room at both ends}
Grothendieck wants to replace every map by an inclusion along which gluing respects the chosen weak equivalences. This will make a homotopy cofibre accessible to the same relative methods that have already produced homotopy fibres. The interval is therefore being chosen for a relative property of its two endpoints. Its contractibility supplies only the other half of the construction.

We recall precisely the relevant meaning of cofibration. A map $i:A\to B$ is a $W$-cofibration when the pushout functor
\[
 A\backslash\mathcal M\longrightarrow B\backslash\mathcal M,
 \qquad (A\to T)\longmapsto(B\to B\amalg_A T),
\]
preserves morphisms belonging to $W$. This is the definition on manuscript pp.~109--110 \citep[Definition~2(c)]{Grothendieck1990XII}. It includes every subsequent change of the object under $A$. Associativity of pushouts immediately makes this class stable under pushout and composition. A Quillen cofibration is defined by a lifting property; the relation between the two notions requires a properness or gluing theorem. In the present argument, the required gluing property is part of the definition of $\operatorname{Cof}_W$.

The geometry of a Dwyer map provides a tractable supply of such inclusions. Write an open immersion as a full sieve $A\subseteq B$: an arrow of $B$ whose target belongs to $A$ has its source in $A$. It is Dwyer when the inclusion of $A$ into its generated cosieve $Z_BA$ has a right adjoint $r$ with $r|_A=1_A$. The counit gives arrows $r(b)\to b$, coherently in $b$. Thus objects immediately outside $A$ have canonical points of attachment in $A$; this is the categorical structure used when the inclusion is pushed out \citep[Definition~4.1]{Thomason1980}.

For the fence $I=\Phi_n$, the cosieve generated by its endpoints is
\[
 \{0\longrightarrow1\}\ \amalg\
 \{2n\longrightarrow2n-1\},\qquad n\geq2.
\]
The right adjoint sends the first component to $0$ and the second to $2n$. The unit is the identity, and the two displayed arrows give its counit. Taking its product with $X$ gives the required adjoint for $X\amalg X\subseteq X\times I$. The same argument also explains the lower bound on $n$. For $n=1$, the vertex $1$ receives arrows from both $0$ and $2$. If a right adjoint existed, its value at $1$ would have to represent the functor on the discrete two-point category which takes both objects to a singleton. Neither object represents that functor. The two endpoint neighbourhoods need a vertex between them before the Dwyer condition can hold.

The notation $W_\omega$ belongs to the earlier construction of a minimal localizer from saturation, the relative Theorem~A condition, and the terminal-object axiom \citep[manuscript p.~14]{Grothendieck1990XIIa}. The elementary contraction of this finite fence already belongs to that small class. One can remove the outermost endpoint, then its newly terminal neighbour, and repeat; each removal is witnessed by an adjunction between the remaining posets. Products preserve these adjunctions. Consequently $X\times I\to X$ is a homotopy equivalence for every $X$, and remains one after an arbitrary change of its base $X$.
\end{dcomment}


Il s'agit donc d'abord de factoriser
\[
X \amalg X \xrightarrow{\ \mathrm{codiag}_X\ } X,
\]
et pour ceci je choisis un intervalle fixe
\[
I \ = \ \Phi_n \ = \ (0 \longrightarrow 1 \longleftarrow 2 \longrightarrow \cdots \longrightarrow 2n-1 \longleftarrow 2n)
\]
avec
\[
n \geq 2
\]
de sorte que l'inclusion
\[
\varepsilon = \text{catégorie discrète } \{0,1\} \xrightarrow{\partial} I, \qquad \partial(0) = 0, \quad \partial(1) = 2n,
\]
soit une immersion ouverte de Dwyer. Cette propriété est stable par changement de base Cat-fibrant, en particulier en prenant un produit
\[
\underbrace{X \times \varepsilon}_{\simeq X \coprod X} \xrightarrow{i_X} X \times I,
\]
% continuation of a footnote begun on the preceding page of the edition (no footnote mark on this page)
\begingroup\renewcommand{\thefootnote}{}\footnotetext{%
\[
\begin{tikzcd}[ampersand replacement=\&]
U \arrow[r, hook, "i"] \arrow[d, "p"'] \& X \arrow[d, "q"] \\
U' \arrow[r, hook, "i'"] \& X'
\end{tikzcd}
\]
avec $i$ immersion ouverte de Dwyer, alors
\[
\begin{aligned}
i \in W \quad &\Longleftrightarrow \quad i' \in W \\
p \in W \quad &\Longleftrightarrow \quad q \in W,
\end{aligned}
\]
j'ai donc besoin de W(7 bis) \textit{et} même de W(7) pour la deuxième implication.
\begin{enumerate}
\item[c)] Les immersions ouvertes de Dwyer \texttt{[sont]} stables par cochangement de base et par composition, contiennent les isomorphismes.
\end{enumerate}}\endgroup
\printedpage{146}
% head: A. Grothendieck : Dérivateurs, Ch. XII, Édition des Universités de Montpellier II et Paris VII
% foot: Version 30 mars 2003
\msp{page 229}
d'autre part
\[
X \times I \xrightarrow{\mathrm{pr}_1} X
\]
est $W_\omega$-parfait, d'où la factoriastion cherchée % [sic] factoriastion -> factorisation, misprint
\[
X \amalg X \xrightarrow{i_X} X \times I \ \xrightarrow{\ p_X = \mathrm{pr}_1^X\ } X
\]
- il n'y a pas plus simple! Comment l'utiliser pour construire la factorisation d'une flèche $f : X \longrightarrow Y$ quelconque, en dualisant les constructions relatives aux espaces de chemins? On définit $Z(f)$ par le carré cocartésien dans Cat
\[
\begin{tikzcd}[column sep=5em, row sep=4em]
X \arrow[r, dotted, "\mathrm{inj}_1^X"'] & X \times I \arrow[r, "f'"] & Z(f) \arrow[d, dotted, bend left=45, "\,p_f"] \\
& X \arrow[u, hook, "\mathrm{inj}_0^X"'] \arrow[r, "f"] & Y, \arrow[u, hook, "W"']
\end{tikzcd}
\]
on a
\[
Z(f) \xrightarrow{\alpha} Y \qquad \text{induit par } \mathrm{pr}_1 : X \times I \longrightarrow X,
\]
et la factorisation cherchée de $f$ est
\[
X \ \xrightarrow{\ i_f = f' \circ \mathrm{inj}_1^X\ } Z(f) \xrightarrow{p_f} Y.
\]

\dgrblock{dgrb571}{Conservative gluing is an additional hypothesis}
\begin{dcomment}{development}{p228}{Conservative gluing is an additional hypothesis}
The continued footnote~124 is stronger than the gluing theorem usually used in homotopy theory. It prints two biconditionals for a Dwyer pushout. Its stated hypotheses include the conservative axioms W(7 bis), and for the second assertion W(7), inherited from earlier pages \citep[Chapter XII, manuscript pp.~76--78]{Grothendieck1990XII}. The conditional assertion should not be silently replaced by a theorem about $W_\infty$. In fact the reverse implications fail in the ordinary homotopy theory, even for finite posets and an open cover.

\textbf{Example (a conservative converse fails).}\ours\
There is a pushout square of finite posets
\[
\begin{tikzcd}
 A \arrow[r,hook,"i"] \arrow[d,hook,"p"'] & B \arrow[d,hook,"q"]\\
 C \arrow[r,hook,"i'"'] & D
\end{tikzcd}
\]
in which $i$ and $p$ are Dwyer maps, $i'$ and $q$ belong to $W_\infty$, and neither $i$ nor $p$ belongs to $W_\infty$.

\textit{Construction and proof.} Start with a connected finite acyclic complex $K$ with nontrivial fundamental group. For completeness, the presentation complex with two generators and two relations
\[
 \langle a,b\mid a^5b^{-3},\ b^3(ab)^{-2}\rangle
\]
has cellular boundary matrix $\left(\begin{smallmatrix}5&-2\\-3&1\end{smallmatrix}\right)$, of determinant $-1$, so its reduced integral homology vanishes. Sending $a$ to $(12345)$ and $b$ to $(253)$ defines a nontrivial homomorphism of its fundamental group to $A_5$: $a^5=b^3=(ab)^2=1$, since $ab=(12)(45)$. Hence it is not contractible. This is the acyclic complex of \citet[Example 2.38]{Hatcher2002}. Choose a finite simplicial model of the same homotopy type. Such a model can be obtained by replacing the two attaching maps by simplicial maps between subdivided finite graphs and using their finite triangulated mapping cylinders before adjoining the two discs.

Let $P(L)$ be the poset of nonempty simplices of a finite simplicial complex $L$, ordered by inclusion. Set
\[
 A=P(K),\quad B=P(C_+K),\quad C=P(C_-K),\quad D=P(\Sigma K).
\]
The two cones are subcomplexes of the suspension and meet in $K$. Their face posets are sieves, their union is $D$, and there are no arrows between the two complements. Their ordinary pushout is therefore $D$. For $A\hookrightarrow B$, the cosieve generated by $A$ consists of all simplices except the one-vertex apex simplex. On this cosieve, deleting the apex defines a right adjoint to the inclusion of $A$: a simplex of $K$ is a face of $\sigma$ if and only if it is a face of $\sigma$ with the apex removed. This verifies the Dwyer condition. The proof for $A\hookrightarrow C$ is identical.

The realization of $NP(L)$ is the barycentric subdivision of $L$. Both $B$ and $C$ are therefore contractible. The suspension of a connected space is simply connected, and its reduced homology is the shifted reduced homology of the space. Thus $\Sigma K$ is simply connected and acyclic. If it had a first nonzero homotopy group, the Hurewicz theorem would identify that group with a nonzero homology group. It has none, and Whitehead's theorem makes $\Sigma K$ contractible. Consequently $i'$ and $q$ are weak equivalences, whereas $i$ and $p$ cannot be, because $|NA|$ has nontrivial $\pi_1$.\hfill$\square$

This application of a classical acyclic example is included to expose the strength of the assumed axiom, not as a claim of novelty. We checked the original Dwyer pushout theorem, its retract correction, and the later categorical refinement; they assert preservation, not these converse implications. The earlier W(7 bis), formulated for a union of open sieves, already fails on the displayed union. The claim in Chapter XII, manuscript p.~78, that $W_\infty$ satisfies the stronger W(7), cannot supply it. The present proof of the mapping cylinder and Proposition~16 uses only the forward gluing implications, which are valid for $W_\infty$. Keeping this distinction allows the construction to survive with its natural hypotheses.
\end{dcomment}


Je dis que
\[
\left\{
\begin{aligned}
& i_f \text{ est immersion ouverte de Dwyer, donc une } W\text{-cofibration.} \\
& p_f \in W.
\end{aligned}
\right.
\]

\msp{page 230}
On voit sur le diagramme plus complet
\[
\begin{tikzcd}[column sep=6.5em, row sep=4.5em]
& X \times I \arrow[r, "f'"] \arrow[dd, dotted, bend left=45, "\mathrm{pr}_1"{pos=.60}] \arrow[dr, phantom, "\substack{\text{cocart.}\\ \text{et cart.}}"] & Z(f) \arrow[dd, dotted, bend left=60, "p_f"] \\
X \arrow[ur, dotted, "\mathrm{inj}_1^X"] \arrow[r, hook, "i_1^X"] & X \amalg X \arrow[u, hook, "i_X"'] \arrow[r, "f \coprod \mathrm{id}_X"] \arrow[dr, phantom, "\substack{\text{cocart.}\\ \text{et cart.}}"] & Y \amalg X \arrow[u, hook, "j"'] \\
& X \arrow[u, hook, "i_0^X"'] \arrow[r, "f"] & Y, \arrow[u, hook]
\end{tikzcd}
\]
\printedpage{147}
% head: A. Grothendieck : Dérivateurs, Ch. XII, Édition des Universités de Montpellier II et Paris VII
% foot: Version 30 mars 2003
\[
i_f = f' \circ \mathrm{inj}_1^X = f' i_X i_1^X = j \ \ \underbrace{(f \amalg \mathrm{id}_X) i_1^X}_{\substack{\text{inclusion canonique}\\ \text{de } X \text{ dans}\\ Y \coprod X}} \ \ ,
\]
c'est \texttt{[une]} immersion ouverte de Dwyer, comme composé de $Y \longrightarrow Y \amalg X$ et $j : Y \amalg X \longrightarrow Z(f)$ qui le sont ($j$ obtenu par cochangement de base à partir de $i_X$, qui est immersion ouverte de Dwyer).

D'autre part, $X \xrightarrow{\mathrm{inj}_0^X} X \times I$ est \texttt{[une]} immersion ouverte de Dwyer et dans $W$, donc coüniversellement dans $W$, donc $Y \longrightarrow Z(f)$ est elle aussi dans $W$, donc aussi la rétraction $p_f$.

\dgrblock{dgrb572}{A cylinder and a circle in finite categories}
\begin{dcomment}{examples}{p230:donc aussi la rétraction}{A cylinder and a circle in finite categories}
The mapping cylinder puts the preceding two properties to work simultaneously. Its construction is worth following through a small example, because the example also shows what a homotopy cofibre retains that a strict categorical quotient can lose. First write its two structural maps without the changing names in the display:
\[
 Z(f)=(X\times I)\amalg_{X\times\{0\}}Y,
 \qquad s_f:Y\longrightarrow Z(f),
 \qquad p_f:Z(f)\longrightarrow Y.
\]
The map $p_f$ is induced by $f\operatorname{pr}_1$ on $X\times I$ and by $1_Y$ on $Y$. Hence $p_fs_f=1_Y$, while the free endpoint defines $i_f:X\to Z(f)$ with $p_fi_f=f$. These equations also show that the construction is functorial in a commutative square of functors: the two maps on the pieces agree on the identified endpoint and therefore induce a unique map on their pushout.

To see that the free endpoint is still a Dwyer inclusion, separate both endpoints before making the attachment. The pushout of $X\amalg X\hookrightarrow X\times I$ along $f\amalg1_X$ is
\[
 j:Y\amalg X\longrightarrow Z(f).
\]
The desired inclusion is $j$ composed with the $X$-summand inclusion. Both are Dwyer maps. The printed sentence has $Y\to Y\amalg X$ in this position; its preceding formula and its source $X$ specify the intended summand. The other endpoint inclusion $X\to X\times I$ is a weak equivalence, and its pushout $s_f$ is one by the forward gluing property. Two-out-of-three applied to $p_fs_f=1_Y$ proves $p_f\in W$. No fibration property of $p_f$ is required for this finite construction.

Now take $X=\{a,b\}$ discrete, $Y=\{y\}$, and let $f$ be the unique functor. Choose $I=\Phi_2$. The category $Z(f)$ has the nine objects
\[
 y,\quad a_1,a_2,a_3,a_4,\quad b_1,b_2,b_3,b_4.
\]
Its nonidentity arrows are exactly, for $x=a,b$,
\[
 y\longrightarrow x_1\longleftarrow x_2
       \longrightarrow x_3\longleftarrow x_4.
\]
There are no composable pairs of nonidentity arrows. Thus its nerve has precisely these vertices and eight nondegenerate edges, and no higher nondegenerate simplices. Its realization is a tree with two arms. The contraction to $y$ verifies directly that $p_f$ is a weak homotopy equivalence. The free endpoints $a_4,b_4$ form a Dwyer subcategory: their generated cosieve consists of the two separate arrows $a_4\to a_3$ and $b_4\to b_3$.

We may consequently form the homotopy cofibre by the strict pushout
\[
 C(f)=Z(f)\amalg_{\{a_4,b_4\}}\{t\}.
\]
It has eight objects and eight nonidentity arrows. Its nerve is the graph obtained from the tree by identifying its two free endpoints. This graph is an eight-edge circle. In particular,
\[
 H_1(NC(f);\mathbb Z)=\mathbb Z,
 \qquad \pi_1(NC(f),t)\cong\mathbb Z.
\]
The direct pushout $\{y\}\amalg_{\{a,b\}}\{t\}$ in $\mathrm{Cat}$ is instead the terminal category: both maps identify the two objects at once. The cylinder has retained the two paths by which the gluing takes place. This is the finite categorical model of the familiar homotopy pushout $*\amalg^h_{S^0}*\simeq S^1$, and it is the reason the question on p.~232 requires a proof.
\end{dcomment}


\dgrblock{dgrb573}{What the appeal to Thomason supplies}
\begin{dcomment}{later}{p230:donc aussi la rétraction}{What the appeal to Thomason supplies}
Grothendieck's appeal to Thomason supplies a stronger factorization in the ordinary spatial case. The cellular construction in the proof of the model-category theorem factors a functor as a Dwyer inclusion followed by a trivial fibration for the Thomason model structure \citep[proof of CM5 and Theorem~4.9]{Thomason1980}. The cylinder just constructed gives a finite and transparent substitute for the weaker factorization needed here.

A later correction matters when using that historical reference. \citet{Cisinski1999Dwyer} exhibited a retract of a Dwyer map which is a Thomason cofibration but is not itself Dwyer, and repaired the properness argument. Thomason's assertion that all cofibrations are Dwyer therefore cannot be used to justify the first paragraph. The particular cellular factorization survives, as do the pushout and composition properties used in this passage. Grothendieck's explicit cylinder is also independent of closure under retracts. The distinction protects exactly the argument being made, without changing its conclusion.
\end{dcomment}


D'autre part, la proposition duale de
\[
\underline{\mathrm{Fib}}_W S \cap W_S = W^{\mathrm{u}}_{S\,\mathrm{univ}}
\]
est valable pour les $W$-cofibrations :

\msp{page 231}
\textbf{Proposition 16} (\footnote[125]{N'a rien à voir avec Cat - une chose générale quand il y a assez de $W$-cofibrations.}). \textit{Sur la sous-catégorie pleine $\underline{\mathrm{Cofib}}_W(S)$ de $S\backslash\mathrm{Cat}$, le localiseur ${}_SW$ est égal au $({}_SW)^{\mathrm{univ}}$ ($W$-équivalences couniverselles), i.e. on a ceci : Soient}
\[
S \longrightarrow X, \qquad S \longrightarrow Y
\]
\textit{des $W$-cofibrations, et soit $f : X \longrightarrow Y$ un morphisme de $S\backslash\mathrm{Cat}$ tel que $f \in W$. Alors pour tout cochangement de base $S \longrightarrow S'$, la flèche correspondante}
\[
f' : X' = X \amalg_S S' \longrightarrow Y' = Y \amalg_S S'
\]
\textit{est dans $W$.}

\dgrblock{dgrb574}{The mapping cylinder, with the two endpoints distinguished}
\begin{dcomment}{exposition}{p230}{The mapping cylinder, with the two endpoints distinguished}
The construction now turns the endpoint geometry into a factorization of an arbitrary functor. Let $a_0,a_1:X\to X\times I$ denote the two endpoint inclusions and form
\[
 Z(f)=Y\amalg_{X,a_0}(X\times I).
\]
Write $k:Y\to Z(f)$ and $f':X\times I\to Z(f)$ for the coprojections. The maps $f\,\mathrm{pr}_X$ and $\id_Y$ agree on the identified copy of $X$, hence induce a unique $p_f:Z(f)\to Y$. Put $i_f=f'a_1$. Then $p_fi_f=f$ and $p_fk=\id_Y$. The displayed $\alpha$ in the source is this same $p_f$; the projection used to define it is followed by $f$ on the cylinder summand.

The endpoint $a_0$ is a Dwyer map and a weak equivalence. Its pushout $k$ is consequently a weak equivalence by the forward gluing property. Two-out-of-three applied to $p_fk=\id_Y$ gives $p_f\in W$. To prove the other assertion, use both endpoints at once. The map
\[
 j:Y\amalg X\longrightarrow Z(f)
\]
is the pushout of $a_0\amalg a_1:X\amalg X\to X\times I$ along $f\amalg\id_X$, hence is Dwyer. The inclusion of the second summand $X\hookrightarrow Y\amalg X$ is Dwyer, and its composite with $j$ is $i_f$. Closure under composition proves the required assertion. The edition says $Y\to Y\amalg X$ in this sentence; the formula immediately preceding it requires the second summand $X$.

The construction is functorial in $f$. A commutative square $bf=f'a$ gives maps $b$ on the $Y$ summand and $a\times\id_I$ on the cylinder summand, compatible on the attaching copy of $X$. The universal property gives the induced map of mapping cylinders, and uniqueness proves compatibility with composition and identities. This supplies an actual replacement functor on the arrow category, not merely objectwise existence of a factorization. Nothing in this argument says that the displayed $p_f$ is a Thomason fibration; the source only claims $p_f\in W$ for this explicit factorization. The stronger model-category factorization mentioned earlier is a different choice.

A small computation shows what the construction buys. Let $X$ be the discrete two-object category and let $f:X\to e$ be the unique functor. With $I=\Phi_2$, $Z(f)$ is made of two four-edge intervals sharing their attached endpoint $y$. Form the further pushout $e\amalg_X Z(f)$, identifying their two free endpoints to a second point $z$. The resulting poset has eight vertices: $y,z$, and three interior vertices on each arm. Every nonidentity arrow runs from an even vertex to an adjacent odd vertex, so there are no nondegenerate simplices of its nerve in dimensions above one. Its realization is the graph of two four-edge paths from $y$ to $z$, hence a circle. Its cellular differential has rank seven, giving $H_1=\mathbf Z$ and $H_q=0$ for $q>1$; a spanning tree leaves one generator of $\pi_1$. This is the homotopy cofiber of the collapse $S^0\to *$, obtained by an ordinary categorical pushout only after the mapping-cylinder replacement.
\end{dcomment}


\textsc{Démonstration}. On factorise $S \longrightarrow S'$ en
\[
S \xrightarrow{i} S'_0 \xrightarrow{j} S',
\]
avec $i$ immersion ouverte de Dwyer (donc $W$-cofibration), et $j \in W$. On a donc le diagramme à carrés cocartésiens
\[
\begin{tikzcd}
S \arrow[r] \arrow[d, hook, "i"'] & X \arrow[r, "f"] \arrow[d, hook] & Y \arrow[d, hook] \\
S'_0 \arrow[r] \arrow[d, "\in W"'] & X'_0 \arrow[r, "f'_0", "W"'] \arrow[d, "W"] & Y'_0 \arrow[d, "W"] \\
S' \arrow[r] & X' \arrow[r, "f'"] & Y'.
\end{tikzcd}
\]
Comme $i$ est une $W$-cofibration et $f \in W$, on a $f'_0 \in W$. D'autre part, comme $S \longrightarrow X$, donc $S'_0 \longrightarrow X'_0$ est
\printedpage{148}
% head: A. Grothendieck : Dérivateurs, Ch. XII, Édition des Universités de Montpellier II et Paris VII
% foot: Version 30 mars 2003
\msp{page 232}
une $W$-cofibration, $X'_0 \longrightarrow X'$ est dans $W$, et pour la même raison, $Y'_0 \longrightarrow Y$ l'est aussi. Donc $f'$ l'est aussi. % [sic] Y'_0 \longrightarrow Y -> Y'_0 \longrightarrow Y', the diagram has Y'_0 \to Y'

\textbf{NB} Ici, c'est un argument formel, valable (sous forme duale) dans tout $(\mathcal{M}, W)$ où on a un énoncé de factorisation pour une flèche quelconque $f : X \longrightarrow Y$ en $X \xrightarrow{i_f} \bar{X} \xrightarrow{\alpha} Y$, avec $i_f \in W$ et $\alpha$ une $W$-fibration. Alors, si $X$, $Y$ sont $W$-fibrés sur $S$, toute $S$-flèche de $X$ dans $Y$ qui est dans $W$, est dans $W^{\mathrm{u}}_S$. Et dualement, si toute flèche $X \longrightarrow Y$ se factorise en $X \longrightarrow \tilde{Y} \longrightarrow Y$, avec $X \longrightarrow \tilde{Y}$ une $W$-cofibration et $\tilde{Y} \longrightarrow Y$ dans $W$.

\dgrblock{dgrb575}{Weak equivalences between objects under a fixed base}
\begin{dcomment}{exposition}{p232:Ici, c'est un argument formel, valable}{Weak equivalences between objects under a fixed base}
Proposition~16 provides the independence of replacement that a relative homotopy construction needs. A map under $S$ can be a weak equivalence while its image after a badly behaved attachment is not. The proposition locates a class of objects under $S$ for which every attachment behaves correctly: both structure maps from $S$ must be $W$-cofibrations. It does not require the map between those objects to be a cofibration.

The argument can be expressed in any category with the pushouts in question. Let $a:S\to X$ and $b:S\to Y$ be $W$-cofibrations, and let $f:X\to Y$ satisfy $fa=b$ and $f\in W$. To change $S$ into an arbitrary $T$, factor the change of base as
\[
 S\xrightarrow{k}T_0\xrightarrow{q}T,
 \qquad k\in\operatorname{Cof}_W,\quad q\in W.
\]
Set $X_0=X\amalg_ST_0$, $Y_0=Y\amalg_ST_0$, and set $X_T=X\amalg_ST$, $Y_T=Y\amalg_ST$. The map $f_0:X_0\to Y_0$ is in $W$ because $k$ is a $W$-cofibration and it pushes out the map $f$ in $S\backslash\mathcal M$. On the other hand, $T_0\to X_0$ and $T_0\to Y_0$ are pushouts of $a$ and $b$, so they are $W$-cofibrations. Pushing out $q$ along them gives weak equivalences
\[
 u:X_0\longrightarrow X_T,
 \qquad v:Y_0\longrightarrow Y_T.
\]
The equality $f_Tu=vf_0$ now proves $f_T\in W$ by two-out-of-three. These are three separate applications of the cofibration condition, each with its own fixed base. They account for every marked $W$ in the displayed proof.

There is a useful immediate consequence. If $a:S\to X$ is both a $W$-cofibration and a weak equivalence, apply the proposition to the two objects $1_S:S\to S$ and $a:S\to X$ of $S\backslash\mathcal M$. Every pushout of $a$ belongs to $W$. Thus, once these factorizations exist, a trivial $W$-cofibration is exactly a couniversal weak equivalence. This is the abstract version of the special Dwyer fact used to construct the cylinder. The order of proof is sound: the special Dwyer statement constructs enough cofibrations first, and Proposition~16 then gives the general statement.

The dual formulation on p.~232 is obtained by taking opposite categories throughout. A weak equivalence between two $W$-fibrations over $S$ survives every pullback along $T\to S$. The symmetry is a symmetry of the relative arguments. It does not assert that every object of a slice or coslice is already suitable for calculating a derived change of base.
\end{dcomment}


Il me vient ici une (petite) inquiétude : ai-je au moins prouvé que si
\[
S \longrightarrow X
\]
est une $W$-cofibration, alors pour toute flèche $S' \longrightarrow S$ dans Cat, le carré cocartésien dans Cat % [sic] S' \longrightarrow S -> S \longrightarrow S', the square below has S \to S'
\[
\begin{tikzcd}
S \arrow[r] \arrow[d] & X \arrow[d] \\
S' \arrow[r] & X'
\end{tikzcd}
\]
est aussi $W$-cocartésien (\footnote[126]{dans le sens ancien, spécial à Cat (déf. 3, p. 131). Pour le comparer avec la théorie générale, cf. prop. 20, p. 262.})? Autrement,
\msp{page 233}
à quoi ce carré-là serait-il bon?! Pour le prouver, factorisons encore $S \longrightarrow S'$ en
\[
\underset{\substack{\text{immersion}\\ \text{ouverte de}\\ \text{Dwyer}}}{S \hookrightarrow} S'_0 \underset{\in W}{\longrightarrow} S',
\]
d'où le diagramme de 2 carrés cocartésiens
\[
\begin{tikzcd}
S \arrow[r, "\mathrm{Cof}_W"] \arrow[d, hook, "D_W"'] & X \arrow[d, hook, "D_W"] \\
S'_0 \arrow[r, "\mathrm{Cof}_W"] \arrow[d, "\in W"'] & X'_0 \arrow[d, "\in W"] \\
S' \arrow[r, "\mathrm{Cof}_W"] & X',
\end{tikzcd}
\]
où $X'_0 \longrightarrow X'$ est $\in W$ d'après l'hypothèse sur $X$ sous $S$. Il faut donc vérifier :
\printedpage{149}
% head: A. Grothendieck : Dérivateurs, Ch. XII, Édition des Universités de Montpellier II et Paris VII
% foot: Version 30 mars 2003
\begin{enumerate}
\item[a)] Un carré cocartésien construit sur une flèche qui est \texttt{[une]} immersion de Dwyer, est $W$-cocartésien.
\item[b)] Un carré commutatif dans lequel deux côtés opposés sont dans $W$, est $W$-cocartésien.
\item[c)] Le composé de deux carrés $W$-cocartésiens est $W$-cocartésien.
\end{enumerate}

\dgrblock{dgrb576}{Why Proposition 16 is a formal gluing lemma}
\begin{dcomment}{exposition}{p232}{Why Proposition 16 is a formal gluing lemma}
Grothendieck needs weak equivalences between objects under $S$ to survive an arbitrary change of the attaching object. The hypothesis that both objects are $W$-cofibrant under $S$ is exactly what controls the two comparisons introduced by a factorization of that change. It cannot be replaced by a statement about the underlying objects alone.

Assume an ambient category has pushouts, a weakly saturated class $W$, and a class of $W$-cofibrations stable under pushout, with the property that pushout along such a map preserves $W$-equivalences. Suppose every map factors as a $W$-cofibration followed by a member of $W$. Let $S\to X$ and $S\to Y$ be $W$-cofibrations, and let $f:X\to Y$ be a $W$-equivalence under $S$. Factor the chosen $S\to S'$ as
\[
 S\xrightarrow{i}S_0'\xrightarrow{j}S',\qquad i\in\mathrm{Cof}_W,\quad j\in W.
\]
Put $X_0'=X\amalg_SS_0'$ and $Y_0'=Y\amalg_SS_0'$. The map $X\to X_0'$ is a pushout of $i$, hence a $W$-cofibration. Pushing out $f$ along this map gives $f_0':X_0'\to Y_0'$ in $W$; its target is $Y\amalg_X X_0'\cong Y\amalg_SS_0'$. The maps $S_0'\to X_0'$ and $S_0'\to Y_0'$ are pushouts of the two original cofibrations. Pushing out $j$ along them gives weak equivalences $v_X:X_0'\to X'$ and $v_Y:Y_0'\to Y'$. In the square
\[
 f'v_X=v_Yf_0',
\]
the right-hand composite belongs to $W$, and $v_X$ belongs to $W$. Two-out-of-three proves $f'\in W$.

This is the whole argument. It uses neither a lifting axiom nor the conservative converses of footnote~124. It also identifies the dual hypotheses precisely: pullbacks, fibrations whose pullback functors preserve weak equivalences, stability under pullback, and a factorization as a weak equivalence followed by a fibration. Under those hypotheses a weak equivalence between fibrations over the same base is universally a weak equivalence over that base. Thus the parallel with the preceding fibre theory is mathematical, not just terminological.

The small slips in the final prose are harmless once the diagram is read: the comparison is $Y_0'\to Y'$, not $Y_0'\to Y$, and the subsequent cobase-change arrow is $S\to S'$, not $S'\to S$. Both are printed slips and remain in the source text. They are not gaps in Proposition~16.
\end{dcomment}


\dgrblock{dgrb577}{Comparing the two meanings of homotopy pushout}
\begin{dcomment}{exposition}{p233:Il faut donc vérifier :}{Comparing the two meanings of homotopy pushout}
Grothendieck now asks whether the replacement method agrees with the earlier construction special to small categories. The earlier definition starts with a span $X_1\leftarrow X_0\to X_2$, integrates that diagram into a category $\mathcal X$, and calls a cocone to $Y$ homotopy cocartesian when its comparison $\mathcal X\to Y$ belongs to $W$. The new method starts with an ordinary pushout and asks whether a $W$-cofibration makes that pushout homotopically correct. The two descriptions have different starting data; their agreement is the content of the corollary on p.~239, and later of Proposition~20.

The proof has three stages, and keeping their dependence visible avoids a misleading impression of circularity. A pushout along a Dwyer inclusion is first reduced to a pushout whose two incoming maps are Dwyer inclusions, followed by a square with two opposite sides in $W$. In the first square, both categories embed as open subcategories of their union, so the earlier open-cover calculation applies. In the second, the definition of $W$-cofibration supplies the second weak equivalence. The only remaining operation is composition of homotopy-cocartesian squares.

That composition rule is then proved independently, by replacing the comparison maps with inclusions between categories of elements. Its proof uses restrictions of a single diagram over a six-object poset and a union of closed subcategories. It therefore does not presuppose that all Dwyer pushouts have already been identified as homotopy pushouts. Finally, an arbitrary $W$-cofibration is handled by factoring the other incoming map through a Dwyer inclusion and a weak equivalence. The two resulting squares fall under the established cases.

The assertion about two opposite weak equivalences has an especially short explanation. In the source's notation, if $i_2:X_0\to X_2$ is in $W$, the other leg $j_1:X_1\to\mathcal X$ of the integrated span is in $W$. The cocone satisfies $\varphi j_1=f_1$. If $f_1$ also belongs to $W$, two-out-of-three gives $\varphi\in W$. This argument concerns the actual coherent square, including its commutativity. It will reappear in the general formalism as the rule that a homotopy pushout of an equivalence is an equivalence.
\end{dcomment}


Admettons provisoirement b) et c), et prouvons a), pour un carré cocartésien
\[
\begin{tikzcd}
U \arrow[r, hook, "i"] \arrow[d, "p"'] & X \arrow[d, "q"] \\
U' \arrow[r, hook, "i'"] & X',
\end{tikzcd}
\]
où $i$ est immersion ouverte de Dwyer. On procède encore comme tantôt, en factorisant $p$ en $U \overset{j}{\hookrightarrow} U'_0 \xrightarrow{p'} U'$, où $j$ est \texttt{[une]} immersion ouverte de Dwyer, $p' \in W$. On a donc un diagramme de deux carrés cocartésiens accolés
\msp{page 234}
\[
\begin{tikzcd}
U \arrow[r, hook, "i \in D_W"] \arrow[d, hook, "D_W"'] & X \arrow[d, hook] \\
U'_0 \arrow[r, hook, "D_W"] \arrow[d, "p' \in W"'] & X'_0 \arrow[d, "q' \in W"] \\
U' \arrow[r, hook] & X',
\end{tikzcd}
\]
où le premier carré est formé d'immersions ouverts. On sait alors qu'il est $W$-cocartésien. Dans le deuxième, on a $p' \in W$, donc $q' \in W$ puisque $i$ est une immersion de Dwyer, donc une $W$-cofibration. Donc il est également $W$-cocartésien. On conclut alors grâce à c) ci-dessus. % [sic] immersions ouverts -> immersions ouvertes, agreement

\noindent Prouvons b), pour
\begin{tikzcd}[baseline=(\tikzcdmatrixname-1-1.base)]
X_0 \arrow[r, "i_1"] \arrow[d, "i_2 \in W"'] & X_1, \arrow[d, "f_1 \in W"] \\
X_2 \arrow[r, "f_2"] & Y
\end{tikzcd}
en l'insérant dans
\begin{tikzcd}[baseline=(\tikzcdmatrixname-1-1.base)]
X_0 \arrow[r, "i_1"] \arrow[d, "i_2 \in W"'] & X_1, \arrow[d, "j_1"] \arrow[ddr, "f_1"] & \\
X_2 \arrow[r, "j_2"] \arrow[drr, "f_2"'] & \mathcal{X} \arrow[dr, "\varphi"] & \\
& & Y
\end{tikzcd}
où $\mathcal{X} = \displaystyle\int X_0$
\begin{tikzcd}[baseline=(\tikzcdmatrixname-2-1.base), row sep=small]
& X_1 \\
{} \arrow[ur] \arrow[dr] & \\
& X_2
\end{tikzcd}
.

Comme $i_2 \in W$, on sait qu'il en résulte $j_1 \in W$, et comme $f_1 \in W$, il en résulte que $\varphi \in W$, q.e.d.


\dgrblock{dgrb578}{Completing the pasting step}
\begin{dcomment}{exposition}{p234}{Completing the pasting step}
At the end of manuscript page~234, the proof that a pushout along a $W$-cofibration is $W$-cocartesian still invokes the composition of homotopy pushout squares. This is the deferred item~c), not a defect in item~b), whose proof has just been given. The composition statement is a standard consequence of coherent homotopy colimits; in derivator language it is the composition part of \citet[Proposition 3.13, published version]{Groth2013}. We spell out the argument to make clear which compatibility is required.

Consider a coherent commutative diagram
\[
\begin{tikzcd}
 A\arrow[r]\arrow[d]&B\arrow[d]\\
 C\arrow[r]\arrow[d]&D\arrow[d]\\
 E\arrow[r]&F .
\end{tikzcd}
\]
Suppose each small square is a homotopy pushout. Its two universal properties give
\[
 D\simeq B\amalg_A^{h}C,\qquad
 F\simeq D\amalg_C^{h}E.
\]
For any test object $T$, write $M_U=\operatorname{Map}(U,T)$. These mapping spaces are derived mapping spaces, so that homotopy colimits become homotopy limits. Then
\[
\begin{aligned}
 \operatorname{Map}(F,T)
 &\simeq (M_B\times^h_{M_A}M_C)\times^h_{M_C}M_E\\
 &\simeq M_B\times^h_{M_A}M_E.
\end{aligned}
\]
The second equivalence is the associative law for homotopy pullbacks: both sides are limits of the same diagram after eliminating the repeated $M_C$. The natural equivalences identify the maps to $M_B$ and $M_E$ and their common homotopy over $M_A$. The mapping-space universal property therefore identifies $F$ with $B\amalg_A^h E$, including its cocone. This proves that the outside rectangle is a homotopy pushout.

For Grothendieck's $W$-cocartesian comparison, we give the categorical calculation explicitly. Suppose $W$ is a fundamental localizer. First, integrating an objectwise $W$-equivalence $\theta:H\to H'$ of covariant diagrams on $J$ gives a $W$-equivalence. Indeed, the over-category $(\int_J H)/j$ of the projection to $J$ has an evaluation functor
\[
 (i,x,\alpha:i\to j)\longmapsto H(\alpha)x.
\]
It is left adjoint to $y\mapsto(j,y,\id_j)$, with identity counit. Its nerve is therefore a homotopy equivalence. The analogous evaluation for $H'$ commutes with $\theta$; two-out-of-three shows that the induced maps of over-categories belong to $W$. The relative Theorem A gives $\int\theta\in W$. Here and below an ordinary weak equivalence belongs to $W$ by \citet[Theorem 2.2.11]{Cisinski2004Minimal}.

Write $H_1=\int(B\leftarrow A\to C)$ and $H_2=\int(D\leftarrow C\to E)$. The two small squares being $W$-cocartesian means that their canonical maps $H_1\to D$ and $H_2\to F$ lie in $W$. Objectwise invariance now shows that
\[
 H_{12}=\int(H_1\leftarrow C\to E)\longrightarrow H_2\longrightarrow F
\]
belongs to $W$. To compare this iterated integral with the outside span, introduce the five-object poset $K$ whose only nonidentity arrows are
\[
 a\to b,\quad a\to c_1,\quad c_0\to c_1,\quad c_0\to e.
\]
Assign $A,B,C,C,E$ to $a,b,c_1,c_0,e$, with the identity on the arrow between the two copies of $C$. Flattening the two Grothendieck constructions identifies its integral with $H_{12}$: the outer copy of $C$ is $c_0$, and the copy inside $H_1$ is $c_1$.

Let $K'$ have objects $a,b,c,e$ and arrows generated by $a\to b$ and $a\to c\to e$. The functor $r:K\to K'$ identifies $c_0,c_1$ with $c$. Let $G:K'\to\Cat$ be the original diagram $A\to B$, $A\to C\to E$. Thus the flattened diagram is $r^*G$. Let $L\subset K'$ be the full subcategory on $a,b,e$; its integral is the outside-span integral $H_0=\int(B\leftarrow A\to E)$.

Both $r$ and $L\hookrightarrow K'$ are homotopy final. This can be checked without an abstract pasting theorem. For $r$, the comma category $a/r$ is the tree-shaped poset $K$; $c/r$ is the cospan $c_1\leftarrow c_0\to e$; and $b/r,e/r$ are singletons. All have contractible nerves. For the inclusion of $L$, the comma categories at $a,b,c,e$ are respectively a span with initial object $a$, a singleton, the singleton $e$, and a singleton. Homotopy cofinality for homotopy colimits, together with the nerve comparison for integration \citep[Theorem 1.2]{Thomason1979}, therefore gives ordinary weak equivalences
\[
 H_{12}\longrightarrow \int_{K'}G\longleftarrow H_0.
\]
All these functors respect the cocone to $F$. Since $H_{12}\to F$ lies in $W$, two-out-of-three first gives $\int_{K'}G\to F$ in $W$ and then $H_0\to F$ in $W$. This proves the required pasting statement for every fundamental localizer. It supplies the comparison and its compatibility with the cocone, rather than merely an equality of unstructured homotopy types.

Item~b) has an equally transparent interpretation. If $A\to C$ is a weak equivalence, then $B\to B\amalg_A^hC$ is a weak equivalence. If also $B\to D$ is a weak equivalence, the comparison from that homotopy pushout to $D$ is one by two-out-of-three. The final integral diagram in the source is a categorical model for exactly this proof. For a general $W$, use the objectwise invariance of integration just proved instead of appealing to ordinary spaces.

These two facts complete the reduction. Factor the cobase-change map through a Dwyer inclusion followed by a weak equivalence. The first pushout is $W$-cocartesian by the Dwyer gluing theorem, and the second has the two opposite weak equivalences supplied by the $W$-cofibration hypothesis. Item~b) handles the second square; pasting handles their composite. This proof establishes the intended statement without adding a converse detection axiom for individual edges.
\end{dcomment}


\dgrblock{dgrb579}{What survives beyond homotopy types}
\begin{dcomment}{later}{p234}{What survives beyond homotopy types}
The Dwyer condition ultimately controls more than the weak homotopy type of the nerve. \citet[Theorems 1.5 and 1.6; arXiv version 3]{HackneyOzornovaRiehlRovelli2024} prove that pushouts along Dwyer maps are preserved when ordinary categories are viewed as $(\infty,1)$-categories. Equivalently, the comparison from the pushout of nerves to the nerve of the ordinary pushout is a categorical equivalence, a stronger assertion than Thomason's weak homotopy equivalence. Their formulation also covers suitable generalized Dwyer maps.

This later theorem concerns the fidelity of the gluing construction. It supplies no converse of the kind excluded by the acyclic example above. It is compatible with Grothendieck's strategy of finding strict representatives on which homotopy operations become ordinary operations, while explaining why the coherence of the original diagram must be retained. In the pointed derivator formulation, a cofiber is correspondingly functorial on the category of coherent arrows, rather than determined by an arrow in the ordinary homotopy category alone \citep[Definition 3.18]{Groth2013}.
\end{dcomment}
\begin{otherlanguage}{english}
\dgroverview{XIIc-a}{Chapter XII \textperiodcentered\ \textup{XIIc-a} \textperiodcentered\ Overview}
\selectlanguage{english}

These pages begin by asking what the categorical loop constructions actually detect, and end by constructing their dual, a categorical mapping cylinder \citep[manuscript pp.~171--234]{Grothendieck1990XIIc}. Between these two constructions Grothendieck makes a decisive change of viewpoint. A path category is first a concrete device for computing homotopy groups. It then becomes one possible presentation of an intrinsically defined operation on localized categories. The problem is to retain enough of the family of relative homotopy types that changing the presentation does not change the operation or its compatibility with a further change of base.

The first section tests this programme against a classical comparison between cohomology and homotopy. For a small category $X$, write $BX=|NX|$. Constant coefficients detect connected components, while locally constant coefficients detect transport around the fundamental groupoid. Grothendieck's several formulations of $W_0$ and $W_1$ are therefore exact descriptions of the information that a zero-type and a one-type retain. The nonabelian degree-one condition matters: an abelian cohomology test alone can leave a perfect fundamental group undetected. Our commentary spells out the small group-theoretic argument showing why surjectivity on constant nonabelian $H^1$, together with the local-system $H^0$ test, already determines the fundamental groupoid.

The attempted extension to higher degrees is more delicate. With his notation, $a_n$ means that every homotopy fibre is $n$-connected; $b_n$ means that the map induces isomorphisms on homotopy groups through degree $n$, including components; and $c_n$ means that it induces the corresponding cohomology isomorphisms through degree $n$ for all permitted locally constant coefficients. The range of $i$ is missing from the printed statement of $c_n$, but is fixed by the proof. The corrected theorem establishes
\[
 a_n\ \Longrightarrow\ b_n\ \Longrightarrow\ c_n\
 \Longrightarrow\ a_{n-1}.
\]
The homotopy exact sequence explains the first loss of one degree. The next implication uses the cohomological spectral sequence of the homotopy fibre, with its monodromy retained. The last implication tests the first possible nonzero fibre group against itself as a coefficient module. These arguments are sound once the stated verifications concerning twisted coefficients and the degree-one fibre group are supplied.

The withdrawn equivalence $b_n\Leftrightarrow c_n$ really does fail, but the footnote's suggested degree-two failure does not occur. The inclusion $S^2\to\mathbf{CP}^{\infty}=K(\mathbf Z,2)$ is a cohomology isomorphism through degree three with every abelian coefficient group and has trivial fundamental-group data; it kills $\pi_3(S^2)$. In degree two, by contrast, after identifying the fundamental groups, the coinduced local system
\[
 \operatorname{Hom}_{\mathbf Z}(\mathbf Z[\pi],\mathbf Q/\mathbf Z)
\]
dually detects the second homology of the universal cover. The natural Hurewicz isomorphism then detects $\pi_2$. The commentary gives the complete argument and distinguishes it from a claim of historical priority. It also completes the universal-cover spectral-sequence approach that Grothendieck leaves aside: a finite-range truncation argument for derived invariants supplies the needed comparison without assuming an isomorphism in the next degree.

There is a second correction at the end of the section. The class $a_n$ is not a localizer, because it fails two-out-of-three. For a nonzero abelian group $A$, the composite
\[
 *\longrightarrow K(A,n+1)\longrightarrow *
\]
is the identity, and the second map has $n$-connected fibre, whereas the first does not. The class $b_n$, the equivalences detected by Postnikov $n$-truncation, supplies the intended weak equivalences. There is nevertheless a precise sense in which $a_n$ suffices: its saturation generates the same localization. The two conditions already agree for maps to a point, so this correction leaves Theorem~2 intact. That theorem recognizes weak contractibility by connectedness of all iterated loop categories. At the degree-two endpoint its reverse induction needs the abelianness of $\pi_1\Omega X$; we supply the Eckmann--Hilton verification. The modern identification of the minimal fundamental localizer with $W_\infty$ and the later loop-based characterization are established by \citet[Theorems~2.2.11 and 2.3.6]{Cisinski2004Minimal}.

The question on p.~191 is consequently no longer whether the path objects work, but in what sense their answers are canonical. Fix a base $S$. A replacement consists of an endofunctor $C$ of $\Cat/S$, a natural map $i:1\to C$ whose components belong to $W$, and the requirement that $C(X)$ be $W$-fibred over $S$. Two replacements compare through their composite. The compatibility of this comparison with the units is exactly naturality, and two-out-of-three shows that the comparison maps are weak equivalences. These elementary observations already contain the solution to the coherence problem posed a few pages later.

Indeed, fixing one replacement $C_0$, composition gives natural transformations from the identity functor on the category of replacements to $C\mapsto C C_0$, and from the constant functor at $C_0$ to the same functor. Passing to nerves contracts the identity to a constant map. Thus this category is contractible whenever it is nonempty, and its fundamental groupoid is equivalent to the terminal groupoid. The manuscript's assertion of filteredness is not justified by its common-target construction: equalizing parallel arrows is a separate requirement. Contractibility provides exactly the comparison that is needed without that assertion. The localization argument on pp.~204--207 gives a second formulation: the inclusion of suitable fibred representatives induces an equivalence of localized categories, and the unit of replacement supplies the specified comparison to its quasi-inverse.

This concern with operations independent of their computation is continuous with the programme of \citet[\S69, printed pp.~1--3]{Grothendieck1983}. There the derived categories of diagrams, together with their change-of-index functors and adjoints, replace the information lost on keeping only one derived category. The letter to Thomason makes the distinction between the relative category $(\mathcal M,W)$ and supplementary choices of fibrations and cofibrations explicit \citep[points~1--2]{Grothendieck1991Letter}. The present pages examine that distinction in a concrete and demanding case. An ordinary functor $\Pi S\to\mathrm{Hot}$ records the homotopy types of fibres and their first transport, but it does not by itself retain all higher coherence of a locally constant family. The later theories of derived diagrams and homotopy Kan extensions preserve the additional structure rather than identifying these two levels of information.

The fourth section asks how much of this picture can be stated for a category with pullbacks and weak equivalences. Its intrinsic $W$-fibrations are maps whose pullback functors preserve weak equivalences. These are the maps later called sharp in a model-category setting; they should not be identified indiscriminately with an independently chosen class of Quillen fibrations. Grothendieck compares the localization of the entire slice at underlying weak equivalences with the localization of the subcategory of such fibrations at universally preserved weak equivalences. The fibre axiom says that these localizations agree. The adjoining axiom says that the two weak-equivalence classes already agree on fibrant representatives, before localization.

Some care with the hypotheses is essential here. The stated two-out-of-three assumption does not imply closure under retracts. Our finite-set example shows that the corresponding claimed retract closure of intrinsic fibrations also fails at precisely that level of generality; imposing retract closure on $W$ repairs it. On the other hand, a functorial fibrant replacement with a weak-equivalence unit gives the desired localization comparison for every full intermediate subcategory containing its image. This proves the intermediate-subcategory assertion in the situations for which the manuscript actually has a replacement construction. The bare logical implication from the two fibre axioms to that assertion is not settled here; this restricted conclusion is recorded as such, not presented as a newly identified open problem. The subsequent manuscript discussion at pp.~240--241 adds a factorization hypothesis and its own caution about fibrancy over a terminal object.

Non-locally-constant relative types require more structure than the preceding slice localization. Grothendieck introduces a pullback-stable class $\mathcal L$ of allowed tests and defines $W_S^{\mathcal L}$ by preservation under these tests. In $\Cat$ the two comma-category variances give $W_S^{\mathrm g}$ and $W_S^{\mathrm d}$. Smooth and proper functors provide complementary ways of changing the base and choosing representatives. Our commentary proves the formal universal-saturation construction and identifies the saturation asked about on p.~222 with $W$-smooth functors. The key is to test a $W$-aspherical map as an arrow to the terminal object of its slice and use the universal pullback characterization in \citet[Theorem~1.2(c)]{Maltsiniotis2005Asphericite}. A calculation with a degree-two self-map of a circle then shows why a genuinely varying diagram and a locally constant family are different objects. The abstract existence of both adjoints to restriction requires further hypotheses; in the combinatorial model-category setting the resulting strong derivator is established by \citet[Proposition~1.30]{Groth2013}.

Finally the interval $\Phi_n$, with $n\geq2$, gives a Dwyer inclusion of its two endpoints. Multiplying by $X$ produces a cylinder, and attaching its first endpoint to $Y$ along $f:X\to Y$ gives a factorization of $f$ into a Dwyer inclusion and a weak equivalence. The commentary identifies the required adjoints explicitly, explains why the shorter interval fails for the two-endpoint inclusion, and verifies the induced pushout and comparison maps. Proposition~16 then follows by factoring a change of cobase and applying preservation of weak equivalences along cofibrations. The final pasting step is completed with its hypotheses stated; its later derivator formulation is \citet[Proposition~3.13]{Groth2013}.

The two biconditionals in footnote~124 need a different reading. They use the manuscript's stronger conservative pushout axioms, whose converse implications are not properties of ordinary weak homotopy equivalences. A finite acyclic, noncontractible complex and the two cones in its suspension give a counterexample by face posets, even with Dwyer inclusions. This does not invalidate the mapping-cylinder construction or Proposition~16: their proofs use forward preservation, which is the part available for $W_\infty$. Keeping that distinction visible lets the successful construction retain its force while preventing an extra axiom from silently becoming a claim about classical homotopy theory.
\end{otherlanguage}
\dgrchunkend

% DER_XIIc-b: Grothendieck's text with the commentary set in, then the overview, inside a group that applies the chunk's own definitions
\dgrpart{1}{XIIc-b}{Section 6, part 2 of 2, manuscript pages 235--275}{XIIc-b\enspace\textperiodcentered\enspace Chapter XII, section 6, part 2 of 2\enspace\textperiodcentered\enspace manuscript pages 235--275}{Chapter XII \textperiodcentered\ \textup{XIIc-b}}
\dgrchunkbegin
% DER_XIIc-b: the definitions of the chunk's preamble, applied inside its group
\tikzcdset{every diagram/.append style={nodes in empty cells}}
\fancyhf{}
\fancyhead[L]{\footnotesize\itshape Les D\'erivateurs -- XIIc-b}
\fancyhead[R]{\footnotesize Manuscript 228--275}
\fancyfoot[C]{\footnotesize\thepage}
\providecommand\headrulewidth{}\renewcommand{\headrulewidth}{0.3pt}
\let\printedpage\relax\newcommand{\printedpage}[1]{\marginpar{\footnotesize p.~#1}}
\let\msp\relax\newcommand{\msp}[1]{\marginpar{\footnotesize [#1]}}
\let\cX\relax\newcommand{\cX}{\mathscr{X}}
\let\Ob\relax\newcommand{\Ob}{\operatorname{Ob}}
\let\cercle\relax\newcommand{\cercle}[1]{\tikz[baseline=(c.base)]\node[draw,circle,inner sep=1pt](c){#1};}
\let\Fibu\relax\newcommand{\Fibu}{\underline{\mathrm{Fib}}}
\let\Homu\relax\newcommand{\Homu}{\underline{\mathrm{Hom}}}
\let\Ho\relax\newcommand{\Ho}{\mathrm{Ho}}
\let\ho\relax\newcommand{\ho}{\mathrm{ho}}
\let\lc\relax\newcommand{\lc}{\mathrm{lc}}
\let\uFib\relax\newcommand{\uFib}{\underline{\mathrm{Fib}}}
\let\Fib\relax\newcommand{\Fib}{\mathrm{Fib}}
\let\Cof\relax\newcommand{\Cof}{\mathrm{Cof}}
\let\Asph\relax\newcommand{\Asph}{\mathrm{Asph}}

% ---- from Der-ChXIIc-30-03-03_p036-040.tex
\printedpage{150}
% head: A. Grothendieck : Dérivateurs, Ch. XII, Édition des Universités de Montpellier II et Paris VII
% foot: Version 30 mars 2003

Reste à prouver c). Je vais le déduire d'une

\textbf{Proposition 17.} \textit{Soit un carré commutatif}
\[
Q \quad = \quad
\begin{tikzcd}[ampersand replacement=\&]
X_0 \arrow[r, "i_1"] \arrow[d, "i_2"'] \& X_1 \arrow[d, "f_1"] \\
X_2 \arrow[r, "f_2"] \& Y
\end{tikzcd}
\]
\textit{dans} Cat\textit{, et $\cX_Q$ la catégorie fibrée sur}
\[
\tilde{\Phi} \quad (\simeq \mathbf{\Delta}^1 \times \mathbf{\Delta}^1) \quad = \quad
\begin{tikzcd}[ampersand replacement=\&]
0 \& 1 \arrow[l] \\
2 \arrow[u] \& 3 \arrow[l] \arrow[u]
\end{tikzcd}
\]
\textit{qui l'intègre. Soit $\cX$ l'image inverse du fermé}
\[
\begin{tikzcd}[ampersand replacement=\&]
0 \& 1 \arrow[l] \\
2 \arrow[u] \&
\end{tikzcd}
\]
\msp{page 235}
\textit{dans $\cX_Q$, qui intègre le diagramme}
\[
\begin{tikzcd}[ampersand replacement=\&, row sep=small]
 \& X_1 \\
X_0 \arrow[ru] \arrow[rd] \& \\
 \& X_2
\end{tikzcd}
\]
\textit{. Pour que $Q$ soit $W$-cocartésien, il faut et il suffit que l'inclusion}
\[
i : \cX \longrightarrow \cX_Q
\]
\textit{soit dans $W$.}

Comme $\cX_Q$ est lisse sur $\tilde{\Phi}$, et 3 est objet initial de $\tilde{\Phi}$, il s'ensuit que l'inclusion de la fibre de $\cX_Q$ en le sommet 3 de $\tilde{\Phi}$ est dans $W$,
\[
Y \xrightarrow{i_Y} \cX_Q \qquad\qquad \in W.
\]

Considérons alors le diagramme
\[
\begin{tikzcd}[ampersand replacement=\&, column sep=small]
\cX \arrow[rr, "\varphi_Q"] \arrow[rd, "i"', ""{name=direct,pos=.55,above right}] \& \& Y \arrow[ld, "i_Y\in W", ""{name=via,pos=.30,below left}] \\
 \& \cX_Q. \& \arrow[from=via,to=direct,dotted,shorten <=1pt,shorten >=1pt,"\alpha"']
\end{tikzcd}
\]

\printedpage{151}
% head: A. Grothendieck : Dérivateurs, Ch. XII, Édition des Universités de Montpellier II et Paris VII
% foot: Version 30 mars 2003

Ce diagramme n'est pas commutatif, mais il y a une flèche de commutativité
\[
\alpha : i_Y\varphi_Q \longrightarrow i.
\]

\dgrblock{dgrb580}{The integral criterion with its natural transformation}
\begin{dcomment}{exposition}{p235:Ce diagramme n'est pas commutatif}{The integral criterion with its natural transformation}
Proposition~17 turns the comparison with a homotopy pushout into a test on one inclusion. This is what makes the pasting proof manageable: several comparison maps can then be placed in the same category and related by literal inclusions. The natural transformation in the triangle is essential to that change of viewpoint.

Here is an explicit model, fixing the variance. Regard the square as a contravariant diagram $F:\widetilde\Phi^{\mathrm{op}}\to\mathrm{Cat}$, with $F(0)=X_0$, $F(1)=X_1$, $F(2)=X_2$, and $F(3)=Y$. The fibred Grothendieck construction $E=\int F$ has objects $(b,x)$ for $x\in F(b)$. A morphism $(b,x)\to(c,y)$ consists of an arrow $u:b\to c$ in the base together with an arrow $x\to F(u)(y)$ in $F(b)$. Composition uses the contravariance of $F$.

The vertex $3$ is initial in $\widetilde\Phi$. Its unique arrows $3\to b$ therefore define a functor
\[
 r:E\longrightarrow Y,
 \qquad r(b,x)=F(3\to b)(x).
\]
For $y\in Y$ there is a natural bijection
\[
 \operatorname{Hom}_E\bigl((3,y),(b,x)\bigr)
   =\operatorname{Hom}_Y\bigl(y,r(b,x)\bigr).
\]
Thus the fibre inclusion $i_Y:Y\to E$ is left adjoint to $r$, with $ri_Y=1_Y$. The counit $i_Yr\Rightarrow1_E$ is the arrow over $3\to b$ represented by the identity of $F(3\to b)x$. This displays the homotopy equivalence of $Y$ with $E$ directly.

Now restrict to $\mathcal X=E|_{\{0,1,2\}}$. The restricted functor $r|_{\mathcal X}$ is exactly the source's comparison $\varphi_Q:\mathcal X\to Y$. Restricting the counit gives
\[
 \alpha:i_Y\varphi_Q\Longrightarrow i,
 \qquad i:\mathcal X\hookrightarrow E.
\]
This is an arrow between functors with source $\mathcal X$ and target $E$. In the edition it is drawn inside the triangle. The transcription had instead attached it to the vertices $Y$ and $\mathcal X$; the correction register restores the interior arrow.

To finish the criterion, a natural transformation $g\Rightarrow h$ determines a functor $\mathcal X\times[1]\to E$. The projection $\mathcal X\times[1]\to\mathcal X$ and its two sections belong to the localizer. Consequently $g$ belongs to $W$ exactly when $h$ does. Applying this to $\alpha$, and then using $i_Y\in W$, gives
\[
 i\in W
 \ \Longleftrightarrow\ i_Y\varphi_Q\in W
 \ \Longleftrightarrow\ \varphi_Q\in W.
\]
The triangle need not commute strictly. Its specified transformation is precisely enough because the localizer already identifies naturally homotopic functors.
\end{dcomment}


On sait que ceci implique que
\[
i \in W \qquad\qquad \Longleftrightarrow \qquad\qquad (i_Y\varphi_Q) \in W
\]
(cf. lemme p. 2), et comme $i_Y \in W$,
\[
(i_Y\varphi_Q) \in W \qquad\qquad \Longleftrightarrow \qquad\qquad \varphi_Q \in W,
\]
d'où
\[
i \in W \qquad\qquad \Longleftrightarrow \qquad\qquad \varphi_Q \in W,
\]

\hfill q.e.d.

Nous pouvons maintenant prouver

\msp{page 236}

\textbf{Proposition 18} (\footnote[127]{devrait être vrai dans toute catégorie de modèles ayant assez de cofibrations.}). (Moyennant W(7 bis).) \textit{Considérons un diagramme commutatif}
\[
\begin{tikzcd}[ampersand replacement=\&]
X_0 \arrow[r] \arrow[d] \arrow[rd, phantom, "1"] \& X_1 \arrow[r] \arrow[d] \arrow[rd, phantom, "2"] \& X_2 \arrow[d] \\
X'_0 \arrow[r] \& X'_1 \arrow[r] \& X'_2
\end{tikzcd}
\tag*{$(\ast)$}
\]
\textit{d'où un carré composé}
\[
\begin{tikzcd}[ampersand replacement=\&]
X_0 \arrow[r] \arrow[d] \arrow[rd, phantom, "3"] \& X_2 \arrow[d] \\
X'_0 \arrow[r] \& X'_2.
\end{tikzcd}
\]
\textit{Supposons que le carré} 1 \textit{soit $W$-cocartésien. Alors le carré} 2 \textit{est $W$-cocartésien si et seulement si le carré} 3 \textit{l'est.}

\textsc{Démonstration.}

\texttt{[Maintenant, }$X_i$\texttt{ et }$X'_i$\texttt{ échangent leurs rôles.]}

\printedpage{152}
% head: A. Grothendieck : Dérivateurs, Ch. XII, Édition des Universités de Montpellier II et Paris VII
% foot: Version 30 mars 2003

Soit
\[
I \quad = \quad
\left(
\begin{tikzcd}[ampersand replacement=\&]
0 \arrow[d] \& 1 \arrow[l] \arrow[d] \& 2 \arrow[l] \arrow[d] \\
0' \& 1' \arrow[l] \& 2' \arrow[l]
\end{tikzcd}
\right)
\]
et soit $\cX$ l'intégrale du diagramme $(\ast)$, catégorie fibrée sur $\cX$. On désigne par $\cX_{010'}$ la sous-catégorie pleine de $\cX$ image inverse de la sous-catégorie pleine % [sic] fibrée sur $\cX$ -> fibrée sur $I$, a category fibred over itself; the base is the category $I$ just defined
\[
\underbrace{I_{010'}}_{=\{0,1,0'\}} \ \subseteq I,
\]
et de même pour toute autre partie $A$ de $\Ob I = \{0,1,2,0',1',2'\}$, donnant une sous-catégorie pleine $\cX_A$ de $\cX$. On pose
\[
U \quad = \quad
\begin{tikzcd}[ampersand replacement=\&]
0 \arrow[d] \& 1 \arrow[l] \& 2 \arrow[l] \\
0' \& \&
\end{tikzcd}
, \qquad
V \quad = \quad
\begin{tikzcd}[ampersand replacement=\&]
0 \arrow[d] \& 1 \arrow[l] \arrow[d] \& 2 \arrow[l] \\
0' \& 1' \arrow[l] \&
\end{tikzcd}
,
\]
d'où
\[
\cX_U = \cX_{0120'}, \qquad\qquad \cX_V = \cX_{0120'1'}.
\]
Les flèches entre les $\cX_A$ sont les inclusions

\msp{page 237}

déduites d'inclusions $A \hookrightarrow A'$ dans $\mathfrak{P}(\Ob I)$. Ainsi, le caractère $W$-cocartésien des carrés 1, 2, 3 s'exprime respectivement par la condition que les inclusions
\[
\begin{aligned}
\cX_{010'} \quad &\xrightarrow{\varphi_1} \quad \cX_{010'1'} \\
\cX_{121'} \quad &\xrightarrow{\varphi_2} \quad \cX_{121'2'} \\
\cX_{020'} \quad &\xrightarrow{\varphi_1} \quad \cX_{020'2'} % [sic] $\varphi_1$ -> $\varphi_3$, this inclusion is called $\varphi_3$ in the diagram that follows
\end{aligned}
\]
soient dans $W$. Considérons alors le diagramme de flèches d'inclusion
\printedpage{153}
% head: A. Grothendieck : Dérivateurs, Ch. XII, Édition des Universités de Montpellier II et Paris VII
% foot: Version 30 mars 2003
\[
\begin{tikzcd}[ampersand replacement=\&]
\cX_{020'} \arrow[r, "\varphi_3"] \arrow[d, "\alpha"'] \& \cX_{020'2'} \arrow[d, "\beta"] \\
\cX_U \arrow[r, "\psi_3"] \arrow[d, "\psi_1"'] \& \cX \arrow[d, "\mathrm{id}"] \\
\cX_V \arrow[r, "\psi_2"] \& \cX \\
\cX_{121'} \arrow[r, "\varphi_2"] \arrow[u, "\alpha'"] \& \cX_{121'2'}. \arrow[u, "\beta'"']
\end{tikzcd}
\]

Je dis que dans ce diagramme on a
\[
\alpha, \beta, \alpha', \beta' \in W
\tag*{$(\ast)$}
\]
et
\[
\varphi_1 \in W \ \ \text{(i.e. carré 1 $W$-cocartésien)} \qquad\qquad \Longleftrightarrow \qquad\qquad \psi_1 \in W.
\tag*{$(\ast\ast)$}
\]

Les relations $(\ast)$ impliquent déja que % [sic] déja -> déjà, missing grave accent
\[
\begin{aligned}
\varphi_2 \in W \quad &\Longleftrightarrow \quad \psi_2 \in W \\
\varphi_3 \in W \quad &\Longleftrightarrow \quad \psi_3 \in W,
\end{aligned}
\]
et sous l'hypothèse $W$-cocartésienne $(\ast\ast)$, impliquant $\psi_1 \in W$,

\msp{page 238}

on trouve
\[
\psi_2 \in W \qquad\qquad \Longleftrightarrow \qquad\qquad \psi_3 \in W,
\]
d'où
\[
\varphi_2 \in W \qquad\qquad \Longleftrightarrow \qquad\qquad \varphi_3 \in W,
\]
ce qu'il fallait prouver.

Il reste à établir $(\ast)$ et $(\ast\ast)$. Les relations $(\ast)$ proviennent de la remarque que si $i : A \longrightarrow A'$ est une inclusion de sous-catégories pleines de $I$, telle que $i$ soit $W$-\textit{asphérique} (i.e. $A/j$ $W$-asphérique pour $j \in \Ob A'\backslash\Ob A$), alors $\cX_A \longrightarrow \cX_{A'}$ l'est aussi (changement de base lisse), donc est \texttt{[dans]} $W$. On doit vérifier la $W$-asphéricité pour les quatre inclusions
\[
\begin{array}{rcll}
\{0,2,0'\} & \longrightarrow & U = \{0,1,2,0'\} & \qquad\qquad \mathbf{\Delta}^0 \\[1ex]
\{0,2,0',2'\} & \longrightarrow & I & \qquad\qquad \mathbf{\Delta}^0/\mathbf{\Delta}^1 \\[1ex]
\{1,2,1'\} & \longrightarrow & V = \{0,1,2,0',1'\} & \qquad\qquad \mathbf{\Delta}^1,\ \Phi \\[1ex]
\{1,2,1',2'\} & \longrightarrow & I & \qquad\qquad \mathbf{\Delta}^1,\ \mathbf{\Delta}^1 \times \mathbf{\Delta}^1.
\end{array}
\]

\dgrblock{dgrb581}{The seven slices behind the pasting theorem}
\begin{dcomment}{examples}{p238:On doit vérifier la $W$-asphéricité}{The seven slices behind the pasting theorem}
The six-object poset is the bookkeeping device that puts the left square, right square and outer rectangle in one ambient diagram. Grothendieck reduces the substantial part of pasting to seven very small slice categories. Computing them explicitly both explains the proof and corrects one harmless shape label in its margin.

Use the arrows of $I$ exactly as printed: the top row is $2\to1\to0$, the bottom row is $2'\to1'\to0'$, and the vertical arrows are $k\to k'$. A slice $A/j$ has as objects the elements of $A$ admitting an arrow to $j$. Since $I$ is a poset, its morphisms are just the induced order. For the inclusions underlying $\alpha,\beta,\alpha',\beta'$ we obtain
\[
\begin{array}{c|c|l}
 \text{inclusion}&j&A/j\\ \hline
 \alpha&1&\{2\}\\
 \beta&1&\{2\}\\
 \beta&1'&(2\longrightarrow2')\\
 \alpha'&0&(2\longrightarrow1)\\
 \alpha'&0'&(2\longrightarrow1\longrightarrow1')\\
 \beta'&0&(2\longrightarrow1)\\
 \beta'&0'&\{1,2,1',2'\}\cong[1]\times[1].
\end{array}
\]
Each displayed category has initial object $2$. For an object $j$ already in $A$, the slice $A/j$ has terminal object $1_j$, so no separate check is necessary. In the fifth line the composite $2\to1'$ exists: the inclusion is full, and the defining rectangle is commutative. The slice is consequently $[2]=\Delta^2$. The edition prints $\Phi$ in this place. The asphericity argument uses the initial object and is unaffected by this label.

We can also make the passage from base slices to total-category slices explicit. Let $F:I^{\mathrm{op}}\to\mathrm{Cat}$ be the integrated diagram and let $e=(j,y)$ be an object of its total category. There is a projection
\[
 (\mathcal X_A)/e\longrightarrow A/j.
\]
An object above $u:a\to j$ is an arrow $x\to F(u)y$ in $F(a)$. The choice $x=F(u)y$, with the identity arrow, defines a section. For an object $z$ of the total slice, every arrow from its image in $A/j$ to $u$ has a unique lift from $z$ to this chosen object over $u$. Hence the projection is left adjoint to the section. The two categories are homotopy equivalent. Asphericity of $A/j$ therefore proves asphericity of $(\mathcal X_A)/e$; the localizer's Theorem~A condition gives $\mathcal X_A\to\mathcal X_{A'}\in W$. This is the ``smooth change of base'' step in concrete form.

The last ingredient is gluing. The two subsets
\[
 U=\{0,1,2,0'\},\qquad
 U'=\{0,1,0',1'\}
\]
are cosieves in $V=\{0,1,2,0',1'\}$: following an arrow from either subset never leaves that subset. Their total categories are therefore closed subcategories covering $\mathcal X_V$, with intersection $\mathcal X_{010'}$. The forward part of the closed-cover axiom sends $\varphi_1\in W$ to $\psi_1\in W$. As $\psi_3=\psi_2\psi_1$, two-out-of-three then says that $\psi_2\in W$ exactly when $\psi_3\in W$. The four slice equivalences transport this assertion back to $\varphi_2$ and $\varphi_3$.

The proof thus needs only the forward part of W(7 bis). The source writes the stronger equivalence $\varphi_1\in W\Longleftrightarrow\psi_1\in W$, which its full hypothesis permits, but its reverse direction is not used. Cancellation of the right square against the outer rectangle is produced by a known weak equivalence $\psi_1$ and two-out-of-three. It should be kept separate from the stronger conservative gluing condition examined in the appendix.
\end{dcomment}


\printedpage{154}
% head: A. Grothendieck : Dérivateurs, Ch. XII, Édition des Universités de Montpellier II et Paris VII
% foot: Version 30 mars 2003

(NB J'ai marqué en marge les catégories qu'on obtient comme $A/j$, à isomorphisme près, dans les quatre cas envisagés - elles ont toutes un objet initial, donc sont $W$-asphériques.)

Reste à prouver $(\ast\ast)$. Soit
\[
U' = \{0,1,0',1'\}, \qquad\qquad \text{donc } V = U \cup U',\ \ U \cap U' = \{0,1,0'\},
\]
où $U$, $U'$ sont des \textit{fermés} de $V$ qui recouvrent $V$. Donc on aura de même
\[
\cX_V = \cX_U \cup \cX_{U'}, \qquad\qquad \cX_U \cap \cX_{U'} = \cX_{010'},
\]

\msp{page 239}

\[
\begin{tikzcd}[ampersand replacement=\&]
\cX_{U\cap U'} \arrow[r] \arrow[d, "\varphi_1"] \& \cX_U \arrow[d, "\psi_1"] \\
\cX_{U'} \arrow[r] \& \cX_V.
\end{tikzcd}
\]

Donc par W(7 bis) (cas dual des immersions fermées), on a
\[
\varphi_1 \in W \qquad\qquad \Longleftrightarrow \qquad\qquad \psi_1 \in W,
\]

\hfill q.e.d.

Notons pour mémoire le corollaire qui nous avait motivé.

\textbf{Corollaire.} \textit{Supposons} W(7 bis) (comme dans toute cette section). \textit{Soit un carré cocartésien dans} Cat
\[
\begin{tikzcd}[ampersand replacement=\&]
X_0 \arrow[r, "i_1"] \arrow[d, "i_2"'] \& X_1 \arrow[d] \\
X_2 \arrow[r] \& X_3,
\end{tikzcd}
\]
\textit{avec $i_1$ (ou $i_2$) une $W$-cofibration (p.ex. une immersion de Dwyer). Alors ce carré est $W$-cocartésien.}

\dgrblock{dgrb582}{From categorical gluing to coherent cocartesian squares}
\begin{dcomment}{later}{p239:Alors ce carré est $W$-cocartésien.}{From categorical gluing to coherent cocartesian squares}
The corollary finishes the comparison that motivated this detour. Factor the second leg of a strict pushout as a Dwyer map followed by a weak equivalence. The first resulting square is homotopy cocartesian by the Dwyer case; the second is so because its opposite sides are weak equivalences, using the first leg's $W$-cofibration property. Proposition~18 identifies their composite as homotopy cocartesian. The categorical integral and the replacement construction therefore give the same answer on these pushouts.

In the modern derivator language, put $\square=[1]\times[1]$ and let $\ulcorner\subset\square$ omit $(1,1)$. A coherent square $Q\in\mathbb D(\square)$ is cocartesian when the counit
\[
 \iota_!\iota^*Q\longrightarrow Q,
 \qquad \iota:\ulcorner\hookrightarrow\square,
\]
is an isomorphism. The composition and cancellation rule for such squares is \citet[Proposition~3.13]{Groth2013}. It requires neither stability nor a zero object. This identifies the modern form of the assertion in footnote~127: the coherence belongs to an object of $\mathbb D([2]\times[1])$, and the left square controls equivalence of the two remaining cocartesianness tests.

For Dwyer maps the later result is stronger still. \citet[Theorems~1.5--1.6, arXiv version~3]{HackneyOzornovaRiehlRovelli2024} show that if $D=B\amalg_A C$ and $A\to B$ is Dwyer, then
\[
 NB\amalg_{NA}NC\longrightarrow ND
\]
is a weak categorical equivalence in the Joyal model structure. Thus the ordinary pushout represents the pushout of the corresponding $(\infty,1)$-categories, retaining categorical information before passage to the homotopy types of their classifying spaces. The hypotheses concern the attaching map, so the result applies directly to the finite cylinders above.

The source's final sentence on cofibres should be read at this finite level. It has exhibited functorial replacements, proved the relative invariance needed to compare them, and established the basic gluing calculus. Those constructions supply the material for the abstract homotopy-square theory in the next section. A claim that all homotopy Kan extensions exist for an arbitrary relative category would require further hypotheses; no such claim is established by this closing factorization alone.
\end{dcomment}


Cette digression (p. 232 ff.) n'est pas nécessaire d'ailleurs pour formuler dès maintenant la conclusion de nos cogitations sur les cofibres $W$-homotopiques : \textit{Il y a une telle théorie dans} Cat, pourvu seulement que $W$ satisfasse à W(7 bis). En effet, tout résulte du fait que toute flèche $f : X \longrightarrow Y$ dans Cat se factorise en
\[
X \xrightarrow{i_f} \tilde{Y} \xrightarrow{p_f} Y,
\]
avec $i_f$ une $W$-cofibration (et même une imm-

% ---- from Der-ChXIIc-30-03-03_p041-045.tex
\printedpage{155}
% head: A. Grothendieck : Dérivateurs, Ch. XII, Édition des Universités de Montpellier II et Paris VII
% foot: Version 30 mars 2003

\msp{page 240}
ersion ouverte de Dwyer), et $p_f \in W$. % beginning of a word continued from the preceding page (outside this range), as printed

\section*{7\quad Retour sur l'axiomatique de la fibre $W$-homotopique.}

Je reviens sur la réflexion du n$^\circ$ 4. Finalement, pour
\[
(\mathcal{M}, W \subseteq \mathrm{Fl}(\mathcal{M}))
\]
donnés avec $W$ faiblement saturé, il apparaît qu'une excellente théorie des fibres $W$-homotopiques, satisfaisant aux conditions \cercle{a}\cercle{b}\cercle{c} du n$^\circ$ 4 (pages 215-218), existe dès que l'axiome d'apparence anodine suivant est satisfait (en plus de l'existence de l'objet final et des produits fibrés (\footnote[128]{et encore! il me semble qu'il suffit de supposer que les $X \longrightarrow Y$ qui sont $W$-fibrants sont `quarrables'.})).

\textbf{Factorisation gauche :} (\footnote[129]{Sous ces conditions, $\mathcal{M}$ muni de $(W, F_W)$ est une catégorie à cofibrations de K.S. \textsc{Brown-Anderson}, et `tout est O.K.', mais il faudrait vérifier a) b) c) ci-dessus (p. 215, 218), je ne suis pas sûr de l'avoir fait.},\footnote[130]{\textbf{NB} L'axiome de factorisation gauche n'est suffisant pour l'existence d'une bonne théorie des fibres homotopiques (semble-t-il) que dans le cas où tout $X$ dans $\mathcal{M}$ est $W$-fibrant, i.e. est dans $\Fibu_W e$.}). Toute flèche $f : X \longrightarrow Y$ de $\mathcal{M}$ se factorise en
\[
X \xrightarrow{i_f} \bar{X} \xrightarrow{p_f} Y,
\]
avec
\[
i_f \in W, \qquad p_f \in \mathrm{Fib}_W.
\]

L'axiome dual donne lieu à une bonne théorie des cofibres $W$-homotopiques.

\textbf{Factorisation droite :} Toute flèche $f : X \longrightarrow Y$ dans $\mathcal{M}$ se factorise en
\[
X \xrightarrow{j_f} \tilde{Y} \xrightarrow{q_f} Y,
\]
avec
\[
q_f \in W, \qquad j_f \in \mathrm{Cofib}_W.
\]

\msp{page 241}
Je voudrais à présent développer quelque peu le formalisme des fibres $W$-homotopiques, sous l'hypothèse (disons) que $W$ satisfait l'axiome de factorisation gauche.

\dgrblock{dgrb583}{The factorization axiom and Brown's hypothesis}
\begin{dcomment}{context}{p240}{The factorization axiom and Brown's hypothesis}
Grothendieck now wants to detach the theory of homotopy fibres from the particular geometry of categories. The interval and mapping-cylinder arguments of the preceding pages have supplied a mechanism; he asks which part of that mechanism belongs to an arbitrary category with weak equivalences. His phrase ``catégorie de modèles'' here means a category equipped with a class $W$ and the properties under discussion. It does not silently import all the axioms of a Quillen model category.

The relevant meaning of $W$-fibration is likewise intrinsic to $W$. Precisely, for $p:X\to Y$, the pullback functor $p^*:\mathcal M/Y\to\mathcal M/X$ must preserve the arrows whose underlying maps belong to $W$. This also implies the same assertion after another base change. This is the property used below, together with the stability of these maps under composition and pullback. In particular, one must not replace $\mathrm{Fib}_W$ by the fibrations of an arbitrary model structure without checking the corresponding properness statement. Here ``weakly saturated'' means that $W$ contains the isomorphisms and satisfies two-out-of-three; it does not mean that every arrow inverted by localization already belongs to $W$. A factorization $f=pi$ with $i\in W$ and $p\in\mathrm{Fib}_W$ is sufficient for the replacement operations that are actually written down; it is not, by its form alone, a pair of lifting axioms.

The two footnotes beneath the axiom record a real issue in its scope. Brown's category of fibrant objects requires every map to the terminal object to be a fibration. Brown also requires pullbacks of fibrations and of trivial fibrations to exist and to remain of the same kind, and a path-object factorization of every diagonal. His factorization lemma produces a factorization of every map into a weak equivalence followed by a fibration; conversely, applying such a factorization to the diagonal supplies the path-object axiom \citep[Part~I, \S1, pp.~420--422]{Brown1973}. Thus the all-object-fibrancy qualification in footnote~130 gives the clean connection with Brown. The word ``cofibrations'' in footnote~129 is a slip for the fibration version of this theory.

Here is why a weakened existence condition on pullbacks is reasonable. If $S$ is fixed and $A\to S$, $B\to S$ are $W$-fibrations, their product over $S$ is formed by pulling back one of those fibrations. The projection to the other factor is again a fibration, and the composite to $S$ is a fibration. The same observation constructs a product with any finite number of factors; the empty product is $S\to S$. A factorization of a diagonal then provides the path objects required in this relative category. Arbitrary pullbacks in $\mathcal M$ are consequently stronger than what these constructions use. Footnote~128 is a proposed economy of axioms, rather than an obstruction to the argument.

The distinction between Brown's category and a category merely admitting fibrant replacement matters for the footnote's present status. In the language of left derivable categories, the written factorization and pullback hypotheses already suffice for the finite diagram theorem. We verify this below, including the required stability of trivial fibrations; no global-fibrancy hypothesis has to be added. Functorial factorization gives a shorter separate proof of the two replacement lemmas, but is also unnecessary in the derivable-category argument. Thus the source's cautious necessity suggestion is not part of the modern theorem.
\end{dcomment}


Considérons un diagramme cartésien
\[
\begin{tikzcd}[ampersand replacement=\&]
X_1 \arrow[d, "p_1"'] \& Y \arrow[l, "q_2"'] \arrow[d, "q_1"] \\
S \& X_2 \arrow[l, "p_2"]
\end{tikzcd}
\]
\printedpage{156}
% head: A. Grothendieck : Dérivateurs, Ch. XII, Édition des Universités de Montpellier II et Paris VII
% foot: Version 30 mars 2003
dans $\mathcal{M}$. Si $p_1$ est dans $\mathrm{Fib}_W$, alors $Y$ sur $X_2$ peut s'interpréter comme la fibre $W$-homotopique de $X_1/S$ en $X_2/S$. Plus précisément, il en est ainsi de la classe $\ho^{\lc}_{W,X_2}(Y)$,
\[
\ho^{\lc}_{W,X_2}(Y) \simeq p_2^*(\ho^{\lc}_{W,S}(X_1/S)). \tag*{$(*)$}
\]
Mais en tant qu'objet sur $X_1$, $Y$ peut s'interpréter également comme fibre $W$-homotopique de $X_2/S$ en $X_1/S$ :
\[
\ho^{\lc}_{W,X_1}(Y/X_1) \simeq p_1^*(\ho^{\lc}_{W,S}(X_2/S)). \tag*{$(**)$}
\]
À dire vrai, il faudrait le prouver. Pour ceci, considérons une factorisation
\[
X_2 \xrightarrow{i_2} \bar{X}_2 \xrightarrow{\bar{p}_2} S \qquad i_2 \in W\ p_2 \in \mathrm{Fib}_W, % [sic] p_2 \in \mathrm{Fib}_W -> \bar{p}_2 \in \mathrm{Fib}_W, the factorisation just written is through \bar{p}_2
\]
de sorte qu'on trouve un diagramme de carrés cartésiens
\[
\begin{tikzcd}[ampersand replacement=\&]
X_1 \arrow[d] \& \bar{Y} \arrow[l, "\bar{q}_2"'] \arrow[d] \& Y \arrow[l, "j_2"', "W"] \arrow[d] \\
S \& \bar{X}_2 \arrow[l, "\bar{q}_2"'] \& X_2. \arrow[l, "i_2"', "W"] % [sic] \bar{q}_2 (bottom left arrow) -> \bar{p}_2, the arrow \bar{X}_2 \to S is \bar{p}_2 of the factorisation above
\end{tikzcd}
\]

\msp{page 242}
Par définition de $p_1^*$, on a
\[
p_1^*(\ho^{\lc}_{W,S}(X_2/S)) \simeq \ho^{\lc}_{W,X_1}(\bar{Y}/X_1),
\]
d'autre part, comme $X_1 \longrightarrow S$ est dans $\mathrm{Fib}_W$, et $i_2 \in W$, on a $j_2 \in W$, donc
\[
\ho^{\lc}_{W,X_1}(Y/X_1) \xrightarrow{\sim} \ho^{\lc}_{W,X_1}(\bar{Y}/X_1),
\]
d'où la formule $(**)$.

Ainsi, l'objet $Y$ a deux interprétations homotopiques différentes, suivant qu'on le regarde comme étant sur $X_1$, ou sur $X_2$. On aimerait alors une interprétation qui englobe les deux à la fois. On va considérer le diagramme
\[
\begin{tikzcd}[ampersand replacement=\&]
X_1 \arrow[d] \& \\
S \& X_2 \arrow[l]
\end{tikzcd}
\]
de type $\Psi^\circ$, ou
\[
\Psi \quad = \quad
\begin{tikzcd}[ampersand replacement=\&]
1 \& \\
0 \arrow[u] \arrow[r] \& 2
\end{tikzcd}
\qquad \text{ou} \qquad
\begin{tikzcd}[ampersand replacement=\&, row sep=small, column sep=small]
 \& 1 \\
0 \arrow[ur] \arrow[dr] \& \\
 \& 2
\end{tikzcd}
\quad ,
\]
% the last "2" of the first diagram is printed overlapping the page number 156
\printedpage{157}
% head: A. Grothendieck : Dérivateurs, Ch. XII, Édition des Universités de Montpellier II et Paris VII
% foot: Version 30 mars 2003
dont $Y$ se déduit par produit fibré, en complétant ce diagramme en un diagramme de type $\tilde{\Psi}^\circ$, où
\[
\tilde{\Psi} =
\begin{tikzcd}[ampersand replacement=\&, row sep=small, column sep=small]
 \& 1 \arrow[dr] \& \\
0 \arrow[ur] \arrow[dr] \& \& 3 \\
 \& 2 \arrow[ur] \&
\end{tikzcd}
\qquad (\simeq \boldsymbol{\Delta}^1 \times \boldsymbol{\Delta}^1).
\]

\msp{page 243}
Pour toute catégorie $I$, soit
\[
\Ho_{\mathcal{M},W}(I) \text{ \ ou \ } \Ho_W(I) \overset{\text{déf}}{=} \mathcal{M}(I)W(I)^{-1},
\]
où
\[
\left\{
\begin{array}{rcl}
\mathcal{M}(I) & = & \Homu(I^\circ, \mathcal{M}) \\[1ex]
W(I) & \subseteq & \mathrm{Fl}(\mathcal{M}(I)) \text{ est formé des } u : F \longrightarrow G \text{ telles que} \\[1ex]
 & & u_i : F(i) \longrightarrow G(i) \text{ dans } W \text{ pour tout } i \text{ dans } \mathrm{Ob}\, I.
\end{array}
\right.
\]

\dgrblock{dgrb584}{One pullback with two relative meanings}
\begin{dcomment}{exposition}{p242}{One pullback with two relative meanings}
The symmetry that Grothendieck is seeking concerns the whole square. A pullback $P=A\times_S B$ has a map to $A$ and a map to $B$; a construction that remembers only its isomorphism class in a homotopy category has forgotten part of the answer. Suppose $A\to S$ is a $W$-fibration and factor $B\to S$ as $B\xrightarrow{w}\bar B\to S$, with $w\in W$ and $\bar B\to S$ a $W$-fibration. Then
\[
 A\times_S B\longrightarrow A\times_S\bar B
\]
is the pullback of $w$ along the base change of $A\to S$. It belongs to $W$. It is also a morphism over $A$, so it identifies the class of the original pullback over $A$ with the result of deriving the base-change operation on $B$. This proves the second interpretation in the text, while the first interpretation is already computed by the good representative $A\to S$. No second fibration hypothesis on $B\to S$ has been introduced.

The natural domain of this simultaneous construction is therefore the category of cospans before localization. With the manuscript's contravariant convention,
\[
 \Psi=(1\leftarrow0\to2),\qquad
 \mathcal M(\Psi)=\operatorname{Fun}(\Psi^{\mathrm{op}},\mathcal M),
\]
so an object of $\mathcal M(\Psi)$ is $A\to S\leftarrow B$. The additional vertex $3$ of $\tilde\Psi$ points, after taking the opposite, to both $1$ and $2$. A right Kan extension along $\Psi^{\mathrm{op}}\hookrightarrow\tilde\Psi^{\mathrm{op}}$ consequently puts a limit at vertex~$3$. This orientation explains the star in $i_*$: it is a right adjoint to restriction.

One should retain the distinction already emphasized in the commentary on Chapter~VI, manuscript p.~87. The category
\[
 \operatorname{Fun}(\Psi^{\mathrm{op}},\mathcal M)[W(\Psi)^{-1}]
\]
retains how the maps and their weak-equivalence comparisons occur together. It is not being replaced by $\operatorname{Fun}(\Psi^{\mathrm{op}},\mathcal M[W^{-1}])$. In the latter, a commutativity equation is only an equality of homotopy classes. The former is the diagram category on which a canonical homotopy completion can be defined. This is precisely the reason that the later derivator formalism forms diagram categories first and localizes afterwards \citep[pp.~313--315]{Groth2013}.

The coming proof accordingly has two tasks. It must show that the chosen good cospans present the localization, and then show that taking their strict pullbacks respects weak equivalences. The first task concerns all morphisms and relations in a localized category; replacing one object at a time does not settle it.
\end{dcomment}


Ainsi, le diagramme $\vcenter{\hbox{\begin{tikzcd}[ampersand replacement=\&, row sep=small, column sep=small]
 \& X_1 \arrow[dl] \\
S \& \\
 \& X_2 \arrow[ul]
\end{tikzcd}}}$ définit un objet $\xi$ de $\mathcal{M}(\Psi)$, donc de $\Ho_W(\Psi)$, et le diagramme $\vcenter{\hbox{\begin{tikzcd}[ampersand replacement=\&, row sep=small, column sep=small]
 \& X_1 \arrow[dl] \& \\
S \& \& Y \arrow[ul] \arrow[dl] \\
 \& X_2 \arrow[ul] \&
\end{tikzcd}}}$ peut de même s'interpréter comme un objet $\tilde{\xi}$ de $\Ho_W(\tilde{\Psi})$, dont la restriction à $\Psi$ est $\xi$.

Ceci posé, j'aimerais définir un foncteur
\[
i_* : \Ho_W(\Psi) \longrightarrow \Ho_W(\tilde{\Psi})
\]
(où
\[
i : \Psi \longrightarrow \tilde{\Psi}
\]
est l'inclusion, qui est une immersion ouverte), lequel doit correspondre, dans le cas d'un objet $\vcenter{\hbox{\begin{tikzcd}[ampersand replacement=\&, row sep=small, column sep=small]
 \& X_1 \arrow[dl] \\
S \& \\
 \& X_2 \arrow[ul]
\end{tikzcd}}}$ de $\Ho_W(\Psi)$ avec $p_1 \in \mathrm{Fib}_W$ \texttt{[}$:$ $X_1 \longrightarrow S$\texttt{]}, à l'opération de produit fibré (\footnote[131]{\textbf{NB} Cette construction, qui se fait ici assez lourdement, devient immédiate, et aussi ses principales propriétés, si on admet que la catégorie à fibrations $(\mathcal{M}, W, \mathrm{Fib}_W)$ définit un \textit{dérivateur} sur le domaine des ensembles ordonnés finis, ce qui se prouve en effet assez simplement (avec pas plus de mal que je ne me donne ici).}).

\printedpage{158}
% head: A. Grothendieck : Dérivateurs, Ch. XII, Édition des Universités de Montpellier II et Paris VII
% foot: Version 30 mars 2003

\msp{page 244}
La construction de ce foncteur $i_*$ devrait être essentiellement sorital, avec ce dont nous disposons. Nous allons définir d'abord, pour $S$ \textit{fixé,} un foncteur canonique
\[
\Ho^{\lc}_W(S) \times \Ho^{\lc}_W(S) \longrightarrow \Ho_{W,S}(\Psi), \tag*{$(*)$}
\]
où le deuxième membre est défini, par analogie avec les $\Ho_W(I)$, par
\[
\Ho_{W,S}(I) \overset{\text{déf}}{=} ((\mathcal{M}/S)(I))(W_S(I))^{-1}
\]
- i.e. c'est la même définition, sauf que $(\mathcal{M}, W)$ est remplacé par $(\mathcal{M}/S, W_S)$. Quant à $\Phi = \vcenter{\hbox{\begin{tikzcd}[ampersand replacement=\&, row sep=small, column sep=small]
1 \arrow[dr] \& \\
 \& 3 \\
2 \arrow[ur] \&
\end{tikzcd}}}$ , on le considère comme une sous-catégorie fermée de $\tilde{\Psi}$, qui s'en déduit en lui rajoutant un objet initial 0 (tout comme $\tilde{\Psi}$ se déduit de $\Psi$ en lui rajoutant un objet final 3). D'autre part, on a un foncteur tautologique
\[
\Ho_{W,S}(\Phi) \longrightarrow \Ho_W(\tilde{\Psi}), \tag*{$(*')$}
\]
qui sur les objet consiste à compléter un diagramme $\vcenter{\hbox{\begin{tikzcd}[ampersand replacement=\&, row sep=small, column sep=small] % [sic] les objet -> les objets, plural agreement
X_1 \& \\
 \& Y \arrow[ul] \arrow[dl] \\
X_2 \&
\end{tikzcd}}}$ de $\mathcal{M}/S$ en le diagramme carré $\vcenter{\hbox{\begin{tikzcd}[ampersand replacement=\&, row sep=small, column sep=small]
 \& X_1 \arrow[dl] \& \\
S \& \& Y \arrow[ul] \arrow[dl] \\
 \& X_2 \arrow[ul] \&
\end{tikzcd}}}$ de $\mathcal{M}$.

\msp{page 245}
Le composé de $(*)$ et de $(*')$ donnera donc un foncteur
\[
\Ho^{\lc}_W(S) \times \Ho^{\lc}_W(S) \longrightarrow \Ho_W(\tilde{\Psi}), \tag*{$(**)$}
\]
très proche déjà du foncteur $i_*$ à construire. Pour construire $(*)$, on interprète (à équivalence de catégories près) le premier membre comme
\[
(\Fibu_W S \times \Fibu_W S)(W_S \times W_S)^{-1},
\]
donc il faut construire un foncteur
\[
\Fibu_W S \times \Fibu_W S \longrightarrow \Ho_{W,S}(\Phi),
\]
qui transforme flèches dans $W_S \times W_S$ en isomorphismes. On prendra le foncteur
\[
(X_1, X_2) \ \longmapsto \
\left(
\begin{tikzcd}[ampersand replacement=\&]
X_1 \& X_1 \times_S X_2 \arrow[l] \arrow[d] \\
 \& X_2
\end{tikzcd}
\right),
\]
\printedpage{159}
% head: A. Grothendieck : Dérivateurs, Ch. XII, Édition des Universités de Montpellier II et Paris VII
% foot: Version 30 mars 2003
et il faut seulement vérifier que si
\[
X_1 \xrightarrow{\alpha_1} X_1', \qquad X_2 \xrightarrow{\alpha_2} X_2'
\]
sont deux flèches dans $\Fibu_W S$, d'où
\[
X_1 \times_S X_2 \xrightarrow{\alpha = \alpha_1 \times_S \alpha_2} X_1' \times_S X_2',
\]
alors
\[
\alpha_1 \in W,\ \alpha_2 \in W \qquad \Longrightarrow \qquad \alpha \in W,
\]
ce qui est immédiat sur la factorisation
\[
X_1 \times_S X_2 \quad \underbrace{\longrightarrow}_{\substack{\in W \\ \text{car } X_2 \in \Fibu_W S}} \quad X_1' \times_S X_2 \quad \underbrace{\longrightarrow}_{\substack{\in W \\ \text{car } X_1' \in \Fibu_W S}} \quad X_1' \times_S X_2'.
\]

\msp{page 246}
On a prouvé la

\textbf{Proposition 18.} \textit{Soit $S$ dans $\mathcal{M}$. Il existe, à isomorphisme unique près, un foncteur et un seul}
\[
\Ho^{\lc}_W(S) \times \Ho^{\lc}_W(S) \xrightarrow{\pi} \Ho_{W,S}(\Phi)
\]
\textit{donnant lieu à un carré commutatif à isomorphisme près}
\[
\resizebox{\linewidth}{!}{$\begin{tikzcd}[ampersand replacement=\&, column sep=large]
 \& {\scriptstyle (X_1,X_2)\ \longmapsto\ \begin{array}{@{}c@{}c@{}c@{}} \scriptstyle X_1 & & \\ & \nwarrow & \\ & & \scriptstyle X_1 \times_S X_2 \\ & \swarrow & \\ \scriptstyle X_2 & & \end{array}} \& \\[-3ex]
\Fibu_W S \times \Fibu_W S \arrow[rr] \arrow[d] \& \& (\mathcal{M}/S)(\Phi) \arrow[d] \\
\Ho^{\lc}_W(S) \times \Ho^{\lc}_W(S) \arrow[rr] \& \& \Ho_{W,S}(\Phi) \\[-4ex]
(= (\mathcal{M}/S)(W_S)^{-1} \times (\mathcal{M}/S)(W_S)^{-1}) \& \& (= \mathcal{M}/S(\Phi)(W_S(\Phi))^{-1}).
\end{tikzcd}$}
\]
\textit{Pour $X_1$, $X_2$ dans $\mathcal{M}/S$, le diagramme}
\[
\pi(X_1, X_2) \overset{\text{déf}}{=} \pi(\ho^{\lc}_{W,S}(X_1), \ho^{\lc}_{W,S}(X_2))
\]
\textit{se `calcule' ainsi, à isomorphisme canonique près : on choisit}
\[
X_1 \xrightarrow{\alpha_1} \bar{X}_1, \qquad X_2 \xrightarrow{\alpha_2} \bar{X}_2 \qquad \textit{flèches dans } \mathcal{M}/S
\]
\textit{avec}
\[
\alpha_1, \alpha_2 \in W_S, \qquad \bar{X}_1, \bar{X}_2 \in \mathrm{Ob}\, \Fibu_W S,
\]
\textit{ce qui définit}
\[
\pi(\alpha_1, \alpha_2) : \pi(X_1, X_2) \xrightarrow{\sim} \pi(\bar{X}_1, \bar{X}_2) = \ho_{S,W,\Phi}
\left(
\begin{tikzcd}[ampersand replacement=\&]
\bar{X}_1 \& \bar{X}_1 \times_S \bar{X}_2 \arrow[l] \arrow[d] \\
 \& \bar{X}_2
\end{tikzcd}
\right).
\]

% ---- from Der-ChXIIc-30-03-03_p046-050.tex
\setcounter{footnote}{131}

\printedpage{160}
% head: A. Grothendieck : Dérivateurs, Ch. XII, Édition des Universités de Montpellier II et Paris VII
% foot: Version 30 mars 2003

\msp{page 247}

\dgrblock{dgrb585}{A product together with its projections}
\begin{dcomment}{exposition}{p246}{A product together with its projections}
The intermediate construction over a fixed base is more than a calculation of the object $A\times_S B$. Grothendieck keeps both projections in the value of $\pi$. Thus $\pi$ lands in the localized category of $\Phi^{\mathrm{op}}$-diagrams in $\mathcal M/S$, where $\Phi=(1\to3\leftarrow2)$. The occurrence of $\Psi$ in the first display of ms.~244 is an edition slip: the next functor, and Proposition~18 itself, require $\Phi$.

Write $\mathcal C_S$ for the full subcategory of $\mathcal M/S$ on the $W$-fibrations. If $a:A\to A'$ and $b:B\to B'$ are weak equivalences in this category, the comparison of products has the factorization
\[
 A\times_S B\xrightarrow{a\times_S1}
 A'\times_S B\xrightarrow{1\times_S b}
 A'\times_S B'.
\]
The first arrow is a pullback of $a$ along a base change of $B\to S$; the second is a pullback of $b$ along a base change of $A'\to S$. Both belong to $W$. The same morphism of spans has $a$ and $b$ at its other vertices, so it is an objectwise weak equivalence. This proves exactly the compatibility needed to descend the entire span functor.

There is also a small categorical point behind the source's product of localizations. For relative categories $(\mathcal C,U)$ and $(\mathcal D,V)$, with identities in $U,V$, there is a canonical equivalence
\[
 (\mathcal C\times\mathcal D)[(U\times V)^{-1}]
 \simeq\mathcal C[U^{-1}]\times\mathcal D[V^{-1}].
\]
Indeed a functor out of $\mathcal C\times\mathcal D$ inverts $U\times V$ precisely when it inverts $(u,1)$ and $(1,v)$ separately. Localizing successively in the two variables gives the displayed universal property, including natural transformations. Applied to $\mathcal C_S$ twice, this justifies the domain used in Proposition~18. The localization equivalence between $\mathcal C_S$ and the relative homotopy category is an earlier input to this paragraph, not a consequence of the displayed product calculation.

The phrase ``unique up to unique isomorphism'' is understood relative to the specified comparison with the functor on good representatives. Without that comparison, a functor can have automorphisms and the uniqueness assertion would be too strong. With it, full faithfulness of restriction along a localization supplies the required uniqueness. A pair of replacement arrows $A\to\bar A$, $B\to\bar B$ therefore gives the canonical comparison written by Grothendieck; it does not require a preferred strict map between every two independently chosen replacements.

This proposition is numbered~18 a second time in the edition. We distinguish it from the pasting proposition on ms.~236 by its page and by the name ``fixed-base product proposition''. Its content will also be used when the base itself varies, where a further argument is needed.
\end{dcomment}


\textbf{Corollaire.} {\itshape Description du foncteur
\[
(\xi_1,\xi_2)\longmapsto\xi_1\times_S\xi_2 : \Ho^{\lc}_W(S)\times\Ho^{\lc}_W(S)\longrightarrow\Ho^{\lc}_W(S).
\]
Ce foncteur est bel et bien un produit cartésien dans} $\Ho^{\lc}_W(S)$ (\footnote{mais c'est un cas particulier des cas de $\Ho_W$ (si $\mathcal{M}$ a des produits binaires), en remplaçant $\mathcal{M}$ par $\mathcal{M}/S$.}), {\itshape où les produits finis existent} (\footnote{Dire aussi que les foncteurs $f^* : \Ho^{\lc}_W(S)\longrightarrow\Ho^{\lc}_W(S')$ (pour $f : S'\longrightarrow S$) commutent aux produits finis.}).

Pour prouver ce dernier point, posant
\[
\mathcal{H}=\Ho^{\lc}_W(S),
\]
il faut montrer que les foncteurs
\[
\begin{tikzcd}[ampersand replacement=\&]
\mathcal{H} \arrow[r, shift left, "\delta"] \& \mathcal{H}\times\mathcal{H} \arrow[l, shift left, "\pi"]
\end{tikzcd}
\]
sont adjoints, avec
\[
\begin{array}{rclcrcl}
\delta & = & \mathrm{diag}_{\mathcal{H}} & : & \xi & \longmapsto & (\xi,\xi)\\
\pi & = & \text{`produit'} & : & (\xi_1,\xi_2) & \longmapsto & \xi_1\times_S\xi_2.
\end{array}
\]
On va définir les homomorphismes d'adjonction
\[
\begin{array}{rcllrcl}
\mathrm{id}_{\mathcal{H}} & \xrightarrow{\alpha} & \pi\delta & \text{i.e.} & \xi & \xrightarrow{\alpha_\xi} & \xi\times_S\xi\\
\mathrm{id}_{\mathcal{H}\times\mathcal{H}} & \xleftarrow{\beta} & \delta\pi & \text{i.e.} & (\xi_1,\xi_2) & \xleftarrow{\beta_{\xi_1,\xi_2}} & (\xi_1\times_S\xi_2,\xi_1\times_S\xi_2).
\end{array}
\]
Pour ceci, on décrit $\mathcal{H}$ comme
\[
\mathcal{H}=(\uFib_W S)(W_S)^{-1}
\]
et on note que $\delta$, $\pi$ proviennent de foncteurs adjoints au niveau de $\mathcal{H}_0=\uFib_W S$
\[
\begin{tikzcd}[ampersand replacement=\&]
\mathcal{H}_0 \arrow[r, shift left, "\delta_0"] \& \mathcal{H}_0\times\mathcal{H}_0 \arrow[l, shift left, "\pi_0"]
\end{tikzcd},
\]
lesquels donnent lieu à $\alpha_0$, $\beta_0$
\[
\begin{array}{rcl}
\mathrm{id}_{\mathcal{H}_0} & \xrightarrow{\alpha_0} & \pi_0\delta_0\\
\mathrm{id}_{\mathcal{H}_0\times\mathcal{H}_0} & \xleftarrow{\beta_0} & \delta_0\pi_0.
\end{array}
\]

\printedpage{161}
% head: A. Grothendieck : Dérivateurs, Ch. XII, Édition des Universités de Montpellier II et Paris VII
% foot: Version 30 mars 2003

\msp{page 248}

On définit $\alpha$, $\beta$ par passage aux catégories de fractions pour $\alpha_0$, $\beta_0$. Les 2 compatibilités, étant vérifiées pour $\alpha_0$, $\beta_0$, le sont aussi pour $\alpha$, $\beta$. Q.e.d.

Je dégage ici le

\textbf{Lemme.} {\itshape Soit \begin{tikzcd}[ampersand replacement=\&]
\mathcal{X} \arrow[r, shift left, "f"] \& \mathcal{Y} \arrow[l, shift left, "g"]
\end{tikzcd} un couple de foncteurs adjoint, et soient $\Sigma$, $\Sigma'$ des localiseurs dans $\mathcal{X}$, $\mathcal{Y}$ tels que $f(\Sigma)\subseteq\Sigma'$, $g(\Sigma')\subseteq\Sigma$, de sorte qu'on a un couple de foncteurs
\[
\begin{tikzcd}[ampersand replacement=\&]
\mathcal{X}\Sigma^{-1} \arrow[r, shift left, "\bar{f}"] \& \mathcal{Y}\Sigma'^{-1} \arrow[l, shift left, "\bar{g}"]
\end{tikzcd}.
\]
Alors $\bar{f}$, $\bar{g}$ sont adjoints, les flèches d'adjonction étant déduites de celles pour $f$, $g$.}

\textbf{Corollaire.} {\itshape Soit $(\mathcal{M},W)$ une catégorie de modèles, supposons que dans $\mathcal{M}$ les produits binaires (resp. les produits finis) existent, et que pour $u_1 : X_1\longrightarrow X'_1$ et $u_2 : X_2\longrightarrow X'_2$, si $u_1,u_2\in W$, alors $u_1\times u_2 : X_1\times X_2\longrightarrow X'_1\times X'_2$ est dans $W$. Alors dans $\mathcal{M}W^{-1}$ les produits binaires (resp. produits finis) existent, et le foncteur localisation $\mathcal{M}\longrightarrow\mathcal{M}W^{-1}$ y commute.}

On applique le lemme à
\[
\begin{tikzcd}[ampersand replacement=\&]
\mathcal{M} \arrow[r, shift left, "\delta"] \& \mathcal{M}\times\mathcal{M} \arrow[l, shift left, "\pi"]
\end{tikzcd}
\qquad \text{et } W,\ W\times W,
\]
et pour l'objet final,
\[
\begin{tikzcd}[ampersand replacement=\&]
\mathcal{M} \arrow[r, shift left, "e"] \& e \arrow[l, shift left, "\bullet"]
\end{tikzcd}
\qquad\qquad \text{et } W,\ W_0=\{\mathrm{id}_e\}.
\]

\msp{page 249}

\dgrblock{dgrb586}{Why the product remains a categorical product}
\begin{dcomment}{exposition}{p248}{Why the product remains a categorical product}
The assertion about products in $\Ho_W^{\mathrm{lc}}(S)$ is stronger than the existence of a homotopically well-defined binary operation. Grothendieck proves the universal property by descending an adjunction. This is an efficient place to make the small localization argument completely explicit.

Suppose $F:\mathcal C\rightleftarrows\mathcal D:G$ is an adjunction, with $F$ left adjoint to $G$, and both functors preserve the chosen weak equivalences. Its unit $\eta:1\to GF$ and counit $\epsilon:FG\to1$ have images in the localizations. Their naturality for an inverted weak equivalence follows from their naturality for that weak equivalence: multiply the naturality square by the inverse image of the equivalence. Naturality then holds for any composite of ordinary arrows and formal inverses, which are the morphisms of the localization. The identities
\[
 \epsilon_{F X}F(\eta_X)=1_{F X},\qquad
 G(\epsilon_Y)\eta_{G Y}=1_{G Y}
\]
are equations already true before localization, hence remain true afterwards. The descended unit and counit establish $\bar F\dashv\bar G$. No lifting property and no fullness assertion about the localization functor are needed.

Apply this argument to the diagonal functor on $\mathcal C_S$ and its right adjoint, the ordinary product over $S$. The preceding block proves that this right adjoint preserves weak equivalences. The unit is the diagonal $A\to A\times_S A$; the counit consists of the two projections. Consequently the operation constructed from good representatives is the categorical product in the localized relative category. Applying the same argument to the unique functor $\mathcal C_S\to e$ and the functor selecting $S\to S$ proves that this object remains terminal. This supplies all finite products.

The statement in footnote~133 has an equally concrete proof. Pullback along $f:S'\to S$ sends good objects to good objects and, by the preceding relative invariance theorem, sends weak equivalences between them to weak equivalences. The canonical strict isomorphism
\[
 f^*(A\times_S B)\cong f^*A\times_{S'}f^*B
\]
is compatible with both projections; also $f^*(S\to S)=(S'\to S')$. Their images after localization are therefore the comparison isomorphisms for a functor preserving finite products. This proof concerns the product in the relative homotopy category. It does not assert that an arbitrary strict pullback in $\mathcal M$ is already a homotopy pullback.
\end{dcomment}


Je vais maintenant construire le foncteur
\[
i_* : \Ho_W(\Psi)\longrightarrow\Ho_W(\tilde{\Psi}).
\]
Il me faut d'abord prouver le

\textbf{Lemme 1.} {\itshape Soit $\mathcal{F}$ la sous-catégorie pleine de $\mathcal{M}(\Psi)$, catégorie des diagrammes \begin{tikzcd}[ampersand replacement=\&, row sep=small, column sep=small]
 \& X_1 \arrow[dl, "p_1"'] \\
S \& \\
 \& X_2 \arrow[ul, "p_2"]
\end{tikzcd} dans $S$ \texttt{[plutôt dans }$\mathcal{M}$\texttt{, sur }$S$\texttt{]}, formée de ceux pour lesquels $p_1$, $p_2$ sont dans Fib$_W$. Soit $W_{\mathcal{F}}$ le localiseur induit sur $\mathcal{F}$ par $W(\Psi)$ ($W$-équivalences terme à terme). Alors
\[
\mathcal{F}W_{\mathcal{F}}^{-1}\xrightarrow{\sim}\mathcal{M}(\Psi)W(\Psi)^{-1}\ (\overset{\text{déf}}{=}\Ho_W(\Psi)).
\]}

\printedpage{162}
% head: A. Grothendieck : Dérivateurs, Ch. XII, Édition des Universités de Montpellier II et Paris VII
% foot: Version 30 mars 2003

Nous allons prouver ce lemme plus bas. En l'admettant, pour construire $i_*$, il suffit de construire
\[
\mathcal{F}\longrightarrow\Ho_W(\tilde{\Psi})
\]
qui applique $W_{\mathcal{F}}$ en des isomorphismes. Je prendrai, bien sûr,
\[
\begin{gathered}
\begin{tikzcd}[ampersand replacement=\&, row sep=small, column sep=small]
 \& X_1 \arrow[dl] \& \& \& \& X_1 \arrow[dl] \& \\
S \& \& \longmapsto \& \& S \& \& X_1\times_S X_2 \arrow[ul] \arrow[dl] \\
 \& X_2 \arrow[ul] \& \& \& \& X_2 \arrow[ul] \&
\end{tikzcd}\\
\mathcal{F}\quad\longrightarrow\quad\Ho_W(\Psi).
\end{gathered}
\] % [sic] \Ho_W(\Psi) -> \Ho_W(\tilde{\Psi}), the image is a diagram of type \tilde{\Psi} and the functor to be constructed is \mathcal{F}\longrightarrow\Ho_W(\tilde{\Psi})
Vérifions qu'il transforme flèches de $W_{\mathcal{F}}$ en isomorphismes.

\msp{page 250}

\dgrblock{dgrb587}{The deferred lemmas with functorial factorization}
\begin{dcomment}{exposition}{p249}{The deferred lemmas with functorial factorization}
Grothendieck has deliberately separated the existence of useful representatives from the assertion that they compute the whole localization. Lemmas~1 and~2 assert the latter. Their proofs are deferred, and the promise on ms.~257 does not complete them in the pages supplied here. There is a particularly short completion when the factorization axiom is available functorially. We state that extra hypothesis explicitly, since the axiom on ms.~240 does not include it.

Let $D=\mathcal M(\Psi)$, and suppose that the assignment
\[
 (A\xrightarrow p S)\longmapsto
 (A\xrightarrow{j_p}R(p)\xrightarrow{r_p}S)
\]
is a functor on the arrow category, with $j_p\in W$, $r_p\in\mathrm{Fib}_W$, and $r_pj_p=p$. For $d=(A\xrightarrow p S\xleftarrow q B)$, set
\[
 R_Dd=(R(p)\xrightarrow{r_p}S\xleftarrow{r_q}R(q)).
\]
This gives a functor $R_D:D\to\mathcal F$ and a natural transformation $j:1_D\to J R_D$, where $J:\mathcal F\hookrightarrow D$ is inclusion. Every component of $j$ is an objectwise weak equivalence. If $a:d\to d'$ is in $W(\Psi)$, naturality gives
\[
 (J R_Da)j_d=j_{d'}a.
\]
At each vertex the two maps $j_d,j_{d'}$ and the map $a$ belong to $W$; two-out-of-three therefore implies that $R_Da$ is in $W_{\mathcal F}$. The functor $R_D$ consequently descends to localization. Restricted to $\mathcal F$, the same natural transformation has both source and target in $\mathcal F$. After localization it supplies an isomorphism $1\cong \bar R_D\bar J$, while its unrestricted version supplies $1\cong\bar J\bar R_D$. These are precisely the two isomorphisms that show $\bar J$ is an equivalence. This proves Lemma~1 with the additional hypothesis. Over a fixed base $S$, the same argument applies to every full intermediate category between $\mathrm{Fib}_W(S)$ and $\mathcal M/S$: the replacement always lands in $\mathrm{Fib}_W(S)$ and the natural comparison stays in the intermediate category. It therefore proves the strengthened replacement assertion~(c) recalled on ms.~240 in this functorial setting.

For Lemma~2 use the identical construction on the three vertices of a square belonging to $\Psi$. Keep its fourth vertex $T$ unchanged, and replace its maps to $A$ and $B$ by their composites with $j_p$ and $j_q$. Their composites to $S$ still agree, so this is a strictly commutative square. It defines a functor $R_E:\mathcal M(\tilde\Psi)\to\mathcal G$ with a natural objectwise weak equivalence $1\to K R_E$, whose component at $T$ is the identity. The same two-out-of-three argument makes $R_E$ homotopical, and its restriction to $\mathcal G$ gives the other inverse comparison. Hence $\mathcal G[W_{\mathcal G}^{-1}]\to\Ho_W(\tilde\Psi)$ is an equivalence.

This argument also specifies what the arbitrary choice of a factorization would fail to provide: not the value on an object, but the functor and the natural transformation used in the proof. It uses neither an assumption that every object is fibrant nor a saturation theorem identifying $W$ with all maps inverted by localization. The required input is a functorial factorization, which is a separate sufficient hypothesis.
\end{dcomment}


\dgrblock{dgrb588}{A nonfunctorial repair by derivability}
\begin{dcomment}{later}{p249}{A nonfunctorial repair by derivability}
The passage to diagrams can be made without functorial choices, and the later theory also removes the need to add global fibrancy to the written hypotheses. The appropriate structure is a left derivable category in the sense of \citet[Definition~1.1]{Cisinski2010}. We check its one potentially hidden axiom. Let $p:X\to Y$ belong to $W\cap\mathrm{Fib}_W$, and let $u:Y'\to Y$ be arbitrary. Factor $u$ as $Y'\xrightarrow{j}Z\xrightarrow{r}Y$, where $j\in W$ and $r\in\mathrm{Fib}_W$. Since $r^*$ preserves $W$, the map $X\times_Y Z\to Z$ is in $W$. The map $X\times_Y Y'\to X\times_Y Z$ is the pullback of $j$ along a base change of $p$, hence is also in $W$. The commutative square with bottom map $j$ and two-out-of-three now give $X\times_Y Y'\to Y'\in W$. Thus trivial fibrations remain trivial under arbitrary pullback.

The other derivability axioms are immediate from the standing assumptions: a terminal object exists, weak equivalences contain isomorphisms and satisfy two-out-of-three, fibrations compose and are stable under the available pullbacks, and the required factorization is the axiom of ms.~240. These properties are stronger than the restricted pullback and factorization requirements of Definition~1.1. The precise later results are the finite-diagram structures and localization comparison of Theorem~1.30 and Corollaries~1.31--1.32, followed by the right-Kan-extension theorem, Proposition~2.6; Theorem~2.21 packages the construction as a one-sided derivator on finite direct categories.

This also settles the strengthened fixed-base condition~(c) recalled from ms.~218. The slice $\mathcal M/S$, with the inherited fibrations and weak equivalences, is left derivable and has $\mathrm{Fib}_W(S)$ as its full subcategory of fibrant objects. Every full intermediate category $\mathcal E$ containing these objects is left derivable with the same fibrant objects: the pullbacks demanded by the axioms take place among fibrant objects, and the middle object in a factorization with fibrant target is again fibrant. Proposition~1.8 identifies the localization of each such category with that of its fibrant objects. It follows that both inclusions in condition~(c) are localization equivalences. Condition~(b), the equality of ordinary and universal weak equivalences between good objects, is the relative invariance argument dual to Proposition~16; together they give~(a).

For clarity, the particular inclusions $\mathcal F\subset D$ and $\mathcal G\subset E$ can be checked by that paper's approximation theorem, Theorem~3.12. We give the verification, since $\mathcal G$ imposes no fibration condition on its fourth vertex. On $\mathcal F$ take a morphism $(a_0,a_1,a_2)$ to be a fibration when the three maps
\[
 S\to S',\qquad
 A\to S\times_{S'}A',\qquad
 B\to S\times_{S'}B'
\]
are fibrations. On $\mathcal G$ use these same conditions and additionally require the map $T\to T'$ between fourth vertices to be a fibration. Take weak equivalences objectwise in both cases. An object of $\mathcal F$ is fibrant precisely when its base $S$ is fibrant in $\mathcal M$. An object of $\mathcal G$ is fibrant precisely when both $S$ and its fourth vertex $T$ are fibrant. Its remaining vertices are then fibrant because their maps to $S$ are fibrations.

Here is the factorization needed both for these structures and for approximation. Given a map from an arbitrary cospan to a fibrant good cospan $(A'\to S'\leftarrow B')$, first factor its base map as
\[
 S\xrightarrow{\sim}S_1\twoheadrightarrow S'.
\]
Factor the resulting maps $A\to S_1\times_{S'}A'$ and $B\to S_1\times_{S'}B'$ as weak equivalences followed by fibrations. Call the middle objects $A_1,B_1$. Their maps to $S_1$ are fibrations, and the resulting cospan maps to the target by a fibration of the displayed kind. This constructs an objectwise weak equivalence from the original cospan to an object of $\mathcal F$, compatibly with the given map to the target.

For a map of squares with fibrant target in $\mathcal G$, make those three factorizations and put
\[
 P_1=A_1\times_{S_1}B_1,\qquad P'=A'\times_{S'}B'.
\]
The map $P_1\to P'$ is a fibration. Indeed it factors through $S_1\times_{S'}P'$; the first map is a composite of pullbacks of the two relative leg fibrations, and the second is a pullback of $S_1\to S'$. The source corner $T$ maps canonically to $P_1\times_{P'}T'$. Factor this map into $T\xrightarrow{\sim}T_1$ followed by a fibration. Since $P_1\times_{P'}T'\to T'$ is itself a fibration, so is $T_1\to T'$. The middle square lies in $\mathcal G$ and is fibrant: its base maps fibrantly to the fibrant target base, and $T_1$ maps fibrantly to the fibrant object $T'$. Thus the factorization has all the required properties.

The remaining derivability axioms can be checked on these same relative maps. Composition is a composite of a relative map and a pullback of the next one. A pullback of a fibration is computed vertex by vertex; at either leg its relative map is the corresponding pullback of the original relative map, so the resulting object remains good. For a trivial fibration, the relative leg maps are also weak equivalences, by two-out-of-three applied to the weak base map and the good target leg. Their pullbacks remain weak equivalences, as do the fourth-vertex maps. This verifies the trivial-fibration axiom. Factoring a diagonal of a fibrant object by the construction above supplies its path object. The ambient categories of cospans and squares carry the objectwise fibration structures of Cisinski's Corollary~1.31; the nontrivial factorization axiom there is supplied by the finite inverse diagram construction. The inclusions preserve fibrations, trivial fibrations, terminal objects and their pullbacks; they are exact and reflect weak equivalences.

Finally, Cisinski's approximation condition AP2 asks, for a map from an ambient fibrant object to an included fibrant object, for a weak-equivalence comparison with an object of the smaller category over the same target. The explicit factorizations just constructed supply it, with the preliminary weak equivalence of the ambient source taken to be its identity. AP1 is reflection of weak equivalences. Theorem~3.12 now proves both localization equivalences. This proves the two deferred lemmas under the written factorization and pullback hypotheses. The derivable-category argument controls the morphisms without either global fibrancy or functorial choices of replacements.
\end{dcomment}


\[
\begin{tikzcd}[ampersand replacement=\&, row sep=small, column sep=small]
 \& X_1\times_S X_2 \arrow[dl] \arrow[dd] \& \& X'_1\times_{S'}X'_2 \arrow[ll, "f_3"'] \arrow[dl] \arrow[dd] \\
X_1 \arrow[dd] \& \& X'_1 \arrow[ll, crossing over, "f_1"'] \& \\
 \& X_2 \arrow[dl] \& \& X'_2 \arrow[ll, "f_2"'] \arrow[dl] \\
S \& \& S' \arrow[ll, "f_0"'] \arrow[from=uu, crossing over] \&
\end{tikzcd}
\]

Il faut prouver que si $f_0,f_1,f_2\in W$, alors $f_3\in W$. Mais considérons l'image inverse du carré cartésien $Q$ par $f_0 : S'\longrightarrow S$, d'où un carré cartésien $Q_{S'}$ et une factorisation
\[
\begin{tikzcd}[ampersand replacement=\&]
Q \& Q_{S'} \arrow[l, "\tilde{f}"'] \& Q' \arrow[l, "g"'] \arrow[ll, bend left, "f"]
\end{tikzcd}
\]
de $Q'\longrightarrow Q$. Comme $f_0\in W$, et que $X_1$, $X_2$, $X_1\times_S X_2$ sont dans $\uFib_W S$, les flèches horizontales réalisant $Q_{S'}\longrightarrow Q$ sont toutes dans $W$. Ceci, plus le fait que $f_1,f_2\in W$, impliquent que dans le morphisme $g$, les flèches $g_1$, $g_2$ sont dans $W$. Il en résulte que la flèche
\[
g_3 : X'_1\times_{S'}X'_2\longrightarrow(X_1\times_S X_2)\times_S S'
\]
l'est aussi (théorie du produit dans $\uFib_W S'$ et $\Ho^{\lc}_W(S)$), donc $f_3=\tilde{f}_3g_3\in W$, q.e.d.

\msp{page 251}

\dgrblock{dgrb589}{The base is allowed to move}
\begin{dcomment}{exposition}{p250}{The base is allowed to move}
Having obtained good representatives, Grothendieck must show that their pullbacks vary homotopically in the base as well as in the two inputs. This is the extra content of the argument on ms.~250. The fixed-base product proposition cannot by itself be applied to a morphism of cospans whose middle objects differ.

Suppose a morphism of good cospans has components
\[
 a:A'\to A,\qquad b:B'\to B,\qquad s:S'\to S
\]
in $W$. Pull the target cospan back along $s$, obtaining
\[
 \tilde A=S'\times_S A,\quad
 \tilde B=S'\times_S B,\quad
 \tilde P=S'\times_S(A\times_S B).
\]
The projections $\tilde A\to A$ and $\tilde B\to B$ are weak equivalences because $A\to S$ and $B\to S$ are $W$-fibrations. The map $A\times_S B\to S$ is a composite of $W$-fibrations, so $\tilde P\to A\times_S B$ is also in $W$. The given maps factor through $\tilde A$ and $\tilde B$; two-out-of-three shows that $A'\to\tilde A$ and $B'\to\tilde B$ lie in $W$. These are now maps over the same base $S'$. The product calculation gives
\[
 A'\times_{S'}B'\longrightarrow
 \tilde A\times_{S'}\tilde B\cong\tilde P
 \quad\text{in }W.
\]
Composition proves that the original map of pullbacks is a weak equivalence. All four components of the completed square have therefore been checked. The final parenthesis of the source should refer to $\Ho_W^{\mathrm{lc}}(S')$, since this intermediate product is taken over $S'$; the printed $S$ is retained as an edition slip.

The calculation explains why one can descend the strict completion functor on $\mathcal F$. It also shows that its homotopicality is a theorem about a whole diagram, not an assertion that the pullback functor preserves arbitrary weak equivalences in each uncontrolled variable. The good-map assumptions are used at the three projections above and at the two fixed-base product maps. In Brown's setting the corresponding pullback invariance is already the content of \citet[Part~I, \S4, Lemmas~1--2, pp.~428--429]{Brown1973}.
\end{dcomment}


Notons qu'on a un isomorphisme canonique
\[
\alpha : i^*i_*\xrightarrow{\sim}\mathrm{id}_{\Ho_W(\Phi)}. \tag*{$(*)$}
\]
J'ai envie de prouver que les foncteurs
\[
\begin{tikzcd}[ampersand replacement=\&]
\Ho_W(\Phi) \arrow[r, shift right, "i_*"'] \& \Ho_W(\tilde{\Phi}) \arrow[l, shift right, "i^*"']
\end{tikzcd}
\]
\printedpage{163}
% head: A. Grothendieck : Dérivateurs, Ch. XII, Édition des Universités de Montpellier II et Paris VII
% foot: Version 30 mars 2003
sont adjoints, et que $(*)$ est une des flèches d'adjonction. Il faut définir l'autre flèche d'adjonction
\[
\beta : i_*i^*\longleftarrow\mathrm{id}_{\Ho_W(\tilde{\Phi})}. \tag*{$(**)$}
\]
Mais pour ceci, je dois utiliser le

\textbf{Lemme 2.} Soit $\mathcal{G}$ la sous-catégorie pleine de $\mathcal{M}(\tilde{\Phi})$, formée des \begin{tikzcd}[ampersand replacement=\&]
X_1 \arrow[d, "p_1"'] \& X_3 \arrow[l, "q_2"'] \arrow[d, "q_1"] \\
S \& X_2 \arrow[l, "p_2"]
\end{tikzcd} tels que $p_1,p_2\in\mathrm{Fib}_W$ et soit $W_{\mathcal{G}}$ le localiseur induit par $W(\tilde{\Phi})$ ($W$-équivalences terme à terme). Alors l'inclusion $\mathcal{G}\longrightarrow\mathcal{M}(\tilde{\Phi})$ induit une équivalence
\[
\mathcal{G}W_{\mathcal{G}}^{-1}\xrightarrow{\sim}\mathcal{M}(\tilde{\Phi})W(\tilde{\Phi})^{-1}\overset{\text{déf}}{=}\Ho_W(\tilde{\Phi}).
\]

Je vais pour le moment admettre ce

\msp{page 252}

lemme, tout comme le lemme 1 (p. 249). Je vais considérer $i^*$, $i_*$ comme des flèches entre les categories $\mathcal{F}W_{\mathcal{F}}^{-1}$, $\mathcal{G}W_{\mathcal{G}}^{-1}$. Elles sont induites par les flèches % [sic] categories -> catégories, missing accent
\[
\begin{tikzcd}[ampersand replacement=\&]
\mathcal{F} \arrow[r, shift right, "i_{0\,*}"'] \& \mathcal{G} \arrow[l, shift right, "i_0^*"']
\end{tikzcd}
\]
entre $\mathcal{F}$ et $\mathcal{G}$ eux mêmes, % [sic] eux mêmes -> eux-mêmes, missing hyphen
\[
\begin{tikzcd}[ampersand replacement=\&]
 \& X_1 \arrow[dd] \& X_3 \arrow[l] \arrow[dd] \& \& X_1 \arrow[dd] \& \\
i_0^* :\quad \mathcal{G}\quad\longrightarrow\quad\mathcal{F} : \& \& \& \longmapsto \& \& \\
 \& S \& X_2 \arrow[l] \& \& S \& X_2 \arrow[l]
\end{tikzcd}
\]
\[
\begin{tikzcd}[ampersand replacement=\&]
 \& X_1 \arrow[dd] \& \& \& X_1 \arrow[dd] \& X_1\times_S X_2 \arrow[l] \arrow[dd] \\
i_{0\,*} :\quad \mathcal{F}\quad\longrightarrow\quad\mathcal{G} : \& \& \& \longmapsto \& \& \\
 \& S \& X_2 \arrow[l] \& \& S \& X_2. \arrow[l]
\end{tikzcd}
\]
Ces deux foncteurs sont adjoints l'un de l'autre, et compatibles aux localiseurs $W_{\mathcal{F}}$, $W_{\mathcal{G}}$. Donc les foncteurs déduits de $i_0^*$, $i_{0\,*}$ par passage aux catégories de fractions sont également adjoints (cf. lemme p. 248).

Ainsi, modulo la démonstration des lemmes 1 et 2, on a prouvé ceci :

\printedpage{164}
% head: A. Grothendieck : Dérivateurs, Ch. XII, Édition des Universités de Montpellier II et Paris VII
% foot: Version 30 mars 2003

\msp{page 253}

\textbf{Théorème.} {\itshape Soit $(\mathcal{M},W)$ une catégorie de modèles ($W\subseteq\mathit{Fl}\,\mathcal{M}$) satisfaisant l'axiome de factorisation gauche (p. 240), donnant donc lieu à une bonne théorie des fibres $W$-homotopiques. Soient
\[
\begin{tikzcd}[ampersand replacement=\&, row sep=small, column sep=small]
 \& 1 \& \& \& \& 1 \arrow[dr] \& \\
\Psi=0 \arrow[ur] \arrow[dr] \& \& {} \arrow[r, "i"] \& {} \& \tilde{\Psi}=0 \arrow[ur] \arrow[dr] \& \& 3 \\
 \& 2 \& \& \& \& 2 \arrow[ur] \&
\end{tikzcd}
\]
(immersions ouvertes dans} $\mathrm{Cat}${\itshape ), et considérons le foncteur restriction
\[
i^* : \Ho_W(\tilde{\Psi})\longrightarrow\Ho_W(\Psi).
\]
Ce foncteur admet un adjoint à droite
\[
i_* : \Ho_W(\Psi)\longrightarrow\Ho_W(\tilde{\Psi}).
\]
Ce foncteur $i_*$ se `calcule' ainsi : sa valeur en \begin{tikzcd}[ampersand replacement=\&]
X_1 \arrow[d, "p_1"'] \& \\
S \& X_2 \arrow[l, "p_2"]
\end{tikzcd} est canoniquement isomorphe (dans} $\Ho_W(\tilde{\Psi})${\itshape ) à \begin{tikzcd}[ampersand replacement=\&]
X_1 \arrow[d] \& X_1\times_S X_2 \arrow[l] \arrow[d] \\
S \& X_2 \arrow[l]
\end{tikzcd}, si on suppose que $p_1$ ou $p_2$ est une $W$-fibration. Dans le cas général il se calcule en factorisant $p_1$ en
\[
X_1\xrightarrow{i_1}\bar{X}_1\xrightarrow{p'_1}S,\qquad i_1\in W,\ p'_1\in\mathit{Fib}_W,
\]
et en prenant \begin{tikzcd}[ampersand replacement=\&]
\bar{X}_1 \arrow[d] \& \bar{X}_1\times_S X_2 \arrow[l] \arrow[d] \\
S \& X_2 \arrow[l]
\end{tikzcd}, ou en

\msp{page 254}

procédant de façon symétrique, ou en factorisant simultanement $p_1$ et $p_2$, et prenant % [sic] simultanement -> simultanément, missing accent
\[
\begin{tikzcd}[ampersand replacement=\&]
\bar{X}_1 \arrow[d] \& \bar{X}_1\times_S\bar{X}_2 \arrow[l] \arrow[d] \\
S \& \bar{X}_2. \arrow[l]
\end{tikzcd}
\]}

\dgrblock{dgrb590}{The missing corner can be an entire group}
\begin{dcomment}{examples}{p254:procédant de façon symétrique}{The missing corner can be an entire group}
A groupoid calculation makes the difference between a strict pullback and its homotopy replacement visible without any topological machinery. Work in groupoids with equivalences as $W$, and let $G$ be a group. Write $BG$ for the one-object groupoid with automorphism group $G$. Consider the cospan
\[
 *\longrightarrow BG\longleftarrow *.
\]
Its strict pullback is the terminal groupoid: there is one pair of objects, and the only pair of morphisms is the pair of identities. The homotopy pullback will instead be the discrete groupoid on the set $G$.

To compute it, let $EG$ be the groupoid whose objects are the elements of $G$ and with exactly one morphism from each object to every other object. Define $p:EG\to BG$ by sending the unique morphism $g\to h$ to $hg^{-1}$. The composition formula
\[
 (kh^{-1})(hg^{-1})=kg^{-1}
\]
shows that this is a functor. The inclusion $*\to EG$ selecting the identity element is fully faithful and essentially surjective, hence is an equivalence; its composite with $p$ is the specified map $*\to BG$.

The functor $p$ is an isofibration. Given an object $g$ and an arrow $a$ in $BG$, the arrow $g\to ag$ lies above $a$. Isofibrations of groupoids have the required $W$-fibration property. For completeness, pull back an equivalence $u:C\to D$ over a base along an isofibration $E\to S$. The induced functor on pullbacks is fully faithful because $u$ is. Given an object $(d,e)$ of the target pullback, choose an isomorphism $u(c)\cong d$; its image in $S$ lifts through the isofibration, and the lift changes $e$ to an object lying strictly over the image of $c$. This produces an object in the source pullback isomorphic to $(d,e)$, proving essential surjectivity. The same reasoning applies after any base change.

We can therefore replace the left leg by $*\xrightarrow{\sim}EG\xrightarrow p BG$. In the strict pullback $EG\times_{BG}*$, an object is an element $g\in G$. A morphism from $g$ to $h$ exists precisely when its image $hg^{-1}$ in $BG$ is the identity, which is equivalent to $g=h$. Thus
\[
 EG\times_{BG}*\cong G_{\mathrm{disc}}.
\]
For $G=\mathbb Z/2$, the strict pullback has one connected component while the homotopy pullback has two. For a general group, its elements record the possible isomorphisms between the two images of the chosen object of $*$. The derived completion preserves those comparisons.

This is the cospan version of the groupoid computations already used in the commentary on Chapter~VI, manuscript p.~107. Its purpose here is narrower: it exhibits the actual fourth vertex produced by the theorem and explains why the maps $*\to BG$ must be replaced. The whole completed square, including its two projections and its map to the original cospan after localization, is the value of $i_*$.
\end{dcomment}


% ---- from Der-ChXIIc-30-03-03_p051-055.tex
\printedpage{165}
% head: A. Grothendieck : Dérivateurs, Ch. XII, Édition des Universités de Montpellier II et Paris VII
% foot: Version 30 mars 2003

\textit{La loi fonctorielle s'obtient à l'aide du lemme 1 p. 249, qui permet de se borner aux arguments $\begin{tikzcd}[ampersand replacement=\&] X_1 \arrow[d, "p_1"'] \& \\ S \& X_2 \arrow[l, "p_2"] \end{tikzcd}$ tels que $p_1,p_2\in \mathit{Fib}_W$, de sorte que l'on trouve un foncteur}
\[
\begin{tikzcd}[ampersand replacement=\&]
X_1 \arrow[d] \& \\
S \& X_2 \arrow[l]
\end{tikzcd}
\qquad\longmapsto\qquad
\begin{tikzcd}[ampersand replacement=\&]
X_1 \arrow[d] \& X_1\times_S X_2 \arrow[l] \arrow[d] \\
S \& X_2, \arrow[l]
\end{tikzcd}
\]
\textit{compatible aux localiseurs, et passant aux catégories de fractions} $\Ho_W(\Psi)$ \textit{et} $\Ho_W(\tilde{\Psi})$ \textit{respectivement.}

\textbf{NB} En procédant de façon similaire, plus simple (sans besoin d'un recours au lemme 1), on trouve l'existence d'un \textit{adjoint à gauche} $i_!$ de $i^* : \Ho_W(\Psi)\longrightarrow \Ho_W(\tilde{\Psi})$, sous réserve de l'existence d'un objet initial $\varphi$ dans $\mathcal{M}$. Cet foncteur $i_!$ est donné par % [sic] Cet foncteur -> Ce foncteur, "foncteur" begins with a consonant

\msp{page 255}
\[
i_!\left(
\begin{tikzcd}[ampersand replacement=\&]
X_1 \arrow[d] \& \\
S \& X_2 \arrow[l]
\end{tikzcd}
\right)
\qquad = \qquad
\left(
\begin{tikzcd}[ampersand replacement=\&]
X_1 \arrow[d] \& \varphi \arrow[l] \arrow[d] \\
S \& X_2 \arrow[l]
\end{tikzcd}
\right),
\]
la loi fonctorielle étant évidente.

\dgrblock{dgrb591}{The adjunction determines the completion}
\begin{dcomment}{exposition}{p254}{The adjunction determines the completion}
The adjunction is what makes Grothendieck's construction a canonical completion of a cospan. On the good-diagram categories, restriction $r:\mathcal G\to\mathcal F$ has right adjoint $c$, where $c$ supplies the strict pullback. A morphism from a square $Q$ to $c(d)$ is uniquely determined by its three components on the cospan $r(Q)$: the fourth component exists uniquely by the pullback universal property. This gives the adjunction bijection before localization.

The counit $rc\to1_{\mathcal F}$ is the identity on the original three vertices. The unit at a square with fourth vertex $T$ is the map to its completion, whose fourth component is
\[
 T\longrightarrow A\times_S B.
\]
At a completed square this map is the identity, and restricting the unit to its three original vertices also gives the identity. These observations verify both triangle identities directly. Restriction preserves objectwise weak equivalences; the preceding base-change calculation proves that completion does too. The localization lemma on ms.~248 therefore descends the adjunction. Once Lemmas~1 and~2 identify the localized good-diagram categories with the required diagram categories, it is exactly $i^*\dashv i_*$, with invertible counit $i^*i_*\to1$.

The formula using only one fibrant leg follows by an explicit comparison, not by a symmetry convention. If $A\to S$ is already a $W$-fibration, replace only $B$ by $\bar B$. The map
\[
 A\times_S B\longrightarrow A\times_S\bar B
\]
is in $W$, as in the block at ms.~242, and forms an objectwise weak equivalence between the completed square and one with two good legs. The latter computes $i_*$. For an arbitrary cospan, first replace $A$; the same argument compares the one-leg and two-leg formulas. Interchanging $A$ and $B$ gives the other formula. These comparisons concern objects of $\Ho_W(\tilde\Psi)$ and retain all the projection maps.

There is also the left adjoint mentioned in the concluding NB. If $\mathcal M$ has an initial object $\varphi$, extending a cospan by putting $\varphi$ at vertex~$3$ is left adjoint to restriction already on ordinary diagram categories. The fourth component of any extension of a map of cospans is the unique map $\varphi\to T$. This extension preserves objectwise weak equivalences, since its fourth component on every morphism is $1_\varphi$, so the adjunction descends by the same lemma. It needs no fibrant replacement. The displayed type of $i^*$ in that NB is reversed in the edition; its correct type remains $\Ho_W(\tilde\Psi)\to\Ho_W(\Psi)$. Likewise the $\Phi,\tilde\Phi$ used on ms.~251--252 denote the $\Psi,\tilde\Psi$ of the theorem, not a second pair of diagram shapes.
\end{dcomment}


\dgrblock{dgrb592}{Two extensions of the same cospan}
\begin{dcomment}{exposition}{p255:la loi fonctorielle étant évidente.}{Two extensions of the same cospan}
Grothendieck has now reached the point at which a homotopy fibre can be characterized by its relation to all diagrams, rather than by a particular construction of paths. The two extensions of the three-vertex diagram make this distinction visible. Restriction has the direction
\[
 i^*: \Ho_W(\tilde\Psi)\longrightarrow\Ho_W(\Psi).
\]
The reversed direction in the preceding NB is an edition slip. Its left adjoint, discussed in the preceding block, fills the fourth vertex with an initial object; its right adjoint fills that vertex with a homotopy pullback. These operations answer different universal problems on the same cospan.

For the right adjoint, the fourth vertex must instead receive a compatible pair of maps. Once the two legs to $S$ are $W$-fibrations, that universal object is $X_1\times_S X_2$. The extra work in Lemmas 1 and 2 is what allows this calculation on suitably replaced diagrams to determine a functor and an adjunction on the original localization. The homotopy limit of the cospan is then obtained by evaluating this completed square at vertex 3.

The indexing convention explains the apparently reversed endpoint in the following lemma. An object of $\mathcal M(I)$ is a functor $I^{\mathrm{op}}\to\mathcal M$. If $t$ is terminal in $I$, each arrow $i\to t$ induces $F(t)\to F(i)$. A natural transformation from the constant diagram on $A$ to $F$ is consequently determined by its component $A\to F(t)$, and every such component gives a natural transformation. Thus
\[
 \operatorname{Hom}_{\mathcal M(I)}(p^*A,F)
 \cong \operatorname{Hom}_{\mathcal M}(A,F(t)).
\]
Evaluation at $t$ is a right adjoint. The displayed formula survives localization because both functors preserve the designated weak equivalences. In particular, the object constructed here is universal for maps of coherent cospans. There is no assertion that taking an ordinary pullback in $\mathcal M[W^{-1}]$ would recover it.
\end{dcomment}


Ce n'est pas encore la fin du sorite! Il faut encore interpréter le foncteur composé
\[
\Ho_W(\Psi)\xrightarrow{\ i_*\ }\Ho_W(\tilde{\Psi})\xrightarrow{\ \varepsilon_3^*\ }\Ho_W(e)=\Ho_W,
\]
où
\[
\varepsilon_3 : e\xrightarrow{\ \sim\ }\{3\}
\qquad\longrightarrow\qquad
\begin{tikzcd}[ampersand replacement=\&, column sep=small, row sep=small]
\& 1 \arrow[dl] \& \\
3 \& \& 0 \arrow[ul] \arrow[dl] \\
\& 2 \arrow[ul] \&
\end{tikzcd}
\]
est l'inclusion du sommet 3 dans $\tilde{\Psi}$. Si $\Ho_W(\Psi)$ est interprété, conformément au lemme 1, comme $\mathcal{F}W_{\mathcal{F}}^{-1}$, ce foncteur composé n'est autre que
\[
\begin{tikzcd}[ampersand replacement=\&]
X_1 \arrow[d] \& \\
S \& X_2 \arrow[l]
\end{tikzcd}
\qquad\longmapsto\qquad X_1\times_S X_2
\]
\printedpage{166}
% head: A. Grothendieck : Dérivateurs, Ch. XII, Édition des Universités de Montpellier II et Paris VII
% foot: Version 30 mars 2003
(pour $\left(\begin{tikzcd}[ampersand replacement=\&] X_1 \arrow[d, "p_1"'] \& \\ S \& X_2 \arrow[l, "p_2"] \end{tikzcd}\right)\in \mathrm{Ob}\,\mathcal{F}$, i.e. $p_1,p_2\in\Fib_W$.)

On veut interpréter ce foncteur comme un $\varprojlim$ $W$-homotopique, i.e. vu comme un `produit fibré $W$-homotopique'. De façon précise, on a ceci :

\msp{page 256}
\textbf{Corollaire 1.} \textit{Considérons $p_\Psi : \Psi\longrightarrow e$, et le foncteur}
\[
p_\Psi^* : \Ho_W(e)=\Ho_W\longrightarrow\Ho_W(\Psi),
\]
\textit{déduit par passage aux catégories de fractions du foncteur}
\[
\begin{array}{ccc}
\mathcal{M} & \longrightarrow & \mathcal{M}(\Psi) \\
X & \longrightarrow & \left(
\begin{tikzcd}[ampersand replacement=\&]
X \arrow[d] \& \\
X \& X \arrow[l]
\end{tikzcd}
\right).
\end{array}
\]
\textit{Ce foncteur admet le foncteur $\varepsilon_3^*i_*$ précédent comme adjoint à droite, i.e.}
\[
p_{\Psi\,*}=\varepsilon_3^*i_* : \Ho_W(\Psi)\longrightarrow\Ho_W.
\]

\textsc{Démonstration}. On peut regarder $p_\psi$ comme le composé % [sic] p_\psi -> p_\Psi, the subscript is the capital \Psi everywhere else
\[
\Psi\xrightarrow{\ i\ }\tilde{\Psi}\xrightarrow{\ p_{\tilde{\Psi}}\ }e,
\]
et on a donc
\[
p_\Psi^*\simeq i^*(p_{\tilde{\Psi}})^*.
\]
Donc l'existence et la détermination d'un adjoint à droite pour $p_\Psi^*$ résultera de l'existence et de la détermination d'adjoints à droite $i_*$, $(p_{\tilde{\Psi}})_*$ de $i^*$, $(p_{\tilde{\Psi}})^*$, via
\[
(p_\Psi)_*\simeq(p_{\tilde{\Psi}})_*i_*.
\]

\msp{page 257}
Donc le corollaire 1 résulte de la formule
\[
p_{\tilde{\Psi}\,*}=\varepsilon_3^*,
\]
qui elle-même est un cas particulier du

\printedpage{167}
% head: A. Grothendieck : Dérivateurs, Ch. XII, Édition des Universités de Montpellier II et Paris VII
% foot: Version 30 mars 2003

\textbf{Lemme.} \textit{Soit $I$ dans} Cat\textit{, avec objet final $e$, considérons l'inclusion}
\[
\varepsilon : e\longrightarrow I,\qquad \textit{et } p : I\longrightarrow e.
\]
\textit{Soit $(\mathcal{M},W)$ une catégorie de modèles, et considérons}
\[
\Ho_W(I)\xleftarrow{\ p^*\ }\Ho_W.
\]
\textit{Ce foncteur admet un adjoint à gauche} \texttt{[plutôt à droite]} \textit{$p_*$, lequel est isomorphe à $\varepsilon^*$.}

\textsc{Démonstration}. Considérons le couple de foncteurs adjoints
\[
\begin{tikzcd}[ampersand replacement=\&, column sep=large]
\mathcal{M}(I) \arrow[r, shift right, "p_*=\varprojlim\limits_I"'] \& \mathcal{M} \arrow[l, shift right, "p^*"']
\end{tikzcd}\ ;
\]
comme $I$ admet l'objet final $e$, le foncteur $p_*$ est isomorphe à $\varepsilon^*$. Ces foncteurs $p^*$, $\varepsilon^*$ sont compatibles aux localiseurs $W(I)$, $W$, donc on conclut par le lemme page 248.

Il ne reste plus qu'à prouver les lemmes 1 et 2. Je le ferai demain.

\msp{page 258}
Aujourd'hui, j'ai envie encore de donner les sorites des carrés $W$-cartésiens dans une

\textbf{Proposition 19 :} \textit{Soit $\begin{tikzcd}[ampersand replacement=\&] X_1 \arrow[d, "p_1"'] \& X_3 \arrow[l, "f_1"'] \arrow[d, "f_2"] \\ S \& X_2 \arrow[l, "p_2"] \end{tikzcd}$ un carré commutatif dans $\mathcal{M}$. Considérons des factorisations de $p_1$, $p_2$ en}
\[
X_1\xrightarrow{\ i_1\ }\bar{X}_1\xrightarrow{\ \bar{p}_1\ }S,\qquad X_2\xrightarrow{\ i_2\ }\bar{X}_2\xrightarrow{\ \bar{p}_2\ }S,
\]
\textit{$i_1,i_2\in W$, $\bar{p}_1,\bar{p}_2\in F_W$} \texttt{[$=$ Fib$_W$]}\textit{, et considérons les flèches induites par $X_3\longrightarrow X_1\times_S X_2$ (déduite de $Q$) via $\begin{tikzcd}[ampersand replacement=\&, row sep=small] \& \bar{X}_1\times_S X_2 \\ X_1\times_S X_2 \arrow[ur] \arrow[r] \arrow[dr] \& \bar{X}_1\times_S\bar{X}_2 \\ \& X_1\times_S\bar{X}_2 \end{tikzcd}$, de sorte qu'on trouve un diagramme commutatif}
\[
\begin{tikzcd}[ampersand replacement=\&, row sep=small]
\& \bar{X}_1\times_S X_2 \arrow[d, "j_2"'] \\
X_3 \arrow[ur, "\pi_1"] \arrow[r, "\pi_0"] \arrow[dr, "\pi_2"] \& \bar{X}_1\times_S\bar{X}_2 \\
\& X_1\times_S\bar{X}_2, \arrow[u, "j_1"]
\end{tikzcd}
\]
\textit{où $j_1=i_1\times_S\mathrm{id}_{\bar{X}_2}$ et $j_2=\mathrm{id}_{\bar{X}_1}\times i_2$ sont dans $W$ (car $i_1$ et $i_2$ le sont, et $\bar{X}_1$ et $\bar{X}_2$ sont dans $\underline{\mathrm{Fib}}_W S$). On a donc}
\[
(\pi_1\in W)\qquad\iff\qquad(\pi_2\in W)\qquad\qquad(\iff\ \pi_0\in W),
\]
\textit{ce qui prouve que la condition $\pi_1\in W$, qui semble dépendre du choix de la factorisation par $i_1$, n'en dépend pas, et idem pour $\pi_2\in W$, et pour $\pi_0\in W$ (qui semble dépendre à la fois de $i_1$, et de $i_2$).}

\dgrblock{dgrb593}{Independence of replacements and invariance of the test}
\begin{dcomment}{exposition}{p258:ce qui prouve que la condition}{Independence of replacements and invariance of the test}
Grothendieck wants the property of a square to be independent of the auxiliary factorization used to recognize it. Proposition 19 achieves this with a particularly economical comparison: one changes the first leg while holding the second fixed, then exchanges their roles. We spell out the argument because it also supplies the invariance used in the next pasting proof.

Write a square as $D\to A$, $D\to B$ over the cospan $A\to S\leftarrow B$. Choose $A\xrightarrow{a}\bar A\to S$ and $B\xrightarrow{b}\bar B\to S$ as in the proposition. The two arrows
\[
 \bar A\times_S B\longrightarrow\bar A\times_S\bar B,
 \qquad
 A\times_S\bar B\longrightarrow\bar A\times_S\bar B
\]
are in $W$, since each pulls back a member of $W$ along a $W$-fibration. Thus the three comparison maps from $D$ are simultaneously in $W$. Now choose a second factorization of $A\to S$, retaining the chosen factorization of $B\to S$. Both first-leg tests are equivalent to the very same test $D\to A\times_S\bar B$. They therefore agree. The argument with $A$ and $B$ interchanged proves independence of the other choice. No map between the two factorizations of $A$ has been assumed.

For the pasting argument we also need the following verification. Suppose a map of squares $Q\to Q'$ is in $W$ at all four vertices. Denote the cospans by $A\to S\leftarrow B$ and $A'\to S'\leftarrow B'$, and choose a factorization
\[
 A'\xrightarrow{a'}\widehat A'\xrightarrow{q'}S',
 \qquad a'\in W,\quad q'\in\Fib_W.
\]
Set $T=S\times_{S'}\widehat A'$. The projection $T\to\widehat A'$ is in $W$, being the pullback of $S\to S'$ along $q'$. The induced $A\to T$ is in $W$ by two-out-of-three, since its composite to $\widehat A'$ is $A\to A'\xrightarrow{a'}\widehat A'$. Moreover $T\to S$ is a $W$-fibration. This is an admissible replacement of the first leg of $Q$.

With these choices the comparison of the two pullback objects is
\[
 T\times_S B\cong\widehat A'\times_{S'}B
 \longrightarrow\widehat A'\times_{S'}B'.
\]
It is in $W$, since it is obtained by pulling back $B\to B'$ along the pullback of $q'$. The square relating these objects to $D\to D'$ commutes. Its vertical arrows are in $W$, so two-out-of-three shows that the comparison for $Q$ is in $W$ if and only if that for $Q'$ is. This proves invariance under objectwise weak equivalences, and hence under zigzags of them, using only the factorization hypothesis and the intrinsic definition of $\Fib_W$. Strong saturation, functorial choices of factorization, and the deferred adjunction lemmas are unnecessary for this particular assertion.
\end{dcomment}


\printedpage{168}
% head: A. Grothendieck : Dérivateurs, Ch. XII, Édition des Universités de Montpellier II et Paris VII
% foot: Version 30 mars 2003

\msp{page 259}
\textbf{Définition.} \textit{Quand ces conditions sont satisfaites, on dit que le carré $Q$ est $W$-cartésien.}

\textbf{Corollaire 1.} \textit{Si $Q$ est $W$-cartésien, alors l'image de $Q$ dans} $\Ho_W(\tilde{\Psi})$ \textit{est dans l'image essentielle de}
\[
i_* : \Ho_W(\Psi)\longrightarrow\Ho_W(\tilde{\Psi}),
\]
\textit{i.e. la flèche canonique dans} $\Ho_W(\tilde{\Psi})$
\[
Q\longrightarrow i_*i^*Q
\]
\textit{est} \texttt{[un]} \textit{isomorphisme. La réciproque est vraie si $W$ est fortement saturé.}

\dgrblock{dgrb594}{The unit and the role of strong saturation}
\begin{dcomment}{exposition}{p259:La réciproque est vraie si}{The unit and the role of strong saturation}
Grothendieck now compares his concrete test with the intrinsic language of homotopy Kan extension. This comparison separates two tasks which coincide in the familiar examples: constructing an isomorphism after localization, and recognizing that a specified original arrow belongs to $W$.

Let $Q^r$ be obtained by replacing $X_1\to S$ and $X_2\to S$ by $\bar X_1\to S$ and $\bar X_2\to S$, while keeping $X_3$ and $S$. There is an objectwise weak equivalence $Q\to Q^r$. The unit on $Q^r$ is represented by the map to the completed pullback square which is the identity on $\bar X_1$, $\bar X_2$, and $S$, and is
\[
 \pi_0:X_3\longrightarrow\bar X_1\times_S\bar X_2
\]
at the fourth vertex. If $\pi_0\in W$, this map of squares becomes an isomorphism. Naturality of the unit then gives the assertion for $Q$.

Conversely, if the unit on $Q$ is invertible, the unit on $Q^r$ is invertible. Evaluation at vertex 3 shows that $\pi_0$ is inverted by $\mathcal M\to\mathcal M[W^{-1}]$. Strong saturation says exactly that this forces $\pi_0\in W$. Two-out-of-three by itself makes no such assertion about every map inverted by localization. This is the reason for the qualification in Corollary 1, rather than an incidental technical precaution.

Since $i^*i_*\cong\mathrm{id}$, the right adjoint is fully faithful, and the usual adjunction argument identifies its essential image with the squares whose unit is invertible. In modern derivator terminology this essential-image condition defines cartesian squares; the concrete $W$ test agrees with it under the saturation hypothesis just explained. See \citet[Proposition 1.20 and Lemma 1.21]{Groth2013} for the fully faithful Kan-extension formulation. That formulation presupposes the construction of the adjunction, so it does not by itself discharge the two deferred localization lemmas in this manuscript.
\end{dcomment}


\textbf{Corollaire 2.} \textit{Considérons un diagramme commutatif dans $\mathcal{M}$}
\[
\begin{tikzcd}[ampersand replacement=\&]
X_0' \arrow[d] \& X_1' \arrow[l] \arrow[d] \arrow[dl, phantom, "\textit{1}"] \& X_2' \arrow[l] \arrow[d] \arrow[dl, phantom, "\textit{2}"] \\
X_0 \& X_1 \arrow[l] \& X_2, \arrow[l]
\end{tikzcd}
\tag*{(D)}
\]
\textit{d'où deux carrés commutatifs 1, 2, et un carré composé}
\[
\begin{tikzcd}[ampersand replacement=\&]
X_0' \arrow[d] \& X_2' \arrow[l] \arrow[d] \arrow[dl, phantom, "\textit{3}"] \\
X_0 \& X_2. \arrow[l]
\end{tikzcd}
\]
\textit{Supposons le carré 1 $W$-cartésien. Alors le carré 2 est $W$-cartésien si et seulement si le carré 3 l'est.}

\textsc{Démonstration}. Considérons une factorisation de $\begin{tikzcd}[ampersand replacement=\&, row sep=small] X_0' \arrow[d] \\ X_0 \end{tikzcd}$ en
\[
\begin{tikzcd}[ampersand replacement=\&, column sep=small, row sep=small]
\bar{X}_0' \arrow[dr] \& X_0' \arrow[l] \arrow[d] \\
\& X_0
\end{tikzcd}
\]
avec $(X_0'\longrightarrow\bar{X}_0')\in W$, et $(\bar{X}_0'\longrightarrow X_0)\in\Fib_W$. On peut alors remplacer dans (D) $X_0'$ par $\bar{X}_0'$, obtenant un diagramme $(\bar{D})$, et on voit aussitôt
\printedpage{169}
% head: A. Grothendieck : Dérivateurs, Ch. XII, Édition des Universités de Montpellier II et Paris VII
% foot: Version 30 mars 2003
\msp{page 260}
que chacun des carrés 1(D), 2(D), 3(D) est $W$-cartésien si et seulement si le carré correspondant $1(\bar{D})$, $2(\bar{D})$ ou $3(\bar{D})$ l'est. Cela montre qu'on peut supposer $X_0\longrightarrow S$ dans $\Fib_W S$. Considérons alors le diagramme commutatif déduit de (D)
\[
\begin{tikzcd}[ampersand replacement=\&, column sep=small]
X_0' \arrow[d] \& \tilde{X}_1'\overset{\text{déf}}{=}X_0'\times_{X_0}X_1 \arrow[l] \arrow[d] \& X_1' \arrow[l, "\alpha"'] \& X_2' \arrow[l] \arrow[d] \\
X_0 \& X_1 \arrow[l] \& \& X_2, \arrow[ll]
\end{tikzcd}
\]
où $\alpha\in W$ en vertu de l'hypothèse $W$-cartésienne sur le carré 1. On voit encore que le carré 2 est $W$-cartésien si et seulement si le carré correspondant avec $X_1'$ remplacé par $\tilde{X}_1'$ l'est, quant aux carrés 1, 3, ils ne changent pas par cette modification. Cela montre qu'on peut supposer le carré 1 \textit{cartésien.} Considérons alors le diagramme
\[
\begin{tikzcd}[ampersand replacement=\&]
X_0' \arrow[d, "\mathrm{Fib}_W"'] \& X_1' \arrow[l] \arrow[d, "\mathrm{Fib}_W"] \arrow[dl, phantom, "\text{cart.}"] \& \tilde{X}_2'\overset{\text{déf}}{=}X_1'\times_{X_1}X_2 \arrow[l] \arrow[d, "\mathrm{Fib}_W"] \arrow[dl, phantom, "\text{cart.}"] \& X_2' \arrow[l, "\beta"'] \\
X_0 \& X_1 \arrow[l] \& X_2. \arrow[l]
\end{tikzcd}
\]
Alors la proposition 19 et la définition montrent que le carré 2 est $W$-cartésien si et seulenment si $\beta\in W$, et de même pour le carré 3. Ce qui achève la démonstration. % [sic] seulenment -> seulement, misspelling

\dgrblock{dgrb595}{Pasting reduces two questions to the same arrow}
\begin{dcomment}{exposition}{p260:Ce qui achève la démonstration.}{Pasting reduces two questions to the same arrow}
Grothendieck's proof makes pasting effective by arranging that the first square is an actual pullback along a $W$-fibration. The apparently immediate replacements in the proof are justified by the objectwise invariance established in our preceding block. The printed phrase $X_0\to S\in\Fib_W S$ should be read as $X'_0\to X_0\in\Fib_W$; it is the latter arrow that has just been factored.

Write $E\to X_0$ for a fibrant replacement of $X'_0\to X_0$. Replacing $X'_0$ by $E$ changes neither relevant cartesianness test. Because the first square is $W$-cartesian, its comparison
\[
 X'_1\longrightarrow E\times_{X_0}X_1
\]
is in $W$. We may therefore replace $X'_1$ by $E_1=E\times_{X_0}X_1$, preserving the test for the second square. The first square is now strictly cartesian, and $E_1\to X_1$ is a $W$-fibration.

Both remaining conditions are consequently tested by the one arrow
\[
 X'_2\longrightarrow E_1\times_{X_1}X_2
       \cong E\times_{X_0}X_2.
\]
Using the expression on the left tests the second square; using the expression on the right tests the outer rectangle. The two conditions are identical. This proves both the composition implication and its indicated cancellation, without assuming that arbitrary base change preserves $W$.

The result became a basic part of the calculus of derivators \citep[Proposition 3.13]{Groth2013}. Its orientation matters: in the manuscript the horizontal arrows point towards the left, so its left-hand square is the square adjacent to the common target. One cannot replace the hypothesis by either small square indiscriminately in an unstable theory.
\end{dcomment}


\textbf{Corollaire 2.} \textit{Soit $\begin{tikzcd}[ampersand replacement=\&] X_1 \arrow[d, "p_1"'] \& X_3 \arrow[l, "f_2"'] \arrow[d, "f_1"] \\ S \& X_2 \arrow[l, "p_2"] \end{tikzcd}$ un carré $W$-cartésien. Alors}
\[
\begin{aligned}
p_1\in W\quad&\Longrightarrow\quad f_1\in W\\
p_2\in W\quad&\Longrightarrow\quad f_2\in W,
\end{aligned}
\]
\textit{plus précisément, si $p_i\in W$, alors $Q$ est $W$-cartésien si et seulement si $f_i\in W$.}

\msp{page 261}
Supposons $p_1\in W$. Alors dans la factorisation de $p_1$ en $X_1\longrightarrow\bar{X}_1\longrightarrow S$, on peut prendre $\bar{X}_1=S$. Donc le carré est $W$-cartésien si et seulement si
\[
(X_3\longrightarrow\bar{X}_1\times_S X_2=X_2)\in W,
\]
i.e. $f_1\in W$, q.e.d.

Bien sûr, on ne peut en général affirmer les implications inverses : $f_i\in W$ n'implique nullement $p_i\in W$.

\dgrblock{dgrb596}{A nonempty counterexample to the converse}
\begin{dcomment}{examples}{p261:Bien sûr, on ne peut en général affirmer}{A nonempty counterexample to the converse}
Grothendieck's final caution about cartesian squares concerns a loss of information under restriction to part of the base. A finite example already shows the issue, without empty objects or any higher homotopy calculation.

Take $\mathcal M=\mathrm{Set}$ and let $W$ be the bijections. Every map is an intrinsic $W$-fibration, and the identity followed by the given map provides the required factorization. Put
\[
 S=\{0,1\},\quad X_1=\{a,b,c\},\quad X_2=\{0\},
 \qquad p_1(a)=0,\quad p_1(b)=p_1(c)=1,
\]
and let $p_2$ be the inclusion of $\{0\}$ in $S$. The pullback is $X_3=\{a\}$. Its projection $f_1:X_3\to X_2$ is a bijection, whereas $p_1:X_1\to S$ is not. The square is strictly cartesian and hence $W$-cartesian.

Only the fibre over 0 has been tested. Nothing in the upper projection records the two elements over 1. The forward assertion in the source has no such problem: if $p_1\in W$, use the factorization $X_1\xrightarrow{p_1}S\xrightarrow{\mathrm{id}}S$, and the cartesianness comparison is precisely $f_1$. The later appendix will ask for a converse to the analogous pushout statement. That converse is an additional conservativity condition, rather than a formal consequence of the existence of homotopy pushouts.
\end{dcomment}


% ---- from Der-ChXIIc-30-03-03_p056-060.tex
\printedpage{170}
% head: A. Grothendieck : Dérivateurs, Ch. XII, Édition des Universités de Montpellier II et Paris VII

Tous ces résultats, et la notion de carré $W$-cartésien, se dualisent (moyennant l'axiome idoine de factorisation à droite, cf. page 240) en une théorie des sommes amalgamées $W$-homotopiques, et des carrés $W$-cocartésiens. Cette fois on travaille avec la catégorie
\[
\Phi \;=\; \left(
\begin{tikzcd}[ampersand replacement=\&, row sep=small, column sep=small]
 \& 1 \arrow[dl] \\
0 \& \\
 \& 2 \arrow[ul]
\end{tikzcd}
\right)
\qquad\quad \overset{j}{\hookrightarrow} \qquad\quad
\tilde{\Phi} \;=\; \left(
\begin{tikzcd}[ampersand replacement=\&, row sep=small, column sep=small]
 \& 1 \arrow[dl] \& \\
0 \& \& 3 \arrow[ul] \arrow[dl] \\
 \& 2 \arrow[ul] \&
\end{tikzcd}
\right),
\]
$j$ étant une immersion fermée. On a un foncteur restriction
\[
j^* : \Ho_W(\tilde{\Phi}) \longrightarrow \Ho_W(\Phi),
\]
dont on construit un adjoint à gauche
\[
j_! : \Ho_W(\Phi) \longrightarrow \Ho_W(\tilde{\Phi}).
\]
Et on trouve une notion naturelle de
\msp{page 262}
carré $W$-cocartésien, qui paraît irréprochable à tout points de vue, et d'une simplicité parfaite. % [sic] tout points -> tous points, agreement of "tout"

Cela amène à revoir la notion en question dans Cat, où elle avait l'air très astucieuse, profonde et un peu obscure, avec des intégrales et tout ça. C'est le moment de se rendre compte que `ce n'est que ça'. Donc de prouver la

\textbf{Proposition 20.} \textit{Supposons} $\mathcal{M} = \mathrm{Cat}$\textit{, $W$ un localiseur fondamental satisfaisant W(7 bis). Considérons un carré commutatif dans} Cat
\[
Q = \left(
\begin{tikzcd}[ampersand replacement=\&]
S \arrow[r] \arrow[d] \& X_1 \arrow[d] \\
X_2 \arrow[r] \& X_3
\end{tikzcd}
\right).
\]
\textit{Pour qu'il soit $W$-cartésien} \texttt{[plutôt} $W$\texttt{-cocartésien]} \textit{au sens de la théorie développée ici (pour des catégories de modèles générales), il faut et il suffit que}
\[
\varphi_Q : \quad \mathcal{X}_Q = \left(
\begin{tikzcd}[ampersand replacement=\&, row sep=small, column sep=small]
 \& X_1 \arrow[dl] \\
\int S \& \\
 \& X_2 \arrow[ul]
\end{tikzcd}
\right) \qquad\quad \longrightarrow \qquad\quad X_3
\]

\texttt{[plutôt}
\[
\varphi_Q : \quad \mathcal{X}_Q = \left(
\begin{tikzcd}[ampersand replacement=\&, row sep=small, column sep=small]
 \& X_1 \\
\int S \arrow[ur] \arrow[dr] \& \\
 \& X_2
\end{tikzcd}
\right) \qquad\quad \longrightarrow \qquad\quad X_3\text{\texttt{]}}
\]
\printedpage{171}
% head: A. Grothendieck : Dérivateurs, Ch. XII, Édition des Universités de Montpellier II et Paris VII
\textit{soit dans $W$.} (NB L'intérêt de ce critère est qu'il est `intrinsèque' à $Q$, ne faisant pas appel à une factorisation de $\beta_1$ ou de $\beta_2$ \texttt{[où} $\beta_i : S \longrightarrow X_i$\texttt{]} en $S \xrightarrow{\;\in \Cof_W\;} \tilde{X}_i \xrightarrow{\in W} X_i$. Cf. application intéressante dans section V 3.4., p. 146 ff., th. 4 et corollaires.)

\dgrblock{dgrb597}{The integral as a categorical homotopy pushout}
\begin{dcomment}{later}{p262:L'intérêt de ce critère est qu'il est}{The integral as a categorical homotopy pushout}
Grothendieck wants to identify the abstract operation just constructed with the categorical integral used earlier in the chapter. Proposition 20 says that the older construction was already providing a model for the same homotopy pushout. Its advantage is that the comparison can be read directly from the original square.

Here $\Phi$ has arrows $1\to0\leftarrow2$, and the diagram of categories is contravariant: its data are $X_1\leftarrow S\to X_2$. The associated fibred Grothendieck construction has objects $(i,x)$, with $x$ an object of the fibre at $i$. For $\alpha:i\to j$, a morphism over $\alpha$ is a map $x\to F(\alpha)(y)$ in the fibre at $i$. A compatible cocone to $X_3$ gives the displayed functor $\varphi_Q$.

For ordinary weak homotopy equivalences, this interpretation is the content of Thomason's homotopy-colimit theorem: the classifying space of the integral models the homotopy colimit of the classifying spaces of the diagram \citep[Theorem 1.2]{Thomason1979}. We checked its statement and a later proof in \citet[Theorem 26.8 and Corollary 26.9]{ChacholskiScherer2001}; the verification note distinguishes this from inspection of the original 1979 proof. In the present shape it gives
\[
 \left|N\mathcal X_Q\right|
 \simeq |NX_1|\amalg^{h}_{|NS|}|NX_2|.
\]
The fibred version follows from the usual covariant theorem by applying it to the opposite fibre categories on $\Phi^{\mathrm{op}}$ and then taking the opposite total category. Taking opposites leaves classifying spaces canonically homeomorphic. The equivalence explains the notation, but the manuscript's argument addresses a general $W$ under its stated axioms, rather than only $W_\infty$.

Its proof has one directional slip worth resolving before following the reduction. The transformation $\alpha:\tilde Q\to Q$ induces
\[
 \lambda:\mathcal X_{\tilde Q}\longrightarrow\mathcal X_Q,
 \qquad \varphi_{\tilde Q}=\varphi_Q\lambda.
\]
The edition reverses the horizontal arrow in this triangle. With the correct direction, the fibrewise weak-equivalence assertion gives $\lambda\in W$, and two-out-of-three gives the required equivalence of tests. Footnote 134 expressly identifies the strong localization axiom used for this step; it should not be silently dropped when $W$ is left arbitrary.

After both legs have been replaced by Dwyer immersions, put $P=X_1\amalg_SX_2$. The earlier result on pushouts of these immersions supplies $\mathcal X_Q\to P$ in $W$. The map to $X_3$ factors as $\mathcal X_Q\to P\xrightarrow{\pi}X_3$, so $\varphi_Q\in W$ if and only if $\pi\in W$. The integral and the abstract replacement criterion therefore agree. None of this requires the reverse assertion that a homotopy pushout reflects weak equivalences in an attached leg; that separate condition enters the appendix.
\end{dcomment}


\msp{page 263}
\textsc{Démonstration.} Considérons une factorisation de $S \longrightarrow X_1$ en
\[
S \xrightarrow{i_1} \tilde{X}_1 \xrightarrow{p_1} X, % [sic] X -> X_1, the factorisation is of S -> X_1
\]
avec $i_1$ immersion ouverte de Dwyer, $p_1 \in W$. Donc on a un carré commutatif
\[
\tilde{Q} = \left(
\begin{tikzcd}[ampersand replacement=\&]
S \arrow[r, hook, "\mathrm{Dw}"] \arrow[d] \& \tilde{X}_1 \arrow[d] \\
X_2 \arrow[r] \& X_3
\end{tikzcd}
\right)
\]
et un homomorphisme $\tilde{Q} \xrightarrow{\alpha} Q$ qui est l'identité sur les trois sommets distincts de $\tilde{X}_1$, et qui est $p_1$ en $\tilde{X}_1$. On en déduit un diagramme commutatif
\[
\begin{tikzcd}[ampersand replacement=\&, row sep=small, column sep=small]
\mathcal{X}_Q \arrow[rr, "\mathcal{X}_\alpha = \lambda"] \arrow[dr, "\varphi_Q"'] \& \& \mathcal{X}_{\tilde{Q}} \arrow[dl, "\varphi_{\tilde{Q}}"] \\
 \& X_3, \&
\end{tikzcd}
\]
et on a $\lambda \in W$ (car c'est une $\Phi$-flèche de catégories fibrées sur $\Phi$, qui induit une $W$-équivalence sur chacune de trois fibres.) Donc $\varphi_Q \in W \iff \varphi_{\tilde{Q}} \in W$, d'autre part il est clair que $Q$ est $W$-cocartésien si et seulement si $\tilde{Q}$ l'est. Donc on peut supposer que $S \hookrightarrow X_1$ est \texttt{[une]} immersion ouverte (de Dwyer) (\footnote[134]{On utilise ici W(3) (forme forte), i.e. Loc(3 bis) avec les notations de XVI. Mais on peut s'en passer, cf. plus bas.}). Appliquant ceci
\msp{page 264}
en changeant les rôles de $X_1$ et $X_2$, on voit qu'on peut supposer de plus que $S \longrightarrow X_2$ est également \texttt{[une]} immersion ouverte de Dwyer. Considérons alors
\[
\begin{tikzcd}[ampersand replacement=\&]
S \arrow[r, hook] \arrow[d, hook] \arrow[dr, phantom, "Q'"] \& X_1 \arrow[d, hook] \& \\
X_2 \arrow[r, hook] \& X_1 \amalg_S X_2 \arrow[r, "\pi"] \& X_3 \\
 \& \mathcal{X}_Q \arrow[u, "\varphi_{Q'}"] \arrow[ur, "\varphi_Q"'] \&
\end{tikzcd}
\]
\printedpage{172}
% head: A. Grothendieck : Dérivateurs, Ch. XII, Édition des Universités de Montpellier II et Paris VII
On a $\varphi_Q = \pi\varphi_{Q'}$, d'autre part on sait que $\varphi_{Q'} \in W$ (cor. 3, p. 149). Donc
\[
\pi \in W \qquad\quad \iff \qquad\quad \varphi_Q \in W,
\]
où $\pi \in W$ signifie que $Q$ est $W$-cocartésien, q.e.d.

\msp{page 265}

\dgrblock{dgrb598}{A four-object circle behind a strict pushout}
\begin{dcomment}{examples}{p264}{A four-object circle behind a strict pushout}
Return to the two-point span whose cylinder cofiber was calculated at ms.~230. The integral gives a smaller categorical model of the same circle, and explains why the word ``homotopy'' cannot be omitted from Proposition 20. Let $S$ be the discrete category on two objects $a,b$, and take $X_1=X_2=e$, both maps $S\to e$ being the unique functor. The ordinary pushout in $\mathrm{Cat}$ is $e$.

The fibred integral of this span has four objects, which we denote $u,v,a,b$. Its only nonidentity arrows are
\[
 u\longrightarrow a,\qquad u\longrightarrow b,
 \qquad v\longrightarrow a,\qquad v\longrightarrow b.
\]
There are no composable pairs of nonidentity arrows. Thus its nerve has four vertices, four nondegenerate edges, and no nondegenerate simplices of dimension at least two. Its realization is the four-edge cycle, homeomorphic to a circle. Equivalently, the cellular boundary map has rank three, so its first homology is $\mathbb Z$.

The comparison $\mathcal X_Q\to e$ therefore fails to be in $W_\infty$. The strict pushout square is not a homotopy pushout, and its actual homotopy pushout has the homotopy type of $S^1$. This four-edge circle and the earlier eight-edge cylinder cofiber both model the suspension of the two-point space. Neither original leg $S\to e$ is a Dwyer immersion: both identify the two objects. Replacing these legs before taking the ordinary pushout is precisely what preserves the missing circle. The example also fixes the variance of the integral without relying on a pictorial convention for its arrows.
\end{dcomment}


\section*{8\quad Appendice : Retour sur les axiomes W(7), W(7 bis) (pages 75, 76).}

(\footnote[135]{Reflexion provisoire, supplanté par XVI.}). On a vu (prop. 6, p. 78) que W(7 bis) implique W(7). Montrons une réciproque partielle. % [sic] Reflexion -> Réflexion, missing accent; supplanté -> supplantée, agreement (in footnote 135)

\dgrblock{dgrb599}{What the two converse gluing assertions require}
\begin{dcomment}{exposition}{p265:Montrons une réciproque partielle.}{What the two converse gluing assertions require}
Grothendieck wants to recognize a weak equivalence from the pieces of a decomposition. The familiar direction starts with equivalences on the common part and the two pieces and concludes that their homotopy pushouts are equivalent. His appendix asks for two converses: can the common part be recovered from the pieces and their union, and can the pieces be recovered from the common part and their union? These are substantive additions to homotopy-colimit invariance.

Write a square as $A\to B$, $A\to C$, with fourth vertex $P$, and write
\[
   C_e(u)=e\amalg^{\mathrm h}_{\operatorname{dom}(u)}
                    \operatorname{cod}(u)
\]
for the homotopy cofiber formed using the terminal object $e$. Whenever the square is a homotopy pushout, an equivalence $A\to B$ gives an equivalence $C\to P$. Grothendieck's conservative axiom asks for the converse as well. Applying that converse to the defining square of $C_e(u)$ says that a terminal homotopy cofiber detects an equivalence. This detection property is the extra ingredient in the proposition, through the full hypothesis W(7 bis).

It is useful to keep the three levels separate as the proof unfolds. Invariance of a homotopy pushout is formal once a coherent homotopy-colimit theory exists. The fact that an equivalence has terminal homotopy cofiber is equally formal. The converse, that terminal cofiber forces equivalence, is generally false for spaces; a complete example appears below. Thus the proposition under W(7 bis) and the later theorem for a general category of models must be read with different hypotheses. The last theorem drops an assumption that its third assertion still needs.
\end{dcomment}


\textbf{Proposition.} \textit{Supposons} W(7 bis) (axiome du carré cocartésien) \textit{satisfait, et plaçons nous sous les conditions de l'énoncé de} W(7), \textit{et plus généralement, dans le cas d'un diagramme cubique}
\[
\begin{tikzcd}[ampersand replacement=\&, row sep=small, column sep=small]
 \& X_2 \arrow[rr] \arrow[dd, "f_2" near end] \& \& X_3 \arrow[dd, "f_3"] \\
X_0 \arrow[ur] \arrow[rr, crossing over] \arrow[dd, "f_0"'] \& \& X_1 \arrow[ur] \& \\
 \& Y_2 \arrow[rr] \& \& Y_3 \\
Y_0 \arrow[ur] \arrow[rr] \& \& Y_1 \arrow[ur] \arrow[from=uu, crossing over, "f_1"' near end] \&
\end{tikzcd}
\]
\textit{où les deux carrés $X$ et $Y$ (face supérieure et inférieure du cube) sont $W$-cocartésiens. Considérons les conditions suivantes sur ce diagramme}
\begin{enumerate}[label={}]
\item[\textit{a)}] \textit{$f_0 \in W$.}
\item[\textit{b)}] \textit{$f_1$ et $f_2 \in W$.}
\item[\textit{c)}] \textit{$f_3 \in W$.}
\end{enumerate}
\textit{Alors on a les implications}
\[
a \text{ \textit{\&} } b \Longrightarrow c, \qquad\quad b \text{ \textit{\&} } c \Longrightarrow a.
\]
\textit{De plus, si $W$ est stable par facteurs directs (p. ex. $W$ fortement saturé), alors on a aussi}
\[
a \text{ \textit{\&} } c \Longrightarrow b,
\]
\textit{donc la conjonction} \texttt{[de]} \textit{deux quelconques parmi les conditions a) b) c) implique la troisième.}

\textsc{Démonstration.} Utilisant la factorisation
\printedpage{173}
% head: A. Grothendieck : Dérivateurs, Ch. XII, Édition des Universités de Montpellier II et Paris VII
\msp{page 266}
fonctorielle d'une flèche quelconque (ici les $f_i$) en $pi$, avec $i$ immersion ouverte de Dwyer et $p \in W$ (p. 229), on est remené au cas où les $f_i$ sont des immersions ouvertes de Dwyer. Mais pour toute flèche $g : S \longrightarrow T$ qui est dans $\Cof_W$, soient $Z(g)$ son mapping-cone % [sic] remené -> ramené, spelling
\[
\begin{tikzcd}[ampersand replacement=\&, row sep=small]
S \arrow[rr, "g"] \arrow[dr] \& \& T \arrow[dr, "\alpha(g)"] \& \\
 \& e \arrow[rr, "\bar{g}"] \& \& Z(g),
\end{tikzcd}
\]
et
\[
\alpha(g) : T = \mathrm{but}(g) \longrightarrow Z(g)
\]
la flèche canonique (\footnote[136]{C'est un peu vif, car il n'est pas clair que $X_1' \amalg_{X_0} X_2' \longrightarrow X_1 \amalg_{X_0} X_2$ soit \texttt{[dans]} $W$, sans utiliser déjà un des trois ingrédients de W(7) (savoir $f_0, f_1, f_2 \in W \implies f_3 \in W$). Il semble donc que notre argument ne marche que moyennant une hypothèse supplémentaire de cet ordre \dots Mais il y a le \textit{lemme de Brown,} cf. XVI \S3, notamment p. 20.}). Celle-ci est fonctorielle en $g$, de sorte que le diagramme cubique donné se prolonge en un double cube
\[
\begin{tikzcd}[ampersand replacement=\&, row sep=small, column sep=small]
 \& X_2 \arrow[rr] \arrow[dd, "f_2" near end] \& \& X_3 \arrow[dd, "f_3"] \\
X_0 \arrow[ur] \arrow[rr, crossing over] \arrow[dd, "f_0"'] \& \& X_1 \arrow[ur] \& \\
 \& Y_2 \arrow[rr] \arrow[dd, "\alpha_2" near end] \& \& Y_3 \arrow[dd, "\alpha_3"] \\
Y_0 \arrow[ur] \arrow[rr, crossing over] \arrow[dd, "\alpha_0"'] \& \& Y_1 \arrow[ur] \arrow[from=uu, crossing over, "f_1"' near end] \& \\
 \& Z_2 \arrow[rr] \& \& Z_3 \\
Z_0 \arrow[ur] \arrow[rr] \& \& Z_1 \arrow[ur] \arrow[from=uu, crossing over, "\alpha_1"' near end] \&
\end{tikzcd}
\]
avec
\[
Z_i = Z(f_i),\ \alpha_i = \alpha(f_i) \qquad\quad \text{pour } i \in \{0,1,2,3\}.
\]
Notons que W(7 bis) implique que si $g \in \Cof_W$,
\msp{page 267}
on a
\[
g \in W \qquad\quad \iff \qquad\quad \bar{g} \in W, \qquad\quad \text{i.e. } Z(g)\ W\text{-asphérique},
\]
en particulier
\[
f_i \in W \qquad\quad \iff \qquad\quad Z_i\ W\text{-asphérique}.
\]
Ainsi, l'énoncé à prouver sur le cube initial se ramène à un énoncé sur le carré
\printedpage{174}
% head: A. Grothendieck : Dérivateurs, Ch. XII, Édition des Universités de Montpellier II et Paris VII
\[
\begin{tikzcd}[ampersand replacement=\&, row sep=small, column sep=small]
 \& Z_1 \arrow[dr] \& \\
Z_0 \arrow[ur] \arrow[dr] \& \& Z_3, \\
 \& Z_2 \arrow[ur] \&
\end{tikzcd}
\]
savoir que les trois conditions
\begin{enumerate}[label={}]
\item[a$'$)] $Z_0$ $W$-asphérique.
\item[b$'$)] $Z_1$ et $Z_2$ $W$-asphériques.
\item[c$'$)] $Z_3$ $W$-asphérique.
\end{enumerate}
donnant lieu aux implications
\[
\mathrm{a}' \text{ \& } \mathrm{b}' \Longrightarrow \mathrm{c}', \qquad\quad \mathrm{b}' \text{ \& } \mathrm{c}' \Longrightarrow \mathrm{a}',
\]
et que si $W$ est stable par facteurs directs, alors
\[
\mathrm{a}' \text{ \& } \mathrm{c}' \Longrightarrow \mathrm{b}'.
\]

C'est à peu de choses près le corollaire de la page 76, complété à la page 77 (corollaire), à cela près qu'on y suppose les flèches des immersions ouvertes, et le carré cartésien (donc aussi cocartésien) dans Cat, i.e. $Z_3 = Z_1 \cup Z_2$, $Z_0 = Z_1 \cap Z_2$. Mais
\msp{page 268}
ce qui importe pour la validité de cet énoncé, c'est (on s'en doute) que $Z_*$ soit $W$-cocartésien. Donc il faut d'abord prouver ceci.

\dgrblock{dgrb600}{The gluing step that Brown supplies}
\begin{dcomment}{context}{p267:C'est à peu de choses près le corollaire de la page 76}{The gluing step that Brown supplies}
The factorization into a Dwyer immersion and a weak equivalence gives Grothendieck an explicit replacement of every vertical edge. To use the replacement on an entire cube, however, one must know that the induced replacement of a pushout is again a weak equivalence. Footnote~136 identifies this dependence. It is a gap in the chosen proof, not a counterexample to the conditional proposition.

The later discussion is unusually precise about the dependency. In Chapter~XVI, Commentaire~3.2, Grothendieck observes that the appendix had used the fact that Dwyer immersions are $W$-cofibrations at a point where his earlier proof of that fact already used a gluing assertion. Lemma~3.2.1 and Corollary~3.2.2 isolate the required Brown lemma; Proposition~3.3 then records the intended equivalence of the strong axioms. The chapter identifies the repair but also says explicitly that its author does not have a proof of that lemma there \citep[\S3, ms.~pp.~18--23]{GrothendieckDerivateursXVI}.

Here is the precise external result that closes this dependency. Suppose a class $C$ contains the isomorphisms and all maps from an initial object, is stable under composition and pushout, and every map factors as a member of $C$ followed by a weak equivalence. Suppose also that pushout preserves $C\cap W$. These are the relevant axioms of a category of cofibrant objects. Brown's gluing lemma implies that pushing any weak equivalence out along a member of $C$ preserves weak equivalences. In Brown's original formulation this is the dual of Part~I, \S4, Lemma~2; the proof uses his factorization lemma and the invariance result in Lemma~1 \citep[pp.~428--429]{Brown1973}. With $C$ the Dwyer immersions, this yields $C\subseteq\Cof_W$ and legitimizes the replacement of the cube. No converse detection property is supplied by Brown's lemma; that remains the role of W(7 bis).
\end{dcomment}


\textbf{Lemme.} \textit{Si le diagramme cubique p. 265 est tel que les faces $X$, $Y$ soient $W$-cocartésiennes, et les $f_i$ ($i \in \{0,1,2,3\}$) des $W$-cofibrations, alors le carré $Z$ des cofibres des $f_i$ est lui aussi $W$-cocartésien.}

Ce lemme résulte de la théorie du `déploiement des diagrammes cubiques', relatif à une catégorie de modèles $(\mathcal{M}, W)$ admettant une théorie des cofibres $W$-homotopiques, qu'il me faudra bien écrire! (Mais l'existence pour $(\mathrm{Cat}, W)$ d'une telle théorie des cofibres $W$-homotopiques, et par suite de celle des déploiements des diagrammes dans $\mathcal{M}$ de type $(\boldsymbol{\Delta}^1)^I$, $I$ ensemble d'indices fini, ne dépend que de W(7 bis) - on n'a pas à faire appel à W(7).)

\dgrblock{dgrb601}{Taking cofibers commutes with the pushout}
\begin{dcomment}{exposition}{p268:ne dépend que de W(7 bis)}{Taking cofibers commutes with the pushout}
The useful part of Grothendieck's large diagram is the compatibility of two successive homotopy pushouts. We can prove the lemma on p.~268 directly from this compatibility, without using the claim that a terminal cofiber detects an equivalence.

Let $J=(1\leftarrow0\to2)$, and let $X\to Y$ be a coherent morphism of homotopy-pushout squares in a homotopy theory with a terminal object and finite homotopy colimits. Put $Z_j=C_e(X_j\to Y_j)$. Computing the colimit in the two orders gives
\[
\begin{aligned}
 \mathop{\mathrm{hocolim}}_{j\in J} Z_j
 &\simeq
 \left(\mathop{\mathrm{hocolim}}_{j\in J}Y_j\right)
 \amalg^{\mathrm h}_{\mathop{\mathrm{hocolim}}_{j\in J}X_j}
 \left(\mathop{\mathrm{hocolim}}_{j\in J} e\right)\\
 &\simeq Y_3\amalg^{\mathrm h}_{X_3}e=Z_3.
\end{aligned}
\]
The first equivalence is the interchange of colimits of the single diagram indexed by $J$ and the cofiber span. The constant span $e\leftarrow e\to e$ has homotopy pushout $e$, which proves the second equivalence. Under these identifications the comparison map is the one in the given square $Z$, so this proves the lemma, including its compatibility with the prescribed edges. In derivator language the same calculation is the composition of left Kan extensions, with the parameter retained; the pointwise criterion is \citet[Corollary~3.14]{Groth2013}.

This proof also locates the excessive assertion on p.~269. From $u\in W$ one may conclude $C_e(u)\simeq e$ in any such theory. Going back from $C_e(u)\simeq e$ to $u\in W$ uses conservativity. The full W(7 bis) hypothesis permits that return, but the statement that the whole reduction used only its forward half is false. The pointwise-cofiber lemma survives without the reverse half; the reduction of all three cube implications does not.
\end{dcomment}


\msp{page 269}
Ce `deploiement' est un diagramme infini dont toutes les suites infinies, dans les trois directions des axes de coordonnées, sont des `suites de suspension'. Je reproduis la partie utile du déploiement. % [sic] deploiement -> déploiement, missing accent

% ---- from Der-ChXIIc-30-03-03_p061-064.tex
\printedpage{175}
% head: A. Grothendieck : Dérivateurs, Ch. XII, Édition des Universités de Montpellier II et Paris VII
% foot: Version 30 mars 2003

% dangling arrows (from X_2, X_3, Y_2, Y_3, Z_2, Z_3) as printed
\[
\begin{tikzcd}[ampersand replacement=\&, column sep=small, row sep=small]
 \& \& {} \& \& {} \& \& U \arrow[dd, "\mu"] \\
 \& X_2 \arrow[rr] \arrow[dd] \arrow[ur, shorten >=1ex] \& \& X_3 \arrow[rr] \arrow[dd] \arrow[ur, shorten >=1ex] \& \& T \arrow[ur] \arrow[dd] \& \\
X_0 \arrow[rr, crossing over] \arrow[ur] \arrow[dd] \& \& X_1 \arrow[rr, crossing over] \arrow[ur] \arrow[dd, crossing over] \& \& S \arrow[ur, "\lambda"] \arrow[dd, crossing over] \& \& U' \arrow[dd] \\
 \& Y_2 \arrow[rr] \arrow[dd] \arrow[ur, shorten >=1.5ex] \& \& Y_3 \arrow[rr] \arrow[dd] \arrow[ur, shorten >=1.5ex] \& \& T' \arrow[ur] \arrow[dd] \& \\
Y_0 \arrow[rr, crossing over] \arrow[ur] \arrow[dd] \& \& Y_1 \arrow[rr, crossing over] \arrow[ur] \arrow[dd, crossing over] \& \& S' \arrow[ur, "\lambda'"] \arrow[dd, crossing over] \& \& U'' \\
 \& Z_2 \arrow[rr] \arrow[ur, shorten >=1.5ex] \& \& Z_3 \arrow[rr] \arrow[ur, shorten >=1.5ex] \& \& T'' \arrow[ur] \& \\
Z_0 \arrow[rr] \arrow[ur] \& \& Z_1 \arrow[rr] \arrow[ur] \& \& S'' \arrow[ur, "\lambda''"] \& \&
\end{tikzcd}
\]

La condition que $X$ soit $W$-cocartésien s'exprime par $\lambda \in W$, ou encore par $U \in \Asph_W$ (car $U$ s'identifie à la cofibre homotopique $Z(\lambda)$ de $\lambda$ \dots\,), de même $Y$ $W$-cocartésien s'exprime par $U' \in \Asph_W$, et $Z$ $W$-cocartésien par $U'' \in \Asph_W$. Or $U'' \sim Z(\mu)$, où $\mu : U \longrightarrow U'$ est dans $W$ si $U$, $U'$ sont $W$-asphériques, donc $U'' = Z(\mu)$ est alors $W$-asphérique. Cela prouve le lemme.

Cela montre que, moyennant W(7 bis) (en fait, on a utilisé seulement la partie ‘universelle’ de W(7 bis)), savoir que dans
\begin{tikzcd}[ampersand replacement=\&, column sep=small, row sep=small, baseline=(\tikzcdmatrixname-2-1.base)]
T_0 \arrow[rr] \arrow[dr, "i"'] \& \& T_2 \arrow[dr, "i'"] \& \\
 \& T_1 \arrow[rr] \& \& T_3
\end{tikzcd}
diagramme cocartésien d'immersions ouvertes, $i \in W \Longrightarrow i' \in W$ (au lieu de $\Longleftrightarrow$),

\msp{page 270}

la validité de W(7) se ramène au cas où le carré $Y$ est le carré trivial
\begin{tikzcd}[ampersand replacement=\&, column sep=small, row sep=small, baseline=(\tikzcdmatrixname-2-1.base)]
 \& e \arrow[dr] \& \\
e \arrow[ur] \arrow[dr] \& \& e, \\
 \& e \arrow[ur] \&
\end{tikzcd}
$e$ l'objet final de Cat.

On peut maintenant conclure : On a (moyennant la partie directe de W(7 bis))
\[
\mathrm{a}' \ \& \ \mathrm{b}' \Longrightarrow \mathrm{c}',
\]
car si $Z_0$, $Z_1$ sont $\in \Asph_W$, alors $i_1 \in W$, donc $j_1 \in W$, et si de plus $Z_2 \in \Asph_W$, on a donc $Z_3 \in \Asph_W$.
\[
\begin{tikzcd}[ampersand replacement=\&, column sep=small, row sep=small]
 \& Z_1 \arrow[dr, "j_2"] \& \\
Z_0 \arrow[ur, "i_1"] \arrow[dr, "i_2"'] \& \& Z_3 \\
 \& Z_2 \arrow[ur, "j_1"'] \&
\end{tikzcd}
\]

Si on a aussi la partie inverse (‘inhabituelle’) de W(7 bis), alors
\[
\mathrm{b}' \ \& \ \mathrm{c}' \Longrightarrow \mathrm{a}',
\]
car si $Z_1$ et $Z_3$ sont $W$-asphériques, alors $j_2 \in W$, donc $i_2 \in W$, et si de plus $Z_2$ est $W$-asphérique, on conclut $Z_0$ $W$-asphérique.

\printedpage{176}
% head: A. Grothendieck : Dérivateurs, Ch. XII, Édition des Universités de Montpellier II et Paris VII
% foot: Version 30 mars 2003

Enfin, supposons que $W$ soit \textit{stable par facteurs directs,} et W(7 bis) satisfait sous forme directe. Je dis qu'alors
\[
\mathrm{c}' \ \& \ \mathrm{a}' \Longrightarrow \mathrm{b}',
\]
i.e.
\[
Z_0 \text{ et } Z_3 \ \in \Asph_W \qquad\quad \Longrightarrow \qquad\quad Z_1 \text{ et } Z_2 \ \in \Asph_W.
\]
Factorisant $i_1$ et $i_2$ en
\[
\begin{aligned}
Z_0 &\xrightarrow{i'_1} Z'_1 \xrightarrow{\alpha_1} Z_1, \\
Z_0 &\xrightarrow{i'_2} Z'_2 \xrightarrow{\alpha_2} Z_2,
\end{aligned}
\]
avec $i'_1$, $i'_2$ des immersions ouvertes de Dwyer, et $\alpha_1, \alpha_2 \in W$, on est ramené au cas où

\msp{page 271}

$i_1$, $i_2$ sont des \textit{immersions ouvertes de Dwyer,} et où le carré est cocartésien (ce qui implique qu'il est $W$-cocartésien). Faisant alors le cochangement de base par $Z_0 \xrightarrow{p_0} Z'_0 = e$, on trouve un homomorphisme de carrés cocartésiens
\[
\begin{tikzcd}[ampersand replacement=\&, column sep=small, row sep=small]
 \& Z_1 \arrow[rr] \arrow[dd, "p_1" near end] \& \& Z_3 \arrow[dd, "p_3"] \\
Z_0 \arrow[ur] \arrow[rr, crossing over] \arrow[dd, "p_0"'] \& \& Z_2 \arrow[ur] \& \\
 \& Z'_1 \arrow[rr] \& \& Z'_3, \\
e = Z'_0 \arrow[ur] \arrow[rr] \& \& Z'_2 \arrow[ur] \arrow[from=uu, crossing over, "p_2"' near end] \&
\end{tikzcd}
\]
où toutes les faces du cube sont $W$-cocartésiennes, et comme $p_0 \in W$ ($Z_0$ étant $W$-asphérique), les autres flèches verticales $p_1$, $p_2$, $p_3$ sont \texttt{[dans]} $W$, donc le carré $Z'_*$ satisfait aussi $Z'_0$ et $Z'_3$ \texttt{[}$W$\texttt{-]}asphériques (et de plus $Z'_0 = e$), et il faut prouver $Z'_1$, $Z'_2$ $W$-asphériques (d'où itou pour $Z_1$, $Z_2$). Donc, on est ramené au cas particulier où $Z_0 = e$.

Mais dans ce cas il existe une rétraction (triviale!) de $Z_1$ sur $Z_0$, donc aussi de $Z_3$ sur $Z_2$. Ainsi, $Z_2$ est un facteur

\msp{page 272}

direct d'un objet $W$-asphérique, ce qui implique, par l'hypothèse sur $W$, que $Z_2$ est $W$-asphérique. Par symétrie, $Z_1$ l'est aussi. On a gagné!

\dgrblock{dgrb602}{The retract argument is valid at the level of objects}
\begin{dcomment}{exposition}{p272:Par symétrie, $Z_1$ l'est aussi. On a gagné!}{The retract argument is valid at the level of objects}
Grothendieck's reduction to a common part equal to $e$ is sound and deserves to be retained. If
\[
 Z_3\simeq Z_1\amalg^{\mathrm h}_{Z_0}Z_2,
 \qquad Z_0\simeq e,
\]
then homotopy-colimit invariance identifies $Z_3$ with the homotopy coproduct of two objects under $e$. Each summand is a retract of that coproduct: the retraction onto $Z_1$ is induced by the identity of $Z_1$ and the composite $Z_2\to e\to Z_1$. Both agree on $e$, so the universal property supplies the retraction; its composite with the summand inclusion is the identity. The same construction applies to $Z_2$. Consequently, if $Z_3$ is $W$-aspherical and $W$-aspherical objects are closed under retracts, both summands are $W$-aspherical. This proves the object statement without any conservative axiom.

For the source's weaker hypothesis of closure under strict retracts, realize this argument using cofibration representatives $e\to Z_i^c$ in the category under $e$ and their ordinary pushout $P^c=Z_1^c\amalg_e Z_2^c$. These representatives compute the displayed homotopy pushout. The identity of $Z_1^c$ and the composite $Z_2^c\to e\to Z_1^c$ agree strictly on $e$, so the pushout retraction and its splitting are literal maps, with composite the identity. The comparison with $Z_3$ makes $P^c$ aspherical; strict retract closure applies to $Z_i^c$, and the replacement weak equivalences give the assertion for $Z_i$.

The phrase \emph{facteur direct} here means a retract, as the explicit construction shows. Strong saturation of $W$ also suffices because localization carries a retract diagram to a retract diagram, and a retract of an isomorphism is an isomorphism.

What this argument does not supply is the return from the objects $Z_i=C_e(f_i)$ to the maps $f_i$. Terminal cofibers need not detect equivalences. The object statement just proved therefore cannot by itself justify assertion~3\,\textdegree{} of the following general theorem, or the equivalence asserted in Corollary~1. Under the full conservative hypothesis that return is valid, and the final Corollary~2 survives.
\end{dcomment}


\textbf{NB} En fait, ces arguments sont généraux, et donnent ceci :

\setcounter{footnote}{136}
\textbf{Théorème.} \textit{Soit $(\mathcal{M}, W)$ une catégorie des modèles, admettant une théorie de la cofibre $W$-homotopique \textup{(}\footnote{Mais on ne sait vérifier l'existence d'une telle théorie dans $(\mathrm{Cat}, W)$ que si on admet la forme faible de W(7), i.e. $f_0, f_1, f_2 \in W \implies f_3 \in W$, ou quelque cas particulier de celle-ci. Mais le \textit{lemme de Brown} sauve la situation, cf. XV \S3.}\textup{)} (i.e. on a une factorisation de toute flèche $f : X \longrightarrow Y$ en $pi$, avec $i \in \Cof_W$, $p \in W$.) On a donc une notion de carré $W$-cocartésien dans $\mathcal{M}$ \textup{(}\footnote{à vérifier que la factorisation suffit, même si tous les objets ne sont pas fibrants sur $e$. Mais on peut alors remplacer $\mathcal{M}$ par $\mathcal{M}_f$, et développer ce formalisme des fibres homotopiques en conséquence \dots}\textup{)}. Soit
\printedpage{177}
% head: A. Grothendieck : Dérivateurs, Ch. XII, Édition des Universités de Montpellier II et Paris VII
% foot: Version 30 mars 2003
alors $f : X \longrightarrow Y$ un homomorphisme d'un carré $W$-cocartésien $X$ dans un autre $Y$, $f = (f_0, f_1, f_2, f_3)$.}
\[
\begin{tikzcd}[ampersand replacement=\&, column sep=small, row sep=small]
 \& X_2 \arrow[rr] \arrow[dd, "f_2" near end] \& \& X_3 \arrow[dd, "f_3"] \\
X_0 \arrow[ur] \arrow[rr, crossing over] \arrow[dd, "f_0"'] \& \& X_1 \arrow[ur] \& \\
 \& Y_2 \arrow[rr] \& \& Y_3 \\
Y_0 \arrow[ur] \arrow[rr] \& \& Y_1 \arrow[ur] \arrow[from=uu, crossing over, "f_1"' near end] \&
\end{tikzcd}
\]

\textit{Considérons les conditions sur $f$}

\begin{enumerate}
\item[\textit{a)}] \textit{$f_0 \in W$.}
\item[\textit{b)}] \textit{$f_1$ et $f_2 \in W$.}
\item[\textit{c)}] \textit{$f_3 \in W$.}
\end{enumerate}

\textit{Alors}

\begin{enumerate}
\item[\textit{1°)}] \textit{a \& b $\Longrightarrow$ c.}

\msp{page 273}

\item[\textit{2°)}] \textit{Supposons que $\mathcal{M}$ ait un objet final $e$. Alors l'implication}
\[
\textit{b \& c} \Longrightarrow \textit{a}
\]
\textit{pour}~tout \textit{homomorphisme $f : X \longrightarrow Y$ de carrés $W$-cocartésiens équivaut à l'axiome suivant sur $W$ :}

Axiome des carrés $W$-cocartésiens conservatifs (\footnote{cf. argument page 78 qui montre que l'implication b \& c $\Longrightarrow$ a entraine l'axiome conservatif.}) : \textit{Si}
\[
\begin{tikzcd}[ampersand replacement=\&, column sep=small, row sep=small]
 \& Z_1 \arrow[dr, "j_2"] \& \\
Z_0 \arrow[ur, "i_1"] \arrow[dr, "i_2"'] \& \& Z_3 \\
 \& Z_2 \arrow[ur, "j_1"'] \&
\end{tikzcd}
\]
\textit{est $W$-cocartésien, alors $j_1 \in W \Longrightarrow i_1 \in W$ (donc en fait $i_1 \in W \Longleftrightarrow j_1 \in W$ et $i_2 \in W \Longleftrightarrow j_2 \in W$).}

\textit{Cet axiome est équivalent à la conjonction des deux axiomes suivants.}

\dgrblock{dgrb603}{A direct test for conservative homotopy pushouts}
\begin{dcomment}{exposition}{p273:Cet axiome est équivalent}{A direct test for conservative homotopy pushouts}
Grothendieck now isolates the property needed by the preceding reduction. Once homotopy pushouts are available coherently, there is a short proof of the test that also clarifies Cons~1 and Cons~2. We work in a homotopy theory with a terminal object and finite homotopy colimits, and use ``equivalence'' for a morphism detected as an isomorphism by its localization.

For every homotopy-pushout square
\[
\begin{tikzcd}[ampersand replacement=\&]
 A\arrow[r,"u"]\arrow[d] \& B\arrow[d]\\
 C\arrow[r,"v"'] \& P
\end{tikzcd}
\]
associativity gives a canonical equivalence
\[
 C_e(v)=e\amalg_C^{\mathrm h}P
 \simeq e\amalg_C^{\mathrm h}(C\amalg_A^{\mathrm h}B)
 \simeq e\amalg_A^{\mathrm h}B=C_e(u).
\]
Thus, if terminal homotopy cofiber detects equivalences, $v$ being an equivalence forces $u$ to be one. Conversely, apply the conservative-pushout axiom to the defining cofiber square: if $C_e(u)\simeq e$, its lower edge $e\to C_e(u)$ is an equivalence, and the axiom forces $u$ to be an equivalence. This proves the equivalence of conservativity and Cons~1, with the latter interpreted for a cofibration replacement of every map. The intended cocartesian Cons~2 follows: if $B,C,P$ are terminal up to equivalence, the edge $B\to P$ is an equivalence, hence so is $A\to C$; consequently $A\simeq e$. It is therefore redundant under this coherent formulation.

These observations complete the proof of the two converse cube implications under the appropriate hypothesis. Given a morphism $f:X\to Y$ of homotopy-pushout squares, its pointwise cofibers $Z_i=C_e(f_i)$ form a homotopy-pushout square by the preceding lemma. If $f_1,f_2,f_3$ are equivalences, then $Z_1,Z_2,Z_3\simeq e$; conservativity gives $Z_0\simeq e$, and Cons~1 gives $f_0\in W$. Conversely, the implication $b\mathbin{\&}c\Rightarrow a$ applied to the map from a square $(A,B,C,P)$ to $(B,B,P,P)$ with components $(u,1_B,v,1_P)$ proves conservativity immediately.

If instead $f_0,f_3$ are equivalences and aspherical objects are closed under retracts, the object argument on pp.~271--272 gives $Z_1,Z_2\simeq e$; Cons~1 then gives $f_1,f_2\in W$. Together with ordinary pushout invariance this proves the final Corollary~2. It also identifies exactly what must be added to the printed assertion~3\,\textdegree{}: a cofiber-detection hypothesis, equivalently conservative homotopy pushouts, in addition to the retract condition.
\end{dcomment}


\textbf{Cons 1.} \textit{Si on a un carré cocartésien}
\printedpage{178}
% head: A. Grothendieck : Dérivateurs, Ch. XII, Édition des Universités de Montpellier II et Paris VII
% foot: Version 30 mars 2003
\[
\begin{tikzcd}[ampersand replacement=\&, column sep=small, row sep=small]
S \arrow[rr] \arrow[dr] \& \& T \arrow[dr] \& \\
 \& e \arrow[rr] \& \& U,
\end{tikzcd}
\]
\textit{avec $f \in \Cof_W$, alors $U \in \Asph_W \Longrightarrow f \in W$ (donc $U \in \Asph_W \Longleftrightarrow f \in W$).}

\textbf{Cons 2.} \textit{Si on a un carré cartésien,}
\[
\begin{tikzcd}[ampersand replacement=\&, column sep=small, row sep=small]
 \& Z_1 \arrow[dr] \& \\
Z_0 \arrow[ur, "i_1"] \arrow[dr, "i_2"'] \& \& Z_3, \\
 \& Z_2 \arrow[ur] \&
\end{tikzcd}
\]
\textit{avec $i_1, i_2 \in \Cof_W$, alors : Si $Z_1$, $Z_2$, $Z_3$ sont $W$-asphériques, $Z_0$ l'est aussi.}

\dgrblock{dgrb604}{Why Cons 2 must concern pushouts}
\begin{dcomment}{examples}{p273:Si on a un carré cartésien}{Why Cons 2 must concern pushouts}
The edition prints \emph{cartésien} in Cons~2. Taken literally, this gives a different and false consequence of conservativity. The issue can be seen inside a stable homotopy theory, where conservative homotopy pushouts do hold.

Let $k$ be a field and use homologically graded chain complexes. Put $A=ka$ in degree zero, and let
\[
 B=(kb\xrightarrow{d}ka),\qquad
 C=(kc\xrightarrow{d}ka),\qquad d(b)=d(c)=a.
\]
Embed both in the complex
\[
 D=(kd\longrightarrow kb\oplus kc\longrightarrow ka),
 \qquad \partial d=b-c,\quad \partial b=\partial c=a.
\]
Here $d$ as a basis symbol lies in degree two. Both $B$ and $C$ are acyclic. In $D$ the first differential is injective, its image is the kernel of the second, and the second is surjective, so $D$ is acyclic too. Their ordinary pullback over $D$ is the intersection $A$, whose zeroth homology is $k$. All four inclusions are degreewise monomorphisms. Such a map is a $W$-cofibration for quasi-isomorphisms: pushout leaves its cokernel unchanged, and the long exact homology sequences show that a weak equivalence of sources induces one of the pushouts. Over a field these maps are also cofibrations in the usual model structure.

Thus a strict cartesian square with the printed cofibration hypotheses need not satisfy Cons~2, even when the surrounding derived theory is stable. Replacing \emph{cartésien} by \emph{cocartésien} restores the statement used in the argument. The source wording is retained in the text and recorded as an edition slip.
\end{dcomment}


\msp{page 274}

\item[\textit{3°)}] \textit{Supposons que tout facteur direct d'un objet $W$-asphérique soit $W$-asphérique, ce qui est le cas notamment si $W$ est lui-même stable par facteurs directs, a fortiori si $W$ est fortement saturé. Supposons de plus que $\mathcal{M}$ ait} \texttt{[un]} \textit{objet final. Alors on a (sous les conditions générales du début)}
\[
\textit{c \& a} \Longrightarrow \textit{b,}
\]
\textit{i.e. si $f_0$, $f_3$ sont} \texttt{[dans]} \textit{$W$, de même $f_1$, $f_2$.}
\end{enumerate}

Cette dernière partie peut se déployer ainsi :

\dgrblock{dgrb605}{An explicit cube testing the generalization}
\begin{dcomment}{development}{p274:Cette dernière partie peut se déployer ainsi}{An explicit cube testing the generalization}
Grothendieck's retraction proof invites a test of the stronger assertion about arbitrary maps of squares. We give the test in full, since it distinguishes a missing hypothesis from a missing step in a proof. The acyclic complex used below is classical; the assembly into the two diagnostic cubes is our application to this passage. We checked the source's later Chapter~XVI discussion, Brown's gluing lemmas, and the cofiber and Mayer--Vietoris treatments cited here; we did not find this particular counterexample to the printed generalization in those sources. No priority is claimed.

\textbf{A cube with $a$ and $c$, but not $b$.}\ours\ 
In simplicial sets with weak homotopy equivalences, there is a map between homotopy-pushout squares whose zeroth, second and third components are equivalences and whose first component is not.

\emph{Proof.} Begin with the finite CW complex $A$ having one vertex, two edges $a,b$, and two two-cells attached by $a^5b^{-3}$ and $b^3(ab)^{-2}$ \citep[Example~2.38]{Hatcher2002}. Its cellular boundary in positive dimensions is
\[
 0\longrightarrow\mathbb Z^2
 \xrightarrow{\left(\begin{smallmatrix}5&-2\\-3&1\end{smallmatrix}\right)}
 \mathbb Z^2\longrightarrow0.
\]
The determinant is $-1$, so $A$ is acyclic. Its fundamental group is nontrivial: sending $a$ to $(12345)$ and $b$ to $(153)$ respects the relations, since their product is $(23)(45)$, and gives a nontrivial permutation representation. The suspension $\Sigma A$ is simply connected by van Kampen and acyclic by the suspension isomorphism in reduced homology. Hurewicz and Whitehead therefore make it contractible \citep[Theorem~4.32 and Corollary~4.33]{Hatcher2002}.

Choose a simplicial set $K$ of the homotopy type of $A$, let $CK=K\star\Delta^0$, and write $e=\Delta^0$. The inclusion $K\hookrightarrow CK$ is a cofibration, $CK$ is contractible, and $CK\amalg_K e$ has the homotopy type of $\Sigma A$, hence is weakly contractible. Consider
\[
 X=\left(\begin{tikzcd}[ampersand replacement=\&, column sep=small,row sep=small]
 K\arrow[r,equal]\arrow[d,hook] \& K\arrow[d,hook]\\
 CK\arrow[r,equal] \& CK
 \end{tikzcd}\right),\qquad
 Y=\left(\begin{tikzcd}[ampersand replacement=\&, column sep=small,row sep=small]
 K\arrow[r]\arrow[d,hook] \& e\arrow[d]\\
 CK\arrow[r] \& CK\amalg_K e
 \end{tikzcd}\right).
\]
Both strict pushouts compute homotopy pushouts, since the common vertical leg is a cofibration. There is a strict map $X\to Y$ with components
\[
 (f_0,f_1,f_2,f_3)=(1_K,\ K\to e,\ 1_{CK},\ CK\to CK\amalg_K e).
\]
Its zeroth and second components are identities; its third connects weakly contractible objects; its first is not a weak equivalence because $\pi_1(K)\ne1$. This proves the assertion. All simplicial sets are cofibrant, factorizations exist, and weak equivalences are strongly saturated and closed under retracts. Thus the failure is present under the printed hypotheses of assertion~3\,\textdegree{}, not a failure of replacement machinery.

The same $K$ tests the other converse. Map the homotopy-pushout square with vertices $(K,CK,CK,CK\amalg_K CK)$ to the constant square at $e$. Its two cone vertices and its suspension vertex are weakly contractible, while $K$ is not. Hence $b\mathbin{\&}c\Rightarrow a$ fails for spaces, as the conservative axiom predicts. Finally, the object statement in Corollary~1 remains valid for spaces by the retract proof above. Its proposed equivalence with $a\mathbin{\&}c\Rightarrow b$ therefore fails as printed. Adding cofiber detection repairs both assertions.
\end{dcomment}


\textbf{Corollaire 1.} \textit{Sous l'hypothèse supplémentaire} \textup{3\,°)} \textit{sur $W$, la validité de c \& a $\Longrightarrow$ b (pour toute flèche $f$ entre carrés $W$-cocartésiens) équivaut à chacun des deux autres axiomes suivants sur $W$ (qui sont donc équivalents) :}

\begin{enumerate}
\item[\textit{1)}] \textit{Si on a un carré $W$-cocartésien}
\[
\begin{tikzcd}[ampersand replacement=\&, column sep=small, row sep=small]
 \& Z_1 \arrow[dr] \& \\
Z_0 \arrow[ur] \arrow[dr] \& \& Z_3, \\
 \& Z_2 \arrow[ur] \&
\end{tikzcd}
\]
\textit{alors $Z_0$ et $Z_3$ $W$-asphériques $\Longrightarrow$ $Z_1$ et $Z_2$ $W$-asphériques.}
\item[\textit{1$'$)}] \textit{Comme 1), avec $Z_0 = e$ (dans le cas où on suppose que $\mathcal{M}$ admet} \texttt{[un]} \textit{objet final $e$).}
\end{enumerate}

\msp{page 275}

\textbf{Corollaire 2.} \textit{Supposons que $W$ satisfasse l'axiome des carrés $W$-cocartésiens conservatifs, que tout facteur direct d'un objet $W$-asphérique soit $W$-asphérique (p. ex. $W$ stable par facteurs directs), enfin que $\mathcal{M}$ admette un objet final. Alors pout tout homomorphisme $f : X \longrightarrow Y$ de carrés $W$-cocartésiens, la conjonction de deux parmi les conditions a) b) c) de p. 272 entraîne la troisième.} % [sic] pout -> pour, typographical error


\dgrblock{dgrb606}{The stable form of the cube calculus}
\begin{dcomment}{later}{p275}{The stable form of the cube calculus}
The last corollary has a particularly transparent home in a stable derivator. Coherent cocartesian squares then are also cartesian, and their geometry gives a Mayer--Vietoris triangle
\[
 X_0\xrightarrow{(i_1,-i_2)}X_1\oplus X_2
 \xrightarrow{(j_2,j_1)}X_3\longrightarrow\Sigma X_0.
\]
The existence of this triangle is \citet[Theorem~6.1, arXiv v2]{GrothPontoShulman2014}; their Theorem~7.1 compares stability with the invertibility of suspension and of the cofiber/fiber adjunction. This gives an established realization of the suspension calculus suggested by the large diagram on p.~269.

For a morphism of two such squares, the induced map of triangles makes the three implications immediate but still instructive. Equivalences on the first and middle terms give one on the third; equivalences on the middle and third give one on the first. Equivalences on the first and third give an equivalence
$f_1\oplus f_2$. Each $f_i$ is a retract of that direct-sum map in the arrow category, so each is an equivalence. This is the additive form of the retraction argument, with cofiber detection now supplied by stability.

Conservativity should nevertheless not be identified with stability. The former is the particular detection property isolated on p.~273; stability also requires a pointed theory and the agreement of cartesian and cocartesian squares. The source's general constructions of coherent squares and iterated cofibers remain useful in unstable theories, while the converse arguments require their stated additional hypotheses.
\end{dcomment}
\begin{otherlanguage}{english}
\dgroverview{XIIc-b}{Chapter XII \textperiodcentered\ \textup{XIIc-b} \textperiodcentered\ Overview}
% Overview by GPT-6 Astra (OpenAI, ChatGPT Pro mode).
\begingroup\selectlanguage{english}
\noindent\textbf{From replacements to a calculus of coherent squares.}
The pages gathered here form a single mathematical movement. Grothendieck first constructs categorical cylinders and uses them to establish the gluing rules for homotopy pushouts. He then asks which part of this geometry can be recovered from a category $\mathcal M$ and a class $W$ of weak equivalences alone. The answer takes the form of a homotopy completion of a cospan, characterized by an adjunction. Finally he returns to pushouts and investigates how far weak equivalences can be recovered from the pieces of a decomposition. The first two movements build the calculus; the appendix tests a considerably stronger conservativity property.

These pages are manuscript pages 228--275, printed pages 145--178 of the edition \citep{Grothendieck1990XIIc}: the second part of section~VI of Chapter~XII, that is, the end of its section~6 followed by sections~7 and~8. The audit covers that entire range, including the fragment of a footnote continued at the beginning. Page labels in the margins reproduce the two numbering systems of the transcription; citations in the commentary use manuscript pages unless another system is specified.

\medskip\noindent\textbf{Why the diagrams must precede localization.}
The motivation is older than the technical apparatus of these pages. In section 69 of \emph{Pursuing Stacks}, Grothendieck traces the need for derivators to the trace formalism of SGA~5: a single derived category did not retain the diagrammatic information needed for the desired compatibilities \citep[\S69, pp.~1--2]{Grothendieck1983}. His letter to Thomason of 2 April 1991 returns to this history and to the system of localized diagram categories \citep[\S2, pp.~5--6]{Grothendieck1991Letter}. These are historical statements checked in those texts; we do not infer that every later definition was already present there in its final form.

Here this motivation becomes a concrete distinction. With the manuscript's contravariant convention, the relevant category is
\[
 \mathbb D(I)=\operatorname{Fun}(I^{\mathrm{op}},\mathcal M)
                  [W(I)^{-1}],
\]
where $W(I)$ consists of the objectwise weak equivalences. It is not obtained by first passing to $\mathcal M[W^{-1}]$ and then forming its ordinary diagram category. A cospan remembers two maps to a common base together; its homotopy completion must preserve that relationship and its naturality in all three vertices. This explains both the attention given to the whole completed square and the two localization lemmas whose proofs the manuscript postpones. The modern introduction in \citet[pp.~313--315]{Groth2013} makes the same structural distinction.

\medskip\noindent\textbf{The constructive part and its later justification.}
The fence interval $\Phi_n$, with $n\geq2$, lets both endpoints enter by a Dwyer immersion. The associated mapping cylinder gives a functorial factorization into a $W$-cofibration followed by a weak equivalence. The commentary proves the endpoint condition and follows the construction on the map from two discrete objects to a point. Taking its categorical cofiber produces an eight-edge circle. Later, the Grothendieck integral of the same span gives a four-edge circle. These are two explicit models of the homotopy pushout; the ordinary categorical pushout is only a point.

The appeal to Thomason is valid in the needed cellular form. It must not be strengthened to the assertion that every Thomason cofibration is a Dwyer map: the retract assertion used to obtain that conclusion is false, and \citet{Cisinski1999Dwyer} gives the counterexample and repair. The recent comparison of Dwyer pushouts with $(\infty,1)$-categorical pushouts is stronger in a different direction, identifying what the nerve preserves rather than adding the appendix's conservative axiom \citep[Theorems~1.5--1.6, arXiv v3]{HackneyOzornovaRiehlRovelli2024}.

In section 7, a $W$-fibration means a map whose pullback functor preserves the designated weak equivalences. The factorization assumption says that every arrow is a weak equivalence followed by such a fibration. It is essential to retain this intrinsic meaning: the manuscript has not assumed a Quillen model structure with arbitrary lifting classes. Under its stated terminal-object and pullback hypotheses, these data form a left derivable category. The commentary verifies the nontrivial axiom, stability of trivial fibrations under arbitrary base change, by a two-step factorization and two-out-of-three argument.

This observation gives a complete modern repair of the deferred Lemmas 1 and 2. We construct the appropriate fibration structures on the good cospans and good squares, check the factorization and pullback conditions, and verify approximation for their inclusions in the ambient diagram categories. Cisinski's Proposition~1.8 also supplies the strengthened replacement assertion for every full intermediate subcategory of a slice containing its fibrant objects; his Theorem~3.12 establishes the two diagram-localization equivalences \citep{Cisinski2010}. No functorial factorization or global assumption that all objects are fibrant is required. A separate, shorter deformation argument explains what functorial factorization would give. The audit still labels the source proofs as deferred or sketched: the later completion does not change what is written on these pages.

With those equivalences available, restriction from squares to cospans has a fully faithful right adjoint. On a cospan $A\to S\leftarrow B$ with one good leg, its missing corner is represented by $A\times_S B$. The groupoid example $*\to BG\leftarrow *$ gives the discrete groupoid $G$ at that corner, whereas the strict pullback without replacement is a point. The adjunction retains the two projections and their compatibility, as well as the corner's homotopy type. Proposition 19 then gives an intrinsic test for cartesianness, and we prove its independence of choices, invariance under objectwise weak equivalences, and pasting directly from the factorization axiom. Strong saturation enters precisely when an isomorphism after localization must be recognized as an original member of $W$.

\medskip\noindent\textbf{What the appendix adds, and where its generalization fails.}
The ordinary gluing implication says that equivalences on the common part and the two pieces induce an equivalence on their homotopy pushouts. The appendix asks for converses. Its full W(7 bis) hypothesis is substantially stronger: a homotopy pushout must also reflect weak equivalences in an attached leg. In a coherent homotopy-pushout theory this is equivalent to detection by the terminal homotopy cofiber
\[
 C_e(u)=e\amalg^{\mathrm h}_{\operatorname{dom}(u)}
                         \operatorname{cod}(u).
\]
The forward implication $u\in W\Rightarrow C_e(u)\simeq e$ is formal. The reverse implication is the additional detection property. We prove the equivalence with conservativity by pasting cofiber squares.

This distinction exposes an actual error in the general theorem on manuscript page 274. Retract closure of aspherical objects alone does not imply that equivalences at the initial and final vertices of a map of homotopy-pushout squares force equivalences at the two middle vertices. Our counterexample uses Hatcher's finite acyclic noncontractible complex $A$. Its suspension is contractible; the map from the square $(A,A,CA,CA)$ to $(A,*,CA,CA/A)$ has three weak-equivalence components and one component $A\to *$ that is not a weak equivalence. Both squares are homotopy pushouts. We give the presentation of $A$, its cellular homology calculation, a nontrivial permutation representation of its fundamental group, and the suspension argument \citep[Example~2.38]{Hatcher2002}. The related equivalence asserted in Corollary 1 therefore fails as printed. Adding cofiber detection repairs both statements, and the final Corollary 2 already has the required conservative hypothesis.

There is also a source slip in Cons~2: the edition prints \emph{cart\'esien} where its argument requires \emph{cocart\'esien}. A finite chain-complex example shows that the literal cartesian assertion fails even in a stable theory. The proof of the intended criterion shows that its cocartesian version follows from Cons~1 under the coherent formulation. In stable derivators the entire corrected cube calculus has a familiar expression through the Mayer--Vietoris triangle \citep[Theorems~6.1 and 7.1, arXiv v2]{GrothPontoShulman2014}.

The appendix's earlier circularity is a different issue. Chapter XVI identifies the dependence on the Dwyer-cofibration gluing lemma but explicitly does not give its proof there \citep[\S3, ms.~pp.~18--23]{GrothendieckDerivateursXVI}. Brown's published gluing lemma, dualized, supplies the missing result \citep[Part~I, \S4, Lemma~2]{Brown1973}. Thus the initial conditional proposition survives, while the later hypothesis-free generalization does not. These are distinct diagnoses, with distinct repairs.

\medskip\noindent\textbf{Editorial and verification limits.}
The text preserves slips printed in the edition. Its one mathematical transcription correction restores the interior natural transformation $\alpha$ on manuscript page 235, which the transcription had drawn as an arrow between the wrong objects. The separate registers record the source's proof status, the present repair, and the evidence for each reading. Bibliography notes identify exactly which statements, proofs, or metadata were checked; in particular, the original proof of Thomason's 1979 theorem was not inspected, although its statement and a later proof were checked.

\endgroup
\end{otherlanguage}
\dgrchunkend

% DER_XIIapp: Grothendieck's text with the commentary set in, then the overview, inside a group that applies the chunk's own definitions
\dgrchapter{}{Chapter XII, appendix}{Appendice au chapitre XII}
\dgrpart{0}{XIIapp}{Manuscript pages 1--18}{XIIapp\enspace\textperiodcentered\enspace Chapter XII, appendix\enspace\textperiodcentered\enspace manuscript pages 1--18}{Chapter XII, appendix \textperiodcentered\ \textup{XIIapp}}
\dgrchunkbegin
% DER_XIIapp: the definitions of the chunk's preamble, applied inside its group
\let\printedpage\relax\newcommand{\printedpage}[1]{\if\relax\detokenize{#1}\relax\else\marginnote{\footnotesize\strut ed.~#1}[-14pt]\fi}
\def\msnumber page #1\relax{#1}
\let\msp\relax\newcommand{\msp}[1]{\marginnote{\footnotesize\strut [ms.~\msnumber#1\relax]}}
\let\Hom\relax\newcommand{\Hom}{\mathrm{Hom}}
\let\Fl\relax\newcommand{\Fl}{\mathrm{Fl}}
\let\Ob\relax\newcommand{\Ob}{\mathrm{Ob}}
\let\Cat\relax\newcommand{\Cat}{\mathrm{Cat}}
\let\Ens\relax\newcommand{\Ens}{\mathrm{Ens}}
\let\pr\relax\newcommand{\pr}{\mathrm{pr}}
\let\Ker\relax\newcommand{\Ker}{\mathrm{Ker}}
\let\Coker\relax\newcommand{\Coker}{\mathrm{Coker}}
\let\mon\relax\newcommand{\mon}{\mathrm{mon}}
\let\deq\relax\newcommand{\deq}{\mathrel{\overset{\text{déf}}{=}}}
\let\id\relax\newcommand{\id}{\mathrm{id}}
\let\Winf\relax\newcommand{\Winf}{\mathbf{W}_{\infty}}
\let\dotto\relax\newcommand{\dotto}{\mathrel{\tikz[baseline=-0.5ex]\draw[dotted,-latex] (0,0)--(1.6em,0);}}

% ---- from Der-ChXIIapp-19-06-03_p001-005.tex
\printedpage{}
% no page number printed (title page)
% head: (none)

\begin{center}
{\LARGE\textsf{LES DÉRIVATEURS}}

\bigskip

\textbf{ALEXANDRE GROTHENDIECK}

\bigskip\bigskip

{\large\textbf{Appendice au Chapitre XII}}
\end{center}

\newpage

\printedpage{1}
% head: A. Grothendieck : Dérivateurs, Ch. XII, Édition des Universités de Montpellier II et Paris VII
% foot: Version 20 juin 2003

\msp{page 1}
\textbf{Proposition} (?) (\footnote{Cf. pages 11 ff. pour les contrexemples-clef.}). \textit{Soit $X$ dans} $\Cat$\textit{, avec $b$, $c$ dans $X$ tels que}

\begin{itemize}
\item[\textit{a)}] \textit{$b\times c$ est dans $X$.}
\item[\textit{b)}] \textit{$\forall$ $x$ dans $X$, $b(x)$ ($\deq$} $\Hom_X(x,b)$\textit{) et $c(x)$ sont $\neq\emptyset$, i.e. $X/b \longrightarrow X$ et $X/c \longrightarrow X$ sont surjectifs sur les objets.}
\end{itemize}

\textit{Alors $X$ est $\mathbf{W}_\infty$-asphérique.} (\footnote{Faux, cf. plus bas. $X$ est 1-connexe, mais pas nécessairement 2-connexe.})

Le fait que $b$ soit un plus grand élément de $\Ob X$ assure déjà la 0-connexité.

\dgrblock{dgrb607}{What an ordinary pushout forgets}
\begin{dcomment}{context}{p1:Le fait que $b$ soit un plus grand élément}{What an ordinary pushout forgets}
Grothendieck is looking for categorical gluings that preserve a homotopy type. The little span that appears throughout this appendix is a particularly severe test. Its indexing category
\[
  \Psi=(b_0\longleftarrow a_0\longrightarrow c_0)
\]
has an initial object, hence a contractible nerve. Nevertheless, taking the ordinary colimit of its representable presheaves can introduce higher homotopy into the category of elements. The purpose of the examples is to locate this discrepancy precisely enough to decide which inclusions can be used in a homotopy theory on \(\Cat\).

We write \(BC=|NC|\), and interpret \(\Winf\) as the functors inducing weak homotopy equivalences on these classifying spaces. In this terminology ``aspherical'' means weakly contractible, not merely a space with vanishing higher homotopy groups. For a presheaf \(F:C^{\mathrm{op}}\to\Ens\), the notation \(C/F\) denotes its category of elements: its objects are \((c,x)\), and an arrow \((c,x)\to(c',x')\) is an arrow \(u:c\to c'\) with \(F(u)(x')=x\). Thus a represented presheaf has contractible category of elements \(C/c\). The central phenomenon is that a colimit of represented presheaves need not retain that property.

This is already the tension described in \citet[\S69]{Grothendieck1983}: categorical integration gives homotopy colimits, whereas actual amalgamated sums require additional care. The letter to Thomason of 2 April 1991 explains why the two adjoints to restriction, together with their exactness properties, should be treated as the essential structure \citep[pp.~8--9]{Grothendieck1991Letter}. The appendix examines the small categorical configurations behind that program. Its final gluing theorem belongs to the theory of homotopy colimits; its earlier counterexamples concern ordinary pushouts in \(\Cat\). Keeping the two operations distinct makes the apparent reversal between the examples and the closing conjecture intelligible.

There are also two different uses of asphericity for an inclusion. A functor \(U\to X\) is aspherical in the usual sense if every comma category \(U/x\) is weakly contractible. An open immersion in these pages is a full sieve, and the local condition used for it is imposed on \(U\to\overline U\), where \(\overline U\) is its cosieve closure. Objects outside \(\overline U\) have empty comma categories. This distinction will matter for the boundary of a subdivided square: that boundary has the homotopy type of a circle, even though the indicated local condition holds. We refer to the precise conventions of Chapter XII, \S V.2 \citep{Grothendieck1990XII}, and to the commentary on Chapter~II, manuscript p.~36 (``The quantifier in total asphericity'') for the universal quantifier in that second notion.
\end{dcomment}


\dgrblock{dgrb608}{What the product actually contracts}
\begin{dcomment}{exposition}{p1:Le fait que $b$ soit un plus grand élément}{What the product actually contracts}
Grothendieck begins by asking how much of a terminal object can be replaced by the existence of maps. The objects $b$ and $c$ are \emph{weakly terminal}: every object admits a map to each of them, without any assertion of uniqueness. Their product $a=b\times c$ is weakly terminal as well. The argument establishes the exact conclusion needed at the first stage: the nerve $NX$ is connected and simply connected. Here and below $BX=|NX|$, and ``aspherical'' in the terminology of these pages means weakly contractible, not merely a $K(\pi,1)$.

The cancellation argument can be compressed without losing its content. For two maps $s,t:x\to a$, insert the mixed pair $(p_bs,p_ct)$. Cancellation against the now invertible projections shows that all three maps have the same image in the groupoid completion $\Pi X$. Choose one such image $\alpha_x:x\to a$ for each object, taking $\alpha_a=1_a$. For every arrow $u:x\to y$, both $\alpha_yu$ and $\alpha_x$ come from arrows $x\to a$ and therefore agree. The same identity then holds for inverse arrows and their composites. Every morphism $x\to a$ in $\Pi X$ is consequently $\alpha_x$. Thus $a$ is terminal in this groupoid, which is equivalent to the terminal groupoid. The longer zigzag induction in the text is a complete proof of this fact; its undefined $a_0$ and $u_0$ are printing slips, not a missing argument.

There is a significant distinction between this proposition and the later four-object example. In that example $b\amalg_a c=1$ in the presheaf category, but $a$ is \emph{not} $b\times c$: at the added object there are two maps to $a$ and only one pair of maps to $(b,c)$. It cannot by itself refute the opening product proposition. The finite two-object category below does refute that proposition, although its first obstruction occurs in degree three. We give the calculation because the provisional degree-two calculation in the manuscript does not survive verification.
\end{dcomment}


Pour la 1-connexité, prouvons que dans le groupoïde fondamental $\Pi(X)$, $a \deq b\times c$ devient un élément final.
\[
\begin{tikzcd}
b\times c \arrow[r] \arrow[d] & b \\
c &
\end{tikzcd}
\]

Pour ceci, prouvons d'abord que $\forall$ $x$ dans $X$, tous les éléments de $a(x) = b(x)\times c(x)$ sont égaux dans $\Fl\Pi(X)$. En effet, comme $b\times c \longrightarrow b$ est \texttt{[un]} isomorphisme, donc \texttt{[un]} monomorphisme dans $\Pi(X)$, cela montre déjà que si $\beta,\beta'\in b(x)$, $\gamma,\gamma'\in c(x)$, d'où $(\beta,\gamma)$, $(\beta,\gamma')$, $(\beta',\gamma')\in b(x)\times c(x)$, on a $(\beta,\gamma) = (\beta,\gamma')$ dans $\Pi(X)$ (puisque $\pr_1(\beta,\gamma) = \pr_1(\beta,\gamma')$, et de même (en utilisant que $\pr_2$ est \texttt{[un]} isomorphisme, donc \texttt{[un]} monomorphisme) $(\beta,\gamma') = (\beta',\gamma')$). Donc tous les $(\beta,\gamma)\in b(x)\times c(x) = a(x)$ sont égaux dans $\Fl\Pi(X)$. Pour achever de prouver que $a$ est objet final de $\Pi(X)$, il suffit de voir que toute flèche \textit{dans} $\Pi(X)$ d'un $x\in\Ob X$ dans $a$, provient d'une flèche $x\longrightarrow a$ dans $X$. Mais une flèche dans $\Pi(X)$ est obtenue par un chemin dans $X$
\[
\begin{tikzcd}
x = x_0 \arrow[r, "u_1"] \arrow[dr, "\alpha_0"] & x_1 \arrow[d, "\alpha_1"'] & x_2 \arrow[l, "u_2"'] \arrow[r] \arrow[dl, "\alpha_2"'] & {} & \cdots & {} & x_{2N} = a \arrow[l, "u_{2N}"'] \\
& a. & & & & &
\end{tikzcd}
\]

Prenons une flèche $\alpha_0 : x\longrightarrow a_0$ dans $X$, je dis que la % [sic] a_0 -> a, the arrows α_i go to a (no a_0 is defined)
\msp{page 2}
classe d'homotopie du chemin $c$ (i.e. son image dans $\Fl(\Pi(X))$ est égale à celle de $u_0$. % [sic] u_0 -> \alpha_0 (?), no u_0 is defined; also the parenthesis opened at "(i.e." is not closed
Pour le voir, choisissons pour tout $i\in[0,2N]$ une flèche $\alpha_i : x_i\longrightarrow a$. (Ça existe grâce à l'hypothèse b)). Je dis que pour chaque $i\in[0,2N]$, le triangle
\printedpage{2}
% head: A. Grothendieck : Dérivateurs, Ch. XII, Édition des Universités de Montpellier II et Paris VII
% foot: Version 20 juin 2003
\[
\begin{tikzcd}[column sep=small]
x_0 \arrow[rr, "c_i"] \arrow[dr, "\alpha_0"] & & x_i \arrow[dl, "\alpha_i"'] \\
& a &
\end{tikzcd}
\]
dans $\Pi(X)$ est commutatif, où $c_i$ désigne la classe d'homotopie du chemin de $x_0$ à $x_i$ définie par les $u_i$ ($i\in[1,i]$) \texttt{[plutôt} ($i\in[1,2N]$)\texttt{]}. Nous le faisons par récurrence sur $i$, l'assertion est triviale si $i = 0$ ($c_0$ représente l'application identitique). Le diagramme % [sic] identitique -> identique, non-existent word
\[
\begin{tikzcd}[column sep=small]
x_0 \arrow[rr, "c_i"] \arrow[dr, "\alpha_0"] \arrow[rrrr, bend left=32, "c_{i+1}"] & & x_i \arrow[dl, "\alpha_i"'] \arrow[rr, "u_{i+1}"] & & x_{i+1} \arrow[dlll, "\alpha_{i+1}"'] \\ % verified: solid curved arrow; bend increased for label clearance
& a & & &
\end{tikzcd}
\]
montre qu'il suffit de prouver que pour $i\in[0,2N-1]$ le triangle
\[
\begin{tikzcd}[column sep=small]
x_i \arrow[rr, "u_{i+1}"] \arrow[dr, "\alpha_i"'] & & x_{i+1} \arrow[dl, "\alpha_{i+1}"] \\
& a &
\end{tikzcd}
\]
est commutatif dans $\Pi(X)$. (\textbf{NB} $u_{i+1}$ est en sens directe ou rétrograde, suivant que % [sic] directe -> direct, agreement with "sens"
$i$ est paire ou impaire.) Mais cela résulte du fait que l'image de $\Hom_X(x_i,a)$ (resp. $\Hom_X(x_{i+1},a)$) % [sic] paire ou impaire -> pair ou impair, agreement with "i"
\msp{page 3}
dans $\Fl\Pi(X)$ est réduite à un seul élément.

Ainsi on trouve que $X$ \textit{est} 1-\textit{connexe.} Il y a donc de l'espoir! Il reste à prouver
\[
H^i(X,M) = 0, \qquad i\geq 2, \quad M \text{ coefficient constant.}
\]

J'ai regardé le cas lorsqu'on a ceci : $b = c$, $\Ob X = \{b, b\times b\}$, $\Hom_X(b\times b)$ \texttt{[plutôt} $\Hom(b\times b, b)$\texttt{]} $= \{$ réunion disjointe des ensembles $u\circ\pr_1$ et $u\circ\pr_2$, avec $u\in\Hom_X(b,b) = G$ un monoïde quelconque $\}$. Alors $X^\circ$ est fibrée, via $I\longmapsto b^I$, sur la sous-catégorie pleine $S$ de $\Ens$ formée de $\{0\}$ et $\{0,1\}$, et donc $X$ est cofibrée (en particulier propre) sur $S^\circ \deq T$. (NB $S^\circ = T$ a un objet initial, mais non un objet final. Sa cohomologie à valeurs dans un préfaisceau quelconque n'est pas toujours triviale, cf. ci-dessous.)

On trouve \texttt{[`On trouve' barré]}
\[
H^n(X,M) \Longleftarrow H^p(T, \underbrace{R^q f_*(M_X)}_{= \mathcal{H}^q}), \qquad \text{où } f : X \longrightarrow S^\circ
\]
\texttt{[}$n = p+q$\texttt{]}. J'admet ici (provisoirement) ceci : % [sic] J'admet -> J'admets, verb ending
\[
\begin{aligned}
H^p(T,F) &= 0 \ , &&\text{pour } p\geq 2, \text{ tout } F\in T^\wedge \ ,\\
H^0(T,F) &= \Ker\delta \ , &&\\
H^2(T,F) &= \Coker\delta \ , && % [sic] H^2(T,F) -> H^1(T,F), the first line gives H^p(T,F) = 0 for p >= 2
\end{aligned}
\]
\printedpage{3}
% head: A. Grothendieck : Dérivateurs, Ch. XII, Édition des Universités de Montpellier II et Paris VII
% foot: Version 20 juin 2003
où $\delta : F_0 \longrightarrow F_1$ est défini comme la différence $\delta = p_1^* - p_2^*$, où $p_1$, $p_2$ sont les deux projections de $\{0,1\}\simeq\{0\}^2$ dans $T = S^\circ$. On a de plus
\[
(\mathcal{H}^q)_0 = H^q(G,M), \qquad (\mathcal{H}^q)_1 = H^q(G\times G, M),
\]
et
\[
p_1^*,\ p_2^*\ :\ H^q(G,M) \rightrightarrows H^q(G{\times}G, M) \qquad \text{sont les transposés de } \pr_1,\ \pr_2 : G\times G\longrightarrow G.
\]

\dgrblock{dgrb609}{The finite-power category and the missing symmetries}
\begin{dcomment}{exposition}{p3:J'admet ici (provisoirement) ceci}{The finite-power category and the missing symmetries}
Grothendieck now makes the product condition explicit in a category small enough to calculate. Write $X_G$ for the category with objects $1,2$, representing $b,b^2$, and
\[
 \Hom_{X_G}(m,n)=\bigl(\{1,\ldots,m\}\times G\bigr)^n.
\]
A morphism lists, for each output coordinate, one input coordinate and one element of the monoid $G$; composition substitutes these expressions. Thus $(i,g):b^m\to b$ means $g\circ p_i$. This convention fixes multiplication even when $G$ is not commutative. The evident maps $2\to1$ exhibit $2$ as the product of $1$ with itself. Both objects map to $1$, so the hypotheses of the opening proposition hold. For a group $G$, the opposite category is the full category of free $G$-sets of ranks one and two.

Let $S$ be the full subcategory of finite sets on $\mathbf1,\mathbf2$, and put $T=S^{\mathrm{op}}$. The functor $X_G\to T$ forgets the monoid labels. A map $v:m\to n$ of $T$, equivalently a function $n\to m$, induces the coordinate-selection homomorphism $G^m\to G^n$. The arrows with all labels equal to $1_G$ are chosen cocartesian lifts, so $X_G$ is the Grothendieck construction of $n\mapsto BG^n$, where here $BG^n$ denotes the one-object category. This proves the asserted opfibration and justifies the cohomological spectral sequence, with constant abelian coefficients $M$:
\[
 E_2^{p,q}=R^p\!\varprojlim_{S} F_q
 \ \Longrightarrow\ H^{p+q}(BX_G;M),
 \qquad F_q(n)=H^q(BG^n;M).
\]
The functor $F_q$ is covariant on $S$. A modern treatment of this spectral sequence for Grothendieck constructions is \citet[Corollaries 3--4]{Cegarra2020}; the underlying homotopy-colimit comparison is \citet[Theorem 1.2]{Thomason1979}.

The provisional formula is already contradicted by a constant diagram $F=M$. The category $T$ has an initial object, hence contractible nerve, and $H^{>0}(T,M)=0$. Yet the proposed differential is $1_M-1_M=0$, whose cokernel is $M$. Changing the printed $H^2$ to the evidently intended $H^1$ does not repair this. The full category $S$ contains two inclusions, their common retraction, and the transposition of $\mathbf2$; retaining only a parallel pair loses relations and higher derived limits. On the other hand the formula for $H^0$ is correct: a compatible pair is determined by its value at $\mathbf1$, and the equality of the two induced values at $\mathbf2$ implies compatibility with the retraction and transposition.
\end{dcomment}


\dgrblock{dgrb610}{A small resolution for the actual base category}
\begin{dcomment}{development}{p3:sont les transposés de}{A small resolution for the actual base category}
The point of repairing the base calculation is to retain Grothendieck's finite example, while locating its obstruction in the correct degree. The following elementary computation supplies all the derived limits used below. We did not find this calculation for the appendix's two-object category in the literature checked; the argument is included in full, without a claim of priority.

\textbf{Derived-limit calculation.}\ours
Let $F:S\to\mathrm{Ab}$, let $i_1,i_2:\mathbf1\to\mathbf2$ be the inclusions, let $r:\mathbf2\to\mathbf1$ be the unique map, and let $\tau$ transpose $\mathbf2$. Put
\[
 A=F(\mathbf1),\quad B=F(\mathbf2),\quad
 L=\ker(F(r):B\to A),\quad
 L^- =\{v\in L:F(\tau)v=-v\}.
\]
Then
\[
 \begin{aligned}
 R^0\!\varprojlim_S F&=\ker(F(i_1)-F(i_2):A\to B),\\
 R^1\!\varprojlim_S F&=L^-/(F(i_1)-F(i_2))A,\\
 R^p\!\varprojlim_S F&=\operatorname{Ext}^{p-1}_{\mathbb Z[C_2]}
       (\mathbb Z_{\mathrm{sgn}},L)\qquad(p\ge2).
 \end{aligned}
\]
The action of $C_2$ on $L$ is induced by $F(\tau)$.

\emph{Proof.} Work in the abelian category of covariant functors on $S$. The functors $P_j=\mathbb Z[S(\mathbf j,-)]$ are projective, since $\operatorname{Hom}(P_j,F)=F(\mathbf j)$ and evaluation is exact. The augmentation $P_1\to\underline{\mathbb Z}$ gives
\[
 0\longrightarrow K\longrightarrow P_1
   \longrightarrow\underline{\mathbb Z}\longrightarrow0.
\]
Here $K(\mathbf1)=0$, and $K(\mathbf2)=\mathbb Z(i_1-i_2)$ is the sign representation; the two constant endomorphisms act as zero. This immediately gives the formula for degree zero and identifies the degree-one group as the stated cokernel, because a map $K\to F$ is precisely an element of $L^-$.

For completeness the higher Ext groups can be computed inside the same tiny category. Precomposition with $r$ gives a monomorphism $P_1\to P_2$. It has the natural retraction induced by $i_1$, because $ri_1=1$. Consequently $Q=P_2/P_1$ is projective. It is zero at $\mathbf1$ and is the regular representation $\mathbb Z[C_2]$ at $\mathbf2$, with basis the classes of $1$ and $\tau$; both constant maps act as zero. Any free $\mathbb Z[C_2]$-resolution of the sign representation therefore gives a projective resolution of $K$, by replacing the regular module with $Q$. Applying $\operatorname{Hom}(-,F)$ gives exactly the complex computing $\operatorname{Ext}^*_{\mathbb Z[C_2]}(\mathbb Z_{\mathrm{sgn}},L)$: maps out of $Q$ are elements annihilated by $F(r)$. Finally dimension shifting through the displayed short exact sequence proves the formula for $p\ge2$. This also proves that no omitted higher relation is being assumed away.\hfill$\square$

For example, take $F(\mathbf1)=0$, $F(\mathbf2)=\mathbb Z$ with $\tau$ acting as $-1$. The formula gives $R^3\!\varprojlim F\cong\mathbb Z/2$, from the usual alternating resolution of $C_2$. Thus even the claimed vanishing in all degrees at least two is false for the full base category.
\end{dcomment}


\msp{page 4}
On vérifie aisément que
\[
H^0(T,\mathcal{H}^q) = 0 \qquad \text{si } q\geq 1,
\]
et il reste
\[
H^n(X,M) \simeq H^1(S,\mathcal{H}^{n-1}) = \Coker(\delta : H^{n-1}(G) \longrightarrow H^{n-1}(G\times G)). % [sic] H^1(S, ...) -> H^1(T, ...), the cohomology is taken on T = S°
\]
Cela donne, pour $n = 2$,
\[
\begin{aligned}
H^2(X,M) &\simeq \Coker\Big(\delta : \Hom_\mon(G,M) \longrightarrow \underbrace{\mathit{Hom}_\mon(G\times G, M)}_{\Hom_\mon(G,M)\times\Hom_\mon(G,M)} \ \Big)\\
&\simeq \Coker\Big(u\longmapsto (u,-u) : H \longrightarrow H\times H\Big),
\end{aligned}
\]
où $H$ est le groupe abélien $\Hom_\mon(G,M)$.

D'où
\[
H^2(X,M) \ \simeq\ H = \Hom_\mon(G,M).
\]
Donc
\begin{quote}
$X$ est 2-connexe si et seulement si le groupe \textit{abélien} enveloppant du monoïde $G$ est nul.
\end{quote}

Ainsi, si $G$ est un groupe abélien $\neq 0$, $X$ \textit{n'est pas} 2-connexe.

\dgrblock{dgrb611}{The counterexample survives, one degree later}
\begin{dcomment}{development}{p4:Ainsi, si $G$ est un groupe abélien}{The counterexample survives, one degree later}
Grothendieck's proposed category is indeed nonaspherical for $G=C_2$. Its degree-two cohomology, however, is zero. The distinction matters twice: it changes the audit of the first four pages, and it supplies a sound finite example for the pushouts that follow.

\textbf{The actual low-degree calculation.}\ours
For every monoid $G$, the space $BX_G$ is $2$-connected. For $G=C_2$ one moreover has
\[
 H^3(BX_{C_2};\mathbb F_2)\cong\mathbb F_2.
\]
In particular the opening weak-contractibility assertion is false under its exact product hypotheses.

\emph{Proof.} The groupoid argument already proves simple connectivity. Use the spectral sequence and notation of the preceding two comments. The $q=0$ row is the constant diagram $M$, so its terms vanish for $p>0$. For $q>0$ the two coordinate pullbacks
\[
 H^q(BG;M)\rightrightarrows H^q(BG^2;M)
\]
have zero equalizer: restriction along $g\mapsto(g,1)$ sends the first pullback to the identity and the second to zero. Thus $E_2^{0,q}=0$ for $q>0$.

Set $A=H^1(BG;M)=\Hom_{\mon}(G,M)$, with $M$ regarded as an additive group. A normalized degree-one cocycle on the monoid is exactly an additive homomorphism, which proves this identification directly. A homomorphism $G^2\to M$ is uniquely the sum of its restrictions to the two factors, so the degree-one diagram is
\[
 F_1(\mathbf1)=A,\qquad F_1(\mathbf2)=A\oplus A,
\]
with $i_1,i_2$ inducing the two inclusions, $r$ inducing addition, and $\tau$ interchanging the factors. Hence $L=\{(a,-a):a\in A\}$, $L^-=L$, and the map $A\to L$ in the degree-one formula is an isomorphism. Thus $E_2^{1,1}=0$. All three entries of total degree two vanish; consequently $H^2(BX_G;M)=0$ for every abelian $M$. Simple connectivity and the universal coefficient theorem give $H^2(BX_G;M)=\Hom(H_2(BX_G;\mathbb Z),M)$. Taking $M=H_2(BX_G;\mathbb Z)$ forces this group to be zero. The Hurewicz theorem now gives $\pi_2(BX_G)=0$; see \citet[Theorems 3.2 and 4.32]{Hatcher2002}.

Now put $G=C_2$ and $M=\mathbb F_2=k$. The same degree-one diagram has $L=k$ with trivial $C_2$-action. The base calculation gives
\[
 E_2^{2,1}=\operatorname{Ext}^1_{\mathbb Z[C_2]}(\mathbb Z_{\mathrm{sgn}},k)=k.
\]
Indeed the free resolution of the sign module has differentials alternating between $1+\tau$ and $1-\tau$, both zero after applying $\operatorname{Hom}(-,k)$. No differential can kill this term: its sole possible incoming differential is from $E_2^{0,2}=0$ and its sole possible outgoing differential lands in $E_2^{4,0}=0$.

To identify the whole degree-three group, the only remaining possible entry is $E_2^{1,2}$. Write $H^*(BC_2;k)=k[t]$ and $H^*(B(C_2^2);k)=k[x,y]$, with all three generators in degree one; these are the usual projective-space and K\"unneth calculations \citep[Theorems 3.15 and 3.19]{Hatcher2002}. In degree two, $F_2(r)$ sends each of $x^2,xy,y^2$ to $t^2$. The $\tau$-fixed part of its kernel is therefore $k(x^2+y^2)$, exactly the image of $F_2(i_1)-F_2(i_2)$. The degree-one formula gives $E_2^{1,2}=0$. All other entries of total degree three vanish, proving the displayed isomorphism.\hfill$\square$

This proof also explains the discrepancy in the manuscript. The abstract cokernel of $a\mapsto(a,-a)$ is indeed $A$, but it is not the derived limit being computed. The transposition and the common retraction replace that cokernel by the preceding resolution. For $G=C_2$ the category has two objects and $2+4+4+16=26$ arrows. It supplies the finite counterexample promised in the text, with a degree-three obstruction. The stronger assertion in footnote~2 that the general product hypotheses need not imply $2$-connectivity requires a different example. The next comment verifies it with an actual product and a two-sphere.
\end{dcomment}


\dgrblock{dgrb612}{A product whose category is a two-sphere}
\begin{dcomment}{development}{p4:Ainsi, si $G$ est un groupe abélien}{A product whose category is a two-sphere}
The correction of the monoid calculation leaves one assertion in footnote~2 to check: the exact product hypotheses may fail to force $2$-connectivity. That assertion is true. One extra object supplies a smaller example in object count, and reveals geometrically the second homotopy group that the proposed monoid calculation did not produce. The construction below is our synthesis of the product condition and the elementary nerve-of-a-cone calculation; we did not locate this particular finite example in the literature searched, and make no claim of priority.

\textbf{A family realizing wedges of two-spheres.}\ours
For every integer $r\ge1$ there is a category $Y_r$ with three objects $d,b,a$, where $b$ is weakly terminal and $a=b\times b$, such that
\[
 BY_r\simeq\bigvee^{(r-1)(r-2)/2}S^2.
\]
The wedge with zero summands means a point. In particular $Y_3$ has $21$ arrows, satisfies the opening hypotheses with $c=b$, and has the homotopy type of $S^2$.

\emph{Proof.} Let $D$ be an $r$-element set. On the two objects $b=b^1$ and $a=b^2$ use the category $T=S^{\mathrm{op}}$ already encountered, with
\[
 \Hom_T(b^m,b^n)=\{1,\ldots,m\}^{\{1,\ldots,n\}},
 \qquad m,n\in\{1,2\}.
\]
These are coordinate-selection maps, composed by substitution. Add the object $d$, its identity, the maps
\[
 \Hom(d,b)=D,\qquad \Hom(d,a)=D^2,
\]
and no maps from $a$ or $b$ to $d$. Composition with old arrows again selects coordinates. Substitution is associative, so this defines a category. The product bijection
\[
 \Hom(z,a)\xrightarrow{\ (p_1,p_2)\ }
 \Hom(z,b)\times\Hom(z,b)
\]
holds for $z=a,b$ by the definition of $T$, and for $z=d$ it is the identity on $D^2$. Hence $a=b\times b$ in the enlarged category. As $D$ is nonempty, all three objects map to $b$.

To see the homotopy type, consider $E=d\backslash T$. An object of $E$ is a list of length one or two with entries in $D$; a morphism selects coordinates to obtain its target list. Let $P$ be the poset of nonempty subsets of $D$ of size at most two, ordered by reverse inclusion. Sending a list to its support defines $s:E\to P$. Choose one ordering of $D$ and let $j:P\to E$ list each subset without repetition in that order. Coordinate selection makes $j$ a functor and $sj=1_P$. There is a natural transformation $js\to1_E$: from the distinct list of a support to any given list, repeat the indicated coordinates. This map is unique because each label occurs once in its source. Its uniqueness also proves naturality for every coordinate-selection arrow. Therefore $BE\simeq BP$.

For clarity, the last gluing step can be read directly on the nerve. Every nondegenerate simplex of $NY_r$ is either wholly in $NT$, or starts at $d$ and continues in $T$. A simplex of the second kind is precisely the cone on a simplex of $N(d\backslash T)$. Its face opposite $d$ is obtained by forgetting the initial arrow from $d$. Thus $NY_r$ is the simplicial mapping cone of $NE\to NT$. After realization it is the mapping cone of $BE\to BT$. The category $T$ has initial object $b$, so $BT$ is contractible and
\[
 BY_r\simeq\Sigma BE\simeq\Sigma BP.
\]
The space $BP$ is the barycentric subdivision of the complete graph on the $r$ vertices of $D$. For $r\ge1$ that graph is connected and has $\binom r2-r+1=(r-1)(r-2)/2$ independent cycles: collapse a spanning tree to obtain the wedge of that many circles. Suspending gives the asserted wedge of two-spheres. Finally $T$ has $8$ arrows; adding $d$ contributes $1+r+r^2$ arrows. For $r=3$ the total is $21$.\hfill$\square$

This example verifies footnote~2 under the actual product assumptions, including its failure of $2$-connectivity. It also keeps the two repairs distinct. The two-object category $X_{C_2}$ first detects nonasphericity in degree three; the three-object category $Y_3$ detects it in degree two. The four-object category considered next in the manuscript has the latter homotopy type but satisfies only the terminal-presheaf-span condition.
\end{dcomment}


\textbf{Corollaire.} $\Psi = {}$\begin{tikzcd}[ampersand replacement=\&, row sep=tiny, column sep=small, baseline=(A.base)]
\& b_0 \\
|[alias=A]| a_0 \arrow[ur] \arrow[dr] \& \\
\& c_0
\end{tikzcd}
\textit{n'est} pas \textit{totalement} $\mathbf{W}_\infty$\textit{-asphérique; i.e. si on a un diagramme} \begin{tikzcd}[ampersand replacement=\&, row sep=tiny, column sep=small, baseline=(B.base)]
\& b \\
|[alias=B]| a \arrow[ur] \arrow[dr] \& \\
\& c
\end{tikzcd}
\textit{dans une catégorie} $X\in\Ob\Cat$\textit{, tel que sa somme amalgamée dans $X^\wedge$ soit égale à $e_X$, il ne s'ensuit pas que $X$ soit $\mathbf{W}_\infty$-asphérique.} (Mais peut-être voudra bien être 1-connexe?)

\printedpage{4}
% head: A. Grothendieck : Dérivateurs, Ch. XII, Édition des Universités de Montpellier II et Paris VII
% foot: Version 20 juin 2003

\msp{page 5}
\textbf{Question.} \textit{Y-a-t-il d'autres objets totalement} $\mathbf{W}_\infty$\textit{-asphériques de} $\Cat$\textit{, que les catégories ayant un objet final?} Je suspecte à présent que non -- donc que l'on ne peut améliorer substantiellement sur la notion déjà connue d'immersion ouverte de Dwyer.

\dgrblock{dgrb613}{Total asphericity without a terminal object}
\begin{dcomment}{status}{p5:Je suspecte à présent que non}{Total asphericity without a terminal object}
The question concerns how far the useful gluing condition extends beyond representability. In the sense used here, a category \(A\) is totally \(\Winf\)-aspherical when, for every functor \(f:A\to C\), the presheaf \(f_!1_A\) has weakly contractible category of elements. This is the universal left-Kan-extension condition of Chapter XII, ms.~88--90; the editors explicitly distinguish it from the total asphericity of \emph{Pursuing Stacks}. The suspicion that a terminal object is necessary is false. Filtered categories provide counterexamples; indeed Grothendieck records this correction himself in Chapter XII, ms.~109, footnote~78 \citep{Grothendieck1990XII}. Chapter II, Corollaries~2.17--2.18 (ms.~34--35), gives the underlying filtered argument \citep{Grothendieck1990II}. We recall its mechanism because it answers the question on this page, while the commentary on Chapter~II, manuscript pp.~35--36, already discusses its broader setting.

For any small \(A\), the canonical presentation of the terminal presheaf gives
\[
  f_!1_A\cong\mathop{\mathrm{colim}}_{a\in A}C(-,f(a)).
\]
Suppose that \(A\) is filtered, and denote this presheaf by \(F\). We verify directly that \(C/F\) is filtered. It is nonempty, since the identity of any \(f(a)\) supplies an object. Given \((c,x)\) and \((d,y)\), choose representatives \(u:c\to f(a)\) and \(v:d\to f(b)\). A common successor \(k\) of \(a,b\) produces arrows from both objects to \((f(k),[\id_{f(k)}])\). Finally, let \(r,s:(c,x)\rightrightarrows(d,y)\) be parallel arrows. Represent \(y\) by \(v:d\to f(b)\). The equality \(F(r)y=F(s)y\) says that \(vr\) and \(vs\) become equal at some later stage of the filtered colimit. Thus there is an arrow \(b\to k\) such that its image, composed with \(v\), equalizes \(r\) and \(s\). This is an equalizer after postcomposition inside \(C/F\).

The nerve of a filtered category is weakly contractible. One way to see this is to take the finitely many objects and arrows occurring in a simplicial representative of a homotopy class, choose a common target, and then equalize the finitely many competing composites. The resulting compatible cone gives a null homotopy. Equivalently, the usual finite-diagram characterization of filteredness cones off every finite simplicial diagram. Consequently \(C/F\) is weakly contractible for every \(f\), as required. The elementary eventual-equality fact used above is the standard filtered-colimit argument \citep[\S4.19]{StacksFiltered}.

For a concrete infinite example take \(A=(\mathbb N,\leq)\). It has no greatest element and hence no terminal object. The phenomenon is not a consequence of infinitude: the one-object category with endomorphism monoid \(\{1,e\}\), \(e^2=e\), is filtered, since postcomposition by \(e\) equalizes any pair of endomorphisms. Its unique object is not terminal, because it has two endomorphisms. For a functor out of this category, \(f_!1_A\) is the image of the induced idempotent on the represented presheaf \(C(-,f(*))\). Taking \(C=A\) and \(f=\id_A\) makes this image the terminal presheaf, with one element at the unique object, whereas the sole representable has two elements there. Thus the gluing weight need not be representable. These examples settle the question as asked; they do not assert an exhaustive classification of totally aspherical categories.
\end{dcomment}


Dans le diagramme
\[
\begin{tikzcd}
U \arrow[r, hook, "i"] \arrow[d, "p"'] \arrow[dr, phantom, "\substack{\text{cocart.}\\ \text{et cart.}}"] & X \arrow[d, "q"'] & T \arrow[l, hook', dotted, "j"'] \arrow[d, dotted, "\text{iso}"] \\
U' \arrow[r, hook, "i'"] & X' & T', \arrow[l, hook', dotted, "j'"']
\end{tikzcd}
\tag*{$(*)$}
\]
prenons
\[
\begin{aligned}
U = \Psi, \quad & T = \{t\} = e \text{ catégorie finale,}\\
& t \text{ objet final de } X,\\
& \text{i.e. } H_t = (u\longmapsto \Hom_X(u,t)) = e_U \in \Ob U^\wedge \ \mbox{(\footnotemark)}.
\end{aligned}
\]
\footnotetext{i.e. $X$ déduit de $U = \Psi$ en lui rajoutant un objet final $t$ -- donc $X\simeq\boldsymbol{\Delta}^1\times\boldsymbol{\Delta}^1$,
\[
\begin{tikzcd}[ampersand replacement=\&,
execute at end picture={
\draw[rounded corners] ([xshift=-10pt,yshift=10pt]A.center) -- ([xshift=-10pt,yshift=-10pt]C.center) -- ([xshift=10pt,yshift=-10pt]D.center) -- ([xshift=10pt,yshift=3pt]D.center) -- ([xshift=3pt,yshift=3pt]C.center) -- ([xshift=3pt,yshift=10pt]A.center) -- cycle;
\node (P) at ([xshift=-35pt,yshift=-22pt]C.center) {$\Psi$};
\draw[-stealth] (P) -- ([xshift=-11pt,yshift=-4pt]C.center);
}]
|[alias=A]| {} \arrow[r] \& {} \\
|[alias=C]| {} \arrow[u] \arrow[r] \& |[alias=D]| {} \arrow[u]
\end{tikzcd}
\ .
\]}
Prenons pour $U'$ la catégorie dont il est question dans la `proposition' hypothétique (et fausse) p. 1 avec $b = c$, plus précisément une catégorie avec deux objets $b$ et $b\times b$, un groupe (voire un monoïde) $G$ d'automorphismes \texttt{[(voire d'endomorphismes)]} de $b$, et les autres flèches celles qui s'en déduisent, en exigeant que $b\times b$ soit bel et bien le carré cartésien de $b$. On trouve une catégorie cofibrée (donc propre) sur la catégorie $T$ opposée de celle des ensembles finis de la forme $\{0\}\cup\{0,1\}$ (deux objets) avec comme fibres $B_G$ et $B_{G\times G}$ respectivement ($B_\pi =$ catégorie avec un seul objet, et monoïde d'endomorphismes $\pi$). Le foncteur
\[
U = \Psi \longrightarrow U'
\]
est celui donné par le diagramme

\msp{page 6}
\[
\begin{tikzcd}[row sep=small]
& b \\
b\times b \arrow[ur, "\pr_1"] \arrow[dr, "\pr_2"'] & \\
& b
\end{tikzcd}
\]

% ---- from Der-ChXIIapp-19-06-03_p006-010.tex
\printedpage{5}
% head: A. Grothendieck : Dérivateurs, Ch. XII, Édition des Universités de Montpellier II et Paris VII
% foot: Version 20 juin 2003

dans $U'$, et les hypothèses a), b) de la page 1 assurent que la somme amalgamée $\varinjlim$ de ce diagramme dans $U'^{\wedge}$ est $e_{U'}$, i.e. que
\[ p_!(e_U) = e_{U'}. \]
Il s'ensuit que dans la somme amalgamée, $X'$ envisagée dans $(*)$, l'unique élément $t'$ de $T'$ est objet final de $X'$. On a vu que $U'$ est 1-connexe, mais non 2-connexe si le groupe abélien enveloppant de $G$ (i.e. $G/$(groupe des commutateurs) dans le cas où $G$ est un groupe) n'est pas réduit à 1. L'exemple le plus économique s'obtient en prenant $G = \{\pm 1\}$ le groupe à deux éléments. C'est donc un cas où dans $(*)$ on a
\[ i \in W \ \ \text{(et même $i$ asphérique),} \qquad \text{mais } i' \notin W, \]
et qui a l'avantage sur l'exemple plus sophistiqué de la page 103 (prop. 10) de ne mettre en jeu que des catégories finies -- et même fort petites !

\dgrblock{dgrb614}{Adjoining a terminal object after the pushout}
\begin{dcomment}{examples}{p6:ne mettre en jeu que des catégories finies}{Adjoining a terminal object after the pushout}
This construction uses the failure of asphericity to test whether an ordinary pushout preserves weak equivalences. Its mechanism is transparent on representables. For any object $x$ of a category satisfying the opening hypotheses, put $B_x=\Hom(x,b)$ and $C_x=\Hom(x,c)$. Both sets are nonempty, and
\[
 B_x\amalg_{B_x\times C_x}C_x=\{*\}.
\]
Every element of $B_x$ is identified with every element of $C_x$, so there is exactly one equivalence class. As presheaf colimits are computed objectwise, the span $b\leftarrow b\times c\to c$ has terminal colimit. This verifies $p_!1_\Psi=1_{U'}$ without a homotopy assertion.

Write $\Psi^{\triangleright}$ for $\Psi$ with a new terminal object. In the pushout $U'\amalg_\Psi\Psi^{\triangleright}$, a map from $x\in U'$ to the new object is represented by a map from $x$ to an object of the span, modulo its two identifications. Thus the presheaf of maps to that object is precisely $p_!1_\Psi$, and the pushout is $(U')^{\triangleright}$. The original inclusion $\Psi\to\Psi^{\triangleright}$ is aspherical: the comma category over the new vertex is $\Psi$, with initial object, and the comma categories over old vertices have terminal objects. After the pushout, $U'\to(U')^{\triangleright}$ is a weak equivalence exactly when $BU'$ is weakly contractible.

Taking $U'=X_{C_2}$ therefore validates the announced failure. Both the original span and its cone are contractible, whereas $H^3(BU';\mathbb F_2)=\mathbb F_2$. The source's asserted obstruction in $\pi_2$ must be replaced by the degree-three calculation above; the categorical pushout and its terminal object require no alteration. This is also the point where the distinction from a homotopy pushout becomes visible: a homotopy pushout along the equivalence $\Psi\to\Psi^{\triangleright}$ would retain the homotopy type of $U'$, while the ordinary categorical pushout has become a cone.
\end{dcomment}


Procédons maintenant comme dans la proposition 11, en désignant par $U''$ la catégorie déduite de $U'$ en lui ajoutant un objet
\msp{page 7}
final, et soit
\[ p' : U' \longrightarrow U'' \]
l'inclusion, de sorte que
\[
\begin{aligned}
p'_!(e_{U'}) \quad &= \quad (p'p)_!(e_U) = U' \quad && \text{en tant qu'objet dans } U''^{\wedge}\\
& && \text{(sous-préfaisceau du préfaisceau final } e_{U''}\text{).}
\end{aligned}
\]
Considérons
\[
\begin{tikzcd}
U \arrow[r, hook, "i"] \arrow[d, "{p''=p'p}"'] \arrow[dr, phantom, "\begin{smallmatrix}\text{cocart.}\\ \text{et cart.}\end{smallmatrix}"] & X \arrow[d, "q''"'] & T \arrow[l, hook', dotted, "j"'] \arrow[d, dotted, "\text{iso}"] \\
U'' \arrow[r, hook, "i''"] & X'' & T''. \arrow[l, hook', dotted, "j''"']
\end{tikzcd}
\]
Cette fois $U$, $X$ \textit{et} $U''$ sont $\Winf$-asphériques (et même contractiles toutes trois), donc $i$, $p''$ sont dans $W$, pourtant $X''$ n'est pas $\Winf$-asphérique (donc $i'', q'' \notin \Winf$), car pour qu'elle le soit, i.e. pour que $i'' \in \Winf$, il faudrait (prop. 7, cor. page 85 (\footnote[4]{pas encore démontré !})) que $i''$ soit $\Winf$-asphérique, i.e. que $U''/t''$ soit $\Winf$-asphérique. Mais par construction $U''/t'' \simeq U'$, qui n'est \textit{pas} $\Winf$-asphérique.

\dgrblock{dgrb615}{Two cone vertices and a suspension}
\begin{dcomment}{exposition}{p7:et même contractiles toutes trois}{Two cone vertices and a suspension}
Grothendieck strengthens the example by making the other map out of $\Psi$ a weak equivalence as well. Let $C=U'$, let $e$ be the terminal object adjoined to form $U''=C^{\triangleright}$, and retain $t$ for the cone vertex coming from $X=\Psi^{\triangleright}$. The left Kan extension $p'_!1_C$ is the presheaf with value one on the old objects of $C$ and value zero at $e$: no map runs from the new terminal object back to $C$. This is the subterminal presheaf called $U'$ in the displayed formula.

Consequently the pushout $X''$ consists of $C$ with two incomparable new terminal vertices $e,t$, and one map from each old object to each vertex. There are no arrows between $e$ and $t$. Every nondegenerate nerve simplex lies in one of the two cones; their intersection is exactly $NC$. Hence
\[
 BX''=C(BC)\cup_{BC}C(BC)\simeq\Sigma BC.
\]
This gives the omitted nonasphericity argument directly. For the finite choice $C=X_{C_2}$, suspension in reduced cohomology yields
\[
 \widetilde H^4(BX'';\mathbb F_2)
   \cong\widetilde H^3(BC;\mathbb F_2)
   \cong\mathbb F_2.
\]
Thus $X''$ is not weakly contractible, although $\Psi$, $X$ and $U''$ are contractible. Both maps to $X''$ from the latter two categories fail to be weak equivalences. No converse to Quillen's Theorem A is required. In general a noncontractible space may have contractible suspension if it is connected and acyclic, so a bare appeal to the nonasphericity of a comma category would not suffice. Here the explicit surviving cohomology class proves exactly what is needed.
\end{dcomment}


Voilà un exemple encore plus simple d'un foncteur $f : \Psi \longrightarrow X$ tel que $f_!(e_{\Psi})$ ne soit pas $\Winf$-asphérique, i.e. d'un diagramme
\[
\begin{tikzcd}
& b \\
a \arrow[ur, "\beta_0"] \arrow[dr, "\gamma_0"'] & \\
& c
\end{tikzcd}
\]
\printedpage{6}
% head: A. Grothendieck : Dérivateurs, Ch. XII, Édition des Universités de Montpellier II et Paris VII
% foot: Version 20 juin 2003
\msp{page 8}
dans une catégorie $V$, telle que la limite inductive $L = b \amalg_a c$ \textit{dans} $V^{\wedge}$ soit un objet non asphérique de $V^{\wedge}$, i.e. $V/L$ soit non asphérique. Je sais qu'il faut, pour un tel contre-exemple, que $a \longrightarrow b$ et $a \longrightarrow c$ ne soient pas des monomorphismes. L'idée, c'est de prendre un exemple universel où ces flèches ne soient pas \texttt{[des]} monomorphismes, en prenant
\msp{page 9}
\[
V = \left(
\begin{tikzcd}
& & b \\
0 \arrow[r, shift left, "\lambda"] \arrow[r, shift right, "\mu"'] & a \arrow[ur, "\beta"] \arrow[dr, "\gamma"'] & \\
& & c
\end{tikzcd}
\right)
\]
avec
\[ \beta\lambda = \beta\mu \qquad \text{et} \qquad \gamma\lambda = \gamma\mu \]
les deux seules flèches non identiques, en plus des quatre flèches marquées $\beta$, $\gamma$, $\lambda$, $\mu$.

\dgrblock{dgrb616}{A presheaf pushout remembers only components}
\begin{dcomment}{exposition}{p8:Je sais qu'il faut, pour un tel contre-exemple}{A presheaf pushout remembers only components}
Grothendieck now seeks a finite obstruction that can be inserted into a much larger category. The decisive distinction is between the ordinary pushout of three representable presheaves and the homotopy carried by its category of elements. We make that distinction explicit before using the example. Write $y(a)=\Hom_C(-,a)$, and let
\[
 L=y(b)\amalg_{y(a)}y(c)
\]
for a span $b\xleftarrow{\beta}a\xrightarrow{\gamma}c$ in $C$. For $x\in C$, form the bipartite multigraph $K_x$ with vertices $b(x)\amalg c(x)$ and one edge for every $u\in a(x)$, joining $\beta u$ to $\gamma u$. Then $L(x)=\pi_0K_x$. An ordinary pushout has forgotten the cycles in this graph.

This description proves the assertion about monomorphisms, including the case where $L$ is not terminal. Let $E$ have objects $(v,x,u)$, with $v\in\Psi$ and $u:x\to f(v)$ for the given span $f:\Psi\to C$. A morphism to $(v',x',u')$ consists of $h:v\to v'$ and $r:x\to x'$ satisfying $f(h)u=u'r$. The projection $E\to\Psi$ has right adjoint $v\mapsto(v,f(v),1)$, so $BE$ is contractible. There is a canonical functor
\[
 P:E\longrightarrow C/L,\qquad (v,x,u)\longmapsto(x,[u]).
\]
Here $C/L$ denotes the category of elements of the presheaf $L$. For $(x,l)\in C/L$, the comma category $(x,l)\backslash E$ retracts up to a natural transformation onto the incidence category of the component $l$ of $K_x$. Indeed an object of this comma category includes $r:x\to x'$; send it to $(v,x,u'r)$ and use the morphism given by $r$ for the comparison with the original object. The incidence category has one object for each vertex and edge, and two arrows from each edge to its endpoints.

If $\beta$ is monic, every vertex in $b(x)$ has degree at most one. Every connected component of $K_x$ is consequently a star, possibly a single vertex, and its incidence category is contractible. Quillen's Theorem A, in its comma-category form, gives $BP:BE\simeq B(C/L)$. Thus $L$ is aspherical. The same argument applies with $b$ and $c$ interchanged. This supplies the proof promised in the source; the only external ingredient is the comma-category criterion \citep[\S1, Theorem A]{Quillen1973}.

There is also a useful conclusion without a monomorphism hypothesis. If $L=e_C$, every $K_x$ is nonempty and connected, and $C$ is simply connected. To see this directly, work in $\Pi C$. Put $d=[\beta][\gamma]^{-1}:c\to b$. Send a vertex $v\in b(x)$ to $[v]:x\to b$ and a vertex $w\in c(x)$ to $d[w]$. The two endpoints of an edge $u\in a(x)$ have the same image, since $d[\gamma u]=[\beta u]$. Connectedness therefore gives a single arrow $s_x:x\to b$. For every $r:x\to y$, precomposition of a vertex of $K_y$ proves $s_y[r]=s_x$. These identities extend to all zigzags and force every arrow from $x$ to $y$ in $\Pi C$ to equal $s_y^{-1}s_x$. Thus $\Pi C$ is equivalent to the terminal groupoid. The correction on manuscript p.~10 is justified: the first possible obstruction is in degree two.
\end{dcomment}


On trouve que
\[ b \amalg_a c \xrightarrow{\ \sim\ } e_X \qquad \text{(somme amalgamée dans } V^{\wedge}\text{).} \] % [sic] e_X -> e_V, the category considered here is V

Pourtant $V$ n'est pas asphérique, ni même 1-connexe : il est homotope à la sous-catégorie ouverte $V_0 = (0 \overset{\lambda}{\underset{\mu}{\rightrightarrows}} a)$, car on a une rétraction $V \xrightarrow{\ r\ } V_0$ (envoyant $b$ et $c$ sur $a$), et une flèche $ir \longrightarrow \id_V$ (où $i : X_0 \hookrightarrow X$ est l'inclusion). Donc $V$ a le même type d'homotopie \textit{que} $V_0$, qui est un cercle. Donc on a intérêt à prendre ce $V$, où $\tilde{V}$ déduit de $V$ en rajoutant un élément final $e$, % [sic] i : X_0 \hookrightarrow X -> i : V_0 \hookrightarrow V, the inclusion of V_0 in V is meant
\[
\tilde{V} = \left(
\begin{tikzcd}
& & b \arrow[dr] & \\
0 \arrow[r, shift left] \arrow[r, shift right] & a \arrow[ur] \arrow[dr] & & e \\
& & c \arrow[ur] &
\end{tikzcd}
\right),
\]
comme $U'$, pour les contre-exemples pages 6 et 7.

\dgrblock{dgrb617}{The missing retraction and the two-sphere}
\begin{dcomment}{examples}{p9:comme $U'$, pour les contre-exemples pages 6 et 7.}{The missing retraction and the two-sphere}
The four-object category is useful because its relations can be seen as attaching two disks to the same circle. We use the notation $Z$ of the following pages, suppressing subscripts: its objects are $0,a,b,c$, its generating arrows are $\lambda,\mu:0\rightrightarrows a$, $\beta:a\to b$, and $\gamma:a\to c$, and its relations identify the two composites to each of $b,c$. The commentary on Chapter~II, manuscript p.~36 (``The retracted proposition and a visible sphere'') already records the general obstruction to total asphericity. Here we need the precise attaching maps, because the next pushout will suspend this particular sphere once more.

Let $D=(0\rightrightarrows a)$. A retraction $Z\to D$ would fix $a$, so it must send $b$ to $a$ and $\beta$ to $1_a$. Applying it to $\beta\lambda=\beta\mu$ would give $\lambda=\mu$, a contradiction. This is the exact point at which the proposed retraction fails. The full subcategories on $\{0,a,b\}$ and $\{0,a,c\}$ are instead the two terminal cones on $D$. Every composable string in $Z$ belongs to one of them, and their intersection is $D$. Consequently
\[
 NZ=N(D^{\triangleright_b})\amalg_{ND}N(D^{\triangleright_c}),
 \qquad BZ\cong\Sigma BD\cong S^2.
\]
The first identity is an actual identity of simplicial sets. Realization preserves this union, and $BD$ is the circle with two vertices and two edges. This proves the correction already supplied by footnote~7, rather than replacing one guessed homotopy type by another.

The calculation is also visible in normalized chains. There are four vertices, six nonidentity arrows, four nondegenerate $2$-simplices, and no higher nondegenerate simplices. With $[u,v]$ denoting a composable pair, a fundamental cycle is
\[
 [\lambda,\beta]-[\mu,\beta]-[\lambda,\gamma]+[\mu,\gamma].
\]
Its boundary vanishes by the two defining relations; it generates $H_2(BZ;\mathbf Z)\cong\mathbf Z$. In particular $\pi_1(BZ)=0$ and $\pi_2(BZ)\cong\mathbf Z$.

The presheaf pushout nevertheless is terminal. Its pointwise calculation is
\[
\begin{array}{c|cccc}
 x&0&a&b&c\\\hline
 |a(x)|&2&1&0&0\\
 |b(x)|&1&1&1&0\\
 |c(x)|&1&1&0&1\\
 |L(x)|&1&1&1&1
\end{array}
\]
and all restriction maps of $L$ are the unique maps between singletons. Thus $Z/L\cong Z$, while $L=e_Z$. The same data establish that $\Psi$ is not totally $\Winf$-aspherical. Both arrows out of $a$ fail to be monic, and the same pair $\lambda,\mu$ witnesses both failures. The word ``universal'' refers to the free category with these specified generators and relations; arbitrary failures of the two monomorphisms need not have a common witnessing pair.
\end{dcomment}


\msp{page 10}
\subsection*{2.6. Contre-exemples pour la totale asphéricité.}

Je me suis finalement convaincu que
\[
\Psi = \left\{
\begin{tikzcd}
& b_0 \\
a_0 \arrow[ur] \arrow[dr] & \\
& c_0
\end{tikzcd}
\right\}
\]
\printedpage{7}
% head: A. Grothendieck : Dérivateurs, Ch. XII, Édition des Universités de Montpellier II et Paris VII
% foot: Version 20 juin 2003
n'est \textit{pas} totalement asphérique (pour le localiseur $\Winf$ -- dans le présent numéro, on supposera $W = \Winf$), ou ce qui revient au même, qu'on peut avoir un diagramme de type $\Psi$
\[
\begin{tikzcd}
& b \\
a \arrow[ur, "\beta"] \arrow[dr, "\gamma"'] & \\
& c
\end{tikzcd}
\tag*{$(*)$}
\]
dans une catégorie $Z \in \mathrm{Ob}\,\mathrm{Cat}$, dont la limite inductive $b \amalg_a c$ dans $Z^{\wedge}$ soit le préfaisceau final $e_Z$, i.e. tel que l'on ait
\[ b(z) \amalg_{a(z)} c(z) \xrightarrow{\ \sim\ } \text{ensemble ponctuel} \]
pour tout objet $z$ de $Z$ (où $a(z) = \Hom_Z(a,z)$ \texttt{[plutôt} $\Hom_Z(z,a)$\texttt{]}) sans que pour autant $Z$ soit asphérique, ni même 1-connexe (\footnote[5]{faux -- il est 1-connexe forcément (\footnotemark[6]), mais non 2-connexe.}).\footnotetext[6]{cf. commentaire XII, p. 101.} L'exemple le plus simple que j'ai trouvé, est l'exemple ‘universel’ d'un diagramme $(*)$ où ni l'une ni l'autre des deux flèches $a \longrightarrow b$, $a \longrightarrow c$ n'est \texttt{[un]} monomorphisme (car si elle l'est, alors on peut prouver que $Z$ est bel et bien asphérique, comme nous verrons par la suite). C'est la catégorie
\[
Z = \left\{
\begin{tikzcd}
& & b_0 \\
0_0 \arrow[r, shift left, "\lambda_0"] \arrow[r, shift right, "\mu_0"'] & a_0 \arrow[ur, "\beta_0"] \arrow[dr, "\gamma_0"'] & \\
& & c_0
\end{tikzcd}
\right\}
\]
\msp{page 11}
déduite de $\Psi$ en lui rajoutant un objet $0_0$, relié à $a_0$ par deux flèches distinctes $\lambda_0$, $\mu_0$, telles que
\[ \beta_0\lambda_0 = \beta_0\mu_0 \qquad \text{et} \qquad \gamma_0\lambda_0 = \gamma_0\mu_0, \]
ces deux flèches composées étant, avec les quatre flèches $\beta_0$, $\gamma_0$, $\lambda_0$, $\mu_0$ marquées sur le diagramme, les seules (six) flèches non identiques de $Z$. $\Psi$ immergée dans $Z$ comme un fermé, de complémentaire la catégorie ponctuelle réduite à $0_0$. (NB L'immersion fermée $\Psi \longrightarrow Z$ n'est pas immersion de Dwyer, i.e. (comme $Z$ est le plus petit ouvert de $Z$ contenant $\Psi$) n'est pas coasphérique, car la catégorie $0\backslash\Psi$, cofibrée sur $\Psi$ avec fibres discrètes, de cardinal 2 en $a_0$ et 1 en $b_0$ et $c_0$, est la catégorie
\[
\begin{tikzcd}
& \tilde{b}_0 \\
a'_0 \arrow[ur] \arrow[ddr] & \\
a''_0 \arrow[uur, crossing over] \arrow[dr] & \\
& \tilde{c}_0
\end{tikzcd}
\]
\printedpage{8}
% head: A. Grothendieck : Dérivateurs, Ch. XII, Édition des Universités de Montpellier II et Paris VII
% foot: Version 20 juin 2003
i.e.
\[
\begin{tikzcd}
& {} & \\
{} \arrow[ur] \arrow[dr] & & {} \arrow[ul] \arrow[dl] \\
& {} &
\end{tikzcd}
\quad ,
\]
qui a le type d'homotopie du cercle.) Ni l'immersion ouverte $\{0_0\} \longrightarrow Z$, car $\{0_0\}/a_0$ est la catégorie discrète de cardinalité 2.

\dgrblock{dgrb618}{The link detects the failure of Dwyer's condition}
\begin{dcomment}{exposition}{p11:qui a le type d'homotopie du cercle.)}{The link detects the failure of Dwyer's condition}
Grothendieck is locating the obstruction where it will survive a change of ambient category. For the closed inclusion $\Psi\hookrightarrow Z$, the relevant link is $0\backslash\Psi$. Its four objects are the two arrows $\lambda,\mu$ to $a$ and the unique arrows to $b,c$. Each of the first two objects maps to each of the last two, and there are no other nonidentity arrows. Its nerve is the square graph, hence its classifying space is $S^1$. Projection to $\Psi$ is the discrete opfibration classified by the covariant functor $\Hom_Z(0,-)|_\Psi$, with fibre cardinalities $2,1,1$.

This link has neither an initial object nor a contractible nerve. The inclusion is therefore neither a closed Dwyer inclusion nor coaspherical. These are separate conclusions: Dwyer's condition asks for a left adjoint to the inclusion into the smallest open neighbourhood, equivalently for initial objects in the relevant comma categories; weak contractibility alone would be a weaker requirement. For the complementary open inclusion $\{0\}\hookrightarrow Z$, the slice over $a$ is the discrete two-object category, so this inclusion also fails Dwyer's condition. The larger category can be simply connected while these local links retain nontrivial information. That local information, rather than the nonexistent retraction, will drive the finite pushout counterexample.
\end{dcomment}


L'inclusion $i : Z_0 \longrightarrow Z$ de
\[ Z_0 \ = \ (\ 0_0 \overset{\lambda_0}{\underset{\mu_0}{\rightrightarrows}} a_0\ ) \]
dans $Z$ est un homotopisme (\footnote[7]{faux, il n'y a pas de rétraction du tout! En fait, $Z$ est la \textit{suspension} de la $\rightrightarrows$, qui a le type d'homotopie (topologique) du cercle, et il a donc le type d'homotopie de la 2-sphère.}) (car il existe une rétraction évidente $r$ de $Z$ sur $Z_0$ (\footnotemark[7]), et une flèche $ir \longrightarrow \id_Z$), donc $Z$ est homotope à $Z_0$, qui a le type d'homotopie du cercle, donc n'est pas 1-connexe.

\msp{page 12}
D'autre part, la vérification de
\[ b_0 \amalg_{a_0} c_0 \ = \ e_Z \qquad \text{dans } Z^{\wedge} \]
est immédiate.

On va utiliser cette flèche
\[ \Psi \longrightarrow Z \]
pour construire des exemples intéressants (et déchirants) de sommes amalgamées
\[
\begin{tikzcd}
U \arrow[r, hook, "i"] \arrow[d, "p"'] & X \arrow[d, "q"'] \\
U' \arrow[r, hook, "i'"] & X',
\end{tikzcd}
\]
avec $i$ immersion ouverte, telle qu'on ait même
\[ i :\ U \longrightarrow \overline{U} \qquad \text{est asphérique, donc dans } W, \]
alors que pourtant on aura
\[ i' :\ U' \longrightarrow \overline{U'} \qquad \notin W, \]
et ce qui mieux (ou pis) est,
\[ p \in W \qquad \text{mais} \qquad q \notin W \]
\printedpage{9}
% head: A. Grothendieck : Dérivateurs, Ch. XII, Édition des Universités de Montpellier II et Paris VII
% foot: Version 20 juin 2003
(et éventuellement, en même temps, $i \in W$ mais $i' \notin W$). De plus (contrairement aux constructions des propositions 10, 11 (p. 101-107)), les catégories envisagées pourront toutes être finies. Je suis particulièrement intéressé par le cas où $X$ est un ‘carré’
\[ X = \Phi_n \times \Phi_m \qquad (n,m \geq 2), \]
et où $U$ est son ‘bord’
\[ U = \dot{X} \overset{\text{déf}}{=} (\dot{\Phi}_n \times \Phi_m) \cup (\Phi_n \times \dot{\Phi}_m). \]

\dgrblock{dgrb619}{The boundary and its categorical closure}
\begin{dcomment}{exposition}{p12:et où $U$ est son ‘bord’}{The boundary and its categorical closure}
Grothendieck wants to know whether ordinary subdivisions of a square provide attaching maps compatible with weak equivalences. The notation carries two geometries: the realization of the nerve is a genuine square, while the open and closed subsets of the category are determined by the directions of its arrows. Write
\[
 \Phi_n=(0\longrightarrow1\longleftarrow2\longrightarrow\cdots
                \longrightarrow2n-1\longleftarrow2n),
 \qquad \dot\Phi_n=\{0,2n\}.
\]
This is the convention in Chapter XII, manuscript p.~114 \citep{Grothendieck1990XII}; the commentary on directed intervals in Chapter~VII, manuscript pp.~1--19, explains its relation to the underlying unoriented interval. The boundary $U$ consists of the pairs with at least one coordinate an endpoint. It is a sieve, and
\[
 BX\cong[0,1]^2,\qquad BU\cong\partial[0,1]^2\cong S^1.
\]
In particular $U\to X$ is not a weak equivalence. The expression ``immersion ouverte $W$-asph\'erique'' in the proposition must be read in the local sense $U\to\overline U$: it is this induced functor to the categorical closure that is aspherical.

The slice calculation makes the distinction exact. Put
\[
 E_n=\{0,2n\},\qquad E_n^+=\{0,1,2n-1,2n\}.
\]
The cosieve generated by $U$ is
\[
 \overline U=(E_n^+\times\Phi_m)\cup(\Phi_n\times E_m^+).
\]
For $t\notin\overline U$, the slice $U/t$ is empty. For $t\in U$ it has terminal object $t$. Suppose that $t\in\overline U\setminus U$. If exactly one coordinate is adjacent to an endpoint and the other is an interior even vertex, then $U/t$ is a point. If that other coordinate is an odd vertex not adjacent to an endpoint, $U/t$ is the three-object interval with a terminal middle vertex. Finally, when both coordinates are adjacent to endpoints, $U/t$ is the five-object category $\widetilde\Psi$ drawn in the source. There are exactly four such $t$, because $n,m\geq2$.

The nerve of $\widetilde\Psi$ is the path with successive vertices $a',a,0,b,b'$. It is contractible, although neither $a$ nor $b$ receives a map from the other. Thus every slice over $\overline U$ is contractible, proving the claimed local asphericity, while the four exceptional slices have no terminal object. For a sieve of posets, Dwyer's condition is precisely the existence of these terminal objects over its generated cosieve: they assemble into the right adjoint to the inclusion. This proves both that $U\hookrightarrow X$ is not a Dwyer inclusion and that deleting the four exceptional vertices makes it one. Contractible slices suffice for a weak equivalence to the closure; an adjoint supplies the stronger structure that behaves well under arbitrary pushout.
\end{dcomment}


\msp{page 13}
\[
\begin{tikzcd}
0 \arrow[r] \arrow[d] & a \arrow[d] & a' \arrow[l] \arrow[r] \arrow[d] & \cdots \\
b \arrow[r] & t_1 & \cdots \arrow[l] & \\
b' \arrow[u] \arrow[r] \arrow[d] & {} \arrow[u] & \cdots \arrow[l] \arrow[u] & \\
\cdots & & &
\end{tikzcd}
\]

\textbf{Proposition.} \textit{Soient $n,m \geq 2$, et considérons l'inclusion $U = \partial X \overset{i}{\hookrightarrow} X$ du bord du ‘carré’ $X = \Phi_n \times \Phi_m$ dans $X$. C'est une immersion ouverte $W$-asphérique, de façon précise, les $U/t$ ($t \in \dot{U} = \bar{U}\backslash U$) sont isomorphes à $e$ (catégorie ponctuelle) ou à $\Phi =$ \begin{tikzcd}[ampersand replacement=\&, row sep=small, column sep=small] \& {} \arrow[dl] \\ {} \& \\ \& {} \arrow[ul] \end{tikzcd} (catégorie avec élément final), sauf pour les quatre éléments $t_i$ des quatre ‘coins’ de $\dot{U}$, où c'est la catégorie}
\[
\tilde{\Psi} = \left(
\begin{tikzcd}
& & a' \arrow[dl] \\
& a & \\
0 \arrow[ur] \arrow[dr] & & \\
& b & \\
& & b' \arrow[ul]
\end{tikzcd}
\right)
\]
\textit{n'ayant pas d'élément final. Donc $U \xrightarrow{\ i\ } X$ n'est pas immersion ouverte de Dwyer (mais $U \longrightarrow X\backslash\{t_1,t_2,t_3,t_4\}$ l'est \ldots). Ceci posé, je dis que $i : U \longrightarrow X$ n'est pas une $W$-cofibration.}

% ---- from Der-ChXIIapp-19-06-03_p011-014.tex
\printedpage{10}
% head: A. Grothendieck : Dérivateurs, Ch. XII, Édition des Universités de Montpellier II et Paris VII
% foot: Version 20 juin 2003
\textsc{Démonstration}. Soit
\[ F \hookrightarrow U \]
l'inclusion dans $U = \partial X$ de la catégorie-coin \texttt{[}$V =$\texttt{]} $\vcenter{\hbox{\begin{tikzcd}[row sep=small, column sep=small]
0 \arrow[r] \arrow[d] & a \\
b &
\end{tikzcd}}}$ (cf. figure), isomorphe à $\Psi$. On voit aussitôt que c'est une immersion fermée de Dwyer (le plus petit ouvert de $U$ contenant $\Psi$ est justement la catégorie $\tilde{\Psi} = U/t_1$ de tantôt), donc une $W$-cofibration. Si $U \longrightarrow X$ était une $W$-cofibration, il en serait de même du composé $V \longrightarrow X$. On va voir que ce n'est pas le cas.

Soit $\overline{V} = \vcenter{\hbox{\begin{tikzcd}[row sep=small, column sep=small]
0 \arrow[r] \arrow[d] & a \arrow[d] \\
b \arrow[r] & t_1
\end{tikzcd}}}$ l'adhérence de $V$ dans $X$,
\msp{page 14}et considérons l'inclusion
\[ \overline{V} \longrightarrow X, \]
je dis que \texttt{[c]}'est une immersion fermée de Dwyer, i.e. que pour tout $s \in W\backslash\overline{V}$, où $W$ est le plus petit ouvert de $X$ contenant $\overline{V}$, i.e.
\[
W = \quad
\begin{tikzcd}[row sep=small, column sep=small]
0 \arrow[r] \arrow[d] & a \arrow[d] & a' \arrow[l] \arrow[d] \\
b \arrow[r] & {} & a'' \arrow[l] \\
b' \arrow[u] \arrow[r] & b'' \arrow[u] & c, \arrow[l] \arrow[u]
\end{tikzcd}
\]
la catégorie $s\backslash\overline{V}$ a un objet initial. (Les catégories $s\backslash\overline{V}$ qu'on obtient, à partir des cinq choix possibles pour $s$, sont $\boldsymbol{\Delta}^1$ (si $s = a'$ ou $b'$) ou $e$ (pour $s = a''$, $b''$ ou $c$)). Donc $\overline{V} \longrightarrow X$ est une $W$-cofibration.

Ceci posé, on sait qu'il existe, pour [?] l'inclusion (immersion ouverte)
\[ V \hookrightarrow \overline{V} \qquad \text{i.e. } \Psi \hookrightarrow \boldsymbol{\Delta}^1 \times \boldsymbol{\Delta}^1, \]
un foncteur $V \longrightarrow V'$, donnant lieu à un carré cocartésien
\[
\begin{tikzcd}
V \arrow[r, hook, "i"] \arrow[d, "p"'] & \overline{V} \arrow[d, "\overline{p}"] \\
V' \arrow[r, hook, "i'"] & \overline{V'}
\end{tikzcd}
\tag*{$(*)$}
\]
avec
\[ i, p \in W \ \text{(car $V$, $V'$, $\overline{V}$ $W$-asphérique), mais } i', \overline{p} \notin W \ \text{(i.e. $\overline{V'}$ \textit{pas} $W$-asphérique).} \]

Je rappelle qu'on prend d'abord
\printedpage{11}
% head: A. Grothendieck : Dérivateurs, Ch. XII, Édition des Universités de Montpellier II et Paris VII
% foot: Version 20 juin 2003
\[
V'_0 = \left(
\begin{tikzcd}[row sep=small]
 & & a \\
0' \arrow[r, shift left, "\lambda"] \arrow[r, shift right, "\mu"'] & 0 \arrow[ur, "\alpha"] \arrow[dr, "\beta"'] & \\
 & & b
\end{tikzcd}
\right)
\qquad \text{avec} \quad
\left\{
\begin{aligned}
\alpha\lambda &= \alpha\mu \\
\beta\lambda &= \beta\mu
\end{aligned}
\right.
\]
et l'immersion fermée ‘identique’ de $V$ dans $V'$ \texttt{[plutôt} $V'_0$\texttt{]}
\[ p_0 : V \longrightarrow V'_0 \]
\msp{page 15}(qui est telle que
\[ p_{0!}(e_V) = e_{V'_0} \ \ ), \]
puis on prend pour $V'$ la catégorie obtenue en ajoutant un objet final à $V'_0$
\[
V' = \left(
\begin{tikzcd}[row sep=small, column sep=small]
 & & a \arrow[dr, "\rho"] & \\
0' \arrow[r, shift left, "\lambda"] \arrow[r, shift right, "\mu"'] & 0 \arrow[ur, "\alpha"] \arrow[dr, "\beta"'] & & e \\
 & & b \arrow[ur, "\sigma"'] &
\end{tikzcd}
\right)
\qquad \text{avec} \quad
\left\{
\begin{aligned}
\alpha\lambda &= \alpha\mu \\
\beta\lambda &= \beta\mu \\
\rho\alpha &= \sigma\beta
\end{aligned}
\right.
\]
et pour $p$ le foncteur composé
\[
\begin{tikzcd}
V \arrow[r, "p_0"] \arrow[rr, bend right, "p"'] & V'_0 \arrow[r, "p_1"] & V',
\end{tikzcd}
\]
où $p_1$ est l'inclusion (ouverte). On trouve alors le carré $(*)$, avec $\overline{V'}$ non $W$-asphérique, ou ce qui revient au même, $i' \notin W$. Car comme $\overline{V'}\backslash V' \simeq \overline{V}\backslash V$ est la catégorie ponctuelle, $i' \in W$ équivaut à $V'/t$ $W$-asphérique (\footnote[8]{moyennant (W 7 bis).}) (cf. cor. 4, p. 149 dans XII). Or par construction $V'/t \simeq V'_0$, qui n'est pas $W$-asphérique (dans le cas, disons, où $W \subseteq W_1$, car $\pi_1(V'_0) \simeq \mathbf{Z}$).

\dgrblock{dgrb620}{The pushout is a three-sphere}
\begin{dcomment}{examples}{p15:dans le cas, disons, où $W \subseteq W_1$}{The pushout is a three-sphere}
The purpose of adjoining the new terminal object $e$ is to make the replacement $V'$ contractible without removing the obstruction in the link of the other cone point $t$. It is possible to calculate the entire pushout, so the argument need not depend on the converse excision principle mentioned in footnote~8. Rename $V'_0$ as $Z$ and the original three-object corner as $V$. Thus $Z$ is the four-object category just analysed, $V'=Z^{\triangleright_e}$, and $\overline V=V^{\triangleright_t}$.

Let $Q=V'\amalg_V\overline V$. Its objects are the four objects of $Z$ and two further objects $e,t$. It contains $Z$ as a full subcategory. For each $z\in Z$ there is exactly one arrow $z\to e$ and exactly one arrow $z\to t$; there are no arrows between $e$ and $t$. To check the only nontrivial assertion, the cocone over $V$ supplies the maps from its three objects to $t$. The two composites from the extra object $0'$ to $t$ agree, because the relation $\alpha\lambda=\alpha\mu$ already identifies their composites through $a$. Hence there is a unique arrow $0'\to t$. This description has the pushout universal property: a functor out of $Q$ is exactly a functor out of $Z$ together with one cocone over $Z$ with vertex the image of $e$ and one with vertex the image of $t$; the latter is uniquely determined by its restriction to $V$.

Every simplex of $NQ$ lies in one of the two cones on $NZ$, and their intersection is $NZ$. Therefore
\[
 BQ\cong\Sigma BZ\cong S^3,
 \qquad \widetilde H_k(BQ;\mathbf Z)=
 \begin{cases}\mathbf Z,&k=3,\\0,&k\ne3.\end{cases}
\]
Since $BV$, $B\overline V$, and $BV'$ are contractible, the maps $i$ and $p$ are weak equivalences, but $i'$ and $\overline p$ are not. This is a completely finite example: $Q$ has six objects, fourteen nonidentity arrows, sixteen nondegenerate $2$-simplices, and eight nondegenerate $3$-simplices. Its nerve has no higher nondegenerate simplices.

The sentence $\pi_1(V'_0)\cong\mathbf Z$ is present in the edition and must remain in the historical text. Its correct replacement in the argument is
\[
 \pi_1(V'_0)=0,\qquad \pi_2(V'_0)\cong\mathbf Z,
 \qquad \pi_3(Q)\cong\mathbf Z.
\]
The example proves the required assertion for $W=\Winf$. It does not prove the printed extension to every $W\subseteq W_1$: both $V'\to Q$ and $\overline V\to Q$ are equivalences on $2$-types. Here $W_k$ means equivalence on homotopy groups in degrees at most $k$ and on components.

There is a second, logically independent issue. For an arbitrary open inclusion with one-point complement, membership in $\Winf$ does not force its comma category over that point to be contractible. The referenced Chapter XII corollary explicitly assumes the inverse part of $W(7\,\mathrm{bis})$. The acyclic-link construction in the commentary on manuscript pp.~17--18 exhibits the failure of that inverse implication for spaces. In the present example both the link and the pushout are known spheres, so the conclusion follows from the displayed nerve calculation with no appeal to this extra axiom.
\end{dcomment}


\msp{page 16}Considérons alors le diagramme de carrés cocartésiens
\[
\begin{tikzcd}
V \arrow[r, hook, "i"] \arrow[d, "p"'] & \overline{V} \arrow[r, hook] \arrow[d, "\overline{p}"] & X \arrow[d, "q"] \\
V' \arrow[r, hook, "i'"] & \overline{V'} \arrow[r, hook] & X'
\end{tikzcd}
\qquad .
\]

\printedpage{12}
% head: A. Grothendieck : Dérivateurs, Ch. XII, Édition des Universités de Montpellier II et Paris VII
% foot: Version 20 juin 2003
Comme $\overline{V} \hookrightarrow X$ est \texttt{[une]} immersion fermée de Dwyer, le carré de droite est $W$-cocartésien, donc $(\overline{p} \in W) \iff (q \in W)$, donc (comme $\overline{p} \notin W$) on a
\[ p \in W, \qquad \text{mais} \qquad q \notin W, \]
ce qui montre bien que $V \longrightarrow X$ n'est pas une $W$-cofibration. Factorisant $V \longrightarrow X$ par $U$ (au lieu de $\overline{V}$), on trouve
\[
\begin{tikzcd}
V \arrow[r] \arrow[d, "p\in W"'] & U \arrow[r] \arrow[d, "q_U\in W"] & X \arrow[d, "q\notin W"] \\
V' \arrow[r] & U' \arrow[r] & X'
\end{tikzcd}
\qquad ,
\]
où $q_U \in W$ puisque $p \in W$ et $V \longrightarrow U$ est \texttt{[une]} immersion de Dwyer. Cela rend explicite que $U \longrightarrow X$ n'est pas une $W$-cofibration, puisque
\[ q_U \in W, \qquad \text{mais} \qquad q \notin W, \]
\hfill q.e.d.

\dgrblock{dgrb621}{Carrying the local obstruction to the boundary}
\begin{dcomment}{exposition}{p16:ce qui montre bien que $V \longrightarrow X$}{Carrying the local obstruction to the boundary}
The last step places the six-object defect inside the subdivided rectangle. Grothendieck has chosen the corner so that the two surrounding inclusions are closed Dwyer inclusions. We record their adjoints, because they explain exactly which pushouts preserve homotopy and complete the logical step in the proof.

Use coordinates $0=(0,0)$, $a=(1,0)$, $b=(0,1)$, and $t_1=(1,1)$. The smallest open neighbourhood of $V=\{0,a,b\}$ in $U$ consists of these three points and $a'=(2,0), b'=(0,2)$. Its inclusion of $V$ has the left adjoint fixing $V$ and sending $a'$ to $a$ and $b'$ to $b$. Indeed these images are the least points of $V$ above the corresponding objects. The inclusion $V\hookrightarrow U$ is a cosieve, so it is a closed Dwyer inclusion.

Similarly $\overline V=\{0,1\}\times\{0,1\}$ is a cosieve of $X$. Its smallest open neighbourhood is $\{0,1,2\}^2$. Define $r_0(0)=0$, $r_0(1)=r_0(2)=1$; then $r_0\times r_0$ is left adjoint to the inclusion of $\overline V$ into this neighbourhood. Thus $\overline V\hookrightarrow X$ is also a closed Dwyer inclusion. This gives the source's five comma categories explicitly: $a'\backslash\overline V$ and $b'\backslash\overline V$ are two-object chains, whereas the other three are points.

Thomason's pushout theorem, applied after taking opposite categories, now identifies the nerves of the relevant categorical pushouts with homotopy pushouts \citep[Proposition~4.3 and Corollary~4.4]{Thomason1980}. One additional observation is needed for the reverse implication written in the source: $\overline V\to X$ is itself a weak equivalence, since both classifying spaces are contractible. Consequently its pushout $Q\to X'$ is a weak equivalence. In the commutative square
\[
\begin{tikzcd}
\overline V\arrow[r,"\sim"]\arrow[d,"\overline p"']&X\arrow[d,"q"]\\
Q\arrow[r,"\sim"]&X'
\end{tikzcd}
\]
two-out-of-three gives $\overline p\in\Winf$ if and only if $q\in\Winf$. Since $BQ\simeq S^3$, we obtain $BX'\simeq S^3$ and $q\notin\Winf$. A homotopy pushout by itself would not justify the converse: the two horizontal equivalences are the missing reason in this particular square.

Finally let $U'=U\amalg_VV'$. The other closed Dwyer inclusion, $V\to U$, ensures that $q_U:U\to U'$ is a weak equivalence because $p:V\to V'$ is one. Associativity of ordinary pushout identifies $X\amalg_UU'$ with the same $X'$. Thus the cobase-change functor along $U\to X$ takes the weak equivalence $q_U$ to the map $q$ which is not a weak equivalence. This proves exactly the definition of failure of a $W$-cofibration used in Chapter XII: the functor $U\backslash\Cat\to X\backslash\Cat$ fails to preserve weak equivalences. In modern terminology, this inclusion is not an $h$-cofibration. The target replacement still satisfies $BU'\simeq S^1$; the change in the ambient category is detected by its new third homology group.
\end{dcomment}


\dgrblock{dgrb622}{Dwyer maps, homotopy pushouts, and products}
\begin{dcomment}{later}{p16:Cela rend explicite que $U \longrightarrow X$}{Dwyer maps, homotopy pushouts, and products}
Grothendieck's distinction between contractible links and adjoint neighbourhoods remains useful in the modern theory. Thomason's original Dwyer condition guarantees that the comparison from a pushout of nerves to the nerve of a categorical pushout is a weak homotopy equivalence. It therefore controls the entire change-of-base operation, rather than merely the homotopy type of the initial inclusion. One should not identify Dwyer maps with all cofibrations in the Thomason model structure: \citet[Lemma~4 and Proposition~2]{Cisinski1999Dwyer} corrected the claim that the original Dwyer class is closed under retracts. His larger class of pseudo-Dwyer maps retains the required pushout property.

The later result of \citet[Theorems~1.5--1.6]{HackneyOzornovaRiehlRovelli2024} strengthens the comparison to an equivalence of $(\infty,1)$-categories. Thus pushouts along Dwyer maps preserve more than the classifying-space homotopy type. Our two uses of closed Dwyer maps are obtained by passing to opposites; the reversal changes the direction of the neighbourhood adjunction but does not change the classifying-space conclusion.

The negative product phenomenon also has a precise later form. \citet[Definition~2.6]{Pavlov2024Thomason} uses $h$-cofibration for the preservation condition Grothendieck calls $W$-cofibration. In the same paper, Proposition~4.3 proves that no model structure on $\Cat$ with Thomason weak equivalences can be cartesian. Its proof uses a different finite cube construction, credited in Proposition~4.2 to Maltsiniotis. The present appendix supplies a particularly transparent earlier obstruction: the boundary of a product of alternating intervals can have contractible links in its closure and still fail the preservation condition. These are results about strict categorical pushouts; homotopy pushouts of spaces retain their usual invariance.
\end{dcomment}


\dgrblock{dgrb623}{The same defect on higher cubical boundaries}
\begin{dcomment}{development}{p16:Cela rend explicite que $U \longrightarrow X$}{The same defect on higher cubical boundaries}
The rectangle argument is local enough to survive in every higher dimension. This gives a concrete continuation of Grothendieck's construction: the defect need not grow with the dimension of the ambient cube. The literature checked for this extension comprised Thomason's pushout theorem, Cisinski's correction, the Dwyer pushout comparison of Hackney--Ozornova--Riehl--Rovelli, and Pavlov's distinct cube obstruction. We did not locate the following precise family there; the proposition is an elementary extension of the appendix, without a claim of priority.

\textbf{Proposition (one fixed obstruction in every dimension).}\ours
Let $d\geq2$, let $n_j\geq2$, and put
\[
 X_d=\prod_{j=1}^d\Phi_{n_j},\qquad
 U_d=\bigcup_{j=1}^d
   \Phi_{n_1}\times\cdots\times\dot\Phi_{n_j}\times\cdots\times\Phi_{n_d}.
\]
There is a finite weak equivalence $U_d\to U'_d$ whose pushout along $U_d\hookrightarrow X_d$ has target $X'_d$ with $BX'_d\simeq S^3$. In particular every inclusion $U_d\hookrightarrow X_d$ fails to be a $\Winf$-cofibration.

\emph{Proof.} Fix the sink $1$ in each of the last $d-2$ factors. Let $V$ be the three-object full subposet with first two coordinates $(0,0),(1,0),(0,1)$ and all other coordinates $1$. Let $S$ add the point $(1,1,1,\ldots,1)$. Thus $V\cong\Psi$ and $S\cong[1]\times[1]$. Every point of $V$ lies in $U_d$, while the fourth point of $S$ does not. Since each coordinate $1$ is maximal in its alternating interval, all successors of $S$ remain in $S$; all successors of $V$ that remain in $U_d$ remain in $V$. Hence $S\hookrightarrow X_d$ and $V\hookrightarrow U_d$ are closed inclusions.

Their Dwyer structures can be given without induction. The smallest open neighbourhood of $S$ is $\{0,1,2\}^d$. Map its first two coordinates by $r_0(0)=0$, $r_0(1)=r_0(2)=1$, and map every remaining coordinate to $1$. This sends an object to the least element of $S$ above it, and is therefore left adjoint to the inclusion of $S$. The smallest open neighbourhood of $V$ in $U_d$ is
\[
 \bigl(\{0,1,2\}\times\{0\}\;\cup\;
       \{0\}\times\{0,1,2\}\bigr)\times\{0,1,2\}^{d-2}.
\]
The same map lands in $V$ and is its left adjoint. Both inclusions are closed Dwyer inclusions.

Use the finite weak equivalence $p:V\to V'=Z^{\triangleright_e}$ constructed above. Put
\[
 U'_d=U_d\amalg_VV',\quad Q=S\amalg_VV',\quad
 X'_d=X_d\amalg_VV'.
\]
The first Dwyer pushout gives $U_d\simeq U'_d$. Both $BS$ and $BX_d$ are contractible, so $S\to X_d$ is a weak equivalence as well as a closed Dwyer inclusion. Its pushout gives a weak equivalence $Q\to X'_d$. The calculation on manuscript p.~15 identifies $BQ\simeq S^3$, so $BX'_d\simeq S^3$. Finally associativity of pushouts gives $X'_d=X_d\amalg_{U_d}U'_d$. The resulting map from the contractible $BX_d$ to $BX'_d$ is not a weak equivalence, proving the assertion. All categories involved are finite, since the displayed pushouts are obtained by the same finite gluing constructions.\hfill$\square$
\end{dcomment}


\msp{page 17}\textbf{Recollement dans} Cat \textbf{d'un ouvert et d'un fermé complémentaire.}

\texttt{[Supposons donné]}
\[ X_0 \underbrace{\hookrightarrow}_{\text{ouvert}} X \ \underbrace{\hookleftarrow}_{\substack{\text{fermé}\\ \text{complé-}\\ \text{mentaire}}} X_1 \]
dans Cat, \texttt{[et soit]}
\[ \tilde{X}_{01} \subseteq \mathrm{Fl}(X) \]
\texttt{[la]} sous-catégorie pleine des flèches de source dans $X_0$, but dans $X_1$.

\dgrblock{dgrb624}{The arrows which join the two pieces}
\begin{dcomment}{exposition}{p17:sous-catégorie pleine des flèches de source}{The arrows which join the two pieces}
Grothendieck now turns from unsuccessful pushouts to a decomposition which retains the attaching data. An open subcategory $A=X_0$ is a full sieve; its complementary closed subcategory $B=X_1$ is a full cosieve. Thus every object belongs to exactly one piece, and there are no arrows from $B$ to $A$. The additional information in $X$ is the functor
\[
 M:A^{\mathrm{op}}\times B\longrightarrow\Ens,
 \qquad M(a,b)=\Hom_X(a,b).
\]
In the language of distributors, $X$ is the collage of $M$; this description, and its simplicial counterpart, are given in \citet[\S\S7.1--7.2, pp.~309--311]{Joyal2008Quasi}. The direction of these attaching arrows is the feature that permits the argument below.

Write $E=\widetilde X_{01}$. Its objects are arrows $f:a\to b$ with $a\in A$ and $b\in B$; a morphism $f\to f'$ is a commuting square
\[
 (u,v):f\longrightarrow f',\qquad
 u:a\to a',\quad v:b\to b',\quad vf=f'u.
\]
Both vertical sides point in their ordinary directions. In particular $E$ is the ordinary crossing-arrow category, not a twisted-arrow category. The projection $p:E\to A$ is a split Grothendieck fibration: pullback along $u:a\to a'$ sends $f':a'\to b$ to $f'u:a\to b$. Its fibre is $a\backslash B$. Dually, $q:E\to B$ is a split opfibration, with fibre $A/b$ and transport given by postcomposition. The arrows $f$ themselves constitute a natural transformation $ip\Rightarrow jq$.

After realization this natural transformation provides a specified homotopy between the two maps $BE\rightrightarrows BX$, where $BC=|NC|$. Consequently it defines a map from the double mapping cylinder of $BA\leftarrow BE\to BB$ to $BX$. This is the comparison whose weak equivalence is at issue. The ordinary square does not commute strictly, and replacing its specified homotopy by an arbitrary one would change the question.

The easy cases stated in the manuscript have a useful explanation. If $i$ is aspherical, every $A/b$ is contractible, hence $q$ is a weak equivalence; the homotopy pushout then has the homotopy type of $BA$, which maps equivalently to $BX$. If $j$ is coaspherical, the dual argument uses the fibres $a\backslash B$ of $p$. The general argument must retain these fibres rather than assume their contractibility.
\end{dcomment}


\textbf{Question.} \textit{Ce diagramme}
\[
\begin{tikzcd}[column sep=small]
 & & \tilde{X}_{01} \arrow[dll, "p"' near start, "\text{fibrant, fibres les $x_0\backslash X_1$}"' near end] \arrow[drr, "q" near start, "\text{cofibrant, fibres les $X_0/x_1$}" near end] & & \\
X_0 \arrow[ddrr, "i"'] & & & & X_1 \arrow[ddll, "j"] \\
 & {} \arrow[rr, bend left] & & {} & \\
 & & X, & &
\end{tikzcd}
\tag*{$(*)$}
\]
\printedpage{13}
% head: A. Grothendieck : Dérivateurs, Ch. XII, Édition des Universités de Montpellier II et Paris VII
% foot: Version 20 juin 2003
\textit{est-il $W$-cocartésien?}

OK si $i$ $W$-asphérique ou j $W$-coasphérique, i.e. $X \longrightarrow \boldsymbol{\Delta}^1$ lisse ou propre -- c'est un cas où on a soit $X_0 \xrightarrow{W} X$, soit $X_1 \xrightarrow{W} X$. Mais même en supposant $i$ ou $j$ dans $W$, et de plus $X$ asphérique (donc $X_0$ et $X$ asphériques, ou $X_1$ et $X$ asphériques), et $W = \mathbf{W}_\infty$, je n'ai pu prouver que le diagramme est cocartésien, i.e. $q \in W$ resp. $p \in W$ (\footnote[9]{Je suppose que $W$ satisfait Loc (6).}). (Pas non plus si $X_0 \simeq e$, $X$, $X_1$ $W$-asphérique $\Longrightarrow$ $e\backslash X$ $W$-asphérique, i.e. $X_1 \longrightarrow X$ $W$-coasphérique $\ldots$\,) Mais j'ai prouvé du moins, dans tous les cas, que % [sic] j (upright) -> $j$, letter denoting a functor printed in roman type
\[
\begin{tikzcd}[row sep=small]
 & X_0 & & \\
\mathcal{X} \overset{\text{déf}}{=} \displaystyle\int \tilde{X}_{01} \arrow[ur] \arrow[dr] & & {} \arrow[r, "\varphi"] & X \\
 & X_1 & &
\end{tikzcd}
\tag*{$(**)$} % verified: CMSY10 glyph, \mathcal{X}; corrections.csv, PDF leaf 14
\]
est dans $\mathbf{W}_1$.

Chose \textit{fausse} (sauf si $i$ ou $j$ dans $W = \mathbf{W}_\infty$) si on remplace $\tilde{X}_{01}$ par $\tilde{X}'_{01}$, formée des triples $(x_0, x_1, c)$, avec $x_0 \in \mathrm{Ob}\,X_0$, $x_1 \in \mathrm{Ob}\,X_1$, $c : x_0 \dotto x_1$ un \textit{chemin} de $x_0$ à $x_1$ dans $X$! On a
\msp{page 18}donc une situation ‘topologique’ tout à fait spéciale à la catégorie de modèles Cat, qui ne semble pas avoir d'analogue dans des catégories telles que Esp ou $\boldsymbol{\Delta}^\wedge$, où pourtant les notions d'ouvert et de fermé complémentaire (avec grain de sel dans $\boldsymbol{\Delta}^\wedge$) ont un sens.

\dgrblock{dgrb625}{Acyclic links and a false converse}
\begin{dcomment}{development}{p17}{Acyclic links and a false converse}
The proposed $\mathbf W_\infty$ gluing formula and the opposite-side reflection invoked by the manuscript's ``i.e.'' have different strengths. The formula is true. Footnote~9 adds Loc~$(6)$, a numbered hypothesis which requires qualification. If it uses the full reversible convention of Chapter~XVI, Corollary~3.7, then cocartesianness together with $i\in W$ does imply $q\in W$, and dually for $j,p$: this is the additional condition Loc~$(6\mathrm C)$ \citep[\S3, Corollary~3.7, ms.~26]{GrothendieckDerivateursXVI}. Under that interpretation the conditional inference is valid, but its specialization to $\mathbf W_\infty$ is incompatible with the axiom. The following construction proves the failure for ordinary weak homotopy equivalences.

\textbf{Proposition (a contractible union with a noncontractible link).}\ours
There is a finite poset $X$, a one-object sieve $A=\{a\}$, and its complementary cosieve $B$, such that $BA$, $BB$, and $BX$ are contractible, whereas $B\widetilde X_{01}$ is not contractible. In particular neither $p$ nor $q$ need be a weak equivalence, even when both $i$ and $j$ are weak equivalences. Moreover, $\mathbf W_\infty$ fails Loc~$(6\mathrm b)$ in the convention of Chapter~XVI.

\emph{Proof.} We first choose a connected acyclic finite complex $P$ with nontrivial fundamental group. A standard source of such complexes is \citet[Example~2.38]{Hatcher2002}. For explicitness, take the presentation complex
\[
 \langle s,t\mid s^3(st)^{-2},\ t^5(st)^{-2}\rangle.
\]
It has one vertex, two edges and two two-cells. Its cellular differential in positive degrees is
\[
 0\longrightarrow\mathbb Z^2
 \xrightarrow{\left(\begin{smallmatrix}1&-2\\-2&3\end{smallmatrix}\right)}
 \mathbb Z^2\longrightarrow0.
\]
The determinant is $-1$, so its reduced integral homology vanishes. The assignments $s\mapsto(1\,2\,3)$ and $t\mapsto(1\,4\,3\,2\,5)$ define a nontrivial homomorphism to $A_5$: their product is $(1\,4)(2\,5)$, and both relations hold. Thus $P$ is not simply connected. By the finite form of the simplicial replacement theorem \citep[Theorem~2C.5]{Hatcher2002}, choose a finite simplicial complex $L\simeq P$ and let $K$ be its poset of nonempty faces. Its order complex is the barycentric subdivision of $L$, so $BK\simeq P$.

Adjoin two incomparable new elements $a,b$ to $K$, each smaller than every element of $K$, and call the resulting poset $X$. Set $A=\{a\}$ and $B=\{b\}\cup K$. Then $A$ is a sieve and $B$ is its complementary cosieve; $b$ is initial in $B$. The nerve of $X$ is the union of two cones on $NK$, along $NK$ itself. Hence
\[
 BX\cong\Sigma BK.
\]
This suspension is simply connected because $BK$ is connected, and its reduced homology vanishes by the suspension isomorphism. The simply connected homology form of Whitehead's theorem therefore makes it contractible \citep[Corollary~4.33]{Hatcher2002}.

The crossing arrows are exactly the unique arrows $a\to k$, for $k\in K$. Their commuting squares are exactly the order relations of $K$. Thus $E\cong K$, $p:K\to\{a\}$, and $q:K\hookrightarrow\{b\}\cup K$. Neither map is a weak equivalence, since $\pi_1(BK)$ is nontrivial.

For the last assertion, consider the full sieves
\[
 U=(\{a\}\cup K)^{\mathrm{op}},\qquad
 V=(\{b\}\cup K)^{\mathrm{op}}
 \quad\text{in }X^{\mathrm{op}}.
\]
They cover $X^{\mathrm{op}}$ and intersect in $K^{\mathrm{op}}$. Map their Mayer--Vietoris square to the square whose four vertices are the terminal category. The component maps from $U,V,X^{\mathrm{op}}$ are weak equivalences: their classifying spaces are two cones and the contractible suspension. The component from $K^{\mathrm{op}}$ is not. This violates precisely the inverse cube implication Loc~$(6\mathrm b)$; the full axiom includes that implication \citep[\S3, ms.~$16'$--$17'$, footnote~13]{GrothendieckDerivateursXVI}.\hfill$\square$

For $\mathbf W_\infty$, this also answers negatively the subsidiary question with $X_0\simeq e$ and $X,X_1$ aspherical. Its intended comma category is $a\backslash X_1=K$; the printed $a\backslash X$ would contain the initial object $\id_a$ and would be automatically contractible. The example does not contradict the gluing conjecture: its homotopy pushout is exactly the contractible suspension $\Sigma BK$. It disproves opposite-side reflection for $\mathbf W_\infty$, while leaving the stronger conditional principle intact.

The axiom labels cannot be identified solely by their numbers. Chapter~XV, \S9, ms.~117, calls a fibration axiom Loc~$(6)$ \citep{Grothendieck1990XV} and refers to Chapter~XII, \S IV, ms.~53--59; there, variant $\mathrm a'$ is a pullback condition for $W_\omega$-perfect functors \citep{Grothendieck1990XIIa}. Chapter~XVI instead writes the forward Mayer--Vietoris condition as Loc~$(6'\mathrm a)$, whereas the appendix's last line prints Loc~$(6\mathrm a')$ \citep[\S3, ms.~$17'$]{GrothendieckDerivateursXVI}. We retain these distinct labels and cannot certify that footnote~9 invokes the later reversible convention. We checked the standard acyclic-complex construction and the homotopy-colimit sources cited below; we did not locate this precise finite-poset application there. The marker records our worked application, without a claim of priority.
\end{dcomment}


L'argument catégorique quasi-trivial qui établit que $(**)$ est dans $\mathbf{W}_1$, i.e. induit une équivalence de catégories entre $\mathrm{Rev}(X)$ et
\[ \mathrm{Rev}(\mathcal{X}) \ \simeq \ \mathrm{Rev}(X_0) \overset{2}{\times}_{\mathrm{Rev}(\tilde{X}_{01})} \mathrm{Rev}(X_1), \]
me semble devoir s'étendre aux dimensions supérieures, à coups de $n$-gerbes sur $X$ -- mais les fondements sur la théorie des $n$-gerbes manquent. (Et je n'ai pas eu le courage d'expliciter l'argument même pour $n = 2$ (où pourtant \textsc{Giraud} fournit ce qu'il faut) (\footnote[10]{Mais j'ai fait quand même une esquisse de vérification pour $n = 2$, qui me convainc qu'on a bien $\varphi \in \mathbf{W}_2$. Et je n'ai pas guère de doute qu'on a même $\varphi \in \mathbf{W}_\infty$.}).) Aussi je conjecture que $\varphi$ est tout au moins dans $\mathbf{W}_\infty$, i.e. que le diagramme $(*)$ est $W$-cocartésien si $W \supseteq \mathbf{W}_\infty$ (et satisfait Loc (6 a$'$)). % [sic] je n'ai pas guère (footnote 10) -> je n'ai guère, double negation pas + guère


\dgrblock{dgrb626}{The open–closed gluing formula}
\begin{dcomment}{development}{p18}{The open--closed gluing formula}
The final conjecture $\varphi\in\mathbf W_\infty$ has an affirmative answer. Its proof comes from the fact that a simplex in $NX$ crosses from $A$ to $B$ at most once. We give the precise specialization of Thomason's homotopy-colimit comparison that settles this question. We did not locate this exact formulation in the checked sources; the result below is presented as a consequence of an established theorem, without a claim of priority.

\textbf{Proposition (directed gluing).}\ours
For every small category $X$ with a full sieve $A$ and complementary full cosieve $B$, let $E$ be its crossing-arrow category. The natural transformation $ip\Rightarrow jq$ induces a weak homotopy equivalence
\[
 \operatorname{hocolim}(BA\longleftarrow BE\longrightarrow BB)
 \longrightarrow BX.
\]
Consequently the comparison $\varphi:\mathcal X\to X$ of the manuscript belongs to $\mathbf W_\infty$.

\emph{Proof.} To keep track of the unique crossing, introduce the bisimplicial set
\[
 S_{r,s}=\left\{
 a_0\longrightarrow\cdots\longrightarrow a_r
 \xrightarrow{f}b_0\longrightarrow\cdots\longrightarrow b_s
 \ \middle|\ a_i\in A,\ b_j\in B
 \right\},\qquad r,s\geq0.
\]
Its faces compose or omit arrows on the indicated side, and its degeneracies insert identities. The last horizontal face precomposes $f$; the first vertical face postcomposes $f$. Thus $S$ is the simplicial replacement of the contravariant functor
\[
 A^{\mathrm{op}}\longrightarrow\Cat,\qquad
 a\longmapsto a\backslash B.
\]
Its Grothendieck construction, with cartesian variance, is exactly $E$: a morphism over $u:a\to a'$ is a map $f\to f'u$ in $a\backslash B$, namely a $v$ with $vf=f'u$.

There is a comparison $\theta:\operatorname{diag}S\to NE$. On an $n$-simplex it gives the chain of crossing arrows
\[
 a_i\longrightarrow a_n\xrightarrow{f}b_0
 \longrightarrow b_i,\qquad 0\leq i\leq n,
\]
whose successive morphisms have sides $a_i\to a_{i+1}$ and $b_i\to b_{i+1}$. This is Thomason's comparison for the preceding contravariant functor. It is a weak equivalence; the precise variance and this formula are verified in \citet[\S3.1, equation~(9), and Theorem~3]{Cegarra2020}, which gives an independent proof of the comparison originally due to \citet[Theorem~1.2]{Thomason1979}. Both endpoint maps commute with $\theta$: they simply forget the $b$-chain or the $a$-chain.

It remains to identify what the mixed simplices assemble to. Let $R=\lVert S\rVert$ be the double geometric realization. It has maps $r_A:R\to BA$ and $r_B:R\to BB$ obtained by forgetting one side. The natural diagonal-to-double-realization homeomorphism identifies $R$ with $|\operatorname{diag}S|$. Form
\[
 D(R)=\bigl(BA\amalg(R\times[0,1])\amalg BB\bigr)
       /\bigl((z,0)\sim r_A(z),\ (z,1)\sim r_B(z)\bigr).
\]
For a mixed simplex as above, send $(u,v,t)\in\Delta^r\times\Delta^s\times[0,1]$ to the point of $\Delta^{r+s+1}$ with coordinates
\[
 ((1-t)u_0,\ldots,(1-t)u_r,tv_0,\ldots,tv_s).
\]
This respects faces and degeneracies, and respects the endpoint identifications. It supplies a homeomorphism $D(R)\cong BX$. Indeed, in a point with positive total weight on both pieces, $t$ is the total weight on the $B$ vertices; dividing the two sets of coordinates by $1-t$ and $t$ recovers $u,v$. At $t=0,1$ precisely the identifications defining $D(R)$ occur. This description is compatible with the realization quotient topology because it is the usual simplex-join homeomorphism on every simplex. Pure simplices give the two endpoint spaces.

The double mapping cylinder computes a homotopy pushout. Replacing its middle space $R$ by the weakly equivalent CW complex $BE$ therefore preserves its homotopy type, by homotopy invariance of this construction. The resulting map to $BX$ is the canonical one: the homotopy obtained from $ip\Rightarrow jq$, after composition with $\theta$, is homotopic relative to its endpoints to the displayed barycentric homotopy. On each prism, both lie in the same mixed simplex, and linear interpolation supplies compatible homotopies on its faces. Thus the canonical comparison, and not merely an abstract map between the same homotopy types, is a weak equivalence.

Finally, let $J$ be the span category with arrows $*\to0$ and $*\to1$, and take its covariant diagram with values $E,A,B$ and structure maps $p,q$. Its Grothendieck construction models $\mathcal X$. A comparison to $X$ sends its $A$- and $B$-objects to themselves and its $E$-object $f:a\to b$ to $a$; the arrows towards the $B$-part use $f$. Thomason's theorem identifies its classifying space with the preceding homotopy pushout, compatibly with this cocone. Hence $B\varphi$ is a weak equivalence.\hfill$\square$

The cases $A=\varnothing$ or $B=\varnothing$ are included: then $E$ is empty and the assertion reduces to the other piece. For every localizer $W$ containing $\mathbf W_\infty$, the comparison just proved also lies in $W$. If the span Grothendieck construction represents the $W$-homotopy pushout, this proves $W$-cocartesianness as well. For a fundamental localizer, Chapter~XVI explicitly supplies that interpretation under its forward axiom Loc~$(6'\mathrm a)$ \citep[\S3, ms.~24, footnote~15]{GrothendieckDerivateursXVI}. This comparison hypothesis requires no opposite-side reflection. The manuscript cites Loc~$(6\mathrm a')$ for its general-$W$ interpretation; because of the numbering ambiguity just documented, we do not silently identify that label with the distinct forward or full axioms of Chapter~XVI. The theorem for $\mathbf W_\infty$ is independent of that identification.

The proof also clarifies the preceding observation about a phenomenon peculiar to categories. Its realization is an explicit topological double mapping cylinder, and the mixed-simplex construction has a directed simplicial formulation. What a bare pair of complementary topological subsets does not supply is the specified category of attaching arrows.
\end{dcomment}


\dgrblock{dgrb627}{Why replacing arrows by all paths changes the answer}
\begin{dcomment}{examples}{p18}{Why replacing arrows by all paths changes the answer}
Grothendieck's caution about paths concerns the attaching information, rather than an insufficiency of path spaces as models of homotopy types. A bridge of directed arrows records how the two pieces were assembled. A bridge of all paths in the already assembled space usually records a homotopy pullback instead. These two bridges can differ even in the smallest example with a loop.

Let $X$ have objects $a,b$ and precisely two nonidentity arrows $\lambda,\mu:a\rightrightarrows b$. Take $A=\{a\}$ and $B=\{b\}$. There are no nonidentity commuting squares between its two crossing arrows, so $E$ is the discrete two-point category. The nerve of $X$ has two vertices and two edges and no nondegenerate higher simplices. Thus
\[
 BX\cong S^1,\qquad
 \operatorname{hocolim}(*\longleftarrow\{\lambda,\mu\}\longrightarrow*)
 \cong S^1.
\]

Now replace the bridge by a model of the space of all paths from $a$ to $b$ in $BX$. Choosing one reference path identifies this space up to homotopy with $\Omega S^1$, whose components are contractible and indexed by $\mathbb Z$. The new homotopy pushout is therefore the graph with two vertices and one edge $e_k$ for each $k\in\mathbb Z$. Collapsing $e_0$ gives
\[
 \bigvee_{k\in\mathbb Z\setminus\{0\}}S^1.
\]
The comparison to $S^1$ sends the loop represented by $e_ke_0^{-1}$ to the integer $k$. Its map on fundamental groups is the homomorphism from the free group on these generators to $\mathbb Z$ which sends the generator indexed by $k$ to $k$. It is surjective but not injective; for example the word consisting of the generator indexed by $2$ followed by the inverse square of the generator indexed by $1$ lies in its kernel and is a nonempty reduced word. Thus even the $\mathbf W_1$ conclusion fails.

More generally, when a path category models $BA\times^h_{BX}BB$, the exception stated in the manuscript is valid: if $BA\to BX$ is a weak equivalence, its projection to $BB$ is a weak equivalence, so the associated homotopy pushout maps equivalently to $BX$; the other case is dual. This is a sufficient condition, not a characterization of every exceptional case. The manuscript specifies the objects of $\widetilde X'_{01}$ but does not spell out its morphisms. Our computation concerns the usual homotopically correct all-paths interpretation; it does not assert an equivalence for an unspecified raw category of paths.
\end{dcomment}


\dgrblock{dgrb628}{From covering descent to higher local systems}
\begin{dcomment}{later}{p18}{From covering descent to higher local systems}
The last paragraph identifies the right source of confidence in the conjecture: a local system can be reconstructed from its restrictions to the two pieces and its transport across every attaching arrow. We can make the argument for coverings completely explicit before passing to higher coefficients.

Use covariant transport to identify $\operatorname{Rev}(C)$ with the category of functors $C\to\Ens$ which send every arrow to a bijection. An object of
\[
 \operatorname{Rev}(A)
 \mathbin{\mathop{\times}\limits^{2}_{\operatorname{Rev}(E)}}
 \operatorname{Rev}(B)
\]
is a pair $L_A,L_B$ with an isomorphism $\alpha:p^*L_A\xrightarrow{\sim}q^*L_B$. Define $L$ on the two full subcategories by $L_A,L_B$, and on a crossing arrow $f:a\to b$ by $L(f)=\alpha_f$. To check functoriality, there are only two mixed cases. The squares $(u,\id):fu\to f$ and $(\id,v):f\to vf$ in $E$ give
\[
 \alpha_{fu}=\alpha_f L_A(u),\qquad
 \alpha_{vf}=L_B(v)\alpha_f.
\]
They are exactly the two composition laws needed in $X$, since a composable chain cannot cross from $A$ to $B$ twice. All transports are bijective. A compatible pair of morphisms between such descent data gives a natural transformation of the reconstructed functors, and restriction is inverse to this construction. This proves the equivalence of covering categories in the manuscript; applying the same gluing description to the categorical mapping cylinder gives its displayed formula for $\operatorname{Rev}(\mathcal X)$.

This argument also identifies the fundamental groupoid directly. Its generators are the arrows internal to the two pieces and the crossing arrows, now formally invertible; the equations just verified supply all mixed composition relations. Consequently the comparison $\varphi$ induces an equivalence of fundamental groupoids. The source offers a sketch of this $\mathbf W_1$ assertion and reports a further sketch for $\mathbf W_2$; the preceding mixed-simplex proof establishes all degrees at once.

The contemporary interpretation of the proposed higher-gerbe argument uses local systems valued in spaces or in $n$-groupoids. If $\mathcal S$ denotes the $\infty$-category of spaces, define $\operatorname{Loc}(C)=\operatorname{Fun}(BC,\mathcal S)$, viewing $BC$ as its homotopy type. The proved pushout formula, by the universal property of a colimit, yields
\[
 \operatorname{Loc}(X)\simeq
 \operatorname{Loc}(A)\times^h_{\operatorname{Loc}(E)}
 \operatorname{Loc}(B).
\]
One can restrict the values to $n$-truncated spaces to obtain every finite stage. Straightening gives the precise realization of space-valued coefficients, and higher descent formalizes such pullbacks of coefficient categories; see \citet[Theorem~2.2.1.2, Corollary~3.3.4.6, and Theorem~6.1.3.9]{Lurie2009HTT}. Giraud's theory supplied the stack and gerbe language for the first nonabelian stages \citep{Giraud1971}; the claim here is about locally constant higher coefficients, not an identification of all stacks with local systems. The categorical question of the appendix is therefore settled without requiring a separate construction and comparison of $n$-gerbes for each $n$.
\end{dcomment}
\begin{otherlanguage}{english}
\dgroverview{XIIapp}{Chapter XII, appendix \textperiodcentered\ \textup{XIIapp} \textperiodcentered\ Overview}
\selectlanguage{english}
Grothendieck's appendix to Chapter XII examines a difficulty that becomes visible as soon as one tries to use small categories as models for spaces: an ordinary categorical pushout can discard the homotopy information that a homotopy pushout must retain. The investigation begins with a proposed contractibility criterion, passes through explicit finite counterexamples, and ends with a conjectural gluing theorem. Its exploratory order matters. An unsuccessful cohomology calculation and an impossible retraction are part of the route to the examples; correcting them changes some intermediate claims without destroying the main obstruction. The electronic edition already marks several of these revisions \citep{Grothendieck1990XIIapp}.

Throughout these pages \(BC=|NC|\) denotes the classifying space of a small category. The localizer \(\mathbf W_\infty\) consists of the functors inducing weak homotopy equivalences. The word ``aspherical'' means weakly contractible in this usage. A presheaf \(F:C^{\mathrm{op}}\to\Ens\) has a homotopy type through its category of elements \(C/F\). In particular, a representable presheaf is aspherical, since its category of elements is a slice \(C/c\) with a terminal object. These conventions connect the appendix to the integration of homotopy types described in \emph{Pursuing Stacks}, \S69, and to the adjoints-to-restriction viewpoint of the 1991 letter to Thomason \citep{Grothendieck1983,Grothendieck1991Letter}.

The opening hypothesis asks for two weakly terminal objects \(b,c\), meaning that every object admits a map to each, and for their categorical product \(a=b\times c\). The source proves that \(BC\) is simply connected. Its proof is sound: in the groupoid completion, cancellation against the two projections identifies all maps from an object to \(a\); those unique classes are compatible with every arrow and its inverse. The completion is therefore equivalent to the terminal groupoid. The proposed stronger conclusion, weak contractibility, is false, but the first example in the manuscript needs a different calculation from the one printed.

That example is the category \(X_G\) of the first two formal powers of an object with endomorphism monoid \(G\):
\[
 \Hom_{X_G}(b^m,b^n)=\bigl(\{1,\ldots,m\}\times G\bigr)^n,
 \qquad m,n\in\{1,2\}.
\]
It is an opfibration over the opposite of the full finite-set category on a singleton and a two-element set, with fibres the one-object categories of \(G\) and \(G^2\). The resulting cohomological spectral sequence is legitimate. The provisional derived-limit formula used to evaluate it is not: the transposition of the two-element set and the common retraction of its two inclusions supply relations that the proposed two-term complex has omitted. Even a constant coefficient diagram contradicts the displayed cokernel formula.

The commentary gives an explicit projective resolution for the actual indexing category. It proves that \(BX_G\) is \(2\)-connected for every monoid \(G\), while
\[
 H^3(BX_{C_2};\mathbb F_2)\cong\mathbb F_2.
\]
Thus the two-object, twenty-six-arrow category remains a finite counterexample to the opening proposition. Its obstruction occurs one degree later than the manuscript asserts. The repaired computation is supplied as a commentators' development, with complete proofs and a separate normalized-nerve verification. The standard spectral-sequence and homotopy-colimit ingredients are identified explicitly \citep{Cegarra2020,Thomason1979}.

A second development supplies the degree-two counterexample promised by footnote~2 under the exact product hypotheses. For an integer \(r\geq1\), adjoin a new source \(d\) to the category of the first two powers for the trivial monoid, with \(r\) maps \(d\to b\) and the corresponding \(r^2\) maps \(d\to b^2\). Composition selects coordinates. The product remains an actual categorical product, and its new source link retracts up to homotopy onto the poset of nonempty subsets of an \(r\)-element set of cardinality at most two. Its nerve is the subdivided complete graph. The resulting three-object category \(Y_r\) satisfies
\[
 BY_r\simeq\bigvee^{(r-1)(r-2)/2}S^2.
\]
For \(r=3\) this gives a twenty-one-arrow model of \(S^2\). It verifies the footnote without confusing the opening product condition with the weaker presheaf condition used subsequently.

The persistent indexing shape is the span \(\Psi=(b_0\leftarrow a_0\to c_0)\). For a span of representables in \(C^\wedge\), its value at \(x\) is the set of components of a bipartite graph: the vertices are the maps \(x\to b\) and \(x\to c\), and the edges are the maps \(x\to a\). A terminal ordinary pushout says that all these graphs are connected. It does not say that their cycles vanish. The commentary proves both useful consequences stated in the source: a terminal pushout forces simple connectivity of \(BC\), and if either arrow out of \(a\) is monic, the presheaf pushout is aspherical. The latter follows because the graph components are then stars, and Quillen's Theorem A applies to their incidence categories \citep{Quillen1973}.

The four-object example \(Z\) makes the missing cycles visible. It has two parallel arrows \(0\rightrightarrows a\), followed by arrows \(a\to b,c\) which equalize them. Its presheaf pushout \(b\amalg_a c\) is terminal. Its nerve, however, is two cones on the parallel-arrow circle, and hence \(BZ\cong S^2\). The proposed retraction onto the circle cannot be a functor: applying it to either equalizing relation would identify the two distinct arrows. Footnote 7 corrects the homotopy type, but the later statement \(\pi_1(Z)\cong\mathbb Z\) repeats the earlier error. We retain both printed readings and explain their mathematical status.

Adding two incomparable cone vertices to \(Z\) gives the six-object category \(Q\) with \(BQ\cong S^3\). This is the reliable finite obstruction in the later pushout square: its three input categories are contractible, while its ordinary categorical pushout is not. The assertion is established directly from the nerve, without an unrestricted converse to a comma-category theorem. Such a converse would be false: a connected acyclic noncontractible link can have contractible suspension. The appendix's extension to arbitrary \(W\subseteq W_1\) therefore does not follow from this example; the obstruction for this pushout lies in degree three.

Grothendieck then inserts this local defect into the boundary of a product \(\Phi_n\times\Phi_m\) of alternating directed intervals. The boundary has classifying space \(S^1\), while the product is contractible. Its slices over its categorical closure are contractible, but four of them have no terminal object. Accordingly the boundary inclusion satisfies the local asphericity condition of Chapter XII and fails the stronger Dwyer neighbourhood-adjoint condition. Two surrounding closed Dwyer inclusions allow the six-object defect to survive the enlargement. The necessary pushout comparison is Thomason's; Cisinski's later correction distinguishes Dwyer maps from their retracts \citep{Thomason1980,Cisinski1999Dwyer}. The resulting inclusion fails to preserve weak equivalences under cobase change. In contemporary terminology it is not an \(h\)-cofibration.

The same construction extends, with proof here, to the boundary of every product of at least two such intervals: fixing interior sinks in the remaining coordinates carries the same \(S^3\) defect into higher dimensions. Later theory gives two complementary perspectives. Pushouts along Dwyer maps are preserved even as \((\infty,1)\)-categorical pushouts \citep{HackneyOzornovaRiehlRovelli2024}; no model structure on \(\Cat\) with Thomason weak equivalences can be cartesian \citep[Proposition~4.3]{Pavlov2024Thomason}. These results explain the continuing significance of the appendix's distinction between contractible links and actual adjoint neighbourhoods.

The question whether total asphericity forces a terminal object has a simpler answer. Here ``total'' quantifies over every left Kan extension \(f_!1_A\), a convention expressly distinguished in the edition from the total asphericity of \emph{Pursuing Stacks}. Filtered categories satisfy this condition, because the category of elements of a filtered colimit of representables is filtered. Grothendieck already records the correction elsewhere in the manuscript \citep[Chapter II, Corollaries~2.17--2.18]{Grothendieck1990II}. Both \((\mathbb N,\leq)\) and the one-object category of an unsplit idempotent give examples without a terminal object. The surrounding classification is discussed in the commentary on Chapter~II, manuscript pages~35--36, and is not repeated here.

The closing conjecture concerns a different, homotopically correct gluing. If \(A\) is a full sieve of \(X\), \(B\) its complementary full cosieve, and \(E\) the category of arrows from \(A\) to \(B\) with ordinary commuting squares, then
\[
 BX\simeq\operatorname{hocolim}\bigl(BA\longleftarrow BE\longrightarrow BB\bigr).
\]
The comparison is defined by the natural transformation supplied by the crossing arrows. We prove the formula using mixed simplices of the nerve: each crosses from \(A\) to \(B\) once, and its barycentric coordinates split into a point on each side and an interval coordinate. Thomason's comparison identifies the resulting bisimplicial bridge with \(BE\). This settles the \(\mathbf W_\infty\) conjecture and supplies the whole hierarchy of local-system descent statements at once. The exact specialization is worked out here; the general comparison theorem predates the manuscript, and no claim about priority for this formulation is made.

Two qualifications are mathematically necessary. First, for \(\mathbf W_\infty\), the true gluing formula does not imply the subsidiary assertion that one of its bridge projections must be a weak equivalence whenever a boundary inclusion is. The acyclic-link example disproves that implication even for finite posets with contractible pieces and union. A stronger reversible localizer axiom would supply the implication, but \(\mathbf W_\infty\) does not satisfy that axiom. The numbered references in the appendix require caution: Chapters XV and XVI use different Loc~\((6)\) conventions, and the positions of the primes in their variants differ. The detailed commentary records the checked formulations without certifying which convention the appendix intended \citep{Grothendieck1990XV,GrothendieckDerivateursXVI}. Second, replacing crossing arrows by all paths in the completed space changes the bridge: for the parallel-arrow circle, two crossing arrows give \(S^1\), whereas an all-paths bridge gives a countable wedge of circles. The source's covering-descent argument is fully justified by gluing locally constant set-valued functors across commuting squares; higher local systems give its contemporary extension \citep{Lurie2009HTT}. The explicit directed attaching data are what make the final formula work.
\end{otherlanguage}
\dgrchunkend

% DER_XIII-a: Grothendieck's text with the commentary set in, then the overview, inside a group that applies the chunk's own definitions
\dgrchapter{XIII}{Chapter XIII}{Catégories de modèles (1)}
\dgrpart{0}{XIII-a}{Part 1 of 2, manuscript pages 1--60}{XIII-a\enspace\textperiodcentered\enspace Chapter XIII, part 1 of 2\enspace\textperiodcentered\enspace manuscript pages 1--60}{Chapter XIII \textperiodcentered\ \textup{XIII-a}}
\dgrchunkbegin
% DER_XIII-a: the definitions of the chunk's preamble, applied inside its group
\let\printedpage\relax\newcommand{\printedpage}[1]{\marginpar{\footnotesize p.~#1}}
\let\msp\relax\newcommand{\msp}[1]{\marginpar{\footnotesize [#1]}}
\let\Fl\relax\newcommand{\Fl}{\mathrm{Fl}}
\let\M\relax\newcommand{\M}{\mathcal{M}}
\let\cM\relax\newcommand{\cM}{\mathcal{M}}
\let\Cof\relax\newcommand{\Cof}{\mathrm{Cof}}
\let\Fib\relax\newcommand{\Fib}{\mathrm{Fib}}
\let\Ob\relax\newcommand{\Ob}{\mathrm{Ob}}
\let\circled\relax\newcommand{\circled}[1]{\tikz[baseline=(c.base)]\node[draw,circle,inner sep=1pt](c){#1};}
\let\univ\relax\newcommand{\univ}{\mathrm{univ}}
\let\cercle\relax\newcommand{\cercle}[1]{\mbox{\tikz[baseline=(c.base)]\node[draw,circle,inner sep=1pt](c){#1};}}
\let\clarrow\relax\newcommand{\clarrow}{\mathrel{\leftarrow\mkern-4mu{+}\mkern-4mu\rightarrow}}
\let\limind\relax\newcommand{\limind}{\mathop{\underrightarrow{\mathrm{lim}}}}
\let\limproj\relax\newcommand{\limproj}{\mathop{\underleftarrow{\mathrm{lim}}}}
\let\HO\relax\newcommand{\HO}{\mathrm{HO}}

% ---- from Der-ChXIII-25-01-04_p001-005.tex
\printedpage{--} % no page number printed on this page
% head: (none)

\begin{center}
\vspace*{5cm}
{\LARGE\sffamily LES DÉRIVATEURS}

\vspace{1.5cm}
{\large\textbf{ALEXANDRE GROTHENDIECK}}

\vspace{2.5cm}
{\large\textbf{Chapitre XIII}}

\vspace{0.5cm}
{\large\textbf{Catégories de modèles (1)}}
\end{center}

\clearpage

\printedpage{--} % no page number printed on this page
% head: (none)

\vspace*{6cm}
\noindent Ce texte a été déchiffré et transcrit en \LaTeXe{} par M. Künzer. Il a été édité par M. Künzer, J. Malgoire, et G. Maltsiniotis. La transcription est aussi fidèle que possible au manuscrit. Pour les quelques corrections évidentes, ou rares commentaires des éditeurs, ainsi que pour la numérotation originale des pages du manuscrit, les caractères de machine à écrire \texttt{[typewriter]} entre crochets sont utilisés. Un point d'interrogation entre crochets signifie que l'on n'est pas sûr du mot qui précède. Les index, la liste des axiomes, le glossaire, et les références bibliographiques à la fin du chapitre ont été compilés par les éditeurs. Les numéros des pages dans les index sont ceux de l'original.

\noindent Cette édition est provisoire. Les remarques, commentaires et corrections sont bienvenus. Envoyer un message à :

\noindent\textbf{maltsin@math.jussieu.fr}

\hfill G. Maltsiniotis\hspace*{2cm}

\clearpage

\printedpage{1}
% head: A. Grothendieck : Dérivateurs, Ch. XIII, Édition des Universités de Montpellier II et Paris VII
% foot: Version 25 janvier 2004

\msp{page 1}

\section*{1\quad Couples de Quillen $(\Phi, \Psi)$}

Je veux d'abord donner une variante affaiblie de la notion de triple de Quillen (``model category'' de Quillen). Ce ne sera pas une notion tout-à-fait autoduale -- ainsi il y aura \textit{deux} variantes affaiblies, duales l'une de l'autre.

\dgrblock{dgrb666}{Two calculi behind a model structure}
\begin{dcomment}{exposition}{p1:Je veux d'abord donner}{Two calculi behind a model structure}
Grothendieck begins by separating the machinery that constructs replacements from the equivalences those replacements are meant to respect. This separation will make it possible to ask whether the homotopy theory belongs to the pair $(\mathcal M,W)$ itself. The first section therefore contains no weak equivalences: its subject is the elementary calculus of lifting and factorization.

For arrows $i:A\to B$ and $p:X\to Y$, write $i\mathrel{\square}p$ when every commutative square with these vertical sides admits a diagonal $h:B\to X$, satisfying $hi=a$ and $ph=b$, where $a:A\to X$ and $b:B\to Y$ are its horizontal sides. No uniqueness is required. In current notation Grothendieck's $\Phi_*$ is $\Phi^{\square}$, and $\Psi^*$ is ${}^{\square}\Psi$. Consequently
\[
 \Phi\subseteq{}^{\square}\Psi
 \quad\Longleftrightarrow\quad
 \Psi\subseteq\Phi^{\square}.
\]
This is an order-reversing Galois connection between classes of arrows. The arrows in the two displayed boxes operate on classes, rather than on individual morphisms; the editors' correction is mathematically essential, while the missing parenthesis in that correction is merely typographical. For a large category we work relative to universes, or speak of classes instead of forming an unrestricted set of all its subclasses.

There are already two different closure notions in view. A class is \emph{left lifting closed} when it has the form ${}^{\square}\Psi$. Later, a class will be enlarged by isomorphisms, pushouts, transfinite composites and retracts; this is its \emph{weak saturation} in current model-category language. These operations agree for the closure generated by a set in a locally presentable category, by the small object argument. Their identification requires that theorem; it is not a consequence of the Galois connection alone. Conversely, Grothendieck's phrase ``$W$ faiblement saturée'' on manuscript p.~7 means two-out-of-three together with the isomorphisms. Keeping these uses separate prevents a genuine ambiguity in the conjectural argument below.
\end{dcomment}


Mais avant encore, un sorite sur les LLP (left lifting property) et RLP (right lifting property) (\footnote[1]{donner définition en forme, avec renvoi à \textsc{Quillen}.}). Soit $\M$ une catégorie quelconque (le plus souvent, $\M$ sera grande). On considère des parties
\[
\Phi,\ \Psi\ \subset\ \Fl(\M)
\]
et on définit des applications
\[
\boxed{\begin{array}{ccc}
\Phi & \longmapsto & \Phi_* \\[1ex]
\Psi & \longmapsto & \Psi^*
\end{array}}
\qquad
\Fl(\M) \longrightarrow \Fl(\M)
\]
\texttt{[plutôt} $\mathfrak{P}(\Fl(\M) \longrightarrow \mathfrak{P}(\Fl(\M))$\texttt{,}\newline\texttt{où pour tout ensemble} $E$\texttt{,} $\mathfrak{P}(E)$ \texttt{désigne l'ensemble des parties de} $E$\,\texttt{]} par % [sic] \mathfrak{P}(\Fl(\M) -> \mathfrak{P}(\Fl(\M)), a closing parenthesis is missing in the first term
\[
\begin{aligned}
\Phi_* \ &= \ \{\text{ensemble des } p \in \Fl(\M) \text{ qui ont la RLP par rapport à } \Phi\} \\
\Psi^* \ &= \ \{\text{ensemble des } i \in \Fl(\M) \text{ qui ont la LLP par rapport à } \Psi\},
\end{aligned}
\]
et ensuite
\[
\boxed{\begin{array}{ccc}
\Phi & \longmapsto & \widetilde{\Phi} \\[1ex]
\Psi & \longmapsto & \overline{\Psi}
\end{array}}
\qquad
\Fl(\M) \longrightarrow \Fl(\M)
\]
\texttt{[plutôt} $\mathfrak{P}(\Fl(\M) \longrightarrow \mathfrak{P}(\Fl(\M))$\texttt{]} par % [sic] \mathfrak{P}(\Fl(\M) -> \mathfrak{P}(\Fl(\M)), a closing parenthesis is missing in the first term
\[
\begin{aligned}
\widetilde{\Phi} \ &= \ (\Phi_*)^* \\
\overline{\Psi} \ &= \ (\Psi^*)_* \ .
\end{aligned}
\]

\msp{page 2}

Je résume les propriétés les plus immédiates de ces opérations dans une

\noindent\textbf{Proposition 1.}

\begin{enumerate}[label={}]
\item[a)] $\Phi \longmapsto \Phi_*$, $\Psi \longmapsto \Psi^*$ \textit{sont des applications décroissantes de} $\Fl(\M)$ \texttt{[plutôt} $\mathfrak{P}(\Fl(\M))$\texttt{]} \textit{dans lui-même, pour la relation d'ordre} $\subset$\textit{, et} $\Phi \longmapsto \widetilde{\Phi}$\textit{,} $\Psi \longmapsto \overline{\Psi}$ \textit{des applications croissantes.}

\printedpage{2}
% head: A. Grothendieck : Dérivateurs, Ch. XIII, Édition des Universités de Montpellier II et Paris VII
% foot: Version 25 janvier 2004

\item[b)] \textit{On a}
\[
\Phi\ \subset\ \widetilde{\Phi},
\]
\textit{et pour qu'on ait $\Phi = \widetilde{\Phi}$, il faut et il suffit que $\Phi$ soit de la forme $\Psi^*$. On dit alors que $\Phi$ est} \textup{clos à gauche} (\footnote[2]{Donner définition en forme.})\textit{. L'ensemble des parties de} $\Fl(\M)$ \textit{closes à gauche est stable par intersections quelconques} (\footnote[3]{de façon précise, $\bigcap_i(\Psi_i^*) = (\bigcup_i \Psi_i)^*$}), \textit{et $\widetilde{\Phi}$ est la plus petite partie close de} $\Fl(\M)$ \textit{qui contienne $\Phi$ et la plus grande partie $\Phi'$ de} $\Fl(\M)$ \textit{telle que $\Phi'_* = \Phi_*$ (ou telle que $\Phi'_* \supset \Phi_*$).}

\textit{Énoncés symétriques pour $\Psi$ : on a}
\[
\Psi\ \subset\ \overline{\Psi},
\]
\textit{on a $\Psi = \overline{\Psi}$ si et seulement si $\Psi$} \texttt{[est]} \textit{de la forme $\Phi_*$, on dit alors que $\Psi$ est} \textup{clos à droite.} \textit{L'ensemble des parties de} $\Fl(\M)$ \textit{closes à droite est stable par intersections quelconques} (\footnote[4]{de façon précise, on a $\bigcap_i(\Phi_{i\,*}) = (\bigcup_i \Phi_i)_*$}), \textit{et $\overline{\Psi}$ est la plus petite partie close à droite contenant $\Psi$, et aussi la plus grande partie $\Psi'$ de} $\Fl(\M)$ \textit{telle que $\Psi'^* = \Psi^*$ (ou telle que $\Psi'^* \supset \Psi^*$).}

\item[] \hspace*{-2.5em}\msp{page 3}

\item[c)] \textit{Les applications}
\[
\Phi \ \longmapsto\ \Phi_*, \qquad\quad \Psi\ \longmapsto\ \Psi^*
\]
\textit{établissent des bijections inverses l'une de l'autre entre l'ensemble des parties $\Phi$ closes à gauche de} $\Fl(\M)$\textit{, et l'ensemble des parties $\Psi$ closes à droite. C'est un anti-isomorphisme d'ensembles ordonnés. On a alors}
\[
(\Phi \cap \Phi')_* \ = \quad
\begin{array}{l}
\textit{partie close à droite de } \Fl(\M) \ \textit{``engendrée par } \Phi_*,\ \Phi'_*\textit{'',} \\
\text{i.e. \texttt{[la]} \textit{plus petite partie close à droite contenant} } \Phi_*, \\
\Phi'_*, \text{ i.e. } \overline{\Phi_* \cup \Phi'_*}
\end{array}
\]
\texttt{[pour} $\Phi$\texttt{,} $\Phi'$ \texttt{parties closes à gauche]}\textit{, plus généralement}
\[
(\bigcap_i \Phi_i)_* \ = \ \overline{\bigcup_i (\Phi_{i\,*})}
\]
\textit{pour toute famille $(\Phi_i)_{i \in I}$ de parties closes à gauche de} $\Fl(\M)$\textit{. Formules duales} \texttt{[pour des parties closes à droite]}
\[
\left\{
\begin{aligned}
(\Psi \cap \Psi')^* \ &= \ \widetilde{\Psi^* \cup \Psi'^*} \\
(\bigcap_j \Psi_j)^* \ &= \ \widetilde{\bigcup_j (\Psi_j^*)}\ .
\end{aligned}
\right.
\]

\item[d)] \textit{Si} $\Phi, \Psi \subset \Fl(\M)$\textit{, alors on a}
\[
\Phi \subset \Psi^* \iff \Psi \subset \Phi_*, \qquad\quad \textit{on écrit alors } \Phi \longleftrightarrow \Psi,
\]
\textit{cela signifie que $\Phi$ a la LLP pour $\Psi$, ou que $\Psi$ a la RLP pour $\Phi$.}
\end{enumerate}

\printedpage{3}
% head: A. Grothendieck : Dérivateurs, Ch. XIII, Édition des Universités de Montpellier II et Paris VII
% foot: Version 25 janvier 2004

\msp{page 4}

\dgrblock{dgrb667}{Closure, transfinite lifting and the retract argument}
\begin{dcomment}{exposition}{p3}{Closure, transfinite lifting and the retract argument}
Grothendieck wants the formal part of Quillen's method to be available before choosing any examples. Proposition~1 provides exactly this preparation. To see its content directly, let $L(R)={}^{\square}R$ and $R(L)=L^{\square}$. Antitonicity follows by adding lifting problems. Every arrow of $L$ lifts against every arrow of $R(L)$, so $L\subseteq L(R(L))$. Applying antitonicity gives $R(L(R(L)))\subseteq R(L)$, while the reverse inclusion is the same elementary observation with the two roles exchanged. Thus $RLR=R$ and $LRL=L$. These identities prove idempotence, describe the fixed points and show that the two fixed-point posets are oppositely isomorphic. For instance,
\[
 R\left(\bigcap_iL_i\right)
   =\overline{\bigcup_iR(L_i)}
\]
when each $L_i$ is left lifting closed. The hypothesis matters: the right side is the join in the lattice of right lifting closed classes, not their ordinary union.

The stability assertions in Proposition~2 express how to transport a chosen solution. Isomorphisms lift by inversion. A lifting problem for a composite is solved successively. A lifting problem for a pushout is equivalent, by its universal property, to one for the original map. For a continuous ordinal chain, begin with the prescribed map on its initial object, extend at each successor, and take the unique induced map at each limit. Transfinite recursion then supplies the diagonal against the composite. The existence of the colimits appearing in the chain is the only colimit hypothesis used here; no presentability has yet entered.

The retract assertion deserves an explicit formulation, since it is the bridge from a constructed factorization to a lifting characterization. Suppose $f:A\to B$ factors as $A\xrightarrow{i}E\xrightarrow{p}B$, and suppose $f\mathrel{\square}p$. Apply this lifting property to the square with top $i$ and bottom $1_B$. The resulting $s:B\to E$ satisfies $sf=i$ and $ps=1_B$. In the arrow category, the morphisms $(1_A,s):f\to i$ and $(1_A,p):i\to f$ have composite $1_f$. Hence $f$ is a retract of $i$. In particular, a factorizing pair $(L,R)$ with $L\mathrel{\square}R$ becomes a weak factorization system as soon as both classes are retract closed. This is the content of the standard retract argument; compare \citet[Lemmas~11.1.4 and~11.2.3]{Riehl2014}. It is this system of two classes of arrows that the manuscript calls a \emph{couple de Quillen}. A modern \emph{Quillen adjunction} is a different notion, involving two model categories and an adjoint pair of functors.
\end{dcomment}


\noindent\textbf{Proposition 2.} \textit{Soit $\Phi$ une partie close à gauche de} $\Fl(\M)$\textit{. Alors $\Phi$ satisfait les conditions} a) b) c) \textit{ci-dessous :}

\begin{enumerate}[label={}]
\item[a)] \textit{$\Phi$ contient les isomorphismes, et est stable par composition, et par cochangement de base (}i.e. \textit{si on a}
$\vcenter{\hbox{\begin{tikzcd}
A \arrow[r, "i"] \arrow[d] & B \arrow[d] \\
A' \arrow[r, "i'"] & B'
\end{tikzcd}}}$
\textit{carré cocartésien dans $\M$, avec $i \in \Phi$, on a $i' \in \Phi$).}

\item[b)] \textit{$\Phi$ stable par facteurs directs (en tant que partie de l'ensemble des objets de} $\underline{\Fl}(\M)$\textit{).}

\item[c)] \textit{$\Phi$ ``stable par limites inductives ordinales'', dans le sens suivant : Soit $I$ un ensemble bien ordonné, $(A_i)_{i \in I}$ un système inductif dans $\M$, tel que les $u_{ji} : A_i \longrightarrow A_j$ soient dans $\Phi$, et $A_{i_0} \xleftarrow{\ \sim\ } \varinjlim\limits_{i<i_0} A_i$ si $i_0 \in I$ est un élément limite. Supposons que $A = \varinjlim\limits_{I} A_i$ existe. Alors $\forall\ i_0 \in I$, la flèche canonique $A_{i_0} \longrightarrow A$ est dans $\Phi$.}
\end{enumerate}

Il y a un énoncé dual pour une partie $\Psi$ close à droite; elle contient les isomorphismes, est stable par composition et par changement de base, par facteurs directs, et enfin stable par limites projectives ordinales.

\msp{page 5}

On aura une réciproque partielle dans les catégories accessibles. C'est pour l'énoncé de cette réciproque que j'ai mis à part la condition \textit{b}) de stabilité par facteurs directs, au lieu de la bloquer avec \textit{a}).

\noindent\textbf{Définition 1.}

\begin{enumerate}[label={}]
\item[\textit{a})] (Pour mémoire.) RLP, LLP, la relation $\Phi \longleftrightarrow \Psi$ entre parties de $\Fl(\M)$ (relation non symétrique!).

\item[\textit{b})] (Pour mémoire.) Partie de $\Fl(\M)$ close à gauche, close à droite.

\item[\textit{c})] Soit $\Phi \longleftrightarrow \Psi$. On dit que le couple $(\Phi, \Psi)$ a la \textit{propriété de factorisation de Quillen}, si \textit{toute} flèche $f$ de $\M$ se factorise en
\[
f = pi, \qquad\quad i \in \Phi,\ p \in \Psi.
\]
On dit que $(\Phi, \Psi)$ est un \textit{couple de Quillen} si c'est un couple associé par \texttt{[les]} propriétés LLP et RLP ($\Phi \longleftrightarrow \Psi$), et \texttt{[si]} elles \texttt{[les parties} $\Phi$ \texttt{et} $\Psi$\texttt{]} sont closes (\textit{i.e.} $\Phi = \Psi^*$, $\Psi = \Phi_*$), et si elles ont la propriété de factorisation de Quillen.
\end{enumerate}

J'énonce dès maintenant le résultat technique essentiel, assurant l'existence de beaucoup de couples de Quillen :

\dgrblock{dgrb668}{A weak factorization system in sets}
\begin{dcomment}{examples}{p5:J'énonce dès maintenant}{A weak factorization system in sets}
The elementary example of injections and surjections makes the preceding abstraction concrete. We work in sets with the axiom of choice. Every injection $i:A\hookrightarrow B$ lifts against every surjection $p:X\twoheadrightarrow Y$: retain the given map on $A$, and choose a preimage in $X$ of the prescribed image of each point of $B\setminus A$. Conversely, lifting against $\varnothing\to\{*\}$ characterizes surjections. If a map $i$ has the left lifting property against all surjections, factor it as
\[
 A\xrightarrow{\iota} A\amalg B
     \xrightarrow{[i,1_B]} B.
\]
The second arrow is surjective, so the retract argument makes $i$ a retract of the injection $\iota$, hence an injection. Thus $(\mathrm{Mono},\mathrm{Surj})$ is a weak factorization system, generated by the single arrow $\varnothing\to\{*\}$.

The factorization displayed above is functorial: a square $(u,v):i\to i'$ acts on its middle object by $u\amalg v$. The diagonal in an arbitrary lifting problem need not be unique or functorially chosen. This is why functorial factorization and a specified coherent family of fillers are additional kinds of structure, even though both are useful in later developments. Nor should this factorization be confused with the familiar image factorization, whose left class is surjections and whose right class is injections. The same two classes occur in opposite order, with a different lifting relation. The asymmetry noted in Definition~1(a) is already visible in this example.
\end{dcomment}


% ---- from Der-ChXIII-25-01-04_p006-010.tex
\setcounter{footnote}{4}

\printedpage{4}
% head: A. Grothendieck : Dérivateurs, Ch. XIII, Édition des Universités de Montpellier II et Paris VII
% foot: Version 25 janvier 2004
\msp{page 6}\textbf{Théorème 1.} \textit{Supposons que $\cM$ soit} accessible, \textit{et stable par (petites) limites inductives. (Il doit suffire stabilité par sommes amalgamées, et par $\varinjlim$ filtrantes, voire ordinales \dots).} \textit{Soit $\Phi_0$ une} petite \textit{partie de} $\Fl(\cM)$\textit{, $\Psi=(\Phi_0)_*$, et $\Phi$ une partie de} $\Fl(\cM)$ \textit{telle que}
\[
\Phi_0\ \subset\ \Phi\ \subset\Psi^*\ (=\widetilde{\Phi}_0).
\]
\textit{Alors :}

\begin{itemize}[leftmargin=2.5em]
\item[a)] \textit{Pour que $\Phi$ soit close à gauche,} i.e. $\Phi=\widetilde{\Phi}_0$\textit{, il faut et il suffit qu'elle satisfasse les conditions} a) b) c) \textit{de la proposition 2. Quand il en est ainsi, le couple $(\Phi,\Psi)$ est un couple de Quillen.}

\item[b)] \textit{Supposons seulement que $\Phi$ satisfasse aux conditions} a), c) \textit{de la proposition 2 (à l'exclusion de la condition} b) \textit{des facteurs directs) – donc on ne peut affirmer que $\Phi$ soit close à gauche. Néanmoins, $(\Phi,\underbrace{\Psi}_{=\,\Phi_*})$ satisfait à la propriété de factorisation de Quillen.}
\end{itemize}

\dgrblock{dgrb669}{The small object theorem and its precise input}
\begin{dcomment}{exposition}{p6:on ne peut affirmer que}{The small object theorem and its precise input}
Theorem~1 is sound under its stated cocompleteness and accessibility hypotheses, understood in the usual locally small sense: an accessible cocomplete category is locally presentable. Its proof explains why the small generating family and the separate retract condition have been singled out. Let $I=\Phi_0$, choose a regular cardinal $\kappa$ such that the domains of all arrows in $I$ are $\kappa$-presentable, and fix $f:X\to Y$. Start with $X_0=X$. At stage $\alpha$, take every square
\[
\begin{tikzcd}
A\arrow[r]\arrow[d,"j"'] &X_\alpha\arrow[d]\\
B\arrow[r]&Y
\end{tikzcd}
\qquad(j\in I)
\]
and attach one copy of $B$ along its specified map from $A$. This is a set of squares because $I$ is small and the category is locally small. Their coproduct and pushout give $X_\alpha\to X_{\alpha+1}$ over $Y$. At limit stages take colimits. After $\kappa$ stages put $E=\mathop{\rm colim}_{\alpha<\kappa}X_\alpha$.

Every map $A\to E$ with $A$ a chosen domain factors through some $X_\alpha$, since the ordinal $\kappa$ is $\kappa$-filtered. A square $j\to(E\to Y)$ is therefore one of the squares attached at the next stage; its attaching cell supplies a filler into $E$. Hence $E\to Y$ belongs to $I^{\square}=\Psi$. The arrow $X\to E$ is an $I$-cell complex. The construction is functorial if every square is attached, without making choices of representatives depending on $f$.

For the theorem's slightly unusual formulation, coproducts of attaching maps cause no problem. A simultaneous attachment indexed by a set can be replaced, after well ordering that set, by successive pushouts. Thus closure under pushouts and continuous transfinite composition already puts every $I$-cell complex in $\Phi$. This proves part~(b). Moreover
\[
 I\subseteq\Phi\subseteq{}^{\square}(I^{\square})
 \quad\Longrightarrow\quad \Phi^{\square}=I^{\square}.
\]
If $\Phi$ is also retract closed, any $g\in{}^{\square}(I^{\square})$ is a retract of the left factor of its cellular factorization. Consequently $g\in\Phi$ and $\Phi={} ^{\square}(I^{\square})$, proving part~(a). This proof supplies the theorem only stated here; a modern account is \citet[Theorem~12.2.2 and Corollary~12.2.4]{Riehl2014}.

The parenthetical weakening of cocompleteness can also be made precise. For this construction it suffices to have the needed pushouts, ordinal colimits, local smallness and one regular cardinal giving the required smallness of the generating domains relative to the cellular chains. Pushouts together with ordinal colimits construct each set of attachments successively from the existing object $X_\alpha$, so a separately postulated empty coproduct is unnecessary. What cannot be dropped is the smallness condition guaranteeing that a lifting problem has appeared before the construction ends.
\end{dcomment}


\textbf{Remarque.} Dans la démonstration, l'existence de $\Phi_0\subset\Phi$, qui joue le rôle d'un petit ensemble de “générateurs” (à gauche) pour $\Phi$, joue un rôle essentiel. Il y aurait lieu de développer la notion de \textit{parties accessibles}
\msp{page 7}
de $\Fl(\cM)$, qui impliquerait l'existence d'un petit ensemble $\Phi_0\subset\Phi$ tel que $\Phi\subset\widetilde{\Phi}_0$, et qui aurait la propriété que
\[
\begin{aligned}
&\Phi\ \text{accessible}\quad\Longrightarrow\quad\Phi_*\ \text{accessible}\\
&\Psi\ \text{accessible}\quad\Longrightarrow\quad\Psi^*\ \text{accessible.}
\end{aligned}
\]
Tout couple $(\Phi,\Psi)$ de parties closes associées ($\Phi=\Psi^*$, $\Psi=\Phi_*$) et accessible, serait donc un couple de Quillen.

\dgrblock{dgrb670}{Accessibility does not supply generating arrows}
\begin{dcomment}{examples}{p7:serait donc un couple de Quillen.}{Accessibility does not supply generating arrows}
Grothendieck's proposed notion of an accessible class is meant to make the small object theorem reusable without displaying its generators. Modern terminology separates three properties: accessibility of a full subcategory of the arrow category, accessibility of a chosen factorization functor, and generation of its right lifting class by a set of arrows. The following known example shows why these cannot be substituted for one another; see \citet[Introduction and Examples~3.7]{Rosicky2021Remarks}.

Let $\mathcal M=\mathrm{Pos}$, and let $L$ consist of the order embeddings. This is a left lifting closed class. Indeed, a monotone map from a subposet $A\subseteq B$ to the two-element chain $2=\{0<1\}$ extends to $B$ by taking the upward closure of the inverse image of $1$. Thus every order embedding lifts against $2\to1$. Conversely, if $u:A\to B$ is not an order embedding, there exist $a,a'\in A$ with $u(a)\leq u(a')$ but $a\not\leq a'$. This also detects failure of injectivity by choosing the appropriate ordered pair of distinct points with the same image. The characteristic function of the principal upper set $\{x\mid a\leq x\}$ takes $a$ to $1$ and $a'$ to $0$, and cannot extend over $u$. Hence
\[
 L={} ^{\square}\{2\to1\}.
\]
One can even construct the associated factorization. Embed $A$ into the complete lattice $2^{\operatorname{Up}(A)}$ by recording membership in all its upper sets, and factor $f:A\to B$ through its graph in $B\times2^{\operatorname{Up}(A)}$. Projection to $B$ lifts against order embeddings, because a monotone map into a complete lattice extends by taking joins of the images of elements below the new point. Thus $(L,L^{\square})$ is a weak factorization system.

As a full subcategory of $\mathrm{Pos}^{[1]}$, $L$ is accessible and accessibly embedded. Filtered colimits preserve order reflection: any relation between images is witnessed at a finite stage, where it already reflects. Every embedding is a filtered union of its finite subembeddings; such finite arrows are finitely presentable in the arrow category and remain so in this full subcategory.

Nevertheless no set $I\subseteq L$ satisfies $I^{\square}=L^{\square}$. Choose an infinite regular $\kappa$ larger than the cardinalities of all domains and codomains in $I$, and let $P$ be the poset of subsets of $\kappa$ of cardinality less than $\kappa$. A map $h:A\to P$ extends along every $i:A\hookrightarrow B$ in $I$ by
\[
 \bar h(b)=\bigcup_{i(a)\leq b}h(a).
\]
Regularity makes this union an element of $P$, and order reflection gives $\bar h(i(a))=h(a)$. Therefore $P\to1$ belongs to $I^{\square}$. Now embed the discrete poset on $\kappa$ into the poset obtained by adjoining a top element. The map $a\mapsto\{a\}$ cannot extend to that top, since its image would have to contain all of $\kappa$. Thus $P\to1\notin L^{\square}$. The proposed generating-set conclusion fails even for an accessible, left lifting closed class having factorizations. This does not refute the possibility of defining the stronger notion Grothendieck wants; it identifies what such a definition would have to add.
\end{dcomment}


\section*{2\quad Catégories de Quillen faibles, et construction d'icelles (Théorème-scholie conjectural)}

Dans toute la suite, on se donne une \textit{catégorie de modèles}
\[
\left\{
\begin{array}{ll}
\mbox{$(\cM,W)$, où $W\subset\Fl(\cM)$} & \\[3ex]
\mbox{\textit{$W$ faiblement saturée} (\footnotemark) :} &
\begin{array}[c]{@{}l@{\quad}p{6.5cm}@{}}
a) & $W$ contient les isomorphismes.\\[1ex]
b) & Si deux parmi $f$, $g$, $gf$ \texttt{[sont]} dans $W$, le troisième aussi.
\end{array}
\end{array}
\right.
\]
\footnotetext{C'est l'axiome M 5 de \textsc{Quillen}.}

\printedpage{5}
% head: A. Grothendieck : Dérivateurs, Ch. XIII, Édition des Universités de Montpellier II et Paris VII
% foot: Version 25 janvier 2004
On veut décrire des données supplémentaires un peu plus faibles que celles d'une catégorie de modèles de Quillen. On suppose données des parties
\[
C,\ F,\ TC,\ TF\ \subset\ \Fl(\cM),
\]
avec les propriétés suivantes. Comme axiome préliminaire, je pose

\begin{itemize}[leftmargin=3em,labelsep=0.6em]
\item[\textbf{M 0}] Stabilité de $\cM$ par $\varprojlim$ finies et $\varprojlim$ \texttt{[plutôt $\varinjlim$]} finies.

\item[\textbf{M 1}]\msp{page 8}Propriété de relèvement-prolongement
\[
C\longleftrightarrow TF,\qquad TC\longleftrightarrow F.
\]

\item[\textbf{M 2}] Propriété de factorisation pour $(C,TF)$, $(TC,F)$.

\item[\textbf{M 3}] Stabilité de $F$ par composition et par changement de base, de $C$ par composition et cochangement de base, $F$ et $C$ contiennent les isomorphismes.

\item[\textbf{M 4}] Stabilité de $TF$ par composition et changement de base, de $TC$ par composition et cochangement de base, et $TF$ et $TC$ contiennent les isomorphismes.

\item[\textbf{M}$'$] $TC\subset C\cap W$, $TF\subset F\cap W$.

\item[\textbf{M}$''$] $C\subset\ \underbrace{\Cof_W}_{\text{$W$-cofibrations}}\ ,\ F\subset\ \underbrace{\Fib_W}_{\text{$W$-fibrations}}\ .$
\end{itemize}

\texttt{[Un morphisme $p:X\longrightarrow Y$ de $\cM$ est une $W$-fibration si pour tout diagramme de carrés cartésiens de la forme :}
\[
\begin{tikzcd}
X \ar[d,"p"'] & X' \ar[l] \ar[d] & X'' \ar[l,"t"'] \ar[d] \\
Y & Y' \ar[l] & Y'' \ar[l,"s"] \ ,
\end{tikzcd}
\]
\texttt{si $s\in W$, alors $t\in W$. La notion de $W$-cofibration est la notion duale.]}

Dans le cas de \textsc{Quillen} (Q-\textit{catégories de modèles}), M$'$ est renforcé en
\[
TC=C\cap W,\qquad TF=F\cap W,
\]
et il n'y a pas d'axiome M$''$. L'axiome M$''$ équivaut à la propriété de “propriety”, correspondant aux Q-catégories de modèles \textit{propres.}

\dgrblock{dgrb671}{Relative equivalences and properness}
\begin{dcomment}{exposition}{p8:correspondant aux Q-catégories de modèles}{Relative equivalences and properness}
The additional four classes serve two purposes: they furnish replacements and make ordinary base change behave homotopically. In modern language $(\mathcal M,W)$ is first a relative category with two-out-of-three. It is not yet a Quillen model category. Axiom M~2 supplies both factorizations, and M~1 supplies the corresponding lifting relations; without retract closure these pairs need not yet be weak factorization systems. The strict inclusions allowed in M$'$ deliberately leave open whether the chosen trivial cofibrations exhaust $C\cap W$, and similarly for fibrations.

The intrinsic class $\mathrm{Fib}_W$ has a stronger universal quantifier than the assertion that one pullback preserves weak equivalences. A map $p:X\to Y$ belongs to it precisely when, after every base change $Y'\to Y$, the functor
\[
 (p')^*:\mathcal M/Y'\longrightarrow\mathcal M/(X\times_Y Y')
\]
takes weak equivalences to weak equivalences. Its dual is $\mathrm{Cof}_W$. For an ordinary model structure, pullbacks of fibrations are again fibrations. Therefore right properness says exactly that $F\subseteq\mathrm{Fib}_W$; left properness similarly gives $C\subseteq\mathrm{Cof}_W$. The identification in the manuscript is correct, with the customary completeness and retract conventions supplied when translating to a modern closed model category.

There is a useful unconditional inclusion
\[
 W_{\mathrm{univ}}\subseteq\mathrm{Fib}_W\cap W,
 \qquad
 W^{\mathrm{univ}}\subseteq\mathrm{Cof}_W\cap W.
\]
For the first, put a universally weak equivalence vertically in a cartesian square. Both vertical sides are in $W$. If the bottom horizontal arrow is in $W$, two-out-of-three puts the top arrow in $W$; apply the same argument after any further base change. The reverse inclusions require the factorization argument of Proposition~4. This logical dependence explains why Proposition~3 first invokes the displayed equality and then proves the result which justifies it. There is no circularity: Proposition~4 uses factorization and the intrinsic fibrations only, not Proposition~3's saturation construction.
\end{dcomment}


\textbf{Proposition 3.} \textit{Supposons les conditions} M 0 \textit{(pour mémoire),} M 1, M 2, M$'$, M$''$ \textit{satisfaites. Alors en remplaçant $C$, $TC$ par les saturés pour les propriétés de stabilité de contenir les isomorphismes, d'être stable par composition et par cochangement de base, et $F$, $TF$ par celles de contenir les isomorphismes et d'être stable par composition et par changement de base, on trouve $C'$, $TC'$, $F'$, $TF'$ satisfaisant}
\printedpage{6}
% head: A. Grothendieck : Dérivateurs, Ch. XIII, Édition des Universités de Montpellier II et Paris VII
% foot: Version 25 janvier 2004
\msp{page 9}\textit{à tous les axiomes} M 1 à M 4, M$'$, $M''$. % [sic] $M''$ (italic M) -> M$''$ (roman M), the axiom is printed with roman M everywhere else

\textsc{Démonstration}. La proposition 2 nous assure que M 1 continue à être satisfait, et pour M 2 c'est trivial. M 3 et M 4 seront valables par construction. D'autre part, M$''$ reste valable, car $\Cof_W$ a les stabilités requises pour $C'$, et $\Fib_W$ celles requises pour $F'$. Reste à prouver M$'$, \textit{i.e.}
\[
TC'\ \subset\ C'\cap W,\qquad TF'\subset F'\cap W,
\]
et pour ceci il suffit de voir que $C'\cap W$ resp. $F'\cap W$ a les stabilités requises pour $TC'$ resp. $TF'$. Cela résulte du

\textbf{Lemme.} \textit{Soit $C$ une partie de} $\Cof_W$\textit{. Si $C$ contient les isomorphismes, il en est de même de $C\cap W$. Si de plus $C$ est stable par composition resp. par cochangement de base, il en est de même} \texttt{[de]} \textit{$C\cap W$. Énoncé dual pour une partie $F$ de} $\Fib_W$\textit{.}

Cela provient du fait que
\[
\boxed{\ \Cof_W\cap W=W^{\mathrm{univ}}\ (\mbox{\footnotemark})\ }
\stackrel{\text{déf}}{=}
\begin{array}{l}
\text{ensemble des flèches $A\longrightarrow B$ telles que pour tout}\\
A\longrightarrow A',\ A'\longrightarrow B'=A'\vee_A B\ \ \text{est dans $W$.}
\end{array}
\]
\footnotetext{flèches \textit{coüniversellement} dans $W$, ou $W$-\textit{équivalences coüniverselles}.}
est lui-même stable par composition et par cochangement de base, et que l'on a bien sûr, si $C\subset\Cof_W$,
\[
C\cap W\ =\ C\cap(\Cof_W\cap W)\ =\ C\cap W^{\mathrm{univ}}.
\]

\msp{page 10}Donc il reste à prouver le

\textbf{Corollaire 1.} \textit{Supposons qu'il existe des parties $C$, $F$, $TC$, $TF$ $\subset$} $\Fl(\cM)$ \textit{satisfaisant} M 2, M$'$ (\footnote{On n'a pas besoin de tout M$'$, seulement de $TC\subset W$, $TF\subset W$.}), M$''$. \textit{Alors}
\[
\begin{aligned}
\Cof_W\cap W\ &=\ W^{\mathrm{univ}}\\
\Fib_W\cap W\ &=\ W_{\mathrm{univ}}\ \ (\mbox{\footnotemark}).
\end{aligned}
\]
\footnotetext{Flèches $f:X\longrightarrow Y$ \textit{universellement} dans $W$, \textit{i.e.} telles que pour tout $Y'\longrightarrow Y$, $f':X'=X\times_Y Y'\longrightarrow Y'$ soit dans $W$.}

Pour voir ceci, montrons p. ex. le deuxième (le premier s'en déduira par dualité).
\[
\begin{tikzcd}
X \ar[d,"f"'] & X' \ar[l] \ar[d,"f'"] \\
Y & Y' \ar[l]
\end{tikzcd}
\qquad\qquad\qquad
f\in\Fib_W,\ f\in W\Longrightarrow f'\in W.
\]

\printedpage{7}
% head: A. Grothendieck : Dérivateurs, Ch. XIII, Édition des Universités de Montpellier II et Paris VII
% foot: Version 25 janvier 2004
On a mieux :

\textbf{Proposition 4.} \textit{Supposons que toute flèche de $\cM$ se factorise en $pi$, avec $i\in W$ et $p\in$} $\Fib_W$\textit{. Alors, si $X$, $Y$ sont dans} $\underline{\Fib}_W S\subset\cM/S$ \texttt{[sous-catégorie pleine de $\cM/S$ formée des $r:Z\longrightarrow S$, où $r\in\Fib_W$]}\textit{, alors toute $f:X\longrightarrow Y$ dans $\cM/S$ qui est dans} $W_S$ \texttt{[formé des flèches au-dessus de $S$ qui sont dans $W$]} \textit{est dans $W^{\mathrm{u}}_S$,} i.e. \textit{elle reste dans $W$ après tout changement de base $S'\longrightarrow S$.}
\[
\begin{tikzcd}[column sep=small]
X \ar[rr,"f"] \ar[dr,"p"'] & & Y \ar[dl,"q"] \\
& S &
\end{tikzcd}
\qquad\qquad\qquad
p,q\in\Fib_W
\]

\textit{Énoncé dual pour les flèches de} $S\backslash\underline{\Cof}_W\subset S\backslash\cM$ \texttt{[sous-catégorie pleine de $S\backslash\cM$ formée des $i:S\longrightarrow Z$, où $i\in\Cof_W$]}\textit{, une telle flèche qui est dans $W$ est $S$-coüniversellement dans $W$, si toute flèche de $\cM$ se factorise en $qj$, avec $j\in$} $\Cof_W$ \textit{et $q\in W$.}

\msp{page 11}\textsc{Démonstration}.
\[
\begin{tikzcd}[column sep=small]
X \ar[dr,"f\in W"] \ar[dd] & & \overline{X'} \ar[ll] \ar[dr,"\overline{f'}\in W"] \ar[dd] & & X' \ar[ll,"i_X"'] \ar[dr,"f'"] \ar[dd] & \\
& Y \ar[dl] & & \overline{Y'} \ar[ll,crossing over] \ar[dl] & & Y' \ar[ll,"i_Y"',crossing over] \ar[dl] \\
S & & \overline{S'} \ar[ll,"F"] & & S' \ar[ll,"i"',"TC"] &
\end{tikzcd}
\]
On factorise $S'\longrightarrow S$ en $S'\overset{i}{\longrightarrow}\overline{S'}\overset{p}{\longrightarrow}S$, avec $i\in W$, $p\in\Fib_W$, par M 2. Donc $f\in W\Longrightarrow\overline{f'}\in W$. Comme $X,Y\in\Ob\,\underline{\Fib}_W S$, on a $\overline{X'},\overline{Y'}\in\Ob\,\underline{\Fib}_W\overline{S'}$, et comme $i:S'\longrightarrow\overline{S'}$ est dans $W$, donc aussi $i_X$, $i_Y$, on a $f'\in W$ par M 5 (saturation faible de $W$).

\dgrblock{dgrb672}{Why weak equivalences between fibrations survive base change}
\begin{dcomment}{exposition}{p11:par M 5 (saturation faible}{Why weak equivalences between fibrations survive base change}
Proposition~4 turns the factorization machinery into a statement about $W$ alone. Its proof is short, but each use of the intrinsic fibration condition has a different role. Given a weak equivalence $f:X\to Y$ over $S$, with $X\to S$ and $Y\to S$ in $\mathrm{Fib}_W$, factor a proposed base change as
\[
 S'\xrightarrow{i}\overline S\xrightarrow{q}S,
 \qquad i\in W,\quad q\in\mathrm{Fib}_W.
\]
Pulling $f$ back along $q$ gives $\bar f\in W$: it is a pullback of the weak equivalence $f$ along the base change $Y\times_S\overline S\to Y$ of $q$. Next, the maps $X\times_S\overline S\to\overline S$ and $Y\times_S\overline S\to\overline S$ remain intrinsic fibrations. Pulling $i$ back along them gives weak equivalences
\[
 X\times_SS'\longrightarrow X\times_S\overline S,
 \qquad
 Y\times_SS'\longrightarrow Y\times_S\overline S.
\]
The commutative square with these sides and $\bar f$ then places $f\times_SS'$ in $W$ by two-out-of-three. No lifting property, functorial choice of factorization, or retract argument is used.

For Corollary~1 apply this assertion in the slice over the codomain of $f\in\mathrm{Fib}_W\cap W$, comparing $f$ with the identity fibration. The converse inclusion was proved above. Duality gives the cofibration equality. Now Proposition~3 can be completed: adjoining the specified composition and base-change operations preserves lifting by Proposition~2 and preserves available factorizations because the four classes only grow. The intrinsic classes have the same finite closure properties. Finally $C'\cap W=C'\cap W^{\mathrm{univ}}$ is stable under pushout and composition, and therefore contains the closure of $TC$; the fibration side is dual. The strengthening by retracts on manuscript p.~11 is valid if $W$ is retract closed, since pullback and pushout transport retract diagrams to retract diagrams.
\end{dcomment}


\dgrblock{dgrb673}{The factorization hypothesis seen on three objects}
\begin{dcomment}{development}{p11:par M 5 (saturation faible}{The factorization hypothesis seen on three objects}
Grothendieck's intrinsic classes make it possible to isolate the role of factorization without any topology. The following small calculation shows exactly what would go wrong if Corollary~1 were read as a formal property of two-out-of-three alone. We checked the comparisons with relative-category and factorization-system treatments in \citet{Riehl2014,Rosicky2021Remarks}; we did not find this particular three-object calculation there. Its purpose is to clarify the hypothesis, not to claim a new general phenomenon.

\textbf{A weak equivalence which is an intrinsic fibration but is not universal.}\ours
Let $\mathcal M$ be the chain $0<1<2$, regarded as a category, and let $W$ consist of the identities and the arrow $w:0\to2$. Then $W$ has two-out-of-three, $w\in\mathrm{Fib}_W$, but $w\notin W_{\mathrm{univ}}$.

\emph{Proof.} The only composite of two nonidentity arrows is $0\to1\to2$. Only its composite is in $W$, so two-out-of-three imposes no further weak equivalence; composites involving identities satisfy the condition. Pullback in this chain is minimum. Every base change of $w$ is $0\to y$ for $y\in\{0,1,2\}$. Below $y=0$ or $1$, the only weak equivalences are identities. Below $2$, the sole additional test is $0\to2$, whose pullback along $0\to2$ is $1_0$. Thus all base changes of $w$ preserve weak equivalences and $w\in\mathrm{Fib}_W$. Its pullback along $1\to2$ is $0\to1$, which is not in $W$. This proves the assertion.

The missing factorization can also be located. The map $1\to2$ is not in $\mathrm{Fib}_W$, since it pulls $w$ back to $0\to1$. A factorization of $1\to2$ with first factor in $W$ must have first factor $1_1$, leaving precisely this nonfibration as second factor. Thus the hypothesis of Proposition~4 fails for a specified arrow. This finite category has all small limits and colimits; filtered diagrams of arrows have colimits attained in their finite directed image. Hence $W$ is even stable under filtered colimits. Those size and closure properties do not replace factorization in the corollary.\hfill$\square$
\end{dcomment}


\textbf{Commentaire sur Proposition 3.} On peut être encore plus exigeant, en demandant dans M 3, M 4 également la stabilité par facteurs directs, et (quand seuls M 1, M2, M$'$, M$''$ % [sic] M2 -> M 2, printed without the space used everywhere else
sont satisfaits) en saturant $C$, $TC$, $F$, $TF$ également pour cette stabilité supplémentaire. La proposition 2 assure encore que M 1 continue à être vérifié (et M 2, M 3, M 4 trivialement), on voit qu'il en est ainsi encore de M$'$, M$''$ à condition qu'on renforce l'axiome de saturation faible M 5 sur $W$ par
\[
\textit{stabilité de $W$ par facteurs directs,}
\]
ce qui implique en effet que $\Cof_W$, $\Fib_W$ sont stables par facteurs directs, ainsi que $W^{\mathrm{univ}}$ et $W_{\mathrm{univ}}$.

\msp{page 12}Il ne semble pas que les axiomes M 0 à M 4, plus M$'$ et M$''$ suffisent à développer les suites exactes de suspension et de cosuspension. Il faut un petit chouia de plus :

\printedpage{8}
% head: A. Grothendieck : Dérivateurs, Ch. XIII, Édition des Universités de Montpellier II et Paris VII
% foot: Version 25 janvier 2004
\textbf{Définition 2.} J'appelle \textit{pré-Q-catégorie de modèles à gauche,} la donnée de $\cM$ et de $W$, $C$, $F$, $TC$, $TF$ $\subset$ $\Fl(\cM)$, satisfaisant les axiomes M 0 à M 5, M$'$, M$''$, plus la forme renforcée de la première inclusion dans M$'$
\[
\text{M}'_0\ \text{(g)}\qquad\qquad\qquad TC=C\cap W\qquad\quad(\text{non seulement } TC\subset C\cap W).
\]
De même, on définit la \textit{pré-Q-catégorie de modèles à droite,} par la condition
\[
\text{M}'_0\ \text{(d)}\qquad\qquad\qquad\qquad\qquad\qquad TF=F\cap W.
\]

NB Quand les deux conditions sont satisfaites à la fois, alors la donnée de $TC$, $TF$ résulte déjà de celle de $W$, $C$, $F$, et on a simplement une Q-catégorie de modèles “propre” (au sens de \textsc{Quillen}).

\dgrblock{dgrb674}{The purpose of the one-sided equalities}
\begin{dcomment}{exposition}{p12:Quand les deux conditions sont satisfaites à la fois}{The purpose of the one-sided equalities}
The one-sided conditions are intended to recover exactness of homotopy fibres or cofibres while leaving the auxiliary choices flexible. They should not be read as a claim that every relative category has such a structure, or that a left structure and a right structure must use the same four classes. Grothendieck explicitly allows different choices. The assertion about a proper Quillen category when both equalities hold uses Quillen's original conventions; for a modern closed model category one also requires the retract conditions, or replaces the classes by their retract closures when the stated hypotheses allow it.

The historical motivation is coherent with the earlier reduction to categories of fibrant objects. Brown's formulation develops homotopy using weak equivalences, fibrations and path objects, without specifying a full model structure; see \citet[\S1, pp.~420--421]{Brown1973}. Grothendieck is pursuing a further independence statement: the resulting homotopy operations should be attached to $W$, with $(C,TC,F,TF)$ used to prove that they exist and have exactness properties. Later in these pages, the same issue reappears for all diagram categories, where objectwise replacement alone does not automatically construct a coherent derived Kan extension.

The claims on manuscript p.~13 about exactness are explicitly provisional: the source says that they still require checking. They are therefore recorded as conjectural assertions in the audit. The completion proved in the commentary on manuscript p.~32 settles their conclusion under the printed hypotheses: either one-sided equality implies exact factorization of weak equivalences, after which the two designated trivial classes can be enlarged to $C\cap W$ and $F\cap W$. The resulting structure is biexact with the same $W,C,F$, though it need not be closed. Thus one-sided structures retain flexibility in the designated trivial maps, but their existence already yields a possibly nonclosed Quillen structure. This is a later verification of the assertion, not a proof supplied on p.~13. The cross-reference to ``proposition~3'' on p.~13 points in fact to Proposition~4: the property cited there is the universal base-change property just proved, together with Corollary~1.
\end{dcomment}


\textbf{Scholie.} Pour moi, l'introduction de $C$, $TC$, $F$, $TF$ est purement auxiliaire – ce qui m'importe, ce sont les propriétés de $W$. Ainsi, l'existence de $C$, $TC$, $F$, $TF$ satisfaisant aux axiomes faibles M 0 à M 4, M$'$, M$''$, du fait qu'elle donne les factorisations
\msp{page 13}
dont il est question dans \texttt{[la]} proposition 3, entraîne la propriété remarquable de $W$ énoncée dans cette proposition, et dans le corollaire qui la précède. De même, ces propriétés de factorisation (\footnote{sous une forme légèrement renforcée il est vrai – il faut une factorisation \textit{fonctorielle} pour $f$ \dots}) assurent l'existence d'une bonne théorie des fibres et cofibres $W$-homotopiques sans besoin de référence dans cette théorie à $C$, $TC$, $F$, $TF$ (dont l'introduction donne tout au plus des petits compléments \dots). Mais je ne vois pas comment, à partir des seules hypothèses de factorisation, prouver la propriété cruciale d'exactitude pour les carrés $W$-cartésiens, et $W$-cocartésiens, nécessaire pour pouvoir développer la suite \textit{exacte} de cosuspension et de suspension. C'est pour avoir cette exactitude que l'introduction auxiliaire, et transitoire, de l'attirail $(C,TC,F,TF)$ est utile, l'axiome supplémentaire à gauche $TC=C\cap W$ assurant (sauf erreur) l'exactitude des carrés $W$-cartésiens, et l'axiome à droite $TF=F\cap W$ celle des carrés $W$-cocartésiens. (Chose à vérifier avec plus de soin que je ne l'ai fait à présent.) Une fois ceci obtenu, et éventuellement itou pour les $(\cM(I),W(I))$ pour toute catégorie-diagramme $I$, on peut
\msp{page 14}
totalement oublier ces ensembles auxiliaires – j'ai envie en tous cas de les oublier.

\dgrblock{dgrb675}{The auxiliary status of the Quillen classes}
\begin{dcomment}{context}{p14:totalement oublier ces ensembles auxiliaires}{The auxiliary status of the Quillen classes}
Grothendieck wants the eventual constructions to belong to the pair
$(\mathcal M,W)$, even when their existence is established by choosing
cofibrations and fibrations. His wish to forget those choices is therefore a
mathematical requirement: two suitable presentations should give the same
homotopy theory of diagrams, including its Kan extensions and their comparison
maps. It does not assert that the choices are dispensable in a proof.

The letter to Thomason of 2 April 1991 places this passage in its immediate
setting. Grothendieck begins with the accessibility theory he had been
developing from SGA~4, then distinguishes the intrinsic localizer from the
additional classes used to calculate with it. Concerning closed Quillen
triples he writes, ``J'avais cru pouvoir m'en passer, mais finalement n'y suis
pas parvenu'' \citep[pp.~3--4]{Grothendieck1991Letter}. The contrast he draws
there with Brown's one-sided methods explains why the present chapter keeps
trying weaker combinations of axioms. The asymmetry is deliberate: the
existence of homotopy fibres and that of homotopy cofibres need not initially
be established by the same auxiliary structure.

There are two distinct tasks in such a programme. One must first construct
the localized categories of diagrams, with sufficient control over their
morphisms to form derived operations. One must then prove that the operations
satisfy comparison laws as the indexing category varies. An equivalence
between two ordinary homotopy categories settles neither task by itself.
This is why the apparently elementary question whether a lifting problem
can be solved naturally in a diagram will become decisive on manuscript
page~41. A family of componentwise solutions need not be a natural
transformation.

The later literature realizes these two tasks in several settings, with
hypotheses that must be retained. Cisinski's construction of homotopy right
Kan extensions starts from a closed Quillen model category with small
limits; it does not require an injective model structure on every diagram
category \citep[Theorem~6.11]{Cisinski2003}. His subsequent theory of derivable
categories provides a further framework for one-sided and exact-category
examples \citep[Introduction and \S1]{Cisinski2010}. These results give a
precise continuation of the programme, but they should not be substituted
for a proof of the particular conjectural recognition criterion now being
formulated. Here the outstanding issue is the smallness of the classes from
which the factorization systems are to be built.
\end{dcomment}


Mon espoir, c'est que pour une $W$ donnée, même si on n'arrive pas à trouver une structure de triple de Quillen $(W,C,F)$ associée (p. ex. parce qu'il n'y en aurait pas), on peut trouver à la fois des quintuplets $(W,C,F,TC,TF)$ qui soient des pré-Quillen à gauche, et des quintuplets (d'\textit{autres} quintuplets!) qui soient des pré-Quillen à droite.

Voici le genre d'énoncé en forme, pour le moment conjectural voire totalement heuristique, auquel je songe :

\textbf{Théorème-scholie conjectural :} \textit{On suppose que $\cM$ est stable par $\varinjlim$ “quelconques”, et par $\varprojlim$ finies, et que $W$ est stable par $\varinjlim$ filtrantes (en plus de la saturation faible).}

% ---- from Der-ChXIII-25-01-04_p011-015.tex
\printedpage{9}
% head: A. Grothendieck : Dérivateurs, Ch. XIII, Édition des Universités de Montpellier II et Paris VII

\begin{itemize}[leftmargin=2.5em]
\item[\circled{1}] \textit{Alors} $\Cof_W$ \textit{et} $W^{\univ}$ \textit{ont les propriétés de stabilité} a) \textit{et} c) \textit{de la proposition} 2 (p.~4), \textit{et aussi la propriété de stabilité} b) \textit{par facteurs directs, pourvu que} $W$ \textit{lui-même l'a.}

\item[\circled{2}] \textit{Supposons} $\mathcal{M}$ accessible, \textit{de sorte qu'on dispose du théorème} 1. \textit{Conditions équivalentes sur} $W$ \textit{:}

\msp{page 15}

\begin{itemize}[leftmargin=2.5em]
\item[a)] \textit{Il existe une partie petite} $C_0$ \textit{de} $\Cof_W$\textit{, et} $TC_0$ \textit{de} $W^{\univ}$\textit{, telles que}
\[ TF \overset{\text{déf}}{=} C_{0*} \subset W, \qquad (TC_0)_* \overset{\text{déf}}{=} F \subset \Fib_W. \]

\item[a$'$)] \textit{Il existe un quadruplet} $(C,TC,F,TF)$ \textit{de parties de} $\Fl(\mathcal{M})$\textit{, satisfaisant les conditions} M~1 \textit{à} M~4, M$'$, M$''$ \textit{de la page} 6 (\footnote[10]{avec $C$, $TC$ clos à gauche, $F$, $TF$ clos à droite.}).

\item[a$''$)] \textit{Comme} a$'$)\textit{, en supposant} M$'$ \textit{sous forme renforcée avec} M$'_0$ (g), i.e. \textit{on a une structure de pré-Q-catégorie de modèles à gauche (}$TC = C \cap W$\textit{).}

\item[b)] \textit{On a}
\[ (\Cof_W)_* \ \subset\ W, \qquad (W^{\univ})_* \ \subset\ \Fib_W \]
\textit{(donc} $(\Cof_W)_* \subset \Fib_W \cap W$ \textit{puisque} $(\Cof_W)_* \subset (W^{\univ})_*$ (\footnote[11]{$W^{\univ} \subset \Cof_W$.})\textit{).}

\item[b$'$)] \textit{On a}
\[ (\Cof_W)_* \ \subset\ W_{\univ}, \qquad (W^{\univ})_* \ \subset\ \Fib_W. \]
\end{itemize}

\item[\circled{3}] \textit{Pour qu'il existe un quadruplet qui, avec} $W$\textit{, définisse une pré-Q-catégorie de modèles à droite (}$TF = F \cap W$\textit{)} (\footnote[12]{On suppose toujours $\mathcal{M}$ accessible dans 3° et 4°, avec en plus $F$, $TF$ (et $C$, $TC$) closes à droite (resp. à gauche).})\textit{, il faut et il suffit qu'on ait} b$'$) \textit{ci-dessus, ainsi que}
\[ (W^{\univ})_* \cap W_{\univ} \ \ \text{\itshape est clos à droite.} \]

\msp{page 16}

\item[\circled{4}] \textit{Les hypothèses préliminaires sur} $(\mathcal{M}, W)$\textit{, y inclus l'accessibilité, ainsi que les conditions équivalentes dans}\circled{2}\textit{, sont stables par passage de} $(\mathcal{M}, W)$ \textit{à} $(\mathcal{M}(I), W(I))$\textit{. Mais qu'en est-il de la condition dans}\circled{3} \textit{? Est-elle conséquence des conditions dans} \circled{2} \textit{(}i.e. \textit{leurs est-elle équivalente) ?} % [sic] leurs est-elle -> lui est-elle, grammar
\end{itemize}

Je vais esquisser la démonstration du théorème-scholie, pour cerner le travail technique qui reste à faire.

\dgrblock{dgrb676}{The closure part of the conjectural theorem}
\begin{dcomment}{exposition}{p16:Je vais esquisser la démonstration}{The closure part of the conjectural theorem}
The theorem-scholie combines a correct closure calculation with a much more substantial assertion that small generating data exist. It helps to finish the closure calculation before attempting the existence argument. Assume the colimit and weak-equivalence hypotheses stated on manuscript p.~14. Then $W^{\mathrm{univ}}$ and $\mathrm{Cof}_W$ contain isomorphisms and are stable under pushout and composition, by pasting pushout squares and using two-out-of-three where appropriate.

For transfinite composition, let $(A_\alpha)_{\alpha\leq\lambda}$ be continuous. Suppose first that every successor map is a universal weak equivalence. Push the chain out along an arbitrary $A_0\to T$, obtaining $T_\alpha=T\amalg_{A_0}A_\alpha$. Each successor is in $W$. Transfinite induction shows that $T\to T_\beta$ is in $W$: the successor step uses composition, while a limit step is the filtered colimit of the arrows $T\to T_\alpha$. Thus $A_0\to A_\lambda$ lies in $W^{\mathrm{univ}}$.

For $\mathrm{Cof}_W$, start with a weak equivalence $T\to T'$ under $A_0$. The two pushed-out chains satisfy $T_0\to T'_0\in W$. The successor step preserves this relation because the corresponding base change of $A_\alpha\to A_{\alpha+1}$ is a $W$-cofibration; at a limit step use filtered-colimit closure of $W$. It follows that pushing out along $A_0\to A_\lambda$ preserves weak equivalences. Repeating after arbitrary cobase change establishes the universal quantifier in the definition of $\mathrm{Cof}_W$.

Retract closure is available here without a further assumption on $W$. If $f$ is a retract of $g\in W$ in the arrow category, write the retract maps as $s:f\to g$ and $r:g\to f$, with $rs=1_f$. Then $f$ is the colimit of
\[
 g\xrightarrow{sr}g\xrightarrow{sr}g\xrightarrow{sr}\cdots.
\]
Indeed a compatible cocone $h:g\to z$ satisfies $h(sr)=h$ and factors uniquely as $hs:f\to z$ followed by $r$. Since this is a filtered diagram of arrows in $W$, its colimit $f$ belongs to $W$. Pullback and pushout preserve retract diagrams, so the intrinsic and universal classes inherit retract closure. This is the elementary mechanism behind \citet[Proposition~1.19]{Beke2000}; the stronger accessibility hypothesis is not needed when filtered-colimit closure is already assumed.

These arguments establish item~\circled{1}. They do not establish that either whole class is defined by a lifting property, nor do they produce a small generating family. The distinction becomes decisive in the next paragraph of the manuscript.
\end{dcomment}


1°) Est de la vérification de routine.

2°) Voici les implication $\pm$ tautologiques entre les conditions envisagées : % [sic] les implication -> les implications, agreement
\[ a'' \Longrightarrow a' \overset{?}{\Longrightarrow} a \Longrightarrow b \Longleftarrow b', \tag*{$(*)$} \]
% foot: Version 25 janvier 2004
\printedpage{10}
% head: A. Grothendieck : Dérivateurs, Ch. XIII, Édition des Universités de Montpellier II et Paris VII
mais je vois que $a' \Longrightarrow a$ n'est nullement tautologique, et peut-être même faux après tout~! (\footnote[13]{C'est pourquoi je l'ai modifié, en renforçant $a'$) par \texttt{[la condition de la note 10]}.}). Il faudrait supposer que dans \texttt{[le]} quadruplet $(C,F,TC,TF)$, $F$ et $TF$ sont clos à droite, et plus précisément, que $F = (TC)_*$, $TF = C_*$, ce qui sera le cas p.~ex. si on a une structure de Q-catégorie de modèles “close” (“closed model category”), au sens de \textsc{Quillen}. Les autres trois implications dans $(*)$ sont bel et bien tautologiques.

$a \Longrightarrow a'$. On peut dans $a$) supposer
\[ TC_0 \ \subset\ C_0 \]
(quitte à remplacer $C_0$ par $C_0 \cup TC_0$), puis, grâce à
\msp{page 17}
1°) ($W^{\univ}$ et $\Cof_W$ sont clos à gauche), posant
\[ \overline{TC_0} = TC, \qquad \overline{C_0} = C, \qquad \text{\texttt{[plutôt}}\ \ \widetilde{TC_0} = TC, \qquad \widetilde{C_0} = C, \text{\texttt{]}} \]
on aura $TC \subset W^{\univ}$, $C \subset \Cof_W$. Passant de $TC_0$ à $TC$, de $C_0$ à $C$, on a perdu seulement que ces parties de $\Fl(\mathcal{M})$ soient petites. Mais par le théorème 1, posant
\[ F \ =\ (TC)_*, \qquad TF \ =\ C_*, \]
les couples de parties closes (à gauche -- à droite)
\[ (C,\, TF), \qquad (TC,\, F) \]
satisfont à la condition de factorisation de Quillen. On aura alors les propriétés M~1 à M~4, plus les inclusions $TC \ \subset\ C \cap W$ (car $TC \ \subset\ W^{\univ} \ \subset\ W$, et $TC \ \subset\ C$ par construction), et $C \subset \Cof_W$ de M$'$ et M$''$. Il reste à vérifier $TF \subset F \cap W$, $F \subset \Fib_W$, mais $F$ $(= (TC)_*)$ $\ \subset \Fib_W$ par hypothèse, et $TF$ $(= C_*)$ $\ \subset W$ aussi~; d'autre part $TF \subset F$ résulte de $TC \subset C$.

\dgrblock{dgrb677}{From generators to the two factorizations}
\begin{dcomment}{exposition}{p17:par construction), et}{From generators to the two factorizations}
The implication $(a)\Rightarrow(a')$ can be proved completely without the unsupported assertion that $\mathrm{Cof}_W$ and $W^{\mathrm{univ}}$ themselves are left lifting closed. Replace $C_0$ by $C_0\cup TC_0$; this preserves the required inclusions and makes $TC_0\subseteq C_0$. Put
\[
 C=\operatorname{cof}(C_0),\qquad
 TC=\operatorname{cof}(TC_0),\qquad
 TF=C_0^{\square},\qquad F=TC_0^{\square},
\]
where $\operatorname{cof}$ means pushouts, transfinite composites and retracts. The small object theorem identifies these left classes with the corresponding double lifting closures and supplies both factorizations. The closure calculation in the preceding block gives
\[
 C\subseteq\mathrm{Cof}_W,\qquad
 TC\subseteq W^{\mathrm{univ}},\qquad TC\subseteq C.
\]
Condition $(a)$ supplies $TF\subseteq W$ and $F\subseteq\mathrm{Fib}_W$; antitonicity gives $TF\subseteq F$. All of M~1--M~4, M$'$ and M$''$ follow.

This repairs the proof, not by assuming the desired global closure statement, but by applying the theorem to the actual small sets that condition $(a)$ provides. The reverse implication remains a separate issue: a right lifting closed class need not be generated by a set on the left. The order-embedding example above explains why the footnote adding lifting closedness to $(a')$ does not, by itself, supply that missing smallness. Thus $(a'')\Rightarrow(a')$, $(a)\Rightarrow(b)$ and $(b')\Rightarrow(b)$ are immediate, and $(a)\Rightarrow(a')$ is now proved; the remaining arrows in the proposed equivalence require additional arguments.
\end{dcomment}


$a' \Longrightarrow a''$. On veut donc, en plus des conditions M~1 à M~4 et M$'$, M$''$, la condition
\[ TC \ =\ C \cap W\ (= C \cap W^{\univ}). \]
L'idée, c'est de remplacer le couple
\msp{page 18}
\[ C,\ TC \qquad (C \subset \Cof_W,\ TC \subset C \cap W), \]
donnant naissance à
\[ F = (TC)_*, \qquad TF = C_*, \]
par
\[ C,\ C \cap W\ (\supset TC), \]
% foot: Version 25 janvier 2004
\printedpage{11}
% head: A. Grothendieck : Dérivateurs, Ch. XIII, Édition des Universités de Montpellier II et Paris VII
ce qui ne change pas le couple $(C, TF)$, mais bien le couple $(TC, F)$, remplacé par
\[ \underbrace{(TC)'}_{\supset TC} \ =\ \underbrace{C \cap W}_{= C\cap W^{\univ}}, \qquad F' \ =\ \underbrace{(TC')_*}_{\subset F} \]
(on a remplacé $F$ par un $F'$ plus petit) (\footnote[14]{NB $TC'$ sera encore clos à gauche comme intersection de deux parties closes à gauche.}). Le seul problème, c'est si $(TC', F')$ satisfait encore à la condition de factorisation de Quillen. Ce sont des parties closes (à gauche pour $TC'$, à droite pour $F'$) “duales” l'une de l'autre,
\[ TC' = F'^*, \qquad F' = TC'_* \ , \]
mais comme on ne sait pas si $TC'$ est de la forme $\widetilde{TC'_0}$, avec $TC'_0$ petite, on n'est pas sûr que la conclusion du théorème 1 s'applique. Mais si on suppose dans $a'$) que $C$ est une partie \textit{accessible} de $\Fl(\mathcal{M})$ (\textit{i.e.}, étant déjà close à gauche, que $C = \widetilde{C_0}$, $C_0$ petit),
\msp{page 19}
alors je présume qu'on pourra prouver qu'il en est de même de $C \cap W$ ($W$ elle-même étant accessible, grâce à la condition de stabilité par petites $\varinjlim$ filtrantes), donc le théorème 1 s'applique bel et bien à $TC' = C \cap W$. Mais si on garde $a'$) telle quelle, alors il faudrait sans doute prouver que $a'$) \textit{implique} $a'$) sous la forme plus forte, avec $C$ \textit{accessible.} Mais comme on a $a' \Longrightarrow a$, et qu'on vient d'établir que $a \Longrightarrow a'$ sous cette forme plus forte, on gagne (modulo le fait conjectural admis sur les parties admissibles \texttt{[plutôt accessibles]} de $\Fl(\mathcal{M})$).

\dgrblock{dgrb678}{The solution-set repair and Smith's theorem}
\begin{dcomment}{later}{p19:donc le théorème 1 s'applique}{The solution-set repair and Smith's theorem}
Grothendieck is trying to enlarge $TC$ to $C\cap W$ while keeping a factorization theorem. The successful modern recognition theorem places the smallness assumption on the weak equivalences, where it can be checked independently. Smith's theorem, published with proof by \citet[Theorem~1.7]{Beke2000}, assumes a locally presentable category, a set $I$ of arrows, two-out-of-three and retract closure for $W$, the inclusion $I^{\square}\subseteq W$, closure of $\operatorname{cof}(I)\cap W$ under pushouts and transfinite composition, and a solution-set condition for maps from each $i\in I$ into $W$ in the arrow category. It then gives a cofibrantly generated model structure with cofibrations $\operatorname{cof}(I)$ and fibrations $(\operatorname{cof}(I)\cap W)^{\square}$. Accessibility and accessible embeddedness of the full weak-equivalence subcategory imply the required solution-set condition; this is Beke's Proposition~1.15.

Here is the exact consequence for these pages. Retain their hypotheses and condition $(a)$, and assume in addition that the full subcategory of weak equivalences is accessible and accessibly embedded in $\mathcal M^{[1]}$. Use $I=C_0\cup TC_0$ and the two factorizations constructed in the preceding block. Corollary~1 gives $\mathrm{Cof}_W\cap W=W^{\mathrm{univ}}$. Hence
\[
 C\cap W=C\cap W^{\mathrm{univ}}
\]
is closed under pushouts, transfinite composites and retracts. Smith's hypotheses hold, so $C\cap W$ is generated by a set and the desired replacement of $TC$ has factorizations. In fact the result is stronger than the proposed left pre-Q structure: it is a full combinatorial model structure. It is proper because $C\subseteq\mathrm{Cof}_W$ and
\[
 (C\cap W)^{\square}\subseteq TC_0^{\square}
       \subseteq\mathrm{Fib}_W.
\]
The first inclusion uses $TC_0\subseteq C\cap W$.

This application is a repair under an explicit additional hypothesis, not a proof of the original theorem-scholie. Closure under filtered colimits says that colimits of diagrams already in $W$ stay in $W$; accessibility additionally supplies a set of sufficiently presentable arrows from which all of $W$ is built. Nor is the manuscript's parenthetical definition $C=\widetilde C_0$ with $C_0$ small equivalent to accessibility of the full arrow subcategory. The example in $\mathrm{Pos}$ separates those notions even when $C$ is left lifting closed.
\end{dcomment}


Il reste donc à prouver
\[ b \ \Longrightarrow\ a, \qquad b \ \Longrightarrow\ b', \]
et il suffit de prouver la première implication, car on sait que $a$) $\Longrightarrow$ $a'$), et que $a'$) implique que
\[ \Fib_W \cap W \ =\ W_{\univ} \]
(cor.~1, page 10), donc $(\Cof_W)_* \subset W$ et $(W^{\univ})_* \subset \Fib_W$ (\textit{i.e.} $b$)) impliquent, comme $(\Cof_W)_* \subset (W^{\univ})_*$ (car $W^{\univ} \subset \Cof_W$), $(\Cof_W)_* \subset W \cap \Fib_W = W_{\univ}$, ce qui est la partie à établir de $b'$) (l'autre $(W^{\univ})_* \subset \Fib_W$ étant tautologiquement contenue dans $b$).

\msp{page 20}

Notons maintenant que $b$) n'est autre que la variante affaiblie de $a$), quand on omet de spécifier que $C_0$, $TC_0$ doivent être petits. Donc le seul problème, pour $b \Longrightarrow a$, c'est de voir ceci :

\textbf{Lemme conjectural} (\footnote[15]{Ce lemme a l'air un peu faux, mais ça doit pouvoir s'arranger \dots}).

\begin{itemize}[leftmargin=2.5em]
\item[\circled{1}] \textit{Si on a une partie close à gauche} $\Phi$ \textit{de} $\Fl(\mathcal{M})$\textit{, telle que} $\Phi_* \subset \Fib_W$ \textit{(resp.} $\Phi_* \subset W$\textit{), alors il existe une partie} petite $\Phi_0 \subset \Phi$\textit{, telle que} $\Phi_{0*} \subset \Fib_W$ \textit{(resp.} $\Phi_0 \subset W$\textit{).}

\item[\circled{2}] \textit{Si} $C$\textit{,} $C'$ \textit{sont deux parties closes à gauche de} $\Fl(\mathcal{M})$\textit{, si} $C$ \textit{et} $C'$ \textit{sont accessibles, il en est de même de} $C \cap C'$\textit{.}
\end{itemize}
% foot: Version 25 janvier 2004

\printedpage{12}
% head: A. Grothendieck : Dérivateurs, Ch. XIII, Édition des Universités de Montpellier II et Paris VII
Je crois que c'est bel et bien vrai sous les hypothèses faites sur $\mathcal{M}$, $W$ -- qui doivent impliquer que dans un sens convenable, $W$ est une partie accessible de l'ensemble $\Fl(\mathcal{M})$ des flèches de la catégorie accessible $\mathcal{M}$. Il doit en résulter qu'une partie de $\mathcal{M}$ construite explicitement en termes de $W$, telle $\Fib_W$, est elle-même accessible, et le lemme doit sortir de là.

Ainsi, 2° est prouvé, modulo ce lemme d'accessibilité, qu'il me faudra prouver $\pm$ en même temps que le théorème 1 (que je n'ai fait ici qu'énoncer).

\msp{page 21}

\dgrblock{dgrb679}{The intersection lemma has a later affirmative answer}
\begin{dcomment}{later}{p20}{The intersection lemma has a later affirmative answer}
The second part of the conjectural lemma has a precise modern resolution when its word ``accessible'' is read in the sense specified on manuscript p.~18, namely set-generated under left lifting closure. In a locally presentable category, if
\[
 C=\operatorname{cof}(I),\qquad C'=\operatorname{cof}(J)
\]
for sets $I,J$, then $C\cap C'$ is again $\operatorname{cof}(K)$ for a set $K$. This follows from \citet[Theorem~3.2 and Remark~3.8]{MakkaiRosicky2014Cellular}: combinatorial cellular structures are stable under the relevant pseudopullbacks. Apply their result to the two identity functors from $(\mathcal M,C)$ and $(\mathcal M,C')$ into $\mathcal M$ with every arrow cellular. The resulting structure has exactly the intersection as its cellular class. This resolves the intersection assertion; it does not imply that the entire intrinsic class $W^{\mathrm{univ}}$ has such generators.

The first part of the lemma asks for a different kind of compactness: from $\Phi^{\square}\subseteq T$, where $T$ is $W$ or $\mathrm{Fib}_W$, it seeks a set $\Phi_0\subseteq\Phi$ with $\Phi_0^{\square}\subseteq T$. Neither the Galois connection nor the preceding intersection theorem provides this reduction. The final parenthesis in the edition has $\Phi_0\subseteq W$; its intended expression is $\Phi_{0*}\subseteq W$, as the surrounding reduction from $(b)$ to $(a)$ requires. The printed slip is retained and recorded separately.

We found the applicable positive results in the small object theorem, Smith's recognition theorem, and the cellular-category theorem just cited. They do not prove this small-subfamily assertion under precisely the preliminary assumptions of the manuscript. We therefore leave that assertion unresolved in this audit rather than attach a claim of current openness to a loosely specified accessibility principle. In a small category the selection is immediate by taking $\Phi_0=\Phi$; its entire content here is the passage from a proper class to a set. Modern accessible weak factorization systems also require care: accessibility refers to a factorization functor, and does not guarantee that its full left or right class is an accessible subcategory. The corrected statements and the distinction from set generation are recorded in \citet[\S3, especially Remarks~3.2, 3.5 and Theorem~3.9]{Rosicky2021Remarks}.
\end{dcomment}


Prouvons 3°. Supposons d'abord qu'il existe un quadruplet “clos” $(C, TC, F, TF)$, s'insérant dans un “pré-quintuplet de Quillen clos” qui soit une Q-catégorie de modèles à droite, \textit{i.e.}
\[ TF \ =\ F \cap W \ =\ F \cap W_{\univ}. \]
Prouvons alors \texttt{[que]}
\[ \underbrace{(W^{\univ})_* \cap W_{\univ}}_{A_W} \qquad \text{est clos à droite.} \tag*{$(*)$} \]
En effet, on a
\[ TC \ \subset\ W^{\univ}, \]
d'où
\[ (W^{\univ})_* \ \subset\ (TC)_* \ =\ F, \]
donc
\[ A = (W^{\univ})_* \cap W_{\univ} \ \subset\ F \cap W_{\univ} \ \overset{\text{(par hyp.)}}{=}\ TF, \]
or $TF$ est clos à droite et contenu dans $W_{\univ}$, d'où
\[ \overline{A} \ \subset\ TF \ \subset\ W_{\univ}, \]
et comme $(W^{\univ})_*$ est clos à droite, on a aussi
\[ \overline{A} \ \subset\ (W^{\univ})_*, \]
d'où $\overline{A} \subset W_{\univ} \cap (W^{\univ})_* = A$, donc $A$ est clos à droite, q.e.d.

Prouvons la réciproque, \textit{i.e.} supposons satisfaites les conditions de 2°, \textit{plus} la condition $(*)$ ci-dessus. On doit construire le quadruplet.

\msp{page 22}

\dgrblock{dgrb680}{A closure obstruction before the existence problem}
\begin{dcomment}{exposition}{p21}{A closure obstruction before the existence problem}
Grothendieck is trying to extract a property of the weak equivalences which is independent of the auxiliary Quillen structure. The class
\[
 A_W=(W^{\univ})_*\cap W_{\univ}
\]
is designed for this purpose: its first condition is a lifting condition, its second says that every pullback is a weak equivalence. There is a useful separation here between a necessary condition, which is established without an accessibility argument, and a construction of factorizations, which still needs one.

Suppose that a closed weak structure is right exact. Since $TC\subset W^{\univ}$, we have $(W^{\univ})_*\subset F$. Consequently
\[
 A_W= (W^{\univ})_*\cap TF.
\]
Indeed, the forward inclusion uses $TF=F\cap W_{\univ}$, while the reverse inclusion uses $TF\subset W_{\univ}$. This writes $A_W$ as an intersection of two right lifting classes. It is therefore right closed. This short proof is equivalent to the closure computation in the manuscript and shows exactly why no generation hypothesis enters the necessary condition.

The converse is more delicate. Set $L=W^{\univ}$, $F=L_*$, $T=F\cap W_{\univ}$ and $C=T^*$. If $T$ is right closed, then $C_*=T$, and $L\subset C$. These are statements about orthogonality; they do not factor an arbitrary arrow. To obtain the proposed structure, one must additionally know that $(L,L_*)$ and $(T^*,T)$ admit factorizations and that $T^*\subset\Cof_W$. If also $L$ is left closed, these conditions give the two weak factorization systems. Since $L\subset C\cap W$, $T\subset F\cap W$, and $F\subset\Fib_W$, they give a closed weak structure; the right exactness follows from $F\cap W=F\cap W_{\univ}$, using the universal-equivalence criterion already proved on manuscript pp.~10--11.

The sentence on the next page which attributes left lifting closure of $\Cof_W$ and $W^{\univ}$ to part~1 of the conjectural theorem passes over precisely this distinction. Part~1 establishes stability under certain constructions. Left lifting closure is a stronger assertion unless a generation theorem applies. The earlier block \emph{From generators to the two factorizations} supplies the appropriate repair when generators have been specified. We read the diagrams on pp.~22--24 as the formal organization of a proposed construction, with the remaining factorization hypotheses still visible; they are not an additional proof of those hypotheses.
\end{dcomment}


Notons d'abord, pour nous y reconnaître, que les relations $b'$) de 2°), et les relations transposées (compte tenu que $\Cof_W$ et $W^{\univ}$ sont clos à gauche, par 1°), savoir
\[ \left\{ \begin{array}{lllllll}
(\Cof_W)_* & \subset & W_{\univ} & , & (W^{\univ})_* & \subset & \Fib_W \\
(W_{\univ})^* & \subset & \Cof_W & , & (\Fib_W)^* & \subset & W^{\univ},
\end{array} \right. \]
donnant (\footnote[16]{compte tenu de $W^{\univ} \subset \Cof_W$, $W_{\univ} \subset \Fib_W$ et des inclusions transposées.}) les deux diagrammes carrés d'inclusions
% foot: Version 25 janvier 2004
\printedpage{13}
% head: A. Grothendieck : Dérivateurs, Ch. XIII, Édition des Universités de Montpellier II et Paris VII
\[
\begin{tikzcd}[column sep=small, row sep=small]
 & W^{\univ} \arrow[dr, hook] & \\
(\Fib_W)^* \arrow[ur, hook] \arrow[dr, hook] & & \Cof_W \\
 & (W_{\univ})^* \arrow[ur, hook] &
\end{tikzcd}
\qquad\qquad (\text{\footnotemark[17]})
\tag*{(I)}
\]
\footnotetext[17]{\textbf{NB} On a ces carrés d'inclusions, chaque fois qu'on a un système de Quillen $(W, C, F, TC, TF)$ \textit{clos.}}
et
\[
\begin{tikzcd}[column sep=small, row sep=small]
 & (W^{\univ})_* \arrow[dl, hook'] & \\
\Fib_W & & (\Cof_W)_*\ , \arrow[ul, hook'] \arrow[dl, hook'] \\
 & W_{\univ} \arrow[ul, hook', "\alpha"'] &
\end{tikzcd}
\tag*{(II)}
\]
dont le premier est transposé (par $\Phi \longmapsto \Phi^*$) du second, et où les 8 ensembles de flèches envisagés, à l'exception de $\Fib_W$ et $W_{\univ}$ (reliés par $\alpha$ dans le carré II) sont clos (les quatre ensembles de I clos à gauche, et les deux ensembles de II non intéressés par $\alpha$, clos à droite).

\msp{page 23}

Les “pré-quintuplets de Quillen clos”, associés à $W$, correspondent au choix de parties closes à gauche $TC$ et $C$, s'insérant dans le diagramme (I) comme montré dans I$'$ ci-dessous, et donnent lieu par transposition à $F$ et $TF$ comme dans II$'$ ci-dessous,
\[
\begin{tikzcd}[column sep=small, row sep=small]
 & & W^{\univ} \arrow[dr, hook] & \\
 & TC \arrow[ur, hook] \arrow[dr, hook] & & \Cof_W \\
(\Fib_W)^* \arrow[ur, hook] \arrow[dr, hook] & & C \arrow[ur, hook] & \\
 & (W_{\univ})^* \arrow[ur, hook] & &
\end{tikzcd}
\tag*{(I$'$)}
\]
\[
\begin{tikzcd}[column sep=small, row sep=small]
 & & (W^{\univ})_* \arrow[dl, hook'] & \\
 & F \arrow[dl, hook'] & & (\Cof_W)_* \arrow[ul, hook'] \arrow[dl, hook'] \\
\Fib_W & & TF \arrow[ul, hook'] \arrow[dl, hook'] & \\
 & W_{\univ}\ , \arrow[ul, hook'] & &
\end{tikzcd}
\tag*{(II$'$)}
\]
le “hic” étant donc de choisir $TC$ et $C$ assez grandes pour qu'on ait
\[ \underbrace{(TC)_*}_{= F} \ \subset\ \Fib_W, \qquad \underbrace{C_*}_{= TF} \ \subset\ W_{\univ} \qquad \text{--} \]
% foot: Version 25 janvier 2004

% ---- from Der-ChXIII-25-01-04_p016-020.tex
\printedpage{14}
% head: A. Grothendieck : Dérivateurs, Ch. XIII, Édition des Universités de Montpellier II et Paris VII
% foot: Version 25 janvier 2004

\msp{page 24}
\noindent et les inclusions stipulées dans 2$^\circ$, $b'$) assurent qu'on peut bel et bien trouver $TC$ et $C$, quitte à les prendre les plus grandes possibles, $TC = W^{\univ}$, $C = \Cof_W$ -- auquel cas, d'ailleurs, on aura
\[
TC \;=\; C \cap W^{\univ},
\]
\textit{i.e.} on a une pré-Q-catégorie de modèles à gauche. Mais il s'agit d'assurer maintenant, par contre, que
\[
TF \;=\; F \cap W_{\univ},
\]
ce qui implique, on vient de le voir, (en notant $(W^{\univ})_* \subset F$) que
\[
\overline{(W^{\univ})_* \cap W_{\univ}} \;\subset\; \overline{F \cap W_{\univ}} = \overline{TF} = TF \subset W_{\univ}. \tag*{$(*)$}
\]
Inversement, si $(*)$ est satisfait, prenons
\[
F = (W^{\univ})_*
\]
le plus petit possible, \textit{i.e.} $C = W^{\univ}$ \texttt{[plutôt} $TC = W^{\mathtt{univ}}$\texttt{]} le plus grand possible. On doit donc prendre
\[
TF \;=\; F \cap W_{\univ} \;=\; (W^{\univ})_* \cap W_{\univ},
\]
et cet ensemble est bien clos à droite par hypothèse.

\textbf{Corollaire.} \textit{Supposons la condition} b$'$) \textit{du théorème-scholie satisfaite. Les conditions suivantes} a), b) \textit{sur $W$ sont équivalentes (et impliquent la condition de} 3$^\circ$ \textit{du théorème-scholie).}

\msp{page 25}
\begin{itemize}[leftmargin=3em]
\item[\cercle{a}] $(W^{\univ})_* \cap \overline{W_{\univ}} = (\Cof_W)_*$.

\item[\cercle{b}] $\widetilde{W^{\univ} \cup (W_{\univ})}^{*} = \Cof_W$, i.e. $\Cof_W$ \textit{est la partie close à gauche engendrée par} $W^{\univ}$ \textit{et} $(W_{\univ})^*$\textit{,} % doubt resolved: printed tilde ends before the final star; preserve the edition reading, see corrections.csv (PDF leaf 16).
\end{itemize}

\noindent\textit{et ces conditions impliquent}

\begin{itemize}[leftmargin=3em]
\item[\cercle{c}] \textit{qu'il existe une structure de} triple clos de Quillen $(W, C, F)$ \textit{incluant} $W$ \textit{(donc avec} $TC = W \cap C$, $TF = W \cap F$\textit{) (lequel sera clos au sens de Quillen si et seulement si} $W$ \texttt{[est]} \textit{stable par facteurs directs -- car $C$ et $F$ sont déjà stables spontanément).}
\end{itemize}

\noindent Il suffit en effet de prendre alors
\[
TC \;=\; W^{\univ}, \qquad\quad C \;=\; \Cof_W .
\]
Mais je doute que l'existence d'un triple clos de Quillen $(W, C, F)$ suffise à impliquer les conditions \textit{a}), \textit{b}) ci-dessus, ni même la condition plus faible
\[
(W^{\univ})_* \cap W_{\univ} \;=\; (\Cof_W)_*
\]
qui signifie \textit{exactement} que $C = \Cof_W$, $TC = W^{\univ}$, et leurs transposés $TF = (\Cof_W)_*$, $F = (W^{\univ})_*$ forment une catégorie de Quillen close.

\clearpage
\printedpage{15}
% head: A. Grothendieck : Dérivateurs, Ch. XIII, Édition des Universités de Montpellier II et Paris VII
% foot: Version 25 janvier 2004

\msp{page 26}
\noindent Conditions sur $(\mathcal{M}, W)$ avec
\begin{itemize}[leftmargin=3em]
\item[\textit{a})] $\mathcal{M}$ stable par $\varprojlim$ finies et $\varinjlim$ (petites),
\item[\textit{b})] $W$ stable par $\varinjlim$ filtrantes,
\item[\textit{c})] $\mathcal{M}$ accessible.
\end{itemize}

\dgrblock{dgrb681}{The two closures in the maximal construction}
\begin{dcomment}{exposition}{p25}{The two closures in the maximal construction}
The corollary seeks a particularly economical structure: take every universal weak cofibration as a trivial cofibration and every $W$-cofibration as a cofibration. Its strength comes from this maximal choice, and the two different closures in the formulas measure its cost.

For arbitrary classes $L,R$ of arrows, the lifting calculus gives
\[
 \bigl(L\cup R^*\bigr)_*=L_*\cap(R^*)_*
       =L_*\cap\overline R.
\]
Thus the right orthogonal of $W^{\univ}\cup(W_{\univ})^*$ is
$(W^{\univ})_*\cap\overline{W_{\univ}}$, with a bar on the second factor. It is not, without another hypothesis, $(W^{\univ})_*\cap W_{\univ}$. If $\Cof_W$ is left closed, applying the Galois correspondence proves the equivalence
\[
 (W^{\univ})_*\cap\overline{W_{\univ}}=(\Cof_W)_*
 \quad\Longleftrightarrow\quad
 \widetilde{\,W^{\univ}\cup(W_{\univ})^*\,}=\Cof_W.
\]
The superscript star must belong to the second summand inside the tilde. The printed formula on manuscript p.~25 places it outside the visible tilde; the accompanying prose and the summary on p.~26 identify the intended expression. We preserve the edition's formula and record the discrepancy separately.

Under condition $b'$ one has $(\Cof_W)_*\subset W_{\univ}$. Hence the stronger displayed equality implies
\[
 (W^{\univ})_*\cap W_{\univ}=(\Cof_W)_*.
\]
For the inclusion from right to left, use $W^{\univ}\subset\Cof_W$ and $b'$; for the reverse inclusion, enlarge $W_{\univ}$ to its right closure and apply the stronger equality. The latter equality is exactly the compatibility needed by the maximal two pairs, once their factorizations exist. It does not itself supply factorizations.

This explains the hierarchy I--IV on p.~26. At the level of formal implications, IV imposes a closure identity stronger than the identity in III, and a maximal structure in III is a special instance of a structure in II. The alleged equivalences inside I and the construction clauses inside II retain the accessibility and factorization problems from pp.~17--20. The hierarchy is a research programme with several proved implications, not four unconditional recognition theorems. In particular, the correct formula in II(c) is $C=(TF)^*$; the printed subscript star is a variance slip.
\end{dcomment}


\[
\mathrm{I}\ \left\{
\begin{array}{@{}l@{\quad}p{11cm}@{}}
\cercle{a} & existence d'un système de Quillen faible $(W, C, TC, F, TF)$ (clos si on y tient) \\[1ex]
\Updownarrow & \\[1ex]
\cercle{b} & existence d'un système de Quillen faible exact à gauche ($TC = C \cap W$). \\[1ex]
\Updownarrow & \\[1ex]
\cercle{c} & $\underbrace{(\Cof_W)_* \subset W,\ (W^{\univ})_* \subset \Fib_W}_{(*)}$ ($\Longleftrightarrow$ $(\Cof_W)_* \subset W_{\univ}$, $(W^{\univ})_* \subset \Fib_W$).
\end{array}\right.
\]

\noindent$\Uparrow$
\[
\mathrm{II}\ \left\{
\begin{array}{@{}l@{\quad}p{11cm}@{}}
\cercle{a} & existence d'un système de Quillen faible exact à droite ($TF = F \cap W$) \\[1ex]
\Updownarrow & \\[1ex]
\cercle{d} & On a $(*)$, et de plus $(W^{\univ})_* \cap W_{\univ}$ clos à droite. \\[1ex]
\Updownarrow & \\[1ex]
\cercle{b} & existence d'un système de Quillen clos (closed model category), $TC = C \cap W$, $TF = F \cap W$, $F = (TC)_*$, $TF = C_*$, $TC = F^*$, $C = (TF)^*$. \\[1ex]
\Updownarrow & \\[1ex]
\cercle{c} & Posant $TC = W^{\univ}$, $F = (TC)_*$, $TF = F \cap W_{\univ}$, $C = (TF)_*$, on trouve avec $W$ un système de Quillen clos. % [sic] $C = (TF)_*$ -> $C = (TF)^*$, cf. the line b) above where $C = (TF)^*$
\end{array}\right.
\]

\noindent$\Uparrow$
\[
\mathrm{III}\ \left\{
\begin{array}{@{}p{11.5cm}@{}}
$C = \Cof_W$ et $TC = W^{\univ}$ font partie (avec $W$) d'un système de Quillen clos (donc $F = (W^{\univ})_*$, $TF = (\Cof_W)_*$, $W = TF \circ TC$). \\[1ex]
$\Updownarrow$ \\[1ex]
\texttt{[On a]} $(*)$, et de plus $(W^{\univ})_* \cap W_{\univ} = (\Cof_W)_*$.
\end{array}\right.
\]

\noindent$\Uparrow$
\[
\mathrm{IV}\ \left\{
\begin{array}{@{}p{11.5cm}@{}}
On a $(*)$ et $\Cof_W = \widetilde{(W_{\univ})^* \cup W^{\univ}}$\ \ . \\[1ex]
$\Updownarrow$ \\[1ex]
On a $(*)$ et $(W^{\univ})_* \cap \overline{W_{\univ}} = (\Cof_W)_*$.
\end{array}\right.
\]

\clearpage
\printedpage{16}
% head: A. Grothendieck : Dérivateurs, Ch. XIII, Édition des Universités de Montpellier II et Paris VII
% foot: Version 25 janvier 2004

\noindent Cas $\mathcal{M} = \mathrm{Cat}$, $W = \mathbf{W}_\infty$. On sait que II est satisfait, par \textsc{Thomason}. Donc on a
\[
(\Cof_W)_* \subset W_{\univ}, \qquad\quad (W^{\univ})_* \subset \Fib_W, \qquad\quad (W^{\univ})_* \cap W_{\univ} \ \text{ clos à droite.}
\]
J'ai pu prouver directement les deux premières inclusions (\textit{i.e.} I), mais pas que $(W^{\univ})_* \cap W_{\univ}$ soit clos à droite. (Je devrais au moins essayer!) (\footnote[18]{Il suffirait de savoir que $\overline{W_{\univ}} \cap \Fib_W = W_{\univ}$, ce qui impliquerait que III et IV sont équivalents (\footnotemark[19]).}\footnotetext[19]{Mais c'est faux, p. ex. si $X$ \texttt{[est une]} catégorie discrète non vide, $X \longrightarrow e$ est dans $\overline{W_{\univ}} \cap \Fib_W$.}). On n'a pas IV, \textit{i.e.} on a
\[
(W^{\univ})_* \cap \overline{W_{\univ}} \;\neq\; (\Cof_W)_*
\]
car
\[
(W^{\univ})_* \cap \overline{W_{\univ}} \not\subset W_{\univ}
\]
(contient les $X \longrightarrow e$ avec $X$ discrète), donc on a
\[
(W^{\univ})_* \cap \overline{W_{\univ}} \;\;\underset{\neq}{\supset}\;\; \underbrace{(W^{\univ})_* \cap W_{\univ}}_{\text{n'est \textit{pas} l'orthogonal de } W^{\univ} \cup (W_{\univ})^*} \quad . % doubt resolved: printed strict-superset sign confirmed; see corrections.csv (PDF leaf 18).
\]

\msp{page 27}

\dgrblock{dgrb682}{Discrete categories detect the extra closure}
\begin{dcomment}{development}{p26}{Discrete categories detect the extra closure}
Grothendieck's example deserves a proof because it distinguishes two conditions that otherwise look nearly identical. We take $W$ in $\mathrm{Cat}$ to be the functors inducing weak homotopy equivalences of nerves. Thomason's theorem gives a proper model structure with this $W$; its existence does not imply the maximal closure condition IV \citep[Theorem~4.9 and Corollary~5.5]{Thomason1980}.

\noindent\textbf{A discrete obstruction to condition IV.}\ours\quad
For every nonempty discrete category $D$, its terminal map belongs to
\[
 (W^{\univ})_*\cap\overline{W_{\univ}}\cap\Fib_W.
\]
If $D$ has at least two objects, this map does not belong to $W_{\univ}$.

\emph{Proof.} We first examine a functor $i:A\to B$ having the left lifting property against every universally weakly equivalent functor. We claim that $\pi_0(i):\pi_0A\to\pi_0B$ is injective. Suppose that two distinct components of $A$ have images connected in $B$ by a zigzag of length $m$. Choose a finite poset $J$ whose Hasse diagram is an alternating linear zigzag with more than $m$ edges. Its nerve is an interval subdivided into edges, so $J\to e$ is a weak equivalence. It is universally so: its pullback over any category $E$ is the projection $J\times E\to E$, and the nerve preserves products; a product with this finite contractible simplicial set is a weak homotopy equivalence.

Send the first chosen component of $A$ to one endpoint of $J$, the second to the other endpoint, and every other component to either endpoint. This defines a functor $a:A\to J$, constant on components. A lift against $J\to e$ would extend $a$ to $B\to J$. The zigzag of length $m$ in $B$ would then connect the two endpoints of $J$ by at most $m$ edges. Every nonidentity arrow of the alternating poset joins adjacent vertices, so this is impossible. This proves the claim.

Now a functor $A\to D$ is just a function $\pi_0A\to D$. If $i\in(W_{\univ})^*$, injectivity of $\pi_0(i)$ lets us extend this function to $\pi_0B$: use its prescribed values on the image and one fixed object of $D$ on all remaining components. Thus $D\to e$ has the right lifting property against $(W_{\univ})^*$, which is exactly membership in $\overline{W_{\univ}}$.

Every functor in $W^{\univ}$ is in $W$, so induces a bijection on components. The same extension argument, now with a unique possible value on every component, gives $D\to e\in(W^{\univ})_*$. Finally, any base change of $D\to e$ is a coproduct of copies of an identity map. Pulling a weak equivalence back along it gives the corresponding coproduct of that weak equivalence; the nerve and geometric realization verify that the result is again a weak equivalence. Therefore $D\to e\in\Fib_W$. If $|D|\geq2$, the map fails already to induce a bijection on components, so is not in $W$, and a fortiori not in $W_{\univ}$.\hfill$\square$

This also explains the qualification required in the footnote: a nonempty discrete category supplies the displayed membership, but a counterexample requires at least two components. The empty category is not included in the lifting argument because extending a function to a newly appearing component would then be impossible. The proof expands the manuscript's example; we checked Thomason's construction and the cited modern treatments of lifting and localization, without finding this precise zigzag proof there. The marker records a proof supplied for this commentary, not a claim of priority.
\end{dcomment}


\section*{3\quad Catégorie de Quillen (deuxième mouture)}

On se donne toujours un couple
\[
(\mathcal{M}, W), \qquad\quad W \subset \mathrm{Fl}(\mathcal{M}), \qquad\quad W \text{ faiblement saturé,}
\]
$\mathcal{M}$ satisfait M 0 (stabilité par $\varprojlim$ finies et $\varinjlim$ finies).

\noindent\textbf{Définition 1 :} Une \textit{structure de Quillen faible} associée à la catégorie de modèles $(\mathcal{M}, W)$ est la donnée de quatre parties $(C, TC, F, TF)$ de $\mathrm{Fl}(\mathcal{M})$, satisfaisant les conditions suivantes.

\begin{itemize}[leftmargin=3em,labelsep=0.5em,align=left]
\item[\textbf{M 1}] $C \longleftrightarrow TF$, $TC \longleftrightarrow F$ (\textit{i.e.} $TF \subset C_*$, $F \subset (TC)_*$).
\item[\textbf{M 2}] Propriété de factorisation de Quillen pour $(C, TF)$ et $(TC, F)$.
\item[\textbf{M 3}] Stabilité de $C$ par composition et par cochangement de base, de $F$ par composition et changement de base -- $C$ et $F$ contiennent les isomorphismes.
\item[\textbf{M 4}] Stabilité similaire pour $TC$ et $TF$.
\item[$\mathbf{M}'$] $TC \subset C \cap W$, $TF \subset F \cap W$.
\item[$\mathbf{M}''$] $C \subset \Cof_W$, $F \subset \Fib_W$.
\end{itemize}

\noindent On dit que la structure de Quillen est \textit{exacte} si elle satisfait à l'axiome supplémentaire

\clearpage
\printedpage{17}
% head: A. Grothendieck : Dérivateurs, Ch. XIII, Édition des Universités de Montpellier II et Paris VII
% foot: Version 25 janvier 2004

\noindent\textbf{M 2 bis}\qquad $W = TF \circ TC$.

\noindent On dit qu'elle est \textit{exacte à gauche} si elle satisfait l'axiome supplémentaire

\noindent\textbf{M$'_0$ (g)}\qquad $TC = C \cap W$ (au lieu de $TC \subset C \cap W$ seulement),

\noindent\textit{exacte à droite} si elle satisfait

\noindent\textbf{M$'_0$ (d)}\qquad $TF = F \cap W$ (au lieu de $TF \subset F \cap W$ seulement) (\footnote[20]{\textbf{NB} Si $C$, $F$, $W$ \texttt{[sont]} stables par facteurs directs (donc aussi $C \cap W$ et $F \cap W$), alors l'exactitude à gauche $TC = C \cap W$ \textit{équivaut} à l'exactitude à droite $TF = F \cap W$, et implique que $C$, $TC$, $F$, $TF$ sont \textit{clos} (à gauche pour $C$, $TC$, à droite pour $F$, $TF$).}).

\msp{page 28}
\noindent\textbf{Remarque.}
\begin{itemize}[leftmargin=3em]
\item[a)] Une structure de Quillen exacte à gauche ou exacte à droite est exacte.
\item[b)] Si on a une donnée $(C, TC, F, TF)$ satisfaisant M 1, M 2, M$'$, $M''$, quand on sature $C$, $TC$ à gauche et $F$, $TF$ à droite, pour les stabilités précisées dans M 2, M 3, on trouve une structure de Quillen. (Tous les axiomes sont évidents, sauf M$'$, pour lequel je renvoie à p. 9, 10, notamment prop. 4.) Si on avait M 2 bis, cela sera encore vérifié (a fortiori) après saturation, \textit{i.e.} on a une structure de Quillen exacte.
\end{itemize}

\dgrblock{dgrb683}{Exactness belongs to the chosen classes}
\begin{dcomment}{exposition}{p27}{Exactness belongs to the chosen classes}
The second definition rearranges the earlier axioms around the task they have to perform. The inclusions $TC\subset C\cap W$ and $TF\subset F\cap W$ distinguish designated trivial maps from all weak equivalences among the designated cofibrations and fibrations. This distinction is essential before retract closure has been imposed. In modern terminology, the closed case consists of two nested weak factorization systems; Grothendieck's $\text{M}''$ adds properness relative to the specified weak equivalences.

There are three different assertions. Exactness says that every weak equivalence has a factorization using the two designated trivial classes. Left exactness identifies the designated trivial cofibrations with $C\cap W$; right exactness identifies the designated trivial fibrations with $F\cap W$. Either identification implies exactness. For example, if $w=pi$ with $i\in C$ and $p\in TF$, then $w,p\in W$ imply $i\in W$, and left exactness gives $i\in TC$. The dual argument uses the other factorization.

Neither this argument nor the two-out-of-three axiom proves that the original chosen class $TF$ is retract closed. Footnote~20 tries to infer too much from retract closure of $C,F,W$. Its conclusion becomes correct after replacing the missing trivial class by its retract closure, but the chosen class need not already equal that closure. The next example makes the distinction concrete. For closed structures, by contrast, exactness alone forces both equalities by the retract argument developed later in the commentary on manuscript p.~57.
\end{dcomment}


\noindent\textbf{Définition 1 bis.} La structure de Quillen faible est dite \textit{structure de Quillen} (tout court) si elle est exacte à gauche et à droite, ou (comme on dit) “biexacte”, alors $(C, F)$ determine déjà $(C, TC, F, TF)$. (Ça correspond exactement aux “model categories” de \textsc{Quillen} \texttt{[propres]}.) On dit que la structure de Quillen faible est \textit{close} si on a % [sic] determine -> détermine, missing accent
\[
C \clarrow TF\ , \qquad\quad TC \clarrow F\ ,
\]
\textit{i.e.} $TF = C_*$, $C = (TF)^*$ et $F = (TC)_*$, $TC = F^*$ (\footnote[21]{Cela signifie aussi que $C$, $TF$, $TC$, $F$ sont stables par facteurs directs. On dit que la structure de Quillen faible est \textit{close à gauche} (resp. \textit{à droite}) si $C$, $TC$ (resp. $F$, $TF$) sont clos à gauche (resp. à droite).}). Une structure de Quillen faible à la fois biexacte (\textit{i.e.} pas faible du tout!) et close est appelée structure de Quillen \textit{parfaite}. (Ça correspond exactement à la “closed model category” \texttt{[propre]} de \textsc{Quillen}.)

\msp{page 29}
\noindent\textbf{Remarque.} Pour qu'une structure de Quillen faible soit parfaite, il suffit qu'elle soit biexacte et qu'on ait
\[
C = (TF)^*, \qquad\quad F = (TC)_*
\]
(ce qui implique déjà que $TC = F^*$, $TF = C_*$, \textit{i.e.} que $TC$ est clos à gauche et $TF$ clos à droite). C'est prouvé dans Quillen. Ça implique aussi que $W$ est fortement saturé -- c'est le pays de cocagne !

\dgrblock{dgrb684}{A retract condition missing from footnote 20}
\begin{dcomment}{development}{p28}{A retract condition missing from footnote 20}
Grothendieck wants the ordinary model-category identities to reappear as soon as retracts behave well. They do reappear after completing the designated classes, but not necessarily before. The following example isolates the missing step without involving any topological pathology.

\noindent\textbf{Left exactness without right exactness.}\ours\quad
Assume the axiom of choice. In $\mathrm{Set}$ put
\[
 W=\Fl(\mathrm{Set}),\qquad C=TC=\{\text{injections}\},
 \qquad F=\{\text{surjections}\},
\]
and let $TF$ consist of those surjections all of whose fibres are either singletons or infinite sets. These classes satisfy all the weak Quillen axioms on manuscript p.~27. The classes $C,F,W$ are retract closed, and the structure is left exact, but it is not right exact.

\emph{Proof.} Injections have the left lifting property against surjections: on the complement of the injected subset, choose a preimage in each relevant nonempty fibre. Both required lifting properties follow. For an arbitrary map $f:A\to B$, use
\[
 A\ \longrightarrow\ A\amalg(B\times\mathbb N)
 \ \xrightarrow{\ p\ }\ B,
 \qquad p|_A=f,\quad p(b,n)=b.
\]
The first map is injective and every fibre of the second is infinite; when $B$ is empty, necessarily $A$ is empty and the assertion is vacuous. This one factorization serves both required pairs.

Injections contain the isomorphisms and are stable under composition and pushout. Surjections have the corresponding pullback and composition properties. For the special class $TF$, pullback preserves each fibre up to bijection. A fibre of a composite is a disjoint union of nonempty fibres of the first map indexed by a fibre of the second. If the indexing fibre is a singleton, the result is a permitted fibre of the first map; if it is infinite, the disjoint union is infinite. Thus $TF$ also has the required stability. The inclusions $TC\subset C\cap W$ and $TF\subset F\cap W$ hold. Since all maps are weak equivalences, $\Cof_W$ and $\Fib_W$ are both the class of all maps, so $\text{M}''$ holds as well.

The usual elementwise arguments show that injections and surjections are retract closed in the arrow category, while $W$ trivially is. We have $TC=C\cap W$. However, the map from a two-element set to a singleton belongs to $F\cap W$ and does not belong to $TF$. It is a retract of $\mathbb N\to\{*\}$, which does belong to $TF$. This proves both the failure of right exactness and the failure of retract closure of the designated class.\hfill$\square$

The repaired assertion can be stated precisely. Suppose $C,F,W$ are retract closed and $TC=C\cap W$. Then $TC$ is retract closed. The factorization and lifting axioms give $TC={}^\perp F$ and $F=TC^\perp$. If $p\in F\cap W$, factor $p=qi$ with $i\in C$ and $q\in TF$. Two-out-of-three gives $i\in TC$, and lifting $i$ against $p$ exhibits $p$ as a retract of $q$. Conversely every retract of a map in $TF$ lies in $F\cap W$. Hence
\[
 \operatorname{Retr}(TF)=F\cap W.
\]
Replacing $TF$ by this class preserves lifting, factorization and properness and produces a perfect structure. The standard retract argument also gives $C={}^\perp(F\cap W)$ and $C^\perp=F\cap W$. The dual repair applies to a right exact structure. We checked the standard model-category retract arguments and Cisinski's treatments; the particular fibre-cardinality example is supplied here to settle this exact footnote, without a priority claim.
\end{dcomment}


\clearpage
\printedpage{18}
% head: A. Grothendieck : Dérivateurs, Ch. XIII, Édition des Universités de Montpellier II et Paris VII
% foot: Version 25 janvier 2004

\noindent\textbf{Scholie :} Le couple $TC \longleftrightarrow F$, satisfaisant à la condition de factorisation de Quillen, et satisfaisant $TC \subset W$, $F \subset \Fib_W$, permet déjà de développer une théorie satisfaisante des fibres $W$-homotopiques dans $\mathcal{M}$ -- mais sans propriété d'exactitude pour les carrés $W$-cartésiens. Dualement, la donnée de $C \longleftrightarrow TF$, avec factorisation, et $TF \subset W$, $C \subset \mathrm{Cofib}_W$, donne une bonne théorie des cofibres $W$-homotopiques, mais sans propriété d'exactitude pour les carrés $W$-cocartésiens. La structure de Quillen faible réunit ces deux données duales l'une de l'autre, avec en plus le lien supplémentaire entre ces deux donnés, par $TC \subset C$, $TF \subset F$ (stipulé dans $M'$). Cette donnée assure déjà que % [sic] donnés -> données, feminine noun as in "ces deux données duales"
\[
\underbrace{\mathcal{M}_{\mathrm{cf}}}_{\substack{\text{sous-catégorie pleine}\\ \text{de $\mathcal{M}$ formée des}\\ \text{objets fibrants et}\\ \text{cofibrants à la fois}}} W_{\mathrm{cf}}^{-1} \;\xrightarrow{\ \sim\ }\; \mathcal{M}W^{-1} = \mathrm{Ho}_W
\]
\msp{page 30}
est une équivalence, mais non pas que le foncteur localisation
\[
\mathcal{M}_{\mathrm{cf}} \longrightarrow \mathcal{M}_{\mathrm{cf}}W_{\mathrm{cf}}^{-1}
\]
est un simple passage au quotient sur les flèches, \textit{i.e.} que si $X$, $Y$ sont dans $\mathcal{M}_{\mathrm{cf}}$, alors
\[
\mathrm{Hom}_{\mathcal{M}}(X, Y) \;\longrightarrow\; \underbrace{\mathrm{Hom}_W(X, Y)}_{\overset{\text{déf}}{=}\,\mathrm{Hom}_{\mathrm{Ho}_W}(X, Y)} \tag*{$(*)$}
\]
soit surjectif. Ceci est assuré par l'axiome d'exactitude M 2 bis,
\[
W \;=\; TF \circ TC.
\]
Avec ce seul axiome, on n'a pas de caractérisation sympathique de la relation d'équivalence dans $\mathrm{Hom}_{\mathcal{M}}(X, Y)$ définie par la flèche $(*)$ (pour $X$, $Y$ dans $\mathcal{M}_{\mathrm{cf}}$).

\dgrblock{dgrb685}{Replacement controls arrows as well as objects}
\begin{dcomment}{exposition}{p29}{Replacement controls arrows as well as objects}
The passage to fibrant-cofibrant objects is the first point where the lifting axioms recover information which would be lost by merely choosing a weakly equivalent object. Grothendieck's assertion of an equivalence of localizations is correct under the weak axioms. The mechanism is worth making explicit because exactness will be needed only for the next assertion, that ordinary arrows represent all localized arrows.

Choose $q_X:QX\to X$ in $TF$ with $QX$ cofibrant. Given $f:X\to Y$, lift against $q_Y$ to obtain $Qf:QX\to QY$ with $q_YQf=fq_X$. This lift need not be unique as an arrow. It is unique in the localization of the cofibrant objects. Indeed, for two lifts $a,b:QX\to QY$, form $P=QY\times_YQY$. Its projections belong to $TF$ by pullback stability. Choose a cofibrant replacement $q_P:QP\to P$ and lift $(a,b)$ through it. The diagonal $QY\to P$ likewise lifts through $q_P$ because $QY$ is cofibrant. Call this lift $d:QY\to QP$. The two composites
\[
 QY\xrightarrow{d}QP\xrightarrow{\operatorname{pr}_j q_P}QY
\]
are identities. Each $\operatorname{pr}_j q_P$ belongs to $TF$, so $d\in W$. After localization both projections are therefore the inverse of $d$, and $a=b$ there. This uniqueness proves compatibility with identities and composition. Two-out-of-three shows that a weak equivalence $f$ gives a weak equivalence $Qf$. Thus $Q$ defines an inverse, after localization, to the inclusion of the cofibrant objects.

For the second step choose $r_X:X\to RX$ in $TC$, with $RX$ fibrant, when $X$ is cofibrant. Then $RX$ remains cofibrant. An arrow $X\to Y$ extends to $RX\to RY$ by lifting $r_X$ against $RY\to e$. To compare two extensions, form $P=RX\amalg_XRX$. The two inclusions $RX\to P$ are in $TC$. Choose $P\to RP$ in $TC$ with $RP$ fibrant. The fold $P\to RX$ extends to $d:RP\to RX$ because $RX$ is fibrant. If $v_0,v_1:RX\to RP$ are the two induced maps, then $dv_0=dv_1=1$, and $d,v_0,v_1$ are weak equivalences. The two extensions consequently agree after localization, by first extending their common pushout map $P\to RY$ to $RP\to RY$. This gives the inverse equivalence for the inclusion of fibrant-cofibrant objects into cofibrant objects.

No equality $TC=C\cap W$ or $TF=F\cap W$ entered these two arguments. What is indispensable is the ability to compare choices inside the relevant full subcategory. This is the same distinction already emphasized in the commentary on Chapter~VI, manuscript p.~87. For the usual model-category statement, see \citet[Theorem~2.2.15 and Corollary~2.2.18]{Cisinski2019HigherCategories}; the argument above records the weaker hypotheses actually used here.
\end{dcomment}


Si on a exactitude à gauche,
\[
TC = C \cap W \qquad\quad (\Longleftrightarrow\ TC = C \cap W^{\univ}),
\]
alors cette relation d'équivalence est celle engendrée par la relation de \textit{Q-homotopie à droite} (\footnote[22]{Il semblerait que la relation de Q-homotopie à droite est une relation d'équivalence, dès que $Y$ \texttt{[est]} dans $\mathcal{M}_{\mathrm{f}}$ et sans supposer la structure de Quillen faible exacte à droite ou à gauche.}, \footnote[23]{\textbf{NB} D'après \textsc{Baues}, il semblerait que sur $\mathcal{M}_{\mathrm{cf}}$, cette relation d'homotopie à droite équivaut à la relation d'homotopie à gauche, et que celle-ci ne dépend pas du cylindre choisi \ldots}) (impliquée par la Q-homotopie à gauche), et la catégorie quotient ainsi définie
\[
\pi_{\mathrm{d}}\mathcal{M}_{\mathrm{cf}} \;=\; \pi\mathcal{M}_{\mathrm{cf}} \;\simeq\; \mathrm{Ho}_W,
\]
s'envoie dans $\pi_{\mathrm{d}}\mathcal{M}_{\mathrm{c}}$,
\[
\pi\mathcal{M}_{\mathrm{cf}} \longrightarrow \pi_{\mathrm{d}}\mathcal{M}_{\mathrm{c}},
\]
et ce foncteur peut être interprété (sauf erreur)

\dgrblock{dgrb686}{Two path conventions and a nontransitive homotopy}
\begin{dcomment}{examples}{p30:alors cette relation d'équivalence est celle engendrée}{Two path conventions and a nontransitive homotopy}
Footnote~22 requires a distinction which the text does not explicitly make. A diagonal factorization
\[
 Y\xrightarrow{s}PY\xrightarrow{p=(d_0,d_1)}Y\times Y
\]
may be required to have $s\in TC$ and $p\in F$, or only $s\in W$ and $p\in F$. Call these, respectively, a \emph{designated} path object and a \emph{broad} path object. In standard model-category terminology the first is a very good path object and the second a good path object. The raw relation associated with a fixed path object declares $f,g:X\to Y$ related when $(f,g)=pH$ for some $H:X\to PY$. Taking the union of these relations over all broad path objects is another operation; it must not be silently identified with the relation from a designated one.

Every designated path object gives a reflexive and symmetric relation, for every source. Reflexivity uses $sf$. For symmetry, lift $s$ against $p$, with top map $s$ and bottom map $\tau p$, where $\tau$ interchanges the two factors of $Y^2$. The filler $r:PY\to PY$ satisfies $pr=\tau p$, and $rH$ reverses a homotopy. Moreover, if $Y\xrightarrow{s'}P'Y\xrightarrow{p'}Y^2$ is another designated path object, the square with vertical maps $s,p'$ and horizontal maps $s',p$ has a filler $PY\to P'Y$ preserving the endpoints. Interchanging the objects gives the converse comparison. Thus \emph{one fixed designated path object and all designated path objects define the same raw relation}, without fibrancy, cofibrancy or retract-closure assumptions. Transitivity is the additional issue.

\noindent\textbf{A counterexample with every object fibrant and cofibrant.}
The following is a dual application of the graph example of \citet[Definition~3.1, Observation~3.3 and Example~3.6]{KurzRosicky2005}. Let $\mathcal G$ be the category of reflexive directed multigraphs. An object consists of sets of vertices and edges, maps $s,t:E\rightrightarrows V$, and a distinguished loop $e:V\to E$ satisfying $se=te=1_V$. Morphisms preserve these data. This is a functor category, hence has all small limits and colimits. Let
\[
 L=\operatorname{Mono}(\mathcal G),\qquad R=L^{\square}.
\]
Here monomorphisms are injective on both vertices and edges. For precision, $R$ is defined by the literal lifting property, not by an abbreviated assertion about whether two image vertices are adjacent: in a multigraph a lifting problem specifies the target edge as well as its endpoints.

The pair $(L,R)$ is a weak factorization system. One generating map adds a vertex with its distinguished loop to the empty graph. The other adds a directed edge between the two vertices of a discrete two-vertex reflexive graph. In a pushout the two endpoints may coincide, so the latter generator also adds a nondistinguished loop. Every subgraph inclusion is a composite of pushouts of these maps: first adjoin its missing vertices, then its missing edges. Conversely such pushouts, their composites and their retracts are monomorphisms. The finite generating graphs are finitely presentable, so the small object theorem proved earlier gives exactly $(L,L^{\square})$.

Pass to $\mathcal M=\mathcal G^{\mathrm{op}}$ and put
\[
 W=C=\operatorname{Mor}(\mathcal M),\qquad TF=\operatorname{Iso}(\mathcal M),
 \qquad TC=R^{\mathrm{op}},\qquad F=L^{\mathrm{op}}.
\]
Both $(C,TF)$ and $(TC,F)$ are weak factorization systems. In particular their classes contain the isomorphisms and have the required composition, pushout or pullback stability, and even retract stability. The inclusions $TC\subset C\cap W$ and $TF\subset F\cap W$ hold. Since every arrow is weak, the intrinsic classes $\Cof_W$ and $\Fib_W$ also contain every arrow. Thus all the printed weak axioms, including $\text{M}''$, hold.

Every object is cofibrant because $C$ consists of all arrows. Every object is fibrant because its map to the terminal object of $\mathcal M$ is opposite to an inclusion of the empty graph, hence belongs to $L^{\mathrm{op}}$. These observations make the example independent of any failure to replace its source or target.

Let $V$ denote the graph with one vertex and its single distinguished loop. Let $D$ have two vertices $0,1$ and exactly one directed edge for each ordered pair of vertices, with the two diagonal edges distinguished. The inclusion $j:V\amalg V\to D$ belongs to $L$. The unique map $q:D\to V$ belongs to $R$: a graph map from a subgraph to $D$ extends by sending each additional vertex to $0$, and then sending each edge to the unique edge with the prescribed endpoints. This extends every lifting problem against a monomorphism. The equality $qj=\nabla_V$ therefore becomes a designated path factorization in the opposite category,
\[
 Y\xrightarrow{q^{\mathrm{op}}}D
     \xrightarrow{j^{\mathrm{op}}}Y\times Y,
 \qquad Y=V\quad\text{viewed in }\mathcal M.
\]

Now let $X$ be the reflexive graph with vertices $a,b,c$, distinguished loops, and four further edges
\[
 a\rightleftarrows b\rightleftarrows c,
\]
with no edge from $a$ to $c$ or from $c$ to $a$, again viewed as an object of $\mathcal M$. The three vertex maps $V\to X$ in $\mathcal G$ are three arrows $f_a,f_b,f_c:X\to Y$ in $\mathcal M$. A right homotopy through the displayed path object is precisely a graph map $D\to X$ with the specified vertex images. Such maps exist for $(a,b)$ and $(b,c)$, whereas none exists for $(a,c)$. Consequently
\[
 f_a\sim_r f_b,\qquad f_b\sim_r f_c,
 \qquad f_a\not\sim_r f_c.
\]
The comparison argument above shows that allowing every designated path object does not change this calculation. Raw designated right homotopy is not transitive even though every object is fibrant and cofibrant and all four lifting classes are retract closed.

This weak structure is not exact. Indeed $TF\circ TC=TC$, and the opposite of a vertex inclusion $V\to X$ is not in $TC$: the vertex inclusion fails the right lifting property against the map adding a vertex. Thus $TC\ne W$. There is also a useful exact variant: replace $C$ by $TC$ and $TF$ by $F$, retaining $W,TC,F$. This is a perfect model structure with all arrows weak. The same path calculation still fails transitivity for the same source, which is now noncofibrant. To see the latter directly, $X\to V$ in $\mathcal G$ does not lift the edge-adding generator with endpoints mapped to $a,c$. Hence even exactness does not make the raw relation from very good paths transitive for unrestricted sources.

\noindent\textbf{The broad relation for a fibrant target.}
Under the printed weak axioms, the relation defined by existence of a homotopy through \emph{some} broad path object is an equivalence relation whenever $Y$ is fibrant, with arbitrary source $X$. Here is the direct argument. Reflexivity uses the constant section of any designated, hence broad, path object. Symmetry uses the same broad object with its endpoints interchanged. For transitivity choose broad path objects
\[
 Y\xrightarrow{s_P}P\xrightarrow{(d_0,d_1)}Y^2,
 \qquad
 Y\xrightarrow{s_Q}Q\xrightarrow{(e_0,e_1)}Y^2
\]
and homotopies from $f$ to $g$ through $P$ and from $g$ to $h$ through $Q$. Form
\[
 K=P\mathbin{\times_{Y}}Q
 \quad\text{using }d_1\text{ and }e_0,
 \qquad
 t=(s_P,s_Q):Y\longrightarrow K.
\]
Each endpoint map $d_i,e_i$ belongs to $F$, since it is the composite of a path fibration with a projection $Y^2\to Y$, and the latter is a pullback of $Y\to e$. Each endpoint is also in $W$ by two-out-of-three applied to its composite with the constant section. The projection $K\to P$ is the pullback of $e_0\in W$ along $d_1\in F\subset\Fib_W$, so it is in $W$. Its composite with $t$ is $s_P\in W$, and hence $t\in W$.

The map $K\to Y^3$ recording the two outer endpoints and their common middle endpoint is the pullback of
\[
 P\times Q\longrightarrow Y^2\times Y^2
\]
along $(y_0,y_1,y_2)\mapsto(y_0,y_1,y_1,y_2)$. Products and pullbacks of $F$-maps are in $F$, so this map is in $F$. The projection $Y^3\to Y^2$ retaining the outer coordinates is again a pullback of $Y\to e$. Their composite $K\to Y^2$ is therefore in $F$. Thus $Y\xrightarrow{t}K\to Y^2$ is a broad path object. The two given homotopies induce $X\to K$, supplying the homotopy from $f$ to $h$. No exactness or retract closure has entered.

The missing implication for designated paths is now visible: this argument proves $t\in W$, but does not prove $t\in TC$. The graph example shows that the omission cannot be repaired from the weak axioms alone. In that example the broad path $Y\xrightarrow{\Delta}Y^2\xrightarrow{1}Y^2$ makes every pair related, whereas every designated path retains the nontransitive relation just calculated. Consequently footnote~22 has a positive answer for the existential broad convention and a negative answer for the designated convention used in the replacement constructions. Its status is a question of these precise conventions, not an unresolved general homotopy theorem.
\end{dcomment}


% ---- from Der-ChXIII-25-01-04_p021-025.tex
\setcounter{footnote}{23}

\printedpage{19}
% head: A. Grothendieck : Dérivateurs, Ch. XIII, Édition des Universités de Montpellier II et Paris VII
% foot: Version 25 janvier 2004
\msp{page 31}
comme un \textit{adjoint à droite} (pleinement fidèle) du foncteur de localisation
\[
\pi_{\mathrm{d}}\cM_{\mathrm{c}} \longrightarrow (\pi_{\mathrm{d}}\cM_{\mathrm{c}})W_{\mathrm{c}}^{-1} \;\simeq\; \mathrm{Ho}_W,
\]
et $\mathrm{Ho}_W$ se déduit de $\pi_{\mathrm{d}}\cM_{\mathrm{c}}$ par un \textit{calcul de fractions à gauche.} Dualement, si on a exactitude à droite, \textit{i.e.} $TF = F\cap W$ ($\Longleftrightarrow$ $TF = F\cap W_{\mathrm{univ}}$), alors la relation d'équivalence pertinente sur $\mathrm{Hom}_{\cM}(X,Y)$, pour $X$, $Y$ dans $\cM_{\mathrm{cf}}$, est la Q-homotopie à gauche (impliquée par la Q-homotopie à droite), et cette fois la flèche naturelle de $\pi_{\mathrm{g}}\cM_{\mathrm{cf}} = \pi\cM_{\mathrm{cf}} \simeq \mathrm{Ho}_W$ dans $\pi_{\mathrm{d}}\cM_{\mathrm{f}}$ \texttt{[plutôt }$\pi_{\mathtt{g}}\cM_{\mathtt{f}}$\texttt{]},
\[
\underbrace{\pi\cM_{\mathrm{cf}}}_{\simeq\,\mathrm{Ho}_W} \longrightarrow \pi_{\mathrm{d}}\cM_{\mathrm{f}} \qquad \texttt{[plutôt }\pi_{\mathtt{g}}\cM_{\mathtt{f}}\texttt{]}
\]
est un \textit{adjoint à gauche} (pleinement fidèle) du foncteur de localisation
\[
\pi_{\mathrm{d}}\cM_f \longrightarrow \cM_{\mathrm{f}}W_{\mathrm{f}}^{-1} \;\simeq\; \mathrm{Ho}_W \qquad \texttt{[plutôt }\pi_{\mathtt{g}}\cM_f \longrightarrow \cM_{\mathtt{f}}W_{\mathtt{f}}^{-1} \;\simeq\; \mathrm{Ho}_W\texttt{]}, % [sic] \cM_f (italic f, twice) -> \cM_{\mathrm{f}}, elsewhere the index f is upright
\]
et la localisée s'explicite par un calcul de fractions à droite.

\dgrblock{dgrb687}{Exactness and the congruence generated by both homotopies}
\begin{dcomment}{development}{p30}{Exactness and the congruence generated by both homotopies}
Grothendieck has now isolated the hypothesis which allows a weak equivalence between good objects to be inverted by an ordinary arrow. His observation admits a precise completion even before one-sided exactness identifies a preferred homotopy relation.

A chosen cylinder of $X$ is a factorization
\[
 X\amalg X\xrightarrow{(j_0,j_1)}IX\xrightarrow{q}X,
 \qquad (j_0,j_1)\in C,\quad q\in TF.
\]
A chosen path object of $Y$ is a factorization
\[
 Y\xrightarrow{s}PY\xrightarrow{(p_0,p_1)}Y\times Y,
 \qquad s\in TC,\quad(p_0,p_1)\in F.
\]
Left and right Q-homotopies are maps $IX\to Y$ and $X\to PY$ with the prescribed endpoints. We allow all such factorizations in these definitions.

\noindent\textbf{A quotient description under exactness alone.}\ours\quad
Assume the weak Quillen axioms and $W=TF\circ TC$. Let $\equiv$ be the smallest congruence on $\mathcal M_{\mathrm{cf}}$ containing both left and right Q-homotopy. Then
\[
 \mathcal M_{\mathrm{cf}}/{\equiv}\ \simeq\ \mathcal M[W^{-1}].
\]
In particular every morphism between fibrant-cofibrant objects in the localization is represented by an ordinary morphism.

\emph{Proof.} Both kinds of homotopy become equalities after localization: for a cylinder, $qj_0=qj_1=1$ and $q\in W$; for a path object, $p_0s=p_1s=1$ and $s\in W$. It therefore suffices to prove that the quotient inverts every weak equivalence between fibrant-cofibrant objects.

Let $w:X\to Y$ be such a weak equivalence. Exactness gives $w=pi$ with $i:X\to Z$ in $TC$ and $p:Z\to Y$ in $TF$. The intermediate object $Z$ is cofibrant because $TC\subset C$ and $X$ is cofibrant; it is fibrant because $TF\subset F$ and $Y$ is fibrant. Lifting $i$ against $X\to e$ gives $r:Z\to X$ with $ri=1_X$. Lifting $\varnothing\to Y$ against $p$ gives $t:Y\to Z$ with $pt=1_Y$.

To compare $ir$ with $1_Z$, choose a path object $Z\xrightarrow{s}PZ\xrightarrow{v}Z\times Z$. The square with left edge $i$, right edge $v$, top edge $si$ and bottom edge $(ir,1_Z)$ commutes. The lifting axiom $TC\perp F$ supplies a right homotopy from $ir$ to $1_Z$. To compare $tp$ with $1_Z$, choose a cylinder $Z\amalg Z\xrightarrow{j}IZ\xrightarrow{q}Z$. The square with left edge $j$, right edge $p$, top edge $(tp,1_Z)$ and bottom edge $pq$ commutes. The lifting axiom $C\perp TF$ supplies a left homotopy from $tp$ to $1_Z$.

Consequently $i$ and $p$ are invertible in the quotient, with inverses $r$ and $t$, and $w$ has inverse $rt$. The universal property of localization now identifies the quotient with $\mathcal M_{\mathrm{cf}}[W_{\mathrm{cf}}^{-1}]$. The replacement argument in the preceding block identifies that category with $\mathcal M[W^{-1}]$.\hfill$\square$

The congruence formulation deliberately keeps the two homotopies together. Under left exactness a left homotopy implies right homotopy: for a cofibrant source, either cylinder endpoint is in $C\cap W=TC$, so the familiar endpoint lifting square against a path fibration produces the right homotopy. Dually, right exactness converts a right homotopy into a left homotopy. Thus the proposition also proves the manuscript's description by the relation \emph{generated} by right homotopy under left exactness, and its dual. We have checked Cisinski's model-category and derivable-category treatments; their standard statements identify a single homotopy relation under their own hypotheses. The joint-congruence formulation is supplied here for the exact weak axioms on these pages, without claiming that it is new.
\end{dcomment}


Dans le cas biexacte (\textit{i.e.} celui des véritables structures de Quillen), les deux relations de Q-homotopie dans $\cM_{\mathrm{cf}}$ (plus généralement dans $\mathrm{Hom}(X,Y)$, si $X$ dans $\cM_{\mathrm{c}}$ et $Y$ dans $\cM_{\mathrm{f}}$)
\msp{page 32}
coïncident, et deviennent même particulièrement simples. (Car si $f, g : X \rightrightarrows Y$, avec $X$ cofibrant et $Y$ fibrant, alors $f\sim g$ peut se tester sur \textit{n'importe quel} “objet cylindre” pour $X$, ou sur \textit{n'importe quel} “objet chemins” pour $Y$.)

\dgrblock{dgrb688}{Why the localization has a fully faithful adjoint}
\begin{dcomment}{exposition}{p31}{Why the localization has a fully faithful adjoint}
The proposed adjoint is a way to recognize fibrant replacement from a universal property in the homotopy quotient. Assume left exactness, and write $\pi_{\mathrm d}\mathcal M_{\mathrm c}$ for the category of cofibrant objects modulo the equivalence relation generated by right Q-homotopies. Right homotopy is compatible with both compositions: precomposition is immediate, while postcomposition follows by extending a map between the constant paths across $s:Y\to PY$ using $TC\perp F$. Thus taking its generated equivalence relation really defines a category.

Choose $r_X:X\to RX$ in $TC$ with $RX$ fibrant. For any fibrant-cofibrant $Y$, restriction along $r_X$ gives a bijection
\[
 \operatorname{Hom}_{\pi_{\mathrm d}\mathcal M_{\mathrm c}}(RX,Y)
 \ \xrightarrow{\ \sim\ }\
 \operatorname{Hom}_{\pi_{\mathrm d}\mathcal M_{\mathrm c}}(X,Y).
\]
Surjectivity follows by lifting $r_X$ against $Y\to e$. For injectivity, a right homotopy between $ur_X$ and $vr_X$ extends across $r_X$ by lifting against the path fibration $PY\to Y^2$, giving a right homotopy between $u$ and $v$. For a chain of homotopies, choose extensions of all intermediate maps and apply this argument successively. The same proof works for chains with either orientation.

The previous block identifies the homotopy quotient on fibrant-cofibrant objects with $\mathrm{Ho}_W$. Moreover its inclusion in the larger quotient is fully faithful: any relation in the larger quotient gives equality after localization, and on fibrant-cofibrant objects such equality is already the generated right-homotopy relation. The displayed bijection therefore says exactly that the localization
\[
 L:\pi_{\mathrm d}\mathcal M_{\mathrm c}\longrightarrow\mathrm{Ho}_W
\]
is left adjoint to the inclusion of the fibrant-cofibrant homotopy category. The unit is represented by $r_X$. This proves the suggested fully faithful right adjoint without any claim that a chosen lift is unique before homotopy.

It also proves the stated calculus of fractions. Write $S$ for the images of the arrows of $W$ in this quotient; identities and compositions belong to $S$. Precomposition by $s:X\to X'$ in $S$ is bijective on maps into any fibrant-cofibrant object, since the displayed bijection identifies these Hom sets with those in the localization. Given $f:X\to Y$ and $s:X\to X'$ in $S$, choose the denominator $r_Y:Y\to RY$. That bijectivity supplies a map $g:X'\to RY$ with
\[
 gs=r_Yf,
\]
which is the common-denominator condition. For cancellation, suppose $f,g:X\to Y$ and $s:X'\to X$ in $S$ satisfy $fs=gs$. The same bijectivity, now with target $RY$, gives $r_Yf=r_Yg$. Thus the denominator $r_Y$ also supplies the cancellation condition. Every localized arrow consequently has a representative $X\to RY\leftarrow Y$, with the right-hand arrow in $S$; this specifies the manuscript's left-fraction convention.

Applying the whole argument to the opposite category proves the right-exact statement and the fully faithful left adjoint on the quotient of fibrant objects. The editors' changes from $\pi_{\mathrm d}$ to $\pi_{\mathrm g}$ on this page are therefore mathematically required.

For the cofibrant sources and fibrant targets considered here, the biexact case gives a single homotopy relation, computable using any cylinder or path object \citep[Lemma~2.2.12]{Cisinski2019HigherCategories}. Footnote~22 omits the cofibrancy assumption on the source and leaves the path convention implicit. The block \emph{Two path conventions and a nontransitive homotopy} resolves that distinction. For designated $TC$--$F$ paths, raw right homotopy can fail transitivity even when every object is fibrant and cofibrant if the weak structure is not exact; it can also fail in a perfect structure for a noncofibrant source. Both failures come from the graph example of \citet[Example~3.6]{KurzRosicky2005}. In contrast, the existential relation obtained by allowing all $W$--$F$ path objects is an equivalence relation for a fibrant target and an arbitrary source, by direct concatenation. The generated relation used in the adjoint argument above is therefore a deliberate and sufficient formulation.
\end{dcomment}


Dans le cas où la structure de Quillen faible est exacte à gauche, \textit{i.e.} où la relation d'équivalence pertinente dans $\mathrm{Fl}\,\cM_{\mathrm{cf}}$ est la Q-homotopie à droite, alors l'axiome d'exactitude est vérifié pour les carrés $W$-cartésiens, mais je doute qu'il en soit de même pour les carrés $W$-cocartésiens (\footnote{Il semble pourtant que c'est le cas, d'après \textsc{Baues}, cf. annotation marginale correspondante page 30.}). (Si on ne fait une hypothèse supplémentaire, p. ex. l'exactitude à droite de la structure de Quillen.) Si on n'a que l'exactitude (mais ni à gauche ni à droite), il semble qu'on ne peut en conclure l'exactitude ni des carrés $W$-cartésiens, ni des carrés $W$-cocartésiens, donc pas celle des suites de cosuspension (loop sequences), ni des suites de suspension. On pourrait donc penser que la condition d'exactitude, à elle seule, ne sert pas à grand
\msp{page 33}
chose pour le formalisme homotopique dans les catégories de fractions $\mathrm{Ho}_W$ etc. Mais on a pourtant ceci : soit $\Psi = {}$\begin{tikzcd}[row sep=tiny, column sep=small, baseline=(\tikzcdmatrixname-2-1.base)]
 & 1 \\
0 \arrow[ru] \arrow[rd] & \\
 & 2
\end{tikzcd}, et supposons que $(\cM,W)$ \textit{et} $(\cM(\Psi),W(\Psi))$ \texttt{[où pour toute petite catégorie }$I$\texttt{, }$\cM(I)$\texttt{ désigne la catégorie des préfaisceaux sur }$I$\texttt{ à valeurs dans }$\cM$\texttt{, et }$W(I)$\texttt{ est la partie de }\texttt{Fl}$(\cM(I))$\texttt{ formée des morphismes}
\printedpage{20}
% head: A. Grothendieck : Dérivateurs, Ch. XIII, Édition des Universités de Montpellier II et Paris VII
% foot: Version 25 janvier 2004
\texttt{de préfaisceaux qui sont dans }$W$\texttt{ "argument par argument"]} admettent une structure de Quillen faible exacte. Alors dans $\cM$ les carrés $W$-cartésiens sont $W$-exactes, donc on a dans $\cM$ des suites exactes longues de cosuspension (exact loop sequences) (\footnote{il semble que je me suis canulé ici, et je doute que le résultat évoqué soit exact.}). (L'hypothèse sur $(\cM(\Psi),W(\Psi))$ sert uniquement à assurer que toute flèche de $\mathrm{Ho}_{W(\Psi)} = \cM(\Psi)W(\Psi)^{-1}$ est isomorphe à une flèche provenant d'une flèche de $\cM(\Psi)$ lui-même. Mais c'est là une propriété apparemment très délicate, et je ne connais d'autre façon de l'obtenir que via l'existence d'une structure de Quillen faible exacte.) Dualement, si $(\cM,W)$ et $(\cM(\Phi),W(\Phi))$ admettent des structures de Quillen faibles exactes, alors les carrés $W$-cocartésiens dans $\cM$ sont $W$-exacts, et on trouve (si $\cM$ est ponctué) des suites exactes longues de suspension. % [sic] W-exactes -> W-exacts, agreement with "carrés" (cf. "W-exacts" below)

\dgrblock{dgrb689}{Exactness permits completion of the trivial classes}
\begin{dcomment}{development}{p32}{Exactness permits completion of the trivial classes}
Grothendieck's hesitation about exactness can be resolved by changing the designated trivial classes while keeping the weak equivalences, cofibrations and fibrations. This is a different operation from asserting that the original designated classes were already complete. It is therefore compatible with the counterexample to footnote~20.

\noindent\textbf{Completion under exactness.}\ours\quad
Suppose $(C,TC,F,TF)$ is an exact weak Quillen structure. Then
\[
 \bigl(C,\ C\cap W,\ F,\ F\cap W\bigr)
\]
is a biexact Quillen structure, with the same $C,F,W$. No retract-closure hypothesis on these classes is needed for this assertion. If all four original lifting classes were closed, the completion makes no change.

\emph{Proof.} The only new lifting problems involve $C\cap W$ against $F$, and $C$ against $F\cap W$. Let $j\in C\cap W$. Exactness gives $j=qi$ with $i\in TC$ and $q\in TF$. Since $j\in C$ and $q\in TF$, the lifting axiom exhibits $j$ as a retract of $i$ in the arrow category. Lifting properties are preserved by retracts, so $j$ has the left lifting property against every map in $F$. Dually, if $p\in F\cap W$, factor $p=qi$ with $i\in TC$ and $q\in TF$. Lifting $i$ against $p$ exhibits $p$ as a retract of $q$, so it has the right lifting property against $C$.

The original factorizations still serve, because $TC\subset C\cap W$ and $TF\subset F\cap W$. The classes $C$ and $F$ have not changed. The classes $C\cap W$ and $F\cap W$ have the required pushout or pullback stability by $\text{M}''$ and the universal-equivalence criterion of manuscript pp.~10--11; composition and identities follow from two-out-of-three. Thus every axiom holds, and both exactness equalities hold by definition. When $TC$ and $TF$ are retract closed, the two retract arguments also show $C\cap W\subset TC$ and $F\cap W\subset TF$, giving equality with the original classes.\hfill$\square$

This completion is a possibly nonclosed model structure in Quillen's original terminology. It suffices for the ordinary homotopy theory of cofibrant and fibrant objects: the standard endpoint arguments now use $C\cap W\perp F$ and $C\perp F\cap W$. In particular both homotopies classify the localized maps from a cofibrant object to a fibrant object, and any good chosen cylinder or path object from the original weak structure remains available. The existence of different designated classes therefore does not obstruct the usual homotopy calculation once exactness holds. This also explains the applicability of the nonclosed Quillen case in \citet[Example~5.2]{Cisinski2010}.

We checked the quoted Quillen and derivable-category frameworks and did not locate this completion stated using the exact four-class axioms of the manuscript. The proof is a direct application of the standard retract argument; its purpose is to settle the hypotheses on these pages, not to claim a new model-category theorem.
\end{dcomment}


\msp{page 34}
Ceci montre que si, pour une catégorie de modèles $(\cM,W)$ donnée, on s'intéresse à “faire de l'algèbre homotopique” pour tous les couples $(\cM(I),W(I))$ à la fois, avec $I$ dans Cat, il suffit (si on veut avoir une bonne théorie des fibres et cofibres $W$-homotopiques dans ces catégories, y inclus la propriété de $W$-exactitude pour les carrés $W$-cartésiens et pour les carrés $W$-cocartésiens) de disposer pour chaque $(\cM(I),W(I))$ d'une structure de Quillen faible exacte sans plus.

\dgrblock{dgrb690}{Exactness of squares after completing the classes}
\begin{dcomment}{exposition}{p33}{Exactness of squares after completing the classes}
Grothendieck introduces $\mathcal M(\Psi)$ to recover relations among maps while representatives remain available. Exactness in that diagram category indeed says that every localized arrow is isomorphic, in the arrow category, to an actual natural transformation: replace its source and target by fibrant-cofibrant objects and apply the quotient description above. His footnote correctly warns that this representability argument has not, by itself, established compatibility with the chosen homotopy pullback.

There is a more direct repair. The completion just proved turns an exact weak structure on $\mathcal M$ into a possibly nonclosed proper Quillen structure with the same $W,C,F$. Both homotopy-cartesian and homotopy-cocartesian exactness follow from that structure. For this conclusion the additional exact structure on $\mathcal M(\Psi)$ is unnecessary. Thus the asserted conclusion on p.~33 is valid, while the proposed route through arrow representability leaves a step unproved.

Here is the local lifting calculation behind the cartesian assertion. After taking the usual fibrant representatives, write the homotopy pullback as the ordinary pullback $P=E\times_SY$, where $p:E\to S$ is a fibration and $E,Y,S$ are fibrant. Let $T$ be cofibrant. Suppose maps $a:T\to E$ and $b:T\to Y$ have homotopic composites to $S$. Represent the homotopy by $H:IT\to S$, with endpoints $pa$ and $qb$. The endpoint inclusion $j_0:T\to IT$ is in $C\cap W$: it is a cofibration because $T$ is cofibrant, and a weak equivalence because the cylinder projection is one. The completed lifting property against $p$ gives a lift $\widetilde H:IT\to E$. Its other endpoint $a'=\widetilde H j_1$ is homotopic to $a$ and satisfies $pa'=qb$. The pair $(a',b)$ therefore gives a map $T\to P$. This proves the required surjectivity onto pairs of homotopy classes with equal composites. The dual argument gives cocartesian exactness and the corresponding suspension sequences in the pointed case.

The geometry of the calculation is visible in the cospan $*\to S^1\leftarrow *$. Its homotopy pullback is $\Omega S^1$, not the one-point strict pullback. Lifting a based loop to $\mathbb R$ with initial point $0$ assigns its endpoint in $\mathbb Z$; two loops are homotopic relative endpoints precisely when those integers agree. Hence $\pi_0\Omega S^1=\mathbb Z$. The path between the two composites is part of the homotopy-pullback data. The lifting argument preserves that datum while making the square strictly commute. This is the compatibility that mere representability of an arrow does not explain.
\end{dcomment}


Ceci nous montre donc l'intérêt de développer des critères, pour une $(\cM,W)$, assurant qu'il en est bien ainsi. À ce sujet, je note que je ne sais comment vérifier l'exactitude d'une structure de Quillen \textit{que si} je dispose de l'exactitude à gauche, ou à droite. C'est dire que je dois dégager en fait des conditions sur $(\cM,W)$ qui assurent (disons) que les $(\cM(I),W(I))$ admettent des structures de Quillen faibles exactes à gauche.

\msp{page 35}
Les structures de Quillen faibles pour les $(\cM(I),W(I))$ seront déduites d'une structure de Quillen faible pour $(\cM,W)$, de la façon suivante : elles seront toutes \textit{“closes à droite”,} \textit{i.e.}
\[
F(I) \;=\; (TC(I))_*, \qquad TF(I) \;=\; (C(I))_*, \tag*{(1)}
\]
en particulier (pour $I=e$)
\[
F \;=\; (TC)_*, \qquad TF \;=\; C_*,
\]
ce qui implique déjà (puisque $C\subset \mathrm{Cof}_W$, $TC\subset W^{\mathrm{univ}}$)
\[
(\mathrm{Cof}_W)_* \;\subset\; W, \qquad (W^{\mathrm{univ}})_* \;\subset\; \mathrm{Fib}_W \tag*{(2)}
\]
(d'où $(\mathrm{Cof}_W)_*\subset W_{\mathrm{univ}}$). En vertu de (1), pour définir les structures de Quillen faibles pour les $(\cM(I),W(I))$, il suffit de définir $C(I)$, $TC(I)$, et même, comme on veut l'exactitude à gauche, \textit{i.e.}
\[
TC(I) \;=\; C(I)\cap W(I) \tag*{(3)}
\]
\printedpage{21}
% head: A. Grothendieck : Dérivateurs, Ch. XIII, Édition des Universités de Montpellier II et Paris VII
% foot: Version 25 janvier 2004
(en particulier
\[
TC \;=\; C\cap W \quad ), \tag*{(4)}
\]
il reste à définir $C$. On définira $C(I)$ “argument par argument”, \textit{i.e.} pour $f\in\mathrm{Fl}(\cM(I))$
\[
f\in C(I) \qquad\Longleftrightarrow\qquad \alpha_s^*(f)\in C \ \text{ pour tout } s\in\mathrm{Ob}\,I, \tag*{(5)}
\]
où
\[
\alpha_s : e\longrightarrow I \qquad \texttt{[est]}\ \text{défini par } \alpha_s(e)=s.
\]
Il en résulte d'ailleurs aussitôt, par (3),
\[
f\in TC(I) \qquad\Longleftrightarrow\qquad \alpha_s^*(f)\in TC \ \text{ pour tout } s\in\mathrm{Ob}\,I. \tag*{(6)}
\]

\msp{page 36}
On suppose donc que $(\cM,W)$ admet une structure de Quillen faible exacte à gauche et close à droite $(C,TC,F,TF)$, définie en termes de
\[
C\;\subset\;\mathrm{Fl}(\cM),
\]
par
\[
TC \;=\; C\cap W, \qquad F \;=\; (TC)_*, \qquad TF = C_*.
\]
Pour que cela donne bel et bien une structure de Quillen faible (forcément exacte à gauche et close à droite), il faut simplement vérifier les conditions suivantes.

\dgrblock{dgrb691}{The injective choice of diagram structure}
\begin{dcomment}{exposition}{p35}{The injective choice of diagram structure}
Grothendieck now makes a choice which determines the variance of the whole construction. Cofibrations and weak equivalences will be tested at each object of the indexing category. Fibrations will be characterized by lifting against these objectwise classes. In modern terminology this is the injective choice. The asymmetry is productive: every restriction of diagrams preserves the classes specified objectwise, so its right adjoint becomes a candidate for a right Quillen functor.

Keep the variance explicit: $\mathcal M(I)=\operatorname{Fun}(I^{\mathrm{op}},\mathcal M)$. If $C(I)$ is the objectwise class associated with $C$, then
\[
 TC(I)=C(I)\cap W(I)=(C\cap W)(I)
\]
is an identity of classes, not an existence theorem. Once $F(I)=TC(I)_*$ and $TF(I)=C(I)_*$ are defined, lifting is automatic and the right classes have their required composition and pullback stability. Objectwise pushouts show the stability of $C(I)$; the corresponding statement for $TC(I)$ uses the universal weak-equivalence property already obtained in the base category.

The remaining two tasks have different characters. The inclusions $TF(I)\subset W(I)$ and $F(I)\subset\Fib_{W(I)}$ can be checked by evaluation and its adjoints; pp.~37--39 give that argument. Factorization in the diagram category requires a construction of coherent fillers and cannot be inferred merely by factoring every component. Even functorial componentwise factorizations do not make both lifting classes objectwise. The walking-arrow example below shows the distinction in the smallest nontrivial indexing category.
\end{dcomment}


\begin{itemize}[leftmargin=3em]
\item[\circled{i}] Factorisation pour $(C, TF=C_*)$ et $(TC=C\cap W, F=(TC)_*)$ – on reviendra sur ce point.
\item[\circled{ii}] Stabilité de $C$ par composition et cochangement de base, et $C$ contient les isomorphismes (d'où résultera la stabilité similaire de $TC = C\cap W$, compte tenu que $C\subset\mathrm{Cof}_W$) (NB Les stabilités idoines pour $F$, $TF$ sont triviales par définition $F=(TC)_*$, $TF=C_*$.)
\item[\circled{iii}] $C\subset\mathrm{Cof}_W$, $\underbrace{C_*}_{=\,TF}\subset W$, $(C\cap W)_*\subset\mathrm{Fib}_W$.
\end{itemize}

Pour vérifier, de même, que $(C(I),TC(I),F(I),TF(I))$ soit une structure de Quillen faible pour $W(I)$, il faut vérifier les conditions similaires à \circled{i}\,\circled{ii}\,\circled{iii} . Or (\textit{ii}) pour $C(I)$ résulte aussitôt de la condition (\textit{ii}) pour $C$. Il reste à examiner (\textit{iii}) et (\textit{i}).

\printedpage{22}
% head: A. Grothendieck : Dérivateurs, Ch. XIII, Édition des Universités de Montpellier II et Paris VII
% foot: Version 25 janvier 2004
\msp{page 37}
Les inclusions de (\textit{iii}).

L'inclusion
\[
C\;\subset\;\mathrm{Cof}_W
\]
implique que toute $f\in C(I)$ est dans $\mathrm{Cof}_W$ argument par argument. Cela implique que $f$ est dans $\mathrm{Cof}_{W(I)}$, comme on voit aussitôt sur la définition – je l'explicite en le

\textbf{Lemme 1 :} \textit{Soit $f : X_\bullet\longrightarrow Y_\bullet$ une flèche de $\cM(I)$. Pour que $f$ soit dans} $\mathrm{Cof}_{W(I)}$ (resp. \textit{dans} $\mathrm{Fib}_{W(I)}$), \textit{il} suffit \textit{que pour tout $s$ dans} $\mathrm{Ob}\,I$, \textit{$\alpha_s^*(f)$ soit dans} $\mathrm{Cof}_W$ (resp. \textit{dans} $\mathrm{Fib}_W$)\textit{. Itou (pour mémoire) pour une condition pour que $f\in W(I)$ – là il faut et il suffit que $\alpha_s^*(f)\in W\ \forall\, s\in\mathrm{Ob}\,I$.}

Nous allons appliquer ce même lemme pour établir les inclusions
\[
C(I)_* \;\subset\; W(I), \qquad (C(I)\cap W(I))_* \;\subset\; \mathrm{Fib}_{W(I)}.
\]
Donc il suffit de voir que
\[
\begin{aligned}
\alpha_s^*(C(I)_*) \;&\subset\; C_* &&\qquad (\subset\; W)\\
\alpha_s^*((\underbrace{C(I)\cap W(I)}_{(C\cap W)(I)})_*) \;&\subset\; (C\cap W)_* &&\qquad (\subset\; \mathrm{Fib}_W),
\end{aligned}
\]
et plus généralement
\[
\alpha_s^*(\Phi(I)_*) \;\subset\; \Phi_* \tag*{(7)}
\]
pour toute ( ?) partie $\Phi$ de $\mathrm{Fl}(\cM)$, définissant donc $\Phi(I)\subset\mathrm{Fl}(\cM(I))$ terme à terme.

\msp{page 38}
En explicitant la chose, on trouve la condition
\[
\alpha_{s\,!}(\Phi)\subset\Phi(I) \qquad \text{pour tout } s\in\mathrm{Ob}\,I,
\]
pour que ça marche, \textit{i.e.} on veut que
\[
i\in\Phi \qquad\Longrightarrow\qquad \alpha_{s\,!}(i)\in\Phi(I), \quad \textit{i.e.}\ \ \alpha_t^*(\alpha_{s\,!}(i))\in\Phi\ \ \forall\, t\in\mathrm{Ob}\,I.
\]
Or le foncteur $\alpha_t^*\alpha_{s\,!}$ s'explicite comme
\[
A \longmapsto \underbrace{\mathrm{Hom}_I(t,s)\times A}_{\substack{\text{somme de } \mathrm{Hom}_I(t,s)\\ \text{copies de } A \text{ dans } \cM}},
\]
donc il faut supposer $\Phi$ \textit{stable par sommes directes} (du type envisagé $\mathrm{Hom}_I(t,s)\times A$). C'est bien le cas pour $\Phi = C$ et $\Phi = C\cap W$ et pour des sommes directes \textit{finies,} grâce à (\textit{ii}). On trouve donc le

\printedpage{23}
% head: A. Grothendieck : Dérivateurs, Ch. XIII, Édition des Universités de Montpellier II et Paris VII
% foot: Version 25 janvier 2004
\textbf{Lemme 2.}

\begin{itemize}[leftmargin=3em]
\item[a)] \textit{Soit $I$ dans} Cat \textit{telle que les sommes directes de type} $\mathrm{Hom}_I(t,s)$ \textit{existent dans $\cM$. (P. ex. on suppose que $I$ est} finie, \textit{plus généralement que les} $\mathrm{Hom}_I(t,s)$ \textit{sont des ensembles finis ; ou que les “petites” sommes directes existent dans $\cM$.) Soit $\Phi$ une partie de} $\mathrm{Fl}(\cM)$\textit{, stable par ce type de sommes directes, et soit} $\Phi(I)\subset\mathrm{Fl}(\cM(I))$ \textit{la partie correspondante. Alors pour tout $s\in\mathrm{Ob}\,I$, on a}

\msp{page 39}
\[
\alpha_s^*(\Phi(I)_*) \;\subset\; \Phi_*.
\]
\item[b)] \textit{Supposons que les} $\mathrm{Hom}_I(t,s)$ \textit{soient finis (p. ex. $I$ finie). Alors l'hypothèse de} a) \textit{s'applique pour $\Phi=C$ et pour $\Phi=C{\cap}W$, donc la condition} (iii) \textit{p.} 36 \textit{est satisfaite pour la structure $(W(I),C(I),TC(I),F(I),TF(I))$ dans $\cM(I)$. Même conclusion sans hypothèse de finitude sur les} $\mathrm{Hom}_I(t,s)$\textit{, quand on suppose $\cM$ stable par “petites” sommes directes, et $C$ et $W$ stables également par de telles sommes directes.}
\end{itemize}

\dgrblock{dgrb692}{Evaluation transports the lifting problem}
\begin{dcomment}{exposition}{p38}{Evaluation transports the lifting problem}
The copower formula is the essential calculation in the two lemmas. Let $s$ be an object of $I$. The left adjoint to evaluation at $s$ is
\[
 (\alpha_{s!}A)_t=\coprod_{\operatorname{Hom}_I(t,s)}A.
\]
For an arrow $t'\to t$, precomposition sends $\operatorname{Hom}_I(t,s)$ to $\operatorname{Hom}_I(t',s)$ and defines the presheaf restriction map. A natural transformation $\alpha_{s!}A\to X$ is uniquely determined by the map from the summand indexed by $1_s$ to $X_s$; all other summands are obtained by restriction. This proves the adjunction and fixes the direction of the Hom set.

Suppose $\Phi$ is stable under the copowers appearing in this formula. Then $\alpha_{s!}(i)\in\Phi(I)$ for every $i\in\Phi$. For $p\in\Phi(I)_*$, a lifting square from $i$ to $p_s$ transposes to a square from $\alpha_{s!}i$ to $p$. Its diagonal transposes back to the required lift. Hence
\[
 \alpha_s^*(\Phi(I)_*)\subset\Phi_*.
\]
The statement on p.~37 for an arbitrary class, accompanied by a question mark, therefore needs exactly the additional closure used in Lemma~2. We are not assuming that evaluation itself preserves right orthogonals without an adjoint test.

Finite copowers require no new hypothesis for $C$ or $TC$. For maps $i:A\to B$ and $j:A'\to B'$, the coproduct map factors as
\[
 A\amalg A'\longrightarrow B\amalg A'\longrightarrow B\amalg B',
\]
where each factor is a pushout of one of the original maps. The empty copower is an identity on the initial object. Induction proves the finite statement. For arbitrary copowers one must either assume their stability or obtain it from transfinite composition by adjoining the summands in a well-order. The required colimits must exist.

Finally, finite limits and colimits of presheaves are computed objectwise. A double pushout diagram testing a $W(I)$-cofibration evaluates to the corresponding test in $\mathcal M$, and an objectwise weak equivalence remains one after all these tests. This proves Lemma~1 and its dual. Combining it with the lifting argument proves all the inclusions on p.~36 under the finite-Hom hypothesis, while deliberately leaving factorization for the next stage.
\end{dcomment}


\textbf{NB} Vu le scholie plus haut, il est parfaitement raisonnable de se borner au besoin aux catégories $I$ \textit{finies} – vu que la catégorie de ces $I$ est stable par les produits $I\longmapsto I\times\Psi$, $I\longmapsto I\times\Phi$.

Il reste, pour être heureux, à vérifier la condition de factorisation \circled{i} pour $(C,TF)$, $(TC,F)$, mais aussi et surtout pour les $(C(I),TF(I),TC(I),F(I))$. C'est dire que de partir même d'une structure de Quillen close (pays de cocagne)
\msp{page 40}
ne nous avance guère à lui tout seul, faute de savoir que l'axiome de factorisation se conserve, en passant de $(\cM,W,\ldots)$ à $(\cM(I),W(I),\ldots)$.

\dgrblock{dgrb693}{A complete calculation for the walking arrow}
\begin{dcomment}{examples}{p39}{A complete calculation for the walking arrow}
The formulas become especially transparent when $I=[1]$, with its arrow $0\to1$. A presheaf is a map $X_1\to X_0$. Let $\Phi$ contain identities and be stable under finite coproducts, and put $R=\Phi_*$. Write $x:X_1\to X_0$ and $y:Y_1\to Y_0$ for the structural maps. For a natural transformation $p:X\to Y$, consider the matching map
\[
 m_p:X_1\longrightarrow Y_1\times_{Y_0}X_0.
\]
Then $p$ has the right lifting property against every objectwise $\Phi$-map if and only if
\[
 p_0:X_0\to Y_0\ \in R,
 \qquad m_p\in R.
\]

To prove sufficiency, take a square from an objectwise $\Phi$-map $A\to B$ to $p$. First lift its component at $0$ against $p_0$, obtaining $B_0\to X_0$. The component at $1$, together with this already chosen map composed with $B_1\to B_0$, gives a lifting square against $m_p$. A lift provides $B_1\to X_1$ and enforces precisely the naturality relation. Thus the two component lifts fit together.

For necessity, evaluation at $0$ and its left adjoint show $p_0\in R$. To test $m_p$ against $i:A\to B$ in $\Phi$, write a square as maps $a:A\to X_1$, $b_1:B\to Y_1$, and $b_0:B\to X_0$ with
\[
 p_1a=b_1i,\qquad xa=b_0i,\qquad p_0b_0=yb_1.
\]
Use the map of arrow diagrams
\[
 (A\xrightarrow{i}B)\longrightarrow(B\xrightarrow{1}B),
\]
whose components are $i$ at $1$ and $1_B$ at $0$. This is objectwise in $\Phi$. A lift of its induced square against $p$ is exactly a solution to the original square against $m_p$. This proves necessity.

Apply the calculation in sets with $\Phi$ the injections and $R$ the surjections, assuming choice. Let $E_1=\{a\}$, $E_0=\{0,1\}$, and let $E_1\to E_0$ send $a$ to $0$. The transformation from $E$ to the terminal arrow diagram is surjective at both objects. Its matching map is the non-surjective map $\{a\}\to\{0,1\}$, so it is not in $\Phi(I)_*$. Explicitly, a proposed lift choosing the element $1$ at vertex $0$ cannot be extended to vertex $1$. This is a failure of naturality, despite the existence of lifts at each object separately.

For the injective model structure, taking $\Phi=C$ gives the characterization of trivial fibrations, and taking $\Phi=TC$ gives the characterization of fibrations. It follows in particular that an injectively fibrant arrow has fibrant target and a fibration as its structural arrow. The general admissible-edge result later in the manuscript is an extension of this relative lifting calculation, rather than a consequence of pointwise fibrancy alone.
\end{dcomment}


C'est ici que le théorème de factorisation du no.~1 doit s'avérer providentiel. Mais pour pouvoir l'appliquer aux couples $(C(I),C(I)_*)$, $(TC(I),TC(I)_*)$, il faut (en plus de l'accessibilité de $\cM$, ce qui est une condition fort anodine) des conditions de stabilité très fortes pour $C(I)$, $TC(I)$ (\footnote{Plus les \textit{petits} systèmes de générateurs pour $C$, $TC$.}). À vrai dire, il reste la stabilité de $C$ et de $C\cap W$ par limites inductives ordinales (ce qui impliquera la même stabilité pour les $C(I)$, $(C{\cap}W)(I) = C(I){\cap}W(I)$).

À dire vrai, je vois peu de chances de trouver une structure de Quillen faible $(C,TC,F,TF)$ associé à $W$, et satisfaisant toutes les conditions dites, y compris celle de stabilité de $C$ et de $C\cap W$ par $\varinjlim$ ordinales, sans supposer que $W$ elle-même ne soit déjà stable par (petites) limites inductives filtrantes. % [sic] associé -> associée, agreement with "structure"

\msp{page 41}
On peut songer à utiliser une factorisation \textit{fonctorielle} des flèches dans $\cM$ en
\[
\begin{aligned}
f = pi \quad \text{avec}\quad & i\in TC, && p\in F\\
\text{ou}\quad & i\in C, && p\in TF,
\end{aligned}
\]

\dgrblock{dgrb694}{The factorization problem remains a coherence problem}
\begin{dcomment}{later}{p40}{The factorization problem remains a coherence problem}
The final paragraph should be read as a choice of a sufficient strategy, not as a necessity theorem about weak equivalences. Transfinite closure of cofibrations and trivial cofibrations is part of the small-object method; closure of every weak equivalence under every filtered colimit is a convenient, stronger condition from which Grothendieck hopes to obtain it. Nothing on this page proves that such closure of $W$ is necessary for diagram model structures.

The later injective existence theorem in the combinatorial setting answers the construction problem with explicit generation hypotheses; its precise application belongs to the discussion of manuscript pp.~42--44. There is also a different route through weaker diagram-compatible structures: \citet[Theorem~6.22]{Cisinski2010} proves stability under presheaf categories for homotopically complete model categories in his sense; the theorem also covers homotopically complete and cocomplete weak model categories. These are two ways to obtain the coherence which the walking-arrow calculation detects. Neither licenses treating the pointwise right class as the required right lifting class.
\end{dcomment}


% ---- from Der-ChXIII-25-01-04_p026-030.tex
\printedpage{24}
% head: A. Grothendieck : Dérivateurs, Ch. XIII, Édition des Universités de Montpellier II et Paris VII
pour déduire une factorisation similaire pour les flèches $f$ de $\M(I)$, en
\[
\begin{aligned}
f = pi \quad \text{avec}\quad & i \in TC(I), && p \in F(I)\\
\text{ou}\quad & i \in C(I), && p \in TF(I),
\end{aligned}
\]
où $C(I)$, $TC(I)$, $F(I)$, $TF(I)$ sont ici définis via $C$, $TC$, $F$, $TF$ “argument par argument”. Mais l'ennui, c'est que en procédant de cette façon naïve, on n'aura plus nécessairement $C(I) \longleftrightarrow TF(I)$, $TC(I) \longleftrightarrow F(I)$. Il faut considérablement rapetisser $TF(I)$ et $F(I)$, pour avoir une propriété RLP par rapport à $C(I)$ resp. $TC(I)$.

Donc je ne vois guère d'autre moyen de s'en tirer, dans ce contexte général, qu'en supposant la stabilité de $W$ par $\limind$ filtrantes (ce qui n'est après tout pas très méchant).

\noindent\msp{page 42}

\dgrblock{dgrb695}{Natural factorizations and natural fillers}
\begin{dcomment}{exposition}{p41}{Natural factorizations and natural fillers}
Grothendieck has reached the obstruction which makes a theory of diagrams more than a repetition of the theory of objects. A functorial factorization in $\mathcal M$ does produce a factorization of every natural transformation. Its middle objects and their transition maps fit together because the original factorization is functorial on the arrow category. But a choice of a filler for each component of a lifting problem need not commute with the transition maps. The missing datum is a compatible family of fillers, and the factorization functor supplies no such family.

The walking-arrow calculation in the commentary at manuscript pp.~38--39 makes this failure completely explicit. For diagrams $X_1\to X_0$ and $Y_1\to Y_0$, right lifting against all objectwise maps in a left class $L$ requires both
\[
 X_0\longrightarrow Y_0\in R,
 \qquad X_1\longrightarrow Y_1\mathbin{\times}_{Y_0}X_0\in R,
 \qquad R=L_*.
\]
The second condition remembers the compatibility of the lift with the arrow. Requiring only $X_1\to Y_1\in R$ forgets it. This is the first matching-object condition of the later Reedy formalism; the general condition appears in \citet[Definition~5.2.2 and Theorem~5.2.5]{Hovey1999}.

There are consequently two different ways to organize the model structure of a diagram category. In the injective structure the cofibrations and weak equivalences are objectwise, and fibrations are defined by a global right lifting property. In the projective structure the fibrations and weak equivalences are objectwise, and the cofibrations carry the compatibility condition. The manuscript is pursuing the injective structure. Its decision to keep $C(I)$ and $TC(I)$ objectwise therefore forces the replacement of the naive classes of fibrations by the smaller orthogonal classes $(TC(I))_*$ and $C(I)_*$.

This choice has a purpose beyond obtaining factorizations. Restriction along any functor preserves the two objectwise left classes. Its right adjoint will therefore preserve the two right classes, which is exactly the input needed to construct homotopy right Kan extensions. Thus the apparent asymmetry of the construction already prepares the three variances of the next section.
\end{dcomment}


Bon gré, mal gré, on revient à la situation envisagée dans le paragraphe 2. On va réénoncer l'essentiel de ce qui a été pressenti alors :

\textbf{Théorème 2} (encore un peu conjectural). {\itshape Soit $(\M,W)$ une catégorie de modèles, avec $W$ (comme toujours) faiblement saturée, mais de plus on suppose $\M$ stable par petites $\limind$ et par $\limproj$ finies, et $W$ stable par $\limind$ filtrantes, enfin $\M$ accessible. Considérons les conditions suivantes :

\begin{enumerate}
\item[\textup{a)}] Il existe une structure de Quillen faible pour $W$, exacte à gauche, close à droite, $(C,TC,F,TF)$.
\item[\textup{a$'$)}] Comme \textup{a)}, mais avec de plus $C$, $TC$ clos à gauche (\textup{i.e.} on a une structure de Quillen faible close), et $C$, $TC$ de la forme $\widetilde{C_0}$, $\widetilde{TC_0}$, où $C_0$ et $TC_0$ sont petits.
\item[\textup{a$''$)}] Il existe $C_0 \subset \Cof_W$, $TC_0 \subset W^{\univ}$, parties petites, telles que $C_{0*} \subset W$, $(TC_0)_* \subset \Fib_W$.
\item[\textup{b)}] On a
\[
(\Cof_W)_* \;\subset\; W_{\univ}, \qquad\quad (W^{\univ})_* \;\subset\; \Fib_W
\]
(d'où résulte, $\Cof_W$ et $W^{\univ}$ étant clos à gauche, que
\[
(W_{\univ})^* \subset \Cof_W, \qquad\quad (\Fib_W)^* \subset W^{\univ} \qquad\quad \text{\itshape ).}
\]
\end{enumerate}

\noindent\msp{page 43}

Ces conditions impliquent que pour toute $I$ dans \textup{Cat}, il existe une structure de Quillen faible exacte à gauche et close pour $W(I)$ dans $\M(I)$. Une telle structure se définit, en termes d'une structure similaire $C$, $TC = C \cap W$, $F = (TC)_*$, $TF = C_*$ pour $W$, par}
\[
C(I) \;=\; \Bigl\{ f \in \Fl(\M(I)) \Bigm| \alpha_s^*(f) \in C \ \ \forall\, s \in \Ob I \Bigr\},
\]

\printedpage{25}
% head: A. Grothendieck : Dérivateurs, Ch. XIII, Édition des Universités de Montpellier II et Paris VII
\noindent\textit{puis}
\[
\begin{aligned}
TC(I) \;&=\; C(I)\cap W(I) \;=\; \Bigl\{ f \in \Fl(\M(I)) \Bigm| \alpha_s^*(f) \in TC \ \ \forall\, s \in \Ob I \Bigr\}\\
F(I) \;&=\; (TC(I))_*\\
TF(I) \;&=\; C(I)_*\,.
\end{aligned}
\]

\textbf{Corollaire-Scholie.} \textit{Ces conditions impliquent que pour tout $I$ dans} Cat, \textit{$(\M(I),W(I))$ donne lieu à une bonne théorie des fibres et cofibres $W(I)$-homotopiques, avec carrés $W(I)$-cartésiens $W(I)$-exacts, et carrés $W(I)$-cocartésiens $W(I)$-exacts – donc avec des suites} exactes \textit{longues de cosuspension et de suspension, dans le cas où $\M$, donc aussi $\M(I)$, est une catégorie avec objet neutre.}

\noindent\msp{page 44}

\dgrblock{dgrb696}{The injective existence theorem and its precise hypothesis}
\begin{dcomment}{later}{p43}{The injective existence theorem and its precise hypothesis}
Grothendieck wants the factorization theorem to work uniformly after adding any small indexing category. One part of this proposal is now a standard existence theorem: a combinatorial model category has injective and projective model structures on each small diagram category. In the terminology of the present pages, a closed structure exact on the left is already an ordinary closed model structure, and condition $a'$ supplies the two sets of generators. Since an accessible cocomplete category is locally presentable, $a'$ places us in the combinatorial situation.

For the injective construction, the set of generators on the diagram category is the substantive point. One cannot simply use the free diagrams on the generators at each object: those give the projective left class. A theorem of \citet[Corollary~3.4, author preprint of 21 November 2013]{MakkaiRosicky2014Cellular} instead says that the pointwise left class associated with a set-generated, retract-closed cellular structure on a locally presentable category is again set-generated. Apply it separately to $C$ and $TC$. The two resulting small-object factorizations have left classes $C(I)$ and $TC(I)$ and right classes $C(I)_*$ and $(TC(I))_*$.

Here is why the compatibility of these two factorizations survives. Evaluation takes either right class into its corresponding right class in $\mathcal M$: its left adjoint is the copower by the representable presheaf, and copowers preserve $C$ and $TC$. Consequently $C(I)_*\subset W(I)$, while an element of $(TC(I))_*$ which belongs to $W(I)$ is objectwise in $TF$. To obtain the reverse lifting assertion without confusing componentwise and natural fillers, factor such a map $f$ using $(C(I),C(I)_*)$, say $f=pi$. The maps $p$ and $f$ are weak equivalences, so $i\in C(I)\cap W(I)=TC(I)$. The lifting property of $f$ against $i$ exhibits $f$ as a retract of $p$. Hence $f\in C(I)_*$. This proves the second trivial-class equality, and the model axioms follow.

Properness also passes to these structures. Pushouts and pullbacks are objectwise; injective cofibrations are objectwise cofibrations, and injective fibrations are objectwise fibrations by the preceding evaluation argument. The base properness axioms therefore apply at every object. These conclusions are recorded in \citet[Theorems~1.17 and 1.19, Propositions~1.21--1.22]{Barwick2007}; the numbering here is that of the 2007 preprint.

This settles the diagram construction under $a'$. It does not establish every proposed recognition implication involving $a$, $a''$ or $b$. Closure under filtered colimits alone is not the solution-set condition in a model-structure recognition theorem. The standard Smith theorem requires an accessible, accessibly embedded class of weak equivalences in the arrow category, as well as a specified set of generating cofibrations; see \citet[Proposition~1.7]{Barwick2007}. No assertion that the whole conjectural theorem has been proved is justified by citing the injective existence theorem. Conversely, once a combinatorial model structure is given, its diagram structures do not require preservation of \emph{every} filtered colimit by the weak equivalences. That stronger hypothesis belongs to Grothendieck's proposed route to recognition.
\end{dcomment}


Je m'occuperai de la démonstration du théorème précédent, ainsi que du théorème de factorisation du paragraphe 1, quand j'aurai mes notes sur les catégories et foncteurs accessibles et filtrations cardinales. En attendant, j'ai envie déjà d'expliciter un peu plus le formalisme auquel on parvient moyennant les hypothèses faites ci-dessus. (Je suis à peu près sûr tout au moins que la condition $a''$) suffit à impliquer toutes les autres. Et c'est bien sous la forme $a''$) que ces conditions vont se vérifier en pratique, je présume, notamment dans le cas où $\M = $ Cat.)

\dgrblock{dgrb697}{The variance is fixed by presheaves}
\begin{dcomment}{exposition}{p44:En attendant, j'ai envie déjà d'expliciter un peu plus}{The variance is fixed by presheaves}
Grothendieck now turns the diagram model structures into operations. Throughout these pages
\[
 \mathcal M(I)=\operatorname{Fun}(I^{\mathrm{op}},\mathcal M),
 \qquad u:I'\longrightarrow I.
\]
This convention determines the comma categories in both pointwise Kan formulas. If $X':(I')^{\mathrm{op}}\to\mathcal M$, then
\[
 \begin{aligned}
 (u_!X')(i)&=\operatorname*{colim}_{(i\downarrow u)^{\mathrm{op}}}X'(i'),\\
 (u_*X')(i)&=\operatorname*{lim}_{(u\downarrow i)^{\mathrm{op}}}X'(i').
 \end{aligned}
\]
The manuscript writes $i\backslash I'$ and $I'/i$, with the contravariant convention understood. Making the opposites explicit explains the otherwise surprising occurrence of an \emph{initial} object in footnote~27: it becomes terminal in the indexing category of the colimit. It also prevents interchanging the two comma categories when passing to modern covariant notation.

The lifting argument behind Proposition~2 is a precise form of the principle that adjunction transfers extension problems. Suppose $L:\mathcal A\rightleftarrows\mathcal B:R$ and let $i:A\to A'$ and $p:B\to B'$. By the adjunction, the two horizontal maps in a square from $Li$ to $p$ correspond to the two horizontal maps in a square from $i$ to $Rp$. Naturality of the adjunction bijection identifies the commutativity conditions. A diagonal $LA'\to B$ corresponds to a diagonal $A'\to RB$, and each of the two triangle identities required of a filler is carried to the corresponding identity. The sets of fillers are therefore naturally bijective. In particular
\[
 Li\mathrel{\perp}p\quad\Longleftrightarrow\quad i\mathrel{\perp}Rp.
\]
Applying this first to $u^*\dashv u_*$ proves the preservation of $F$ and $TF$ by $u_*$. Applying it to $u_!\dashv u^*$ proves the two equivalences in (3). This is the same elementary mechanism as \citet[Lemma~1.3.4]{Hovey1999}.

For a useful special case, suppose every component of $i\downarrow u$ has an initial object $e_\alpha$. The left Kan formula reduces to
\[
 (u_!X')(i)\cong\coprod_\alpha X'(e_\alpha).
\]
A left orthogonal class is closed under coproducts: a square out of a coproduct is a family of squares, and a family of fillers assembles to a filler. Thus $u_!$ preserves $C$ and $TC$ under this hypothesis, and $u^*$ preserves $F$ and $TF$. No claim about preservation of general colimits by the weak equivalences is involved. The comma-category hypothesis has replaced that issue by a coproduct calculation.
\end{dcomment}


\section*{4\quad Les variances $u^*$, $u_*$, $u_{!}$ pour $(\M(I),W(I))$ (I variable)}

Je suppose les conditions du paragraphe précédent vérifiées pour $(\M,W)$ (cf. théorème p. 42), et choisis une structure de Quillen faible close et exacte à gauche $(C,TC,F,TF)$ pour $W$, d'où pour tout $I$ dans Cat

\noindent\msp{page 45}

\noindent une structure de Quillen faible close et exacte à gauche $(C(I),TC(I),F(I),TF(I))$ pour $W(I)$, dans $\M(I)$, décrite page 43. Soit maintenant
\[
u \;:\; I' \longrightarrow I
\]
une flèche dans Cat, d'où une flèche
\[
u^* \;:\; (\M(I),W(I)) \;\longrightarrow\; (\M(I'),W(I')).
\]
On sait, vu que $\M$ est stable par $\limind$ quelconques et aussi (étant de plus accessible) par $\limproj$ quelconques, que $u^*$ admet des foncteurs adjoints à gauche et à droite,
\[
u_{!},\ u_* \;:\; \M(I') \;\longrightarrow\; \M(I),
\]
\printedpage{26}
% head: A. Grothendieck : Dérivateurs, Ch. XIII, Édition des Universités de Montpellier II et Paris VII
et je me propose ici de résumer les propriétés les plus immédiates de $u^*$ et de $u_{!}$, $u_*$ par rapport aux structures de Quillen faibles envisagées.

\textbf{Proposition 1 :} \textit{Avec les notations précédentes, on a les relations :}
\[
u^*(C(I)) \subset C(I'), \qquad u^*(TC(I)) \subset TC(I'). \tag*{(1)}
\]
\[
u_*(F(I')) \subset F(I), \qquad u_*(TF(I')) \subset TF(I). \tag*{(2)}
\]
\[
\begin{aligned}
u^*(F(I)) \subset F(I') &\iff u_{!}(TC(I')) \subset TC(I)\\
u^*(TF(I)) \subset TF(I') &\iff u_{!}(C(I')) \subset C(I).
\end{aligned} \tag*{(3)}
\]

D'autre part, je ne sais rien dire de positif pour les inclusions
\[
\begin{aligned}
u_*(C(I')) \;&\subset\; C(I), \;\; u_*(TC(I')) \;\subset\; TC(I)\\
u_{!}(F(I')) \;&\subset\; F(I), \;\; u_{!}(TF(I')) \;\subset\; TF(I).
\end{aligned}
\]
Si ce n'est que pour un $u$ tant soit peu général ($u$ pas un isomorphisme, disons), $u_*$ n'a aucune envie

\noindent\msp{page 46}

\noindent d'être compatible avec les structures $C$, $TC$ (propriété d'exactitude homotopique à droite, moralement, vu que $u_*$ est exact à gauche et en général pas du tout à droite), ni $u_{!}$ avec les structures $F$, $TF$ (commentaire dual à donner, $u_{!}$ est exacte à \textit{droite}, et n'a pas envie d'être “homotopiquement exact à gauche”). % [sic] exacte -> exact, accord avec « foncteur » / « exact à gauche »

\textsc{Démonstration} de (1), (2), (3). Les relations (1) sont évidentes par définition. Les relations (2), (3) le sont à peine moins, vu la propriété soritale générale suivante, dont la place serait au paragraphe 1 :

\textbf{Proposition 2.} \textit{Soit}
\[
\begin{tikzcd}
C' \arrow[r, shift left, "u_{!}"] & C \arrow[l, shift left, "u^*"]
\end{tikzcd}
\]
\textit{un couple de foncteurs adjoints $(u_{!},u^*)$, et soit $\Phi' \subset \Fl(C')$, $\Psi \subset \Fl(C)$. Alors}
\[
(u_{!}(\Phi') \longleftrightarrow \Psi) \qquad\quad \iff \qquad\quad (\Phi' \longleftrightarrow u^*(\Psi)).
\]

\textbf{Corollaire.} (C'est le même énoncé, avec seulement des notations changées.) \textit{Supposons que $u^* : C \longrightarrow C'$ admette un adjoint à} droite \textit{$u_*$ (pas forcément un adjoint à gauche)}
\[
\begin{tikzcd}
C' \arrow[r, shift right, "u_*"'] & C, \arrow[l, shift right, "u^*"']
\end{tikzcd}
\]

\printedpage{27}
% head: A. Grothendieck : Dérivateurs, Ch. XIII, Édition des Universités de Montpellier II et Paris VII
\noindent\textit{soit $\Phi \subset \Fl(C)$, $\Psi' \subset \Fl(C')$, alors}
\[
(u^*(\Phi) \longleftrightarrow \Psi') \qquad\quad \iff \qquad\quad (\Phi \longleftrightarrow u_*(\Psi')).
\]

C'est la proposition 2 pour \begin{tikzcd} C'_0 \arrow[r, shift left, "u_{0\,!}"] & C_0 \arrow[l, shift left, "u_0^*"] \end{tikzcd}, avec $C'_0 = C$, $C_0 = C'$, $(u_{0\,!} : C \longrightarrow C') = u^*$, et $(u_0^* : C' \longrightarrow C) = u_*$.

\noindent\msp{page 47}

Pour les inclusions (2), on utilise le corollaire de la proposition 2 pour les couples $\Phi = TC(I)$, $\Psi' = F(I')$, resp. $\Phi = C(I)$, $\Psi' = F(I')$ \texttt{[plutôt} $\Psi' = TF(I')$\texttt{]}, et on applique (1) qui donne $u^*(\Phi) \longleftrightarrow \Psi'$, d'où $\Phi \longleftrightarrow u_*(\Psi')$, ce qui n'est autre que (2). Pour les équivalences (3), on utilise la proposition 2, en faisant $\Psi = F(I)$, $\Phi' = TC(I')$ resp. $\Psi = TF(I)$, $\Phi' = C(I')$.

\textbf{Corollaire 1} (de la proposition 1) : \textit{Supposons $u : I' \longrightarrow I$ telle que pour tout $i \in \Ob I$, la catégorie $i\backslash I'$ soit discrète} (\footnote[27]{Plus généralement, il suffit que toute composante connexe $Z_\alpha$ de $i\backslash I'$ ait un objet initial $e_\alpha$, de sorte que $\limind\limits_{i\backslash I'} X'(i') \simeq \coprod_{\alpha\in\pi_0(i\backslash I')} X'(e_\alpha)$. C'est le cas notamment si $I' \longrightarrow I$ est un isomorphisme local, p. ex. une immersion ouverte.}) \textit{(p. ex. $I'$ discrète). Alors on a}
\[
\left\{
\begin{aligned}
u_{!}(C(I')) \;&\subset\; C(I), \;\; u_{!}(TC(I')) \;\subset\; TC(I)\\
u^*(F(I)) \;&\subset\; F(I'), \;\; u^*(TF(I)) \;\subset\; TF(I').
\end{aligned}
\right.
\]

En effet, les relations de la deuxième ligne équivalent (à cause de (3)) à celles de la première. Il faut prouver que si $f$ dans $\Fl\M(I')$ est dans $C(I')$ resp. $TC(I')$, $u_{!}(f)$ est dans $C(I)$ resp. $TC(I)$, ce qui signifie aussi que $\alpha_s^*(u_{!}(f))$ est dans $C$ resp. $TC$, pour tout $s \in \Ob I$. Or le foncteur
\[
\alpha_s^* u_{!} \;:\; \M(I') \;\longrightarrow\; \M
\]
s'explicite par la formule bien connue
\[
\alpha_s^*(u_{!}(X')) \;\overset{\text{déf}}{=}\; u_{!}(X')(s) \;\simeq\; \limind\limits_{i'\in s\backslash I'} X'(i').
\]

Comme

\noindent\msp{page 48}

\noindent $s\backslash I'$ est discrète, on a donc l'isomorphisme fonctoriel en $X'$ dans $\M(I')$
\[
\alpha_s^*(u_{!}(X')) \;\simeq \coprod_{(s',\, x \,:\, s \,\longrightarrow\, u(s')) \text{ dans } \Ob\,(s\backslash I')} X'(s'),
\]
donc si $f \,:\, X' \longrightarrow Y'$, $\alpha_s^*(u_{!}(f))$ s'explicite comme une somme directe dans $\underline{\mathrm{Fl}}(\M)$ de flèches de la forme
\[
f_{s'} \;:\; X'(s') \;\longrightarrow\; Y'(s'),
\]
\printedpage{28}
% head: A. Grothendieck : Dérivateurs, Ch. XIII, Édition des Universités de Montpellier II et Paris VII
lesquelles sont dans $C$ resp. $TC$ puisque $f$ est dans $C(I')$ resp. $TC(I')$. Comme $C$, $TC$ sont stables par sommes directes quelconques ($C$, $TC$ étant clos à gauche), on gagne.

\textbf{Corollaire 2.} \textit{Soit $I$ dans} Cat, \textit{et $f \in \Fl(\M(I))$. Pour que $f$ soit dans $F(I)$} (resp. \textit{dans $TF(I)$\,), il} faut \textit{que $\forall\, s \in \Ob I$, $f_s : X(s) \longrightarrow Y(s)$ le soit.} (Mais cette condition n'a pas du tout l'air d'être suffisante !) \textit{En particulier, soit $X \in \Ob\M(I)$, alors $X \in \Ob\M(I)_f \Longrightarrow X(s) \in \Ob\M_f \ \ \forall\, s \in \Ob I$.}

\dgrblock{dgrb698}{Ordinary limits and the two homotopical exactness properties}
\begin{dcomment}{examples}{p48:Mais cette condition n'a pas du tout l'air}{Ordinary limits and the two homotopical exactness properties}
Grothendieck's hesitation on pp.~49--50 identifies a necessary distinction. A right Kan extension can preserve fibrations and trivial fibrations without preserving weak equivalences between arbitrary diagrams. Its derived functor is computed on fibrant diagrams, and it is this functor which preserves homotopy limits. The distinction already matters for a finite poset.

Work in simplicial sets, and let $K$ be the nerve of the groupoid with two objects and exactly one morphism between every ordered pair. The simplicial set $K$ is Kan and contractible. Consider the cospan
\[
 *\xrightarrow{0}K\xleftarrow{1}*,
\]
where the two maps select distinct vertices. Its strict pullback is empty. Collapsing $K$ to a point gives an objectwise weak equivalence from this cospan to the constant one-point cospan, whose strict pullback is a point. Hence the ordinary limit functor does not preserve objectwise weak equivalences, even though every value of both diagrams is fibrant. The square formed by this weak equivalence over the identity of the terminal diagram is homotopy cartesian; its image under the strict limit is not, since it would require the map from the empty space to the point to be a weak equivalence. Fibrancy of values has not imposed the matching condition. The homotopy pullback of the first cospan is the space of paths from $0$ to $1$ in $K$, which is contractible, as it must be.

There is an equally concrete colimit example. Form the simplicial span
\[
 \Delta^1\longleftarrow\partial\Delta^1\longrightarrow\Delta^1
\]
using the two boundary inclusions. Its pushout is a simplicial circle. Collapsing both intervals gives an objectwise weak equivalence to $*\leftarrow\partial\Delta^1\to *$, whose strict pushout is a point. The first span consists of cofibrations and calculates the homotopy pushout; the second does not. Thus objectwise cofibrancy of all simplicial sets does not make every span a cofibrant diagram for computing colimits.

Even after deriving the right Kan extension, its preservation of homotopy pushouts is a different assertion and is generally false. Take $u$ from the discrete two-object category to the terminal category and work with pointed spaces. Then $\mathbf R u_*$ is the product. In the category of pairs, the coproduct of $(S^1,*)$ and $(*,S^1)$ is $(S^1,S^1)$. Applying the product gives the torus $S^1\times S^1$, whereas taking the coproduct after applying the product gives $S^1\vee S^1$. The comparison induces
\[
 \mathbb Z*\mathbb Z\longrightarrow\mathbb Z\times\mathbb Z
\]
on fundamental groups and is not an isomorphism. Since a pointed homotopy coproduct is a homotopy pushout over the zero object, this supplies the finite-poset example requested in the manuscript.

Restriction has a different behavior. Once the diagram localizations form a derivator, restrictions between shifted derivators preserve homotopy Kan extensions in the independent diagram variable, and hence both cartesian and cocartesian squares; this is \citet[Proposition~2.5]{Groth2013}. The two claims about restriction are therefore compatible. The asymmetry belongs to a general left or right Kan extension, and disappears in stable situations for \emph{finite} homotopy limits and colimits, where the two notions coincide.
\end{dcomment}


On applique le corollaire 1 à $\alpha_s : e \longrightarrow I$.

\textbf{Scholie.} La proposition 1 (relation (1)) montre que le foncteur $u^*$ transforme carrés $W$-cocartésiens en itou, suites exactes longues de suspension en itou – c'est “l'exactitude à droite homotopique” de $u^*$. On s'attend à ce qu'il en soit de même pour les carrés cartésiens et les suites exactes

\noindent\msp{page 49}

\noindent longues de cosuspension, mais pour le prouver il faudra un travail, qu'on fera dans la suite.

Les relations (2) montrent que le foncteur $u_*$ transforme carrés $W$-cartésiens en itou, donc suites exactes longues de cosuspension en itou. Je ne m'attends \textit{pas} à ce qu'il en soit de même en général pour les carrés $W$-cocartésiens – \textit{i.e.} que $u_*$ soit “homotopiquement exact à droite”, et non seulement à gauche. Mais il faudrait au moins que je fasse des exemples (avec $I'$, $I$ des catégories finies ordonnées).

Quant à $u_{!}$, la proposition ne dit rien de bien palpable à son sujet. Néanmoins on s'attend à ce que $u_{!}$ transforme carrés $W$-cocartésiens en itou – c'est la moindre des choses ! Ou peut-être je fais erreur, en confondant les foncteurs $u_*$, $u_{!}$ envisagés ici, avec les foncteurs similaires
\[
\begin{tikzcd}
\HO_W(I') \arrow[r, shift left, "u_{!}^W"] \arrow[r, shift right, "u_*^W"'] & \HO_W(I),
\end{tikzcd}
\]
adjoints à gauche resp. à droite de
\[
\begin{tikzcd}
\HO_W(I') & \HO_W(I), \arrow[l, "u_W^*"']
\end{tikzcd}
\]
induit par
\[
u^* \;:\; (\M(I),W(I)) \;\longrightarrow\; (\M(I'),W(I')).
\]
Mais en général $u_*$, $u_{!}$ ne sont pas compatibles

\noindent\msp{page 50}

\noindent avec les localiseurs $W(I')$, $W(I)$, donc ils ne doivent pas être confondus avec les foncteurs $u_{!}^W$, $u_*^W$. Donc j'ai été peut-être hâtif, dans ce que j'ai dit au sujet de $u_*$ et de $u_{!}$. Je serai obligé d'étudier la situation de très près, quand j'essayerai de construire les foncteurs $u_*^W$ et $u_{!}^W$. Pour le moment, je n'ai encore aucune assurance que les conditions envisagées pour $(\M,W)$ suffisent à assurer l'existence de $u_*^W$, $u_{!}^W$ – pas même pour $I$, $I'$ des ensembles ordonnés finis !

\textbf{Corollaire 1.} \textit{Soit $X \in \M(I)_{\mathrm{f}}$, alors pour toute $\alpha : s \longrightarrow t$ dans $I$ qui est} \texttt{[un]} monomorphisme, \textit{$X(\alpha) : X(t) \longrightarrow X(s)$ est dans $F$.}

% foot: Version 25 janvier 2004

% ---- from Der-ChXIII-25-01-04_p031-035.tex
\printedpage{29}
% head: A. Grothendieck : Dérivateurs, Ch. XIII, Édition des Universités de Montpellier II et Paris VII
(On applique \texttt{[le]} corollaire 1 (p. 47) à $\tilde{\alpha} : \boldsymbol{\Delta}_1 \longrightarrow I$.)

\textbf{Proposition 3.} (\footnote[28]{Comparer avec le lemme p. 63 et sa démonstration (plus exactement, le corollaire du lemme).}) \textit{Soit $X$ dans $\mathcal{M}(\boldsymbol{\Delta}_1)$, $X = (X_0 \xleftarrow{\alpha} X_1)$. Alors $X$ est dans $\mathcal{M}(\boldsymbol{\Delta}_1)_{\mathrm{f}}$ si et seulement si $X_0$, $X_1$} \texttt{[sont]} \textit{dans $\mathcal{M}_{\mathrm{f}}$ et $\alpha \in F$.}

\textbf{Corollaire 2.} \textit{Soit $I' \overset{u}{\hookrightarrow} I$} \texttt{[une]} immersion ouverte. \textit{Alors pour $X \in \operatorname{Ob}\mathcal{M}(I)_{\mathrm{f}}$, $\varprojlim\limits_{I} X(i) \longrightarrow \varprojlim\limits_{I'} X(i)$ est dans $F$.}

\section*{5\quad Catégories de modèles quillénisables}

Je reviens au cas général d'une catégorie de modèles $(\mathcal{M},W)$.

\dgrblock{dgrb699}{The two finite lifting arguments}
\begin{dcomment}{exposition}{p50:Je reviens au cas général d'une catégorie de modèles}{The two finite lifting arguments}
Grothendieck isolates consequences of diagram fibrancy which can be proved before any large existence theorem. The walking-arrow criterion of Proposition~3 is the special case $Y=*$ of the matching-map calculation discussed above: $X=(X_1\to X_0)$ is injectively fibrant precisely when $X_0$ is fibrant and $X_1\to X_0$ is a fibration. The fibrancy of $X_1$ then follows by composition. This proves the proposition without an appeal to the later lemma cited in footnote~28.

The monomorphism hypothesis in Corollary~1 has an exact combinatorial role. Let $\alpha:s\to t$ be monic and let $a:[1]\to I$ classify it. For any object $i$, the category $i\downarrow a$ has objects given by maps $i\to s$ or $i\to t$. A map of the first kind is joined to its composite with $\alpha$; a map of the second kind has at most one such predecessor because $\alpha$ is monic. Thus every component is either a singleton or a single arrow, and has an initial object. Footnote~27 shows that $a^*$ preserves injective fibrancy. Applying Proposition~3 to $a^*X$ gives $X(t)\to X(s)\in F$. No assumption that $I$ is itself a poset is needed.

For the restriction of limits, let $j:J\hookrightarrow I$ be a sieve, the presheaf convention for the manuscript's open immersion, and let $X$ be injectively fibrant. We prove directly that
\[
 \lim_I X\longrightarrow\lim_J j^*X
\]
has the right lifting property against $TC$. Take a square against $a:A\to B\in TC$. The upper map is a cone from $A$ to $X$ over $I$, while the lower map is a cone from $B$ over $J$ extending it. Since $j$ is a sieve, left extension $j_!$ is extension by the initial object outside $J$. The two cones combine to a map into $X$ from
\[
 T=(\operatorname{const}_I A)
 \mathbin{\amalg}_{j_!\operatorname{const}_J A}
 j_!\operatorname{const}_J B.
\]
The canonical map $T\to\operatorname{const}_I B$ is the identity at objects of $J$ and is $a$ elsewhere. It belongs to $TC(I)$. The fibrancy of $X$ supplies a lift $\operatorname{const}_I B\to X$, which is exactly the missing cone. This proves the corollary whenever the displayed limits and the required finite colimits exist. The construction exhibits the compatibility information in the lifting problem instead of suppressing it in a pointwise calculation.
\end{dcomment}


% marginal vertical bar printed beside the following Définition 1
\textbf{Définition 1.} On dit que $(\mathcal{M},W)$ est \textit{quillénisable} s'il existe une structure de Quillen faible \textit{exacte} $(C,TC,F,TF)$ pour $\mathcal{M}$.

C'est là une condition autoduale, mais que je ne sais vérifier en pratique que moyennant l'exactitude soit à gauche, soit à droite -- donc par une condition plus forte qui détruit la symétrie. C'est pour la sauvegarder dans la mesure du possible que j'introduis la définition précédente, ainsi que la

\textbf{Définition 2.} Soit, en plus de $(\mathcal{M},W)$, Diag $\subset$ Cat une sous-catégorie pleine de Cat. On dit que
\msp{page 51}$W$ est \textit{quillénisable sur} Diag, si pour tout $I$ dans Diag, $(\mathcal{M}(I),W(I))$ est quillénisable.

En pratique, j'utiliserai cette définition dans le cas où Diag $=$ Cat, ou alors, qu'il satisfasse au moins les propriétés de stabilité utiles pour le domaine d'un dérivateur : stable par $\varprojlim$ finies, par sommes finies -- eventuellement par $\varinjlim$ finies, etc (\footnote[29]{Je supposerai que Diag contient tout au moins les catégories provenant d'ensembles ordonnés finis.}). % [sic] eventuellement -> éventuellement, missing accent

\dgrblock{dgrb700}{What quillenizability promises}
\begin{dcomment}{exposition}{p51:En pratique, j'utiliserai cette définition}{What quillenizability promises}
Grothendieck wants a condition on the relative category $(\mathcal M,W)$, independent of the auxiliary choices of cofibrations and fibrations. Quillenizability says that some exact weak structure exists. Quillenizability on a domain of diagrams repeats this existence statement for every $I$ in the domain. It does not itself choose compatible structures. The exact-completion argument in the commentary to pp.~32--34 does allow one to enlarge the two trivial classes to $C\cap W$ and $F\cap W$, recovering the original, possibly nonclosed Quillen axioms without changing $C$, $F$ or $W$. Thus exactness has real force even before closure is imposed. The adjective ``weak'' must nevertheless retain the definition from p.~27. It is not a synonym for the later notions of a semi-model category or of a weak model category.

There are three separate questions. First, can one form the localization $\mathcal M(I)[W(I)^{-1}]$? Second, does it carry the fibre and cofibre operations constructed by replacement? Third, do restriction and Kan extension relate these categories with the exactness and pointwise formulas required of a derivator? Existence of localizations answers the first question, an exact weak structure addresses the second, and the uniform constructions aim at the third. Treating these three tasks separately explains why the author returns to examples from ordinary homological algebra rather than simply assuming arbitrary colimits from the outset.

His quotation marks around ``triangulée'' are also essential. In a pointed unstable theory there are suspension and loop functors, and long exact sequences obtained from fibres or cofibres. The suspension need not be an equivalence, the homotopy category need not be additive, and these data are not a triangulated category in Verdier's sense. Stable derivators do yield canonical triangulations, with exact restriction and Kan extension functors; the stability hypothesis is the additional mathematical content in \citet[Theorem~4.16]{Groth2013}. For complexes, shifting is invertible on the unbounded derived category, whereas nonnegative chain complexes provide a pointed but generally unstable example. This is precisely the distinction the next example needs.
\end{dcomment}


\dgrblock{dgrb701}{Quillenisability and the strength of its qualifiers}
\begin{dcomment}{exposition}{p51:En pratique, j'utiliserai cette définition}{Quillenisability and the strength of its qualifiers}
Grothendieck is looking for hypotheses on a category with weak equivalences which survive passage to diagrams. The point is not merely to construct another presentation of its homotopy category. A derivator must retain enough of the original category to restrict a diagram, replace it, and eventually take its homotopy limit. His succession of adjectives records different parts of this programme; reading them all as the modern expression ``model category'' would erase the question he is asking.

Throughout these pages we use his variance:
\[
 \mathcal M(I)=\operatorname{Fun}(I^{\mathrm{op}},\mathcal M),
 \qquad W(I)=\{f\mid f_i\in W\text{ for every }i\}.
\]
Thus the ordered category $0\to1$ indexes an arrow $X_1\to X_0$ in $\mathcal M$. A terminal object of $I$ computes the ordinary limit of such a diagram. This explains several formulas which would have the opposite orientation with covariant diagrams.

The glossary supplied with the chapter makes the terminology precise. A weak Quillen structure has the two factorizations $(C,TF)$ and $(TC,F)$, their lifting properties, and $TC\subset C\cap W$, $TF\subset F\cap W$. The axiom called exactness is
\[
                         W=TF\circ TC.
\]
Exactness on the left is the stronger-looking equality $TC=C\cap W$; exactness on the right is $TF=F\cap W$. Closedness says that the four classes are determined by their respective lifting properties. The additional axiom $\mathbf M''$ concerns preservation of weak equivalences under base change along fibrations and cobase change along cofibrations. In a usual model structure these are the properness requirements. Finally, $\mathbf M0$ asks only for finite limits and finite colimits. Small completeness must therefore be supplied when an infinite Kan extension is used.

There is a useful simplification precisely when \emph{closedness and exactness are both present}, as they will be in \S6. It is the retract argument, not an extra smallness theorem. Suppose $f\in C\cap W$. Exactness supplies $f=pi$ with $i\in TC$ and $p\in TF$. The left lifting property of $f$ against $p$ makes $f$ a retract of $i$ in the arrow category. A class defined by a left lifting property is closed under retracts, so $f\in TC$. For $f\in F\cap W$, use the same factorization and lift $i$ against $f$; this exhibits $f$ as a retract of $p$, and hence puts it in $TF$. Consequently
\[
              TC=C\cap W,\qquad TF=F\cap W
\]
under the closed exact hypotheses. This application of the familiar retract argument is particularly consequential here: see \citet[Lemma 1.1.9]{Hovey1999}. The proof uses neither arbitrary products nor a prior assertion that $W$ is closed under retracts.

Quillenisability by itself, in Definition 1, does not include closedness. We therefore do not substitute the stronger hypothesis into \S5. Once a uniform closed exact quillenisability has been assumed in \S6, however, the preceding equalities are available in every diagram category. They will allow the apparently separate questions about finite posets, products, towers and group actions to be treated by one argument. Keeping this distinction visible preserves both the breadth of the definition and the actual force of the theorem's hypotheses.
\end{dcomment}


\textbf{Exemple 1 :} Celui des paragraphes 3, 4, où on a dégagé des conditions sur $(\mathcal{M},W)$, avec $\mathcal{M}$ une ``grosse'' catégorie et $\mathcal{M}$ accessible, pour que $W$ soit quillénisable sur Cat tout entier. C'est le cas qui me paraît à présent le plus intéressant.

\textbf{Exemple 2 :} Soient $\mathcal{A}$ une catégorie abélienne, et $\mathcal{M} = K^{\bullet}(A)$ la catégorie des complexes de cochaînes à valeurs dans $\mathcal{A}$, $W = W_{\mathcal{A}}$ l'ensemble des quasi-isomorphismes, donc % [sic] K^\bullet(A) -> K^\bullet(\mathcal{A}), plain italic A where calligraphic A is used everywhere else
\[ \mathcal{M}W^{-1} \;=\; K^{\bullet}(\mathcal{A})W_{\mathcal{A}}^{-1} \;\simeq\; D^{\bullet}(\mathcal{A}), \qquad \text{catégorie dérivée de $\mathcal{A}$.} \]
On définit aussi $\mathcal{M}^{+}$, $\mathcal{M}^{-}$, $\mathcal{M}^{\mathrm{b}}$, avec les localiseurs induits $W_{\mathcal{A}}^{+}$, $W_{\mathcal{A}}^{-}$, $W_{\mathcal{A}}^{\mathrm{b}}$, et les catégories de fractions $D^{+}(\mathcal{A})$, $D^{-}(\mathcal{A})$, $D^{\mathrm{b}}(\mathcal{A})$. Ou
\msp{page 52}encore, pour $I$ dans Cat,
\[ \mathcal{M}(I) \;=\; K^{\bullet}(\mathcal{A}(I)) \]
% foot: Version 25 janvier 2004
\printedpage{30}
% head: A. Grothendieck : Dérivateurs, Ch. XIII, Édition des Universités de Montpellier II et Paris VII
\[ W(I) \;=\; \underbrace{W_{\mathcal{A}(I)}}_{\text{quasi-isomorphismes}} \;\subset\; \mathrm{Fl}(K^{\bullet}(\mathcal{A}(I))), \]
et de même pour $\mathcal{M}^{\alpha}(I)$, $W^{\alpha}(I)$ $(\alpha \in \{+,-,\mathrm{b}\})$, du moins \textit{si $I$ est finie.} Il y a aussi les sous-catégories $K^{\geq 0}(\mathcal{A})$ et $K^{\leq 0}(\mathcal{A})$ des complexes de cochaînes resp. de chaînes au sens strict (avec $K^{i} = 0$ si $i \geq 0$, resp. si $i \leq 0$ \texttt{[plutôt si $i < 0$, resp. si $i > 0$]}). On aura pour $I$ quelconque
\[ \left\{ \begin{aligned} K^{\geq 0}(\mathcal{A})(I) &\;\simeq\; K^{\geq 0}(\mathcal{A}(I)) \\ W_{\mathcal{A}}^{\geq 0}(I) &\;=\; W_{\mathcal{A}(I)}^{\geq 0}, \end{aligned} \right. \]
et de même pour $K^{\leq 0}$, $W^{\leq 0}$.

D'après \textsc{Quillen} page 1.2, $(K^{+}(\mathcal{A}),W_{\mathcal{A}}^{+})$ est quillénisable si $\mathcal{A}$ a assez d'injectifs, $(K^{-}(\mathcal{A}),W_{\mathcal{A}}^{-})$ l'est si $\mathcal{A}$ a assez de projectifs. Il ne dit rien pour $(K^{\bullet}(\mathcal{A}),W_{\mathcal{A}})$ sans plus -- pourtant on fait de l'algèbre cohomologique avec $D^{\bullet}(\mathcal{A})$ sans plus, et même sans faire d'hypothèse d'existence d'injectifs ou de projectifs. J'ai un peu oublié ce que \textsc{Verdier} a fait à ce sujet -- je devrais quand même regarder sa thèse (!), ou bien l'exposé de \textsc{Hartshorne}.

\dgrblock{dgrb702}{Two meanings of a bounded diagram of complexes}
\begin{dcomment}{examples}{p52:je devrais quand même regarder sa thèse}{Two meanings of a bounded diagram of complexes}
Grothendieck's qualification ``du moins si $I$ est finie'' protects an important distinction. We use $\operatorname{Ch}(\mathcal A)$ for his $K^\bullet(\mathcal A)$, which here denotes complexes and chain maps, before quotienting by chain homotopy. For unbounded complexes there is a genuine identification
\[
 \operatorname{Fun}(I^{\mathrm{op}},\operatorname{Ch}(\mathcal A))
 \cong\operatorname{Ch}(\operatorname{Fun}(I^{\mathrm{op}},\mathcal A)).
\]
Both sides specify objects $X_i^n$, differentials, and structure maps satisfying the same commutation relations. Quasi-isomorphisms are objectwise because kernels, cokernels, and hence cohomology in an abelian functor category are computed objectwise. The same argument works for a fixed half-line of degrees, such as $n\geq0$.

For $\operatorname{Ch}^+$ there are two possible bounds. An object of $\operatorname{Fun}(I^{\mathrm{op}},\operatorname{Ch}^+(\mathcal A))$ has a lower bound depending on $i$. An object of $\operatorname{Ch}^+(\mathcal A(I))$ has one lower bound valid for every $i$. A finite object set allows us to take the minimum of the finitely many lower bounds, so the two categories agree. An infinite object set need not. For $I=\mathbb N$ discrete and a nonzero abelian group $A$, the family $A[n]$, $n\geq0$, has one nonzero term in degree $-n$ and gives an immediate counterexample. There are analogous upper-bound and two-sided-bound distinctions.

This also explains footnote~31. In $\operatorname{Ch}^+(\mathbf{Ab})$ take the filtered system
\[
 X_m=\bigoplus_{0\leq n\leq m}\mathbb Z[n],\qquad
 X_m\longrightarrow X_{m+1}
\]
with the evident inclusions and zero differentials. Its colimit in unbounded complexes has a copy of $\mathbb Z$ in every nonpositive degree. It does not belong to $\operatorname{Ch}^+$. Thus the assertion ``$K^\bullet(\mathcal A)$ etc. have the same properties'' cannot include bounded-below or bounded-above complexes without a common bound; the manuscript itself flags this point.

More seriously, the difference can destroy an adjoint. Suppose the right adjoint to the constant-diagram functor
\[
 D^+(\mathbf{Ab})\longrightarrow
 \prod_{n\geq0}D^+(\mathbf{Ab})
\]
existed at the family $(\mathbb Z[n])_{n\geq0}$, and call its value $P$. Applying the adjunction to $\mathbb Z[k]$ gives
\[
 \begin{aligned}
 H^{-k}(P)&\cong\operatorname{Hom}_{D(\mathbf{Ab})}(\mathbb Z[k],P)\\
 &\cong\prod_{n\geq0}
 \operatorname{Hom}_{D(\mathbf{Ab})}(\mathbb Z[k],\mathbb Z[n])\\
 &\cong\mathbb Z\qquad(k\geq0).
 \end{aligned}
\]
The last equality uses the projectivity of $\mathbb Z$: only the factor $n=k$ is nonzero. No bounded-below complex can have these cohomology groups. Therefore this right adjoint does not exist. For a discrete indexing category, the localization of the objectwise bounded complexes is indeed the product of the bounded derived categories: choose a bounded-below injective resolution separately in each component and represent morphisms by chain maps between these resolutions.

Accordingly, the familiar $\mathbf R u_*:D^+(\mathcal A(I'))\to D^+(\mathcal A(I))$ refers to a \emph{uniformly} bounded-below input. It must not be used without comment for the generally larger localization of $\operatorname{Fun}((I')^{\mathrm{op}},\operatorname{Ch}^+(\mathcal A))$. This failure concerns the requested adjoint, not a proof that the latter relative category has no possible weak quillenization. The commentary on Chapter~VI, manuscript p.~131, already records a uniform-bound obstruction; here the obstruction is located directly in the proposed domain of a derived Kan extension.
\end{dcomment}


\msp{page 53}Toujours est-il qu'on a là des exemples de type tout différent de ceux envisagés dans 1°. Ainsi, il n'y a pas à supposer que $\mathcal{M}$ soit stable par petites $\varinjlim$ -- $\mathcal{M}$ peut même être une petite catégorie, elle l'est si et seulement si $\mathcal{A}$ elle même l'est. Si $\mathcal{A}$ est stable par (petites) $\varinjlim$ filtrantes, et est accessible (p. ex. a une petite famille génératrice, et les $\varinjlim$ filtrantes sont exactes), alors $K^{\bullet}(\mathcal{A})$ etc. (\footnote[30]{!!}) ont les mêmes propriétés, et $(K^{+}(\mathcal{A}),W_{\mathcal{A}}^{+})$ est quillénisable sur Cat tout entier (\footnote[31]{attention, $K^{+}(\mathcal{A})$ n'est \textit{pas} stable par $\varinjlim$ filtrantes infinies -- donc on ne peut appliquer le tapis des paragraphes 2, 3 directement à $K^{+}$.}). (C'est le cas où on sait que $\mathcal{A}$ a assez d'injectifs, et l'algèbre cohomologique se fait surtout à coups de résolutions injectives \dots) Dans ce cas on sait que les % [sic] elle même -> elle-même, missing hyphen
\[ u_{*} \;:\; D^{+}(I',\mathcal{A}) \longrightarrow D^{+}(I,\mathcal{A}) \]
existent comme adjoints à droite des
\[ u^{*} \;:\; D^{+}(I,\mathcal{A}) \longrightarrow D^{+}(I',\mathcal{A}), \]
mais je ne suis plus trop sûr si on sait définir $u_{*}$ aussi sur $D(I',\mathcal{A})$ tout entier, et je doute que $u_{!}$ existe sans hypothèse sur $u$ (genre : dimension cohomologique finie). Sans doute on doit toujours avoir
\[ u_{!} \;:\; D^{-}(I',\mathcal{A}) \longrightarrow D^{-}(I,\mathcal{A}), \]
\msp{page 54}adjoint à gauche de
\[ u^{*} \;:\; D^{-}(I,\mathcal{A}) \longrightarrow D^{-}(I',\mathcal{A}). \]
% foot: Version 25 janvier 2004

\dgrblock{dgrb703}{A bound on each component is not a bound on the diagram}
\begin{dcomment}{examples}{p53}{A bound on each component is not a bound on the diagram}
Grothendieck turns to complexes because they offer a familiar theory of derived coefficients without the large-colimit hypotheses of the preceding sections. His notation $K^\bullet(\mathcal A)$ denotes complexes and chain maps here. To avoid conflict with the now customary use of $K(\mathcal A)$ for the chain-homotopy category, we write $\operatorname{Ch}(\mathcal A)$ for that category of complexes. Localization at quasi-isomorphisms gives the derived category in either convention.

For unbounded complexes the exchange of the two diagram directions is literal:
\[
 \operatorname{Fun}(I^{\mathrm{op}},\operatorname{Ch}(\mathcal A))
 \simeq \operatorname{Ch}(\operatorname{Fun}(I^{\mathrm{op}},\mathcal A)).
\]
Both sides are the same collection of objects, differentials, transition maps and commutativities. Quasi-isomorphisms agree because kernels, cokernels and cohomology in the functor category are computed at each object. The fixed truncation conditions $X^n=0$ for $n<0$, or for $n>0$, also commute with arbitrary diagram categories. The editor's correction to the reversed inequalities on manuscript p.~52 is essential.

The movable bound in $\operatorname{Ch}^+$ behaves differently. A functor into $\operatorname{Ch}^+(\mathcal A)$ gives a lower bound $b_i$ for each component. An object of $\operatorname{Ch}^+(\mathcal A(I))$ requires one bound valid for all components. These conditions agree for a finite set of objects, since one can take $\min_i b_i$, but need not agree otherwise. The same distinction applies to upper bounds and to bounded complexes. This is why the words ``si $I$ est finie'' belong to the argument, rather than merely to its notation.

Here is the obstruction in its simplest form. Let $k$ be a field and put $X_n=k[n]$, with our cochain convention $X[n]^j=X^{j+n}$. Thus $X_n$ has a single nonzero term in degree $-n$. Each $X_n$ is bounded below. The degreewise product
\[
                         \prod_{n\geq0} k[n]
\]
has nonzero cohomology in every degree $0,-1,-2,\ldots$ and is not bounded below. This is not merely a defect of that chosen complex. If a product $P$ of the family existed in $D^+(k)$, the defining universal property, tested on $k[m]$, would give
\[
 H^{-m}(P)=\operatorname{Hom}_{D(k)}(k[m],P)
 \cong\prod_{n\geq0}\operatorname{Hom}_{D(k)}(k[m],k[n])
 \cong k
\]
for every $m\geq0$. No object of $D^+(k)$ has this property. Thus the projection from an infinite discrete category does not have a right adjoint on the theory of \emph{objectwise} bounded-below diagrams with values in $D^+(k)$.

The filtered-colimit warning in footnote 31 has an equally concrete witness. Let $P_r=\bigoplus_{0\leq n\leq r}k[n]$ and include $P_r$ into $P_{r+1}$. Each $P_r$ is bounded, whereas their degreewise filtered colimit has nonzero terms arbitrarily far to the left. Hence $\operatorname{Ch}^+$, and likewise $\operatorname{Ch}^{\mathrm b}$, is not closed under arbitrary filtered colimits in $\operatorname{Ch}$. The dual example uses $k[-n]$ to treat $\operatorname{Ch}^-$. The unbounded category and the two fixed half-line categories do not have this problem.

Enough injectives do give the familiar bounded-below resolutions, and enough projectives give bounded-above resolutions. What the argument on these pages does not establish is quillenisability on every infinite diagram shape for the category allowing a different bound at every object. The accessibility theorem quoted from the preceding sections cannot be applied directly to that category. We record this as a missing justification, not as a counterexample to quillenisability itself: the product calculation concerns a Kan extension, which a finite-limit model structure is not obliged to possess. Passing to unbounded complexes is the clean way to obtain the unrestricted diagram theory.
\end{dcomment}


\printedpage{31}
% head: A. Grothendieck : Dérivateurs, Ch. XIII, Édition des Universités de Montpellier II et Paris VII
Mais je ne sais plus trop s'il faut s'attendre à pouvoir définir $u_{!}$ sur $D(I',\mathcal{A})$ tout entier, en dehors d'hypothèses genre dimension cohomologique finie. La triste vérité, c'est que j'ai $\pm$ oublié la théorie des foncteurs dérivés ! (Il faudrait que je demande encore à \textsc{Jean Malgoire} de m'envoyer des textes canoniques.)

\dgrblock{dgrb704}{Unbounded resolutions remove the dimensional restriction}
\begin{dcomment}{later}{p54:La triste vérité, c'est que}{Unbounded resolutions remove the dimensional restriction}
Grothendieck's uncertainty about unbounded derived functors has a precise modern answer for Grothendieck abelian categories. An unbounded complex needs a resolution which is injective as a complex, in the homotopical sense, rather than merely injective in each degree. A complex $J$ is $K$-injective when $\operatorname{Hom}^\bullet(A,J)$ is acyclic for every acyclic complex $A$. This condition makes homotopy classes of maps into $J$ invariant under quasi-isomorphisms in the source.

\citet[Theorems~C and D]{Spaltenstein1988} constructed unbounded resolutions for modules and for sheaves of modules. For an arbitrary Grothendieck abelian category, the general existence result is stated and proved in \citet[Theorem~19.12.6, tag 079P]{StacksKInjectives}. One may also equip $\operatorname{Ch}(\mathcal A)$ with the proper injective model structure whose weak equivalences are quasi-isomorphisms and whose cofibrations are monomorphisms; see \citet[Theorem~1.2]{CisinskiDeglise2009}, where the existence theorem is attributed to Beke.

For a small category $I$, the abelian category $\mathcal A(I)$ is again Grothendieck. Exactness of filtered colimits is objectwise. A generating family is obtained by applying the left adjoints to evaluation to a generating family of $\mathcal A$; these left adjoints are coproducts indexed by hom-sets. Applying the injective and projective diagram constructions to unbounded complexes consequently gives
\[
 \mathbf L u_!\dashv u^*\dashv\mathbf R u_*
 \quad\text{on}\quad
 D(\mathcal A(I'))\ \text{and}\ D(\mathcal A(I))
\]
for every functor of small indexing categories. Here ``projective diagram structure'' refers to the diagram model structure and does not assert that $\mathcal A$ has enough projective objects. No finite cohomological dimension of $u$ is required for existence on the unbounded categories.

Finite dimensionality becomes relevant when one asks these functors to preserve a chosen bounded subcategory. This can be seen without any sophisticated example. Let $G$ be cyclic of prime order $p$, let $k=\mathbb F_p$, and consider constant coefficients $k$ over $BG$. Write $\epsilon=g-1$ in $k[G]\cong k[\epsilon]/(\epsilon^p)$. The periodic free resolution of the trivial module has differentials alternating multiplication by $\epsilon$ and $\epsilon^{p-1}$:
\[
 \cdots\xrightarrow{\epsilon}k[G]
 \xrightarrow{\epsilon^{p-1}}k[G]
 \xrightarrow{\epsilon}k[G]\longrightarrow k\longrightarrow0.
\]
It is exact because the kernel of multiplication by $\epsilon$ is the image of multiplication by $\epsilon^{p-1}$, and conversely. After tensoring with $k$, or applying $\operatorname{Hom}_{k[G]}(-,k)$, every differential is zero. Thus
\[
 H^{-n}(\mathbf L\operatorname{colim}_{BG}k)\cong k,
 \qquad
 H^n(\mathbf R\lim_{BG}k)\cong k
 \quad(n\geq0).
\]
The left derived colimit of this degree-zero input is not bounded below, and the right derived limit is not bounded above. Both nevertheless exist as unbounded complexes. The source's worries about existence and about boundedness therefore have different resolutions. For an arbitrary abelian category, existence of the Verdier localization alone still does not supply either unbounded resolutions or all these Kan extensions.
\end{dcomment}


\dgrblock{dgrb705}{From a presentation to a homotopy theory of diagrams}
\begin{dcomment}{later}{p54:Mais je ne sais plus trop s'il faut}{From a presentation to a homotopy theory of diagrams}
The uncertainty about derived direct images brings Grothendieck back to the
reason for keeping all indexing categories at once. A construction on one
localized category can conceal the choices needed to make it functorial.
Passing first to diagrams and then localizing preserves the information on
which the construction depends. In particular, a coherent square is an
object of the localization of the category of squares; it is not merely a
commutative square in the ordinary homotopy category.

Renaudin makes the comparison with Quillen presentations precise for
combinatorial model categories. In the published version, the bicategorical
localization of left proper combinatorial model categories at Quillen
equivalences maps locally equivalently to the 2-category of derivators with
adjunctions as morphisms; restricting to derivators of small presentation
gives a biequivalence \citep[Theorems~3.3.2 and~3.4.4]{Renaudin2009}.
The choice of morphisms is part of this statement. It is not an assertion
that every prederivator, or every transformation of prederivators, comes from
an arbitrary relative category in this way.

The analogous construction for an infinity-category $\mathcal C$ has the
particularly transparent form
\[
 I\longmapsto
 \operatorname{Ho}\operatorname{Fun}(NI^{\mathrm{op}},\mathcal C).
\]
Here $N$ is the nerve, and the opposite records the presheaf convention of
this chapter. Limits and colimits in $\mathcal C$ give the homotopy Kan
extensions of this system. Riehl and Verity express the relevant comparison
law as a Beck--Chevalley theorem for comma squares, and explain the resulting
derivator of a complete and cocomplete quasi-category
\citep[Proposition~5.3.9 and Remark~5.3.10]{RiehlVerity2017}.

This description also clarifies the role of the classes $C$ and $F$ in the
preceding calculations. They supply replacements that compute these
intrinsic operations. The right replacement for a presheaf is generally
more than a list of fibrant values: the transition maps carry compatibility
conditions. Grothendieck's arrow example on manuscript page~50 already
exhibits this fact. On the other hand, forgetting to derive an ordinary Kan
extension destroys precisely the homotopy invariance that the intrinsic
description requires. The withdrawal on manuscript pages~49--50 corrects
the construction while preserving its objective.
\end{dcomment}


\textbf{Scholie.} Si $(\mathcal{M},W)$ est quillénisable sur Diag $\subset$ Cat (et même sans qu'il soit quillénisable), il définit un prédérivateur $\mathbf{D}_{\mathcal{M},W}$ sur Diag. On posera parfois
\[ \mathrm{HO}_{\mathcal{M},W}(I) \;\text{ ou }\; \mathrm{HO}_{W}(I) \;=\; \mathbf{D}_{\mathcal{M},W}(I). \]
L'hypothèse quillénisable implique, sur chacune de ces catégories homotopiques, la structure (relativement faible) ``triangulée'', provenant des suites \textit{exactes} de suspension et de cosuspension (\footnote[32]{Dans le cas où $\mathcal{M}$ est ponctuée par \texttt{[un]} objet neutre. Sinon, il y a la structure associée aux carrés cartésiens et cocartésiens. Il serait assez intéressant d'axiomatiser une telle espèce de structure.}). (Le caractère quillénisable assure l'exactitude de ces suites \dots) Les foncteurs $u^{*} : \mathrm{HO}_{W}(I) \longrightarrow \mathrm{HO}_{W}(I')$
\msp{page 55}sont ``homotopiquement exacts'' dans le sens des ces structures ``triangulées'' (non commutatives), \textit{i.e.} transforment suites exactes de suspension resp. de cosuspension en itou. % [sic] des ces -> de ces, grammar

\dgrblock{dgrb706}{An answer already present in Joyal's letter}
\begin{dcomment}{context}{p54}{An answer already present in Joyal's letter}
The questions about complexes and uniform diagram structures have a particularly close antecedent in the correspondence surrounding \emph{Pursuing Stacks}. In the letter begun on 11 April 1984, Joyal describes Heller's two model structures on simplicial presheaves and explicitly connects them with derived right and left Kan extensions for an arbitrary functor of indexing categories. In the continuation dated 15 April, Theorem 2 states the injective model structure on unbounded complexes in an arbitrary Grothendieck abelian category, with monomorphisms as cofibrations and homology isomorphisms as weak equivalences \citep[pp.~3 and 10]{Joyal1984Letter}.

This locates the present hesitation accurately. The manuscript records a return to an earlier set of questions, not the historical first appearance of the unrestricted complex example. The later proofs by Beke and the constructions of $K$-injective resolutions make those statements accessible in systematic published forms. The letter is also unusually direct evidence for the role of Heller in the diagram-category approach; it supports that precise connection without requiring us to infer a chronology from similarities between later axiomatizations.
\end{dcomment}


\dgrblock{dgrb707}{Unbounded resolutions answer the derived-image question}
\begin{dcomment}{later}{p54}{Unbounded resolutions answer the derived-image question}
Grothendieck's hesitation about unbounded derived images concerns the existence of resolutions on which an ordinary functor can safely be evaluated. A complex of injective objects is not automatically suitable when it is unbounded: what is needed is a condition on the whole complex. A complex $J$ is $K$-injective when $\operatorname{Hom}^\bullet(A,J)$ is acyclic for every acyclic complex $A$. Equivalently, maps into $J$ in the homotopy category already compute maps into its image in the derived category.

Spaltenstein's work was already available when these manuscript pages were written. His Theorem C treats unbounded complexes of modules, Theorem D treats sheaves of modules, and Theorem 4.5 gives the latter existence construction; these results should not be retrospectively attributed to the present text \citep{Spaltenstein1988}. For an arbitrary Grothendieck abelian category the general result is associated with \citet{Serpe2003}. A detailed modern statement, including a functorial monomorphic quasi-isomorphism into a termwise injective $K$-injective complex, is Theorem 19.12.6, Tag 079P, of the Stacks Project \citep{StacksKInjective079P}.

The application to right Kan extensions is short once this existence theorem has been supplied. Let $u:I'\to I$ be a functor of small categories and let $\mathcal A$ be Grothendieck abelian. The category $\mathcal A(I')$ is again Grothendieck: exactness of filtered colimits is checked componentwise, and free diagrams on a set of generators give a set of generators. Restriction $u^*$ is exact, and its ordinary right adjoint $u_*$ exists. If $J$ is $K$-injective in $\operatorname{Ch}(\mathcal A(I'))$, then for any acyclic $A$ in $\operatorname{Ch}(\mathcal A(I))$ the ordinary adjunction gives
\[
 \operatorname{Hom}^\bullet(A,u_*J)
   \cong\operatorname{Hom}^\bullet(u^*A,J).
\]
The right-hand side is acyclic because $u^*A$ is acyclic. Thus $u_*J$ is $K$-injective. Replacing $X'$ by such a $J$ defines $\mathbf Ru_*X'=u_*J$ on the entire unbounded derived category, and the same Hom identity proves $u^*\dashv\mathbf Ru_*$ there.

The corresponding unbounded left Kan extensions also exist for Grothendieck coefficients. The precise model-theoretic input is \citet[Proposition 3.13]{Beke2000}: unbounded complexes have a cofibrantly generated model structure with monomorphisms as cofibrations and quasi-isomorphisms as weak equivalences. The underlying category is locally presentable, so this is a combinatorial model; its projective diagram structures and the Brown argument supply all left Kan extensions. Neither construction requires a finite cohomological dimension for the indexing functor. Such dimension restrictions instead govern preservation of bounded ranges. Since restriction is exact for the standard degreewise $t$-structures, its right adjoint preserves lower bounds and its left adjoint preserves upper bounds. Thus the usual statements on globally bounded-below and globally bounded-above diagrams remain meaningful, while a bounded object can acquire cohomology in arbitrarily many degrees on the other side.

For this reason the letter $u_!$ here must be read as the left adjoint to restriction of diagrams. It is not a direct image with proper supports on a space. Those functors have their own geometric hypotheses; importing them would introduce difficulties unrelated to the question asked on pp.~53--54. The example of a finite group below will display the genuine boundedness issue without any such ambiguity.
\end{dcomment}


Quels sont les axiomes des dérivateurs qui sont vérifiés sous la seule hypothèse quillénisable ? \textbf{Der} 1 \texttt{[voir chapitre I, p. 43]} l'est sans hypothèse aucune sur la catégorie de modèles $(\mathcal{M},W)$ :
\[ \mathrm{HO}_{W}(\coprod_{\alpha} I_{\alpha}) \;\simeq\; \prod_{\alpha} \mathrm{HO}_{W}(I_{\alpha}) \]
pour une famille \textit{finie} d'objets $I_{\alpha}$ dans Diag. L'axiome de localisation \textbf{Der} 2 \texttt{[voir chapitre I, p. 43]} l'est sous l'hypothèse que $W$ soit \textit{fortement saturé,} et quillénisable. On utilise le fait que toute flèche de $\mathrm{HO}_{W}(I)$ provient, à isomorphisme près, d'une flèche dans $\mathcal{M}(I)$. (C'était là un des points qui restait en suspens dans VI, dans la théorie de $\mathrm{HOT}_{W}$, $W \subset \mathrm{Fl}(\mathrm{Cat})$ un localiseur fondamental \dots)

\dgrblock{dgrb708}{The two elementary derivator axioms}
\begin{dcomment}{exposition}{p55:On utilise le fait que toute flèche de}{The two elementary derivator axioms}
Grothendieck is now testing how much of the derivator is already present before constructing homotopy Kan extensions. The two assertions about finite coproducts of shapes and detection of isomorphisms have different proofs, and the latter uses a homotopical input which cannot be replaced by saturation alone.

For a finite family of shapes there is an isomorphism of relative categories
\[
 \mathcal M\Bigl(\coprod_{r=1}^n I_r\Bigr)
 \cong\prod_{r=1}^n\mathcal M(I_r)
\]
with componentwise weak equivalences. Localization commutes with this finite product. One way to see the point is to separate each arrow $(f_1,\ldots,f_n)$ into the $n$ coordinate arrows, inserting identities in the other positions. Coordinate arrows commute, and this commutation remains true after inverting weak equivalences. Finite zigzags in the individual coordinates can therefore be padded with identities and assembled into a finite zigzag in the product. The same operation on the elementary relations between zigzags gives full faithfulness as well as essential surjectivity. The empty family gives the terminal category. Infinite products require an additional argument: one cannot pad zigzags of unbounded length to a single finite length. This is why the statement on the page is deliberately finite.

For detection of isomorphisms, use the representation property of the exact weak structure proved in the commentary to pp.~29--34: every arrow $\xi$ in $\mathcal M(I)[W(I)^{-1}]$ is isomorphic in the arrow category to the localization of an actual map $f:X\to Y$. If every evaluation of $\xi$ is invertible, the same is true of every $f_i$ in $\mathcal M[W^{-1}]$. Strong saturation of $W$ then gives $f_i\in W$ for all $i$, hence $f\in W(I)$, and $\xi$ is invertible. The role of representability is to put all components into one honest natural transformation before applying saturation. The commentary on Chapter~VI, manuscript p.~130, already explains why saturation by itself is insufficient.

These arguments establish the two stated axioms with their stated qualifications. They do not produce the pointwise Kan-extension formula. In the modern model-category treatment, that remaining structure is available: \citet[Proposition~1.30]{Groth2013} proves the strong derivator theorem in the combinatorial case, and his Theorem~1.32 records Cisinski's theorem for model categories in general. The word ``model'' in that latter statement has the ordinary model-category meaning. The exact-completion lemma above first recovers Quillen's original axioms from the weak structure; closedness, size conventions and the required small limits and colimits must still be tracked when invoking a modern formulation.
\end{dcomment}


L'axiome \textbf{Der} 3, nouvelle mouture \texttt{[voir chapitre I, pp. 111-112]}, implique un localiseur fondamental $\mathbf{W} \subset \mathrm{Fl}(\mathrm{Cat})$ (pas forcément $\mathbf{W}_{\infty}$) (\footnote[33]{Il est naturel ici de supposer que $\mathbf{W}$ satisfait à tous les ``bons axiomes'' sur les localiseurs fondamentaux W (1, 2, 3, 4, 5, 6, 7) \texttt{[voir chapitre XII]}.}). La forme la plus forte serait celle-ci : si
\[ u \;:\; I' \longrightarrow I \]
\msp{page 56}est dans $\mathbf{W}$, on veut que le foncteur $u^{*}$ induise une \textit{équivalence}
\[ u_{\mathrm{lc}}^{*} \;:\; \mathrm{HO}_{W}^{\mathrm{lc}}(I) \xrightarrow{\;\sim\;} \mathrm{HO}_{W}^{\mathrm{lc}}(I'). \]
Mais je doute qu'on soit armé pour vérifier un axiome aussi fort, sans avoir au moins une théorie des opérations $u_{!}$ ou $u_{*}$. On pourrait à tout hasard poser cet axiome en prenant $\mathbf{W} = \mathbf{W}_{\alpha} \overset{\text{déf}}{=}$ le plus petit localiseur fondamental satisfaisant W (1 à 7) (le plus grand de
% foot: Version 25 janvier 2004
\printedpage{32}
% head: A. Grothendieck : Dérivateurs, Ch. XIII, Édition des Universités de Montpellier II et Paris VII
ces localiseurs étant $\mathbf{W}_{\infty}$). Pour l'instant, la façon de vérifier cet axiome (en termes de la donnée $(\mathcal{M},W)$) est pour moi un mystère $\pm$ total. (Mais je sais le vérifier pour $\mathrm{HOT}_{W}$, quand $W \subset \mathrm{Fl}(\mathrm{Cat})$ satisfait W (1, 2, 3, 4, 5, \dots).)

\dgrblock{dgrb709}{Localization axioms, stability, and locally constant coefficients}
\begin{dcomment}{exposition}{p55}{Localization axioms, stability, and locally constant coefficients}
Grothendieck now asks which parts of the derivator formalism are already forced by a good calculus of weak equivalences. Two separate gains should be distinguished. Localizing categories of diagrams preserves their coherent origin; obtaining homotopy Kan extensions requires additional functors between these localized categories. The prederivator exists before the second gain has been made.

Finite disjoint unions are formal. A diagram on $I_1\amalg\cdots\amalg I_r$ is a tuple of diagrams, and weak equivalences are componentwise. Localization commutes with this finite product: finite zigzags may be padded by identities, and the relations may be checked in one coordinate at a time. This explains the absence of a quillenisability hypothesis in the stated finite version of $\mathbf{Der}1$. The same argument is not an argument for arbitrary products of unrelated relative categories, since there may be no uniform finite zigzag length. No such infinite assertion is needed here.

For $\mathbf{Der}2$, the point is to represent a localized arrow by a genuine arrow after replacing its endpoints. Let $f$ be an arrow of $\mathrm{HO}_W(I)$ whose evaluations are all invertible. Under the stated homotopy calculus, choose a representative $a:A\to B$ in $\mathcal M(I)$, with the endpoint identifications weak equivalences. Every $a_i$ is invertible after localization in $\mathcal M$. Strong saturation of $W$ implies $a_i\in W$, so $a\in W(I)$ and $f$ is invertible. The source is correctly using both ingredients: componentwise detection in the original relative category and representability of arrows in the localized diagram category. Strong saturation alone would not supply the latter.

The word ``triangulée'' has a deliberately broader meaning in this passage than the Verdier definition. In a pointed homotopy theory one has cofiber sequences, fiber sequences, suspension and loop objects. Suspension need not be invertible, and there is no reason for the homotopy category to be additive. Pointed spaces give the basic example. An actual triangulated category emerges in the stable case, where suspension and loop are inverse equivalences and cocartesian squares agree with cartesian squares. For strong stable derivators, the canonical triangulations and exactness of the induced functors are established in \citet[Theorem 4.16 and Proposition 4.18]{Groth2013}. This is a later completion of the stable part of the programme, not an assumption hidden in the quoted scholium.

The locally constant condition on p.~56 is likewise a condition on a coherent diagram. Its transition arrows become isomorphisms after localization, but the object still belongs to $\mathrm{HO}_W(I)$, not merely to $\operatorname{Fun}(I^{\mathrm{op}},\mathrm{HO}_W(e))$. The higher coherence carries information. For ordinary model categories, the precise later answer is \citet[Theorem 1.3]{Cisinski2009LocallyConstant}: a functor inducing a weak equivalence of nerves induces an equivalence between the corresponding categories of locally constant diagrams. One can interpret this as dependence only on the infinity-groupoid obtained by inverting all arrows of the indexing category. It does not assert the result for every relative category merely satisfying the definition of quillenisability in \S5; the necessary homotopy theory must still be available.

The axiom numbers are historical. The present $\mathbf{Der}3$ concerns a localizer and locally constant objects, whereas $\mathbf{Der}4$ and $\mathbf{Der}5$ below concern existence and pointwise computation of right images. Modern texts commonly number the existence and pointwise axioms as Der3 and Der4, with strength as an additional Der5. A reference to a modern numbered axiom must therefore be translated by its content.
\end{dcomment}


\dgrblock{dgrb710}{Locally constant diagrams depend on the homotopy type}
\begin{dcomment}{later}{p56:Pour l'instant, la façon de vérifier cet axiome}{Locally constant diagrams depend on the homotopy type}
Grothendieck wants a homotopy-invariance theorem for coefficients, stronger than invariance of a single homotopy limit. A locally constant diagram sends every arrow of the indexing category to a weak equivalence. Its transition maps consequently remember transport along paths, together with the higher compatibility data carried by the original diagram before localization. This is why the whole homotopy type of the indexing category, rather than only its set of connected components, enters.

For an ordinary model category $\mathcal V$, \citet[Theorem~1.3]{Cisinski2009LocallyConstant} proves the precise statement anticipated here: if $N(u)$ is a weak homotopy equivalence, restriction induces an equivalence between the homotopy categories of locally constant $\mathcal V$-valued diagrams. Corollary~1.28 gives the converse when all coefficient model categories are allowed; simplicial sets already suffice. Apply the theorem to the opposite indexing categories to match the presheaf convention of these pages. This is an established answer for the usual localizer of nerve weak equivalences, not a conjecture awaiting the construction of $u_!$ or $u_*$.

In the language of higher categories, the reason is that a locally constant diagram factors through the groupoid completion of the index. That completion is the homotopy type represented by its nerve. An equivalence of those homotopy types induces an equivalence of their categories of functors into any fixed homotopy theory. The model-category theorem makes this explanation precise while retaining strict diagrams and their objectwise weak equivalences.

Two qualifications keep this comparison faithful to the text. The source's ``Der~3, nouvelle mouture'' is an axiom about local systems and a chosen fundamental localizer; it is not the modern axiom usually numbered (Der3), which asserts existence of homotopy Kan extensions. Also, a fixed coefficient theory may detect fewer distinctions between shapes than all coefficient theories together. Cisinski's converse uses the latter quantifier. We refer to the commentary on Chapter~VI, manuscript pp.~161--165, for the detailed comparison of those two forms of detection, rather than repeating it here.
\end{dcomment}


Je vais m'en occuper, une fois que j'aurai une maîtrise raisonnable sur le formalisme de $u_{*}$, $u_{!}$ relatif aux $u : I' \longrightarrow I$ dans Diag. L'hypothèse quillénisable seule semble trop grossière pour assurer l'existence de $u_{!}$ ou de $u_{\star}$ dans des cas intéressants. % [sic] u_\star -> u_*, star symbol where u_* is used elsewhere

\noindent\msp{page 57}

\section*{6\quad Catégorie de modèles quillénisables, et foncteurs $u_{!}$, $u_{*}$}

\textbf{Définition 1 :}

\begin{itemize}[leftmargin=2.5em]
\item[\circled{a}] On dit que $(\mathcal{M},W)$ est \textit{quillénisable uniformément à gauche} sur Diag, si on peut trouver une structure de Quillen faible $(C,TC,F,TF)$ pour $\mathcal{M}$ qui soit \textit{close,} et qui de plus ait la propriété suivante : pour $I$ dans Cat, soient $C(I)$, $TC(I) \subset \mathrm{Fl}(\mathcal{M}(I))$ définis par
\[ \begin{aligned} &(u : X \longrightarrow Y) \in C(I) \text{ (resp. } TC(I)) \text{ si et seulement si } \ \forall\, s \in \operatorname{Ob} I, \\ &u_{s} : X_{s} \longrightarrow Y_{s} \text{ est dans } C(I) \text{ (resp. } TC(I)). \end{aligned} \]
Soit de plus
\[ F(I) \;=\; TC(I)_{*}, \qquad TF(I) \;=\; C(I)_{*}. \]
La condition envisagée, c'est que pour tout $I$ \textit{dans} Diag, $(C(I),TC(I),F(I),TF(I))$ est une structure de Quillen faible \textit{close exacte} pour $W(I)$ (\footnote[34]{$(C,TC,F,TF)$ est appelé alors une \textit{quillénisation uniforme à gauche} de $(\mathcal{M},W)$ sur Diag.}).

\item[\circled{b}] On dit que $(\mathcal{M},W)$ est \textit{quillénisable sur Diag uniformément et exactement} \texttt{[à]} \textit{gauche}, si on a de plus $TC = C \cap W$, \textit{i.e.} si $(C,TC,F,TF)$ est exacte à gauche, donc défini par $C$ seul, via
\[ TC = C \cap W, \qquad F = (TC)_{*}, \qquad TF = C_{*} \qquad \text{(\footnotemark[35])} \]
\footnotetext[35]{On dit que $(C,\ldots)$, ou simplement $C$, définit une \textit{quillénisation uniforme et exacte à gauche} de $(\mathcal{M},W)$ sur Diag.}
\end{itemize}

On a vu au paragraphe 3 que si $\mathcal{M}$ est stable par petites $\varinjlim$ et accessible, et $W$ stable par petites $\varinjlim$ filtrantes, enfin si
% foot: Version 25 janvier 2004
\printedpage{33}
% head: A. Grothendieck : Dérivateurs, Ch. XIII, Édition des Universités de Montpellier II et Paris VII
\msp{page 58}$W$ admet une structure de Quillen faible $(C,TC,F,TF)$ qui soit close et exacte à gauche -- alors les conditions de la définition 1 sont satisfaites.

\dgrblock{dgrb711}{Why uniform left structures produce right images}
\begin{dcomment}{exposition}{p57}{Why uniform left structures produce right images}
Grothendieck's word ``uniform'' puts the same replacement mechanism on every diagram category. With cofibrations and trivial cofibrations defined objectwise, restriction preserves both classes for the elementary reason that it merely selects components. This makes restriction the left-hand member of the relevant Quillen adjunction, even though its geometric name is ``inverse image''. Its right adjoint is therefore the functor on which fibrant replacement should be performed.

Write $u:I'\to I$ and assume that the ordinary right Kan extension exists. The lifting adjunction immediately gives
\[
 u_*F(I')\subset F(I),\qquad u_*TF(I')\subset TF(I).
\]
For example, lifting a map of $\mathcal M(I)$ in $C(I)$ against $u_*p$ is equivalent to lifting its restriction, which lies in $C(I')$, against $p\in TF(I')$. The argument for $F$ uses $TC$ instead. Since $u_*$ preserves the terminal object, it also preserves fibrant objects. These are the facts summarized before the theorem-scholie; they do not yet appeal to a formula for a homotopy limit.

The pointwise formula determines which ordinary limits must exist:
\[
 (u_*X')(s)=\varprojlim_{(i',\,u(i')\to s)\in I'/s}X'(i').
\]
Here the displayed limit is the limit of the contravariant diagram on $I'/s$. A finite category has finitely many morphisms as well as finitely many objects, so these are finite limits when $I$ and $I'$ are finite. The phrase ``$\varinjlim$ finies'' on p.~59 is a source slip: the relevant half of $\mathbf M0$ is finite \emph{limits}. For an infinite comma category, neither $\mathbf M0$ nor the word quillenisable creates the missing limit.

There are two ways to continue. The manuscript isolates preservation of weak equivalences between fibrant diagrams as a new condition called admissibility, then proves it by finite gluing. This is a sound and instructive construction of finite homotopy limits. The closed exact hypotheses also permit the general Brown argument given below, in the block ``The Brown argument makes admissibility automatic''. It shows that admissibility is automatic whenever the ordinary right Kan extension exists. Consequently the remaining substantive issue is existence of appropriate diagram structures and ordinary limits, rather than an independent homotopy-invariance property of their right adjoints.

For comparison, a combinatorial model category has both the injective and projective model structures on every small diagram category. The injective structure is the modern counterpart of the objectwise-cofibration construction used here, and the projective structure of its dual. The existence statement and its derivator application are respectively Theorem 1.28 and Proposition 1.30 of \citet{Groth2013}. The original finite-limit formulation is still useful when arbitrary small limits do not exist, as in bounded complexes; it then yields only the Kan extensions for which the underlying limit problem can be solved.
\end{dcomment}


D'autre part, notons que la proposition 1 page 45 et ses corollaires, y inclus la proposition 3 (page 50) et ses corollaires, sont valables dès lors que les quillénisations des $(\mathcal{M}(I),W(I))$ envisagés sont définies ``uniformément à gauche'' par une quillénisation uniforme à gauche (pas nécessairement exacte à gauche) de $(\mathcal{M},W)$ sur Diag.

\dgrblock{dgrb712}{Closure and exactness already imply both trivial-class equalities}
\begin{dcomment}{exposition}{p58:D'autre part, notons que la proposition}{Closure and exactness already imply both trivial-class equalities}
Grothendieck introduces uniform quillenization to remove the dependence of the cofibration classes on the shape. The exact definition on p.~57 has a useful consequence which shortens the later discussion. In each diagram category it asks for a structure which is both \emph{closed} and \emph{exact}. Together these hypotheses force the two equalities which define biexactness.

Indeed, suppose $i\in C\cap W$. Exactness factors this very map as
\[
 i=pj,\qquad j\in TC,\quad p\in TF.
\]
Since $i\in C$ and $p\in TF$, the square with top map $j$ and bottom map the identity has a lift. The usual retract argument makes $i$ a retract of $j$ in the arrow category. Closedness of $TC$ gives $i\in TC$. Thus $C\cap W=TC$. Dually, factoring $p\in F\cap W$ as $qj$ with $j\in TC$, $q\in TF$ makes $p$ a retract of $q$, so $F\cap W=TF$. These are applications of the ordinary retract argument, as in \citet[Lemma~1.1.9]{Hovey1999}; no transfinite construction is involved.

Consequently a closed exact weak structure is perfect in the manuscript's terminology. In particular, if the domain contains the terminal category, as stipulated on p.~51, Definition~1(a) already implies the additional equality in Definition~1(b). The advertised distinction between uniform quillenizability and uniform \emph{exact} left quillenizability is therefore redundant under the printed hypotheses. This is a logical consequence of those hypotheses, not a proposed emendation of the French text. If ``exact'' were removed from Definition~1(a), there would be a different and genuinely weaker notion, but that is not the definition supplied here.

The formula in Definition~1(a) has a separate typographical error: a component $u_s:X_s\to Y_s$ belongs to $C$ or $TC$, not to $C(I)$ or $TC(I)$. The latter classes contain natural transformations, so the printed condition is not well-typed. The correction ledger records the error as a slip of the edition and preserves the printed text. After reading the component condition correctly, the uniform structures are precisely injective model structures with the same underlying cofibration classes, subject here also to the properness axiom M$''$ and to the available finite or small limits.
\end{dcomment}


La quillénisation uniforme gauche nous servira à établir l'existence des images directes homotopiques $u_{*}$, pour certains $u$. La quillénisation uniforme droite, dualement, aura comme conséquence principale l'existence des $u_{!}$ homotopiques pour certains $u : I' \longrightarrow I$ dans Diag.

Je vais résumer ce que j'ai pu établir jusqu'à présent dans ce sens, dans le

\msp{page 59}\textbf{Théorème-scholie :} \textit{Soit $(C,TC,F,TF)$ une quillénisation uniforme gauche de $W$ sur} Diag $\subset$ Cat.

\begin{itemize}[leftmargin=3em]
\item[1)] \textit{Soit $u : I' \longrightarrow I$ une flèche dans} Diag. \textit{On suppose} Diag \textit{accordé à $\mathcal{M}$ de façon que pour $u$, flèche dans} Diag, \textit{les $\varprojlim$ dans $\mathcal{M}$ (de type $I'/i$) requises pour assurer l'existence de}
\[ u_{*} \;:\; \mathcal{M}(I') \longrightarrow \mathcal{M}(I) \]
\textit{existent dans $\mathcal{M}$. (C'est O.K. si $I$, $I'$} \texttt{[sont]} \textit{finies, $\mathcal{M}$ étant supposée stable par $\varinjlim$ finies de toutes façons.) On sait} (prop. 1, page 45) \textit{que $u_{*}$ induit}
\[ u_{*}^{\mathrm{f}} \;:\; \mathcal{M}(I')_{\mathrm{f}} \longrightarrow \mathcal{M}(I)_{\mathrm{f}}\,. \tag*{$(*)$} \]
\textit{On dit que $u$ est} admissible (à gauche) pour la quillénisation uniforme gauche donnée de $W$ sur Diag, \textit{si $u_{*}^{\mathrm{f}}$ est compatible avec les localiseurs $W(I')_{\mathrm{f}}$, $W(I)_{\mathrm{f}}$,} i.e. \textit{si $f' : X' \longrightarrow Y'$ dans $\mathcal{M}(I')$, $X'$, $Y'$ fibrants, $f' \in W(I')$ implique $u_{*}(f') \in W(I)$. On dit que $Z$ dans} Diag \textit{est} admissible (pour la quillénisation uniforme gauche donnée de $W$ sur Diag \dots) \textit{si le foncteur $Z \longrightarrow e$ est admissible. Ceci posé :}
\begin{itemize}[leftmargin=2.5em]
\item[a)] \textit{Pour que $u$ soit admissible, il suffit que les $I'/s$ pour $s \in \operatorname{Ob} I$, le soient. NB il y a une réciproque, quand $I$ est ordonné, donc les $I'/s \longrightarrow I'$ des immersions ouvertes.}
\item[b)] \msp{page 60}\textit{Tout ensemble ordonné fini est admissible. Plus généralement (sauf erreur) tout ensemble ordonné $I$ tel que l'espace topologique associé à $I^{\circ}$ soit noethérien (}i.e. \textit{toute suite croissante d'ouverts de $I^{\circ}$ est stationnaire, ou encore toute suite croissante de fermées de $I$ est stationnaire), et si de plus toute partie fermée irréductible de $I^{\circ}$ admet un point générique,} i.e. \textit{$I^{\circ}$ est (non seulement un espace noethérien, mais aussi) un espace topologique} sobre. \textit{(Pour $I$, cela signifie que toute partie ouverte filtrante croissante de $I$ admet un plus grand élément.)} % [sic] fermées -> fermés, agreement with masculine (ouverts / fermés)
\end{itemize}
\end{itemize}
% foot: Version 25 janvier 2004


\dgrblock{dgrb713}{The Brown argument makes admissibility automatic}
\begin{dcomment}{later}{p59}{The Brown argument makes admissibility automatic}
The purpose of admissibility is to ensure that a limit of a suitably fibrant diagram depends only on its homotopy class. Under the closed exact hypotheses actually imposed in Definition 1, this is already a consequence of the lifting and factorization axioms. We give the argument explicitly, since it both strengthens the theorem-scholie and settles the later questions without assuming the existence of limits which the source has not provided.

First recall the consequence of closedness and exactness proved above: $F\cap W=TF$. Let $T:\mathcal C\to\mathcal E$ be a functor, where $\mathcal C$ has the source's closed exact structure and $\mathcal E$ has a class $V$ containing the identity maps and satisfying two-out-of-three. Assume that $T$ carries trivial fibrations between fibrant objects into $V$. If $f:X\to Y$ is a weak equivalence between fibrant objects, factor its graph as
\[
 X\xrightarrow{i} Z\xrightarrow{p}X\times Y,
 \qquad i\in TC,\quad p\in F,\quad pi=(1_X,f).
\]
Put $p_X=\operatorname{pr}_Xp$ and $p_Y=\operatorname{pr}_Yp$. The two projections are fibrations because $X$ and $Y$ are fibrant. Consequently $Z$ is fibrant and $p_X,p_Y\in F$. The identities $p_Xi=1_X$ and $p_Yi=f$, together with $i,f\in W$, imply $p_X,p_Y\in W$. Both maps therefore lie in $TF$. Their images under $T$ belong to $V$. Since $T(p_X)T(i)=1$, two-out-of-three first gives $T(i)\in V$ and then gives
\[
                         T(f)=T(p_Y)T(i)\in V.
\]
This is the dual form of Ken Brown's lemma \citep[Lemma 1.1.12]{Hovey1999}. We have displayed its proof to identify the exact hypotheses used: finite products, factorization, composition and pullback stability of $F$, two-out-of-three, and the equality $TF=F\cap W$. There is no appeal to smallness, accessibility, or filtered colimits.

Apply this to $T=u_*$. Uniformity makes $u^*$ preserve $C(I)$, so the lifting adjunction makes $u_*$ preserve $TF(I')$. It also preserves fibrant objects. The lemma gives precisely
\[
 f\in W(I'),\quad X',Y'\in\mathcal M(I')_{\mathrm f}
 \quad\Longrightarrow\quad u_*f\in W(I).
\]
Thus every functor $u$ for which the ordinary right Kan extension exists is admissible under the stated uniform closed exact assumptions. The finite-poset theorem remains correct, but its restriction on shapes is unnecessary for this particular homotopy-invariance assertion.

The conclusion should not be enlarged in a different direction. It does not manufacture an ordinary right adjoint in a category lacking the required limits. Nor does it, by itself, prove the pointwise homotopy formula: that formula additionally needs fibrant restriction to the comma categories, as discussed at p.~72. These are genuine existence and compatibility questions. The question of preserving weak equivalences \emph{after the relevant fibrant replacements exist} is answered by the short argument above. We attribute it to the established Brown lemma, rather than present the application as a new result.
\end{dcomment}


\dgrblock{dgrb714}{Every available right Kan extension is admissible}
\begin{dcomment}{later}{p60}{Every available right Kan extension is admissible}
Grothendieck's admissibility condition asks for the exact homotopical property needed to derive $u_*$. Under the uniform hypotheses actually printed on p.~57, it holds for every functor for which the ordinary right Kan extension exists. There is no restriction to finite posets. The reason is the dual of Ken Brown's lemma, whose standard statement is \citet[Lemma~1.1.12]{Hovey1999}.

We give the short proof in the present notation. By the preceding retract argument, each uniform structure is perfect. Proposition~1 says that $u_*$ preserves $F$ and $TF$; as a right adjoint it preserves the terminal object, and hence fibrant objects. Let $f:X\to Y$ be a weak equivalence between fibrant diagrams. Factor its graph as
\[
 X\xrightarrow{j}Z\xrightarrow{q}X\times Y,
 \qquad j\in TC,\quad q\in F.
\]
The composites $q_X:Z\to X$ and $q_Y:Z\to Y$ are fibrations because $X$ and $Y$ are fibrant. Their identities $q_Xj=1_X$ and $q_Yj=f$, together with $j,f\in W$, show that both belong to $W$. Biexactness gives $q_X,q_Y\in TF$. Their images under $u_*$ therefore lie in $W$. The equality $(u_*q_X)(u_*j)=1$ and two-out-of-three show $u_*j\in W$; hence $u_*f=(u_*q_Y)(u_*j)\in W$. This is precisely admissibility.

Choosing a fibrant replacement $X\to RX$ now defines
\[
 \mathbf R u_*[X]=[u_*RX].
\]
Brown's lemma ensures independence of the replacement in the localized category. The adjunction with the derived left functor of $u^*$ is the usual derived Quillen adjunction; since restriction already preserves all objectwise weak equivalences, that left derived functor is $u^*$ itself. Thus $u^*\dashv\mathbf R u_*$. The dual argument with uniformly right-defined structures gives $\mathbf L u_!\dashv u^*$. This is exactly the construction developed in \citet[Definition~1.3.6 and Lemma~1.3.10]{Hovey1999}.

The pointwise criterion in part~(a) remains useful for calculations. In the right Kan formula, the projection $I'/s\to I'$ is a discrete fibration. When the requisite fibre coproducts exist, in particular for finite indexing categories, the left Kan extension of presheaves along such a projection is computed by those coproducts: each relevant comma component has an initial object obtained by lifting the arrow to the fibre. It therefore preserves $TC$, so restriction preserves fibrant diagrams. The admissibility of the slice then makes the pointwise limit preserve weak equivalences. If $I$ is a poset, the slice projection identifies its domain with a sieve, as the author says. Under the full uniform hypotheses, the global Brown argument has already established both sides of the proposed converse, without requiring these auxiliary coproducts.

This proof uses finite products and factorizations in the diagram categories, plus the existence of the indicated right adjoint. The parenthesis on p.~59 should invoke finite \emph{limits}; it prints finite colimits. Both are included in M~0, so this slip does not invalidate the finite-index assertion. Up to manuscript p.~60 the theorem is only stated; the proof above supplies its admissibility conclusion independently of the continuation.
\end{dcomment}


\dgrblock{dgrb715}{The noetherian sober poset condition is finite}
\begin{dcomment}{exposition}{p60}{The noetherian sober poset condition is finite}
Grothendieck tries to enlarge the finite-poset case by replacing finite induction with topological noetherian induction. For the usual Alexandrov topology attached to a poset, however, the two hypotheses he writes down already force finiteness. The observation concerns this topology, in which arbitrary intersections of open sets are open; it would not be true for arbitrary noetherian sober spaces.

To make the order convention immaterial, let $P$ be the specialization poset and take its Alexandrov opens to be the upper sets. Closed sets are lower sets. An irreducible nonempty closed set is then exactly a directed lower set. For one direction, if $D$ is directed and meets two upper open sets, a common upper bound in $D$ lies in their intersection. Conversely, if a lower set $D$ is irreducible and $x,y\in D$, the open sets $\mathord\uparrow x$ and $\mathord\uparrow y$ both meet $D$, so their intersection meets $D$; that is a common upper bound in $D$. Sobriety says that each such directed lower set is $\mathord\downarrow z$ for a unique $z$.

It follows that $P$ has no strictly increasing infinite sequence. Otherwise, from $x_0<x_1<\cdots$ form the directed lower set $D=\bigcup_n\mathord\downarrow x_n$. If $D=\mathord\downarrow z$, then $z\leq x_N$ for some $N$, whereas $x_{N+1}\leq z$, a contradiction. This is the sobriety criterion of \citet[Chapter~II, Proposition~6.4.4]{PicadoPultr2012}.

Every nonempty subset of $P$ consequently has a maximal element: starting from an element with no maximal element above it would produce such an increasing sequence. A maximal element $m$ of a subset $S$ is isolated in $S$, since $(\mathord\uparrow m)\cap S=\{m\}$. The Alexandrov space is therefore scattered, meaning that every nonempty subspace has an isolated point.

Now use noetherianity only once. Among the finite open subsets choose a maximal one, $U$; it exists because a strictly increasing sequence of finite open subsets is forbidden. If the complement $P\setminus U$ were nonempty, scatteredness would give an isolated point $x$ in that subspace. There would be an open $V$ of $P$ with $V\cap(P\setminus U)=\{x\}$. Then $U\cup V=U\cup\{x\}$ would be a strictly larger finite open set, contradicting maximality. Hence $P=U$ is finite.

Applied to $I^{\mathrm{op}}$, the proof shows that the class on p.~60 consists exactly of finite posets. Both chain conditions are used: the order condition controls sobriety, and topological noetherianity controls opens. The sequence arguments use classical choice.
\end{dcomment}
\begin{otherlanguage}{english}
\dgroverview{XIII-a}{Chapter XIII \textperiodcentered\ \textup{XIII-a} \textperiodcentered\ Overview}
\selectlanguage{english}

\noindent\textbf{Chapter XIII, first part: building the machinery that can later be forgotten.}

These sixty manuscript pages turn from the intrinsic calculus of a
localizer to the auxiliary structures that make it calculable. Grothendieck
calls a pair $(\mathcal M,W)$ a category of models when $W$ contains the
isomorphisms and satisfies two-out-of-three. This is much less than what
``model category'' ordinarily means today. He seeks classes of cofibrations
and fibrations that provide replacements, homotopy fibres and cofibres, and
eventually the two derived Kan extensions. Their purpose is to establish
properties of the localized diagram categories
\[
 \operatorname{HO}_W(I)
   =\mathcal M^{I^{\mathrm{op}}}[W(I)^{-1}],
\]
where $W(I)$ is defined objectwise. Once these constructions and their
universal properties have been established, the auxiliary classes should
cease to be part of the intrinsic answer
\citep[manuscript pp.~12--14]{Grothendieck1990XIII}.

The chapter opens with the algebra of lifting properties. In current
notation, Grothendieck's $\Phi_*$ is $\Phi^{\square}$, his $\Psi^*$ is
${}^{\square}\Psi$, and his ``couple de Quillen'' is a weak factorization
system. It is not what is now called a Quillen adjunction. The Galois
connection explains why left classes are stable under pushouts, retracts,
and continuous transfinite compositions, and why right classes have the
dual properties. The small-object theorem then supplies factorizations
from a set of generating arrows in a locally presentable category. Our
commentary proves the construction with its cardinal bound and separates
the cellular class from its retract closure. This distinction matters
because factorization can exist before the two classes are closed under
retracts.

The intrinsic classes $\operatorname{Cof}_W$ and
$\operatorname{Fib}_W$ express preservation of weak equivalences under
cobase change and base change. They are not defined by a lifting property.
Their relation to the universally trivial arrows is established using the
factorization hypotheses, not just two-out-of-three. The base-change
argument on manuscript pages~10--11 is particularly economical: factor the
change of base, apply the fibration condition along its weak-equivalence
part, and use two-out-of-three. The commentary expands this proof and gives
a three-element-poset example showing why the factorization assumption
cannot simply be omitted.

The conjectural theorem of manuscript pages~14--26 asks whether the
intrinsic classes contain small collections of tests sufficient to build
the desired structures. Here three different forms of smallness enter:
closure under filtered colimits, accessibility of a full subcategory of the
arrow category, and generation of a lifting class by a set. They do not
mean the same thing. In particular, even an accessible left lifting class
can fail to be generated by a set in the ordinary small-object sense; the
order-embedding example is worked out in the commentary. Conversely, one
specific conjecture in these pages has a precise later answer: intersections
of two set-generated left classes in a locally presentable category are
again set-generated, by the limit theorem for combinatorial cellular
categories \citep[Theorem~3.2 and Remark~3.8]{MakkaiRosicky2014Cellular}.

Smith's recognition theorem gives a useful repair of the larger programme.
One fixes a set of generating cofibrations, assumes that its right
orthogonal consists of weak equivalences, requires the trivial cofibrations
to be stable under pushouts and transfinite composition, and imposes a
solution-set condition on $W$. Under these hypotheses one obtains an
ordinary cofibrantly generated model structure
\citep[Theorem~1.7]{Beke2000}. With the intrinsic inclusions envisaged here,
the resulting structure is proper. This is a theorem with explicit extra
hypotheses, not a proof that the manuscript's filtered-colimit assumption
alone yields every assertion in its conjectural equivalence. The audit
records the individual implications and their dependencies rather than
assigning one status to the whole discussion.

The second formulation, beginning on manuscript page~27, isolates a
different issue. A weak Quillen structure is called exact when
$W=TF\circ TC$; it is left exact when $TC=C\cap W$ and right exact when
$TF=F\cap W$. These distinctions are meaningful for classes that have not
been closed under retracts. The marginal claim that retract closure of
$C,F,W$ alone makes the two one-sided exactness conditions equivalent is
false. Our example in sets takes $C=TC$ to be the monomorphisms, $F$ the
surjections, $W$ all maps, and $TF$ the surjections whose fibres are singletons
or infinite. It satisfies the weak axioms and left exactness, but excludes
the map from a two-element set to a singleton from $TF$. The missing
retract condition is visible directly.

Exactness nevertheless permits a useful completion without any retract
closure assumption. If $j\in C\cap W$, factor $j=qi$ with $i\in TC$ and
$q\in TF$. The lifting axiom makes $j$ a retract of $i$, so $j$ has the left
lifting property against every fibration. The dual argument applies to
$F\cap W$. Replacing the two designated trivial classes by $C\cap W$ and
$F\cap W$ therefore gives a biexact Quillen structure with the same
$C,F,W$. Properness supplies the required stability under base and cobase
change. This is a completion of the designated classes; the counterexample
shows why it cannot be replaced by an assertion that they were equal
already. When all four original classes are closed, the completion makes
no change, and a \emph{closed exact} weak structure is already perfect.
We use this observation again at the end of these pages to remove an
unnecessary admissibility restriction on right Quillen functors.

The discussion of homotopy on manuscript pages~29--34 concerns the actual
morphisms in the localization. Replacing objects by objects that are both
cofibrant and fibrant is not yet the same as proving that all localized
morphisms are represented by single arrows. Exactness supplies such
representatives. We first prove a description by the congruence generated
jointly by left and right homotopies. The completion then identifies each
of the usual homotopies with equality after localization for a cofibrant
source and a fibrant target. It also resolves the doubts on manuscript
pages~32--34: homotopy-cartesian and homotopy-cocartesian squares are exact
under the weak exact axioms themselves. An endpoint-lifting argument makes
the first assertion explicit; the extra hypothesis about the diagram
category used in the source's proposed proof is unnecessary for this
conclusion. Pointedness gives the suspension and loop sequences. Stability
is the additional condition needed for a triangulated structure. The
separate suggestion about raw right homotopy depends on the path-object
convention. A graph example, adapted from \citet[Example~3.6]{KurzRosicky2005},
shows that paths supplied by the designated factorization system need not
give a transitive relation without exactness, even on fibrant-cofibrant
objects. Allowing some weak-equivalence--fibration path object gives an
equivalence relation for a fibrant target by the concatenation argument
proved in the commentary.

The passage to presheaves supplies the chapter's central technical test.
Objectwise cofibrations lead to an injective diagram structure. Its
fibrations are defined by lifting against all objectwise trivial
cofibrations, and need not be merely objectwise fibrations. The calculation
\[
 (\alpha_{s!}A)(t)=\coprod_{\operatorname{Hom}_I(t,s)} A
\]
explains the coproduct hypothesis in the evaluation lemma. The subsequent
failure of componentwise factorization to solve natural lifting problems
is exhibited by explicit examples. Modern existence results for injective
structures apply in the combinatorial setting, while the cellular-category
theorem also accounts for the set-generation of the pointwise left
classes \citep[Corollary~3.4]{MakkaiRosicky2014Cellular}.

The variances then become transparent. Restriction preserves the
objectwise cofibrations and trivial cofibrations, so its right adjoint is
right Quillen. Its left adjoint is not automatically left Quillen for these
same injective structures. Above all, an ordinary Kan extension and its
derived functor are different constructions. Grothendieck explicitly
withdraws the premature identification on manuscript pages~49--50. The
commentary supplies counterexamples and the appropriate derived statements.
Under the uniform closed exact hypotheses finally adopted on manuscript
page~57, every existing ordinary right Kan extension preserves weak
equivalences between injectively fibrant diagrams by Brown's lemma.
Consequently the final admissibility conclusion holds for every such
indexing functor, beyond the finite posets singled out in the text.

Complexes give a complementary warning about size. For an infinite
indexing category, a diagram of individually bounded-below complexes need
not have a common lower bound. It therefore need not be a bounded-below
complex in the abelian category of diagrams. A family with nonzero
cohomology in successively lower degrees makes the resulting failure of a
product in $D^+(\mathbf{Ab})$ explicit. Unbounded derived operations in
Grothendieck abelian categories are instead available through suitable
resolutions and model structures; the boundedness issue concerns the
chosen subcategory and cannot be removed by renaming it.

The later derivator constructions realize the intended independence from
presentation. Cisinski supplies homotopy right Kan extensions for general
closed model categories with small limits
\citep[Theorem~6.11]{Cisinski2003}; Renaudin compares the localized
2-category of suitable combinatorial presentations with derivators of
small presentation \citep[Theorems~3.3.2 and~3.4.4]{Renaudin2009}; and the
infinity-categorical treatment recovers the same pointwise comparison
laws from Kan extensions in the underlying infinity-category
\citep[Proposition~5.3.9 and Remark~5.3.10]{RiehlVerity2017}. These results
give precise answers in their respective settings. The final
order-topological extension also becomes precise: a noetherian sober
Alexandrov $T_0$ space is finite, so under the ordinary poset topology the
apparently broader class of shapes on manuscript page~60 adds no new cases.
For the weakest recognition assertions of the manuscript, what was
established here is kept apart from what was not located in the checked
literature. An unlocated proof is not certified as an open problem.
\end{otherlanguage}
\dgrchunkend

% DER_XIII-b: Grothendieck's text with the commentary set in, then the overview, inside a group that applies the chunk's own definitions
\dgrpart{1}{XIII-b}{Part 2 of 2, manuscript pages 61--98}{XIII-b\enspace\textperiodcentered\enspace Chapter XIII, part 2 of 2\enspace\textperiodcentered\enspace manuscript pages 61--98}{Chapter XIII \textperiodcentered\ \textup{XIII-b}}
\dgrchunkbegin
% DER_XIII-b: the definitions of the chunk's preamble, applied inside its group
\let\printedpage\relax\newcommand{\printedpage}[1]{\marginpar{\footnotesize p.~#1}}
\let\msp\relax\newcommand{\msp}[1]{\marginpar{\footnotesize [#1]}}
\let\circled\relax\newcommand{\circled}[1]{\tikz[baseline=(c.base)]\node[draw,circle,inner sep=1pt](c){#1};}
\let\cM\relax\newcommand{\cM}{\mathcal{M}}
\let\HO\relax\newcommand{\HO}{\mathrm{HO}}
\let\Hom\relax\newcommand{\Hom}{\mathrm{Hom}}
\let\Ob\relax\newcommand{\Ob}{\mathrm{Ob}}
\let\rf\relax\newcommand{\rf}{\mathrm{f}}
\let\rc\relax\newcommand{\rc}{\mathrm{c}}
\let\M\relax\newcommand{\M}{\mathcal{M}}
\let\Fl\relax\newcommand{\Fl}{\mathrm{Fl}\,}
\let\Filt\relax\newcommand{\Filt}{\mathrm{Filt}}
\let\card\relax\newcommand{\card}{\mathrm{card}\,}
\let\barlongleftrightarrow\relax\newcommand{\barlongleftrightarrow}{\mathrel{\text{\ooalign{$\longleftrightarrow$\cr\hidewidth$\mid$\hidewidth\cr}}}}
\let\idx\relax\newcommand{\idx}[2]{\noindent #1\dotfill #2\par}
\let\entry\relax\newcommand{\entry}[2]{\noindent #1\dotfill\ #2\par}
\let\Mo\relax\newcommand{\Mo}{\mathcal{M}}
\let\Mzg\relax\newcommand{\Mzg}{$\mathbf{M}'_{\mathbf{0}}\mathbf{(g)}$}
\let\Mzd\relax\newcommand{\Mzd}{$\mathbf{M}'_{\mathbf{0}}\mathbf{(d)}$}
\providecommand\thefootnote{}\renewcommand{\thefootnote}{\arabic{footnote}}

% ---- from Der-ChXIII-25-01-04_p036-040.tex
\printedpage{34}
% head: A. Grothendieck : Dérivateurs, Ch. XIII, Édition des Universités de Montpellier II et Paris VII

\begin{itemize}[leftmargin=3em]
\item[c)] \textit{En particulier, si $I'$, $I$ sont des ensembles ordonnés finis, alors tout $u : I' \longrightarrow I$ est admissible. Il en est de même, plus généralement, si $I$, $I'$ sont des ensembles ordonnés et $I'^{\circ}$ est un espace noethérien sobre.}
\end{itemize}

\begin{itemize}[leftmargin=2em]
\item[2)] \textit{Supposons $u : I' \longrightarrow I$ admissible. Considérons}
\[
u_*^W \ : \ \underbrace{\cM(I')_{\rf}W(I')_{\rf}^{-1}}_{\simeq\, \HO_W(I')} \ \longrightarrow \ \underbrace{\cM(I)_{\rf}W(I)_{\rf}^{-1}}_{\HO_W(I)},
\]
\textit{induit par $u_*$. Ce foncteur est adjoint à droite de $u_W^* : \HO_W(I) \longrightarrow \HO_W(I')$.}

\msp{page 61}
\textit{De plus, l'homomorphisme canonique (pour $s \in \Ob I$, $X' \in \Ob \HO_W(I')$)}
\[
u_*^W(X')(s) \ \simeq \ p_{s\,*}(X'|I'/s) \ \overset{\text{déf}}{=} \ \underset{I'/s}{\underleftarrow{\mathrm{holim}}}{}_W(X'|I'/s)
\]
\textit{(où $p_s : I'/s \longrightarrow e$) est un isomorphisme, si on suppose que les $I'/s$ sont admissibles (on dira alors que $u$ est} strictement \textit{admissible).}
\end{itemize}

\textbf{Corollaire.} \textit{Sur le domaine formé des ensembles ordonnés finis, le prédérivateur} $\HO_W$ \textit{satisfait (en plus des axiomes} \textbf{Der}~1, \textbf{Der}~2\textit{) aux axiomes} \textbf{Der}~4 \textit{et} \textbf{Der}~5 \textit{(existence des images directes $u_*$, adjoints à droite des foncteurs $u^*$, et calcul standard des fibres d'une image directe).}

\textsc{Démonstration} du théorème-scholie.

$1^\circ$) \textit{a}) est immédiat. Pour \textit{b}), on procède par récurrence sur le cardinal $n$ de l'ensemble ordonné fini $I$, les cas $n \in \{0,1\}$ étant triviaux. Supposons donc $n \geq 2$, et le théorème prouvé pour les $n' < n$. Je distingue deux cas : si $I$ admet un plus grand élément $a$, alors, désignant par $\alpha : e \longrightarrow I$ le foncteur de valeur $a$, par $p : I \longrightarrow e$ la projection, on sait que l'on a
\[
p_* \ \simeq \ \alpha^*, \qquad \textit{i.e.} \ \varprojlim_{I} X_i \simeq X_a,
\]
d'où le fait que sur $\cM(I)$ tout entier
\msp{page 62}
(pas seulement sur $\cM(I)_{\rf}$), $p_*$ est compatible aux localiseurs $W(I)$, $W$.

\dgrblock{dgrb716}{Finite limits assembled from a two-piece cover}
\begin{dcomment}{exposition}{p61}{Finite limits assembled from a two-piece cover}
Grothendieck's finite-poset induction makes the geometry of a homotopy limit visible. Instead of treating a large compatibility diagram at once, he cuts the indexing poset into two lower sets and glues the two resulting limits over their overlap. This argument remains useful even though the Brown lemma gives a shorter proof of admissibility under the stronger standing hypotheses.

Let $I=I_1\cup I_2$ with $I_1,I_2$ lower subsets, and put $I_0=I_1\cap I_2$. Define $v:I\to\Psi=(1\leftarrow0\to2)$ by sending the overlap to $0$, the points belonging only to $I_1$ to $1$, and those belonging only to $I_2$ to $2$. The lower-set condition makes this assignment order preserving: an arrow cannot travel from one exclusive part into the other. Its comma categories over $0,1,2$ are precisely $I_0,I_1,I_2$. Right Kan extension along $v$ therefore produces the span
\[
 \varprojlim_{I_1}X\longrightarrow\varprojlim_{I_0}X
 \longleftarrow\varprojlim_{I_2}X.
\]
Taking its pullback gives $\varprojlim_I X$. Indeed, a cone to $X$ is exactly a pair of cones to its restrictions whose two cones on $I_0$ agree. This proves the ordinary limit identity by its universal property, including the cases in which the overlap is empty.

For a finite nonempty poset without a greatest element, partition its maximal elements into two nonempty sets $M_1,M_2$ and take $I_j=\bigcup_{m\in M_j}\{i:i\leq m\}$. Every point lies below a maximal element, so the lower sets cover $I$, and neither is all of $I$. Their overlap is also smaller. The inductive argument is therefore well founded. If $I$ has a greatest element $a$, the limit is $X_a$ in the source's contravariant convention; the empty poset has the terminal object as its limit.

The homotopical step is to ensure that the span just constructed is fibrant. Right Kan extension preserves fibrancy, and a fibrant $\Psi$-diagram consists of fibrant objects $A_0,A_1,A_2$ with both arrows $A_j\to A_0$ fibrations. Its ordinary pullback then represents its homotopy pullback. The induction replaces ``take an arbitrary pullback'' by ``take a pullback whose legs satisfy a lifting condition''. The latter is the operation stable under weak equivalences of the input span.

Composition of admissible functors is also transparent at this point. If $u$ and $v$ preserve weak equivalences between fibrant diagrams after right Kan extension, then $u_*$ first yields another fibrant diagram, to which $v_*$ can be applied. The canonical identification $(vu)_*\cong v_*u_*$ proves admissibility of the composite. No new comparison of replacements is required. This explains why the elementary shape $\Psi$ controls the whole finite induction.
\end{dcomment}


Si $I$ n'a pas de plus grand élément, alors l'ensemble $M$ des ses éléments maximaux est de cardinal $\geq 2$, on peut donc l'écrire $M = M_1 \amalg M_2$, et on aura % [sic] des ses -> de ses, typographical slip
\[
I \ = \ I_1 \cup I_2, \qquad \text{où } I_1 = I/M_1,\ I_2 = I/M_2
\]
sont deux ouverts de $I$, différents de $I$, donc de cardinaux $n_1$, $n_2 \ < \ n$. Considérons aussi $I_0 = I_1 \cap I_2$, c'est un ouvert de cardinal $< n$. Le diagramme
\[
\begin{tikzcd}[row sep=small]
 & I_1 \arrow[dr, dotted] & \\
I_0 \arrow[ur] \arrow[dr] & & I = I_1 \cup I_2 \\
 & I_2 \arrow[ur, dotted] &
\end{tikzcd}
\]
% foot: Version 25 janvier 2004
\printedpage{35}
% head: A. Grothendieck : Dérivateurs, Ch. XIII, Édition des Universités de Montpellier II et Paris VII
correspond à une flèche
\[
\begin{tikzcd}[row sep=small, column sep=small]
 & & & 1 & \\
u \ : \ I \arrow[r] & \Psi = 0 \arrow[urr, start anchor=east, end anchor=south west] \arrow[drr, start anchor=east, end anchor=north west] & & & , \\
 & & & 2 &
\end{tikzcd}
\]
avec
\[
I_1 \ = \ I/1, \qquad I_2 \ = \ I/2, \qquad I_0 \ = \ I/0.
\]
L'hypothèse de récurrence et \textit{a}) impliquent que $u$ est admissible. Comme un composé de foncteurs admissibles est admissible (tautologiquement, via $(vu)_* \simeq v_*u_*$), on est réduit finalement au
\msp{page 63}

\textbf{Lemme.} \textit{La catégorie} \begin{tikzcd}[row sep=small, column sep=small, baseline=(\tikzcdmatrixname-2-1.base)]
 & 1 \\
\Psi = 0 \arrow[ur] \arrow[dr] & \\
 & 2
\end{tikzcd} \quad \textit{est admissible.}

Ceci va résulter du

\textbf{Corollaire} (\footnote[36]{\textbf{NB} Il y a une variante évidente, caractérisant les $X$ tels que $X \longrightarrow e$ soit dans $TF(\Psi)$ comme ceux pour lesquels $X_0$, $X_1$, $X_2$ sont fibrants \textit{triviaux,} et $\alpha_1$, $\alpha_2 \in TF$.}). \textit{Soit} \begin{tikzcd}[row sep=small, column sep=small, baseline=(\tikzcdmatrixname-2-1.base)]
 & X_1 \arrow[dl, "\alpha_1"'] \\
X_\bullet = X_0 & \\
 & X_2 \arrow[ul, "\alpha_2"]
\end{tikzcd} \quad \textit{dans $\cM(\Psi)$. Pour que $X_\bullet$ soit fibrant, il faut et il suffit que l'on ait}

\begin{itemize}[leftmargin=3em]
\item[a)] \textit{$X_0$, $X_1$, $X_2$ sont fibrants, et}
\item[b)] \textit{$\alpha_1$, $\alpha_2 \in F$.}
\end{itemize}

La nécessité de \textit{a}) est cas particulier du corollaire 2, page 48, celle de \textit{b}) est cas particulier du corollaire 1 page 50. Pour la suffisance (dont on n'aura pas besoin pour le théorème), on sait montrer que dans le diagramme commutatif à traits pleins
\[
% Audit resolved: PDF leaf 7 confirms interrupted vertical sides A_0 -> X_0 and A_2 -> X_2, obscured by the displaced B labels; the printed geometry is retained. See corrections.csv.
\begin{tikzcd}
A_1 \arrow[r] \arrow[dr, "TC"] \arrow[dd] & A_0 \arrow[dr, "TC"] \arrow[d, dash] & A_2 \arrow[l] \arrow[dr, "TC"] \arrow[d, dash] & \\
 & B_1 \arrow[r] \arrow[dl, dotted, "h_1"'] \arrow[d] & B_0 \arrow[dl, dotted, "h_0"'] \arrow[d] & B_2 \arrow[l] \arrow[dl, dotted, "h_2"'] \\
X_1 \arrow[r, "\alpha_1"'] & X_0 & X_2 \ , \arrow[l, "\alpha_2"] &
\end{tikzcd}
\]
avec trois flèches marquées $TC$, on peut trouver $h_1$, $h_0$, $h_2$ laissant le diagramme commutatif (trois triangles et deux carrés commutatifs à respecter). Pour ceci, on commence par choisir $h_0$ de façon à rendre le triangle vertical médian commutatif -- c'est possible car $X_0$ est fibrant, $A_0 \longrightarrow B_0$ dans $TC$.

% foot: Version 25 janvier 2004
\printedpage{36}
% head: A. Grothendieck : Dérivateurs, Ch. XIII, Édition des Universités de Montpellier II et Paris VII
\msp{page 64}

\dgrblock{dgrb717}{Matching objects explain the lifting diagram}
\begin{dcomment}{examples}{p63}{Matching objects explain the lifting diagram}
Grothendieck's three lifts express a recursive compatibility condition. The lift at $X_0$ must be chosen first, because the two remaining lifts are required to induce that same map after composition with $X_1\to X_0$ and $X_2\to X_0$. Choosing three unrelated lifts from objectwise fibrancy would not enforce the squares. The fibrations on the two legs are exactly what turns each final choice into a lifting problem with a prescribed target.

The calculation is the smallest instance of the matching-object description for an inverse diagram. Let $I$ be a finite poset, still indexing contravariant diagrams. For $i\in I$ put
\[
                M_iX=\varprojlim_{j<i}X_j.
\]
An injectively fibrant diagram is characterized by the requirement that each natural map $X_i\to M_iX$ be a fibration. When $i$ has no smaller element, the matching object is terminal, so this says that $X_i$ is fibrant. For $I=\Psi$, the matching objects are $M_0X=e$ and $M_1X=M_2X=X_0$, giving exactly the corollary. This is the inverse-category case of the familiar Reedy description \citep[Theorem 5.1.3]{Hovey1999}.

The lifting criterion can be proved without invoking the terminology. Choose a linear ordering of the finite set of vertices in which smaller vertices occur first. Given an objectwise trivial cofibration $A\to B$ and a map $A\to X$, suppose the lift $B_j\to X_j$ has been chosen for every $j<i$. These maps induce $B_i\to M_iX$. The desired map $B_i\to X_i$ is now a lift in the square whose left side is $A_i\to B_i$ and whose right side is $X_i\to M_iX$. Proceeding through the vertices gives a natural transformation $B\to X$.

For necessity, start with a lifting problem $A\to X_i$, $B\to M_iX$, where $A\to B$ is a trivial cofibration. Form two diagrams which equal $X_j$ for $j<i$, take the values $A$ and $B$ at $i$, respectively, and equal the initial object outside the lower subset $\{j:j\leq i\}$. The given maps to the matching object specify their transition maps. Extension by the initial object is consistent because this subset is lower. Their comparison is objectwise a trivial cofibration, and the first diagram maps to $X$ by the identity below $i$ and the given map at $i$. A lift against $X\to e$ supplies the required $B\to X_i$. This proves the converse. The same finite recursion, factoring $X_i\to M_i(RX)$ at each stage as a trivial cofibration followed by a fibration, constructs a fibrant replacement $X\to RX$.

The condition becomes more informative at a branching vertex. Consider the diamond $0<1<3$, $0<2<3$. Its top matching object is
\[
                           M_3X=X_1\times_{X_0}X_2.
\]
In cochain complexes of vector spaces over a field $k$, take $X_0=0$ and $X_1=X_2=X_3=k[0]$, with both maps from $X_3$ equal to the identity. All objects are fibrant in the usual model structure, and every nonidentity arrow of this diagram is a fibration: the maps are degreewise surjective. Yet
\[
                  X_3\longrightarrow M_3X=k[0]\oplus k[0]
\]
is the diagonal and is not surjective. The diagram is not injectively fibrant. Thus the two separate legs happen to suffice for $\Psi$, whereas a general poset requires compatibility with the entire lower matching diagram. This distinction is why replacements of diagrams carry more information than separate replacements of their components.
\end{dcomment}


Ensuite, on voit séparément pour $h_1$, $h_2$, qu'on peut les trouver en satisfaisant chacune aux deux contraintes de commutativité (1 triangle vertical et un carré penché). Pour $h_1$ p.~ex. l'existence se voit sur le diagramme carré auxiliaire (traits pleins)
\[
\begin{tikzcd}
A_1 \arrow[r, "TC"] \arrow[dd] & B_1 \arrow[dd] \arrow[ddl, dotted, "h_1"'] \arrow[dr, dotted, "\alpha_1^B"] & \\
 & & B_0 \arrow[dl, dotted, "h_0"] \\
X_1 \arrow[r, "F", "\alpha_1^X"'] & X_0 \ \ . &
\end{tikzcd}
\]

\textbf{NB} C'est la démonstration de la proposition 3, page 50.

Le corollaire (partie “il faut”) étant vu, le lemme en résulte, car on sait que si on a un homomorphisme $f$ de carrés cartésiens $X$, $Y$
\[
\begin{tikzcd}[row sep=small, column sep=small]
 & X_1 \arrow[dl, "p_1^X"'] \arrow[dd, "f_1"'] & & X_3 \arrow[ll] \arrow[dl] \arrow[dd, "f_3"] \\
X_0 \arrow[dd, "f_0"'] & & X_2 \arrow[ll, crossing over, "p_2^X"'] & \\
 & Y_1 \arrow[dl, "p_1^Y"'] & & Y_3 \arrow[ll] \arrow[dl] \\
Y_0 & & Y_2 \ \ , \arrow[ll, "p_2^Y"'] \arrow[from=uu, crossing over, "f_2"] &
\end{tikzcd}
\]
avec
\[
p_1^X,\ p_2^X,\ p_1^Y,\ p_2^Y \ \ \text{fibrants \texttt{[appartenant à} $F$\texttt{]} (il suffit même $p_1^X$ et $p_1^Y$ fibrant),}
\]
et
\[
f_0,\ f_1,\ f_2 \ \in \ W,
\]
alors $f_3 \in W$. (Ça a été vu dans la théorie des produits fibrés $W$-homotopiques \texttt{[voir chapitre XII, p. 253]}, et la vérification
\msp{page 65}
directe est immédiate.) De l'hypothèse initiale pertinente que \begin{tikzcd}[row sep=small, column sep=small, baseline=(\tikzcdmatrixname-2-1.base)]
 & X_1 \arrow[dl] \\
X_0 & \\
 & X_2 \arrow[ul]
\end{tikzcd} et \begin{tikzcd}[row sep=small, column sep=small, baseline=(\tikzcdmatrixname-2-1.base)]
 & Y_1 \arrow[dl] \\
Y_0 & \\
 & Y_2 \arrow[ul]
\end{tikzcd}
sont fibrants, on n'a besoin d'utiliser qu'un petit bout, p.~ex. que $p_1^X$ et $p_1^Y$ sont fibrants, qui assure que les deux carrés $X$, $Y$ sont $W$-homotopiquement cartésiens, ce qui suffit à impliquer alors que $f_0$, $f_1$, $f_2 \in W \Longrightarrow f_3 \in W$.

\dgrblock{dgrb718}{A homotopy pullback computed as a complex}
\begin{dcomment}{examples}{p64}{A homotopy pullback computed as a complex}
Grothendieck uses the pullback lemma to make finite gluing invariant under weak equivalences. A concrete calculation shows what fibrant replacement adds to an ordinary pullback. Work with cochain complexes over a field and consider maps $a:X_1\to X_0$ and $b:X_2\to X_0$. A model for their homotopy pullback is
\[
 H=\operatorname{Cone}(a-b:X_1\oplus X_2\to X_0)[-1].
\]
With coordinates chosen so that
\[
 H^n=X_1^n\oplus X_2^n\oplus X_0^{n-1},
 \quad d(x_1,x_2,h)=(dx_1,dx_2,-dh+a(x_1)-b(x_2)),
\]
the equation $d^2=0$ follows from the fact that $a$ and $b$ are chain maps. The third coordinate records the homotopy witnessing equality of the first two images. A strict pullback asks those images to be equal without this additional coordinate.

To verify that this complex computes the derived pullback, replace $a$ by the surjective map $\widetilde X_1\to X_0$, where
\[
 \widetilde X_1^n=X_1^n\oplus X_0^n\oplus X_0^{n-1},
 \quad d(x,z,h)=(dx,dz,-dh+a(x)-z),
\]
and the map to $X_0$ is $(x,z,h)\mapsto z$. The inclusion $x\mapsto(x,a(x),0)$ is a quasi-isomorphism: the projection back to $X_1$ has kernel the contractible complex with differential $(z,h)\mapsto(dz,-dh-z)$, contracted by $(z,h)\mapsto(-h,0)$. The new map is degreewise surjective and hence a fibration. Its ordinary pullback along $b$ imposes $z=b(x_2)$ and is exactly the displayed complex $H$.

There is a short exact sequence of complexes
\[
 0\longrightarrow X_0[-1]\longrightarrow H
   \longrightarrow X_1\oplus X_2\longrightarrow0.
\]
Its connecting map on cohomology is induced by $a-b$. This gives the usual exact sequence for a homotopy pullback and also proves invariance directly: a componentwise quasi-isomorphism between two such spans induces an isomorphism on the two outside terms of the associated long exact sequences, and hence on the middle term. If $a-b$ is degreewise surjective, its ordinary kernel maps quasi-isomorphically to $H$. In particular this occurs when either $a$ or $b$ is degreewise surjective, exactly the fibration condition relevant to this model.

For the span $0\to k[0]\leftarrow0$, the ordinary pullback is zero, while $H=k[-1]$ has $H^1(H)=k$. A fibration replacement of one leg restores this missing class. This is the algebraic version of the familiar fact that the homotopy pullback of the basepoint maps $*\to Y\leftarrow *$ is the loop object of $Y$. It explains both why the manuscript requires fibrant legs and why its gluing step cannot be replaced by taking limits in the homotopy category with no coherent diagram behind them.

In the source's general framework, the corresponding invariance statement uses the $W$-fibration property included in $\mathbf M''$, or the equivalent derived-functor argument under the closed exact hypotheses. No failure of properness is being overlooked in the proof on pp.~64--65: the relevant base-change compatibility is already among its assumptions. The statement that one fibrant leg suffices concerns this particular pullback computation, not a claim that all transition maps in an arbitrary diagram can be ignored.
\end{dcomment}


Je vais essayer de généraliser ce raisonnement, pour traiter le cas d'un ensemble ordonné $I$ qui peut être infini, mais tel que $I^\circ$ soit noethérien et sobre, donc toute suite croissante de parties fermées de $I$, et toute suite décroissante de parties ouvertes, est stationnaire, et de plus toute partie ouverte filtrante croissante a un plus grand élément.

% foot: Version 25 janvier 2004
\printedpage{37}
% head: A. Grothendieck : Dérivateurs, Ch. XIII, Édition des Universités de Montpellier II et Paris VII

% margin: a vertical bar is printed in the left margin beside the lines "prouver ceci : (*) Si I est tel que ... alors I aussi."
On veut prouver que $I$ est admissible, et par “récurrence noethérienne” on est ramené à prouver ceci : $(*)$ Si $I$ est tel que pour tout ouvert $I' \subset I$, avec $I' \neq I$, $I'$ soit admissible, alors $I$ aussi.

$(*)$ admis, prouvons que $I = I_0$ est admissible.

\msp{page 66}
S'il ne l'était, par $(*)$ il existerait $I_1 \subset I_0$, $I_1 \neq I_0$, tel que $I_1$ non admissible. De même, il y aurait $I_2 \subset I_1$, $I_2 \neq I_1$, non admissible etc. -- on aurait construit une suite strictement décroissante d'ouverts $I_n$ de $I$, contrairement au caractère noethérien de $I^\circ$.

Prouvons $(*)$, et distinguons encore deux cas.

\textit{a}) $I^\circ$ est irréductible, donc admet un point générique, \textit{i.e.} un plus petit élément, \textit{i.e.} $I$ admet un plus grand élément. Ce cas a déjà été traité tantôt -- toute catégorie avec objet final est admissible.

\textit{b}) $I^\circ$ est réductible, donc réunion de deux fermés $I_1^\circ$, $I_2^\circ$ distincts de $I^\circ$, donc $I$ est réunion de deux ouverts $I_1$, $I_2$ distincts de $I$. Alors par l'hypothèse dans $(*)$, $I_1$, $I_2$ et $I_0 = I_1 \cap I_2$ sont admissibles. On termine alors comme tantôt, étant ramené à l'admissibilité de \begin{tikzcd}[row sep=small, column sep=small, baseline=(\tikzcdmatrixname-2-1.base)]
 & 1 & \\
\Psi = 0 \arrow[ur] \arrow[dr] & & ! \\
 & 2 &
\end{tikzcd}

\msp{page 67}
Prouvons la partie \textit{c}) du théorème-scholie, partie $1^\circ$). Cela résulte de \textit{a}), une fois vu que pour $s \in I$, $I'/s$ satisfait aux mêmes hypothèses que $I'$. Or $I$ étant ordonné, $I/s$ est une sous-catégorie ouverte de $I$, donc $I'/s$ une sous-catégorie ouverte de $I'$, dont $(I'/s)^\circ$ un sous-ensemble ordonné fermé de $I'^{\circ}$. Ce dernier étant noethérien et sobre, il en est de même de tout fermé, d'où la conclusion.

\dgrblock{dgrb719}{The proposed noetherian enlargement is actually finite}
\begin{dcomment}{exposition}{p66}{The proposed noetherian enlargement is actually finite}
Grothendieck's noetherian induction is patterned on the decomposition of a closed space into irreducible pieces. A generic point makes one piece have a greatest indexing object, and a reducible piece can be glued from two smaller ones. The induction is valid. For the specific Alexandrov topology of a poset, however, the combination ``noetherian and sober'' is stronger than it first appears: it forces the poset to be finite. Thus this paragraph does not in fact add infinite diagram shapes.

We verify the assertion in the precise order convention used here. Opens of $I$ are lower subsets, so opens of $I^{\mathrm{op}}$ are upper subsets of $I$. Sobriety of $I^{\mathrm{op}}$ says that every nonempty directed lower subset of $I$ has a greatest element. If $x_0<x_1<\cdots$ were a strictly increasing sequence, its downward closure would be directed. It could not have a greatest element: any alleged greatest element belongs below some $x_n$, whereas $x_{n+1}$ is strictly larger. Consequently every nonempty subset of $I$ has a maximal element; otherwise dependent choice would construct such a sequence. This is the order-theoretic sobriety criterion, with the order reversed, in \citet[Proposition B.5]{Lindenhovius2014}.

Let $X$ denote $I$ with its upper topology, that is, the topology of $I^{\mathrm{op}}$. Every nonempty subspace $S\subset X$ has an isolated point: take a maximal element $m$ of $S$, so that $S\cap\{x:m\leq x\}=\{m\}$. We now use noetherianity. The union $U$ of all finite open subsets of $X$ is open. Every open subset of a noetherian space is quasi-compact, so the cover of $U$ by its finite open subsets has a finite subcover. Hence $U$ itself is finite. If $X\setminus U$ were nonempty, choose an isolated point $x$ of that subspace. There would be an open $V\subset X$ with $V\cap(X\setminus U)=\{x\}$. Then $U\cup V=U\cup\{x\}$ would be a finite open subset not contained in $U$, contradicting the definition of $U$. Therefore $X=U$ is finite.

Conversely, every finite poset has a noetherian sober Alexandrov topology. Noetherianity is immediate. A nonempty irreducible closed subset has finitely many minimal elements in the lower-topology convention, and irreducibility forces exactly one of them; the closed subset is its closure. Thus the two hypotheses in the proposed enlargement are equivalent to finiteness for posets.

This observation should not be confused with the false assertion that every noetherian sober topological space is finite. Spectra supply abundant infinite examples. The essential extra hypothesis is the Alexandrov topology, in which arbitrary intersections of opens remain open. Nor is it necessary to insert a defect into the manuscript's induction: its irreducible and reducible cases work. The correction is to its claimed extra scope. Infinite towers and group actions on the following pages require another argument, which the Brown lemma supplies whenever the necessary diagram structures and limits exist.
\end{dcomment}


Prouvons la partie $2^\circ$) du théorème-scholie. Considérons
\[
\begin{tikzcd}[row sep=large]
\cM(I')_{\rc} \arrow[r, hook, dotted] & \cM(I') & \cM(I')_{\rf} \arrow[l, hook', dotted] \arrow[d, "u_*^{\rf}"'] \\
\cM(I)_{\rc} \arrow[r, hook, dotted] \arrow[u, "u_{\rc}^*"] & \cM(I) \arrow[u, "u^*"] & \cM(I)_{\rf} \ , \arrow[l, hook', dotted]
\end{tikzcd}
\]
où $u^*$ et a fortiori $u_{\rc}^*$ sont compatibles avec $W(I)$, $W(I')$, et par hypothèse $u_*$ est compatible avec $W(I')_{\rf}$, $W(I)_{\rf}$. En passant aux catégories de fractions, on trouve

% foot: Version 25 janvier 2004
\printedpage{38}
% head: A. Grothendieck : Dérivateurs, Ch. XIII, Édition des Universités de Montpellier II et Paris VII
\msp{page 68}
\[
\begin{tikzcd}[row sep=large]
\cM(I')_{\rc}W(I')_{\rc}^{-1} \arrow[r, dotted, "\sim"] & \HO_W(I') & \cM(I')_{\rf} \arrow[l, dotted, "\sim"'] \arrow[d, "u_*^W"'] \\
\cM(I)_{\rc}W(I)_{\rc}^{-1} \arrow[r, dotted, "\sim"] \arrow[u, "u_W^*"] & \HO_W(I) & \cM(I)_{\rf} \arrow[l, dotted, "\sim"']
\end{tikzcd}
\]
\texttt{[plutôt} $\cM(I')_{\rf}W(I')_{\rf}^{-1}$ \texttt{(resp.} $\cM(I)_{\rf}W(I)_{\rf}^{-1}$\texttt{) au lieu de} $\cM(I')_{\rf}$ \texttt{(resp.} $\cM(I)_{\rf}$\texttt{)]} et on veut montrer que $u_W^*$ et $u_*^W$ sont adjoints, de façon précise trouver un isomorphisme bifonctoriel, pour $F$ \textit{dans} $\cM(I)_{\rc}$ et $G'$ \textit{dans} $\cM(I')_{\rf}$
\[
\Hom_W(u^*F, G') \xrightarrow{\ \sim\ } \Hom_W(F, u_*G') \tag*{$(*)$}
\]
(où $\Hom_W$ désigne le Hom dans $\HO_W(I')$ resp. $\HO_W(I)$). Comme $F$ est dans $\cM(I)_{\rc}$, $u^*F$ est dans $\cM(I')_{\rc}$, et comme $G'$ est dans $\cM(I')_{\rf}$, $u_*G'$ est dans $\cM(I)_{\rf}$. Il s'ensuit que les deux membres de $(*)$ sont des quotients des deux membres dans la bijection canonique
\[
\Hom_{\cM(I')}(u^*F, G') \xrightarrow[\ \sim\ ]{\alpha} \Hom_{\cM(I)}(F, u_*G'). \tag*{$(**)$}
\]
Il faut simplement montrer que cette dernière transforme la relation d'équivalence pertinente $R'$ du premier membre,
\msp{page 69}
en celle $R$ du second.

\dgrblock{dgrb720}{The adjunction after localization, with both inverses}
\begin{dcomment}{exposition}{p68}{The adjunction after localization, with both inverses}
Grothendieck wants the limit functor to retain its universal property after replacement. Knowing that it sends weak equivalences to weak equivalences makes it descend to the homotopy category; it does not alone identify its adjoint. The argument on pp.~68--71 checks that the ordinary adjunction identifies the two homotopy relations. We can make its two directions explicit without choosing a globally functorial replacement.

Choose $F$ cofibrant in $\mathcal M(I)$ and $G$ fibrant in $\mathcal M(I')$. Then $u^*F$ is cofibrant and $u_*G$ is fibrant. Every morphism between these objects in the respective homotopy categories is represented by an ordinary map. Restriction followed by the ordinary counit defines
\[
 \beta:\operatorname{Hom}_{\mathrm{HO}(I)}(F,u_*G)
 \longrightarrow\operatorname{Hom}_{\mathrm{HO}(I')}(u^*F,G).
\]
It is well defined because restriction preserves all weak equivalences. It is surjective because every ordinary map $u^*F\to G$ has a strict adjoint $F\to u_*G$. This is the one direction available before admissibility has been used.

To construct the inverse, choose a trivial cofibration $i':u^*F\to F'$ with $F'$ fibrant, and let $i:F\to u_*F'$ be its strict adjoint. The object $F'$ is also cofibrant. Composition with $i'$ induces a bijection
\[
 \operatorname{Hom}_{\mathrm{HO}(I')}(F',G)
    \xrightarrow{\sim}\operatorname{Hom}_{\mathrm{HO}(I')}(u^*F,G).
\]
For a class $g$ on the right, let $a$ be its unique inverse image and define
\[
                            \gamma(g)=u_*^W(a)\circ[i].
\]
This is where admissibility enters: $u_*^W$ is a functor on the localization of fibrant objects, so it acts on the class $a$, independently of its representative. Naturality of the ordinary counit, together with the fact that $i'$ is adjoint to $i$, gives $\beta\gamma(g)=a[i']=g$.

For the other composite, represent a class $b:F\to u_*G$ by an ordinary map and let $g:u^*F\to G$ be its ordinary adjoint. Since $i'$ is a trivial cofibration and $G$ is fibrant, $g$ extends to an ordinary map $a:F'\to G$ satisfying $ai'=g$. Its class is the inverse image used in defining $\gamma$. The ordinary adjunction then gives $u_*a\circ i=b$, proving $\gamma\beta([b])=[b]$. Thus $\beta$ is bijective. Its naturality, already built into restriction and the counit, gives the required bifunctorial adjunction.

This is the derived-adjunction construction; compare \citet[Lemma 1.3.10]{Hovey1999}. Independence of the chosen $F'$ follows from uniqueness of the inverse of $\beta$, rather than from an additional comparison of arbitrary replacements. It also explains the manuscript's remark that functoriality in $G$ appears a posteriori. The adjoint obtained is
\[
                                u_W^*\dashv u_*^W.
\]
The two occurrences of ``adjoint \`a gauche'' on p.~71 are slips in the source and must read ``adjoint \`a droite''. The Hom formula itself already fixes the orientation, so the proof is unaffected. Likewise the target containing $u^*F'$ on p.~70 is ill typed: $F'$ is already an $I'$-diagram, and the intended object is $u^*F$.
\end{dcomment}


Je dis que l'on a de toute façon
\[
R \ \leq \ \alpha(R')
\]
\textit{sans} hypothèse d'admissibilité sur $u$, \textit{i.e.} que la bijection invese de $(**)$ donne par passage aux quotients un homomorphisme % [sic] invese -> inverse, missing letter
\[
\Hom_W(u^*F, G') \longleftarrow \Hom_W(F, u_*G'), \tag*{$(***)$}
\]
forcément surjectif, et de plus fonctoriel en $F$ \texttt{[,} $G'$ \texttt{barré]}, variables dans $\HO_W(I)$ et $\HO_W(I')$ respectivement. En effet, on définit la flèche $(*$ bis$)$ \texttt{[plutôt} $(***)$\texttt{]} comme composé
\[
\begin{tikzcd}[column sep=large]
\Hom_W(F, u_*G') \arrow[r, "\underbrace{\hspace{2.6cm}}_{\text{\footnotesize\begin{tabular}{@{}c@{}}car $u^*$ est\\ compatible avec\\ $W(I)$, $W(I')$ sur\\ $\cM(I)$ tout entier,\\ et on interprète\\ les deux membres\\ comme des Hom\\ dans $\cM(I)W(I)^{-1}$\\ et $\cM(I')W(I')^{-1}$.\end{tabular}}}"'] & \Hom_W(u^*F, u^*u_*G') \arrow[dr, "\text{\footnotesize\begin{tabular}{@{}l@{}}composition\\ avec l'homomorphisme\\ d'adjonction $u^*u_*G' \longrightarrow G'$\end{tabular}}"] & \\
 & & \Hom_W(u^*F, G').
\end{tikzcd}
\]
% foot: Version 25 janvier 2004

% ---- from Der-ChXIII-25-01-04_p041-045.tex
\printedpage{39}
% head: A. Grothendieck : Dérivateurs, Ch. XIII, Édition des Universités de Montpellier II et Paris VII
% foot: Version 25 janvier 2004

Cela implique donc que le foncteur sur $\HO_W(I)$ $F \longmapsto \Hom_{\HO_{W}(I')}(u^*F,G')$ est, sinon représentable, du moins quotient terme à terme du foncteur représenté par $u_*G'$ (\footnote[37]{Mais il n'est pas clair, sans hypothèse supplémentaire sur la situation, que ce foncteur représentable puisse être choisi en dépendance fonctorielle de $G'$, envisagé dans $\HO_W(I')$.}).

\msp{page 70}

Mais supposons maintenant que $u$ soit admissible. Alors je vais définir la flèche en sens initial $(*)$ (p. 68) de la façon suivante : Je choisis une flèche
\[
i' \ : \ u^*F \longrightarrow F' \qquad \text{avec } i' \in TC \subset C \cap W,\ F' \in \Ob \M(I')_{\rf}
\]
(donc en fait $F' \in \Ob \M(I')_{\mathrm{cf}}$, car $u^*F$ est cofibrant), et par adjonction la flèche $i'$ correspond à
\[
i \ : \ F \longrightarrow u_*F'.
\]
Je considère le diagramme (commutatif!)
\[
\begin{tikzcd}[column sep=large, row sep=large]
\Hom(F',G') \arrow[r, "\alpha"] \arrow[d, "{\scriptstyle \lambda:f'\,\longmapsto\, f'\circ i'}"] & \Hom(u_*F',u_*G') \arrow[d, "{\scriptstyle \mu:f\,\longmapsto\, f\circ i}"] \\
\Hom(u^*F,G') \arrow[r, "\text{adj.}"] & \Hom(F,u_*G')\ .
\end{tikzcd}
\]

On note que l'on a encore $F'$ cofibrant (car $u^*F$ l'est et $i' \in TC \subset C$), donc $\Hom_W(F',G')$ est un quotient de $\Hom(F',G')$, et $\lambda$ induit un isomorphisme
\[
\Hom_W(F',G') \xrightarrow{\ \sim\ } \Hom_W(u^*F',G') % [sic] u^*F' -> u^*F, the target of lambda is Hom(u^*F,G')
\]
(car $i' \in TC \subset W$). Comme $F'$, $G' \in \Ob \M(I')_{\rf}$ et que $u_*^{\rf}$ est compatible avec $W(I')_{\rf}$, $W(I)_{\rf}$, il s'ensuit que la flèche horizontale $\alpha$
\msp{page 71}
est compatible avec les relations d'équivalence définies par
\[
\Hom(F',G') \longrightarrow \Hom_{W(I')}(F',G') \quad \text{resp.} \ \ \Hom(u_*F',u_*G') \longrightarrow \Hom_{W(I)}(u_*F',u_*G')\ .
\]
Et il en est de même de la flèche $\mu : f \longmapsto f \circ i$, trivialement. Donc la flèche d'adjonction $\Hom(u^*F,G') \longrightarrow \Hom(F,u_*G')$ est compatible avec les relations $R$, $R'$, ce qui achève de prouver que $(**)$, page 68, induit une bijection $(*)$, page 68.

Ceci montre donc

\begin{enumerate}
\item[1°)] que pour un objet fixé $G'$ dans $\HO_W(I')$ (qu'on peut toujours supposer, à isomorphisme près, provenir de $\M(I')_{\rf}$), le foncteur
\[
F \ \longmapsto\ \Hom_{\HO_{W}(I')}(u^*_W F,G')
\]
est représentable -- \textit{i.e.} le foncteur $u^*_W$ admet un adjoint à gauche, et
\item[2°)] que cet adjoint à gauche se calcule comme $u_*(G')$ (quand l'argument $G'$ est pris fibrant) -- vu que l'homomorphisme $({*}{*}{*})$, page 69, est \textit{fonctoriel en $F$}. (A posteriori il s'ensuit qu'il est aussi fonctoriel en $G'$ \dots)
\end{enumerate}

\printedpage{40}
% head: A. Grothendieck : Dérivateurs, Ch. XIII, Édition des Universités de Montpellier II et Paris VII
% foot: Version 25 janvier 2004

\msp{page 72}

\dgrblock{dgrb721}{Pointwise right images require fibrant restriction}
\begin{dcomment}{exposition}{p71}{Pointwise right images require fibrant restriction}
Grothendieck now wants more than an abstract right adjoint: he wants to compute its value at an object by a homotopy limit over a comma category. The comparison is canonical, but it becomes an isomorphism only after the same replacement is known to be suitable on both sides. This is the content of the invocation of local isomorphisms on p.~72.

Fix $s\in I$, let $q:I'/s\to I'$ be the projection, and write $p_s:I'/s\to e$. If $X'\to RX'$ is fibrant replacement and $q^*$ preserves fibrant objects, then $q^*RX'$ is a fibrant replacement of $q^*X'$, because restriction always preserves the objectwise weak equivalences. The ordinary pointwise formula therefore identifies
\[
 \bigl(\mathbf Ru_*X'\bigr)_s
   =\bigl(u_*RX'\bigr)_s
   \cong\varprojlim_{I'/s}q^*RX'
   \cong\mathbf R(p_s)_*q^*X'.
\]
The last step uses admissibility of $p_s$. Thus the comparison is computed on one common representative, rather than by comparing two unrelated fibrant replacements.

For clarity, there is a simple sufficient condition for the fibrancy assertion. The functor $q$ is a discrete fibration: an arrow $j\to j'$ in $I'$ and an object $(j',u(j')\to s)$ have a unique lift with source $(j,u(j)\to u(j')\to s)$. In our contravariant convention, when the necessary coproducts exist, its left Kan extension has the formula
\[
                           (q_!A)_j\cong\coprod_{x\in q^{-1}(j)}A_x.
\]
Coproducts of trivial cofibrations have the left lifting property against every fibration, so $q_!$ preserves the objectwise trivial cofibrations. Its right adjoint $q^*$ consequently preserves fibrations and fibrant objects. This is the contravariant version of the discrete-opfibration argument in \citet[Lemma 1.31]{Groth2013}.

If $I$ is a poset, each fiber has at most one element, so only the initial object is needed for this coproduct formula. In the large cocomplete model examples, all small fibers are allowed. In a setting with only finite limits and finite colimits, an arbitrary infinite fiber requires an additional existence hypothesis; it must not be supplied by notation. This is why the general admissibility result and the pointwise formula have been stated separately in the commentary. The finite-poset corollary satisfies the needed conditions, and the usual combinatorial model setting satisfies them for every small functor.

Modern derivator theory packages this calculation as the pointwise Kan-extension axiom. The square of shapes contains the canonical natural transformation $uq\Rightarrow s p_s$; its mate is exactly the displayed comparison. Thinking in terms of that mate explains its compatibility with composition and base change, while the fibrant-representative calculation proves that it is invertible in the present models.

The same proof gives a useful class of base-change squares, with a condition that can be checked directly. Let $v:I_1\to I$ be a discrete fibration, form the strict pullback $I'_1=I'\times_I I_1$, and denote its projections by $v':I'_1\to I'$ and $u_1:I'_1\to I_1$. Assume that the required right Kan extensions and the coproducts over the fibers exist. For each $t\in I_1$ the natural functor
\[
                         I'_1/t\longrightarrow I'/v(t)
\]
is an isomorphism. Indeed, an object of the target is an arrow $u(i')\to v(t)$. Its unique lift along $v$ is an arrow $t'\to t$ with $v(t')=u(i')$, and therefore supplies the object $((i',t'),t'\to t)$ of the source. The uniqueness of lifts gives the inverse also on morphisms. This proves the ordinary base-change identity by the pointwise limit formula.

The pullback $v'$ is again a discrete fibration, so the preceding coproduct argument shows that $(v')^*$ preserves fibrant objects. Evaluate both sides of the ordinary identity on a fibrant replacement of $X'$. They then compute the two derived functors, proving the canonical isomorphism
\[
                      v^*\mathbf Ru_*X'
                     \cong\mathbf R(u_1)_*(v')^*X'.
\]
This is the same discrete-fibration mechanism used in the comma calculation, applied to a family of objects at once; it is part of the standard pointwise derivator calculus cited above.

A strict pullback square without this hypothesis need not have such an isomorphism. Take $I=(0\to1)$, let $I'=\{0\}$ and $I_1=\{1\}$ be the two endpoint inclusions, and use complexes over a field. Their pullback is empty. For $X'=k[0]$, the left side evaluated at $1$ is $k[0]$, because $\{0\}/1$ is terminal. The right side is zero, the limit of the empty diagram in complexes. Thus the base-change map cannot be an isomorphism. The pointwise axiom guarantees the designated comma squares; additional squares require the corresponding exactness argument.
\end{dcomment}


Le dernier énoncé dans 2° du théorème-scholie résulte du calcul de $u_*^W$, compte tenu que $I'/s \longrightarrow I'$ est un isomorphisme local, donc le foncteur image inverse envoie
\[
\M(I')_{\rf} \longrightarrow \M(I'/s)_{\rf},
\]
donc si $G'$ est dans $\M(I')_{\rf}$,
\[
u_*(G')(s) \ \simeq\ p_{s\,*}(G'|I'/s)
\]
est un calcul de $p^W_{s\,*}(G'|I'/s)$.

Cela achève la démonstration du théorème-scholie sur les images directes.

\textbf{Commentaires.}

1) Ce théorème me paraît loin d'être satisfaisant, c'est un pis-aller. J'avais espéré, avec les hypothèses du paragraphe 3, qu'on aurait l'existence des $u_*^W$ pour \textit{toute} flèche $u : I' \longrightarrow I$ dans Cat. Je doute à présent que ce soit le cas, mais n'ai pas encore réussi de faire un contre-exemple. Les cas test les plus importants pour la question d'admissibilité d'une catégorie $I$ me paraissent les suivants :

\begin{enumerate}
\item[\textit{a})] $I$ catégorie discrète infinie (disons dénombrable). La question revient ici à celle-ci : soient $f_i : X_i \longrightarrow Y_i$ dans $W$, avec $X_i$, $Y_i$ fibrants, \textit{i.e.} dans $F$ \texttt{[\textit{i.e.}} $X_i \longrightarrow e$\texttt{,} $Y_i \longrightarrow e$ \texttt{dans} $F$\texttt{]}. Alors $f = \prod f_i : \prod X_i \longrightarrow \prod Y_i$ est-il dans $W$ ?

\msp{page 73}

\dgrblock{dgrb722}{Products and towers have a positive answer}
\begin{dcomment}{status}{p72}{Products and towers have a positive answer}
Grothendieck's test cases isolate two ways in which an infinite limit could lose a weak equivalence. The first takes infinitely many independent factors; the second imposes infinitely many compatibility equations. Under the closed model hypotheses he explicitly mentions, both questions have affirmative answers. Their difficulty is the availability of suitable replacements, not a failure of weak equivalence to survive the limit of a fibrant diagram.

For a set $S$, equip $\mathcal M^S$ with the product structure. The diagonal functor preserves cofibrations and trivial cofibrations componentwise. Its right adjoint, whenever it exists, is $\prod_S$. The Brown argument therefore shows that a product of weak equivalences between fibrant objects is a weak equivalence. This remains true for uncountable $S$ whenever the required products exist. No assumption about exact filtered colimits is used. The derived product is consequently computed by the product of fibrant replacements; the same conclusion appears explicitly in \citet[Example 1.3.11]{Hovey1999}.

For the tower shape, use the matching-object description already seen in finite posets. A tower
\[
                         X_0\longleftarrow X_1\longleftarrow X_2\longleftarrow\cdots
\]
is fibrant precisely when $X_0$ is fibrant and every bonding map is a fibration; all other $X_n$ are then fibrant by composition. The recursive lifting proof has one step for each natural number and defines a compatible lift on the whole tower. Taking the limit is right adjoint to the constant-tower functor, which preserves the objectwise cofibrations and trivial cofibrations. Brown's lemma again gives invariance under levelwise weak equivalences between these towers. In the ordinary complete model setting this is also the inverse Reedy construction.

The manuscript correctly observes that the tower assertion implies the countable product assertion. For a countable family $(X_n)$ of fibrant objects, put $P_r=\prod_{0\leq n\leq r}X_n$. The projection $P_{r+1}\to P_r$ is a pullback of $X_{r+1}\to e$, hence a fibration. A family of weak equivalences induces weak equivalences between the finite products, by the finite case; the inverse limit of the resulting map of towers is the original map of infinite products. This reduction proves the stated implication, although the direct product adjunction is shorter once the general lemma is recognized.

The affirmative answer does not mean that ordinary inverse limits preserve all levelwise weak equivalences. The following example shows exactly where the bonding-map condition matters. It is also a useful safeguard when replacing an abstract existence theorem by an explicit tower of complexes.
\end{dcomment}


\item[\textit{b})] I catégorie % [sic] I (upright roman) -> $I$, italic everywhere else for the category
\[
0 \longrightarrow 1 \longrightarrow 2 \longrightarrow 3 \longrightarrow \cdots,
\]
\textit{i.e.} la catégorie ordonnée associée à l'ensemble ordonné $\mathbf{N}$, pour la relation d'ordre ordinaire $m \leq n$. Donc $\M(I)$ est formée des systèmes \textit{projectifs} dans $\M$, indexés par $\mathbf{N}$. On prouve facilement (en procédant comme dans le corollaire du lemme p. 63) que $\M(I)_{\rf}$ est formée des systèmes projectifs
\[
X_0 \longleftarrow X_1 \longleftarrow X_2 \longleftarrow \cdots \longleftarrow X_n \xleftarrow{\ p_n\ } X_{n+1} \longleftarrow \cdots
\]
tels que les $X_i$ soient dans $\M_{\rf}$, et les $p_n$ dans $F$. La question est donc : si $X_\bullet$, $Y_\bullet$ sont fibrants et $f_\bullet : X_\bullet \longrightarrow Y_\bullet$ est telle que $f_i \in W$ pour $i \in \mathbf{N}$, alors a-t-on $\varprojlim f_i : X_\infty \longrightarrow Y_\infty$ dans $W$ ?

Si la reponse est affirmative pour $\mathbf{N}$ (pour un $W$ et une quillénisation donnée), elle % [sic] reponse -> réponse, accent missing
l'est aussi pour toute catégorie discrète dénombrable, comme on le voit aussitôt. Mais je ne vois pas comment déduire même \textit{a}) des axiomes, disons pour un triplet de Quillen clos (closed model category) tel que $\M$ soit de plus stable par $\varinjlim$ filtrantes, que celles-ci soient exactes, et qu'il y ait
\printedpage{41}
% head: A. Grothendieck : Dérivateurs, Ch. XIII, Édition des Universités de Montpellier II et Paris VII
% foot: Version 25 janvier 2004
\msp{page 74}
un petit ensemble de générateurs dans $\M$ (donc $\M$ est accessible), que $W$ soit stable par $\varinjlim$ filtrantes. Je devrais demander à \textsc{Thomason} s'il connait un contre-exemple. % [sic] connait -> connaît, circumflex missing

\dgrblock{dgrb723}{An inverse limit of acyclic complexes can have cohomology}
\begin{dcomment}{examples}{p73}{An inverse limit of acyclic complexes can have cohomology}
The insistence on fibrant towers is mathematically necessary, even in the most familiar algebraic model. Fix a prime $p$ and, for $n\geq1$, form the cochain complex
\[
 A_n=\bigl[\mathbb Z\xrightarrow{p^n}\mathbb Z
                   \xrightarrow{\bmod p^n}\mathbb Z/p^n\bigr]
\]
in degrees $0,1,2$. It is acyclic: the first map is injective, the image of that map is the kernel of the second, and the second is surjective. Define $A_{n+1}\to A_n$ by multiplication by $p$ in degree $0$, the identity in degree $1$, and reduction modulo $p^n$ in degree $2$. These maps commute with both differentials. The maps of towers $A_n\to0$ are therefore quasi-isomorphisms at every level.

Their ordinary inverse limit is readily computed. A compatible sequence in degree $0$ satisfies $a_n=p a_{n+1}$. Its first coordinate is divisible by every power of $p$, and hence vanishes; then every coordinate vanishes. The limit in degree $1$ is $\mathbb Z$, and the limit in degree $2$ is $\mathbb Z_p$. Thus
\[
               \varprojlim_n A_n=\bigl[0\longrightarrow\mathbb Z
                                          \longrightarrow\mathbb Z_p\bigr].
\]
The map $\mathbb Z\to\mathbb Z_p$ is injective but not surjective. For example, $\mathbb Z_p$ is uncountable, since a sequence of digits in $\{0,\ldots,p-1\}$ uniquely determines a $p$-adic integer, while $\mathbb Z$ is countable. Consequently
\[
              H^2\!\left(\varprojlim_n A_n\right)=\mathbb Z_p/\mathbb Z\ne0.
\]
The inverse limit of the levelwise quasi-isomorphism $A_\bullet\to0$ is not a quasi-isomorphism.

There is no conflict with the positive result on the preceding page. In the projective model structure on unbounded complexes of abelian groups every object is fibrant, but a fibration is degreewise surjective. The bonding map of this tower is multiplication by $p$ in degree $0$, so it is not a fibration. Individual fibrancy has not made the \emph{diagram} fibrant. A fibrant tower replacement would have an acyclic inverse limit by the Brown argument.

This example separates three statements which can otherwise look deceptively similar: ordinary inverse limit is not exact; inverse limit need not preserve levelwise quasi-isomorphisms of arbitrary towers; inverse limit does preserve weak equivalences between fibrant towers in the model structure under discussion. Only the third statement is the admissibility problem posed in the manuscript. Its positive answer is compatible with the failure exhibited here.
\end{dcomment}


\item[\textit{c})] Les deux exemples précédents étaient des catégories ordonnées. Comme exemple typique de catégorie non-ordonnée, où je ne vois pas ce qui se passe, je signale $I = B_G$, $G$ un groupe (donc un seul objet, avec un groupe d'automorphismes $G$). Aussi le cas de $\boldsymbol{\Delta}[1] = \{\boldsymbol{\Delta}_0, \boldsymbol{\Delta}_1\}$ (sous-catégorie pleine de $\boldsymbol{\Delta}$) n'est pas clair pour moi.
\end{enumerate}

\dgrblock{dgrb724}{A group action detects the difference between fixed and homotopy fixed points}
\begin{dcomment}{examples}{p74:Les deux exemples précédents étaient des catégories ordonnées.}{A group action detects the difference between fixed and homotopy fixed points}
Grothendieck introduces $B_G$ because automorphisms are absent from the finite-poset argument. The limit of a $G$-diagram is a fixed-point object. Its derived version measures the coherence of invariance under the whole action, and a single fibrant value does not guarantee that the action diagram is injectively fibrant. The algebraic example below makes the difference explicit and at the same time answers the boundedness question from \S5.

Let $G=C_p=\langle g\mid g^p=1\rangle$, let $k$ be a field of characteristic $p$, and give $k$ the trivial action. Using inversion to change a right action into a left action if necessary, diagrams of complexes on $B_G$ are complexes of $R=k[G]$-modules. With $t=g-1$, one has $R\cong k[t]/(t^p)$. The trivial module has the free resolution
\[
 \cdots\xrightarrow{t}R\xrightarrow{t^{p-1}}R
       \xrightarrow{t}R\longrightarrow k\longrightarrow0.
\]
Its exactness follows immediately from
\[
       \ker(t:R\to R)=(t^{p-1}),\qquad
       \ker(t^{p-1}:R\to R)=(t).
\]
Applying $\operatorname{Hom}_R(-,k)$ makes every differential zero, since $t$ acts by zero on $k$. Tensoring with $k$ does the same. Therefore
\[
 \operatorname{Ext}^n_R(k,k)\cong k,
 \qquad \operatorname{Tor}^R_n(k,k)\cong k
 \qquad(n\geq0).
\]
For the projection $p:B_G\to e$ this gives
\[
 \mathbf Rp_*k=\mathbf R\operatorname{Hom}_{k[G]}(k,k),
 \qquad \mathbf Lp_!k=k\otimes^{\mathbf L}_{k[G]}k.
\]
The first object has nonzero cohomology in every nonnegative degree, and the second in every nonpositive degree. Both exist in $D(k)$. Neither belongs to $D^{\mathrm b}(k)$, although the input is concentrated in degree zero. Finite cohomological dimension is therefore a boundedness condition in this example, not a prerequisite for the unbounded operations.

Ordinary invariants of the original diagram give only $k[0]$. Moreover, over a field its underlying complex is already fibrant. The additional cohomology arises because injective fibrancy in the diagram category is a condition on the $k[G]$-complex, not merely on the underlying $k$-complex. Replacing that diagram by a $K$-injective $k[G]$-complex and taking ordinary invariants computes $\mathbf Rp_*$. The Brown lemma proves invariance between such fibrant diagrams, so the presence of a nontrivial automorphism group creates no exception to admissibility.

Exactly the same general argument applies to the full subcategory of $\boldsymbol\Delta$ on $[0]$ and $[1]$ mentioned next in the manuscript. It is not the walking arrow: it contains all maps between those two ordinals, including the two sections and the retraction. It has two objects and seven morphisms: the four Hom sets have respectively $1,2,1,3$ elements. Thus its ordinary limit uses only the finite limits of $\mathbf M0$. Its relations matter when constructing a coherent diagram, but the assumed uniform structure and the lifting-adjunction proof already give admissibility; no further infinite-completeness assumption is needed for this test.
\end{dcomment}


2) On aimerait pouvoir traiter de façon duale l'existence des $u_!^W$, à défaut de mieux. Dans l'hypothèse de départ du paragraphe 4 sur $(\M,W)$, cela amène donc à changer de tactique. Cette fois, on devra supposer l'existence d'une quillénisation close \textit{exacte à droite} pour $W$ (ce qui est plus exigeant, semble-t-il, qu'une telle structure exacte à gauche, dans le cas envisagé), donc
\[
TF \ = \ F \cap C,
\]
\msp{page 75}
et de plus encore $C$, $TC$ tels qu'on ait
\[
C \ = \ \widetilde{C_0}, \qquad TC \ = \ \widetilde{TC_0}
\]
avec $C_0$, $TC_0$ petits (\footnote[38]{On dit alors que $C$, $TC$ (en tant que parties closes à droite \texttt{[plutôt à gauche]}) sont “accessibles”.}). On posera maintenant
\[
\begin{aligned}
F(I) \ &= \ \{f \in \Fl \M(I) \ \mid \ f_s \in F \quad\ \ \forall\, s \in \Ob I\}\\
TF(I) \ &= \ \{f \in \Fl \M(I) \ \mid \ f_s \in TF \ \ \forall\, s \in \Ob I\},
\end{aligned}
\]
puis
\[
\begin{aligned}
C(I) \ &= \ TF(I)^*\\
TC(I) \ &= \ F(I)^*\ .
\end{aligned}
\]
L'espoir (que je devrai vérifier ultérieurement), c'est qu'on obtient ainsi une quillénisation uniforme à droite de $W$, sur Cat tout entier. Le seul problème, c'est la validité des axiomes de factorisation de $f \in \Fl \M(I)$ en $f = pi$, avec $i \in TC$, $f \in F$, et aussi avec $i \in C$, $f \in TF$. % [sic] f \in F, f \in TF -> p \in F, p \in TF, the factorisation is f = pi
J'ai bon espoir que c'est vrai -- cela se réduit essentiellement à voir que $(C(I),TF(I))$ et $(TC(I),F(I))$, soit $(\Phi(I),\Psi(I))$, satisfont en plus de $\Phi(I) = \Psi(I)^*$, à $\Psi(I) = \Phi(I)_*$ (\textit{i.e.} $\Psi(I)$ est clos à droite), et que $\Phi(I)$ est accessible (tout comme $\Phi = \Phi(e)$). Si ça marche, on aura donc l'existence
\printedpage{42}
% head: A. Grothendieck : Dérivateurs, Ch. XIII, Édition des Universités de Montpellier II et Paris VII
% foot: Version 25 janvier 2004
\msp{page 76}
de $u_!$ pour $u$ un foncteur entre catégories ordonnées finies, et sous des conditions légèrement plus générales. (P. ex. $u : I' \longrightarrow I$ avec $I$ ordonnée, $I'$ ordonnée et à espace sous-jacent noethérien et sobre.) Peut-être que cette fois, l'hypothèse que $W$ est stable par $\varinjlim$ filtrantes va quand même servir! Ainsi, le cas des catégories discrètes, et des catégories filtrantes à gauche comme $\mathbf{N}^\circ$, ne va plus poser de problèmes à cause de cela. Il y a donc de l'espoir! Tout ensemble ordonné qui est $\varinjlim$ filtrante de sous-ensembles ordonnés \textit{fermés} noethériens et sobres, est alors $W$-admissible à droite. Mais qu'en est-il de catégories de type $\Gamma = B_G$ p. ex. ? ? ?

\dgrblock{dgrb725}{The projective construction supplies the dual theory}
\begin{dcomment}{later}{p75}{The projective construction supplies the dual theory}
Grothendieck's change of tactic is the right one for left images. To derive a colimit, fibrations should be defined objectwise, so restriction becomes a right Quillen functor and its left adjoint acts on cofibrant diagrams. The extra small generating sets are needed to construct these projective diagram structures. Once they have been constructed, the dual Brown argument supplies homotopy invariance for every existing left Kan extension, rather than only for finite posets.

Here is the construction with its size hypothesis visible. Suppose that $\mathcal M$ is a locally presentable, cofibrantly generated model category, with generating cofibrations $C_0$ and generating trivial cofibrations $T_0$. For $i\in I$, evaluation has a left adjoint $L_i$, given in the contravariant convention by
\[
                         (L_iA)_j=\coprod_{I(j,i)}A.
\]
Use the sets $\{L_ic:i\in I,c\in C_0\}$ and $\{L_it:i\in I,t\in T_0\}$. By the evaluation adjunction, their right lifting classes are exactly the objectwise trivial fibrations and the objectwise fibrations, respectively. The sets really are small: $I$ and the two generating collections are small. For a sufficiently large regular cardinal the domains are presentable, because evaluation preserves filtered colimits and
\[
 \operatorname{Hom}_{\mathcal M(I)}(L_iA,X)
                        \cong\operatorname{Hom}_{\mathcal M}(A,X_i).
\]
The small object argument therefore gives both desired factorizations.

One also checks the weak equivalences in the transferred structure. A relative cell complex built from the $L_it$ is componentwise a composite of pushouts of coproducts of trivial cofibrations. The base trivial cofibrations form a left lifting class and are closed under these operations, so the resulting map is an objectwise trivial cofibration and hence an objectwise weak equivalence. Taking retracts gives the same conclusion for every map in $TC(I)$. Moreover $TC(I)\subset C(I)$, because the objectwise trivial fibrations are objectwise fibrations. If $f\in C(I)\cap W(I)$, factor it as $f=pj$ with $j\in TC(I)$ and $p\in F(I)$. Two-out-of-three gives $p\in W(I)$, so $p\in TF(I)$ by the objectwise definitions. The lifting property of $f$ against $p$ makes $f$ a retract of $j$, hence $f\in TC(I)$. This identifies the required acyclic class and completes the model-structure verification. It is the projective transfer argument sought in the paragraph following the two formulas on p.~75. See \citet[Theorem 1.28 and Lemma 1.29]{Groth2013}; the small object argument itself is recorded in \citet[Theorem 2.1.14]{Hovey1999}.

For $u:I'\to I$, restriction now preserves the objectwise fibrations and trivial fibrations. Adjointness makes $u_!$ preserve cofibrations and trivial cofibrations. The dual Brown lemma shows that it preserves weak equivalences between cofibrant diagrams, and the derived adjunction is
\[
                              \mathbf Lu_!\dashv u^*.
\]
The hypotheses of local presentability and cofibrant generation make all small diagram shapes available, including $B_G$. An objectwise construction of fibrations does not mean that cofibrant diagrams are objectwise cofibrant with no further condition; the latter would miss precisely the relations that the free diagrams $L_iA$ impose.

Two slips in the source obscure this otherwise correct programme. Exactness on the right is $TF=F\cap W$, not $F\cap C$. In the factorizations $f=pi$, the membership assertions must refer to $p\in F$ or $p\in TF$, not to $f$. They are retained as printed and explained in the correction register. The role of filtered colimits also deserves precision: in the projective transfer just described, closure of trivial cofibrations under transfinite composites follows from their lifting property. A separate assertion that all weak equivalences are closed under filtered colimits is not what proves that step.
\end{dcomment}


\msp{page 77}

\dgrblock{dgrb726}{Why the homotopy problem returns to SGA 4}
\begin{dcomment}{context}{p76}{Why the homotopy problem returns to SGA 4}
Grothendieck turns to cardinal filtrations because a construction on every diagram category needs a controlled supply of small pieces. The issue is already visible in the proposed right uniform quillenisation: a lifting condition must be generated by a set of maps, and the factorisations built from that set must remain available after changing the indexing category. The return to SGA~4, I, \S9 is therefore part of the homotopical argument. It supplies a way of bringing a large category within reach of set-indexed constructions.

The later letter to Thomason makes this connection explicit. In its opening pages Grothendieck describes his renewed work on accessible categories and accessible parts as a rebuilding of the provisional account in SGA~4. In the next section of the letter, he distinguishes the relative category $(\mathcal M,W)$, which specifies the homotopy theory, from the cofibrations and fibrations used to construct its operations \citep[pp.~2--4 of the electronic edition]{Grothendieck1991Letter}. The present chapter places these two concerns next to one another: first find enough replacements to derive direct images; then explain the smallness that makes such replacements possible. This is also why the size of a \emph{diagram}, rather than just the sizes of its values, becomes decisive near the end of \S7.

There are two historical cautions. The word ``accessible'' here follows the cardinal conventions of SGA~4: ``grand devant $\pi$'' bounds every subset of size at most $\pi$, and thus corresponds to $\pi^+$-directedness in the regular-cardinal terminology now customary. Also, the closing definition tests membership in a part against \emph{every} perfectly adapted presentation. It is stronger than merely exhibiting the part as a category of sufficiently directed colimits of small objects. The distinction is mathematically effective: the counterexamples below separate those conditions even in a powerset ordered by inclusion. The later development of accessibility does not license us to identify formulations whose quantifiers differ.
\end{dcomment}


\dgrblock{dgrb727}{What has been settled before the cardinal detour}
\begin{dcomment}{status}{p76}{What has been settled before the cardinal detour}
The final questions of \S6 can now be separated by the task they perform. For the uniform closed exact structures assumed in the theorem-scholie, the Brown lemma settles admissibility of every functor with the requisite ordinary right adjoint; its dual settles left admissibility. Products, towers, group actions and the truncated simplex category are all included. The finite-poset proof remains a direct construction of the finite calculus, while the noetherian sober Alexandrov condition adds no infinite posets.

The next question is existence of those uniform structures. In the usual combinatorial setting, injective and projective diagram model structures, derived adjunctions, and the pointwise comma-category formulas furnish a strong derivator on all small categories. This is the later theorem corresponding most directly to Grothendieck's programme. In smaller categories such as bounded complexes, one retains finite operations and must check each infinite ordinary limit or colimit separately. The bounded-product example above prevents a blanket completeness assertion.

Section 7 addresses the remaining constructive burden: how a large lifting class or a large collection of objects can be controlled by a small family at a sufficiently large cardinal. That question is substantial even though the preceding homotopy-invariance doubts have a formal answer. It is also where the older SGA~4 language of accessibility and the later language of presentability need to be compared with particular care. The change of subject is therefore motivated by the production of factorizations, rather than by a need to repair Brown's argument.
\end{dcomment}


\section*{7\quad Filtrations cardinales et parties accessibles d'une grosse catégorie}

\dgrblock{dgrb728}{A size filtration and a presentability rank}
\begin{dcomment}{exposition}{p77:Filtrations cardinales et parties accessibles d'une grosse catégorie}{A size filtration and a presentability rank}
Grothendieck returns to cardinal estimates because the constructions of model structures need a controlled supply of small objects. A replacement construction may be performed object by object and still fail to be controlled uniformly over a category. The filtration is intended to provide that control: a large object is assembled from bounded pieces, while bounded constructions on bounded pieces remain bounded. The categorical purpose of the estimates should remain visible through the set theory.

There is a one-cardinal shift in the language. Throughout these pages, ``grand devant $\pi$'' means that every subset of cardinal at most $\pi$ has an upper bound. In modern terminology such a nonempty poset is $\pi^+$-directed, where $\pi^+$ is the successor cardinal. Write $\kappa=\pi^+$ when translating the assertions. An object called $\pi$-accessible in the sense of SGA~4 is correspondingly $\kappa$-presentable: its covariant Hom functor preserves these directed colimits. This meaning of ``accessible object'' must also be kept distinct from the new meaning of ``accessible part'' introduced on manuscript page~86. The latter concerns a full subcategory and its position inside the ambient category.

The three filtration conditions do not say that every object of $\Filt^\pi(E)$ is $\pi^+$-presentable. They say that this stage is essentially small, is closed under the specified bounded colimits, and generates $E$ under sufficiently directed colimits. Preservation of such colimits by Hom is an additional hypothesis. This explains why the hypothesis $\Filt^\pi(E)\subset E_\pi$, initially left tacit, becomes necessary on pages~90--94. Under suitable accessibility hypotheses, SGA~4 supplies useful cardinals at which size and presentability do coincide; see \citet[I, 9.16--9.18]{SGA4I1972}. It does not identify the two notions at every index of an arbitrary filtration.

The universe convention is also doing real work. ``Essentially small'' means that a skeleton belongs to the fixed universe; it does not mean that there are finitely or countably many objects. Once a skeleton is chosen, the relevant comma categories are small because $E$ is a $\mathfrak U$-category. All diagrams below remain small relative to this universe. No assertion that the whole large category has a set of objects in that same universe is needed.

Finally, the filtration is exhaustive even though exhaustiveness is not listed separately. Choose the presentation from condition~(c) at $\pi_0$, and choose $\rho\geq\pi_0$ larger than the cardinality of its indexing poset. Its terms lie in $\Filt^\rho(E)$, and condition~(b) places its colimit in the same stage. Thus every object eventually enters the filtration. This short argument is the content of footnote~39.
\end{dcomment}


Je me réfère à la théorie des foncteurs et catégories accessibles de SGA 4, I, §9 (p. 138-178). Je me donne une grosse catégorie (toujours une $\mathfrak{U}$-catégorie) $E$ avec une \textit{filtration cardinale} (croissante)
\[
\Bigl(\Filt^{\pi}(E)\Bigr)_{\pi \geq \pi_0}, \qquad \pi_0 \text{ un cardinal infini fixé,}
\]
par les sous-catégories strictement pleines $\Filt^{\pi}(E)$. On a \texttt{[SGA 4, I, 9.12]}

\begin{enumerate}
\item[\textit{a})] $\Filt^{\pi}(E)$ essentiellement petite, \textit{i.e.} équivalente à une petite catégorie (\textit{i.e.} à une catégorie appartenant à $\mathfrak{U}$).
\item[\textit{b})] $E$ satisfait $L_{\pi_0}$ (\textit{i.e.} $E$ est stable par $\varinjlim$ de type $I$, quand $I$ est un ensemble ordonné grand devant $\pi_0$), et $\Filt^{\pi}(E)$ est stable par les limites inductives filtrantes indexées par $I$ grand devant $\pi_0$, et tel que $\card I \leq \pi$.
\item[\textit{c})] Pour tout $\pi \geq \pi_0$, et pour tout $X \in \Ob E$, on a
\[
X \ \simeq\ \varinjlim_{i \in I} X_i,
\]
où les $X_i$ \texttt{[sont]} dans $\Filt^{\pi}(E)$, et où $I$ est un ensemble ordonné grand devant $\pi$. De plus, si $X \in \Ob \Filt^{\pi'}(E)$, avec $\pi' > \pi$, on peut prendre $I$ tel que $\card I \leq \pi'^{\pi}$ (donc si $\pi' = 2^{c'}$ avec $c' \geq \pi$, d'où $(\pi')^{\pi} = (2^{c'})^{\pi} = 2^{c'\pi} = 2^{c'} = \pi'$, on aura $\card I \leq \pi'$). (\textbf{NB} Si on admet l'hypothèse du continu généralisée, on aura toujours $\pi'^{\pi} = \pi'$, donc $\card I \leq \pi'$.) (\footnote[39]{NB \textit{b}) et \textit{c}) impliquent que $E = \bigcup_{\pi \geq \pi_0} \Filt^{\pi}(E)$.})
\end{enumerate}

\dgrblock{dgrb729}{The useful exponent bound and the singular-cardinal exception}
\begin{dcomment}{examples}{p77:Si on admet l'hypothèse du continu généralisée}{The useful exponent bound and the singular-cardinal exception}
The bound $(\pi')^\pi$ measures how many small pieces may be needed to describe an object of size $\pi'$. It should be retained until an actual cardinal identity permits its replacement by $\pi'$. The calculation for a power-set cardinal is correct: if $\pi'=2^{c'}$ and $c'\geq\pi$ are infinite, then
\[
 (\pi')^\pi=(2^{c'})^\pi=2^{c'\cdot\pi}=2^{c'}=\pi'.
\]
Thus arbitrarily large useful choices are available without a continuum hypothesis.

The subsequent claim about the generalized continuum hypothesis is false as written. If $\mu$ is singular and $\operatorname{cf}(\mu)\leq\pi<\mu$, then $\mu^\pi>\mu$, independently of GCH. Here is the elementary obstruction. Choose increasing infinite cardinals $\mu_\alpha<\mu$, for $\alpha<\delta=\operatorname{cf}(\mu)$, with supremum $\mu$. K\"onig's diagonal argument gives
\[
 \mu=\sum_{\alpha<\delta}\mu_\alpha
       <\prod_{\alpha<\delta}\mu_\alpha^+
       \leq\mu^\delta\leq\mu^\pi.
\]
For completeness, the strict inequality follows by partitioning a purported list of all elements of the product into pieces $S_\alpha$ of size $\mu_\alpha$. In coordinate $\alpha$, choose a value outside the $\alpha$-coordinates of $S_\alpha$; the resulting element is absent from every piece of the list. Taking $\mu=\aleph_\omega$ and $\pi=\aleph_0$ already disproves the printed assertion.

Under GCH the conclusion does hold when $\pi<\operatorname{cf}(\pi')$. In that case each map $\pi\to\pi'$ has bounded image. For every infinite $\nu<\pi'$, GCH gives
\[
 \nu^\pi\leq(2^\nu)^\pi=2^{\max(\nu,\pi)}
 =\max(\nu,\pi)^+\leq\pi',
\]
and the union of $\pi'$ such sets has size at most $\pi'$. The constant functions give the reverse inequality. The missing cofinality condition matters precisely at singular cardinals. Nothing later in the construction requires the erroneous unconditional simplification: the bound $(\pi')^\pi$, and the unconditional power-set choices, suffice.
\end{dcomment}


\printedpage{43}
% head: A. Grothendieck : Dérivateurs, Ch. XIII, Édition des Universités de Montpellier II et Paris VII
% foot: Version 25 janvier 2004

\msp{page 78 barrée}

\msp{page 79}

Considérons la représentation de \textit{c}),
\[
X \ = \ \varinjlim_{I} X_i \qquad\qquad
\left\{
\begin{aligned}
& X_i \in \Filt^{\pi},\\
& I \text{ grand devant } \pi.
\end{aligned}
\right.
\]
Soit
\[
\tilde{I} \ = \ \{I' \subset I \mid \card I' \leq \pi,\ \ I' \text{ grand devant } \pi_0\},
\]
muni de la relation d'ordre d'inclusion. Pour
\[
I' \ \in\ \tilde{I}, \qquad \text{soit}\ \ X_{I'} = \varinjlim_{I'} X_i\ ,
\]
cette limite existe par $L_{\pi_0}$ (hypothèse \textit{b}), elle est dans $\Filt^{\pi}$ par \textit{b}), donc on a un système inductif
\[
(X_{I'})_{I' \in \tilde{I}}, \qquad X_{I'} \in \Filt^{\pi},
\]
avec
\[
\card \tilde{I} \ \leq\ (\card I)^{\pi},
\]
et si $X \in \Filt^{\pi'}$ $(\pi' > \pi)$, prenant
\[
\card I \ \leq\ \pi'^{\pi},
\]
on aura $\card \tilde{I} \leq (\pi'^{\pi})^{\pi} = \pi'^{\pi^2} = \pi'^{\pi}$, donc
\[
\card \tilde{I} \ \leq\ \pi'^{\pi}.
\]
Je dis que $\tilde{I}$ est lui aussi \textit{grand devant $\pi$,} tout comme $I$. En effet,
\msp{page 80}
soit
\[
(I_\alpha)_{\alpha \in A}, \qquad I_\alpha \in \tilde{I}, \qquad \card A \ \leq\ \pi,
\]
prouvons qu'il existe $I' \in \tilde{I}$ qui majore les $I_\alpha$, \textit{i.e.} qui contient $I'_0 = \bigcup_{\alpha \in A} I_\alpha$. On aura $\card I'_0 \leq \pi^2 = \pi$ (car $\card A \leq \pi$ et $\card I_\alpha \leq \pi$ pour tout $\alpha$), et il suffit de voir ceci :

\texttt{[En fait, l'assertion est évidente : comme} $\card I'_0 \leq \pi$\texttt{, et comme} $I$ \texttt{est grand devant} $\pi$\texttt{,} $I'_0$ \texttt{admet un majorant} $i'_0$ \texttt{dans} $I$\texttt{, et on peut prendre} $I' = I'_0 \cup \{i'_0\}$\texttt{.]}

\textbf{Lemme 1 :} \textit{Soient $\pi_0$, $\pi$ deux cardinaux, avec $\pi \geq \pi_0$, $I$ un ensemble ordonné grand devant $\pi$, et $I'_0$ une partie de $I$ de cardinal $\leq \pi$. Alors il existe une partie $I'$ de $I$ satisfaisant les conditions :} (\footnote[40]{Pourvu qu'on suppose $\pi > \pi_0$, et de plus $\pi^{\pi_0} = \pi$ (p. ex. $\pi = 2^c$ avec $c \geq \pi_0$).})

\begin{enumerate}
\item[a)] $I' \supset I'_0$.
\item[b)] $\card I' \leq \pi$.
\end{enumerate}

% ---- from Der-ChXIII-25-01-04_p046-050.tex
\printedpage{44}
% head: A. Grothendieck : Dérivateurs, Ch. XIII, Édition des Universités de Montpellier II et Paris VII

\begin{itemize}
\item[c)] \textit{$I'$ grand devant $\pi_0$.}
\end{itemize}

\texttt{[Le lemme tel qu'il est énoncé est évident pour les mêmes raisons que ci-dessus (même sans l'hypothèse de la note de bas de page). La démonstration qui suit prouve en fait un résultat plus fin :}

\texttt{\textsc{Lemme} 1$'$ : \textit{Soient $\pi_0$, $\pi$ deux cardinaux, avec $\pi_0<\pi$ et $\pi^{\pi_0}=\pi$, $I$ un ensemble ordonné grand devant $\pi_0$, et $I'_0$ une partie de $I$ de cardinal $\le\pi$. Alors il existe une partie $I'$ de $I$ satisfaisant les conditions :}}

\begin{itemize}
\item[\texttt{a)}] $I'\supset I'_0$\texttt{.}
\item[\texttt{b)}] $\card I'\le\pi$\texttt{.}
\item[\texttt{c)}] \texttt{\textit{$I'$ grand devant $\pi_0$. ]}}
\end{itemize}

On va essayer de construire $I'$ par récurrence transfinie, en définissant les $I'_\alpha\subset I$ ainsi :
\[
\begin{aligned}
I'_0 &= I'_0\;, \\
I'_\alpha &= \textstyle\bigcup_{\beta<\alpha} I'_\beta \quad \text{si $\alpha$ est un ordinal limite,} \\
I'_{\alpha+1} &= \left\{\parbox{0.6\textwidth}{déduit de $I'_\alpha$ ainsi : pour toute partie $J$ de $I'_\alpha$ telle que $\card J\le\pi_0$, qui n'est pas majorée dans $I'_\alpha$, on choisit un majorant $i_J$ dans $I$ (\texttt{[qui]} existe, car $I$ \texttt{[est]} grand devant $\pi$, donc devant $\pi_0$) et $I'_{\alpha+1}$ s'obtient en adjoignant ces $i_J$ à $I'_\alpha$.}\right.
\end{aligned}
\]

\msp{page 81}
Soit alors $\alpha_0$ le premier ordinal qui soit de cardinal $>\pi_0$, et considérons
\[
I' \;=\; I'_{\alpha_0} \;=\; \bigcup\limits_{\substack{\alpha<\alpha_0\\ \textit{i.e.}\ \card\alpha\le\pi_0}} I'_\alpha .
\]
Alors $I'$ contient $I'_0$, il est grand devant $\pi_0$, mais quel est son cardinal ? Prouvons par récurrence transfinie que
\[
\card I'_\alpha \;\le\pi \qquad\quad \text{si } \alpha<\alpha_0,\ \ \textit{i.e.}\ \ \card\alpha\le\pi_0 .
\]
C'est O.K. pour $\alpha=0$ par hypothèse sur $I'_0$. Si c'est vrai pour $I'_\alpha$, est-ce vrai pour $I'_{\alpha+1}$ ? On adjoint à $I'_\alpha$ un ensemble d'éléments $i_J$ qui est de cardinal $\le\card\mathfrak{P}_{\le\pi_0}(I'_\alpha)$, donc de cardinal $\le\card(I'_\alpha)^{\pi_0}\le\pi^{\pi_0}$, et si on suppose
\[
\pi^{\pi_0} \;=\; \pi,
\]
on gagne. Supposons enfin que $\alpha$ soit un ordinal limite (de cardinal $\le\pi_0$). Donc on a
\[
I'_\alpha \;=\; \bigcup_{\beta<\alpha} I'_\beta ,
\]
donc
\printedpage{45}
% head: A. Grothendieck : Dérivateurs, Ch. XIII, Édition des Universités de Montpellier II et Paris VII
\msp{page 82}
\[
\card I'_\alpha \;\le\; \underbrace{(\card\alpha)\pi}_{\le\pi\pi=\pi^2=\pi} \qquad\quad \text{puisque les } \card I'_\beta\le\pi,
\]
on gagne encore.

On trouve enfin, par le même argument,
\[
\card I'_{\alpha_0} \;\le\; \underbrace{(\card\alpha_0)\pi\le\pi^2}_{\substack{\text{si }\card\alpha_0\le\pi,\\ \textit{i.e.}\text{ si }\pi_0<\pi}} = \pi,
\]
donc on gagne encore.

\dgrblock{dgrb730}{Why the transfinite closure has to run far enough}
\begin{dcomment}{exposition}{p82:donc on gagne encore.}{Why the transfinite closure has to run far enough}
The editors distinguish two different closure problems. If $I$ is already grand devant $\pi$ and $S\subset I$ has size at most $\pi$, a single upper bound $s$ gives the required subposet $S\cup\{s\}$. It has a largest element, so it is grand devant every cardinal. This proves the stated Lemma~1 without any exponent assumption. The longer proof is useful when the ambient poset is only grand devant $\pi_0$; that is the corrected Lemma~$1'$.

Put $\lambda=\pi_0^+$. Starting with $S_0=S$, adjoin at each successor stage an upper bound for every subset of $S_\alpha$ having cardinal at most $\pi_0$, and take unions at limit stages. Suppose $\pi>\pi_0$ and $\pi^{\pi_0}=\pi$. Each successor step adds at most $\pi^{\pi_0}=\pi$ elements. A union of at most $\pi$ sets of size at most $\pi$ again has size at most $\pi$, so every $S_\alpha$, including $S_\lambda$, satisfies the desired bound.

The omitted verification is that $S_\lambda$ is grand devant $\pi_0$. Given $A\subset S_\lambda$ with $|A|\leq\pi_0$, choose a stage containing each of its elements. Since $\lambda=\pi_0^+$ is regular, these at most $\pi_0$ stages are bounded below $\lambda$. Hence $A\subset S_\alpha$ for some $\alpha<\lambda$, and the next stage contains an upper bound for $A$. This is why merely iterating the operation finitely many times, or through a sequence of cofinality at most $\pi_0$, would not prove the assertion.

The exponent hypothesis is sufficient rather than intrinsic. For example, when the desired bound is exactly $\pi_0^+$, enumerate the initial subset in a sequence of length $\pi_0^+$ and recursively choose an upper bound for the elements enumerated so far and the previously chosen bounds. At each proper stage there are at most $\pi_0$ requirements. Adjoining the resulting increasing sequence to the initial subset produces a subposet grand devant $\pi_0$, of cardinal at most $\pi_0^+$, without assuming $(\pi_0^+)^{\pi_0}=\pi_0^+$. We will use this elementary successor-cardinal improvement when returning to the unfinished final lemma.

These are two applications of the same method: control the number of requirements at a successor step, then choose a length whose cofinality catches every small family of requirements. Modern accessibility arguments still use this mechanism. The explicit stronger size assumption in \citet[Lemma~5.4.2.8]{Lurie2009HTT} serves exactly this purpose in the setting of $\infty$-categories.
\end{dcomment}


Ainsi, on a trouvé un deuxième système inductif
\[
(X_{\tilde\imath})_{\tilde\imath\in\tilde I}
\]
du même type que le système initial, ayant comme limite inductive $X$ (qu'il ait même limite inductive est immédiat). Mais qu'a-t-on gagné ? Ceci : dans $\tilde I$ toute partie de cardinal $\le\pi$ et grande devant $\pi_0$ est non seulement \textit{majorée} dans $\tilde I$, mais même y admet une borne supérieure. Cela résulte du % margin: double vertical rule beside the four lines "inductive est immédiat" ... "supérieure. Cela résulte du"

\msp{page 83}
\textbf{Lemme 2 :} \textit{Soit $I$ grand devant $\pi>\pi_0$, avec $\pi^{\pi_0}=\pi$} \texttt{[il suffit $\pi\ge\pi_0$]} \textit{(cf. lemme 1). Soit $(I'_\alpha)_{\alpha\in A}$ un ensemble de parties de $I$ qui sont dans $\tilde I$,} i.e.

\begin{itemize}
\item[a)] \textit{de cardinal $\le\pi$, et}
\item[b)] \textit{grandes devant $\pi_0$,}
\end{itemize}

\textit{et supposons que} $\card A\le\pi$ \textit{et que $A$, ordonné par inclusion, soit grand devant $\pi_0$. Alors $I'=\bigcup_{\alpha\in A} I'_\alpha$ est encore dans $\tilde I$,} i.e. \textit{$I'$ est de cardinal $\le\pi$ (car de cardinal $\le\pi^2=\pi$), et de plus, grande devant $\pi_0$.} (Vérification immédiate.)

Ainsi, on a prouvé ceci :

% margin: single vertical rule beside Proposition 1, from here to the end of the page
\textbf{Proposition 1.} \textit{Soit $E$} \texttt{[une]} \textit{grosse catégorie, avec filtration cardinale} $(\Filt^\pi)_{\pi\ge\pi_0}$. \textit{Soit $\pi$ un cardinal avec}
\[
\pi>\pi_0, \qquad\quad \pi^{\pi_0}=\pi \qquad\quad \textit{(p. ex. $\pi=2^c$, avec $c\ge\pi_0$)} \quad \texttt{[}\pi\ge\pi_0\ \texttt{suffit].}
\]
\textit{Alors pour tout $X$ dans $E$, on peut trouver un système inductif $(X_i)_{i\in I}$, de limite $X$, avec}

\begin{itemize}
\item[a)] \[ \left\{ \begin{aligned} & X_i\in\Filt^\pi \ \ \textit{pour tout $i\in I$,} \\ & \ I \ \textit{grand devant $\pi$,} \end{aligned} \right. \]
\end{itemize}

\textit{satisfaisant aux conditions supplémentaires suivantes :}

\begin{itemize}
\item[b)] \textit{Toute partie $J$ de $I$ de cardinal $\le\pi$, et grande devant $\pi_0$, admet une borne supérieure $i_J$ dans $I$.}
\item[c)] \textit{On a, sous ces conditions}
\end{itemize}

\printedpage{46}
% head: A. Grothendieck : Dérivateurs, Ch. XIII, Édition des Universités de Montpellier II et Paris VII
\msp{page 84}
\[
X_{i_J} \xleftarrow{\ \sim\ } \varinjlim_{i\in J} X_i .
\]
\textit{De plus, si $X\in\Filt^{\pi'}$ avec $\pi'>\pi$, on peut prendre $I$ de cardinal $\le\pi'^{\pi}$ (donc de cardinal $\le\pi'$ si $\pi'^{\pi}=\pi'$, p. ex. si $\pi'=2^{c'}$, avec $c'\ge\pi$).}

\textbf{Définition.} On appelle un système inductif $(X_i)_{i\in I}$, satisfaisant les conditions \textit{a}), \textit{b}), \textit{c}) ci-dessus, un \textit{système inductif bien adapté} (à la filtration cardinale donnée, pour les cardinaux $\pi_0$, $\pi$ donnés), et la représentation correspondante
\[
X \xleftarrow{\ \sim\ } \varinjlim X_i
\]
une \textit{représentation bien adaptée} de $X$ en termes d'objets de $\Filt^\pi$.

\dgrblock{dgrb731}{A presentation that is continuous on its small pieces}
\begin{dcomment}{exposition}{p84:représentation bien adaptée}{A presentation that is continuous on its small pieces}
Grothendieck wants the indexing set to remember small colimits already performed. An arbitrary directed presentation has upper bounds, but its value at an upper bound need not equal the colimit of the stages below it. Replacing indices by small directed subsets builds this continuity into the presentation.

Take $X=\varinjlim_{i\in I}X_i$ as in condition~(c), and let $P$ consist of the subposets $S\subset I$ of size at most $\pi$ which are grand devant $\pi_0$. Order $P$ by inclusion and put $Z_S=\varinjlim_{s\in S}X_s$. The ambient hypothesis on $I$ gives an upper bound for the union of at most $\pi$ members of $P$; adjoining that element produces another member of $P$. Thus $P$ is grand devant $\pi$. Condition~(b) gives $Z_S\in\Filt^\pi(E)$, while the elementary bound on subsets gives $|P|\leq|I|^\pi$.

The equality of the two colimits can be verified directly, without a cofinality convention. A compatible family of maps $Z_S\to Y$ gives maps $X_i\to Y$ by using the singleton subset $\{i\}$. For $i\leq j$, the two-element subset $\{i,j\}$, which has largest element $j$, ensures compatibility. Conversely a cocone on the $X_i$ induces one on every $Z_S$, and these two operations are inverse. Hence $\varinjlim_P Z_S\simeq X$.

Now let $A\subset P$ be grand devant $\pi_0$ and have cardinal at most $\pi$. Its union $S=\bigcup_{T\in A}T$ has size at most $\pi$. For a subset $B\subset S$ of size at most $\pi_0$, choose one $T_b\in A$ containing each $b$. An upper bound $T\in A$ for these members contains $B$, and an upper bound for $B$ inside $T$ belongs to $S$. Thus $S\in P$, and $S$ is the supremum of $A$. A cocone on the diagrams over all $T\in A$ is exactly a cocone over their union, so
\[
 Z_{\sup A}=Z_S\simeq\varinjlim_{T\in A}Z_T.
\]
This proves the continuity clause in Proposition~1, not merely the existence of suprema in its indexing poset.

When all filtered colimits exist and the filtration is strict in the bounded sense of footnote~41, the same proof uses ordinary directed subposets $S$ and ordinary directed families $A$. A union of a directed family of directed subposets is directed. A subset of $I$ of size at most $\pi$ can be enlarged to a directed subset of that same size by adding upper bounds for finite subsets through countably many stages. This yields Proposition~1 bis. The word ``perfectly'' records the stronger continuity, with no change to the cardinality bound on the partial diagrams. The inverse-limit sign in the preceding display on manuscript page~85 is a source slip: the construction uses a colimit.
\end{dcomment}


\textbf{NB} Quand $E$ est stable par petites $\varinjlim$ filtrantes, la situation se simplifie considérablement. On peut alors prendre
\[
\tilde I \;=\; \{I'\subset I \mid \card I'\le\pi,\ I' \text{ filtrant}\}
\]
(sans s'encombrer de la condition $I'$ grand devant $\pi_0$) et définir encore le système inductif $(X_{I'})_{I'\in\tilde I}$ par
\[
X_{I'} \;=\; \varinjlim_{i\in I'} X_i \;.
\]
On a encore
\msp{page 85}
\[
\card\tilde I \;\le\; (\card I)^\pi \;\le\; \pi'^{\pi} \qquad\quad \text{si }\ \card I\le\pi'^{\pi}
\]
(notamment, si $X\in\Filt^{\pi'}$, et qu'on prenne $I$ comme indiqué dans la condition $c$ des filtrations cardinales). De plus, cette fois, le fait que $\tilde I$ soit grand devant $\pi$ est trivial, on a même mieux, à part que $I$ soit filtrante :

\begin{itemize}
\item[$(*)$] \textit{Toute} partie $(I'_\alpha)_{\alpha\in A}$ de $\tilde I$ filtrante et de cardinal $\le\pi$ admet une borne supérieure $\tilde\imath_A$,
\end{itemize}

savoir la réunion (qui est bien filtrante). De plus, le système inductif $(X_{\tilde\imath})_{\tilde\imath\in\tilde I}$ a la propriété plus forte que
\[
X_{\tilde\imath_A} \xleftarrow{\ \sim\ } \varprojlim_{\tilde\imath\in A} X_{\tilde\imath} \;. % [sic] \varprojlim (arrow under lim pointing left) -> \varinjlim, inductive limit of the inductive system X_ĩ over the filtrant part A
\]
Ainsi on a la

\textbf{Proposition 1 bis.} \textit{Supposons $E$ stable par petites $\varinjlim$ filtrantes, et} $\Filt^\pi$ \textit{stable par celles-ci, où $\pi\ge\pi_0$ est donné. Alors pour tout $X$ dans $E$, on peut trouver un système inductif $(X_i)_{i\in I}$ de limite inductive $X$, avec}

\begin{itemize}
\item[a)] $X_i\in\Filt^\pi$ \textit{pour tout $i\in I$, $I$ ensemble ordonné grand devant $\pi$.}
\item[b)] \textit{Toute partie $J$ de $I$ filtrante de cardinal $\le\pi$ admet une borne supérieure $i_J$ dans $I$.}
\item[c)] \textit{Sous ces conditions}
\end{itemize}

\printedpage{47}
% head: A. Grothendieck : Dérivateurs, Ch. XIII, Édition des Universités de Montpellier II et Paris VII
\msp{page 86}
\[
X_{i_J} \;=\; \varinjlim_{i\in J} X_i \;.
\]
De tels systèmes inductifs seront dits \textit{“parfaitement adaptés”} (à la filtration cardinale donnée, pour le cardinal $\pi$ donné – $\pi_0$ n'y figure plus, si ce n'est comme début de la filtration cardinale sans plus).

\dgrblock{dgrb732}{Sets and modules as perfectly adapted presentations}
\begin{dcomment}{examples}{p86:De tels systèmes inductifs seront dits}{Sets and modules as perfectly adapted presentations}
The basic example makes every condition concrete. In $E=\mathrm{Set}$ put $\Filt^\pi(E)=\{S:|S|\leq\pi\}$. A skeleton is a set of cardinals, so the stage is essentially small. A colimit of a diagram with at most $\pi$ objects, each of size at most $\pi$, is a quotient of a disjoint union of size at most $\pi\cdot\pi=\pi$. Thus the required bounded closure holds, even for arbitrary colimits of that size.

For a set $X$ consider
\[
 I=[X]^{\leq\pi},\qquad X_A=A\quad(A\in I),
\]
with inclusion maps. The union of at most $\pi$ members has size at most $\pi$, so $I$ is grand devant $\pi$. Every nonempty directed family $J\subset I$ of size at most $\pi$ has supremum $\bigcup J$, and its colimit in sets is this same union. The colimit of the entire diagram is $X$, since the singletons cover it. If $|X|\leq\pi'$, then $|I|\leq(\pi')^\pi$ for infinite $X$; finite $X$ causes no difficulty. This is a perfectly adapted representation with the precise bound appearing in the source.

A set $A$ of size at most $\pi$ is $\pi^+$-presentable. To factor a map $A\to\varinjlim_IY_i$ through one stage, choose representatives for its at most $\pi$ values, and then one upper bound for their indices. If two such factorizations induce the same map, each of their at most $\pi$ component equalities holds at some later stage; another upper bound makes the maps equal. Both surjectivity and injectivity of the Hom comparison are used. They are the categorical counterparts of choosing generators and then checking relations.

For modules over a ring $R$, fix an infinite $\pi\geq|R|$. Use modules whose underlying sets have size at most $\pi$. They form an essentially small class, and bounded filtered colimits preserve the size bound. Every module $X$ is the union of its submodules of size at most $\pi$. Any set of at most $\pi$ elements generates such a submodule: a linear combination uses finitely many coefficients, so there are at most $\pi$ combinations. The poset of these submodules is grand devant $\pi$, and directed unions of at most $\pi$ of them again have size at most $\pi$. This gives the same perfectly adapted presentation.

The module presentability assertion also includes relations. A module of size at most $\pi$ has a surjection from a free module on at most $\pi$ generators. That free module and its kernel have size at most $\pi$, so the kernel is generated by at most $\pi$ elements. Factoring a homomorphism through a $\pi^+$-directed colimit requires at most $\pi$ choices for generators and at most $\pi$ equations for relations; a common later stage handles them all. Equality of two factorizations is checked on the generators. Thus the chosen stage consists of $\pi^+$-presentable modules.

These examples also locate an omission in the glossary. There, ``strict filtration'' is stated as closure of each stage under all filtered colimits, with no bound on the indexing cardinal. Together with condition~(c), that statement would force $E=\Filt^\pi(E)$ for every $\pi$, hence would make $E$ essentially small. The definition actually used here is the bounded one in footnote~41: closure under filtered diagrams of cardinal at most $\pi$. For instance, the union of $\pi^+$ many successively larger subsets of a set of cardinal $\pi^+$ shows why the bound cannot be discarded.
\end{dcomment}


Comme pratiquement, on aura à manier surtout des $E$ stables par petites $\varinjlim$ filtrantes, et des filtrations cardinales stables en conséquence (renforcement de la condition \textit{b}) des filtrations cardinales), c'est le cas, et la notion de \textit{représentation parfaitement adaptée} de $X$ comme $\varinjlim\limits_{I} X_i$ ($X_i\in\Filt^\pi$), que je mettrai en avant. Le cas “général”, où on suppose seulement \texttt{[que]} $E$ satisfait $L_{\pi_0}$, est plus alambiqué et se traite moins joliment (\footnote[41]{On parlera de filtration cardinale \textit{stricte,} si les $\underline{\Filt}^\pi$ sont stables par $\varinjlim$ filtrantes de cardinal $\le\pi$ (et que ces limites filtrantes petites existent dans $E$).}).

\textbf{Définition 2.} Soit $E$ munie d'une filtration cardinale stricte $(\underline{\Filt}^\pi)_{\pi\ge\pi_0}$. Soit $M$ une partie de $\Ob E$, $\underline{M}$ la sous-catégorie pleine de $E$ qu'elle définit, et soit $\pi$ un cardinal $\ge\pi_0$. On dit que $M$ est \textit{$\pi$-accessible,} où que $\underline{M}$ est \textit{$\pi$-accessible}, % [sic] où que -> ou que, disjunction "ou" not the relative "où"
\msp{page 87}
si les conditions suivantes sont satisfaites. Pour tout système inductif
\[
(X_i)_{i\in I}
\]
parfaitement adapté à la filtration cardinale pour $\pi$ (donc $I$ ensemble ordonné filtrant tel que toute partie filtrante $J$ de cardinal $\le\pi$ ait une borne supérieure $i_J$, et $X_{i_J}\xleftarrow{\ \sim\ }\varinjlim\limits_{j\in J} X_j$) ; soit $I_M$ l'ensemble des $i\in I$ tels que $X_i\in M$. Ceci posé, on veut ceci : pour que $X$ ($=\varinjlim X_i$) soit dans $M$, il faut et il suffit que $I_M$ soit cofinal dans $I$ (de sorte que $X$ apparaît comme $\varinjlim$ filtrante, avec ensemble d'indices $I_M$ grand devant $\pi$, d'éléments de $M\cap\underline{\Filt}^\pi\overset{\text{déf}}{=}M^\pi$).

\msp{page 88}
J'ai envie de reprendre la définition de la $\pi$-accessibilité de $M$, pour une filtration cardinale donnée, à coups de conditions équivalentes. Pour tout système inductif $(X_i)_{i\in I}$ dans $M$ \texttt{[plutôt dans $E$]}, je garde la notation
\[
I_M \;=\; \{i\in I \mid X_i\in M\}.
\]
Je pose, pour tout cardinal $\pi\ge\pi_0$,
\[
M^\pi \;=\; M\cap\Filt^\pi .
\]

\begin{itemize}
\item[\circled{a$_\pi$}] $M$ est l'ensemble des objets de la forme
\[
\varinjlim_{i\in I} X_i \;,
\]
où les $X_i$ sont dans $M^\pi$, et $I$ est grand devant $\pi$.
\end{itemize}

\printedpage{48}
% head: A. Grothendieck : Dérivateurs, Ch. XIII, Édition des Universités de Montpellier II et Paris VII

\begin{itemize}
\item[\circled{b$_\pi$}] Soit $(X_i)_{i\in I}$ un système inductif dans $M$ \texttt{[plutôt dans $E$]} parfaitement adapté à la filtration cardinale pour $\pi$, \textit{i.e.} $I$ filtrante et toute partie filtrante $J$ de $I$ de cardinal $\le\pi$ admet une borne supérieure $i_J$, et $X_{i_J}\xleftarrow{\ \sim\ }\varinjlim\limits_{J} X_i$, et enfin $X_i\in\Filt^\pi$ pour tout $i\in I$. Soit $X=\varinjlim X_i$. Pour que l'on ait $X\in M$, il faut et il suffit que $I_M$ soit cofinal dans $I$.
\item[\circled{c$_\pi$}] $\underline{M}^\pi=\underline{M}\cap\underline{\Filt}^\pi$ est une sous-catégorie de $\underline{M}$ génératrice par épimorphismes stricts, $\underline{M}$ est stable par $\varinjlim$ grandes devant $\pi$, enfin pour tout $X\in M$, $\underline{M}^\pi/X$ est une catégorie filtrante.
\end{itemize}

\msp{page 89}
Je présume que ces conditions sont équivalentes.

\dgrblock{dgrb733}{Three different ways to recognize a part}
\begin{dcomment}{exposition}{p89:Je présume que ces conditions sont équivalentes.}{Three different ways to recognize a part}
The purpose of the new definition is to recognize membership in a class $M$ from small approximations. This would be particularly useful when $E$ is an arrow category and $M$ is a class of weak equivalences. Grothendieck proposes three tests, but their logical strengths differ.

The first, $a_\pi$, says that $M$ is exactly the class of $\pi^+$-directed colimits of objects of $M^\pi=M\cap\Filt^\pi(E)$. It includes a closure assertion as well as an existence assertion. The second, $b_\pi$, says that \emph{every} perfectly adapted presentation detects membership by a cofinal set of $M$-stages. It therefore imposes a condition on the individual stages of a chosen presentation, beyond the existence of some good presentation. The third, $c_\pi$, combines strict generation, closure under $\pi^+$-directed colimits, and ordinary filteredness of $M^\pi/X$.

At a useful cardinal, with $M^\pi$ consisting of $\pi^+$-presentable objects, $a_\pi$ is the familiar description by an $\operatorname{Ind}_{\pi^+}$-completion. Indeed the canonical colimit functor out of that completion is fully faithful; the Hom calculation is the one in \citet[I, 9.18(a)]{SGA4I1972}. It follows in particular that the resulting full subcategory is closed under $\pi^+$-directed colimits. This explains the modern accessibility of the inclusion, but it does not make $b_\pi$ automatic.

There are two distinct issues in the proposed equivalence. Ordinary filteredness controls finite diagrams, whereas $a_\pi$ asks for compatibility of up to $\pi$ constraints at once. Separately, two cofinal collections of small approximations need not meet: one may pick only successor stages and another only limit stages. Continuity of the second presentation does not force membership to pass to its countable colimits. The following examples isolate these two failures while keeping all objects in the small stage genuinely presentable.
\end{dcomment}


\dgrblock{dgrb734}{Two counterexamples to the proposed equivalence}
\begin{dcomment}{development}{p89:Je présume que ces conditions sont équivalentes.}{Two counterexamples to the proposed equivalence}
The distinction just made can be checked in a category with no nontrivial parallel arrows. This removes any possible ambiguity about coherence and leaves the cardinal issue exposed. We give the examples as our own verification of the conjecture on page~89. We did not find these particular examples in the relevant portions inspected of SGA~4 I, \S9, \emph{Les D\'erivateurs} XVIII, \S4.13 and Lemma~4.16.7.11, and the accessible-category references listed below; no priority is asserted.

\textbf{Proposition: the implication $a_\pi\Rightarrow b_\pi$ fails at the same cardinal.}\ours
Take $\pi=\aleph_0$ and let $E$ be the poset $\mathcal P(\omega_1)$, viewed as a category under inclusion. Give it the strict filtration by cardinality, $\Filt^\rho(E)=\{S\subset\omega_1:|S|\leq\rho\}$. Let
\[
 M=\{\alpha+1:\alpha<\omega_1\}\cup\{\omega_1\},
\]
where ordinals denote their sets of predecessors. Then $a_\pi$ and $c_\pi$ hold, while $b_\pi$ fails. Every object of $\Filt^\pi(E)$ is $\pi^+$-presentable in $E$.

\emph{Proof.} Colimits in $E$ are unions. Each filtration stage is small, since it is a subset of the universe-small set $\mathcal P(\omega_1)$. A union of at most $\rho$ sets of size at most $\rho$ again has size at most $\rho$, proving the bounded closure axiom for every infinite $\rho$. Each $S\subset\omega_1$ is the union of its subsets of size at most $\rho$; their poset is $\rho^+$-directed because a union of at most $\rho$ such subsets is another one. If $|S|\leq\rho'$ and $\rho'>\rho$, the number of indices is at most $(\rho')^\rho$; the finite cases satisfy this bound as well. This verifies the full filtration axiom~(c), including its size estimate. Finally, in a $\pi^+$-directed union every countable subset of the union is contained in one stage. Since Hom sets in $E$ have at most one element, this factorization assertion proves presentability of all the countable objects.

The objects of $M^\pi$ are exactly the countable successor ordinals. Each is a constant directed colimit of such objects, and $\omega_1$ is the union of all countable successor ordinals; their ordered set is countably directed. Conversely, consider a countably directed family of countable successor ordinals. If its union $\beta$ is countable and is not attained, choose a countable cofinal subfamily. A single member bounding this subfamily would contain its union, contradicting nonattainment. Therefore a countable union arising this way is an attained successor ordinal. An uncountable union is $\omega_1$. This proves $a_\pi$ and closure of $M$ under countably directed colimits.

For $X\in M$, the slice $M^\pi/X$ is a nonempty linearly ordered set, hence is filtered. Its union is $X$. In a poset this says exactly that the canonical family of arrows from $M^\pi$ is strictly epimorphic: any common upper bound of that family in $M$ is an upper bound of $X$, and the required factorization is unique. Thus $c_\pi$ holds as well.

Now index a diagram by the nonzero countable limit ordinals $\delta<\omega_1$, and set $X_\delta=\delta$. The index is countably directed: the supremum of any countable family is still countable, and a larger countable limit ordinal is an upper bound. Every nonempty directed subset of this index of cardinal at most $\aleph_0$ has its supremum in the index; a supremum of nonzero limit ordinals is again a nonzero limit ordinal. At that supremum the diagram value is the union of the earlier values. The presentation is therefore perfectly adapted, and its colimit is $\omega_1\in M$. Nevertheless no stage belongs to $M$. Its set $I_M$ is empty. This disproves $b_\pi$.\hfill$\square$

\textbf{Proposition: the implication $c_\pi\Rightarrow a_\pi$ also fails.}\ours
In the same filtered category $E$, replace $M$ by
\[
 N=\{S\subset\omega_1:S\text{ is finite}\}\cup\{\omega_1\}.
\]
Then $c_\pi$ holds for $N$, but $a_\pi$ does not.

\emph{Proof.} Here $N^\pi$ consists of the finite subsets. For a finite $X$, the slice $N^\pi/X$ has largest element $X$; for $X=\omega_1$, it is the directed set of finite subsets, whose union is $\omega_1$. These unions prove strict generation, as in the first example. A countably directed family of finite subsets has finite union. Otherwise choose countably many distinct elements of the union and a stage containing each; an upper bound for these countably many stages would be a finite set containing infinitely many elements. A countably directed diagram in $N$ that contains $\omega_1$ has union $\omega_1$. Thus $N$ is closed under the required colimits, proving $c_\pi$.

But no countably directed colimit of objects of $N^\pi$ can equal $\omega_1$, by the same finite-union argument. Hence $a_\pi$ fails.\hfill$\square$

The examples are legitimate $\mathfrak U$-categories. If one requires the ambient category itself to be non-small, take $E\times\mathrm{Set}$ with
\[
 \Filt^\rho(E\times\mathrm{Set})=
 \{(S,T): |S|\leq\rho,\ |T|\leq\rho\}.
\]
A small skeleton is obtained by retaining the actual subsets $S\subset\omega_1$ and replacing $T$ by a cardinal. Colimits are computed in each coordinate, so bounded closure follows from the union calculation above and the set-colimit calculation in the block ``Sets and modules as perfectly adapted presentations.'' For condition~(c), use the product poset $[S]^{\leq\rho}\times[T]^{\leq\rho}$. It is $\rho^+$-directed and has colimit $(S,T)$. If both coordinates have size at most $\rho'>\rho$, its cardinal is at most $(\rho')^\rho\cdot(\rho')^\rho=(\rho')^\rho$. Countable objects are $\aleph_1$-presentable by the same simultaneous choice of representatives and equations in the two coordinates. Restrict the second coordinate of $M$ or $N$ to the empty set. All presentations, strict generating families and counterexample diagrams used above then have empty second coordinate, so every preceding verification remains valid. The repair to $c_\pi$ is substantial: one needs a $\pi^+$-filtered canonical slice with the asserted canonical colimit, rather than merely an ordinary filtered slice. The first example shows that even this repaired generation condition does not recover the universal-presentation condition $b_\pi$ at the same rank.
\end{dcomment}


$b_\pi\Longrightarrow a_\pi$. On suppose \textit{b}) vérifié, on veut prouver \textit{a}). D'abord que tout $X\in M$ est de la forme dite $\varinjlim\limits_{I} X_i$, avec les $X_i\in M^\pi$ et $I$ grand devant $\pi$. Par \texttt{[la]} proposition 1, on sait qu'on a une présentation parfaite adaptée pour $\pi$, $X=\varinjlim\limits_{i\in I} X_i$. Par \textit{b}) “il faut”, la partie $I_M$ de $I$ est cofinale dans $I$, donc comme $I$ est grand devant $\pi$, $I_M$ aussi, et on aura (posant $J=I_M$)
\[
X \;=\; \varinjlim_{i\in J} X_i, \qquad\quad X_i\in M^\pi,\ \ J \text{ grand devant } \pi,
\]
O.K. Inversement, prouvons que si $X$ se met sous la forme dite dans \textit{a}), il est dans $M$. Cela va résulter du

\textbf{Lemme 1 :} \textit{Si la condition} b) (\textit{ou} a)) \textit{est satisfaite, $M$ est stable par limites inductives grandes devant $\pi$.}

Il faut prouver que si
\msp{page 90}
\[
X \;=\; \varinjlim_{i\in I} X_i \;,
\]
avec $I$ grand devant $\pi$, les $X_i$ dans $M$, alors $X\in M$. La condition \textit{b}) “il suffit” l'assure quand on suppose de plus les $X_i$ dans $\Filt^\pi$ (donc dans $M^\pi$), de sorte que l'on aura $I_M=I$, donc $I_M$ cofinal, donc $X\in M$ (\footnote[42]{Pas tout à fait – on y suppose plus, sur le système inductif. On ne peut se dispenser, même dans le cas $X_i\in M^\pi$ $\forall\, i\in I$, faire le raisonnement des pages suivantes.}). Si on ne suppose pas les $X_i$ dans $\Filt^\pi$, il faut travailler un peu plus.

Passer tout de suite à la partie de la démonstration ci-dessous, page 96, qui prouve \textit{b}) $\Longrightarrow$ \textit{a}). Le reste de la démonstration va prouver le

\textbf{Lemme 2.} \textit{Si tout élément de Filt$^\pi$ est $\pi$-accessible, alors} b$_\pi$) \textit{(et même} a$_\pi$\textit{) implique que $M$ est stable par $\varinjlim$ grandes devant $\pi$.}

% ---- from Der-ChXIII-25-01-04_p051-055.tex
\printedpage{49}
% head: A. Grothendieck : Dérivateurs, Ch. XIII, Édition des Universités de Montpellier II et Paris VII
% foot: Version 25 janvier 2004

\msp{page 91}Soit
\[ \mathcal{C} \;=\; \underline{M}^\pi/\mathcal{X} \]
la catégorie suivante. Ses objets sont les triples
\[ (Z,i,\alpha) \quad \text{tels que} \quad Z\in M^\pi,\ i\in I,\ \alpha : Z\longrightarrow X_i \ , \]
les flèches
\[ (Z,i,\alpha) \;\longrightarrow\; (Z',i',\alpha') \]
sont les couples $(u,\lambda)$ avec
\[ u : Z\longrightarrow Z', \qquad \lambda : i\longrightarrow i', \qquad (\textit{i.e.}\ \text{on a } i\le i' \text{ et } \lambda \text{ est l'unique flèche de } i \text{ dans } i'), \]
tels que le diagramme
\[
\begin{tikzcd}[row sep=large, column sep=large]
Z \arrow[r,"\alpha"] \arrow[d,"u"'] & X_i \arrow[d,"\text{transition } \lambda_{i',i}"] \\
Z' \arrow[r,"\alpha'"] & X_{i'}
\end{tikzcd}
\]
commute. On peut donc définir $\mathcal{C}$ comme le produit fibré dans le carré cartésien
\[
\begin{tikzcd}[row sep=large, column sep=large]
 & \mathcal{C} \arrow[r,"\text{cofibré}"] \arrow[d] \arrow[dl,dotted,"\mathcal{Z}"'] & I \arrow[d,"\mathcal{X} : X \longmapsto X_i"] \\ % [sic] X \longmapsto X_i -> i \longmapsto X_i, the functor \mathcal{X} goes from I, so the variable should be i
\underline{M}^\pi & \underline{M}^\pi/E \arrow[l,dotted,"\text{source}"'] \arrow[r,"\text{but}","\text{cofibré}"'] & E.
\end{tikzcd}
\]
La fibre $\mathcal{C}_i$ de $\mathcal{C}$ en $i\in I$ est $\mathrm{Filt}^\pi/X_i$ \texttt{[plutôt} $M^\pi/X_i$\texttt{]}. On a, comme \textit{b}) $\Longrightarrow$ première moitié de \textit{a}), que $\underline{M}^\pi$ est un
\msp{page 92}système de générateurs stricts de $\underline{M}$ \texttt{[cela n'est pas clair sans autres hypothèses sur} $\underline{M}$\texttt{]}, et $X_i\in\operatorname{Ob}\underline{M}$ par hypothèse,
\[ X_i \;=\; \varinjlim_{Z\in\operatorname{Ob}\underline{M}^\pi/X_i} Z \;=\; \varinjlim_{\mathcal{C}_i}\mathcal{Z}|\mathcal{C}_i, \]
donc le système inductif donné $(X_i)_{i\in I}$ sur $I$ à valeurs dans $X$, se déduit du système inductif (non ordonné en général) sur $\mathcal{C}$, à valeurs dans $\underline{M}^\pi$, en ``intégrant sur les fibres'' de $\mathcal{C}\longrightarrow I$. D'ailleurs, comme $\underline{M}^\pi/E\longrightarrow E$ est cofibrée, donc propre, il en est de même de $\mathcal{C}\longrightarrow I$, donc les $\varinjlim$ de $\mathcal{Z}$ sur les fibres $\mathcal{C}_i$ ($=$ $\underline{M}^\pi/X_i$) forment bien un foncteur covariant $I\longrightarrow E$, dont la limite inductive est celle de $\mathcal{Z}$ (\footnote[43]{\textbf{NB} Si on admet la validité de $a$), alors il en résulte déjà que $\varinjlim\limits_{i\in I} X_i \simeq \varinjlim\limits_{\alpha\in C} Z_\alpha$ est dans $M$. (Stabilité de $M$ par $\varinjlim$ grandes devant $\pi$).}).

\printedpage{50}
% head: A. Grothendieck : Dérivateurs, Ch. XIII, Édition des Universités de Montpellier II et Paris VII
% foot: Version 25 janvier 2004

\msp{page 93}Soit, pour tout $i\in I$, $\tilde{\mathcal{C}}_i$ l'ensemble ordonné
\[ \tilde{\mathcal{C}}_i = \{\text{sous-catégories } \mathcal{C}'_i \text{ de } \mathcal{C}_i = \underline{M}^\pi/X_i \mid \operatorname{card}\mathcal{Z}(\mathcal{C}'_i)\le\pi,\ \ \mathcal{C}'_i \text{ filtrante}\}. \]
J'admets pour l'instant le

\textbf{Lemme} (\footnote[44]{Utilise l'hypothèse tacite que les éléments de $\mathrm{Filt}^\pi$ sont $\pi$-accessibles.}). \textit{Pour tout $X\in M$, la catégorie $\underline{M}^\pi/X$ est filtrante. (Sous l'hypothèse} b), \textit{ou que tout objet de $M$ est de la forme dite dans} a).\textit{)}

\dgrblock{dgrb735}{Presentability supplies the strict generators}
\begin{dcomment}{exposition}{p93:est filtrante. (Sous l'hypothèse}{Presentability supplies the strict generators}
Grothendieck now tries to replace a diagram of large $M$-objects by one diagram of small $M$-objects. The comma category is the right construction: it records a small object together with the map by which it participates in a particular large object. The editors' reservation on page~92 concerns the claim that a chosen colimit presentation already makes \emph{all} maps from small objects into a canonical colimit presentation. Presentability supplies the missing comparison.

Here is the precise assertion. Let $\kappa=\pi^+$, let $A\subset E$ be an essentially small full subcategory of $\kappa$-presentable objects, and suppose
\[
 X\simeq\varinjlim_{j\in J}Y_j,
 \qquad J\ \kappa\text{-directed},\qquad Y_j\in A.
\]
Then $A/X$ is $\kappa$-filtered, and its canonical diagram has colimit $X$. This is the canonical-diagram form of the fully faithful $\operatorname{Ind}_\kappa$ construction in \citet[I, 9.18(a)]{SGA4I1972}.

To prove the assertion, take any diagram of fewer than $\kappa$ morphisms in $A/X$. Each of its objects is a map $T\to X$ with $T$ $\kappa$-presentable, so it factors through a $Y_j$. Choose one common upper bound for the indices used. For every arrow in the given diagram, the two resulting maps into this common stage agree after passage to $X$. Presentability of its source makes that equality hold at some later stage. There are fewer than $\kappa$ equalities, so another common upper bound makes them all hold together. The resulting object $Y_{j'}\to X$ is a cocone in $A/X$. This proves $\kappa$-filteredness, including the case of the empty diagram.

To identify the colimit, let a compatible family $T\to Z$ be given for all objects $(T\to X)$ of $A/X$. Restrict it to the chosen presentation $Y_j\to X$, producing a unique map $h:X\to Z$. Any $T\to X$ factors through a $Y_j$, and compatibility with that factorization says that the prescribed map $T\to Z$ is $h$ composed with $T\to X$. Thus $h$ induces the whole family, uniquely. This proves the universal property of the canonical colimit. In particular its canonical family is strictly epimorphic in the ambient category, hence also in any full subcategory containing $X$ and $A$.

Apply this with $A=M^\pi$. The first half of $a_\pi$, already obtained from $b_\pi$, supplies the required presentations. The tacit hypothesis in footnote~44 supplies presentability. Under those hypotheses the lemma on page~93 is correct, and in fact its slice is $\pi^+$-filtered, a stronger conclusion than printed. The precise missing step is factorization and eventual equality for \emph{arbitrary} maps from $M^\pi$, not an additional assumption about epimorphisms in $M$.

The same argument explains why this construction is stable under a sufficiently directed change of the large object. For a $\kappa$-directed diagram $(X_i)$ in $M$, form the category $C$ of triples $(T,i,T\to X_i)$. A diagram of fewer than $\kappa$ arrows in $C$ can first be moved into one fibre over $I$, and then given a cocone in that fibre by the preceding paragraph. Thus $C$ is $\kappa$-filtered. Its colimit is $\varinjlim_I X_i$, because giving a cocone on $C$ is giving compatible cocones on all the fibres. After the usual cofinal directed-poset replacement, this is a presentation by objects of $M^\pi$. If $a_\pi$ holds, it proves closure of $M$ under $\kappa$-directed colimits. The directed-poset replacement is the ordinary-category case of \citet[Proposition~5.3.1.18]{Lurie2009HTT}; it changes the index, not the colimit.
\end{dcomment}


Il en résulte, comme $X_i\in M$, que $\mathcal{C}_i$ est filtrante, donc la réunion des $\mathcal{C}'_i\in\tilde{\mathcal{C}}_i$ est $\mathcal{C}_i$ ; l'ensemble ordonné $\tilde{\mathcal{C}}_i$ est filtrant, et toute partie $(\mathcal{C}'_{i_\alpha})_{\alpha\in A}$ filtrante et de cardinal $\le\pi$ de $\tilde{\mathcal{C}}_i$ y admet une borne supérieure $\mathcal{C}'_{i_A}$, savoir la réunion des $\mathcal{C}'_{i_\alpha}$. $\tilde{\mathcal{C}}_i$ est donc un ensemble ordonné $\pi$-adapté.

On peut considérer les ensembles ordonnés $\tilde{\mathcal{C}}_i$, pour $i\in I$ variable, comme les fibres d'une application croissante d'ensembles ordonnés
\[ p \;:\; \tilde{\mathcal{C}} \;\longrightarrow\; \mathcal{C}, \]
en écrivant
\[ \left\{ \begin{array}{l}
(\mathcal{C}'_i\subset\mathcal{C}_i) \;\le\; (\mathcal{C}'_j\subset\mathcal{C}_j) \\
\text{si et seulement si } i\le j, \\
\text{et le foncteur canonique de transition } \mathcal{C}_i\longrightarrow\mathcal{C}_j \text{ applique } \mathcal{C}'_i \text{ dans } \mathcal{C}'_j.
\end{array}\right. \]
On voit donc que $\tilde{\mathcal{C}}$ est lui aussi filtrant et grand devant $\pi$.

\msp{page 94}On définit un système inductif sur $\tilde{\mathcal{C}}$, à valeurs dans $\underline{\mathrm{Filt}}^\pi$, en posant
\[ Z_{\mathcal{C}'_i} \;=\; \varinjlim_{\alpha\,\in\,\mathcal{C}'_i\ (\subset\,\mathcal{C}_i\,=\,\underline{M}^\pi/X_i)} Z_\alpha \]
(où un $\alpha\in\mathcal{C}_i$ est formé, je le rappelle, de $Z_\alpha\xrightarrow{u_\alpha}X_i$, avec $Z_\alpha\in M^\pi\subset\underline{\mathrm{Filt}}^\pi$). Ces $Z_{\mathcal{C}'_i}$ ne sont plus forcément dans $M^\pi$, mais seulement dans $\mathrm{Filt}^\pi$. (Alors à quoi ça a servi de définir les $\mathcal{C}_i$ eux-mêmes comme $\underline{M}^\pi/X_i$, au lieu de $\underline{\mathrm{Filt}}^\pi/X_i$ ?)

Le système inductif
\[ (Z_{\mathcal{C}'_i})_{\mathcal{C}'_i\in\tilde{\mathcal{C}}_i} \]
est parfaitement bien adapté à la filtration cardinale pour le cardinal $\pi$, et il est immédiat que sa $\varinjlim$ est isomorphe à $\varinjlim\limits_{\alpha\in\operatorname{Ob}\mathcal{C}_i=\underline{M}^\pi/X_i} Z_\alpha = X_i\in M$ . Donc par \textit{b}), l'ensemble $\tilde{\mathcal{C}}_M$ des $\mathcal{C}'_i\in\tilde{\mathcal{C}}_i$ tels que $Z_{\mathcal{C}'_i}\in M$ (donc $Z_{\mathcal{C}'_i}\in M^\pi$) est cofinal dans $\tilde{\mathcal{C}}_i$. Il s'ensuit donc que l'ensemble $\tilde{\mathcal{C}}_M$ est lui aussi cofinal dans $\tilde{\mathcal{C}}$.

\printedpage{51}
% head: A. Grothendieck : Dérivateurs, Ch. XIII, Édition des Universités de Montpellier II et Paris VII
% foot: Version 25 janvier 2004

\msp{page 95}On vérifie aussi immédiatement que
\[ \varinjlim_{\alpha\in\tilde{\mathcal{C}}} Z_\alpha \;=\; \varinjlim X_i \;=\; X, \]
donc, $\tilde{\mathcal{C}}_M$ étant cofinal,
\[ \varinjlim_{\alpha\in\tilde{\mathcal{C}}_M} Z_\alpha \;=\; X. \]
Au total, on est arrivé à représenter $X$ sous la forme
\[ X \;=\; \varinjlim_{\alpha\in A} X_\alpha, \]
où cette fois les $X_\alpha$ sont \textit{dans} $M^\pi$ (pas seulement dans $M$), et $A$ est grand devant $\pi$. Mais on ne sait pas que toute partie filtrante $B$ de cardinal $\le\pi$ de $A$ admet une borne supérieure $\alpha_B$ dans $A$, et encore moins que $X_{\alpha_B} = \varinjlim\limits_{\alpha\in B} X_\alpha$ (\footnote[45]{Si on suppose $(a_\pi)$, cet argument prouve du moins que $M$ est stable par $\varinjlim$ grandes devant $\pi$.}). Donc on ne peut encore appliquer directement \textit{b}). S'il ne s'agissait que de prouver la deuxième moitié de \textit{a}), où on suppose déjà les $X_i\in M^\pi$ et $I$ grand devant $\pi$, on n'a rien gagné du tout par nos labeurs !

\dgrblock{dgrb736}{The size of an indexing category cannot be replaced by the size of its image}
\begin{dcomment}{development}{p95:Donc on ne peut encore appliquer directement}{The size of an indexing category cannot be replaced by the size of its image}
The partial-colimit construction is meant to keep all intermediate objects within $\Filt^\pi$. For that purpose it must bound the indexing category. The literal condition on page~93, $\operatorname{card}\mathcal Z(C_i')\leq\pi$, bounds only the image of the diagram. That is insufficient, even if its cardinal counts image morphisms as well as image objects.

\textbf{A filtered counterexample to the image bound.}\ours
Let $E=M=\mathrm{Set}$ with the usual cardinal filtration, let $\pi=\aleph_0$, and choose an uncountable set $X$. Inside $M^\pi/X$, take the category $C'$ whose objects are injections $a:n\hookrightarrow X$ from the standard finite sets $n=\{0,\ldots,n-1\}$. A morphism from $a:n\hookrightarrow X$ to $b:m\hookrightarrow X$ is a map $u:n\to m$ satisfying $bu=a$.

This category is filtered. Two objects map into an injection whose image is the union of their two finite images. Parallel arrows are equal, since their composites with the injective target map are equal. The source functor $\mathcal Z:C'\to\mathrm{Set}$ has image inside the countable category of standard finite sets and maps; hence $\operatorname{card}\mathcal Z(C')\leq\aleph_0$. Yet
\[
 \varinjlim_{C'}\mathcal Z\simeq X.
\]
Indeed every element of $X$ occurs through a one-element injection, and two representatives determine the same element precisely when they become equal in a finite union of their images. Equivalently the compatible injections give the usual directed union of the finite subsets of $X$. This colimit is not in $\Filt^{\aleph_0}(E)$.\hfill$\square$

The corrected construction uses filtered \emph{indexing categories} $D\subset C_i$ with at most $\pi$ morphisms, or directed indexing posets of that size mapping to $C_i$. This correction does not require a bound on every Hom set of $C_i$. The subcategories need not be full. Start with at most $\pi$ specified objects and arrows; close under identities and composition, adjoin cones for finite diagrams, and repeat countably often. Each step adds at most $\pi$ data, since $\pi$ is infinite. The resulting filtered subcategory still has at most $\pi$ morphisms. In particular the poset of these small filtered subcategories is grand devant $\pi$: perform this construction on the union of any at most $\pi$ of them.

A directed union of at most $\pi$ such subcategories is again a filtered subcategory of size at most $\pi$ and is their supremum. Colimits over this union are the colimits of the partial colimits. To apply the source's filtration axiom, which is phrased for posets, a filtered category with at most $\pi$ arrows can be replaced cofinally by a directed poset of size at most $\pi$: the usual finite-diagram construction adjoins finitely many compatibility data at each of countably many stages. Thus the corrected partial colimit lies in $\Filt^\pi$, and its presentation over the poset of small subcategories is perfectly adapted.

This also repairs the rest of the argument on pages~93--95. The map from the total poset of pairs $(i,D\subset C_i)$ has target $I$, as its displayed order relation specifies; the printed target $C$ is a slip. To bound at most $\pi$ such pairs, move their indices to one index of $I$ and enlarge the union of their images to a small filtered subcategory of the corresponding fibre. Cofinality of the good partial colimits in every fibre then gives their cofinality in the total poset. The final diagram so obtained need not itself be continuous at all small directed suprema, exactly as Grothendieck observes on page~95. The reindexing on page~96 is therefore still needed.

The counterexample concerns this particular image-size formulation. The corrected small-subcategory construction is the familiar filtered-diagram argument, already used in SGA~4 I, 9.13.1. We did not find the specific counterexample above in the literature checked; it records our verification of the printed hypothesis, not a claim of priority.
\end{dcomment}


\msp{page 96}Mais revenons à ce cas
\[ X \;=\; \varinjlim X_i \;\left|\; \begin{array}{l} X_i\in M^\pi \\[1.5ex] I \text{ grand devant } \pi, \end{array}\right. \]
et introduisons à nouveau
\[ \tilde{I} \;=\; \{I'\subset I \mid I' \text{ filtrant, de cardinal } \le\pi\}, \]
donc cette fois $\tilde{I}$ est bien adapté à $\pi$ (filtrant, et toute partie filtrante de cardinal $\le\pi$ a une borne supérieure), et si on pose
\[ X_{\tilde{\imath}} \;=\; \varinjlim_{i\in I'_{\tilde{\imath}}} X_i, \]
on voit que le système inductif
\[ (X_{\tilde{\imath}})_{\tilde{\imath}\in\tilde{I}} \]
est parfaitement adapté à $\pi$ et la filtration cardinale. De plus, je dis que l'ensemble $\tilde{I}_M$ est \textit{cofinal} dans $\tilde{I}$. En effet, $\tilde{I}_M$ contient le sous-ensemble $\tilde{I}_0$ formé des $I'\in\tilde{I}$ qui ont un plus grand élément -- et $\tilde{I}_0$ est bien cofinal ! Maintenant, on peut appliquer
\printedpage{52}
% head: A. Grothendieck : Dérivateurs, Ch. XIII, Édition des Universités de Montpellier II et Paris VII
% foot: Version 25 janvier 2004
\msp{page 97}le ``il suffit'' dans \textit{b}), qui conclue que $X = \varinjlim X_i\in M$, q.e.d. % [sic] conclue -> conclut, verb form (indicative of conclure)

On a prouvé $b\Longrightarrow a$, sans utiliser l'hypothèse que tout objet de $\mathrm{Filt}^\pi$ est $\pi$-accessible. Mais on a utilisé cette hypothèse pour prouver que $b$ (ou $a$) implique que $M$ est stable par $\varinjlim$ grandes devant $\pi$ (lemme 2, p. 90).

\dgrblock{dgrb737}{The implication that does not need presentable objects}
\begin{dcomment}{exposition}{p97:On a prouvé $b\Longrightarrow a$}{The implication that does not need presentable objects}
The short argument on pages~96--97 is more robust than the intervening attempt to prove closure under large colimits. It proves $b_\pi\Rightarrow a_\pi$ directly, using only the strict filtration and the existence of perfectly adapted presentations. In particular it should not inherit the tacit presentability hypothesis needed for the preceding comma-category construction.

First suppose $X\in M$. Choose a perfectly adapted presentation $X=\varinjlim_I X_i$ at $\pi$. By $b_\pi$, the subset $I_M$ is cofinal. A cofinal subset of a $\pi^+$-directed poset is again $\pi^+$-directed: take an upper bound in the original poset for any at most $\pi$ elements, and then move above it into the cofinal subset. Restricting to $I_M$ therefore presents $X$ as a $\pi^+$-directed colimit of objects of $M^\pi$.

For the other implication in $a_\pi$, suppose $X=\varinjlim_{i\in I}Y_i$ with $Y_i\in M^\pi$ and $I$ grand devant $\pi$. Let $P$ be the poset of nonempty directed subsets $S\subset I$ of cardinal at most $\pi$, and put $Z_S=\varinjlim_S Y_i$. Strictness places each $Z_S$ in $\Filt^\pi$. The union construction makes $(Z_S)_{S\in P}$ perfectly adapted, with colimit $X$.

The key cofinal family consists of the subsets having a largest element. For any $S\in P$, choose an upper bound $s$ of all of $S$ in $I$ and put $S^+=S\cup\{s\}$. Then $S^+\in P$, it contains $S$, and
\[
 Z_{S^+}\simeq Y_s\in M^\pi.
\]
Hence $P_M$ is cofinal in $P$. The sufficiency direction in $b_\pi$ gives $X\in M$. This completes both directions of $a_\pi$.

The earlier inference on page~90 applied $b_\pi$ directly to an arbitrary directed presentation and omitted its continuity hypothesis; footnote~42 explicitly retracts that inference. The argument just given is the source's own repair. After it is established, the additional hypothesis that $M^\pi$ consists of $\pi^+$-presentable objects permits the closure argument through canonical comma categories. These are two different uses of the hypotheses, and retaining that separation prevents an artificial circularity in the proof.
\end{dcomment}


$a\overset{?}{\Longrightarrow} b$. Le ``il suffit'' dans \textit{b}) est conséquence immédiate de \textit{a}). Il reste à prouver que si, dans le système inductif $\pi$-adapté $(X_i)_{i\in I}$ de limite $X$, on a $X\in M$, alors $I_M$ est cofinal dans $I$. Mais par hypothèse \textit{a}), on aura $X = \varinjlim\limits_J Y_j$, avec $J$ grand devant $\pi$, les $Y_j$ dans $M^\pi$. Admettons que $\mathrm{Filt}^\pi\subset E_\pi$, \textit{i.e.} tout élément de $\mathrm{Filt}^\pi$, en particulier les $X_i$, sont $\pi$-accessibles.

Je pense que je n'arriverai pas à prouver $a_\pi\Longrightarrow b_\pi$, mais seulement ceci : Si $\pi'$ est un cardinal $>\pi$ (disons le successeur de $\pi$ dans la suite ordinale des cardinaux), alors $a_\pi\Longrightarrow b_\pi$.

\msp{page 98}Ceci va résulter du

\textbf{Lemme 3.} \textit{Soient $X\simeq\varinjlim\limits_I X_i$, $X\simeq\varinjlim\limits_J Y_j$ deux représentations de $X$ sous forme des $\varinjlim$ filtrantes, avec $I$ grand devant $\pi'$, $J$ grand devant $\pi$, avec $\pi$, $\pi'$ deux cardinaux tels que $\pi'>\pi$. On suppose $(X_i)_{i\in I}$ $\pi'$-adapté,} i.e. \textit{pour toute partie filtrante $I'$ de $I$ de cardinal $\le\pi'$, $I'$ a une borne supérieure $i_{I'}$, et $X_{i_{I'}} = \varinjlim\limits_{i\in I'} X_i$. Alors pour tout $i\in I$, on peut trouver une partie $I'$ de $I$, contenant $i$,} grande devant $\pi$, \textit{de cardinal $\le\pi'$, et une partie $J'$ de $J$, enfin un isomorphisme $I'\xrightarrow[\sim]{\lambda}J'$ et une application strictement croissante $\mu : J'\longrightarrow I'$ et des homomorphismes} \texttt{[fin du chapitre]}

\dgrblock{dgrb738}{Where the manuscript ends}
\begin{dcomment}{exposition}{p98:et une application strictement croissante}{Where the manuscript ends}
Grothendieck has reached the comparison of two presentations when the manuscript breaks off. The final words announce homomorphisms without specifying them. We therefore have a programme for a comparison lemma, not a completed statement whose omitted proof can simply be supplied. The separate larger-cardinal repair is a theorem with its own stated hypotheses; it does not turn this fragment retrospectively into a finished argument.

The material that follows belongs to the edition's apparatus. Its notation and terminology indices refer to manuscript page numbers. Thus the reference ``77'' beside cardinal filtrations points to manuscript page~77, on printed page~42 of this edition. The glossary then restates definitions from the whole chapter, including sections outside these pages. Its axiom list is particularly useful here: it records finite limits and colimits in \textbf{M~0}, two factorizations in \textbf{M~2}, the properness conditions in $\mathbf M''$, and the separate requirement $W=TF\circ TC$ in \textbf{M~2\,bis}. These distinctions explain why the names of the various Quillen structures cannot be replaced indiscriminately by the single modern phrase ``model category''.

The glossary is a guide to those definitions, and should be checked against their explicit formulation in the manuscript. In its account of a strict cardinal filtration, the qualification ``of cardinality at most $\pi$'' has disappeared; manuscript page~86, footnote~41, supplies it. Its final paragraph also says ``well adapted'' where the immediately preceding definition requires ``perfectly adapted''. The audit registers these two discrepancies separately, retaining the printed wording. The bound matters: by the representation axiom every object of $E$ is a filtered colimit of objects in any fixed stage, so unrestricted filtered-colimit closure would identify that stage with all of $E$. The printed glossary would then express a substantially different definition.

Finally, the bibliography after the glossary is the edition's bibliography. It includes the 1996 publication of Verdier's thesis; its presence is not evidence that Grothendieck had consulted that publication while writing these pages in 1990--1991. Our later references and our repairs are placed in the commentary and its own bibliography, where the date and the responsibility for each statement remain visible.
\end{dcomment}


\dgrblock{dgrb739}{A complete comparison after one successor step}
\begin{dcomment}{development}{p98:et des homomorphismes}{A complete comparison after one successor step}
The chapter stops while Grothendieck is constructing a comparison between two presentations. The final Lemma~3 has no completed conclusion or proof, and we do not supply words to its unfinished sentence. His preceding proposal does, however, suggest a precise statement which can be proved under explicit presentability assumptions. The printed implication on page~97 repeats $b_\pi$; the surrounding reference to a larger cardinal indicates $b_{\pi'}$ as the intended target.

\textbf{Successor comparison.}\ours
Let $E$ carry a strict cardinal filtration and let $\pi\geq\pi_0$ be infinite. Put $\kappa=\pi^+$ and $\mu=\kappa^+$. Assume that objects of $\Filt^\pi(E)$ are $\kappa$-presentable and that objects of $\Filt^\kappa(E)$ are $\mu$-presentable. Let $M$ be a replete full subcategory satisfying $a_\pi$. Then $M$ satisfies $b_\kappa$, that is, $b_{\pi^+}$ in the manuscript's notation.

We prove the statement rather than infer it from a general assertion that increasing a cardinal preserves every accessibility property. The literature checked includes SGA~4 I, 9.18, Grothendieck's later Chapter~XVIII, \S4.13 and Lemma~4.16.7.11, and \citet[\S8]{LiebermanRosickyVasey2020}. The argument below is a formulation for the particular continuous presentations used here; its use of alternating factorizations belongs to the standard comparison method for filtrations. We make no priority claim.

\emph{Proof.} The canonical-slice argument above shows first that $M$ is closed under $\kappa$-directed colimits. Indeed every object of $M$ has a presentation by objects of $M^\pi$, these objects are $\kappa$-presentable, and a $\kappa$-directed diagram of such presentations can be combined into one $\kappa$-filtered diagram of their small terms. The closure half of $a_\pi$ then applies.

We next obtain a more highly directed presentation of each $X\in M$, while keeping its terms in $M$. Write $X=\varinjlim_{j\in J}Y_j$ as in $a_\pi$. Any subset $S\subset J$ of cardinal at most $\kappa$ lies in a $\kappa$-directed subposet $K\subset J$ of size at most $\kappa$. To see this, first adjoin any element of the nonempty poset $J$ if $S$ is empty, then enumerate $S$ in length at most $\kappa$, repeating elements if necessary. At stage $\alpha<\kappa$, choose an upper bound in $J$ for the current enumerated element and all previously chosen bounds. There are at most $\pi$ of them. The increasing sequence of bounds, together with $S$, is $\kappa$-directed: any fewer than $\kappa$ of its elements is accounted for before some stage, since $\kappa$ is regular.

Let $H$ be the poset of all these $\kappa$-directed subposets of $J$ of size at most $\kappa$. It is $\mu$-directed. The union of at most $\kappa$ of its members has size at most $\kappa$, and the construction just given places that union in another member of $H$. Put
\[
 Z_K=\varinjlim_{j\in K}Y_j\qquad(K\in H).
\]
The closure already proved gives $Z_K\in M$, and strictness of the filtration gives $Z_K\in\Filt^\kappa(E)$. Thus $Z_K$ is $\mu$-presentable by hypothesis. The union-of-subdiagrams argument gives $X\simeq\varinjlim_{K\in H}Z_K$.

Now take a perfectly adapted presentation $X=\varinjlim_{i\in I}X_i$ at cardinal $\kappa$. Its index $I$ is $\mu$-directed and its objects are $\mu$-presentable. Fix $i_0\in I$. We will construct increasing sequences $(i_\alpha)_{\alpha<\kappa}$ in $I$ and $(K_\alpha)_{\alpha<\kappa}$ in $H$, together with maps over $X$
\[
 X_{i_\alpha}\xrightarrow{u_\alpha}Z_{K_\alpha}
       \xrightarrow{v_\alpha}X_{i_{\alpha+1}},
\]
such that
\[
 v_\alpha u_\alpha=X(i_\alpha\leq i_{\alpha+1}),
\]
and, for $\alpha<\beta<\kappa$,
\[
 u_\beta X(i_{\alpha+1}\leq i_\beta)v_\alpha
       =Z(K_\alpha\subset K_\beta).                 \tag{SC}
\]
The first equality makes a trip through the other presentation agree with a transition in the $X$-diagram; the second does the same for the $Z$-diagram. They will give inverse maps on colimits.

Here are all the inductive choices. At a limit stage $\beta<\kappa$, choose $i_\beta$ above the earlier $i_\alpha$ and $i_{\alpha+1}$; at a successor it has already been chosen. Since $X_{i_\beta}$ is $\mu$-presentable, its map to $X=\varinjlim_H Z_K$ factors through some $Z_K$. Enlarge $K$ above all the previously chosen $K_\alpha$. For each $\alpha<\beta$, the two maps from $Z_{K_\alpha}$ to this stage occurring in (SC) agree after passage to $X$. Presentability of $Z_{K_\alpha}$ makes that equality hold at some later $H$-stage. There are fewer than $\kappa$, hence fewer than $\mu$, such equalities. The $\mu$-directedness of $H$ permits one enlargement making all of them hold. Call the resulting index $K_\beta$ and the factorization $u_\beta$.

Next factor $Z_{K_\beta}\to X$ through an $X$-stage, using its $\mu$-presentability. Enlarge that stage above $i_\beta$. The composite with $u_\beta$ and the transition out of $X_{i_\beta}$ have equal images in $X$; presentability of $X_{i_\beta}$ allows a further enlargement on which they are equal. This gives $i_{\beta+1}$ and $v_\beta$. The same process starts at $\beta=0$, when there are no equations (SC) from earlier stages. Every choice takes place in a small diagram, and the regularity and directedness just used justify every limit stage.

Equations (SC) imply that the $u_\alpha$ define a map between the two chain colimits: compose (SC) with $u_\alpha$ and use $v_\alpha u_\alpha=X(i_\alpha\leq i_{\alpha+1})$. They also imply that the $v_\alpha$, followed by the structural maps to the $X$-chain colimit, are compatible. The two resulting maps are inverse, by the same displayed equalities. Consequently
\[
 \varinjlim_{\alpha<\kappa}X_{i_\alpha}
 \simeq\varinjlim_{\alpha<\kappa}Z_{K_\alpha}\in M.
\]
The membership assertion uses closure of $M$ under $\kappa$-directed colimits, since the regular ordinal $\kappa$ is a $\kappa$-directed poset.

The perfectly adapted $X$-diagram has a supremum $i_* =\sup_{\alpha<\kappa}i_\alpha$ and is continuous there: the image of this chain is a directed subset of $I$ with at most $\kappa$ elements. Therefore $X_{i_*}\simeq\varinjlim_{\alpha<\kappa}X_{i_\alpha}$ belongs to $M$, and $i_*\geq i_0$. Thus $I_M$ is cofinal whenever $X\in M$. Conversely, if $I_M$ is cofinal then it is $\mu$-directed, hence $\kappa$-directed, and closure of $M$ gives $X=\varinjlim_{I_M}X_i\in M$. This proves $b_\kappa$.\hfill$\square$

The cardinal shift has a definite function. It permits a comparison chain of length $\pi^+$ whose final colimit remains a single stage of the larger adapted presentation. At the original cardinal, that chain is too large for the continuity clause. The counterexample with successor and limit ordinals shows that this obstruction can be real. No continuum hypothesis entered the proof, but presentability at both filtration levels did; the filtration axioms alone do not provide it.
\end{dcomment}


\dgrblock{dgrb740}{Canonical slices, later accessibility, and weak equivalences}
\begin{dcomment}{later}{p98:et des homomorphismes}{Canonical slices, later accessibility, and weak equivalences}
Grothendieck returned to these questions at much greater length in Chapter~XVIII. Definition~4.13.1 there asks that a full subcategory be replete, accessible, and accessibly included in the ambient accessible category. Lemma~4.16.7.11 tests membership by cofinality of the small canonical slice inside the ambient small-object slice. That is an intrinsic form of the generation condition, with presentability built into the indexing category. It does not assert that every continuous presentation already has cofinally many stages in the subcategory at the original rank. The two powerset examples above therefore leave the later criterion intact. The edition consulted is dated 17~December~2025; the cited lemma lies on manuscript pages~488--490 \citep{Grothendieck1990XVIII}.

In present terminology, accessibility and accessible inclusion are distinct conditions, and a rank belongs to the statement. The definition and the description as a completion under directed colimits are given by \citet[Definition~2.1 and Theorem~2.26]{AdamekRosicky1994}. The earlier monograph of \citet{MakkaiPare1989} situates these constructions within categorical model theory. For the higher-categorical formulation, \citet[Proposition~5.4.2.2]{Lurie2009HTT} characterizes accessible $\infty$-categories by small generating subcategories of $\kappa$-compact objects for a suitable regular $\kappa$. Its nearby rank-enlargement results retain explicit cardinal conditions; a bare change from one regular cardinal to a larger one is not a general theorem.

The modern theory of filtrations also distinguishes a directed presentation from a continuous chain. \citet[\S8, especially Lemma~8.4 and the following paragraph]{LiebermanRosickyVasey2020} discusses continuous filtrations and their agreement on a closed unbounded set of indices when the size has the appropriate regularity. This helps explain the example above: the successor-ordinal presentation of $\omega_1$ is not continuous at its countable limit stages. The comparison theorem for two smooth filtrations therefore cannot turn its stages into stages of the continuous limit-ordinal presentation. We checked that discussion and Theorem~8.8 in the authors' 5~June~2019 version; we do not identify its notion of filtration with Grothendieck's cardinal filtration of the whole category.

For homotopical algebra, the relevant application is to a full subcategory of the arrow category. \citet[Proposition~2.2 in author version~2; Proposition~2.3 in the journal]{RaptisRosicky2015} gives a controlled accessibility rank for weak equivalences in a combinatorial model category after the replacement functors preserve the specified presentable objects. Their Theorem~4.2 supplies better bounds under additional hypotheses, and Corollary~5.1 proves finite accessibility for weak equivalences of simplicial sets. These are genuine refinements of the small-generation problem motivating this chapter. They do not show that weak equivalences in every combinatorial model category are finitely accessible, nor do they establish the source's universal-presentation condition for an arbitrary filtration. The present status of the local questions is thus precise: $b_\pi\Rightarrow a_\pi$ is proved in the manuscript; the three-way equivalence as stated is false; the unfinished comparison has the explicit conditional completion proved above.
\end{dcomment}


\newpage
\printedpage{53}
% head: A. Grothendieck : Dérivateurs, Ch. XIII, Édition des Universités de Montpellier II et Paris VII
% foot: Version 25 janvier 2004

\section*{Index des notations}

\idx{LLP, RLP}{1}
\idx{$\Phi_*$, $\Psi^*$}{1}
\idx{$\widetilde{\Phi}$, $\overline{\Psi}$}{1}
\idx{$\Phi \longleftrightarrow \Psi$}{3}
\idx{$\mathrm{Cof}_W$, $\mathrm{Fib}_W$}{8}
\idx{$W^{\mathrm{univ}}$}{9}
\idx{$W_{\mathrm{univ}}$}{10}
\idx{$\underline{\mathrm{Fib}}_W S\subset\mathcal{M}/S$}{10}
\idx{$W_S$}{10}
\idx{$W_S^{\mathrm{u}}$}{10}
\idx{$S\backslash\underline{\mathrm{Cof}}_W$}{10}
\idx{$\Phi \barlongleftrightarrow \Psi$}{28}
\idx{$\mathcal{M}_{\mathrm{cf}}$, $W_{\mathrm{cf}}$}{29}
\idx{$\mathrm{Ho}_W$}{29}
\idx{$\mathrm{Hom}_W$}{30}
\idx{$\pi_{\mathrm{d}}$, $\pi$}{30}
\idx{$\mathcal{M}_{\mathrm{c}}$}{30}
\idx{$W_{\mathrm{c}}$}{31}
\idx{$\pi_{\mathrm{g}}$}{31}
\idx{$\mathcal{M}_{\mathrm{f}}$}{31}
\idx{$f\sim g$}{32}
\idx{$\Psi,\ \Phi,\ \mathcal{M}(\Psi),\ W(\Psi),\ \mathcal{M}(\Phi),\ W(\Phi)$}{33}
\idx{$\mathcal{M}(I)$, $W(I)$}{34}
\idx{$F(I),\ TF(I),\ C(I),\ TC(I)$}{35}
\idx{$\alpha_s$}{35}
\idx{$u^*$, $u_*$, $u_!$}{45}
\idx{$u^*_W$, $u^W_*$, $u^W_!$}{49}
\idx{$\mathrm{HO}_W(I)$}{49}
\idx{Diag}{50}
\idx{$\mathrm{HO}_{\mathcal{M},W}(I)$, $\mathbf{D}_{\mathcal{M},W}(I)$}{54}
\idx{$\mathrm{HO}^{\mathrm{lc}}_W(I)$, $u^*_{\mathrm{lc}}$}{56}
\idx{$u^{\mathrm{f}}_*$}{59}
\idx{$u^*_{\mathrm{c}}$}{67}
\idx{$\mathfrak{U}$}{77}
\idx{$\mathrm{Filt}^\pi(E)$}{77}

% ---- from Der-ChXIII-25-01-04_p056-060.tex
\printedpage{54}
% head: A. Grothendieck : Dérivateurs, Ch. XIII, Édition des Universités de Montpellier II et Paris VII
% foot: Version 25 janvier 2004

\entry{$L_\pi$}{77}
\entry{$\tilde{I}$}{79}
\entry{$\underline{M}$}{86}
\entry{$I_M$}{87}
\entry{$M^\pi$}{87}

\clearpage

\printedpage{55}
% head: A. Grothendieck : Dérivateurs, Ch. XIII, Édition des Universités de Montpellier II et Paris VII
% foot: Version 25 janvier 2004

\section*{Index terminologique}

\entry{accessible (partie close à gauche)}{75}
\entry{admissible (à gauche)}{59}
\entry{catégorie de modèles}{7}
\entry{catégorie de Quillen close}{25}
\entry{clos à gauche, à droite}{2}
\entry{couple de Quillen}{5}
\entry{faiblement saturée}{7}
\entry{filtration cardinale}{77}
\entry{filtration cardinale stricte}{86}
\entry{flèche coüniversellement dans $W$}{9}
\entry{flèche universellement dans $W$}{10}
\entry{grand devant $\pi$}{77}
\entry{objet chemin}{32}
\entry{objet cylindre}{32}
\entry{$\pi$-accessible (partie d'une catégorie munie d'une filtration cardinale stricte)}{86}
\entry{pré-Q-catégorie de modèles à gauche, à droite}{12}
\entry{pré-Quillen à gauche, à droite}{14}
\entry{pré-quintuplet de Quillen clos}{21}
\entry{propre}{8}
\entry{propriété de factorisation de Quillen}{5}
\entry{Q-catégorie de modèles}{8}
\entry{Q-catégorie de modèles à droite}{21}
\entry{Q-catégorie de modèles “close”}{16}
\entry{Q-homotopie à gauche, à droite}{30}
\entry{quillénisable}{50}
\entry{quillénisable sur Diag}{51}
\entry{quillénisable sur Diag uniformément et exactement à gauche}{57}
\entry{quillénisable uniformément à gauche sur Diag}{57}
\entry{quillénisation uniforme à droite sur Diag}{58}
\entry{quillénisation uniforme à gauche sur Diag}{57}
\entry{quillénisation uniforme et exacte à gauche sur Diag}{57}
\entry{représentation bien adaptée}{84}
\entry{représentation parfaitement adaptée}{86}
\entry{strictement admissible}{61}
\entry{structure de Quillen}{28}

\clearpage

\printedpage{56}
% head: A. Grothendieck : Dérivateurs, Ch. XIII, Édition des Universités de Montpellier II et Paris VII
% foot: Version 25 janvier 2004

\entry{structure de Quillen exacte}{27}
\entry{structure de Quillen exacte à gauche, à droite}{27}
\entry{structure de Quillen faible}{27}
\entry{structure de Quillen faible biexacte}{28}
\entry{structure de Quillen faible close, close à gauche, à droite}{28}
\entry{structure de Quillen parfaite}{28}
\entry{système de Quillen clos}{26}
\entry{système de Quillen faible}{26}
\entry{système de Quillen faible exact à gauche, à droite}{26}
\entry{système inductif bien adapté}{84}
\entry{système inductif parfaitement adapté}{86}
\entry{triple clos de Quillen}{25}
\entry{triple de Quillen}{14}
\entry{$W$-cofibration}{8}
\entry{$W$-équivalence coüniverselle}{9}
\entry{$W$-fibration}{8}

\clearpage

\printedpage{57}
% head: A. Grothendieck : Dérivateurs, Ch. XIII, Édition des Universités de Montpellier II et Paris VII
% foot: Version 25 janvier 2004

\section*{Liste des axiomes}

\noindent $\mathcal{M}$ catégorie

\noindent $W \ \subset\ \Fl(\mathcal{M})$ , \qquad $C$, $F$, $TC$, $TF \ \subset\ \Fl(\mathcal{M})$

\begin{description}[leftmargin=4.5em,labelwidth=4em]
\item[M 0] La catégorie $\mathcal{M}$ admet des limites projectives finies et des limites inductives finies.

\item[M 1] Les flèches appartenant à $C$ ont la propriété de relèvement à gauche relativement à celles appartenant à $TF$, et les flèches appartenant à $F$ ont la propriété de relèvement à droite relativement à celles appartenant à $TC$.

\item[M 2] Toute flèche de $\mathcal{M}$ se décompose en une flèche de $C$ suivie d'une flèche de $TF$, et en une flèche de $TC$ suivie d'une flèche de $F$.

\item[M 3] La partie $F$ de $\Fl(\mathcal{M})$ est stable par composition et changement de base, $C$ est stable par composition et cochangement de base, $F$ et $C$ contiennent les isomorphismes.

\item[M 4] La partie $TF$ de $\Fl(\mathcal{M})$ est stable par composition et changement de base, $TC$ est stable par composition et cochangement de base, $TF$ et $TC$ contiennent les isomorphismes.

\item[M 5] Si dans un triangle commutatif de $\mathcal{M}$ deux des trois flèches sont dans $W$, il en est de même de la troisième ; $W$ contient les isomorphismes.

\item[M$'$] $TC \ \subset\ C \cap W$, $TF \ \subset\ F \cap W$.

\item[M$''$] Les flèches appartenant à $C$ sont des $W$-cofibrations, et les flèches appartenant à $F$ sont des $W$-fibrations.

\item[M$'_{\mathbf{0}}$(g)] $TC = C \cap W$.

\item[M$'_{\mathbf{0}}$(d)] $TF = F \cap W$.

\item[M 2\,bis] $W = TF \circ TC$.
\end{description}

\clearpage

\printedpage{58}
% head: A. Grothendieck : Dérivateurs, Ch. XIII, Édition des Universités de Montpellier II et Paris VII
% foot: Version 25 janvier 2004

\setcounter{footnote}{0}
\renewcommand{\thefootnote}{\fnsymbol{footnote}}
\section*{Glossaire}

\subsection*{Section 1}

\begin{itemize}
\item $\Phi \ \subset\ \Fl(\mathcal{M})$

$\Phi$ est \textit{clos à gauche} : il existe \ $\Psi \ \subset\ \Fl(\mathcal{M})$ tel que $\Phi$ soit formé des flèches ayant la propriété de relèvement à gauche relativement aux flèches de $\Psi$.

\item $\Psi \ \subset\ \Fl(\mathcal{M})$

$\Psi$ est \textit{clos à droite} : il existe \ $\Phi \ \subset\ \Fl(\mathcal{M})$ tel que $\Psi$ soit formé des flèches ayant la propriété de relèvement à droite relativement aux flèches de $\Phi$.

\item $\Phi$, $\Psi \ \subset\ \Fl(\mathcal{M})$

$(\Phi, \Psi)$ a la \textit{propriété de factorisation de Quillen} :
\begin{itemize}
\item[a)] les flèches de $\Phi$ ont la propriété de relèvement à gauche relativement aux flèches de $\Psi$ ;
\item[b)] toute flèche $f$ de $\mathcal{M}$ se factorise en
\[
f = pi \quad , \qquad i \in \Phi \ , \quad p \in \Psi \ .
\]
\end{itemize}

\item $(\Phi, \Psi)$ \textit{couple de Quillen} :
\begin{itemize}
\item[a)] $(\Phi, \Psi)$ a la propriété de factorisation de Quillen ;
\item[b)] $\Phi$ est clos à gauche et $\Psi$ clos à droite.
\end{itemize}
\end{itemize}

\subsection*{Section 2}

\begin{itemize}
\item $W \ \subset\ \Fl(\mathcal{M})$

\textit{$W$ faiblement saturé} :
\begin{itemize}
\item[a)] $W$ contient les isomorphismes ;
\item[b)] si deux parmi les trois flèches d'un triangle commutatif sont dans $W$, la troisième aussi.
\end{itemize}

\item $W \ \subset\ \Fl(\mathcal{M})$

Un morphisme $p : X \longrightarrow Y$ de $\mathcal{M}$ est une \textit{$W$-fibration} ($p \in \mathrm{Fib}_W$) si pour tout diagramme de carrés cartésiens
\end{itemize}

% ---- from Der-ChXIII-25-01-04_p061-065.tex
\printedpage{59}
% head: A. Grothendieck : Dérivateurs, Ch. XIII, Édition des Universités de Montpellier II et Paris VII

\[
\begin{tikzcd}
X \arrow[d, "p"'] & X' \arrow[l] \arrow[d] & X'' \arrow[l, "t"'] \arrow[d] \\
Y & Y' \arrow[l] & Y'' \arrow[l, "s"]
\end{tikzcd}
\quad ,
\]

si $s \in W$, alors $t \in W$.

\begin{itemize}
\item[\textbullet] $W \ \subset\ \Fl(\Mo)$

Un morphisme $i : A \longrightarrow B$ de $\Mo$ est une $W$-\textit{cofibration} ($i \in \mathrm{Cof}_W$) si pour tout diagramme de carrés cocartésiens
\[
\begin{tikzcd}
A'' \arrow[d] & A' \arrow[l, "s"'] \arrow[d] & A \arrow[l] \arrow[d, "i"] \\
B'' & B' \arrow[l, "t"] & B \arrow[l]
\end{tikzcd}
\quad ,
\]
si $s \in W$, alors $t \in W$.

\item[\textbullet] $W \ \subset\ \Fl(\Mo)$

Un morphisme $p : X \longrightarrow Y$ de $\Mo$ est \textit{universellement dans} $W$, ou une $W$-\textit{équivalence universelle} ($p \in W_{\mathrm{univ}}$) si pour tout carré cartésien
\[
\begin{tikzcd}
X \arrow[d, "p"'] & X' \arrow[l] \arrow[d, "p'"] \\
Y & Y' \arrow[l]
\end{tikzcd}
\quad ,
\]
la flèche $p'$ appartient à $W$.

\item[\textbullet] $W \ \subset\ \Fl(\Mo)$

Un morphisme $i : A \longrightarrow B$ de $\Mo$ est \textit{coüniversellement dans} $W$, ou une $W$-\textit{équivalence coüniverselle} ($i \in W^{\mathrm{univ}}$) si pour tout carré cocartésien
\[
\begin{tikzcd}
A' \arrow[d, "i'"'] & A \arrow[l] \arrow[d, "i"] \\
B' & B \arrow[l]
\end{tikzcd}
\quad ,
\]
la flèche $i'$ appartient à $W$.
\end{itemize}
% foot: Version 25 janvier 2004

\printedpage{60}
% head: A. Grothendieck : Dérivateurs, Ch. XIII, Édition des Universités de Montpellier II et Paris VII

\begin{itemize}
\item[\textbullet] \textit{Pré-Q-catégorie de modèles à gauche} (resp. \textit{à droite}) :

$\Mo$ et $W$, $C$, $F$, $TC$, $TF$ $\subset$ $\Fl(\Mo)$ satisfaisant aux axiomes \textbf{M~0} à \textbf{M~5}, $\mathbf{M}'$, $\mathbf{M}''$ et \Mzg{} (resp. \Mzd).

\item[\textbullet] \textit{Q-catégorie de modèles}, \textit{triple de Quillen} :

catégorie de modèles au sens de Quillen.

\item[\textbullet] \textit{Q-catégorie de modèles “close”}, \textit{triple clos de Quillen}, \textit{catégorie de Quillen close}, \textit{système de Quillen clos} :

catégorie de modèles fermée au sens de Quillen.

\item[\textbullet] \textit{Pré-Quillen à gauche} (resp. \textit{à droite}) :

pré-Q-catégorie de modèles à gauche (resp. à droite).

\item[\textbullet] \textit{Pré-quintuplet de Quillen clos} :

$W$, $C$, $F$, $TC$, $TF$ $\subset$ $\Fl(\Mo)$ satisfaisant aux axiomes \textbf{M~0} à \textbf{M~5}, $\mathbf{M}'$, $\mathbf{M}''$, $C$, $TC$ étant clos à gauche et $F$, $TF$ clos à droite.

\item[\textbullet] \textit{Système de Quillen faible} :

$W$, $C$, $F$, $TC$, $TF$ $\subset$ $\Fl(\Mo)$ satisfaisant aux axiomes \textbf{M~0} à \textbf{M~5}, $\mathbf{M}'$, $\mathbf{M}''$.

\item[\textbullet] \textit{Système de Quillen faible exact à gauche} (resp. \textit{à droite}) :

$W$, $C$, $F$, $TC$, $TF$ $\subset$ $\Fl(\Mo)$ satisfaisant aux axiomes \textbf{M~0} à \textbf{M~5}, $\mathbf{M}'$, $\mathbf{M}''$ et \Mzg{} (resp. \Mzd).
\end{itemize}

\section*{Section 3}

\begin{itemize}
\item[\textbullet] $(\Mo,W)$ \textit{catégorie de modèles} :

$\Mo$ catégorie satisfaisant à l'axiome \textbf{M~0} ;

$W$ $\subset$ $\Fl(\Mo)$ faiblement saturé.

\item[\textbullet] \textit{Structure de Quillen faible associée à la catégorie de modèles} $(\Mo,W)$ :

quadruplet $(C,F,TC,TF)$ de parties de $\Fl(\Mo)$ satisfaisant aux axiomes \textbf{M~1} à \textbf{M~5}, $\mathbf{M}'$, $\mathbf{M}''$.

\item[\textbullet] \textit{Structure de Quillen faible exacte} :

structure de Quillen faible satisfaisant en outre à l'axiome \textbf{M~2\,bis}.

\item[\textbullet] \textit{Structure de Quillen faible exacte à gauche} (resp. \textit{à droite}) :

structure de Quillen faible satisfaisant en outre à l'axiome \Mzg{} (resp. \Mzd).
\end{itemize}
% foot: Version 25 janvier 2004

\printedpage{61}
% head: A. Grothendieck : Dérivateurs, Ch. XIII, Édition des Universités de Montpellier II et Paris VII

\begin{itemize}
\item[\textbullet] \textit{Structure de Quillen faible close à gauche} (resp. \textit{à droite}) :

structure de Quillen faible $(C,F,TC,TF)$ telle que $C$, $TC$ (resp. $F$, $TF$) soient clos à gauche (resp. à droite).

\item[\textbullet] \textit{Structure de Quillen faible close} :

structure de Quillen faible à la fois close à gauche et à droite.

\item[\textbullet] \textit{Structure de Quillen} :

structure de Quillen faible “biexacte” :

structure de Quillen faible à la fois exacte à gauche et à droite

$=$ catégorie de modèles de Quillen propre.

\item[\textbullet] \textit{Structure de Quillen parfaite} :

structure de Quillen faible “biexacte” et close

$=$ catégorie de modèles fermée de Quillen propre.
\end{itemize}

\section*{Section 5}

\begin{itemize}
\item[\textbullet] Catégorie de modèles $(\Mo,W)$ \textit{quillénisable} :

il existe une structure de Quillen faible \textit{exacte} $(C,F,TC,TF)$ associée à $(\Mo,W)$.

\item[\textbullet] $(\Mo,W)$ \textit{quillénisable sur} Diag $\subset$ Cat :

pour tout $I$ dans Diag, la catégorie des modèles $(\Mo(I),\, W(I))$ est quillénisable.
\end{itemize}

\section*{Section 6}

\begin{itemize}
\item[\textbullet] Catégorie de modèles $(\Mo,W)$ \textit{quillénisable uniformément à gauche sur} Diag :

il existe une structure de Quillen faible \textit{close} $(C,F,TC,TF)$ pour $(\Mo,\, W)$ telle que pour tout $I$ dans Diag, si l'on définit $C(I)$ et $TC(I)$ “argument par argument” (\footnote{Un morphisme de foncteurs $u : X \longrightarrow Y$ de $I^{\circ}$ vers $\Mo$ étant dans $C(I)$ (resp. $TC(I)$) si et seulement si pour tout objet $i$ de $I$, $u_i$ est dans $C$ (resp. $TC$).}), et $F(I)$ et $TF(I)$ par la propriété de relèvement à droite relativement à $TC(I)$ et $C(I)$ respectivement (\footnote{Un morphisme de $\Mo(I)$ étant dans $F(I)$ (resp. $TF(I)$) si et seulement si il a la propriété de relèvement à droite relativement aux flèches appartenant à $TC(I)$ (resp. $C(I)$).}), le quadruplet $(C(I),TC(I),F(I),TF(I))$ soit une structure de Quillen faible \textit{close exacte} pour $(\Mo(I),\, W(I))$. On dit alors que $(C,F,TC,TF)$ est une \textit{quillénisation uniforme à gauche} de $(\Mo,\, W)$ sur Diag).
\end{itemize}
% foot: Version 25 janvier 2004

\printedpage{62}
% head: A. Grothendieck : Dérivateurs, Ch. XIII, Édition des Universités de Montpellier II et Paris VII

\begin{itemize}
\item[\textbullet] Catégorie de modèles $(\Mo,W)$ \textit{quillénisable sur} Diag \textit{uniformément et exactement à gauche} :

$(\Mo,W)$ admet une quillénisation uniforme à gauche sur Diag $(C,F,TC,TF)$ telle que $TC = C \cap W$ (autrement dit, telle que $(C,F,TC,TF)$ soit une strucure de Quillen faible \textit{exacte à gauche}). On dit alors que $(C,F,TC,TF)$ est une \textit{quillénisation uniforme et exacte à gauche} de $(\Mo,\, W)$ sur Diag. % [sic] strucure -> structure, letter t omitted
\end{itemize}

\section*{Section 7}

\begin{itemize}
\item[\textbullet] $I$ ensemble ordonné filtrant, \quad $\pi$ cardinal

$I$ \textit{grand devant} $\pi$ :

toute partie de $I$ de cardinal inférieur ou égal à $\pi$ admet un majorant dans $I$.

\item[\textbullet] $E$ catégorie

Une \textit{filtration cardinale} de $E$ est une filtration croissante
\[
\Big(\Filt^{\pi}(E)\Big)_{\pi\geq\pi_0} \qquad , \qquad\qquad \pi_0 \text{ un cardinal infini fixé,}
\]
de $E$ par des sous-catégories strictement pleines $\Filt^{\pi}(E)$, indexée par les cardinaux $\pi \geq \pi_0$ appartenant à un univers donné, et satisfaisant aux conditions suivantes :
\begin{itemize}
\item[\textit{a})] pour tout $\pi \geq \pi_0$, $\Filt^{\pi}(E)$ est équivalente à une petite catégorie ;
\item[\textit{b})] la catégorie $E$ admet des limites inductives filtrantes indexées par des ensembles ordonnés $I$ grands devant $\pi_0$, et pour tout $\pi \geq \pi_0$, $\Filt^{\pi}(E)$ est stable par les dites limites si $\mathrm{card}\, I \leq \pi$ ;
\item[\textit{c})] Pour tout $\pi \geq \pi_0$, et pour tout $X \in \mathrm{Ob}\, E$, il existe un isomorphisme
\[
X \ \simeq\ \varinjlim_{i\in I} X_i \ ,
\]
où $(X_i)_{i\in I}$ est un système inductif filtrant dans $E$, indexé par un ensemble ordonné $I$ grand devant $\pi$, et où les $X_i$ sont des objets de $\Filt^{\pi}(E)$. De plus, si $X \in \mathrm{Ob}\,\Filt^{\pi'}(E)$, $\pi' \geq \pi$, on peut prendre $I$ tel que $\mathrm{card}\, I \leq \pi'^{\pi}$.
\end{itemize}

\item[\textbullet] $E$ catégorie, \quad $(\Filt^{\pi}(E))_{\pi\geq\pi_0}$ filtration cardinale de $E$, \quad $\pi \geq \pi_0$ cardinal

Un \textit{système inductif bien adapté} (à la filtration cardinale donnée, pour les cardinaux $\pi_0$, $\pi$ donnés) est un système inductif filtrant $(X_i)_{i\in I}$ satisfaisant aux conditions suivantes :
\begin{itemize}
\item[\textit{a})] $I$ est grand devant $\pi$, et pour tout $i \in I$, $X_i \in \Filt^{\pi}(E)$ ;
\item[\textit{b})] Toute partie $J$ de $I$ de cardinal $\leq \pi$, et grande devant $\pi_0$, admet une borne supérieure $i_J$, et on a
\[
X_{i_J} \ \xleftarrow{\ \sim\ }\ \varinjlim_{i\in J} X_i \ .
\]
% foot: Version 25 janvier 2004
\end{itemize}

\printedpage{63}
% head: A. Grothendieck : Dérivateurs, Ch. XIII, Édition des Universités de Montpellier II et Paris VII

Si $X$ est un objet de $E$, une \textit{représentation bien adaptée} de $X$, en termes d'objets de $\Filt^{\pi}(E)$, est la donnée d'un système inductif bien adapté $(X_i)_{i\in I}$, à la filtration cardinale donnée, pour les cardinaux $\pi_0$, $\pi$, et d'un isomorphisme
\[
X \ \xleftarrow{\ \sim\ }\ \varinjlim_{i\in I} X_i \ .
\]

\item[\textbullet] $E$ catégorie admettant des petites limites inductives filtrantes

Une \textit{filtration cardinale stricte} de $E$ est une filtration cardinale $(\Filt^{\pi}(E))_{\pi\geq\pi_0}$ telle que pour tout cardinal $\pi \geq \pi_0$, $\Filt^{\pi}(E)$ soit stable par limites inductives filtrantes.

\item[\textbullet] $E$ catégorie admettant des petites limites inductives filtrantes

$(\Filt^{\pi}(E))_{\pi\geq\pi_0}$ filtration cardinale stricte de $E$, \quad $\pi \geq \pi_0$ cardinal

Un \textit{système inductif parfaitement adapté} (à la filtration cardinale donnée, pour le cardinal $\pi$) est un système inductif filtrant $(X_i)_{i\in I}$ satisfaisant aux conditions suivantes :
\begin{itemize}
\item[\textit{a})] $I$ est grand devant $\pi$, et pour tout $i \in I$, $X_i \in \Filt^{\pi}(E)$ ;
\item[\textit{b})] Toute partie $J$ de $I$ filtrante de cardinal $\leq \pi$, admet une borne supérieure $i_J$, et on a
\[
X_{i_J} \ \xleftarrow{\ \sim\ }\ \varinjlim_{i\in J} X_i \ .
\]
\end{itemize}
Si $X$ est un objet de $E$, une \textit{représentation parfaitement adaptée} de $X$, en termes d'objets de $\Filt^{\pi}(E)$, est la donnée d'un système inductif bien adapté $(X_i)_{i\in I}$, à la filtration cardinale donnée, pour le cardinal $\pi$, et d'un isomorphisme
\[
X \ \xleftarrow{\ \sim\ }\ \varinjlim_{i\in I} X_i \ .
\]
\end{itemize}
% foot: Version 25 janvier 2004

% ---- from Der-ChXIII-25-01-04_p066-066.tex
\printedpage{64}
% head: A. Grothendieck : Dérivateurs, Ch. XIII, Édition des Universités de Montpellier II et Paris VII
% foot: Version 25 janvier 2004

\section*{Références}

\noindent H. J. Baues, ``Algebraic Homotopy'', Cambridge studies in advanced mathematics, vol. 15 (1989).

\noindent R. Hartshorne, ``Residues and Duality'', Lecture Notes in Mathematics, vol. 20, Springer-Verlag (1966).

\noindent D. G. Quillen, ``Homotopical Algebra'', Lecture Notes in Mathematics, vol. 43, Springer-Verlag (1967).

\noindent R. W. Thomason, \textit{Cat as a Closed Model Category}, Cahiers de Topologie et Géométrie Différentielle, Vol. XXI-3, pp. 305-324 (1980).

\noindent J.-L. Verdier, ``Des catégories dérivées des catégories abéliennes'', Astérisque, vol. 239 (1996).

\noindent SGA 4, I, ``Théorie des Topos et Cohomologie Étale des Schémas'', Séminaire de Géométrie Algébrique du Bois-Marie, dirigé par M. Artin, A. Grothendieck, J.-L. Verdier, Lecture Notes in Mathematics, vol. 269, Springer-Verlag (1972).
\begin{otherlanguage}{english}
\dgroverview{XIII-b}{Chapter XIII \textperiodcentered\ \textup{XIII-b} \textperiodcentered\ Overview}
\selectlanguage{english}

The second part of Chapter XIII joins two constructions that Grothendieck needs for the same purpose. He wants a homotopy theory to provide a coherent calculus of diagrams and their direct images. He also wants the replacements used in that calculus to be obtainable from a controlled set of small data. Sections~5 and~6 address the first problem through quillenisation; section~7 returns to the cardinal filtrations of SGA~4 to address the second. The manuscript ends during the statement of its final comparison lemma. The break is part of the surviving source, and the commentary does not supply a purported continuation of Grothendieck's text \citep{GrothendieckDerivateursXIII}.

Throughout these pages, $\mathcal M(I)$ means the category of contravariant diagrams $I^{\mathrm{op}}\to\mathcal M$, and $W(I)$ consists of objectwise weak equivalences. This variance explains both the right Kan extension formula over $I'/s$ and the fact that a greatest element of a poset computes the ordinary limit by evaluation. Grothendieck's phrase ``cat\'egorie de mod\`eles'' initially specifies much less than a modern Quillen model category: it is a category with finite limits and colimits and a two-out-of-three class of morphisms. Cofibrations, fibrations, closedness, and exactness are additional data or properties. Keeping those qualifications visible is necessary both to understand the proofs and to avoid attributing modern theorems to a weaker hypothesis.

A closed exact structure has more force than the subsequent doubts suggest. Exactness says that every weak equivalence factors as a trivial cofibration followed by a trivial fibration. Closedness supplies closure under retracts for the lifting classes. The elementary retract argument therefore gives $TC=C\cap W$ and $TF=F\cap W$. In a uniformly left quillenised diagram theory, restriction preserves the objectwise cofibrations and trivial cofibrations, so its ordinary right adjoint preserves trivial fibrations. The graph-factorisation proof of Ken Brown's lemma then shows that this right adjoint preserves weak equivalences between fibrant diagrams. Consequently every $u$ for which the ordinary $u_*$ exists is admissible under those uniform closed exact hypotheses. This answers the product, tower, and group-indexing questions in the scope where the required ordinary operations exist. The argument uses only finite products, the specified factorizations, lifting, and two-out-of-three; it does not infer the existence of missing infinite limits \citep[Lemmas~1.1.9 and~1.1.12]{Hovey1999}.

Grothendieck nevertheless gives a useful geometric proof for finite posets. A downward-closed cover reduces a limit to a pullback, and the fibrancy of the cospan makes that pullback homotopy invariant. The construction exposes the local geometry behind pointwise derived direct images. His subsequent noetherian induction is correct, but its proposed larger class has an unexpected limitation: a sober noetherian Alexandrov space associated to a poset is finite. The commentary proves the assertion directly. Pointwise computation for a general comma category also deserves its own hypothesis: restriction to that comma category must preserve the chosen fibrant diagrams. This holds for the finite-poset situation and in the usual cocomplete diagram-model setting, where the left adjoint to restriction along a discrete fibration is computed by fibrewise coproducts. Admissibility and the pointwise comparison are thus separately justified.

The derived-category examples illuminate a different boundary. For a finite shape, a diagram of bounded-below complexes has a common lower bound. For an infinite discrete shape, its components may be bounded below at degrees that tend to minus infinity. Accordingly, $\operatorname{Ch}^+(\mathcal A(I))$ and $\operatorname{Fun}(I^{\mathrm{op}},\operatorname{Ch}^+(\mathcal A))$ need not agree, and products need not stay in $D^+(\mathcal A)$. An explicit family of shifted copies of a field $k$ makes the obstruction visible. Unbounded resolutions answer the substantive existence question in the standard module and sheaf settings: $K$-injective resolutions allow right derived functors to be defined on the full unbounded derived category. Spaltenstein's work predates the present manuscript, and the general Grothendieck-category theorem is recorded precisely in the references. Beke's construction of the combinatorial injective model structure on unbounded complexes, combined with the diagram-model theorem, gives the all-small-category derivator for any Grothendieck abelian category \citep[Proposition~3.13]{Beke2000}. As a direct later algebraic treatment under the additional assumption of enough projectives, Alonso Tarr\'io, \'Alvarez D\'iaz, and Jerem\'ias L\'opez prove the all-small-category strong derivator theorem for a Grothendieck category with enough projectives \citep[Theorem~3.4, arXiv version~3]{AlonsoAlvarezJeremias2023}.

The modern comparison also separates pointedness from stability. Fibre and cofiber sequences are available in pointed homotopy theories without making their homotopy categories additive or their suspension functors invertible. The manuscript's deliberately broad ``triangul\'ee'' language must be read accordingly. Strong stable derivators give actual canonical triangulations, with exact diagram operations, as proved by \citet[Theorem~4.16 and Proposition~4.18]{Groth2013}. His Theorem~1.28 and Proposition~1.30 provide the combinatorial model-category construction of the full derivator. These results explain the successful modern form of the programme without replacing the original finite-limit framework by an unstated cocompleteness assumption.

Section~7 changes the technical language but keeps the construction problem in view. A cardinal filtration $(\Filt^\pi(E))$ consists of essentially small stages, bounded colimit-closure conditions, and presentations of every object by sufficiently directed diagrams in each stage. ``Grand devant $\pi$'' means that every subset of cardinality at most $\pi$ has an upper bound, hence means $\pi^+$-directed in the customary regular-cardinal convention. Membership in $\Filt^\pi(E)$ is a size condition; preservation of these colimits by $\operatorname{Hom}(A,-)$ is a presentability condition. The filtration axioms do not identify the two. SGA~4's comparison results explain why suitable cardinals permit one to use both, and why the extra hypothesis becomes visible in the discussion of canonical slice categories \citep[I, 9.12 and 9.16--9.18]{SGA4I1972}.

The adapted-presentation construction itself survives the audit well. Replace an indexing poset by its small sufficiently directed subposets and take the colimit over each piece. The new index admits the relevant small suprema, and the new diagram preserves them. The editors correctly observe that the first enlargement lemma is immediate under its stated hypotheses: one may simply adjoin an upper bound to the given small subset. The transfinite proof does something more useful when only the weaker directedness of the ambient poset is assumed. Its cardinal equation controls the number of newly imposed bounds. The commentary proves both versions and checks their different assumptions. The separate assertion that generalized continuum hypothesis always makes $(\pi')^\pi=\pi'$ is false at singular $\pi'$ with small cofinality; the displayed power-of-two case remains valid without that hypothesis.

The proposed equivalence of the three conditions on a part $M$ needs a more substantial correction. Condition $a_\pi$ describes $M$ as the objects obtainable by $\pi^+$-directed colimits from $M\cap\Filt^\pi(E)$. Condition $b_\pi$ asks, for every perfectly adapted presentation, that its colimit belong to $M$ if and only if its $M$-stages are cofinal. Condition $c_\pi$ combines strict generation, closure under sufficiently directed colimits, and ordinary filteredness of a canonical slice. These quantify over different data. The commentary gives two fully checked counterexamples in $E=\mathcal P(\omega_1)$ with $\pi=\aleph_0$. Taking $M$ to be the countable successor initial segments together with $\omega_1$ satisfies $a_\pi$, but a continuous diagram of countable limit initial segments violates $b_\pi$. Taking $M$ to be the finite subsets together with $\omega_1$ satisfies $c_\pi$ but violates $a_\pi$. Both examples use the ordinary cardinality filtration, whose stages have the required presentability. They therefore cannot be removed by adding only the compactness hypothesis mentioned in the manuscript.

What does remain is substantial. Grothendieck's proof of $b_\pi\Rightarrow a_\pi$ is completed without any presentability assumption, using a cofinal family of small subposets with greatest elements. With presentable small stages, the canonical slice construction gives the stronger directedness needed to prove colimit closure. Near page~94 there is a distinct size error: bounding the image of an indexing category does not bound that category itself. A filtered diagram of finite sets can have an uncountable colimit even though only countably many finite sets and maps occur as values. The repair bounds the indexing morphisms. Finally, the enlarged-cardinal idea is developed as a separate conditional result, with every object of $\Filt^\pi(E)$ assumed $\pi^+$-presentable and every object of $\Filt^{\pi^+}(E)$ assumed $\pi^{++}$-presentable, and with a complete alternating-chain proof. It is not presented as the missing end of Lemma~3. The powerset counterexamples and this conditional completion are marked as the commentators' work, with no claim of priority.

The accompanying registers preserve the distinction between a transcription mistake, a slip already present in the edition, and a mathematical claim requiring a proof or a correction. Every doubtful reading and every \texttt{[sic]} marker of the transcription is accounted for. The editorial indices, axiom list, and glossary have also been checked, including the missing cardinal bound in the glossary's definition of strict filtration. References report the portions actually verified. The purpose throughout is to let the reader follow the mathematical argument while retaining an exact account of what this fragment establishes, what later theory supplies, and what the present commentary proves separately.
\end{otherlanguage}
\dgrchunkend

% DER_XIV: Grothendieck's text with the commentary set in, then the overview, inside a group that applies the chunk's own definitions
\dgrchapter{XIV}{Chapter XIV}{Carrés h-cartésiens et h-cocartésiens}
\dgrpart{0}{XIV}{Manuscript pages 1--8}{XIV\enspace\textperiodcentered\enspace Chapter XIV\enspace\textperiodcentered\enspace manuscript pages 1--8}{Chapter XIV \textperiodcentered\ \textup{XIV}}
\dgrchunkbegin
% DER_XIV: the definitions of the chunk's preamble, applied inside its group
\let\printedpage\relax\newcommand{\printedpage}[1]{\if\relax\detokenize{#1}\relax\else\begingroup\reversemarginpar\marginnote{\footnotesize p.~#1}\endgroup\fi}
\let\msp\relax\newcommand{\msp}[1]{\marginnote{\footnotesize [#1]}}

% footnote marks as printed: "(1)" in the text, "1" before the footnote text
\makeatletter
\providecommand\@makefnmark{}\renewcommand\@makefnmark{\hbox{\textsuperscript{\normalfont(\@thefnmark)}}}
\providecommand\@makefntext{}\renewcommand\@makefntext[1]{\parindent 1em\noindent\hb@xt@1.8em{\hss\textsuperscript{\normalfont\@thefnmark}}#1}
\makeatother
\let\uHom\relax\newcommand{\uHom}{\underline{\mathrm{Hom}}}
\let\bDelta\relax\newcommand{\bDelta}{\mathbf{\Delta}}
\let\ho\relax\newcommand{\ho}{\mathrm{ho}}
\let\cM\relax\newcommand{\cM}{\mathcal{M}}

% ---- from Der-ChXIV-07-10-10_p001-005.tex
% ===================== PDF page 1 =====================
\printedpage{}% no page number printed on this page
% head: (none)

\begin{center}
\vspace*{6cm}

{\LARGE \textsf{LES DÉRIVATEURS}}

\vspace{1cm}

{\large \textbf{ALEXANDRE GROTHENDIECK}}

\vspace{2cm}

{\large \textbf{Chapitre XIV}}

\vspace{0.3cm}

{\large \textbf{Carrés h-cartésiens et h-cocartésiens}}
\end{center}

\clearpage

% ===================== PDF page 2 =====================
\printedpage{}% no page number printed on this page
% head: (none)

\vspace*{6cm}

\noindent Ce texte a été déchiffré et transcrit en \LaTeXe{} par M. Künzer. Il a été édité par M. Künzer, J. Malgoire, et G. Maltsiniotis. La transcription est aussi fidèle que possible au manuscrit. Pour les quelques corrections évidentes, ou rares commentaires des éditeurs, ainsi que pour la numérotation originale des pages du manuscrit, les caractères de machine à écrire \texttt{[typewriter]} entre crochets sont utilisés. Un point d'interrogation entre crochets signifie que l'on n'est pas sûr du mot qui précède.

\noindent Ce chapitre est inachevé. Dans la première section, Grothendieck introduit une notion de “catégorie quadrangulée” à gauche (resp. à droite). Le texte s'interrompt brutalement au début de la deuxième section, consacrée aux exemples, quand Grothendieck réalise, sans doute, que l'exemple principal qu'il a en vue (la catégorie homotopique d'une catégorie de modèles munie des carrés homotopiquement cartésiens (resp. cocartésiens)) ne satisfait \textit{pas} à l'axiomatique des catégories quadrangulées. Cela est brièvement expliqué dans une postface des éditeurs.

\noindent Cette édition est provisoire. Les remarques, commentaires et corrections sont bienvenus. Envoyer un message à :

\noindent \textbf{maltsin@math.jussieu.fr}

\hfill G. Maltsiniotis\hspace*{2cm}

\clearpage

% ===================== PDF page 3 =====================
\printedpage{1}
% head: A. Grothendieck : Dérivateurs, Ch. XIV, Édition des Universités de Montpellier II et Paris VII
% foot: Version 7 octobre 2010

\msp{page 1}

\section*{1\quad Axiomes de structures quadrangulées (une tentative avortée)}

Soit $(\mathcal{M},W)$ une catégorie de modèles admettant une théorie de la fibre homotopique. Alors on a dans $\mathrm{Ho}_W=\mathcal{M}W^{-1}$ une notion de \textit{carré $W$-cartésien} (ou $W$-homotopiquement cartésien), par quoi on entend un diagramme carré
\[
\begin{tikzcd}
X_1 \arrow[d] & X_3 \arrow[l] \arrow[d] \\
X_0 & X_2 \arrow[l]
\end{tikzcd}
\]
dans $\mathrm{Ho}_W$ provenant, à isomorphisme près, d'un carré $W$-cartésien dans $\mathcal{M}$ – donc isomorphe, dans $\underline{\mathrm{Hom}}(\boldsymbol{\Delta}_1\times\boldsymbol{\Delta}_1,\mathrm{Ho}_W)$, a un carré \textit{cartésien} \texttt{[dans $\mathcal{M}$]} % [sic] a un carré -> à un carré, preposition printed without its grave accent
\[
\begin{tikzcd}
Y_1 \arrow[d, "p_1"'] & Y_3\ (\simeq Y_1\times_{Y_0}Y_2) \arrow[l] \arrow[d] & \\
Y_0 & Y_2 \arrow[l, "p_2"] & ,
\end{tikzcd}
\]
avec $p_1,\ p_2\in\mathrm{Fib}_W$. (En fait, il suffit de $p_1$ ou $p_2$ dans $\mathrm{Fib}_W$ pour assurer le caractère $W$-cartésien du carré cartésien envisagé.) Je voudrais axiomatiser cette situation, histoire surtout de faire la liste des propriétés serviables essentielles de cette notion, et de la notion duale de carré $W$-cocartésien (ou $W$-homotopiquement cocartésien), dans le cas où $(\mathcal{M},W)$ admet une théorie de la cofibre homotopique. L'espèce de structure, savoir celle des “catégories quadrangulées” (à gauche, à droite, bilatères) est une variante, dans un cas \msp{page 2} non additif, de la notion de catégorie triangulée dégagée par \textsc{Verdier}. Comme celle-ci, cette notion est en fait trop faible – ce n'est \textit{pas} “la bonne” – mais elle a un intérêt pratique provisoire. La “bonne” notion (s'il en est une à part celle de dérivateur) devrait travailler avec une notion plus fine de carré et de carré $W$-cartésien, en travaillant systématiquement dans $\mathrm{HO}_W(\boldsymbol{\Delta}_1\times\boldsymbol{\Delta}_1)=\mathcal{M}(\boldsymbol{\Delta}_1\times\boldsymbol{\Delta}_1)W(\boldsymbol{\Delta}_1\times\boldsymbol{\Delta}_1)^{-1}$ (où vivent les “vrais” carrés $W$-cartésiens !), et catégories similaires, telles que $\mathrm{HO}_W(\Psi)$, $\mathrm{HO}_W(\Phi)$. Mais les inclure dans une axiomatique ne faisant pas mention d'une catégorie de modèles $(\mathcal{M},W)$, revient manifestement à mettre le doigt dans l'engrenage de la notion de dérivateur, et ce n'est pas mon propos pour l'instant.

\dgrblock{dgrb741}{What a square must remember}
\begin{dcomment}{exposition}{p1}{What a square must remember}
Grothendieck wants a usable calculus of homotopy pullbacks on the category of homotopy types. The attraction is clear: four objects and four arrows are much easier to manipulate than a model category. The difficulty is already identified in his next paragraph. The comparison
\[
 \operatorname{dia}_{\square}:
 \mathcal M^{\square}(W^{\square})^{-1}
 \longrightarrow (\mathcal MW^{-1})^{\square}
\]
forgets part of a diagram. Here $\square$ has arrows $3\to1\to0$ and $3\to2\to0$, with the two composites equal; $\Phi\subset\square$ omits $3$. We use covariant diagrams, following the arrows actually drawn. An object on the left is a \emph{coherent square}; on the right it is an ordinary commutative square in the localized category. Calling an object coherent does not mean that its chosen representative has stopped commuting strictly. It means that the whole diagram was localized together.

Write $\operatorname{Cart}_{\mathrm{im}}$ for the replete image, on the right, of the homotopy-cartesian objects on the left. Membership says that \emph{some} coherent realization is cartesian. Proposition~1 will assert, incorrectly, that \emph{every} strict realization is then cartesian. An isomorphism of underlying squares need not lift to an isomorphism of coherent squares. Likewise, four arrows forming a cube after localization need not form a morphism of the chosen coherent squares. These are the two losses that the postface makes explicit.

The opening fibration formula belongs to the assumed theory of homotopy fibres. For an arbitrary Quillen model category, a safe version is that $Y_1\times_{Y_0}Y_2$ computes the homotopy pullback when $Y_1\to Y_0$ is a fibration and $Y_0,Y_2$ are fibrant; right properness permits the familiar version without these fibrancy assumptions \citep[Theorem~5.2.5 and the mixed Reedy argument of Lemma~5.2.6, dualized]{Hovey1999ModelCategories}. Thus the opening paragraph should not be detached from its hypotheses or used to diagnose a false formula in the transcription. The editors state the precise Quillen version below.

The commentary on Chapter~I, manuscript p.~7 (``Coherent squares and what the axioms remember''), and on Chapter~IV, manuscript p.~75 (``Constructing the loop square without choosing homotopies''), supplies the general background. The new issue here is reflection of a universal property through $\operatorname{dia}_{\square}$, and the chapter itself supplies a counterexample.
\end{dcomment}


\noindent\textbf{Définition 1.} J'appelle \textit{“catégorie à carrés h-cartésiens”,} ou \textit{“catégorie quadrangulée à gauche”,} une catégorie $\mathcal{H}$, munie d'un sous-ensemble Cart de l'ensemble $\mathrm{Hom}(\boldsymbol{\Delta}_1\times\boldsymbol{\Delta}_1,\mathcal{H})$
\printedpage{2}
% head: A. Grothendieck : Dérivateurs, Ch. XIV, Édition des Universités de Montpellier II et Paris VII
% foot: Version 7 octobre 2010
des diagrammes carrés \texttt{[par carré on entend carré commutatif]}
\[
\begin{tikzcd}
X_1 \arrow[d] & X_3 \arrow[l] \arrow[d] \\
X_0 & X_2 \arrow[l]
\end{tikzcd}
\]
dans $\mathcal{H}$, carrés appelés \textit{“h-cartésiens”} (h comme “homotopiquement”), cette \msp{page 3} donnée étant soumise aux axiomes qui suivent, Cart (1) à Cart (\ldots).

\dgrblock{dgrb742}{The earlier reason for keeping diagram categories}
\begin{dcomment}{context}{p2}{The earlier reason for keeping diagram categories}
The retreat to diagram categories on manuscript page~2 is part of a long-standing programme. In \emph{Pursuing Stacks}, \S69, Grothendieck relates it to the difficulty of formulating trace additivity and determinant multiplicativity using only a derived category. His proposed object retains the derived categories of diagram categories, including filtered complexes of varying lengths \citep[pp.~1--2 of the edited excerpt]{Grothendieck1983}. The letter to Thomason of 2~April~1991 returns to the mapping-cone problem and the need for this whole family \citep[\S2, pp.~5--6]{Grothendieck1991Letter}. These passages explain why an apparently modest square calculus leads back to derivators.

Heller's \emph{Homotopy Theories} already organizes homotopy Kan extensions on diagram categories \citep[Introduction, pp.~1--2]{Heller1988}. This is a historical parallel, not a warrant for all the proposed square rules. The point of the present fragment is precisely to test how much of that structure can survive on the single category $\mathcal H$.

Two different judgments in the source must therefore be distinguished. The ordinary square is too impoverished to carry the desired coherent universal property; nevertheless, some of the axioms imposed on such squares are too strong for their intended examples. There is no contradiction between these diagnoses. Nor should the editors' suggestion that Grothendieck stopped when he noticed the failure be converted into a documented account of his decision: the editors' postface proves the mathematical obstruction, not that biographical explanation.
\end{dcomment}


\noindent Comme je ne suis pas sûr de connaître d'avance toutes les propriétés essentielles, le nombre de ces axiomes reste pour le moment ouvert.

\noindent\textbf{Cart (1) (Symétrie).} Si un carré
\[
\begin{tikzcd}
X_1 \arrow[d] & X_3 \arrow[l] \arrow[d] \\
X_0 & X_2 \arrow[l]
\end{tikzcd}
\]
est h-cartésien, de même le carré symétrique (obtenu en échangeant le rôle de $X_1$ et de $X_2$).

\noindent\textbf{Cart (2) (Carrés triviaux).} Soit
\[
\begin{tikzcd}
X_1 \arrow[d, "p_1"'] & X_3 \arrow[l, "q_2"'] \arrow[d, "q_1"] \\
X_0 & X_2 \arrow[l, "p_2"]
\end{tikzcd}
\]
un carré dans $\mathcal{H}$, avec $p_1$ un isomorphisme. Alors le carré est h-cartésien si et seulement si $q_1$ est un isomorphisme.

\noindent\textbf{Cart (3) (Carrés composés).} Soit un diagramme commutatif
\[
\begin{tikzcd}
X'_0 \arrow[d] \arrow[dr, phantom, "\mathrm{I}"] & X'_1 \arrow[l] \arrow[d] \arrow[dr, phantom, "\mathrm{II}"] & X'_2 \arrow[l] \arrow[d] & \\
X_0 & X_1 \arrow[l] & X_2 \arrow[l] & ,
\end{tikzcd}
\]
d'où les carrés I, II et III :
\[
\begin{tikzcd}
X'_0 \arrow[d] & X'_2 \arrow[l] \arrow[d] \\
X_0 & X_2 \arrow[l]
\end{tikzcd}
\]
\printedpage{3}
% head: A. Grothendieck : Dérivateurs, Ch. XIV, Édition des Universités de Montpellier II et Paris VII
% foot: Version 7 octobre 2010
(qu'on appelle \textit{carré composé} (“horizontal”) des carrés I et II). Supposons I h-cartésien. Alors II est h-cartésien si et seulement si III l'est \texttt{[voir postface des éditeurs]}.

\noindent\textbf{Corollaire.} \textit{Soit un diagramme commutatif cubique, qu'on peut considérer comme homomorphisme d'un carré} (\textit{face supérieure du cube}) \textit{dans un autre} (\textit{face inférieure}) :
\[
\begin{tikzcd}
 & {} \arrow[dl] \arrow[dd] & & {} \arrow[ll] \arrow[dl] \arrow[dd] \\
{} \arrow[dd] & & {} \arrow[ll, crossing over] & \\
 & {} \arrow[dl] & & {} \arrow[ll] \arrow[dl] \\
{} & & {} \arrow[ll] \arrow[from=uu, crossing over] & .
\end{tikzcd}
\]
\textit{Si toutes les faces du cube}\footnote{En fait, il suffit de la face inférieure, et des deux faces parallèles latérales (p. ex. celle de devant et celle de derrière).} \textit{sauf la face supérieure sont} \msp{page 4} \textit{des carrés h-cartésiens (savoir la face inférieure et les quatre faces latérales), il en est de même de la face supérieure.}\footnote{\textbf{Corollaire} (de Cart (3) et Cart (2)) : \textit{tout carré isomorphe à un carré h-cartésien est h-cartésien.}}

\dgrblock{dgrb743}{Pasting, cancellation, and the cube}
\begin{dcomment}{exposition}{p3}{Pasting, cancellation, and the cube}
The first three axioms try to isolate the part of pullback calculus that can be used without constructing a pullback. Cart~(2) has exactly the usual orientation: if the left vertical arrow is invertible, a square is distinguished precisely when the right vertical arrow is invertible. Cart~(1) supplies the transposed assertion. Cart~(3) then includes both composition and cancellation, with the left square held cartesian. Its failure for $\operatorname{Cart}_{\mathrm{im}}$ will be a failure of cancellation; the postface should not be read as disproving each implication separately.

We can fill in the cube argument entirely formally. Regard the cube as two horizontal squares $S_X,S_Y$ and name the opposite lateral faces
\[
 B=\left(\begin{matrix}X_1&\leftarrow&X_3\\\downarrow&&\downarrow\\Y_1&\leftarrow&Y_3\end{matrix}\right),
 \qquad
 F=\left(\begin{matrix}X_0&\leftarrow&X_2\\\downarrow&&\downarrow\\Y_0&\leftarrow&Y_2\end{matrix}\right).
\]
If $B$ and $S_Y$ are distinguished, their vertical composite $R$ is distinguished. The same rectangle is obtained by composing $S_X$ with $F$. If $F$ is distinguished, vertical cancellation implies that $S_X$ is distinguished. Vertical pasting is horizontal pasting after Cart~(1). This proves exactly the stronger three-face version in the first footnote; the other two lateral faces are unnecessary. An isomorphism of squares makes $B,F$ distinguished by Cart~(2), so the second footnote follows as well.

There is also an economy in the later list: Cart~(5)(b) is already a consequence of Cart~(1)--(3). For a cube between distinguished $S_X,S_Y$, suppose $p_0,p_1,p_2$ are isomorphisms. Then $F$ is distinguished. Composing $S_X$ and $F$ makes $R$ distinguished; cancelling $S_Y$ makes $B$ distinguished. Applying Cart~(2) to $B$ and $p_1$ forces $p_3$ to be an isomorphism. This deduction concerns the abstract axioms. The editors' two counterexamples will show why the intended class cannot satisfy them.
\end{dcomment}


\noindent Ces axiomes sont dans la nature d'axiomes de stabilité, nous donnant des critères affirmant que certains carrés sont h-cartésiens, si telles conditions sont remplies (p. ex. que certains autres carrés sont déjà h-cartésiens). J'ai l'impression d'avoir épuisé ici les propriétés pertinentes de ce type (\textit{i.e.} que les autres devraient pouvoir s'en déduire). À présent, je m'occupe de propriétés d'existence et d'unicité de carrés h-cartésien. % [sic] carrés h-cartésien -> carrés h-cartésiens, missing plural agreement

\noindent\textbf{Cart (4) (Axiome du produit fibré).} Tout diagramme
\[
\begin{tikzcd}
X_1 \arrow[d] & \\
X_0 & X_2 \arrow[l]
\end{tikzcd}
\]
dans $\mathcal{H}$ peut se compléter en un carré h-cartésien
\[
\begin{tikzcd}
X_1 \arrow[d] & X_3 \arrow[l] \arrow[d] & \\
X_0 & X_2 \arrow[l] & .
\end{tikzcd}
\]
L'axiome suivant (\textit{cf.} son corollaire) assurera entre autres que le carré h-cartésien en question est \textit{unique à isomorphisme près} – mais on fera attention qu'il n'est en général pas unique “à isomorphisme unique près”.

% ---- from Der-ChXIV-07-10-10_p006-010.tex
\printedpage{4}
% head: A. Grothendieck : Dérivateurs, Ch. XIV, Édition des Universités de Montpellier II et Paris VII

\textbf{Cart (5) (Axiome d'exactitude)} Considérons deux carrés h-cartésiens
\[
X = \left(\begin{tikzcd}
X_1 \ar[d] & X_3 \ar[l] \ar[d] \\
X_0 & X_2 \ar[l]
\end{tikzcd}\right)
\qquad \text{et} \qquad
Y = \left(\begin{tikzcd}
Y_1 \ar[d] & Y_3 \ar[l] \ar[d] \\
Y_0 & Y_2 \ar[l]
\end{tikzcd}\right)
\]
et un homomorphisme $p = (p_0, p_1, p_2)$ entre les diagrammes de type
\[
\Phi \ = \
\begin{tikzcd}
1 \ar[d] & \\
0 & 2 \ar[l]
\end{tikzcd}
\]
associés (\textit{cf.} diagramme en traits pleins ci-dessous) :

\msp{page 5}
\[
\begin{tikzcd}[row sep=small, column sep=small]
& X_1 \ar[dl] \ar[dd, "p_1" near start] & & X_3 \ar[ll] \ar[dl] \ar[dd, dotted, "p_3"] \\
X_0 \ar[dd, "p_0"'] & & X_2 \ar[ll, crossing over] & \\
& Y_1 \ar[dl] & & Y_3 \ar[ll] \ar[dl] \\
Y_0 & & Y_2 \ar[ll] \ar[from=uu, crossing over, "p_2"' near end] & .
\end{tikzcd}
\]

Alors :
\begin{itemize}[leftmargin=2.5em]
\item[a)] Il existe un homomorphisme $\bar{p} : X \longrightarrow Y$ induisant $p$, \textit{i.e.} il existe une flèche $p_3$ de façon à compléter le diagramme cubique commutatif (en respectant donc la contrainte de commutativité pour les deux carrés adjacents à l'arête $p_3$).
\item[b)] Si de plus $p_0$, $p_1$, $p_2$ sont des isomorphismes, il en est de même de $p_3$ \texttt{[voir postface des éditeurs]}.
\end{itemize}

Prenant le cas où $p_0$, $p_1$, $p_2$ sont des identités, on trouve :

\textbf{Corollaire 1.} \textit{Sous les conditions de} Cart (4)\textit{, si on a} deux \textit{carrés h-cartésiens $X$, $X'$ qui complètent}
\[
\begin{tikzcd}[row sep=small]
& X_1 \ar[dl] & \\
X_0 & & , \\
& X_2 \ar[ul] &
\end{tikzcd}
\]
\textit{il existe un homomorphisme de $X$ dans $X'$ qui induise l'identité sur les sommets d'indices} 0, 1, 2\textit{, et un tel homomorphisme est toujours un isomorphisme.}

Mais cet isomorphisme n'est pas forcément unique, en d'autres termes, un automorphisme d'un carré h-cartésien peut induire l'identité sur les sommets $X_0$, $X_1$, $X_2$, sans être l'identité.
% foot: Version 7 octobre 2010

\dgrblock{dgrb744}{Existence of a comparison is not a chosen comparison}
\begin{dcomment}{exposition}{p5:sans être l'identité.}{Existence of a comparison is not a chosen comparison}
Cart~(4) supplies an object over each cospan. Cart~(5)(a) asks that a map of cospans extend to some map of these objects; it does not specify a choice compatible with composition. Cart~(5)(b), with its universal quantifier on the already chosen $p_3$, says that \emph{every} such extension over an isomorphism of cospans is invertible. Replacing ``every'' by ``there exists an invertible extension'' would be a different axiom.

For a fixed cospan $A\to C\leftarrow B$, consider the category of cones whose objects are triples $(T,a:T\to A,b:T\to B)$ satisfying $fa=gb$, and whose arrows preserve $a,b$. The first corollary says that the full subcategory consisting of distinguished completions is nonempty, that any two of its objects admit a comparison, and that every comparison is invertible. This is a connected groupoid, not a proof that it is a terminal category. Existence follows by Cart~(5)(a) with three identities; invertibility follows by (b). The source's nonuniqueness remark records the distinction between these conclusions. It supplies no separate example satisfying all five axioms with a nontrivial automorphism group.

For the intended homotopy class, the distinction is quite visible: the square with vertices $\Omega X,0,0,X$ and all edges zero permits every automorphism of $\Omega X$ over the fixed cospan. It also permits every \emph{noninvertible} endomorphism, including zero. Thus this example illustrates nonuniqueness but simultaneously invalidates (b); it cannot validate the literal axiom system. In an additive category satisfying the whole system, the development below proves that no such nonuniqueness is possible.

In the coherent repair the relevant base is itself a category of coherent cospans. Restriction from cartesian coherent squares to that category is an equivalence. This gives an honest functorial construction there, while forgetting coherence can again identify distinct maps. Neither a noncanonical choice of apex nor an arbitrary commuting cube restores the information that was forgotten.
\end{dcomment}


\printedpage{5}
% head: A. Grothendieck : Dérivateurs, Ch. XIV, Édition des Universités de Montpellier II et Paris VII

\msp{page 6}
\textbf{Corollaire 2.} \textit{Considérons un carré commutatif}
\[
Y = \left(\begin{tikzcd}
Y_1 \ar[d] & Y_3 \ar[l] \ar[d] \\
Y_0 & Y_2 \ar[l]
\end{tikzcd}\right) \quad ,
\]
\textit{et un objet $T$ de $\mathcal{H}$, d'où (posant} $\mathrm{Hom}(T,X) = X(T)$\textit{) un carré commutatif dans} (Ens)
\[
\begin{tikzcd}
Y_1(T) \ar[d] & Y_3(T) \ar[l] \ar[d] \\
Y_0(T) & Y_2(T) \ar[l]
\end{tikzcd} \quad ,
\]
\textit{et par suite}
\[
Y_3(T) \longrightarrow Y_1(T) \times_{Y_0(T)} Y_2(T) \quad . \tag*{$(*)$}
\]
\textit{Si $Y$ est h-cartésien, cette application} $(*)$ \textit{est} surjective \textit{quel que soit $T$.}

On applique Cart (5) à $Y$ et au carré constant $X$ de valeur $T$
\[
X = \quad
\begin{tikzcd}
T \ar[d, "\mathrm{id}"'] & T \ar[l, "\mathrm{id}"'] \ar[d, "\mathrm{id}"] \\
T & T \ar[l, "\mathrm{id}"']
\end{tikzcd} \quad ,
\]
en notant que la donnée de $p = (p_0, p_1, p_2)$ de
\[
\begin{tikzcd}
T \ar[d] & \\
T & T \ar[l]
\end{tikzcd}
\qquad\qquad \text{dans} \qquad\qquad
\begin{tikzcd}
Y_1 \ar[d] & \\
Y_0 & Y_2 \ar[l]
\end{tikzcd}
\]
revient à la donnée d'un élément du deuxième membre de $(*)$,
\[
\begin{tikzcd}
Y_1 \ar[d] & T \ar[l, "p_1"'] \ar[dl, "p_0"] \ar[d, "p_2"] \\
Y_0 & Y_2 \ar[l]
\end{tikzcd} \quad ,
\]
et que les éléments de $Y_3(T)$ qui s'envoient (par $(*)$) dans cet élément ne sont autres que les $p_3$ dont l'axiome Cart (5) affirme l'existence.

\textbf{N.B.} On notera que l'application surjective $(*)$ n'est en général \textit{pas} bijective, \textit{i.e.} le carré h-cartésien n'est en général pas cartésien.
% foot: Version 7 octobre 2010

\dgrblock{dgrb745}{Weak universality and the information in a homotopy}
\begin{dcomment}{exposition}{p6:le carré h-cartésien n'est en général pas cartésien.}{Weak universality and the information in a homotopy}
The second corollary identifies the positive content of the existence axiom. A distinguished square is a \emph{weak pullback}: every compatible pair of maps from $T$ factors through its apex, without uniqueness. The constant square on $T$ is distinguished by Cart~(2), and applying Cart~(5)(a) to it gives exactly the surjection displayed in the text. Conversely, if every distinguished square is a weak pullback, apply its weak property to $T=X_3$ and the two composites $X_3\to X_i\to Y_i$. This supplies Cart~(5)(a). Thus, given the constant squares, the Hom-set assertion and that part of Cart~(5) are equivalent.

For an actual homotopy pullback the surjection has a geometric explanation. In an enhanced homotopy theory, put
\[
 M_i=\operatorname{Map}(T,Y_i).
 \qquad
 \operatorname{Map}(T,Y_3)\simeq M_1\mathbin{\times^h_{M_0}}M_2.
\]
A point of the homotopy pullback consists of two maps and a path between their composites. Equality in $\pi_0M_0$ ensures that such a path exists, which proves surjectivity after taking components. Different choices of path can give different components. For the loop square, taking $T=S^0$ in pointed spaces and $X=S^1$ gives the map $\pi_0\Omega S^1\cong\mathbb Z\to\{*\}$, a fully explicit failure of injectivity. This statement is about the intended homotopy class, not an example of all five literal axioms. Mapping spaces supply the universal property before $\pi_0$ is taken; compare \citet[\S1.2.13]{Lurie2009HTT}.

There is a useful rigidity argument whenever an ordinary pullback $P=A\times_C B$ already exists in $\mathcal H$. Let $Q$ be a distinguished completion. Its canonical map $q:Q\to P$ has a section $s:P\to Q$ by weak universality. The endomorphism $sq$ of $Q$ is over the identity cospan, so Cart~(5)(b) makes it invertible. But $(sq)^2=sq$; an invertible idempotent is the identity. Hence $q$ is an isomorphism, and $Q$ is an ordinary pullback. Conversely, Cart~(4) and invariance under isomorphisms show that every existing ordinary pullback square is distinguished. Consequently, if $\mathcal H$ has all ordinary pullbacks, the proposed class is forced to be precisely that class. This is the familiar minimality argument for weak universal objects; its weak-kernel form is explicitly recorded by \citet[\S1, ``Weak kernels'', p.~6 of the 2014 manuscript]{Krause2016FlatCovers}.
\end{dcomment}


\printedpage{6}
% head: A. Grothendieck : Dérivateurs, Ch. XIV, Édition des Universités de Montpellier II et Paris VII

\msp{page 7}
\textbf{N.B.} Il y a une notion duale de \textit{catégorie à carrés h-cocartésiens,} ou \textit{catégorie quadrangulée à droite.} Le plus souvent, les deux structures sont présentes simultanément. Il y a lieu de dégager des propriétés qui relient les deux structures duales sur une même catégorie. Ainsi, sauf erreur, dans le cas $\mathcal{H}$ additive, les deux espèces de carrés, les h-cartésiens et les h-cocartésiens, devraient \textit{coïncider.} (Il me faudra revenir là-dessus par la suite.) Dans le cas $\mathcal{H}$ ``ponctuée'', il y a l'adjonction des foncteurs $\Sigma$, $\Omega$ liés l'un à la structure quadrangulée droite, l'autre à la structure quadrangulée gauche. Et il y a, quand il y a des \underline{Hom} internes ou externes, des propriétés éventuelles de $\uHom(X,Y)$ par rapport aux quadrangles h-cocartésiens en $X$, h-cartésiens en $Y$ -- mais c'est à peine un lien entre les deux structures \ldots

\dgrblock{dgrb746}{Additivity, stability, and the two variables of Hom}
\begin{dcomment}{later}{p7:mais c'est à peine un lien entre les deux structures}{Additivity, stability, and the two variables of Hom}
Grothendieck is looking for a common structure behind fibres and cofibres. Additivity alone cannot supply it. Take $\mathcal H=\mathrm{Vect}_k$, with ordinary pullback squares on the left and ordinary pushout squares on the right. Each structure satisfies the corresponding five axioms. For $V\ne0$, the square
\[
 \begin{tikzcd}[column sep=large]
 0\ar[d]&0\ar[l]\ar[d]\\
 V&0\ar[l]
 \end{tikzcd}
\]
is a pullback, since $0\times_V0=0$, but is not a pushout, since $0\amalg_00=0$. This refutes the additive speculation within the literal framework, without using any defective homotopy example.

For coherent diagrams, the corrected hypothesis is \emph{stability}. A pointed derivator is stable when its cartesian and cocartesian coherent squares coincide. The suspension--loop adjunction exists already in the pointed case; the equivalence of that adjunction characterizes stability \citetext{\citealp[Proposition~3.17]{Groth2013}; \citealp[Theorem~7.1, arXiv v2]{GrothPontoShulman2014}}. The corresponding higher-categorical statement appears in \citet[Proposition~1.1.3.4]{Lurie2017HA}. These theorems concern coherent squares, and do not rescue Cart~(3) on the ordinary homotopy category.

The Hom remark also has a precise scope. For fixed $T$, a derived mapping-space functor sends homotopy limits in its second variable to limits of spaces. In the first variable, mapping out of a homotopy colimit produces a homotopy limit. In a pointed enhancement this gives
\[
 \operatorname{Map}(\Sigma X,Y)
 \simeq\Omega\operatorname{Map}(X,Y)
 \simeq\operatorname{Map}(X,\Omega Y),
\]
where loops in the mapping space are based at the zero map; taking components gives the adjunction. Cisinski constructs the external derived Hom for model-derived derivators and proves its continuity \citep[Proposition~6.13, pp.~239--240]{Cisinski2003}. Internal Hom requires a specified closed monoidal structure and the appropriate derived adjunctions. Nothing on this page supplies such structure to an arbitrary category $\mathcal H$.
\end{dcomment}


\dgrblock{dgrb747}{Binary products force the literal structure to be ordinary}
\begin{dcomment}{development}{p7:mais c'est à peine un lien entre les deux structures}{Binary products force the literal structure to be ordinary}
The proposed axioms are meant to retain a weak universal property while allowing nonunique comparisons. The preceding rigidity argument shows that this freedom disappears wherever an ordinary pullback already exists. Pasting makes the restriction stronger: binary products alone force all the missing pullbacks to exist. This applies in particular to additive categories, without any assumption of kernels or of an abelian structure.

We assume Cart~(1)--(5) literally throughout this block. The weak-kernel minimality argument has the antecedent cited above. We checked the editors' postface, the coherent-square discussion of \citet[\S3]{Groth2013}, and \citet[\S1]{Krause2016FlatCovers}, together with targeted literature searches. We did not find the following deductions from these particular axioms there. The marking identifies proofs supplied by us, not a claim of priority.

\textbf{Pullback along a monomorphism.}\ours\quad
A distinguished square with a monomorphism among its two base arrows is an ordinary pullback.

\emph{Proof.} Write the square as the left square below, with $f$ monic. To test uniqueness of its lifting property, apply Cart~(4) to $P\xrightarrow qB\xleftarrow qP$ and obtain the right square:
\[
\begin{tikzcd}[column sep=large]
A\ar[d,"f"'] & P\ar[l,"p"']\ar[d,"q"]
 & R\ar[l,"r"']\ar[d,"s"]\\
C & B\ar[l,"g"] & P\ar[l,"q"] .
\end{tikzcd}
\]
Cart~(3) makes the outer rectangle distinguished. Its base cospan is $A\xrightarrow f C\xleftarrow{fp}P$, because $gq=fp$. Since $f$ is monic, this cospan has ordinary pullback $P$ with projections $p$ and $1_P$. The rigidity argument proved in the preceding exposition makes the canonical comparison $s:R\to P$ an isomorphism.

Weak universality of the right square supplies a diagonal $d:P\to R$ with $rd=sd=1_P$. Thus $d=s^{-1}$ and $r=s$. If $u,v:T\to P$ satisfy $qu=qv$, the same weak property gives $w:T\to R$ with $rw=u$ and $sw=v$; hence $u=v$. Therefore $q$ is monic. The original square is a weak pullback, and its factorization maps are unique because their composites with $q$ agree. It is an ordinary pullback. Cart~(1) treats the other base arrow.\hfill$\square$

\textbf{Collapse under binary products.}\ours\quad
If $\mathcal H$ has binary products and satisfies Cart~(1)--(5), it has all pullbacks, and its distinguished squares are exactly its ordinary pullback squares.

\emph{Proof.} Given $A\xrightarrow fC\xleftarrow gB$, form the cospan
\[
C\xrightarrow{\Delta}C\times C
  \xleftarrow{\,f\times g\,}A\times B.
\]
The diagonal is monic. A distinguished completion exists by Cart~(4), and the preceding result makes it an ordinary pullback. Its apex is $A\times_CB$: a pair $a:T\to A$, $b:T\to B$ with $fa=gb$ gives the unique map $(a,b)$ to $A\times B$ and the compatible map $fa$ to $C$, so the displayed pullback property gives a unique map to the apex. All pullbacks therefore exist. Rigidity shows that each distinguished square is ordinary; Cart~(4) and isomorphism invariance show that every ordinary pullback square is distinguished.\hfill$\square$

No terminal object was used. Additivity supplies the required binary products as biproducts, so the theorem is stronger than an additive statement. It does not classify categories lacking binary products, and we draw no general nonexistence conclusion in that setting.

\textbf{Pointed coincidence obstruction.}\ours\quad
If a pointed category carries both literal structures and their distinguished classes coincide, every object is isomorphic to zero.

\emph{Proof.} For any $X$, the cospan $0\to X\leftarrow0$ has ordinary pullback $0$. Cart~(4), rigidity, and isomorphism invariance designate the square with base $X$ and its other three vertices zero. Coincidence makes it distinguished cocartesian. Its upper span $0\leftarrow0\to0$ has ordinary pushout $0$; dual rigidity therefore gives $X\cong0$.\hfill$\square$

Even without coincidence, the literal loop construction is always zero, by the same pullback argument, and the literal suspension is zero if the right structure is present. Their adjunction would be the vacuous zero-functor adjunction. The meaningful suspension and loop functors in homotopy theory belong to the coherent setting anticipated on manuscript page~2.
\end{dcomment}


\section*{2\quad Exemples}

Un premier exemple trivial, est celui où $\mathcal{H}$ est stable par produits fibrés, et où les carrés h-cartésiens sont définis comme les carrés cartésiens. Les axiomes Cart (1) à Cart (5) sont alors des propriétés bien connues, $\pm$ tautologiques de la notion de carré cartésien. Et
\msp{page 8}
en un sens, le cas général doit se ramener à ce cas particulier banal.

Le seul type d'autre cas que je connaisse, est celui où $\mathcal{H} = \mathcal{M}W^{-1}$, où $W$ est un localiseur admettant une théorie de la fibre $W$-homotopique (\textit{cf.} chapitre XII). Dans ce cas, on l'a dit, on prend comme carrés h-cartésiens dans $\mathcal{H}$ ceux qui sont isomorphes, dans $\uHom(\bDelta_1 \times \bDelta_1, \mathcal{H})$, à un carré $W$-cartésien dans $\mathcal{M}$. Encore faut-il vérifier les axiomes Cart (1) à Cart (5). Je pense que je n'ai guère de chance d'y arriver que moyennant une

\textbf{Proposition 1.} \textit{Moyennant des hypothèses sur $(\mathcal{M}, W)$ qui vont être dégagées ci-dessous, vérifiées en tous cas si $(\mathcal{M}, W)$ est quillénisable, sur la catégorie des ensembles ordonnés finis} (\textit{cf.} XIII, p. 51)\textit{, on a ceci : soit}
\[
X = \left(\begin{tikzcd}
X_1 \ar[d] & X_3 \ar[l] \ar[d] \\
X_0 & X_2 \ar[l]
\end{tikzcd}\right)
\]
\textit{un carré dans $\mathcal{M}$, et soit} $\ho_W(X)$ \textit{son image dans} $\uHom(\bDelta_1 \times \bDelta_1, \mathcal{H})$\textit{. Pour que} $\ho_W(X)$ \textit{soit h-cartésien, il} faut (\textit{et il suffit, bien sûr}) \textit{que $X$ soit $W$-cartésien} \texttt{[voir postface des éditeurs]}\textit{.}
% foot: Version 7 octobre 2010

\dgrblock{dgrb748}{The precise assertion that fails in Proposition 1}
\begin{dcomment}{exposition}{p8:voir postface des éditeurs}{The precise assertion that fails in Proposition 1}
Grothendieck hopes to verify the axioms by transporting a homotopy-cartesian criterion through localization. The sufficient direction of Proposition~1 is the definition of $\operatorname{Cart}_{\mathrm{im}}$: a homotopy-cartesian strict square has a distinguished underlying image. The necessary direction would make homotopy cartesianness a property of the ordinary underlying square alone. The editors disprove precisely this reflection statement.

It matters that the obstruction occurs in Quillen model categories with entirely familiar diagram model structures. It cannot be removed by the promised availability of models for finite ordered indexing categories. The walking square and the walking cube are already finite posets; the example below lives in complexes over a field, where all their diagram model structures are available. An additional hypothesis could restrict the homotopy theory drastically, but Quillen presentability of diagrams does not establish the proposition.

The preliminary example, ordinary pullbacks in an ordinary category, remains valid. Symmetry, the isomorphism test, and pasting follow directly from the pullback universal property; every cospan has a completion by hypothesis; a map of cospans induces a unique map of pullbacks; an invertible base map induces an invertible comparison. These statements also prove the stronger uniqueness omitted from the proposed homotopical axioms. The preceding development gives a precise circumstance in which the general axiom system really does reduce to this example: the existence of binary products. It does not vindicate the intended model-category example.
\end{dcomment}


\printedpage{7}
% head: A. Grothendieck : Dérivateurs, Ch. XIV, Édition des Universités de Montpellier II et Paris VII

\section*{Postface des éditeurs}

La proposition ci-dessus est fausse. On rappelle que si $\mathcal{M}$ est une catégorie de modèles fermée de Quillen de classe d'équivalences faibles $W$, un carré commutatif
\[
X = \left(\begin{tikzcd}
X_1 \ar[d] & X_3 \ar[l] \ar[d] \\
X_0 & X_2 \ar[l]
\end{tikzcd}\right)
\]
de $\mathcal{M}$ est $W$-cartésien si et seulement si il existe un carré \textit{cartésien}
\[
Y = \left(\begin{tikzcd}
Y_1 \ar[d] & Y_3 \ar[l] \ar[d] \\
Y_0 & Y_2 \ar[l]
\end{tikzcd}\right)
\]
avec $Y_1 \longrightarrow Y_0$ une fibration, et $Y_0$ et $Y_2$ fibrants (cette dernière condition étant superflue si la catégorie de modèles $\mathcal{M}$ est propre à droite), et un morphisme de carrés $s : X \longrightarrow Y$, $s = (s_0, s_1, s_2, s_3)$, \textit{i.e.} un cube commutatif
\[
\begin{tikzcd}[row sep=small, column sep=small]
& X_1 \ar[dl] \ar[dd, "s_1" near start] & & X_3 \ar[ll] \ar[dl] \ar[dd, "s_3"] \\
X_0 \ar[dd, "s_0"'] & & X_2 \ar[ll, crossing over] & \\
& Y_1 \ar[dl] & & Y_3 \ar[ll] \ar[dl] \\
Y_0 & & Y_2 \ar[ll] \ar[from=uu, crossing over, "s_2"' near end] & ,
\end{tikzcd}
\]
tel que $s$ soit une équivalence faible argument par argument, \textit{i.e.} $s_0, s_1, s_2, s_3 \in W$. Il est équivalent de demander que pour \textit{tout} diagramme commutatif à flèches verticales des équivalences faibles
\[
\begin{tikzcd}
X_1 \ar[r] \ar[d] & X_0 \ar[d] & X_2 \ar[l] \ar[d] \\
Y_1 \ar[r] & Y_0 & Y_2 \ar[l]
\end{tikzcd}
\]
tel que $Y_1 \longrightarrow Y_0$ soit une fibration, et $Y_0$, $Y_2$ des objets fibrants, la flèche canonique $X_3 \longrightarrow Y_1 \times_{Y_0} Y_2$ soit une équivalence faible.

\dgrblock{dgrb749}{Postface: why one strict comparison cube suffices}
\begin{dcomment}{exposition}{p8:Il est équivalent de demander}{Postface: why one strict comparison cube suffices}
\emph{Editors' postface, printed p.~7.}
The editors replace the cospan before taking its limit. This makes the missing comparison map canonical, and explains the difference between a weak equivalence of whole squares and a mere isomorphism after applying $\operatorname{dia}_{\square}$.

Here is a verification with exactly their fibrancy hypotheses. Give $\Phi=(1\to0\leftarrow2)$ the Reedy degrees $\deg(2)=0$, $\deg(0)=1$, $\deg(1)=2$, with $2\to0$ direct and $1\to0$ inverse. A diagram $Y$ is Reedy fibrant precisely when $Y_0,Y_2$ are fibrant and $Y_1\to Y_0$ is a fibration. The constant-diagram functor is left Quillen: for a (trivial) cofibration its relative latching maps at $1,2$ are that map, and at $0$ the identity. Hence its right adjoint, the pullback functor, is right Quillen. This is the mixed Reedy method of \citet[Theorem~5.2.5 and the proof of Lemma~5.2.6]{Hovey1999ModelCategories}, dualized.

For a cospan $K$, choose a Reedy trivial cofibration $K\to RK$ with fibrant target. Any objectwise weak equivalence $K\to Y$ to a cospan with the printed fibrancy conditions extends to a map $RK\to Y$ by lifting against $Y\to *$. It is a weak equivalence by two-out-of-three. Applying the right Quillen pullback functor preserves it, since both diagrams are fibrant. Consequently the two comparison maps from the original fourth vertex are weak equivalences simultaneously. This proves ``some'' if and only if ``every''. Taking $Y_3=\lim Y$ makes the fourth component of the strict cube the unique induced cone map, so it also proves the single-cube criterion.

For completeness, right properness removes both objectwise fibrancy conditions. If $K_1\to K_0$ is a fibration and $K\to R$ is a fibrant replacement of the cospan, put $K'_1=R_1\times_{R_0}K_0$. Right properness and two-out-of-three make $K_1\to K'_1$ a weak equivalence between fibrant objects of the slice $\mathcal M/K_0$. Pullback to $K_2$ is right Quillen, so
\[
K_1\times_{K_0}K_2\longrightarrow
R_1\times_{R_0}K_2\longrightarrow R_1\times_{R_0}R_2
\]
are weak equivalences: the first by that slice argument, the second by right properness. This verifies the parenthesis in the source, without adding a properness assumption to the fibrant version.
\end{dcomment}


Supposons maintenant que la catégorie de modèles $\mathcal{M}$ soit ponctuée, autrement dit qu'elle admette un objet $*$ à la fois initial et final, et soit $X$ un objet fibrant de $\mathcal{M}$. Choisissons une décomposition de la flèche $* \longrightarrow X$, en une équivalence faible suivie d'une fibration
\[
* \longrightarrow C \longrightarrow X \quad ,
\]
et formons le carré cartésien
% foot: Version 7 octobre 2010
\printedpage{8}
% head: A. Grothendieck : Dérivateurs, Ch. XIV, Édition des Universités de Montpellier II et Paris VII
\[
\begin{tikzcd}
C \ar[d] & \Omega X \ar[l] \ar[d] \\
X & * \ar[l]
\end{tikzcd} \quad ,
\]
où par définition $\Omega X = C \times_X *$ est l'\textit{espace des lacets} de $X$. Ce carré est $W$-cartésien. Par ailleurs, il est isomorphe dans $\uHom(\bDelta_1 \times \bDelta_1,\ \mathcal{M}W^{-1})$ au carré trivial
\[
\begin{tikzcd}
* \ar[d] & \Omega X \ar[l] \ar[d] \\
X & * \ar[l]
\end{tikzcd} \quad .
\]

Néanmoins, ce dernier n'est \textit{pas} en général $W$-cartésien. En effet, s'il l'était, l'unique flèche $s_3 : \Omega X \longrightarrow \Omega X$ de $\mathcal{M}$ rendant commutatif le cube
\[
\begin{tikzcd}[row sep=small, column sep=small]
& * \ar[dl] \ar[dd] & & \Omega X \ar[ll] \ar[dl] \ar[dd, dotted, "s_3"] \\
X \ar[dd, "\parallel"'] & & * \ar[ll, crossing over] & \\ % Audit resolved: vertical identity mark on X; see corrections.csv, PDF leaf 10.
& C \ar[dl] & & \Omega X \ar[ll] \ar[dl] \\
X & & * \ar[ll] \ar[from=uu, crossing over] &
\end{tikzcd}
\]
serait une équivalence faible. Or, comme la flèche $s_3$ se factorise par $* \times_X * \simeq *$, cela impliquerait que $\Omega X \to *$ est une équivalence faible. Cette condition est très restrictive. Si, par exemple, $\mathcal{M}$ est la catégorie de modèles des espaces topologiques pointés, et $X$ un CW-complexe pointé, elle équivaut à la contractibilité de la composante connexe de $X$ contenant le point base. De même, si $\mathcal{M}$ est la catégorie de modèles des complexes d'une catégorie abélienne de Grothendieck et $X$ un complexe, elle signifie que $X$ est acyclique.

\dgrblock{dgrb750}{Postface: a two-term complex and the missing filling}
\begin{dcomment}{examples}{p8:Cette condition est très restrictive.}{Postface: a two-term complex and the missing filling}
\emph{Editors' postface, printed p.~8.}
The counterexample is especially transparent when every map can be written as a scalar. Let $k$ be a nonzero field and use unbounded cochain complexes, with quasi-isomorphisms as weak equivalences and degreewise surjections as fibrations \citep[Theorem~2.3.11, with degrees reversed]{Hovey1999ModelCategories}. Our shift convention is $(V[r])^n=V^{n+r}$, with differential $(-1)^r d_V$.

Take $X=k[0]$ and the acyclic complex $C=(k\xrightarrow{1}k)$ in degrees $0,1$. Let $p:C\to X$ be the identity in degree $0$ and zero in degree $1$. This is a surjective chain map. Its kernel $P$ is $k[-1]$, concentrated in degree $1$. Thus $0\xrightarrow{\sim}C\xrightarrow pX$ is the editors' factorization and
\[
\begin{tikzcd}[column sep=large]
C\ar[d,"p"'] & P\ar[l,"i"']\ar[d]\\
X & 0\ar[l]
\end{tikzcd}
\]
is both an ordinary and a homotopy pullback. The contraction $s^1:C^1\to C^0$ is the identity, and $d_Cs+sd_C=1_C$. In $D(k)$, therefore, $C\cong0$, and the underlying square is isomorphic to the square with apex $P$, base $X$, and all edges zero.

That latter square, taken \emph{strictly} in complexes, has comparison $s_3=0_P$ to the displayed fibrational replacement: commutativity requires $is_3=0$, and $i$ is injective. The comparison induces zero on $H^1(P)=k$, so it is not a quasi-isomorphism. This computes the dotted arrow in the postface and disproves the necessary direction of Proposition~1.

To locate the information lost, write down a homotopy-pullback model for $f:A\to Z\leftarrow B:g$:
\[
\begin{split}
H(f,g)^n&=A^n\oplus B^n\oplus Z^{n-1},\\
d(a,b,z)&=(d_Aa,d_Bb,fa-gb-d_Zz).
\end{split}
\]
The chain-map identities imply $d^2=0$. This is a shifted mapping cone of $f-g$, with the last coordinate negated relative to one common convention. A cone with maps $u:T\to A$, $v:T\to B$ and a specified degree $-1$ homotopy $h$ satisfying
\[
 d_Zh+hd_T=fu-gv
\]
has comparison $c_h=(u,v,h):T\to H(f,g)$. The displayed identity is exactly its chain-map condition. The coherent square is cartesian precisely when this comparison is a quasi-isomorphism.

For $A=B=0$, $Z=X=k[0]$, and $T=P=k[-1]$, the four edges are zero and $H(0,0)=P$. Every scalar $\lambda\in k$ specifies a filling $h_\lambda^1=\lambda$, and $c_{h_\lambda}$ is multiplication by $\lambda$ on $P$. It is a quasi-isomorphism exactly when $\lambda\ne0$. Its mapping cone is the two-term complex
\[
 k\xrightarrow{\lambda}k
 \quad\hbox{in degrees }0,1,
\]
which is acyclic for $\lambda\ne0$ and has $k$ in both cohomology degrees for $\lambda=0$. Thus identical vertices and edges support both a cartesian and a noncartesian coherent filling.

We can check the transport sign as well. Put $t=ps$, so $d_Xt+td_C=p$. The map
\[
 H(p,0)\longrightarrow X[-1],\qquad (c,z)\longmapsto z-tc
\]
is a chain map with acyclic kernel, the diagonal copy of $C$. Composing with $P\to H(p,0)$, $a\mapsto(i(a),0)$, gives $-1_P$. The strict pullback transported through this contraction has filling $h_{-1}$, whereas the strict zero-edge square has filling $h_0$. The sign depends on the displayed cone convention; its being nonzero does not. With the boundary and its identifications fixed, the classes of fillings are indexed by $k$, because $\underline{\operatorname{Hom}}(P,X)$ is concentrated in degree $-1$ and has no degree $-2$ boundaries. Allowing scalar automorphisms of the apex identifies all nonzero fillings, while the zero filling remains separate.

The editors' other interpretations agree with this computation. In pointed spaces, $X=S^1$ gives $\pi_0\Omega X=\mathbb Z$, so the zero self-map cannot be a weak equivalence. For a CW complex in general, the homotopy groups of its loop space detect all positive homotopy groups of the base component; Whitehead's theorem then makes exactly that component contractible \citep[Theorem~4.5 and \S4.3]{Hatcher2002}. Other components are unrestricted. For unbounded complexes in a Grothendieck abelian category, an acyclic middle term in $0\to P\to C\to X\to0$ gives $H^n(P)\cong H^{n-1}(X)$, so $P$ is acyclic exactly when $X$ is. The standard injective model structure supplies this setting \citep[Theorem~2.2 and Proposition~2.12 of the preprint]{Hovey2001Sheaves}; no assertion that acyclic complexes are chain-contractible is needed.
\end{dcomment}


Des considérations analogues montrent que les carrés $h$-cartésiens de $\mathcal{M}W^{-1}$, définis dans la section 2 comme étant les carrés commutatifs de $\mathcal{M}W^{-1}$ qui sont isomorphes dans $\uHom(\bDelta_1 \times \bDelta_1,\ \mathcal{M}W^{-1})$ à un carré $W$-cartésien, ne satisfont pas à toutes les conditions de Cart (3) et Cart (5) de la section 1. En effet, supposons toujours que $\mathcal{M}$ soit ponctuée, et gardons les notations ci-dessus. Considérons le diagramme commutatif dans $\mathcal{M}W^{-1}$
\[
\begin{tikzcd}
* \ar[d] & \Omega X \ar[l] \ar[d] & \Omega X \ar[l, "*_{\Omega X}"'] \ar[d] \\
X & * \ar[l] & * \ar[l]
\end{tikzcd} \quad ,
\]
où $*_{\Omega X}$ désigne l'unique endomorphisme de $\Omega X$ se factorisant par $*$. Il résulte de ce qui précède que le carré de gauche et le carré composé sont $h$-cartésiens. En revanche, le carré de droite ne l'est pas en général. En effet, si ce carré était isomorphe dans $\uHom(\bDelta_1 \times \bDelta_1,\ \mathcal{M}W^{-1})$ à un carré $W$-cartésien, la flèche $*_{\Omega X}$ serait un isomorphisme
% foot: Version 7 octobre 2010

% ---- from Der-ChXIV-07-10-10_p011-011.tex
\printedpage{9}
% head: A. Grothendieck : Dérivateurs, Ch. XIV, Édition des Universités de Montpellier II et Paris VII
% foot: Version 7 octobre 2010
de $\cM W^{-1}$, donc une équivalence faible de $\cM$, ce qui impliquerait les mêmes conditions très restrictives que ci-dessus. Ainsi l'une des implications de l'axiome Cart 3 n'est pas satisfaite.

De même le cube commutatif dans $\cM W^{-1}$
\[
\begin{tikzcd}[row sep=small, column sep=small]
 & * \arrow[dl] \arrow[dd] & & \Omega X \arrow[ll] \arrow[dl] \arrow[dd, "*_{\Omega X}"] \\
X \arrow[dd, "\shortparallel"'] & & * \arrow[ll, crossing over] & \\
 & * \arrow[dl] & & \Omega X \arrow[ll] \arrow[dl] \\
X & & * \arrow[ll] \arrow[from=uu, crossing over] & \mathllap{,}
\end{tikzcd}
\]
dont les faces horizontales sont $h$-cartésiennes, montre que la condition $(b)$ de l'axiome Cart (5) n'est pas satisfaite en général.


\dgrblock{dgrb751}{Postface: the two failed quantifiers}
\begin{dcomment}{status}{p8:dont les faces horizontales sont}{Postface: the two failed quantifiers}
\emph{Editors' postface, printed pp.~8--9.}
The two final counterexamples use the same object $P=\Omega X$ and the same underlying zero-edge square $Q$. They apply whenever $P$ is not isomorphic to zero, a condition realized by both worked examples above. In the horizontal pair, the left square and outer rectangle are $Q$, hence belong to $\operatorname{Cart}_{\mathrm{im}}$. In the right square the bottom edge is $1_0$ and the top edge is $0_P$. A homotopy pullback of an equivalence has an equivalence as opposite edge, so this right square cannot belong to $\operatorname{Cart}_{\mathrm{im}}$. This is exactly the implication ``I and III cartesian imply II cartesian'' that fails.

For Cart~(5)(b), use the endomorphism of $Q$ that is the identity at the three cospan vertices and zero at $P$. It is a commuting cube in the ordinary homotopy category, but its fourth component is not invertible. The positive existence statement in (a), and the failure of uniqueness already noticed in Corollary~2, should not be conflated with this false assertion about \emph{every} fourth component. As the formal deduction above shows, (b) would have followed from the full pasting axiom anyway.

The implication $0_P\in W\Rightarrow P\to0\in W$ used in the postface needs no cofibrancy assumption. Writing $r:P\to0$ and $s:0\to P$, the arrow $r$ is a retract, in the arrow category, of $sr=0_P$, using components $(1_P,s)$ and $(1_P,r)$. Weak equivalences are closed under retracts. The converse follows from two-out-of-three. Equivalently, an invertible zero endomorphism makes the object zero in the homotopy category. These are settled counterexamples, not conjectures awaiting a stronger model-category theorem.
\end{dcomment}


\dgrblock{dgrb752}{Why composition nevertheless survives}
\begin{dcomment}{exposition}{p8:dont les faces horizontales sont}{Why composition nevertheless survives}
The failure of cancellation leaves a useful positive statement: for every Quillen model category, $\operatorname{Cart}_{\mathrm{im}}$ is closed under composition. Two adjacent cartesian witnesses can be made compatible along their shared edge. A left square and an outer rectangle, as in the counterexample, are different data. We give the rectification argument to keep this distinction precise.

For any model category $\mathcal A$, the arrow comparison
\[
U_{\mathcal A}:\operatorname{Ho}(\mathcal A^{[1]})
 \longrightarrow\operatorname{Ho}(\mathcal A)^{[1]}
\]
is full, essentially surjective, and conservative. The first two properties are the arrow-shape lifting property used in derivator strength \citep[Lemma~1.9]{Groth2013}. To verify fullness directly, take a Reedy-cofibrant source arrow $f:X_0\to X_1$ and a Reedy-fibrant target arrow $g:Y_0\to Y_1$. Thus $X_0$ is cofibrant, $f$ is a cofibration, and $Y_0,Y_1$ are fibrant. Represent a square in the homotopy category by maps $a_0,a_1$ and a right homotopy between $a_1f$ and $ga_0$. The endpoint map $\mathrm{ev}_0:\operatorname{Path}(Y_1)\to Y_1$ is a trivial fibration. Lifting against $f$ extends that homotopy from $X_0$ to $X_1$, replacing $a_1$ by a homotopic $a'_1$ with $a'_1f=ga_0$ strictly. This gives the required coherent-arrow morphism. Any arrow in the homotopy category has an actual representative between cofibrant-fibrant objects, proving essential surjectivity. Finally, pointwise weak equivalences and saturation imply conservativity. In particular, an isomorphism of underlying arrows lifts to an isomorphism of coherent arrows.

Now choose cartesian strict lifts of the two adjacent squares. Their common columns may differ, but the given underlying identifications specify an isomorphism between these columns in $\operatorname{Ho}(\mathcal M)^{[1]}$. Apply the preceding result for $\mathcal A=\mathcal M$ to lift it. The two squares can then be regarded as composable morphisms in $\operatorname{Ho}(\mathcal M^{[1]})$. Choose bifibrant representatives for the three column objects in the arrow model category, with a single shared middle representative, and represent the two morphisms by actual maps. They form one strict rectangle in $\mathcal M$.

Apply the arrow comparison a second time, now to $\mathcal A=\mathcal M^{[1]}$. Each new small square has an underlying arrow in $\operatorname{Ho}(\mathcal A)$ isomorphic to that of its original witness. Fullness lifts this isomorphism and conservativity makes the lift invertible in $\operatorname{Ho}(\mathcal M^{\square})$. Thus the two new squares are cartesian as coherent squares. Coherent pasting makes their outer rectangle cartesian, and its underlying square is isomorphic to the required composite. This proves the assertion without a combinatorial or stability assumption, and without lifting an arbitrary incoherent cube.

In the stable setting, \citet[Proposition~6.11, pp.~17--18 of the 7~June~2021 author version]{ChristensenFrankland2022} prove the corresponding pasting statement even for homotopy-cartesian squares in an arbitrary triangulated category. The valid remnant of the proposed calculus is therefore substantial: symmetry, the isomorphism test, completion of cospans, weak universality, and positive composition survive for the homotopy-image class. Cancellation and Cart~(5)(b) do not.
\end{dcomment}


\dgrblock{dgrb753}{The coherent replacement of the square axioms}
\begin{dcomment}{later}{p8:dont les faces horizontales sont}{The coherent replacement of the square axioms}
\emph{After the editors' postface.}
The construction Grothendieck names on manuscript page~2 gives the appropriate replacement. Let $\mathbb D$ be a derivator and $j:\Phi\hookrightarrow\square$ the full inclusion omitting the prospective pullback vertex. A coherent square $Q\in\mathbb D(\square)$ is cartesian when the adjunction unit
\[
 Q\longrightarrow j_*j^*Q
\]
is invertible. Right Kan extension along $j$ is fully faithful; consequently
\[
 j^*:\mathbb D(\square)_{\mathrm{cart}}
       \simeq\mathbb D(\Phi)
\]
is an equivalence, with inverse $j_*$ \citep[Proposition~1.20 and Definition~3.9]{Groth2013}. A coherent cospan therefore has a cartesian completion, and a morphism of coherent cospans has a unique extension between these completions once their identifications are fixed. If its three components are isomorphisms, pointwise conservativity makes it an isomorphism in $\mathbb D(\Phi)$; the equivalence makes the extended square morphism invertible. This proves the correct fourth-vertex assertion. Its input is a morphism in the coherent category, not an arbitrary cube in $\mathbb D(e)$.

The full pasting-and-cancellation law is a statement about one coherent rectangle. To see the mechanism in a model or higher-category enhancement, write it as
\[
\begin{tikzcd}
A\ar[d]&B\ar[l]\ar[d]&C\ar[l]\ar[d]\\
D&E\ar[l]&F\ar[l].
\end{tikzcd}
\]
If the left square is cartesian, $B\simeq A\times^h_D E$. Associativity of the homotopy limit gives
\[
 B\times^h_E F\simeq
 (A\times^h_D E)\times^h_E F\simeq A\times^h_D F.
\]
Coherence identifies the comparison from $C$ to the first expression with its comparison to the last. One is an equivalence exactly when the other is. For derivators this is the right-Kan-extension pasting statement of \citet[Proposition~3.13]{Groth2013}; the higher-categorical proof is given, dually, by \citet[Lemma~4.4.2.1]{Lurie2009HTT}. The editors' rectangle loses the compatibility needed to make those two comparison maps agree.

In a stable enhancement there is a useful translation into the underlying triangulated category. A coherent cartesian square with projections $u:P\to A$, $v:P\to B$ and base maps $f:A\to C$, $g:B\to C$ supplies a distinguished triangle
\[
 P\xrightarrow{(u,-v)}A\oplus B
 \xrightarrow{(f,g)}C\xrightarrow{\delta}\Sigma P
\]
\citep[Theorem~6.1, arXiv v2]{GrothPontoShulman2014}. The connecting morphism is part of this conclusion. In the zero-edge loop example it is an isomorphism $X\to\Sigma P$, with the sign dictated by conventions. The cube with fourth component $0_P$ and base component $1_X$ cannot preserve that morphism, since it would require $\delta=(\Sigma0_P)\delta=0$. The triangle therefore diagnoses exactly why that cube is incoherent, even though its four visible faces commute in the homotopy category.

The relation with higher categories is a comparison of enhanced diagram theories. For a suitable model presentation, localizing diagram categories gives the homotopy categories of coherent diagrams; precise comparisons appear in \citet[Propositions~1.3.4.24--1.3.4.25]{Lurie2017HA}. It is not necessary here to identify every abstract derivator with an infinity category. The correct conclusion is narrower and sufficient: the familiar homotopy theories retain the diagram information on which the failed inference depended.
\end{dcomment}
\begin{otherlanguage}{english}
\dgroverview{XIV}{Chapter XIV \textperiodcentered\ \textup{XIV} \textperiodcentered\ Overview}
\selectlanguage{english}
Chapter XIV is a short experiment in reducing homotopy theory to a calculus of squares in a single category. Its eight manuscript pages introduce five proposed axioms and begin the examples. The editors' postface, which has no manuscript pagination, explains why the intended model-category example fails. The fragment is especially instructive because Grothendieck identifies the necessary finer setting before attempting the reduction: square diagrams must be localized as diagrams, together with their maps, rather than reconstructed from four arrows in the localized category.

The mathematical object under discussion can be described without ambiguity. For $\mathcal H=\mathcal MW^{-1}$, let $\operatorname{Cart}_{\mathrm{im}}$ consist of those commutative squares in $\mathcal H$ that are isomorphic to the underlying images of homotopy-cartesian strict squares in $\mathcal M$. With $\square$ denoting the walking commutative square, the relevant functor is
\[
\operatorname{dia}_{\square}:
\mathcal M^{\square}(W^{\square})^{-1}
\longrightarrow\mathcal H^{\square}.
\]
The domain retains coherent diagrams. The target retains ordinary commutativity after localization. Membership in $\operatorname{Cart}_{\mathrm{im}}$ is an existential statement about a coherent realization. Proposition~1 would promote it to a statement about every strict realization, and this is the false inference.

There is a specific historical reason for wanting the larger diagram theory. In \emph{Pursuing Stacks}, \S69, Grothendieck connects it with the trace and determinant problems that motivated him in the period of SGA~5; the letter to Thomason of 2~April~1991 again stresses derived diagram categories and mapping cones \citep{Grothendieck1983,Grothendieck1991Letter}. The present chapter tests whether one can retain a useful part of their calculus in a simpler language. The editors' preliminary notice suggests that recognition of the failed example explains the abrupt ending. That suggestion is conjectural; the postface establishes the mathematical failure, not the author's reason for interrupting the manuscript.

The formal axioms are consistent: any category with pullbacks, equipped with its ordinary pullback squares, satisfies them. Cart~(1) gives symmetry, Cart~(2) tests trivial squares against isomorphisms, and Cart~(3) combines pasting with cancellation. Cart~(4) requires a completion of every cospan. Cart~(5)(a) extends a map between the three-vertex boundaries of two distinguished squares; (b) asserts that every such extension is invertible whenever the boundary map is invertible. The commentary supplies the cube proof, including the stronger three-face hypothesis in the footnote, and proves that (b) already follows from Cart~(1)--(3).

The source's second corollary isolates the weak universal property that does survive in homotopy theory. For a homotopy pullback $P$ of $A\to C\leftarrow B$, there is a surjection
\[
[T,P]\longrightarrow[T,A]\times_{[T,C]}[T,B].
\]
A compatible pair in the target admits a homotopy between its composites; a lift remembers a choice of that homotopy. Passing to components forgets the choice, so the map need not be injective. In the loop square for $S^1$, with $T=S^0$ in pointed spaces, it is the map $\mathbb Z\to\{*\}$. The surjection is equivalent to Cart~(5)(a) for a class containing the constant squares. The much stronger assertion in (b) does not follow from this weak property alone.

Positive composition does hold for the homotopy-image class of every Quillen model category. The commentary proves this by applying fullness and conservativity of the arrow comparison twice, first to the model category and then to its arrow category, so that two cartesian witnesses acquire a single shared middle column. This is compatible with the postface: its counterexample concerns cancellation.

The editors' example distinguishes the underlying square from its realization. Factor the basepoint of a fibrant object as
\[
0\xrightarrow{\sim}C\twoheadrightarrow X,
\qquad P=C\times_X0.
\]
The resulting strict pullback is homotopy cartesian. Its image in $\mathcal H^{\square}$ is isomorphic to the square with apex $P$, base $X$, the other vertices zero, and all edges zero. But that zero-edge square, regarded strictly in $\mathcal M$, has comparison $0_P:P\to P$ to the fibrational replacement. It is homotopy cartesian exactly when $P$ is weakly zero. This proves the failure of Proposition~1 whenever the loop object is nontrivial. The same underlying square gives the two final counterexamples: a horizontal rectangle with left and outer squares distinguished but right square not distinguished, and a cube whose first three components are identities while the fourth is zero.

The worked algebraic example takes $X=k[0]$ and $C=(k\xrightarrow1k)$ in degrees $0,1$. Then $P=k[-1]$. A specified filling of the fixed zero-edge square is a scalar $\lambda$; the homotopy-pullback comparison is multiplication by $\lambda$ on $P$. Its cone is $k\xrightarrow\lambda k$, so the coherent square is cartesian precisely for $\lambda\ne0$. The commentary computes the transport through a contraction and its sign explicitly. It also checks the editors' topological qualification: for a CW complex, weak triviality of the loop object means contractibility of the base component, with no restriction on other components. For unbounded complexes in a Grothendieck abelian category it means acyclicity.

The literal axioms have a stronger limitation than the displayed counterexample alone reveals. Our development proves that, if $\mathcal H$ has binary products, the axioms force all ordinary pullbacks to exist and force the distinguished class to be exactly their class. The proof first compares a distinguished weak pullback with an existing ordinary one: a section and the invertibility of a cone-preserving idempotent make the comparison an isomorphism. This is the minimality argument familiar for weak kernels \citep[\S1, p.~6 of the 2014 manuscript]{Krause2016FlatCovers}. Pasting a distinguished square with a weak kernel pair then shows that pullbacks along monomorphisms are ordinary. Finally one pulls back the diagonal $C\to C\times C$ along $A\times B\to C\times C$. The complete argument is given in the marked development block; its application to these axioms was not located in the bounded literature search, and no priority is claimed.

Additive categories satisfy the binary-products hypothesis. Moreover, even without products, the literal loop object in a pointed category must be zero, because $0\to X\leftarrow0$ already has ordinary pullback zero. If the dual structure also exists, its suspension is zero. If the two classes coincide, the square with base $X$ and three zero vertices makes $X$ itself zero. Thus Grothendieck's additive coincidence speculation is false for the literal structures: ordinary pullbacks and pushouts in vector spaces give an immediate counterexample.

The modern positive statement belongs to coherent squares. For the inclusion $j:\Phi\hookrightarrow\square$ of the cospan, cartesian coherent squares are the essential image of the fully faithful right Kan extension $j_*$. Restriction consequently identifies their category with $\mathbb D(\Phi)$. This makes the completion functorial and proves the correct fourth-vertex assertion for coherent morphisms. Coherent rectangles admit pasting and cancellation. In a pointed derivator the suspension--loop adjunction exists; stability, rather than additivity, is the condition under which cartesian and cocartesian coherent squares agree, equivalently that adjunction is an equivalence \citetext{\citealp[Proposition~3.17]{Groth2013}; \citealp[Theorem~7.1, arXiv v2]{GrothPontoShulman2014}}. External derived Hom has the expected continuity, with the usual enrichment hypotheses for an internal version \citep[Proposition~6.13]{Cisinski2003}. None of these assertions makes the forgetful functor reflect the property disproved in the postface.

The proposed main example and Proposition~1 are therefore settled negatively, while the coherent replacements are established theorems. The manuscript's broad request for further compatibility axioms does not by itself specify a present open problem. This commentary does not classify all literal quadrangulated categories without binary products or produce a nonordinary example in that generality. Its audit keeps that limit separate from the explicit counterexamples and the proved structural deductions.
\end{otherlanguage}
\dgrchunkend

% DER_XV-a: Grothendieck's text with the commentary set in, then the overview, inside a group that applies the chunk's own definitions
\dgrchapter{XV}{Chapter XV}{Théorèmes de factorisation (Modèles (2))}
\dgrpart{0}{XV-a}{Part 1 of 2, manuscript pages 1--66}{XV-a\enspace\textperiodcentered\enspace Chapter XV, part 1 of 2\enspace\textperiodcentered\enspace manuscript pages 1--66}{Chapter XV \textperiodcentered\ \textup{XV-a}}
\dgrchunkbegin
% DER_XV-a: the definitions of the chunk's preamble, applied inside its group
\let\printedpage\relax\newcommand{\printedpage}[1]{\marginpar{\footnotesize p.~#1}}
\let\msp\relax\newcommand{\msp}[1]{\marginpar{\footnotesize [#1]}}
\let\M\relax\newcommand{\M}{\mathcal{M}}
\let\Fl\relax\newcommand{\Fl}{\mathrm{Fl}}
\let\uFl\relax\newcommand{\uFl}{\underline{\mathrm{Fl}}}
\let\uNerf\relax\newcommand{\uNerf}{\underline{\mathrm{Nerf}}}
\let\cM\relax\newcommand{\cM}{\mathcal{M}}
\let\Flu\relax\newcommand{\Flu}{\underline{\mathrm{Fl}}}
\let\fU\relax\newcommand{\fU}{\mathfrak{U}}
\let\fV\relax\newcommand{\fV}{\mathfrak{V}}
\let\Filt\relax\newcommand{\Filt}{\mathrm{Filt}}
\let\card\relax\newcommand{\card}{\operatorname{card}}
\let\circled\relax\newcommand{\circled}[1]{\tikz[baseline=(c.base)]\node[draw,circle,inner sep=0.5pt](c){#1};}
\let\pfn\relax\newcommand{\pfn}[2]{\textsuperscript{(}\footnote[#1]{#2}\textsuperscript{)}}
\let\uFilt\relax\newcommand{\uFilt}{\underline{\mathrm{Filt}}}
\let\Hom\relax\newcommand{\Hom}{\operatorname{Hom}}
\let\Ob\relax\newcommand{\Ob}{\operatorname{Ob}}
\let\wavy\relax\newcommand{\wavy}[1]{\widetilde{#1}}
\let\id\relax\newcommand{\id}{\mathrm{id}}
\let\qpre\relax\newcommand{\qpre}{\mathrel{\leftarrow\!\!\!\!\mapsto}}
\let\Cof\relax\newcommand{\Cof}{\mathrm{Cof}}
\let\deq\relax\newcommand{\deq}{\overset{\text{déf}}{=}}
\let\univ\relax\newcommand{\univ}{\mathrm{univ}}
\let\cC\relax\newcommand{\cC}{\mathcal{C}}
\let\Ss\relax\newcommand{\Ss}{\mathrm{Ss}}
\let\Img\relax\newcommand{\Img}{\mathrm{Im}\,}
\let\isoto\relax\newcommand{\isoto}{\xrightarrow{\ \sim\ }}
\let\hookto\relax\newcommand{\hookto}{\lhook\joinrel\longrightarrow}
\let\Ci\relax\newcommand{\Ci}{\mathcal{C}}

% ---- from Der-ChXV-05-11-10_p001-005.tex
% ===================== PDF page 1 =====================
\printedpage{---}
% no page number printed on this page
% head: (none)

\begin{center}
{\Large\textsf{LES DÉRIVATEURS}}

\bigskip

\textbf{ALEXANDRE GROTHENDIECK}

\vspace{2cm}

\textbf{Chapitre XV}

\medskip

\textbf{Théorèmes de factorisation}

\textbf{(Modèles (2))}
\end{center}

\clearpage

% ===================== PDF page 2 =====================
\printedpage{---}
% no page number printed on this page
% head: (none)

Ce texte a été déchiffré et transcrit en \LaTeXe{} par M. Künzer. Il a été édité par M. Künzer, J. Malgoire, et G. Maltsiniotis. La transcription est aussi fidèle que possible au manuscrit. Seules quelques corrections mineures évidentes ont été effectuées. Pour les rares ajouts ou commentaires des éditeurs, ainsi que pour la numérotation originale des pages du manuscrit, les caractères de machine à écrire \texttt{[typewriter]} entre crochets sont utilisés.

Cette édition est provisoire. Les remarques, commentaires et corrections sont bienvenus. Envoyer un message à :

\textbf{maltsin@math.jussieu.fr}

\hfill G. Maltsiniotis\hspace*{2cm}

\clearpage

% ===================== PDF page 3 =====================
\printedpage{1}
% head: A. Grothendieck : Dérivateurs, Ch. XV, Édition des Universités de Montpellier II et Paris VII
% foot: Version 5 novembre 2010

\msp{page 1}

\section*{1\quad Construction fondamentale : la suite transfinie de factorisation d'une flèche $\boldsymbol{f}$.}

Soit $\M$ une $\mathfrak{U}$-catégorie, stable par “petites” (\textit{i.e.} $\mathfrak{U}$-petites) $\varinjlim$ filtrantes. Soit
\[
C_0 \;\subset\; \Fl(\M)
\tag*{(1.1)}
\]
un \textit{petit} ensemble de flèches, on se propose de construire pour toute flèche $f$ de $\M$, une factorisation
\[
f \;=\; p\,i \,, \qquad i \;\in\; \widetilde{C_0} = ((C_0)_*)^* \;, \;\; p \;\in\; (C_0)_* \;.
\tag*{(1.2)}
\]
Pour ceci, on va construire une suite ordinale de foncteurs
\[
\left\{
\begin{array}{rcl}
f & \longmapsto & X_\alpha(f) \;=\; X_\alpha(C_0,f) \\
X_\alpha : \uFl(\M) & \longrightarrow & \M
\end{array}
\right.
\tag*{(1.3)}
\]
et des homomorphismes fonctoriels
\[
i_\alpha(f) : \mathrm{source}\, f \;\longrightarrow\; X_\alpha(f) \;,
\tag*{(1.4)}
\]
\[
p_\alpha(f) : X_\alpha(f) \;\longrightarrow\; \mathrm{but}\, f \;,
\tag*{(1.5)}
\]
de sorte qu'on aura le diagramme (où $X = \mathrm{source}\, f$, $Y = \mathrm{but}\, f$)
\[
\begin{tikzcd}[column sep=small]
X \arrow[rr, "f"] \arrow[dr, "i_\alpha(f)"'] & & Y \\
& X_\alpha(f) \arrow[ur, "p_\alpha(f)"'] &
\end{tikzcd}
\qquad ,
\tag*{(1.6)}
\]
ce diagramme étant commutatif
\[
f \;=\; p_\alpha(f) i_\alpha(f) \;.
\tag*{(1.7)}
\]
On peut dire aussi qu'on a une section du foncteur canonique de “composition”, transposé de $\boldsymbol{\Delta}_1 \longrightarrow \boldsymbol{\Delta}_2$, donné par $0 \longmapsto 0$, $1 \longmapsto 2$ :
\[
\left\{
\begin{array}{l}
\begin{tikzcd}[ampersand replacement=\&]
\uNerf_2(\M) \arrow[r, shift left, "\gamma"] \& \uNerf_1(\M) \arrow[l, dotted, shift left, "\Phi_\alpha"]
\end{tikzcd}
\;=\; \uFl(\M) \;, \\[1ex]
\\[1ex]
\\
\end{array}
\right.
\tag*{(1.8)}
\]
$\Phi$ comme “factorisation”,
\printedpage{2}
% head: A. Grothendieck : Dérivateurs, Ch. XV, Édition des Universités de Montpellier II et Paris VII
% foot: Version 5 novembre 2010
\msp{page 2}
où
\[
\Phi_\alpha(f) \;=\; \{\; X \xrightarrow{i_\alpha(f)} X_\alpha(f) \xrightarrow{p_\alpha(f)} Y \;\}
\tag*{(1.9)}
\]
avec $p_\alpha(f)i_\alpha(f) = f$. La fonctorialité de 1.4 et 1.5, ou celle de $\Phi_\alpha$, \textit{i.e.} de 1.9, s'exprime par le fait que tout carré commutatif
\[
\begin{tikzcd}[row sep=large]
X \arrow[r, "f"] \arrow[d, "u_X"'] & Y \arrow[d, "u_Y"] \\
X' \arrow[r, "{f'}"] & Y'
\end{tikzcd}
\tag*{(1.10)}
\]
dans $\M$, interprété comme flèche $u = (u_X, u_Y)$ de $f$ dans $f'$ dans $\uFl(\M)$, donne naissance au double carré commutatif
\[
\begin{tikzcd}[row sep=large]
X \arrow[r, "i_\alpha(f)"] \arrow[d, "u_X"] & X_\alpha(f) \arrow[r, "p_\alpha(f)"] \arrow[d, "X_\alpha(u)"] & Y \arrow[d, "u_Y"] \\
X' \arrow[r, "{i_\alpha(f')}"] & X_\alpha(f') \arrow[r, "{p_\alpha(f')}"] & Y' \;.
\end{tikzcd}
\tag*{(1.11)}
\]
\textbf{N.B.} Quand on voudra expliciter la dépendance de ces constructions par rapport à $C_0$, on écrira
\[
X_\alpha(C_0,f) \;, \qquad i_\alpha(C_0,f) \;, \qquad p_\alpha(C_0,f) \;, \qquad \Phi_\alpha(C_0,f) \;,
\]
ou
\[
X_\alpha(C_0,-) \;\;\text{ou}\;\; X_\alpha^{C_0} \;, \qquad i_\alpha(C_0,-) \;\;\text{ou}\;\; i_\alpha^{C_0} \qquad \textit{etc.} \;,
\]
au lieu de $X_\alpha(f)$, $i_\alpha(f)$, \dots

\dgrblock{dgrb717}{Why this construction belongs in a theory of derivators}
\begin{dcomment}{context}{p1}{Why this construction belongs in a theory of derivators}
Grothendieck needs factorizations which can be used uniformly across a category of coefficients. A construction for isolated arrows would not supply the functorial replacements needed when those coefficients occur in diagrams. The emphasis on an endofunctor of the arrow category, and then on an entire transfinite sequence of such endofunctors, should be read with that purpose in view.

Two strands of the chapter have different origins. Its language of lifting and factorization is that of homotopical algebra; its treatment of size continues the cardinal filtrations and accessibility developed in \citep[Expos\'e~I, \S9]{SGA4I1972}. A filtration does more than assign a cardinal label to an object: it specifies presentations, closure under controlled colimits, and the bounds within which represented functors commute with those colimits. Here that apparatus will turn an unbounded repetition of a local operation into a uniform existence theorem.

The commentary on Chapter~VI already discusses the two forms of the replacement axiom, by a functorial replacement and by a factorization ``à la Quillen'', at manuscript p.~108, and the distinction between replacements at manuscript p.~180. The present pages address the additional existence problem: how to manufacture the replacements from a set of maps. They should therefore be read as preparation for the diagram theory, rather than as another formulation of its axioms. The modern account of the same small object mechanism in \citep[Chapter~12]{Riehl2014} makes this separation of construction and use especially clear.
\end{dcomment}


\dgrblock{dgrb718}{Factorization by adjoining solutions}
\begin{dcomment}{exposition}{p1}{Factorization by adjoining solutions}
Grothendieck wants a construction which can be used uniformly for every arrow and subsequently compared as its generators vary. The object to be constructed is therefore more than an isolated factorization. It is a section of the composition functor
\[
 \mathcal M^{[2]}\longrightarrow\mathcal M^{[1]},
 \qquad (X\longrightarrow E\longrightarrow Y)\longmapsto(X\longrightarrow Y),
\]
where $[n]=\{0<\cdots<n\}$. This explains the apparently leisurely passage through the arrow category: naturality in a commutative square will eventually supply natural replacements and comparisons between them.

Write $I=C_0$. In modern lifting notation the two desired classes are
\[
 R=I^\perp,\qquad L={}^{\perp}R.
\]
Here $j\perp p$ means that every commutative square with $j$ on the left and $p$ on the right admits a diagonal filler; uniqueness is not required. Thus Grothendieck's $(C_0)_*$ is $R$, and $\widetilde{C_0}$ is $L$. The construction will first produce a much more concrete subclass of $L$, consisting of maps built by attaching copies of the generators. Retracts enter only when one identifies all of $L$.

There are two independent tasks. Adjoining a cell for every lifting problem makes all problems present at one stage solvable at the next. Smallness of the domains then ensures that a problem appearing at a sufficiently long limit was already present earlier. Neither task asserts that one application solves every problem in its newly enlarged domain. This distinction is the reason for the transfinite sequence.

The resulting argument is the classical small object argument, whose modern formulation and retract conclusion are given in \citet[Theorem~12.2.2 and Corollary~12.2.4]{Riehl2014}. Grothendieck isolates a useful economy in its hypotheses: arbitrary coproducts in $\mathcal M$ are unnecessary. Pushouts of the generating arrows, together with filtered colimits, suffice to form all the amalgamated sums under the fixed source. The word ``accessible'' concerns objects as tested by their covariant representable functors; in the proof below, only factorization of maps from the generating domains through stages is used. The categorical accessibility assumptions later in this chapter are a way of obtaining a uniform bound for that property.
\end{dcomment}


Voici les propriétés essentielles du système inductif ordinal des $X_\alpha(f)$ ou $\Phi_\alpha(f)$.

\msp{page 3}

\textbf{1.} $X_0(f) = X$, $i_0(f) = \mathrm{id}_X$, $p_0(f) = f$ (\textit{i.e.} la factorisation $\Phi_0(f)$ de $f$ est la factorisation triviale $X \xrightarrow{\mathrm{id}_X} X \xrightarrow{f} Y$).

\textbf{2.}
\[
\begin{aligned}
X_{\alpha+1}(f) &\;=\; X_1(p_\alpha(f)) \\
i_{\alpha+1}(f) &\;=\; i_1(p_\alpha(f)) \circ i_\alpha(f) \\
p_{\alpha+1}(f) &\;=\; p_1(p_\alpha(f))
\end{aligned}
\]
\[
\begin{tikzcd}
\underbrace{X}_{= \,\mathrm{source}\, f} \arrow[r, "i_\alpha(f)"] \arrow[rrr, bend left=20, "f"] & X_\alpha(f) \arrow[rr, "p_\alpha(f)"] \arrow[dr, bend right=40, "i_1(p_\alpha(f))"'] & & Y \\
& & \underbrace{X_1(p_\alpha(f))}_{= X_{\alpha+1}(f)} \arrow[ur, bend right=40] &
\end{tikzcd}
\]

\clearpage

% ===================== PDF page 5 =====================
\printedpage{3}
% head: A. Grothendieck : Dérivateurs, Ch. XV, Édition des Universités de Montpellier II et Paris VII
% foot: Version 5 novembre 2010

ou mieux
\[
\begin{tikzcd}[row sep=large]
X \arrow[rrr, "f"] \arrow[d, "\mathrm{id}"] & & & Y \arrow[d, "\mathrm{id}"] \\
X \arrow[r, "i_\alpha(f)"] \arrow[d, "\mathrm{id}"] & X_\alpha(f) \arrow[rr, "p_\alpha(f)"] \arrow[d, "\mathrm{id}"] & & Y \arrow[d, "\mathrm{id}"] \\
X \arrow[r] & X_\alpha(f) \arrow[r, "i_1(p_\alpha(f))"] & \underbrace{X_1(p_\alpha(f))}_{= X_{\alpha+1}(f)} \arrow[r, "p_1(p_\alpha(f))"] & Y \;.
\end{tikzcd}
\tag*{(1.12)}
\]
Le morphisme de transition $X_\alpha(f) \longrightarrow X_{\alpha+1}(f)$ est
\[
X_\alpha(f) \;\; \xrightarrow{\;\;i_1(p_\alpha(f))\;\;} \;\; X_1(p_\alpha(f)) \;=\; X_{\alpha+1}(f) \;.
\]

\textbf{3.} Si $\alpha$ est un ordinal limite, on a
\[
\left\{
\begin{array}{rcl}
X_\alpha(f) & = & \varinjlim\limits_{\alpha'<\alpha} X_{\alpha'}(f). \\[3ex]
i_\alpha(f) & = & \parbox[t]{8cm}{morphisme canonique de $X = X_0(f)$ dans la limite inductive.} \\[5ex]
p_\alpha(f) & = & \text{limite projective des } p_{\alpha'}(f) : X_{\alpha'}(f) \longrightarrow Y.
\end{array}
\right.
\]
(On peut considérer la limite inductive définissant $X_\alpha(f)$ comme une limite dans $X\backslash\M/Y\dots$)

Les morphismes de transition
\[
X_{\alpha'}(f) \;\longrightarrow\; X_\alpha(f) \qquad (\; \alpha' \;<\; \alpha \;)
\]
sont les morphismes canoniques relatifs au système inductif des $X_{\alpha'}(f)$.

\msp{page 4}

On peut donc dire que la suite ordinale des factorisations fonctorielles $\Phi_\alpha(f)$ de $f$ variable, se déduit de la factorisation fonctorielle $\Phi_1(f)$ de $f$ en
\[
\begin{tikzcd}
X \arrow[r, "i_1(f)"] \arrow[rr, bend left=20, "f"] & X_1(f) \arrow[r, "p_1(f)"] & Y
\end{tikzcd}
\]
par “itération transfinie à droite”, \textit{i.e.} en regardant les opérations $\Phi_\alpha$ comme des foncteurs
% footnote mark (1) printed beside the following display; \footnotemark/\footnotetext used because \footnote does not work inside math
\[
\left\{
\begin{array}{rcl}
p_\alpha :\; \uFl(\M) & \longrightarrow & \uFl(\M) \\
f & \longmapsto & p_\alpha(f) \;,
\end{array}
\right.
\qquad\qquad \text{\footnotemark}
\]
\footnotetext{\textbf{N.B.} $p_\alpha$ respecte les structures de catégories \textit{sur} $\M$, via les foncteurs but : $\uFl(\M) \longrightarrow \M$, \textit{i.e.} $\uFl(\M)$ est considérée comme objet de $\mathrm{CAT}/\M$.}
munis d'homomorphismes fonctoriels (au-dessus de $\M$)
\[
\mathrm{id}_{\uFl(\M)} \;\longrightarrow\; p_\alpha \;.
\]

% ---- from Der-ChXV-05-11-10_p006-010.tex
\printedpage{4}
% head: A. Grothendieck : Dérivateurs, Ch. XV, Édition des Universités de Montpellier II et Paris VII
% foot: Version 5 novembre 2010

Ainsi, si on considère
\[
\left\{
\begin{array}{rl}
\Flu(\cM) & \text{comme objet de CAT}/\cM \text{ par le foncteur } f \longmapsto \text{but}\, f\\[1ex]
\text{où} & \text{CAT est la catégorie des catégories \textit{éléments de} } \fV\text{\textit{,} univers}\\
 & \text{tel que } \fU \in \fV\text{, d'où Cat } \overset{\text{(déf)}}{=} \mathrm{Cat}^{\fU},\ \cM,\ \Flu(\cM) \in \text{CAT} ,
\end{array}
\right.
\]
on a un endofoncteur $p_1$ de $\Flu(\cM)$ sur $\cM$, et un homomorphisme fonctoriel
\[
i_1 \ : \ \mathrm{id}_{\Flu(\cM)/\cM} \longrightarrow p_1 \ , \qquad \textit{i.e.} \quad
\begin{tikzcd}[column sep=small]
X \arrow[rr, "f"] \arrow[dr, "i_1(f)"'] & & Y \\
 & \underbrace{X_1(f)}_{=\,\text{source}\, p_1(f)} \arrow[ur, "p_1(f)"'] &
\end{tikzcd} \ ,
\tag*{(1.13)}
\]
et la description de la suite ordinale des endofoncteurs $p_\alpha$ de $\Flu\,\cM/\cM$ résulte de cette seule donnée, compte tenu que les petites limites filtrantes existent dans les fibres de $\Flu\,\cM \xrightarrow{\text{but}} \cM$.

\msp{page 5}

\[
\begin{tikzcd}[column sep=normal]
X \arrow[r, "i_1(f)"] \arrow[dr, "f"'] & X_1(f) \arrow[r, "i_1(p_1(f))"] \arrow[d, "p_1(f)"] & X_2(f) \arrow[r] \arrow[dl, "p_2(f)"] & \cdots \arrow[r] & X_\alpha(f) \arrow[r, "i_1(p_\alpha(f))"] \arrow[dlll, "p_\alpha(f)"] & X_{\alpha+1}(f) \arrow[r] \arrow[dllll, "p_{\alpha+1}(f) = p_1(p_\alpha(f))"] & \cdots \\
 & Y & & & & &
\end{tikzcd}
\tag*{(1.14)}
\]

Donc, pour achever de définir la suite ordinale des $X_\alpha(f)$ (munie des $i_\alpha(f)$, $p_\alpha(f)$), ou des $p_\alpha(f)$ (munie des $i_\alpha(f)$), il reste à décrire $X_1(f)$, $i_1(f)$, $p_1(f)$.

La propriété essentielle de ce triple dont nous aurons besoin est la suivante :

\begin{itemize}[leftmargin=5em, labelwidth=4em, labelsep=1em, align=left]
\item[(1.15)] \textbf{Fact}$(C_0)$ (propriété de factorisation par rapport à $C_0$).

Pour tout diagramme commutatif
\[
D : \qquad
\begin{tikzcd}
A \arrow[r] \arrow[d, "i"'] & X \arrow[d, "f"] \\
B \arrow[r] & Y
\end{tikzcd}
\qquad \text{avec } i \in C_0 \ ,
\]
d'où un diagramme traits pleins
\[
\begin{tikzcd}
A \arrow[r] \arrow[d, "i"'] & X \arrow[d, "i_D"] \arrow[dr, dotted, "i_1(f)"] \arrow[dd, dotted, bend right=60, "f"'] & \\
B \arrow[r] & X_D \overset{\text{déf}}{=} X \amalg_A B \arrow[r, dotted, "u_D"] \arrow[d, "p_D"] & X_1(f) \arrow[dl, dotted, "p_1(f)"] \\
 & Y \ , &
\end{tikzcd}
\]
\printedpage{5}
% head: A. Grothendieck : Dérivateurs, Ch. XV, Édition des Universités de Montpellier II et Paris VII
% foot: Version 5 novembre 2010
il existe un homomorphisme de la factorisation $f = p_D i_D$ de $f$ dans la factorisation $f = p_1(f)i_1(f)$.
\end{itemize}

Pour formuler cette condition, j'ai utilisé l'existence des sommes amalgamées $X_D$, \textit{i.e.} la condition sur $C_0$ :

\msp{page 6}

\dgrblock{dgrb719}{A single stage and the colimits it needs}
\begin{dcomment}{exposition}{p5}{A single stage and the colimits it needs}
The first stage is arranged to solve every prescribed extension problem without choosing one problem in preference to another. For $f:X\to Y$, its indexing set can be written explicitly as
\[
 \Delta(I,f)=\coprod_{(j:A\to B)\in I}
 \{(a,b)\in\operatorname{Hom}(A,X)\times
                   \operatorname{Hom}(B,Y)\mid fa=bj\}.
\]
This is a small set: $I$ is small and the category is locally small in the chosen universe. Each triple $D=(j,a,b)$ gives the pushout $X_D=X\amalg_A B$, together with its induced map to $Y$. In the category of objects under $X$, take their coproduct. This gives exactly (1.23), and the inclusion of the $D$th cell supplies the map $B\to X_1(f)$ required by (1.17).

The existence argument deserves to be explicit because the stated assumptions do not include all small colimits. A pushout of a pushout of $j$ is again a pushout of $j$: this is the pasting property of pushout squares. Consequently every map $X\to X_D$ is itself pushoutable along every map out of $X$. Any finite family of these objects therefore has an amalgamated sum under $X$, built by repeated binary pushouts. For a finite subset $F\subseteq\Delta(I,f)$, write this sum as $X_F$, and set $X_{\varnothing}=X$. Chosen colimit cocones give transition maps for inclusions $F\subseteq G$. The poset of finite subsets is filtered, and
\[
 X_1(f)=\mathop{\mathrm{colim}}_{F\subseteq\Delta(I,f),\;F\text{ finite}} X_F.
\]
To verify its universal property, a map from this colimit to an object under $X$ is a compatible family of maps from the finite amalgamations. Restriction to singleton subsets identifies such a family with precisely one compatible map from each $X_D$. This proves the required coproduct property in the undercategory, including the case $\Delta(I,f)=\varnothing$.

If ordinary coproducts also exist, the same object is the pushout
\[
\begin{tikzcd}
 \coprod_{D\in\Delta(I,f)} A_D \arrow[r]\arrow[d] & X\arrow[d,"i_1(f)"]\\
 \coprod_{D\in\Delta(I,f)} B_D \arrow[r] & X_1(f).
\end{tikzcd}
\]
This familiar presentation is convenient for calculations; the preceding finite-subset construction explains why it is not an additional assumption in Theorem~1. The construction of a stage uses the codomains $B_D$, but the eventual lifting argument will require smallness only of the domains $A_D$.
\end{dcomment}


\begin{itemize}[leftmargin=5em, labelwidth=4em, labelsep=1em, align=left]
\item[(1.16)] Les flèches $i \in C_0$ sont coquarrables, \textit{i.e.} pour un diagramme
\[
\begin{tikzcd}
A \arrow[r, "u"] \arrow[d, "i"'] & X \\
B &
\end{tikzcd}
\]
avec $i \in C_0$, la somme amalgamée existe.
\end{itemize}

Mais on peut formuler 1.15 aussi ainsi, indépendamment de l'hypothèse (1.16) :

\begin{itemize}[leftmargin=5em, labelwidth=4em, labelsep=1em, align=left]
\item[(1.17)] Pour tout carré commutatif
\[
\begin{tikzcd}
A \arrow[r, "u"] \arrow[d, "i"'] & X \arrow[d, "f"] \\
B \arrow[r, "v"] & Y
\end{tikzcd}
\qquad ,
\]
avec $i \in C_0$, il existe un carré commutatif traits pleins
\[
\begin{tikzcd}
A \arrow[r, "u"] \arrow[d, "i"'] & X \arrow[d, "i_1(f)"] \\
B \arrow[r, "v_1"] \arrow[dr, dotted, "v"] & X_1(f) \arrow[d, dotted, "p_1(f)"] \\
 & Y
\end{tikzcd}
\]
tel que $v = p_1(f)v_1$.
\end{itemize}

Ceci a pour conséquence le

\textbf{Corollaire 1} (de 1.17). \textit{Soit} $f \in \Fl(\cM)$\textit{, et} $\alpha$ \textit{un ordinal} $(\in \fU)$. \textit{Pour tout carré commutatif}
\[
\begin{tikzcd}
A \arrow[r, "u_\alpha"] \arrow[d, "i"'] & X_\alpha(f) \arrow[d, "p_\alpha(f)"] \\
B \arrow[r, "v"] & Y
\end{tikzcd}
\]
\printedpage{6}
% head: A. Grothendieck : Dérivateurs, Ch. XV, Édition des Universités de Montpellier II et Paris VII
% foot: Version 5 novembre 2010
\textit{il existe un carré commutatif}
\[
\begin{tikzcd}
A \arrow[r, "u_\alpha"] \arrow[d, "i"'] & X_\alpha(f) \arrow[d, "\text{trans.}"] \arrow[dd, dotted, bend left=60, "p_\alpha(f)"] \\
B \arrow[r, "v_1"] \arrow[dr, dotted, "v"] & X_{\alpha+1}(f) \arrow[d, dotted, "p_{\alpha+1}(f)"] \\
 & Y
\end{tikzcd}
\]
\textit{tel que}
\[
v \ = \ p_{\alpha+1}(f)\ v_1 \ .
\]

\msp{page 7}

\textbf{Corollaire 2.} \textit{Soit $\pi$ un cardinal} $(\in \fU)$ \textit{tel que pour toute flèche $i : A \longrightarrow B$ dans $C_0$, la source $A$ soit $\pi$-accessible. Soit $\alpha$ un ordinal grand devant $\pi$} (i.e. \textit{tel que l'ensemble ordonné $I_\alpha$ des ordinaux $\alpha' < \alpha$ soit grand devant $\pi$}) -- \textit{p. ex. $\alpha$ le premier ordinal dont le cardinal soit $> \pi$. Alors pour toute flèche $f \in \Fl(\cM)$, on a $p_\alpha(f) \in (C_0)_*$.}

\textsc{Démonstration}. Cela signifie que pour tout carré commutatif
\[
\begin{tikzcd}
A \arrow[r, "u"] \arrow[d, "i"'] & X_\alpha(f) \arrow[d, "p_\alpha(f)"] \\
B \arrow[r, "v"] \arrow[ur, dotted, "w"] & Y \ ,
\end{tikzcd}
\tag*{$(*)$}
\]
avec $i \in C_0$, la flèche pointillée $w$ (avec $wi = u$, $p_\alpha(f)w = v$) existe. Or
\[
X_\alpha(f) = \varinjlim_{\alpha' \in I_\alpha} X_{\alpha'}(f) \ ,
\]
donc $A$ étant $\pi$-accessible et $I_\alpha$ grand devant $\pi$, $u$ provient d'une flèche $u_{\alpha'} : A \longrightarrow X_{\alpha'}(f)$, et (quitte à remplacer $\alpha'$ par $\alpha'' > \alpha'$ et $u_{\alpha'}$ par le composé $u_{\alpha''} : A \longrightarrow X_{\alpha'}(f) \longrightarrow X_{\alpha''}(f)$) on peut supposer $p_{\alpha'}(f)u_{\alpha'} = vi$, \textit{i.e.} que $(*)$ provient d'un carré commutatif
\[
\begin{tikzcd}
A \arrow[r, "u_{\alpha'}"] \arrow[d, "i"'] \arrow[rr, dotted, bend left=30, "u"] & X_{\alpha'}(f) \arrow[r, dotted, "\text{trans.}"'] \arrow[d, "p_{\alpha'}(f)"'] & X_\alpha(f) \arrow[dl, dotted, "p_\alpha(f)"] \\
B \arrow[r, "v"] & Y \ , &
\end{tikzcd}
\]

\printedpage{7}
% head: A. Grothendieck : Dérivateurs, Ch. XV, Édition des Universités de Montpellier II et Paris VII
% foot: Version 5 novembre 2010
et appliquant le corollaire 1, on trouve
\[
\begin{tikzcd}
A \arrow[r, "u_{\alpha'}"] \arrow[d, "i"'] & X_{\alpha'}(f) \arrow[d, "\text{trans.}"] & \\
B \arrow[r, "v_1"] \arrow[dr, dotted, "v"] & X_{\alpha'+1}(f) \arrow[r, dotted, "\text{trans.}"] \arrow[d, "p_{\alpha+1}(f)"] & X_\alpha(f) \arrow[dl, dotted, "p_\alpha(f)"] \\ % [sic] p_{\alpha+1}(f) -> p_{\alpha'+1}(f), the arrow starts at X_{\alpha'+1}(f); ?? uncertain: subscript of the label, tightly set -- \alpha+1 or \alpha'+1
 & Y &
\end{tikzcd}
\]
avec $v = p_{\alpha+1}(f)v_1$, et on prendra % [sic] p_{\alpha+1}(f)v_1 -> p_{\alpha'+1}(f)v_1, corollaire 1 is applied to \alpha' (arrow from X_{\alpha'+1}(f))
\[
w \ = \ \mathrm{trans}_{\alpha,\alpha'+1} \circ v_1 \ .
\]

\msp{page 8}

\dgrblock{dgrb720}{Why the cofinality bound solves every problem}
\begin{dcomment}{exposition}{p7}{Why the cofinality bound solves every problem}
Grothendieck's Corollary~2 supplies the decisive passage from a succession of partial solutions to an actual right lifting property. The statement needs the additional word \emph{limit}. For a limit ordinal, the condition that every subset of cardinal at most $\pi$ has an upper bound is equivalent to cofinality greater than $\pi$. More invariantly, choose a regular cardinal $\kappa$ for which every generating domain is $\kappa$-presentable, and use a limit ordinal $\lambda$ with $\operatorname{cf}(\lambda)\geq\kappa$. The ordered set of stages $\beta<\lambda$ is then $\kappa$-filtered.

Consider a square $j:A\to B$ against $p_\lambda:X_\lambda\to Y$, with top map $a:A\to X_\lambda$ and bottom map $b:B\to Y$. Since
\[
 X_\lambda=\mathop{\mathrm{colim}}_{\beta<\lambda}X_\beta,
\]
smallness gives $\beta<\lambda$ and $a_\beta:A\to X_\beta$ representing $a$. The compatibility of the maps to $Y$ gives
\[
 p_\beta a_\beta=p_\lambda a=bj.
\]
Thus the square already commutes at this stage. The parenthetical suggestion to enlarge $\beta$ is harmless but unnecessary: the target $Y$ has remained fixed throughout. The corresponding cell at stage $\beta+1$ provides $h:B\to X_{\beta+1}$ with $hj=t_{\beta+1,\beta}a_\beta$ and $p_{\beta+1}h=b$. Because $\lambda$ is a limit, $\beta+1<\lambda$, and $t_{\lambda,\beta+1}h$ is the desired filler.

No smallness assertion about $B$ was used. In fact the proof needs only the surjectivity of the map from the colimit of the relevant Hom sets, rather than its injectivity. This is why the usual relative-smallness hypothesis, restricted to the cellular chains actually constructed, is sufficient as well. The bound must be common to all generating domains, and the stopping ordinal must be common to all input arrows if one wants a single factorization functor. Since the generators form a set, individual accessibility bounds can be replaced by a common larger one.

The omitted limit condition is a real, although easily repaired, boundary error. In the definition used in \citet[Expos\'e~I, Definition~9.1]{SGA4I1972}, a poset with a maximum is ``grand devant $\pi$'' for every $\pi$. Thus $I_1=\{0\}$ satisfies the printed hypothesis. Take sets, generators $\varnothing\to 1$ and the fold $2\to 1$, and $f=1_1$. The first generator adds one point; the square for the fold identifies no distinct old points. Hence $X_1(f)=2$ and $p_1(f):2\to1$ is not injective. Since the right class for these two generators consists of bijections, Corollary~2 would be false at $\alpha=1$. Its proof itself shows the repair: choose a limit ordinal with the indicated cofinality, as in the successor-cardinal example given in the text. The same qualification is required in (1.22) and Theorem~1(b).

The two printed labels $p_{\alpha+1}$ at the end of the proof, on printed page~7, omit the prime on $\alpha$. Their intended value is $p_{\alpha'+1}$. The displayed source of that arrow determines the correction; the edition's readings are retained in the French text and registered in the corrections table.
\end{dcomment}


Mais on veut une factorisation de $f$, flèche quelconque dans $\cM$, en
\[
f \ = \ p\ i \ , \qquad \text{avec} \ \ p \in (C_0)_* \ \ \textit{et} \ \ i \in \widetilde{C_0} \ ,
\]
donc il faut encore énoncer une propriété, assurant que les $i_\alpha(f)$ sont dans $\widetilde{C_0}$. On aura même quelque chose de plus précis. Introduisons

\begin{itemize}[leftmargin=10em, labelwidth=9em, labelsep=1em, align=left]
\item[(1.18)\hfill $C'_0 \ = $] plus petit sous-ensemble de $\Fl(\cM)$ contenant $C_0$, stable par cochangement de base (chaque fois qu'il existe).
\item[\hfill $C''_0 \ = $] ensemble des flèches qui sont soit dans $C'_0$, soit des sommes amalgamées d'une famille \textit{non vide} de telles flèches $A \xrightarrow{i_\alpha} B_\alpha$ de même source $A$, $i_\alpha \in C'_0$ ; la flèche $A \longrightarrow B = \displaystyle\coprod_{\alpha}{}_{(A)} B_\alpha$ est dans $C''_0$.
\end{itemize}

(La condition 1.16 implique que les éléments de $C'_0$ sont coquarrables, et cela suffit, avec la condition de stabilité de $\cM$ par $\varinjlim$ filtrantes, à assurer que les sommes amalgamées infinies envisagées existent dans $\cM$.)

\textbf{N.B.} $C''_0$ est \textit{égal} au plus petit ensemble de flèches contenant $C_0$ et stable par cochangement de base et par petites sommes amalgamées. \texttt{[Si les petites sommes existent dans $\cM$, alors $C''_0$ est aussi l'ensemble des flèches obtenues par cochangement de base à partir de sommes de flèches de $C_0$.]}

\setcounter{footnote}{1}% the edition numbers these footnotes 2 and 3
\begin{itemize}[leftmargin=10em, labelwidth=9em, labelsep=1em, align=left]
\item[(1.19)\hfill $C'''_0 \ = $] plus petit sous-ensemble de $\Fl(\cM)$ contenant $C''_0$, et stable de plus par limites inductives ordinales strictes, \textit{i.e.} pour un système inductif ordinal $(A_\alpha)_{\alpha<\alpha_0}$, tel que pour $\alpha$ ordinal limite on ait $A_\alpha \xleftarrow{\ \sim\ } \varinjlim\limits_{\alpha' \in I_\alpha} A_{\alpha'}$, si les morphismes de transition $i_{\beta,\alpha} : A_\alpha \longrightarrow A_\beta$ sont dans $C'''_0$, de même les morphismes canoniques
\[
A_\alpha \xrightarrow{\ i_\alpha\ } \varinjlim_{\beta \in I_{\alpha_0}} A_\beta \ . \quad \footnotemark
\]
\end{itemize}
\footnotetext{$C_0 \subset C'_0 \subset C''_0 \subset C'''_0$}% footnote 2: mark printed in the display, text printed at the foot of p. 8

\printedpage{8}
% head: A. Grothendieck : Dérivateurs, Ch. XV, Édition des Universités de Montpellier II et Paris VII
% foot: Version 5 novembre 2010
\msp{page 9}

\dgrblock{dgrb721}{Empty families and a complete calculation in sets}
\begin{dcomment}{examples}{p8}{Empty families and a complete calculation in sets}
Grothendieck has now separated attachment from lifting, and the simplest generator makes both operations visible. Work in sets and let $I=\{\varnothing\to\{*\}\}$. For $f:X\to Y$, a lifting problem is exactly an element of $Y$. Hence
\[
 X_1(f)=X\amalg Y,\qquad i_1(f)=\mathrm{in}_X,
 \qquad p_1(f)=[f,1_Y].
\]
The right class is the class of surjections, because a lift against $\varnothing\to\{*\}$ is a preimage of one specified element. With the usual axiom of choice, the left class is the class of injections. An injection lifts against a surjection by extending the prescribed map over its complement. Conversely, if $g:A\to B$ is not injective, choose distinct $a,a'$ with $g(a)=g(a')$, map $A$ to $\{0,1\}$ with different values at these points, and use the surjection $\{0,1\}\to\{*\}$. A filler would factor that map through $g$, which is impossible.

The whole transfinite construction can be calculated, with the obvious compatible choices of coproducts:
\[
 X_\alpha(f)=X\amalg(Y\times\alpha),\qquad
 p_\alpha(x)=f(x),\qquad p_\alpha(y,\beta)=y.
\]
At a successor one adds another copy of $Y$; at a limit one takes the union of the earlier copies. For every $\alpha>0$, $p_\alpha$ is already a surjection, but if $Y\ne\varnothing$ every transition to the next stage adds elements. Success in the lifting problem therefore does not mean that the unmodified construction has stabilized.

There is also a small literal omission in (1.18)--(1.20). The definition of $C_0''$ expressly allows only \emph{nonempty} amalgamated families. In this example $C_0'$ consists of the injections which add one element, and $C_0''$ consists of the injections with nonempty complement. For $f:\varnothing\to\varnothing$, however, $\Delta(I,f)$ is empty and (1.23) gives $i_1(f)=1_{\varnothing}$. Thus (1.20), read with the printed definition, is false in this boundary case. The minimal repair is to permit the empty amalgamated family, whose structural map is the identity, or explicitly to adjoin identities to $C_0''$. Under the printed convention the always valid assertion is $i_1(f)\in C_0''\cup\mathrm{Iso}(\mathcal M)$. This is a convention repair in the source, not a transcription correction, and it does not affect the resulting weak factorization theorem. Zero-length cellular compositions must in any event include identities.
\end{dcomment}


La proposition 2, p. 4 de XIII assure que l'on a
\[
C'''_0 \ \subset \ \widetilde{C_0} \ .
\]
Ceci posé, voici la deuxième propriété essentielle de la factorisation $f = p_1(f)i_1(f)$ (la première étant 1.17) :

\begin{itemize}[leftmargin=5em, labelwidth=4em, labelsep=1em, align=left]
\item[(1.20)] \qquad\qquad Pour toute flèche $f \in \Fl(\cM)$, on a $i_1(f) \in C''_0$.
\end{itemize}

Il en résulte aussitôt, par récurrence ordinale :

\dgrblock{dgrb722}{The cellular closure and its successive steps}
\begin{dcomment}{exposition}{p9:Il en résulte aussitôt, par récurrence ordinale}{The cellular closure and its successive steps}
The finer class $C_0'''$ records how a map was constructed, before retracts erase that information. Its intended modern description is the class of relative $I$-cell maps: continuous transfinite composites in which every successor map is a pushout of an amalgamated family of copies of generators. Identities are allowed, and all diagrams are small in the chosen universe. The separate stages $C_0'$ and $C_0''$ distinguish one attachment from simultaneous attachment of a family.

The assertion that these maps lie in ${}^{\perp}(I^\perp)$ has a direct proof. Fix $q:U\to V$ in $I^\perp$. A lifting problem for a pushout of $j\in I$ reduces, by the pushout property, to one for $j$. For an amalgamated family of such maps, choose a filler for each member with the common prescribed value on the source; the resulting maps agree there and therefore induce a filler on the amalgamation. For a continuous chain $(Z_\beta)_{\beta\leq\lambda}$, construct the lift on $Z_\beta$ by transfinite recursion. At a successor extend the lift already constructed using the lifting property of $Z_\beta\to Z_{\beta+1}$. At a limit the compatible lifts determine a map from the colimit. This proves the lifting property for $Z_0\to Z_\lambda$. It uses existence of fillers, together with the ordinary choice principle, and no uniqueness.

The word ``successor'' is essential in this formulation. In the explanatory sentence of (1.19), the displayed symbols $i_{\beta,\alpha}$ can be read as demanding that \emph{every} transition already belong to the class. That condition by itself would not imply closure under finite composition. For example, in the poset category $0<1<2$, take $I=\{0\to1,1\to2\}$. Pushouts and amalgamations give these two arrows and identities, but not $0\to2$. A chain all of whose transitions lie in this class cannot contain both $0$ and $2$; its colimit introduces no missing arrow. Nevertheless two successive cell attachments factor $0\to2$ through $1$. Thus the induction in (1.21) requires the usual successorwise interpretation of ``limites inductives ordinales strictes'', or an explicit additional closure under composition. We adopt that intended interpretation throughout; the claim about the lifting class is valid independently of this notational compression.
\end{dcomment}


\begin{itemize}[leftmargin=5em, labelwidth=4em, labelsep=1em, align=left]
\item[(1.21)] \qquad Pour toute flèche $f \in \Fl(\cM)$ et tout ordinal $\alpha \in \fU$, on a $i_\alpha(f) \in C'''_0$, a fortiori $i_\alpha(f) \in \widetilde{C_0}$.
\end{itemize}

Donc si $\alpha$ est grand devant $\pi$ (\textit{cf.} cor. 2, page 7), la factorisation $f = p_\alpha(f)i_\alpha(f)$ satisfait
\[
\left\{
\begin{array}{l}
f = p_\alpha(f)i_\alpha(f)\text{, avec } i_\alpha(f) \in C'''_0 \subset \widetilde{C_0} = \Phi \text{ \texttt{[à ne pas}}\\
\text{\texttt{confondre avec }}\Phi_\alpha\text{\texttt{]}, } p_\alpha(f) \in C_{0*} = \Psi \text{, lorsque } \alpha \text{ est}\\
\text{grand devant } \pi\text{, les sources des } i \in C_0 \ \pi\text{-accessibles.}
\end{array}
\right.
\tag*{(1.22)}
\]

Il reste à construire le foncteur $f \longmapsto X_1(f)$, avec $i_1(f)$, $p_1(f)$, satisfaisant les conditions (1.17) et (1.20), sous l'hypothèse (1.16) que les $i \in C_0$ sont coquarrables. On considère alors, pour $f \in \Fl(\cM)$ fixé, l'ensemble $\Delta(f)$ (ou $\Delta(C_0, f)$) de \textit{tous} les diagrammes $D$ de la forme (1.15), avec $i \in C_0$, et on pose
\[
X_1(f) \ = \ \coprod_{D \in \Delta(f)}{}_{(X)} X_D
\tag*{(1.23)}
\]
(somme amalgamée sous $X$, \textit{i.e.} somme indexée
\msp{page 10}
par $\Delta(f)$ dans $X\backslash\cM$), le morphisme $X \xrightarrow{i_1(f)} X_1(f)$ étant le morphisme costructural, et le morphisme $X_1(f) \xrightarrow{p_1(f)} Y$ étant celui ayant pour composantes les $p_D$. En somme, on fait “au plus économique” (et au plus naturel) pour assurer la validité de (1.17), sous la forme équivalente (1.15). D'autre part, par définition de $C'_0$, $C''_0$, les $i_D$ sont dans $C'_0$, et $i_1(f)$ est dans $C''_0$, ce qui est (1.20)\footnote{Il faudrait quand même expliciter la dépendance fonctorielle de $X_1(f)$ par rapport à $f$\,!!!}.

\dgrblock{dgrb723}{The omitted functoriality in the input arrow}
\begin{dcomment}{exposition}{p10:Il faudrait quand même expliciter la dépendance fonctorielle}{The omitted functoriality in the input arrow}
The footnote asks for the functorial dependence which makes Theorem~1 usable. Its proof is short once the indexing squares are retained as data. A morphism $(u,v):f\to f'$ in the arrow category consists of $u:X\to X'$ and $v:Y\to Y'$ with $vf=f'u$. It sends
\[
 (j,a,b)\longmapsto(j,ua,vb)
\]
to a well-defined member of $\Delta(I,f')$. For this image $D'$, the pair of maps $X\xrightarrow{u}X'$ and $1_B$ induces a map
\[
 X_D=X\amalg_A B\longrightarrow X'\amalg_A B=X'_{D'}.
\]
These maps respect the structural maps from $X$ and the maps to the target. Composing them with the inclusions into $X_1(f')$ gives a cocone over the family defining $X_1(f)$, hence a unique map $X_1(u,v):X_1(f)\to X_1(f')$.

Identity and composition laws follow by restricting to $X$ and to each copy of $B$: there the asserted equalities are the identity and composition laws for $u$ and $v$, and the universal property makes their extensions unique. The same argument proves the naturality of $i_1$ and $p_1$. At successor ordinals one applies this functor to the already constructed square between $p_\alpha(f)$ and $p_\alpha(f')$; at limits one takes the map of colimits. Transfinite induction gives all of (1.10)--(1.14), including compatibility with every transition map.

One must fix the relevant colimit cocones to obtain literal functors. Another choice gives canonically isomorphic functors by the same universal properties. There is no need to choose fillers in arbitrary lifting problems: the attached cell itself supplies the particular solution required at each stage. This distinction explains why the factorization is functorial even when the fillers subsequently used to characterize the left or right class are highly nonunique.
\end{dcomment}


Pour résumer :

\textbf{Théorème 1.} \textit{Soient $\cM$ une $\fU$-catégorie stable par $\fU$-limites inductives filtrantes, $C_0 \subset \Fl(\cM)$ un petit ensemble de flèches de $\cM$. On suppose les flèches $i \in C_0$ coquarrables. Cela assure la définition d'un endofoncteur $f \longmapsto p_1(f)$ de $\Flu(\cM)$, et d'un homomorphisme fonctoriel $i_1 : \mathrm{id}_{\Flu(\cM)} \longrightarrow p_1$,} i.e. \textit{$i_1(f) : f \longrightarrow p_1(f)$, induisant l'identité sur les objets buts de $f$ et de $p_1(f)$,}
\[
\begin{tikzcd}[column sep=small]
X \arrow[rr, "i_1(f)"] \arrow[dr, "f"'] & & X_1(f) \arrow[dl, "p_1(f)"] \\
 & Y \ , &
\end{tikzcd}
\]

% ---- from Der-ChXV-05-11-10_p011-015.tex
\printedpage{9}
% head: A. Grothendieck : Dérivateurs, Ch. XV, Édition des Universités de Montpellier II et Paris VII
% foot: Version 5 novembre 2010

{\itshape d'où par itération un système inductif ordinal d'endofoncteurs $f\longmapsto p_\alpha(f)$ de $\uFl(\cM)$, et d'homomorphismes $i_\alpha : \mathrm{id}_{\uFl(\cM)} \longrightarrow p_\alpha$, \textup{i.e.}}
\[
\begin{tikzcd}[column sep=small]
X \arrow[rr, "i_\alpha(f)"] \arrow[dr, "f"'] & & X_\alpha(f) \arrow[dl, "p_\alpha(f)"] \\
& Y\,, &
\end{tikzcd}
\]
{\itshape formant un système inductif ordinal de factorisations de $f$ fonctorielles en $f$.}

\msp{page 11}

{\itshape Ceci posé, on a :}
\begin{enumerate}[label={}, leftmargin=*]
\item[a)] {\itshape Pour tout ordinal $\alpha$, on a $i_\alpha(f)\in\Phi$.}
\item[b)] {\itshape Supposons que les sources des $i\in C_0$ soient accessibles et soit alors $\pi$ un cardinal dans $\mathfrak{U}$ tel qu'elles soient $\pi$-accessibles. Alors pour tout nombre ordinal $\alpha$ grand devant $\pi$, on a $p_\alpha(f)\in C_{0*}$, donc}
\[
f \;=\; p_\alpha(f)\, i_\alpha(f) \qquad\quad \textit{avec} \qquad\quad i_\alpha(f)\in\Phi\,, \qquad\quad p_\alpha(f)\in\Psi\,.
\]
\end{enumerate}

On a un résultat légèrement plus précis, que voici :

\textbf{Corollaire 1.} \textit{Soit $\Phi_0$ une partie de} $\Fl(\cM)$ \textit{telle que}
\[
C_0''' \;\subset\; \Phi_0 \;\subset\; \Phi \qquad\quad (\textit{d'où}\ \ C_0\subset\Phi_0\subset\Phi)\,,
\tag*{(1.24)}
\]
\textit{où $C_0'''\subset\Phi$ est définie dans} 1.19, \textit{p.} 8. \textit{Par exemple, si $\Phi_0$ est une partie de} $\Fl(\cM)$ \textit{satisfaisant}
\[
C_0 \;\subset\; \Phi_0 \;\subset\; \Phi\,,
\]
\textit{et stable par cochangement de base, par petites sommes amalgamées} (\textit{cf.} (1.18), p.~8), \textit{par $\varinjlim$ ordinales strictes} (\textit{cf.} (1.19), p.~8), \textit{elle satisfait} (1.24). (\textbf{N.B.} \textit{Si $\cM$ admet des sommes amalgamées simples, de tels $\Phi_0$ existent trivialement, il suffit de prendre $\Phi_0=\Phi$.})

\msp{page 12}

\textit{Ceci posé, pour toute flèche $f$ dans $\cM$, et tout ordinal $\alpha$, on a}
\[
i_\alpha(f) \;\in\; \Phi_0\,.
\]

\textbf{Corollaire 2.} \textit{Supposons que les sources des flèches $i\in C_0$ soient accessibles, et soit $\Phi_0$ comme dans le corollaire} 1 (p.~ex. $\Phi_0=C_0'''$). \textit{Alors le couple $(\Phi_0,\Psi)$ est un précouple de Quillen,} i.e.
\begin{enumerate}[label={}, leftmargin=*]
\item[a)] $\Phi_0\longleftrightarrow\Psi$, i.e. $\Psi\subset\Phi_{0*}$ (\textit{ou ce qui revient au même, $\Phi_0\subset\Psi^*$}), \textit{et}
\item[b)] \textit{on a la propriété de factorisation de Quillen.}
\end{enumerate}
\textit{Cela implique donc que $\Phi$ $(=\widetilde{C_0}=\widetilde{\Phi_0})$ est formé des flèches qui sont facteurs directs de flèches de $\Phi_0$. Donc si $\Phi_0$ est stable par facteurs directs, on a $\Phi_0=\Phi$,} i.e. \textit{$\Phi_0$ est clos à gauche.}

\dgrblock{dgrb724}{Why the retract conclusion is exact}
\begin{dcomment}{exposition}{p12:Donc si $\Phi_0$ est stable par facteurs directs}{Why the retract conclusion is exact}
The corollary identifies the amount of closure required after the cellular construction. Put $R=I^\perp$ and $L={}^{\perp}R$. If $I\subseteq\Phi_0\subseteq L$, then $\Phi_0^\perp=R$: the inclusion $I\subseteq\Phi_0$ gives one containment, and $\Phi_0\subseteq{}^{\perp}R$ gives the other. Assume, as in the corollary, that every map factors as a map in $\Phi_0$ followed by a map in $R$.

Take $g:A\to B$ in $L$ and choose a factorization
\[
 A\xrightarrow{j}E\xrightarrow{q}B,
 \qquad j\in\Phi_0,\quad q\in R,\quad qj=g.
\]
The square with left map $g$, top map $j$, right map $q$ and bottom map $1_B$ admits a filler $s:B\to E$. It satisfies $sg=j$ and $qs=1_B$. In the arrow category the maps
\[
 g\xrightarrow{(1_A,s)}j\xrightarrow{(1_A,q)}g
\]
compose to the identity. Thus $g$ is a retract of a member of $\Phi_0$, with the source held fixed. Conversely a retract of a map with the left lifting property has that property: insert the retract maps into a lifting square, solve the resulting square, and compose back. Therefore $L$ is exactly the retract closure of $\Phi_0$.

The cellular class and its retract closure consequently have different jobs. Cellular maps are the output of the explicit construction; the retract closure is the class characterized purely by lifting. In particular, if $\Phi_0$ is retract closed then it is already $L$. The stronger description using $C_0'''$ follows from the transfinite attachment construction, with identities included as explained above. This is the content, rather than an additional hypothesis, of the assertion that $(\Phi_0,R)$ is a Quillen precouple. Section~4 will isolate the same retract mechanism without any reference to the particular generating set.
\end{dcomment}


\dgrblock{dgrb725}{From a pointed endofunctor to algebraic factorization}
\begin{dcomment}{later}{p12:Donc si $\Phi_0$ est stable par facteurs directs}{From a pointed endofunctor to algebraic factorization}
Grothendieck's emphasis on $1\to p_1$ makes the construction especially close to the later algebraic treatment. It should not, however, be read as already providing a monad: a pointed endofunctor comes with a unit but has no specified multiplication or associativity law. The set example above exhibits the difference concretely, since the ordinary iteration keeps adding solutions to problems which have already been solved.

\citet[Theorem~4.4]{Garner2009} constructs an algebraically free natural weak factorization system on a small category of generating arrows under specified cocompleteness and smallness hypotheses, in particular for a locally presentable category. Its right factor carries a monad and its left factor a comonad. For a discrete generating category, the algebras on a right map encode chosen solutions of the generating lifting problems; morphisms preserve those choices. The added algebraic structure gives a universal property to the construction. It is stronger than the existence of the functorial factorization obtained here.

This places the chapter's two concerns in their subsequent setting. The transfinite argument proves that the required lifting classes form a weak factorization system, while the algebraic refinement controls the choices of solutions and their compatibility. A functorial section of composition, an algebraic weak factorization system, and an orthogonal factorization system with unique fillers are three distinct structures. Nothing in Theorem~1 requires unique fillers.
\end{dcomment}


\printedpage{10}
% head: A. Grothendieck : Dérivateurs, Ch. XV, Édition des Universités de Montpellier II et Paris VII
% foot: Version 5 novembre 2010

\section*{2\quad Variation des factorisations $f=p_\alpha i_\alpha$ avec la partie $C_0$ génératrice}

Je dis que le système inductif ordinal de foncteurs $p_\alpha=p_\alpha^{C_0} : \uFl(\cM)\longrightarrow\uFl(\cM)$,
\msp{page 13}
muni des flèches
\[
i_\alpha^{C_0} \;:\; \mathrm{id}_{\uFl(\cM)} \longrightarrow p_\alpha^{C_0}\,,
\]
dépend fonctoriellement de $C_0$, quand on considère $C_0$ comme un élément de l'ensemble ordonné des \textit{petites} parties de $\Fl(\cM)$, formées de flèches coquarrables (donc de l'ensemble ordonné
\[
\mathfrak{P}_{\mathfrak{U}\text{-petites}}(\Fl_{\text{coquar.}}(\cM)) \quad ).
\]
Si on a une inclusion
\[
C_0 \;\hookrightarrow\; C_1\,,
\tag*{(2.1)}
\]
on en déduit pour tout ordinal $\alpha$ un homomorphisme fonctoriel
\[
u_\alpha \;=\; u_\alpha^{C_1,C_0} \;:\; X_\alpha^{C_0} \;\longrightarrow\; X_\alpha^{C_1}\,,
\tag*{(2.2)}
\]
rendant commutatifs les carrés déduits de $f : X\longrightarrow Y$ dans $\cM$,
\[
\begin{tikzcd}
X \arrow[r, "i_\alpha^{C_0}(f)"] \arrow[d, "\mathrm{id}"'] & X_\alpha^{C_0}(f) \arrow[r, "p_\alpha^{C_0}(f)"] \arrow[d, "u_\alpha(f)"] & Y \arrow[d, "\mathrm{id}"] \\
X \arrow[r, "i_\alpha^{C_1}(f)"] & X_\alpha^{C_1}(f) \arrow[r, "p_\alpha^{C_1}(f)"] & Y\,,
\end{tikzcd}
\tag*{(2.3)}
\]
et compatible avec les morphismes de transition pour $\alpha$ variable,
\[
\begin{tikzcd}
X_\alpha^{C_0} \arrow[r, "{\mathrm{trans}_{\beta,\alpha}^{C_0}}"] \arrow[d, "u_\alpha"'] & X_\beta^{C_0} \arrow[d, "u_\beta"] \\
X_\alpha^{C_1} \arrow[r, "{\mathrm{trans}_{\beta,\alpha}^{C_1}}"] & X_\beta^{C_1}\,.
\end{tikzcd}
\tag*{(2.4)}
\]

\msp{page 14}

Pour définir les $u_\alpha$, satisfaisant aux conditions précédentes, il suffit de les définir pour $\alpha=1$ ``et d'itérer'' :
\[
\begin{tikzcd}[column sep=small]
& X_1^{C_0}(f) \arrow[dd, "u_1(f)"] \arrow[dr, "p_1^{C_0}(f)"] & \\
X \arrow[ur, "i_1^{C_0}(f)"] \arrow[dr, "i_1^{C_1}(f)"'] & & Y \\
& X_1^{C_1}(f) \arrow[ur, "p_1^{C_1}(f)"'] & {}.
\end{tikzcd}
\tag*{(2.5)}
\]
Or $X_1^{C_0}(f)$ est défini comme une somme directe dans $X\backslash\cM/Y$, relative à un ensemble d'indices $\Delta(C_0,f)$ \textit{contenu} dans $\Delta(C_1,f)$ -- et la famille d'objets dont on prend la somme directe étant induite d'une famille (non petite) dont l'ensemble d'indices serait $\Delta(\Fl_{\text{coquar.}}(\cM),f)$, savoir formé de \textit{tous} les diagrammes carrés
\printedpage{11}
% head: A. Grothendieck : Dérivateurs, Ch. XV, Édition des Universités de Montpellier II et Paris VII
% foot: Version 5 novembre 2010
\[
D : \qquad
\begin{tikzcd}
A \arrow[r, "u"] \arrow[d, "i"'] & X \arrow[d, "f"] \\
B \arrow[r, "v"] & Y
\end{tikzcd}
\qquad\qquad \text{avec $i$ coquarrable sans plus ,}
\]
la définition des $X_D$, $i_D$, $p_D$ associés à $D$ ne dépendant que de $D$, et non de l'ensemble $C_0$ ou $C_1$ de flèches coquarrables qui contient $i$. Ceci dit, la flèche $u_1(f)$
\msp{page 15}
dans (2.5) n'est autre que l'inclusion canonique d'une somme amalgamée partielle, indexée par $\Delta(C_0,f)$, dans la somme amalgamée correspondante à l'ensemble d'indices plus grand. La fonctorialité de cette flèche $u_1(f)$, pour $f$ variable, ``est immédiate'', elle résulte de la commutativité du carré, associé à un morphisme $v=(v_X,v_Y)$ de flèches
\[
\begin{tikzcd}
X \arrow[r, "f"] \arrow[d, "v_X"'] & Y \arrow[d, "v_Y"] \\
X' \arrow[r, "f'"] & Y'\,,
\end{tikzcd}
\tag*{(2.6)}
\]
savoir le diagramme d'applications entre ensembles d'indices
\[
\begin{tikzcd}
\Delta(C_0,f) \arrow[r, hook, "u_1(f)"] \arrow[d, "{\begin{array}{r}\text{dépendance fonctorielle}\\ \text{en $f$ de $\Delta(C_0,f)$}\end{array}}"'] & \Delta(C_1,f) \arrow[d, "{\begin{array}{l}\text{dépendance fonctorielle}\\ \text{en $f$ de $\Delta(C_1,f)$}\end{array}}"] \\
\Delta(C_0,f') \arrow[r, hook, "u_1(f')"] & \Delta(C_1,f')
\end{tikzcd}
\tag*{(2.7)}
\]
(commutativité qui revient à dire que $\Delta(C_0,f)$ peut être considéré comme un foncteur en le couple $(C_0,f)$ comme objet de la catégorie produit $\underline{\mathfrak{P}_{\text{pet.}}(\Fl_{\text{coquar.}}(\cM))}\times\uFl(\cM)$ \ldots). Le carré commutatif (2.7) est la partie ``indices'' d'un carré commutatif entre familles d'objets $\mathscr{X}(C_0,f)$, $\mathscr{X}(C_0,f')$, $\mathscr{X}(C_1,f)$, $\mathscr{X}(C_1,f')$,
\msp{page 16}
indexés par ces ensembles, savoir les $X_D^{C_0}$ et les $X_D^{C_1}$ relatifs aux diagrammes $D$ impliquant, soit la flèche $f$, soit la flèche $f'$ \ldots

\dgrblock{dgrb726}{Functoriality in the generators, including limit stages}
\begin{dcomment}{exposition}{p15:commutativité qui revient à dire que}{Functoriality in the generators, including limit stages}
Grothendieck next keeps track of what happens when more lifting problems are declared relevant. The one-stage comparison is induced by the inclusion of indexing sets, but at later stages the input to the attachment functor has also changed. Writing the recursive formula makes the joint naturality precise.

Let $I\subseteq J$, put $P_I=p_1^I$ and $P_J=p_1^J$, and let $\tau:P_I\to P_J$ be the natural transformation constructed from the inclusion of squares. Both endofunctors lie over the codomain functor, and their units satisfy $\tau\eta_I=\eta_J$. Let $u_\alpha:P_I^\alpha\to P_J^\alpha$ denote the comparison at stage $\alpha$, where the superscript means the continuous iteration used in the text, not ordinary iteration at a limit. Then
\[
 u_0=1,\qquad
 u_{\alpha+1}=P_J(u_\alpha)\,\tau_{P_I^\alpha},\qquad
 u_\lambda=\mathop{\mathrm{colim}}_{\alpha<\lambda}u_\alpha
 \quad(\lambda\text{ a limit}).
\]
The successor formula starts with $P_I(P_I^\alpha f)$, changes the outer attachment functor to $P_J$, and then changes its input to $P_J^\alpha f$. Naturality of $\tau$ allows these two operations in the opposite order as well.

Here is the compatibility with a successor transition, which is the point needed to make the limit formula meaningful. Naturality of $\eta_J$ and the equation $\tau\eta_I=\eta_J$ give
\[
 \begin{split}
 u_{\alpha+1}\eta_{I,P_I^\alpha}
 &=P_J(u_\alpha)\eta_{J,P_I^\alpha}\\
 &=\eta_{J,P_J^\alpha}u_\alpha.
 \end{split}
\]
Compatibility with all transitions follows by composition and the defining colimit property at limits. Each $u_\alpha$ is natural in $f$ by the same induction. For $I\subseteq J\subseteq K$, the equation $u_\alpha^{K,I}=u_\alpha^{K,J}u_\alpha^{J,I}$ holds at stage one because both maps restrict to the same inclusions of cells. At successors it follows by naturality of the one-stage transformations, and at limits by equality on every colimit component. Identity inclusions are handled identically. This proves the functoriality on the product of the poset of generating sets with the arrow category that the final paragraph of Section~2 leaves unexpanded.

The term ``inclusion'' describes the indexing sets. It does not by itself assert that the resulting morphism of objects is a monomorphism in an arbitrary category. The following elementary refinement explains when the same word can also be used for the actual stages.
\end{dcomment}


Or la construction canonique du système inductif ordinal $p_\alpha$, avec les $i_\alpha : \mathrm{id}_{\uFl(\cM)}\longrightarrow p_\alpha$, est fonctorielle par rapport à la donnée initiale $p_1$, $i_1$ -- je n'explicite pas le détail de la fonctorialité, bien que je devrais \ldots

\dgrblock{dgrb727}{Comparisons which are monomorphisms}
\begin{dcomment}{development}{p16:je n'explicite pas le détail de la fonctorialité}{Comparisons which are monomorphisms}
The parameter dependence in Section~2 is useful when a family of generators is enlarged gradually. For presheaves, monomorphic generators give a stronger conclusion: an enlargement can be seen stage by stage inside the resulting objects. We record the precise two-variable statement, with its proof, for this construction. We checked the small object construction in \citet[Section~12.2]{Riehl2014} and the algebraic comparison in \citet[Sections~3--4]{Garner2009}; we did not find this particular formulation there. The statement below is an elementary elaboration of the displayed construction, with no claim of priority.

\textbf{Proposition (monomorphic comparison).}\ours
Let $\mathcal E=\mathbf{Set}^{\mathcal A^{\mathrm{op}}}$ for a small category $\mathcal A$. Let $I\subseteq J$ be small sets of monomorphisms in $\mathcal E$. Suppose $(u,v):f\to f'$ is a commutative square whose source map $u:X\to X'$ and target map $v:Y\to Y'$ are both monomorphisms. For every ordinal $\alpha$ for which the construction is performed, the jointly natural comparison
\[
 X_\alpha^I(f)\longrightarrow X_\alpha^J(f')
\]
is a monomorphism. In particular, at a fixed input $f$, enlargement of a set of monomorphic generators gives monomorphisms at every stage.

\textit{Proof.} First consider one stage. The map
\[
 \Delta(I,f)\longrightarrow\Delta(J,f'),\qquad
 (j,a,b)\longmapsto(j,ua,vb)
\]
is injective. Indeed equal images have the same generator $j$, and cancellation of the monomorphisms $u$ and $v$ gives equality of both maps $a$ and $b$. Evaluate the pushout formula at an object $c\in\mathcal A$. Evaluation preserves colimits, and every $j:A_j\to B_j$ is pointwise injective. Thus, as sets,
\[
 X_1^I(f)(c)\cong X(c)\amalg
 \coprod_{(j,a,b)\in\Delta(I,f)}
 \bigl(B_j(c)\setminus j(A_j(c))\bigr).
\]
This decomposition is only used at the level of sets; the complements need not form subpresheaves. On the first summand the comparison is $u(c)$, which is injective. On each remaining summand it is the identity on the same complement, placed in the summand with image index. Injectivity of the indexing map shows that distinct summands remain distinct. Hence the one-stage comparison is pointwise injective.

At stage zero the assertion is the hypothesis on $u$. If it holds at stage $\alpha$, apply the one-stage result to the square between $p_\alpha^I(f)$ and $p_\alpha^J(f')$; its source map is monic by induction and its target map is still $v$. This proves the successor step. At a limit $\lambda$, evaluate at $c$ and use the two compatible diagrams of sets. If two representatives in the first colimit have equal images in the second, they become equal at a common later stage of the second diagram. The injectivity of the comparison at that stage makes their images in the first diagram equal there as well. They therefore represent the same element of the first colimit. This proves the limit step and the proposition.\hfill$\square$

The assumption on the target map carries real information. With $I=J=\{\varnothing\to\{*\}\}$, take $f:\varnothing\to\{0,1\}$ and $f':\varnothing\to\{*\}$, the identity on the empty source and the constant map on targets. The induced map between first stages is $\{0,1\}\to\{*\}$, so it is not monic. The construction distinguishes lifting problems by their target component as well as their boundary map.
\end{dcomment}


\section*{3\quad Majoration des cardinaux}

Je suppose à présent que $\cM$ est muni d'une \textit{filtration cardinale}
\[
(\Filt^\pi)_{\pi\geq\pi_0} \qquad\quad (\pi_0,\ \pi \text{ cardinaux dans } \mathfrak{U})
\tag*{(3.1)}
\]
(SGA 4 I 9.12, p.~149), les $\Filt^\pi$ stables par sommes amalgamées simples (chaque fois qu'elles existent), donc aussi par sommes amalgamées infinies indexées par des ensembles
\printedpage{12}
% head: A. Grothendieck : Dérivateurs, Ch. XV, Édition des Universités de Montpellier II et Paris VII
% foot: Version 5 novembre 2010
d'indices de cardinal $\leq\pi$. On suppose de plus que tout objet de $\Filt^{\pi_0}$ soit accessible, et on désigne par \circled{$\pi_1$} le plus petit cardinal tel que $\pi_1\geq\pi_0$ et que les objets de $\Filt^{\pi_0}$ soient $\pi_1$-accessibles. Désignant, pour tout cardinal, par $\cM_\pi$ le sous-ensemble de $\mathrm{Ob}\,\cM$ formé des objets $\pi$-accessibles, on aura alors pour tout $\pi\geq\pi_1$ (\textit{loc. cit.} 9.16)
\[
\cM_\pi \;\subset\; \Filt^\pi \;\subset\; \cM_{(\pi^{\pi_0})} \qquad\quad (\text{si } \pi\geq\pi_1)\,,
\tag*{(3.2)}
\]

\msp{page 17}

\dgrblock{dgrb728}{A filtration measures the cost of a construction}
\begin{dcomment}{exposition}{p16}{A filtration measures the cost of a construction}
Grothendieck now asks for more than the existence of a factorization. He wants to perform the construction inside a bounded part of the category, uniformly for every arrow in that part. This is what will eventually let a small amount of data control a model structure. The three quantities to keep separate are the size of the domains and codomains of the generators, the number of generators actually used, and the length of the construction. Isomorphic generators have the same lifting property, but the unmodified construction attaches a cell for each of them.

The terminology comes from \citet[Expos\'e I, Definitions 9.1--9.3 and 9.12]{SGA4I1972}. A partially ordered set is \emph{large relative to} an infinite cardinal $\pi$ if each of its subsets of cardinality at most $\pi$ has an upper bound. Thus the modern directedness index is $\pi^+$. An object called $\pi$-accessible here has a representable functor preserving such colimits; in the usual setting this is $\pi^+$-presentability, not presentability at the cardinal $\pi$. The filtration is additional structure: its levels are essentially small, are increasing, admit bounded presentations of objects at larger levels, and have specified closure properties. Membership in a filtration level should therefore not be replaced by an informal assertion that an object has at most $\pi$ elements.

For the calculation ahead we need the following concrete consequences. Put $\mathcal M^\pi=\Filt^\pi$. If $A\in\mathcal M^r$, it has a filtered presentation by objects of $\mathcal M^{\pi_0}$ with at most $r^{\pi_0}$ indices. If $X\in\mathcal M^s$ and $s\geq\pi_1$, it has a presentation by objects of $\mathcal M^{\pi_1}$ over a $\pi_1^+$-directed poset with at most $s^{\pi_1}$ indices. These are the presentations used on ms.~20. The stronger closure condition (3.13), which is explicitly assumed below, includes \emph{all} filtered diagrams of cardinality at most the level, not only those large relative to $\pi_0$.

This stronger closure also supplies the retract argument in (3.2). If $X$ is a retract of $Y\in\mathcal M^\pi$, with idempotent $e:Y\to Y$, the countable chain
\[
Y\xrightarrow{e}Y\xrightarrow{e}Y\xrightarrow{e}\cdots
\]
has colimit $X$, so $X\in\mathcal M^\pi$. If $X$ is $\pi$-accessible, factor its identity through a stage of its $\pi^+$-directed presentation by $\mathcal M^\pi$-objects. This makes it such a retract. In the other direction, a presentation with at most $\pi^{\pi_0}$ arrows expresses $\Hom(X,-)$ as a limit of that many representables. Directed colimits of sets with index larger than this bound commute with those limits. This proves the inclusions in (3.2) under the closure hypotheses used here and explains why their equality at suitable cardinals is useful.
\end{dcomment}


donc si $\pi^{\pi_0}=\pi$, p.~ex. si $\pi=2^c$ avec $c\geq\pi_0$, on aura
\[
\Filt^\pi \;=\; \cM_\pi\,.
\tag*{(3.2 bis)}
\]
(Si on admet l'hypothèse du continu généralisée, on a $\pi^{\pi_0}=\pi$ pour tout $\pi>\pi_0$, donc $\Filt^\pi=\cM_\pi$ si $\pi\geq\pi_1$ et $\pi>\pi_0$.)

Je pose, pour tout cardinal $\pi\geq\pi_0$,
\[
\Fl^\pi(\cM) \;=\; \{\ f\in\Fl(\cM) \mid \text{source et but de $f$ sont dans } \Filt^\pi\ \}
\tag*{(3.3)}
\]
-- ce qui sans doute donne une filtration cardinale sur $\uFl(\cM)$ (mais ce n'est pas là notre souci \ldots). Mon but à présent, c'est de dégager des conditions sur $C_0\subset\Fl_{\text{coquar.}}(\cM)$ (petite partie de $\Fl(\cM)$) et sur le cardinal $\pi$, et l'ordinal $\alpha$, pour que l'on ait
\[
\begin{aligned}
f\in\Fl^\pi(\cM) \;\;\overset{?}{\Longrightarrow}\;\; & X_\alpha(f)\in\Filt^\pi \\
& (\textit{i.e.}\ p_\alpha(f)\in\Fl^\pi(\cM),\ \text{ou encore } i_\alpha(f)\in\Fl^\pi(\cM))\,.
\end{aligned}
\tag*{(3.4)}
\]
Plus généralement, je me propose de donner une estimation de la ``cardinalité'' de $X_\alpha(f)$, pour $C_0$ fixé, en fonction de $C_0$, de $f$, de $\alpha$, plus exactement, en fonction des trois cardinaux $\pi(C_0)$, $\pi(f)$, $\pi(\alpha)$ définis comme
\[
\left\{
\begin{array}{rcl}
\pi(C_0) & = & \text{plus petit cardinal tel que } C_0\subset\Fl^\pi(\cM) \\[1ex]
\pi(f) & = & \pi(\{f\}) \;=\; \text{plus petit cardinal tel que } f\in\Fl^\pi(\cM) \\[1ex]
\pi(\alpha) & = & \card\alpha\,.
\end{array}
\right.
\tag*{(3.5)}
\]
Donc je cherche une fonction
\msp{page 18}
cardinale de variables cardinales $\varphi$, telle que l'on ait
\[
X_\alpha(f) \;\in\; \Filt^\pi \qquad\quad \text{si} \qquad\quad \pi \;=\; \varphi(\underbrace{\pi(C_0)}_{r},\underbrace{\pi(f)}_{s},\underbrace{\pi(\alpha)}_{t})\,.
\tag*{(3.6)}
\]
L'estimation cruciale concernera sûrement le cas $\alpha=1$. Notons que $X_1^{C_0}(f)=X_1(f)$ est défini comme une somme amalgamée (en général infinie) d'objets $X_D$, pour un ensemble $\Delta(C_0,f)=\Delta(f)$ de diagrammes, savoir les diagrammes commutatifs
\[
\begin{tikzcd}
A \arrow[r] \arrow[d, "i"'] & X \arrow[d, "f"] \\
B \arrow[r] & Y
\end{tikzcd}
\]
\printedpage{13}
% head: A. Grothendieck : Dérivateurs, Ch. XV, Édition des Universités de Montpellier II et Paris VII
% foot: Version 5 novembre 2010
avec $i\in C_0$. Pour un $i$ fixé, et un $f$ fixé, le cardinal de l'ensemble des diagrammes de ce type est majoré par celui de $\mathrm{Hom}(A,X)\times\mathrm{Hom}(B,Y)$. On a $A,B\in\Filt^r$, $X,Y\in\Filt^s$, donc il nous faut une majoration de
\[
\card\mathrm{Hom}_{\cM}(A,X)\,, \qquad\quad A\;\in\;\Filt^r\,, \qquad\quad X\;\in\;\Filt^s\,.
\]

\dgrblock{dgrb729}{The cofinality qualification in the cardinal arithmetic}
\begin{dcomment}{exposition}{p17}{The cofinality qualification in the cardinal arithmetic}
Grothendieck looks for cardinals at which the estimate stops growing. The unconditional choice $\pi=2^c$, $c\geq\pi_0$, works because
\[
(2^c)^{\pi_0}=2^{c\pi_0}=2^c.
\]
The parenthetical appeal to the generalized continuum hypothesis has a different status. Even under GCH, $\pi>\pi_0$ does not imply $\pi^{\pi_0}=\pi$. For example, $\pi=\aleph_\omega$ and $\pi_0=\aleph_0$ give $\pi^{\pi_0}>\pi$. More generally, K\"onig's theorem gives $\pi^{\operatorname{cf}\pi}>\pi$ for every infinite $\pi$. A direct proof of the case just used is instructive: if $(f_\xi)_{\xi<\aleph_\omega}$ enumerated all maps $\mathbb N\to\aleph_\omega$, choose $g(n)$ outside the set $\{f_\xi(n):\xi<\aleph_n\}$. For every $\xi$ some $n$ satisfies $\xi<\aleph_n$, and then $g(n)\ne f_\xi(n)$, a contradiction.

Under GCH the additional condition $\pi_0<\operatorname{cf}\pi$ is sufficient. The general exponentiation rule and its cofinality qualification are recorded in \citet[Theorems 10.12 and 11.7, Exercise 11.8]{Marks2026}. Nothing in the factorization theorem requires GCH: the fixed exponent identity above is all that is needed. The same correction applies to the repeated GCH assertion on ms.~19.

The sentence about the arrow category immediately before (3.4) should also be read with its size threshold in view. Pointwise filtered colimits and essential smallness of the bounded arrow categories are immediate. Uniformly bounded presentations of arrows require control of the Hom sets between small objects and simultaneous realization of the arrow in two presentations. At a sufficiently enlarged initial level this is the usual accessibility theorem for a finite diagram category. The text does not supply that extra argument at the original $\pi_0$, and does not use the assertion in its proof.
\end{dcomment}


\msp{page 19}

Voici les majorations essentielles dans ce sens :

\textbf{Proposition 3.1.} \textit{Soit $\pi_0$ le cardinal initial de la filtration cardinale} $(\Filt^\pi)_{\pi\geq\pi_0}$ \textit{envisagé, $\pi_1$ le plus petit cardinal $\geq\pi_0$ tel que les objets de} $\Filt^{\pi_0}$ \textit{soient $\pi_1$-accessibles, $C_1=$} card $\underbrace{\widetilde{\uFl^{\pi_1}(\cM)}}_{=\uFl(\cM^{\pi_1})}$\pfn{4}{\textit{cf.} plus bas la notation $\widetilde{\mathcal{C}}$.}. \textit{Pour tout cardinal $\pi\geq\pi_0$, je désigne par $\cM^\pi$ la sous-catégorie strictement pleine de $\cM$ dont l'ensemble des objets est} $\Filt^\pi$. \textit{Enfin, si $\mathcal{C}$ est une catégorie, je désigne par $\widetilde{\mathcal{C}}$ l'ensemble des classes d'isomorphie d'objets de $\mathcal{C}$.}\pfn{5}{Signalons qu'on a toujours $\card\widetilde{\uFl(\mathcal{C})}\geq\card\widetilde{\mathcal{C}}$ car $\widetilde{\mathcal{C}}\hookrightarrow\widetilde{\uFl(\mathcal{C})}$ via $\mathrm{cl}(X)\longmapsto\mathrm{cl}(\mathrm{id}_X)$.}
% ?? uncertain: straight underlines under "le plus petit" and under "pleine" -- underlined / not underlined (a wavy tilde and an underbrace stand close by)

$1^\circ$) \textit{Si $r\geq\sup(\pi_1,C_1)$, on a}
\[
\card\,(\underbrace{\widetilde{\uFl^r(\cM)}}_{=\uFl(\cM^r)}) \;\leq\; 2^{(r^{\pi_0})} \quad\ (=2^r,\ \textit{si } r=2^c,\ \textit{avec } c\geq\pi_0)\,.
\tag*{(3.7)}
\]

$2^\circ$) \textit{Si $r,s\geq\sup(\pi_1,C_1)$, et si $A\in\Filt^r$, $X\in\Filt^s$, on a}
% (3.8) and (3.9) share one brace on the right
\[
\begin{array}{l}
\text{(3.8)}\\[3ex]
\text{(3.9)}
\end{array}
\qquad\qquad
\begin{array}{r@{\;\;\leq\;\;}l}
\card\mathrm{Hom}_{\cM}(A,X) & s^{(r^{\pi_0})}\\[3ex]
\card\widetilde{\uFl(\cM^r/X)} & s^{(r^{\pi_0})}
\end{array}
\qquad\quad
\left\{
\begin{array}{lll}
= & s^r & \textit{si } r=2^c,\ \textit{avec } c\geq\pi_0\,,\\[3ex]
= & s & \textit{si } s=2^{c'},\ \textit{avec } c'\geq r^{\pi_0}\,.
\end{array}
\right.
\]

\textbf{N.B.} Si on admet l'hypothèse du continu généralisée, qui implique que $\pi^{\pi'}=\pi$ si $\pi>\pi'$ ($\pi$ cardinal infini, $\pi'\geq1$), alors prenant $r$ tel que $r>\pi_0$, les formules précédentes donnent (avec de plus $r\geq\sup(\pi_1,C_1)$) :
\[
\card\widetilde{\uFl^r(\cM)} \;\;\leq\;\; 2^r\,,
\tag*{(3.7 bis)}
\]
\[
\card\mathrm{Hom}_{\cM}(A,X) \;\leq\; s^r \qquad\quad (\text{si } A\in\Filt^r,\ X\in\Filt^s)\,,
\tag*{(3.8 bis)}
\]
\[
\card\widetilde{\uFl(\cM^r/X)} \;\leq\; s^r \qquad\quad (\text{si } X\in\Filt^s)\,.
\tag*{(3.9 bis)}
\]
Sans hypothèse du continu, ces formules sont aussi valables pour $r$ de la forme $2^c$, avec $c\geq\pi_0$. Ce sont ces dernières formules qui on l'air les plus naturelles. Pour (3.8 bis) et % [sic] qui on l'air -> qui ont l'air, verb form
(3.9 bis), je signale que le deuxième membre $s^r$ est le cardinal de l'ensemble $\mathrm{Hom}_{\mathrm{Ens}}(A,X)$, où $A$, $X$ sont deux ensembles tels que $\card A=r$, $\card X=s$. On ne peut donc espérer améliorer (3.8 bis)\,! Pour la

% ---- from Der-ChXV-05-11-10_p016-020.tex
\printedpage{14}
% head: A. Grothendieck : Dérivateurs, Ch. XV, Édition des Universités de Montpellier II et Paris VII
% foot: Version 5 novembre 2010
\msp{page 20}
première formule (3.7 bis), notons que l'on a égalité lorsque $\M$ est la catégorie des couples $(A,B)$, $A$ un ensemble et $B$ une partie de $A$, avec les flèches évidentes, et qu'on définit $\M^\pi$ (pour $\pi\geq\pi_0$) comme la sous-catégorie strictement pleine formée des $(A,B)$ tels que $\card A\geq\pi$.

\dgrblock{dgrb730}{An orbit count is not a bound on a Hom set}
\begin{dcomment}{examples}{p19}{An orbit count is not a bound on a Hom set}
The purpose of $C_1$ is to pay for every choice of a map between the small objects. As printed, however, it counts isomorphism classes of objects of the arrow category. These two numbers can be very different. This is a mathematical issue in the edition, not an error introduced by the transcription.

Let $G$ be any small group and let $BG$ have one object $*$, with $\Hom(*,*)=G$. Give $BG$ the constant filtration $\mathcal M^\pi=BG$, with $\pi_0=\pi_1=\aleph_0$. Every filtered diagram in $BG$ is isomorphic to the constant diagram: transport from a fixed index to another through a common upper bound; a further common upper bound shows that transport does not depend on the choices. Consequently filtered colimits exist, and $\Hom(*,-)$ preserves them. Pushouts of spans of isomorphisms exist as well. The constant filtration has all the presentation and closure properties needed here.

Every two arrows $g,h$ are isomorphic as objects of $\Fl(BG)$: take the square with vertical maps $1$ and $hg^{-1}$. Thus the printed constant is $C_1=1$, whereas $|\Hom(*,*)|=|G|$. Take $r=s=\aleph_0$ and a group of cardinality greater than $\aleph_0^{(\aleph_0^{\aleph_0})}$. Formula (3.8), with the printed $C_1$, is false. The example also identifies precisely the invalid step on ms.~20: the indicated bound $|\Hom(A_i,X_j)|\leq C_1$ does not follow from the definition of $C_1$.

There is a straightforward repair. Choose a small full skeleton $K_1$ of $\mathcal M^{\pi_1}$, choosing its representatives compatibly with those at level $\pi_0$, and use
\[
\chi=\max\{\aleph_0,\pi_1,|\mathrm{Mor}(K_1)|\}
\]
as the threshold for $r$ and $s$. It counts actual maps between representatives, including automorphisms. This change is also the natural reading of the lowercase $c_1$ in Theorem~2: the entire strictly full level need not be a small category, so its raw collection of arrows cannot serve as a small cardinal. The choice of skeleton affects a numerical bound, not the factorization or the lifting classes.
\end{dcomment}


\textsc{Démonstration de} 3.1. Prouvons d'abord la formule (3.8),
\[
\card\Hom(A,X)\ \leq\ s^{(r^{\pi_0})}\qquad\quad
\left\{\begin{array}{l}
\text{si } A\in\Filt^r,\ X\in\Filt^s,\\[1.5ex]
r\geq\pi_1,\ s\geq\sup(\pi_1,C_1).
\end{array}\right.
\]

(\textbf{N.B.} Je n'utiliserai pas $r\geq C_1$, postulé dans 3.1 $1^\circ$.) J'écris \texttt{[<<J'écris>> rayé]}
\[
\begin{aligned}
A\ &=\ \varinjlim_I A_i\ ,\quad \card I\ \leq\ r^{\pi_0}\ ,\quad A_i\ \in\ \Filt^{\pi_0}\ ,\\
X\ &=\ \varinjlim_J X_j\ ,\quad \card J\ \leq\ s^{\pi_1}\ ,\quad X_j\ \in\ \Filt^{\pi_1}\ ,\qquad\qquad J\text{ grand devant }\pi_1\ .
\end{aligned}
\]
On aura
\[
\Hom(A,X)\ \simeq\ \varprojlim_{i\in I}\ \underbrace{\Hom(A_i,\,\varinjlim_J X_j)}_{\overset{(\text{\footnotemark[6]})}{\simeq}\ \varinjlim\limits_{j\in J}\ \underbrace{\Hom(A_i,X_j)}_{\text{de cardinal }\leq C_1}}\ ,
\]
\footnotetext[6]{car les $A\in\Filt^{\pi_0}$ sont $\pi_1$-accessibles.}
donc pour tout $i\in I$
\[
\Hom(A_i,X)\ \leq\ \card(J)\times C_1\ \leq\ s^{\pi_1}\times C_1\ =\ s^{\pi_1}\ ,
\]
car $s\geq C_1$, donc $s^{\pi_1}\geq C_1$, d'où
\[
\card(\,\varprojlim_{j\in I}\Hom(A_i,X))\ \leq\ (s^{\pi_1})^{\card I}\ \leq\ \underbrace{(s^{\pi_1})^{(r^{\pi_0})}}_{s^{\pi_1 r^{\pi_0}}}\ . % [sic] j\in I -> i\in I, the limit is over the index i of A_i
\]
Comme $r\geq\pi_1$, on a $r^{\pi_0}\geq\pi_1$, donc $r^{\pi_0}\pi_1=r^{\pi_0}$, d'où
\[
\card\Hom(A,X)\ \leq\ s^{(r^{\pi_0})}\ ,
\]
c'est bien la formule (3.8). Pour en tirer une estimation pour $\M^r/X$ (avec $X\in\Filt^s$), \msp{page 21} je vais d'abord majorer le cardinal des classes d'isomorphie d'objets de $\M^r/X$, \textit{i.e.} des flèches à isomorphisme près (dans $\M^r$)
\[
A\ \longrightarrow\ X\ ,\qquad\quad X\text{ fixé dans }\Filt^s,\ A\text{ variable dans }\M^r\ .
\]
On aura
\[
A\ =\ \varinjlim_{i\in I} A_i\qquad\quad\text{avec } A_i\in\Filt^{\pi_0},\ I\text{ ensemble ordonné de cardinal}\leq r^{\pi_0}\ .
\]

\printedpage{15}
% head: A. Grothendieck : Dérivateurs, Ch. XV, Édition des Universités de Montpellier II et Paris VII
% foot: Version 5 novembre 2010
Pour $I$ fixé, le cardinal de l'ensemble des systèmes inductifs indexés par $I$, dans $\M^{\pi_0}/X$, à isomorphisme près, peut se majorer ainsi : pour tout $i\in I$, il faut choisir $A_i$ dans $\M^{\pi_0}$ (mais seule sa classe dans $\wavy{\M^{\pi_0}}$ importe), donc c'est majoré par card $\wavy{\M^{\pi_0}}\ \leq C_0$~\textsuperscript{(}\footnote[7]{On a posé ici $C_0=\card\wavy{\uFl^{\pi_0}}$ , donc $C_0\leq C_1$.}\textsuperscript{)}, et pour tous les $i$, par $\card I\times C_0\leq r^{\pi_0}\times C_0$, et si $r\geq C_0$, donc $r^{\pi_0}\geq C_0$, on trouve une majoration des choix possibles du système des \textit{objets} $A_i$ dans $\M^\pi$, à isomorphisme près, par $r^{\pi_0}$. Si on se donne de plus les morphismes de transition $A_i\longrightarrow A_j$ pour $i<j$, cela fait pour chaque paire comparable, élément de $I\times I$, un choix majoré par $C_0$, donc encore une majoration par $\ \underbrace{\card(I\times I)}_{=(\card I)^2=\card I}\cdot C_0\leq r^{\pi_0}\cdot C_0=r^{\pi_0}$. Donc l'ensemble des systèmes inductifs en question, à isomorphisme près, pour $I$ fixé, est de cardinal $\leq r^{\pi_0}\cdot r^{\pi_0}=r^{\pi_0}$. Enfin, les \textit{homomorphismes} d'un système inductif donné dans $X$ sont majorés par $\card I\times s^{(r^{\pi_0})}$ (par (3.8) déjà prouvé), donc par
\[
r^{\pi_0}\cdot s^{(r^{\pi_0})}\ =\ s^{(r^{\pi_0})}\qquad\quad(\text{puisque } r^{\pi_0}\leq s^{(r^{\pi_0})})\ .
\]

Au total, pour $I$ fixé, l'ensemble des systèmes inductifs à isomorphisme près, indexés par $I$, à valeurs dans $\M^r/X$, est \msp{page 22} de cardinal majoré par $\quad\underbrace{r^{\pi_0}}_{\text{\scriptsize\begin{tabular}{@{}c@{}}nombre des\\ systèmes inductifs\\ dans $\M^\pi$, à\\ isomorphisme près\end{tabular}}}\quad\cdot\quad\underbrace{s^{(r^{\pi_0})}}_{\text{\scriptsize\begin{tabular}{@{}c@{}}nombre des\\ flèches d'un\\ tel système inductif\\ dans $X$\end{tabular}}}\quad=s^{(r^{\pi_0})}$.

Mais il faut maintenant faire varier $I$, à isomorphisme près, dans la catégorie des ensembles ordonnés (filtrants \textit{etc.}) de cardinal $\leq\lambda=r^{\pi_0}$. On voit de suite que cet ensemble de classes d'isomorphie (qui peut s'interpréter comme un ensemble quotient d'un certain ensemble de parties d'un ensemble fixe $I_0$ de cardinal $\lambda$) est de cardinal $\leq 2^\lambda$. On trouve donc finalement
\[
\card\wavy{\M^r/X}\ \ \leq\ s^{(r^{\pi_0})}\cdot 2^{(r^{\pi_0})}\ =\ s^{(r^{\pi_0})}\ ,
\]
sous condition que
\[
r\ \geq\ \sup(\pi_1,C_0)\ ,\qquad\quad s\ \geq\ \sup(\pi_1,C_1)\ .
\]
Il faut prouver l'inégalité similaire pour card $\wavy{\uFl^r\,\M}$ \ (3.9). Or, comme pour $A$, $B$ dans $\M^r$, on a, si $r\geq\sup(\pi_1,C_1)$,
\[
\card\Hom_{\M}(A,B)\ \leq\ r^{(r^{\pi_0})}
\]
en vertu de (3.8) où on fait $s=r$, on trouve que
\[
\begin{aligned}
\card\wavy{\uFl(\M^r/X)}\ &\leq\ \Bigl(\card\wavy{\M^r/X}\,\Bigr)^2\cdot r^{(r^{\pi_0})}\\
&\leq\ s^{(r^{\pi_0})}\cdot r^{(r^{\pi_0})}\\
&=\ s^{(r^{\pi_0})}\qquad\text{si } s\geq r\ .
\end{aligned}
\]
(On va regarder plus bas les cas $s\leq r$~\textsuperscript{(}\footnote[8]{\textit{cf.} cor. 2, page 23.}\textsuperscript{)}). C'est bien la formule (3.9).

\printedpage{16}
% head: A. Grothendieck : Dérivateurs, Ch. XV, Édition des Universités de Montpellier II et Paris VII
% foot: Version 5 novembre 2010
\msp{page 23}
Venons-en à la catégorie $\M^r$ elle-même. L'argument de la page 21 montre que l'ensemble des classes d'isomorphie d'objets de $\M^r$ indexés par un $I$ fixé, de cardinal $\leq r^{\pi_0}$, comme $A=\varinjlim\limits_i A_i$ avec $A_i\in\Filt^{\pi_0}$, est de cardinal majoré par $r^{\pi_0}$ (si $r\geq C_0$). Si on fait varier $I$, on trouve une majoration par $r^{\pi_0}\cdot 2^{(r^{\pi_0})}=2^{(r^{\pi_0})}$. On en conclut comme tantôt que l'ensemble $\ \wavy{\Fl\,\M^r}$ est de cardinal majoré par
\[
(2^{(r^{\pi_0})})^2\times r^{(r^{\pi_0})}\qquad\quad(\text{si } r\geq\sup(\pi_1,C_1))\ ,
\]
donc par $r^{(r^{\pi_0})}$. Je dis qu'on a
\[
r^{(r^{\pi_0})}\ =\ 2^{(r^{\pi_0})}\ ,
\]
en effet, on a $\geq$, et pour prouver $\leq$, on note que $r<2^r$, donc $r^{(r^{\pi_0})}\leq(2^r)^{(r^{\pi_0})}=2^{r\cdot r^{\pi_0}}=2^{(r^{\pi_0})}$, O.K.! Cela établit donc 3.7. Cet argument prouve en fait le

\dgrblock{dgrb731}{The estimates survive the corrected counting argument}
\begin{dcomment}{exposition}{p22}{The estimates survive the corrected counting argument}
Grothendieck's desired bound has enough room to absorb the choices that were omitted in the displayed counting. We give the calculation with the corrected threshold $r,s\geq\chi$ of the preceding block. Write $\mu=r^{\pi_0}$. For the presentations on ms.~20, accessibility gives
\[
\Hom(A_i,X)\cong\varinjlim_j\Hom(A_i,X_j),
\qquad |\Hom(A_i,X)|\leq s^{\pi_1}.
\]
Taking a limit over the presentation of $A$ then yields
\[
|\Hom(A,X)|\leq (s^{\pi_1})^{|I|}
\leq s^{\pi_1\mu}=s^\mu.
\]
This proves the repaired (3.8). The last equality uses $\pi_1\leq r\leq\mu$ and infinite cardinal arithmetic.

The choices of all the objects $A_i$ must be counted by a product, not by $|I|$ times the number of possibilities at one index. The same applies to the transition arrows. A poset with at most $\mu$ objects has at most $\mu$ arrows; a diagram on it with values in a fixed small skeleton at level $\pi_0$ can therefore be specified in at most $\chi^\mu\leq 2^\mu$ ways. The requirements of identity and composition merely reduce this number. There are at most $2^\mu$ isomorphism types of posets of size at most $\mu$, since both their underlying sets and their order relations can be coded inside a fixed set of size $\mu$ and its square. Taking colimits gives
\[
|\operatorname{IsoOb}(\mathcal M^r)|\leq 2^\mu.
\]
This replaces the unsupported fixed-index bound $r^{\pi_0}$ on ms.~21 and~23 by $2^\mu$; the eventual bound is unchanged.

For a fixed $X\in\mathcal M^s$, choose one representative of each object at level $r$. There are at most $s^\mu$ arrows from any such representative to $X$, hence at most $2^\mu s^\mu=s^\mu$ objects of $\mathcal M^r/X$ up to isomorphism. Between two of them there are at most $r^\mu=2^\mu$ morphisms, because a morphism in the slice is a particular underlying morphism in $\mathcal M$. Consequently
\[
|\operatorname{IsoOb}(\Fl(\mathcal M^r/X))|
\leq(s^\mu)^2 2^\mu=s^\mu.
\]
This is (3.9), including the case $s<r$. Finally, applying the same argument to the arrow category without a fixed target gives (3.7). In particular the appearance of products where exponential counts are required is a false intermediate inference, but it does not force a worse final estimate once the small skeleton has been counted correctly.

The elementary identity used here is $a^b=2^b$ for infinite $b$ and $2\leq a\leq b$: one has $2^b\leq a^b\leq(2^b)^b=2^b$. These are the hypotheses under which Lemma~3.2 is used. Its unqualified wording should not be extended to arbitrary finite cardinals, for which, for example, $3^3\ne2^3$.
\end{dcomment}


\textbf{Lemme 3.2.} \textit{Soient $r$ un cardinal, $r'$ un cardinal $\geq r$. Alors on a}
\[
r^{r'}\ =\ 2^{r'}\qquad\quad(=s^{r'}\ \textit{pour tout}\ 2\leq s\leq r)\ .
\]

\textbf{Corollaire 1.} \textit{Soient $r\geq\sup(\pi_1,C_1)$, et $A,X\in\Filt^\pi$. Alors}
\[
\card\Hom_{\M}(A,X)\ \leq\ 2^{(r^{\pi_0})}\qquad\quad(=2^r\ \textit{si}\ r=2^c,\ \textit{avec}\ c\geq\pi_0)\ . \tag*{(3.10)}
\]

On applique (3.8), qui donne la majoration $r^{(r^{\pi_0})}$, et le lemme 3.2.

On notera qu'on trouve la même majoration pour le cardinal de $\wavy{\uFl(\M^r)}$ tout entier (3.7)! De même, (3.9) donne :

\textbf{Corollaire 2.} \textit{Soit $r\geq\sup(\pi_1,C_1)$, et $X\in\Ob\M^r$. Alors}
\[
\card\wavy{\uFl(\M^r/X)}\ \leq\ 2^{(r^{\pi_0})}\qquad\quad(=2^r\ \textit{si}\ r=2^c,\ \textit{avec}\ c\geq\pi_0)\ . \tag*{(3.11)}
\]
c'est donc la formule (3.9) pour $s\leq r$, laissée en suspens.

\msp{page 24}
Revenons à la majoration cardinale des $X_\alpha(f)$, où $f : X\longrightarrow Y$ est dans $\Fl^s(\M)$, \textit{i.e.} $X,Y\in\Filt^s$. On suppose de plus $C_0\subset\Fl^r(\M)$,
\[
\begin{aligned}
&C_0\ \subset\ \Fl^r(\M)\ =\ \Fl(\M^r)\ ,\qquad\quad(f:X\longrightarrow Y)\in\Fl^s(\M)\ ,\ \ \textit{i.e.}\ X,Y\in\Filt^s\ ,\\
&\text{avec } s\geq r\ .
\end{aligned} \tag*{(3.10)}
\]

\dgrblock{dgrb732}{What the proposed sharpness example does show}
\begin{dcomment}{examples}{p23}{What the proposed sharpness example does show}
Grothendieck rightly observes that the Hom bound cannot improve on the elementary set-theoretic exponent $s^r$: sets of cardinalities $r$ and $s$ have exactly that many functions. The claim that the category of pairs $(A,B)$, $B\subseteq A$, attains the arrow-isomorphism bound $2^r$ at every allowed $r$ needs a separate check. First, the printed inequality $|A|\geq\pi$ on ms.~20 must be $|A|\leq\pi$ to describe the intended increasing bounded filtration; the original sign remains in the source text and is recorded in the corrections register.

Even after that change, equality does not hold at every cardinal. Assume GCH and take $r=\aleph_1$, $\pi_0=\pi_1=\aleph_0$. A map of such pairs is a function between sets with two colours, required to preserve the distinguished colour. Its isomorphism type is determined by the colour of a target element, the two cardinalities of its source fibre, and the multiplicity of each such fibre type. Among cardinalities at most $\aleph_1$ there are only countably many choices. Thus there are countably many fibre types, and at most
\[
\aleph_1^{\aleph_0}=\aleph_1
\]
possible multiplicity lists. The count of arrow-isomorphism types is therefore at most $\aleph_1$, strictly less than $2^{\aleph_1}=\aleph_2$. The threshold involving the printed $C_1$ is met: arrows between countable coloured sets have at most $2^{\aleph_0}=\aleph_1$ isomorphism types.

The same example does attain $2^{\aleph_0}$ at the countable level. For each subset $S$ of the positive integers, take one fibre of size $n$ for each $n\in S$, and add a countably infinite fibre so that the source remains countable. Distinct $S$ give nonisomorphic arrows. This is an orbit-count example at a level below the proposition's threshold $C_1=2^{\aleph_0}$, rather than an admissible equality case for that proposition. These calculations retain the useful point of the example while separating a genuine sharp Hom estimate from a stronger, false, universal sharpness assertion about arrow orbits.
\end{dcomment}


\texttt{[Deux fois (3.10).]} Il y a une difficulté idiote, c'est que $\M^r$ est équivalente à une petite catégorie, mais n'est pas elle-même petite. Même si $C_0$ ne contenait qu'une seule classe d'isomorphie de flèches de $\M^r$, son cardinal pourrait être (pour $r$ fixé) arbitrairement élevé. Aussi, dans ce qui suit, je dois supposer la condition de “non pléthoricité cardinale” pour $C_0$ :
\[
\parbox{0.72\linewidth}{$\forall\,i\in C_0$, désignant par $C_0(i)$ l'ensemble des $i'\in C_0$ telles que $i\simeq i'$ dans $\Fl(\M)$, j'impose que $\card C_0(i)\leq 2^{(r^{\pi_0})}$ (ce qui est aussi la valeur de la majoration pour le cardinal de $\wavy{\uFl^r(\M)}=\wavy{\uFl(\M^r)}$\,).} \tag*{(3.11)}
\]

\printedpage{17}
% head: A. Grothendieck : Dérivateurs, Ch. XV, Édition des Universités de Montpellier II et Paris VII
% foot: Version 5 novembre 2010
\texttt{[Deux fois (3.11).]} Cela implique donc que l'on a $\card C_0\leq(2^{(r^{\pi_0})})^2=2^{(r^{\pi_0})}$,
\[
\card C_0\ \leq\ 2^{(r^{\pi_0})}\ , \tag*{(3.12)}
\]
qui est équivalente à (3.11) (moyennant $C_0\subset\Fl^r(\M)$, postulé dans (3.10)).

Il s'ensuit par (3.8) que l'on a
\[
\card\Delta(C_0,f)\ \leq\ \underbrace{\card(C_0)}_{\leq 2^{(r^{\pi_0})}}\cdot s^{(r^{\pi_0})}\cdot s^{(r^{\pi_0})}\ \leq\ s^{(r^{\pi_0})}\ .
\]

D'autre part, les $X_D$ dont on prend la somme amalgamée sous $X$, pour obtenir $X_1(f)$, sont dans $\M^s$, si on fait l'hypothèse suivante. \msp{page 25}
\[
\left|\ \parbox{0.72\linewidth}{Les $\M^\pi=\uFilt^\pi$ sont stables par sommes amalgamées simples, et par $\varinjlim$ filtrantes de cardinal $\leq\pi$ (pas nécessairement grandes devant $\pi_0$) ,}\right. \tag*{(3.13)}
\]
compte tenu que $r\leq s$, donc $\M^r\subset\M^s$, donc $A$, $B$, $X$ dans $X_D=B\amalg_A X$ sont dans $\M^s$. Les sommes amalgamées partielles \textit{finies} des $X_D$ sont donc aussi dans $\M^s$, par récurrence sur le cardinal de l'ensemble d'indices. La somme amalgamée infinie apparaît ainsi comme $\varinjlim$ filtrante d'objets de $\M^s$, l'ensemble d'indices étant l'ensemble $\mathfrak{P}^*_{\mathfrak{f}}\Delta(C_0,f)$ des parties finies non vides de $\Delta(C_0,f)$. Cet ensemble est de cardinal majoré par $s'$, si $s'$ est un cardinal \textit{infini} tel que $\card\Delta(C_0,f)\leq s'$. On peut donc prendre $s'=s^{(r^{\pi_0})}$. On a donc prouvé, compte tenu de (3.13) appliqué à $\M^{s'}$, le % ?? uncertain: symbol in X_D = B \amalg_A X -- \amalg or a doubled \sqcup; subscript of \mathfrak{P}^* -- fraktur f or roman f

\textbf{Lemme 3.3.} \textit{On suppose} \textup{(3.10), i.e.} $C_0\subset\Fl^r(\M)$, $f\in\Fl^s(\M)$, \textit{avec} $s\geq r\geq\sup(\pi_1,C_1)$, \textit{et de plus que $C_0$ satisfait la condition de non-pléthoricité} \textup{(3.11)} \textup{(}\textit{ce qu'on peut toujours supposer, quitte à réduire $C_0$ sans changer $\widetilde{C_0}=\Phi$ et $C_{0*}=\Psi$}\textup{)}\textit{, et que la filtration cardinale satisfait la condition de stabilité (anodine!)} \textup{(3.13)}\textit{. Posons}
\[
s'\ =\ s^{(r^{\pi_0})}\ . \tag*{(3.14)}
\]
\textit{Alors $X_1(C_0,f)$ est dans $\M^{s'}$. Plus généralement et plus joliment, pour $C_0$ et $r$ fixés} \textup{(}\textit{avec $C_0\subset\Fl^r(\M)$ et satisfaisant la condition de non-pléthoricité} \textup{(3.11)} \textit{relativement à $r$}\textup{)}\textit{, pour tout cardinal} \msp{page 26} \textit{de la forme}
\[
s'\ =\ s^t\ ,\qquad\quad\textit{avec}\ t\geq r^{\pi_0},\ s\geq 2 \tag*{(3.15)}
\]
\textit{(d'où $s'\geq 2^{(r^{\pi_0})}>r^{\pi_0}\geq r$), $\uFl^{s'}(\M)$ est stable sous le foncteur $p_1^{C_0}:\uFl(\M)\longrightarrow\uFl(\M)$.}

En effet, si on applique la première partie du lemme au couple $(r,s')$, on trouve que $X_1(f)$ est dans $\M^{s''}$, avec $s''=s'^{(r^{\pi_0})}$. Mais
\[
s''\ =\ s'^{(r^{\pi_0})}\ =\ (s^t)^{(r^{\pi_0})}\ =\ s^{tr^{\pi_0}}\ \underset{\substack{\vert\\ \text{\scriptsize car $t\geq r^{\pi_0}$,}\\ \text{\scriptsize donc $tr^{\pi_0}=t$}}}{=}\ s^t\ =\ s'\ ,
\]
\printedpage{18}
% head: A. Grothendieck : Dérivateurs, Ch. XV, Édition des Universités de Montpellier II et Paris VII
% foot: Version 5 novembre 2010
donc on voit que $f\in\Fl^{s'}$ implique $X_1(f)\in\M^{s'}$, ou ce qui revient au même, $p_1(f)\in\Fl^{s'}$.

\textbf{Corollaire.} \textit{Soient $C_0$, $r$, $s'=s^t$ $(t\geq r^{\pi_0},\ s\geq 2)$ comme dessus, et soit $\alpha$ un nombre ordinal tel que}
\[
\card\alpha\ \leq\ s'\ =\ s^t\ .
\]
\textit{Alors $\uFl^{s'}(\M)$ est stable par le foncteur $p_\alpha=p_\alpha^{C_0}$.}

On le voit aussitôt par récurrence transfinie, compte tenu du lemme 3.3, et de l'hypothèse de stabilité (3.13).

Pour résumer :

\textbf{Théorème 2.} \textit{Soit $\M$ une catégorie stable par petites $\varinjlim$ filtrantes, et munie d'une filtration cardinale} $(\Filt^\pi)_{\pi\geq\pi_0}$\textit{, telle que les} $\Filt^\pi$ \textit{soient stables par $\varinjlim$ filtrantes de cardinal $\leq\pi$ (pas forcément grandes devant $\pi_0$), et par les sommes amalgamées simples qui existent dans $\M$. On suppose que les éléments de} $\Filt^{\pi_0}$ \textit{soient accessibles, et on désigne par $\pi_1$ le plus petit cardinal $\geq\pi_0$, tel que les objets de} $\Filt^{\pi_0}$ \textit{soient $\pi_1$-accessibles. On pose} \msp{page 27}
\[
c_1\ =\ \card\Fl(\M^{\pi_1})\ =\ \card\Fl^{\pi_1}(\M)
\]
\textit{(où pour tout $\pi\geq\pi_0$, $\M^\pi$ désigne la sous-catégorie strictement pleine de $\M$ définie par} $\Filt^\pi$\textit{, et} $\Fl^\pi(\M)=\Fl(\M^\pi)$\textit{).}

\dgrblock{dgrb733}{One fixed exponent controls every stage}
\begin{dcomment}{exposition}{p26}{One fixed exponent controls every stage}
The geometric calculation has now become very simple. If $C_0\subseteq\Fl(\mathcal M^r)$ and there are at most $2^\mu$ generators, where $\mu=r^{\pi_0}$ and $r\geq\chi$, then for an arrow with endpoints in $\mathcal M^s$, $s\geq r$, the set of attaching squares has cardinality at most $s^\mu$. Each individual pushout stays in $\mathcal M^s$. Their common-source amalgam can be obtained as a filtered colimit of finite partial amalgams, using at most $s^\mu$ indices. This proves Lemma~3.3 with the corrected counting constant. If there are no squares, the amalgam is simply the original source; this case is included by allowing the empty finite subset, or by treating it separately.

What matters for iteration is an exponent fixed point. If $S=c^t$ with $c\geq2$ and $t\geq\mu$, then $S^\mu=c^{t\mu}=c^t=S$. Thus the successor construction preserves $\mathcal M^S$. A limit stage indexed by an ordinal of cardinality at most $S$ stays there by (3.13). Transfinite induction proves that every stage up to any ordinal $\alpha$ with $|\alpha|\leq S$ has its middle object in $\mathcal M^S$.

To obtain the lifting property at the final stage, the ordinal must also be a limit of cofinality greater than $\mu$. The initial ordinal of $\mu^+$ is a uniform choice: $\mu^+\leq2^\mu\leq S$. The assertion in Theorem~2 that the initial ordinal of \emph{any} cardinal $r'>\mu$ will do is too broad as a justification of the cofinality condition; a singular $r'$ may have cofinality at most $\mu$. Choosing the successor cardinal, as the theorem itself suggests, repairs this point. The size estimate and the cofinality test perform different jobs and must both be retained.
\end{dcomment}


\textit{Soit $C_0$ une petite partie de} $\Fl(\M)$\textit{, formée de flèches coquarrables, et posons}
\[
\Phi\ =\ \widetilde{C_0}\ =\ \Psi^*\ ,\qquad\quad\textit{où}\ \ \Psi=C_{0*}=\Phi_*\ .
\]
\textit{Soit $r$ un cardinal} $\geq\sup(\pi_1,c_1)$\textit{, et tel que l'on ait $C_0\subset\Fl^r(\M)$. On suppose que $C_0$ est irrédondant (}i.e.\textit{ deux objets de $C_0$ isomorphes sont égaux), ou plus généralement, que $C_0$ satisfait la très anodine condition de non-pléthoricité} \textup{(3.11)} \textit{relativement à $r$. On considère la suite transfinie des foncteurs $p_\alpha=p_\alpha^{C_0}$ sur $\uFl(\M)$} \textup{(}\textit{cf.} \textup{th. 1)}\textit{, donnant lieu aux factorisations fonctorielles d'une $f\in\uFl(\M)$, $f:X\longrightarrow Y$, en}
\[
X\ \xrightarrow{i_\alpha(f)}\ X_\alpha(f)\ \xrightarrow{p_\alpha(f)}\ Y\ . \tag*{(3.16)}
\]

\textit{On a alors ce qui suit :}

$1^\circ$) \textit{On a (pour mémoire) $i_\alpha(f)\in\Phi$ pour tout $\alpha$, et de plus,} si $\alpha$ est grand devant $r^{\pi_0}$, \textit{on a aussi $p_\alpha(f)\in\Psi$. Si $r'$ est un cardinal $>r^{\pi_0}$ (p. ex. le cardinal successeur de $r^{\pi_0}$), on peut prendre pour $\alpha$ le premier ordinal dont le cardinal soit égal à $r'$.}

$2^\circ$) \textit{Pour tout cardinal $s$ de la forme}
\[
s\ =\ c^t\qquad\quad\textit{avec}\ \ t\geq r^{\pi_0},\ c\geq 2
\]
\textit{(donc $s=2^t$ si on prend $c\leq t$,} cf. \textit{lemme 3.2), $\uFl^s(\M)$ est stable sous les $p_\alpha$ tels que} $\card\alpha\leq s$\textit{. Si on prend p. ex. pour $\alpha$ le premier ordinal dont le cardinal soit $>r^{\pi_0}$, alors $\uFl^s(\M)$ est stable par $p_\alpha$, qui par $1^\circ$) réalise donc une factorisation fonctorielle} \textup{(3.16)} \textit{avec $i_\alpha(f)\in\Phi$, $p_\alpha(f)\in\Psi$.}

% ---- from Der-ChXV-05-11-10_p021-025.tex
\printedpage{19}
% head: A. Grothendieck : Dérivateurs, Ch. XV, Édition des Universités de Montpellier II et Paris VII
% foot: Version 5 novembre 2010

\msp{page 28}
\textsc{Démonstration}. Tout a été prouvé, si on se rappelle que les objets de $\mathcal{M}^r$ sont $r^{\pi_0}$-accessibles (\textit{cf.} (3.2)).

\dgrblock{dgrb734}{Keeping the multiplicity of generators in the estimate}
\begin{dcomment}{development}{p27}{Keeping the multiplicity of generators in the estimate}
The non-plethoricity hypothesis is convenient when one is free to choose a generating set. In a family of constructions, however, discarding isomorphic copies can obscure the given maps between generating data. We can keep those copies and make their cost explicit. The following is a quantitative completion of the calculation, with no claim of priority. We checked the cardinal-filtration argument in SGA~4 and the accessibility bounds for functorial factorizations in \citet[Theorem 4.4]{Low2016Heart}; neither is cited as asserting this exact formula with the raw multiplicity parameter.

\textbf{A uniform bound with multiplicities.}\ours
Assume the filtered-colimit, pushout and cardinal-filtration hypotheses of Section~3, including (3.13). Choose $\chi$ by counting the actual arrows of a small full skeleton of $\mathcal M^{\pi_1}$. Let $C_0$ be any small set of coquarrable arrows whose endpoints lie in $\mathcal M^r$, where $r\geq\chi$. Let $f$ have endpoints in $\mathcal M^s$, and let $\alpha$ be any small ordinal. Put
\[
\mu=r^{\pi_0},\qquad
q=\max\{2,r,s,|C_0|,|\alpha|\},\qquad B=q^\mu.
\]
Then $X_\beta(C_0,f)\in\mathcal M^B$ for every $\beta\leq\alpha$. If $\alpha$ is a limit with $\operatorname{cf}(\alpha)>\mu$, the terminal factor $p_\alpha(f)$ has the right lifting property against $C_0$.

\emph{Proof.} The cardinal $\mu$ is infinite, $B\geq q\geq r,s$, and
\[
B^\mu=q^{\mu\mu}=q^\mu=B.
\]
At stage zero the assertion follows from monotonicity of the filtration. Suppose it is true at $\beta$. The repaired Hom estimate applies to each endpoint of a generator and to each endpoint of $p_\beta(f)$, and gives at most $B^\mu=B$ possibilities for either side of an attaching square. There are therefore at most $|C_0|B^2\leq B$ squares. Each of their pushouts belongs to $\mathcal M^B$. A finite partial amalgam also belongs to $\mathcal M^B$, by induction on the number of its terms; the empty one is $X_\beta$. The poset of finite subsets of the set of squares is filtered and has cardinality at most $B$. Condition (3.13) puts its colimit, which is $X_{\beta+1}$, in the same level.

If $0<\lambda\leq\alpha$ is a limit, its diagram has cardinality $|\lambda|\leq|\alpha|\leq B$. Again (3.13) puts $X_\lambda=\varinjlim_{\beta<\lambda}X_\beta$ in $\mathcal M^B$. This completes the induction. Finally, a generator domain is $\mu$-accessible by (3.2), so any map from it to $X_\alpha$ factors through some $X_\beta$, $\beta<\alpha$. The prescribed map from the generator codomain to the fixed target then determines a square at that stage, solved at $\beta+1<\alpha$. Its image in $X_\alpha$ is the required lift.\hfill$\square$

In particular take $\alpha$ to be the initial ordinal of $\mu^+$. The formula gives an explicit answer to the size question (3.6) for the unaltered generating set. The extra parameter $|C_0|$ is necessary in this formulation: in $\mathrm{Set}$, arbitrarily many distinct copies of the isomorphism class $\varnothing\to\{*\}$ attach arbitrarily many points to $\varnothing\to\{*\}$ at the first stage, while the filtration levels of the generator endpoints remain fixed.
\end{dcomment}


\dgrblock{dgrb735}{From bounded factorizations to small model structures}
\begin{dcomment}{later}{p28:Tout a été prouvé, si on se rappelle}{From bounded factorizations to small model structures}
Theorem~2 addresses a question that remains useful long after the existence of the small object argument is understood: can a chosen factorization be performed inside a controlled small part of the ambient category? In the locally presentable setting, \citet[Theorem 4.4]{Low2016Heart} gives a functorial weak factorization system whose functors preserve specified presentable objects under explicit bounds on the generators, their Hom sets and two regular cardinals. His Theorem~5.9 then tests when two such systems form a model structure by restricting them to an appropriate small full subcategory; Proposition~5.12 gives the corresponding statement for an arbitrary combinatorial model category.

These results use regular presentability indices and the relation of one cardinal being sharply below another. The numerical dictionary is therefore not obtained by replacing every $\Filt^\pi$ with the modern class of $\pi$-presentable objects. What is common is the argument's architecture: control one stage, choose a cardinal stable under its operations, and keep the entire factorization within that bound. This is the immediate background for the model-structure criteria that Grothendieck develops in Section~4.
\end{dcomment}


\section*{4\quad Compléments sur les structures de Quillen}

\textbf{Définition 4.1.} Soit $\mathcal{M}$ une catégorie. On appelle \textit{précouple de Quillen} dans $\mathcal{M}$ un couple $(\Phi,\Psi)$ de parties de $\Fl(\mathcal{M})$, satisfaisant les deux conditions suivantes.
\begin{itemize}[leftmargin=2.5em]
\item[\textit{a})] On a $\Phi\longleftrightarrow\Psi$, \textit{i.e.} $\Psi\subset\Phi_*$, ou ce qui revient au même, $\Phi\subset\Psi^*$ -- soit encore : pour tout carré commutatif dans $\mathcal{M}$
\[
\begin{tikzcd}[row sep=large, column sep=large]
A \arrow[r, "u"] \arrow[d, "i"'] & X \arrow[d, "p"] \\
B \arrow[ur, dotted, "w" description] \arrow[r, "v"] & Y \;,
\end{tikzcd}
\tag*{(4.1.1)}
\]
avec $i\in\Phi$, $p\in\Psi$, il existe $w$ telle que $wi=u$, $pw=v$.
\item[\textit{b})] Toute $f\in\Fl(\mathcal{M})$ se factorise en
\[
f \;=\; p\,i\;, \qquad \text{avec } i\in\Phi,\, p\in\Psi\ .
\]
\end{itemize}
On écrit
\[
\Phi \;\qpre\; \Psi
\tag*{(4.1.2)}
\] % ?? uncertain: the printed relation sign -- a double arrow with a small vertical bar in the middle (rendered by \qpre)
pour indiquer que $(\Phi,\Psi)$ est un précouple de Quillen.

On dit que le précouple de Quillen est \textit{clos}, si $\Phi$ est clos à gauche, $\Psi$ clos à droit. On va voir plus bas (cor. 2, page 31) que cela implique (donc équivaut à) % [sic] clos à droit -> clos à droite, gender agreement with « à droite »
\[
\Phi \;=\; \Psi^*\;, \qquad \Psi \;=\; \Phi_*\ .
\tag*{(4.1.3)}
\]

\msp{page 29}
Un précouple de Quillen clos est aussi appelé un couple de Quillen.

\dgrblock{dgrb736}{From a factorization procedure to its two lifting classes}
\begin{dcomment}{exposition}{p28}{From a factorization procedure to its two lifting classes}
Grothendieck has constructed factorizations from a set of generators. He now asks which parts of that construction survive when the generators and their cardinal estimates are forgotten. This change of viewpoint explains the apparently elementary opening of Section~4: the next existence argument will need to recognize a model structure from two factorization procedures, and a factorization procedure does not by itself determine the classes of maps which ought to be called cofibrations and fibrations.

Write $L^{\square}$ for the maps having the right lifting property against every member of $L$, and ${}^{\square}R$ for the corresponding left lifting class. Grothendieck's notation is $L_*=L^{\square}$ and $R^*={}^{\square}R$. A \emph{pr\'ecouple} supplies factorizations and the inclusion $L\subset{}^{\square}R$. A \emph{couple} supplies the equalities
\[
 L={}^{\square}R,\qquad R=L^{\square}.
\]
Thus a couple is what is now called a weak factorization system. Neither uniqueness of the diagonal filler nor a functorial choice of factorization belongs to this definition. The constructions of Section~1 provide functoriality under their additional smallness hypotheses. Modern definitions and the retract argument can be compared directly in \citet[Definition~11.2.1 and Lemma~11.2.3]{Riehl2014}.

The word \emph{clos} concerns equality with a lifting class. It must be distinguished from closure under the operations listed later in (4.4.3). Those operations preserve a left lifting class, but the converse requires an existence theorem. Section~1 furnishes that converse for classes generated by a small set of sufficiently small maps. This distinction will become decisive when the class $C\cap W$ is known to have the expected stability properties but no generating set has yet been found.

The expression ``facteur direct'' also has a categorical meaning here. It is a retract in the arrow category, without any additive direct-sum decomposition. For arrows $j:A\to B$ and $i:A'\to B'$, it means that there are morphisms of arrows $j\to i\to j$ whose composite is the identity of $j$. Proposition~4.2 shows that a factorization of $j$ lets us keep its source fixed. That stronger form is why one diagonal filler is enough to turn a lifting condition into a retract statement.
\end{dcomment}


\textbf{Proposition 4.2.} \textit{Soit} $(\Phi,\Psi)$ \textit{un précouple de Quillen dans} $\mathcal{M}$\textit{, et soit} $j\in\Fl(\mathcal{M})$\footnote[9]{\textbf{N.B.} On utilise seulement que $\Phi\longleftrightarrow\Psi$, et que $j$ peut se mettre sous la forme $pi$, avec $i\in\Phi$, $p\in\Psi$.}. % footnote mark printed as "(9)"
\textit{On a alors les conditions équivalentes sur} $j$ \textit{:}
\begin{itemize}[leftmargin=2.5em]
\item[a)] $j\in\widetilde{\Phi}$.
\item[b)] $j\in\Psi^*$.
\item[c)] $j$ \textit{est facteur direct d'un} $i\in\Phi$.
\item[c$'$)] $j$ \textit{est facteur direct d'un} $i\in\Phi$ \textit{ayant même source que} $j$\textit{, plus précisément, il existe un diagramme}
\[
\begin{tikzcd}[row sep=large, column sep=large]
A \arrow[r, "j"] \arrow[d, shift right=1ex, "\id_A"'] & B \arrow[d, shift right=1ex, "q"'] \\
A \arrow[u, dotted, shift right=1ex, "\id_A"'] \arrow[r, "i"] & B' \;, \arrow[u, dotted, shift right=1ex, "p"']
\end{tikzcd}
\tag*{(4.2.1)}
\]
\printedpage{20}
% head: A. Grothendieck : Dérivateurs, Ch. XV, Édition des Universités de Montpellier II et Paris VII
% foot: Version 5 novembre 2010
\textit{avec} $i\in\Phi$\textit{, et}
\[
qj \;=\; i\;, \qquad pi \;=\; j\;, \qquad pq \;=\; \id_B\;,
\tag*{(4.2.2)}
\]
\textit{de sorte que} $j$ \textit{est un facteur direct de} $i\in\Phi$\textit{, par} $(\id_A,q)$\textit{, admettant la rétraction} $(\id_A,p)$.
\item[d)] \textit{Pour tout diagramme commutatif traits pleins}
\[
\begin{tikzcd}[column sep=small]
A \arrow[rr, "j"] \arrow[dr, "i"'] & & B \arrow[dl, dotted, bend left=40, "q"] \\
& B' \arrow[ur, "p"'] &
\end{tikzcd}
\;,
\tag*{(4.2.3)}
\]
i.e. \textit{toute factorisation de} $j$ \textit{en} $pi$\textit{, si on a} $p\in\Psi$\textit{, il existe une section} $q$ \textit{de} $B'$ \textit{sur} $B$ \textit{dans} $A\backslash\mathcal{M}$\textit{,} i.e. \textit{une flèche} $q:B\longrightarrow B'$ \textit{telle que}
\[
pq \;=\; \id_B \qquad \textit{et} \qquad qj \;=\; i\ .
\]
\end{itemize}

\textsc{Démonstration} de la proposition. Elle se fait par la suite d'implications suivante, toutes $(\pm)$ tautologiques :
\[
(c') \;\Longrightarrow\; (c) \;\Longrightarrow\; (a) \;\Longrightarrow\; (b) \;\Longrightarrow\; (d) \;\Longrightarrow\; (c')\;,
\]
\msp{page 30}
dont les quatre premières utilisent seulement l'hypothèse $\Phi\longleftrightarrow\Psi$ (de $(a)$ de la définition 4.1), et la dernière seulement, $(d)\Longrightarrow(c')$, fait appel à l'hypothèse de factorisation de Quillen ($(b)$ dans 4.1) ; et encore il suffit de la supposer pour $j$ seulement !

$(c')\Longrightarrow(c)$ tautologie pure, $(c)\Longrightarrow(a)$ car on sait que $\widetilde{\Phi}$ est stable par facteurs directs (et contient $\Phi$), $(a)\Longrightarrow(b)$ car $\widetilde{\Phi}\subset\Psi^*$ ($\widetilde{\Phi}$ étant par définition la plus petite partie close à gauche de $\Fl(\mathcal{M})$ contenant $\Phi$, et $\Psi^*$ étant close à gauche), $(b)\Longrightarrow(d)$ car on écrit le diagramme donné (4.2.2) (traits pleins) comme % [sic] (4.2.2) -> (4.2.3), the diagram given in condition d) is (4.2.3)
\[
\begin{tikzcd}[row sep=large, column sep=large]
A \arrow[r, "u"] \arrow[d, "j\in\Psi^*"'] & X \arrow[d, "p\in\Psi"] \\
B \arrow[ur, dotted, "q" description] \arrow[r, "v"] & Y \;,
\end{tikzcd}
\tag*{$(*)$}
\]
en y faisant $Y=B$, $v=\id_B$, et en rebaptisant $X$ par $B'$, et $u$ par $i$ : l'hypothèse $(b)$, $j\in\Psi^*$, signifie que dans un diagramme traits pleins $(*)$, il existe toujours $q$ telle que $qj=u$, $pq=v$, et en l'exprimant dans le cas particulier où $v=\id_B$, on trouve $(d)$.

Reste à prouver $(d)\Longrightarrow(c')$, en utilisant cette fois l'hypothèse de factorisation sur $j$, \textit{i.e.} qu'on peut écrire $j$ sous la forme $pi$ :
\[
\begin{tikzcd}[column sep=small]
A \arrow[rr, "j"] \arrow[dr, "i"'] & & B \arrow[dl, dotted, bend left=40, "q"] \\
& B' \arrow[ur, "p"'] &
\end{tikzcd}
\;,
\]
avec $i\in\Phi$, $p\in\Psi$. En vertu de $(d)$, comme $p\in\Psi$, il existe donc $q:B\longrightarrow B'$ satisfaisant
\printedpage{21}
% head: A. Grothendieck : Dérivateurs, Ch. XV, Édition des Universités de Montpellier II et Paris VII
% foot: Version 5 novembre 2010
\msp{page 31}
\[
pq \;=\; \id_B\;, \qquad qj \;=\; i\;,
\]
ce qui prouve $(c')$, puisque $i\in\Phi$.

\dgrblock{dgrb737}{The retraction which completes both classes}
\begin{dcomment}{exposition}{p30}{The retraction which completes both classes}
The role of Proposition~4.2 is to identify exactly what can be missing from a pr\'ecouple. Suppose $j:A\to B$ lifts against $R$, and choose a factorization $j=pi$ with $i:A\to Z$ in $L$ and $p:Z\to B$ in $R$. Apply the lifting property to
\[
\begin{tikzcd}[row sep=large,column sep=large]
A \arrow[r,"i"] \arrow[d,"j"'] & Z \arrow[d,"p"]\\
B \arrow[r,"\id_B"'] \arrow[ur,dashed,"q"] & B .
\end{tikzcd}
\]
The equations are $qj=i$ and $pq=\id_B$. Hence $(\id_A,q):j\to i$ and $(\id_A,p):i\to j$ exhibit $j$ as a retract of $i$. Conversely, a lifting property passes to retracts: extend a proposed lifting square along the section, solve it for the larger arrow, and compose the resulting filler with the retraction. These two observations prove
\[
{}^{\square}R=\operatorname{Retr}(L),\qquad
L^{\square}=\operatorname{Retr}(R)
\]
whenever $(L,R)$ is a pr\'ecouple. The dual assertion uses a factorization with target fixed.

This proof requires a factorization only of the particular arrow $j$ under consideration, as footnote~9 correctly emphasizes. There is no induction hidden in the argument and no completeness hypothesis on the category. Once all arrows have factorizations, closing the two classes under retracts therefore produces a weak factorization system without changing any chosen factorization. This is the precise content of the two corollaries on manuscript page~31.

It also explains a useful economy in model-category arguments. One first constructs maps in classes with convenient explicit descriptions, such as relative cell complexes. One then enlarges those classes by retracts. The enlargement does not require repeating the construction: lifting supplies the necessary retraction. The enlargement is nevertheless real, and omitting it from a recognition statement can change that statement's truth value.
\end{dcomment}


\textbf{Corollaire 1.} \textit{Soit} $(\Phi,\Psi)$ \textit{un précouple de Quillen. Alors on a}
\[
\begin{aligned}
\widetilde{\Phi} \;\;=\;\; \Psi^* \;\;=\;\; &\textit{ensemble des objets de } \underline{\Fl}(\mathcal{M}) \textit{ qui sont fac-}\\
&\textit{teurs directs d'éléments de } \Phi\ ,
\end{aligned}
\tag*{(4.3)}
\] % printed line break of the display falls inside "fac-teurs"; kept as printed
\textit{et dualement}
\[
\begin{aligned}
\overline{\Psi} \;\;=\;\; \Phi_* \;\;=\;\; &\textit{ensemble des objets de } \underline{\Fl}(\mathcal{M}) \textit{ qui sont fac-}\\
&\textit{teurs directs d'éléments de } \Psi\ .
\end{aligned}
\tag*{(4.4)}
\] % printed line break of the display falls inside "fac-teurs"; kept as printed

\textbf{Corollaire 2.} \textit{Les conditions suivantes sur le précouple de Quillen} $(\Phi,\Psi)$ \textit{sont équivalentes.}
\begin{itemize}[leftmargin=2.5em]
\item[a)] \textit{Le précouple est} clos, i.e. $\Phi=\widetilde{\Phi}$\textit{,} $\Psi=\overline{\Psi}$\textit{.}
\item[b)] $\Phi=\Psi^*$\textit{,} $\Psi=\Phi_*$\textit{.}
\item[c)] $\Phi$ \textit{et} $\Psi$ \textit{sont stables par facteurs directs.}
\end{itemize}
Soit maintenant
\[
W \;\subset\; \Fl(\mathcal{M})
\tag*{(4.5)}
\]
un localiseur, satisfaisant les conditions habituelles
\[
\left\{
\begin{aligned}
&\textit{a})\quad W \text{ contient les isomorphismes.}\\[1ex]
&\textit{b})\quad \text{Si } X\xrightarrow{\;f\;}Y\xrightarrow{\;g\;}Z \text{ dans } \mathcal{M}\text{, si deux parmi } f,\ g,\ gf \text{ sont}\\
&\phantom{\textit{b})}\quad \text{dans } W\text{, alors le troisième aussi.}
\end{aligned}
\right.
\tag*{(4.6)}
\]
($W$ modérément saturé). Supposons de plus données quatre autres parties
\[
C\ ,\ TC\ , \qquad F\ ,\ TF
\tag*{(4.7)}
\]
de $\Fl(\mathcal{M})$, satisfaisant les conditions suivantes.
\[
\left\{
\begin{aligned}
&\textit{a})\quad (C,TF) \text{ et } (TC,F) \text{ sont des précouples de Quillen.}\\[1ex]
&\textit{b})\quad TC\subset C\cap W,\ TF\subset F\cap W.
\end{aligned}
\right.
\tag*{(4.8)}
\]

\msp{page 32}
Ces propriétés s'explicitent ainsi :
\begin{itemize}[leftmargin=2.5em]
\item[\textbf{M\,1}] Propriétés de relèvement-prolongement, $C\longleftrightarrow TF$ et $TC\longleftrightarrow F$.
\item[\textbf{M\,2}] Propriété de factorisation de toute flèche $f$ de $\mathcal{M}$ en $f=pi$, avec, au choix, $i\in C$, $p\in TF$, ou $i\in TC$, $p\in F$.
\item[\textbf{M$'$}] $TC\subset C\cap W$, $TF\subset F\cap W$.
\end{itemize}

\printedpage{22}
% head: A. Grothendieck : Dérivateurs, Ch. XV, Édition des Universités de Montpellier II et Paris VII
% foot: Version 5 novembre 2010
C'est donc une partie des axiomes d'une structure de Quillen faible pour $W$ (XIII, p. 27, déf. 1) -- il manque les axiomes de stabilité M\,3, M\,4, et les axiomes de “propreté” (propriety) M$''$, savoir $F\subset\mathrm{Fib}_W$, $C\subset\mathrm{Cof}_W$.

\textbf{Proposition 4.3.} \textit{Soient} $\mathcal{M}$ \textit{une catégorie,} $(C,TF)$ \textit{et} $(TC,F)$ \textit{deux précouples de Quillen dans} $\mathcal{M}$\textit{, satisfaisant}
\[
TC \;\subset\; C\;, \qquad TF \;\subset\; F\;,
\tag*{(4.3.1)}
\]
\textit{et soit} $W$ \textit{une partie de} $\Fl(\mathcal{M})$ \textit{dont on ne suppose pour l'instant que} 4.6, (b)\footnote[10]{On utilise seulement que si $X\xrightarrow{\;f\;}Y\xrightarrow{\;g\;}Z$, et si $gf\in W$, alors $f\in W\Longleftrightarrow g\in W$. On n'a pas besoin que $W$ contienne les isomorphismes, ni qu'il soit stable par composition.}\textit{, et} % footnote mark printed as "(10)"
\[
W \;\supset\; TC\ ,\ TF\ ,
\tag*{(4.3.2)}
\]
\textit{donc on a}
\[
TC \;\subset\; C\cap W\ , \qquad TF \;\subset\; F\cap W\ .
\tag*{(4.3.3)}
\]
\textit{Considérons les conditions de prolongement-relèvement suivantes} (a), (a$'$)\textit{, et les conditions de factorisation} (b), (b$'$), (c), (d) (\textit{en se rappelant que} $\widetilde{TC}=F^*$\textit{,} $\overline{TF}=C_*$\textit{, en vertu de la proposition} 4.2 (\textit{cor.} 1) \textit{appliquée aux précouples de Quillen} $(TC,F)$\textit{,} $(C,TF)$) :
\begin{itemize}[leftmargin=2.5em]
\item[a)] $C\cap W\subset\widetilde{TC}$\quad (\textit{donc} $TC\subset C\cap W\subset\widetilde{TC}$).
\item[a$'$)] $F\cap W\subset\overline{TF}$\quad (\textit{donc} $TF\subset F\cap W\subset\overline{TF}$).
\item[b)] $W\subset TF\circ\widetilde{TC}$.
\item[b$'$)] $W\subset\overline{TF}\circ TC$.
\item[c)] $W\subset\overline{TF}\circ\widetilde{TC}$.
\item[d)] $W\subset TF{\circ}TC$ (\textit{donc, par} (4.3.2)\textit{,} $W=TF{\circ}TC$ \textit{si} $W$ \textit{est stable par composition}).\footnote[11]{\textbf{N.B.} La condition $(a)$ signifie aussi que $(C\cap W,F)$ est encore un précouple de Quillen, resp. $(b)$ signifie que $(C,F\cap W)$ en est un -- et alors cette paire de précouples de Quillen satisfait cette fois à $(d)$, \textit{i.e.} $W=TF'\circ TC'$.} % footnote mark printed as "(11)"
\end{itemize}

\msp{page 33}
\textit{Alors toutes les conditions sauf la dernière sont équivalentes} (\textit{et elles sont donc impliquées par cette dernière}). \textit{Donc lorsque} $TC$ \textit{ou} $TF$ \textit{sont stables par facteurs directs} (\textit{p. ex. si} $TC$ \textit{est clos à gauche, ou} $TF$ \textit{clos à droite}), \textit{toutes les conditions envisagées sont équivalentes} (\textit{puisque par le corollaire} 2 \textit{à la proposition} 4.2\textit{, on a alors} $TC=\widetilde{TC}$ \textit{resp.} $TF=\overline{TF}$\textit{, donc} (b$'$) \textit{resp.} (b) \textit{équivaut alors à} (d)).

\textsc{Démonstration}. Elle se fait par le diagramme d'implications
\[
\begin{tikzcd}[column sep=small, row sep=small]
& (d) \arrow[dl, Rightarrow] \arrow[dr, Rightarrow] & \\
(b) \arrow[dr, Rightarrow] & \mathrm{Q} & (b') \arrow[dl, Rightarrow] \\
& (c) \arrow[dl, Rightarrow] \arrow[dr, Rightarrow] & \\
(a) \arrow[uu, Rightarrow] & & (a') \arrow[uu, Rightarrow]
\end{tikzcd}
\;,
\]
\printedpage{23}
% head: A. Grothendieck : Dérivateurs, Ch. XV, Édition des Universités de Montpellier II et Paris VII
% foot: Version 5 novembre 2010
où les quatre implications du carré supérieur Q sont tautologiques. Il suffit donc d'établir les implications
\[
(c) \;\Longrightarrow\; (a) \;\Longrightarrow\; (b)\;,
\]
les implications $(c)\Longrightarrow(a')\Longrightarrow(b')$ en résultent par dualité.

Supposons donc $(c)$, \textit{i.e.}
\[
W \;\subset\; \overline{TF}\circ\widetilde{TC}\;,
\tag*{$(*)$}
\]
et prouvons $(a)$, \textit{i.e.} que
\[
C\cap W \;\subset\; \widetilde{TC}\ .
\]
Mais si $j\in C\cap W$, pour prouver que $j\in\widetilde{TC}$, écrivons $j=pi$, $i\in\widetilde{TC}$, $p\in\overline{TF}$, ce qui est possible par l'hypothèse $(*)$ et $j\in W$.

\msp{page 34}
\[
\begin{tikzcd}[column sep=small]
A \arrow[rr, "j\in C"] \arrow[dr, "i\in\widetilde{TC}"'] & & B \arrow[dl, dotted, bend left=60, "q"] \\
& B' \arrow[ur, "p\in\overline{TF}"'] &
\end{tikzcd}
\;.
\]
Utilisons maintenant l'hypothèse $j\in C$ et $C\longleftrightarrow\overline{TF}$, on trouve donc (prop. 3.2, $(a)\Longrightarrow(d)$) une flèche $q:B\longrightarrow B'$ telle que $pq=\id_B$, $qj=i$, donc $j$ apparaît comme facteur direct de $i$. Comme $i\in\widetilde{TC}$ et que $\widetilde{TC}$ est stable par facteurs directs, on en conclut $j\in\widetilde{TC}$, \quad \textit{q.e.d.}

$(a)\Longrightarrow(b)$. On suppose donc
\[
C\cap W \;\subset\; \widetilde{TC}\;,
\tag*{$(*)$}
\]
et on veut prouver $(b)$, \textit{i.e.}
\[
W \;\subset\; TF\circ\widetilde{TC}\ .
\]
Soit donc $f\in W$, et factorisons $f$ en $pi$, avec $i\in C$, $p\in TF$. (O.K. à cause du précouple de Quillen $(C,TF)$). Comme $TF\subset W$, on a $p\in W$, et comme $pi\in W$, on en conclut $i\in W$ (saturation modérée de $W$ -- on doit quand même l'utiliser !). On a donc $i\in C{\cap}W$, donc $i\in\widetilde{TC}$ par $(*)$, d'où $f\in TF\circ\widetilde{TC}$, \quad \textit{q.e.d.}

\textbf{Corollaire 1.} \textit{Supposons} $TC$ \textit{stable par facteurs directs. Alors toutes les conditions de la proposition} 4.3 \textit{sont équivalentes, et équivalent à la relation}
\[
C\cap W \;=\; TC\ .
\tag*{(4.3.4)}
\]
\textit{Dualement, si} $TF$ \textit{est stable par facteurs directs, les conditions de la proposition} 4.3 \textit{sont toutes équivalentes et équivalent à}
\[
F\cap W \;=\; TF\ .
\tag*{(4.3.4$'$)}
\]

% ---- from Der-ChXV-05-11-10_p026-030.tex
\printedpage{24}
% head: A. Grothendieck : Dérivateurs, Ch. XV, Édition des Universités de Montpellier II et Paris VII
% foot: Version 5 novembre 2010

\msp{page 35}
En effet, par le corollaire 2 de la proposition 4.2, l'hypothèse sur $TC$ implique que $TC = \widetilde{TC}$, donc comme $TC \subset C{\cap}W$, l'hypothèse $C{\cap}W \subset \widetilde{TC}$ $(= TC)$ équivaut à $C{\cap}W = TC$.

\dgrblock{dgrb738}{Recognition by two compatible factorizations}
\begin{dcomment}{exposition}{p34}{Recognition by two compatible factorizations}
Grothendieck wants to describe weak equivalences through their two possible factorizations. The names $TC$ and $TF$ are provisional at this stage: the assumptions only give $TC\subset C\cap W$ and $TF\subset F\cap W$. Proposition~4.3 measures what must be added before those inclusions become the desired equalities.

The proof has two different mechanisms. A factorization
\[
j=pi,\qquad j\in C\cap W,\quad
i\in\widetilde{TC},\quad p\in\overline{TF}
\]
exhibits $j$ as a retract of $i$, because $C$ lifts against $\overline{TF}=C_*$. This gives $C\cap W\subset\widetilde{TC}$ without using a property of $W$ beyond membership of $j$ in it. In the other direction, factor an arbitrary $w\in W$ using $(C,TF)$. If $w=pi$, then $p\in TF\subset W$; cancellation in $W$ forces $i\in C\cap W$. This is the sole use of the two-out-of-three axiom in that half of the argument. The dual proof exchanges the two factorizations.

Consequently, the equivalences among (a), (a$'$), (b), (b$'$), and (c) need only the restricted cancellation condition stated in footnote~10: whenever $gf\in W$, membership of $f$ in $W$ is equivalent to membership of $g$ in $W$. A displayed equality $W=TF\circ TC$ additionally uses closure of $W$ under composition for the reverse inclusion. The source itself makes this qualification beside (d). There is no need to strengthen the main proposition to a saturation condition involving localization.

The subsequent corollaries must be read with their respective closure hypotheses. If $TC$ is retract closed, then $\widetilde{TC}=TC$, so the recognition condition becomes $C\cap W=TC$. The analogous statement for $TF$ needs retract closure of $TF$. If $C$ and $W$ are retract closed, there is an even shorter proof of Corollary~2 than the printed one: $TC\subset C\cap W$ gives
\[
\widetilde{TC}=\operatorname{Retr}(TC)\subset C\cap W,
\]
and condition (a) supplies the reverse inclusion. The printed relation $C\cap W\longleftrightarrow TC$ in this proof is a slip. The lifting relation available here is $C\cap W\longleftrightarrow F$, and the relevant pr\'ecouple is $(C\cap W,F)$.
\end{dcomment}


\textbf{Corollaire 2.} \textit{Supposons $C$ et $W$ stables par facteurs directs. Alors les conditions équivalentes de la proposition} 4.3 \textit{équivalent à}
\[
C\cap W \;=\; \widetilde{TC}\ . \tag*{(4.3.5)}
\]
\textit{Dualement, si $F$, $W$ sont stables par facteurs directs, les conditions équivalentes de la proposition} 4.3 \textit{équivalent à}
\[
F\cap W \;=\; \overline{TF}\ . \tag*{(4.3.5$'$)}
\]

En effet, si $C$ et $W$ sont stables par facteurs directs, il en est de même de $C\cap W$, et dès lors si $C\cap W\subset\widetilde{TC}$, comme $C\cap W$ contient $TC$, on en conclut $\widetilde{C\cap W}=\widetilde{TC}$. Mais comme $C\cap W\longleftrightarrow TC$ et $C\cap W$ est stable par facteurs directs, $C\cap W$ est clos à gauche par corollaire 2 de la proposition 4.2, \textit{i.e.} $C\cap W=\widetilde{C\cap W}$, d'où $C\cap W=\widetilde{TC}$, \quad \textit{q.e.d.}

\textbf{Corollaire 3.} \textit{Supposons les conditions équivalentes} (a) \textit{à} (c) \textit{de} 4.3 \textit{satisfaites, et de plus} $W$ stable par facteurs directs\footnote[12]{En fait, il suffit de supposer que l'on a $W\supset\widetilde{TC}$, $\overline{TF}$, \textit{i.e.} que tout facteur direct d'un élément de $W$ qui est soit dans $C$, soit dans $F$, est encore dans $W$. Alors $(\widetilde{C},\widetilde{TC},\bar{F},\overline{TF})$ définissent une structure de Quillen close, et on sait que cela implique qu'en fait, $W$ lui aussi est stable par facteurs directs.}. \textit{Alors les conditions préliminaires} (4.3.1) \textit{et} (4.3.2) \textit{restent satisfaites quand on remplace $(C,TF)$ et $(TC,F)$ par leurs clôtures $(C',TF')=(\widetilde{C},\overline{TF})$ et $(TC',F')=(\widetilde{TC},\bar{F})$, et les conditions équivalentes de la proposition} 4.3 \textit{sont encore satisfaites pour ces couples, sous les formes fortes}
\[
\left\{
\begin{aligned}
C'\cap W &\;=\; TC' && (=\widetilde{TC'}\,)\\
F'\cap W &\;=\; TF' && (=\overline{TF'}\,)\\
W &\;=\; TF'\circ TC'\ .
\end{aligned}
\right.
\tag*{(4.3.6)}
\]

\msp{page 36}
\textit{Le triple $(W,C',F')$ est un triple de Quillen clos $($correspondant à une “closed model category” de Quillen, du moins si on suppose $\mathcal{M}$ stable par $\varinjlim$ finies et $\varprojlim$ finies, et que $W$ est modérément saturé$)$.}\footnote[13]{Il en résulte le\par
\textbf{Corollaire 4.} \textit{Soient $W$ un localiseur modérément saturé, et $(C,TF)$, $(TC,F)$ deux précouples de Quillen tels que $TC\subset C\cap W$, $TF\subset F\cap W$. Alors les conditions suivantes sont équivalentes.}
\begin{enumerate}[leftmargin=*]
\item[a)] $TC=C\cap W$.
\item[b)] $TF=F\cap W$.
\item[c)] $W=TF\circ TC$.
\item[d)] $(W,C,F)$ \textit{est un triple de Quillen clos.}
\end{enumerate}}

\textsc{Démonstration.} Les relations (4.3.1) impliquent
\[
\widetilde{TC}\;\subset\;\widetilde{C}\qquad\text{et}\qquad\overline{TF}\;\subset\;\bar{F}\ ,
\]
\printedpage{25}
% head: A. Grothendieck : Dérivateurs, Ch. XV, Édition des Universités de Montpellier II et Paris VII
% foot: Version 5 novembre 2010
\textit{i.e.}
\[
TC'\;\subset\;C'\ ,\qquad TF'\;\subset\;F'\ ,
\]
\textit{i.e.} (4.3.1) pour les nouveaux couples de Quillen. La relation 4.3.2, $W\supset TC,TF$, implique
\[
W\;\supset\;\widetilde{TC}\ ,\ \overline{TF}\ ,\qquad \textit{i.e.}\qquad W\;\supset\;TC'\ ,\ TF'\ ,
\]
en vertu du corollaire 2 à la proposition 4.2 et du fait que $W$ est stable par facteurs directs ; donc on a (4.3.2) pour les nouveaux couples. Comme par hypothèse on a
\[
W\;=\;\overline{TF}\circ\widetilde{TC}\ ,
\]
ceci s'écrit
\[
W\;=\;TF'\circ TC'\ ,
\]
donc les nouveaux couples satisfont aux mêmes conditions, qui cette fois, les couples étant clos, prennent la forme forte (4.3.6).

Prouvons que le triple $(W,C',F')$ satisfait aux axiomes de Quillen d'un triple de Quillen clos. On a bien, par (4.3.6),
\[
\left\{
\begin{aligned}
C' &\;=\;(TF')^*\;=\;(F'\cap W)^*\ , &\qquad F' &\;=\;(TC')_*\;=\;(C'\cap W)_*\ ,\\
W &\;=\;TF'\circ TC'\;=\;C'_*\circ F'^*\ ,
\end{aligned}
\right.
\]
qui correspond à la définition des catégories de modèles de Quillen \textit{clos} : les axiomes M\,1, M\,2, M\,3, M\,4 de Quillen sont réalisés de façon triviale (M\,3 et M\,4 parce que
\msp{page 37}
$C$ et $F$, $TC$ et $TF$ sont clos à gauche resp. à droite). L'axiome M\,5 est la saturation modérée de $W$.

\dgrblock{dgrb739}{The missing closure hypothesis in the footnote}
\begin{dcomment}{development}{p36}{The missing closure hypothesis in the footnote}
Grothendieck's main closure result, Corollary~3, deliberately replaces the four provisional classes by their retract closures. Footnote~13 then proposes a shorter Corollary~4 in which that replacement disappears. As printed, its equivalence with a \emph{closed} Quillen structure is false. The following example isolates the defect while retaining invariance of the classes under isomorphism; it is not an objection based merely on a choice of representatives.

\textbf{A pr\'ecouple whose two classes have not both been closed.}\ours
In the category of sets, with choice, let $W$ be all maps. Let $L$ consist of the injections $i:A\hookrightarrow B$ for which $B\setminus i(A)$ is empty or infinite, and let $R$ consist of all surjections. Set
\[
C=TC=L,\qquad F=TF=R.
\]
Then both required pairs are pr\'ecouples, and conditions (a), (b), and (c) of the printed Corollary~4 hold, but $(W,C,F)$ is not closed.

To verify the lifting condition, consider an injection $i:A\hookrightarrow B$ and a surjection $p:X\twoheadrightarrow Y$. A commutative square specifies the filler on $i(A)$; on each remaining element of $B$, choose an inverse image under $p$ of its prescribed image in $Y$. This proves $L\longleftrightarrow R$. For factorization, take any $f:A\to B$ and form
\[
A\ \longrightarrow\ A\amalg(B\times\mathbb N)
  \ \xrightarrow{\ [f,\operatorname{pr}_B]\ }\ B.
\]
The first map has empty complement when $B=\varnothing$ and infinite complement otherwise; the second is surjective. Thus it is an $(L,R)$-factorization. Since $W$ contains every map, the equalities $TC=C\cap W$ and $TF=F\cap W$ hold, and the factorization just given proves $W=TF\circ TC$.

On the other hand, $\varnothing\to\{0\}$ is a retract of $\varnothing\to\mathbb N$: use the inclusion selecting $0$ and the unique map $\mathbb N\to\{0\}$. The larger map belongs to $L$, whereas the smaller one does not. Therefore $L$ is not retract closed and cannot equal ${}^{\square}R$. This proves the failure of (d).

The appropriate repair of the footnote is to require the two input pairs to be \emph{couples}, or equivalently to require all four classes to be retract closed. Proposition~4.3 then gives the equivalence of (a), (b), and (c), and the resulting two weak factorization systems give (d), subject to the required existence of limits and colimits. Another valid conclusion under the weaker hypotheses is the one actually stated in Corollary~3: take retract closures first. The general modern formulation, including automatic retract closure of $W$ once the two weak factorization systems exist, is in \citet[Definition~11.3.1 and Lemma~11.3.3]{Riehl2014}.

We checked those general criteria and the recognition theorems cited below. We did not find this particular isomorphism invariant example in those sources; the claim here is the explicit counterexample and its verification, not a claim of priority for the retract obstruction.
\end{dcomment}


\textbf{4.4.} Voici l'idée pour construire plus ou moins “toutes” les structures de Quillen closes $(C,TC,F,TF)$ associées à un localiseur donné $W\subset\mathcal{M}$. On suppose que $\mathcal{M}$ satisfait aux conditions du paragraphe précédent, concernant l'existence de $\varinjlim$ filtrantes, d'une filtration cardinale, l'accessibilité dans $\mathcal{M}$ (celle des éléments de $\mathrm{Filt}^{\pi_0}$). On part d'une partie petite
\[
C_0\;\subset\;\Fl(\mathcal{M})
\]
telle que
\[
C_{0*}\;\subset\;W\ , \tag*{(4.4.1)}
\]
et on pose
\[
\left\{
\begin{aligned}
C &\;=\;\widetilde{C_0}\ , &\qquad TC &\;=\;C\cap W\ ,\\
TF &\;=\;C_*\;=\;C_{0*}\ , &\qquad F &\;=\;(TC)_*\ .
\end{aligned}
\right.
\tag*{(4.4.2)}
\]
Il y a un problème seulement, vu que l'on sait déjà que $(C,TF)$ est un couple de Quillen par le théorème de factorisation dans §1 (p. 10) : il faut la factorisation pour le couple $(TC,F\;=\;(TC)_*)$. À ce moment là, les conditions préliminaires (4.3.1), (4.3.2) de la proposition 4.3 sont satisfaites :
\[
TC\;\subset C\ ,\qquad TF\;\subset F\ ,\qquad W\;\supset\;TC\ ,\ TF\ ,
\]
ainsi que la condition $(a)$, $C\cap W\subset\widetilde{TC}$, puisque $C\cap W=TC$. Donc on trouve, par l'équivalence des conditions $(a)$, $(a')$, $(c)$, et vu que $TC$ et $TF$ sont clos,
\[
TF\;=\;F\cap W\ ,\qquad W\;=\;TF\circ TC\ ,
\]
\printedpage{26}
% head: A. Grothendieck : Dérivateurs, Ch. XV, Édition des Universités de Montpellier II et Paris VII
% foot: Version 5 novembre 2010
\msp{page 38}
d'autre part l'hypothèse sur $W$ implique que $W$ est stable par facteurs directs, $C$ l'est aussi, étant close à gauche, donc $TC=C{\cap}W$ aussi, donc par proposition 4.1, corollaire 2, on a $TC\;=\;\widetilde{TC}$, $TC$ est close à gauche. On gagne donc, par la dernière assertion du corollaire 3, proposition 4.3 (p. 35).

En fait, indépendamment des conditions particulières sur $\mathcal{M}$, $W$, on vient de prouver ceci :

\textbf{Lemme 4.4.1} (corollaire de la proposition 4.3.). \textit{Soient $\mathcal{M}$ une catégorie, $W$ un localiseur dans $\mathcal{M}$ modérément saturé et stable par facteurs directs. Alors les structures de Quillen closes $(C,F)$ associées à $W$ s'obtiennent ainsi : on prend une partie $C_0$ de} $\Fl(\mathcal{M})$ \textit{telle que}
\[
TF\;\deq\;C_{0*}\;\subset\;W\ ,
\]
\textit{on pose}
\[
C\;=\;\widetilde{C_0}\;=\;TF^*\ ,\qquad TC\;=\;C\cap W\ ,\qquad F\;=\;(TC)_*\ ,
\]
\textit{et on doit se débrouiller pour que $(C,TF)$ et $(TC,F)$ satisfassent à la condition de factorisation de Quillen,} i.e. \textit{soient des précouples de Quillen.} (En fait, $C$, $TF$ et $F$ étant clos par construction, et $TC=C\cap W$ étant stable par facteurs directs ($C$ et $W$ l'étant), ce seront même des \textit{couples} de Quillen.)

\msp{page 39}
Les hypothèses supplémentaires faites tantôt sur $\mathcal{M}$ servaient à assurer que si je prends $C_0$ petite, alors la factorisation marche pour $(C,TF)$ automatiquement. Il reste le problème pour $TC$, $F$.

Notons que si elle marche, alors $TC=\widetilde{TC}=F^*$, donc $TC$ est clos à gauche, et \textit{a fortiori} satisfait les conditions de stabilité correspondantes :
\[
\left\{
\begin{tabular}{rp{0.6\textwidth}}
\textit{a}) & Stable par cochangement de base, par composition, par facteurs directs (ces deux dernières sont assurées automatiquement, car $C$, $W$ sont stables).\\[1ex]
\textit{b}) & Stable par $\varinjlim$ ordinales strictes.
\end{tabular}
\right.
\tag*{(4.4.3)}
\]

L'hypothèse de stabilité sur $W$ sert (entre autres) à garantir la stabilité $(b)$ pour $TC$ (en fait, il suffirait pour cela seulement que $W$ soit stable par $\varinjlim$ ordinales strictes). Il reste la stabilité de $TC=C{\cap}W$ par cochangement de base, qui est problématique car $W$ n'est en général \textit{pas} stable par cochangement de base. C'est pour assurer cette stabilité que je me vois pratiquement contraint ici à imposer
\[
C_0\;\subset\;\Cof_W\ , \tag*{(4.4.4)}
\]
je dis qu'avec l'hypothèse de stabilité faite sur $W$, cela implique
\[
C\;\deq\;\widetilde{C_0}\subset\Cof_W\ . \tag*{(4.4.5)}
\]
\printedpage{27}
% head: A. Grothendieck : Dérivateurs, Ch. XV, Édition des Universités de Montpellier II et Paris VII
% foot: Version 5 novembre 2010
\msp{page 40}
En effet, on a :

\dgrblock{dgrb740}{Why preservation under pushout is the right condition}
\begin{dcomment}{exposition}{p39}{Why preservation under pushout is the right condition}
The second factorization cannot be obtained merely by declaring $TC=C\cap W$. Its left class must survive the operations used in the small object construction. Composition and retracts are already available, but an arbitrary weak equivalence need not remain a weak equivalence after pushout. Grothendieck's choice $C_0\subset\Cof_W$ makes the missing stability a property of the cofibrations along which one pushes out.

With the definition recalled in Chapter~XIII, manuscript page~8, a map $i:A\to B$ lies in $\Cof_W$ when its pushout functor
\[
i_!:A\backslash\M\longrightarrow B\backslash\M,
\qquad (A\to X)\longmapsto(B\to B\amalg_A X),
\]
preserves weak equivalences. All pushouts mentioned here are assumed to exist. This is often called the condition of being an $h$-cofibration relative to $W$. The property automatically passes to cobase changes, by the pasting law for pushout squares. It is defined from $W$, before a model structure has been chosen.

The distinction between preserving weak equivalences under pushout \emph{along} $i$ and having every pushout \emph{of} $i$ belong to $W$ is essential. The latter condition defines $W^{\univ}$. Two-out-of-three gives $W^{\univ}\subset\Cof_W\cap W$: in the square comparing the cobase changes of $i$ to $X$ and to $Y$, the two cobase changes and the given map $X\to Y$ are weak equivalences. The reverse inclusion needs a factorization hypothesis. Grothendieck had explicitly used that hypothesis in Chapter~XIII, Corollary~1 and Proposition~4, manuscript pages~10--11; it has been dropped from Corollary~4.4.3 here. The next block records the resulting counterexample and restores the hypothesis exactly where it is needed.

The filtered-colimit assertion of Proposition~4.4.2 itself is sound and has a short explicit proof. Let $i_t:A_t\to B_t$ be a filtered diagram in $\Cof_W$, with colimit $i:A\to B$. Given a weak equivalence $u:X\to Y$ under $A$, restrict it under each $A_t$. Each map
\[
B_t\amalg_{A_t}X\longrightarrow B_t\amalg_{A_t}Y
\]
lies in $W$. Their filtered colimit is $B\amalg_A X\to B\amalg_A Y$, since colimits commute with colimits. Closure of $W$ under filtered colimits proves the required preservation. It also gives retract closure of $W$: a retract is the colimit of the countable sequence which repeats the associated idempotent. Thus a retract of a map in $\Cof_W$ again preserves $W$, because its pushout comparisons are retracts of the corresponding comparisons for the larger map. Together with pushout and composition, these facts prove the asserted stability under continuous transfinite composition by induction on the ordinal.
\end{dcomment}


\textbf{Proposition 4.4.2.} \textit{Soit $\mathcal{M}$ stable par petites $\varinjlim$ inductives filtrantes, et $W\subset$} $\Fl(\mathcal{M})$ \textit{un localiseur $($modérément saturé$)$} stable par lesdites limites. \textit{Alors} $\Cof_W$ \textit{satisfait les conditions de stabilité} (4.4.3) (“\textit{stabilité de Quillen gauche}”), \textit{et est stable aussi par $\varinjlim$ filtrantes quelconques.}

En effet, la condition $(a)$ est satisfaite sans condition supplémentaire sur $W$. D'autre part, il est immédiat que cette condition par rapport aux $\varinjlim$ filtrantes implique que $\Cof_W$ \textit{est stable par limites inductives filtrantes} dans $\underline{\Fl}(\mathcal{M})$.

\textbf{Corollaire 4.4.3.} \textit{Sous les conditions de la proposition} 4.4.2, $\Cof_W\cap W=W^{\univ}$ $($\textit{ensemble des $W$-équivalences co-universelles dans $\mathcal{M}$}$)$ \textit{est stable par limites inductives filtrantes, et donc “stable à gauche” en sens de} (4.4.3).

\textbf{Corollaire 4.4.4.} \textit{Sous les conditions de la proposition} 4.4.2, \textit{supposons de plus que les sommes amalgamées simples existent dans $\mathcal{M}$, et soit $C_0$ une petite partie de} $\Cof_W$ \textit{dont les éléments soient accessibles, et considérons}
\[
C\;=\;\widetilde{C_0}\ ,\qquad TC\;\deq\;C\cap W\ .
\]
\textit{Alors on a}
\[
C\;=\;\Cof_W\ ,\qquad TC\;=\;C\cap W^{\univ}\;\subset\;W^{\univ}\ ,
\]
\msp{page 41}
\textit{et $C$ et $TC$ sont $Q$-stables à gauche.}

\dgrblock{dgrb741}{A missing hypothesis in the couniversal-equivalence assertion}
\begin{dcomment}{development}{p40}{A missing hypothesis in the couniversal-equivalence assertion}
The intended application will have a useful factorization available before the second small object argument starts. Remembering that factorization repairs Corollary~4.4.3 and preserves the eventual recognition theorem. Forgetting it makes the intermediate assertion false even in a finite lattice.

\textbf{A finite counterexample.}\ours
Let $\M$ be the diamond lattice, viewed as a category, with $0<a<1$ and $0<b<1$, and with $a,b$ incomparable. Let $W$ consist of the identity maps and $i:0\to a$. Then $W$ satisfies two-out-of-three and is closed under all filtered colimits; all objects of $\M$ are finitely presentable. Nevertheless,
\[
i\in\Cof_W\cap W,\qquad i\notin W^{\univ}.
\]
Indeed, a filtered diagram in a finite poset has a greatest element in its image, and that element is its colimit. The same observation applies to the finite arrow poset, proving closure of $W$ under filtered colimits. It also shows directly that each representable preserves filtered colimits. To check two-out-of-three, the only nonidentity composable pairs are the two chains from $0$ to $1$ through $a$ and $b$; neither chain has two of its three maps in $W$.

Pushout along $i$ is the operation $x\mapsto a\vee x$. Among the arrows under $0$, the only nonidentity weak equivalence is $0\to a$, and it is sent to $\id_a$. Hence $i\in\Cof_W$. But its pushout along $0\to b$ is $b\to1$, which does not belong to $W$. This verifies every claimed property and disproves the equality in Corollary~4.4.3 under its printed hypotheses. If $C_0=\{i\}$, then $C=\widetilde{C_0}$ contains $i$ and its pushout $b\to1$. Thus $TC=C\cap W$ contains $i$ but not that pushout, disproving also the asserted pushout stability of $TC$ in Corollary~4.4.4. The other identities involving $W^{\univ}$ in that corollary require the same repair.

\textbf{The factorization that repairs the argument.}
Assume that every map factors as a member of $\Cof_W$ followed by a member of $W$. Then $\Cof_W\cap W=W^{\univ}$. This is the dual form of \citet[Chapter~XIII, Proposition~4, ms.~pp.~10--11]{GrothendieckDerivateursXIII}, and we give the complete argument in the notation used here.

Take $i:A\to B$ in $\Cof_W\cap W$ and any $u:A\to X$. Factor $u=qj$ with $j:A\to Z$ in $\Cof_W$ and $q:Z\to X$ in $W$. Form $P=B\amalg_A Z$ and $Q=B\amalg_A X$. The map $k:Z\to P$ is weak because it is the pushout of the weak equivalence $i$ along the $W$-cofibration $j$. It is also a $W$-cofibration, being a pushout of $i$. Therefore the induced map $q':P\to Q$, the pushout of $q$ along $k$, is weak. The equation $q'k=i'q$, where $i':X\to Q$ is the desired cobase change, and two-out-of-three imply $i'\in W$. The reverse inclusion was proved in the preceding block.

Proposition~4.4.6 supplies exactly this hypothesis: its first factorization has left map in $C\subset\Cof_W$ and right map in $TF\subset W$. Once that first factorization is constructed, the proof above establishes the pushout stability of $TC$, and all subsequent uses of Corollary~4.4.4 are justified. The false intermediate generalization therefore does not invalidate the final recognition proposition. The class $W^{\univ}$ itself remains closed under filtered colimits and the stated left stability operations even without the extra factorization.

Our literature check covered the earlier manuscript argument and the modern comparison in \citet[Lemma~1.10]{BataninBerger2017}, which distinguishes couniversal weak equivalences from trivial $h$-cofibrations and gives equality in left proper model categories. We did not find this finite lattice example in those sources.
\end{dcomment}


En effet, par le théorème 1 du paragraphe 1 (p. 10) et son corollaire, $\widetilde{C_0}$ est la plus petite partie de $\Fl(\mathcal{M})$ contenant $C_0$ et Q-stable à gauche. Comme $\Cof_W$ est Q-stable à gauche et contient $C_0$, on a donc $C=\widetilde{C_0}\subset\Cof_W$. Donc $TC\deq C{\cap}W\subset\Cof_W{\cap}W\deq W^{\univ}$, donc $TC=C\cap W^{\univ}$. Comme $W^{\univ}$ est lui aussi $Q$-stable à gauche (corollaire précédent), il en est de même de $C\cap W$, \quad \textit{q.e.d.}

\textbf{Corollaire 4.4.5.} \textit{Supposons que tout objet de $\mathcal{M}$ soit accessible. Sous les conditions du corollaire précédent, posant}
\[
TF\;=\;C_{0*}\;=\;C_*\ ,\qquad F\;=\;(TC)_*\ ,
\]
\textit{le couple $(C,TF)$ est un couple de Quillen. Pour qu'il en soit de même du couple $(TC,F)$, il suffit qu'il existe une} petite \textit{partie $TC_0$ de $C\cap W$, telle que $\widetilde{TC_0}=TC$ $(\deq C\cap W)$.}

La première assertion est, on l'a vu tantôt, le théorème 1 du paragraphe 1. La deuxième résulte de ce même théorème. On peut préciser dans ce sens :

\dgrblock{dgrb742}{An inclusion which cannot be an equality}
\begin{dcomment}{examples}{p41:La première assertion est, on l'a vu tantôt}{An inclusion which cannot be an equality}
Grothendieck has established enough closure to apply the small object theorem to a chosen set $C_0\subset\Cof_W$. The conclusion is that the class it generates stays inside $\Cof_W$. The display in Corollary~4.4.4 prints $C=\Cof_W$, although the proof on the following manuscript page proves $C\subset\Cof_W$. This is a slip in the edition, not in the transcription.

A small example shows why the stronger assertion cannot follow. Take $\M=\mathrm{Set}$, let $W$ consist of the bijections, and take $C_0=\{\id_{\{0\}}\}$. Every map belongs to $\Cof_W$, since every pushout functor preserves isomorphisms. Every map also has the right lifting property against $\id_{\{0\}}$, so
\[
C_0^{\square}=\operatorname{Mor}(\mathrm{Set}),\qquad
{}^{\square}(C_0^{\square})=\operatorname{Iso}(\mathrm{Set}).
\]
For the second equality, a map which lifts against every map lifts against itself in the square with identity horizontal maps, and the resulting filler is a two-sided inverse. Thus $C$ consists of the bijections, whereas $\Cof_W$ contains, for example, $\varnothing\to\{0\}$. All objects of $\mathrm{Set}$ are presentable for some regular cardinal, and the required colimits exist, so the example meets the hypotheses of the corollary.

Under the extra factorization hypothesis explained in the preceding block, the rest of the corollary has its intended force:
\[
TC=C\cap W=C\cap W^{\univ}.
\]
Both factors in the last intersection have the required stability under pushout, retracts, and continuous transfinite composition. Hence so does $TC$. This establishes the necessary closure of the proposed trivial cofibrations; the existence of a small set generating them remains the next problem.
\end{dcomment}


\textbf{Proposition 4.4.6.} \textit{Soit $\mathcal{M}$ une catégorie stable par petites $\varinjlim$ filtrantes, par sommes amalgamées, et dont tous les objets sont accessibles, $W$ un localiseur}
\msp{page 42}
\textit{modérément saturé dans $\mathcal{M}$, stable par $\varinjlim$ filtrantes, $C_0$ une} petite \textit{partie de} $\Fl(\mathcal{M})$, \textit{telle que}
\[
C_0\;\subset\;\Cof_W\ ,\qquad \underbrace{C_{0*}}_{(=TF)}\;\subset\;W\ ,
\]
\printedpage{28}
% head: A. Grothendieck : Dérivateurs, Ch. XV, Édition des Universités de Montpellier II et Paris VII
% foot: Version 5 novembre 2010
\textit{posons}
\[
C\;=\;\widetilde{C_0}\qquad TC\;=\;C\cap W\ ,\qquad TF\;=\;C_*\ ,\qquad F\;=\;(TC)_*\ ,
\]
\textit{et soit $TC_0$ une petite partie de $TC=C\cap W$. On a alors}
\[
\left\{
\begin{aligned}
C &\;\subset\;\Cof_W\ , &\qquad TC &\;=\;C\cap W^{\univ}\quad(\textit{pour mémoire})\\
\widetilde{TC_0} &\;\subset\;TC\ , & F_0 &\supset F\quad(\textit{ce dernier trivial})\ ,
\end{aligned}
\right.
\tag*{(4.4.6)}
\]
\textit{et les conditions suivantes sont équivalentes :}
\begin{enumerate}[leftmargin=*]
\item[a)] $\widetilde{TC_0}=TC\quad(\textit{on a}\subset)$\,;
\item[a$'$)] $F_0=F\quad(\textit{on a}\supset)$\,;
\item[b)] $TF=F_0\cap W\quad(\textit{on a}\subset)$\footnote[14]{\textbf{N.B.} On voit facilement que $\boxed{TF=F\cap W,}$ mais ça ne suffit pas pour avoir $TF=F_0\cap W$ pour $TC_0$ convenable.}\,;
\item[c)] $W=TF\circ\widetilde{TC_0}\quad(\textit{on a}\supset)$\footnote[15]{\textbf{N.B.} On sait que $\boxed{W=TF\circ TC,}$ ce qui n'est nullement suffisant pour impliquer $(c)$\,!}\,;
\end{enumerate}
\textit{et impliquent que $(W,C,F)$ est un triple de Quillen clos.}

\textsc{Démonstration.} L'inclusion $\widetilde{TC_0}\;\subset\;TC$ résulte de la Q-stabilité à gauche de $TC$ (cor. 4.4.4) et du fait que $TC_0\subset TC$ est petite (on applique alors le théorème 1 tout comme pour $C\subset\Cof_W\ldots$). L'équivalence de $(a)$, $(b)$, $(c)$ n'est autre que la proposition 4.3 appliquée à $W$ et aux deux couples de Quillen $(C,TF)$, $(\widetilde{TC_0},F_0)$ (compte tenu du théorème 1 du paragraphe 1). Qu'on a alors un triple de Quillen clos a été vu dans le corollaire 3 de la proposition 4.3 (p. 35). Il reste à prouver l'équivalence de $(a)$ et de $(a')$. On a $(a)\Longrightarrow(a')$, puisque $(TC_0)_*=(\widetilde{TC_0})_*=F_0$. Inversement, $(a')$ implique $\widetilde{TC_0}=\widetilde{TC}$, donc le théorème de factorisation pour $(TC,F)$ en vertu du théorème de factorisation, corollaire 1 (p. 11) appliqué à $TC_0\subset TC\subset\widetilde{TC_0}$ et à $F_0=F=(TC_0)_*$.

\dgrblock{dgrb743}{The two identities already available before the second factorization}
\begin{dcomment}{exposition}{p42:Il reste à prouver l'équivalence}{The two identities already available before the second factorization}
Grothendieck's two footnotes on this page separate a formal recognition problem from a problem of size. Put $C=\widetilde{C_0}$, $TF=C^{\square}\subset W$, $TC=C\cap W$, and $F=TC^{\square}$. The first weak factorization system $(C,TF)$ already gives
\[
W=TF\circ TC,\qquad TF=F\cap W,
\]
without assuming a factorization with left map in $TC$ for every arrow. The missing quantifier concerns arbitrary arrows, not only weak equivalences.

Here are complete proofs. Given $w\in W$, factor it as $w=pi$ using $(C,TF)$. Since $p\in TF\subset W$, two-out-of-three gives $i\in C\cap W=TC$. Conversely, a composite of a member of $TC$ and a member of $TF$ belongs to $W$. For the second identity, $TF\subset F\cap W$ follows from $TC\subset C$. If $f\in F\cap W$, factor $f=pi$ with $i\in C$ and $p\in TF$. Again $i\in TC$. The square with left arrow $i$, right arrow $f$, top arrow the identity of the source of $f$, and bottom arrow $p$ has a diagonal $s$ satisfying $si=\id$ and $fs=p$. Consequently $f$ is a retract of $p$ in the arrow category. The right lifting class $TF$ is retract closed, so $f\in TF$.

For a candidate set $TC_0\subset TC$, the notation left implicit in the statement of Proposition~4.4.6 is $F_0=TC_0^{\square}$. Smallness produces a weak factorization system $(\widetilde{TC_0},F_0)$, but only the inclusions
\[
\widetilde{TC_0}\subset TC,\qquad F\subset F_0
\]
are automatic. Proposition~4.4.6 says that the candidate succeeds precisely when $F_0\cap W=TF$, equivalently when every weak equivalence has the more restrictive factorization with left map in $\widetilde{TC_0}$. If $F_0=F$, then $\widetilde{TC_0}={}^{\square}F=\widetilde{TC}$; combined with $\widetilde{TC_0}\subset TC\subset\widetilde{TC}$, this gives equality immediately. Thus the final implication does not assume the desired second factorization in disguise.

To see the distinction concretely, take $\M=\mathrm{Set}$ and $W$ all maps. Let $C_0$ contain $\varnothing\to\{0\}$ and the fold $\{0,1\}\to\{0\}$. Right lifting against the first map is surjectivity, and right lifting against the second is injectivity. Hence $TF$ consists of bijections and $C$ consists of all maps. Then $TC$ is all maps and $F$ consists of bijections. If the candidate $TC_0$ contains only $\varnothing\to\{0\}$, its right class $F_0$ consists of surjections. The projection $\{0,1\}\to\{0\}$ belongs to $F_0\cap W$ but not to $TF$. Both identities of the preceding paragraph still hold; the candidate has simply omitted a necessary generating map.
\end{dcomment}


\msp{page 43}
Ainsi, le seul problème, une fois trouvé $C_0$ (petite) satisfaisant $C_0\subset\Cof_W$, $C_{0*}\subset W$, c'est de trouver une \textit{petite} partie $TC_0\subset TC$ qui “engendre $TC\;=\;C\cap W$ ”, \textit{i.e.} qui satisfasse aux conditions équivalentes $(a)$, $(a')$, $(b)$, $(c)$ de la proposition précédente. J'ai passé des jours à essayer de prouver l'existence de $TC_0$, en supposant $\mathcal{M}$ munie d'une filtration cardinale, et $W$ une partie accessible de $\mathcal{M}$, sans y être arrivé. Bien sûr, tout ce qui nous importe vraiment, c'est que $(TC,F)$ satisfasse à la condition de factorisation. Mais je ne vois aucun moyen de le prouver, si ce n'est en “engendrant” à gauche $TC$ par un \textit{petit} sous-ensemble $TC_0$.

\dgrblock{dgrb744}{The generating set supplied by Smith's theorem}
\begin{dcomment}{later}{p43:J'ai passé des jours à essayer}{The generating set supplied by Smith's theorem}
The question at the end of Section~4 became a standard recognition theorem for combinatorial model categories. In the locally presentable setting, Smith's theorem, published with proof by \citet[Theorem~1.7]{Beke2000}, gives the missing set of generating trivial cofibrations. Its hypotheses are precisely organized around the difficulty seen here: $W$ has two-out-of-three and is retract closed, $C_0^{\square}\subset W$, $C\cap W$ is stable under pushout and transfinite composition, and $W$ satisfies a solution set condition at each map of $C_0$. Accessibility of the full subcategory $W\subset\M^{[1]}$ supplies the last condition by \citet[Proposition~1.15]{Beke2000}.

The solution set condition says that, for each $i\in C_0$, a set of maps $i\to w_s$ with $w_s\in W$ is weakly initial among all maps from $i$ to members of $W$ in the arrow category. This is the categorical form of the bounded choices Grothendieck is seeking. Smith's argument converts these choices into a set $J\subset C\cap W$ with $\operatorname{cof}(J)=C\cap W$; Beke's Lemmas~1.8 and~1.9 supply the construction and proof.

This answers the existence question under the stated modern hypotheses. A cardinal filtration on its own should not silently be substituted for local presentability, and the claim does not cover every category allowed by the more flexible assumptions on the preceding pages. Nor does the theorem require $C$ itself to be stable under all filtered colimits, the additional route Grothendieck begins to explore in Section~5.
\end{dcomment}


\dgrblock{dgrb745}{The form closest to the proposed localizer construction}
\begin{dcomment}{status}{p43:J'ai passé des jours à essayer}{The form closest to the proposed localizer construction}
Grothendieck's use of $\Cof_W$ also survives in a recognition theorem whose hypotheses closely follow these pages. In the updated online text of \citet[Proposition~A.2.6.15]{Lurie2009HTT}, a locally presentable category, a perfect class $W$, and a set of generating cofibrations satisfying the pushout preservation condition and $C_0^{\square}\subset W$ produce a left proper combinatorial model structure. Here ``perfect'' means that $W$ contains the isomorphisms, satisfies two-out-of-three, is closed under filtered colimits, and is generated under those colimits by a set. The proof first establishes the required stability of $C\cap W$, then applies the small-generation result of Lemma~A.2.6.9 and the general recognition Proposition~A.2.6.10. These are the numbers in the online revision consulted, not the earlier printed numbering.

The later theorem makes precise both parts of the intended conclusion: the second factorization exists, and the chosen cofibrations preserve weak equivalences under pushout. The latter is left properness. It remains separate from the formal two-factorization axioms studied at the beginning of Section~4; an arbitrary model structure need not have it.
\end{dcomment}


\dgrblock{dgrb746}{The passage from models to diagram homotopy theories}
\begin{dcomment}{later}{p43:J'ai passé des jours à essayer}{The passage from models to diagram homotopy theories}
The point of obtaining a model structure is that it can support coherent operations on diagrams. In the combinatorial setting, the later passage is precise: \citet[Theorem~1.28 and Proposition~1.30]{Groth2013} constructs projective and injective model structures on diagram categories and proves that $J\mapsto\operatorname{Ho}(\mathcal M^J)$ is a strong derivator. The two diagram model structures supply the appropriate derived Kan extensions.

The existence of these diagram model structures is a convenient sufficient route. The broader construction of homotopy right Kan extensions for model categories with small limits is treated by \citet[Theorem~6.11]{Cisinski2003}, in that paper's left-derivator convention. This later theorem explains why the goal is larger than a particular generating-set argument. Chapter~XV builds and controls models; the derivator records their homotopy theories of diagrams and the functors between them. Neither the cardinal estimates nor the two lifting systems alone constitute the passage between the two.
\end{dcomment}


\section*{5\quad Résultats techniques préliminaires sur les parties accessibles de $\mathrm{Ob}\,\mathcal{M}$, stables par limites inductives filtrantes}

Je crois pourtant que je peux arriver à quelque chose d'utilisable, en travaillant systématiquement avec des $C$ ($\subset\Cof_W$, Q-closes à gauche, Q-engendrées à gauche par une \textit{petite} partie $C_0$ de $C$) qui sont stables par toutes limites inductives filtrantes.

% ---- from Der-ChXV-05-11-10_p031-035.tex
\printedpage{29}
% head: A. Grothendieck : Dérivateurs, Ch. XV, Édition des Universités de Montpellier II et Paris VII
% foot: Version 5 novembre 2010

Je suppose à nouveau sur $\cM$ les conditions suivantes :
\msp{page 44}

\begin{enumerate}[nosep]
\item[\textit{a})] $\cM$ stable par petites limites inductives (et pas seulement filtrantes).
\item[\textit{b})] $\cM$ muni d'une filtration cardinale $(\Filt^{\pi})_{\pi\geq\pi_0}$.
\item[\textit{c})] Les objets de $\cM$ sont accessibles (\textbf{N.B.} Il suffit de le supposer pour les objets de $\Filt^{\pi_0}$).
\end{enumerate}

La condition ($b$) implique que $\cM$ contient une petite famille génératrice par épimorphismes stricts (p. ex. une sous-catégorie $\cC$ petite de $\uFilt^{\pi_0} = \cM^{\pi_0}$ telle que l'inclusion $\cC \longrightarrow \uFilt^{\pi_0}$ soit une équivalence de catégories). Par SGA 4 I 8.12.7, il en résulte que $\cM$ est stable par petites limites projectives. Il s'ensuit le

\textbf{Lemme 5.1} (sous les conditions ($a$), ($b$)). \textit{Pour tout objet $X$ de $\cM$, l'ensemble} $\Ss(X)$ \textit{des sous-objets de $X$, muni de la relation d'ordre naturelle ``d'inclusion'', admet des} inf \textit{et des} sup \textit{quelconques.}

Des inf, car ceux-ci s'expriment comme des limites projectives dans $\cM$. Or l'existence des inf implique celle des sup, puisque $\sup_J = \inf_{J'}$, où $J'$ est l'ensemble des majorants de $J$ dans $\Ss(X)$.

\textbf{Corollaire 5.2.} \textit{Pour toute flèche $f : Z \longrightarrow X$ dans $\cM$, il existe, parmi les sous-objets $X'$ de $X$ par lesquelles $f$ se factorise, un plus petit} % [sic] lesquelles -> lesquels, "sous-objets" is masculine
\msp{page 45}
\textit{(qu'on notera} $\Img f$\textit{). Le morphisme $Z \longrightarrow \Img f$ est un épimorphisme.}

\dgrblock{dgrb747}{Images and the meaning of a small subobject}
\begin{dcomment}{exposition}{p44}{Images and the meaning of a small subobject}
Grothendieck now wants the size of an object to be visible inside the object itself. A presentation by a filtered diagram can repeat objects, identify elements, and carry many choices. The ordered collection of its small subobjects is a more geometric presentation: its transition maps are inclusions, and its small directed suprema can be calculated in the ambient category. This change of presentation is the preparation for recognizing a class of objects from its bounded part.

We keep the three conditions on the filtration separate. Essential smallness says that a bounded collection has a set of isomorphism classes; generation says that every object has an appropriately directed presentation by bounded objects; preservation of such colimits by a represented functor is an additional condition on the representing object. In the terminology of these pages, ``grand devant $\pi$'' means that every subset of cardinal at most $\pi$ has an upper bound. For infinite $\pi$, this is the usual $\pi^+$-directedness. A $\pi$-accessible object is consequently a $\pi^+$-presentable object in current terminology. It is not an object with exactly $\pi$ elements, and the bound is not assumed regular.

The set-sized assertion in Lemma~5.1 also has content. Let $G$ run through a small set of representatives of the base level. A subobject $U\hookrightarrow X$ determines the subsets of $\Hom(G,X)$ consisting of the maps which factor through $U$. These subsets determine $U$: a generating presentation of its domain lifts uniquely to any other subobject having the same subsets, and these unique lifts assemble. Thus there is a set of subobjects of $X$.

The same base level is a conservative generating family. Suppose $f:X\to Y$ induces bijections on maps from its objects. Write $Y=\varinjlim_iG_i$ using the filtration. Each structural map $G_i\to Y$ lifts uniquely to $X$; uniqueness makes the lifts compatible, so they induce $g:Y\to X$ with $fg=\mathrm{id}_Y$. For a presentation $X=\varinjlim_jH_j$, injectivity of $\Hom(H_j,f)$ gives $gf h_j=h_j$ at every structural map $h_j$, hence $gf=\mathrm{id}_X$. No accessibility rank is needed. The dual of \citep[Expos\'e~I, 8.12.7]{SGA4I1972} now supplies completeness: the functor of cones $Z\mapsto\varprojlim_j\Hom(Z,X_j)$ takes colimits to limits and is represented by the desired limit. This argument uses conservative generation, without requiring the stronger density assertion in the source.

Once intersections exist, the image used here is the intersection of \emph{all} subobjects through which the map factors. If $Z\to\operatorname{Im}(f)$ were not epic, two different maps out of its codomain would agree after composition; their equalizer would give a smaller subobject through which $f$ factors. This proves Corollary~5.2. In fact the first map in this factorization is extremal epic: any factorization through a monomorphism would contradict the minimality of the image unless that monomorphism were invertible.

In Lemma~5.3 the adjective \emph{strict} has the categorical meaning of SGA~4, not merely ``different from the whole object''; the latter requirement is written separately. A strict monomorphism which is also epic is invertible. Conversely, failure to be epic is detected by a proper equalizer, which is a strict monomorphism. This proves the stated criterion, after replacing the accidental name $g$ by $f$. The distinction matters, for example, for the epimorphism $\mathbb Z\to\mathbb Q$ in commutative rings: it factors through a proper monomorphism, namely itself, but not through a proper strict monomorphism. The terminology is fixed in \citep[Expos\'e~I, 10.2--10.5]{SGA4I1972}.
\end{dcomment}


Cette dernière assertion provient de celle-ci :

\textbf{Lemme 5.3.} \textit{Soient $\cM$ une catégorie où les noyaux des doubles flèches existent. Soit $f : Z \longrightarrow Y$ une flèche de $\cM$. Alors $g$ est un épimorphisme si et seulement si $g$ ne se factorise pas par un sous-objet} strict $Y_0$ \textit{de $Y$ différent de $Y$.} % [sic] g -> f, the arrow has been named f

On va faire une hypothèse supplémentaire (pratiquement anodine) sur les $\Filt^{\pi}$ :

\begin{enumerate}[nosep]
\item[\textit{d})] $\Filt^{\pi}$ est stable par limites inductives indexées par des catégories d'indices $I$ telles que $\card \Fl(I) \leq \pi$\footnote[16]{Cette condition sera nommée ($d_0$) par la suite.}. (Il suffit de supposer la stabilité par objet initial, sommes amalgamées simples, et $\varinjlim$ filtrantes indexées par un ensemble ordonné $I$ de cardinal $\leq \pi$.) De plus, si $f : Z \longrightarrow X$ est une flèche de $\cM$ telle que $Z$ ou $X$ est dans $\Filt^{\pi}$, alors $\Img f \in \Filt^{\pi}$.
\end{enumerate}

Cette dernière condition sur $\Img f$ se décompose en deux :
\msp{page 46}

\begin{enumerate}[nosep]
\item[$d_1$)] Tout sous-objet d'un objet de $\Filt^{\pi}$ est dans $\Filt^{\pi}$.
\item[$d_2$)] Si $Z$ est dans $\Filt^{\pi}$, alors tout quotient $Z'$ de $Z$ tel que $\Img p = Z'$ (où $p : Z \longrightarrow Z'$ est le morphisme canonique) est dans $\Filt^{\pi}$.\footnote[17]{Il semblerait qu'on n'aura pas besoin de ($d_1$), seulement de ($d_2$). Mais en fait, ($d_1$) doit être conséquence de ($d_2$), et ($d_2$) est conséquence de la première partie de ($d$) (stabilité par $\varinjlim$ finies) si tout épimorphisme dans $\cM$ est strict, donc effectif (si les $\Filt^{\pi}$ sont stables par $\varprojlim$ finies aussi \ldots).}
\end{enumerate}

\printedpage{30}
% head: A. Grothendieck : Dérivateurs, Ch. XV, Édition des Universités de Montpellier II et Paris VII
% foot: Version 5 novembre 2010

Bien sûr, pratiquement il sera toujours vrai que $\Filt^{\pi}$ est stable par passage aux sous-objets et aux quotients d'un objet (pas seulement les quotients particuliers dont il est question dans ($d_2$). % [sic] closing parenthesis of "(pas seulement" missing, parentheses unbalanced as printed
D'ailleurs, lorsque tout bimorphisme (\textit{i.e.} monomorphisme \textit{et} épimorphisme) de $\cM$ est un isomorphisme, alors ($d_2$) équivaut bien à la stabilité de $\Filt^{\pi}$ par quotients d'objets.

\textbf{Proposition 5.3.} \texttt{[Deux fois 5.3.]} \textit{Supposons les conditions} (a), (b), (d) \textit{vérifiées pour $\cM$. Soient $\pi$ un cardinal $\geq \pi_0$, et $X$ un objet de $\cM$. Soit} $\Ss_{\pi}(X)$ \textit{l'ensemble des sous-objets de $X$ qui sont dans} $\Filt^{\pi}$\textit{. Cet ensemble, ordonné par la relation d'ordre induite par} $\Ss(X)$\textit{, est grand devant $\pi$ (et en particulier filtrant). Si de plus, les limites inductives filtrantes de monomorphismes dans $\cM$ sont des monomorphismes}\footnote[18]{C'est une \textit{condition d'exactitude} de $\cM$, que je vais noter comme :\par
\noindent\circled{e}\quad Les monomorphismes dans $\cM$ sont stables par $\varinjlim$ filtrantes dans $\uFl(\cM)$.} \textit{(p. ex. si les limites inductives filtrantes sont exactes à gauche), ou si les objets de} $\Filt^{\pi}$ \textit{sont $\pi$-accessibles, on a $\varinjlim\limits_{\alpha\in\Ss_{\pi}(X)} Z_{\alpha} \isoto X$}\footnote[19]{Dire aussi que $\Ss_{\pi}(X)$ est $\pi$-adapté, et le système inductif $(Z_{\alpha})_{\alpha}$ aussi.}.

\msp{page 47}
\textsc{Démonstration}. Soit $(Z_{\alpha})_{\alpha\in A}$ une famille de sous-objets de $X$ tels que $Z_{\alpha} \in \Filt^{\pi}$, et $\card A \leq \pi$. Soit $Z = \coprod_{\alpha\in A} Z_{\alpha}$, on a $Z \in \Filt^{\pi}$ par ($c$). Le plus petit sous-objet de $X$ majorant les $Z_{\alpha}$ n'est autre que $\Img (Z \longrightarrow X)$, qui est dans $\Filt^{\pi}$ par ($d_1$). Prouvons que $\varinjlim\limits_{\alpha\in\Ss_{\pi}(X)} Z_{\alpha} \isoto X$ est un isomorphisme, quand on sait que c'est un monomorphisme. Or
\[
X \;=\; \varinjlim_{I} X_i \ ,
\]
où $(X_i)_{i\in I}$ est un système inductif filtrant, $I$ ensemble ordonné grand devant $\pi$, les $X_i$ dans $\Filt^{\pi}$. Soit, pour tout $i \in I$, $X'_i$ l'image de $X_i$ dans $X$. On a donc $X'_i \in \Filt^{\pi}$ par ($d_2$). Utilisant le fait que les $X_i \longrightarrow X'_i$ sont des épimorphismes, on voit de suite que $X = \varinjlim\limits_{I} X'_i$. Considérons l'application
\[
\begin{aligned}
I \quad &\longrightarrow \quad \Ss_{\pi}(X)\\
i \quad &\longmapsto \quad X'_i \ .
\end{aligned}
\]
C'est une application croissante, soit $J$ son image, partie filtrante grande devant $\pi$ de $\Ss_{\pi}(X)$, telle que
\[
\varinjlim_{\alpha\in J} Z_{\alpha} \ \isoto \ X \ .
\]
\msp{page 48}
On a donc
\[
\underbrace{\varinjlim_{\alpha\in J} Z_{\alpha}}_{\simeq X} \ \longrightarrow \ \varinjlim_{\alpha\in\Ss_{\pi}X} Z_{\alpha} \ \xrightarrow{\ \varphi\ } \ X \ ,
\]
le composé étant l'identité. Si on sait que $\varphi$ est un monomorphisme, il en résulte que c'est un isomorphisme (monomorphisme ayant une section). Or $\varphi$ sera un monomorphisme si toute $\varinjlim$ filtrante de monomorphismes est un monomorphisme, et alors on gagne.

\dgrblock{dgrb748}{Quotient bounds do not bound subobjects}
\begin{dcomment}{development}{p46}{Quotient bounds do not bound subobjects}
The two image conditions serve different purposes. The second controls the image of a bounded source; the first controls an arbitrary subobject of a bounded target. The suggestion in footnote~17 that the second should imply the first is false even in a module category with exact filtered colimits. Giving the example explicitly explains why the two conditions must be distinguished, even in arguments where only one of them is used.

\textbf{A filtration separating $(d_1)$ and $(d_2)$.}\ours
Fix an infinite cardinal $\pi$, a field $k$, and a set $S$ of cardinal greater than $\pi$. Put
\[
 R=k[x_s\mid s\in S],\qquad
 \Filt^{\kappa}(R\text{-}\mathrm{Mod})
 =\{N\mid N\text{ is generated by at most }\kappa\text{ elements}\},
 \quad \kappa\geq\aleph_0.
\]
This is a cardinal filtration satisfying $(a)$, $(b)$, $(c)$, $(d_0)$, $(d_2)$ and $(e)$, but $(d_1)$ fails at $\pi$.

\emph{Proof.} Modules generated by at most $\kappa$ elements are quotients of $R^{(\kappa)}$, so they have a set of isomorphism classes. A colimit of a diagram with at most $\kappa$ arrows and $\kappa$-generated values is a quotient of the sum of at most $\kappa$ such modules; it is again $\kappa$-generated. This proves $(d_0)$ and the required closure part of the filtration, and quotient closure proves $(d_2)$. Every module is presentable for some sufficiently large regular cardinal: choose a presentation, and bound both its generators and relations; factorization and eventual equality in a sufficiently directed colimit are then checked on those sets. This proves $(c)$. Filtered colimits of modules are exact, giving $(e)$.

For the quantitative generation clause, let $N$ have a chosen generating set $T$ of size at most $\kappa'$, where $\kappa'\geq\kappa$. The submodules generated by subsets of $T$ of size at most $\kappa$ form a $\kappa^+$-directed diagram with colimit $N$. Its indexing set has cardinal at most $(\kappa')^{\kappa}$. This is precisely the bound required of a cardinal filtration, including when the ring itself is much larger than $\kappa$.

Finally $R$ is cyclic, whereas its ideal $\mathfrak a=(x_s\mid s\in S)$ is not $\pi$-generated. Indeed $R/\mathfrak a=k$ and
\[
 \mathfrak a/\mathfrak a^2\cong\bigoplus_{s\in S}k\,\overline{x_s}.
\]
Any generating set of $\mathfrak a$ maps to a spanning set of this $k$-vector space and therefore has cardinal at least $|S|$. Thus $\mathfrak a\hookrightarrow R$ disproves $(d_1)$.\hfill$\square$

The example does not challenge the propositions which follow: they assume $(d)$, hence both image conditions. It also explains why increasing the cardinal can remove the difficulty in familiar concrete categories. If $\kappa\geq|R|$, every submodule of a $\kappa$-generated module has at most $\kappa$ underlying elements and is itself $\kappa$-generated. We checked the filtration definition and the relevant size arguments in \citep[Expos\'e~I, 9.12--9.16]{SGA4I1972}; the particular counterexample above is our explicit verification of the remark's failure, with no claim of priority.
\end{dcomment}


\printedpage{31}
% head: A. Grothendieck : Dérivateurs, Ch. XV, Édition des Universités de Montpellier II et Paris VII
% foot: Version 5 novembre 2010

Si on suppose que les éléments de $\Filt^{\pi}$ sont $\pi$-accessibles, on gagne aussi, car on trouve :

\textbf{Corollaire 5.4.} \textit{Si les objets de} $\Filt^{\pi}$ \textit{sont $\pi$-accessibles (p. ex. si $\pi \geq \pi_1$ et $\pi = \pi^{\pi_0}$, p. ex. $\pi \geq \pi_1$ et $\pi = 2^{c}$, avec $c \geq \pi_0$,} cf. (3.2) \textit{p. 16), alors pour toute partie filtrante $J$ de} $\Ss(X)$ \textit{grande devant $\pi$, telle que $\varinjlim\limits_{\alpha\in J} Z_{\alpha} \isoto X$, $J$ est cofinale dans} $\Ss_{\pi}(X)$ \textit{(donc $\varinjlim\limits_{\alpha\in J} Z_{\alpha} \isoto \varinjlim\limits_{\alpha\in\Ss(X)} Z_{\alpha}$).}

C'est immédiat.

\msp{page 49}
\textbf{Scholie 5.5.} Quand on fait sur $\cM$ l'hypothèse d'exactitude assez anodine que les $\varinjlim$ filtrantes de monomorphismes sont des monomorphismes, \textit{ou} qu'on se borne aux cardinaux $\pi$ tels que les objets de $\Filt^{\pi}$ soient $\pi$-accessibles, on peut \textit{canoniser} la représentation d'un objet $X$ comme $\varinjlim\limits_{i\in I} X_i$, avec $X_i \in \Filt^{\pi}$ et $I$ grand devant $\pi$, en prenant $I = \Ss_{\pi}(X)$, et le système inductif correspondant. On a de plus l'agréable propriété supplémentaire que voici : toute partie filtrante $J$ de $I$, de cardinal $\leq \pi$, admet une borne supérieure $i_J$, et (si on suppose que $\cM$ satisfait la condition ($e$), page 46, note 18), on aura
\[
X_{i_J} \;=\; \varinjlim_{i\in J} X_i \ ;
\]
en d'autres termes, on a une représentation de $X$ par un système inductif (parfaitement) \textit{$\pi$-adapté} (\textit{cf.} XIII, p. 85, 86).

\msp{page 50}
Soit maintenant
\[
M \;\subset\; \Ob\cM
\]
une partie de $\Ob\cM$, stable par la relation d'isomorphie. Si $M$ est $\pi$-accessible (suivant la condition $\mathrm{b}_{\pi}$ de XIII, p. 88), alors en posant
\[
M^{\pi} \;=\; M \cap \Filt^{\pi} \ ,
\]
$M^{\pi}$ détermine $M$ par la condition suivante sur $X \in \Ob\cM$.
\[
\begin{aligned}
&X \in M \iff \text{le sous-ensemble } \Ss^{M}_{\pi}(X) \text{ de } \Ss_{\pi}(X) \text{ formé des } X_{\alpha} \in \Ss_{\pi}(X)\\
&\text{tels que } X_{\alpha} \in M \ (\iff X_{\alpha} \in M^{\pi}) \text{ est cofinal dans } \Ss_{\pi}(X).
\end{aligned}
\tag*{(5.1)}
\]
Notons que lorsque $X \in \Filt^{\pi}$, ce critère est trivialement valable, puisque alors $\Ss_{\pi}(X)$ admet $X$ comme plus grand élément, et les sous-ensembles cofinaux sont ceux qui contiennent cet objet final. Donc (5.1) est intéressant seulement pour les $X \in \Ob\cM \backslash \Filt^{\pi}$. Notons que pour \textit{toute} partie $M^{\pi}$ de $\Filt^{\pi}$, on peut définir une partie $M$ de $\Ob\cM$ par la condition (5.1), et on aura trivialement
\[
M^{\pi} \;=\; M \cap \Filt^{\pi} \ . \tag*{(5.2)}
\]
Mais sans doute, bien que $M$ satisfasse (5.1), il n'est pas pour autant nécessairement $\pi$-accessible. En effet, pour qu'il le soit, il faut qu'il soit stable par $\varinjlim$
\printedpage{32}
% head: A. Grothendieck : Dérivateurs, Ch. XV, Édition des Universités de Montpellier II et Paris VII
% foot: Version 5 novembre 2010
\msp{page 51}
grandes devant $\pi$ d'éléments de $M^{\pi}$, ce qui implique que tout facteur direct d'un élément de $M^{\pi}$ est dans $M$, donc aussi dans $M^{\pi}$, d'après la condition de stabilité ($d$) sur $\Filt^{\pi}$. C'est là une condition particulière sur la partie $M^{\pi}$ de $\Filt^{\pi}$. Aussi, on devra avoir une condition sur $\cM$ plus forte que la validité de (5.1), savoir :
\[
\parbox{0.72\textwidth}{Pour tout $X \in \Ob\cM$ et pour toute partie $J$ de $\Ss_{\pi}(X)$, telle que
\begin{enumerate}[nosep]
\item[\textit{a})] $J$ est cofinale dans $\Ss_{\pi}(X)$ et
\item[\textit{b})] Si $J' \subset J$ est une partie filtrante de $J$, de cardinal $\leq \pi$, alors sa borne supérieure $\alpha_{J'}$ dans $\Ss_{\pi}(X)$ appartient à $J$,
\end{enumerate}
on a $X \in M$ si et seulement si $J^{M} \overset{\text{déf}}{=} J \cap \Ss^{M}_{\pi}(X)$ est cofinal dans $J$ (ou ce qui revient au même, dans $\Ss_{\pi}(X)$)\,\footnotemark[20].}
\tag*{(5.3)}
\]
\footnotetext[20]{Au lieu de $J$ cofinale dans $\Ss_{\pi}(X)$, il suffit même que l'on ait $\varinjlim\limits_{\alpha\in J} X_{\alpha} \isoto X$ -- la condition correspondante sur $M$ devient donc plus exigeante. Mais si les objets de $\Filt^{\pi}$ sont $\pi$-accessibles, cette condition est en fait équivalente à la condition plus faible.}

\dgrblock{dgrb749}{The canonical directed union}
\begin{dcomment}{exposition}{p49}{The canonical directed union}
The gain in Proposition~5.3 is a canonical presentation whose indexing order has actual suprema. Write $\operatorname{Sub}_{\pi}(X)$ for Grothendieck's $\Ss_{\pi}(X)$. Given at most $\pi$ members $U_a\hookrightarrow X$, form their coproduct and then its image in $X$. The coproduct belongs to $\Filt^\pi$ by $(d_0)$, and its image belongs to $\Filt^\pi$ by $(d_2)$. It is their least upper bound in the full subobject lattice. These are the two assumptions used in the proof; the printed references to $(c)$ and $(d_1)$ on manuscript page~47 are misplaced.

Here is the colimit argument in a form that also explains the role of epimorphisms. Choose a generating presentation $X=\varinjlim_iX_i$ by $\Filt^\pi$-objects, and factor each map to $X$ as
\[
 X_i\twoheadrightarrow U_i\hookrightarrow X.
\]
A compatible family $U_i\to T$ gives, by composition, a compatible family $X_i\to T$ and hence a unique map $X\to T$. Its restriction to $U_i$ is the given map because $X_i\to U_i$ is epic. Therefore $\varinjlim_iU_i=X$. Passing to the image of the indexing order in $\operatorname{Sub}_{\pi}(X)$ does not change this colimit: above any chosen $U_i$, the original directed order supplies a common upper bound for every finite family of representatives.

There is consequently a section of the canonical map
\[
 q:\varinjlim_{U\in\operatorname{Sub}_{\pi}(X)}U\longrightarrow X.
\]
Under $(e)$ this map is monic, being the filtered colimit of the monomorphisms $U\to X$; a monomorphism with a section is invertible. Alternatively, if every $\Filt^\pi$-object is $\pi$-accessible, any small subobject $V\hookrightarrow X$ factors through some $U_i$, and the factorization is a map of subobjects. The $U_i$ are then cofinal in $\operatorname{Sub}_{\pi}(X)$, proving the same result without $(e)$.

For the stronger continuity assertion of Scholie~5.5, $(e)$ is used again. If $K\subset\operatorname{Sub}_{\pi}(X)$ is directed and $|K|\leq\pi$, its colimit belongs to $\Filt^\pi$ by $(d_0)$, and its map to $X$ is monic by $(e)$. It is the supremum of $K$. Thus the canonical diagram carries small directed suprema to colimits. This is the precise content of a \emph{perfectly $\pi$-adapted} representation. It is stronger than the mere requirement that its index be $\pi^+$-directed.

In a presheaf category on a small category $A$, this construction has a concrete form once $\pi$ bounds $|\operatorname{Fl}(A)|$. Use the filtration by presheaves whose total set of elements has cardinal at most $\pi$. Closing at most $\pi$ chosen elements under all restriction maps still gives at most $\pi$ elements. The poset of such subpresheaves is $\pi^+$-directed, its small directed suprema are unions, and their union is the whole presheaf. This familiar example is the geometric picture underlying the abstract image argument.
\end{dcomment}


Si on admet déjà le ``il suffit'' de (5.1), alors il suffit dans (5.3) de supposer $X \in M$ et exiger qu'alors $J^{M}$ soit cofinal dans $J$.

Mais notons la
\msp{page 52}

\dgrblock{dgrb750}{Accessibility tests every adapted presentation}
\begin{dcomment}{exposition}{p51}{Accessibility tests every adapted presentation}
The word ``accessible'' now applies to a \emph{class of objects}. Its meaning must be recovered from the definition to which Grothendieck refers, since it is more specific than simply saying that the corresponding full subcategory is accessible in the modern sense. The relevant condition $b_\pi$ in \citep[Chapter~XIII, manuscript pp.~88--89]{GrothendieckDerivateursXIII} is this: for every perfectly $\pi$-adapted presentation
\[
 X=\varinjlim_{i\in I}X_i,\qquad X_i\in\Filt^\pi,
\]
one has $X\in M$ if and only if $\{i\mid X_i\in M\}$ is cofinal in $I$. Thus the test concerns all such presentations, including presentations whose maps identify elements.

Formula~(5.1) applies this test to the canonical subobject presentation. It records a necessary condition, and gives a way to reconstruct a candidate class from any chosen bounded part $M^\pi$. On bounded objects the recovery formula is automatic: $X$ itself is the largest element of $\operatorname{Sub}_{\pi}(X)$, so a cofinal subset must contain it. This explains (5.2), but supplies no invariance under changing the presentation.

Condition~(5.3) asks for that invariance among the continuous cofinal subdiagrams of the canonical diagram. The closure requirement on $J$ says that it contains the supremum of every directed subset of itself of cardinal at most $\pi$. It makes the induced diagram perfectly adapted. It does not turn a nonmonic presentation into a subobject presentation, and this is where the first version of Proposition~5.6 overreaches. The source itself identifies the obstruction and inserts the correction in footnote~21.

The repair on manuscript page~55 asks that $M^\pi$ be closed under filtered colimits of diagrams of cardinal at most $\pi$. Then the good indices in a perfectly adapted diagram are themselves closed under the same directed suprema. They form another perfectly adapted presentation whenever they are cofinal. Lemma~5.7 can therefore be applied to that presentation, reducing its colimit to a continuous cofinal family of monomorphic stages. This is why the extra closure condition repairs the converse: it allows the comparison lemma to be used a second time, on the good part of the diagram.
\end{dcomment}


\textbf{Proposition 5.6.}\footnote[21]{Faux tel quel, vrai si on suppose que $M^{\pi}$ est stable par $\varinjlim$ filtrantes.} \textit{Supposons $\pi$ choisi de telle façon que les objets de} $\Filt^{\pi}$ \textit{soient $\pi$-accessibles (p. ex. $\pi = 2^{c}$, avec $c \geq \pi_1$,} cf. (3.2), \textit{p. 16). Alors, pour une partie $M$ de} $\Ob\cM$\textit{, $M$ est $\pi$-accessible si et seulement si elle satisfait au critère} (5.3).

On a vu que c'est nécessaire (tautologiquement, compte tenu de la proposition 5.3 et du scholie 5.5, bas de la page 49). Prouvons que c'est suffisant, en établissant d'abord le

\textbf{Lemme 5.7.} \textit{Sous les conditions de} 5.6 \textit{sur $\pi$, et la condition} (e) \textit{sur $\cM$} (p. 46), \textit{pour toute représentation $\pi$-adaptée $X = \varinjlim\limits_{i\in I} X_i$ d'un objet $X$ de $\cM$, le sous-ensemble $I'$ des $i \in I$ tels que $X_i \longrightarrow X$ soit un monomorphisme, est cofinal dans $I$. De plus, $I'$ est stable par bornes supérieures dans $I$ de parties filtrantes $J$ de $I'$ de cardinal $\leq \pi$.}

Montrons comment ce lemme 5.7 implique le ``il suffit'' de la proposition 5.6. Soit donc $X = \varinjlim\limits_{i\in I} X_i$ une représentation $\pi$-adaptée de $X$ ($X_i \in \Filt^{\pi}$, $I$ grand devant $\pi$, \textit{etc.}). Il faut prouver que, si 5.3 est valable, alors $X \in M$ si et seulement si l'ensemble $I_M$ est cofinal dans $I$.

\msp{page 53}
Supposons $X \in M$. Considérons le sous-ensemble $I' \subset I$ décrit dans le lemme. On a une application $I' \longrightarrow \Ss_{\pi}(X)$, soit $i \longmapsto X'_i = \Img (X_i \hookto X)$, et on a
\[
X \simeq \varinjlim_{i\in I} X_i \ \simeq \ \varinjlim_{i\in I'} X_i \ \simeq \ \varinjlim_{i\in I'} X'_i \ \simeq \ \varinjlim_{\alpha\in J} Z_{\alpha} \ ,
\]
où $J$ est l'image de $I'$ dans $\Ss_{\pi}(X)$. Comme les objets de $\Filt^{\pi}$ sont $\pi$-accessibles, on en conclut que $J$ est cofinal dans $\Ss_{\pi}(X)$. Par 5.3, il s'ensuit que l'ensemble $J_M$ des $\alpha \in J$
\printedpage{33}
% head: A. Grothendieck : Dérivateurs, Ch. XV, Édition des Universités de Montpellier II et Paris VII
% foot: Version 5 novembre 2010
tels que $Z_{\alpha} \in M$ est cofinal dans $J$, donc $\neq \emptyset$. Donc $I'_M$, qui est l'image inverse de $J_M$ dans $I'$, est en tous cas non vide. Appliquant le même résultat en remplaçant $I'$ par le sous-ensemble $I'_{\geq i'_0}$, on voit que $I'_M$ est cofinal dans $I'$, donc aussi dans $I$, puisque $I'$ est cofinal dans $I$ ; donc $I_M \supset I'_M$ est cofinal dans $I$, \quad \textit{q.e.d.}

Inversement, supposons que $I_M$ soit cofinal dans $I$, et prouvons $X \in M$. La difficulté ici provient du fait que la représentation $X = \varinjlim\limits_{i\in I_M} X_i$ n'est plus adaptée à $\pi$ -- $I_M$ n'est pas nécessairement stable dans $I$ par bornes supérieures (dans $I$) de parties filtrantes dans $I_M$, $I'$ de cardinal $\leq \pi$. En somme, il faut
\msp{page 54}
prouver que $M$ est stable par $\varinjlim$ grandes devant $\pi$ sans plus, et je doute maintenant que ça résulte des seules hypothèses faites ici, qui ne concernent que les systèmes inductifs où les morphismes de transition sont des monomorphismes.

Voici, je crois, un contre-exemple. Prenons
\[
M \quad = \quad \parbox[t]{0.6\textwidth}{ensemble des objets de $\cM$ isomorphes à un objet fixe $X_0 \in \Filt^{\pi}$, tel que $X_0$ n'ait pas d'endomorphisme qui soit un monomorphisme autre que l'identité.}
\tag*{(5.4)}
\]
Alors la condition 5.3 est satisfaite. Car :
\begin{enumerate}[nosep]
\item[\textit{a})] Si $X \in M$, \textit{i.e.} $X \simeq X_0$, alors $X \in \Filt^{\pi}$, donc $\Ss_{\pi}(X)$ admet $X$ comme objet final, et $J$, étant cofinal, contient cet élément, donc $J_M$ (qui le contient aussi) est cofinal.
\end{enumerate}
Et :
\begin{enumerate}[nosep]
\item[\textit{b})] Si $J_M$ est cofinal, l'hypothèse (5.4) implique que l'ensemble des morphismes de transition entre les $Z_{\alpha}$ ($\alpha \in J_M$) est formé d'isomorphismes, donc la limite inductive $X$ des $Z_{\alpha}$ en question est isomorphe à chacun d'eux, donc $X \simeq X_0$, \textit{i.e.} $X \in M$.
\end{enumerate}
Pourtant, $M$ n'est $\pi$-accessible que s'il est stable par $\varinjlim$ grandes devant $\pi$, ce qui implique qu'il est stable par \textit{facteurs directs.} Mais l'hypothèse
\msp{page 55}
(5.4) n'exclut nullement que $X_0$ admette un projecteur $p$ ($p^{2} = p$) non trivial. Celui-ci aura une image $X'$ dans $\cM$, qui sera facteur direct de $X_0$ non isomorphe à $X_0$, \textit{i.e.} non dans $M$.

Ces difficultés disparaissent si on suppose que $M^{\pi}$ est stable par $\varinjlim$ filtrantes indexées par des ensembles ordonnés $I$ de cardinal $\leq \pi$. En effet, dans ce cas, dans la démonstration précédente, le sous-ensemble $I_M$ est stable par bornes supérieures (dans $I$) de parties $J$ de $I_M$ filtrantes et de cardinal $\leq \pi$. Donc le système inductif $(X_i)_{i\in I_M}$ est encore $\pi$-adapté, et on peut lui appliquer le lemme 5.7, pour trouver un sous-ensemble cofinal $I'_M \subset I_M$, tel que les $X_i \longrightarrow X$ pour $i \in I'_M$ soient des monomorphismes. L'image $J$ de $I'_M$ dans $\Ss_{\pi}(X)$ est alors une partie de $\Ss_{\pi}(X)$ telle que $\varinjlim\limits_{\alpha\in J} Z_{\alpha} = X$, donc (par l'hypothèse faite sur $\pi$) cofinale dans $\Ss_{\pi}(X)$ (5.4), avec $Z_{\alpha} \in M^{\pi}$ pour tout $\alpha \in J$. Donc par 5.3 on a $X \in M$, \textit{q.e.d.}

\msp{page 56}
Il reste à prouver le lemme 5.7. Avec les hypothèses et notations du lemme, soit donc $i_0 \in I$, il faut trouver $i' \in I'$ tel que $i' \geq i_0$. Considérons l'image $X'_i$ de $X_i$ dans $X$, et

\dgrblock{dgrb751}{A two-point model of the retract obstruction}
\begin{dcomment}{development}{p55}{A two-point model of the retract obstruction}
Grothendieck's proposed counterexample asks for an object which has no proper monic self-map, but does have a nontrivial retract. The walking arrow provides such an object in an elementary ambient category, and also makes the failure visible in a perfectly adapted presentation. This gives a concrete completion of the mechanism proposed in (5.4).

\textbf{A class satisfying (5.3) but failing $b_\pi$.}\ours
Let $\mathcal M=\mathrm{Pos}$ be the category of small partially ordered sets and monotone maps. For an infinite cardinal $\pi$, use the filtration by underlying cardinality at most $\pi$, and let $M$ consist of the objects isomorphic to the two-element chain $P=\{0<1\}$. Then $M$ satisfies (5.3) but is not $\pi$-accessible.

\emph{Proof.} Monomorphisms in $\mathrm{Pos}$ are the injective monotone maps, so every monic endomorphism of $P$ is the identity. The constant map $P\to P$ with value $0$ is nevertheless an idempotent whose image is the one-point poset $1$. The ambient hypotheses are satisfied: a colimit is obtained from a disjoint union by generating the order relations and quotienting by the resulting equivalence relation, so a diagram of cardinal at most $\pi$ with $\pi$-small values has a $\pi$-small colimit. Subobjects and image quotients have the same cardinal bound. Filtered colimits have underlying set the filtered colimit of the underlying sets, and an order relation between two representatives holds exactly when it holds at some later common stage. This description proves $(e)$. A poset of cardinal at most $\pi$ has at most $\pi$ order relations; choosing representatives for its elements and making those relations hold at one common stage proves its $\pi$-accessibility. The eventual-equality part follows by the same argument on the at most $\pi$ elements.

To check (5.3), first suppose $X\cong P$. Every cofinal subset of $\operatorname{Sub}_{\pi}(X)$ contains its largest member $X$, so its good members are cofinal. Conversely, suppose $J^M$ is cofinal in a continuous cofinal subset $J$ of $\operatorname{Sub}_{\pi}(X)$. Then $X$ is the filtered colimit of the subobjects indexed by $J^M$. Each has type $P$, and every transition between them is an isomorphism. The index is nonempty and directed, so its colimit is again $P$. This establishes both directions of (5.3).

For a direct failure of the actual definition $b_\pi$, put $\lambda=\pi^+$ and index a diagram by the ordinals $\alpha<\lambda$. Take $X_\alpha=P$ at $0$ and at every successor ordinal, and $X_\alpha=1$ at every nonzero limit ordinal. Give each $X_\alpha$ its distinguished least element. For $\alpha<\beta$, let the transition map be the constant map to that element of $X_\beta$. These maps compose as required. The colimit is $1$, since every element is identified with a distinguished element at a later stage.

This diagram is perfectly $\pi$-adapted. The ordinal $\lambda$ is $\pi^+$-directed. Every nonempty directed subset of cardinal at most $\pi$ has a supremum below $\lambda$; if it has a maximum, the colimit is its final stage, and otherwise its supremum is a nonzero limit ordinal and both the colimit and that stage are $1$. Yet the good indices, namely $0$ and the successor ordinals, are cofinal, while the colimit is not in $M$. Thus $b_\pi$ fails.\hfill$\square$

The same phenomenon can be described as splitting an idempotent by a long filtered diagram. The continuous version above shows why replacing arbitrary presentations by subobject presentations loses precisely the needed information. We checked the definition in Chapter~XIII, manuscript pp.~85--97, and the treatment of idempotents in \citep[Expos\'e~I, 8.6.5 and 10.6]{SGA4I1972}. We did not find this specified two-point diagram in those sources; it is an elaboration of the manuscript's counterexample, not a claim of priority.
\end{dcomment}


% ---- from Der-ChXV-05-11-10_p036-040.tex
\printedpage{34}
% head: A. Grothendieck : Dérivateurs, Ch. XV, Édition des Universités de Montpellier II et Paris VII
% foot: Version 5 novembre 2010
l'inclusion
\[
\underbrace{X'_{i_0}}_{\in\, \Filt^\pi} \;\longrightarrow\; X \;\simeq\; \varinjlim X_i \;.
\]
Comme $X'_{i_0} \in \Filt^\pi$, donc est $\pi$-accessible, il existe un $i_1 > i_0$ tel que $X'_{i_0} \longrightarrow X$ se factorise par
\[
X'_{i_0} \;\longrightarrow\; X_{i_1} \;.
\]
Appliquant ce même résultat à $X_{i_1}$, $X'_{i_1}$ \textit{etc.}, on trouve une suite infinie de morphismes
\[
\begin{tikzcd}[column sep=small]
X_{i_0} \arrow[dr] & & X_{i_1} \arrow[dr] & & X_{i_2} \arrow[dr] & \\
& X'_{i_0} \arrow[ur] & & X'_{i_1} \arrow[ur] & & X'_{i_2} \quad \cdots
\end{tikzcd}
\]
telle que les morphismes composés de deux en deux soient, sur la ligne supérieure, les transitions
\[
X_{i_m} \longrightarrow X_{i_n} \qquad (\; i_m \;<\; i_n \;) \;,
\]
et sur la ligne inférieure, les inclusions
\[
X'_{i_m} \hookrightarrow X'_{i_n} \qquad (\; \hookrightarrow X \;) \;.
\]
Soit $i_\omega = \sup_n i_n$ dans $I$, donc
\[
X_{i_\omega} \;=\; \varinjlim_{n} X_{i_n} \;\simeq\; \varinjlim_{n} X'_{i_n} \;\simeq\; X'_{i_\omega} \;,
\]
donc $X_{i_\omega} \longrightarrow X$ est un monomorphisme.

\msp{page 57}

La dernière assertion de 5.7 résulte aussitôt du fait qu'une limite inductive filtrante de monomorphismes dans $\M$ est un monomorphisme (condition $(e)$, p.~46).

On peut préciser la proposition 5.6 obtenue ainsi :

\textbf{Proposition 5.7}. \texttt{[Deux fois 5.7.]} \textit{Soient $\M$ une catégorie satisfaisant les conditions} (a), (b), (c), (d) (p.~44, 45) et (e) (p.~46). \textit{Soit $\M$ une partie de} $\Ob\M$\textit{, stable par isomorphie, $\pi$ un cardinal tel que les objets de} $\Filt^\pi$ \textit{soient $\pi$-accessibles. Les conditions suivantes sont équivalentes :} % [sic] Soit $\M$ (calligraphique) -> Soit $M$, la partie de Ob M est notée M partout ailleurs

\begin{enumerate}[label={}]
\item[a)] \textit{$M$ est $\pi$-accessible, et stable par $\varinjlim$ filtrantes.}
\item[b)] \textit{$M$ est $\pi$-accessible, et $M \cap \Filt^\pi$ est stable par $\varinjlim$ filtrante de cardinal $\leq \pi$.}
\item[c)] \textit{Il existe une partie $M^\pi$ de} $\Filt^\pi$\textit{, stable par $\varinjlim$ filtrantes de cardinal $\leq \pi$ et telle que pour tout $X \in \Ob\M$, on ait $X \in M$ si et seulement si l'ensemble $J_M$ des $Z_\alpha \in J$ tels que $Z_\alpha \in M^\pi$ est cofinal dans $J$.}
\end{enumerate}

(\textit{De plus, sous la condition} (c), \textit{on a}
\printedpage{35}
% head: A. Grothendieck : Dérivateurs, Ch. XV, Édition des Universités de Montpellier II et Paris VII
% foot: Version 5 novembre 2010
\msp{page 58}
$M^\pi = M \cap \Filt^\pi$.)
\begin{enumerate}[label={}]
\item[c$'$)] \textit{Il existe une partie $M_0$ de} $\Filt^\pi$ \textit{telle que $M$ soit la plus petite partie de $\M$ stable par limites inductives filtrantes.} (\textit{De plus, on peut alors prendre $M_0$} petite.)\footnote[22]{\textbf{N.B.} On a
\[
\begin{tikzcd}[ampersand replacement=\&]
(a) \arrow[r, Leftrightarrow] \arrow[d, Rightarrow] \& (b) \arrow[d, Rightarrow] \\
(c') \arrow[r, Rightarrow, "?"] \& (c)
\end{tikzcd}
\quad ,
\]
sans conditions sur $\pi$ (sauf $\pi \geq \pi_0$).}.
\end{enumerate}

\dgrblock{dgrb752}{Making the image factorization a transition}
\begin{dcomment}{exposition}{p57}{Making the image factorization a transition}
Lemma~5.7 is the comparison that makes the repaired recognition theorem work. Its diagonal construction is sound once one supplies an eventual-equality step. That step uses the accessibility of $X_i$ as well as the accessibility of its image; factorization of the image alone does not identify the resulting composite with a transition of the original diagram.

Write $q_i:X_i\twoheadrightarrow U_i$ for the image factorization of the structural map to $X$. Given $i_0$, accessibility of $U_{i_0}\in\Filt^\pi$ gives $j\geq i_0$ and a factorization
\[
 U_{i_0}\xrightarrow{a}X_j\longrightarrow X.
\]
The two maps $a q_{i_0}$ and $X(i_0\leq j)$ from $X_{i_0}$ to $X_j$ have the same composite to $X$. Accessibility of $X_{i_0}$ therefore gives a further stage $i_1\geq j$ at which these maps agree. After replacing $a$ by its composite to $X_{i_1}$, the upper composite in Grothendieck's zigzag really is the transition $X_{i_0}\to X_{i_1}$. Its lower composite is the inclusion $U_{i_0}\hookrightarrow U_{i_1}$, since both maps are over $X$ and the latter image maps monomorphically to $X$.

Repeating gives an interleaving of the chains $(X_{i_n})$ and $(U_{i_n})$. A cocone on either chain extends uniquely to a cocone on the other, so their colimits agree. Let $i_\omega=\sup_n i_n$. Perfect adaptation identifies $X_{i_\omega}$ with that colimit, while $(e)$ makes the map $\varinjlim_nU_{i_n}\to X$ monic. Thus $i_\omega$ is a monomorphic stage above $i_0$. The same use of $(e)$ and continuity proves that the monomorphic stages are closed under directed suprema of cardinal at most $\pi$. The choices may be non-strictly increasing; the proof does not require that the index have no maximal element.

There is a second detail in applying this lemma on manuscript page~53. The image of a continuous cofinal map of directed posets need not itself be closed under directed suprema, as the example below shows. Consequently, the image $J$ of the monomorphic stage indices cannot simply be inserted into (5.3) without checking its closure requirement.

The correction already adopted on manuscript page~55 supplies the repair. If $X\in M$, the good small subobjects form a cofinal subset closed under small directed unions, because $M^\pi$ has that closure. Starting with any monomorphic stage, choose a good small subobject containing it, then a later monomorphic stage containing that subobject, and continue alternately. At the countable supremum the stage and the union of the good subobjects agree, so the resulting stage is good. Hence good stages are cofinal. In the converse direction, if good stages are cofinal, their closure under small directed suprema makes their restricted diagram perfectly adapted. Applying Lemma~5.7 to it produces cofinally many good monomorphic stages; their subobjects are cofinal in $\operatorname{Sub}_{\pi}(X)$, so (5.1) gives $X\in M$. This proves the corrected assertion without imposing a closure property on the raw image of an indexing map.
\end{dcomment}


\dgrblock{dgrb753}{A continuous image need not be closed}
\begin{dcomment}{development}{p57}{A continuous image need not be closed}
The distinction just used is easy to miss because the indexing map preserves every supremum which can be formed upstairs. The difficulty is that a directed family of image values need not have coherently ordered representatives. The following example isolates that point in sets.

\textbf{A perfectly adapted diagram with nonclosed image.}\ours
Let $I$ be the countable subsets of the successor ordinals below $\omega_1$, ordered by inclusion. Choose $e\notin\omega_1$, put $X=\omega_1\cup\{e\}$, and define subobjects of $X$ by
\[
 X_S=\bigcup_{\alpha\in S}\alpha
       \ \cup\ \{e\mid |S|\geq2\}.
\]
With inclusion maps this is a perfectly $\aleph_0$-adapted presentation of $X$, and its image is cofinal among the countable subobjects of $X$, but that image is not closed under countable directed suprema.

\emph{Proof.} The union of countably many countable subsets is countable, so $I$ is $\aleph_1$-directed, and countable directed suprema in $I$ are unions. If $\mathcal K\subset I$ is directed, its union contains two distinct ordinals exactly when some member of $\mathcal K$ contains both: direct the two witnesses to a common member. It follows that
\[
 X_{\bigcup\mathcal K}=\bigcup_{S\in\mathcal K}X_S.
\]
Every $X_S$ is countable, and the union of the whole diagram is $X$, proving perfect adaptation. Any countable subset of $X$ is contained in $X_S$ for a two-element $S$ containing a successor ordinal above its ordinal elements; this proves cofinality. Finally every positive finite ordinal $n$ is $X_{\{n\}}$. Their directed supremum is $\omega$ without $e$. No $S$ with two or more elements gives this value, because its value contains $e$; a singleton gives a successor ordinal, and the empty set gives the empty set. Thus $\omega$ is not in the image.\hfill$\square$

The literature check here was bounded to the adapted-presentation discussion in Chapter~XIII, manuscript pp.~85--97, and SGA~4, Expos\'e~I, 9.1--9.3 and 9.12--9.16. This explicit example supplies a missing verification in the present argument; we make no assertion that the observation is new.
\end{dcomment}


\textbf{Corollaire 5.8.} \textit{Les applications}
\[
M \;\longmapsto\; M^\pi = M \cap \Filt^\pi \;,
\]
\[
M^\pi \quad\longmapsto\quad \parbox[t]{0.5\textwidth}{\textit{l'ensemble des $X \in \Ob\M$ tels que l'ensemble $\Ss_\pi^{M^\pi}(X)$ des $Z_\alpha \in \Ss_\pi(X)$ qui sont dans $M^\pi$ soit cofinal} (= \textit{plus petit sous-ensemble de $\Ob\M$ stable par $\varinjlim$ filtrantes})}
\]
\textit{établissent des bijections, inverses l'une de l'autre, entre l'ensemble des $M \subset \Ob\M$ satisfaisant les conditions équivalentes de la proposition} 5.7\textit{, et l'ensemble des parties $M^\pi$ de} $\Filt^\pi$ \textit{stables par $\varinjlim$ filtrantes de cardinal $\leq \pi$.}

\textsc{Démonstration de} 5.7. L'implication $(b) \Longrightarrow (c)$ est tautologique, compte tenu de la définition de la $\pi$-accessibilité de $M$, et de la proposition 5.3. Et $(a) \Longrightarrow (b)$ provient de la stabilité de $\Filt^\pi$ par $\varinjlim$ filtrantes de cardinal $\leq \pi$ (condition $(d)$).

Il reste donc à prouver $(c) \Longrightarrow (a)$. (C'est ici qu'on doit utiliser l'hypothèse particulière sur $\pi$.) Une fois ceci vu, on saura l'équivalence de $(a)$, $(b)$, $(c)$, et d'autre part, il est immédiat que $(a) \Longrightarrow (c')$, en
\msp{page 59}
prenant $M_0 = M \cap \Filt^\pi$. D'autre part, $(c') \Longrightarrow (c)$ (sans hypothèses sur $\pi\,$?), car on peut supposer que la partie $M_0$ de $\Filt^\pi$ est stable par $\varinjlim$ filtrantes de cardinal $\leq \pi$. Mais alors, connaissant l'équivalence de $(a)$, $(b)$ $(c)$, d'où aussitôt le corollaire, il s'ensuit que l'enveloppe stable par $\varinjlim$ filtrantes de $M_0$ dans $\M$ n'est autre que la partie accessible déduite de $M_0$ dont il est question dans $(c)$, et on gagne. % [sic] $(b)$ $(c)$ -> $(b)$, $(c)$, virgule manquante

Donc tout revient à prouver que, si tous les objets de $\Filt^\pi$ sont accessibles, alors $(c)$ implique $(a)$. Montrons d'abord que $(c)$ implique le critère d'appartenance (5.3) (p.~51), donc par la proposition 5.6 “revue et corrigée” dans la note 21, il s'ensuit déjà que $(c) \Longrightarrow (b)$, et il restera à établir $(b) \Longrightarrow (a)$, \textit{i.e.} que $(b)$ implique la stabilité de $M$ par $\varinjlim$ filtrantes quelconques (pas seulement grandes devant $\pi$).

Que $(c)$ implique (5.3) va résulter du lemme suivant, appliqué à $I = \Ss_\pi(X)$ :

\textbf{Lemme 5.8.} \texttt{[Deux fois 5.8.]} \textit{Soit $I$ un ensemble ordonné} $\pi$-adapté, i.e. \textit{filtrant et tel que toute partie filtrante $K$ de $I$, de cardinal $\leq \pi$, ait une borne supérieure $i_K$ dans $I$. Je dis qu'une partie $J$ de $I$ est} $\pi$-cofinale, \textit{si elle est cofinale, et si pour toute partie filtrante $K$ de $J$ de cardinal $\leq \pi$, sa borne supérieure $i_K$ dans $I$ est}
\printedpage{36}
% head: A. Grothendieck : Dérivateurs, Ch. XV, Édition des Universités de Montpellier II et Paris VII
% foot: Version 5 novembre 2010
\msp{page 60}
\textit{dans $J$.} (\textit{Il s'ensuit donc que $J$ est lui aussi $\pi$-adapté, pour la structure d'ordre induite.}) \textit{Ceci dit, l'intersection de deux parties $\pi$-cofinales de $I$ est encore $\pi$-cofinale.}

Il suffit de voir qu'elle est cofinale, car sa stabilité par bornes supérieures dans $I$ de parties filtrantes de cardinal $\leq \pi$ est tautologique (si elle est vraie pour des $J_\alpha$, elle est vraie pour $\bigcap_\alpha J_\alpha$). Soit $i_0 \in I$, et construisons un élément de $J \cap J'$ qui majore $i_0$. On construit de proche en proche une double suite
\[
\begin{tikzcd}
i_1 \arrow[dr] & & i_2 \arrow[dr] & & \cdots \\
& i'_1 \arrow[ur] & & i'_2 \arrow[ur] &
\end{tikzcd}
\]
avec les $i_\alpha$ dans $J$, les $i'_\alpha$ dans $J'$, et $i_1 \geq i_0$. On a donc
\[
\sup i_\alpha \;=\; \sup i'_\alpha \;,
\]
et si $i$ est la valeur commune des deux membres, on a $i \in J \cap J'$.

Ceci vu, pour prouver la validité de (5.3), on note que la partie $J' = \Ss_\pi^{(M^\pi)}(X)$ de $I = \Ss_\pi(X)$, formée des $Z_\alpha \subset X$ tels que $Z_\alpha \in M^\pi$, est stable par bornes supérieures de parties filtrantes de cardinal $\leq \pi$, vu que $M^\pi$ est stable
\msp{page 61}
par $\varinjlim$ de cardinal $\leq \pi$. Donc si $J$ est une partie $\pi$-cofinale de $\Ss_\pi(X)$, $J \cap J'$ est cofinale dans $J$, \textit{i.e.} dans $I$, si et seulement si $J'$ est cofinale dans $I$ (\textit{i.e.} si et seulement si $X \in M$, par la condition $(c)$). Car $J \cap J'$ cofinale $\Longrightarrow$ $J'$ cofinale tautologiquement, et $J'$ cofinale implique $J'$ $\pi$-cofinale, donc $J \cap J'$ $\pi$-cofinale par le lemme 5.8, donc cofinale.

\dgrblock{dgrb754}{Closed cofinal families and the four descriptions}
\begin{dcomment}{exposition}{p60}{Closed cofinal families and the four descriptions}
The alternating construction in Lemma~5.8 is a small version of the closed-unbounded argument familiar in set theory. Its purpose here is representation independence. Two canonical-looking systems of small subobjects need not meet cofinally merely because each is cofinal. Closure under countable directed suprema is what allows an alternating sequence to finish in both systems.

For example, the even and odd natural numbers are cofinal in $\mathbb N$ and disjoint. That order lacks the required countable suprema. In a $\pi$-adapted order, with $\pi$ infinite, start above any prescribed element and alternate choices in two closed cofinal subsets. The supremum of the even subsequence equals that of the odd subsequence, and belongs to both. The intersection is therefore cofinal; closure of the intersection is immediate. This gives exactly the assertion used in the manuscript, without any assumption that $\pi$ itself is regular.

We can now organize the corrected Proposition~5.7. Let $S\subset\Filt^\pi$ be invariant under isomorphism and closed under filtered colimits with at most $\pi$ arrows. Define $M_S$ by cofinality of its $S$-subobjects in $\operatorname{Sub}_{\pi}(X)$. Under $(e)$ these good subobjects are closed under small directed suprema. Their intersection with any other closed cofinal family is cofinal by the preceding argument. Thus $M_S$ satisfies (5.3), and the corrected Proposition~5.6 identifies it as a $\pi$-accessible class when the $\Filt^\pi$-objects are $\pi$-accessible. Proposition~5.9, in this same range of cardinals, supplies closure under all filtered colimits.

Conversely, a $\pi$-accessible class closed under all filtered colimits has bounded part $S=M\cap\Filt^\pi$ with the required small-colimit closure, by $(d_0)$, and (5.1) reconstructs $M$. These two operations are inverse because on a bounded object cofinality is equivalent to containing the largest subobject, namely the object itself. This is Corollary~5.8.

The description by an arbitrary generating set $M_0\subset\Filt^\pi$ follows by first closing $M_0$ under the filtered colimits of cardinal at most $\pi$ which exist inside the bounded level, and under isomorphism. Call the resulting class $S$. It remains essentially small because it is contained in $\Filt^\pi$. The class $M_S$ is filtered-colimit closed and contains $M_0$. Any filtered-colimit closed class containing $M_0$ contains $S$, and then contains $M_S$, since every object of $M_S$ is a filtered colimit of its good small subobjects. Hence $M_S$ is precisely the filtered-colimit closure of $M_0$.

The printed condition $(c')$, and the matching parenthesis in Corollary~5.8, omit the words saying that the smallest closed class must \emph{contain} the proposed generators. Those words are indispensable: without them the smallest class closed under nonempty filtered colimits is the empty class. The intended statement and its proof are unambiguous from the surrounding discussion. The stronger question in footnote~22 concerns retaining the same rank without the special assumption on the bounded objects; it is not resolved merely by the proof at the admissible ranks used in the proposition.
\end{dcomment}


Reste à prouver seulement que $(b) \Longrightarrow (a)$, \textit{i.e.} la stabilité de $M$ par limites inductives filtrantes quelconques, moyennant $(b)$. Ici à nouveau, cela ne doit pas utiliser l'hypothèse particulière sur $\pi$, ni même l'hypothèse d'exactitude $(e)$ sur $\M$, mais être une histoire générale d'accessibilité :

\textbf{Proposition 5.9.} \textit{Soit $\M$ une catégorie satisfaisant} (a), (b), (c), (d$_0$) (pages 44, 45), \textit{$\pi \geq \pi_0$ un cardinal, $M$ une partie $\pi$-accessible de} $\Ob\M$, i.e. \textit{satisfaisant la condition $b_\pi$ de} XIII, p.~88. \textit{Alors $M$ est stable par $\varinjlim$ filtrantes si et seulement si $M^\pi \overset{\text{déf}}{=} M \cap \Filt^\pi$ est stable par $\varinjlim$ filtrantes de cardinal $\leq \pi$.}

Je n'ai pas envie d'expliciter ici la démonstration, ayant l'impression pénible de refaire à perpette la même démonstration (ici et dans XIII). Je devrai revenir là-dessus dans une situation plus décantée.

\msp{page 62}

Je vais quand même esquisser la démonstration. Soit $X = \varinjlim\limits_{I} X_i$, $I$ filtrant, les $X_i$ dans $M$, il faut prouver que $X \in M$ (moyennant l'hypothèse de stabilité sur $M^\pi = M \cap \Filt^\pi$).
\printedpage{37}
% head: A. Grothendieck : Dérivateurs, Ch. XV, Édition des Universités de Montpellier II et Paris VII
% foot: Version 5 novembre 2010
\[
\begin{tikzcd}
& \Ci \arrow[r] \arrow[d] \arrow[dl, dotted, "\mathfrak{z}"'] \arrow[dr, phantom, "\text{\footnotesize cart.}"] & I \arrow[d, "\varphi"] \\
\underline{M}^\pi & \underline{M}^\pi/\underline{M} \arrow[l, dotted, "\text{\footnotesize “source”}"] \arrow[r, "\text{\footnotesize “but”}"'] & \underline{M}
\end{tikzcd}
\] % ?? uncertain: label of the dotted arrow -- printed glyph resembles a Fraktur z (\mathfrak{z}); possibly meant \exists

$\Ci$ est une catégorie cofibrée sur $I$, dont les fibres sont
\[
\Ci_i \;\simeq M^\pi/X_i \;,
\]
or comme $X_i \in M$, on voit, à cause de la condition de $\pi$-accessibilité, que $M^\pi/X_i$ est \textit{filtrante.} On va le dégager en

\textbf{Lemme 5.10.} \textit{Soit $M$ une partie $\pi$-accessible de $\M$. Alors pour tout $Y \in M$, $M^\pi/Y$ est une catégorie} filtrante, \textit{et $\Ss_\pi(Y)$ y est cofinale.}

\msp{page 63}

Considérons le foncteur canonique
\[
I \;=\; \Ss_\pi(Y) \xrightarrow{\;\varphi\;} \underline{M}^\pi/Y \;=\; I' \;,
\]
où la catégorie source $I$ est filtrante. Par SGA 4 I 8.1.3, ce foncteur est cofinal si et seulement s'il satisfait deux conditions F\,1, F\,2, et alors $I'$ est également filtrante. La condition F\,1 est la condition habituelle que tout $i' \in \Ob I'$ est “majoré” par un $\varphi(i)$, \textit{i.e.} tout $Z \longrightarrow Y$, avec $Z \in M^\pi$, est “majoré” par un \textit{monomorphisme} $Z_\alpha \hookrightarrow Y$, avec $Z_\alpha \in M^\pi$, ce qui résulte du fait que c'est majoré par $Z_\alpha \hookrightarrow Y$ avec $Z_\alpha \in \Filt^\pi$, \textit{i.e.} $Z_\alpha \in \Ss_\pi(Y)$, et que parmi ces $Z_\alpha$ ceux qui sont dans $M^\pi$ sont cofinaux (vu que $Y \in M$). L'autre condition\footnote[23]{Comme $I$ est filtrante.} est qu'une double flèche $i' \overset{f'}{\underset{g'}{\rightrightarrows}} \varphi(i)$ peut
\msp{page 64}
s'égaliser par une flèche $\varphi(h)$, avec $h : i \longrightarrow j$. Mais en fait, pour un $i' \in \Ob I'$, \textit{i.e.} un $Z \longrightarrow Y$ donné ($Z \in M^\pi$), et un $i \in \Ob I$, \textit{i.e.} un $Z_\alpha \hookrightarrow Y$ donné ($Z_\alpha \in M^\pi$), il existe au plus \textit{une} flèche dans $I' = \underline{M}^\pi/Y$ de $Z \longrightarrow Y$ dans $Z_\alpha \hookrightarrow Y$. Donc la condition F\,2 est trivialement satisfaite (on prend $h = \mathrm{id}_i$).

\textbf{Corollaire 5.11.} \textit{Si $(X_i)_{i \in I}$ est un système inductif filtrant, avec les $X_i$ dans $M^\pi$, alors la catégorie $\Ci = \underline{M}^\pi/\underline{M} \times_{\underline{M}} (I,\varphi)$, construite page} 62, \textit{est filtrante, et le foncteur canonique} (\textit{fibrant}) \textit{$\Ci \longrightarrow I$ est cofinal.}

Que $\Ci$ soit filtrante et $\Ci \longrightarrow I$ cofinal résulte formellement du fait que $I$ l'est, que les fibres de $\Ci$ le sont, et que $\Ci$ est cofibrée sur $I$. (Mérite un lemme, à toutes fins utiles.)

\textbf{N.B.} On peut remplacer $\Ci$ par $\Ci_0$ dont les fibres sont les
\[
\Ss_\pi^{M^\pi}(X_i) \;\subset\; \underline{M}^\pi \;; \qquad \text{cofinale dans } \Ci_i = \underline{M}^\pi/X_i \;,
\]
\msp{page 65}
on trouve une sous-catégorie de $\Ci$ cofinale dans $\Ci$, et qui est cette fois ordonnée tout comme $I$.

\dgrblock{dgrb755}{The slice category that must be integrated}
\begin{dcomment}{exposition}{p64}{The slice category that must be integrated}
The final argument returns from subobjects to maps, because maps assemble functorially when the target moves. Its fibre over $Y\in M$ is $M^\pi/Y$, the category of maps $A\to Y$ with $A\in M^\pi$. The poset naturally included in this fibre is
\[
 \operatorname{Sub}_{\pi}^{M}(Y)
 =\{U\hookrightarrow Y\mid U\in M\cap\Filt^\pi\}.
\]
The unrestricted $\Ss_\pi(Y)$ printed in Lemma~5.10 and its first display must be restricted in this way. An arbitrary small subobject of an object of $M$ need not itself belong to $M$, so the printed unrestricted functor is not defined in general.

Under the image and continuity hypotheses used earlier in this section, the corrected assertion follows directly from (5.1). Given $A\to Y$, its image is small by $(d_2)$, and is contained in a good small subobject by cofinality. Hence every object of $M^\pi/Y$ maps to an object of $\operatorname{Sub}_{\pi}^{M}(Y)$. There is at most one map over $Y$ into any fixed subobject. The good subobjects form a directed poset, and their inclusion in the slice is final: the relevant comma category is nonempty, and two of its objects have a common upper bound. This also proves that $M^\pi/Y$ is filtered. Its source diagram has colimit $Y$, because the final subobject diagram does.

For a filtered diagram $i\mapsto Y_i$ in $M$, the total category $\mathcal E$ has objects $(i,A,a)$, where $A\in M^\pi$ and $a:A\to Y_i$. A morphism to $(j,B,b)$ is a pair $(u,v)$ with $u:i\to j$, $v:A\to B$, and $bv=Y(u)a$. Postcomposition gives an opfibration $\mathcal E\to I$. The adjective ``fibrant'' in Corollary~5.11 has the wrong variance for this construction; no pullback closure of $M^\pi$ is assumed.

The total category is filtered. Two objects can first be moved to a common base index and then given a common successor in that fibre. For two parallel arrows, first equalize their base arrows in $I$, then equalize the resulting two vertical arrows in the target fibre. The same argument shows that $\mathcal E\to I$ is final. More particularly, a cocone on the source objects $A$ of $\mathcal E$ is exactly a compatible family of cocones on the fibres. Taking their colimits gives
\[
 \varinjlim_{(i,A,a)\in\mathcal E}A
 \cong \varinjlim_{i\in I}
          \left(\varinjlim_{(A\to Y_i)\in M^\pi/Y_i}A\right)
 \cong \varinjlim_iY_i.
\]
This is an ordinary colimit calculation; no homotopy-cofinality theorem is being used.

Corollary~5.11 prints $Y_i\in M^\pi$, a stronger hypothesis than the argument needs. The same proof works for $Y_i\in M$, and this is the form required on manuscript page~62. If the base index is a poset, the full subcategory consisting of good subobjects in each fibre is again a poset and is final in $\mathcal E$. Its projection need not remain an opfibration: the image of a good subobject after postcomposition need not be good. The proof requires only finality of this subcategory, which the preceding factorization through good subobjects establishes.
\end{dcomment}


\printedpage{38}
% head: A. Grothendieck : Dérivateurs, Ch. XV, Édition des Universités de Montpellier II et Paris VII
% foot: Version 5 novembre 2010
On a
\[
X_i \;=\; \varinjlim_{Z_\alpha \in \Ci_i} Z_\alpha \;,
\]
donc
\[
X \;=\; \varinjlim_{\alpha \in \Ci} Z_\alpha \;,
\]
au total on est arrivé à représenter $X$ comme $\varinjlim$ filtrante d'éléments $Z_\alpha$ qui sont (non seulement dans $M$, mais même) dans $M^\pi$.

Changeons la notation, en prenant
\[
X \;=\; \varinjlim_{i \in I} X_i \;, \qquad I \text{ ensemble ordonné filtrant, } X_i \in M^\pi \;\; \forall\, i \in I \;,
\]
soit $I'$ l'ensemble des parties \textit{filtrantes} $J \subset I$ telles que $\mathrm{card}\, J \leq \pi$, $I'$ ordonné par inclusion. Alors $I'$ est grand devant $\pi$, et on a un système inductif
\[
(X_J)_{J \in I'} \;, \qquad X_J \;=\; \varinjlim_{i \in J} X_i \;.
\]
Par l'hypothèse de stabilité de $M^\pi$ (qui n'a pas encore servi), les $X_J$ sont dans $M^\pi$. On a donc
\[
X \;=\; \varinjlim_{J \in I'} X_J \;, \qquad \text{les } X_J \text{ dans } M^\pi\text{, } I' \text{ grande devant } \pi \;,
\]
donc $X \in M$ puisque $M$ est $\pi$-admissible, \quad \textit{q.e.d.}

\msp{page 66}

\dgrblock{dgrb756}{Completing the filtered-colimit closure argument}
\begin{dcomment}{exposition}{p65}{Completing the filtered-colimit closure argument}
The last reindexing in Proposition~5.9 is the right way to pass from bounded filtered colimits to arbitrary ones. There is, however, a scope issue in the preceding reduction. The proposition drops the image hypotheses $(d_1),(d_2)$ and the exactness hypothesis $(e)$, whereas the proof of Lemma~5.10 as printed uses the earlier subobject presentation. We give a complete proof under the additional assumption that every $\Filt^\pi$-object is $\pi$-accessible. This is the assumption already present in Proposition~5.7 and suffices for every application to the four equivalent descriptions above. The printed general assertion for an arbitrary rank is not certified by the image proof alone.

Put $A=M\cap\Filt^\pi$, and suppose that $A$ is closed under filtered colimits indexed by categories with at most $\pi$ arrows. When forming slices, replace this essentially small class by a small full set of representatives; this does not change their colimits. The necessary direction of Proposition~5.9 is immediate from $(d_0)$, so consider a filtered diagram $Y_i\in M$ with colimit $Y$. For each $Y_i$, choose a perfectly $\pi$-adapted presentation by objects of $\Filt^\pi$. Such presentations exist under $(b),(d_0)$: apply the reindexing by small filtered subcategories given below to any bounded generating presentation. Condition $b_\pi$ makes its $A$-stages cofinal. Thus $Y_i$ is a $\pi^+$-directed colimit of objects of $A$.

Smallness of the domains now supplies the slice argument without images. Any finite diagram in $A/Y_i$ factors through one stage of the chosen $A$-presentation: factor the finitely many structural maps, choose a common stage, and then pass further until its finitely many commutativity equations hold. The same reasoning works for a diagram with at most $\pi$ arrows. Hence $A/Y_i$ is $\pi^+$-filtered. The chosen presentation is final in it, since every map from an $A$-object factors through a stage and two factorizations become equal at a later stage. In particular
\[
 Y_i\cong\varinjlim_{A/Y_i}A.
\]
The opfibration argument in the preceding block then expresses $Y$ as a filtered colimit of $A$-objects, say $Y=\varinjlim_{e\in E}A_e$ with $E$ small and filtered.

To make this presentation perfectly adapted, let $P$ be the poset of filtered subcategories $J\subset E$ with at most $\pi$ arrows. Every set of at most $\pi$ arrows of $E$ is contained in such a subcategory: repeatedly adjoin identities, composites, cocones for finite diagrams and equalizers of parallel pairs. There are at most $\pi$ choices at each step, and countably many repetitions suffice, since $\pi$ is infinite. This also shows that $P$ is $\pi^+$-directed. A directed family of at most $\pi$ members of $P$ has its union as supremum; that union is still filtered and has at most $\pi$ arrows.

Define $A_J=\varinjlim_{e\in J}A_e$. By the assumed bounded closure, $A_J\in A$. Colimits commute with this directed union of indexing categories: a cocone on the union is exactly a compatible family of cocones on the pieces. Consequently
\[
 Y\cong\varinjlim_{J\in P}A_J,
 \qquad
 A_{\bigcup K}\cong\varinjlim_{J\in K}A_J
 \quad(K\subset P\text{ directed},\ |K|\leq\pi).
\]
This is a perfectly $\pi$-adapted presentation with every stage in $M$. Applying $b_\pi$ gives $Y\in M$, which completes the proof in the stated range.

The proof also makes the remaining issue precise. To remove the extra presentability assumption one needs, under the proposition's weaker hypotheses, a replacement for the slice factorization and density argument. The statement on manuscript page~44 about generation by strict epimorphisms would provide such density if established at the required bounded level, but it should not be inferred merely from a choice of colimit presentations. Density and generation under colimits are distinct notions. The earlier discussion already flags this difficulty: Chapter~XIII, manuscript p.~92, questions the asserted strict generation, and footnote~44 on manuscript p.~93 explicitly identifies the tacit hypothesis that the bounded objects are $\pi$-accessible \citep{GrothendieckDerivateursXIII}. We have not settled that unrestricted fixed-rank step here; this is a limit of the present verification, not a claim that it is an open problem in the literature.
\end{dcomment}


\textbf{Définition 5.12.}
\begin{enumerate}[label={}]
\item[1)] Soit $\M$ une catégorie satisfaisant les conditions $(a)$, $(b)$, $(c)$, $(d)$ (p.~44-46). Soit $\pi \geq \pi_0$ un cardinal. Une partie $M$ de $\Ob\M$ est dite \textit{L-$\pi$-accessible,} si elle est $\pi$-accessible et stable par petites $\varinjlim$ filtrantes\footnote[24]{On dit alors que $M$ est \textit{L-stable.}}, ou ce qui revient au même (5.9), si elle est $\pi$-accessible et si $M^\pi = M \cap \Filt^\pi$ est stable par $\varinjlim$ filtrantes de cardinal $\leq \pi$. (\textbf{N.B.} Si les objets de $\Filt^\pi$ sont $\pi$-accessibles, 5.8 nous dit que cela signifie aussi qu'il suffit pour cela qu'il existe une partie $M_0$ de $\Filt^\pi$ telle que $M$ soit la L-enveloppe (enveloppe stable par $\varinjlim$ filtrantes) de $M_0$ -- condition qui est de toutes façons toujours nécessaire, en prenant $M_0 = M^\pi$.)
\item[2)] Soit $M$ une partie de $\Ob\M$. On dit qu'elle est \textit{L-accessible,} si elle satisfait à l'une des quatre conditions équivalentes suivantes :
\begin{enumerate}[label={}]
\item[\textit{a})] Il existe un cardinal $\pi$ tel que $M$ soit L-$\pi$-accessible.
\item[\textit{a}$'$)] $M$ est accessible et L-stable.
\item[\textit{b})] Il existe une petite partie $M_0$ de $\Ob\M$, telle que $M$ soit la L-enveloppe de $M_0$\footnote[25]{Cette condition ne dépend pas de la filtration cardinale choisie de $\M$.}.
\item[\textit{b}$'$)] Il existe une partie $M_0$ de $\M$ contenue dans une partie $\Filt^\pi$ (ou ce qui revient au même, telle que la sous-catégorie pleine $\underline{M}_0$ de $\M$ soit essentiellement petite, \textit{i.e.} équivalente à une petite catégorie, telle que $M$ soit la \textit{L}-enveloppe de $M_0$).
\end{enumerate}
\end{enumerate}


\dgrblock{dgrb757}{Filtered-colimit classes and modern accessible categories}
\begin{dcomment}{later}{p66}{Filtered-colimit classes and modern accessible categories}
Definition~5.12 isolates the feature needed for the intended application: a class is generated by bounded objects and is closed under \emph{all} filtered colimits. At the good cardinals used in Proposition~5.7, its full subcategory is $\pi^+$-accessible in the current sense, and its inclusion into $\mathcal M$ preserves filtered colimits. Indeed its bounded objects are $\pi^+$-presentable even in the ambient category, every object has a $\pi^+$-directed presentation by them, and the required colimits exist in the full subcategory because they are computed in $\mathcal M$. This is the familiar small-generation description of accessible categories; see \citep[Definition~2.1 and Theorem~2.26]{AdamekRosicky1994}. The manuscript's $b_\pi$ also controls every continuous bounded presentation, so one should not identify its fixed-rank definition with the modern definition without the comparison just proved.

A simple example shows that L-accessibility does not mean local presentability of the full subcategory. In $\mathrm{Set}$, let $M$ be the nonempty sets. Its bounded part consists of the nonempty sets of cardinal at most $\pi$ and is closed under filtered colimits of cardinal at most $\pi$. For a nonempty $X$, every small subset can be enlarged to a nonempty small subset; for $X=\varnothing$ there is no good subobject at all. Thus the small-subobject criterion recovers exactly $M$. Filtered colimits of nonempty sets are nonempty, so this is an L-accessible class. Its full subcategory has no initial object: a singleton has more than one map to a two-element set, and a larger set does not have a unique map either. It is therefore accessible but not cocomplete.

Under the hypotheses of Proposition~5.7, the statement independent of a chosen filtration is the existence of an essentially small family whose filtered-colimit closure is the given class. To pass from a small family to a bounded one, bound the filtration indices of its representatives. To pass back, take representatives of the bounded level. At a sufficiently large cardinal satisfying the accessibility requirement, Proposition~5.7 identifies this closure with the class detected by cofinally many good small subobjects. The cardinal filtration is then a means of making the construction uniform; the resulting notion of L-accessibility does not retain that choice.

There is a final scope convention to keep visible. Definition~5.12 restates $(a),(b),(c),(d)$ but omits $(e)$. The definition of L-$\pi$-accessibility itself makes sense with those hypotheses. Its cited classification through Proposition~5.7, however, was proved using $(e)$, and our small-subobject argument retains that hypothesis. We have not supplied a proof removing it from the classification, just as we have not removed the rank restriction from Proposition~5.9. These qualifications concern the claimed equivalences under weaker hypotheses, not the definition or the preceding example.

This is the part of the chapter which is closest to the later theory of accessible and accessibly embedded subcategories. It supplies a method for obtaining a small amount of generating data, while the earlier factorization argument turns small generating \emph{maps} into lifting factorizations. These are different passages: existence of a small generating family under filtered colimits does not by itself assert that the same maps generate a weak factorization system. The additional closure and lifting arguments remain necessary.
\end{dcomment}
\begin{otherlanguage}{english}
\dgroverview{XV-a}{Chapter XV \textperiodcentered\ \textup{XV-a} \textperiodcentered\ Overview}
\selectlanguage{english}
Grothendieck devotes these pages to the passage from a small amount of generating data to a factorization theory on a large category. The point is not merely to factor an individual map. The construction must vary with that map, remain compatible with enlargement of the generators, and admit a bound which is independent of the particular object being treated. Only then can it contribute to the construction of categories of models and, through them, to the homotopy theories of diagrams which the manuscript seeks to organize.

The chapter brings together two kinds of control. Lifting properties say what a factorization must accomplish; cardinal filtrations say when the construction has accomplished enough and how large its intermediate objects can become. The former belongs to the language of model categories, while the latter continues the systematic treatment of generation and accessibility in SGA~4, Expos\'e~I. Their interaction is the mathematical subject of this part of Chapter~XV. The later discussion of accessible classes is therefore not a digression. It addresses the problem of replacing a class of prospective generators by a set to which the preceding construction can be applied.

\medskip\noindent\textbf{Attaching solutions and taking a sufficiently long limit.}
For a set $C_0$ of arrows and a map $f:X\to Y$, the first stage attaches, simultaneously, one solution for every square from an element of $C_0$ to $f$. Grothendieck realizes this simultaneous attachment as the filtered colimit of the finite amalgamations. This permits him to work with the relevant pushouts and filtered colimits, without assuming arbitrary coproducts at the outset. The residual map to $Y$ is then subjected to the same operation, and limit stages are formed by filtered colimits. The result is the classical small object construction, expressed as a transfinite sequence of functorial factorizations; compare \citet[Chapter~12]{Riehl2014}.

The decisive step is a factorization of a map from the domain of a generator through an earlier stage. If these domains are $\pi$-accessible in the manuscript's convention, a limit ordinal of cofinality greater than $\pi$ suffices. A square reaching the limit then already occurs at some stage and receives its solution at the next. The word \emph{limit} is necessary: an ordinal with a last predecessor can satisfy the order-theoretic definition of ``grand devant $\pi$'' while providing no further stage in which to attach the required solution. The printed statements omit this qualification. Their usual successor-cardinal choice of stopping ordinal has the required property and repairs them.

The cellular left class is compared with the left lifting class by the familiar retract argument. Factoring a map which already has the left lifting property exhibits it as a retract of its cellular factor. The commentary also supplies the functoriality arguments left to the reader, including compatibility with inclusions of generating sets. One should not confuse completion of the lifting task with stabilization of the sequence: for the sole generator $\varnothing\to1$ in sets, the stages have the form $X\amalg(Y\times\alpha)$. They continue growing after the residual map has become surjective. Garner's algebraic refinement treats the additional structure of chosen fillers and its composition laws \citep{Garner2009}; that structure is stronger than the functoriality established here.

\medskip\noindent\textbf{The arithmetic behind a uniform factorization.}
A cardinal filtration presents a large object by sufficiently directed diagrams of smaller ones. The two base cardinals have distinct roles: $\pi_0$ bounds the complexity of these presentations, while $\pi_1$ provides a uniform accessibility bound for the first bounded level. The estimates then depend on a small category of representatives and on its actual arrows. Counting only isomorphism classes in the arrow category loses the automorphisms and cannot bound individual Hom sets. The one-object category of an arbitrarily large small group gives a particularly short counterexample: its arrows have one isomorphism class, although its unique Hom set has the cardinality of the group.

Choose a small full skeleton $K_1$ of the $\pi_1$-level and put
\[
 \chi=\max\bigl(\aleph_0,\pi_1,|\operatorname{Mor}(K_1)|\bigr),
 \qquad \mu=r^{\pi_0},\qquad r,s\geq\chi.
\]
The corrected counting argument gives the intended forms of the bounds: at most $2^\mu$ isomorphism classes of objects and arrows at level $r$, and at most $s^\mu$ maps from an $r$-bounded object to an $s$-bounded one. The number of choices of a diagram must be counted by exponentiation, not by the printed products. With that change the final estimates still follow. The asserted universal sharpness of the two-coloured-set example does not: under GCH its arrow-orbit count at $r=\aleph_1$ is smaller than $2^r$.

The fixed-point device is robust. Taking a sufficiently large $\rho=2^c$ makes $\rho^{\pi_0}=\rho$ and absorbs the repeated attachment bounds. The manuscript's alternative GCH shortcut needs a cofinality condition: GCH does not imply $\rho^{\pi_0}=\rho$ at singular cardinals of cofinality at most $\pi_0$. The commentary separates these two choices and proves a bound retaining the actual number of generators. If their endpoints lie at level $r$, the input endpoints at level $s$, and the construction is followed through $\alpha$, then
\[
 B=\bigl(\max(2,r,s,|C_0|,|\alpha|)\bigr)^{r^{\pi_0}}
\]
bounds every stage, under the stated closure hypotheses. Thus the principal uniform-factorization conclusion survives the correction of its numerical threshold. Low's later analysis identifies regular-cardinal conditions under which factorization functors preserve presentable objects and a small subcategory detects a combinatorial model structure \citep[Theorem~4.4, Theorem~5.9, Proposition~5.12]{Low2016Heart}.

\medskip\noindent\textbf{Recognizing a category of models.}
Section~4 separates the elementary logic of lifting from the smallness needed to produce the factorizations. A precouple supplies lifting and factorization; closing its two classes under retracts gives the corresponding weak factorization system. Two compatible systems, together with the cancellation properties of $W$, provide the model-category identities. This distinction explains both the strength of Proposition~4.2 and the defect in the later footnote which treats unclosed precouples as though they were already closed. A set-theoretic counterexample in the commentary satisfies the displayed equalities but fails retract closure.

The class $\Cof_W$ records preservation of weak equivalences by pushout. An arrow in $W\cap\Cof_W$ need not, under the hypotheses printed at that point, remain in $W$ after every cobase change. A four-object diamond lattice makes the difference explicit. The omitted sufficient hypothesis is already present in Chapter~XIII: every map must factor as a $W$-cofibration followed by a weak equivalence. The commentary proves the resulting identification with couniversal weak equivalences. The eventual recognition theorem, Proposition~4.4.6, supplies this factorization through its first weak factorization system, so its principal conclusion is preserved.

For the existence of the second small generating family, modern accessibility hypotheses give a precise answer. Smith's theorem as published by \citet[Theorem~1.7 and Proposition~1.15]{Beke2000} supplies the relevant solution-set argument in a locally presentable category. The formulation in \citet[Proposition~A.2.6.15, updated online version]{Lurie2009HTT} is especially close to the proposed construction from $W$ and $\Cof_W$. These results resolve the problem in their stated setting; they do not justify deleting hypotheses from an arbitrary cardinally filtered category. Once such a combinatorial model category is obtained, its diagram homotopy categories form a strong derivator \citep[Proposition~1.30]{Groth2013}. This identifies the later destination of the machinery without treating this portion of the manuscript as a proof of every derivator axiom.

\medskip\noindent\textbf{Recovering a class from its bounded part.}
Section~5 replaces presentations by maps with presentations by small subobjects. Images, closure of bounded levels under suitable colimits, and exactness of filtered colimits make the poset of small subobjects into a continuous, sufficiently directed presentation of the object. The hypotheses on images are independent in the direction questioned by the source: a cyclic module may contain a submodule needing arbitrarily many generators. The commentary gives the complete polynomial-ring example, and distinguishes a strict monomorphism from an arbitrary proper monomorphism in the preceding image lemmas.

The definition of an accessible class used here tests \emph{every} perfectly adapted presentation. Recognizing membership from cofinally many good small subobjects is therefore initially only a necessary condition. The first version of Proposition~5.6 misses the information carried by nonmonic presentations; Grothendieck himself corrects it in footnote~21. A continuous diagram of two-point and one-point posets makes the retract obstruction explicit. Closure of the bounded part under small filtered colimits repairs the assertion. The remaining comparison is a closed-cofinal-set argument: two cofinal families meet cofinally when both contain the relevant directed suprema. The supplied proofs include the eventual-equality step in Lemma~5.7 and avoid assuming that the image of a continuous cofinal indexing map is itself closed.

At cardinals where the bounded objects are themselves $\pi$-accessible, these arguments give the four descriptions in the corrected Proposition~5.7, the recovery bijection of Corollary~5.8, and filtered-colimit closure. Proposition~5.9 is printed in wider generality. Its proof invokes a slice argument whose supplied justification uses hypotheses no longer retained. We prove the proposition at the cardinals required for the preceding recognition theorem, but do not certify the unrestricted fixed-rank assertion. This is a stated limit of the present verification, not a designation of an open problem. In the verified range, the resulting full subcategory is accessible and accessibly embedded; it need not be cocomplete, as the nonempty sets show \citep[Definition~2.1]{AdamekRosicky1994}.

The commentary keeps these distinctions visible in its apparatus. The edition remains the authority for the historical text; genuine transcription corrections are applied, while slips visible in the edition are retained and explained. The claim register records the source's proof status separately from the repair. Results developed here carry the marker \ours{} in their commentary blocks and no claim of priority.
\end{otherlanguage}
\dgrchunkend

\addtocontents{toc}{\protect\dgrtocmissing{XV-b}{Part 2 of 2}}
\addtocontents{toc}{\protect\dgrtocmissingchapter{XVI}{Les localiseurs fondamentaux}}
\addtocontents{toc}{\protect\dgrtocmissingchapter{XVII}{Catégories à fibrations et à cofibrations (Modèles (3))}}
\addtocontents{toc}{\protect\dgrtocmissingchapter{XVIII}{Catégories et ensembles accessibles}}
% DER_XIX: Grothendieck's text with the commentary set in, then the overview, inside a group that applies the chunk's own definitions
\dgrchapter{XIX}{Chapter XIX}{Modèles (4)}
\dgrpart{0}{XIX}{Manuscript pages 1--50}{XIX\enspace\textperiodcentered\enspace Chapter XIX\enspace\textperiodcentered\enspace manuscript pages 1--50}{Chapter XIX \textperiodcentered\ \textup{XIX}}
\dgrchunkbegin
% DER_XIX: the definitions of the chunk's preamble, applied inside its group
\let\printedpage\relax\newcommand{\printedpage}[1]{\marginpar{\footnotesize p.~#1}}
\let\msp\relax\newcommand{\msp}[1]{\marginpar{\footnotesize [#1]}}

% Q-orthogonality arrows (half arrowheads = Q-orthogonal, full arrowheads = equality)
\let\Qo\relax\newcommand{\Qo}{\mathrel{\leftharpoondown\joinrel\relbar\joinrel\rightharpoonup}}   % lower half-head left, upper half-head right
\let\Qorev\relax\newcommand{\Qorev}{\mathrel{\leftharpoonup\joinrel\relbar\joinrel\rightharpoondown}} % mirrored form of (1.2)
\let\Qog\relax\newcommand{\Qog}{\mathrel{\leftarrow\joinrel\relbar\joinrel\rightharpoonup}}         % full head left
\let\Qod\relax\newcommand{\Qod}{\mathrel{\leftharpoondown\joinrel\relbar\joinrel\rightarrow}}       % full head right
\let\Qogd\relax\newcommand{\Qogd}{\mathrel{\leftarrow\joinrel\relbar\joinrel\rightarrow}}            % full heads both sides
% the same arrows crossed by a vertical bar (Quillen factorisation)
\let\Qx\relax\newcommand{\Qx}[2]{\mathrel{#1\joinrel\mapstochar#2}}
% crossed by two vertical bars (functorial factorisation)
\let\Qxx\relax\newcommand{\Qxx}[2]{\mathrel{#1\joinrel\mapstochar\mkern2mu\mapstochar#2}}
% Q-closure marks
\let\tld\relax\newcommand{\tld}{\tilde{}}   % left Q-closure mark (small tilde in exponent)
\let\rcl\relax\newcommand{\rcl}{\bar{}}       % right Q-closure mark (horizontal bar in exponent)
\let\Qorth\relax\newcommand{\Qorth}{\mathrel{\leftharpoondown\mkern-6mu\relbar\mkern-6mu\rightharpoonup}}
\let\Qorthx\relax\newcommand{\Qorthx}{\mathrel{\leftharpoondown\mkern-9mu{+}\mkern-9mu\rightharpoonup}}
\let\Qorthxfull\relax\newcommand{\Qorthxfull}{\mathrel{\leftharpoondown\mkern-9mu{+}\mkern-9mu\rightarrow}}
\let\Qcouple\relax\newcommand{\Qcouple}{\mathrel{\leftarrow\mkern-9mu{+}\mkern-9mu\rightarrow}}
\let\clos\relax\newcommand{\clos}{\longleftrightarrow}
\let\fl\relax\newcommand{\fl}{\longrightarrow}
\let\M\relax\newcommand{\M}{\mathcal{M}}
\let\Fl\relax\newcommand{\Fl}{\mathrm{Fl}}
\let\Ho\relax\newcommand{\Ho}{\mathrm{Ho}}
\let\cc\relax\newcommand{\cc}{\mathrm{c}}
\let\ddarr\relax\newcommand{\ddarr}{\mathrel{\twoheadleftarrow\joinrel\twoheadrightarrow}}
\let\harp\relax\newcommand{\harp}{\mathrel{\leftharpoondown\joinrel\relbar\joinrel\rightharpoonup}}
\let\deq\relax\newcommand{\deq}{\overset{\text{déf}}{=}}
\let\Mono\relax\newcommand{\Mono}{\mathrm{Mono}}
\let\Filt\relax\newcommand{\Filt}{\mathrm{Filt}}
\let\Cof\relax\newcommand{\Cof}{\mathrm{Cof}}
\let\qpair\relax\newcommand{\qpair}{\mathrel{\leftharpoondown\joinrel\relbar\joinrel\rightharpoonup}}
\let\Cat\relax\newcommand{\Cat}{\mathrm{Cat}}
\let\Diag\relax\newcommand{\Diag}{\mathrm{Diag}}
\let\id\relax\newcommand{\id}{\mathrm{id}}
\let\Qs\relax\newcommand{\Qs}{\tilde{\ }}
\let\fact\relax\newcommand{\fact}{\mathrel{\leftarrow\mkern-11mu\twoheadrightarrow}}
\let\limind\relax\newcommand{\limind}{\varinjlim}
\let\limproj\relax\newcommand{\limproj}{\varprojlim}
\let\Is\relax\newcommand{\Is}{\mathrm{Is}}

% ---- from der-ch19_p001-005.tex
\printedpage{(non numérotée)}
% head: (none; title page)

\begin{center}
{\Large\textsf{LES DÉRIVATEURS}}

\bigskip

\textbf{ALEXANDRE GROTHENDIECK}

\bigskip\bigskip

\textbf{Chapitre XIX}

\medskip

\textbf{Modèles (4)}
\end{center}

\newpage

\printedpage{(non numérotée)}
% head: (none; editorial note)

Ce texte a été déchiffré et transcrit en \LaTeXe\ par M. Künzer. Il a été édité par M. Künzer, J. Malgoire, et G. Maltsiniotis. La transcription est aussi fidèle que possible au manuscrit. Seules quelques corrections mineures évidentes ont été effectuées. Pour les rares ajouts ou commentaires des éditeurs, ainsi que pour la numérotation originale des pages du manuscrit, les caractères de machine à écrire \texttt{[typewriter]} entre crochets sont utilisés.

Cette édition est provisoire. Les remarques, commentaires et corrections sont bienvenus. Envoyer un message à :

\textbf{georges.maltsiniotis@imj-prg.fr}

\hfill G. Maltsiniotis\hspace*{2cm}

\newpage

\printedpage{1}
% head: A. Grothendieck : Dérivateurs, Ch. XIX, Édition des Universités de Montpellier II et Paris VII
% foot: Version 11 février 2019

\msp{page 1}

\section*{1\quad Rappels sur les couples de Quillen (couples de Quillen spéciaux)}

Je vais reprendre les notations et la terminologie pour les relations de LLP \texttt{[left lifting property]} et RLP \texttt{[right lifting property]} de Quillen. J'écris (pour une catégorie $\mathcal{M}$, et $\Phi, \Psi \subset \mathrm{Fl}(\mathcal{M})$)
\[
\Phi \Qo \Psi \qquad \iff \qquad (\Phi \subset \Psi^{*}) \qquad \iff \qquad (\Psi \subset \Phi_{*})\ ,
\tag*{(1.1)}
\]
i.e. $\Phi$ a la LLP par rapport à $\Psi$, ou encore $\Psi$ la RLP par rapport à $\Phi$. Je dis aussi que $(\Phi, \Psi)$ est un \textit{couple} Q-\textit{orthogonal.} Cette notion n'est pas symétrique, mais je me permets parfois (exceptionnellement) d'écrire
\[
\Psi \Qorev \Phi \qquad \text{comme synonyme de} \qquad \Phi \Qo \Psi\ .
\tag*{(1.2)}
\]
$\Phi$ est Q-\textit{orthogonal à gauche} de $\Psi$, $\Psi$ est Q-\textit{orthogonal à droite} à $\Phi$, et la Q-orthogonalité à gauche (resp. à droite) est notée par le demi-sommet de flèche pointant vers le bas (resp. vers le haut). J'écris
\[
\left\{
\begin{aligned}
&\Phi \quad \Qog \quad \Psi \quad \iff \quad (\Phi = \Psi^{*})\\
&\Phi \quad \Qod \quad \Psi \quad \iff \quad (\Psi = \Phi_{*})\\
&\Phi \quad \Qogd \quad \Psi \quad \iff \quad (\Phi = \Psi^{*} \ \text{ et } \ \Psi = \Phi_{*})\ .
\end{aligned}
\right.
\tag*{(1.3)}
\]
Ces trois cas de Q-orthogonalité renforcée sont donc ceux où on a de plus $\Phi = \Phi^{\tld}$
resp. $\Psi = \Psi^{\rcl}$ resp. $\Phi = \Phi^{\tld}$ et $\Psi = \Psi^{\rcl}$. On dit que $\Phi$ est Q-\textit{clos à gauche} si $\Phi = \Phi^{\tld}$, que % reading resolved: DER_XIX_corrections.csv
$\Psi$ est Q-\textit{clos à droite} si $\Psi = \Psi^{\rcl}$, que le couple Q-orthogonal $(\Phi, \Psi)$ est Q-\textit{clos à gauche} (resp. \textit{à droite}) si $\Phi$ (resp. $\Psi$) est Q-clos à gauche (resp. à droite). On dit qu'il est Q-\textit{clos} si les deux conditions % reading resolved: DER_XIX_corrections.csv

\msp{page 2}

sont satisfaites simultanément. Quand aucune confusion n'est à craindre, il m'arrivera de dire simplement que $\Phi$ est Q-clos, ou que $\Psi$ est Q-clos, sans spécifier si c'est à gauche où % [sic] où -> ou, conjunction "or" intended
à droite. Mais pour un couple Q-orthogonal $(\Phi, \Psi)$ (ou je préfère l'écrire $\Phi \Qo \Psi$), Q-clos % [sic] ou -> où, relative adverb intended
signifie toujours Q-clos à gauche et à droite.

\dgrblock{dgrb1158}{Orthogonality, closure, and factorization}
\begin{dcomment}{exposition}{p2:signifie toujours Q-clos à gauche et à droite.}{Orthogonality, closure, and factorization}
Grothendieck begins by separating the ingredients that will later have to survive passage to a category of diagrams. A lifting property is easy to transport along an adjunction. A factorization requires a construction. The purpose of the elaborate arrows is to prevent the existence of that construction from being smuggled into the lifting notation.

For classes of arrows $L,R$ in a locally small category, write $L\perp R$ when every commutative square with left side in $L$ and right side in $R$ has a diagonal filler. The filler need not be unique. In modern notation,
\[
 {}^{\square}R=\{i:i\perp R\},\qquad
 L^{\square}=\{p:L\perp p\}.
\]
These are the source's $R^*$ and $L_*$. They form an antitone Galois connection: $L\subseteq{}^{\square}R$ if and only if $R\subseteq L^{\square}$. Consequently
\[
 \overline L={}^{\square}(L^{\square}),\qquad
 \overline R=({}^{\square}R)^{\square}
\]
are extensive, monotone, idempotent operations. Grothendieck writes a small tilde for the first closure and an overbar for the second. An arrow with two complete heads means the stronger pair of identities $L={}^{\square}R$ and $R=L^{\square}$. A transverse stroke adds the assertion that every $f$ factors as $pi$ with $i\in L$ and $p\in R$. Together these are precisely the data of a weak factorization system, without a chosen functorial factorization.

There is a genuine logical slip in the sentence following (1.3). For an arbitrary orthogonal pair, separate closedness of $L$ and $R$ does not imply that they are one another's orthogonals. Indeed, in $\mathrm{Set}$ take both classes to be the isomorphisms. Each is closed under its appropriate double-orthogonal operation, because
\[
 {}^{\square}\mathrm{Fl}=\mathrm{Is},\quad
 \mathrm{Is}^{\square}=\mathrm{Fl},\quad
 {}^{\square}\mathrm{Is}=\mathrm{Fl},\quad
 \mathrm{Fl}^{\square}=\mathrm{Is}.
\]
For example, an arrow having the left lifting property against itself has a two-sided inverse, by using the square with identity horizontal sides; this proves the first identity. The other three follow in the same way or directly from invertibility. Thus $(\mathrm{Is},\mathrm{Is})$ is separately closed and orthogonal, but is not a mutually maximal orthogonal pair unless the category is a groupoid.

The repair is to retain the explicit equalities of (1.3) as the meaning of a closed \emph{pair}. There is no difficulty with the later terminology for a Quillen pair: once factorization is available, Proposition 1.3 shows that separate retract-closedness forces these equalities. The distinction matters precisely before factorization, where this chapter's questions lie. We therefore retain the printed sentence, record the defect in the audit, and use the mutually orthogonal meaning whenever a two-headed arrow occurs below.
\end{dcomment}


Quand un couple Q-orthogonal $\Phi \Qo \Psi$ a la propriété de factorisation de Quillen, je l'indique par la barre verticale croisant l'horizontale, ainsi
\[
\Phi \ \Qx{\leftharpoondown}{\rightharpoonup} \ \Psi\ ,
\tag*{(1.4)}
\]
et on définit de même le sens des graphismes
\[
\left\{
\begin{aligned}
&\Phi \quad \Qx{\leftharpoondown}{\rightharpoonup} \quad \Psi\\
&\Phi \quad \Qx{\leftarrow}{\rightharpoonup} \quad \Psi\\
&\Phi \quad \Qx{\leftharpoondown}{\rightarrow} \quad \Psi\\
&\Phi \quad \Qx{\leftarrow}{\rightarrow} \quad \Psi\ .
\end{aligned}
\right.
\tag*{(1.3)} % [sic] (1.3) -> (1.5), number already used above and (1.5) is missing
\]

\newpage

\printedpage{2}
% head: A. Grothendieck : Dérivateurs, Ch. XIX, Édition des Universités de Montpellier II et Paris VII
% foot: Version 11 février 2019

Un couple Q-\textit{orthogonal} Q-\textit{factorisant} $\Phi \Qx{\leftharpoondown}{\rightharpoonup} \Psi$ est appelé un \textit{précouple de Quillen.} S'il est de plus Q-clos, i.e. si on a $\Phi \Qx{\leftarrow}{\rightarrow} \Psi$, on dit que c'est un \textit{couple de Quillen.}

$\Phi$ est dit Q-\textit{saturé à gauche} s'il contient les isomorphismes et est stable par cochangement de base, par (petite) composition transfinie, et par facteurs directs. Dualement $\Psi$ est dit Q-\textit{saturé à droite} s'il

\msp{page 3}

contient les isomorphismes, est stable par changement de base, par composition transfinie, et par facteurs directs. Si $\Phi$ est Q-saturé à gauche, il est stable par (petites) sommes directes, i.e. si les $f_i : A_i \longrightarrow B_i$ ($i \in I$, $I$ petit) sont dans $\Phi$, $f = \coprod f_i : \coprod A_i \longrightarrow \coprod B_i$ aussi, et de même pour (petites) sommes amalgamées finies ou infinies (quand on se place dans un $A\backslash\mathcal{M}$ au lieu de $\mathcal{M}$). On a l'énoncé dual pour $\Psi$, quand $\Psi$ est Q-saturé à droite.

En fait, il y a des notions plus faibles de Q-saturation que celle que je viens de dire : Q-\textit{saturation simple} (à gauche : $\Phi$ contient les isomorphismes, stable par composition binaire, par cochangement de base), Q-\textit{saturation transfinie} (on exige de plus stabilité par composition transfinie), enfin Q-\textit{saturation totale,} ou Q-saturation tout court (quand on exige de plus la stabilité par facteurs directs). On a les implications
\[
\left\{
\begin{aligned}
&\text{Q-clos à gauche}\\
&\Downarrow\\
&\text{Q-saturé (total)}\\
&\Downarrow\\
&\text{Q-saturé transfini}\\
&\Downarrow\\
&\text{Q-saturé simple.}
\end{aligned}
\right.
\tag*{(1.6)}
\]
Pour $\Phi \subset \mathrm{Fl}(\mathcal{M})$, on a sa Q-clôture (à gauche) $\Phi^{\tld}$ (qui est la plus petite partie Q-close à gauche de $\mathrm{Fl}(\mathcal{M})$ contenant $\Phi$) et de même son Q-saturé (total) à gauche, son Q-saturé transfini, son Q-saturé simple

\msp{page 4}

avec les inclusions tautologiques
\[
\Phi \ \subset \
\underbrace{\Phi'}_{\substack{\text{Q-saturé}\\ \text{à gauche simple}}}
\ \subset \
\underbrace{\Phi''}_{\substack{\text{Q-saturé}\\ \text{à gauche transfini}}}
\ \subset \
\underbrace{\Phi'''}_{\substack{\text{Q-saturé}\\ \text{à gauche total}}}
\ \subset \
\underbrace{\Phi^{\tld}}_{\substack{\text{Q-clôture}\\ \text{à gauche}}}
\ ,
\tag*{(1.7)}
\]
et dualement pour $\Psi$. Entre $\Phi$ et $\Phi'$ il y a encore un $\Phi'_0$,
\[
\Phi \ \subset \ \Phi'_0 \ \subset \ \Phi'\ ,
\tag*{(1.8)}
\]
qui est le Q-\textit{saturé pour cochangement de base} (contenant aussi, par définition, tous les isomorphismes). Pour construire tous ces saturés, on procède ainsi :

\dgrblock{dgrb1159}{What cellular closure does and does not supply}
\begin{dcomment}{exposition}{p4:Pour construire tous ces saturés, on procède ainsi}{What cellular closure does and does not supply}
The four closures organize a constructive question: how much of a left lifting class can be built by attaching specified arrows? The elementary operations are chosen because every left lifting class survives them. For a pushout of $i$, a lifting square is equivalent, by the pushout property, to a lifting square against $i$. For a composite $i_2i_1$, solve the lifting problem first against $i_1$ and then against $i_2$. For a continuous ordinal sequence $(X_\alpha)_{\alpha\leq\lambda}$, compatible fillers are constructed at successor stages and are assembled at a limit stage by the colimit property. The resulting map $X_\lambda\to E$ is the required filler for the transfinite composite. Finally, a lifting problem against a retract of an arrow is a retract of a lifting problem against that arrow, so a filler restricts along the retraction.

These arguments prove the implications in (1.6), including the dual assertions, whenever the indicated colimits or limits exist. They also explain why small coproducts need not be added as an independent closure operation. Well-order the indexing set $S$ of a family $(A_s\to B_s)$. Starting with $\coprod_sA_s$, replace the summand $A_s$ by $B_s$ at stage $s$, using a pushout. At a limit stage take the colimit. The final map is $\coprod_sA_s\to\coprod_sB_s$. The same construction in an undercategory treats amalgamated coproducts. This use of a well-order belongs to the usual set-theoretic choice convention.

The formula for $\Phi''$ on the next page rests on a useful flattening argument. A pushout of a transfinite composite is the transfinite composite of the successive pushouts. A transfinite composite of transfinite composites can be indexed by the ordinal sum of the component ordinals. Thus the transfinite composites of pushouts of the original generators already form a class closed under pushout and transfinite composition. Identity maps allow stages at which nothing is attached. This proves the asserted description of $\Phi''$ without first adding an unspecified hierarchy of repeated cellular operations.

What remains at this point is the relation with retracts. Closure under retracts alone gives no factorization, and the notation $\Phi'''$ initially means the smallest class stable under \emph{all} the listed operations. Grothendieck will identify it with the retracts of $\Phi''$ under small-generation hypotheses, and simultaneously identify both with the entire lifting class. The first of these equalities in fact has a general proof, given in the next block; the second is what the small object argument adds. This distinction prevents a positive answer to the elementary closure question from being mistaken for a factorization theorem. A modern reference for the latter is \citet[Theorem~2.1.14 and Corollary~2.1.15]{Hovey1999}.
\end{dcomment}


\newpage

\printedpage{3}
% head: A. Grothendieck : Dérivateurs, Ch. XIX, Édition des Universités de Montpellier II et Paris VII
% foot: Version 11 février 2019

\begin{enumerate}[leftmargin=*]
\item[1)] $\Phi'_0 \quad = \quad$ ensemble des flèches $i : A \longrightarrow B$ s'insérant dans un carré cocartésien
\[
\begin{tikzcd}
A_0 \arrow[r, "i_0"] \arrow[d] & B_0 \arrow[d]\\
A \arrow[r, "i"] & B\ ,
\end{tikzcd}
\]
avec $i_0 \in \Phi$ ou $i_0$ une identité (auquel cas on trouve pour $i$ les isomorphismes).

\item[2)] $\Phi' \quad = \quad$ ensemble des flèches composées multiples
\[
i_n \, i_{n-1} \ \cdots \ i_1\ ,
\]
avec $i_\alpha \in \Phi'_0$.

\item[3)] $\Phi'' \quad = \quad$ ensemble des composés transfinis de flèches $\in \Phi'_0$.
\end{enumerate}

On a envie de poser
\[
\Phi''' \ = \ \Phi^{\prime\prime\natural} \ \overset{\text{déf}}{=} \ \text{ensemble des facteurs directs d'éléments de } \Phi'',
\]
et c'est un fait que $\Phi^{\prime\prime\natural} \subset \Phi'''$ trivialement, mais il n'est par clair en général qu'on ait % [sic] par -> pas, typo
\msp{page 5}
égalité, faute de pouvoir montrer que $\Phi^{\prime\prime\natural}$ est stable, non seulement par facteurs directs, mais aussi par cochangement de base et par composition transfinie. Cependant, rappelons :

\dgrblock{dgrb1160}{Retracts of cellular arrows are closed under transfinite composition}
\begin{dcomment}{later}{p5:égalité, faute de pouvoir montrer}{Retracts of cellular arrows are closed under transfinite composition}
The uncertainty expressed here has a positive answer without a small generating set. If $\mathcal M$ is cocomplete and $\mathcal C$ contains the isomorphisms and is stable under pushout and transfinite composition, then $\operatorname{Retr}(\mathcal C)$ has the same properties. This is \citet[Lemma~2.2]{MakkaiRosicky2014}. We give the explicit compatible-retraction construction, since it isolates exactly the issue Grothendieck raises.

First every retract of a map in $\mathcal C$ can be presented as a retract with its domain fixed. Suppose $f:A\to B$ is a retract of $g:X\to Y$, with maps $a:A\to X$, $b:X\to A$, $c:B\to Y$, $d:Y\to B$ satisfying
\[
ba=1_A,\quad dc=1_B,\quad ga=cf,\quad fb=dg.
\]
Push $g$ out along $b$, obtaining $k:A\to P=A\amalg_XY$ in $\mathcal C$. If $q:Y\to P$ is the other pushout map, put $s=qc:B\to P$. The maps $f:A\to B$ and $d:Y\to B$ induce $r:P\to B$. Then $sf=k$, $rk=f$, and $rs=1_B$. Thus $f$ is a retract of $k$ in $A\backslash\mathcal M$.

Now let $(X_\alpha)_{\alpha\leq\lambda}$ be a continuous chain with each $f_\alpha:X_\alpha\to X_{\alpha+1}$ in $\operatorname{Retr}(\mathcal C)$. Choose fixed-domain presentations
\[
 X_\alpha\xrightarrow{h_\alpha}Z_\alpha,
 \quad s_\alpha':X_{\alpha+1}\to Z_\alpha,
 \quad r_\alpha':Z_\alpha\to X_{\alpha+1},
\]
with $h_\alpha\in\mathcal C$, $s_\alpha'f_\alpha=h_\alpha$, $r_\alpha'h_\alpha=f_\alpha$, and $r_\alpha's_\alpha'=1$. Construct a second continuous chain $Y_\alpha$ and compatible retractions $X_\alpha\xrightarrow{s_\alpha}Y_\alpha\xrightarrow{r_\alpha}X_\alpha$, starting with $Y_0=X_0$. At a successor stage set
\[
Y_{\alpha+1}=Y_\alpha\amalg_{X_\alpha}Z_\alpha,
\]
using $s_\alpha$ and $h_\alpha$ for the pushout. Its structure map $g_\alpha:Y_\alpha\to Y_{\alpha+1}$ lies in $\mathcal C$. Define $s_{\alpha+1}$ by the inclusion of $Z_\alpha$ composed with $s_\alpha'$. Define $r_{\alpha+1}$ by $f_\alpha r_\alpha$ on $Y_\alpha$ and $r_\alpha'$ on $Z_\alpha$; these agree on $X_\alpha$. The equations
\[
r_{\alpha+1}s_{\alpha+1}=1,
\quad g_\alpha s_\alpha=s_{\alpha+1}f_\alpha,
\quad r_{\alpha+1}g_\alpha=f_\alpha r_\alpha
\]
follow directly. At a limit stage take colimits of both the chains and these equations. The composite $X_0\to X_\lambda$ is consequently a retract, with domain fixed, of $Y_0\to Y_\lambda\in\mathcal C$.

Pushout stability follows by pushing out a fixed-domain retraction. Stability under retracts follows by composing retraction diagrams. Taking $\mathcal C=\Phi''$ therefore proves $\Phi'''=\operatorname{Retr}(\Phi'')$ in every cocomplete category, even for a proper class of generators. It does not prove $\Phi'''={}^{\square}(\Phi^{\square})$: that last equality still needs a mechanism, such as Theorem 1.1, for constructing the appropriate factorizations.
\end{dcomment}


\textbf{Théorème 1.1.} \textit{Supposons $\mathcal{M}$ stable par petites $\varinjlim$ , et que ses éléments soient accessibles. Supposons de plus que $\Phi$ soit petit, plus généralement, qu'il existe une partie $\Phi_0 \subset \Phi$ petite telle que $\Phi \subset \Phi_0{}^{\tld}$, donc}
\[
\Psi \ \overset{\text{déf}}{=} \ \Phi_{*} \ = \ (\Phi_0)_{*} \qquad \text{\footnotemark}.
\tag*{(1.1.1)} % footnote mark printed as superscript (1)
\]
\footnotetext{\textbf{N.B.} On a donc $\Phi \Qod \Psi$, a fortiori $\Phi_0 \Qod \Psi$.}
\textit{Sous ces conditions, si $\Phi''_0$ désigne le \textup{Q}-saturé à gauche transfini de $\Phi_0$, alors on a $\Phi''_0 \Qx{\leftharpoondown}{\rightarrow} \Psi$, i.e. toute flèche $f$ de $\mathcal{M}$ se factorise en $f = pi$, avec $i \in \Phi''_0$, $p \in \Psi$\,\footnote{Il y a même factorisation \textit{fonctorielle} en $f$, propriété qui mérite une notation, p.ex. $\Phi''_0 \Qxx{\leftharpoondown}{\rightarrow} \Psi$.}. A fortiori, si $\Phi''$ désigne le \textup{Q}-saturé à gauche transfini de $\Phi$, on a} % footnote mark printed as superscript (2)
\[
\Phi'' \ \Qx{\leftharpoondown}{\rightarrow} \ \Psi\ .
\]

\dgrblock{dgrb1161}{The small object argument and the retract that finishes it}
\begin{dcomment}{exposition}{p5:A fortiori, si $\Phi''$ désigne}{The small object argument and the retract that finishes it}
Theorem 1.1 is the engine of the chapter. It recalls the transfinite construction developed in \citet[Chapter~XV, \S1, Theorem~1]{Grothendieck1990XV}. Its force is not merely that sufficiently many cells exist, but that a \emph{set} of generators admits one common bound on the stages needed to see all its lifting problems. Read ``accessible object'' here as an object whose maps into sufficiently filtered colimits are detected at a stage; for this argument the weaker condition of smallness relative to the relevant cellular sequences suffices. Local smallness is understood in the chosen universe. These hypotheses do not assert that the entire category is locally presentable.

Put $S=\Phi_0$ and let $f:X\to Y$. Choose an infinite regular cardinal $\kappa$ such that every domain of an arrow of $S$ is $\kappa$-small for the sequences being used. Such a common bound exists because $S$ is a set. Starting from $X_0=X$, form at stage $\alpha+1$ the pushout
\[
\begin{tikzcd}[column sep=large]
\displaystyle\coprod_{\sigma} A_\sigma\arrow[r]\arrow[d] & X_\alpha\arrow[d]\\
\displaystyle\coprod_{\sigma} B_\sigma\arrow[r] & X_{\alpha+1},
\end{tikzcd}
\]
where $\sigma$ runs over all commutative squares $s:A_\sigma\to B_\sigma$ against $X_\alpha\to Y$, with $s\in S$. There is a set of such squares, by local smallness and the smallness of $S$. Their lower maps induce $X_{\alpha+1}\to Y$. At a limit stage take the colimit. Choose the length $\lambda$ so that the sequence is sufficiently filtered for every domain in $S$, for example a regular cardinal larger than a common smallness bound.

The composite $i:X\to X_\lambda$ lies in $S''$: attaching a coproduct of cells can be serialized as on the preceding page. The induced $p:X_\lambda\to Y$ has the right lifting property against $S$. Indeed, in a square from $s:A\to B$ to $p$, the map $A\to X_\lambda$ factors through some $X_\alpha$. Together with the fixed map $B\to Y$, it is one of the squares attached at the next stage. The corresponding copy of $B$ in $X_{\alpha+1}$ gives the filler after composition with $X_{\alpha+1}\to X_\lambda$. The equation with $Y$ already holds at stage $\alpha$, since all maps are over $Y$. This proves $f=pi$. Attaching \emph{all} squares makes the construction functorial once colimit choices are fixed; a morphism of arrows sends each indexed square to an indexed square at the next stage.

The retract argument is independent of this construction and deserves to be explicit. Suppose $f:A\to B$ has the left lifting property against $p:E\to B$ and $f=pi$ for $i:A\to E$. A filler in
\[
\begin{tikzcd}
A\arrow[r,"i"]\arrow[d,"f"'] & E\arrow[d,"p"]\\
B\arrow[r,"1_B"']\arrow[ur,dashed,"s"] & B
\end{tikzcd}
\]
satisfies $sf=i$ and $ps=1_B$. Hence $(1_A,s):f\to i$ and $(1_A,p):i\to f$ exhibit $f$ as a retract of $i$ in the arrow category. In particular, for any factorizing orthogonal pair $(L,R)$,
\[
{}^{\square}R=\operatorname{Retr}(L),\qquad
L^{\square}=\operatorname{Retr}(R).
\]
The reverse inclusions follow from stability of lifting under retracts. The double-orthogonal closures equal these classes because orthogonality is unchanged by taking retracts. This proves Proposition 1.3 and, applied to $(S'',S^{\square})$, Corollary 1.2. It also proves Corollary 1.4: factor an arrow of $L$ as a map in $W$ followed by a map in $R$, then use the same square. If $W$ is retract-closed and $(L,R)$ itself factorizes, the converse follows from $L\subseteq W$.
\end{dcomment}


\bigskip

\textbf{Corollaire 1.2.} \textit{Sous les conditions du théorème, on a}
\[
\Phi''' \ \text{(Q-saturé total à gauche de $\Phi$)} \ = \
\underbrace{\Phi^{\prime\prime\natural}}_{\substack{\text{ensemble des}\\ \text{facteurs directs d'éléments}\\ \text{de $\Phi''$}}}
\ = \ \Phi^{\tld}\ .
\tag*{(1.2.1)}
\]
\textit{Plus généralement (en apparence seulement!), si $\Phi_1$ est une partie de \textup{$\mathrm{Fl}(\mathcal{M})$} qui est \textup{Q}-saturée à gauche transfinie, et telle que}
\[
\Phi_0 \ \subset \ \Phi_1 \ \subset \ \Phi^{\tld} \ = \ \Phi_0{}^{\tld}\ ,
\]

% ---- from der-ch19_p006-010.tex
\printedpage{4}
% head: A. Grothendieck : Dérivateurs, Ch. XIX, Édition des Universités de Montpellier II et Paris VII
% foot: Version 11 février 2019

\msp{page 6}
\textit{alors}
\begin{enumerate}[label={}]
\item[a)] \textit{on a $\Phi_1 \Qorthxfull \Psi$ et} % reading resolved: DER_XIX_corrections.csv
\item[b)] $\Phi_1^{\sim} = \Phi^{\sim} = \Phi_1^{\natural}$.
\end{enumerate}
Ce corollaire résulte du théorème en vertu de XV prop. 4.2, dont je rappelle ici l'essentiel :

\textbf{Proposition 1.3.} \textit{Soit $\Phi \Qorthx \Psi$. Alors $\Phi^{\sim} = \Phi^{\natural} = \Psi^{*}$, et dualement on a $\Psi^{-} = \Psi^{\natural} = \Phi_{*}$.}

Notons aussi, dans cet ordre d'idées :

\textbf{Corollaire 1.4.} \textit{Soit $\Phi \Qorth \Psi$, et soit $W$ une partie quelconque de $\mathrm{Fl}(\mathcal{M})$. Pour qu'on ait $\Phi \subset W^{\natural}$, il suffit que toute flèche de $\mathcal{M}$ se factorise en $f = pi$, avec $i \in W$, $p \in \Psi$. (Et cette condition est évidemment aussi nécessaire si $\Phi \Qorthx \Psi$, et $W$ stable par facteurs directs.)}

On dit qu'une partie $\Phi_0$ de $\Phi$ est Q-\textit{génératrice à gauche} de $\Phi$ si elle satisfait la condition
\[
\Phi \ \subset \Phi_0^{\sim}\ , \tag*{(1.9)}
\]
et on dit que $\Phi$ est une partie \textit{spéciale} si $\Phi$ admet une \textit{petite} partie Q-génératrice à gauche. Si $\Phi_0 \subset \Phi$ est petit, et si $\mathcal{M}$ satisfait les conditions du théorème 1.1, alors $\Phi_0$ est une partie Q-génératrice (à gauche) si et seulement si son Q-saturé total (à gauche) $\Phi_0'''$ contient $\Phi$ (condition en apparence plus forte que $\Phi \subset \Phi_0^{\sim}$, puisque en général on a $\Phi_0''' \subset \Phi_0^{\sim}$). Cela signifie donc aussi, via 1.2, que tout \msp{page 7} élément de $\Phi$ est facteur direct d'un composé transfini de flèches déduites de
\[
\Phi_0 \cup (\text{Identités de } \mathcal{M})
\]
par cochangement de base.

Je rappelle encore le sorite $\pm$ tautologique des foncteurs adjoints :

\textbf{Proposition 1.5.} \textit{Soit}
\[
\begin{tikzcd}
\mathcal{M}' \arrow[r, shift right=1, "v^{*}"'] & \mathcal{M} \arrow[l, shift right=1, "v_{!}"']
\end{tikzcd}
\tag*{(1.5.1)}
\]
\textit{un couple de foncteurs adjoints, soient}
\[
\Phi \ \subset\ \mathrm{Fl}(\mathcal{M})\ , \qquad \Psi' \ \subset\ \mathrm{Fl}(\mathcal{M}')\ ,
\]
\textit{alors on a}
\[
(v_{!}(\Phi) \Qorth \Psi') \qquad \Longleftrightarrow \qquad (\Phi \Qorth v^{*}(\Psi'))\ . \tag*{(1.5.2)}
\]

\textbf{Corollaire 1.6.} \textit{Pour toute partie $\Phi$ de} $\mathrm{Fl}(\mathcal{M})$,
\[
(v_{!}(\Phi))_{*} \ =\ (v^{*})^{-1}(\Phi_{*})\ . \tag*{(1.6.1)}
\]

\printedpage{5}
% head: A. Grothendieck : Dérivateurs, Ch. XIX, Édition des Universités de Montpellier II et Paris VII
% foot: Version 11 février 2019

En effet, appliquant 1.5 à $\Phi$ et à $\{f'\}$, où $f' \in \mathrm{Fl}(\mathcal{M}')$, on trouve
\[
f' \in (v_{!}(\Phi))_{*} \qquad \Longleftrightarrow \qquad v^{*}(f') \in \Phi_{*}\ ,
\]
ce qui équivaut à (1.6.1). Dualement :

\textbf{Corollaire 1.7.} \textit{Pour toute partie $\Psi'$ de $\mathcal{M}'$, on a}
\[
(v^{*}(\Psi'))^{*} \ =\ (v_{!})^{-1}(\Psi'^{*}). \tag*{(1.7.1)}
\]

En effet, si $f \in \mathrm{Fl}(\mathcal{M})$, on a par (1.5.2)
\[
f \in (v^{*}(\Psi'))^{*} \qquad \Longleftrightarrow \qquad v_{!}(f) \in \Psi^{*}\ . % [sic] \Psi^{*} -> \Psi'^{*}, cf. (1.7.1)
\]
\textsuperscript{(}\footnote[3]{donner ici 1.9 page 9.}\textsuperscript{)}.

\msp{page 8}
Un couple Q-orthogonal $\Phi \Qorth \Psi$ est dit \textit{spécial} s'il satisfait les conditions suivantes.
\[
\left\{
\begin{array}{ll}
\text{a)} & \parbox[t]{0.7\textwidth}{$\mathcal{M}$ est stable par petites $\varinjlim$ , les objets de $\mathcal{M}$ sont accessibles. (On dit alors que $\mathcal{M}$ est “spéciale”.)} \\[3ex]
\text{b)} & \text{Le couple est clos : } \Phi \clos \Psi. \\[1ex]
\text{c)} & \Phi \text{ est \textit{spéciale.}}
\end{array}
\right.
\tag*{(1.10)}
\]
Cela implique donc, par 1.1, que $\Phi \Qcouple \Psi$ est un couple de Quillen. Si $\mathcal{M}$ satisfait a), les couples Q-orthogonaux spéciaux dans $\mathcal{M}$, ou encore les couples de Quillen spéciaux, correspondent donc biunivoquement aux parties $\Phi$ de $\mathrm{Fl}(\mathcal{M})$ qui sont \textit{spéciales} et closes à gauche (ou ce qui revient au même, Q-saturées à gauche), en associant à une telle $\Phi$ le couple $\Phi \Qcouple \Phi_{*}$.

\dgrblock{dgrb1162}{An adjunction transports the whole lifting problem}
\begin{dcomment}{exposition}{p7}{An adjunction transports the whole lifting problem}
The next passage prepares the construction of diagrams by proving a statement more economical than a transfer theorem. If $L=v_!$ is left adjoint to $R=v^*$, a commutative square from $Li$ to $p'$ corresponds to a commutative square from $i$ to $Rp'$. This correspondence includes the diagonal fillers: a map from the codomain of $Li$ to the domain of $p'$ corresponds to a map from the codomain of $i$ to the domain of $Rp'$, and the two triangle equations correspond by naturality of the adjunction. Thus
\[
 Li\perp p'\quad\Longleftrightarrow\quad i\perp Rp'.
\]
There is no need for unique fillers, and no factorization hypothesis enters this equivalence.

Applying it one arrow at a time gives
\[
 (L\Phi)^{\square}=R^{-1}(\Phi^{\square}),\qquad
 {}^{\square}(R\Psi')=L^{-1}({}^{\square}\Psi').
\]
Here the inverse image means inverse image of a class of arrows, rather than an inverse functor. The unprimed $\Psi^*$ in the proof of Corollary 1.7 is a printed slip: its ambient category forces $\Psi'{}^*$.

The same calculation explains why a left adjoint preserves generation by a set even when it does not literally preserve a chosen list of cellular presentations. If $\Phi\subseteq{}^{\square}(\Phi_0^{\square})$, then $\Phi_0^{\square}\subseteq\Phi^{\square}$; inverse image under $R$ gives
\[
 (L\Phi_0)^{\square}\subseteq(L\Phi)^{\square},
\]
and taking left orthogonals gives $L\Phi\subseteq{}^{\square}((L\Phi_0)^{\square})$. This is Corollary 1.9. It needs neither accessibility nor cocompleteness. The extra hypotheses of Corollary 1.8 are used only after this calculation: the small object argument in the target produces the factorizations and identifies the left class with the retracts of the relative $L\Phi_0$-cell complexes.

This order of reasoning is helpful throughout the chapter. First identify the desired right class as an inverse image. Then identify a small family that defines it. Only then apply a factorization theorem in the target. In particular, the source category's ``specialness'' in Corollary 1.8 is stronger than is required for the transport of generators itself; Grothendieck explicitly observes this when extracting Corollary 1.9. The later transfer of a \emph{model structure} will require an additional condition on weak equivalences. That condition is absent here because a single weak factorization system has no class of weak equivalences to preserve.
\end{dcomment}


Ceci posé :

\marginpar{$\Vert$}% double vertical bar printed in the left margin beside Corollaire 1.8
\textbf{Corollaire 1.8.} \textit{Supposons que $\mathcal{M}$, $\mathcal{M}'$ satisfassent la condition} a) \textit{de} (1.10), \textit{et soit $(\Phi, \Psi)$ un couple de Quillen spécial dans $\mathcal{M}$. Alors $\Bigl(\Phi' = v_{!}(\Phi)^{\sim}, \Psi' = (v^{*})^{-1}(\Psi)\Bigr)$ est un couple de Quillen spécial dans $\mathcal{M}'$.}

En effet, par le corollaire 1.6 on a
\[
\Phi'_{*} \ =\ (v_{!}(\Phi))_{*} \ =\ (v^{*})^{-1}(\Psi) \ =\ \Psi'\ ,
\]
donc $\Phi' \clos \Psi'$. Il reste à prouver que $\Phi'$ est spécial. Soit $\Phi_0$ une partie petite de $\Phi$ telle que $\Phi \subset \Phi_0^{\sim}$, posons $\Phi'_0 \ =\ v_{!}(\Phi_0) \subset \Phi'$, \msp{page 9} il suffit de voir qu'on a
\[
v_{!}(\Phi) \ \subset\ v_{!}(\Phi_0)^{\sim}\ , \tag*{($*$)}
\]
d'où résultera qu'on a $\Phi' \overset{\text{déf}}{=} (v_{!}(\Phi))^{\sim} \subset v_{!}(\Phi_0)^{\sim}$. Mais on a
\[
v_{!}(\Phi_0)^{\sim} \ =\ (v_{!}(\Phi_0)_{*})^{*}\ ,
\]
\printedpage{6}
% head: A. Grothendieck : Dérivateurs, Ch. XIX, Édition des Universités de Montpellier II et Paris VII
% foot: Version 11 février 2019
et
\[
v_{!}(\Phi_0)_{*} \ =\ (v^{*})^{-1}(\underbrace{\Phi_{0\,*}}_{\Psi}) \ =\ (v^{*})^{-1}(\Psi) \ \overset{\text{déf}}{=}\ \Psi'\ ,
\]
donc
\[
v_{!}(\Phi_0)^{\sim} \ =\ \Psi'^{*} \ =\ \Phi' \ =\ v_{!}(\Phi)^{\sim}\ ,
\]
d'où ($*$). Cet argument ne fait pas appel à l'hypothèse “spéciale” sur $\mathcal{M}$, et peut s'expliciter en le

\textbf{Corollaire 1.9.} \textsuperscript{(}\footnote[4]{devrait passer avant 1.8.}\textsuperscript{)} \textit{Sous les conditions générales de} 1.5 \textit{($v_{!} : \mathcal{M} \fl \mathcal{M}'$ étant un foncteur admettant un adjoint à droite), soient $\Phi_0$, $\Phi$ deux parties de} $\mathrm{Fl}(\mathcal{M})$ \textit{telles que $\Phi \subset \Phi_0^{\sim}$ (ou ce qui revient au même, $\Phi^{\sim} \subset \Phi_0^{\sim}$). Alors $v_{!}(\Phi) \subset v_{!}(\Phi_0)^{\sim}$, i.e. $v_{!}(\Phi)^{\sim} \subset v_{!}(\Phi_0)^{\sim}$.}

En effet, cette relation équivaut aussi à $v_{!}(\Phi_0)_{*} \ \subset\ v_{!}(\Phi)_{*}$, c'est-à-dire (par 1.6) $(v^{*})^{-1}(\Phi_{0\,*}) \ \subset\ (v^{*})^{-1}(\Phi_{*})$, et est impliqué par $\Phi_{0\,*} \ \subset\ \Phi_{*}$. Cette dernière d'autre part équivaut à l'hypothèse $\Phi^{\sim} \subset \Phi_0^{\sim}$, q.e.d.

\msp{page 10}
Si $\Psi$ est une partie de $\mathrm{Fl}(\mathcal{M})$, on en déduit pour toute $I$ dans Cat une partie
\[
\left\{
\begin{array}{l}
\Psi(I) \subset \mathrm{Fl}(\mathcal{M}(I)), \quad \text{avec}\ \ \mathcal{M}(I) = \underline{\mathrm{Hom}}(I^{\circ}, \mathcal{M})\ , \\[1ex]
f \in \Psi(I) \iff f(i) \in \Psi \ \ \forall\, i \in \mathrm{Ob}\, I\ .
\end{array}
\right.
\tag*{(1.11)}
\]
Si on introduit la catégorie discrète $I_0$ sous-jacente à $\mathrm{Ob}\, I$, et le foncteur canonique
\[
v_I = v \ :\ I_0 \fl I\ , \tag*{(1.12)}
\]
on trouve un couple de foncteurs adjoints
\[
\begin{tikzcd}[row sep=small]
\mathcal{M}(I) \arrow[r, shift right=1, "v^{*}"'] & \mathcal{M}(I_0) \arrow[l, shift right=1, "v_{!}"'] \\
( & ) \arrow[l, dotted, "v_{*}"]
\end{tikzcd}
\ \simeq\ \mathcal{M}^{I_0}\ ,
\tag*{(1.13)}
\]
et on a
\[
\Psi(I) \ =\ (v^{*})^{-1}(\Psi(I_0))\ .
\]
Supposons maintenant qu'on ait
\[
\left\{
\begin{array}{rcl}
\Psi & = & \Phi_{*} \\[1ex]
\text{d'où}\ \ \Psi(I_0) & = & \Phi(I_0)_{*}
\end{array}
\right.
\tag*{(1.14)}
\]
(vérification immédiate), on conclut alors par 1.6
\[
\Psi(I) \ =\ (v_{!}(\Phi(I_0)))_{*} \ =\ (\Phi\langle I\rangle)_{*}\ , \tag*{(1.15)}
\]
avec
\[
\Phi\langle I\rangle \ \overset{\text{déf}}{=}\ v_{!}(\Phi(I_0))^{\sim} \ =\ \Psi(I)^{*}\ . \tag*{(1.16)}
\]
Si on a un foncteur
\[
u \ :\ I \fl I'\ , \tag*{(1.17)}
\]
d'où un triple de foncteurs adjoints
\printedpage{7}
% head: A. Grothendieck : Dérivateurs, Ch. XIX, Édition des Universités de Montpellier II et Paris VII
% foot: Version 11 février 2019
\msp{page 11}
\[
\left\{
\begin{tikzcd}[column sep=large]
\mathcal{M}(I) \arrow[r, shift left=3, "u_{!}"] \arrow[r, shift right=3, "u_{*}"] & \mathcal{M}(I') \arrow[l, "u^{*}"']
\end{tikzcd}
\right.
,
\tag*{(1.18)}
\]
on aura de façon $\pm$ tautologique
\[
\left\{
\begin{array}{rcl}
u_{!}(\Phi\langle I\rangle) & \subset & \Phi\langle I'\rangle \\
u^{*}(\Psi(I')) & \subset & \Psi(I)\ .
\end{array}
\right.
\tag*{(1.19)}
\]
Le corollaire 1.8 donne :

\dgrblock{dgrb1163}{Free diagrams and the projective construction}
\begin{dcomment}{exposition}{p11:on aura de façon $\pm$ tautologique}{Free diagrams and the projective construction}
The reason for passing through the discrete category $I_0$ is that independent lifting problems can then be solved component by component. The left Kan extension freely inserts the naturality required by the arrows of $I$. Throughout this chapter $\mathcal M(I)=\operatorname{Fun}(I^{\mathrm{op}},\mathcal M)$. For an object $i$ of $I$, the left adjoint to evaluation at $i$ is
\[
 i_!A(j)=\coprod_{I(j,i)} A.
\]
Thus the free diagram on a family $(A_i)$ has value $\coprod_i\coprod_{I(j,i)}A_i$ at $j$. This formula fixes the variance: it is the representable presheaf $I(-,i)$, not $I(i,-)$, that occurs.

Suppose $(\Phi,\Psi)$ is a mutually orthogonal pair and $S\subseteq\Phi$ generates it. The set
\[
 S_I=\{i_!s\mid i\in\operatorname{Ob}I,\ s\in S\}
\]
generates the projective pair. Indeed, $p$ has the right lifting property against $i_!s$ precisely when $p(i)$ has the right lifting property against $s$. Hence
\[
 S_I^{\square}=\Psi(I),\qquad
 {}^{\square}(S_I^{\square})=\Phi\langle I\rangle.
\]
This makes the smallness in Proposition 1.10 concrete. One need not use the much larger class of all objectwise $\Phi$-arrows on the discrete category. If the domain $A$ of $s$ is small, then $i_!A$ is small for pointwise colimits because
\[
\operatorname{Hom}(i_!A,X)=\operatorname{Hom}(A,X(i)).
\]
The same common cardinal as in the original small object argument therefore suffices for these generators.

There is a small scope error in the preliminary equation (1.14). Its equality $\Psi(I_0)=\Phi(I_0)^{\square}$ is valid for the closed left classes used in Proposition 1.10: identities in $\Phi$ allow one to isolate a chosen coordinate, using identity arrows on the domains of the right-hand maps at all other coordinates. It is not valid for an arbitrary $\Phi$ merely satisfying $\Psi=\Phi^{\square}$. For instance, let $\Phi=\{\varnothing\to1\}$ in $\mathrm{Set}$ and let $I_0$ have two objects. The arrow $(\varnothing\to\varnothing,\varnothing\to1)$ has the right lifting property against the sole member $(\varnothing\to1,\varnothing\to1)$ of $\Phi(I_0)$, because no square can be formed in the first coordinate. Its second component is not surjective. Adding identities to $\Phi$ repairs this preliminary assertion; it is automatic in all the subsequent applications to a closed pair.

The source also asserts that the diagram category itself is special. Cocompleteness is pointwise. To check its object-smallness condition, fix $X:I^{\mathrm{op}}\to\mathcal M$ and choose one regular cardinal dominating the ranks of its set of values and the cardinality of the set of arrows of $I$. A natural transformation out of $X$ is a family of component maps subject to one equation for each arrow of $I$. In a sufficiently filtered colimit, all these component maps occur at one stage, and all their naturality equations hold together at a later stage. Equality of two transformations is detected similarly. This proves smallness of $X$. If one uses the stronger filtered-presentability convention for accessible objects, the same argument is expressed by the equalizer of the two products computing the set of natural transformations.

Restriction plainly preserves the objectwise class $\Psi(I')$. Its left adjoint consequently preserves $\Phi\langle I\rangle$, by the preceding lifting correspondence. The dual construction uses objectwise left arrows $\Phi(I)$ and defines $\Psi[I]=\Phi(I)^{\square}$. This is the injective pair; its factorization is a different question. Any displayed right Kan extension is understood where its required limits exist. Cocompleteness and smallness of individual objects alone do not imply completeness: the ordered class of ordinals, with arrows $\alpha\to\beta$ for $\alpha\leq\beta$, has all small suprema and small individual objects, but no terminal object. This observation keeps the existence of the right adjoints separate from the lifting calculations that use them.
\end{dcomment}


\textbf{Proposition 1.10.} \textit{Supposons que $\Phi \Qcouple \Psi$ soit un couple de Quillen} spécial \textit{dans $\mathcal{M}$, i.e. de la forme $(\Phi, \Phi_{*})$, avec $\Phi$ partie Q-close à gauche spéciale dans $\mathcal{M}$ ($\mathcal{M}$ satisfaisant la condition} (1.10. a))\textit{. Alors pour toute $I$ dans} Cat, \textit{$(\Phi\langle I\rangle, \Psi(I))$ est un couple spécial dans $\mathcal{M}(I)$, en particulier il satisfait au théorème de factorisation.}

Considérons la construction duale, qui nous fait introduire, à partir d'un couple $\Phi \clos \Psi$ dans $\mathcal{M}$, le couple
\[
\Phi(I) \ \clos\ \Psi[I]\ , \tag*{(1.20)}
\]
où $\Phi(I)$ est défini par la formule générale (1.11) (où $\Psi$ est remplacé par $\Phi$), et où $\Psi[I]$ est défini soit comme $\Phi(I)_{*}$, soit comme
\[
\Psi[I] \ =\ v_{*}(\Psi(I_0))^{-} \ =\ \Phi(I)_{*}\ . \tag*{(1.21)}
\]

\msp{page 12}
On a alors, pour $u : I \fl I'$
\[
\left\{
\begin{array}{rcl}
u^{*}(\Phi(I')) & \subset & \Phi(I) \\
u_{*}(\Psi[I]) & \subset & \Psi[I']\ .
\end{array}
\right.
\tag*{(1.22)}
\]
Il se pose encore la question si le couple $(\Phi(I), \Psi[I])$ satisfait au théorème de factorisation, sous des hypothèses convenables sur $(\Phi, \Psi)$. Je peux seulement hasarder la

\textbf{Question-conjecture 1.11.} \textit{Supposons que $(\Phi, \Psi)$ soit un couple de Quillen} spécial, \textit{en est-il de même de $(\Phi(I), \Psi[I])$ pour toute $I$ dans} (Cat) \textit{? En d'autres termes, si $\mathcal{M}$ est spéciale, condition} (1.10. a))\textit{, et si $\Phi$ est une partie Q-close à gauche} spéciale \textit{de $\mathcal{M}$, i.e. admettant une petite partie Q-génératrice, est-il vrai que pour toute $I$ dans} Cat, \textit{$\Phi(I)$ est également spéciale ? Je conjecture que oui, en faisant au besoin des hypothèses supplémentaires sur $\mathcal{M}$ (existence d'une filtration cardinale ?). P.ex. si $\Phi_0$ est Q-génératrice dans $\Phi$, et petite, alors $\Phi_0(I)$ est-elle Q-génératrice dans $\Phi(I)$ ?}

\dgrblock{dgrb1164}{The injective question under local presentability}
\begin{dcomment}{later}{p12:Je conjecture que oui, en faisant au besoin}{The injective question under local presentability}
Grothendieck's first question in 1.11 has a precise affirmative answer after the additional hypothesis of local presentability. In a locally presentable category, a weakly saturated class generated by a set remains generated by a set when imposed objectwise on a small diagram category. This is exactly \citet[Lemma~A.2.8.3]{Lurie2009HTT}; its proof is supplied in the more general Lemma~A.3.3.3. The corresponding statement for combinatorial categories, without introducing weak equivalences, is \citet[Corollary~3.4]{MakkaiRosicky2014}, using the numbering of the 2013 arXiv version. These statements apply directly to a cofibrantly generated weak factorization system and prove that $(\Phi(I),\Psi[I])$ is again cofibrantly generated.

The additional hypothesis is substantive. A category being cocomplete with individually accessible objects does not, by itself, specify a set of presentable objects generating the whole category under filtered colimits. The cited theorems therefore settle the locally presentable form of the question, not its entire literal generality. Our literature check did not establish a theorem with exactly that weaker generality. The final sentence of 1.11 asks for more than existence of \emph{some} generating set: it proposes the particular set $\Phi_0(I)$. That proposal fails, as the following calculation shows.
\end{dcomment}


\dgrblock{dgrb1165}{Fixed points detect the missing generators}
\begin{dcomment}{examples}{p12:Je conjecture que oui, en faisant au besoin}{Fixed points detect the missing generators}
The difficulty with the proposed generators can be seen without homotopy theory. Work in $\mathrm{Set}$ with its usual choice convention. The single map $\varnothing\to 1$ generates the weak factorization system of injections and surjections: its right lifting property is surjectivity, injections lift against surjections by choosing preimages on their complements, and the standard factorization is
\[
 X\longrightarrow X\amalg Y\xrightarrow{[f,1_Y]}Y.
\]
Take $I=BG$ for a nontrivial group $G$, say $G=C_2$. The diagram category consists of right $G$-sets. The class $\Phi_0(I)$ contains only the map from the empty $G$-set to the one-point trivial $G$-set. A $G$-map $p:X\to Y$ has its right lifting property if and only if $X^G\to Y^G$ is surjective: a map from the trivial point chooses a fixed point.

Now consider $p:\varnothing\to G$, where $G$ acts on itself freely. Both fixed-point sets are empty, so $p$ passes the proposed test. But the objectwise monomorphism $\varnothing\to G$ does not lift against $p$: in the square whose lower horizontal map is $1_G$, a filler would be a map $G\to\varnothing$. Thus $\Phi_0(I)$ does not Q-generate $\Phi(I)$.

The correct generating family in this example is
\[
 \{\varnothing\to G/H\mid H\leq G\},
\]
with right-coset or left-coset notation adjusted to the chosen action convention. An equivariant injection is obtained by attaching the orbits of its invariant complement; conversely, such attachments are injections, and injections are preserved under retracts. The right class consists of maps surjective on $H$-fixed points for every subgroup $H$. To see this directly, a map from a transitive orbit to $Y$ is a point fixed by the stabilizer of that orbit, and a lift is a fixed preimage. Choosing such lifts orbit by orbit solves an arbitrary equivariant lifting problem against an injection. The single proposed test sees only $H=G$ and misses the free orbit $H=\{1\}$. This fully accounts for the failure; it does not contradict existence of the injective pair.
\end{dcomment}


\dgrblock{dgrb1166}{Exactly when the original singleton still generates}
\begin{dcomment}{development}{p12:Je conjecture que oui, en faisant au besoin}{Exactly when the original singleton still generates}
The preceding example suggests a sharper answer to Grothendieck's last sentence: one can characterize all diagram shapes for which his proposed singleton works. The following elementary result is our development of that sentence. We checked the generator constructions of \citet[Lemma~A.2.8.3]{Lurie2009HTT} and \citet[Corollary~3.4]{MakkaiRosicky2014}; we did not locate this precise characterization there. It is a calculation with constant presheaves, not a claim of priority.

\textbf{Proposition (the constant-generator test).}\ours
Let $I$ be a small category, let $\Delta$ denote the constant-presheaf functor, and let $\Phi_0=\{\varnothing\to1\}$ in $\mathrm{Set}$. Then $\Phi_0(I)$ Q-generates all monomorphisms in $\mathrm{Set}^{I^{\mathrm{op}}}$ if and only if $I$ is empty or equivalent to the terminal category.

\textit{Proof.} The class $\Phi_0(I)$ consists of the single map $\varnothing\to\Delta1$: all its component domains and codomains are prescribed, and their transition maps are forced. A pushout of a coproduct of copies of this map adds a constant summand. Consequently every relative cell complex starting from the initial presheaf is $\varnothing\to\Delta S$ for some set $S$. This includes transfinite cell complexes, since colimits of the successive constant summands are computed pointwise. The small object argument applies, since the initial presheaf is small. Therefore every object $P$ for which $\varnothing\to P$ belongs to the generated left class is a retract of some $\Delta S$.

Suppose now that this left class contains every monomorphism and that $I$ is nonempty. For every $i\in I$, the representable $h_i=I(-,i)$ is then a retract of a constant presheaf. Choose natural transformations
\[
 h_i\xrightarrow{a}\Delta S_i\xrightarrow{b}h_i,
 \qquad ba=1_{h_i}.
\]
For any arrow $\alpha:j\to k$ of $I$, naturality gives
\[
 a_jh_i(\alpha)=a_k,\qquad
 h_i(\alpha)b_k=b_j.
\]
It follows that $h_i(\alpha)$ has inverse $b_ka_j$: composing in either order and using the two displayed equations yields the appropriate identity. Since this holds for every $i$, every map
\[
 I(k,i)\longrightarrow I(j,i),\qquad t\longmapsto t\alpha
\]
is bijective. This forces $\alpha$ to be invertible. Explicitly, surjectivity with $i=j$ gives $\beta:k\to j$ with $\beta\alpha=1_j$; injectivity with $i=k$ then gives $\alpha\beta=1_k$, since $(\alpha\beta)\alpha=\alpha=1_k\alpha$. Thus $I$ is a groupoid.

Moreover, $S_i$ is nonempty because $h_i(i)$ contains $1_i$ and maps into $S_i$. For every $j$, the map $b_j:S_i\to I(j,i)$ therefore shows that $I(j,i)$ is nonempty. The groupoid is connected. Finally, an endomorphism $\gamma:i\to i$ acts trivially on $h_i(i)$: naturality gives $a_ih_i(\gamma)=a_i$, and $a_i$ is injective because it has a left inverse. Applying this identity to $1_i$ gives $\gamma=1_i$. Every isotropy group is therefore trivial. A connected groupoid with trivial isotropy has exactly one arrow between any two objects, so it is equivalent to the terminal category.

Conversely, if $I$ is equivalent to the terminal category, evaluation at any object is an equivalence $\mathrm{Set}^{I^{\mathrm{op}}}\simeq\mathrm{Set}$ and takes $\varnothing\to\Delta1$ to $\varnothing\to1$. Equivalences preserve and reflect lifting properties and monomorphisms, so the stated generation follows from the set calculation. If $I$ is empty, the presheaf category has one object and one arrow; both the generated left class and its class of monomorphisms consist of that arrow. This proves all cases.\hfill$\square$

The statement concerns the proposed \emph{unchanged} generator. Even a discrete category with two objects fails the test: it asks one attachment to add a point simultaneously in both coordinates. The affirmative injective-generation theorem supplies a different set, large enough to express the independent and compatible attachments required by the chosen shape.
\end{dcomment}


\textbf{Proposition 1.12} \textsuperscript{(}\footnote[5]{cf. XVII 6.4 (p. 117), mais où les notations pour $\Phi\langle I\rangle$, $\Psi(I)$ sont inverses.}\textsuperscript{)}. \textit{Soit $\Phi \clos \Psi$ dans $\mathcal{M}$, soit $u : I \fl I'$ une flèche dans} Cat, \textit{d'où}
\[
\left\{
\begin{tikzcd}[column sep=large]
\mathcal{M}(I) \arrow[r, shift left=3, "u_{!}"] \arrow[r, shift right=3, "u_{*}"] & \mathcal{M}(I') \arrow[l, "u^{*}"']
\end{tikzcd}
\right.
.
\tag*{(1.12.1)}
\]
\printedpage{8}
% head: A. Grothendieck : Dérivateurs, Ch. XIX, Édition des Universités de Montpellier II et Paris VII
% foot: Version 11 février 2019
\textit{On a alors} (1.19), \textit{i.e.}
\[
\left\{
\begin{array}{rcl}
u_{!}(\Phi\langle I\rangle) & \subset & \Phi\langle I'\rangle \\
u^{*}(\Psi(I')) & \subset & \Psi(I)\ ,
\end{array}
\right.
\tag*{(1.13.1)}
\]
\textit{et de plus l'équivalence}
\[
u^{*}(\Phi(I')) \subset \Phi\langle I'\rangle \qquad \Longleftrightarrow \qquad u_{*}(\Psi(I)) \subset \Psi(I')\ . \tag*{(1.13.2)}
\]
\textit{Ces relations équivalentes sont satisfaites dans chacun des trois cas}~\textsuperscript{(}\footnote[6]{\textbf{N.B.} On a $1^{\circ} \Longrightarrow 3^{\circ}$, $2^{\circ} \Longrightarrow 3^{\circ}$.}\textsuperscript{)} \textit{:}
\begin{enumerate}[label={}]
\item[1°)] \textit{$u$ est coétale.}
\item[2°)] \textit{$I$ est discrète.}
\item[3°)] $\forall\, i' \in \mathrm{Ob}\, I'$\textit{, toute composante connexe de $I/i'$ admet un objet final.}
\end{enumerate}
\msp{page 13}

\dgrblock{dgrb1167}{The comma-category test for preserving projective arrows}
\begin{dcomment}{exposition}{p12:Ces relations équivalentes sont satisfaites}{The comma-category test for preserving projective arrows}
Proposition 1.12 identifies shapes for which restriction preserves the projective left class. This is a useful exception to the usual asymmetry: restriction always preserves objectwise right arrows, but it need not preserve the left arrows defined by a lifting property against them.

The printed formula (1.13.2) contains two slips that have to be corrected in interpreting the statement. Its left side should read
\[
 u^*\bigl(\Phi\langle I'\rangle\bigr)
       \subseteq\Phi\langle I\rangle,
 \tag{*}
\]
not $u^*(\Phi(I'))\subseteq\Phi\langle I'\rangle$. The latter has the wrong ambient category on the right and the wrong, objectwise, class on the left. The intended equivalence follows by applying the lifting correspondence to the adjunction $u^*\dashv u_*$:
\[
\begin{aligned}
u^*\Phi\langle I'\rangle\subseteq{}^{\square}\Psi(I)
&\Longleftrightarrow
\Phi\langle I'\rangle\perp u_*\Psi(I)\\
&\Longleftrightarrow
u_*\Psi(I)\subseteq\Psi(I').
\end{aligned}
\]
The source's established mutual-orthogonality identities are exactly what the last step needs.

To prove the sufficient condition, let $K_{i'}=(u\downarrow i')$: its objects are pairs $(i,u(i)\to i')$. For presheaves, the pointwise right Kan extension is
\[
 (u_*X)(i')=\lim_{K_{i'}^{\mathrm{op}}}X(i).
\]
Suppose each connected component $c$ of $K_{i'}$ has a terminal object $t_c=(i_c,u(i_c)\to i')$. Then $t_c$ is initial in the opposite category, and a limit over that component is evaluation at $i_c$. Therefore
\[
 (u_*X)(i')\cong\prod_cX(i_c).
\]
For an objectwise arrow in $\Psi(I)$, the induced arrow is a product of arrows in $\Psi$. Any right lifting class is closed under products: project a lifting square to each factor, choose its fillers, and combine them by the product property. The empty product gives an identity of a terminal object. Thus $u_*$ preserves $\Psi(I)$ and (*) follows. This supplies the proof left implicit in the statement.

If $I$ is discrete, every $K_{i'}$ is discrete and each nonempty component is a singleton. For a coétale functor in the source's convention, lifting a map into $i'$ identifies a terminal lift in each component of $K_{i'}$. These are the two advertised special cases of the third condition. All formulas use the relevant products or right Kan extensions where they exist. No factorization theorem and no accessibility hypothesis are needed for this argument; those enter only when one wants the orthogonal pairs themselves to give replacement factorizations.
\end{dcomment}


\section*{2\quad Rappels sur les triples de Quillen clos et les triples de Quillen spéciaux}

Soit $(\mathcal{M}, W)$ une \textit{catégorie de modèles,} i.e. $\mathcal{M}$ munie d'un localiseur $W \subset \mathrm{Fl}(\mathcal{M})$, i.e. une partie satisfaisant
\[
\parbox{0.75\textwidth}{$W$ contient les isomorphismes, et si $f$, $g$ sont des flèches composables, si deux parmi $f$, $g$, $gf$ sont dans $W$, la troisième aussi \textit{(saturation faible).}}
\tag*{(2.1)}
\]
Je rappelle que la donnée d'un \textit{triple de Quillen clos} $(W, C, F)$\textit{,} ou d'une structure de Quillen close, sur $\mathcal{M}$, ayant $W$ comme ensemble de quasi-isomorphismes, revient à la donnée de deux couples de Quillen
\[
C \Qcouple TF\ , \qquad TC \ \Qcouple\ F\ , \tag*{(2.2)}
\]
satisfaisant
\[
TC \ \subset\ C \cap W\ , \qquad TF \ \subset\ F \cap W\ , \tag*{(2.3)}
\]
plus une des trois conditions équivalentes :
\[
\left\{
\begin{array}{lrcll}
\text{a)} & TC & = & C \cap W & \qquad (\text{i.e. } C \cap W \subset TC) \\[1ex]
\text{b)} & TF & = & F \cap W & \qquad (\text{i.e. } F \cap W \subset TF) \\[1ex]
\text{c)} & W & = & TF \circ TC\ . &
\end{array}
\right.
\tag*{(2.4)}
\]
Rappelons aussi que ceci implique que $W$ \textit{est fortement saturée,} et a fortiori que $W = W^{\natural}$, i.e. $W$ stable par facteurs directs.

Rappelons aussi qu'en vertu de (1.4), on a :

% ---- from der-ch19_p011-015.tex
\printedpage{9}
% head: A. Grothendieck : Dérivateurs, Ch. XIX, Édition des Universités de Montpellier II et Paris VII
% foot: Version 11 février 2019

\msp{page 14}

\dgrblock{dgrb1168}{How the two factorizations recognize weak equivalences}
\begin{dcomment}{exposition}{p13}{How the two factorizations recognize weak equivalences}
The purpose of the second section is to add compatibility with a localizer to the machinery already established for one weak factorization system. Grothendieck's ``category of models'' at (2.1) means a category with a specified two-out-of-three class. It should not be silently identified with a modern model category. The model structure enters through the two factorizing pairs
\[
 (C,TF),\qquad (TC,F),
\]
with $TC\subseteq C$ and $TF\subseteq F$. The letters $T$ stand for the arrows intended to be trivial, and that intention has to be checked against $W$. For the equivalence in (2.4), the source explicitly assumes $TC\subseteq C\cap W$ and $TF\subseteq F\cap W$ before asking whether these inclusions are equalities.

The equivalence of the three conditions (2.4) is a short but instructive use of both factorizations. Suppose first that $TC=C\cap W$. If $f\in W$, factor $f=pi$ using $(C,TF)$. Since $p\in W$, two-out-of-three gives $i\in C\cap W=TC$, hence $W\subseteq TF\circ TC$. The reverse inclusion follows because both factors belong to $W$. Conversely, if $W=TF\circ TC$ and $f\in C\cap W$, choose such a factorization. Since $f\perp p$, the retract argument makes $f$ a retract of $i\in TC$, so $f\in TC$. This proves equivalence of (a) and (c). Using the other factorization and the dual retract argument proves equivalence with (b).

Formula (2.5) is another application of the same idea. To test $TC\subseteq W$ when $W$ is retract-closed, it suffices to factor each arrow of $TC$ as a $W$-arrow followed by an $F$-arrow; the left lifting property makes the original arrow a retract of the first factor. If $TC\subseteq W$ already holds, the $(TC,F)$-factorization supplies this for all arrows. The dual test is a factorization $f=pi$ with $i\in C$ and $p\in W$. The printed last clause ``$i\in W$, $f\in C$'' is incorrect and is retained as an edition slip, with its intended dual stated here.

The ``special'' condition adds small generators for both left classes, together with the ambient smallness hypotheses that make their small object arguments work. It is closely related to cofibrant generation in contemporary terminology, but stronger ambient hypotheses are explicitly included here. It is not synonymous with ``combinatorial'': local presentability is not part of (2.6). The recalled strong saturation of $W$ uses the ambient model-category convention: \citet[Chapter~XIII, \S3, manuscript pp.~27--29]{Grothendieck1990XIII} explicitly assumes M0, existence of finite limits and finite colimits, and invokes Quillen's saturation result. Strong saturation means that precisely the arrows of $W$ become invertible in $\mathcal M[W^{-1}]$; retract-closedness follows since a retract of an isomorphism is an isomorphism. A modern reference is \citet[Theorem~1.2.10(iv)]{Hovey1999}. The audit records this as a recalled theorem with its inherited ambient hypotheses, not as a gap in the present proof.
\end{dcomment}


\begin{itemize}[leftmargin=4em]
\item[(2.5)] Étant donnés les deux couples (2.2), avec $TC \subset C$ (équivalent à $TF \subset F$), pour qu'on ait la première relation (2.3), i.e. $TC \subset W$, il faut, et il suffit si $W = W^{\natural}$, que toute flèche $f$ de $\M$ (ou seulement toute flèche de $TC$) se factorise en $f = pi$, avec $i \in W$, $p \in F$. Dualement, pour qu'on ait $TF \subset W$, il faut et il suffit si $W = W^{\natural}$, que toute flèche $f$ de $\M$ (ou seulement toute flèche de $TF$) se factorise en $f = pi$, avec $i \in W$, $f \in C$. % [sic] i \in W, f \in C -> i \in C, p \in W (?), dual of the preceding factorisation with p \in F
\end{itemize}

On dit que la structure de Quillen $(W,C,F)$ est \textit{spéciale} si les deux couples de Quillen (2.3) sont spéciaux, i.e. si les conditions suivantes sont satisfaites.

\begin{itemize}[leftmargin=4em]
\item[(2.6)]
\begin{itemize}
\item[a)] $\M$ est stable par petites $\varinjlim$, et ses éléments sont accessibles ($\M$ “spéciale”).
\item[b)] Il existe des \textit{petites} parties $C_0 \subset C$, $TC_0 \subset TC$ qui soient Q-génératrices à gauche, i.e. telles que l'on ait
\[
\left\{
\begin{array}{rclcrclcrcl}
C & \subset & C_0^{\sim}\ , & \text{i.e.} & C & = & C_0^{\sim}\ , & \text{i.e.} & TF & = & C_{0*}\ ;\\
TC & \subset & (TC_0)^{\sim}\ , & \text{i.e.} & TC & = & (TC_0)^{\sim}\ , & \text{i.e.} & F & = & (TC_0)_*\ .
\end{array}\right. % reading resolved: DER_XIX_corrections.csv
\]
\end{itemize}
\end{itemize}

Ainsi, la donnée d'un tel triple de Quillen \textit{spécial} dans $\M$, quand $\M$ satisfait (2.6. a), revient à la donnée des deux parties
\msp{page 15}
\[
TC \ \subset\ C
\tag*{(2.7)}
\]
de $\Fl(\M)$, satisfaisant les conditions :
\[
\left\{
\begin{array}{ll}
C,\ TC\ \text{“spéciales”} &
\left\{\begin{array}{ll}
\text{a)} & \text{$C$, $TC$ sont Q-closes à gauche.}\\[1ex]
\text{b)} & \parbox[t]{0.45\textwidth}{$C$, $TC$ sont Q-engendrées à gauche par des \textit{petites} parties $C_0$, $TC_0$ de $C$, $TC$ respectivement.}
\end{array}\right.\\[5ex]
 & \quad\ \begin{array}{ll}
\text{c)} & TC \subset W,\ TF \overset{\text{déf}}{=} C_* \subset W.\\[1ex]
\text{d)} & \parbox[t]{0.45\textwidth}{On a l'une des trois conditions équivalentes supplémentaires de (2.4), i.e. :}\\[5ex]
 & \begin{array}{lrcllrcl}
\alpha) & TC & = & C \cap W\ , & \text{i.e.} & C \cap W & \subset & TC\ .\\
\beta) & TF & = & F \cap W\ , & \text{i.e.} & F \cap W & \subset & TF\ .\\
\gamma) & W & = & TF \circ TC\ . & & & &
\end{array}
\end{array}
\end{array}\right.
\tag*{(2.8)}
\]

\dgrblock{dgrb1169}{The extra condition in transferring a model structure}
\begin{dcomment}{exposition}{p15:de $\Fl(\M)$, satisfaisant les conditions}{The extra condition in transferring a model structure}
Before Proposition 2.1, it helps to identify what can go wrong. The adjunction transfers the two lifting relations formally, and small generation can provide both factorizations. Nevertheless, a map built from the transferred trivial generators might cease to be a weak equivalence after applying the right adjoint. The condition $TC'\subseteq W'$ rules out exactly that possibility.

Write $L=v_!$ and $R=v^*$. With
\[
 W'=R^{-1}W,\quad F'=R^{-1}F,\quad
 TF'=R^{-1}TF,
\]
the adjunction identifies
\[
 C'={}^{\square}TF',\qquad TC'={}^{\square}F'.
\]
Thus both pairs are already mutually orthogonal, although not yet known to factorize. Since $TC\subseteq C$, monotonicity gives $TC'\subseteq C'$, and inverse image gives $TF'\subseteq F'$. The class $W'$ satisfies two-out-of-three because $R$ is a functor. Once the two factorizations exist, the equation $TF'=F'\cap W'$ is built into the definitions; the previous block therefore reduces the entire model-structure question to $TC'\subseteq W'$.

The equivalent condition (2.1.5) has a useful practical interpretation: factor every target arrow into a weak equivalence followed by a member of $F'$, without insisting in advance that the first arrow is in $TC'$. If the arrow being factored already belongs to $TC'$, its left lifting property makes it a retract of that weak equivalence. Since $W'$ is retract-closed, it belongs to $W'$. Conversely, a $(TC',F')$-factorization works once $TC'\subseteq W'$ is known.

The convenient stronger condition is $RL(TC)\subseteq TC$. It tests the images of the original trivial left arrows before taking their left closure. Cases A and B show two mechanisms by which that condition passes to the closure. Case C, completed on manuscript page 19, instead transports a factorization of $Rf'$ through a fully faithful adjunction. This is why the three cases have different hypotheses. A right adjoint generally does not preserve pushouts, so the hypothesis of Case B must not be suppressed in a purported general transfer theorem.
\end{dcomment}


\printedpage{10}
% head: A. Grothendieck : Dérivateurs, Ch. XIX, Édition des Universités de Montpellier II et Paris VII
% foot: Version 11 février 2019

\textbf{Proposition 2.1.} \textit{Soit}
\[
\begin{tikzcd}
\M' \arrow[r, shift right, "v^*"'] & \M \arrow[l, shift right, "v_!"']
\end{tikzcd}
\tag*{(2.11)} % [sic] (2.11) -> (2.1.1), numbering sequence (2.1.2), (2.1.3), ...
\]
\textit{un couple de foncteurs adjoints, $\M$ étant muni d'une structure de Quillen close} $(W,C,F)$. \textit{Posons}
\[
\left\{
\begin{array}{rcl}
W' & = & (v^*)^{-1}(W)\\
F' & = & (v^*)^{-1}(F)\\
TF' & = & F' \cap W' \ = \ (v^*)^{-1}(TF)\\
C' & = & v_!(C)^{\sim}\\
TC' & = & v_!(TC)^{\sim}\ .
\end{array}\right.
\tag*{(2.1.2)}
\]

\begin{itemize}
\item[1°)] \textit{On a alors}
\[
\left\{
\begin{array}{rclcrcl}
C' & \longleftrightarrow & TF'\ , & & TC' & \longleftrightarrow & F'\ ,\\
TC' & \subset & C'\ , & & TF' & \subset & F'\ .
\end{array}\right.
\tag*{(2.1.3)}
\]

\textit{Pour que les données} (2.1.2) \textit{dans $\M'$ y correspondent à un triple de Quillen $(W',C',F')$, il faut et il suffit que les conditions suivantes soient satisfaites :}
\msp{page 16}
\[
\left\{
\begin{array}{ll}
\text{a)} & \parbox[t]{0.7\textwidth}{\textit{Factorisation pour $C' \longleftrightarrow TF'$ et $TC' \longleftrightarrow F'$, i.e. les relations $C' \Qcouple TF'$ et $TC' \Qcouple F'$,}}\\[4ex]
\text{b)} & TC' \subset W',
\end{array}\right. % reading resolved: DER_XIX_corrections.csv
\tag*{(2.1.4)}
\]
\textit{ou encore, au lieu de} b)\textit{, la condition équivalente} (en présence de a)) :
\[
\left\{
\begin{array}{ll}
\text{b$'$)} & \parbox[t]{0.7\textwidth}{\textit{Toute flèche $f'$ dans $\M'$ (non seulement $f' \in TC'$) se factorise en $f' = p'i'$, avec $i' \in W'$, $p' \in F'$.}}
\end{array}\right.
\tag*{(2.1.5)}
\]
\item[2°)] \textit{Supposons} (2.1.4. a) \textit{satisfaite (factorisation). Alors, pour que $(W',C',F')$ soit un triple de Quillen clos, il} suffit \textit{que l'on ait}
\[
\text{b$''$)}\qquad \boxed{\ v^*v_!(TC)\ \subset\ TC\ }\ ,
\tag*{(2.1.6)}
\]
\textit{et que de plus on soit dans l'un des trois cas suivants :}
\begin{itemize}
\item[A)] \textit{$v^*$ admet un adjoint à droite $v_*$.}
\item[B)] \textit{$v^*$ commute aux petites} $\varinjlim$\textit{, $(W,C,F)$ est spécial, et $\M'$ est spéciale, i.e. satisfait à la condition} (2.6. a).
\item[C)] \textit{$v^*$ est pleinement fidèle, et $v^*v_!(F) \subset F$.}
\end{itemize}
\end{itemize}
Notons tout de suite le

\textbf{Corollaire 2.2.} \textit{Supposons que $(W,C,F)$ soit un triple de Quillen spécial, que $\M'$ soit spéciale, que $v^*$ commute aux petites} $\varinjlim$\textit{, enfin que l'on ait} (2.1.6)\textit{, i.e. $v^*v_!(TC) \subset TC$. Alors $(W',C',F')$ est un triple de Quillen spécial.}

\printedpage{11}
% head: A. Grothendieck : Dérivateurs, Ch. XIX, Édition des Universités de Montpellier II et Paris VII
% foot: Version 11 février 2019

\msp{page 17}
En effet, en vertu du 2.1. 2°, cas B, pour voir que $(W',C',F')$ soit un triple de Quillen clos, il est à vérifier (2.1.6. b$''$). Mais cela résulte de 1.8, appliqué aux deux couples de Quillen spéciaux $(C,TF)$, $(TC,F)$ dans $\M$ : les couples $(C',TC')$ et $(TC',F')$ sont des couples de Quillen spéciaux, ce qui prouve en même temps que le triple $(W',C',F')$ est spécial, q.e.d. % [sic] (C',TC') -> (C',TF'), image of the couple (C,TF)

Démontrons 2.1. Les premières deux relations (2.1.3) sont contenues dans 1.6, les relations $TC' \subset C'$, $TF' \subset F'$ sont tautologiques. Le critère (2.1.4) pour que $(W',C',F')$ soit un triple de Quillen clos, est contenu dans le rappel de la page 13, compte tenu que l'on sait déjà que $TF' = F' \cap W'$ (relation (2.4. b)). Que la condition (2.1.4. b), en présence de (2.1.4. a), soit équivalente à la condition (2.1.5. b$'$), est contenu dans la première assertion (2.5), en notant que l'on a $W' = W'^{\natural}$ (ce qui résulte de la définition de $W'$ et de $W = W^{\natural}$).

Il reste à prouver 2°), i.e. que la condition b) est satisfaite moyennant a) et b$''$) (2.1.6) dans chacun des trois cas A), B), C).

\textit{Cas A.} Par définition de $W' = (v^*)^{-1}(W)$, pour prouver que $TC' \subset W'$, il suffit de prouver
\[
v^*(TC')\ \subset\ W'\ ,  % [sic] W' -> W, v^*(TC') lies in M and W' = (v^*)^{-1}(W)
\]
et a fortiori, il suffit de prouver
\[
v^*(TC')\ \subset\ TC\ ,\qquad\quad \text{i.e. } v^*(TC') \harp F\ .
\tag*{(2.1.7)} % reading resolved: DER_XIX_corrections.csv
\]
Or par 1.5, appliqué au couple de foncteurs
\msp{page 18}
adjoints $(v^*,v_*)$, cette relation équivaut à
\[
TC' \harp v_*(F)\ .
\]
Comme $TC' \overset{\text{déf}}{=} v_!(TC)^{\sim}$, ceci équivaut à
\[
v_!(TC) \harp v_*(F)\ ,
\]
ou encore (en réappliquant 1.5 à $(v_!(TC),F)$ au lieu de $(TC',F)$), à
\[
v^*v_!(TC) \harp F\ ,\qquad\quad \text{i.e. } v^*v_!(TC) \subset \underbrace{TC}_{=\,F^*}\ ,
\]
or c'est là l'hypothèse b$''$).

\textit{Cas B.} Il suffit encore de prouver 2.1.7. Or, du fait que $v_!(TC)$ admet un petit ensemble Q-générateur à gauche, $v_!(TC)^{\sim}$ est la clôture de $v_!(TC)$ pour les opérations de cochangement de base, de composition transfinie, de facteurs directs. Comme $v^*$ commute à toutes ces opérations (comme il commute aux petites $\varinjlim$), et comme $TC$ dans $\Fl(\M)$ est stable par ces opérations, on voit donc que pour montrer que $v^*$ applique $v_!(TC)^{\sim}$ dans $TC$, il suffit de voir qu'il applique $v_!(TC)$ dans $TC$, ce qui est encore l'hypothèse b$''$).

\dgrblock{dgrb1170}{Two ways of passing from generators to their closure}
\begin{dcomment}{exposition}{p18:Comme $v^*$ commute à toutes ces opérations}{Two ways of passing from generators to their closure}
The proofs of Cases A and B establish the stronger conclusion $R(TC')\subseteq TC$. They are worth comparing because one uses only lifting properties, whereas the other uses the cellular description supplied by small generation.

In Case A, $R$ has a right adjoint $V=v_*$. For each $p\in F$, the condition $RL(TC)\perp p$ is equivalent to $L(TC)\perp Vp$. Taking the left double-orthogonal closure of $L(TC)$ does not change which arrows have the right lifting property against it; hence $TC'\perp Vp$. Applying the adjunction once more gives $R(TC')\perp p$. Since $TC={}^{\square}F$, this proves the desired containment. This argument does not require $R$ to preserve any colimit and does not require the ambient categories to be locally presentable.

In Case B, the small object argument identifies $TC'$ with the retracts of cellular composites of a small family inside $L(TC)$. The functor $R$ preserves pushouts and transfinite composites by the stated colimit hypothesis; every functor preserves retract diagrams. Since $TC$ is closed under these operations and $RL(TC)\subseteq TC$, it follows that $R(TC')\subseteq TC$. One may equally use the full weakly saturated closure of $L(TC)$: Corollaries 1.8--1.9 show that its small generating family has the same left closure.

Corollary 2.2 combines this argument with Corollary 1.8, which gives both target factorizations and their small generators. Its printed proof contains two local slips. The first transferred pair is $(C',TF')$, not $(C',TC')$. Moreover, the remaining condition supplied by Corollary 1.8 is factorization (2.1.4.a), not (2.1.6.b$''$), which is already an explicit hypothesis of Corollary 2.2. In the subsequent Case A proof, $R(TC')$ is a class in $\mathcal M$, so the intermediate target must be $W$, not $W'$. With these corrections of interpretation, the proof establishes the corollary exactly as stated.
\end{dcomment}


\textit{Cas C.} On va vérifier b) sous la forme équivalente b$'$). Soit donc $f' \in \Fl(\M')$, $f' : X' \longrightarrow Y'$,
\printedpage{12}
% head: A. Grothendieck : Dérivateurs, Ch. XIX, Édition des Universités de Montpellier II et Paris VII
% foot: Version 11 février 2019
\msp{page 19}
posons
\[
f\ =\ v^*(f'),
\]
et factorisons $f$ en $pi$, avec $i \in TC$, $p \in F$, d'où
\[
v_!v^*(f') = v_!(f) = v_!(p)v_!(i)\ .
\]
Mais $v^*$ étant pleinement fidèle, $v_!v^* \xrightarrow{\ \sim\ } \mathrm{id}$, donc on trouve une factorisation de $f'$ en $v_!(p)v_!(i)$. Il suffit donc de voir que
\[
v_!(p)\ \in\ F'\ ,\qquad\quad v_!(i)\ \in\ W'\ ,
\]
ce qui équivaut à
\[
v^*v_!(p) \in F\ ,\qquad\quad v^*v_!(i) \in W\ ,
\]
et résulte des hypothèses que $v^*v_!(F) \subset F$ et $v^*v_!(TC) \subset TC \subset W$, q.e.d.

\textbf{N.B.} J'ai noté les cas A) et C) par acquis de conscience, sans application précise en vue. Le cas le plus utile sera justiciable du cas B), et mérite d'être réexplicité en un

\dgrblock{dgrb1171}{Free diagrams and the projective model structure}
\begin{dcomment}{exposition}{p19:J'ai noté les cas A) et C) par acquis de conscience}{Free diagrams and the projective model structure}
The concluding Case~C can first be made completely explicit. Write $L=v_!$, $R=v^*$ and $\varepsilon:LR\to1$ for the counit. Since $R$ is fully faithful, $\varepsilon$ is an isomorphism. For $f':X'\to Y'$, choose in $\mathcal M$ a factorization
\[
 R X'\xrightarrow{i}Z\xrightarrow{p}R Y',
 \qquad i\in TC,\quad p\in F.
\]
Define
\[
 i'=L(i)\varepsilon_{X'}^{-1}:X'\longrightarrow LZ,
 \qquad p'=\varepsilon_{Y'}L(p):LZ\longrightarrow Y'.
\]
Naturality of the counit gives $p'i'=f'$. Moreover $R(i')=RL(i)R(\varepsilon_{X'}^{-1})$ lies in $TC$, since $RL(TC)\subseteq TC$ and $TC$ is stable under composition with isomorphisms. Likewise $R(p')=R(\varepsilon_{Y'})RL(p)$ lies in $F$, using the other hypothesis $RL(F)\subseteq F$. Thus $i'\in W'=R^{-1}W$ and $p'\in F'=R^{-1}F$, proving the factorization condition (2.1.5). The two factorizing pairs required separately in Proposition~2.1 remain part of its hypotheses; this argument supplies their compatibility with weak equivalences.

Grothendieck now puts the transfer argument to its principal use: a homotopy theory of coefficients should supply a homotopy theory of diagrams of every small shape. The weak equivalences must be tested at the objects of the indexing category, since this is what allows a diagram to retain its homotopy coherence while its individual values are replaced. The freedom lies in choosing the auxiliary cofibrations and fibrations. For homotopy left Kan extensions, the useful choice fixes the fibrations componentwise and lets lifting determine the cofibrations.

Write $\mathcal M(I)=\operatorname{Fun}(I^{\mathrm{op}},\mathcal M)$, as in the source, and use $L^{\square}$ and ${}^{\square}R$ for right and left lifting classes. The intended formulas are
\[
 C_{\mathrm{proj}}(I)={}^{\square}(F(I)\cap W(I)),
 \qquad
 TC_{\mathrm{proj}}(I)={}^{\square}F(I).
\]
The first equality in Theorem 2.3 as printed has $F(I)$ where $TF(I)$ is required. This is visible in the edition and is therefore retained in the transcription. Its correction follows already from (1.16), applied to $(C,TF)$. With the printed reading, the alleged cofibrations would be precisely the trivial cofibrations, and even the case $I=e$ would give the wrong class.

The free-diagram adjunction makes the construction concrete. For $i\in I$, evaluation at $i$ has a left adjoint
\[
 (i_!A)(k)=\coprod_{I(k,i)}A.
\]
A lifting problem against $i_!a$ is exactly a lifting problem against $a$ at $i$. If $S$ and $T$ generate the two lifting pairs in $\mathcal M$, the small families
\[
 \{i_!s\mid i\in\operatorname{Ob}I,\ s\in S\},
 \qquad
 \{i_!t\mid i\in\operatorname{Ob}I,\ t\in T\}
\]
therefore detect the componentwise trivial fibrations and fibrations. The required smallness is inherited because evaluation preserves the colimits used in the small object argument. Thus a domain small for a sufficiently long cellular construction in $\mathcal M$ remains small after applying $i_!$. The common bound may depend on $I$ and on the generating families; no bound independent of every small category is needed.

To see why the second family produces weak equivalences, evaluate a cell attachment at $k$. It becomes a coproduct of pushouts of maps in $T$, indexed by the sets $I(k,i)$. Trivial cofibrations in $\mathcal M$ are closed under coproducts, pushouts, transfinite composites and retracts. The evaluated relative cell map is consequently a trivial cofibration. The transfer criterion applies. This is the projective model structure of present terminology; the modern combinatorial statement appears in \citet[Theorem~1.28]{Groth2013}. The word ``special'' in this chapter includes small generating families and the stated accessibility of objects, but it does not by itself assert that the underlying category is locally presentable.

The rest of Theorem 2.3 is formal once these factorizations exist, but one step deserves to be made explicit. Suppose $L$ preserves cofibrations and trivial cofibrations and that $f:X\to Y$ is a weak equivalence between cofibrant objects. Factor
\[
 (f,1_Y):X\amalg Y\xrightarrow{k}Z\xrightarrow{p}Y
\]
with $k$ a cofibration and $p$ a trivial fibration. Let $i=k|_X$ and $j=k|_Y$. Both are cofibrations because $X$ and $Y$ are cofibrant. Since $pi=f$ and $pj=1_Y$, two-out-of-three makes $i$ and $j$ trivial cofibrations. Now $L(p)L(j)=1$, while $L(j)$ is weak; hence $L(p)$ is weak, and so is $L(f)=L(p)L(i)$. This proves exactly the homotopy invariance needed to define $u_!^W$ on cofibrant diagrams. It also explains why taking the ordinary Kan extension of an arbitrary representative is insufficient.
\end{dcomment}


\textbf{Théorème 2.3}\footnote[7]{cf. scholie 2.10, qui devrait précéder le théorème 2.3.}. \textit{Soit $(W,C,F)$ un triple de Quillen spécial dans $\M$. Pour toute $I$ dans} Cat\textit{, considérons le triple $(W(I),C\langle I\rangle,F(I))$ dans $\M(I) = \underline{\mathrm{Hom}}(I^{\circ},\M)$, où les notations sont celles de page 10, donc $W(I)$, $F(I)$ et $TF(I)$ se définissent fibre par fibre, et $C\langle I\rangle = (F(I))^* = (v_{I\,!}(C(I_0)))^{\sim}$}\footnote[8]{où $v_I : I_0 \longrightarrow I$ ($I_0$ catégorie discrète ayant mêmes objets que I).}\textit{.}
\begin{itemize}
\item[1°)] \textit{Pour toute $I$ dans} Cat\textit{, $(W(I),C\langle I\rangle,F(I))$ est}
\msp{page 20}
\textit{un triple de Quillen spécial dans $\M(I)$.}
\item[2°)] \textit{Si $u : I \longrightarrow I'$ est une flèche dans} Cat\textit{, d'où}
\[
\begin{tikzcd}
\M(I) \arrow[r, shift left, "u_!"] & \M(I'), \arrow[l, shift left, "u^*"]
\end{tikzcd}
\tag*{(2.3.1)}
\]
\textit{alors}
\[
\left\{
\begin{array}{l}
\left.\begin{array}{rcl}
u^*(W(I')) & \subset & W(I)\\
u^*(F(I')) & \subset & F(I)\\
u^*(TF(I')) & \subset & TF(I)
\end{array}\right\}\textit{tautologie}\\
\begin{array}{rcl}
u_!(C\langle I\rangle) & \subset & C\langle I'\rangle\\
u_!(TC\langle I\rangle) & \subset & TC\langle I'\rangle\ ,
\end{array}
\end{array}\right.
\tag*{(2.3.2)}
\]
\textit{enfin, $u_!$ induit un foncteur}
\[
\M(I)_{\cc} \longrightarrow \M(I')_{\cc} \qquad\quad \text{\footnotemark[9]}
\tag*{(2.3.3)}
\]
\footnotetext[9]{Sous-catégories des objets cofibrants.}
compatible aux localiseurs $W(I)_{\cc} = W(I)|\M(I)_{\cc}$ \textit{et} $W(I')_{\cc} = W(I')|\M(I')_{\cc}$\textit{, donc induit un foncteur}
\[
u_!^W\ :\ \Ho_W(I) \longrightarrow \Ho_W(I')\ .
\tag*{(2.3.4)}
\]
\printedpage{13}
% head: A. Grothendieck : Dérivateurs, Ch. XIX, Édition des Universités de Montpellier II et Paris VII
% foot: Version 11 février 2019
\textit{Ce foncteur est un adjoint à gauche du foncteur $u^*_W$,}
\[
u^*_W\ :\ \Ho_W(I') \longrightarrow \Ho_W(I)\ ,
\tag*{(2.3.5)}
\]
\textit{induit par $u^*$} (3.3.1). % [sic] (3.3.1) -> (2.3.1), the adjunction displayed in (2.3.1)
\end{itemize}
\textsc{Démonstration de 2.3.}

\textit{Prouvons} 1°). On considère comme dans 1.10 le foncteur
\[
v = v_I\ :\ I_0 \longrightarrow I
\]
et le couple de foncteurs adjoints
\[
\begin{tikzcd}
\M(I) \arrow[r, shift right, "v^*"'] & \M(I_0)\ , \arrow[l, shift right, "v_!"']
\end{tikzcd}
\tag*{(2.3.6)}
\]
auquel on peut appliquer 2.2. Il reste
\msp{page 21}
à vérifier seulement
\[
v^*v_!(TC)\ \subset\ TC\ ,
\tag*{(2.3.7)}
\]
or on a, pour $X = (X_i)_{i \in I_0} \in \M(I_0)$,
\[
\begin{aligned}
v_!(X)(i) &\ \simeq\ \textstyle\coprod_{j \in I_0} v_{j\,!}(X_j)(i)\\
&\ \simeq\ \textstyle\coprod_{j \in I_0}\ \ \underbrace{X_j^{\mathrm{Hom}(i,j)}}_{\substack{\text{somme directe}\\ \text{de }\mathrm{Hom}(i,j)\text{ copies}}}\qquad ,
\end{aligned}
\]
et la relation (2.3.7) résulte du fait que $TC$, étant Q-clos à gauche, est stable par petites sommes directes.

\dgrblock{dgrb1172}{A constant diagram whose homotopy colimit is a classifying space}
\begin{dcomment}{examples}{p21:et la relation (2.3.7) résulte}{A constant diagram whose homotopy colimit is a classifying space}
The free objects in the preceding proof have a geometric meaning: they introduce the room needed to carry the relations in a diagram through a replacement. A diagram may have cofibrant values and still fail to be cofibrant as a diagram. Group actions give a complete calculation of the difference.

Let $G$ be a nontrivial group and let $I=BG$ be the category with one object and automorphism group $G$. A contravariant diagram of simplicial sets on $BG$ is a simplicial set with a right $G$-action. Give simplicial sets their usual Kan model structure and give these action diagrams the projective structure. The ordinary colimit is the orbit simplicial set. In particular the colimit of the point with trivial action is a point. We calculate its derived colimit directly.

Let $\mathcal E_G$ be the category whose objects are the elements of $G$ and which has exactly one morphism between any ordered pair of objects. Right translation gives a free action on this category. Its nerve $EG$ has
\[
 (EG)_n=G^{n+1},
\]
with faces deleting an entry and degeneracies repeating an entry; $G$ acts diagonally on the right. The category $\mathcal E_G$ has an initial object, indeed every object is initial, so its nerve is contractible. An explicit contraction is obtained by choosing the identity element and inserting it as an extra first vertex. The augmentation $EG\to *$ is therefore a weak equivalence of the underlying simplicial sets. Since $\mathcal E_G$ is a groupoid, its nerve is Kan: a missing face in a horn is supplied by composition and inversion. Thus the augmentation is also a fibration.

It is also a projective cofibrant replacement. To verify cofibrancy without assuming an equivariant model theorem, filter $EG$ by its skeleta. The action on its nondegenerate simplices is free: an element fixing the tuple $(g_0,\ldots,g_n)$ must fix $g_0$, hence is the identity. After choosing one representative in every orbit of nondegenerate $n$-simplices, the $n$-skeleton is obtained from the preceding skeleton by attaching copies of
\[
 G\times\partial\Delta^n\longrightarrow G\times\Delta^n.
\]
These are the free diagrams on the generating cofibrations of simplicial sets. Starting from the empty diagram proves that $EG$ is projectively cofibrant.

Its orbit simplicial set is the ordinary bar construction $BG$. There is one vertex, an edge for every group element, and a two-simplex for every composable pair. Its fundamental group has generators $[g]$, $g\in G$, with relations
\[
 [1]=1,\qquad [g][h]=[gh],
\]
so the evident map identifies it with $G$. Thus the derived colimit of the constant point is
\[
 \operatorname{hocolim}_{BG^{\mathrm{op}}}*\simeq EG/G\simeq BG,
\]
which cannot be contractible for nontrivial $G$. For $G=C_2$ the first homology group is $\mathbb Z/2$, whereas the first homology of the ordinary orbit point vanishes. The calculation also proves that the trivial-action point was not projectively cofibrant: otherwise the weak equivalence $EG\to *$ between cofibrant diagrams would remain a weak equivalence after taking orbits, contrary to the calculation.

Every individual simplicial set in this example is cofibrant. What the replacement corrects is therefore the action, including its coherent relations, rather than the values of the diagram. This is the concrete content of the angle brackets in $C\langle I\rangle$. The same distinction will reappear with arrows reversed when injectively fibrant diagrams are used to compute right Kan extensions. The general diagram conventions and their Kan adjunctions are recorded in \citet[Definitions~A.2.8.1 and Proposition~A.2.8.7]{Lurie2009HTT}; the calculation above gives the elementary free resolution behind this instance.
\end{dcomment}


\textit{Prouvons} 2°). Les formules (2.3.2) sont essentiellement tautologiques, les deux dernières résultent de la troisième et de la deuxième, en vertu de 1.5 appliqué au couple de foncteurs adjoints $(u_!,u^*)$. La relation $u_!(C\langle I\rangle) \subset C\langle I'\rangle$, et le fait que $u_!$ transforme objet initial en objet initial, impliquent (2.3.3). Que $u_!$ restreint à $\M(I)_{\cc}$ soit compatible aux localiseurs provient de $u_!(TC\langle I\rangle) \subset TC\langle I'\rangle$, et du fait que toute flèche $f$ dans $\M(I)_{\cc}$ qui est dans $W(I)$, se décompose en $f = pi$, avec $i \in TC$ et $p$ inverse à gauche d'une $j \in TC\langle I\rangle$. Ceci vu, joint aux deux premières formules (2.3.2) montre qu'on est sous les conditions
\msp{page 22}
du théorème d'adjonction de Quillen, d'où l'adjonction de $u_!^W$ et de $u^*_W$, q.e.d.

\textbf{Corollaire 2.4.} \textit{Soit $W$ un localiseur dans $\M$, tel qu'il existe un triple de Quillen spécial $(W,C,F)$ incluant $W$. Alors le prédérivateur} $\Ho_W$ \textit{sur} Cat \textit{satisfait aux axiomes Dér 1 (décomposition), Dér 2 (localisation), Dér 3 (objet final), Der 4 bis (existence des $u_!^W$), Dér 5 bis (calcul des fibres des $u_!^W$), Dér 6 bis (exactitude à gauche, i.e. exactitude à gauche des carrés $W$-cocartésiens dans les} $\Ho_W(I)$\textit{), Dér 6 (exactitude à droite des carrés $W$-cartésiens).} % [sic] Der 4 bis -> Dér 4 bis, accent missing as printed

On vient de prouver Dér 4 bis, et Dér 5 bis est pratiquement contenu dans le mode d'obtenir de $u_!$, compte tenu de proposition 1.12, qui implique que pour $X$ dans $\M(I)_{\cc}$, et $i'$ dans $I'$, l'image inverse de $X$ dans $\M(i'\backslash I)$ est dans $\M(i'\backslash I)_{\cc}$ \dots

% ---- from der-ch19_p016-020.tex
\printedpage{14}
% head: A. Grothendieck : Dérivateurs, Ch. XIX, Édition des Universités de Montpellier II et Paris VII
% foot: Version 11 février 2019

\msp{page 23}
L'axiome de décomposition est satisfait pour tout prédérivateur de la forme $\mathbf{D}_{\mathcal{M},W}$. L'axiome d'exactitude à gauche Dér 5 bis est valable, pourvu que les $(\mathcal{M}(I),W(I))$ s'insèrent dans une structure à cofibrations, ce qui est le cas. De même pour l'axiome Dér 2 de localisation, pourvu que de plus $W$ soit fortement saturé, ce qui est également le cas. Enfin, l'axiome de l'objet final dit que si $I$ admet un objet final, alors
\[ \Ho_W(e) \longrightarrow \Ho_W(I) \]
est pleinement fidèle, ou encore que le morphisme d'adjonction
\[ u_!u^*(\xi) \longrightarrow \xi \]
est un isomorphisme. Mais en vertu de 1.12, cas 3°, dans ce cas $u^*$ applique $\mathcal{M}(e)_{\mathrm{c}} = \mathcal{M}_{\mathrm{c}}$ dans $\mathcal{M}(I)_{\mathrm{c}}$, donc le calcul du couple de foncteurs adjoints $u_!^W$, $u_W^*$ peut se faire en utilisant les restrictions de $u_!$, $u^*$ à $\mathcal{M}(I)_{\mathrm{c}}$ et $\mathcal{M}_{\mathrm{c}}$ respectivement. Comme le morphisme d'adjonction $u_!^{\mathrm{c}}u_{\mathrm{c}}^* \longrightarrow \mathrm{id}$ est un isomorphisme, il en est donc de même de $u_!^W u_W^* \longrightarrow \mathrm{id}$, qui s'en déduit en localisant, q.e.d.

\dgrblock{dgrb1173}{The comma formula and the first derivator axioms}
\begin{dcomment}{exposition}{p23:Comme le morphisme d'adjonction}{The comma formula and the first derivator axioms}
Grothendieck wants more than an adjoint on each localized diagram category. He wants its values to be calculable from the part of the diagram lying over a chosen object. This local calculation is what makes the construction a derivator and permits the subsequent base-change arguments. The opposites in his convention must be retained throughout the calculation.

Let $u:I\to J$, $j\in J$, and let $q:j\downarrow u\to I$ forget the arrow $j\to u(i)$. For a contravariant diagram $X$, the ordinary formula is
\[
 (u_!X)(j)=\operatorname{colim}_{(j\downarrow u)^{\mathrm{op}}}q^*X.
\]
To derive this equality, one must know that the same representatives compute both sides. Put $E=j\downarrow u$. The functor $q:E\to I$ is a discrete opfibration: an arrow $i\to i'$ carries $(i,a:j\to u(i))$ to $(i',u(i\to i')a)$, uniquely. For every $i\in I$ and object $A\in\mathcal M$, this unique-lifting property gives
\[
 q^*(i_!A)\cong\coprod_{e\in q^{-1}(i)} e_!A.
\]
Indeed, at $e'\in E$, the indexing sets on the two sides are identified by the bijection
\[
 \coprod_{e\in q^{-1}(i)} E(e',e)
       \xrightarrow{\ \sim\ } I(q(e'),i).
\]
Each arrow on the right has exactly one lift with source $e'$, whose target is in the indicated fibre. Uniqueness also makes the bijection compatible with precomposition, so it defines an isomorphism of diagrams, natural in $A$. The same formula holds for arrows $a$ in place of objects $A$.

Apply it to the two generating families in the preceding block. Restriction sends each generating projective cofibration to a coproduct of generating projective cofibrations in $\mathcal M(E)$, and likewise for trivial cofibrations. It preserves colimits, computed pointwise, and every functor preserves retract diagrams. It consequently preserves the cellular and retract closures of both families, hence all projective cofibrations and trivial cofibrations. It also preserves the initial object, so $q^*QX$ is projectively cofibrant whenever $QX$ is. This proof uses precisely the cocompleteness and generation hypotheses already available here.

Apply the ordinary formula to a projectively cofibrant replacement $QX\to X$. On the left it computes $u_!^WX$; on the right it computes the homotopy colimit of $q^*X$. Hence
\[
 (u_!^WX)_j\cong
 \operatorname{hocolim}_{(j\downarrow u)^{\mathrm{op}}}q^*X.
\]
The isomorphism is the canonical comparison obtained from the comma square, since the ordinary comparison was used before localization. This last point matters: an abstract isomorphism between the two objects would not establish the base-change axiom.

The argument for a terminal object in $I$ is similarly precise. If $t$ is terminal, then $I(i,t)$ is a singleton for each $i$, so $t_!A$ is the constant diagram on $A$. Thus the constant-diagram functor takes cofibrant objects to projectively cofibrant diagrams. Since $I$ is nonempty and connected, the colimit of a constant diagram is $A$, and the counit is an isomorphism already before localization. Its derived counit is consequently an isomorphism. This establishes full faithfulness of constant coefficients in the form claimed by Dér~3.

There is one overstatement in the paragraph being commented on. Decomposition for arbitrary small disjoint unions does not follow for an entirely arbitrary relative category merely from the formula $\mathbf D_{\mathcal M,W}(I)=\mathcal M(I)[W(I)^{-1}]$. Localization need not commute with infinite products: a single morphism in a product localization has a finite zigzag presentation whose length is uniform over all components. The commentary on Chapter~VI, manuscript p.~131, addresses that obstruction, and should be consulted here. In the present model-category situation the conclusion is valid. Componentwise replacements and the description of maps between cofibrant-fibrant objects by homotopy classes supply uniform representatives and uniform homotopies. The same description proves that evaluation detects isomorphisms: a representative with weak-equivalence components is a weak equivalence in the diagram model structure.

Finally, the two exactness axioms concern finite homotopy pushouts and pullbacks. They do not require that arbitrary right Kan extensions have already been constructed. The cofibration and fibration structures provide the required finite calculus. The reference to ``Dér~5 bis'' for left exactness on manuscript page~23 should read ``Dér~6 bis''; the statement on page~22 has the right numbering. The author supplies the omitted right-exactness argument himself on page~24. These are local repairs in the exposition, not failures of Corollary~2.4.
\end{dcomment}


\textbf{2.5.} Pour achever de construire $\Ho_W$ comme un dérivateur sur Cat tout entier, il faut \msp{page 24} donc établir Dér 4, Dér 5 : existence des $u_*^W$, et calcul des fibres de $u_*^W$. Quant à Dér 6, il est acquis automatiquement, du fait que chaque $(\mathcal{M}(I),W(I))$ s'insère dans une structure de fibration $(W(I),F(I))$. (Oubli dans la démonstration de 2.4.) L'idée naturelle est d'utiliser à présent la structure de triple
\[ (W(I),C(I),\ \underbrace{F[I]}_{\deq TC(I)_*}\ )\ , \tag*{(2.5.1)} \]
mais la difficulté est de prouver la factorisation pour un couple
\[ C(I) \longleftrightarrow TF[I]\ , \qquad \underbrace{TC(I)}_{\deq C(I)\cap W(I)} \quad \longleftrightarrow F[I]\ , \tag*{(2.5.2)} \]
faute de pouvoir montrer, dans les conditions générales de 2.3 seulement (triple de Quillen spécial) que $C(I)$ et $TC(I)$ sont également spéciaux. (Cf. question 1.11.) Mais moyennant des conditions convenables sur $\mathcal{M}$, on arrivera à s'en tirer si $C = \mathit{mono}$. L'énoncé pertinent, un peu technique, se trouve dans XV cor. 7.4 (p. 98-99), que je vais rappeler ici.

\dgrblock{dgrb1174}{Why monomorphisms enter the construction}
\begin{dcomment}{context}{p24:faute de pouvoir montrer}{Why monomorphisms enter the construction}
The obstruction to right Kan extensions is now a problem about generating cofibrations. The source already knows how to generate the projective classes by free diagrams. The injective classes are prescribed componentwise, and a small generating family for them is not produced by simply taking the old generators componentwise. Grothendieck returns to monomorphisms because subobjects provide a way to control the size of successive attachments. This is the same geometry of unions and injective objects that underlies the older theory of sheaves and abelian categories.

The six conditions in Theorem~2.6 separate this size control from the exactness needed to glue subobjects. Cocompleteness and accessibility of individual objects supply the small object argument. The cardinal filtration adds a supply of objects bounded in size; it must not be silently absorbed into the word ``special.'' Left exactness of filtered colimits allows intersections to pass through increasing unions. Condition~(e) says that the union of two subobjects can be formed by their pushout over the intersection. Condition~(f) ensures that attaching such a piece preserves monicity even when the attaching map is not itself monic. These hypotheses explain why the subsequent proof can remain inside the category of subobjects of a fixed target.

The direction of the limit in (2.6.1d) requires an editorial correction. The edition prints a projective-limit arrow. The cited antecedent, XV, Theorem~6.1(d), manuscript page~76, has inductive limits, and XVIII, Proposition~4.6.7, manuscript page~157, likewise gives closure of the cardinal stages under inductive limits of bounded size \citep{Grothendieck1990XV,Grothendieck1990XVIII}. This is also what the intended examples require. For the ordinary filtration of sets by cardinality at most $\pi$, a coproduct of at most $\pi$ sets of that size still has size at most $\pi$. By contrast, the product of $\pi$ copies of a two-element set has size $2^\pi>\pi$. We preserve the printed sign and use the intended inductive-limit condition when explaining the theorem.

The examples in the footnotes give a useful test of the exactness assumptions. In an abelian category the pushout of subobjects $A_1,A_2\subset B$ over their intersection is their sum $A_1+A_2$, which embeds in $B$. If $i:A_0\hookrightarrow A_1$ is monic and $f:A_0\to A_2$ is arbitrary, the pushout is
\[
 (A_1\oplus A_2)/\operatorname{im}(i,-f).
\]
The map from $A_2$ into this quotient is monic: an element of the intersection of $0\oplus A_2$ with $\operatorname{im}(i,-f)$ must come from an element killed by $i$, hence from zero. The same proof can be phrased with kernels and images in any abelian category. Thus these conditions extend beyond topoi.

For groups the effective-union condition fails visibly. Inside $C_2\times C_2$, take the two coordinate subgroups. Their intersection is trivial, while their pushout in groups is $C_2*C_2$. The canonical map to $C_2\times C_2$ kills the commutator of the two nonidentity generators. That commutator is a nonempty reduced word in the free product, so the map is not injective. The appeal to the absence of enough injectives in groups is therefore not an arbitrary warning: the precise gluing mechanism used by the theorem has already broken. The cited SGA passage establishes that there are groups which embed in no injective group \citep[Exposé~I, Example~8.12.9]{SGA4I1972}.

The tentative footnote about removing condition~(d) should be read alongside Grothendieck's own refinement in XVIII, 4.6.7. There he replaces the excessive demand on all quotients by factorization into an effective epimorphism followed by a monomorphism, recovering the weaker quotient condition actually used in XV. This is a specified replacement of hypotheses; it does not justify the projective-limit sign printed here, nor does it establish a small generating set for every accessible class of monomorphisms.
\end{dcomment}


\printedpage{15}
% head: A. Grothendieck : Dérivateurs, Ch. XIX, Édition des Universités de Montpellier II et Paris VII
% foot: Version 11 février 2019

\msp{page 25}
\textbf{Théorème 2.6.} \textit{Soit $\mathcal{M}$ une catégorie satisfaisant les conditions suivantes.}
\[ \left\{ \begin{minipage}{0.72\textwidth}\itshape
\begin{itemize}[leftmargin=2.5em]
\item[\upshape a)] $\mathcal{M}$ est “spéciale”, i.e. stable par petites $\varinjlim$ , et ses éléments sont accessibles.
\item[\upshape b)] Les limites inductives filtrantes dans $\mathcal{M}$ sont exactes à gauche (et dans $\mathcal{M}$ les $\varprojlim$ finies sont représentables – résulte de \textup{a} et \textup{c}).
\item[\upshape c)] $\mathcal{M}$ admet une filtration cardinale $(\Filt^\pi)_{\pi\geq\pi_0}$.
\item[\upshape d)] Les $\Filt^\pi$ sont stables par $\varprojlim$ indexées par des catégories $I$ telles que $\operatorname{card}\Fl(I) \leq \pi$, et par quotients ($Z \in \Filt^\pi$ et $Z \longrightarrow Z'$ épimorphisme implique $Z' \in \Filt^\pi$).
\item[\upshape e)] Si $A_1$, $A_2$ sont deux sous-objets d'un objet $B$ de $\mathcal{M}$, et $A_0 = A_1 \cap A_2$, alors $A_1 \amalg_{A_0} A_2 \longrightarrow B$ est un monomorphisme.
\item[\upshape f)] Si on a un carré cocartésien
\begin{center}
\begin{tikzcd}
A_0 \arrow[r, "i"] \arrow[d] & A_1 \arrow[d] \\
A_2 \arrow[r, "i'"] & A_3
\end{tikzcd}
\end{center}
avec $i$ monomorphisme, alors $i'$ est un monomorphisme.
\end{itemize}
\end{minipage}\right. \tag*{(2.6.1)} \]
$^{(10,11,12,13)}$\footnotetext[10]{\textbf{N.B.} a + c signifie aussi que $\mathcal{M}$ est une “catégorie $\mathfrak{U}$-paratopos”, et implique l'existence d'une filtration cardinale satisfaisant d) (sauf l'assertion concernant les quotients $Z'$ de $Z$). Donc il s'agit d'un $\mathfrak{U}$-paratopos satisfaisant de plus b) e) f). (Est-ce alors un topos\,? \textit{Non,} car e), f) sont O.K. dans les \textit{catégories abéliennes\,!})}\footnotetext[11]{\textbf{N.B.} La condition d) est sans doute inutile, du moins si les épimorphismes sont effectifs, car d) sera satisfaite en tous cas pour $\pi = 2^c$, avec $c \geq \pi_0$.}\footnotetext[12]{\textbf{N.B.} La condition e) n'est \textit{pas} satisfaite pour la catégorie (Gr), et le théorème est en défaut pour cette catégorie, laquelle n'admet \textit{pas} “assez d'injectifs”, cf. SGA 4 I 8.12.9. Par contre, e) et f) sont satisfaites par les catégories \textit{abéliennes.}}\footnotetext[13]{Cf. XVIII 4.6.7 (p. 157) pour une formulation plus sympathique de l'ensemble des conditions a) à f) pour $\mathcal{M}$.}. \textit{(On dira alors que $\mathcal{M}$ est} “très spéciale”\textit{). Soit de plus}
\[ W \ \subset \Fl(\mathcal{M}) \tag*{(2.6.2)} \]
\textit{un localiseur dans $\mathcal{M}$, qui satisfait les conditions}
\[ \left\{ \begin{minipage}{0.72\textwidth}\itshape
\begin{itemize}[leftmargin=2.5em]
\item[\upshape a)] $W$ est stable par $\varinjlim$ filtrantes $^{\textup{(14)}}$.
\item[\upshape b)] $W$ est de plus $\mathrm{L}_1$-accessible, i.e. il existe une \textup{petite} partie $W_0 \subset W$ telle que $W$ soit la $\mathrm{L}_1$-enveloppe de $W$, i.e. que $W$ soit la plus petite partie de $\Fl(\mathcal{M})$ contenant $W_0$, et stable par $\varinjlim$ filtrantes. % [sic] la L_1-enveloppe de W -> la L_1-enveloppe de W_0, W_0 being the small generating part just introduced
\end{itemize}
\end{minipage}\right. \tag*{(2.6.3)} \]
\footnotetext[14]{\textbf{N.B.} a) implique déjà que $W = W^\natural$.}

\printedpage{16}
% head: A. Grothendieck : Dérivateurs, Ch. XIX, Édition des Universités de Montpellier II et Paris VII
% foot: Version 11 février 2019

\msp{page 26}
\textit{Posons alors}
\[ C \ = \ \Mono(\mathcal{M}), \qquad TC \ = \ C\cap W\ , \qquad TF \ = \ C_*\ , \qquad F = (TC)_*\ . \tag*{(2.6.4)} \]
\textit{Sous ces conditions :}
\begin{itemize}[leftmargin=2.5em]
\item[1°)] \textit{$C \Qcouple TF$, plus précisément $C \qpair TF$ est un couple de Quillen spécial.} % reading resolved: DER_XIX_corrections.csv
\item[2°)] \textit{Pour que $TC \qpair F$ soit un couple de Quillen, il faut et il suffit que $TC$ soit stable par cochangement de base, et alors c'est un couple de Quillen spécial.} % reading resolved: DER_XIX_corrections.csv
\item[3°)] \textit{Pour que $(W,C,F)$ soit un triple de Quillen (lequel sera nécessairement spécial par \textup{1°, 2°}, et en particulier clos), il faut et il suffit que les conditions suivantes soient satisfaites.}
\[ \left\{ \begin{array}{ll}
\text{a)} & TC\ (\deq W\cap\Mono)\ \text{\itshape est stable par cochangement de base.}\\[1ex]
\text{b)} & TF\ (\deq (\Mono)_*) \ = \ \text{\itshape “applications de type injectif”}\ \subset W.
\end{array}\right. \tag*{(2.6.5)~$^{(15)}$} \]
\footnotetext[15]{\textbf{N.B.} Ces conditions (2.6.5) équivalent à dire que $(W,C)$ est une structure de catégorie à cofibrations sur $\mathcal{M}$.}
\textit{Et la condition} b) \textit{équivaut encore à :}
\[ \begin{minipage}{0.7\textwidth}
\begin{itemize}[leftmargin=2.5em]
\item[b$'$)] \textit{Toute flèche $f$ de $\mathcal{M}$ (ou seulement de $TF$) se factorise en $f = pi$, avec $i$ monomorphisme, $p \in W$.}
\end{itemize}
\end{minipage} \tag*{(2.6.6)} \]
\end{itemize}

\textbf{Corollaire 2.7.} \textit{$\mathcal{M}$ et $W$ étant comme dans} (2.6.1), (2.6.3), \textit{pour que $(W,C,F)$ soit un triple de Quillen} propre à gauche \textit{(i.e. $C \subset \Cof_W$), il faut et il suffit que $W$ satisfasse la condition} (2.6.5. b) \textit{(i.e. $W$ contient les flèches de type injectif), et de plus}
\msp{page 27}

\begin{itemize}[leftmargin=7em]
\item[] $C \subset \Cof_W$, \textit{i.e. la condition :}
\item[a$'$)] \textit{Pour tout diagramme cocartésien}
\[ \begin{tikzcd}
A_0 \arrow[r, "i"] \arrow[d, "f"'] & A_1 \arrow[d, "f'"'] \\
A_2 \arrow[r, "i'"] & A_3
\end{tikzcd} \tag*{(2.7.7)} \]
\textit{avec $i$ monomorphisme, si $f \in W$, alors $f' \in W$. (Stabilité de $W$ par cochangement de base par monomorphismes.)}
\end{itemize}

\textit{Sous ces conditions, ce triple de Quillen est spécial, et a fortiori clos.}

Ce corollaire résulte de 2.6. 3°), puisque $C \subset \Cof_W$ implique $TC \deq C\cap W = C\cap W^{\mathrm{univ}}$ (car $\Cof_W \cap W \ = \ W^{\mathrm{univ}}$), donc que $TC$ est stable par cochangement de base, d'où la condition (2.6.5. a).

\dgrblock{dgrb1175}{The recognition argument and its essential cancellation}
\begin{dcomment}{exposition}{p27:Ce corollaire résulte de 2.6. 3°)}{The recognition argument and its essential cancellation}
Grothendieck's criterion is designed so that the difficult set-theoretic work produces two factorizations, after which the model axioms follow from a short argument. Keeping these stages distinct also prevents the recalled theorem from being confused with a broader assertion about arbitrary accessible classes of monomorphisms. The latter is not what Theorem~2.6 needs.

Set $C=\operatorname{Mono}$, $TC=C\cap W$, $TF=C^{\square}$ and $F=(TC)^{\square}$. Suppose that the two asserted weak factorization systems have been obtained and that $TF\subset W$. Then $TF\subset F$ because $TC\subset C$. We verify the remaining equality $F\cap W=TF$. Given $f\in F\cap W$, use the first factorization system to write
\[
 f=pi,\qquad i\in C,\quad p\in TF.
\]
Since $p,f\in W$, two-out-of-three puts $i$ in $TC$. Since $f$ has the right lifting property against $i$, the usual diagonal makes $f$ a retract of $p$ in the arrow category. Right lifting classes are retract-closed, so $f\in TF$. The reverse inclusion has already been observed. Hence the two weak factorization systems have exactly the intersections required for a closed Quillen triple.

The alternative condition (2.6.6) is the same retract argument in a different position. If $TF\subset W$, factorization with respect to $(C,TF)$ supplies a monomorphism followed by a weak equivalence. Conversely suppose each $f\in TF$ factors as $f=pi$ with $i$ monic and $p\in W$. The right lifting property of $f$ against $i$ makes $f$ a retract of $p$, so $f\in W$, provided $W$ is retract-closed. That last property is already forced by (2.6.3a): if $f$ is a retract of $g\in W$, split the associated idempotent $e:g\to g$ by the chain
\[
 g\xrightarrow e g\xrightarrow e g\xrightarrow e\cdots .
\]
Its colimit in the arrow category is $f$. Filtered-colimit closure therefore gives $f\in W$.

The historical proof of the generating-set assertion should be located accurately. XV, Theorem~6.1 carries the author's warning that its proof is invalid. The surviving argument is the separate Lemma~7.3 and its application in~7.1--7.4. Starting with $A\hookrightarrow B$ in $TC$, one chooses a bounded subobject $B_0\subset B$ for which $A\cap B_0\hookrightarrow B_0$ lies in $TC$ and adds it to $A$. The density argument permits $B_0$ to contain a prescribed bounded piece of $B$; when $A\ne B$, choose that piece not contained in $A$. Thus the union strictly enlarges $A$. Effective unions keep $A\cup B_0\to B$ monic; pushout stability makes $A\to A\cup B_0$ weak. Since $A\to B$ is weak, cancellation makes $A\cup B_0\to B$ weak. Repeat this construction, taking unions at limit stages. Filtered-colimit closure keeps the resulting map to $B$ in $TC$ at such stages. Since $B$ has only a set of subobjects, strict enlargement cannot continue indefinitely: the process reaches $B$. The cancellation step is the extra ingredient missing for an arbitrary accessible mono class \citep[§7, especially Lemma~7.3 and Corollary~7.4]{Grothendieck1990XV}.

A later independent formulation is Smith's recognition theorem, proved in \citet[Theorem~1.7]{Beke2000}: with a locally presentable underlying category and a chosen set generating the cofibrations, accessibility supplies the solution-set condition for weak equivalences, and pushout/transfinite closure of the trivial cofibrations completes the construction. Beke's Proposition~1.12 identifies hypotheses ensuring that all monomorphisms admit such generators. This provides a modern route to the intended very-special case without appealing to the withdrawn general claim.

Corollary~2.7 records an additional condition, left properness. Its pushout square tests whether a weak equivalence remains weak after attachment along a monomorphism. When the attached map itself is both a monomorphism and a weak equivalence, the two-out-of-three argument in the source identifies it with a universally weak map, whose every cobase change is weak. Hence left properness supplies the pushout stability of $TC$ required by Theorem~2.6. The recognition theorem needs that stability; the stronger properness property is singled out for applications rather than inserted into the definition of a model structure.
\end{dcomment}


\printedpage{17}
% head: A. Grothendieck : Dérivateurs, Ch. XIX, Édition des Universités de Montpellier II et Paris VII
% foot: Version 11 février 2019

\textbf{2.8.} On dit qu'une \textit{structure de Quillen} $(W,C,F)$ sur $\mathcal{M}$ est \textit{très spéciale,} si elle est obtenue comme au théorème 2.1., i.e. si on a les conditions :
\[ \left\{ \begin{minipage}{0.72\textwidth}
\begin{itemize}[leftmargin=2.5em]
\item[a)] $\mathcal{M}$ est très spéciale, i.e. satisfait aux conditions (2.6.1).
\item[b)] $W$ est stable par petites $\varinjlim$ filtrantes, et est accessible (conditions (2.6.3)).
\item[c)] $C = \Mono$.
\end{itemize}
\end{minipage}\right. \tag*{(2.8.1)} \]
Donc pour qu'un localiseur $W$ dans $\mathcal{M}$ corresponde à une structure de Quillen très spéciale \msp{page 28} dans $\mathcal{M}$ (laquelle sera forcément unique), il faut et il suffit que $\mathcal{M}$ et $W$ satisfassent les conditions (2.8.1. a)\,b)), et que de plus on ait :
\[ \left\{ \begin{minipage}{0.65\textwidth}
\begin{itemize}[leftmargin=2.5em]
\item[a)] $TC \deq \Mono\cap W$ est stable par cochangement de base.
\item[b)] $TF \deq (\Mono)_* \subset W$, i.e. $W$ contient les flèches de type injectif dans $\mathcal{M}$.
\end{itemize}
\end{minipage}\right. \tag*{(2.8.2)} \]
\footnote[16]{On dit aussi que le localiseur $W$ est très spécial, s'il définit une structure de Quillen très spéciale.}.

\dgrblock{dgrb1176}{Injective structures and the role of accessibility}
\begin{dcomment}{later}{p28:laquelle sera forcément unique}{Injective structures and the role of accessibility}
Grothendieck has reached a class of coefficient categories for which the second half of the diagram theory can be constructed. Theorem~2.9 uses componentwise cofibrations, so restriction along every functor preserves both cofibrations and trivial cofibrations. The right Kan adjoint is then a right Quillen functor. This is the injective model structure; the square brackets in $F[I]$ indicate that its fibrations have been determined by a lifting property on whole diagrams.

There are two sorts of permanence hidden in the phrase ``par AQT.'' The exactness conditions are evaluated at objects of $I$: limits, colimits, intersections, pushouts and the maps from unions to ambient objects are computed there. Monicity is also componentwise. For the converse direction of that assertion, apply the defining monomorphism test to free diagrams $i_!A$: monicity of $X\to Y$ in the diagram category says that $\mathcal M(A,X_i)\to\mathcal M(A,Y_i)$ is injective for every $A$, hence $X_i\to Y_i$ is monic. The existence of the free adjoint is what makes this argument valid in the stated generality.

The size statements require a common bound. One chooses a cardinal larger than the arrow set of $I$ and enlarges the bounds used for the coefficients and for $W$. A small diagram then involves only a bounded family of objects, structure maps and equations. A sufficiently large accessible presentation can carry all these data at one stage. Equalities of composites are imposed at a later common stage, and the filtered system of such choices recovers the original diagram. Thus the bounds are allowed to grow with the shape. The source's cardinal-filtration hypothesis is doing real work here; accessibility of each individual object alone would not justify this conclusion. A precise antecedent for the diagram-filtration step is \citet[Exposé~I, Corollary~9.24]{SGA4I1972}, with the indexing category held fixed and the cardinal enlarged as prescribed there.

In the later combinatorial formulation, \citet[Lemma~A.2.8.3]{Lurie2009HTT} proves that the componentwise class associated to a set-generated weakly saturated class is again set-generated in a small diagram category. Proposition~A.2.8.2 then gives both projective and injective model structures. This settles the obstacle of page~24 when local presentability is imposed; it should not be quoted as a theorem about every historical ``special'' category with that hypothesis omitted. For presheaves with an accessible localizer and monomorphism cofibrations, the corresponding statement is \citet[Corollary~1.4.22]{Cisinski2006}. These references supply exact later versions of the passage from coefficients to diagrams.

The derived right adjunction itself can be checked without new size arguments. Since $u^*$ preserves the componentwise trivial cofibrations, adjunction shows that $u_*$ preserves injective fibrations. Applying the same reasoning to all componentwise cofibrations gives preservation of trivial fibrations. It preserves the terminal object, hence fibrant objects. The dual of the Brown argument above shows that it preserves weak equivalences between those objects. Localizing its restriction to fibrant objects defines $u_*^W$, and the derived unit and counit give
\[
 u_W^*\dashv u_*^W.
\]
The replacement needed for the pointwise formula can also be specified. Let $q:E=(u\downarrow j)\to I$ forget the arrow $u(i)\to j$. This projection is a discrete fibration: an arrow $i'\to i$ uniquely lifts into an object $(i,u(i)\to j)$ by precomposition. Consequently its left Kan extension on contravariant diagrams is
\[
 (q_!Y)(i)\cong\coprod_{e\in q^{-1}(i)}Y(e).
\]
Coproducts preserve cofibrations and trivial cofibrations, so $q_!$ preserves both componentwise left classes. The adjunction $q_!\dashv q^*$ therefore makes $q^*$ a right Quillen functor for the injective structures. In particular, if $X\to RX$ is an injective fibrant replacement, then $q^*RX$ is injectively fibrant and $q^*X\to q^*RX$ is a weak equivalence. Applying the ordinary right Kan formula to $RX$ now gives
\[
 (u_*^WX)_j\cong
 \operatorname{holim}_{(u\downarrow j)^{\mathrm{op}}}q^*X.
\]
Here the comma category has objects $u(i)\to j$, opposite to the comma category used for the left extension. This accounts for the two different fibre formulas in the derivator axioms.

The two model structures have the same weak equivalences and therefore produce the same localized diagram category. Their purpose is computational: one makes the left adjoints homotopical on cofibrant objects, the other makes the right adjoints homotopical on fibrant objects. The special monomorphism hypotheses guarantee these choices in the source's proof. They are not an invariant property that every presentation of the resulting derivator must share.
\end{dcomment}


On trouve maintenant le

\textbf{Théorème 2.9.} \textit{Soit} $(W, C = \Mono, F = (TC)_* = (W\cap\Mono)_*)$ \textit{un triple de Quillen très spécial dans une catégorie $\mathcal{M}$} (2.8). \textit{Alors :}
\begin{itemize}[leftmargin=2.5em]
\item[1°)] \textit{Pour toute $I$ dans} Cat, $(W(I),C(I),F[I])$ \textit{est un triple de Quillen très spécial dans $\mathcal{M}(I)$, et on a en particulier}
\[ C(I) \ = \ \Mono(\mathcal{M}(I))\ , \tag*{(2.9.1)} \]
\textit{et de plus}
\[ F[I]\cap W(I) \ = \ (TF)[I]\ . \tag*{(2.9.2)} \]
\item[2°)] \textit{Pour toute flèche $u : I \longrightarrow I'$ dans} Cat, \textit{d'où le couple de foncteurs adjoints}
\[ \begin{tikzcd}
\mathcal{M}(I) \arrow[r, shift right, "u_*"'] & \mathcal{M}(I') \arrow[l, shift right, "u^*"']
\end{tikzcd}\ , \tag*{(2.9.3)} \]
\textit{on a}
\[ \left\{ \begin{array}{rcl}
u^*(W(I')) & \subset & W(I)\\
u^*(C(I')) & \subset & C(I)\\
u^*(TC(I')) & \subset & TC(I)\\
u_*(F[I]) & \subset & F[I']\\
u_*(TF[I]) & \subset & TF[I']\ ,
\end{array}\right. \tag*{(2.9.4)} \]
\textit{enfin $u_*$ induit un foncteur}
\printedpage{18}
% head: A. Grothendieck : Dérivateurs, Ch. XIX, Édition des Universités de Montpellier II et Paris VII
% foot: Version 11 février 2019
\msp{page 29}
\[ u_*^{\mathrm{f}} \ :\ \mathcal{M}(I)_{\mathrm{f}} \longrightarrow \mathcal{M}(I')_{\mathrm{f}} \qquad\quad ^{(17)} \tag*{(2.9.5)} \]
\footnotetext[17]{sous-catégories pleines des objets fibrants.}
compatible aux localiseurs $W(I)_{\mathrm{f}}$ \textit{et} $W(I')_{\mathrm{f}}$ \textit{induits par $W(I)$, $W(I')$ respectivement, donc induit un foncteur}
\[ u_*^W \ :\ \Ho_W(I) \longrightarrow \mathit{Ho}_W(I')\ . \tag*{(2.9.6)} \] % [sic] italic Ho_W(I') -> upright Ho_W(I'), inconsistent with the upright Ho_W elsewhere
\textit{Ce foncteur est un} adjoint à droite \textit{du foncteur}
\[ u_W^* \ :\ \mathit{Ho}_W(I') \longrightarrow \Ho_W(I) \tag*{(2.9.7)} \] % [sic] italic Ho_W(I') -> upright Ho_W(I'), inconsistent with the upright Ho_W elsewhere
\textit{induit par $u^*$} (2.9.3).
\end{itemize}
\textsc{Démonstration de 2.9.} Pour prouver 1°), il faut vérifier d'abord que les conditions préliminaires (2.8.1. a)\,b)) sur $\mathcal{M}$, $W$ sont stables par passage de $\mathcal{M}$, $W$ à $\mathcal{M}(I)$, $W(I)$, ce qui se fait par AQT \texttt{[âne qui trotte]}. D'ailleurs on vérifie tout de suite que
\[ C(I) \ = \ \Mono(\mathcal{M}(I))\ , \]
et le théorème 2.6 (1°, 2°) implique alors que
\[ \begin{array}{ccccc}
C(I) & \qpair & C(I)_* & \deq & TF[I]\\
\underbrace{TC(I)}_{\deq C(I)\cap W(I)} & \qpair & (TC(I))_* & \deq & F[I]
\end{array} \] % reading resolved: DER_XIX_corrections.csv
sont des couples de Quillen spéciaux. (En notant pour le deuxième, que $TC(I)$ est bien \msp{page 30} stable par cochangement de base, $TC$ l'étant.) Pour vérifier que l'on a un triple de Quillen très spécial, il reste en vertu de 2.8 à vérifier que
\[ TF[I] \ \subset\ W(I)\ , \]
ce qui s'écrit
\[ v^*(TF[I]) \ \subset\ W(I_0)\ , \]
où
\[ v \ :\ I_0 \longrightarrow I \]
est défini comme dans (1.12) (p. 10). En fait, on aura même
\[ v^*(TF[I]) \ = \ v^*v_*(TF(I_0)) \ \subset\ \underbrace{TF(I_0)}_{= F(I_0)\cap W(I_0)}\quad , \]
et la démonstration est duale de celle de
\[ v^*v_!(TC(I_0)) \ \subset\ TC(I_0) \]
(cf. (2.3.7), p. 21), elle ne dépend pas des questions de factorisation dans $\mathcal{M}(I)$, et provient ici de la stabilité de $TF$ par petits produits. Cela prouve 1°.

\dgrblock{dgrb1177}{The cofree image is smaller than the lifting class}
\begin{dcomment}{development}{p30:elle ne dépend pas des questions de factorisation}{The cofree image is smaller than the lifting class}
Grothendieck needs the inclusion $TF[I]\subset W(I)$ to finish the injective model structure. The intended conclusion is correct, but the equality involving cofree diagrams on this page is stronger than the adjunction permits. We spell out both the direct repair and a finite example distinguishing the two classes.

First there is a short proof of the necessary inclusion. For $v:I_0\to I$, the free-diagram functor $v_!$ takes componentwise monomorphisms to monomorphisms. Indeed, at $i$ it is a coproduct of the given component maps, and these coproducts are monic under the hypotheses in use. One can also see the coproduct assertion by expressing a coproduct of monomorphisms as a transfinite succession of pushouts, using closure of monomorphisms under filtered colimits at limit stages. If $p\in TF[I]=\operatorname{Mono}(\mathcal M(I))^{\square}$, adjunction now gives
\[
 v^*p\in\operatorname{Mono}(\mathcal M(I_0))^{\square}=TF(I_0).
\]
Since $TF\subset W$, this proves the required inclusion. No factorization of $p$ and no assertion that it is cofree have been used.

\textbf{Finite separation of the two image classes.}\ours
In $\mathrm{Set}$, with weak equivalences all maps and cofibrations the injections, take $I=(0\to1)$. Then
\[
 v^*(TF[I])\ne v^*v_*(TF(I_0)).
\]
Here the model structure is generated by the usual injection--surjection weak factorization system; it satisfies the intended, inductive-limit version of the very-special hypotheses.

\emph{Proof.} A contravariant $I$-diagram is a map $X_1\to X_0$. Let $X_0$ have two elements and $X_1$ three, and choose a surjection $r:X_1\twoheadrightarrow X_0$. Let $p:X\to1$ be the map to the terminal diagram. We show directly that it has the right lifting property against every componentwise injection $A\hookrightarrow B$. Regard $A_i$ as a subset of $B_i$. Given a map $A\to X$, first extend $A_0\to X_0$ to $h_0:B_0\to X_0$ by sending all new elements to one fixed element of $X_0$. Choose a section $s:X_0\to X_1$ of $r$. On $B_1\setminus A_1$, define
\[
 h_1(b)=s\bigl(h_0(r_B(b))\bigr),
\]
and on $A_1$ retain the specified map. The two definitions have disjoint domains and agree with the required compatibility on $A_1$. Thus $rh_1=h_0r_B$, giving a lift of diagrams. Hence $p\in TF[I]$.

For a family $(E_0,E_1)$, the cofree diagram is
\[
 (v_*E)_0=E_0,\qquad (v_*E)_1=E_0\times E_1,
\]
with the projection as its structure map. Consequently the restriction of a cofree arrow has component source sizes $|E_0|$ and $|E_0||E_1|$. If it were isomorphic, after restriction, to $p$, these sizes would be $2$ and $3$. The equation $2|E_1|=3$ has no solution. Therefore $v^*p$ is not in the cofree image class, even up to isomorphism. This proves the assertion.\hfill$\square$

The product argument in the source does establish
\[
 v^*v_*(TF(I_0))\subset TF(I_0),
\]
because right lifting classes are closed under products. It does not replace the direct lifting proof for the larger class $TF[I]$. In a closure formulation one must retain the appropriate right-orthogonal saturation; deleting it creates exactly the mistaken equality.

We checked the diagram-model discussions of \citet[§A.2.8]{Lurie2009HTT} and \citet[Corollary~1.4.22]{Cisinski2006}; both use lifting to establish the relevant preservation statements. We did not find this particular finite illustration there. The example is recorded as a local development of the printed argument, with no claim of priority for the general distinction between cofree objects and a right lifting class.
\end{dcomment}


La démonstration de 2° est maintenant la duale de celle de (2.3. 2°). Comme cela nous fait répéter deux fois la même démonstration, cela montre qu'on n'a pas “monté le sorite” tout à fait avec assez de soin. Donc à titre de rattrapage, j'ajoute un

% ---- from der-ch19_p021-025.tex
\printedpage{19}
% head: A. Grothendieck : Dérivateurs, Ch. XIX, Édition des Universités de Montpellier II et Paris VII
% foot: Version 11 février 2019

\msp{page 31}
\textbf{Scholie 2.10.} \textit{Soit $(W,C,F)$ un triple de Quillen clos dans $\M$. Pour $I$ dans \textup{Cat}, pour que $(W(I),C\langle I\rangle,F(I))$ soit un triple de Quillen clos pour $\M(I)$, il faut et il suffit que les couples $C\langle I\rangle\longleftrightarrow TF(I)$ et $TC\langle I\rangle\longleftrightarrow F(I)$ soient factorisants. Supposons que cette condition soit satisfaite pour $I$ et pour $I'$, et soit $u : I\longrightarrow I'$ une flèche de $I$ dans $I'$. Alors les assertions de \textup{2.3. 2°} sont valables, en particulier on trouve un foncteur $u^{W}_{!} : \Ho_W(I)\longrightarrow\Ho_W(I')$, adjoint à gauche de $u^{*}_{W}$\footnote[18]{Dire aussi que le calcul des fibres de $u^{W}_{!}(X)$, prévu dans Dér 5 bis, est OK.}. Ces résultats se dualisent aussitôt (par $\M\longmapsto\M^{\circ}$, $I\longmapsto I^{\circ}$) pour les triples $(W(I),C(I),F[I])$, et pour la construction de $u^{W}_{*}$, adjoint à droite de $u^{*}_{W}$. En particulier, pour que le triple précédent soit un triple de Quillen (forcément clos), il faut et il suffit que les deux couples $C(I)\longleftrightarrow TF[I]$ et $TC(I)\longleftrightarrow F[I]$ soient factorisants ; et quand cette condition est satisfaite pour $I$ et pour $I'$ (pour une flèche $u : I\longrightarrow I'$ dans \textup{Cat} donnée), alors $u^{W}_{*}$ existe et se construit dualement de la construction de $u^{W}_{!}$ explicitée dans \textup{2.3. 2°}.}

\msp{page 32}
% a vertical bar is printed in the left margin alongside Corollaire 2.11
\textbf{Corollaire 2.11.} \textit{Le prédérivateur $\mathbf{D}_{\M,W}$ défini pour un localiseur \textup{très spécial} $W$ dans $\M$ est un dérivateur en plein sens du terme, i.e. il satisfait à tous les axiomes Dér 1 à Dér 6 (y inclus Dér 4 bis et Dér 5 bis et Dér 6 bis) des dérivateurs.}

\textbf{N.B.} À ma propre surprise, je n'ai pas été capable de trouver un cas d'un localiseur $W$ qui définisse un dérivateur sur Cat tout entier (avec existence de $u^{W}_{!}$, $u^{W}_{*}$), \textit{que} pour $W$ un localiseur très spécial, donc s'insérant dans un triple de Quillen très spécial -- ou quand on peut s'y ramener de la façon standard. (Par exemple dans le cas où $\M=\Cat$, $W$ un localiseur fondamental satisfaisant Loc 6 a) et Loc 8, 8 bis)\footnote[19]{axiome des carrés cocartésiens (version directe), axiome des limites inductives, accessibilité.} -- les monomorphismes dans $W$ se refusent formellement à être stable par cochangement de base !) J'ignore totalement si le résultat mirobolant ``prouvé'' par \textsc{Anderson} est vrai ou faux. Ça a l'air presque trop beau pour être vrai \ldots{}D'ailleurs, je ne connais pas de structure de Quillen spéciale qui ne soit très spéciale -- désolé ! % [sic] stable -> stables, agreement with « les monomorphismes »

\dgrblock{dgrb1178}{The existence question after Chapter XIX}
\begin{dcomment}{status}{p32:J'ignore totalement si le résultat mirobolant}{The existence question after Chapter XIX}
The final note measures the distance between the construction just completed and the general statement Grothendieck hoped for. Theorem~2.11 provides both Kan extensions by imposing enough smallness and exactness to construct two model structures on every diagram category. The subsequent uncertainty concerns how much of this apparatus is necessary. That existence question now has a definitive answer for ordinary complete and cocomplete Quillen model categories.

\citet[Theorem~6.11]{Cisinski2003} proves the existence and pointwise calculation of homotopy right Kan extensions for a Quillen model category admitting all small limits; applying the dual result supplies the left extensions when all small colimits exist. This gives the associated derivator on all small categories, without requiring the original category to be combinatorial or the diagram categories to possess the two projective and injective model structures used here. The same general statement is recorded in \citet[Theorem~1.32]{Groth2013}. Groth's Proposition~1.30 gives a shorter proof for combinatorial models and establishes the usual strength axiom in that setting. These later results settle the mathematical question; they do not certify the details of the historical proof attributed to Anderson.

There is already a very small distinction between ``special'' and ``very special.'' On $\mathrm{Set}$ take
\[
 W=\operatorname{Iso},\qquad C=F=\operatorname{Fl}(\mathrm{Set}).
\]
The two weak factorization systems are $(\operatorname{Fl},\operatorname{Iso})$ and $(\operatorname{Iso},\operatorname{Fl})$. The first is generated by $\varnothing\to1$ and $2\to1$: right lifting against the first means surjectivity, and right lifting against the second means injectivity. The empty generating family gives the second system. All sets are accessible at a suitable cardinal and all small colimits exist, so this is special in the chapter's sense. It is not very special because $2\to1$ is a cofibration but not a monomorphism. Its associated derivator is the represented diagram theory of sets. This elementary example separates a chosen model presentation from the special construction with all monomorphisms as cofibrations.

The warning about $\mathrm{Cat}$ points to a different, genuine obstruction. Here is a complete finite construction for the ordinary weak-homotopy localizer. Let $P$ be the poset of nonempty proper faces of a tetrahedron. Its nerve is the barycentric subdivision of the boundary sphere, so $|NP|$ is homeomorphic to $S^2$. For each nonidentity arrow $a$ of $P$, take a copy of the walking arrow $[1]$ and its inclusion in the groupoid $J$ with two objects and one isomorphism between them. Form the disjoint unions
\[
 A=\coprod_a[1],\qquad B=\coprod_a J,
\]
and map $A\to P$ by the chosen arrows. The map $A\to B$ is a monomorphism of categories: it is injective on objects and faithful. It is also a weak homotopy equivalence, since it is a disjoint union of maps between contractible nerves.

The pushout $P\amalg_A B$ is the ordinary groupoid completion of $P$, by the universal property of freely adjoining inverses to all arrows. The fundamental groupoid of the nerve of a category is its groupoid completion: its generators are the one-simplices, and its two-simplex relations are exactly composition. Since $S^2$ is connected and simply connected, this groupoid is equivalent to the terminal groupoid. Thus the pushout map
\[
 P\longrightarrow P\amalg_A B
\]
induces $S^2\to *$ on homotopy types and is not a weak equivalence, as its effect on $H_2(-;\mathbb Z)$ already shows. All categories used to specify the pushout are finite. The claim on page~32 that weak monomorphisms in $\mathrm{Cat}$ need not survive pushout therefore has a concrete content; one cannot obtain the usual homotopy theory of categories by simply declaring every monomorphism to be a cofibration.

Scholie~2.10 isolates what the direct proof actually requires: the two lifting pairs must factor at the diagram shapes being used. For its pointwise assertion one must also justify the corresponding comma-shape calculation, as done above when all the projective structures are available. Cisinski's theorem supplies the general alternative when that direct model-structure construction is unavailable. The categorical limits and colimits remain essential hypotheses; this is the distinction the examples in the following section must continue to respect.
\end{dcomment}


\msp{page 33}

\section*{3\quad Cas des catégories à cofibrations, des catégories à fibrations}

Je rappelle les axiomes d'une \textit{structure à cofibrations} $(W,C)$ sur une catégorie $\M$.

\begin{itemize}[leftmargin=4em,labelwidth=3.5em,align=left]
\item[Cof$_0$ :] $W$ est un localiseur.
\item[Cof$_1$ :] $C$ contient les isomorphismes, est stable par cochangement de base, par composition.
\item[Cof$_2$ :] $C\cap W$ est stable par cochangement de base (et aussi par composition, et contient les isomorphismes, en vertu de Cof$_0$, Cof$_1$).
\item[Cof$_3$ :] Toute flèche $f$ dans $\M$ se factorise en $f=pi$, avec $i\in C$, $p\in W$.
\end{itemize}

(Axiomes de K. S. \textsc{Brown}, légèrement généralisés par \textsc{Anderson}). Il y a un ``axiome
\printedpage{20}
% head: A. Grothendieck : Dérivateurs, Ch. XIX, Édition des Universités de Montpellier II et Paris VII
% foot: Version 11 février 2019
de continuité'' ou \textit{``axiome transfini''} facultatif, qui est le suivant :

\begin{itemize}[leftmargin=4em,labelwidth=3.5em,align=left]
\item[Cof$_4$ :] $\M$ stable par petites $\varinjlim$ filtrantes, et $C$ et $TC = C\cap W$ sont stables par composition transfinie.
\end{itemize}

Ainsi, Cof$_1$, Cof$_2$, Cof$_4$ peuvent se résumer en disant que $C$ et $TC$ sont Q-saturés transfiniment à gauche.

La structure de cofibration est dite \textit{spéciale} si elle satisfait Cof$_4$, si $\M$ est spéciale, et si $C$ admet une \textit{petite} partie $C_0$ qui soit Q-génératrice à gauche, i.e. telle que
\[
C \ \subset\ C_0\Qs\,, \qquad \text{i.e.}\quad \tilde{C}=C_0\Qs. \tag*{(3.1)}
\]

\msp{page 34}
Alors, par le théorème de factorisation 1.1, on a
\[
C \ \fact\ TF \ \deq\ C_* \ ,
\]
et l'axiome de factorisation Cof$_3$ implique (cf. 1.4) $TF\subset W^{\natural}$. Plus généralement, on a :

\textbf{Proposition 3.1.} \textit{Supposons que $\M$ soit spéciale (stable par petites $\varinjlim$ , et éléments accessibles), que $(W,C)$ satisfasse à Cof$_0$, Cof$_1$, et que $C$ soit stable par composition transfinie (ce qui, joint à Cof$_1$, s'énonce en disant que $C$ est Q-saturé transfini à gauche), enfin que $C$ admette une petite partie $C_0$ Q-génératrice à gauche, de sorte que $C\fact TF\deq C_*$. Alors pour que l'axiome de factorisation Cof$_3$ soit satisfait, il suffit qu'on ait}
\[
TF \ \subset\ W \ , \tag*{(3.1.1)}
\]
\textit{et il faut qu'on ait $TF\subset W^{\natural}$, donc si $W=W^{\natural}$, l'axiome Cof$_3$ équivaut à} (3.1.1).

\dgrblock{dgrb1179}{What one factorization can accomplish}
\begin{dcomment}{exposition}{p34:l'axiome Cof$_3$ équivaut}{What one factorization can accomplish}
Grothendieck now removes the fibrations from the starting data. His purpose is to isolate what is actually needed to construct homotopy colimits: the ability to replace a map by a cofibration, together with the invariance of gluing along an acyclic cofibration. The resulting structure is deliberately asymmetric. It controls homotopy cofibres without committing to a second factorization, and therefore without assuming that every lifting class already belongs to a model structure.

Write $TC=C\cap W$ and, in modern notation, write $S^\perp$ for the maps having the right lifting property against $S$. The symbol $TF=C_* =C^\perp$ initially names a lifting class; its membership in $W$ is a separate assertion. Proposition 3.1 keeps this distinction correctly. If the small object argument gives a factorization $f=pi$ with $i\in C$ and $p\in C^\perp$, then $C^\perp\subset W$ immediately supplies Cof$_3$. Conversely, apply Cof$_3$ to a map $f\in C^\perp$, obtaining $f=pi$, $i\in C$, $p\in W$. The square with left side $i$, right side $f$, and identity on the common source has a diagonal $s$ satisfying $si=1$ and $fs=p$. In the arrow category, $(i,1)$ and $(s,1)$ exhibit $f$ as a retract of $p$. This proves exactly $C^\perp\subset W^{\natural}$; replacing the latter class by $W$ requires retract closure. No additional lifting axiom has entered the argument.

The two closure requirements on $C$ and $TC$ have different tasks. Closure under cobase change gives the elementary gluing operation. Closure under transfinite composition permits a diagram to be assembled through its successive stages. Closure under retracts then turns a factorization into a characterization by lifting. In the terminology of these pages, Q-closed means equality with ${}^\perp(S^\perp)$, whereas Q-saturated describes specified operations on arrows. Theorem 1.1 identifies them under a small generation hypothesis; section 3 ends by asking how far that hypothesis can be removed.

The historical starting point is Brown's theory of categories of fibrant objects, used here in its dual form \citep{Brown1973}. Theorem-Scholie~2 of Chapter~VI, manuscript p.~108, already places this asymmetry in the development of the manuscript. What is new here is its systematic application to categories of diagrams. Grothendieck's assertion that a special cofibration structure has functorial factorizations uses the functorial version of the small object construction from section 1. It is a conclusion of that construction, not part of the definition of an arbitrary cofibration category.
\end{dcomment}


La structure de cofibration est dite \textit{semi-close} si $C$ est Q-clos à gauche, i.e. si on a $C=C\Qs$ (ce qui implique que $C$ est stable par composition transfinie), \textit{close} si de plus $TC=(TC)\Qs$ (ce qui implique l'axiome transfini Cof$_4$).

\msp{page 35}
Je ne reviens pas ici sur la théorie de K. S. \textsc{Brown} des catégories à cofibrations. C'est à ma connaissance l'ensemble d'axiomes le plus économique connu, pour assurer d'une part une bonne prise sur $\Ho_W$ (par une notion correspondante d'homotopie d'applications, et un calcul de fractions dans $\Pi\M_{\mathrm{c}}$), et d'autre part l'existence de la suite \textit{exacte} de cofibrations, donc une théorie de cofibres $W$-homotopiques. (\textbf{N.B.} Quand la structure est \textit{propre,} i.e. $C\subset\Cof_W$, alors l'existence de cette structure revient à dire simplement que ``$W$ admet assez de $W$-cofibrations'', i.e. que toute $f\in\Fl(\M)$ se factorise en $pi$, avec $i\in\Cof_W$, $p\in W$ ; et on peut donc prendre $C=\Cof_W$. Dans le cas où $W$ est stable par petites $\varinjlim$ filtrantes, $C=\Cof_W$ et $TC=W^{\mathrm{univ}}$ satisfont de plus l'axiome transfini Cof$_4$.) Enfin, K. S. \textsc{Brown} montre que toute flèche $i : X\longrightarrow Y$ dans $C$, \textit{avec $X$, $Y$ cofibrants} (i.e. $\emptyset_{\M}\longrightarrow X$, $\emptyset_{\M}\longrightarrow Y$ dans $C$) est dans $\Cof_W$ (``Lemme de Brown'').

\dgrblock{dgrb1180}{Brown excision and the scope of properness}
\begin{dcomment}{examples}{p35:avec $X$, $Y$ cofibrants}{Brown excision and the scope of properness}
Grothendieck wants gluing to depend only on the weak equivalences, which explains the intrinsic class $\Cof_W$. One qualification is necessary in the final sentence on this page. Brown excision gives homotopy invariance of a pushout along a cofibration when the weak-equivalence comparison is between cofibrant objects. It does not say that a cofibration with cofibrant endpoints belongs to the ambient, universally defined class $\Cof_W$.

This is a distinction in the manuscript's own terminology. Chapter XIII, ms.~p.~8, defines $\Cof_W$ by two consecutive pushout squares: after any initial cobase change, pushing out a weak equivalence must again give a weak equivalence \citep[Ch.~XIII, ms.~p.~8]{Grothendieck1990XIII}. Thus it requires invariance on the entire undercategory. The safe form of Brown's assertion is obtained by restricting the coefficient category and the comparison maps to cofibrant objects. Alternatively, one may assume $C\subset\Cof_W$, the properness hypothesis already distinguished in the preceding parenthesis.

A concrete obstruction makes the difference visible. In simplicial associative unital $\mathbb Z$-algebras, let $A=\mathbb Z/p$ be constant and choose a simplicial free algebra resolution $B\longrightarrow A$. Adjoining a free noncommuting generator has underlying simplicial abelian group
\[
 B\langle x\rangle\cong\bigoplus_{n\geq1}B^{\otimes_{\mathbb Z} n}.
\]
The summand for $n=2$ has first homotopy group
$\operatorname{Tor}^{\mathbb Z}_1(A,A)\cong\mathbb Z/p$, whereas
$A\langle x\rangle$ is constant. Hence
$B\langle x\rangle\longrightarrow A\langle x\rangle$ is not a weak equivalence. But $\mathbb Z\longrightarrow\mathbb Z\langle x\rangle$ is a free cofibration with cofibrant source and target. Its pushout functor is precisely adjoining $x$, so it is not a $W$-cofibration. This is the mechanism of \citet[Example 2.11, arXiv version]{Rezk2002Proper}. The error concerns the unrestricted formulation of the sentence, not Brown's excision theorem and not the finite-diagram constructions that follow.
\end{dcomment}


Je m'intéresse ici à la question d'étendre la structure à cofibrations à une catégorie de la forme $\M(I)$, et d'utiliser cette extension pour construire
\msp{page 36}
les foncteurs $u^{W}_{!}$. Notons déjà que la théorie de K. S. \textsc{Brown}, avec les compléments d'\textsc{Anderson}\footnote[20]{revus et arrangés par mes soins.}, donne essentiellement ceci :

\printedpage{21}
% head: A. Grothendieck : Dérivateurs, Ch. XIX, Édition des Universités de Montpellier II et Paris VII
% foot: Version 11 février 2019

\textbf{Proposition 3.2.} \textit{Soit $W\subset\Fl(\M)$ un localiseur associé à une structure à cofibrations $(C,W)$ dans $\M$.}
\begin{itemize}[leftmargin=3em,labelwidth=2.5em,align=right]
\item[1°)] \textit{Sur la catégorie \textup{$\Diag_0$} des ensembles ordonnés \textup{finis}, ce prédérivateur satisfait aux axiomes Dér 1 (décomposition), Dér 2 (localisation), Dér 3 (objet final), Dér 4 bis (existence de $u^{W}_{!}$), Dér 5 bis (calcul des fibres de $u^{W}_{!}$), Dér 6 bis (exactitude à gauche).}
\item[2°)] \textit{C'est encore vrai sur la sous-catégorie \textup{$\Diag_1$} dans \textup{Cat} plus grande, formée des ensembles ordonnés $I$ \textup{noethériens} tels que l'espace topologique associé à $I^{\circ}$ soit noethérien (toute suite croissante d'objets de $I$, ou d'ouverts de $I^{\circ}$, est stationnaire).}
\item[3°)] \textit{Enfin, si l'axiome transfini Cof$_4$ est satisfait, i.e. si $C$ et $TC$ sont stables par composition transfinie, alors la conclusion de \textup{1°} reste valable sur la sous-catégorie pleine \textup{$\Diag_2$} de \textup{Cat} formée des ensembles ordonnés qui sont noethériens sans plus.}
\end{itemize}

\textbf{Corollaire 3.3.} \textit{Supposons que l'on ait Cof$_3$ sous la forme plus forte d'une factorisation
\msp{page 37}
$f=pi$ \textup{fonctorielle} en $f$, ce qui est le cas en particulier si la structure à cofibrations est spéciale. Alors pour toute $I$ dans \textup{Cat}, $\M(I)$ muni de $(W(I),C(I))$ est une catégorie à cofibrations, laquelle satisfait l'axiome transfini Cof$_4$ si $(W,C)$ y satisfait. Donc on peut lui appliquer les résultats de \textup{3.2}, qui impliquent notamment que les axiomes Dér 1, 2, 3 ( ?) sont valables sur \textup{Cat} tout entier, et que $u^{W}_{!}$ existe quand $u$ est de la forme $\id_I\times v$, avec $v : J\longrightarrow K$ une flèche dans \textup{$\Diag_1$} (resp. dans \textup{$\Diag_2$}, si l'axiome transfini est satisfait).}

\textbf{3.4.} La clef de la théorie est dans l'extension de la structure à cofibrations $(W,C)$ de $\M$ à une catégorie $\M(I)$, où $I$ est un ensemble ordonné, en $(W(I),C_I)$. Les faits principaux sont énoncés dans XVII, dans les théorèmes p. 26-27, p. 50, p. 60. C'est quand $I$ est \textit{noethérien,} et qu'on suppose soit que $\mathrm{Esp}(I^{\circ})$ est également noethérien (p.ex. $I$ fini), soit que $(W,C)$ satisfait l'axiome transfini Cof$_4$, qu'on prouve que $(W(I),C_I)$ est également une structure à cofibrations sur $\M(I)$.

\dgrblock{dgrb1181}{Attaching a diagram in the order of its dependencies}
\begin{dcomment}{exposition}{p37:Les faits principaux sont énoncés dans XVII}{Attaching a diagram in the order of its dependencies}
Grothendieck's noetherian hypothesis makes an inductive construction possible: the value at an object can be attached after the values on which it depends have been constructed. With his convention $\M(I)=\operatorname{Fun}(I^\circ,\M)$, it is $I^\circ$ that serves as a direct indexing category. In a finite example this is the familiar Reedy construction, although the manuscript allows transfinite well-founded orders as well.

For a direct category $J$ and a diagram $X:J\longrightarrow\M$, let
\[
 L_jX=\mathop{\rm colim}_{(k\to j),\ k\ne j}X_k
\]
be the latching object, when the indicated colimit exists. A map $f:X\longrightarrow Y$ is a Reedy cofibration when
\[
 X_j\amalg_{L_jX}L_jY\longrightarrow Y_j
\]
belongs to $C$ for every $j$. The relative latching object incorporates all the compatibility already imposed at preceding objects. Merely requiring $X_j\longrightarrow Y_j$ to be a cofibration would omit those compatibility conditions.

For the arrow $0\to1$, the two tests are $X_0\to Y_0\in C$ and
$X_1\amalg_{X_0}Y_0\to Y_1\in C$. Given a map of arrows $X\to Y$, first factor its component at $0$ as $X_0\to Z_0\to Y_0$. Next factor the induced map $X_1\amalg_{X_0}Z_0\to Y_1$ as a cofibration followed by a weak equivalence. Its intermediate object defines $Z_1$, and the pushout supplies $Z_0\to Z_1$. The resulting $X\to Z$ is a Reedy cofibration and $Z\to Y$ is objectwise a weak equivalence. The same construction applies successively at every vertex of a finite direct category. Transfinite stages require colimits and the continuity condition Cof$_4$.

The identity $C_I\cap W(I)=(TC)_I$ records a second induction. In the arrow example, if the component at $0$ is acyclic, its pushout $X_1\to X_1\amalg_{X_0}Y_0$ is acyclic. Consequently the relative latching map is a weak equivalence exactly when $X_1\to Y_1$ is one, by two-out-of-three. This proves both directions of the identity at the second vertex. The general induction repeats this argument with the latching diagram; at a limit stage Cof$_4$ supplies the required continuity. This is why the equality, rather than the bare definition of $C_I$, is central to homotopy left Kan extension.

Corollary 3.3 uses a different construction. A functorial factorization in $\M$ can be applied to all components and all transition maps of an arbitrary diagram at once. Cof$_0$, Cof$_1$, and Cof$_2$ then hold objectwise, and functoriality supplies Cof$_3$ in $\M(I)$. This gives finite homotopy-colimit operations inside every shifted diagram category. It does not itself make the ordinary colimit functor preserve objectwise acyclic cofibrations. That additional issue is exactly what the projective classes $C\langle I\rangle$ are meant to address.
\end{dcomment}


\msp{page 38}
Un fait important pour la construction des foncteurs $u^{W}_{!}$, pour $u : I\longrightarrow I'$ une flèche entre catégories ordonnées noethériennes, est la relation
\[
C_I\cap W(I) \ = \ (TC)_I \ , \tag*{(3.4.1)}
\]
d'où résulte que l'on a, pour $u$ comme dessus,
\[
u_!(C_I\cap W(I)) \ \subset\ C_{I'}\cap W(I') \ , \tag*{(3.4.2)}
\]
ingrédient essentiel de la construction du foncteur $u^{W}_{!}$. Pour étendre cette construction au cas d'une flèche $\pm$ quelconque dans Cat, il se pose donc le problème d'étendre la définition de $C_I$, et les formules (3.4.1) et (3.4.2), au cas de catégories $\pm$ quelconques dans Cat. Or on a prouvé :

\textbf{Proposition 3.5.} \textit{Si $\Phi=\Psi^*$ est une partie \textup{Q-close} à gauche de \textup{$\Fl(\M)$}, et $I$ un ensemble ordonné noethérien, $\Phi_I$ l'extension de $\Phi$ en une partie de \textup{$\Fl(\M(I))$} définie par \textup{\textsc{Anderson}}, alors on a}
\[
\Phi_I \ = \ \Psi(I)^* \ = \ v_!(\Phi(I_0))\Qs \ , \tag*{(3.5.1)}
\]
\textit{où $v : I_0\longrightarrow I$ est le foncteur canonique de la catégorie discrète définie par \textup{$\mathrm{Ob}\,I$}.}

\printedpage{22}
% head: A. Grothendieck : Dérivateurs, Ch. XIX, Édition des Universités de Montpellier II et Paris VII
% foot: Version 11 février 2019

C'est (avec des notations différentes) le lemme 6.2.5 (p. 107) de XVII, dont je ne reprends pas ici la démonstration.

Ceci donne la clef d'une définition de $\Phi_I$ comme $\Phi\langle I\rangle$ (avec les notations des paragraphes 1, 2),
\msp{page 39}
pour $I$ quelconque dans Cat,
\[
\Phi_I \ \deq\ \Phi\langle I\rangle \ \deq\ \Psi(I)^* \ = \ v_!(\Phi(I_0))\Qs \ , \tag*{(3.5.2)}
\]
où $\Psi=\Phi_*$, et où on suppose $\Phi$ Q-clos à gauche, i.e. $\Phi=\Psi^*$. Dans le cas où on suppose seulement $\Phi$ Q-saturé transfini sans plus, on aurait envie de poser
\[
\Phi_I \ \deq\ \text{Q-saturé transfini de } v_!(\Phi(I_0)) \ , \tag*{(3.5.3)}
\]
et je présume que pour $I$ ordonné noethérien, cela coïncide avec la définition d'\textsc{Anderson}. Mais je vais me borner pour la suite au cas où $C=\Phi$ est Q-clos à gauche.

Revenons alors au cas d'une structure à cofibrations $(W,C)$ dans $\M$, en supposant $C$ et $TC=C\cap W$ Q-clos à gauche. Définissons pour toute $I$ dans Cat $C_I=C\langle I\rangle$ comme dans (3.5.2). On a alors, pour $u : I\longrightarrow I'$ une flèche dans Cat,
\[
u_!(C\langle I\rangle) \ \subset\ C\langle I'\rangle \ , \tag*{(3.5.4)}
\]
mais on n'y a guère de raison d'espérer avoir aussi (3.4.2), sans disposer de la formule
\[
C\langle I\rangle\cap W(I) \ \overset{?}{=}\ (TC)\langle I\rangle \qquad\qquad \text{\footnote[21]{\textbf{N.B.} On a l'inclusion $\supset$, car $v^*v_!TC(I_0)\subset TC(I_0)$.}}. \tag*{(3.5.5)}
\]
On a prouvé que cette formule est valable si $I$ est ordonné noethérien (XVII th. 2.1, p. 26-27). Cela
\msp{page 40}
nous amène à la

\textbf{Question 3.6.} \textit{Soit $\M$ une catégorie stable par petites limites, et soit $(W,C)$ une structure à cofibrations dans $\M$ telle que $C$ et $TC=C\cap W$ soient Q-clos à gauche. Soit $I$ un objet de \textup{Cat}, a-t-on alors}
\[
C\langle I\rangle\cap W(I) \ = \ (TC)\langle I\rangle \text{ ?} \tag*{(3.6.1)}
\]

\dgrblock{dgrb1182}{Lifting properties retain the coherence}
\begin{dcomment}{exposition}{p40:telle que $C$ et $TC=C\cap W$ soient Q-clos à gauche}{Lifting properties retain the coherence}
Grothendieck replaces a geometric attachment rule by a characterization that makes sense for every small category. Put $TF=C^\perp$, $F=(TC)^\perp$, and
\[
 C\langle I\rangle={}^\perp TF(I),\qquad
 (TC)\langle I\rangle={}^\perp F(I).
\]
The lifting problem here is a problem in the category of diagrams. A diagonal must commute with every arrow of $I$; a collection of independently chosen componentwise diagonals is insufficient. This is the coherence retained by the projective lifting class.

Let $v:I_0\to I$ include the discrete category of objects. The adjunction $v_!\dashv v^*$ gives
\[
 (v_!\Phi(I_0))^\perp
   =\{p:v^*p\in\Phi(I_0)^\perp\}
   =\Psi(I),
\]
where $\Psi=\Phi^\perp$ and $\Phi$ is Q-closed. The last equality follows because a lifting problem in a discrete family is solved component by component; equivalently, one tests a single component and puts identity maps in the others. Taking left orthogonals proves (3.5.2). The same adjunction calculation for $u_!\dashv u^*$ gives (3.5.4), since restriction preserves the objectwise right class. These arguments require existence of the ordinary left Kan extension, not a factorization system on diagrams.

The specification of $\Psi$ matters, just as in the arbitrary-test-class defect in (1.14) discussed earlier. If Proposition 3.5 is read as allowing an arbitrary presentation $\Phi={}^\perp\Psi$, its first equality needs qualification. In sets, take $\Psi=\{\emptyset\to1\}$. Then $\Phi$ consists of the maps with nonempty source together with $1_{\emptyset}$. On the direct arrow category $J=[1]$, the diagram map
\[
 f:(\emptyset\longrightarrow1)\longrightarrow(1\xrightarrow{1}1)
\]
fails the relative latching condition at $0$. Yet it lifts vacuously against $\Psi(J)$: that class contains only the map from the constant empty diagram to the constant singleton diagram, and the source of $f$ has no map to the former. Thus ${}^\perp\Psi(J)$ is too large for this nonclosed $\Psi$. The canonical choice $\Psi=\Phi^\perp$, explicitly imposed in (3.5.2), removes this issue. It leaves the intended construction and the subsequent propositions intact.

For a well-founded direct indexing order, the latching description and this lifting description agree by induction on objects: at stage $j$, the already constructed diagonals determine a map out of the relative latching object, and its lifting property supplies the next component. The converse can be checked explicitly for the posets in Proposition 3.5. Given a lifting problem from the relative latching map at $j$ to $p:A\to B$, use the diagram map which is $1_A$ at strict predecessors of $j$, is $p$ at $j$, and is the identity of the terminal object elsewhere. The map $L_jY\to A$ supplies its predecessors. A coherent diagonal then gives exactly the requested relative lift. The canonical right class $\Phi^\perp$ contains all identities, as this test requires. This proves the comparison under the stated existence assumptions.

There is also a direct cellular verification of the conjecture (3.5.3) whenever the small colimits needed for its transfinite construction exist. Assume that $\Phi$ contains isomorphisms and is closed under pushouts and small transfinite composition. The following is the direct-poset version of the standard Reedy cellular construction; compare \citet[Proposition~6.3 and Corollary~6.7]{RiehlVerity2014Reedy}. Let $J=I^\circ$ be the well-founded poset and well order its objects so that every strict predecessor is earlier. Write $j_!$ for the left adjoint of evaluation, so $(j_!A)_k=A$ if $j\leq k$ and is initial otherwise. A map $j_!a$, $a\in\Phi$, has relative latching map $a$ at $j$ and an identity elsewhere.

The class of maps whose relative latching maps lie in $\Phi$ is closed under pushouts, since forming a relative latching map takes a pushout of diagram maps to a pushout of that relative map. For transfinite closure, let $(X^\alpha)_{\alpha\leq\lambda}$ be a continuous chain of such maps. For fixed $j$ put
\[
 Q_\alpha=X^\alpha_j\amalg_{L_jX^\alpha}L_jX^\lambda.
\]
Each $Q_\alpha\to Q_{\alpha+1}$ is a pushout of the relative latching map of $X^\alpha\to X^{\alpha+1}$. The $Q$-chain is continuous because latching objects are colimits. Its starting object is the relative-latching domain of $X^0\to X^\lambda$, and $Q_\lambda=X^\lambda_j$. Transfinite closure of $\Phi$ therefore proves the assertion. Consequently this relative-latching class contains the transfinite saturation of all the maps $j_!a$.

For the reverse inclusion, take $f:X\to Y$ with all relative latching maps in $\Phi$. Starting from $X$, successively attach the value $Y_j$ at each object $j$. When $j$ is reached, all strict predecessors already equal their values in $Y$, and the current value at $j$ is
\[
 X_j\amalg_{L_jX}L_jY.
\]
Push out the free diagram map $j_!(X_j\amalg_{L_jX}L_jY\to Y_j)$ along its evaluation-adjoint map to the current diagram. This replaces the value at $j$ by $Y_j$ and leaves every earlier value fixed. At limit ordinals take colimits. The resulting composite ends at $Y$, since every component is eventually fixed to $Y_j$. Each step is a pushout of $j_!a$ for some $a\in\Phi$. Finally these one-object free maps belong to $v_!\Phi(I_0)$ by putting identity maps of the initial object in the other components; conversely a whole family is a coproduct of such one-object maps, formed by successive pushouts. This proves precisely the proposed transfinite-saturation description, without adding retract closure. It is a presentation by a class of cells, not an assertion that a small set of them suffices. If the required colimits have not been assumed, that existence qualification must remain explicit.

The easy inclusion in (3.5.5) is structural. Under the small-completeness hypothesis of 3.6, evaluation at an object $j$ has a right adjoint whose values are products indexed by the arrows into $j$ in $I^\circ$. Right lifting classes are closed under products, so this right adjoint sends a map of $F$ to an objectwise $F$-map. If $f\in{}^\perp F(I)$, adjunction therefore implies $f_j\in{}^\perp F=TC$. Hence projective $TC$-cofibrations are objectwise in $TC$, and therefore in $W(I)$; also $TC\subset C$ gives $(TC)\langle I\rangle\subset C\langle I\rangle$. The difficult inclusion says that an objectwise weak equivalence which already satisfies one coherent lifting condition satisfies the stronger one. Existence of abstract homotopy Kan extensions, established later, does not by itself identify these two specifically prescribed classes of strict diagrams.
\end{dcomment}


On a vu au paragraphe précédent qu'il en est ainsi si $(W,C)$ s'insère dans un triple de Quillen clos $(W,C,F=(TC)_*)$, et quand on suppose de plus que les couples
\[
C\langle I\rangle \ \longleftrightarrow\ TF(I) \ , \qquad TC\langle I\rangle \ \longleftrightarrow\ F(I)
\]
sont factorisant, auquel cas $(W(I),C\langle I\rangle,F(I))$ est une structure de Quillen close dans $\M(I)$. D'ailleurs, $(W,C,F)$ est une structure de Quillen (nécessairement close) si et seulement si les couples $C\longleftrightarrow \underbrace{TF}_{\deq C_*}$, $TC\longleftrightarrow \underbrace{\ F\ }_{\deq (TC)_*}$ eux-mêmes sont factorisants, et si % [sic] factorisant -> factorisants, agreement with « les couples »
\printedpage{23}
% head: A. Grothendieck : Dérivateurs, Ch. XIX, Édition des Universités de Montpellier II et Paris VII
% foot: Version 11 février 2019
de plus $TF\subset W$ (ce qui est automatique si $W=W^{\natural}$, car on aura en tous cas, par Cof$_4$, $TF\subset W^{\natural}$).

Mais peut-on s'en tirer avec moins, pour pouvoir conclure à (3.6.1) ? Ne suffirait-il pas de savoir que $TC\langle I\rangle\longleftrightarrow F(I)$ est factorisant (ce qui n'a guère de chances d'être vrai pour une $I$ $\pm$ générale, si on n'a déjà $TC\fact F$, factorisation dans $\M$ lui-même). C'est bien le cas :

\msp{page 41}
\textbf{Proposition 3.7.} \textit{Sous les conditions générales de \textup{(3.6)} sur $\M$, $W$, $C$, $I$, supposons que le couple $TC\langle I\rangle\longleftrightarrow F(I)$ (où $F\deq (TC)_*$) soit factorisant, i.e. $TC\langle I\rangle\fact F(I)$. Alors la formule \textup{(3.6.1)} est vraie.}

En effet, si $f\in C\langle I\rangle\cap W(I)$, on factorise $f$ en $pi$, avec $i\in TC\langle I\rangle$, $p\in F(I)$. Comme $f,i\in W(I)$, on aura $p\in W(I)$, i.e. $p\in TF(I)$, et à cause de $\underbrace{C\langle I\rangle}_{\ni\, f}\harp\underbrace{TF(I)}_{\ni\, p}$, il existe $p : Y\longrightarrow X'$ telle que $pq=\id_Y$, $qf=i$, % reading resolved: DER_XIX_corrections.csv; % [sic] p : Y -> q : Y, the arrow Y -> X' is q (see the diagram and pq = id_Y)
\[
\begin{tikzcd}[column sep=small, row sep=large]
X \arrow[rr, "f\in C\langle I\rangle\cap W(I)"] \arrow[dr, "i\in TC\langle I\rangle"'] & & Y \arrow[dl, dotted, shift right=1ex, "q"'] \\
& X' \arrow[ur, "p\in F(I)"'] \ , &
\end{tikzcd}
\]
donc $f$ est facteur direct de $i\in TC\langle I\rangle$. Comme $TC\langle I\rangle$ est stable par facteurs directs, on gagne.

\textbf{Corollaire 3.8.} \textit{Soit $(W,C,F)$ une structure de Quillen close dans $\M$, et supposons que $TC=C\cap W$ soit une partie Q-spéciale de $\M$, i.e. le couple $TC\fact F$ un couple spécial (i.e. $\M$ est \textup{spéciale}, et il existe une petite partie $TC_0$ Q-génératrice à gauche pour $TC$). Alors pour toute $I$ dans \textup{Cat}, on a}
\[
C\langle I\rangle\cap W(I) \ = \ TC\langle I\rangle \ . \tag*{(3.8.1)}
\]

\dgrblock{dgrb1183}{The lifting lemma inside the retract argument}
\begin{dcomment}{exposition}{p41:Q-génératrice à gauche pour $TC$}{The lifting lemma inside the retract argument}
Grothendieck has found that only one diagram factorization is needed for (3.6.1). The proof is sound, but the line identifying a weak equivalence in $F$ with a map lifting against $C$ uses a short lemma that is useful beyond the present proposition.

Suppose $(W,C)$ satisfies Cof$_0$--Cof$_3$, and set $F=(C\cap W)^\perp$. Then
\[
 F\cap W\subset C^\perp.
\]
To prove this, consider a square from $a:A\to B$ in $C$ to $p:X\to Y$ in $F\cap W$, with top map $h:A\to X$ and bottom map $k:B\to Y$. Form $Z=X\amalg_A B$, write $j:X\to Z$, and let $g:Z\to Y$ be induced by $p$ and $k$. Factor $g=q\ell$ by Cof$_3$, where $\ell:Z\to T$ is in $C$ and $q:T\to Y$ is in $W$. The composite $\ell j$ belongs to $C$, and $q\ell j=p$ shows by two-out-of-three that $\ell j\in TC$. Since $p\in F$, the square with left side $\ell j$, right side $p$, top $1_X$, and bottom $q$ admits $s:T\to X$ satisfying $s\ell j=1_X$ and $ps=q$. If $b:B\to Z$ is the other pushout map, then $s\ell b$ is the desired lift. This proves the inclusion without a second factorization in $\M$, and without retract closure of $W$.

Now take $f:X\to Y$ in $C\langle I\rangle\cap W(I)$ and factor it as $f=pi$ with $i\in(TC)\langle I\rangle$ and $p\in F(I)$. The already established easy inclusion makes $i$ an objectwise weak equivalence, hence $p\in F(I)\cap W(I)$. Applying the preceding lemma at every object gives $p\in TF(I)=C^\perp(I)$. The defining coherent lifting property of $f$ produces a natural map $q:Y\to X'$ with $qf=i$ and $pq=1_Y$. The arrow-category maps $(1_X,q)$ and $(1_X,p)$ exhibit $f$ as a retract of $i$. Left lifting classes are retract closed, so $f\in(TC)\langle I\rangle$.

The printed prose writes $p:Y\to X'$ for the new diagonal; the diagram and both identities show that $q$ is intended. The earlier reference on ms.~p.~40 to Cof$_4$ in the assertion $TF\subset W^{\natural}$ is likewise a reference slip: the relevant axiom is Cof$_3$, as in Proposition 3.1. Neither slip changes the argument. Corollary 3.8 is then exactly the small object theorem applied to $TC$ in diagram categories. It proves the equality without assuming that $C$ has a small generating set; the separate difficulty is obtaining the other factorization needed for a complete diagram model structure.
\end{dcomment}


\msp{page 42}
En apparence, on n'a pas eu à supposer que $C$ soit elle aussi une partie Q-spéciale de $\Fl(\M)$. Mais comme on a besoin de la factorisation pour $C\longleftrightarrow TF$ (pour avoir au moins une structure de cofibrations sur $\M$), on n'a guère le moyen de l'établir sans supposer $C$ spéciale. Sans compter que pour pouvoir expliciter (3.8.1) pour la construction de $u^{W}_{!}$, on a besoin de savoir que $(W(I),C\langle I\rangle)$ définit bien une structure à cofibrations sur $\M(I)$, ce qui ne peut guère se faire (il me semble) sans avoir la factorisation pour $C\langle I\rangle\longleftrightarrow TF(I)$ également ; chose que je n'arrive à établir pour un $I$ général dans Cat, que moyennant l'hypothèse de spécialité sur $C$ également.

\textbf{3.9.} Pour en avoir le cœur net, il faudrait quelques exemples concrets et illuminants, du type suivant :

a) \texttt{[Il n'y a pas de b).]} Structure de Quillen close $(W,C,F)$ sur $\M$, telle que pour le prédérivateur associé, $u^{W}_{!}$ n'existe pas toujours, ou $u^{W}_{*}$ n'existe par toujours. Je pense notamment au cas d'une flèche $u : I\longrightarrow I'$ avec $I'=e$ (présumant que si dans ce cas, $u^{W}_{!}$ resp. $u^{W}_{*}$ existe, alors il existe dans tous les cas). Un tel exemple serait encore plus % [sic] n'existe par toujours -> n'existe pas toujours, typing error

% ---- from der-ch19_p026-028.tex
\printedpage{24}
% head: A. Grothendieck : Dérivateurs, Ch. XIX, Édition des Universités de Montpellier II et Paris VII
% foot: Version 11 février 2019
\msp{page 43}
instructif, bien sûr, si $\mathcal{M}$ est stable par petites $\limind$ , que ses objets soient accessibles etc. et que $W$ soit L-accessible (i.e. stable par petites $\limind$ filtrantes, et de plus accessible \dots). Le premier candidat pour $I$ auquel je pense serait un $B_G$ ($G$ un groupe, p.ex. $G=\pm1$), c'est en effet l'exemple aux antipodes des ensembles ordonnés, lesquels, eux, ne semblent pas faire trop de problème (moyennant l'hypothèse noethérienne tout au moins).

Les exemples suivants ne sont pas aussi convaincants que ça, hélas!

\textbf{Exemple 3.10.} Soit $\mathcal{M}$ une catégorie quelconque, alors les couples $(\Fl(\mathcal{M}),\Is(\mathcal{M}))$ et $(\Is(\mathcal{M}),\Fl(\mathcal{M}))$ sont tous deux des couples de Quillen,
\[
\Fl(\mathcal{M}) \Qcouple \Is(\mathcal{M}) \ , \qquad \Is(\mathcal{M}) \Qcouple \Fl(\mathcal{M}) \ . \tag*{(3.10.1)} % reading resolved: DER_XIX_corrections.csv
\]
Il s'ensuit que si on pose
\[
W \ = \ \Is(\mathcal{M}) \ , \qquad C \ = \ F \ = \ \Fl(\mathcal{M}) \ , \tag*{(3.10.2)}
\]
$(W,C,F)$ est un triple de Quillen clos \textsuperscript{(}\footnote[22]{au moins si on admet que $\mathcal{M}$ est stable par $\limind$ et $\limproj$ \textit{finies.}}\textsuperscript{)}. D'ailleurs pour toute $I$ dans Cat, on a $W(I)=\Is(\mathcal{M}(I))$, donc
\[
\Ho_W(I) \ = \ \mathcal{M}(I) \ . \tag*{(3.10.3)}
\]
Donc la question de l'existence des foncteurs $u_!^W$, $u_*^W$ se réduit à celle des foncteurs $u_!$, $u_*$ habituels, adjoints à gauche et à droite de $u^* : \mathcal{M}(I') \longrightarrow \mathcal{M}(I)$.

\dgrblock{dgrb1184}{Finite completeness and the modern model category theorem}
\begin{dcomment}{status}{p43:Donc la question de l'existence des foncteurs}{Finite completeness and the modern model category theorem}
Example 3.10 isolates the completeness issue. Every map factors as itself followed by an identity, or as an identity followed by itself. A map with the lifting property against every map lifts against itself, and that diagonal is its inverse. Consequently the two displayed pairs really are closed factorization pairs in any category. With finite limits and colimits they give the historical, finitely complete version of Quillen's axioms. Their localization does nothing, so existence of homotopy Kan extensions is precisely existence of ordinary Kan extensions. For a functor between finite categories, the relevant comma categories are finite and the assumed finite limits and colimits suffice. For arbitrary indexing categories those assumptions say nothing about the required infinite constructions. The small-generation calculation for the corresponding example in sets is already given in the commentary to ms.~p.~32.

The terminology now usually includes \emph{all} small limits and colimits in the definition of a model category. Under that convention the requested example in 3.9 does not exist. \citet[Theorem~6.11]{Cisinski2003} proves the derivator assertion for an arbitrary Quillen model category with the required small limits; duality gives the assertion about colimits. Neither cofibrant generation nor local presentability is required. The distinction matters: this theorem supplies homotopy Kan extensions without asserting that the specific projective Quillen structure written in Question 3.6 exists on every diagram category. The bounded-complex example below therefore addresses the older finite-completeness convention, while Question 3.6 retains its more specific lifting-class content.

Grothendieck's suggestion that terminal projections should suffice anticipates the pointwise Kan-extension formulas. In the established derivator framework the value of a Kan extension is a homotopy limit or colimit over a comma category. For an arbitrary prederivator, however, adjoints to terminal projections alone do not supply the coherent pointwise construction or its base-change isomorphisms. Those are additional assertions proved by the model-category theorem, rather than a formal consequence of representability at the terminal category.
\end{dcomment}


\msp{page 44}
On sait bien que l'existence de ces foncteurs pour $u$ donné, fait appel à l'existence de certaines $\limind$ resp. $\limproj$ dans $\mathcal{M}$; p.ex. si $I'=e$, il est nécessaire et suffisant que les $\limind$ (resp. $\limproj$ ) de type $I^\circ$ soient représentables dans $\mathcal{M}$. Mais avec les hypothèses faites sur $\mathcal{M}$, il n'y a pas de problème si $u : I \longrightarrow I'$ est un foncteur entre catégories \textit{finies.} On va donner un exemple où il n'en est plus ainsi.

\textbf{Exemple 3.11.} Soit $k$ un anneau commutatif, et $\mathcal{M}=K^{\mathrm{b}}(k\text{-Ab})$ la catégorie des complexes de cochaînes bornés de $k$-modules, $W$ l'ensemble des quasi-isomorphismes, donc
\[
\Ho_W \ = \ D^{\mathrm{b}}(k\text{-Ab}) \ ,
\]
catégorie dérivée bornée de la catégorie abélienne des $k$-modules. On sait, en utilisant l'existence d'assez de projectifs dans $k$-Ab, que cette structure est associée à une structure de Quillen close $(W,C,F)$, décrite par \textsc{Quillen}
\msp{page 45}
dans Homotopical Algebra -- à cela près qu'il prend $D^-(k)$ au lieu de $D^{\mathrm{b}}(k)$. Peut-être, pour vérifier les axiomes, doit-on supposer que $k$ est de dimension cohomologique finie.

Prenons $I=B_G$, où $G$ est un groupe (p.ex. $G=\{\pm1\}$), et $u : I \longrightarrow e$. Alors $\mathcal{M}(I)$ est la catégorie des complexes bornés de $k[G]$-modules. On voit que $u_!^W$ resp. $u_*^W$ existe si et seulement si pour tout tel complexe borné, ou simplement pour un $k[G]$-module sans plus $M$, le complexe $L\Gamma^G(M)$ resp. $R\Gamma_G(M)$ appartient à $D^{\mathrm{b}}(k)$, ou ce qui revient au même,
\printedpage{25}
% head: A. Grothendieck : Dérivateurs, Ch. XIX, Édition des Universités de Montpellier II et Paris VII
% foot: Version 11 février 2019
si et seulement si on a
\[
\begin{aligned}
H_i(G,M) \ &= \ 0 \quad \text{pour } i \gg 0 \ ,\\
\text{resp.}\quad H^i(G,M) \ &= \ 0 \quad \text{pour } i \gg 0 \ ,
\end{aligned}
\]
pour tout $k[G]$-module $M$. Ce n'est pratiquement jamais le cas; quand on prend p.ex. $k=\mathbf{Z}$ et $G$ groupe fini, cela n'est le cas que si $G=e$. L'exemple le plus économique me paraît $G=\pm1$, $k=\mathbf{F}_2$.

\dgrblock{dgrb1185}{A finite group can force an unbounded resolution}
\begin{dcomment}{examples}{p45:L'exemple le plus économique me paraît}{A finite group can force an unbounded resolution}
The notation $K^{\mathrm b}(k\text{-Ab})$ is explicitly defined here as the category of bounded complexes and chain maps. It is what is now usually denoted $C^{\mathrm b}(k\text{-Mod})$, before quotienting by chain homotopy. The sentence about enough projectives needs the qualification which Grothendieck immediately adds: an unbounded projective resolution does not provide a bounded cofibrant replacement. It is enough for his example to take the proposed field $k=\mathbf F_2$, where every bounded acyclic complex is contractible.

Here is a complete verification of the finite Quillen structure in that case. Put $W=$ quasi-isomorphisms, $C=$ degreewise injections, and $F=$ degreewise surjections. A bounded contractible complex of vector spaces is a finite direct sum of disk complexes and is both projective and injective in the category of complexes. Thus a surjective quasi-isomorphism, whose kernel is contractible, has the right lifting property against every injection; dually an injective quasi-isomorphism has the left lifting property against every surjection. For a bounded complex $X$, the cone $I(X)=\operatorname{Cone}(1_X)$ is contractible and has a natural degreewise injection $X\to I(X)$. Dually $P(Y)=\operatorname{Cone}(1_Y)[-1]$ has a natural degreewise surjection $P(Y)\to Y$. For $f:X\to Y$ the factorizations
\[
 X\xrightarrow{(f,i_X)}Y\oplus I(X)\longrightarrow Y,
 \qquad
 X\longrightarrow X\oplus P(Y)\xrightarrow{(f,p_Y)}Y
\]
have types $(C,F\cap W)$ and $(C\cap W,F)$ respectively. Retract closure and two-out-of-three are immediate from degreewise linear algebra and cohomology. All objects and constructions remain bounded. Finite limits and colimits exist, whereas, for example, a coproduct of $k[n]$ over all integers cannot be bounded: its universal injections and retractions would force nonzero cohomology in every degree.

For $G=C_2$ write $A=k[G]=k[\epsilon]/(\epsilon^2)$, with $\epsilon=g-1$. Since $BG$ has one object, a diagram of bounded complexes is exactly a bounded complex of $A$-modules. There is no uniform-boundedness issue here; for a general category with infinitely many objects, an objectwise bounded diagram need not be bounded uniformly. Thus the relevant restriction functor is
\[
 p^*:D^{\mathrm b}(k)\longrightarrow D^{\mathrm b}(A),
 \qquad V\longmapsto V\text{ with trivial }G\text{-action}.
\]
The exact projective resolution
\[
 \cdots\xrightarrow{\epsilon}A\xrightarrow{\epsilon}A
 \xrightarrow{\epsilon}A\longrightarrow k\longrightarrow0
\]
is exact because $\ker(\epsilon)=\epsilon A$. Tensoring it with $k$, or applying $\operatorname{Hom}_A(-,k)$, kills every differential. Consequently
\[
 \operatorname{Tor}^{A}_n(k,k)=k
 \quad\hbox{and}\quad
 \operatorname{Ext}^{n}_A(k,k)=k\qquad(n\geq0).
\]
This proves nonexistence of the adjoints themselves, not merely failure of one proposed construction. If $R$ were right adjoint to $p^*$, then
\[
 H^n(Rk)=\operatorname{Hom}_{D(k)}(k[-n],Rk)
 \cong\operatorname{Hom}_{D(A)}(k[-n],k)=k
\]
for every $n\geq0$, contradicting boundedness. If $L$ were left adjoint, then
\[
 \operatorname{Hom}_{D(k)}(Lk,k[n])
 \cong\operatorname{Hom}_{D(A)}(k,k[n])=k
\]
for every $n\geq0$. Over a field, the left side vanishes outside a finite range for every bounded $Lk$, another contradiction.

The source's general boundedness criterion can also be justified. The unbounded derived adjunctions are tensoring with the trivial module and derived invariants. If they take every module to a bounded complex, finite truncation triangles show that they do so for every bounded complex. Their restrictions are then the required adjoints. Conversely, the shifts of $k$ detect all cohomology of the proposed right adjoint. For the left adjoint, test maps into shifts of an injective cogenerator of $k$-modules: these detect the homology of the derived tensor product and force its boundedness. For $k=\mathbf Z$ and a nontrivial finite group, choose a subgroup $C_p$ of prime order and induce its trivial $\mathbf F_p$-module to $G$. The cyclic periodic resolution and the induction/coinduction identities give nonzero homology and cohomology in arbitrarily large degrees; see \citet[Lecture~3]{BrownCohomologyLectures}. Thus the stated finite-group obstruction is valid. Its mechanism is the arbitrarily long composable strings in $BG$, despite the finiteness of its object and arrow sets.
\end{dcomment}


\msp{page 46}
Donc ici $u_!^W$ et $u_*^W$ n'existent pas, pour $I \longrightarrow e$ avec $I$ catégorie finie. Mais aussi $\mathcal{M}$ n'est pas stable par petites $\limind$ .

\textbf{Exemple 3.12.} Soit $\mathcal{M}$ une catégorie muni d'un localiseur $W$ fortement saturé, tel que $\mathcal{M} \longrightarrow \underbrace{\mathcal{M}W^{-1}}_{=\mathcal{H}}$ admette un adjoint -- à gauche disons. On a donc une situation de deux foncteurs adjoints % [sic] muni -> munie, agreement with « une catégorie »
\[
\begin{tikzcd}[column sep=huge]
\mathcal{H} \arrow[r, shift left, "v_!\ \text{pleinement fidèle}"] & \mathcal{M} \arrow[l, shift left, "v^*\ \text{localisation}"]
\end{tikzcd}
\ , \tag*{(3.12.1)}
\]
et $W=(v^*)^{-1}(\mathrm{Iso}(\mathcal{H}))$. On trouve que pour toute $I$ dans Cat, on a
\[
\mathrm{HO}_W(I) \ \simeq \ \mathcal{H}(I) \ \overset{\text{déf}}{=} \ \mathrm{Hom}(I^\circ,\mathcal{H}). \tag*{(3.12.2)} % [sic] HO_W -> Ho_W, notation Ho_W used in (3.10.3) and Exemple 3.11
\]
On peut donc songer à faire un exemple où dans $\mathcal{M}$ les petites $\limind$ et $\limproj$ existent, mais non dans $\mathcal{H}$, de sorte que les $u_!^W$ ou $u_*^W$ n'existent pas toujours. Mais ce semble difficile à construire, vu que dans la situation envisagé, on voit facilement que les petites $\limproj$ existent dans $\mathcal{H}$. Donc si dans $\mathcal{H}$ % [sic] envisagé -> envisagée, agreement with « la situation »
\msp{page 47}
il existe une petite famille cogénératrice (et pour ceci, il suffit qu'une telle famille existe dans $\mathcal{M}$), il y aura également des petites $\limind$ , en vertu de SGA 4 I 8.12.7. De même, elles existent si $\mathcal{H}$ est une sous-catégorie pleine de $\mathcal{M}$ \textit{stable par petites} $\limind$ , et il paraît pratiquement impossible de faire un exemple (3.12.1) sans que $\mathcal{H}$ ait cette propriété de stabilité. Sans compter qu'il y a même la perplexité suivante :

\marginpar{! \textbar}\textbf{Question 3.13.} \textit{Existe-t-il une $\mathfrak{U}$-catégorie (\,$\mathfrak{U}$ un univers donné) où les petites $\limind$ soient représentables, sans qu'il en soit de même des petites $\limproj$ ?} \textbf{N.B.} Une telle catégorie n'admet pas de petite famille génératrice, a fortiori elle n'est pas équivalente à une petite catégorie.

\dgrblock{dgrb1186}{A coreflective localization inherits both kinds of limits}
\begin{dcomment}{exposition}{p47:Une telle catégorie n'admet pas de petite famille génératrice}{A coreflective localization inherits both kinds of limits}
The proposed difficulty in Example 3.12 is removed by a formal adjunction argument. Write $L:\mathcal H\to\M$ for the fully faithful left adjoint and $R:\M\to\mathcal H$ for the localization. The unit is an isomorphism. The counit $\varepsilon_X:LRX\to X$ belongs to $W$, since $R\varepsilon_X$ is an isomorphism by the triangle identity. These facts hold objectwise in every diagram category. Therefore $L^I$ and $R^I$ descend to inverse equivalences
\[
 \M(I)[W(I)^{-1}]\simeq\mathcal H(I),
\]
which proves (3.12.2) directly and explains why the example becomes a question about ordinary completeness.

If $\M$ is complete and $D:J\to\mathcal H$ is small, then
\[
 R\bigl(\lim_J LD\bigr)
\]
is its limit in $\mathcal H$: applying $\operatorname{Hom}_{\mathcal H}(H,-)$ and the adjunction identifies the required universal property with the limit universal property in $\M$. If $\M$ is cocomplete, let $X=\operatorname{colim}_J LD$. Every cocone map $LD(j)\to X$ has a unique factorization through $\varepsilon_X:LRX\to X$, because its source lies in the image of $L$. These factorizations form a cocone, hence induce $g:X\to LRX$ with $\varepsilon_Xg=1_X$. Postcomposition with $\varepsilon_X$ is a bijection on maps from any object in the image of $L$. Applying this with source $LRX$ gives $g\varepsilon_X=1_{LRX}$. Thus $\varepsilon_X$ is an isomorphism and $X$ belongs to the essential image of $L$. In particular that image is automatically closed under all small colimits, and $R\operatorname{colim}LD$ is the colimit in $\mathcal H$.

Consequently a full coreflective subcategory of a bicomplete category is bicomplete. No cogenerating family is needed, and the envisaged counterexample cannot have the adjunction (3.12.1). The source's auxiliary observation about cogenerators is nevertheless correct: if $\{C_s\}$ conservatively detects arrows in $\M$ by $\operatorname{Hom}_{\M}(-,C_s)$, then $\{RC_s\}$ does so in $\mathcal H$, by adjunction and full faithfulness of $L$. The cited representability criterion, \citet[Expos\'e~I, 8.12.7]{SGA4I1972}, supplies another sufficient argument when that hypothesis is available; the colimit argument above shows why it is unnecessary here.
\end{dcomment}


\dgrblock{dgrb1187}{All small suprema without a largest ordinal}
\begin{dcomment}{examples}{p47:Une telle catégorie n'admet pas de petite famille génératrice}{All small suprema without a largest ordinal}
Question 3.13 has a direct affirmative answer, with its universe distinction essential. Let $\operatorname{Ord}_{\mathfrak U}$ have the ordinals belonging to $\mathfrak U$ as objects, with one arrow $\alpha\to\beta$ when $\alpha\leq\beta$. It is locally $\mathfrak U$-small: every Hom-set is empty or a singleton. For a $\mathfrak U$-small diagram its colimit is the supremum of the ordinals occurring in the diagram. This supremum belongs to $\mathfrak U$ by the universe axioms; for the empty diagram it is $0$. There is no terminal object, since $\alpha+1$ is again an object and is larger than $\alpha$. Thus all small colimits exist but even the empty limit does not. The category is necessarily large relative to $\mathfrak U$.

The note about generators is correct with the convention used in SGA~4. In \citet[Expos\'e~I, 7.1]{SGA4I1972}, a generating family means that the associated represented functors jointly \emph{reflect isomorphisms}. Given any small family of ordinals $\{\alpha_s\}$, put $\beta=\sup_s\alpha_s$. Every map
\[
 \operatorname{Hom}(\alpha_s,\beta)
 \longrightarrow\operatorname{Hom}(\alpha_s,\beta+1)
\]
is a bijection between singletons, but $\beta\to\beta+1$ is not an isomorphism. Hence there is no small generating family in this conservative sense. A merely faithful family of Hom-functors would be a different convention: every poset has at most one arrow with given endpoints, so faithfulness alone says almost nothing here. It would be incorrect either to reject Grothendieck's note using that weaker convention or to conclude from this example that a small cocomplete category need not be complete.
\end{dcomment}


\msp{page 48}
\textbf{Question 3.14.} \textsc{Anderson} prétend que pour le prédérivateur défini par une catégorie à cofibrations satisfaisant l'axiome Cof$_4$ transfini, les foncteurs $u_!^W$ existent. Sa démonstration est totalement canulée, néanmoins je n'arrive pas à faire de contre-exemple. Le résultat qu'il annonce serait-il vrai? Je ne connais pas même d'exemple d'un localiseur $W$, dans une catégorie $\mathcal{M}$ \textit{avec petites $\limind$ et petites $\limproj$ ,} tel que $u_!^W$ ou $u_*^W$ n'existent pas toujours -- mais je suis persuadé qu'il y a des exemples (même avec $\mathcal{M}$ une petite catégorie?). J'en suis moins convaincu si on suppose $W$ stable par $\limind$ filtrantes, et au besoin, de plus, que les éléments de $\mathcal{M}$ sont accessibles (plus éventuellement d'autres propriétés catégoriques anodines sur $\mathcal{M}$, telles l'existence d'une filtration
\printedpage{26}
% head: A. Grothendieck : Dérivateurs, Ch. XIX, Édition des Universités de Montpellier II et Paris VII
% foot: Version 11 février 2019
\msp{page 49}
cardinale etc.)

\dgrblock{dgrb1188}{The cofibration assertion is true under the stated continuity hypothesis}
\begin{dcomment}{later}{p48:Le résultat qu'il annonce serait-il vrai?}{The cofibration assertion is true under the stated continuity hypothesis}
The first question in 3.14 now has a positive answer. Modern cofibration-category theory separates this assertion from the existence of projective Quillen structures. In the presence of the initial object and pushouts in the standing framework, Cof$_4$ implies the coproduct and sequential-colimit axioms of an Anderson--Brown--Cisinski cofibration category. To see the coproduct assertion, well order a small family of cofibrant objects, attach its members one at a time by pushout from the initial object, and take colimits at limit stages. The resulting transfinite composite is the coproduct. Applied to a family of cofibrations, the same construction proves closure of cofibrations under coproducts; for acyclic cofibrations use the asserted transfinite closure of $TC$. Countable sequential composites are already a special case of Cof$_4$.

\citet[Lemma~1.6.2 and Theorem~10.3.2]{RadulescuBanu2009} gives the corresponding axioms and proves that every such cofibration category defines a strong derivator with all homotopy left Kan extensions. His Theorem~10.3.3 includes arbitrary Quillen model categories. A complementary formulation uses derivable categories: the dual of \citet[Corollary~6.21]{Cisinski2010} supplies the colimit theorem, while that corollary itself treats the limit direction in Cisinski's conventions. These results require neither a small set of generating cofibrations nor functorial factorization. They settle the existence question raised here, without retroactively validating the proof which Grothendieck criticizes.

The broader request in the same paragraph drops the cofibration structure entirely. It has a different answer, even when the underlying category and localizer have every finiteness and accessibility property explicitly mentioned here. The following finite example makes that distinction precise.
\end{dcomment}


\dgrblock{dgrb1189}{A finite complete lattice whose localization loses binary limits}
\begin{dcomment}{development}{p48:Le résultat qu'il annonce serait-il vrai?}{A finite complete lattice whose localization loses binary limits}
\textbf{An eight-object counterexample to the broader request in 3.14.}\ours
There is a finite complete and cocomplete category $P$, with a strongly saturated localizer $W$ stable under filtered colimits and accessible as a full subcategory of the arrow category, for which neither homotopy Kan extension along the discrete two-object category $D\to e$ exists. This elementary construction was prepared for this commentary. A targeted search for localizations of finite and Boolean lattices and failures of products or coproducts did not establish an attribution for this exact presentation. No priority claim is made. This is not a counterexample to the cofibration-category theorem just cited.

Let $P$ be the Boolean lattice of subsets of $\{1,2,3\}$, ordered by inclusion, and abbreviate subsets by their strings of elements. Put $a=1$, $b=12$, and let $W$ consist of the identities and $s:a\to b$. The two-out-of-three property holds: $a<b$ is a cover, so a factorization of its only nonidentity weak equivalence has an identity factor. The underlying category has all small limits and colimits, computed by intersections and unions.

We first compute the localization as a category, not just its object classes. Adjoin $s^{-1}:b\to a$. Any composable word containing two such inverses has the intervening forward path $a\to b$, which is $s$; it therefore reduces to a word with at most one inverse. Thus a morphism $x\to y$ is represented either by an inclusion or by
\[
 x\longrightarrow b\xrightarrow{s^{-1}}a\longrightarrow y.
\]
If both forms occur, then $x\subseteq b$, $a\subseteq y$, and $x\subseteq y$. Either $x\subseteq a$, or $x$ contains $2$ and therefore $b\subseteq y$. In the first case the inverse cancels on the source side; in the second it cancels on the target side. The two forms consequently represent the same arrow. The localization is thin, with relation
\[
 x\preceq y\quad\Longleftrightarrow\quad
 x\subseteq y\ \text{ or }\ (x\subseteq12\text{ and }1\subseteq y).
\]
This also verifies its universal property: for any functor $F$ inverting $s$, the new arrow must be sent to
$F(a\to y)F(s)^{-1}F(x\to b)$; the reductions just proved verify compositions and uniqueness. Only $a$ and $b$ become isomorphic. A skeleton $Q$ therefore has seven objects, and $W$ is exactly the inverse image of its isomorphisms, hence is strongly saturated.

In $Q$ the common upper bounds of $2$ and $3$ are $13$, $23$, and $123$. The first two are incomparable, so no least upper bound exists. Dually the common lower bounds of $13$ and $23$ are $\emptyset$, $2$, and $3$, and there is no greatest one. Thus $Q$ lacks both a binary coproduct and a binary product. Localization of the product relative category $(P\times P,W\times W)$ gives $Q\times Q$: one may first invert $(w,1)$ and then $(1,w)$, and these arrows generate the same inversions as $W\times W$. The restriction along $D\to e$ consequently identifies with the diagonal $Q\to Q\times Q$. Its left and right adjoints would give exactly the missing binary coproducts and products. Both therefore fail to exist.

Finally, all the accessibility qualifications can be checked directly. A filtered diagram in a finite poset has a finite directed image, hence a largest value which occurs at some stage and is its colimit. Every object of $P$ is therefore finitely presentable. The same observation applies to the finite arrow category $P^{[1]}$. A filtered diagram all of whose objects lie in $W$ has its colimit equal to one of those objects, hence in $W$. The full subcategory on $W$ is itself finitely accessible: it has filtered colimits, every object is finitely presentable, and its finite set of objects supplies a generating set. The inclusion preserves filtered colimits. A cardinal filtration is trivial for this finite category. Thus the counterexample persists under the explicit additional hypotheses in the final sentences of 3.14. What fails is the extra homotopical structure supplied by the cofibration axioms, not any smallness or ordinary completeness condition.
\end{dcomment}


\textbf{Question 3.15.} Cela me ramène à la question particulièrement élusive et irritante, de l'existence des petits ensembles Q-générateurs, dans une partie $\Phi$ de $\Fl(\mathcal{M})$, quand on suppose (disons) que $\mathcal{M}$ est stable par petites $\limproj$ , que ses éléments sont accessibles et qu'il admet une filtration cardinale. Ainsi on a la question : une partie L-accessible de $\Fl(\mathcal{M})$, qu'on suppose au besoin Q-saturé à gauche (ou ce qui revient au même, contenant les isomorphismes, et stable par cochangement de base, par composition et par facteurs directs) admet-elle une \textit{petite partie} $\Phi_0$ Q-génératrice à gauche? % [sic] Q-saturé -> Q-saturée, agreement with « une partie »

\dgrblock{dgrb1190}{An accessible class with no small Q-generating subset}
\begin{dcomment}{development}{p49:Q-génératrice à gauche?}{An accessible class with no small Q-generating subset}
\textbf{A negative answer to the Q-saturated form of Question 3.15.}\ours
Let $\mathcal K$ be the category of directed multigraphs, with sets of edges and vertices and two functions $s,t:E\to V$, and let $\Phi$ be its class of pure monomorphisms. Then $\mathcal K$ is bicomplete, every object is accessible, and cardinality supplies a cardinal filtration. The full subcategory of $\mathcal K^{[1]}$ on $\Phi$ is accessible and is closed under all filtered colimits. The class $\Phi$ contains the isomorphisms and is closed under pushouts, transfinite compositions, and retracts, yet it has no small Q-generating subset. The non-generation theorem is due to \citet[Theorem~5.1]{CoxEtAl2026Pure}; the detailed translation to the hypotheses of this manuscript is the development supplied here. We make no priority claim for that translation.

Write $\mathcal K=\mathbf{Set}^{\mathcal E}$, where $\mathcal E$ has objects $e,v$, the distinct arrows $s,t:e\to v$, and identities. Limits and colimits are computed on the two sets, and finite graphs are finitely presentable. More generally a graph of cardinality less than a regular infinite cardinal $\lambda$ is $\lambda$-presentable: maps out of it and equalities between such maps involve fewer than $\lambda$ elements and equations, all of which occur at a common stage of a $\lambda$-filtered diagram. Thus every graph is presentable at some rank, and the subcategories with bounded cardinality are essentially small.

For an inclusion $A\subseteq B$, purity means that every finite system of equations in vertices and edges, with parameters from $A$, which has a solution in $B$, has a solution in $A$. Equivalently, in a commutative square from an arrow $P\to Q$ of finitely presentable graphs to $A\to B$, the top map $P\to A$ extends to $Q\to A$; equality with the given bottom map into $B$ is not required. These are the equivalent descriptions in \citet[Definition~2.7 and Remark~2.8]{CoxEtAl2026Pure}.

We verify accessibility of the \emph{full arrow subcategory}, whose maps are arbitrary commuting squares. This is distinct from the category with graphs as objects and only pure maps as morphisms. Given a pure inclusion $A\subseteq B$ and countable subsets in its domain and codomain, first close them under $s,t$ and ensure that the chosen part of $A$ is included in the chosen part of $B$. Suppose countable subgraphs $A_n\subseteq A$ and $B_n\subseteq B$, with $A_n\subseteq B_n$, have been chosen. There are only countably many finite equation systems with parameters in $A_n$. For each such system that has a solution in $B_n$, choose a solution in $A$, using purity of $A\subseteq B$. Adjoin all these witnesses to $A_n$, close under $s,t$, and adjoin their images to $B_n$. Repeat for $n<\omega$. The unions $A'\subseteq B'$ are countable. Any finite system over $A'$ with a solution in $B'$ has all its parameters and witnesses in a common stage, and was therefore given a solution in $A'$ at the next stage. Hence $A'\to B'$ is pure.

The poset of countable pure subarrows of $A\to B$ is $\aleph_1$-directed: for a countable collection take the unions of their domains and codomains as seeds and apply the preceding closure construction. Its union is the whole arrow, since any chosen vertex or edge can be included in the seeds. There is a set of isomorphism types of countable graph arrows. Each is $\aleph_1$-presentable in $\mathcal K^{[1]}$, by the same element-and-equation argument already used for graphs, with the two components of the arrow included as additional functions. In particular every pure arrow is a filtered colimit of a fixed small family of bounded-size pure arrows, exactly the filtered-colimit-envelope condition used by the manuscript's earlier accessibility language.

We must also check preservation of the relevant colimits. Let $f_i:A_i\to B_i$ be a filtered diagram of pure arrows with arbitrary commuting-square transition maps. A square from finite $P\to Q$ into its colimit factors through a common stage, and its required equality holds at some later stage, because $P,Q$ are finitely presentable. Purity at that stage gives $Q\to A_i$ extending the top map; composing into the colimit proves purity of the colimit arrow. Filtered colimits of sets preserve injections, so it is again a monomorphism. Thus the full subcategory on $\Phi$ has filtered colimits created in $\mathcal K^{[1]}$. Its countable arrows remain $\aleph_1$-presentable in that full subcategory, and the preceding directed presentation proves that it is $\aleph_1$-accessible. This verifies both parts of the manuscript's L-accessibility requirement, rather than inferring them from accessibility of a different category of pure maps.

For completeness, the closure conditions on $\Phi$ are also those required in the question. Composition follows by solving the same finite system first in the intermediate graph and then in the source. A retract of a pure arrow is pure: carry a parameter system and its solution along the section of the retract diagram, solve it in the larger source, and carry the solution back along the retraction. Arbitrary pushouts of pure monomorphisms are pure, as recorded in \citet[Fact~2.9(iii)]{CoxEtAl2026Pure}. Finally transfinite closure follows from composition at successor stages and the filtered-colimit argument at limit stages, applied to the arrows from the fixed starting object. Thus $\Phi$ is Q-saturated in the sense of the question, and in fact has the additional transfinite stability used elsewhere in the chapter.

The cited classification says that pure monomorphisms in $\mathbf{Set}^{\mathcal E}$ are cofibrantly generated precisely when, for every span $y\leftarrow x\to z$ in $\mathcal E$, one leg factors through the other. The span $v\xleftarrow{s}e\xrightarrow{t}v$ fails this test: the only endomorphism of $v$ is its identity, and $s\ne t$. Hence $\Phi$ is not cofibrantly generated.

To connect that conclusion exactly to Grothendieck's terminology, suppose a small subset $S\subseteq\Phi$ were Q-generating. Its defining property is
\[
 \Phi\subseteq{}^\perp(S^\perp).
\]
The small object argument applies in the locally presentable category $\mathcal K$. It identifies ${}^\perp(S^\perp)$ with $\operatorname{cof}(S)$, the class of retracts of transfinite composites of pushouts of maps in $S$. The closure properties just proved give $\operatorname{cof}(S)\subseteq\Phi$. Both inclusions force equality, contradicting the classification theorem. This proves the asserted negative answer with no additional set-theoretic principle. The example concerns Q-\emph{saturation}; it does not silently strengthen the question to the different hypothesis that $\Phi$ is already Q-\emph{closed}. Nor does the argument identify accessibility of an arrow class with combinatoriality of a model structure, since no model structure has been supplied.
\end{dcomment}


\marginpar{! \textbar}D'autre part, je n'ai pas à présent construit un seul exemple d'un couple $\Phi \longleftrightarrow \Psi$, dans une catégorie $\mathcal{M}$, qui ne soit factorisant! Le cas où $\mathcal{M}$ est spéciale, était bien sûr le plus intéressant. Visiblement, dans tout
\msp{page 50}
ça ce sont les exemples qui manquent. Le clef d'une compréhension semble bien être dans les questions de factorisation entre ensembles $\Phi \longleftrightarrow \Psi$. Il faudrait en faire une étude systématique, mais je recule devant la perspective d'une digression qui me prendrait peut-être encore des semaines! % [sic] Le clef -> La clef, « clef » is feminine


\dgrblock{dgrb1191}{Mutual lifting classes need not factor every arrow}
\begin{dcomment}{examples}{p50:Il faudrait en faire une étude systématique}{Mutual lifting classes need not factor every arrow}
The failure of a mutually orthogonal pair to admit factorizations is already illustrated by \citet[Example~5.7.10]{Taylor1999PracticalFoundations}. The following rational-cut presentation gives a small elementary instance, with no priority claim; it does not meet the stronger condition that the ambient category be special. Regard $\mathbf Q$ as a category under its usual order, and fix an irrational real number $r$, for example $\sqrt2$. Define $\mathcal L$ to contain the identities and the nonidentity arrows $a\to b$ with $b<r$. Define $\mathcal R$ to contain the identities and the nonidentity arrows $x\to y$ with $r<x$.

In a poset a lifting problem from $a\to b$ to $x\to y$ has a solution exactly when $b\leq x$, provided the square exists. Nonidentity arrows in the displayed classes have $b<r<x$, so $\mathcal L\perp\mathcal R$. Conversely, if a nonidentity $a\to b$ is not in $\mathcal L$, then $b>r$. Choose a rational $x$ with $\max(a,r)<x<b$. The square from $a\to b$ to $x\to b\in\mathcal R$, with bottom identity, has no lift. Hence ${}^\perp\mathcal R=\mathcal L$. If a nonidentity $x\to y$ is not in $\mathcal R$, then $x<r$. Choose rational $b$ with $x<b<\min(y,r)$. The square from $x\to b\in\mathcal L$ to $x\to y$, with top identity, has no lift. Thus $\mathcal L^\perp=\mathcal R$.

Nevertheless, an arrow $a\to b$ with $a<r<b$ has no factorization of the indicated type. An intermediate object $z$ with both factors nonidentity would have to satisfy simultaneously $z<r$ and $r<z$; choosing either factor to be an identity also fails. This is a pair of mutually closed lifting classes, but not a weak factorization system. The rational poset is not cocomplete, so the example answers the unrestricted request only. It leaves the final, more demanding special-ambient question logically separate from the counterexample.
\end{dcomment}
\begin{otherlanguage}{english}
\dgroverview{XIX}{Chapter XIX \textperiodcentered\ \textup{XIX} \textperiodcentered\ Overview}
\begingroup
\selectlanguage{english}
\noindent\textbf{Chapter XIX: factorization, diagram categories, and the existence of derived Kan extensions.}



\medskip
Grothendieck has reached the point at which his formalism must acquire its
examples. Given a category $\mathcal M$ and a class $W$ of weak equivalences,
the intended coefficient system is
\[
 \mathbb D_{\mathcal M,W}(I)
   =\operatorname{Fun}(I^{\mathrm{op}},\mathcal M)[W(I)^{-1}],
 \qquad u^*:\mathbb D(J)\longrightarrow\mathbb D(I).
\]
The problem is to construct both adjoints to restriction, prove their pointwise
formulas, and explain which hypotheses on the models are doing the work. The
chapter's fifty manuscript pages therefore move from lifting and factorization
to model structures on diagrams, and then return to the places where that
method appears to stop. In the edition the chapter occupies printed pages 1--26,
after two preliminary leaves.

The historical motivation precedes the terminology of this chapter. In
\emph{Pursuing Stacks}, \S69, Grothendieck relates the passage from a single
derived category to a system of derived diagram categories to the inadequacy
of the former for the variance and trace constructions arising in SGA~5.
His letter to Thomason of 2 April 1991 connects the long investigation of
accessible categories in Chapter XVIII with the foundations needed for this
coefficient formalism \citep{Grothendieck1983,Grothendieck1991Letter}. Thus
accessibility enters here for a concrete reason: the homotopical operations
must be constructed using sets of generators and controlled colimits before
they can be recognized intrinsically after localization.

\noindent\textbf{The two levels of closure.}
The opening pages isolate the relation $i\mathbin{\perp}p$ of admitting a
diagonal in every lifting square. It gives an antitone correspondence between
classes of arrows, hence the closures
${}^{\perp}(S^{\perp})$ and $({}^{\perp}T)^{\perp}$. Pushouts, transfinite
compositions, and retracts describe a second closure process. The small
object argument identifies the two when a set of arrows generates the
lifting condition and their domains have the required smallness. The
retract argument is the decisive step: factor an arrow already having the
left lifting property, then use its lift against the right factor to express
it as a retract of the constructed cellular arrow. This is why the distinction
between an orthogonal pair and a factorizing pair must remain visible.
The source's initial identification of separate closedness of two orthogonal
classes with mutual determination of the classes is too strong; the
factorization hypothesis supplies the missing implication later.

Grothendieck's word ``special'' must also retain its stated meaning. His
ambient category has small colimits and each object is presentable at some
cardinal; his special left classes have a small set of lifting generators.
This does not itself supply the small generating subcategory required for
local presentability. Modern statements about combinatorial model categories
may consequently be used only after the additional ambient hypothesis has
been made explicit.

\noindent\textbf{The asymmetry of projective and injective diagrams.}
For contravariant diagrams, evaluation at $i$ has a left adjoint
$i_!A(j)=\coprod_{\operatorname{Hom}_I(j,i)}A$. Applying these adjoints to
generating cofibrations and generating trivial cofibrations constructs the
projective structure. Its fibrations and weak equivalences are pointwise;
its cofibrations encode compatibility with the arrows of the indexing
category. Theorem 2.3 then derives $u_!$, and Proposition 1.12 supplies the
restriction property needed to compute its values on comma categories.
In the variance used here the formulas are
\[
 (\mathbf L u_!X)_j\simeq
   \operatorname*{hocolim}_{(j\downarrow u)^{\mathrm{op}}}X,
 \qquad
 (\mathbf R u_*X)_j\simeq
   \operatorname*{holim}_{(u\downarrow j)^{\mathrm{op}}}X.
\]
Writing the opposite categories is necessary: the manuscript's diagrams
are presheaves, whereas much modern derivator notation uses covariant
diagrams.

For the injective structure the cofibrations are pointwise, and it is the
right class whose compatibility is no longer pointwise. One cannot obtain
small generators merely by placing the old generators at every component.
Question 1.11 has a positive answer for locally presentable ambient
categories, by the generation theorem of \citet[Lemma A.2.8.3]{Lurie2009HTT},
but its final suggested choice of generators already fails for sets with a
nontrivial group action. The commentary gives the sharper elementary
criterion for when the sole generator $\varnothing\to 1$ can be used unchanged
on a diagram category, with a complete proof and an explicit statement of
the literature checked.

Theorems 2.6 and 2.9 secure the injective construction in Grothendieck's
``very special'' situation, in which the cofibrations are all monomorphisms.
The geometric conditions on unions and pushouts, together with accessibility
of $W$, are the mechanism behind the recognition theorem. Here two-out-of-three
does essential work: after building a bounded trivial subextension, it keeps
the remaining extension in $W$. This is a substantially stronger situation
than an arbitrary accessible class of monomorphisms. The proof recalled
from Chapter XV must be read through its corrected argument in \S7, not
through the explicitly abandoned general claim in \S6. The cardinal-stage
condition printed with a projective-limit arrow in XIX is identified in the
apparatus as a source slip: the cited formulations require closure under
the corresponding small inductive limits.

The latching interpretation in section 3 returns to the same care about
testing classes. Proposition 3.5 needs the canonical right class, which the
source explicitly chooses in (3.5.2); an arbitrary presentation by right
tests can give a false equality. The commentary supplies the counterexample
and the valid comparison. It also proves the classwise cellular description
conjectured in (3.5.3) for noetherian posets, with the required colimits stated,
and supplies the short lifting lemma implicit in Proposition 3.7.

\noindent\textbf{What later existence theorems change.}
The existence of a convenient diagram model structure and the existence of
derived Kan extensions are different assertions. Cisinski proves the latter
for an arbitrary Quillen model category with small limits; duality supplies
the colimit statement when small colimits exist
\citep[Theorem 6.11]{Cisinski2003}. Consequently the bicomplete Quillen
case of Question 3.9 admits no counterexample of the requested kind.
For combinatorial models, both diagram model structures give a direct and
symmetric proof; \citet[Theorem 1.28 and Proposition 1.30]{Groth2013} presents
this route. The later existence theorem does not retroactively prove every
factorization assertion contemplated in the manuscript, and it does not
make the adjective ``very special'' necessary.

The final section shows why the size of the permitted diagrams matters.
Its bounded-complex example is effective precisely because bounded derived
categories need not contain the result of a homotopy limit indexed by a
finite category with arbitrarily long composable strings. For
$k=\mathbf F_2$ and $G=C_2$, the resolution of the trivial $k[G]$-module is
periodic and gives nonzero group homology and cohomology in every
nonnegative degree. Testing a hypothetical adjoint to the trivial-action
functor against shifts proves its nonexistence within the bounded category.
This supplies the required obstruction without identifying a bounded
category with a bicomplete modern model category. The distinction between
cofibration-category excision on cofibrant objects and the manuscript's
absolute $W$-cofibrations is equally necessary in reading the Brown lemma
invoked on manuscript page 35.

Some of the closing searches have elementary answers. A fully faithful
coreflective image in a bicomplete category is itself bicomplete: the
coreflector constructs limits, and the image is closed under colimits.
Thus the proposed adjoint-localization construction in Example 3.12 cannot
produce its hoped-for failure. The category of ordinals small relative to
the fixed universe has all small suprema and no terminal object, answering
Question 3.13 with the universe convention made explicit. Its lack of a
small conservative generating family is consistent with the convention of
SGA~4. A rationally ordered category cut at an irrational number supplies
a mutually lifting-closed pair which does not factor arrows across the cut;
it does not have the cocompleteness required of a special ambient category.

The broader localizer question in 3.14 admits a small counterexample.
Localize the Boolean lattice of subsets of $\{1,2,3\}$ at the single cover
$\{1\}\subset\{1,2\}$. The resulting seven-object skeleton has neither
a least upper bound for $\{2\},\{3\}$ nor a greatest lower bound for
$\{1,3\},\{2,3\}$. The localizer is strongly saturated, accessible, and
stable under filtered colimits, and the original lattice is complete and
cocomplete. The commentary computes the localization and proves each
assertion. This answers the broad question even with the additional
smallness hypotheses. The cofibration-category question in the same
paragraph has the opposite answer: the transfinite continuity axiom
supplies the coproduct and sequential-colimit hypotheses of modern
cofibration-category theory, which yields all homotopy left Kan
extensions \citep[Theorem 10.3.2]{RadulescuBanu2009}.

Question 3.15 also benefits from separating the closures introduced at the
start. The class of pure monomorphisms in directed multigraphs is accessible,
closed under filtered colimits, pushouts, transfinite composition, and
retracts. It nevertheless has no small set of Quillen generators. The
non-generation theorem is due to Cox, Feigert, Kamsma, Mazari-Armida, and
Rosick\'y, in their 2026 classification for presheaf categories
\citep[Theorem 5.1]{CoxEtAl2026Pure}; the
commentary verifies the translation into the hypotheses stated here and
supplies the small-subdiagram argument for accessibility. This answers the
question for a Q-saturated class. It does not identify that class with its
double lifting orthogonal, and the stronger question obtained by imposing
that additional condition is not silently substituted for the question
Grothendieck actually wrote.

\medskip
The commentary retains the edition's slips as source evidence, distinguishes
them from errors of transcription, and states repairs at the
point where they are needed. The claim register records the source's proof
status separately from later theorems or counterexamples. Mathematical
dependencies on Chapters XIII, XV, XVII, and XVIII are checked against
the relevant original passages or explicitly identified as recalled results.

\endgroup
\end{otherlanguage}
\dgrchunkend

\backmatter
\dgrbackchapter{Register of claims}
The audit of each part records the definitions, axioms, propositions, claims, conjectures and questions of the text, each with the status the text gives it: proved, sketched, stated, assumed, conjectured, false (the assertion, or the inference as written, is false) or open. The table counts the entries of each part by status; the list after it gives every entry whose status is false, open or conjectured, with its manuscript page and its identifier in the audit.

\begin{center}\small
\begin{tabular}{lrrrrrrrr}
\toprule
Part & proved & sketched & stated & assumed & conjectured & false & open & total \\
\midrule
I-a & 53 & 24 & 136 & 19 & 7 & 14 & 0 & 253 \\
I-b & 72 & 37 & 101 & 22 & 9 & 12 & 0 & 253 \\
II-a & 50 & 21 & 145 & 1 & 5 & 20 & 3 & 245 \\
III-V & 54 & 27 & 128 & 10 & 11 & 9 & 10 & 249 \\
VI-a & 88 & 29 & 137 & 19 & 16 & 23 & 0 & 312 \\
VI-b & 115 & 18 & 59 & 17 & 9 & 15 & 17 & 250 \\
VI-c & 88 & 27 & 72 & 9 & 8 & 5 & 5 & 214 \\
VII-a & 92 & 19 & 74 & 2 & 0 & 12 & 0 & 199 \\
VII-b & 54 & 16 & 82 & 3 & 5 & 17 & 7 & 184 \\
VII-c & 70 & 26 & 62 & 4 & 30 & 35 & 11 & 238 \\
VIII-X & 11 & 19 & 104 & 5 & 24 & 12 & 10 & 185 \\
XI-XIIa & 124 & 35 & 90 & 43 & 12 & 9 & 11 & 324 \\
XIIb-a & 91 & 25 & 70 & 7 & 0 & 16 & 12 & 221 \\
XIIb-b & 74 & 11 & 36 & 4 & 3 & 4 & 1 & 133 \\
XIIc-a & 51 & 21 & 89 & 14 & 11 & 9 & 19 & 214 \\
XIIc-b & 63 & 20 & 31 & 4 & 3 & 5 & 1 & 127 \\
XIIapp & 11 & 31 & 73 & 4 & 4 & 17 & 5 & 145 \\
XIII-a & 56 & 19 & 172 & 27 & 44 & 12 & 17 & 347 \\
XIII-b & 33 & 13 & 63 & 17 & 5 & 11 & 9 & 151 \\
XIV & 18 & 6 & 27 & 8 & 2 & 4 & 0 & 65 \\
XV-a & 88 & 23 & 66 & 17 & 1 & 28 & 2 & 225 \\
XIX & 62 & 17 & 148 & 7 & 4 & 11 & 20 & 269 \\
\midrule
total & 1418 & 484 & 1965 & 263 & 213 & 300 & 160 & 4803 \\
\bottomrule
\end{tabular}
\end{center}

\section*{I-a \textperiodcentered\ Chapter I, part 1 of 2, manuscript pages 1--66}
\pdfbookmark[1]{I-a}{dgrclaims.I-a}
\begingroup\raggedright
% DER_I-a: the definitions of the chunk's preamble, applied inside its group
\let\printedpage\relax\newcommand{\printedpage}[1]{\marginpar{\footnotesize p.~#1}}
\let\msp\relax\newcommand{\msp}[1]{\marginpar{\footnotesize [#1]}}
\let\AMSTeX\relax\newcommand{\AMSTeX}{$\mathcal{A}\kern-.1667em\lower.5ex\hbox{$\mathcal{M}$}\kern-.125em\mathcal{S}$-\TeX}
\let\uHom\relax\newcommand{\uHom}{\underline{\mathsf{Hom}}}
\let\Ob\relax\newcommand{\Ob}{\mathsf{Ob}}
\let\Ab\relax\newcommand{\Ab}{\mathsf{Ab}}
\let\bDelta\relax\newcommand{\bDelta}{\boldsymbol{\Delta}}
\let\Dia\relax\newcommand{\Dia}{\mathsf{Dia}}
\let\Cat\relax\newcommand{\Cat}{\mathcal{C}\mathit{at}}
\let\DD\relax\newcommand{\DD}{\mathbb{D}}
\let\Hom\relax\newcommand{\Hom}{\underline{\mathsf{Hom}}}
\let\id\relax\newcommand{\id}{\mathsf{id}}
\let\Del\relax\newcommand{\Del}{\boldsymbol{\Delta}}
\let\Ens\relax\newcommand{\Ens}{\mathcal{E}ns}
\let\defeq\relax\newcommand{\defeq}{\overset{\text{déf}}{=}}
\let\D\relax\newcommand{\D}{\mathbb{D}}
\let\sHom\relax\newcommand{\sHom}{\mathsf{Hom}}
\let\DER\relax\newcommand{\DER}{\mathsf{DER}}
\let\sIm\relax\newcommand{\sIm}{\mathsf{Im}}
\let\sDia\relax\newcommand{\sDia}{\mathsf{Dia}}
\let\sOb\relax\newcommand{\sOb}{\mathsf{Ob}}
\let\sid\relax\newcommand{\sid}{\mathsf{id}}
\let\HOT\relax\newcommand{\HOT}{\mathsf{HOT}}
\let\cA\relax\newcommand{\cA}{\mathcal{A}}
\let\hot\relax\newcommand{\hot}{\mathsf{hot}}
\let\Hot\relax\newcommand{\Hot}{\mathsf{Hot}}
\let\isoto\relax\newcommand{\isoto}{\overset{\sim}{\to}}
\let\isolongto\relax\newcommand{\isolongto}{\overset{\sim}{\longrightarrow}}

\dgrclaim{conjectured}{3}{I-a-E020}{For a monotone map of finite ordinals, restriction on A-valued presheaves has a second adjoint on each side when A is pointed and has finite limits and colimits.}
\dgrclaim{false}{3}{I-a-E021}{All restriction functors associated with all maps between two fixed finite ordinals occur in the same successive adjunction chain.}
\dgrclaim{conjectured}{3}{I-a-E022}{For [n] and [n+1], the face and degeneracy restrictions form a chain of 2n+3 terms, extendible by two terms at each end to 2n+7.}
\dgrclaim{false}{4}{I-a-E024}{A category in which every arrow is monic and every isomorphism is an identity is a poset.}
\dgrclaim{false}{15}{I-a-E066}{The abelianization morphism preserves all right Kan extensions.}
\dgrclaim{false}{15}{I-a-E067}{A homotopy pullback of spaces along a fibration is sent by abelianization to a homotopy pullback.}
\dgrclaim{false}{25}{I-a-M007}{The mixed interchange map might be invertible in additive derivators.}
\dgrclaim{conjectured}{29}{I-a-M021}{There should be a stable or triangulated derivator theory in which left and right exact morphisms coincide.}
\dgrclaim{false}{32}{I-a-M030}{The withdrawn Der5 asserts that a strict hot-equivalence has fully faithful restriction, equivalently $1\to f_*f^*$ is invertible.}
\dgrclaim{false}{32}{I-a-M032}{The dual withdrawn Der5 asserts $f_!f^*\to1$ invertible for every costrict hot-equivalence.}
\dgrclaim{false}{32}{I-a-M033}{The withdrawn argument would make $f_*$, respectively $f_!$, localization functors.}
\dgrclaim{false}{47}{I-a-M060}{The same cohomology invariance would follow if instead all $(\alpha/x)$ were contractible.}
\dgrclaim{false}{47}{I-a-M062}{The terminal-object condition above would substantially strengthen cohomological cofinality.}
\dgrclaim{false}{47}{I-a-M064}{The condition $\alpha_!\varepsilon_{X'}\simeq\varepsilon_X$ follows from either orientation of cofinality.}
\dgrclaim{false}{47–48}{I-a-M067}{Corollary2 repeats the claim that strict hot-equivalence implies $F\to f_*f^*F$ invertible for arbitrary $F$.}
\dgrclaim{false}{48}{I-a-M071}{The preceding slice calculation proves the withdrawn full-faithfulness corollary by opposite cofinality.}
\dgrclaim{false}{49}{I-a-L003}{The same sufficient base-change criterion holds with $h_{y'}$ homotopy aspherical (all $h_{y'}/(x,u)$ contractible).}
\dgrclaim{conjectured}{58}{I-a-L035}{If right base change for a fixed f holds for every base change and every derivator, then f is proper.}
\dgrclaim{conjectured}{58}{I-a-L036}{If base change along a fixed p commutes with right direct images for every f and every coefficient theory, then p is smooth.}
\dgrclaim{conjectured}{58}{I-a-L037}{Testing HOTAB and HOT should suffice to detect the universal proper/smooth base-change properties.}
\dgrclaim{conjectured}{59}{I-a-L039}{For f, the four conditions universal right base change, universal left-mate base change, coasphericity of X\_y->X/y, and contractibility of every fibre of x\textbackslash{}X->f(x)\textbackslash{}Y are equivalent.}
\par\endgroup

\section*{I-b \textperiodcentered\ Chapter I, part 2 of 2, manuscript pages 67--126}
\pdfbookmark[1]{I-b}{dgrclaims.I-b}
\begingroup\raggedright
% DER_I-b: the definitions of the chunk's preamble, applied inside its group
\csname tagsleft@true\endcsname
\let\printedpage\relax\newcommand{\printedpage}[1]{\marginnote{\scriptsize\sffamily ed.\ #1}}
\let\msp\relax\newcommand{\msp}[1]{%
  \begingroup\def\thismsp{#1}\def\blankmsp{page 114 vide}%
  \ifx\thismsp\blankmsp
    \marginnote{\scriptsize\sffamily [114: vide]}[1.3\baselineskip]%
  \else\def\followingmsp{page 115}%
    \ifx\thismsp\followingmsp
      \marginnote{\scriptsize\sffamily [page 115]}[2.6\baselineskip]%
    \else\marginnote{\scriptsize\sffamily [#1]}[1.3\baselineskip]\fi
  \fi\endgroup}
\let\DD\relax\newcommand{\DD}{\ensuremath{\mathbb{D}}}
\let\cA\relax\newcommand{\cA}{\ensuremath{\mathcal{A}}}
\let\Dia\relax\newcommand{\Dia}{\ensuremath{\mathsf{Dia}}}
\let\Ob\relax\newcommand{\Ob}{\ensuremath{\mathsf{Ob}}}
\let\Hom\relax\newcommand{\Hom}{\ensuremath{\mathsf{Hom}}}
\let\hot\relax\newcommand{\hot}{\ensuremath{\mathsf{hot}}}
\let\Hot\relax\newcommand{\Hot}{\ensuremath{\mathsf{Hot}}}
\let\HOT\relax\newcommand{\HOT}{\ensuremath{\mathsf{HOT}}}
\let\isoto\relax\newcommand{\isoto}{\overset{\sim}{\to}}
\let\isolongto\relax\newcommand{\isolongto}{\overset{\sim}{\longrightarrow}}
\let\D\relax\newcommand{\D}{\ensuremath{\mathbb{D}}}
\let\Fl\relax\newcommand{\Fl}{\ensuremath{\mathsf{Fl}}}
\let\Cat\relax\newcommand{\Cat}{\ensuremath{\mathcal{C}at}}
\let\bs\relax\newcommand{\bs}{\backslash}
\let\RR\relax\newcommand{\RR}{\ensuremath{\mathbb{R}}}
\let\sfup\relax\newcommand{\sfup}[1]{\textsf{\upshape #1}}

\dgrclaim{conjectured}{78}{I-b-A069}{The doubled-arrow construction probably does not give the desired counterexample to base-change stability.}
\dgrclaim{false}{79}{I-b-B079-03}{Every universally cohomological D-fibration preserves all constant-coefficient D-equivalences under pullback.}
\dgrclaim{false}{79}{I-b-B079-04}{Every universally homological D-fibration preserves all constant-coefficient D-equivalences under pullback.}
\dgrclaim{false}{79}{I-b-B079-07}{The implication (b)\ensuremath{\Rightarrow}(a) nevertheless holds for all derivators D\_i of i-truncated homotopy types.}
\dgrclaim{false}{81}{I-b-B081-04}{Universal cohomological/homological D-fibrations as subsequently redefined are closed under composition.}
\dgrclaim{conjectured}{84}{I-b-B084-01}{For a universal D-equivalence f:X\ensuremath{\to}Y, restriction f\textasciicircum{}*:D(Y)\ensuremath{\to}D(X) is fully faithful.}
\dgrclaim{false}{84}{I-b-B084-06}{For D-fibrations over S, a D-equivalence of total categories must be a D-equivalence on each fibre.}
\dgrclaim{false}{84}{I-b-B084-07}{The converse of Proposition6(a) holds whenever S is D-contractible.}
\dgrclaim{false}{87}{I-b-B087-05}{A bundle over the circle is determined up to unique isomorphism by a specified monodromy automorphism of its fibre.}
\dgrclaim{false}{93}{I-b-B093-04}{Conditions(a)–(c) alone imply that the inclusion of a fibre of(E×Z)/G\ensuremath{\to}E/G is not a D-equivalence.}
\dgrclaim{false}{93}{I-b-B093-05}{If additionally H*(Z;k)\textasciicircum{}G=k, then X=(E×Z)/G is D(k)-aspherical.}
\dgrclaim{false}{95}{I-b-B095-01}{The orbit of the six-object poset has the displayed globular relations p1\textasciicircum{}-p0\textasciicircum{}-=p1\textasciicircum{}-p0\textasciicircum{}+ and p1\textasciicircum{}+p0\textasciicircum{}-=p1\textasciicircum{}+p0\textasciicircum{}+.}
\dgrclaim{conjectured}{96}{I-b-B096-04}{Under Der3a\ensuremath{'}, the revised weakly-contractible-base test condition is equivalent to each of the universal one-sided fibration conditions.}
\dgrclaim{false}{102}{Ib-C003}{If f is Hot-proper then both X\_y\ensuremath{\to}X/y and X/y\ensuremath{\to}X/y\ensuremath{'} are Hot-equivalences, universally.}
\dgrclaim{conjectured}{107}{Ib-C021}{Footnote55: for HOT, cohomological equivalence, homological equivalence, and constant-coefficient equivalence coincide.}
\dgrclaim{conjectured}{109–110}{Ib-C026}{Every Hot-equivalence induces an equivalence on locally constant HOT coefficients and hence preserves their cohomology and homology.}
\dgrclaim{conjectured}{115}{Ib-C048}{For HOT all five equivalence notions coincide with nerve weak equivalence.}
\dgrclaim{conjectured}{119}{Ib-C069}{For HOT, D-fibrations do preserve D-equivalences by pullback.}
\dgrclaim{false}{122–123}{Ib-C082}{With compatible internal Homs, every unqualified D-equivalence is 1-complete.}
\dgrclaim{conjectured}{125–126}{Ib-C092}{HOT satisfies c\ensuremath{''} and its equivalences are precisely nerve weak equivalences.}
\dgrclaim{conjectured}{126}{Ib-C096}{There should be a single intrinsic HOT-fibration notion on all small categories; the homotopy derivator must be constructed to justify the coefficient route.}
\par\endgroup

\section*{II-a \textperiodcentered\ Chapter II, part 1 of 2, manuscript pages 1--53}
\pdfbookmark[1]{II-a}{dgrclaims.II-a}
\begingroup\raggedright
% DER_II-a: the definitions of the chunk's preamble, applied inside its group
\let\printedpage\relax\newcommand{\printedpage}[1]{\begingroup\def\dcpagevalue{#1}\def\dcblankpage{--}\ifx\dcpagevalue\dcblankpage\else\marginpar{\scriptsize\raggedright p.~#1}\fi\endgroup}
\def\dcmsnumber page #1\relax{#1}
\newif\ifdcskipmsp
\let\dcpairedpage\relax\newcommand{\dcpairedpage}[2]{\marginpar{\scriptsize\raggedright p.~#1\\\mbox{[#2]}}\global\dcskipmsptrue}
\let\msp\relax\newcommand{\msp}[1]{\ifdcskipmsp\global\dcskipmspfalse\else\marginpar{\scriptsize\raggedright[\dcmsnumber#1\relax]}\fi}
% printed notation: the category C' enclosed in a V-shaped wedge with the arrow name under it
\let\Vcat\relax\newcommand{\Vcat}[2]{\underset{#2}{\underset{\vee}{#1}}}
\let\Homint\relax\newcommand{\Homint}{\operatorname{Hom\,int}}
\let\sX\relax\newcommand{\sX}{\mathcal{X}} % resolved against PDF: CMSY calligraphic X
\let\X\relax\newcommand{\X}{\mathcal{X}}
\let\Hom\relax\newcommand{\Hom}{\operatorname{Hom}}
\let\Ker\relax\newcommand{\Ker}{\operatorname{Ker}}
\let\Ob\relax\newcommand{\Ob}{\operatorname{Ob}}
\let\mspin\relax\newcommand{\mspin}[1]{\texttt{[#1]}}
\let\HOT\relax\newcommand{\HOT}{\mathrm{HOT}}
\let\Cat\relax\newcommand{\Cat}{\mathrm{Cat}}
\let\Fl\relax\newcommand{\Fl}{\mathrm{Fl}}
\let\Ens\relax\newcommand{\Ens}{\mathrm{Ens}}
\let\pl\relax\newcommand{\pl}{\texttt{[plutôt}}
\let\uHom\relax\newcommand{\uHom}{\underline{\operatorname{Hom}}}
\let\bs\relax\newcommand{\bs}{\backslash}
\let\cM\relax\newcommand{\cM}{\mathcal{M}}
\let\dbl\relax\newcommand{\dbl}[2]{\mathrel{\overset{#1}{\underset{#2}{\rightrightarrows}}}}
\let\defeq\relax\newcommand{\defeq}{\overset{\text{déf}}{=}}

\dgrclaim{conjectured}{3–4}{II-A-013}{For homotopical cofinality, replace the connectedness criterion on comma categories by weak contractibility of their nerves.}
\dgrclaim{false}{5}{II-A-017}{The manuscript invokes Der5 to infer $F\to u_*u^*F$ invertible for all F from homotopy strictness, i.e. contractibility of each $u/c$.}
\dgrclaim{false}{5}{II-A-018}{The dual asserted consequence is $u_!u^*F\to F$ invertible for all F when all $c/u$ are contractible.}
\dgrclaim{open}{5}{II-A-019}{If each $c/u$ is weakly contractible, is $H^i(C,F)\to H^i(C\prime,u^*F)$ an isomorphism for every abelian presheaf F and every i?}
\dgrclaim{conjectured}{7}{II-A-025}{For $K_u(c)=\pi_0(c/u)$ there is a natural isomorphism $\lim u^*F\cong\operatorname{Nat}(K_u,F)$.}
\dgrclaim{conjectured}{7–8}{II-A-026}{$\mathbf R\Gamma(C\prime,u^*F)\simeq\mathbf R\operatorname{Hom}_{C;k}(\mathbf Lu_!k_{C\prime},F)$ in D(k).}
\dgrclaim{conjectured}{8}{II-A-028}{For space-valued presheaves, $\mathbf R\Gamma(C\prime,u^*F)\simeq\mathbf R\Gamma(C,\underline{\operatorname{Hom}}(u_!1,F))$.}
\dgrclaim{false}{9}{II-A-032}{The proposed unsectioned identity is $u_*u^*F\simeq\underline{\operatorname{Hom}}(u_!1,F)$.}
\dgrclaim{false}{10}{II-A-034}{The proposed general formula is $u_*\underline{\operatorname{Hom}}(G\prime,u^*F)\simeq\underline{\operatorname{Hom}}(u_!G\prime,F)$.}
\dgrclaim{open}{10}{II-A-036}{There is a canonical projection comparison from adjunction, whose invertibility is in question.}
\dgrclaim{false}{12–13}{II-A-049}{The first list claims equivalence between local irreducibility of Top(X), local irreducibility of Esp(X), irreducibility of components, and irreducibility of all slices X/x.}
\dgrclaim{false}{12–13}{II-A-051}{The second list a–e again includes weak cofiltering of connected components among conditions equivalent to local irreducibility of Top(X).}
\dgrclaim{false}{13}{II-A-054}{The prose identifies local irreducibility of the presheaf topos with PS1 plus PS2, or componentwise cofiltering.}
\dgrclaim{open}{15}{II-A-060}{Find a topos-theoretic condition equivalent to PS1+PS2 and ensuring every connected étale category has weakly contractible nerve.}
\dgrclaim{false}{16}{II-A-066}{The NB claims PS1 is equivalent to componentwise weak cofiltering.}
\dgrclaim{false}{18}{II-B-006}{Total 0-connectedness alone implies W\_infty-asphericity.}
\dgrclaim{false}{19}{II-B-019}{The representable-base conditions (c,c-prime) imply the strong equivalent conditions (a,a-prime,b,d).}
\dgrclaim{false}{21}{II-B-024}{For every pair of connected U,V over an arbitrary object W, U times\_W V is nonempty.}
\dgrclaim{false}{24}{II-B-033}{Condition 2.6(c-prime) implies PS2.}
\dgrclaim{false}{28}{II-B-045}{If all slices X/x are cofiltered, then X must be locally cofiltered.}
\dgrclaim{conjectured}{36;36ter}{II-C-003}{Does T(A) imply that A is filtered?}
\dgrclaim{false}{36}{II-C-006}{Proposition2.19: a nonempty category with cofinal diagonal satisfies T(A).}
\dgrclaim{false}{36}{II-C-008}{A totally0-connected category has weakly contractible nerve.}
\dgrclaim{false}{36}{II-C-009}{Corollary to2.19: Delta\textasciicircum{}op satisfies T despite not being filtered.}
\dgrclaim{false}{36bis--36ter}{II-C-015}{In the proposed proof of2.20, a weak equivalence i:X\ensuremath{\to} E is used to infer that X/H is weakly contractible by a reference to ChapterXII.}
\dgrclaim{false}{51}{II-C-074}{The source asserts that M-connectedness is equivalent to the slice-limit comparison for coefficients on I/i restricted from I.}
\dgrclaim{false}{51}{II-C-076}{The converse to4.8 holds for ordered I without restrictions on M.}
\dgrclaim{false}{52}{II-C-078}{Proposition4.1(d prime) as printed: all coslice functors i\textbackslash{}I prime\ensuremath{\to}i\textbackslash{}I are cofinal is equivalent to full faithful restriction.}
\par\endgroup

\section*{III-V \textperiodcentered\ Chapter III, manuscript pages 1--12; Chapter IV, manuscript pages 73--83; Chapter V, manuscript pages 1--52}
\pdfbookmark[1]{III-V}{dgrclaims.III-V}
\begingroup\raggedright
% DER_III-V: the definitions of the chunk's preamble, applied inside its group
\let\printedpage\relax\newcommand{\printedpage}[1]{\marginpar{\footnotesize p.~#1}}
\let\msp\relax\newcommand{\msp}[1]{\marginpar{\footnotesize [#1]}}
\let\Hot\relax\newcommand{\Hot}{\mathrm{Hot}}
\let\D\relax\newcommand{\D}{\mathbf{D}}
\let\uHom\relax\newcommand{\uHom}{\underline{\mathrm{Hom}}}
\let\hm\relax\newcommand{\hm}{\mathrm{hom}}
\let\Hom\relax\newcommand{\Hom}{\mathrm{Hom}}
\let\homext\relax\newcommand{\homext}{\mathrm{hom}\,\mathrm{ext}}
\let\HOT\relax\newcommand{\HOT}{\mathrm{HOT}}
\let\Ob\relax\newcommand{\Ob}{\mathrm{Ob}\,}
\let\Ltens\relax\newcommand{\Ltens}{\overset{\mathbf{L}}{\otimes}}
\let\RHom\relax\newcommand{\RHom}{\mathbf{R}\mathrm{Hom}}
\let\ab\relax\newcommand{\ab}{\mathrm{ab}}
\let\defeq\relax\newcommand{\defeq}{\overset{\text{déf}}{=}}
\let\cA\relax\newcommand{\cA}{\mathcal{A}}
\let\bDelta\relax\newcommand{\bDelta}{\boldsymbol{\Delta}}
\let\id\relax\newcommand{\id}{\mathrm{id}}
\let\Fl\relax\newcommand{\Fl}{\operatorname{Fl}}
\let\uCh\relax\newcommand{\uCh}{\underline{\mathrm{Ch}}}
\let\bPhi\relax\newcommand{\bPhi}{\boldsymbol{\Phi}}
\let\uFl\relax\newcommand{\uFl}{\underline{\mathrm{Fl}}}
\let\Ch\relax\newcommand{\Ch}{\underline{\mathrm{Ch}}}
\let\card\relax\newcommand{\card}{\mathrm{card}}
\let\cM\relax\newcommand{\cM}{\mathcal{M}}
\let\Chh\relax\newcommand{\Chh}{\underline{\underline{\mathrm{Ch}}}}
\let\KK\relax\newcommand{\KK}{\underline{\underline{K}}}
\let\dfeq\relax\newcommand{\dfeq}{\overset{\text{déf}}{=}}
\let\Cathot\relax\newcommand{\Cathot}{\mathrm{Cat\,hot}}
\let\Ind\relax\newcommand{\Ind}{\operatorname{Ind}}
\let\Ens\relax\newcommand{\Ens}{\mathrm{Ens}}
\let\Cat\relax\newcommand{\Cat}{\mathrm{Cat}}
\let\uuK\relax\newcommand{\uuK}{\underline{\underline{K}}}

\dgrclaim{open}{III:1}{III-002}{Seek a covariant tensor $\mathrm{Hot}\times\mathcal A\to\mathcal A$ and a contravariant-in-space cotensor $\mathrm{Hot}^{op}\times\mathcal A\to\mathcal A$.}
\dgrclaim{conjectured}{III:1}{III-004}{For each space the two operations should extend from $\mathbf D(e)$ to morphisms $\mathbf D\to\mathbf D$.}
\dgrclaim{open}{III:3}{III-010}{The comparison $\mathrm{Hot}_0\to\mathrm{Hot}$ may be fully faithful for a suitable domain of sufficiently small homotopy types.}
\dgrclaim{conjectured}{III:4}{III-013}{A nerve weak equivalence $f:X'\to X$ is a $\mathbf D$-equivalence.}
\dgrclaim{open}{III:4}{III-015}{An equivalence of categories in Diag is a $\mathbf D$-equivalence.}
\dgrclaim{false}{III:5}{III-018}{Tensor and cotensor by the unpointed circle are suspension and based loops.}
\dgrclaim{conjectured}{III:7}{III-026}{The functor $\operatorname{Hom}_{(-)}(M,N)$ appears representable by an object of $\mathrm{Hot}_0$, or at least by a space.}
\dgrclaim{open}{III:8}{III-029}{For a Quillen model presentation, the external mapping space satisfies the representation formula (10).}
\dgrclaim{false}{III:10}{III-034}{Equation (14) specializes (13) to a circle and identifies it with suspension–loop adjunction.}
\dgrclaim{open}{III:11}{III-043}{Footnote 2 requests an explicit verification of the tensor-cotensor adjunction for complexes.}
\dgrclaim{conjectured}{III:11}{III-044}{A candidate external Hom is the space underlying the nonpositive truncation of $\mathbf R\operatorname{Hom}_k(M,N)$.}
\dgrclaim{open}{III:12}{III-045}{Footnote 3 asks whether $[\xi,H]\cong\operatorname{Hom}_{\mathrm{Hotab}}(\xi_{\mathrm{ab}},H)$.}
\dgrclaim{conjectured}{III:12}{III-051}{The better external Hom might take values in $D^-(\mathbb Z)$, retaining more than its nonpositive truncation.}
\dgrclaim{open}{IV:73}{IV-004}{Are coherent diagrams with isomorphic skeletons necessarily isomorphic?}
\dgrclaim{open}{IV:78}{IV-018}{Seek canonical coherent lifts of the unit $G\to f_*f^*G$ and counit $f^*f_*F\to F$, and dually for $f_!\dashv f^*$.}
\dgrclaim{conjectured}{IV:82}{IV-030}{The collage with its structural maps is characterized by the universal triple description, up to compatible canonical isomorphism.}
\dgrclaim{open}{IV:82}{IV-034}{Seek $\mathbf D(K)\to\mathbf D(I\times[1])$ whose underlying arrow is u.}
\dgrclaim{open}{IV:82}{IV-035}{Seek $\mathbf D(K)\to\mathbf D(J\times[1])$ whose underlying arrow is the adjoint v.}
\dgrclaim{false}{V:8--9}{V-P035}{Raw concatenation has shape $\Phi_{m+n}$ whenever $m,n$ have the same parity.}
\dgrclaim{false}{V:13}{V-P047}{Checking only simultaneous increments $(m,n)\mapsto(m+1,n+1)$ suffices for the proposed compatibility.}
\dgrclaim{false}{V:15--16,24}{V-P052}{Fixed-endpoint path categories are full subcategories of the unqualified functor categories.}
\dgrclaim{conjectured}{V:25--26}{V-P081}{A nonempty category of endpoint-preserving functors between generalized intervals is connected.}
\dgrclaim{conjectured}{V:26}{V-P082}{The same nonempty endpoint-preserving functor category is contractible.}
\dgrclaim{conjectured}{V:46}{V-I057}{Essential constancy should almost never occur in practice.}
\dgrclaim{conjectured}{V:49}{V-I064}{The required minimal lengths of stagewise homotopies may tend to infinity, obstructing a global categorical homotopy.}
\dgrclaim{false}{V:49}{V-I065}{For a Thomason fibrant K\ensuremath{'}, are categorically homotopic functors X\ensuremath{\to}K\ensuremath{'} always connected by one natural transformation?}
\dgrclaim{conjectured}{V:49-50}{V-I066}{Different reasonable theories of paths might yield canonically and transitively isomorphic unions in Cathot, as well as isomorphic ind-objects.}
\dgrclaim{false}{V:52}{V-I071}{This restricted Yoneda functor is conservative.}
\dgrclaim{false}{V:52}{V-I072}{Bijectivity on unbased homotopy classes from all spheres already detects weak equivalence.}
\dgrclaim{false}{V:52}{V-I073}{Is the restricted Yoneda functor Hot\ensuremath{\to}Fun(Hot\_0\textasciicircum{}op,Ens) fully faithful?}
\par\endgroup

\section*{VI-a \textperiodcentered\ Chapter VI, part 1 of 3, manuscript pages 1--70}
\pdfbookmark[1]{VI-a}{dgrclaims.VI-a}
\begingroup\raggedright
% DER_VI-a: the definitions of the chunk's preamble, applied inside its group
\let\printedpage\relax\newcommand{\printedpage}[1]{\marginpar{\footnotesize p.~#1}}
\let\msp\relax\newcommand{\msp}[1]{\marginpar{\footnotesize [#1]}}
\let\deq\relax\newcommand{\deq}{\overset{\text{déf}}{=}}
\let\ul\relax\newcommand{\ul}[1]{\underline{\mathrm{#1}}}
\let\xline\relax\newcommand{\xline}[1]{\mathrel{\overset{#1}{\rule[0.6ex]{2em}{0.4pt}}}}
\let\cat\relax\newcommand{\cat}[1]{\underline{\mathrm{#1}}}
\let\circled\relax\newcommand{\circled}[1]{\tikz[baseline=(c.base)]\node[draw,circle,inner sep=1pt](c){#1};}
\let\und\relax\newcommand{\und}[1]{\underline{\mathrm{#1}}}
\let\Cat\relax\newcommand{\Cat}{\mathrm{Cat}}
\let\uu\relax\newcommand{\uu}[1]{\underline{\mathrm{#1}}}
\let\lline\relax\newcommand{\lline}{\mathrel{\relbar\joinrel\relbar\joinrel\relbar}}
\let\uc\relax\newcommand{\uc}[1]{\underline{\mathrm{#1}}}
\let\HOT\relax\newcommand{\HOT}{\mathrm{HOT}}
\let\defeq\relax\newcommand{\defeq}{\overset{\text{déf}}{=}}
\let\cercle\relax\newcommand{\cercle}[1]{\tikz[baseline=(c.base)]\node[draw,circle,inner sep=1pt](c){#1};}
\let\Ob\relax\newcommand{\Ob}{\mathrm{Ob}}
\let\Fl\relax\newcommand{\Fl}{\mathrm{Fl}}
\let\id\relax\newcommand{\id}{\mathrm{id}}
\let\cA\relax\newcommand{\cA}{\mathcal{A}}

\dgrclaim{false}{4, 8}{VIa-039}{The initial comparison theorem asserts that all functors of the thirteen-category diagram induce equivalences after the localizations specified on ms7, identifying strict-fibrewise diagram theory with coarse Cat/S theory.}
\dgrclaim{conjectured}{5}{VIa-040}{The higher coherence discarded by ordinary localization should be recovered in relative internal Hom objects interpreted as \ensuremath{\infty}-groupoid stacks over S.}
\dgrclaim{false}{6, 33}{VIa-048}{Read literally for arbitrary W\ensuremath{\supseteq}W\ensuremath{\infty}, the source identifies transitions being Hot-equivalences with transitions belonging to W (ms6), and identifies Hot\_W-valued local coefficients exactly with the index-0 split categories defined by W\ensuremath{\infty}-transitions (ms33).}
\dgrclaim{conjectured}{14}{VIa-067}{Replacing \ensuremath{\Phi} and \ensuremath{\Psi} by \ensuremath{\Phi}\ensuremath{'}=\ensuremath{\alpha}\ensuremath{'}\ensuremath{\Phi}\ensuremath{_0} and \ensuremath{\Psi}\ensuremath{'}=\ensuremath{\beta}\ensuremath{'}\ensuremath{\Psi}\ensuremath{_0} is proposed as a way to make the initial localized equivalence argument work.}
\dgrclaim{false}{14}{VIa-068}{For an arbitrary W satisfying the initially stated localizer assumptions, the source identifies total and strict-fibrewise W-equivalence on Parf(S), Fibhot(S), Fib\ensuremath{_0}(S) and Cofib\ensuremath{_0}(S).}
\dgrclaim{false}{21}{VIa-097}{Footnote 13 literally states that if X is a homotopy fibration over S then \ensuremath{\Psi}(X) belongs to Fibsc\ensuremath{_0}(S).}
\dgrclaim{conjectured}{22}{VIa-103}{Desired numbered assertion (2): for f:X\ensuremath{\to}Y over S with X,Y perfect, or merely homotopy fibrations, total f\ensuremath{\in}W should imply f\_s\ensuremath{\in}W for every s.}
\dgrclaim{conjectured}{23, 45}{VIa-108}{The source conjectures on ms23 that perfect-base-change stability should hold for the finite homotopy-truncation localizers W\_i, and on ms45 asks whether they satisfy new L5, L6 and L6 bis.}
\dgrclaim{false}{23}{VIa-111}{From a commutative square with vertical weak homotopy equivalences and lower horizontal map f\ensuremath{'} merely in W, the source concludes that the upper fibre map f\ensuremath{'}\_\{s\ensuremath{'}\} is also a weak homotopy equivalence.}
\dgrclaim{false}{26}{VIa-123}{Every Grothendieck fibration over S is isomorphic over S to a split fibration, so Fibsc(S)\ensuremath{\to}Fib(S) is essentially surjective as an ordinary functor; dually for opfibrations.}
\dgrclaim{false}{26}{VIa-124}{Restricting Fib(S) to Cartesian functors makes Fibsc(S)\ensuremath{\to}Fib(S) fully faithful and an equivalence of ordinary categories; the cocartesian dual assertion also holds.}
\dgrclaim{conjectured}{26-27}{VIa-125}{After restricting morphisms from arbitrary functors over S to coherent splitting-preserving ones and simultaneously restricting weak equivalences from coarse to fibrewise ones, the resulting localized categories might still be equivalent.}
\dgrclaim{false}{30-31}{VIa-135}{In the surrounding contravariant \ensuremath{\Phi} construction, its transport square is the printed covariant square with fibres X/s and X/t and transport u\_*:X\_s\ensuremath{\to}X\_t.}
\dgrclaim{conjectured}{32-35}{VIa-141}{The nine nonsplit categories of the p8 diagram, with coarse W-equivalences, should model the full subcategory of locally constant objects of HOT\_W(S); passing from an arbitrary diagram to its total category and then C should give its locally constant replacement.}
\dgrclaim{false}{33}{VIa-143}{The category Cofibsc(S) is the outer opposite category (Fibsc(S\textasciicircum{}op))\textasciicircum{}op.}
\dgrclaim{conjectured}{33, 35, 37}{VIa-145}{The functors Fibsc\_0(S)\ensuremath{\to}Fib\_0(S) and Cofibsc\_0(S)\ensuremath{\to}Cofib\_0(S) should induce equivalences after localization, with quasi-inverses represented by \ensuremath{\Phi}\_0 and \ensuremath{\Psi}\_0, under the standing fibre/coarse identification hypothesis on internal objects.}
\dgrclaim{conjectured}{33-35, 37}{VIa-146}{The inclusions Fibsc\_0(S)\ensuremath{\to}Fibsc(S) and Cofibsc\_0(S)\ensuremath{\to}Cofibsc(S), localized fibrewise, should be fully faithful, with essential images precisely the locally constant objects of the respective contravariant and covariant diagram homotopy categories.}
\dgrclaim{false}{33-34}{VIa-147}{The fine covariant theory and its locally constant inclusion are to be written using (HOT\_W(S\textasciicircum{}op))\textasciicircum{}op, as printed in the p33–34 discussion.}
\dgrclaim{conjectured}{34-35}{VIa-148}{After identifying the coarse localization with local systems, forgetting the strict diagram structure should induce retractions from HOT\_W(S) and HOT\_W(S\textasciicircum{}op) onto their locally constant full subcategories.}
\dgrclaim{false}{38}{VIa-156}{The equation k\ensuremath{\psi}\ensuremath{\cong}id\_B implies that k is full on Hom\_A(a,a\ensuremath{'}) for arbitrary objects a,a\ensuremath{'} of A.}
\dgrclaim{false}{38}{VIa-157}{Under k\ensuremath{\psi}\ensuremath{\cong}id\_B, failure of full faithfulness of \ensuremath{\psi} is equivalent to failure of faithfulness of k.}
\dgrclaim{conjectured}{38}{VIa-160}{The essential image of \ensuremath{\bar{\ensuremath{\psi}}}\ensuremath{'} is expected to consist precisely of the locally constant objects of HOT\_W(S).}
\dgrclaim{conjectured}{41}{VIa-173}{Left-fine W-equivalences should be preserved by smooth base change T\ensuremath{\to}S, extending their stability under discrete-fibration base change.}
\dgrclaim{conjectured}{41}{VIa-174}{Right-fine W-equivalences should be preserved by proper base change T\ensuremath{\to}S.}
\dgrclaim{false}{42, 44}{VIa-180}{Proposition 1(1), which compares X/s\ensuremath{\to}Y/s over an independent base S, is simply the original absolute axiom L4.}
\dgrclaim{conjectured}{43}{VIa-186}{The integral-homology localizer is suspected also to fail the total-to-fibre implication for maps between homotopy fibrations.}
\dgrclaim{false}{43, 45}{VIa-188}{L1, L3 quater, relative L4 bis, and the new L5 alone force W=W\ensuremath{\infty}, with no component-preservation hypothesis.}
\dgrclaim{conjectured}{44}{VIa-191}{For the fundamental localizers under consideration, L6—or even the new L5 alone—should imply L6 bis.}
\dgrclaim{false}{45}{VIa-200}{Saturation, the terminal-object axiom, relative L4 bis, and the new L5 alone imply W\ensuremath{\subseteq}W\ensuremath{\infty}, without any component condition.}
\dgrclaim{false}{47}{VIa-208}{On Liss(S), Fib(S), and Fibsc(S), the coarse W-class equals the right-fine W-class; dually on Prop(S), Cofib(S), and Cofibsc(S), it equals the left-fine W-class.}
\dgrclaim{conjectured}{49, 53-57}{VIa-213}{Does the coarse weak-equivalence class also give the asserted comparison for the split categories Fibsc(S) and Cofibsc(S), whose morphisms must preserve the chosen splitting?}
\dgrclaim{false}{50}{VIa-218}{The common left-fine localization of Cat/S, Prop(S), Cofib(S), and Cofibsc(S) is isomorphic to HOT\_W(S\textasciicircum{}op)\textasciicircum{}op.}
\dgrclaim{false}{51}{VIa-224}{The editorial replacement HOT\_W\textasciicircum{}o(S):=HOT\_W(S\textasciicircum{}op)\textasciicircum{}op is the correct target convention for the covariant comparison functor j.}
\dgrclaim{conjectured}{51-52, 64-69}{VIa-226}{The canonical functors i:HOT\_W\textasciicircum{}lc(S)\ensuremath{\to}HOT\_W(S) and j:HOT\_W\textasciicircum{}lc(S)\ensuremath{\to}HOT\_W(S\textasciicircum{}op) should be fully faithful under the p45 package L1, L3 quater, L4 bis, L7 and the new L5.}
\dgrclaim{false}{55}{VIa-238}{A Cartesian functor between split fibrations is therefore compatible with their specified splittings.}
\dgrclaim{false}{55-57}{VIa-239}{The pullback retractions \ensuremath{\gamma}\_X:\ensuremath{\Phi}(X)\ensuremath{\to}X form a strictly natural transformation on the full category Fib(S), allowing every functor over S and imposing no Cartesian condition on the morphisms.}
\dgrclaim{false}{56}{VIa-242}{For every X\ensuremath{\to}S and split fibration Y\ensuremath{\to}S, restriction along \ensuremath{\eta}\_X gives a bijection Hom\_cart/S(\ensuremath{\Phi}X,Y)\ensuremath{\cong}Hom\_S(X,UY), where the left-hand side includes all Cartesian functors.}
\dgrclaim{false}{65-67}{VIa-267}{To make Cbar a right adjoint to \ensuremath{\rho}, the required adjunction arrows should have the initially drawn directions id\_coarse\ensuremath{\to}\ensuremath{\rho}Cbar and Cbar\ensuremath{\rho}\ensuremath{\to}id\_fine.}
\dgrclaim{false}{65}{VIb-A006}{The attempted unit/counit directions before p.67 do not yet give $\rho\dashv\bar C$.}
\par\endgroup

\section*{VI-b \textperiodcentered\ Chapter VI, part 2 of 3, manuscript pages 71--145}
\pdfbookmark[1]{VI-b}{dgrclaims.VI-b}
\begingroup\raggedright
% DER_VI-b: the definitions of the chunk's preamble, applied inside its group
\let\printedpage\relax\newcommand{\printedpage}[1]{\marginpar{\footnotesize p.~#1}}
\let\msp\relax\newcommand{\msp}[1]{\marginpar{\footnotesize [#1]}}
\let\HOT\relax\newcommand{\HOT}{\mathrm{HOT}}
\let\Cat\relax\newcommand{\Cat}{\mathrm{Cat}}
\let\Ob\relax\newcommand{\Ob}{\mathrm{Ob}}
\let\Fl\relax\newcommand{\Fl}{\mathrm{Fl}}
\let\id\relax\newcommand{\id}{\mathrm{id}}
\let\cA\relax\newcommand{\cA}{\mathcal{A}}
\let\ul\relax\newcommand{\ul}[1]{\underline{\mathrm{#1}}}
\let\circled\relax\newcommand{\circled}[1]{\tikz[baseline=(c.base)]\node[draw,circle,inner sep=1pt](c){#1};}
\let\Hot\relax\newcommand{\Hot}{\mathrm{Hot}}
\let\hot\relax\newcommand{\hot}{\mathrm{hot}}
\let\Hom\relax\newcommand{\Hom}{\operatorname{Hom}}
\let\cercle\relax\newcommand{\cercle}[1]{\tikz[baseline=(c.base)]\node[draw,circle,inner sep=1pt](c){#1};}
\let\ovale\relax\newcommand{\ovale}[1]{\tikz[baseline=(c.base)]\node[draw,rounded corners=0.8ex,inner sep=2pt](c){#1};}
\let\usl\relax\newcommand{\usl}[1]{\underline{\mathrm{#1}}}
\let\encadre\relax\newcommand{\encadre}[1]{\tikz[baseline=(n.base)]\node[draw,rounded corners=3pt,inner sep=1.5pt](n){#1};}
\let\cM\relax\newcommand{\cM}{\mathcal{M}}
\let\cF\relax\newcommand{\cF}{\mathcal{F}}
\let\Fib\relax\newcommand{\Fib}{\underline{\mathrm{Fib}}\,}
\let\Fibhot\relax\newcommand{\Fibhot}{\underline{\mathrm{Fibhot}}\,}
\let\Ho\relax\newcommand{\Ho}{\mathrm{Ho}}
\let\Fibu\relax\newcommand{\Fibu}{\underline{\mathrm{Fib}}\,}
\let\WL\relax\newcommand{\WL}{W_S^{\mathcal{L}}}
\let\defeq\relax\newcommand{\defeq}{\overset{\text{déf}}{=}}
\let\Ch\relax\newcommand{\Ch}{\mathrm{Ch}}
\let\pr\relax\newcommand{\pr}{\mathrm{pr}}
\let\Fibsc\relax\newcommand{\Fibsc}{\underline{\mathrm{Fibsc}}\,}
\let\Cofib\relax\newcommand{\Cofib}{\underline{\mathrm{Cofib}}\,}
\let\Cofibsc\relax\newcommand{\Cofibsc}{\underline{\mathrm{Cofibsc}}\,}
\let\Parf\relax\newcommand{\Parf}{\underline{\mathrm{Parf}}\,}
\let\Liss\relax\newcommand{\Liss}{\underline{\mathrm{Liss}}}
\let\Prop\relax\newcommand{\Prop}{\underline{\mathrm{Prop}}}
\let\M\relax\newcommand{\M}{\mathcal{M}}
\let\und\relax\newcommand{\und}[1]{\underline{\mathrm{#1}}}
\let\uHom\relax\newcommand{\uHom}{\underline{\mathrm{Hom}}}
\let\uFib\relax\newcommand{\uFib}{\underline{\mathrm{Fib}}}
\let\uFibsc\relax\newcommand{\uFibsc}{\underline{\mathrm{Fibsc}}}
\let\uCh\relax\newcommand{\uCh}{\underline{\mathrm{Ch}}}
\let\Wd\relax\newcommand{\Wd}{W^{\mathrm{d}}}

\dgrclaim{false}{71}{VIb-A027}{The editorial outer-opposite correction asserts $\HOT_W^o(S)\simeq(\HOT_W(S^{op}))^{op}$ while retaining the ordinary localization definitions.}
\dgrclaim{false}{71}{VIb-A031}{Footnote46 conjectures that coarse and fibrewise localized theories differ whenever S is not equivalent to a discrete category.}
\dgrclaim{false}{74}{VIb-A041}{The covariant underlying diagram for arbitrary X/S is printed as $s\mapsto\hot_W(s\backslash X)$.}
\dgrclaim{open}{82}{VIb-A070}{Can the coarse comparison be obtained from a Quillen-type replacement without choosing a functorial factorization?}
\dgrclaim{open}{82}{VIb-A071}{Can fine right and left relative equivalences in M/S be defined intrinsically from a general category of models (M,W)?}
\dgrclaim{false}{82}{VIb-R001}{The prose identifies local-isomorphism tests with W\_S\textasciicircum{}d and co-local-isomorphism tests with W\_S\textasciicircum{}g.}
\dgrclaim{conjectured}{83}{VIb-A073}{The suggested smooth and Zariski-like localizations should have the same relevant invariants, and dually proper and co-Zariski-like localizations should have the same cohomology.}
\dgrclaim{false}{84}{VIb-A075}{The compatibility $i_{CX}=C(i_X)$ cannot be formal under the hypotheses already available.}
\dgrclaim{open}{84}{VIb-A077}{Does $i_{CX}=C(i_X)$ hold after inverting universal weak equivalences over S?}
\dgrclaim{open}{84}{VIb-A078}{Does duality reverse the order of the two adjunctions in this comparison?}
\dgrclaim{open}{85}{VIb-A079}{A finer localizer on M/S contained in W\_S gives a localization from the fine quotient to the coarse quotient; constructing a section and adjunction requires suitable fibrant representatives.}
\dgrclaim{open}{87}{VIb-A086}{To construct a quasi-inverse one seeks a functor gamma:M/S\ensuremath{\to}(Fib S)[W\_S\textasciicircum{}\{-1\}] sending X to C(X) and W\_S to isomorphisms.}
\dgrclaim{conjectured}{94}{VIb-B015}{Replacing local-isomorphism tests by all smooth morphisms gives the same left relative localizer.}
\dgrclaim{open}{94}{VIb-B017}{Is closure of the test class $\mathcal L$ under composition needed?}
\dgrclaim{conjectured}{97–98}{VIb-B028}{The section i is expected to be fully faithful.}
\dgrclaim{conjectured}{98}{VIb-B029}{The localization rho is expected to be left adjoint to i.}
\dgrclaim{false}{105}{VIb-B049}{The source attributes $(x,c)\mapsto(x,1_{f(x)},c)$ to $i_{CX}$ and $(x,c)\mapsto(x,c,1_{\beta(c)})$ to $C(i_X)$, respectively.}
\dgrclaim{false}{108}{VIb-B057}{Footnote59 asserts that the stability hypotheses on Fib are used only in II and only in case A/A-prime.}
\dgrclaim{false}{115}{VIb-B085}{The editorial correction inserts an outer opposite in $\mathrm{HOT}_W^o(S)=\mathrm{HOT}_W(S^{op})^{op}$.}
\dgrclaim{false}{120}{VIb-C009}{The path replacement $C_S(X)$ represents the same object of $\HOT_W(S)$ as arbitrary $X\to S$.}
\dgrclaim{false}{120}{VIb-C011}{Corollary 2, second isomorphism: $f^*q_S(X)\cong q_{S^{\prime}}(C_S(X)\times_S S^{\prime})$ for arbitrary $X\to S$.}
\dgrclaim{open}{120}{VIb-C012}{Which changes of base $f:S^{\prime}\to S$ allow strict pullback to calculate derived restriction on every $X\to S$?}
\dgrclaim{conjectured}{122}{VIb-C016}{Proper $f:S^{\prime}\to S$ induces a strict pullback functor preserving $W_S^{\mathrm d}$.}
\dgrclaim{conjectured}{122}{VIb-C017}{Smooth $f:S^{\prime}\to S$ induces a strict pullback functor preserving $W_S^{\mathrm g}$.}
\dgrclaim{open}{123}{VIb-C019}{Is preservation of right-fine weak equivalences by pullback equivalent to $W$-properness of the base functor?}
\dgrclaim{open}{124}{VIb-C024}{If an arrow of a chosen category over $S^{\prime}$ becomes invertible in its coarse localization, must it belong to $W_{S^{\prime}}$?}
\dgrclaim{conjectured}{124}{VIb-C025}{For proper change of base, the same comparison criterion can be formulated on all of $\Cat/S$.}
\dgrclaim{false}{126}{VIb-C028}{An arrow invertible after an arbitrary localization necessarily belonged to the original class of arrows being inverted.}
\dgrclaim{open}{126}{VIb-C031}{Does a Hot-fibration which is not proper furnish a pullback that fails to preserve the right-fine class on all categories over the base?}
\dgrclaim{open}{127}{VIb-C032}{Are there functors other than Hot-fibrations or W-fibrations for which restriction commutes with local-constant reflection?}
\dgrclaim{conjectured}{128}{VIb-C035}{The varying-base assignment $S\mapsto\HOT_W^{\mathrm{lc}}(S)$ is not expected to be a derivator.}
\dgrclaim{open}{128}{VIb-C036}{Do restriction functors between locally constant categories have left and right adjoints?}
\dgrclaim{open}{130}{VIb-C042}{For an arbitrary relative category (M,W), are isomorphisms in localized diagram categories detected by evaluation?}
\dgrclaim{open}{130–132}{VIb-C043}{Does $\HOT_W(\coprod_i S_i)\to\prod_i\HOT_W(S_i)$ remain an equivalence for arbitrary small index sets?}
\dgrclaim{open}{131}{VIb-C045}{For arbitrary relative categories (F\_i,W\_i), is $(\prod_iF_i)[(\prod_iW_i)^{-1}]\to\prod_i(F_i[W_i^{-1}])$ an equivalence?}
\dgrclaim{conjectured}{132}{VIb-C048}{Uniform short presentations of morphisms could prove the infinite coproduct axiom for HOT.}
\dgrclaim{open}{137}{VIb-C061}{Does a Hot-fibration f ensure that direct image preserves locally constant types?}
\dgrclaim{false}{142}{VIb-C070}{Pullback of the right-fine replacement map $i_X$ along arbitrary S-prime\ensuremath{\to}S is again right-fine.}
\dgrclaim{false}{142–143}{VIb-C071}{The counit may always be defined by projection composed with the inverse of the base-changed replacement map.}
\dgrclaim{false}{142}{VIc-A007}{Question: does the pullback i\_X times\_S id\_\{S\ensuremath{\prime}\}:X times\_S S\ensuremath{\prime} -> Phi\_S(X) times\_S S\ensuremath{\prime} necessarily belong to W\_S\textasciicircum{}d?}
\dgrclaim{false}{142}{VIc-A008}{Proposed formula: beta\_X=pr\_1\textasciicircum{}X (i\_X times\_S id\_\{S\ensuremath{\prime}\})\textasciicircum{}\{-1\} defines the counit at every representative X over S.}
\par\endgroup

\section*{VI-c \textperiodcentered\ Chapter VI, part 3 of 3, manuscript pages 146--204}
\pdfbookmark[1]{VI-c}{dgrclaims.VI-c}
\begingroup\raggedright
% DER_VI-c: the definitions of the chunk's preamble, applied inside its group
\let\printedpage\relax\newcommand{\printedpage}[1]{\marginpar{\footnotesize p.~#1}}
\let\msp\relax\newcommand{\msp}[1]{\marginpar{\footnotesize [#1]}}
\let\HOT\relax\newcommand{\HOT}{\mathrm{HOT}}
\let\hot\relax\newcommand{\hot}{\mathrm{hot}}
\let\Cat\relax\newcommand{\Cat}{\mathrm{Cat}}
\let\id\relax\newcommand{\id}{\mathrm{id}}
\let\pr\relax\newcommand{\pr}{\mathrm{pr}}
\let\Ob\relax\newcommand{\Ob}{\mathrm{Ob}}
\let\Wd\relax\newcommand{\Wd}{W^{\mathrm{d}}}
\let\Hom\relax\newcommand{\Hom}{\mathrm{Hom}}
\let\uHom\relax\newcommand{\uHom}{\underline{\mathrm{Hom}}}
\let\Liss\relax\newcommand{\Liss}{\underline{\mathrm{Liss}}\,}
\let\cA\relax\newcommand{\cA}{\mathcal{A}}
\let\dd\relax\newcommand{\dd}{\mathrm{d}}
\let\lc\relax\newcommand{\lc}{\mathrm{lc}}
\let\adj\relax\newcommand{\adj}[2]{\mathrel{\underset{#2}{\overset{#1}{\rightleftarrows}}}}
\let\multhot\relax\newcommand{\multhot}{\mathrm{multhot}}
\let\mult\relax\newcommand{\mult}{\mathrm{mult}}
\let\Hot\relax\newcommand{\Hot}{\mathrm{Hot}}
\let\D\relax\newcommand{\D}{\mathbf{D}}
\let\bs\relax\newcommand{\bs}{\backslash}
\let\ul\relax\newcommand{\ul}[1]{\underline{\mathrm{#1}}}
\let\circled\relax\newcommand{\circled}[1]{\tikz[baseline=(c.base)]\node[draw,rounded corners,inner sep=2pt](c){#1};}
\let\Fib\relax\newcommand{\Fib}{\underline{\mathrm{Fib}}}
\let\Parf\relax\newcommand{\Parf}{\underline{\mathrm{Parf}}}
\let\Prop\relax\newcommand{\Prop}{\underline{\mathrm{Prop}}}
\let\Fibsc\relax\newcommand{\Fibsc}{\underline{\mathrm{Fibsc}}}
\let\Cofibsc\relax\newcommand{\Cofibsc}{\underline{\mathrm{Cofibsc}}}
\let\lcs\relax\newcommand{\lcs}{\mathrm{lc}\,!}
\let\cercle\relax\newcommand{\cercle}[1]{\tikz[baseline=(c.base)]\node[draw,circle,inner sep=1pt](c){#1};}
\let\troisbarres\relax\newcommand{\troisbarres}{\leavevmode\llap{$|||$\quad}}
\let\Bifib\relax\newcommand{\Bifib}{\underline{\mathrm{Bifib}}}
\let\margbar\relax\newcommand{\margbar}{\marginpar{\(\|\)}}
\let\und\relax\newcommand{\und}[1]{\underline{\mathrm{#1}}}
\let\circnum\relax\newcommand{\circnum}[1]{\tikz[baseline=(c.base)]\node[draw,circle,inner sep=1pt](c){#1};}

\dgrclaim{false}{151}{VIc-A021}{Proposed product formula: f\_!f\textasciicircum{}*(xi)=[S\ensuremath{\prime}/S] times\_S xi for every base functor f and every coefficient xi.}
\dgrclaim{false}{152}{VIc-A025}{False Corollary1: if only one of X,Y is smooth, then [X times\_S Y] is necessarily the product [X] times\_S [Y].}
\dgrclaim{false}{152}{VIc-A026}{Every smooth X -> S induces by pullback a functor Cat/S -> Cat/X preserving all right-fine equivalences.}
\dgrclaim{conjectured}{159}{VIc-A052}{The strengthened Der3b\ensuremath{\prime}\_2 says that every f in W induces an equivalence f\_0\textasciicircum{}*:HOT\_W\textasciicircum{}lc(S)->HOT\_W\textasciicircum{}lc(S\ensuremath{\prime}).}
\dgrclaim{conjectured}{159}{VIc-A054}{Full faithfulness of f\_0\textasciicircum{}* is expected to follow from the enhanced cohomological equivalence condition by applying internal Hom and cohomological direct images.}
\dgrclaim{conjectured}{159}{VIc-A055}{The construction of the right adjoints f\_* and of internal Hom is expected to proceed together.}
\dgrclaim{conjectured}{163}{VIc-B011}{Properness should be equivalent to preservation of the right-fine localizers by pullback.}
\dgrclaim{conjectured}{172–173}{VIc-B034}{There may exist substantial homotopy theories HOT\_W for classes W strictly smaller than W\ensuremath{\infty} while W(1,2,3) hold.}
\dgrclaim{false}{180–181}{VIc-B065}{Intermediate full subcategories between Parf\_omega S and Fib\_W S might yield distinct fine localizations.}
\dgrclaim{open}{184}{VIc-C010}{Does the old restricted W(4b\ensuremath{\prime}) alone force $\lambda:\mathrm{HOT}_W^{lc}(S)\to\mathrm{HOT}_W^{lc!}(S)$ to be an equivalence?}
\dgrclaim{conjectured}{185}{VIc-C015}{The new Cartesian fibre-detection axiom should imply the strongest W(4a) preservation condition, at least for the strong saturation of W.}
\dgrclaim{open}{186}{VIc-C019}{Do $\mathrm{Parf}_W S\hookrightarrow\mathrm{Liss}_{0,W}S$ and $\mathrm{Parf}_W S\hookrightarrow\mathrm{Prop}_{0,W}S$ induce equivalences after inverting fibrewise W-equivalences, without W(4b)?}
\dgrclaim{open}{186}{VIc-C020}{The same fine-localization comparison is asked with $\mathrm{Parf}_\omega S$ in place of $\mathrm{Parf}_W S$.}
\dgrclaim{open}{186}{VIc-C022}{Does $(\mathrm{Bifib}S)[(W_S^f)^{-1}]\to(\mathrm{Fib}_WS)[(W_S^f)^{-1}]$ become essentially surjective, or an equivalence, even assuming W(4a)?}
\dgrclaim{open}{187}{VIc-C024}{Does Phi carry W-proper objects, or proper weakly W-fibred objects, to W-proper objects?}
\dgrclaim{conjectured}{189}{VIc-C028}{If both arrows in the span are open embeddings, or both closed embeddings, is the total category W-equivalent to the ordinary amalgamated sum $X_a\amalg_{X_0}X_b$?}
\dgrclaim{conjectured}{201}{VIc-C077}{If X is weakly W-fibred, then $[X\times_S Y]\to[X]\times_S[Y]$ should be invertible for every Y, without smoothness of Y.}
\dgrclaim{false}{202}{VIc-C083}{The general product formula should fail if neither X is W-proper nor Y W-smooth, even when X is weakly W-fibred.}
\par\endgroup

\section*{VII-a \textperiodcentered\ Chapter VII, part 1 of 3, manuscript pages 1--57}
\pdfbookmark[1]{VII-a}{dgrclaims.VII-a}
\begingroup\raggedright
% DER_VII-a: the definitions of the chunk's preamble, applied inside its group
\let\printedpage\relax\newcommand{\printedpage}[1]{\if\relax\detokenize{#1}\relax\else\marginpar{\raggedright\footnotesize p.~#1}\fi}
\let\msp\relax\newcommand{\msp}[1]{\marginpar{\raggedright\footnotesize [#1]}}
\let\Fltr\relax\newcommand{\Fltr}{\operatorname{Fl}\operatorname{tr}}
\let\Ob\relax\newcommand{\Ob}{\operatorname{Ob}}
\let\bDelta\relax\newcommand{\bDelta}{\boldsymbol{\Delta}}
\let\bD\relax\newcommand{\bD}{\boldsymbol{\Delta}}
\let\iddots\relax\newcommand{\iddots}{\mathinner{\mkern1mu\raise1pt\hbox{.}\mkern2mu\raise4pt\hbox{.}\mkern2mu\raise7pt\hbox{.}\mkern1mu}}
\let\Ch\relax\newcommand{\Ch}{\underline{\mathrm{Ch}}}
\let\uHom\relax\newcommand{\uHom}{\underline{\mathrm{Hom}}}
\let\uFl\relax\newcommand{\uFl}{\underline{\mathrm{Fl}}}
\let\opp\relax\newcommand{\opp}{\mathrm{opp}}
\let\Hom\relax\newcommand{\Hom}{\mathrm{Hom}}
\let\uCh\relax\newcommand{\uCh}{\underline{\mathrm{Ch}}}
\let\Ty\relax\newcommand{\Ty}{\mathrm{Ty}}
\let\uTy\relax\newcommand{\uTy}{\underline{\mathrm{Ty}}}
\let\Cat\relax\newcommand{\Cat}{\mathrm{Cat}}
\let\circled\relax\newcommand{\circled}[1]{\text{\textcircled{$\scriptstyle #1$}}}
\let\ul\relax\newcommand{\ul}[1]{\underline{\mathrm{#1}}}
\let\red\relax\newcommand{\red}{\text{\upshape réd}}
\let\Tr\relax\newcommand{\Tr}{\mathrm{Tr}}
\let\Cr\relax\newcommand{\Cr}{\mathrm{Cr}}
\let\lng\relax\newcommand{\lng}{\operatorname{long}}
\let\card\relax\newcommand{\card}{\operatorname{card}}
\let\uTyp\relax\newcommand{\uTyp}{\underline{\mathrm{Typ}}}
\let\ovalbox\relax\newcommand{\ovalbox}[1]{\tikz[baseline=(X.base)]\node[draw,rounded corners=6pt,inner sep=2pt](X){#1};}

\dgrclaim{false}{4-5}{VIIa-A013}{As printed: composable chains correspond to full simplex subcategories, and maximal chains of length i correspond exactly to object-injective functors $\Delta_i\to I$.}
\dgrclaim{false}{5}{VIIa-A018}{Proposition 1(b), literally: every endofunctor of an interval fixing both endpoints is the identity.}
\dgrclaim{false}{10}{VIIa-A031}{Proposition 3 as printed: with $A=(I/s)\setminus\{s\}$ and $B=(s\backslash I)\setminus\{s\}$, s is critical iff $I\setminus\{s\}=A\amalg B$.}
\dgrclaim{false}{11}{VIIa-A037}{Footnote 4, literally: every subcategory of I which is an interval is necessarily full.}
\dgrclaim{false}{13}{VIIa-A046}{As printed, the length n of the path is $\operatorname{card} I$.}
\dgrclaim{false}{33}{VIIa.B030}{A constant functor at every non-endpoint vertex of a target interval factors through exactly two critical intervals.}
\dgrclaim{false}{34}{VIIa.B033}{For increasing f, the full subcategory on the image critical vertices equals the ambient integer interval [f(s\_0),f(s\_r)].}
\dgrclaim{false}{36}{VIIa.B046}{The global source critical intervals contained in f\textasciicircum{}\{-1\}(J') are exactly all critical intervals of the preimage.}
\dgrclaim{false}{42}{VIIa-C009}{Each target transition segment [i-1,i] equals the unique image block [t\_\{j-1\},t\_j] which crosses it.}
\dgrclaim{false}{45}{VIIa-C017}{The object map is constant on every integer-connected component of the set of vertices incident to a collapsed source transition.}
\dgrclaim{false}{47}{VIIa-C030}{The printed object-map table sends the third fibre block to l\_2 and its general block to l\_i.}
\dgrclaim{false}{50}{VIIa-C036}{The same object concatenation extends to a bifunctor $\uTyp\times\uTyp\to\uTyp$, where Typ has all functors of interval categories.}
\par\endgroup

\section*{VII-b \textperiodcentered\ Chapter VII, part 2 of 3, manuscript pages 58--117}
\pdfbookmark[1]{VII-b}{dgrclaims.VII-b}
\begingroup\raggedright
% DER_VII-b: the definitions of the chunk's preamble, applied inside its group
\providecommand\IfDocumentMetadataT[1]{}
\let\printedpage\relax\newcommand{\printedpage}[1]{\marginpar{\footnotesize p.~#1}}
\let\msp\relax\newcommand{\msp}[1]{\marginpar{\vspace*{10pt}\footnotesize [#1]}}
\let\uTy\relax\newcommand{\uTy}{\underline{\mathrm{Ty}}}
\let\Ch\relax\newcommand{\Ch}{\underline{\mathrm{Ch}}}
\let\Ty\relax\newcommand{\Ty}{\underline{\mathrm{Ty}}}
\let\Nerf\relax\newcommand{\Nerf}{\mathrm{Nerf}}
\let\Fl\relax\newcommand{\Fl}{\mathrm{Fl}}
\let\Hom\relax\newcommand{\Hom}{\mathrm{Hom}}
\let\ch\relax\newcommand{\ch}{\mathrm{ch}}
\let\hot\relax\newcommand{\hot}{\mathrm{hot}}
\let\id\relax\newcommand{\id}{\mathrm{id}}
\let\CH\relax\newcommand{\CH}{\underline{\mathrm{CH}}}
\let\red\relax\newcommand{\red}{\mathrm{r\acute{e}d}}
\let\uHom\relax\newcommand{\uHom}{\underline{\mathrm{Hom}}}
\let\ul\relax\newcommand{\ul}[1]{\underline{#1}}
\let\triv\relax\newcommand{\triv}{\mathrm{triv}}
\let\uCh\relax\newcommand{\uCh}{\underline{\mathrm{Ch}}}
\let\Cat\relax\newcommand{\Cat}{\mathrm{Cat}}
\let\Cathot\relax\newcommand{\Cathot}{\mathrm{Cat\,hot}}
\let\Hot\relax\newcommand{\Hot}{\mathrm{Hot}}
\let\Ob\relax\newcommand{\Ob}{\mathrm{Ob}\,}
\let\sHom\relax\newcommand{\sHom}{\underline{\mathcal{H}\mathit{om}}}
\let\To\relax\newcommand{\To}{T_{\mathrm{o}}}
\let\circled\relax\newcommand{\circled}[1]{\tikz[baseline=(c.base)]\node[draw,circle,inner sep=1pt](c){#1};}
\let\Hotz\relax\newcommand{\Hotz}{\mathrm{Hot}^{\wedge}_{0}}
\let\CNo\relax\newcommand{\CNo}{C_{\mathbf{N}^{o}}}
\let\lab\relax\newcommand{\lab}[2][l]{\text{\footnotesize\begin{tabular}{@{}#1@{}}#2\end{tabular}}}

\dgrclaim{false}{60}{VIIb-A032}{The parenthetical restriction of the possibility of assigning an X-morphism to a path to type $\Delta_+$ alone is too narrow.}
\dgrclaim{false}{63}{VIIb-A040}{The displayed uncorrected map from a constant a-path to $(b\leftarrow a\to b)$ is a map in $\underline{Ch}_{-+}(b,b)$.}
\dgrclaim{false}{69}{VIIb-A055}{An equality established only in the groupoid completion Pi(X) implies the needed equality in the relative localization Pi\_Sigma(X).}
\dgrclaim{false}{76}{VIIb-B012}{For a word with r alternating runs beginning with sign epsilon, its last run has sign $(-1)^r\varepsilon$.}
\dgrclaim{false}{77}{VIIb-B015}{The concatenated interval $I_{\tau\tau'}$ is the pushout of $I_\tau$ and $I_{\tau'}$ over their two-element endpoint category $\{0,1\}$.}
\dgrclaim{open}{79--80}{VIIb-B019}{Does the fibrewise/cartesian construction extend to a functor on the total ordinary path categories?}
\dgrclaim{false}{85}{VIIb-B034}{The collapsing remainder in the alpha-th source run has type $\Delta_{\varepsilon_\alpha}^{l_{\alpha-1}}$.}
\dgrclaim{false}{86}{VIIb-B036}{The quickest map is a terminal object because it is numerically ahead of every other map.}
\dgrclaim{conjectured}{86}{VIIb-B037}{The quickest map is the greatest admissible map for pointwise numerical order.}
\dgrclaim{open}{86}{VIIb-B038}{Is every nonempty increasing endpoint-preserving interval mapping category contractible?}
\dgrclaim{false}{93}{VIIb-B052}{Reversal gives an isomorphism $P(\tau,\tau')\to P((\tau')^o,\tau^o)$.}
\dgrclaim{false}{90}{VIIb-B057}{Applying the endpoint-fixed induction to the subinterval $[1,n-1]\subseteq[0,n]$ reduces a map whose active edges have one sign and include the first and last edges to a map with only those two active edges.}
\dgrclaim{false}{90}{VIIb-B058}{A first source run that is not a single edge has length2.}
\dgrclaim{false}{102}{VIIb-C024}{The displayed transition $C_{t_n}\to C_{t_{n-1}}$ has the variance required by the sequential direct system.}
\dgrclaim{open}{102–103}{VIIb-C026}{Compare the total category $C$, the various sequential colimits $C_\infty$, and the ind-object $H_C$, in Cat hot and in Hot; determine dependence on the chosen cofinal chain.}
\dgrclaim{false}{103}{VIIb-C028}{The thin reflection of the slice satisfies $(T/t_0)_o=P/t_0=\{t:t\le t_0\}$.}
\dgrclaim{false}{103}{VIIb-C029}{Under only the stated preorder filtration, passage from $T$ to $T/t_0$ automatically gives another filtered indexing preorder and leaves the same ind-object construction available.}
\dgrclaim{conjectured}{103–104}{VIIb-C030}{If the total category were canonically recovered from the ind-object, restriction to a preorder-cofinal chain would preserve its homotopy type.}
\dgrclaim{open}{104–105}{VIIb-C031}{For a left-cofinal base change $\alpha:T^{\prime}\to T$ preserving the indicated ind-diagram, is $C\times_TT^{\prime}\to C$ a Cat-hot isomorphism or at least a weak equivalence of nerves?}
\dgrclaim{open}{107}{VIIb-C040}{Are the sets $\mathrm{Cat\,hot}(S,C)$ or $\mathrm{Hot}(S,C)$ naturally identified with $H(S)$ beyond S=e?}
\dgrclaim{conjectured}{107}{VIIb-C041}{An isomorphism on the preceding component sets need not extend to higher homotopy groups or to a Hot equivalence of the compared objects.}
\dgrclaim{conjectured}{112–113}{VIIb-C053}{For the motivating path fibration, Grothendieck suspects that $C_\infty\to C$ is generally not a Hot isomorphism.}
\dgrclaim{false}{113}{VIIb-C054}{A loop-space multiplication cannot in general be represented by a strictly associative simplicial monoid, so such strict multiplication would obstruct the proposed comparison.}
\dgrclaim{false}{115}{VIIb-C060}{For an arbitrary chosen cofinal sequence of types, concatenation has target stage n+m: $P_{t_n}(a,b)\times P_{t_m}(b,c)\to P_{t_{n+m}}(a,c)$.}
\dgrclaim{open}{108–109}{VIIc-A003}{For T=N\textasciicircum{}op, does the canonical localization C\ensuremath{\to}colim\_n C\_n induce a weak equivalence?}
\dgrclaim{open}{109}{VIIc-A006}{Is the canonical morphism C\_\ensuremath{\infty}\ensuremath{\to}C always an isomorphism in Hot?}
\dgrclaim{conjectured}{112–113}{VIIc-A015}{The comparison C\_\ensuremath{\infty}\ensuremath{\to}Ch\_\ensuremath{\Sigma}(X;a,b) is conjectured not to be a weak equivalence in general.}
\dgrclaim{false}{113}{VIIc-A016}{The based loop object \ensuremath{\Omega}(BX,a) cannot in general be modeled by a strictly associative simplicial monoid.}
\dgrclaim{false}{114}{VIIc-A020}{The object of the finite-test presheaf category associated with C\_\ensuremath{\infty} equals the filtered colimit of the stage-representable presheaves.}
\par\endgroup

\section*{VII-c \textperiodcentered\ Chapter VII, part 3 of 3, manuscript pages 118--171}
\pdfbookmark[1]{VII-c}{dgrclaims.VII-c}
\begingroup\raggedright
% DER_VII-c: the definitions of the chunk's preamble, applied inside its group
\providecommand\IfDocumentMetadataT[1]{}
\let\printedpage\relax\newcommand{\printedpage}[1]{\marginpar{\footnotesize p.~#1}}
\let\msp\relax\newcommand{\msp}[1]{\marginpar{\footnotesize [#1]}}
\let\Hotz\relax\newcommand{\Hotz}{\mathrm{Hot}^{\wedge}_{0}}
\let\Ch\relax\newcommand{\Ch}{\underline{\mathrm{Ch}}}
\let\Ty\relax\newcommand{\Ty}{\underline{\mathrm{Ty}}}
\let\CNo\relax\newcommand{\CNo}{C_{\mathbf{N}^{o}}}
\let\lab\relax\newcommand{\lab}[2][l]{\text{\footnotesize\begin{tabular}{@{}#1@{}}#2\end{tabular}}}
\let\Hot\relax\newcommand{\Hot}{\mathrm{Hot}}
\let\Hc\relax\newcommand{\Hc}{\underline{\mathrm{Hom\,cof}}}
\providecommand{\No}{}
\providecommand\No{}\renewcommand{\No}{\mathbf{N}^o}
\let\uHom\relax\newcommand{\uHom}{\underline{\mathrm{Hom}}}
\let\comp\relax\newcommand{\comp}{\operatorname{comp}}
\let\uCh\relax\newcommand{\uCh}{\underline{\mathrm{Ch}}}
\let\hot\relax\newcommand{\hot}{\mathrm{hot}}
\let\dashto\relax\newcommand{\dashto}[1]{\mathrel{\overset{#1}{\rule[0.6ex]{2em}{0.4pt}}}}
\let\CCh\relax\newcommand{\CCh}{\underline{\underline{\mathrm{Ch}}}}
\let\Xb\relax\newcommand{\Xb}{\overline{X'}}
\let\Xbb\relax\newcommand{\Xbb}{\overline{\overline{X'}}}
\let\cC\relax\newcommand{\cC}{\mathcal{C}}
\let\cX\relax\newcommand{\cX}{\mathcal{X}}
\let\NN\relax\newcommand{\NN}{\mathbf{N}^{o}}
\let\Ob\relax\newcommand{\Ob}{\operatorname{Ob}}
\let\CN\relax\newcommand{\CN}{\mathcal{C}_{\mathbf{N}^o}}

\dgrclaim{false}{117–118}{VIIc-A029}{Every natural transformation between cofinal sequences induces a weak equivalence C\_\ensuremath{\infty}(t\ensuremath{'})\ensuremath{\to}C\_\ensuremath{\infty}(t) under the general section10 hypotheses.}
\dgrclaim{open}{119}{VIIc-A034}{Does the induced Hot morphism depend only on the endpoints t,t\ensuremath{'} and not on the chosen natural transformation?}
\dgrclaim{false}{119}{VIIc-A035}{On finite-test presheaves every comparison is the canonical identification with the presheaf associated to H\_T.}
\dgrclaim{false}{120}{VIIc-A038}{The map colim\_n Hom\_Hot(\ensuremath{\xi},C\_n)\ensuremath{\to}Hom\_Hot(\ensuremath{\xi},C\_\ensuremath{\infty}) is bijective for every \ensuremath{\xi}\ensuremath{\in}Hot.}
\dgrclaim{false}{121}{VIIc-A042}{Bijectivity on all unpointed finite homotopy test sets implies that a map is an isomorphism in Hot.}
\dgrclaim{false}{121}{VIIc-A043}{Hom\_Hot(Y,\ensuremath{-}) commutes with filtered colimits whenever Y is a finite category.}
\dgrclaim{open}{121,124}{VIIc-A046}{Are all comparison functors invertible already in Cat hot?}
\dgrclaim{false}{123}{VIIc-A050}{The ind-identification implies that all induced comparison maps on C\_\ensuremath{\infty} are isomorphisms in Hot.}
\dgrclaim{false}{123}{VIIc-A053}{The restriction H\_0 of H to the source’s finite-category test subcategory is represented by the presheaf h(C\_\ensuremath{\infty}).}
\dgrclaim{open}{124–125}{VIIc-A056}{Does a binary path composition C\_\ensuremath{\infty}(a,b)×C\_\ensuremath{\infty}(b,c)\ensuremath{\to}C\_\ensuremath{\infty}(a,c) exist canonically in Hot?}
\dgrclaim{false}{126}{VIIc-A063}{Two unrelated interval types can be compared by their categorical product \ensuremath{\tau}\_0×\ensuremath{\tau}\ensuremath{'}\_0 as another interval type.}
\dgrclaim{open}{127}{VIIc-A065}{Is the constructed binary operation associative in Hot?}
\dgrclaim{open}{127}{VIIc-A066}{Is Hot\ensuremath{\to}widehat(Hot\_0) faithful, or even fully faithful?}
\dgrclaim{open}{127}{VIIc-A069}{Are the canonical maps H\ensuremath{\to}C\_\ensuremath{\infty} epimorphisms in Ind Hot, and do they remain so after products?}
\dgrclaim{conjectured}{128}{VIIc-A072}{For a non-Kan simplicial set K, the strict endpoint fibres of K\textasciicircum{}\{\ensuremath{\Delta}\textasciicircum{}1\}\ensuremath{\to}K×K need not be the derived path spaces.}
\dgrclaim{conjectured}{128}{VIIc-A073}{A correct simple path construction for finite simplicial sets would imply a finite computation of all their homotopy groups.}
\dgrclaim{conjectured}{128}{VIIc-A074}{For \ensuremath{\Sigma}=Fl(X), a weak equivalence X\ensuremath{\to}Xbar should induce Ch\_\ensuremath{\infty}(X;a,b)\ensuremath{\simeq}Ch\_\ensuremath{\infty}(Xbar;a,b).}
\dgrclaim{conjectured}{129}{VIIc-B002}{The external-Hom parameter functor F\_\{a,b\}(Y)=Hom\_\{D(Y)\}(a\_Y,b\_Y) should send nerve weak equivalences of Y to bijections.}
\dgrclaim{false}{129}{VIIc-B005}{The source infers under its locality hypothesis that the finite-parameter evaluation of H(a,b) is [Y,Ch\_\{infty,Sigma\}(a,b)] in Hot, using finite-category compactness.}
\dgrclaim{conjectured}{130}{VIIc-B006}{A nerve weak equivalence f:X\ensuremath{\to}Xprime should induce weak equivalences Ch\_infty(X;a,b)\ensuremath{\to}Ch\_infty(Xprime;fa,fb), with all arrows allowed.}
\dgrclaim{false}{130}{VIIc-B009}{For Thomason-fibrant X, fixed-endpoint paths of type Delta\_+ might suffice to model Ch\_infty(X;a,b).}
\dgrclaim{open}{130}{VIIc-B011}{Is Hot\ensuremath{\to}Hot\_0\textasciicircum{}wedge faithful, so that finite tests would prove associativity?}
\dgrclaim{conjectured}{132}{VIIc-B018}{Noncontractible components of Omega(X,a) would make the invariance condition almost never hold.}
\dgrclaim{false}{133}{VIIc-B022}{For a fixed finite circle category Y, degrees realized by functors Y\ensuremath{\to}Ch\_infty(X;a,a) are bounded by the number of its arrows, hence by four in the displayed example.}
\dgrclaim{false}{133}{VIIc-B023}{Footnote28 proposes the sharper bound -2\ensuremath{\leq}degree\ensuremath{\leq}2 for the displayed four-arrow circle.}
\dgrclaim{conjectured}{133}{VIIc-B024}{By enlarging alternating circle models Y\_n, one can realize degrees of arbitrarily large magnitude.}
\dgrclaim{false}{134}{VIIc-B026}{The proposed bound on degrees proves that pi\_0 Fun(Y\_2,Omega)\ensuremath{\to}pi\_0 Fun(Y\_n,Omega) is not bijective.}
\dgrclaim{conjectured}{134}{VIIc-B027}{The same finite-size obstruction should apply whenever higher homotopy groups of X provide infinitely many classes.}
\dgrclaim{conjectured}{134}{VIIc-B028}{Products of alternating intervals with their boundary identified should provide categorical models of S\textasciicircum{}i.}
\dgrclaim{false}{134}{VIIc-B029}{A fixed finite categorical sphere Y\_n should realize only a bounded subset of pi\_i Omega(X,a).}
\dgrclaim{conjectured}{134–135}{VIIc-B030}{Every homotopy class should be represented after choosing a sufficiently large categorical sphere and a degree-one map to a smaller model.}
\dgrclaim{false}{135}{VIIc-B034}{The absolute comparison already fails for Y=[1].}
\dgrclaim{false}{136}{VIIc-B035}{For any category Y and Thomason-fibrant C, BFun(Y,C) represents RMap(BY,BC).}
\dgrclaim{false}{136}{VIIc-B036}{If C is Thomason-fibrant then Fun(Y,C) is Thomason-fibrant for arbitrary Y.}
\dgrclaim{false}{136}{VIIc-B037}{Footnote29 requires Thomason cofibrations to remain cofibrations after product with arbitrary Y.}
\dgrclaim{false}{136}{VIIc-B038}{For fibrant C and arbitrary Z,Y, pi\_0 Fun(Z×Y,C) computes [B(Z×Y),BC].}
\dgrclaim{conjectured}{136}{VIIc-B040}{Categorical path objects should have the correct homotopy type without first making X Thomason-fibrant.}
\dgrclaim{false}{139}{VIIc-B050}{For every f, bar f in this factorization should be both fibrant and a weak equivalence.}
\dgrclaim{false}{139}{VIIc-B051}{The map alpha should be a Thomason cofibration.}
\dgrclaim{false}{139}{VIIc-B052}{The map bar f is expected to be fibrant in the model-category sense.}
\dgrclaim{conjectured}{139}{VIIc-B053}{Item1: p:Ch\_infty(X)\ensuremath{\to}X×X should be smooth, proper and a HOT-fibration.}
\dgrclaim{conjectured}{139}{VIIc-B055}{Item2 additionally asserts contractible fibres for sigma and beta, hence trivial HOT-fibrations.}
\dgrclaim{conjectured}{139}{VIIc-B056}{Item3: bar f:C(f)\ensuremath{\to}X should be smooth, proper and a HOT-fibration for every f.}
\dgrclaim{conjectured}{140–141}{VIIc-B062}{For a path c:a\ensuremath{\to}b, transport should give an isomorphism bar c:hot(C(f)\_a)\ensuremath{\to}hot(C(f)\_b).}
\dgrclaim{conjectured}{140–141}{VIIc-B063}{Transport should be compatible with concatenation of base paths.}
\dgrclaim{conjectured}{141}{VIIc-B066}{Transport should depend only on the homotopy class of c in the path category.}
\dgrclaim{conjectured}{141}{VIIc-B067}{The fibre homotopy types and transport should define a functor Pi(X)\ensuremath{\to}Hot.}
\dgrclaim{conjectured}{141}{VIIc-B068}{The monodromy construction should apply to every HOT-fibration, not just this mapping-path replacement.}
\dgrclaim{conjectured}{141}{VIIc-B069}{For a D-fibration and coefficient A\ensuremath{\in}D(e), fibrewise D-homology and D-cohomology should give two Pi(X)-functors to D(e).}
\dgrclaim{open}{141–142}{VIIc-B070}{Footnotes33–34 ask for the effect of transport on components of the fibres.}
\dgrclaim{conjectured}{142}{VIIc-B072}{The nerve of Ch\_infty(X) over NX×NX should be weakly equivalent in the slice to a simplicial path object.}
\dgrclaim{conjectured}{142}{VIIc-B073}{The fixed-endpoint categorical path categories should have canonically the same homotopy types as simplicial path spaces.}
\dgrclaim{open}{142}{VIIc-B074}{Does a functorial classical simplicial path object exist for every simplicial set?}
\dgrclaim{false}{145}{VIIc-B084}{The significant edge indices of the t\_n stage are [-2i+1,2i], with 4i edges, as printed.}
\dgrclaim{false}{146}{VIIc-B088}{For even n, tau\_\{n,m\}=(+-)\textasciicircum{}\{(m-n)/2\} if m is even, and the printed (+-)\textasciicircum{}\{(m-n)/2-1\}+ if m is odd.}
\dgrclaim{false}{147}{VIIc-B089}{For odd n, tau\_\{n,m\}=(-+)\textasciicircum{}\{(m-n)/2\} if m is odd, and the printed (-+)\textasciicircum{}\{(m-n)/2-1\}- if m is even.}
\dgrclaim{false}{147}{VIIc-B090}{For n<m the initial sign is (-1)\textasciicircum{}n and the terminal sign is printed as (-1)\textasciicircum{}m.}
\dgrclaim{open}{150}{VIIc-C004}{Is the restricted Yoneda functor $\mathrm{Hot}\to\mathrm{Hot}_0^{\wedge}$ faithful, so that equality of the two finite-test images implies $T^2=1$?}
\dgrclaim{conjectured}{151}{VIIc-C016}{The finite-stage composition operations should make the comparison square commute on associated ind-objects.}
\dgrclaim{false}{152}{VIIc-C017}{If the comma transport becomes an isomorphism in the unbased finite-test presheaf category, it must be an isomorphism in Hot.}
\dgrclaim{conjectured}{153}{VIIc-C023}{The mapping-path projection $\bar f:C(f)\to X$ should be smooth, proper, and a homotopy fibration.}
\dgrclaim{conjectured}{155}{VIIc-C031}{The finite-stage $v_*$ are compatible up to categorical homotopy and induce a morphism of ind-objects in Cat hot, then a map of finite-test presheaves.}
\dgrclaim{conjectured}{155}{VIIc-C033}{The stagewise $v^*$ likewise induce an ind-morphism and hence a map on finite-test presheaves.}
\dgrclaim{conjectured}{156}{VIIc-C034}{The two ind-morphisms $v_*,v^*$ are inverse, because the stagewise backtracking operations are homotopic to the identity.}
\dgrclaim{conjectured}{156}{VIIc-C035}{The ind-object comparison square with the comma transport $\bar v$ commutes.}
\dgrclaim{false}{156}{VIIc-C037}{Invertibility of $\bar v$ after restriction to finite unbased tests implies its invertibility in Hot.}
\dgrclaim{open}{156}{VIIc-C038}{$\bar f$ is a trivial homotopy fibration if and only if f is a weak equivalence; a categorical long exact sequence seems necessary to prove it.}
\dgrclaim{conjectured}{157}{VIIc-C039}{The canonical maps $\Ch_\infty(X)\to\Ch(X)$ and $\Ch_\infty(X;a,b)\to\Ch(X;a,b)$ in Hot are isomorphisms.}
\dgrclaim{conjectured}{157}{VIIc-C040}{The total comparison is an isomorphism in the relative homotopy category over $X\times X$, giving the fixed-endpoint comparisons.}
\dgrclaim{false}{157}{VIIc-C042}{$\Ch(X)\to X\times X$ is both smooth and proper.}
\dgrclaim{false}{158}{VIIc-C045}{$\mathcal C_{\mathbf N^o}\to X\times X$ is smooth and proper.}
\dgrclaim{false}{161}{VIIc-C058}{Smoothness of $\Ch(X^o)\to X^o\times X^o$ implies properness of $\Ch(X)\to X\times X$.}
\dgrclaim{false}{166}{VIIc-C073}{The composite $j^{\prime} i^{\prime}:X\to\mathcal C_{\mathbf N^o}$ is aspherical.}
\dgrclaim{false}{168}{VIIc-C078}{The padded formula $\varphi(c^{\prime},w^{\prime})=(\widetilde c^{\prime},\widetilde w^{\prime})$ defines a functor from the sequential lift category to the sequential fixed-endpoint slice.}
\dgrclaim{false}{169}{VIIc-C081}{The proposed sequential endofunctor $\psi\varphi$ admits a natural transformation to the identity.}
\dgrclaim{false}{169}{VIIc-C083}{Because the system of interval types is self-opposite, smoothness for X\textasciicircum{}o proves properness of $\mathcal C_{\mathbf N^o}(X)\to X\times X$.}
\par\endgroup

\section*{VIII-X \textperiodcentered\ Chapter VIII, manuscript pages 1--9; Chapter IX, manuscript pages 1--25; Chapter X, manuscript pages 1--5}
\pdfbookmark[1]{VIII-X}{dgrclaims.VIII-X}
\begingroup\raggedright
% DER_VIII-X: the definitions of the chunk's preamble, applied inside its group
\let\printedpage\relax\newcommand{\printedpage}[1]{\marginpar{\footnotesize p.~#1}}
\let\msp\relax\newcommand{\msp}[1]{\marginpar{\footnotesize [#1]}}
\let\Grp\relax\newcommand{\Grp}{\underline{\mathrm{Grp}}}
\let\uHom\relax\newcommand{\uHom}{\underline{\mathrm{Hom}}}
\let\uGrp\relax\newcommand{\uGrp}{\underline{\mathrm{Grp}}}
\let\uuGrp\relax\newcommand{\uuGrp}{\underline{\underline{\mathrm{Grp}}}}
\let\fl\relax\newcommand{\fl}{\longrightarrow}
\let\Ch\relax\newcommand{\Ch}{\mathrm{Ch}}
\let\Hot\relax\newcommand{\Hot}{\mathrm{Hot}}
\let\Hom\relax\newcommand{\Hom}{\mathrm{Hom}}
\let\HotO\relax\newcommand{\HotO}{\mathrm{Hot}_0^{\wedge}}
\let\MOP\relax\newcommand{\MOP}{\mathcal{M}_{0\,\Pi}}
\let\M\relax\newcommand{\M}{\mathcal{M}}
\let\uCh\relax\newcommand{\uCh}{\underline{\mathrm{Ch}}}
\let\Homext\relax\newcommand{\Homext}{\mathrm{Hom\,ext}}
\let\Lotimes\relax\newcommand{\Lotimes}{\overset{\mathbf{L}}{\otimes}}
\let\cercle\relax\newcommand{\cercle}[1]{\tikz[baseline=(c.base)]\node[draw,circle,inner sep=1pt](c){#1};}
\let\Chu\relax\newcommand{\Chu}{\underline{\mathrm{Ch}}}
\let\D\relax\newcommand{\D}{\mathbf{D}}
\let\Cat\relax\newcommand{\Cat}{\mathrm{Cat}}
\let\Fl\relax\newcommand{\Fl}{\operatorname{Fl}}
\let\Ob\relax\newcommand{\Ob}{\operatorname{Ob}}
\let\bDelta\relax\newcommand{\bDelta}{\boldsymbol{\Delta}}

\dgrclaim{conjectured}{VIII:4}{VIII-016}{For X\ensuremath{\to}S\ensuremath{\to}T and every groupoid fibration \ensuremath{\Gamma} on X, (fp)\_!\ensuremath{\Gamma}\ensuremath{\simeq}f\_!p\_!\ensuremath{\Gamma}, naturally and coherently.}
\dgrclaim{conjectured}{VIII:7}{VIII-024}{For a cartesian square with invertible \ensuremath{\gamma}, the comparison f\ensuremath{'}\_!q*G\ensuremath{\to}p*f\_!G is an equivalence if f is 1-smooth or p is 1-proper.}
\dgrclaim{false}{VIII:7}{VIII-029}{For i\ensuremath{\in}\{0,1,2\}, i-faithfulness of f* is equivalent to i-faithfulness of \ensuremath{\Gamma}(T,H)\ensuremath{\to}\ensuremath{\Gamma}(S,f*H) for every H, taking only G=e\_T.}
\dgrclaim{conjectured}{VIII:8}{VIII-036}{If f is (-1)-proper and (-1)-aspherical, or (-1)-smooth and (-1)-coaspherical, then f* is 0-faithful.}
\dgrclaim{conjectured}{VIII:8}{VIII-037}{If f is 0-proper and 0-aspherical, or 0-smooth and 0-coaspherical, then f* is 1-faithful.}
\dgrclaim{conjectured}{VIII:8}{VIII-038}{If f is 1-proper and 1-aspherical, or 1-smooth and 1-coaspherical, then f* is 2-faithful.}
\dgrclaim{conjectured}{VIII:8–9}{VIII-040}{Under either hypothesis of Theorem 3 with i=2, G on S is in the essential image of f* exactly when G restricted to each strict fibre S\_t is equivalent to a constant groupoid fibration.}
\dgrclaim{false}{VIII:9}{VIII-041}{For 1-connected Z, a groupoid fibration on Z is constant up to fibred equivalence if and only if every transport functor is an equivalence.}
\dgrclaim{false}{VIII:9}{VIII-043}{Under Theorem 3 with i=2, G descends along f if and only if every f-vertical arrow induces an equivalence on its fibres.}
\dgrclaim{false}{IX:4}{IX-011}{Arbitrary arrows u:a-prime->a and v:b->b-prime induce an isomorphism of Hom objects.}
\dgrclaim{false}{IX:4}{IX-013}{The displayed Ind Cat hot bifunctor has domain M\_Pi x M\_Pi.}
\dgrclaim{open}{IX:9-10}{IX-R01}{For the full subcategory Hot\_0 of types of finite categories, is the restricted Yoneda functor psi:Hot -> PSh(Hot\_0) fully faithful (and likewise Ind psi)?}
\dgrclaim{open}{IX:13}{IX-037}{Ch composition is associative in Hot.}
\dgrclaim{conjectured}{IX:14}{IX-040}{Quillen pi\_1(f,g) agrees with arrows between f and g in the fundamental groupoid of Ch\_infty(a,b).}
\dgrclaim{conjectured}{IX:15}{IX-041}{The homotopy groupoid construction extends naturally from actual arrows to all stabilized zigzags.}
\dgrclaim{open}{IX:15}{IX-044}{The path compositions should lift to pointed Hom objects with associative composition and units.}
\dgrclaim{false}{IX:16}{IX-046}{There are canonical arrows in both directions between these zero representatives, identity at every vertex except degree zero.}
\dgrclaim{conjectured}{IX:17}{IX-049}{The resulting exact path pairings are associative in the pointed homotopy category and perhaps in pointed Cat hot.}
\dgrclaim{false}{IX:17}{IX-050}{Cartesian pointed spaces instantiate the earlier monoidal axiom with pi\_0=Hom(unit,-).}
\dgrclaim{false}{IX:17}{IX-051}{Forgetting basepoints preserves the corrected monoidal product.}
\dgrclaim{conjectured}{IX:17}{IX-052}{Pointed CH should exist in Ind Hot\_bullet and Ind Cat hot\_bullet.}
\dgrclaim{false}{IX:18}{IX-054}{The displayed truncated path a->0<-b represents the desired distinguished element.}
\dgrclaim{open}{IX:18}{IX-055}{The zero object should remain a zero object in M[W\textasciicircum{}\{-1\}].}
\dgrclaim{open}{IX:18}{IX-056}{Find conditions under which localization preserves finite coproducts and finite products.}
\dgrclaim{conjectured}{IX:18}{IX-059}{An additional five-lemma-type condition may be needed for the additive localization.}
\dgrclaim{conjectured}{IX:19}{IX-061}{The external Hom of an additive relative category should lift to D\textasciicircum{}\{<=0\}(k-Mod).}
\dgrclaim{conjectured}{IX:19}{IX-065}{Derived Hom objects should admit associative composition over derived tensor and units.}
\dgrclaim{conjectured}{IX:21}{IX-070}{A full theory of strong derived k-bilinear Hom composition is expected.}
\dgrclaim{false}{IX:21}{IX-071}{pi\_1(a,b) is written pi\_1(a,b;0\_a,0\_b).}
\dgrclaim{conjectured}{IX:21}{IX-072}{Quillen pi\_1(a,b) is canonically pi\_1 of the pointed Ch\_infty(a,b).}
\dgrclaim{conjectured}{IX:22}{IX-080}{The two formulations of test (b) are equivalent up to minor qualifications.}
\dgrclaim{conjectured}{IX:22}{IX-081}{Tests (a) and (b) are a decisive criterion for a homotopy theory in M.}
\dgrclaim{false}{IX:23}{IX-082}{The stated tests cannot hold for a finite M except in the trivial cases.}
\dgrclaim{false}{IX:23}{IX-083}{If M is empty, its localization is equivalent to the point.}
\dgrclaim{conjectured}{IX:23}{IX-084}{Test (c) asks for Homext(Sigma a,b)=Homext(a,Omega b)=Omega Homext(a,b) as pointed homotopy types.}
\dgrclaim{conjectured}{IX:23}{IX-086}{External Hom preserves finite products in the second variable and sends finite coproducts in the first variable to products.}
\dgrclaim{conjectured}{IX:23}{IX-087}{Omega should extend to a left exact endomorphism of the derivator, and Sigma dually.}
\dgrclaim{conjectured}{IX:23}{IX-088}{Test (b) is necessary for the loop-space equivalences in test (c).}
\dgrclaim{conjectured}{IX:24}{IX-090}{The essential condition on the path factorization is that p be a fibration in a suitable sense.}
\dgrclaim{open}{IX:25}{IX-091}{A cofibration condition on i may be useful; fibrancy of S is queried.}
\dgrclaim{open}{IX:25}{IX-093}{The constructed object should represent pi\_1 Ch\_infty(T,S) naturally in T after localization.}
\dgrclaim{conjectured}{X:3}{X022}{The categorical assertions of Chapter I that use only weak equivalences and their formal properties should extend from W\_D to abstract localizers W.}
\dgrclaim{open}{X:3}{X023}{Under suitable conditions on W, the pointwise localization A |-> Cat\textasciicircum{}\{A\textasciicircum{}op\}[(W\textasciicircum{}\{A\textasciicircum{}op\})\textasciicircum{}\{-1\}] should be a derivator.}
\dgrclaim{open}{X:3}{X024}{One should be able to realize a given localizer W as W\_D for a derivator constructed from (Cat,W).}
\dgrclaim{conjectured}{X:5}{X038}{The derivators useful in practice seem to satisfy La.}
\dgrclaim{open}{X:5}{X042}{Do Lb and Lc hold for every W\_D under consideration, and what are their concrete conditions on W?}
\par\endgroup

\section*{XI-XIIa \textperiodcentered\ Chapter XI, manuscript pages 127--136; Chapter XII, sections 1--4, manuscript pages 1--58}
\pdfbookmark[1]{XI-XIIa}{dgrclaims.XI-XIIa}
\begingroup\raggedright
% DER_XI-XIIa: the definitions of the chunk's preamble, applied inside its group
\let\printedpage\relax\newcommand{\printedpage}[1]{\marginpar{\footnotesize p.~#1}}
\let\msp\relax\newcommand{\msp}[1]{\marginpar{\footnotesize [#1]}}
\let\Hot\relax\newcommand{\Hot}{\mathrm{Hot}}
\let\Cat\relax\newcommand{\Cat}{\mathrm{Cat}}
\let\Ob\relax\newcommand{\Ob}{\operatorname{Ob}}
\let\Ch\relax\newcommand{\Ch}{\underline{\mathrm{Ch}}}
\let\Dl\relax\newcommand{\Dl}{\boldsymbol{\Delta}}
\let\fl\relax\newcommand{\fl}{\longrightarrow}
\let\hot\relax\newcommand{\hot}{\mathrm{hot}}
\let\uHom\relax\newcommand{\uHom}{\underline{\mathrm{Hom}}}
\let\id\relax\newcommand{\id}{\mathrm{id}}
\let\pr\relax\newcommand{\pr}{\mathrm{pr}}
\let\Fl\relax\newcommand{\Fl}{\mathrm{Fl}}
\let\card\relax\newcommand{\card}{\operatorname{card}}
\let\uFib\relax\newcommand{\uFib}{\underline{\mathrm{Fib}}}
\let\HOT\relax\newcommand{\HOT}{\mathrm{HOT}}
\let\HOTlc\relax\newcommand{\HOTlc}{\mathrm{HOT}^{\mathrm{lc}}_W}
\let\Hom\relax\newcommand{\Hom}{\mathrm{Hom}}
\let\Wd\relax\newcommand{\Wd}{W^{\mathrm{d}}_S}
\let\Wf\relax\newcommand{\Wf}{W^{\mathrm{f}}_S}
\let\dashedbox\relax\newcommand{\dashedbox}[1]{\tikz[baseline=(n.base)]\node[draw,dashed,inner sep=2pt](n){#1};}
\let\Fib\relax\newcommand{\Fib}{\mathrm{Fib}}
\let\Parf\relax\newcommand{\Parf}{\mathrm{Parf}}

\dgrclaim{false}{XI:129}{XI-017}{For an arbitrary locally constant nonabelian F, lifting \ensuremath{\eta}\ensuremath{\in}H0(Y,R\textasciicircum{}1f\_*F) is governed by a fixed-band obstruction in H2(Y,f\_*F), as written.}
\dgrclaim{open}{XII:1}{XII-I-001}{Find intrinsic conditions on W implying W\ensuremath{\subseteq} W\_\ensuremath{\infty}, without defining those conditions through W\_\ensuremath{\infty}.}
\dgrclaim{open}{XII:5}{XII-I-028}{Are Ch(X)\ensuremath{\to} X\ensuremath{\times} X and Ch\_\ensuremath{\infty}(X)\ensuremath{\to} X\ensuremath{\times} X W-fibrations for a general W?}
\dgrclaim{open}{XII:5}{XII-I-032}{Does the left test alone imply preservation over W-aspherical bases without additional axioms?}
\dgrclaim{open}{XII:13}{XII-I-057}{Do W(1)-W(5) already imply W=W\_\ensuremath{\infty}, or is an additional descent/limit axiom necessary?}
\dgrclaim{conjectured}{XII:14}{XII-I-071}{Even the strongest W(4), with W(1)-W(3) and W(5), probably does not force W=W\_infty.}
\dgrclaim{conjectured}{XII:15}{XII-I-076}{W(1)-W(6), using the weakest fibre axiom, imply W=W\_infty.}
\dgrclaim{open}{XII:16}{XII-I-077}{Perhaps closure under retracts must be added to W(1) for the equality conjecture.}
\dgrclaim{false}{XII:17}{XII-I-084}{The formula pi\_n(F)=coker(pi\_\{n+1\}(E)\ensuremath{\to} pi\_\{n+1\}(B)) is asserted without excluding n=0.}
\dgrclaim{false}{XII:18,20}{XII-II-002}{W\_D\ensuremath{\supseteq}W\_infty and W(4b'') imply W\_D=W\_infty without another stated qualification; the analogous inference is repeated for W in the HOT\_W discussion.}
\dgrclaim{false}{XII:19}{XII-II-005}{The left/right conditions are the same as the earlier D-cohomological/D-homological fibrations.}
\dgrclaim{conjectured}{XII:19}{XII-II-007}{Every genuinely useful derivator/localizer should satisfy the indicated constancy condition.}
\dgrclaim{open}{XII:20}{XII-II-023}{Did the preceding construction and properties of f\_! for HOT\_W actually use the working assumption W\_infty\ensuremath{\subseteq}W?}
\dgrclaim{conjectured}{XII:26}{XII-III-A-024}{Aspherical functors should be preserved by smooth base change, and coaspherical functors by proper base change.}
\dgrclaim{false}{XII:36}{XII-III-A-057}{It suffices, to define functorial covariant transport universally, to require only that the terminal-fibre inclusion in every pullback over [1] belongs to W.}
\dgrclaim{false}{XII:36}{XII-III-A-058}{The analogous initial-fibre edge-only condition suffices to define universal contravariant transport.}
\dgrclaim{conjectured}{XII:40}{XII-III-B-003}{Among the provisional fibration notions, only ultrafibrations appear likely to be stable under composition.}
\dgrclaim{open}{XII:40}{XII-III-B-005}{Find minimal hypotheses on W making the stronger fibration variants equivalent.}
\dgrclaim{open}{XII:40}{XII-III-B-006}{Does the interval endpoint condition a0 alone imply the remaining conditions for W\ensuremath{\infty}?}
\dgrclaim{conjectured}{XII:40}{XII-III-B-008}{W-perfect\ensuremath{\Rightarrow}hot-strong W-fibration (circled implication).}
\dgrclaim{conjectured}{XII:40}{XII-III-B-009}{Footnote 48: the hot-strong condition should be equivalent to both the left and right fibration conditions.}
\dgrclaim{open}{XII:42}{XII-III-B-017}{If f is W-perfect and S categorically contractible, is X\_s\ensuremath{\to}X in W?}
\dgrclaim{false}{XII:44}{XII-III-B-025}{The first displayed transport is F\ensuremath{\circ}f\ensuremath{\to}F\ensuremath{\circ}g for \ensuremath{\alpha}:f\ensuremath{\to}g.}
\dgrclaim{false}{XII:45}{XII-III-B-026}{The classifying functor F itself is 'in W fibre by fibre'.}
\dgrclaim{conjectured}{XII:46}{XII-III-B-032}{These results should use at most full W(2) rather than W(2 bis).}
\dgrclaim{conjectured}{XII:48}{XII-III-B-040}{The stronger complement: pullback of an arbitrary categorical homotopy equivalence along the family is a W-map.}
\dgrclaim{open}{XII:54}{XII-IV-011}{Does W(4c\ensuremath{'}) imply W(4c)?}
\dgrclaim{open}{XII:54}{XII-IV-012}{Does W(4c) imply W(4b)?}
\dgrclaim{conjectured}{XII:54}{XII-IV-014}{The weaker split test W(4b\ensuremath{'}) should suffice for the same conclusion.}
\dgrclaim{conjectured}{XII:54}{XII-IV-015}{W(4b\ensuremath{''}), perhaps with W(6 bis), should likewise suffice.}
\dgrclaim{conjectured}{XII:54}{XII-IV-017}{The equivalence should also include W(4b\ensuremath{'}) and W(4b\ensuremath{''}).}
\dgrclaim{false}{XII:58}{XII-IV-030}{Proposition 3: under W(4a), a strict Cartesian map of split fibrations with W-fibrant fibre maps is W-fibrant.}
\par\endgroup

\section*{XIIb-a \textperiodcentered\ Chapter XII, section 5, part 1 of 2, manuscript pages 59--122}
\pdfbookmark[1]{XIIb-a}{dgrclaims.XIIb-a}
\begingroup\raggedright
% DER_XIIb-a: the definitions of the chunk's preamble, applied inside its group
\let\printedpage\relax\newcommand{\printedpage}[1]{\marginpar{\footnotesize p.~#1}}
\let\msp\relax\newcommand{\msp}[1]{\marginpar{\footnotesize [#1]}}
\let\uHom\relax\newcommand{\uHom}{\underline{\mathrm{Hom}}}
\let\Hom\relax\newcommand{\Hom}{\mathrm{Hom}}
\let\Cat\relax\newcommand{\Cat}{\mathrm{Cat}}
\let\Ens\relax\newcommand{\Ens}{\mathrm{Ens}}
\let\id\relax\newcommand{\id}{\mathrm{id}}
\let\Ob\relax\newcommand{\Ob}{\mathrm{Ob}\,}
\let\circled\relax\newcommand{\circled}[1]{\tikz[baseline=(c.base)]\node[draw,circle,inner sep=1pt](c){#1};}
\let\cercle\relax\newcommand{\cercle}[1]{\tikz[baseline=(c.base)]\node[draw,circle,inner sep=1pt](c){#1};}
\let\defeq\relax\newcommand{\defeq}{\overset{\text{déf}}{=}}
\let\upi\relax\newcommand{\upi}{\underline{\pi}}
\let\uCh\relax\newcommand{\uCh}{\underline{\mathrm{Ch}}}
\let\bDelta\relax\newcommand{\bDelta}{\boldsymbol{\Delta}}
\let\hra\relax\newcommand{\hra}{\lhook\joinrel\longrightarrow}

% Reader layout additions only; all edition slips in the body are retained.

\dgrclaim{false}{59}{XIIba-A-003}{With x<=y iff Hom\_X(x,y) is nonempty and opens the lower sets, the source asserts x<=y iff x belongs to the closure of \{y\}.}
\dgrclaim{false}{61}{XIIba-A-012}{The induced triple represents on all triples Y=(Y,V,S) the functor Hom(V,U) x Hom(S,T), as asserted in displayed (*).}
\dgrclaim{false}{61}{XIIba-A-013}{The displayed unqualified functor-category version (**) of the same representability is an isomorphism.}
\dgrclaim{open}{76}{XIIba-A-053}{Does W(7bis) alone imply that W-asphericity of the intersection and union forces W-asphericity of both pieces?}
\dgrclaim{false}{78}{XIIba-A-058}{The ordinary weak homotopy equivalences W\_infty satisfy the full axiom W(7).}
\dgrclaim{false}{79-80}{XIIba-A-062}{Reverse excision is asserted to be a special property of HOT, while the dual derivator fails it.}
\dgrclaim{open}{94–97}{XIIb-a-B035}{Does total W-asphericity of H imply total W-asphericity of U/H?}
\dgrclaim{false}{97–100}{XIIb-a-B045}{Assuming W(7 bis), a union V=V\ensuremath{'}\ensuremath{\cup}V\ensuremath{''} of two sieves is claimed to be totally W-aspherical whenever V\ensuremath{'},V\ensuremath{''},V\ensuremath{'}\ensuremath{\cap}V\ensuremath{''} are totally W-aspherical, by preservation of the subobject square under f!.}
\dgrclaim{false}{98–100}{XIIb-a-B047}{Under the standing W(7 bis) hypothesis, the assertion of Proposition 9(b) is repaired by categorical retractions V\ensuremath{'}\ensuremath{\to}V\ensuremath{'}\ensuremath{\cap}V\ensuremath{''} and V\ensuremath{''}\ensuremath{\to}V\ensuremath{'}\ensuremath{\cap}V\ensuremath{''}.}
\dgrclaim{false}{98}{XIIb-a-B048}{Under the same gluing hypothesis, the analogous union assertion is true for a pushout square of subpresheaves in an arbitrary presheaf category merely because the three smaller objects are totally W-aspherical.}
\dgrclaim{false}{98}{XIIb-a-B049}{Under the same hypotheses, one of the two categorical retractions already suffices to imply the asserted total-asphericity conclusion.}
\dgrclaim{false}{99}{XIIb-a-B051}{The existence of those categorical retractions implies that f! sends the square of the associated subterminal presheaves to a square of monomorphisms.}
\dgrclaim{false}{99}{XIIb-a-B052}{The same retractions imply f!(V\ensuremath{'}\ensuremath{\cap}V\ensuremath{''})=f!V\ensuremath{'}\ensuremath{\cap}f!V\ensuremath{''}.}
\dgrclaim{false}{99–100}{XIIb-a-B053}{Consequently the four categories of elements of these Kan extensions form a Cartesian and cocartesian square of sieve inclusions.}
\dgrclaim{false}{100}{XIIb-a-B054}{Assuming W(7 bis), the outgoing fork \ensuremath{\Psi}=(0\ensuremath{\to}a,0\ensuremath{\to}b) is totally W-aspherical.}
\dgrclaim{false}{100}{XIIb-a-B055}{Assuming W(7 bis), every finite ordered category with an initial object is totally W-aspherical.}
\dgrclaim{false}{100}{XIIb-a-B056}{A category is totally W-aspherical only if it has a final object.}
\dgrclaim{open}{103}{XIIba-C005}{If $Y$ is $W_\infty$-aspherical, must $\underline{\mathrm{Ch}}(Y;y_0,-)$ be totally $W_\infty$-aspherical?}
\dgrclaim{open}{104}{XIIba-C009}{Can one similarly find a sieve pushout with $p\in W$ but $q\notin W$?}
\dgrclaim{open}{111}{XIIba-C035}{Must a sieve that is a W-cofibration, or even a couniversal W-equivalence, be W-dwyerienne?}
\dgrclaim{open}{111}{XIIba-C036}{Must such a sieve even satisfy the original Dwyer right-adjoint condition?}
\dgrclaim{open}{111--112}{XIIba-C039}{Is there a useful Dwyer notion for arbitrary full inclusions which includes both the sieve and cosieve cases?}
\dgrclaim{false}{112}{XIIba-C040}{The proof says that W-dwyerienne stability makes $i'$ W-aspherical and in W.}
\dgrclaim{open}{113}{XIIba-C045}{Are W-dwyerienne and totally W-dwyerienne inclusions stable under composition?}
\dgrclaim{open}{114}{XIIba-C047}{Are strictly W-aspherical functors and totally W-aspherical functors stable under composition?}
\dgrclaim{open}{114}{XIIba-C048}{Seek a class of sieve W-cofibrations stable under the endpoint pushout product with $\{0,2n\}\to\Phi_n$, n>=2.}
\dgrclaim{open}{119}{XIIba-C064}{Request: make the map $(\lambda_\alpha)_!$ in the preceding transition formula explicit.}
\dgrclaim{open}{114 (continued before 115)}{XIIbb-A001}{Find a class of sieve inclusions that are W-cofibrations and is closed under the endpoint pushout-product U×Phi\_n union X×\{0,2n\} -> X×Phi\_n for every n>=2.}
\par\endgroup

\section*{XIIb-b \textperiodcentered\ Chapter XII, section 5, part 2 of 2, manuscript pages 123--169}
\pdfbookmark[1]{XIIb-b}{dgrclaims.XIIb-b}
\begingroup\raggedright
% DER_XIIb-b: the definitions of the chunk's preamble, applied inside its group
\providecommand\IfDocumentMetadataT[1]{\IfDocumentMetadataTF{#1}{}}
\let\printedpage\relax\newcommand{\printedpage}[1]{\marginpar{\footnotesize p.~#1}}
\let\msp\relax\newcommand{\msp}[1]{\marginpar{\footnotesize [#1]}}
\let\Hot\relax\newcommand{\Hot}{\mathrm{Hot}}
\let\HOT\relax\newcommand{\HOT}{\mathrm{HOT}}
\let\Cat\relax\newcommand{\Cat}{\mathrm{Cat}}
\let\cX\relax\newcommand{\cX}{\mathcal{X}}
\let\sX\relax\newcommand{\sX}{\mathscr{X}}
\let\tX\relax\newcommand{\tX}{\tilde{\mathscr{X}}}
\let\Hom\relax\newcommand{\Hom}{\mathrm{Hom}}
\let\Fl\relax\newcommand{\Fl}{\mathrm{Fl}}
\let\tsX\relax\newcommand{\tsX}{\tilde{\mathscr{X}}}
\let\circled\relax\newcommand{\circled}[1]{\tikz[baseline=(char.base)]{\node[shape=circle,draw,inner sep=1pt] (char) {#1};}}
\let\uHom\relax\newcommand{\uHom}{\underline{\mathrm{Hom}}}
\let\tld\relax\newcommand{\tld}{^{\tilde{}}}
\let\hto\relax\newcommand{\hto}[1][]{\overset{#1}{\lhook\joinrel\longrightarrow}}
% Transcription repair: complete the truncated macro; the printed integral diagrams are unchanged.
\let\intdiag\relax\newcommand{\intdiag}[3]{\int\begin{tikzcd}[row sep=small,column sep=small,ampersand replacement=\&,baseline=(\tikzcdmatrixname-2-1.base)] \& #2 \\ #1 \arrow[ur]\arrow[dr] \& \\ \& #3 \end{tikzcd}}
\let\Ob\relax\newcommand{\Ob}{\operatorname{Ob}}
\let\Fibdisc\relax\newcommand{\Fibdisc}{\underline{\mathrm{Fib\ disc}}}
\let\cercle\relax\newcommand{\cercle}[1]{\tikz[baseline=(c.base)]\node[draw,circle,inner sep=1pt](c){#1};}

\dgrclaim{open}{136}{XIIb-b-B009}{Is phi\_Q in W equivalent to tilde(phi)\_Q in W, so that smooth and proper integrals define the same W-homotopy-cocartesian squares?}
\dgrclaim{false}{139}{XIIb-b-B014}{The indexing category Gamma is the ordered square [1]x[1] with an initial object adjoined.}
\dgrclaim{conjectured}{143}{XIIb-b-B025}{The canonical isomorphisms among the four expanded diagrams should form a transitive system.}
\dgrclaim{conjectured}{144}{XIIb-b-B029}{For arbitrary diagram shapes the four expanded comparisons have the same canonical-isomorphism relations as in Proposition 14.}
\dgrclaim{conjectured}{145–146}{XIIb-b-B033}{Can phi\_\{Q\_0\}:int\_lis F->P fail to belong to W while there exists a strict cocone P->Y with composite phi\_Q in W?}
\dgrclaim{false}{155}{XIIbb-C009}{The arrows x\_0\ensuremath{\to}\ensuremath{\alpha}\_i(x\_0) displayed in the source are the cartesian arrows in the smooth integral.}
\dgrclaim{false}{156}{XIIbb-C013}{The displayed diamond A\_0(u)\ensuremath{\to}A\_1(u),A\_2(u)\ensuremath{\to}(E/x\_0)\_u is a commutative W-cocartesian square to which W(7 bis) can be applied directly.}
\dgrclaim{false}{157}{XIIbb-C021}{For x in any of the three strata, local W-smoothness of \ensuremath{\varphi}\_Q is equivalent to local W-smoothness of f\_i without any additional hypotheses on the other strata.}
\par\endgroup

\section*{XIIc-a \textperiodcentered\ Chapter XII, section 6, part 1 of 2, manuscript pages 171--234}
\pdfbookmark[1]{XIIc-a}{dgrclaims.XIIc-a}
\begingroup\raggedright
% DER_XIIc-a: the definitions of the chunk's preamble, applied inside its group
\let\printedpage\relax\newcommand{\printedpage}[1]{{\reversemarginpar\marginnote{\raggedleft\scriptsize ed.~#1}}}
\let\msp\relax\newcommand{\msp}[1]{\marginpar{\raggedright\scriptsize ms.~#1}}
\let\Ch\relax\newcommand{\Ch}{\underline{\mathrm{Ch}}}
\let\Hotp\relax\newcommand{\Hotp}[1]{\mathrm{Hot}_{\bullet\,#1}}
\let\hotp\relax\newcommand{\hotp}[1]{\mathrm{hot}_{\bullet\,#1}}
\let\Ob\relax\newcommand{\Ob}{\mathrm{Ob}\,}
\let\HomEns\relax\newcommand{\HomEns}{\mathrm{Hom}_{\mathrm{Ens}}}
\let\utpi\relax\newcommand{\utpi}{\underset{\sim}{\pi}}
\let\Fib\relax\newcommand{\Fib}{\underline{\mathrm{Fib}}}
\let\Fl\relax\newcommand{\Fl}{\underline{\mathrm{Fl}}}
\let\Hom\relax\newcommand{\Hom}{\underline{\mathrm{Hom}}}
\let\Cat\relax\newcommand{\Cat}{\mathrm{Cat}}
\let\HOT\relax\newcommand{\HOT}{\mathrm{HOT}}
\let\Hot\relax\newcommand{\Hot}{\mathrm{Hot}}
\let\hot\relax\newcommand{\hot}{\mathrm{hot}}
\let\id\relax\newcommand{\id}{\mathrm{id}}
\let\cM\relax\newcommand{\cM}{\mathcal{M}}
\let\cF\relax\newcommand{\cF}{\mathcal{F}}
\let\Ho\relax\newcommand{\Ho}{\mathrm{Ho}}
\let\ho\relax\newcommand{\ho}{\mathrm{ho}}
\let\Wu\relax\newcommand{\Wu}{W^{\mathrm{u}}}
\let\lc\relax\newcommand{\lc}{\mathrm{lc}}
\let\circled\relax\newcommand{\circled}[1]{\tikz[baseline=(c.base)]\node[draw,circle,inner sep=1pt](c){#1};}
\let\deff\relax\newcommand{\deff}{\overset{\text{déf}}{=}}

\dgrclaim{open}{172}{A011}{Compare the homotopy-group theories defined using $\Ch$ and $\Ch_\infty$.}
\dgrclaim{false}{172}{A012}{Theorem1: isomorphisms of homotopy groups through n, either at all basepoints or at one per component, are equivalent to isomorphisms of local-coefficient cohomology through n, with nonabelian H1.}
\dgrclaim{false}{172}{A013}{Footnote112 suggests that the preceding converse already fails for n=2 unless source and target are simply connected.}
\dgrclaim{open}{174}{A019}{Does a W-fibration E->X with a section canonically determine a functor X->Hot\_\{bullet,W\} sending all arrows to isomorphisms?}
\dgrclaim{false}{176-177}{A022}{Under the lower-degree homotopy and cohomology comparison hypotheses, pi\_n(f) is an isomorphism iff H\textasciicircum{}n(f;F) is an isomorphism for all abelian local systems F.}
\dgrclaim{open}{177-178}{A024}{For a homomorphism of the categorical loop monoids, a homotopy-group comparison at the units extends to every component.}
\dgrclaim{false}{179}{A029}{Each arrow in the corrected implication chain is conjectured strict.}
\dgrclaim{open}{183}{A042}{Does the comparison of the universal-cover E2 rows for q<=n suffice for the desired total-cohomology comparison through n?}
\dgrclaim{false}{187}{A050}{The good fundamental localizer W\_n for n-types is defined by a\_n.}
\dgrclaim{open}{179}{A062}{Footnote117 asks whether c\_n implies b\_n even for simply connected X,Y and conjectures counterexamples because only the Hurewicz composite from pi\_nZ is seen.}
\dgrclaim{open}{194}{B006}{In what sense are fixed-base or varying-base functorial factorizations unique?}
\dgrclaim{open}{196-197}{B011}{Is there a construction-independent canonical comparison \ensuremath{\varphi}\_C\ensuremath{\cong}\ensuremath{\varphi}\_C\ensuremath{'}, independent of using CC\ensuremath{'} or C\ensuremath{'}C?}
\dgrclaim{false}{201}{B022}{The existence of these common targets, or of the natural binary comparison maps, proves C\_S is filtered.}
\dgrclaim{conjectured}{202}{B024}{The three binary-to-ternary and three unary-to-ternary comparisons obtained by inserting coaugmentations agree.}
\dgrclaim{conjectured}{202-203}{B025}{For any ordered partial product of a finite list of replacements there is a unique comparison to the full ordered product obtained by the permitted unit insertions and associativity maps.}
\dgrclaim{open}{203}{B026}{Do these coherences imply that C\_S is 1-connected?}
\dgrclaim{false}{204}{B028}{The one-object category B\_G of an arbitrary group G is a strict monoidal category of the indicated kind.}
\dgrclaim{false}{205}{B032}{The functor labelled (*) is an equivalence.}
\dgrclaim{false}{208}{B039}{Under only the assumptions of footnote 119, Fib\_W is stable under retracts in the arrow category.}
\dgrclaim{open}{209}{B042}{Construct inverse images on Ho\_W\textasciicircum{}!(S) for arbitrary base maps, coherently in S.}
\dgrclaim{conjectured}{216}{B056}{Axioms a and b will usually be satisfied and proved together.}
\dgrclaim{conjectured}{216}{B058}{For Cat, W(4a) supplies the fibre-localization comparison; footnote 122 suggests W(4b) should suffice.}
\dgrclaim{open}{217-218}{B062}{Do axioms a and b alone imply the intermediate-full-subcategory assertion c?}
\dgrclaim{conjectured}{218}{B065}{In practical cases verification of a should also establish c.}
\dgrclaim{conjectured}{219}{C004}{For the main models of ordinary homotopy types one expects a natural notion of local isomorphism.}
\dgrclaim{conjectured}{221}{C012}{A satisfactory theory for varying S should have $f^*: \mathrm{Ho}^{\mathcal L}_W(S)\to\mathrm{Ho}^{\mathcal L}_W(S^{\prime})$ for every $f:S^{\prime}\to S$.}
\dgrclaim{conjectured}{221}{C013}{The localization of $\mathrm{Fib}_W(S)$ should identify the locally constant full subcategory of the finer relative theory, and usually not all of it.}
\dgrclaim{open}{222}{C017}{The provisional definition does not visibly imply stability under all pullbacks.}
\dgrclaim{open}{222}{C023}{Is $\widehat{\mathcal L}=\mathrm{Liss}_W$?}
\dgrclaim{open}{224-225}{C030}{Does a smooth/proper complementarity occur in general models of homotopy types?}
\dgrclaim{conjectured}{225}{C032}{The smooth/proper duality may extend to all presheaf models and to topological spaces.}
\dgrclaim{conjectured}{225-226}{C033}{Proposed data: classes Liss and Prop containing isomorphisms and stable under composition and pullback, perhaps with $\mathrm{Liss}\cap\mathrm{Prop}\subseteq\mathrm{Fib}_W$.}
\dgrclaim{conjectured}{226}{C038}{Proposed axiom (c): the two dual adjunction properties from ChapterVI theorem-scholie2 should be added.}
\dgrclaim{open}{226}{C039}{Can the hypotheses for those adjunction formulas be simplified?}
\dgrclaim{open}{226-227}{C040}{Construct $f_!$ and $f_*$ for the locally constant, proper and smooth relative theories; $f_!\dashv f^*$ is expected to be comparatively direct.}
\dgrclaim{open}{227}{C041}{Seek a cofiber theory dual to the homotopy-fibre theory for (Cat,W).}
\dgrclaim{open}{232-233}{C057}{Does every ordinary pushout along a W-cofibration give a W-cocartesian square, in the earlier integration-based definition?}
\dgrclaim{open}{228}{XIIcb-A001}{Every functor $f:X\to Y$ admits a factorization $X\xrightarrow{i_f}\widetilde Y\xrightarrow{p_f}Y$ with $i_f\in\operatorname{Cof}_W$ and $p_f\in W$.}
\dgrclaim{open}{232–233}{XIIcb-A019}{Is every strict pushout along a $W$-cofibration $W$-cocartesian in the categorical-integral sense of XII, definition 3, ms.131?}
\par\endgroup

\section*{XIIc-b \textperiodcentered\ Chapter XII, section 6, part 2 of 2, manuscript pages 235--275}
\pdfbookmark[1]{XIIc-b}{dgrclaims.XIIc-b}
\begingroup\raggedright
% DER_XIIc-b: the definitions of the chunk's preamble, applied inside its group
\tikzcdset{every diagram/.append style={nodes in empty cells}}
\fancyhf{}
\fancyhead[L]{\footnotesize\itshape Les D\'erivateurs -- XIIc-b}
\fancyhead[R]{\footnotesize Manuscript 228--275}
\fancyfoot[C]{\footnotesize\thepage}
\providecommand\headrulewidth{}\renewcommand{\headrulewidth}{0.3pt}
\let\printedpage\relax\newcommand{\printedpage}[1]{\marginpar{\footnotesize p.~#1}}
\let\msp\relax\newcommand{\msp}[1]{\marginpar{\footnotesize [#1]}}
\let\cX\relax\newcommand{\cX}{\mathscr{X}}
\let\Ob\relax\newcommand{\Ob}{\operatorname{Ob}}
\let\cercle\relax\newcommand{\cercle}[1]{\tikz[baseline=(c.base)]\node[draw,circle,inner sep=1pt](c){#1};}
\let\Fibu\relax\newcommand{\Fibu}{\underline{\mathrm{Fib}}}
\let\Homu\relax\newcommand{\Homu}{\underline{\mathrm{Hom}}}
\let\Ho\relax\newcommand{\Ho}{\mathrm{Ho}}
\let\ho\relax\newcommand{\ho}{\mathrm{ho}}
\let\lc\relax\newcommand{\lc}{\mathrm{lc}}
\let\uFib\relax\newcommand{\uFib}{\underline{\mathrm{Fib}}}
\let\Fib\relax\newcommand{\Fib}{\mathrm{Fib}}
\let\Cof\relax\newcommand{\Cof}{\mathrm{Cof}}
\let\Asph\relax\newcommand{\Asph}{\mathrm{Asph}}

\dgrclaim{conjectured}{236 footnote 127}{XIIcb-A029}{The composition-and-cancellation statement should hold in every model setting with sufficient cofibrations.}
\dgrclaim{false}{238}{XIIcb-A038}{For $\{1,2,1^{\prime}\}\hookrightarrow V$, the slices at $0,0^{\prime}$ are annotated as $\Delta^1$ and $\Phi$.}
\dgrclaim{conjectured}{240}{XIIcb-B004}{For the proposed fibre theory, it should suffice that every W-fibration admit arbitrary pullbacks, rather than that every pair of maps admit a pullback.}
\dgrclaim{conjectured}{240}{XIIcb-B006}{The author suspects that the left factorization axiom gives the entire good relative fibre theory only if every X\ensuremath{\to} e is a W-fibration.}
\dgrclaim{false}{269–270}{XIIcb-C020}{The reduction of all three cube implications to the case where Y is the terminal square uses only the direct, pushout-stability half of W(7 bis).}
\dgrclaim{open}{272}{XIIcb-C028}{Does the factorization hypothesis alone suffice to construct the required theory when not every object is fibrant over the terminal object?}
\dgrclaim{false}{273}{XIIcb-C035}{The conservative-pushout axiom is equivalent to the conjunction of Cons1 and literal cartesian Cons2.}
\dgrclaim{false}{274}{XIIcb-C037}{Under only the theorem’s general cofiber/factorization hypotheses, a terminal object and closure of aspherical objects under retracts, f\_0,f\_3\ensuremath{\in} W imply f\_1,f\_2\ensuremath{\in} W for every map of W-cocartesian squares.}
\dgrclaim{false}{274}{XIIcb-C038}{Corollary1: the universal map-of-squares implication f\_0,f\_3\ensuremath{\in} W\ensuremath{\Rightarrow}f\_1,f\_2\ensuremath{\in} W is equivalent to the object condition Z\_0,Z\_3 aspherical\ensuremath{\Rightarrow}Z\_1,Z\_2 aspherical, and to its special case Z\_0=e.}
\par\endgroup

\section*{XIIapp \textperiodcentered\ Chapter XII, appendix, manuscript pages 1--18}
\pdfbookmark[1]{XIIapp}{dgrclaims.XIIapp}
\begingroup\raggedright
% DER_XIIapp: the definitions of the chunk's preamble, applied inside its group
\let\printedpage\relax\newcommand{\printedpage}[1]{\if\relax\detokenize{#1}\relax\else\marginnote{\footnotesize\strut ed.~#1}[-14pt]\fi}
\def\msnumber page #1\relax{#1}
\let\msp\relax\newcommand{\msp}[1]{\marginnote{\footnotesize\strut [ms.~\msnumber#1\relax]}}
\let\Hom\relax\newcommand{\Hom}{\mathrm{Hom}}
\let\Fl\relax\newcommand{\Fl}{\mathrm{Fl}}
\let\Ob\relax\newcommand{\Ob}{\mathrm{Ob}}
\let\Cat\relax\newcommand{\Cat}{\mathrm{Cat}}
\let\Ens\relax\newcommand{\Ens}{\mathrm{Ens}}
\let\pr\relax\newcommand{\pr}{\mathrm{pr}}
\let\Ker\relax\newcommand{\Ker}{\mathrm{Ker}}
\let\Coker\relax\newcommand{\Coker}{\mathrm{Coker}}
\let\mon\relax\newcommand{\mon}{\mathrm{mon}}
\let\deq\relax\newcommand{\deq}{\mathrel{\overset{\text{déf}}{=}}}
\let\id\relax\newcommand{\id}{\mathrm{id}}
\let\Winf\relax\newcommand{\Winf}{\mathbf{W}_{\infty}}
\let\dotto\relax\newcommand{\dotto}{\mathrel{\tikz[baseline=-0.5ex]\draw[dotted,-latex] (0,0)--(1.6em,0);}}

\dgrclaim{false}{1}{XIIapp-O002}{If b and c are weakly terminal and their categorical product b×c exists, BX is weakly contractible.}
\dgrclaim{false}{3}{XIIapp-O015}{H\textasciicircum{}p(T,F)=0 for every abelian presheaf F and every p\ensuremath{\geq}2.}
\dgrclaim{false}{3}{XIIapp-O017}{H\textasciicircum{}2(T,F)=coker(F(i\_1)-F(i\_2)), with H\textasciicircum{}1 evidently intended in the displayed two-term formula.}
\dgrclaim{false}{4}{XIIapp-O021}{H\textasciicircum{}n(X\_G;M)\ensuremath{\simeq}H\textasciicircum{}1(T,R\textasciicircum{}\{n-1\}f\_*M)\ensuremath{\simeq}coker(H\textasciicircum{}\{n-1\}(G;M)\ensuremath{\to}H\textasciicircum{}\{n-1\}(G²;M)).}
\dgrclaim{false}{4}{XIIapp-O025}{H²(BX\_G;M)\ensuremath{\simeq}Hom\_mon(G,M).}
\dgrclaim{false}{4}{XIIapp-O026}{X\_G is 2-connected if and only if the universal abelian group of G is zero.}
\dgrclaim{false}{4}{XIIapp-O027}{If G is a nonzero abelian group, then X\_G is not 2-connected.}
\dgrclaim{open}{4}{XIIapp-O029}{If y\_b\ensuremath{\sqcup}\_\{y\_a\}y\_c=1\_X, must X at least be simply connected?}
\dgrclaim{open}{5}{XIIapp-O030}{Are totally W\ensuremath{\infty}-aspherical categories precisely categories with a terminal object?}
\dgrclaim{false}{6}{XIIapp-O035}{X\_G is simply connected, but not 2-connected whenever its universal abelian group is nontrivial.}
\dgrclaim{false}{5}{XIIapp-R001}{Every totally W\ensuremath{\infty}-aspherical category, in the universal left-Kan-extension sense used here, has a terminal object.}
\dgrclaim{conjectured}{5}{XIIapp-R002}{The suspected terminal-object characterization would mean that the existing notion of open Dwyer immersion cannot be substantially improved.}
\dgrclaim{false}{9}{XIIapp-C007}{There is a retraction r:V\ensuremath{\to} D fixing D and sending b,c to a.}
\dgrclaim{false}{9}{XIIapp-C008}{The alleged retraction gives a natural transformation ir\ensuremath{\Rightarrow} id\_V and a categorical homotopy equivalence D\ensuremath{\to} V.}
\dgrclaim{false}{9}{XIIapp-C010}{BV has the homotopy type of S\textasciicircum{}1 and is not simply connected.}
\dgrclaim{false}{10}{XIIapp-C018}{The singleton-pushout condition can hold while Z is not simply connected.}
\dgrclaim{false}{11}{XIIapp-C029}{The inclusion Z\_0=(0\_0\ensuremath{\rightrightarrows} a\_0)\ensuremath{\to} Z is a homotopy equivalence by an evident retraction.}
\dgrclaim{false}{15}{XIIapp-C062}{pi\_1(V\_0 prime) is isomorphic to Z.}
\dgrclaim{false}{15}{XIIapp-C063}{V\_0 prime is non-W-aspherical whenever W is contained in W\_1, by its alleged nontrivial fundamental group.}
\dgrclaim{open}{17}{XIIapp-G05}{The crossing arrows define a natural transformation $ip\Rightarrow jq$ and hence a homotopy-commutative square with vertices $E,A,B,X$; is it $W$-cocartesian?}
\dgrclaim{open}{17}{XIIapp-G11}{Under $i\in W$ or $j\in W$ and W-asphericity of $X$, is the same square still W-cocartesian, in particular for $W=\mathbf W_\infty$?}
\dgrclaim{false}{17}{XIIapp-G12}{For $W=\mathbf W_\infty$, under $i\in W$ (respectively $j\in W$) and W-asphericity of $X$, W-cocartesianness is equivalent to $q\in W$ (respectively $p\in W$).}
\dgrclaim{open}{17}{XIIapp-G14}{If $A\simeq e$ and $X,B$ are W-aspherical, must the attaching comma category $a\backslash B$ be W-aspherical, equivalently must $j$ be W-coaspherical?}
\dgrclaim{conjectured}{18}{XIIapp-G23}{The covering-descent argument should extend to higher homotopy degrees through n-gerbes.}
\dgrclaim{conjectured}{18}{XIIapp-G26}{For every open--closed decomposition, $\varphi\in\mathbf W_\infty$.}
\dgrclaim{conjectured}{18}{XIIapp-G27}{If $W\supseteq\mathbf W_\infty$ satisfies Loc $(6\mathrm a')$, then the open--closed square is W-cocartesian.}
\par\endgroup

\section*{XIII-a \textperiodcentered\ Chapter XIII, part 1 of 2, manuscript pages 1--60}
\pdfbookmark[1]{XIII-a}{dgrclaims.XIII-a}
\begingroup\raggedright
% DER_XIII-a: the definitions of the chunk's preamble, applied inside its group
\let\printedpage\relax\newcommand{\printedpage}[1]{\marginpar{\footnotesize p.~#1}}
\let\msp\relax\newcommand{\msp}[1]{\marginpar{\footnotesize [#1]}}
\let\Fl\relax\newcommand{\Fl}{\mathrm{Fl}}
\let\M\relax\newcommand{\M}{\mathcal{M}}
\let\cM\relax\newcommand{\cM}{\mathcal{M}}
\let\Cof\relax\newcommand{\Cof}{\mathrm{Cof}}
\let\Fib\relax\newcommand{\Fib}{\mathrm{Fib}}
\let\Ob\relax\newcommand{\Ob}{\mathrm{Ob}}
\let\circled\relax\newcommand{\circled}[1]{\tikz[baseline=(c.base)]\node[draw,circle,inner sep=1pt](c){#1};}
\let\univ\relax\newcommand{\univ}{\mathrm{univ}}
\let\cercle\relax\newcommand{\cercle}[1]{\mbox{\tikz[baseline=(c.base)]\node[draw,circle,inner sep=1pt](c){#1};}}
\let\clarrow\relax\newcommand{\clarrow}{\mathrel{\leftarrow\mkern-4mu{+}\mkern-4mu\rightarrow}}
\let\limind\relax\newcommand{\limind}{\mathop{\underrightarrow{\mathrm{lim}}}}
\let\limproj\relax\newcommand{\limproj}{\mathop{\underleftarrow{\mathrm{lim}}}}
\let\HO\relax\newcommand{\HO}{\mathrm{HO}}

\dgrclaim{conjectured}{6}{XIIIa-A025}{The small object construction should require only pushouts and filtered, or ordinal, colimits instead of full cocompleteness.}
\dgrclaim{conjectured}{6–7}{XIIIa-A026}{There should be a notion of accessible arrow class \ensuremath{\Phi} implying existence of a set I\ensuremath{\subseteq}\ensuremath{\Phi} with \ensuremath{\Phi}\ensuremath{\subseteq}\{\}\textasciicircum{}\ensuremath{\square}(I\textasciicircum{}\ensuremath{\square}).}
\dgrclaim{conjectured}{7}{XIIIa-A027}{The proposed accessible classes should be preserved by \ensuremath{\Phi}\ensuremath{\mapsto}\ensuremath{\Phi}\_* and \ensuremath{\Psi}\ensuremath{\mapsto}\ensuremath{\Psi}\textasciicircum{}*.}
\dgrclaim{conjectured}{7}{XIIIa-A028}{If a left/right lifting closed pair is accessible in a sense supplying a set of left generators, it has Quillen factorization.}
\dgrclaim{conjectured}{12}{XIIIa-A052}{The weak axioms M0–M4,M\ensuremath{\prime},M\ensuremath{\prime}\ensuremath{\prime} may be insufficient to obtain exact suspension and cosuspension sequences.}
\dgrclaim{conjectured}{13}{XIIIa-A058}{The equality TC=C\ensuremath{\cap} W should imply the required exactness of W-cartesian squares and cosuspension sequences.}
\dgrclaim{conjectured}{13}{XIIIa-A059}{The equality TF=F\ensuremath{\cap} W should imply the required exactness of W-cocartesian squares and suspension sequences.}
\dgrclaim{conjectured}{14}{XIIIa-A061}{A relative category may admit distinct left and right pre-Q structures even when no full Quillen structure with those weak equivalences is available.}
\dgrclaim{conjectured}{15}{XIIIa-A069}{Under accessibility of M, conditions (a),(a\ensuremath{\prime}),(a\ensuremath{\prime}\ensuremath{\prime}),(b),(b\ensuremath{\prime}) are proposed to be equivalent.}
\dgrclaim{conjectured}{15}{XIIIa-A071}{Theorem-scholie(3): a closed right pre-Q structure exists iff (b\ensuremath{\prime}) holds and A\_W=(W\textasciicircum{}univ)\textasciicircum{}\ensuremath{\square}\ensuremath{\cap} W\_univ is right lifting closed.}
\dgrclaim{conjectured}{16}{XIIIa-A072}{The preliminary hypotheses, accessibility of M, and the proposed equivalent recognition conditions are stable under (M,W)\ensuremath{\mapsto}(M\textasciicircum{}I,W\textasciicircum{}I) for every small diagram category I.}
\dgrclaim{open}{16}{XIIIa-A073}{Is the extra right-closedness condition in item 3 stable under passage to diagram categories?}
\dgrclaim{open}{16}{XIIIa-A074}{Do the conditions in item 2 imply the extra condition in item 3, making the two criteria equivalent?}
\dgrclaim{conjectured}{16}{XIIIa-A078}{(a\ensuremath{\prime}) should imply (a), after requiring lifting closedness of the quadruple.}
\dgrclaim{false}{17}{XIIIa-A080}{The closure properties proved in item 1 imply that Cof\_W and W\textasciicircum{}univ are left lifting closed.}
\dgrclaim{conjectured}{19}{XIIIa-A085}{Filtered-colimit stability of W should make W accessible in a suitable sense, and then imply small generation of C\ensuremath{\cap} W when C is small-generated.}
\dgrclaim{conjectured}{20}{XIIIa-A089}{Conjectural lemma(1), fibration version: if \ensuremath{\Phi} is left lifting closed and \ensuremath{\Phi}\textasciicircum{}\ensuremath{\square}\ensuremath{\subseteq} Fib\_W, there exists a set \ensuremath{\Phi}0\ensuremath{\subseteq}\ensuremath{\Phi} such that \ensuremath{\Phi}0\textasciicircum{}\ensuremath{\square}\ensuremath{\subseteq} Fib\_W.}
\dgrclaim{conjectured}{20}{XIIIa-A090}{Conjectural lemma(1), weak-equivalence version intended: if \ensuremath{\Phi} is left lifting closed and \ensuremath{\Phi}\textasciicircum{}\ensuremath{\square}\ensuremath{\subseteq} W, there exists a set \ensuremath{\Phi}0\ensuremath{\subseteq}\ensuremath{\Phi} such that \ensuremath{\Phi}0\textasciicircum{}\ensuremath{\square}\ensuremath{\subseteq} W.}
\dgrclaim{conjectured}{20}{XIIIa-A091}{Conjectural lemma(2): the intersection of two set-generated left lifting closed classes in a locally presentable category is again set-generated.}
\dgrclaim{conjectured}{20}{XIIIa-A092}{An intrinsic class such as Fib\_W, defined from an accessible W, should itself be accessible and yield the compactness lemma.}
\dgrclaim{false}{22}{XIIIa-B003}{Part 1 of the conjectural theorem implies that \ensuremath{\Cof}\_W and W\textasciicircum{}{\ensuremath{\univ}} are left lifting classes.}
\dgrclaim{open}{25}{XIIIa-B016}{Does existence of some closed Quillen triple imply either the stronger closure conditions or the maximal compatibility A\_W=(\ensuremath{\Cof}\_W)\_*?}
\dgrclaim{conjectured}{26}{XIIIa-B019}{Group I identifies weak Quillen existence, left exact weak Quillen existence, and the two orthogonality inclusions (*).}
\dgrclaim{conjectured}{26}{XIIIa-B020}{Group II identifies right exact weak Quillen existence, condition (*) plus right closure of A\_W, closed Quillen existence, and the displayed maximal-TC construction.}
\dgrclaim{false}{26}{XIIIa-B021}{With TC=W\textasciicircum{}{\ensuremath{\univ}}, F=TC\_*, TF=F\ensuremath{\cap} W\_{\ensuremath{\univ}}, define C=(TF)\_*.}
\dgrclaim{conjectured}{26}{XIIIa-B022}{Group III identifies the maximal closed model candidate with (*) plus A\_W=(\ensuremath{\Cof}\_W)\_*.}
\dgrclaim{open}{26}{XIIIa-B027}{Could the identity \ensuremath{\overline{W_{\univ}}}\ensuremath{\cap}\ensuremath{\Fib}\_W=W\_{\ensuremath{\univ}} imply III\ensuremath{\Leftrightarrow}IV?}
\dgrclaim{false}{26}{XIIIa-B029}{The displayed intersection (W\textasciicircum{}{\ensuremath{\univ}})\_*\ensuremath{\cap}\ensuremath{\overline{W_{\univ}}} contains D\ensuremath{\to} e for discrete D without qualification.}
\dgrclaim{false}{27}{XIIIa-B041}{If C,F,W are retract closed, left exactness is equivalent to right exactness for the original designated classes and forces all four classes closed.}
\dgrclaim{false}{28}{XIIIa-B044}{The preceding saturation produces a structure de Quillen in the biexact sense introduced immediately afterwards.}
\dgrclaim{conjectured}{30}{XIIIa-B060}{For arbitrary weak structures, right Q-homotopy is an equivalence relation whenever the target is fibrant, without exactness.}
\dgrclaim{conjectured}{30}{XIIIa-B061}{On fibrant-cofibrant objects, the right and left homotopies coincide and are independent of the chosen cylinder, as suggested by Baues.}
\dgrclaim{conjectured}{30–31}{XIIIa-B062}{Under left exactness, the fibrant-cofibrant quotient equivalent to Ho\_W embeds as a fully faithful right adjoint to localization of the right-homotopy quotient of cofibrant objects.}
\dgrclaim{open}{32}{XIIIa-B067}{Does left exactness also suffice for exactness of W-homotopy-cocartesian squares?}
\dgrclaim{false}{32}{XIIIa-B068}{Exactness alone may fail to imply either cartesian or cocartesian homotopy exactness.}
\dgrclaim{conjectured}{33–34}{XIIIa-B069}{If exact weak structures exist on (M,W) and on (M(\ensuremath{\Psi}),W(\ensuremath{\Psi})) for the cospan indexing category, W-homotopy-cartesian squares are W-exact and give loop exact sequences.}
\dgrclaim{conjectured}{33}{XIIIa-B071}{The dual exact-structure hypothesis for M(\ensuremath{\Phi}) gives cocartesian exactness and pointed suspension exact sequences.}
\dgrclaim{conjectured}{34}{XIIIa-B072}{For simultaneous homotopical algebra in all M(I), it suffices to have an exact weak Quillen structure on each relative diagram category.}
\dgrclaim{open}{37}{XIIIa-B086}{For any arbitrary arrow class \ensuremath{\Phi}, evaluation might send \ensuremath{\Phi}(I)\_* into \ensuremath{\Phi}\_*.}
\dgrclaim{conjectured}{40}{XIIIa-B097}{The author expects filtered-colimit stability of W to be practically necessary for this strategy.}
\dgrclaim{conjectured}{41}{XIIIa-C003}{Closure of W under all filtered colimits is proposed as the practical way to establish the required diagram factorizations.}
\dgrclaim{conjectured}{42–43}{XIIIa-C011}{Condition a implies existence, for every small I, of a closed left-exact weak Quillen structure on Fun(I\textasciicircum{}op,M) with objectwise cofibrations.}
\dgrclaim{conjectured}{42–43}{XIIIa-C012}{Condition a-prime implies existence of the prescribed injective structures for all small I.}
\dgrclaim{conjectured}{42–43}{XIIIa-C013}{Condition a-double-prime implies existence of the prescribed closed left-exact weak structures for all small I.}
\dgrclaim{conjectured}{42–43}{XIIIa-C014}{Condition b implies existence of the prescribed closed left-exact weak structures for all small I.}
\dgrclaim{conjectured}{44}{XIIIa-C020}{The author expects a-double-prime to imply all other conditions of Theorem 2.}
\dgrclaim{conjectured}{44}{XIIIa-C021}{The author expects the small-set condition a-double-prime to be the practical verification method, notably for M=Cat.}
\dgrclaim{conjectured}{48}{XIIIa-C041}{The converse componentwise condition for injective fibrations need not hold.}
\dgrclaim{conjectured}{48–49}{XIIIa-C044}{Restriction also preserves homotopy cartesian squares and cosuspension exact sequences.}
\dgrclaim{false}{49}{XIIIa-C045}{The ordinary functor u\_* preserves arbitrary homotopy cartesian squares merely because it preserves F and TF.}
\dgrclaim{conjectured}{49}{XIIIa-C046}{A right homotopy Kan extension generally does not preserve homotopy cocartesian squares.}
\dgrclaim{false}{49}{XIIIa-C047}{The ordinary u\_! should preserve arbitrary homotopy cocartesian squares.}
\dgrclaim{open}{50}{XIIIa-C049}{Do the preceding assumptions ensure the existence of homotopy adjoints u\_!\textasciicircum{}W and u\_*\textasciicircum{}W, even for finite posets?}
\dgrclaim{open}{52}{XIIIa-C068}{What quillenization or derived-functor theory is available for unbounded complexes in a general abelian category?}
\dgrclaim{false}{53}{XIIIa-C071}{The same filtered-colimit closure holds for Ch\textasciicircum{}+(A), Ch\textasciicircum{}-(A), and Ch\textasciicircum{}b(A).}
\dgrclaim{open}{53}{XIIIa-C075}{Does the right derived Kan extension exist on the whole unbounded derived category?}
\dgrclaim{conjectured}{53}{XIIIa-C076}{A left homotopy Kan extension need not preserve D\textasciicircum{}+ without a homological-dimension restriction.}
\dgrclaim{conjectured}{53–54}{XIIIa-C077}{For Grothendieck A, a left derived Kan extension is expected on D\textasciicircum{}-(A(I-prime)) with values in D\textasciicircum{}-(A(I)).}
\dgrclaim{open}{54}{XIIIa-C078}{Does left derived Kan extension exist on all of D(A(I-prime)) without finite cohomological dimension?}
\dgrclaim{open}{55}{XIIIa-C083}{Which derivator axioms follow from quillenizability alone?}
\dgrclaim{conjectured}{55–56}{XIIIa-C087}{The strong local-system axiom asks that u in the chosen fundamental localizer induce an equivalence on homotopy categories of locally constant diagrams.}
\dgrclaim{conjectured}{56}{XIIIa-C092}{Quillenizability alone seems too weak to ensure interesting homotopy Kan extensions.}
\dgrclaim{false}{59}{XIIIa-C101}{The parenthesis invokes finite colimits to justify the finite limits required by finite-index u\_*.}
\dgrclaim{conjectured}{60}{XIIIa-C109}{A poset whose associated opposite Alexandrov space is noetherian and sober is admissible.}
\dgrclaim{open}{52}{XIIIb-Q013}{What hypotheses suffice for quillenisability of unbounded complexes, and what can be done without enough injectives or projectives?}
\dgrclaim{false}{53}{XIIIb-Q015}{If A is accessible and has filtered colimits, all the listed complex categories have the same closure properties.}
\dgrclaim{open}{53–54}{XIIIb-Q019}{Does u\_* extend to a right adjoint on unbounded derived categories of diagrams?}
\dgrclaim{conjectured}{53–54}{XIIIb-Q020}{For globally bounded-above diagrams of Grothendieck coefficients, one expects a left adjoint Lu\_! to restriction preserving upper bounds.}
\dgrclaim{open}{54}{XIIIb-Q021}{Does a left adjoint to restriction exist on the entire unbounded derived diagram category without a finite cohomological-dimension hypothesis?}
\dgrclaim{open}{54 footnote32}{XIIIb-Q024}{Can the square-based structure of an unpointed homotopy theory be axiomatized?}
\dgrclaim{open}{55}{XIIIb-Q026}{Which derivator axioms follow from quillenisability alone?}
\dgrclaim{conjectured}{55–56}{XIIIb-Q030}{A strongest proposed local-constancy axiom says that u in a chosen fundamental localizer induces HO\_W\textasciicircum{}lc(I)\ensuremath{\simeq}HO\_W\textasciicircum{}lc(I\ensuremath{'}).}
\dgrclaim{open}{56}{XIIIb-Q032}{How can the local-constancy equivalence be verified directly from (M,W), before a theory of u\_! or u\_* has been built?}
\par\endgroup

\section*{XIII-b \textperiodcentered\ Chapter XIII, part 2 of 2, manuscript pages 61--98}
\pdfbookmark[1]{XIII-b}{dgrclaims.XIII-b}
\begingroup\raggedright
% DER_XIII-b: the definitions of the chunk's preamble, applied inside its group
\let\printedpage\relax\newcommand{\printedpage}[1]{\marginpar{\footnotesize p.~#1}}
\let\msp\relax\newcommand{\msp}[1]{\marginpar{\footnotesize [#1]}}
\let\circled\relax\newcommand{\circled}[1]{\tikz[baseline=(c.base)]\node[draw,circle,inner sep=1pt](c){#1};}
\let\cM\relax\newcommand{\cM}{\mathcal{M}}
\let\HO\relax\newcommand{\HO}{\mathrm{HO}}
\let\Hom\relax\newcommand{\Hom}{\mathrm{Hom}}
\let\Ob\relax\newcommand{\Ob}{\mathrm{Ob}}
\let\rf\relax\newcommand{\rf}{\mathrm{f}}
\let\rc\relax\newcommand{\rc}{\mathrm{c}}
\let\M\relax\newcommand{\M}{\mathcal{M}}
\let\Fl\relax\newcommand{\Fl}{\mathrm{Fl}\,}
\let\Filt\relax\newcommand{\Filt}{\mathrm{Filt}}
\let\card\relax\newcommand{\card}{\mathrm{card}\,}
\let\barlongleftrightarrow\relax\newcommand{\barlongleftrightarrow}{\mathrel{\text{\ooalign{$\longleftrightarrow$\cr\hidewidth$\mid$\hidewidth\cr}}}}
\let\idx\relax\newcommand{\idx}[2]{\noindent #1\dotfill #2\par}
\let\entry\relax\newcommand{\entry}[2]{\noindent #1\dotfill\ #2\par}
\let\Mo\relax\newcommand{\Mo}{\mathcal{M}}
\let\Mzg\relax\newcommand{\Mzg}{$\mathbf{M}'_{\mathbf{0}}\mathbf{(g)}$}
\let\Mzd\relax\newcommand{\Mzd}{$\mathbf{M}'_{\mathbf{0}}\mathbf{(d)}$}
\providecommand\thefootnote{}\renewcommand{\thefootnote}{\arabic{footnote}}

\dgrclaim{open}{69 footnote37}{XIIIb-Q066}{Can the representing object of the preceding pointwise quotient be chosen functorially in localized G\ensuremath{'} without admissibility?}
\dgrclaim{false}{71}{XIIIb-Q070}{Representability of F\ensuremath{\mapsto}Hom(HO(I\ensuremath{'}))(u*F,G\ensuremath{'}) gives a left adjoint of u*.}
\dgrclaim{false}{71}{XIIIb-Q071}{The left adjoint just obtained is computed by u\_* on fibrant G\ensuremath{'}.}
\dgrclaim{open}{72}{XIIIb-Q073}{Do all functors u admit homotopy right images under the earlier large-category hypotheses; can a counterexample be found?}
\dgrclaim{open}{72–73}{XIIIb-Q074}{Is a small product of W-maps between fibrant objects again a W-map?}
\dgrclaim{open}{73}{XIIIb-Q076}{Does inverse limit preserve levelwise W-maps between such fibrant towers?}
\dgrclaim{open}{73–74}{XIIIb-Q078}{Can the product assertion fail even for a closed model category with the additional accessibility and filtered-colimit exactness hypotheses listed?}
\dgrclaim{open}{74}{XIIIb-Q079}{Is the one-object category BG admissible for a group G?}
\dgrclaim{open}{74}{XIIIb-Q080}{Is the full subcategory of \ensuremath{\Delta} on [0] and [1] admissible?}
\dgrclaim{false}{74–75}{XIIIb-Q081}{For the dual construction assume a closed right-exact structure with TF=F\ensuremath{\cap}C.}
\dgrclaim{conjectured}{75}{XIIIb-Q084}{These classes should define a uniform right quillenisability on all small categories.}
\dgrclaim{conjectured}{76}{XIIIb-Q087}{If the projective construction succeeds, left homotopy Kan extensions exist for maps of finite posets and the indicated noetherian sober extensions.}
\dgrclaim{conjectured}{76}{XIIIb-Q088}{Closure of W under filtered colimits should make discrete and left-filtered categories such as N\textasciicircum{}op admissible for left images.}
\dgrclaim{conjectured}{76}{XIIIb-Q089}{A poset which is a filtered union of closed noetherian sober subposets should be right-admissible under the projected dual hypotheses.}
\dgrclaim{open}{76}{XIIIb-Q090}{What happens to left admissibility for BG?}
\dgrclaim{false}{77}{XIIIb-C07-007}{Under GCH, (pi')\textasciicircum{}pi=pi' for every pair of infinite cardinals pi'>pi.}
\dgrclaim{false}{85}{XIIIb-C07-027}{The printed formula asserts Z\_sup A \ensuremath{\cong} lim\_A Z\_T for the covariant partial-colimit diagram.}
\dgrclaim{false}{89}{XIIIb-C07-037}{The three conditions a\_pi, b\_pi, c\_pi are equivalent.}
\dgrclaim{false}{90}{XIIIb-C07-040}{To deduce membership of an arbitrary pi\textasciicircum{}+-directed colimit of M\textasciicircum{}pi-objects, it is enough to apply b\_pi directly because I\_M=I.}
\dgrclaim{false}{94}{XIIIb-C07-054}{For every C\_i' admitted by the printed image-size definition, Z\_\{C\_i'\}=colim\_\{C\_i'\} Z belongs to Filt\textasciicircum{}pi(E).}
\dgrclaim{false}{94}{XIIIb-C07-055}{With the printed tilde C\_i, the diagram (Z\_\{C\_i'\}) is perfectly adapted at pi.}
\dgrclaim{open}{97}{XIIIb-C07-068}{If X has both a perfectly adapted Filt\textasciicircum{}pi presentation and a pi\textasciicircum{}+-directed presentation by M\textasciicircum{}pi-objects, can one show its adapted M-stages cofinal?}
\dgrclaim{conjectured}{97}{XIIIb-C07-069}{After replacing pi by a larger cardinal pi' (suggested pi'=pi\textasciicircum{}+), a\_pi should imply b\_\{pi'\}.}
\dgrclaim{false}{86 (glossary discrepancy)}{XIIIb-G036}{The glossary states that a strict cardinal filtration has every stage \ensuremath{\mathrm{Filt}}\textasciicircum{}{\ensuremath{\pi}}(E) closed under all filtered colimits.}
\dgrclaim{false}{86 (glossary discrepancy)}{XIIIb-G037}{The final glossary paragraph calls a perfectly adapted representation the colimit of a 'well adapted' system.}
\par\endgroup

\section*{XIV \textperiodcentered\ Chapter XIV, manuscript pages 1--8}
\pdfbookmark[1]{XIV}{dgrclaims.XIV}
\begingroup\raggedright
% DER_XIV: the definitions of the chunk's preamble, applied inside its group
\let\printedpage\relax\newcommand{\printedpage}[1]{\if\relax\detokenize{#1}\relax\else\begingroup\reversemarginpar\marginnote{\footnotesize p.~#1}\endgroup\fi}
\let\msp\relax\newcommand{\msp}[1]{\marginnote{\footnotesize [#1]}}

% footnote marks as printed: "(1)" in the text, "1" before the footnote text
\makeatletter
\providecommand\@makefnmark{}\renewcommand\@makefnmark{\hbox{\textsuperscript{\normalfont(\@thefnmark)}}}
\providecommand\@makefntext{}\renewcommand\@makefntext[1]{\parindent 1em\noindent\hb@xt@1.8em{\hss\textsuperscript{\normalfont\@thefnmark}}#1}
\makeatother
\let\uHom\relax\newcommand{\uHom}{\underline{\mathrm{Hom}}}
\let\bDelta\relax\newcommand{\bDelta}{\mathbf{\Delta}}
\let\ho\relax\newcommand{\ho}{\mathrm{ho}}
\let\cM\relax\newcommand{\cM}{\mathcal{M}}

\dgrclaim{conjectured}{4}{XIV.016}{The preceding stability axioms appear to exhaust the useful stability properties, so further such properties should be consequences.}
\dgrclaim{false}{7}{XIV.030}{If the underlying category is additive, the h-cartesian and h-cocartesian classes should coincide.}
\dgrclaim{conjectured}{8}{XIV.035}{In some sense the general case ought to reduce to the elementary ordinary-pullback case.}
\dgrclaim{false}{8}{XIV.036}{The proposed second family of left quadrangulated examples is the localized category with the homotopy-image class of squares, for localizers admitting homotopy fibres.}
\dgrclaim{false}{8}{XIV.037}{Proposition 1 is expected to hold under further hypotheses to be identified, including every presentation quillenizable on finite poset diagram shapes.}
\dgrclaim{false}{8}{XIV.039}{Proposition 1, necessary direction: if $\mathrm{ho}_W(X)$ belongs to the designated h-cartesian class, the chosen strict square $X$ is $W$-cartesian.}
\par\endgroup

\section*{XV-a \textperiodcentered\ Chapter XV, part 1 of 2, manuscript pages 1--66}
\pdfbookmark[1]{XV-a}{dgrclaims.XV-a}
\begingroup\raggedright
% DER_XV-a: the definitions of the chunk's preamble, applied inside its group
\let\printedpage\relax\newcommand{\printedpage}[1]{\marginpar{\footnotesize p.~#1}}
\let\msp\relax\newcommand{\msp}[1]{\marginpar{\footnotesize [#1]}}
\let\M\relax\newcommand{\M}{\mathcal{M}}
\let\Fl\relax\newcommand{\Fl}{\mathrm{Fl}}
\let\uFl\relax\newcommand{\uFl}{\underline{\mathrm{Fl}}}
\let\uNerf\relax\newcommand{\uNerf}{\underline{\mathrm{Nerf}}}
\let\cM\relax\newcommand{\cM}{\mathcal{M}}
\let\Flu\relax\newcommand{\Flu}{\underline{\mathrm{Fl}}}
\let\fU\relax\newcommand{\fU}{\mathfrak{U}}
\let\fV\relax\newcommand{\fV}{\mathfrak{V}}
\let\Filt\relax\newcommand{\Filt}{\mathrm{Filt}}
\let\card\relax\newcommand{\card}{\operatorname{card}}
\let\circled\relax\newcommand{\circled}[1]{\tikz[baseline=(c.base)]\node[draw,circle,inner sep=0.5pt](c){#1};}
\let\pfn\relax\newcommand{\pfn}[2]{\textsuperscript{(}\footnote[#1]{#2}\textsuperscript{)}}
\let\uFilt\relax\newcommand{\uFilt}{\underline{\mathrm{Filt}}}
\let\Hom\relax\newcommand{\Hom}{\operatorname{Hom}}
\let\Ob\relax\newcommand{\Ob}{\operatorname{Ob}}
\let\wavy\relax\newcommand{\wavy}[1]{\widetilde{#1}}
\let\id\relax\newcommand{\id}{\mathrm{id}}
\let\qpre\relax\newcommand{\qpre}{\mathrel{\leftarrow\!\!\!\!\mapsto}}
\let\Cof\relax\newcommand{\Cof}{\mathrm{Cof}}
\let\deq\relax\newcommand{\deq}{\overset{\text{déf}}{=}}
\let\univ\relax\newcommand{\univ}{\mathrm{univ}}
\let\cC\relax\newcommand{\cC}{\mathcal{C}}
\let\Ss\relax\newcommand{\Ss}{\mathrm{Ss}}
\let\Img\relax\newcommand{\Img}{\mathrm{Im}\,}
\let\isoto\relax\newcommand{\isoto}{\xrightarrow{\ \sim\ }}
\let\hookto\relax\newcommand{\hookto}{\lhook\joinrel\longrightarrow}
\let\Ci\relax\newcommand{\Ci}{\mathcal{C}}

\dgrclaim{false}{7–8}{XV-A-012}{If every domain of a generator is $\pi$-accessible and $I_\alpha=\{\beta<\alpha\}$ is grand devant $\pi$, the source asserts $p_\alpha(f)\in I^\perp$.}
\dgrclaim{false}{9–10}{XV-A-024}{The proposed first stage should satisfy $i_1(f)\in C_0^{\prime\prime}$ for every arrow $f$, and the construction is asserted to do so.}
\dgrclaim{false}{11}{XV-A-033}{If the domains of $I$ are accessible, choose a common $\pi$; the source claims $p_\alpha(f)\in I^\perp$ for every ordinal grand devant $\pi$.}
\dgrclaim{false}{17}{XV3-006}{GCH implies pi\textasciicircum{}pi\_0=pi for every pi>pi\_0.}
\dgrclaim{conjectured}{17}{XV3-008}{The endpoint filtration Fl\textasciicircum{}pi(M) should itself be a cardinal filtration of the arrow category.}
\dgrclaim{false}{19}{XV3-015}{Proposition 3.1(2), (3.8): for r,s>=max(pi\_1,C\_1), A in Filt\textasciicircum{}r and X in Filt\textasciicircum{}s imply |Hom(A,X)|<=s\textasciicircum{}(r\textasciicircum{}pi\_0).}
\dgrclaim{false}{19}{XV3-018}{GCH implies pi\textasciicircum{}pi\_prime=pi for infinite pi>pi\_prime>=1 and hence gives (3.7 bis)-(3.9 bis) for every r>pi\_0.}
\dgrclaim{false}{20}{XV3-020}{The category of pairs (A,B subset A) gives equality in the arrow-isomorphism bound (3.7 bis) for the indicated filtration.}
\dgrclaim{false}{20}{XV3-022}{|Hom(A\_i,X\_j)|<=C\_1 because C\_1 counts arrow-isomorphism classes at level pi\_1.}
\dgrclaim{false}{21}{XV3-024}{For a fixed poset I of size at most mu=r\textasciicircum{}pi\_0, the choices of all objects A\_i up to isomorphism can be bounded by |I| times C\_0<=mu.}
\dgrclaim{false}{21}{XV3-025}{For fixed objects A\_i, choices of all transition maps can be bounded by |I times I| times C\_0<=mu.}
\dgrclaim{false}{21}{XV3-026}{A cocone from a fixed I-diagram to X has at most |I| times s\textasciicircum{}mu choices.}
\dgrclaim{false}{23}{XV3-031}{Lemma 3.2 literally states r\textasciicircum{}rprime=2\textasciicircum{}rprime whenever rprime>=r, without an infinitude or lower-bound qualification.}
\dgrclaim{false}{23}{XV3-032}{Corollary 1 claims |Hom(A,X)|<=2\textasciicircum{}(r\textasciicircum{}pi\_0) for A,X in Filt\textasciicircum{}pi and r>=max(pi\_1,C\_1).}
\dgrclaim{false}{27}{XV3-048}{Theorem 2(1): p\_alpha(f) lies in Psi when alpha is large relative to r\textasciicircum{}pi\_0.}
\dgrclaim{false}{27}{XV3-049}{The initial ordinal of any cardinal rprime>r\textasciicircum{}pi\_0 may be chosen to satisfy the large-cofinality condition.}
\dgrclaim{false}{35}{XV4-025}{In the proof of Corollary2, $C\cap W\longleftrightarrow TC$.}
\dgrclaim{false}{36}{XV4-030}{For two precouples with $TC\subset C\cap W$ and $TF\subset F\cap W$, the equalities $TC=C\cap W$, $TF=F\cap W$, $W=TF\circ TC$ and the assertion that $(W,C,F)$ is closed are all equivalent.}
\dgrclaim{false}{40}{XV4-039}{Under only the hypotheses of Proposition4.4.2, Cof\_W cap W equals W\textasciicircum{}univ, the maps all of whose cobase changes belong to W.}
\dgrclaim{false}{40}{XV4-041}{Under only the Proposition4.4.2 hypotheses, Cof\_W cap W is Q-stable on the left.}
\dgrclaim{false}{40–41}{XV4-042}{Under the hypotheses of Corollary4.4.4, C=tilde(C0) equals Cof\_W.}
\dgrclaim{false}{40–41}{XV4-043}{Under only the hypotheses of Corollary4.4.4, TC=C cap W equals C cap W\textasciicircum{}univ and is contained in W\textasciicircum{}univ.}
\dgrclaim{false}{41}{XV4-045}{TC is Q-stable on the left under only Corollary4.4.4 hypotheses.}
\dgrclaim{open}{43}{XV4-054}{Given suitable small C0 and accessible W, does there exist a small TC0\ensuremath{\subset}C cap W generating C cap W under left lifting closure?}
\dgrclaim{false}{46}{XV5-015}{Footnote 17: $(d_2)$ should imply $(d_1)$.}
\dgrclaim{false}{52}{XV5-033}{Proposition 5.6 as first stated: at a rank where bounded objects are $\pi$-accessible, condition (5.3) alone is equivalent to $b_\pi$.}
\dgrclaim{false}{53}{XV5-036}{The proof applies (5.3) directly to the image J of the monomorphic stage indices in $\operatorname{Sub}_{\pi}(X)$.}
\dgrclaim{false}{58}{XV5-042}{Printed condition (cprime): M is the smallest class closed under filtered colimits, with no stated requirement to contain M\_0.}
\dgrclaim{open}{58–59}{XV5-045}{Footnote 22 asks whether (cprime) implies (c) without the special object-accessibility assumption on pi.}
\dgrclaim{false}{62–63}{XV5-053}{The printed lemma and display assert that the unrestricted $\operatorname{Sub}_{\pi}(Y)$ maps canonically and cofinally into $M^\pi/Y$.}
\dgrclaim{false}{64}{XV5-057}{The projection in Corollary 5.11 is described parenthetically as fibrant.}
\par\endgroup

\section*{XIX \textperiodcentered\ Chapter XIX, manuscript pages 1--50}
\pdfbookmark[1]{XIX}{dgrclaims.XIX}
\begingroup\raggedright
% DER_XIX: the definitions of the chunk's preamble, applied inside its group
\let\printedpage\relax\newcommand{\printedpage}[1]{\marginpar{\footnotesize p.~#1}}
\let\msp\relax\newcommand{\msp}[1]{\marginpar{\footnotesize [#1]}}

% Q-orthogonality arrows (half arrowheads = Q-orthogonal, full arrowheads = equality)
\let\Qo\relax\newcommand{\Qo}{\mathrel{\leftharpoondown\joinrel\relbar\joinrel\rightharpoonup}}   % lower half-head left, upper half-head right
\let\Qorev\relax\newcommand{\Qorev}{\mathrel{\leftharpoonup\joinrel\relbar\joinrel\rightharpoondown}} % mirrored form of (1.2)
\let\Qog\relax\newcommand{\Qog}{\mathrel{\leftarrow\joinrel\relbar\joinrel\rightharpoonup}}         % full head left
\let\Qod\relax\newcommand{\Qod}{\mathrel{\leftharpoondown\joinrel\relbar\joinrel\rightarrow}}       % full head right
\let\Qogd\relax\newcommand{\Qogd}{\mathrel{\leftarrow\joinrel\relbar\joinrel\rightarrow}}            % full heads both sides
% the same arrows crossed by a vertical bar (Quillen factorisation)
\let\Qx\relax\newcommand{\Qx}[2]{\mathrel{#1\joinrel\mapstochar#2}}
% crossed by two vertical bars (functorial factorisation)
\let\Qxx\relax\newcommand{\Qxx}[2]{\mathrel{#1\joinrel\mapstochar\mkern2mu\mapstochar#2}}
% Q-closure marks
\let\tld\relax\newcommand{\tld}{\tilde{}}   % left Q-closure mark (small tilde in exponent)
\let\rcl\relax\newcommand{\rcl}{\bar{}}       % right Q-closure mark (horizontal bar in exponent)
\let\Qorth\relax\newcommand{\Qorth}{\mathrel{\leftharpoondown\mkern-6mu\relbar\mkern-6mu\rightharpoonup}}
\let\Qorthx\relax\newcommand{\Qorthx}{\mathrel{\leftharpoondown\mkern-9mu{+}\mkern-9mu\rightharpoonup}}
\let\Qorthxfull\relax\newcommand{\Qorthxfull}{\mathrel{\leftharpoondown\mkern-9mu{+}\mkern-9mu\rightarrow}}
\let\Qcouple\relax\newcommand{\Qcouple}{\mathrel{\leftarrow\mkern-9mu{+}\mkern-9mu\rightarrow}}
\let\clos\relax\newcommand{\clos}{\longleftrightarrow}
\let\fl\relax\newcommand{\fl}{\longrightarrow}
\let\M\relax\newcommand{\M}{\mathcal{M}}
\let\Fl\relax\newcommand{\Fl}{\mathrm{Fl}}
\let\Ho\relax\newcommand{\Ho}{\mathrm{Ho}}
\let\cc\relax\newcommand{\cc}{\mathrm{c}}
\let\ddarr\relax\newcommand{\ddarr}{\mathrel{\twoheadleftarrow\joinrel\twoheadrightarrow}}
\let\harp\relax\newcommand{\harp}{\mathrel{\leftharpoondown\joinrel\relbar\joinrel\rightharpoonup}}
\let\deq\relax\newcommand{\deq}{\overset{\text{déf}}{=}}
\let\Mono\relax\newcommand{\Mono}{\mathrm{Mono}}
\let\Filt\relax\newcommand{\Filt}{\mathrm{Filt}}
\let\Cof\relax\newcommand{\Cof}{\mathrm{Cof}}
\let\qpair\relax\newcommand{\qpair}{\mathrel{\leftharpoondown\joinrel\relbar\joinrel\rightharpoonup}}
\let\Cat\relax\newcommand{\Cat}{\mathrm{Cat}}
\let\Diag\relax\newcommand{\Diag}{\mathrm{Diag}}
\let\id\relax\newcommand{\id}{\mathrm{id}}
\let\Qs\relax\newcommand{\Qs}{\tilde{\ }}
\let\fact\relax\newcommand{\fact}{\mathrel{\leftarrow\mkern-11mu\twoheadrightarrow}}
\let\limind\relax\newcommand{\limind}{\varinjlim}
\let\limproj\relax\newcommand{\limproj}{\varprojlim}
\let\Is\relax\newcommand{\Is}{\mathrm{Is}}

\dgrclaim{false}{1}{XIX-A007}{The strengthened equalities of (1.3) are claimed to be equivalent to adding separate double-orthogonal closedness of the corresponding classes to $\Phi\perp\Psi$.}
\dgrclaim{open}{4–5}{XIX-A026}{Is $\operatorname{Retr}(\Phi'')$ stable under pushout and transfinite composition, and therefore equal to $\Phi'''$, without small generation?}
\dgrclaim{false}{10}{XIX-A052}{If $\Psi=\Phi^\square$, then $\Psi(I_0)=\Phi(I_0)^\square$.}
\dgrclaim{conjectured}{12}{XIX-A060}{Question-conjecture 1.11: if $(\Phi,\Psi)$ is special, is the injective pair $(\Phi(I),\Psi[I])$ special for every small $I$, perhaps after an additional cardinal-filtration hypothesis?}
\dgrclaim{open}{12}{XIX-A061}{In 1.11, if $\Phi_0$ is a small generating family for $\Phi$, does the particular class $\Phi_0(I)$ generate $\Phi(I)$?}
\dgrclaim{false}{12}{XIX-A062}{Proposition 1.12, formula (1.13.2), prints $u^*(\Phi(I'))\subseteq\Phi\langle I'\rangle\iff u_*\Psi(I)\subseteq\Psi(I')$.}
\dgrclaim{false}{14}{XIX-A074}{The printed dual criterion in (2.5) uses a factorization $f=pi$ with $i\in W$ and $f\in C$.}
\dgrclaim{false}{17}{XIX-A087}{The Case A proof says that $TC'\subseteq W'$ is equivalent to $R(TC')\subseteq W'$.}
\dgrclaim{false}{10}{XIX-A090}{The equalities (1.15)--(1.16) identify $(v_!\Phi(I_0))^\square$ with $\Psi(I)$ and its left orthogonal with $\Phi\langle I\rangle$ under the sole displayed assumption $\Psi=\Phi^\square$.}
\dgrclaim{false}{19}{XIX-B003}{The displayed equality in Theorem 2.3 identifies $C\langle I\rangle$ with ${}^{\square}F(I)$.}
\dgrclaim{false}{23}{XIX-B021}{Decomposition holds for every prederivator of the form $\mathbf D_{\mathcal M,W}$, without hypotheses on the relative category.}
\dgrclaim{open}{24}{XIX-B024}{For a merely special triple, the source has not established small generating sets or factorizations for the componentwise classes $C(I)$ and $TC(I)$.}
\dgrclaim{conjectured}{25}{XIX-B035}{Condition (d) should be redundant, at least when epimorphisms are effective, after choosing suitable cardinals of the form $2^c$.}
\dgrclaim{false}{30}{XIX-B062}{$v^*(TF[I])=v^*v_*(TF(I_0))$ as image classes of arrows.}
\dgrclaim{open}{32}{XIX-B070}{Can special model structures outside the very special class define derivators on all small shapes?}
\dgrclaim{open}{32}{XIX-B072}{The author expresses uncertainty about the all-shapes homotopy-Kan-extension theorem attributed to Anderson.}
\dgrclaim{false}{35}{XIX-C023}{Every cofibration between cofibrant objects belongs to ambient Cof\_W by Brown's lemma.}
\dgrclaim{false}{38}{XIX-C038}{Proposition3.5: if \ensuremath{\Phi}=\textasciicircum{}perp\ensuremath{\Psi} is Q-closed, then for noetherian I the Anderson class \ensuremath{\Phi}\_I equals \textasciicircum{}perp(\ensuremath{\Psi}(I)).}
\dgrclaim{conjectured}{39}{XIX-C043}{For only transfinitely Q-saturated \ensuremath{\Phi}, conjecturally the transfinite Q-saturation of v!\ensuremath{\Phi}(I\_0) agrees with Anderson's noetherian extension.}
\dgrclaim{open}{39}{XIX-C045}{The desired identity is C<I>\ensuremath{\cap}W(I)=(TC)<I>.}
\dgrclaim{open}{40}{XIX-C048}{Question3.6 asks the identity for every small I, under small completeness of M and Q-closedness of C and TC.}
\dgrclaim{open}{40}{XIX-C053}{It may suffice for Question3.6 that only (TC)<I>--F(I) factor.}
\dgrclaim{open}{42}{XIX-C062}{Question3.9 requests a closed Quillen structure whose prederivator lacks some u!\textasciicircum{}W.}
\dgrclaim{open}{42}{XIX-C063}{Question3.9 also requests such a structure lacking some u*\textasciicircum{}W.}
\dgrclaim{conjectured}{42}{XIX-C064}{The author presumes testing u:I\ensuremath{\to}e should suffice for general Kan extensions.}
\dgrclaim{open}{43}{XIX-C065}{The requested counterexample would be stronger with cocomplete accessible ambient M and L-accessible W.}
\dgrclaim{open}{46–47}{XIX-C092}{A counterexample might have bicomplete M but H lack some small colimits.}
\dgrclaim{open}{47}{XIX-C096}{Question3.13: can a locally U-small category have all U-small colimits but lack some U-small limits?}
\dgrclaim{open}{48}{XIX-C099}{Question3.14: a cofibration category satisfying Cof\_4 should have u!\textasciicircum{}W for all small u.}
\dgrclaim{open}{48}{XIX-C100}{The author requests a localizer on bicomplete M whose prederivator lacks some left or right Kan extension, perhaps with M small.}
\dgrclaim{open}{48–49}{XIX-C101}{The author is less convinced such a counterexample exists when W is filtered-colimit stable and M has accessible objects and a cardinal filtration.}
\dgrclaim{open}{49}{XIX-C102}{Question3.15 asks whether an L-accessible arrow class in the stated accessible complete ambient category has a small left Q-generating subset.}
\dgrclaim{open}{49}{XIX-C103}{Question3.15 allows the extra assumption that \ensuremath{\Phi} is left Q-saturated: isomorphisms, cobase change, composition, and retracts.}
\dgrclaim{open}{49–50}{XIX-C105}{The author requests any mutually Q-orthogonal closed pair that does not factor all arrows.}
\dgrclaim{open}{49–50}{XIX-C106}{The same nonfactorization request with special ambient M is of particular interest.}
\par\endgroup

\dgrbackchapter{Transcription corrections}
The corrections to the transcription, counted by kind for each part: readings of the transcription corrected against the edition (transcription), slips of the edition kept in the text with their intended reading recorded (edition slip kept), and doubtful readings settled against the edition (doubt resolved).

\begin{center}\small
\begin{tabular}{lrrrr}
\toprule
Part & transcription & edition slip kept & doubt resolved & total \\
\midrule
I-a & 8 & 29 & 5 & 42 \\
I-b & 3 & 69 & 3 & 75 \\
II-a & 11 & 65 & 1 & 77 \\
III-V & 14 & 114 & 6 & 134 \\
VI-a & 11 & 43 & 5 & 59 \\
VI-b & 2 & 67 & 3 & 72 \\
VI-c & 1 & 43 & 3 & 47 \\
VII-a & 19 & 81 & 2 & 102 \\
VII-b & 6 & 53 & 0 & 59 \\
VII-c & 3 & 45 & 7 & 55 \\
VIII-X & 0 & 29 & 3 & 32 \\
XI-XIIa & 1 & 72 & 6 & 79 \\
XIIb-a & 0 & 48 & 2 & 50 \\
XIIb-b & 1 & 27 & 2 & 30 \\
XIIc-a & 2 & 41 & 4 & 47 \\
XIIc-b & 3 & 38 & 0 & 41 \\
XIIapp & 2 & 21 & 3 & 26 \\
XIII-a & 0 & 33 & 2 & 35 \\
XIII-b & 5 & 28 & 1 & 34 \\
XIV & 0 & 2 & 1 & 3 \\
XV-a & 2 & 37 & 5 & 44 \\
XIX & 6 & 36 & 10 & 52 \\
\midrule
total & 100 & 1021 & 74 & 1195 \\
\bottomrule
\end{tabular}
\end{center}

\nocite{*}
\bibliographystyle{plainnat}
\bibliography{LesDerivateurs_companion}
\end{document}
